\frontmatter

\quellseite{1}
\begin{titelseite}
Dr. Martin Luthers\\
Sämmtliche Schriften,\\
herausgegeben von\\
Dr. Joh. Georg Walch.\\
Zweiter Band.\\
Auslegung des ersten Buches Mose.\\
Zweiter Theil.\\
Neue revidirte Stereotypausgabe.\\
St. Louis, Mo.\\
CONCORDIA PUBLISHING HOUSE.
\end{titelseite}

\begin{titelseite}
Dr. Martin Luthers\\
Auslegung des ersten Buches Mose.\\
Zweiter Theil.\\
In deutscher Sprache herausgegeben\\
von\\
Dr. J. G. Walch.\\
Aufs Neue herausgegeben im Auftrag des Ministeriums der deutschen ev.-luth. Synode von Missouri, Ohio u. a. Staaten.\\
St. Louis, Mo.\\
CONCORDIA PUBLISHING HOUSE.
\end{titelseite}

\quellseite{2}
\chapter{Vorwort}

Bei der Redaction dieses zweiten Theils von Luthers Auslegung des ersten Buches Mose, der hiermit von Neuem der Oeffentlichkeit übergeben wird, sind genau dieselben Grundsätze befolgt worden, wie bei der Herausgabe des ersten Theils. Es genügt daher hier, auf das Vorwort zum ersten Band der revidirten Walch'schen Ausgabe von Luthers Werken zurückzuweisen. Das diesem zweiten Band angefügte ausführliche Register ist von Herrn Christian Körner angefertigt worden. Dasselbe zerfällt in drei Theile. Der erste Theil, das „Hebräische Wort-Register“, enthält alle hebräischen Wörter, welche in den beiden bisher erschienenen Bänden der Sämmtlichen Schriften Luthers erklärt werden. In solchen Fällen, in denen Luther eine abgeleitete Form des Wortes erklärt, ist diese neben der Grundform in Parenthese angegeben. Der zweite Theil, das „Spruch-Register“, enthält zwar nicht alle Schriftstellen, welche Luther bei seiner Erklärung der Genesis anführt, aber doch alle diejenigen, welche er in den genannten beiden ersten Bänden seiner Schriften, ausführlicher oder kürzer, erklärt. Dabei werden nicht, wie dies wohl früher üblich war, die Worte, auf welche sich die Erklärung bezieht, angeführt, weil dies unnöthig schien und zu viel Raum beansprucht hätte. Der dritte Theil, das „Sach-Register“, umfaßt zugleich ein „Autoren-Register“. Es will in möglichster Vollständigkeit, was Luther über irgend einen Gegenstand oder Autor gesagt hat, wiedergeben. Bei der Anordnung der angeführten Stellen ist Bedacht darauf genommen worden, solche derselben, die Verwandtes behandeln, zu gruppiren, ohne daß eine streng logische Eintheilung beabsichtigt worden wäre. Der Uebersicht wegen sind längere Artikel in numerirte Unterabtheilungen zerlegt, während in anderen Fällen durch gesperrten Satz oder alphabetische Anordnung die Uebersichtlichkeit befördert wird. Beim Nachschlagen wolle man stets das Hauptwort aufsuchen, nicht eine nähere Bestimmung desselben; also suche man z. B. „Rothes Meer“, oder „Jüngster Tag“, unter „Meer“, „Tag“ u. s. w. Hin und wieder sind wohlbekannte Abkürzungen des Raumersparnisses wegen gebraucht worden. Wo das Stichwort in einem Abschnitt wiederkehrt, wie das öfters vorkommt, ist es, außer wo die Verständlichkeit es anders verlangte, nur mit den Anfangsbuchstaben angegeben, woran man sich ja leicht gewöhnen wird. In allen Theilen des Registers bezeichnet die römische Zahl den Band, die arabische die Columne.

\quellseite{3}
St. Louis, am Pfingstfeste 1881.

\begin{unterschrift}Die Redaction.\end{unterschrift}

\chapter{Des Uebersetzers Johannes Gudenus Vorrede}

Dem Gestrengen, Ehrbaren und Ehrenfesten Christoph von Steinberg, meinem günstigen lieben Junker und Patron.

Gottes Gnade durch Jesum Christum, unsern einigen Erlöser und Heiland! Gestrenger, Ehrbarer und Ehrenfester, günstiger lieber Junker! Wie unser Herr GOtt bei den Zeiten Josephs, des großen, heiligen und vortrefflichen Patriarchen, da derselbe ist nach Egypten verkauft worden, dem egyptischen Könige, seinen Fürsten, Räthen und Amtleuten, samt vielen Unterthanen im selben Königreiche, ein hörend Ohr, wie es Salomo Spr. 20, 12. nennt, gegeben, daß sie sich von Joseph (welcher durch sonderliche Schickung GOttes zu ihnen nach Egypten gekommen) haben lehren und berichten lassen in der Lehre von der rechten Erkenntniß GOttes, und dadurch also sind bekehrt worden: also hat der liebe, gnädige GOtt in dieser letzten Zeit, innerhalb etlichen und dreißig Jahren, in Deutschland auch solche unaussprechliche Gnade erzeigt, daß er nicht allein ein sehend Auge, das ist, den theuren Mann D. Martin Luther, heiliger Gedächtniß, als einen trefflichen, geistreichen Lehrer erwecket hat, durch welchen der alte Sauerteig falscher verkehrter Menschenlehren, so im Pabstthum allenthalben mit Gewalt eingerissen und überhand genommen hatten, ausgefegt und die Kirche recht ist reformirt worden; sondern hat auch daneben hörende Ohren gegeben, so die Lehre dieses trefflichen Propheten gehört, erwogen und angenommen haben; nicht allein was gemeine Unterthanen sind in Landen und Städten, sondern auch etliche große namhaftige Regenten, als Chur- und Fürsten, insonderheit des löblichen Hauses zu Sachsen, dazu Grafen und viel Ehrliebende und auch Namhaftige vom Adel, zudem auch etliche ehrbare Städte, welche alle sich dermaßen erzeigt und vernehmen lassen, daß sie unter die Zahl derer gehören, davon Jesaia sagt im 49. Capitel v. 23.: „Und die Könige sollen deine Pfleger, und ihre Fürstinnen deine Säugammen sein.“ Denn wiewohl der Teufel (wie vom Anfange seine Art gewesen) durch seine Glieder, den andern großen Haufen, über alle Maßen sehr getobet und gewüthet, und fast auf allen Seiten versucht hat, wie er dieses Licht der reinen Lehre, uns durch gedachten D. Luther offenbart, wiederum auslöschen möchte; so hat gleichwohl GOtt seiner kleinen Heerde und der rechten Kirche, so er sich durch die reine Lehre seines Wortes sammeln läßt, diese Pfleger gegeben, welche sich das haben lassen ernstlich angelegen sein: und noch jetzt, nachdem sie erstlich ihre Thore weit aufgethan und dieselben weit gemacht, daß der König der Ehren, wie der 24. Psalm V. 7. sagt, bei ihnen ist eingezogen, wie sie nun auch können und mögen die christliche Lehre dieses Königes befördern, handhaben und ausbreiten helfen. Mit solchen Pflegern und Säugammen will, sage ich, unser HErr GOtt zu jeder Zeit aus Gnaden seine Kirche versehen, und sie also mehren und erhalten, wenn sich auch der Teufel und seine Schuppen darüber zerreißen sollten. Unter solche getreue Pfleger und Kirchenväter ist billig Euer Ehrenfeste und Euer gleichen etliche andere Ehrliebende vom Adel mehr auch zu rechnen. Denn was ihr von Anfange dabei gethan, damit ja diese reine Lehre ihren Lauf möchte frei behalten, weiß der wohl, vor welchem nichts verborgen ist, der es auch laut seiner Verheißung reichlich vergelten wird. Und dieweil ich Eurer Ehrenfesten ehrliches und christliches Gemüthe also erkannt habe, daß Ihr vor der Welt keinen sonderlichen Ruhm damit wollt gesucht haben, daß Ihr vielen armen Kirchendienern und Schülern große Hülfe und Steuer gethan, damit ja der Kirchen- und Schuldienst im Schwange bleiben möchte, in Ansehung, daß der Teufel die böse und unbändige Welt dazu reizt, daß sie die Schul- und Kirchendiener nur verhungern soll lassen, auf daß also die reine Lehre und rechter Gottesdienst dahin falle und wiederum allerlei Abgötterei und falscher Gottesdienst allgemach aufgerichtet werde (dazu denn die Welt ihr Geld gerne pfleget zu geben); wie im Volke Israel auch geschehen, wie man im Buche Nehemias Cap. 13, V. 4. lesen mag: so will ich davon jetzt nicht sagen; Christus, unser lieber Herr, der da sagt Matth. 10, V. 42.: Es soll auch ein Trunk kaltes Wasser, seiner geringsten Diener einem gegeben in seinem Namen, nicht unvergolten bleiben, derselbe wird am Tage seiner Erscheinung solche reiche Wohlthaten, so Euer Ehrenfeste mir und andern Dienern des Wortes freiwillig erzeigt, öffentlich ausrufen, und davon urtheilen, daß es ihm selbst geschehen, wird es auch unvergolten nicht lassen.

\quellseite{4}
Damit ich aber gleichwohl mich etwas ein wenig dankbar erzeigen möge, will ich Euer Ehrenfesten diese meine Arbeit, so ich durch Euer vielfältiges Anregen, mit Verleihung göttlicher Gnade, nun innerhalb eines Jahres darauf gewendet, daß dieser reiche Schatz, so der theure Mann D. Luther, der Deutschen letzter Prophet, in diesem Commentar über das Erste Buch Mose hinter sich gelassen, aus dem Lateinischen in das Deutsche möchte gebracht werden, gewidmet haben, und unter Eurem und Anderer vom Adel Namen lassen ausgehen, dieweil es unserm lieben Lehrer, dem Herrn Philippus, also gefallen hat, daß man es durch den Druck hat verfertigen sollen. Denn wo es ohne desselben Mannes Bewilligung wäre, sollte ich mir zwar ungern das zuzumessen haben, daß ich diese meine Arbeit sollte haben in Druck gegeben. Was aber nun dieses für ein trefflicher Schatz ist, werden diejenigen mit der That befinden, so diese Auslegung mit Fleiß lesen werden. In Summa, man findet hierin den rechten Kern der Lehre, so unser Herr Gott uns durch D. Luther geoffenbart hat; wie davon der Herr Philippus, D. Jonas, und andere geistliche Männer recht geurtheilt haben. Daß nun Eure Ehrenfesten sammt Andern vom Adel dieses Werk so fleißig und treulich, aus christlichem Eifer gefördert und darauf nicht geringe Kosten gewendet, solches wird ohne Zweifel unserm Herrn Gott zu Ehren und seiner lieben Kirche zur Besserung gerathen. Der liebe gnädige GOtt, Vater unsers Herrn Jesu Christi, wolle durch den Heiligen Geist unser aller Herzen erleuchten, daß wir diese Lehre, so er uns aus großer Gnade in diesen letzten Zeiten durch vielgedachten theuern Mann D. Martin Luther, heiliger Gedächtniß, wiederum geoffenbart und damit aus der großen Finsterniß und Greuel des Papstes erlöset, ernstlich annehmen, dadurch auf vielerlei Weise gebessert mögen werden, und dabei bis an unser Ende beharren und beständig bleiben, auf daß wir endlich durch Christum, den einigen Mittler und Heiland, auf welchen uns allein diese Lehre weiset, zur ewigen Seligkeit mögen erhoben werden.

Datum Braunschweig, am Tage der Erscheinung Christi, im Jahre nach Christi, unseres Herrn, Geburt 1557.

Euer Ehrenfesten williger Johannes Gudenus.

\chapter{Philipp Melanchthons Vorrede}

\begin{center}Philippus Melanchthon.

wünschet dem frommen gottseligen Leser seinen Gruß.\end{center}

Man soll oft bedenken den Rath Gottes, nämlich, daß er haben will, daß man zu allen Zeiten wissen soll den Anfang der Welt, und dazu auch, wie die Kirche erstlich habe angefangen, und folgendes ausgebreitet und immer gemehret sei worden, und daß man ja solche große Dinge im frischen Gedächtniß behalte, und dasselbe nach der Ordnung der Jahre, so dazu eingesezt und beschrieben worden, auf daß man immer zurück an den Anfang denken möge. Und wir sollen GOtt dafür danken, daß er unter so großen Zertrennungen der Regimente und Herrschaften und Zerstörungen der Städte, dazu auch unter solchem Verderben des menschlichen Geschlechts gleichwohl diese seine göttliche Lehre nicht läßt vertilgt werden, und daß er auch schöne klare Zeugnisse sehen läßt, als nämlich, die Auferweckung der Todten und daß etwa die Sonne in ihrem Laufe aufgehalten wird, und viel andere Dinge mehr; auf daß wir ja gewißlich wissen sollen, beide, daß ein GOtt sei und daß er sich eine Kirche lasse sammeln, und dazu auch, daß diese Lehre, so er gegeben hat, die rechte, gewisse, unbewegliche Wahrheit sei.

Derhalben sollen wir diese göttliche Stimme für gewiß annehmen, und beständiglich fliehen und verdammen die heidnischen greulichen Irrthümer, deren Etliche gedichtet haben, daß die Creaturen in der Welt von ohngefähr aus den Atomen, das ist, aus dem kleinen Staube, so in der Sonne fliegt, zusammengefetzt und gemacht sein sollen, und daß viel unzähliger Welten gewesen und sein werden und dazu auch mancherlei Natur aller Dinge. Die Andern aber, wiewohl sie diese scheußlichen und erschrecklichen Worte nicht loben dürfen, und etwas suchen, das zierlicher und besser geschmückt ist, so fehlen sie doch der Wahrheit weit, gleich als die da blind sind, und haben greuliche Irrthümer, damit GOtt auf das höchste gelästert wird, erdacht; wie die Stoiker dichten, daß auch GOtt keine andere Bewegung erregen könne, denn allein dahin er gerissen wird durch das Getriebe der Materie oder Mitteldingen, und daß er ebensowohl eine Ursache sei der greulichen Laster als der guten Werke. Endlich ist ein großer Haufe scheußlicher Irrthümer außerhalb der Kirche und Gemeinde GOttes, als nämlich, vom Anfange der Welt, von der Natur des Menschen, von göttlicher Verheißung, von Ursachen des Bösen, vom Ende des Menschen: ob der Mensch allein zu diesem sterblichen Leben geschaffen sei. Item, ob man auch auf ein ander Leben warten und hoffen solle, in welchem ein Unterschied sein werde zwischen den Gerechten und Ungerechten; und ob auch eine Versammlung oder Haufen Volks sei in solcher großen Uneinigkeit so mancherlei Meinung und Sitten der Menschen, davon man gewißlich halten und sagen möge, daß er GOtt wohlgefalle. Item, wie dieser Haufe versammelt und von andern Menschen unterschieden müsse werden: welche Aemter oder Werke GOtt wohlgefallen mögen, und ob der andere Haufe, darin so eine große Menge ist, von GOtt gar verworfen sei.

Von diesen großen Dingen wird uns allein in der Gemeinde GOttes die rechte nöthige Lehre vorgehalten. Und ist im ersten Buche Mose der rechte Grund und Ursprung, welcher jedermann soll und muß sehr wohl bekannt sein, dieweil in einer jeglichen Lehre daran viel gelegen ist, daß man darin den Anfang und rechten Ursprung wissen und recht verstehen möge; wie man sagt: Principium dimidium est totius, et plus quam dimidium, das ist, wenn man den Anfang oder Grund eines Dinges hat, so hat man es halb, ja, man hat davon mehr als die Hälfte.

\quellseite{5}
Es hat aber GOtt das Predigamt eingesetzt, nicht daß jemand eine neue Lehre erdichten soll, sondern daß die Lehrer eben diese prophetischen und apostolischen Bücher dem Volke erzählen und vorhalten sollen, auf daß sie daraus die reine unverfälschte Lehre von GOtt und der Seligkeit des menschlichen Geschlechts nehmen sollen. Darnach aber, dieweil man das junge Volk und den gemeinen Mann erinnern muß der Form und Weise zu reden, und wie ein Stück ordentlich auf das andere folge, wo man die Schrift liest: so muß man erklären, was an einem jeden Orte für eine Art und Weise zu reden geführt werde. Zudem muß man auch anzeigen, woher die mancherlei Argumente, so hin und wieder gebraucht werden, genommen seien, und muß die Ungelehrten berichten, wie sie in ihrem Gebete, da sie GOtt anrufen, und sonst in ihrem Leben die Sprüche und Exempel der Schrift brauchen sollen. Hierin ist ein großer Unterschied unter denen, so die Schrift auslegen.

Dieweil aber GOtt durch das Amt Luthers die Kirche von mancherlei Irrthum gereinigt, dazu er ihm große treffliche Gaben hat gegeben, so weiset es nun das Werk an sich selbst aus, daß seine Auslegungen die Schriften der andern Schreiber, so die Schrift ausgelegt, weit übertreffen. Und daß man dieselben haben mag und daß sie auch die Nachkommen nach uns haben mögen, ist nütze und gut, und nicht allein daß die Schriften der Propheten und Apostel recht mögen erklärt werden, sondern auch daß die Zeugnisse und schönen Bekenntnisse dieses großen trefflichen Mannes (deß Herz vom Heiligen Geiste regieret worden) von allen Artikeln jedermann mögen bekannt werden, daß damit hernach die Göttlichen mögen gestärkt werden, wo es vorfällt, daß man über der Lehre streitig wird.

Derhalben soll man billig den frommen gelehrten Männern danken, welche so treulich beide, andere Auslegungen Luthers und sonderlich auch diese über das erste Buch Mose (das man ja vornehmlich wissen soll), in Druck haben lassen ausgehen, nämlich D. Caspar Creuziger, Georg Rörer, Veit Dietrich und Hieronymus Besold zu Nürnberg. Denn nach dem Tode Veit Dietrichs ist diese Arbeit auf Hieronymus Besold geerbt (daß ich es also nenne), welcher dem, so D. Creuziger und M. Georg Rörer mit ihren Händen geschrieben, treulich nachgefolgt. Denn er hat gesehen, daß dasselbe mit dem Exemplar, so er mit seiner Hand geschrieben, und auch noch mit einem andern, so M. Johannes Stoltz mit seiner Hand geschrieben (welcher Mann denn auch sehr fleißig und getreu ist) übereingekommen und denselben gleichförmig gewesen.

Nun soll aber der Leser die Lehre noch einander betrachten, als nämlich, von GOtt, dem Schöpfer; vom Anfange der Welt; von der Schöpfung der Menschen; wie sich GOtt vor dem Falle des Menschen geoffenbart und wie die Kirche angefangen; darnach, vom Falle und Erlösung des Menschen, darin sich GOtt abermals geoffenbart, und die Verheißung gegeben von der Erlösung der Menschen und wie die alte Schlange sollte zertreten werden; und wie wiederum aus großer Gnade und Barmherzigkeit eine Kirche aufgerichtet und unsere ersten Eltern mitten aus dem Tode errettet.

Darnach sollst du das auch betrachten, welches und wo die rechte Kirche und Gemeinde GOttes gewesen sei, nämlich, darin die gegebene Verheißung vom künftigen Opfer ist gegeben worden. Zudem sollst du auch das Gesetz und die Verheißung unterscheiden und diese betrübte Zwietracht und Uneinigkeit im menschlichen Geschlechte erkennen, nämlich, daß allezeit zwei Haufen gewesen und jetzt noch sind, dazu auch hinfort bleiben werden, da den einen der Sohn GOttes selbst, von welchem die Verheißung redet, durch das Predigtamt des Evangeliums sammelt, und wahrhaftig bei demselben ist, darin regiert, und bedeckt die Sünde derer, so ihn erkennen, zündet ihre Herzen an mit dem Heiligen Geiste, daß sie GOtt recht anrufen und ihm gehorsam sind, und gibt ihnen das ewige Leben und Seligkeit. Der andere Haufe aber kennt GOtt nicht, geht in Finsterniß um, ist entbrannt vom Ehrgeize, Unzucht und andern sündlichen Begierden, verachtet und haßt das Evangelium. Diesen unsinnigen Haufen reizen und hetzen die Teufel wider die rechte Kirche, als, daß Cain seinen Bruder erwürgt hat und Esau auch gleich denselben Willen gehabt.

Dieser große Kampf, darüber sich alle Menschen sehr verwundern, ist bald im Anfange in der ersten Predigt mit diesen Worten zuvor verkündigt worden, da GOtt 1 Mos. 3, 15. spricht: „Ich will Feindschaft setzen zwischen dir und dem Weibe, und zwischen deinem Samen und ihrem Samen.“ Denn also hat die Kirche müssen zuvor erinnert werden, auf daß sie wissen möchte, daß sie darum von GOtt nicht verachtet oder verworfen wäre, dieweil die Andern das Regiment würden haben, so die göttliche Lehre verachten, und das Häuflein, welches die Lehre des Evangeliums bewahrt, immer verfolgen und unterdrücken werden. Man soll auch betrachten die Ordnung, wie die nach einander gelebt, so die Kirche mit dem Worte regiert haben; und soll bei jeglichem insbesondere betrachten, was er für eine Lehre geführt habe; darnach auch die Zeugnisse der Offenbarung GOttes und der Lehre. Item, wie er sie überwunden, dazu, wie man gesehen, daß GOtt gegenwärtig bei ihnen gewesen und sie errettet hat; auch wie die Kirche ist gesammelt worden; was sie für Sprüche und Exempel geführt, so uns unterweisen und unsern Glauben stärken mögen, auch was sie für Tugenden gehabt, denen wir billig sollten nachfolgen.

Diese Hauptstücke sind nach einander in den Historien der vorigen Bücher beschrieben worden. Als: von Adam und Eva, von Abel, Seth, Henoch, Noah, Sem, Abraham, Lot, Isaak. Diese alle führen einerlei Lehre von GOtt dem Schöpfer, von der Schöpfung aller Dinge, von dem Fall unserer ersten Eltern, von der Sünde im menschlichen Geschlechte, vom Gesetze GOttes, von der Verheißung und vom Mittler, von der Gnade und Vergebung der Sünden, vom Glauben, vom rechten Anrufen, vom neuen Gehorsam, von der Kirche, vom Kreuze, vom zukünftigen Gericht. Aber der Beruf und die Werke ihrer aller sind fast ungleich. Noah wird berufen, daß er die Kirche, so in der Sintfluth übergeblieben, erhalten sollte und darnach dieselbe bei den Nachkommen wiederum aufrichten. Da die Nachkommen Noahs zerstreut und jetzt nun das babylonische Reich angefangen war, dazu auch die Fürsten da unten Heiden immer größer geworden, da richtet GOtt die Kirche wiederum auf mit einem neuen Beruf, wiederholt die Verheißung und thut noch die Beschneidung dazu, auf daß die Menschen, wenn sie an dasselbe Glied denken würden, sich des zukünftigen Samens zu erinnern hätten. Darum haben Isaak und Jakob die Lehre, so sie von Abraham genommen, auch getrieben und dieselbe andere Leute gelehret.

Nun sollst du aber in dieser Historie bedenken, wie eine große Ehre und wie das den Gottseligen so eine große Stärkung gewesen, daß zu einerlei Zeit umher, fast nahe bei einander gewohnt und zugleich mit einander gelebt haben, auch oftmals zusammen gekommen sind, die Lehre auszulegen und sich unter einander zu trösten, nämlich, Sem (von dem man meint, daß er Melchisedek genannt werde), Abraham, Isaak und der König Abimelech. Und wiewohl die Menge der Kirche und Gemeinde GOttes allezeit kleiner ist denn der Haufe der gottlosen Heiden, so haben doch diese großen Lehrer gleichwohl eine eben große Gemeinde gehabt. Denn GOtt will allezeit haben, daß die Kirche im menschlichen Geschlechte soll gehört werden, auf daß der Sohn erkannt werde und daß auch ein Theil des menschlichen Geschlechts Erben werden der himmlischen Herrlichkeit; darum stellt er sie gleich als auf ein hohes Schloß, auf daß sie ja gesehen und gehört werden möge. Und gibt ihr eine eben große Menge Volkes, daß sie Zeugen seien der Lehre, und die großen Wohltaten des Sohnes GOttes empfangen mögen, welchen ja der ewige Vater nicht vergeblich auf Erden gesandt hat.

Darum soll man auch in dieser Historie Jakob vornehmlich die Lehre betrachten, welche zum öftermal ist wiederholt worden. Denn GOtt selbst, da er zu Bethel auf der Leiter steht und mit Jakob redet, zeigt an, daß dieses der rechte wahrhaftige GOtt sei, der sich dem Abraham geoffenbart und die Verheißung gegeben hat; die ist an dem Orte mit denselben Worten wiederholt worden.

Darum will er haben, daß man diesen wahren GOtt erkennen und anrufen soll, der die Verheißung gegeben hat, und will haben, daß man den Mittler ansehen soll. Er will, man soll Vergebung der Sünden und die Versöhnung mit Gott durch den Mittler empfangen. Er will, daß man wissen soll, daß derselbe Mittler auch zur selbigen Zeit den Jakob bewahre und regiere. Und läßt sich den Gottesdienst gefallen, welchen Abraham von Gott selbst, da er mit ihm mit seiner göttlichen Stimme geredet, gelernt hat; davon geschrieben steht Römi. 4, 3.: „Abraham glaubte dem Herrn, und das rechnete er ihm zur Gerechtigkeit.“ Daher hat Jakob auch gewußt, daß dieser Glaube der vornehmste Gottesdienst wäre, und hat es dafür gehalten, daß er dadurch um dieses Mittlers willen Gott gefällig wäre und nicht um seiner eigenen Würdigkeit willen; und daß er durch diesen Mittler regieret würde, und daß die Sünde und der Tod weggenommen und vertilget sind; so werde um desselben Mittlers willen und durch ihn Gerechtigkeit und ewiges Leben wiedergebracht und gegeben werden. Denn er hat zu der Verheißung, welche sein Vater und er gehört hatten, dieses Wort auch noch gesetzt: „Der Same des Weibes wird der Schlange den Kopf zertreten“, 1 Mos. 3, 15.

\quellseite{6}
Dieses Bekenntniß thut er selbst, da er sagt, 1 Mos. 48, 15. 16.: „Der Gott, vor dem meine Väter, Abraham und Isaak, gewandelt haben; der Engel, der mich erlöset hat von allem Uebel“, vom ewigen Tode, und anderm Unfall, „der segne die Knaben“ etc. Und ist daran kein Zweifel, dieser Engel, dessen hier gedacht wird, ist der Sohn Gottes, von welchem Jakob gewußt, daß er zu der Kirche und Gemeinde Gottes gesandt wird, daß er sie regieren, schützen, dazu von Sünde und Tod erlösen soll. Diese Wohlthat gibt keine Gewalt irgend einer Creatur, sondern der Sohn Gottes erzeiget sie der Kirche.

Und ist daran kein Zweifel, es hat Jakob von diesem Herrn auch gepredigt, der Gerechtigkeit und Leben wiederbringt, da er hernach sagt im 49. Cap. V. 10.: „Es wird das Scepter von Juda nicht entwendet werden“ etc., „bis daß der Held komme; und demselben werden die Völker anhangen.“ Wiewohl aber andere das Wort schiloh nicht auf einerlei Weise ausgelegt haben, so zeigt doch das Stück, so im Texte darauf folgt, offenbar an, daß dadurch der Held verstanden werde, der auch die Heiden mit Gott versöhnen und zu ihm versammeln wird, welcher der Schlange den Kopf zertreten, die Sünde und den Tod tilgen und die Leute davon erlösen wird. Und ich halte, daß dies Wort, schiloh, recht verdolmetscht wird, daß es Helfer und Heiland heiße.

Und ist die Etymologie oder Eigenschaft des Wortes lustig und fein; wie dieselbe den Hebräern wohl bekannt ist, welche das Wort auslegen, daß es im Lateinischen soviel heißen soll als, foetus ejus, videlicet matris, das ist, die Frucht der Mutter; auf daß es sich mit dem Spruche reime: „Der Same des Weibes wird der Schlange den Kopf zertreten.“

Und ist herrlich und sehr fein geredet: die Frucht der Mutter. Denn dieser Sohn Gottes nimmt die schwache menschliche Natur an sich, welche mit großem Jammer und Leid beladen ist, dazu Gottes Zorn und den Tod auf sich liegen hat; und hängt dieser Sohn Gottes am Kreuze, ist eine mißgestaltete blutige Frucht geworden, und steht doch gleichwohl vom Tode wiederum auf und wird lebendig. In uns aber ist die Natur ganz unrein, ist eine befleckte und unreine Frucht und der Sünde willen. Wenn aber der Sohn Gottes diese unsere Frucht an sich nimmt und sich darein verkleidet; wenn er uns durch das Wort des Evangeliums und durch seinen Heiligen Geist lebendig macht: so wird diese Haut ein Tempel Gottes und wird von Sünde und Tod erlöset.

Aber eben in dieser Auslegung Luthers werden alle diese Predigten Jakobs und der Andern reichlich erklärt. Und dieses habe ich nur darum erzählt, daß ich gleich zum Eingange die jungen Leute damit möchte erinnern, daß man in diesen Historien nicht allein Exempel des Haus- und Weltregiments suchen soll, sondern man soll vornehmlich die Lehre und Einhelligkeit der ersten Lehrer in der Kirche, wie sie im Glauben und rechter Anrufung einig gewesen, ansehen und betrachten, auf daß wir uns selbst mit ihrer Einhelligkeit auch stärken mögen. Denn Adam, Abraham, Isaak und Jakob sind auch Zuhörer und Zeugen eben desselben Evangeliums, und haben dasselbe auch ausgebreitet, welches jetzt von Gottes Gnaden in den Kirchen gelehrt und gehört wird, so die Lehre Lutheri, welche er uns in dieser Auslegung und in seinen andern Schriften vorgehalten, annehmen.

Derhalben ist es sehr nütze und gut, daß man die Anrufung und Predigten dieser Väter oft und fleißig betrachte; auf daß du wissen mögest, daß du wahrhaftig ein Mitgeselle seiest ihrer Erlösung, ihres Glaubens und Anrufung, dazu auch ein Miterbe der ewigen Herrlichkeit. Und wo du die Lehre betrachtest hast, so siehe auf die Kirche und Gemeinde Jakobs: siehe, wie viel Elendes und Jammers er gelitten, und schaue darin auf die Exempel des Glaubens und anderer Tugenden, dazu auch auf die Gegenwart des Sohnes Gottes, der ihn erlöset; auf daß dadurch die Zeugnisse der Offenbarung an den Tag gegeben und die Kirche auch möchte erhalten werden.

Darnach siehe auch an, wie die Kirche verwundet, betrübet und geplaget wird. Es toben und wüthen die Teufel so sehr, daß sie auch in diesem kleinen Häuflein, als im Hausgesinde Jakobs, etliche vornehme Männer treiben, daß sie große greuliche Sünden begangen haben. Der eine Sohn hat seines Vaters Weib durch Blutschande beschlafen. Darnach unterstehen sich ihrer viel, ihren eigenen Bruder Joseph zu tödten. Wer wollte nicht erschrecken, wenn er diese greulichen Exempel ansieht in dem Hausgesinde, welches wahrhaftig die Kirche und Gemeinde Gottes war, darum zumal auch kein Zweifel gewesen? Aber Gott will haben, daß wir unsere Schwachheit erkennen sollen und wider die hitzigen Anschläge des Teufels immer wacker seien; zudem will er auch haben, daß man die, so gefallen sind, mit den Exempeln solcher großen Leute wiederum aufrichten soll; denn wiewohl dieselbigen auch schwerlich verwundet gewesen, sind sie doch wiederum geheilet worden.

Ueber das, wiewohl Jakob viel Feinde gehabt, ist ihm doch auch sein eigener Bruder Esau sehr feind, der ihm doch von Natur am nächsten verwandt gewesen, und den sein Vater auch auferzogen und unterwiesen hat eben zu derselbigen Religion, wie man Gott glauben und recht anrufen, sich auch in allerlei Tugenden und guten Werken üben soll. Also, dieweil es in der Kirche gebräuchlich und gemein ist, daß die, so einander verwandt und gute Freunde sind, über der Religion pflegen uneinig zu werden; so sollen wir unsere Herzen dazu bereiten und stärken, daß wir diesen großen Schmerz, welchen die Feindschaft derer, so unsere guten Freunde gewesen, gebieret, weißlich dulden und vertragen können. Sophocles sagt: Es sei keine schwerere Wunde und thue einem Menschen nichts so wehe, als wenn ihm ein guter Freund gefällig und uneinig mit ihm wird. Und in den Propheten findet man auch dergleichen Klage. Dieweil aber die Exempel uns lehren und berichten, was die Ursachen sein solcher Trübsale, und dabei auch Trost anzeigen, so lindern dieselben Exempel dem Leser seinen Schmerz. Und ist ein Stolz und Hoffart, daß wir nicht leiden wollen, was Andere gelitten haben, so frömmer gewesen sind denn wir, als, Abel, Ismael, Joseph und Andere mehr. Wir sollen wissen, daß der Gehorsam und wo man in Herzen stille ist, solche Trübsale mit Geduld zu ertragen, Gott wohlgefällt, und daß unterdeß der Sohn Gottes für uns sorgt und daß er unser getreuer und beständiger Freund ist. Mit diesem Trost hat sich David aufgerichtet, da er sagt im 27. Psalm V. 10.: „Mein Vater und meine Mutter verlassen mich; aber der HErr nimmt mich auf.“

Jetzig(er) Zeit gerathen zum öftern die, so gute Freunde mit einander gewesen, in Uneinigkeit und Zank, wird auch der rechten Kirche von mancherlei Feinden fast hart zugesetzt. Die Türken, Pabst und Bischöfe, Könige und Fürsten, welche des Pabstes und seiner Bischöfe getreue Diener sind, unterstehen sich, die arme Kirche mit Schwert und Feuer zu vertilgen. Und die Mönche, und etliche Andere mehr, so an der Fürsten Höfen für gelehrt gehalten werden und deßhalben auch das Regiment mit haben, reizen die Fürsten dazu, daß sie auf uns immer zorniger werden. Und damit dieselben nun immer weiter toben und wüthen mögen, dieweil sie mit großer Gewalt gerüstet, und dazu eine große Menge haben derer, so es mit ihnen halten, haben sie jetzt neulich auf dem Tridentinischen Concil Decrete und Satzungen gemacht, deren viele offenbar falsch und unrecht sind, aber noch viel mehr, die sophistisch sind und die man deuten kann, auf welche Seite man will; und damit sie dieselben bestätigen mögen, werden die Könige dazu sonderliche Gebote lassen ausgehen, mit Blut geschrieben.

Sie haben nun etliche, viele Jahre fast getracht und sich wider uns hoch gerühmt, als nämlich, daß sie den Namen der Kirche führen, sich auf die Concilien berufen und daß es die allgemeine Kirche einhellig mit ihnen halte; damit, sage ich, wollen sie uns gleich hart schrecken. Wiewohl nun dadurch viele Leute etwas bewogen werden, sollen doch die alle, so GOtt recht anrufen, diesen Betrug verstehen, und lernen, daß sie solche Gaukelei, damit die Widersacher umgehen, widerlegen können. Es ist zwar eine klare und gewisse Regel, welche diesen Betrug und Gaukelei gar umstößt, nämlich, da St. Paulus sagt Gal. 1, 9.: „So jemand euch Evangelium predigt anders, denn das ihr empfangen habt, der sei verflucht.“

\quellseite{7}
Es weist es die Sache an sich selbst aus, daß etliche Satzungen des Tridentinischen Concils dem Evangelium offenbar zuwider sind, dazu diesen Predigten Abrahams auch, darin je der Grund rechter Lehre gelegt ist, als da gesagt wird Röm. 4, 3.: „Abraham glaubete dem HErrn, und das rechnete er ihm zur Gerechtigkeit.“ Denselben Spruch zieht St. Paulus darum an, auf daß er damit anzeigen möchte, daß man also in der Kirche für und für einhellig vom Glauben gelehret hat. Da bringt er diese Predigt hervor, welche die erste ist, so ausdrücklich vom Glauben und Gerechtigkeit redet, und sagt klarlich, daß man die Verheißung mit dem Glauben soll annehmen; und wenn sich der Mensch mit diesem Glauben an den Mittler hält, so werde ihm das wahrhaftig zur Gerechtigkeit gerechnet, nämlich, daß ihm Vergebung der Sünden geschenkt wird und daß also der Mensch mit GOtt um des Mittlers willen versöhnt wird, und daß eben dieser Mittler in solchem Menschen auch kräftig sei und durch das Wort der Verheißung und mit seinem Heiligen Geiste ein neu Licht und Gehorsam anzünde, und da im Herzen Friede mache, welches der Anfang ist des ewigen Lebens.

Diese ganze Lehre, so eigentlich die rechte Lehre der Kirche und Gemeinde GOttes ist, welche auch die Zuflucht und höchste Trost ist der Gottseligen, und eben in diesem Buche und darnach auch in den Schriften der Propheten und Apostel ausdrücklich gelehrt und erklärt ist, unterdrücken die Tridentinischen Areopagiten ganz trotzig und freventlich mit einer solchen Lehre und Satzung, so dieser Lehre stracks zuwider und gottlos und verflucht ist. Denn sie lehren und gebieten, daß alle Menschen sollen im Zweifel bleiben; und tilgen damit also gar aus dieses ganze Hauptstück im Glauben, da wir sagen: „Ich glaube Vergebung der Sünden“; ja, sagen sie, ich glaube nicht, sondern zweifle noch immer daran.

Dieser öffentliche und schädliche Irrthum zeigt an, daß die, so solche Satzung gestiftet, Feinde sind der Kirche und Gemeinde GOttes, und daß sie diese Namen und Titel mißbrauchen, nämlich, da sie sich der Kirche, der Conciliën rühmen und daß sie es einhellig mit der allgemeinen Kirche halten; sie halten aber den Leuten die rechte Lehre der Kirche nicht vor, sondern bestätigen noch immer die Meinung, so die Heiden oder auch die Pharisäer gehabt. Denn die Herzen der Menschen, so vom Evangelium nichts wissen und das Gesetz etlichermaßen anschauen, werden in diesem Zweifel ersäuft, welcher ein elendes und betrübtes Toben wider GOtt und dazu auch Verzweiflung gebieret. Ich sage jetzt nicht von solchen Satzungen, so man auf beide Seiten deuten kann, deren sie viel gemacht haben, auf daß sie dazu ihre Zuflucht haben mögen, wenn sie etwas von uns gestraft werden, und bestätigen doch damit in Wahrheit nur ihre Mißbräuche; sondern ich rede jetzt von den öffentlichen und schädlichen Irrthümern. Und wo der Leser solche Irrthümer in ihren Decreten und Satzungen sieht, soll er der Regel folgen Gal. 1, 9.: „So jemand euch Evangelium predigt anders, denn das ihr empfangen habt, der sei verflucht.“

Hilf GOtt, wie sind so viel schädliche Verfälschungen in ihren Decreten und Satzungen von der Buße! Es ist nie in keinem Concil gesagt oder geschlossen worden, daß vermöge göttlichen Rechts vornehmen wäre, in der Beichte alle Sünden einzeln zu erzählen, und dasselbe darum, daß man den Leuten dafür eine Genugthuung oder Buße soll auflegen (wie sie dieselbe genannt haben). Nun haben aber die Gesetzgeber zu Trident noch unverschämt diese Lüge zu den alten Decreten gesetzt und schmähen damit also den Sohn GOttes. Denn da wird nicht bedacht, wie groß die Sünde sei, und wird auch nicht betrachtet, warum der Zorn und die Genugthuung auf den Sohn GOttes gewendet und gelegt sei, da die Menschen überschüttet werden mit Schauspielen der Genugthuung, deren etliche lächerlich, etliche auch schändlich gewesen sind. Endlich ist das Narrenwerk, so sie in dieser Sache getrieben, nicht eine geringe Unsinnigkeit gewesen. Und die Meister solcher Tridentinischer Satzung wissen zwar selbst wohl, daß die verworrene Lehre von der Genugthuung voller Lügen gewesen und noch jetzt falsch und lügenhaftig ist, und daß dieselbe Weise nun Alters halben gar erloschen und abgenommen ist. Und zeigen doch gleichwohl dem Volke den vergeblichen Schatten des Wortes von der Genugthuung, auf daß sie an ihrem großen Namen keinen Schaden oder Abbruch leiden mögen, wenn sie ihre Irrthümer bekennen würden; darnach, daß sie auch mit einem Scheine die Dispensation oder Linderung der alten Canones und Satzungen verkaufen können, welche sie Ablaß nennen.

Vom Opfer der Messe erneuern sie das Narrenwerk der Pharisäer, nämlich, daß man mit demselben Opfer verdienen soll Vergebung der Sünden, der Schuld und Pein denn, der das Opfer thut, und auch andern Leuten. Derhalben vertilgen sie damit das Wort des Evangeliums von Vergebung der Sünden, die uns aus Gnaden um des Sohnes GOttes willen widerfährt und mit dem Glauben muß angenommen werden; ja, sie tilgen damit den Spruch aus, so in der Epistel an die Hebräer am 10. Cap. V. 14. steht: „Mit Einem Opfer hat er in Ewigkeit vollendet, die geheiliget werden.“

Ich will aber jetzt davon nicht mehr sagen; denn es ist dasselbe Gesetzbuch dermaßen noch nicht ganz ausgegangen, und wenn es auch wird an den Tag kommen, so werden beide, jetzt und auch hernach bei den Nachkommen, viel fromme gelehrte Männer sein und die GOtt recht anrufen, welche die Irrthümer des Tridentinischen Concils widerlegen werden und der ganzen Welt anzeigen die heimliche List und Sophisterei in solchen ihren Satzungen, und dazu auch die verkehrte Bosheit derjenigen, so tückischer Weise viel ungewisses Dinges, daß man auf beide Seiten deuten möchte, unter die Satzungen gemacht haben, daß sie mit solchem Schein zugleich das gemeine Volk betrügen, und der rechten Strafe, so sie dagegen hören und leiden werden, entfliehen mögen.

Auf daß aber fromme Herzen vor den greulichen Irrthümern des Tridentinischen Concils, und andern dergleichen, so noch immer von den Papisten vertheidigt, bewahret mögen werden und unversehrt bleiben, sollen sie die Bücher Luthers fleißig lesen; unter welchen diese Auslegungen des ersten Buchs Mose die allerlieblichsten und herrlichsten sind, und wie man im Sprüchwort sagt, Schwanengesang \&c. Dieselben Auslegungen soll man thun zu andern seinen Bekenntnissen, so in andern seinen Büchern erzählt werden von allen Hauptstücken christlicher Lehre, und solcher Bekenntnisse Auslegung ist in diesen längeren Predigten auch enthalten.

Diese Zeugnisse sind nütze, fromme Herzen damit zu lehren und zu befestigen. Dieselben sollen wir, als eine gute Beilage (wie Paulus 2 Tim. 1, 12. 14. seine Zeugnisse nennt) treulich bewahren, und nicht zulassen, daß sie entweder mit falschen und sophistischen Decreten des Papstes unterdrückt, oder aber mit List verkehrt und durch hinterlistige Auslegung in fremden Verstand verändert werden.

In diesem äußersten Alter und Ende der Welt werden mehr seltsame verworrene Meinungen sein, denn zuvor je gewesen sind, und wird die Leichtfertigkeit und Unbeständigkeit der Leute noch größer sein. Viel Menschen werden von Natur ehrgeizig und aufrührerisch sein, und ist auch das Toben des Teufels etwas größer. Denn dieweil er wohl weiß, daß der Tag nicht weit ist, an welchem sich der Sohn GOttes von dem unmenschlichen Geschlechte wiederum wird sehen lassen, einen Triumph führen und seine Kirche mit ewiger Herrlichkeit schmücken, und an welchem auch die große Sünde und Schande der Teufel wird gesehen und ihre Pein und Marter auch gemehrt werden: sind sie nun so viel desto unsinniger, auf den Sohn GOttes ihren Zorn auszuschütten, und toben und wüthen derhalben jetzt noch mehr wider die arme Kirche, die nun Alters halben etwas schwach ist.

Derhalben soll uns diese große Gefahr aufwecken und munter machen, daß wir uns soviel desto mehr befleißigen, über der reinen Lehre zu halten und darin die Einhelligkeit der allgemeinen Kirche bewahren, und dabei beständig bleiben. Dazu ist nun vornöthen, daß man in der Schrift fleißig lese und mit heißem Gebet bei dem lieben gnädigen GOtt und Vater unseres HErrn Jesu Christi immer anhalte, wie der Sohn GOttes sagt Matth. 26, 41.: „Wachet und betet, daß ihr nicht in Anfechtung fallet.“ Daß er sagt, daß man wachen soll, darin wird begriffen, daß man der Lehre fleißig soll nachdenken; mit dem Gebete aber bittet man Gott um Hülfe; denn bei menschlichem Fleiße ist kein Glück, wo er nicht von Gott regiert, und durch seinen Segen fortgesetzt und gefördert wird.

\quellseite{8}
Wie nun der Herr Jesu Christi, unser Hoherpriester, der für uns gekreuzigt und wiederum von den Todten auferstanden ist, in seinem Todeskampf gebeten, daß der ewige Vater seine Kirche erhalten und sie in seiner Wahrheit heiligen wolle, und machen, daß sie alle in ihm eines seien: also thut er noch immer dieses Gebet für uns in dem himmlischen Rathe der göttlichen Majestät, und ist gewiß, daß sein Gebet nicht vergeblich oder umsonst ist.

Zu diesem Gebet unsers Hohenpriesters sollen wir auch unser Seufzen, Thränen und Gebet thun, daß der ewige Vater in dieser elenden betrübten Zeit seinem Sohn für und für auch in diesen Landen eine Kirche und Gemeinde sammeln und uns in seiner Wahrheit heiligen wolle, und machen, daß wir alle in ihm eines seien; dazu auch die Herberge seiner Kirche schützen und bewahren wolle, und nicht zulassen, daß heidnische und scheußliche Verwüstungen geschehen oder einreißen mögen.

Dieses bitte ich mit herzlichem Seufzen, daß uns der ewige Vater unsers Herrn Jesu Christi um der Fürbitte willen seines lieben Sohnes gnädiglich verleihen wolle.

Datum, Nürnberg, am Tage von Pauli Bekehrung im Jahre nach Christi, unsers Herrn, Geburt 1552.

\chapter[Hieronymus Besolds Vorrede]{Hieronymus Besolds Vorrede über den letzten Theil des ersten Buchs Mose, so vom 37. Capitel angeht}

Dem frommen und getreuen Leser Gnade und Friede durch Christum.

Dieweil ich gewußt, daß viel gutherziger Leute schon vorlängst heftig begehrt, daß dieser letzte Theil der Auslegung der Genesis auch im Drucke möchte ausgehen, habe ich es nicht allein unnöthig geachtet, eine lange Vorrede zu stellen, sondern habe auch gesehen, daß es den Gottseligen einen Verzug und Verdruß bringen würde, welche eilen, daß sie die Auslegung selbst je eher je lieber lesen mögen, so ihrer Würde und Nutzens halben sehr hoch berühmt ist, und dazu diejenigen, so diesem Schreiber und seinen Schriften etwas zuwider sind, wegen der wunderbaren Wichtigkeit und Heiligkeit der Dinge, davon hierin gehandelt wird, zu sich rezet. Vornehmlich aber habe ich mich besorgt, ich möchte vielleicht mit dieser meiner Gutwilligkeit keinen Beifall bei denjenigen, welchen mein geringer Verstand wohl bekannt ist, finden, daß ich entweder mit langer Vorrede, oder aber daß ich diese Auslegungen geschrieben, solche Arbeit mir habe lassen auflegen. Denn ich gedenke noch wohl, daß Herr Veit Dietrich, ein trefflicher, gottseliger und gelehrter Mann, über diese schwere Arbeit geklagt und gewünscht hat, daß diese Auslegung hätte mögen von D. Luther, unserm lieben Präceptor, selbst geschrieben und in den Druck gegeben werden, welche er in der Schule zu Wittenberg mit großem Ernste und wohlgelehrten Worten, und dazu mit reicher Materie und wichtigen Dingen, davon er gehandelt, gelehret hat; dieweil es fast schwer wäre, alle Worte unter dem Lesen, wie behende Einer auch schreiben kann, zu erreichen und aufzuschreiben, oder aber auch in so feine gleiche Ordnung zu bringen, wie von ihm, dem Herrn D. Luther selbst, hätte mögen geschrieben und geredet werden. Nun wissen aber alle Gottseligen wohl, daß an Herrn Veit sonderliche Geschicklichkeit gewesen, in diesem Theile der guten Künste zu arbeiten, und bezeugen solches auch die Schriften D. Luthers, seliger Gedächtniß, so er im Druck hat lassen ausgehen, den Nachkommen damit zu dienen. Und ich habe gesehen, daß er nicht allein aus dem, so er selbst, sondern auch, so D. Creutziger und M. Georg Rörer aufgeschrieben gehabt, mit großer Geschicklichkeit viel anderes Dinges und auch den ersten Theil dieser Auslegung der Genesis geschrieben hat. Darum ich auch von wegen der Schwachheit meiner Kräfte so viel desto mehr besorgt geworden, und hätte mich auch solche Schwachheit und Unvermögen gänzlich dahin bewegen, dieses vorgenommene Werk liegen zu lassen, wo ich es mit deren Willen und Erlaubniß hätte thun mögen, die mich mit ihrer Gottseligkeit und großem Ansehen dazu gebracht haben, daß ich diese Last und Arbeit auf mich genommen.

Derhalben habe ich gedacht, ich müßte kürzlich die Ursache dieser Arbeit erzählen, nicht allein daß ein gottseliger und getreuer Leser dieselbe möchte wissen, sondern darum am allermeisten, daß die Nachkommen dessen möchten vergewissert werden, daß diese Auslegung treulich beschrieben worden ist. Denn nachdem Herr Veit Dietrich von diesem Leben abgeschieden, im Jahre nach Christi, unsers Herrn, Geburt 1549, eben in der Nacht, so vor dem Tage, an welchem in der christlichen Kirche das Fest der Verkündigung Mariä oder Empfängniß Christi gehalten wird, hergeht, an welchem Tage auch Isaac geopfert, und Christus den Tod gelitten hat, wie die Alten dafür gehalten haben, ist nicht allein unsere Kirche, sondern auch diese Auslegung ihres Lehrers und getreuen Merkgutes beraubt worden, welche er, so viel ihm seiner Krankheit halben immer möglich gewesen, täglich in Händen gehabt, als seinen einigen Trost und Linderung so viel großer Schmerzen, die er erlitten hat; also daß er kurz vor seinem Abschiede aus diesem Leben, da er die Feder nicht mehr in der Hand hat halten können, gleichwohl die angefangene Arbeit gefördert und einen Andern in die Feder zugelesen hat, und gehofft, er wollte die Historie von Abraham ausmachen und zu Ende bringen, welchen er dazumal (wie er selbst pflegte zu sagen) an den Berg Morija gebracht hatte, daß er mit derselben sein Leben enden und mit Abraham möchte begraben werden. Und solche seine Hoffnung und Verlangen, daß er aus diesem Leben möchte erlöset werden, ist nicht vergeblich gewesen. Denn ehe er denselbigen Theil vom Berge Morija zum Ende gebracht, dahin er in seinem Herzen pflegte wallen zu gehen, ist er zu der ewigen Gemeinschaft Christi, dazu der lieben Patriarchen und Propheten berufen und gefördert worden, da er die Lehre, so aus derselben Schriften gezogen, beide mündlich und schriftlich der Kirche und Gemeinde Gottes mit großer Treue und Fleiß vorgehalten. Welche Wohltat ich und alle Gottseligen mit der Gnade Christi in alle Ewigkeit mit dankbarem Herzen rühmen wollen. Und fällt mir jetzt ein, daß ich nicht ohne großen Schmerz gedenke beide an die Tugenden und Gaben, damit er begnadet gewesen, der Kirche in diesem unserm Vaterlande zu dienen und zu helfen, und auch die guten Künste und die darin studirt, zu befördern und fortzusetzen, und dazu auch an die Gefahr und erbärmliche Verwirrung, so noch bei seinem Leben über die Kirche gekommen ist, und nach seinem Absterben allemal mehr gemehret und größer geworden ist.

Es tobte und wüthete zu der Zeit über ganz Deutschland die greuliche Gewalt der Finsterniß, so vom höllischen Antichrist erregt war, daß Licht des Evangeliums, so uns von Gott erzeigt, auszulöschen, und die Kirchen hin und wieder zu unterdrücken, dazu auch derselbigen Herberge, darin die Stimme des Sohnes Gottes durch fromme gottselige Lehrer und Kirchendiener gelehret worden. Und in dem großen Schrecken, da fromme Herzen also bestürzt gewesen, dieweil alle Dinge sich so gar ungewöhnlich und jämmerlich haben angelassen, ist diese Furcht auch dazu gekommen, daß die Diener des Antichrists nach diesen Schriften forschten, und die Macht und Freiheit, solche Schriften in Druck zu geben, so den Kirchen zuvor von gottseeligen Regenten verliehen, wiederum nehmen möchten; welches mir denn etlichemal gehört, daß sie damit umgingen und uns damit gedroht hätten. Derhalben, da die Freunde und Erben des Herrn Veit in der Furcht gestanden und darum bestürzt gewesen, haben sie dieser Auslegung einen sicheren Ort und Behältnis gesucht, nämlich, bei dem vornehmen Manne, Herrn Hieronymus Baumgärtner, da dieselbe möchte verborgen liegen; wie die hundert Propheten bei Obadja verborgen gelegen. 1 Kön. 18, 4.

\quellseite{9}
Wiewohl aber nun in solchem großen Schrecken und Bestürzung zumal kleine Hoffnung war, daß wir wiederum zu voriger Freiheit kommen möchten, wo nicht Gott selbst den Feinden Christi wehren würde, hat doch gleichwohl das Seufzen der Gottseeligen, so in derselben betrübten Zeit immer groß Verlangen gehabt nach der letzten und edelsten Arbeit ihres lieben Propheten und Pfarrherrn D. Luthers seligen, diesen Commentar oder Auslegung nicht länger also verborgen liegen lassen. Und er, Hieronymus, hat selbst großen Fleiß darauf gewendet, damit je zum wenigsten dasjenige im Drucke möchte ausgehen, was er, Veit, geschrieben hatte, da die Widersacher noch immer getobt und sich mancherlei unterstanden haben. Und dieweil er leichtlich konnte abnehmen, daß jedermann gerne sehen würde, daß die ganze Historie möchte im Drucke ausgehen, hat er von mir begehrt, daß ich den übrigen Theil vollends wollte ausmachen, als der ich in der Schule zu Wittenberga D. Luther auch gehört, da er dasselbe gelesen und gelehrt hat. Und gefiel solches dem Herrn Philippus und M. Rörer auch wohl, und M. Stolz gab mir auch, was er aufgeschrieben hatte: welches denn sehr wohl zu lesen ist. Dazu hatte ich selbst auch viel aufgeschrieben aus D. Luthers Munde, was er gelesen hatte; welches mir denn sehr wohl geholfen, da ich an die fremde Hand oder Schrift noch nicht gewöhnt, und mir sehr lieb, daß alle Exemplare so fein übereinstimmt und einander so gleich waren.

Dieweil ich nun denen, so mir diese Arbeit auflegten, vieler großen Wohltaten halben verpflichtet und sie in mir lieb gewesen; auch durch ihre und Herrn Beurs Einnehmung und Angeben dazu gekommen war, daß ich D. Luthers Hausgenosse geworden, hat es mir mit nichten gezimeten wollen, solchen großen aufsehnlichen Leuten dieses zu versagen. Und war gänzlich der Hoffnung, der fromme getreue Leser würde es zum allerbesten deuten und aufnehmen, wenn er verstehen würde, daß ich weder aus meinem eigenen Gutdünken oder eigenem Willen, sondern aus schuldiger Pflicht und einfältigem guttem Herzen diese Last auf mich geladen, nur daß diese Auslegung also möchte bewahret und erhalten werden.

Da aber die Historie Abrahams war ausgegangen, habe ich die andern auch vorgenommen, eben auf Vermahnung und Anhalten derselbigen Männer, und habe in etwas eine Erleichterung empfunden der verdrießlichen Arbeit des Abschreibens, von wegen der Größe und Nutzen der guten Dinge, Materie und Exempel, welche vornehmlich D. Luther in seinen letzten Worten der Kirche hat wollen vorhalten und befehlen, auf daß damit den Lesern der Weg geöffnet würde, die Bibel so viel desto fleißiger und besser zu verstehen. Denn dieses sind wahrhaftig die rechten Brunnen Israel, daraus alle Propheten, Apostel und frommen gottseligen Lehrer zu allen Zeiten alle himmlische Weisheit geschöpft haben. Welche Brunnen, da sie mit Dreck und Unflat verflossener Zeit überschüttet gewesen, ist es eine große Wohltat, daß Gott den lieben Luther erweckt hat, dieselben wiederum zu reinigen, und die ältesten Schriften der ersten Väter zu erläutern, welchen anfänglich die Verheißung vom Samen des Weibes und dem Sohne Gottes, der da sollte Mensch worden, gegeben, und immer bei ihnen ist wiederholt worden. Wie aber die Propheten und Apostel auf dieselben Verheißungen fleißig gesehen, also ist sie unter dem Papste, dem Antichrist, gar verachtet und verdunkelt worden durch mancherlei Abgötterei, Aberglauben, und die verworrenen Menschen-Satzungen von Klostergelübden, vom Cölibat oder ehelosem Leben der Mönche und Pfaffen etc.

Und ich gedenke noch wohl, daß D. Luther etlichemal in seinem Hause über Tisch, da ich sein Tischgeselle war, mit Seufzen und Schmerzen geredet hat von der Blindheit und Finsternis seiner Zeit, und sagte unter andern, wie die ehelosen Päbstler das Buch der Schöpfung sehr verachtet hätten, und vor demselben eine Scheu gehabt, als vor einem solchen Buche, das voller Unflaths und Unzucht wäre, darum daß es Anderes nichts lehrete, denn nur, wie die Patriarchen Kinder gezeuget und der Haushaltung gewartet hätten. Er setzte aber eine treffliche Empfehlung und schönes Lob dieses Buchs dazu, nämlich, vom Segen des Wortes und Stimme Gottes, von welcher man oftmals hört, daß sie zu den Vätern ist geredet worden. Das, sagt er, macht Genesis köstlich: „Der Herr hat geredet“; „Der Herr hat gesagt“; „der Herr ist erschienen“. Das haben die Mönche nicht geachtet; denn sie haben Gottes Wort nicht gehabt. Das hat kein Heide. So haben es die Jura nicht, Medicina auch nicht. Die ganze Bibel hebet sich also an: Dixit Dominus: „Der Herr hat gesagt“; das Neue Testament auch: „Im Anfang war das Wort“, Joh. 1, 1. Wer das Dixit verliert, der ist verloren. Sie haben das Dixit (Er hat gesagt) nicht, das ist, Vater (der Vater), so haben sie das Fiat (Es werde) nicht, das ist, Filius (der Sohn), das Facta sunt (Es ist geworden) haben sie, wie Paulus sagt, daß der Schöpfer durch die Creatur erkannt werde a posteriori, wie man sagt; aber wir Theologen haben es a priori, wir haben das Dixit und das Fiat.

Und eben dasselbe hat er auch in der Auslegung der Genesis immer fleißig getrieben. Ich habe aber dieses gleichwohl auch erzählen wollen, Andere dadurch zu erwecken, daß sie Lust und Liebe gewinnen möchten zu den Historien in der heiligen Schrift und daß sie sich derselben lernen verwundern: welche zwar nicht also herrlich scheinen vor den Augen der Menschen, wie die Geschichte der Heiden, oder auch der Mönche und Heuchler abergläubiger Gottesdienst thut; läßt sich auch ansehen, daß sie mit feiner zierlichen Rede geschmückt seien. Es sind aber doch die Worte und heiligen Dinge, davon gehandelt wird, sehr wichtig; wie denn solches dieser unser Ausleger mit großem Fleiße erklärt hat, und wird auch dasselbe in dem wahrhaftigen und geistlichen Kampf allererst recht gespüret und verstanden.

Er hat auch oft in gemeinen Gesprächen eine Vergleichung gemacht der Schriften der Propheten und Apostel, welche von Mose gleich als aus einem Brunnen wären hergeflossen, und zeigte nicht allein an, wie sich die Dinge, davon sie schreiben, mit einander reimeten, sondern auch in etlichen den Unterschied der Reden. Von Mose sagte er, daß er wohl und ordentlich hätte lehren können; von Jesaia, daß er große, prächtige und zierliche Worte hätte gebraucht, wie Cicero zierlich geredet hat. Davids Rede, sagt er, wäre gleich wie des Sallust Rede ist, das ist, fein kurz und reinlich. Aber Salomons Rede wäre geschmückt; wiewohl auch Davids Rede wunderbarlich und sehr reich wäre, sonderlich im 19. Psalm, wie er ihn pflegte zu rühmen. Endlich gab er ihnen aus dem Neuen Testamente zu Paulus und Johannes als vortreffliche köstliche Ausleger des Alten Testaments, darum daß dieselben viel tiefer und fleißiger Mose und die Propheten besehen hätten. Und ermahnte uns dabei, daß wir die hebräische und griechische Sprache ja fleißig studiren sollten, auf daß wir so viel desto leichter die Sprüche der Propheten und Apostel mit einander vergleichen und verstehen könnten. Und dieweil ich es dafür halte, daß solche Vermahnung denen, so in der heiligen Schrift studiren, nicht allein lieb und angenehm, sondern auch nöthig und nütze sein würde, will ich sie eben mit denselbigen Worten erzählen, damit er sie ganz freundlich und herzlich geredet hat. Und ich zweifle nicht, der ehrwürdige Herr D. Justus Jonas, der bei dieser Rede auch mit gewesen und dazu Ursach gegeben hat, habe dieselben noch im frischen Gedächtnis behalten. Denn er war ungefähr daselbst einmal aus den sächsischen Salzgruben (von Halle) spazieren gezogen zu D. Luther, und auch die andern Herren Prädicanten zu besuchen; und da er sie in der Schule von mancherlei Dingen künstlich und schriftlich hat hören disputiren, erzählte er über Tische unter dem Essen, wie er von Herrn Philippus gehöret, daß die ganze Schrift Anderes nichts wäre, denn ein Streit oder Kampf zwischen der Schlange und dem Samen des Weibes. Dazu sagte D. Luther also: Glaubet ihr auch, daß Johannes ein Commentar und Auslegung sei über die ganze Bibel? Paulus auch? Es ist kein Wort, Johannes wollte gerne Christum zum Gott machen. Die heilige Schrift bringt viel mehr auf den Sohn als auf den Vater; denn die ganze Schrift ist um des Sohnes willen geschrieben: darum sind auch im Alten Testamente mehr Sprüche oder Zeugnisse vom Sohne als vom Vater.

Unter andern aber ward auch gedacht der Dolmetschung der siebenzig Dolmetscher, davon sagte er damals also: Paulus hätte an vielen Orten dieselbe Dolmetschung behalten. Er hat es nicht verachtet; denn er hat mit den Griechen geredet. Im 19. Psalm V. 5. steht auch denn hebräischen Texte: In omnem terram exivit regula eorum, das ist: „Ihre Schnur gehet aus in alle Lande“, das nennt Paulus Röm. 10, 18.: Sonum, „einen Schall“, und sagt gleich auch daselbst. Also sagt er 1 Cor. 15, 55.: Absorpta est mors in victoriam, das ist: „Der Tod ist verschlungen in den Sieg“; da lautet es aus dem hebräischen Texte auf Lateinisch, in finem, oder in aeternum, etc. Und Jes. 25, V. 8.: Praecipitabit mortem in aeternum: „Der Tod wird gar verschlungen werden“; er wird nicht wiederkommen in victoriam, id est, vita vincet, das Leben wird siegen. Paulus ist reich von Worten. Ein Wort Pauli hat drei Reden Cicerons. Er redet ein Wort, das steht durch einen ganzen Jesaia oder Jeremia. O, es ist ein feiner Prediger! Er heißt nicht vergebens ein auserwähltes Rüstzeug. Unser Herr GOtt sagt: Ich will der Welt einen Prediger geben, der soll köstlich sein. Es ist keiner, der das Alte Testament so wohl versteht als er; Johannes nehme ich aus. (Petrus ist auch köstlich). Matthäus und die Andern schreiben ja wohl die Historien, und die sind hoch vornöthig: aber die Kraft und Worte des Alten Testaments drücken sie so nicht aus. Paulus hat viel Dings aus dem Hebräischen Griechisch übersetzt, das keiner sonst thun könnte. Er redet in einem Capitel, das oft vier oder fünf Propheten auslegt. O, er hat Jesaia und Mose lieb gehabt! Die sind es auch. Die Worte und Materie, davon Paulus handelt, die sind in den Propheten und Mose. Darum sollen die jungen Theologen studiren Hebräisch, auf daß sie das Griechische und Hebräische mit einander vergleichen können. Mose sagt von Abraham: Et imputatum est ei ad justitiam: „Und es ward ihm gerechnet zur Gerechtigkeit“ etc. Das hat sich Paulus so müße gemacht, und steht doch im Hebräischen nur: Et cogitatum est. Wenn ich jung wäre und wollte nun ein hoher Theologe werden, so wollte ich Paulum vergleichen mit dem Alten Testamente. Er ist ein Dialektiker und köstlicher Rhetor oder Redner gewesen.

\quellseite{10}
Dieses hat mich gut gedäucht allhier zu erzählen, dieweil ich ja etwas habe sagen müssen; denn es reimt sich daher nicht übel: und habe gehofft, es würde denen, so je D. Luther selbst gehört haben, auch Lust bringen und ihnen wohlgefallen. Denn ich kann mich zwar niegendes mit so wohl zufrieden geben in diesem jetzigen grimmigen Zorn des Teufels, damit er jetzt so greulich tobet und wüthet wider die Kirche und freie Städte oder Regimenter, welche wohl bestellt und ordentlich gefaßt sind, als daß ich daran denke, was ich von D. Luther, beide öffentlich und insonderheit, gehört habe; damit ich mich aufrechtshalte, so wohl als ich kann, dulde und ertrage so viel desto leichter beide gemeinen Jammer und Unglück: und habe mich bisher auch dazu beflissen, mich mit so vieler Sorgfältigkeit und größerem Fleiße zu hüten vor Verfälschung der Lehre, so der liebe Luther gelehret und hinter sich gelassen hat. Denn unter andern hat er davon auch Erinnerung gethan, daß greuliche Empörungen und Strafen über Deutschland kommen würden, nicht allein von den Feinden des Evangeliums, sondern auch von wegen der Grimmigkeit etlicher scheußlicher und wunderlicher Menschen, welche mit allerlei Vorwendung eines guten Scheines der Freiheit des Vaterlandes heimlich würden nachstellen. Und es ist auch zwar also ergangen. Denn wir haben ja gehört, daß man den Namen und Lehre des Evangeliums zum Scheine vorgewendet, da man unbillige Grausamkeit, Mord, Brennen und Räuberei getrieben hat. Man weiß aber auch wohl, was die Kirche und gottselige Leute für schwere Kämpfe gehalten mit denen, von welchen er zuvor gesagt, daß Etliche würden neue Künste und Betrügereien erdenken, die alte Abgötterei damit wiederum hereinzuführen. Die Andern aber würden aus Vorwitz und Ehrgeiz, und daß sie gerne etwas Neues wollten anrichten, Alles verwirren in den Kirchen, die recht und wohl angerichtet und in gute Ordnung gebracht sind. Und dieselben würden mit großem Ansehen und Namen großer Kunst umgeben sein, und würden sich also unter solchem Namen und erdichteter Hülfe bei vorwitzigen und unbeständigen Menschen rühmen und angeben.

Und daß ich anderer Schäden geschweige, damit jetzt etliche Jahre die Kirche und gemeine Regiment fast hart angefochten und geschwächt worden, so hat uns wahrlich viel Mühe und Zanks geboren die neue Weise zu reden, so Flandern in der Lehre, von der Justification oder Rechtfertigung eingeführt, welcher beide mit Verfälschung vieler Sprüche in den Schriftteil der Propheten und Apostel, und auch etlicher Zeugnisse aus D. Luthers Schriften, die er fälschlich auf seinen Irrthum gezogen, viel Leute betrogen hat. Denn er nimmt das Wort von der wesentlichen Gerechtigkeit aus dem Buche, so im Anfange des Evangeliums D. Luther hat lassen ausgehen von dreierlei Gerechtigkeit: mit welchem Worte er die Würdigkeit der Gerechtigkeit des Glaubens hat erheben und rühmen wollen wider die Gerechtigkeit menschlicher Werke. Wie denn die Auslegung, so dabei gesetzt ist, klarlich anzeigt, nämlich, daß D. Luther durch die wesentliche Gerechtigkeit nicht verstanden habe die Gerechtigkeit Gottes des Vaters, des Sohnes und Heiligen Geistes, oder allein die göttliche Natur in Christo, wie sich die Flandern träumen läßt: sondern alles dasjenige, was Christus für uns gethan und gelitten hat; wie er unter andern Zeugnissen der Schrift den Spruch Pauli anziehet, Röm. 5, 18. 19.: „Wie durch Eines Sünde die Verdammniß über alle Menschen kommen ist; also ist auch durch Eines Gerechtigkeit die Rechtfertigung des Lebens über alle Menschen kommen“, und „durch Eines Gehorsam werden viel Gerechte.“

In der andern Schrift aber, von zweierlei Gerechtigkeit, so daselbst angegeben (welches meines Erachtens sonderlich also von GOtt zugeschickt ist, dem Irrthum zu begegnen und damit dasselbe Wort nicht so fälschlich verstanden werden möchte), hat er mit klaren Worten eine Definition gesetzt, was die wesentliche Gerechtigkeit sei, welche er daselbst also nennt. Denn also lauten die Worte D. Luthers: So wird nun diese Gerechtigkeit den Menschen gegeben in der Taufe, und zu aller Zeit, wenn man rechte Buße thut, also daß der Mensch sich mit aller Zuversicht in Christo rühmen kann und sagen: Das Christus gelebt, gethan und gesagt, daß er gelitten hat und gestorben ist, das ist alles mein, nicht anders, denn als ob ich selbst also gelebt, gethan und gesagt, gelitten hätte und gestorben wäre. Und bald darnach: Daher auch Christus selbst, der da sagt, er sei gekommen, daß er diesen gnädigen Willen seines Vaters thun wollte, ihm ist gehorsam worden, und alles, was er gethan, hat er uns gethan, und hat wollen haben, daß es sollte unser sein, da er sagte, Luc. 22, 27.: „Ich bin mitten unter euch, als Einer, der da dienet.“ Und abermals in den Worten des Abendmahls sagt er, V. 19.: „Das ist mein Leib, der für euch gegeben wird.“ Und Jesaia sagt Cap. 43, 24.: „Ja, mir hast du Arbeit gemacht in deinen Sünden, und hast mir Mühe gemacht in deinen Missethaten“ etc.

Nun zweifele ich gar nicht, daß dieses alle verständigen und gottseligen Menschen wohl wissen, und die Sprüche der Propheten und Apostel, keinen sich auch wohl dazu; und in Luthers Büchern allen wird eben dasselbe mehr als tausendmal wiederholt und jedermann eingebildet. Und dann solches auch nicht geleugnet werden von denen, so der neuen Flandrischen Lehre (welche wenig Leute recht verstehen) anhängig sind; und schämen sich doch nicht, den ansehnlichen Namen Luthers vorzuwenden, ihren Irrthum damit zu beschönigen. Wo sie aber etwa in den Schriften Luthers an solche Oerter kommen, welche so klar sind, daß sie damit leichtlich können überwiesen werden, daselbst dichten sie, es sei durch irgend eine Figur oder Gleichniß geredet; erdenken etliche Figuren und Tropen (sonderliche Arten zu reden), gleich als ob D. Luther dazu Lust gehabt, daß er in diesem vornehmsten Hauptstücke unseres Heils und Seligkeit mit Figuren und verdunkelten Worten also hätten spielen und Scherz treiben wollen. Dieser Mann GOttes hat wohl verstanden aus so vielen schweren Anfechtungen und geistlichen Kämpfen, mit welchem Trost das menschliche Herz, so aus Furcht des Gerichts und Zorns Gottes bestürzt und erschrocken ist, am allerbesten kann zufrieden gestellt werden, nämlich, mit dem Kreuze, Tode und Blute des Sohnes Gottes, unsers Mittlers Jesu Christi; wie die Propheten, Johannes der Täufer, Christus selbst und die Apostel uns auf die Erhöhung und sein Opfer am Kreuze weisen, und das Blut des Bundes, dadurch die Gefangenen aus der Grube gelassen werden, da kein Wasser innen ist, Sach. 9, 11. Ja, sie zeigen uns seine Wunden und Striemen, und sagen: Das sei unsere Gerechtigkeit vor GOtt; wie Paulus klarlich sagt Röm. 5, 9.: „Nachdem wir durch sein Blut gerecht worden sind“ etc. Daselbst legt er aus, was Zacharja mit dem Auslassen der Gefangenen aus der Grube gemeint habe, oder wie an andern Orten die Schrift redet, erlösen und erretten von der Gewalt des Satans. Das,
\quellseite{11}
sage ich, legt Paulus aus durch das Wort, gerecht werden, auf daß er damit anzeige, daß zwischen denselben Worten kein Unterschied sei. Und Röm. 3, 25. sagt er: „Welchen Gott hat vorgestellt zu einem Gnadenstuhl, durch den Glauben in seinem Blut.“ Item Joh. 16, 8. 10.: „Der Heilige Geist wird die Welt strafen um die Gerechtigkeit, denn ich gehe zum Vater.“ Welches man alles fleißig merken soll ums des geringen und kindischen Gespötts oder Betrugs willen, so sie treiben mit der Präposition per (durch), mit welcher Mißhandler dergleichen Sprüche Pauli verkehrt und verwirret. Unter vielen andern Zeugnissen aber Luthers sind sehr schöne und merkliche Sprüche in der letzten Lection dieser Auslegung, nämlich, vom Tode Christi, wider das Wortgezänke, dadurch die Erlösung und Versöhnung abgesondert wird von der Rechtfertigung und christlicher Gerechtigkeit, als, mit dem Unterschied der Jahre oder der Zeit. Denn wiewohl Christus einmal und zu bestimmter Zeit gelitten und so das Werk der Erlösung ausgerichtet hat, so gilt doch gleichwohl die Kraft und Wirkung des Leidens und Erlösung zu allen Zeiten, und dient zur Gerechtigkeit und Seligkeit derer, die dasselbe in einem jeden Augenblicke durch den Glauben an sich bringen. Und der Glaube fragt nichts nach der arithmetischen Jahrrechnung, oder nach irgend einer Zeit, wie lange oder kurz dieselbe auch sein mag; sondern sagt mit D. Luther also: Ich weiß, daß Christus nicht allein nach der Historie gestorben ist, sondern daß sein Tod währet vom Anfang bis an das Ende der Welt, und daß er zu Hülfe kommt allen Heiligen durch die ganze Zeit der Welt, dem Adam, der Eva und Andern mehr. So neu ist mir jetzt Christus, als hätte er sein Blut diese Stunde vergossen. Und wie hätte er doch klärlicher oder eigentlicher reden können? Und wir haben gehört, daß er dieses also geredet nicht ohne sonderlichen Eifer und hitzigen Geist. Denn eben zu derselben Zeit hat er sich mit diesen Worten aufgerichtet, getröstet und gestärkt, da er von ganzem Herzen begehrte aus diesem Leben abzuscheiden; welches denn die letzte Clausel, damit er dieses Buch Genesis beschließen, angezeigt. Denn er hat oft gesagt, daß er mit diesem Buche Genesis sein Leben begehrete zu beschließen, und hat Gott zwar sein Begier und Wunsch auch erfüllt. Und da er von dem Tage, darin er Genesium vollendet, nicht mehr ist in die Schule gekommen und des nächstfolgenden Jahres 1546, den 18. Tag Februarii zum ewigen Leben gefordert worden, hat er etlichemal, da er in den letzten Zügen gelegen, denselben Trost wiederholt, so uns in den Worten Christi vorgehalten wird, Joh. 3, 16.: „Also hat Gott die Welt geliebet, daß er seinen eingebornen Sohn gab, auf daß alle, die an ihn glauben, nicht verloren werden, sondern das ewige Leben haben.“ Von welchem Hingeben des Sohnes Gottes er gewiß dafür gehalten, daß es die Erlösung sei vom ewigen Tode, und dadurch wir absolvirt und losgesprochen werden vor Gottes Gericht; ja, er hat es für seine Gerechtigkeit gehalten, damit ihn Gott bekleidet und begnadet hätte. Die Schwärmer ziehen den Spruch Salomos an, da er sagt, Spr. 13, 2.: „Gerechtigkeit erlöset vom Tode.“ Dieweil aber die Propheten und Apostel so oft dieselbe Kraft dem Blute und Gehorsam Christi zuschreiben, wie sind denn die Schwärmer so unverschämt, daß sie dem Blute und Gehorsam Christi diese Ehre und Namen der Gerechtigkeit nehmen dürfen? Aber wollte Gott, daß sie solches in geistlichen Anfechtungen mit der That erfahren und lernen möchten, welche mit Worten auf den rechten Weg nicht wiederum mögen gebracht werden.

Was sie aber in die Länge davon disputiren, habe ich kürzlich wollen anziehen, auf daß ich damit bei frommen und gutherzigen Leuten anzeigte, daß ich nicht ohngefähr oder ohne Ursache, sondern durch solche heilsame Vermahnung D. Luthers, meines lieben Präceptors, bewogen worden, von denen abzuweichen, welchen ich doch naher Freundschaft halben verwandt gewesen, und daß ich damit den jämmerlichen Schaden unserer Kirche beklagen möchte, welche mit dieser Seuche und Irrthum jämmerlich vergiftet ist; nicht so sehr um der Widersacher willen, als auch von wegen der verkehrten Heuchelei etlicher Heuchler, damit sie auch anfänglich herein geschlichen und unter die Zahl der rechten Lehrer gekommen sind, und eben mit solcher Heuchelei jetzt auch die Kirche und Gemeinde Gottes betrügen, und was Andere recht gesagt und gelehrt haben, verfälschen und verderben. Sie stellen sich, als wären sie dem Irrthum und Rotten der Schwärmer ganz feind und gar zuwider; welchen Irrthum und Rotten sie doch gleichwohl heimlich bei sich haben und auch darüber halten. Damit sie sich aber vor Feindschaft hüten mögen auf beiden Seiten, hört man oftmals von ihnen, daß sie den Spruch anziehen und immer wiederholen: „Christus ist unsere Gerechtigkeit.“ Und wenn du dieselben Worte anziehst, so wirst du sagen, daß solches treulich und christlich also geredet werde. Siehe aber an den Betrug, und wie sie die Leute verblenden können, damit sie also davon kommen und den armen tollen Haufen an sich hängen, auf daß sich derselbe ihr Ding gefallen lasse. Sie behalten den Namen und die Person Christi, welcher Gott und Mensch ist, und halten dieselben gegen seine Werke und Leiden. Christus, sagen sie, ist unsere Gerechtigkeit, nicht die Werke, nicht das Blut, nicht der Tod oder Auferstehung, sondern die Person. Was ist aber das Andere, denn daß sie damit aus den Herzen und Augen der Gottseligen hinwegnehmen die Wohlthat des Mittlers Christi, damit er uns gemacht ist von Gott zur Weisheit und zur Gerechtigkeit, zur Heiligung und zur Erlösung, 1 Cor. 1, 30., und um welcher willen Gott Jesum zu einem Herrn und Christ gemacht hat, wie Petrus sagt Ap. Gesch. 2, 36.? Und was ist es Anderes, Gott und Mensch sein, als erniedriget werden, und dem Vater gehorsam werden bis zum Tode des Kreuzes, Phil. 2, 8.? Wie wollen sie denn nun dasselbe von einander scheiden, es sei denn, daß sie die zwei Naturen in Christo von einander reißen, und dahin fallen wollen (welches sie jetzt noch mit Fleiß verhehlen und heimlich halten), daß allein die göttliche Natur unsere Gerechtigkeit sei? Paulus, 1 Cor. 1, 30., redet von dem gekreuzigten Christo; und Röm. 3, 25. sagt er, daß derselbe uns von Gott vorgestellt sei zu einem Gnadenstuhl in seinem Blute. Christus, Joh. 3, 14. 15., sagt: „Wie Moses in der Wüste eine Schlange erhöhet hat, also muß des Menschen Sohn erhöhet werden, auf daß alle, die an ihn glauben, nicht verloren werden, sondern das ewige Leben haben.“ Daselbst hat er es ja gleichsam mit Fingern gezeigt, woran sich der Glaube halten und worauf er sich verlassen soll. Wollen sie aber haben, daß du allein die Person sollst ansehen, damit du nicht verloren, sondern gerecht und zum ewigen Leben erhalten werdest, und nicht die Erhöhung am Kreuze, so wird alsbald das Herz hin und her gezogen, und der Trost verschwindet, dadurch das Herz, so vor Gottes Gericht erschrocken, aufgerichtet und erhalten sollte werden. Derhalben müssen und können diese Stücke an dem Mittler Christo nicht von einander getheilt oder gesondert werden, und man soll sich fleißig hüten vor solcher heimlichen List, dadurch sich der Teufel untersteht, das Amt und Gehorsam Christi zu verkleinern und zu verwerfen.

Der heilige Augustinus thut eine herrliche und schöne Vermahnung aus großem Verstande und Vorzicht, da er sagt in seinem Enchiridion, oder Handbüchlein, ad Laurentium: Das ewige und rechte eigentliche Fundament des katholischen und allgemeinen Glaubens ist Christus. Denn „einen andern Grund“, sagt Paulus 1 Cor. 3, 11., „kann zwar niemand legen, außer dem, der gelegt ist, welcher ist Jesus Christ.“ Und man muß nicht sagen, daß dieses darum nicht der Grund des katholischen und christlichen Glaubens sei, dieweil man keinen mag, daß etliche Ketzer dieses auch also mit uns halten. Denn wo man fleißig bedenkt wird, was zu Christo gehört, so wird Christus bei allen Ketzern gefunden nach dem Namen, welche auch wollen Christen genannt werden; aber in Wahrheit ist er nicht bei ihnen. Solches aber zu beweisen, würde zu langwierig sein. Denn da muß man alle Ketzereien erzählen, die entweder gewesen sind, oder jetzt noch sind, oder aber so unter dem christlichen Namen haben sein mögen. Und wie wahrhaftig das auch sei, müßte man durch eine eigentliche Ketzerei insonderheit beweisen; zu welcher Disputation viel Bücher gehören, also daß sie auch angesehen wird, daß sie unendlich würde sein.

Und ist sehr nütze, daß man an diese Vermahnung gedenke, auf daß die einfältigen Herzen dadurch gestärkt werden wider die Heuchelei und Schein, so vorgewendet wird von dem Namen und Person Christi, da man doch die Stücke hinweg wirft, so Christo angehören. Denn wiewohl sie sagen, daß sie sich an das Verdienst Christi mit beiden Händen halten, wie sie pflegen zu reden, so nehmen sie doch dem Gehorsam Christi wider das öffentliche Wort Gottes und so vieler Kirchen Bewilligung diesen Namen der Gerechtigkeit, durch welchen die Sünder vor GOtt gerecht werden, daß derselbe nicht soll unsere Gerechtigkeit seyn. Sie werden aber selbst wohl sehen und inne werden, ob GOtt bei solchen Menschen sein und bei ihnen wohnen wolle, welches sie doch so oft und mit vollem Munde von sich selbst sagen und rühmen, und sind doch gar sicher dabei, daß sie der Worte Christi also mißbrauchen und sie verfälschen. Denn Christus sagt Joh. 14, 23., daß er und sein Vater die lieben wollen, welche sein Wort halten, und verheißt, daß sie bei denselben Wohnung machen wollen. Er sagt aber nicht, daß er von denen geliebt werde, die das Wort, so er aus dem Schooße des ewigen Vaters hervorgebracht, nicht halten. Darum fordert er von uns nicht mit vergeblicher Ernste, denn Mose vom Volke Israel thut, da er sagt 5 Mos. 4, 2.: „Ihr sollt nichts dazu thun, das ich euch gebiete, und sollt auch nichts davon thun, auf daß ihr bewahren möget die Gebote des HErrn, eures GOttes, die ich euch gebiete.“ Und derhalben sollten wir die Worte des Apostels Paulus immer vor unsern Ohren haben und dieselben auch inwendig zu Gemüthe führen, in welchen er den Timotheus ermahnt, sich zu befleißigen, daß er ja fest halte an dem heilsamen Worte, und sich zu hüten vor dem listigen Betrug der irrigen Geister, da er sagt 2 Tim. 1, 13.: „Diese gute Beylage bewahre durch den Heiligen Geist, der in uns wohnet“ etc., „des ungeistlichen losen Geschwätzes aber entschlage dich.“ Cap. 2, 16. Und 1 Tim. 6, 20. 21.: „Meide die ungeistlichen losen Geschwätze, und das Gezänke der falsch berühmten Kunst, welche etliche vorgeben, und fehlen des Glaubens.“ Denn darum hat der Sohn GOttes, der zur Rechten des Vaters sitzt, in diesen letzten Zeiten den lieben theuren Mann D. Luther erweckt, und andere dergleichen geistreiche Männer mehr, welche getreue Ausleger der Schrift sind im Hause des HErrn. Deren Auslegung und Bekenntniß der reinen Lehre man behalten soll, auf daß wir durch ihre ernste Vermahnung gestärkt werden, damit wir uns nicht wägen und wiegen lassen von allerlei Wind der Lehre durch Schalkheit der Menschen und Täuscherei, damit sie uns erschleichen zu verführen, Eph. 4, 14.

\quellseite{12}

Darum lässet uns allezeit wachen und beten, und dasselbe so viel desto häufiger, wie viel größere Verwirrung und Irrthum in diesem äußersten und kindischen Alter der Welt und Unbeständigkeit der Menschen in ihrem Verstande zu besorgen sind, auf daß wir würdig geachtet werden, diesem allen zu entfliehen, und zu stehen vor dem Sohne des Menschen.

Hieronymus Besold.

\tableofcontents
\mainmatter
\part[Auslegung des ersten Buches Mose]{D. Martin Luthers gründliche und erbauliche Auslegung des ersten Buches Mose}
\chapter{Das fünfundzwanzigste Capitel}

\section{Zweiter Theil}

\subsection{Vorbericht}

1. Bis hieher haben wir drei Theile dieses ersten Buches Mose ausgelegt, wie wir dasselbe droben getheilt haben. Denn im ersten Theil wird die Historie unsers ersten Vaters Adam angezeigt. Der andere Theil preiset Noah, den Prediger der Gerechtigkeit, und sagt von der greulichen Strafe der Sintfluth, dadurch die erste Welt vertilget ist, also daß nur der einige Noah mit seiner Familie erhalten ist worden. Der dritte Theil wird dem Abraham zugeeignet. Welche Historie vornehmlich wohl zu merken ist, darum daß daselbst die Verheißung vom Samen, der alle Völker segnen würde, wiederholt wird, und um mancherlei Exempel willen des Glaubens und anderer Tugenden.

2. Und wahrlich, wo jemand dies alles wird überlesen und erwägen, der wird verstehen, daß es eine große und unaussprechliche Wohlthat Gottes ist, daß die Historien solcher trefflichen großen Männer in den Büchern der heiligen Schrift beschrieben, und bis auf diesen Tag bewahrt und erhalten sind worden. Denn die Bücher und Schriften der Heidnischen wissen davon gar nichts; allein die Kirche und Gemeinde Gottes mag dieser Wohlthat genießen, und hat eine rechte Erkenntniß und göttlich Zeugniß vom rechten Anfang und wunderbarer Ausbreitung der Lehre, welche jetzt die Kirche noch hat und die GOtt anfänglich durch seine Stimme den ersten Eltern gegeben, und hat haben wollen, daß dieselbige Lehre und Zeugniß im ganzen Stamm und Geschlecht der Väter hat bleiben und getrieben sollen werden.

3. Derhalben wird die Verheißung oftmals wiederholt und erneuert durch die göttliche Stimme, welche anfänglich im Paradeis ist gehört, und hernach dem Abraham auch wiederholt und erklärt worden. Wiewohl GOtt nicht so vielmal mit Isaak und Jakob geredet hat, als er mit Abraham gethan; denn sie haben wohl gewußt, daß sie auf die Verheißung, dem Abraham geschehen, und auf seinen Glauben haben sehen sollen: jedoch, so oft diese heiligen Patriarchen in Anfechtung gewesen und ihr Glaube schwach geworden ist, richtet sie GOtt wiederum auf und stärkt sie, auf daß sie in der Gefahr und Beschwerung, so ihnen begegnet und sie betroffen hat, nicht dahin fallen und verzagen möchten. Denn wir werden in dem Leben beider Patriarchen wunderbare Kämpfe, und ganz schweren Jammer und Trübsal sehen, und in denselbigen Kämpfen wird leuchten und sich sehen lassen die wunderbare Regierung der Heiligen, dadurch sie von GOtt regieret und erhalten werden ohne menschliche Hülfe.

4. Weltliche Menschen verstehen von diesen Dingen nichts, und wenn sie lesen, daß Isaak oder Jakob wie Landstreicher sind umhergezogen, daß sie keinen gewissen Sitz gehabt, und daß es gemeine Leute gewesen, die sich mit häuslicher und knechtlicher gemeiner Arbeit bemüht haben, finden sie solches höchst lächerlich und verächtlich; denn sie sehen daselbst nichts, das sonderlich trefflich und herrlich ausgerichtet wäre nach dem Urtheil, Klugheit und Geschicklichkeit der Vernunft; wie sie denn in den Historien der Heiden solcher trefflichen und herrlichen Thaten viel lesen mit großer Verwunderung. Sie verachten aber diese Geschichten der heiligen Väter darum, weil sie nicht sehen dieser Historien rechten Schmuck und Zierde. Sie sehen nicht die so treffliche Gabe des göttlichen Wortes; sehen nicht den Glauben, die Anrufung und Gebet; sehen nicht die Geduld in der Widerwärtigkeit, dadurch die Heiligen die Welt und Nachstellung des Satans überwunden haben; welche Kämpfe und Siege, so sie gehabt, alle Geschichten und Triumphe aller andern großen Helden weit übertreffen.

\quellseite{13}
\quellseite{14}

5. Wer nun diese Historien nicht vergeblich lesen will, der soll es gewiß dafür halten, daß die heilige Schrift keine menschliche, sondern göttliche Weisheit ist: und alsdann wird er empfinden, daß sein Herz wunderbare große Liebe und Verlangen haben wird nach den Dingen, so in der heiligen Schrift begriffen sind. Denn es ist ein solcher Brunn, welcher, je mehr er ausgeschöpft und getrunken wird, je mehr Durst man darnach hat. Wie die Weisheit sagt im Buch Sirach Cap. 24, 27. 28. 29.: „Wer von mir isset, den hungert immer nach mir. Und wer von mir trinket, den dürstet immer nach mir. Denn meine Predigt ist süßer, denn Honigseim.“ Denn es wird in diesem Buche von solchen Dingen gehandelt, welche „auch die Engel gelüstet zu schauen“, wie 1 Petri 1, 12. gesagt wird; sie sind aber verborgen vor den Augen derer, so nach dem Fleisch weise sind, und sind nur dem Geiste offenbar. Wir wollen aber nun an diesem Orte das vierte Buch anfangen. Denn wiewohl die Hebräer dies Stück von Ismael und Isaak noch zu der Historie Abrahams lesen, ist es doch leichtlich zu sehen, daß Mose selbst hier ein Buch von dem Patriarchen Isaak habe wollen anfangen, nachdem Ismael und Isaak ihren Vater Abraham begraben haben.

\section[Erstes Stück]{Erstes Stück. Von Isaaks Harren auf die Verheißung; item, von Ismael, dessen Nachkommen und Tode}

\subsection{I}

\begin{quote}
\textbf{V. 11. Und nach dem Tod Abrahams segnete Gott Isaak, seinen Sohn. Und er wohnte bei dem Brunnen des Lebendigen und Sehenden.}
\end{quote}

1. Mose nimmt eine kurze Erzählung vor von Isaak, nämlich, daß ihn Gott nach seines Vaters Tode gesegnet habe; aber er bricht bald ab, und macht kürzlich das vollends ab, was noch von Ismael zu sagen ist, den Segen aber läßt er stehen bis an seinen Ort. Denn es hat sich damit noch zwanzig Jahre verzogen, in welchen Jahren Isaak mit Rebekka, seinem Weibe, ohne Kinder gelebt, und hat durch natürliche Kräfte kein Vater können werden, sondern mit Flehen und demütigem Gebet von Gott Kinder und Erben erlangt.

2. Dies ist aber ein wunderbar Ding, wenn man es mit dem vergleicht, so zuvor im 24. Cap. V. 60. gesagt ist. Denn es hat Abrahams Knecht mit großem Gepränge und Herrlichkeit für Isaak um eine Jungfrau geworben, die er möchte zum Weibe nehmen, und ist ihm auch durch göttliche Versehung dieselbige Jungfrau ehelich vertraut worden mit solcher großen Hoffnung und Erwartung, daß sie ist angesehen worden als die unzähligen Kinder Mutter sein würde. Denn wie sollte Isaak nicht darauf gehofft haben, welchem ja durch göttliche Versehung das Weib ist zugeführt worden, und zwar durch so einen frommen, heiligen Vater, dazu auch durch einen sehr getreuen Diener und endlich durch den Dienst des Engels? Doch hat er gleichwohl Rebekka zwanzig Jahre unfruchtbar gehabt. Es ist aber sein Vater noch am Leben gewesen, und werden ohne Zweifel beide herzlich gewünscht, gefleust und um Kinder gebeten haben. Und ich kann zwar nicht wissen, ob solches die Ursache gewesen sei, daß Abraham Ketura zum Weibe genommen hat, da er gesehen, daß es sich so lange verzogen, ehe Isaak Samen bekommen hat. Ich will aber das nicht für gewiß sagen. Denn diese heiligen Männer, und die mit so mancherlei Anfechtung geübet worden, haben sich durch solche lange Uebung daran gewöhnt, daß sie auf die Hand des Herrn geharrt haben: und Gott hat sie uns vorgestellt zum Exempel des rechten, wahren Gottesdienstes, welcher vornehmlich und wahrhaftig nicht in äußerlichen Opfern oder Werken besteht, sondern im Glauben, Hoffnung und in der Liebe Gottes.

3. Die Vernunft läßt sich träumen, man müsse Gott dienen und ihn versöhnen mit leiblichen Opfer oder andern Uebungen, so von Menschen erdacht sind. Aber die Exempel der Väter zeigen an, daß dies der vornehmste und höchste Gottesdienst sei, daß man auf Gott harre. Und das ist auch der rechte Nutzen und die rechte eigentliche Uebung des Glaubens. Denn der Glaube reißt uns erstlich dahin auf das Unsichtbare, wenn er uns vorhält die Dinge anzunehmen, so man vor Augen nicht sieht, und dasselbe können wir zwar etlichermaßen dulden und ertragen; aber das Herz wird nicht allein auf das Unsichtbare geführt, sondern wird oft auch lange hingehalten und wird die Hülfe in die Länge gezogen. Wie Abraham droben fünf und zwanzig Jahre gewartet, ehe ihm der Sohn geboren worden, und Isaak muß zwanzig Jahre ohne Kinder sein. Das Dritte aber und das Allerbeschwerlichste ist dies, wenn auf den Verzug und die lange Zeit das auch noch folgt, daß es sich gar widersinnig schickt und ansehen läßt. Wer alsdann verharren und warten kann, und auf das hoffen, so verzogen wird, und das lieben, so gar widersinnig ist, derselbe wird endlich erfahren, daß Gott wahrhaftig ist und seine Verheißung treulich hält.

4. Dies sind aber keine Werke des Fleisches oder menschlicher Vernunft, oder auch der Mönche und Heuchler, welchen aller Verzug zu lange und verdrießlich ist; denn was sie bitten, dasselbe wollen sie alsbald haben, und vertrauen Gott nicht, denn allein in den Dingen, so gegenwärtig sind, aber nicht in dem, damit es sich noch immer verzieht. Gott aber will haben, daß was er verheißt, unsichtbar und widersinnig sein soll, auf daß wir also versucht und geübt werden, und lernen, daß dies der rechte wahre Gottesdienst sei, und der Gott am allerangenehmsten ist, nämlich, wo man auf ihn harret. Daher kommt es, daß in den Psalmen diese Vermahnungen so gemein sind, als, im 27. Psalm V. 14.: „Harre des Herrn, sei getrost und unverzagt.“ Und die Propheten prägen uns dasselbe auch immer ein, als wollten sie sagen: Harret und glaubet doch.

5. Aber diese Lehre vom Glauben, und daß man auf Gott harren soll, wird von den Weltmenschen verlacht und verspottet; wie ihre spöttischen Worte erzählt werden Jes. 28, 13.: „Gebet hin, gehet her; harre hie, harre da; hie ein wenig, da ein wenig“; denn also verlachen sie dazumal auch die Lehre und den Trost vom Glauben und Geduld als vergeblich und ungewiß Ding. Nun ist es aber doch gewiß, daß das, darum man bittet oder das verheißet wird, alsdann stehet und wahrhaftig vorhanden ist, wenn wir anfangen zu beten und zu vertrauen; wie der Spruch lautet Jes. 65, 24.: „Und soll geschehen, ehe sie rufen, will ich antworten; wenn sie noch reden, will ich hören.“ Das ist wahrlich am gewissensten und wahrhaftigsten.

6. Aber wenn das menschliche Herz in der Anfechtung und Gefahr ist, kann es schwerlich mit diesem Trost zufrieden sein; denn also pflegt es sich für und für zu ängsten und zu klagen: Was will denn werden? Wann wird es geschehen? Wo wird es denn geschehen? Darum antworte ich: Harre, harre. Wo es sich aber noch etwas länger verzieht, und das Herz abermals fragt: Wann es doch endlich einmal geschehen werde? so sollst du sagen: Ich weiß dir keinen andern Rath zu geben, denn daß du Geduld trägst und noch länger harrest ein, zwei, drei Jahre. Es wird der Herr mit seiner Hülfe gewißlich kommen und nicht verziehen.

7. Denn das ist wahrhaftig und ist auch unser Herr Gottes Gebrauch also, daß er pflegt zu verziehen und die Hülfe in die Länge aufzuschieben, aber doch also, daß er in der Noth komme und bald komme; wie das in der wunderbaren Erlösung des Volkes Israel, da sie aus Egypten sind geführt worden, zu sehen ist. Das Volk hatte lange Zeit gewartet, da es mit schwerer Dienstbarkeit beladen war, und der Herr hatte immer mit der Hülfe verzogen; aber da sie in die äußerste Gefahr ihres Leibes und Lebens geführt worden sind, und den Tod vor ihren Augen gesehen haben von allen Seiten her, da sagt der Herr zu Mose: „Was schreist du zu mir?“ 2 Mos. 14, 15.; und alsbald ist die Erlösung da und werden ihnen die großen Thaten Gottes gezeigt. Wo es aber die höchste Noth nicht erfordert, und die Beschwerung so groß nicht ist, daß uns durch menschlichen Rath davon nicht kann geholfen werden, alsdann verzieht Gott mit seinem Werk und mit seiner Hülfe.

\quellseite{15}

8. Diese schönen herrlichen Exempel werden uns an den heiligen Männern vorgehalten, auf daß wir lernen, daß der vornehmste und geistliche Gottesdienst nicht daran gelegen sei, daß man Kirchen baue, viele Ceremonien anrichte; welches alles Kinderspiel ist, damit uns Gott lockt und zu sich reizt: wie er die Juden mit mancherlei solcher äußerlichen Weise und gleich als mit einer Kinderzucht geübt hat. Aber der Gottesdienst der Väter ist, daß sie auf Gott geharrt haben; wie Hiob sagt: „Wenn er mich auch schon tödten würde, will ich doch auf ihn hoffen.“ Und Jakob hernach 1 Mos. 32, 24 ff., da er mit dem Engel ringt, hat fest und mit beständigem Herzen an der Verheißung gehalten. Aber dieselbe Verheißung verzieht sich fast lange und sein Glaube wird auf mancherlei Weise angefochten; jedoch ist sein Herz also gestanden: Ich glaube, daß die Verheißung wahrhaftig sei; wo sie nun nicht bald erfüllet wird, so hoffe ich doch immer darauf; denn wo er auch schon mit mir ringt, will ich ihn doch lieben. Solche Opfer geben den süßesten und lieblichsten Geruch von sich vor Gott, und erfreuen endlich das Herz des Menschen mit ewiger Freude und Seligkeit.

9. Dagegen aber, die in der Noth zaghaftig werden und von Gott abfallen, berauben sich selbst der ewigen Seligkeit. Wie Sirach sagt Cap. 2, 16.: „Wehe denen, so die Geduld verloren haben.“ Wer aber beharrt und das Vertrauen und Hoffnung behält, derselbe überwindet endlich und bringt den Sieg davon, Sir. Cap. 16, 13. Also hat Noah hundert und zwanzig Jahre geharrt und hat unterdeß Spott und Hohn von den Gottlosen leiden müssen, endlich aber ist die ganze Welt durch die Sintfluth ersäuft worden, er aber ist mit seinem Hause wunderbarer Weise erhalten und ein Vater der Gerechten geworden.

10. Darum, sage ich, soll man Gott dienen mit Glauben, Hoffnung und Liebe; die äußerlichen Ceremonien sind nur Nebenübungen der groben einfältigen Leute, dadurch sie an die Werke der Gottseligkeit, so viel höher sind, gewöhnet werden; wie die kleinen unmündigen Kinder erstlich mit Milch und weicher Speise genähret und auferzogen werden. Wir aber, so Andere zu regieren und zu lehren gesetzt sind, sollen lernen, daß der rechte und wahre Gottesdienst nicht sei, wo man unvernünftige Thiere opfert etc., sondern wo man an seiner Verheißung hält und glaubt, daß dieselbe wahrhaftig ist und nicht fehlen kann. Auf solchen Glauben und Vertrauen folgt dann die Hoffnung, die mich dessen erinnert, daß ich harren soll und daß ich einen gnädigen Gott habe. Wenn ich aber also in der Noth beharre, alsdann wird Pharao mit seinem ganzen Heer in der Tiefe des Meeres ersäuft, und wird mit dem, so im Glauben beharrt und beständig geblieben ist, einen solchen Ausgang nehmen, wie er gehofft hat.

11. So bricht nun Mose ab die Erzählung, so er vorgenommen hat vom Segen Isaaks. Denn dieweil er sagt, daß Gott Isaak gesegnet habe, soll man dasselbe verstehen im Glauben und Hoffnung; sintemal daselbst noch das Widerspiel ist gesehen worden und Isaak ganzer zwanzig Jahre auf Kinder gewartet hat, und hat ihm am Alter, an der Natur und auch an Kräften nicht gemangelt. Zudem war auch Rebekka selbst jung, schön und zur Kinderzucht und Geburt ihres Alters wegen wohl geschickt; und hat sich doch gleichwohl mit der Verheißung verzogen, und ist dieselbe unsichtbar geblieben, bis daß sie sicher an ihre sechzig Jahre gekommen sind, und daß nun gewißlich die Gefahr der Unfruchtbarkeit vorhanden gewesen. Da wird Isaak ohne Zweifel in seinem Herzen bedacht und sich erinnert haben, was seine Eltern für Sorge gehabt in gleicher Gefahr der Unfruchtbarkeit, und wird derhalben also gedacht haben: Siehe, es hat eben dieselbe Anfechtung mich und meine Rebekka betroffen, die zuvor meine Eltern, Abraham und Sarah, auch geplagt hat, dieweil sich Alles lässet ansehen, als ob es widervinnig und der Verheißung stracks zuwider sei; und hat doch bei sich beschlossen, an derselbigen fest zu halten nach dem Exempel seines Vaters.

12. Dies geschieht aber darum, auf daß Gott dem rufe, das nicht ist, daß es sei, wie Paulus sagt Röm. 4, 17. Wie heutiges Tages der Türke und Papst mit unmenschlichem Haß und Grausamkeit wider uns toben und wüthen, unwendig aber unsere Gewissen auf mancherlei Weise versucht werden und kämpfen mit Furcht, Unglauben und Schrecken der Sünden. Und wir haben zwar wider dies alles starken Trost, als, die Verheißung des Wortes, die heilige Taufe, das Abendmahl des Herrn, welche man mit festem Glauben ergreifen und sich daran halten soll. Ich sehe aber die Vergebung der Sünden nicht, ich sehe nicht die ewige Seligkeit, ich sehe nicht das Leben; sondern glaube es und halte mich also immer mit der Hoffnung aufrecht. Und wenn der Glaube angefochten und schwach wird, es sei gleich durch Anfechtung, oder durch Sündenfälle und menschliche Gebrechlichkeit, so halte ich mich doch gleichwohl an die Schlüssel. Ich bleibe in der göttlichen Verheißung und wenn auch der Himmel darüber einfallen sollte. Das heißt Gott dienen und das erste Gebot erfüllen; denn wozu dasselbe nütze ist, sieht man erst in der Anfechtung, wo die Verheißung unsichtbar ist; ich werde lange hingehalten und wird endlich die Sache gar widerwillig.

13. Der Teufel aber handelt und wirkt nach einer anderen und entgegengesetzten Regel; denn derselbe pfleget bald bei seinen Priestern und Propheten zu sein, wenn sie nur wollen, und verzieht nicht mit dem, daß sie bitten und begehren. Also erwecken oftmals die Zauberer plötzlich Ungewitter und Donnerschlag. Auf diese Weise bethört und fängt der Teufel die Leute, daß sie ihm gerne dienen, dieweil er so bald hilft und sie erhört; und was das Ungerechte ist, wird dafür gehalten, als geschehe solches alles von Gott. Darum soll man auf solche Ungleichheit der Wirkung Gottes und Satans fleißig Acht geben. Gott will haben, daß wir harren sollen; wie der 130. Psalm V. 5. sagt: „Meine Seele harret, und ich hoffe auf sein Wort.“

14. Wenn du Gott dienen willst, so sollst du glauben, was unsichtbar ist; du sollst hoffen auf das, so verzogen wird; und sollst Gott lieben, wenn er sich auch schon erzeigt, als ob er dein Feind und dir zuwider sei; und sollst also beharren bis an das Ende. Das meint Mose damit, da er es einen Segen nennt, nämlich, Isaak ein Weib geben und ihm doch keine Kinder geben. Isaak hat um ein Weib gebeten, den Samen auszubreiten und zu mehren, und hat dasselbe auch erlangt. Aber er hat zwanzig Jahre geharrt ohne Kinder; wie sein Vater Abraham, auch so viel Jahre geharrt hat, bis Ismael ist geboren worden, welcher doch der Sohn nicht war, so ihm Gott verheißen hat, und er hat auf den Sohn von der Sarah noch viel länger hoffen und harren müssen.

15. Von dem Brunnen des Lebendigen und Sehenden ist droben gesagt worden. Da Hagar von Sarah, ihrer Frau, ist geflohen, ging sie in der Wüste irre bei Bersaba nach Egypten hin, und hat daselbst dem Brunnen einen Namen gegeben, daß bei allen Nachkommen Hagar sollte gerühmt werden als eine treffliche Domina und Matrone, von welcher der Brunnen den Namen erlangt hätte. Denn sie sagt: „Gewißlich, hie hab ich gesehen den, der mich hernach angesehen hat“, 1 Mos. 16, 13. Dasselbe hat Hagar dafür gehalten, daß es ihm auch rühmlich wäre, wiewohl Hagar seine Stiefmutter war: und darum hat er gewohnt bei dem Brunnen des Lebendigen und Sehenden. Jetzt aber wird nun Mose die Historie von Ismael zu Ende führen.

\subsection{II}

\begin{quote}
\textbf{V. 12—16. Dies ist das Geschlecht Ismaels, Abrahams Sohns, den ihm Hagar gebar, die Magd Sarahs aus Egypten; und das sind die Namen der Kinder Ismael, davon ihre Geschlechter genannt sind: Der erstgeborne Sohn Ismaels Nebajoth, Kedar, Abdeel, Mibsam, Misma, Duma, Masa, Hadar, Thema, Jetur, Naphis und Kedma. Dies sind die Kinder Ismaels mit ihren Namen in ihren Höfen und Städten, zwölf Fürsten über ihre Leute.}
\end{quote}

16. Diesen Ismael beschreibt nun Mose auf eine andere und entgegengesetzte Weise, denn er Isaak beschrieben, von dem er gesagt hat, daß ihn Gott gesegnet habe, und hat doch den Segen nicht ausgedrückt. Vorhin (Cap. 21, V. 11. ff.) aber haben wir von Ismael ein sehr betrübt Urtheil gehört, nämlich, daß er mit Schmerzen und Unwillen seines Vaters aus dem Hause verstoßen worden nicht allein durch der Mutter, sondern auch durch göttliche Gewalt, und ist ein Sohn des Fleisches und nicht der Verheißung genannt worden. Das ist sehr hart und beschwerlich. An diesem Orte aber wird er nun gerühmt nicht allein des fleischlichen, sondern auch des geistlichen Segens halben.

\quellseite{16}

17. Und im fleischlichen Segen ist er zwar höher denn sein Bruder Isaak, welches denn den Isaak seine Anfechtung und Bekümmerniß sehr gemehret hat. Denn derselbe hat noch keine Kinder, und hat so viel Jahre auf den Segen warten müssen, bis ihm erst Kinder geboren worden sind, da er achtzig Jahr alt gewesen ist. Mittlerweile werden Ismael zwölf Fürsten gegeben, lange zuvor ehe denn Isaak Kinder gezeuget, und hat doch weder er noch Abraham so viel Fürsten gezeuget. Ismael aber, sein Bruder, der von der Magd geboren war, wird mit großem und reichem Segen gezieret, daß er bald zwölf Fürsten gezeuget, welche sie alle gesehen haben, beide Abraham und Isaak, daß sie gelebet haben und in großer Herrlichkeit gewesen sind. Denn Ismael ist acht und vierzig Jahre nach Abraham gestorben, und die Kinder Ismael sind bald zu großen Männern und Fürsten gerathen, ehe denn der Segen, dem Abraham und Isaak verheißen, ist angefangen worden.

18. Dies ist wahrlich eine schwere Anfechtung des Glaubens, der Hoffnung und der Liebe gewesen. Denn sollte das nicht ein Herz anfechten, wie gottselig und fromm dasselbe auch sein mag, vornehmlich aber das Herz Isaaks, der Gott geglaubt und gehört hat, daß er es vorgezogen worden allen Völkern und seinem Bruder Ismael. Denn also ist gesagt worden 1 Mos. 21, 12.: „In Isaak“, und nicht in Ismael, „soll dir der Same genannt werden“, in welchem Samen die ganze Welt wird gesegnet werden. Er hat wohl Ismael auch gesegnet, da er 1 Mos. 17, 20. sagt: „Ich will ihn zum großen Volk machen, und er wird zwölf Fürsten zeugen“, aber das ist alles nichts gegen den verheißenen Samen zu rechnen.

19. Hier läßt es sich aber nun ansehen, daß solches alles geändert und umgekehret sei. Denn Isaak, der die rechte Verheißung hat, wird hinter Ismael gesetzt, und Ismael, der doch nur eine leibliche Verheißung hat, geht vor dem Andern weit her und wird ein Herr über andere Herren, hat zwölf Fürsten, da Isaak allein und ohne Kinder zu Bersaba lebt wie ein Stock oder Klotz. Darum wird er gedacht haben: O Gott! ist das deine Verheißung und deine Wahrheit? Solche Gedanken sind sehr scharfe Stacheln gewesen, so ihn zur Verzweiflung und Unglauben gereizet haben. Wie das Weib Hiobs ihrem Manne auch aufrückt und vorwirft, da sie sagt: „Hältst du noch fest an deiner Frömmigkeit? Segne Gott und stirb“, Hiob 2, 9. Dasselbe hätte man zwar zu Isaak auch sagen mögen. Denn der, so zuvor verachtet gewesen und aus dem Hause verstoßen ist, erlangt den göttlichen Segen, der viel reicher und größer ist, denn Isaak oder Jakob je erlangt haben.

20. Derselben sind dies wunderbare Historien und gehen den fleischlichen Menschen Ursache zu großen Aergernißen. Denn alle Werke Gottes sind stracks wider die Verheißung, welche doch gleichwohl ganz wahrhaftig und beständig bleibt. Aber dies ärgert die Leute, daß sie unsichtbar ist, daß sie also verzogen wird und immer das Gegentheil geschieht.

21. Darum soll man den wunderbaren Rath Gottes lernen, damit er seine Heiligen regiert, und die Herzen der Gottseeligen sollen sich daran gewöhnen. Wenn du eine göttliche Verheißung hast, so wird dir, je mehr du von Gott geliebet wirst, desto mehr solche Verheißung verborgen sein, wird verzogen und ins Gegentheil verkehrt werden. Denn wo er dich nicht so sehr liebte, würde er mit dir nicht also spielen, das ist, er würde seine Verheißung und Hülfe nicht also verziehen und ins Gegentheil verkehren.

22. Denn das sind die gewissesten Zeichen eines väterlichen Herzens, daß gleichsam entbrannt ist von großer Liebe, die er zu dir trägt. Also thut ein leiblicher Vater auch, der sein Kind liebt: er spielt mit ihm, verheißt ihm ein Geschenk zu geben; aber er verzieht damit, und stellt sich, als wollte er seine Verheißung nicht halten. Daselbige thut er darum, dieweil er das Kind sehr lieb hat, und will es ihm reichlich geben, sondern auch das Kind beharrt und den Verzug in sich frißt und überwindet. Also spielt Gott mit uns auch; aber wir werden bald zornig und ungeduldig selbst um eines kleinen Verzuges willen. Denn wir wollen dasjenige, so uns verheißen wird, entweder alsbald gegenwärtig haben, oder aber wollen es sonst auf eine andere Weise und Wege, was auch das für ein Weg sein mag, erlangen.

23. Darum lehren uns die Exempel der Väter auf diese Weise, welches der rechte Gottesdienst sei, nämlich, ein feiner, reiner Glaube, vollkommene Hoffnung und beständige Liebe, dadurch wir verstehen und empfinden, daß Gott gegenwärtig und gnädig ist, wenn wir uns auch schon dünken lassen, daß er uns zuwider sei.

24. Es läßt sich aber übrigens ansehen, daß Ismael ein großer Theologe gewesen sei, der die rechte Kraft und Würde der Verheißungen, welche wahrlich groß und herrlich sind, fleißig erwogen hat, und das vornehmlich, daß seine Mutter Hagar dessen werth gewesen ist, daß sie den Engel gehört, der sie geheißen hat, sie sollte wiederum zu ihrer Frau Sarah kehren und sich unter ihre Hand demüthigen, 1 Mos. 16, 9. Dies alles wird er gerühmt und herrlich davon geredet haben; wie ein jegliches Land oder Geschlecht vortrefflicher sein will, als die Andern sind, und will allein die Welt regieren. Die Andern aber verachtet es gemeiniglich. Darum hat Ismael sich selbst die Verheißungen, so ihm geschehen waren, fleißig eingepräget, und hat sagen können: „Ich sehe, daß der Herr bei mir ist und mich gesegnet hat auch mit geistlichem Segen.“ Und wird auch dabei zugleich den äußerlichen Gottesdienst, so er in seines Vaters Hause gesehen hat, behalten und gemehret haben; und werden auch etliche von seinen Kindern und Hausgesinde durch solche Ceremonien zur Erkenntniß der Göttlichkeit gekommen sein; wie auch ihrer Viele vom Geschlecht und Nachkommen Kains zu der rechten wahren Kirche Adams versammelt worden sind und heutiges Tages auch immer viel Leute zu unserer Kirche und kleinen Häuflein versammelt werden.

25. Auf dieselbige Weise haben sich auch ihrer Viele aus dem Hause Ismaels zu der Kirche und Gemeinde gehalten, so im Hause Abrahams gewesen ist, vornehmlich Ismael selbst, von dem ich glaube, daß er Buße gethan, und darnach wiederum sich zur rechten wahren Kirche bekehret und derselben auch theilhaftig geworden sei, nicht vom Fleisch, sondern aus dem Geist; denn nach dem Fleisch ist er verworfen worden. Aber im Text wird er auch geistlich gelobt. Item, Mose hat droben V. 9. erzählt, wie er mit Ehrerbietung seinen Vater Abraham begraben habe, welches ein Anzeichen ist, daß er von der Kirche und Gemeinde, so im Hause Isaaks gewesen, nicht abgesondert gewesen ist.

26. Darum halte ich, er sei fromm und gottselig gewesen, wiewohl er ein Kriegsamann war. Er hat auch seinen Kindern herrliche Namen gegeben. Denn Kedar heißt Einer, der da traurig, sauer und schwarz ist, wie die Mönche, als der vor Leide sterben wird. Und vielleicht ist der fromme Mann auch etwas traurig gewesen, da er dem Kinde denselben Namen gegeben hat. Und es lautet auch zugleich, als bedeutete dieser Name eine Ehrerbietung mit Kniebeugen, oder Einen, vor dem man sich geneiget hätte. Was die Juden von diesem Namen für Lügen eingeführt, die lasse ich stehen.

27. Mir scheint es aber, daß man daraus schließen und abnehmen könne, daß Ismael die Lehre von dem Gott Abrahams fleißig getrieben, und seine Kinder und Hausgesinde an den rechten Gottesdienst, welchen er zuvor gesehen, gewöhnet hat. Aber wie es pfleget zu gehen, daß der rechte Gottesdienst allezeit Aberglauben gebieret, also haben auch seine Nachkommen den rechten Gottesdienst fahren lassen, und nur einen äußerlichen Schein und Heuchelei behalten; denn wo das erste Gebot weg ist, da bleiben nur Spuren und Träber übrig. Also, da anfänglich die Leute gesehen haben, daß Abraham auf diese oder andere Weise geopfert, Gott gelobet und gedankt hat, haben sie denselben äußerlichen Gottesdienst angenommen, aber dabei des ersten Gebotes vergessen. Also werden die Leute lauter Affen der Väter; denn sie folgen allein den äußerlichen Werken nach, und lassen den Glauben gar liegen und achten desselben nicht.

28. Daher ist alle Abgötterei gekommen. Die Heiden haben gesehen, daß Abraham und die andern Väter ihr Angesicht nach dem Aufgang der Sonne gewendet und den Schöpfer dieses schönen Lichtes angebetet haben: dasselbe haben sie auch nachthun wollen, haben aber nur die leeren Schalen behalten und den rechten Kern verloren. Also thun auch die Rottengeister, nehmen von uns das Wort, aber nicht wahrhaftig; denn sie mißbrauchen desselben zu eitler Ehre und daß sie nur den gemeinen Pöbel damit an sich hängen.

29. Also haben sich die Ismaeliten Kinder Gottes genannt, und wollen heutiges Tages nicht Hagarener, sondern Saracener genannt werden, allein darum, daß sie dem äußerlichen Gottesdienst nachfolgen und den Schein davon haben, und um des leiblichen Segens willen, welchen die Araber jetzt noch haben, so Ismaeliten sind; wie droben Cap. 16, 12. in der Verheißung ist gesagt worden von Ismael: „Seine Hand wider jedermann“ etc. Denn sie haben das Regiment über die Welt, und rühmen sich, daß ihnen die Beute und der Raub gebühre, und leben auch vom Raube. Denn wiewohl die Türken, welche Seythen sind, Arabien eingenommen und daselbst das Regiment haben, behalten sie doch die arabische Sprache und wird auch dieselbe am türkischen Hofe am meisten gebraucht. Sie sind aber bisher zuvor nie überwunden worden, weder von den Römern, noch von den Persern. Der Türke hat sie gedemüthigt; sie sind aber doch zugleich auch mit im Regimente. Also ist nun Ismael zuerst leiblich gesegnet worden, da zwölf Fürsten von ihm geboren sind; nun wird auch der geistliche Segen folgen.

\quellseite{17}

\begin{quote}
\textbf{V. 17. Und das ist das Alter Ismaels, hundert und sieben und dreißig Jahr; und nahm ab, und starb, und ward gesammelt zu seinem Volk.}
\end{quote}

30. Ismael hat hundert und sieben und dreißig Jahre gelebt, ist geboren worden, da Abraham sechzehn und achtzig Jahre alt gewesen ist, vierzehn Jahre vor der Geburt Isaaks, und ist gestorben acht und vierzig Jahre nach Abraham. Isaak aber hat lange nach Ismael gelebt, und Ismael hat doch gleichwohl gesehen die Kinder Isaaks, Jakob nämlich und Esau. Dieser Text aber zeigt klärlich an den geistlichen Segen Ismaels. Denn Mose gebraucht dieselben Worte, die er droben (V. 8.) auch gebraucht hat, da er den Tod Abrahams beschreibt; es ist nicht ein Tüttel geändert worden. Darum bezeuget er damit, daß er rechtgläubig und gottselig gewesen sei, dieweil er sagt, er sei gesammelt worden zu seinem Volk, das ist, zum heiligen Volk.

31. Das hebräische Wort gava heißt, wenn Einer zu Bette liegt und sein sanft und stille verschiedet; wie droben (V. 8.) von Abraham steht, da sich die Stunde des Todes genahet, ward er schwach und starb in einem ruhigen Alter, ohne Schrecken, sondern wie der Leib pflegt aufgelöst zu werden, auf daß die Seele vom Leibe erlöset werde. Also ist Ismael auch gestorben: der Herr hat ihm ein feines, seliges, und stilles Ende dieses Lebens bescheret, nämlich, daß er schwach geworden, sein Leben beschlossen und sein sanft gestorben und in den Schooß der heiligen Väter aufgenommen worden ist, dahin kurz zuvor Abraham und die andern Väter auch gezogen sind. Dies ist der geistliche Segen.

32. Es ist aber auch eine schöne und liebliche Beschreibung der Unsterblichkeit, da Mose saget: „Er ist zu seinem Volk gesammelt worden.“ Wir leben jetzt unter dem groben Volk dieser Welt, welches wenig nach GOtt fragt, ja, wir leben hier im Reiche des Teufels. Wenn wir aber aus diesem armseligen Leben ziehen werden, so werden wir sein sanft verscheiden, und zu unserem Volke versammelt werden, da kein Unglück, keine Noth, keine Trübsal ist, sondern Friede und Ruhe, und ein feiner sanfter Schlaf im Herrn. Wo aber nun ein ander Volk ist ohne das, mit welchem wir hier leben, so muß auch eine Auferstehung der Todten sein, und ist dies ein gewisses Anzeichen, daß ein GOtt sei und daß die Welt nicht ohngefähre und vergeblich geschaffen sei; denn daß man nach dem Tode leben soll, ist keines Menschen, sondern GOttes Werk.

33. Plinius und die Epicurer spotten und verlachen diese Lehre, und geben uns darin keinen Beifall, daß es ein solcher Gang sein sollte, dahin man von diesem unruhigen Leben zur Ruhe gehen sollte: darum führen sie immer diese Epicurischen Worte und sagen: Post mortem nulla voluptas: Wenn man stirbt, so hat die Freude ein Ende. Item: Pereat, qui crastina curat: GOtt gebe dem dies und jenes, der da sorgt, wo er morgen sein wird. Item, wie Martial sagt, daß ein seliges Leben darin bestehe, wo man des letzten Tages nicht begehrt, sich auch davor nicht fürchtet: Summum nec metuas diem, nec optes etc. Also stärken sie sich gegen die Verachtung des Todes und heben einfach alle Furcht auf, das ist, sie nehmen die Hoffnung der Unsterblichkeit gar hinweg. Wovor sollte ich mich fürchten? sagen sie; worauf sollte ich hoffen, dieweil kein GOtt ist? Warum lasset uns essen, trinken, spielen, guter Dinge sein etc.

34. Wenn aber die Vernunft das könnte zugeben, daß es ein Gang wäre von dem betrübten und geplagten Volk zu einem andern Volk des Friedens, so würde sie von Noth wegen daß auch bekennen, daß nach diesem Leben ein ander Leben sei. Denn „zum Volke gehen“ heißt ja nicht, daß man zu Nichts gehe oder zunichte werde. Und vornehmlich, wenn gefragt wird: „zu seinem Volke gehen“, nämlich, die auch denselbigen Glauben, Hoffnung und Trübsal gehabt haben und dieselbe Sprache, ja, zu unsern Mitbürgern und Landsleuten, welche wahrlich etwas sind. Wir ziehen nicht zu den Feinden, auch nicht zu den bösen Geistern, ja, wir weichen von denselben weg und werden versammelt zu unsern Vätern.

35. Dies sind Zeugnisse der Väter von der Auferstehung, von der Unsterblichkeit und vom ewigen Leben, nämlich, dieweil daselbst Völker sind; darum ist es nothwendig, daß wir leben und auferstehen werden. Die nun diesem Worte glauben, werden sich der Epicurer Narrenwerk nimmermehr bewegen lassen: aber die Vernunft verstehet solches nicht und nimmt es auch nicht an; ja, das noch mehr ist, sie sieht vor Augen, daß die Geburt aus Nichts geschieht, und kann doch daraus nicht schließen, daß sie ein göttlich Werk sei. Ich sehe, daß Isaak von einer unfruchtbaren Mutter geboren wird, welcher Leib erstorben, und so hart gewesen ist, wie die Knochen sind im Schädel. Ich sehe kein Haar, keinen Leib, keine Seele, und kommt doch gleichwohl ein Sohn aus den Lenden Abrahams und vom Leibe Sarahs, der Mutter, der dies alles hat. Kommt denn solches nicht her aus Nichts? Es ist wahrlich der Same ein todt Ding und der Leib auch, und ist noch mehr todt, denn das Korn, so in die Erde fällt. Darum soll die Vernunft denken, daß dies keine natürliche Wirkung sei, wie sie sich träumen läßt. Es ist aber eine Blindheit und Bosheit des Teufels, daß wir dies nicht ansehen oder betrachten, und verstehen also der keines, wo wir der Vernunft folgen, wie sie uns leitet, weder die künftige Auferstehung noch die Geburt des Menschen. Denn wer die Auferstehung der Todten nicht glaubt, derselbe glaubt oder sieht auch nicht das Wunderwerk der Geburt, nämlich, daß ein Mensch vom Menschen, ein Kind vom Kinde geboren wird.

36. Nun ist aber die ganze Welt voll Zeugniß von der Auferstehung. Aus einem Baume und harten Holze wächst eine schöne Blume; es wachsen Blätter, Zweige und schöne liebliche Früchte heraus. Dieweil es aber ein so gemein und alltäglich Ding ist, darum wird es gering geachtet, und sind die Herzen der Menschen so gar erstarrt, daß wenn auch Lazarus alle Tage auferweckt würde, lassen sich doch die Ungläubigen solches gar nicht bewegen. Daraus folgt aber doch nicht, daß man auf kein ander Leben hoffen solle, welches besser und seliger ist denn dies zeitliche Leben. Es wird aber die Einsammlung und die Offenbarung der ewigen Freude verzogen, bis dies feindselige Volk aufhören wird, und wird unser Volk bleiben, welches mit uns die Auferstehung der Todten geglaubt hat. Derhalben, um dieses Zeugnisses willen der heiligen Schrift von Ismael machen wir ihn zum Heiligen und großen Patriarchen; wiewohl seine Kindeskinder und Nachkommen abgefallen und ihrem Vater nicht gleich gewesen.

\begin{quote}
\textbf{V. 18. Und sie wohneten von Hevila an, bis gen Sur gegen Egypten, wenn man gen Assyrien gehet. Er legte sich aber gegen alle seine Brüder.}
\end{quote}

37. Die Kinder Ismael haben heutiges Tages oder je bei den Zeiten des Hieronymus diese Orte auch inne gehabt. Der Erstgeborne, Nebajoth, hat ein groß Theil Arabiens eingenommen, aber doch nicht das ganze Arabien: von welchem auch das ganze Land Nebathea ist genannt worden, welches in den Büchern der griechischen und lateinischen Schreiber gerühmt wird, und zwar um der Steine Ursache willen, daß bei den Einwohnern desselbigen Landes die Undankbarkeit der Kinder gegen die Eltern und Lehrmeister am Leibe ist gestraft worden. Welches eine treffliche Disciplin und Theologie gewesen ist in demselbigen Lande; und was heutiges Tages vielen gestattet wird ungestraft zu thun, das haben diese Heiden so hart und ernstlich gestraft.

38. Der andere Theil Arabiens ist Kedar, welches nicht weniger berühmt ist, und wird das ganze steinichte Arabien mit diesem Namen genannt. Die Einwohner haben sich einen Theil selbstiglich zugeeignet; wie im 120. Psalm V. 5. steht: „Wehe mir, daß ich ein Fremdling bin unter Mesech; ich muß wohnen unter den Hütten Kedars“ etc., das ist, ich bin ein Fremdling unter den Moscovitern und Tartaren. Derselbige Strich von Hevila bis gen Sur ist in Arabien, oder ist die Wüste zwischen Egypten und dem gelobten Lande. Von Sur aber wird gesagt, daß es gegen Egypten liegen soll, darum daß es an dasselbe grenzt. So haben nun die Kinder Ismael gewohnt in dem Strich vom Rothen Meere bis an den Phrath; sie haben beinahe ganz Arabien inne gehabt, ausgenommen den Theil, so die Jdumäer eingenommen haben.

\quellseite{18}

39. Von dem letzten Stücke im Texte: Coram cunctis fratribus suis corruit, welches wir verdolmetscht haben: „Er legte sich aber wider alle seine Brüder“ etc., disputiren die Ausleger mancherlei Weise. Die gemeine Auslegung hält dafür, daß Ismael gestorben sei vor dem Angesichte seiner Brüder, das ist, seiner Kinder, welche vom Vater gerufen worden und also zusammengekommen sind, ihren Vater zu hören und zu trösten, also daß er verschieden sei, daß solches die Brüder angesehen und dabei gegenwärtig gewesen sind.

40. Dies ist aber meine Meinung, wiewohl ich es nicht für gewiß sagen will: Mose redet nicht von Ismael nach der Person, denn er war jetzt schon gestorben: item, er hat zuvor nicht gesagt, habitarit, er hat gewohnt, sondern, habitaverunt, sie haben gewohnt, nämlich, die zwölf Söhne, in dem Striche vom Rothen Meere nach Egypten. Oder, wenn du sagen wolltest, daß Ismael todt war, hat er gewohnt etc., muß solches von seinen Nachkommen verstanden werden. Wie hernach eine gleiche Weise zu reden folgen wird, nämlich, Jakob ist hinabgezogen in Egypten und ist aus Egypten wieder heraufgezogen; und im 114. Psalm V. 2.: „Da ward Juda sein Heiligthum“ etc. Daselbst werden die Eigennamen verändert in Gattungsnamen und Völkernamen, das ist, in gemeine Wörter und die ein sonderlich Volk bedeuten, welches in der heiligen Schrift gemein ist. Darum halte ich, daß eben derselbe Verstand sei mit dem droben im 16. Cap. V. 12.: In facie omnium fratrum suorum habitabit, das ist: „Er wird gegen allen seinen Brüdern wohnen.“ Dasselbe wird an diesem Orte gesagt: Corruit, vel cecidit, vel fixit tabernaculum, er ist niedergelassen oder hat eine Hütte aufgerichtet, daß, da der Vater gestorben ist, und ein Testament hinterlassen hat, darin er seine Kinder und Brüder ermahnt, daß sie Gott fleißig dienen wollten: da sind sie zugehöret, und haben greulich wie die unvernünftigen Thiere mit Gewalt und mit dem Schwert getobt und gewüthet, haben das Land eingenommen und sind nephilim, das ist, Tyrannen geworden; wie an einem andern Ort gesagt wird.

\section[Zweites Stück]{Zweites Stück. Von Isaaks Verheirathung und der Rebecca Unfruchtbarkeit}

\begin{quote}
\textbf{V. 19. 20. Dies ist das Geschlecht Isaaks, Abrahams Sohns: Abraham zeugete Isaak. Isaak aber war vierzig Jahr alt, da er Rebecca zum Weibe nahm, die Tochter Bethuel, des Syrers, von Mesopotamien, Labans, des Syrers, Schwester.}
\end{quote}

41. Ich habe oft ermahnt und muß es auch den Leuten immer eingeprägt werden, nämlich, daß in den Historien der Väter das Allerkräftigste und Lieblichste sei, zu sehen, wie sie beschrieben und abgemalt werden als wahrhaftige Menschen, die da schwach oder gebrechlich und uns allenthalben gleich gewesen sind, und doch wiederum unter solcher menschlichen Schwachheit heilige Engel und auch Kinder Gottes gewesen sind. Denn was sehr wunderlich ist, im Stande der Haushaltung sieht man an ihnen gar keinen ungewöhnlichen oder sonderlichen Schein der Heiligkeit; und wo das Fleisch, das ist, die Weltweisen und Mönche solche Schwachheit sehen, ärgern sie sich daran gar sehr und verachten die heiligen Patriarchen hochmüthig, wie Augustinus von sich selbst bekennt, daß er die Historien Isaaks und der andern Väter belacht habe, da er noch ein Manichäer gewesen. Denn er hat auf nichts Anderes gesehen, denn nur auf den gemeinen Stand, als, daß Einer ein Weib hat, Kinder zeugt, ein wenig Schafe und Rinder hat, und mit andern Bürgern und Nachbarn umgeht und Gemeinschaft hat. Was könntet du daraus Treffliches oder Sonderliches lernen? oder, warum werden doch solche Historien gelesen und den Leuten vorgehalten, so doch anzusehen sind, als ob darin fast geringe und unnütze Dinge begriffen seien?

42. Dies ist zwar ein Jammer, damit das Fleisch beladen ist, daß es also soll hangen bleiben an der gemeinen Schwachheit, damit die Väter andern Menschen gleich sind, und daß sie darum geärgert werden und das gemeine Leben verachten, und ein ander, sonderlich Leben erdenken und erwählen, als, den ehelichen Stand, Möncherei und Pfafferei etc. Denn der Gottlose soll die Ehre und Herrlichkeit Gottes nicht sehen, sondern nur die Schwachheit und was thöricht ist, und (daß ich also sagen möge) die Nichtigkeit Gottes: die Ehre aber und Majestät, die Gewalt und Weisheit Gottes soll er nicht sehen, ob sie ihm wohl vor Augen gestellt werden. Wie Mose allhier auf das allereinfältigste erzählt, daß Isaak, da er vierzig Jahre alt gewesen, ein Weib, ja ein schwaches Weiblein genommen habe aus Mesopotamien Syrien. Was ist doch das? Werden nicht andere Leute auch ehelich mit einander, deren Ehestand diesem gleich oder auch wohl etwas herrlicher ist und größeren Schein hat? Warum wird denn dies gelesen? Antwort: Das Fleisch kann an den Heiligen sehen, daß sie Menschen und schwach gewesen sind, aber die Gottheit und Heiligkeit der Engel kann es an ihnen nicht sehen, auf daß es nur geärgert werde, und sich dünken lasse, es habe nun gute Ursache gefunden, neue Gottesdienste zu erdenken, darein es die Heiligkeit setzen möge.

\subsection{I}

43. Es ist aber nicht vergeblich oder ohne Ursache, wenn die heilige Schrift sagt, daß Isaak, da er vierzig Jahre alt gewesen, Rebecca zum Weibe genommen habe. Denn sie zeigt damit an, daß er in der ersten Brunst seiner Jugend kein Weib genommen habe, sondern daß er eine Zeitlang im Kampf und Siege gestanden sei wider das Fleisch und den Teufel. Denn die Historien und Erfahrung aller Menschen zeugen, wie sich die Unzucht mit so großer Ungeduld pflegt zu regen in der ersten Jugend, da die Kraft und Empfindung des Fleisches erst anfähet, und wird das männliche und weibliche Geschlecht gegen einander entbrannt.

44. Das ist eine gemeine Klage des ganzen menschlichen Geschlechts; und die solcher ersten Hitze nicht widerstehen, und nicht bedenken, daß sie etwas leiden müssen, dieselben fallen in Hurerei, Ehebruch und greuliche Unzucht, oder aber wo sie keuschfertig und ohne Rath Weiber nehmen, bringen sie sich selbst in ewige Marter und Plage. Derhalben hat Isaak diesen Kampf ertragen, und hat mit solcher Flamme und seinem Fleisch ganz hart gestritten; denn er ist ja ein wahrhaftiger und bloßer Mensch gewesen, wie wir auch sind.

45. In dieser letzten Zeit der Welt will unsere Jugend solche Last nicht tragen und will auch nicht eine Zeitlang Geduld haben. Darum wenn sie in der ersten Hitze Weiber nehmen, gönnt der Teufel, der sie zuvor mit Unzucht angezündet und dazu gereizt hat, und macht sie darnach mit einem widerwärtigen Weibe so gar kalt, daß sie dem Weibe auf das äußerste feind werden. Das ist ein recht teuflisch Ding. Darum sollte man zuvor mit den Exempeln der Väter das Herz unterrichten, daß es den ersten Kampf wider das Fleisch annehmen und ertragen könnte. Das reifere und männliche Alter hat darnach auch seine eigenen Kämpfe, die etwas größer sind. In der Jugend fängt man erstlich mit der Liebe an, wie in den Comödien solches an den Junggesellen beschrieben und abgemalt wird. Die Historien aber der heiligen Schrift halten uns solche Exempel vor, darin zugleich der Sieg und Kampf wider das Fleisch begriffen und uns vorgelegt werden. Also wird Isaak auch die Flammen der Unzucht gefühlt haben wie andere Jünglinge; er ist aber von seinem Vater unterrichtet worden, daß man dagegen kämpfen müsse, nämlich dadurch, daß man die heilige Schrift fleißig lese und Gott anrufe, darnach fleißig arbeite, ein fein mäßig, eingezogen Leben führe und faste.

46. Dies sollte die Uebung der Junggesellen sein auf gewisse ein Jahr oder zwo, wo sich etliche nicht länger enthalten könnten, daß sie doch lernen möchten, was das sei, daß man solche Flammen der Unzucht mit Geduld ertragen könne; denn dasselbe ist auch eine Geduld und Marter, wie Etliche vielerlei Art der Marter machen und zählen diese auch darunter, nämlich einen reichen, milden und züchtigen Jüngling. Und derselbe ist auch wahrlich ein rechter Märtyrer; denn er wird täglich mit vielem Leiden seines Fleisches gekreuzigt und gemartert.

47. Derhalben sollen die Junggesellen sich vor Unzucht und unzimlicher Vermischung hüten, und auf daß sie ihre Keuschheit bewahren und erhalten mögen, sollen sie ihre Herzen stärken mit Lesen und Betrachtung der Psalmen und des Wortes GOttes wider die grimmigen Anfechtungen des Fleisches. Wenn du die Flammen fühlst, so nimm einen Psalmen oder ein Capitel oder zwei aus der Bibel und lies; wenn aber die Flamme gestillt ist, alsdann bete fleißig: wo es sich aber nicht so bald stillt, so leide es mit Geduld und halte dich männlich ein Jahr, zwei oder mehr und halte immer an mit dem Gebet.

\quellseite{19}

48. Wo du aber die hitzigen Flammen des Fleisches nicht länger ertragen oder überwinden kannst, so bitte GOtt, daß er dir ein fromm Weib geben wolle, mit welchem du freundlich und in rechter Liebe leben mögest. Ich habe selbst ihrer Viele gesehen, die ihren bösen Lüsten nachgegangen und den Zaum gelassen haben, und sind darüber in schändliche Unzucht gefallen; aber endlich sind sie auch erbärmlich gestraft worden und ist die Strafe auf ihnen liegen geblieben; oder aber, wo sie in ihrer Blindheit zugeplagt und sich zum Ehestande begeben haben, ist es ihnen also geraten, daß sie Weiber bekommen haben, so ungeeignet und ihnen gar nicht gehorsam gewesen. Und ist ihnen daran eben recht geschehen.

49. Denn dies sollten sie ja alle gewußt haben, daß sie dazu berufen sind, daß sie wider das Fleisch streiten sollen, welches der eine Streit oder Kampf ist. Der andere ist wider den Teufel; der dritte wider die Welt. Darum muß man denselben ersten Anreizungen nicht so bald weichen und sich von ihnen überwinden lassen; vornehmlich jetzt zur Zeit, da den jungen Leuten immer die Hoffnung des Ehestandes vorgehalten wird. Dieselbige Hoffnung haben wir im Papstthum nicht gehabt; denn wer daselbst hat Priester werden wollen, ist gezwungen worden, daß er hat müssen ewige Keuschheit geloben. Dieselbige päbstische Tyrannei ist nun jetzt aufgedeckt und geoffenbart worden, und ist die rechte Freiheit wiederum aufgerichtet. Darum lerne, daß du fleißig betest und wider das Fleisch streitest; darnach aber bitte GOtt, daß er dir eine Rebekka geben wolle und keine Hagar, oder Eine, die noch ärger sei. Denn ein fromm Weib bekommt man nicht ohngefähr und ohne göttliche Schickung, sondern es ist eine Gabe GOttes, kommt nicht von unserm eigenen Rath oder Willen, wie die Heiden meinen.

50. Auf solche Weise ist Isaac unterrichtet gewesen, der vierzig Jahre gelebt hat, ehe er ehelich geworden ist. Dasselbe wird ohne Zweifel nicht ohne Plage und Flammen zugegangen sein, so er von seinem Fleisch gehegt; denn das Fleisch streitet immer wider den Geist nicht weniger im Hausregiment, als auch im Welt- und Kirchenregiment. Er hat aber seinem Vater Abraham gehorcht, der ihn unterrichtet und gelehrt hatte, wie er die Gebote GOttes betrachten sollte, und sich rüsten mit der heiligen Schrift zu dem Kampff, so ihm endlich begegnen würde. Darum hat ihm GOtt hernach Rebekka gegeben, mit welcher er in gutem Frieden gelebt hat. Dasselbe will die heilige Schrift also verborgener Weise angezeigt haben, und in solcher Schwachheit hält sie uns Isaac vor als ein trefflich Exempel der Keuschheit, so junge Leute haben sollen. Welches ein groß Ding ist; denn es ist ein Streit, so junge Leute führen müssen. Und in solcher Keuschheit, deren sich Isaac geflissen, wird auch mit begriffen, wie er sei auferzogen und unterrichtet worden; durch welchen Unterricht er gelernt hat, sich höher Gesellschaft zu enthalten. Und wird damit auch angezeigt, daß er fleißig gewesen, GOttes Wort betrachtet, gebetet und sich mit Arbeit geübt hat. Das alles scheint allhier verborgener Weise in den vierzig Jahren, darin er ohne Weib gelebt hat.

51. Es scheinet aber, daß es zu derselben Zeit also gebräuchlich gewesen sei, daß die Junggesellen nicht haben Weiber genommen, ehe sie vierzig Jahre alt geworden. Die Jungfrauen aber haben Männer angenommen, wenn sie zehn Jahre jünger gewesen sind; wie wir droben (Cap. 17, 17.) an Sarah gesehen haben. Und ich halte, daß Rebekka auch dreißig Jahre gewesen sei. Nach der Sintfluth haben schier die Väter alle um ihr dreißigstes oder vierzigstes Jahr Weiber genommen, vor der Sintfluth aber etwas langsamer. Denn da haben sie gewartet bis an die hundert, achtzig und neunzig Jahre; daher, nach der Sintfluth hat GOtt geeilt, daß das menschliche Geschlecht möchte gemehret werden: darum ist daselbst die Zeit verkürzt worden, also daß die Männer im vierzigsten und die Weiber im dreißigsten Jahre sind ehelich geworden. Derhalben ist dieselbige Zeit viel besser und trefflicher gewesen, denn unsere Zeit jetzt ist.

52. Wir meinen, diese Plage solle damit gelindert werden, wenn wir unsere böse Lust befriedigen mit Hurerei und Ehebruch: aber also werden aus den Menschen gar unvernünftige Thiere, und werden zu allen guten Werken ungeschickt und gar untauglich. Wenn sie leichtfertig zwolagen und ohne gewisse Ordnung, von GOtt uns vorgeschrieben, ehelich werden, alsdann findet es sich, daß sie nicht Weiber nehmen, sondern sich selbst damit Strafe und ewige Plage auf den Hals laden; denn sie haben GOtt nicht angerufen und ist bei ihnen keine Gottesfurcht gewesen. Davon sagt aber GOttes Gebot 2 Mos. 20, 7. also: „Du sollst den Namen des Herrn, deines GOttes, nicht mißbrauchen“; item Psalm 50, 15.: „Rufe mich an in der Noth; so will ich dich erretten, so sollst du mich preisen.“

53. Ja, so möchte Einer sagen, es ist aber solcher Verzug verdrießlich und unleidlich. Es ist recht geredet. Darum habe ich droben gesagt, daß es verglichen wird mit andern Uebungen der Geduld, so die Heiligen gehabt haben, wie denn ist, Hunger, Gefängniß, Frost, Krankheit, und da sie Geduld gehabt im Leiden und Marter, und was sie mehr Verdruss oder Beschwerung haben tragen können. Also ist die Anfechtung der Unzucht eine schwere Seuche und Last; aber man muß derselben widerstehen und dawider streiten. Wie hernach, wenn du wirst in der Polizei zween weltlichen Regiment kommen, da wirst du mit anderer Beschwerung umgetrieben und versucht werden, als, mit Dieberei, Räuberei und sonst mit mancherlei Bosheit der Menschen. Im Kirchenregiment wirst du streiten müssen mit den Ketzern, mit dem Teufel, welcher sich wider den Glauben, die Hoffnung und die Liebe GOttes legt. Aber da hast du das Wort, du hast die heilige Schrift, hast deine Studien, deine Uebung und Arbeit in deinem Beruf; daher wird der Glaube gemehret und gestärket werden. Also wird die Anfechtung der Unzucht, wo man sie überwindet mit dem Gebet, auch dienen, daß dadurch der Glaube und das Gebet gemehret wird. Derhalben, sage ich, hält uns die heilige Schrift in diesen (wie es scheint) vergeblichen Worten vor ein groß Exempel der Keuschheit, so Isaac gehabt, und dazu der schönen Zucht, so in der Gemeinde und im Hause Abrahams gewesen ist.

\subsection{II}

\begin{quote}
\textbf{V. 21. Isaac aber bat den HErrn für sein Weib, denn sie war unfruchtbar. Und der HErr ließ sich erbitten, und Rebekka, sein Weib, ward schwanger.}
\end{quote}

54. Dies ist nun eine andere Anfechtung, nachdem die Flamme der Unzucht aufgehört, und Isaac nun ist ehelich geworden und Rebekka zwanzig Jahre zum Weib gehabt. (Denn so lange hat GOtt die Verheißung verzogen, darin er seinem Vater Abraham verheißen und gesagt hatte 1 Mos. 21, 12.: „In Isaac soll dir der Same genannt werden.“) Da folgt nun ein ander Kreuz, welches zwar noch viel schwerer ist, denn die vorige Anfechtung war. Denn Isaac, der Ueberwinder der Unzucht, hat mit seiner Keuschheit den Teufel überwunden, bis er nun in den Ehestand gekommen ist. Im Ehestande begehrt er nun Kinder laut der Verheißung, und hat wahrlich darauf nicht wenig Hoffnung, dieweil er weiß, daß ihm durch göttlichen Rath und Schickung das Weib bescheret und zugeführt ist worden. Rebekka aber gebieret nicht, und hat auch keine sonderliche Verheißung, daß sie eine Mutter sollte werden, wie erstlich Sarah auch nicht gehabt. Das hat ihm ohne Zweifel in seinem Herzen wehe gethan und ist betrübt darüber geworden. Und ist zu solcher Anfechtung das auch noch gekommen, daß er sich geschicket und geforget hat, es würde die Unfruchtbarkeit also für und für währen und auf Rebekka liegen bleiben, welche Unfruchtbarkeit sie für einen Fluch gehalten haben. Denn die Väter haben sehr auf das Wort gedrungen, da GOtt 1 Mos. 1, 28. gesagt hatte: „Wachset und mehret euch“ etc., und haben das Kinderzeugen für einen sonderlichen göttlichen Segen gehalten; und diejenige, so nicht fruchtbar waren, haben sie für verfluchte Leute gehalten und die unter GOttes Zorn wären.

55. Darum kann man leichtlich abnehmen, wie schwerlich und hart Rebekka deshalb sei geplagt worden, und welch ein großer Schmerz den Isaac das gewesen sei, da er gesehen hat, daß sein Weib daran gar verzagt hat, daß sie würde Kinder gebären; denn sie ist dazumal bei fünfzig Jahren gewesen. Da wird sie gedacht haben: Nun werde ich veralten und untüchtig werden, Kinder zu gebären. Isaac mag noch eine Hoffnung gehabt haben, nämlich, wo Rebekka nicht würde gebären, so wollte er ein ander Weib nehmen, wie Abraham, sein Vater, gethan hatte.

\quellseite{20}

56. Diese Hoffnung war der Rebekka genommen, und hat deshalb alle Jahre und Tage mit großen Sorgen gezählt, welche verlaufen waren von der Zeit an, da sie mit Isaak war ehelich geworden, ob sie vielleicht noch der Jahre und Alters halben Hoffnung haben möchte. Darum ist diese Anfechtung viel verdrießlicher gewesen denn die vorige, weil die arme Rebekka der Meinung ist, daß sie auch unter die Weiber mit gezählt sei, so des Segens Gottes beraubt sind.

57. Was soll sie aber nun thun, dieweil sie sieht, daß es vergeblich ist, daß sie sich also gefehlt und Kinder begehrt hat? Sie ist zwanzig Jahre unfruchtbar gewesen und wird nunmehr eine verstorbene Frau; denn es war vorhanden das Jahr und die Zeit, da ihr Leib ersterben und zur Geburt würde untüchtig werden. Darum wird sie ohne Zweifel ihren Mann gebeten haben, daß er wollte für sie bitten. Diese einzige und äußerste Hülfe hat sie erfunden. Sie will ihm traun kein ander Weib geben, will auch der Mutterehre sich nicht berauben lassen, wie Sarah droben 1 Mos. 16, 3. gethan, die ihrem Manne die Magd Hagar zum Weibe gegeben hat. Darum sagt Mose: „Und Isaak bat den Herrn für sein Weib“ etc.

58. Hier wird sich nun abermals das Fleisch ärgern, und dies alles verachten als schlechte, gemeine und geringe Dinge. Denn was ist das so groß Ding, daß der Mann für das Weib bittet, welches unfruchtbar ist! sind doch sonst viel Weiber, die schwanger werden, daß man darum nicht bittet; ja, solches geschieht bisweilen wider Begehren und Willen Vieler, die nicht begehren Kinder zu haben. Aber siehe du mir an die große, treffliche Beständigkeit des Glaubens, die Geduld und Hoffnung, welche alle das Fleisch nicht sieht, alsdann wirst du finden, daß du dich billig verwundern magst; denn sie hat nicht ohne großen Schmerz und Bekümmerniß denken können an die Verheißung, daß von Isaak sollten Kinder geboren werden. Und während die andern Weiber zwar alle gesegnet waren, die weder gebetet noch Verheißungen gehabt haben, sie aber allein ohne alle Hoffnung, daß sie würde Kinder zeugen, hat Leben müssen, und also die Zeit ihres ehelichen Lebens in großer Traurigkeit und mit vielen Thränen zugebracht hat: behält sie doch noch den Glauben, und ermahnt ihren Ehmann mit großer Beständigkeit, daß er den Herrn für sie bitten wolle.

59. Wo unser Einer mit solchem Kreuz beladen wäre, oder aber so viel Jahre mit anderem Unglück, Krankheit, Elend, oder Gefängniß sollte geplagt werden, und er darüber nicht murren und nicht ungeduldig würde, sondern im Glauben und Hoffnung beharren und beständig würde bleiben, derselbe würde sehen, was Rebekka erlitten hat. Das Fleisch sieht nur, was äußerlich ist und was die Haushaltung belanget, nämlich, daß sie häusliche und tägliche Arbeit thut, daß sie bei ihrem Manne schläft: aber es sieht das nicht, daß sie so geduldig gewesen ist, daß sie die ganzen zwanzig Jahre geseufzt und geweint hat.

60. Denn solche treffliche, große Tugenden sind vor den Augen dieser Welt verborgen, als nämlich sind Geduld, Glaube, und daß man mit Geduld harret, wenn es sich mit der göttlichen Verheißung verziehet. Dies aber heißt das Fleisch an und verwundert sich dessen höchlich, wo es einen Mönch in einem grauen Rocke sieht daher gehen, der einen Strick umgürtet hat, der kein Fleisch ißt, hat aber keinen rechten Glauben, keine Geduld, kein recht Kreuz, und, in Summa, der Tugenden keine, so Rebekka gehabt. Warum aber das? Ei darum, daß die Welt blind ist.

61. Wir aber sollen uns gewöhnen an solchen Kampf, da immer die Göttlichen Eines nach dem Andern erwarten müssen, und sollen lernen glauben und Geduld tragen, und daß wir auch beharren, daß wir nicht murren oder von der Verheißung nicht abfallen, sondern daß wir stark und männlich seien, und streiten wider die Ungeduld und feurigen Pfeile des Teufels, welcher unsere Herzen dazu treibt, daß sie ungeduldig werden und mit Gott zürnen sollen, damit er uns den Glauben und die Geduld im Unglück aus dem Herzen treibe. Wir sollen uns das Exempel Isaaks und Rebekkas vorhalten. Die haben beide zwanzig Jahre geharrt, und haben inzwischen gesehen, wie die Gottlosen so glücklich und fruchtbar gewesen sind, welche sie verspottet und gelästert; werden gesagt haben: Warum hat er eben diese genommen? hat eine Fremde wollen zum Weibe haben. Warum hat er nicht etwa aus unsern Geschlechtern eine ehrliche Jungfrau genommen? Rebekka ist verloren und von Gott verworfen. Solche Läster- und Schmähworte wird er ohne Zweifel haben hören müssen, nicht ohne großen Schmerz seines Herzens, und nicht ohne viel Thränen, so er darüber vergossen wird haben. Wie droben Sarah auch ihre Unfruchtbarkeit beweint hat, und 1 Sam. 1, 10., um derselben Ursache willen auch jämmerlich heult und schreit; und hat doch gleichwohl solches mit Geduld und starkem Glauben überwunden.

62. Diese Tugenden sollen wir rühmen und in den Historien der Väter fleißig auf solche Exempel der Geduld sehen. Denn am schwersten ist dieser Kampf wider die göttliche Verheißung, dadurch überwunden wird der Zorn, das Murren und die Ungeduld über das Verziehen Gottes. Denn das ist unseres Herrn Gottes Eigenschaft, und er wird recht genannt expectatus, das ist, auf den man harren soll; wir aber heißen expectantes, das ist, die da harren und warten sollen. Diese Namen sollen immer vor unsern Augen und Gemüth stehen, auf daß wir lernen den ersten Zorn, dazu wir bewogen werden, brechen, und daß wir nicht sobald ungeduldig werden, wenn uns Gott ein, oder zwei, oder mehr Jahre hinhält und harren läßt; sondern sollen bedenken, daß man beharren müsse, und mit starkem Muth alles überwinden, damit unsere Geduld angefochten wird. Wie Rebekka hat lernen verachten, da sie von andern Weibern und vielleicht auch wohl von ihrem eigenen Hausgesinde ist geschmäht worden, bis sie endlich mit ihrem und ihres Mannes Gebet den Herrn überwunden hat.

63. Das hebräische Wort athar hat einen großen Nachdruck; denn es ist ein sonderlich Wort des Gebetes, und heißt so viel, als wenn man gleichsam zur Unzeit und ungeheuer betet, also daß man mit Anklopfen und heftigem, feindseligem Anhalten Gott einen Verdruss bereitet. Wir heißen es im Lateinischen exorare, das ist, erbitten. Denn erstlich muß man bitten; zum Andern, suchen, und zum Dritten, anklopfen. Wenn wir schreien werden: Herr Gott, hilf mir in dieser Noth, erlöse uns von diesem oder anderem Unglück! und die Erlösung folgt nicht alsbald, alsdann soll man suchen alle Exempel der Väter: Siehe doch, lieber himmlischer Vater, wie du zu allen Zeiten bei deinem Volke gestanden und demselben geholfen hast. Wo er alsdann noch verzieht, sollst du doch nicht aufhören zu bitten, sondern sollst sagen: Ich will nicht ablassen, will nicht aufhören zu klopfen, sondern will schreien und anklopfen bis zum Ende meines Lebens. Also wird Rebekka ihren Mann ermahnt haben: Mein lieber Isaak, werde nicht müde, höre nicht auf. Und Isaak hat ihr Weinen und Seufzen angesehen und hat den Herrn erbeten.

64. Daraus soll man nun lernen, wenn wir beten, daß wir auch gewißlich sollen erhöret werden; wie denn wahrlich bis daher die Kirche mit ihrem Gebet Frieden erlanget und den Türken sammt dem Pabst zurückgehalten hat. Wir sollen uns aber allein davor hüten, daß wir nicht bald müde werden, wenn wir einmal angefangen haben zu beten; sondern sollen suchen, und alle Sorge, Unglück und Kreuz auf Gott werfen und ihm die Exempel allerlei Erlösung vorhalten. Und endlich sollen wir mit herzlichem Vertrauen und mit stetigen Schlägen an die Thüre klopfen. Da werden wir erfahren, daß St. Jacobus sagt im 5. Cap. v. 16.: „Des Gerechten Gebet vermag viel, wenn es ernstlich ist“; denn es dringt durch Himmel und Erde, Gott kann alsdann unser Schreien nicht länger leiden, wie Lucä 18, 5. von dem ungerechten Richter und der Wittwe gesagt wird. Man muß aber nicht allein Eine Stunde bitten, sondern man muß schreien und anklopfen; alsdann wirst du ihn zwingen, daß er komme und dir helfe. Also habe ich eigentlich das Vertrauen, wenn wir mit dem Gebet ernstlich und brünstig werden anhalten, daß wir Gott erbitten würden, daß der jüngste Tag kommen würde.

65. Auf dieselbe Weise hat Rebekka auch ihre Zuflucht gehabt zum ernstlichen Gebet, damit man immer anhält, und hat Tag und Nacht ängstlich geseufzt. Isaak hat auch für sie gebeten, und hat Gott Anderes nichts vorgehalten, denn allein diese einige Noth, nämlich, die Unfruchtbarkeit seines Weibes. Daraus wir lernen sollen, daß man Gott allerlei Noth, auch die leibliche, vorhalten soll, vor allen Dingen aber die geistliche Noth. Isaak
\quellseite{21}
wird also gebetet haben: Lieber Herr Gott, wo es zur Heiligung deines Namens dient, wo es zu Erhaltung deines Reichs gehört, so gib der Rebekka Kinder. Wo keine sonderliche Verheißung ist, wie Rebekka keine gehabt, da muß das Gebet solches erstatten und daselbst zu Hülfe kommen.

66. Es ist aber ein schwer Ding und ist eine große Arbeit, und viel schwerer, denn die Predigt des Wortes oder sonst andere Kirchenämter sein mögen. Wenn wir lehren oder predigen, so leiden wir mehr, denn daß wir etwas thun sollten: Gott redet durch uns und ist ein göttlich Werk; aber beten ist das allerschwerste Werk, darum ist es auch so gar selten. So ist es nun ein groß Ding, daß Isaac seine Augen und Hände hat dürfen aufheben zu der göttlichen Majestät, daß er bittet, sucht und anklopft; denn es ist ja ein sehr groß Ding, mit Gott reden. Es ist wohl auch ein groß Ding, wenn er mit uns redet, aber dies ist etwas schwerer; denn es kommt dazu unsere Schwachheit und Unwürdigkeit, die uns aufhält und zurückzieht, daß wir denken: Wer bin ich, der ich meine Augen darf aufheben und dazu meine Hände zu der göttlichen Majestät, da die Engel sind und auf dessen Wink die ganze Welt erzittert? Soll denn ich armer Mensch zu demselben treten, und sagen: Dies will ich haben, und bitte, du wollest mir es geben?

67. Der gemeine Haufe der Mönche und Pfaffen weiß davon nichts, und wissen auch nicht, was beten sei; wiewohl etliche gottselige solche Gedanken etwas leichtlicher überwinden. Aber das Gebet, so recht kräftig und heftig ist, und das durch die Wolken dringen soll, ist wahrlich ein sehr schwer Ding; denn der ich Asche und Staub und voller Sünden bin, soll den lebendigen, wahrhaftigen und ewigen Gott anreden. Darum ist es nicht Wunder, daß der, so da betet, erzittert und zurück läuft. Wie ich ehemals that, da ich noch ein Mönch war, und zum erstenmal im Canon der Seelmesse diese Worte lesen sollte: Wir bitten dich demüthiglich, gnädigster Vater etc.; item: Wir opfern dir, dem lebendigen, wahrhaftigen und ewigen Gott etc.: da entsetzte ich mich, und erschrack ganz und gar derselben Worte halben; denn ich dachte also: Wie komme ich dazu, daß ich die hohe Majestät anreden soll, da sich doch sonst alle Menschen entsetzen müssen, wenn sie irgend einen Fürsten oder König ansehen oder anreden sollen?

68. Der Glaube aber, so sich an Gottes Barmherzigkeit und Wort hält, der steigt und überwindet diesen Schrecken, wie er in Isaac auch gesiegt hat, welcher an aller menschlichen Hülfe hat verzagen müssen; denn niemand kann einem unfruchtbaren Weibe helfen. Darum faßt er einen Muth und richtet ein heißes und heftiges Gebet zu Gott. Solchen trefflichen Muth und großen Glauben sieht das Fleisch nicht. Dies wird aber um unsertwillen geschrieben, auf daß wir auch kühn und getrost seien, und beten lernen. Denn das Gebet der Gläubigen kann nicht vergeblich sein; wie Isaac auch nicht vergeblich gebetet hat, sondern wie Mose sagt: „Der Herr ließ sich erbitten.“ Also wird der Herr unser Seufzen und Schreien auch nicht verachten. Lasset uns nur aufwachen zum Gebet.

69. Die Juden erregen allhier eine Frage von dem Wort noch. Dasselbe heißt bei den Hebräern eigentlich so viel, wie stracks vor sich; wie in 57. Cap. Jesaiä 2. stehet: Der da wandelt, nechocho (wie es im Hebräischen lautet), das ist, der sich weder zur Linken noch zur Rechten gewendet hat. Darum sagen die Juden, Isaac habe stracks vor seinem Weibe gebetet. Und wo dasselbe ein gemeiner Gebrauch ist gewesen zu beten, so wird Rebekka stracks vor ihm gestanden sein, oder ist auf die Kniee gefallen, auf welche er die Hände gelegt, da sie gemeint und \unsicher{gefürset} hat, und haben also zugleich mit einander den Herrn gebeten. Wenn ein solcher Gebrauch oder Weise gewesen wäre zu beten, so könnte man die rechte eigentliche Bedeutung desselben Wortes behalten: ist es aber nicht also gebräuchlich gewesen, so muß man es geistlich auslegen, nämlich, daß er von ganzem Herzen, und da er auf seines Weibes Elend gesehen, gebetet habe; gleich als wenn ich für jemand bete, stelle ich mir denselben in meinem Herzen vor, und sehe oder denke an nichts Anderes, schaue in meinem Herzen nur ihn allein an.

70. Also hat Isaac gebetet und hat sich sein Weib im Herzen vorgestellt. Damit Mose anzeigen will, daß es ein heiß und ernstlich Gebet gewesen sei, darin er nicht gezweifelt habe, und in seinem Herzen und Gedanken nicht umher irre gegangen sei. Gleichwie man sagt vom Gebet Eines, der mit St. Bernhard gewettet, er wollte das Vater unser ohne alle Gedanken, und daß ihm nichts Fremdes einfallen sollte, ausbeten: da sie aber einen guten Hengst aufgesetzt, den es gelten sollte, und das Gebet nun aus war, und ihm derselbe gute Freund frei heraus die Wahrheit sagen sollte, wie sie darüber mit einander eins geworden waren, da hat er bekannt, daß ihm, dieweil er gebetet habe, eingefallen sei, wo er das Pferd gewönne, ob ihm auch Sattel und Zaum dazu gebühren wolle oder nicht. Also soll das Gebet der Gottseligen nicht sein; denn solch Gebet geschieht nicht stracks zu Gott, sondern das Herz geht umher spazieren, jetzt zur Rechten, dann zur Linken. Aber ein recht wahrhaftiges und heißes Gebet trägt Gott seine Sachen vor, und sieht mit großer Begierde und herzlichem Verlangen nur darauf, und läßt sich keinerlei Vermessenheit oder Verzweiflung irre machen; sondern sagt: Lieber Herr Gott, siehe doch dies arme betrübte Weib an und gedenke an deine Verheißung. Denkt an nichts Anderes, bekümmert sich auch um nichts Anderes. Und dies ist das ernstliche, anhaltende Gebet des Gerechten, das so einig und stracks zu Gott gerichtet ist; davon in der Epistel Jacobi im 5. Cap. V. 16. gesagt wird.

\section[Drittes Stück]{Drittes Stück. Von Rebekkas Schwangerschaft, wie sie dabei angefochten worden, den Herrn um Rath gefragt und was sie für Antwort erhalten}

\subsection{I}

\begin{quote}
\textbf{V. 21. Und Rebekka, sein Weib, ward schwanger.}
\end{quote}

71. Dies ist ein sehr trefflicher und merkwürdiger Text, welcher St. Paulus in der Epistel an die Römer Cap. 9. V. 10—13. herrlich handelt, da er sagt: „Nicht allein aber ists mit dem also, sondern auch, da Rebekka von dem einigen Isaac, unserm Vater, schwanger ward; ehe die Kinder geboren waren, und weder Gutes noch Böses gethan hatten, auf daß der Vorsatz Gottes bestünde nach der Wahl, ward zu ihr gesagt, nicht aus Verdienst der Werke, sondern aus Gnaden des Berufers, also: Der Größere soll dienstbar werden dem Kleinern, wie denn geschrieben stehet: Jacob habe ich geliebet; aber Esau habe ich gehasset.“ Alhier wollen wir aber gerne von St. Paulo aufstehen, ihm die Ehre erzeigen und die Meisterschaft zuschreiben, daß er allein diesen Text nach seiner Würde habe können auslegen und erklären; denn wir können keine andere oder bessere Auslegung daher bringen. Und diesmal man noch etwas hierbei sagen muß, so wollen wir wie unvollkommene Schüler dem vollkommensten Meister nachfolgen.

72. Ich hätte zwar diese Worte immermehr also ansehen können. Mose sagt also: „Rebekka ward schwanger.“ Wie aber? Nicht allein leiblich oder durch Kraft der Natur, sondern dadurch, daß der heilige Patriarch solches mit seinem Gebet erlangt hat. Und dies ist keine schlechte oder geringe Gnade; denn dieweil damals so viel Patriarchen gelebt und die Kinder gesehen haben, als nämlich, Abraham, Sem, Salah und Eber, wird doch aller dieser Gebet verschwiegen, und wird allein Isaac diese Ehre zugeschrieben als dem Erben der Verheißung, daß er gebetet habe, das ist, das priesterliche Amt verwaltet, da er Rebekka die Hände aufgelegt nach der Weise, so die Priester gebraucht haben. Darum ist diese Empfängnis nicht fleischlich oder allein natürlich, sondern Gott hat haben wollen, daß sie hat geschehen sollen durch das Gebet des heiligen Patriarchen, aus dem Glauben, Hoffnung und Liebe.

73. Darum lehrt nun dieser Text eben dasselbe, das Paulus in der Epistel an die Römer auch lehrt, da er einen Unterschied macht zwischen der Geburt aus dem Fleische, das ist, aus der Schöpfung, und der geistlichen Geburt. Denn die fleischliche Geburt hat Gott noch erhalten nach dem Fall Adams. Wiewohl die Natur verderbet war durch die Sünde und den Teufel, der Adam geheißen hat, daß er Gott sollte gleich werden, eben in solchem Falle, dadurch Lucifer auch aus dem Himmel gefallen ist: so hat doch Gott die Natur der Geburt oder Kinderzucht nicht verhaßet, sondern hat zugelassen, daß beide die Gottseligen und auch die Gottlosen Kinder gezeugt. Das aber hat er allein anzeigen wollen, daß es nicht genug sei, in diese Welt geboren werden aus dem Fleisch, sondern daß über solche Geburt, so noch in der Natur geblieben ist, vonnöthen sei eine Wiedergeburt und Erneuerung der Wiedergeburt durch den Heiligen Geist. Dasselbe zu beweisen hat der heilige Paulus diesen Text von Esau und Jakob gebraucht, und auch den von Isaak und Ismael. Denn sie alle beide sind geboren nach dem Fleisch und nach der ersten Geburt. Aber es wäre ihrer keiner in das Himmelreich gekommen aus dem menschlichen Reiche, wo nicht Isaak durch die neue Geburt zum Erben wäre gesetzt worden und desgleichen Jakob auch.

\quellseite{22}

74. Gott hat dieser vortrefflichen Personen dazu gebraucht, daß er dieselbigen zu großen, trefflichen Exempeln vorstellen möchte dem hoffärtigsten Volke, welches Abrahams Nachkommen würden sein. Denn er hat wohl gewußt, wie hoffärtig sie waren, daß sie niemand hat zähmen können, und daß sie solche Nacken hätten, die härter wären, denn irgend ein Diamant sein kann, nämlich um dieser Ehre willen, daß sie Kinder und Nachkommen Abrahams waren.

75. Darum hat er dies zur Vorbeugung in den Text setzen wollen, daß er damit den Juden ihr hoffärtiges Maul verstopfen möchte und ihnen damit ihr Argument von der fleischlichen Geburt zerstörete. Denn es ist an derselben mit nichten genug; sondern über solche fleischliche Geburt gehört auch die Wiedergeburt noch dazu. Wie heilig die Väter sind, Abraham, Isaak und Andere mehr, zeugen sie doch keine Kinder des Fleisches durch die erste Geburt, welche erstlich von Anfang der Welt geschaffen und durch den Teufel verderbet ist. Sondern über die Schöpfung muß auch noch der Beruf dazu kommen. Denn also sagt Paulus Röm. 9, 11. 12.: „Ehe die Kinder geboren waren, und weder Gutes noch Böses gethan hatten“ etc., „ward zu ihr gesagt, aus Gnaden des Berufers“, und nicht von dem Schöpfer; denn sie sind alle beide zugleich empfangen und geschaffen. Auf daß aber die ganze Welt und vornehmlich das hartnäckige Volk wissen möchte, daß über die Schöpfung auch noch der Beruf dazu gehöre, welches sonst die Verheißung genannt wird, ist zu ihr gesagt worden: „Nicht aus Verdienst der Werke, sondern aus Gnaden des Berufers: der Größere soll dienstbar werden dem Kleinern.“

76. Dies Argument und Materie ist vom Anfang der Welt für und für getrieben und gehandelt worden, wie der Zank zwischen den Brüdern Esau und Jakob sich in Mutterleibe angefangen hat. Und ist des Krieges noch kein Ende. Denn es ist eben derselbe Streit, so auch zwischen Kain und Abel gewesen ist, und zwischen den Nachkommen der Schlange und des Samens des Weibes. Die erste Geburt ist hoffärtig und erhebt sich über die Maßen sehr, vornehmlich wo das dazu kommt, daß die Eltern auch ein sonderlich Ansehen gehabt, daß sie vom Blut der Väter geboren sind. Darnach hilft auch etwas dazu Reichthum, Gewalt, das Reich dieser Welt, Weisheit, Gerechtigkeit, Religion oder Gottesdienst. Daher kommt die Feindschaft zwischen Kain und Abel, Ismael und Isaak, Esau und Jakob, zwischen unsers Herrn Gottes und des Teufels Kirche. Und hat des Teufels Kirche stets dazu Lust, daß sie herrschen und regieren möge: zum Ersten um des Geblütes und der Väter willen gleich als aus rechtmäßiger Ursache; darnach um des zeitlichen Segens willen. Wir sind, sagen sie, Kinder der Väter; wir sind groß und gemehret worden von wegen des Regiments, so auf uns liegt, dazu von wegen der Religion, Geldes und Gutes, auch vieler herrlichen Siege und großen Wunderwerke. Und endlich, sind wir also erhöhet, daß alle Völker wie der Schwanz sind, wir aber sind wie das Haupt. Das Argument scheint unüberwindlich und unumstößlich. Aber die Schrift sagt das Gegenteil: Nicht die vom Fleische geboren sind, nicht die der Väter Kinder sind, sind eben darum auch Gottes Kinder; denn die Kinder Gottes werden nicht von dem Geblüte, noch von dem Willen des Fleisches, noch von dem Willen des Mannes, sondern von Gott geboren, Joh. 1, 13.

77. Darum ist das ein sehr stark Argument, welches man gegen das, so droben (1 Mos. 17, 21. Cap. 21, 12.) ist gehört worden, halten kann, nämlich, daß Gott einen Unterschied gemacht hat zwischen Isaak und Ismael, zwischen Jakob und Esau, welche ja zugleich von Abraham und Isaak geboren sind: und kann doch gleichwohl Ismael kein Erbe sein, wenn er es auch schon gerne sein wollte, sondern der Same der Verheißung, der den Beruf, und über die erste auch die andere Geburt und die Wiedergeburt hat, derselbe wird vorgezogen.

78. Daher ist aber nun der Streit gekommen, so für und für währet vom Anfang bis an das Ende der Welt; nicht ein Streit über solchem Dinge, daß geringe und nichts werth sei, sondern über dem großen herrlichen Titel, nämlich, über der Kirche, Gottes Volk, über dem Himmelreich und ewigen Leben. Also müssen wir jetziger Zeit auch streiten mit des Papstes Kirche, welche Gottes Volk will sein, und will das Regiment und Priesteramt haben, rühmen sich, daß sie allein die Kirche seien, welche Gott recht als einen Vater erkennt und ihm auch recht dient: uns aber verdammen sie gar und verfolgen uns dazu als Ketzer und die Kirche des Teufels. Und das ist es auch, daß sich diese Kinder, Jakob und Esau, mit einander stoßen, ehe denn sie geboren sind worden. Denn es ist in der Welt zweierlei Kirche vom Anfange, gleichwie zweierlei Samen sind, nämlich, der Schlange Same und des Weibes Same, welche mit einander streiten und sich stoßen um des Titels willen der Kirche.

79. Aber St. Paulus hat uns hier eine sehr scharfe und starke Dialektik gezeigt und gelehrt, da er den Unterschied angezeigt hat zwischen der Geburt und dem Beruf. Wo allein die Geburt ist, da ist die Verdammniß; denn was vom Fleisch geboren ist, das ist Fleisch, Joh. 3, 6. Item, die nicht vom Geblüte geboren sind etc., Joh. 1, 13. Es ist ein Anderes, daß St. Paulus Röm. 9, 11. 12. sagt: „Es ist zu ihr gesagt aus Gnaden des Berufers“, das ist, es gehört Gottes Wort und die Verheißung dazu; man muß über die Creatur den hören, der da regieren und Gottes Sohn sein will, nicht als einen Gott, der da schafft, sondern als einen, der da beruft. Wo es nun an der ersten Geburt genug wäre, wozu bedürften wir denn Gottes?

80. Die Türken sind sehr ehrliche Leute, sehr weise und auf ihren Gottesdienst beflissen. Das Reich oder Regiment, so sie mit großer Mühe und Arbeit zuwege gebracht haben, erhalten sie auch mit ernster Zucht und Ordnung. Wenn nun das Volk Gottes ein solch Volk wäre, wie das türkische Volk ist, so wäre ihm nichts mehr vonnöthen. Aber das werden sie auch noch geziert von Gott mit Reichthum, Weisheit, Ehre, Vernunft und großem herrlichem Siege: aber was ist bei den Türken Anderes denn die erste Geburt? Denn die Vernunft ist vom Weibe geboren, und zu der Vernunft gehört die Weisheit, das weltliche Recht, Disciplin, Gesetze und äußerliche Ordnung; von dem allen sagt die heilige Schrift, daß es vor Gott verloren und verdammt sei.

81. Die Juden haben zu ihrer Zeit zugleich das Geblüt mit dem Ruhm des Fleisches gehabt: der Türke hat allein den fleischlichen Ruhm, die Juden aber rühmten sich des Geblütes halben, das ist, daß sie von den Vätern und Propheten geboren wären. Darnach ist der fleischliche Ruhm auch dazu gekommen, nämlich, das Regiment, die Gewalt, die Religion oder Gottesdienst, und der äußerliche Schein oder Heuchelei. Darum waren sie in allewege ruhmredig, nicht allein um der Ursache willen, daher solcher Ruhm entstanden, sondern auch um deswillen, daß daraus gefolgt ist. Die Ursache war diese, daß sie von sehr heiligen Eltern geboren waren; der Erfolg aber war, daß sie das Regiment, die Ehre, Geld und Gut hatten.

82. Der Türke hat den Ruhm vom Geblüte nicht, denn er ist nicht von den Vätern geboren, sondern hat das Regiment, Weisheit, Ehre: und mit solchen großen Werken untersteht er sich zu beweisen, daß sein Volk soll Gottes Volk sein. Aber hier wird gefragt: Wo ist denn der Beruf? Ei, sagen sie, man bedarf dessen nicht. Es ist genug an der ersten Geburt: wenn der Mensch thut, was an ihm ist, so wird er selig werden. Aber dem ist nicht also, sondern es gehört dazu der Beruf, das ist, das Wort. Wo es an der ersten Geburt genug wäre, wozu bedürften wir denn des Wortes? Warum berufst Gott vom Anfang bis zum Niedergange? Sind wir nicht gerecht, fromm, heilig und feine vernünftige Menschen? Aber wider das alles sagt Gott: Ich will es also haben, daß alles Fleisch sey, und alle seine Herrlichkeit oder Güte wie eine Blume auf dem Felde ist, Jes. 40, 6., auf daß sich niemand des Fleisches und Geblütes halben rühme, sondern sich an die Erneuerung halte, welche geschieht, wenn uns Gott aus Gnaden beruft; sonst wird aller Ruhm vergeblich und nichts sein, so beide, der anfänglichen Ursache halben, und auch von den Werken, so daraus erfolgt, hergekommen ist. Und das Volk, so sich allein auf diesen Ruhm verläßt und darauf seinen Trost setzt, wird ewig verloren und verdammt werden.

83. Wenn aber dies also gelehrt wird, so zürnen und toben die Juden, und werden schier rasend und toll über solcher Lehre. Denn sie sind gar unsinnig über der ersten Geburt, und können so weit nicht kommen, daß sie an die andere Geburt einmal denken, und werden darüber also verdammt. Dergleichen ist der Türke auch hoffärtig und stolz nur einfach der ersten Geburt halben und fragt nach dem Berufe gar nichts. Der Papst ist darin auch ganz und gar ersoffen, so auch, daß er alle Gottseeligkeit, Religion und Gottesdienst auf dieselbe erste Geburt gezogen hat, als, da er gelehrt hat, das Sacrament des Abendmahls sei ein solch äußerlich Werk, damit die Gottseeligen oder Christen ihren Gehorsam beweisen gegen die Kirche. Daselbst ist ja nichts, denn die Natur von der ersten Geburt, welche das äußerliche Werk an sich selbst, das ist, den Gehorsam der Kirche, verstehen kann. Darnach hat er die Welt voll Gesetze gemacht, welche die Früchte sind des bösen Baumes, wiewohl sie einen sehr trefflichen und schönen Schein haben; ich höre aber daselbst Gott nicht, der da berufe und etwas verheiße, sondern ich höre nur einen Menschen, der solche äußerliche Werke thut aus der ersten Geburt.

\quellseite{23}

84. Derhalben stoßen wir uns mit einander, wie in der ganzen Welt vom Anfange bis an das Ende geschehen ist. Wir sagen, daß sie nicht die rechte Kirche sind, dieweil sie sich allein der ersten Geburt rühmen: dagegen aber verdammen sie uns, die wir Gott glauben, der uns aus Gnaden berufen und seine Verheißung gegeben hat, und daher warten und hoffen wir selig zu werden. Wer soll nun hier Richter und Schiedsmann sein? Antwort: Das soll dieser Text thun: Nicht die aus dem Fleisch geboren, sondern die des Glaubens sind. Die Seligkeit zu erlangen, dazu gehört der Beruf und die Verheißung.

85. Wenn man fragt, ob auch die Türken sollen selig werden? soll man antworten: Sie werden nicht selig werden. Ist doch bei ihnen große Weisheit, ein sehr ehrbar Leben, mancherlei Gottesdienst, Ehrerbietung und Gehorsam gegen ihre Obrigkeit, eine treffliche Disciplin und Kriegsordnung? Ja, solches findet man zwar alles bei ihnen. Wohin gehört aber dies alles? Antwort: Es gehört zu der ersten Geburt, darum müssen sie verdammt werden und können nicht Kinder sein des Reiches; denn es mangelt ihnen am Beruf und an der Wiedergeburt, und sie wollen selig werden durch die erste Geburt, welche verderbet ist, wo sie nicht wiederum durch die neue Geburt aufgerichtet und zurechtgebracht wird. Darum ist es gewiß, daß die Türken die Kirche nicht sind, die Juden auch nicht; denn sie haben den Beruf nicht, wiewohl es ihnen weder an der anfänglichen Ursache, noch an den Dingen, so daraus folgen, mangelt, wie droben gesagt ist.

86. Also ist auch des Papistes Kirche nicht die rechte Kirche; denn sie gehen auch daher in der ersten Geburt, und sind so vermessen, daß sie gedenken durch ihre Werke die Seligkeit zu erlangen. Wo aber Gott über die erste Geburt auch den Beruf dazu thut, wie Paulus Röm. 9, 12. sagt: „Es ist zu ihr gesagt worden, aus Gnaden des Berufes“ etc.: daselbst wirst du die Kirche finden. Die Römer haben sehr schöne Rechte und Gesetze gemacht, wie sie der Poet rühmt, da er sagt: Tu regere imperio populos, Romane, memento: das ist: Ihr Römer, gedenket, daß ihr in eurer Herrschaft über die Völker regieret etc. Aber es ist alles nichts Anderes, denn eitel Vernunft, die aus dem Fleische vom Weib und Mann geboren ist; dieselbe ist gar unnütz und ist todt vor Gott. Derhalben wird kein Papist, oder Türke, oder Mönch, oder Jurist, oder Arzt (von sich selbst) selig werden. Denn alles, was in der ganzen Welt ist von den allerschönsten Dingen und wie nütze sie auch sind, das wird alles unter dem Einen Namen verdammt, dieweil es Fleisch und Herrlichkeit oder Güte des Fleisches ist, Jes. 40, 6.

87. Darum lehrt die rechte Kirche also: Wir sind nicht allein geschaffen und geboren vom Fleische, sondern Gott hat uns sein Wort gesandt und beruft uns dadurch, das ist, er predigt uns von Vergebung der Sünden, und daß er uns durch Jesum Christum zu Kindern habe angenommen. Wenn wir das glauben, so sind wir die rechte Kirche. Wie Isaak ein Sohn ist der Verheißung, und hat dieselbe Kindschaft nicht von Abraham und Sarah, denn das ist eine fleischliche Geburt, sondern hat sie von diesem Worte: „Isaak soll der Same sein.“ Gleich auf dieselbige Weise verhält es sich mit Jakob auch. Daß er ein Theil oder Gliedmaß der Kirche ist, hat er daher nicht, daß er vom Isaac geboren ist, sondern daß die Stimme des Berufers sagt: Jakob ist der Kleinste und soll doch gleichwohl der Herr sein. Also ist Abraham zu Babel der Allerweiseste und Schriftlichste gewesen, hat ein eheliches Weib gehabt, mit derselben hat er keusch und züchtig gelebt; aber er hätte mit dem allen müssen verloren und verdammt sein, wo ihn Gott nicht berufen hätte aus Ur in Chaldäa.

88. Also soll man nun die rechte Kirche von der falschen unterscheiden. Denn die falsche Kirche ist vermessen, entweder des Geblüts halben, das ist, von wegen der ersten Ursache der Geburt, oder aber der Dinge halben, so daraus folgen. Ismael argumentirt oder beweist sein Ding mit der anfänglichen Ursache und sagt also: Ich bin von Abraham geboren und bin dazu auch der erstgeborne Sohn: ich werde Gottes Volk sein. Da ist zwar die anfängliche Ursache oder der Ursprung und das Geblüt, aber wo er nicht auch den Beruf haben wird, so ist seine Vermessenheit des Geblütes und der Geburt halben vergeblich und umsonst. Also rühmen alle Papisten und Türken, daß sie Gottes Volk seien; und wiewohl sie sich vom Geblüte der Väter nicht rühmen können als vom Ursprung und der anfänglichen Ursache, rühmen sie sich doch der Dinge halben, so daraus folgen, das ist, sie rühmen sich von wegen ihrer Gewalt, Gerechtigkeit, Religion oder äußerlichen Gottesdiensts, welches alles vom Fleische ist. Wenn du sie aber fragst: Ob sie auch über diese Stücke; davon sie sich rühmen, einen andern Beruf haben, daß sie Gottes Volk seien? darauf können sie nicht antworten, sondern rühmen nur von diesen großen und göttlichen Wohlthaten, welches doch ein vergeblicher unnützer Ruhm ist, und werden damit zugleich die Frommen und Bösen überschüttet.

89. Wenn du den Papist fragst: Warum er mit seinem Hause Gottes Volk sei? so gibt er diese Antwort und sagt: Dieweil ich sitze auf dem Stuhl der Apostel Petri und Pauli und ihr Stuhlerbe bin. Darnach habe ich dessen auch Ursache in der Schrift, als nämlich, aus dem Spruche Christi Matth. 16, 18.: „Du bist Petrus, und auf diesen Fels will ich bauen meine Gemeine“ etc. Nun könnte aber auch wohl ein Hund oder Sau an Petri Stätte sitzen. Aber wo man außer solcher Nachfolge den Beruf, das ist, das Wort, hat, dem du glauben mögest, dasselbe macht die rechte Kirche und Kinder Gottes. Das hat aber der Papist nicht, sondern verfolgt es feindseliger Weise; darum ist er auch kein Kind, sondern ein Feind Gottes und ein Widersacher Christi. Wir aber zweifeln daran gar nicht, daß wir die rechte Kirche sind; denn wir haben das Evangelium, die Taufe, die Schlüssel, die heilige Schrift, die da lehrt, daß der Mensch in der Erbsünde verloren und verdammt ist, und daß vonnöthen sei, daß er durch Christum neu geboren müsse werden.

90. Derhalben ist dies eine nöthige Lehre, so den Unterschied zwischen der rechten und falschen Kirche anzeigt, und ist erstlich von Gott den Juden vorgehalten, auf daß er damit ihre Hoffart, da sie sich des Geblüts halben rühmen, zu Schanden mache, nämlich, mit diesen trefflichen Exempeln in beiden Häusern Abrahams und Isaaks. Denn beide Ismael und Isaak sind von einerlei Geblüt Abrahams geboren, und ist doch gleichwohl Ismael der Erbe nicht; denn er will durch die erste Geburt, und zwar durch dieselbe allein, selig werden. Isaak aber ist der rechte Erbe, darum daß er das Wort und die Verheißung hat. Also sind Esau und Jakob nicht allein von Einem Vater, sondern auch von einer Mutter geboren; und dasselbe ist noch ein stärker Argument oder Beweis denn das vorige. Denn da ist einerlei Vater und Mutter, auch einerlei Geblüt ihrer beiden.

91. Warum wird aber Jakob dem Esau vorgezogen? Antwort: Der Esau ist zu Jakob gekommen; Esau aber ist vermessen und denkt, es müsse das Erbe des Reichs auf ihn kommen von wegen der ersten Geburt. Solcher Vermessenheit ist Gott feind, und will haben, daß die Erneuerung der Natur soll dazu kommen; denn außerhalb derselben hat man sich keiner Gerechtigkeit des Reiches oder Erbes zu trösten. Und aus solchen Orten der Schrift können wir auch unsere Juden gewaltig und kräftig überbinden, und ihren Ruhm vom Geblüt, davon sie sich noch heutiges Tages träumen lassen, widerlegen. Denn es werden aus einerlei Geblüt Zwei geboren, und ist der Eine der Erbe, der Andere aber nicht, und werden unterscheiden Einer vom Andern durch den Beruf. So viel den fleischlichen Ruhm belanget, sind sie allenthalben gleich; aber der Eine hat das Wort und gehört demselben, der Andere aber verachtet es.

\quellseite{24}

92. Deß haben wir auch gar ein gleich Exempel an der päbstischen und unserer Kirche. Denn wir haben beyderseits einerlei Taufe, einerlei Sacrament des Abendmahls, einerlei Schlüssel, auch einerlei Schrift und Wort, und endlich sind wir auch beyderseits von einerlei Aposteln und von einerlei Kirche gleich als von Einer Mutter Rebecca hergekommen. Warum sind wir denn der Sachen mit einander uneins? Antwort: Wir nehmen das Wort an im Sacrament, folgen dem Beruf, handeln dasselbe und gehen damit um nach dem Glauben; sie aber nach dem äußerlichen Ansehen und Schein. Wir sagen: Wo man das heilige Sacrament gebraucht, soll man auf das Wort sehen und soll es mit rechtem Glauben empfangen; sie machen nur ein äußerlich Werk daraus. Darum ist das die rechte Kirche, die am Wort und Glauben hängt, verläßt sich auf seine Werke, sondern hört und folgt Gott, der sie aus Gnaden beruft: der Beruf aber ist das Wort, dadurch die Natur, so vom Teufel und der Sünde verderbet ist, neu geboren wird.

93. Dies ist keine geringe Uneinigkeit und Streit. Es zürnen auch die Widersacher nicht unbillig so feindlich mit uns; denn wir streiten über die höchsten Dinge in der ganzen weiten Welt, nämlich, nicht über das Regiment oder weltliche Herrschaft, nicht über Geld und Gut, über große Herrlichkeit und Gewalt; sondern der Streit ist über ewige Verdammniß und ewiges Leben. Darum ist auch die Feindschaft so bitter und heftig groß, und ist sonst nimmer bitterer oder größer, denn wo man über die Religion uneinig ist. Denn die Sache an sich selbst ist auch sehr groß und schwer, nämlich, ob wir Gottes Kinder, oder des Teufels Kinder seien. Sie wollen traun des Teufels Kinder nicht sein, sondern reißen mit Gewalt an sich die Ehre und den Namen der rechten wahren Kirche, und darum verdammen und tödten sie uns auch. Dagegen verdammen wir sie auch aus dem Worte und sind aus gewissen Zeugnissen der Schrift unserer Sicherheit gewiß, und halten fest über dem Unterschied, welchen Gott gemacht hat zwischen Abraham und den Babyloniern, Ismael und Isaak, Jakob und Esau, nämlich dem Berufe, dadurch zu Abraham ist gesagt worden 1 Mos. 12, 1.: „Gehe aus deinem Vaterland“ etc.; item Cap. 21, 12.: „In Isaak soll dir der Same genannt werden“; item Cap. 25, 23.: „Der Größere wird dem Kleinern dienen.“

94. Wenn man dies nicht hinnimmt, ist die erste Geburt nicht allein dem Abraham, oder Isaak und Jakob, sondern auch allen Menschen ganz und gar unnütze; ja, sie ist ein Brunnen und Ursprung alles Unglücks, welches man an dem Exempel der zwölf Patriarchen, so von Jakob geboren sind, sehen und lernen kann. Denn dieselbigen werden Todschläger, ja Vater- und Brudermörder, Verräther und Mörder ihres eigenen Vaters, dieweil sie den frommen alten Mann also betrüben, daß er begehrt zu sterben, 1 Mos. 37, 18. 35. Und möchte solche ihre Sünde, wo es Einer wollte aufmützen und groß machen, recht und billig ein Vatermord genannt werden. Das, meine ich, sind die rechten Früchte und Werke des Fleisches und Blutes, welche anzeigen, daß der Mensch von seiner Art und Natur ganz und gar Kainisch und mörderisch ist. Und an den Kindern Jakobs hätte zwar auch wohl ein Exempel mögen vorgestellt werden; aber Gott hat sie nicht wollen verwerfen und läßt es bei diesen zwei Zeugnissen bleiben, als Ismaels und Isaaks, Esaus und Jakobs. Und ich halte, daß sie selig geworden sind; wie ich von Ismael und Esau auch die Gedanken habe, daß sie werden Buße gethan haben und zum Beruf getreten sein: und hat doch Gott haben wollen, daß dies Argument und Wahrzeichen da stehen sollte, nämlich, daß wir daran lernen sollten, daß die erste Geburt vor Gott nichts sei, und das Fleisch gar verdammt ist mit allem, was es hat und vermag, es sei gleich wie groß und trefflich es immer sein kann; wie denn das Regiment und weltliche Herrschaft ist, dazu Gesetze und Gerechtigkeit. Wie des Kaisers Augustus Reich oder Regiment sehr schön und herrlich gewesen ist, da Alles in großem Flor gestanden, Weisheit, Gerechtigkeit, Friede, Geld und Gut, und dazu viel herrliche Siege. Aber dies alles gehört zu der ersten Geburt oder Vernunft, zur Schöpfung und zum Fleisch und ist derhalben auch gar verdammt.

95. Darum muß man drauf sehen, ob am Menschen etwas über solche natürliche Gaben sei, mögen nun die Heiden, oder der Türke, oder Papst dieselben besitzen. Wo der Beruf nicht da ist, sollt du es dafür halten, daß einfach Alles verloren und verdammt sei. Denn also hat Gott selbst geurtheilt und hat es vornehmlich diesem Volk an ihren vornehmsten Vätern gezeigt, nämlich, daß dieselbigen außerhalb und über die erste Geburt auch haben müssen berufen werden. Und was die erste Geburt vermocht hat ohne den Beruf, dasselbe zeigen an die greulichen Sünden und Laster der allerhöchsten Patriarchen, welche sich mit Vater- und Brudermord, und dazu mit Blutschande verunreinigt haben. Und dasselbe bezeugen auch die Historien der Könige und Richter, und endlich die große, greuliche Sünde, da sie den Sohn Gottes, ihren Messias, an das Creuz geschlagen haben. Diesen Unterschied der rechten und falschen Kirche lehrt das treffliche und wunderbare Exempel von der Empfängniß der zwei Kinder Jakob und Esau, so sich in Mutterleibe mit einander gestoßen haben.

96. Nun wollen wir weiter auch etwas von den Lügen der Juden daher setzen um derer willen, so dermaleins auch die Auslegungen der Rabbinen lesen werden. Denn dieselben erregen an diesem Ort eine Frage, nämlich: Warum Isaak nicht habe ein ander Weib genommen, nachdem er schier zwanzig Jahre die Unfruchtbarkeit seines Weibes erfahren hatte? Da sie aus dem Exempel Sarahs schließen, daß man zehn Jahre dazu geben sollte, darin man sich erkundigen und erfahren möchte, ob ein Weib unfruchtbar oder fruchtbar wäre. Sie antworten aber auf diese Frage also: Daß es mit Isaak eine andere Gestalt habe als mit Abraham; denn Isaak sei Gott geopfert gewesen, und sei aus Gottes Befehl ein Opfer und Brandopfer geworden; deshalb habe ihn nicht geziemen wollen, ein ander Weib zu nehmen. Wie Paulus 1 Tim. 3, 2. haben will, daß ein Bischof soll Eines Weibes Mann sein.

97. Wie ungereimt und loses, unnützes Geschwätze dies sei, sieht zwar jedermann wohl; und man muß doch solches bisweilen anführen, diejenigen zu ermahnen, so in der hebräischen Sprache studiren, daß sie die Worte und Schriften der Juden mit Verstand lesen und daß sie davon recht urtheilen lernen. Wir erkennen zwar, daß es eine große Wohlthat sei, daß wir diese Sprache von ihnen empfangen haben; aber vor dem Unrath der Rabbinen muß man sich hüten, welche aus der heiligen Schrift gleichsam ein heimlich Gemach haben gemacht, darein sie ihre Schande und närrische Meinungen getragen. Diese Erinnerung thue ich darum, dieweil heutiges Tages unserer Theologen viele sind, die allzuviel von den Rabbinen halten, was den Verstand der Schrift belanget.

98. In der Grammatik kann ich sie leichtlich ertragen und Geduld mit ihnen haben; aber es mangelt ihnen am rechten Verstande der Schrift; nach dem Spruche Jes. 29, 14.: „Die Weisheit der Weisen dieses Volkes wird untergehen, und der Verstand seiner Klugen wird verblendet werden.“ Der Spruch zeigt an, daß bei den Juden kein rechter Verstand der Schrift sein wird, sondern das Buch der heiligen Schrift ist ihnen verschlossen und versiegelt. Gott redet zu diesem Volk mit einer andern Zunge, Jes. 28, 11., und sie können oder wissen Anderes nichts, denn eitel Lästerung auszuspielen auf die christliche Religion.

99. Darum reimet sich dies nicht auf einander, daß sie sagen: Isaak ist geopfert gewesen, darum hat er kein ander Weib nehmen mögen. Abraham ist ebenowohl Gott gewidmet gewesen als Isaak, und noch wohl etwas mehr; denn er war berufen aus Ur in Chaldäa, und war nicht einmal, sondern oft geschlachtet worden: und hat ihn doch solche Heiligung daran nicht gehindert, daß er deshalb nicht ein ander Weib hätte mögen nehmen. Also wäre auch Isaak gar frei gewesen, ein anderes Weib zu nehmen, wo er es hätte thun wollen; aber er hat lieber Alles dulden und das Aeußerste leiden wollen, auf daß er sich damit stellte zum Exempel der Geduld und wie man auf Gottes Hülfe harren müsse. Denn es war daselbst große und gewisse Hoffnung, daß Rebekka eine Mutter werden würde, dieweil sie von Gott dazu erwählt, daß sie sein Weib sein sollte, und hatten ihn dieselbe die Engel zugeführt und die Heirath gemacht, indem sein Vater den Knecht dazu bestellt und darnach ausgesandt. Er wird aber darum also hingehalten bis in das zwanzigste Jahr, weil GOTT diesen heiligen Mann und die heilige Matrone hat wollen versuchen, zum Exempel und Unterweisung aller Gemeinden Gottes.

\subsection{II}

\begin{quote}
\textbf{V. 22. Und die Kinder stießen sich mit einander in ihrem Leibe. Da sprach sie: Da mirs also sollte gehen, warum bin ich schwanger worden?}
\end{quote}

\quellseite{25}

100. Diese Historien der heiligen Väter reden von der Kinderzeugung, vom Beischlafen, und wie sich Mann und Weib ehelich zusammen gehalten haben, darüber in der Welt nichts Geringeres und das lächerlicher sein möchte, angesehen oder gehalten wird. Und werden uns doch in solchen geringen Dingen so treffliche hohe Exempel des Glaubens und andere christliche Tugenden vorgehalten. Dies gute fromme Weib Rebekka, welche Fleisch und Blut gewesen ist nicht anders, denn wie wir auch sind, ist ganz schwerlich angefochten worden so viel Jahre lang. Denn erstlich, war sie gewißlich von GOtt dazu erwählt, daß sie sollte eine Mutter sein. Sie war von den Engeln gesucht, gefunden und zu ihrem Bräutigam geführt worden darum, daß sie eine Mutter werden sollte. Aber daraus wird gar nichts, und sie wird also von GOtt mit einem langen Creuz und Leiden angefochten. Da ist wahrlich nach der Lust und Freude, so sie zuvor gehabt, Traurigkeit, Bitterkeit, Weinen und Jammer gefolgt. Dies, meine ich, sei ja die rechte Freude des Ehestandes. Darum bittet sie und begehrt heftig und brünstig, daß sie doch die Mutter des Samens werden möge. Darnach, da sie nun erhört ist, und GOtt gleichsam mit Gewalt gedrungen worden, die unfruchtbare Rebekka fruchtbar zu machen, kommt bald eine andere Beschwerung darauf. Es war ihr zuvor schwer, daß sie die Hoffnung verlieren sollte, welche sie aus der Erwählung und göttlichem Beruf gehabt, und daß sie unter den Nachbarinnen und ihrem Hausgesinde sollte geschmäht und verachtet werden.

101. Nun hat sie eine Frucht erlangt und ist schwanger geworden, und ist über und wider die Natur die Unfruchtbarkeit überwunden, daß sie durch ein neu Wunderwerk schwanger geworden. Aber da ist nun etwas vorhanden, das sie leiden muß, welches noch viel beschwerlicher ist. Denn da sie nun die Frucht im Leibe hat, und bei sich selbst fröhlich darüber ist, daß sie nach solchem großen Zagen mit so reicher Hoffnung getröstet und aufgerichtet wird, da sie eine Mutter wird zweier Kinder, während sie zuvor verworfen und verstoßen war gewesen und nun bald einem andern Weibe hätte weichen und Raum geben sollen, und die Ehre, daß sie möchte Mutter werden, übergeben (wie droben 1 Mos. 16, 2. Sarah auch gethan): da wird sie nun erst in die allerbetrübteste und äußerste Verzweiflung geworfen. Denn die Kinder stoßen sich mit einander in ihrem Leibe, und sie kann doch nicht eigentlich wissen, ob es Kinder sind oder nicht. Darum sagt sie: „Da mirs also sollte gehen, warum bin ich schwanger worden?“ Sie redet nach dem Urtheil in abgebrochenen Worten, als die erschrocken gewesen, und sich hat gereuen lassen, daß sie zuvor gewünscht und begehrt hat schwanger zu werden. Ach, denkt sie, wollte GOtt, ich hätte nie um Kinder gebeten! Warum habe ich nicht Geduld getragen und einem andern Weibe das Mutterrecht übergeben, sintemal ich nun viel größeren Jammer erfahren und leiden muß denn zuvor!

102. Dies ist eine zwiespältige, ja, wohl eine siebenfältige Anfechtung. Denn jetzt meint sie, es wäre besser gewesen, daß sie für und für unfruchtbar geblieben wäre, denn daß sie so jämmerlich soll geplagt werden, und soll auch noch zugleich das in Zweifel stellen, ob sie eine Mutter werden soll oder nicht; denn sie hat sich nun dessen zum andernmal begeben müssen, darum daß sie jetzt den Tod fühlt, und anfängt, an ihrem Leben zu verzagen der großen Gefahr halben, so sie empfunden hat. Denn sie weiß nicht, was es für ein Unglück sein möge, was für eine Bewegung in ihrem Leibe sei: sie fühlt nur das Stoßen, und kann keine andere Schuld finden, denn daß sie dieselbe nur auf sich selbst legen und sich eine Märtyrin schelten muß, daß sie GOtt so heftig um Kinder gebeten und geheischt hat. Sie sieht, daß sie nun mit größerer Schande wird sterben müssen, denn wo sie wäre unfruchtbar geblieben. Also wird abermals die Hoffnung, daß sie eine Mutter werden sollte, gar erstickt, ja, es wird die Mutter mit der Frucht getödtet. Denn die Vernunft kann keinen andern Verstand haben oder nicht anders denken, so oft GOtt unsere Hoffnung und Gebet schwach werden läßt, denn daß sie sage: Je länger ich bete, je härter ich betrübt und geplagt werde.

103. Soll man aber denn nun darum nicht beten, GOtt nicht vertrauen und auf ihn nicht hoffen? wie jetzt die Epicurer sagen, sie wollen es alles dem Willen und Schickung GOttes befehlen; denn was GOtt beschlossen habe, das müsse nothwendig also geschehen. Nein, man soll das Gebet mit nichten unterlassen; denn GOtt hat uns heißen beten, glauben, ihn anrufen, seinen Namen heiligen und auf ihn hoffen. Er könnte zwar wohl sich eine Kirche sammeln ohne das Wort oder Predigtamt, könnte das politische oder weltliche Regiment wohl verwalten ohne Obrigkeit, Kinder zeugen ohne Eltern und Fische schaffen ohne Wasser: aber er gebietet uns und will haben, daß wir predigen und beten sollen, und daß ein Jeder an seinem bestimmten Orte und in seinem Berufe thun soll, was ihm gebührt. Ist es aber der Fall, daß es nicht allezeit so hinaus geht und gerathen will, wie wir es zuvor bedacht und bestimmt haben, soll doch gleichwohl ein Jeder in seinem Beruf und Amt bleiben, und soll GOtt das Gedeihen, den Ausgang oder Ende lassen befohlen sein.

104. Eine solche Anfechtung ist dies auch gewesen, so Rebekka gehabt. Denn sie wird also gedacht haben: Siehe, du hast GOtt gebeten, daß er dir wolle Kinder geben: es wäre besser gewesen, du hättest darum nicht gebeten. Denn da sie das Stoßen in ihrem Leibe gefühlt, denkt sie nochmals nicht, daß es Kinder seien, sondern denkt, sie trage in ihrem Leibe etwa ein seltsam und scheußlich Wunder. Darum sagt sie: „Da mirs also sollte gehen, warum bin ich schwanger worden?“ Das fromme, heilige Weib zeigt mit diesen Worten an, daß es gereut und sie sich selbst gestraft habe um der vorigen Bitte willen, so sie gethan hatte. Als wollte sie sagen: Da mir es also gehen sollte, wäre es besser gewesen, ich hätte gar nichts gebeten; soll es also zugehen, ach, warum habe ich denn gebetet? warum habe ich begehrt, eine Mutter zu werden wider GOttes Willen? Es sind Worte Einer, die sich ihr voriges Gebet hat gereuen lassen und nun gar verzagt ist.

105. Denn wo die Menschen in großer Verzweiflung stecken, pflegen sie in abgebrochenen Worten zu reden, haben in ihrer Rede mehr Ausrufe als andere Worte. Und solche kurze abgebrochene Reden zeigen an die Größe des Schmerzes, so in den Herzen der Betrübten ist und derer, die schier gar verzagt sind; wie denn Rebekka ein solch betrübt und verzagt Herz gehabt hat. Denn dieser Heiligen vornehmste Sorge und Bekümmerniß, so sie gehabt, ist gewesen von der Verheißung des Samens. Darum, sage ich, verzagt Rebekka gar, und wird gedacht haben: Ich werde in dieser Gemeinde keine Mutter werden, sondern es wird dazu eine Andere erwählt werden nach meinem Tode, die dessen mehr werth ist denn ich: und begibt sich also desselben gar, daß sie würde eine Mutter werden; wie droben (1 Mos. 16, 2.) Sarah gethan, und hernach (Cap. 30, 3.) Rahel auch thut, welche ja wahrhaftig Jakobs Weib war und muß doch gleichwohl einer Andern das Mutterrecht übergeben; denn sie sagte zu Jakob: Gehe hin zu der Magd etc.

106. Aber es sind wenig Leute, die dies in den Historien der Väter betrachten. Wie die Papisten davon urtheilen und denken, die Väter haben nur ihre Lust geübt und Unzucht getrieben, wie sie selbst thun, die ohne Anfechtung, und ohne Mühe und Arbeit dahin leben. Diese heiligen Menschen aber sind immer mit Anfechtung und Betrübniß beladen gewesen, welches ihnen die bösen Gedanken und böse Lust leiblich hat vertreiben können. Sie haben ein groß Gesinde zu regieren gehabt, und unterdeß dabei große Anfechtung leiden müssen, gleich als ob sie aller göttlichen Hülfe beraubet und von GOtt gar verlassen wären.

107. Uns aber wird dies zum Exempel und zur Lehre vorgehalten, erfrischt, daß die fleischliche Geburt, wie scheußlich und verderbt sie auch ist, doch nicht der Natur eigen, sondern eine Gabe GOttes ist. Das soll man an den Vätern lernen. GOtt theilt zwar diese Gabe aus unter die Menge auch der bößesten Menschen, aber die Gottseligen verstehen es recht, daß es eigentlich eine Gabe GOttes sei. Wie der 128. Psalm V. 3. 4. sagt: „Deine Kinder werden sein wie die Oelzweige um deinen Tisch her. Siehe, also wird gesegnet der Mann, der den HErrn fürchtet.“ Darum soll jeder gottseelige Mensch lernen, daß es eine sehr große Gabe sei und eine göttliche Schöpfung, wo man Kinder, Söhne oder Töchter zeuget und gebieret. Denn daher heißen wir GOtt unsern Schöpfer und unsern Vater, und sollen uns die fleischlichen Sünden und Strafe, damit die Geburt beladen ist, nicht anfechten oder beirren lassen, daß wir es nicht dafür halten sollten, daß es ein Segen GOttes sei, und ein solch Werk, das GOtt angenehm und gefällig sei.

108. Die Gottlosen verstehen es nicht, welche auch die andern Creaturen wie den Schöpfer selbst vergeblich führen und unnützlich brauchen, so daß nicht allein die Creaturen, sondern auch der Schöpfer der Eitelkeit unterworfen ist. Man muß aber die scheußliche Mißgestalt unterscheiden und absondern von der Creatur, von der GOtt haben will, daß sie soll herrlich sein und groß geachtet werden; wie solches die heiligen Väter gelehrt haben aus diesen Anfechtungen, nämlich, daß die Geburt oder Kinderzeugung ein göttlich Werk und allein GOttes eigen Werk sei. Und ist gewißlich wahr, wo es nicht mit so scheußlicher Unzucht und Schande vergiftet wäre, so würden wir in Staunen geraten über die Würdigkeit und die Wunderbarkeit solches großen Werkes. Denn was ist wunderbarer, denn daß ein Mensch vom andern Menschen geboren wird? Es ist aber verwundelt, unterdrückt mit der scheußlichen Sünde der bösen Lust. Darum muß es mit solchen Exempeln gerühmt und herausgestrichen werden, auf daß die Leute die Größe und Hoheit dieses Werkes und des göttlichen Segens erkennen lernen.

\quellseite{26}

109. Zum Andern, halten uns diese Exempel vor den Glauben, die Hoffnung und die Liebe, das ist, die Weise, dadurch GOtt mit uns pfleget zu handeln. Denn also handelt er auch mit den heiligen Patriarchen; wie St. Paulus Röm. 11, 33. sagt: „Seine Gerichte sind unbegreiflich, und seine Wege unerforschlich“, nicht allein in seinen Werken, sondern auch in seinen Worten und Verheißungen. Und darum läßt es sich auch in der Welt ansehen, daß nichts Ungewisseres sei, denn GOttes Wort und der Glaube, und nichts, das so eitel und vergeblich sei als die Hoffnung auf GOttes Verheißung, und endlich ist nichts, das da scheinet mehr nichts zu sein als GOtt selbst. Derhalben ist dies die rechte Wissenschaft der Heiligen und Gläubigen, und das Geheimniß, welches vor den Weisen verborgen und den Unmündigen offenbaret ist.

110. Der Pabst und Türke können leichtlich glauben, dieweil bei ihnen so groß Glück ist und darum daß sie in so großem Flor sind mit Gewalt, Weisheit, mit großem Schein der Heiligkeit, Religion und äußerlichem Gottesdienst, daß nichts darüber ist. Der Türke hält es so gar für gewiß, daß er mit seinem Haufen GOttes Volk sei, daß er auch dabei sein Leib und Leben würde dürfen zum Pfande setzen. Denn er hat sein Ding und seinen Gott in Händen; GOtt zieret ihn mit sehr schönen und hohen Gaben. Und mit den Christen geht er gleichsam widerinnig um, gibt ihnen das Gegenteil, wirft sie unter das Kreuz, und gibt sie dahin, daß sie von den Widersachern gewürgt werden wie die Schlachtschafe. Darum trotzt der Türke wider uns mit großem Stolz und Hoffahrt. Wie in einer Schlacht oder Niederlage der Christen soll geschehen sein, da die Türken gesehen, daß hin und wieder die toten Körper gelegen und ihnen die Nasen abgehauen gewesen, da haben sie solchen Jammer derer, so erschlagen waren, beklagt und bitter und giftig Gespött damit getrieben, und diese Worte oftmals wiederholt und gesagt: Jesus Maria! Als wollten sie sagen: Wo ist nun ihr GOtt?

111. Daher kommen solche betrübte Klagen in den Psalmen: Ps. 7, 7.: „Stehe auf, Herr“, schlafe nicht; Ps. 10, 1.: „Warum trittst du so ferne“ etc., und dergleichen Worte mehr der verzagten Herzen, so da zweifeln, ob wir auch einen GOtt haben, und die da fragen, Ps. 42, 11.: „Wo ist nun dein GOtt?“ Soll man an den glauben, der sich also verbirgt und sich von uns wendet? Aber der Türke weiß nicht, was er selbst thut oder was mit uns werden soll. Er denkt und urtheilt also: Die Christen sind von GOtt gar verworfen und verstoßen; und ist vermessen, als ob ihm GOtt günstig sei und ihn erwählet habe; ist derhalben noch mehr verhärtet und verstöcket denn der Teufel selbst. Aber wenn wir ihm in seine Hände werden gegeben und darnieder liegen, alsdann wird der Himmel voll Märtyrer, und die Gottlosen werden versucht, daß sie seufzen und schreien, daß nur der jüngste Tag kommen wolle. Er aber richtet mit seinem Stolz und Uebermuth an, daß über ihn die ewige Verdammniß so viel desto eher kommen wird.

112. Also hat GOtt auch nicht anders gehandelt mit dem Volk Israel zur Zeit der Assyrer und Babylonier. Da war kein Glück oder Sieg und ließ sich auch da kein GOtt sehen, und da war er recht ein verborgener GOtt, der GOtt Israel, der Heiland, wie im Propheten Jesaia steht Cap. 45, 15. Wie es heutiges Tages mit uns auch geht, da Himmel und Erde, der Pabst und Türke, und Alles wider uns ist, und GOtt, an den wir glauben, ist gar nichts; der Fürst und Gott dieser Welt, der ist Gott. Darum fallen auch viel Christen von uns ab zu den Türken, nämlich die, so nicht recht unterrichtet sind und die das Wort nicht haben. Sie sehen, daß dort Alles in Flor und überaus wohl steht, und daraus nehmen sie ab, es müsse GOtt gegenwärtig bei den Türken sein; sintemal sie GOtt also lebendig macht, uns aber tödtet und übergibt er in die Hände des grausamen Feindes.

113. Wir aber haben von GOtt nichts Anderes, denn allein das reine Wort, daß unser Herr Jesus zur Rechten des Vaters sitzt und ein Richter ist der Lebendigen und der Todten, durch welchen wir Könige und Priester sind. Aber wobei kann man dies sehen? Wir sehen es freilich nicht gegenwärtig vor Augen, sondern müssen wünschen, hoffen, harren und dem glauben, was GOtt in seinem Wort gesagt hat. Denn warum er sich also verberge, werden wir an jenem Tage sehen, wenn alle Feinde ihm unter seine Füße geworfen werden. Inzwischen müssen wir glauben und hoffen; denn wenn wir es jetzt gegenwärtig vor Augen sehen könnten, so bedürfte man des Glaubens nicht.

114. Wie ungefähr und falsch nun der Glaube, und wie eitel und vergeblich die Hoffnung scheinen mag, so weiß ich doch, daß es dazu kommen soll, daß wir den Türken unter unsere Füße treten werden, und die jetzt begraben liegen, deren Blut er vergossen hat, die werden ihn zertreten und in den Abgrund der Hölle stoßen. Dasselbe werden auch die anderen Märtyrer alle thun, so vom Pabst, Kaiser, Franzosen und von Andern verbrannt und getödtet sind. Denn dies ist die Kunst und Weisheit der Heiligen, nämlich, daß sie wider die Lüge an die Wahrheit glauben, und wider die öffentliche Wahrheit an die verborgene Wahrheit; ja, sie glauben auf Hoffnung, da nichts zu hoffen ist.

115. Dieselbige Kunst wird uns in diesem Exempel der Rebekka vorgehalten. Denn sie ist erhört worden, und ist erwählt, daß sie sollte eine Mutter werden; sie ist aber noch nichts weniger denn eine Mutter, auch da sie jetzt schon erlangt hat, daß sie ist fruchtbar geworden, und da sie große Hoffnung bekommen hat, daß sie soll eine Mutter werden. Also fällt bald alle Hoffnung und Warten dahin, und liegt gar darnieder. Ist denn nun das die Erfüllung des göttlichen Wortes? Antwort: Man muß einfach an der Kunst oder Erkenntniß, so die Heiligen haben, hängen bleiben, und auf die wunderbare Regierung sehen, wie sie von GOtt geführt werden. Denn also ist es mit Christo selbst auch zugegangen. Da der in die Herrlichkeit des Vaters hat gehen wollen, ist er gestorben und niedergfahren zur Hölle; daselbst verschwindet alle Ehre und Herrlichkeit. Er sagt zu seinen Jüngern Joh. 14, 12.: „Ich gehe zum Vater“, und geht doch zum Grabe.

116. Solche Exempel sollen wir immer in frischem Gedächtniß und vor Augen haben, die heiligen fleißig bedenken und erwägen. Denn wo wir vom Türken erwürget werden mit großer Schmach und Lästerung; wenn die kleinen unmündigen Kindlein, dazu alte Leute und junge Mägdlein geschlachtet und grausamlich getödtet werden, da rufen und schreien wir: Wo ist nun Christus, wo ist unser GOtt? Zudem verspottet sie uns schmählich und sagen: Gehet nun hin und glaubet an euren Christum. Aber da sollen wir ihnen getrost mit Freudigkeit antworten, wenn sie auch dasselbe schon nicht hören, und sollen sagen: Siehe, dies kleine Kindlein, welches du zerhackt hast oder hast es an einen Pfahl gestecket, dasselbe wird im andern und bessern Leben mit so großer Klarheit geziert werden, daß du wünschen wirst, daß dir es möchte so gut werden, daß du nur die Wunden sehen möchtest, so du ihm gemacht hast. Denn der Türke richtet sonst mit seiner Tyrannei nichts aus, denn daß er den Himmel voll Märtyrer macht, dieweil sie alle um Christi willen getödtet werden: und wiederum füllt er die Hölle mit seines und seines Volkes Leibern.

117. Also auch der Pabst (wiewohl derselbe ärger ist denn der Türke, dieweil er gar nichts glaubt) meint von denen, die er tödtet, sie seien nun gar dahin und sei mit ihnen aus, und denkt nicht anders, denn daß ihm frei und erlaubt sei, wider jedermann seines Gefallens zu toben und zu wüthen, darum ihn niemand soll zu strafen haben, und zeigt oder hebt die Monarchen und weltlichen Herren dazu, daß sie dergleichen grausame Tyrannei auch üben sollen. Unterdeß klagen wir, daß wir jämmerlich erwürgt und erdrückt werden, und schreien Ps. 44, 24.: „Herr, warum schläfst du?“ Denn GOtt verbirgt sich eine Zeitlang, und wie im Hohen Liede Salomonis Cap. 2, 9. steht, „er stehet hinter unserer Wand, und siehet durchs Fenster, und gucket durchs Gitter“; darum ist er nicht weit von uns, sondern ist gegenwärtig und sieht Alles an. Und beide Papst und Türke eben damit, daß sie uns tödten, müssen uns dienen, sie wollen oder wollen nicht, auf daß wir die unverwelkliche Krone der Gerechtigkeit empfangen mögen, 1 Petri 5, 4. Wie auch Herodes den kleinen unmündigen Kindern mit Wohlthaten immermehr solchen Dienst hätte thun können, wie er ihnen wohl damit gethan, daß er sie getödtet hat.

\quellseite{27}

118. Also läßt es sich ansehen, daß er uns auch verlassen habe, dieweil er aufgefahren ist gen Himmel und sitzt zur Rechten Gottes des Vaters; und hat uns doch gleichwohl sein Wort und die heiligen Sacramente gelassen. Was ist aber das? wird uns das Fleisch da verwesen und sagen; sehe ich doch nichts, es ist ein bloßes Wort. Ja, es ist zwar ein Wort, aber in demselben Wort ist eine sonderliche und göttliche Kraft, welche du alsdann verstehen und empfinden wirst, wenn du dich daran mit festem Glauben hängen und halten wirst. Wo du schon dein Leben darüber wirst in Gefahr müssen setzen, ist daran wenig gelegen; denn es ist nicht weit dahin, daß du wiederum von den Todten wirst auferstehen und lebendig werden. Seine Wege sind unerforschlich, und seine Gerichte sind verborgen gegen seine Heiligen, Röm. 11, 33., wie das sein eingeborner Sohn bezeuget, welcher das allervollkommenste Exempel ist aller Kinder.

119. Also handelt nun Mose von den wichtigsten Dingen, da er von der Geburt der Menschen und von dem fleischlichen Werk redet. Rebekka verliert die Ehre, daß sie hätte sollen eine Mutter werden, und verliert dazu Gott selbst, der sie erhört und schwanger gemacht hat: derselbe zieht seine Hand ab und verbirgt sein Angesicht, daß sein Werk, daß sie war schwanger geworden, nicht mehr angesehen wird für ein göttlich Werk, sondern wird für ein Gespött und Betrug des Satans gehalten, und sie denkt nicht anders, denn daß sie verworfen und verstoßen sei. Er verzieht aber nicht ewig, bleibt auch nicht allezeit weg von uns. Er regieret zwar alle Dinge also, daß das Wort angesehen wird, als das gar vergeblich und nichts sei. Aber aus demselben, das so gar nichts ist, macht er Alles durch große und wunderbare Veränderung, also daß die Dinge, so zuvor angesehen worden, daß sie Alles sind, verderben und zunichte werden müssen.

Also ist das Verderben des Papstes jetzt schon angefangen und er wird alle Tage je mehr und mehr durch den Geist des Mundes Christi zertreten werden, und wird in kurzer Zeit dazu kommen, daß auch der Türke seine verdiente Strafe bekommen wird.

\subsection{III}

\begin{quote}
\textbf{V. 22. Und sie (Rebekka) ging hin, den Herrn zu fragen.}
\end{quote}

120. Die arme Rebekka versteht nun und hält es dafür, daß Alles verdorben und verloren mit ihr sei. Was soll sie aber nun thun? Sie geht hin, den Herrn zu fragen. Bis daher ist sie angesehen worden, als wäre sie vom Angesicht des Herrn verstoßen und verworfen, und ihre Hoffnung und Glaube sind jetzt in die höchste Noth gekommen, daß es damit Mühe und Arbeit gehabt. Aber dies ist eben der rechte dies criticus, der kritische, entscheidende Tag, da sie wiederum lebendig wird, und hingeht, den Herrn zu fragen, da es schien, als ob sie in der äußersten Noth und Gefahr sei. Denn da pflegt die Saite leichtlich zu brechen, wo sie am härtesten gespannt oder gezogen wird.

121. Die Juden disputiren darüber: Wohin Rebekka gegangen sei, den Herrn zu fragen? Die Rabbinen sind allzusehr mönchisch, da sie sagen, sie sei entwichen an einen Ort besonders allein. Darum wollen wir die Meinung behalten derer, die da sagen, daß sie hingegangen sei zu dem heiligen Patriarchen Sem, welcher noch gelebet hat und mit ihm auch Salah und Eber. Dieselbigen Väter haben die Kirche und Gemeinde Gottes regiert, und haben hin und wieder das Wort und den rechten Gottesdienst erhalten. Denn Sem war ein sehr heiliger Mann, welcher die Kirche regiert und so viel hundert Jahre nach der Sintfluth über dem Wort und Gottesdienst gehalten hat; darum hat man seine Rede und Predigten für göttliche Antwort gehalten.

122. Es hat aber Rebekka nicht allein gefragt nach der Frucht oder Kindern, damit sie schwanger ist gegangen, sondern auch nach ihrer Seligkeit, und ob das, so ihr begegnet, ein Zeichen wäre des Zorns oder aber der Gnade Gottes? Denn ein jeglicher Unfall oder Schaden, auch der leibliche, bringt mit sich Verzweiflung, oder daß man anfängt zu zweifeln an dem Willen Gottes, so er zu uns trägt. Und derhalben muß man in aller Anfechtung das Herz verwahren und stärken, daß es nicht denke, wenn wir betrübt und gestraft werden, daß solches geschehe aus Gottes Zorn, daß er uns ewig verdammen wollte. Als, wenn der Türke und Papst so greulich wider uns toben, so ärgsten wir uns billig und erschrecken davor alle mit einander. Wo du aber das Wort hast und die heiligen Sacramente, so thue die Augen zu und verachte nur getrost ihre grausame Tyrannei, also daß du auch zu ihnen sagest: Tobet und wüthet, und seid nur unsinnig, so viel als euch selbst gelüstet, und zerreißet mich auch in einzelne Stücke: ich weiß doch, daß Christus, mein Erlöser, noch lebt und sitzt zur Rechten des Vaters; darum werde ich auch leben und werde euch wiederum mit Füßen treten.

123. Aber dasselbe Vertrauen und solche geistliche Hoffahrt kommt her nicht vom gegenwärtigen, sondern vom zukünftigen Dinge. Die Vernunft folgt nur dem nach, was sichtbar ist; man muß sie aber allhier tödten, auf daß das Wort und der Glaube Raum haben. Es kann aber die Vernunft nicht getödtet werden denn allein durch Verzweiflung, Mißglauben, Haß und Murren wider Gott, auf daß endlich, wenn alles, was äußerlich vorfallen kann, hintan gesetzt ist, das Herz allein am Worte und den Sacramenten hange, sich daran halte und damit zufrieden sei. Denn Gott ist unbegreiflich, und wir halten ihn für nichts bei allen seinen Thaten und Wunderwerken, darum daß er uns, unserm Gefühl nach, im Leiden und in Trübsal verläßt, und gibt uns nur das Wort, damit er uns, wie ein Fisch mit der Angel gezogen wird, also durch dies Meer der Gefahr und Anfechtungen an das Ufer oder Land bringt.

\begin{quote}
\textbf{V. 23. Und der Herr sprach zu ihr: Zwei Völker sind in deinem Leibe, und zweierlei Leute werden sich scheiden aus deinem Leibe; und ein Volk wird dem andern überlegen sein, und der Größere wird dem Kleinern dienen.}
\end{quote}

124. Bisher ist gesagt worden, wie sich Gott verbirgt, ja, gar verloren hat vor dem heiligen Weiblein, der Rebekka, ja, auch vor Isaac selbst und Abraham, dazu vor allen, so in derselben Gemeinde gewesen sind: also daß sie hat laufen und Zuflucht haben müssen zu der Gemeinde und zum Patriarchen Sem, und daselbst Hülfe und Rath suchen. Nun aber beschreibt Mose, wie dieser verlorne Gott wiederum gefunden werde und sich sehen lasse: uns zum Exempel, auf daß wir in unsern Anfechtungen nicht verzagen, wenn sich das Gebet an die gewisse Verheißung hält; sondern sollen gewiß sein, daß Gott gemächlich müsse erscheinen, sich sehen lassen und Antwort geben, wenn auch schon der Himmel und alle Creaturen den Fall und Untergang drohen sollten.

125. Wiewohl aber nun dies Harren fast beschwerlich ist, sollen wir doch vom Heiligen Geist also unterrichtet und gelehret sein, daß wir wissen, daß Gott unser Gebet gewißlich erhören muß, dieweil es unmöglich ist, daß er lügen muß. Er hatte aber Isaac den Segen verheißen. Aber bei solcher Verheißung sind viel Beschwernungen und Hindernisse vorgefallen, daß ein schwach Herz wohl zweifeln möchte, ob auch die Verheißung sollte wahr sein. Und wiewohl Rebekka seine sonderliche Verheißung gehabt, hat sie doch gleichwohl gebetet. Und wo nun das Gebet recht und wahrhaftig ist, wie denn ohne Zweifel der Rebekka Gebet gewesen ist, so muß es sich doch von Noth wegen an die göttliche Verheißung halten.

126. Denn wer wollte von selbst so kühn sein, daß er seine Hände und Augen dürfte zu Gott aufheben, wo er nicht des Willens Gottes gewiß wäre? Derhalben hält sich alles Gebet an das Gebot und Verheißung Gottes. Denn im andern Gebot steht Gottes Wort, welches uns lehrt beten, nämlich, da Gott 2 Mos. 20, 7. sagt: „Du sollst den Namen des Herrn, deines Gottes, nicht mißbrauchen“; item Ps. 50, 15.: „Rufe mich an in der Noth“ u. s. w. Wenn aber nun das Gebet geschehen ist, sollst du wissen, daß du an der Erhörung eben so wenig zweifeln mußt, als du an der Verheißung nicht zweifeln kannst oder sollst. Denn das Fundament oder Grund des Gebets ist die Verheißung, nicht unser Wille, nicht unsere Würdigkeit oder unser Verdienst. Und derselbigen Lehre haben wir ein trefflich Exempel an Gebet der Rebekka, welche nach so vielen Kinderwünschen zuletzt wunderbarlich erhört wird. Denn im Anfang, da sie schwanger geworden, zweifelt sie daran nicht, sie sei nun erhört und ihr sei von Gott geholfen worden: aber dieser Glaube und Zuversicht wird bald gehindert von wegen einer neuen Gefahr, deren sie sich mit nichten versehen hatte. Denn die Kinder stoßen sich miteinander in ihrem Leibe, also daß sie wiederum in Verzweiflung gebracht wird. Da sie nun auf das höchste erschrocken ist, richtet sie Gott mit dem Wort wiederum auf, und zeigt an, daß sie wahrhaftig erhört sei. Auf daß er damit bezeuge, daß er diejenigen liebe, so im Glauben beten, und im Gebet, das ist, wo das Gebet lange unerhört bleibt, auf ihn harren und seiner Hülfe warten.

\quellseite{28}

127. Solche Exempel erkennen oder verstehen die Papisten nicht; denn ihr Gebet ist Anderes nichts, denn daß sie damit Gott spotten und verlachen. Ein ernstlich und brünstig Gebet aber, welches nicht aufhört und nicht müde wird, sondern wartet und harret bis auf den letzten Augenblick, das bricht endlich durch Himmel und Erde, und ist unmöglich, daß es nicht sollte erhört werden. Denn da ist es Gott das angenehmste Opfer, wenn wir also beten, daß das Gebet unsern Sinn und Verstand übertrifft. Wie St. Paulus sagt Eph. 3, 20.: „Gott ist mächtig und kann überschwänglich thun über alles, das wir bitten oder verstehen“ etc. Wenn die Sache gar verloren und aller Rath und Fleiß vergeblich ist, da sei stark und hüte dich, daß du von Gott nicht abfällest; denn Gott ruft allen Dinge aus den Todten und aus Nichts; wo keine Hülfe oder Hoffnung mehr vorhanden ist, da fähet erst die göttliche Hülfe an. Und das ist ein rechtschaffen vollkommen Gebet.

128. Daneben aber soll man doch das schwache Gebet und der schwachen Menschen Gebet nicht verachten. Dieses Gebet aber, welches so ernstlich und brünstig ist, wird uns im höchsten Grad und als ein vollkommen Exempel vorgehalten, auf daß wir uns befleißigen, demselben nachzufolgen. Denn wie die Kinder in der Schule erstlich das ABC und den Donat lernen, also stammeln wir auch erstlich an dem unvollkommenen Gebet, welches mit nichten soll verachtet werden. Denn wie es geht mit dem Knaben, der nicht erstlich das Alphabet oder Donat lernt, daß derselbe auch nimmermehr den Virgili oder Katechismus wird reden lernen: also geht es mit dem Gebet auch zu. Wer nicht erstlich mit Schwachheit betet, der wird nimmermehr zum vollkommenen Gebet kommen. Darum soll man die Lehre an diesem Ort fleißig merken, welche anzeigt, wie das Gebet seine große Kraft habe, und wie das Gott so ein angenehm Opfer sei, wo man darin beharrt; wie Rebekka auch beharrt bis an das Ende, da kein Rath mehr vorhanden ist denn allein, daß sie die Kirche oder Gemeinde Gottes um Rath frage.

\subsection{IV}

129. Mose aber beschreibt sehr weitläufig, wie sie erhört worden sei. „Der Herr sprach“, nämlich durch Sem. Denn wir haben droben gesagt, daß Gott auch durch Adam geredet habe. Wie es Christus auslegt Matth. 19, 4 f.: „Habt ihr nicht gelesen, daß, der im Anfang den Menschen gemacht hat, der machte, daß ein Mann und Weib sein sollte, und sprach: Darum wird ein Mensch Vater und Mutter verlassen“ etc. Diese Worte hat Gott geredet, und Mose schreibt doch 1 Mos. 2, 24., daß sie Adam geredet habe. Derhalben ist Adam voll Heiliges Geistes und voll Gottes gewesen, welcher durch Adam geredet hat. Darum sind Adams Worte Gottes Worte. Also ist noch jetzt die Kirche und Gemeinde Gottes voll Gottes, und alles, was sie sagt, soll man nicht anders verstehen, denn als wäre es von Gott selbst geredet. Also gebietet Gott, daß wir den Eltern und der Obrigkeit gehorsam sein und von ihnen sein Wort hören sollen. Auf dieselbe Weise hat Gott der Herr auch mit dieser sehr heiligen alten Mutter Rebekka geredet durch Sem, welcher außer dem gemeinen Verstand des Glaubens auch voller Weissagung gewesen ist.

130. Es wird uns aber ein sehr großer und reicher Trost in diesen Worten vorgehalten, da im Text steht: „Zwei Völker sind in deinem Leibe“ etc. Denn Sem will also sagen: Du hast gebetet, und dein Mann und Abraham haben auch mit dir gebetet; meinst du aber, daß euer Gebet nicht sei erhört worden und daß du sterben werdest? Nein, das fürchte am allerwenigsten! Und Sem kommt gleich recht mitten in die Sache und auf den rechten Trost. Denn er sagt: So weit fehlet es, daß du in Gefahr deines Lebens kommen solltest, daß aus deinem Leibe zweierlei Völker kommen sollen. Du bist nur ein wenig angefochten worden, und dein Gebet ist versucht worden, ob du auch beharren wolltest. Darum sollst du wissen, daß mehr erhört oder erlangt, denn gebeten worden ist, und daß dir mehr gegeben sei, als begehrt oder verstanden gewesen ist. Denn du hast nicht verstehen können, was oder wie viel dir Gott geben wird. Er hat dein Gebet besser verstanden denn du selbst.

131. Dies Exempel ist wohl zu merken und wir sollen demselben auch nachfolgen. Denn wir beten auch wider den Antichrist, den Türken und die Rotten. Und wir haben gewiß sehr wichtige und verzweifelte Sachen, so gar, daß wir nicht denken oder verstehen können irgend einen Trost oder Hülfe, denn allein, daß wir Gott um Gnade anrufen, daß er eilen wolle mit der Zukunft des jüngsten Tages. Darum soll man nicht aufhören zu beten, sondern die starken sollen heftig schreien und die Schwachen sollen auch in ihrer Schwachheit zu Gott rufen; wie Christus selbst sagt Luc. 18, 7. von den Gottseeligen, die zu ihm rufen Tag und Nacht.

132. Es ist aber jetzt unsere Noth und Sache größer und beschwerlicher, denn sie jemals zu irgend einer Zeit in der Kirche gewesen ist. Denn wir haben wider uns die Höllen und das Ende der Welt, den äußersten grimmigen Zorn des Pabstes und Türken, welche uns gar verschlingen wollen und auch verschlingen werden. Darum soll man beten, und also beten, daß wir uns halten an das Gebot und Verheißung Christi. Denn wiewohl es sich läßt ansehen, daß in Kürze Alles werde über einen Haufen fallen: was liegt daran? Laß immer hinfallen. Dein Gebet wird doch noch mehr erlangen, als du bittest. Denn er hat es uns verheißen. Und wir bitten in unserem Gebet vor allen Dingen und auch vor unserem Unglück selbst, daß Gottes Name geheiliget und sein Reich erhalten möge werden. Wiewohl nun Alles verloren und eine verzweifelte Sache ist, und ob wir auch hier alle sollten getödtet werden: so weiß ich doch, daß die Widersacher sollen umkommen. Sie sollen herunter, denn das Gebet ist erhört worden.

133. So tröstet nun Sem Rebekka also: „Du liebe Rebekka, dein Gebet ist reichlich und über alle deinen Verstand erhört worden! Du hast den Herrn um einen Sohn gebeten, und begehrst jetzt in dieser Gefahr nur, daß du glücklich gebären und das Leben behalten mögest; ja, du würdest dir daran gerne genügen lassen, wenn du wissen möchtest, daß dein Gebet erhört wäre. Aber sei getrost und habe nur guten Muth der Geburt halben; denn Gott hat dich erhört und er versteht dein Gebet recht wohl. Du hast um einen Sohn oder Tochter gebeten, dieselbige Bitte wird er häufen und reichlich mehren. Wie soll aber das zugehen? Antwort: „Zweierlei Völker sind in deinem Leibe“ etc. Das heißt nun recht, daß das Gebet erhört ist. Denn es sind nicht nur zwei Söhne, die du gebären wirst, sondern zwei sehr große Patriarchen. Darum ist dein Gebet nicht allein darin erhört, daß ein Sohn von dir soll geboren werden, sondern daß du auch Nachkommen oder Erben haben sollst von zwei Söhnen und daß du zweierlei große Völker gebären wirst. Darum disputire nur nicht, verzage nicht, du trägst kein Ungeheuer oder Teufel, wie du dich deswegen hast befürchten müssen, daß sich die Kinder mit einander im Leibe gestoßen haben: sondern du bist jetzt schon eine Mutter zweier großer Völker.

134. Also hat sie allein gebeten um ihr und eines Sohnes Leben, und erlangt damit das Leben zweier Patriarchen und ihrer Nachkommen. Sie hat um einen Heller gebeten und erlangt einen goldenen Berg, welches sie weder zu hoffen noch zu glauben sich unterstanden hatte. Denn sie hat das Gebet bescheiden wollen sein lassen, und daß sie sich an Wenigem genügen ließ; wie wir auch pflegen um kleine und geringe Dinge zu bitten, und bedenken nicht, daß wir mit solcher großen Majestät reden, wenn wir beten. Denn wo er nur kleine und geringe Dinge geben wollte, so hätte er uns solche große und herrliche Form und Weise zu beten nicht vorgeschrieben, da er uns also hat heißen sagen: „Unser Vater, der du bist im Himmel, dein Name werde geheiliget; dein Reich komme“ etc. Gott ist kein karger Euclio, oder ein armer Trus, sondern legt uns vor und bietet uns große Güter an und die allerhöchsten Gaben im Himmel und auf Erden, und will haben, daß du dieselbigen von ihm bitten und darauf warten sollt. Denn in einer jeglichen Bitte, da wir im Vater Unser sagen: Dein Name werde geheiliget; dein Reich komme; dein Wille geschehe; unser täglich Brod gib uns heute etc., da wird Himmel und Erde, und alles, was darin ist, begriffen. Denn was ist es, daß sein Name geheiliget wird, daß sein Reich kommt, und sein Wille geschieht? Antwort: Es ist so viel, als unzählige Teufel darnieder legen und die ganze Welt in Einem Gebet verschlingen. Aber wir haben enge Herzen und sind von schwachem Glauben: darum sollen wir dies Exempel mit Fleiß merken und lernen, wie sich GOtt nicht genügen läßt, uns wenig zu geben, ob wir wohl um ein Weniges bitten; sondern er gibt mehr, denn wir verstehen oder bitten dürfen, Eph. 3, 20.

\quellseite{29}

135. Wir wollen aber nun die Worte des Patriarchen Sem besehen. Denn sie begreifen in sich vier Stücke, welche wir erstlich allein nach dem Gesetz verstehen wollen. Zum Ersten sagt er: „Zweierlei Völker sind in deinem Leibe“, und sind jetzt schon vor GOtt so groß, wie groß sie werden sollen. Wenn sie aber werden geboren und auferwachsen sein, alsdann werden sie getheilt werden. Dies ist der andere Theil: eine Verheißung, daß Edom und Jakob zweierlei Völker sein werden, nach dem Fleisch getheilt, nämlich, die Edomiter und Israeliten; wie die heilige Schrift solches bezeuget. Sie werden aber nicht unter Einem Haupte sein, sondern ein Jeder wird seine Ehe und Herrlichkeit haben wie der Fürst seines Volks. Sie werden auch Haus- und Weltregimente geordnet und angerichtet haben, dazu auch Kirchen: es wird aber doch ein Unterschied darin sein. Denn sie werden getheilt werden in mancherlei Religion, Polizei, weltliche Rechte und Gesetze. Denn es ist wohl glaublich, daß Esau aus dem Hause Abrahams mit sich genommen habe die Beschneidung und sonst etliche Weise zu opfern. Also wird sein Volk eine sonderliche Polizei, Ceremonien und Kirche haben. Jakob aber hat auch seine sonderliche Ordnung empfangen in der Polizei und Kirche, dergleichen in der Welt nicht gewesen ist, und hat zugleich dabei die Succession (Nachfolge) und Erbschaft der Verheißung gehabt. Darum werden diese Völker von einander getheilt werden beide nach der Polizei und nach dem Kirchenregiment. Und dasselbe ist auch leiblich erfüllet worden.

Denn St. Paulus hat in der Epistel an die Römer Cap. 9. 10. f. die geistliche Deutung gesetzt, davon an andern Orten ist gesagt worden. Die letzten zwei Stücke sind, da er sagt: „Und ein Volk wird dem andern überlegen sein, und der Größere wird dem Kleinern dienen.“

136. Edom oder Esau ist sehr gewachsen und gemehret worden, wie wir hernach sehen werden, da er seinem Bruder Jakob begegnet mit vierhundert Mann, da Jakob niemand bei sich gehabt denn vier Weiber. Und im sechunddreißigsten Capitel werden elf Fürsten erzählt, die das Land Edom inne gehabt. Daraus ist abzunehmen, daß Esau sehr große Gewalt gehabt habe. Jakob aber ist noch im Lande Canaan umher gezogen, und darnach auch nach Egypten hinab als ein Fremdling gekommen, da unterdeß von seinem Bruder elf Kindeskinder und große Fürsten sind geboren worden. Darum ist diese Weissagung auch äußerlich nach dem Buchstaben erfüllt worden. Aber zur Zeit Davids ist der Kleinere gemachet und gemahret worden, nämlich Israel, und hat Königreiche eingenommen und hat sich auch Edom unterthänig gemacht; wie im 60. Psalm 3. 10. stehet: „Mein Schuh trete ich über Edom.“ Hernach aber wird Isaac in Segen Esaus diese Verheißung auslegen und mildern, daß solche Dienstbarkeit bis auf eine bestimmte Zeit währen wird. Denn er sagt also, 1 Mos. 27, 40.: „Und es wird geschehen, daß du sein Joch von deinem Halse reißen wirst.“ Und dasselbe ist geschehen zur Zeit Ahab, da Edom ist abgefallen und sich einen eigenen König erwählt hat.

137. So ist nun dies die Summa der Verheißung und so reicher Erhörung. Denn also hat sie Sem getröstet: Siehe, meine liebe Rebekka, du hast für dich gebeten, und hast nur begehrt, daß du und dein Kind möchten beim Leben erhalten werden: ich aber verheiße dir viel größere Dinge. Du wirst nicht allein Einen Sohn haben, sondern zwei Söhne und Fürsten zweier Völker. Und über das zeigt dir GOtt auch an, was für einen Stand und Amt sie alle beide haben werden, auf daß du wissen mögest, was du werdest für Nachkommen und Erben haben auch nach deinem Tode. Einer wird stärker, der Andere etwas schwächer sein in der Polizei und Religion; aber das kleinste Volk wird der rechte Erbe und Herr bleiben. Der Stärkere wird den Schwächeren nicht gar unterdrücken; ja, der Kleinere wird über den Größeren herrschen. Darum hast du nun gehört, daß nicht allein deine Person, und die Kinder, so du gebären wirst, gesund und lebendig bleiben sollen, sondern daß auch hier beschrieben ist die Historie dieser beiden Theile bis auf alle Nachkommen ohne Ende.

138. Dies wird nun also beschrieben uns zur Lehre und Trost, daß wir lernen getrost und mit vollem Vertrauen beten, und daß wir an der Erhörung nicht verzagen oder zweifeln. Denn wiewohl das Gebet wider und über unsern Verstand geschieht, ist es doch nicht vergeblich. Denn GOtt will die Schmach auf sich nicht liegen haben, daß wir es sollten dafür halten, daß er nicht wahrhaftig wäre. Und darum hat er uns heißen beten, und hat solches nicht allein geboten, sondern auch verheißen, daß er unser Gebet hören wolle, auf daß wir es ja gewißlich dafür halten, daß er dasjenige geben wolle, darum wir bitten. Derhalben ist dies der höchste Gottesdienst und Religion, daß man glaube, daß GOtt wahrhaftig sei.

139. Ueber das haben wir auch eine feine, einfältige und gewisse Form, darin er uns vorgeschrieben hat, was und wie man beten soll. Darum, wenn du nun betest: Unser Vater, der du bist im Himmel etc., so denke also: Herr, du hast mich heißen beten, und hast auch verheißen, du wollest den, der da betet, erhören; und sollst an der Treue und Wahrheit GOttes, dazu auch an seinen Verheißungen nicht zweifeln. Denn ob schon der Teufel in Einem Augenblick die ganze Welt würde umkehren und gar über einen Haufen werfen, so wird doch gleichwohl der Name GOttes geheiliget werden. Und wenn auch wir alle sterben müßten, so werden wir doch ohne allen Zweifel behalten und bewahret auf die zukünftige Auferstehung. Ja, über das sehen wir auch in diesem Leben offenbar die Kraft und Wirkung unsers Gebets. Denn gewißlich, durch das Gebet wird das Regiment der Welt erhalten; es wird unser Leben und alle Güter, deren wir gebrauchen in dieser Welt, dadurch erhalten.

140. Es hat ferner diesen Text der heilige Paulus Röm. 9, 10. ff. etwas herrlicher gehandelt, und zieht denselben auf die rechte geistliche und wahrhaftige Historie, da er aus dem Propheten Maleachi Cap. 1. V. 2. 3. anführt, da GOtt spricht: „Jakob habe ich lieb, und hasse Esau“; item Röm. 9, 12.: „Es ward zu ihr gesagt, nicht aus Verdienst der Werke, sondern aus Gnaden des Berufers, also: Der Größere soll dienstbar werden dem Kleinern“ etc. Da berührt er und gedenkt nicht der zeitlichen, sondern der rechten wahrhaftigen Verheißung. Die Juden verstehen dies allein nach dem Fleische, als nämlich, daß sie sind aus Egypten geführet worden und das Land Canaan eingenommen haben. Aber in diesen leiblichen Verheißungen werden auch die geistlichen begriffen. Denn aus demselben schwachen Theil hat der Sohn GOttes sollen geboren werden, auf daß das Heil nicht von den Heiden und auch nicht von den Idumäern, sondern von den Juden herkommen möchte, wie Joh. 4, 22. gesagt wird. Es ist alles um Christi willen geschrieben, welcher aus der Linie des kleineren Volks hergekommen ist.

141. Derhalben handelt St. Paulus diesen Text auf solche Weise, daß er zweierlei Geburt macht, nämlich, des Fleisches und der Verheißung. Esau ist auch ein Sohn der Verheißung gewesen, aber da scheidet es sich von einander und werden getheilt in ein groß und ein klein Volk. Da legt es Paulus nach seinem großen und apostolischen Geiste also aus, daß der Größere nicht werde der Erbe sein, sondern der Kleinere, nämlich Jakob. Also werden die Andern auch alle Erben, die der Verheißung folgen und sich daran halten, wie droben davon weitläufig ist gesagt worden. Darum gehört die leibliche Verheißung dem Esau, aber nach dem Geist gehört sie dem Jakob, welcher der Steckel kleiner Sohn war nach dem Fleisch; aber nach dem Geist ist er der Herr und der Größere gewesen.

142. Wie ich aber von der Kainischen und Ismaelithischen Kirche gesagt habe, nämlich, daß sie sind verworfen worden, aber doch also, daß solche Verwerfung dazu hat dienen sollen, daß sie gedemüthigt würden, auf daß sie auf das Erbe verzichteten, daß sie sich hatten vermessen, daß sie dasselbe haben würden von wegen der fleischlichen Geburt; sie sind aber noch selig geworden durch Buße und Glauben an die Verheißung: also bezeugt auch der Text, daß Viele, so von Edom gekommen, selig sind geworden, nicht darum, daß sie Kinder Edoms waren (denn dieselbe Linie ist verworfen), sondern weil sie mit dem Glauben die Verheißung ergriffen und sich daran gehalten haben: nach dem Spruche Pauli, da er Röm. 9, 8. sagt: „Nicht sind das Gottes Kinder, die nach dem Fleisch Kinder sind; sondern die Kinder der Verheißung werden für Samen gerechnet.“ Derhalben haben sie sich zu den Ismaeliten geschlagen und gesagt: Sarah ist unsere Großmutter und Abraham unser Vater, aber dadurch werden wir nicht selig werden; denn die leibliche Geburt ist kein nütze; sondern wir glauben an den Samen, welcher den Vätern verheißen ist, wie Isaak und Jakob geglaubt haben und selig geworden sind. Und alle, die diesen Glauben gehabt, haben das Erbe des ewigen Lebens erlangt. Aber die aus diesem Geschlechte oder Nachkommen in solcher Verheißung geblieben sind, daß sie die Größeren wären von wegen der leiblichen Geburt, die sind alle verloren und verdammt worden.

\quellseite{30}

143. Also können wir Heiden uns nicht rühmen, wie sich der Türke rühmt von wegen seiner vielen Königreiche, oder wie sich der Papst rühmt der ordentlichen Succession und höchsten Gewalt in der Kirche. Aber deß rühmen wir uns, daß Joh. 1, 12. 13. geschrieben steht: „Wie viele ihn aber aufnahmen, denen gab er Macht, Gottes Kinder zu werden, die an seinen Namen glauben. Welche nicht von dem Geblüt, noch von dem Willen des Fleisches, noch von dem Willen eines Mannes, sondern von Gott geboren sind.“ Zu dieser Erkenntniß müssen wir alle kommen, wie Paulus Röm. 4, 12. Abraham einen „Vater“ nennt „der Beschneidung, nicht allein derer, die von der Beschneidung sind, sondern auch derer, die wandeln in den Fußstapfen des Glaubens, welcher war in der Vorhaut unsers Vaters Abraham.“ Also sagt auch Christus Joh. 8, 39.: „Wenn ihr Abrahams Kinder wäret, so thätet ihr Abrahams Werke.“ Denn derselbe hat Gott geglaubt und hat durch den verheißenen Samen Gerechtigkeit und Seligkeit erlangt. Also, die Gott glauben, werden auch Abrahams und Gottes Kinder.

144. Darum soll man aus diesem Text lernen, daß alle Ehre oder Ruhm des Fleisches von Gott verworfen und verdammt sei, und daß wir dadurch nicht können gerecht werden. Wie auch die Väter durch solchen Ruhm des Fleisches oder Geblütes nicht sind gerecht geworden, welches dies Exempel gewaltiglich lehrt und beweist, und ist ein sehr trefflich Exempel, das auch werth ist, daß man es mit goldenen Buchstaben schreiben sollte.

145. Es haben sich aber diese Brüder mit einander gestoßen im Mutterleibe, und da sie geboren sind, haben sie sich darnach auch die ganze Zeit ihres Lebens gezankt. Also müssen wir auch immer mit dem Papst und Türken streiten, und können mit ihnen nimmermehr Gemeinschaft oder Friede haben. Ismael und Esau, das ist, dieselben Größeren und Stärkeren, die uns weit übertreffen mit Geld und Gut, mit Gewalt, mit großer Menge und auch mit Weisheit, dieselben sind wider uns, die wir die Jüngeren und Schwächeren sind. Aber dies ist unser Trost, sintemal wir uns ja mit einander müssen stoßen, und mit den Edomitern solchen stetigen Krieg führen, welcher nicht vertragen oder aufgehoben kann werden, nämlich, daß wir des Sieges gewiß sind. Denn das Wort oder die Verheißung wird ohne allen Zweifel gewinnen müssen.

146. Wo aber etliche von den Widersachern selig werden, soll man wissen, daß sie nicht darum selig werden, daß sie Papisten, Mönche oder Türken sind, sondern weil sie sich zu uns begeben, das ist, zu Abraham, Isaak und Jakob. Und wiewohl sie von unserem Samen nicht sind, so sind sie doch eingepfropfet. Also wird nicht der ganze Teig verworfen, sondern nur der Stolz und die Hoffahrt des Teiges, so sich der Väter, der Ceremonien und des Gesetzes rühmt; und wo er mit uns will eingepfropfet sein, so wird es ein heiliger Teig mit uns sein und mit dem ganzen Geschlechte der frommen gottseligen Menschen. Aber solches geschieht aus Gnaden, nicht aus Gerechtigkeit des Fleisches oder durch Verdienst der Werke, darin die Welt ihr Heil und Gerechtigkeit sucht.

\section[Viertes Stück]{Viertes Stück. Von der Geburt Esau’s und Jakobs; von ihrer Lebensart, und von dem Verhalten Isaaks und Rebekkas gegen diese ihre Söhne}

\subsection{I}

\begin{quote}
\textbf{V. 24—26. Da nun die Zeit kam, daß sie gebären sollte, siehe, da waren Zwillinge in ihrem Leibe. Der erste, der heraus kam, war röthlich, ganz rund wie ein Fell; und sie nannten ihn Esau. Darnach kam hernach sein Bruder, der hielt mit seiner Hand die Ferse des Esau; und hießen ihn Jakob.}
\end{quote}

147. Die wunderbare und gefährliche Geburt wird allhier beschrieben, dadurch die Mutter ohne Zweifel sehr wird erschrocken sein, da sich die Kinder also in ihrem Leibe mit einander gestoßen haben. Denn es wird hier eine solche Geburt beschrieben, daß wo die Verheißung Rebekka nicht hätte erhalten, so wäre sie alsbald in der Geburt gestorben. Und ich verstehe es wahrlich nicht, habe auch nie kein Exempel weder gelesen oder gehört, das diesem gleich sei. Denn also lautet der Text: „Da nun die Zeit kam, daß sie gebären sollte, da waren Zwillinge in ihrem Leibe.“ Denn die Weiber, so dabei waren, haben nur Hoffnung gehabt auf Ein Kind; bald aber, da sie zusehen, verstehen sie und werden gewahr, daß es Zwillinge sind.

148. Nun ist dies wohl gewöhnlich, daß Zwillinge geboren werden, daß darüber die Mutter in keine Gefahr kommt: aber hier ist große Gefahr, da der Erste, nämlich Esau, geboren worden, daß daselbst auch bald Jakob heraus kommt, und da daselbst also, daß er mit der Hand die Ferse seines Bruders gehalten hat. Wir wissen aber, daß es eine unglückselige, schwere Geburt ist, wo ein Kindlein nicht recht ordentlich geboren wird, also daß erst das Haupt heraus komme: sondern wenn die eine Hand oder Fuß erst hervor kommt, daselbst kommt beide, Mutter und Kind, in Gefahr Leibes und Lebens. Es wird auch oft in solcher Noth die Frucht zerrissen und in einzelnen Stücken herausgezogen. Also wird diese Geburt auch beschrieben, daß sie wunderbar und voller Gefahr gewesen sei, so viel die Mutter und Jakob betroffen hat. Der Erste kommt frisch und gesund heraus, aber der Andere kommt in Gefahr seines Lebens. Die Verheißung hat aber Rebekka erhalten, darin sie gehört hatte, daß sie zweierlei Völker oder zweier Völker Häupter gebären würde. Darum ist sie in solcher großen Noth und Gefahr nicht verzagt.

149. Es beschreibt aber Mose die Kinder beide: „Der erste“, sagt er, „der heraus kam, war röthlich.“ Da steht im Hebräischen admoni. Davon hat Adam seinen Namen; denn er ist von röthlicher Erde gemacht worden, davon Plinius und Josephus sagen, daß es die beste Erde sei soll. Nach derselben wird die schwarze Erde gelobt; aber die gelbe und falbe Erde ist etwas geringer. Edom heißt röthlich. Es wird aber eine sonderliche Röthe an ihm müssen gewesen sein. Denn es werden sonst auch alle Kinder roth geboren und nach der Geburt sind sie auch röthlich, Esau aber wird sehr röthlich gewesen sein, mehr als seine Brüder oder andere Kinder.

150. Darnach hat er auch noch ein ander Zeichen, dadurch er auch von seinem Bruder ist unterschieden worden. Denn er ist rauh wie ein Fell nicht allein am Haupte, wie andere Kinder sind und vornehmlich wie die Knäblein viel Haare pflegen zu haben; sondern über das, daß Esau röthlich gewesen, hat er auch gar eine rauhe Haut wie ein Fell. In Summa, er ist an Händen und Füßen und am ganzen Leibe rauh. Denn Mose setzt dies Gleichniß dazu und sagt: wie andrereth, wie das hebräische Wort lautet, das ist, wie ein Fell. Es ist wahrlich eine wunderbare Beschreibung eines Kindes, das so röthlich und rauh gewesen ist, dergleichen man kaum in den Historien der Heiden finden soll, darin doch oftmals von ungewöhnlichen Geburten Meldung geschieht.

151. Das hebräische Wort heißt ein Fell oder ein rauher Mantel; wie im andern Buche der Könige im 2. Capitel V. 8. das Kleid Elias beschrieben wird, welches ein rauher Mantel gewesen ist, damit er in das Wasser des Jordans geschlagen, daß es sich auf beide Seiten theilete etc. Und im Buch Josua im 7. Capitel V. 21.: Vidi inter spolia addereth schineaar, welches hier im Hebräischen sear genannt wird, das ist: „Ich sahe unter dem Raube einen köstlichen babylonischen Mantel“ etc. Ob sie aber zu der Zeit aus Wolle Kleider gemacht haben, kann ich nicht wissen. Dieser Art und Gestalt, halte ich, sei gewesen das lederne Kleid, schön und lang, wie ein Mantel oder ein langer Rock sein mag, welches der Propheten Kleidung gewesen ist. Dies kann aber auch wohl verstanden werden, daß nach der Figur, Hyperbel genannt, geredet sei, da man ein Ding so gar groß macht; wie wir in unserer deutschen Sprache der Gleichnisse auch haben, als, da wir von denen, so
\quellseite{31}
lange, scheußliche, verworrene Haare haben, sagen, daß sie gestaltet seien wie ein Igel oder eine Haareule. Also haben die Weiber vom Esau auch reden mögen, wie man etwa im Hause pflegt zu reden, und ein Ding größer zu machen, denn es an sich ist.

152. Da nun Esau geboren ist, wird auch die Geburt Jakobs beschrieben. Derselbe ist in Gefahr gewesen beide seines eigenen und auch seiner Mutter Lebens halben. Denn er folgt seinem Bruder so gar bald nach, daß er auch seine Ferse hat ergriffen und sich daran gehalten. Und man mag denken, daß sie schier in einem Augenblicke zur Welt geboren seien. Darum sind dies wunderbare Dinge, dergleichen man in keinen Historien mehr findet. Denn es ist schlechthin und wahrhaftig ein göttlich Werk, das Gott allhier wirkt, und darum hat auch Mose diese Dinge so gar genau und eigentlich beschreiben wollen, auf daß ja fleißig eine jegliche Geschichte möchte gemerket werden, und daß der Leser erweckt und erinnert werde, die Größe dieser Dinge recht zu betrachten.

153. Denn allhier wird bestätigt das Fundament und der Grund der ganzen christlichen Lehre, und Gott will durch die Geburt dieser Zwillinge der ganzen Welt im voraus das Urteil sprechen, ja, alle Gerechtigkeit des Fleisches damit abschneiden und unterdrücken. Er will damit lehren, daß alle Weisheit und Herrlichkeit und Vorzug des Fleisches verloren und vergeblich sei; das will er, sage ich, durch dies allein gelehrt und angezeigt haben, daß er gesagt hat: „Der Größere wird dem Kleinern dienen.“ Darum ist dies keine Disputation von Königreichen und Herrschaften der Welt, vielweniger von andern geringen und unnützen Dingen, daran nichts gelegen, oder von altvettelischen Fabeln; sondern diese Worte legen den rechten Grund unserer christlichen Lehre. Und wird derselbe auch dadurch gewaltiglich bestätigt, nämlich, daß vor Gott das Fleisch todt und verdammt sei, aber der Geist lebendig gemacht wird. Und wird damit angezeigt die angeborene Sünde und verderbte Natur, auf daß sich ja niemand auf das Fleisch verlassen, sondern jedermann sein Vertrauen auf den Herrn und auf seine Verheißung setzen soll.

154. Dasselbe kann nicht genugsam gesagt oder den Leuten eingeschärft werden. Denn es sind die höchsten Gerichte und die unerforschlichen Wege Gottes. Ich rede jetzt nicht von den Gerichten Gottes a priori, wie man es nennt, welche er bei sich selbst in seinem innerlichen und heimlichen Rath hat, nämlich, warum er auf diese oder eine andere Weise rathe, dies oder jenes thue, regiere, helfe, verderbe etc.: sondern ich rede von Gott, wie er die Leute aus Gnaden beruft, mit ihnen redet und sich offenbart a posteriori; wie er zu Mose sagte: „Du wirst mir hinten nach sehen“, 2 Mos. 33, 23. Denn wir sehen zwar allhier solche Wege und solche Gerichte, welche kein Mensch verstehen oder begreifen kann, wie gesagt wird 1 Cor. 2, 14.: „Der natürliche Mensch vernimmt nichts vom Geist Gottes; es ist ihm eine Thorheit, und kann es nicht erkennen“, also daß auch die Heiligen nicht verstehen können, was sie bitten. So gar sind uns diese Gerichte und Wege Gottes unerforschlich, a posteriori, wie man es nennt, das ist, des Gottes, der uns beruft, und sich uns offenbart im Worte und äußerlichen Werke. Können wir nun diese Dinge nicht begreifen, wie sollten wir denn begreifen können, daß Gott in sich selbst und in seiner Gottheit vor uns verborgen ist? Darum hast du damit genug zu thun, daß du dich dessen befleißigest, daß du verstehen mögest, wie Gott wirkt und uns beruft. Es ist nicht vonnöthen, daß du suchest oder forschest die großen Geheimnisse und solche Dinge, die deinem Verstande zu hoch sind.

155. Denn ist das nicht die größte Blindheit des Fleisches des Menschen, daß es sich so schwerlich ärgert an diesem Exempel des göttlichen Gerichtes, nämlich, daß Esau, der Erstgeborne, nicht soll der Erbe sein, sondern Jakob, welcher der Kleinere und nach ihm geboren ist? Dasselbe wirst du mit menschlicher Vernunft nimmermehr verstehen, wiewohl du hörst, daß ihn Gott beruft, und diese seine Worte selbst klarlich redet, da er sagt: „Der Größere wird dem Kleinern dienen“, ja, daß er auch um derselben Worte willen die Könige und Königreiche dieser Welt zerstört: so wird doch gleichwohl durch die Vernunft dies Gericht Gottes verdammt werden, da er sagt: „Der Größere wird dem Kleinern dienen.“ Du wirst den Türken dessen nimmermehr überreden, daß er vor Gott verdammt sei; den Papst auch nicht, daß der Stuhl zu Rom des Teufels Stuhl sei; wirst dessen auch den Kaiser Karl den Fünften, Ferdinand und den König von Frankreich nicht überreden. Denn das Fleisch, das ist, der Größere, stolzirt, legt und ist gleichsam Alles in der Welt. Dagegen aber ist der Geist betrübt, ist todt und gar nichts. Wie Esau denkt: Ich bin der Erstgeborne, darum bin ich auch die Kirche, bin der Sohn und Volk Gottes; Jakob aber ist verworfen, ist ohne Ehre und hat gar keinen Namen. Wie wir zu unsern Zeiten erfahren, daß unsere Widersacher auch den Titel oder Namen der Kirche mit Gewalt an sich reißen, so sie doch nichts weniger sind denn die rechte Kirche.

156. Derhalben ist Esau der Erstgeborne, Jakob ist der Letztere. Aber Gott kehrt diese Ordnung um und sagt: Das Fleisch ist stolz und hoffärtig; der Geist ist traurig und trägt Leid: jedoch ist es mit dem Fleische eitel Lüge, der Geist aber ist die Wahrheit. Denn alle Menschen sind nicht allein Sünder, sondern haben auch sich selbst allzu lieb, überreden sich allezeit dessen selbst und sind sehr vermessen, als ob sie es immer besser wollen wissen, und verachten und hassen auch dies Urteil Gottes und diese Lehre. Hebe an vom Anfang der Welt und betrachte nach einander die Exempel zu allen Zeiten. Kain, der Erstgeborne, ist Herr, Abel ist Knecht. Das Königreich Egypten ist wie eine schöne Blume in der Welt und ist gleichwohl des Teufels Dreck. Darnach das Königreich der Assyrer, Babylonier und der Griechen, dazu auch der Römer, herrscht; die Israeliten aber und das Volk Gottes auch im Neuen Testamente wird fast hart mit schwerer Dienstbarkeit gedrückt.

157. Wider diese unüberwindliche Hoffart des Fleisches wird dies geschrieben, auf daß wir wissen, wie wir uns halten sollen in solcher großen Ungleichheit des Glückes, wie wir ein rechtes gewisses Urteil, dazu auch einen beständigen und starken Muth haben und behalten mögen, solche widerwärtige Dinge zu ertragen. Als wenn ich sehe, daß der Türke oder Papst in der größten Majestät und Herrlichkeit sitzen, soll ich sagen: Du hast kein Gottes Wort, bei dir ist das Wort nicht, darum magst du dich rühmen, so viel du kannst, du magst auch gewaltiger sein denn alle Engel, so bist du doch gleichwohl verdammt. Denn das Fleisch oder der Größere soll dem Kleineren dienen.

158. Aber der Türke und die Päpstischen sagen dagegen also: Was Wort! Was Elia, Elia! Was Luther oder Philippus! Ich habe Königreich, Geld und Gut, Gewalt, groß Vermögen, schöne Gestalt und dazu auch Weisheit, dessen sich die ganze Welt rühmt und damit über Andere herrscht. Darum, sagen sie, bin ich selig und bin Gottes Volk; ihr aber seid arme Bettler und werdet geplagt, und derhalben seid ihr auch von Gott verworfen.

159. Diesem so großen Aergerniß setzen wir dies entgegen und sagen: Wir gehören dem Gott zu, der uns aus Gnaden beruft, der sich offenbar und klärlich heraus sagt: „Der Größere wird der Kleinste sein“, das ist, das Fleisch, so da herrscht, soll dienstbar werden, und der Geist, so dienstbar ist, soll Herr werden. Dies ist unser Ruhm, welcher zwar von der Welt verlacht und von den Gewaltigen dieser Welt verspottet wird; wie bald im Anfange Kain geherrscht hat, Abel aber dienstbar gewesen ist. Doch ist Kain verworfen worden, Abel aber lebt. Also, die Assyrer und Babylonier herrschten auch und verderbten das jüdische Land, Jerusalem und den Tempel. Also herrschten die Perser auch, die Juden aber waren gefangen. Doch sind die Juden geblieben, jene aber sind gar in Grund verdorben und umgekommen. Also auch, die Römer haben mit großer Hoffahrt das Regiment inne gehabt, dagegen aber haben die Apostel im Elende müssen umher ziehen und flüchtig werden, Hunger und Durst leiden, waren nackend und gleichsam zum Tode verurtheilt, wie 2 Cor. 4, 11. geschrieben steht. Jeglicher Zeit herrscht und regieret der Türke, die Christen aber sind dienstbar; der Papst ist stolz und hoffärtig, die Kirche ist elend und betrübt. Doch wird endlich der Papst fallen und mit dem Geiste des Mundes Christi umgebracht werden. Es ist wohl wahr, spricht Gott, die Welt ist größer, die Kirche und mein Volk ist kleiner, ich aber urtheile also, daß der Größere wird der Kleinere und der Kleinere der Größere werden.

160. Dies ist eine Lehre vom Glauben, Hoffnung, Trost und Liebe gegen Gott und Menschen. Es ist kein anderer Unterschied, denn allein, daß hier in der Kirche Gott ist, welcher uns durch sein Wort beruft, dort aber in der Welt und bei den Gottlosen schweiget Gott stille. Der Türke hat auch nicht einen Buchstaben, damit er beweisen könne, daß recht sei, was er thut. Da ist kein GOtt, der da berufe. Also hören der Papst und Kaiser auch nichts, es ist bei ihnen gar stille: wir aber haben einen GOtt, der uns beruft. Denn er hat sich uns geoffenbart, daß wir an ihm einen solchen GOtt haben, den wir sehen, fühlen und ergreifen mögen. Wir haben das Wort, die Taufe, die Schlüssel und müssen doch dabei Verfolgung leiden. Nun, es ist gut, lasset uns leiden, was wir sollen, es werden doch endlich die Knechte Herren und wiederum die Herren Knechte werden. Wir werden bleiben, sie aber werden untergehen. Unsere Nachkommen, die wahrhaftig die Unsern sind und uns angehören, werden bleiben auch hier auf Erden; wie immer von Anfange der Welt etliche Propheten geblieben sind, und haben uns, gleichsam als erblich, die Stimme und Wort GOttes hinterlassen und gelehret.

\quellseite{32}

161. Dies Argument hat GOtt der ganzen Welt wollen vorhalten, und mit dieser Historie beweisen und bekräftigen. Darum sollen wir dieselbe nicht so schläfrig oder unfließig lesen, wie sonst etwa eine weltliche oder heidnische Historie. Denn wo man sie nicht mit Fleiß lieset, ist sie gar kalt, schlecht und todt, gleich als wenn ich eine Erzählung lesen würde von der Geburt oder von den Thaten eines türkischen Kaisers. GOtt hat aus sonderlichem Rathe seine Lehre mit diesem Exempel anzeigen und sein Werk offenbaren wollen in diesen zwei Brüdern, auf daß wir sehen sollten, wie der Größere würde dem Kleineren dienen müssen. Esau ist der Erstgeborne, und ist derhalben sehr stolz, fährt hoch her: er ist ein chaldäischer, babylonischer, persischer, ein griechischer und römischer Kaiser, und in Summa, er ist der Papst und ein Herr aller Herren. Er hat aber keinen GOtt, der ihn berufe, sondern der ihn verdammt mit diesem Worte, da er sagt: „Der Größere soll dienstbar sein“; wiewohl er dieser Stimme nicht glaubt. Jakob aber hat GOtt, der ihn beruft und ihm die Herrschaft gibt über den Größern: er ist aber nicht gewaltig, weise oder heilig. Jedoch, dieweil er berufen ist, wird er das Regiment haben, wie die Propheten dies fleißig dargelegt, daraus auch vielfältigen Trost genommen haben; wie im 41. Cap. Jesaias V. 14. steht: „So fürchte dich nun nicht, du Würmlein Jakob, ihr armer Haufe Israel“ 2c.

162. Dasjenige hat GOtt alsobald in dieser Geburt anzeigen wollen. Esau kommt zuerst heraus, ist ein rechter Edom, ist röthlich und ganz rauh wie ein Fell. Der Andere aber hat eine glatte Haut, und ist nicht röthlicher, denn sonst Kindlein sind, wenn sie geboren werden. Es scheint an ihm nichts, das sonderlich vortrefflich wäre. Schau aber, welcher der Erste ist, soll vortrefflich sein vor seinem Bruder zweier großer Zeichen halben, und haben ihn nicht allein die Eltern, sondern auch die Freunde und Nachbarn, so dabei sind, wie mit seinem Munde Esau genannt, dazu sie durch diese zwei Zeichen sind bewogen worden. Denn menschliche Weisheit kann anders nicht urtheilen, denn daß Esau werde der Erbe sein, darum daß er der Erstgeborne und gezeichnet ist mit zwei trefflichen Wahrzeichen, nämlich, daß er so röthlich und rauh ist.

163. Er hat aber den Namen vom Wirken oder vom Thun. Als wollten sie sagen: Dies wird der Mann sein, welcher Alles herrlich ausrichten wird, und wie man pfleget zu sagen, er wird Factotum sein, der Alles gar allein thun wird. Wie Kain seinen Namen daher bekommen hat, daß seine Mutter meinte, sie hätte den Mann, den Herrn bekommen, 1 Mos. 4, 1., von dem GOtt verheissen hatte, daß er die Sünde und die Schlange überwinden würde. Also ist nun die Welt beschaffen, wenn sie von so hohen Dingen urtheilt. Darum sollen wir besser urtheilen lernen von den göttlichen Urtheilen, welche vor der Welt Thorheit sind. Denn sie ist der rechte Esau und Factotum, wie der Türke und Papst auch Esau ist. Wir sind verachtet und verworfen. Aber im Wort und Beruf GOttes kann dies alles sein unterschieden werden.

164. Sie haben aber drei Gründe, warum sie Esau so hoch rühmen. Als zum Ersten, daß er der Erstgeborne ist; darnach, daß die Eltern mit großer Begierde und Verlangen ganzer zwanzig Jahre auf das Kind gehofft und gewartet, und dasselbe mit großer Bitte endlich erlanget haben; zum Dritten, daß er röthlich und rauh ist. Welches alles dazu beigetragen hat, daß sie alle sich von diesem Kinde etwas Sonderliches versprochen haben, und dasselbe wider das göttliche Urtheil.

165. Also pflegen wir nämlich die Zeichen, so wider uns sind und damit GOtt den Gottlosen drohet, auf das Gegentheil zu deuten, gleich als ob sie für uns seien. Wir klopfen uns selbst und trösten uns, gleich als ob uns nichts Widerwärtiges widerfahren könnte; wie die Juden vor der Zerstörung Jerusalems auch gethan, da der Tempel in der Nacht ohne Menschen Hände geöffnet worden, welches ein Zeichen war, daß die Stadt mit dem Tempel würde zerstöret werden und daß die Heiden ohne alles Hinderniß in den Tempel gehen würden: dasselbe haben sie zum allerbesten und für ein gut Zeichen angenommen. Der Herr, sagten sie, wird jetzt bei uns sein, er geht von selbst zu uns herein. So gar hat sich unsere Vernunft selbst lieb, und alles, was Gutes gesagt wird, zieht sie auf sich selbst. Als, wenn Einer gelobt wird, meint er alsbald, es sei wahr, und sei er dessen wohl werth, was von ihm gesagt werde: so gar ist der Mensch der ärgste Heuchler, der sich selbst so fein schmeicheln kann. Aus bösen Zeichen nehmen wir etwas Gutes ab, und was böse ist, ziehen wir immer auf andere Leute.

166. Also sind diese zwei Zeichen, von denen sie meinen, daß sie dazu gegeben seien, daß dadurch der Erstgeborne geehrt und gepriesen würde, sehr böse Zeichen, nämlich, röthlich und rauh sein. Denn die Röthe bedeutet in der heiligen Schrift die blutdürstige Gerechtigkeit, rauh sein aber bedeutet Tyrannei oder Grausamkeit; wie auch wiederum weiß und hell sein etwas Gutes bedeutet. So ist nun allhier in dieser Geburt bedeutet und angezeigt worden, daß Alle, die vom Geblüt sind, unversöhnlich, blutdürstig und greulich sind, ohne alle Barmherzigkeit. Darnach, die rauhe Haut ist betrüglich und bedeutet Heuchelei. Dies sind wahrlich sehr böse Zeichen und sind wider die Vernunft, darum zieht sie dieselben auf ihren Nutzen und deutet sie zum allerbesten.

167. Wir sollen aber merken, daß die rauhen und abschäulichen Heuchler am Leibe Satans sitzen. Denn die rauhe Haut ist eine falsche, lügenhaftige Haut; die natürliche Haut aber soll nicht rauh sein. Darnach gehört die röthliche Farbe den Todtschlägern zu. Das sind zwei greuliche Zeichen an der Kirche und am Leibe des Teufels, nämlich, die Lügen und Mord. Sie aber erkennen diese Zeichen nicht, sondern rühmen sich von wegen ihrer Gerechtigkeit und Barmherzigkeit; wie Christus sagt

\begin{quote}
Joh. 16, 2.: „Wer euch tödtet, wird meinen, er thue GOtt einen Dienst daran.“
\end{quote}

Wird denn nun der, so Christum und die Seinen tödtet, angesehen, daß er rauh und röthlich sei? Ja, er wird angesehen, daß er ganz weiß sei wie die Sonne und heller scheine denn alle Sterne. An den Heuchlern scheint äußerlich nichts Rauhes, sondern die höchste Gerechtigkeit, Gottesdienst und Wahrheit. Darum soll man fleißig merken, daß Esau die rechten, wahrhaftigen Zeichen mit sich bringt, so eigentlich zu des Teufels Kirche gehören, da er beschrieben wird, daß er derselbigen Teufelskirche Haupt sei. Der andere aber, Jakob, ist nicht röthlich, sondern fein glatt und weiß.

168. Auf solche Weise hat GOtt diese Historie abgemalt nicht allein mit Verheißungen, sondern auch mit Zeichen, diese Meinung und Verstand damit zu befestigen, nämlich, daß das Fleisch dem Tode und Teufel unterworfen sei, wie gewaltig, herrlich und tausendmal größer es auch ist, denn der Türke und sonst alle Königreiche. Jakob aber und die rechten göttlichen Menschen sind in der Welt gar nichts, jene sind Alles. Aber das ist der Unterschied: bei uns ist die Stimme und das Wort, auf jener Seite ist das Stillschweigen. GOtt aber ruft denn, das nichts ist, daß es sei, und zerbricht, was etwas ist, daß es nichts sei. Denn GOttes Wort ist größer denn unzählige viele Welten. Darum hat uns GOtt nicht ohne sonderlichen Rath diese zwei sehr klaren Zeichen vorgestellt, welche allen Menschen vor Augen und Ohren stehen sollten. Esau ist das Thun und Werk, so GOtt gethan hat: wir sind Esau nicht, wir sind, wir sind gar nichts. Diese Historie aber lehrt gar das Gegentheil, bestätiget und vergewissert uns des gewaltig, nämlich, daß alles, was in der Welt und am Fleisch ist, Nichts sei, und daß allein GOttes Wort Alles sei.

169. An diesem Ort soll man das auch noch merken, dieweil Isaak die Zeichen und Geburt seiner Kinder nicht so bald nach dem rechten Verstande der Verheißung versteht, daß die Heiligen oder Gläubigen nicht sobald auf einmal Alles lernen und verstehen können; wie wir droben auch gehört haben, daß Abraham und Sarah auch haben müssen zunehmen und immer fortschreiten, und daß sie die göttliche Verheißung nicht sobald verstanden haben.

\quellseite{33}

Derhalben müssen die Heiligen durch gemeinen Gebrauch oder tägliche Uebung immer zunehmen von Klarheit zu Klarheit. Abraham meint, er würde von Sarah Kinder und Erben haben; aber da er sieht, daß sie veraltet und unfruchtbar war des Alters halben, da nimmt er Hagar zum Weibe und hält dieselbe für die Mutter des Samens, bis ihm aufs neue durch des Heiligen Geistes Auslegung die Verheißung ist gedeutet worden, nämlich, daß Sarah die Mutter werde sein des Samens, so ihm verheißen war. Also sind die Väter von Tage zu Tage immer mehr und mehr geübt worden, daß sie also zugenommen haben; wie St. Petrus vermahnet 2 Epist. 3, 18.: „Wachset aber in der Gnade und Erkenntniß unsers Herrn und Heilandes Jesu Christi.“ Sie sind nicht sobald zu vollkommener und völliger Erkenntniß gekommen, sondern mancherlei Uebung, Erfahrung, Anfechtung und Gefährlichkeit haben sie unterwiesen und gelehrt.

170. Also haben allhier Isaac und Rebekka gehört, daß zweierlei Völker sollen werden und daß der Größere dem Kleineren dienen werde; da haben sie also gedacht: Der, so der Erstgeborne ist, wird ohne Zweifel der Größere sein; und meinen, der Streit habe nun in Mutterleibe aufgehört. Darum nennen sie den Ersten Esau, das ist, einen solchen Poeten, der Alles ausrichten werde. Als wollten sie sagen: Jakob sollte der Größere gewesen sein, aber er ist in Mutterleibe von seinem Bruder durch Schickung GOttes überwunden worden, welcher zuerst heraus ist gekommen und dazu mit diesen zwei Zeichen geziert ist. Darum wird von Jakob nicht gesagt: „Sie nannten ihn“ etc., wie der Erstgeborne durch ihrer aller Bewilligung Esau genannt wird, sondern der Vater hat ihm den Namen gegeben.

171. Und zwar hat er den Namen von der Fußsohle, als wenn man sagen wollte, ein Untertreter. Doch ist diese Namengebung mehr ein Wunsch, als eine bestimmte Aussage. Esaus Name dagegen ist nur bestimmte Aussage, kein Wunsch. Derselbe wird solche Dinge ausrichten. Dieweil aber Esau, so im Mutterleibe der Größere gewesen, in der Geburt der Kleinere geworden ist und vom Großen überwunden ist, darum hat der Vater also gedacht: GOtt gebe, daß du ein rechter Untertreter mögest werden, wie noch an äußerlichen Gebärden zu sehen ist, daß du den Größern hast untertreten wollen, sintemal du seine Ferse ergriffen hast; aber die göttliche Stimme hat es umgekehrt. Derhalben wünscht er, daß er das sein möge, das er hat sein wollen, nämlich, ein Untertreter, und dasselbe zwar nicht seines Bruders, des Erstgebornen und des Ueberwinders, sondern daß er ein Untertreter sollte sein der Heiden und anderer Völker.

172. Also, meine ich, habe Isaac die Verheißung und Geburt seiner zwei Söhne anfänglich verstanden. Denn die Väter sind so heilig und vollkommen nicht gewesen, wie die Mönche von St. Augustinus und andern Heiligen träumen, von welchen sie dichten, daß sie ganz vollkommen gewesen seien. Und dies ist ein Exempel, das da zeuget, daß die heiligen Väter in der Erkenntniß immer gewachsen und zugenommen haben; wie denn alle gottseligen für und für auch wachsen müssen. Denn darum hat sie GOtt fallen und irren lassen, daß wir sollen wissen, daß an ihnen Erkenntniß, Glaube und Liebe unvollkommen gewesen ist, wie hernach ein sehr schwerer Fall Isaaks beschrieben wird.

173. Durch solche Exempel werden wir gelehrt, daß man keinem Menschen, wie heilig er auch sein möge, glauben soll, sondern man soll nur auf GOtt sehen, der uns beruft, und auf sein Wort. Daß man aber das Wort studirt und lernt, soll nicht allein ein oder zwei Jahre währen; denn es ist GOttes Wort ein unendlich Ding.

174. Die Widersacher schreien: Ambrosius, Augustinus haben also oder anders gelehrt. Aber die gehen uns nichts an, es sei denn, daß sie das Wort Christi bringen. Welcher aus ihnen das Wort und die Stimme GOttes klarer hat, demselbigen will ich lieber folgen, denn Augustinus oder sonst einem Andern, wer derselbe auch sein mag. Wie Augustinus etwas klarer und deutlicher lehrt und die Schrift besser auslegt denn Ambrosius, und ist derhalben auch nützlicher, man lese den Augustinus als den Ambrosius.

175. Aber ingemein sollen wir von Allen dies halten, wie der 119. Psalm V. 105. sagt: „Dein Wort ist meines Fußes Leuchte, und ein Licht auf meinem Wege“; und 2 Petri 1, 19.: „Wir haben ein festes prophetisches Wort, und ihr thut wohl, daß ihr darauf achtet, als auf ein Licht, das da scheinet in einem dunkeln Ort, bis der Tag anbreche und der Morgenstern aufgehe in euren Herzen.“ Denn man kann Alles so vollkommen und gänzlich nicht lernen, wie es an sich vollkommen und rein ist, dieweil uns noch immer anhängt, daß wir fleischlich gesinnet sind: dasselbe aber wird an uns ausgefegt durch Anfechtung, so wir leiden müssen, auf daß also das Licht in uns immer klarer möge werden.

\subsection{II}

\begin{quote}
\textbf{V. 26. 27. Sechzig Jahr alt war Isaac, da sie geboren wurden. Und da nun die Knaben groß wurden, ward Esau ein Jäger und ein Ackermann, Jakob aber ein frommer Mann, und blieb in den Hütten.}
\end{quote}

176. Darum erzählt Mose, wie alt Isaac gewesen, auf daß damit angezeigt und gerühmt werde seine sehr reine Keuschheit, darin er zwanzig Jahre mit Rebekka gelebt hat. Darnach wird damit auch angezeigt, wie sie die lange Zeit über auf die Kinder gehofft und gewartet haben, bis sie auch fast daran verzagt. Derhalben sind diese so heiligen Eheleute mit vielen Anfechtungen versucht und geübt worden. Sie sind zwar nicht so müßig und sicher gewesen, daß sie sich so gar unsinniger Weise auf die Unzucht oder fleischliche Lust wollten begeben haben, wie heutiges Tages die faulen Menschen in dieser letzten Zeit pflegen zu thun. Dergleichen sind sie auch Fremdlinge gewesen und haben nichts Eigenes gehabt, sondern sind hin und her gezogen mit großer Geduld und Demuth, wie droben von Abraham ist gesagt worden, welches Alles daher gehöret.

177. Bisher hat Abraham noch gelebet und hat fünfzehn Jahre lang mit seinen Kindeskindern gelebt, jedoch ist sein Leben viel kürzer gewesen als der andern Patriarchen. Ich zweifele aber nicht, es werden Kinder gewesen sein von einer guten Art und die sich wohl haben angelassen. Denn sie sind beide wohl erzogen und von Abraham selbst unterwiesen worden im Gehorsam, Zucht und gebührlicher Ehrerbietung nicht allein gegen ihre Eltern, sondern auch gegen die Patriarchen. Und vielleicht hat Abraham diese Verheißung auch ausgelegt und hat es mit den Andern dafür auch gehalten, daß Esau wäre der Größere geworden dadurch, daß ihm GOtt in Mutterleibe hätte den Sieg gegeben, da er der Kleinere war. Mehr hat er davon nicht verstanden oder gesehen. Denn er ist gestorben, ehe er Beide Sitten und Leben, dazu auch, wes sie sich befleißigen haben, gesehen hat.

178. Nun aber hebt die heilige Schrift an und eilt, den Unterschied dieser zwei Brüder zu beschreiben, und zu beweisen, welcher unter ihnen wahrhaftig der Größere und der Herr sei, auf daß diese Verheißung recht möge verstanden werden: „Ein Volk wird dem andern überlegen werden“; item: „Der Größere wird dem Kleineren dienen.“ Da sie aber erwachsen und groß geworden sind, hebt dieser Unterschied an, daß sie sich von einander theilen, und da wird Jeder erkannt an den Werken, deren er sich vornehmlich befliessen hat; wie das alte Sprüchwort lautet: Curvum se praebet, quod in uncum crescere debet, das ist: Was ein Häkchen werden will, krümmt sich bei Zeiten. Denn daran ist dazumal kein Zweifel gewesen, daß Esau nicht sollte der Erstgeborne sein, und daß die Verheißung müßte auf ihn gezogen werden, daß er der Größere wäre und das Regiment haben sollte, Jakob aber als der Kleinere würde dienen müssen.

179. Aber dies ist die Beschreibung der zwei Brüder, wie man sagt per antithesin, das ist, da einer mit dem andern verglichen wird, darin Esau also beschrieben wird, daß er sich durch seine eigene Schuld und Verdienst der Größe und der Herrschaft beraubt habe, ja, daß er gar bald der Kleinere soll werden. Dagegen aber soll Jakob die Herrschaft und die Größe bekommen; denn Esau ist ein Jäger und ein Ackermann. Ja, er ist der Herr, der rechte Poet, das Factotum, der da sicher und aufgeblasen ist von wegen der Erzfegheit, und fällt allemal ab von der Verheißung und vom Geist auf die äußerliche und fleischliche Ehre; denn er hört, daß er der Priester und der Herr sei: daraus erwächst bei ihm bald Verachtung seines Bruders und geistlicher Gaben, und folgt darauf die Hoffarth des Fleisches. Dagegen aber ist Jakob fein erzogen, demüthig, stille und einfältig, überhebt sich nicht über die Maaße mehr, denn sich gebührt, und sucht seine Ehre nicht in Summa, er ist der Sackträger und Schattenträger.

180. Und ist dies der Anfang des greulichen Falles, so Esau gethan, und ist ein recht Bild der ganzen Welt, welche also dahin fällt vom Geiste auf das Fleisch. Denn die Leute verachten, daß wenig, gering, kläglich und gefährlich, dazu auch der Strafe und vieler Mühe unterworfen ist, und suchen nur große Ehre, suchen Frieden und gute Tage, Reichthum und Freundschaft großer, gewaltiger Leute.

\quellseite{34}

181. Es hat aber Esau anfänglich seinen Fall nicht gefühlt, bis er des Segens beraubt wird; unterdeß ist er stolz und hoffärtig, fragt nichts nach seinem Vater, hält ihn auch nicht in Ehren, nimmt ein Weib aus dem Geschlecht der Cananiter, und ist allzu trotzig und verläßt sich darauf, als könnte er solcher großen Güter nicht beraubt werden. Welche Weise für und für alle Heuchler, und die, so sehr vermessen sind, pflegen zu haben; wie im 10. Psalm V. 6—11. steht: „Der Gottlose spricht in seinem Herzen: Ich werde nimmermehr darnieder liegen“ etc. „Gott hats vergessen, er hat sein Antlitz verborgen, er wird es nimmermehr sehen.“ Von den Gottlosen aber sagen sie: Es sind Ketzer, sie sind verflucht und verdammt; wir aber sind die Kinder des Reichs, wir sind die Kirche und Gottes Volk.

182. Also sind zu allen Zeiten alle Ketzer, alle Widersacher des Wortes ganz gewiß und sicher, daß ihr Ding Gott solle wohl gefallen. Ja, dieses Gebrechen breitet sich so gar weit aus, daß man es mit Worten nicht genugsam ausreden kann. Also schrieen im Anfange die Sacramentsschwärmer auch: Was wir lehren, ist die gewisseste Wahrheit. Denn diese alle werden aufgeblasen und stolz vom Glücke nach dem Fleische, und derer Dinge, welche das Fleisch vornehmlich sucht, als da ist, große Ehre und Majestät, Gewalt, und daß ihnen der Pöbel beifällt und sie hoch rühmt, dagegen aber ärgern sie sich an der Thorheit und Schwachheit des Geistes.

183. Darum hat sich Esau auch das Bisthum angemaßt und daß er in der Kirche oder Gemeinde Abrahams das oberste Regiment wollte haben; hat den Gehorsam gegen seine Mutter nicht viel geachtet, wie Jakob sich gehorsam erzeigt hat. Denn es ließ sich ansehen, als wäre es fast ein knechtisch Ding, der Mutter in solchen geringen und häuslichen Werken gehorsam zu sein, welche den Knechten gebühren auszurichten und auch dem Jakob. Ich, hat er gedacht, bin der Herr im Hause, Jakob muß Wasser holen; ich bin der Herr in beiden Ständen, Jakob ist mein Diener. Darum ist er ganz sicher, und erwählt einen solchen Stand, darin er leben will, welcher nicht mühselig oder schwer, sondern lustig und herrlich ist. Esau, sagt Mose, ist ein Jäger und ein Ackersmann, nicht daß er eigentlich ein rechter Ackersmann im Felde wäre, wiewohl er ohne Zweifel Aecker und andere liegende Gründe gehabt hat, sondern es ist eine Vergleichung der Arbeit, welcher sich die beiden Brüder befliessen haben. Esau hat sich im Felde geübt mit Jagen, mit Reiten und Kriegen. Was hat aber Jakob gethan? Antwort: Der war ein einfältiger, frommer und aufrichtiger Mann, und wohnte in den Hütten, in summa, er war ein Mann, der sich sein stille hielt.

184. Es ist aber unmöglich, daß Esau auch nicht in viel mehr andere Laster sollte gefallen sein. Denn ein Jäger, wenn er schon sich vor andern Sünden und Lastern hütet, so versündigt er sich doch oftmals mit Zorn und gottslästerlichem Fluchen, wenn es ihm nicht allenthalben nach seinem Willen geht; und wenn wir nur äußerlich nach dem Fleisch große Ehre haben, so stehen wir in gewisser Gefahr. Darum ermahnt der 62. Psalm V. 11. die Reichen und sagt also: „Fällt auch Reichthum zu, so hänget das Herz nicht dran“; und Hiob sagt auch im 31. Cap. V. 24. von sich selbst: „Hab ich das Gold zu meiner Zuversicht gestellt, und zu dem Goldklumpen gesagt: Mein Trost?“; das ist, ich war nicht stolz und hoffärtig von wegen meines vielen Geldes und Gutes, und verachtete auch die Andern nicht. Item daselbst sagt Hiob V. 17.: „Hab ich meinen Bissen allein gegessen, und nicht der Waise auch davon gegeben?“ etc.

185. Darum ist es gefährlich, wo man groß Geld und Gut hat, und wird im Evangelium der Reichthum Dornen und unrecht Gut genannt, und wo der Heilige Geist nicht dabei ist, so ist der reiche Mann nicht in der Zahl derer, die da Leid tragen, sondern geht einher in seinem Reichthum, und fällt gewißlich in Hoffart und Verachtung Gottes und der Menschen, Luc. 8, 7., und hält Fleisch für seinen Arm, Jer. 17, 5., hält das Gut für seinen Gott und betet es gleichsam an, und also wird der Reichthum ein rechter Abgott. Darum sagt auch Christus Matth. 6, 24.: „Ihr könnet nicht Gott dienen und dem Mammon.“ Und ist bei den Deutschen ein altes Sprüchwort, daß man sagt: Gut macht Muth. Dies geschieht aber nicht, als sollte der Fehler am Reichthum sein, sondern kommt von unserm sündlichen Fleisch her, welches so stolz und hoffärtig ist, und immer die Andern verachtet.

186. Darnach zum Andern, so treibt das Fleisch auch noch über diese Hoffahrt und Abgötterei einen Diebstahl und Räuberei, nämlich alsdann, wo es dem Exempel Hiobs nicht nachfolgt, daß es sein Gut den Armen nicht mittheilt, die Nackenden nicht bekleidet und andere Werke der Liebe auch nicht erzeigt. Denn das Fleisch hat das Geld und Gut so gar sehr lieb und verwahret es, daß auch die Reichen selbst ihres Reichthums nicht recht gebrauchen oder genießen.

187. Derhalben werden erstlich die Reichen abgöttisch wider die erste Tafel; zum Andern werden sie auch ungerecht und sündigen wider den Nächsten, werden Räuber und Diebe. Denn sie sollten den Armen geben von ihrem Gut und geben ihnen doch nichts. Es sind zwar etliche, die da etwas geben, aber unter der Bedingung, daß diejenigen, welche von ihnen etwas empfangen, ihnen dagegen auch sollen unterworfen sein, daß ist, sie geben häßlich und nur aus Hoffahrt, nicht mit einfältigen Augen oder Herzen, sondern nur daß sie dagegen wollen angebetet und hoch gerühmt werden.

188. Darum ist es sehr schwer, daß ein Reicher gerecht und gottselig sei, weil dem Reichthum das immer anhänget, daß Abgötterei und Verachtung des Nächsten daraus folgt. Ja, was soll man davon sagen, daß sie sich oftmals auch öffentlicher Sünde und Laster nicht enthalten, nämlich, daß sie anderer Leute Gut und Nahrung heimlich nachstellen, den armen Waisen Gewalt und Unrecht thun, und berauben nicht allein die Armen der Wohlthaten, so sie ihnen schuldig sind zu erzeigen, sondern reißen auch noch dasjenige, so sie übrig haben, mit Gewalt an sich? Daher kommen die Geizhälse und Wucherer; welche Laster auch in weltlichen Rechten sträflich sind.

189. Ich rede aber allein von den Reichen, die sich äußerlich des Stehlens enthalten, daß sie Andern das Ihre mit Gewalt nicht nehmen, und begehen doch damit einen Diebstahl, daß sie den Armen von ihrem Gut nichts geben, und erkennen nicht, daß ihr Reichthum eine Gabe Gottes sei, sind Gott dafür nicht dankbar. Dasselbe ist die Art und Eigenschaft aller Reichen, wo nicht der Heilige Geist dabei ist; wie er in Hiob, David und andern frommen, gottseeligen Königen gewesen ist, welche alle Geld und Gut wohl gebraucht haben. Aber David sagt im 62. Psalm V. 11.: „Fällt euch Reichthum zu, so hänget das Herz nicht daran.“

190. Es ist aber eine gemeine Gefahr, darin die Reichen stehen, nämlich, die Sünde, daß sie unterlassen oder versäumen, den Armen und Nothdürftigen mit ihrem Gut zu helfen, so doch Gott ernstlich geboten hat Jes. 58, 7.: „Brich dem Hungrigen dein Brod, und die, so im Elend sind, führe ins Haus“ etc. „Willst du das nicht thun, so wirst du das traurige Urtheil hören Luc. 12, 48.: „Welchem viel befohlen ist, von dem wird man viel fordern.“ Darum ist es sehr gefährlich, reich zu sein. Und Jes. 53, 9. werden die Reichen an der Gottlosen Statt gesetzt, darum daß zwischen einem Reichen und Gottlosen kein Unterschied ist, wo sie ohne den Heiligen Geist, und nicht gottselig und gläubig sind. Ein Gottloser raubt Gott seine Ehre, dankt Gott nicht für seine Gaben, und darnach weigert er sich auch, dem Armen Hülfe zu erzeigen. Dasselbe thut ein Reicher auch, welcher ohne Glauben ist; darum ist er auch gottlos.

191. So hat nun Esau immer eine heftige Begierde und Verlangen nach fleischlicher Ehre, nach großer Gewalt und Weisheit um der Ehrbarkeit willen. Er will sich an dem geringen knechtischen Stande nicht genügen lassen, wie sein Bruder Jakob thut. Er wählt sich den Jäger- und Reiterstand, darin er nicht ohne Sünde hat leben können, auch nicht ohne Begehungssünde, vornehmlich bei der Jagd.

192. Wie unsere Fürsten nicht allein damit sündigen, daß sie ihrem Amte nicht genug thun und sich der armen Unterthanen nicht annehmen, sondern sündigen auch ganz schwerlich, daß sie mit ihrem vielen unmäßigen Jagen die Armen beschweren, den armen Bauern und Ackerleuten die Früchte verderben, machen ihnen den Acker gar wüste. Und man darf keinerlei Weise das Wild aus den Gärten oder Aeckern wegtreiben, sondern es darf frei Schaden thun, und den Acker, so mit großer Mühe und Arbeit gebauet und besäet ist, verderben und die Früchte abfressen. Daselbst liegt nicht allein der Schutz darneben, daß sie den Unterthanen keine Hülfe thun, sondern man thut ihnen auch großen Schaden, welchen sie doch helfen sollten. Derhalben wird endlich der Türke oder ein anderer Jäger kommen, der den deutschen Fürsten beides, ihr Netz und Spieß, so sie in der Jagd gebrauchen, mit Gewalt aus der Hand nehmen wird.

\quellseite{35}

193. Dies sage ich darum, daß wir wissen sollen, daß die Jagd nicht ohne große Sünde sei, und ist von Vielen davon disputirt worden: Ob auch die Jagd ein ordentlich Werk sei? Es ist zwar an sich selbst nicht böse, und daß man damit umgeht, kann wohl christlich und gut sein; wie wir an dem Durchlauchtigen, unserm gnädigsten Herrn, Herzog Friedrich von Sachsen, seliger Gedächtniß, gesehen haben. Derselbe hat also gejagt, daß er damit niemand Schaden gethan, ja Vielen genützt hat. Wo er etwa vernommen, daß auch nur ein kleiner Schaden geschehen war, so hat er solches doppelt lassen bezahlen. Er hat auch oftmals den armen Bauern etliche Scheffel Getreides lassen austheilen, daß sie dafür haben sollten, wo ihnen das Wild etwas würde abfressen. Daselbst steht einem frommen, löblichen Fürsten wohl an. Denn wir wollen den Fürsten ihre fürstlichen Regalien und Gerechtigkeit, so sie haben, nicht nehmen; wie sich dessen die Bauern in dem Aufruhr im Jahre 1525 unterstanden haben.

194. Darnach soll dies die endliche Ursache sein, daß die Jagd dahin gerichtet werde, daß die greulichen und schädlichen wilden Thiere, als, Wölfe, Bären und wilde Schweine vertrieben mögen werden, und daß beide, Menschen und Vieh, von denselben befreiet und sicher sein mögen. Die Jagd soll also gethan sein, daß dadurch die armen Schäflein und andere unschuldige Thiere geschützt werden. Denn dasselbe ist ein Fürst von Amtes wegen schuldig zu thun. Welches auch der heidnische Poet gesagt hat: Parcere subjectis et debellare superbos, das ist: Er soll der armen Unterthanen verschonen, dieselben schützen und schirmen, und die Unruhigen und Stolzen, von denen die Unterthanen beschädigt werden, strafen und austrotten. Darnach ist die Lust und Ergötzung, so man an der Jagd hat, an sich selbst auch ehrlich, welche denen erlaubt ist, die derselben ohne Schaden und Nachtheil der Unterthanen gebrauchen. Wo es aber anders zugeht, so ist es mit der Jagd ein sehr böses und schädliches Ding.

195. Ich habe aber Sorge, dieweil die Vermuthung da ist, daß Esau ganz im Fleisch erwachsen gewesen, er werde seiner darin oft vergessen und viel Sünde begangen haben, also daß er beide, Schaden gethan, und auch vergessen habe, Andere zu schützen. Dieweil er aber der Erstgeborne ist und der Vater ihn sonderlich lieb gehabt, darum sündigt er nicht und thut seines Trachtens nichts Böses oder Unrechtes auch in den Dingen, so an sich selbst unrecht sind. Wiewohl wo er etwas gesündigt und unrecht gethan hat, wird er ohne Zweifel solches vor dem Vater verholen haben. Und dasselbe hat er zwar wohl thun können, dieweil er im Hause der oberste Fürst oder Regent war, daß er über das Hausgesinde regieret hat.

196. Nun wollen wir auch den andern Theil der Vergleichung besehen, da diese zwei Brüder werden gegen einander gehalten. „Jakob war ein frommer Mann, und blieb in den Hütten.“ Sie haben das hebräische Wort tham alle vermalmescht im Lateinischen, simplex, auf Deutsch, albern: so doch der Text im Hebräischen das nicht meint. Denn das Wort tham heißt nicht Einen, der albern ist von Verstand, das ist, an welchem gar keine Geschicklichkeit ist, sondern heißt Einen, der dem Willen nach einfältig, fromm und aufrichtig ist. Und eben dasselbe Wort steht im 119. Psalm, bald im Anfang: Beati immaculati etc., das ist: „Wohl denen, die ohne Wandel leben“ etc. Denn Jakob ist nicht albern oder närrisch gewesen von Verstand. Wie dasselbe hernach zu sehen ist, da von seinem Bruder die Erstgeburt gekauft hat.

197. Ein solcher frommer, aufrichtiger Mann ist Esau nicht gewesen, sondern ist dem Willen nach viel unbeständiger und tückischer gewesen; denn er heidnische Weiber zur Ehe genommen und sich also unter die Cananiter gemischt. Die Jäger sind nicht aufrichtig, sie richten oftmals ihr Ding, damit sie umgehen, nach dem die Zeit, die Sachen und Personen sind; Jakob aber ist fromm, aufrichtig und beständig, läßt sich keine Person bewegen, ändert sich auch in seinem Gemüth nicht, wenn die Sachen anders vorfallen, und ist doch gleichwohl nicht närrisch. Dies ist eine schöne Vergleichung der zwei Brüder.

198. Zum Andern ist Jakob kein Kriegsmann. Er befleißigte sich keiner weltlichen oder fleischlichen Händel. Wiewohl man das weltliche Regiment, und was dazu gehört, auch wohl christlich und recht verwalten kann. Jakob aber hat sich ganz und gar befliessen, daß er in den Hütten ist geblieben, das ist, er ist daheim bei Vater und Mutter geblieben und hat denselben gedienet. Und Lyra erzählt eine Meinung der Juden von den Hütten, welche mir sehr wohl gefällt, weil sie von den Vätern ist hergekommen. Denn sie sagen, daß die Hütten nicht allein sollen verstanden werden von denen, darin sie gewohnt, als in ihren Häusern, sondern daß damit auch die Hütten gemeint sind, darin sie ihre Kirchen gehabt, wie die Hütte Mosis genannt wird, dadurch auch eine heilige Hütte verstanden wird. Wir heißen es Kirchen und Schulen. Zu der Zeit aber war kein Unterschied zwischen den heiligen Hütten und den andern, darin sie sonst ihre Wohnung gehabt haben. Abraham hat in seiner Hütte ein Gotteshaus und Kirche gehabt; wie jetziger Zeit ein jeglicher frommer und gottseliger Hausvater seine Kinder und Gesinde in der Gottseligkeit unterrichtet. Darum ist ein solch Haus eine rechte Schule und Kirche, und ein Hausvater ist in seinem Hause Bischof, Priester und Pfarrherr. Also waren die Hütten, darin Jakob gewohnt hat, auch heilig, und er hat darin zum Ersten das Reich GOttes gesucht.

199. Dagegen aber hat Esau das Reich der Welt gesucht, darum hat er auch beiderlei Reich, das weltliche und geistliche, verloren. Denn das ist eine gemeine Regel, die für und für währet, nämlich, die am Ersten das Reich GOttes suchen, denselben wird das Andere auch zufallen. Wiederum aber, die am Ersten nach ihrem eigenen Nutzen und dem täglichen Brod trachten, ehe sie sich um das Reich GOttes annehmen, denen wird das Gegentheil begegnen, nämlich, daß sie das Reich GOttes mit dem äußerlichen Reich zugleich verlieren werden.

200. Darum ist Jakob eine Zeitlang kleiner und geringer gewesen denn sein Bruder. Aber sein Herz und Zuversicht hat gar an GOtt gehangen, und ist ein guter, frommer, gottseliger Jüngling gewesen, ist fleißig und gern in die Kirche und Schule gegangen, hat GOttes Wort und gute Sitten gelernt, hat die Predigt des Patriarchen Sem und der Andern gehört; wie die Hebräer anzeigen und solches von den Vätern genommen haben, nämlich, daß dieser Sem solle verstanden werden von den Hütten Sems und Eber, welcher achtzig Jahre mit Jakob gelebt hat.

201. Es sind aber sehr heilige und fromme Männer gewesen, die gar entbrannt gewesen sind von Liebe, so sie zu GOttes Wort gehabt, und haben viel größere Achtung gehabt auf die Religion und Gottesdienst, denn wir haben. Und wir mögen doch unsere Häuser und Kirchen auch wohl mit Recht Hütten nennen, wenn jemand eine Kirche oder Gemeinde in seinem Hause hat, und darin seine Kinder und Gesinde unterrichtet, lehrt sie die rechte Gottseligkeit und gute Tugenden. Mit je brünstigerem und fröhlicherem Geiste man aber das Wort handelt, je mehr Frucht bringt es. Darum lasse ich mir an diesem Orte der Juden Glosse und Auslegung wohl gefallen, und ist dies ein trefflich Exempel eines frommen und gottseligen Jünglings, der sich am Ersten vor allen Dingen in weltliche Händel mengt, befleißigt sich der Dinge nicht, damit sich das Fleisch vornehmlich pfleget zu bekümmern.

202. Er hat wohl das Weltliche nicht gar liegen lassen; aber er hat sich vornehmlich das lassen angelegen sein, wie er den Vätern Sem, Salah und Eber in ihrer Lehre anhangen und folgen möchte. Und aus der Jahrrechnung ist abzunehmen, daß Jakob mit Eber achtzig Jahre gelebt hat, mit Sem sechs und vierzig und mit Salah drei und dreißig Jahre. Zu denselben Vätern wird sich Jakob ohne Zweifel fleißig und mit gebührlicher Ehrerbietung gehalten haben, vornehmlich zu Eber. Denn da Sem alt geworden, hat er den Knaben Jakob dem Patriarchen Eber befohlen. Denn Sem ist der Vater gewesen aller Gemeinden, wie droben Cap. 10, V. 21. gesagt ist: „Sem, der ein Vater ist aller Kinder von Eber“; darum hat alles Volk und die ganze Gemeinde in seinem äußersten Alter auf ihn gesehen. Eber aber hat die Gemeinde versorgt und gelehrt; denn derselbe hat die Last der Regierung tragen können. Derhalben thue ich den Worten der Juden oder vielmehr der Väter gerne Beifall, da sie bei diesem Texte sagen, daß Jakob nicht allein daheim in der Hütte gewesen sei, sondern auch in fremden Hütten, und vornehmlich bei Eber und Sem.

\quellseite{36}

203. Jakob aber hat sich des Ehestandes enthalten beinahe in die siebenzig Jahre. Esau, als der Herr und gar ein Weltmensch, hat Canaaninnen zu Weibern genommen, da er vierzig Jahre gewesen ist. Diese Umstände alle zeigen an, daß Esau um der Erstgeburt willen über die Maßen stolz gewesen sei, daß er seinen Bruder verachtet, der Hütten auch gar nichts geachtet, und daß er sich nur ganz und gar des Reitens und der Jagd beflissen habe.

204. Dies ist aber gleichwohl nicht also zu verstehen, als ob Esau nie in den Hütten wäre gewesen, sondern er hat bisweilen dieselben auch besucht mit seinem Bruder und mit Andern mehr. Gleichwie die Heuchler mit Fleiß darauf sehen, und wollen traurig für heiliger und frömmer gehalten werden, denn die, so recht heilig, fromm und gottselig sind. Darum ist er auch ein Zuhörer gewesen in seines Vaters Kirche, und hat bisweilen aus freiem Antrieb die andern Kirchen oder gemeine Versammlung auch besucht, darin Eber und Salah gelehrt haben; aber heimlich und in der Stille hat er derselben nicht groß geachtet: wie bei uns der gemeine Haufe der Christen das Predigtamt und den Gottesdienst verachtet. Warum, sagen sie, sollte ich in die Kirche gehen? kann ich doch wohl daheim lesen; und stehen also allgemein an das Wort zu verachten, bis sie in die weltlichen Händel geraten, und gar weltlich und gottlos werden.

205. Aber fromme und gottselige Leute sollen immer daran mit Fleiß denken: Wer das Reich Gottes sucht, der soll beides finden; wer aber das Reich Gottes verachtet, derselbe verliert auch beiderlei. Darum sollen sie sich diesen Jakob vorstellen, daß sie demselben nachfolgen mögen, und daß sie ja Gottes Wort fleißig hören und lernen; an welchem Ort dasselbe gelehrt wird, daselbst sind die Hütten Gottes. Und ein Jeder soll es dafür gewißlich halten, wenn er zum allerersten das Wort hat, so werde ihm darnach auch Alles zufallen. Denn Gott war, ehe denn die Welt ist geschaffen worden, und hat Alles durch das Wort geschaffen; derhalben thut und bringt auch das Wort Alles.

206. So ist nun dies die Beschreibung der zwei Brüder. Die heilige Schrift lobt Jakob, und zieht ihn hervor in der Religion und Gottseligkeit, Esau aber in weltlicher Polizei; denn derselbe hat das Reich der Welt gesucht, welches billig an ihm gestraft wird. Und daselbst mag man sehen, wie er gefallen sei und seine Erstgeburt verloren habe, darum daß er so gar weltlich und gottlos gewesen; und wie Jakob dieselbe bekommen habe, dieweil er sich der Religion beflissen und gottselig gewesen ist. Daraus wir lernen sollen, daß wir am allerersten nach dem Reich Gottes trachten und daß wir dasselbe nicht nachsetzen sollen den Bemühungen und Begierden dieser Welt.

\subsection{III}

\begin{quote}
\textbf{V. 28. Und Isaak hatte Esau lieb, und aß gern von seinem Wildwerk, Rebekka aber hatte Jakob lieb.}
\end{quote}

207. Dieser Theil der Historie ist sehr dunkel, schwer und verborgen; denn Mose setzt allhier einen wunderlichen Unterschied zwischen diesen zwei Brüdern: der Vater hat Esau lieb und die Mutter Jakob. Und dasselbe kommt daher, daß sie alle beide nicht einerlei Natur gehabt, sich auch nicht einerlei Werks beflissen haben. Denn alle Mütter haben dies von Natur an sich, daß sie Lust und Gefallen haben an den Kindern, die von seiner stiller Art und guten züchtigen Sitten sind; denn die Weiber, welche fromm und ehrbar sind, dieselben sind auch von Natur furchtsam und sanftmüthig. Wie St. Petrus 1. Ep. 3, 6. die Weiber als die schwachen Werkzeuge stärkt, und ermahnt, daß sie nicht sollen furchtsam sein und vor nichts erschrecken. Und ist von Natur ein Unterschied zwischen der Art der Söhne und Töchter. Die Knaben haben eine etwas unbändige Art an sich; die Töchter aber sind lieblicher und halten sich freundlich zu den Eltern.

208. Darum hatte Rebekka den Jakob von Natur lieb; denn er hatte seine freundliche und stille Sitten; er war von guter, züchtiger und schamhaftiger Art, war dazu auch etwas schüchtern und furchtsam; denn er fürchtete seine Mutter mit aller Ehrerbietung, er hörte das Wort und die Predigten, hielt die heiligen Väter in Ehren. Welche Stücke alle an den Jünglingen den Müttern überaus angenehm und wohlgefällig sind und haben sie daran sonderliche Lust.

209. Dagegen aber hat sich Rebekka etwas geärgert an der groben Art und harten, unfreundlichen Sitten, so Esau an sich gehabt. Wie aber die Mütter die Söhne mehr lieben, so von stiller und freundlicher Art sind, denn die, so etwas unbändig und keck sind: also haben die Väter die Söhne lieb, so etwas hurtig und muthig sind, und dieselben werden deshalb angesehen, daß sie geschickt und hurtig werden wollen, etwas auszurichten. Und es gebührt sich auch gänzlich also, daß die Söhne hurtiger sein müssen als die Töchter. Ein solcher war nun Esau auch. Und Isaak hatte den Verstand, daß zum Regiment ein wackerer und frischer Muth gehöre, darum ist er dem Esau günstig gewesen.

210. Ueber das aber, so dieser Brüder Art und Natur gewesen, haben sie beide etwas an sich gehabt, darum sie sonderlich sind gelobt worden. Jakob war aufrichtig und ohne Falsch, war heilig und befließigte sich der Gottseligkeit, war gar entbrannt von großer Liebe und Verlangen nach dem Reich Gottes; dasselbe gefällt seiner Mutter Rebekka wohl, denn sie ist ein frommes gottseliges Weib. Darum wünscht und begehrt sie, daß Jakob möchte der Erstgeborne sein, und was das Vornehmste war, ist ihr Begehr, daß die Verheißung möchte erfüllt werden, daß der Größere dem Kleineren dienen möchte. Ueber das aber, daß Esau einen feinen scharfen Verstand gehabt, hat er auch eine sonderliche Gabe, daß er ein guter Schütze gewesen ist, wie droben (Cap. 21, V. 20.) von Ismael ist gesagt worden. Dasselbe gefällt nun dem Vater wohl, und lässet er sich dünken, es sei dem erstgebornen Sohne ein schöner Schmuck, als der das Regiment haben und verwalten sollte: dieweil auch das Gesetz der Natur, so allen Heiden bekannt ist, dem Erstgebornen die Herrschaft zueignet.

211. Jetzt ist es in weltlichen Rechten anders verordnet. Esau aber hat sich über das Gesetz der Natur auch auf die Gewalt der Väter und auf das göttliche Gesetz etwas verlassen, nach welchem der erstgeborne Sohn der Herr war, die andern Brüder aber waren Knechte. Dies ist der Unterschied der Dinge, welcher sich diese zwei Brüder beflissen, und auch der ungleichen Liebe an den Eltern gegen diese Zwillinge. Daß Jakob gottseelig war; hat ihm bei der Mutter Gunst gebracht, der Vater aber hat seine Ursache aus Gottes Gesetz, daß er dem Esau günstig und gewogen ist um der Erstgeburt willen. Wiewohl die Mutter nicht so sehr auf die Erstgeburt sieht, als darauf, daß er zur Linie der frommen Väter gehöret und ein Erbe der Verheißung sein möchte: da 1 Mos. 22, 18. gesagt wird: „Durch deinen Samen sollen alle Völker auf Erden gesegnet werden“; dies Erbe begehrt sie vornehmlich auf Jakob zu wenden.

212. In diesem Ort fällt aber nun eine schwere Frage ein, die ich droben ein wenig berührt habe. Die Verheißung sagt klärlich: „Der Größere wird dem Kleinern dienen.“ Nun weiß aber jedermann wohl, daß Esau der Größere, wie auch Jakob der Kleinere ist; wie das jedermann hat erkennen und sagen müssen. Wo nun Isaak die Verheißung gewußt, nach welcher Jakob der Größere sein soll, so hat er ja gar mit Vorsatz, wissentlich und boshaft wider die offenbare Verheißung Gottes gehandelt. Was wollen wir nun hier sagen? Denn wiewohl er Gottes Gesetz für sich hat, welches alle Erstgebornen zu Herren macht, so hat doch Gott Macht als ein Stifter des Gesetzes, das Gesetz aufzuheben, und davon auszunehmen, welche er will. Darum ist Isaak damit nicht entschuldiget, daß er das Gesetz für sich hat und sich darnach richtet, sondern er soll Jakob die Erstgeburt geben und nicht dem Esau, laut der Verheißung; wie er hernach darüber erschrickt, da Jakob dem Esau die Erstgeburt hat entzogen, als der sich der Sünde schuldig weiß.

213. Darnach fällt noch eine andere Frage ein, nämlich: Ob Isaak diese Verheißung gewußt oder nicht gewußt habe? Denn sie legen beinahe alle diesen Text also aus, daß es sich ansehen läßt, als seien sie der Meinung, daß Rebekka diese Verheißung bei sich behalten und vor ihrem Manne verhehlt habe. Denn darum eignet Isaak so beständig und von ganzem Herzen die Erstgeburt seinem Sohne Esau zu und nicht dem Jakob.

214. Mit diesen kann ich es aber nicht wohl halten; denn wir haben droben (V. 21.) gehört, wie Isaak für Rebekka gebeten und sie zu dem Patriarchen Sem gesandt habe. Darum ist es der Wahrheit nicht gemäß, daß Rebekka die Antwort, so sie von Sem gehört, sollte heimlich gehalten haben, vornehmlich dieweil sie dazumal nicht gewußt, welcher zuerst herkommen würde. Darum zweifle ich daran nicht, sie werde mit ihrem Manne von der Antwort, so sie vom Patriarchen Sem gehört, geredet haben, welcher gesagt hat, daß der Größere dem Kleineren dienen sollte. Da sie aber geboren sind, fährt Isaac zu, und wendet die Erstgeburt und den Segen, der seinem Sohne Jakob gebührt, auf den gottlosen Esau: dasselbe ist sehr kümmerlich und fast schwer, und steht noch die Frage da; und ich weiß nicht, wie wir dieselbe sollen auflösen. Ich wollte aber gern dieser Meinung folgen, daß Isaac, nachdem er die Antwort, wie sich die zwei Kindlein im Mutterleibe stoßen würden, gehört, gedacht habe, daß dies also zu verstehen wäre, daß der zuerst würde herauskommen, in Mutterleibe den Größern überwunden hätte, dieweil er der Kleinere war, und daß also der Streit nun sein Ende hätte und die Verheißung damit erfüllt wäre. Jakob ist in Mutterleibe der Größere gewesen; aber die göttliche Verheißung hat diese Ordnung geändert, also daß Jakob, der Zertreter, der Kleinere und der Letzte in der Geburt geworden ist.

\quellseite{37}

215. Dies ist mir also eingefallen, auf diese Frage zu antworten; wiewohl ich nichts Gewisses dazu sagen darf und will es auch nicht thun. Es läßt sich aber ansehen, daß es sehr glaublich und der Wahrheit fast gemäß sei. Und mich beweget das Exempel Ismaels, von welchem Abraham auch gemeint, daß er würde der Erbe sein, da er noch von Isaac nichts gewußt; er hat an dem Spruch gehangen, der droben im 16. Cap. 3. 10. steht: „Ich will deinen Samen also mehren, daß er von großer Menge nicht soll gezählet werden“; daselbst hat er die Verheißung nicht anders verstanden, denn daß Ismael würde der Erbe sein. Denn die Verheißung ist gar insgemein geredet, und man kann daraus keinen klaren und gewissen Verstand nehmen, bis GOtt selbst sein Wort hat ausgelegt. Die Verheißung, sagt er, soll nicht von Ismael verstanden werden, wie du meinst, sondern von Sarah und ihrem Sohn.

216. Es läßt sich also ansehen, als sei Isaac ganz und gar der Meinung gewesen, daß die Erstgeburt seinem Sohne Esau gebühre, und hat ihn nicht angesehen oder bewogen, daß Esau von Art und auch seiner Sitten halben fast unbändig und wild war, welche Unbändigkeit und wildes Wesen die Mutter an ihm pflegte zu strafen: sondern hat gemeint, daß dasselbe dem Erstgebornen als dem künftigen Herrn und Regenten wohl anstehe. Und wo etwa ein Fehler an ihm wäre, hat er die Hoffnung, wenn er würde älter werden und sich bemühen, sein Amt fleißig auszurichten, so würde solcher Fehler an ihm gebessert werden und dahin fallen, und er würde sich daran gewöhnen können, daß er etwas sanftmüthiger werden möchte.

217. Dies lasse ich mir zwar gefallen, dieser That Isaacs also nachzudenken oder dies davon zu erachten; ist aber kein Herz und Meinung anders gewesen, und daß er mit Wissen und Willen den Segen auf Esau gewendet hat, so ist er nicht wohl zu entschuldigen. Wie man Rebekka auch nicht wohl entschuldigen kann, wo sie die Antwort, so sie vom Patriarchen Sem gehört, vor ihrem Mann verhehlet hat. Darum muß man auf der Mittelstraße bleiben.

218. Dieweil sie aber beiderseits so fest auf ihrer Meinung bestehen bleiben (denn die Meinung, so der Mann gefasset, ist gar dahin gerichtet, daß er es also haben will, aber der Mutter Meinung ist, daß sie wünscht und begehrt, daß es also gehen möchte, wie sie es gern haben wollte, und dieselbe Meinung ist wahrhaftiger), darum disputiren sie noch weiter, daß Rebekka durch ein sonderliches und göttliches Eingeben diesen Verstand gehabt habe, daß der Kleinere der Größere würde werden. Hier gehe ich auch zu, daß es also sei. Denn also ist der Sarah auch von heiligen Geist eingegeben worden, daß der Sohn Ismael sollte ausgetrieben werden. Denn es ist ja an keinem Orte geschrieben oder öffentlich geboten, daß er sollte ausgetrieben werden, sondern Sarah hat solches aus eigener Meinung und Geist also beschlossen, weil sie an Ismael viel findet, das alles der Erstgeburt stracks zuwider war, und siehet, daß er den verfolgt, welcher nach der Verheißung geboren war. Also wird Rebekka auch ein sonderliches Eingeben gehabt haben, welches Isaac nicht gehabt.

219. Man soll aber nicht leichtlich und ohne gewisse Ursache ein jegliches Eingeben annehmen; denn es ist damit gefährlich. Der Papst sagt, daß Eingeben der Natur und des Teufels solle bisweilen dem Eingeben der Gnade gleich sein, und setzt des zum Exempel die Historie von den Versuchungen Christi, Matth. 4, V. 1. ff. Aber diese Meinung ist besser, daß sich der Teufel in einen Engel des Lichts verstellen und in die göttliche Majestät verändern kann, 2 Cor. 11, 14. Darum muß zu solchem Eingeben GOttes Wort kommen; wie der Sarah Eingeben daher entstanden ist, daß sie weiß, daß sie eine Verheißung hat, und daß sie nach derselben Verheißung die Mutter ist des rechten Samens, welche Verheißung Ismael nicht gehabt. Darnach urtheilt sie von dem Baum aus den Früchten, nämlich, dieweil Ismael ihren Sohn und den rechten Erben verfolgt, sucht nur, was fleischlich ist, und lehrt Abgötterei. Dies Eingeben der Sarah ist gewiß.

220. Auf solche Weise hat es wohl sein mögen, daß Rebekka von den Vätern mag gehört haben: Dein Sohn Esau ist ein wilder, ungezähmter und grausamer Mensch, darum wird er den Segen nicht erben; er geht ja in den Tempel oder Hütte, er läßt sich aber die Gottseligkeit nicht ernstlich zu Herzen gehen; aber Jakob ist fromm und gottselig, darum wird es geschehen, daß er wird der Größere werden. Dazu kommt auch das Exempel und die Erfahrung an sich selbst, dieweil Esau gottlos und immer ärger wird; dagegen trachtet Jakob nach dem Himmelreich, wächst und nimmt immer zu in Gottseligkeit und in andern Tugenden. Dasselbe Eingeben ist auch gut und gewiß.

221. Also muß man sich nun immer befleißigen, daß wir ein gewisses Wort haben, oder, wo ja das Gebot GOttes nicht allezeit vorgebracht, so sieht man doch nach der That, was es für ein Eingeben gewesen ist. Wie David den Löwen, den Bären und Goliath, den Philister, erwürgt; aber da ist der Geist auf ihn gefallen, das ist das Eingeben gewesen, so man erst nach der That gesehen hat. Und was solche Eingebungen sind, die haben zum Grunde GOttes Wort, oder wo sie anfänglich ohne Wort sind, so werden sie doch hernach erkannt und dafür gehalten, daß sie von GOtt sind hergekommen. Wie es uns auch ist gegangen. Wir haben den Papst ohne Wort angegriffen; jetzt aber sehen wir, daß es ein göttlich Eingeben gewesen, ohne unsere Gedanken und Rath. Über davon ein ander mal mehr.

222. Jetzt rathschlagt Rebekka weiter, wie Jakob die Erstgeburt bekommen möge. Esau ist gar sicher und denkt, er sei nun der Gefahr gar entronnen, sintemal er der Erstgeborne und der Herr ist; es folgt aber nun eine Gelegenheit, dadurch er die Erstgeburt gar verliert.

\section[Fünftes Stück]{Fünftes Stück. Wie Esau seine Erstgeburt verkauft, und wie Jakob dieselbe an sich gebracht}

\begin{quote}
\textbf{V. 29. 30. Und Jakob kochte ein Gerichte. Da kam Esau vom Felde, und war müde, und sprach zu Jakob: Laß mich kosten das rothe Gericht, denn ich bin müde. Daher heißt er Edom.}
\end{quote}

223. Mose hat bisher beschrieben, wie an den zwei Brüdern ungleiche Werke gespürt worden, deren sie sich beflissen haben, nämlich, daß Esau ein Ackermann und ein Jäger gewesen, das ist, daß er nach der Gottseligkeit, nach der Predigt und Lehre der Väter, so dazumal noch gelebet haben, wenig gefragt hat. Es ist ihm aber gegangen, wie es Allen pfleget zu gehen, wenn sie das Kleine verachten, daß sie darüber auch allgemach das Große pflegen zu verlieren. Denn er hat das für ein klein gering Ding gehalten, daß er in Sachen, die Gottseligkeit belangend, etwas nachlässiger und saumseliger gewesen denn Jakob; daß er zu göttlichen Dingen, zur Erkenntniß Gottes und zum Gottesdienst nicht so brünstig gewesen. Aber darüber ist er allgemach kälter und nachlässiger geworden, und ist dahin gerathen, daß er auch die Erstgeburt verachtet.

224. Darnach ist er in solche Sicherheit und Stolz gefallen, daß er meint, er möchte ungestraft seinen Bruder damit wohl betrügen, als wäre es sein Ernst nicht gewesen, daß er ihm die Erstgeburt verkauft hätte. Denn er gedachte, es könnte oder möchte nicht geschehen, daß die Erstgeburt auf seinen Bruder kommen sollte, sondern dieselbe gebührete ihm, und hoffte auch, er werde für immer den Segen behalten. Er spottet also seines Bruders, und wird also gar ein gottloser Mensch gegen GOtt, gegen die Menschen und seinen eigenen Bruder.

\quellseite{38}

225. Ein solch gottloses Wesen und Sicherheit pflegt daraus zu folgen: wo man GOttes Wort verachtet und sich darin nicht fleißig übt, da werden die Leute Atheisten, Epicurer und gar thöricht, daß sie weder Sinn noch Verstand mehr haben. Dasselbe bezeugen die Exempel auch der größten Leute. David war ein sehr heiliger Mann, war sehr brünstig im Gottesdienst, aber siehe, wie leichtlich und schnell er in Ehebruch, in Todtschlag und Gotteslästerung gefallen. Denn wo man in GOttes Wort schläfrig und unfleißig, dazu auch sicher ist, das ist eben so viel, als wo man dem Teufel Thür und Fenster aufthut. Darum ist uns geboten, daß wir wachen sollen, wie geschrieben steht 1 Petr. 5, 8.: „Seid nüchtern und wachet, denn euer Widersacher, der Teufel, gehet umher wie ein brüllender Löwe, und suchet, welchen er verschlinge“; item Eph. 5, 15. 16.: „So sehet nun zu, wie ihr vorsichtiglich wandelt, nicht als die Unweisen, sondern als die Weisen, und schicket euch in die Zeit“. Aus Sicherheit und daß man in göttlichen Sachen nachlässig und unfleißig ist, sind oftmals große schwere Fälle und greuliche Sünden gekommen.

226. Ja, es ist auch in weltlichen Sachen daran viel gelegen, daß man wacker sei, und ob man die Gelegenheit wahrnimmt oder versäumt. Wie die Türken darauf beflissen sind, daß sie sich immer nach Gelegenheit in die Zeit wissen zu schicken. Darum haben sie in kurzer Zeit so manchen Sieg wider uns erhalten. Unser Kaiser hat im Jahr 1532 auch gute Gelegenheit gehabt, die Sachen wider den Türken herrlich anzugreifen, nämlich, vor Wien in Oestreich. Aber da er dieselbe versäumt hat, suchen wir jetzt vergeblich andere Gelegenheit.

227. Derhalben sollen wir unsere Herzen daran gewöhnen, daß wir die erste Tafel von ganzem Herzen lieb gewinnen. Wo man soll auf das Feld gehen, oder sonst ein ander Amt im Haus- oder Weltregiment annehmen zu verwalten, da soll das die allererste und vornehmste Sorge sein, daß man sich mit Ernst um das Reich GOttes bekümmere; wo man aber das versäumt, so sehen wir wohl, wie jämmerlich Esau dahin fällt, daß er auch nimmermehr wiederum aufsteht. Dagegen aber wächst Jakob und wird der Erstgeborne durch sehr seine Gelegenheit. Esau hat den Namen und den Ruhm, daß er der Erstgeborne sei: aber unterwegs fällt er bald und greulich dahin und wird aller Ehre und Herrlichkeit beraubt. Darum gedenke, daß du nicht schläfrig seiest, sondern auf allerlei Gelegenheit gute Achtung gebest, auf daß du durch Nachlässigkeit deine Gaben und das Reich GOttes nicht verlierest. Denn darum hat uns die Schrift dieses Exempel vorgehalten, nicht als eine kalte und todte vergebliche Historie, sondern daß sie uns lehren und erinnern soll, wie es jetzt und zu allen Zeiten um alle Sachen und Personen bestellt sei.

228. Der heilige Geist will haben, daß wir diese zwei Brüder stets sollen vor Augen haben und ein tägliches Sprüchwort von ihnen machen. Denn zu allen Zeiten sind wir entweder Esauiten oder Jakobiten. Esau ist dem äußerlichen Ansehen nach vor den Andern allen heilig und geistlich, und läßt sich ansehen, daß er ein Erbe sei des Himmelreichs, und dasselbe schon inne habe und besitze: und ist doch nichts mehr denn ein rechter Esau, das ist, ein Werktheiliger und Heuchler, der sich die Religion nicht mit Ernst läßt angelegen sein, ist nur allein äußerlich fromm um des Genusses und Bauches willen und kann man ihm den Namen oder Titel der Kirche nicht nehmen. Wie wir vor der Welt den Namen der Kirche nicht haben, sondern werden für Ketzer gehalten, und für die, so hernach geboren sind. Die Päpstlichen aber wissen sich zu rühmen, sie seien die Erstgeborenen. Und wenn wir auch schon den Pabst mit seinem Anhange überwinden, und mit dieser unserer Lehre, welche eigentlich die Lehre ist der rechten Kirche, siegen würden, so würden doch gleichwohl unter uns aufstehen Papisten oder Türken, die sich den Titel und Namen der Kirche würden anmaßen, wie bisher die Wiedertäufer und Sacrament-Schwärmer gewesen sind; und kann auch nicht anders zugehen. Also, da die Leiden, so die Kirche verfolgt, sind überwunden worden, sind bald die Arianer aufgestanden, die den Namen der Kirche mit Gewalt haben an sich gerissen, und ist die rechte Kirche gar verachtet gewesen. Also kommen allezeit Kain oder Abel, Esau oder Jakob, und zwar von einerlei Stamme und auch von einerlei heilsamen Evangelium.

229. So ist nun Esau der Bischof, das Factotum, und ist die Kirche und gebührt ihm doch davon nichts. Denn es wird der ganze Segen aus göttlichem Befehl auf Jakob gewendet, welcher der rechte Erbe ist aller Ehren und des Himmelreichs; bisher aber ist er des Namens und der Herrlichkeit, so ihm gebührt, beraubt gewesen, bis ihm sein Bruder Esau selbst wunderbare Gelegenheit gibt, diese Herrlichkeit zu erlangen; welche Gelegenheit alle Hoffnung und Gedanken Jakobs übertrifft. Und in Summa: Es hebt sich schlecht über dem Brei. Jakob gedenkt nicht, daß er seinen Bruder wollte betrügen. Aber es geschieht alles durch göttliche Ordnung und Rath, da keiner, weder Jakob noch Esau, daran gedacht oder des Willens gewesen.

230. Jakob thut gemeine gewöhnliche Werke, er ist mit seiner Mutter Rebekka in der Hütte und kocht einen Brei; diese Einfalt wendet GOtt in eine wunderbare Gelegenheit. Denn Esau kommt wieder vom Felde und ist müde, daß er sich auf der Jagd so bemüht und so emsig gewesen war. Derselbe Fleiß gefiel seinem Vater Isaac wohl; und läßt sich ansehen, als habe er nichts mitgebracht zu essen, dieweil er nicht allein müde, sondern auch hungrig und durstig war. Und darum bittet er mit solcher Begierde, daß ihn Jakob wollte das rothe Gericht kosten lassen. Denn er hat es zweimal wiederholt: „Laß mich“, sagte er, „das rothe Gericht, ja, eben das rothe Gericht kosten.“ Damit wird angezeigt, wie er nach der Speise so große Begierde und Verlangen gehabt. Als wenn ein Hungriger sagen würde: Laß mich mit dir essen vom Karpfen, und eben von dem Karpfen, daß mich gelüstet. Mose hat an ihm wollen anzeigen, daß er nicht allein hungrig und müde gewesen, sondern auch große Lust und Begierde zu essen gehabt, da er das rothe Gericht mit sonderlicher Lust angeschaut und desselben begehrt hat.

231. Ich will aber diesen Text den hebräischen Grammatikern befohlen haben, daß sie denselben etwas fleißiger auslegen. Denn es ist unmöglich, daß man eine Sprache also verdolmetschen sollte, daß man an allen Orten den Nachdruck, dazu auch die Figuren in allen Worten und Sätzen behalten könnte.

232. Das hebräische Wort, welches Esau gebraucht, da er Jakob anredet, halitheni, steht sonst nirgends denn an diesem Orte. Was es aber bedeute, wissen die Hebräer selbst nicht und ich kann es auch nicht wissen. Die Conjugation und Construction, wie man es in der Grammatik nennt, zeigt etlichermaßen die Bedeutung an. Er will sagen: Ich bitte dich, speise mich, weide mich, gib mir zu essen. Die heilige Schrift hat nicht ohne Ursache so eigentlich reden wollen; sie hätte wohl ein ander Wort gebrauchen können, das Speisen oder Nähren, oder desgleichen etwas bedeutet, aber sie hat ein sonderlich Wort gebrauchen wollen. Rabbi Salomo dichtet, daß Esau so müde gewesen sei, daß er die Hände nicht habe können zum Munde bringen und selber die Speise hineinstecken. Ich weiß aber nicht, ob es sich auch mit dem, so hernach folgt, recht reimen wolle, dieweil der Text sagt: „Er hat gegessen, getrunken, und ist aufgestanden und davon gegangen.“ Dies ist aber doch die Meinung und der rechte Verstand, daß Esau hat haben wollen, daß Jakob ihm das Gericht geben sollte, das er für sich selbst gekocht, und hat zu ihm gesagt: Speise mich doch oder erquicke mich mit diesem rothen Gerichte.

233. Daß aber im Texte noch hinzu gesetzt wird: „Daher heißt er Edom“, ist mit nichten vergeblich oder umsonst also hinzu gesetzt, wiewohl es sich ansehen läßt, daß die Ursache dieses Namens lächerlich und ungereimt genug sei, nämlich, daß er den Namen haben soll vom Linsenmus oder vom rothen Gerichte. Es wird aber ohne Zweifel etwas Größeres darunter verborgen sein. Denn droben (V. 25.), da er ganz und gar röthlich geboren worden, hat man ihn um der röthlichen Haut willen nicht Edom genannt, welches ja eine wichtigere und rechtmäßigere Ursache wäre desselbigen Zunamens; aber daselbst heißt sein Name Esau: jetzt, da er nun vom rothen Gerichte ist satt geworden, wird er Edom genannt.

234. Er hat aber drei Namen. Esau ist sein rechter Name, den er, da er ist geboren worden oder aber in der Beschneidung empfangen hat. Der Zuname ist Edom von dem rothen Gerichte. Der dritte Name heißt Seir. Droben werden die zwei Namen zusammen gesetzt, Edom und Seir, und er wird doch Esau genannt. Er hat aber den Namen Seir bekommen entweder von dem Berge, darauf er gewohnt hat, welcher zuvor Seir ist genannt worden, oder von der Rauhheit, darum daß er röthlich und rauh war, und ist vielleicht auch das Land, so er inne gehabt, hart und rauh gewesen. Daher werden die Jdumäer das Volk Seir genannt eben von Esau, der den Berg Seir inne gehabt, oder daß im selbigen Lande gewohnt haben die seirim, das ist, die fauni und satyri pilosi, welche wir die wilden rauhen Männer nennen. Ob aber in dem Lande solche wilde rauhe Männer gewohnt haben oder nicht, kann ich nicht wissen. Es werden sonst in der Bibel die seirim verstanden von den Teufeln, so in der Gestalt der wilden Männer pflegen zu erscheinen.

\quellseite{39}

\begin{quote}
\textbf{V. 31—34. Aber Jakob sprach: Verkaufe mir heute deine Erstgeburt. Esau antwortete: Siehe, ich muß doch sterben, was soll mir dann die Erstgeburt? Jakob sprach: So schwöre mir heute. Und er schwur ihm, und verkaufte also Jakob seine Erstgeburt. Da gab ihm Jakob Brod und das Linsengericht, und er aß und trank, und stund auf, und ging davon. Also verachtete Esau seine Erstgeburt.}
\end{quote}

235. Alhier fällt eine große Frage ein vom Kaufen und Verkaufen der Erstgeburt; denn es ist ein geistlich Gut, das man weder kaufen noch verkaufen muß. Derhalben wird gefragt: Ob beide gesündigt haben, der Eine mit Verkaufen, der Andere mit Kaufen. Lyra sagt, Esau habe gesündigt, daß er die Erstgeburt verkauft; von Jakob aber sagt er, daß er damit nicht gesündigt habe, daß er sie gekauft hat, darum daß er nach göttlicher Ordnung der Erstgeburne gewesen sei; wie der Spruch lautet: „Der Größere wird dem Kleinern dienen.“

236. Esau hat die Erstgeburt inne gehabt, die ihm doch nicht gebührt; denn wiewohl er sich dieselbe von wegen der Geburt anmaßte, so hat er sie doch nach GOttes Willen nicht inne gehabt, sondern wie man pfleget zu sagen, welcher ein Gut besitzt und inne hat, der hat es damit so viel desto besser. Darum hat ihn Jakob nicht in seinem Besitz stören können, und Esau hat sich auch geschützt mit weltlichem und göttlichem Recht, und dazu auch mit der Geburt; welche Stücke alle dem erstgebornen Zwillinge den Besitz und das Eigenthum zueignen bei allen Völkern. Ueber das gibt ihm das Gesetz auch zwei Theile an des Vaters Gütern. Derhalben wo GOtt dies Recht nicht hätte geändert und widerrufen, so wäre dem Esau die Erstgeburt unverrückt geblieben. Dieweil aber Jakob aus der göttlichen Stimme weiß, daß sie ihm gebührt, so sündigt er daran nicht, daß er sein Recht in Anspruch nimmt, sondern er hat daran wohl gethan, daß er auf allerlei Gelegenheit aufmerksam hat, auf daß er die Erstgeburt erlangen möchte, so ihm von Rechts wegen gebührt hat.

237. Dies hat aber auf keine andere und bessere Weise geschehen können, denn daß sie ihm sein Bruder verkauft hat. Denn wiewohl das Gericht, so er gehabt hatte, oder die Suppe, nicht so viel werth ist, die Erstgeburt damit zu bezahlen, so dient es doch dazu, daß er derhalben nicht hat können vor Gericht beschuldigt oder geplagt dürfen werden. Denn das geschieht oftmals, daß Einer unbilliger Weise von den Sophokanten und Verleumdern geplagt wird, die ihn vor Gericht verklagen, und es versuchen, entweder durch falsche Zeugen, die sie herbeibringen, oder aber durch andere listige Tücke, so sie erfunden und erdacht haben, ihm sein Gut und Habe zu nehmen, unter dem Schein, als habe er dasselbe nicht mit Recht besessen und inne gehabt. Die nun daselbst gerne wollen Frieden haben, wiewohl sie ihr Recht vertheidigen und erhalten könnten, geben doch solchen Sophokanten, die sie so fälschlich verklagt haben, viel lieber noch Geld zu, denn daß sie mit ihnen lange wollten vor Gericht hadern und zanken, auf daß sie also ihrer falschen Anklage los werden. Daselbst gilt das gemeine Sprüchwort, daß man sagt: Man muß dem Teufel zwei Kerzen aufstecken. Auf solche Weise hat Jakob das rothe Gericht gegeben, und sich damit errettet, daß ihn sein Bruder vor Gericht nicht hat plagen dürfen, welcher sich der Erstgeburt rühmt, auf daß er hernach sein Gut und sein gebührlich Recht mit gutem Frieden besitze. Dies ist es etwa, so Lyra an diesem Ort wohl und göttlich genug geredet in Anbetracht der Zeit, darin er gelebt hat.

238. Ich habe aber droben in der Historie Abrahams von der Simonie gesagt, welcher einen Acker zum Begräbniß gekauft hat, 1 Mos. 23. V. 16., nämlich, ob Abraham derhalben soll ein Simonist genannt werden, weil das Begräbniß auch mit unter die geistlichen Güter gerechnet wird, und wo man solche geistliche Güter kauft oder verkauft, solches Simonie genannt wird? Wir haben aber gehört, daß derselbe gekaufte Acker ein weltlich Gut gewesen, und das seinem leiblichen Herrn unterworfen war. Und wo ein Acker verkauft wird, ist solches gar ein bürgerlicher Handel, wie ein Bürger mit dem andern kauft und verkauft. Derhalben hat Abraham nicht gesündigt, da er den Acker gekauft und daselbst ein Begräbniß gemacht hat; denn sonst würde es alles eitel Simonie sein, was die Christen zu ihrer täglichen Nothdurft kaufen, als Brod, Fleisch, Wein, welches alles die Heiligen oder Gläubigen zu ihrem Gebrauch verwenden; welches gar ungereimt Ding ist, wenn man davon also urtheilen wollte.

239. Die Rechte und Canones nennen aber das Simonie, wo man geistliche Güter um Geld kauft. Also, wenn der Gottesacker oder Begräbniß verkauft wird, sagen sie, es sei Simonie; denn es ist bei den Christen ein geistlich Gut. Item, wo Einer ein Theil etwa von einer Pfarre verkauft, oder liegende Güter von den geistlichen Pfründen oder Lehen, welche sie auch geistliche Kirchengüter heißen.

240. Das war die große Sünde, darauf vor dieser Zeit die Theologen und Canonisten einmüthig so heftig gescholten haben, und verdammten die Bischöfe und römischen Curtisanen (Höflinge), so die Bisthümer kauften und verkauften. Aber sie haben mit ihrem großen Geschrei nichts können ausrichten und werden jetzt vom ganzen Papstthum verspottet; welches Papstthum mit seiner Substanz und Wesen voller Simonie ist, und wenn du aus dem Papstthum die Simonie wirst wegnehmen, so hast du schon die Sonne heraus genommen. Denn es haben die Curtisanen sonst Anderes nichts zu thun, denn daß sie geistliche Lehen kaufen und verkaufen.

241. Derhalben, wo man die römische Kirche von der Simonie reformiren wollte, das würde Anderes nichts sein, denn dieselbe in den Grund zerstören und gar vertilgen. Wo wird die Herrlichkeit und große Pracht der Cardinäle am Hofe zu Rom bleiben, wenn du die Simonie daselbst wirst wegnehmen? Ja, alsdann wird der ganze Hof zu Rom gar zu Boden gehen und dahin fallen. Denn es sind der Cardinäle gar zu viel, die jetzt alle geistlichen Kirchengüter getroffen und verzehrt haben.

242. Es hat Johann Huß und Andere mehr, so zu der Zeit das Lehramt geführt, dies Laster fast hart gestraft, aber es ist alles vergeblich und umsonst gewesen. Denn es ist um den Hof zu Rom ein verdammt Ding, vor welchem jedermann sich hüten und denselben verfluchen soll; was sollte er denn Gutes können lehren vom Glauben und guten Werken? Das Papstthum ist eine Versammlung der Teufel und der allerschändlichsten Menschen. Sie verdammen die Simonie und leben doch selbst davon, haben daher ihren Unterhalt. Darum vermähnt uns St. Johannes in seiner Offenbarung am 18. Cap. V. 4. 5. recht, da er von Babylon sagt: „Gehet aus von ihr, mein Volk, daß ihr nicht theilhaftig werdet ihrer Sünden, auf daß ihr nicht empfahet etwas von ihren Plagen; denn ihre Sünden reichen bis in den Himmel, und GOtt denkt an ihren Frevel“ etc. Es sind alle des Teufels Stifte oder Kirchen, darum sollen wir davor fliehen, daß wir ihrer Plagen nicht theilhaftig werden.

243. Wir könnten aber diesen ihren Diebstahl wohl entschuldigen mit einer Glosse, daß es nicht Simonie genannt werde; denn es ist ein weltlicher Raub, und ist kein Unterschied zwischen solcher Simonie und wo man sonst weltlicher Weise ein Gut mit Unrecht einnimmt: nicht anders, denn wie wir jetzt zu unserer Zeit sehen, daß die vom Adel mit einander zanken, und die liegenden Güter, Aecker, Fuhrwerke und Dörfer einander einnehmen, zu dem so sind unsere Lehen und Besoldungen, deren wir gebrauchen, auch durch Simonie erworben. Sie können von Hurenlohn gesammelt, wie im Propheten Micha am 1. Cap. V. 7. steht. Wenn sie aber die Kirchengüter an sich reißen, wenden sie dieselben nicht dazu an, daß davon fromme gottseelige Priester oder Kirchendiener unterhalten werden, sondern wenden dieselbigen nur zum Ueberfluß auf ihre Pracht, und daß sie davon im Saufe leben. Wir wenden sie ja zum rechten Gebrauch, nämlich, daß davon die Kirchendiener genähret, und die Schulen bestellt und erhalten werden. Derhalben achten wir der Simonie nicht, sondern dieweil sie sich mit ihrem geistlichen Recht selbst beschweren, so sagen wir zu ihnen: „Aus deinem Munde richte ich dich“ 2c., Luc. 19, 22.

\quellseite{40}

244. Das ist aber die rechte Simonie, wo man geistliche Dinge für Geld kauft und verkauft. Als, wenn Einer wollte Gottes Gnade, den Glauben, Liebe, Wunderwerke, die Gaben und Kraft des Heiligen Geistes, damit die Kirche und Gemeinde Gottes gezieret wird, kaufen und damit Hantierung treiben, wie sich das Simon unterstanden hat in den Geschichten der Apostel am 8. Cap. V. 19.; nicht, daß solch geistlich Ding an sich selbst Einem könnte gegeben werden, wo man solch Heiligthum kauft oder verkauft. Denn der Heilige Geist gibt seine Gaben nicht für Geld, und Gott nimmt auch weder Gold noch Silber für die Vergebung der Sünden, für das Leben, für die Augen, Sinne und Vernunft, und für alle anderen Wohlthaten, so er uns verliehen hat: sondern nach seiner göttlichen Güte streut er und theilt seine Gaben aus beide Dankbaren und Undankbaren, fordert dagegen von uns nur Danksagung, daß wir sagen: Gelobet sei Gott der Herr in seinen Gaben; daß ihm also die Gottheit zugeschrieben werde, das ist, daß wir ihn für unsern Gott erkennen, den wir auch mit ganzem Herzen, mit Mund und mit Werken bekennen sollen, daß er der Schöpfer sei. Denn wir können ihm sonst für solche große Gaben Anderes nichts geben oder ihm dieselben bezahlen.

245. Man soll aber darauf sehen, was das Vorhaben und der Wille sei derer, so die geistlichen Gaben kaufen oder verkaufen. Das Vorhaben Simons in den Geschichten der Apostel ist Simonie. Denn er begehrt, daß ihm widerfahren möchte, daß er für sein Geld den Heiligen Geist empfangen möge. Solches ist auch der Verstand der Worte, wenn wir bitten, daß Gott geehrt und gepreiset werde, und daß sein Reich erhalten werde. Gott kann nicht geschändet oder geschmähet werden; denn er ist wesentlich die Ehre, die Gewalt und Gütigkeit selbst. Unser Vorhaben aber und Wille ist also gethan, daß wir begehren, daß kein Gott wäre und daß seine Ehre nur möchte zu Schanden werden. Und dasselbe beweisen wir auch gegen sein Wort und Sacramente. Wenn wir Gott also beleidigen oder schmähen könnten, wie wir sein Wort und Sacramente schmähen, wir würden uns auch dessen alsdann nicht enthalten.

246. So ist nun das die rechte Simonie, nämlich, ein solches Kaufen und Verkaufen, so man mit Gottes Wort und den Sacramenten, mit der Kirche und Heiligen übt und treibt. Diese Simonie hat der Papst nicht verstanden und versteht es auch noch heutiges Tages nicht; wiewohl das ganze Papstthum darin gar ersoffen ist, und ist mit der Simonie behaftet gewesen lange zuvor, ehe denn es nach den geistlichen Rechten darein gerathen ist. Denn sie haben den Himmel, Vergebung der Sünden, Gottes Gnade und den Heiligen Geist für ihre Werke verkauft. Daher sind so viel Kirchen, Klöster, Messen, Ablässe und Mönchorden gekommen, mit welchen Werken wir Gottes Gnade haben verdienen wollen.

247. Auf diese Weise haben sie aus Gott einen Kaufmann gemacht, der das Himmelreich nicht umsonst und aus Gnaden, sondern für Geld und Menschenverdienste geben wolle. Denn wiewohl er es aus Gnaden zu geben anbietet, so werfen sie doch solche gnädige und unverdiente Wohlthat hinweg, und wenden hundert oder tausend Gulden darauf, daß sie damit Klöster und Kirchen bauen, und dergleichen Dinges mehr anrichten, damit sie die Gnade und Wohlthaten Gottes lösen und an sich bringen mögen. Heißt das nicht, die Gnade und Vergebung der Sünden kaufen und damit Hantierung treiben? Und ein solcher Geselle ist der Papst gewesen, da er auch am allerbesten war, und mit solcher Simonie hat er die ganze Welt erfüllt, so gar, daß die Welt jetzt nichts Anderes ist, denn ein solcher unschlächtiger Ort, der voller rechter Simonie ist. Denn was thun sonst die Stifte und Klöster Anderes, denn daß sie Geld nehmen und mit der Gottlosigkeit Hantierung treiben.

248. Aber eigentlich davon zu reden, ist es unmöglich, daß man Gottes Gnade kaufen könne; und wenn du dich dessen unterstehen würdest, so würde es eben so viel sein, als wolltest du einen Bock melken, das ist, du würdest gar vergeblich und nur eine verlorne Arbeit thun. Der Papst stellt den Bock, das Volk hält das Sieb unter, wie man im Sprüchwort sagt. Wenn es nun möglich wäre, daß der Papst reformirt und die gar grobe Simonie möchte ausgefegt werden, so können sie doch von der rechten Simonie nicht gereiniget werden. Denn das ist des Papstes Herz und seine Seele. Wie die Seele die Substanz und Wesen des Leibes ist, also ist der Papst mit Leib und Seele ein Simonist. Wenn man aber die Simonie würde hinwegnehmen und abschaffen, so würde bald alle Pracht, Ehre und Herrlichkeit des Papstes dahin fallen, und er würde nur ein Bischof zu Rom bleiben.

249. Darum mag man nun sehen, wie viel großer Gefahr und bösen Dinges im Papstthum ist, welch ein Schlund es sei, und wie die ganze Welt so greulich verschlungen ist durch den Primat des Papstes, da er allein hat wollen der oberste Bischof und das Haupt der Kirche sein. Darum sollen wir diese Bestie mit vollkommenem Ernst hassen. Und sie kämpfen dennoch noch mehr darüber, daß der Papst seinen Primat, den er sich angemaßet, ganz und vollkommen behalten möge. Was wäre aber das Anderes, denn daß man der Simonie damit ihre Seele und Leben wieder gäbe? Es ist wahrlich eine greuliche Lästerung und schändlicher Mißbrauch des Namens Gottes, welcher mit der ewigen Verdammniß wird gestraft werden. Denn dahin haben sie das Wort, den Namen Gottes und die Kirche mit Gewalt gezogen und ist doch nichts gekauft worden. Sie haben nur gedichtet, daß Gott ein solcher Gott wäre, der die Gnade verkaufen wollte. Ihr Vorhaben und Wille ist Simonisch gewesen. Aber eigentlich ist es an sich selber alles nur der Bock und das Sieb, das ist, eitel verloren und vergeblich Ding, damit sie sind umgegangen.

250. Es ist aber doch dieser Fehler von Natur in aller Menschen Herzen gepflanzet, nämlich, wenn Gott seine Gnade verkaufen wollte, so würden wir dieselbe viel eher und lieber annehmen, denn da er uns dieselbe umsonst und ohne unser Verdienst zu geben anbietet. Zuvor, da die Simonie regiert und im vollen Schwange gegangen ist, waren alle Menschen sehr begierig und bestissen, unzählige Klöster und Kirchen zu bauen. Denn man hat damit das Himmelreich kaufen und verkaufen wollen. Jetzt, da es den Leuten umsonst und ohne Verdienst angeboten wird, und da man sagt: Gott hat seinen Sohn in das Fleisch gesandt und Mensch lassen werden, auf daß er allen, die an ihn glauben, das ewige Leben möchte geben; du hast gar nichts verdient, diese Güter werden dir aus Gnaden gegeben, nämlich, Errettung vom Tode und von Sünden, der Heilige Geist und das Himmelreich: so verachtet die Welt diese Gnade und milde Güte Gottes. Zuvor, da man uns lehrte, wie wir mit unsern Werken die Gnade Gottes verdienen könnten, da liefen die Leute häufig herzu, und jedermann hat darauf alles, was er nur vermocht, gewendet. Wie diese kleine Stadt Wittenberg jährlich den Mönchen tausend Gulden hat gegeben.

251. Es wurde allenthalben der Spruch des Ambrosius, welcher sehr schädlich ist, hoch gerühmt: Regnum coelorum est pauperum, das ist: Das Himmelreich ist der Armen, nämlich, aus dem Spruche Christi Luc. 16, 9.: „Auf daß sie euch aufnehmen in die ewigen Hütten.“ Der Spruch ist allen Menschen durch das Herz gegangen wie ein scharfes Schwerdmesser. Da haben sie gerufen und geschrieen: Wo die Reichen wollen selig werden, so müssen sie es von den Armen kaufen; die Armen aber sind die Franciscaner oder Barfüßermönche (fratres minores) und die andern geistlichen Brüder, die haben das Himmelreich in ihren Händen; von denselben mußt du es kaufen.

252. Was ist denn nun das für ein Scheusal und Ungeheuer, daß die Welt die Erlösung von der Sünde, vom Tode und Hölle, so uns durch Christum aus Gnaden angeboten wird, so schändlich verachtet? Wenn die Simonie wiederum regieren und in Schwang kommen würde, so würden die Leute ganz freigebig alle ihr Gut darauf wenden; aber wenn wir jetzt lehren aus Gottes Wort und Befehl Christi und sagen: Lieber Mensch, nimm die Wohlthaten Gottes aus Gnaden an, und gib den Pfarrherren, Kirchen- und Schuldienern auch wiederum umsonst, daß du dafür nichts wieder begehrest: das will niemand hören, da sind die Leute gar taub, ja, sie berauben die, welchen sie billig helfen und die sie in Ehren halten sollten. Ist nicht das ein erschrecklicher Wahnsinn der Christenheit? Also sind nun alle Menschen mit dem Laster der Simonie behaftet, und man kann doch gleichwohl solche Dinge nicht verkaufen. Denn Gott hat nie einem
\quellseite{41}
Menschen, auch nicht den Mönchen für ihre Klostergelübde etwas Anderes gegeben, denn höllisch Feuer und ewige Pein und Marter.

253. Da mag man aber nun sehen, welch ein jämmerlich und gefährlich Ding das sei, wo man von der rechten Hauptlehre, von der Rechtfertigung des Glaubens und wie man ohne Werke selig müsse werden, abgefallen ist. Denn wo dieselbe Lehre verfinstert und verlöschen ist, so wird die ganze Welt mit Simonie behaftet. Also unterstehen sich die Türken auch, mit ihren Werken GOtt zu versöhnen, und die geistlichen Gaben für Geld und eigen Verdienst zu kaufen. Darum ist es dem Namen nach eine solche Simonie, wie diejenige ist, so im geistlichen Rechte des Papstes beschrieben wird. Denn sie lösen oder erlangen mit Gelde oder eigenem Verdienst gar nichts, sondern das Vorhaben und der Wille ist im Herzen des Menschen so tief eingewurzelt, dadurch alle Menschen von Natur mit dem Laster der Simonie behaftet sind, daß sie wollten, daß GOtt ein solcher GOtt wäre, der sich mit Werken und Verdienst der Menschen versöhnen ließe, dem sie könnten aufreden und vorwerfen: Ich habe viel Gutes gethan, habe viel gefastet, das mußt du ansehen und mir dafür die Seligkeit und das ewige Leben verkaufen. Aber GOtt kennt oder hört solche Uebelthäter nicht, und hält uns seinen Willen vor in der apostolischen Lehre. Als, da St. Paulus sagt (Eph. 2, 8. 9.): „Aus Gnaden seid ihr“ gerecht und „selig worden durch den Glauben; und dasselbige nicht aus euch, GOttes Gabe ist es; nicht aus den Werken, auf daß sich nicht jemand rühme.“ Wo diese Gnade aufgehoben wird, so ist kein Unterschied mehr zwischen Türken und dem Papst, zwischen Heiden und Juden, so sind sie alle mit der Simonie behaftet, sind rechte caparmulgi, Bockmelker.

254. Die Simonie, davon die Canonisten schreiben, ist etwas leichter und nur leiblich, welche man auf äußerliche und weltliche Räuberei ziehen soll: wenn aber der Papst von den Königen und andern Monarchen Kronen, Städte und Land nimmt, daß er ihrer in der Messe dafür gedenken wolle; wenn er den Schein der Göttlichkeit vorwendet, und will für den angesehen und gehalten sein, als der der Allerfrömmste und Heiligste sei: daselbst ist er der Allerbößeste und ist ganz und gar ein Simonist. Denn dort thut er nur mit seinem Stempel Schaden und bekennt die Simonie mit äußerlicher oder weltlicher Räuberei; hier aber bekennet er die Sünde nicht, sondern lobt und rühmt seine Gerechtigkeit und Religion.

255. Derhalben habe ich davon wollen Erinnerung thun, daß man verstehen könne, welches die rechte Simonie sei, und daß wir ja bei der reinen Lehre bleiben, die uns lehrt, wie uns GOtt seinen Sohn habe gegeben, und durch ihn alle himmlischen und geistlichen Güter, welche die Welt für Geld und Verdienst der Werke kaufen, und nicht umsonst und aus Gnaden empfangen und annehmen will.

256. Nun wollen wir wiederum auf die Historie kommen. So läßt sich ansehen, als habe Jakob auf den Kauf gedrungen, also, daß man ihn nicht wohl entschuldigen und lossprechen kann von dem Laster der Simonie. Denn er sagt ja klärlich: „Verkaufe mir deine Erstgeburt.“ Mit welchem Recht aber fordert er dasselbe und begeht so eine große Sünde, daß er sich unterstehen darf, den Segen zu kaufen, der da geistlich ist, und gehört zu Christo und zukünftigen Dingen, welche alle heilig sind und das zukünftige Leben betreffen?

257. In diesem Leben ist alles weltlich, unheilig und befleckt; die Sacramente aber und andere heilige Dinge sind nicht um dieses gegenwärtigen Lebens willen eingesetzt und verordnet, sondern darum, daß sie eine heilige Vorbereitung sein sollen auf die zukünftigen Dinge. Denn durch das Evangelium werde ich zubereitet, wiedergeboren und verändert auf das zukünftige Leben; es ist mir nicht dazu gegeben, daß ich davon soll reich werden. Wer aber nun ein Pfarrherr oder Prediger ist, und das Evangelium darum predigt, daß er Reichthum und Ehre dieser Welt damit erlangen möge, der ist mit dem Laster der Simonie behaftet. Also ist dies der Gebrauch und Nutz der Taufe, daß ich dadurch aus dem Schooße meiner Mutter und aus dem Grabe versetzet und geleget werde in das Paradies, und aus dem Tode in das Leben. Derhalben suche ich daselbst keine Lust und Freude dieser Welt, sondern brauche der geistlichen Gaben also, daß wenn die leiblichen aufhören, ich alsdann aus diesem Leben in das ewige Leben und Unsterblichkeit geführet möge werden. Wer nun darum ein Pfarrherr ist, darum absolutirt oder Sacramente reicht, daß er davon Ehre und Gewalt will erlangen, der ist mit der Simonie behaftet nach der groben Definition, wie die Canonisten die Simonie beschrieben haben; er ist aber nicht in der Zahl derer, die noch etwas subtiler und behender sünd, und das ewige Leben mit Werken verdienen wollen.

258. Warum begehrt aber Jakob, daß ihm die Erstgeburt möge verkauft werden? Und dieweil er dieselbe kaufen will, so folget ja, daß es eine Hantierung ist. Darauf antworte ich also: Jakob hat Achtung auf die Gelegenheit, dadurch er die Erstgeburt erlangen möchte, und dieselbige Gelegenheit hat er gebrauchen wollen.

259. Darnach wird sein Bruder Esau unter die gerechnet, die gar grob und äußerlich mit der Simonie beladen sind, das ist, er hat im geistlichen Segen Anderes nichts gesucht, denn nur den Bauch damit zu füllen. Denn er ist nicht in die Hütte gegangen, sondern hat den Gottesdienst verachtet. Er hat erstlich das Kleine verachtet, bis er darüber in Verachtung der größeren Güter gefallen ist; wie Kain, der Erstgeborne, auch gethan hat. Esau hat nicht verstanden, wie groß der Segen gewesen ist, darum hat er ihn verkauft für solche Güter und Lust, so nur diese gegenwärtige Welt betreffen, also daß er vor GOtt wahrhaftig der Erstgeburt und aller Güter beraubt ist worden. Darum hat Jakob mit Fleiß Achtung gegeben auf die gegenwärtige Gelegenheit, und ist mit seinem Bruder eins geworden, daß er ihm gewisse Bezahlung gibt, damit ihm hernachmals sein Bruder um der Erstgeburt willen keinen Lärm oder sonst andere Unlust machen oder erregen könnte, und daß er ihm auch derhalben nichts aufreden oder vorwerfen möchte.

260. Also hat heutiges Tages der Papst den Namen der Kirche und wir haben ihn nicht: wir wissen aber doch, daß wir wahrhaftig die Kirche sind. Denn wir haben das Wort, die heiligen Sacramente und die Schlüssel, so Christus hinterlassen und zu gebrauchen befohlen hat, nicht daß sie zu unserer Gewalt und Lüsten dieses Lebens dienen sollen, sondern daß sie uns bereiten sollen auf die Zukunft des Sohnes GOttes. Darum sind wir die rechte Kirche.

261. Dagegen aber eignet sich der Papst diesen Titel und Namen an, weil er im Besitz ist, und rühmt sich, er sei der Nachfolger und Stuhlerbe Petri und Pauli. Wenn ich nun mit dem Papst einen solchen Vertrag machen wollte und zu ihm also sagen: Ich will dir jährlich so viel tausend Gulden geben, daß du mich wollest das Evangelium frei lauter und rein predigen lassen, wäre das auch Simonie? Antwort: Nein, denn ich kaufe das nicht, so ich schon zuvor habe, nämlich, den rechten Gebrauch der heiligen Sacramente, die reine Lehre des Evangeliums, des Glaubens, der Hoffnung, da ich aufs zukünftige Leben warte, welches mir alles schon erlaubt ist auch wider des Papstes Willen und ohne seine Macht oder Gewalt. Ich suche aber damit nur das, daß ich ihm die tausend Gulden zu geben anbiete, daß ich mich von seiner Quälerei loskaufe. Derhalben kaufe ich von ihm nichts, sondern werfe nur dem Widersacher eine Suppe und ein roth Linsengericht vor, auf daß ich also den Papst berauben möge des Namens, den er führt, daß er will die Kirche sein: ich achte sonst seiner Bullen und Ablaßbriefe gar nicht. Derhalben wollte ich den Fürsten unserm Theils rathen, daß sie ihm jährlich etliche tausend Kronen gäben um der Ursache willen, damit doch das Blutvergießen und anderes Aergerniß mehr aufhöhren möchte und das Evangelium im Frieden möchte ausgebreitet werden. Dies ist keine Simonie, sondern ist eine redemtionem vexationis, daß man es ihm abkauft, daß er uns nicht also plagen soll.

262. Darum da Jakob zu Esau sagt: „Verkaufe mir deine Erstgeburt“, dies ist in Wahrheit kein rechter Kauf. Denn die Erstgeburt hatte schon zuvor dem Jakob gebührt und Esau ist derselben beraubt worden. Bei Jakob ist der Besitz, er hat aber den Namen nicht. Bei Esau ist der Name und nicht der rechte Besitz. Dieweil aber nun Jakob sieht, daß Esau keinen andern Verstand hat, oder in dem Segen Anderes nichts sucht, denn daß er den Bauch füllen möge; wie der Papst auch unter dem Schein, daß er den Namen der Kirche führt, nur seine Ehre und Gewalt sucht: darum wirft er dem Esau eine Suppe in das Maul, auf daß er ihm hernachmals keine Mühe damit machen soll.

263. Auf solche Weise wollte ich zwar wünschen, daß unsere Fürsten dem Kaiser jährlich zwanzigtausend Gulden geben wollten, auf daß die Lehre des Evangeliums im Frieden und stiller Ruhe möchte gepredigt werden. Ja, möchte Einer sagen, so höre ich wohl, ihr wollt für solche geistliche Gaben Geld geben? Antwort: Mit nichten; denn was mir jetzt schon haben, daßelbe kaufen wir nicht und sie können uns das auch nicht nehmen, sondern wir wollen das Geld darum ausgeben, daß sie uns sollen unbemüht lassen und daß uns der Pabst und Kaiser Frieden geben wollten.

\quellseite{42}

264. Aber, lieber Pabst und Cardinäle, ihr sollt gleichwohl wissen, daß diese geistlichen Güter nicht euer, sondern unser sind. Und wenn schon der Pabst nimmermehr beiderlei Gestalt im Sacrament würde geben, so haben wir dieselben doch von GOtt, und sie können die nicht verkaufen; oder, wenn sie das schon thun wollten, so werden sie den Bock melken und das Sieb unterhalten. Wir haben beiderlei Gestalt wider des Pabstes Willen und mit GOttes Erlaubniß. Wenn sie aber nicht wollen aufhören, uns zu verfolgen, zu würgen und zu berauben, so wollen wir mit ihnen einen solchen Vertrag machen, daß sie das rothe Gericht fressen; die Ehre sollen sie aber nicht behalten, daß sie uns beiderlei Gestalt des Sacramentes und andere dergleichen geistliche Gaben gegeben oder geschenkt hätten. Denn das wäre so viel, als wenn wir in ihre Gewalt, die sie sich anmaßen, auch mitwilligen und sie bestätigen wollten. Und dieses ist ein göttliches und heiliges Gut, das ich für kein Geld kaufen kann. Darum wollen wir es auch vom Pabst weder gegeben noch geschenkt haben.

265. Also wird Esau (wie droben 1 Mos. 21, 9. Ismael auch) seinem Bruder Jakob ohne Zweifel viel zum Verdruss gethan und wird trotzig über ihn geherrscht haben, darum daß er der Herr und Erstgeborne wäre; Jakob aber hat für einen Knecht und Narren müssen gehalten werden. Von solcher Verfolgung, und daß ihm sein Bruder solchen Verdruß bereitet, hat sich Jakob befreien und erretten wollen, auf daß er den geistlichen Segen sicher behalten möchte.

266. Also findet man heutiges Tages viel fromme Christen, die um des Bekenntnisses willen des Evangeliums aus den Städten weichen, auf daß sie die Tyrannei der Päpisten nicht leiden müssen. Und St. Paulus sagt auch also Röm. 12, 18: „Ist es möglich, so viel an euch ist, so habt mit allen Menschen Friede.“ Wollen die Widersacher nicht Friede halten, so gib ihnen ein roth Gericht, daß sie dir Friede lassen, das ist, daß sie dich der Güter, so du schon zuvor hast, mit Frieden gebrauchen lassen.

267. Nun folgt die Antwort Esaus auf seines Bruders Bitte. „Siehe“, sagt er, „ich muß doch sterben, was soll mir dann die Erstgeburt?“ Da vergewissert der Bösewicht, und hat sich mit seinem eigenen Zeugniß verrathen, daß er ein gottlos Herz gehabt. Soll man also mit der Religion, mit der Gottseligkeit und Erstgeburt umgehen? Soll man davon nicht mehr halten? Denn jetzt bekennt er öffentlich, daß er die Erstgeburt nicht verzeihe, und daß er sie nicht für ein Erbe des zukünftigen Lebens und für eine heilige Sache halte.

268. Wenn aber die Erstgeburt zu nichts Anderem nütze wäre, denn nur zu großer Ehre und Reichthum dieses Lebens, wo würde dann der Glaube, wo würde die Hoffnung und das Warten auf die ewigen Güter bleiben? So lasset uns nun mit den Gottlosen und Epicurern sagen, wie sie pflegen zu sagen. Wenn sie hören, daß man den Himmel nennt, so antworten sie und sagen: Was Himmel! hätten wir hier Mehl. Wie der Pabst auch nichts mehr von der Kirche und Schlüssel hält: Was sagst du mir viel von Schlüsseln? ich muß sie gebrauchen, daß ich damit Geld und Gut sammle, und die Könige und Fürsten bezwingen könne. Dieses sind die Gauklerischen Worte der Gottlosen, die da meinen, daß die Gottseligkeit nur darin bestehe, daß man davon großen Gewinn und Genuß haben möge; und würden Christo und seinem Evangelium bald gute Nacht sagen und fahren lassen, wenn sie nicht zeitlich davon groß Glück und viel Geldes und Gutes haben könnten.

269. Derhalben ist dies ein sehr tapfer Zeugniß aus der heiligen Schrift, daß Esau jetzt die Erstgeburt schon verloren hat. Denn er wird aus seinem eigenen Munde gerichtet; wiewohl er ohne Zweifel also wird gedacht haben: Ich bin der Herr; ich will zwar die Erstgeburt wohl verkaufen, aber ich will damit meines Bruders nur spotten. Denn wiewohl ich sie ihm werde verkaufen, kann ich sie doch wohl wieder an mich bringen. Also hat er seines Bruders gespottet und ihn verachtet. Wie der Pabst mit seinen Larven die ganze Welt verspottet und sie getäuscht hat, da er seine guten Werke und Ablaß verkauft; und hat zwar gemeint, es sollte ihm also dabei bleiben und würde ihm Alles frei ausgehen, und müßte ihn niemand darum strafen, oder aber, er würde noch endlich wiederum können Buße thun.

270. Und heutiges Tages sind derer auch Viele, so der Lehre des Evangeliums schändlich mißbrauchen, und können unterdeß sich fein liebkosen und sagen: Wiewohl ich jetzt in Sünden liege und gottlos bin, so will ich doch endlich dermaßen wiederum Buße thun und zum Leben kommen. Wie auch Etliche vom Herzog Georg sagen: Ei, er hat noch wohl an seinem letzten Ende, da er jetzt hat aus diesem Leben abschieden sollen, bekehrt können werden. Aber diese alle werden zuletzt dies betrübte Urtheil hören aus dem Spruch, so in der Epistel an die Hebräer Cap. 12, 17. von Esau geschrieben steht: „Wisset aber, daß er hernach, da er den Segen ererben wollte, verworfen ist; denn er fand keinen Raum zur Buße, wiewohl er sie mit Thränen suchte.“

271. Also genießen die Antinomer auch der Güter dieser Welt, daß es ihnen nur in diesem Leben möge wohl gehen, und sagen, sie wollen sich noch zu seiner Zeit wohl bekehren und Buße thun; verachten also ihren Segen, die Kirche, Taufe, die Schlüssel, Vergebung der Sünden, ewiges Leben und GOttes Gnade; empfahen sie alle vergeblich; verschwenden die angenehme Zeit und den Tag des Heils, davon St. Paulus sagt 2 Cor. 6, 2. Sie werden aber hernach solche Gelegenheit, so sie jetzt versäumen und verachten haben, vergeblich suchen und werden alsdann zu lange geharrt haben; wie die Braut im Hohenliede Salomonis klagt im 5. Cap. V. 6.: „Da ich meinem Freunde aufgethan hatte, war er weg und hingegangen.“

272. Derhalben, da Esau sagt: „Siehe, ich muß doch sterben, was soll mir dann die Erstgeburt?“ zeigt er damit an, daß er nur auf den Bauch gedacht habe, wie er den füllen möchte, und daß er die Verheißung verworfen habe, als die vergeblich und kein nütze sei zu dem zukünftigen Leben; ganz und gar wie heutiges Tages die Epicurer auch sagen: Wozu ist es vonnöthen, daß man das Evangelium soll hören, der Schlüssel oder des Abendmahls des Herrn gebrauchen?

273. Jakob aber ist viel anders gesinnt, und weiß, daß solche heilige Sachen zur Vorbereitung dienen auf das zukünftige Leben. Darum achtet er sie groß, und ist ganz brünstig vor großer Liebe, so er zu solchen göttlichen Dingen trägt. Die aber nun derselben nur zu diesem Leben gebrauchen wollen, dieselbigen werden derselben aus göttlichem Recht und Gewalt beraubt. Denn daselbst ist es ihr freier Wille, daß sie mit gutem Wissen und Willen solche göttliche Güter verkaufen oder verwerfen.

274. Nachdem aber nun Jakob sieht, daß die rechte gelegene Zeit eben jetzt vorhanden ist, da hält er mit seiner Bitte an, und will haben, daß der Vertrag nun soll bestätigt und vollzogen werden. „So schwöre mir“, sagt er, „heute“, oder jetzt, bald in dieser Stunde. Denn man soll solche Gelegenheit nicht versäumen. „Da schwur ihm Esau, und verkaufte also Jakob seine Erstgeburt.“ Also sollte man mit den Päpisten auch handeln, wenn sie schwören wollten und Bürgschaft setzen, dazu auch Brief und Siegel geben, daß sie keinen Krieg wider uns anfangen wollten, so wollten wir ihnen alsbald ein roth Zinsgericht geben.

275. Nun ist aber noch eine andere Frage übrig, so von Andern allhier angeregt wird. Von derselben wollen wir also handeln, damit sie den Leser nicht irre mache, nämlich: Ob das rothe Gericht der Preis gewesen sei, dafür Jakob die Erstgeburt von seinem Bruder gekauft hat? Darauf antworten sie und sagen, daß es nicht der eigentliche Kaufpreis gewesen sei, so Jakob für die Erstgeburt gegeben habe, vielmehr die Erstgeburt ein sehr groß und köstlich Ding gewesen. Darum, sagen sie, sei es der Wahrheit nicht gemäß, daß Esau sollte dies schlechte Gericht genommen haben als Bezahlung der Erstgeburt. Derhalben wollen sie haben, daß dies Gericht nur ein Pfand soll gewesen sein, sein Kauf damit zu bestätigen, welches Esau zum Wahrzeichen empfangen habe, daß er Jakob die Erstgeburt verkauft hatte; wie man bei uns zwischen dem Käufer und Verkäufer pflegt eine Kanne Weins oder zwei aufzusetzen zum Zeichen, daß Einer dem Andern etwas hat abgekauft. Die Deutschen heißen es einen Gleichkauf.

276. Diese Meinung ist aus den Büchern der jüdischen Rabbinen hergekommen, und wir verwerfen dieselbe und nehmen sie nicht an; wie man es mit allen ihren Meinungen halten muß, nämlich, daß man das Gute, so darin ist, behalte, aber was böse und falsch ist, verwerfe und sich davor hüte. Denn diese Meinung allhier ist straks wider den klaren Text in der Epistel an die Hebräer Cap. 12, 16., da ausdrücklich gesagt wird, daß Esau um einer Speise willen seine Erstgeburt verkauft habe.

\quellseite{43}

277. Denn die heilige Schrift hat damit wollen abmalen die Art und Sitten der falschen und heuchlerischen Kirche, welche man an ihren Sitten erkennen und verstehen muß lernen, deren Vorbild und Figur Esau ist gewesen. Denn dieselbe pflegt sich des Namens Gottes fast hoch zu rühmen und will mit Gewalt die Kirche und für die Kirche gehalten sein, und ist es auch vor der Welt, wie Esau die Erstgeburt hat; aber er verachtet dieselbe und durch solche Verachtung hat er sie auch verloren. Denn er wird wiederum verachtet, verworfen und beraubt. Darum ist es nichts mit seinem Rühmen; wie es mit der Papisten Rühmen auch nichts ist, die sich dessen hoch rühmen dürfen, daß sie die Kirche und Nachfolger der Apostel seien. Denn wer die Kirche sein will, derselbe muß wahrlich die Erstgeburt groß achten, nicht allein hier im Leben, sondern auch nach diesem Leben.

278. Und dieser Zank währt in der Welt für und für. Wie Esau mit Jakob gezankt hat, also haben wir mit den Papisten auch einen solchen Zank, und vor uns haben die Juden mit den Aposteln auch also gezankt. Die Juden haben darüber heftig gestritten, daß sie Gottes Volk wären, und St. Paulus Röm. 2, 17. und 9, 24. bekennt auch, daß sie die Kirche seien. Dieselben haben allein den bloßen Namen behalten, die rechte wahre Kirche aber an sich selbst verachtet haben, darum sind sie verworfen und die Heiden berufen und angenommen worden; wie im Propheten Hosea im 2. Cap. V. 23. steht: „Ich will sagen zu dem, das nicht mein Volk war: Du bist mein Volk; und es wird sagen: Du bist mein Gott“; und 5 Mos. 32, 21.: „Mit ihrer Abgötterei haben sie mich erzürnet. Und ich will sie wieder reizen an dem, das nicht mein Volk ist; an einem närrischen Volk will ich sie erzürnen.“

279. Derhalben ist dies das erste Wahrzeichen, ja, eine sehr scheußliche böse Farbe der heuchlerischen Kirche, nämlich, daß sie sehr stolz und hoffärtig ist, und trotzt heftig mit dem Namen Gottes, der Kirche, der Schlüssel, des Kirchenamtes, der Sacramente, und kann nicht mehr, denn daß sie sich nur Gottes rühmt.

280. Das andere Wahrzeichen aber ist, daß sie in Wahrheit die Dinge verachtet und versäumt, so zu der rechten Kirche gehören; sie fragt nicht viel nach ihrer Erstgeburt. Wie der Papst und seine Bischöfe weder lehren noch taufen, und keinerlei Kirchenamt verwalten. Sie haben nur den Namen, die Frucht aber und das Werk beweisen sie nicht. Darum sagt Christus Math. 7, 16.: „An ihren Früchten sollt ihr sie erkennen.“ Dieweil aber der Papst die Kirche verachtet, derhalben verliert er auch die Kirche, und gilt der vergebliche Ruhm gar nichts, daß sie sagen: Wir sind die rechten Christen, wir sind die Kirche, sind Bischöfe. Wenn die Wahrheit nicht auch dabei ist, so wirst du den Namen nicht lange behalten, ja, du bist jetzt schon verworfen, und was du hast, das hast du doch nicht. Also werden die Ketzer, so nach uns kommen werden, die Kirche auch verlieren. Denn sie ist ein Pfeiler und Grundfeste der Wahrheit, 1 Tim. 3, 15., welche sie verwerfen werden, und hilft sie ihr Rühmen gar nichts und daß sie den Namen haben ohne Wahrheit.

281. Wenn sie nun schreien und sagen: Wir sind die Kirche, wir haben Macht und Gewalt zu gebieten und zu verbieten, und wollen uns zwingen, daß wir ihren Statuten und Satzungen beifallen sollen, so wollen wir ihnen darauf also antworten: Liebe Herren, ihr sollt wissen, daß zweierlei Kirchen sind, nämlich, die eine, welche man gewöhnlich dafür hält, und die andere, welcher der Name mit Recht gebührt; item, die eine, so den Schein hat der Kirche, und die andere, so in Wahrheit die Kirche ist; oder, die eine dem Namen nach, die andere dem Wesen nach, wie sie selbst lehren. Esau ist nur mit dem Namen und nach dem äußerlichen Schein die Kirche, hat aber dabei das Wesen nicht, wie alle Abgöttischen und Ketzer, die sich auch den Namen der Kirche anmaßen.

282. Das kann ich zwar wohl leiden und zugeben, daß Einer nur dem Namen nach ein Christ sei und doch unterdeß äußerlich das Amt der Kirche verwalte. Wenn auch Einer schon heimlich gottlos und ein boshafter Mensch ist, wo er nur tauft, die Kranken besucht, und die armen betrübten Gewissen tröstet und stärkt, so hindert das nicht, daß er äußerlich unter die Christen und zur Kirche gerechnet werde, aber nur nicht von Rechts wegen. Dasselbe thut aber der Papst und seine Bischöfe nicht. Darum kann man sie auch nicht einmal mit den Schriften vergleichen, die da heimlich gottlos sind; denn der Papst und seine Bischöfe thun, das dem Kirchenamt stracks zuwider ist: sie lästern das Wort, verfolgen die gottseligen frommen Christen, das ist, die rechte wahre Kirche.

283. In unserer Kirche sind zwar viel Diener, nur dem Namen nach, aber nicht dem Recht nach. Aber der Papst ist öffentlich ein abgesagter Feind der Kirche, er unterdrückt die reine Lehre vom Glauben, erfüllt die Welt mit eitel Abgötterei, Simonie und vielen unzähligen scheußlichen Dingen mehr. Das heißt recht, unsinnig sein, und wider die sichtbare und unsichtbare Kirche wüthen und toben, wie die Türken thun. Darum gilt der Papisten Argument nicht, wenn sie sagen: Wir sind die äußerliche Kirche, deshalb sollt ihr uns dulden und leiden. Wenn sie solche Leute sein wollten, wie sie wollen genannt sein und die Kirche heißen, so müßten sie mit uns in der Lehre eins sein: sie müßten lehren und das Kirchenamt recht verwalten, welches rechten und frommen Bischöfen gebührt. Aber sie sind Esaus Nachkommen, sind Widersacher und Zerstörer der Kirche, wie die Ketzer sind.

284. Die sich aber stellen, als ob sie gottselig und fromm seien und doch Heuchler sind, wie deren Viele unter der Versammlung der Gottseligen gefunden werden: dieselbigen dulden wir, daß sie auch mit unter die Gläubigen gezählt werden, sofern sie gleichwohl nichts lehren, das dem Bekenntniß der rechten Lehre zuwider sei.

285. Das dritte Wahrzeichen der falschen heuchlerischen Kirche hat Mose in diesem Vorbilde auch abgemalt. Denn dieweil Esau so gar sicher gewesen, und mit großem Trotz die geistlichen Güter oder Gaben verachtet hat, als nämlich, den Glauben, die Liebe, Geduld, Freundlichkeit, welche in der Kirche sind, ja, dies alles gar verachtet hat: so folgt nun weiter, daß er Lust und eine hitzige Begierde gewinnt zu dem rothen Gerichte. Das ist so viel gesagt: Dieweil er nach der Erkenntniß Christi nichts fragt, achtet dessen auch nicht, daß bei ihm die geistlichen Gaben möchten gemehret werden, achtet des ewigen Lebens nicht: da wird er nun auch ein Geizhals, und ein solcher Mensch, der nur nach Wollust, Ehre und Herrlichkeit dieses Lebens strebt.

286. Ja, das ist es, daß er von dem Linsengerichte mehr hält denn von der Erstgeburt. Das ist so viel gesagt: Er gehört zur äußerlichen Kirche; er folgt aber Fleisch und Blut nach, sucht bei der Gottseligkeit nur leiblichen Gewinn und Genuß. Also stellt sich die falsche Kirche, als ob sie zumal gottselig sei, nur darum, daß sie der Wollust und Ehre dieser Welt genießen möge. Aber die rechte Kirche trachtet nach dem ewigen Leben durch Geduld und rechten Glauben, jene aber sind rechte Mauldiener und der Bauch ist ihr Gott, Phil. 3, 19.

287. Dies ist das Geheimniß mit dem rothen Gerichte, welches uns an diesem Ort die heilige Schrift vorhält, damit der Zank zwischen diesen Brüdern von der Erstgeburt entschieden wird. Die Linsen aber sind nicht roth, sondern sind vielmehr gelb. Derhalben meine ich, sie werden mit Safran oder anderer Würze zugerichtet gewesen sein, auf daß sie eine Begierde machen könnten zum Essen. Also ist es aber um die gegenwärtigen Güter dieser Welt bestellt, daß dazu die Menschen sehr große Lust haben, daß sie viel lieber der Wollerei, Saufen und Fressen nachgehen, denn daß sie sollten das Fleisch und Blut des Sohnes Gottes essen und trinken, das ist, glauben und nach den himmlischen Gütern trachten. Sie suchen nur ihren Gewinn mehr denn Gottes und die Ehre dieser Welt; schlechte Dinge so gar angenehm und lustig sind, daß sie den Menschen nicht allein Hunger, sondern auch Lust und Begierde machen. Darum hat Esau dreimal wiederholt, daß er von seinem Bruder begehrt und sagt: Lieber, speise mich doch, laß mich kosten das rothe Gerichte, ja, das rothe Linsengericht.

288. Darum will hiermit die heilige Schrift anzeigen, daß alle Menschen immer geiziger, hoffärtiger und fleischlicher werden, wo sie den Segen verlieren, denn sie zuvor gewesen sind; denn sie haben alle Fleisch und Blut. Darum suchen sie auch nichts Anderes, denn was fleischlich ist, und mit greulicher Lust begehren sie das rothe Linsengericht. Aber das geschieht alles unter dem Namen Gottes und Christi, unter dem Titel der Kirche. Jakob ist zwar auch von dem rothen Gerichte, aber ist dessen nicht so begierig; er wiederholt „das rothe“ nicht zweimal, wie Esau thut. Das ist so viel gesagt: Die Gottseligen gebrauchen der gegenwärtigen Güter nicht als zur Seligkeit, sondern nur zu ihrer Nothdurft.

\quellseite{44}

289. Zuletzt setzt Mose noch hinzu, wie es Esau so gering geschätzt, daß er die Erstgeburt verkauft hatte. Daselbst wird auch der große Trotz und Hoffart angezeigt und gestraft. Denn er hat gedacht, er wäre durch menschliche Hülfe wohl gesichert, er könnte sein Recht wider seinen Bruder leichtlich erhalten, angesehen, daß er mit ihm einen Contract und Kauf gemacht hätte.

290. Also pflegen die Heuchler und falschen Brüder zu thun: sie tragen den Baum auf beiden Achseln, und sind tückisch und listig; jetzt stellen sie sich, als wollten sie weichen und etwas nachgeben, bald reißen sie wiederum an sich, das ihnen mit keinem Recht gebührt, und haben doch ihre Glossen, damit sie ihre Heuchelei bedecken und sich entschuldigen können, daß sie zugleich mit der Erstgeburt auch das Linsenmus behalten mögen. Denn sie wollen zugleich GOtt und dem Bauch oder dem Mammon dienen.

291. Also hat Esau auch gedacht: Wiewohl ich die Erstgeburt verkauft habe, item, geschworen und die Suppe gegessen habe, und es sich ansehen läßt, als habe mein Bruder, der gar verachtet gewesen, die Erstgeburt hinweg genommen: so will ich sie doch behalten und ihm damit seinen Glauben halten. Denn er hat sie im Ernst nicht verkauft, sondern hat seines Bruders spotten und ihn betrügen wollen, als hätte er es ihm verkauft, und wollte doch gleichwohl sein Recht erhalten.

292. Aber die Epistel an die Hebräer Cap. 12, 17. zeigt an, daß er hernach verworfen sei. Und wiewohl er sich unterstanden, die Erstgeburt wieder an sich zu bringen, darum daß es ihm gereuet, daß er sie verkauft hatte, so habe er doch mit solcher Buße nichts ausgerichtet. Nicht sage ich, daß er nicht solle selig sein geworden, sondern daß er den Segen, so er einmal verloren, mit seinen Thränen habe wiederum erlangen können. Also verstehe ich es nun, daß er für diese Suppe die Erstgeburt verkauft habe, und nehme der jüdischen Rabbinen Glosse nicht an. Denn sie sind fleischliche Menschen, die in der heiligen Schrift keinen Verstand haben.

\chapter{Das sechsundzwanzigste Capitel}

\section[Erster Theil]{Erster Theil. Von Isaaks Reise nach Gerar, und wie ihm der HErr erschienen und ihn mit einer neuen Verheißung aufgerichtet hat}

\begin{quote}
\textbf{V. 1. Es kam aber eine Theurung ins Land über die vorige, so zu Abrahams Zeiten war. Und Isaak zog zu Abimelech, der Philister König, gen Gerar.}
\end{quote}

1. Es ist nun oftmals gesagt worden, daß in diesem ganzen Buche die Historien der Väter beschrieben werden wie auf einer gar schlechten Rinde oder Blatt, und daß sie uns vorgehalten werden ohne allen Schein oder äußerliches Gepränge der Religion, Gerechtigkeit und Weisheit, ja, in einer sehr verachteten Gestalt der Haushaltung und dieses leiblichen Lebens. Denn was erzählt Mose von Isaak Anderes, denn daß er von seinem Vater Abraham geboren sei, daß er Kinder gezeuget, das Vieh gemehret, und bald in diesem, bald in einem andern Lande umher gezogen? Da wird gar wenig oder nichts vom Gebet gelehrt und von den abenteuerlichen Religionen der Mönche. Was geht aber mich das an, daß er ein Ehemann gewesen und bei seinem Weibe geschlafen hat? Soll man das in der Kirche lehren? Denn also hängt das Fleisch an dieser äußerlichen und sehr geringen verachteten Gestalt, wenn es das Leben der Väter ansieht, und wird daraus gar nichts gebessert, sondern nur geärgert.

2. Danach aber denkt es, wie es das Leben des Bernhardus, Antonius, und dergleichen anderer Mönche mehr möge verstehen lernen, darin wunderbare und unglaubliche Werke gefunden werden, wie sie sich etlicher Speisen und anderer Dinge mehr enthalten, wie sie gefastet und gewacht haben; da man von keinen Weibern, Knechten und viel weniger vom Vieh findet, wie man damit wäre umgegangen. Diese gemeinen und verachteten Dinge in der Haushaltung Abrahams, Isaaks oder Jakobs verachtet, sage ich, das Fleisch und treibt damit sein Gespötte, und folgt dieweil andern Werken, die einen großen Schein haben und prächtig sind. Oder wo man diese Historien etwa in der Kirche lieset, verwundert sich niemand derselben; denn niemand sieht auf ihre rechte Würde und den rechten Schmuck, davon solche Historien ihre Ehre haben.

3. Derhalben lehren wir also und das soll man auch den Leuten fleißig einprägen, nämlich, daß man in den Exempeln der heiligen Väter dies als das rechte Hauptstück und höchstes Lob ansehen soll, daß GOtt mit ihnen geredet hat, daß sie GOttes Wort gehabt. Das ist das rechte Stück, so diese Historien zieret, und den rechten Verstand gibt, was und wie groß diese heißen und scheinbar verächtliche Gestalt der heiligen Männer sind. Denn wo GOttes Wort ist, da ist auch der rechte Glaube und da sind auch rechte gute Werke; denn da geschieht Alles im Wort und unter dem Wort. Die andern Dinge, so außerhalb und ohne das Wort geschehen, allein nach unserm eigenen Willen und Gutdünken, sind wahrhaftig nichts Anderes als Hefen und nur eitel Dreck vor GOtt.

4. Darum ist aller Mönche Leben, wie köstlich es auch scheint vor den Augen des Fleisches, doch gar Nichts. Es ist ein großer Schein der Heiligkeit an Antonius, Hilarion, und vielen Andern mehr gewesen, deren etliche mit Fasten, etliche mit wunderbarem Wachen ihr Leben bis in das siebenzigste Jahr zugebracht haben; und zu solchen wunderbaren Werken hat das Fleisch, das Herz und Auge des Menschen Lust. Siehe aber du zu, ob auch GOttes Wort dabei sei. Frage Antonius, ob er auch ein Wort habe, das ihn habe in die Wüste heißen gehen und sein Fleisch kasteien? Da wird er Nein sagen; sondern es hat mich also gedünkt gut und GOtt angenehm zu sein. Ja, lieber Antonius, das ist das rechte Hauptstück, welches in deinen Werken vornehmlich sich sollte sehen lassen, ohne welches das ganze Leben, so du führst, der Tod ist, und nur eine Wahl des Fleisches, ein eigen gut Bedünken, und eine lautere Prahlerei und Bethörung fleischlicher Menschen.

5. Darum lobt nun die heilige Schrift in den Historien der Väter vornehmlich den Glauben an das Wort; denn das Wort GOttes heiligt alle Dinge, sintemal es heilig ist, ja, es ist die Heiligkeit, Wahrheit und Weisheit selbst. Und das Leben, so durch das Wort regiert wird, ist ein wahrhaftiges, gerechtes, weises und ewiges Leben. Wo es aber das Wort nicht hat, so hat es auch keine Wahrheit, kein Licht, keine Weisheit vor GOtt, und alle seine Werke sind eitel Werke der Finsterniß.

6. Also, wiewohl der Ehestand ein unreiner Stand ist (weil die Vereinigung des Mannes und Weibes nicht ohne fleischliche Unreinigkeit geschehen kann), auch ein unflätig Ding ist, des Viehes Hüten; dazu auch der Obrigkeit und Unterthanen Leben fast unrein und voller Gebrechen ist, daselbst viel Sünden zusammen kommen, dazu auch sonderliche Uebertretungen und Gebrechen, so an den Personen sind, noch über die gemeinen Gebrechen: so hat doch GOtt dies alles herrlich geziert und in seinem Wort verordnet, und wenn du dem Worte anhängst, so bist du schon gereiniget von aller deiner Unreinigkeit.

7. Endlich ist kein Mensch, der ohne Sünde lebe; aber das Wort hat solche große Kraft, daß es dies alles verschlingt, daß du sagen magst: Ich lebe im Ehestande mit meinem Weibe und Kindern in gutem Frieden, in Gottesfurcht und Vertrauen auf ihn, und also weiß ich, daß Alles recht ist und wohl steht; denn also sagt St. Paulus vom Leben der Eheleute 1 Tim. 2, 15.: „Das Weib wird selig durch Kinderzeugen.“ Wie aber? „So sie bleibet im Glauben“ etc. Dies ist die erste Kraft des Wortes.

8. Danach gibt das Wort nicht dies allein, daß der Stand heilig sei und GOtt wohlgefalle, sondern es erweckt dich auch zu allerlei Tugenden und zu lieblichem guten Werken. Denn es feiert nicht, wo es nur in deinem Herzen ist nach dem rechten Verstande, sondern es wird dich erinnern, daß du gedenkest, wie du GOtt wollest anrufen und loben; es wird dich zum Priester und Propheten GOttes machen, daß Opfer GOtt sehr angenehm sein werden; denn deine Augen sehen auf den Glauben. Und diese Frucht und Kraft des Wortes sehen die heuchlichen Mönche auch nicht.

\quellseite{45}

9. Zum dritten, wo der Teufel fühlen und merken wird, daß du das Wort hast, und dich darauf verlässest, daß dein Leben Gott angenehm und wohlgefällig ist um des Wortes willen. So wird er nicht ruhen, sondern wird allerlei Anfechtung und Kreuz auch in den geringsten Dingen lassen vorfallen. In der Haushaltung wirst du erfahren die große Untreue des Hausgesindes, Haß und Feindschaft der Nachbarn, item, daß dir die Kinder oder das Weib dahin stirbt. Dies wird dir alles darum begegnen, daß dadurch dein Glaube geübt und geprüft werde. Wo aber nun dabei das Wort nicht ist, so folgt Ungeduld und Zorn um dieses Standes willen, der so verächtlich und mühselig ist. Wie wir deren Viele hören, die da schreien, daß sie nicht Gott, sondern der Teufel geführt habe, da sie sich in den Ehestand begeben haben.

10. Daselbst begegnet auch den Regenten. Wenn sie sehen die ungezähmte Bosheit der Unterthanen, dazu den Wucher, Geiz und Unzucht, und von Amts wegen solches gern bessern und strafen wollten, müssen sie damit Haß und Feindschaft der Leute auf sich laden. Da kommt es denn, daß sie auch toben und fluchen, nicht weniger als im Ehestande auch geschieht. Und wenn solche Leute Papisten sind, die von Gottes Wort nichts wissen, so werden sie davon laufen und den Ehestand oder das Regiment verlassen, und etwa in eine Wüste sich begeben, daß sie daselbst Gott in Frieden und Ruhe dienen mögen, aber Alles ohne die Gnade und den Willen Gottes. Denn Gott hat allen Menschen sein Wort vorgehalten und geboten, daß ein Jeder seines Berufes warten und ihm darin dienen soll.

11. Derhalben, auf daß du Beschwerung und Mühseligkeit im Ehestande oder im Amte des weltlichen Regiments überwinden mögest, sollst du dich am allerersten befestigen, daß du ja vornehmlich an das Wort Gottes denkest, damit die Obrigkeit reichlich bestätiget ist, Röm. 12, 8. und Cap. 13, 1—4. Desselbigen gleichen auch im Ehestande, ein Knecht oder eine Magd, ein Schulmeister oder ein Schüler sind alle ihres Standes und des Willens Gottes gewiß. Nimm du nur das Wort an und thue rechtschaffene Frucht des Wortes, so wirst du sehen, daß bald das Kreuz und Anfechtung folgen wird. Darauf aber folgt das Gebet, dem Gebet folgt die Erlösung, der Erlösung folgt das Dankopfer, damit man Gott preiset. Also kannst du zugleich das Kreuz tragen und auch das Dankopfer opfern, welches die Mönche nicht thun wollen und können, es auch nicht thun, die nur äußerlichen Frieden und gute Tage suchen, und daß sie den Bauch füllen und in Wollust leben mögen.

12. Auf solche Weise beschreibt Mose die Legenden dieses heiligen Vaters Isaak nach dem Einen Hauptstück, nämlich, daß er sein Leben in vieler Trübsal zugebracht habe. Er schreibt nicht von Grübeleien, beschaulichem Leben und Heuchelei der selbst erwählten Werke, sondern er beschreibt an ihm große vortreffliche Tugenden des Glaubens und rechter Gottseligkeit, welche die Gottlosen und fleischlichen Menschen nicht sehen.

13. Und man soll diesen Gegensatz des Lebens, so die Heuchler führen, und des Lebens dieser Väter wohl und fleißig merken. Die Heuchler erwählen äußerliche Werke, die einen großen Schein haben, trinken keinen Wein und essen kein Fleisch, gehen und lassen den Kopf hängen, haben sonderliche Kleidung vor andern Leuten, und hüten sich vor der Mühe und Arbeit im Haus- und Weltregiment. Dagegen aber warten die Väter der Haushaltungen, gehen mit ihren Weibern und Kindern, mit dem Gesinde und dem Vieh um. Da ist gar kein äußerlicher Schein der Religion oder Geistlichkeit, sondern nur ein großer Sack des Lebens der gemeinen Haushaltung. Die Heuchler ziehen vergoldete Kleider an, damit sie vor der Welt scheinen, wenn sie sich des gemeinen Lebens und der Gemeinschaft anderer Leute entziehen. Aber unter dem groben Sack der gemeinen Haushaltung, darin die Väter gelebet haben, glänzt Sonne, Mond und Sterne, das ist, die allertrefflichsten Tugenden. Dagegen unter dem sehr herrlichen Kleide der Heuchler, welches äußerlich so einen großen Schein hat, liegen verborgen die heuchlerischen Greuel des Unglaubens, geistlicher Hoffart, Neides, Unreinigkeit, und hat doch gleichwohl den Namen und Schein eines geistlichen Lebens.

\subsection{I}

14. Lasset uns aber die geistlichen Augen auftun, und von ihrer geistlichen Schalkheit urtheilen nach der Regel und Vorbild dieser Historie. Denn Mose gedenkt keines Fastens oder Wachens und dergleichen. Das hört du allein, daß Isaak von Hebron nach Gerar zieht, daß er Hunger leidet, sucht einen Ort, da er wohnen möge. Aber siehe, welch ein großer Glaube da ist, welch eine unschätzbare Geduld! Welch eine unglaubliche Langmüthigkeit, Gütigkeit und Freundlichkeit ist an ihm gesehen und gespüret worden! Welche Tugenden alle scheinen wie Sonne und Mond. Es wäre zwar mir unmöglich, solchen großen Gehorsam zu leisten. Wenn ich zwei Söhne und ein Weib hätte, dazu Hausgesinde, und so eine große Menge Viehes und Knechte, und hätte nicht einen Fuß breit Landes, da ich einen Fuß möchte hinsetzen, was sollte ich da wohl thun? Ich würde wahrlich davon laufen und das Regiment und die Haushaltung fahren lassen. Denn wo Einer also in der Irre und auf das Ungewisse gehen und sitzen soll, und bleibt doch gleichwohl noch bei Weib und Kindern, und bei seinem Hausgesinde und Vieh, dazu gehört ein wunderbar großer Glaube, der aus Steinen könnte Brod machen. Das sehen aber die Gottlosen nicht.

15. Es ist aber wahrlich eine schwere Anfechtung gewesen, daß er täglich an einen andern Ort hat ziehen müssen und immer eine neue Herberge suchen. Denn das heißt in der Welt wohnen, und doch in der Welt elend und fremd sein. Und Isaak hätte wohl mit Christo das Büchlein singen mögen: „Des Menschen Sohn hat nicht, da er sein Haupt hinlege“, Luc. 9, 58. Derhalben kann es nicht bald ein Jeder sehen, der nur schnell darüberhin liefet, wie so ein großer Glaube an dem Patriarchen Isaak gerühmt wird; und wenn man mit dem Glauben aller Papisten Kappen, Orden und Werke vergleichen sollte, müssen sie davor gar verschwinden. Denn es ist ihrer keiner, der eines Tages seines Bettes, seiner Wohnung oder Speisen entrathen wollte. Ei, wie ist mir aber das so eine schöne Religion! Des Patriarchen Isaak Stand aber ist so viel desto härter, weil er nicht allein ist, sondern ist an sein Weib, Kinder und Gesinde gebunden, welche er in fremdem Lande mit Nahrung hat versorgen müssen. Da ist wahrlich sein Glaube fast hart angefochten und geprüft worden.

16. Darum lasset uns lernen, daß wir die Historien der Patriarchen groß machen und hoch preisen wider die abenteuerlichen Werke der Heuchler und Mönche, welche sagen: Es sei ein Leben der Wollust, wenn man Weib und Kinder hat, und darum könne es kein geistlich oder heilig Leben sein. Denn siehe doch, ist nicht dem Isaak allerlei Ursache der Freude und leiblicher Wollust abgeschnitten und genommen, da er alle Stunden dessen ungewis gewesen ist, wo er seinen Fuß hat hinsetzen sollen? Dem wird kein Mönch nachfolgen, ja, sie führen in dieser Welt ein gut geruhig Leben und das voller Wollust ist, sie haben ihre auserlesenen Häuser, sind unter dem Schutz und Schirm des Papstes und sind Herren über die ganze Welt.

17. Darum macht diese Pilgerschaft alle Religion und vermeinten Gottesdienste der Papisten zu Schanden, und leuchtet darin der höchste Gottesdienst, welcher ist der rechte Glaube an Gott. Er zieht umher in fremden Landen, und hat keine gewisse Herberge, hat keine Weide und nicht ein Tröpflein Wassers. Ja, wo soll ich es nehmen? würde ein Anderer sagen. Antwort: „Ich glaube an Gott den Vater“ 2c. Denn also hat Isaak gedacht: Gott wird mich mit Herberge versorgen, der ich also umher wallen muß; er wird meinem Viehe Weide geben, er wird meinem Hausgesinde Essen und Trinken schaffen. Da lebt Isaak also, daß er täglich Gott das höchste Opfer des Glaubens thut und lebt in der Welt ohne Welt und außerhalb der Welt.

18. Dieses Glaubens, welcher in den Vätern sehr groß ist gewesen, achten die Gottlosen nicht, sie sehen nur auf das Weiblein, Gesinde und Vieh; den Glauben sehen sie nicht, der dies alles erhält. Wenn zu unsern Zeiten Einer also sollte umher ziehen, der würde weder Knechte noch Mägde haben können, die ihm folgeten, da er also von einem Ort zum andern zöge. Aber der Glaube thut Wunderwerke, der macht, daß das Gesinde so großen Unfall, Trübsal und Armuth mit Geduld leidet, und sich in allen Dingen gehorsam erzeigt. Ich könnte zwar solche Knechte und die so fromm wären, nicht bekommen; denn ich habe so einen großen Glauben nicht. Und ist kein Heiliger im Neuen Testament, der diese großen herrlichen Tugenden, so an den Vätern gewesen, erreichen kann; sie sind ganz kleine und unmündige Kinder, wenn man sie mit diesen Altvätern vergleicht. Sie haben keine Heiligkeit mehr, denn nur ihren Ehelosen oder ehelosen Stand. Aber auch denselbigen übertrifft Isaak mit seiner Pilgerschaft, ja, sein Ehestand ist auch noch weit besser denn aller eheloser Stand der Mönche; wie das Licht der Sonne viel heller und klarer ist denn sonst alle anderen Lichter.

\quellseite{46}

19. Derhalben ist ihr ganzes Leben voller Wunderwerke. Denn sie leben einfach von der Hand GOttes und schier von nichts; sie hangen nur an der Güte und Verheißung GOttes, und halten fest an dieser Hoffnung, daß sie sagen können: Wenn ich heute nichts habe oder auch nicht sehe, davon ich leben kann, so werde ich es doch morgen oder übermorgen gewißlich bekommen. Solche Tugenden soll man in der Historie Isaaks fleißig betrachten, vornehmlich aber, daß er voll Glaubens gewesen zu GOtt, der ihn ernährt und erhalten hat.

20. Darnach hat er ohne Zweifel sein Herz auch gewöhnt zur Geduld, Langmuth und Wohlthat gegen den Nächsten. Er hat gelernt Unrecht leiden, es geschehe gleich ihm selber oder aber seinem Hausgesinde. Er hat die, so ihn gehaßt, nicht wieder gehaßt, auch nicht begehrt, sich an ihnen zu rächen. Denn die heiligen Väter sind gastfrei und barmherzige Leute gewesen, und wohlthätig gegen die Freunde sowohl, als gegen die Feinde. Denn es hat in ihnen ein trefflicher Glaube geleuchtet, welcher solche herrliche Tugenden zuwege bringt; von welchen die Sophisten ungesichtet und ungöttlich disputiren, nämlich, ob es Gebote oder Rathschläge seien, und schließen oder halten es dafür, daß es nur Rathschläge seien. Die Theologen der Sorbonne zu Paris sagen: Die Menschen können wohl selig werden, wiewohl sie einander nichts Gutes gethan haben. Und ein solch Leben ist allerdings der Mönche Leben, nämlich, nur eine äußerliche Heuchelei, welche diese Tugenden verachtet, als wären es bloße Rathschläge.

21. Aber diese Exempel lehren, daß es der höchste Gottesdienst sei, und der höchst nöthig ist, nämlich, daß man an der Verheißung und Vorsehung GOttes mit rechtem Glauben hangen bleibe. Denn er hat verheißen, daß er unser Vater sein wolle, und daß man auf ihn hoffe und Hülfe von ihm erwarte. Das thut kein Mönch, sofern er ein rechter Mönch ist. Nach dem Glauben aber folgt die Liebe gegen die Guten und Bösen, und daß man wohlthätig ist gegen die Dankbaren und Undankbaren. Diese Tugenden, sage ich, folgen aus dem Glauben, wie die guten Früchte von einem guten Baume folgen.

22. Es erzählt aber Mose eine treffliche Gelegenheit, dadurch der Glaube dieses sehr heiligen Mannes geprüft und bewährt wird. Denn es kommt eine Theurung in das Land, welche nicht allein Isaak mit seinen zwei Söhnen, mit seinem Weibe und ganzem Hausgesinde betrifft, sondern geht auch über das ganze Land bei Hebron. In diesem großen und gemeinen Unglück wird sein Glaube geprüft, und so hoch gelobt, daß wir uns über denselbigen alle verwundern, aber demselben können wir nicht leichtlich nachfolgen.

23. Auch soll man vornehmlich merken, daß gemeiniglich, wo das Wort im Flor und im Schwang geht, wenn der Herr die geistliche Speise reichlich gibt, alsbald die leibliche Theurung hereinfällt; denn der Teufel entzieht der Kirche die Nahrung und will sie durch Hunger tödten. Und dies muß wahrlich eine große Theurung gewesen sein, dieweil sie Mose vergleicht mit der vorigen Theurung, so zu Abrahams Zeiten gewesen war, welcher Gedächtniß dazumal schier verloschen war. Denn es waren ungefähr hundert Jahre oder mehr vergangen seit jener Theurung. Derhalben erinnert sie nun Mose solcher vorigen Theurung wieder und vergleicht sie mit dieser jetzigen gegenwärtigen Theurung, auf daß er damit anzeige, daß großer Kummer und Mangel aller Dinge gewesen sei.

24. Wie fromm es aber, daß solche heilige Menschen von GOtt nicht erlangt haben nothdürftige Nahrung für sich selbst und auch für andere Leute? Denn es hat nicht allein Abraham und Isaak, sondern auch andere hohe Patriarchen und Propheten, als Jacob, Joseph, Elia und Elisa, und endlich auch der heilige Paulus und andere Gottfelige mehr solchen gemeinen Schaden der Theurung zugleich mit andern Menschen tragen müssen. Antwort: GOtt sendet Theurung, Krieg, Pestilenz und dergleichen andere Plagen und Unglück mehr erstlich, daß die Gottfeligen dadurch versucht und geprüft werden, daß sie es lernen gewiß dafür halten, daß sie auch in der Zeit der Theurung sollen ernährt werden; wiewohl sie mancherlei Beschwerung erfahren, dazu auch unbekannte und ungewisse Wohnung suchen müssen.

25. Darnach läßt GOtt solche Plagen kommen zum Aergerniß und Strafe der Gottlosen. Denn wo das Wort reichlich offenbar wird, werden die Leute undankbar, ja, sie verfolgen und hassen das Wort. Die Anderen, mit denen es sich ansehen läßt, daß sie es annehmen, gerathen doch dahin, daß sie desselben bald überdrüssig werden und einen Ekel gewinnen an dieser losen Speise, verachten und plagen die Diener des Wortes. Mit welcher Verachtung und Haß sie GOtt zum Zorn reizen, daß er sagt: Willst du dich nicht speisen und sättigen lassen mit der Speise meines Wortes und mit dem geistlichen Leben, so will ich dir auch dies leibliche Leben nehmen und dich Hungers sterben lassen. Derhalben, wenn die Leute den unermesslichen Schatz des Wortes verachten, so verlieren sie darüber auch die Wohlfahrt und den Nutzen dieses gegenwärtigen zeitlichen Lebens; und wie die Gottfeligen in der Theurung erhalten werden, also werden die Gottlosen von wegen des Unglücks und Verzweiflung unterdrückt, und dieweil sie ohne Wort sind, sterben sie also dahin.

26. Daher nimmt denn der Teufel Ursache, das Wort zu lästern und demselben übel nachzureden; welcher an diesem Aergerniß Lust hat, und sich dasselbe zu Nutz machen kann, daß er damit die Herzen der Menschen vom Worte abwende. Denn was können wir denken, daß die Cananiter Anderes gedacht haben, welche, ehe Abraham zu ihnen gekommen, allerlei reichen Segen gehabt, aber nun mit ihm haben müssen Hunger leiden und die Beschwerung der Theurung tragen: was haben sie, sage ich, Anderes denken können, denn daß dieser Schadbäder alles dieses Unglücks und Jammers Ursache sein müßte? Dies ist die andere Ursache, nämlich, daß der Teufel und die Gottlosen Gelegenheit haben, das Evangelium zu lästern, auf daß sie immer ärger werden.

27. Wie wir zu unsern Zeiten auch hin und wieder dergleichen Klage hören und leiden müssen, nämlich, daß vor dieser Zeit alles Dinges übrig genug und Alles wohl gerathen, dazu allenthalben gut Glück gewesen; jetzt aber sei das Getreide so viel theurer und stehe allenthalben sehr übel. Wiewohl ich es nicht dafür halte, daß solches daher kommen sollte, daß jetzt der Frucht oder des Getreides weniger sein sollte; sondern daß es vielmehr vom Geiz und Bosheit der Menschen kommen müsse, welche mit allen Dingen Theurung machen nach ihrem eigenen Gefallen. Es ist aber doch nicht ein geringes Unglück, damit die Armen und auch die Diener des Worts hart gedrückt und beschweret werden. Die Anderen, so sehr reich sind, haben weniger Noth. Derhalben findet man nun deren Viele, die da wünschen und begehren, daß es wiederum in den alten vorigen Stand kommen möchte, und daß sie dabei solches Glück haben möchten, wie man zuvor gehabt; und setzen noch diese Lästerung dazu, daß sie sagen: Es sei aus dieser Lehre des Evangeliums nichts Gutes gekommen; und über das, so seien die Leute auch viel ärger geworden, denn sie vor der Zeit gewesen sind.

28. Also wird der Lehre des Evangeliums Schuld gegeben, daß sie alles Unglückes Ursache sein soll. Zur Zeit Augustini, da die Gothen Italien verwüsteten, wurde zu Rom alle Schuld auf die Apostel gelegt, nämlich, auf St. Petrus und Paulus. Denn die gottlosen Menschen thun aus den Augen hinweg und wollen nicht sehen auf die Sünde der Welt, und muß also das Wort GOttes, welches ganz rein und heilig ist, unbilliger Weise die Schuld tragen, daß davon alle Sünden sollen herkommen. Denn es lehrt ja nicht, daß man wuchern, geizen, Schwelgerei treiben, und andere Sünden, Schande und Trügerei, damit die Welt umgeht, üben solle; sondern streitet und ruft wider diese Sünden alle. Warum wird denn nun das Evangelium mit solchen greulichen Lästerungen beschweret? Antwort: Darum daß diese die Bosheit ist des Teufels, dadurch er sich freut das Evangelium zu lästern und wider dasselbe allerlei Schmachworte aufzubringen, woher er nur kann und mag.

29. Hilf GOtt, wie ist das so ein erschrecklich Ding, und ist zu beweinen, daß dieweil dies Licht des Evangeliums scheint, die Leute noch so überaus böse und muthwillig sein sollen! Du aber sollst dich aus dem Worte lehren und dessen berichten lassen, daß das Evangelium nicht erlaubt, Muthwillen zu treiben mit Wucher, Räuberei, Fressen und Saufen, oder auch mit Spielen, und dergleichen andern Sünden und Schaden mehr. Wo du aber diese Schuld auf das Evangelium legst, so wirfst du jetzt aus deinem eigenen Munde gerichtet und übermunden, daß du mußt vom Teufel besessen sein, welcher ein Stifter ist solcher Lästerung. Denn du siehst ja wohl, daß ja das Widerspiel gelehrt wird, und daß die, so der Lehre gehorchen, das Heil und Seligkeit erkennen und empfinden, ja, daß sie auch anfangen, ein gottseliges und christliches Leben zu führen. Und du kannst ja nicht leugnen, daß es die Wahrheit sei, und willst doch gleichwohl alle Sünde und Schande, so du an den Epicurern und Gottlosen siehst, der Lehre des Evangeliums zuschreiben. Das thust du aus lauter Bosheit des Teufels, auf daß du und alle, so der Wahrheit Feind sind, gerichtet werden.

\quellseite{47}

30. Augustinus erzählt viele greuliche Sünden und Laster der Römer, als Mord, Unzucht, Wucher und andere mehr. Sollten nun solche Sünden nicht billig von GOtt gestraft werden? Derhalben ist es gar ein teuflisch Ding und kann ohne greuliche Sünde nicht einmal gedacht werden, daß die Leute dieses alles vergessen können und die Ursache alles Unglücks und der Strafen nur auf das Wort legen wollen.

31. Dies ist aber alles uns zur Lehre geschrieben, daß wir uns an solcher Lästerung nicht ärgern sollten. Denn also haben zu allen Zeiten alle Gottseligen Theurung und Lästerung der Gottlosen leiden und tragen müssen, und je heiliger sie gewesen sind, je mehr Strafen und Unglück in der Welt gewesen ist; nicht, daß sie solches verursachet hätten und die Welt verwirreten, sondern weil die Welt das Licht des Evangeliums verachtet und der Abgötterei nachgefolgt hat; wie Elia dem Könige Ahab antwortete 1 Kön. 18, 18.: „Ich verwirre Israel nicht, sondern du und deines Vaters Haus, damit, daß ihr des HErrn Gebote verlassen habt, und wandelt Baalim nach“ etc.

32. Auf solche Weise wollen wir zu unserer Zeit auf die lästerlichen Klagen auch antworten und zu den Leuten also sagen: Da das Evangelium euch vorleuchtet und den Weg zur Seligkeit weist, da seid ihr undankbar, verachtet und verfolgt diese reine Lehre; darum kommen die Strafen über die Welt, daher kommt den Türken und allerlei Unglück.

33. Unterdeß aber werden die Gottseligen auch mit gestraft, wie Isaac Mangel und Theurung zu seiner Zeit auch gelitten und getragen hat; er war fremd und elend, und dachte, er müßte das gegenwärtige Unglück mit Geduld tragen. Man soll aber fleißig merken, was die endliche Ursache sei solcher Trübsale, so die Gottseligen leiden müssen. Es trifft zwar das gemeine Unglück die Heiligen und Propheten auch mit; aber das geschieht ihnen nicht zur Strafe oder aus Zorn, wie es den Gottlosen und Undankbaren geschieht, sondern ihnen zur Seligkeit, und daß sie dadurch bewähret und geprüfet werden im Glauben, in der Liebe und Geduld, auf daß die Gottseligen lernen in der Haushaltung die Hand Gottes dulden und ertragen. Denn GOtt hat ihnen verheißen, daß er sie in der Theurung ernähren wolle. Wie Ps. 37, 25. geschrieben steht: „Ich habe nie gesehen den Gerechten verlassen“; item V. 19.: „In der Theurung werden sie genug haben.“ Diese Verheißung bestätigt GOtt mit solchen Exempeln der Heiligen, und mit solcher häuslichen Trübsal und Noth unterrichtet er die Gottseligen im Wort, im Glauben, in geistlicher Demuth, in der Liebe und andern Tugenden.

34. Die Gottlosen aber werden geplagt zur Strafe und Aergerniß, daß sie verhärtet und immer ärger werden; denn durch das Gute werden sie nicht gebessert, durch das Böse aber werden sie schlimmer. Darum richtet man damit nichts aus und hilft bei ihnen gar nichts, man hinge ihnen gleich süß oder sauer, man rede freundlich oder hart mit ihnen, wie Christus auch klagt Matth. 11, 17.

35. Nun siehe und betrachte diese geringe und verächtliche Gestalt der Haushaltung Isaaks, und halte sie gegen alle Werke der Mönche, des Antonius oder Hilarion. Denn was ist alles Fasten und harte Kasteiung der Mönche gegen diese einige Theurung? Es scheint, als sei es eine kleine und geringe Anfechtung, wenn du sie nur äußerlich und obenhin ansiehst; aber gesetzt den Fall, daß du in Isaaks Hause wärest, welcher Weib und Kinder, item, eine große Menge Knechte und Vieh hat, in solchem großen Mangel, da es ihm an Speise und Futter gebricht hat. Gegen den Fall, daß du hören mußt, wie das Gesinde klagt, welches begehrt und fordert, daß man ihm Brod oder dem Viehe Futter schaffen und geben soll; und wo du daselbst antworten wirst: Ich habe weder Brod noch Futter, ich will davon laufen und das Gesinde sitzen lassen, so liegt dir GOttes Gebot im Wege, der von dir haben will, daß du in deinem Berufe, und an dem Orte, dahin er dich bestellt hat, bleiben sollst.

36. Ja, sagst du, womit soll ich denn das Gesinde erhalten und das Vieh füttern? Antwort: Du sollst und mußt verstehen lernen, was dies für ein Fasten sei, da ein armer Hausvater nicht allein für sich muß Hunger leiden, sondern muß auch vor seinen Augen sehen, daß sein Weib und lieben Kinderlein schier müssen Hungers sterben. Ja, thue das auch noch dazu, daß Isaac an dem Orte elend und fremd war, und keine gewisse Herberge gehabt hat. Er wäre ohne Zweifel wohl gern zu Hebron geblieben: dieweil aber die Leute an dem Ort selbst auch mit der Theurung sind geplagt worden, haben sie ihn heißen wegziehen, darum daß sie ihn mit einem solchen großen Gesinde nicht haben ernähren können, und weil sie selbst kaum so viel gehabt, davon sie sich und ihr Vieh erhalten möchten.

37. Derhalben muß er weichen und wird das Unglück nun noch einmal so groß; denn zu der Theurung kommt nun auch, daß er muß in das Elend ziehen. Da wird wahrlich Isaac das Herz im Leibe geschlagen und geschrieen haben, weil er ja sich der Hausorge nicht hat begeben oder dieselbige fahren lassen können von wegen des Gebotes Gottes, 1 Tim. 5, 8.: „So jemand die Seinen, sonderlich seine Hausgenossen, nicht versorget, der hat den Glauben verleugnet, und ist ärger denn ein Heide.“ Und hat doch gleichwohl nichts, davon er mit ihnen hätte leben und sich erhalten mögen.

38. Gib du mir jetzt ein Exempel, daß dieser Geduld möchte gleich sein! Bei den Carthäuser-Mönchen wirst du zwar keines finden oder auch bei den andern Heuchlern; die würden kurz und gut sagen, es wäre ein unmöglich Ding. Warum sagen sie aber denn, daß es ein verächtlich und unnützig Ding sein soll, daß ein Ehemann bei seinem Weibe liegt und sein Hausgesinde regiert? Ja, siehe uns selbst auch an, wie wir uns in solcher Noth und Gefahr pflegen zu halten; wie verzagt wir in unserm Herzen sind, wenn wir nicht Getreide oder Geld vollauf haben: hilf GOtt, wie leichtlich fallen wir dann dahin in Verzweiflung oder Ungeduld. Da ist nun Isaac in der höchsten Noth und äußersten Gefahr gewesen, doch läßt er nicht davon und läßt gleichwohl den Glauben und Hoffnung nicht fahren.

39. Und daraus kann man nun sehen, was das sei, daß Einer ein frommer gottseliger Ehemann ist. Isaac kämpft nicht allein mit der Theurung und Armuth, sondern ringt auch mit Verzweiflung und Ungeduld; aber da leuchtet hervor sein starker und unüberwindlicher Glaube. Denn er bleibt ein Hausvater, vermißt sein Weib nicht, verläßt sein Hausgesinde nicht, sondern zieht mit ihnen an einen andern Ort, daß er daselbst seine Nahrung suchen möge. Wiewohl er nun auch nicht ohne Gefahr und Schaden gewesen ist, richtet er sich doch gleichwohl auf und hält sich an den Trost, daß er gedacht hat: GOtt hat mich gewiesen, daß ich sollte ein Ehemann werden, und hab mir zwei Söhne gegeben; ich weiß, daß solches sein Werk ist, darum wird er mich nicht verlassen.

40. Also lehren nun diese häuslichen Werke nicht, daß dabei eitel Wollust und fleischliche Freude sei, sondern es werden uns darin vorgehalten Exempel eines heiligen Kampfes und Streites, so diese Väter geführt haben wider den Unglauben, Mißtrauen, dazu wider die Verachtung und Ungeduld, wo man wider GOtt murret. Von welchem schweren Kampfe die Mönche gar nichts wissen und suchen mit ihrer Heiligkeit nur einen vollen Bauch; sie wollen und können auch GOtt in der Noth und Widerwärtigkeit nicht vertrauen.

41. Alhier aber ist ein trefflich Exempel eines solchen Glaubens, der im Kampfe liegt und zu GOtt schreit; denn wo das Wort ist, da folgt auch die Anrufung. Und wie groß dieser Kampf gewesen sei, zeigt die Reise an, so er vornimmt, da er durch Palästina nach Egypten zieht. Darum hat sich GOtt sein erbarmt und stärkt ihn mit einem neuen Wort. Es pflegt aber Gottes Wort nicht zu kommen, vornehmlich aufs neue, denn nur zu denen, so sehr betrübt sind und die sonderlichen Trostes bedürfen. Es kommt zu denen nicht, die da satt und überdrüssig sind; denn es hat bei ihnen keine Statt.

42. Ueber das werden ihm die zwei Söhne und Rebekka, sein Weib, den Schmerz auch etwas genehret haben. Esau und Jakob waren damals ungefähr achtzehn Jahre alt, welchen ohne Zweifel die Reise in das unbekannte Land und in solcher großen Beschwerung, so auf ihnen gelegen, nicht wird gefallen haben, und werden zum Vater gesagt haben: Lieber Vater, wo willst du hinziehen? Also wird Rebekka auch gesagt haben: Mein lieber Isaac, siehe, was du thust. Wie das Weib Hiobs ihn sein Unglück vorrückt und ihn darüber lästert; denn es sind wohl heilige Weiber, sie sind aber doch nicht ohne Anfechtung. Darum hat er ihre Herzen stärken müssen, daß sie ein wenig harren und warten sollten, und dies Kreuz mit Geduld tragen. Und ist diese Theurung eine Ursache gewesen, daher viel treffliche schöne Predigten gekommen sind, welche Isaac damals gethan hat, sein Gesinde, dazu auch sein Weib und Kinder in Glauben zu erhalten.

\quellseite{48}

43. Aber gleich im selbigen Augenblick, da die Kinder, das Weib und Hausgesinde schier der Noth und Trübsal halben verzagt worden, und nun angefangen zu murren, oder aber gedenken davon zu laufen, da thut sich der Himmel auf und kommt eine neue und herrliche Verheißung. Da er gedenkt, durch Palästina nach Egypten zu ziehen und in Gottes Namen in das Elend gehen will, und Gott dasselbe auch also geschehen läßt, daß er sich rüstet auf die Reise; da die Schnur oder Strick so hart gespannt ist, daß sie jetzt schier brechen will: da kommt ein mächtiger und trefflicher Trost vom Himmel herab, daß Gott zu ihm spricht: „Zeuch nicht hinab in Egypten.“

\subsection{II}

\begin{quote}
\textbf{V. 2—5. Da erschien ihm der HErr, und sprach: Zeuch nicht hinab in Egypten, sondern bleibe in dem Lande, das ich dir sage. Sei ein Fremdling in diesem Lande, und ich will mit dir sein, und dich segnen; denn dir und deinem Samen will ich alle diese Länder geben, und will meinen Eid bestätigen, den ich deinem Vater Abraham geschworen habe, und will deinem Samen mehren, wie die Sterne am Himmel, und will deinem Samen alle diese Länder geben. Und durch deinen Samen sollen alle Völker auf Erden gesegnet werden; darum, daß Abraham meiner Stimme gehorsam gewesen ist, und hat gehalten meine Rechte, meine Gebote, meine Weise und meine Gesetze.}
\end{quote}

44. Dies muß wahrlich eine große Noth und Trübsal gewesen sein; welches auch daraus gespüret wird, weil Isaac sich nicht hat können genügen lassen an der vorigen Verheißung, welche er von seinem Vater Abraham und dem Patriarchen Sem gehabt: „Der Größere wird dem Kleinern dienen“ etc. Denn sein Herz schlägt ihm noch im Leibe und der Glaube ist in Angst, daß ihn Gott muß mit einem neuen Trost aufrichten.

45. Derhalben ist dies Exempel nicht allein uns, sondern auch den Aposteln und Heiligen selbst vorgehalten, daß wir lernen unser Herz zur Geduld stärken und in keinerlei Anfechtung wider Gott murren, in welchem Stande wir auch leben sollen. Es ist wahrlich der Stand des Kirchen-, des Haus- oder Weltregiments fast erbärmlich, vornehmlich, wenn du willst fromm und gottselig sein; denn „alle, die gottselig leben wollen in Christo Jesu, müssen Verfolgung leiden“, 2 Tim. 3, 12. In der Haushaltung muß man stets mit dem Gesinde zanken und hadern. Im weltlichen Regiment zankt die Ehre und Gewalt hervor, und darum läßt es sich ansehen, daß es ein wohlhäusig Leben sei: aber wirst du ein frommer Richter, ein gottseliger Regent sein, so wirst du empfinden, welch eine große Last dieser äußerliche Schein der Ehre, Herrlichkeit, und des großen Gutes, so an den Regenten äußerlich gesehen wird, auf sich liegen habe. Wirst du ein frommer gottseliger Prediger sein, auch bei gottseligen und dankbaren Zuhörern, so wirst du doch von den Letzern große Undankbarkeit leiden müssen, desgleichen auch Haß und unendliche Stücke des Teufels, damit er dir heimlich nachstellen wird.

46. Derhalben laß dich deinen Stand nicht gereuen, davon du weißt, daß er dir von Gott zugeeignet und bescheret ist, er sei gleich so gering und verachtet, als er immer sein kann: sondern rühme dich Gottes des HErrn, der dich berufen hat, der dich regieret, und daß du auch das Wort gegenwärtig hast. Denn „ist Gott für uns“, wie Paulus sagt Röm. 8, 31., „wer mag wider uns sein?“ Wenn irgend eine Bitterkeit oder Beschwerung vorfallen wird, so gieße du guten Wein darauf und vertreibe damit die Bitterkeit, welcher Wein Gottes Wort ist, auf daß dadurch dein Glaube gemehret oder gestärkt werde, und daß du sagen könnest: Ich bin gewiß, daß der Stand, darin ich jetzt lebe, gottselig und christlich ist, und ich will darin bleiben und verharren, daß ich Gott loben und preisen, und andere Leute lehren möge.

47. Wir sehen aber, daß allhier wiederholet und zusammengefaßt werden schier alle Verheißungen, welche Gott dem Abraham an mancherlei Orten gegeben hatte: dieselbigen alle werden hier in eine Summa gebracht. Denn Gott hat mit Abraham oftmals geredet, mit Isaac aber kaum zwei- oder dreimal, und daran ist es auch genug. Denn er bestätiget hier in einer Summa alle Verheißungen, auf daß der heilige Patriarch durch des Teufels Anreizung nicht möchte anfangen zu zweifeln an Gottes Willen. Denn der Satan hört nicht auf, auch die allerheiligsten Männer und die ganz vollkommenen sind, mit seinen feurigen und vergifteten Pfeilen anzufechten und ihnen damit hart zuzusetzen.

48. Diese Verheißung aber hat zwei Theile: der erste ist zeitlich, nämlich, vom Erbe oder Besitz dieser Länder. „Ich will mit dir sein“, spricht Gott, fürchte dich nicht, du wirst in das Elend getrieben und wirst mit Theurung beschweret werden, es wird aber keine Noth oder Gefahr haben; du sollst nicht Hungers sterben, es soll dir auch der Feind in fremden Landen keinen Schaden thun. Also tröstet ihn Gott in der sehr schweren Anfechtung, die Herzen seines Weibes, der Kinder und des Gesindes im Hause damit zu stärken, welchen Isaac diesen Trost wird vorgehalten und gesagt haben: Gott der HErr ist mit mir, und hat mir verheißen, mich aus aller Noth und Unglück zu erlösen; darum werde ich Frieden haben, auch Schutz und Schirm, und mein täglich Brod mitten unter den Feinden und in der Theurung. Denn dieser Verheißung hat er fest geglaubt, wiewohl er zuvor angefochten und gereizt worden ist, daß er zweifeln sollte. Hernach aber werden wir hören, wie er sich gefürchtet habe vor dem Tode, und wie er verleugnet habe, daß er ein Weib hätte. Denn die heiligen Männer haben Fleisch und Blut gehabt wie wir auch; darum haben sie bisweilen auch Schwachheit des Glaubens gefühlt wie wir.

49. Er wiederholet aber, das er zuvor gesagt hat: „Ich will dir und deinem Samen alle diese Länder geben“, und läßt hier das Wörtlein „dir“ aus; denn er sagt nur: „Ich will deinem Samen alle diese Länder geben“, und legt die Verheißung aus. Denn es hat weder Abraham, noch Isaac oder Jakob irgend einen Theil dieses Landes inne gehabt oder besessen; wie droben (Cap. 23, V. 16, 17.) gesagt ist, daß Abraham um sein Geld einen Acker gekauft habe zum Begräbniß der Sarah. Sonst hat ihm Gott auch nicht eines Fußes breit in dem Lande gegeben, so ihm verheißen war, bis die Nachkommen da hinein sind geführt worden. Und die erste Verheißung dieses Landes ist dem Abraham geschehen um des gebenedeieten Samens willen, durch welchen alle Völker sollten gesegnet werden, auf daß ein gewisses und bestimmtes Land und Ort wäre, daher man den HErrn Christum erwarten sollte, der die ganze Welt und alle Völker segnen würde: damit das Evangelium nicht ohne gewisses Zeugniß in die Welt ausginge, sondern die ganze Welt dessen vergewissert möchte werden, daß dies der rechte Heiland wäre, der den rechten Segen bringen würde, nämlich, der nach der Verheißung in diesem Lande geboren wäre.

50. Es wird aber doch dadurch, daß zuvor Abraham und der Same zusammen gesetzt worden, den Gläubigen angezeigt das ewige Leben und Auferstehung der Todten, nämlich, daß Abraham, Isaac und Jakob dieses Landes Erben sind, wiewohl sie darin auch nicht eines Fußes breit inne gehabt. Denn ob sie wohl gestorben sind, leben sie doch gleichwohl noch, und darum gehört ihnen auch dies Erbe, sintemal Abraham nicht gestorben ist, sondern lebt. Darum wird das auch angezeigt, da er sagt: „dir und deinem Samen“, nämlich, daß die Nachkommen das Land nicht besitzen könnten, wo die Väter die Verheißung nicht empfangen hätten; und im Glauben der Väter haben die Nachkommen den Besitz des Landes erlangt.

51. Der andere Theil der Verheißung ist geistlich, um welches willen die leibliche Verheißung ist gegeben worden, wie eben gesagt. Es sind aber eben dieselbigen Worte und auch derselbe Verstand, wie droben im 22. Cap. V. 18., und was man hier sagen kann, diese Verheißung auszulegen und zu erklären, ist alles droben gesagt worden, wie auch daselbst die Bedeutung des Wortes „Segen“ erklärt ist. Denn dies Wort gehört zur vierten Ordnung der Zeitwörter (wie man nach der Grammatik redet), welche Bedeutung die Lateiner mit solcher Art und Weise zu reden pflegen zu geben, daß sie sagen: gloriari in Domino, laudari in Domino, sich im HErrn rühmen, im HErrn gelobt werden. Wenn ich mich selbst im HErrn segne, so mache ich nicht, daß du Andere segnest, sondern ich erhebe und lobe mich selbst, nicht in mir, sondern im Samen Abrahams; denn ich bin verdammt und verflucht in Adam, bin unter die Gewalt des Teufels geraten, der mich gefangen hält unter dem Joch der Sünde, des Todes und Verdammniß; da kann sich niemand segnen oder rühmen, sondern ist eitel Seufzen, Schreien und großer unendlicher Jammer mit uns.

\quellseite{49}

52. Wenn aber der Same Abrahams kommen wird, alsdann werden alle Völker andere Menschen werden. Sie werden zwar an sich selber nichts haben, davon sie rühmen könnten, sondern werden von sich selbst sagen müssen, daß sie verfluchte und elende Leute sind: sie sollen aber doch gleichwohl herrschen, fröhlich und seliglich siegen, nicht in sich selber, sondern in diesem Samen. Und also geschieht es heutiges Tages von allen, die an Christum glauben. Also ist die Schrift erfüllt worden. Denn also halten und sagen alle Christen: Wenn ich schon in mir selber verdammt bin durch Adam, so bin ich doch gerecht und heilig in dem HErrn Jesu Christo, dem Sohne Gottes, der ist mein Leben, Gerechtigkeit, Heiligung und Erlösung. Also segne ich mich selbst und eines andern Dinges willen und durch einen Andern, nicht mit meinem, sondern mit dem Segen des Samens Abrahams. Und wenn ich denselbigen durch das Wort ergreife und an ihn glaube, so kann ich rühmen, wie im 34. Psalm V. 2. geschrieben steht: „Ich will den HErrn loben, und meine Seele soll sich rühmen des HErrn“, nämlich, daß ich wider den Tod das Leben, wider die Verdammniß die Seligkeit und wider den Tyrannen und meinen Feind den Teufel Gott zum Vater habe.

53. Dies sind große und herrliche Dinge, und man kann sie nicht genugsam wiederholen und den Leuten einprägen; denn es sind Worte des Trostes und des ewigen Lebens: wie Christus eben dasselbe seinen Jüngern auch oftmals und fleißig einprägt, als, Joh. 14, 1.: „Euer Herz erschrecke nicht. Glaubet ihr an Gott, so glaubet ihr auch an mich“; item V. 19.: „Ich lebe, und ihr sollt auch leben“; item Joh. 16, 33.: „Seid getrost, ich habe die Welt“ und den Teufel „überwunden.“ Wodurch aber? Durch meinen Sieg, der euer ist. Dies ist die höchste und allerobernehmste Lehre in der ganzen heiligen Schrift, nämlich, von der Verheißung und Glauben Abrahams, wie dieselbige Christus auch rühmt Joh. 8, 56.: „Abraham ward froh, daß er meinen Tag sehen sollte; und er sahe ihn, und freuete sich.“ Denn er hat wohl verstanden, daß dies der Segen sein würde, dadurch alle Völker in der ganzen Welt sollten gesegnet werden, daß dieser Ruhm nicht allein die Nachkommen des Fleisches angehe, sondern alle Völker über die ganze Welt. Das ist ein menschenwerthes Wort, welches die Propheten und Apostel fleißig gelehrt haben.

54. Was mehr in diesen Worten begriffen ist, das ist alles droben ausgelegt worden, nämlich, die Auferstehung der Todten, der Sieg oder Überwindung der Sünde, des Todes, des Teufels und das ewige Reich. Item, daß dieser Same Gottes Sohn und Mensch sei. Denn wo er nicht Gott wäre, so könnte er uns verfluchte Menschen nicht segnen: und wenn er nicht Mensch wäre, so würde der Segen nicht zu uns kommen.

55. Es ist das Stück auch erklärt worden, da im Text steht: „Darum, daß Abraham meiner Stimme gehorsam gewesen ist“ etc., droben im 22. Cap. V. 16. Denn dieser Segen ist dem Abraham nicht gegeben worden, daß er dadurch sollte gerecht werden; sondern da er schon gerecht war durch den Glauben, hat er diesen Segen empfangen als eine sehr treffliche Belohnung. Er ist gerecht, gehorsam und heilig. Und weil er so gehorsam ist, darum wird er so hoch erhoben werden, daß Christus von seinem Samen seine Menschheit hat angenommen, wie Röm. 9, 5. steht. Das ist wahrlich eine große Ehre, daß der von Abrahams Samen herkommt, welcher Gottes Sohn ist, der die Hölle zerstört, den Tod überwunden, das Gesetz aufhebt und das ewige Leben wiederbringt.

56. Es werden aber im Text vier Worte gesetzt, dieselbigen wollen wir erklären. Denn die Juden wollen daraus schließen, daß auch vor der Zeit Moses das Gesetz Moses gewesen sei; und ist zwar schier alles, was in Mose steht, aus den Vätern und ihren Büchern genommen. Es ist aber doch ein Unterschied darunter. Das Wort mischmar, auf Lateinisch observantia, das wir gegeben haben, Rechte, ist ein gemeines Wort, wird oftmals im dritten Buch Mose gebraucht. Denn von den Leviten wird gesagt, daß sie die Rechte des HErrn halten, wie wir auf Deutsch sagen: „Warte, was ich dich heißen werde“: da befehle ich Einem, daß er bereit sei und die Befehle erwarte. Also stunden die Leviten in der Hütte, und warteten, was sie der Hohepriester heißen würde. Auf dieselbige Weise wird hier auch gesagt: Abraham wartete auf mich, that, was ich ihn hieß. Also, da ihm geboten ward, seinen Sohn zu schlachten, war er gehorsam. Das hebräische Wort mizva heißt auf Lateinisch eigentlich praecepta, Gebote, und das man für und für halten soll, wie die Zehen Gebote sind.

57. Das Wort chok heißt die Ordnung der Weise, die zum Gottesdienst gehört, wenn Gott etwas zum Gottesdienst geordnet. Das ist kein Gebot, sondern ein eingesetzter Gottesdienst. Als, im 2. Psalm V. 7.: Praedicabo statutum etc.: „Ich will von einer solchen Weise predigen“ etc., das ist, von einem solchen Gottesdienste; er will sagen, Gott wird durch seinen Sohn einen andern Gottesdienst anrichten und einsetzen. Torah heißt, was zu guten Sitten dient. Dazu gehört auch mischmar, Lateinisch jura, Rechte. Also hat Abraham immer gewartet, was ihn der HErr heißen würde, auch was gute Sitten belangt hat: er hat zugleich die Zehen Gebote, die Weise des Sabbathes und das Gesetz der Beschneidung gehalten.

58. Dies ist aber der Unterschied unter dem Gesetze Moses und der Väter: Mose hat ein gewiß Gebot gehabt, das Levitische Priesterthum einzusetzen, welches auch ein gewiß Volk, Person, Stätte und eine bestimmte Zeit bis auf Christum anging; das ist das vornehmste Stück im Gesetze Moses. Darnach hat er auch das Königreich und die zwölf Fürsten im weltlichen Regiment verordnet. Die Zehen Gebote aber und andere Gesetze, so von den Vätern herkommen, sind nicht eigentlich Moses eigene Gesetze. Allein wo die Ceremonien belanget, so auf gewisse Personen gestellt waren, das ist Moses Gesetz. Da aber die Zeit vollendet war, so zu diesen Dingen allen bestimmt ist gewesen, da haben auch die Opfer und alle Ordnungen Moses ihr Ende, daß sie nicht mehr gelten: da hat das Priesterthum mit seinen Opfern und auch das Königreich nicht mehr gegolten; denn es ist das Ende Moses da gewesen. Darauf ist aber gefolgt das Königreich und Priesterthum Christi ohne Stätte und Zeit, oder Person, ohne alle äußerliche Geberden, wie Christus sagt Luc. 17, 20.

59. Endlich müssen wir auch etwas sagen von der göttlichen Erscheinung. Es ist in der heiligen Schrift zweierlei Erscheinung, nämlich im Traum und im Gesicht oder sichtbarer Form und Weise. Es wird auch von Gott gesagt, daß er erschienen sei, wenn er geredet hat durch Sem oder Eber, die dazumal gelebt haben und Hohepriester gewesen sind. Dieweil aber allhier nichts dazu gesetzt, sondern nur gesagt wird: „Da erschien ihm der HErr“, so wird hier verstanden, daß die Erscheinung nicht im Traume, sondern in sichtbarer Gestalt geschehen sei, das ist, es ist ein Engel gewesen, der ihm in der Gestalt eines Menschen erschienen ist; wie zu der Jungfrau Maria Luc. 1, 26. 27., zu St. Petro Ap. Gesch. 12, 7., und zu Abraham die Engel gekommen sind, wie droben im 18. Cap. V. 1. von Abraham steht. Und ist wahrlich diese Erscheinung dem Isaac eine große treffliche Ehre gewesen. Denn wiewohl er seines Vaters Predigen gehört hat von den Verheißungen, so hat doch gleichwohl sein Herz in der Noth und Trübsal etwas gezittert und ist erschrocken gewesen. Derhalben kommt nun Gott mit einer neuen Bestätigung der Verheißungen.

60. Und es ist wahrlich ein großes Ding, daß Gott einem Menschen erscheint und seine Verheißung also sonderlich auf einen Menschen insonderheit richtet. Darum wird in dem Theil von Vielen gehalten, daß die heiligen Väter viel seliger gewesen sind als wir, sintemal sie so gewissen und sonderlichen Trost und Erscheinung von Gott gehabt durch den Dienst der Engel. Denn es möchte jetzt Einer sagen: Wenn Gott mir auch etwa in einer menschlichen Gestalt erschiene, o wie sollte das meinem Herzen so eine große Freude bringen! Ja, alsdann wollte ich mich nicht mehr fürchten, allerlei Gefahr und Widerwärtigkeit um Gottes willen zu leiden. Aber so gut wird mir es nicht. Ich höre nur die Predigten, lese die Schrift und gebrauche die Sacramente; ich habe keine Erscheinung der Engel.

\quellseite{50}

61. Darauf antworte ich also: Du hast dich nicht zu beklagen, daß du weniger heimgesucht seiest, als Abraham oder Isaak gewesen sind. Du hast auch Erscheinungen, und die etlichermaßen stärker und klarer sind, und hast deren auch mehr, als sie gehabt haben, wenn du nur die Augen und das Herz könnest aufthun und dieselbigen annehmen. Du hast die heilige Taufe, das Abendmahl des Herrn, da Brod und Wein die Gestalt, Figur und Formen sind, darin und unter welchen GOtt gegenwärtig dir in die Ohren, Augen und Herz redet und wirkt. Darnach hast du das Predigtamt, und solche Lehrer, durch welche GOtt mit dir redet; du hast das Amt und den Gebrauch der Schlüssel, dadurch er dich absolvirt und tröstet. Er sagt zu dir auch, Jes. 41, 10.: „Fürchte dich nicht, ich bin mit dir“; er erscheint dir in der Taufe, und ist es selber, der dich tauft und anredet; er sagt nicht allein: „Ich bin mit dir“, sondern: „Ich vergebe dir deine Sünde, ich bringe dir Hülfe, dadurch dir vom Tode geholfen wird, dazu Erlösung von allem Schrecken und Gewalt des Teufels und der Hölle; und ich bin nicht allein mit dir, sondern auch alle Engel neben mir. Was wolltest du mehr begehren? Es ist alles voll göttlicher Erscheinung und Gespräche, so er mit dir hält.

62. Da sollten wir aber nun klagen und seufzen über unser vergiftet Fleisch, das in Sünden gar ersoffen ist, welches uns nicht zuläßt, daß wir glauben und solche große Wohlthaten annehmen, und überdies uns heißt disputiren und noch in Zweifel ziehen, ob dies alles also wahr sei. Ich rede jetzt von uns, die wir rechte Christen sind, die wir dies selbst lehren und glauben. Ich glaube nicht allein an Christum, sondern ich weiß auch, daß er zur Rechten des Vaters sitzt, daß er unser Mittler sei und uns vertreten möge. Ich weiß, daß das Brod und Wein im Abendmahl des Herrn der wahre Leib und Blut Christi ist, daß das Wort des Pfarrherrn, wenn er prediget oder absolviret, Gottes Wort ist; und ist doch gleichwohl das Fleisch mit Zweifel beschweret, daß es dies nicht glaubt. Das ist ein großer Jammer und schwerer als der Tod selbst; ja, der Tod ist darum bitter und schwer, dieweil wir nicht glauben und des Hindernisses wollen unseres Fleisches. Sonst würde, wenn wir glaubten, die Trübsal Freude und der Tod ein Schlaf sein.

63. Dies Unglück, so uns von der Erbsünde durch Adam angehängt ist worden, sollten wir beklagen, und bitten, daß Gott den Glauben in uns mehren und stärken wolle, und uns unter dem Himmel der Vergebung der Sünden erhalten, wie uns Christus hat beten gelehrt: „Vergib uns unsere Schuld.“ Denn es zumal ein groß Unglück, daß wir mit so vielen und großen herrlichen Erscheinungen, Gesprächen und mancherlei Gestalt, darin sich Gott uns offenbaret, überladen und überschüttet sind, und doch gleichwohl nicht glauben; und mögen wir billig mit dem heiligen Paulo sagen, Röm. 7, 22. 23.: „Ich habe Lust an Gottes Gesetz nach dem inwendigen Menschen. Ich sehe aber ein ander Gesetz in meinen Gliedern, das da widerstreitet dem Gesetz in meinem Gemüthe“; item R. 19.: „Das Gute, das ich will, das thue ich nicht, sondern das Böse, das ich nicht will, das thue ich.“ Ich wollte gern glauben, daß wahr ist, und also glauben, wie es in Wahrheit ist; ich wollte gern, daß mein Glaube so gewiß und stark wäre, als das Ding selber ist; aber die Sünde im Fleisch widerstrebt dem Geist, daß ich ja so stark nicht glauben kann. Wer sich aber der heiligen Taufe nach ihrer Würde recht verwundern und dieselbe in seinem Herzen hoch genug erheben könnte, der würde des Todes und aller Teufel spotten; der Türke würde ihm cammoz (wie das hebräische Wort lautet) sein, das ist, wie lauter „Spreu“, wie der 1. Psalm V. 4. sagt, „die der Wind zerstreuet“.

64. Darum soll man dies stets lesen und betrachten, und bitten, daß unser Glaube von Tage zu Tage gemehret werden möge. Dasjenige, so wir glauben, ist gewiß, und kann uns nicht fehlen. Darum dürfen wir uns der sonderlichen Erscheinungen oder Gesichte, so die Väter gehabt, nicht so hoch verwundern oder derselben auch für uns begehren. Abraham wird am jüngsten Tage sagen: Wenn ich wäre ein Christ gewesen im Neuen Testament, so wollte ich dies mit viel einem stärkeren Glauben geglaubt haben. Ich habe nur der einigen Verheißung geglaubt; ich bin das Exempel und der Erste gewesen, du hast unzählige Exempel, du hast deine Eltern, deine Brüder, die dich absolvieren und dir diese sichtlichen Gestalten darreichen.

65. Deshalb ist es eine Schande, daß wir im Glauben noch schwächer sind, als die Patriarchen gewesen, welche auf eine Verheißung Könige und Königreiche überwunden haben. Sie haben Einem Worte geglaubt, wir aber, denen mit so vielen Exempeln der vorigen und auch jetzt gegenwärtigen Heiligen geholfen wird, je länger wir predigen, je fauler wir, leider! werden. Und widerfährt uns eigentlich eben das, so Pythagoras gedichtet hat vom Laufe des Himmels, er habe dasselbe von sich selber gehabt, oder aber von den Vätern, nämlich, daß die Bewegung oder Lauf der Sterne stets einen lieblichen Gesang von sich gebe; aber die Herzen und Ohren der Menschen sollen dadurch, daß solches für und für geschieht, erstarrt sein, daß sie solchen Gesang nicht mehr hören.

66. Also sehen wir GOtt taufen, absolviren, trösten und das Abendmahl reichen; wer ist aber, der zuhöre oder sich dessen verwundere? Wir lassen uns seine Drohungen, Verheißungen, Trost und Sacramente nichts oder gar wenig bewegen. Derhalben sind wir billig zu strafen, die wir nach dem so leiblichen und schönen Gesange des Himmelreichs nichts fragen. Und wenn wir auch die Engel in ihrer Majestät predigen hörten, so würden wir doch denen nicht mehr achten, als wir es jetzt achten, da wir dem Pfarrherrn oder andere Kirchendiener hören. Wenn wir es aber gewißlich dafür hielten, daß es Gottes Wort wäre, das wir hören, so würden wir nicht also schnarchen, und so gar sicher und faul sein. Aber bisweilen wir meinen, daß nur Menschen mit uns reden, und daß es eines Menschen Wort sei, das wir hören, werden wir gar zu unvernünftigen Thieren.

67. Hier sollten wir nun seufzen und klagen, daß die Erbsünde, leider! so kräftig und gewaltig ist auch in denen, so neu geboren sind. Es sind ihrer sonst genug und übrig genug, die das Wort lästern und verfolgen. Wiewohl wir auch nicht gut sind, wie Christus sagt Luc. 11, 13.: „Ihr, die ihr arg seid“, so sollen wir uns doch befleißigen, daß wir ja das Wort nicht verfolgen, und sollen die Bosheit unseres Fleisches erkennen; wie St. Paulus Röm. 7, 24. auch schreit: „Ich elender Mensch, wer wird mich erlösen von dem Leibe dieses Todes?“ Der Leib ist sterblich, vergiftet, verderbt und ganz und gar höllisch. Wenn es uns wohl geht, so können wir uns leichtlich rühmen und uns im Herrn segnen; wie im 30. Psalm V. 7. steht: „Ich aber sprach, da mirs wohl ging: Ich werde nimmermehr darnieder liegen“; dagegen V. 8.: „Aber da du dein Antlitz verbargest, erschrak ich.“ Da kommt das Fleisch wieder und vertreibt die Freude alle.

68. Darum, sage ich, wird dies alles uns zur Lehre geschrieben, daß wir nicht allein der Widersacher Argumente strafen und widerlegen, sondern auch dem Fleische widerstehen können, welches uns mit der schweren Last der Sünde und der Furcht niederdrückt, daß wir nicht sagen können: In Samen Abrahams bin ich gerecht, selig, ein Herr des Todes, des Teufels und Türken. Derhalben sollen wir das Wort lernen und dasselbe lieb haben, darin gelehrt wird, daß man an den verheißenen Samen glauben soll. Darnach aber soll man auch das Fleisch zähmen und zwingen, welches dem Gesetze im Gemüth widerstreitet und uns gefangen nimmt. Und wenn wir etwa überwunden sind, soll sich bald der Glaube wiederum aufrichten und stärken, damit das Fleisch den Geist nicht gar unterdrücke.

69. Wo wir aber ja so bald nicht können zur Vollkommenheit kommen, daß wir gar keine Sünde, Tod und Hölle fühlen oder empfinden, und daß wir keine vollkommene Dialectiker und Rhetoriker werden können: so lasset uns unterdeß ABC-Schüler oder aber Donatisten sein, bis wir zu der Dialectik und Rhetorik kommen mögen. Denn dies ist ein Leben der Erstlinge und nicht der Zehnten. Darum klagt St. Paulus Röm. 7, 23. selbst auch, wie zuvor gesagt ist: „Ich sehe ein ander Gesetz in meinen Gliedern, das da widerstreitet dem Gesetz in meinem Gemüthe, und nimmt mich gefangen“; das Gute wollte ich wohl gerne thun, wollte wohl gerne ein solcher Mensch sein, der da nicht zweifelte, sich nicht bekümmerte irgend einer Noth oder Trübsal halben. Dasjenige, so wir glauben sollen, ist zwar an sich gewiß und die Schrift ist wahr, aber das Fleisch liegt uns im Wege, hindert uns allenthalben.

70. Was soll ich aber thun? Wenn ich es in meinem Herzen so stark nicht glauben und ergreifen kann, so zürne ich mit dem Fleische, und sage mit St. Paulo Röm. 7, 24.: „Wer wird mich erlösen von dem Leibe dieses Todes?“ Denn dieser Leib und Fleisch macht den Tod bitter und schrecklich, das Leben aber, so in
\quellseite{51}
Sünden und vielen Dornen ganz bitter ist, wird durch den Glauben süß. Und damit geht die heilige Schrift allenthalben um, daß sie uns lehrt, wie wir den Samen Gottes und uns selbst erkennen sollen, das ist, daß wir erkennen den alten und angebornen Fluch von Adam und den angefangenen Segen durch Christum.

\section[Zweiter Theil]{Zweiter Theil. Wie Isaak sein Weib verleugnet, mit ihr scherzt, und Abimelech diesen Scherz gewahret}

\begin{quote}
\textbf{V. 6. 7. Also wohnte Isaak zu Gerar. Und wenn die Leute am selben Orte fragten von seinem Weibe, so sprach er: Sie ist meine Schwester. Denn er fürchtete sich zu sagen: Sie ist mein Weib; sie möchten mich erwürgen um Rebekka willen, denn sie war schön von Angesicht.}
\end{quote}

71. An diesem Orte mag man abermal fragen: Warum doch der heilige Geist solche unnütze Dinge schreiben möge? Denn wozu dient es, daß man der Gemeinde Gottes solch ungereimtes und närrisch Ding soll vorhalten? Soll man damit das Volk Gottes unterrichten, daß Rebekka eine Schwester Isaaks ist genannt worden und nicht sein Weib? Darauf antworte ich also: Auf solche Weise macht Gott die Welt mit aller ihrer Weisheit zur Närrin. Wenn irgend ein Grieche oder Römer dies lesen sollte, so würde er dessen spotten als eines solchen Dinges, das nicht werth wäre, daß man es hören sollte. Aber Gott will damit erklären, daß die Weisheit der Menschen lautere Thorheit sei und daß sie also zu Schanden werde.

\subsection{I}

72. Bisher hat Mose beschrieben, wie Isaak ein göttlich Gespräch und himmlische Verheißung gehabt, und hat ihn stark gemacht und gar unüberwindlich wider die höllischen Pforten und alles Unglück; hier sieht man an ihm so große Schwachheit und solch schändlich Unvermögen, daß man auch nichts Schändlicheres erdenken könnte. Denn ist das nicht ein schändlich Exempel, daß er leugnet, daß Rebekka sein Weib sei?

73. Die Theologen disputiren davon: Ob er daran nicht gesündigt habe, daß er sein Weib verleugnet, und lügt, daß sie seine Schwester sei? Es ist wahrlich eine schändliche Schwachheit. Denn er sagt: Ich will sagen, sie sei meine Schwester; denn sie möchten mich sonst erwürgen. Als wollte er sagen: Es mag das Weib hinnehmen und schänden, wer da will, wenn ich nur leben mag; wenn ich sage, sie sei mein Weib, so werden sie denken, sie können mir dieselbe nicht nehmen, es sei denn, daß sie mich zuvor erwürgen. Ist aber das nicht ein närrisch, ungeschickt und unbillig Ding von einem solchen großen Manne? Sollte er nicht frei heraus gesagt haben: Sie ist mein Weib, ihr mögt mich erwürgen oder nicht? „Isaak fürchtete sich“, sagt der Text. Ei, daß ein solcher Mann in ein so schändlich Fürchten fallen soll, daß er sich vor dem Tode fürchtet! Also hatte Elia 1 Kön. 19, 1—3. achthundert falsche Propheten mit großem Eifer und starkem Muth erwürgt, und war niemandes Gewalt so groß, davor er sich gefürchtet hätte; aber da ihm die Isebel drohte, wird er so gar erschrocken, daß er davon flieht; zuvor hat er sich vor dem Könige nicht gefürchtet, jetzt flieht er vor einem Weibe.

74. Dies ist nun alles, wie es äußerlich scheint, fast närrisch, aber doch sehr fein und heilsam. Denn es ist geschrieben der Kirche und Gemeinde Gottes zum Trost, auf daß wir wissen sollen, wie gnädig und barmherzig Gott sei. Wir mögen zwar arg und schwach sein, allein, daß wir uns unter denen nicht finden lassen, so Gott verfolgen, ihn hassen und lästern; mit unserer Schwachheit will Gott gerne Geduld tragen.

75. Ich kann zwar die Väter nicht entschuldigen, wie die Andern thun, ich will es auch nicht thun, ja, ich höre gern, daß die Heiligen etwa gefallen und schwach geworden sind: nicht, daß ich es an ihnen loben wollte, als wäre es wohl gethan und als wären es Tugenden; wie ich die Apostel auch nicht entschuldige, daß sie von Christo sind gelaufen und daß St. Petrus ihn verleugnet hat, und dazu andere Schwachheit, Thorheit und Ungeschicklichkeit derselben. Und dies wird nicht geschrieben um der harten, hoffärtigen und verstockten Menschen willen, sondern auf daß damit angezeigt werde, wie es um das Reich Christi bestellt sei, der in seiner kleinen Heerde arme, elende und schwache Gewissen hat, welche leichtlich verletzt werden und sich schwerlich wiederum trösten lassen: er ist ein König der Starken und zugleich auch der Schwachen und Blöden; den Hoffärtigen ist er feind, den Starken sagt er ab, will gegen sie streiten. Die Pharisäer und sicheren Menschen straft er, aber die erschrockenen, kleinmüthigen, betrübten und bekümmerten Herzen will er nicht zerbrechen oder zu Schanden machen: er will das glimmende Docht nicht auslöschen, Jes. 42, 3. Dies ist seine Weise und Gebrauch für und für.

76. Also hat er vom Anfange der Welt gethan bis an das Ende. Er ist stark und allmächtig, aber das glimmende Docht wird er nicht auslöschen, und gebietet uns auch, daß wir den Schwachen im Glauben sollen aufnehmen. Denn es ist ein großes Ding, daß man glauben soll, daß die Christen durch den Tod in das Leben eingehen. Das Fleisch erschrickt, wo man des Todes, des Gerichtes und ewigen Feuers gedenkt; sonst ist es stark und unerschrocken, wo es ohne große Sturmwinde ist, die ihm zuwider sind.

77. Derhalben ist das Reich Christi ein solch Reich, daß zugleich Leute hat, welche immer oder doch zeitweilig schwach oder stark sind; jedoch sind deren wenige, die allezeit stark und beherzt sind. Die aber etwas schwächer sind, denn solche Starke, mit denen soll man nicht zu hart umgehen. Denn dies Reich ist ein Reich des Trostes, der Elenden und Betrübten. Es ist nicht dazu gesetzt, daß man darin herrschen und schrecken soll, es sei denn die hoffärtigen und verhärteten Köpfe, welche man allewege schrecken muß; die aber erschrocken sind, denen soll man Trost mittheilen und zu ihnen sagen: Du bist ein glimmendes Docht, Gott will nicht haben, daß ich dich auslöschen soll. Denn Christus ist ein König der armen, elenden, betrübten und schwachen Herzen.

78. Hierher gehören die Exempel der heiligen Schrift, welche die sehr großen und heiligen Männer, Abraham und Isaak, anführt, als die den Tod, Teufel und Hölle überwunden haben, und welche sie zu Freunden und Propheten Gottes macht, wie im 105. Psalm V. 6. 8. 9. steht. Und dennoch beschreibt sie dieselbigen, wie sie auch so gar schwach gewesen und sich schändlich gefürchtet haben. Also sagt David von sich selber, im 30. Psalm V. 7. 8.: „Ich aber sprach: Ich werde nimmermehr darnieder liegen“, der Tod und Teufel sind nichts bei mir; „denn, Herr, durch dein Wohlgefallen hast du meinen Berg stark gemacht; aber da du dein Antlitz verbargest, erschrak ich.“ Wenn uns Gott stärkt und Muth gibt, so hat man gut stark und unerschrocken sein; auf daß wir aber nicht stolz werden und solche Stärke unserer eigenen Kraft zuschreiben, läßt er uns etwa sinken. Eben also ist es dem Patriarchen Isaak auch ergangen.

79. Dies ist nun nütze und gut, die Kirche und Gemeinde Gottes damit zu lehren, nämlich, daß solche sehr große und heilige Männer, so die göttlichen Verheißungen gehabt, und wie der die Sünde und Tod triumphirt und gesiegt haben, so jämmerlich dahin fallen, und werden darnach gar zu Kindern und flüchtig. Den Tod fürchten, ist an sich selber böse und eine Sünde; wenn aber das dermaßen geschieht, daß darüber dein Weib in Gefahr ihrer Zucht und Keuschheit kommt, so wird es noch einmal so böse. Und diese Lehre hält uns der Heilige Geist vor unter solcher geringen und verächtlichen Gestalt.

80. Derhalben sollen wir dies alles auf uns ziehen. Denn also stärkten wir heutiges Tages die Gottseligen auch, so des Fürsten halben in Gefahr stehen, und sagen zu ihnen: Wenn dich schon der Türke erwürgt und verschlingt, was thust er denn? Was kann er mehr thun denn würgen? „Fürchtet euch nicht vor denen“, spricht Christus Matth. 10, 28., „die den Leib tödten“ etc. Wenn er dir schon dein Leib, dein Haus und Anderes mehr mit Gewalt nimmt, was liegt daran? Das alles muß man doch sonst auch verlassen, und ich wollte es dafür halten, daß mir daran recht wohl geschähe, wenn mich der Teufel um Gottes Sohnes willen erwürgte. Ja, also sind wir gesinnet, wenn wir beherzt sind und uns der Heilige Geist stärkt. Aber in der andern Stunde kann mir oder auch einem jeglichen anderen Menschen eben dasselbe widerfahren, das Isaak auch widerfahren ist.

81. Darum, sage ich, ist dies uns zur Lehre geschrieben, auf daß wir durch Geduld und Trost der Schrift Hoffnung haben, Röm. 15, 4., daß wir nicht denken: Ich bin so heilig oder stark im Glauben nicht, wie Isaak oder Jakob gewesen sind. Sie sind zwar stark und beständig gewesen, wenn sie GOtt gehalten hat und ihren Berg stark gemacht; aber da er sein Antlitz verborgen hat, sind sie erschrocken, Ps. 30, 8. Und wir mögen von ihnen eben dasselbe auch sagen, daß St. Petrus zu Cornelius sagt in den Geschichten der Apostel am 10. Cap. V. 26.: „Ich bin auch ein Mensch.“ Denn es hat weder Petrus noch Paulus, ob sie wohl heilig gewesen, ein besser Fleisch gehabt als wir: sie sind auch Adams Kinder gewesen, wie wir sind. St. Paulus rühmt sich bisweilen und ist so hoffärtig, gleich als hätte er jetzt schon alles Unglück überwunden; bisweilen klagt er auch ganz erbärmlich über den Schrecken, so er innerlich gehabt, und über den Streit oder Kampf, so er auswendig hat führen müssen. Wo war da der große unüberwindliche Held und der etwa die ganze Hölle überwunden hat?

\quellseite{52}

82. Derhalben sollen wir am bösen Tage, und wenn Unglück und Widerwärtigkeit vorhanden ist, die Hoffnung nicht fallen lassen und verzagen: und am guten Tage, wenn es uns glücklich und wohl geht, sollen wir auch nicht stolz und hoffärtig werden. Denn die, so im Glauben am stärksten gewesen, als Abraham und Isaac, sind in solche schändliche Schwachheit gefallen, uns zur Hoffnung und Trost, auf daß wir lernen, was die Weise und Gebrauch sei in diesem Reiche und mit diesem Könige. Denn da sieht man eine unüberwindliche Stärke und wiederum eine höchst überwindliche Schwachheit, daß ich also rede, welches alles dahin geht, daß allein GOtt geehrt und gerühmt soll werden. Denn wir rühmen uns unserer eigenen Kräfte nicht, sondern weil unser König zur Rechten GOttes des Vaters sitzet, so rühmen wir uns in ihm und erheben uns hoch, seine Kraft ist in unserer Schwachheit mächtig, 2 Cor. 12, 9., wie man das hin und wieder in allen Exempeln der Väter sehen mag.

83. Von dieser Weisheit aber weiß die Welt nichts, es hat sie der Heilige Geist allein und hält sie in der Kirche vor unter einer sehr schlechten und geringen Decke. Und wir sollen die Historien der Väter also lesen, daß uns die Exempel des Glaubens stärken und beherzt machen; die Exempel aber der Schwachheit sollen uns trösten, daß wir nicht verzagen und auch nicht hoffärtig werden, und daß wir Andere auch trösten. Wenn du erschrocken bist, so thue ein Unglück nicht zum andern, fliehe vor dem Herrn Christo nicht, verleugne ihn auch nicht. Denn wiewohl ich schwach bin, sollst du sagen, so wollte ich doch nicht gerne, daß Christus, mein Herr, sollte gelästert oder geschmäht werden, sondern ich will, daß seine Ehre vorkommen, und sein Reich fest und unbeweglich sein und bleiben soll: ich bin zwar gefallen, was soll ich nun thun? soll ich darum verzagen? Mit nichten, ist doch eben dasselbe auch den so großen Erzvätern widerfahren, von welchen Christus ist hergekommen.

84. Mose hat fleißig ausgedrückt, daß sich Isaac vor dem Tode gefürchtet hat. „Er sprach“, sagt er, „sie ist meine Schwester.“ Warum sagt er aber das? „Sie möchten mich erwürgen um Rebekka willen.“ Derhalben hat er aus Furcht vor dem Tode sein Weib in die Gefahr gesetzet, daß sie hätte im Ehebruch kommen mögen, auf daß nur er frisch und gesund und am Leben bleiben möchte.

85. Wer nun diese Sünde Isaacs wollte hoch aufmüßen und sehr groß und schwer machen, der würde viel Gründe dafür finden können. Denn da du leugnest, daß Rebekka dein Weib sei, damit verleugnest du GOtt, der sie dir zugeschickt hat und dich dazu mit zwei Söhnen gesegnet hat; darnach, über die Sünde, die du mit der Lüge thust, führst du dein Weib in die Gefahr und verräthst sie, daß sie darüber möchte um ihre Keuschheit kommen. Das ist des Teufels Rhetorik. Laßt uns aber in der Rhetorik des Heiligen Geistes bleiben.

86. Jedoch würde diese Schwachheit Isaaks zuwege gebracht haben, wo es GOtt nicht gewehret und gehindert hätte, daß sein Weib geschwächt worden und in Ehebruch gefallen wäre, und daß sonst viel mehr Böses hätte geschehen mögen. Aber GOtt läßt seine Heiligen also fallen, daß er sie doch gleichwohl noch erhält.

87. Was die Sitten der Bewohner Palästinas gewesen und was sie für ein Leben geführt haben, weiß ich nicht, und verstehe auch den Umstand nicht, davon Mose handelt, daß Abimelech sei fromm gewesen und den Ehebruch bei Lebensstrafe verboten habe. Warum fürchtet sich denn nun Isaac vor ihm? Er wird müssen gesehen haben (1 Mos. 12, 11. ff. Cap. 20, 2.) am Exempel Abrahams, wie sie einander die Weiber mit Gewalt genommen und entführt haben; wo jemand mäßiger gelebt und sind auch vom Leibe stärker gewesen denn heutiges Tages. Nun folgt ein Spiel oder Scherz mit Rebekka, so über die Maßen ungeschickt und ungereimt ist.

88. Es ist wohl glaublich, daß er dieser Sünde etliche Exempel zuvor wird gesehen oder aber von Andern gehört haben; wo sie aber andere Sitten gehabt und ein solch Leben nicht geführt haben, und Isaac gefürchtet hat, daß sie solche Leute sein möchten, da er doch nirgends ein Exempel gehabt: so kann man doch daraus abnehmen, daß eine große unglaubliche Schwachheit an ihm muß gewesen sein. Er wird vielleicht allein gedacht oder sich eingebildet haben: Die Einwohner dieses Landes sind gottlose Leute und fragen nach GOtt nichts; darum verachten sie auch die Menschen und gehen unbillig mit ihnen um.

89. Daher ist aber diese Sorge und Bekümmerniß gekommen, weil Rebekka schön war von Angesichte. Denn also steht im Texte, und dünkt es weltliche Menschen sehr lächerlich. Denn Isaac ist zum wenigsten achtzig oder neunzig Jahre alt gewesen. Im sechzigsten Jahre seines Alters sind ihm zwei Söhne geboren worden, fünfzehn Jahre darnach ist sein Vater Abraham gestorben, sie aber, Rebekka, ist bei siebenzig Jahren alt gewesen. Denn die Weiber sind zehn Jahre jünger gewesen, als die Männer; wie droben (Cap. 17. V. 17.) von Sarah ist gesagt worden, welche auch, da sie schon neunzig Jahre alt war, doch gleichwohl noch schön von Angesicht gewesen.

90. Derhalben sind die Menschen zu der Zeit nicht allein schöner, sondern auch stärker vom Leibe gewesen denn jetzt, sintemal sie in solchem Alter Kinder zeugen und die Weiber haben schwanger können werden, in welchem Alter bei uns keine Hoffnung mehr ist, daß sie könnten fruchtbar sein. Rebekka aber, ein alt Weib von siebenzig Jahren, ist noch so schön von Angesichte, daß sie es den Andern hat zuvor gethan und daß sie um der Schönheit willen in Gefahr ihrer Zucht und Keuschheit gekommen ist. Damit wird die Zeit, so dazumal gewesen ist, gelobt, nämlich, daß die Leute

\subsection{II}

\begin{quote}
\textbf{V. 8. Als er nun eine Zeitlang da war, sahe Abimelech, der Philister König, durchs Fenster, und ward gewahr, daß Isaac scherzte mit seinem Weibe Rebekka.}
\end{quote}

91. Ein solcher großer Mann, und der schon achtzig Jahre ist, scherzt noch mit seinem Weibe, die auch fast siebenzig Jahre alt ist. Hat denn der Heilige Geist dazu Lust, daß er sich Narrnwerk schreiben mag, und sich also hernieder lasse, von solchem geringen Scherz, so im Ehestande geschieht, zu reden? Warum erzählt er nicht lieber das Fasten der vierzig Tage, wie von Elia geschrieben steht? 1 Kön. 19, 8. Erstlich aber wollen wir die Meinung der Juden, die gar grob und unverschämt ist, nach der Weise der cynischen Philosophen, verwerfen. Denn dieselbigen groben unflätigen Säue sagen, Isaac soll ehelichen Scherz mit Rebekka gepflogen haben frei öffentlich, daß Abimelech habe zugesehen; als ob dieser so heilige Vater ein Cyniker oder Diogenes gewesen wäre, der öffentlich auf der Gasse bei seinem Weibe soll gelegen haben.

92. Der Heilige Geist pflegt das Beischlafen der Eheleute zu bedecken und zu beschönigen und bringt es nicht also öffentlich an den Tag; und Isaac und seine Rebekka sind schamhaftiger gewesen, denn daß sie eine solche große Schande hätten begangen haben; denn sie sind in der Furcht GOttes auferzogen worden und daß sie GOtt mit gebührlicher Ehrfurcht geehrt haben, und haben dazu die Erbsünde auch wohl verstanden, damit das Werk der Kinderzucht entstellt ist. Derhalben es GOtt hat verborgen haben wollen und hat es mit Kleidung und der hochzeitlichen Fröhlichkeit geziert.

93. Der Text aber zeigt heimlich an, wenn du die Umstände fleißig betrachtest, daß dieser Scherz hergekommen sei mehr aus Traurigkeit des Herzens und daß sie etwas erschrocken gewesen, denn aus Geilheit oder Unzucht, vornehmlich an Rebette. Da diese von so vielem Unglück und Jammer schier verzagt worden, mag sie vielleicht zu ihrem Manne gesagt haben: Ach! mein lieber Isaak, siehe doch, wohin wir gekommen sind, wie jämmerlich und elend wir im Lande umher ziehen. Wie das Weib Hiobs ihrem Manne auch seinen schweren Unfall vorrückt, da sie zu ihm sagte Hiob 2, 9.: „Segne Gott und stirb“ etc. Siehe, wir haben zwei Söhne, dazu auch Schafe und ander Vieh mit unserem Hausgesinde; wir sind Fremdlinge und im Elende; und über das wird uns nun auch verboten, nach Egypten zu ziehen, da muß ich nun auch den Namen eines Eheweibes verlieren.

\quellseite{53}

94. Dies, sage ich, hat wahrlich dem armen Weibe können einfallen, daß sie also wird gedacht und zu ihrem Manne gesagt haben; dieweil sie nicht sehr begehrt war, und mit dem Zorn und Ungeduld, so sie in solchem Unglück gefühlt, gekämpft hat. Denn wer könnte doch solches dulden und ertragen ohne Seufzen? Darum, auf daß Isaak sie stärken möchte, hat er sie vielleicht um den Hals genommen und geküßt, und sie heißen zufrieden und getrost sein; wird zu ihr gesagt haben: Meine liebe Rebette, trage doch dies Unglück mit Geduld; Gott ist mir erschienen, und hat mir verheißen, daß er mich erretten, und dazu schützen und schirmen wolle. Also könnten wir Isaak vertheidigen, oder was er gethan hat, entschuldigen, daß wir aus ihm mit den groben unflätigen Säuen, den Juden, keinen Cyniker oder Diogenes machen.

95. Ich will aber diese Gedanken unentschieden lassen, will darauf nicht hart dringen und will sie auch nicht gar verwerfen. Es läßt sich aber gleichwohl ansehen, daß sie der Wahrheit fast ähnlich seien, dieweil ja sonst niemand Lust bezeigt, daß er sein Weib in fremdem Lande, da er unter den Feinden ist, küssen sollte. Es war ihr aber Trost vonnöthen, und daß ihr möchte angezeigt werden, daß Gott gegenwärtig bei ihnen wäre, und ihnen verheißen, sie zu erretten und zu schützen; und dasselbe hat mit solchem Küssen und Scherzen geschehen können.

96. Das ist ein ehrlicher Scherz, so einem frommen ehrlichen Ehemann mit seinem Weibe wohl ansteht. Wenn er mit seiner Schwester oder Gesinde im Hause dermaßen scherzen wollte, das würde ihm nicht wohl anstehen. Denn daselbst gehört es sich, daß man sie heiße, was sie thun und lassen sollen, und daß dabei Ernst sei, auch wenn man sie tröstet. Aber mit der, die mir Gott zugefügt hat, mag ich scherzen, spielen und freundlich reden, auf daß ich mit Vernunft und Bescheidenheit bei ihr leben möge. Wie St. Petrus 1. Ep. 3, 7. die Männer ermahnt, daß sie dem weiblichen, als dem schwächsten Werkzeuge seine Ehre geben sollen, als auch Miterben der Gnade des Lebens, und die da Mütter und der andere Theil des menschlichen Geschlechtes sind, wiewohl der schwächste Theil, so viel den Leib betrifft.

97. Wir Männer sind etwas stärker, darum sind uns auch in unserm Amte größere Werke befohlen, darin oftmals an uns gleichwohl ein weibisches Herz gespüret wird. Es ist aber doch ein Unterschied zwischen einem männlichen und weiblichen Gemüth. Die Weiber sind mit der Hausarbeit beladen, dazu mit der Geburt und daß sie die Kinder müssen aufziehen, haben dazu auch die Hausorge auf sich liegen; wir aber haben auf uns das Welt- und Kirchenregiment, dieweil wir etwas stärker sind. Wie wir uns aber im Kirchen- und Weltregiment halten, sehen wir zwar wohl. Es wird oftmals in der Regierung, so die Männer zu verwalten haben, nicht weniger Schwachheit gespüret.

98. Derhalben will ich diesen Gedanken folgen, und will zugeben, daß sie Isaak freundlich geküßt habe, nicht allein um des Trostes willen, sondern daß er damit anzeigen möchte die Geberde eines frommen Ehemannes, damit er sich gegen sein ehrliches Weib scherzhaft und freundlich pflegt zu erzeigen; wie man denn dergleichen Geberde oft sehen mag an Ehemännern, die da etwas freundlich sind.

99. Der König ist zwar in Weise nicht gewesen, daß er solche wichtige Ursachen, wie ich sie zuvor erzählt, hätte verstehen können, warum Isaak sein Weib möchte geküßt haben; sondern hat verstanden, daß er dazu die gemeinen Ursachen der Ehemänner gehabt hätte, die man wohl öffentlich zeigen mag.

100. Und ist sehr fein, wo ein Ehemann auch in äußerlichen Sitten sich gegen sein Weib freundlich erzeiget. Dagegen aber ist es schändlich und ärgerlich, wo etliche Männer sind, die sich seltsam, wunderlich, hart und unfreundlich gegen ihre Weiber erzeigen, und geben weder mit Worten noch mit Werken gar kein Anzeichen von sich, daran man spüren möchte, daß sie ihre Weiber lieb hätten und ihnen freundlich gesinnet wären. Gegen die Fremden oder auch gegen das Gesinde soll man ernst sein. Also hören wir oftmals, daß von den Weibern solches gelobt wird und daß sie sich darüber freuen, wenn sie sehen, daß der Mann mit seinem Weibe freundlich und in guter Einigkeit lebt. O eine gute Ehe ist das, da ist rechte Liebe, so Eines zum Andern hat, und eine selige Einigkeit; davon pflegen die Männer auch zu rühmen. Woher wollten sie aber das wissen, wenn sie nicht gesehen hätten, daß solcher Scherz und eheliche Freundlichkeit bei vielen Eheleuten fast selten ist?

101. Und ist aber keine härtere und greulichere Uneinigkeit und Zwietracht, denn zwischen Eheleuten und Brüdern; dagegen aber wenn sie einander lieb haben, mit einander scherzen und freundlich leben, solche Ehe hat man lieb und wird hin und wieder von jedermann gerühmt. Denn es ist ein selten Ding um des Teufels willen, welchen der Ehestand für und für zum Feinde hat, der sich untersteht, die, so Gott zusammen gefügt hat, zu verwirren, wie er nur kann und mag. Wiewohl es die Welt nicht versteht, sondern läßt sich träumen, daß der Ehestand Anderes nichts sei denn nur das Beischlafen, und daß Mann und Weib einander küssen und äußerlich solcher Liebe pflegen.

102. Es ist aber um den Ehestand nicht also bestellt, daß sie allein nach dem Leibe zusammen sollten gefügt sein, sondern auch mit den Herzen; welches sich an etlichen gewissen Anzeichen und Wahrzeichen sehen läßt. Und wiewohl solche äußerliche Anzeichen ungeschickt und lächerlich scheinen, so stehen sie doch einer guten Ehe wohl an. Ich habe oft von frommen ehrlichen Matronen gehört, wenn sie andere Ehemänner lobten von wegen ihrer Freundlichkeit, Gelassenheit, und daß sie so gar stille und fromm wären, daß sie solches als ein sonderlich Wunderwerk rühmten, wenn ein Weib einen solchen frommen Mann oder ein Mann ein solch fromm Weib bekommen hätte.

103. Derhalben will ich solchem Verstande nach diesen Scherz Isaaks deuten und auslegen; wiewohl ich die vorigen Gedanken nicht verwerfe, und will sie auch nicht vertheidigt haben; denn allein den schändlichen Juden zuwider. Dies sind aber gewiß feine und göttliche Gedanken, nämlich, daß ein Ehemann nicht allein in der Schlafkammer, sondern auch sonst öffentlich dem Ehestande zu Lob sich fein freundlich gegen sein Weib halten soll: er soll nicht wunderlich, zornig und ungebärdig sein. Denn man kann die Exempel der Uneinigkeit, und wo ein Theil den anderen etwa verletzt oder erzürnt, sehr leichtlich sehen und ärgern sie auch sehr; vornehmlich, wo die Eifersucht dazu kommt, ist die Hölle selbst.

104. Derhalben ist es nütze und gut, daß solche Exempel der Freundlichkeit und Liebe unter den Eheleuten vorhanden seien, auf daß sich Andere auch gewöhnen an Freundlichkeit, Liebe und Geduld, wenn irgend ein Anstoß oder Unglück vorfällt. Denn der Ehestand ist eine göttliche Ordnung. Und von dem Weibe ist gesagt 1 Mos. 2, 18.: „Wir wollen dem Menschen eine Gehülfin machen, die um ihn sei.“ Derhalben, so muß ja ein äußerliches Anzeichen der Gewogenheit da sein um der Andern willen, vornehmlich dieweil wir solche treffliche Heilige, wie Isaak ist, in unserer Gesellschaft haben. Denn dies Leben ist voll unsäglicher Anstöße, damit leichtlich Eines das Andere erzürnen kann, wie man das in der Haushaltung vornehmlich sehen mag. Darum sollen sich fromme Eheleute dazu schicken, daß sie einander verzeihen und dessen bald vergessen mögen, wo Eines das Andere verletzt oder erzürnt hätte, wie hart oder schwer die Verletzung auch sein möchte, und sollen wiederum zu ihrer gewöhnlichen Freundlichkeit greifen auch äußerlich.

105. Denn darum verschmähet es der heilige Geist nicht, solch gering Ding zu beschreiben und zu rühmen, daß er damit erstlich den groben unflätigen Säuen begegnen möge, die nur fleischliche Lust suchen, aber gegen ihre Weiber sich fast rauh und ungebärdig bezeigen. Darnach hat er damit auch begegnen wollen den Feinden des Ehestandes, als den Papisten, welche den Ehestand verboten haben, als sollte er von Gott verdammt sein, auf daß wir ein Exempel haben, daran wir sehen, daß auch ein gewisser leichter Sinn im Ehestande Gott nicht mißfalle; welcher Stand sonst mit der Unreinigkeit der Erbsünde entstellt ist, daher die gewöhnliche und hitzige Brunst kommt, so der Mann zum Weibe oder auch das Weib zum Manne hat. Aber GOtt gebraucht derselben, den Mann und das Weib also zusammen zu fügen, und nennt das Weib eine Gehülfin um den Mann. Derhalben fügt sie GOtt zusammen, wie Christus Matth. 19, 6. sagt: „Was GOtt zusammen gefüget hat, das soll der Mensch nicht scheiden.“

\quellseite{54}

106. Und es ist ein großer Trost, daß wir wissen, daß GOtt ein Mißfallen daran hat, wo sich Eheleute von einander scheiden, und daß er haben will, daß sie Liebe und Freundlichkeit unter einander üben, sich freundlich küssen und mit einander scherzen, dem Teufel zuwider, welcher ein Stifter aller Uneinigkeit ist.

107. Da werden uns aber die, so in dem unreinen ehelosen Stande sind, vorwerfen, was St. Paulus sagt 1 Cor. 7, 1.: „Es ist dem Menschen gut, daß er kein Weib berühre“: und der Ehestand, sagen sie, ist ein unreiner Stand. Item, die Andern sagen: Sollte ich mit dem Weibe so freundlich umgehen und ihre seltsamen, wunderlichen Sitten alle dulden und ertragen? Ich will kurzum von ihr gefürchtet sein, will mit Gewalt über sie herrschen und regieren, daß sie sehe, daß ich ein Mann sei. Ei, wie wirst du auch so große Ehre einlegen, wenn du das arme schwache Werkzeug zerbrechen und das Gewissen beschweren wirst mit unbilliger Bitterkeit, so du geübt hast an deiner Schwester und Miterbin des Reichs GOttes, als die auch in derselben Gemeinschaft der Taufe, dazu aller Wohlthaten GOttes und der ganzen christlichen Kirche mit ist? Bedenke doch, was ein Weib sei, wer du seiest; oder wo du ja streiten willst, warum streitest du nicht wider deines Gleichen? Was meinst du, daß du damit für große Ehre einlegen werdest, wenn du so tyrannisch mit dem schwachen Werkzeuge umgehen wirst? Es ist wahrlich unbillig gethan von einem jeden Menschen und viel mehr von einem Christen, und ist zwar eine sehr geringe, ja, gar keine Ehre, ein solch schwach Geschlecht zu bezwingen. Darum bin ich denen sehr feind, die gegen die armen Weiber so beherzt sind und, wie man zu sagen pflegt: Domi leones, foris lepores, das ist, die im Hause Löwen, und draußen so blöde und verzagt sind wie die Hasen. Von den bösen und schändlichen Weibern rede ich nicht, die man auf feinere Weise bessern oder fromm machen kann.

108. Wir sollen aber GOtt dafür danken, daß wir wissen, daß der Ehestand von GOtt eingesetzt sei, und daß er GOtt, den Engeln und allen Creaturen wohlgefalle. Es gefällt GOtt wohl, daß ich gegen mein Weib freundlich bin, und dagegen gefällt es ihm gar nicht, wo ich seltsam und wunderlich bin, dazu mich mit Bitterkeit gegen sie vernehmen lasse; viel weniger aber gefallen GOtt die, welche den Ehestand verbieten und verdammen. Und solch Gebot, damit sie den Ehestand verboten haben, ist zwar damit greulich gestraft worden, da der Teufel anstatt der göttlichen Ordnung das wilde Wesen mit allerlei Unzucht und unzähliger Vermischung, dazu auch mit unsäglichen sodomitischen Sünden hat eingeführt. Und dasselbe ist also geschehen durch GOttes gerechten Gericht, die weil sie dem heiligen Stande Abbruch gethan und denselben geschmäht haben, welchen doch GOtt selbst geordnet und auch gesegnet hat.

109. Es hat zwar dieser Stand seine eigenen Gebrechen oder Unreinigkeit, aber GOtt duldet solche Gebrechen im Ehestande und thut, als sehe er es nicht. Wo die rechte natürliche Ordnung gehalten wird, Kinder zu zeugen, alsdann thut GOtt, als sehe er es nicht, und verzeiht frommen Eheleuten solche Gebrechen und Unreinigkeit; oder wie Lucas in den Geschichten der Apostel ein sehr fein griechisch Wort braucht, da er sagt im 13. Cap. V. 18., wie GOtt bei vierzig Jahr lang der Juden Weise in der Wüsten geduldet habe. Also will er hier im Ehestande sich diese Weise auch gefallen lassen und sie dulden, auf daß also das menschliche Geschlecht möge erhalten werden. Ich weiß wohl, sagt GOtt zu dir, daß du in Sünden empfangen bist, ich will dir es aber verzeihen, und dir zulassen, daß du dein eigen Eheweib haben magst zur Gehülfin; an der allein sollst du deine Lust haben, und sollt es daneben auch dulden und fein in der Stille ertragen, wenn irgend eine Traurigkeit oder Müßhelligkeit vorfallen wird. Ich will, sagt er, damit auch Geduld tragen, du greifest entweder dein Weib freundlich an bei nächtlicher Weile, wenn sie nackend ist, oder aber am Tage, wenn sie bekleidet ist.

110. Wenn GOtt nicht aus sonderlicher Gnade dies eingesetzt hätte, daß also im Ehestande ein Mann und ein Weib sollten zusammen gefügt werden, ei, wie sollte die ganze Welt so großes Verlangen darnach haben, auf daß sie durch dies Mittel und Hülfe von unziemlicher Vermischung und Unreinigkeit möchte erlöst werden! Jetzt aber verachten wir diese rechtmäßige Verbindung, welche er mit einem stetigen und unbeweglichen Gebot bestätiget hat, dadurch er die Hurerei und andere unziemliche Vermischung verboten haben will. Er gebietet, daß du dir Eine wählen sollst, die dir wohlgefällt, mit welcher du dein Leben mögest zubringen. Es ist ja wohl Sünde; aber ich will dir es zulassen, spricht GOtt, daß du scherzest und lächest mit deinem eigenen Weibe und dich fremder Weiber enthaltst. Wie geschrieben steht in den Sprüchen Salomos am 5. Cap. V. 18. 19.: „Freue dich des Weibes deiner Jugend. Sie ist lieblich wie eine Hindin, und holdselig wie ein Rehe. Laß dich ihre Liebe allezeit sättigen, und ergötze dich allewege in ihrer Liebe.“

111. Also tröstet uns dieser Text, darin geschrieben steht, wie Isaac im Alter mit seiner lieben Rebekka geschert habe nicht allein von Noth und Trostes wegen, sondern auch mit ehelichem Scherz, so diesem Stande wohl ansteht, auf daß er mit solcher äußerlichen Freundlichkeit anzeigen möchte, daß sie einander herzlich lieb gehabt. Und wird dies geschrieben den Eheleuten zum Trost, und zum Verdruß, Verachtung und Widerlegung der schändlichen heillosen Schwärmer, so den Ehestand verdammen.

112. Ich gedenke noch wohl, wie vor dieser Zeit, ehe denn das Evangelium wiederum an den Tag gekommen ist, die Ehemänner oftmals in der Beichte jämmerlich geklagt haben über solchen ehelichen Scherz, gleich als über eine große schwere Sünde. Ja, das doch sehr unbillig ist und damit auch der Schöpfer geschmäht wird, wenn irgend eine ehrliche Matrone oder fromme Jungfrau über den Kirchhof der Barfüßer gegangen war, haben die Mönche alsbald den Kirchhof mit Besen gefegt, und die heilige Stätte mit Feuer wiederum gereiniget. Lieber Herr GOtt, sind es doch unsere Mütter, und man kann ja die Flamme oder Brunst der Unzucht, so der Mensch durch die Erbsünde empfangen hat, mit keinem Besen oder Feuer wegkehren oder ausfegen, sondern wird durch solche Weise immer mehr und mehr angezündet und gelockt; und wenn man kein Weib haben kann, wird dadurch die Brunst noch immer gemehret. Aber diese Heuchler haben die Erbsünde nicht verstanden, darum haben sie auch nichts gewußt von der Hülfe und Arznei, so GOtt selbst uns dagegen gezeiget und gegeben hat.

113. Christus will haben, daß wir erstlich sollen getauft werden, auf daß das Herz gebessert und gesegnet werde; darnach wo du die Gabe der jungfräulichen Keuschheit nicht haben kannst, so magst du die Keuschheit wählen, so im Wittwen- oder Ehestande ist; denn es ist beides an männlichem und weiblichem Geschlecht gleich: derhalben haben sie auch Samen, und empfinden bei sich die Brunst und Begierde, Kinder zu zeugen. Was soll man denn nun hier rathen? Antwort: Dies ist der Rath, den man folgen soll: Hast du die Gabe der Keuschheit, so bleibe dabei; wo aber nicht, so nimm ein Weib, auf daß du mit unziemlicher Vermischung und cardinalischer Unzucht nicht befleckt, und darüber ewig verloren und verdammt werdest.

114. Das hebräische Wort zachak ist seiner Bedeutung nach nicht weniger allgemein, als in der deutschen Sprache das Wort lachen und in der lateinischen das Wort ridere auch ist. Denn das Wort ridere, lachen, wird bisweilen in seiner eigentlichen Bedeutung gebraucht, ein andermal aber wird es durch die Figur, antiphrasis genannt, gebraucht für ein bitter Lachen (das heißt im lateinischen Sprüchwort risus Sardonius) welches bitterer ist, denn alles Weinen sein kann. Und noch ein andermal heißt es Fröhlichkeit oder fröhliche Geberde und wenn man einen fröhlich ansieht; und wird solches auch verstanden von den Creaturen, so kein Leben haben. Als wenn wir sagen: Der Acker, die Wiese oder der Wald lacht, so ist das eben so viel gesagt, als daß es lustig und lieblich anzusehen ist; also sagen wir auf Deutsch auch: Er lacht, daß er sich schüttelt, das ist, er ist fröhlich; und durch die Figur, antiphrasis genannt, sagt man: Er lacht meiner dazu in die Faust; item: Wie höhnisch lacht er mein! Damit wird ein sehr bitteres und saures Lachen angezeigt.

115. Eben in solcher Bedeutung steht dies hebräische Wort im 2. Psalm V. 4. auch: „Aber der im Himmel wohnet, lachet ihrer, und der HErr spottet ihrer.“ Item, also sprach Abner zu Joab 2 Sam. 2, 14.: „Laß sich die Knaben aufmachen, und vor uns spielen“, als wollte er sagen: Laß sie einen guten Muth haben (per antiphrasin); das heißt auch gescherzt, wenn dir das Schwerdt durch den Bauch geht, das ist zumal ein greulich Spielen oder Lachen. Derhalben heißt an diesem Orte das Wort ridere, lachen oder scherzen, nicht, bei dem Weibe schlafen, sondern heißt eben so viel, als was droben 1 Mos. 18, 12. gesagt wird, da Sarah auch gelacht hat. Item, also hat Ismael den Isaac auch verlacht 1 Mos. 21, 9., das ist, er hat seiner gespottet. Also steht 2 Mos. 32, 6.: „Darnach setzte sich das Volk zu essen und zu trinken, und stunden auf zu spielen“, das heißt nicht, Abgötterei treiben, wie sich die Juden träumen lassen, sondern heißt, lachen und fröhlich sein. Es ist aber im Hebräischen zumal eine feine Anspielung, die fein zutrifft und auf den Namen Isaak deutet, welche man in keiner Sprache wiedergeben kann, als, wenn man auf Lateinisch sagen wollte, Isaac isaacabat, Isaac, das ist, der da lacht, der freundlich ist, thät freundlich gegen sein Weib; der da lieblich und liebenswert ist, liebte sein Weib; er stellte sich wie ein rechter Isaak, und dasselbe that er mit solcher Sicherheit und Vertrauen, wie einem Ehemanne geziemt.
\quellseite{55}

116. Derhalben wird in diesem Exempel christlichen Eheleuten eine treffliche Lehre und Trost vorgehalten, und darnach auch eine Widerlegung der schändlichen und schädlichen Papisten und Ketzer, welche aus allen Worten und Werken der Eheleute Todsünden gemacht haben. Und ich selber, da ich noch ein Mönch war, bin auch derselben Meinung gewesen, nämlich, daß der Ehestand ein verdammter Stand wäre. Wir pflegten davon zu disputiren: Ob es sich auch geziemen möchte, daß Einer eine fromme ehrliche Jungfrau liebe und sie zum Eheweibe begehre? Item: Ob es auch Sünde wäre, daß der Mann mit seinem Weibe scherze? Und ich verwunderte mich des Spruches Bonaventura, welcher unter den Mönchen der heiligste gewesen, sehr, da er sagt: Das ist keine Sünde, wo Einer irgend um eine Jungfrau werben läßt und sie zum Weibe zu haben begehrt: ja, man mag solches mit Recht wohl thun. Derhalben soll man davon nicht disputiren, soll auch daran nicht zweifeln. Item, er sagt auch, daß ein Ehemann mit seinem Eheweibe scherzen möge. Ich hätte mich viel eines andern Spruches oder Meinung zu ihm versehen, und die sich mit seiner Mönchsprofession besser gereimt hätte; denn ich verstand dazu mal dies nicht anders, denn wie es die Juden auslegen.

117. Darum bedenkt, ihr Jünglinge, die ihr jetzt daher wachset, daß ihr werdet zum Ehestand kommen und zu gemeinen Aemtern gezogen, daß ihr Gott danket für dieses göttliche Licht und Segen, daß wir nun wissen, daß der Ehestand ein heilig Ding sei, und daß wir mit unsern Eheweibern lachen, scherzen, und uns freundlich zu ihnen halten mögen, sie seien gleich nackend oder bekleidet, allein, daß wir uns fremder Weiber enthalten. Dies bringt uns einen großen Trost, und kommt dazu auch noch das Zeugniß des heiligen Paulus. Ob er wohl gerne haben wollte, daß wir uns des Ehestandes ganz und gar sollten enthalten, so läßt er es uns doch gleichwohl zu, dieweil es Gott selbst zuläßt, 1 Cor. 7, 2. 9.

118. Man kann aber hieraus auch eine Vermahnung nehmen, wie sich Eheleute mäßig und keusch halten sollen. Denn Etliche sind so gar grobe unflätige Säue, daß sie meinen, sie haben im Ehestande Macht, mit ihren Weibern zu thun alles, was sie nur gelüstet. Dieselbigen irren sehr weit, und sollen wissen, daß ihnen nur dieselbe elende Lust, und daß sie sich freundlich zu den Weibern halten mögen, erlaubt ist; nicht, als wäre solche Lust und Beischlafen rein im Fleische, denn die Eheleute sind von beiden Theilen mit der Erbsünde, Krankheit und Lustseuche vergiftet, sondern dieweil Gott gesagt hat: Ich will dir diese Unreinigkeit nicht zurechnen, sofern du es bei der Gebühr bleiben läßt, die um dich ist.

119. Derhalben müssen wir fein ehrlich und mit Gottesfurcht in diesem Stande leben. Ich habe zwar durch göttliche Erlaubniß ein Eheweib; es ist aber doch in unser beider Fleisch Sünde: die schändliche Lust und Unzucht ist nicht gut. Denn wiewohl von Natur ein Ehemann mit besserem Gewissen sich zu seinem Eheweibe hält, als wenn sonst der Same vergeblich ausgeschüttet würde, oder aber wenn er sich der Hurerei und Ehebruchs befleißigen wollte, so geschieht doch gleichwohl solches nicht ohne Unreinigkeit. Derhalben sind die Männer zu fragen, die mit ihren Weibern in der Schlafkammer nach ihrem Muthwillen schändliche und unsägliche Dinge üben. Denn GOtt hat ihnen eine Sicherheit vergönnt und zugelassen, aber das ist nur eine Erlaubniß und dieser muß man mit nichten mißbrauchen.

120. Und also muß man verstehen den Spruch Augustini, da er sagt: Wer sein Weib liebt, der mag frei sicher des jüngsten Tages dabei erwarten. Wie aber? Antwort: Aus sonderlicher Erlaubniß; wo das nicht wäre, so wäre es ein Ehebruch und Unreinigkeit. Dieweil es aber GOtt also verordnet und Mann und Weib zusammengefüget hat, darum rechnet es ihnen GOtt nicht zu, was daselbst Schändliches oder Unreines geschieht. Man soll aber gleichwohl das Böse im Ehestande und in der ehelichen Beiwohnung nicht vertheidigen, als ob es etwas Gutes sei. Die Heiden haben gesagt, daß es ein natürlich Ding wäre und daß es nicht unehrlich sei, daß etwa ein Jüngling eine Jungfrau lieb habe; wie jener Alte sagt in der Komödie. Aber GOtt will von uns haben, daß wir im Ehestande einen Verstand haben sollen, und daß wir bei Sinnen seien und wissen, daß wir wohl nicht verdammt oder beschuldigt werden sollen am jüngsten Gericht: wir sollen aber doch die böse Lust und Schande, so im Ehestande ist, nicht entschuldigen. Wir sollen nicht sagen: Es ist wohl gethan, daß ich bei meinem Weibe geschlafen habe; sondern wir sollen unsere Unreinigkeit erkennen, und sollen es doch gleichwohl dafür halten, daß im Ehestande Alles rein sei, darum daß GOtt Mann und Weib zusammen gefügt hat, auf daß wir unsere Geschäfte in Furcht und gebührlicher Ehre besitzen mögen. Und wo etwas Böses begangen ist, daß doch solches in den Grenzen des Ehestandes bleiben möge.

121. Also werden wir das Reich Christi recht verstehen und ansehen. Denn es ist ein Anderes, durch das Gesetz gerecht werden, und ein Anderes, aus Gnaden gerecht werden. Ich gebrauche des Gesetzes und Ordnung des Ehestandes, nicht daß ich es werth bin; ja, ich bin nicht werth des Brodes, das ich esse, dazu des Geldes und aller Güter, die ich gebrauche. Denn die Creatur ist unterworfen der Eitelkeit, Röm. 8, 20. Das Weib dient mir als Einem, der dessen nicht werth ist, in allen praedicamentis, wie man in den Schulen nach der Dialectik pflegt zu reden, substantiae, quantitatis, qualitatis, das ist, es sei gleich nach dem Wesen, nach der Größe, oder auch nach der Geschicklichkeit: aber in praedicamento relationis da bin ich es würdig, daß ich bin nicht werth, daß ich ein Weib oder Kinder haben soll, nach dem Wesen und nach dem Dinge an sich selbst; denn ich bin ein armer Sünder: dann werde ich aber dessen würdig, wenn GOtt sagt: Ich will dir es erlauben und will es also geschehen lassen, weil ich es also geordnet habe; du sollst aber dich selbst erkennen und in Gottesfurcht und mit gebührlicher Ehrerbietung in diesem Stande leben.

122. Also sind wir alle nicht werth weder des Ehestandes, noch des Welt- oder Kirchenregimentes. Ich bin nicht werth, daß ich eine Lectio oder Predigt thun soll. Wir haben aber einen gnädigen und barmherzigen GOtt, der uns unsere Sünde und Unwürdigkeit verzeiht, der uns aus dem Reich des Zornes in das Reich der Gnaden versetzt hat, und uns geheißen, daß wir, unter dieser Decke der Gnaden beschützt, zu ihm treten sollen, wie unwürdig wir auch sind, und sollen also der göttlichen Güter durch seine Gnade gebrauchen und derselben genießen.

\section[Dritter Theil]{Dritter Theil. Wie Abimelech den Isaak zur Rede setzt, mit Worten hart straft, und einen Befehl Isaaks wegen ergehen läßt}

\begin{quote}
\textbf{V. 9. Da rief Abimelech dem Isaak, und sprach: Siehe, es ist dein Weib. Wie hast du denn gesagt: Sie ist meine Schwester? Isaak antwortete ihm: Ich gedachte, ich möchte vielleicht sterben müssen um ihretwillen.}
\end{quote}

123. Dies alles ist nur solch Ding, das zum Haus- und Weltregiment gehört, und wird hier nichts von großen und rechten geistlichen Dingen gelehrt, davon das Evangelium predigt und lehrt, als, vom Glauben an Christum, von der Dreifaltigkeit, von der Auferstehung der Todten, von der Ueberwindung der Sünden, des Todes und des Teufels. Darum wird nun dies verachtet von den Heuchlern, die einen großen Geist und Heiligkeit zu haben vermeinen. Was, sagen sie, gehen uns diese fleischlichen Dinge an? Gleich als ob sie ohne Fleisch lebten; gleich als ob sie nicht trinken, nicht essen, nicht schlafen, dazu auch kein männlich oder weiblich Geschlecht an sich haben; endlich gleich als bedürften sie der fleischlichen Dinge gar nicht, oder als ob sie darnach ganz und gar nichts fragten. Aber sie verstehen nicht und wissen auch nicht, was Geist oder was Fleisch sei.

\quellseite{56}

124. Um solcher schwärmerischen Meinungen willen schreibt nun der Heilige Geist solch kindisch Ding und das zum Hausregiment gehört, und zeigt damit an, daß er auch die leiblichen Sachen regiere, ohne welche dieses Leben nicht kann regiert oder erhalten werden. Denn wiewohl die Enthusiasten oder himmlischen Propheten selbst nach fleischlichen Dingen gar nichts zu fragen scheinen, bedürfen sie doch gleichwohl Essens, Trinkens und Kleidung. Warum sollten wir denn davon nicht disputiren und reden, wie man essen und leben soll in rechter Gottseligkeit und mit gutem Gewissen? Denn in der Kirche und Gemeinde Gottes soll man nicht allein vom künftigen Leben lehren; wie davon das Evangelium lehrt, unangesehen dieses gegenwärtigen Lebens, wiewohl es gleichwohl dies Leben auch nicht gar verachtet: sondern wir sollen dasselbe nach dem Gesetz anrichten, und die Leute daraus lehren, wie sie ein gottseeliges und ehrliches Leben führen sollen bis an ihr letztes Ende. Denn wir werden täglich geboren und sterben auch täglich.

125. Darum muß man die Leute lehren und berichten, wie sich das Weib gegen den Mann und der Knecht gegen seinen Herrn halten soll, daß man nicht stehlen, nicht ehebrechen etc. soll. Ja, werden die Enthusiasten und sauren Heiligen sagen: Dies ist aber eitel fleischlich Ding. Antwort: Es ist recht geredet; wo dir es nun nicht gefällt, solcher fleischlichen Dinge zu gebrauchen, so lasse auch Essen, Trinken, Kleider und dergleichen andere Dinge mehr fahren und enthalte dich derselben auch. Diese alle sind Ordnungen und Creaturen Gottes, so er uns zu diesem Leben hat gegeben, welches wir ehrlich und gottseelig führen und zubringen sollen bis in den Tod. Dies muß man sagen um der feindseligen Leute willen, die sich von einem anderen und gleichsam überhimmlischen Leben, hoch über die Wolken hinüber, träumen lassen, so sie doch dieses gegenwärtigen Lebens gar nicht entbehren können.

126. Lasset uns aber nun sehen, was darauf gefolgt sei, daß Isaak vernommen hat, daß Rebekka sein Weib wäre. Abimelech ruft ihn zu sich, und spricht: Ich habe gesehen, daß du mit Rebekka gescherzt hast; daraus kann ich abnehmen und schließen, daß sie ja nicht deine Schwester, sondern dein Weib sein muß; denn es pflegt ja ein Bruder mit einer Schwester also nicht zu scherzen. Warum hast du mir denn vorgelogen und gesagt, es sei deine Schwester? Isaak antwortete ihm: „Ich gedachte, ich möchte vielleicht sterben müssen.“

127. An diesem Ort muß man von der Lüge reden, wie wir droben (im 12. Capitel) gehört, davon Augustinus wider Hieronymus scharf disputirt um der Worte Pauli willen an die Galater im 1. Cap. V. 20., da er sagt: „Was ich euch aber schreibe, siehe, Gott weiß, ich lüge nicht.“ Da sagt er, Paulus habe Petrum nicht mit Ernst gestraft, sondern habe sich nur äußerlich also gestellt, als ob er ihn strafte. Er unterscheidet aber dreierlei Lügen. Eine schädliche Lüge ist, wenn man mit solchem Vorsatz, daß man Einen betrügen will, Lügen redet, dem Nächsten zu schaden entweder an seinem Gute, oder Leumunde, oder aber am Leben. Und ist das noch das Allergärigste, wenn unter dem Namen Gottes Lügen und falsche Lehre ausgebreitet werden, welches im achten Gebote Gott verboten hat, da er sagt: „Du sollst kein falsch Zeugniß reden“ etc., 2 Mos. 20, 16.

128. Die andere Art der Lüge heißt mendacium officiosum, eine Lüge, da jemand mitlügt Michal 1 Sam. 19, 17. ihrem Vater auch vor, daß David sei weggegangen; das thut sie aber dem David zu gute, ihn damit zu retten. Derhalben, wo es ohne der Obrigkeit oder der Eltern Schaden geschehen kann, so mag man die, welche sie suchen oder angreifen wollen, verbergen oder sie schützen. Der gleichen Lüge thut Rahab auch im Buche Josua im 2. Cap. V. 5. Darum ist das eine Nothlüge (mendacium officiosum) damit man Andern hilft und dient, daß sie an ihrem guten Namen oder Leumunde, oder auch an Leib und Leben nicht verletzt oder beschädigt werden.

129. Hingegen die schädliche Lüge schädigt das alles, grade wie die Nothlüge es vertheidigt, und daher nicht eigentlich, sondern uneigentlich und mißbräuchlich Lüge genannt wird; denn sie ist ein sehr schöner Schutz wider die Gefahr Leibes und Lebens, dazu auch des Gutes. Darum ist es eine ehrliche und gute Lüge, und man sollte es vielmehr ein Werk der Liebe nennen; wiewohl es Augustinus eine Lüge nennt, und lindert es doch gleichwohl mit dem Zusatze, da er es mendacium officiosum nennt, das ist, eine Nothlüge. Denn es wird damit der Verfolger betrogen, auf daß der Teufel und die Bosheit des Verfolgers gehindert und der Unschuldige vertheidigt werde. Das heißt die Gebote Gottes halten und nicht übertreten. Ja, möchte Einer sagen, man soll gleichwohl die Wahrheit nicht verschweigen. Darauf antworte ich also, daß man in solchen Fällen die Wahrheit nicht sagen soll, es sei denn, daß du dazu mit Gewalt genöthigt und gezwungen werdest.

130. Die dritte Art der Lüge heißt auf Lateinisch mendacium jocosum, eine Scherzlüge, wenn Einer mit dem Andern scherzt, daß doch gleichwohl damit die Gottseligkeit und der Glaube unverletzt bleibe, und daß immer noch ein solcher Ernst dabei gespürt werde, wie einem Christen geziemt. Wie dieser Scherz ist, so Rebekka und Isaak mit einander gehabt; oder wenn ein Ehemann mit seinem Weibe oder Söhnlein etwa Scherz treibt und sie täuscht, daß sie darüber lachen und fröhlich werden. Wenn solcher Betrug aufgedeckt und offenbar wird, da hört dann die Lüge auf, und ist nichts, denn nur allein ein Scherz und ein solch Ding, darüber man fröhlich und guter Dinge ist. Und dies ist auch eine nützliche Lüge, vornehmlich unter Bekannten, und die einander verwandt und gute Freunde sind.

131. Derhalben wird nun gefragt von dieser That Isaaks: Ob es eine Sünde sei, das er gethan hat? Darauf antworte ich also: Es ist keine Sünde, sondern ist eine Nothlüge, damit er verhütet hat, daß er nicht ermürget würde von denen, bei welchen er sich dazumal aufgehalten hat, wenn er gesagt hätte, daß Rebekka sein Weib wäre; wiewohl es eine Schwachheit des Glaubens ist, daß er es nicht frei öffentlich und beständig bekannt und den Tod verachtet hat. Denn das wäre eine schöne herrliche Tugend und Beständigkeit gewesen, die einem großen tapferen Helden wohl hätte angestanden.

132. Aber GOtt hat haben wollen, daß er etwas schwach sein sollte, auf daß man an ihm ein Exempel hätte, daraus die Kirche und Gemeinde GOttes möge unterrichtet und gelehret werden, nämlich, daß GOtt damit nicht erzürnt wird, du bekennest gleich mit großer Beständigkeit, welches einem tapferen Helden gebührt, oder aber daß du etwas schwach seiest. Denn GOtt will damit Geduld tragen und es dir verzeihen. Und daraus sehen wir endlich, daß wir einen gnädigen GOtt haben, der uns unsere Schwachheit verzeihen, dieselbe übersehen und unsere Sünden vergeben kann, allein, daß wir nicht schädlicher Weise lügen.

133. Es hat aber Isaak nicht geringe Ursache gehabt, daß er sich also gefürchtet. Denn wiewohl er in dem Lande einen König gefunden, der gerecht, heilig und fromm genug war, hat er doch also gedacht: Siehe, ich werde der Theurung halben genöthigt, daß ich unter den Heiden werde müssen ein Fremdling sein, welche diese Verheißung nicht haben, die ich habe, nämlich, daß der Herr mein GOtt sein wolle, und daß ich glaube, und im Glauben an ihn bleiben und verharren soll. Diese Leute haben weder Wort noch Glauben. Derhalben ist wohl zu vermuthen, oder aber man muß es ja zum wenigsten befürchten, daß es gottlose Leute sind, und die mit Mord, Ehebruch und Hurerei umgehen und mit solchen Sachen behaftet sein müssen. Und wenn sie schon jetzt mit der That nicht Ehebrecher und Mörder sind, so möchte doch leichtlich eine Gelegenheit sich bieten, dadurch sie in solche Sünden zu fallen könnten bewogen werden. Nun ist aber mein Weib schön von Angesicht, und könnte irgend einen gottlosen Mann damit reizen, daß er sie ihrer Schöne halben lieb gewinnen möchte; darum stehe ich in großer Gefahr.

\quellseite{57}

134. Dies ist eine rechtmäßige Furcht, und eine solche Furcht, die einem ehrlichen beständigen Manne wohl anstehen kann, wie die Juristen oder Rechtsverständigen pflegen zu reden. Wenn sich ein verzagter und erschrockener Mann fürchtet, daß ist eine unnütze Furcht; aber die Furcht, so etwa ein tapferer Held oder beherzter Mann hat, ist rechtmäßig und vernünftig. Derhalben wird eine solche Furcht gelobt von den Juristen und auch von den Theologen. Denn das ist nicht das Zeichen eines tapferen Helden oder starken beherzten Mannes, daß er sich gar vor keinerlei Gefahr fürchten oder besorgen, und sich tollkühn in unnöthige Gefahr begeben wollte. Die das thun, sind tollkühne und verwegene Leute und nicht tapfer oder beherzt. Wer recht stark und beherzt ist, der weiß, wie er die Kühnheit und die Furcht in Schranken und damit das rechte Maß halten soll. In nöthiger Gefahr fürchtet er sich nicht, da es die Noth, die Ehrbarkeit und gemeiner Nutzen erfordert, daß man einen starken und unverzagten Muth haben soll; aber außerhalb der Noth bleibt er sich nicht so leichtlich oder tollkühn in Gefahr. Und also mag sich ein frommer ehrlicher Mann wohl billig fürchten. Denn das steht einem beherzten Manne und der beständig ist, nicht übel an. Auf solche Weise fürchtet sich nun Isaac auch und ist schwach, und wird doch gleichwohl entschuldigt, dieweil er furchtsam ist und ein beständiger Mann.

135. Ich lasse mir es aber gefallen, daß wir aus diesem Text Ursache nehmen, zu disputiren vom Zweifel, da man an GOtt und seinem Willen zweifelt. Denn ich höre, daß hin und wieder unter denen vom Adel und andern großen Herren etwa böse Worte fallen und ausgebreitet sollen werden von der Prädestination oder Versehung Gottes. Denn also sollen sie reden: Wenn ich versehen bin, so mag ich Gutes oder Böses thun, ich werde doch selig werden; bin ich aber nicht versehen, so muß ich verdammt werden, unangesehen meiner Werke.

136. Wider solche gottlose Worte wollte ich gerne lange disputiren, wenn ich es meiner schwankenden Gesundheit halben thun könnte. Denn wenn die Worte wahr sind, wie sie sich dünken lassen, so wird damit die Menschwerdung des Sohnes Gottes, sein Leiden und Auferstehung, und alles, was er gethan hat um der Welt Heil und Seligkeit willen, gar aufgehoben und hinweg genommen. Was werden uns dann die Propheten und die ganze heilige Schrift nütze sein? Wozu dienen uns dann die heiligen Sacramente? Derhalben lasset uns nur dies alles verwerfen und mit Füßen treten.

137. Dieß sind teuflische und vergiftete Pfeile, und ist eben die Erbsünde selbst, damit der Teufel unsere ersten Eltern verführt hat, da er sagte: „Ihr werdet sein, wie GOtt“, 1 Mos. 3, 5. Denn sie ließen sich mit der Gottheit, so ihnen geoffenbart war, durch welcher Erkenntniß sie selig waren, nicht genügen, sondern sie wollten in die Tiefe der Gottheit eindringen. Denn sie haben also vernünftelt: Es müßte irgend eine heimliche Ursache, warum ihnen GOtt verboten hätte, daß sie von der Frucht des Baumes, so mitten im Paradies war, nicht essen sollten, darunter verborgen sein; dieselbige Ursache wollten sie wissen.

138. Wie diese Leute jetziger Zeit auch davon reden: Was GOtt versehen hat, das muß geschehen, darum ist es alles ungewis und vergeblich, daß wir uns um die Religion oder der Seelen Seligkeit viel bekümmern sollten. Es ist aber dir nicht befohlen, daß du davon urtheilen sollst; denn das Urtheil oder Gericht Gottes ist unerschöpflich. Warum zweifelst du oder verwirfst den Glauben, den dir GOtt geboten hat? Denn wozu dient es, daß GOtt seinen Sohn gesandt hat, daß er für uns leiden und gekreuziget sollte werden? Wozu ist es nütze gewesen, daß er die heiligen Sacramente eingesetzt hat, wenn es alles ungewiß und gar vergeblich ist zu unserer Seligkeit? Denn sonst, wo jemand wäre versehen worden, der wäre ohne den Sohn und ohne die Sacramente oder die heilige Schrift selig geworden. Derhalben so muß GOtt nach dieser Leute Lästerung ein greulicher Narr gewesen sein, daß er seinen Sohn gesandt, das Gesetz und Evangelium gegeben und die Apostel gesandt hat, wenn er nur das hat haben wollen, daß wir ungewiß sein und daran noch zweifeln sollten, ob wir selig oder verdammt werden.

139. Aber dies ist des Teufels Gespenst und Betrug, dadurch er sich untersteht, uns zweifelhaftig und ungläubig zu machen; so doch Christus darum in diese Welt gekommen ist, daß er uns der Seligkeit hat wollen ganz gewiß machen. Denn es muß endlich auf solche lästerliche Meinung entweder Verzweiflung folgen, oder Verachtung Gottes, der heiligen Bibel, der Taufe und aller göttlichen Wohlthaten, damit er uns wider den Zweifel hat stärken wollen und daß wir ja unserer Seligkeit nicht ungewiß sein sollten. Denn die Lästerer werden doch zuletzt mit den Epicureern sagen: Lasset uns im Saufe leben, essen und trinken, wir werden doch morgen sterben müssen. Sie werden, wie die Türken pflegen zu thun, tollkühn in das Schwert und Feuer fallen, sintemal, wie sie meinen, die Stunde schon bestimmt ist, darin du entweder niedergelegt und erschlagen werdest, oder aber davon kommen mögest.

140. Gegen diese Gedanken aber soll man die wahre und gewisse Erkenntniß Christi halten; wie ich oft vermahne, daß es vornehmlich nütze und nöthig sei, daß ja die Erkenntniß Gottes in uns ganz gewiß sein möge, daß wir es im Herzen gewiß fassen und fest daran hangen; denn sonst wird unser Glaube vergeblich und umsonst sein. Denn wo GOtt seine Verheißungen nicht gewiß hält, so ist es mit unserer Seligkeit gar verloren. Da hingegen unser Trost ist: daß ob wir uns wohl ändern, wir zu dem unsere Zuflucht haben, der sich nicht ändert, sondern immer beständig bleibt. Denn also sagt er von sich selber im Propheten Maleachi am 3. Cap. V. 6.: „Ich bin der HErr, der nicht lüget.“ Und St. Paulus sagt Röm. 11, 29.: „Gottes Gaben und Berufung mögen ihn nicht gereuen.“

141. Darum habe ich in meinem Buche De servo arbitrio (daß der freie Wille nichts sei) wider Erasmus, und an andern Orten mehr also gelehrt, daß man unterscheiden soll, wenn man von der Erkenntniß oder vielmehr von dem Wesen der Gottheit handelt. Denn man muß entweder disputiren von dem verborgenen GOtt, oder aber von dem geoffenbarten GOtt. Von GOtt, sofern er nicht geoffenbart ist, ist kein Glaube und keine Erkenntniß und man kann von solchem GOtt nichts wissen, und da muß man sich nach dem Spruche halten: Quae supra nos, nihil ad nos: Was über uns ist, das geht uns nichts an. Denn solche Gedanken, die über oder außerhalb der Offenbarung Gottes etwas Höheres forschen wollen, sind gar teuflische Gedanken, damit man nichts mehr ausrichten kann, denn daß wir uns selber in das Verderben hinein stürzen; denn sie halten uns einen solchen Gegenstand vor, der unerforschlich ist, nämlich, GOtt, der nicht geoffenbart ist. Man lasse viel lieber GOtt seine Schlüsse und Geheimnisse im Verborgenen behalten. Wir dürfen uns darum, daß uns dieselbigen sollen geoffenbart werden, so sehr nicht bemühen.

142. Mose begehrte auch 2 Mos. 33, 18., daß ihn GOtt sein Angesicht oder Herrlichkeit wollte sehen lassen; aber der HErr antwortete ihm darauf also, V. 20.: „Du wirst mir hinten nachsehen, aber mein Angesicht kann man nicht sehen.“ Denn dieser Vorwitz ist die Erbsünde selbst, dadurch wir getrieben und gereizt werden, daß wir durch natürliche Speculation einen Weg suchen zu GOtt. Es ist aber eine große Sünde und ist ein unnütz und vergeblich Ding, daß man sich dessen unterstehen will. Denn also spricht Christus Joh. 6, 65. Cap. 14, 6.: „Niemand kommt zum Vater, denn durch mich.“ Darum wenn wir zu dem GOtt, der sich nicht geoffenbart hat, treten, so ist da kein Glaube, kein Wort oder gar keine Erkenntniß. Denn es ist ein unsichtbarer GOtt, den wirst du nicht sichtbar machen.

143. Darnach hat GOtt auch ganz ernstlich verboten, daß man sich dermaßen nicht soll gelüsten lassen, seine Gottheit zu erkennen. Wie Christus zu den Aposteln sagt, Ap. Gesch. 1, 7., da sie fragten: „Herr, ist es nicht also versehen, daß auf diese Zeit das Reich Israel soll wieder aufgerichtet werden?“ – „Es gebühret euch nicht zu wissen Zeit oder Stunde“ \&c. Laß mich, spricht GOtt, verborgen sein, da ich mich dir nicht geoffenbart habe, oder du wirst dir selbst eine Ursache sein deines eigenen Verderbens, wie Adam greulich gefallen ist. „Denn wer schwere Dinge forschet, dem wird es zu schwer“, Spr. 25, 27.

144. Und anfänglich hat GOtt zwar alsbald diesem Vorwize wollen zuvorkommen. Denn also hat er uns seinen Willen und Rath vorgehalten und sagt nämlich also: Siehe, Mensch, ich will dir die Versehung und Prädestination herrlich offenbar; aber nicht auf dem Wege deiner Vernunft und fleischlicher Weisheit, wie du dir träumen lässest und denkst. Ich will also thun: aus einem GOtt, der nicht geoffenbart ist, will ich ein geoffenbarter GOtt werden, und will doch derselbige GOtt bleiben. Ich will Mensch werden, oder will meinen Sohn senden, der soll für deine Sünde sterben und wieder vom Tode auferstehen; und also will ich deine Begierde erfüllen, auf daß du wissen mögest, ob du verstehen seiest oder nicht. „Siehe, das ist mein Sohn, den sollst du hören“, Matth. 17, 5., den siehe an, wie er in der Krippe liegt und auf der Mutter Schooße, dazu auch, wie er am Kreuze hängt; siehe, was derselbe thue, was er rede. Da wirst du mich gewißlich ergreifen. Denn „wer mich siehet“, spricht Christus Joh. 14, 9., „der siehet den Vater.“ Wo du diesen hören und in seinem Namen getauft werden, dazu sein Wort lieben wirst, alsdann bist du gewißlich versehen und deiner Seligkeit ganz gewiß. Wenn du aber das Wort lästerst oder verachtest, so bist du verdammt: „Denn wer nicht glaubet, der wird verdammt werden“, Marc. 16, 16.
\quellseite{58}

145. Die andern Gedanken und Wege, so deine Vernunft oder Fleisch dir vorgibt, sollst du tödten. Denn GOtt ist denselben feind. Das allein lasse dir angelegen sein, daß du meinen Sohn annehmest, daß dir Christus in deinem Herzen wohlgefalle in seiner Geburt, Wunderwerken und Kreuz. Denn da ist das Buch des Lebens, darin du geschrieben bist.

146. Und dieses ist der einzige und gewisse Rath wider diese greuliche Seuche, damit die Menschen immer nach ihrer Speculation und hohen Gedanken, GOtt in seiner hohen Majestät zu erforschen, wollen fortfahren, und endlich darüber in Verzweiflung oder Verachtung GOttes fallen. Willst du nun der Verzweiflung, der Feindschaft und Lästerung GOttes entfliehen, so laß die Speculation und hohe Gedanken von dem verborgenen GOtt fahren, und höre auf, dessen vergeblich zu begehren und darnach zu streben, daß du das Angesicht oder die Herrlichkeit GOttes sehen wollest; sonst wirst du für und für im Unglauben und Verdammniß hängen bleiben und verloren werden. Denn wer da zweifelt, der glaubt nicht, und wer nicht glaubt, wird verdammt werden.

147. Darum sollen wir diesen schändlichen bösen Worten feind sein und uns davor hüten, welche die Epicurer führen, wenn sie sagen: Wo dieses ja nothwendig also gesehen muß, so geschehe es immerhin. Denn GOtt ist nicht darum vom Himmel herab gekommen, daß er sich der Versehung wolle ungewiß machen, und daß er dich lehrete die Sacramente, Absolution und andere göttliche Ordnungen mehr verachten; ja, er hat dieses alles darum eingesetzt, daß er dich damit wollte ganz gewiß machen, und aus deinem Herzen den großen Mangel und Fehler des Zweifels wegnehmen, auf daß du nicht allein im Herzen glauben, sondern auch mit leiblichen Augen sehen und dazu mit den Händen greifen mögest. Warum verwirfst du nun dieses alles, und klagst, daß du nicht wissen könnest, ob du zur Seligkeit versehen seiest? Du hast ja das Evangelium, bist getauft, hast die Absolution, bist ein Christ und zweifelst doch noch und sagst, du wissest nicht, ob du glaubest oder nicht glaubest, ob du das auch für wahrhaftig haltest, was dir im Wort und Sacramenten von Christo gesagt und gepredigt wird.

148. Aber du wirst vielleicht sagen: Ich kann nicht glauben; wie deren viele sind, die mit dieser Anfechtung geplagt werden. Und erinnere ich mich, daß zu Torgau ein armes Weiblein zu mir gekommen, und mit bitteren Thränen geklagt hat, daß sie nicht glauben könnte. Da ich nun die Artikel des Glaubens nach einander hersagte, und sie um einen jeglichen Artikel fragte: Ob sie es auch dafür hielte, daß dies alles wahr und also auch geschehen wäre, oder nicht? antwortete sie mir und sagte: Ich halte allerdings dafür, daß es wahr sei; ich kann es aber ja nicht glauben. Das war ein teuflischer Betrug. Darum sagte ich zu ihr: Wenn du dies alles für wahr hältst, so hast du keinen Grund, über deinen Unglauben zu klagen; denn wo du daran nicht zweifelst, daß der Sohn GOttes für dich gestorben ist, so glaubst du ja gewißlich. Denn Glauben ist Anderes nichts, als daß man dies für gewisse unzweifelhafte Wahrheit halte.

149. GOtt sagt zu dir: Siehe, da hast du meinen Sohn, den höre und nimm ihn an; wenn du das thust, so bist du jetzt schon deines Glaubens und deiner Seligkeit gewiß. Ja, sagst du, ich weiß aber nicht, ob ich im Glauben bleiben kann? Ei, so nimm doch gleichwohl die gegenwärtige Verheißung und Versehung an, und hüte dich, daß du nicht vorsätzlich oder zu genau nach den heimlichen Rathschlüssen GOttes forschest. Wenn du an den geoffenbarten GOtt glaubst und sein Wort annimmst, so wird dir allemal auch der verborgene GOtt geoffenbart werden. Denn „wer mich siehet“, spricht Christus Joh. 14, 9., „der siehet den Vater.“ Wer aber den Sohn verwirft, der verleugnet mit dem geoffenbarten GOtt auch den verborgenen GOtt, der sich nicht geoffenbart hat. Wirst du aber mit starkem Glauben dem geoffenbarten GOtt anhangen, also daß du in deinem Herzen gesinnet seiest, du wollest Christum nicht verlieren, wenn du auch sonst alles, was du hast, solltest beraubet werden, so bist du gewißlich versehen, und wirst den verborgenen GOtt verstehen, ja, du verstehst ihn jetzt schon allebereit: wenn du den Sohn erkennst und seinen Willen, daß er sich dir offenbaren und dein Herr und Heiland sein wolle, so bist du dessen gewiß, daß GOtt auch dein Herr und dein Vater sei.

150. Siehe doch, wie dich GOtt so freundlich und gnädig von dieser greulichen Anfechtung erlöset, welche der Satan jetziger Zeit über die Maßen hart treibt, auf daß er die Leute zweifelhaft und ungewiß mache, dazu endlich auch von GOttes Wort gar abwende. Denn warum wolltest du das Evangelium hören, sagen die Epicurer, dieweil es doch alles an der Versehung gelegen ist? Also nimmt uns der Satan mit Gewalt hinweg die Versehung, deren wir vergewissert sind durch den Sohn GOttes und durch die heiligen Sacramente, und macht uns ungewiß, da wir doch zuvor ganz gewiß sind: und wenn er die armen erschrockenen Gewissen mit dieser Anfechtung angreift, so sterben sie in der Verzweiflung dahin; wie mir auch schier geschehen wäre, wo mich Staupitz nicht errettet hätte, da ich eben dieselbige Anfechtung hatte. Wenn es aber Verächter sind, die also angefochten werden, so werden die allergärigsten und schändlichsten Epicurer daraus.

151. Darum sollen wir viel lieber diese Sprüche in unser Herz prägen, als, Joh. 6, 44., da Christus spricht: „Es kann niemand zu mir kommen, es sei denn, daß ihn ziehe der Vater, der mich gesandt hat.“ Durch wen ziehet er ihn aber? Antwort: Durch mich. „Denn wer mich siehet, der siehet den Vater.“ Und zu Mose spricht GOtt: „Mein Angesicht kannst du nicht sehen; denn kein Mensch wird leben, der mich siehet“, 2 Mos. 33, 20. Item, Ap. Gesch. 1, 7. spricht Christus zu den Jüngern: „Es gebühret euch nicht zu wissen Zeit oder Stunde, welche der Vater seiner Macht vorbehalten hat“; gehet ihr aber hin, und richtet aus, was ich euch befohlen habe. Item, Sirach spricht am 3. Cap. V. 22. 23. 24.: „Stehe nicht nach höherm Stande, und denke nicht über dein Vermögen; sondern was GOtt dir befohlen hat, des nimm dich stets an“ etc., „und was deines Amts nicht ist, da laß deinen Vorwitz.“ Höre meinen Sohn, sagt GOtt, der Mensch geworden ist, so wird die Versehung von sich selber kommen.

152. D. Staupitz pflegte mich mit diesen Worten zu trösten und sagte zu mir also: Lieber, warum plagst du dich also mit diesen Speculationen und hohen Gedanken; schaue an die Wunden Christi, und sein Blut, das er für dich vergossen hat, daraus wird die Versehung hervorscheinen. Derhalben soll man den Sohn GOttes hören, der in das Fleisch gesandt, Mensch geworden, und darum erschienen ist, daß er die Werke des Teufels zerstöre, 1 Joh. 3, 8., und dich der Versehung gewiß mache. Und darum sagt er auch zu dir: Du bist mein Schäflein; denn du hörst meine Stimme; und niemand wird dich mir aus meiner Hand reißen, Joh. 10, 29.

153. Es sind deren Viele, die nicht auf solche Weise dieser Anfechtung widerstrebet, darüber in Verderben und ewige Verdammniß gestürzet worden. Derhalben muß man die Herzen gottseliger Leute fleißig stärken, daß sie dagegen immer gerüstet sein mögen. Wie ein Grenit oder Einsiedler in dem Buche „Leben der Väter“ ermahnt, daß sie sich solcher Speculation und hohen Gedanken entschlagen und enthalten sollen, und sagt also: Wenn du siehst, daß irgend Einer seinen Fuß in den Himmel gesetzt hat, so ziehe ihn wiederum zurück. Denn also pflegen die Heiligen oder Christen, die noch Neulinge sind, außerhalb Christo von GOtt zu denken; und die sind es, die sich in den Himmel zu steigen unterstehen und daselbst beide Füße hinzusetzen; aber sie werden bald in die Hölle geworfen und versenkt.

\quellseite{59}

154. Darum sollen sich die Gottseeligen davor hüten, und sich dessen allein befleißigen, daß sie lernen dem Kindlein und dem Sohne Gottes, Jesu, anhangen, welcher dein Gott ist und um deinetwillen Mensch geworden ist; denselbigen sollst du erkennen und hören, dazu deine Lust an ihm haben und ihm auch dafür danken. Wenn du den hast, so hast du auch den verborgenen Gott zugleich mit dem geoffenbarten. Und das ist der einige Weg, die Wahrheit und das Leben: außerhalb desselbigen Weges, Wahrheit und Lebens wirst du nichts denn eitel Verderben, Verdammniß und Tod finden.

155. Darum hat er sich im Fleische geoffenbart, auf daß er uns aus dem Tode, aus dem Fleische und aus des Teufels Gewalt reißen und erlösen möchte. Aus solcher Erkenntniß muß gewißlich große Freude und Lust kommen, daß Gott unwandelbar ist und daß er nach unwandelbarer Nothwendigkeit wirkt, und sich selbst nicht leugnen kann, sondern seine Verheißung treulich hält.

156. Darum ist es uns nicht frei, daß wir mit solchen hohen Gedanken umgehen und zweifeln an der Verheißung; sondern dieselbigen Gedanken sind gottlos, böse und teuflisch. Darum wenn dich der Teufel damit anficht, so sage nur: „Ich glaube an Jesum Christum, unsern Herrn“, an dem ich keinen Zweifel habe, daß er Mensch geworden, gelitten hat und für mich gestorben ist, in dessen Tod ich getauft bin. Mit dieser Antwort wird die Anfechtung verschwinden und der Satan wird dir den Rücken zukehren. Wie ich an andern Orten oft ein merklich Exempel von einer Nonne erzählt habe, welche auch eben dieselbige Anfechtung gehabt hat. Denn unter dem Papstthum sind auch viele gottselige Leute gewesen, so diese geistlichen Anfechtungen gefühlt haben, welches rechte höllische und der verdammten Menschen Gedanken sind; denn es ist kein Unterschied zwischen einem, der da zweifelt, und einem Verdammten. Derhalben, so oft dieselbige Nonne gefühlt hat, daß sie mit den feurigen Pfeilen Satans ist angegriffen worden, hat sie Anderes nichts gesagt denn: Ich bin eine Christin.

157. Also müssen wir auch thun. Man muß das Disputiren unterwegen lassen und sagen: Ich bin ein Christ, das ist, der Sohn Gottes ist Mensch geworden und auf diese Welt gebohren, der hat mich erlöset und sitzt zur Rechten des Vaters und ist mein Heiland. Also treibe den Teufel von dir weg, mit so wenig Worten, als du immer kannst, und sprich: „Heb dich weg von mir, Satan“, Matth. 4, 10., mache mir keinen Zweifel; der Sohn Gottes ist in diese Welt gekommen, daß er dein Werk und Zweifel zerstöre; da hört denn die Anfechtung auf, und das Herz kommt wieder zum Frieden, Ruhe und Liebe Gottes.

158. Wo man aber zweifelt an irgend eines Menschen Wollen, das ist keine Sünde; wie Isaak zweifelt, ob er auch lebendig bleiben würde oder einen guten frommen Wirth haben. An einem Menschen mag ich und soll auch zweifeln; denn er ist mein Heiland nicht. Und im 146. Psalm V. 3. steht geschrieben: „Verlasset euch nicht auf Fürsten.“ Denn alle Menschen sind Lügner, und können nicht helfen. Aber mit Gott kann man nicht so im Zweifel handeln. Denn er will und kann auch nicht wandelbar oder ein Lügner sein; sondern der höchste Gottesdienst, den er fordert und haben will, ist, daß du ihn für wahrhaftig haltest. Denn darum hat er so starke Argumente und Wahrzeichen seiner Wahrheit gegeben und daß Alles mit ihm eitel gewiß Ding ist. Er hat seinen Sohn in das Fleisch und in den Tod gegeben, hat dazu die Sacramente eingesetzt, daß du wissen sollst, daß er kein Lügner, sondern wahrhaftig ist.

159. Und das beweist und bestätiget er nicht mit geistlichen, sondern mit greiflichen Argumenten und Wahrzeichen. Denn ich sehe ja das Wasser in der Taufe; ich sehe das Brod und Wein im Abendmahl; ich sehe den Diener des Wortes: welches ja alles leiblich ist; in welchen leiblichen Figuren oder Bildern er sich offenbart. Wenn man mit Menschen handeln soll, da mag man zweifeln, was oder wie weit du Einem glauben mögest und wie Andere gegen dich mögen gesinnet sein; aber von Gott sollst du das gewiß und ungezweifelt halten, daß er dir um Christi willen gnädig sei, und daß du durch das theure Blut des Sohnes Gottes erlöst und geheiliget seiest; und also wirst du deiner Versöhnung auch gewiß sein; du wirst alle vorwitzigen und gefährlichen Fragen von den heimlichen Ratschlüssen Gottes fahren lassen, zu welchen der Teufel sich unterstehet uns zu treiben, wie er dazu unsern ersten Vater Adam getrieben und gebracht hat.

160. O wie wäre er aber so selig gewesen, wenn er Gottes Wort fleißig vor Augen gehabt, und von allen andern Bäumen gegessen hätte, den allein ausgenommen, davon ihm verboten war zu essen! Aber er wollte darnach forschen, was Gott damit gemeint hätte, daß er ihm verboten hatte, des einzigen Baumes nicht zu gebrauchen und davon nicht zu essen. Dazu ist auch gekommen der schalkhaftige Meister, der Satan, welcher den Vorwitz gemehret und dazu weiblich geholfen hat. Also ist er darüber in Sünde und Tod gestürzet worden.

161. Also offenbart uns Gott seinen Willen durch Christum und das Evangelium. Das verachten wir aber, und lassen uns nach dem Exempel Adams vor allen andern Bäumen des verbotenen Baumes gelüsten. Dieses Gebrechen haben wir von Natur alle: wenn das Paradies und der Himmel zugeschlossen, und der Engel daselbst hingestellt ist, dasselbe zu bewahren, so unternehmen wir es vergeblich, dahin einzukommen. Denn Christus hat recht gesagt Joh. 1, 18.: „Niemand hat Gott je gesehen.“ Und dennoch hat sich Gott uns aus unermesslicher Gnade geoffenbart, auf daß er unsere Begierde erfüllen und derselbigen Genug thun möchte. Er hat uns ein sichtbarlich Bild dargestellt, und sagt: „Siehe, da hast du meinen Sohn; wer den hört und getauft wird, der ist in das Buch des Lebens geschrieben; das offenbare ich dir durch meinen Sohn, welchen du mit den Händen kannst angreifen und mit den Augen sehen.“

162. Dies habe ich also gerne und mit Fleiß vernommen wollen. Denn nach meinem Tode werden Viele meine Bücher hervorbringen und dieselben anführen, und werden daraus allerlei Irrthümer und ihre eigene Phantasei beweisen und bestätigen wollen. Nun habe ich aber unter andern geschrieben, daß Alles nothwendig sei und geschehen müsse; ich habe aber zugleich das auch dazu gesetzt, daß man den Gott, der sich geoffenbart hat, ansehen soll, wie wir im 46. Psalm singen: „Er heißt Jesus Christ, der Herr Zebaoth, und ist kein ander Gott“, und sonst an vielen andern Orten. Aber sie werden an den Orten allen vorübergehen, und nur die annehmen, da von dem verborgenen Gott gehandelt wird.

163. Derhalben sollt ihr daran denken, die ihr auch jetzt hört, daß ich dies also gelehrt habe, nämlich, daß man nicht forschen solle nach der Versöhnung des verborgenen Gottes, sondern daß man sich an derselben Versöhnung solle genügen lassen, so durch die Berufung und durch das Predigtamt offenbart wird; denn daselbst kannst du deines Glaubens und Seligkeit gewiß sein, und sagen: Ich glaube an den Sohn Gottes, der da gesagt hat: „Wer an den Sohn glaubet, der hat das ewige Leben“, Joh. 3, 36. Darum ist an ihm keine Verdammniß oder Zorn, sondern ein Wohlgefallen Gottes des Vaters. Und eben dasselbige habe ich auch also in meinen andern Büchern öffentlich gelehrt und lehre es jetzt auch noch mit lauter Stimme. Darum will ich entschuldiget sein.

\begin{quote}
\textbf{V. 10. 11. Abimelech sprach: Warum hast du denn uns das gethan? Es wäre leicht geschehen, daß jemand vom Volk sich zu deinem Weibe gelegt hätte; und hättest also eine Schuld auf uns gebracht. Da gebot Abimelech allem Volk, und sprach: Wer diesen Mann oder sein Weib antastet, der soll des Todes sterben.}
\end{quote}

164. Der König zeigt sich hier als ein sonderlich frommer Mann, der die Gerechtigkeit und Gottseligkeit lieb hat. Darum sagt Petrus Ap. Gesch. 10, 34. 35.: „Gott stehet die Person nicht an, sondern in allerlei Volk, wer ihn fürchtet und recht thut, der ist ihm angenehm.“ Denn wiewohl dieser König Abrahams Sohn nicht ist, so ist er dennoch heilig und hofft auf die Gerechtigkeit Gottes, wie Cornelius auch gethan hat. Darum ist Isaak zu ihm gekommen, wie zu dem vorigen Abimelech Abraham gekommen ist, und er ist von Isaak in rechter Gottseligkeit und Erkenntniß Gottes unterrichtet worden; denn der wird ohne Zweifel daselbst Gottes Wort gepredigt haben. Es läßt sich aber ansehen, daß der Name Abimelech der Königsname gewesen sei und nicht der Name seiner Person; wie der Name Pharao auch ein gemeiner Name gewesen ist aller Könige in Egypten.

165. Siehe aber, wie heftig feind der König dem Ehebruch ist. Und ist derhalben dies ein trefflicher Text, da er sagt: „Warum hast du denn uns das gethan?“ Als wollte er sagen: Wolltest du mich und das ganze Königreich also betrügen, daß mir unwissend mit solcher schweren Sünde GOTT erzürnen sollten? Denn nicht allein habe ich, sondern auch mein ganz Königreich in Gefahr gestanden, wo jemand von den Unterthanen sich zu ihr gelegt hätte. Dies wird er wahrlich mit großem Zorn gesagt haben; und ich halte dafür, daß er von solchem Beischlafen rede, das ehrlich ist, und von der Bewohnung Mannes und Weibes, und nicht von unziemlicher Vermischung, davor sich Isaak vornehmlich gefürchtet hat. Ich halte aber dafür, daß dieses des Königs Meinung gewesen sei, daß er also habe sagen wollen: Es hätte irgend Einer dein Weib wegnehmen mögen, als ob sie frei wäre, und hätte sie mögen heimführen und sich zu ihr legen, wie sich ein Ehemann zu seinem Weibe pfleget zu halten; so wäre doch daselbst eine heimliche Sünde des Ehebruchs: und ob er wohl ehrlich bei ihr wohnete wie ein Ehemann, so würde er dennoch erfunden und möchte auch billig für einen Ehebrecher gehalten werden.

\quellseite{60}

166. Das sind Worte eines Mannes, der GOTT fürchtet und gottselig ist. Und steht solches einem Fürsten und Könige wohl an, der auch dem heimlichen Ehebruche feind ist, und ist glaublich, daß er noch viel mehr den öffentlichen Ehebruch und unziemliche Vermischung werde gestraft haben. Und wenn sie schon einem Andern ohne des Königes Wissen wäre zugefüget worden, so hätte er doch daran keine Schuld gehabt.

167. An unsern Königen und Fürsten wirst du zwar jetziger Zeit nichts finden, daß diesem Exempel möchte gleich sein. Die sollten wohl sagen: Sie mag seine Schwester oder sein Weib sein, man mag ehrlich oder unehrlich mit ihr umgehen, was geht es mich an? Aber an diesem Abimelech spürt man, daß er die Furcht GOTTes gehabt hat; darum fürchtete er sich, daß in seinem Königreiche auch kein heimlicher Ehebruch möchte begangen werden. Also werden ohne Zweifel die andern Fürsten und Regenten im Lande auch gesinnet gewesen sein.

168. Derhalben ist es eine große Gnade und Wohltat gewesen, daß zu derselbigen Zeit GOTT so fromme und gottselige Herren gegeben hat, so die Kirche und Gemeinde GOTTes beherbergt haben, wie droben (Cap. 14, V. 13.) Abraham Escol, Aner und Abimelech gehabt hat. Denn es muß ja die Kirche und Apostel irgend einen Winkel haben, da sie leben und sich aufhalten können. Ein solcher König ist dieser Abimelech gewesen, der den heiligen Apostel und Propheten GOTTes, Isaak, aufgenommen und gepflegt hat. Darum halte ich dafür, daß er der frommen heiligen Könige einer gewesen sei. Denn wiewohl er kein Feind der Verheißung gewesen, so ist er dennoch derselbigen auch theilhaftig geworden, wie zu derselbigen Zeit auch viel Heiden sind selig geworden.

169. Das hebräische Wort ascham hat Hieronymus im Lateinischen verdolmetscht: peccatum grande, eine große Schuld. Es wird aber damit angezeigt, daß zu der Zeit der Ehebruch für keinen Scherz oder Kinderspiel ist gehalten worden; wie heutiges Tages unter den Christen der Ehestand für einen Scherz gehalten wird, vornehmlich bei den großen Herren. Und die Papisten halten den Ehebruch und die Hurerei für lauter Scherz und treiben ihr Gespött damit. Man verstehe es nun gleich vom Ehebruche oder aber von Hurerei, so kann man wahrlich daraus abnehmen, daß des Königs Herz voll Gottseligkeit und Furcht GOTTes muß gewesen sein, der dieser Sünden keine hat können ungestraft hingehen lassen; denn er wird sehr zornig darüber: er hält es wahrlich für keinen Scherz, sondern hält es für eine Schuld, da unser Herr GOTT einen Zuspruch zu hat (auf Hebräisch ascham, auf Lateinisch reatus); wie wir im Vater Unser bitten: „Vergib uns unsere Schuld“, und im 34. Psalm V. 22.: „Die den Gerechten hassen, werden Schuld haben.“ Daher kommt es, daß hostia so viel heißt als delictum, Schuldopfer; und im 5. Psalm V. 11.: Judica illos, Deus: „Schuldige sie, GOTT“; da hat es Hieronymus im Lateinischen verdolmetscht: Condemna eos, non justifica eos: Verdamme sie, mache sie nicht gerecht, sondern halte sie für Verdammte, die da schuldig und verloren sind. Es lautet schier eben also, wie das Wort, das St. Paulus 1 Cor. 4, 13. gebraucht, auf Griechisch katharma, auf Deutsch Sündopfer.

170. Derhalben sagt nun der König also: „Es wäre leicht geschehen, daß du hättest also eine Schuld auf uns gebracht.“ Er scheut und fürchtet sich nicht allein vor GOTT, sondern erschrickt auch vor solcher Gefahr, welche daher über ihn und das ganze Königreich gekommen wäre, wenn sie ein Anderer weggeführt oder zum Eheweibe genommen hätte.

171. Darnach läßt er es dabei noch nicht bleiben, daß er Isaak hart gestraft oder gescholten hat, sondern bezeugt es auch noch mit einem sonderlichen Edict oder Befehl, daß er darum ein Mißfallen gehabt, und daß er es ernstlich strafen wolle, die diesen fremden Gästen werden Gewalt thun. Um dieser Fremdlinge willen macht er ein eigen Gesetz, und setzt auch eine Leibesstrafe dazu, so über die gehen soll, welche das Gebot übertreten werden; vermehret noch dazu das Gebot, da er sagt: „Wer nicht allein das Weib, sondern auch den Mann wird antasten“; wer ihn an seinem Leibe, Hausgesinde, an seiner Habe und gutem Namen verletzen oder beschädigen wird, „der soll des Todes sterben“; denn ich will haben, daß er in meinem Königreiche von allerlei Gefahr frei und sicher sein soll.

172. Da siehst du nun aber wohl, wie gütig und gnädig sich der Herr gegen seine Heiligen pfleget zu erzeigen: er versuchet sie zwar wohl, schickt sie in das Elend und läßt sie in Gefahr kommen ihres guten Namens und Lebens halben, läßt sie die Theurung und sonst allerlei Unglück treffen; aber er bereitet ihnen doch eine feine, friedsame und sichere Herberge und gibt ihnen Frieden auch mitten unter den Feinden.

173. Dieses alles wird geschrieben, unsern Glauben damit zu stärken, wenn schon allerlei Gefahr und Unglück vorfällt, wie Isaak in der größten Gefahr ist: daß dennoch GOTT der sei, der uns regieren und erhalten wolle. Denn es muß die Kirche hier auf Erden eine Stätte und Nest haben, da sie sich aufhalten möge. Wenn dieser Fürst oder auch ein anderer uns nicht schützen will, so wird doch GOTT einen andern geben, der uns freundlich und gern beherbergen wird; und er pflegt gemeiniglich ohne unsern Rath oder Gedanken, ja, ohne und wider alle unsere Hoffnung und Erwartung einen Wirth zu erwehlen, dabei die Seinen ihre Herberge mögen haben.

174. Auf solche Weise hat Isaak in demselbigen Lande Sicherheit, Frieden und Ruhe erlangt, und ist hernach sehr reich geworden, daß ihm auch die Einwohner im Lande solch groß Glück mißgönnt und ihn darüber beneidet haben, und daß er wiederum von dannen an einen andern Ort hat ziehen müssen. Dieses wird darum gelehrt, daß wir an GOTT glauben sollen, der uns regiert und erhält. Denn wenn wir das Wort haben, so werden wir zwar versucht mit mancherlei Gefahr, aber doch also, daß wir in der Versuchung nicht verderben oder verzweifeln, sondern daß wir uns in Frieden freuen und GOTT danken sollen.

\section[Vierter Theil]{Vierter Theil. Wie Gott Isaak in Gerar segnet; wie Isaak deswegen von den Philistern beneidet wird, Gerar verläßt, drei Brunnen gräbt und gen Bersaba zieht}

\subsection{I}

\begin{quote}
\textbf{V. 12—14. Und Isaak säete in dem Lande, und kriegte desselben Jahrs hundertfältig; denn der HErr segnete ihn. Und er ward ein großer Mann, ging und nahm zu, bis er fast groß ward, daß er viel Guts hatte an kleinem und großem Vieh, und ein groß Gesinde.}
\end{quote}

175. Dieses alles, was Mose von Isaak da erzählt, läßt sich ansehen, als habe er damit übertreibend (hyperbolisch) geredet und könnte nicht wohl möglich sein, dieweil solches gesagt wird von einem armen elenden Manne, und der ein Fremdling war, welcher auch nicht einen Fuß breit Eigenes gehabt in dem Lande, weder Acker noch Scheune, auch nicht eine arme Schäferhütte.

176. Wir haben aber oftmals gesagt, daß diese so heiligen Väter also beschrieben werden, daß wir an ihnen ein Exempel haben sollen eines christlichen und gottseligen Lebens, wie wir dasselbe sein nüchtern, züchtig, gerecht und gottseliglich führen sollen um des zukünftigen Lebens willen. Denn also sehen wir, wie dieser so heilige Patriarch Isaak so mancherlei Weise versucht, und dennoch nur sein einfältig im Worte GOttes gelebt hat, welches ihm in den Verheißungen reichlich gegeben war, nämlich, daß GOtt zu ihm gesagt hat: „Ich will mit dir sein, ich will dir alle diese Länder geben, und durch deinen Samen sollen alle Völker auf Erden gesegnet werden“; dieselbigen Worte hat er wohl bedacht und in seinem Herzen fest gehalten. Wir hören es schläfrig und mit tauben Ohren; wir achten des Wortes so hoch und
\quellseite{61}
groß nicht, wie die Väter gethan haben. Darum hat uns GOtt anzeigen wollen, wie diese so heiligen Männer gelebt haben, „nicht allein vom Brod, sondern von einem jeglichen Wort, das durch den Mund GOttes gehet“, Matth. 4, 4. Denn es ist unmöglich gewesen, daß Fleisch und Blut sollte solche große Mühseligkeit und Beschwerung dieses Lebens haben dulden und ertragen können, so allhier beschrieben werden.

177. Wer nun dieses Ding nicht auch gegenwärtig versucht hat und dieses lieset oder nach dem Fleisch beurtheilt, demselbigen dünkt es alles schlecht, gering und zeitlich Ding zu sein, als, von der Nahrung, vom Weibe, Gesinde, von Weide und Vieh. Was ist das? Da ist ja nichts Herrliches oder Köstliches, das einen sonderlichen Schein hätte. Aber wer es recht erwägt und bedenkt, wie groß diese Dinge sind, dem werden sie in Wahrheit ganz schwer, wichtig und groß zu sein dünken. Denn nach der so reichen Verheißung ist die Versicherung erfolgt, und daß er und sein Weib, sie ihrer Zucht und er seines Lebens halben, in Gefahr gekommen sind. Unterdeß aber hält er sich aufrecht mit dem Wort der Verheißung, daß ihm droben (V. 3. 4.) ist gegeben worden, und ist ein Fremdling in dem Lande, so ihm von GOtt verheißen und sein eigen war, gleichwie ein Exulant und Fremdling, und leidet mit Geduld, daß Fremde darin herrschen und regieren. Nun betrachte aber die Umstände, und setze den Fall, daß es jetzt mit dir auch dermaßen zuginge. Denn was wolltest du wohl thun, wenn du in deinem Hause einen fremden Herrn leiden solltest und wenn du aus allen deinen liegenden und andern Gütern verstoßen würdest? Du würdest wahrlich solches nimmermehr leiden.

178. Nun muß aber Isaak elend und fremde sein in dem Lande, so ihm von GOtt gegeben war, und also, daß er sich hat besorgen müssen, daß sein Weib in Gefahr kommen möchte ihrer Zucht und Keuschheit halben, das ist, daß ihm sein ganz Haus verwirret werden und er selbst darüber auch umkommen möchte. Denn wo er wäre erwürget worden, so wäre sein Gesind in große Gefahr und Noth gekommen. Wir sind zumal zarte Märtyrer, wir sitzen müßig und haben guten Frieden und Ruhe, gebrauchen unserer Güter, essen und trinken in Frieden: das ist aber zumal eine große Anfechtung, nämlich, wo Einer mit seinem Weibe und zwei Söhnen, dazu mit so einem großen Gesinde, mit so viel kleinem und großem Viehe in fremden Landen wohnen soll, da er das Futter und alle Nothdurft dieses Lebens um Geld kaufen muß. Es ist kein größerer Jammer, denn also in der ganzen Welt umher ziehen und nirgends eine gewisse Wohnung oder eigen Haus haben; wie das die Sprüche bei allen Völkern bezeugen, und sonderlich auch dieses Sprüchwort, so die Sachsen pflegen zu gebrauchen: Eigen Wat (Gewand), gut is dat. Wo man aber der Gnade oder Gunst der Freunde und Feinde, die fast ungewisß ist, leben soll, das ist eine sehr schwere Anfechtung.

179. Also ist aber das Leben der heiligen Väter gewesen, nämlich voller Anfechtungen, die ihnen auch äußerlich und in leiblichen Dingen widerfahren sind, auf daß ihr Glaube an die göttliche Verheißung damit geübet und geprüft werden möchte. Denn Isaak bleibt an himmlischen Dingen hangen, so er doch noch hier auf Erden ist; aber ohne Glauben hätte er ein solch hartes Leben nicht führen können. Mich nimmt Wunder, woher er doch Gesinde und Knechte möge bekommen haben; es sei denn, daß er dieselben aus Arabien bekommen habe, welches Land viel Knechte gehabt hat. Es wird aber ohne Zweifel viel Zankes und Haderes, viel Mühseligkeit in der Haushaltung, desgleichen viel Dieberei mit umgelaufen sein, und daß immer die Hirten und das Gesinde einander heimlich werden nachgestellt haben. Isaak aber hat damit Geduld und wartet immer auf Hülfe; denn er hat einen GOtt, der ihm Hülfe verheißen hat, und der zwar unsichtbar ist und eine Zeitlang die Seinen pfleget zu verlassen, aber der doch zu rechter Zeit gewißlich hilft.

180. Auf diese Anfechtung aber folgt nun der Trost, auf daß es nicht dafür gehalten werde, als hätte Isaak vergeblich und umsonst auf den HErrn geharrt und gehofft; hernach wird er wiederum versucht werden, wird aber auch abermals Trost empfangen. Denn diese Dinge sind stets also untereinander gemengt, daß eins immer auf das andere folgt, nämlich, Anfechtung und Trost, wie immer Nacht und Tag auf einander folgen; denn also wird das Wort der Verheißung und der Glaube geübt und geprüft. Derhalben so säet nun Isaak in dem fremden Lande, welches doch sein eigen ist, und bekommt desselbigen Jahres hundertfältig.

181. Es leuchtet aber an dem Könige Abimelech eine sonderliche Frömmigkeit; denn er nimmt nicht allein Isaak in seinen Schutz und Schirm, und läßt bei Leibesstrafe verbieten, daß ihm oder seinem Weibe jemand soll Gewalt thun: sondern er verleiht ihm auch ein Theil im Lande, Aecker und Wiesen für sein Vieh und Gesinde. Und wiewohl er ihm das um eine gewisse Summe Geldes geliehen hat, so ist doch solches auch ehrlich und wohl gethan; denn das wäre Isaak vielleicht an andern Orten nicht gestattet worden. Derhalben rühmt Mose diesen König billig, daß er fromm und gottselig gewesen sei, und ich zweifle nicht, er werde selig geworden sein wie sein Vater; wofern es nicht eben derselbe Abimelech ist, der zur Zeit Abrahams gelebet hat, daß ich gleichwohl ihn nicht halte; denn Isaak ist jetzt schier hundert Jahre alt.

182. Dies dient aber dazu, daß wir daraus GOttes Güte und Gnade erkennen lernen, der die Seinen pfleget zu trösten: wenn er sie betrübt und ihren Glauben und Beständigkeit geprüft hat, so läßt er wiederum über ihnen aufgehen und scheinen das Licht und die Sonne seiner Gnade und Barmherzigkeit; er gibt ihnen eine gute Herberge, daß sie nicht allein im Frieden leben mögen, sondern daß sie auch groß Gut bekommen.

183. Die Juden sagen, es sei dies Jahr das ganze Land sehr unfruchtbar gewesen, und habe nur Isaaks Acker, den er um den Zins gedingt hatte, so reiche Frucht gebracht. Es bedarf aber des jüdischen Zusatzes nicht, daß man dieses damit groß mache, es ist sonst an sich selber groß genug, daß Ein Scheffel hundert Scheffel wiedergegeben hat, und wird ohne Zweifel ein groß Wunder und in dem Lande ungewöhnlich gewesen sein, daß der Acker so reiche Frucht gebracht hat. Hundert Scheffel aber gehen tausend, und so immerfort, ohne Ende und Zahl. Ich will der andern Güter schweigen, so vom Ackerbau und Viehzucht kommen, als, Butter, Milch, Eier und dergleichen. In unsern Landen ist es ein unglaublich und unbegreiflich Ding. Das hält man für eine reiche Frucht, wenn Ein Scheffel fünf oder sechs Scheffel bringt; das ist aber noch keine dreißigfältige oder sechzigfältige Frucht, davon im Evangelium gemeldet wird im Gleichniß vom Säemann, Luc. 8, 8.

184. Daraus lasset uns nun lernen, daß wir im Glauben beharren und beständig bleiben, daß wir nicht zweifeln oder wanken, wenn irgend eine Anfechtung vorfällt. Denn siehe, wie mildiglich GOtt dem Isaak die vorige Anfechtung vergolten und erstattet hat, auf daß wir wissen, daß er nicht ewig zürnen wolle, und wenn wir seine Verheißung haben, sollen wir in der Anfechtung beharren und beständig bleiben, und also schließen und sagen: Der Herr, der zu mir gesagt hat: Glaube du mir, der wird meine Verheißung gewißlich halten; ich will unterdeß sein warten, und auf ihn harren. Wie im 27. Psalm V. 14. steht: „Harre des HErrn, sei getrost und unverzagt“; item im 31. Psalm V. 25.: „Seyd getrost und unverzagt, alle, die ihr des HErrn harret.“ Die Gottseligen sollen harren und beständig bleiben auch in der allergrößten Gefahr und Noth: denn er hat verheißen, daß er für uns sorgen wolle; wie St. Petrus sagt 1. Epist. 5, 7.: „Alle eure Sorge werfet auf ihn; denn er sorget für euch“ etc.

185. Diese Uebung und Prüfung des Glaubens thut den Gottseligen noth; denn wo sie ohne dieselbige wären, so würde der Glaube matt und kalt werden, ja, er würde endlich gar verlöschen. Wenn sie aber durch Anfechtung also geprüft werden, so lernen sie, was der Glaube sei; und wenn sie versucht werden, so wachsen sie in der Erkenntniß des Sohnes GOttes, und werden so stark und kräftig, daß sie in der Noth und Widerwärtigkeit sich nicht weniger freuen und rühmen können, als im Glück und wenn es ihnen wohl geht; ja, daß sie eine jede Anfechtung nur für eine Wolke oder Nebel halten, welche bald verschwinden und vorübergehen werde.

186. Wie viel Scheffel aber Isaak gesäet habe, erzählt Mose nicht; er sagt nur, daß Ein Scheffel hundert gebracht habe; wenn er hundert gesäet hat, so hat er zehenmal tausend Scheffel wiederbekommen. Das ist ein sonderlicher Segen, nach der Verheißung, so ihm droben (V. 3.) geschehen ist, und Mose beharrt noch immer dabei, daß er dieses Wunderwerk groß macht und sagt, daß Isaak sehr viel Gut bekommen habe. Und wo diese reiche Frucht und Einkommen zwei oder drei Jahre gewähret hätte, so wäre er über alle Maßen groß und reich davon geworden. Er hätte sich zwar daran wohl genügen lassen, wenn Ein Scheffel hätte zehn oder noch weniger gebracht; aber GOtt sättigt ihn nicht allein, sondern überschüttet ihn mit reichem Segen, Psalm 21, 4., und gibt ihm viel mehr, als ihm vonnöthen ist für sein Haus und Gesinde, wie viel er dessen auch gehabt hat; denn er hat kaum den zwanzigsten Theil bedurft. Derhalben kann GOtt die Seinen leidlich reich machen, wenn er den Ackermann für tausend Scheffel, daran er zuvor reichlich genug gehabt, zehnmal tausend Scheffel gibt; und Isaac hat doch dermaßen noch nichts Eigenes im Lande, er lebt in seinem eigenen Lande, gleich als wäre es nicht sein eigen.
\quellseite{62}

187. Darum ist dies ein großer Trost, wenn wir ihn nur recht verstehen und annehmen. Denn GOtt gibt dieses alles, auf daß wir lernen glauben und in der Noth und Widerwärtigkeit auf die Hülfe des Herrn warten, im Glück aber und wenn es uns wohl geht, GOtt für seine Wohlthat danken; denn darin besteht das Leben der Heiligen. Der 104. Psalm V. 8. sagt: „Die Berge gehen hoch hervor, und die Breiten setzen sich herunter“ etc. Jetzt ist es Tag, bald wird es Nacht. Heute haben wir Ehre, morgen Schande, 2 Cor. 6, 8. Allezeit aber sollen wir unsere Augen und Herz richten auf die Verheißung und uns daran halten, die wird gewißlich kommen und nicht verziehen. Nun folgt wieder eine andere Anfechtung, damit der Glaube Isaacs aufs neue angefochten und geprüft wird.

\subsection{II}

\begin{quote}
\textbf{V. 14--16. Darum neideten ihn die Philister, und verstopften alle Brunnen, die seines Vaters Knechte gegraben hatten, zur Zeit Abrahams, seines Vaters, und füllten sie mit Erde; daß auch Abimelech zu ihm sprach: Weich von uns, denn du bist uns zu mächtig worden.}
\end{quote}

188. Nach dem göttlichen Trost wird der so heilige Patriarch wiederum in mancherlei und schwere Noth und Trübsal geworfen. Denn also müssen gottselige Menschen ihr Leben anrichten, daß sie zwar dieser Welt, des Ehestandes, Essens, Trinkens und aller andern Nothdurft dieses Lebens gebrauchen, aber doch gleichwohl mit dem Herzen an den himmlischen Gütern hangen. Die das nicht thun, bleiben Bauchdiener und des Mammons Knechte. Der Teufel aber läßt uns in diesem Leben keine Ruhe, es ist unmöglich, daß wir mit ihm können Frieden haben, sintemal er GOttes abgesagter Feind ist und weil auch Feindschaft gesetzt ist zwischen den Kindern der Schlange und dem Samen des Weibes.

189. Die erste Anfechtung aber ist diese, daß die Philister aus Haß und Neid alle Brunnen verstopfen und mit Erde füllen; dadurch wollen sie ihn bewegen und verursachen, daß er von selber davon ziehen soll. Isaac aber hat dies alles mit Geduld ertragen, und gebraucht unterdeß Wasser, so er von Andern um sein Geld gekauft hat.

190. Weil sie ihn auf diese Weise nicht vertreiben können, fahren sie zu und überreden den König, daß er ihm durch ein sonderlich Gebot sagen und gebieten wolle, aus dem Lande hinweg zu ziehen; da muß er weichen und davon ziehen, und nicht ein fremdes, sondern sein eigenes Land verlassen. Derhalben ist diese Anfechtung härter und betrübter, als sie äußerlich scheint, wenn sie Einer nur obenhin ansieht. Origenes bringt an diesem Ort, ich weiß nicht was für Weisheit vor, von den verstopften Brunnen, welcher Weisheit sich Erasmus hoch verwundert; unterdeß aber verliert man die Lehre, wie man ein christlich Leben anrichten soll, und wird die Historie nicht gehandelt, darin wir seine christliche Lehren haben.

191. Isaac ist reich und hat viel Getreide, er lehrt sein Gesinde Geduld tragen und sich der Gottseligkeit befleißigen: was widerfährt aber diesem heiligen Manne und Priester GOttes? Da es sich nun ansehen läßt, daß er alles überflüssig gesegnet hatte, was ihm nur vonnöthen gewesen, da werden die Philister aus Haß und Neid bewogen und nehmen ihm das Wasser mit Gewalt, also daß sie auch selber desselben nicht haben genießen mögen. Wenn sie ihm das Wasser nur verboten hätten, so hätte es etwa heimlich nehmen und gebrauchen können; aber sie fahren zu und verstopfen die Brunnen gar und füllen sie mit Erde, daß sie alle zugleich des Wassers nicht gebrauchen können, ohne welches man gleichwohl dies Leben nicht kann erhalten, vornehmlich bei dem Vieh. Wo man Einen des Wassers beraubt, ist es eben so viel, als beraubte man ihn des Lebens; denn davon müssen beide Menschen und Vieh sterben.

192. Derhalben ist dies eine sehr böse und schwere Anfechtung gewesen bei so einer großen Menge Gesindes und Viehes, und ist eine außerordentlich große Boßheit, Einem das Wasser verbieten in solcher großen Noth. Was hat Isaak da Anderes können denken, denn daß er mit Vieh und Gesinde würde Durstes sterben müssen? Er hat nur diesen einigen Trost gehabt, daß er seine Diener gen Gerar hat senden können, und daher in Kannen und Flaschen Wasser holen lassen, das er um sein Geld gekauft hat, wie wir den Wein kaufen: und hat ihm gleichwohl das seinen Schmerz auch gemehret, daß er sein eigen Wasser hat kaufen müssen; denn er war eigentlich der rechte Herr des Wassers, des Brunnens und des ganzen Landes, und sein Vater Abraham hatte dieselbigen Brunnen gegraben.

193. Darum lasset uns nun hieraus lernen, daß unsere Anfechtungen viel geringer und leichter sind, als der heiligen Patriarchen Noth und Trübsal gewesen. Wir sind ja zumal zarte Märtyrer, die wir alle Dinge vollauf haben. Derhalben haben diese Historien gar kein gleich Exempel, und ist heutiges Tages keiner auch unter den allerheiligsten Menschen, der dies würde leiden und ertragen können. Diese guten Leute haben müssen hin und wieder in der Irre ziehen von einem Dorf zum andern bis gen Hebron, vielleicht auch wohl bis in die Wüste zu dem Brunnen des Lebendigen und Sehenden (denn Mose schreibt nicht, woher sie das Wasser bekommen haben), und von bannen für sich und für das Vieh Trank holen. Welchem Menschen sollte doch solche große Unbilligkeit nicht zu Herzen gehen? Es ist wohl glaublich, daß bisweilen Rebekka aus Ungeduld wird gesagt haben: Sind wir doch Herren dieses Landes, wo bleibt denn nun die göttliche Verheißung? Wird sie also erfüllt, daß Fremde unser Güter genießen und wir sollen vertrieben werden? oder aber daß uns sonst alles genommen wird, so wir zu unserer Nothdurft haben und genießen sollten?

194. Derhalben ist dies, so hier vom Wasser und vom Brunnen erzählt wird, nicht so lächerlich Ding, wie das Fleisch davon urtheilt; sondern es hat Alles dazu gedient, daß damit der Glaube hat sollen geprüft werden. Denn Isaak hat ja eine göttliche Verheißung gehabt, da GOtt zu ihm gesagt hat: Ich will mit dir sein; jetzt aber begegnet ihm das Gegentheil, gleich als ob er von GOtt, der ihm die Verheißung gethan, nun gar verlassen wäre.

195. Warum sollen wir auch lernen auf den Herrn harren, wenn wir seine Verheißung haben, es gerathe gleich mit dem Anderen, wie es wolle, es komme gleich eine Verfolgung über uns an deß Statt, da wir gedacht haben, daß GOtt bei uns sein würde laut seiner Verheißung, oder aber, daß wir anstatt der Gnade seinen Zorn fühlen und empfinden; daß wir dennoch sagen: Ich glaube, ich bin getauft, ich bin absolvirt, ich habe die göttliche Verheißung von seiner Gnade und Barmherzigkeit; ich habe daran genug; es mag mir nun Tag oder Nacht, Trübsal oder Fröhlichkeit vorfallen, so will ich doch seine Gnade und Barmherzigkeit nicht fahren lassen, will auch in meinem Herzen nicht verzagen.

196. Die andere Anfechtung ist diese. Nachdem sie ihm das Wasser genommen und ihn desselben beraubt haben, da kommen die Scharrhansen, und versuchen den König, daß er sein Herz auch von ihm wendet, und den Vertrag widerruft, so er mit Isaak gemacht, daß er ihm die Wiesen und Aecker um Zins vermiethet hatte. Und diese Scharrhansen maßen sich nun die Aecker und Wiesen alle an; sie verstehen nicht, daß es ein sonderlich Wunderwerk und Segen GOttes sei, sondern meinen, der Acker sei sonst von selber so fruchtbar, daß er so viel reiche Frucht bringe: darum wollen sie lieber selbst für sich denselbigen Acker gebrauchen, denn daß sie ihn einen Fremden gebrauchen lassen wollten, welcher etwa bald hernach, wenn er würde reich werden, mit Gewalt oder Geld alle Aecker und Wiesen der Einwohner und darnach allgemach das ganze Land unter sich in seine Gewalt würde bringen können. Aber dadurch haben sie den Fluch über sich selbst geführt. Denn es will sich nicht reimen, daß eben solche Fruchtbarkeit, welche Isaak zuvor durch göttlichen Segen gehabt hat, ihnen auch verliehen worden sei.

197. Vielleicht mögen die Philister auch wohl irgend aus den Worten Eines unter seinem Gesinde erfahren und verstanden haben, daß dieser Fremdling sich sollte rühmen, daß ihm die Herrschaft desselbigen Landes verheissen wäre. Das hat sie sehr unbillig zu sein gedünkt, und haben gedacht, sie möchten ihm darüber billig feind sein und ihn in das Elend treiben. Derhalben wird er erstlich durch Privatgewalt, nun aber durch die königliche Gewalt und durch ein öffentlich Edict vertrieben. Denn das war ihr Plan, als sie die Brunnen verstopft hatten, daß er derhalben, da es ihm an Wasser mangeln würde, von selber aus dem Lande ziehen sollte, wenn ihn auch niemand dazu zwingen würde.
\quellseite{63}

198. Aber die Anfechtung ist noch viel schwerer, daß er durch den Befehl des Königes vertrieben wird, von dem er bisher gemeint hat, daß er ihm günstig und wohlgelegen wäre, durch welches Schutz und Schirm er wider allerlei Gewalt sicher sein möchte, und muß nun jetzt von demselbigen hören, daß er zu ihm sagt: „Weich von uns, denn du bist uns zu mächtig“, zu stark und zu reich geworden, du hast mehr Knechte und mehr Vieh als wir. Also wird Isaak vertrieben, dieweil sie ihm den reichen Segen mißgönnen und ihn darum neiden, und können nicht leiden, daß er daselbst reich werde, verbieten ihm derhalben erstlich das Wasser und darnach das Land.

199. Dieses ist die andere Anfechtung, so stracks wider die Verheißung ist; aber Isaak bleibt doch gleichwohl an derselben hangen und hält sich daran fest, und denkt also: Gott hat mich vom Tode errettet, er hat mein Weib bei ihrer Zucht und Keuschheit erhalten, er hat mir desgleichen auch hundertfältige Früchte gegeben von dem, was ich gesäet hatte: darum will ich mich nicht beschweren, diese Wohnung zu verlassen, und will dennoch in der Nähe bleiben; denn der Herr hat mich heißen ein Pilger oder Fremdling sein in diesem Lande.

\subsection{III}

\begin{quote}
\textbf{V. 17. 18. Da zog Isaak von dannen, und schlug sein Gezelt auf im Grunde Gerar, und wohnte allda, und ließ die Wasserbrunnen wieder aufgraben, die sie zu Abrahams Zeiten, seines Vaters, gegraben hatten, welche die Philister verstopfet hatten nach Abrahams Tod, und nannte sie mit denselben Namen, da sie sein Vater mit genannt hatte.}
\end{quote}

200. Isaak weicht und zieht an einen andern Ort, der nicht weit von dem vorigen gelegen war, ohne Zweifel nicht ohne großen Schmerz und Anfechtung des Fleisches. Was er aber für einen Acker gehabt habe, zeigt die Schrift nicht an, der Text sagt nur: „Er schlug sein Gezelt auf im Grunde Gerar“, in demselben Winkel. Es haben aber an diesen drei Orten, nämlich zu Hebron, Bersaba und Gerar, drei Patriarchen gewohnt, Abraham, Isaak und Jakob: daraus denn erhellt, daß Gott an denselbigen Orten sonderliche Lust müsse gehabt haben. Und diese Orte sind gelegen im Stamm Juda gegen Mittag, sie sind vier oder fünf Meilen breit gewesen. Daselbst sind die heiligen Patriarchen umher gezogen, haben bisweilen zu Gerar, bisweilen auch wohl zu Bersaba gewohnt, und Mose nennt es einen Grund; denn der Acker daselbst ist nicht gebauet oder fruchtbar gewesen in dem Thale Gerar; denn in den Gründen bei den Wasserbächen pflegt das Land trocken und dürr zu sein.

201. Derhalben ist es eine außerordentlich große Anfechtung gewesen, in welcher Isaak gelernt hat, nicht allein vom Brod zu leben, sondern von einem jeglichen Wort, das durch den Mund Gottes gehet, Matth. 4, 4. Um des Ortes willen aber haben ihn die Philister nicht beneidet, da er keinen guten Acker gehabt, da er hätte säen und ernten mögen. Denn die Wasser pflegen die Aecker wüste zu machen mit ihrem Ueberlaufen. Womit hat er denn nun sich und sein Gesinde erhalten? Hat er Geld gehabt, so hat er Proviant zu seiner Nahrung kaufen mögen. Darnach graben sie andere Wasserbrunnen, welche hernach die Philister abermals mit Erde füllen und verstopfen werden, wie sie zuvor auch gethan haben, und Isaak wird doch gleichwohl solche Feindschaft und Neid überwinden, und unüberwindlich verharren und beständig bleiben.

202. Wiewohl da immer eine Anfechtung über die andere auf den Trost, so er gehabt, gefolgt ist; wie Mose mit einem nachdrücklichen und emphatischen Worte beschreibt, wie er auf dieser Reise so einen starken und beständigen Muth gehabt habe: Isaak zieht von dannen, wird vertrieben und verjagt aus seinem Lande, und schlägt dennoch ein Gezelt auf. Das ist auch ein kriegerisch Wort, das auch im 34. Psalm V. 8. gebraucht wird:

„Der Engel des Herrn lagert sich“ etc. Das hebräische Wort ohel heißt auf Lateinisch tabernaculum, Hütte; chanah aber heißt, castra figere, Lager aufschlagen. Darnach wird es auch gebraucht, daß es eine Wohnung heißt; aber eigentlich versteht man es von einer solchen Hütte oder Gezelt, wie man im Kriege hat.

203. Darum hat Mose anzeigen wollen, daß der Patriarch Isaak ein unverzagtes Herz gehabt habe in so vieler großer Trübsal, und daß er dessen gar gewiß ist, an welchem Orte er nun wohnte, daselbst wäre ein Schloß oder feste Burg, darin er von Gott würde bewahret werden. So kräftig und stark sind die göttlichen Verheißungen, daß es sich mit Isaak ansehen läßt, als habe er des Teufels und seiner Glieder gespottet, die ihn aus seinem sehr schönen Sitze vertrieben haben. Denn er hat also geschlossen: Wenn ich an diesem Orte nicht kann hundertfältige Frucht haben, sondern muß verlassen, was ich hier mit gutem Rechte inne habe, so werde ich dennoch gleichwohl mein Auskommen und Nothdurft haben, ja, ich werde sicher leben und in meinem Gezelt wohnen; wenn ich nicht bei den fruchtbaren Aeckern zu Gerar wohnen kann, so will ich im Grunde wohnen, welcher mir ein Gezelt oder Lager sein wird, nicht mit Holz oder Steinen bewahret oder befestiget, sondern mit der Heerschaar der heiligen Engel.

204. Also beschreibt Mose Isaak, wie er fast erschrocken und fast in höchster Trübsal und Angst gewesen ist, und beschreibt ihn doch gleichwohl dermaßen, daß er noch immer einen starken Muth gehabt wie ein Held, und daß er die Welt verachtet und fest an der Hoffnung zu Gott gehalten, dadurch er also gesinnet gewesen, daß er sagt: Ich werde nicht allein meine Nahrung zur Nothdurft haben, sondern auch dabei Schutz und sichere Wohnung.

205. Dieses wird uns zum Exempel und Trost vorgehalten. Denn wo wir eine göttliche Verheißung haben, so muß das gewißlich folgen und kann nicht fehlen, daß auch die heiligen Engel bei und um uns sind. Und daher kommt der Vers im 34. Psalm V. 8.: „Der Engel des Herrn lagert sich um die her, so ihn fürchten“ etc., und im 91. Psalm V. 11.: „Er hat seinen Engeln befohlen über dir“ etc. Das haben die Väter fest und ohne allen Zweifel geglaubt. Derhalben, wo wir göttliche Leute sind, sollen wir auch an die Verheißung des Glauben, der nicht lügen kann, alsdann sind wir gewißlich unter seinem Schutz und Schirm; dazu ist es auch gewiß, daß die heiligen Engel bei uns sind.

206. Wo uns aber über oder wider solch Vertrauen und Schutz irgend ein Unfall begegnet, so geschieht dasselbe aus sonderlichem Rath Gottes, der uns verborgen ist und davon vornehmlich die Widersacher nichts wissen. Also sollen sich aber die Gottseligen trösten und sagen: Ich weiß, daß ich die heiligen Engel zu Hütern habe; daß ich aber noch gleichwohl Widerwärtigkeit leiden muß, dasselbe befehle ich dem Willen Gottes; denn ich bin in dem Heerlager der heiligen Engel: Gott ist kein Lügner, darum wird er mich nicht verlassen.

207. Darum will uns Mose mit diesem Exempel ermahnt haben, daß wir hochfärtig sein sollen der göttlichen Verheißung halben, die wir auch reichlich im Evangelium haben. Denn sintemal ich weiß, daß ich durch das Predigtamt von meinen Sünden absolvirt, auch getauft bin, da soll ich auch einen starken unverzagten Muth haben in allerlei Noth und Gefahr, es müsse gleich mein Leib oder Seele die äußerste Noth leiden.

208. Ja, sagst du, ich kann aber das nicht glauben, kann so stark nicht sein. Antwort: Das gebe ich zu; denn das Fleisch widerstrebet der Verheißung; wie St. Paulus von sich selber auch sagt Röm. 7, 23.: „Ich sehe ein ander Gesetz in meinen Gliedern, das da widerstreitet dem Gesetz in meinem Gemüthe.“ Da redet er von der allgemeinen Sünde und vom Verderben der ganzen Natur, davon es herkommt, daß wir nicht so stark und beständiglich dem Gesetz in unsern Gliedern widerstreben und dem Gesetz in unserm Gemüth anhangen. Und eben darüber klagen alle Gottseligen, dieweil sie noch in diesem Fleische leben. Wir haben zwar diese Weisheit und Verstand, daß wir etlichermaßen und schüchtern die göttlichen Verheißungen ergreifen, und wir seufzen täglich, daß unser Fleisch, so den Glauben hindert, zu Asche werden möge; man muß sich aber befleißigen, daß wir in derselben starken Erkenntniß wachsen und zunehmen mögen, bis wir endlich vollkommen werden.

209. Aber wenn wir ohne Anfechtung sind, so lernen wir nichts, nehmen auch gar nicht zu. Denn das ist die Ritterschaft und Uebung der Christen, dadurch wir erkennen, daß wir unter dem Schutz der heiligen Engel sind, und wiewohl wir mit harten und schweren Anfechtungen geübt und versucht werden, daß sie uns doch gleichwohl nicht schaden. Das ist unsere Theologie, die man nicht so leichtlich und bald lernen kann, sondern man muß an das Gesetz des Herrn denken, davon Tag und Nacht reden, man muß immer zu Felde liegen und wider den Teufel streiten, der versucht, uns vom Worte abzuziehen und wie er unsern Glauben schwächen möge. Da Isaak auf solche Weise mit Gewalt von den Philistern geplagt worden ist, richtet er sich doch wiederum auf, und schließt bei sich selbst also: Ich werde in diesem Lande wohnen, wenn es auch denen allen leid sein sollte, die mich da heraus wollen verstoßen haben; denn das Land ist mein nach der Verheißung: das wird mir ein königlicher Sitz sein im Grunde bei diesem Bach. Ich halte aber dafür, daß es ein solches Thal gewesen sei wie in Sachsen zwischen Düben und Eilenburg bei dem Wasser, die Mulde genannt. Wiewohl er daselbst auch noch keinen Sitz oder Wohnung haben wird, so hat er dennoch also gesagt im Glauben an die Verheißung und hat geglaubt, daß er im Heerlager der Engel wohnte. Denn Gott lügt nicht, darum werden wir, hat er gedacht, nicht verderben, intemal er mir dies Land verheißen hat. Derhalben wenn wir auch schon verderben müßten, so werden wir doch nicht verderben. Dies hat aber Isaak von seinem Vater Abraham gelernt, dem eben solches auch begegnet ist und der auch mit gleich starkem Glauben Alles überwunden hat.
\quellseite{64}

210. Hier wird aber der Wasserbrunnen gedacht, so die Knechte seines Vaters Abraham gegraben und welche die Philister nach Abrahams Tode verstopft hatten. Dieses gehört noch in das 21. Cap. 3. 25., dahin soll man es ziehen; denn hiermit wird ausgelegt und erklärt, was daselbst nur von Einem Brunnen gesagt wird, welchen die Knechte des Königs Abimelech mit Gewalt genommen haben, da sich Abimelech vor Abraham entschuldigt, daß er darum nichts gewußt habe, und darnach einen Bund mit ihm macht. Da wird vielleicht Isaak eingefallen sein, daß er gedacht hat, daß er gleich auf dieselbige Weise von den Einwohnern des Landes angeklagt würde wie zuvor Abraham. Wiewohl das härter und beschwerlicher ist, daß ihn der König selbst heißt aus dem Lande ziehen. Es thun ihm nicht allein die Knechte Gewalt, sondern der König selbst, da er sich durch die falsche Anklage der Hoffschranzen hat verführen lassen, vertreibt ihn aus dem Lande. Denn wiewohl wir es dafür halten, daß er fromm gewesen sei, so hat er doch dem Rath der Fürsten in seinem Königreich folgen müssen.

\subsection{IV}

\begin{quote}
\textbf{V. 19--21. Auch gruben Isaaks Knechte im Grunde, und fanden daselbst einen Brunnen lebendigen Wassers. Aber die Hirten von Gerar zankten mit den Hirten Isaaks, und sprachen: Das Wasser ist unser. Da hieß er den Brunnen Esek, darum, daß sie ihm da Unrecht gethan hatten. Da gruben sie einen andern Brunnen, da zankten sie auch über; darum hieß er ihn Sitna.}
\end{quote}

211. Dies ist nun ein neuer Segen, damit Gott Isaak aus sonderlicher Gnade überschüttet hat nach der vorigen Anfechtung, so er gehabt; denn er gibt ihm nicht allein solche Wasserbrunnen, die wiederum aufgeräumt waren, sondern auch einen neuen Brunnen lebendigen Wassers, den man zuvor in dem Thale nicht gesehen hatte. Denselbigen Brunnen aber hat Isaak sehr lieb, und lobt ihn und nennt ihn „einen Brunnen lebendigen Wassers“. Denn es ist ein sehr schöner Trost, wo man einen Brunnen hat, der da für und für fließt; und er wird vielleicht aus Bergen geflossen sein.

212. Augustinus macht einen Unterschied zwischen den beiden lateinischen Wörtern puteus und fons (wir heißen es beides einen Brunnen), und sagt, puteus heiße ein solcher Brunnen, den wir mit Eimern pflegen zu schöpfen, darin das Wasser gleichsam todt ist und nicht fließt; fontes aber, sagt er, sind die Brunnen, so da fließen, und mit dem Fließen anzeigen, daß sie leben, und sind den Thieren gleich, so da leben.

213. Also hat nun Gott dies lebendige Wasser Isaak zum Troste gegeben; es ist ihm aber eben damit gegangen, wie man im Sprüchwort sagt: Etliche sind, die da arbeiten, die Andern aber genießen der Arbeit: Etliche säen, die Andern aber ernten. Denn da die Philister sehen, daß Isaak den Brunnen des lebendigen Wassers und die schöne lustige Brunnenader gefunden hat, da verstehen sie bald, wie großen Nutzen man davon haben könne; denn er nicht allein den Bürgern Wasser gibt und Wasser im Brunnen drinnen bleibt, sondern es auch heraus fließt und das ganze Land feuchtet. Darum gefällt er den Philistern so wohl, wiewohl sie ihn nicht gefunden noch gegraben haben; ja, sie haben ihn in ihrem Lande nicht einmal gesehen, und dennoch wollen sie denselbigen sich zueignen.

214. Mose schweigt von dem Könige und den Fürsten fein stille und sagt nur von den Dorfhirten, gleich als hätten dieselbigen den Zank über den Brunnen erregt. Es ist aber ein schwerer Zank gewesen, weil sie sich derselben mit Recht anmaßen, und ziehen den Text an, daß sie sagen: Das Land ist unser; derhalben alles, was es gemacht hat und das Land in sich begreift, steht uns zu.

215. Was soll nun Isaak thun? Er ist ein guter, frommer und aufrichtiger Mann, ist nicht zänkisch; derhalben trägt und leidet er diese Gewalt auch mit Geduld. Wiewohl es ein schwer Leiden ist, daß dieser Leute Bosheit so groß, daß sie ihm auch den Trank mißgönnen und mit Gewalt nehmen, den er von dem Wasser gehabt hat, so ihm von Gott verliehen war. Denn es war ihnen nicht gegeben, so hatten sie es auch nicht gefunden. Sondern es war durch sonderlichen Segen diesem heiligen Manne gegeben. Die Juristen hätten an diesem Ort Ursache gehabt, ganzer Hundert Jahre zu zanken; Isaak aber weicht der Zeit und Person, und läßt die Sache auch fahren und verläßt den Brunnen. Ich würde zwar gebeten haben, daß Gott die Ader und den Brunnen, der ohne ihre Arbeit erfunden war, mit einander vergeben lassen sollte.

216. Isaak aber rächt diese Gewalt anders nicht, denn daß er dem Brunnen einen Namen gibt und heißt ihn Esek zum ewigen Gedächtnis der großen Bosheit; auf daß er damit anzeige, daß das Unrecht von ihnen hergekommen wäre, bei ihm aber eitel unschuldige Geduld gespüret worden sei, welches eine sehr heilige Tugend ist, und hält sich aufrecht mit einem guten Gewissen, das sich unschuldig weiß, und nichts Unrechtes gethan hat, verletzet oder beschädigt den Nächsten nicht, viel weniger Gott in seinen Heiligen. Darum auch Isaak die Philister nicht schmäht; es ist genug, daß er den schändlichen Namen Esek hinter sich läßt zum Zeugniß der großen Bosheit dieses Volkes.

217. Denn auf solchen ungerechten Raub pfleget gemeiniglich von selber die Rache und Strafe zu folgen; wie das der Spruch Augustinus bezeuget, da er sagt: Lucrum in arca, damnum in conscientia: Gewinn in der Kiste, Schade im Gewissen. Kein ungerechter Gewinn geht ohne sehr ungerechten Schaden ab. Denn wer einem Andern sein Gut nimmt, hat schon den Glauben verloren; mit dem Gute bekommt man den Teufel und verliert Gott mit dem Glauben und mit der Gerechtigkeit. Wenn du mir das Gut nimmst und mir den Glauben und ein gut Gewissen bleiben lässest, da hast du die Spreu, ich aber habe das gute Korn. Was du nun mit Recht oder Unrecht rauben wirst, so hast du damit schon deinen Glauben verloren und hast Gott zum Feinde, und wirst darnach mit der Spreu in die ewige Verdammnis gestürzt werden.

218. Dies wird zwar oft gesagt und den Leuten eingeschärft, aber man hört und glaubt es nicht. Es ist bei allen Völkern ein gemeiner Spruch: Male parta male dilabuntur; item: Male partatum male disperit; item: De male quaesitis non gaudet tertius haeres. Diese Sprüchwörter lauten bei den Deutschen also: Wie gewonnen, so zerronnen; oder: Uebel gewonnen, häßlich zerronnen; item: Unrecht Gut währt nicht (das heißt, gedeihet nicht) und kommt nicht auf den dritten Erben. Und wir hören doch nicht auf zu rauben und zu stehlen mit falschem Maß, unrechtem Gewicht und sonst mit allerlei List und Trügerei.

219. Ich selbst, nachdem ich mein zwanzigstes Jahr vollendet hatte (denn vor der Zeit steht der Verstand der Mensch in gemeinen Händeln nichts), habe darnach innerhalb vierzig Jahren viel Stuhlräuber, Geizhälse und Wucherer gesehen, welche, da sie großes Geld und Gut zusammengescharrt, dasselbe dennoch kaum dreißig Jahre haben behalten können.

220. Eben auf diese Weise halte ich, daß dieser König und seine großen Scharhansen wenig Nutzen oder Lust von diesem ungerechten Gewinne werden gehabt haben. Derhalben soll man den Spruch im 37. Psalm V. 16. wohl behalten: „Das Wenige, das ein Gerechter hat, ist besser, denn das große Gut vieler Gottlosen.“ Denn die Ungerechtigkeit läßt das Gewissen nicht fröhlich oder friedsam sein. Wer groß Gut sammelt, der kann sich dessen nicht freuen, und wer solch unrecht Gut erbt, wird arm und verderbt darüber. Denn im Hause des Gottlosen zerrinnt es alles. Man scharrt daselbst wohl großen Reichthum zusammen, aber er muß doch zerrinnen und vergehen; denn GOtt bläset drein, und der es sammelt, der sammelt es in einen löchrichten Beutel, wie die Poeten eine Fabel haben von dem Siebe der Danaiden.
\quellseite{65}

221. Darum, wer gottseelig und fromm ist, soll sich an dem Wenigen genügen lassen nach dem Spruche im 37. Psalm V. 16., und gebrauche desselben mit seinem Weibe und Kindern und enthalte sich vom unrechten Gewinne. Dasselbe lehren auch die Historien und Bücher aller Völker, nämlich, daß es viel besser sei, daß du von einem oder zweien Gulden, welche du mit Ehren erworben hast, lebest, denn sonst von vielen tausend Gulden, so mit anderer Leute Schaden schändlich erworben sind.

222. Und ich habe eine Historie gehört, die sehr merkwürdig ist, von einem frommen und ehrlichen Bürger, welcher, da er seine Tochter verlobet und derselbigen nun ihre Aussteuer geben wollte, zu seinem Schwiegersohne sagte: Siehe, da hast du dreißig Gulden, die ich mit GOtt und mit Ehren erworben habe ohne irgend eines Menschen Nachtheil und Schaden. Diese Aussteuer wollte ich genommen haben für dreitausend Gulden. Denn GOtt segnet das wohlerworbene Gut also, daß ein Heller mehr bringt, denn ein ganzer Gulden eines Räubers oder Wucherers.

223. Derhalben ist dies eine schwere Sünde der Leute zu Gerar, damit sie den Fluch GOttes auf sich laden, darum daß sie einem solchen frommen Gast Gewalt und Unrecht thun. Er rächt sich solcher Gewalt halben an ihnen nicht, flucht ihnen auch nicht, ja, er wollte sie wohl viel lieber segnen. Und er möchte ihnen zwar schier für den Brunnen danken, welchen sie ihm genommen haben; aber diese Danksagung ist fast unselig und ist kein Glück dabei; es ist ein unglückselig Balet, wenn solche Menschen, die voll Heiligen Geistes sind, hinweg ziehen, und sagen zum Lande, daraus sie ziehen, Essig und Sitna. Da folgt in Kürze der Fluch. Denn das will Isaak sagen: Hier ist nichts als Esel und Sitna. Er wollte gerne etwas Anderes sagen, aber er kann die Gewalt, so die zu Gerar an ihm geübt hatten, nicht billigen oder rechtfertigen. Derhalben nimmt er Abschied von ihnen, indem er ihnen die Sünde und Gewalt vorwirft, welche er bei ihnen erlitten hatte. Er wünscht ihnen nichts Böses; er weissagt aber Böses, so über sie kommen werde, nämlich, daß daselbst die Brunnen sind Esel und Sitna.

224. Dieses hat der Heilige Geist uns zur Lehre schreiben wollen, auf daß wir wissen, daß man sich der Gewalt halben nicht selbst rächen soll. Klagen mag man darüber wohl, aber rächen mag man sich nicht. Darum rächt sich auch Isaak nicht; er zeigt nur seinen Schmerz an, und läßt zum Zeugniß der Gewalt die Namen hinter sich, nämlich, daß sie Männer sind Esel und Sitna, das ist, daß sie ungerechte und zänkische Leute sind. Dem Hebräischen lautet es nicht wohl. Esel, Zank, zänkische Leute, die nach seinem Frieden stehen, sind nicht Kinder der Liebe, sondern des Zankes und Haders, zu des Nächsten Nachtheil und Schaden. Es heißt, einem Andern Gewalt und Unrecht thun, aber mit einem Schein des Rechtes. „Das Land“, sagen sie, „ist unser“: das ist der Schein des Rechtes. „Darum ist der Brunnen auch unser“: das ist die Gewalt und das Unrecht. Und ist dieses eine gemeine und sehr böse Weise in der Welt, daß wer dem Andern mit Gewalt keinen Schaden zufügen kann, der sieht, wie er ihm heimlich mit List und Betrug Schaden thun möge. Die Deutschen heißen es und sagen davon also: Gewalt und Unrecht thun, das ist Essig. Und ist dieses ein böses Balet und bringt denen zu Gerar kein Glück. GOtt behüte und bewahre mich, daß ich dergleichen Segen oder Balet von einem frommen und gottseeligen Menschen nicht hören dürfe.

225. Nun dem Worte sitna kommt Satan her, und will Isaak so viel sagen: Sie vertreiben mich nicht allein und nehmen mir den Brunnen mit Gewalt unter dem Scheine des Rechtes, sondern sie verfolgen mich auch feindseliger Weise und sind meine Feinde: sie wollen nicht allein den Brunnen haben, und die neue Brunnenader, so ich zuvor gefunden und die ich ihnen jetzt auch eingeräumt habe, sondern sie nehmen mir diesen Brunnen auch mit Gewalt. Sie sind mir gram, sie hassen mich. Nun bedürfen sie ja dieses Brunnens nicht; sie können sich wohl genügen lassen an dem schönen Wasserflusse aus des neuen Brunnens Ader; warum nehmen sie mir denn diesen Brunnen? Sie gönnen mir das nicht, daß sie nicht bedürfen noch Lust daran haben. Ich habe zwar ihnen keine Ursache gegeben zu Feindschaft und sie gebrauchen selbst diesen Brunnen nicht, sondern es ist nur Gewalt; es ist ein verstockter Muthwille und die äußerste Bosheit. Derhalben sagt er zu ihnen: Ade zu guter Nacht und seid nur immerhin Esel und Sitna; nehmt den Brunnen hin unter dem feinen Scheine, deß ihr doch selbst nicht bedürft, auch keine Lust daran habt, ihn dazu auch nicht gebraucht.

226. Das sind die aller schändlichsten Leute, denen weder Vergnügen noch Nutzen Anlaß zu Raub und Sünde gibt. Denn sie hätten diesen Brunnen nie begehrt oder etwas darnach gefragt, wo ihnen jemand denselbigen auch umsonst hätte angeboten; aber dieweil sein ein Anderer gebraucht, vornehmlich aber Isaak, da werden sie gar entbrannt von Haß und Neid. Isaak aber leidet solches mit Geduld und sagt: Behaltet den Esel, weil ihr mir so gram seid.

227. Und das ist uns zum Exempel geschehen, auf daß wir auch lernen Geduld tragen, wenn uns von Andern Gewalt und Unrecht geschieht. Denn er gibt ihnen nun das Balet und sagt: Behaltet immerhin den Esel und Sitna, ich kann in diesem Thale nicht bleiben; und zieht also davon und ist in seinem Herzen wohl zufrieden. Das ist zumal eine große Geduld, die keines Gleichen hat, und dazu ein unüberwindlicher Glaube, von welchem Glauben weder die Rechte, noch die Künste, und auch die Medicin oder Arznei nichts lehren. Allein die heilige Schrift hält uns diese Exempel und Lehre vor.

\begin{quote}
\textbf{V. 22. 23. Da machte er sich von dannen, und grub einen andern Brunnen, da zankten sie sich nicht über; darum hieß er ihn Rehoboth, und sprach: Nun hat uns der HErr Raum gemacht, und uns wachsen lassen im Lande. Danach zog er von dannen gen Bersaba.}
\end{quote}

228. Nun zieht Isaak an einen andern Ort, und gräbt einen andern Brunnen, über den sie sich nicht gezankt haben. Da tröstet ihn der Herr nun wieder, da er ihm einen neuen Brunnen gibt, welchen er Rehoboth nennt, und sagt: „Nun hat uns der HErr Raum gemacht.“ Diese Worte sind hergekommen aus dem Glauben an die Verheißung. Denn er hat also gedacht: Ich habe müssen sehr gute Brunnen, Esel und Sitna, verlassen; ich will aber darum den Glauben nicht fallen lassen, sondern man muß verharren, und fest an der Verheißung GOttes halten, welcher ja nicht lügt. Nun wollen wir zuletzt noch reich werden, sagt er, dieweil wir einen eigenen Brunnen haben.

229. Siehe aber, wie GOtt seine Heiligen pflege zu führen und zu prüfen; bald macht er es ihnen süß, bald sauer. Wer wahrhaftig ein Christ ist, der wird oben solchen Wechsel erfahren, wie er in unseren Psalmen beschrieben wird: Angst, Trost, Trübsal und wieder Trost darauf. Er gibt aber diesem Brunnen einen Namen vom Troste, als wollte er sagen: Er soll heißen Trostbrunnen, ein Lustbrunnen, wie die andern Brunnen der Trübsal, des Hasses und Neides gewesen sind. Also haben wir an diesem Brunnen nicht allein leiblichen, sondern auch geistlichen Trost gefunden, auf daß wir GOtt loben mögen, der uns nicht verlässt. Jetzt werden wir nun wachsen, groß und reich werden in diesem Lande.

230. Aber diese Freude währt auch nicht lange; denn es folgt nun eine andere Reise, daß er gen Bersaba gezogen ist. Wiewohl Mose dessen hier nicht gedenkt, daß er diese Wohnung deswegen verändert hätte, daß ihn die Philister verfolgt hätten; sondern er wird vielleicht von selbst weggezogen sein, und wird entweder von Andern dazu gereizt, oder aber sonst durch andere Veranlassung bewogen worden sein, daß er einen besseren Ort gesucht hat. Da wird sich aber ohne Zweifel sein Gesinde von wegen der großen Bosheit der Leute zu Gerar beklagt haben, daß sie den so frommen und heiligen Mann vertrieben hätten. Es rühmt aber GOtt die große Geduld Isaaks, welche er uns allen zum Exempel vorgehalten hat; denn er versorgt ihn wiederum mit anderen Wirthen, die ihn freundlich und gern aufnehmen und beherbergen.

231. Und bisher haben wir gehört von den drei Brunnen und äußerlichen Anfechtungen des heiligen Patriarchen Isaak, damit ihn der HErr versucht und geprüft hat. Es leuchtet aber an ihm ein brünstiger und unüberwindlicher Glaube, dadurch er beständiglich verharrt und die schwersten Anfechtungen überwunden hat, uns zum Exempel, daß wir ja in keinerlei Noth oder Gefahr verzagen sollen. Denn dieweil der HErr lebt, der uns seine Verheißung gegeben hat, so werden wir nicht verlassen werden; wo er aber stirbt, so wird auch die Verheißung und Alles mit einander zu Trümmern gehen. So lange er lebt, dürfen wir nicht verzagen, und wir sollen uns befleißigen, daß unser Glaube nicht matt werde und unsere Hoffnung nicht aufhöre. Das gefällt ihm wohl und ist uns nütze und heilsam, und kein Unglück oder Anfechtung, wie schwer sie auch ist, kann so lange währen, daß sie Gott tödten könnte.
\quellseite{66}

232. In solchem Vertrauen haben alle Heiligen ihre Hoffnung gehabt und sind also beständig geblieben. Also sollen wir heutiges Tages auch thun. Es komme der Türke oder sonst ein ander Unglück, was es auch sein mag, sollen wir also denken: Der uns ausgesondert und in seine Kirche gesetzt hat, daß wir an ihn glauben und auf ihn hoffen sollen, der wird das wohl halten, was er uns verheißen hat, und wird uns bis an jenen Tag erhalten. Also hat Noah und alle Patriarchen geglaubt, und sind auch durch den Glauben zur Seligkeit erhalten worden.

233. Nun hat aber die äußerliche Anfechtung ein Ende, und läßt es sich ansehen, als habe Isaak mit seiner Geduld die Welt bezwungen und die äußerlichen Feinde überwunden; aber bald hernach wird er mit häuslicher Anfechtung von seinem Sohne Esau geplagt werden, da er seine zwei Weiber heimführen wird. Aber vor derselbigen Anfechtung geht abermal eine große herrliche Verheißung her.

\section[Fünfter Theil]{Fünfter Theil. Wie der HErr dem Isaak erschienen, und wie Isaak einen Altar baut und des HErrn Namen predigt}

\begin{quote}
\textbf{V. 24. 25. Und der HErr erschien ihm in derselben Nacht, und sprach: Ich bin deines Vaters Abrahams Gott. Fürchte dich nicht, denn ich bin mit dir, und will dich segnen, und deinen Samen mehren, um meines Knechts Abrahams willen. Da bauete er einen Altar daselbst, und predigte von dem Namen des HErrn, und richtete daselbst seine Hütte auf; und seine Knechte gruben daselbst einen Brunnen.}
\end{quote}

234. Ich habe oftmals ermahnt, daß man in den Historien und Legenden der heiligen Patriarchen vornehmlich auf das Wort Achtung geben solle, und daß man fleißig unterscheiden müsse die Werke der Patriarchen und das Wort Gottes als das vornehmste Stück in dem Leben der Väter, nämlich, daß wir vornehmlich darauf merken und sehen, was und wie oft Gott mit ihnen redet; denn Gottes Wort ist größer und mehr, denn aller Menschen, und dazu der allerheiligsten, ja, auch der Engel Werke sein können.

235. Das ist in den Historien der Väter der einige und der allerschönste Edelstein, nämlich, daß Gott mit ihnen und mit uns redet. Diesen Edelstein sieht die Welt nicht, versteht auch nicht, wie theuer und werth dieser Schatz ist, und wir selbst haben auch nicht genugsam Acht darauf, nämlich, daß im ganzen Leben das Wort das rechte Maß und das Allerwichtigste ist, welches unser Leben regiert, daß du sagen kannst: Dieses thue ich in Gottes Wort; dieses hat mich der HErr heißen thun; dieses gefällt Gott wohl; wie wir dessen können gewiß sein vom höchsten Stande an bis zu dem allergeringsten: Gott hat es befohlen, Gott hat es gesagt.

236. Wo aber das Wort nicht ist, da fliehen, wer nur fliehen kann. Wenn auch das Leben äußerlich anzusehen ist, als ob es ein engelisch Leben wäre. Wenn auch dasjenige, so man vornimmt, schon anzusehen ist, als sollte es Gott angenehm und ein sehr geistlich und heilig Ding sein, so muß das äußerliche Ansehen oder Gutdünken immer weg, wo du des Wortes nicht gewiß bist. Denn warum willst du dich vergeblich also plagen und die Zeit so jämmerlich und übel zubringen, wie die Mönche im Pabstthum, item, die Türken und Juden in dem menschlichen Wähnen einhergehen, daß sie meinen, alsdann gefallen sie Gott wohl, wenn sie sich mit Fasten und Kasteiung des Fleisches zu Tode martern; wie von St. Bernhard gesagt wird, daß er seinen Leib mit Fasten also gemartert habe, daß ihn die andern Mönche im Chore nicht haben bei sich leiden können um seines stinkenden Odems willen. Also hat er nämlich die böse Lust zur Unzucht im Zaume halten und auslöschen wollen. Und also sind auch alle Werke in diesem Leben der Menschen, so ohne Gottes Wort vorgenommen werden.

237. Derhalben thue ich diese Vermahnung so fleißig und zum öftern, und man muß solches den Leuten allezeit einschärfen, daß wir uns ja unsern eigenen Wahn oder Gedanken nicht führen und regieren lassen, wenn sie auch schon göttlich, englisch und himmlisch wären; wie St. Paulus Col. 2, 18. auch ermahnt, da er sagt: „Lasset euch niemand das Ziel verrücken, der nach eigener Wahl einhergehet in Demuth und Geistlichkeit der Engel, daß er nie seines gesehen hat, und ist ohne Sache aufgeblasen“ etc. Denn solche Gedanken sind ohne Wort. Denn darum redet Gott mit uns und handelt mit uns durch die Diener des Wortes, durch unsere Eltern und durch die Obrigkeit, auf daß wir uns nicht wägen und wiegen lassen von allerlei Wind der Lehre. Die Kinder sollen ihre Eltern, die Bürger und Unterthanen ihre Obrigkeit, ein Christ seinen Pfarrherrn und die Diener des Wortes, ein Schüler seinen Schulmeister hören. Außerhalb des Wortes ist alles Leben verdammt, und ist mit allen Secten und Orden verloren: wo aber das Wort ist, alsdann habe ich einen gewissen Trost, ich sei Vater oder Mutter, oder ein Kind im Hause, da höre ich das Wort, und weiß, was ich glauben und thun soll. Denn Gott redet mit mir auch in dem Stande, darin ich lebe.

238. Dagegen aber, wo man das Wort verfälscht und verkehret, das ist ein greulicher Zorn Gottes, und ist eine schwere Verachtung und Strafe über die, so das Wort verachtet haben. Denn wo du Gott nicht hören willst, da er dir die gewisse Wahrheit sagt, ei, so magst du die Lüge unter dem Schein der Wahrheit hören; wie St. Paulus 2 Thess. 2, 10. 11. von dieser Strafe der Verachtung des Wortes auch redet, und beweisen solches auch die elenden und betrübten Exempel, so man davon hat. Griechenland muß heutiges Tages den Mahomet hören; die, so gegen Abend oder Niedergang der Sonne wohnen, müssen den Pabst hören.

239. Man soll zwar die Werke und Exempel in den Legenden der Patriarchen auch loben und sich derselben verwundern; denn sie auf das Wort mehr Achtung geben, welches das Leben der Patriarchen führt und regiert. Das ist die rechte Regel und Richtschnur, daraus ich also schließen kann: Ich bin ein Prediger, ein Ehemann, ein Gewerb, ein Knecht, eine Magd oder ein Kind im Hause: dieses hat mir Gott befohlen zu thun und darum gefällt es ihm auch wohl, nicht um der Werke willen an sich selbst, denn wir sind Fleisch, sondern um deswillen, der uns führt und regiert, welches nämlich das Wort ist, darunter ich einher gehe. Wenn ich sterbe, so sterbe ich nicht in der Regel oder dem Orden des Franciscus, sondern in der Regel und im Worte des Heiligen Geistes.

\subsection{I}

240. Nun wollen wir auch die Worte, wie sie lauten, besehen. Ego, dii Abraham: Ich, die Götter Abrahams. Denn es steht im Hebräischen das Wort elohim, welches in der Mehrzahl Götter heißt; wie droben im 1. Cap. V. 1. steht: „Im Anfang schuf Gott Himmel und Erde“ etc. Daselbst gebraucht Mose auch das Wort elohim. Wir Christen sollen und müssen an der Einigkeit Gottes, daß nur Ein Gott sei, nicht zweifeln; was aber dasselbe für eine Einigkeit sei, ist unserer Vernunft und menschlicher Weisheit unerforschlich; denn es ist eine Offenbarung des Heiligen Geistes und des Glaubens.

241. Die Juden und Türken spotten unser und dichten auf uns eine schändliche und unverschämte Lüge, als sollten wir viel Götter annehmen, und glauben, daß drei Götter wären. Derhalben rühmen sie sich, daß sie Gottes Volk seien, daß sie mit vielen herrlichen Siegen und Reichthum der ganzen Welt gesegnet und begnadet werden, darum daß sie an einen Gott glauben; wir aber, sagen sie, haben so viel Niederlagen erlitten und sind so manchmal geschlagen worden, darum daß wir dichten, daß viel Götter sein sollen; und darum nennen sie uns abgöttische Leute, und lästern den Sohn Gottes greulich, wie das alle ihre Historien bezeugen.

242. Derhalben müssen wir uns in unsern Herzen rüsten gegen solche Lästerworte, und wider den Ruhm, der (wie es vor der Vernunft (scheint) gar zu wahr und greifbar ist. Denn das ist ja gewiß, daß wir glauben, daß nur ein einiger Gott sei, der da ist lauter und rein von aller Zusammensetzung und Vermengung oder Mischung; daß aber der Türke unsern Glauben oder Lehre nicht recht versteht, und nicht hört, wer dieser einige Gott, dieses lautere, reine und unvermengte göttliche Wesen sei, das ist seine eigene und nicht unsere Schuld.
\quellseite{67}

243. Denn wir lehren und glauben nicht allein, daß Ein Gott sei, sondern daß er durchaus lauter und rein sei von aller Zusammensetzung und Vermengung, und so vollkommen einig sei, daß außer ihm kein Gott sein kann; wir theilen oder trennen diese Drei nicht von einander, den Vater, den Sohn und den Heiligen Geist; wir machen nicht unterschiedliche Götter: sondern glauben, daß in rechter Einigkeit ein einiger Gott, ein lauteres, reines und unvermengtes göttliches Wesen, sei.

244. Ja, möchte irgend Einer sagen, ich verstehe das Ding nicht. Antwort: Das ist recht, du sollst es auch nicht verstehen. Wir glauben, daß Ein Gott sei, welcher doch von sich selber pflegt zu reden und sich zu zählen, als wären seiner Viele; wie droben, so oft gesagt ist: „durch deinen Samen“ etc., daselbst ist der Same von Noth wegen und wahrhaftig rechter Gott, dieweil er ohne Sünde ist, und Macht und Gewalt hat, die Sünde zu zerstören und abzuschaffen. Eben auf solche Weise wird an diesem Ort Gott auch genannt „Götter Abrahams“. Und wir fragen nach den groben Ochsen und Eseln, den Türken und Juden, gar nichts, die da sagen, sie können es mit menschlicher Vernunft nicht verstehen; und sie sollen es auch damit nicht verstehen, da solche große Dinge nur aus göttlicher Offenbarung bekannt sind.

245. Darnach, wenn wir diese Worte etwas fleißig ansehen werden, so werden wir sehen, daß man daraus auch schließen und abnehmen kann, daß eine Auferstehung der Todten sei. Denn wie man im gemeinen Sprüchwort sagt: Man muß der Fürsten Briefe dreimal lesen: also muß man viel mehr die heilige Schrift auch wohl tausendmal lesen und immer wiederholen. Denn wir wissen, daß Abraham gestorben ist und daß ihn seine Söhne Isaak und Ismael begraben haben. Nun sagt aber hier Gott: „Ich bin Abrahams Gott“, und Christus sagt Matth. 22, 32.: „Gott ist nicht ein Gott der Todten“; derhalben macht er hier Abraham lebendig und erweckt ihn von den Todten. Dies ist der Schluß und Beweis, so Christus führt Matth. 22, 32., da er sagt: „Ich bin der Gott Abrahams und der Gott Isaaks, und der Gott Jakobs: Gott aber ist nicht ein Gott der Todten, sondern der Lebendigen.“ Gott ist desjenigen Gott nicht, das an sich selbst nicht ist. Keiner und Niemand beten Gott nicht an und Gott regiert über sie nicht. Abraham ist gestorben und Gott ist Abrahams Gott: so folgt ja, daß Abraham leben muß. Er ist zwar gestorben und begraben, aber mir, spricht Gott, der ich Gott bin, lebt Abraham, und hat die Dreifaltigkeit der Person und Christum, seinen Samen, erkannt. Denn Gott ist nicht ein Gott des, das da nichts ist.

246. Dieses Argument haben die Propheten durch Erleuchtung des Heiligen Geistes herrlich gesehen und wohl verstanden. Der Türke sieht diese nicht, und der Papst und seine großen Doctoren haben es auch nie verstanden; denn sie gehen in die Schrift nach dem Verstande eines Moffes und Mäusels, wie der 32. Psalm V. 9. sagt. Wenn sie etwa gelesen haben Matth. 1, 2.: „Abraham zeugete Isaak“ etc., so haben sie gedacht, dieselbigen seien gute, schlichte, fromme Männer und Eheleute gewesen, an denen sei nichts, deß man sich hoch verwundern möchte. Wenn aber der Heilige Geist dazu kommt, der die Herzen erleuchtet, da wird alsdann das Licht hervorleuchten von der Unsterblichkeit und von ewigen Leben. Denn wo Abraham einen Gott und wiederum Gott einen Abraham hat, so folgt nothwendig wiederum das auch, daß Gott und Abraham zugleich leben müssen; denn diese zwei stehen und fallen mit einander, sintemal Gott mit den Todten nichts zu thun hat.

247. Diese Lehre, daß mehr denn Eine Person in der Gottheit ist, und von der Auferstehung der Todten, hat Noah verstanden, und ohne Zweifel auch gelehrt und gepredigt. Denn es ist das Wort, so in die Kirche gehört, und ist die Predigt des Evangeliums. Er hat erstlich gebetet, und darnach gepredigt, und wird also gesagt haben: Mir ist von Gott geöffnet worden, daß mein Vater Abraham in Gottes Gnade und Barmherzigkeit lebe; derhalben so werden wir auch nicht sterben. Denn ob wir wohl vor der Welt sterben, so werden wir doch von den Gottlosen abgesondert werden, welche in dieser Welt das Ansehen haben, daß sie leben und daß sie einen Gott haben, der ihnen günstig und gewogen sei: aber sie haben einen zornigen Gott; wir aber haben hier und werden dort auch einen gnädigen Gott haben.

248. Daher gehört das auch, daß er hinzu setzt: „Um meines Knechtes willen“ 2c. Denn wie ist Abraham nach dem Tode ein Knecht Gottes? Ob denn Gott des Abraham nicht endlich wird vergessen können? Er dient wahrlich noch heutiges Tages Gott, wie Adam, Abel und Noah Gott dienen. Und dies soll man fleißig merken, denn es ist die göttliche Wahrheit, daß Abraham mit Gott lebt, ihm dient und auch mit ihm regiert. Was aber das für ein Leben sei, ob er schlafe oder wache, das ist eine andere Frage.

249. Wie die Seele ruhe, sollen wir nicht wissen, es ist aber gewiß, daß sie lebt. Siehe die Menschen an, die entzückt sind oder schlafen. Von denselbigen kann niemand sagen, was sie sind, auch wenn sie leiblich schlafen; denn sie sind ohne Empfindung und als wären sie todt, und kann doch nicht gesagt werden: Dieser ist todt, sondern er lebt im Schlaf oder im Traum, wie das die ganze Welt bezeugen muß. Ich fühle aber doch nicht, daß ich lebe, wenn ich schlafe; denn da sind alle Sinne und die Vernunft selbst auch gar stille, und richten der Dinge keines aus, das ihr Amt ist: da weiß niemand, wo er ist. Wenn wir im Schlaf lebten, so könnten wir im Schlaf auch denken: Ich bin in diesem Hause, in dieser Kammer. Nun geschieht es aber oftmals, daß ich mich im Schlaf oder Traum dünken lasse, daß ich in der Hölle, im Himmel, zu Venedig oder sonst an andern Orten sei. Derhalben ist das ja ein groß Zeichen, daß ich lebe und doch gleichwohl nicht lebe.

250. Wie können wir nun meinen, daß es um die Seele stehe, wie sie ruhen oder leben möge? Sie wird ohne Zweifel ihre eigene Weise haben, wie sie schläft, welche ich nicht weiß, wie ich auch den leiblichen Schlaf nicht verstehen kann. Denn ich habe oftmals mit Fleiß merken wollen auf den Augenblick, da ich eingeschlafen bin, und auf den Augenblick, da ich wiederum würde aufwachen; aber ehe ich darauf Achtung gehabt, bin ich schon aufgewacht. Solche Exempel unseres Schlafes, welcher auch eine Art Tod ist, beweisen, daß die Seelen der Heiligen ruhen, wie Jer. 26, 20. und 57, 2. gesagt wird, und dasselbe noch viel stiller und friedlicher, als die da schlafen. Dieselbigen ruhen zwar auch, und um der Ruhe willen leben die Menschen niemals mehr, als wenn sie schlafen. Denn das Leben derjenigen, die da wachen, ist voller Sorge, Jammer, Bekümmerniß, Mühe und Arbeit: der leibliche Schlaf aber ändert und überwindet die Krankheit; wie die Jünger sagen Joh. 11, 12. von Lazarus: „Herr, schläft er, so wirds besser mit ihm.“ Derhalben sagen wir recht, daß wir dann am allermeisten leben, wenn wir schlafen und sterben; denn der lebendige Odem ist dann am kräftigsten. Eben aber auf solche Weise schlafen die Seelen auch; wie aber das zugehe, verstehen wir nicht.

251. Und das ist kein Wunder, dieweil wir andere Dinge auch nicht verstehen, die uns doch bekannter sein sollten, darin wir geboren und auferzogen sind. Denn siehe deine Kindheit an, und bedenke, ob du dich dessen auch wissest zu erinnern, daß du in deiner Mutter Leibe gewesen seiest, daß du in der Wiege gelegen, dazu deiner Mutter Brüste gesogen, geschrien und Brei gegessen habest, und wie du gewachsen seiest 2c. Nun leben wir wahrlich auch im ersten Jahre, da die Frucht in Mutterleibe getragen wird; wie wir aber gelebt haben, davon wissen wir gar nichts. Also auch nach der Geburt weiß ein jährlich Kindlein nichts vom Leben, es weiß nicht, daß es lebt, es kann das Leben nicht überlegen, wir sehen an solchen Kindlein das Leben ohne Leben. Also habe ich sechzig Jahre gelebt und habe auch in Mutterleibe gelebt; aber von diesem Leben habe ich hernach nie etwas gewußt. Und das dennoch die Frucht im Mutterleibe gewißlich und kräftig leben müsse, wird dadurch angezeigt, daß sie sich zum erstenmal bewegt, zudem, daß auch das Kindlein schreit, wenn es aus dem Mutterleibe kommt. Da wir aber nun dies mit unsern Gedanken nicht erreichen können, so werden wir viel weniger das begreifen, wie es um das Leben nach dem Tode stehe, und ist doch gleichwohl gewiß, daß wir leben. Darum soll man das Forschen unterwegen lassen und es Gott befehlen; wie man dem vorigen Artikel von der Dreifaltigkeit auch nicht nachforschen soll, nämlich, wie Gott zugleich einig und nach den Personen dreifaltig sein möge. Euch aber also: höre das Wort, dasselbe laß dich führen und regieren, du aber thue die Augen zu und folge immer dem Wort nach, alsdann wirst du finden, wie Gott einig und dreifaltig sei.
\quellseite{68}

253. Auf solche Weise höre ich hier auch, daß im Tode und nach dem Tode das Leben und daß eine Auferstehung der Todten sei; wie aber dasselbe werde zugehen, damit bekümmere du dich nicht, dieweil du nicht weißt, daß du lebst, wenn du schläfst, und lebst dennoch allezeit; also hast du auch gelebt in Mutterleibe und da du ein kleines unmündiges Kind warest, daß du davon auch selbst nichts gewußt hast. Denn sollte unser Herr Gott nicht mehr Weisen haben zu leben als diese zweierlei allein, daß ein Leben ohne Leben und Verstand sei? Also kann die Seele auch ihre eigene Weise haben zu leben, welche über unsern Verstand ist.

254. Unterdeß sollen wir uns an diesen Trost halten und darauf mit festem Vertrauen warten, welchen Gottes Wort vom Himmel uns vorhält, daß er sagt: „Ich bin dein Gott, ich bin der Gott Abrahams, ihr seid meine Knechte etc. Also steht Jes. 26, 19.: „Wachet auf und rühmet, die ihr lieget unter der Erde; denn dein Thau ist ein Thau des grünen Feldes.“ Daselbst redet Gott mit den Todten nicht anders, denn als ob sie lebten.

255. Und dasselbe Wort ist ein sehr kräftiges Zeugniß, daß wir nicht sterblich sind, sondern unsterblich auch im Tode. Und ist dies die Ursache, weil Gott mit uns redet auch in unserer Sprache und wie wir pflegen zu reden. Gott weiß, daß dies Leben vergänglich ist; warum wollte er aber mit uns reden, und dasselbe also, daß er auch unserer Sprache gebrauchte, wenn wir nicht ewig leben? denn sonst redete er sein Wort ja vergeblich, nur um einer kurzen Zeit und Augenblickes willen. Nun redet er aber nicht vergeblich, und nimmt ja der Ochsen oder Esel Stimme nicht an sich, daß er damit redet: er brüllt ja mit ihnen nicht; er redet allein mit dem Menschen. Wo nun und mit wem Gott redet, der sei gleich wer er wolle und er rede im Zorn oder in Gnaden mit ihm, derselbe ist gewißlich unsterblich. Die Person Gottes, der da redet, und das Wort zeigen an, daß wir solche Creaturen sind, mit denen Gott bis in Ewigkeit und unsterblicher Weise reden wolle. Einen solchen Gott oder Götter (wie hier nach dem hebräischen Texte steht) hat Abraham, und wer der Verheißung Abrahams anhängt, der hat auch denselbigen Gott und ist ein Knecht Gottes, und wird endlich, wenn er gestorben ist, auch im Schlafe leben.

256. Es wiederholt aber Gott den Segen darum, daß dieser heilige Mann Isaak so mancherlei Weise ist geplagt worden, und ohne Zweifel in der Anfechtung wird gedacht haben: Ach ich armer elender Mann, was soll ich doch thun, wohin soll ich mich wenden? Denn die vergangenen Anfechtungen machen uns ängstlich und bekümmern, daß wir denken: wie lange werde ich wohl diesen Frieden haben? Vielleicht wird er etwa Einen Monat oder auf das längste zwei währen? Denn er hat den Brunnen Rehoboth auch verlassen müssen. Und vielleicht mag Rebekka angefangen haben zu klagen und zu zürnen wegen des Unglücks, und daß es sich mit ihnen so oft und auf so mancherlei Weise geändert hat. Sie wird gesagt haben: Wer wird dies auf die Länge dulden und leiden können, so lange in der Irre umher zu ziehen und keinen gewissen Sitz oder Wohnung zu haben? Denn ob dies wohl fleißig Ding ist, wird doch gleichwohl der Glaube dadurch heftig angefochten.

257. Darum halte ich dafür, daß Isaak schier würde müde geworden sein von so großem und vielem Unglück, und daß Rebekka auch aus Ungeduld werde gemurret haben; wie das Weib Tobias ihm seine Almosen aufrückte, Tob. 2, 15. 22. 23., da sie sagt: „Da siehet man, daß dein Vertrauen nichts, und deine Almosen verloren sind. Mit solchen und andern Worten mehr warf sie ihm sein Elend für.“ Also halte ich dafür, daß hier auch solcher gestalt Rebekka sammt Isaak werden matt und müde geworden sein, und dasselbe wird durch das Klagen und Geschrei des Gesindes immer gemehrt worden sein. Derhalben tröstet ihn der Herr und sein ganz Gesinde auch und spricht: „Fürchte dich nicht.“ Daraus folgt, daß Isaak furchtsam und erschrocken gewesen sei und sein Glaube angefochten worden ist. Denn der Teufel pflegt die Gottseligen müde zu machen nicht allein mit der Größe und Schwere der Anfechtungen, sondern daß er damit auch immer anhält und ihnen für und für zusetzt, läßt ihnen keine Ruhe. Derhalben spricht nun Gott zu ihm: Mein lieber Isaak, sei du stark, fürchte dich nicht; und also richtet er ihn auf und erquickt ihn, dieweil er im Geist fast müde war. Das thut er aber mit einem sehr schönen Trost, nämlich, mit der Offenbarung der Unsterblichkeit, und sagt: Du wirst leben und wirst in Ewigkeit leben; denn ich bin deines Vaters Abraham Gott, der lebt, darum wirst du auch leben.

258. Diese Worte soll man fleißig merken; denn Gott redet nicht vergeblich. Er sagt zum Ersten: „Ich bin mit dir, aber doch heimlich und verborgener Weise. Ja, möchte Isaak sagen: Empfinde ich es doch nicht; du stehest zumal ferne von mir. Wie Jeremia auch klagt im 12. Cap. V. 2. von den Gottlosen: „Du lässest sie viel von dir rühmen, und züchtigest sie nicht“, als wollte er sagen: Die Andern sind fröhlich, springen und jauchzen, du bist fast nahe bei ihnen.

259. Wie jetzt der Türke hoch rühmt, daß Gott auf seiner Seite stehe und es mit ihm halte; wie er mit so vielen Siegen und so großem Reichthum von Gott gesegnet sei worden. Denn er hat einen großen reichen Segen Gottes. Unterdeß läßt es sich ansehen, daß der Teufel nahe bei den Christen sei, die von den Türken gefangen und jämmerlich erwürgt werden, und daß Gott sehr ferne von ihnen sei. Bei den Gottlosen ist der Teufel gestorben und, wie sie sich dünken lassen, regiert Gott unter ihnen allein. Aber er ist eigentlich ganz ferne von ihnen, und sie sind, wie im Propheten Jeremia im 46. Cap. V. 20. 21. steht, „gemästete Kälber“ auf den Tag, da sie sollen geschlachtet und geopfert werden. Eben also wird es auch dem Türken und Pabst gehen, ob sie wohl ihre glücklichen Verrichtungen, so sie haben, dahin deuten, daß es Zeichen der Gnade und nicht des Zorns Gottes seien.

260. Du aber, wenn dich Gott versucht und stäupt, glaube, und laß dir an dem Worte genügen, daß du hast, da Gott spricht im 91. Psalm V. 15.: „Ich bin bei dir in der Noth, ich will dich heraus reißen“, nicht allein aus der Noth, sondern auch aus dem Tode, Krankheit, Schande und Schmach etc., ich will für dich sorgen. Dasselbe sehen zwar die Augen nicht, es greifen es auch die Hände nicht, aber den Gläubigen sind alle Dinge möglich.

261. Und was nun an diesem Orte dem Patriarchen Isaak gesagt ist, eben dasselbe wird uns auch gesagt, nämlich, da Gott spricht: Ich bin mit dir, und will halten, was ich verheißen habe; du sollst von mir nicht allein leibliche, sondern auch ewige Güter zu erwarten haben. Ja, möchte Isaak sagen: Ich sehe aber das Gegentheil; es läßt sich ansehen, als ob du mich verlassest; ich muß schier vor Durst sterben mit meinem Weibe, mit meinen Kindern und ganzem Hausgesinde: ist das ein göttlicher Segen? Da antwortet Gott, und spricht: „Fürchte dich nicht“, du sollst nicht Durstes sterben. Und sagt noch weiter: „Ich will deinen Samen mehren, um meines Knechtes Abrahams willen.“ Ich will dich nicht allein jetzt gegenwärtig mehren, sondern auch deinen Samen und deine Nachkommen.

262. Dies ist eine Wiederholung des Segens mit einer Vermehrung desselbigen, dadurch er ihn stärkt und des ewigen Lebens und der Unsterblichkeit gewiß macht. Warum sagt er aber „um meines Knechtes Abrahams“, und nicht, um Isaaks willen? Antwort: Er will damit das Exempel seines Vaters Abraham loben und ihm dasselbe vorhalten, daß er demselben nachfolgen soll. Als wollte er sagen: Ob du es wohl nicht verdienst hast, daß ich mit dir reden sollte, so will ich doch um deines Vaters willen, welchen ich so sehr lieb habe, mit dir reden. Dies ist ein sehr lieblicher Trost, und ist viel angenehmer, denn wo er das Exempel Abrahams nicht angezogen und gesagt hätte: Um deinetwillen, Isaak, will ich deinen Samen mehren etc. Und dies ist die letzte Verheißung, damit sich Isaak bis an das Ende seines Lebens aufrecht gehalten und getröstet hat.

\subsection{II}

263. Derhalben baut er daselbst einen Altar und ist bis an seinen Tod zu Bersaba geblieben, und in der Höhle, davon im 23. Cap. V. 17. 19. gemeldet wird, begraben worden. Er hat aber daselbst einen Altar gebaut nicht um des Opfers, sondern um der Predigt des Wortes willen. Denn an allen Orten, da gedacht wird, daß sie Altäre aufgerichtet und Hütten gemacht haben, daselbst wird angezeigt, daß sie an solchen Orten ein Kirchlein angerichtet haben, da die Leute sind zusammengekommen, daß man daselbst gelehrt, Gottes Wort gehört, gebetet, Gott gelobt und geopfert hat. Also sind unsere Kirchen solche Orte auch, da man zum Gebet, und Gott zu loben und zu preisen zusammen kommt. Denn ein Altar ist kein Ort oder Stätte im Hause, sondern in der Kirche, und die Werke, so man in der Kirche verrichten soll, sind, daß man Gottes Wort lehre und höre, dazu Gott lobe und ihm danke etc.
\quellseite{69}

264. Daselbst nun hat Isaak viel schöne Predigten gethan fast siebenzig oder achtzig Jahre. Da hat er von dem Geheimniß der Dreifaltigkeit, item, von der Menschwerdung Christi, von der Unsterblichkeit, und von allen Stücken, so jetziger Zeit im Evangelium gelesen und gelehrt werden, gepredigt. Daselbst wird aber ein fein Kirchlein gewesen sein, etwa unter einem Baum aufgerichtet, unter welchem er, sein Gesinde und Nachbarn umher sind zusammen gekommen. Daselbst hätte ich gewiß gerne ansehen mögen; denn es ist mit keinem Golde, Silber oder Edelsteinen geschmückt gewesen, sondern Gott hat daselbst mit lebendiger Stimme geredet. Die Hütte ist auch nicht köstlich, daß sie einen großen Schein gehabt hätte, sondern von Bäumen aufgebaut und aufgerichtet gewesen. Aber sie ist derhalben herrlich und trefflich groß, weil daselbst ein vortrefflicher großer Prophet und eine Prophetin, Rebekka, gewohnt haben, zudem auch viel heilige Menschen von seinem Gesinde, die zu Isaak gekommen sind, da er gelehrt und gepredigt hat, und haben ihn mit großer Ehrerbietung gehört, haben Gott für diese Lehrer und Prediger gedankt.

265. Also thut nun Mose: nachdem er die leiblichen und äußerlichen Dinge ausgerichtet hat, die dazu gehören, daß man damit den Leib erhalten muß, darinn gleichwohl viel schöne Exempel sind des Glaubens, der Liebe, der Hoffnung, Demuth und Geduld: setzt er auch nun in die Historie die Stücke und Werke, so zu der Religion und Glaubenssachen gehören, als, die Verheißungen, die Predigten, und wie sie Gott gelobt und geistliche Psalmen gesungen haben, damit die Heiligen sich in ihrem Glauben gestärkt und unter einander in der Noth getröstet haben. Denn das ist das Ende, dahin alle Predigten und die ganze himmlische Lehre soll gerichtet werden, nämlich, daß der Glaube gemehret, die Verheißung von Gottes Gnade gerühmt, ausgebreitet und in die Herzen der Menschen gepflanzt werden möge; zudem, daß die Geduld und andere Früchte des Glaubens auch mehr mögen gemehret werden.

266. Wir behalten aber an diesem Orte den gemeinen Unterschied zwischen diesen zwei Stücken, als, „den Namen des Herrn anrufen“ und „im Namen des Herrn anrufen“; denn das Erste heißt anrufen oder beten, als, im 50. Psalm V. 15.: „Rufe mich an in der Noth.“ Aber „im Namen des Herrn anrufen“ heißt so viel als, predigen im Namen des Herrn; als, droben im 4. Cap. V. 26. dieses ersten Buches Mose: „Zu der Zeit fing man an zu predigen von des Herrn Namen“ etc. Das hebräische Wort kara heißt eigentlich, rufen, nennen, aus dem Buche lesen, predigen. Bisweilen heißt es auch wohl, begegnen. Ich halte aber dafür, daß Mahomet seinem Buche daher den Titel und Namen gegeben habe, daß er es den Alkoran nennt, nämlich, daß es ein Lehrbuch, darin seine Lehre zusammengetragen, und gleichsam seine Bibel ist, wie der Papst seine Decretalien nennt. Derhalben, so verstehen wir das Wort an diesem Ort also, daß es heiße, lehren oder öffentlich in der gemeinen Predigt etwas lesen; wo sie anders Bücher gehabt haben.

267. Es wird aber zugleich hier auch beschrieben, was es für eine Predigt gewesen sei, nämlich, daß er im Namen des Herrn gepredigt, daß man von unserm Herrn Gott gepredigt hat; es ist eine Lehre gewesen von Gott. Es ist nicht eine lose unnütze Predigt von Menschenlehren gewesen, wie des Papstes Lehre ist, sondern es ist eine rechte, wahre und reine Lehre gewesen, welche lauter und rein vom Herrn und seinem Namen aus der Verheißung gelehrt hat, auf daß damit der Glaube geübet würde, welcher ein Werk ist der ersten Tafel.

268. Darnach, in der andern Tafel, ist die Lehre gewesen von guten Werken. Denn wir haben zweierlei Lehren: die eine ist von den göttlichen Verheißungen, so in die erste Tafel gehören, da Gott mit uns redet, und uns verheißt, daß er unser Gott sein wolle, auf daß wir daraus Glauben, Liebe und Hoffnung zu Gott fassen mögen; und diese Lehre erstreckt sich sehr weit und geht durch alle Werke. Der andere Theil der Lehre ist von guten Werken und Früchten des Glaubens. Also sind dies die vornehmsten Stücke der christlichen Lehre, nämlich, die Verheißung und das Gesetz.

269. Zum Dritten kommen dazu auch die Ceremonien, als, die Beschneidung, welche Isaak von Gott gehabt; dieselbige hat er auch seinen Zuhörern vorgetragen. Und er wird ohne Zweifel ein sehr reiner Lehrer gewesen sein, der die Kirche und Gemeinde Gottes fleißig gelehrt und unterrichtet hat im Worte des Glaubens und der Wahrheit.

270. Und dasselbe hat er recht und wohl thun können zur Zeit des Friedens. Denn wo man nicht Friede und Ruhe hat, da ist es nicht Zeit und hat auch keine Statt, Hütten zu bauen oder Altäre aufzurichten, es können auch die Leute nicht gelehret, und weder Welt- noch Hausregiment recht verwaltet werden. Denn es gehört dazu ein stiller Ort und eine friedsame Zeit. Zur Zeit aber der Anfechtung sollen wir deß gebrauchen, das wir im Frieden gelernt haben. Wie Isaak gethan hat: da er sieht, daß er von Gottes Gnade Frieden und Ruhe gehabt, da ist er nicht träge und faul, läßt sich den Müßiggang nicht verderben, befleißigt sich nicht, wie er viel Geld und Gut sammeln möge; sondern gebraucht die Zeit also, daß er in der Ruhe am allermeisten zu arbeiten hat; denn er versteht es dahin, daß ihm die Ruhe dazu gegeben sei, daß darinn die Gemeinde solle gelehrt und unterrichtet werden um der Nachkommen willen.

271. Dasselbe sollen wir auch thun, da die Mönche und Dompfaffen faule Schelme sind, und nur fressen und saufen, und des Friedens und der guten Tage, des Rehoboth und Ueberflusses aller Dinge genießen. Sie mißbrauchen aber aller der Dinge schändlich, nämlich, daß sie dabei nur immer im Hause leben, ihre Wollust und allerlei teuflische Sünde und Schande treiben: so doch die Schätze der Kaiser und Könige dazu nicht den Kirchen und Stiften gegeben sind, sondern daß man die Werke, so zur Gottseligkeit gehören, damit anrichten und erhalten soll. Als, wenn du Friede und Ruhe hast, sollst du, wie Isaak gethan hat, derselben dazu gebrauchen, daß der Gottesdienst gebessert und immer weiter ausgebreitet möge werden. Also sollten die Dompfaffen die heilige Schrift lesen und lernen, auf daß sie selbst darnach Andere auch von der Religion lehren könnten, daß sie beteten, dem Altare dieneten und im Namen des Herrn predigten.

\section[Sechster Theil]{Sechster Theil. Von dem Bunde Abimelechs mit Isaak; von dem neuen Brunnen Isaaks, und von Esaus Verheirathung}

\subsection{I}

\begin{quote}
\textbf{V. 26. Und Abimelech ging zu ihm von Gerar, und Ahusath, sein Freund, und Phichol, sein Feldhauptmann.}
\end{quote}

272. Diese Historie ist der gleich, die droben im 21. Cap. V. 25. steht, daß es fast scheint, als sei es ein und dieselbe Historie. Denn daselbst zankten Abimelech und Abraham auch über den Brunnen, und Abimelech entschuldigt sich, daß er davon nichts gewußt und daß der Brunnen dem Abraham ohne sein Wissen sei genommen worden. Aber diese Historie ist noch etwas klarer und erklärt die vorige; denn es wird die Glosse dazu gesetzt, daß der Brunnen viele gewesen sind.

273. Dies dient aber nun auch zum Trost Isaaks. Denn es hat ihn Gott schon mit der Verheißung aufgerichtet und hat ihm Frieden und Ruhe gegeben, daß er sicher und stille sein möchte, den Altar und die Hütte aufzurichten. Derselbigen Ruhe gebraucht er nun also, daß er gleichwohl dabei seine Ruhe und gute Tage haben und im Müßiggang leben wollte, sondern daß er thun möge, das einem Propheten und Priester Gottes gebührt zu thun: er arbeitet im Wort und in der Lehre, ist lehrhaftig, ist geschäftig und fertig, das Lehramt zu verwalten; darum folgt nun eine Vermehrung des Trostes, nämlich, daß Abimelech selbst, der ihn zuvor vertrieben hatte, nun von selbst und freiwillig Frieden und einen Bund mit ihm macht. Ich halte aber gänzlich dafür, daß dieser Abimelech des Sohn sei, dessen droben (Cap. 20.) ist gedacht worden. Er ist aber listiger als sein Vater; denn er vertreibt Isaac, und stellt sich, als wisse er davon nichts, daß seine Knechte Isaaks Brunnen genommen hätten.
\quellseite{70}

274. Derhalben muß man erstlich erkennen, daß es eine Wohlthat und Gnade Gottes ist, dadurch Isaac mit dem Könige versöhnt wird. Denn also sagt Salomo Spr. 16, 7.: „Wenn jemandes Wege dem HErrn wohlgefallen, so macht er auch seine Feinde mit ihm zufrieden.“

275. Darnach wird mit diesem Exempel klarlich bewiesen, wie große Kraft die Geduld habe, die Herzen der Menschen zu erweichen; denn es ist eine allmächtige Kraft, die aus einem Feinde den allerbesten Freund machen kann. Dieser Abimelech, der zuvor Isaac auf das äußerste feind war, ihn auch in seinem Königreich nicht hat leiden können, der kommt von selbst, ungebeten, mit seinen Fürsten und andern großen Herren wiederum zu Isaac, nicht mit dem gemeinen Pöbel, sondern mit seinem Freunde Ahusath und Phichol, seinem Feldhauptmann, bieten ihm Frieden an und begehren mit Isaac Freundschaft zu haben, und versprechen ihm Sicherheit und allerlei Wohlthat, damit sie ihm dienen wollen. Das ist eine große Wohlthat und Ehre, damit Gott Isaac begnadet, auf daß er ja reichlich Trost, Frieden, Schutz und Schirm haben möge nicht allein vom Könige, der nun den Verleumdern und den bösen schalkhaftigen Hirten nicht mehr zuhört, sondern er hat auch die Obersten und Vornehmsten im Königreiche zu Freunden, also daß dieser Schutz nun gewiß und stark genug ist, ohne alle Furcht und Gefahr.

276. Dies kommt alles daher, weil Isaac in seiner Geduld geblieben, und nicht aufgehört hat, den undankbaren Gutes zu thun, und er selbst um der Undankbarkeit und Bosheit der Andern willen nicht ärger geworden ist. Wie ein Weinberg oder Feigenbaum, wenn er schon zerrissen und übel mit ihm umgegangen wird, doch gleichwohl keine Dornhecke wird, sondern bleibt dennoch ein guter Baum: also ist Isaac auch wie ein guter Baum, der angepflanzet ist an den Wasserbächen, der seine Frucht reichlich und glücklichlich bringt, Ps. 1, 3. Denn dieweil er so geduldig ist, und das Unrecht, so ihm begegnet, hat ertragen können, öffnet er damit den Nachbarn die Augen, bewegt und erweicht ihre Herzen, daß sie ihre Sünde erkennen und bekennen, so sie an ihm begangen, und bedenken, wie große Gewalt und Unrecht sie ihm gethan haben, wie viel und wie unbillig er bisher von ihnen habe leiden müssen, ohne alle sein oder seines frommen ehrlichen Weibes und des ganzen Gesindes Schuld, in welchen allen sie zu jeder Zeit sonderliche Freundlichkeit, Mildigkeit und Gastfreiheit gespürt, dazu an seinen Knechten und Mägden auch gesehen, daß sie gute Sitten gehabt, und recht und wohl gelehrt und unterrichtet gewesen sind; derhalben erkennen sie, daß bei solchem ehrlichen und gottseligen Gesinde Gott mit seinem Segen gegenwärtig ist. Also bringt die Geduld und Freundlichkeit Isaaks seine Feinde zur Buße und zur Erkenntniß ihrer selbst, daß sie nun sich selbst strafen, und sich Schuld geben, daß sie den frommen heiligen Mann mit seinem Gesinde vertrieben hätten.

277. Und stimmt also dies sein überein mit dem Spruch Pauli Röm. 12, 20., den er aus den Sprüchen Salomos Cap. 25, V. 21. 22. genommen hat, da er sagt: „So nun deinen Feind hungert, so speise ihn; dürstet ihn, so tränke ihn: so wirst du feurige Kohlen auf sein Haupt sammeln.“ Denn wie sehr des Menschen Herz erbittert ist, so kann es doch erweicht und versöhnt werden, wo es in der Bosheit nicht gar verhärtet und verstockt ist, wie das Herz Pharaos war. Wo es aber aus menschlicher Schwachheit, oder irgend einer Unwissenheit und durch falsches Anbringen anderer Leute, oder auch sonst aus rechten Ursachen zum Zorn gereizt und bewogen ist, so kann es doch durch Wohlthat und Geduld also erweicht werden, daß es in sich selbst erbitzt wird und entbrennt, und die Güte und Freundlichkeit deß, den es für seinen Feind hält, erkennt, und sich selbst straft und also denkt: Warum habt du doch diesen frommen Mann verletzt und ihm Schaden gethan? Warum hast du ihn, der so unschuldig ist, verfolgt? Dies sind die Kohlen, die von solcher Geduld, Freundlichkeit und Wohlthat erhitzen und anfangen zu brennen.

278. Auf solche Weise bekehrt und überwindet Isaac den König auch, da er in Glauben, Hoffnung und Liebe und daß er gastfrei ist, beharrt, daß er sich mit ihm nicht allein versöhnt hat, sondern auch dadurch zum rechten Gottesdienst und zur wahren Erkenntniß Gottes bekehrt wird. Denn er hat ohne Zweifel die Predigten gehört, ist in die Hütte Isaaks gegangen, daß er von ihm Gottes Wort lernen möchte. Also hat Gott den König und viele andere Menschen mehr durch das Wort, so Isaac gelehrt hat, gewonnen und zur rechten Erkenntniß gebracht. Denn er hat uns sein Wort nicht allein um unsert-, sondern auch um Anderer willen gegeben, nämlich, daß wir es anderen Leuten mittheilen sollen.

279. Auf diese Weise ist der gute Geruch beides, der reinen gesunden Lehre Isaaks und auch seines unschuldigen Lebens, bei dem Könige und den Fürsten ausgebreitet worden, daß sie werden gesagt haben: Ach, warum haben wir diesen frommen heiligen Mann vertrieben, welcher recht und herrlich lehrt und auch dabei ein heilig Leben führt; deß Weib eine sehr fromme Matrone, dazu wohlthätig, milde und mit allerlei Tugenden geziert ist! Und also sind die Herzen der Menschen im ganzen Königreich und in der Nachbarschaft bekehrt worden. Darum sei geduldig, du seiest, wer du wollest, der du von andern Leuten Gewalt und Unrecht leiden mußt, und harre des HErrn. Denn dies wird alles von Gott geschickt und regiert, da Isaac selbst stille sitzt und wider den König gar nichts vornimmt, auch mit ihm keinen Vertrag macht; denn er leidet das nur mit Geduld und in der Stille, das ihm geboten war, er predigt, richtet eine Hütte auf und treibt das Wort fleißig.

280. Also sind wir ja jetzt auch stille, streiten nicht: es sind aber ihrer Viele zu uns gekommen auch von den Feinden und sind bekehrt worden, dazu ich sie weder gezwungen, noch jemals daran gedacht habe. Derhalben ist dies ein sehr lieblicher Trost, der sich fein reimt mit dem Spruch Salomos in seinem Prediger am 7. Cap. V. 9.: „Ein geduldiger Geist ist besser, denn ein hoher Geist.“ Darum lasset uns unsere Herzen auch daran gewöhnen, daß wir mögen geduldig sein, und dabei immer beharren, daß wir unser Amt ausrichten, die Laster strafen und das Wort immer fleißig treiben. Alsdann werden es die Nachbarn und Feinde hören und der gute Geruch Christi wird zu ihnen kommen, wird sie fangen und herzuziehen: wie allhier Abimelech von selbst zu Isaac kommt und wird ein heiliger König, wie auch sein Vater war, wiewohl er nicht beschnitten ist.

281. Denn die Beschneidung geht die Heiden nichts an, sondern gehört nur zum Hause Abrahams, zu den Nachkommen und zu seinem Gesinde: denselben, die im Volk Abrahams sind, ist sie nur geboten. Daß aber die Juden auch die Judengenossen (Proselyten) beschnitten haben, daran haben sie nicht recht gethan; denn sie haben davon kein Gebot oder Befehl gehabt. Denn es sind auch zur Zeit der Beschneidung viel Heiden selig geworden allein durch das Haus Abrahams, darin das Wort, die rechte Religion und Altar Gottes gewesen ist, da man das Wort im Namen des HErrn gehört hat, und da das Wort der Verheißung, „durch deinen Samen“, gelesen hat. Darum aber hat Gott die Beschneidung dem Hause Abrahams gegeben, auf daß gewisse Nachkommen und ein gewisser Ort sein möchte, da Gottes Wort und Christus der HErr auch von den Heiden möchte gefunden werden. Er hat aber 1 Mos. 22, 18. nicht gesagt: „Durch deinen Samen werden alle Völker beschnitten“, sondern „gesegnet werden“; darum haben sie den Segen ohne die Beschneidung empfangen können.

282. Derhalben ist diese ganze Historie eine Bestätigung des Trostes, den die Kirche und Gemeinde Gottes haben muß. Denn sie muß Schutz und Herberge haben in diesem Leben. Darum wird Abimelech wiederum durch das Wort dahin gebracht, daß er sich des vornehmsten und höchsten Amtes eines Königes annimmt, welches ist, daß er das Gesetz handhaben, Frieden erhalten, dazu auch Gerechtigkeit und Gericht üben und darüber halten soll; vornehmlich aber soll er über dem Gottesdienst halten, die Propheten und Lehrer schützen, Isaac nähren und unterhalten, auf daß in seinem Lande die Erkenntniß Gottes bleiben möge.

283. Das ist das allerhöchste und herrlichste Werk der Menschen, welches die Herzen mit rechter Freude erfüllet, dazu auch mit Gnade bei Gott und den Menschen; wie David im 60. Psalm V. 8. darüber auch fröhlich ist, da er sagt: „Ich bin froh, und will theilen Sichem, und abmessen das Thal Suchoth.“ Warum aber? Antwort: Darum, daß „Gott in seinem Heiligthum redet, deß bin ich froh“, sagt er. Das ist so viel gesagt: Wir hören in meinem Königreiche Gottes Wort; die Kirche und Gemeinde Gottes hat in meinem Lande ihre

\quellseite{71}
Hütte, Gottesdienst, Priester \&c. Darum ist das vornehmste und rechte königliche Amt eines Königes, daß er das Wort handhabe, und dafür sorge, daß dasselbe auch auf die Nachkommen erben und bei ihnen ausgebreitet werden möge.

284. Aber solcher Könige und Fürsten sind fast wenig. Denn der größte Theil der Könige und Fürsten in der Welt hassen und verfolgen das Wort allezeit, zerstören die Kirchen und plagen die Göttseligen. Derhalben wird uns dies zum Exempel und zum Trost vorgehalten, daß wir wissen sollen, daß Gott aus solcher großen Anzahl der Fürsten, die da Feinde des Wortes sind, den Einen oder Andern erwählt, die sich der Kirche annehmen, das Wort Gottes lieb haben und lernen. Also wird Abimelech im Heiligen Geiste getauft und mit geistlicher Beschneidung beschnitten, und glaubt an den Gott, von dem gepredigt und der geehrt wird im Hause Isaaks; wird ein Schüler, ein Freund und Schirmherr Isaaks. Denn das ist der Segen der Göttseligen, welchen Gott verheißen hat; wie im 1. Buche Samuelis steht im 2. Cap. V. 30.: „Wer mich ehret, den will ich auch ehren“ \&c.

285. Die Hebräer disputiren von dem Namen achusath, ob es ein Eigenname eines Mannes oder aber ein Gattungswort sei? Hieronymus hält dafür, daß es ein Gattungswort sei, ein collectivum, wie man es in der Grammatik nennt, und daß es nicht nur Einen Menschen, sondern einen ganzen Haufen Volks guter Freunde bedeuten soll, die um den König her gestanden und mit ihm gekommen sind. Das hebräische Wort achas heißt ergreifen; wie gute Freunde unter einander pflegen die Hände zu geben, einander dabei halten und einander Handreichung und Hülfe leisten; daß dies der Verstand und die Meinung sei: Abimelech ist gekommen, und ein ganzer Haufe seiner Freunde, die einander bei den Händen gehalten, wie ein Haufen Volks oder eine ganze Menge, das ist, Abimelech ist zu Isaak mit der ganzen Gesellschaft seiner guten Freunde, so entweder im Lande oder aber außerhalb des Landes gewesen, gekommen. Denn es werden ohne Zweifel viele von den Nachbarn umher, die nicht unter Abimelech in seinem Lande gesessen, sondern etwa ihre Verwandte oder sonst gute Freunde gewesen sind, zum Hause Isaaks gekommen sein. Zu denselbigen allen ist der gute liebliche Geruch Christi gekommen, daß sie sich unter einander ermahnt haben und gesagt: Lieber, lasset uns hingehen zu dem heiligen Manne, lasset uns seine Predigt hören \&c.

286. Da mag man aber sehen, was ein Prophet für eine nützliche Person in der Welt sei; wie die Worte Christi Joh. 15, 16. anzeigen, da er zu seinen Jüngern sagt: „Ich habe euch erwählet, und gesetzt, daß ihr hingehet, und Frucht bringet, und eure Frucht bleibe.“ Denn hier kommen die großen Herren und Fürsten im selbigen Königreiche zu Isaak mit ihren Blutsverwandten, Schwägern und Freunden. Ein solch groß Ansehen hat er bei ihnen, und also wird die Geduld des heiligen Mannes gekrönt und belohnt, dieweil er das Böse mit Gutem überwunden, nach der Lehre Pauli Röm. 12, 21.

\begin{quote}
\textbf{V. 27—31. Aber Isaak sprach zu ihnen: Warum kommt ihr zu mir? hasset ihr mich doch, und habt mich von euch getrieben. Sie sprachen: Wir sehen mit sehenden Augen, daß der HErr mit dir ist. Darum sprachen wir: Es soll ein Eid zwischen uns und dir sein, und wollen einen Bund mit dir machen, daß du uns keinen Schaden thust, gleichwie wir dich nicht angetastet haben, und wie wir dir nichts, denn alles Gute gethan haben, und dich mit Frieden ziehen lassen. Du aber bist nun der Gesegnete des HErrn. Da machte er ihnen ein Mahl, und sie aßen und tranken. Und des Morgens frühe stunden sie auf, und schwur einer dem andern; und Isaak ließ sie gehen, und sie zogen von ihm mit Frieden.}
\end{quote}

287. Isaak straft sie und wirft ihnen ihre Sünde vor, damit sie ihn und sein Gesinde hart verletzt und erzürnt hatten. Aus dieser Strafe mag man eine solche Lehre, so zu guten Sitten dient, nehmen, nämlich, daß man also im Glauben, in der Göttlichkeit und Geduld leben soll, daß wir gleichwohl den großen Haufen nicht schmeicheln und ihre Sünde nicht zu decken und verhehlen. Denn wo das geschieht, so beschweren und beladen wir uns selbst mit fremden Sünden. Die Alten soll man allerdings nicht mit bittern Worten strafen, sondern sie ermahnen, auf daß der Jüngere wider den Alten nicht stolz sei, wie Jes. 3, 5. steht. Und zu unsern Zeiten schleichen eben solche böse verderbte Sitten allgemach herein. Es soll ein Unterschied zwischen den Obersten und Untersten, zwischen Eltern und Kindern sein. Der Vater ist der Gewalt und Noth halben mehr und größer als der Sohn; darum soll sich der Sohn mit nichten gegen den Vater halten, wie sich der Vater gegen den Sohn hält; und Gottes Wort sagt 3 Mos. 19, 32: „Vor einem grauen Haupt sollst du aufstehen, und die Alten ehren.“

288. Also verhehlt Isaak die Sünde des Königs und seiner Freunde nicht, er heuchelt ihnen nicht, sondern straft sie mit gebührlicher Ehrerbietung und sagt: „Warum kommet ihr zu mir, so ihr mich doch hasset?“ Wißt ihr nicht, daß ihr gesündigt habt? Denn warum habt ihr mich von euch getrieben? Ich will mich zwar nicht weigern, mit euch Gespräche zu halten; ich weiß aber, daß ihr voll Hasses und Neides seid.

289. Auf solche Weise soll man die Sünde nicht rechtfertigen, man soll auch den nicht loben, der gesündigt hat; wo es aber irgend ein Alter ist, der eine Sünde begangen hat, soll man ihn freundlich ermahnen. Als, wenn der Sohn seine Eltern um irgend einer Sünde willen straft, soll er sagen: Mein lieber Vater, meine liebe Mutter, dies läßt sich ansehen, als daß frommen ehrlichen Leuten nicht wohl anstehe. Einem alten Mann steht es zumal übel an, daß er Hurerei, Ehebruch oder andere dergleichen Sünden und Schande treibt. Einem Propheten aber oder Prediger gebührt, daß er bisweilen etwas härter und schärfer strafe, wie allhier an dem Exempel Isaaks zu sehen ist.

290. Dieweil aber der König und seine Freunde sehen, daß die Strafe billig und recht ist, bekennen sie ihre Sünde und nennen ihn einen Gesegneten des HErrn. Als wollten sie sagen: Wir erkennen, daß wir unrecht gethan, dieweil wir dich unschuldiger Weise gehaßt haben; denn wir haben nicht gewußt, daß du ein Gesegneter des HErrn wärest. Wir haben zwar wohl gesehen, daß du gewachsen und reich geworden bist; wir haben aber nicht gedacht, daß man solches dem göttlichen Segen müßte zuschreiben: nun aber sind uns die Augen, dazu auch unsere Herzen und Ohren aufgethan, und erkennen gerne, daß wir dich unbilliger Weise beleidigt und vertrieben haben.

291. Sie entschuldigen und leugnen zwar die Sünde nicht, sondern lindern sie etwas; bekennen aber doch genugsam, daß sie Schuld haben, da sie sagen: „Wir sehen mit sehenden Augen, daß der HErr mit dir ist.“ Also begreifen sie zugleich miteinander und fassen zusammen das Bekenntniß ihrer Sünde und die Erkenntniß der göttlichen Güte und Gnade an Isaak; denn er ist noch der Meinung gewesen, daß sie mit ihm noch nicht versöhnet, sondern dermalen noch Feinde wären. Darum sagt Abimelech: Wir sind dir nun nicht übel geneigt, sintemal wir nun mit sehenden Augen sehen, daß Gott mit dir ist. Und wir wollen auch zugleich mit dir sein in der Gemeinschaft einer Kirche und eines Glaubens, wollen deinen Gott erkennen und ehren: darum soll ein Eid zwischen dir und uns sein. Denn wiewohl wir nicht beschnitten sind, so wollen wir dennoch den rechten Glauben von dir lernen und wollen des zeitlichen Friedens und Gemeinschaft theilhaftig sein; sei du nur unser Freund und thue uns keinen Schaden, wir wollen dich nimmermehr beleidigen; wie wir dir bisher ja nichts denn Gutes gethan haben.

292. Also entschuldigt sich Abimelech so viel, als er kann. Er hat zwar Isaak sein Gut nicht genommen, hat ihm nichts mit Gewalt geraubt, so auch, daß diese Verkleinerung seiner Schuld billig und recht zu sein scheint; es war aber doch durch die Diener sein Herz also erbittert, daß er Isaak gebot, aus dem Lande zu ziehen. Sonst, sagt Abimelech, dieweil du in unserm Lande fremd wärest, haben wir dich freundlich und mit gutem Frieden ziehen lassen, haben dir gar kein Leid gethan, und so dir die Hirten die Wasserbrunnen haben genommen, so ist dir doch solch Unrecht und Schaden durch den göttlichen Segen reichlich erstattet worden. Darum bitten wir, du wollest uns verzeihen, was geschehen ist.

293. Dies ist die Demuth und Bekehrung des Königes Abimelech: dagegen zankt Isaak auch nicht, rechnet nicht mit ihm und streitet nicht über irgend einer Sache, sondern freut sich von wegen des guten Willens des Königes und seiner Freunde und richtet ein Mahl an, auf daß der Eid, der Friede und Schutz zwischen ihnen beständig und gewiß sein möge. Und also, da Isaak ist geprüft worden, wird er nun erhöht über die, so ihn zuvor gehaßt und verfolgt hatten.
\quellseite{72}
\subsection{II}

\begin{quote}
\textbf{V. 32. 33. Desjenigen Tages kamen Isaaks Knechte, und sagten ihm um den Brunnen, den sie gegraben hatten, und sprachen zu ihm: Wir haben Wasser funden. Und er nannte ihn Saba; daher heißt die Stadt Bersaba, bis auf den heutigen Tag.}
\end{quote}

294. Alhier kommt noch ein anderer Segen. Nachdem er mit Abimelech Frieden gemacht, gibt ihm GOtt einen neuen Brunnen. Durch welch Exempel wir ermahnt werden, daß man auf den Trost harren soll, der nach der Trübsal kommen wird. Denn diese Veränderung währet in dem Leben der Heiligen für und für, nämlich, daß nach der Trübsal Trost folgt und wiederum nach dem Troste eine andere Trübsal; so lange aber GOtt lebt, ist es gewiß, daß wir auch leben werden: „Denn unser keiner lebt ihm selber, und keiner stirbt ihm selber. Leben wir“, sagt St. Paulus Röm. 14, 7. 8., „so leben wir dem HErrn; sterben wir, so sterben wir dem HErrn.“ Es sei gleich Tag oder Nacht, Kreuz oder Freude, Alles ist GOttes: niemand ist ihm selber krank, sondern dem HErrn; wer da schläft, der schläft dem HErrn; wer da ist, der ist dem HErrn; es ist Alles gut, wenn du GOttes Wort hörst.

295. Von dem hebräischen Wort scheba ist droben im 21. Capitel gesagt worden, nämlich, daß es zweierlei Bedeutung hat, als, von der Zahl, die sieben bedeutet, und vom Eide; und das lateinische Wort septem läßt sich ansehen, als sei es von dem hebräischen hergekommen; wie wir Deutschen dasselbe Wort auch aussprechen, sieben. Hieronymus will haben, daß man es an diesem Orte auslegen soll, daß es so viel heiße, als saturitas, Fülle, denn dasselbe heißt scheba auch; daß es hier so viel heißen soll, als wenn man sagt, puteus saturitatis, ein Brunnen der Fülle oder ein voller Brunnen. Aber die Rabbinen der Hebräer halten hart darüber, daß man es nach seiner rechten eigentlichen Bedeutung verdolmetschen soll, daß es heiße, puteus juramenti (Brunnen des Eides). Und dasselbe gefällt mir auch besser, daß es ein Brunnen des Eides genannt werde und nicht der Fülle. Wie er droben ist genannt worden ein Schwörborn; welchen Namen man zu Erfurt auch einem Brunnen gegeben hat.

296. Zuletzt sollte man auch an diesem Orte die heimlichen Deutungen der Brunnen suchen. Ich habe aber dazu nicht so große Lust wie Origenes und Hieronymus. Ich frage nichts darnach, denn allein sofern sie den Verstand der Historie zieren, welchen man aus der einfältigen Historie nehmen kann. Und daselbst sind die Allegorien oder heimlichen Deutungen wie Blumen, so unter den einfältigen Verstand der Historien gestreuet werden; man kann aber damit nichts beweisen, welches Augustinus von den Figuren auch gesagt hat.

297. Wo aber ja jemand eine heimliche Deutung haben wollte, der mag die drei Brunnen ziehen auf die drei Häupstücke der heiligen Schrift, nämlich, auf das Gesetz, auf die Propheten und auf das Evangelium. Denn das Gesetz ist Esek, das ist, wie man in der deutschen Sprache sagt, es ist Eßig: ja, in der Seele oder Gewissen ist es Eßig; denn „das Gesetz richtet nur Zorn an“, Röm. 4, 15. Die Propheten sind sitna; denn wo sie auf das Gesetz bringen, so ist da der Zorn und Feindschaft GOttes; nicht daß solches vom Gesetz herkommen sollte, sondern um der Philister willen; aus dem Guten wird Böses. Das Evangelium aber ist rechoboth: wo dasselbe in der ganzen Welt ausgebreitet wird, macht es die Gewissen weit und bringt ihnen Trost. Dies mag die heimliche Deutung sein dieses Textes. Wer nun will, der mag dieselbige gebrauchen.

\subsection{III}

\begin{quote}
\textbf{V. 34. 35. Da Esau vierzig Jahr alt war, nahm er zum Weibe Judith, die Tochter Beri, des Hethiters; und Basemath, die Tochter Elon, des Hethiters. Die machten beide Isaac und Rebekka eitel Herzeleid.}
\end{quote}

298. Dies sind Exempel der heiligen Patriarchen, an welchen wir nicht allein sehen, wie sie so einen trefflichen Glauben gehabt, sondern auch, wie sie sich so höchlich befliessen haben, die Werke der Liebe, Geduld und allerlei Tugenden zu üben. Bisher haben wir gesehen, wie viel schwere Trübsal Isaac erlitten habe, da er mit Anfechtung, so ihm von den Fremden begegnet, geplagt ist; wie er unter den Heiden als ein Fremdling umhergezogen, welche Heiden, wo sie nicht öffentlich seine Feinde gewesen sind, so haben sie ihm doch mit heimlicher Feindschaft, mit bitterem Haß, und daß sie ihm heimlich nachgestellt, viel Mühe und Arbeit gemacht. Und das sind noch äußerliche Anfechtungen gewesen.

299. Hier aber folget nun der Anfang eines viel greulicheren Unglückes, nachdem er alle Anfechtungen überwunden hat, damit er sammt seinem Weibe, Kindern und Gesinde sehr hart geplagt worden, davon er schon vollkommene Geduld gelernt hat. Und sind die Heiden, unter welchen er gewohnt, auch nun mit ihm ausgesöhnt, und die Kirche oder Gemeinde GOttes ist in gutem Frieden und Ruhe gewachsen und hat zugenommen, also daß immer mehr und mehr Zuhörer von alenthalben herzugelaufen sind, über das ist er auch mit leiblichem Segen begnadigt. Denn GOtt gibt ihm Früchte der Erde vollauf, dazu seine lustige Brunnen und allerlei Nothdurft. In Summa, er hat Frieden und Ruhe und Alles vollauf.

300. Nun aber fangen die Anfechtungen in seinem Hause an, und die viel schwerer sind als die vorigen. In Summa, nun kommt eitel Jungfer, Pfeffer, Kellerhals, Jammer und Herzeleid, da er nun alt ist, und hofft, die Anfechtungen sollen ein Ende nehmen, dieweil er von so vielem Jammer und Elende, so er erlitten hat, ermüdet und erschöpft ist. Denn er ist nun schier ein Hundertjähriger; welches man daraus abnehmen kann, weil Esau und Jakob geboren sind, da Isaac im sechzigsten Jahr seines Alters gewesen, und Esau, da er vierzig Jahre alt war, ein Weib nimmt: darnach, sage ich, ist Isaac hundert Jahre alt. Derselbe alte Mann wird nun aufs neue mit schwerer Anfechtung in seinem eigenen Hause geplagt und gezüchtigt; da es sich ansehen läßt, daß nun alles Unglück vorüber und überwunden sei, da hebt sich erst ein neu Kreuz und Angst an.

301. Dieweil aber uns eben dasselbe auch pflegt zu begegnen, sollen wir bedenken, daß dies uns zum Exempel geschrieben ist, auf daß wir sehen, wie GOtt seine Heiligen, denen er das Wort gegeben und das ewige Leben verheißen hat, auf so mancherlei Weise und so hart pflege zu versuchen in diesem elenden Leben, daß der Spruch des heiligen Paulus im Buch von der Apostel Geschichten am 14. Cap. V. 22. wahr sei, da er sagt, „daß wir durch viel Trübsal müssen in das Reich GOttes gehen.“

302. Denn also leidet Isaak in seinem Hause viel schweres Unglück vom hundertsten Jahre seines Alters bis auf das achtzigste über hundert. Das hat er alles leiden und in sich fressen müssen bis an die Zeit, da Joseph, seines Sohnes Jakobs Kind, nach Egypten ist vertrieben worden. Denn in demselben Jahre oder um die Zeit ist er gestorben, und hat denselbigen Jammer alle gesehen und erlitten, der hernach folgen wird an Jakob, seinem Sohne, an Joseph, Jakobs Sohne, an Esau, der abtrünnig geworden ist, an Dina, so beschlafen worden ist, an Ruben, der auf seines Vaters Jacob Bett gestiegen. Dieselbigen vielen, großen und schweren Anfechtungen, damit er diese achtzig Jahre ist geplagt worden, haben ihn endlich müssen aufreiben.

303. Derhalben sind diese Patriarchen rechte Heilige, und wenn wir mit ihnen sollten verglichen werden, so sind wir gar nichts; ja, es haben sich alle Bischöfe, Märtyrer und Apostel ihrer hoch verwundert und gedacht: Wir sind zwar in dem ersten Grade und müssen viel leiden, aber das alles ist nichts gegen die Anfechtungen der Patriarchen zu rechnen; und darum haben sie dieselbigen auch also geehrt, daß sie sich ihrer selbst, mit den Patriarchen verglichen, geschämt haben. Denn es ist einem Märtyrer um eine oder eine halbe Stunde zu thun, hier aber ist der Anfechtung, der Mühseligkeit und täglichen Jammers kein Ende. Denn, lieber GOtt, wie ist das so ein großer Jammer und Elend, daß er ganz achtzig Jahre in seinem Hause mit höchsten Schmerzen solch Unglück hat müssen ansehen und leiden. Und ist das kein gemeiner Jammer und Unglück, sondern ein sonderlicher und der höchste Jammer und Elend gewesen.

304. Es ist aber niemand so beredt, der die Geduld Isaaks genugsam könnte ausreden, wie die Geduld Abrahams und Jakobs auch hoch gerühmt wird. Daraus genugsam zu sehen ist, daß diese drei Patriarchen GOtt vornehmlich lieb gewesen, und darum sie auch des Namens der Patriarchen und der allerheiligsten Märtyrer wohl werth sind. Und wir möchten sie wohl billig unter allen andern heiligen eine Dreifaltigkeit der Heiligen (wenn man also reden möchte) nennen; denn sie haben sich weder die Größe, noch die Menge der Anfechtungen, dazu auch, daß sie für und für gewähret, überwinden lassen, sondern haben mit großem starkem Muthe diese Anfechtungen alle männlich und ritterlich überwunden und darüber triumphirt; und wir werden endlich am jüngsten Tage die Ehre solcher Ueberwindung sehen, und uns derselben verwundern; jetzt mögen wir sie etlichermaßen bedenken und anschauen.
\quellseite{73}

305. Vor den Papisten aber ist dies alles gar verborgen um des einigen Aergernisses willen, daß sie Scheinmänner gewesen, Weib und Kind gehabt etc., keinen ehelosen Stand geführt haben. Denn das ist im Papstthum allein die Heiligkeit gewesen, von Einer in ehelosen Stande gelebt hat; darum können sie diese schönen Lichter, ja, die Sonne und Mond der allerschönsten Tugenden an den heiligen Patriarchen nicht sehen.

306. Derhalben werden wir dadurch gelehrt und unterwiesen, nicht allein, daß wir solchen Exempeln der Patriarchen sollen nachfolgen, sondern daß wir auch schamroth werden und uns selbst schämen lernen, daß wir, mit diesen Patriarchen verglichen, so gar nichts sind. Denn wir müssen uns hier schämen und sagen: Augustinus und Ambrosius sind nichts, welche zwar ihre Gaben gehabt haben, aber mit diesen heiligen Vätern nicht einmal zu vergleichen sind, welche Gott allenthalben mit so großer Geduld und Demuth gehorsam gewesen. Darum spricht Gott nicht vergeblich, er sei ein Gott Abrahams, Isaaks und Jacobs, 2 Mos. 3, 6.

307. So läßt uns nun die Anfechtung, so Isaak in seinem Hause gehabt, überlaufen. „Da Esau vierzig Jahr alt war“, sagt der Text, „nahm er zwei Weiber, die beide sehr bitter gegen Isaak und Rebekka waren.“ Dies sehr schwere Kreuz hat Mose mit kurzen Worten beschrieben. Er zeigt an das Alter Esau’s; denn zu der Zeit haben sie etwas eher Weiber genommen als vor der Sintfluth: nach der Sintfluth ist es gebräuchlich und auch geboten gewesen, daß sie haben Weiber genommen, wenn sie an ihre vierzig Jahre gekommen sind.

308. Danach schreibt auch Mose, daß Esau zwei Weiber genommen, die fremde oder Heidinnen gewesen sind; daraus wir abnehmen können, daß er nicht habe nach Syrien ziehen und eine von seinen Blutsverwandten nehmen wollen, das vor ihm sein Vater Isaak gethan und hernach Jakob auch thun wird.

309. Es läßt sich ansehen, als habe Mose wollen anzeigen, daß Esau mit den Männern und Jünglingen, den Hethitern, einen Bund gemacht, und auf daß derselbe möchte kräftiger sein, so habe er derhalben ihre Töchter oder Schwestern zur Ehe genommen, so doch die Heiden vorlängst von Gott verdammt gewesen; wie er droben diese selben Völker aufzählt und sie verdammt hat, da er sagt im 15. Capitel dieses ersten Buchs V. 16.: „Die Mißethat der Amoriter ist noch nicht alle“, welches ohne Zweifel Esau gewußt hat. Nicht daß sie Gott gänzlich verworfen hätte; denn Juda, einer von den zwölf Patriarchen, ist ein Vater Christi, von des Samen Christi gekommen ist; dieser hat eine Cananitin zum Weibe genommen, die Tochter des Suach. Und Thamar hat durch Hurerei und Blutschande Perez und Serah gezeugt. Derhalben hat Gott diese Heiden verdammt, aber doch gleichwohl aus dem Geschlechte und Geblüte solcher verdammten Völker etliche Mütter erwählt.

310. Und kann wohl sein, daß es Isaak anfänglich nicht so sehr verdrossen hat, daß Esau Hethiter zu Weibern genommen hat und derselbigen zwar zwei zugleich. Er hat vielleicht gesehen, daß dies bei den selbigen Völkern ein gemeiner Gebrauch gewesen ist. Wie Abimelech auch viel Weiber gehabt und im ganzen Morgenlande eben derselbige Gebrauch auch gewesen ist. Darum halte ich dafür, daß der zwei Ursachen keine das Herz Isaaks und Rebekkas erbittert und sie betrübt gemacht habe. Aber daher hat sich das Kreuz und die Trübsal dieser sehr heiligen Eltern erhoben, daß sie diese zwei Weiber ihres Sohnes in ihr Haus haben führen müssen, von welchen darnach aller Lärm und allerlei Verletzung und Widerwille entstanden ist.

311. Die Mönche im ehelosen Stande wissen von dem Kreuz und Mühseligkeit im Hausregiment gar nichts. Es ist aber wahrlich beiden Eltern ein sehr hart und schwer Kreuz gewesen, vornehmlich aber der Mutter, welche bis an ihr neunzigstes Jahr mit großer schwerer Mühe und Arbeit das Haus allenthalben regiert hatte. Nun aber, da sie Alters und vieler Sorgen halben, so sie über der Haushaltung gehabt, fast entkräftet ist, muß sie noch die jungen Weiber in ihr eigen Haus nehmen, dieselbigen von dem Ihren nähren und unterhalten, daß sie die alte Schwiegermutter nur wohl martern und plagen, und nach ihrem eigenen Muthwillen Alles im Hause verwirren und das unterste zu oberst kehren. Denn also sagt Mose: „Sie waren beide sehr bitter, und betrübten Isaak und Rebekka.“

312. Welcher Art aber die Anfechtungen gewesen seien, deß gedenkt die Schrift nicht ausdrücklich; und ich halte dafür, daß die Eltern nicht gleich um irgend einer geringen Schuld willen erzürnt und erbittert worden sind; wie im Hausregiment viel geringe Verletzungen sich zu begeben pflegen. Denn beide, Isaak und Rebekka, sind fast wohl geübt gewesen, Unglück und Trübsal zu leiden, und haben bei sich die Ueberwindung und den Sieg des Geistes gehabt. Solche Leute lassen sich ja nicht leichtlich erzürnen.

313. Derhalben wird es irgend eine schwere und besondere Ursache müssen gewesen sein. Dieselbige können wir aber einigermaßen wohl erforschen und verstehen lernen, wenn wir die Zehn Gebote werden fleißig ansehen; dieselbigen werden wohl anzeigen, welches die Ursachen gewesen seien dieser Anfechtungen. Darnach läßt uns hierher ziehen das Exempel Hagars, welches droben, 1 Mos. 16, 4. ist erzählt worden. Denn dieweil sich dieselbige wider ihre Frau aufzulehnen und sie zu betrüben gewagt hat, was, meinest du, sollten wohl diese Weiber nicht gethan haben? Denn die werden sich ohne Zweifel gerühmt und gesagt haben: Wir sind des erstgebornen Sohnes Weiber, die Erstgeburt ist unser. Denn Jakob ist gar verachtet und verworfen gewesen, Esau aber hat sich gehalten wie ein Priester und wie einem Fürsten und rechten Erben in der Kirche und im Hause zusteht. Derhalben haben auch die fremden Weiber nicht wollen Mägde, sondern selbst Hausmütter sein, und solche Ehre im Hause haben und erhalten als Weiber des erstgebornen Sohnes.

314. Diese Ehre und Herrlichkeit hat sie sehr stolz und hoffärtig gemacht, daher es gekommen ist, daß Isaak und Rebekka allerlei Unbilligkeit haben leiden müssen. Und wie sollten sie sich anders gehalten haben, dieweil droben gesagt ist, daß Esau selbst auch der Art gewesen, daß er sich also gerühmt und so wilde Sitten gehabt hat? Wie, meinest du, wird er mit so großem Stolz bei seinen Weibern gerühmt und gesagt haben: Ich bin der Erstgeborne; das Erbe gehört mir zu; ich habe die göttlichen Verheißungen, so meinem Vater geschehen sind; die gehören alle aus göttlichem Recht mir zu; ich werde der Vater sein des künftigen Samens; ich werde das Land besitzen. Und so oft Isaak vor seinem Gesinde von den geistlichen und leiblichen Verheißungen gepredigt hat, so oft hat Esau sich dieselbigen zugeschrieben: Ich bin der Erstgeborne und der rechte Erbe dieser Verheißung, und ihr seid Weiber des Erstgebornen; ihr seid die rechten Hausmütter und Frauen über das ganze Land; unsere Kinder werden dasselbe erben und besitzen; Jakob ist viel geringer denn ich, ist verachtet und verworfen.

315. Das ist die Eine Sünde und die Eine Bitterkeit gewesen, nämlich, die große Hoffahrt der zwei Weiber, Judith und Basemath; welche Hoffahrt diese zwei heiligen Eheleute haben leiden müssen. Denn also wird diese Hoffahrt Esau’s und seiner Weiber recht verglichen mit dem vorigen Exempel Ismaels in Abrahams Hause, da sich Ismael auch dessen gerühmt hat, daß er der Erstgeborne wäre, wie sein Vater selbst bezeugt hatte. Da ist Sarah zornig geworden, daß er wollte der Erbe sein und ihren Sohn Isaak verstoßen. Nein, sagte sie, also hat Gott nicht geredet, sondern: „In Isaak soll dir der Same genannt werden.“ Wiewohl nun Ismael nach dem Fleisch der Erstgeborne ist, so ist er doch nach dem Geist der Erstgeborne nicht.

316. Eben dasselbe hat sich in diesem Hause auch also zugetragen. Esau und seine Weiber haben sich des Priesterthums und Regimentes gerühmt, und daß Jakob auf nichts mehr hoffen oder warten sollte, denn nur, daß ihm ein klein Theil vom Erbe oder aber irgend eine geringe Gabe werden möchte, daran er sich sollte genügen lassen, und würde also seines Vaters Haus verlassen und in das Elend ziehen müssen; Esau aber würde im Hause bleiben, Factotum sein, würde dazu der Fürst und Priester, und Vater Christi sein. Wie aber Sarah über das Rühmen Ismaels ungeduldig und zornig geworden ist, also ist Rebekka auch heftig zornig geworden, und hat auf Mittel und Wege gedacht, wie Jakob die Erstgeburt erhalten möchte. Und dasselbe hat sie auch endlich erlangt.

317. Es ist aber noch fast dreißig Jahre verzogen worden. Denn da dieselben vorüber gewesen, ist Jakob erst an seines Vaters Statt gekommen und der Erbe geworden beide in leiblichen und auch in geistlichen Dingen, und bleibt in seines Vaters Hause; Esau aber hat dasselbe mit Weib und Kindern verlassen. Das ist also unseres Herrn GOttes Gericht und seine sonderliche Regierung, welcher, wie die Jungfrau Maria Luc. 1, 52. singt, die Gewaltigen vom Stuhl stößet und die Elenden erhebet: Denn er will und kann die Hoffahrt nicht leiden, ja, er verstößt und stürzet der Hoffahrt halben diejenigen, welche er kurz zuvor auf den Stuhl gesetzt und hoch erhoben hat; denn er thut beides, er richtet die Elenden auf und setzt sie auf den Stuhl, und stößt ebendieselbigen hernach wieder vom Stuhle.
\quellseite{74}

318. Warum thut aber GOtt das und warum ist er wider sich selbst? Antwort: Die er deswegen, daß sie elend sind, erhöht, dieselben stürzt er hernach wieder, wenn sie hoffärtig werden. Und ist kein Kaiser so mächtig und gewaltig, den er nicht demüthigen und vom Stuhle stoßen könnte; ja, es ist ihm viel leichter, irgend einen großen Monarchen zu stürzen, denn mir, eine Fliege zu tödten. Denn also sagt er zu Saul 1 Sam. 15, 17.: „Ists nicht also, da du klein warest vor deinen Augen, wurdest du das Haupt unter den Stämmen Israel?“ Du bist aber nun stolz und hoffärtig, und willst meiner Stimme nicht gehorchen, darum stoße ich dich vom Stuhle.

319. Darum, sage ich, ist das GOttes Werk, daß er die Elenden erhöhet, und wiederum die Hoffärtigen stürzt. Dasselbe bezeugen auch die Historien bei allen Völkern. „Er setzt Könige ab, und setzt Könige ein“, spricht Daniel im 2. Cap. V. 21. Dieweil die Könige demüthig sind, schützt er sie und erhält sie bei ihrem Königreiche; wenn sie aber stolz werden und ihm seine Heiligen oder Gläubigen verfolgen, so setzt er sie ab.

320. Es ist wahrlich das Exempel mit Ismael 1 Mos. 21, 9. ff. fast erschrecklich, welcher um seiner Hoffahrt willen vertrieben und der Erstgeburt beraubt wird, vornehmlich soviel die geistliche Verheißung belanget. Also wird Esau auch stolz und hoffärtig, darum wird er verworfen, und ist nicht werth, daß er soll ein Erbe sein des Landes Canaan und ein Vater des Messias, der in dem Lande regieren sollte. Derhalben ist die Erbitterung Isaaks und der Rebekka aus der ersten und andern Tafel gekommen. Aus der ersten um des Priesterthums willen, aus der andern aber um des weltlichen Regiments willen. Denn die zwei Weiber Esaus haben gerühmt und gesagt: Unser Mann wird den Gottesdienst und die Gemeinde GOttes haben und Christus wird von unserm Manne geboren werden. In der ersten Tafel soll er das Lehramt führen, in der andern soll er auch das Regiment verwalten. Das ist das Recht des Erstgebornen, daß er sich mit großem Stolze gerühmt hat.

321. Droben aber (Cap. 25. V. 23.) ist in der Verheißung das Gegentheil gesetzt und beschlossen worden, da also gesagt wird: „Der Größere wird dem Kleinern dienen.“ Denn GOtt setzt seine Verheißung also, wenn er Gnade und Segen verheißt, daß er GOtt bleibe. Wenn die Verheißung fest sein und bleiben soll, so ist vonnöthen, daß GOtt selbst auch bleibe und erkannt werde, als der den Segen gibt und verheißt. Ist es aber der Fall, daß ich stolz und hoffärtig werde, so verliere ich die Verheißung. Denn die Verheißung schließet die Hoffahrt aus, und begreift in sich die Demuth und die Erkenntniß GOttes, der die Verheißung thut.

322. Es haben aber Rebekka und Jsaak diese Erbitterung tragen und dulden müssen bis in das dreißigste Jahr. Denn so lange hat Jakob ohne Weib gelebt, und ist auch vom Gesinde im Hause gegen Esau verachtet und verworfen gewesen, welches seiner Mutter großen und gerechten Schmerz gemacht hat. Denn das man gemeiniglich pflegt zu sagen, daß alle Schwiegermütter ihrer Söhne Weiber pflegen zu hassen, soll man von dieser heiligen Matrone Rebekka nicht sagen. Es ist hier viel ein ander Ding; denn sie bekümmert sich um die Erstgeburt Jakobs, wie Sarah auch that. Da muß man nichts verschweigen, da muß man kein Unrecht leiden, und man muß daselbst von seinem Rechte nicht weichen oder davon abstehen, weil es eine göttliche Verheißung ist. Wo man aber das Wort verachten oder verwerfen, und die göttliche Verheißung verleugnen will, das ist keine Geduld, sondern ist eine Faulheit, ja, eine Verachtung GOttes. Denn ich könnte alles leiden, was ich sollte, allein, daß man mir meinen Herrn Christum nicht nehme.

323. Ja, sagst du, es soll ein Christ stille schweigen und Geduld haben mit aller Demuth, soll dazu verzeihen und vergeben denen, so ihn beleidigen. Antwort: Alsdann hört aber die Geduld, Demuth und alle andern Werke der Liebe auf, wenn ich den verlieren soll, um deß willen ich lebe: wo man GOtt verlieren, sein Wort und den rechten Gottesdienst verleugnen soll, da muß man keine Geduld haben; da müssen wir der göttlichen Verheißung gewiß sein, die uns angeht, daß wir uns dieselbe auf keinerlei Weise nehmen lassen.

324. Also, da Rebekka sieht, daß Esau sich mit seinen Weibern hoch erhebt, als sei er der König und sie die Königinnen, ja, daß sie die rechten Erben seien beide zum weltlichen und geistlichen Seegen, da hört die Geduld auf, und denkt Rebekka also: Ich werde daraufhin arbeiten müssen, daß meinem Sohne Jakob der Segen nicht genommen und entzogen werde, sondern daß er möge der Erbe sein nach der Verheißung. Und der Heilige Geist gibt der Rebekka den Rath, davon wir hernach hören werden, daß Jsaak, der Vater, da er nicht wohl sehen konnte, unwißentlich die Erstgeburt von Esau auf Jakob gewendet hat. Und wie Abraham die Stimme der Sarah hat müssen hören, also muß Jsaak auch wider seinen Willen der Rebekka gehorchen. Und daselbst richtet Rebekka aus durch einen heilsamen und gottseeligen Betrug, auf daß das Erbe von dem falschen unrechten Erstgebornen genommen, und Jsaak, dem rechten Erstgebornen, wiederum zugestellt werde.

325. Dieses ist die eine Ursache der Verbitterung; und dieselbige ist wahrlich größer gewesen, denn wir denken können, die wir dies Ding schläfrig und unfleißig lesen und auch diesen Kampf nicht verstehen. Zu diesem Jammer und Unglück ist das andere auch gekommen, nämlich, die Abgötterei. Denn diese zwei Weiber ihres Sohnes, welche in der göttlichen Lehre nicht unterrichtet waren, haben dem rechten Gottesdienst und der wahren Religion widersprochen. Vielleicht mögen sie auch wohl die Satzung ihrer Väter, der Heiden, in das Haus ihres Schwiegervaters gebracht, und mit demselbigen Sauerteige die Kirche und Gemeinde GOttes, so Jsaak gehabt, vergiftet haben. Ja, das hat erst dieser heiligen Seelen Leute Herzen erbittert, da die Weiber an dem Aberglauben ihrer Väter fest gehalten und denselbigen vertheidigt haben, daß davon Kezerei und Secten entstanden sind.

326. Denn das ist das Allergrößte und Betrübteste: wenn die Kirche wohl angerichtet und bestellt ist, und die Lehrer und Gläubigen einig sind, reden aus Einem Munde, haben Ein Herz, Eine Feder, wenn wir alle einerlei Lehre führen und einträchtig schreiben: daß alsdann Einer kommt, der das alles verwirret und die Leute irre macht, der Doctor und Meister sein und die ganze Gemeinde auf seine Seite ziehen will; wie Arius aufgestanden und die ganze Kirche zu Alexandrien zerstört hat. Denn diese sind gleichsam zwei Schwestern und Zwillinge, nämlich, der Kain und Habel über die Erstgeburt und die Ketzerei.

327. Die Ketzer sehen, daß sie nicht können groß und berühmt werden, wenn sie in der Gemeinschaft und Einigkeit der Kirche bleiben. Was soll ich thun? sagen sie, mein Name wird gar verdunkelt, wird in andern Herrschaften und Landen nicht hoch gerühmt; derhalben muß ich einen Weg treffen, dadurch ich mich selbst etwas erheben und hervorthun möge, daß doch die Leute sehen, was ich ausrichten könne. Da erhebt sich denn ein Streit über der Erstgeburt und über dem Titel oder Namen der Kirche.

328. Die rechte Kirche aber wird von solchen ehrgeizigen Geistern gar verachtet und gleichsam mit Füßen getreten; wie Jsaak und Jakob von Esau und seinen Weibern allhier für die rechte Kirche nicht sind gehalten worden. Es sind Narren, werden sie gesagt haben, sie haben zumal einen kleinen Geist; man muß höher fahren. Denn solche Worte hat man zu unserer Zeit von Münzer gehört: D. der Luther und Andere haben wohl das Evangelium angefangen, haben es aber nicht hinausgeführt etc.

329. Auf solche Weise haben diese zwei stolzen Weiber den Esau auch angereizt, die Kirche zu reformiren, und zu ihm gesagt: Du bist der Herr, du magst nach deinem Gefallen Alles anrichten und regieren, du wirst nicht irren oder fehlen können; wohl, laß uns unsern Vater Religion hier auch einführen und in Schwang bringen. Was! dein Vater Jsaak ist ein grober Mann, der bleibt immer bei seiner Einfalt, lehrt uns nur beten. Das ist alles einfältig Ding, hat gar keinen Schein oder Ansehen, bewegt oder erweckt die Herzen der Menschen nicht. Man muß in die Kirche was bringen, das einen Schein habe und vor den Augen der Menschen köstlich daher leuchte. Unsere Väter, die Hethiter, beten Sonne und Mond an; derhalben lasset uns auch zugleich die Religion und Ceremonien der Hethiter mit halten.
\quellseite{75}

330. Das ist der große verderbliche Schade gewesen derselbigen Kirche. Denn die Menschenfatzungen und Ceremonien haben einen großen Schein und fallen den Leuten in die Augen, bewegen den gemeinen Pöbel, daß sie sich der äußerlichen Pracht und des großen Scheines hoch verwundern. Auf die Weise haben diese zwei heidnischen Weiber ihrer Viele unter dem Gesinde und auch etliche Nachbarn zu sich gezogen, und werden gesagt haben: Isaak ist ein alter kindischer Mann, er kann nichts ausrichten; Rebekka melkt die Ziegen und kühne; derhalben wollen wir einen andern und viel größeren Gottesdienst anrichten.

331. Also hat die falsche Kirche wegen der äußerlichen Pracht angefangen zu herrschen, die Verheißung aber und die reine göttliche Lehre ist verachtet worden. Denselbigen Jammer haben Isaak und Rebekka müssen ansehen und dulden. Esau ist mit seinen Weibern vorgezogen worden und hat seine Eltern verachtet; und haben also weder GOtt recht gedienet, noch ihre Eltern in Ehren gehalten, sondern haben mit ihrem Anhang die guten frommen Eltern erbittert und hart erzürnt.

332. Und eben also geht es uns jetzt auch noch, die wir sehen, daß so viel Spaltungen und Secten aufkommen, um welcher Ursache willen die Kirche in Angst und fast sehr betrübt ist. Aber es kann ja dies Unglück nicht verhütet werden. Wenn aber die Erstgeburt hinweg genommen ist, alsdann folgt gewißlich die Abgötterei auch nach. Derhalben hat die Kirche, so Isaak gehabt, in großer Gefahr gestanden, welche Gefahr von dem Erstgebornen hat angefangen, dazu auch von seinen Weibern, welche beide das Herz und die Seele Isaaks und Rebekkas hart betrübet und ganz erbittert haben. Und Isaak hat seinen Schmerz nicht verschwiegen, oder verschweigen, oder verhehlen können; denn es ist eine heftige Betrübniß gewesen und die für und für gewähret hat. Denn sie hatten Isaak und Rebekka nicht allein erzürnt und beleidigt, wie in der gemeinen lateinischen Uebersetzung steht, sondern sie haben sie noch immer ohne Aufhören verletzt und erbittert; haben nicht nachgelassen, die beiden, Isaak und Rebekka, zu erbitterten. Solches ist nicht allein einmal, zwei- oder dreimal geschehen, sondern die Erbitterung hat lange Zeit gewähret bis auf den harten Kampf, der nun folgt.

\chapter{Das siebenundzwanzigste Capitel}

\section[Erster Theil]{Erster Theil. Wie Isaak den Beruf faßt, Esau zu segnen, Rebekka aber sich bemüht, den Segen von Esau auf Jakob zu bringen}

\subsection{I}

\begin{quote}
\textbf{V. 1–4. Und es begab sich, da Isaak war alt worden, daß seine Augen dunkel wurden zu sehen; rief er Esau, seinem größern Sohn, und sprach zu ihm: Mein Sohn! Er aber antwortete ihm: hie bin ich. Und er sprach: Siehe, ich bin alt worden, und weiß nicht, wann ich sterben soll. So nimm nun deinen Zeug, Köcher und Bogen, und gehe aufs Feld, und fahre mir ein Wildpret, und mache mir ein Essen, wie ichs gerne habe, und bringe mirs herein, daß ich esse, daß dich meine Seele segne, ehe ich sterbe.}
\end{quote}\footnote{1. In diesem Capitel wird davon gehandelt, wie der Segen der Erstgeburt so wunderbarer Weise von Esau auf Jakob gewendet ist, und wird hier beschrieben das Gericht und Urtheil GOttes, welches bisher verborgen gewesen, aber endlich geoffenbart wird. Vor allen Dingen aber soll man merken und Acht haben auf die Zeitrechnung, auf daß wir wissen, in welchem Jahre diese Historie geschehen sei. Denn Esau und Jakob sind zu der Zeit sieben und siebenzig Jahre alt gewesen. Denn Jakob ist nach Mesopotamien gezogen, da er sieben und siebenzig Jahre alt gewesen ist. Daran ist kein Zweifel, die Zeitrechnung beweiset solches klarlich. Wenn man nun die vorigen sechzig Jahre des Alters Isaaks, ehe denn seine Söhne geboren worden, zu den sieben und siebenzig Jahren thut, so kommen hundert und sieben und dreißig Jahre heraus. Esau hat zwei Weiber genommen, da Isaak hundert Jahre alt war; über diese hundert Jahre hat er darnach gelebet sechs oder sieben und dreißig Jahre; derhalben ist dies, so jetzt folgt, geschehen im hundert und sieben und dreißigsten Jahre des Alters Isaaks.}
\footnote{2. Da aber Jakob in Mesopotamien sieben Jahre gedienet hat, wird er erst im vier und achtzigsten Jahre seines Alters ein Bräutigam. Esau hat vor ihm vier und vierzig Jahre seine zwei Weiber genommen und hat schon viel Kinder gezeuget, und wird von denselbigen ohne Zweifel auch Kindeskinder gehabt haben; denn er war nun sieben und dreißig Jahre ein Ehemann gewesen, in welchen man viel Kinder zeugen kann, vornehmlich von zwei Weibern. Darum hat er ein groß Gesinde und eine eigene Hütte mit seinem Weibe und Kindern, wie wir hernach sehen werden. Er ist der Fürst und Herr im Hause Isaaks, hat zwei Weiber, die sind Frauen im Hause; hat dazu viel Kinder und Kindeskinder, und ist der Hoffnung, er wolle ohne alles Hinderniß an seines Vaters Statt kommen und allein der Erbe sein. Jakob aber ist elend und verachtet, hat seine sonderliche Hoffnung, wartet nur etwa auf ein Geschenk: wenn er das empfangen hat, wird er aus seines Vaters Hause hinweg geschickt werden und alles Andere seinem Bruder Esau lassen.}
\footnote{3. Endlich aber hat Isaak gedacht: Siehe, du bist hundert und sieben und dreißig Jahre alt, du mußt nun auch an deinen Abschied aus diesem Leben denken, und wie du es mit deinen Kindern und Nachkommen ordnen und bestellen willst. Denn nach der Sintfluth haben die Väter nicht mehr zweihundert Jahre erreicht. Und waren sie damals schier alle gestorben, Eber ausgenommen, der zwei Jahre darnach gestorben ist. Die ganze Majestät und Herrlichkeit der Patriarchen war begraben, so daß Isaak und Eber noch übrig waren, welcher Eber zwei Jahr hernach stirbt. Derhalben ist Isaak allein mit seinen Söhnen der Patriarch und Papst in der ganzen Welt. Darum hat er also gedacht: Die andern Väter haben nicht zweihundert Jahre gelebt, einen oder zwei ausgenommen; so will ich auch auf kein länger Leben hoffen. „Und es begab sich“, sagt Mose, „da Isaak war alt worden“, nämlich, im hundert und sieben und dreißigsten Jahre, „daß seine Augen dunkel wurden“, vielleicht nicht so sehr vom Alter, als sonst irgend von wegen eines andern Unfalls.}
\footnote{4. Aber sonst im Alter pfleget nicht allein das Gesicht abzunehmen, sondern auch das Schmecken, Greifen, dazu alle Sinne und Glieder; denn das Alter ist ein Vorspiel des Todes. Wie auch die kleinen unmündigen Kinder sind, als wären sie todt, sehen und sehen doch nicht, schmecken und schmecken doch nicht; also werden die Alten auch wieder zu Kindern und werden vom Alter gar kindisch; die Augen werden ihnen dunkel, das Gehör nimmt auch ab, daß sie schwer hören. Derhalben, sobald der Mensch geboren ist, hebt es sich an, daß die Sinne abnehmen, und wenn der Tod herzukommt, da werden wiederum alle Kräfte und Sinne der Menschen matt und schwach. Da nun Isaak die Augen dunkel wurden entweder Alters halben oder aber sonst von wegen einer zufälligen Krankheit, daß er nicht mehr hat sehen können, da rief er Esau, seinem größeren Sohne, auf daß er ihn segnen möchte, ehe denn er stürbe.}
\footnote{5. Hier erhebt sich aber eine schwere und hohe Frage, nämlich: Warum Isaak nach der göttlichen Antwort, so droben (1 Mos. 25, 23.) der Rebekka ist gegeben worden, noch gleichwohl fortfährt und den größeren Sohn segnen will. Denn das ist ja gewiß, daß er gehört hat, daß zu Rebekka durch eine göttliche Antwort ist gesagt worden: „Der Größere wird dem Kleinern dienen.“ Nem: „Zweierlei Leute werden sich scheiden aus deinem Leibe; und ein Volk wird dem andern überlegen sein.“ Und darum hält Rebekka immer auf den Kleinern, nämlich, auf Jakob; Isaak aber dagegen hält auf Esau.}
\footnote{6. Die Juden bringen hier unnützes Geschwätz vor und sagen: Isaak habe von der göttlichen Antwort nichts gewußt; denn Rebekka habe ihm solches verborgen. Aber das reimt sich nicht, und ist wider die Art und Gebrauch der Eheleute, die solche Dinge einander pflegen zu sagen. Und Rebekka hat den Schmerz nicht verhehlen können, den sie gefühlt und empfunden hat, da sich die Kinder in ihrem Leibe mit einander gestoßen haben. Sie wird ohne Zweifel geklagt haben, daß sie in Gefahr ihres Leibes und Lebens stehe, und wird auf Befehl ihres Mannes hingegangen sein zu Eber oder zu Sem, auf daß sie seinen Rath oder Meinung hören möchte. Da hat ihr derselbe Patriarch die Antwort gegeben: „Du wirst nicht sterben; sondern zwei Völker sind in deinem Leibe“, und du wirst eine Mutter werden zweier Zwillinge, ja, zweier Völker; und die Kinder werden in ihrer Kindheit auch nicht sterben, sondern sie werden wachsen und groß werden, und die Verheißung des Lebens haben, auf daß sie wissen, daß sie Patriarchen werden sollen sehr großer Völker. Und da sie das gehört, ist sie wieder zu ihrem Manne gekommen und hat ihm die Antwort erzählt mit Freuden, wird gesagt haben: Sem hat mir gesagt, ich soll Zwillinge gebären, die da leben und zwei Völker zeugen sollen.}
\quellseite{76}

7. Worauf denkt aber nun Isaak nach der göttlichen Antwort, daß der Größere der Kleinere werden soll? Es ist wahrlich eine schwere Frage und ist dies ganze Capitel fast wunderlich. Die Auslegung der andern, so dies erste Buch Mose erklärt haben, ist gar kalt und eitel unnütz Ding. Ich will göttlich also denken, daß zwischen diesen Eheleuten ein freundlicher Streit gewesen und daß sie sich gütlich mit einander beredet haben über den Verstand der göttlichen Antwort. Rebekka hat zum Vater sagen mögen: Es läßt sich wahrlich ansehen, daß Jakob soll der Erstgeborne werden. Isaak wird dagegen gesagt haben: Du verstehest die Antwort nicht recht; mich dünkt, diese Prophezeiung oder Weissagung sei schon erfüllt, also daß Esau den Jakob, der im Mutterleibe der Größere gewesen, überwunden hat, da er den Kleineren hat wollen unterdrücken; Esau aber ist der Größere geworden und soll auch nun der Größere bleiben. Rebekka hat das Gegentheil gehalten und ist fest auf ihrer Meinung geblieben, bis daß hernach das vorgefallen, daß der Segen wunderbarer Weise auf Jakob gewendet worden ist.

8. Mittlerzeit aber ist derhalben zwischen den Eheleuten ein freundlicher Streit gewesen. Isaak wird versucht haben, Rebekka von ihrer Meinung zu bringen, wird gesagt haben: Meine liebe Rebekka, du fehlest, du verstehst die Antwort nicht recht. Rebekka aber wird ihm geantwortet haben: Esau wird der Erstgeborne nicht sein. Und hat dessen sehr wichtige Argumente und gute Ursache gehabt. Denn Esau war ein Jäger, hat Bündnis gemacht mit den Heiden und führt in seines Vaters Haus zwei sehr böse Weiber; hat seine Kinder nicht gelehrt, wie sie ihren Eltern sollen gehorsam sein: daraus hat sie geschlossen, Esau würde verworfen werden. Und das hat ihr der Heilige Geist eingegeben, der sie dessen erinnert, daß Esau, der zuerst geboren war, der Größere sei und der Andere der Kleinere.

9. Also nennen beide Eltern jedes seinen Sohn, den es lieb hat, den Größern und Kleinern. Solche freundliche oder liebliche Irrthümer können zwischen Eheleuten wohl täglich vorfallen. Daselbst ist hier auch also geschehen über eine nicht geringe Sache, und hat dieser Streit zwischen ihnen gewährt, bis daß es sich allhier nun umkehrt. Rebekka hält Esau für den Größern, und schließt, daß er derhalben der Erstgeborne nicht könne werden laut der göttlichen Antwort; darnach, so hatte er auch die Erstgeburt verachtet und verkauft. Dies hat sie alles fleißig gemerkt und behalten; wie denn die Weiber auch ihre Gedanken haben.

10. Isaak hat aber eine Meinung gefaßt, die dieser stracks zuwider ist, und hat für sich die Regel oder das Gesetz, nämlich: Esau ist der Erstgeborne, derhalben so gebührt ihm auch das Recht der Erstgeburt. Gleich als wenn man in der Grammatik sagt: Die Hauptwörter, die sich auf den Buchstaben a endigen, sind weiblichen Geschlechtes; es werden aber doch etliche ausgenommen. Also macht Rebekka auch eine Ausnahme; der Vater bleibt stracks an der Regel hängen ohne alle Ausnahme. Und kommt das auch dazu, daß Esau im Besitz ist: er regiert das Haus und ist Herr im Hause; Jakob aber bleibt immer ein Knecht, lebt ohne Weib und ist wie sonst ein Knecht im Hause, hat gar kein Regiment diese ganzen sieben und dreißig Jahre.

11. Diese Beschreibung der Personen und wenn man sie also gegen einander hält, hilft viel und erklärt auch die Historie. Esau führt nicht allein das Haus- und Regiment, sondern verwaltet auch allenthalben das Kirchenregiment mit der Religion und Gottesdienst; bei Abwesenheit seines Vaters hat er gepredigt, gebetet, geopfert, Vieh geschlachtet zum Opfer; er ist König und Papst in derselbigen Kirche und Gemeinde. Seine Weiber sind stolze hoffärtige Herrinnen und Königinnen, die werden sich ohne Zweifel gerühmt haben, daß sie im Besitze seien, und daß Esau, ihr Mann, das Priestertum und Regiment inne habe. Dagegen aber ist Jakob gar verworfen, ist gleichsam todt und begraben, ist wie ein fauler Klotz und „eine Wurzel aus dürrem Erdreich“, wie Jesaia Cap. 53, 2. von Christo sagt. Das ist aber ein Werk Gottes, der aus einem dürren Klotze einen feinen grünen Baum machen kann.

12. Unterdeß aber, dieweil die Sache also gestanden und ein Jedes seine eigene Meinung gehabt hat, wird Esau für den Auserwählten, Jakob aber für den Verworfenen gehalten, nach der Regel und Urtheil aller, die im Hause waren, und dazu auch nach dem Ansehen und Gewalt Isaaks, des Vaters. Allein Rebekka, die Mutter, ist dawider, welche diese sieben und dreißig Jahre wird Geduld gehabt haben, und viel hat leiden müssen, so der Sohn und seine Weiber unbillig und mit großem Stolze gethan und angerichtet haben. Sie hat Esau im Besitze müssen sitzen sehen, daß er das Regiment gehabt, gepredigt und das priesterliche Amt verwaltet hat, und daß Jakob für ein armes Knechtlein ist gehalten worden, der Wasser und Feuer hat müssen zutragen zum Opfer, und andere häusliche Arbeit, wie die Knechte pflegen, hat thun müssen.

13. Dies läßt sich ansehen, daß es sich nicht übel reime, daß man es an diesem Orte sage. Denn in diesem Leben fällt oftmals solcher Streit vor über Dinge, so ungewiß sind und daran man zweifelt. Wie zwischen Augustinus und seiner Mutter Monica auch ein Streit, diesem fast gleich, gewesen ist: die Mutter hat Tag und Nacht für ihren Sohn gebetet, daß er möchte von der Secte und Ketzerei der Manichäer errettet werden, welcher Secte er neun Jahre ist zugethan und anhängig gewesen, und darin also gefangen war, daß ihn niemand hat erretten oder ihm widerstehen können, wenn man mit ihm darüber disputirt hat; die Mutter hat Viele dazu gezogen, die mit ihm davon sollten reden und ihn von diesem Irrthume bringen, aber es hat ihm niemand weder überwinden, noch davon bringen können, also daß seine Mutter darüber jämmerlich in Angst gerathen und seinen Trost mehr gehabt hat. Da trägt es sich zu, daß sie in der Nacht einen Traum gehabt, da ihr ein Jüngling erschienen, der in einem Kreis stand und freundlich mit ihr redete: Liebes Weib, warum weinst du also? Sie sagte: Ich bin betrübt, daß mein Sohn soll verloren werden. Da antwortete ihr der Jüngling: Wo du bist, da ist er auch. Da erwacht die Mutter aus dem Schlafe, ist fröhlich, und erzählt ihrem Sohne den Traum und sagt: Lieber Sohn, ich habe endlich Trost gefunden, und habe die Verheißung, daß du sollst zu meiner Regel oder Glauben bekehrt werden; und erzählte ihm die Worte des Jünglings. Da fährt Augustinus den Spruch um und sagt zu seiner Mutter: Lieber Mutter, du wirst nicht recht gehört haben, sondern der Jüngling hat also gesagt: Wo ich bin (und hat das Wörtlein „ich“ auf sich und seine Lehre gezogen), da wirst du auch sein. Die Mutter hat sich das nicht lassen anfechten, sondern ist beständig bei ihrem Verstande geblieben, und Augustinus bekennt, daß ihn solche Beständigkeit seiner Mutter sehr bewegt habe, daß er desto leichter dem Sinn und Meinung seiner Mutter beigepflichtet.

14. Solcher freundlicher Zank oder Streit ist gemein beide zwischen Eheleuten und sonst auch zwischen guten Freunden, als zwischen dem Pomponius und Cicero und hier zwischen Isaak und Rebekka. Und sie haben nicht allein über den Verstand der Antwort disputirt, sondern auch beider, Esaus und Jakobs, Art und Sitten fleißig gegen einander gehalten. Meine liebe Rebekka, wird Isaak gesagt haben, du siehst, daß Jakob ein schlichter einfältiger Mensch und verzagt ist, und daß er weder nach dem Regimente noch Priesterthum etwas fragt; Esau aber führt Bogen, Pfeil und Schwert, der wird hurtiger und geschickter sein zum Regimente. Wie wir uns pflegen Gutes zu versehen zu denen, an denen wir sehen, daß sie ein wenig fertiger und geübter sind denn Andere, ihre Geschäfte auszurichten.

15. Danach hängt Isaak an dem äußerlichen Verstande, so die Antwort haben soll nach dem Buchstaben, wie die Juden an dem fleischlichen Verstande der Schrift hängen; denn derselbigen Ebenbild ist Isaak, da er so dunkel sieht und alt geworden ist. Dagegen aber, da die heilige Matrone Rebekka durch die Bitterkeit der Weiber Haus und die große Hoffart ihres Sohnes verletzt und betrübt worden ist, hält sie durch Erleuchtung des Heiligen Geistes auf Jakob und sagt zu ihrem Manne: Esau wird mitnichten der Herr werden, seine Sitten und sein Leben gefallen mir gar nicht, ob er wohl dafür angesehen wird, daß er zum Regimente geschickt sei; Jakob demütigt sich, ist gehorsam und wird von seinem Bruder gar unter die Füße getreten; darum ist die Verheißung noch nicht erfüllt; denn die Demütigen werden erhöht, und die Hoffärtigen werden gedemütigt und gestürzt. Denn dies ist eine gewisse Regel in der Schrift, die nicht fehlet, wie St. Petrus sagt 1. Epist. 5, 5.: „Gott widersteht den Hoffärtigen, aber den Demütigen gibt er Gnade.“
\quellseite{77}

16. Dieweil nun die Regierung unsers Herrn Gottes damit übereinstimmt, so bringt Rebekka hart auf diese Regel. Dieweil Esau merkt, daß er von seinem Vater hoch in Ehren gehalten und vorgezogen wird, und derselben von wegen der Gunst, so er bei dem Gesinde und gemeinen Volk gehabt, stolz und aufgeblasen wird; darum wird er gedemütigt werden. Und Gott setzt auch hernach das Gegenteil, also daß der Andere, nämlich Jakob, der jetzt verworfen ist und im Staube liegt, erhöhet wird.

17. Dies sind die Gerichte Gottes. Und das sage ich darum, daß man so viel desto leichter verstehen möge, warum doch Isaak so wunderbarer Weise betrogen worden sei, also daß es schier anzusehen ist, als sei diese Geschichte erlogen oder erdichtet. Denn man muß es dafür gewiß halten, daß er die göttliche Antwort wohl gewußt, und daß er sehr hart über dem Verstande gehalten habe, daß Esau der Größere wäre, der zuvor in Mutterleibe von Jakob getreten und überwunden wäre worden, und der würde nun wiederum Jakob, seinen Größern, unter die Füße treten und überwinden. Denn sonst, wenn man es also nicht versteht, wird man den Text nicht ausgleichen können, daß es sich mit einander reimen möge.

18. Der Vater hat am Buchstaben gehangen, wie Abraham auch that, der Ismael auch für den Erstgebornen gehalten hat; aber Rebekka hat stärkere Argumente und bessere Ursachen ihres Verstandes gehabt. Darum duldet sie wohl die Hoffahrt und Gewalt Esaus, und befleißigt sich doch gleichwohl in ihrem Herzen, daß sie auf allerlei Gelegenheiten Achtung gibt, auf daß sie die Erstgeburt von Esau auf Jakob wenden möchte. Und mit demselbigen Rath ist sie nicht erst in dem Jahre oder Tage umgegangen, sondern hat denselbigen lange Zeit in ihrem Herzen auf mancherlei Weise erwogen und also gedacht: Ich will sehen, wie ich meinem Sohne den Segen möge wiedererlangen.

19. Und ist wohl glaublich, daß sie dies nicht von sich selbst oder aus ihrem eigenen Verstande erdichtet, sondern daß sie dessen der Heilige Geist erinnert habe durch den Patriarchen Eber, nämlich, daß sie an die Ausnahme von der gemeinen Regel gedacht hat: der wird sie ermahnt und erinnert haben, daß der Segen durch göttlichen Befehl wie mit einer sonderlichen Ausnahme von der gemeinen Regel Jakob und nicht Esau zugehöre. Wiewohl sich Esau auf den Buchstaben des Gesetzes und auf die Regel verlassen und sich damit geschützt hat. Aber, wie St. Paulus Gal. 3, 18. sagt, die Kinder des Gesetzes sind nicht die rechten Erben, sondern die Kinder der Verheißungen, die vom Gesetze ausgenommen werden. Darum weicht Rebekka von ihrer Meinung nicht, und verläßt sich darauf gewißlich, Jakob solle der Erbe und Herr werden. Dagegen will Isaak auch von der Regel und vom Gesetze nicht weichen, bis auf den Zeitpunkt, da die Mutter und der Sohn anfangen, den Vater mit Lügen zu hintergehen und zu betrügen.

20. So ruft nun Isaak seinem größeren Sohne Esau, und begehrt von ihm, er wollte ihm ein Essen machen, dazu er wüßte, daß der Vater Lust hätte. Daraus man ja klarlich abnehmen kann, daß er es in seinem Herzen gewiß dafür gehalten hat, daß Esau der Größere wäre. Denn er redet dies frei heraus mit vollem Vertrauen; wie die frommen heiligen Menschen pflegen zu thun, die sind nicht zweizüngig, haben auch kein doppelt Herz, sondern wo sie hin geneigt sind, dahin sind sie mit ganzem Herzen geneigt: sie sind entweder nichts oder sind es gar. Dagegen sind die betrüglichen Menschen mit sich selbst nicht eins, sind unbeständig und haben zweierlei Zungen im Munde. Aber hier ist das Herz des Vaters Isaak unzerteilt; er ist ganz fest davon überzeugt, daß Esau der Erstgeborne sein soll; darum wird er auch hernach so gar verwirrt und hart erzürnt, da er merkt, daß er in solcher Einfalt und gewissem Vertrauen betrogen worden ist. Daselbst kommt von der Einfalt, welche das Herz nicht theilt in mancherlei Meinung. Gleich als wenn an einen Gott glaubt, der kann nicht zwei Götter haben; denn da will Gott das Herz haben, da er 5 Mos. 6, 5. sagt: „Du sollst den Herrn, deinen Gott, lieb haben von ganzem Herzen.“ Derhalben, daß er betrogen wird, kommt von der Ursache her, daß er sich auf seine Meinung, so er gefaßt hatte, allzusehr verlassen hat.

21. Esau aber wird hier beschrieben, als der sehr gehorsam und fromm sei; darum setzt auch Isaak all sein Vertrauen auf ihn, daß er werde der Erbe sein, der ihm folgen werde; und redet auch mit ihm nicht anders, als mit Einem, der an seiner Statt das Erbe besitzen werde. Als wollte er sagen: Mein lieber Sohn, es will nun fast Zeit sein, daß die Verheißung erfüllet werde. Du bist der Größere und der Erstgeborne, derhalben so wollen wir nun verschaffen, daß du mir im Regiment und Priesterthum am nächsten nachfolgen mögest; denn ich besorge mich, es möchte zwischen dir und deinem Bruder ein Zank werden, von dem ich jetzt länger denn dreißig Jahre etliche Anfänge gesehen habe. Derhalben, auf daß solche Uneinigkeit zwischen euch Brüdern nun verhütet werden möge, habe ich mir vorgenommen und bei mir beschlossen, daß ich einen gewissen Erben hinter mir lassen und verordnen wolle, der beiden, dem Regimente und auch der Kirche, möge vorstehen. Dazu erwählet er nun den Esau; denn der ist ihm sehr gehorsam. Und ist Isaak noch der Meinung, daß Rebekka, die Mutter, irre, welche den Jakob lieb hatte, und hoffte, der solle der Erstgeborne werden. Darum gehen diese Worte Isaaks aus einem solchen Herzen, das darauf hart gestanden und dessen wollte gewiß sein, daß Esau würden müssen für den rechten Erben erklärt werden. „Gehe hin“, sagt er, „aufs Feld“, nimm deinen Köcher und Bogen etc., „fange mir ein Wildpret, und mache mir ein Essen“ etc., „daß dich meine Seele segne, ehe ich sterbe.“ Da ist gar kein Argwohn, daß es anders gehen könnte; und wird doch gleichwohl betrogen werden.

22. Es scheint aber aus dieser Rede Isaaks, daß sich Esau fein hat wissen einzuschmeicheln bei dem Vater. Die Schrift sagt zwar das nicht von ihm, daß er seinen Vater habe in Ehren gehalten, sondern daß er seinen Vater geheuchelt und weiblich geschmeichelt hat. Klappern gehört zum Handwerk. Er hat sich, sage ich, können einschmeicheln und sein dazu schicken, wohin er gesehen hat, daß der Vater geneigt gewesen; darnach hat er sich stellen können, als ob er sonderliche Sorge trüge für den Vater, mehr als sein Bruder. Darum verzagt Isaak an Jakob, der immer um die Mutter, und dazu ganz bescheiden und stille ist. Es läßt sich mit ihm ansehen, als sei er dazu geboren, daß er solle ein Knecht sein. Esau aber ist der Herr im Hause, der da pflegt zu jagen und seinem Vater von der Jagd gut auserlesen Essen zu bringen: nicht gemeine Speise, als etwa eine Henne, Wild oder Käse, wie ihm Rebekka und Jakob pflegten zu geben, sondern Wildpret. Und also beweist er viel größere Fürsorge und Eifer um den Vater, als Jakob oder die Mutter. Und mit solchem Dienst und Fleiß macht er, daß ihm der Vater immer günstiger und je länger je mehr angenehm geworden ist, der ihm sonst ohne das sehr günstig war. Ja, also muß man mit denen umgehen, die wir hintergehen und betrügen wollen.

23. Es haben ihm aber auch seine Weiber fleißig dazu geholfen, welche den alten Mann mit ihren Schmeichelworten sein haben schmieren können, und unterdes die arme Mutter über die Maßen sehr betrübt und erbittert, dieweil sie wußten, daß dieselbe nicht damit zufrieden war, daß Esau die Ehre haben und Jakob also sollte verachtet werden. Aber Rebekka hat gedacht: O Esau, du ehrst deinen Vater nicht um seinetwillen, als wärest du so fromm, sondern um deines eigenen Nutzens willen: non propter te, sed propter tuum, wie man gemeiniglich auf Lateinisch pflegen zu sagen. Du meinst damit nicht deinen Vater, sondern dich selbst. Denn er geht nicht damit um, daß er seinen Vater aus reinem Herzen ehren wolle, sondern er sucht damit für sich, seine Weiber und Kinder das Regiment und Priesterthum zu erlangen. Das hat Rebekka gemerkt.
\quellseite{78}

24. Denn also sind alle Heuchler gesinnt. Es sind Wölfe in Schafskleidern, damit sie ihre Eltern, Herren und Brüder etc. betrügen. Isaak hat zwar darauf keine Achtung gehabt; aber Rebekka hat es sein merken können und wird diese Hinterlist durch einen wunderlichen Rath zu Schanden machen. Derhalben ist Esau ein gehorsamer Sohn gewesen, aber Gott hat davon nichts gewußt. Dagegen läßt es sich mit Jakob ansehen, daß er seinen Vater gar mit nichts geehrt habe, aber nach dem Urtheil Gottes ehrt er seinen Vater allein. Denn Gott sieht in das Herz. Esau bringt Hasen und Wildpret und schmeichelt damit dem alten Vater, darum wird er desto mehr geliebt; aber es ist lauter Heuchelei. Daher wird Isaak immer mehr in seinem Wahn bestärkt, da er sich mit solcher Heuchelei betrügen läßt.

25. Es ist aber zumal ein böses Exempel, und die Gottseeligen sollen lernen, daß sie solcher Heuchelei feind sein und den Herrn fürchten. Denn Gott fragt darnach nichts, wenn Einer groß oder viel Ding thut; wie Esau das vierte Gebot erfüllt im höchsten Grad, wie es äußerlich scheint: aber er thut in Wahrheit nichts weniger; denn sein Herz ist nicht rechtschaffen und rein, welches doch Gott vor allen Dingen haben will. Darnach achtet Esau seines Vaters nicht, sondern ist ihm um seine Erstgeburt zu thun; welches man hernach mehr sehen wird, da er den Segen verloren hat. Denn wo er seinen Vater recht und von Herzen geliebt hätte, so hätte er auch den Fluch mit Geduld ertragen können und gesagt: Lieber Vater, wiewohl es dir und mir anders geschienen hat, so will ich doch meiner Mutter nun auch gehorchen, und will damit Geduld tragen und vorlieb nehmen, daß der Segen dermaßen von mir ist gewendet worden. Aber sein Herz steht viel anders; denn er sagt: „Es wird die Zeit bald kommen, daß mein Vater Leid tragen muß; denn ich will meinen Bruder Jakob erwürgen.“ Soll man aber also die Eltern ehren? Derhalben hat Gott das teuflische und mörderische Herz gesehen, als der die Herzen prüft. Esau hat es wohl eine Zeitlang können verbergen, daß es der Vater nicht hat merken sollen, aber darnach fährt er mit Gewalt heraus.

26. Das fördert aber die Ehre, so die Kinder den Eltern zu erzeigen schuldig sind, daß wenn sie schon etwas wider deinen Willen verordnen würden, daß du ihnen nichts desto weniger gehorsam sein und bleiben sollst. Darum beschreibet Mose Esau also, daß er ein großer Heuchler gewesen, der sich mit sonderlichen Diensten bei dem Vater verdient gemacht und ihn doch mit solcher Heuchelei schändlich betrogen hat.

27. Wiewohl nun Isaak, der da blind und gar alt ist, von dem Sohn, der dem äußerlichen Schein nach der allergehorsamste war, betrogen worden ist: so urtheilt doch Gott recht und wendet den Segen von dem bösen Heuchler Esau auf den frommen und einfältigen Jakob. Und das wird uns zur Lehre und Exempel vorgehalten, auf daß wir lernen Gott fürchten, den wir mit unserer Heuchelei nicht betrügen können, sondern er sieht in das Herz. Darum ist dies das Abenteuerte, daß man sein einfältig handle, was man nur zu handeln hat, und mit aufrichtigem Herzen lebe. Und man soll mit Gott nicht scherzen weder in der Religion, noch im Weltregiment, noch im Hausregiment, darum daß solcher Scherz und Heuchelei nimmer pflegt ungestraft zu bleiben.

28. Nachdem nun der Heuchler sich mit seinen zwei Weibern und Kindern verbunden und vereinigt hat, daß sie den alten Vater betrügen wollen, und daß er, Esau, möchte zum Erben gesetzt werden: so folgt nun der Rath Rebekkas, welcher jenem Anschlag stracks zuwider ist. Denn Rebekka hat eine große Klugheit begangen, und kehrt um und verhindert alle Anschläge, so Isaak, ihr Mann, und Esau, ihr Sohn, gemacht hatten.

29. 5–10. Rebekka aber hörte solche Worte, die Isaak zu seinem Sohne Esau sagte. Und Esau ging hin aufs Feld, daß er ein Wildpret jagte, und heim brächte. Da sprach Rebekka zu Jakob, ihrem Sohn: Siehe, ich habe gehört deinen Vater reden mit Esau, deinem Bruder, und sagen: Bringe mir ein Wildpret, und mache mir ein Essen, daß ich esse, und dich segne vor dem Herrn, ehe ich sterbe. So höre nun, mein Sohn, meine Stimme, was ich dich heiße. Gehe hin zu der Heerde, und hole mir zwei gute Böcklein, daß ich deinem Vater ein Essen mache, wie ers gerne hat. Das sollst du deinem Vater hinein tragen, daß er esse, auf daß er dich segne vor seinem Tode.

29. Hierher gehört die Beschreibung der Weise, so die Väter gebraucht haben, wenn sie Andere gesegnet oder aber den Segen empfangen haben. Aber Mose erzählt nur nach der Historie, was der Erstgeborne habe thun müssen, wie er habe müssen gekleidet sein, die Lehengüter des Segens zu empfangen. Ueber das setzt er nichts mehr dazu, so sie doch etliche gewisse und sonderliche Ceremonien dazu gebraucht haben. Wie jetziger Zeit die Fürsten, wenn sie vom Kaiser die Lehen empfangen, nicht gekleidet sind, wie sie sonst gemeiniglich ihre gewöhnlichen Kleider tragen; sondern ziehen königliche Kleider an mit trefflichen Helmen, Schild und Waffen. Das sind äußerliche Ceremonien, und desselbigen Gepränge gebraucht die Welt, dasjenige, so man vorhat, bekannt zu machen und zu zieren. Also hat Mose auf Josua seine Hände gelegt, 4 Mos. 27, 20. 23., und die Schrift sagt: Pones super eum laudes tuas, das ist: Du sollst ihn vor dem Volk mit einer kurzen Rede loben, daß sie wissen, wenn er von dir also gelobt ist, daß er dessen werth sei, daß sie ihn hören und ihm folgen sollen. Wie wir auch thun, wenn die Schüler zusammen kommen und die Doctores oder Magister öffentlich abkündiget werden. Dasselbe geschieht auch also mit den Hochzeiten: da gehen die jungen Eheleute in die Kirche, haben sich mit neuen Kleidern geschmückt, die Nachbarn werden geladen, daß sie auch wollen mitgehen und solch hochzeitlich Gepränge und gewöhnliche Ceremonien helfen zieren; da bereitet man ein Mahl, dem Bräutigam und der Braut zu Ehren, daß die Leute, so geladen sind, mit einander essen und trinken, und fröhlich sein sollen.

30. Eben auf dieselbe Weise haben sie auch etliche sonderliche Ceremonien gebraucht, wenn sie den Segen der Erstgeburt empfangen haben. Erstlich hat man dem Vater müssen ein Essen bereiten. Dasselbe ist ohne Opfer nicht geschehen: sie haben beide Gott ein Opfer gethan und haben es also mit dem Opfer angefangen. Wie wir auch, wo wir öffentlich zusammen kommen, mit dem Gebet oder Danksagung pflegen anzufangen. Und wie wir beten, wenn wir denen die Hände auflegen, so zum Kirchen- oder Predigtamt ordinirt und bestätigt werden: also hat Esau auch anfänglich ein Opfer gethan, darnach hat er dem Vater davon auch etwas geben wollen.

31. Die andere Ceremonie ist diese gewesen: der den Segen empfangen sollte, hat ein priesterlich Kleid, das sich zu dem priesterlichen Gepränge und zu dem Segen der Erstgeburt gereimt hat, anziehen müssen. Und da zeigt Mose an, daß Rebekka im Besitz der priesterlichen Kleider gewesen sei. Jakob hat derselben nicht geachtet und hat es mit der Erstgeburt fast verloren gegeben, hat für seine Person daran verzaget; sie aber bewahrt die Kleider mit großer Sorgfalt. Sie sind aber dazumal beide, nämlich Esau und Jakob, sieben und siebenzig Jahre alt gewesen, und sind doch diese Kleider noch immer in der Verwahrung bei Vater und Mutter geblieben, von welchen man sie hat fordern müssen, wenn man sie hat gebrauchen sollen. Und ich halte dafür, daß Rebekka dabei eine sonderliche Absicht gehabt, daß sie die Kleider in ihrer Gewalt hätte, als die eine Krone, das Scepter, das Königskleid, und was sonst mehr zu der Ehre und Herrlichkeit der Erstgeburt gehört hat. Dieses können wir also sicher aus Mose abnehmen und schließen von denselbigen Ceremonien.

\subsection{II}

32. Da kommt aber nun eine sehr schwere Frage, und ein solcher Knoten, der nicht wohl aufzulösen ist, nämlich: Ob es sich auch geziemt habe, daß Rebekka und ihr Sohn Jakob haben lügen, und den heiligen Patriarchen, da er blind und alt war, in einer solchen großen und wichtigen Sache betrügen sollen und also Esau, seinen Bruder, den Segen und die Erstgeburt nehmen? Denn, erstlich, sündigen sie wider das Gesetz und die gemeine Regel; darnach, wider die göttliche Antwort, welche Isaak anders verstanden hat, als es Rebekka verstehen wollte; zum Dritten, wider den Willen und Gewalt des Vaters. Und Rebekka reißt die königlichen und priesterlichen Kleider an sich, damit sie Jakob schmückt und ihn also zum Vater führt, an welchen der Vater nie gedacht hatte, daß er ihn hätte wollen zum Erben machen.

33. Dies ist ja ein großer Frevel und Kühnheit, dazu ein sehr großer Betrug und schwerer Schaden. Denn dieses ist ja keine Scherz- oder Dienstlüge, sondern ist sehr schädlich. Denn die Erstgeburt ist ein solch Ding, davon man in Religions- und Kirchensachen pflegt zu streiten und die das ewige zugleich mit dem zeitlichen Leben an sich hangen hat. Derhalben ist es ja eine große Sünde, davon Einer nicht meinen sollte, daß sie irgend ein Mensch, wenn er auch der allergottloseste sein sollte, geschweige denn solche heilige Menschen, wie Rebekka und Jakob, begehen sollten. Und es läßt sich ansehen, daß man keine größere Sünde erdenken oder von keiner größeren Bosheit sagen könnte, auch nach dem Urtheil der Welt. Es ist ein sehr böser Betrug in solchen wichtigen Sachen, wider die Antwort und göttliches Gesetz. Darnach, so hat Esau den Besitz und das Regiment beide im Hause und auch in weltlichen und Kirchensachen; und wird doch gleichwohl wider dies alles und wider den Willen des Vaters von seinem Bruder Jakob davon gestoßen.
\quellseite{79}

34. Ich weiß aber nichts aus den Schriften der Väter, Augustins und der andern, zu finden, daß mir helfen möchte diese Frage auszulegen. Denn sie gehen alle an diesem Texte vorüber und sagen von diesem Stück nichts. Derhalben werden wir es müssen errathen. Und ist wahrlich ein wunderlich Ding, daß dies Weiblein so kühn ist, daß sie ein solch groß Ding hat thun dürfen und mit so großer Gefahr ihren Sohn vorstellen in Abwesenheit seines Bruders, und dazu wider des Vaters Willen, welcher auch den Betrug schier merkt und erschrickt, da er die Stimme Esaus hört und seine Hände angreift. Denn er sagt V. 22. also: „Die Stimme ist Jakobs Stimme, aber die Hände sind Esaus Hände“; item V. 20.: „Wie hast du so bald funden?“ Und Jakob erschrickt selbst auch, da er den Anschlag seiner Mutter hört, und sagt V. 12. voller Furcht zu der Mutter: „Ich fürchte, daß ich über mich einen Fluch, und nicht einen Segen bringen werde.“ Aber Rebekka folgt stracks dem Geist, der sie geführt und getrieben hat: und darum hat sie diese Gefahr und \unsicher{Verwegenheit} alle mit großem Geist und starkem Muth aus den Augen geschlagen.

35. Und es läßt sich ansehen, daß sie diese Sache aus ihrem eignen Rath nicht also gehandelt, sondern daß sie den \unsicher{Befehl} vorlängst durch anderer Leute Erinnerung in ihrem Herzen also gefaßt habe. Sie hat aber ihren Mann nicht überreden können, daß er den Segen wider die Regel auf Jakob hätte wenden wollen. Darum ist sie oft zu Eber gekommen, und wird gesagt haben, daß Isaak von der Regel nicht weichen wolle und Esau segnen. Da hat sie von Eber gehört, daß er zu ihr gesagt hat: Thue, was du nur thun kannst, Isaak zu betrügen und das Recht der Erstgeburt hinweg zu reißen: versuche alles, was du nur kannst. Denn der Glaube ist nicht ohne Wort. Darum wird sie ohne Zweifel das mündliche Wort gehört haben und hat dasselbe in ihrem Herzen erwogen, sonst hätte sie dieses nicht anfangen dürfen.

36. Denn der Heilige Geist kommt nicht ohne Wort, sondern er will durch die Harfe kommen, das ist, durch das Wort, wo man davon redet und dasselbe betrachtet, oder durch die mündliche Stimme des Vaters, der Mutter oder sonst Anderer. Sonst kommt der Teufel. Aber wie David, Jesaia und die ganze Schrift bezeugt, so kommt der Heilige Geist mit dem Wort und durch das Wort, nach dem Spruch im ersten Psalm V. 2.: „Wohl dem, der Lust zum Gesetz des Herrn hat, und redet von seinem Gesetz Tag und Nacht“ etc.

37. Darnach, so pflegt Gott die Regel bisweilen zu übertreten, darum daß die, so nach der Regel und nach dem Gesetz Kinder sind, ihrer Regel auch pflegen zu mißbrauchen. Und insofern ist ihnen das Gesetz nichts nütze. Denn wo du anfängst dasselbe zu mißbrauchen, so hebt es Gott auf, daß es nicht mehr gelten muß. Esau war stolz und hoffärtig wider die göttliche Antwort, und gedachte bei sich selbst also: Jakob ist verworfen, ist gestorben und begraben; und hat also seinen armen Bruder unterdrückt. Zudem hat er auch dulden können den Muthwillen und Bitterkeit, so seine Weiber gegen seine Eltern geübt haben.

38. Derhalben, wenn die Hoffahrt kommt, alsdann hört die Regel auf. Denn Gott hat die Königreiche und Herrschaften nicht eingesetzt, gibt auch den Menschen die andern Gaben nicht darum, daß sie davon sollen hoffärtig und stolz werden, und wider die Armen und Elenden toben und wüthen; sondern alles, was er den Menschen gibt, es sei Reichthum, oder Stärke, oder Schöne, das gibt er zu seiner Ehre und dem Nächsten zu Nutz.

39. So ist nun die endliche Ursache der Gaben Gottes nicht, daß diejenigen, so die Gaben haben, damit Wollust treiben und Tyrannei oder Muthwillen üben sollen; sondern der rechte Gebrauch der Gaben soll zu Gottes Ehre, Nutzen und Wohlfahrt des Nächsten gerichtet sein. Aber die Menschen empfangen den Segen Gottes, Königreiche, Priesterthum, Stärke, Kraft, Klugheit, und haben zwar wohl die Gaben von Gott, aber fragen nichts nach dem Zweck, wozu sie die Gaben gebrauchen sollen. Warum bist du aber ein König? Warum bist du ein Fürst, Priester, Vater oder Mutter? Ei, darum, wirst du sagen, daß ich es in diesem Leben soll gut haben, soll meine Lust haben, in Wollust leben, alle meine Begierde erfüllen; darum bin ich gelehrt und reich, daß ich dadurch einen großen Namen und große Ehre unter den Leuten erlangen möge. Ja, daselbst wird die Regel, davon du dich so hoch rühmst, gar aufgehoben und muß nicht mehr gelten.

40. Denn Gott will seinen Segen zu seinem andern Ende ausgeschüttet und gegeben haben, denn zu seiner Ehre, zum Lobe deß, von welchem die Gaben herkommen, und zur Wohlfahrt und Besserung der Kirche. Die Obrigkeit hat darum ihre Ehre, daß sie der Gemeinde und den Unterthanen möge nütze sein. Mann und Weib werden darum in den Ehestand zusammengefügt, daß sie Kinder zeugen und dieselbigen zu Nutz dem Haus- und gemeinen Weltregimente aufziehen sollen. Aber die Welt schneidet den Endzweck in allen Gaben Gottes ab.

41. Das hat Rebekka gemerkt, und hat die göttliche Antwort etwas fleißiger angesehen und betrachtet als Isaak, ihr Mann, der nur an der Regel ist hängen geblieben, nämlich also: Esau ist zuerst geboren und Jakob zuletzt; darum so ist Esau auch der Erstgeborne. Dagegen aber hat sich Rebekka an die Verheißung gehalten: sie hat nicht sowohl die Sache, als den Zweck der Sache für sich. Denn das Leben und die Sitten Esaus reimen sich mit der Antwort nicht. Darum hat sie also gedacht: Mein Sohn Esau hält sich also, daß ich sehe, daß er der Erstgeborne nicht werden wird. Darnach hat er seine Erstgeburt verkauft, und hat den Endzweck der Erstgeburt, dahin sie eigentlich gerichtet ist, verachtet und gesagt: Ich will nur jetzt das rothe Gericht hinnehmen; das ist so viel gesagt: Wenn ich nur der Herr und Fürst und Priester sein mag, so frage ich nach dem Endzweck, was den Nutzen der Kirche belanget, gar nichts; ich denke nur darauf, daß ich nach meiner eigenen Begierde leben und große Ehre haben möge. Also hält nun Isaak über der Regel, aber Rebekka bringt auf die Ausnahme. Denn Esau hält sich also, daß er der Erstgeborne nicht ist und auch nicht sein kann. Darnach, so stimmt die Antwort mit dem überein, so ihr von Gott eingegeben war wider die Regel.

42. Das ist es, worum es sich bei dem ganzen Streite handelt: Rebekka und Jakob fahren immer fort und beharren bei der Ausnahme und bei dem Geiste, darum daß Esau die Erstgeburt verkauft hatte, und dazu um des Mißbrauchs willen der Erstgeburt nicht werth gehalten wird, daß er sie haben sollte. Derhalben kann man auf die Frage also antworten, als nämlich für Rebekka und Jakob: Was mir von Gott gegeben ist, und davon ich weiß, daß es mir zugehört, dasselbe mag ich mir auch zueignen und mich desselben anmaßen, es sei gleich mit welcherlei List oder Betrug solches geschehen möge; denn ich bin schuldig, bei einer Todsünde zu denken, mich zu stellen, oder auch mein Vornehmen zu verhehlen, wie ich nur kann und mag, damit dasjenige geschehen möge, so mir von Gott befohlen ist.

43. Ezra sagt, daß Jakob und Rebekka mitnichten zu entschuldigen sind. Und wir können sie zwar auch von der Sünde nicht gar lossprechen; denn es sind Menschen gewesen, die fleischliche Bewegungen gehabt wie wir auch. Es ist aber nicht vonnöthen, daß man sie deswegen, als sollten sie gelogen haben, oder daß man ihnen wollte irgend einen Betrug aufrücken. Denn dieweil Gott sagt: Ich will haben, daß Jakob soll der Erstgeborne sein, so ist daselbst das Gesetz aufgehoben: „Wo aber das Gesetz nicht ist, da ist auch keine Uebertretung“, Röm. 4, 15.

44. Dieweil nun Jakob und Rebekka deß gewiß waren, daß die Erstgeburt dem Jakob gehörte, nämlich, aus der göttlichen Antwort und dieweil sie an Esau so böse Früchte und Sitten gesehen: so haben sie billig das Gesetz und die Regel verachten sollen, und der Ausnahme folgen, dadurch GOtt die Erstgeburt von Esau auf Jakob gewendet hat. Darum hat Rebekka darauf gedacht, wie sie Isaak, ihren Mann, und Esau, ihren Sohn, dazu alle, die im Hause waren, betrügen möchte. Denn sie gehorcht jetzt nicht der Regel oder dem Gesetze, sondern GOtt, der die Erstgeburt also verwendet und damit bisweilen wider die Regel: darum hat sie nicht gesündigt.
\quellseite{80}

45. Und dies, halte ich, sei eine einfältige und rechte Antwort auf die Frage von dem Betrüge, Lüge und Schaden, so Esau, und allen, die ihm angehangen haben, das ist, der ganzen Kirche im Hause Isaaks, widerfahren ist. Denn sie waren ihm alle gehorsam und hielten ihn für den Erstgebornen, sintemal er dasselbe Recht und Herrschaft länger als sieben und dreißig Jahre inne gehabt. Darum zweifelten sie daran gar nicht, er würde wohl der Erstgeborne bleiben. Aber da der Segen der Erstgeburt jetzt sollte bestätigt werden, ist es alles plötzlich und unversehens umgekehrt worden wider der ganzen Kirche und des Patriarchen Isaak Rath und Willen. Denn Rebekka folgt stracks dem nach, daß GOtt die Erstgeburt von Esau gewendet hat, darum daß er verworfen war nach der göttlichen Antwort und von wegen der Verachtung der Erstgeburt, dazu um der Frucht und des Lebens willen, so dieser Herrlichkeit nicht werth gewesen. Und da stimmen sein mit einander überein die Ursache und die Wirkung und kommt auch dazu die Bestätigung des heiligen Geistes: Isaak aber wird betrogen. Welches zwar an sich selbst eine schwere und greuliche Sünde ist wider den Mann, wider den Vater und Nächsten, ja, wider seinen Bruder und wider die Regel. Aber wie es dem Herrn gefallen hat, also ist es geschehen.

46. Denn der Gehorsam der Eltern, die Brüderschaft und Liebe sind der ersten Tafel unterworfen: die andere Tafel muß weichen und gilt nichts, wenn sie sich mit der ersten stößt. Wo man den Bruder oder Vater lieb hat mit Verachtung Gottes, das ist eine Feindschaft wider Gott. „So jemand zu mir kommt“, spricht Christus Luc. 14, 26., „und hasset nicht seinen Vater, Mutter, Weib, Kinder, Brüder, Schwestern, auch dazu sein eigen Leben, der kann nicht mein Jünger sein.“ Da hebt GOtt das Gesetz auf, daß ich nicht schuldig bin, dem selbigen in der andern Tafel gehorsam zu sein.

Aber im Fall, da die erste und andere Tafel nicht wider einander sind, da bleibt das Gesetz und der Gehorsam bestehen.

47. Wo die Obrigkeit von mir dulden kann, daß ich GOttes Wort lehre, da ehre ich sie und halte eine solche Obrigkeit für meinen Fürsten mit aller Ehrerbietung. Wenn sie aber zu mir sagt: Du sollst GOtt verleugnen, das Wort fahren lassen; alsdann erkenne ich sie nicht mehr für meine Obrigkeit. Auf dieselbige Weise soll man auch den Eltern Gehorsam leisten. Aber wenn sie sagen würden: Ich will haben, daß du ein Mönch oder Meßpfaffe sollst werden nach päpstlicher Abgötterei, da soll man ihnen mit nichten gehorchen. Denn also sagt Mose 5 Mos. 33, 9.: „Wer zu seinem Vater und zu seiner Mutter spricht: Ich sehe ihn nicht; und zu seinem Bruder: Ich kenne ihn nicht; und zu seinem Sohn: Ich weiß nicht; die halten deine Rede und bewahren deinen Bund.“ GOtt will haben, daß wir uns und unser Leben verleugnen sollen in der andern Tafel, wo sie wider die erste ist. Wenn sie aber zugleich mit einander bestehen können, alsdann ist die Ehre, damit man die Eltern ehrt, GOttes Ehre: wo sie aber wider einander sind, so ist vonnöthen, daß man eine Ausnahme mache.

48. Ja, sagst du, daran sagt man gar recht, und ist auch vonnöthen, daß man also lehre, daß man der Obrigkeit, den Eltern und aller ordentlichen Gewalt gehorsam sein soll. Das ist recht geredet, ich gebe die Regel zu. Warum hältst du aber denn die Regel nicht, sprichst du, wo entweder die Eltern oder die Obrigkeit gebieten, daß du ihrer Religion folgen sollst? Antwort: Da ist eine Ausnahme: es soll die erste Tafel der andern vorgezogen werden. Also, wenn die Eltern etwas ordnen oder gebieten, das wider GOtt ist, da wird das vierte Gebot aufgehoben, welches zuvor gewiß und kräftig war. Denn im ersten Gebot wird gesagt, daß man GOtt über alle Dinge lieben und ehren soll; und in den Geschichten der Apostel am 5. Cap. 29. spricht St. Petrus also: „Man muß GOtt mehr gehorchen, denn den Menschen.“

49. Auf solche Weise schließen und erkennen wir auch in der Sache des Evangeliums wider die Herrschaft oder Gewalt des Kaisers und Papstes, ja, wider das ganze gottlose Wesen des Papstes ohne alles vorgehende Recht. Wir haben den Rath nicht verlassen, haben es auch nicht thun können; denn da war kein Richter: wir haben den Papst, unsere Eltern und den Kaiser geehrt. Dieweil aber Christus Joh. 10, 27. sagt: „Meine Schafe hören meine Stimme; einem Fremden aber folgen sie nicht, sondern fliehen von ihm“ etc.: da haben wir nicht warten müssen, bis nach menschlicher Weise erkannt würde, ob wir wohl oder übel daran thäten, wenn wir uns vom Papste absondern. Denn wo man GOttes Willen erkannt hat, da soll man nicht lange disputiren vom Rechte, oder von Herkommen, Brauch oder dergleichen andern Dingen mehr; sondern man soll ohne alles Bedenken dem Gebot und Befehl GOttes gehorchen. Denn es führen weder Papst, noch Eltern, noch auch der Kaiser diesen Titel: „Ich bin der Herr, dein Gott“ etc., 2 Mos. 20, 2.

50. Also kann man auch leichtlich auf diese Frage, die sonst schwer genug ist, antworten von dem, daß Rebekka und Jakob gethan haben. Das Gesetz und die Regel hat bestimmt, daß Esau der Erstgeborne sein sollte. Dieweil aber GOtt in seiner ersten Tafel solches abgewendet, ja, dies Gesetz geändert, und also geschlossen und geboten hat: Esau habe ich nicht wollen haben, Jakob aber wollte ich haben; derhalben haben Rebekka und Jakob nicht gesündigt, sondern heilig und gottseelig gehandelt, und haben mit allen Rechten den Esau beraubt und ihm das Lehengut der Erstgeburt genommen.

51. Sie haben es mit Gewalt nicht können erlangen; denn dasselbe hätte ohne Aufruhr nicht geschehen können. Denn so viel die Regierung und das Regiment belangt hat, da hätte sich die Kirche und das weltliche Regiment, dazu auch das ganze Haus Isaaks wider Jakob gelegt. Denn Esau hat es ihm der Gewalt und des Rechten halben weit zuvor gethan. Jakob war viel geringer, ja, er war gar verachtet und unterdrückt; darum hat Rebekka die Sache fein künstlich und behende, dazu mit einem sehr feinen Betrüge angegriffen; und dasselbe nach GOttes Willen, der auch in solcher wichtigen Sache und die voller großer Gefahr war, einen glückseligen Ausgang gegeben hat, daß es alles wohl gerathen ist.

52. Das hebräische Wort mathamim heißt so viel als schmecken, und reimt sich mit der Art und Weise zu reden, die Lucas in der Apostel Geschichte am 10. Cap. V. 10. auch gebraucht, da er von Petro sagt: „Und als er hungrig ward, wollte er anbeißen“ (voluit gustare cibum), das ist, er wollte zu Mittage essen. Damit wird angezeigt die Mäßigkeit, so sie mit Essen und Trinken gehalten haben, daß es mehr ein Anbiß sei, denn daß sie sollten gefressen und gesoffen haben, auf daß der Leib damit ernährt und erhalten, und nicht beschweret oder beladen werde. Es ist aber eine sehr feine Weise gewesen, das Sehen und den Segen der Erstgeburt zu empfangen, nämlich, da der, welcher hat sollen gesegnet werden, dem Vater, der den Segen hat geben sollen, ein gut Essen hat müssen bringen.

53. Es hat demnach die Mutter nach ihrer Gemahls Befehl gehorsam: „Höre, mein Sohn, meine Stimme“, zum Trost des Sohnes und daß sie ihm Muth möchte machen, diese schwere und wichtige Sache auszurichten. Sie hat vor sich die göttliche Antwort, darum fährt sie stracks fort im Glauben. Und Mose hat sonderlich das Wort gebraucht, daß die Mutter dem Sohne solches befohlen habe zu thun, auf daß er damit anzeigen möchte, daß die Mutter bemüht habe, den Sohn zu überreden, daß er dies hat thun sollen. Denn er weigert sich, es zu thun, und disputirt mit der Mutter über das Gebot.

\section[Zweiter Theil]{Zweiter Theil. Wie Jakob sich anfänglich weigert, in der Rebekka Rath zu willigen, wie ihn aber Rebekka endlich dahin bringt, daß er ihr folgt}

\begin{quote}
\textbf{V. 11—14. Jakob aber sprach zu seiner Mutter Rebekka: Siehe, mein Bruder Esau ist rauh, und ich glatt; so möchte vielleicht mein Vater mich begreifen, und würde vor ihm geachtet, als ob ich ihn betrügen wollte, und brächte über mich einen Fluch, und nicht einen Segen. Da sprach seine Mutter zu ihm: Der Fluch sei auf mir, mein Sohn; gehorche nur meiner Stimme, gehe und hole mir. Da ging er hin und holte und brachte seiner Mutter. Da machte seine Mutter ein Essen, wie sein Vater gerne hatte.}
\end{quote}

54. Das hebräische Wort chalak heißt auf Deutsch glatt, das keine Haare hat. Dieses ist aber eine wichtige Entschuldigung, damit er das Gebot seiner Mutter umzustoßen versucht. Er
\quellseite{81}
erinnert sie, wie die Sache, damit sie umging, an sich selbst so groß und schwer sei. Als wollte er sagen: Siehe zu, liebe Mutter, daß du mich nicht in große Noth und Gefahr bringest, daß ich über mich selbst einen Fluch und nicht einen Segen bringen möge. Er ist wahrlich ein guter Dialektiker; er hält seiner Mutter ein Argument vor, das stark genug ist, damit er versucht, ihr Herz von ihrem Vornehmen abzuwenden. Wie? wenn mein Vater sagen würde: Du hast gemacht, daß ich mit dem Segen geirrt habe, ich will den Segen widerrufen haben; da werde ich alsdann über mich den Fluch bringen. Du und mein Vater werdet leichtlich eine Entschuldigung finden können; ich aber werde mit diesem meinem Gehorsam und daß ich dir gehorcht habe, in große Noth und Gefahr kommen.

55. Alhier kommt nun eine andere Frage von diesem frommen heiligen Weibe, welche zwar vom Heiligen Geiste getrieben wird und immer in ihrem Glauben fortfährt, thut aber gleichwohl viel, das etwas frevelhaft scheint, und bringt sich und ihren Sohn in vielerlei Gefahr. Es ist zumal ein sehr feiner Rath, daß sie so listig dem größeren Sohne zuvor kommt: sie lockt die Böcklein, hintergehet den Vater mit einem großen Betrüge, welches schier über den Verstand eines Weibes ist; verhüllt dem Sohne die Hände, bedeckt ihm seine glatte Haut mit den Bockfellen. Aber das ist zumal ein närrischer Gedanke, daß sie nicht merkt, daß der Vater seinen Sohn auf eine andere Weise sollte erkennen können, denn mit den Augen, die ihm nun dunkel waren. Sie sieht das große Fenster nicht, nämlich die Stimme, welche den Menschen am allermeisten zu erkennen gibt, obwohl die andern Sinne davon nicht urtheilen können. Und sie kann seinerlei Weise verhehlet oder verändert werden. Ja, das noch mehr ist, es ist das Urtheil der Stimme noch etwas gewisser, als wenn du auch das Angesicht selbst anschaust.

56. Darum würde Rebekka ihren Sohn in die größte Noth und Gefahr gebracht haben, wo nicht GOtt dies Ding alles wunderbärlich regiert hätte. Denn höre, was der Vater antwortet. Da Jakob sagt: „Ich bin Esau, dein erstgeborner Sohn“, da erkennt er alsobald die Stimme Jakobs und spricht: „Die Stimme ist Jakobs Stimme, aber die Hände sind Esaus Hände.“ Da wird er ohne Zweifel erschrocken sein und gedacht haben, er habe es nicht alles recht bedacht und erwogen. Ich wäre zwar bald davon gelaufen, wenn ich die Antwort des Vaters gehört hätte: „Die Stimme ist Jakobs Stimme“ etc.; denn dieser Handel ist voller Gefahr gewesen.

57. Dies gehört aber alles zum Trost der Heiligen oder Gläubigen, welcher Rath und Vornehmen glücklich sein muß, dieweil sie es im Glauben anfangen. Denn wer da glaubet, dem muß Alles zum Besten dienen, Röm. 8, 28. Der Glaube ruht nicht, und ist ihm Alles möglich, auch einen heiligen Mann zu betrüben und etlichermaßen auch wohl den Heiligen Geist. Also fährt Rebekka immer fort, und denkt daran nicht, wie sie sich und ihren Sohn Jakob in große Gefahr möchte bringen. Denn Isaak hätte sagen mögen: Warum willst du deinen Vater betrügen und deinen Bruder berauben? Meinst du, daß ich die Stimme nicht höre und erkenne? Da ist es wahrlich in der Wage gestanden. Durch solche Worte hätte Jakob mögen erschrocken sein und denken: O liebe Mutter, was hast du gethan? der Vater hat Acht auf die Stimme, welches zwar uns nicht ist eingefallen, daß wir daran gedacht und uns vor dem Wahrzeichen hätten hüten mögen.

58. Also ist Jakob in größter Angst gewesen: aber er fährt doch fort, und gehorcht seiner Mutter, die ihm gesagt hatte: „Der Fluch sei auf mir, mein Sohn.“ Darum ist es ein großer Glaube gewesen, damit sie sich wider Jakob gelegt, da er sich geweigert und beschweret hat zu thun, was sie ihm befohlen hatte. Lieber Sohn, sagt sie, disputire nicht lange; ich befehle und gebiete dir, daß du thust, was ich gesagt habe. Wo daraus etwas Böses wird entstehen, will ich es lassen über mich gehen; es sei gleich ein Fluch oder Segen, fahre du nur immer fort in meinem Namen, auf mein Heißen und Befehl. Wo du wirst verflucht werden, so will ich den Fluch tragen.

59. Also kann niemand reden, es sei denn, daß er voll Glaubens und Heiligkeit Geistes sei. Denn daß Einer seinen Sohn also soll lassen hinziehen, und in gewisse Gefahr bringen, daß er möchte verflucht werden, und dennoch dabei also glauben, wo es übel gerathen würde, daß er allen Schaden tragen und den Fluch selbst auch über sich gehen lassen wollte, das kann nicht bald ein Jeder thun, sondern ist ein Werk eines sehr brünstigen Glaubens. Darnach wird sie auch gebetet haben: O mein lieber himmlischer Vater, gib, daß ja mein Sohn Alles möge glücklich und wohl ausrichten!

60. Daran hat sie aber nicht gedacht, daß der Vater die Stimme kennen würde. Darum, wer da lügen will, der muß vorsichtig sein, muß ein sehr gut Gedächtnis haben, muß flug und listig sein. Ich würde in solcher Sache demmaßen nichts ausrichten können, ich würde bald Alles verderben und umstoßen. Und Rebekka ist gleichwohl auch nicht vorsichtig genug gewesen: daß sie am allermeisten hätte sollen verhüten, daran hat sie am wenigsten gedacht, also daß schier die Sache dahin gerathen ist, daß es sich hat lassen ansehen, als wäre es gar verdorben und als würden sie verjagen müssen.

61. Aber der Glaube behält immer den Sieg, und überwindet alle Fehler und Gefahr, so den Heiligen begegnen. Wie Rebekka geglaubt und gebetet, und wie Jakob gehorcht, dazu wie endlich die Antwort auch gelautet hat: also hat die Sache auch einen Ausgang genommen. Denn der Glaube macht die Hörenden taub und die Sehenden blind, und wiederum auch die Tauben hörend und die Blinden sehend. Kurz, er verfehlt und verderbt nichts, oder wo er ja bisweilen aus Thorheit oder Unvorsichtigkeit sich stößt oder fehlet, so pflegt doch GOtt solches vorher zu sehen und zu bessern, auf daß die Fehler zugedeckt werden und endlich dennoch einen glückseligen guten Ausgang gewinnen mögen.

62. Ich weiß, daß ich auch oftmals viel Dinges närrisch und fast unvorsichtig ausgerichtet habe, so gar, daß ich wohl gedacht habe: Warum hat mich GOtt zum Predigtamte berufen, dieweil zwar an mir so viel Kunst, so viel Rathes und Verstandes oder Klugheit nicht ist, als zu dem großen Amte gehört? Denn wiewohl ich Alles aus gutem und einfältigem Herzen, dazu auch mit Fleiß und gutem Eifer ausgerichtet habe, so ist dennoch zum öftern viel ungereimtes Dinges und etliche Fehler vorgefallen und mit untergelaufen, daß es sich ansehen ließ, als wollte der Himmel und die ganze Welt einfallen. Da mußte ich auf meine Kniee fallen, und Hülfe und Rath von GOtt begehren, der da mächtig ist, und machen kann, daß auch ein harter greulicher Anstoß oder Unfall einen fröhlichen Ausgang gewinnen muß, wenn wir schlafen. Wie er Eva geschaffen hat, dieweil Adam schlief. Denn da nimmt er Adam im Schlafe seiner Ribben eine und schließt die Stätte zu mit Fleisch, und baut ein Weib aus der Ribbe, 1 Mos. 2, 21. 22. Da möchte Einer sagen, daß GOtt selten Finger hätte, da er ein solch groß Werk so leicht und ohne Mühe ausrichten kann.

63. Gleich nun auf dieselbige Weise regiert er seine Heiligen auch. Wiewohl sie etwa schwerlich geirret oder gefehlet, und unvorsichtig eine große Thorheit begangen haben, davon unzählig viel Unglück kommen möchte, so kann er doch diesem allen ein fröhlich und glückselig Ende machen. So kräftig ist der Glaube und das Gebet, ja, der Glaube ist in Wahrheit allmächtig.

64. Derhalben wird Isaak gleichsam wissentlich, vorsätzlich und zusehens betrogen; denn er sagt selbst: „Die Stimme ist Jakobs Stimme, aber die Hände sind Esaus Hände“, und läßt sich doch gleichwohl betrügen. Wie in den Büchern der Heiden zu sehen, daß die Herren von den Knechten betrogen werden; wie Jener sagt in der Komödie: Wenn ich nicht so gar unverständig oder grob wäre wie ein Stein oder Klotz, so könnte ich den Betrug merken. Also, sage ich, handelt GOtt mit seinen Heiligen oder Gläubigen, daß sie nicht einmal merken oder verstehen können, wo irgend eine närrische oder unvorsichtige That geschehen ist.

65. Darum sollen wir an dies Exempel fleißig denken, dadurch wir ermahnt werden, daß der Glaube solche große Kraft habe, daß er die, so auf das schärfste sehen und hören können, blind und taub macht, daß sie hörend und sehend betrogen werden.

66. Und also mag man auf die Frage antworten, daß Rebekka wohl unvorsichtig gehandelt habe; es hat sie aber die göttliche Vorsichtigkeit regiert um ihres trefflichen großen Glaubens willen, also daß wiewohl Isaak, der Vater, die Stimme Jakobs gehört, er doch, weil er so ein groß Vertrauen gehabt und so gar sicher gewesen, betrogen worden ist.

67. Ein ungläubiger Mensch wird nicht leichtlich betrogen, aber wer leichtlich einem Jeden traut, der wird auch leichtlich betrogen; wie Salomo sagt in seinen Sprüchen Cap. 14, 15.: „Ein Alberner glaubt alles, aber ein Witziger merkt auf seinen Gang.“ Derhalben muß man die, welche man betrügen will, nicht angreifen zu der Zeit, wenn wir wissen, daß sie wachen und auf ihr Ding Achtung haben; sondern man muß erstlich alles Mißtrauen von ihnen wenden, und muß es dahin zu bringen suchen, daß sie uns trauen. Man muß die Herzen immer hingehalten mit Hoffnung, und sie überreden, daß sie sich zu uns nichts Böses zu versehen haben: wie die Knechte Davus oder Cyrus in den Komödien ihre Herren so fein listig wissen zu betrügen, daß sie die Augen zuthun, schlafen und ihnen vertrauen. Denn also zweifelt Isaak an seines Weibes Rebekka und an seines Sohnes Jakob Falsch Treue gar nicht, versieht sich zu ihnen keiner Untreue: er denkt, sie seien so fromm und getreu, daß sie nicht listig oder betrüglich wider ihn handeln können. Darum da er auch schon erinnert wird durch die Stimme Jakobs, hat er doch keinen bösen Argwohn auf ihn haben können, sondern gemeint, Esau würde irgend durch einen Scherz die Stimme geändert haben, nicht den Vater damit zu betrügen.

\quellseite{82}
68. Und dies ist also die gemeine Weise des ganzen Lebens, nicht allein in den Heiligen, sondern auch im gemeinen Leben bei allen Menschen, wie das die Exempel in den Komödien bezeugen. Wer leichtlich glaubt, der wird betrogen; wiederum aber, wer nicht leichtlich glaubt, da ist es vergeblich, daß du versuchen willst, ihn zu betrügen. Darum geschieht aber des Dinges viel im Leben der Heiligen, daß wir sehen sollen, wie Gott seine Heiligen regiere und aus ihrer Thorheit und Verwegenheit die allergrößte Weisheit machen kann; dazu gibt er zu ihrem Rath und Anschlägen, die sich anfänglich sehr närrisch ansehen lassen, Glück, daß sie einen schönen Ausgang gewinnen müssen. Es kommt oftmals vor, daß heilige Menschen unvorsichtig genug handeln, wie die Rebekka hier auch schmählich fehlgreift, aber sie bricht dennoch hindurch, dieweil sie gläubig ist.

69. Derhalben soll man dies im Anfang dieses Capitels fleißig merken, nämlich, zum Ersten, daß die erste Tafel der andern soll vorgezogen werden; darnach, daß Gott seinen Heiligen ihre Thorheit und Verwegenheit verzeiht, ja, er läßt auch dieselbige wohl gerathen und gibt ihnen dazu gut Glück. Derhalben, wo du ja etwas thust, das an sich selbst nicht gut ist, es sei in welcherlei Stande es wolle, so verzage darum nicht, sondern erkenne deinen Fehler oder Irrthum mit aller Demuth, und denke also: Gott ist allmächtig, der hat Rebekka und Jakob aus solcher großen Noth und Gefahr führen und erretten können, darum will ich auch nicht verzagen, sondern will ihm vertrauen, er werde mich aus diesem Unfall auch erlösen. Der Rath und Anschlag der Rebekka war zwar fast verwegen, ist aber doch wohl gerathen; denn Gott macht den Rath der Gottlosen zunichte, aber er ehrt und hilft seinen Heiligen.

70. Darum soll man sich vor beiderlei hüten, nämlich, daß du nicht hoffärtig und stolz werdest von wegen deiner Weisheit und klugen Rathes, und darnach wo etwa ein Rath oder Anschlag nicht wohl gerathen will, daß du daselbst gleichwohl nicht verzagest. Denn Gott entschuldigt und begünstigt die Irrthümer der Frommen. Ich habe oftmals die größten Thorheiten und Verwegenheiten begangen, aber ich habe solches aus gutem Herzen gethan, nicht daß ich damit jemand hätte wollen Schaden thun; habe es auch nicht wissentlich, sondern unwissentlich und aus Thorheit gethan; habe damit begehrt, Andern treulich zu rathen und zu helfen. Da habe ich wahrlich beten müssen, daß Gott meinen Irrthum und Fehler zudecken und verbessern wolle. Und in großen Händeln, die etwas verworren sind, kann es nicht anders zugehen. Denn es geschieht oft, daß große und sehr fromme Männer im gemeinen Regiment großen Schaden thun mit ihrem närrischen und verwegenen Rath und Anschlägen. Und wenn sich Gott ihrer nicht erbarmete, so würde allenthalben das Unterste nach oben und das Oberste herunter gekehrt, und wohl Alles über einen Haufen fallen.

71. Also steht es um unser Leben, welches wahrlich ein arm elend Leben ist. Da wir uns vermessen, sehr weise zu sein, und den Sachen sehr wohl rathen und helfen wollen, da richten wir oft den größten Schaden an; so, daß wo nicht solche unsere Fehler durch göttliche Gnade und Vorsehung gebessert würden, Alles gründlich verderben würde. Denn es widerfährt uns eben dasselbe, das jenem Bauern widerfahren ist, der seinen Wagen hat schwerlich vorwärtsbringen können deswegen (wie er sich dünken ließ), weil die Räder am Wagen etwas zu breit waren: derselbige Bauer hat, damit die Pferde den Wagen leichter ziehen könnten, die Reifen der Räder abgespitzt; aber mit diesem Anschlag richtete er so wenig aus, daß der Wagen mit den Rädern nur desto tiefer in den Koth einsank und gar nicht von der Stelle gebracht werden konnte.

72. Also geht es mit unserer Weisheit auch, deren wir uns anmaßen, nach welcher wir uns dünken lassen, daß auf Erden kein leichter Ding sei, denn das gemeine Regiment zu verwalten. Wie jener Narr in der Komödie sagt: Me regem esse oportuit: Ich sollte billig ein König sein, und wie die Poeten solche Vermessenheit der Menschen in der Fabel von Phaethon fein abgemalt haben; denn solcher Leute hörst du jetziger Zeit auch viele, die da schreien können: Wenn ich an D. Martinus oder Philippus Statt wäre, ich wollte jenen wohl einen besseren Rath geben; ja, wahrlich, du würdest es recht treffen etc. Also sehen wir, daß oftmals große Fürsten anlaufen und fehlen, und daß die allerbesten und frömmsten Bischöfe sich auch zum öftern am allerärgsten stellen.

73. Wie nun? Soll man denn gar nichts thun und soll man ganz und gar vor allem Regimente fliehen? Antwort: Mit nichten, es soll vielmehr ein Jeder sein Amt fleißig und treulich ausrichten, das ihm von Gott befohlen ist; davor aber soll er sich hüten, daß er sich auf seine Kräfte oder Weisheit nicht verlasse, und daß er es dafür nicht halte, daß er so ein großer Mann sei, mit deß Rathe Alles müßte regieret werden. Denn solchem Frevel und Vermessenheit ist nicht zu helfen, und ist verdammlich, wenn ich mir das zuschreibe, daß ich ein solcher und so ein vortrefflicher Mann sei, daß ich allenthalben weislich, recht und wohl das gemeine Regiment verwalten, dazu auch ein Haus und die Kirche wohl regieren könne. Bist du aber ein Richter, ein Bischof, oder ein Fürst, so schäme dich nicht, auf deine Kniee zu fallen und zu sagen: Lieber Herre Gott, du hast mich zu einem Fürsten, Richter, Bischof, Hausvater, Pfarrherrn oder Kirchendiener gesetzt, darum regiere und lehre du mich, gib mir Rath, Weisheit, Stärke und Kraft, daß ich mein befohlen Amt fleißig und wohl ausrichten möge.

74. Also that Salomo auch, da er im Traume des Nachts vom Herrn gehört hatte, daß er zu ihm sagte, 1 Kön. 3, 5.: „Bitte, was ich dir geben soll.“ Da bittet er also V. 7. ff.: „HErr, mein Gott, du hast mir ein groß mächtig Volk gegeben, gib mir nun auch Weisheit, daß dein Knecht dein Volk richten möge, und verstehen, was gut und böse ist.“ Und die Schrift zeiget an, daß solch Gebet dem HErrn sehr wohl gefallen habe; denn es ist ein sehr schön Bekenntniß und Beichte. Er war ein sehr weiser und gottseeliger König, und nennt sich dennoch einen kleinen Knaben; denn er sagt also V. 7.: „Ich bin ein kleiner Knabe, weiß nicht weder meinen Ausgang noch Eingang.“ Darum wird daselbst im Texte V. 10. 11. 12. hinzugefügt: „Das gefiel dem HErrn wohl, daß Salomo um ein solches bat.“ Und Gott sprach zu ihm: „Weil du solches bittest, und bittest nicht um langes Leben, noch um Reichthum, noch um deiner Feinde Seele, sondern um Verstand, Gericht zu hören; siehe, so habe ich gethan nach deinen Worten. Siehe, ich habe dir ein weises und verständiges Herz gegeben, daß deines gleichen vor dir nicht gewesen ist, und nach dir nicht aufkommen wird“ etc.

75. Derhalben soll ein Jeder insonderheit lernen, daß er mit aller Demuth seine Schwachheit und Unvermögen erkenne und Gott um Weisheit und Rath bitte. Denn die Leute werden nicht darum zum Regimente berufen, dasselbe zu verwalten, daß sie sich eine vollkommene Erkenntniß aller Dinge anmaßen sollen, sondern daß sie dadurch sollen unterrichtet werden und lernen, was Gott sei, was er ausrichte durch die Obrigkeit und Regenten, welches Werkzeuge sind der Werke Gottes, dadurch Gott die Menschen regiert. Da werden sie dann recht verständig und haben Glück in ihrem Regimente; wo sie aber ihrem eigenen Rathe und ihren Gedanken folgen, so können sie nichts mehr ausrichten, sondern verwirren und mengen Alles unter einander.

76. Derhalben, sage ich, soll man zum Gebet fliehen, und Gott darin vortragen, wie groß und schwer das Amt sei, so uns befohlen ist, und soll sagen: „Vater unser, der du bist im Himmel“ etc., „gib mir die Weisheit, die stets um deinen Thron ist“ etc., wie im Buch der Weisheit Cap. 9, 4. steht. Vornehmlich aber soll jeder Pfarrherr oder Kirchendiener also beten: HErr Gott, du hast mich in der Kirche zu einem Bischof und Pfarrherrn gesetzt; Du siehst, wie ich so ungeschickt bin, solch großes und schweres Amt recht auszurichten, und wo es ohne deinen Rath gewesen wäre, so hätte ich schon vorlängst Alles miteinander verderbet. Darum rufe ich dich an. Ich will zwar gerne meinen Mund und mein Herz dazu leihen und neigen; ich will das Volk lehren; ich will selber auch immer lernen, und mit deinem Worte umgehen und demselben fleißig nachdenken; brauche du mich als deines Werkzeuges. Lieber HErr, verlaß du mich nur nicht; denn wo ich werde allein sein, so werde ich leichtlich Alles miteinander verderben.

\quellseite{83}
77. Da thun nun aber die Secten und Rottengeister gar das Gegenteil. Denn dieselbigen schreiben sich selbst große Weisheit zu und daß sie sehr wohl regieren und lehren können; darum brechen sie verwegen in die Kirche ein. Sie beten nicht, glauben dazu auch nicht, daß es eine Gabe GOttes sei, das Kirchen- oder Weltregiment zu verwalten; sondern drängen sich selbst hinein als große treffliche Meister und Regenten. Darum kommt es endlich dahin, daß sie Alles verwirren und verhindern, so von Andern nützlich und wohl gebaut und angerichtet ist.

78. Eben dasselbe begegnet im Weltregiment denen, die sich auf ihren eigenen Rath und Klugheit verlassen und darauf zum Regiment greifen. Derhalben, wo du im Regiment bist, so bedenke und hüte dich, daß du dich auf deine Weisheit nicht verlässest; hüte dich, sage ich, vor dir selbst und bete mit gefalteten Händen heimlich also: Du, mein lieber himmlischer Vater, stehe mir bei, hilf mir, regiere und führe du mich. Wie Mose die ganze Sache, das Volk aus Egypten zu führen, auf GOtt legt, und weigert sich den göttlichen Befehl auszurichten, wo nicht GOtt selbst dabei sein, und sie führen und regieren werde. Denn er sagt also, 2 Mos. 33, 15.: „Wo nicht dein Angesicht gehe, so führe uns nicht von dannen hinauf“ etc. Wenn du nun auf solche Weise gebet hast, so fahre nur in deinem Amt fort, habe einen starken unverzagten Muth, und zweifle daran gar nicht, deine Sache wird wohl vor sich gehen, es wird Alles wohl gerathen und einen glücklichen Ausgang nehmen.

79. Also habe ich auch oftmals, wie man im deutschen Sprichwort sagt, die Schlüssel dem HErrn vor die Füße geworfen, das ist, ich habe mein befohlen Amt auf ihn gelegt, darum ist es auch viel anders gerathen, denn ich selbst zuvor gesehen oder gedacht hatte. Und wo ich irgend gefehlet hatte, da hat er es nach seiner Güte und Weisheit gebessert.

80. Auf dieselbige Weise sind auch Vater und Mutter im Hausregiment Werkzeuge, dadurch das Haus und Hauswesen regieret wird. Sie sollen aber auch erkennen, daß sie mit ihren Kräften und durch ihren Fleiß oder Arbeit die Kinder nimmermehr recht und wohl aufziehen können. Derhalben sollen sie zu GOtt schreien und also beten: HErr GOtt, himmlischer Vater, hilf uns, daß die Kinder wohl gerathen mögen! Gib, daß das Weib in Zucht und Ehrbarkeit lebe und daß sie in Erkenntniß und Furcht GOttes beharren möge.

81. Du wirst aber deren Viele finden, die diese oberste Gewalt und Weisheit im Regiment nicht erkennen, sondern wo irgend etliche Hindernisse vorfallen, da meinen sie, daß sie es am besten ausrichten und zurechtbringen wollen, wenn sie mit der Strafe desto größeren Ernst anwenden, daß die Unterthanen durch Furcht der Strafe gezwungen und auch wider ihren Willen zum Gehorsam getrieben werden. Nun ist es wahr, es thut Zucht und zwar strenge Zucht in diesen unsern Zeiten sehr noth; das aber gar zu gewiß, daß du ohne das Gebet nimmermehr etwas ausrichten wirst. Denn das Regiment ist eine göttliche Kraft und darum nennt auch GOtt alle Obrigkeit Götter, nicht um der Schöpfung willen, sondern um des Regiments willen, welches allein GOtt gebühret. Derhalben, wer im Regiment ist, der ist gleichsam ein Gott im Fleisch.

82. Wo aber die Leute aus Frevel und gleichsam mit ungewachsenen Händen, wie man sagt, das ist, ganz ungeschickt und unvorbereitet in das Regiment fallen, es sei im Kirchen- oder Weltregiment, oder auch wohl in der Haushaltung, und GOtt ausschließen, nicht beten und GOtt nicht um Rath fragen, sondern alles durch ihren eigenen Rath und Kräfte regieren und ausrichten wollen: da wird es endlich in der Haushaltung also kommen, daß aus einem ehrbaren und züchtigen Weibe die ärgste Hure werden wird, daß die Kinder übel geraten und dem Henker in die Hände kommen. Im Weltregiment wird die Gemeinde verwirret durch Aufruhr, Krieg, Mord und viel andere unzählige Gefahr. In der Kirche aber entstehen Ketzereien, es wird eine epicurische Verachtung des Wortes; item, daß die heiligen Sacramente gelästert und schändlich gemißbraucht werden etc. Ei Lieber, warum? Antwort: Dieweil ein solcher Hausvater, Fürst oder Pfarrherr nicht erkennen will, daß GOtt allein der sei, von dem Rath und Regiment herkommt; sondern mit seiner Vermessenheit und daß er von sich selbst so viel hält, verderbt er sich und Andere, welchen er vorsteht.

83. Dies, habe ich gedacht, bei diesem Text zum Exempel und zur Lehre sagen zu müssen, als nämlich zum Ersten, daß wir wissen, daß die erste Tafel die andere regieren soll nach dem Spruch im ersten Gebot: „Ich bin der HErr, dein GOtt“; darnach, daß wir GOttes Gnade und Barmherzigkeit erkennen lernen, der seine Heiligen wunderlich führt auch in ihrer Verwegenheit und Thorheit, auf daß wir GOtt von ganzem Herzen vertrauen, und lernen, daß der Glaube Alles verschlingt, nicht allein die Sünde, sondern auch die Vermessenheit und Thorheit der Heiligen, und wo sie etwas unbedachtsam vorgenommen haben. Das ist die Sache, um die es sich in diesem Capitel handelt. Nun wollen wir den Text nacheinander lesen, und wollen dieser wunderbaren Historie zusehen.

\section[Dritter Theil]{Dritter Theil. Wie Rebekka den Jakob zubereitet, und Jakob also zu Isaac geht, den Segen zu erhalten, und von ihm nicht erkannt wird}

\begin{quote}
\textbf{V. 15. Und (Rebekka) nahm Esau's, ihres größern Sohns, köstliche Kleider, die sie bei sich im Hause hatte, und zog sie Jakob an, ihrem kleinern Sohn.}
\end{quote}

84. Mose setzt nun die andere Ceremonie hinzu, nämlich die mit der schönen und köstlichen Kleidung. Denn so viel heißt das hebräische Wort chamudoth, auf Lateinisch desiderabilis, pretiosus, das köstlich ist, und wird dies Wort oftmals in der heiligen Schrift gebraucht. Als, da der Engel zu Daniel sagt: Es vir chamudoth, auf Lateinisch, vir desideriorum, das ist, ein lieber und werther Mann, Dan. 9, 23. Und im Evangelium Lucä Cap. 7, 4. wird gesagt vom Hauptmann zu Capernaum: Erat ei servus chamudoth, auf Lateinisch: carus et pretiosus, das ist: Er hatte einen Knecht, den er lieb und werth hielt. Also hat Rebekka auch hervorgebracht die lieben, schönen, herrlichen Kleider, die priesterlichen Kleider, die ohne Zweifel sehr schön gewesen sind, wie sie im Mose beschrieben werden. Denn sie haben zum Opfer und Gottesdienst sonderlicher Kleider gebraucht, wie wir in unsern Kirchen auch thun, welches wir den Ornat zu nennen pflegen; und dieselben Kleider hat Rebekka bei sich im Hause.

85. Daraus denn erhellt, daß Esau noch nicht ist bestätigt gewesen, wiewohl er den Besitz gehabt und die Kirche oder Gemeinde regieret hat; er ist Licentiat gewesen zum Pabstthum und hat des Priesteramtes gepflegt, ist aber nicht confirmirt oder bestätigt gewesen, und hat die Lehen noch nicht empfangen. Darum hat er die Kleider in seiner Gewalt nicht; welches ohne Zweifel durch sonderlichen Rath Rebekka behütet und ihm zuvorgekommen ist. Denn sie hat die Kleider an sich genommen und hat sie behalten, wiewohl es dem Vater Isaac gebührte, sie zu haben. Denn sie hat auf die Gelegenheit gewartet, daß sie die Kleider dem geben möchte, der die Lehen vom Vater empfangen und zum Erben sollte bestätigt werden.

\begin{quote}
\textbf{V. 16–18. Aber die Felle von den Böcklein that sie ihm um seine Hände, und wo er glatt war am Halse. Und gab also das Essen mit Brod, wie sie es gemacht hatte, in Jakobs Hand, ihres Sohns. Und er ging hinein zu seinem Vater, und sprach: „Mein Vater!“ Er antwortete: „Sie bin ich. Wer bist du, mein Sohn?“}
\end{quote}

86. Es ist ein Wunder, daß die Mutter Jakob nur am Halse mit den Fellen bedeckt und nicht auch über das ganze Angesicht. Sie hat nur daran gedacht, daß er dem Vater würde müssen die Hände reichen, und wenn er dem Vater hätte sollen einen Kuß geben, so würde wiederum der Vater ihn um den Hals nehmen, und ihn also umfangen oder herzen. Darum hat sie diese beiden Theile vornehmlich wollen bedeckt haben. Sie meint aber, man könne das glatte Angesicht so leichtlich nicht merken, dieweil Isaac einen Bart gehabt und den Fuß auf den Boden genommen hat.

\quellseite{84}
87. Vielleicht mag sie ihn dessen auch erinnert haben, während sie ihn mit der Kleidung geschmückt hat, daß er, so gut er könnte, die Stimme Esau nachahmen sollte, auf daß der Vater keinen Unterschied machen könnte zwischen der Stimme Esaus und Jakobs; wie er ihn bald fragt: „Wer bist du, mein Sohn?“ Bist du Esau, den ich auf die Jagd gesandt habe? Mich dünkt, ich höre die Stimme Jakobs etc.

\begin{quote}
\textbf{V. 19. 20. Jakob sprach zu seinem Vater: Ich bin Esau, dein erstgeborner Sohn; ich habe gethan, wie du mir gesagt hast; stehe auf, sitze dich, und iß von meinem Wildpret, auf daß mich deine Seele segne. Isaak aber sprach zu seinem Sohn: Mein Sohn, wie hast du so bald funden? Er antwortete: Der HErr, dein GOtt, bescherte mirs.}
\end{quote}

88. Jakob hält eine lange Rede, damit er seinem Vater antwortet, und dieselbe hört der Vater mit Fleiß, und wird endlich daraus schließen, daß hernach folgt, da er sagt: „Die Stimme ist Jakobs Stimme, aber die Hände sind Esaus Hände.“ Das ist aber eine offenbarbare Lüge, daß Jakob sagt, daß er der Erstgeborne sei. Darnach hat er auch etwas unvorsichtig geeilt, denn die Mutter hat das nicht vorhergesehen oder bedacht, daß solch Eilen bei dem Vater würde verdächtig sein. Denn er sagt: „Wie hast du so bald funden?“ Dies ist ein Anstoß, der sehr hart ist, und er wird dadurch des Betruges erinnert, damit sie umgegangen sind. Wie kommst du so bald von der Jagd, sagt er, hast du so bald jagen und Wildpret fangen können? Wenn du gegen den Abend oder Morgen wieder gekommen wärest, so wärest du bald genug wieder gekommen. Jakob antwortete: „GOtt bescherte mirs.“ Es ist wahr, GOtt muß es auf sich nehmen. Er lügt gewaltig; sagt, er habe Wildpret gefangen, so er doch im Stalle zwei Ziegenböcklein genommen hatte. Aber Isaak verbeißt es, und thut, als wisse er davon nichts; denn er ist sicher und voll Vertrauens, nämlich, daß er von seinem Sohne oder von der Mutter nicht könne betrogen werden.

\begin{quote}
\textbf{V. 21. 22. Da sprach Isaak zu Jakob: Tritt herzu, mein Sohn, daß ich dich begreife, ob du seist mein Sohn Esau, oder nicht. Also trat Jakob zu seinem Vater; und da er ihn begriffen hatte, sprach er: Die Stimme ist Jakobs Stimme, aber die Hände sind Esaus Hände.}
\end{quote}

89. Dies ist ein anderer und viel härterer Anstoß; da will es böse werden. Der Vater heißt ihn herzutreten, daß er ihn begreifen möge. Da wird er heiß geschwitzt haben und hat stillschweigend der Mutter Rath und Anschlag gefraget. O meine liebe Mutter, wird er heimlich gesagt haben, was hast du gethan? Wozu hast du mich gebracht? Da er nun hinzugetreten ist, hört er, daß der Vater sagt: „Die Stimme ist Jakobs Stimme.“ Ach die Stimme, die Stimme! Da hätte ich die Schüssel fallen lassen, und wäre gelaufen, als hätte mir der Kopf gebrannt. Denn Isaak ist da auf dem rechten Wege, den Betrug zu merken, da er die Stimme Esaus und Jakobs unterscheidet, und es hätte ihn dieser Anstoß wohl aufwecken sollen.

90. Aber er fällt wieder in den Schlaf, und meint, Esau stelle sich also, daß er seines Bruders Stimme nachahme. Jakob tritt getrost hinzu, und hört, daß der Vater die Worte redet. Derhalben hört nun der Rath Rebekkas und Jakobs da auf, ja, er läuft damit an. Aber GOtt macht, daß Isaak auf sein Amt beflißen ist, daß er ihn segnen soll, und daß er diesen Anstoß aus dem Herzen fahren läßt.

91. Nun muß man aber dies alles auf unsere Lehre ziehen, daß wir daraus gelehret und unterrichtet werden mögen, wie kräftig der Glaube sei, und daß dem Gläubigen alle Dinge möglich sind. Denn der Glaube macht aus dem, das nichts ist, daß es sei; und aus den Dingen, so unmöglich sind, macht er Alles möglich. Du hast gesehen die Thorheit und Verwegenheit der Rebekka, wie sie ihren Sohn Jakob gezieret und gerüstet hat, daß er den Vater betrügen möchte: sie hat ihm die Felle von den Böcklein um seine Hände und um den Hals, daß der Vater meinen sollte, er wäre in Wahrheit recht rauh.

92. Dies ist ja ein gefährlich Ding; wie Jakob auch zittert und darob erschrickt, will dem Rathe der Mutter anfänglich nicht gerne folgen, und sagt: Es möchte vielleicht mein Vater mich begreifen und den Betrug merken, so würde ich über mich selbst einen Fluch und nicht einen Segen bringen. Aber der Glaube dieses sehr frommen Weibes ist so groß, daß er durch dies alles durchbricht, Jakob ist zum Vater hinein gegangen in der Gestalt seines Sohnes Esau, und hat rauhe Hände, die also mit List und mit sonderlichem Betrug zugerichtet waren, aber er kann gleichwohl den Vater nicht in allen Stücken betrügen; denn er kennt dennoch die Stimme Jakobs. Da siehst du nun, daß GOtt diese Verwegenheit zum Besten wendet durch den Glauben des Weibes, um welches Glaubens willen der Rath und Anschlag, so fast närrisch ist, sehr gut und glücklich ablaufen muß.

93. Also beschreibt nun die heilige Schrift die Historien der Heiligen, daß sie darin die große Kraft des Glaubens rühmt und lobt. Derhalben, wer GOttes Wort hat, der soll sich für selig halten, und soll die Augen von diesen gegenwärtigen und zeitlichen Dingen abwenden und sie auf die zukünftigen unsichtbaren Güter richten. Denn das Wort Gottes und vornehmlich die Verheißung redet nicht von gegenwärtigen, sondern von den zukünftigen Dingen und die niemand erfahren hat. Der Glaube hängt sich straks an das Ding, das noch lauter Nichts ist, und wartet darauf, bis daß daraus Alles werde. Er ist eine Kunst und Weisheit der Finsterniß und der Nichtigkeit, das ist, der Dinge, die er nicht erfahren hat, die man nicht sieht und die schier unmöglich sind. Und wer ein Christ sein will, der muß das in seinem Herzen wohl bedenken und daran fest halten. Denn die andern Künste alle werden gelehrt aus Schlüssen, aus bewährten Gründen und aus der Erfahrung: sie haben aber ihren Grund oder Anfang nicht in dem, das nicht ist, und sonderlich, das man nicht sieht, das unmöglich, ungereimt und närrisch ist. Aber der Glaube, der die Verheißung ergreift und sich daran hält, haftet mit dem Herzen an dem, das gar ungereimt und unmöglich ist, und ist mit dem Wort und der göttlichen Verheißung zufrieden.

94. Denn also thut Jakob auch. Die Mutter heißt ihn, er soll den Segen vom Vater empfangen; da liegt ihm das im Wege, das lauter Nichts ist; denn der Segen gebührte Esau und nicht Jakob. Es lag ihm, sage ich, im Wege, daß die Dinge gar unmöglich waren. Denn er hat müssen denken: Wie sollte ich können gesegnet werden? Es gebührt ja der Segen meinem Bruder. Ueber das, so ist mein Bruder rauh, ich aber bin glatt: es ist die Gefahr dabei, daß ich über mich einen Fluch und nicht einen Segen möchte bringen. Warum ist es unmöglich, so viel die Sache an sich selbst betrifft; denn der Segen ist nicht mein. Der Vater gedenkt Esau zu segnen, bekümmert sich in dem Falle um mich gar nichts. Darnach ist es wohl möglich, daß der Fluch auf mich komme und auf mir bleibe, der ich an der Haut glatt und weich bin, denn der Vater den Betrug erkennen und merken würde. Darum, meine liebe Mutter, werden wir gewißlich anlaufen.

95. Aber wider dies alles sagt die Mutter also: Wo ein Fluch kommen wird, der soll mir Schaden und auf mir ruhen. Das ist ein Weib, und kein schlecht Weib. Es ist ein Mensch von weiblichem Geiste, aber es ist so großer Eifer in ihr, daß sie auch die allerstärksten Männer überwindet. Denn sie denkt, sie könne hindurchbrechen, daß der Segen dem verliehen und gegeben werde, dem er doch nach dem Gesetz und der Regel nicht gebühret hat. Sie glaubt, was einfach unmöglich und ungereimt, närrisch und gefährlich ist. Sie zweifelt daran nicht, Jakob werde müssen gesegnet werden; wiewohl sich Alles dawider gelegt und demselben stracks gar zuwider war.

96. Also ist der Glaube und seine Kraft ein Gedanke des Herzens oder (daß ich es also nenne) eine solche Meinung, damit er sich hängt und fest hält an dem, das nichts, das unmöglich und gar ungereimt ist. Er hat nur das göttliche Wort, das da sagt: „Der Größere wird dem Kleineren dienen.“ Sie aber, als die Mutter, hat verstanden, welcher unter den beiden der Größere wäre. Und ist die Erleuchtung des Heiligen Geistes auch dazu gekommen, daß sie dasselbe Wort recht verstanden hat. Isaak aber dagegen besteht noch immer stracks und fest wie ein unbeweglicher Steinfels, und hängt nur am Gesetz und an der Regel, nämlich, daß Esau der Erstgeborne sei; aber die Mutter läßt sich das alles nicht anfechten.

97. Auf solche Weise sollen wir uns auch in die Sache schicken lernen. Wenn wir das Wort haben, alsdann sollen wir allerlei Disputation wider das Wort fahren lassen und verwerfen, und sollen uns das gar nichts anfechten oder innen lassen, es werde uns gleich vorgeworfen, das entweder närrisch, oder unmöglich, oder endlich auch dem Gesetz der Natur und Moses stracks zuwider ist; sondern wir sollen stracks dem folgen, das Gott gesagt hat. Wo aber die Vernunft dazu kommt und vom Wort disputirt, so ist es verloren. Und daher kommt der Kampf, davon St. Paulus Gal. 5, 17. sagt, daß das Fleisch gelüste wider den Geist, und den Geist wider das Fleisch. Denn da redet Paulus von jenen hohen Affecten, als da ist, daß du glauben sollst, das du nicht verstehst und das nach der Vernunft unmöglich ist.

\quellseite{85}
98. Also muß man im Berufe und endlich auch im Tode thun: wenn mich Gott zum Herrn setzt wider die höllischen Pforten, wider das Toben und grimmigen Zorn der ganzen Welt, wider die Schwachheit meines Fleisches, daß ich da durchbrechen soll, wie soll uns da zu Muthe sein? Denn wer bin ich, der ich streiten soll, erstlich wider so viel Heereshaufen aller Teufel und Menschen; darnach auch wider die Schwachheit meines Fleisches? Antwort: Da muß man auf Gottes Gebot und Befehl sehen.

99. Der Herr sagt: Gehe hin, du sollst die Oberhand behalten und siegen; ich will mit dir sein. Wie er Mose heißt den egyptischen König angreifen und das Volk Israel herauszuführen aus der Tyrannei und dem eisernen Feuerofen. Läßt sich aber das nicht ansehen, daß es eitel närrisch und unmöglich Ding sei, nämlich, ein so groß Volk aus Egypten führen, und es also herauszuführen, daß auch nicht eine Klaue dahinten bliebe? wie Mose zu Pharao spricht 2 Mos. 10, 26. Da hätte ich aber nach der Vernunft gesagt: Du wirst auch nicht eine Klaue herausführen, geschweige denn ein solch groß Volk; denn der König ist mächtig und ein greulicher Tyrann. Aber Mose redet mit Pharao nicht anders, denn als wenn er jetzt schon das große Volk herausgeführt hätte; ob es wohl äußerlich scheint eine Fabel zu sein, ein närrischer Gedanke und gar ein unmöglich Ding, wenn man auf die Vernunft sieht. Denn alles, was der Glaube angibt und das Wort verheißt, das muß gewißlich geschehen; denn Gott ist das Wort und das Wort ist Gott: wer das Wort hat, der hat eine göttliche Kraft, Röm. 1, 16.

100. Es scheint zwar äußerlich nichts Anderes als eine menschliche Stimme, und eine Thorheit des Menschen, der da redet; wie St. Petrus sagt in den Geschichten der Apostel am 10. Capitel §. 26. zu Cornelius: „Ich bin auch ein Mensch“, wie du bist etc. Petrus hat auch Wort und in den Ohren des Cornelius klingt die Stimme eines Menschen, sie thut aber doch allerlei große Wunder.

101. Also sollen wir beherzt und getrost sein wider allerlei Gefahr, so beide vom Türken und andern Feinden kommen kann, ja, wider den Tod selbst. Denn wo wir durch den Mund irgend eines Bruders oder Kirchendieners absolvirt sind, soll man den Menschen, der da redet, nicht ansehen, soll auch die Augen nicht richten auf die Gefahr oder den Tod; man soll nicht auf das Grab, auf die Verwesung, auf die Würmer und Fräule sehen, so uns vorgeworfen werden; denn die Vernunft denkt ohne Glauben also: Ist es auch möglich, daß ein todter Leib leben soll, der da begraben, von Würmern gefressen, oder verbrannt und zu Asche gemacht wird? Dazu sagt sie aber einfach Nein, und spricht: Es ist nicht mehr. Der Geist aber, der da widerstreitet, sagt: Es ist wahr; der ich jetzt zum Tode verordnet bin, lasse mich doch gewißlich dünken, daß ich lebe, und aus dem Grabe wiederum herrlich auferstehen werde und mit einem verklärkten Leibe.

102. Darum, sage ich, ist der Glaube ein allmächtiger Kampf, der wider den Streit des Fleisches stracks in einem und einerlei Gesicht am Wort und an der Verheißung hängt; denn wir reden hier nicht vom historischen Glauben, welcher eine bloße Erkenntniß ist, als, wenn ich weiß, daß David den Goliath und die Philister überwunden hat, 1 Sam. 17, 49. Wiewohl dasselbe ein Exempel des Glaubens ist, so ist es doch mir allein eine Historie, und wo ich den Goliath erlegen und schlagen, dazu den Löwen zerreißen sollte, da würde mir die Historie zu nichts dienen, denn allein, daß sie mir eine Vermahnung wäre, die mich erwecken möchte zu gleicher Stärke. Wenn ich aber also gedächte: David und Simson haben schöne herrliche Siege erlangt, darum werde ich mich auch also halten müssen, daß ich solch Ding möge ausrichten; das wäre mir unmöglich. Denn es hat David auch nicht ohne großen schweren Kampf solche herrliche Thaten gethan; denn er hat auch im Fleisch gestreitet wie wir, er hat ein solch Fleisch gehabt, das ihm immer widersprochen, wider ihn gestritten, ja, ihn gefangen gehalten hat nicht anders, als wir auch haben; er hat viel mehr Unglauben gefühlt denn Glauben. Denn der Glaube scheint und ist auch schwach, dieweil das Fleisch nicht allein widerstreitet, sondern auch gefangen nimmt, wie St. Paulus sagt Röm. 7, 14.

103. Ich habe zwar vor Zeiten anders nicht gewußt, denn daß Paulus nicht gedacht hätte an den Kampf und Streit wider das Fleisch; er sagt aber selbst 2 Cor. 12, 7.: „Mir ist gegeben ein Pfahl ins Fleisch, nämlich des Satans Engel, der mich mit Fäusten schlage.“ Ja, der lehrt mich, was der Glaube sei. Ich sehe wohl die großen herrlichen Siege der Märtyrer und der Heiligen, aber mein Glaube kann solches nicht ausrichten; und sie haben es sich auch blutsauer lassen werden und haben den Tod überwunden nicht ohne großen Kampf. Also haben wir auch dieselbige Taufe, denselbigen Gott, der uns tröstet, den Tröster, den Heiligen Geist, der uns im Wort viel reichen Trost vorhält. Und ich höre zwar die Historie von dem Kampf und Sieg der Heiligen gern: ich kann aber dasselbe nicht thun und ausrichten, das sie gethan haben. Darum muß man immer das Wort ergreifen und sich daran halten; wie Ps. 27, 14. steht: „Harre des Herrn, sei getrost und unverzagt, und harre des Herrn.“

104. Gott verwirft auch den Schwachen im Glauben nicht; ja, die Apostel selbst und die Propheten sind auch im Glauben so gar stark nicht gewesen, vornehmlich wenn sie durch den Glauben haben große Dinge thun sollen. Mose zittert und schreit am Rothen Meer; denn das Fleisch hat wider ihn nicht allein gestritten, sondern ihn auch gefangen genommen. Denn du siehst wohl, wird ihm sein Fleisch eingegeben haben, daß an beiden Seiten Gebirge sind; dahinten ist Egypten, vor uns ist das Meer, dadurch werden wir mit den Kindern und armen Weibern, dazu auch mit unserm Vieh nicht gehen können. Also hat er schier verzagen müssen. Denn das Gesetz in seinem Fleisch und Gliedern hat ihm dies also eingegeben, daß er wird gedacht haben: Siehe, wir sind doch gar ohne alle Hilfe; was habe ich gethan, daß ich das Volk herausgeführet habe? Wir werden alle sterben müssen. Da fängt dann an das unaussprechliche Seufzen, wie es Paulus nennt Röm. 8, 26. Dies Geschrei kann nicht beschrieben noch mit Worten erzählt werden; denn dies Wort: das Seufzen des Herzens, ist sehr groß, welches den Himmel erfüllt, und durch die Wolken dringt und bis an die Thron der göttlichen Majestät reicht, daß Gott endlich antworten muß: „Was schreiet du?“ 2 Mos. 14, 15. Ja, hat er doch nicht geschrien, sondern ist erschrocken gewesen, daß ihm die Haare zu Berge gestanden und er kein Wort hat reden können. Aber der die Herzen prüft, der versteht wohl, was der Geist haben will, der da seufzt.

105. Dies wiederhole ich nun darum und schärfe es euch so fleißig ein, weil der Glaube vor der Vernunft scheint wie ein geringes Ding und eitler Wahn, den ein Jeder im Herzen fassen könnte. Die Leute denken, es sei nur ein historisch Ding, das ganz leicht sei. Wenn man aber vor das rothe Meer kommt, ja, wenn der Tod kommt, und die Ueberwindung der Sünde, des Todes und der Hölle, so durch uns geschehen soll, alsdann wird die Kraft Gottes mit der menschlichen Schwachheit verknüpft und verbunden, und wird die Allmacht mit der Nichtigkeit und mit der äußersten Schwachheit zusammengefügt, und bringt endlich den Schwachen dahin, daß er solche Dinge thut, die unmöglich und unglaublich sind; wie Christus sagt Joh. 14, 12.: „Wer an mich glaubt, der wird größere Werke thun, denn diese“ etc. Denn also thut die Kirche auch Alles, wiewohl sie schwach ist: sie betet, sie trägt und verschlingt alle Gewalt und grimmigen Zorn des Teufels und der Menschen; sie zieht und nimmt auch hinweg Schwachheit, Sünde und Tod etc.

106. Diese Kraft und Macht des Glaubens beweist Rebekka mit der allergefährlichsten That, bei der es sich um ewigen Fluch oder ewigen Segen handelt. Denn der Segen ist das ewige Leben, der Fluch ist der Tod. Sie haben darüber mit Worten gestritten: Ob Jakob, der Sohn Isaaks, ewig sollte verflucht sein oder nicht, und mit ihm auch selbst die Mutter? In solchem großen und harten Streite hält Rebekka sich selbst vor, wie das Himmelreich offen stehe, und setzt sich zu einer Mutter des Segens.

107. Lieber, woher nimmt sie das? Antwort: Aus dem Worte. Denn niemand kann etwas Gutes und Glückseliges thun oder recht glauben ohne das Wort. Der Text 1 Mos. 25, 23.: „Der Größere wird dem Kleinern dienen“, hat ihr Muth gemacht; daraus hat sie geschlossen, daß der Segen dem Kleinern und nicht dem Größern gebühre. Isaac hat es zwar nicht verstanden, aber die Mutter hat es fleißig erwogen und der heilige Geist hat auch mitgewirkt; welcher niemand den Glauben gibt durch bloße Speculation oder Gedanken, sondern nur durch das Wort. Denn das Fleisch treibet immer wider den Geist und auf die Gedanken oder Speculation des Fleisches folgt nichts. Denn dieweil der Dinge keines gesehen wird, so dem Heiligen Geiste gebühren, so läßt es sich auch ansehen, daß Alles widerwärtig und todt Ding sei. Wenn aber das Herz das Wort ergreift, da folgt die Erleuchtung des Geistes, dazu die Kraft und Macht, wunderliche Dinge auszurichten.

\quellseite{86}
\begin{quote}
\textbf{V. 23. Und er kannte ihn nicht, denn seine Hände waren rauh, wie Esaus, seines Bruders, Hände.}
\end{quote}

108. Dies hat auch wohl natürlich und nach gemeiner Weise geschehen mögen, daß Isaac seinen Sohn nicht gekannt hat, wiewohl er seine Stimme hört, nämlich, weil er so gar fleißig und emsig war, den Segen zu geben. Denn ein Mensch, der sich nur Einer Sache befleißigt und sich dazu ganz ergeben hat, der achtet anderer Dinge nicht, wenn er dieselben auch schon sieht und mit allen Sinnen begreift. Und dies geschiehet nicht allein geistlich, sondern pfleget auch wohl leiblich zu geschehen. Wie man an den melancholischen Menschen sehen mag: wenn schon andere Leute reden, trinken, ein- oder ausgehen, so hören oder sehen sie doch nichts; denn die Gedanken des Herzens sind von den Sinnen gar abgezogen. Darum wo das Herz eines melancholischen Menschen Gedanken hat auf andere Dinge, so hat es gar keine Acht auf die Dinge, so ihm doch öffentlich vor die Sinne kommen. Ein solcher Mensch ist wohl bei andern Menschen, so mit einander reden, Historien und Fabeln erzählen, und hört doch der keines; denn das Herz ist gleich als wäre es nicht daheim, sondern gar in fremden Landen. Es kommt oftmals vor, daß ein melancholischer Mensch ißt und trinkt, und doch selbst nicht weiß, was er esse oder trinke, ob es Bier oder Wein sei. Man sagt von Bernhards, daß er Oel für Wein soll getrunken haben, da er in tiefen Gedanken gesessen hat.

109. Und solches geschiehet viel mehr, wenn eine geistliche Ueberhaftung dazu kommt und wo ein geistlich Herz ganz und gar auf geistliche Dinge beflissen ist. Also sollen wir denken, daß eben dasselbe Isaac auch widerfahren sei, der die Stimme Jakobs hört und kennt, und verschweiget solches auch nicht; denn er sagt: „Die Stimme ist Jakobs Stimme“: und läßt sich dennoch anderswohin führen. Denn es liegt ihm dieser Gedanke im Sinn, wie er keinen erstgebornen Sohn segnen wolle, und betrachtet zugleich dabei, wie es so groß Ding sei um diesen Segen, und hat zusammengefaßt alle Beschreibung, so Abraham geschehen waren.

110. Derhalben ist Isaac nicht allein mit natürlichen Gedanken bekümmert gewesen, wie die Buhler und Melancholiker pflegen, sondern auch mit geistlichen Gedanken. Es hat aber Rebecca auch zuvor solche Melancholie an ihm gemerkt; denn sie hat oftmals ihrem Manne Essen vorgesetzt, daß er es nicht ist gewohnt worden. Endlich ist Isaac dessen bei sich selbst ganz gewiß, daß die Mutter mit ihrem Sohne ihn nicht betrügen könne und auch nicht wolle; darum sagt er: „Die Stimme ist Jakobs Stimme“, und vergisst doch wieder diesen sinnlichen Eindruck und wird gar in den geistlichen Gedanken verschlungen. Die Juden haben an diesem Orte keine rechten Gedanken oder Meinung, daß Isaac den Betrug gewußt und gemerkt, doch soll er sich verstellt und dabei gethan haben, als hätte er ihn nicht gemerkt; denn sie wissen die Kraft des Geistes nicht, welche stärker ist als die melancholischen Gedanken. Wenn damit das Herz des Menschen bekümmert ist, machen sie, daß es gar erstarrt und entzückt wird.

\begin{quote}
\textbf{Und segnete ihn.}
\end{quote}

111. Der Segen hebt sich hier noch nicht an; denn es folgt bald im Text hernach, daß er zu ihm sagt: „Bist du mein Sohn Esau?“ Aber das meint Mose mit diesen Worten: Da Isaac Jakobs Hände angerührt und begriffen hatte, ist er gar entzückt und erstarrt worden, ist schier von Sinnen gekommen, und hat in seinem Herzen diesen Segen bestätigt und beschlossen, daß er beständig und gewiß sein sollte. Als wollte er sagen: Der Segen ist schon geschehen und ist gewiß; wie er hernach zu Esau sagen wird: „Er wird auch gesegnet bleiben.“ Denn es ist der Heilige Geist, der durch Isaac den Segen gibt; darum geziemet es sich nicht, daß man darin etwas widerrufen oder ändern wollte. Dieses ist ein sonderlicher Trieb und Bewegung des Heiligen Geistes gewesen, dadurch er bei sich selbst beschlossen hat, da er Jakobs Hände und Hals begriffen hatte, er wollte seinen Sohn segnen; und das hat er nicht geändert, sondern ist stracks bei der Meinung geblieben, wiewohl ihm viel gewichtige Gründe entgegenstanden.

112. Das meint Mose mit dieser Vorrede: „Und segnete ihn.“ Als wollte er sagen: Der Segen, so Jakob zugeordnet, ist jetzt schon gewiß und bestätigt; wer denselbigen hernach suchen oder begehren wird, dem wird es fehlen. Denn der Heilige Geist widerruft seine Anregung nicht; wie Maleachi Cap. 3 V. 6 steht: „Ich bin der HErr, der nicht lüget“; und 4 Mos. 23, 19.: „Gott ist nicht ein Mensch, daß er lüge, noch ein Menschenkind, daß ihn etwas gereue.“ Wenn Gott irgend ein Urtheil gefällt hat, so ändert oder widerruft er es nicht, wie die Menschen pflegen zu thun. Also hat an diesem Ort der Heilige Geist Isaaks Herz gerühret und in seinem Herzen gesagt: Jetzt segne ich diesen. Nun lassen sich aber Gottes Gaben und Berufung nicht ändern, wie St. Paulus sagt Röm. 11, 29. Darum sagt Mose: „Und segnete ihn“, das ist: Es war geschehen.

\section[Vierter Theil]{Vierter Theil. Wie Isaac den Jakob segnet}

\begin{quote}
\textbf{V. 24—27. Und er sprach zu ihm: Bist du mein Sohn Esau? Er antwortete: Ja, ich bin's. Da sprach er: So bringe mir her, mein Sohn, zu essen von deinem Wildpret, daß dich meine Seele segne. Da brachte er's ihm, und er aß; und trug ihm auch Wein hinein, und er trank. Und Isaac, sein Vater, sprach zu ihm: Komm her, und küsse mich, mein Sohn. Er trat hinzu, und küssete ihn. Da roch er den Geruch seiner Kleider, und segnete ihn, und sprach: Siehe, der Geruch meines Sohns ist wie ein Geruch des Feldes, das der HErr gesegnet hat.}
\end{quote}

113. Die Seele ist der Geist oder das Leben der Menschen in den äußerlichen Sinnen: die Seele sieht, hört, redet, meint und lacht. So viel will nun Isaac sagen: Inwendig im Herzen ist es also beschlossen, daß ich dich segnen will, nun wird dich meine Seele segnen, das ist, ich will dich mit den äußerlichen Sinnen segnen: du bist inwendig nach dem Geiste im Herzen schon gesegnet, nun will ich dich auch mit der Seele segnen.

114. Darum kommen nun auch die Ceremonien dazu; denn es können auch die geistlichen Dinge, die da äußerlich sind, ohne äußerliche Ceremonien nicht verwaltet werden. Die fünf Sinne und der ganze Leib haben ihre Geberde und Weise, unter welchen der Leib wie unter etlichen äußerlichen Larven leben muß; darum hat er ihn nicht allein im Herzen, sondern auch mit den Sinnen und äußerlichen Ceremonien gesegnet. Es haben die Söhne, so gesegnet werden sollten, dem Vater ein gut niedlich Essen und Wein vorbringen müssen; darnach sind sie herzu getreten und haben den Vater geküßt. Dieses sind rechte weltliche Ceremonien, und werden heutiges Tages noch also unter Königen und Fürsten gehalten, wo man die Lehen gibt. Und wir halten sie in den Schulen auch, wenn wir Doctorem der Theologie machen.

115. Da nun die Ceremonien gehalten sind, sagt Mose abermal: „Und segnete ihn“, nämlich äußerlich, da er den Geruch seiner Kleider gerochen hatte. Denn es ist allezeit ein äußerlicher Sinn dabei, der den Heiligen Geist beweget. Und der Heilige Geist regt und beweget auch die Herzen durch äußerliche Dinge, als nämlich, durch das Wort, durch Ceremonien und durch solche Gegenstände, so irgend durch einen äußerlichen Sinn das Herz bewegen. Nach solchem Anregen aber und Bewegung des Heiligen Geistes faßt Isaac Muth, und ist nun, gleich als wäre er entzückt, dessen ganz gewiß, daß dieser der erstgeborne Sohn sei; denn durch den Geruch der Kleider wird er in der Meinung gar bestärkt.

116. Dies sind aber die priesterlichen Kleider Abrahams und Isaaks gewesen, welche sie mit großem Fleiß bewahrt, haben Esau und andere wohlriechende Dinge dazu gelegt, daß sie die Motten nicht fressen sollten; wie wir jetzt Narden oder andere dergleichen wohlriechende Kräuter zu den Kleidern legen. Derhalben hat er nun gedacht: Dies ist Esau, mein erstgeborner Sohn; denn er hat ja die priesterlichen Kleider angezogen. Also ist er innerlich im Geist betrogen und wird nun auch mit dem äußerlichen Sinn betrogen.

\quellseite{87}
117. Er gebraucht aber zumal ein fein Gleichniß, da er sagt: „Der Geruch meines Sohnes ist wie ein Geruch des Feldes, das der Herr gesegnet hat.“ Welch(er)lei aber und wie lieblich der Geruch des Feldes sei, wenn bald die Zeit der Ernte kommen wird, das wissen die Ackerleute sehr wohl; desgleichen auch der Geruch der Weinberge, wenn die Weinernte nahe ist: das ist ein überaus lieblicher und gar ein lebendig machender Geruch. Derhalben frohlockt nun Isaak bei sich selbst und erfreut seine Seele, darum daß er seinen Sohn segnen soll. Nun habe ich einen Erben, denkt er, und Einen, der an meiner Statt das Regiment führen wird; damit bin ich so wohl zufrieden und habe daran solche Lust, gleich als ob ich irgendwo auf dem Felde oder in einem wohlriechenden Weinberge spazieren ginge, daß Geruch meinen Leib und Seele erquicken möchte. Nun will ich gern und mit gutem Frieden sterben; denn ich lasse nun hinter mir einen Herrn und Lehrer in meinem Hause, einen Vater der zukünftigen Nachkommen und einen Priester.

118. Also beschreibt Mose den Isaak, daß er freudenvoll und gleichsam trunken gewesen sei von guten fröhlichen Gedanken über den Sohn, der sein Erbe sein sollte. Er trinkt Wein und ist fröhlich; aber doch ist er mehr vom Heiligen Geiste und geistlichen Gedanken (daß ich also sagen möge) trunken gewesen. Denn das ist sein einiger und vornehmster Trost, daß er weiß, daß er nun einen Erben des Segens hat, auf welchen man dazumal warten sollte, daß er durch Christum kommen würde. Und wird ohne Zweifel das auch noch dazu gesetzt haben, nämlich, daß die Verheißung, so Abraham geschehen ist, nun vor sich gehe und bis an den dritten Erben kommen werde.

119. Dieses ist ein viel anderer und höherer Segen, denn das Weihwasser, davon die Papisten viel Lügen erdichten. Dieses sind Segen gewesen vom ewigen Leben wider den ewigen Tod, sind priesterliche und königliche Segen gewesen, die bis in das künftige Leben reichen; wiewohl dieselbigen ohne dies Leben nicht können verwaltet oder gegeben werden, und es ist nöthig, daß wir auch leiblichen Segen haben. Denn wir können des ewigen Segens nicht genießen ohne diesen zeitlichen Segen; Gott muß den Acker segnen, Brod, Fleisch und alles, was zur Nahrung vonnöthen ist, geben: aber der Mensch lebt nicht vom Brod allein, und der Segen solcher leiblichen Güter wird um des ewigen Segens willen gegeben. Darum begreifen allezeit die geistlichen Verheißungen auch die zeitlichen in sich.

120. Also sagt Christus zu seinen Jüngern Luc. 10, 7., da er sie aussendet, das Evangelium zu predigen: „In demselbigen Hause aber bleibet, esset und trinket, was sie haben“; item V. 8.: „Esset, was euch wird vorgetragen.“ Denn der Herr gibt den leiblichen Segen, als, Brod und Wein; sintemal dieses Leben des Segens der leiblichen Güter nicht entbehren kann. Das Evangelium ist der vornehmste Segen; das Andere aber alles fällt denen zu, so das Reich Gottes suchen. Der erste Segen gehört zum ewigen Leben; der andere ist leiblich, davon Christus sagt Matth. 6, 31. 32.: „Darum sollt ihr nicht sorgen und sagen: Was werden wir essen? was werden wir trinken? womit werden wir uns kleiden? Denn euer himmlischer Vater weiß, daß ihr des alles bedürfet.“ Darum werden wir hier auch sehen, wie Isaak den äußerlichen Segen zu dem Segen des ewigen Lebens gesetzt hat, welchen ewigen Segen man ohne den leiblichen nicht haben kann.

\begin{quote}
\textbf{V. 28. 29. Gott gebe dir vom Thau des Himmels, und von der Fettigkeit der Erde, und Korn und Weins die Fülle. Völker müssen dir dienen, und Leute müssen dir zu Füße fallen. Sei ein Herr über deine Brüder, und deiner Mutter Kinder müssen dir zu Füße fallen. Verflucht sei, wer dir fluchet; gesegnet sei, wer dich segnet.}
\end{quote}

121. Dieses ist die Form und Weise des Segens, und gehört der erste Theil dieses Segens auf die leibliche Nahrung, dieweil wir ohne dieselbige nicht leben können, auch im Reiche Gottes, so viel dieses Leben belanget; denn der Leib muß genähret sein, wenn sie lehren und die Kirche regieren sollen. Derhalben gehört nun der erste Theil zu der Haushaltung, und zu den Gütern, so man im Hause haben muß, daß Weib und Kind und auch das Gesinde haben möge, was ihnen täglich zu diesem Leben vonnöthen ist, welches im Vater Unser das tägliche Brod genannt wird, das ist, alles, was zur Unterhaltung des Leibes im Hause erfordert wird. „Gott gebe dir“, sagt er, „von Thau des Himmels, und von der Fettigkeit der Erde, und Korn und Weins die Fülle.“ Das sind die zeitlichen Güter.

122. Hier auf Erden bedürfen wir des Thaues vom Himmel, das ist, daß Gott vom Himmel herab Regen gebe, und bedürfen auch der Fettigkeit der Erde, daß ist, daß Gott ein fruchtbar Erdreich gebe. Denn das ist vergeblich Ding, wo man auf einen Steinfels, in das Wasser, oder in einem Walde säen soll; man muß auf ein gut fruchtbar Land säen. Danach thut auch eines guten Regens vonnöthen. Und diesen Theil des Segens macht er sehr reich und herrlich; denn er sagt: „Er wird Korns und Weins die Fülle haben.“ Er soll nicht allein vom Wasser leben, daß es ihm hart und schwer gehen würde; sondern soll ein reiches Auskommen und vollauf haben. Da wirst, sagt er, so viel haben, daß du deines Leibes wohl wirst warten können; nicht daß es mit deines Leibes Nahrung werde schmal zugehen, sondern du sollst Fett genug und alles Dinges die Fülle haben; wie Mose vom Volk Israel sagt 5 Mos. 32, 15.: „Er ist fett und dick und stark worden“ etc. Also ist nun Jakob für sich und seine Nachkommen gewiß, daß er für sein Haus Nahrung haben soll, und dieselbige soll nicht schmal und gering, sondern reich und niedlich sein. Und dasselbe mag man zwar sehen in den Büchern der Könige, wie diese Verheißung erfüllt worden sei.

123. Derhalben sollen die Gottseligen erkennen, daß sie die leiblichen Güter durch Gottes Gabe und Segen haben; und sollen sich nicht träumen lassen, wie die Heiden und Ungläubigen thun, daß uns in diesem Leben beides, Gutes und Böses, ungefähr widerfahre; sondern sollen wissen, daß solche große Gaben von Gott herkommen, auf daß sie für solche Wohlthaten Gott auch dankbar seien, wie die Apostel predigen in den Geschichten der Apostel Cap. 14, 17.: „Gott hat uns viel Gutes gethan, und vom Himmel Regen und fruchtbare Zeiten gegeben, unsere Herzen erfüllet mit Speise und Freuden.“

124. Die Mönche und andere Ungelehrte haben vor dieser Zeit gelehrt, daß man die zeitlichen Güter, Weingärten und Acker verachten soll; und sie selbst haben doch am allermeisten darnach getrachtet, haben das Allerbeste und Niedlichste gegessen und getrunken. Wider solche Träume der Mönche sollen wir lernen, daß auch die leiblichen Wohlthaten und Gaben Gottes Segen sind.

125. Der andere Theil des Segens ist weltlich und gehört zum Regiment. Denn er wird gesetzt zum Herrn über Völker und viel Leute, und daß seine Nachkommen werden Fürsten und Könige sein, und nicht allein Hausväter. Denn die Völker werden ihm dienen, wenn sie nennen nicht als Hausvätern, sondern als Fürsten und Königen werden unterthänig sein. Und wird ihm nicht allein Ein Volk, sondern viel Völker und viel Leute dienen. Dieses ist auch alles erfüllt worden zur Zeit der Richter und Könige, da die Kinder Israel das Land Canaan eingenommen haben, da sie nicht allein die Idumäer unter sich gebracht, sondern auch alle Völker und Könige desselben Landes gar vertilgt haben.

126. Der dritte Theil des Segens ist geistlich und gehört zum Priesterthum. „Deine Brüder und deiner Mutter Kinder müssen dir zu Füße fallen“, die vielleicht im Welt- und Hausregiment eben dieselbige Ehre und Herrlichkeit haben mögen, die du hast. Die priesterliche Ehre sollst du allein haben. Und dieses ist der vornehmste Theil des Segens. Zuvor hat er gesagt: „Völker müssen dir dienen, und Leute müssen dir zu Füße fallen.“ Das Fußfallen gehört zum weltlichen Regiment; dieses ist ein Anderes und geht auf die Brüder und seiner Mutter Kinder. Wiewohl es auf das weltliche Regiment möchte gezogen werden, so ist es doch besser, daß man es von dem dritten Theil des Segens verstehe, dieweil er des Regimentes zuvor gedacht hat.

127. So sind nun diese drei heiligen Regimente oder Stände, davon wir oft sagen und die wir den Leuten einschärfen, nämlich, das Hausregiment, das weltliche Regiment und das Priesterthum, oder das Haus, die Stadt und die Kirche. Das Haus hat das tägliche Brod und ist wie ein tägliches Reich. Die Polizei oder Weltregiment hat auch zeitliche Güter und ist mehr als ein täglich Regiment; denn es währet die ganze Zeit für und für, die Ewigkeit ausgenommen. Aber das Priestertum ist über das Haus und Stadt, gehört in die Kirche, ist himmlisch und ewig. Auf diese Weise ist nun Jakob, Isaaks Sohn, wohl versorgt; denn er ist zum Erben gesetzt, daß er mit seinen Nachkommen sein Haus, das Regiment und die Kirche haben soll; und ist ein reicher und herrlicher Segen, dadurch er zum künftigen Erbe bestätigt worden ist, darauf er ohne alle Widersprache gewißlich hat warten und hoffen mögen.

\quellseite{88}
128. Dieser Segen aber ist nicht ein leerer Schall von Worten oder ein Glückwunsch, damit Einer dem Andern etwas Gutes pflegt zu wünschen. Als, wenn ich sage: Gott gebe dir feine und gehorsame Kinder; das sind nur solche Worte, damit man Einem etwas Gutes wünscht, damit ich einem Andern nichts gebe, sondern allein etwas wünsche; und ist ein solcher Segen, der ungewiss ist und noch vom Erfolg abhängt. Dieser Segen aber des Patriarchen Isaak zeigt auch ein gegenwärtiges Gut und ist für immer gewiß. Es ist kein Wunsch, sondern er gibt ihm damit das Gut, und sagt damit zu ihm also: Siehe, nimm die Gaben hin, die ich dir mit Worten verspreche. Denn das ist ein Anderes, wenn ich sage: Ich wollte dir wünschen, daß du einen starken und gesunden Leib möchtest haben, daß du einen guten Verstand hättest; da das Wort „haben“ eben nicht folgt. Es ist aber ein Anderes, wenn ich dir einen Sack mit Geld darbiete, und sage: Siehe, nimm hin, da hast du tausend Gulden, die will ich dir schenken; oder da Christus sagt zu dem Bräutigam Matth. 9, 6.: „Stehe auf, heb dein Bette auf und gehe heim“ etc. Nach gemeinem Segen, damit Einer dem Andern Gutes wünscht, hätte er gesagt: Ich wollte Gott, daß du möchtest gesund und stark sein; damit würde aber die Krankheit nicht abgeschafft, und würde darauf nicht folgen, daß der Kranke wieder zu Kräften käme. Darum ist das nur ein Wortsegen.

129. In der heiligen Schrift aber sind thatsächliche Segen: nicht allein Segenswünsche, sondern wirkliche Segen, die das wirklich schenken und mit sich bringen, was die Worte sagen. Wie wir denn im Neuen Testament auch solche Segen haben durch das Priestertum Christi, welches unser Segen ist, wenn ich sage: Nimm hin die Absolution deiner Sünde. Wenn ich aber also sagte: Wollte Gott, daß dir deine Sünden vergeben wären; ach, daß du fromm und in Gottes Gnade wärest; oder: Ich wünsche dir von Gott Gnade und Barmherzigkeit, das ewige Reich und Erlösung von deinen Sünden: das möchte man einen Segen der Liebe heißen. Aber der Segen der Verheißung und des Glaubens und der gegenwärtigen Gaben lautet also: Ich absolviere dich von deinen Sünden im Namen des Vaters, und des Sohnes und des Heiligen Geistes, das ist: Ich versöhne deine Seele mit Gott, nehme von dir den Zorn und Ungnade Gottes und setze dich in Gottes Gnade, ich gebe dir das Erbe des ewigen Lebens und das Himmelreich.

130. Diese Dinge alle haben die Kraft und Gewalt, daß sie dir gegenwärtig und wahrhaftig gegeben werden, wenn du glaubst. Denn es sind nicht unsere Werke, sondern sind Gottes Werke durch unser Amt und Dienst. Derhalben sind es nicht solche Segen, die nur etwas wünschen, sondern die es auch mittheilen. Wenn ich dich taufe im Namen des Vaters und des Sohnes und des Heiligen Geistes, ist es eben so viel, als wenn ich sagte: Ich reiße dich dem Teufel aus den Händen und bringe dich zu Gott, und daselbste wahrhaftig und mit der That.

131. Auf dieselbige Weise haben die Patriarchen auch in ihren Händen Macht und Gewalt gehabt zu segnen, das ist, daß sie in einer gewissen Verheißung haben geben können denen, welche sie gesegnet haben, daß die Nachkommen würden Nahrung, Regiment und Priestertum haben. Nicht anders, denn als ob Isaak sagte: Ich gebe dir Korn, ich gebe dir das Reich und Priestertum. Die Juden behandeln diese Segen allzu kalt; denn sie verstehen dieselben nur menschlicher Weise und daß dadurch nur Gutes gewünscht werde, und nicht daß denen, so gesegnet werden, die Güter gewißlich zuerkannt und gegenwärtig überreicht werden. Der Segen aber ist ein solcher endlicher Spruch oder Urtheil, der da gewißlich bestimmt und schließt, daß die Güter, so im Segen verheißen werden, auch mit der That folgen müssen. Und eine solche Gewalt ist wahrlich ein groß Ding; denn sie bringt und gibt wahrhaftig leibliche Güter für das Haus, zeitliche Güter zum Regiment, ewige Güter zum Priestertum. Für solche Macht und Gewalt haben sie Gott gelobt, der sie den Menschen gegeben hätte, durch welche er segnete, und allerlei Gaben und Wohlthaten gäbe und erzeigte.

132. Aber das ist ein großer und erbärmlicher Jammer, daß wir sehen, wie die heiligen Patriarchen diese Güter mit so gewissem Glauben Andern geben und selbst auch empfangen, daß aber wir mit unserm Segen im Neuen Testament so kalt und schläfrig umgehen und desselben so wenig achten. Ich schäme mich wahrlich meiner selbst sehr und habe an mir einen Verdruss, so oft ich mich mit ihnen vergleiche. Denn, Lieber, siehe doch dieses Weib, Rebekka, an, die auch im Fleisch und Blut lebt, wie wir thun, desgleichen Isaak und Jakob, wie sie so einen gewissen, lebendigen und starken Glauben an die zukünftigen Güter haben, daß sie des Fleisches, so gegenwärtig bei ihnen ist, nichts achten. Ja, es läßt sich ansehen, als ob sie in diesem zeitlichen Leben schlafen und gleichsam schnarchen im Hinblick auf die zukünftigen Güter, die ihnen versprochen werden.

133. Wir haben eine viel reichere Gabe oder haben sie wenigstens nicht in geringerem Maß; aber wir haben ja solchen Glauben nicht, wir sind fast schläfrig, sind halb todt, unsere Augen sind dunkel, die Ohren hören schwer, unsere Herzen wanken auch und sind sehr unbeständig. Sie haben ja große Gaben, und achten doch des nicht, das sie haben. Denn wo man Einem die Absolution spricht, die Sacramente der Taufe und des Abendmahls des Herrn reicht, aus dem Evangelium Vergebung der Sünde verkündigt, das ist ein viel größer Ding, denn wo Isaak den Jakob segnet. Denn es ist gleich so viel, als wenn ich sagte: Ich gebe dir das Himmelreich, das Regiment über den Teufel; und wenn du schon sterben mußt, so will ich dich doch erhalten und bewahren, daß du nicht sollst verloren werden. Das thun wir zwar nicht aus unsrer Kraft, sondern thun es von Gottes wegen und aus göttlichem Befehl, der den Menschen diese Macht gegeben hat, daß Einer den Andern mit dem Priestertum Christi zum ewigen Leben bringen möge.

134. Also, sage ich, haben wir zwar viel reichere Verheißungen, denn sie gehabt haben; wir sind aber schläfriger und verachten diesen Schatz des Reichthums und der Güte Gottes. Denn wir sind, leider, sonst nicht dankbar, wir sind von wegen dieser Seligkeit nicht fröhlich, freuen uns derselben nicht von ganzem Herzen. Wir haben die Gaben in großer Fülle; aber wir achten derselben nicht, ja, verachten sie wohl ganz und gar, so sie doch alle so gewiß sind, als diese Segen der Patriarchen gewesen sind; ja, das noch das \unsicher{Allerbetrüblichste} ist, je reichlicher wir es haben, je mehr die Welt darüber tobt und solchen Schatz verfolgt.

135. Derhalben soll das unsere Herzen erwecken und uns munter machen, soll uns die Schlafsucht vertreiben, daß wir sehen, daß die lieben Väter so fleißig und wacker gewesen sind, die mit größerem Geist und Glauben die Verheißung und Segen angenommen haben, als wir thun. Darnach sollen wir unsere Gaben groß machen, welche ja ihren Gaben gleich oder noch wohl etwas größer sind; wiewohl dieselbigen auch sehr groß und trefflich gewesen sind. Wir haben nicht die Väter, daß die mit uns redeten; sondern wir haben den Sohn Gottes selbst, der redet jetzt mit uns, wie in der Epistel an die Hebräer steht, Cap. 1, 2., desselbigen hören wir sagen: Ich absolviere dich, ich gebe dir die Schlüssel zum Himmelreich, die Macht zu taufen, selig zu machen, Teufel und Hölle unter die Füße zu treten; diese göttliche Macht gebe ich euch, daß ihr eben dieselbigen Werke auch thun sollt, die ich thue, und sollt größere thun denn diese.

136. Aber wir gähnen und schlafen, und denken also: Christus ist wahrer Gott und Mensch, darum ist es sein Wunder, daß er solche Gaben den Menschen mittheilt. Daß aber die Menschen, Einer dem Andern, Gaben dieses und des zukünftigen Lebens geben sollen, das ist zumal ungereimt Ding, und läßt sich ansehen, als sei es unmöglich und unglaublich. Also sind aber die heiligen Väter nicht gesinnet gewesen, sondern haben die göttlichen Segen hoch und werth gehalten, sind fröhlich darüber gewesen und haben Gott dafür gedankt. Darum sollen wir uns, mit ihnen vergleichen, schämen, und unsre Trägheit, und daß wir so gar schläfrig sind, an uns strafen lernen und uns bessern.

137. Nun folgt in dem letzten Theile das liebe heilige Kreuz, aber doch zugleich auch dabei der Sieg durch das Kreuz und im Kreuz. Denn also sagt Isaak: „Verflucht sei, wer dir fluchet.“ Diese Worte sind genommen aus Cap. 12, 3. und 22, 17. droben, da der Herr zu Abraham also sagt: Du wirst verflucht werden, aber ich will den Fluch in einen Segen wenden. Dieses ist aber eine große Gewalt, daß Isaak einen solchen Spruch und Urtheil sprechen darf, vor welchem die Hölle mit allen Teufeln und die ganze Welt mit aller ihrer Macht und Gewalt erschrecken und erzittern muß. Denn er will so viel sagen: Ich weiß, daß diesem Segen der Teufel, die Welt und das Fleisch werden feind sein; dieses weiß ich, daß es gewiß geschehen wird. Lieber Jakob, ich überschütte dich zwar mit großen Gaben, ich erhöhe und ehre dich; denn du wirst ein Vater sein ohne Armuth, ein König ohne Hinderniß, ein Priester und Heiland der Seelen auch wider die Pforten der Hölle. Du mußt aber bedenken, daß dieses alles in der Verheißung steht, und daß es dermalen noch nicht erfüllt ist mit vollem Siege, auf welchen man noch in der Hoffnung warten muß.

\quellseite{89}
138. Derhalben wirst du deine Güter also haben, daß es sich doch gleichwohl mit dir wird ansehen lassen, als habest du davon gar nichts. Denn du wirst angefochten werden in deiner Haushaltung, im Welt- und Kirchenregiment: um dieses alles willen werden dich die Gottlosen neiden und du wirst von ihnen verflucht werden, also daß der Fluch schon bereitet ist bei dem Segen. Ich segne dich zwar reichlich; es wird aber der Teufel und die Welt kommen, es werden deine Brüder kommen, und werden dir fluchen, dich verfolgen, und versuchen, deinen Segen zu verwirren und umzustoßen.

139. Und endlich ist es auch also gekommen, wie diese Weissagung lautet. Denn wie viel hat Jakob wohl gelitten vom Anfang dieses Segens! Sein Bruder Esau droht ihm, er wolle ihn erwürgen; da wird ihm alsbald das Kreuz im Hausregiment aufgelegt. Er ist ein Vater der Nachkommen, so von ihm sollen geboren werden; aber er muß Tag und Nacht in Gefahr Leibes und Lebens stehen, muß erwarten, daß ihn sein eigener Bruder erwürgen werde, und wird mit großen Sorgen und Beschwerung von seinen Eltern heimlich weggebracht, daß er sich vor solcher Gefahr hat hüten mögen. Er muß in Syrien fast an zwanzig Jahre im Elende sein. Ei, welch ein schöner Segen ist mir das! Er hat zwar denken mögen: Lieber Vater, wie reimt sich das, daß du gesagt hast vom Thau des Himmels und Fettigkeit der Erde, mit solcher großen Beschwerung? Denn er hat auch nicht einen Brosamen Brodes, da er nach Syrien zieht. Er ist elend und arm, muß vierzehn Jahre dienen, und in demselbigen Dienst hat er viel Unbilligkeit und Unrecht leiden müssen: darum erinnert ihn der Vater, daß es also kommen werde, daß er anstatt des Segens den Fluch empfinden werde.

140. Hernach aber, da die Theurung im Lande Canaan kommen wird, wird er auch in große Noth und Gefahr kommen mit seinem ganzen Gesinde, schier bis an die Gefahr seines Leibes und Lebens. Und endlich, wie greulich und übel ist man in Egypten mit den Kindern Israel umgegangen! Wo bleibt denn nun der Segen? Antwort: „Der Mensch lebet nicht vom Brod alleine, sondern von einem jeglichen Wort, das durch den Mund Gottes gehet“, Matth. 4, 4. 5 Mos. 8, 3. Mit dem Wort wird der Mensch aufgerichtet und wird auch dadurch erhalten. Wiewohl er nun jetzt gegenwärtig zu seiner Haushaltung nicht Brod hat, so stirbt er dennoch nicht, wird nicht verlassen; sondern der Glaube an Gottes Wort ernährt ihn mitten in Armuth, beide an Leib und Seele, und endlich verändert er auch die Theurung in eine Fülle und daß er muß genug haben. „Denn die den Herrn suchen“, wie im 34. Psalm V. 10. 11. steht, „haben keinen Mangel an irgend einem Gut.“ Dagegen aber sagt derselbe Psalm von den Reichen: „Die Reichen aber müssen darben und hungern.“ Item im 37. Psalm V. 2.: „Wie das grüne Kraut werden sie verwelken“ c.

141. Man muß aber dies im Glauben verstehen und darauf harren, wie Jakob auch im Glauben es dafür gewißlich gehalten hat, daß ihm nimmermehr in seinem Hause etwas mangeln würde. Und dem hat auch der Ausgang sehr schön entsprochen. Denn da der Segen folgt, ist er reicher geworden als Laban, sein Schwiegervater. Aber das ist ohne Anfechtung nicht zugegangen; denn er ist versucht, aber doch nicht verlassen worden. Denn der Segen wird wohl angefochten, aber doch nicht umgestoßen; er wird angegriffen, aber doch nicht umgeworfen; wie im 118. Psalm V. 13. 14. steht: „Man stößet mich, daß ich fallen soll, aber der Herr hilft mir. Der Herr ist meine Macht, und mein Psalm, und ist mein Heil.“

142. Der Herr gibt einen solchen Segen, der mit Geduld gemischt und mit dem Heiligthum des heiligen Kreuzes geziert ist, auf daß wir in der Anfechtung unterrichtet werden und lernen, daß unser Leben nicht allein im Brod stehe, sondern in einem jeglichen Wort Gottes. Das Brod aber gibt endlich der Herr auch genüglich und läßt es daran nicht fehlen, wenn wir im Glauben geprüft und versucht sind, ob wir auch Gott in seinen Verheißungen glauben wollen. Denn er gibt wohl die Verheißung, aber er versucht uns gleichwohl und entzieht uns den Segen, gleich als ob wir auf keinen Segen warten dürften. Aber gerade dann sinnt und denkt er auf den Segen, wenn wir den Fluch empfinden. Derhalben, sage ich, kann der Segen angefochten, aber unterdrückt und umgestoßen kann er nicht werden; wie man auch von der Wahrheit pflegt zu sagen.

143. Und dies ist der Fluch oder das Kreuz im Hause gewesen, welches den Patriarchen Jakob hart genug geplagt hat. Siehe aber nun auch auf das weltliche Regiment bei seinen Nachkommen. Siehe David an, welcher der erste König gewesen ist nach dem Segen. Lieber, wie oft ist sein Reich und Regiment angefochten und in Gefahr gekommen, als sollte es jetzt dahinfallen! Erstlich zu den Zeiten des Königs Saul; darnach durch seinen eigenen Sohn Absalom und Ahithophel, den schalkhaftigen untreuen Rath; und ist noch zur Zeit Absaloms das böse Gewissen auch dazu gekommen. Jedoch hat er sich noch immer mit der göttlichen Verheißung aufgerichtet und erhalten; wie im 21. Psalm V. 2. steht: „Herr, der König freuet sich in deiner Kraft“; item, im 63. Psalm V. 12.: „Der König freuet sich in Gott. Wer bei ihm schwöret, wird gerühmet werden.“ Und 2 Sam. 15, 25. 26. sagt er: „Werde ich Gnade finden vor dem Herrn, so wird er mich wieder holen; spricht er aber also: Ich habe nicht Lust zu dir; siehe, hie bin ich, er mache es mit mir, wie es ihm wohl gefällt.“ Denn das Königreich war ihm verheißen und bestätigt, aber nicht ohne sehr große Mühe und Anfechtung, so er darüber erlitten hat.

144. Auf dieselbige Weise steht man in der Kirche auch mehr Fluch denn Segen. Denn dasselbe lehrt uns auch die Erfahrung, die wir im selbigen Amte sind, welchen die Taufe, die Gewalt der Schlüssel und das Sacrament des Altars gegeben ist. Aber wie viel sind derer, auch der sehr frommen Leute, die der Satan also hart anficht und bethört, wenn er ihnen den himmlischen Segen vor den Augen wegnimmt, daß es sich mit ihnen ansehen lässet, als behielten sie von diesem Segen gar nichts; ja, sie fühlen und empfinden vielmehr Fluch und Angst der Hölle, denn göttlichen und himmlischen Segen? Die Andern fallen in offenbare Sünde und Laster, und gerathen dahin, daß sie den Segen offenbar verachten und gar verwerfen. Diejenigen aber, so ihn noch behalten und lieben, sind so gar schwach, daß ihnen wohl noth thut, daß man sie viel und stets ermahne und den Segen immer wiederhole, daß sie dadurch also erweckt und das Wort Gottes in ihnen geschärft werden möge, wie Mose sagt 5 Mos. 6, 7.

145. Unterdeß aber geht gleichwohl dem Segen nichts ab, sondern er bleibt beständig, kräftig und sehr reich, und größer, als wir verstehen können. St. Paulus klagt eben darüber auch, daß ihm ein Pfahl in das Fleisch gegeben sei, 2 Cor. 12, 7., und daß wir durch viel Trübsal müssen in das Reich Gottes gehen, Ap. Gesch. 14, 22. Derhalben wird es am Fluche nicht mangeln. Aber du mußt dich dawider legen, sei nur stark, und bleibe an dem Segen hangen, wenn es sich schon ansehen läßt, daß alles voll Fluches sei. Denn also sollen wir schließen: Es ist gewiß, daß ich getauft bin; ich habe das Wort gehört aus dem Munde des Dieners; ich habe das Sacrament des Altars gebraucht: das ist die göttliche Wahrheit, die nicht kann geändert werden. Wenn ich schon schwach bin, so ist dennoch dieselbe Wahrheit gewiß und beständig, die Güter sind sehr gewaltig und reich, das Herz aber ist fast schläfrig und wanckt, wenn es die Güter ergreifen und annehmen soll.

146. Davor aber sollen wir uns allein hüten, daß wir ja nicht verlegen; und wo wir ja nicht mit großem Geschrei bekennen können, daß wir dennoch leise zu Gott murmeln. Können wir nicht singen, wenn wir Gott loben, so lasset uns doch zum wenigsten seufzen, daß wir in dem Segen, darein uns der Sohn Gottes gesetzt hat, beharren mögen, welchen man ohne großen Kampf und mancherlei Anfechtung nicht behalten kann. Denn also haben die Väter die Verheißung auch gewiß und stark gehabt, aber nicht ohne Anfechtung. Und um der Ursache willen ermahnt uns Christus auch so fleißig dazu, daß wir beständig bleiben und beharren sollen. „Fasset eure Seelen mit Geduld“, sagt er Luc. 21, 19. Ihr seid Kinder des Reichs, die Sünden sind euch vergeben, der Teufel ist überwunden und unter eure Füße gelegt, die Sünde und Tod werden euch nichts schaden, sondern ihr seid nun unschuldig; derhalben so leidet auch mit geduldigem Herzen den widermärtigen Fluch, so euch begegnen wird.

\quellseite{90}
147. Nachdem er nun den Fluch, das ist, Kreuz und Anfechtung gesetzt hat, so den trefflichen und reichen Segen folgen, thut er nun etwas mehr dazu und sagt weiter: „Gesegnet sey, wer dich segnet.“ Als wollte er sagen: Sie werden euch nicht alle fluchen oder anfechten, sondern es werden ihrer Viele herzu kommen, die euch segnen und eures Segens auch wollen theilhaftig werden; das wird die Frucht sein der Anfechtungen, wenn ihr im Glauben beharren werdet. Wie Christus sagt Joh. 12, 24.: „Wahrlich, wahrlich, ich sage euch, es sei denn, daß das Weizenkorn in die Erde falle und ersterbe, so bleibts alleine; wo es aber erstirbt, so bringts viel Früchte.“

148. Ein Christ, der versucht und angefochten worden, nützt mehr, denn sonst hundert, die nicht versucht worden sind. Denn in der Anfechtung wächst der Segen, daß er viel Leute lehren, trösten und ihnen in leiblichen und geistlichen Sachen mit gutem Rathe helfen kann. Also, will Christus sagen, seid ihr in der Welt versucht, so seid ihr doch zugleich auch mit himmlischem Segen erfüllt.

149. Nun wollen wir auch die Auslegung etlicher Worte nach der Grammatik besehen. Adorare heißt, wenn man sich mit dem Leibe neigt und einem Andern Ehrerbietung erzeigt. Dasselbe Wort wird auch in weiterem Sinne gebraucht; denn man neigt sich vor Königen, Brüdern und vor dem Nächsten. Droben (1 Mos. 23, 7.) neigt sich Abraham vor den Kindern Heth. Darnach wird es auch gezogen auf das geistliche Neigen, welches ist, wo man sich neigt oder bückt vor dem Schöpfer. Denn man neigt sich anders vor dem Schöpfer und wieder anders vor der Creatur.

150. Warum sagt er aber: „Sey ein Herr über deine Brüder“, so er doch nur Einen Bruder hat, nämlich Esau? Antwort: Das ist eine sonderliche Art und Weise zu reden in der heiligen Schrift, da alle Kindeskinder werden „Kinder“ genannt, und Bruderkinder werden „Brüder“ genannt; die von Rebekka und ihren Töchtern geboren sind, dieselben sind seine Kinder genannt worden. Also ist Esau ein Vater seiner Kinder und auch seiner Kindeskinder. „Brüder“ aber werden in der heiligen Schrift genannt nicht allein die von einer Mutter geboren sind, sondern auch die Nachkommen und Kindeskinder.

\section[Fünfter Theil]{Fünfter Theil. Wie Esau den Segen Isaaks verlangt; wie Isaak sich weigert, und Esau sich über Jakob beschwert, bei Isaak beständig anhält um den Segen und endlich ein Stück des Segens erhält}

\begin{quote}
\textbf{V. 30. 31. Als nun Isaak vollendet hatte den Segen über Jakob, und Jakob kaum hinaus gegangen war von dem Vater Isaak, da kam Esau, sein Bruder, von seiner Jagd, und machte auch ein Essen, und trugs hinein zu seinem Vater, und sprach zu ihm: Stehe auf, mein Vater, und iß von dem Wildpret deines Sohns, daß mich deine Seele segne.}
\end{quote}

151. Da der Segen vollendet war, folgt nun darauf ein Klagelied und der entgegengesetzte Segen Esaus. Und das ist der andere Theil von dieser Historie. Man soll aber die Worte in der Rede Esaus fleißig betrachten. Denn Jakob hatte nicht also geredet in der dritten Person: Surgat pater: Mein Vater wolle aufstehen; sondern er sagte: Surge quaeso: Lieber Vater, stehe auf. Aber dieser Heuchler braucht eine sonderliche Rhetorik. Wie solche Heilige pflegen zu thun, nämlich, daß sie solche Worte brauchen, die nur zum Schmuck dienen, und dahin gerichtet sind, daß sie einen großen Schein haben. Esau hat seine Bitte fein einfältig und schlicht vorgetragen, wie die Dialectiker pflegen zu thun; dieser aber macht seine Rede groß nach Art der Rhetorik, darauf Jakob nimmermehr hätte denken können. „Mein Vater“, sagt er, „wolle aufstehen und essen von der Jagd seines Sohnes.“

152. Mit solcher Hoffahrt weiß er seine Rede zu schmücken, gleich als ob er allein und dazu ein eingeborner Sohn wäre, Jakob aber kein Sohn sei. Denn also pflegen die Heuchler zu thun: sie nennen sich die Kirche und haben auch äußerlich eine größere Larve, und einen solchen Schmuck, der einen köstlicheren Schein hat, als die haben, so recht gottselig sind. An den Gottseligen aber ist kaum so viel Kühnheit und auch solcher großer Schein nicht; wie Jakob sich nicht untersteht, seinen Vater mit solchem Stolze anzureden. Mose gedenkt aber hier der Kleider gar nicht, die er doch ohne Zweifel auch wird gebraucht haben. Dieselben wird vielleicht Rebekka heimlich wiederum an ihren Ort gelegt haben, hat sich darum nun nicht mehr bekümmert, wie sie dieselben bewahren möchte; wird gedacht haben: Er mag immerhin hingehen, sich bekleiden und schmücken, so schön als er selbst will, mein Sohn hat nun den Segen schon erlangt.

\begin{quote}
\textbf{V. 32. 33. Da antwortete ihm Isaak, sein Vater: Wer bist du? Er sprach: Ich bin Esau, dein erstgeborner Sohn. Da entsetzte sich Isaak über die Maße sehr, und sprach: Wer? wo ist denn der Jäger, der mir gebracht hat, und ich habe von allem gegessen, ehe du kamest, und habe ihn gesegnet? Er wird auch gesegnet bleiben.}
\end{quote}

153. Alhier fängt er an tapfer zu lügen, und gebraucht zwar prächtige Worte, die einen großen Schein haben: „Ich bin“, sagt er, „dein erstgeborner Sohn“, macht seine Rede herrlich und groß; wie es denn einem solchen Redner gebühret. Jakob ist nur ein einfacher Dialectiker, der ist vor den Augen Esaus gar verachtet, gilt vor ihm nichts, ja, er ist gar verflucht, ist nicht werth, daß er soll Isaaks Sohn genannt werden, vielweniger der Erstgeborne Isaaks. Der Heilige Geist aber hat schon das Gegenteil beschlossen. Denn Jakob hat nun den Segen hinweg und du Esau bist mit nichten der Erstgeborne.

154. Hast du denn vergessen, was du jenesmal gethan hast, da du sagtest, 1 Mos. 25, 32.: „Siehe, ich muß doch sterben, was soll mir dann die Erstgeburt?“ Gib mir her das Linsengericht? Also hat Esau die Erstgeburt verkauft, und da er gegessen und getrunken hatte, stand er auf und ging davon, und verachtete also seine Erstgeburt. Das alles hat der Heilige Geist noch nicht vergessen, wie du dir träumen lässest. Hast du die Erstgeburt deinem Bruder verkauft, so bist du mit Recht verdammt, daß du der Erstgeburt beraubt sollst sein. Darum ist es erlogen, daß du sagst, du seiest mein erstgeborner Sohn; denn du hast die Erstgeburt verschleumet mit dem Linsengerichte, und hast sie verworfen in großer Verachtung. Gott will für seine Wohlthaten geehrt sein, und wer ihn ehrt, denselben ehrt er wiederum auch; wer ihn aber verachtet, wird verschmäht sein, wie 1 Sam. 2, 30. steht. Desselbigen Esauches ist dieses ein greulich Exempel.

155. Isaak aber erschrickt und entsetzt sich über die Maße sehr, da er die Worte Esaus hört, und besorgt sich, daß vielleicht durch Gottes Verhängniß irgend ein Fremder den Segen möchte weggenommen haben und daß also seine Nachkommen desselben sollten beraubet werden. Dieses ist nicht ein geringes Entsetzen gewesen; denn er hat Sorge gehabt, daß diese Ehre und Herrlichkeit des Segens, des Regiments und Priestertums irgend durch Gottes Rath von seinem Hause möchte entwendet werden. Wie auch eben dieselbe Gefahr eingetreten ist, da Esau droht, daß er seinen Bruder Jakob erwürgen wolle: da war die Gefahr vorhanden, wo Jakob wäre umgekommen, daß der Segen Fremden hätte mögen gegeben werden.

156. Jedoch läßt er sich das nicht anfechten, daß er dadurch sollte bewogen worden sein, den Segen zu ändern, den er einmal gegeben hatte; sondern er sei, sagt er, wer er sei, so den Segen hingekommen hat, der soll gesegnet sein und bleiben. Denn die Gaben Gottes können nicht widerrufen werden. Darnach aber, da er sich genauer der Sache erinnert, hat er gedacht: Siehe, ich habe die Stimme Jakobs gehört, der hat mich betrogen, da er die Hände und den Hals mit Bocksfellen überzogen hatte; ja, der ist es, den ich gesegnet habe.

\begin{quote}
\textbf{V. 34. 35. Als Esau diese Rede seines Vaters hörte, schrie er laut, und ward über die Maße sehr betrübt; und sprach zu seinem Vater: Segne mich auch, mein Vater. Er aber sprach: Dein Bruder ist kommen mit List, und hat deinen Segen hinweg.}
\end{quote}

\quellseite{91}

157. Hier wird gefragt: Warum Isaak den Segen nicht widerrufen habe, weil Jakob denselben mit List hinweggenommen hat? Esra erzählt, was der Juden Meinung sei, nämlich, daß Isaak, da er die Klage Esaus gehört, den Segen habe widerrufen wollen; aber er soll daselbst nach Gottes Willen die Hölle haben offen gesehen, und daß sie ihm bereitet wäre, wo er den Segen widerrufen würde. Darum soll er sich also entschlossen haben, und hat den Segen nicht widerrufen, sondern denselben vielmehr bestätigt. Dieses zeigt Esra an nach dem Verstande der Juden; aber sie legen die Worte der Schrift nicht aus, sondern verdunkeln sie vielmehr. Darum soll man es dafür nicht halten, daß Isaak, wiewohl er sich über die Maße sehr entsetzt hat, sollte gedacht haben, den Segen zu widerrufen; denn er hat wohl gewußt, daß dieser Segen ein Werk und Gabe Gottes ist, die gänzlich unbegreiflich und unveränderlich muß bleiben. Als, wenn ich irgend einem die Taufe gebe, so ist mein Herz und Wille ganz gewiß, daß ich ihn wahrhaftig taufen will; wo aber der, so getauft wird, mit List umgeht, so habe ich doch gleichwohl ihm die rechte Taufe gereicht, welche nicht mein ist, sondern ist wahrhaftig ein göttliches Werk.

158. Auf solche Weise hat Isaak auch gesagt: „Ich habe ihn gesegnet, er wird auch gesegnet bleiben.“ Und dasselbe hat er zuvor ernstlich bei sich selbst also beschlossen, und hat damit nicht ohne sonderlichen Rath und Bedenken verzogen bis an das Ende seines Lebens; darum ist er dessen gewiß gewesen, da er Jakob gesegnet, daß er damit einen endlichen Spruch oder Urtheil gesprochen habe, so aus göttlichem Befehl gesprochen und bestätigt sei. Und ist eben derselbe Segen gewesen, den er erblich von den Vätern Adam, Noah, Abraham und den Andern empfangen hatte. Solche Sprüche können und sollen auch nicht geändert werden; denn Gott ändert seine Gaben nicht. Er widerruft die Taufe nicht, die Absolution auch nicht, und andere Gaben, die er uns durch sein Wort mittheilt: wenn er mir meine Sünden vergibt, so sind sie mir wahrhaftig vergeben.

159. Sie disputiren hier auch von der List: Ob die heiligen Väter auch listig gehandelt, und ob sie auch daran gesündigt, wenn sie also mit List sind umgegangen? Denn wir haben oftmals gehört, daß sie tapfer gelogen haben, nicht allein aus Noth, zum Schein, sondern auch in Wirklichkeit. Aber in dieser That ist keine Sünde. Im Werk zwar und vor den Augen der Menschen ist es wohl ein Betrug und List; denn Jakob hat den Vater betrogen damit, daß er die Hände und den Hals mit Fellen überzogen hatte; aber vor Gott ist es kein Betrug. Denn die Erstgeburt und der Segen hat ihm gebühret, welche er beide von seinem Bruder gekauft, und hat sie auch sonst zuvor gehabt, daß sie ihm von Gott zugeordnet war nach der göttlichen Antwort: „Der Größere wird dem Kleinern dienen.“

160. Warum man von einem heimlich nachstellt, und mit List einem Andern das wegnimmt, so Gott dir gegeben hat, ist es keine Sünde. Wie der Raub, so die Israeliten den Egyptern genommen, wohl ein rechter Raub ist nach dem Urtheil der Menschen; aber die Israeliten haben damit nicht gesündigt, daß sie die Egypter beraubt haben. Denn Gott hat sie das geheißen 2 Mos. 11, 2., daß sie sollten zu den Egyptern sagen: Lieber, leihe mir ein silbern Geräthe und Kleider, das Fest des Herrn damit zu schmücken, da sie gedachten, davon zu fliehen. Sie hatten ein öffentlich ausdrückliches Gebot Gottes, daß er sagte: Ich will es also haben, daß du die Egypter dermaßen betrügen und berauben sollst. Denn die Egypter waren den Israeliten ihren Lohn schuldig für ihren Dienst und schwere Tyrannei, so sie bei ihnen hatten leiden müssen; wiewohl es zwar eine schlechte und geringe Erstattung war dafür, daß sie das Volk Israel so lange unterdrückt und beschwert, ihnen dazu auch ihre Kindlein getödtet hatten.

161. Auf solche Weise soll man diese List Jakobs auch beurtheilen. Denn wo die Heiligen irgend mit List umgehen und dessen von Gott Befehl haben, alsdann ist es wohl vor den Augen der Menschen ein Betrug oder List, aber es ist eine heilige, billige und geistliche List. Dergleichen ist es nicht vonnöthen, daß man davon disputiren und viel Fragens machen wollte, wie oder ob Jakob gesündigt habe? sondern dahin soll man sehen, daß ihm zuvor aus Gottes Befehl solches erlaubet gewesen, was er seinem Bruder mit List genommen hat. Also haben die Heiligen im Kriege oftmals die Feinde hintergangen und betrogen; aber das sind solche Lügen, die man wohl gebrauchen mag im göttlichen Amt wider den Teufel und die Feinde Gottes.

162. Also betrügt der Fischer den Fisch auch mit solcher List, wenn er an die Angel Speise bindet; welch Gleichniß die Väter nicht übel auf Christum gezogen haben. Denn der ist in die Welt gekommen und hat sich in unser Fleisch gekleidet, ist in das Wasser geworfen worden wie eine Angel: da aber denselbigen der Teufel gebissen hat, ist er alsbald von Gott wiederum aus dem Wasser gezogen, und auf das Trockene geworfen und zermalmt worden, das ist, Christus hat dem Teufel seine schwache Menschheit vorgehalten, welche die ewige und unüberwindliche Majestät bedeckt hat. Da ist der Teufel an der Angel der Gottheit angelaufen und hat sich daran gestoßen, durch welche Gottheit alle seine Macht, dazu auch die Gewalt des Todes und der Hölle ist überwunden worden; wie St. Paulus Col. 2, 15. sagt: „Er hat ausgezogen die Fürstenthümer und die Gewaltigen, und sie Schau getragen öffentlich, und einen Triumph aus ihnen gemacht durch sich selbst.“ Es hätte sich aber der Teufel nicht unbillig beklagen mögen, daß er auch unbilliger Weise hintergangen und betrogen wäre, sinemal er gemeint hätte, er würde einen Menschen erwürgen, und er wäre nun selbst von ihm mit List erwürget worden. Aber es ist durch Gottes wunderbaren Rath eben dasselbe widerfahren, das man gemeiniglich pfleget zu sagen: Ars ut artem falleret: Daß eine List die andere betrügen möchte.

\begin{quote}
\textbf{V. 36. Da sprach er (Esau): Er heißt wohl Jakob; denn er hat mich nun zweimal untertreten. Meine Erstgeburt hat er dahin; und siehe, nun nimmt er auch meinen Segen. Und sprach: Hast du mir denn keinen Segen vorbehalten?}
\end{quote}

163. Alhier hebt sich die Feindschaft an wider seinen Bruder: „Er heißt wohl Jakob“, sagt er, „denn er hat mich nun zweimal untertreten.“ Es ist aber beides erlogen. Du Esau lügst; weißt du nicht, was du zuvor gethan hast? Du hast selbst willig die Erstgeburt verkauft, da du sagtest 1 Mos. 25, 32.: „Ich muß doch sterben, was soll mir denn die Erstgeburt?“; den Segen hast du verachtet, welcher auch nach dem Tode kräftig ist. Darnach hast du auch den Befehl Gottes nicht geachtet, darin gesagt ist: „Der Größere wird dem Kleinern dienen“; darum bist du nicht mehr der Erstgeborne, sondern der Käufer ist der Erstgeborne, und hat den Segen mit allem Recht, den man von ihm nicht darf zurückfordern, dieweil es eine Gabe Gottes ist, die nicht geändert wird: du aber lügst und zeihst deinen Bruder fälschlich, daß er beiderlei soll genommen oder geraubt haben, nämlich, die Erstgeburt und den Segen.

164. Und daraus kann man sehen, wie der Text in der Epistel an die Hebräer am 12. Cap. V. 16. 17. zu verstehen sei, der also lautet: „Daß nicht jemand sei ein Hurer, oder ein Gottloser, wie Esau, der um einer Speise willen seine Erstgeburt verkaufte. Wisset aber, daß er hernach, da er den Segen erwerben wollte, verworfen ist; denn er fand keinen Raum zur Buße, wiewohl er sie mit Thränen suchte.“ Denn die Ursache, daß er mit seiner Buße oder Thränen nichts hat können ausrichten, ist diese: weil es keine rechte wahrhaftige Buße gewesen ist. Das ist wohl wahr und steht fest, daß Gott, der nicht fehlen oder lügen kann, allen Menschen, die wahrhaftig Buße thun, seine Gnade und Barmherzigkeit vorgehalten hat, und wo man von wegen der Sünde wahrhaftige Buße thut, findet solche Buße allezeit Statt bei Gott. Darüber soll man fest halten und dasselbe vertheidigen; wie Gott solches mit vielen Exempeln und in den Büchern der heiligen Schrift bezeuget hat.

165. Es ist aber eine andere Buße, die nicht wahrhaftig, sondern falsch und erdichtet ist, welche die Deutschen eine Galgenreue nennen, nämlich, wenn ich also Buße thue, daß ich mich nicht schäme, daß ich Gott erzürnt habe, sondern daß ich mir selbst habe Schaden gethan. Solche Buße ist sehr gemein, und ich habe selbst auch oftmals also Buße gethan, und habe mir es lassen leid sein, daß ich irgend etwas närrisch, unmöglich und mit Schaden ausgerichtet hatte. Ich schämte mich der Thorheit und des Schadens mehr als der Sünde, als der Schuld und daß ich Gott erzürnt hatte. Wo man aber allein über den Schaden, so man gethan hat, Leid trägt, das ist eine solche Buße, davon Gott nichts weiß; ja, unsere eigenen Herzen wissen davon auch nichts, wie das an Esau wohl zu sehen ist. Denn er sagt nicht: Nun verstehe ich es, daß ich gesündigt habe, warum habe ich Gott erzürnt, da ich die Erstgeburt verkauft? ich will des Segens nun gerne entrathen, allein, daß mir Gott die Sünde möge verzeihen. Das wäre eine rechte wahrhaftige Buße gewesen, dadurch er sich bekümmert hätte, wie er mit Gott versöhnt werden möchte von wegen der begangenen Sünde.

\quellseite{92}

166. Denn die wahrhaftige Buße steht auf den Zorn Gottes um der Sünde willen, mit dem wollte sie gern versöhnet sein, flieht vor dem Zorn Gottes; sie gebietet nicht allein Schmerz um des Schadens willen, gebietet auch seinen Zorn und Haß wider den Bruder oder Nächsten; sondern sagt also: Wenn mir nur Gott wollte gnädig sein, so wollte ich gern allerlei Schaden und Unglück leiden und damit vorlieb nehmen.

167. Deren hörest du keines von Esau. Er hat eine Buße von wegen der Strafe, und nicht von wegen der Sünde, so er begangen hatte, darum hat er auch keine Buße der Strafe gefunden; denn er sucht nicht die Buße der Sünden, sondern ist noch immer halsstarrig in der Sünde. Wie Saul 1 Sam. 15, 3. ff. auch that, da er den Befehl Gottes übertreten hatte, damit ihn Gott geheißen hatte, Amalek mit Mann und Weib zu erwürgen, dazu auch mit allem großen und kleinen Vieh; er aber verschonte des Königes und des großen Viehes, und da er von Samuel gestraft ward, antwortete er ihm also: „Das Volk verschonete der besten Schafe und Rinder, um des Opfers willen des Herrn, deines Gottes.“ Das hieß nicht die Sünde recht erkennen, sondern dieselbe unter dem Schein der Heiligkeit entschuldigen und vertheidigen. Darum wird der Prophet sehr zornig und verkündigt ihm ein sehr trauriges Urtheil und sagt: „Weil du des Herrn Wort verworfen hast, hat er dich auch verworfen, daß du nicht König seiest“ etc., 1 Sam. 15, 23.

168. Also unterstand sich Adam auch, seine Sünde zu entschuldigen, und legte die Schuld nicht allein auf das Weib, sondern auch auf Gott: „Das Weib“, sagt er 1 Mos. 3, 12, „das du mir zugesellet hast, gab mir von dem Baum, und ich aß“ etc. Also entschuldigt sich Eva auch, und beschuldigt die Schlange, die sie betrogen hätte. Das soll mit nichten eine Buße genennt werden. Denn wo rechte Buße ist, da empfindet der Mensch, daß er in seinem Herzen also gesinnet ist, und bei sich selbst klagt: Ach, warum habe ich Gott erzürnt? Warum habe ich seinen Zorn und schweres Gericht wider mich erregt? Er mag mich Andern zum Exempel strafen, wie er nur will, allein daß er mir die Sünde verzeihe und vergebe. Wie die Kinder Israel im Buch der Richter am 10. Cap. V. 15. bitten: „Wir haben gesündiget, mache du es nur mit uns, wie dirs gefällt; allein errette uns zu dieser Zeit.“ Ein solch Gebet gebührt sich denen zu thun, die wahrhaftig und mit Ernst Buße thun, daß sie sagen: Der Herr mag zwar mit mir machen, was er will, allein, daß er mir gnädig sei. Und da ist auch Gott ganz bereit, Gnade zu erzeigen und die Sünde zu verzeihen, und bietet solchen bußfertigen Sündern alsbald an, daß er ihnen die Sünde vergeben wolle.

169. Wie weit aber Esau noch von solcher wahrhaftigen Demuth und Erkenntniß der Sünden ist, zeigen seine Worte an, in welchen er noch ganz unverschämt lügt, daß er der Erstgeborne sei, so er es doch nicht ist, und sagt, daß die Erstgeburt sein sei, die er doch zuvor aus freien Stücken verworfen hatte. Zuletzt, das noch viel schwerer ist, so beschuldigt und verdammt er seinen unschuldigen Bruder, entschuldigt und lobt sich selbst, da er doch schuldig war: „Er heißt“, sagt er, „wohl Jakob; denn er hat mich untertreten.“ Er gedenkt aber seiner Sünde auch nicht mit einem Worte. Er sagt nicht: Ich habe gesündigt, ich habe bisher nicht gefühlt oder verstanden, daß ich gesündigt habe; nun lerne ich es an der Strafe und am Gerichte Gottes; nun ist es Zeit, daß ich meine Sünde erkenne; denn ich empfinde, daß Gott greulich erzürnt ist dadurch, daß ich die Erstgeburt verkauft und dieselbe verachtet habe; ich will sie aber meinem Bruder nicht mißgönnen, sondern will sie ihm von Herzen gerne gönnen, und ich freue mich dessen, daß er den Segen erlangt hat. Ich habe gesündigt und ist mir herzlich leid, nicht so sehr darum, daß ich den Segen verloren habe, als um des Zorns Gottes willen, den ich wider mich erregt habe.

170. Ja wohl, er sagt oder denkt nichts weniger denn also; ja, er rühmt sich noch seiner Bosheit. Ich bin gerecht, mein Bruder ist gottlos, der hat mir mit List das Meine genommen und heißt fremdes Gut. Daß es ihm übel gerathe! Also verdammt er die Gerechtigkeit eines Andern und seine eigene Sünde vertheidigt er noch immer. Endlich faßt er einen greulichen Haß wider seinen Bruder, gönnt ihm die Gnade nicht, die ihm Gott gegeben hat und die er mit Recht und durch seine eigene Schuld verloren hat; über das droht er ihm, daß er ihn erwürgen wolle. Ei, das ist mir zumal eine feine Buße.

171. Darum ist es Anderes nichts, denn daß er darüber Leid trägt, daß er das Gut verloren und Schaden empfangen hat. Also tragen auch die Diebe und Uebelthäter über ihre Sünde Leid, nicht, daß sie mit Ernst Buße thäten oder den Sünden sollten feind sein; sondern weil sie von der Furcht und dem Gefühl der Strafe gemartert werden. Wie wir im Pabstthum für die Sünde Buße gethan haben, da wir alle Jahre zur Beichte sind gezwungen worden. Denn dazumal hätten die Leute auch lieber den Sünden nachgehangen, wo sie es hätten thun mögen; und bald nach Ostern thaten sie doch wiederum, wie ihre Art war, und freuten sich, daß sie nur der Marter, so sie mit der Beichte gehabt, los waren.

172. Darum, sage ich, heißt Buße thun, wo man, erstlich, den Zorn Gottes um der Sünde willen fühlt; also daß dem Sünder im Herzen angst werde, und er sich bekümmere und Verlangen gewinne nach der Hülfe und Gnade Gottes. Einen solchen bußfertigen Sünder aber soll aufrichten und trösten, wer da kann, daß er nicht verzweifele und verloren möge werden, wie Judas verzweifelt ist. Denn ein solch Herz ist nicht anders gesinnt, denn daß es gern an alle seinem Gute wollte Schaden leiden, allein, daß es nur möge Trost finden, daß es möge Hülfe haben auch von dem aller kleinsten Kinde, das ihm die Verheißung der Seligkeit verkündige. Denn es befindt: Wenn ich nur vom Zorn Gottes, damit ich jämmerlich gemartert werde, erlöset werden möchte, so wollt ich mich darum gar nicht bekümmern, wie ich die Güter, so ich verloren habe, möchte wieder bekommen. Habe ich natürlich gehandelt, ei, so mag ich darum meinen Lohn nehmen, oder um meiner Thorheit willen die Strafe leiden; allein, daß mir nur Gott gnädig sei.

173. Ja nun steht aber also: Wenn ich meine Erstgeburt hätte, so wollte ich mich darum nicht sehr bekümmern, ob mir Gott gnädig oder ungnädig wäre. Derhalben soll man nun den Text in der Epistel an die Hebräer am 12. Capitel v. 17. also verstehen: „Er fand keinen Raum zur Buße“; denn es war keine rechte wahrhafte Buße, sondern es gereuete ihn nur der Schade, den er empfangen, und seine Thorheit, dadurch er die Erstgeburt und seine Herrlichkeit verloren hatte. Darum verfiel er in greulichen Haß gegen seinen Bruder und Gott, welcher Jakob durch seinen Vater Isaak gesegnet hat.

174. Und dies ist ein rechtes Bild und Exempel eines solchen Menschen, der fälschlich Buße thut. Welches man darum merken soll, daß wir lernen, die rechte Buße von der falschen Christen Buße zu unterscheiden, welche wir eine Galgenreue nennen. Denn ein Dieb, wenn er sieht, daß man ihn mit dem Galgen strafen will, trägt auch Leid über die Sünde und wollte gerne länger leben: wo ihm aber Gott die Gnade und ein solch Licht gibt, daß er seine Sünde und den Zorn Gottes wahrhaftig erkennt, so bekümmert er sich nicht mehr um dieses Leben; allein, daß ihm nur möge geholfen werden, daß er nicht in Gefahr komme seiner Seelenseligkeit halben. Und da soll man ihn lehren und aufrichten mit Gottes Gnade und Barmherzigkeit, so er uns in Christo erzeigt hat, und soll ihn nicht mit der Ohrenbeichte oder päpstlichen Genugthuung trösten; denn sonst würde es mit ihm gar verloren sein.

\begin{quote}
\textbf{V. 37. Isaak antwortete, und sprach zu ihm: Ich habe ihn zum Herrn über dich gesetzt, und alle seine Brüder hab ich ihm zu Knechten gemacht, mit Korn und Wein hab ich ihn versehen; was soll ich doch dir nun thun, mein Sohn?}
\end{quote}

175. Der Segen Jacobs begreift die drei heiligen Stände in sich, nämlich, das Hausregiment, das Weltregiment und das Priesterthum. Diese drei Stände oder Aemter zu verwalten, gibt er ihm auch die Güter dieser Welt. In der Polizei hat er das Reich und Regiment über das Volk. In der Kirche hat er das Amt, die Vergebung der Sünden und ewiges Leben zu verkündigen und auszutheilen; darum ist Jacob ein Bischof, er hat die Weissagung, hat das Wort, Vergebung der Sünden und den Gottesdienst. Darnach setzt er noch dazu: „Ich hab ihm alle seine Brüder“ im Weltregimente „zu Knechten gemacht“, und er wird sich die Völker selbst auch unterthänig machen. Welches unter David geschehen ist, der zuerst den Segen erfüllt hat, da er die Syrer, Philister, Araber, Edomiter unter sich gebracht hat. Was soll ich aber nun thun, sagt Isaak, was soll ich über diese dreierlei Segen dir geben? Er hat Alles hinweg. Und das ist ohne Zweifel aus sonderlichem Rath und Willen Gottes geschehen, dessen dieser Segen gewesen ist, darum kann ich ihn nicht ändern. Esau aber hört noch nicht auf, ja, er hält noch immer hart an bei dem Vater mit Bitten und Flehen, daß er ihn auch segnen wolle.

\quellseite{93}

\begin{quote}
\textbf{V. 38. Esau sprach zu seinem Vater: Hast du denn nur Einen Segen, mein Vater? Segne mich auch, mein Vater; und hub auf seine Stimme, und weinete.}
\end{quote}

176. Daher kommt, daß in der Epistel an die Hebräer Cap. 12, 17. gesagt wird: „Wiewohl er sie mit Thränen suchte.“ Denn dieses ist das ängstliche und heiße Flehen mit Thränen, damit er vielmals und mit großer Angst nachsucht; aber er hat zu lange geharrt und ist Alles vergeblich. Denn wo Gott einmal weggezogen, und wo das Wort einmal hinweggenommen ist und die Gnade Gottes, so kann man es nicht leichtlich wieder finden. Zu Rom ist vor Zeiten eine sehr schöne Kirche und Gemeinde Gottes gewesen; es sind daselbst mehr fromme Christen und Märtyrer, so die Wahrheit bekannt haben, gewesen, denn sonst irgend an andern Orten in der ganzen Welt: aber da dieselbige Kirche und Gemeinde hinweggenommen ist, hilf Gott, welche erschreckliche Finsterniß und Greuel sind darnach gefolgt! Und wenn sie sich jetzt auch zu Tode schreien oder heulen würden, oder gar mit Esau sterben sollten, so richteten sie doch damit nichts aus, das vorige Licht und Gnade Gottes wieder zu erlangen.

177. Wir haben vor dieser Zeit im Papsthume geschrien um die ewige Seligkeit und das Reich Gottes; wir haben uns fast wehe gethan, ja, wir haben unsere Leiber schier zu Tode gemartert, nicht mit dem Schwerdte oder anderer äußerlichen Wehre, sondern mit Fasten und Kasteiung des Leibes: da haben wir gesucht, haben angeklopft Tag und Nacht. Und ich selbst, wo ich nicht durch den Trost des Evangeliums Christi wäre erlöset worden, so hätte ich nicht zwei Jahre leben können; also zermarterte ich mich und floh vor dem Zorn Gottes, und mangelte auch an Thränen und Seufzen nicht: wir richteten aber damit nichts aus. Derhalben ermahnt der heilige Paulus 2 Cor. 6, 1. 2. nicht umsonst so fleißig, da er sagt: „Daß ihr nicht vergeblich die Gnade Gottes empfanget. Sehet, jetzt ist die angenehme Zeit; jetzt ist der Tag des Heils.“ Darum lasset uns der Gnade, die uns gegeben ist, gebrauchen, dieweil wir es können; lasset uns den Mund und das Herz aufthun, daß der Segen eingegossen werde; denn wo derselbe sollte weggenommen werden, so ist es mit uns gar verloren.

178. Die Deutschen pflegen in einem Sprüchwort einen Fluch zu brauchen, der nicht böse ist, nämlich, da sie sagen: Gott gebe faulen Händen die Drüse (Beule, Pestbeule). Denn sie ermahnen uns damit, daß wir nicht sollen faul sein, sondern die Gelegenheit, so sich darbietet, wohl gebrauchen. Also, wenn das Evangelium gelehrt wird, so sollen wir es hören und lernen mit dankbarem Herzen. Wie Christus sagt Joh. 12, 35. 36.: „Es ist das Licht noch eine kleine Zeit bei euch. Wandelt, dieweil ihr das Licht habt, daß euch die Finsternisse nicht überfallen. Es ist das Licht noch eine kleine Zeit bei euch, glaubet an das Licht, dieweil ihrs habt, auf daß ihr des Lichtes Kinder seid.“

179. Und das Exempel der Juden bezeuget gewißlich genugsam, wie greuliche Finsterniß sie überfallen habe, nachdem sie an den nicht geglaubt haben, der so große Zeichen vor ihnen gethan hat. Denn sie erfüllen nur herrlich diese Figur oder Vorbild Esaus, schreyen Tag und Nacht, fasten den Leib mit Fasten und Beten nun in die anderthalb tausend Jahre, beten ganz eifrig und sagen: Herr Gott, sende doch den Messias um deines Namens, um deines Wortes, um deines Reichs willen. Von welchem Gebete und Klagen auch Steine und harte Felsen möchten bewegt werden; sie finden aber keinen Raum zur Buße. Nicht daß die Buße keinen Raum oder Statt finden könnte, sondern darum, daß sie auf einem verkehrten Wege in den Himmel kommen und mit ihrem Dienste den Segen erlangen wollen, welches unmöglich ist. Sie erkennen ihre Sünde nicht, sondern rechtfertigen sich selbst, machen sich fromm, wie Esau auch that. Wir, sagen sie, sind das rechte Israel; wie Esau sagt: Ich bin dein erstgeborner Sohn. O lieber Gott, die Heiden haben die Erstgeburt und den Segen mit Gewalt weggenommen etc. Sie fühlen zwar den Zorn Gottes wohl, sie wollen aber die Sünde nicht erkennen.

180. Wo man aber die Sünde und den Zorn Gottes um der Sünde willen fühlt, das ist eine sehr große Gnade, und solchen Sündern ist die Seligkeit nahe, daß sie leichtlich können zur Buße gebracht werden. Aber wenn man die Sünde vertheidigt und entschuldigt, das heißt Gott richten, das ist, verdammen in seinen Worten, wie im 51. Psalm V. 6. steht. Darum ist Esau ein Exempel aller Juden.

181. Und da wir Mönche waren, haben wir eben darum auch mit unserm Kasteien nichts ausgerichtet. Denn wir wollten unsere Sünde und gottlos Wesen nicht erkennen; ja, wir wußten von der Erbsünde nichts, und haben nicht verstanden, daß der Unglaube Sünde wäre. Ja, wir hielten und lehrten auch, daß man an Gottes Gnade und Barmherzigkeit zweifeln müsse. Derhalben, je mehr ich lief und begehrte zu Christo zu kommen, je weiter mich er von mir. Nach der Beichte und wenn ich Messe gehalten hatte, konnte ich in meinem Herzen nimmer zufrieden sein; denn das Gewissen kann keinen rechten gewissen Trost haben von den Werken.

182. Darum lasset uns des Segens, so jetzt gegenwärtig ist, gebrauchen, und der Gnade, so uns gegeben ist, nachdem nun das Licht des Evangeliums wiederum an den Tag gekommen, und lasset uns nicht unfleißig oder undankbar sein. Denn wenn der Segen einmal weggenommen ist, so steht es darnach in unserm Willen nicht, daß wir ihn wieder an uns bringen mögen, sondern ist allein an Gottes Gabe und Gnade gelegen, und auch also, daß er sich keines Menschen Weinen, Schreien, Mühe oder Arbeit bewegen läßt.

183. Das sahe ich zuerst wohl ein, daß die freie Gnade Gottes durchaus nöthig wäre, um das Licht und ewige Leben zu erlangen, und ich arbeitete fleißig und ängstlich, wie ich doch den Spruch Pauli Rom. 1, 17. verstehen sollte, da er sagt: „Die Gerechtigkeit Gottes wird im Evangelio geoffenbaret.“ Da suchte ich lange und klopfte immer an; denn das Wort justitia Dei, „die Gerechtigkeit Gottes“, lag mir im Wege, welches man nach gemeinem Gebrauche also auszulegen pflegte: Die Gerechtigkeit Gottes ist eine solche Tugend, dadurch er für sich gerecht ist und die Sünder verdammt. Also hatten alle Doctores diesen Spruch ausgelegt, Augustin ausgenommen, daß sie sagten: Die Gerechtigkeit Gottes, das ist, der Zorn Gottes. So oft aber ich den Spruch las, wünschte ich allezeit, daß Gott das Evangelium niemals möchte offenbart haben; denn wer könnte den Gott lieben, der da zürnt, richtet und verdammt? bis ich endlich durch Erleuchtung des Heiligen Geistes den Spruch des Propheten Habakuk etwas fleißiger erwogen habe, da er sagt im 2. Cap. V. 4.: „Der Gerechte lebt seines Glaubens.“ Daraus habe ich abgenommen und geschlossen, daß das Leben aus dem Glauben muß herkommen, und zog also das abstractum in das concretum, wie man es in den Schulen zu nennen pflegt, das ist, ich zog das Wort „Gerechtigkeit“ auf das Wort „gerecht“, nämlich, daß der Mensch vor Gott gerecht würde durch den Glauben etc. Da wurde mir die ganze heilige Schrift und der Himmel selbst auch geöffnet.

184. Aber jetzt zur Zeit sehen wir nur dieses große Licht ganz klärlich, und wir mögen desselben reichlich gebrauchen. Wir verachten aber, leider! diesen edelsten und solchen himmlischen Schatz und werden desselben überdrüssig. Wenn uns aber dereinst dieser Schatz wiederum genommen werden sollte, alsdann werden wir abermals schreien und anklopfen, wie Christus sagt im Gleichniß von den thörichten Jungfrauen, Matth. 25, 11. Aber das wird alles vergeblich und umsonst sein. Darum lasset uns Gott fürchten und dankbar sein. Ihr sollt euch aber vornehmlich bewegen lassen mein und Anderer Exempel, die wir im Tode und in der Hölle gelebt haben, und den Segen so reichlich nicht gehabt, wie ihr jetzt habt. Derhalben sehet, daß ihr euch fleißig übet in der Lehre von dem Segen und derselben recht nachdenket, daß ihr sie für euch selbst behalten und auch bei andern Leuten ausbreiten möget. Wir haben gethan, was wir haben thun sollen.

\begin{quote}
\textbf{V. 39. 40. Da antwortete Isaak, sein Vater, und sprach zu ihm: Siehe da, du wirst eine fette Wohnung haben auf Erden, und vom Thau des Himmels von oben her. Deines Schwerts wirst du dich nähren, und deinem Bruder dienen. Und es wird geschehen, daß du auch ein Herr, und sein Joch von deinem Halse reißen wirst.}
\end{quote}

\quellseite{94}

185. Er gibt ihm Antheil am häuslichen Segen; denn was das Welt- und Kirchenregiment belanget, steht Alles bei seinem Bruder Jakob. Dies ist aber kein Segen. Denn er sagt nur: „Du wirst eine fette Wohnung haben“, das ist, du wirst deine Nahrung und Kleidung haben; auf daß das fleischliche Volk der Juden auch etwas haben möchte, daß sie nicht sagen dürften, daß sie gar nichts hätten, so gibt er ihm die leiblichen Verheißungen. Darnach setzt er noch hinzu: „Deines Schwertes wirst du dich nähren“, das ist, du wirst dich mit deinem Schwerte schützen und vertheidigen müssen; du wirst aber nicht gewinnen oder das Regiment haben. Dies ist alles ohne Segen zu verstehen. Gott wird dich nicht zum Herrn über andere Völker setzen, viel weniger über deinen Bruder, welchen er zuvor ausgenommen hat, und hat Jakob zum Herrn gesetzt über seine Brüder. Du wirst dir müssen genügen lassen mit deiner Wohnung im fetten Lande, und wirst dich deines Schwertes nähren müssen, daß du vor den Feinden sicher seiest, doch also, daß du deinem Bruder unterthänig seiest und ihm dienst.

186. Nun halte man diese Zwei gegen einander: Jakob ist der Doctor in der Kirche und ist auch der Herr im Weltregiment über seinen Bruder Esau, und sein Segen wird mit Recht ein Segen genannt. Aber der Segen dieses Esau wird kein Segen genannt; denn der Text sagt nicht: Dies soll dein Segen sein, sondern sagt also: „Du wirst eine fette Wohnung haben auf Erden.“ Und ist im Hebräischen eben dasselbe Wort, so auch im 1. Psalm V. 1. steht: In sede derisorum: „Da die Spötter sitzen.“ Die heilige Schrift hat das fleißig verhütet, daß Esau nicht hat müssen gesegnet heißen; denn der Segen streckt sich in die Breite, in die Länge, auf das Fette und alle Güter; und wird auch im Text nicht gesagt: Der Herr, Gott wird dir es geben. Denn wo das Wort „Herr“ hinzugesetzt wird, so wird es ein Segen. Derhalben will Esau so viel sagen: Gott wird dich in Kraft dieses Segens nicht segnen; du wirst aber zufällig etwas haben, daß dir Gott geben wird, wie er denen in Arabien und den Egyptern auch gibt, die zufällig ihre Nahrung und Schutz gehabt, haben es aber für keinen göttlichen Segen erkannt. Das soll man fleißig merken. Diese zwei Worte „Gott“ und „Segen“ merken im Text nicht gesetzt; darum so ist es auch kein Segen.

187. Er thut aber noch mehr dazu, und sagt: „Und es wird geschehen, daß du auch ein Herr, und sein Joch von deinem Halse reißen wirst“ etc. Dieses ist auch ein sonderlich Stück. Denn Esau möchte ihm gern mehr geben, und sucht auch fleißig, was er ihm doch mehr geben möchte. Dieses gehört aber auf Herodes und sein Geschlecht. Die Juden legen es aus von der Zeit des Königs Joram, 2 Kön. 8, 20., unter dem Propheten Elisa, von welchem die Idumäer sind abgefallen, darum daß der König gottlos war. Und ist nicht ein übler Verstand; denn da hat Edom das Joch zerrissen und hat sich los gemacht von der Herrschaft Israels. Ich halte aber gleichwohl, daß man es noch besser von Herodes verstehen möchte, welcher ein Sohn Antipaters gewesen ist, des Idumäers, der unter Julius Cäsar ein großes Ansehen gehabt und in großen Ehren ist gehalten worden, und das Königreich Israel geplagt hat. Derhalben ist es genommen, daß er das Joch von seinem Halse geworfen und auch geherrscht hat.

188. Daraus entsteht aber nun eine Frage: Wie die Weissagung oder Verheißung erfüllt worden sei, daß Jakob gesegnet sei, und der Größere dem Kleineren dienen werde und aber der Kleinere herrschen werde: und hernach wird folgen, daß sich Jakob neiget vor Esau, beide er selbst und sein Gesinde, dazu seine Kinder und Weiber. Er wird daselbst vier Haufen verordnen, die alle Esau zu Füße fallen und sich vor ihm neigen. Darnach herrschen die Idumäer eine lange Zeit in Israel, und Herodes mit seinen Nachkommen regiert über Israel gewaltig, schier bis auf die Zerstörung Jerusalems; und Jakob hat doch aus diesem Text den Segen genommen, so hernach im 49. Capitel V. 10. 11. folgen wird. Wie kann man nun diese zwei Dinge ausgleichen, daß sie nicht wider einander seien?

189. Antwort: Wo man nur eine Stunde oder einen Tag herrscht, das wird nicht mit Recht eine Herrschaft oder Regiment genannt. Also, da das Volk Israel gefangen gen Babel ist geführt worden, hätte man auch sagen mögen, daß sie das Regiment verloren hätten: sie haben es aber in Wahrheit nicht verloren. Denn sie sind nur gezüchtigt worden. Das Regiment war darum nicht verloren, sondern verzogen und aufgeschoben. Auf solche Weise neigt sich Jakob auch vor Esau. Er ringt mit dem Engel, also daß er auch schier am Regiment und am Segen hätte zergangen mögen; denn er hat in dem Kampfe Anderes nichts gefühlt, denn daß er seinem Bruder würde müssen unterworfen sein und daß sein Bruder ihm sein Gesinde und Weib erwürgen würde: aber nach dem Kampfe wird er vom Engel Israel genannt, und wird eben dazumal in dem Kampfe am allermeisten zum Könige und Regenten gesetzt, da er unterliegt. Denn „du hast“, sagt der Engel 1 Mos. 32, 28., „mit Gott und mit Menschen gekämpft, und bist obgelegen.“

190. Das heißt nicht, den Segen oder das Regiment verlieren, wenn jemand im Glauben angefochten wird über das Regiment und Segen. Wie Abraham eine Stunde ist versucht worden, daß er seinen Sohn schlachten sollte, da Gott zu ihm spricht 1 Mos. 22, 2.: „Opfere mir deinen Sohn.“ Darum hat er aber den Sohn nicht verloren, sondern hat ihn reichlich und mit Frucht und Nutzen wieder bekommen; wie der Engel überdies den Segen Jakobs auch bestätigt auf seine Bitte, da er sagt: „Ich lasse dich nicht, du segnest mich denn.“ Derhalben ist der Segen nicht dahingefallen oder aufgehoben, sondern er ist damit versucht und angefochten worden, auf daß der Segen desto stärker und gewisser würde.

191. Was soll man aber von Herodes sagen, der wahrlich über Jerusalem und Juda mit großem Stolz und Grausamkeit regiert hat? Antwort: Da ist die Prophezeiung oder Weissagung Jakobs erfüllt worden: „Es wird das Scepter von Juda nicht entwendet werden“ etc. Denn Maria, die Mutter Christi, war dazumal schon geboren und haben auch die Eltern Johannes des Täufers schon gelebt, als Herodes das Regiment sechs oder auf das höchste zehn Jahre gehabt hatte; wiewohl es vor der Geburt Christi nicht bestätigt gewesen ist. Wie Philo schreibt, daß er sechs Jahre ordentlich und ein und dreißig Jahre tyrannischer Weise regiert haben soll, welches mit nichten werth ist, daß man es ein Regiment nennen soll, sein Vater aber, Antipater, ist vor ihm nur ein Oberster oder Hauptmann im jüdischen Lande gewesen. Darnach hat er auch kein still oder geruhig Regiment gehabt; denn die Juden haben ihm stets widerstrebt um der Verheißung willen, darüber sie sehr hart hielten, und ist für unbillig angesehen und gehalten worden, daß sie, die Juden, ihrem Knechte, Herodes dienen sollten, dieweil doch zu Esau gesagt war: „Du sollst deinem Bruder Jakob dienen.“ Derhalben hat wohl Herodes über die Juden regiert, aber mit keinem ordentlichen Regiment.

192. Endlich kann man auch wohl also antworten: Daß er das Regiment erlangt habe, nicht aus göttlichem Befehl, dadurch zuvor in Israel und in Juda die Könige eingesetzt und bestellt waren; sondern durch die Gewalt und Tyrannei der Römer, aus Gottes Zorn, daß er auch dasselbe Königreich zerreißen sollte, mit dem es schon fast am Ende war, und das ohnedem allgemach getrennt und zerstört ward. Darum ist seine Herrschaft oder Regiment nicht wider die Verheißung; denn eben dieselbige Zeit ist das Ende gewesen des Königreichs, darin Christus sollte geboren werden.

193. Und dieses sind die zweierlei entgegengesetzten Verheißungen der zwei Brüder. Jakob hat den Segen im Hause, dazu auch im Welt- und Kirchenregiment erlangt. Esau hat deren gar keines, sondern hat nur eine solche Verheißung, wie die andern Völker auch gehabt; als, die Syrer, welchen Herrschaft die Juden bisweilen haben entragen und dulden müssen, das Königreich aber oder Regiment haben sie nicht verloren, denn nur als es am Ende gewesen, auf daß die Prophezeiung Jacobs möchte erfüllt werden, darin er zuvor gesagt hatte, 1 Mos. 49, 10.: Wenn der Silo oder Held komme, würden, sollte das Königreich mit dem Priesterthum und Gesetz sein Ende haben.

\section[Sechster Theil]{Sechster Theil. Wie Esau um des Segens willen zornig wird auf Jakob, und den Vorsatz faßt, ihn umzubringen}

% Spalten 326–329 ergänzt nach Google Books ftFBAQAAMAAJ (Ohio State, St. Louis 1881):
% https://books.google.com/books?id=ftFBAQAAMAAJ&pg=PA327 und &pg=PA329.
\begin{quote}
\textbf{V. 41. Und Esau war Jakob gram um des Segens willen, damit ihn sein Vater gesegnet hatte; und sprach in seinem Herzen: Es wird die Zeit bald kommen, daß mein Vater Leid tragen muß; denn ich will meinen Bruder Jakob erwürgen.}
\end{quote}

194. Wir haben gesagt von der Buße des gottlosen Menschen Esau, daß er geweint habe, und wie es Hieronymus im Lateinischen gegeben hat: Irrugit clamore valido: Er hat heftig geschrieen wie ein Löwe, darum daß ihm der Segen genommen war, den er doch zuvor selbst freiwillig verkauft hatte. Nun folgen auch die Worte, die einem solchen Bußfertigen geziemen. Denn wie die Buße ist, ebenso ist auch die Frucht der Buße und die Genugthuung. Denn er sagt also: „Es wird die Zeit bald kommen, daß mein Vater Leid tragen muß.“ Ei, warum das? Antwort: „Denn ich will meinen Bruder erwürgen.“

195. Ei, wie ist mir das so eine feine Buße. Er zürnt nicht allein mit seinem Bruder, sondern auch mit seinen Eltern und mit GOtt selbst, dessen, wie er wohl weiß, der Segen ist und von dem man auch allein den Segen erwarten soll. Ja, also soll man Buße thun! Zuvor weint er und schreit: „O mein lieber Vater Isaak, segne mich auch“, der ich ja dein Sohn bin; jetzt aber ist er gar unsinnig geworden und gedenkt, dreierlei Mord zu begehen. Denn erstlich will er seinen Bruder erwürgen. Ja, das heißt GOtt für die Sünde genug thun mit der Buße. Darnach aber will er seine beiden Eltern in großes schweres Herzeleid bringen, dadurch sie auch getödtet werden. Den Bruder will er mit einem Knüttel oder Schwert erwürgen; die Eltern mit schwerem Leide und Traurigkeit tödten. Das weiß der gottlose Mann Esau wohl, daß es also kommen wird; darum sagt er vorsätzlich und wissentlich, er werde seine Eltern mit solchem Schmerz tödten.

196. Ja, das heißt mit teuflischer Unsinnigkeit und Bosheit behaftet sein; dieweil er also denkt: Wenn ich den Segen nicht haben soll, so will ich verschaffen, daß die Andern alle desselben auch sollen beraubt werden. Es soll ihn weder mein Bruder, noch meine Eltern, noch die ganze Kirche behalten. Denn es sind der Brüder nur zwei, welchen der Segen gehört oder die Erbschaft des Regiments und Priesterthums. Und wenn Jakob wäre erwürgt worden, so wäre doch das Erbe des Segens auf Esau nicht gekommen, denn dieser wäre entweder am Leibe gestraft, oder aber in das Elend getrieben worden wie Kain. Derhalben ist die Kirche oder Gemeinde und der Segen im Hause Isaaks in großer Gefahr gestanden.

197. Und da siehst du, was Esau für eine Buße gehabt hat. Es ist einfach eine teuflische Unsinnigkeit, dadurch er aus dem Herzen hinweg wirft alle Ehrerbietung und Furcht beide GOttes und seiner Eltern, und vergißt ganz und gar aller Gottseligkeit und Ehrbarkeit; also daß wo er es thun könnte, so wollte er, daß der Segen mit Vater und Mutter und mit der ganzen Kirche gar vertilgt wäre. Denn wo seine Unsinnigkeit sollte vor sich gegangen sein und daß Jakob erwürget wäre, so wäre Isaak wahrlich, der ein sehr zartes Herz gehabt, und Rebekka auch, vor erbärmlicher und schwerer Traurigkeit dahin gestorben; wie es unmöglich gewesen wäre, daß Adam hätte können lebendig bleiben um des Mordes willen seines Sohnes, so Kain begangen hatte, wo ihn nicht GOtt sonderlich und wunderbarlich erhalten hätte. Also hätten solche schwere Schmerzen diese Eltern auch getödtet.

198. Derhalben sehet ihr, wie der Teufel im Hausregiment und Gesinde auch der heiligen Leute tobt und wüthet. Wir erinnern nicht vergeblich so oft, vermahnen oft und schreien, daß ihr ja fleißig und ohne Unterlaß beten wollet, darum daß der Teufel nicht weit von uns, sondern mitten unter uns ist. Siehe, welch einen Lärm er angerichtet hat in der heiligen Kirche und Gemeinde der Väter, im Hause Isaaks, da Esau vorhat, seinen Bruder zu erwürgen, seine Eltern zu tödten, die ganze Kirche zu verderben und zu zerstören!

199. Nun ist aber Esau billig betrübt, und hat deß gute Ursache, daß er so heftig zornig ist; aber es ist gleichwohl seine eigene Schuld. Denn er hat daran gesündigt, da er dasjenige, so er jetzt heftig begehrt, nicht groß geachtet, ja, gar verachtet hat, und jetzt, da er nun den Zorn GOttes fühlt, wird er vor großem Schmerz und Ungeduld gar unsinnig.

200. Diese große Bosheit des Teufels sollen wir erkennen lernen, dadurch er zuwege gebracht, daß er in den Häusern der allerheiligsten Menschen seine eigenen Kinder hat, die ihren Brüdern und ihren eigenen Eltern nach dem Leben stehen. Isaak ist ein sehr heiliger Patriarch, ein Vater der Verheißung; Rebekka ist ein sehr heiliges Weib und Mutter derselbigen Verheißung: nun ist aber von ihrem Fleisch und Blut Esau geboren, der ihren Tod begehrt, und selbst darnach trachtet, wie er sie tödten möge. Lieber GOtt, wovor sollten denn wir uns nicht zu fürchten haben?

201. Es ist aber der Schmerz der Eltern viel heftiger und bitterer, als der Kinder, Brüder oder der angebornen Freunde. Denn es sind sehr große und heftige Neigungen, welche GOtt in der ganzen Welt geschaffen und den Eltern eingepflanzt hat gegen ihre Kinder, und wenn ihre Herzen mit Schmerz und Traurigkeit verwundet werden von wegen irgend eines Unfalls ihrer Kinder, das ist ihr gewisser Tod und ihrem Leben ein Gift. Darum werden die Eltern leichtlich getödtet, wo nicht mit dem Schwert, so geschieht es doch durch Traurigkeit und Schmerzen. Ich habe selbst viel fromme ehrliche Eltern gesehen, die von wegen der Bosheit ihrer Kinder vor Traurigkeit des Herzens gestorben sind.

202. Die Jugend denkt daran nicht, versteht es auch nicht: man soll aber die Kinder lehren und ermahnen, daß sie nicht Mutter- und Vatermörder werden; denn es wird ein sehr schrecklich Gericht und göttliche Strafe über sie gehen, wie St. Paulus sagt 1 Tim. 1, 9.: „Das Gesetz ist gegeben den Vatermördern und Muttermördern.“ Es kommt oft vor, daß die Kinder gar sicher sind und also in mancherlei Sünde und Laster fallen, und betrachten darin die Ehre gegen ihre Eltern gar nicht. Die Töchter lassen sich zu Schanden machen und schmähen damit ihre frommen und ehrlichen Eltern; aber mit solchen Schanden und Lastern tödten sie Vater und Mutter. Denn die haben eine sehr zarte und brünstige Neigung zu ihren Kindern, welche die Kinder nicht so empfinden, ja, sie verstehen solche Neigung nicht einmal.

203. Derhalben soll sich die Jugend wohl vorsehen und lernen, daß sie ihre Eltern in Ehren halten und ihnen gebührliche Ehrerbietung erzeigen, und daß sie diese Worte oder Namen, „Vater“ und „Mutter“, für sehr groß, köstlich Heiligthum halten; ja, sie sollten viel lieber sterben wollen, ehe sie ihre lieben Eltern erzürnen wollten. Denn die Eltern haben in diesem Leben nichts, dadurch sie mehr bewegt werden, als die Liebe zu ihren Kindern; darum können sie so leichtlich verletzt werden. Wiewohl du dasselbe nicht fühlst oder verstehst, sondern bist in der Sünde und im Fleisch ersoffen, und achtest ihre Sorge und Kümmerniß, so sie deinetthalben haben, gering. Aber, wehe dir, wo du sie nicht wiederum auch in Ehren und vor Augen halten wirst. Denn also wirst du bald, ehe denn du daran denkst, ein Vater- oder Muttermörder werden, und wirst darnach vergeblich auf den Segen warten; ja, du wirst vielmehr dich des Galgens und Schwertes besorgen müssen und wirst einen greulichen Fluch auf dich laden. Denn also sagt GOttes Wort 5 Mos. 27, 16.: „Verflucht sei, wer seinem Vater oder Mutter fluchet“; item, 2 Mos. 21, 17.: „Wer Vater oder Mutter flucht, der soll des Todes sterben“; item, 3 Mos. 20, 9.: „Sein Blut sei auf ihm, daß er seinem Vater oder Mutter gefluchet hat.“

204. Esau flucht eigentlich hier seinem Vater; thut er es nicht mit Worten, so thut er es doch mit der That. Denn das meint er, da er sagt: „Es wird die Zeit bald kommen, daß mein Vater Leid tragen muß.“ In der gemeinen lateinischen Dolmetschung steht: Venient dies: „Die Zeit wird kommen.“ Aber das hebräische Wort heißt eigentlich so viel, als wenn man auf Lateinisch sagt: prope sunt: Die Zeit ist nahe, das ist, es wird nicht lange ausbleiben, daß ich die Freude meines Vaters und Mutter, dazu auch meines Bruders Jakob zerstören werde, wie sie mir meine Freude auch zerstört haben. Ich will ihnen Galle und Wermuth darunter mengen, wie man pflegt zu sagen. Das sind wahrlich Worte eines greulichen Fluches, die da genugsam anzeigen, daß er vom Teufel und höllischen Furien getrieben werde.

205. So ist nun das mit das größte Unglück im menschlichen Geschlecht, nämlich, die Thorheit und Bosheit der Kinder, damit sie ihre Eltern tödten, ehe sie verstehen, wie groß diese ihre Sünde sei. Denn sie wissen nicht, wie große Schmerzen und Bekümmernisse die Eltern haben, von wegen der Bosheit ihrer Kinder. Die Eltern aber fühlen und empfinden es wohl, und vor großem Herzeleid vergehen sie und werden darüber gar verzehret. Darum hat GOtt im vierten Gebot nicht vergeblich das Wort „ehren“ gebraucht. Denn
\quellseite{95}
er sagt nicht: Du sollst deinen Eltern gehorchen, sollst sie hören, sondern: Du sollst sie ehren. Er will den Namen der Eltern herrlich und ehrlich gehalten haben. Denn Gott kennt an uns die Bosheit der Erbsünde wohl, die so überaus kräftig ist, daß sie die Menschen treibt, daß sie ihre Eltern tödten und an ihrem eigenen Blut grausame Tyrannei üben. David ist mit diesem Unglück hart geplagt worden, und hat gelernt, welch große Schmerzen und Angst die Eltern haben von der Bosheit ihrer Kinder. Denn er hat mit der That empfunden und erfahren, daß sein Sohn Absalom sein Feind war, der ihm nach Leib und Leben, dazu auch nach ewiger Wohlfahrt und Seligkeit gestanden hat. Aber er ist darüber in einen greulichen Fluch gefallen, daß er eines solchen Todes gestorben ist, wie er es mit seiner bösen That verschuldet hat.

206. Darum vermahne ich und bitte alle Jünglinge, daß sie sich vor dieser Sünde hüten und derselbigen Feind seien, und daß sie ihre Herzen daran gewöhnen, daß sie ihre Eltern in Ehren halten, und dazu in ihrem Gebet Gott täglich und fleißig um Hülfe anrufen. Ist das in Abrahams Hause und in Abrahams Hause geschehen, zwischen Ismael und Isaak, im Hause Isaaks zwischen Jakob und Esau (ich will jetzt von vielen Exempeln der Heiden nicht sagen): was haben denn wir Eltern nicht zu erwarten und was habt ihr Kinder euch nicht zu besorgen? Denn ihr seht, daß der Teufel mitten unter uns ist. Ihr seht auch die betrübtesten Exempel der Straßen zu allen Zeiten, so über die ungehorsamen muthwilligen Kinder gehen, die ihre Eltern schmähen, und sie mit ihrem Ungehorsam martern und betrüben. Wie jämmerlich und plötzlich ist Esau dahin gefallen, ehe denn er solches betrübten Unfalls gewahr geworden, oder daran zuvor möchte gedacht haben! Und wie sollte er daran denken, dieweil er allein in seines Vaters Hause beide das Welt- und Kirchenregiment gehabt hat, ja, dieweil er allein der Herr und Fürst im Hause ist?

207. Denn dieses ist das Bild des Hauses und Gesindes Isaaks. Der alte Vater hat sich des ganzen Amtes, so dem Hausvater obliegt, und des ganzen Regimentes im Hause, dazu auch der weltlichen und Kirchenregierung gar begeben. Denn er ist nun blind, und daßelbe, meines Erachtens, nicht Alters halben, sondern irgend eines anderen Unfalls oder Anfechtung halben, auf daß solches auch noch zu seiner Trübsal und Kreuz kommen möchte. Denn er ist ein heiliger Mann gewesen, und den Gott lieb hatte, und derhalben ist er auch mit so vielen Anfechtungen geprüft und versucht worden. Dieweil er nun nicht mehr hat sehen können, so muß er vom ganzen Regiment gar abstehen, es sei denn, daß vielleicht irgend mit Worten etwas zu vermählen oder zu gebieten sei. Rebekka ist die Hausmutter, hat aber doch keine Macht oder Gewalt im Hause; desgleichen ist Jakob auch verschmäht und verachtet worden. Esau ist allein der Herr, der da zwei Weiber und viel Kinder hat. Denn Jakob wird erst im siebenunddreißigsten Jahre nach der Hochzeit Esaus, so er im vierzigsten Jahr seines Alters gehalten hat, gesegnet, und sie alle beide sind jetzt siebenundsiebzig Jahre; mittler Zeit aber hat Esau viel Kinder gezeugt.

208. Wiewohl nun Rebekka auch zum Theil mit regiert hat, wie eine Hausmutter, so hat sie doch daßelbe mit großer Beschwerung und mit großer Mühseligkeit gethan. Denn die zwei Weiber ihres Sohnes Esau waren heidnisch und gottlos, welche ihre Schwiegermutter verachtet haben als ein alt kindisch Weib; und ist sonst doch also gebräuchlich, daß alle Sohnesweiber nicht wohl leiden können, daß ihre Schwiegermütter über sie regieren. Diese zwei Weiber aber waren nicht Fremde, sondern waren daselbst im Lande daheim, welche einen fremden Mann zur Ehe genommen, nämlich Esau, der arm und verachtet war; darum haben sie sich hochfahrig und mit großer Verachtung wider Rebekka gesetzt.

209. Also hat Rebekka in großer Angst und höchster Verachtung leben müssen. Wie droben im 26. Cap. V. 35. gesagt ist: „Sie machten beide Isaak und Rebekka eitel Herzeleid“, sie betrübten und plagten für und für das gute fromme Weib. Denn sie waren die Königinnen im Hause; Esau war ein Junker, war denen vom Adel und andern großen Herren im Lande durch Verschwägerung verwandt: Jakob war elend und wie ein Fremder, mußte alle Tage erwarten, daß er würde gar verstoßen und vertrieben werden. Denn Esau gedachte, seine Kinder zu Erben zu setzen, und dazu wollten ihn ohne Zweifel die Schwäger und seine zwei Weiber auch angereizt haben. Darum ist er des Regiments und Priesterthums ganz gewiß. Denn er regieret das Haus anstatt des Vaters, als des Vaters Stellvertreter. Aber die heilige Mutter Rebekka wandelt im Schlamm großer Wasser und leidet viel Unbilliges von ihres Sohnes Weibern; denn der Vater war alt und seines Gesichtes beraubt, daß er sich des ganzen Regimentes begeben hat.

210. Lieber Gott, wie viel hat unterdes Rebekka leiden müssen! wie viel Gewalt und Schmach hat sie ertragen und in sich fressen müssen! Darum hat sie nicht ohne Ursache über Jakob gehalten, und ihn auch geschützt, daß er den Segen sollte erlangen, da Esau mit solchem Stolz und Hoffahrt im Hause regiert, und daß die Schwäger sich auch mit so großem Stolz gefreut und hoch gerühmt haben, daß sie, die Hetiter, selbst auch durch diesen Fürsten Esau regieren sollten. Da hat ohne Zweifel Rebekka vor Gott ängstiglich und brünstig ihr Gebet mit Thränen gethan, da sie gesehen hat, daß die Ehre und Herrlichkeit der Verheißung auf die Heiden kommen sollte, nämlich die Hetiter. O lieber Gott! wird sie geschrieben haben, du wollest ja mehren und verhüten, daß ihr Vornehmen nicht vor sich gehen und geraten möge.

211. Da nun das Regiment in solchem Flor gestanden und Esau ein solcher hoffärtiger Fürst, Jakob aber mit seiner frommen heiligen Mutter verachtet und verworfen war, hat Esau nicht anders gewußt und bei sich selbst gedacht, denn daß er allem göttlichen und menschlichen Recht nach König wäre. Und dazu kommt auch der Fleiß, die Kunst und Wohltat seines Vaters Isaak, welches eine Bestätigung und ein gewiß Argument ist, darauf er sich verläßt, und bekümmert sich darum nicht mehr, daß er die Erstgeburt verkauft und verworfen hatte, ist alles Dinges halben ganz sicher, gleich als ob er im Himmel und Paradiese seinen Sitz hätte.

212. Siehe aber, wie gar bald und unversehens ihn der Zorn und Gericht Gottes überfällt. Denn der Vater heißt ihn hinaus auf die Jagd ziehen und daß er ihm ein niedlich Essen bringen soll, da er hat sollen gesegnet werden. Da wird Esau ohne Zweifel seine Schwäger und andere gute Freunde zu Gefährten gehabt haben, die mit ihm auf die Jagd sind gezogen, und zieht also hinaus mit Hunden, und mit den Jägern, die sich mit ihm gefreut und ihm Glück gewünscht haben, und das Gesinde im Hause wird auch gerühmt haben: Morgen wird Isaac unsern Fürsten und Priester Esau segnen. Alles ist eitel Freude und Wonne.

213. Was geschieht aber? Esau kommt wieder heim mit großem Gepränge, mit großem Warten und Hoffnung, bringt dem Vater das Essen, so er ihm bereitet und zugerichtet hatte. Da fällt bald die ganze Sache und alle Hoffnung, dazu auch alle Freude gar dahin, und er selbst auch fällt plötzlich nieder, als hätte ihn der Donner und Blitz vom Himmel herabgeschlagen. Darum hat er große Ursache zu zürnen gehabt. Denn sollte eine solche große Freude so bald und plötzlich geändert werden? Hilf Gott, wie hat er da seinem Bruder so greulich geflucht und ihn gelästert? Wie wird sich da ein Weinen, Schrecken und Klagen erhoben haben? Wie bestürzt wird das ganze Gesinde, seine Weiber und Schwäger geworden sein? Denn der des Morgens um sechs Uhr ein Herr, König und Priester gewesen war, ist am Abend ein Knecht geworden. Also ist ihnen allen plötzlich diese große Hoffnung gar dahin gefallen, daß sie nimmermehr gehofft und gedacht hätten.

214. In solchem Fall kann das Fleisch nicht geduldig sein. Was soll ich nun thun? hat Esau gedacht; die Schmach und Schande ist so groß, daß ich sie weder dulden noch leiden kann, dieweil mir so große Unehre geschehen ist, der ich in meines Vaters Hause König und Priester bin; und mein Bruder, der bisher gar verachtet und verworfen gewesen, soll nun mein Herr sein? Da hat Esau über solchen Unfall greulich geflucht; wird gesagt haben: Ei, da schlage der Teufel in der Hölle zu und das höllische Feuer. Denn daß er also unversehens von solcher großen Hoffnung hat sollen gestoßen werden und ihm dieselbe so gar dahin ist gefallen, hat ihn von Noth wegen sehr heftig erzürnen müssen.

215. Da siehst du nun, wie das so ein groß und greulich Unglück ist, wo man also sicher dahin sündigt und nicht rechte Buße thut. Denn dasselbe hat zu diesem Unfall Esaus, der durch Gottes Zorn alles himmlischen Segens beraubt worden, Ursache gegeben: wie Lucifer vom Himmel gefallen ist, ehe er sich dessen versehen, oder daran gedacht hatte. Denn da Esau den Segen verloren hatte, wie er damals war, hat er auch zugleich damit Vater und Mutter, dazu das ganze Erbe mit dem Regiment und Priesterthum verloren. Darum wird er so gar unsinnig, zürnt, tobt und wüthet, da er sagt: „Es wird die Zeit bald kommen, daß mein Vater Leid tragen muß.“ Denn also hat er zu seinen Schwägern geredet, mit denen er gerathschlagt hat, und zum Gesinde im Hause, welches gesehen hat, daß er ganz zornig und grimmig war.

\quellseite{96}

\quellseite{97}

216. Dazu ist auch gekommen, daß seine Kinder geweint und geschrieen haben, welche gesehen, daß sie nun der großen Hoffnung, davon sie alle so stolz waren, beraubt sind. Diese haben ihn zu solchem grimmigen Zorn gereizt, werden gesagt haben: „Willst du das leiden, daß dir so viel große herrliche Güter so gar plötzlich sollen genommen werden?“ Darauf hat er gesagt: Das will ich mit nichten leiden, sondern will ihnen die Freude zerstören. Solcher Worte wird er sich ohne Zweifel gerühmt haben bei seinen Schwägern und zwei Weibern, so der Rebekka in ihren Augen haben wehe gethan, daß sie dieselben nicht wohl hat sehen mögen, welche den Esau, der allbereit unsinnig war von großem Zorn, noch mehr erhitzt und zum Zorn bewogen haben.

217. Und so hat es damals in dem Hause Isaaks, des heiligen Patriarchen, gestanden; dieses Haus ist ganz in Verwirrung gerathen durch den Rath des Heiligen Geistes, den Rebekka ausgeführt hat, nämlich, daß der Segen von dem größeren Sohne auf den kleinern möchte gewendet werden. Denn es hat eine große Verwirrung und Lärm daraus folgen müssen, da die große Hoffnung und aller Anschlag, so Esau gehabt, über einen Haufen gefallen und zunichte geworden ist; denn er hatte schon das Priesterthum und die Herrschaft im Besitz, und brauchte nicht zu befürchten, daß er daraus verstoßen werden möchte.

218. Aber dies ist ein trefflich Exempel der göttlichen Strafe, dadurch der große Fürst und Regent, der so überaus hoffärtig war, gefallen ist, also daß er schier in einem Augenblicke gar hat verzagen müssen, der doch kurz zuvor sich selbst verheissen hatte, daß er ewig regieren würde. Dasselbe sollen wir uns fleißig vorstellen. Denn so gar böse und verkehrt ist aller Menschen Art: dieweil wir sündigen, sind wir gar sicher; wenn wir aber den Zorn und die Strafe Gottes fühlen, so sind wir arme, elende und verzagte Leute, wie im Hause Esaus sehr groß Heulen und Schreien gewesen ist. Da haben die Weiber, die Kinder und das ganze Gesinde jämmerlich geheult und geweint: er aber geräth dahin, daß er gar verzagt und unsinnig wird; denn der Zorn Gottes ist ein groß und greulich Ding, daß er darüber so gar erhitzt wird, daß er gern seine Eltern tödten und die Kirche gar vertilgen wollte.

\section[Siebenter Theil]{Siebenter Theil. Wie Rebekka Esau’s Vorsatz erfährt und zu verhindern sucht, und deswegen Jakob aus seinem Hause zu ihrer Freundschaft nach Mesopotamien schickt}

\begin{quote}
\textbf{V. 42. Da wurden Rebekka angesagt diese Worte ihres größern Sohns Esau: und schickte hin, und ließ Jakob, ihrem kleinern Sohn, rufen, und sprach zu ihm: Siehe, dein Bruder Esau bräuet dir, daß er dich erwürgen will.}
\end{quote}

219. Esau hat diese Drohung nicht heimlich gehalten, sondern da seine Kinder und Weiber geweint und die Schwäger vor großem Grimm und Zorn gar unsinnig gewesen sind, hat er öffentlich gesagt: Lasset uns nur ein wenig zusehen, die Freude Jakobs und meiner Eltern soll nicht lange währen: ich will machen, daß sie auch mit uns trauern und Leid tragen sollen. Da diese Worte irgend einer unter den Knechten, der etwas frömmer und getreuer gewesen als die andern, gehört hat, ist er bald zu Rebekka gelaufen und hat ihr solches angesagt.

220. Und sie hat zwar selbst auch anders nicht denken können, denn daß er heftig würde erzürnt werden und daß er auch sehr gute Ursache hätte zu zürnen, sintemal er ganze siebenundsiebenzig Jahre in gewisser Hoffnung gelebt, daß er den Segen besitzen würde, und dieweil der Vater ihm denselben auch mit hohem Fleiß und sonderlicher Gunst zugeordnet und beschieden hatte; welches er nun alles verloren hat. Aber das ist der Rebekka dieses auch eingefallen, daß sie gedacht hat, daß sie an all diesem Unglück und greulichen Lärm in ihrem Hause Schuld wäre und wäre von ihr hergekommen; wiewohl sie in der Wahrheit daran nicht Schuld gewesen, sondern Esau hat solches mit seiner Sünde verursacht, und ist vornehmlich alles von der göttlichen Antwort hergekommen, in welcher gesagt war: „Der Größere wird dem Kleinern dienen.“ Derhalben hat sich Esau selbst durch seine eigene Schuld dieses Unglück gemacht, dieweil er Gott erzürnt und verachtet hat, dazu auch die Herrlichkeit und Ehre der Erstgeburt gar gering geachtet, und um ein Geringes, nämlich, um ein Linsengericht, verkauft hat.

221. Da stehe nun, wie viel Unglück von dieser Einen Sünde hergekommen sei; denn nach der göttlichen Antwort ist die Sünde Esaus eine Ursache dieses ganzen Lärms. Und mit diesen zwei Stücken hat sich auch Rebekka getröstet: Erstlich, mit der göttlichen Antwort, und darnach, mit der Sünde Esaus; welcher, da er auch zur Erkenntniß gekommen ist und daß er den Zorn GOttes gefühlt hat, anhebt unsinnig zu werden, aber nicht Buße zu thun.

222. Wo man nun sündigt und die Strafe fühlt, und aber doch nicht erkennen will, daß die Strafe von GOttes gerechtem Urtheile herkomme, sondern nur ungeduldig und zornig sein will von wegen des Gutes, so man verloren hat: das heißt nicht Buße thun, sondern heißt wider GOtt toben und mit ihm zürnen. Wie wir jetzt zur Zeit auch thun: wenn wir mit dem Türkenkriege, mit Pestilenz, mit Theurung gestraft, oder sonst auf andere Weise vom Teufel geplagt werden, so klagen wir alle über den großen Jammer und Unglück. Aber du hörst niemand, der da sage: Wir haben gesündigt, wir haben übel gethan, HErr GOtt, erbarme dich unser, gedenke an deine Barmherzigkeit, die von Ewigkeit währet. Wir bekehren uns nicht zu GOtt, der uns straft; wie darüber der Prophet Jesaia am 9. Cap. V. 13. auch klagt: „Das Volk kehret sich nicht zu dem, der es schlägt; und fragt nichts nach dem HErrn Zebaoth.“

223. Derhalben schreien und klagen wir alle jämmerlich über die Strafe, über die Grausamkeit und Tyrannei des barbarischen Feindes, und über ander Unglück, damit wir überfallen werden; aber wir lassen uns die Sünde nicht leid sein. Darum ist unsere Buße eben eine solche Buße, wie Esau gehabt hat. Endlich aber werden wir auch hören müssen, daß in den Sprüchen Salomos Cap. 1, 24–28. geschrieben steht: „Weil ich denn rufe und ihr weigert euch; ich reckte meine Hand aus, und niemand achtet drauf, und laßt fahren allen meinen Rath, und wollt meiner Strafe nicht: so will ich auch lachen in eurem Unfall, und euer spotten, wenn da kommt, wie ein Sturm, das ihr fürchtet, und euer Unfall als ein Wetter, wenn über euch Angst und Noth kommt. Dann werden sie mir rufen, aber ich werde nicht antworten“ 2c.

224. Denn was ist unbilliger, als daß wir von GOtt Hülfe und Errettung bitten und begehren, dazu auch Segen in leiblichen und geistlichen Dingen, und daß wir gleichwohl noch unterdes Abgötterei treiben, unsere Brüder und Eltern tödten, und dazu allerlei bösen Lüsten ihren Zaum lassen und denselben nachhängen? Das wird wahrlich ohne gewisse und greuliche Strafe nicht abgehen. Denn es steht eine solche Drohung auch im Propheten Jesaia Cap. 65, V. 12. 13.: „Darum, daß ich rief, und ihr antwortet nicht; sondern thatet, was mir übel gefiel, und erwählet, das mir nicht gefiel. Darum spricht der HErr HERR also: Siehe, meine Knechte sollen essen, ihr aber sollt hungern; siehe, meine Knechte sollen trinken, ihr aber sollt dürsten“ 2c.

225. Es haben aber auch die gottlosen und unschuldigen Menschen ihre Trübsal, Unglück und Jammer, und gerathen auch zum öftern mit in das gemeine Unglück. Denn es müssen ja allezeit unter so viel tausend Menschen oder in so großer Menge Volkes etliche Gottlose überbleiben. Wenn dieselbigen aber beten, so werden sie nicht verworfen oder verlassen, wie den Gottlosen geschieht, sondern das widerfährt ihnen, wie St. Petrus sagt 1. Epist. 4, 17., daß am Hause GOttes das Gericht anfähet. Die Gottlosen aber müssen darnach die Höllen auslaufen, Ps. 75, 9. Denn es ist ein gewisses Zeichen, wenn die Leute, die da beten und GOtt von Herzen dienen, daß den gottlosen und sichern Menschen der Untergang bevorstehe: und die Gottseeligen schmecken zwar nur ein wenig aus dem Kelche des Zornes; was will es aber für ein Ende mit denen nehmen, die dem Evangelium nicht glauben?

226. Wie wir zu unsern Zeiten auch in der Versammlung sind derer, so die Strafe und das Verderben verdient haben, und erfahren jetzt, daß das Gericht an dem Hause GOttes aufsteht. Denn wir tragen auch den gemeinen Jammer mit andern Menschen, der da schwer genug ist; wie Daniel, und die andern Göttseligen, so zu Babel in Elende waren, die gemeinen Strafen auch mitgetragen haben: aber den Tyrannen und Feinden wird viel eine größere und greulichere Strafe verkündigt; wie GOttes Wort im Propheten Jeremia am 49. Cap. V. 12. zeuget wider die Amoriter, Moabiter und Andere mehr: „Denn so spricht der HErr: Siehe, diejenigen, so es nicht verschuldet hatten, den Kelch zu trinken, müssen trinken, und du solltest ungestraft bleiben? Du sollst nicht ungestraft bleiben, sondern du mußt auch trinken“ etc.

\quellseite{98}

227. In solchem Jammer aber und Unglück soll man sich im Herzen zur Geduld schicken, daß wir das Gericht und die Strafe GOttes gut aufnehmen; denn es wird uns, die wir Buße thun, und uns der Sünden der Gottlosen nicht theilhaftig machen und ernstlich beten, nichts schaden. Aber wehe denen, so die Hefen aussaufen müssen; denn dieselbigen werden von der Strafe, so sie tragen müssen, immer ärger, und erkennen die Sünde nicht, so sie begangen haben: sie schreyen und zürnen von wegen des Unglückes, das sie leiden; die Sünde aber lassen sie sich nicht leid sein. Wenn du um deiner Sünde willen schreyest, so würde dich GOtt erhören; dieweil du aber die Sünde verhelfst, entschuldigst und dazu noch vertheidigst, so will ich, spricht der HErr, dein Geschrey, da du über die Strafe klagst, nicht hören. Dieses soll man bedenken von der Ursache des sehr betrübten und plötzlichen Falles Esaus, auf daß wir die Ursachen der Strafen verstehen und uns davor hüten lernen, und daß wir in gemeiner Noth und Gefahr immer Hoffnung behalten und GOtt anrufen, daß wir von wegen der Strafe und des Unglücks nicht verzagt werden und zu den Gottlosen fallen.

228. Nachdem nun Rebekka die Botschaft erhalten, daß Esau so überaus zornig wäre und daß er schnauete mit Morden wider seinen Bruder, da redet sie ihren Sohn Jakob an und spricht: „Dein Bruder dräuet dir, daß er dich erwürgen will“, er sucht, wie er sich an dir rächen möge; welches zwar ein sehr böser Trost

229. Darnach hat sie auch gebeten, daß sein Herz möchte zufrieden gestellt werden, und wird ohne Zweifel auch mit freundlichen Worten und Geschenken bei den Schwägern und Weibern angehalten und gesucht haben, daß sie Esau, der so überaus zornig war, wiederum versöhnen und zufrieden stellen wollten. Zudem hat sie Esaus Herz selbst auch mit freundlichen Worten angegangen und zu erweichen versucht, daß er doch so sehr nicht wüthen wollte; wird gesagt haben: Lieber Sohn, laß dir doch an dem genügen, so dir dein Vater zugeeignet hat; du hast ja daran genug, GOtt wird dir etwas Anderes und Größeres geben. Solche Anschläge und freundliche Worte hat sie gebraucht, diesen Zorn zu schlichten; denn es ist ein sehr listig und verständig Weib gewesen, aber doch Alles mit Gottseligkeit und Gottesfurcht.

\begin{quote}
\textbf{V. 43–45. Und nun höre meine Stimme, mein Sohn: Mach dich auf, und fleuch zu meinem Bruder Laban in Haran, und bleib eine Weile bei ihm, bis sich der Grimm deines Bruders wende, und bis sich sein Zorn wider dich von dir wende, und vergeße, was du an ihm gethan hast; so will ich darnach schicken, und dich von dannen holen lassen. Warum sollte ich euer beider beraubet werden auf Einen Tag?}
\end{quote}

230. Das hebräische Wort \unsicher{cheham} heißt, ein heftiger Zorn, welcher ist wie eine große Hitze. Jacham heißt: erwärmen, erhitzen; wie die Poeten pflegen den Zorn zu beschreiben: Ignescunt irae, et duris dolor ossibus ardet etc., das ist: Der Zorn brennt wie ein Feuer und der Schmerz brennt auch in den harten Gebeinen. Also ist Esau auch hitzig geworden vom Zorn. Derhalben meint Rebekka, ihr Sohn Jakob müsse vor ihm fliehen, auf daß sie nicht angesehen würde, als wollte sie ihn so leichtfertig in Gefahr bringen, wo sie ihn bei sich behalten hätte.

231. Denn ein Herz, das vom Zorn so gar erhitzt ist, ist kühn, allerlei böse Thaten anzufangen und zu begehen; wie man pflegt zu sagen: Ira furor brevis est etc.: Der Zorn ist wie eine Unsinnigkeit, die sich in einem Huy zuträgt und begehen läßt. Darum muß man solcher Unsinnigkeit weichen, daß er in zornigem Muthe nicht wider dich wüthe, und solchen Schaden zufüge, den man darnach nimmermehr heilen könnte. Dieses hat Rebekka also ganz klüglich bedacht und geredet. Wir wollen, sagt sie, GOtt nicht versuchen, und sagen: der dich gesegnet hat, wird dich auch wohl erhalten; GOtt thut wohl, was er vorgenommen hat, aber er thut es durch Mittel. So erlaubt nun Rebekka auch, daß die Kirche und der Segen wohl geschützt wird werden: sie gebraucht aber doch die Mittel, so GOtt erlaubt hat, die Gefahr dadurch zu verhüten.

232. Dieses Exempel soll man fleißig merken um derer willen, die Alles auf die Versehung ziehen wollen, und also damit alle Werke und Mittel, die GOtt verordnet hat, aufheben. Denn sie sagen also: Wenn dieses geschehen soll, so wird es nothwendig geschehen müssen auch ohne mein Zuthun. Oder wenn sie sich auch in unnöthige Gefahr begeben, versprechen sie sich doch selbst Schutz und Schirm in solcher Gefahr; denn GOtt werde sie von Noth wegen schützen und bewahren müssen laut seiner Verheißung. Dieses sind böse und gottlose Gedanken. Denn GOtt will haben, daß du der Mittel sollst gebrauchen, die du gebrauchen kannst. Er will haben, daß du die Gelegenheit, die sich darbietet, annehmen und derselben gebrauchen sollst, sintemal er durch dich das ausrichten will, was er verordnet hat. Denn also hat er es haben wollen, daß dich dein Vater zeugen und deine Mutter nähren sollte, so er dich doch wohl ohne Eltern hätte schaffen und nähren können. Also soll man in diesem gemeinen Leben die nöthige Arbeit thun: man soll säen, pflanzen, man soll sich Nahrung suchen etc., darnach wird GOtt auch thun, was er will. Wenn du aber sagen wirst: Ich will dem Kindlein nicht Milch zu saugen geben; denn wo es leben soll, so wird es ohne das wohl leben etc., so wirst du damit betrogen und dich daran schwerlich versündigen. Denn darum hat GOtt der Mutter die Brüste gegeben, daß sie damit das Kind säugen soll. Er könnte zwar das Kindlein wohl ohne Milch ernähren, er wills es aber nicht thun. Derhalben soll man die Mittel gebrauchen, so von ihm verordnet sind.

233. Also hat Rebekka gewißlich schließen und es dafür halten können: Mein Sohn Jakob wird nicht erwürget und der Segen auch nicht widerrufen werden; sie versäumt aber doch nicht, was ihr Amts halben gebührt zu thun, sondern sagt zu Jakob: „Fleuch zu meinem Bruder Laban, bis sich der Grimm deines Bruders wende.“ Es könnte dich zwar GOtt wohl schützen und vor ihm vertheidigen: man muß aber dem Teufel nicht Gelegenheit oder Raum geben, sondern man muß ihm die Gelegenheit abschneiden. Du siehst, wie jetzt dein Bruder im Hause umhergeht wie ein unsinniger Mensch, und wo er dich ergriffe, so würde er dich ohne Zweifel erwürgen; oder wenn schon der Vater ihm sein Herz erweichen würde, so muß man sich doch nicht weniger solcher Gefahr besorgen von seinen Schwägern, Weibern und Kindern, die man so leichtlich nicht wird versöhnen können; und wenn schon Esau zufrieden gestellt wäre, so werden sie doch versuchen dir Schaden zu thun, dieweil sie wissen, daß sie damit ihrem Vater Esau einen angenehmen Dienst thun. Dieweil du nun zu meinem Bruder fliehen kannst und also diese Gefahr verhüten, so gebrauche diese Gelegenheit und diesen Rath; mittler Zeit wollen wir sehen, daß wir beide, die Schwäger und unseres Sohnes Weiber, wiederum versöhnen und zufrieden stellen. Wir wollen sie den göttlichen Willen lehren und berichten, daß GOtt dieses also versehen und verordnet hat. Weiche du nur eine kleine Weile. Ich will allen möglichen Fleiß thun, daß die Herzen mögen erweicht werden. Ich will zu ihnen gehen, will mit ihnen reden und will GOtt um Gnade bitten.

234. Dieses ist eine große treffliche Klugheit an dem sehr frommen heiligen Weibe, dadurch sie in dem großen Kärm, den sie erregt hatte, und in solcher großen Noth und Gefahr sich selbst und ihrem Sohne so fein hat rathen und helfen können; denn es ist beide, das Haus und die Kirche Isaaks, sehr heftig betrübt und verwirrt, und die Herzen, so von rechtschaffigem Schmerz erbittert gewesen, haben nicht so bald können erweicht und versöhnt werden. Das sieht Rebekka wohl, und versucht doch gleichwohl GOtt nicht, verzagt auch nicht. Denn man muß beiderlei thun, nämlich, daß man GOtt in der Noth und Gefahr nicht versuche und auch an der Hülfe nicht verzage.

\quellseite{99}

235. Dieses gehört und dient zur Lehre und zum Trost der Kirche, deren Bild uns hier vorgehalten wird. Denn sie steht allezeit in sehr großer Gefahr, und muß immer solches Unglück erwarten, daraus sie nach menschlicher Vernunft nicht kommen kann. Und dieses ist ihr gemein Wort und rechter Titel, daß sie sagt: „HErr, hilf uns, wir verderben!“ Matth. 8, 25. Sie verdirbt aber doch nicht. Sie wird stets von den Wasserwogen und dem Ungestüm des Meeres hin und her getrieben, daß sie vor ihren Augen Anderes nichts sieht denn eitel Verderben, und daß wir schreien: „Wir verderben“, und ist doch gleichwohl bei GOtt noch immer Hülfe. Wie Rebekka und Jakob in solcher großen Beschwerung auch an der Verheißung GOttes hangen und endlich sich herausringen und den Sieg behalten.

236. Danach soll man aus diesem Exempel auch das lernen, daß wir nicht alsobald verzagen an den Menschen, die da gesündigt haben, sondern daß wir immer an die Worte Christi Joh. 11, 8. 9. denken, damit er seinen Jüngern antwortet, da sie ihm widerrathen wollten, daß er nicht wieder nach Judäa ziehen sollte, darum daß ihn kurz zuvor die Juden hatten steinigen wollen: „Sind nicht des Tages zwölf Stunden?“ das ist, die Herzen können immer mit der Zeit gelindert und erweicht werden, daß sie ihre Bosheit fallen lassen und sich zur Besserung schicken: mit der Zeit vergeht der große Grimm und der hitzige Zorn.

237. Und ist zwar wohl glaublich, daß Esau viel gelinder geworden sei in den zwölf oder mehr Jahren, während Jakob in Mesopotamien geblieben ist. Und ist auch dazu gekommen, daß er zeitlich gut Glück gehabt hat, daran ihm nichts abgeprochen ist; ja, es ist seine gottlose Vermessenheit immer gemehrt worden, daß er nichts desto weniger den Segen behalten würde, sintemal gar das Gegentheil von dem geschehen, so der Segen Jakobs mitgebracht. Denn derselbe war wohl zum Erben des Regimentes und Priesterthumes gesetzt: er hat aber wenig Glück dabei; denn er wird in das Elend getrieben, und sein Glück und ganzes Leben läßt sich ansehen, daß es einem Fluche ähnlicher sei als einem Segen. Esau aber bleibt im Besitz, dieweil Isaak, da er seines Gesichtes beraubt war, sich des Regiments über das Gesinde im Hause und über die Kirche begeben hatte. Darum bleibt Esau der Herr im Hause in Abwesenheit seines Bruders Jakob, und richtet noch immer seine gewöhnliche Arbeit aus, wie er bisher gethan hat.

238. Da dieses die zwei Schwiegertöchter und Schwäger gesehen haben, da sind sie auch leichtlich zufrieden gewesen, und haben gedacht: Laß den Segensaten immer hingehen; wir bleiben im Hause, er aber muß flüchtig werden. Und ist dieselbe Zuversicht sehr gemehret worden, da Jakob so viele Jahre ausgeblieben ist. Vielleicht, werden sie gesagt haben, ist der alte Isaak gar kindisch geworden, daß er nicht gewußt, was er gethan hat, wird von der Närrin, der Rebekka, also mit List betrogen worden sein. Dieser Segen wird GOtt nicht gefallen haben. Und ist ein großes Zeichen, daß er nicht gesegnet muß sein, dieweil er im Elende umher ziehen muß. Wir sind die Weiber des Fürsten und des Priesters, Isaaks Sohnes Weiber, denn die Verheißung geschehen ist; derhalben fragen wir nach dem nichts, der hinweg gelaufen ist. Also werden sie ohne Zweifel alle gedacht haben, was Esaus Familie gewesen ist. Denn wo Glück vorhanden ist und daß dazu ein falscher Wahn kommt des Segens halben, so sagt man alsbald: GOtt ist mir gnädig; das Glück aber und ein gottloser Wahn machen die Leute hoffärtig und sicher. Also kann es wohl zugegangen sein, daß Esau auch mit seinen Weibern und Kindern ist etwas erweitert und gelindert worden, da er gedacht, ja, dessen sei sich selbst ganz gewiß gewesen ist, daß der Segen Jakobs vergeblich wäre; sintemal Jakob nicht durch ein sonderlich Wunderwerk GOttes ist gesegnet worden, sondern GOtt hat ihn verlassen, ließ ihn aus dem Hause gestoßen werden auch durch seine eigene Mutter und Vater selbst. Darum, sage ich, hat Esau nicht anders schließen können, denn daß diesen Segen GOtt damit aufgehoben und zunichte gemacht hätte.

239. Wir sollen aber nicht meinen, daß dieses ohne schwere Anfechtung und sehr großen Schmerz der Eltern geschehen sei. Es ist wahrlich der Glaube Rebekkas abermals angefochten worden. Denn das hat ihr sehr wehe gethan, daß Esau in Abwesenheit ihres Sohnes Jakob bei der vorigen Herrschaft geblieben ist; wiewohl solcher frommen gottseligen Eltern Affecten oder Bewegungen in ihrem Herzen nicht genugsam können ausgelegt oder bedacht werden. Mein Glaube würde wahrlich so viel nicht vertragen können. Rebekka ist jetzt eine Mutter von sieben und siebenzig Jahren; denn so viel Jahre sind es, daß die göttliche Antwort ist eröffnet worden, da sie die Zwillinge noch in ihrem Leibe gehabt: „Der Größere wird dem Kleinern dienen.“ Diese Jahre über hat sie mit großen Sorgen auf den Segen gewartet und Jakob mit ihr, der darum im ehelosen Stande ohne Weib und Kind gelebt hat als Einer, der gar verworfen und verstoßen war. Unterdeß hat Esau im vierzigsten Jahre seines Alters zwei Weiber genommen und hat jetzt Kinder und Kindeskinder; er ist Herr und Fürst, und Priester im Hause. Dieses alles hat Rebekka täglich leiden und ansehen müssen eine solche lange Zeit, und bleibt dennoch beständig, da auch die zwei Schwiegertöchter das Regiment gehabt, welche sie alle beide erbittert und den alten Isaak und die rechte Hausmutter wohl geplagt haben. Denn Rebekka hat ihnen müssen die Schlüssel und das ganze Regiment übergeben. Derhalben siehe, ob dieses nicht ein trefflicher Glaube und Hoffnung sei der Rebekka, die mit solcher Langmuth auf den Segen gewartet hat.

240. Wie geht es aber Jakob? Dem ist neben anderen Anfechtungen und Plagen ein solcher Jammer und Kreuz auferlegt worden, daß schier unleidlich gewesen ist, nämlich, daß er im Alter seines Gesichts beraubt worden ist. Der alte Mann hat in solchem ehrlichen Alter die Sonne nicht sehen können, und andere Dinge, so ihm beide nöthig und lustig wären anzusehen. Er hört bisweilen Klage von Rebekka, seinem Weibe, von Jakob, seinem Sohne, oder aber sonst vom andern Gesinde über das hoffärtige Regiment und Gewalt der Esauiten: darüber mag er sich betrüben und klagen, aber bessern kann er es nicht. Und aus dem, so vorhergeht, kann man schließen, was hernach hat folgen müssen, nämlich, daß Esau, sein Sohn, gar unartig geworden sei, und daß er mehr Lust gehabt zu der Religion und Leben der Heiden, als zu der Gottseligkeit und Lehre seines Vaters, und daß die Schwiegertöchter, weil sie Heiden waren, auch der rechten Lehre abgeneigt, und dazu sehr stolz und frech gewesen sind von wegen der Herrschaft ihres Mannes.

241. Derhalben ist der Glaube dieser frommen heiligen Eltern also angefochten und geübt worden, daß es unsere Herzen weder begreifen noch mit Gedanken erreichen können; und dasselbe hat nicht zehn oder zwanzig, sondern mehr denn hundert Jahre gewähret. Lieber GOtt, wie ist das so ein großer Jammer, daß ein solcher großer Mann, den GOtt also geliebt, daß er ihn gesegnet und ihm das Erbe des Himmelreichs und des ewigen Lebens verheißen hat, so viel Jahre soll in Finsterniß sitzen, Leid tragen und allerlei Plage leiden: daß, sage ich, GOtt diesen Mann läßt seines Gesichtes beraubt werden, und daß er dazu mit so viel und großer Beschwerung muß gequälet und geplagt werden. Warum wollen wir denn nun murren und ungeduldig sein? Warum wollen wir mit GOtt zürnen, wenn uns irgend ein Unglück, ein Schade oder Unfall widerfährt?

242. Es ist zwar das Leiden der Patriarchen nichts, wenn man es mit dem Leiden Christi vergleichen und dagegen halten wollte: es bewegt aber doch unsere Herzen mehr, als das Leiden Christi thut. Denn das Leiden Christi erkennen und verstehen wir so vollkommen nicht; wir können auch nicht glauben, daß er so gar ein armer, elender und verachteter Mensch gewesen sei, wie er doch in Wahrheit gewesen ist. Darnach, so denkt jedermann, dieweil Christus wahrer GOtt ist, habe er leichtlich Alles leiden und überwinden können. Nun ist aber kein Mensch, der sich mit Isaak vergleichen könnte und der solche Geduld und Beständigkeit haben könnte im Unglück und Widerwärtigkeit. Es ist wahrlich der Glaube dieser Eltern wohl geprüft und bewähret worden und ist köstlicher und reiner erfunden denn Gold; und in solchem Glauben und Hoffnung haben sie endlich erlangt, was sie nur begehrt haben.

243. Das im lateinischen Text steht: Diebus paucis, und wir verdeutscht haben: „eine Weile“ etc., das lautet nach dem Hebräischen, als wenn man im Lateinischen sagen wollte: diebus unis, auf Deutsch: ein oder zwei Tage. Am Ende setzt Rebekka die Ursache auch hinzu, warum sie ihn heiße hinweg ziehen, da sie sagt: „Warum sollte ich euer beider beraubet werden auf Einen Tag?“ Denn dieweil sie ein sehr klug Weib war, hat sie wohl gesehen, wo Esau seinen Bruder erwürgen würde, daß er alsdann auch am Leibe würde gestraft werden nach dem Gesetze, so insgemein von Anfang gegeben war, als ein Todtschläger: daß er in das Elend getrieben würde wie Kain, oder aber durch einen Menschen, das ist, durch die Obrigkeit, sollte getödtet werden.

\quellseite{100}

244. Es läßt sich aber ansehen an diesem Ort, als ob Rebekka an der Wahrheit und Verheißung des göttlichen Segens gezweifelt habe, dieweil sie ihren Sohn aus Furcht vor seinem zornigen Bruder wegschickt. Denn warum behält sie ihn nicht bei sich, dieweil sie doch weiß, daß er gewißlich zu erwarten hat, daß ihn GOtt erhalten werde?

245. Darauf antworte ich, wie droben: Es hätte zwar Rebekka wohl sagen können: Du sollst bei mir bleiben, sollst dich vor deinem zornigen Bruder Esau nicht fürchten; denn der Herr wird für dich sorgen, der wird dich bewahren, wie der Segen lautet. Aber das thut sie nicht. Denn sie will uns mit ihrem Exempel lehren, daß man GOtt in Allem glauben und auf ihn hoffen soll; aber versuchen soll man GOtt nicht. Der Glaube ist sehr gewiß und fest um des Wortes und Verheißung willen, die er hat, und sie hätte also schließen können: Mein Sohn hat den Segen hinweg, es mag nun geschehen, was geschehen soll; er wird nicht verderben können, ich will ihn daheim bleiben lassen. Wie zuvor gesagt ist von den gottlosen und heidnischen Menschen, die alle Dinge auf eine nothwendige Ordnung und Schickung Gottes legen wollen. Aber das fromme, gottselige und verständige Weib gebraucht die Mittel und Hülfe, so vorhanden sind und die GOtt verordnet hat.

246. Wie man nicht sagen soll: Ich will nicht essen, ich will nicht trinken; soll ich leben, so werde ich doch wohl leben: wo ich aber nicht leben soll, so wird es mich auch nicht helfen, wenn ich schon esse und trinke. Item: ich bin ein Mann, derhalben werde ich auch müssen ein Vater werden, wenn ich schon kein Weib nehme. Welcher also gesinnt ist, wer wollte den nicht für einen unsinnigen Menschen halten? Denn du sollst der Gaben Gottes, die vorhanden sind, gebrauchen und sollst es nicht auf die Versehung oder Verheißung legen. Wir sollen nicht disputiren ohne Wort, was für einen Ausgang oder Ende die Dinge nehmen werden, und sollen aber doch an der Verheißung nicht zweifeln. Wo keine Verheißung ist, da soll man auch nichts anfangen, sondern man soll dem Exempel der Rebekka folgen, die sich an diese Verheißung hat halten und darauf verlassen können: Jakob, mein Sohn, wird leben; denn er soll ein Vater werden der Nachkommen. Sie verachtet aber das Mittel nicht, so vorhanden war, ja, sie gebraucht desselben, auf daß ihr Sohn der Gefahr entfliehen möge. Ja, das heißt recht glauben und aber doch GOtt nicht versuchen.

247. Wo man aber der Verheißung Gottes nicht glaubt und Gott versucht, da ist Eines wie das Andere und man sündigt an beiderlei gleich schwerlich. Denn GOtt will haben, daß wir der Creatur gebrauchen sollen, die er gegeben hat, und er hat sie auch darum gegeben, daß wir sie gebrauchen sollen. Derhalben soll man ja Gott nicht versuchen, wie die Juden in der Wüste Gott versuchten und wurden darüber greulich gestraft; welch Exempel vom heiligen Paulus auch angezogen wird 1 Cor. 10, 9. Und man soll die Unnüzigkeit der Schwärmer, die sehr schädlich ist, verdammen, welche doch heutiges Tages weit um sich greift, wie zuvor gesagt ist. Denn solche Schwärmer verkehren und verderben die ganze Gottseligkeit, dieweil sie ohne Wort Alles auf die Verheißung legen wollen.

248. Es versendet aber Rebekka Jakob nicht allein darum, daß er dem grimmigen Zorne seines Bruders entfliehen möge, sondern auch darum, weil sie gedenket, den Sohn mit einem frommen Eheweibe zu versorgen, wie hernach folgen wird. Denn sie denkt nicht also, wie die Schwärmer pflegen: Er hat den Segen hinweg, er bleibe nun im Lande oder nicht, ist gleich viel; er nehme ein Weib oder nicht, so wird er doch ein Vater des Samens sein und wird für und für Nachkommen haben. Oder wenn er schon eine Cananitin zum Weibe nehmen wird, welche ich dieser Ehe nicht werth achte, so wird doch der Herr wohl verschaffen, daß sie eine Mutter werde, die des Kindes werth sei. Also, sage ich, denkt sie nicht, sondern sie umgeht die Cananitinnen, welche sie allezeit verabscheut hat, und geht damit um, so viel als sie nur vor Staat kann und vermag, daß Jakob ein Weib nehmen soll aus ihrem Geschlecht.

249. Derhalben soll uns dieses Exempel Rebekkas ermahnen, daß wir in allen unsern Sachen und Werken thun, was an uns ist, das ist, daß wir der ordentlichen Mittel bei der Verheißung gebrauchen und darnach GOtt den Ausgang und die Versehung befehlen. Denn du hast Gottes Wort und Gebot, da sollst du wissen, daß dir gebühren will, nach demselben dein Leben zu richten, daß du nicht irgend anlaufest. Die aber außerhalb der Verheißung umher schweifen und sich selbst aus göttlicher Verheißung einen sonderlichen Ausgang dichten, die verlieren beide, das Gebot und die Verheißung Gottes. Das heißt aber Gott verachten, der das Gebot und die Verheißung gibt.

250. Auf daß nun Rebekka diese große Sünde meide und sich davor hüten möchte (denn es war ein wohlgeleibtes Weib), wiewohl sie jetzt schon der Verheißung gewiß war, so gebraucht sie doch die Mittel, die man hat gebrauchen sollen, dieweil sie dafür Sorge trägt, was es für einen Ausgang nehmen möchte. Sie geht zu ihrem Manne, unterrichtet ihn, woher ihr Sohn ein Weib nehmen solle. Da entsendet Isaak seinen Sohn, daß er also Gelegenheit hätte, sich vor den Cananitinnen zu hüten, und ein ander Weib zu nehmen, so aus seinem Stamm geboren war. Daß sie ihm aber verheißt, sie wolle darnach schicken und Jakob von dannen lassen wieder holen, desselben gedenkt Mose hernach nicht mehr, und läßt sich ansehen, als ob sie innerhalb der zwanzig Jahre, darin Jakob dem Laban gedient hat, gestorben sei; denn allhier wird der Rebekka zum letztenmal gedacht in dieser ganzen Historie.

\begin{quote}
\textbf{V. 46. Und Rebekka sprach zu Isaak: Mich verdrießt zu leben vor den Töchtern Heth. Wo Jakob ein Weib nimmt von den Töchtern Heth, die da sind wie die Töchter dieses Landes; was soll mir das Leben?}
\end{quote}

251. Allhier magst du abermal sehen, wie klug und listig dieses fromme gottselige Weib gewesen ist; denn sie zeigt dem Vater die Ursache nicht an, warum sie ihren Sohn wegschicken wolle; sie zeigt ihm nicht an, wie im Hause so ein groß Unglück und Betrübniß sei, nämlich, daß der eine Bruder den andern erwürgen wolle: aber mit gottseliger List verhehlt sie diesen Zorn vor ihrem Mann. Daraus soll man aber eine Lehre nehmen, die äußerlich zum guten Leben dient, welche wahrlich sehr gut und schön ist, nämlich, daß wir lernen friedsam sein, und daß wir, was andere Leute übel geredet und gethan haben, kein einfältig, treulich und aufrichtig auslegen, dasselbe beschönigen, zudecken, das Beste zu den Sachen reden, daß man eine Sache, die auch sehr böse ist, nicht verschlimmere, sondern dieselbige entschuldige und lindere; wie diese Sache ist, da ein Bruder dem andern nach Leib und Leben steht. Die Sache nimmt Rebekka auf sich allein, legt sie auf ihren Hals, und hilft derselben also, daß beide der Anschlag und das gottlose Vornehmen verhindert und der gute, fromme, alte Isaac in seinem Alter nicht betrübet werde. Denn wo es Isaac sollte erfahren haben, hätte er vor Traurigkeit des Herzens verzehret und getödtet werden mögen; denn es ist glaublich, wie wir droben aus der Historie angezeigt haben, daß er in so viel Jahren große Unbilligkeit gelitten habe von seinem Sohn Esau, der mit so großem Stolz regieret hat. Wenn nun dazu dieses große Unglück gekommen wäre, so hätte er vom Leide und Bekümmerniß seines Herzens gar müssen verzehret werden.

252. Darum zeigt ihm Rebekka eine andere Ursache an, warum Jakob soll hinweg ziehen, nämlich, der Heirath halben, daß sie gerne haben wollte, daß er sich nicht in dem Lande mit den Töchtern Heth verheirathen sollte, sondern mit einer andern, die aus ihrem Geschlechte geboren wäre. Denn ein gottseliger und gläubiger Mensch soll eine heilsame Zunge haben, wie in den Sprüchen Salomos steht, am 12. Cap. V. 18.: „Wer unvorsichtig heraus fährt, sticht wie ein Schwert; aber die Zunge der Weisen ist heilsam.“ Derhalben sollen wir uns daran gewöhnen, daß wir dem Exempel der Rebekka nach in bösen Sachen gute Mittler sein mögen, und daß wir dem, der da gesündigt hat und der da schuldig ist, nicht allzuviel Schaden thun; wie Rebekka ihres Sohnes Esau fleißig verschont. Sie will ihn vor dem Vater nicht verklagen, und sie schont auch damit zugleich ihres frommen Mannes, der alt und blind war, daß sie ihn nicht irre mache. Sie will viel lieber selbst dieses Unglück tragen und in sich fressen.

\quellseite{101}

253. Wenn sie wäre ein böses Weib gewesen, so hätte sie nur flugs in das Feuer geblasen und hätte Oel und Pech in das Feuer gegossen. Denn also pflegen die Leute zu thun, die so gar böser Natur und zu lästern und fluchen so gar geneigt sind. Aber die Zunge des Gerechten ist eine Zunge des Lebens; denn er redet Gutes von GOtt, und gebraucht seine Zunge, Hader und Uneinigkeit zu schlichten, den Zorn zu lindern und zu stillen, die betrübten Herzen aufzurichten und zu stärken. Wie ein schön Exempel gesagt wird von Monica, des heiligen Augustins Mutter. Wenn etwa unter den Weibern in der Nachbarschaft Zank oder Hader entstanden war, und sie hörte, daß auf beiden Seiten viel böse und bittere Worte geredet wurden, so hat sie dieselbigen nicht ausgebreitet, sie auch nicht von einem Theile zum andern gebracht, sondern mit Fleiß verhehlet, und unterdes erdichtet und sich vernehmen lassen, wo sie zu derselbigen Weibern einer gekommen ist, die sich unter einander gehabt, wie die Andere ihrer in Ehren und in allem Guten gedacht hätte. Also ist dieselbige etwas gelinder und sanftmüthiger geworden, davon die Andere gar nichts gewußt hat, daß sie allen Zorn und Haß aus dem Herzen hat fallen lassen.

254. Also sollen auch frommer gottseliger Menschen Zungen heilsam sein: es sollen heilende, lindernde und barmherzige Zungen sein, wie der Prediger auf Lateinisch sagt: Linguae curationis, mitigationis et misericordiae etc. Aber eine bittere und lästerliche Zunge erweckt oft aus einem Tröpflein und aus einem Fünklein ein großes Meer und ein großes Feuer. Darum lassen uns dieses Laster fliehen, und lasset uns dem Exempel Rebekkas folgen und der Monica, auf daß wenn wir auch gleich das Böse hören, wir doch dazu das Beste reden und zum Besten deuten. Es sei denn, daß es also vorfallen würde, daß Einer in Gefahr Leibes und Lebens kommen sollte von wegen der Feindschaft, so ein Anderer gegen ihn hat, der sich nicht will versöhnen lassen: da soll man den, so in Gefahr ist, warnen, daß er sich vor der Gewalt oder heimlichen List seines Widersachers hüten könne. Wie der sehr wohl gethan hat, der Esau bei seiner Mutter Rebekka angegeben hat dieses Streites halben, daß er vorgenommen hätte, seinen Bruder zu erwürgen: und ist eben damit eine heilsame Zunge gewesen, dieweil er den hat heißen aus dem Wege ziehen und weichen, dem Gefahr bevorstand.

255. Derhalben soll man nicht verschweigen, wo Einem also heimlich nachgestellt wird, daß man ihn erwürgen will, sondern man soll solch Ding offenbaren. Dieser Bote, welcher der Rebekka angesagt hat, daß sich Jakob vor seinem Bruder zu befahren hätte, hat nicht sagen sollen: Esau ist wohl zufrieden, hat nichts Böses im Sinn, er hat seinen Bruder lieb etc.; sondern er hat den Teufel völlig beschuldigen und aufdecken sollen, der am andern so fein Fürsten angereizt und getrieben hat, daß er den Mord hat begehen sollen. Denn wo er solches nicht angezeigt hätte, so hätte er dem Mörder geholfen; dieweil er es aber gesagt hat, so hat er damit den Mord verhütet und verhindert.

256. Dieses ist eine Lehre, so zum äußerlichen Leben dient, davon an diesem Ort nützlich geredet wird und die wohl zu merken ist. Denn wir sehen in diesem Exempel beiderlei, nämlich, daß verhütet werden die heimlichen Nachstellungen und der Tod, so Esau zugedacht war, und wird dazu auch dieses alles durch wunderbare Klugheit zugedeckt und heimlich gehalten, daß es nicht vor den alten Vater gekommen ist, daß Esau seinem Bruder gedroht hätte, ihn zu erwürgen. Und ist wahrlich eine sehr gute Ursache, welche Rebekka anführt. Denn Jakob war schon zum Erben gesetzt, der Segen war ihm verliehen und die ganze Kirche war ihm befohlen. Nun ist noch übrig, hat sie zu Esau gesagt, daß man ihm einen Stecken in die Hand thue, daran er sich halten möge, das ist, daß er möge ein Ehebund in sich haben, die seine Gehülfin sei. Das ist ehrlich und wohl geredet, nicht allein, daß solches Jakob nütze und ehrlich, sondern auch nöthig sein wolle. Denn wo er den Segen gebrauchen soll, so muß er ein Weib haben. Er hat bisher ohne Weib gelebt, da er vermodern und nicht mit in der Zahl war; nun aber ist Esau verworfen; derhalben muß man Jakob ein Weib geben, welcher den Segen erlangt hat und hernach das ganze Regiment haben wird.

257. Mit solcher Rhetorik überredet sie Isaak, daß er den Sohn ziehen läßt. Denn es ist ein Argument, damit sie bewährt, daß es nöthig sei, daß ihr Sohn hinweg ziehen müsse; welches alle anderen Argumente übertrifft, welche da hergenommen sind ab utili et honesto, oder vom Nützlichen und Wohlständigen. Wie im Gellius Metellus sehr treffend geredet hat: Wo wir wollen Bürger haben, so müssen wir auch Weiber haben etc. Denn um der Nachkommen willen muß man den Ehestand erhalten und denselben mit nichten fallen lassen.

258. Ja, möchte Einer zu Rebekka gesagt haben, warum vertraust du ihm nicht eine von den cananitischen Jungfrauen? Darauf antwortet sie und sagt: Die gefallen mir nicht. Und gebraucht nun eine andere Rhetorik ab utili et honesto, nämlich, daß sie sagt: Sie müsse auch Achtung geben auf dasjenige, so ihm nütze und geziemend sei. Ich will nicht haben, daß er eine solche nehme, wie diese unsere Schwiegertöchter sind, die zwei Weiber Esau’s. Mein lieber Isaak, du hast wohl gesehen, wie viel wir diese dreißig Jahr lang gelitten haben von den Weibern Esau’s, wie sie nicht unter die Füße getreten, sich aber verachtet haben, und haben sich zwar auf das allerhoffährtigste, ja, auch tyrannisch genug gehalten. Derhalben, wo du mich willst lebendig behalten, so verbiete ihm, daß er ein heidnisch Weib nehme von den Töchtern dieses Landes oder andern Cananitinnen; denn sie sind uns feind, dieweil sie hören, daß wir die Verheißung dieses Landes haben: das können sie nicht leiden, und begehren derhalben, uns mit Gewalt und allerlei Unrecht zu vertreiben.

259. Die Schrift gedenkt ferner Labans, ihres Bruders, nicht mit Namen, zu dem sie ihren Sohn Jakob schicken wollte; vielleicht mag aber Isaak gefragt haben: Woher wollen wir denn unserem Sohn eine Jungfrau nehmen? Da hat sie geantwortet: Laß ihn zu meinem Bruder Laban ziehen. Demselben Rath hat der alte Isaak gefolgt, und hat ihr Alles ganz und gar heimgestellt, daß sie es nach ihrem Willen machte. Denn Rebekka regiert jetzt das Haus, aber mit großer Beschwerung und Gefahr, da Alles so verwirrt war und übel stand.

260. Also wird nun nach solchem gehaltenen Rathschlage Jakob in das Elend getrieben. Der, so schon gesegnet und zum Erben gesetzt ist, wird aus dem Hause gestoßen und flieht vor seinem zornigen Bruder. Ei, wie ist mir das so ein schöner Segen und wie haben die großen Verheißungen so einen schönen Anfang! Denn Jakob hat mögen denken: Ich bin zum Erben gesetzt und werde doch verstoßen, und der, so verstoßen ist, bleibt und hat im Hause das Regiment über Alles; wie reimt sich das mit der Verheißung und mit solchem reichen Segen?

261. Derhalben ist dieses der wunderbaren Exempel eines der göttlichen Regierung, damit GOtt anzeigt, daß er fordere und haben will, daß man seinem Worte und Verheißung vertrauen soll, wenn schon das Gegentheil von dem geschieht, so in der Verheißung begriffen wird, auf daß wir uns daran gewöhnen, daß wir GOtt vertrauen in den Dingen, die nicht gegenwärtig vorhanden und noch weit vor unsern Augen verborgen sind. Denn also hat Jakob den verheißenden Segen, aber nach dem Glauben, welcher eine gewisse Zuversicht ist des, das man hoffet, und nicht zweifelt an dem, das man nicht siehet. Also glaube ich GOtt, der mir verheißt, daß er mich liebe und Achtung auf mich habe, daß er für mich sorge und mich erhöre: und dasselbe habe ich gegenwärtig und ist Alles vorhanden, wiewohl man es nicht siehet. Derhalben so lebt Jakob allein im Glauben, wird jämmerlich verstoßen, muß einsam und arm sein, und hat nichts in seiner Hand oder Gewalt, denn nur einen Stecken oder Stab und ein Stück Brodes im Beutel.

262. Das ist der Anfang des Segens; denn was im Glauben angefangen wird, dasselbe hat man noch nicht, sondern man hoffet es. Wie uns GOtt das ewige Leben verheißen hat, hat uns die Absolution und die Taufe gegeben. Diese Gnade habe ich durch Christum gegenwärtig; ich warte aber auf das ewige Leben, so mir im Wort verheißen wird. Die von dem Wort leben, sind heilig und selig; die Gottlosen aber leben nur vom Brode und nicht vom Worte, darum glauben sie auch kein ewiges Leben, warten oder hoffen auch darauf nicht. Jakob hat wohl sieben und siebenzig Jahre auf den künftigen Segen gewartet; da er denselben nun erlangt hat, muß er in das Elend ziehen, und hebt sich also sein Regiment und Priesterthum an mit einem sehr großen Kreuze, mit größestem Jammer und äußerster Armuth. Er muß seiner freundlichen lieben Eltern entrathen, und die armen Eltern müssen auch ihres allerliebsten Sohnes so eine lange Zeit entrathen und ihn nicht bei sich haben.

\quellseite{102}

263. Wenn aber Einer dieses nur obenhin ansieht und hört, meint er, es sei leicht und gering Ding; aber in der Erfahrung lernt man es, wie es so ein schwer Ding und so voller Anfechtung sei, wo man Vater und Mutter, Segen und Erbe verlassen soll, und in das Elend, da eitel Jammer und Armuth ist, fliehen. Denn dies ist die wunderbare Regierung Gottes, welche das Fleisch mit nichten vertragen kann; denn sie besteht im Glauben.

264. Dieses wird aber uns zum Exempel geschrieben, daß wir lernen an dem unsichtbaren Gott hangen, und uns daran genügen lassen, daß wir allein dieses unsichtbaren und unbegreiflichen Gottes begreifliches Wort haben, und daß wir unser Leben also anrichten, daß wir von unserem unsichtbaren Schöpfer nichts haben, denn sein Wort und die Sacramente; deßgleichen auch unsere Eltern und Obrigkeit, durch welche dieses Leben nach dem Worte regiert wird; und sollen in Hoffnung und Langmüthigkeit auf die Verheißung warten. Denn GOtt wird nicht lügen und uns nicht betrügen. Das Fleisch glaubt zwar schwerlich; denn es hat sich gewöhnt an die Dinge, die da gegenwärtig sind, und läßt sich solche Dinge bewegen, die es fühlt und vor Augen sieht. Es muß aber das Fleisch gekreuziget und getödtet, und von den Dingen, die man mit äußerlichen Sinnen empfindet, abgezogen werden, und muß lernen, daß es leben und seinen Wandel führen könne nach den Dingen, die man nicht sieht und die man mit äußerlichen Sinnen nicht begreifen kann. Das ist die rechte Tödtung des Fleisches, welches immer fleischlich gesinnet ist, und will einfach ganz und gar sicher sein, will sich nur an die Dinge halten, so gegenwärtig und sichtbar sind: darum, wo es das Gegentheil fühlt, so ist ihm angst und bange dabei.

265. Darum soll man den Leuten dies Exempel vor Augen stellen, wie Jakob zum Könige und Priester gesetzt und wie er zum Regiment und Priestertum geweihet werde; denn dieses ist zumal ein erbärmlich Gepränge und Weise, einen König zu salben und einzuführen. Man ziehet ihn kein Feierkleid an, man schmückt ihn mit keinem königlichen Hute oder Krone, man gibt ihm kein Scepter in die Hand: sondern man rüstet ihn mit einem Beutel und Stabe und wird damit also in das Elend getrieben; der Segen aber bleibt bei seinem Bruder Esau, welchem er doch mitnichten gebührt hat. Endlich aber, da Jakob also durch den Glauben an den unsichtbaren Gott getödtet ist, folgt auch das unsichtbare Gut und besitzt der Same Jakobs das Land, und Christus wird von dem Samen geboren, der ewige König und Priester, deß Reich und Priestertum in diesem Segen begriffen wird.

266. Also ward David auch zum Könige gesalbt über Israel, da Saul von GOtt abgesetzt und verworfen worden ist, 1 Sam. 16, 1. 13. ff. Aber dem allen zuwider bleibt gleichwohl Saul noch im Regimente wie zuvor: er richtet seine Sachen daheim in Frieden und draußen wider die Feinde glücklich aus; ja, er verfolgt noch über das den David ganz greulich als einen Aufrührer und der ihm mit List heimlich nach dem Königreich stünde, und David muß unterdeß in der Irre gehen wie Einer, der da flüchtig und gar elend ist in dem Lande, darin er von GOtt zum Könige gesetzt und erklärt war. Wie die Historie das bezeuget, daß er ganze zehn Jahre in der Irre umher gezogen sei, da ihn Saul täglich verfolgt und ihm heimlich nachgestellt hat. Dünket dich das ein Königreich zu sein, wo Einer in seinem eigenen Lande ganze zehn Jahre lang fremd und elend sein muß? Was könnte wunderlicher und unbilliger sein, denn daß Einer zum Könige soll erklärt und bestellt werden ohne Königreich, ohne Scepter, ohne gewisse Wohnung; ja, daß er nicht allein des Landes und Königreiches, sondern auch seines eigenen Hauses, Weibes und Kinder beraubt wird?

267. Aber das Exempel Christi übertrifft dieses alles. Ist der nicht der Allerverachtetste und Unwertheste geworden? wie der Prophet Jesaia Cap. 53, 3. sagt. Da er will gen Himmel fahren und in seine Herrlichkeit gehen, da er den Tod, die Sünde und den Teufel überwinden will, wird er an das Kreuz geschlagen; er stirbt und wird begraben als der Verachtetste unter allen Menschen und Teufeln. Heißt das in die Herrlichkeit gehen, oder den Tod überwinden und wider ihn triumphiren? Ja, das heißt in die Herrlichkeit gehen. Denn das sind unseres HErrn GOttes heimliche und verborgene Wege, die man verstehen muß nicht nach dem Fleisch und menschlichen Verstande, sondern nach der Schrift, und wie sie Christus selbst versteht und deutet, da er zu seinen Jüngern sagt Luc. 24, 26.: „Mußte nicht Christus solches leiden, und zu seiner Herrlichkeit eingehen?“

268. Eben auf dieselbe Weise hat Jakob auch zwanzig Jahre müssen im Elend sein und darnach erst ein König und Priester werden. David mußte zehn Jahre lang hin und her getrieben werden, und mancherlei Winkel suchen, da er etwa eine Nacht möchte sicher liegen, bisweilen auch wohl nur eine Stunde: ist das ein König, von GOtt zum Königreich gesetzt und bestätigt? Ja, also hört du ihn singen und von seinem Königreich hoch rühmen im 63. Psalm V. 12.: „Aber der König freuet sich in GOtt. Wer bei ihm schwöret, wird gerühmet werden; denn die Lügenmäuler sollen verstopft werden“, das ist, die werden sich freuen, die da verstehen, daß er von GOtt zum Könige gesetzt und gesalbt ist. Wo waren aber die? Denn das ganze Regiment und alle Gewalt war bei Saul, da David auch nicht einen Fuß breit Eigenes im Lande gehabt hat.

269. Also sind auch wir, die dem Worte GOttes glauben, die Kirche, haben die allergewisseste Verheißung, dazu wir berufen und getauft sind, dadurch wir ernährt und erhalten werden; wir haben das Sacrament des Altars und das Amt der Schlüssel; wir sind aber nicht Christi Christen und sind auch darum nicht getauft, daß wir dieses Erdreich besitzen sollen, wir sind auch auf dieses Leben nicht getauft und neu geboren, sondern zum ewigen Leben. Was geschieht aber mit uns auch? Antwort: Es geht mit uns also zu, daß wenn die Kirche und Gemeinde GOttes zur Herrlichkeit und zu der ewigen Freude, darauf sie wartet im Worte und Hoffnung, erhoben werden soll, alsdann wird sie unterworfen viel unzähliger Verfolgung der Tyrannen und Teufel; sie wird von den falschen Brüdern wohl geplagt und gelästert und dasselbe auf vielerlei Weise, daß zu erbarmen ist. Heißt aber das in das ewige Leben führen? Ja, das heißt recht in den ewigen Jammer stoßen.

270. Man muß aber die Herzen aufrichten und stärken wider diesen Weg des Kreuzes. Denn dieweil wir das Wort und die Verheißung haben, darum wird die Herrlichkeit, so uns verheißen ist, auch gewißlich folgen, und die Kirche lebt unterdeß und wird erhalten durch den Glauben, der es gewißlich dafür hält, daß GOtt nicht lügen kann, und lernt diese wunderbare Weisheit, so vor dem Fleisch und Vernunft verborgen ist, nämlich, daß GOtt seine Heiligen wunderlich führet, Psalm 4, 4., daß er in seinem Heiligthum wundersam ist, wie der 68. Psalm V. 36. sagt, und daß sein Rath wunderlich ist; daher auch unser HErr und König JEsus Christus den Namen hat, daß er „Wunderbar“ genannt wird, Jes. 9, 6.

271. Dieses gehört und dient nun dazu, daß wir daraus sollen gelehret werden, wenn wir gottselig leben wollen, daß wir alsdann unser Leben auf eine andere Weise anrichten, denn die Welt und das Fleisch pfleget zu thun. Denn wir müssen einfach an dem unsichtbaren Gott hangen, und sollen GOtt mit Freuden danken, daß wir Gottes Wort haben, darin er uns seine Verheißung gibt, davon St. Petrus sagt 2. Epist. 1, 19.: „Wir haben ein festes prophetisches Wort, und ihr thut wohl, daß ihr darauf achtet, als auf ein Licht, das da scheinet in einem dunkeln Orte, bis der Tag anbreche, und der Morgenstern aufgehe in euren Herzen.“ Denn das Wort ist das Licht unsers Lebens, sonst haben wir nichts von der Herrlichkeit. Ich weiß, daß ich getauft bin, ich weiß, daß ich den Leib Christi gegessen und sein Blut getrunken habe, daß ich absolviret, berufen und mit dem Wort des Evangeliums gelehret bin, sonst habe ich vom ewigen Leben nichts mehr; ich habe noch keinen verklärten Leib, der klarer sei, als die Sonne und Sterne sind, sondern ich habe ein solch Herz, das dermalen noch sehr beschweret ist mit vielem großen Unfall und Schrecken; ich trage noch einen solchen Leib umher, der vieler Schwachheit und dem Tode unterworfen ist. Derhalben sieht man beide an unserem Leib und Seele nichts weniger denn das ewige Leben; die Verheißung aber wird uns nicht fehlen. Darum sollen wir am Glauben und an der Hoffnung hangen und beständig bleiben, und uns am Worte genügen lassen, darin wir die Verheißung haben. Danach haben wir dieses äußerliche Leben und Gemeinschaft; wir haben unsere Eltern, Obrigkeit, das Predigtamt, und sonst auch äußerliche Güter, so zu diesem Leben nöthig sind, welche alle eine Zubereitung und Vorstufen sind zu dem künftigen Leben.

272. Dieses ist die Lehre, so eigentlich und vornehmlich in die Kirche gehört, welche der Heilige Geist lehrt, die der Welt und dem Fleische unbekannt ist, darin wir gelehret werden, daß wir Herren und Erben sind des ewigen Lebens in keiner andern Weise oder Gestalt, denn wie Jacob ein Erbe gewesen ist des Segens, welcher, da er den Segen erlangt hatte, aus dem Lande und aus dem Hause seines Vaters in das Elend ist gesandt worden. Denn dieses ist der Proceß, welchen die göttliche Majestät hält mit ihren Heiligen, und dieses ist der Glaube der Heiligen, davon wir bisher gesagt haben. Nun folgen solche Lehren, die zum äußerlichen Leben oder guten Sitten dienen und von den Früchten des Glaubens handeln.

\quellseite{103}

\chapter{Das achtundzwanzigste Capitel}

\section[Erster Theil]{Erster Theil. Wie Isaak dem Jakob befiehlt, kein Weib zu nehmen von den Cananitern; wie er ihn segnet und nach Mesopotamien schicket, und wie sich Esau diese Verheißung zu Nutz machet}

\subsection{I}

\begin{quote}
\textbf{V. 1. 2. Da rief Isaak seinem Sohn Jakob, und segnete ihn, und gebot ihm, und sprach zu ihm: Nimm nicht ein Weib von den Töchtern Canaan; sondern mach dich auf, und zeuch in Mesopotamien zu Bethuels, deiner Mutter Vaters, Haus, und nimm dir ein Weib daselbst von den Töchtern Laban, deiner Mutter Bruders.}
\end{quote}

1. Das erste Stück in diesem Capitel ist nicht theologisch; denn es sind darin keine Exempel oder geistliche Lehren vom Glauben und andern geistlichen Gottesdiensten: wiewohl was hier gehandelt wird, beides Glauben und Gottesfurcht haben will und in sich begreift. Es ist aber eine Lehre von guten Sitten, nämlich, vom Ehestande; welche Lehre man in der Kirche behalten und fleißig treiben soll wegen der Nothwendigkeit und Würde des Ehestandes. Denn nach der Lehre des Evangeliums und des Glaubens, welches eigentlich die rechte Lehre der Kirche ist, soll vornehmlich der Ehestand geehrt und gelobt werden; und dasselbe auch um der Ursache willen, weil die Welt und das Fleisch nicht verstehen, was der Ehestand sei, oder wie viel man davon halten solle.

2. Es wird zwar der Ehestand dem gemeinen Gebrauche nach also beschrieben, daß sie sagen: Der Ehestand ist, wo Mann und Weib zusammen gefüget werden und bei einander leben, also daß sie sich nicht dürfen von einander scheiden in diesem Leben. Aber diese Definition oder Beschreibung des Ehestandes ist nicht vollkommen oder genug; denn es ist noch nicht dabei die Endursache, noch auch die Ursache, so den Ehestand macht. Es ist nur die causa materialis, wie man es in den Schulen nennt, darin angeführt; denn daß Mann und Weib zusammen gefüget werden, das ist das materiale. Diese Definition aber ist besser und vollkommener, wenn man also sagt: Der Ehestand ist, wo Ein Mann und Ein Weib ordentlich und göttlich zusammen gefüget werden, welches dazu verordnet ist, daß man Gott in dem Stande solle anrufen, und Kinder erhalten und aufziehen, das Kirchen- und Weltregiment zu verwalten.

3. Derhalben ist in der christlichen Lehre, darin wir nach der Lehre des Evangeliums und des Glaubens lehren, wie man sich in diesem Leben göttlich und ehrlich halten solle, der Ehestand das erste und das vornehmste Stück. Denn es ist der Anfang und Ursprung des ganzen Lebens, und der Teufel setzt diesem Stande nicht weniger zu als der Kirche; welches daran wohl zu sehen ist, weil es so ein seltsam Ding ist um einen guten und heilsamen Ehestand, und um die Liebe, so Mann und Weib unter einander haben sollen, und auch soviel die Kinder und Nachbarn belanget. Denn dieser Stand ist mit Dornen und Disteln umgeben. Welches zwar anfänglich die jungen Leute nicht sehen können, wenn sie durch häufige Liebe und Lust getrieben werden, daß sie sich leichtfertig in den Ehestand begeben; darnach aber, wenn sie die Mühseligkeit und allerhand Schwierigkeiten fühlen und erfahren, alsdann gereuet sie es zu spät und vergeblich.

4. Darum sollen sich die Christen also in die Sache schicken und ihr Leben also einrichten, daß sie bedenken, der Ehestand sei nicht ein leichtfertig oder zufällig Ding, welches von unserm Gutdünken und dessen Ausgang vom Zufall abhänge; sondern daß in dem Stande Mann und Weib ordentlich und göttlich zusammen gefüget werden; welches dadurch klärlich bewiesen wird, weil Gott Mann und Weib geschaffen hat, und der Mann nicht allein, oder auch das Weib nicht allein geboren wird, sondern beide Mann und Weib. Derhalben hat dieses Zusammenfügen bald in der ersten Geburt seinen Ursprung, und ist darum recht ordentlich und göttlich.

5. Darnach hat Gott den Ehestand nicht um Wollust und fleischlicher Lust willen eingesetzt, das ist nicht die Endursache; sondern es hat der Ehestand zweierlei Endursachen, darum er eingesetzt ist. Die erste ist, daß er eine Arznei sein soll, damit der Unzucht gewehret werden möge. Die andere und vornehmste Ursache ist, daß ein Anfang und Ursprung sein soll des menschlichen Geschlechtes, auf daß Kinder geboren und das menschliche Geschlecht ausgebreitet werden möge, oder wie die Juristen sagen, daß die Stadt voll Bürger werden möge. Aus der heiligen Schrift aber soll man diese Endursache auch noch hinzu setzen, daß in dem Ehestande die Kinder in der Zucht und Furcht des Herrn auferzogen werden, auf daß sie zum Kirchen- und Weltregiment geschickt werden mögen.

6. Das Hauptstück sollen nun die Gottseligen fleißig merken und behalten, nämlich, daß beide Mann und Weib von Gott dazu geschaffen sind, daß sie sollen zusammengefüget werden, und im Ehestande bei einander leben und sich nicht scheiden. Das ist so viel gesagt: Es ist Gottes Wille, dadurch er Mann und Weib will ordentlich zusammengefüget haben, daß wir Kinder aufziehen, die der Kirche und Gott dienen mögen. Wenn wir das fest behalten werden, so werden wir auch um so viel leichter ertragen und überwinden alle Mühseligkeit und Beschwerung, so in diesem elenden Leben den Eheleuten zustoßen kann.

7. Denn nachdem der Mensch durch die Erbsünde dem Tode und der Gewalt des Teufels unterworfen ist, so feiert der Teufel damit nicht, daß er die Eheleute an Leib und Seele greulich plage und betrübe. Und eben das ist es, daß diesen Stand so sehr entstellt, und denselbigen so gar mühselig, verdrießlich und verhaßt macht, daß die Natur vor dem Ehestande erschrickt und flieht, nicht anders als vor dem Kreuz. Wie das die Exempel und schändlichen Worte der Heiden bezeugen. Und heutiges Tages wirst du deren auch Viele hören, die mit Fleiß allerlei Unannehmlichkeiten aufsuchen, wodurch sie sich verführen lassen, daß sie sich nicht zum Ehestande entschließen können.

8. Es ist aber unsere Meinung auch nicht, daß man jemand mit Gewalt und wider seinen eigenen Willen zum Ehestande zwingen sollte. Denn wo es dir gegeben ist, daß du keusch leben und dich ohne Sünde des Ehestandes enthalten kannst, ei, so magst du dich des wohl enthalten, wenn du das ohne Sünde thun kannst; wo du aber ohne Sünde dich dessen nicht enthalten kannst, daß du dich mit einem Weibe vereinigest, so bedenke und gebrauche die Arznei, so dir Gott gezeigt hat; und wo du ja das nicht suchst, wie du mögest Kinder zeugen, so suche doch gleichwohl die Arznei wider die Sünde, nämlich, die Hurerei und Ehebruch zu meiden, zu vermeiden Befleckungen und schändliche Lüste. Denn es ist Jammers mehr denn genug, daß wir mit allerlei Unfall, Sünde und Tod beladen sind, daß du nicht darüber noch eine Sünde über die andere thust und dir selbst mit andern Sünden dein böses Gewissen vermehrest.

9. Ja, sagen sie, es ist aber zumal ein elender und verdrießlicher Stand. Antwort: Es ist zweierlei Leben, das eine ist ein Leben der Sünde, das andere ein Leben der Strafe. Nun betrachte du aber selbst, welches das beste sein möge: daß man in der Strafe lebe ohne Sünde, oder aber dagegen in Sünden lebe ohne Strafe. Die Vernunft ist so gar verderbet, daß sie die Strafen der Sünde nicht mit Geduld ertragen kann; darum hat sie Lust zu den Dingen, die da lieblich, fröhlich und lustig sind, aber vor dem Kreuz und Mühseligkeit fliehet sie immer. Wie man von Socrates liebet, da er gefragt worden ist: Welches besser wäre, ein Weib nehmen oder nicht? daß er darauf geantwortet hat: Du thust gleich welches du willst unter den Zweien, so wird es dich doch gereuen. Wirst du Eine nehmen, die mißgestaltet ist, so hast du eine Strafe; wo du aber ein schön Weib nehmen wirst, so werden sie Andere auch wollen lieb haben 2c. Der hat nur auf die Strafe und den Mißbrauch gesehen.

\quellseite{104}

10. Ein Christ aber soll in seinem Herzen also denken: Man muß die Strafe verachten und alle Müßigkeit von den Augen wegthun, und muß es wagen im Namen des Herrn. Denn man muß auf das Leben denken, so ohne Sünde, ohne Unreinigkeit, ohne Befleckung und Makel ist, auf daß du vor GOtt ein gut Gewissen haben mögest; denn man soll in diesem Leben des Himmelreiches, der Seligkeit oder des ewigen Lebens nicht hoffen; die Erbsünde, die Schwachheit des Fleisches und der Teufel lassen solches nicht zu. Wenn man aber ja Jammer und Elend leiden muß, so sollen wir es lieber mit GOtt denn mit dem Teufel leiden. Denn die vor der Beschwerung des Ehestandes fliehen, die leben zwar wohl in einem sanften und lustigen Leben, sie haben ihre Hurerey, die sie bald ziehen lassen und von sich treiben, bald wieder zu sich nehmen, wie es ihnen selbst wohlgefällt: was meinst du aber, daß solche Leute für ein Gewissen haben mögen? Sie müssen wahrlich ein sehr böses Gewissen haben, und das noch das Allerbeschwerlichste ist, sie müssen doch auch oben dießelbigen, ja, noch wohl mehr Müßigkeit und Beschwerungen tragen, damit sie von den Huren auf unwürdigere Weise geplagt werden denn von den Eheweibern. Also haben sie zweierlei Strafe zu tragen, nämlich, die gemeine Strafe, und sie selbst thun die ewige Strafe auch dazu.

11. Derhalben soll man nach der Lehre, die vornehmlich in der Kirche muß getrieben werden, dieses vom Ehestande in der Kirche auch lehren und immer fleißig einschärfen, daß die Leute wissen mögen, wie sie in diesem Leben ohne Sünde leben mögen. Darnach soll man auch die Herzen der Gottseligen aufrichten und ermahnen, daß sie die Müßigkeit und Beschwerung des Ehestandes lernen verachten. Denn ein Ehemann muß sich männlich halten nicht allein bei der Nacht, sondern auch die Strafen und den Willen GOttes zu dulden, daß er die Bosheit des Gesindes, Untüchtigkeit der Nachbarn und was sonst mehr Unglücks vorfallen kann, vertragen und in sich fressen könne.

12. Darum soll man zum Ehestande fromme Männer und fromme Weiber erwählen; wie Abraham und die andern heiligen Väter stark vom Weibe gewesen und auch einen männlichen Muth gehabt haben: die Andern, die sich vor der Last und Beschwerung des Ehestandes fürchten und sich derselben darum enthalten, thun eben, wie man im Sprüchwort sagt: Sie gehen aus dem Rauch in das Feuer. Denn sie suchen ein ander Leben, denn das, so in der Sünde geboren ist: daselbst finden sie aber den Teufel und die Hölle; denn sie können der bösen Lust, der Ehrschande, der Sünde und dem Tode, so im Fleisch und Herzen steckt, nicht entfliehen.

13. Darum fasse Muth, und denke, daß dieses Leben nichts Anderes sei denn das Elend selbst. Desselben Jammer, Mühe und Beschwerung wirst du also überwinden, wenn du wirst auf GOtt, deinen Schöpfer und Vater, sehen, dessen Willen und Ordnung du dich mit Demuth und Geduld unterwerfen sollst, und also bei dir selbst schließen und sagen: Ich will an GOtt glauben, der mich zu einem Mann geschaffen hat; ich will ihm danken, daß ich sein Wort habe, daß ihm also wohlgefallen hat, daß ich soll ein Ehemann, ein Eheweib sein, daß ich Kinder aufziehen und das Hausgesinde regieren soll, daß ich durch das Evangelium die Verheißung des ewigen Lebens und Trost in diesem gegenwärtigen Leben habe. Wie Paulus sagt 1 Tim. 4, 8.: „Die Gottseligkeit ist zu allen Dingen nütz, und hat die Verheißung dieses und des zukünftigen Lebens“; und Röm. 15, 4.: „Was zuvor geschrieben ist, das ist uns zur Lehre geschrieben, auf daß wir durch Geduld und Trost der Schrift Hoffnung haben.“

14. Auf die Weise sollen wir lernen, mehr vom Ehestande halten und besser davon urtheilen, denn das Fleisch und die Welt pfleget zu thun; und wir sollen die jungen Leute ermahnen, daß sie mit Geduld leiden und ertragen wollen, wenn irgend ein Unglück oder Beschwerung vorfällt. Sie sollen sich ans Gebet gewöhnen und also sagen: Lieber HErr GOtt, ich bin deine Creatur, bin von dir ein Mann geschaffen, und zu diesem Stande verordnet, darin ich jetzt stecke, und bin in viel Unlust und Beschwerung gerathen; aber gib du mir Gnade, daß ich wahrhaftig erkennen möge, daß ich deine Creatur sei, daß du mein Vater und Schöpfer seiest und daß ich von dir Hülfe und Schutz erwarten möge. Ein solches Gebet thut allen Eheleuten vonnöthen. Denn wir erfahren alle, daß uns in dem Stande das Kreuz und Mühseligkeit auferlegt ist.

15. Und wollte GOtt, daß dermaßen der Cölibat oder ehelose Stand in den Stiften der Pfaffen und Bischöfe möchte aufgehoben werden, und daß man ihnen den Ehestand erlauben oder gebieten möchte anzunehmen; wie unsere Kirchendiener gottselig und ehrlich im Ehestande leben: so würden sie nach den Römthürmern und Canonisten oder Präbenden nicht also laufen und darnach so großes Verlangen haben. Denn die Stifte würden alsdann nicht so gewaltig und reich sein. Aber es ist vergeblich, daß wir solches wünschen und begehren: sie sind dumm wie die Ochsen und schmutzig wie Bodensatz.

16. Dieses hat man im Anfang dieses Capitels sagen müssen vom Ehestande, dieweil Mose erzählt, wie Isaak seinem Sohn Esau gerufen und mit ihm geredet habe, wie er sein Haus sollte anfangen zu regieren und ein Weib nehmen. Und der Anfang dieses Ehestandes ist nach GOttes Kraft, Ordnung und Beruf gegangen. Was er aber für Elend und Beschwerung im Ehestande haben werde, werden wir in nachfolgenden Capiteln sehen. Denn er hat im ganzen Ehestande nichts Eigenes, keine Lust oder Freude, ausgenommen die Kinder; sonst ist er ein sehr elender Mensch, lebt im Elende, und dient seinem Schwiegervater Laban, der ganz unbillig mit ihm umgehet.

17. Dieses ist aber die Ordnung und der Beruf zum Ehestande, daß GOtt dem Jakob durch den Mund seines Vaters gebietet, daß er ein Weib nehmen soll. Es setzt aber der Vater noch ein ander Gebot dazu und sagt: Er soll sich kein Weib nehmen von den Töchtern Canaan, sondern soll vielmehr von den Töchtern Labans eine zum Weibe nehmen. Derhalben gehört dieser Text zu der Lehre, daß wo man eine Heirath eingehen will, man wissen soll, daß dazu die Gewalt und Bewilligung der Eltern vonnöthen ist. Denn wir haben die heimlichen Verlöbnisse und Heirathen verworfen und verdammt, weil wir nun, nachdem das helle Licht des Evangeliums wiederum hervor gekommen, wissen, daß der Ehestand ein heilig Ding und jedermann erlaubt, dazu auch eine göttliche Ordnung ist.

18. Es ist keine Schande oder Unehre, ehelich werden, wie wir es vor dieser Zeit für unehrlich gehalten haben, da uns die Mönche in den Irrthum geführt hatten; sondern es ist ein ehrlich und heilig Ding. Wir wissen zwar wohl, daß dieser Stand greulich entstellt ist durch die böse Lust, um welcher Ursache willen es viele Leute für unehrlich gehalten haben, wo man eine Jungfrau zur Ehe begehrt und Heirath eingegangen. Das, sage ich, haben sie gemeint, sei ein schändlich unrein Ding. Denn das Werk der Kinderzeugung ist nicht unterschieden worden von andern Sünden, als Hurerei und Ehebruch. Jetzt aber haben wir nun gelernt, und sind dessen von GOttes Gnade gewiß, daß die Ehe ehrlich ist. Wie in der Epistel an die Hebräer am 13. Cap. V. 4. steht: „Die Ehe soll ehrlich gehalten werden bei allen, und das Ehebette unbefleckt“ 2c., und daß es GOttes Wille und Einsetzung sei, daß ein Jeder sein eigen Gemahl habe, die ihm ordentlich zugefügt worden sei.

19. Darum haben wir dessen keine Ursache, daß wir uns vor diesem Stande scheuen wollten, oder uns die Scham und Zucht schrecken lassen, daß wir uns heimlich und listiger Weise darein begeben und denselben im Winkel anfangen wider die Ehrbarkeit und den Willen unserer Eltern. Eine Jungfrau soll sich nicht scheuen vor ihre Eltern zu treten, und sie zu bitten, daß sie ihr einen frommen ehrlichen Jüngling wollen zur Ehe geben; und die Eltern sind, GOtt Lob! nun auch recht unterrichtet worden, und sind willig und geneigt, ihren Kindern zu rathen und zu helfen. So sind jetzt auch fromme gottesfürchtige Pfarrherren und Regenten, die vermöge ihrer Gewalt, so sie haben, die Eltern überreden können, daß sie sich in solchem Fall nicht hart und schwierig erzeigen, sondern gutwillig sind, ihren Kindern zu Ehren zu helfen.

\quellseite{105}

20. Derhalben soll man sich vor den heimlichen Verlöbnissen hüten und soll diesen sehr heiligen Bund des Ehestandes zu Gottes Ehre, und mit gebührlicher Ehrbietung gegen die Eltern, anfangen. Derselben Lehre wird uns an diesem Ort ein schön Exempel vorgehalten. Denn hier geht die Gewalt und der Befehl des Vaters Isaak vorher, der seinem Sohn gebietet, daß er kein Weib von den Töchtern Canaans nehmen soll, und Jakob ist ganz willig, seinem Vater darin zu gehorchen. Esau aber hat das Gegenteil gethan, hat zwei Weiber genommen, die alle beide Töchter der Hethiter waren, wider seiner Eltern Willen, damit er sie auch hart erzürnt hat, daß beide Isaak und Rebekka mit Schmerzen geplagt haben, wie sie von denselbigen Weibern ihres Sohnes geärgert und jämmerlich geplagt worden sind. Das ist eine schwere Sünde und Laster gewesen.

21. Nun ist es aber an dem, daß wir sonst im Ehestande Mühseligkeit und Beschwerung mehr denn genug zu tragen haben, warum wollen wir uns denn über dies durch die Sünde und Verachtung unserer Eltern noch anderes und schwereres Unglück und dazu den Zorn Gottes auf den Hals laden? Lasset uns viel lieber im Namen des Herrn in den Ehestand treten, wie Jakob gethan hat. Es wird doch an Beschwerungen und allerlei Jammer nicht mangeln. Dieselben Beschwerungen aber wirst du um so viel leichter ertragen können, wenn du ein solches Gewissen in den Ehestand bringen wirst, daß du sagen kannst: Es hat meinen Eltern, Freunden, und Vormündern oder Pflegern, dazu auch Gott also wohlgefallen. Denn du bist in den Stand nicht getreten, daß du dich auf deine Vernunft oder Weisheit verlassen hättest, und die Gewalt und Bewilligung derjenigen hintangesetzt, von denen Gott haben will, daß du sie in Ehren halten sollst. Darum bist du nun des guten gnädigen Willens Gottes und seines Schutzes gewiß. Wo du aber dasselbe nicht thun wirst, so wird zu dem andern gemeinen Unglück auch das böse Gewissen kommen, welches alle Mühseligkeit mehren und schwerer machen wird. Was ist nun das Andere, denn ein Unglück mit dem andern häufen und größer machen, und daß du damit die greuliche Ruthe des Teufels wider dich immer mehrest?

22. Wir haben aber von diesem Stück droben auch und an andern Orten mehr etwas weitläufiger gehandelt, und man soll es den Leuten fleißig einschärfen, nämlich, daß die heimlichen Verlöbnisse und Winkelgehen verboten sind nicht allein in kaiserlichen Rechten, sondern auch in den Exempeln und Sprüchen der heiligen Schrift, um der Rechtserben willen, die des Papstes Decrete der heiligen Schrift und den kaiserlichen Rechten vorzuziehen, und sie säumen sich zu bekennen und zu erkennen, daß sie weiland ungereimt Ding gelehrt haben. Nun ist aber die Lehre der heiligen Schrift viel eine höhere Lehre; und wir müssen doch erkennen, daß die Scholastiker unzählig viele grobe greifliche Irrtümer in dieser Lehre der heiligen Schrift eingeführt haben. Aber die verkehrten Juristen unterstehen sich noch wissentlich und vorsätzlich, den offenbaren Irrthum zu erhalten und zu vertheidigen; welches über die Maße ein unbillig Ding ist. Denn die kaiserlichen Rechte sind offenbar dawider und zeugen klarlich, daß die Söhne, wenn sie wollen zum Ehestande greifen, dazu die Bewilligung ihrer Eltern haben müssen. So sind auch die Exempel der heiligen Schrift stracks dawider. Was geht uns denn der närrische und schändliche Pabst an mit seinen Decreten?

23. Isaak verbietet seinem Sohne, daß er keine Heidin zum Weibe nehmen soll, und heißt ihn zu seiner Mutter Bruder, Laban, ziehen und daselbst ein Weib nehmen. Er zwingt wohl Jakob nicht, daß er die nehmen sollte, die dem Vater wohlgefallen möchte; er will aber doch haben, daß er eine ehrliche Jungfrau erwählen solle, die er selbst lieben und zum Weibe nehmen möge. Denn der Befehl des Vaters soll vorher gehen, der Sohn aber soll denselben gehorchen und folgen. Und der Vater soll seinen Sohn nicht zwingen, daß er Eine nehmen solle, der er feind ist; sondern die Braut nehmen solle, der er selbst lieben und zum Weibe nicht hindern oder verbieten, sondern fördern und helfen. Dieses ist die rechte Ordnung und Weise, Heirath oder Ehen zu stiften. Dieselbige Ordnung soll man nicht verkehren oder umwenden; denn „es ist genug, daß ein jeglicher Tag seine eigene Plage habe“, Matth. 6, 34., das ist, es ist genug an dem Elend und Jammer des Ehestandes, daß du dann genug haben wirst, wenn du in diesen Stand kommen wirst.

24. Derhalben hüte dich, so viel du immer kannst und magst, daß du deine Eltern und Gott nicht erzürnest, und fange viel lieber diesen Stand an mit guter Gunst und Willen Gottes und der Eltern, auf daß du dein ganzes Leben, so lange dasselbe währen wird, wo nicht gar ohne Strafe und Mühseligkeit, doch gleichwohl ohne Schuld zubringen und beschließen mögest; wo nicht mit eitler Freude, doch gleichwohl mit Unschuld deines Herzens, und mit einem guten Gewissen und Gottes gnädigem Wohlgefallen. Es ist genug an der Plage und Anfechtung des Fleisches, und Gott wolle uns ja vor der Plage und Anfechtung des Geistes gnädiglich behüten und bewahren.

25. Andererseits aber, die gar vor dem Ehestand fliehen entweder unter dem Schein der Heiligkeit, oder aber daß sie Macht haben, ein so viel freier Leben zu führen: die soll man vermahnen zu diesem heiligen Stande, daß sie sich mit häuslichen Lüsten nicht beflecken, wie die Dompfaffen und Cardinäle in ihren Stiften pflegen zu thun. Hast du die Gabe nicht, keusch zu leben, so nimm in Gottes Namen ein Weib; kannst du nicht so bald reich werden, so lasse dir genügen an der täglichen Nahrung; kannst du kein König und Herr sein, so sei ein Knecht und gemeiner Unterthan; glaube an Gott und warte des ewigen Lebens in diesem armen und elenden Leben, welches nicht ewig, sondern sehr kurz ist. Oder wenn du auch an dein hundertstes Jahr kommen würdest, was ist doch so eine kurze Zeit, wenn man sie mit der Ewigkeit vergleichen wollte? Wir Christen wissen ja, daß nach diesem Leben ein ander und besser Leben folgen wird, darauf wir warten, darauf wir allezeit sehen sollen. Und dieweil inzwischen dieser Leib nicht keusch und züchtig leben kann, so müssen wir zwar die Dornen und Disteln im Ehestande nach Gelegenheit desselben Standes dulden und also damit vorlieb nehmen, und sollen uns dessen freuen, daß Gott mit Gnaden auf uns Achtung hat, uns schützt und beschirmt. Denn es gefällt Gott wohl, daß du arbeitest, daß du es dir im Ehestande unter den Dornen lässest sauer werden. Laß dir genügen an seiner Gnade und laß dich die Beschwerung nicht abschrecken. Zudem wisse dich zu erinnern und suche die große Mühe, Jammer und Elend auf, so die heiligen Patriarchen gehabt; so wirst du sehen, wie viel großer Trübsal auch Jakob allein im Ehestande erlitten habe.

\subsection{II}

\begin{quote}
\textbf{V. 3—5. Aber der allmächtige GOtt segne dich, und mache dich fruchtbar und mehre dich, daß du werdest ein Haufen Völker; und gebe dir den Segen Abrahams, dir und deinem Samen mit dir, daß du besitzest das Land, da du Fremdling innen bist, das GOtt Abraham gegeben hat. Also fertigte Isaak den Jakob, daß er in Mesopotamien zog zu Laban, Bethuels Sohn, in Syrien, dem Bruder Rebekka, seiner und Esaus Mutter.}
\end{quote}

26. Droben (Cap. 27, 28. 29.) hat Isaak seinem Sohne etliche Segen gegeben; dieselbigen wiederholt er an diesem Orte und wünscht ihm dazu Glück und Heil. Es ist aber ein solcher Segen (wie wir auch droben vermahnt haben), damit das Gut an sich selber gegenwärtig überreicht und gegeben wird; wie mir die Taufe gegenwärtig gegeben wird. Mir wird Vergebung der Sünden gegeben; denn ich erwarte nicht erst auf die Vergebung der Sünden, sondern ich habe dieselbe im Glauben gegenwärtig. Ich glaube nicht, daß Christus noch für mich leiden werde, sondern ich bin dessen durch den Glauben gewiß, daß er jetzt schon für meine Sünde gelitten hat und um meiner Gerechtigkeit willen auferwecket ist. Derhalben ist es keine bloße Bitte oder Wunsch; sondern damit ich dir gegenwärtig durch die Kraft der Schlüssel Vergebung der Sünde, Gottes Gnade und Gunst gebe, daß du es gewiß dafür halten kannst, daß du einen solchen Gott habest, der an dir Wohlgefallen hat. Das gebe ich dir als ein gewisses Gut oder Erbe. Also hat die Kirche gegenwärtig die göttliche Majestät, so ihr günstig und gewogen ist, nämlich, Gott Vater, Sohn und Heiligen Geist; sie hat Christum und die heiligen Engel, die sich mit ihr freuen; sie hat alle Creaturen, die da Gefallen haben und warten auf ihre Erlösung, wie St. Paulus sagt Röm. 8, 19.

27. Darum ist es viel ein ander Ding um den Segen, als um einen solchen Wunsch, damit man einem Glück wünscht mit Worten; welcher Wunsch doch gleichwohl an diesem Orte auch wird dazu gesetzt: „GOtt“, sagt er, „segne dich“, das ist, er wolle dich mehren und fruchtbar machen etc. Droben ist der Segen gesprochen; hier wird nun das Gebet auch dazu gethan.

\quellseite{106}

28. Denn es ist nun oft gesagt worden, daß zweierlei Priesterämter sind und von Anfang der Welt allezeit auch gewesen sind. Das eine ist, lehren, welches auf der Kanzel oder Predigtstuhl geschieht, oder auch in der Absolution, darin ich dir Gottes Gnade gebe, die mit dem Blute Christi bestätiget ist. Das andere priesterliche Amt ist, für sich und Andere beten, welches nach der Predigt pflegt zu geschehen; denn auf eine jede gute Predigt soll auch ein gut Vater Unser folgen. Denn man muß bitten um Mehrung des Glaubens, daß wir im Segen immer wachsen und zunehmen mögen, daß uns der Teufel denselben nicht nehme. Derhalben muß man beiderlei thun, man muß lehren und beten; denn dieses sind die zwei priesterlichen Ämter, nämlich, Gott reden hören, und mit Gott reden, der uns hört, daß man herab und hinauf steige. Durch den Segen, durch die Predigt und Austheilung der heiligen Sacramente steigt Gott hernieder zu uns und redet mit mir; da höre ich ihm zu. Und wiederum steige ich hinauf und rede Gott in die Ohren, der mein Gebet hört.

29. Man soll aber an diesem Orte auch den sonderlichen Namen Gottes merken, auf Hebräisch schaddai, von dem Worte schad, das ist, Brust. Mit diesem Namen will Gott genannt werden, und reimt sich sein mit dem griechischen Namen polymastos, auf Lateinisch mammosus, der viel Brüste hat. Er will daher gerühmt sein, daß er alle Creaturen ernährt und erhält: er ist nicht allein ein Schöpfer, sondern nährt und erhält auch alles, was er geschaffen hat. Droben hat Mose auch denselben Namen gebraucht, da Gott zu Abraham sagte, 1 Mos. 17, 1.: „Ich bin der allmächtige Gott“ etc., ich mache, daß Alles fruchtbar und gemehret wird, daß Alles lebt und ernährt wird.

30. Derhalben rufst nun Isaak den Gott an, der nicht allein ein Schöpfer ist, sondern der auch alles nährt und erhält, daß er seinem Sohn das verheißene Land geben wolle. Das Land ist dein, will er sagen, aber du bist ein Gast und Fremdling darin; ein Anderer hat dasselbe inne, wiewohl du es für dein Land hältst und dir auch nach göttlichem Rechte gebührt. Darum besitzt Jakob das Land, so ihm gebührt, nicht; welches wahrlich schwer und kümmerlich genug ist. Uns aber würde solches gar ein unleidlich Ding sein: wenn ich ein eigen Haus hätte und ein Anderer sollte darin wohnen, ich aber müßte als ein Gast Brod, Wein und alle andere Nothdurft um mein Geld kaufen, daß ich den Fremdling nährete und kleidete.

31. Derhalben haben die Patriarchen einen großen Glauben gehabt; welches diese Worte Isaaks, die voll Glaubens sind, anzeigen und bezeugen; und sie haben ohne Zweifel verstanden, daß in den leiblichen Verheißungen des Landes Canaan die ewigen und geistlichen Verheißungen mitbegriffen wären. Darum haben sie auf ein ander Vaterland gehofft, sintemal die zeitlichen Güter verheißen und aber doch Andern gegeben werden. Und bis daher soll nun vom Glauben und Verheißung des Segens, desgleichen auch vom Ehestande genug gesagt sein.

\begin{quote}
\textbf{V. 6—9. Als nun Esau sahe, daß Isaak Jakob gesegnet hatte, und abgefertiget in Mesopotamien, daß er daselbst ein Weib nähme: und daß, indem er ihn gesegnet, ihm gebot, und sprach: Du sollst nicht ein Weib nehmen von den Töchtern Canaan; und daß Jakob seinem Vater und seiner Mutter gehorchte, und in Mesopotamien zog; sahe auch, daß Isaak, sein Vater, nicht gerne sahe die Töchter Canaan: ging er hin zu Ismael, und nahm über die Weiber, die er zuvor hatte, Mahalath, die Tochter Ismaels, des Sohns Abrahams, die Schwester Nebajoth, zum Weibe.}
\end{quote}

32. Es läßt sich ansehen, daß Esau, der zuvor grimmig und böse war, nun etwas gelinder geworden sei, da er sieht, daß sein Bruder Jakob aus des Vaters Hause gewichen und weggezogen ist. Denn er hat also gedacht: Ich habe nun genug, sofern ich nur im Hause und Erbe sitzen bleibe; dieser mein gesegneter Bruder ist weggegangen und hat den Segen hinter sich gelassen. Und ist der Meinung, daß es nur ein zufälliger Segen gewesen sei, damit ihn der Vater gesegnet hat, daß er ein Theilchen haben möchte, davon er sich im Elende und Pilgerschaft erhalten könne. Und es ist ihm sehr lieb und angenehm gewesen, daß Jakob seinen Eltern nicht widerstrebt, sondern ihnen so leichtlich gehorcht hat, und sich an einem Wenigen genügen lassen und in das Elend gewichen ist. Er hätte bei dem Vater, hat Esau gedacht, hart mögen anhalten, daß er ihm einen größeren Segen mittheile, wie ich mit Drängen und Anhalten erlangt habe, daß ich drüber wohl jetzt im Hause bleibe. Aber es sind gute, schlichte und einfältige Leute, beide Vater und Mutter und auch mein Bruder Jakob. Die Eltern schicken den Sohn in das Elend: der Sohn gehorcht ihnen mit gutem Willen; denn er zittert und fürchtet sich vor meinem Drohen und vor meinem Zorn; darum thut er klüglich, daß er davon flieht. Denn ich bin der Herr und Fürst, Jakob aber ist Knecht, wie er auch bisher allezeit gewesen ist.

33. Dieses ist wahrlich ein fein Exempel der Gottlosen, die sich selbst also pflegen zu schmeicheln, und sich zu trösten auch mit der heiligen Schrift, die doch auf das höchste wider sie ist. Wie wiederum die gottseeligen und gottesfürchtigen Menschen die Sprüche und Worte in der Schrift, so für sie und ihnen selbst zum Leben und Troste geredet sind, umkehren und dieselben dermaßen verstehen, als ob sie ihnen zum Verderben geredet wären. Ein gottloser Mensch fürchtet sich vor Gottes Zorn, da er auf die Gnade und Barmherzigkeit Gottes hoffen sollte; dagegen aber die Gottlosen ziehen die Sprüche von Gottes Gnade und Barmherzigkeit auf sich. Wie die Papisten und Türken sich auch träumen lassen, daß sie allein Gottes liebe und angenehme Kinder seien: sie lassen sich träumen, daß sie Gott dem Vater im Schooße sitzen. Die Exempel aber der Strafe und die Drohungen wollen sie ferne von sich gewendet haben, als gingen sie dieselben gar nichts an. Davor aber fürchten wir uns, und meinen, daß sie auf uns gehen und zu unserem Verderben gehören, so uns doch Gott damit nicht schrecken, sondern trösten, stärken und erfreuen will.

34. Also hat Esau gesehen, und hat mit sonderlicher Lust und Gefallen gesehen, daß die Eltern aus Furcht vor ihrem zornigen Sohne den Jakob abfertigen und wegschicken; und ihm hat nichts Lieberes widerfahren können. Derhalben will er nun dem Vater wiederum auch Gutes thun und sich bei ihm verdient machen, auf daß er nicht angesehen werde, als sei er seinen Eltern undankbar oder ungehorsam. Denn da er hört, daß der Vater an den heidnischen Weibern, den Töchtern der Hethiter, sein Gefallen hat, denkt er: Siehe, ich will nun Eine nehmen, die meinem Vater gefallen soll; und geht hin zu seinem Oheim Ismael und nimmt desselben Tochter zum Weibe.

35. Ob aber der Vater an solcher Heuchelei und List einen Gefallen gehabt oder nicht, zeigt die Schrift nicht an. Ich glaube es aber nicht; denn Isaak hat es dafür nicht gehalten, daß der Segen auf Esau ruhen müßte, sondern er hat wohl gewußt, daß er Jakob gebühre, und hat dazu auch wohl gesehen, daß Esau Alles, was er gethan, nur äußerlich zum Schein und betrüglich oder mit List gethan hat. Dieses ist aber ein trefflich und sehr schön Exempel der gottlosen Menschen und Heuchler, die in diesem Leben groß Glück haben und im Flor stehen, haben alles Dinges vollauf: sie werden herrlich und mit großem Gepränge in ihre Königreiche und Güter dieser Welt gesetzt und bestätigt, wie unsere papistischen Bischöfe; aber am Ende zeigt es sich, nach welcher Melodie.

36. Bisher hat nun Mose die Historie Isaaks und Rebekkas kürzlich verfaßt und beschrieben, darin er sehr schwere Anfechtungen und Kämpfe erzählt hat, die sie beide daheim im Hause und auch draußen mit großem starkem Muthe im Worte und Glauben geduldet und ertragen haben. Denn da Isaak die Verheißung empfangen hat, daß sein Same sollte gemehret werden, hat er noch ganze zwanzig Jahre mit unfruchtbarem Weibe leben und zubringen müssen: und da Rebekka von Gott erhört worden, wird sie schwanger und gebieret Zwillinge nicht ohne Anfechtung und Gefahr ihres Leibes und Lebens. Darnach, da sie Kinder erlangt haben, müssen sie um der Theurung willen einen neuen Sitz oder Wohnung suchen. Da muß der gute Isaak mit seinen Kindern und Gesinde unter den Philistern ein Fremdling sein, da Rebekka ihrer Zucht oder Keuschheit halben in Gefahr kommt; und muß hernach auch den Zank und Feindschaft der Hirten, ja, auch des Königs selbst dulden und leiden. Aber in diesem allen richten sich die heiligen Leute auf und werden erhalten durch das Wort und die Verheißung. Zudem werden auch die Herzen des Königes und der Wideracher wiederum zufrieden gestellt, daß sie allenthalben mit ihm versöhnt werden.

\quellseite{107}

37. Endlich aber, da sie nun der äußerlichen Anfechtung sind los geworden, kommt darüber auch noch das Hauskreuz, von wegen der Bosheit und Ungehorsams ihres Sohnes Esau und seiner zwei Weiber, von denen wir gehört, daß sie die alten Eltern schwer geärgert und wohl geplagt haben. Und hat vornehmlich der Rebekka solches wehe gethan, da die Sorge dazu gekommen, daß die Erstgeburt und der Segen auf Esau möchte gewendet werden wider den Verstand der göttlichen Antwort, damit denn der Vater selbst ist umgegangen und dessen sich Esau auch mit großer Hoffart angemast hat. Aber daselbst leuchtet hervor eine wunderbare und göttliche Kraft des Wortes in der Verheißung; welches Wort unbeweglich und fest hat bestehen müssen, da auch Isaak selbst sich aus Irrthum dawider gesetzt. Derhalben wird er durch Rath und List der Rebekka, die fest über dem rechten Verstande der Verheißung gehalten, fein hintergangen und betrogen, daß er ohne alle Gedanken und unwissend den Segen Jakob gibt, dem er gebührt hat; Esau aber hat den Segen verloren und ist nach seinem Verdienst gestraft worden von wegen seines gottlosen Wesens, und daß er die Erstgeburt zuvor verachtet hatte.

38. Daß man nun diese Exempel der wunderbaren Kämpfe und des Jammers und Elendes, daß das ganze Leben dieses Patriarchen voll ist, in der Kirche haben mag, ist nütze und gut, auf daß dadurch die Herzen der gottseligen Menschen erweckt werden, daß sie auch solche Geduld und dergleichen Hoffnung und Glauben haben mögen und Gott in der Anfechtung und Noth lernen anrufen.

39. Nun folgt aber weiter die Historie des Patriarchen Jakob. Derselbigen wollen wir eben recht den fünften Theil dieses ersten Buchs Mose zueignen. Denn nachdem er den Segen von seinem Vater erlangt und seinen Rath und Vermahnung, was für ein Weib er nehmen solle, gehört hat, da wird nun hinfort Isaaks und Rebekkas nicht mehr gedacht werden, bis in das fünfunddreißigste Capitel, da Mose erzählen wird, wie Isaak alt und des Lebens satt gewesen, und also gestorben und zu seinem Volke sei versammelt worden. Wiewohl aber nun Esau inzwischen in der Kirche und im Hause Isaaks das Regiment hat, und seinen Bruder, der flüchtig hat werden müssen, feindseliger Weise verfolgt: so ist dennoch Jakob um der Verheißung und des Segens willen allein der rechte Patriarch oder Erzvater, und daß ich es also nennen möge, er ist das Licht der Welt, da die andern alle gestorben und sein Vater Alters halben zum Kirchenregiment nicht mehr tauglich gewesen ist. Derhalben wollen wir nun sehen, wie Gott diesen Patriarchen auch wunderbarlich geführt und versucht hat durch große Gefahr und Beschwerung, darin er erfahren wird, daß Gott gegenwärtig bei ihm ist und ihn schützt, nicht anders, denn Abraham und Isaak solches auch erfahren, und alles Unglück, Jammer und Trübsal mit starkem Glauben auf die Verheißung Gottes überwunden und den Sieg behalten haben.

\section[Zweiter Theil]{Zweiter Theil. Wie Jakob seine Pilgrimischafft und Wanderschafft nach Mesopotamien antritt und auf freiem Felde übernachten muß}

\begin{quote}
\textbf{V. 10. Aber Jakob zog aus von Bersaba, und reiste gen Haran.}
\end{quote}

40. In diesem Text ist das erste Stück sehr trefflich, und ist werth, daß man davon fleißig handle. Denn es begreift in sich eine vielfältige und reiche Lehre, die schier reicher ist, denn daß wir sie mit unserer Auslegung erreichen könnten. Denn Jakob zieht hin in das Elend, und muß vor dem grimmigen Zorn seines Bruders fliehen, und dem Brudermörder den Segen, so doch ihm gebührt, lassen. Wiewohl derselbe damit zufrieden gestellt ist, weil der Besitzer des Segens flieht und ihn so viel Jahre lang im vollen Besitz des Erbes bleiben läßt.

41. Dieses ist nun wohl in der göttlichen Regierung vor unsern Augen närrisch und ungereimt Ding, daß der Erbe, welcher der rechte Erbsitz werden soll, von den Bauleuten verworfen wird. Er war dazu verordnet, daß er sollte Herr im Hause sein, und ist ihm von Gottes wegen befohlen, das Regiment im Hause und in der Kirche zu haben. Das heißt aber nicht, ein Haus besitzen oder das Gesinde und die Kirche regieren, wo man von diesem allen flieht und dasselbe alles verlassen muß. Darum schließt die Vernunft, es seien bloße und eitle Worte in der göttlichen Verheißung, dieweil das Gut, so verheißen wird, nicht bald mit der That folgt, sondern gar aufs Gegentheil hingeführt wird.

42. Dieses gehört aber nun zu unserer Lehre. Denn die heilige Schrift ist nicht, wie die Philosophie, welche auf die Dinge sieht, so gegenwärtig sind, und vom Glück aufgeblasen und stolz wird; wenn aber das Glück aufhört, wird sie verzagt. Und dieses ist das Argument der Vernunft und aller Epicurer, nämlich, daß Gott angesehen wird als ein Lügner, dieweil er so gnädige Verheißungen und gute Worte promit, aber doch so böse Dinge gibt. Er gibt Jakob den Segen und läßt ihn denselben wiederum genommen werden: dagegen aber bleibt der verfluchte Esau im Hause mit seinen Kindern, Weibern und der ganzen Freundschaft und regieret Alles; wie er zuvor auch Hausvater und Priester in der Kirche ist. Jakob geht in das Elend und verläßt seinen Segen. Ja, das ist die Regierung göttlicher Dinge.

43. Der Christ muß verworfen werden, auf daß wir lernen einen Unterschied machen zwischen göttlicher und menschlicher Regierung, und daß der Segen Jakobs nicht verloren, sondern daß es nur eine Versuchung sei, ob er auch mit dem Glauben daran fest halten und darauf harren wolle. Denn also wird Abraham auch geboten, daß er seinen Sohn opfern soll, 1 Mos. 22, 2., als ob er ihn verlieren werde: aber bald bekommt er ihn wieder mit großer Herrlichkeit und Frucht, 1 Mos. 22, 12. Also ist es mit David auch gegangen: da er zum Könige gesalbt war, 1 Sam. 16, 13., muß er noch zehn Jahre im Elende umher ziehen; er ist ein Knecht anderer Knechte, ist ein König allein dem Namen nach und von leeren Worten, die er von Samuel gehört hat; er verharrt aber dennoch und bleibt beständig im Glauben, und harrt, bis es mit ihm einen solchen reichen Ausgang nimmt, wie die Verheißung gelautet hat.

44. Dieses ist also der beständige Lauf in der Kirche zu allen Zeiten, nämlich, daß Gott Verheißung gibt, und darnach mit denen, so der Verheißung glauben, also umgehet, daß sie auf das Unsichtbare warten, und was sie nicht sehen, glauben, dazu auch hoffen müssen, das man nicht sieht. Wer das nicht thut, der ist kein Christ. Denn Christus ist selbst also in seine Herrlichkeit gegangen, daß er erst hinunter in die Hölle gefahren; da er regieren will, wird er erst gekreuzigt; da er zur Herrlichkeit soll erhoben werden, wird er erst beleidet: denn man muß erst leiden, und darnach zur Herrlichkeit kommen.

45. Das thut aber Gott darum, daß er unsere Herzen versuchen oder prüfen möge, ob wir auch die verheißenen Güter eine Zeit lang entbehren wollen. Ewig werden wir dieselben nicht entbehren, das ist gewiß; und wo uns Gott nicht versuchen, und dasjenige, so er verheißen hat, verziehen würde, so könnten wir ihn nicht von ganzem Herzen lieben. Denn wo er uns alsbald alles geben würde, was er uns verheißt, so würden wir nicht glauben, sondern gar ersäuft werden in den gegenwärtigen Gütern und Gottes vergessen. Derhalben läßt er die Kirche angefochten und betrübt werden, läßt sie Kummer und Armuth leiden, daß sie lerne, daß sie nicht allein vom Brod, sondern auch vom Wort leben müsse; 5 Mos. 8, 3., und daß in den göttlichen der Glaube und die Hoffnung genähret werde und daß sie immer auf die göttliche Hülfe warten, denn das Wort ist unser Leben und Seligkeit. Und so in diesem Leben von den Kindern Adams gesagt und geglaubt kann werden, daß man gemeinschaftlich pflegt zu sagen: Quod differtur, non aufertur: Aufgeschoben ist nicht aufgehoben: warum sollen denn wir Christen eben dasselbe nicht auch also halten und glauben?

46. Es ist zwar sehr schwer und verdrießlich, daß man sich an diese Hülfe und Trost halten soll; wie Salomo sagt in seinen Sprüchen am 13. Cap. V. 12.: „Die Hoffnung, die sich verzeucht, ängstet das Herz“; das Gut aber, so verheißen wird, tröstet und erhält dennoch das Herz. Darum muß man sich fleißig rüsten und stärken wider die Ungeduld des Fleisches, das in solchen Warten wider den Glauben und Geist streitet; wie dasselbe Murren des Fleisches beschrieben wird in den Exempeln der zwei Weiber, des Tobias und Hiobs. Denn also sagt das Weib des Tobias im 2. Cap. V. 22.: „Da siehet man, daß dein Vertrauen nichts ist, und deine Almosen verloren sind.“ Und das Weib Hiobs sprach auch zu ihm Cap. 2, 9.: „Hältst du noch fest an deiner Frömmigkeit? Ja, segne Gott und stirb“ etc. Und seine Freunde haben ihm dasselbe auch aufgerückt und vorgeworfen: Wo ist nun dein Glaube? Hiob 4, 6. Wenn GOtt an dir ein Wohlgefallen hätte, so würde er dich nicht also plagen lassen etc. Das sind die feurigen Pfeile des Teufels, Eph. 6, 16., damit er versucht uns zu bewegen, daß wir verzweifeln und von GOtt abfallen sollen.

\quellseite{108}

47. Man muß aber dagegen halten den Trost, den wir nun oftmals wiederholt haben; nämlich, daß wir sagen können: Es ist ein Betrug und kein Verlust; wir werden also getödtet, aber doch nicht verderbet, sondern erbaut und gebessert. Der Segen wird verzogen, aber harre und bleibe beständig im Glauben, und sollst nicht sinken, das im 115. Psalm V. 2. steht: „Warum sollen die Heiden sagen: Wo ist nun ihr GOtt?“, sondern du sollst lieber also denken und sagen, wie im 42. Psalm V. 6. steht: „Was betrübest du dich, meine Seele, und bist so unruhig in mir? Harre auf GOtt; denn ich werde ihm noch danken“ etc. Es wird noch gut werden. Denn da streitet David wider sich selbst, da er sagt: „Meine Seele, du bist unruhig in mir“, dieweil das nicht gegenwärtig ist, das du gerne haben wolltest; harre doch nur und sei zufrieden, denn es ist GOtt und sein Wort gegenwärtig bei uns.

48. Wenn sich nun fleischliche Menschen trösten und aufrichten können mit diesem gemeinen Spruch, daß man sagt: Quod differtur, non aufertur: Aufgeschoben ist nicht aufgehoben; warum sollten denn wir das nicht auch thun? Und solcher Stempel ist das Regiment auch in weltlichen Sachen voll. Denn es kommt oft vor, daß irgend ein Fürst durch Brief und Siegel verheißt, daß er Einem einen Acker, Schloß oder dergleichen geben wolle: da wird solches auch nicht sobald gegeben, sondern man muß auf die Gelegenheit warten, auf die Städte, Person und Zeit. Wie viel mehr sollen denn wir auf die Güter harren, so GOtt verheißen hat, welcher ja nicht eitel ist, noch lügen kann, sondern hält und thut gewißlich, was er verheißen hat! wie Jacob seiner Verheißung fest geglaubt, und hat doch unterdeß Geduld, da es sich damit also verzogen. Er weiß, daß er ein Herr über die Kirche und den ganzen Erdboden geworden ist; denn der Segen ist vom Hause Abrahams auf das Haus Jakobs gekommen: und zieht doch gleichwohl in das Elend, und wohnt bei denen, die abgöttisch sind.

49. Mit solchen Exempeln sind wir wie mit einer Wolke umgeben; denn die ganze Kirche und ein jeder Christ, ja, Christus, unser HErr selbst, trägt und duldet diesen Willen und Regierung GOttes: welche Regierung ist, daß er die Seinen versucht, aber doch nicht verdammt; daß er mit seinen Verheißungen verzieht, aber dieselbigen gleichwohl nicht aufhebt. Dieses ist demnach der Christen eigene Lehre, die am Worte hängt, und schließt gewißlich, daß GOtt wahrhaftig geben werde, was er verheißen hat. Aber diese wunderbare Regierung kann der Teufel und das Fleisch nicht vertragen. Wir können es viel leichter dulden und ertragen, wo es sich bei den Menschen verzieht mit demjenigen, so uns verheißen ist, als wenn GOtt mit seiner Verheißung etwas verzieht; denn es liegt uns die Erbsünde im Wege, und hindert uns, daß wir in göttlichen Sachen das nicht leiden können, was doch in menschlichen Sachen nicht so gar schwer ist zu leiden.

50. Ist es denn nun nicht eine große Schande, daß da wir so gewisse und starke Verheißung haben vom ewigen Leben und Segen, wir gleichwohl noch zweifeln an dem Willen GOttes, oder aber doch einen sehr schwachen Glauben haben? Ich bin getauft und habe die Hoffnung zum ewigen Leben, und kann doch so fest nicht glauben, als wenn mir irgend ein Fürst einen Bierkof verheißt und seine Verheißung mit Brief und Siegel bestätigt.

51. So viel desto mehr ist nun zu loben der starke und ritterliche Glaube an diesem unserm Patriarchen Jacob. Denn da derselbe das Regiment im Hause will annehmen, welches er aus göttlicher Einsetzung und Segen hat, flieht er davon, und läßt Esau, seinen Widersacher und Feind, im Regiment und Erbe sitzen: der gebraucht und genießt alles dasjenige, so Jacob gehört. Jedoch läßt sich sein Glaube das nicht anfechten, daß er deshalb wanken oder abfallen sollte, wiewohl es sich läßt ansehen, daß der Segen gar vergeblich und nichts sei.

52. Wiewohl ich daran auch nicht zweifele, ja, ich denke demselben gern also nach, daß diese heiligen Patriarchen Menschen gewesen seien, Fleisch und Blut gehabt haben, wie wir auch haben. Denn wo Petrus, Paulus und die andern Apostel Fleisch und Blut gehabt haben: so werden ja diese Heiligen auch etwas Menschliches gefühlt und gelitten haben; denn Fleisch und Blut streitet wider den Geist. Derhalben ist daran kein Zweifel, Jacob werde zuweilen auch gefühlt haben, daß er gezweifelt und daß er damit ist angefochten worden, ob der Segen vergeblich, oder aber beständig und gewiß sein möchte; denn sein Fleisch ist wahrhaftig Fleisch gewesen und der Teufel war sein Feind. Darum hat er nicht allein äußerliche Anfechtung gefühlt, da er hat müssen in das Elend ziehen und da er aus seinem eigenen Hause und Erbe vertrieben ward; sondern es ist auch die innerliche Anfechtung dazu gekommen, die viel schwerer und höher gewesen ist, da er gedacht hat: Wie, wenn dich deine Mutter betrogen hätte? Denn du siehst ja wohl, daß dein Bruder in dem Besitze bleibt, das Erbe inne behält und des Segens gebraucht. Wie, wenn GOtt anderes Rathes geworden wäre? Denn Esau hat das Gut, ich habe nichts, denn nur bloße und eitle Worte. Und um solcher zweierlei Anfechtung willen wird ihn GOtt hernach mit einem neuen Worte trösten.

53. So soll man nun aus den Exempeln der Väter lernen Geduld haben, welche Väter allenthalben wohl sind versucht worden, aber doch nicht ohne Sünde; davon Christus allein frei gewesen ist. Jacob ist zwar nicht ohne Sünde versucht worden; er ist Christus selbst nicht gewesen, hat auch nicht ein solch Fleisch gehabt, das ohne Sünde gewesen wäre: und ist angefochten und versucht worden uns zum Exempel, daß wir lernen, was der Glaube und was ein christlich Leben sei. Denn das ist ein solch Leben, das eine gewisse Verheißung hat, die nicht allein verzogen wird, sondern gar ins Gegentheil umschlägt. Jacob soll nach göttlicher Ordnung König und Priester sein, aber es geschieht etwas weit Verschiedenes von dem, so ihm ist verheißen worden. Da wird er gedacht haben: Es wird nimmermehr etwas daraus werden; es wird nicht so hinaus gehen, wie die Verheißung lautet.

54. Aber der Glaube hat wider das Fleisch gestritten und gesagt: Du Fleisch und Satan lügst; denn GOtt hat es gesagt und verheißen, der wird ja nicht lügen, es geschehe das Gegentheil, oder ich sterbe gleich unterdeß dahin. Also hat der Glaube gesagt und aus der Finsterniß das Licht heißen hervorscheinen.

55. Darum ist der Glaube nicht eine lächerliche und kalte Eigenschaft oder Geschicklichkeit, die im Herzen faul, schläfrig und müßig sei; sondern wird umgetrieben und geplagt mit greulichen Anfechtungen, nämlich, daß es mit den göttlichen Verheißungen nichts und eitel vergeblich Ding sei. Denn ich glaube an Christum, den ich nicht sehe; ich habe aber keine Taufe, das Sacrament des Altars, den Trost im Worte und die Absolution. Ja, ich sehe aber gleichwohl der Dinge keines, die er mir verheißt, und fühle noch wohl an meinem Fleische das Gegentheil? Darum muß man nun daselbst streiten und kämpfen wider den Unglauben und Zweifel.

56. Aber der Glaube, den die Scholastiker nennen fides informis (der tote Kopf- und Maulglaube), taugt gar nichts, und kann auch die harten Stöße und schweren Anfechtungen nicht vertragen und leiden. Der wahrhaftige und lebendige Glaube aber, welcher den Zweifel überwindet, ist eine recht lebendige Eigenschaft, wenn er noch also zu nennen ist, und ist in unserm Herzen ein solches Gut, das nicht ruhen kann, welches der Teufel alle Augenblicke anficht und darüber er streitet. Wie man das vornehmlich in den letzten Zügen, wenn der Tod daher geht, lernt und empfindet, ob es eines informis sei, wie sie ihn genannt haben, oder aber ein wahrhaftiger und lebendiger Glaube. Denn ein Christ stirbt, wird begraben, wird von den Würmern verzehrt, kurz, er wird gar zu nichts; welches wahrlich allen göttlichen Verheißungen stracks zuwider ist. Was soll er aber nun thun? Da erfährt er erst, daß der Glaube keine bloße Eigenschaft ist, sintemal er den Schrecken des Todes überwindet, und sagt: Wenn du schon nicht allein den Tod, sondern auch tausend Teufel und die Hölle selbst auf deinem Halse tragen solltest, so will ich dennoch GOtt nicht zum Lügner machen. Denn also sind gottselige Menschen in ihrem Herzen gesinnet, wenn die letzten Züge dahergehen, daß man mit dem Tode kämpfen muß: Und wenn auch die Hölle alle ihre Feuerflammen ausschüttet und alles Unglück über mich ausspieet, so bleibe ich doch in dem Glauben, daß ich sage: Ich glaube an JEsum Christum, den einigen Sohn GOttes; lieber HErr GOtt, hilf mir durch Christum, daß diese Lästerung bei mir nicht entspringe, daß ich zu dir sagen sollte, du lüget; ich fühle und empfinde zwar das Gegenteil wohl, aber ich habe dein Wort, welches mir nicht fehlen kann, ja, das ist mir Alles in Allem.

\quellseite{109}

57. Auf solche Weise lobt die heilige Schrift die Väter, und hält uns Exempel ihres Glaubens vor, die da hell und klar sind, und ist sonst in der Welt keine andere Lehre, ausgenommen die heilige Schrift, die da lehrt, daß allein der Glaube die Welt überwindet; welcher Glaube ist, daß man das Wort im Herzen habe und am Wort nicht zweifele, und daß man unterdessen nur das Gegenteil leide, daß man nicht verzage und nicht müde werde bis an das Ende, das ist, bis daß die Verheißung kommt. Hiob beharrte fest, da ihn sein Weib lästerte; denn er antwortete ihr also Cap. 2, 10.: „Du redest, wie die närrischen Weiber reden: haben wir Gutes empfangen von Gott; und sollten das Böse nicht auch annehmen?“ Haben wir das Leben gehabt; haben wir bisher gegessen und getrunken, da uns Gott ernährt und erhalten hat; warum sollten wir denn nicht auch irgend ein Unglück und Widerwärtigkeit tragen? Des Tobias Weib ist auch dahin gefallen, Tob. 2, 22. Aber David hält fest und bleibt beständig ganze zehn Jahre lang, da er zwar ein König war ohne Krone, ohne Land und ohne Volk, ja, da er von den Sauliten erklärt und ausgerufen war als des Königs Feind und ein Aufrührer. Darnach aber, da ihn Absalom aus dem Lande jagte, sagte er nur also, 2 Sam. 15, 25. 26.: „Werde ich Gnade finden vor dem Herrn, so wird er mich wieder holen“ etc. Spricht er aber also: „Ich habe nicht Lust zu dir; siehe, hie bin ich, er mache es mit mir, wie es ihm wohlgefället.“ Ja, das heißt fest an der Verheißung halten: „Der Mensch lebet nicht allein vom Brode“ etc. 5 Mos. 8, 3. Das sei gesagt vom Glauben und den Verheißungen, welches beides auf diesem Orte und sonst an andern Orten auch oft muß wiederholt und den Leuten fleißig eingeschärft werden, sintemal das Fleisch allezeit widerstrebt.

58. Hier entsteht ferner eine Frage: Warum die frommen heiligen Eltern, Isaak und Rebekka, ihren Sohn zu den abgöttischen Leuten gesandt haben? Denn wir werden hernach hören, daß Laban silberne Götzen gehabt habe etc. Derhalben stoßen sie ja ihren Sohn offenbar in Gefahr, und es läßt sich ansehen, daß sie damit GOtt also versuchen. Denn wo ich mich in gewisse offenbare Gefahr begebe, alsdann mag man billig von mir sagen, daß ich GOtt versuche. Wie kann aber nun der Glaube dabei bestehen, da sie GOtt also versuchen? Mit welcher Versuchung sie ihren Sohn in Gefahr der Abgötterei und anderer sehr böser Exempel stoßen; sintemal Laban nicht allein den silbernen Götzen gedient, sondern auch geizig, schändlich karg und sonst voll anderer Laster gewesen ist, wie wir hernach sehen werden. Derhalben hat Jakob seinethalben, wo er mit ihm umgegangen ist, verderbet und verführt werden können: er hat, sage ich, mit vielen Sünden und Gebrechen, so wider den Glauben und gute Sitten sind, befleckt werden können.

59. Antwort: Wo sollte er hingehen? Er hat müssen ein Weib nehmen, er hat sich müssen schicken zum Regiment im Hause und in der Kirche. Die canaanäischen Weiber aber waren gar abgöttisch, hatten gar kein Wort Gottes; dort aber, im Hause Nahor, war noch irgend eine Erkenntniß der Gottseligkeit übriggeblieben. Denn die Linie oder der Stamm der Väter von Nahor bis auf Abraham, darin Nahor, der Bruder Abrahams, auch mitgezählt wird, hat das Wort gehabt. Denn wiewohl die Kirche auf das Haus Abrahams gekommen, so ist doch auch im Geschlechte Nahors das Wort geblieben und daß sie sich der wahren Gottseligkeit befliessen haben. Die Schrift gedenket des Patriarchen Lot nicht mehr; aber Nahor hat ein Geschlecht oder Nachkommen gehabt, in welchem Geschlecht Hiob geboren, welcher aus dem Lande Uz gewesen ist, Hiob 1, 1. Etliche wollen haben, er soll von Edom hergekommen sein; aber Hieronymus und die Andern rechnen also, daß er seinen Ursprung habe von Nahor, Abrahams Bruder. Derhalben ist Hiob Nahors Vetter, dieweil er vom Geschlecht Nahor geboren ist. Sie sagen, Bileam soll auch aus demselben Geschlecht geboren sein, welcher erstlich ein frommer heiliger Mann und gottseliger Priester gewesen, ist aber hernach verderbet und verführt worden durch das Geschenk des Königes der Moabiter, daß er gefallen ist. Und sie wollen haben, daß Bileam der sein soll, der in der Historie Hiobs Elihu genannt wird. Aber das lassen wir alles im Zweifel stehen und halten es für ungewisse Dinge. Das reimt sich aber nicht übel, daß diese Linie und Geschlecht der Patriarchen die Erkenntniß Gottes behalten habe auch außerhalb der Kirche Abrahams und also bei der reinen und rechten Religion geblieben sei. Also mag man zum Ersten auf diese Frage antworten.

60. Darnach, wiewohl man das Bekenntniß der reinen gesunden Lehre behält, so laufen doch immer die Heuchler auch mit unter, und die rechten gottseligen Menschen müssen in ihrer Versammlung diejenigen auch leiden, die dem Mammon dienen, oder sonst zu andern Sünden geneigt sind: und sind aber doch mit den Gottseligen eins in der Lehre und im Gebrauch der heiligen Sacramente. Es sind Christen dem Namen nach; denn sie streiten wider die Lehre nicht, verfolgen sie auch nicht, sondern haben ihre fleischliche Weisheit und sind gemeiniglich Bauchknechte. Wo man nun dieselbigen alle meiden und allein bei frommen gottseligen Leuten wohnen sollte, da gar keine Heuchler wären, so müßten wir aus der Welt gehen, wie Paulus sagt.

61. Die Kirche kann das nicht verhüten oder wehren, daß nicht in ihrer Versammlung Heuchler, das ist, falsche Brüder sein sollten, sofern sie gleichwohl wider die Lehre nicht streiten und uns nicht verdammen, wie zur Zeit Abrahams die Babyloniter waren: wenn sie uns nur bei sich leiden und die Wahrheit lehren lassen. Laß sie immerhin Mammonsdiener sein, laß sie ihre eigene Meinung haben, sofern sie nur still bleiben und den Glauben und meinen Frieden halten. Wo es rechte Feinde sind, so werden sie herausfahren, daß sie uns entweder nicht mehr werden leiden wollen, oder aber daß wir bei ihnen nicht mehr werden wohnen können. Solche Heuchler können wir nicht verhüten, und ist noch Hoffnung, daß sie sich dermaleinst möchten bessern. Wer da geizig oder sonst zu andern Sünden geneigt ist, und seine Sünde fühlt, dieselbe nicht vertheidigt, zu dem wollen wir die Hoffnung haben, daß er noch könne fromm werden. Es ist also eine Thorheit oder Schwachheit dieses Lebens und menschlicher Natur, und wir können nicht allenthalben vollkommen sein; denn nach dem Fleisch thun wir das nicht, was der Geist haben will. Und eben darum beten wir im Vater Unser: „Vergib uns unsere Schulden“, Matth. 6, 12. Wir vertheidigen die Sünde nicht, wir lehren keinen Irrthum oder falsche Meinungen anstatt der rechten Lehre. Es wird aber ein jeder Christ vom Teufel und von seinem eigenen Fleisch geplagt und erkennt seine Bosheit und verderbte Natur; wie St. Paulus selbst darüber auch klagt zu den Römern Cap. 7, 19., da er sagt: „Das Böse, das ich nicht will, das thue ich“ etc.

62. Also hat Nahor die Verheißung, den Vätern geschehen, behalten: wiewohl nicht so gar rein, wie Abraham, der von den Andern abgesondert war, daß er ein Vater der Kirche sein sollte. Er hat aber doch ein Licht gehabt. Derhalben hat nun Jakob keinen bessern Ort finden können, dahin er sich begeben möchte, und hat sich müssen an das gemeine Sprichwort halten, da man pfleget zu sagen: Wo man nicht Pferde hat, da muß man mit Ochsen pflügen. Denn wir sind Menschen; so müssen wir auch andere Menschen um uns leiden, welche, wenn sie schon sündigen, sich doch auch ermahnen und strafen lassen. Darum heißt das nicht GOtt versuchen, wo man mit solchen Leuten umgeht, wie Laban Einer gewesen ist, bei dem Jakob gewohnt hat. Denn es werden ohne Zweifel ihrer Viele aus dem Geschlechte Nahors selig geworden sein, welche unter den andern Heiden wie ein Licht gewesen sind mit Lehren und daß sie die rechte Erkenntniß Gottes ausgebreitet haben; wie Hiob ein trefflicher Mann gewesen ist, desgleichen seine Freunde auch; und Bileam ist ein sehr trefflicher Prophet und ehrlich recht gottselig gewesen, und hat besondere Gaben des Geistes gehabt. Er hat wahrhaftig und mit Recht im Anfang das Volk Israel gesegnet; hernach aber hat ihn der Teufel verführt und verderbt. Denn also kann es wohl geschehen, daß die reine gesunde Lehre bleibe, wenn sich schon die verführt werden, so zuvor rein gewesen sind.

63. Derhalben können wir in der Welt den äußerlichen Wandel nach die Gemeinschaft auch derer, so offenbar gottlos sind, nicht gar meiden. Die Kirche hat sich vor der Tyrannei der Römer nicht können hüten; Abraham hat Ur in Chaldäa nicht meiden können. Wenn ich aber irgend einen Geizhals weiß und der offenbar andern Leuten Unrecht thut, so soll man soviel dazu thun, daß man ihn nicht in der Kirche und zur Communion zulasse: davon kann und soll er ausgeschlossen werden.

\quellseite{110}

64. Also mag man nun auf die Frage: Warum Isaak und Rebekka ihren Sohn Jakob gen Haran gesandt haben? antworten, nämlich: Dieweil die Leute des Ortes annoch in der Lehre eines Theils mit ihnen sind einig gewesen; wiewohl sie ihre Gebrechen dabei auch gehabt haben. Und ist für und für eine beständige Freundschaft gewesen zwischen Nahor und Abraham und ihren Nachkommen. Sie sind zu der Zeit wie Lichter in derselbigen Welt gewesen, bei welchen auch die Verheißung und das Wort gewesen ist und nicht bei den Chaldäern.

65. Lasset uns aber nun die Worte und die Grammatik besehen. Denn das ist mit großer Bewegung und Anzeigung einer großen Traurigkeit gesagt worden, daß Jakob, der zum Herrn und Erben gesetzt war, auszieht von Bersaba und weithin reist bis gen Haran. Ein jeglicher mag es aber bei sich selbst abnehmen, wie hoch dieser fromme und heilige Patriarch betrübt gewesen sei. Denn er wird ohne Zweifel diese lange Reise mit vielen Thränen und öftern herzlichen Seufzen gethan haben. Denn er zieht heimlich davon, daß er sich vor dem grimmigen Zorn seines Bruders Esau verborgen möge, daß ihm der nicht nachfolgen, und auf dem Wege überfallen und ihm Gewalt thun möchte. Derhalben reist er allein, hat keinen Knecht oder Gefährten bei sich, auch niemand, der ihn geleitet hätte. Das ist wahrlich ein großer Jammer gewesen, daß er also hat müssen in das Elend und Finsterniß gehen, und von Vater und Mutter ziehen, die freundliche Gemeinschaft seiner Eltern verlassen, und vergönnen, daß sein wüthender Bruder mit seinen Weibern in seiner Abwesenheit, dieweil er in fremden Landen gewesen, seines Gefallens im Hause regiert. Und Jakob ist ein Mensch gewesen und hat menschliche Affekten und Neigungen gehabt nicht anders denn wir auch. Ja, je heiliger er gewesen ist, je mehr ihn die Bosheit der argen Menschen bewegt und ihm sein Herz gemartert hat, dazu auch seine eigene Sünde und Jammer. Denn die heiligen Menschen sind sehr zart, und werden mehr und heftiger bewegt, denn die groben Klötze und Stücke, die Mönche und Werktheiligen.

66. Derhalben ist dieses eine solche Reise, die betrübt und verdrießlich genug gewesen ist, angenommen die Gefahr, so ihm von seinem Bruder hat begegnen können. Es ist kein lustiger Spaziergang gewesen; er hat auch nicht geruht, ist nicht langsam fortgegangen, wie die pflegen zu thun, die sorglos und sicher sind, daß ihnen niemand heimlich nachstellen werde; sondern er hat geeilt, ist gelaufen, und die Eltern haben ihm darum keinen Gefährten mitgegeben, daß er desto leichter vor seinem Bruder seinen Abschied verhehlen möchte. Denn er hätte Jakob den Weg mögen verlaufen, und wie er vorgenommen hatte, ihn erwürgen. Mose hat aber hier die Figur gebraucht, welche genannt wird hysteron proteron; denn Jakob ist eher geflohen, denn Esau erfahren hat, daß er gezogen war.

67. Endlich ist zu dieser großen Bekümmerniß, so er in seinem Herzen gehabt, der Teufel mit seinen feurigen Pfeilen gekommen, damit er ihn aufgesetzt, daß er gedacht hat: Siehe, was habe ich gethan? Ich habe den Segen meines Bruders weggenommen, ich habe das Haus unruhig gemacht, ich habe meinen Bruder und seine Familie und Schwägerschaft erzürnt. An dieser Anfechtung hat es Jakob ohne Zweifel nicht gemangelt. Und wiewohl er solche Anfechtung überwunden hat, haben sie ihm doch sein Herz sehr gemartert und geängstigt. In solcher Traurigkeit, Schrecken, Sorgen und Thränen ist nun der heilige Patriarch immer fortgezogen. Denn alle Umstände sind derartig, daß sie ihm nicht ohne Grund Thränen aus den Augen gepreßt haben.

68. Haran ist die Stadt, in welche Tharah, Abrahams Vater, von Ur aus Chaldäa gezogen ist und da er gewohnt hat, und daher ist Abraham in das Land Canaan berufen worden. Nahor aber, Abrahams Bruder, und seine Nachkommen sind daselbst geblieben. Hieronymus sagt, diese Stadt sei Charran genannt worden, und ist in Mesopotamien zwischen Tigris und Euphrat gelegen, an welchem Ort die Römer eine große Niederlage erlitten haben unter Crassus, der ihr Feldhauptmann war wider die Parther, welche den Römern die Adler und ihre Fähnlein genommen und Crassus selbst jämmerlich erwürgt haben.

\begin{quote}
\textbf{V. 11. Und kam an einen Ort, da blieb er über Nacht, denn die Sonne war untergegangen. Und er nahm einen Stein des Orts, und legte ihn zu seinen Häupten, und legte sich an denselbigen Ort schlafen.}
\end{quote}

69. Mose hat im Hebräischen ein sonderlich Wort gebraucht, nämlich paga, welches man auf Deutsch nicht vollkommen kann. Denn es lautet eben, als wenn wir sagen: Er traf eben an, begegnete. Das hebräische Wort heißt auch so viel, als, Einen vertreten oder für ihn bitten; wie droben Cap. 23. V. 8. 9. steht, da Abraham sagt zu den Kindern Heth: „Bittet für mich gegen Ephron, daß er mir den Acker verkaufe“ etc. Derhalben heißt es nicht eigentlich, kommen, sondern heißt, er trifft einen Ort an, er begegnet einem Ort. Daher wird es gezogen auf die Fürbitte, darum daß der, so für Einen bittet, demjenigen pflegt zu begegnen, den er versöhnen soll. Wir sagen auf Deutsch: Er kam ohngefähr an den Ort. Darum ist es eben so viel, als wollte Mose sagen: Jakob war betrübt und gar verirrt, und hatte daran nicht gedacht, daß er an dem Ort ruhen wollte; aber es ist ohngefähr also gekommen. Er hat aber daran noch viel weniger gedacht, daß er in der Nacht würde ein Gesicht haben, das ihm erscheinen würde; sondern dieweil er auf der Flucht war, hat er entweder noch wollen weiter gehen, oder aber in der Nähe jenes Ortes bleiben; er hat sich nicht vorgenommen, die Nacht daselbst zu bleiben, sondern es gereicht also, da er nichts weniger gedachte, denn daß er daselbst die Nacht über bleiben und das Gesicht sehen würde; dieweil aber die Sonne untergegangen war, ehe denn er darauf Acht gehabt, so muß er daselbst bleiben.

70. So will nun Mose anzeigen, daß Jakob erschrocken und bekümmert gewesen sei seiner Gefahr halben; dazu auch, wie es um seine Eltern, die er verlassen müssen, stehen möchte, ob sie auch noch gesund und am Leben sein möchten: und daß er also in solcher Traurigkeit geeilt habe, auf daß er bald an den Ort kommen möchte, darnach ihn verlangt hat; und daß er doch an diesem Ort geblieben sei, darum daß ihn die Nacht überfallen hat. Die aber müde sind von Traurigkeit, von der Flucht, Bekümmerniß, werden deswegen, daß der Leib auch müde ist, letztlich schläfrig und schlafen sanft ein; wie Christus seine Jünger im Garten schlafend hat gefunden vor großer Traurigkeit. Andere können bisweilen vor Traurigkeit gar nicht schlafen: aber dieser unser Patriarch ist eingeschlafen, da er von Sorgen und auch von der Reise müde gewesen.

71. Was hat er aber daselbst für ein Bett oder Lager gehabt? Antwort: Er liegt auf der Erde, und legt einen Stein unter das Haupt, den er an dem Ort gefunden hat, das ist sein Hauptkissen. Es läßt sich ansehen, daß er sich keiner Stadt oder Dorfe habe vertrauen dürfen, da er hätte sicher sein mögen; denn er hat gegen alle Menschen Verdacht gehabt: darum bleibt er die Nacht über auf dem Felde, an einem wüsten und dunkeln Orte, da die Sonne untergegangen war. Er hat weder Heu noch Stroh gehabt, daß er hätte mögen unter sich legen; wie die Jünger im Garten auch weder Federn noch Kissen gehabt haben: sondern wie ein Leder auf die Erde gefallen, also hat er geruht, liegend oder sitzend, und hat nicht weniger süß und sanft geschlafen, als wenn er ein Federkissen oder Heu unter sich hätte liegen gehabt. Also hat ganz und gar der Patriarch Jakob in der Wüste auch geruht, da er in großer Gefahr gewesen und großen Schmerz gehabt hat, darum daß er hat müssen flüchtig werden und sich also von seinen lieben Eltern reißen lassen. Denn die Anfechtungen sind sehr groß und schwer; und wo das Herz durch solche Anfechtungen betrübt und beschweret wird, ist ihm die Ruhe vonnöthen, daß es sich ein wenig erhole und der Schmerz gelindert werden möge.

72. Das hebräische Wort meraaschah heißt, was man unter das Haupt legt, und ist ein Wort, das seinen Namen hat vom Haupte. Die Juden sagen (wie Lyra erzählt), Jakob habe drei Steine unter das Haupt gelegt, und da er erwacht, soll aus den drei Steinen Einer geworden sein. Denn an diesem Ort setzt Mose das Wort in der Mehrzahl, als seien der Steine mehr gewesen denn Einer; hernach aber setzt er es in der Einzahl, als ob es nur Ein Stein sei. Ob nun die Juden dasselbe von den Vätern haben oder nicht, kann ich nicht wissen. Es wäre zumal ein feiner Gedanke, wenn er von den Vätern hergekommen wäre. Denn es wollte sich ansehen lassen, als hätte er etliche Anzeichen wollen hinter sich lassen des künftigen Glaubens an Christum. Aber den Traum Jakobs, der da folgt, verstehen die Juden gar nicht.

\quellseite{111}
\section[Dritter Theil]{Dritter Theil. Von dem Traume, den Jakob gehabt, und der Predigt, so ihm in diesem Traume gehalten worden}

\begin{quote}
\textbf{V. 12—14. Und ihm träumete, und siehe, eine Leiter stund auf Erden, die rührete mit der Spitze an den Himmel, und siehe, die Engel GOttes stiegen daran auf und nieder; und der HErr stund oben drauf, und sprach: Ich bin der HErr, Abrahams, deines Vaters, GOtt, und Isaaks GOtt; das Land, da du aufliegst, will ich dir und deinem Samen geben. Und dein Same soll werden, wie der Staub auf Erden.}
\end{quote}

73. Dieses ist eine sehr schöne Predigt und ein trefflicher Edelstein in dieser ganzen Historie, die man genau und fleißig ansehen soll darum, wie wir zum öftermahl gesagt haben, daß man in den Legenden oder Historien der Väter und Heiligen vornehmlich darauf sehen soll, daß GOtt mit ihnen redet; um welcher Ursache willen sie auch heilig sind und heilig genannt werden.

74. Denn es ist zweierlei Heiligkeit: die erste ist, dadurch wir im Wort geheiliget werden. Die andere ist, da wir durch die Werke und Leben heilig sind. Aber diese zweierlei Heiligkeit muß man mit großem Fleiß unterscheiden. Denn die erste und die allerreinste Heiligkeit ist das Wort, darin kein Mangel, kein Flecken, keine Sünde ist, sondern es ist so gar heilig, daß ihm keiner Vergebung der Sünden vonnöthen thut; denn es ist die Wahrheit GOttes, wie Christus Joh. 17, 17. sagt: „Heilige sie in deiner Wahrheit; dein Wort ist die Wahrheit.“ Derselben Heiligkeit rühmen wir uns, die wir durch das Wort berufen sind, und dieselbe ist außerhalb uns und ist nicht unser Werk. Es ist keine wesentliche Gerechtigkeit, sondern ist eine himmlische Heiligkeit, so uns durch das Wort ist mitgetheilt worden, und dasselbe durch das mündliche Wort. Dieselbe Gerechtigkeit rühmen wir nun hoch, und setzen sie wider alle Gerechtigkeit und Heiligkeit des Papstes und aller Heuchler oder Werkheiligen; denn es ist eine unbefleckte Heiligkeit. Ich habe das Wort, ich bin heilig, fromm und rein, ohne alle Sünde und Anklage, sofern ich das Wort habe; wie Christus selbst sagt Joh. 15, 3.: „Ihr seid jetzt rein um des Worts willen, das ich zu euch geredet habe.“

75. Aber von dieser Heiligkeit weiß der Papst nichts, und alle Werktheiligen verachten dieselbe und hängen nur an der Gerechtigkeit des Gesetzes. Sie wollen nicht hinaus steigen zu der himmlischen Heiligkeit, dadurch wir GOtt angenehm und vor ihm gerecht sind um seines Wortes willen. Darum soll man die Legenden der heiligen Patriarchen vornehmlich in dem Theil merken, wenn GOtt mit ihnen redet. Denn aus dem Wort kannst du erkennen, wie groß die Heiligen sind, wenn sie auch schon nimmermehr Wunderzeichen thäten; wiewohl daselbste unmöglich ist. Ehe denn aber das Fleisch etwas thut, sind wir schon heilig durch das Wort. Derhalben schließe ich also: Das Wort ist nicht mein Werk, darum wenn ich mich meines Werkes rühme, so verliere ich das Wort; wiederum aber, wenn ich mich des Wortes rühme, so verachtet mein Werk.

76. Deß kann niemand die Papisten überreden, die stets mit der alten Geige daher kommen und sagen: Ratio deprecatur ad optima: Die Vernunft strebt immer nach dem Besten. Wir aber sollen wissen, daß zweierlei Heiligkeit ist: die eine ist das Wort, welches die Heiligkeit selbst ist. Die aber das Wort haben, denen wird diese Heiligkeit zugerechnet und die Person wird einfach für heilig gehalten; nicht um unsträflichen oder von wegen unserer Werke, sondern um des Wortes willen; und also wird die ganze Person gerecht. Darum wird die Kirche heilig genannt und wir werden auch heilig genannt; denn wir haben eine unsträfliche Heiligkeit, nicht von uns selbst, sondern vom Himmel herab. Und diese Heiligkeit ist nicht zu verachten, und wir sollen uns deß nicht schämen, daß wir heilig genannt werden. Denn wo wir uns dieser Heiligkeit nicht rühmen, so thun wir dem rechten wahren GOtt Gewalt und Unrecht, der uns heilig mit seinem Wort. Ja, sagst du, ich bin aber ein armer Sünder. Antwort: Ich weiß wohl, spricht GOtt, daß du ein Sünder bist, und wenn du das nicht wärest, so wollte ich dich nicht heiligen; du bedürfest auch des Wortes nicht. Darum aber, daß du ein Sünder bist, heilige ich dich.

77. Die andere Heiligkeit ist eine Werkheiligkeit, welche ist die Liebe, die da thut, was GOtt wohlgefällt. Da redet GOtt nicht allein, sondern ich befleißige mich auch, daß ich GOtt, der mit mir redet, folgen und gehorchen möge. Dieweil uns aber noch immer Schwachheit anhängt, darum ist dieselbige Gerechtigkeit nicht rein, sondern das Vater Unser muß da regieren und das Beste thun, und man muß beten: „Dein Name werde geheiliget“, Matth. 6, 9. Dasselbe gehört auf unsere und die Werkgerechtigkeit, welche wesentliche Heiligkeit ist und in die Heiligkeit der Zehen Gebote und in das Vater Unser gehört. Aber die erste Heiligkeit gehört in das Credo, in die Artikel des Glaubens. Denn mit den Zehen Geboten ergreife ich die Verheißung des Wortes nicht, auch mit dem Vater Unser nicht; sondern damit ergreife ich meine Liebe und meine Werke: mit dem Glauben aber ergreife ich das Wort, das ist, die Reinigkeit selbst.

78. Von diesen Dingen kann man nicht genug sagen; man kann solches den Leuten nicht genug einschärfen, und ist doch ein leichter Unterschied, nämlich; zwischen dem Gebot und der Verheißung. Es ist ein ander Wort, das ohne die Liebe und ohne meine Gerechtigkeit den Gläubigen gerecht macht; und ist ein anderes, wenn ich die Gebote GOttes ergreife, daß ich nicht stehle, nicht ehebreche etc. Aber die Papisten sind in ihrer Finsterniß also ersoffen und damit also bedeckt, daß sie diese Lehre nicht hören, ob sie dieselbe schon hören. Und wir selbst halten auch so gar fest nicht daran, als wir billig thun sollten. Derhalben lernt aus diesen Historien, wenn ihr dieselben leset, daß wir in unserm Leben allezeit über diesem Stücke etwas länger abbleiben sind, nämlich, wo GOtt mit den Patriarchen redet. Denn da ist die Legende am allerbesten und köstlichsten.

79. Nun wollen wir die Predigt selbst besehen. Und ist auch daraus wohl abzunehmen, wie Jakob so große Traurigkeit und Angst in seinem Herzen gehabt hat. Denn er ist, so zu sagen, in der äußersten Finsterniß, ist aus seinem Hause und Vaterlande vertrieben worden, ist verlassen und einsam, weiß nicht, wo er sich möchte verborgen liegen. Und ist über das der Teufel auch dazu gekommen, der auf wunderliche und auf tausendfache Weise die betrübten Herzen zu martern pflegt. Also daß man wahrhaftig sagen mag, wie das Sprüchwort lautet: Ein Unglück kommt selten allein. Denn „der Teufel gehet umher wie ein brüllender Löwe“, 1 Petri 5, 8., und sucht, an welchem Orte er am leichtesten über den Zaun steigen möge und mit welcherlei listigen Anschlägen er den Wagen, der schon auf die eine Seite hängt, vollends umstoßen möge. Wo der Zaun am niedrigsten ist, da steigt er hinüber; und wenn der Wagen wankt, so stößt er ihn gar um. Also macht er den betrübten angefochtenen Herzen immer mehr Betrübniß und häuft ihnen die Anfechtung, auf daß er sie zu Verzweiflung, in Gotteslästerung oder Ungeduld führen und stürzen möge.

80. Dies sind des Teufels Werke, also pflegt er für und für den Gottseligen heimlich nachzustellen. Darum wird Jakob ohne Zweifel außer dem Kreuz, so er äußerlich gehabt, und außer dem Gräul, mit den feurigen Pfeilen des Teufels auch getroffen worden sein, und vielleicht hat er gedacht, wie er die Erstgeburt gehabt und seinen Vater betrogen habe. Denn also pflegt der Teufel aus dem allerbesten Werk eine schändliche und große Sünde zu machen. Und daß GOtt mit ihm redet, ist ein Zeichen, daß er diese schwere Anfechtung gehabt hat.

81. Denn GOtt pflegt seine Worte nicht vergeblich zu reden, er redet nur, wenn eine große und nothwendige Ursache vorhanden ist, die ihn dazu bewegt, daß er reden muß. Und er pflegt die auch nicht anzureden oder zu trösten, die sein spotten, die sich freuen und wider ihn trotzen in Wollust oder fleischlicher Weisheit, die in aller Sicherheit, ohne Gottesfurcht und Ehrerbietung dahin leben. „Die Weisheit“, spricht Hiob im 28. Cap. V. 13., „wird nicht funden im Lande der Lebendigen“ etc., sondern unter dem Kreuz derer, so gedrückt werden und die mit geistlichen Anfechtungen kämpfen. Da ist beides, Ursache und Raum, zum Trost; da ist GOtt gegenwärtig und tröstet die Betrübten, „auf daß die Gerechtigkeit ihre Hand nicht ausstrecken zur Ungerechtigkeit“, wie der 125. Psalm V. 3. sagt. Item: „Er wird mit seinen Knechten freundlich reden, wird seinen Freunden Friede wünschen“, wie im 122. Psalm V. 8. steht. Denn wo er allzulange ausbleiben würde, so würde niemand auf die Länge solche schwere Anfechtung und teuflische Plage ertragen und beständig bleiben können. Derhalben ist dieses ein großer Trost in dem so großen jämmerlichen Schrecken, so Jakob gehabt, davon er mehr entschlafen ist, wie es scheint, denn sonst, daß er äußerlich von der Reise müde gewesen. Denn der Teufel ist auch zu seiner Flucht und seinem Elende gekommen und hat ihn innerlich im Herzen erschreckt.

\quellseite{112}

82. Dieses ist aber der Traum, so Jakob gehabt. Da ist eine Leiter auf der Erden gestanden, die rührete mit der Spitze an den Himmel, und die Engel GOttes stiegen daran auf und nieder: und der HErr selbst stand oben darauf und redete die Verheißung zu diesem dritten Patriarchen. Er redete durch keinen Menschen, sondern durch sich selbst; welches in den Historien der Väter fleißig zu merken ist, wie wir gesagt haben.

83. Die Leiter aber ist ein Gemälde oder Bild, das nothwendig etwas bedeuten muß. Denn die Engel sind Geister und Feuer, wie im 104. Psalm V. 4. steht: „Du machest deine Engel zu Winden, und deine Diener zu Feuerflammen.“ Darum bedürfen sie keiner Leiter, daran sie auf- und niedersteigen; viel weniger bedarf dessen GOtt, daß er auf der Leiter stehe, da er zu Jakob, dem Erben der Verheißung, reden will. Es ist aber auf mancherlei Weise ausgelegt worden, was diese Leiter darstellt und abbildet; und ist der Mühe nicht werth, wenn man diese Auslegungen alle zusammenstellen und aufzählen wollte.

84. Lyra sagt: Durch die Stufen oder Stiegen sollen die Patriarchen verstanden werden, die im Anfange des Evangeliums Matthäi im Geschlechtsregister Christi aufgezählt werden. Denn beiderlei Seiten bedeuten, daß Christus sowohl von Sündern, als von Gerechten oder Frommen hergekommen ist. Die Engel, sagt er, sollen bedeuten die Offenbarung der Menschwerdung Christi, so durch die Väter, Propheten und Apostel geschehen ist. Das Aufsteigen legt er aus, daß es bedeuten soll die Andacht der Heiligen oder Gläubigen, wenn sie beten. Das sind nicht böse Gedanken, so Lyra gehabt von dieser Leiter: aber es läßt sich gleichwohl ansehen, daß es die rechte Auslegung oder Allegorie nicht sei.

85. Die Glossæ ordinaria, wie man sie nennt, legt es also aus, daß die aufsteigenden Engel die seligen Geister seien, die im Himmel GOtt dienen, dann die herabsteigenden, daß sie den Menschen dienen, wie geschrieben steht in der Epistel an die Hebräer am 1. Cap. V. 14.: „Sind sie nicht allzumal dienstbare Geister, ausgesandt zum Dienst, um der willen, die ererben sollen die Seligkeit?“ Und Daniel am 7. Cap. V. 10.: „Tausendmal tausend dieneten ihm, und zehn tausendmal zehn tausend stunden vor ihm.“

86. Gregorius nennt die Engel die Prediger, welche aufsteigen und Christum betrachten, und darnach niedersteigen zu der Kirche und Gemeinde, daß sie den Gliedern der Kirche dienen mögen. Wer wollte aber alle Speculationen aufzählen? Sie sind wohl an sich selbst gottselig und gut; sind aber nicht zu rechter Zeit, oder da es Statt hat, geredet worden; wie dieses Dinges viel ist bei den Vätern. Es ist wahr, daß ein Prediger durch das Gebet erst soll aufsteigen, daß er von GOtt das Wort und die rechte Lehre empfangen möge; desgleichen soll er auch studiren, lernen, lesen und der Schrift fleißig nachdenken. Darnach soll er auch hernieder steigen und Andere lehren. Dies ist zweierlei Amt der Priester, nämlich, daß sie sich zu GOtt wenden mit dem Gebet, zum Volke aber mit guter Lehre. Dieses soll man aber an dem Orte bleiben lassen, dahin es gehört.

87. Dieweil aber im Evangelium Johannis Cap. 1. dieser Leiter gedacht wird, wollen wir lieber denselbigen Text ansehen. Denn es läßt sich ansehen, daß am selbigen Ort der HErr selbst dieses Bild ausgelegt habe. Da Philippus den Nathanael zu Christo führt, spricht Christus von ihm: „Siehe, ein rechter Israeliter, in welchem kein Falsch ist“. Joh. 1, 47. Daselbst, wie Augustinus sagt, will er uns dieser Leiter Jakobs erinnert haben, welcher auch Israel genannt wird. Christus sagt am selbigen Orte V. 50. also: „Du glaubest, dieweil ich dir gesagt habe, daß ich dich gesehen habe unter dem Feigenbaum; du wirst noch Größeres denn das sehen.“ Und spricht weiter zu ihm, V. 51.: „Wahrlich, wahrlich, ich sage euch, von nun an werdet ihr den Himmel offen sehen, und die Engel GOttes hinauf und herab fahren auf des Menschen Sohn.“ Dieser Auslegung unseres Heilandes Jesu Christi müssen wir glauben und uns daran genügen lassen. Denn er verstehet es besser, denn die andern Ausleger alle, welche doch darin alle sehr wohl übereinstimmen, daß dieser Traum das unendliche, unaussprechliche und wunderbare Geheimniß der Menschwerdung Christi solle bedeutet haben, welcher von dem Patriarchen Jakob herkommen sollte, wie GOtt spricht: „Durch deinen Samen“ etc.

88. Derhalben hat er dem Jakob selbst geoffenbaret, daß er ein Vater Christi sein würde und daß des Menschen Sohn von seinem Samen sollte geboren werden. Das hat GOtt nicht vergeblich geredet, ja, er hat dieses Bild der Leiter also gemalt, Jakob damit zu stärken und zu trösten im Glauben des künftigen Segens. Wie er droben Abraham (Cap. 22, V. 18.) und Isaak (Cap. 26, V. 4.) auch dieselbe Verheißung gegeben hat, daß sie dieselbe lehren und auf die Nachkommen ausbreiten sollten, als die gewiß wäre und nicht fehlen würde, und daß sie also den Heiland von ihrem Fleische erwarten sollten. Also stärkt GOtt den elenden und betrübten Jakob, der in fremdem Lande wie ein unnützer Kloß war, und wendet mit diesem neuen Bilde allen Segen auf ihn, daß er ihn desselben gewiß mache, daß er der Patriarch sei, von welchem der Same, der Adam verheißen war, kommen werde.

89. Derhalben müssen wir die Engel in ihrer eigentlichen Bedeutung verstehen, wie sie Christus nennt Joh. 1, 51. „die Engel GOttes“, das ist, die seligen Engel; dieselbigen steigen auf und nieder, in Christo oder über Christo. Die Leiter bedeutet das Auf- und Niedersteigen, welches geschieht auf der Leiter oder den Stufen. Wenn du die Leiter wegthust, bedeutet es Anderes nichts, denn das Auf- und Niedersteigen. Die Engel aber brauchen keine äußerliche noch auch gedachte Leiter, und ist aber doch gleichwohl ein Auf- und Niedersteigen, das ist, es ist eine englische Leiter. Dieses ist die beste Erklärung; wie Christus selbst, ohne der Leiter zu gedenken, das Nieder- und Aufsteigen der Engel ausdrückt, wie sie auf des Menschen Sohn hinauf und hinab fahren.

90. Was ist aber das für ein Auf- und Niedersteigen? Antwort: Es ist eben das Geheimniß, daß in der einen und selben Person wahrer GOtt und Mensch ist. Derhalben erfüllet die Einigkeit der Person dieses Geheimniß, und wir, die da glauben, erfüllen das Wort Christi, da er Joh. 1, 51. sagt: „Ihr werdet sehen die Engel GOttes hinauf und herab fahren“ etc. Denn wir glauben an den Einen HErrn, den eingebornen Sohn GOttes, geboren von Maria der Jungfrau, der wahrer GOtt und Mensch ist. Dieses Geheimniß ist so groß, so hoch und unaussprechlich, daß sich dessen auch die Engel selbst nicht genugsam verwundern, viel weniger begreifen können; sondern, wie St. Petrus 1. Epistel 1, 12. sagt: „Welches auch die Engel für und für gelüstet zu schauen.“ Denn die Engel können sich nicht genug freuen und verwundern dieser unaussprechlichen Vereinigung und Einigkeit der Naturen, so doch einander sehr ungleich sind, welche sie nicht erfassen und begreifen weder mit ihrem Auf- noch mit ihrem Niedersteigen. Wenn sie ihre Augen hinauf richten, so sehen sie die unbegreifliche Majestät GOttes über sich: wenn sie aber unter sich sehen, so sehen sie, wie GOtt und die göttliche Majestät den Teufeln und allen Creaturen unterworfen ist.

91. Das sind, sage ich, wunderbare Dinge: einen Menschen und die allergeringste Creatur sehen, die sich unter die andern alle erniedrigt hat, und dieselbige auch zur Rechten GOttes des Vaters sitzen sehen, erhaben über alle Engel: den Sohn GOttes im Schooße des Vaters sehen und ihn bald dem Teufel unterworfen sehen; wie im 8. Psalm V. 6. steht: „Du wirst ihn lassen eine kleine Zeit von GOtt verlassen sein.“ Item Eph. 4, 9.: „Er ist hinunter gefahren in die untersten Oerter der Erde.“ Dieses ist ein wunderbares Auf- und Niedersteigen der Engel, nämlich, in einer und einerlei Person das Höchste und Unterste mit einander vereinigt sehen, daß der höchste GOtt in der Krippe liegt. Darum beten ihn daselbst die Engel an, freuen sich und singen, Luc. 2, 14.: „Ehre sei GOtt in der Höhe“ etc. Und wiederum, wenn sie die Niedrigkeit der menschlichen Natur betrachten, so steigen sie herab und singen: „Und Friede auf Erden.“

92. Und wenn wir nun im andern Leben dasselbe auch sehen, werden wir viel anders gesinnet sein und reden denn jetzt; dieweil es schon jetzt solche große Dinge sind, welche auch die Engel nicht begreifen und deren sie auch nicht satt werden können, ja, sie begehren immer für und für, zu schauen diese unaussprechliche Güte, Weisheit, Gnade und Barmherzigkeit, so über uns ausgegossen ist, sintemal die Person, so die allerhöchste, und über alle Creatur herrlich und erschrecklich ist in ihrer Majestät, die allerniedrigste und verächtlichste Person wird. Dieses wunderbare Schauspiel werden wir in jenem Leben beschauen, und das wird eine ewige Freude sein der Seligen; wie das der Engel einige Lust und Freude ist, daß sie den Herrn über Alles sehen, welcher auch in derselben Person nichts, das ist, gar der Unterste ist.

\quellseite{113}

93. Wir fleischlichen und groben Menschen verstehen solches nicht und achten auch diese großen Dinge nicht hoch; wir haben kaum Milch und keine harte Speise geschmeckt von dieser unaussprechlichen Vereinigung und Gemeinschaft der göttlichen und menschlichen Natur, welche also gethan ist, daß nicht allein die Menschheit angenommen ist, sondern dieselbige Menschheit auch dem Tode und der Hölle ist unterworfen worden, und hat doch in solcher Niedrigkeit den Teufel, die Hölle und Alles in sich selbst verschlungen. Das ist die communicatio idiomatum, die Gemeinschaft der Eigenschaften in Christo. Gott, der alle Dinge geschaffen hat und der über Alles ist, der ist der Höchste und der Unterste, daß wir sagen müssen: Der Mensch, der gegeißelt ist, der unter dem Tode, unter dem Zorn Gottes, unter der Sünde und allerlei Unglück, ja, endlich auch in der untersten Hölle ist, der ist der höchste Gott. Ei, warum das? Antwort: Dieweil es einerlei Person ist. Es ist wohl zweierlei Natur; aber die Person ist nicht zertrennt. Darum ist beides wahr, wenn ich sage: Die höchste Gottheit ist die unterste Creatur, und ist aller Menschen Knecht geworden, ja, sie ist dem Teufel selbst unterworfen worden; und wiederum: Die unterste Creatur, die Menschheit oder der Mensch sitzt zur Rechten des Vaters, ist die allerhöchste geworden und macht sich die Engel unterthänig; nicht um der menschlichen Natur willen, sondern von wegen der wunderbaren Vereinigung und Einigkeit, die aus den zwei Naturen, so stracks wider einander sind und sonst in Einer Person nicht können zusammen gefügt werden, gemacht ist.

94. Dieses ist nun der Artikel, daran sich die ganze Welt, Vernunft und Teufel ärgert. Denn da sind in einerlei Person solche Dinge, die aufs höchste wider einander sind. Der, welcher der Allerhöchste ist, daß ihn auch die Engel nicht begreifen können, wird nicht allein begriffen, sondern ist auch also umfangen, ist also endlich, daß es kein begrenzteres und beschränkteres Ding gibt. Und ebenso umgekehrt. Er wird aber nicht begriffen denn allein im Worte, wie in den Brüsten, darin die Milch eingegriffen ist und dargereicht wird; welches Wort der Glaube ergreift, nämlich, da wir sagen: „Ich glaube an den Sohn Gottes, unsern Herrn Jesum Christum, der empfangen ist von dem Heiligen Geist, geboren aus Maria der Jungfrau (das sind die Brüste), gelitten unter Pontio Pilato, gekreuziget, gestorben und begraben, niedergefahren zur Höllen, am dritten Tage wieder auferstanden von den Todten, aufgefahren gen Himmel, sitzend zur Rechten Gottes, des allmächtigen Vaters“, und sind ihm unterthan alle Engel. Da ist Gott und Mensch, der Höchste und der Unterste, der Unendliche und Endliche in Einer Person, der sich alles Dinges entäußert und Alles erfüllt.

95. Dieses ist nun das Auf- und Niedersteigen der Engel Gottes und der Seligen, die solches anschauen, darauf beflissen sind, daßselbe auch rühmen: wie man am Tage der Geburt Christi sehen mag, da sie herab fahren, als ob droben im Himmel kein Gott sei. Sie kommen gen Bethlehem, sagen zu den Hirten Luc. 2, 10. 11.: „Siehe, ich verkündige euch große Freude, denn euch ist heute der Herr geboren.“ Und in der Epistel an die Hebräer im 1. Cap. V. 6. steht also: „Da er einführet den Erstgebornen in die Welt, spricht er: Und es sollen ihn alle Engel Gottes anbeten“, nämlich, da er jetzt in der Krippe liegt, an den Brüsten seiner Mutter, da beten sie ihn an; ja, sie beten ihn an, und unterwerfen sich für und für diesem Untersten, am Kreuze und da er niederfährt zur Hölle, da er unter der Sünde und Hölle liegt und da er alle Sünde der ganzen Welt trägt. Also steigen nun die Engel auf, und sehen den Sohn Gottes, so von Ewigkeit geboren ist; und dagegen fahren sie auch herab, da sie ihn sehen, wie er in der Zeit von der Jungfrau Maria geboren ist, und sie steigen auf oder nieder, so beten sie ihn doch immer an.

96. Auf diese Weise legt Christus diese Leiter aus. Und ich halte dafür, daß dieses die vornehmste und eigentliche Auslegung dieses Textes sei. Und dieses ist die große und unaussprechliche Herrlichkeit des menschlichen Geschlechtes, die niemand ausreden kann, nämlich, daß Gott durch diese wunderbare Vereinigung die menschliche Natur mit sich selbst verbunden hat.

97. Ambrosius und vornehmlich Bernhardus haben sonderlich große Freude an diesem Text, der über die Maßen lieblich ist, und an dem großen Werke der Menschwerdung Christi. Und sie thun recht und wohl daran; denn diese Lust wird eine Freude über alle Freude und die ewige Seligkeit sein, wenn wir da wahrhaftig unser Fleisch anschauen werden, welches uns allenthalben gleich und zugleich am höchsten und untersten Orte ist. Denn das hat er alles für uns gethan, ist niedergefahren zu der Hölle und wieder aufgefahren gen Himmel. Dieses schauen die Engel im Himmel für uns an. Und das ist es, daß Christus sagt Matth. 18, 10.: „Ihre Engel im Himmel sehen allezeit das Angesicht meines Vaters im Himmel“, sie schauen stets die Gottheit an. Und jetzt, da Gott ist Mensch geworden, fahren sie herab vom Himmel und sehen Christum an, verwundern sich des großen Werkes der Menschwerdung; sie sehen, daß er ist Mensch geworden, erniedrigt, in den Schooß seiner Mutter gelegt, und beten also den Menschen an, der gekreuzigt und verworfen ist, und erkennen ihn für Gottes Sohn.

98. Bernhardus hat die Menschwerdung Christi sehr lieb gehabt, desgleichen Bonaventura; und diese zwei Lehrer lobe ich sehr um desselben Artikels willen, deß sie so gern und so herrlich gedenken, üben sich auch darin bei sich selbst mit großer Freude und Gottseligkeit. Bernhardus meint, und hat diese Gedanken, die gottselig genug sind, nämlich, daß der Teufel gefallen sei von wegen des Neides, daß er den Menschen diese große Herrlichkeit mißgönnt hat, daß Gott würde Mensch werden. Denn der Satan, da er noch ein guter Engel war, soll gesehen haben am Angesicht Gottes, daß etwa geschehen würde, daß die Gottheit würde hernieder fahren und dieses jämmerliche und sterbliche Fleisch an sich nehmen, und würde nicht die Engel an sich nehmen. Durch solche Unbilligkeit und Neid soll der Satan bewogen worden sein, daß er wider Gott getobt hat, darum er aus dem Himmel gestoßen worden. Dieses sind Bernhards Gedanken, die nicht ungeschickt sind; denn sie fließen und kommen her aus Verwunderung der großen Liebe, Gnade und Barmherzigkeit Gottes.

99. Der Teufel ist ein sehr schöner Engel gewesen und eine sehr köstliche edle Creatur; da er aber gesehen hat, daß es also beschlossen gewesen ist, daß Gott nicht englische, sondern menschliche Natur hat annehmen wollen, ist er neidisch, zornig und unwillig geworden wider Gott, daß er ihn, der ein sehr schöner Geist war, nicht wollte an sich nehmen, und daß er der Gottheit und solcher hohen Majestät nicht hat theilhaftig werden können; da hat ihm das wehe gethan, daß das arme, elende, menschliche Fleisch ihm hat sollen vorgezogen werden, da doch die Gottheit sich mit ihm besser gereimt hätte, denn mit diesem sündlichen Fleisch, so dem Tod und allem Unglück unterworfen ist. Und es ist am allerhöchsten zu verwundern, daß diese Meinung in den Alcoran auch heimlich gekommen ist, wer sie auch hinein gebracht oder durch welcherlei Gelegenheit es gekommen sei. Es läßt sich wahrlich ansehen, daß der Teufel selbst dem Meister, der den Alcoran gestellt und geschrieben hat, solches müsse eingegeben haben, nämlich, daß darum aus dem guten Engeln Teufel geworden seien, dieweil sie Adam nicht haben wollen anbeten. Diese seine eigene Sünde hat der Satan nicht können verhehlen oder heimlich halten, darum hat er es diesem seinem Werkzeuge also eingegeben, daß er Gott damit möchte verhaßt machen, und hat die rechte Ursache des Falles verfälscht: gleich als wären die Engel dazu gezwungen worden, daß sie Adam, das ist, eine Creatur sollten anbeten, und da sie sich dessen geweigert, sie aus dem Himmel verstoßen und Teufel geworden sein sollen.

100. Dieses reimt sich schier mit den Gedanken, so Bernhard gehabt; und der Teufel verräth sich also sein selbst, womit er gesündigt habe. Er hat Gott wollen gleich sein. Da er gesehen hat, daß sich Gott dermaleinst also herunterlassen und erniedrigen würde, daß er menschliche Natur an sich nehmen würde, hat er gemeint, diese Ehre sollte ihm von Rechts wegen gebühren. Wie die alten Lehrer den Text des Propheten Jesaia im 14. Cap. verstehen und denselbigen auf diesen Fall und Sünde des Teufels ziehen, da der Prophet V. 13. sagt: „Bedachtest du doch in deinem Herzen: Ich will in den Himmel steigen, und meinen Stuhl über die Sterne Gottes erhöhen“; denn damit wäre er wahrlich Gott gleich geworden, wenn Gott ihn hätte an sich genommen in Einigkeit der Person, wie er den Menschen hat an sich genommen. Denn daß er jetzt die Menschheit an sich genommen hat, daßelbe macht diesen Menschen zum Sohn Gottes, dieweil es Eine Person ist. Dieser Mensch, geboren von Maria der Jungfrau, ist Gott selbst, der Himmel und Erde geschaffen hat. Eben solche Ehre und Herrlichkeit würden die Engel auch gehabt haben, wo der Sohn Gottes ein Engel geworden wäre und diesen sehr schönen Geist an sich genommen hätte; denn so hätte man gesagt: Dieser Lucifer ist der wahre Gott, Schöpfer Himmels und der Erden.

\quellseite{114}

101. Das wird er begehrt haben, sagt Bernhard; da ihm aber das ist abgeschlagen worden, ist er mit großem Haß, Zorn und Neid wider Gott entbrannt, daß er die menschliche Natur also hoch ehrete mit der göttlichen Natur, und daß er nun die menschliche Natur in der Gottheit wußte anbeten. Daher kommt der bittere Haß und grimmige Zorn des Teufels und der Welt, damit er unserer Natur heimlich nachstellt, und damit umgehet, dieselbige zu verderben mit vielerlei Pfeilen oder Geschossen er nur kann und mag. Denn es ist der größeste Haß und Feindschaft, so er gefaßt hat wider den Samen des Weibes, den Sohn Gottes; um den ist es zu thun. Und das ist ein alter Haß, den er im Himmel gefaßt, und der bei ihm daselbst gewurzelt und angefangen hat, also daß er nimmermehr kann ausgerottet werden. Derhalben ist diese Leiter die wunderbare Vereinigung der Gottheit mit unserm Fleisch, da die Engel auf- und niedersteigen und dessen sie sich in Ewigkeit nicht genugsam verwundern können. Das ist nun der einfältige Verstand nach der Historie und nach dem Buchstaben.

102. Danach ist eine andere Vereinigung, da wir mit Christo vereiniget werden; wie es St. Johannes sehr fein zusammen setzt, da Christus Joh. 14, 10. sagt: „Ich bin im Vater, und der Vater ist in mir“; das ist das Erste. Das Andere ist, daß er sagt: „Ihr in mir, und ich in euch.“ Dieses ist die Leiter nach der Allegorie oder heimlichen Deutung. Die Allegorie aber soll dazu dienen, daß sie den Glauben stärke; soll nicht von unserm Ding oder von unseren Werken lehren. So werden wir nun durch den Glauben hineingerissen und werden mit Christo Ein Fleisch; wie er sagt Joh. 17, 21.: „Auf daß sie alle Eines seien, gleichwie du, Vater, in mir, und ich in dir, daß auch sie in uns Eines seien.“ Also steigen wir auf in ihm, und werden durch das Wort und den Heiligen Geist hineingerissen und hangen an ihm, daß wir durch den Glauben Ein Leib mit ihm seien und er mit uns: er ist das Haupt, wir sind die Glieder. Wiederum aber fährt er herab zu uns durch das Wort und die heiligen Sacramente, da er uns lehrt und übt in seiner Erkenntniß. Derhalben, so ist die erste Vereinigung, da der Vater und der Sohn in der Gottheit mit einander vereiniget sind; die andere ist, da die Gottheit und Menschheit in Christo vereiniget wird; die dritte aber ist, wie die Kirche und Christus vereiniget werden.

103. Dieses ist der rechte Verstand, den der Teufel feind ist; und er hat ihn vielleicht im Himmel zuvor gesehen, darum hört er nicht auf, wider die Leiter und dieses Auf- und Niedersteigen zu streiten, daß er uns davon möge abziehen. Das versucht er durch alle Secten und Ketzerei, nämlich, daß er die Leute von der Erkenntniß Christi, von seiner Gottheit und Menschheit abführen, und die ganze Kirche und die Glieder von Christo abreißen möge.

104. Darum hat Gott diesem sehr heiligen Manne im Dunkeln anzeigen wollen dieses große Geheimniß von der Menschwerdung seines Sohnes, nicht allein ihn damit zu trösten, daß er in solchem großen Elende und Trübsal eine Erquickung oder Labsal haben möchte in seinem Glauben, sondern daß es auch eine Prophezeiung oder Weissagung sein sollte auf das Zukünftige für alle Nachkommen. Denn was Mose von den Patriarchen geschrieben hat, hat er nicht vornehmlich um ihretwillen geschrieben; denn sie haben dieser Schrift nicht bedurft und sind nun gestorben und zu ihren Vätern versammelt worden: sondern daß die Kirche und Gemeinde Gottes damit unterrichtet und gestärkt werden möchte bis an das Ende der Welt.

105. Denn wir können diese große Gnade Gottes und die Ehre und Herrlichkeit, so uns damit widerfahren ist, daß Gott der himmlische Vater seinen Sohn in unser Fleisch gesandt hat, nicht fleißig genug treiben und lehren, noch auch durch einiges Forschen ergreifen. Ja, daß noch mehr ist, die Engel können dessen nicht satt werden, wie gesagt ist, daß sie aufsteigen und im Himmel den Sohn Gottes anschauen, und daß sie wiederum herabfahren, da sie ihn sehen in der Krippe liegen, da er gekreuzigt, gestorben und in die unterste Hölle gefahren, und aber wiederum aus der Hölle erhoben ist und zur Rechten des Vaters sitzt.

106. Kurz, das soll man lernen, daß Gott und Mensch Eine Person ist. Wer diesen Artikel glaubt, der weiß jetzt, daß Alles begriffen und verborgen ist in dieser Person. Wie Paulus sagt Col. 2, 3. 9.: „In welchem verborgen liegen alle Schätze der Weisheit und der Erkenntniß. Und in ihm wohnet die ganze Fülle der Gottheit leibhaftig.“ Daß nun der Satan uns diese Erkenntniß und unsere Ehre und Herrlichkeit mißgönnt, uns darum neidet und immer zornig ist, das ist kein Wunder; denn es ist ein sehr hoffärtiger Geist, der diese Niedrigkeit nicht ansehen kann, wie die guten Engel thun, sondern er sieht immer auf in den Himmel und will dem Allerhöchsten gleich sein; darum ist er auch so greulich gefallen.

107. Was ferner in dieser Predigt begriffen wird, davon ist droben oftmals gesagt und gehandelt worden, und wird auch in dem, so hernach folgt, noch immer mehr und mehr wiederholt und bestätigt. Das Erste ist, daß Gott hier in der Mehrzahl genannt wird; denn er heißt elohim, das lautet, als ob man von Vielen redete; darnach, daß Gott nicht ein Gott ist der Todten, sondern der Lebendigen. Diese drei Patriarchen, Abraham, Isaak und Jakob, leben noch, da sie auch schon todt sind. Darum werden wir auch leben, wenn wir schon gestorben sind; sintemal wir eben denselbigen Gott haben. Dieses wird allzeit wiederholt und bestätigt.

108. Endlich wird das Land Canaan dem dritten Patriarchen verheissen, so sie doch auch nicht eines Fußes breit darin gehabt haben, wie Stephanus sagt im Buch von der Apostel Geschichten am 7. Capitel V. 5. Aber mit der Verheißung haben sie ein ander Land empfangen und darauf gehofft, und haben verstanden, daß Gott, da er mit Abraham und den andern Vätern geredet in diesen leiblichen Verheißungen, auch den Todten Verheißung gegeben habe, die nicht allein dieses Land, sondern auch das künftige Vaterland im Glauben haben würden. Vom Samen sagt er: „Dein Same soll werden, wie der Staub auf Erden.“ Droben (Cap. 26. V. 4.), da er mit Isaak redet, hat er ihn den Sternen des Himmels verglichen, item, dem Sande am Meer; hier aber hat er den Staub der Erde genannt. Davon ist droben gesagt worden etc.

\section[Vierter Theil]{Vierter Theil. Von der Verheißung, so dem Jakob in seinem Traum ist gegeben worden}

\begin{quote}
\textbf{V. 14. 15. Und du sollst ausgebreitet werden gegen den Abend, Morgen, Mitternacht und Mittag; und durch dich und deinen Samen sollen alle Geschlechter auf Erden gesegnet werden. Und siehe, ich bin mit dir, und will dich behüten, wo du hinziehst, und will dich wieder herbringen in dies Land; denn ich will dich nicht lassen, bis daß ich thue alles, was ich dir geredet habe.}
\end{quote}

109. Hier setzt Mose ein neu Wort, welches zuvor in dieser Verheißung nicht ist gebraucht worden, nämlich, parazta, und Mose braucht es hernach im 38. Capitel V. 29. auch. Da Perez, der Sohn Juda, geboren wird, da sagt die Wehmutter: „Warum hast du um deinetwillen solchen Riß gerissen?“ da steht in Hebräischen parazta. Und im 2. Buch Samuels am 5. Cap. V. 20., da David die Philister geschlagen hatte, sprach er: „Der Herr hat meine Feinde vor mir von einander gerissen, wie die Wasser reißen. Daher hieß man denselben Ort Baal Prazim.“ Daher kommt das Wort pariz, auf Lateinisch latro, grassator, ein Mörder, der um sich reißt, im 17. Psalm V. 4. Und daher ist auch gekommen der Name des Vogels, der auf Lateinisch parix heißt, eine Meise; denn es ist ein mörderischer Vogel. Darum hat dieses Wort eine wunderbare Bedeutung, da Mose sagt: „Du sollst ausgebreitet werden“, du sollst ausgebreitet, du wirst mächtig einhergehen und wüthen in alle vier Enden der Welt. Das ist nun anders geredet und vollständiger, denn er droben (V. 13.) gesagt: „Ich will dir dieses Land geben“; und du sollst, sagt er, nicht allein ausgebreitet werden, sondern du und dein Same soll ausgebreitet werden und hervorbrechen aus diesem Winkel in alle vier Himmelsgegenden. Er nennet aber die Himmelsgegenden nicht in rechter Ordnung, nennt erst den Abend, denn das bedeutet das Meer; darnach den Morgen, und zuletzt den Mittag. Damit wird angezeigt und verstanden nicht allein das Erbe des Landes Canaan und des künftigen Segens aus der Verheißung, sondern daß Jakob auch regieren würde in den vier Himmelsgegenden. Und daselbst wird er hernach auslegen.

\quellseite{115}

110. Er setzt aber ein treffliches Stück hinein, da er sagt: „Durch dich sollen alle Geschlechter auf Erden gesegnet werden.“ Und will sagen: Du wirst nicht allein das Land besitzen, darauf du liegst, sondern wirst also herrschen, daß der gebenedeite Same ohne allen Widerstand vor sich gehen und mit Gewalt ausgebreitet werden wird über die ganze Welt. Also redet Gott mit Jakob ganz vertraulich. Als wollte er sagen: Die Schlange wird sich unterfangen, daß sie sich wider diesen Segen setze und ihm widerstrebe. Und wird im Lande Canaan nicht allein Einen Ort einnehmen, welcher das Regiment und den leiblichen Segen haben wird, sondern wird sich wider diesen Theil auch setzen und denselben widerstreben, welcher über die ganze Welt ausgebreitet werden soll, in alle vier Himmelsgegenden. Da wird sie mit höchster Macht wider den geistlichen Segen streiten. Sie wird nicht nachlassen oder aufhören; denn ihre Feindschaft hat kein Ende: davon droben in der ersten Verheißung ist gesagt worden. Aber das arme elende Häuflein, so den geistlichen Segen hat, ist schwach und betrübt; es ist eine kleine Heerde, wie Christus sagt Luc. 12, 32., ist verworfen und verachtet; es ist nicht werth, daß es für die Kirche, ja, daß es irgend für ein Volk, viel weniger für Gottes Volk, sollte gehalten werden. Was soll das kleine arme Häuflein ausrichten können, wird der Teufel bedenken, ich will sie alle in einem Augenblick aufreiben und verschlingen. Aber laß ihn nur immer hintreiben, würgen, verfolgen, widerstreben und verhindern (spricht der Herr und hält ihm wiederum vor das Wort von dem Ausbreiten), ich will wiederum auch ausreißen, ich werde sein, wie es im Hebräischen lautet, pariz wider pariz.

111. Auf solche Weise tröstet Gott Jakob und mit ihm die Kirche, daß er seiner Nachkommen gewiß sei. Denn wiewohl Esau das Regiment hat, da Jakob verworfen und in das Elend verstoßen ist, und mag zwar unterdeß König und Priester sein, also daß es sich ansehen läßt, als ob es allenthalben gar verloren, und nichts weniger glaublich sei, denn daß Jakob sollte können ein Erbe und Regent werden im Hause oder in der Kirche: so ist doch daran wenig gelegen, will der Herr sagen, sei du nur stark und beständig. Denn du sollst nicht allein im Hause ein Erbe werden, und dir soll nicht allein der Segen Abrahams und Erbe dieses Landes zugeeignet werden, sondern ich sage dir das auch zu, daß du ein Patriarch über die ganze Welt und ein Vater des gebenedeiten Samens werden sollst, durch welchen alle Völker sollen gesegnet werden. Und daselbst will ich alsdann ausrichten, wenn ich wider den, der mit Gewalt hereindringt, auch mit Gewalt herausbrechen werde.

112. So will nun Gott hiermit so viel zu verstehen geben, daß er durch die Stärke und Kraft Gottes und des Heiligen Geistes durchbrechen wolle wider alle Weisheit und Gewalt des Fürsten dieser Welt. Und wir wissen, daß dieses also erfüllt ist, und erfahren es, daß es auch noch immer erfüllt wird. Es hat der Teufel allezeit mit erschrecklichem grimmigem Zorn wider das Volk Gottes getobt von der Zeit der Richter, der Könige und Apostel an bis zum Ende der Welt, so gar, daß man nichts Greulicheres oder Erschrecklicheres erdenken kann, denn der Fürst und Gott dieser Welt ist. Aber ich, spricht Gott, will wiederum auch ein rechter pariz sein: ich will wiederum hindurchreißen und ausbrechen, daß er empfinden soll, daß in diesem pariz eine unüberwindliche Kraft und Gewalt ist. Und daselbst richtet Gott aus durch seine Kraft, und nicht durch unsere Weisheit und Gewalt; denn alles, was an uns ist, das ist schwach und ist nichts. Aber doch in solcher Schwachheit und Nichtigkeit beweiset Gott seine Kraft; wie in der 2. Epistel an die Corinther steht am 12. Capitel V. 9.: „Meine Kraft ist in den Schwachen mächtig.“ Wenn ihr meinen werdet, daß ihr gar verschlungen und verderbet seid, so will ich hindurchreißen, und will machen, daß ihr aufstehen und regieren sollt. Wie Christus, da er gestorben und am allerschmählichsten geworden war, am allerkürzesten gewesen ist; denn in solcher Schwachheit hat er die Welt und den Fürsten dieser Welt verdammt mit aller ihrer Gewalt und Weisheit.

113. Derhalben soll man dieses Wort fleißig merken; und Mose wird es ohne Zweifel auch wohl verstanden haben und hat dazu große Lust gehabt. Desgleichen haben die Patriarchen diese trefflichen, herrlichen und nachdrücklichen Worte auch fleißig angesehen. Er hätte wohl können ein einfältige Wort brauchen, wie er zuvor einfältig gesagt hat, dilataberis: Du sollst ausgebreitet werden; aber er hat lieber wollen ein stärkeres und kriegerisch Wort gebrauchen. Und St. Paulus hat auf dieses Wort auch gesehen, da er sagt 2 Cor. 10, 4.: „Die Waffen unserer Ritterschaft sind nicht fleischlich, sondern mächtig vor Gott, zu zerstören die Befestungen“ etc. Und alle Diener des Wortes werden auch Kriegsleute und Heerführer genannt; und Gott wird auch in der Schrift ein Gott der Heerschaaren, der Kriegsleute und derer, die da streiten, genannt (auf Hebräisch parizim): und doch ohne Gewalt des Fleisches, ohne menschliche Kraft und Weisheit; sondern aus Kraft des Heiligen Geistes, in höchster Schwachheit, Niedrigkeit und Demuth.

114. Was wird aber der Ausgang und das Ende sein dieses Streites? Antwort, wie im Texte folget: „Und durch dich und deinen Samen sollen alle Geschlechter auf Erden gesegnet werden.“ Es wird keine fleischliche Tyrannei sein; ich werde nicht also herausbrechen, daß ich damit jemand wolte Schaden thun und die Leute verderben (welches die Gewaltigen in der Welt pflegen zu thun, die mit ihrer Tyrannei hervorbrechen durch das Schwert und grimmigen Zorn, nur darum, daß sie über die Andern herrschen und grausame Tyrannei üben gegen die, so unter ihnen sind); sondern mein Wüthen, spricht Gott, meine Kraft und mein Sieg soll heilsam sein, geeignet, gütig, gnädig und das da lebendig mache, das den Menschen keinen Schaden oder Leid thue, sondern jedermann segene, allen Menschen helfe, und von Sünden, von Tod und Teufel erlöse. Und durch solch Wüthen will ich die Tyrannei und Wüthen des Teufels zerstören; welcher Sieg der ganzen Welt durch deinen Samen wird nütze sein. Also erklärt nun Gott dieses Ausreißen durch diese Worte: „Durch dich sollen gesegnet werden“ etc.

115. Droben aber (Cap. 26. V. 4.) hat er gesagt: „alle Völker“, hier sagt er: „alle Geschlechter“; welches Wort sich so weit nicht erstreckt, als der Name, „alle Völker“. Denn ein jeglich Volk hat viel Geschlechter. Derhalben soll dieser Segen nicht allein insgemein unter alle Völker ausgebreitet werden, sondern soll auch insonderheit auf die Geschlechter gehen unter den Völkern. Nicht, daß sie alle aus allen Geschlechtern den Segen empfangen werden; sondern es werden ihn aus allen Geschlechtern Etliche annehmen. Darum ist an diesem Ort der Segen etwas klarer und ausdrücklicher; er ist mit stärkeren und nachdrücklicheren Worten ausgedrückt. Es heißt hier nicht: „Sie werden sich selbst segnen“, sondern passivisch: „Sie sollen gesegnet werden.“ Das Wort, das er droben gebraucht hat, ist etwas stärker und kräftiger. Und dennoch, dieweil er hier sagt, sie sollen gesegnet werden, so hat er damit genugsam ausgedrückt, daß er hervorbrechen werde in die vier Himmelsgegenden. Denn daß er sagt: „Sie sollen gesegnet werden“, heißt eben so viel, als wollte er sagen: Sie werden die Predigt hören. Es ist aber ein Unterschied zwischen den zweien, die Predigt hören, und derselbigen Predigt, die man gehört hat, glauben und sie annehmen. Wie das Wort droben nicht allein die Predigt bedeutet, die man gehört hat, sondern auch diejenigen, so die gehörte Predigt annehmen, wie Ps. 34. V. 3. stehet: „Meine Seele soll sich rühmen des Herrn“, das ist, ich werde nicht allein gelobt von dem Herrn, sondern ich, Seele, die ich gelobt bin, werde ermutigt, oder, meine Seele rühmt sich, wenn sie die Predigt vom Segen hört, sie gibt den Segen zu und nimmt ihn an.

116. So will er nun damit zu verstehen geben, daß dieser Segen soll ausgebreitet werden, welcher von Jakob herkommt. Denn das Evangelium von dem Heiland Christo kann nicht gepredigt werden, es werde denn das auch dazu gesetzt, daß er von Abraham und von Jakob geboren sei. Denn wir müssen ja einen gewissen Beweis haben, daß er wahrhaftig ein Sohn des Menschen gewesen sei, ein rechter natürlicher Mensch, und kein Gespenst, kein erdichteter Mensch, wie der Ketzer Marcion ausgebracht hat. Darum haben wir seine Voreltern, welches rechte wahrhaftige Menschen gewesen sind. Und er kann nicht genannt werden ohne diese Väter, ohne Abraham, Isaak, Jakob und Adam; auf daß er ein gewiß Geschlecht, einen wahrhaftigen Vater und rechte Mutter habe, und daß es gewiß sei, daß er vom menschlichen Samen hergekommen sei; daß er nicht die Engel, nicht irgend eine andere Creatur, sondern menschliche Natur an sich genommen habe.

\quellseite{116}

117. Darum ist den Vätern und Propheten ein gewisser Ort, da Christus hat sollen geboren werden, bestimmt worden; es sind gewisse Personen, von welchen er sollte herkommen, ernannt worden, daß wir daran nicht zweifeln können, daß er wahrhaftig unser Fleisch und Blut sei, Bein von unsern Beinen. Derhalben, wo von ihm gepredigt und gerühmt wird, da werden zugleich die Eltern auch genannt, von welchen er das Fleisch an sich genommen hat. Und wir werden damit auch zugleich begriffen, um derentwillen er der Same Abrahams, Isaaks und Jakobs, und nicht der Engel Same geworden ist. Darum prägt er uns dieses so fleißig ein und wiederholt es so oft: „durch deinen Samen“ etc.; als wollte er sagen: Du sollst gewiß wissen, daß du wirst Nachkommen haben. Du bist jetzt einsam, ohne Weib, ohne Kinder, bist in das Elend vertrieben; aber hernach sollst du Kinderkinder haben, ja, du sollst den Sohn haben, der Kinder Gottes machen wird.

118. Derhalben ist dieses eine große herrliche Verheißung, nicht allein dieses gegenwärtigen Lebens, wiewohl sie dasselbe auch begreift, sondern auch des künftigen Lebens. Denn wo er die vier himmlischen Enden erfüllen soll, so muß er mit nichten gestorben sein, oder wenn er doch stirbt, so wird er ja wiederum müssen auferweckt werden. Darum ist es eine große Verheißung und ist sehr hoch, und gar viel größer, denn daß sie Jakob sollte haben begreifen können; wie wir auch in unsern Herzen deren Größe nicht begreifen können. Denn das menschliche Herz ist zu enge, denn daß es solche große Dinge verstehen oder fassen könnte. Und ich halte dafür, daß der Patriarch Jakob, da er vom Schlafe aufgewacht und über diese Verheißung nachgedacht hat, also freudig und fröhlich gewesen sei, daß er gleichsam darüber entzückt und gar bestürzt worden.

119. Dieses soll man aber nun dahin ziehen, daß wir dadurch auch mögen gebessert werden, auf daß wir die große Gnade Gottes, so uns durch das Evangelium geoffenbart und geschenkt ist, erkennen. Denn wir werden durch das Wort absolvirt, und aus dem Befehl des Herrn Christi wird zu uns gesagt: Ich taufe dich; ich gebe dir den Leib und das Blut Christi, und deine Seele reiße ich mit Gewalt dem Teufel aus dem Rachen; ich erlöse dich vom ewigen Tod und Verdammniß, und mache dich zu einem Kinde Gottes und Erben des ewigen Lebens. Diese Worte, die Gott mit uns redet, sind so groß, ewig und unendlich, daß wir sie nicht begreifen können. Denn meine Natur ist zu schwach, als daß sie solche Worte sollte ertragen können. Und dieselbe Schwachheit zieht uns zurück und hindert uns, daß wir so viel Freude und Wonne nicht empfinden, wie in Wahrheit diese Worte und göttlichen Verheißungen mitbringen. Denn wo ich sie also annehmen und ergreifen könnte, wie ich wünsche und begehre, so wäre es kein Wunder, wenn ich vor Freude plötzlich stürbe. Denn lieber, bedenke doch, wie das so ein groß Ding ist, daß uns durch das Wort eines Menschen, durch Auslegung der Hände das Reich Gottes verkündigt und gegeben wird; dazu Überwindung und Sieg wider Teufel, Sünde und Tod; daß wir unter die Versammlung der Engel und in das Erbe der himmlischen und ewigen Güter gesetzt werden.

120. Aber wer ist unter uns allen, der dieses Glauben glaubt, wie er wohl billig thun sollte? Vielweniger aber können wir es begreifen, sondern wir ergreifen es nur, wie wir am besten können. Und ist doch dieses alles ganz gewiß und wahrhaftig. Und wir müssen wahrlich in diesem Glauben sterben, oder wir werden unserer Seligkeit verlustig werden, und auch nicht einen Augenblick wider den Teufel unverletzt und beständig bleiben können; denn er tobt greulich, ist grimmig und böse, wenn er hört, daß dieses also gelehrt wird, und verfolgt diejenigen, die es predigen, oder auch, die diese Lehre hören und annehmen.

121. Wir aber sollen dieses mit dankbarem Herzen und Freuden hören, und sollen doch zum wenigsten schwächlich glauben. Nur daß wir nicht daneben streiten und lästern, nicht verfolgen, nicht verwerfen oder verleugnen. Daß wir es nicht so fest und vollkommenlich begreifen, davon ist unser armes elendes Fleisch die Ursache; zudem ist auch unser Herz so gar enge, daß es die große unbegreifliche Herrlichkeit nicht fassen kann. Wie St. Paulus redet 2 Cor. 9, 15.: „Gott aber sei Dank für seine unaussprechliche Gabe.“ Wir bemühen uns aber und streben darnach, daß wir es endlich ergreifen mögen, wie wir ergriffen sind, Phil. 3, 12., und mit unserm schwachen Glauben fassen wir es zwar, so viel wir können; und wir müssen uns an den Trost halten, daß Gott gebietet, man soll den Schwachen im Glauben aufnehmen, Röm. 14, 1. Denn das Gesetz in unsern Gliedern widerstrebt, streitet und nimmt mich gefangen, läßt nicht zu, daß ich die vielen und großen Güter stärker ergreifen und fassen möchte. Denn Gott verheißt und schenkt uns einen unaussprechlichen Schatz, nämlich, Vergebung der Sünden, das ewige Leben, seine Gnade, das Erbe der himmlischen Güter, dazu den Sieg und Gewalt, den Teufel unter die Füße zu treten.

122. So soll nun ich armer elender Mensch und Wurm wider so viel Teufel und mit Fürsten und Gewaltigen, mit den Herren, die in der Finsterniß herrschen, kämpfen, Eph. 6, 12., und dasselbe ohne äußerliche Wehre, ohne Kraft, ohne Weisheit und ohne Stärke? Wie aber? Antwort: Ich glaube an den Samen Abrahams, den Sohn Gottes. Dieser Glaube ist unser Sieg, der wider die Welt siegt, der den Teufel überwindet und die Pforten der Hölle zerstört. Aber es geht noch mit uns fast schwach zu, daß wir solches fassen und ergreifen, und wir haben nur ein klein Fünklein Glaubens; darum, da wir noch so gar schwach sind, solches zu ergreifen, sieht man es noch nicht, welch ein groß Ding das sei, das wir glauben. Der Herr wolle uns Gnade verleihen, daß wir im Glauben beharren und beständig bleiben mögen bis an das Ende.

123. Dieses ist nun eine zweifältige Verheißung vom Samen Jakobs, durch welchen alle Völker sollen gesegnet werden, das ist, von Christo, den er in der zeitlichen Verheißung, als die etwas geringer ist, mit begreift; wie ein kleines Kindlein in Tüchlein gewickelt wird und wie man einen Edelstein in Gold setzt, oder wie ein köstlicher Schatz in einem irdenen Gefäße gehalten und bewahrt wird. Der zeitliche Segen ist groß und reich genug; aber dieser ist größer, daß er soll ausreißen und über die ganze Welt ausgebreitet werden. Er wiederholt aber den zeitlichen Segen, und bindet es in ein Bündlein mit dieser zeitlichen Verheißung von den Nachkommen.

123. Denn er sagt also: „Siehe, ich bin mit dir“. etc. Als wollte er sagen: Du hast nun die ewige und zeitliche Verheißung von den Nachkommen; nun will ich dieses dritte Stück auch dazu thun, nämlich, daß ich auf dieser Reise und in deinem Elende mit dir sein will; so viel deine Person belanget, will ich dich selbst auch schützen und vertheidigen. Denn damit diese Dinge, so ich dir verheißen habe, erfüllt werden mögen, so mußt du leben und geschützt werden. Dein Bruder Esau stellt dir heimlich nach; der Teufel ist dir feind, du schläfest oder wachest: aber ich will mit dir wandeln und will dein Gefährte sein auf der Reise. O wie ist das so ein gewünschter und starker Gefährte. Ich will nicht allein mit dir gehen, daß ich dich geleite; sondern ich will deine Mauer und dein Schutz sein, ich will für dich streiten, der ich dein Sohn werden soll von deinem Samen, durch welchen alle Geschlechter der Erde sollen gesegnet werden. Ich will mit dir sein, will dich schützen und beschirmen.

124. Dieses ist wahrlich ein großer Trost, welchen Jakob fast hat und glaubt, aber mit einem solchen Glauben, der schwach genug ist; wie wir hernach hören werden, da er mit sehr schweren Anfechtungen geplagt werden und daß es sich ganz anders gestalten wird, als ihm Gott verheissen hat. Denn er wird müssen leiden und ertragen, daß ihm Laban unrecht thut und ihn nur wohl plagt zwanzig Jahre lang. Heißt das Tag und Nacht Stehen gehabt? Dennoch hat Gott seine Verheißung gehalten, und es ist endlich dahin gerathen, wie die Verheißung gelautet hat; wiewohl es sich nach dem Fleisch ansehen ließ, als ob gar das Gegentheil geschehe.

125. Derhalben ist es eine treffliche Verheißung, und, wie ich zuvor gesagt, soll man sie fleißig betrachten. Denn dieses Stück ist vornehmlich in den Legenden der Väter zu merken, da sie Gott hören reden. Denn unser Leben besteht in der innerlichen Heiligkeit, welche das Wort ist, oder der Segen, der uns zugeeignet wird; es steht nicht in unsern Werken, wie auch die, so unter den Widerwärtigen die Allerbesten sind, streiten. Wiewohl wir nicht leugnen, daß man soll Gutes thun, daß man den Leib kasteien und den Lüsten des Fleisches widerstreben soll. Die Welt aber hat zu solcher Heuchelei und zu dem großen Schein der Werke Lust, daß sie nur das Wort Gottes von den Augen hinweg thun möge. Denn die Thorheit der Menschen kann zwischen Gottes Wort und unserm Werk keinen Unterschied machen, welches doch ein leichter Unterschied und auch sehr nöthig ist; den Heuchlern aber ist es schwer und unmöglich, nicht zu verstehen, sondern, daß sie es sollen annehmen.

\quellseite{117}

126. Unsere Widersacher schreien heftig, es sei allerdings vonnöthen, daß man eine äußerliche Werkgerechtigkeit habe. Dessen verwundert sich die Welt und lobt es hoch. Ist es aber nun daran genug, wenn du mir auch einen Mann rühmest, daß er der Allerheiligste sei? Es haben wahrlich die Heiligen große und löbliche Werke gethan; und wir bekennen auch, daß die Liebe nöthig sei, und daß man sich der fleischlichen Lüste enthalten und das Fleisch tödten soll 2c. Ja, sagen die Widersacher, ihr thut aber solches nicht. Antwort: Wenn wir nun dieses alles thäten, wären wir denn dadurch gerecht? Und warum hebst du diese Dinge so hoch und verwunderst dich dessen so gar sehr, daß es so gar köstlich Ding sein soll? Das thust du darum, auf daß du damit die andere göttliche Heiligkeit unterdrücken mögest.

127. Wir geben zwar das gerne zu, daß man die Werkgerechtigkeit haben muß; aber doch also, daß sie der vorigen Gerechtigkeit nicht vorgezogen werde, und daß man einen Unterschied mache zwischen dem, das kostbar, und dem, das gering ist, zwischen dem Großen und Kleinen; welcher Unterschied nöthig ist, beide im ganzen Leben und sonderlich auch in dieser Lehre; ja, verstehen doch die unvernünftigen Thiere, daß ein Unterschied sei unter dem Futter; ein Hund versteht, daß Fleisch besser sei denn Brod.

128. Derhalben laß die Werkheiligkeit etwas sein und laß sie auch nöthig sein: wir wollen aber haben, daß man über sie erheben und ihr vorziehen soll die Heiligkeit des Wortes, das ist, die Verheißung, nämlich, daß wir vor Gott leben nicht durch die wesentliche und innere Heiligkeit, sondern durch die Heiligkeit des Wortes und des Glaubens.

129. Darum soll man einen rechten Unterschied machen, nämlich, daß eine andere Heiligkeit sei die der Werke und eine andere die des Wortes oder Glaubens. Diese überwindet den Teufel und macht mich zum Kinde Gottes, welches nicht geschieht durch die Heiligkeit des Fleisches oder der Werke, auch der Zehen Gebote, darin wir diese Heiligkeit nicht einmal begreifen; denn wir haben nur etlichermaßen einen Anfang davon. Derhalben soll dieses die erste und vornehmste Predigt sein, von der Heiligkeit des Wortes und Glaubens, und dieselbe soll man nicht gering achten oder verstellen; wie die Widersacher die Heiligkeit der Werke mit List treiben, auf daß sie damit die Heiligkeit des Wortes oder Glaubens verdunkeln mögen. Wir loben die Tugenden, so die Patriarchen gehabt haben, als, ihre Bescheidenheit, Geduld, Demuth, Keuschheit, Liebe, und daß sie die Pilgrimschafften mit rechtem göttlichen Gehorsam angenommen und sich in andern Werken der Gottseligkeit immer geübt haben: dasselbe lehren wir und folgen demselben auch nach. Aber dem allen müssen wir dennoch dieses weit vorziehen, wenn Gott mit ihnen redet. Daß man Gottes Wort hört, der die Verheißung und den Segen gibt, ist viel ein Anderes, denn wo du ein Gebot hörst, darin dir geboten wird, etwas zu thun oder zu leiden.

130. Das ist klar und offenbar genug, daß ein Unterschied ist zwischen dem Wort, darin Gott etwas verheißt, und dem Gebot, darin er gebietet, etwas zu thun. Derhalben soll man auch die Legenden der Heiligen und das ganze Leben in diese zwei Stücke theilen, nämlich, in Gottes Wort und in unsere Werke. Jenes ist allein Gottes; dieses aber ist unser, nämlich, die Liebe, Geduld, und daß wir das Fleisch fasten und kreuzigen. Dieses ist ein Leben, aber ein Werkleben, das die Heiligkeit des Wortes nicht ergreift, welches das Leben ist der Seele. Darum geht es gar nothwendig vorher; denn es ist das Wort und die Verheißung Gottes, der den Segen und die Verheißung gibt, der uns zu Gnaden aufnimmt und uns unsere Sünden vergibt. Ja, das soll und muß vorher gehen: „Mein Sohn, dir sind deine Sünden vergeben“. Matth. 9, 2.; ich bin dein Gott; fürchte dich nicht vor dem Teufel; ich bin mit dir; ich will dich bewahren; ich will dich nicht verlassen. Nach der Verheißung aber der Sünden folgt dann weiter: „Stehe auf, heb dein Bette auf und gehe heim“ 2c., Matth. 9, 6. Das geht vorher: „Mein Sohn, sei getrost“, das ist, du sollst zu allererst erkennen, daß ich dir gnädig und mit mir versöhnet bin, daß ich dir günstig und gewogen bin, daß ich dich von deinen Sünden absolviere und losspreche. Du sollst zuerst meine Wohlthat und Gaben annehmen, daß du von Sünde und Tod erlöset werdest. Danach, wenn du bist gesund geworden, so heb dein Bette auf und gehe hin, lehre und thue gute Werke. Also, da Jakob erstlich im Glauben gestärkt ist von wegen der Gnade und des Segens, geht er darnach auch hin, thut Gutes, und leidet, was ihm auferlegt wird; wie im Text folgt: „Da Jakob von seinem Schlaf aufwachte, sprach er“ 2c.

131. So haben wir nun bisher gehört die große herrliche Verheißung, so dem Patriarchen Jakob ist gegeben worden, nicht allein von dem Segen der Nachkommen und von Christo, sondern auch, daß ihn Gott gegenwärtig führen und in seinem Elend schützen wolle, daß er seines Lebens und Nahrung gewiß sein möge, dazu auch, daß er wiederum aus dem Elende kommen, und alles Leiden und Trübsal, so ihm künftig begegnen würde, überwinden soll. Denn wir werden hernach sehen, wie er mit so großem Unglück hat kämpfen müssen, ehe diese Verheißung ist erfüllt worden. Und er soll dennoch durch diese mancherlei Todesgefahren und so viel unendliche Unglücksfälle wiederum zu dem Erbe gebracht werden, soll der Gesegnete bleiben, und der rechte Erbe der Nachkommen und des künftigen Erbes sein.

132. Derhalben sind dieses sehr merkwürdige Exempel; denn sie zeigen nicht allein an, daß Gott mit den Vätern geredet habe, sondern bezeugen auch, daß Gott in seinem Wort nicht lüge, wenn sich Alles ansehen lässet, als ob es gar widersinnig hinaus gehen und geschehen werde, daß die Welt unsinnig ist, das Fleisch, die Widersacher und der Teufel auch mit allen Pforten der Hölle sehr toben und wüthen. Denn er gebietet, so wird es geschaffen. Es ist aber ganz schwer, daß wir Gläubigen sollen hindurch brechen; denn der Teufel richtet seinen grimmigen Zorn nur ganz und gar auf dies Einzige, daß er uns vom Wort abreiße, daß wir müde und verzagt werden sollen, entweder von wegen der vielen Trübsal, oder aber, daß es damit so lange währet, und daß wir also das Wort verlassen und verwerfen.

Derhalben werden nun diese zwei Stücke zu einander gesetzt, als nämlich die Allmächtigkeit Gottes, und die Schwachheit oder, wenn man also sagen möchte, die Unmächtigkeit des menschlichen Herzens, welches so lange dulden, leiden und beharren soll, bis das Wort erfüllt werde. Also redet Bernhardus vom Glauben der Jungfrau Maria, da ihr der Engel verkündigt hatte, daß sie sollte die Mutter Christi werden, daß die Kraft des Glaubens der Jungfrau, die den Worten des Engels hat glauben können, nicht ein geringeres Wunderwerk gewesen sei, denn die Menschwerdung des Wortes oder des Sohnes Gottes selbst. Derhalben sind in den Historien der Heiligen die größten Dinge die Worte, welche Gott mit den Heiligen redet. Man soll zwar ihre Tugenden und ihre Thaten rühmen, und Gott will dieselbigen auch haben: aber man soll sie an den untersten Ort stellen wie die Füße. Das Haupt aber in dem Leben der Heiligen ist das Wort Gottes selbst.

133. Daß aber Alles, wie es scheint, auf das Gegentheil von dem hinaus kommt, so Gott im Wort verheißt, das geschieht durch die Gewalt des Teufels und Schwachheit unsers Fleisches, welches immer noch zweifelt an dem verheißenen Wort, daß es ewig und wahrhaftig sein sollte. Wie wir hernach hören werden, daß diese gewaltigen und reichen Verheißungen, da Gott zu Jakob sagt: „Ich will mit dir sein, will dich schützen und bewahren“ 2c., mit großer Gewalt angefochten werden.

134. Wir, die in solcher Gefahr und Beschwerung nicht sind, damit Jakob gekämpft hat, lesen dieses Ding schläfrig und nur obenhin. Wir denken, es sei leicht und gering Ding; aber wenn wir an seiner Statt wären, so würden wir viel anders gesinnet sein. Denn wir, die Gott in seinen Verheißungen glauben und durch das Blut Christi erlöset sind, sehen, daß bei den Christen Alles also zugeht, als ob es eitel verloren und unmöglich Ding sei, also, daß der Mensch nur schlechthin muß getödtet werden mit dem Fleisch, und mit aller Weisheit und Vernunft, und muß allein am Wort hangen; wie das an den heiligen Patriarchen auch wohl zu sehen ist. Und wer das Wort nicht vorzieht, und mehr davon hält, als von seinem eigenen Rath und Weisheit, ja, auch von seinem eigenen Leben, der mag immer hinfahren, und mag ein Heide und Epicurer werden; denn er wird nimmermehr Frucht bringen oder nütze sein, weder sich selbst, noch andern Leuten. Denn er will die Dinge, so unmöglich und unbegreiflich sind, mit seiner Weisheit und Rath regieren; darum ist es mit ihm verloren. Derhalben wird es nicht alles nach unserm Willen, und wie wir es bedacht haben, hinaus gehen. Das muß man schlechterdings lernen und für gewiß halten.

135. Also, wenn wir von Sünden absolvirt werden, haben wir das Wort, welches gegenwärtig ist und fest steht in den Verheißungen Gottes. Müssen wir aber nicht viel leiden, das dem stracks zuwider ist, und großen Streit führen wider den Teufel, Tod, Hölle und unsere Vernunft oder Verstand? Und dasselbe thut dem Fleisch sehr wehe. Jedoch erfüllt Gott unterdeß gleichwohl seine Verheißung treulich denen, so da glauben und in Geduld auf ihn harren; wie in der Epistel an die Hebräer steht im 11. Cap. V. 33.: „Die Heiligen haben durch den Glauben Königreiche bezwungen, Gerechtigkeit gewirket, die Verheißung erlanget.“ Also ist dem Patriarchen Jakob vornehmlich Trost vonnöthen gewesen, beide in diesem Elende, und auch hernach, da ihn Laban hart betrübet hat. Denselbigen Trost aber hat er ohne großen Kampf nicht behalten können, wiewohl ihm GOtt klare und gewisse Verheißungen gegeben hat, da er sagt: Ich will dich nicht lassen, bis daß ich gethan habe, was ich zu dir geredet habe. Wie die Jünger zur Zeit des Leidens der Worte Christi und alles Trostes gar vergessen haben: Er hatte ihnen eine lange Predigt gethan, wie dieselbige Predigt St. Johannes in drei Capiteln beschrieben hat, da er sagt, Joh. 16, 1.: „Solches habe ich zu euch geredet, daß ihr euch nicht ärgert“; wenn die Stunde kommen wird etc. Was aber solcher trefflicher Trost und Vermahnung bei ihnen für Nutzen geschafft, hat man wohl gesehen, da er ist gekreuzigt worden. Da haben sie an keines gedacht, was sie zuvor gehört und gesehen hatten, das Christus geredet und gethan hatte. Dieses wird uns aber darum vorgehalten, daß wir lernen sollen, welch ein groß Ding der Glaube sei, daß er keine faule und müßige Eigenschaft oder Geschicklichkeit sei, wie sich die Papisten träumen lassen.

\quellseite{118}

\section[Fünfter Theil]{Fünfter Theil. Was die Erscheinung, dem Jakob geschehen, gewirket}

\begin{quote}
\textbf{V. 16. Da nun Jakob von seinem Schlaf aufwachte, sprach er: Gewißlich ist der HErr an diesem Ort, und ich wußte es nicht.}
\end{quote}

136. Diese Erscheinung oder Gesicht ist Jakob nicht widerfahren, da er gemacht hat, wie sonst gemeiniglich die Erscheinungen der Heiligen geschehen sind; wie droben (1 Mos. 18, V. 1.) der HErr dem Abraham zu Mamre durch die Engel erschienen ist. Das sind Erscheinungen derer, die da wachen, und sind etwas gewisser, denn die, so im Traume geschehen. Und dennoch, wenn jene nicht auch ihre rechte Beziehung, oder Bewegung der Herzen, und den Glauben haben, sind sie nicht recht oder wahrhaftig.

137. Von den Träumen aber kann ich nicht so leichtlich disputiren; denn es haben beide die Gottseeligen und auch die Gottlosen Träume; wie auch in den heidnischen Historien Träume beschrieben werden, als des Julius Cäsar, des Galba, Brutus, Cassius und Anderer mehr, die solche Träume gehabt haben, daß auch die Wahrheit und die That selbst hernach gefolgt, und es nicht allein so hinausgegangen, wie die Träume gelautet haben, sondern sie haben auch die Träume verstanden, da sie vom Schlafe aufgewacht sind. Und denselbigen Träumen sind die gleich, so Pharao 1 Mos. 41, 1. ff. und Nebucadnezar Dan. 2, 1. ff. gehabt haben, welche in der heiligen Schrift stehen.

138. Derhalben wollen wir einen solchen Unterschied machen zwischen den Träumen. Etliche sind weltliche oder sonderliche Träume, das ist, die da sonderliche oder aber öffentliche Angelegenheiten betreffen. Dieselbigen Träume haben ihre Art und ihre eigene Gestalt; darum soll man sie ziehen, dahin sie gehören und da sie ihre Ordnung und Statt haben. Und wir wollen an diesem Orte von denselben Träumen nicht sagen. Etliche Träume aber sind, davon die heilige Schrift redet, die da gehören zu den Kirchensachen und zur Seligkeit. Das sind die rechten Träume, davon an diesem Orte zu sagen ist.

139. Der Traum Pharaos bezog sich auf die unfruchtbaren und fruchtbaren Jahre; derhalben ist es ein weltlicher und sonderlicher Traum. Also war Nebucadnezars Traum auch ein leiblicher Traum und bezog sich auf jenes Königreich. Solcher Träume sind viel, die GOtt unter die Heiden austheilt, wie er viel andere Gaben auch gibt, wie in der Apostelgeschichte stehet im 14. Capitel V. 17.: „Er thut uns Gutes, gibt vom Himmel Regen, und fruchtbare Zeiten, erfüllt unsere Herzen mit Speise und Freuden.“ Also hat er die Königreiche und Kronen ausgetheilt unter die Gottseligen sowohl als unter die Gottlosen.

140. Man soll aber insgemein von den weltlichen Träumen merken, daß auch sie mit der gegenwärtigen Lage in Beziehung stehen, und auch einen solchen großen Eindruck im Herzen machen müssen, daß der den Traum hat, ganz bewegt und bestürzt werde. Es dürfen nicht melancholische Träume sein, die sich nicht zur Sache reimen, davon gesagt wird im Buch Jesu Sirach am 34. Capitel V. 4.: „Was unrein ist, wie kann das rein sein? Und was falsch ist, wie kann das wahr sein?“ Da werden alle Träume unrein genannt, und wer denselbigen glaubt, der wird betrogen. Und Cato sagt auch: Somnia ne cures etc.: Du sollst die Träume nicht achten.

141. Man hat aber allezeit von den Träumen disputirt: Ob man denselben glauben soll, oder nicht? Die falschen Propheten pflegen ihre Träume hoch zu rühmen; darum verbietet GOtt oftmals, daß wir denselben nicht glauben sollen. Denn wenn der Traum nicht der Wirklichkeit entspricht, so wird der betrogen, der ihm glaubt, und der Traum ist eitel und falsch. Nichts desto weniger aber weiß man gewiß, daß auch bösen gottlosen Menschen Träume eingegeben werden, die da wahr sind, und solche Beziehung haben, die sich mit ihnen reimet. Von solchen Träumen haben sie etliche sonderliche Regeln gemacht. Als, zum Ersten, sollen es nicht ungewisse Bilder oder Gedanken sein, die bald vergehen oder verschwinden, sondern sollen ihre Beziehung haben auf das Ding, so gegenwärtig ist: darnach, sollen sie das Herz also bewegen, daß der den Traum hat, davon geängstet und erschrocken werde; wie Daniel zum König Nebucadnezar sagt, Dan. 2, 29.: „Du König dachtest auf deinem Bette, wie es doch hernach gehen würde“ etc. Er zeigt damit an, daß der König ängstlich und bekümmert gewesen ist im Traume um das Königreich und wie es nach ihm damit gehen würde. Derhalben muß eine Bekümmerniß dabei sein, und darnach eine sehr große Bewegung oder Eindruck, dadurch der, so den Traum hat, beweget und erweckt wird.

142. Aber davon disputire ich nicht; denn wir müssen auf die Träume sehen und Achtung geben, von welchen Mose redet, so er doch die Träume an andern Orten verdammt. Denn also spricht der HErr zu Aaron im 4. Buch Mose am 12. Cap. V. 6. 7. 8.: „Ist jemand unter euch ein Prophet des HErrn, dem will ich mich kund machen in einem Gesicht, oder will mit ihm reden in einem Traum. Aber nicht also mein Knecht Mose, der in meinem ganzen Hause treu ist. Mündlich rede ich mit ihm“, das ist, wo jemand das Kirchenregiment zu verwalten hat, mit dem will ich durch mein Wort, oder durch ein Gesicht, oder durch einen Traum reden. Er sagt aber, er wolle mit Mose auf eine viel andere und höhere Weise reden, nicht durch irgend einen Prediger oder Jünger der Propheten, auch nicht in einem Gesicht oder Gleichniß, auch nicht in einem Traum: sondern, ich rede, sagt er, „mündlich mit ihm.“ Daselbst lobt GOtt gleichwohl seine Offenbarung; als zum Ersten, da er sich mit seinem Worte offenbart. Zum Andern, in einem Gesichte wie im Propheten Jesaia im 6. Cap. V. 1. ff. ein Gesicht beschrieben wird. Und zum Dritten, durch einen Traum.

143. Diese drei Arten lobt GOtt, aber an den Propheten, das ist, an denen, die im Kirchenamte sind, und vornehmlich, wenn die Träume sich auf das Wort beziehen und sich mit dem Worte reimen. Da muß man Achtung darauf geben, ob sie auch den, der den Traum hat, bewegen, und ob sie ihre Beziehung haben, und daß sie nicht wider das Wort seien. Und gilt gleich so viel, als wo das Wort selbst entweder im Gesicht, oder aber durch die Stimme derer, die da machen, vorgetragen wird. Wie Jakob das Wort und das Gut auch gewiß hat; denn er hat den Segen empfangen, und ist von GOtt berufen, daß er soll ein Patriarch sein: ist zum Bischof und zum Herrn gesetzt worden, die Kirche und Gemeinde GOttes, dazu auch das Haus seines Vaters zu regieren; und das ist geschehen durch GOttes Wort und durch der Segen seines Vaters Saat. Wo man nun den Grund hat, da soll man den Traum auch ansehen, ob er sich auch mit demselben Grunde reime oder nicht; man soll sehen, ob es Stoppeln oder Gold sei, das auf den Grund gebaut wird, wie Paulus redet 1 Cor. 3, 12.

\quellseite{119}

144. Eben auf dieselbe Weise ist auch im Amte der weltlichen Regierung der Grund gelegt. Wiewohl die Heiden das Wort nicht haben, so haben sie dennoch die Sache und den Beruf, dadurch sie zum Regimente und Herrschaft im Kriege und auch im Frieden verordnet und gesetzt sind. Wenn nun irgend einem heidnischen Könige oder Hauptmann ein Traum eingegeben wird, der seine rechte Beziehung hat, das ist, der sich mit der That an sich selbst und mit seinen Gedanken reimet und übereinkommt, denselbigen Traum soll man mit nichten für vergeblich und unnütz halten.

145. Und darum soll man vielmehr den Traum Jakobs für wahr halten und daß er eine Bedeutung haben müsse; denn es ist ja ein öffentlicher Beruf Gottes vorhergegangen, daß er zum Priesterthum und Erstgeburt verordnet und gesetzt ist, und der rechte Grund ist die Verheißung. Er ist aber darüber bekümmert, wie der ganze Bau auf den Grund gesetzt und in künftiger Zeit möge vollbracht werden; sintemal ihm viel Widerwärtigkeit begegnet, daß es scheint, als wolle der ganze Bau dahin fallen. Da hebt sein Herz an zu wanken und zu zweifeln an dem glücklichen Fortgang und an der Erfüllung der Verheißung. Darum ist nun GOtt da, und gibt ihm einen sehr klaren Traum, der sich auch sein reimet mit dem Dinge, so ihm verheißen war, und der eine solche Bedeutung hat, die auch mit dem Fundament überein kommt. Ueber das, so ist auch die Bewegung dazu genommen, wie der Text sagt: „Er fürchtete sich“, das ist, er ist zur Ehrerbietung bewogen worden, da er gesehen hat, daß sich der Traum mit dem Grunde so fein reimte.

146. Das sind rechte und bedeutende Träume, die nimmer ohne Wirkung abgehen, wie auch das Wort selbst ohne Frucht nicht abgehet; und sie bringen mit sich den rechten Grund und Bewegung des Herzens, der mit dem Grunde übereinstimmt; und kommt zugleich der Glaube auch dazu, der sich mit allen beiden reimet.

147. Wenn es aber vergebliche, unnütze Träume sind, wie sich solcher Träume die falschen Propheten rühmen, so soll man auf die rechte Regel sehen, ob auch dein Herz erschrocken und bewegt worden sei; wiewohl der Satan solches auch kann ausrichten und zuwege bringen, so sollst du dennoch den Traum mit dem Wort vergleichen. Wenn es dir anders träumt, als das Wort lautet, sollst du wissen, daß der Traum falsch und eitel ist. Dieser Traum aber, den Jakob gehabt, kommt sehr fein überein mit dem göttlichen Worte, welches er hört vom Himmel herab kommen: „Ich bin der HErr“ etc.

148. Derhalben irren die Gottlosen in der Auslegung und Verstand der Träume, wie sie sich auch sehr verstoßen, wenn sie die Zeichen und Wunder auslegen; denn sie haben nicht Acht auf das Wort, ja, sie haben kein Wort. Wie es den Wiedertäufern zu Münster ist ergangen, die einen Bogen in den Wolken gesehen hatten und eine blutige Hand dabei: dasselbe haben sie auf sich gezogen, als sollte es ein Zeichen sein des Sieges, den sie haben würden, so es doch ihnen gedroht hat, daß sie würden verderbet werden; wie das Ende solches auch gezeigt hat. Sie haben aber in der Auslegung geirrt, weil sie auf den Grund und Beruf keine Achtung gehabt; denn sie hatten kein Wort oder Gewalt, das Schwert zu führen, sintemal sie nach göttlicher Ordnung und Beruf nicht Regenten gewesen sind. Wenn dergleichen etwas mit Zeichen, Wundern oder Träumen geschehen wäre, da Alexander oder Julius Cäsar regiert haben, so hätte man es leichter auslegen können, warum daß sie GOttes Wort und Ordnung hatten. Derhalben sind die weltlichen Träume auch falsch und vergeblich, die da ungewiß und ohne Grund sind, das ist, wenn es nicht eine öffentliche Person oder die im Regimente ist, und wenn keine Bewegung und keine rechte Beziehung nach der Regel dazu kommt.

149. Aber die Träume, so die Kirche oder Religion belangen, die das Wort, die den rechten Grund und Regel haben, die sind gewisser; wie sich ein Christ vornehmlich darauf verläßt, daß er getauft ist. Wenn nun der Traum mit dem Grunde übereinstimmt, daß er zur Stärkung der Gottseligen und Schrecken der Gottlosen dient, einen solchen Traum kann man ohne Gefahr annehmen. Denn also wird Jakob durch diesen Traum auch gestärkt, der mit seinem Segen in Beziehung steht. Er weiß, daß er gesegnet ist, und daß GOtt sein Vater und sein Schutz und Schirm sein werde, wie das Wort der Verheißung lautet; hiebei aber sein Herz bekümmert und erschrocken ist aus Furcht vor seinem Bruder, der ihm heimlich nachstellt, und gedroht hat, ihn zu erwürgen, darum erscheint ihm der HErr in solcher großen Traurigkeit und redet mit ihm im Traume. Welcher Traum mit dem Dinge, so ihm verheißen war, übereinstimmt und hat die rechte Beziehung nach der Regel: denn eben derselbige HErr redet mit ihm, und redet auch eben dasselbe, so er zuvor im Segen gehört hatte. Es ist eben derselbe HErr, dasselbe Wort, das er in seines Vaters Hause gehört hat.

150. Aber das Gespräch, so Jakob mit GOtt gehabt, ist viel herrlicher, als das er mit seinem Vater gehabt hat. Derhalben, da er aus seinem Schlaf aufgewacht, verwundert er sich und spricht: „Ich hatte gemeint, der HErr wäre allein in meines Vaters Hause; darum habe ich bisher in der Gemeinde meines Vaters GOttes Wort gehört: ich finde aber an diesem Orte GOtt auch. Denn eben derselbe HErr, den ich in meines Vaters Hause habe rühmen und von dem ich habe predigen hören, redet hier auch mit mir. Ja, dieses Gespräch ist viel trefflicher, also daß es sich ansehen läßt, daß viel mehr an diesem Orte das Haus GOttes sei als in meines Vaters Hause; und ist doch gleichwohl dieser Ort mitten unter den Feinden gelegen.

151. Darum lernt Jakob, daß die Kirche und Gemeinde GOttes auch mitten unter den Feinden ist; wie in 110. Psalm V. 2. steht: „Herrsche unter deinen Feinden.“ Das ist wahrlich trefflich. Denn da soll die Gemeinde GOttes sein, da wir meinen, daß unzählig viel Teufel sind. Wie der Ort, da Jakob schläft, bei Jerusalem ist, da die Cananiter und die Gewalt des Teufels regieret hat, daß man nichts weniger sollte gedacht haben, denn daß daselbst das Haus GOttes wäre. Und Jakob hat nur so gedacht: Wenn ich nur eine Nacht allhier sicher und unverletzt bleiben könnte, so will ich mir daran genügen lassen; morgen früh will ich aufstehen und davon laufen. Er hat nicht anders denken können, denn daß er daselbst mitten unter den Feinden schliefe, da der Fürst und GOtt dieser Welt regierete. Siehe aber, „gewißlich ist der HErr an diesem Ort“, spricht er, „und ich wußte es nicht.“ Denn ich meinte, es wäre außerhalb des Hauses meines Vaters kein GOtt, da er der Pfarrherr ist, da mir der Segen ist gegeben worden: ich sehe aber nun in Wahrheit und erfahre es auch, daß GOtt seine Verheißung wunderbarlich erfüllt, darin er gesagt hat: „Dieses Land will ich dir geben.“ Und ich verstehe nun, daß ich ein Herr werden soll dieses Landes, darin ich jetzt liege; ich hätte nicht denken oder hoffen können in solchem großen Schrecken und Traurigkeit, daß dies geschehen würde. Auf solche Weise tröstet GOtt den Jakob durch diesen Traum reichlich, und Jakob wird also wunderbarerweise gestärkt, da er zuvor durch den Traum erschreckt und bewegt worden, der seine rechte Beziehung nach der Regel gehabt und mit der That an sich selbst überein gekommen ist. Derhalben ist es ein rechter wahrhaftiger Traum.

152. Sie disputiren aber auch von dem Orte Bethel, wo derselbe möge gewesen sein, diesmal desselben in den Propheten und in der ganzen heiligen Schrift, vornehmlich aber im Hosea, zum öfteren gedacht wird; als, Hos. 12, 5.: „Desselbst hat er ihn ja zu Bethel funden“ etc. Und darnach hat ihn das Volk Israel stets gerühmt und geehrt, also auch, daß Jerobeam, der ein sehr weiser und kluger Mann war, diesen Ort dazu erwählete, daß er daselbst das eine goldene Kalb aufrichtete bei Jerusalem, wie er das andere zu Dan hat setzen lassen gegen Mitternacht, 1 Kön. 12, 28. 29. Denn diese zwei Götter oder Götzen hat der König aus klugem Rathe nach seiner Vernunft erdichtet und angerichtet, und sie an die zwei äußersten Orte des Landes gesetzt; dazu sind darnach die Könige in Samaria auch gekommen. Bethel aber ist gelegen an der äußersten Grenze des Königreichs Israel, in Ephraim; und der Tempel zu Jerusalem, so von Salomo ist gebaut worden, war zwölf Marksteine von Bethel, das ist, drei kleine Meilen. Derhalben hat Jerobeam sein listig also gedacht: Ich will an zwei Orten Abgötterei anrichten, auf daß das Volk in meinem Königreiche an beiden Orten einen gewissen Gottesdienst und Stätte haben möge; damit es also innerhalb des Landes bleiben muß.

\quellseite{120}

153. Ein solch groß Aergerniß richtet er an, und thut dasselbe an einem solchen gelegenen Ort nahe bei Jerusalem, an welchen Ort die Menge des Volkes häufig gelaufen ist, und haben beide das Wort und den Tempel GOttes fahren lassen; wiewohl es auch Viele unter den Israeliten verachtet haben. Die Propheten aber haben für und für einen Streit gehabt mit denselben Götzendienern, bis das Land gar zerstöret und verderbet worden ist. „Gehet nicht hin gen Gilgal, gen Bethel“, schrieen die Propheten, Hos. 4, 15., Amos 5, 5., „gen Beth Aven“ (denn also haben sie Bethel genannt); denn daselbst suchten sie sonderliche Gottesdienste, Opfer, Räuchwerk, Anrufung oder Gebet, nur darum, daß der Tempel und Gottesdienst zu Jerusalem möchte verachtet werden, davon GOtt gesagt hatte: „Hier will ich wohnen“ etc., 2 Chron. 7, 16.

154. Sie haben aber solche ihre Abgötterei mit einem wunderbaren Schein der Heiligkeit vertheidigt. Denn sie wußten hoch zu rühmen, daß an diesem Ort eine solche herrliche Erscheinung geschehen wäre dem Patriarchen Jakob; hier habe man die Leiter gesehen; an diesem Ort habe GOtt mit Jakob geredet und nicht zu Jerusalem. Derhalben haben die Propheten, und sonderlich Hosea, heftig dawider geredet und diese Abgötterei verdammt. Sie haben aber damit nichts ausgerichtet; wie wir auch nicht viel ausrichten können bei den Widersachern. Wenn wir die Lehre vom Glauben, daß derselbe allein gerecht mache, so oft und so klarlich den Leuten einprägen: da werfen uns die Widersacher auch vor ihr Beth Aven, ihre Kirche, das Ansehen der Väter; welches wahrlich keine geringe Anfechtung ist, dadurch viele Leute von der rechten Erkenntniß und Liebe des Evangeliums abgeführt werden.

155. Bethel aber heißt eben so viel, als wenn man sagt, das Haus GOttes, hier wohnt GOtt, GOttes Haus, Gotta. Der große herrliche Titel und der treffliche Name hat das Volk über die Maßen sehr beweget, daß sie gesagt haben: Du siehst, daß dieses in Wahrheit also geschehen ist, welches Mose selbst bezeuget; dieser Ort ist das rechte Bethel und GOttes Haus; Jerusalem ist nicht der rechte Ort zum Gottesdienste; denn hier ist GOtt erschienen und hat mit Jakob geredet. Darum soll man ihn auch allhier anbeten; wie ihn Jakob angebetet und ein Gelübde gethan hat, daß er daselbst einen Gottesdienst wolle anrichten und GOtt den Zehnten geben. Das ist ein großes Aergerniß gewesen, darüber die Propheten, und vornehmlich Hosea, sehr klagen. Wir haben im Papstthum vergleichen keines, daß einen solchen großen Schein haben möchte.

156. Es sind aber fast alle die, so dieses Buch ausgelegt haben, wider diesen Verstand, und sagen, daß dieser Ort Bethel Jerusalem und der Tempel selbst sein soll. Wie sie auch vom Berge Moria, davon droben im 22. Capitel gesagt worden, sagen, daß eben dieser Ort Bethel sein soll, da Sem und die andern Väter angebetet und geopfert haben, und da hernach auch der Tempel ist gebaut worden. Also daß vor dem Patriarchen Jakob Bethel und das Gotteshaus gewesen ist, da die Väter angebetet und geopfert haben, da Abraham seinen Sohn Isaak hat opfern wollen: und ist daran kein Zweifel, dieser Ort wird dazu verordnet gewesen sein, daß daselbst die Gemeinde zum Gottesdienste ist zusammen gekommen. Daß aber dieser Ort die Stadt Jerusalem selbst sein sollte, deß werden sie mich nicht leichtlich überreden; und so viel die Historien und Propheten anzeigen, ist dieser Ort außerhalb Jerusalems gelegen gewesen und ist von der Stadt auf zwölf Marksteine weit gelegen. Lyra führt zwei Verse an, darin acht Namen der Stadt Jerusalem enthalten sind, die lauten also:

Solyma, Luza, Bethel, Hierosolyma, Hiebus, Helia, Urbs sacra Jerusalem dicitur atque Salem.

157. Aber ich zweifle noch an den zwei Namen Luza und Bethel, und ist dieses meine Meinung von dieser Disputation. Es sind die Veränderungen der Städte wunderbar und mancherlei. Und wo die Meinung derjenigen wahr ist, die der Stadt Jerusalem diesen Namen zeigen, so will ich es zwar zugeben, daß diese Stadt eine Zeit lang Bethel sei genannt worden; aber doch also, daß sie diesen Namen nicht behalten hat, sondern daß er auf die Stadt Jerusalem sei gelegt worden, oder aber, daß sie nicht mit dem rechten Eigennamen, sondern mit dem gemeinen Namen also sei genannt worden. Denn die Beschreibung der Orte bei Josua streitet dawider, welcher Cap. 8. V. 1. 9. anzeigt, daß Ai und Bethel Städte gewesen seien, allernächst bei Jerusalem gelegen, wo anders diese Stadt Bethel bis auf die Zeit Josua geblieben ist; und läßt es sich ansehen, daß es ein kleiner Flecken gewesen sei bei der Stadt Ai; wie droben (1 Mos. 13, 3.) gesagt ist, daß sich Abraham zwischen Bethel und Ai gelagert habe.

158. Wo nun Jakob die Stadt Jerusalem hat Bethel genannt, so hat er es, sage ich, gethan, daß es ein gemeiner, und nicht ein Eigenname sein sollte. Als wollte er sagen: Hier wohnt GOtt. Sonst kann ich die Meinung, so die Andern haben, mit den Propheten nicht ausgleichen oder zusammenreimen. Die Historien im Buche der Richter und Josua stimmen damit auch überein, daß der Berg Moria ein Ort gewesen sei zu Bethel, daß es sich mit den Worten Jakobs auch reimt, da er sagt: „Hier wohnet gewißlich GOtt.“ Sonst ist es gewiß, daß die Namen der Städte auf mancherlei Weise und oftmals geändert werden.

159. Aber das lasse ich fahren. Ich wollte wohl gerne also Gedanken machen, daß dieser Ort der Berg Calvaria sein sollte. Und ob diese Gedanken nicht so gar gewiß und beständig sind, so sind sie dennoch göttselig, als an welchem Ort der HErr selbst geschlafen hat. Und man sagt, daß am selbigen Ort auch soll gestanden haben der Baum, von dessen Früchten zu essen Adam und Eva verboten war. Darum hat GOtt haben wollen, daß daselbst Christus hat gekreuzigt sollen werden und sterben; und der Ort, da Jakob die Leiter gesehen, sollte eben derselbe Ort sein, da Christus, der rechte Jakob, im Grabe geschlafen und auferstanden ist, und da die Engel herab gefahren und hinauf gefahren sind. Also sollte dieser Ort recht und wahrhaftig Bethel genannt werden. Dieses ist aber nur eine Einbildung; und wo solche Gedanken ja nicht wahr sind, so ist es doch ohne alle Gefahr, dieselbigen zu haben, und ist auch dabei gar keine Gefahr, wenn wir schon keine gewisse Namen der Orte haben. Jerusalem und Bethel haben zur Zeit Abrahams und Josuas drei Meilen von einander gelegen: ob nun die Eigennamen in die gemeinen Namen verändert sind worden oder nicht, daran ist wenig gelegen. Die acht Namen, welche Lyra aufzählt, bedeuten eben dasselbe; dieweil daselbst der HErr ist gesehen worden in einem solchen Gesichte, darin man seinen Frieden hat spüren mögen.

160. Aber dieses wollen wir den Grammatikern überlassen. Wir haben daran genug, daß wir wissen, daß Jakob allhier gestärket wird im Glauben der göttlichen Verheißung, daß er eben denselbigen HErrn und GOtt hier gesehen, eben dasselbe Wort gehört, auch dieselbe Kirche und Gemeinde im Schlafe gesehen, die er daheim gesehen hat, wiewohl er allein hier ist und sonst niemand; auf daß wir lernen, daß da die Kirche und Gemeinde GOttes sei, wo das Wort gelehrt und gehört wird, es sei gleich mitten in der Türkei, oder im Papstthum, oder auch in der Hölle. Denn GOttes Wort ist es, das die Kirche macht; das ist der HErr über alle Orte: an welchem Ort nur dasselbe gehört wird, wo die Taufe, das Sacrament des Altars und die Absolution gereicht wird, da sollst du es gewiß dafür halten, schließen und sagen: Hier ist gewißlich GOttes Haus, hier steht der Himmel offen. Wie aber das Wort an keinen Ort gebunden ist, also ist die Kirche auch nicht gebunden an irgend einen Ort. Man soll nicht sagen: Der Papst ist zu Rom, darum ist daselbst auch die Kirche; sondern wo GOtt redet, wo die Leiter Jakobs ist, wo die Engel auf- und niedersteigen, da ist die Kirche, da wird das Himmelreich aufgethan.

\begin{quote}
\textbf{V. 17. Und fürchtete sich und sprach: Wie heilig ist diese Stätte! hie ist nichts anders, denn GOttes Haus, und hie ist die Pforte des Himmels.}
\end{quote}

161. Dieses ist eine sehr wunderbare Rede, dergleichen in der ganzen heiligen Schrift nicht ist, weder im Alten noch im Neuen Testamente. Denn diese Stätte wird GOttes Haus und die Pforte des Himmels genannt, welches an diesem Ort allein und sonst an keinem Orte mehr gesagt wird. Und dieses ist Anderes nichts, denn das Himmelreich und den Himmel selbst nennen. Denn GOttes Haus ist ein solcher Ort, da GOtt wohnt. Wo wohnt aber GOtt? Ist es nicht wahr, daß er im Himmel wohnt? Derhalben setzt er die Erde zum Himmel und den Himmel zur Erde, und bindet die zwei also
\quellseite{121}
zusammen. Er setzt die himmlische Wohnung und das Reich GOttes an den Ort des Landes, und sagt: Hier ist das Himmelreich, hier ist die Pforte des Himmels, da der Weg in den Himmel geht und da man aus dem Hause GOttes geht. Denn er macht einen Unterschied zwischen Himmel und Erde, das ist, zwischen dem gegenwärtigen und zukünftigen Leben; denn er verheißt sich selbst nach diesem Leben das ewige Leben und Himmelreich, welches in diesem gegenwärtigen Leben angefangen wird. Ja, Jakob empfindet, daß er jetzt schon im Himmelreich ist und darin lebt. Derhalben haben die Propheten aus dieser Rede ohne Zweifel sehr viel treffliche Offenbarungen genommen. Da heraus sind gekommen die Verheißungen im Propheten Daniel und an andern Orten, da Jerusalem mit dem Himmelreiche verglichen wird. Denn also steht Dan. 8, 10.: „Und es wuchs bis an des Himmels Heer, und warf etliche davon, und von den Sternen zur Erde, und zertrat sie.“ Daselbst wird nichts Anderes bedeutet, denn die Kirche, und die Lehrer und Zuhörer darin.

162. Darnach wird allhier die rechte Kirche beschrieben, was sie sei und wo sie sei. Denn wo GOtt wohnt, da ist die Kirche und sonst nirgend. Denn die Kirche ist GOttes Haus und die Pforte des Himmels, da der Eingang ist zum ewigen Leben, und da man aus diesem zeitlichen Leben in das himmlische Leben sieht; wo aber GOtt nicht geredet oder gewohnt hat, daselbst ist auch nie eine Kirche gewesen.

163. Wir haben heutiges Tages einen großen Streit mit den sehr verdammten Leuten, den Papisten, über die Kirche, welchen Namen und Titel sie sich mit großem Trotze zuschreiben, und rühmen sich dessen hoch, daß bei ihnen die Kirche sei, und dasselbe in ihrer Lehre und Ceremonien. Wiewohl wir aber nun bekennen, daß bei ihnen eine Kirche sei; denn sie haben ja die Taufe, die Absolution, den Text des Evangeliums, und sind auch viel gottselige Leute unter ihnen: jedoch, wenn sie das wollen dazu gesetzt haben, daß der Pabst und ihre großen Kirchengepränge die Kirche sein sollen, das werden wir ihnen keinesweges zugeben, und halten ihnen diesen Text vor, der wider sie ist. Das Haus GOttes und die Kirche GOttes ist ein Ding, wie Christus Joh. 14, 23. sagt: „Wer mich liebet, der wird mein Wort halten, und mein Vater wird ihn lieben, und wir werden zu ihm kommen und Wohnung bei ihm machen.“ Da muß GOttes Wort sein und das muß man lieben; denn GOtt macht keine Wohnung bei uns, es sei denn, daß wir sein Wort haben. Wenn die Kirche soll GOttes Haus sein, so muß sie auch GOttes Wort haben, und daß in diesem Hause GOtt allein der Hausvater sei.

164. Wir nehmen aber diesen Zusatz der Papisten nicht an, die da schreien: Ihr habt von uns die Schlüssel, das Sacrament des Altars, den Beruf und das Amt, welches ja die rechten Zeichen sind, dabei man die rechte Kirche erkennt. Es mangelt aber noch an Einem Stück, sagen sie, nämlich, daß ihr den Pabst und seine Cardinal[e] nicht für die Kirche halten und erkennen wollt. Darüber ist nun der Streit zwischen uns und den Papisten. Wir wollen den Feind neben dem Hausvater nicht zulassen; denn Christus und der Satan können nicht zugleich Ein Haus regieren. Denn „wie stimmt Christus mit Belial?“ 2 Cor. 6, 15. Wenn wir nun den Pabst ehren sollen, so müssen wir Christum verleugnen; denn er leidet nicht, daß der Satan in seinem Hause regieren soll. Wollen wir eine Kirche haben, so muß sie nicht befleckt und vermischt werden mit irgend einer Teufelslehre. Die papistischen Menschensatzungen aber sind gottlos und stracks wider GOttes Wort; um welches Wortes willen GOtt allein will Wohnung bei uns machen.

165. Darum soll man solche Menschensatzungen fliehen und sich davor hüten; denn Belial und Christus können nicht zugleich in Einem Hause Herren und Hausvater sein. Darum soll man diesen Text und andere dergleichen mehr fleißig merken, daß wir daraus lernen, daß der Pabst und sein Anhang die Kirche nicht sind. Wenn man uns aber vorwirft: Haben sie doch die Taufe, das Abendmahl des HErrn etc., darauf antworte ich also: Die das reine Wort und die Taufe haben, gehören zu uns und zu der rechten Kirche; die aber das Kirchengepränge menschlicher Satzungen dabei haben, sind nicht die Kirche. Ob sie wohl die Taufe und den Text des Evangeliums haben, so haben sie es doch vergeblich. Denn Christus stimmt nicht mit Belial überein. Denn das Bett ist zu enge, also daß der Eine hinaus fällt; und der Mantel ist zu kurz, er kann sie beide nicht bedecken, wie Jes. 28, 20. steht: „Das Bette ist so enge, daß nichts übrig ist; und die Decke so kurz, daß man sich drein schmiegen muß“ etc.

166. Nun gehen sie aber damit einzig und allein um, daß sie ihre Greuel nebeneinführen und vertheidigen wollen. Darum soll man ihnen allezeit das wiederum vorhalten, so im ersten Buch der Könige steht am 18. Cap. V. 21., da Elia spricht: „Wie lange hinket ihr auf beiden Seiten? Ist der HErr GOtt, so wandelt ihm nach; ists aber Baal, so wandelt ihm nach.“ Der HErr und Baal stimmen nicht überein, welche doch Ahab zu vereinigen befohlen war: er wollte GOtt mit dem reinen Worte dienen und wollte doch gleichwohl auch die Menschensatzungen mit untermenge[n]. Nun können aber die Dinge, so stracks wider einander sind, nicht zugleich bestehen; sondern es hebt immer Eins das Andere auf.

167. Derhalben sagt nun Jacob also: Diese Stätte, da ich schlafe, ist das Haus und die Kirche GOttes. Allhier hat GOtt selbst eine Kanzel aufgerichtet, und er ist der Erste, der da predigt von den Nachkommen, und daß diese Kirche für und für währen und bestehen soll; Jacob aber, mit den Nachkommen in seinen Lenden, ist der Zuhörer und die Engel im Himmel. Denn wo auch nur Eine Person ist, die das Wort hört, sammt den Engeln, die zugleich mit dabei sind, daran ist es genug. Er beschreibt aber nun die Herrlichkeit dieser Kirche sehr prächtig, nämlich, daß daselbst die Pforte und der Eingang sei zum Himmelreich. Denn also führt und regiert uns GOtt, daß allenthalben hier auf Erden, wo er mit uns redet, die Thür zum Himmelreich offen steht. Dies ist wahrlich ein großer trefflicher Trost; wo wir nur das Wort hören und getauft werden, da gehen wir in das ewige Leben.

168. Wo findet man aber denselbigen Ort? Antwort: Hier auf Erden, da die Leiter steht, die an den Himmel rührt, da die Engel herab fahren und wieder hinauf steigen, da Jacob schläft. Es ist eine leibliche Stätte; aber man steigt daselbst hinauf in den Himmel ohne irdische Leiter, ohne Flügel und Federn. Der Glaube sagt also: Ich gehe hin an den Ort, da das Wort gelehrt, da das Sacrament und die Taufe gereicht wird; und alles, was da geschieht vor meinen Augen, an dem leiblichen Ort, findet eitel himmlische, göttliche Worte und Werke; die Stätte ist nicht allein Erde oder Land, sondern ist etwas Größeres und Herrlicheres, nämlich, das Reich GOttes und die Pforte des Himmels. Hic itur ad astra: Hier geht man zu den Sternen, wie der Poet sagt. Du darfst nicht zu St. Jacob laufen, oder irgend in einen Winkel kriechen, oder sonst in ein Kloster ziehen. Hüte dich und suche keinen neuen und närrischen Eingang zum Himmel, sondern stehe mit dem Glauben auf den Ort, da das Wort und die Sacramente sind; dahin richte deinen Gang, wo das Wort schallt und die Sacramente verwaltet werden, und schreibe daselbst den Titel hin: Die Pforte GOttes oder des Himmels. Solches geschehe entweder in der Kirche und da die Gemeinde zusammen kommt, oder aber in der Kammer, wenn wir die Kranken trösten und aufrichten; oder wenn wir den, so mit uns über Tische sitzt, absolviren: da ist die Pforte des Himmels. Wie Christus sagt Matth. 18, 20.: „Wo zween oder drei versammelt sind in meinem Namen, da bin ich mitten unter ihnen.“ Da ist das Haus GOttes und die Pforte des Himmels durch die ganze Welt, wo nur das Wort sammt den Sacramenten lauter und rein gelehret wird.

169. Wir sollen aber die äußerliche Stätte nicht allein mit leiblichen Augen, und wie die unvernünftigen Thiere pflegen, ansehen; wir sollen auch nicht meinen, daß das Wort nur ein vergeblicher Laut sei, der nur äußerlich in die Ohren klinge. Es ist zwar eine menschliche Stimme, und der das Wort redet, ist auch ein Mensch; die Kirche ist von Steinen und Holz gebaut, und unsere Kirche, wenn man darin nicht zusammenkommt, ist kein Tempel oder Kirche GOttes; es sei denn, daß sie beziehungsweise (relativ) also genannt werde. Wenn man aber darin predigt, die Sacramente reicht, die Diener zum Lehramte ordinirt und bestätigt, alsdann sollst du sagen: Da ist GOttes Haus und die Pforte des Himmels; denn GOtt redet daselbst, wie St. Petrus sagt 1. Ep. 4, 11.: „So jemand redet, daß ers rede als GOttes Wort. So jemand ein Amt hat, daß ers thue als aus dem Vermögen, das GOtt darreicht.“

170. Aber die Blindheit und Unverstand unserer Herzen liegt vor unsern Augen wie eine finstere Wolke, daß wir die große Herrlichkeit nicht sehen können. Es ist wahrlich eine große Herrlichkeit und Majestät, wenn man sagt: Das ist Gottes Wort. Ich höre die Stimme eines Menschen; ich sehe menschliche Gebärden; Brod und Wein im Abendmahl sind leibliche Dinge; in der Ordination der Kirchendiener legen die Menschen, so Fleisch und Blut sind, den Dienern die Hände auf; in der Taufe ist das Wasser Wasser: denn also und anders nicht urtheilt von diesen Dingen allen das Fleisch. Wenn du aber mit geistlichen Augen diesen Zusatz wirst ansehen, nämlich, wes Wort es sei, das daselbst geredet und gehört wird; zwar keines Menschen (denn wenn es eines Menschen Wort ist, so redet es der Teufel), sondern Gottes Wort: da wirst du verstehen, daß es Gottes Haus und die Pforte des Himmels sei. Das hölzerne Haus oder das Land an sich selbst verdient diesen Namen billig nicht; aber das Land oder Stätte, da man Gottes Wort klingen hört, ist Gottes Land, nicht des Schöpfers, sondern des, der seine Kirche regieret, der seine Heiligen auf solche Weise führt, daß sie in das Himmelreich gehen mögen. Also, das Wasser, so in der Taufe geschöpft und über das Kind gegossen wird, ist auch Wasser, nicht des Schöpfers, sondern Gottes des Heilandes.
\quellseite{122}

171. „Ich sehe aber ein ander Gesetz in meinen Gliedern, das da widerstreitet dem Glauben“, Röm. 7, 23. Denn ich weiß wohl, daß solches wahr ist, und kann demselben doch nicht so gewißlich und mit vollem Vertrauen beifallen: ich kann so viel und so stark nicht glauben, als ich gerne wollte. Also ist es wahrlich um mich bestellt. Ja, es können es auch die Engel selbst nicht fassen und genugsam durchschauen, welch ein groß Ding das sei, davon hier gehandelt und gesagt wird, also daß sie des Anschauens und der Betrachtung nicht satt werden können. Ich wünsche aber, daß ich es so viel ergreifen könnte, als ich jetzt gern wollte. Dasselbe kann aber auch nicht geschehen; denn Fleisch und Blut, dazu die schwere Last, und die Sünde, so uns umfangen hat, streitet immer dawider; von welcher Sünde in der Epistel an die Hebräer steht am 12. Capitel V. 1., daß sie uns immer anklebt und träge macht. Sie unterdrückt den Geist, der dieses gern fassen und vollkommenlich schöpfen wollte, daß er davon möchte trunken werden; aber er muß sich an einem Wenigen genügen lassen, nur daß die Zunge erquickt und gekühlt werden möge.

172. Unterdeß aber müssen wir uns daran gewöhnen, daß wir diesen Zusatz machen können, welchen an diesem Ort Jakob und die Engel im Himmel machen; denn das Fleisch richtet die Augen nur auf das Wasser, Brod und Wein, und auf das Land, da Jakob geschlafen hat. Aber der Geist soll das Wasser, die Hand und das Wort Gottes sehen, und Gott im Wasser. Das Fleisch sieht so scharf, daß es urtheilen kann, Wasser sei Wasser, und daß es Gott ausfließet; wie die Sacramentenschwärmer und Wiedertäufer thun. Derhalben soll man lernen, wider das Anschauen des Fleisches, daß es kein schlecht Wort und nur ein vergeblicher Laut sei, der nur in die Ohren klinge, sondern daß es ein Wort des Schöpfers Himmels und der Erden sei. Also ist das Auslegen der Hände keine Menschenfassung, sondern Gott macht und ordinirt die Diener; und der Pfarrherr ist es nicht, der dich absolvirt, sondern der Mund und die Hand des Dieners ist der Mund und die Hand Gottes.

173. Derhalben sollen wir diese unermeßliche Ehre und Herrlichkeit Gottes erkennen und hoch achten, dadurch er sich uns in seiner Kirche offenbart. Denn die Kirche ist nicht ein solches Haus, darin er schafft, wie er anfänglich Alles aus Nichts geschaffen hat; sondern darin er mit uns redet, bei uns wohnt, uns nährt und versorgt, wir schlafen oder wachen. Wie viel sind aber derer wohl, die das glauben? Es ist ein wahrhaftiges Ding, und ein sehr großes Wort, daran die Engel im Himmel ihre Lust haben, und wenn sie es hören, werden sie darüber freudig und fröhlich, um dieser Erkenntniß willen der Kirche, und können doch nicht satt werden, solches anzuschauen, zu betrachten und sich desselben zu verwundern. Denn die Kirche begreift Gott in sich, der mit uns wandelt und bei uns wohnt; also daß er uns lebendig macht, bewahrt und heilt, und durch solche Beiwohnung wirkt er, daß in diesem Leben das Haus Gottes und die Pforte des Himmels sei.

174. Dieses sind wahrlich unermeßliche und wunderbare Dinge; darum sollen wir lernen, die Ehre und Herrlichkeit der Kirche nach ihrer Würde groß machen und viel davon halten.

Wir sollen sie nicht gering achten, wir sollen auch die rechte Würdigkeit der Kirche nicht in solche schimpfliche närrische Dinge setzen, die nichts werth sind; wie die Kirchenräuber, die Papisten, zu thun pflegen, die sich träumen lassen, daß die Kirche eine leibliche Versammlung sei, die im Trank, Speise und Kleidung etc. bestehe; wiewohl unsere Kirche auch viel Leibliches und Aeußerliches hat, als da ist, Brod, Wasser, Erde etc. Es ist aber ein großer Unterschied zwischen den Dingen, so die Papisten rühmen, und den leiblichen Dingen, die wir haben. Denn jene sind wider Gott angenommen und gebraucht worden, und nicht Gott zu Ehren oder nach Gottes Wort. Darum ist ihre Kirche nur eine vergebliche Heuchelei und daß sie sich etwas davon träumen lassen. Sie ist ein falscher Schein. Dagegen aber ist es ein sehr großes Ding, daß Jakob so eigentlich anzeigt, daß das Gottes Haus sei, wo Gott bei uns wohnt, da wir das Hausgesinde, ja, die Söhne und Töchter sind, und er unser Vater ist, der mit uns redet, mit uns umgeht, und dasselbe auf das allvertrauteste wirkt und ausrichtet, daß die Kirche die Pforte des Himmels sei. Denn darum wohnt er bei uns, daß wir in das Himmelreich eingehen sollen, und, was das Allerlieblichste ist, er ist zuerst gekommen und uns auf der Leiter erschienen, steigt hernieder und hat Gemeinschaft mit uns, er redet und wirkt in uns.

175. Also wird nun die Kirche unter den Leuten gemacht und aufgerichtet, wenn Gott bei den Menschen wohnt, zu dem Ende, daß eine Pforte des Himmels sei und daß wir aus diesem irdischen Leben in das ewige und himmlische Leben gehen mögen. Wer kann sich nun dessen genugsam verwundern oder solches begreifen, daß Gott bei den Menschen wohnt? Ja, dies ist das neue Jerusalem, so von Gott aus dem Himmel herab fährt und die Klarheit Gottes hat etc., wie in der Offenbarung Johannis am 21. Cap. V. 2. steht.

176. Und dieses ist die rechte eigentliche Beschreibung der Kirche, was sie sei nach ihrem Wesen, nämlich, daß man sage: Die Kirche ist ein solcher Ort oder Volk, da Gott wohnt, darum daß er machen will, daß wir in das Himmelreich gehen mögen; denn es ist eine Pforte des Himmels. Und daraus folget sehr fein, daß derhalben in der Kirche nichts soll gehört oder gesehen werden, denn allein was Gott thut; wie St. Petrus sagt 1. Epist. 4, V. 11.: „So jemand redet, daß ers rede als Gottes Wort. So jemand ein Amt hat, daß ers thue als aus dem Vermögen, das Gott darreicht.“ Wo ich aber des Wortes oder der göttlichen Kraft ungewiß bin, so soll ich stille schweigen. So oft ich aber mein Amt verwalte, das ist, taufe oder absolviere, soll ich dessen gewiß sein, daß mein Werk nicht mein, sondern Gottes ist, der durch mich wirkt. Die Taufe ist eine Wirkung Gottes; denn sie ist nicht mein, wiewohl ich die Hände und den Mund dazu leihe als Werkzeuge. Also, wenn ich absolviere, oder zum Kirchenamt berufe und dir die Hände auflege, sollt du daran nicht zweifeln, daß es aus dem Vermögen geschehe, so Gott gibt; wie der Spruch Petri lautet. Und so ist dies die richtige Beschreibung der Kirche, daß sie eine Wohnung Gottes auf Erden ist; nicht daß wir auf Erden bleiben sollen, sondern darum werden die Sacramente gereicht, darum wird das Wort gepredigt, daß wir in das Himmelreich geführt werden und durch die Kirche in den Himmel gehen.

177. Das hat Jakob gesehen und seine Nachkommen haben es gesehen, und wir sehen es auch, und Alle, die jetzt die Kirche sind, oder nach uns die Kirche werden sollen, nämlich, daß sie Gottes Haus ist, die uns von der Erde in den Himmel führt. Und die Kirche hat ihre Stätte im Tempel, in der Schule, im Hause, in der Schlafkammer. Wo zwei oder drei im Namen Christi zusammenkommen, da wohnt Gott, Matth. 18, 20., ja, wenn jemand mit sich selber redet und Gottes Wort betrachtet, da ist Gott mit den Engeln dabei, und wirkt und redet also, daß daselbst die Thür offen steht zum Himmelreich.

178. Dagegen aber wohnt Gott in den Menschensatzungen nicht; denn er redet daselbst nicht. Denn man muß allezeit darauf sehen, daß Jakob Gott reden und wirken hört. Das sind die rechten Werke, so in die Kirche gehören und die uns in das Himmelreich führen. Die päpstliche Kirche aber ist das Himmelreich nicht; denn sie trachtet darnach, wie sie nur Gold, Silber, Königreiche und die Kronen dieser Welt haben möge. Das sind ihre Schlüssel. Was du auf Erden binden willst, das ist, wo du die Könige absetzen und die Königreiche mit Gewalt einnehmen wirst etc. Da ist der rechte Weg zur Hölle. Derhalben lasset uns mit dem Patriarchen Jakob hören nicht die Stimme des Psalmes oder der Menschen, sondern GOttes Wort. Denn wenn der redet, alsdann steht der Weg und die Pforte des Himmels offen.
\quellseite{123}

\section[Sechster Theil]{Sechster Theil. Wie Jakob einen Stein aufrichtet, denselben salbt und dabei ein Gelübde thut}

\subsection{I}

\begin{quote}
\textbf{V. 18. 19. Und Jakob stund des Morgens frühe auf, und nahm den Stein, den er zu seinem Häupten gelegt hatte, und richtete ihn auf zu einem Mal, und goß Oel oben drauf, und hieß die Stätte Bethel; vorher hieß sonst die Stadt Lus.}
\end{quote}

179. Das hebräische Wort mazebah, auf Lateinisch in statuam, sive erectionem, das ist: Er richtete den Stein auf, ist eben dasselbe Wort, so im 2. Psalm V. 2. auch steht, da man im Lateinischen Text liest: Adstiterunt reges terrae: „Die Könige im Lande lehnen sich auf.“ Er hat den Stein aufgerichtet, daß er stehen sollte und nicht liegen, wie ein Altar oder Säule pflegt zu stehen. Ob aber Jakob in der Stadt Lus, oder aber außerhalb der Stadt in Felde geschlafen habe, kann ich nicht gewiß sagen. Etliche wollen haben, er solle auf dem Felde geschlafen haben, die Andern wollen, er soll in der Stadt Lus gelegen haben; wie es sich ansehen lässet, daß fast der Text also lautet.

180. Bethel aber heißt die Stätte, so droben Morija ist genannt worden, und hat vom Anfange diesen Namen gehabt, daß sie GOttes Haus ist genannt worden. Darum ist es wohl glaublich, daß sie zum öftern vom Teufel zerstört und auch oft wiederum gebaut sei. Denn zwischen GOtt und dem Teufel ist immer ein stetiger Krieg: jetzt stürzet GOtt den Teufel und treibt ihn weg, dann kommt der Teufel wieder, und wenn die Lehre verfälscht ist, setzt er an die Statt der Kirche eine Mördergrube, und besudelt die Stätte, da der Tempel GOttes ist. Zur Zeit Davids hat Arauna, der Jebusiter, diese Stätte inne gehabt, und hatte daselbst eine Scheune oder Tenne gebaut, da er den Weizen gemahlen hat; welches also durch des Teufels Treiben geschehen ist aus Haß und Verachtung dieser Stätte, wie er dieselbe zuvor auch öfters verunreinigt hatte. David aber hat dieselbige Stätte gekauft nach dem großen Sterben, so unter dem Volke gewesen war, und hat daselbst einen Tempel oder Altar gebaut und dem HErrn geopfert, darum daß er war erhöret worden; wie im 2. Buch Samuels steht am 24. Cap. V. 21. Denn ich halte dafür, es sei die Stätte gewesen, da Jakob geschlafen hat: und GOtt hat dieselbige Stätte sonderlich lieb gehabt um der Väter willen, um Sems willen, der daselbst gepredigt und Wunderzeichen gethan hat, und auch um der Andern willen vor und nach Sem, als Abraham, Isaaks, Jakobs und um Christi willen selbst.

181. Der Teufel aber ist ein stolzer hoffärtiger Geist, darum nimmt er die heiligen und reinen Stätten gern ein, wie diese Stätte gewesen ist. Wie die Worte Jakobs auch bezeugen, da er sagt: Das muß eine heilige Stätte sein. Als wollte er sagen: Morija ist eine heilige Stätte um GOttes willen, der da herrlich und heilig ist, der an diesem Orte seinen Gottesdienst gehabt und noch hat. Und um der Ursache willen hat er diese Stätte Bethel geheißen, welche vorhin Morija ist genannt worden, darum daß sich GOtt daselbst geoffenbart hat.

182. Mittlerzeit aber, dieweil Jakob in Egypten wohnt mit seinem Hausgesinde, und auch hernach die ganze Zeit der Richter haben die Heiden und der Teufel diese Stätte inne gehabt, bis sie David wiederum recht aufgerichtet und Salomo den Tempel daselbst gebaut. Wiewohl es hat geschehen mögen, daß eben diese Stätte, darauf Jakob gelegen und geschlafen hat, ganz und gar verwüstet und verstört worden ist; sintemal Hosea und das Buch der Richter öffentlich anzeigen, daß die Stadt Bethel von Jerusalem zwölf Marksteine gelegen gewesen. Wo nun Morija die Stätte gewesen ist, da Jakob geschlafen hat, so hat sie den Namen verloren, und ist hernach eine andere Stätte gewesen, die Bethel ist genannt worden. Wiewohl ich das Wort Bethel also verstehe, daß es ein Gattungswort sei, auf Deutsch: Hier wohnt GOtt.

183. Woher hat er aber das Oel genommen, so er oben auf den Stein gegossen hat? Antwort: Wo daselbst eine Stätte gewesen ist, wie der Text anzeigt, so hat er leichtlich Leute finden können, von denen er es gekauft hat. Er hat aber damit, daß er das Oel darauf gegossen, nichts Anderes anzeigen wollen, denn daß er diese Stätte dazu geeignet und geweihet hat, daß es hinfort ein Haus und Wohnung GOttes sein sollte, darum daß daselbst der HErr geredet hat und die Engel seine Zuhörer gewesen sind. Dieweil aber keine Leute dabei gewesen, die dasselbe auch hören und erkennen möchten, so hat er gedacht, er wollte, wenn er wieder aus Mesopotamien kommen würde, eine Kirche oder Altar daselbst aufrichten, und diese herrliche und treffliche Erscheinung GOttes rühmen und preisen. Das Oel bedeutet das Salben und Weihen, damit man die Könige und Priester gesalbt, desgleichen auch die Gefäße und Kleider, auf daß sie dadurch vom weltlichen Gebrauch abgesondert und zum heiligen Gebrauch geweiht würden. Daher ist es gekommen, daß unsere Affen im Papstthum auch Oel gebraucht haben, aber ohne GOttes Wort. Jakob aber hat von sich selber diese Stätte nicht erwählet, sondern ist ohngefehr dahin gekommen und hat daselbst GOttes Wort gehört.

184. Derhalben ist die Stätte von GOtt selbst geweihet worden, dem es also wohlgefallen hat, daß er daselbst hat wohnen wollen. Darum ruft es nun Jakob aus, daß diese Stätte heilig sei. Und auf daß er den Nachkommen sein Zeugniß möchte hinterlassen, so gießt er Oel oben darauf, und verheißt, er wolle die Stätte zum Hause GOttes weihen, wenn er seine Pilgerschaft glücklich beendet hätte. Er ist dazu nicht aus eigenem Eingeben, auch nicht durch Menschenanfahung oder Rath bewogen worden, sondern weil er sie daselbst aufgerichtet, das ist, GOtt gedient hat an der heiligen Stätte um des Wortes willen des HErrn, so er daselbst gehört hatte. Die Allegorie oder heimliche Deutung, von der geistlichen Weihe, will ich bei Seite lassen.

\subsection{II}

\begin{quote}
\textbf{V. 20—22. Und Jakob that ein Gelübde, und sprach: So GOtt wird mit mir sein, und mich behüten auf dem Wege, den ich reise, und Brod zu essen geben, und Kleider anzuziehen, und mich mit Frieden wieder heim zu meinem Vater bringen; so soll der HErr mein GOtt sein; und dieser Stein, den ich aufgerichtet habe zu einem Mal, soll ein Gotteshaus werden; und alles, was du mir gibst, deß will ich dir den Zehnten geben.}
\end{quote}

185. Das letzte Stück in diesem Capitel ist von dem Gelübde, so Jacob gethan hat, davon Lyra in die Länge disputirt. Zum Allerernsten aber wird gefragt: Warum der sehr heilige Patriarch also zweifelt, oder ja zum wenigsten also redet, wie Einer pflegt zu reden, der da zweifelt: „So mir GOtt wird Brod zu essen geben; wo er mich mit Frieden wieder heim zu meinem Vater bringen wird“ etc. Denn hätte er es nicht sollen ganz und gar für gewiß und wahr halten, da er die Verheißung empfangen hatte, daß er dieses alles haben und ihm an der Dinge keinem je etwas mangeln würde? Denn GOtt, der ihm Kinder und Nachkommen verheißen hat, kann ja nicht lügen. Wo er nun ein Vater werden soll vieler Völker, so muß er ja leben; denn die gestorben sind, können nicht Eltern sein oder werden. Soll er nun leben, so muß er auch Nahrung und andere Nothdurft dieses Lebens haben. Warum redet er aber nun also, gleich als ob er an diesem allem gezweifelt hätte? Ist er doch über das durch die herrliche Erscheinung gestärkt und bestätigt worden, da GOtt die Verheißung wiederholt hat und gesagt: „Ich will dich nicht lassen, bis daß ich thue alles, das ich geredet habe.“ Woher kommt denn dieser große Zweifel nach solchen gewissen Verheißungen GOttes?

186. Ich will das bei Seite lassen, was zu hoch ist und hier angeführt werden könnte, und mich bei dem bleiben, das niedrig ist. Denn ich habe droben und sonst an anderen Orten oftmals gesagt, daß ich mir sehr wohl gefallen lasse und daß uns auch nütze und heilsam sei, daß wir hören die Schwachheit der Heiligen; denn solche Exempel von der Schwachheit der Heiligen sind uns nöthiger und bringen mehr Trost, als die Exempel der großen vortrefflichen Stärke und anderer Tugenden, so die Heiligen gehabt haben; wie ich dadurch, daß David den Goliath, den Bären und Löwen erwürgt hat, nicht viel gebessert werden kann etc.

187. Denn in solchen ritterlichen Thaten kann ich ihm nicht nachfolgen; sintemal solches alle meine Kräfte und Gedanken übertrifft. Denn durch solche Exempel der großen Thaten werden die Heiligen gerühmt von wegen ihrer Kraft und großen Stärke, so sie als tapfere Helden gehabt haben; dieselben geben aber uns nichts an; denn sie sind zu hoch, als daß wir sie erreichen und ihnen darin nachfolgen könnten. Wenn uns aber die Exempel der Schwachheit, der Sünden, des Schreckens und Anfechtung, so die Heiligen gehabt, vorgehalten werden, als, wenn ich lese die Klagen, das Seufzen, Schrecken und Zagen, so David gehabt, dasselbe richtet mich über die Maßen sehr auf und gibt mir einen großen Trost. Denn da sehe ich, wie sie in ihrem Zagen und Schrecken nicht verdorben oder umgekommen sind, sondern wie sie sich aufgerichtet und getröstet haben mit den Verheißungen, die sie empfangen hatten: daraus schließe ich, daß ich auch nicht verzagen solle. Denn wo sie im Kampf stehen mit der Hölle und im Gewissen auch erschrocken sind und damit kämpfen müssen, so sind sie also gesinnet und reden auch also, als wenn sie gar keine Verheißung hätten; und dennoch werden sie endlich durch das Wort erhalten, daß sie nicht gar dahinfallen.
\quellseite{124}

188. Denn siehe Paulus an, der 1 Cor. 2, 3. von sich selber sagt: „Ich war bei euch mit Schwachheit, und mit Furcht, und mit großem Zittern.“ Item 2 Cor. 7, 5.: „Auswendig Streit, inwendig Furcht.“ Wie redest du also, Paulus? Gesinnet das einem solchen auserwählten Rüstzeuge, der die Verheißung hat, daß er den Namen Christi tragen soll vor den Heiden? Wo gehst du nun hin, lieber Paulus, in den Kerker der Hölle, des Schreckens und der Verzweiflung? Wo sollen wir bleiben, wenn du zweifeln und schier verzagen willst an deinem gewissen Beruf? Also aber muß es zugehen auch mit den größten Heiligen. Denn die göttlichen Verheißungen werden nicht also gegeben, daß sie uns sollten sicher machen, sondern also, wie St. Paulus 2 Cor. 12, 7. sagt: „Mir ist gegeben ein Pfahl ins Fleisch, nämlich, des Satans Engel, der mich mit Fäusten schlage.“ Ei, warum das? Antwort: „Auf daß ich mich nicht überhebe“ der großen Gaben, Gnade und Barmherzigkeit Gottes.

189. Darum schickt Gott den Seinen solchen Kampf, und Anfechtung und Streit zu, daß wir von Tage zu Tage die Verheißungen Gottes immer klarer und gewisser verstehen und behalten mögen; welches nicht geschehen könnte, wenn die Heiligen allezeit solche ritterliche Stärke üben würden; ja, sie würden endlich sicher und die Verheißungen und alle Hoffnung verlieren. Darum müssen sie also versucht werden, auf daß sie den Glauben und die Hoffnung behalten und auf die Verheißung warten. Und daraus werden wir gebessert und getröstet, wenn wir sehen, daß die Patriarchen und Propheten uns gleich gewesen sind, daß sie mit Schwachheit und Zweifel angefochten worden sind, und dahin geraten, daß sie schier verzagt und den Glauben hätten verlieren sollen.

190. Was kann uns Nützlicheres und Tröstlicheres vorgehalten werden als das Exempel Petri; der ging auf dem Wasser, daß er zu Jesu käme, Matth. 14, 29. Und zuerst, da er aus dem Schiffe trat, ging er zwar auf dem Wasser, daß er zu Jesu käme, wie der Evangelist sagt; er lief mit einem großen Sturme, mit einem starken Muth und gewissen Geiste, dieweil er wußte, daß Christus da war, und hatte das Wort und die Verheißung im Worte auf seine Bitte, da er zum Herrn gesagt hatte: „Bist du es, so heiß mich zu dir kommen auf dem Wasser“, Matth. 14, 28. Aber bald, da sich ein kleiner Wind erhoben hatte, entfiel ihm der Muth, daß er anhub zu sinken. Wie nun? Wo ist der große Muth und Geist geblieben? Warum zweifelst du, lieber Petrus? Aber das hat dem Herrn Christo also wohlgefallen, daß Petrus auf diese Weise sollte versucht werden; denn wo er nicht wäre versucht worden, so wäre er stolz und aufgeblasen geworden. Nun ist es aber besser, versucht werden, denn stolz und aufgeblasen werden; denn also behält man die Verheißung und also lernen wir verstehen das Seufzen der Heiligen. Als, im 6. Psalm V. 1.: „Ach, Herr, strafe mich nicht in deinem Zorn“ etc. Denn David mag auch ein solcher großer Mann, daß ihm Gott das Zeugniß gab und sprach 1 Sam. 13, 14.: „Ich habe gefunden David, den Sohn Jesse, einen Mann nach meinem Herzen, der soll thun allen meinen Willen“; und dennoch betet er also und kämpft mit den schweren Anfechtungen des Unglaubens und der Verzweiflung.

191. Also sind wir auch berufen, und haben solche Verheißungen, die viel klarer und herrlicher sind, denn die Väter gehabt haben. Wie St. Petrus diese heilsame Gnade rühmt, da er sagt 2. Epist. 1. V. 19.: „Wir haben ein festes prophetisches Wort, und ihr thut wohl, daß ihr darauf achtet, als auf ein Licht, das da scheinet in einem dunkeln Ort.“ Die Gnade und das ewige Leben ist uns viel herrlicher verheißen und angeboten als ihnen. Denn der Sohn ist gekommen und alle Verheißungen sind erfüllt. Wir hören den Sohn selbst, wir haben die Sacramente, die Absolution, und das Evangelium verkündigt uns Tag und Nacht und sagt: Du bist heilig, dir sind deine Sünden vergeben, du bist selig etc. Was thun wir aber? Wir zittern noch immer und bleiben in unserer Schwachheit liegen, dieweil dieses Leben währt. Warum werden wir aber nicht erweckt durch das Exempel der Patriarchen, die da ganz vollkommen geglaubt haben? Antwort: Sie sind auch schwach gewesen, wie wir sind; wiewohl wir reichere Verheißungen haben denn sie. Aber also geht es, wie das Wort Gottes lautet, so zu Paulo gesagt worden ist, 2 Cor. 12, 9.: „Meine Kraft ist in den Schwachen mächtig.“ Denn sonst könnte Gott seine Verheißung in uns nicht erhalten und erfüllen, wo er in uns das grobe, hoffärtige und sichere Fleisch nicht tödtete.

192. Eben zu dem Ende ist uns droben (1 Mos. 27, 45. Cap. 28, 1. 2.) das Exempel Isaaks und Rebekkas vorgehalten worden, die sich auch gefürchtet haben und um Furcht willen ihren Sohn in das Elend haben schicken müssen, welches doch wider die Verheißung und den Segen gewesen, so ihnen gegeben war: sie haben aber Gott nicht versuchen und den Sohn nicht ohne Ursache in Gefahr setzen dürfen, obwohl die Verheißungen fest und gewiß waren. Denn also bin ich auch durch das Wort absolviert worden, habe das Sacrament des Altars gebraucht. Soll ich denn nun darob sagen: Ich will nicht arbeiten, ich will gar müßig sitzen; wo ich leben soll, so werde ich doch wohl leben? etc. Das hieße Gott versuchen; ja, du sollst der Mittel und Güter gebrauchen, so Gott dir durch seine Gnade verliehen hat. Du sollst regieren, arbeiten, und darnach trachten, daß du deine Nahrung haben mögest, damit du Gott nicht versuchest. Du sollst dich nicht in Gefahr begeben, sollst nicht in die Elbe gehen, und die Brücke liegen lassen und verachten, darum daß du die Verheißung hast, daß dich Gott nicht lassen will. Denn also haben Isaac und Rebekka den Herrn auch nicht versuchen wollen, sondern da sie eine gewisse Verheißung haben, folgen sie gleichwohl dem Rath, so ihnen GOTT eingegeben hat; denn GOTT erfüllet seine Verheißung durch gewisse Mittel. Man soll mit den losen Schwärmern (die es dafür halten, daß Alles zuvor versehen sei) nicht sagen: Wo ich versehen bin, so werde ich selig werden; wo nicht, so muß ich verdammt werden. Soll ich sterben, so wird es mich nicht helfen, wenn ich schon meines Leibes warte und auf mein Leben gute Achtung gebe; soll ich gelehrt werden, so kann ich es auch wohl ohne Bücher werden etc. Es hat aber GOTT seine Verheißung nicht in solcher Ordnung gegeben. So nämlich heißt es Röm. 8.: „Welche er aber verordnet hat, die hat er auch berufen; welche er aber berufen hat, die hat er auch gerecht gemacht“ etc. Er will das Ende seiner Verheißung nicht ohne Mittel erfüllen, sondern er will es durch Mittel thun; er hat uns die Creaturen gegeben, derselbigen sollen die Christen gebrauchen, bis wir zum Ende der göttlichen Verheißung kommen; wie Jacob, der viele Verheißungen hat, doch gleichwohl der Stätte, Zeit und Person gebraucht.

193. Dies sage ich also zum Exempel, damit Einer vielleicht den sehr heiligen Patriarchen entschuldigen möchte, daß er Gott nicht hätte versuchen wollen. Ich will ihn aber nicht entschuldigen, sowohl als andere Heilige auch nicht, als hätten sie nie gezweifelt oder gesündigt; denn es ist uns ein heilsam und sehr tröstlich Exempel, und viel angenehmer, denn so uns die großen ritterlichen Thaten Simsons oder Davids vorgehalten würden, davon sie keine Verheißung gehabt, sondern daß sie einen gewissen ritterlichen Geist des Glaubens hatten.

194. Nun müssen wir auch von dem Gelübde sagen, so Jacob gethan hat. Denn die Papisten führen diesen Text an und gebrauchen ihn wider uns, die wir ihre Klostergelübde strafen und verwerfen. Jacob, sagen sie, hat ein Gelübde gethan und hat es gehalten; derhalben folget, daß wer ein Gelübde thut, dasselbe auch zu halten schuldig ist. Und sie führen darneben auch an den Text 4 Mos. 30, 3., da Mose also sagt: „Wenn jemand dem Herrn ein Gelübde thut, oder einen Eid schwöret, daß er seine Seele verbindet, der soll sein Wort nicht schwächen, sondern alles thun, wie es zu seinem Munde ist ausgegangen.“ Item 5 Mos.
\quellseite{125}

23, 21. 22.: „Wenn du deinem Herrn, deinem Gott, ein Gelübde thust, sollt du es nicht verziehen zu halten; denn der Herr, dein Gott, wirds von dir fordern, und wird Sünde sein. Wenn du das Geloben unterwegen lässest, so ist dirs keine Sünde“ etc. Diese und viel andere Sprüche und Exempel führen die Mönche wider uns an, auf daß sie damit den Greuel ihrer Gelübde bestätigen mögen.

195. Auf daß wir aber ihre Greuel und ihre Irrthümer widerlegen mögen, soll man zuerst diese Regel fleißig merken, nämlich, daß die Rechtfertigung, Vergebung der Sünde, Gnade oder Barmherzigkeit das Erste und gleichsam der rechte Eckstein ist; und dieselbe kommt einfach aus lauter Gnade und kann mit nichten durch Werke erworben werden. Diese Regel sollen wir allezeit vor Augen haben, daß wir weder durch Werke, noch durch das Gesetz gerecht werden; sondern ohne Verdienst, also daß Gottes Gnade vor allen unsern Werken und Verdienst hergeht. Wie im Propheten Jesaia Cap. 65, 1. steht: „Ich werde funden von denen, die mich nicht suchten.“ Dieses soll der Eckstein sein, und Gott soll der Anfang sein unsrer Seligkeit, der sich uns offenbart, daß wir ihn erkennen lernen. Dieses ist das Hauptstück und der Grund, so uns in der ganzen heiligen Schrift vorgehalten wird.

196. Das Allererste ist Gottes Wort selbst, wie auch die Creaturen selbst das mündliche Wort sind, dadurch alle Heiden Gott erkennen sollten, wie St. Paulus Röm. 1, 19. ff. sagt. Es muß etwas sein, es sei entweder das Wort, oder sonst irgend eine That, so da vorhergeht und uns bewegt, und diese erste Bewegung muß von Gott kommen. Das ist gewiß: Die Person wird erstlich angenehm. Ja, sage ich, sie wird angenehm durch die Gnade, die da gerecht macht, und durch die Gaben des Heiligen Geistes, dadurch der Mensch Gott erkennt, daß er ein solcher Heiland sei. Dieses ist recht die erste Gnade, da wir nichts thun, sondern nur leiden: wir hören, daß Gott sein Wort redet, und empfinden, daß er in uns wirkt durch das mündliche Wort und die Sacramente, dadurch er in uns die Erkenntniß erweckt, daß wir ihn erkennen lernen.

197. Wenn nun das feststeht, welches gewiß und das Erste ist, so folgt dann der andere Satz, nämlich, daß die guten Werke von wegen Anderes nichts sind denn eine Danksagung, und daß dadurch Gott soll gepriesen werden; es seien gleich, was für Werke es sein mögen, Almosen, oder Gelübde, oder Leiden, oder endlich auch der Tod selbst: hie soll man alle zählen unter das Dankopfer. Denn wir können Gott nichts geben, wir können auch nicht verdienen oder zuwege bringen, daß unsere Personen angenehm werden. Nun sind aber die Klostergelübde nur äußerliche Werke; denn ich und alle Mönche geloben Gott Keuschheit, Armuth etc. Warum thun wir aber das? Antwort: Daß wir damit Vergebung der Sünde verdienen und dadurch selig werden wollen. Derhalben sind solche Gelübde stracks wider den Artikel der Rechtfertigung als die allerschändlichsten Greuel und Gotteslästerungen; und es sind alle Klostergelübde gottlos, damit Gott seine Ehre geraubt wird.

198. Denn die Mönche gehören unter die Zahl derer, die da Gott suchen, welche Gott verwirft und verdammt, und sagt dagegen durch den Propheten Jesaia Cap. 65, 1. also: „Ich werde funden von denen, die mich nicht suchten.“ Item: Nehmt eure Opfer hinweg und esset sie, ich will eure Opfer nicht haben, ich frage nach euren Gelübden nichts: „Weisest du, daß ich Ochsenfleisch essen wolle?“ etc., wie im 50. Psalm V. 13. steht, und im Propheten Jesaia am 66. Capitel V. 1.: „Was ists denn für ein Haus, das ihr mir bauen wollt?“ etc. Wo habe ich vorhin gewohnt, ehe denn ich Himmel und Erde gemacht habe? Sondern dieses habe ich euch gesagt: Höret meine Stimme, alles Volk in Juda soll das Wort des Herrn hören; mein Wort soll bei euch sein, daß ich der Anfang eurer Seligkeit sei. Denn ich bin gefunden von denen, die mich nicht gesucht haben; auf daß du wissest, was die Rechtfertigung ohne Verdienst, oder Vergebung der Sünde sei, dadurch die Person Gott angenehm ist, nicht um irgend eines Werks willen, es sei welcherley Werk es immer sein möge. Wie sich die Juden heutiges Tages über die Maßen sehr bemühen und fastieren, daß sie bei Gott Gnade finden und den Zorn Gottes stillen mögen; aber sie richten damit nichts aus. Das ist nun gewiß und beständig genug.

199. Derhalben ist nun ein großer Unterschied zwischen dem Gelübde Jacobs und den Klostergelübden. Das Gelübde, so Jacob gelobte, gethan hat, ist kein Gelübde von der Rechtfertigung oder Vergebung der Sünden; ja, wo nicht zuvor die Person durch die Vergebung der Sünden angenehm und ein Kind Gottes wäre, so würde er durch das Gelübde nicht gerecht, und was noch mehr ist, er würde auch kein Gelübde thun. Das Gelübde macht keinen Baum, wie die Frucht des Baumes auch der Baum nicht ist; sondern das Gelübde ist eine Frucht des Glaubens.

200. Also verwerfen und verdammen wir nun überall die Klostergelübde und alle Opfer, auch die im Gesetz und Mose vorgeschrieben waren. Denn von den denselbigen sagt Gott auch: Ich will eure Opfer nicht haben. Ei, warum das? Antwort: Ihr wollt mir die Figur hysternon proteron daraus machen; ihr wollt mir das Hinterste zuvörderst kehren; ihr wollt mich mit eurem Opfer suchen und versöhnen. Das will ich kurzum nicht leiden, daß ihr etwas dazu thun sollt, daß ihr mich damit versöhnen möchtet; sondern erkennt mich für den Gerechten, nachher und der euch Alles aus Gnaden gibt. Alsdann wird mir euer Werk gefallen, wenn es entweder mir zur Danksagung, oder dem Nächsten zu Nutz und gutem Exempel geschieht. Wenn ihr aber mich mit Ochsen- oder Bockfleisch, oder aber sonst mit eurem Opfer speisen wollt, das wird gar vergeblich und umsonst sein; denn ich will euer Fleisch nicht essen, ich will euer Opfer und Gottesdienst nicht annehmen, damit ihr das sucht und verdienen wollt, daß ich euch meine Gnade verleihen soll. Denn die Gnade soll ohne Verdienst und umsonst sein, sonst heißt es nicht (gratia) Gnade.

201. Darnach will ich eure Gottesdienste, euer Brandopfer und andere Opfer annehmen, dieweil du mir nicht um der Ursache willen opferst, daß du mich damit versöhnen wollest, sondern mich für den erkennst, der ich bin, Gnade erzeige, dich gerecht mache, und der ich dir Alles aus Gnaden, ohne Verdienst und umsonst gegeben habe. So will ich deine Opfer annehmen als eine Danksagung, und daß ich gepriesen werde; wie St. Paulus Röm. 1, 21. von den Heiden sagt: „Dieweil sie wußten, daß ein Gott ist, haben sie ihn doch nicht gepriesen als einen Gott.“ Gott hat es ihnen geoffenbart und gegeben; da sie ihm aber danken und loben sollten, da haben sie ein Anderes gethan, nämlich, daß „sie sind in ihrem Dichten eitel worden, und ihr unverständiges Herz ist verfinstert“ etc.

202. Die erste Gabe ist, daß man Gott erkenne; darauf soll die Danksagung und das Lob folgen; aber sie sind in ihrem Rathe zu Narren geworden, und haben sich andere, abgöttische Götter gemacht, denen sie geben und opfern und von denen sie nichts nehmen oder empfangen wollen. Das heißt einen Abgott machen und den rechten wahren Gott verlassen, und sich selbst preisen und Gott nicht. Denn sie denken also: Gott wird mir das ewige Leben geben um meines Verdienstes, um meiner Kappe und übrigen Werke willen. Das ist eine greuliche Lästerung Gottes.

203. So zeigt nun dieser Text an, daß Jacob ein solch närrisch und gottlos Gelübde nicht gethan hat, sondern es ist ein Gelübde der Danksagung gewesen; denn er sagt also: „So mich Gott mit Frieden wird wieder heimbringen“ etc. Als wollte er sagen: Ich will mich jetzt dazu verpflichten, daß ich schuldig bin, und will dasselbe auch halten und bezahlen, wenn ich wieder heimkomme. Was soll es denn sein? Antwort: Ich will daselbst eine Schule und Kirche bauen, und dieselbigen zu bestellen und zu erhalten will ich den Zehnten geben von Ochsen, Ziegen, Aepfeln, desgleichen den zehnten Theil von Milch und Butter. Desselbigen allen bedarf zwar Gott nicht; denn er ist kein Brod, trinkt auch keinen Wein, wie der 50. Psalm V. 13. sagt: „Meinst du, daß ich Ochsenfleisch essen wolle?“ etc. Was thust du denn nun? Antwort: Jacob ist jetzt schon gerecht, derhalben thut er das Gelübde nicht, daß er damit Gott versöhnen wolle; sondern er will dies thun, Gott damit zu danken und ihn zu preisen, daß er ihn hinführen und wieder heimbringen wolle laut der Verheißung. Wer wird aber den Zehnten nehmen? Antwort: Die Engel, Sonne und Mond werden ihn nicht nehmen, sondern er will das herrliche und vortreffliche Werk thun, das Predigtamt damit zu erhalten, daß man davon Schulen und Kirchen stiften möge, da die Propheten und der Propheten Kinder sich erhalten und genähret werden mögen, die sonst Hunger und Durst leiden müssen, die nackt, arm und elend sind.

204. Wie man heutiges Tages sehen mag, wie deren so wenig sind, die sich solche Armuth und Elend zu Herzen gehen lassen, und wie böslich und untreulich die, so da reich sind und Geldes und Gutes Vollauf haben, etwas dazu geben oder steuern, daß die Jugend, welche an derjenigen Statt kommen soll, die jetzt das Lehramt verwalten, möge erzogen und gelehret werden.
\quellseite{126}

205. Darum hat mir sehr wohl gefallen, daß der Durchlauchtige Herr und Fürst, Herzog Heinrich von Sachsen, der nach Absterben seines Bruders, Herzog Georgs, an dessen Statt zum Regimente gekommen, gesagt und das vornehmlich befohlen hat, daß man zu den Schulen und Kirchen ehrliche Besoldungen machen sollte; denn dieses sind seine Worte gewesen: Ich muß dennoch meinen Pfaffen zu essen und zu trinken geben. Deren sind aber fast wenige, es seien gleich Fürsten oder sonst gemeine Leute, die solche löbliche Worte reden oder solchen Exempel nachfolgen. Darum stehen die Schulen und Kirchen in großer Gefahr, sintemal es bald hernach an solchen Leuten mangeln wird, so die Schulen und Kirchen nähren und erhalten sollen.

206. Derhalben ist dieses Gelübde, so Jakob gethan hat, nicht allein hoch zu loben, sondern auch sehr nöthig gewesen. Und also soll man Mose verstehen, da er 5 Mos. 23, 21. sagt: „Wenn du dem Herrn, deinem Gott, ein Gelübde thust, so sollt du es nicht verziehen zu halten“ etc. Was heißt Gott ein Gelübde thun, oder wen nennt er also, da er sagt: „Dem Herrn, deinem Gott“? Antwort: Es heißt den Zehnten geben den Schülern, Kindern und Eheweibern der Propheten, die da elend, arm und nothdürftig sind; die soll man ernähren und unterhalten von den Zehnten, die Gott gegeben sind; dadurch wird Gott ernährt. Denn was du geben wirst, spricht Gott, den Kindern der Priester und Propheten, durch welche die Lehre ausgebreitet wird, davon sollst du denken, es sei mir gegeben; nicht daß ich dich um des Werkes willen gerecht wolle machen, sondern ich will es annehmen als eine Danksagung und Dankopfer dafür, daß ich dich gerecht gemacht und dich in dem Zeitlichen auch gesegnet habe, dich auch hinfort mehr segnen will, wo du den armen Kindern der Propheten helfen, sie nähren und erhalten wirst.

207. Es ist ein groß Ding, daß die Schrift in Mose, in den Büchern Salomons und in den Psalmen so oft wiederholet: Wenn du dem Herrn ein Gelübde thust, so sollt du es nicht verziehen zu halten. Und dasselbe Gebot betrifft am allermeisten die Priester, auf daß das Predigtamt möge erhalten werden. Wenn du gelobt oder verheißest, die Schulen zu schützen und zu erhalten, so siehe zu, daß du dasselbe auch haltest. Und Christus spricht auch Matth. 26, 11.: „Ihr habt allezeit Arme bei euch“ etc., vornehmlich zu der Zeit, wenn das Evangelium geoffenbart wird.

208. Denn es geht dem Worte Gottes, und denen, die es lehren und predigen, wie es sonst den andern Künsten allen auch geht, wie man im deutschen Sprüchwort sagt: Die Kunst geht nach Brod. Derhalben lautet die göttliche Verheißung im 41. Psalm V. 1. also: „Wohl dem, der sich des Dürftigen annimmt.“ Denn du hast die erste Gnade hinweg, und bist schon aus Gnaden und umsonst gerecht worden; dasselbe sollst du erkennen und Gott dafür danken; alsdann will ich dich, spricht er, also segnen, dich also reich machen und mit Wohltat überschütten, daß du empfinden sollt, dir sei übrig genug geschehen. Und dieses sollte uns vornehmlich dazu ermahnen, daß wir beiden, Schülern und Lehrern, gerne sollten helfen; denn solche Almosen werden an diesem Ort Zehnten genannt, die nicht ihnen, sondern Gott gegeben sind. Und Christus sagt selbst auch Matth. 25, 40.: „Wahrlich, ich sage euch, was ihr gethan habt einem unter diesen meinen geringsten Brüdern, das habt ihr mir gethan.“

209. Die Leviten und Priester im Gesetze Mosis hatten mit dem Volke Israel keinen Theil und Erbe im Lande, und empfingen nichts von dem Erbe ihrer Brüder: „denn ich“, spricht der Herr 4 Mos. 18, 20., „bin dein Theil und dein Erbgut unter den Kindern Israel.“ Was hatte aber der Herr, oder wie hat er können ihr Theil sein? Antwort: Die Gelübde, Opfer und Zehnten, dieselbigen, sagt er, will ich euch geben, daß ihr damit euch und euer Gesinde erhalten und nähren möget, daß ihr dabei meines Amtes und Dienstes warten könnet, nämlich, daß ihr lehret und lernet, und mein Wort also erhaltet. Derhalben hatten die Priester kein eigen Erbe, hatten auch kein Theil im Lande unter den Kindern Israel; sie hatten von ihren Aeckern oder Erbe nichts zu essen, sondern sie aßen mit unserm Herrn Gott von seinen Zehnten.

210. Wo du nun armen Schülern etwas gibst, so hast du es Gott selbst gegeben, und sollst wissen, daß du Gott damit einen sehr angenehmen Dienst und ein solch Opfer gethan hast, das diesen herrlichen und hohen Titel hat, daß es ein Gelübde Gottes genannt wird. Wer aber das glaubt und thut, der soll wissen, daß es eine sonderliche Gabe Gottes ist. Denn die Welt achtet dessen nicht, und darum kommt auch der Türke und nimmt das alles mit Gewalt hinweg, was wir zusammen gescharrt, und denen verweigert haben zu geben, denen wir es billig sollten gegeben haben. Denn Gott will mit uns essen, oder will uns nichts von den Gütern lassen, die wir haben. Ist aber das nicht eine herrliche Ehre, darnach uns alle billig verlangen sollte, daß ich Gott selbst nähre und speise, wenn ich irgend einen Lehrer oder Schüler nähre? Denn so gibt du den Zehnten zum Hause Gottes, welches ohne die Schulen und Schüler nicht kann erhalten werden.

211. Wie Jakob auf dieser Reise die Kinder Sem hat wohnen sehen bei dem Berge Moria, welche weiter auch Söhne, Töchter, Sohnes- und Töchterkinder gehabt haben, dazu sonst auch Blutsfreundinnen, die alle fast elend und arm gewesen sind. Denn dieweil derselbige Berg vorhin durch die Tyrannen verwüstet war, wird ihn vielleicht Eber wiederum aufgerichtet haben; und nach Eber, wo Abraham und Isaac gesehen, daß er verwüstet gewesen ist, haben sie ihn auch wiederum aufgerichtet und gebaut. Also sieht nun Jakob, daß die Kinder, beide Söhne und Töchter, der Patriarchen an demselben Ortte Hungers und Noth leiden; darum thut er ein Gelübde, wo ihn der Herr wiederum gesund heimbringen werde, so wolle er ihm den Zehnten geben, das ist, er wolle die Schule der Patriarchen wiederum aufrichten, und dasjenige dazu geben, was vonnöthen sein wird, daß sie möge erhalten werden.

212. Die Gelübde im Gesetze, dergleichen das der Nazaräer war, waren außerhalb der gemeinen Ordnung aufgerichtet und waren überflüssig. Das ist aber ein recht Gelübde, das auf die Rechtfertigung folgt, und mag ein Gelübde der Liebe genannt werden, wie dieses Gelübde des Patriarchen Jakob war. Denn das höchste und vornehmste Gelübde ist im ersten Gebot, da wir mit Gott einen Bund machen, daß er unser GOtt sein wolle und daß wir an GOtt den Vater glauben wollen: das ist der Bund in der ersten Tafel, der durch die Leviten'schen Gelübde und Opfer bedeutet worden ist, und ist ein Gelübde der Gnade und Gerechtigkeit. Dieses aber sind Gelübde der Liebe und sind natürlich, die für und für gebraucht werden; davon im 5. Buch Mose am 23. Cap. V. 22. 23. gesagt wird: „Wenn du das Geloben unterwegen lässest, so ist es dir keine Sünde“; hast du aber ein Gelübde gethan, so gedenke und halte es auch. Wir sind nicht verpflichtet oder gezwungen, auf eine gewisse Weise oder Regel Gelübde zu thun.

213. Wo aber die Kirchendiener nicht genähret oder erhalten werden, alsdann wird dieses Gelübde ein natürlich und sittlich Gebot, so zu diesem Leben nöthig ist, nämlich, daß ein Jeder unter uns an seinem Ort das Predigt- und Kirchenamt helfe erhalten. Wie der Durchlauchtige Fürst, Herzog Hans Friedrich, Churfürst zu Sachsen, alljährlich dreitausend Gulden gibt, die hohe Schule zu Wittenberg allhier damit zu erhalten. Das ist wahrhaftig das rechte Gelübde Jakobs. Denn wie wiewohl er mit ausdrücklichen sonderlichen Worten nicht sagt, daß er es Gott gebe oder ihm ein Gelübde thue, so ist es doch in Wahrheit ein göttlich Gelübde, damit er es sich so viel kosten läßt, daß die Pfarrherren, die Magister und die Professoren der Künste ernährt und erhalten werden mögen; und ist Seine Churfürstliche Gnaden schuldig, solches also zu halten. Wiewohl es eben mit der That kein Gelübde ist, so ist es doch in Rechten ein Gelübde; denn wir sind alle schuldig, vermöge unserer christlichen Religion, dermaßen zu geloben und den Schulen zu geben, wenn wir schon mit der That kein Gelübde thun.

214. Syrach sagt, was zur Nothdurft gehöre, daßselbe solle dem Gelübde nicht unterworfen sein, sondern ein Gelübde sei freiwillig. Wie die Münche Gehorsam, Keuschheit und Armuth geloben; welche Stücke zwar vor dem Gelübde nicht nöthig sind, aber nach dem Gelübde werden sie sehr nöthig. Ein solch Gelübde ist das Gelübde der Liebe, damit ich irgend einem Armen gehen Gulden zu geben verspreche, deren ihm vonnöthen thut; daselbst bin ich schuldig zu halten, was ich verheißen habe. Ja, ich bin es auch schuldig, wenn ich sehe, daß mein Nächster Noth leidet, daß ich ihm soll zu Hülfe kommen in seiner Armuth; denn ich soll ja nicht leiden, daß Christus Hunger leide, daß er nackend oder arm sei.
\quellseite{127}

215. Auf die Weise gelobt nun Jakob, er wolle zu Morija eine Schule aufrichten, und den Schülern und Kindern der Propheten den nöthigen Unterhalt geben. Wenn ich wieder heimkommen werde, sagt er, so will ich auch Kindern Eber den Zehnten geben um des Hauses GOttes willen, der mich gesegnet und mir dieses Gesicht und die Pforte des Himmels gezeigt hat.

216. Daher kommt es, daß in den Psalmen hin und wieder die Gelübde Danksagungen genannt werden, welche sonst allen Menschen nöthig und geboten sind; und sind nicht allein Werke von Worten, sondern es gehört auch die rechte That dazu. Wie in den Sprüchen Salomos am 3. Cap. V. 9. steht: „Ehre den Herrn von deinem Gut, und von den Erstlingen alle deines Einkommens“ etc. Welches nicht allein von Worten und mündlicher Ehre, oder von Kniebeugen, sondern von der That an sich selbst zu verstehen ist; wie auch die Ehre, so man den Eltern zu erzeigen schuldig ist, nicht das allein fordert und haben will, daß wir vor ihnen den Hut abziehen, sie freundlich anreden und gehorsam sein, sondern daß wir sie auch ernähren und also mit der That ehren sollen. Und St. Paulus sagt auch 1 Tim. 5, 17.: „Die Aeltesten, die wohl vorstehen, die halte man zwiefacher Ehre werth“; item Gal. 6, 6.: „Der aber unterrichtet wird mit dem Wort, der theile mit allerlei Gutes dem, der ihn unterrichtet“ etc.

217. Derhalben, wo GOtt fordert, daß man ihn ehren soll, so will er damit gemeint haben, daß man ihm opfern solle nicht allein das mündliche Opfer, sondern auch Ehrbietung mit dem Herzen; ja, er will auch die That selbst haben. Er will haben, daß ein Jeder unter uns das Predigtamt soll helfen fördern, und dazu geben und steuern, daß die Schüler bei ihrem Studiren erhalten und die Lehre also ausgebreitet werden möge. Da wird dann GOtt recht und wahrhaftig gelobt und gepreiset, da sollt du ein Gelübde thun und sagen: Ich verheiß und gelobe an, daß ich etwas dazu geben und steuern wolle, damit der Kirche möge geholfen werden. Das ist das Gelübde des heiligen Patriarchen Jakob.

218. So ist nun zweierlei Gelübde durch die Gelübde im Gesetz bedeutet worden. Das eine ist, das da gehört zum Glauben gegen GOtt, der mit uns einen Bund macht, daß er unser GOtt sein wolle; wie wir droben vom Gelübde des ersten Gebotes gesagt haben. Das andere ist ein Gelübde der Liebe, welches des Gebotes und Wortes halber zwar frei ist, nämlich, wenn ich nicht weiß, daß mein Nächster Noth leidet, oder wenn er meiner Hülfe nicht bedarf; es wird aber nöthig, wenn er meiner Hülfe bedarf und Andere ihm nicht helfen wollen. Wiewohl ich es nun dem Herrn eben nicht gelobet habe, so soll ich in solchem Fall dem Nächsten dennoch helfen. Denn Christus lehret und gebietet uns solches, da er sagt Luc. 6, 30.: „Wer dich bittet, dem gib.“

219. Die Welt aber verachtet nicht allein solche Werke und Gelübde, sondern sie beraubt auch die Kirchen und Schulen der Zehnten, so andere fromme Leute dazu gegeben haben, und geht nur damit um, daß wir nur wohl hungern und allenthalben darben sollen.

220. Darum soll man den Leuten dieses Exempel fleißig einprägen, daß Jakob gelobet den Zehnten; nicht daß GOtt selbst davon essen oder reich werden solle, sondern er gibt ihn den armen Kirchendienern, und denen, die allzeit die geringsten Brüder oder Jünger Christi auf Erden sind. Denn alle anderen Künste haben ihren Gewinn und man kann damit etwas erwerben; nur diese Profession derer, so das Kirchenamt zu verwalten haben, bedarf, daß man ihr Hülfe und Steuer thue. Wir müssen vom Altar leben, wie St. Paulus sagt. Wer nun gottseelig ist und GOttes Wort lieb hat, der gibt etwas; wer aber der Lehre sammt den Lehrern und Schülern feind ist, der raubt und reißt an sich, wie er nur kann und mag, auch das von Andern gegeben ist.

221. Es hat es im Uebrigen GOtt also geordnet und befohlen, daß wir leben und unseren Unterhalt haben sollen von den Gelübden, Opfern und Almosen, welche nach göttlichem Recht uns gebühren; und wenn gottseelige Leute die Armuth und Nothdurft der Kirchen- und Schuldiener sehen, geloben sie und sind auch schuldig, den Zehnten zu geben; wie ihn Jakob gegeben hat, der die Größe und Herrlichkeit dieses Amtes recht wohl verstanden hat. Denn er hat gesehen, daß die Kirche dazumal betrübet und verlassen war, und gar in der Irre gegangen ist, dazu auch zum öftern von den Tyrannen sehr hart verfolgt und zerrissen worden, da Sem und auch die anderen Patriarchen daselbst die Verheißung gepredigt haben von dem künftigen Samen.

222. Darnach hat Jakob wohl gewußt, daß dieses die Linie und Ordnung wäre, davon Christus hat sollen herkommen, wie droben, Cap. 10. V. 21., Sem ein Vater aller Kinder Eber ist genannt worden, darum daß von den Nachkommen Ebers Christus sollte geboren werden. Denn Eber hat zwei Söhne gehabt, Peleg und Jaktetan; von Jaktetan aber sind dreizehn Söhne geboren worden, und ist Sem deren aller Vater genannt worden. Von diesen Kindeskindern Ebers sind viele Priester geworden, welche in den Schulen und Kirchen die Lehre vom künftigen Samen gelehret haben; wiewohl solche Schulen und Kirchen oftmals zerstört worden sind, wie das am Exempel dieser Patriarchen, Abraham, Isaak und Jakob, welche zum öftern umziehen und andere Wohnungen suchen müssen, und bald zu Gerar, bald in Egypten, bald im Lande Canaan Fremdlinge gewesen sind, wohl zu sehen ist. Wie allezeit die Kirchen und Schulen in der Welt müssen in der Irre gehen und von einem Ort zum andern ziehen, wie etwa die Völker, Nomaden genannt, gethan haben. Darum haben sie der Opfer, Zehnten, Almosen und anderer thätlicher Handreichung bedurft. Und da Jakob gesehen hat, daß es daran in derselbigen Versammlung gemangelt, thut er ein Gelübde und verheißt, daß er den Zehnten daselbst geben wolle, wenn er wiederum heimkommen werde; und hernach im 35. Cap. V. 1. hat ihm GOtt geboten, wiederum an diesen Ort zu ziehen und daselbst einen Altar aufzurichten.

\chapter{Das neunundzwanzigste Capitel}

\section[Erster Theil]{Erster Theil. Von Jacobs Reise nach Haran}

\begin{quote}
\textbf{V. 1—3. Da hub Jakob seine Füße auf, und ging in das Land, daß gegen Morgen liegt; und sahe sich um, und siehe, da war ein Brunnen auf dem Felde, und siehe, drei Heerden Schafe lagen dabei, denn von dem Brunnen pflegten sie die Heerden zu tränken, und lag ein großer Stein vor dem Loch des Brunnens. Und sie pflegten die Heerden alle daselbst zu versammeln, und den Stein von dem Brunnenloch zu wälzen und die Schafe zu tränken, und thaten alsdann den Stein wieder vor das Loch, an seine Stätte.}
\end{quote}

1. Bisher hat Mose beschrieben, wie der heilige Patriarch Jakob sich in den hohen und recht geistlichen Werken des Glaubens und des göttlichen Wortes geübt hat; welches Stück in allen Historien der Heiligen werth ist, daß man es fleißig merken soll, nämlich, wo sie GOttes Wort hören, demselben glauben und im Glauben durch mancherlei Trübsal und Leiden versucht werden. Denn wiewohl sie da etwas schwach sind im Glauben, sind sie dennoch ganz göttliche und himmlische Menschen, sind ganz rein und heilig; und endlich leben und wandeln sie nicht vor den Menschen, sondern vor GOtt.

2. Dieses ist an den heiligen Vätern die rechte Würde und Hoheit. Wie auch in unserm Leben, wenn wir uns in der Kirche mit GOttes Wort üben, die heiligen Sacramente gebrauchen, wir mit mancherlei Anfechtungen versucht werden und unser Glaube wie das Gold im Feuerofen geprüft wird. Das ist die rechte Heiligkeit, davon wir auch heilig genannt werden und uns heilig sind. Denn der heilige Geist macht uns heilig durch das Wort, wo wir dasselbe durch den Glauben ergreifen und annehmen, und das Fleisch durch das Leiden und Trübsal tödten; auf daß also die Heiligen lebendig gemacht werden, und ihre Leiber begeben zum Opfer, daß da lebendig, heilig und Gott wohlgefällig sey etc., wie St. Paulus sagt Röm. 12, 1. Dieses ist der höchste Grad des Lebens der Heiligen.
\quellseite{128}

3. Damit wir aber in unserm Herzen nicht verzagen möchten, wo wir von den heiligen Vätern hören würden, daß sie uns nicht anders denn nur in solchem hohen Grade und heiligen Leben vorgehalten würden, darum führt sie Gott wiederum herab vom Himmel auf die Erde, und beschreibt sie also, wie die allergeringsten Menschen sein mögen, die gar verachtet sind, daß man auch nichts Geringeres oder Verächtlicheres nennen möchte, ausgenommen, daß gleichwohl ihre Sünden nicht gelobt werden.

4. Sie werden im Uebrigen gar gering beschrieben, gleich als die im Staube und Mühseligkeit der Haushaltung und des Weltregimentes immer daher kriechen und niedrig genug sind. Denn sie gehen mit solchen Werken um, die gar nichts werth und ohne alle Heiligkeit zu sein scheinen; so doch die Papisten dagegen nur nach solchem Leben trachten und sich dessen auch hoch verwundern, das schlechterdings von weltlichen Geschäften abgesondert, und gar ohne alle Sorge ist beider des Haus- und Weltregimentes. Ein solch Leben, das dermassen von allen leiblichen und weltlichen Geschäften und Aemtern abgesondert ist, nennen sie Heiligkeit und Gerechtigkeit: aber sie irren und fehlen damit weit. Denn also sollen wir den Glauben und die Gerechtigkeit suchen, daß wir dennoch die Münze, Zinn und Kümmel verzehnten; dieses sollte man thun, und jenes nicht lassen, wie Christus spricht Matth. 23, 23., das ist, dieweil wir im Fleische wohnen und leben, so müssen wir derhalben auch des Leibes und Fleisches warten; aber doch die Sünde ausschließen.

5. Man muß das Welt- und Hausregiment verwalten; sintemal wir noch nicht im Paradiese und den Engeln auch nicht gleich sind, sondern wir leben noch im Fleische, in diesem natürlichen Leben, dem Essen, Trinken, Kleidung, Haus, Kinder und Ackerwerk vonnöthen sind: man bedarf auch des weltlichen Regimentes, und daß wir Schutz haben mögen wider die bösen Buben, die uns beleidigen. Darum ist vonnöthen, daß man die zwei Aemter, so in diesem Leben gehen müssen, im Schwange behalte, nämlich, nähren und wehren: im Hause nährt man die Kinder und Gesinde; das weltliche Regiment aber schützt und beschirmt dieses alles.

6. So werden nun, sage ich, die heiligen Väter niedrig und fleischlich beschrieben, in dem niedrigen Stande dieses Lebens, da vor den Augen der Papisten nichts Unreineres oder Verächtlicheres ist. Wie sie sagen, es werde hier nichts Anderes vorgehalten, denn daß sie Weiber genommen, Kinder gezeugt, die Kühe und Ziegen gemolken haben etc., welches ganz weltliche und heidnische Werke sind. Aber der Herr hat uns nun den rechten Verstand gegeben, dafür sollen wir ihm dankbar sein. Denn wir können dieses alles eigentlich und genauer ansehen, als die Papisten, welche Anderes nichts sehen, denn die fleischlichen Werke, die doch wohl so gar fleischlich nicht sind, als ihre Werke, so vom Teufel regieret werden und ohne den Heiligen Geist geschehen. Denn wenn sie schon fasten, sich des Ehestandes enthalten, in ihren Kirchen ein Gemurmel machen, so ist doch in ihrem Gebet kein Geist, keine herzliche Andacht, und ist kein rechter Gottesdienst, damit sie umgehen: es ist alles voll Geizes, Abgötterei, eitler Ehre und Verachtung Gottes.

7. Wir lehren aber also, daß man zu allererst auf die Person sehen soll, ob dieselbe gerecht und gottselig sei, welches geschieht durch das Wort, dem sie glaubt. Darnach aber richtet sie in der Kirche oder Gemeinde ihr Amt aus, das ist, sie lehrt, vermähnt, betet, lernt und hört das Wort; sie trägt das Kreuz um des Wortes willen, sie wird nach dem Fleische getödtet. Das ist eine heilige lebendige Person und die Gott wohlgefällig ist, welche fortgeht und andere Werke auch thut, wenn sie das Wort gehört und demselben geglaubt, dazu gebetet und ihr Kirchenamt ausgerichtet hat.

8. Wie David, wenn er des Gottesdienstes gepflegt hat, sich des Regimentes annimmt und dasselbe verwaltet, das Recht spricht, Krieg führt, sein Kriegsvolk zum Streit rüstet, die Feinde angreift, sie würgt und ihr Blut vergießt. Das sind nicht eigentlich geistliche, sondern weltliche Werke, die zum Weltregiment gehören. Wenn nun Einer fragen wollte: David ist dann nicht heilig, dieweil er ein Kriegsmann ist und das Schwert führt: der urtheilt zu grob von der Sache, wie die Papisten thun. Denn wer ist der David, der da Blut vergießt und Krieg führt? Antwort: Er ist eine solche Person, die in der Kirche durch das Wort und den Glauben gerechtfertigt ist; darnach aber hat sie ein Amt im weltlichen Regiment, darum richtet sie, verdammt und rechtfertigt, verwaltet das Regiment, straft die da schuldig sind, führt Krieg, und bleibt nichtsdestoweniger ein frommer gläubiger Mensch und ein guter Baum. Solch Blutvergießen aber gefällt Gott wohl, wiewohl die Welt und die Mönche sammt allen andern Heuchlern sich daran sehr ärgern, darum daß sie nur die Larve oder den Schein der äußerlichen Werke ansehen. Sie sehen das Wort und den Glauben, und auch den Geist Gottes nicht, und wie ihn Gott zu solchen Werken getrieben hat, der die Person nicht allein in den hohen, sondern auch in den geringen Werken, das Haus- und Weltregiment belanget, regieret. Denn dazu ist David von Gott berufen, wie die Schrift von ihm sagt, daß er allen meinen Willen thun soll, welcher Wille Gottes ihm das auch befohlen hat, daß er die Philister, Damascener, Amalekiter und Ammoniter etc. hat demüthigen und schlagen sollen.

9. Was ist aber das? sagen die Papisten, er sollte derweil gebetet, geopfert und Gott im Tempel gedient haben. Es ist recht gebetet, und er thut auch zu gelegener Zeit und nach gebührlicher Ordnung, was sich in den Kirchen- und Religionssachen gebührt zu thun. Des Morgens hat er gebetet, von Gottes Wort geredet, demselbigen geglaubt, Psalmen gesungen, und sich dessen fleißig angenommen, was zur Kirche und Religion gehört. Darnach aber nimmt er sich auch des Welt- und Hausregiments an, zeugt Kinder, versorgt sein Gesinde, ißt und trinkt. Also führt David einen göttlichen und heiligen Wandel in allen drei göttlichen Ständen oder Regimenten, als, im Kirchenregiment und dazu auch im Welt- und Hausregiment.

10. Auf solche Weise sollen wir unser Leben auch anrichten, daß wir uns finden lassen in einem solchen Stande, der Gott wohlgefällig sei nach seinem Wort. Vor allen Dingen gedenke, daß du dem Worte glaubest und dasselbe öffentlich bekennest, und daß du bereit seiest, um des Wortes willen zu leiden und zu sterben. Darnach, du seiest ein Regent oder Hausvater, so warte deines Berufs an deinem Orte, dahin du bestellt bist. Ein solch Leben gefällt Gott wohl, und er gibt dazu viel große Belohnung und läßt Alles durch seinen Segen wohl gerathen.

11. Also soll man nun die Exempel der Väter ansehen und davon handeln in dem hohen Grade, daß man gleichwohl die geringen fleischlichen und unreinen Exempel, wofür sie von den Papisten angesehen werden, des Haus- und Weltregimentes nicht verachte. Denn dann finden sie nicht, sind auch nicht gering oder verächtlich, wenn sie von einer solchen Person geschehen, die da gläubig, Gott angenehm, heilig und göttlich ist, die da weiß, daß Gott Alles wohlgefällig ist, was sie thut; aber doch also, daß die rechte Ordnung nicht umgekehrt werde, sondern daß sie im Glauben bleibe und lebe: darnach gefallen auch Gott die Werke wohl, so eine solche Person in ihrem Beruf thut.

12. Auf diese Weise beschreibt Mose den Patriarchen Jakob, wie er gen Haran gekommen sei, und daselbst seines Vetters Tochter gefunden, die Jungfrau lieb gewonnen, sie zum Weibe genommen, Kinder gezeugt und der Sache gehütet habe; welches alles närrische und fleischliche Werke sind, daß sie auch bei den Heiden nicht fleischlicher könnten gefunden werden. Denn niemand sieht den rechten wesentlichen Unterschied, welcher sehr groß ist, zwischen einem Heiden und diesem Patriarchen Jakob. Denn Esau und Ismael bauen auch den Acker, hüten der Schafe, melken die Kühe, haben ihr Gesinde, geben ihrem Vieh sein Futter; welche Werke alle den Werken der Heiligen gleich sind und sind aber doch nicht heilig. Ei warum das? Antwort: Darum, daß ein großer und unbegreiflicher Unterschied ist zwischen den beiderlei Werken. Hier sei Jakob ist der Glaube und das Wort, dort aber ist kein Wort, sondern nur Unglaube. Derhalben sind die Werke Jakobs von den Werken Ismaels oder Esaus so weit unterschieden wie Himmel und Erde, wiewohl es dem äußerlichen Ansehen nach einerlei Werke sind.

13. Also rühmt Erasmus die Tugenden der Heiden sehr hoch, als, des Socrates, Cicero, Atticus etc., und macht einen Vergleich und sagt: Du wirst kaum bei den Christen solche Leute finden, die das gethan hätten, das Pomponius Atticus oder andere Heiden gethan haben; ja, bei den Christen werden deren Viele gefunden, die offenbar böse und schändliche Leute sind, gegen welche diese Heiden viel besser gewesen. Man soll aber darauf also antworten: Nach der Philosophie und in der Materie ist es gleich, das ist, so viel den Stand oder das äußerliche Leben belanget; aber wenn man die Species und den Unterschied ansieht, nicht. Denn wenn schon Cicero und Socrates Blut geschwitzt hätten, so würden sie dennoch um deswillen Gott nicht gefallen können. Und ist das nicht die Frage, was und wie große Werke Alexander der Große, Julius Cäsar und Scipio gethan, von welchen man weiß, daß sie größere Thaten ausgerichtet haben, denn jemals ein Christ gethan hat; denn du wirst solche Stärke im Kriege, item, solche Geduld, und wie sie allerlei Unglück und Mühseligkeit ertragen haben, nicht leichtlich an irgend einem Könige bei den Christen finden; ja, auch nicht einmal bei den Königen im Volke Israel, als, bei David und den andern. Warum willst du sie denn nicht hoch heben, oder allen Königen bei den Christen, dem David und den andern vorziehen?
\quellseite{129}

14. Antwort: Wenn mir die Wahl gegeben würde, so wollte ich das Werk eines Bauern oder einer Dienstmagd, die da Christen sind, wie unschätzig und bäurisch es auch sein möchte, lieber haben, denn alle großen Siege und Triumphe des großen Alexander, Julius Cäsar und anderer Heiden mehr. Ei warum das? Antwort: Darum, daß hier Gott ist, dort aber ist der Teufel, welches der rechte wesentliche Unterschied ist (differentia essentialis). Die Werke haben zwar einerlei Materie, aber wenn man die Species und den Unterschied ansieht, hilft Gott! wie sind sie dann so gar ungleich! Denn Gott spricht: Die Werke dieses Weibleins oder Magd, und die Hausarbeit, die sie thut, daß sie das Haus kehrt und der Hausmutter gehorcht, gefallen mir wohl. Denn er sieht an die Niedrigkeit seiner Magd, da auch keine großen herrlichen Werke sind, denn allein, daß sie daheim Magdarbeit thut, sei es in der Küche, oder aber bei dem Vieh. Solche Magde waren diese zwei auch, nämlich, Lea und Rahel, welche die Schafe ihres Vaters gemeinet, sie zum Wasser getrieben, da sie getrunken haben, zudem haben sie die Kühe und Ziegen gemolken. Diese Werke gefallen Gott wohl. Aber Hannibal, Alexander, Scipio und Cicero gefallen ihm nicht. Warum aber das? Antwort: Darum, daß sie zwar in genere, wie man in den Schulen pflegt zu reden, das ist, so viel die Werke an sich selbst belanget, einander gleich sind, oder wohl vielmehr jene Werke übertreffen; aber wenn man die Species und den rechten Unterschied ansieht, so übertreffen jene Werke diese weit.

15. Das können aber nicht alle Menschen sehen und Erasmus hat es selber auch nicht gesehen: allein die Gläubigen sehen, wie köstlich und wichtig die Werke der Christen sind. Der Glaube aber und das Wort machen die Werke wichtig, köstlich, theuer und groß; denn da ist Gott selbst und der Heilige Geist in dem, der das Werk thut.

16. Die Menschen aber, die da sind wie Rosse und Mäuler, lassen sich den äußerlichen Schein sehr bewegen. Und ich, da ich noch ein Mönch im Kloster war, war ich viel heiliger, als ich jetzt bin, so viel die äußerliche Larve und den Schein belangt. Ich betete mehr, ich wachte, fastete, fastete mein Fleisch, kurz, mein ganzes Leben hatte vor anderer Leute Augen einen großen Schein; wiewohl es solchen Schein vor meinen Augen nicht hatte; denn ich hatte fast einen zerbrochenen Geist und war immer betrübt. Jetzt aber esse ich und kleide mich nach gemeiner Weise, da sieht man an meinem Leben nichts Sonderliches und Vornehmliches vor andern Leuten. Zu jener Zeit aber, da ich noch ein Mönch war, that ich Anderes nichts, denn daß ich die Zeit übel zubrachte, verderbte meine Gesundheit; ja, das noch mehr ist, ich habe auch mein Gewissen mit der Werkgerechtigkeit verwundet und beschweret, daß es sich noch jetzt kümmerlich will heilen lassen. Denn über die Art und Natur, welcher das angehören und eingepflanzet ist, daß sie sich immer der Werke rühmen will, habe ich mir auch einen habitus, wie man es in den Schulen nennt, und eine Gewohnheit daraus gemacht, daß ich immer auf meine Werke und Würdigkeit sehen will. Das weiß ich aber gewiß, daß jetzt Eine Lection und Ein Pater Unser mehr gilt und auch Gott wohlgefälliger ist, denn alle das Gebet, so ich im Kloster ganze fünfzehn Jahre gemurmelt habe; denn ich weiß jetzt, daß ich gewißlich erhöret werde; und ich bedarf keines Wachens, oder sonderlichen Fastens und Kasteiung; denn Gott hat mir einen Engel des Satans gegeben mit vieler andern Beschwerung und Kreuz in dieser Welt, die mir mehr zu schaffen machen, denn jenes alles miteinander.

17. Und bisher haben wir nun von dem Stück gehandelt, daß der heilige Patriarch so gar niedrig und gering beschrieben wird, daß auch zwischen ihm und dem Allergeringsten und verächtlichsten Heiden kein Unterschied ist. Aber daher soll der Leser den Unterschied lernen zwischen einem Christen und einem Heiden, und ihrer beider Werk. Denn ist Einer ein Christ, so siehe, ob er auch das Wort in Ehren halte, dasselbe höre und liebe, und daß ihm solches ein Ernst sei; zu dem, ob er auch mit dem Kreuz und Trübsal beschweret und geplagt werde. Und wo dem eigentlich also ist, so wird er in die Kirche kommen, wird das Wort gern hören, wird seine Lust daran haben, wird demselben von Herzen glauben, wird beten, Gott danken, er wird ein gut Gewissen haben. Wenn du das sehen wirst, sollst du gewißlich schließen und es dafür halten, daß er ein heiliger Mensch sei, der da Gott wohlgefällig ist; und sollst ihn als einen heiligen Menschen loben und dich sein verwundern, und von allem, was er darnach thun wird, es sei gleich in Welt- oder Hausregiment, sollst du sagen, es sei Gott angenehm und wohlgefällig. Und Gott wird es auch belohnen nicht allein in diesem, sondern auch im zukünftigen Leben; wie in der Offenbarung Johannis am 14. Capitel V. 13. steht: „Ihre Werke folgen ihnen nach“, nicht allein die Werke, so die Kirche und Religion belangen, sondern auch die zum Welt- und Hausregiment gehören. Aber die Werke der Gottlosen und der Heiden gefallen Gott nicht, darum folgen sie ihnen nicht nach; sondern die sterben dahin ohne Ehre, ohne Gedächtnis, ohne Verwunderung. Unsere Werke aber haben die Ehre, daß sie Gott und die Engel ansehen, und folgen uns nach in jenes Leben.

18. Wir wollen aber nun den Text besehen, „Er hub seine Füße auf.“ Es ist eine eigene Art zu reden in der hebräischen Sprache. In der gemeinen lateinischen Uebersetzung steht: Profectus est: Er ist hingezogen. Siehe aber, wie ungleich die Auslegung auf diese Weise jener Gesandtschaft sei, davon droben (Cap. 24. V. 2. ff.) Mose gesagt hat, da Abraham seinen Knecht abfertigt, daß er Isaak, seinem Sohn und Erben der Verheißung, ein Weib nehmen und bringen sollte. Denn da wird der Knecht also gesandt, daß er Gold, Silber und andern Schmuck mitgenommen, dazu etliche Kameele und ehrliche Gesellschaft mit sich gehabt hat. Warum wird denn sein Sohn Jakob nicht auch dermaßen gerüstet, so er doch auch nach empfangener Verheißung und Segen zum Erben eingesetzt worden? Er muß zu Fuße gehen ohne Gefährten, ohne Kameele, ohne alle Kosten, er hat kaum Zehrung auf den Weg. Darum soll man diese Art zu reden in der hebräischen Sprache wohl merken um der Vergleichung willen. Er hat seine Füße aufgehoben, er ist zu Fuße gegangen, und hat sich auf dem Wege nicht gesäumt, sondern ist eilends fortgezogen, „endlich“, wie von der Jungfrau Maria gesagt wird Luc. 1, 39., da sie über das Gebirge gegangen ist etc.

19. Also malt Mose den heiligen Patriarchen Jakob und den Erben des Segens ganz verächtlich und niedrig ab, welcher doch jetzt gleichsam allein Papst in der Welt ist. Denn Er ist alt geworden, der noch das vorige Jahr gelebet hat; Isaac ist blind und gar untüchtig zum Regimente; sein Sohn Jakob ist allein der Bischof und Fürst im Hause; nun läßt ihn aber mit höchster Verachtung, mit Gefahr und großem Elend hinweg ziehen, also daß er zu Fuße und ohne Gefährten davon läuft. Es hat zwar an Knechten, Kameelen, Gold und Silber nicht gemangelt: er muß aber so gar elend und arm sein, und muß nichts im Beutel haben, ohne allein kümmerlich die Zehrung auf den Weg.

20. Was ist aber die Ursache? Antwort: Dieses ist die Ursache, die wir droben Cap. 27, V. 43. ff. gehört haben; Rebekka, seine Mutter, hat zu ihm gesagt: „Höre meine Stimme, mein Sohn: Mach dich auf, und fleuch.“ Da siehst du die Ursache, nämlich, daß er hat fliehen müssen, und daß sein Bruder Esau zornig war und gedroht hatte, ihn zu erwürgen. Darum hat Rebekka weislich gethan, daß sie ihn heimlich gerüstet und mit Zehrung versehen hat, auf daß er ohne des Bruders und des ganzen Gesindes Wissen heimlich davon ziehen möchte. Er hat sich müssen heimlich wegstehlen. Schweige still, wird sie gesagt haben, und gehe heimlich hinweg, fliehe so schnell, als du immer kannst. Solche Furcht und Gefahr hat nicht leiden wollen, daß er sich auf die Reise öffentlich gerüstet und viel Gefährten hätte mögen mitnehmen, auf daß er nicht aufgehalten und gehindert würde, weil man mit dem Vieh und Kameelen langsam ziehen muß, und also von seinem zornigen Bruder wäre erhascht worden.
\quellseite{130}

21. Die Juden bringen ohne Ursache an diesem Ort ihre Lügen und unnützen Geschwätze herbei, daß sie vorgeben, es sei nicht wohl glaublich, daß Isaak seinen Sohn Jakob mit weniger Ehre und Herrlichkeit abgefertigt haben sollte, denn Abraham droben, Cap. 24. V. 10., seinen Knecht abgefertigt hat: darum dichten sie, daß Eliphas, der ältere Sohn Esaus, mit einem Haufen gerüsteter Männer dem Jakob nachgefolgt sei, und da er ihn ergriffen, ihn alles seines Gutes und seiner Habe beraubt haben solle.

22. Das läßt sich Lyra gefallen. Aber es sind lauter jüdische Fabeln und Lügen. Denn Jakob ist von seinen Eltern nicht dermaßen abgefertigt worden, daß er hätte Gefährten mit sich gehabt, wie Abraham's Knecht, der da sicher und ohne alle Gefahr war, und nicht hat fliehen müssen. Jakob aber flieht heimlich davon, da alle seine Schwäger nichts darum wissen; dazu auch ohne Vorwissen seines Bruders, von dem er keinen Abschied genommen, der ihm feind war und ihn verfolgte. Darum sagt der Text: „Er hub seine Füße auf“ und floh davon etc.

23. Und es ist wahrlich ein sehr kluger Rath seiner Mutter gewesen, die sich mehr darum bekümmert und bemüht hat, wie der Sohn möchte am Leben und unverletzt bleiben, denn daß sie ihn mit großem Geschenke und Gefährten gerüstet und abgefertigt haben sollte, damit sie ihn nicht dadurch ohne Noth in Gefahr setzte. Und Jakob gehorcht der Mutter sein, widerstrebt ihr nicht, säumet sich auch nicht lange, sondern zieht immer fort und eilt auf der Reise.

24. Dieses, wie wir droben gesagt haben, soll man hart treiben wider die verwegenen und frevelen Köpfe, die da sagen: Wenn ich die Verheißung habe, so habe ich auch dasjenige, so mir verheißen ist, wenn ich schon dazu gar nichts thue. Diese Gedanken soll man verwerfen und verdammen mit diesem Exempel Jakobs, der seiner Mutter gehorcht und flieht, ob er wohl die Verheißung hat. Er sagt nicht: Ich habe die Verheißung, darum werde ich sicher sein und unverletzt bleiben, wenn ich schon dem Rath meiner Mutter nicht folge.

25. Denn die Verheißungen werden nicht darum gegeben, daß wir dabei schlafen, faulenzen und schlafen sollen, oder daß wir thun sollen, das frevels wider die Verheißung ist; sondern wir sollen dabei arbeiten, wacker und wachsam sein, und Frucht bringen. Also werde ich nicht darum getauft, gebrauche auch das Abendmahl des Herrn nicht darum und werde darum auch nicht absolvirt, daß ich dabei faul und schnarchen und faulenzen soll; sondern wenn du die Verheißung, die Taufe und Absolution hast, so bedenke, daß du berufen bist, daß du wachen, und dasjenige, so zum Glauben und deinem Beruf gehört, fleißig und mit großen Sorgen ausrichten sollst. „Wie sollten wir in der Sünde wollen leben“, sagt St. Paulus Röm. 6, 2., „der wir abgestorben sind?“ Wir werden nicht darum von unsern Sünden absolvirt und losgesprochen, daß wir hinfort darin leben und denselben dienen sollten; sondern daß wir ihnen widerstreben und fest in der Verheißung beharren sollen, daß ich mein Fleisch kästen und tödten, und mit Geduld leiden soll, wenn mir Gott das Kreuz auflegt, auf daß wir gesegnet werden und reichere Frucht bringen. „Darinnen“, spricht Christus Joh. 15, 8., „wird mein himmlischer Vater geehrt, wenn ihr meine Jünger werdet“, das ist, wenn ihr euch leidet, wie ich gelitten habe, wenn ihr mir gleichförmig werdet. Denn wer sein crucianus ist (daß ich also reden möge), der ist auch kein christianus, das ist, wer nicht sein Kreuz trägt, der ist auch kein Christ; denn er ist nicht gleichförmig seinem Meister Christo.

26. Das ist nun das Erste, daß dieser sehr große Patriarch, der dazumal allein der Bischof und ein brennend Licht in der ganzen Welt war, der den Segen, die Verheißung und das Wort hat, doch also lebt und wandelt, als hätte er gar nichts. Er, warum das? Antwort: Darum, daß daselbst nicht gar ein geistlicher Mensch ist, sondern ein Hausvater, und der recht elend ist, der gemeine Hausarbeit thut, davon Gott in seinen Verheißungen nichts Sonderliches gesagt oder bestimmt hat; wie er auch nicht gesagt hat, auf welcherlei Weise oder Maße er ihm helfen und ihn regieren wolle, auch nicht, was der Ausgang sein soll: wie er uns auch keine Verheißung hat gegeben, daß eben dieses Jahr soll Friede werden und das Getreide wohl gerathen.

27. Derhalben soll ich nicht sagen: Ich weiß nicht, was werden will, darum will ich auch nichts thun. Ja, Gott spricht vielmehr also: Thue du, was dir in deinem Amte gebühret zu thun, und laß mich das Andere machen. Er hat nicht gesagt: Es wird alles gerathen; sondern du sollst thun, was dir Amts halben gebühret zu thun; du sollst nicht wissen, wie es gerathen, oder was geschehen werde. Du bist gerechtfertigt; derhalben gehe nun hin und übe deinen Glauben mit den Werken des Haus- und Weltregimentes. Und für diese Erkenntniß des Willens Gottes, und für den Beruf, dazu ein Jeder berufen ist, soll man Gott danken, daß ein geistlicher Mensch, das ist, der das Wort und Glauben hat, weiß, daß er Gott wohlgefällt auch in dem niedrigen Grade, im Welt- und Hausregimente, er sei ein Knecht oder eine Magd, ein Regent oder Unterthan: wenn er nur ein Gliedmaß im Welt- und Hausregiment sein mag, so soll er Gott danken, und wissen, daß er einen gnädigen Gott hat, der ihm gütig und wohl gewogen ist.

28. Darnach dient dies auch uns zum Troste, wenn uns die Exempel der sehr heiligen Patriarchen vorgehalten werden, nicht allein was die hohen und heroischen Tugenden belanget, sondern auch in den schlechten und geringen Werken, ja, in den unschätzigen und verächtlichen Werken der Haushaltung; auf daß wir nicht verzagen oder denken, wenn wir mit diesen häuslichen Werken beladen werden, daß wir alsdann von Gott verachtet und verschmähet seien; sondern wir sollen wissen, daß Alles durch das Wort und den Glauben geheiliget werde.

29. Wiewohl die Welt diese Heiligkeit nicht sieht, sondern wo sie hört, daß von solchen schlechten gemeinen Werken erzählt wird, daß auch die allerheiligsten Menschen sie gethan haben, so meint sie, daß es eitel verlorne Arbeit sei, und daß man die gute Zeit damit übel zubringt, wenn man solche Legenden lieset. Denn sie ist nicht werth, daß sie die Ehre und Herrlichkeit Gottes sehen soll; wie geschrieben steht: Hinweg mit dem Gottlosen, daß er die Herrlichkeit Gottes nicht sehe. Es sehen allein die Gläubigen und verstehen die Werke Gottes; darum sind dies vor unsern, ja, vor Gottes Augen köstliche Werke.

30. Und du darfst nicht denken oder dich dessen verwundern, warum der Heilige Geist Lust habe, solche schlechte und verächtliche Werke zu beschreiben; sondern höre, was der heilige Paulus sagt Röm. 15, 4.: „Was aber zuvor geschrieben ist, das ist uns zur Lehre geschrieben, auf daß wir durch Geduld und Trost der Schrift Hoffnung haben.“ Wenn wir fest glauben, wie ich, wiewohl ich schwächlich glaube, daß der Heilige Geist selbst, und Gott, der Schöpfer aller Dinge, der rechte Meister dieses Buchs sei, und solcher schlechten verächtlichen Dinge, wie sie dem Fleische schlecht und gering zu sein dünken: alsdann würden wir den größten Trost davon haben, wie St. Paulus sagt; ja, unser Herz würde sich dessen rühmen und damit hoffärtig sein können, daß sich Gott dazu herabläßt, daß er dieser Patriarchen noch also gedenkt und daß er ihrer nicht vergessen hat, und nicht allein haben wollen, daß ihre ritterlichen Tugenden sollten gerühmet werden, sondern auch die unschätzigen geringen Werke, welche mit dieser Beschreibung wie mit Gold und Edelsteinen sollten geziert und der ganzen Welt vorgehalten werden, daß sie also ausgebreitet, gelesen und erkannt würden. Darum dienen den Gläubigen alle Dinge zum Besten; ja, was auch zuvor gering, unschätzig und verachtet war, muß ihnen zur Ehre dienen, dieweil sie sehen, daß Gott zur solchen geringen Dingen Lust hat, daß er nicht allein singen und rühmen will von den großen und vortrefflichen Tugenden der Heiligen, sondern auch von den geringsten Werklein. Denn es sind Gottes Werke; darum hat Gott Lust an seinen Werken, sie seien gleich groß oder klein: und dasselbe uns zu großem Troste, wie der Spruch lautet im 147. Psalm V. 11.: „Der Herr hat Gefallen an denen, die ihn fürchten“ etc.; item, wie David im 56. Psalm V. 9. sagt: „Zähle meine Flucht, fasse meine Thränen in deinen Sack. Ohne Zweifel, du zählest sie.“

31. Wie nun? Hat denn Gott nichts Anderes
\quellseite{131}
zu thun, denn daß er die Thränen und die Flucht Davids zählen muß? Hat er nicht sonst zu schaffen genug, daß er die Welt regiere? item, den Chören der Engel zuhöre, die ihn loben und preisen ohne Aufhören? Was kann doch Wunderbareres gesagt werden? Und ist doch gleichwohl die Wahrheit, und Gott sorgt dafür auch, daß er die Thränen und die Flucht Davids zähle; wie im 9. Psalm V. 13. steht: „Denn er gedenkt und fragt nach ihrem Blut; er vergisset nicht des Schreiens der Armen“; item, im 116. Psalm V. 15.: „Der Tod seiner Heiligen ist werth gehalten vor dem Herrn.“ Ja, Mose spricht zu Pharao 2 Mos. 10, 26.: „Unser Vieh soll mit uns ziehen, und nicht eine Klaue dahinten bleiben.“ Es sollen nicht allein die Männer, Weiber, Kinder und das Vieh aus Egypten ziehen, sondern alles, was wir haben, auch die allergeringsten Klauen wollen wir nicht dahinten bleiben lassen.

32. Derhalben, sage ich, sind nicht allein die großen ritterlichen Tugenden vor dem Herrn köstlich und werth, dazu auch die großen herrlichen Werke, die er durch uns wirkt, desgleichen das Blut, der Tod und die großen schweren Kämpfe der Heiligen; sondern die allergeringsten Klauen sind auch werth vor ihm. Ja, höre Christum, der macht es noch besser, Matth. 10, 30.: „Nun aber sind auch eure Haare auf dem Haupte alle gezählet“, als wollte er sagen, ihr sollt nicht ein Haar verlieren. Lieber, was ist doch geringer und verächtlicher am ganzen Leibe des Menschen denn ein Haar, oder ein Nagel? Und dieselbigen sind doch auch gezählt und der Vater im Himmel sorgt dafür.

33. Auf solche Weise soll man handeln von den Exempeln der geringen und schlechten Werke der Heiligen, auf daß wir daraus gelehret und getröstet werden mögen. Denn wir lernen nicht, daß sie gering und verächtlich sind, wenn wir nur das Werkzeug ansehen, dadurch sie geschehen: man muß aber auf Gott sehen, der selbst der rechte Thäter ist. Denn du siehst gleich die Unreinigkeit oder aber die Edelsteine an den Heiligen und ihren Werken, so gefallen sie doch Gott wohl, als der beide, die geringen und auch die schönen trefflichen Werke, durch sie wirkt; denn es sind Gottes Werke und Gott wirkt mit ihnen. Derhalben ist dieses ein großer unermesslicher Trost denen, die da glauben. Und dieses alles wird darum beschrieben, daß wir sehen sollen, wie gar herzlich uns Gott liebt und wie wohl er uns gewogen ist, wie große ängstliche Sorge er für uns trage. So genau sieht er auf mich, daß er Sorge hat, ich verliere ein Haar. Wo er nun die Haare zählt und dafür sorgt, so sorgt er ja viel mehr für meinen Leib, Seele, Blut und alles Leiden, so mir begegnen mag.

34. Diese Dinge sind aber allzu hoch, darum glauben wir es nicht, und je geringer und unfläthiger die Werke sind, je weniger wir auch glauben. Die Unreinigkeit der Werke und daß sie so gar schlecht und gering sind, hindert uns am Glauben; sonst würden wir die Gnade und Barmherzigkeit Gottes in diesen sehr schlechten und geringen Werken über die Maßen hoch heben und groß machen, und unser Glaube würde daraus sehr gestärket werden. Denn so große Sorge und Bekümmerniß hat Gott für uns, daß er auch nicht eines Härleins vergessen kann, auch nicht eines Thränleins oder kleinen Bekümmernüsses. Und das meint der Heilige Geist, wenn er so gar niedrig eingehet, wo er seine Heiligen beschreibt, daß auch die allergeringsten Werke der Heiligen Gott wohlgefallen. Es ist ein köstlich Ding um einen Christenmenschen: es ist nichts so gering an ihm, es gefällt Gott. Blut vergießen, sterben, schwitzen, streiten und kämpfen wider den Teufel ist in Wahrheit ein groß Ding und gefällt Gott sehr wohl. Du mußt aber also schließen: Wenn du nun gläubig bist, alsdann gefallen Gott auch die natürlichen, fleischlichen und leiblichen Werke wohl, du essest oder trinkest, du wachest oder schläfest; welches ja lauter leibliche und natürliche Werke sind. Ein so groß Ding und Sache ist der Glaube.

35. Derhalben siehe zu allererst darauf, daß du ein Christ werdest, und daß die Person durch das Wort, durch die Taufe und Sacramente Gott angenehm und gefällig werde. Wenn die Person gläubig ist und dem Worte anhängt, dasselbe nicht verfolgt, sondern Gott dafür dankt: alsdann sollst du Anderes nichts thun, denn das Salomo sagt in seinem Prediger am 9. Capitel V. 7. 8. 9. 10.: „So gehe hin und iß dein Brod mit Freuden, trink deinen Wein mit gutem Muth; denn dein Werk gefällt Gott. Laß deine Kleider immer weiß sein, und laß bei deinem Haupte Salbe nicht mangeln. Gebrauche des Lebens mit deinem Weibe, das du lieb hast, so lange du das eitle Leben hast, das dir Gott unter der Sonne gegeben hat, so lange dein eitel Leben währet, denn das ist dein Theil im Leben und in deiner Arbeit, die du thust unter der Sonne. Alles, was dir vorhanden kommt, das thue frisch“ etc. Was willst du mehr haben? Und was könnte lieblicher, freundlicher und klarer gesagt werden?

36. Das ist wohl wahr, daß Gott auch an den Gottlosen wohlgefällt, daß sie fleißig und geschickt sind, ihr Amt auszurichten; es liegt ihnen aber der Unglaube und die eitle Ehre im Wege, daß sie ihre Werke nicht auf Gottes Ehre richten können: der Mangel und die Sünde ist an der Person, welche Gott nicht gefällt. Derhalben verdienen auch wohl die Gottlosen mit ihren guten Werken Belohnung in diesem Leben; es werden aber ihre Werke nicht gezählt, unser HErr Gott faßt sie nicht in seinen Sack, wie der 56. Psalm V. 9. sagt von den Thränen der Gottseligen. Aber die Thränen, die Flucht, Anfechtung, beide kleine und große Werke der Heiligen, die werden gezählt, und sollen in Ewigkeit gelobt und gerühmt werden.

37. Derhalben ist das ein sehr lieblicher Trost, welchen der heilige Paulus anzeigt, da er Röm. 15, 4. sagt, daß in der heiligen Schrift nichts vergeblich geschrieben. Denn es ist gewiß, daß die sehr niedrigen und sehr geringen Werke uns darum vorgestellet werden, auf daß damit angezeigt werde, daß Gott an seinen Heiligen ein Gefallen habe; denn es bleibt allezeit die Vergebung der Sünde und daß sie Gott zu Gnaden annimmt; und so lange sie in solcher Gnade sind, leben sie immer unter der Wolke und Schatten seiner Flügel, und Gott schützt und beschirmt sie.

38. Und dieses sollen wir uns auch zueignen. Denn wo wir Christen und rechte gottselige Leute sind, so wissen wir, daß wir den großen Heiligen gleich sind, wo nicht in dem höchsten Grade der großen Tugenden, so sind wir ihnen doch gleich in diesen niedrigen und unflätigen Werken dieses Lebens, und wissen, so viel die Sorge und den Schutz Gottes belanget, daß wir nicht weniger von ihm geliebt werden, als sie geliebt worden sind. Und derselbigen zarten und brünstigen Liebe Gottes haben wir ein gewisses Pfand, nämlich, seinen lieben Sohn, um deß willen uns der Vater liebt, und sammt Christo in das himmlische Wesen gesetzt hat, wie St. Paulus sagt Eph. 2, 4. 6. Zugleich sollen wir wissen, daß uns das auch angeht, so im 37. Psalm V. 23. steht: „Von dem Herrn wird solches Mannes Gang gefördert, und hat Lust an seinen Wegen. Fällt er, so wird er nicht weggeworfen; denn der Herr erhält ihn bei der Hand“ etc. Allein lasset uns glauben und diese Verheißung annehmen. Denn wie die Eltern ihre Kindlein sehr fleißig und sorgfältig bewahren, daß sie nicht irgend fallen, sich auch nirgend anstoßen oder verletzet werden mögen; und wenn sie sehen, daß sie einen Flecken oder Schmuz im Gesicht haben, sie flugs da sind, und es ihnen abwischen und sie rein machen, welches ein Feind und Fremder nicht thut; wo ihnen irgend eine Feder an den Haaren hängt, so kämmen und schmücken sie die Kindlein, und halten sie sauber und schön: solche große Sorge, und solche Liebe und rechte Vaterschaft trägt Gott zu Allen, die an ihn glauben.

39. Dieses sei also in einer Summa von der Lehre in diesem Capitel gesagt. Was nun hernach folgt, gehört schlechterdings zum Hausregimente: das wollen wir auch überlaufen, und was zur Grammatik gehört, erklären. „Er ging“, sagt der Text, „in das Land der Kinder des Morgens“, zu Deutsch: in das Land der Morgenländer, oder: „das gegen Morgen liegt.“ Dieses ist eine gemeine Art und Weise zu reden in der hebräischen Sprache. Denn also wird auch gesagt, ein Kind des Todes, nämlich, den man tödten soll; item, ein Kind des Rächers, das ist, ein Pfeil; ein Kind der Hofzucht, das ist, ein wohlfähriger Mensch. Ein Kind der Stärke, ist so viel gesagt, als, seid stark und unverzagt. Ein Kind Belials heißt ein böser Schalk.

40. So ist nun der sehr heilige Patriarch Jakob in das Land, das gegen Morgen liegt, gezogen: mit welchem gemeinen Namen des ganzen Landes die Schrift heimlich anzeigt, daß wiewohl diese Reise sich leicht und gering ansehen lasse, so sei es doch ein hohes und schweres Werk gewesen, so aus dem Worte und Glauben hergekommen ist. Denn er hat in einem unbekannten Lande umherziehen müssen, dessen Name auch nicht ausgedrückt wird. Denn er nennt es nicht Haran, sondern das Land gegen Morgen, welches die Araber und Syrer damals noch nicht gekannt oder bewohnt haben, da ist nichts gewesen, ausgenommen Babel. Darum sagt Mose einfach, er sei in das Land gegen Morgen gegangen, das ist, in eine unbekannte Wüste. Wie droben (Cap. 11. V. 31.) Abraham aus Chaldäa gezogen ist, daß er selbst nicht gewußt, wo er hinging oder wo er seinen Fuß hinsetzen sollte: er wagte es auf unsern Herrn GOtt.
\quellseite{132}

41. Derhalben ist wohl die Reise leicht und gering anzusehen gewesen; aber in Wahrheit ist es ein Werck eines großen brünstigen Glaubens. Denn er weiß nicht, an welchen wüsten Ort der Leute oder Städte er kommen werde; er kennt die Leute nicht, weiß weder die Städte noch den Weg, dahin er gehen soll; er weiß nicht, wo er eine jede Nacht gelegene Herberge haben werde. Also ist es über die Maßen eine gefährliche und beschwerliche Reise gewesen, und er hat unterdessen die ganze Reise über im Glauben mit dem Tode und Teufel gekämpft, wußte nicht, wohin er gehen sollte, oder wie es ihm auf der Reise gehen würde, ehe denn er an den Ort kommen würde, dahin er ziehen sollte.

42. Ohngefähr aber sieht er einen Brunnen auf dem Felde, zu demselben hat er sich begeben. Das ist nun auch schlecht und gering Ding; allein, daß es sich allhier anhebt, daß ihm der Trost gezeigt wird, da er an den Ort kommt, an welchem er nun etwas sicherer ist denn zuvor, und höret über das auch sagen von seinem Vetter Laban.

43. Hier weiß ich nun nichts mehr zu sagen, darum will ich nichts hinzu thun. Denn zu dem Allegorien oder der heimlichen Deutung habe ich nicht Lust; denn der Text an sich selber nach dem Buchstaben gibt uns bessern Bericht, denn die Allegorie thun kann. Es ist daran genug, daß wir wissen, daß diese schlechten und geringen Wercke der Heiligen GOtt wohlgefallen; daher wir auch nicht zweifeln, wo wir in unserm Berufe sind, was wir thun oder leiden werden, solches gefalle GOtt auch wohl, also auch, daß er dasselbe in dieses Buch hat wollen schreiben lassen, und haben wollen, daß alle diese Wercke der Heiligen dermaßen sollten gerühmt werden. Er hat keinesweges leiden wollen, daß solche Wercke sollten vergessen werden. Der Text ist nach dem Buchstaben an sich selber klar und deutlich genug. Nach der Allegorie wollen sie durch die drei Heerden Schafe bei dem Einen Brunnen die Dreifaltigkeit verstehen und durch den großen und schweren Stein verstehen sie Mose. Aber das geht mich nichts an, ich will dasselbe Andern lassen befohlen sein.

\section[Zweiter Theil]{Zweiter Theil. Von Jakobs Ankunft in Haran, wie er von Laban freundlich auf- und in seinen Dienst aufgenommen wird, und was Jakob für einen Lohn verlangt}

\subsection{I}

\begin{quote}
\textbf{V. 4–8. Und Jakob sprach zu ihnen: Lieben Brüder, wo seid ihr her? Sie antworteten: Wir sind von Haran. Er sprach zu ihnen: Kennet ihr auch Laban, den Sohn Nahors? Sie antworteten: Wir kennen ihn wohl. Er sprach: Gehet es ihm auch wohl? Sie antworteten: Es gehet ihm wohl; und siehe, da kommt seine Tochter Rachel mit den Schafen. Er sprach: Es ist noch hoch Tag, und ist noch nicht Zeit das Vieh einzutreiben; tränket die Schafe, und gehet hin, und weidet sie. Sie antworteten: Wir können nicht, bis daß alle Heerden zusammen gebracht werden, und wir den Stein von des Brunnens Loch wälzen, und also die Schafe tränken.}
\end{quote}

44. Haran ist eine Stadt in Mesopotamien gelegen, da hernachmals Crassus eine große Niederlage erlitten und von den Parthern überwunden worden ist. Laban aber ist Nahors Sohnes Kind gewesen, und nicht sein Sohn: und derselbe war nun gestorben. Bethuel, der Sohn Nahors, war Labans und Rebekkas Vater. Nahor wird aber darum genannt, denn er ist der Vater und das Haupt desselbigen Geschlechtes. Jakob aber fragt: Hat er Friede? Das ist nach Art der hebräischen Sprache geredet, als wenn wir sagen: Ist er auch noch gesund und stark?

45. Von dem Steine vor dem Loche des Brunnens erheben die Grammatikter eine Frage: Ob das ein Gesetz gewesen sei, daß Einer allein den Stein nicht hätte abwälzen dürfen und seine Heerde tränken, ehe denn sie alle zusammen kämen? oder aber: Ob der Stein so groß und schwer gewesen sei, daß vonnöthen gewesen, daß ihrer Zwei oder Drei einander haben helfen müssen den Stein abzuwälzen? Ich halte es zwar nicht dafür, daß ein sonderlich Gebot gewesen sei, darin es verboten gewesen wäre, sondern daß der Stein so groß und schwer gewesen sei, daß ihn Einer oder Zwei nicht haben können abwälzen.

46. Man hat aber in Mesopotamien die Brunnen mit Fleiß verwahret und mit großen Steinen zugedeckt, dieweil man im Lande nicht viel Wasser gehabt. Darnach sind auch die Hirten, so die Heerde geweidet, nicht alle starke Männer gewesen, sondern man hat Knaben und junge Mädchen, wie Rebekka war, dazu genommen, denen man die Heerde zu weiden und zu bewahren befohlen hat. Dieselbigen sind bei dem Brunnen zusammen gekommen, und haben gewartet, bis Stärkere kommen würden, die den Stein könnten abwälzen; wie Jakob hinzu tritt und wälzt den Stein allein ab. Dieses halte ich für die Ursache, darum ihrer Viele haben müssen zusammen kommen und einander helfen. Was aber Andere halten, lasse ich ihnen frei stehen.

\begin{quote}
\textbf{V. 9–12. Als er noch mit ihnen redete, kam Rachel mit den Schafen ihres Vaters, denn sie hütete der Schafe. Da aber Jakob sahe Rachel, die Tochter Labans, seiner Mutter Bruders, und die Schafe Labans, seiner Mutter Bruders; trat er hinzu, und wälzte den Stein von dem Loch des Brunnens, und tränkte die Schafe Labans, seiner Mutter Bruders; und küssete Rachel, und weinete laut, und sagte ihr an, daß er ihres Vaters Bruder wäre, und Rebekkas Sohn. Da lief sie, und sagte es ihrem Vater an.}
\end{quote}

47. In den Commentaren der Hebräer ist angezeigt worden, und es läßt sich ansehen, daß der Text eben dasselbe auch mitbringe, nämlich, daß der Heilige Geist auf den Patriarchen Jakob gefallen, und er davon also gestärkt worden sei, daß er durch solchen Trieb und Anregung des Heiligen Geistes den Stein allein habe können abwälzen. Darnach, dieweil er daselbst in seinem Elende Rachel, seine Blutsfreundin, gefunden, faßt er Muth und hat große Hoffnung bekommen, daß er endlich erlangen werde, das er im Sinne gehabt und das ihm sein Vater befohlen hatte, daselbst ein Weib zu nehmen. Darum sobald er Rachel sahe, wird er von Liebe entbrannt, und die natürliche Neigung gegen seine Blutsfreundin thut sich hervor, also daß zweierlei Triebe, des Glaubens und der Liebe, sein Leib und Herz etwas stärker gemacht. Denn er hat sich wollen männlich erzeigen mit Stärke und daß er hurtig wäre, auf daß er damit das Herz der Jungfrau möchte einnehmen und sie anreizen, ihn lieb zu gewinnen.

48. Und dieses sind auch nur natürliche Dinge; es hat sie aber der Heilige Geist darum schreiben lassen, daß niemand denken möchte, daß es schändliche oder ungebührliche Dinge wären. Denn es ist ein christlich und gottselig Ding, daß du eine Jungfrau lieb habest, die du mit Ehren zum Weibe nehmen mögest. Denn solches ist eine natürliche Neigung, so die beiden Geschlechter zu einander haben; und wiewohl nun dieselbige Neigung nicht gar ohne Sünde ist, so will doch GOtt nicht haben, daß man sie verachten soll, als ob sie an sich selber unehrlich wäre. Denn es ist ein Werck GOttes, so er in der Natur des Menschen geschaffen hat, welches man nicht allein nicht verachten oder lästern soll, sondern man soll es auch in Ehren halten. Denn GOtt will in allen Werken, beide klein und groß, geehrt sein. Er will nicht verachtet sein, auch nicht in dem allergeringsten Wercke; wie Christus spricht Matth. 25, 40.: „Was ihr gethan habt einem unter diesen meinen geringsten Brüdern, das habt ihr mir gethan.“ Also sollen wir auch die allergeringsten Gebote und Ordnungen GOttes nicht verwerfen oder unter die Füße treten. Willst du Wasser trinken, so ist es gut, so trinke es und danke GOtt dafür: gelüstet dich es aber nicht zu trinken, sollt du es gleichwohl nicht verachten.

49. So sind nun, sage ich, dieses fleischliche, sündliche und weibische, lauter verächtliche und unreine Dinge, und werden doch vom Heiligen Geiste erzählt; zwar nicht, die Leute damit zur Unzucht zu reizen, nicht, schändliche und unziemliche Liebe damit in den Herzen der Jünglinge anzuzünden; sondern die Hoffnung und Würde des Ehestandes damit zu erhalten, und daß der Heilige Geist anzeige und bezeuge, daß der Ehestand GOtt nicht mißfällt und dazu die Liebe des Bräutigams und der Braut auch nicht; sintemal GOtt diesen Stand geschaffen und eingesetzt hat. Derhalben sollst du ihn nicht schmähen oder lästern, als ob es ein unehrlicher und unziemlicher Stand wäre; wie die Mönche gesagt haben, daß der Ehestand ein gefährlicher und unreiner Stand sei.
\quellseite{133}

50. Derhalben halte ich dafür, daß Jakob in der Stunde, da er seine Blutsfreundin angesehen, etlichermaßen entzückt und also stark im Geist geworden sei, daß er den Stein allein hat abwälzen können, dazu man sonst vier oder fünf starker Knaben bedurft hätte. Und er hat der Jungfrau diese Freundschaft und Dienst thun wollen, daß sie ihre Heerde in der Zeit tränken konnte. Sie aber weiß noch nicht, daß er ihr verwandt ist.

51. Da er nun den Stein abgewälzt, ist er hinzu getreten und hat die Jungfrau geküßt als ein Unbekannter und Fremder, ehe denn er sie gegrüßt. Welches auch für unehrlich angesehen wird und jetziger Zeit nicht gebräuchlich ist, und steht auch nicht wohl: er hat aber zuvor von den andern Hirten gehört, die da sagten: „Siehe, da kommt seine Tochter Rahel“ etc. Darum hat er gewußt, daß sie seine Blutsverwandte war. Und dieselbigen Leute haben den Gebrauch gehabt, daß wenn sie zusammen gekommen und einander gegrüßt, sie auch einander geküßt haben; wie dasselbe noch heutiges Tages in Belgien und auch in andern Ländern gebräuchlich ist. Unsere Leute geben einander nur die Hand und nehmen die Jungfrauen und ehrlichen Matronen fein züchtig in den Arm. So manches Land, so manche Sitten. Derhalben ist dieses ein Anzeichen der Liebe gegen seine angeborne Freundin und sein künftiges Eheweib, welche Liebe er mit diesem Kuß beweisen will, gleich als hätte er ihr die Hand oder sonst irgend ein Geschenk gegeben.

52. Und daher ist die Stärke gekommen, den Stein abzuwälzen, nämlich von der natürlichen Neigung, so er zu seiner Blutsfreundin gehabt. Darnach auch von der Liebe, so er zu seiner Braut gehabt, welche, wie er gesehen hat, ihm von Gott zugeschickt und gegeben war. Zum Dritten ist er auch darüber sonderlich fröhlich geworden, darum daß er gesehen hat, wie er nun nach solcher schweren und gefährlichen Reise den Hafen erlangt hätte. Dieses aber alles hat der heilige Geist in ihn gemehret und bestätiget. Zuletzt wird Jakob ein Bruder genannt des Vaters der Rahel nach hebräischer Weise; denn die Hebräer nennen der Schwester Sohn einen Bruder, wie das Mose selbst erklären wird, da er sagen wird: daß Laban von Jakob, seiner Schwester Sohn, gehört habe etc.

\begin{quote}
\textbf{V. 13. Da aber Laban hörte von Jakob, seiner Schwester Sohn, lief er ihm entgegen, und herzte und küssete ihn, und führete ihn in sein Haus. Da erzählte er dem Laban alle diese Sache.}
\end{quote}

53. Allhier siehst du, daß diese Dinge ganz und gar also beschrieben sind, wie man pflegt von Hirten, wie sich dieselben mit einander unterreden, zu schreiben. Und der Heilige Geist macht doch gleichwohl so viel Worte davon, daß daran der Leser, der ein Weltmensch ist, einen Ekel und Ueberdruß hat. Andere Historien von großen wichtigen Dingen, als, von der Zerstörung Sodoms und Gomorras, item, daß Abraham den Isaac geopfert, hat er mit wenig Worten kaum in fünf oder sechs Versen zusammengefaßt; wenn er aber auf diese unschlächtigen, fleischlichen und närrischen Dinge kommt, so macht er über die Maßen viel Worte davon, nämlich darum, daß wir sollen wissen, daß der Herr Gefallen hat an denen, die ihn fürchten. Denn wo wir gläubig und der Gnade Gottes gewiß sind, so er uns umsonst erzeigt hat, so dürfen wir daran nicht zweifeln, daß nicht alles, was wir thun, Gott sehr wohl gefallen sollte, und daß er auch die Haare auf unserm Haupt gezählt hat, ja, daß ihm wohlgefällt das Küssen und Herzen, und daß Jakob den Stein abwälzt; welches alles an diesem Ort erzählt wird als sehr hohe und köstliche Werke vor den Augen Gottes und der heiligen Engel. Gott hat derselben nicht vergessen können, sondern hat haben wollen, daß sie sollten beschrieben werden uns zur Lehre und Trost.

54. Wiewohl wir nun in diesem Leben sehr hart geplagt werden, und vielem Jammer und Unglück, endlich auch dem Tode selbst unterworfen sind, sollen wir doch wissen, daß es alles eitel köstlich Ding ist; wie im 116. Psalm V. 15. steht: „Der Tod seiner Heiligen ist werth gehalten vor dem Herrn.“ Dasselbe sollen wir auch von andern Werken verstehen, wir essen oder trinken, schlafen oder wachen; wie St. Paulus Col. 3, 17. sagt: „Alles, was ihr thut mit Worten oder mit Werken, das thut alles in dem Namen Jesu, und danket Gott und dem Vater durch ihn.“ Und in solchen geringen Werken sollen wir uns üben, daß wir allgemach fortkommen und uns an höhere Dinge gewöhnen mögen. Der Glaube aber, der große und ritterliche Dinge wirkt, läßt sich hernieder auch zu den häuslichen und geringen Werken, er weidet die Heerde der Schafe, er tränkt sie, er melkt die Ziegen. Denn Gott hat beide geschaffen, den Thoren und Weisen, den Kleinen und Großen, dazu was köstlich und was gering ist. Derhalben gefallen ihm die Werke des Knechtes oder der Magd eben so wohl als des Herrn, des Weibes so wohl als des Mannes. Sei nur gläubig und halte vor allen Dingen in den Kirchen- und Religionssachen Gott in Ehren; mache es also, daß du bei seinem Wort bleibest. „So ihr in mir bleibet“, spricht Christus, Joh. 15, 7. 8. 10., „und meine Worte in euch bleiben, so werdet ihr viel Frucht bringen, und eure Frucht wird bleiben“ etc.

55. Derhalben soll man auf den Glauben sehen, und auf die Person, die bei Gott in Gnaden ist. Denn Gott hat so große Güte über uns ausgeschüttet, daß er auch unsere allergeringsten Werke zählt und Achtung darauf hat. Es seien gleich welcherlei Werke es sein mögen, so sind sie doch alle löblich und Gott wohlgefällig, dieweil ihm die Person gefällt.

56. Und gilt hier nicht, was die Vernunft davon urtheilt, die da sagt: Es gebühret einem hebräischen Schreiber, solche Dinge zu beschreiben, als die nicht werth sind, daß sie in den Büchern der heiligen Schrift stehen sollten; ja, es sollten wohl die griechischen oder lateinischen Schreiber nicht einmal solche geringe Werke schreiben von ihren großen und tapfern Helden, deren Geschichte sie beschrieben haben. Denn sie verstehen nicht, was ein recht gottselig Leben sei, was Gott gefällig oder nicht gefällig sei; wie der Spruch lautet: Hinweg mit den Gottlosen, daß er die Ehre und Herrlichkeit Gottes nicht sehe.

57. Darum sollen wir nun dieses mit gebührlicher Ehrbietung lesen, und Gott dafür danken, daß wir wissen, daß wir Gott gefallen, auch wenn wir schlafen, essen, trinken, freien; wo Mann, Weib und Kind bei einander wohnen; wenn wir das Gesinde regieren, die Ziegen melken etc. In solchen Werken sollen wir uns üben; denn sie sind eben so groß und wunderbar, als die hohen Werke, darum daß sie von einer großen Person geschehen, die da glaubt und ein Erbe des Himmelreichs ist. Wenn es Gott gefallen wird, so wird auch irgend ein Tag oder Stunde kommen, darin wir auch große ritterliche Werke thun werden.

58. Und diesen seinen guten Willen und Wohlgefallen hat Gott in der Neigung, so die Eltern zu ihren Kindern haben, auch abgemalet. Denn wir sehen im Hausregiment, daß Vater und Mutter mehr Lust und Gefallen haben, wo das Söhnlein oder Töchterlein irgend ein Blümlein, oder sonst etwa ein gering Ding daher bringt, denn wo der Knecht oder die Magd einen großen Sack oder Ballen tragen: die geringen Werke sind an den Kindern lieblich und den Eltern angenehm, welche sonst an den Knechten und Mägdgen würden verachtet werden.

59. Darum ist es viel ein Anderes, die Legenden der Heuchler und Mönche zu beschreiben, als der rechten Heiligen. Jene Heuchler verlachen und verspotten die Werke und das ganze Leben dieser rechten Heiligen als ein unreines Leben, und lassen sich derweil träumen, daß sie Gott mit ihren selbstgemachten und ungeheuren Werken versöhnen können. Aber solche ihre Werke sind von Gott verdammt, dieweil sie nicht geschehen von einer solchen Person, die Gott angenehm und gefällig ist. Derhalben lasset uns nun das Andere auch besehen, welches nicht weniger närrisch und fleischlich ist, wie die Heuchler davon urtheilen. Und folgt nun ein neues Stück, nämlich, wie Jakob ein Bräutigam geworden sei.

\subsection{II}

\begin{quote}
\textbf{V. 14. Da sprach Laban zu ihm: Wohlan, du bist mein Bein und Fleisch. Und er (Jacob) blieb bei ihm einen Mond lang.}
\end{quote}

60. Die Sache läßt sich anfänglich wohl an; wie sich denn die Heuchler anfänglich pflegen zu erzeigen. Denn es wird durch die ganze Historie Laban beschrieben, daß er geizig und eigennützig gewesen sei. Und der Name Laban reimet sich fein damit; denn es heißt so viel, als weiß (albus). Denn die Heuchler haben einen großen Schein der Gottseligkeit, der Zucht und guter Sitten, also daß es sich ansehen läßt, als sei niemand heiliger, ehrlicher und geistlicher, denn sie sind. Er sollte aber billiger Nabal genannt worden sein, daß man das Wort hätte umgekehrt. So stellt er sich nun äußerlich gegen seinem Vetter sehr gottseelig und als habe er ihn sehr lieb. Unterdeß aber denkt er: Siehe, nun habe ich einen sehr guten und nützlichen Knecht bekommen, welcher nach meinem Willen und Gefallen Alles wird thun und leiden müssen, deweil er flüchtig und elend ist, der nirgend einen Ort hat finden können, da er sich hätte mögen sicher aufhalten. Er ist arm, hungrig und nackend: in solchem Gefängniß und Banden, die stark genug sind, soll er in meinem Hause gefangen gehalten werden. Denn wo er ohne mich wäre, so würde er müssen Hungers sterben.
\quellseite{134}

61. Anfänglich aber gibt er ihm gute Worte, liebkoset ihn, wie die Heuchler pflegen zu thun. Du bist, sagt er, mein Bein und mein Fleisch, läufst ihm entgegen, herzt und küßt ihn, gleich als ob er ihn herzlich und von Herzen lieb hätte. Aber er wird bald sein verborgen und untreues Herz offenbaren und an den Tag geben. Wie man sonst auch pflegt zu sagen: Post tres saepe dies visitat piscis et hospes, das ist: Wie man der Fische nicht mehr achtet, wenn man sie drei Tage genossen, also achtet man auch eines Gastes nicht viel mehr, wenn man ihn drei Tage gehabt hat. Man wird der Gäste bald müde und überdrüssig.

62. Darum wird Laban dermaßen beschrieben und abgemalt, daß er ein solcher Heuchler gewesen sei, der einen großen Schein gehabt habe, so viel die äußerlichen Werke und Worte belanget; aber der inwendig voll Abgötterei, Hoffart, Geiz und Verachtung des Nächsten gewesen und nur auf seinen eigenen Nutzen sieht. Mit einem solchen Heuchler muß der fromme heilige Mann Jakob umgehen. Siehe aber doch, wie die große Ehre und Herrlichkeit dieses so großen Patriarchen so gar einen geringen Anfang gehabt hat; von welchem Patriarchen so viele Könige, so viele Propheten, Christus selbst und die Apostel hergekommen sind; bei welchem die ganze Majestät des ewigen Lebens; der ein König und Priester ist, und der den ganzen Segen von seinem Vater erlangt hat: der ist so gar arm und verachtet, daß er auch nicht so viel Raums hat, da er möchte mit einem Fuße auftreten. Ja, das heißt aus Nichts einen Menschen schaffen.

\begin{quote}
\textbf{V. 14. 15. Und da er nun einen Monat lang bei ihm gewesen war, sprach Laban zu Jakob: Wiewohl du mein Bruder bist, solltest du mir darum umsonst dienen?}
\end{quote}

63. Jakob ist einen Monat lang bei Laban gewesen, dieselbige Zeit wird er ihn aber nicht umsonst genähret haben. Denn er wird ohne Zweifel nicht müßig gewesen sein, sondern hat gethan, was einem gottseeligen und getreuen Menschen gebühret zu thun. Und dazu hat ihn nicht allein die Armut und Mangel getrieben, sondern er hat solches gethan, weil der Heilige Geist in ihm war, durch welchen die heiligen Menschen angereizt und getrieben werden, allerlei ehrliche und gottseelige Werke zu thun: welche Menschen Gott auch segnet, daß sie vielen Leuten nütze sein können. Wie Joseph, da er von seinem Herrn nichts hatte, denn nur das Brod und die Kleidung, hat er ihn doch reich gemacht, und hat dagegen keine Belohnung von ihm erwarten dürfen: ja, solche Leute sind es, davon St. Paulus sagt Col. 3, 23.: „Alles, was ihr thut, das thut von Herzen, als dem Herrn und nicht den Menschen“ 2c. Solche Knechte sind fast selten, die also treulich dienen, nicht daß sie den Menschen gefallen mögen, sondern Gott. Jakob aber ist den Monat lang ein solcher Knecht gewesen, und die ganze Zeit über, dieweil er Laban gedienet hat; er ist nicht müßig gewesen, sondern hat selbst willig und mit großem Fleiß die häusliche Arbeit auf sich genommen und dieselbe verrichtet, er hat das Vieh gefüttert und getränket: und alle andere dergleichen Hausarbeit hat er auch gethan, so treulich und fleißig, als er immer gekonnt hat. Und das hat Laban, dem Geizhalse, sehr wohlgefallen, da er gesehen hat, daß er so fleißig und lustig war, solche geringe knechtische Arbeit auszurichten. Darum macht er bald einen Vertrag mit ihm des Lohnes halben, dadurch denn der Fleiß und die Geschicklichkeit Jakobs sammt seiner großen Treue hoch gelobt wird.

64. Und ist wahrlich eine treffliche Tugend, die sehr selten und rar ist, vornehmlich unter den Blutsfreunden, da du ihrer sehr wenig finden wirst, die da denken, wie sie ihren angebornen Freunden treulich dienen mögen. Denn sie lassen sich dünken, sie haben sonderlich Recht dazu, daß sie nur verschwinden, an sich reißen und verschwenden, was ihrer angebornen Freunde ist. Die Schwäger aber oder Blutsfreunde sind fast selten, die dem Jakob und Joseph gleich sind, oder wie Naeman war bei dem Könige in Syrien, dem du befehlen und vertrauen möchtest die Sorge und Verwaltung deines Hauses.

65. Derhalben hat Laban gesehen, daß er einen Knecht bekommen hatte nach seinem Gefallen, der da fleißig und emsig war, der das Vieh wohl hat versorgen können, der das Gesinde fein in Ordnung gebracht und im Hause eine sehr schöne Zucht angerichtet, dabei alles, was im Hause war, gar sehr gemehret und gebessert worden ist. Dasselbe bezeugen die Worte, die hernach folgen werden, da Jakob sagt: „Du hattest wenig, ehe ich herkam“ 2c. Und ist zugleich der Segen Gottes auch dazu gekommen um des frommen gläubigen Knechtes willen, welcher in Wahrheit das Heil, und wie eine Säule war, daran sich das Hausgesinde gehalten hat. Derhalben sagt Laban: Ich will nicht haben, daß du mir umsonst dienen sollt. Denn er sieht wohl, daß er kein faules Vieh ist, wie wir heutiges Tages Knechte und Mägd haben, sondern daß er mehr thut, denn ihm befohlen war, Alles seinem Vetter zu Nutz und Frommen. Darum folgt nun der Lohn für solchen treuen Dienst.

\subsection{III}

\begin{quote}
\textbf{V. 15–20. Sage an, was soll dein Lohn sein? Laban aber hatte zwo Töchter; die älteste hieß Lea, und die jüngste hieß Rahel. Aber Lea hatte ein blödes Gesicht, Rahel war hübsch und schön. Und Jakob gewann die Rahel lieb, und sprach: Ich will dir sieben Jahr um Rahel, deine jüngste Tochter, dienen. Laban antwortete: Es ist besser, ich gebe sie dir, denn einem andern; bleibe bei mir. Also dienete Jakob um Rahel sieben Jahr, und däuchten ihn, als wären es einzelne Tage, so lieb hatte er sie.}
\end{quote}

66. Dieses ist ja nur kindisch und fleischlich Ding. Denn was ist es doch, daß der sehr heilige Mann, der so große Verheißung von Gott hat, dem Gott selbst erschienen ist in einem sehr schönen Gesicht, in solchem Alter noch jung und schier gar wiederum zum Kinde wird? Er sieht zwei Jungfrauen, und daß die eine schöner ist denn die andere; von denselben gewinnt er die lieb, welche die schönste ist von Gestalt, die andere hat er nicht lieb. Dünkt dich nun, daß dem Jakob zu ziemen, der ein achtzigjähriger Mann war? Denn es ist gewiß, daß er dazumal an sein achtzigstes oder ja zum wenigsten an das achtundsiebenzigste Jahr gekommen war. Nun hebt er aber in solchem Alter erst an kindischer Weise zu scherzen und die Jungfrau lieb zu gewinnen. Siehe, warum schreibt doch der Heilige Geist solch Ding, und hält uns dasselbe vor, daß wir es lesen sollen, gleich als ob wir aus solcher schändlichen That könnten gebessert werden? Denn das noch mehr ist, hernach wird im Text gesagt werden, daß Jakob seinen Sohn Joseph geliebt habe, darum daß er ihm in seinem Alter geboren war; dieselbe Zeit ist nicht mehr als vierzehn Jahr gewesen, nachdem dieses geschehen war, und wird dennoch das Alter Jakob genannt.

67. Darum ziemt es sich ja nicht, daß ein solcher heiliger Mann also närrisch und geil sein soll, daß er die Schönheit einer Jungfrau ansehen wollte, und eine schöne Jungfrau lieber haben, denn sonst eine, die da häßlich und nicht schön ist. Denn man kann nicht sagen, daß er es darum gethan hätte, daß er gern hätte Kinder mit ihr zeugen wollen; sintemal die Weiber, so nicht schön sind, gemeiniglich am allermeisten pflegen fruchtbar zu sein, wie das an Lea zu sehen ist: und dieselbe hat doch mit keiner Geburt erlangen können, daß Jakob sie so heftig lieb gehabt hätte, wie er Rahel hatte. Derhalben kann man ihn nicht entschuldigen mit dem Schein, den man möchte vorwenden, als hätte er es um der Kinderzeugung willen gethan; sondern dieses ist ein gar schimpflich und lächerlich Ding, nämlich, eine Jungfrau lieb haben um desselben, daß sie schön ist, und nicht vornehmlich um der Kinderzucht willen.

68. Das sollte zwar billig geschehen, daß wir das weibliche Geschlecht an sich selber lieb hätten, nur um der Kinderzeugung willen, dazu dieses Geschlecht von Gott geschaffen ist, und nicht sein allein zur Lust und Unzucht mißbrauchen. Denn daß die Weiber dazu geschaffen sind, daß sie Kinder gebären sollen, zeigt die Gestalt ihres ganzen Leibes an, der sein eigen Werkzeug und Glieder hat, da er die Frucht empfangen, nähren und tragen kann. Über deren sind wenig, die das betrachten, und Jakob fehlet des Endes auch. Denn er hat Rahel lieb, darum daß sie schön ist; Lea aber hat er nicht lieb, dieweil sie so schön nicht war; welches doch alten Männern weniger geziemt als Jünglingen, daß sie also die Schönheit ansehen sollten.
\quellseite{135}

69. Dieses soll man aber in der Kirche darum lehren und treiben, auf daß wir uns vor dem Ehestande nicht scheuen; wie die Papisten denselben mit ihrem Narrenwerk geschändet und verdammt haben. Denn darum schreibt dieses der Heilige Geist, daß er damit anzeigen will, daß Gott auch die nicht verwerfen und verdammen wolle, welche an ihren Eheweibern die Schönheit ansehen, und daß solche Wahl ohne Sünde geschehen möge. Derhalben soll niemand denken, daß er daran sündige, wenn er ein schönes Weib lieber hat, als eine, die häßlich gestaltet ist. Also können wir daraus auch gelehrt und gebessert werden.

70. Darnach hat Jakob die Hoffnung gehabt, daß er mit der Jungfrau, die da schön und stark war von Leibe, Kinder zeugen würde. Denn wo man Ehe stiftet, soll man ja billig darauf Acht geben, daß man es also mache, daß Stärke mit Stärke, Geschickte mit Geschickten erwählt und zusammengefügt werden möchten. Wie solcher Unterschied mit den unvernünftigen Thieren gehalten wird, als, mit Ochsen, Pferden, Schafen, Hähnen und allen andern Thieren. Unter den Menschen geschieht es aber nicht also um der verderbten Natur willen, so von der Erbsünde hergekommen ist. Denn daher kommt es, daß Etliche stark, Etliche schwach geboren werden, und bisweilen auch Etliche aussätzig, Etliche blind, Etliche närrisch, welches der Eltern Schuld ist. Denn um der Erbsünde willen muß man das dulden, daß starke und schwache Personen im Ehestande zusammengefügt werden, dieweil demselbigen Schaden der Erbsünde gesteuert und gewehret muß werden mit der Hülfe, die uns Gott im Ehestande gezeigt hat, wie auch die Leute sein mögen oder welcherlei Leute darin zusammengefügt werden.

71. Also haben auch alle Thiere ihre gewisse Zeit zur Zeugung; allein der Mensch hat solch irrige und ungewisse Bewegung oder Anreizung von wegen des Giftes der Erbsünde. Solcher Anreizung hat man anders nicht rathen können, denn daß sie also zusammengefügt werden. Und wer da ein schön und stark Weib haben kann, der mag es haben; wer aber die Wahl nicht haben kann, derselbe mag die nehmen, die ihm werden kann, sie sei gleich schön oder nicht schön.

72. Darnach ist der Ehestand nach der Erbsünde dem Fleische sehr verdrießlich und beschwerlich, darum daß die Eheleute nicht können oder mögen geschieden werden. Denn wiewohl die Lust und die Liebe bleibt, so Mann und Weib Eines zum Andern haben, so ist doch gleichwohl das Band, davon Mose sagt, daß das Weib um den Mann sein soll als eine Gehülfin seines Lebens, die nicht von ihm soll getrennt werden, fast hart und schwer. Denn welcherlei Unfall sich zutragen mag an ihrem Leibe, oder Gut, oder an den Kindern, so bleibt doch das Bündniß fest bestehen und kann oder mag nicht getrennt werden, und wird alle Wahl aufgehoben, wird ihnen auch dazu kein Wechsel gestattet, und darf kurzum Eines das Andere nicht verstoßen oder verlassen, dieweil sie leben.

73. Ueber das kommt der Teufel auch dazu, der aller Ordnung Gottes feind ist, welcher die Eheleute auf mancherlei Weise bekümmert und plagt. Und um so viel und so großer Mühseligkeit und Beschwerung willen, so dem Ehestande anhangen, ist es nicht böse, wo Einer ein schön und stark Weib erwählt, auf daß er dieses Band des Ehestandes nebst allen Beschwerlichkeiten um so viel länger und leichter dulden und tragen könne. Denn das ist nicht zu rathen, daß Einer wollte eine solche wählen und nehmen, die ihm bald im Anfange nicht gefiele, dazu widerlich und verdrießlich wäre; dieweil der Teufel auch die, so fast schön sind und einander sehr lieb haben, uneinig zu machen und von einander zu reißen pflegt.

74. Und der Heilige Geist lehrt an diesem Ort, daß Gott solche Wahl nicht verboten oder verdammt hat. Denn er lobt dieselbe, oder läßt es diesem heiligen Patriarchen vielmehr zu, daß er, wie die jungen Leute pflegen, die schöne lieber hat, als die nicht schön ist. Er läßt es geschehen. Denn der Ehestand ist ein solcher Stand, der der Vergebung der Sünden bedarf und der Gnade Gottes, daß er an ihm übersehe, nicht allein was solche fleischliche Lust belanget, sofern dieselbe ordentlich im Ehestande geschieht, sondern auch was sonst die Fehler belanget, da man mit dem Gesinde, mit den Kindern, und Mann und Weib untereinander auch zürnen, wie solches in der Haushaltung vorzufallen pflegt. Denn wie das Weltregiment und die Kirche immer unter der Vergebung der Sünde sein muß: also ist das Leben der Eheleute auch unter Gottes Gnade und Vergebung der Sünde.

75. Gott wil die Stücke im Ehestande nicht verdammt haben, als, daß der Mann sich ein schön Weib erwählt und dieselbe lieb hat. Welches zwar im Anfange nicht schwer ist, wenn die Liebe noch etwas hitzig ist, damit der Bräutigam und die Braut gegen einander entbrannt sind; aber sie wird allemal kalt. Ja, es wird bisweilen über einen Monat gar ein bitterer Haß und Feindschaft daraus durch des Teufels Anreizung, und wenn er die Herzen der Eheleute bezaubert; wie er denn mancherlei und wunderbarer Weise pflegt zu thun.

76. So ist nun dieses zwar keine Entschuldigung dieser Sünde, welche die Heuchler an dem heiligen Patriarchen Jacobus nachher groß machen; sondern es ist ein Ruhm und Preis der göttlichen Gnade, nämlich, daß Gott nicht verdammt, daß Jakob Lea verwirft, als die nicht so schön ist, und die Andere wählt, die da schöner war. Die Papisten deuten es als Geilheit, aber sie sehen nicht, was ihn selbst eine Ursache zum Ehestande getrieben. Denn dieweil Jakob die Verheißung gehabt, liegt ihm die Sorge auf dem Halse, daß er gern Kinder und Nachkommen haben wollte; darum sorgt er dafür, wie er möge ein solch Weib bekommen, die ihm genehm und gefällig sei. Dafür, sage ich, hat er mehr Sorge, denn wo er etwa eine gewisse Wohnung und ein Haus kriegen möchte, welches der Weise und Lehre der Philosophen und Weltweisen stracks zuwider ist. Denn also sagt Hesiod: Du sollst dir erstlich ein Haus und ein Weib, dazu auch einen Ochsen schaffen, damit du den Acker pflügen könnest etc. Der Dinge achtet dieser unser Patriarch Jakob keines, weder Haus, noch Ochsen, auch nicht andere Nothdurft; sondern läßt sich das zu allererst angelegen sein, daß er sich ein Weib auserwählt, die ihm gefällig ist.

77. Darnach sehen die Papisten das auch nicht, daß unter der kindischen Liebe des alten Patriarchen eine treffliche Keuschheit verborgen ist. Denn dieweil der Heilige Geist die Liebe an sich, so Bräutigam und Braut Eines zum Andern haben, nicht verwirft, so will er damit zugleich angezeigt haben, wie keusch Jakob gewesen sei. Denn ist das nicht eine treffliche Keuschheit, wo ein Mann bis an sein achtzigstes oder vierundachtzigstes Jahr keusch lebt? Denn in vierundachtzigsten Jahr ist Jakob erst ein Ehemann geworden; so lange hat er außerhalb der Ehe keusch gelebt nach der ersten Blüthe seiner Jugend, bis die Jahre des Alters heran gekommen sind; denn vierzehn Jahre hernach wird er alt genannt werden.

78. Wenn nun die Papisten das sehen können, daß er eine schöne Jungfrau auserwählt hat, die er lieb gewonnen, warum sehen sie denn dieses nicht auch, daß er keusch gelebt hat bis an sein vierundachtzigstes Jahr? Und ist zu der Zeit die Natur noch etwas stärker und vollkommener gewesen, also daß eine Mannsperson tüchtig gewesen ist, Kinder zu zeugen mit fünfzehn oder sechzehn Jahren. Nun ist es wahrlich ein groß Wunder, von demselben Jahre an und recht in der blühenden Jugend die Rohheit des Fleisches und das Gesetz in den Gliedern, welches die Lustseuche genannt wird, tragen und überwinden, und ein sehr harter Kampf wider das Fleisch, den ihrer Wenige ertragen haben.

79. Darum irren und fehlen die Papisten weit, wenn sie denken, daß die Väter ein fleischlich Leben geführt haben, wie sie es zwar dafür halten. Denn sie sind der Meinung, daß niemand keusch lebe, denn die im ehelosen Stande sind, als Mönche oder Nonnen. Und sie ärgern sich daran sehr, daß die Patriarchen ehelich gewesen, und verstehen von derselben vortrefflichen Keuschheit gar nichts. Es wäre aber kein Wunder gewesen, wenn Jakob in so langer Zeit, da er keusch gelebt hat, sein Fleisch gar getödtet und die angeborne Kraft Kinder zu zeugen ausgelöscht hätte. Denn es sind schon die Jahre seiner blühenden Jugend und er erwachsener gewesen, dazu auch keines männlichen Alters verlaufen, und er hat jetzt schon vier und achtzig Jahre auf dem Halse, hat mit seinem Fleische acht und sechzig Jahre gekämpft: und ist noch gleichwohl an ihm eine natürliche Neigung geblieben zu dem weiblichen Geschlechte, zu einer Jungfrau, und zwar zu der schönsten, die er der vorgezogen hat, welche nicht so gar schön gewesen, da er jetzt schon ein alter und grauer Mann war.
\quellseite{136}

80. Derhalben übertreffen diese Exempel der Keuschheit alle unsere Keuschheit, und vornehmlich den sehr schändlichen unreinen Ehestand und das ehelose Leben in den Klöstern, welches jetzt lauter Hurenhäuser sind. Oder wenn sie auch jemals Keuschheit gehalten haben, ist doch dieselbe mit der Väter Keuschheit nicht zu vergleichen, in welchen das Fleisch also ist getödtet worden mit großem Glauben und Geiste, daß doch gleichwohl die natürliche Neigung und Liebe zu dem weiblichen Geschlecht dadurch nicht ist ausgelöscht worden. Wie Jakob seine Braut recht herzlich lieb hat, also daß ihm die sieben Jahre kaum drei oder sieben Tage lang zu sein däuchten. Darnach, was noch mehr ist, da er ein Bräutigam geworden, wartet er noch sieben Jahre, bis ihm die Braut ist gegeben worden. Wenn jetziger Zeit einer so viel Jahre warten sollte auf die Hochzeit, und sollte unterdessen dienen, sollte arm und ein Bettler dabei sein, er würde wahrlich beide Braut und Schwiegervater fahren lassen.

81. Aber dieses betrachten die groben Säue, die Papisten, nicht, und sehen auch das große treffliche Exempel der Keuschheit nicht, dabei er auch zugleich große Geduld gehabt hat, daß er ganze sieben Jahre den schweren Dienst getragen, und den langen Verzug geduldet, der sonst einem Andern, wer er auch sein möchte, in Wahrheit unerträglich sein würde: und ist noch bei solchem Dienste ganz arm und dürftig, hat gar nichts Eigenes, auch nicht einen Nagel oder Faden. Und das thut er alles aus Liebe, die er zu der Jungfrau gehabt, damit er so gar eingenommen war, daß ihm die sieben Jahre däuchten, als wären es einzelne Tage, das ist, als wären es nur sieben Tage gewesen.

82. Augustinus verwundert sich dieser Geduld und dieses Wartens auch, und fragt: Was doch möge die Ursache gewesen sein, daß er solche Geduld hat haben und so lange warten können, dieweil solches wider und über die Natur des Menschen sei? Denn das bezeugen alle Beschreibungen der Liebe, und derjenigen, die mit großer Liebe behaftet sind, daß es wahr sei, wie jener sagt: Impatienter amo:

Ich werde fast ungeduldig über der Liebe: und vornehmlich, wo die Liebe lange hingehalten wird, bekümmert es die Menschen sehr und macht sie ungeduldig. In den Buhlliedern wird geklagt, daß Tag und Nacht allzulang seien; wie jener auch bei dem Poeten sagt: Si mihi non haec lux toto jam longior anno est etc.: Dieser Tag währet mir länger denn sonst ein ganzes Jahr. Augustinus aber antwortet darauf also: Daß ihm diese Zeit und Jahre so kurz gedäucht, sei daher gekommen, daß ihm die Arbeit nicht schwer oder verdrießlich gewesen sei. Denn die Liebe ist allmächtig. Wie wiederum auch, der da haßt, gar nichts vertragen kann. Die Liebe thut das umsonst, das sonst Einer um seines Geldes oder Bitte willen thun würde. Also kommt es oft vor, daß die Huren oder Ehebrecher mehr leiden, als wohl ein ehelich Weib oder Ehemann leiden würde. Und so viel vermag die schändliche unsinnige Buhlliebe; darum thut und leidet die reine eheliche Liebe noch viel mehr.

83. Es ist aber an dieser Antwort nicht genug. Denn man muß das noch hinzu setzen, daß hier angezeigt wird, nämlich, wie diese Brautliebe vom Geiste des Glaubens in Jakob sei regieret worden: wo aber dieser Regent ist, da ist Alles gut und leicht. Denn wo Gott bei uns ist und mit uns umgeht in allem, was wir thun, so ist das unser vornehmster und höchster Trost. „Ich vermag alles“, spricht St. Paulus Phil. 4, 13., „durch den, der mich mächtig macht, Christus.“ Ob nun wohl diese Dinge fleischlich und kindisch sind, so wird es doch Alles gelobt an diesem sehr heiligen Manne, uns zum Exempel und Troste. Denn er hat die Verheißung gehabt, daß er Kinder und Nachkommen haben würde. Darauf hat er vornehmlich gesehen; darum war vonnöthen, daß er mußte ein Weib nehmen.

84. Nun wollen wir auch die Stücke besehen und erklären, so zur Grammatik gehören. Der Text sagt, daß Lea ein blödes Gesicht gehabt habe; da steht im Hebräischen rakkoth; und ich halte dafür, wir haben es ziemlich wohl verdolmetscht. Denn daß die Andern sagen, daß sie klägliche Augen soll gehabt haben, ist eine Fabel. Mich dünkt, es sei ein ehrlich frommes Mägdlein gewesen, von guten Sitten, auch stark von Leibe; aber nicht schön, sonderlich in den Augen, welche den Menschen sehr pflegen zu zieren. Wie Plinius die Augen den König nennt unter den Gliedern, so der Mensch am Leibe hat; und sie sind auch wahrlich der beste Theil am Leibe. Darum ist es eine seltsame Gabe am Menschen, schöne und liebliche Augen haben.

85. Es werden aber vornehmlich die Augen gelobt, die man gemeiniglich Falkenaugen nennt; wie man von den Mägdchen sagt, sie sieht wie ein Falke; denn solche Augen sind sehr scharf und lebendig, und leuchten schier wie der Blitz. Sie sind fein rein und hell, und man kann beiderlei an den Augen abnehmen, nämlich, das Laster und auch die Tugend, so im Herzen verborgen liegt. Darum sind die schönen hellen Augen wie ein Wunderwerk im menschlichen Geschlechte; denn sie bedeuten einen vortrefflichen Verstand und eine sonderliche gute Art. Hernach wird vom Messias gesagt im 49. Capitel dieses Buches V. 12.: „Seine Augen sind röthlicher, denn Wein“; sie sind wie der schwarze rothe Wein. Solche Augen zieren das Angesicht sehr und geben ihm gleichsam Leben und sonderliche Fröhlichkeit.

86. Dagegen aber kleine und blöde Augen, die nicht scharf und lebendig sind, verdunkeln und entstellen das ganze Angesicht. Solche Augen hat Lea gehabt, damit sie nicht scharf gesehen hat; darum hat sie keine Gunst gehabt um der zarten oder blöden Augen willen. Denn das hebräische Wort rakkoth heißt, das da zart oder \unsicher{feinlich} ist; wie droben im 18. Cap. V. 7. von Abraham steht: „Er holete ein gut Kalb“ 2c. Das Geheimniß, wovon die Andern an diesem Orte handeln, lasse ich fahren, und folge dem Verstande, so die Historie mitbringt; denn diese Ursache ist an sich selber einfältig und klar, warum Jakob Lea nicht so lieb gehabt, nämlich, dieweil sie nicht scharfe und schöne Augen gehabt.

87. Rahel aber war schön, und die Bedeutung des hebräischen Wortes thoar bezieht sich eigentlich auf die Form, Gestalt und Proportion. Die Augen der Rahel haben ihre rechte Gestalt gehabt, desgleichen die Stirne und die Backen, und der ganze Leib hat seine ordentliche Gestalt und die Glieder haben ihre Geschicklichkeit auch fein gehabt. Dies ist die wahre Schönheit und Zierde, wenn in dem Gesichte eine geschickte und passende Proportion der Augen, der Stirn, der Wangen und der andern Theile anzutreffen.

88. Und bisher hat Mose diesen Bräutigam und Liebhaber Jakob kindisch genug beschrieben, der um der Liebe willen zu einer solchen schönen hübschen Jungfrau so einen langen und schweren Dienst getragen hat. Daraus man wohl abnehmen kann, wie groß die Liebe gewesen sei. Und ist zu merken, daß die heilige Schrift diese Liebe an dem Bräutigam nicht verdammt, sondern dieselbe rühmt. Nun folgt aber ein schön herrlich Exempel der Geduld, so Jakob in diesem Dienste und in der Liebe zu der Jungfrau gehabt hat.

\section[Dritter Theil]{Dritter Theil. Wie Jakob sein Weib verlangt, daß den Lohn seiner Arbeit, und wie er von Laban schändlich betrogen wird}

\subsection{I}

\begin{quote}
\textbf{V. 21. Und Jakob sprach zu Laban: Gib mir mein Weib; denn die Zeit ist hie, daß ich beiliege.}
\end{quote}

89. Jakob nennt Rahel sein Weib, so sie doch eine Jungfrau ist. Denn die Schrift macht keinen Unterschied zwischen einem Weibe und einer Braut, so einem Manne vertrauet ist: wie das in Mose und in Matthäus zu sehen ist, da die Jungfrauen, so vertrauet sind, Weiber genannt werden. Derhalben hält Jakob Rahel auch für sein ehelich Weib, und hat sie recht herzlich lieb, wie ein Mann sein Weib lieben soll, und wie die Jünglinge und Bräutigame pflegen ihre Bräute lieb zu haben.

90. Es ist aber die Liebe des Bräutigams am allerheftigsten, wenn die Zeit kommt, daß er beiliegen soll. Aber die reine, keusche und heiße Liebe, so dieser unser Patriarch zu seinem Weibe Rahel gehabt, wird an ihm ganz unschuldig Weise verhindert, und er wird darüber gemachet und übel betrogen, wie im Text folget.

\begin{quote}
\textbf{V. 22. Da lud Laban alle Leute des Orts, und machte ein Hochzeitsmahl.}
\end{quote}

91. Der Schalk wäre werth, daß man ihn in Abgrund der Hölle verfluchte. Mose sagt nicht, daß er habe lassen zur Hochzeit laden, welches sein ordentlich und ehrlich pfleget zu geschehen; sondern den guten, frommen und heiligen Mann zu hintergehen und zu betrügen, auf daß, wenn er ihn mit List in seinen Strick gebracht, derjenige behalten müßte und keineswegs wieder verstoßen und umtauschen dürfte, welche er die Nacht erkannt hätte, wiewohl er damit hintergangen und betrogen war. Darum rafft er bald in der Eile etliche Zeugen zusammen. Und es läßt sich ansehen, als wolle die Schrift damit zu verstehen geben, daß Laban kein gut Gericht gehabt habe bei frommen und ehrlichen Nachbarn, welche ohne Zweifel seine bösen Stücke, Gewalt und Unrecht verstanden haben, so er die lieben Jahre über dem frommen und getreuen Knechte Jakob gethan hat. Darum hat er lose leichtfertige Leute zusammengefaßt, welche diesen Betrug nicht strafen würden. So wird damit angezeigt, daß er ein böser Schalk und im Geiz gar ersoffen gewesen sei.
\quellseite{137}

92. Jakob aber war fröhlich, und weil er Rahel so gar heftig lieb gehabt, hat er keinen bösen Argwohn gehabt, daß ihm würde unrecht geschehen, sondern hat gute Hoffnung, ist freudenvoll, und wartet also auf seine allerliebste Braut, daß er mit der Freude, darauf er lange gehofft und gewartet hatte, ergötzt werden möge.

\subsection{II}

\begin{quote}
\textbf{V. 23–25. Des Abends aber nahm er seine Tochter Lea, und brachte sie zu ihm hinein; und er lag bei ihr. Und Laban gab seiner Tochter Lea Silpa zur Magd. Des Morgens aber, siehe, da war es Lea. Und er sprach zu Laban: Warum hast du mir das gethan? habe ich dir nicht um Rahel gedienet? warum hast du mich denn betrogen?}
\end{quote}

93. Dieses ist nun das Exempel der Geduld an Jakob; derhalben wer da kann, der lerne Geduld. Ich kann es nicht; denn diese Geduld ist unerträglich und unnachahmlich. Der fromme heilige Mann hat die Jungfrau so heftig lieb, daß er die ganzen sieben Jahre alles thut und leidet, was er nur kann und mag; ja, daß er den Geizhalß und den alten Schalk Laban möge reich machen mit der Bedingung, daß er ihm dagegen seine Tochter Rahel geben soll. Aber für solchen treuen Dienst hat er ihm also gedankt, daß er ihm sein Weib heimlich nimmt, nämlich die Jungfrau, welche er lieb hatte, und darnach er so viel Jahre Verlangen gehabt und auf die er gewartet hatte. Und das thut er eben in der Stunde, da er die Brautfreude, darnach ihn sehr verlanget, haben soll; und er nimmt sie ihm nicht allein, welches zwar an sich sehr verdrießlich ist, sondern bringt ihm auch wider seinen Willen eine Andere auf, welche er nicht lieb hat, und legt ihm also eine solche Last auf seinen Hals, die er für und für tragen muß.

94. Dieses ist wahrlich ein greulicher und sehr unbilliger Betrug. Jakob weiß anders nicht, denn daß er seine allerliebste Braut in Armen habe, und siehe, da hat er die, so ihm anstatt seiner geliebten Braut beigelegt war. Wenn Einer dem Andern sein Geld, Gold, Silber oder Vieh raubt oder mit Gewalt nimmt, das ist ein geringer Schade; aber wenn dir Einer eine Jungfrau, ein schönes geliebtes Weib nimmt, mit welcher du von ganzem Herzen hoffst Kinder zu zeugen, von der du hoffst Samen zu bekommen der Nachkommen, so dir verheißen sind: das ist wahrlich eine solche Gewalt und Schmach, die größer ist, denn sonst irgend eine Gewalt sein kann. Denn da ist alle Hoffnung und alles Warten Jakobs plötzlich dahin gefallen, und er hat anders nicht denken können, denn daß er nun stets der Liebe Rahels, seiner Braut, würde müssen beraubt sein und sie würde nunmehr sein Weib nicht werden.

95. Also wird er in seiner guten Hoffnung betrogen, da er lange auf die freundliche Liebe, wie er sich zu seiner Braut halten wollte, gewartet hat; welche Hoffnung ganze sieben Jahre gewähret, darüber sie alle beide große Freude gehabt. Daß sie aber diese Gewalt gelitten, ist ohne Thränen und große Schmerzen nicht abgegangen. Denn das Unglück ist allzu greulich und groß, und die Geduld ist auch unglaublich groß. Ich hätte es zwar nicht gelitten, sondern hätte ihn darüber beschuldiget, und vor Gericht mit ihm gehandelt, daß er mit Recht wäre gezwungen worden, mir meine Braut zu geben, und hätte die Andere verstoßen und sie ihrem Vater wiederum nach Hause geschickt. Denn es war ihm nicht Lea, sondern Rahel vertraut worden; und dasselbe ist ohne Zweifel jedermann durch die ganze Stadt öffentlich angesagt und bekannt gewesen.

96. Ich halte aber dafür, daß die Leute, so Laban geladen hat, nichtswürdige und leichtfertige Leute gewesen seien: „lose Fähndichen“, die dem Laban zu Gefallen diese böse That haben loben und ihm darin entschuldigen müssen. Daß aber Jakob und seine liebe Braut dermaßen haben sollen geschieden werden, ist fast ein schwer und unleidlich Ding gewesen. „Ich könnte es wahrlich nicht leiden, wenn ich auch schon älter wäre von Jahren, denn Methusalem gewesen ist.“ Denn solch Scheiden ist wider die Natur und alle menschliche Vernunft. Wie großen Lärm haben in allen Historien angerichtet die, so Andern ihre Jungfrauen und Weiber mit Gewalt entführt haben! Wie oft sind etliche Königreiche, die in großer Gewalt gestanden, mit greulichem Kriege an einander gerathen, und sind um solcher Gewalt willen umgekehrt und zerstört worden!

97. Die Griechen haben schier die ganze Welt erregt um der Helena willen, die auch mit Gewalt weggeführt war. Die Sabiner haben einen Krieg wider die Römer angefangen, darum daß ihnen ihre Töchter waren entführt worden. Und hernach wird im 34. Cap. 23. 25. 26. dieses Buches folgen, wie die Sichemiten von den Söhnen Jakobs sind geschlagen und erwürgt worden, darum daß Dina geraubt war. Denn die herzliche Liebe, so ein Bräutigam hat, ist zu zart und kann nicht viel leiden, vornehmlich wenn er jetzt soll Hochzeit halten und mit der Braut beiliegen, da nun die Hochzeitsfreude erst angeht. Darum kann er nichts weniger dulden oder leiden, denn daß ihm die Braut soll genommen werden, die er heftig lieb hat, daß er auch sein Gold oder Silber, ja, seine Augen und sein eigen Leben so lieb nicht hat.

98. Darnach hatte Jakob diese Braut gesucht und dieselbe erwählt auf Befehl und nach dem Willen seines Vaters, um der Erben und Nachkommen willen, so ihm verheißen waren, die er gehofft hat von der Braut zu bekommen, die da schön und jung war; aber dieselbe Hoffnung und Freude, darauf er so lange gewartet hatte, ist in einem Augenblick gar dahin gefallen. Derhalben, wer es nur kann, mag diese große und unerhörte Bosheit groß machen und verfluchen in den Abgrund der Hölle; denn sie ist größer, denn daß man sie mit Worten genugsam sollte ausreden können. Ja, es läßt sich ansehen, daß im Texte auch angezeigt werde, daß solche Bosheit darin heimlich gestraft sei. Denn es wird von Laban nicht gesagt, daß er Lea zu Jakob gebracht habe als sein Weib, wie hernach von Rahel gesagt wird, daß ihm Laban seine Tochter Rahel zum Weibe gegeben habe. Mose will Lea nicht ein Weib nennen, sondern sagt nur also: „Er nahm seine Tochter Lea und brachte sie zu ihm hinein.“

99. Es hat aber den Geizhalß Laban zu solcher großen Sünde der ehelichen und schändliche Geiz getrieben, und ist ihm nicht so sehr darum zu thun gewesen, daß er dem Jakob seine älteste Tochter wider seinen Willen aufdringen wollte, als daß er ihn noch andere sieben Jahre in seinem Dienste behalten möchte. Denn er hat empfunden und erfahren, daß ihm sein Ding sehr glücklich fortgegangen und sein Gut gemehret worden ist. Derhalben ist es sehr unbillig und greulich, daß ein solcher frommer Mann durch solche große Sünde soll hingehalten werden, daß nur dem Geize dieses alten Schalks möge genug geschehen. Aber er muß den Hohn leiden, der doch gar zu hart ist.

100. Wie Rahel möge gesünnet gewesen sein, wird im Text nicht erzählt; es ist aber wohl anzunehmen, daß sie nicht weniger Unmuthes gewesen sei. Denn sie wird ohne Zweifel wohl gewußt haben, daß sie die Braut war, und die ganze Nachbarschaft hat wohl gewußt, daß Jakob um Rahel sieben Jahre gedient hatte: darum wird Rahel eben so wohl um dieses Unrechts willen zornig geworden sein als der Bräutigam selbst; ja, sie hat es viel weniger ertragen und überwinden können. Denn das weibliche Geschlecht und die Natur an den Weibern ist etwas schwächer als an den Männern; sie war dessen gewiß, daß sie die Braut war und jetzt sollte dem Bräutigam beigelegt werden: aber gleich in derselben Stunde, da sie meinte, daß sie nun der Liebe und ihres erwünschten Bräutigams genießen wolle, wird sie unbilliger Weise betrogen, nicht ohne große Schmerzen und viel bittere Thränen, die sie darüber wird vergossen haben: sie wird jämmerlich geheult und geweint haben.

101. Derhalben ist dieses eine betrübte und traurige Hochzeit. Denn der Bräutigam wird jämmerlich betrogen; Rahel aber, die Braut, wird verstoßen, welche ohne Zweifel ihren Schmerz und Zorn, so sie darüber billig gehabt, nicht verhehlen können, sondern hat denselben mit Weinen und Schreien angezeiget und bezeuget. Zudem wird sie mit ihrem Vater geredet und ihn öffentlich beschuldigt haben. Warum, wird sie gefragt haben, nimmt du mir meinen lieben Bräutigam, den du mir doch selbst vertrauet hast? Aber der grausame Vater ist dadurch gar nicht bewogen worden; vielleicht hat er sie auch wohl hart bedroht, daß sie hat schweigen müssen.
\quellseite{138}

102. Lea hat wohl gewußt, daß sie die Braut nicht war, sondern der Braut Schwester, und wird vielleicht dem Vater etwas widerstrebt haben: aber doch hat sie endlich aus menschlichem Affect nachgegeben, daß sie also ist hingerissen und zu solchem Unrecht ist gezwungen worden. Ja, sie wird sich dieser Gewalt erfreut haben, daß sie also eine Mutter werden möchte der Erben und Nachkommen, so Jacob verheißen waren; denn sie hatte gehört, daß Jacob auf die Erben gewartet hat. Und hernach im 30. Capitel V. 15. läßt es sich ansehen, daß Rahel etlichermaßen an den Tag gegeben hat, daß sie auf ihre Schwester Lea eben um dieser Ursache willen Unmuthes gewesen und sie darum gehaßt habe.

103. Darum, sage ich, ist dieses dem Jacob und der Rahel eine unglückselige Hochzeit gewesen um des Betruges willen, daran sie beiderseits wenig gedacht hatten. Und wird uns hier vorgehalten ein Exempel großer trefflicher Keuschheit an dem heiligen Patriarchen Jacob, und auch eines schweren Kreuzes und sehr großer Geduld. Denn ich weiß nicht, ob irgend ein Heiliger, geschweige ein Heide, dieses Kreuz mit solcher großen Geduld hätte tragen können. Es ist zu hoch und zu schrecklich.

104. Es wird aber gefragt: Was doch diese Leute für einen Gebrauch mögen gehabt haben auf ihren Hochzeiten, weil Jacob mit Lea allein in der Schlafkammer und in Einem Bette schläft, und wird doch nicht gewahr, daß ihm Lea anstatt Rahel beigelegt war; denn er hat ja ganze sieben Jahre in demselben Hause gewohnt, da er ohne Zweifel die Braut etlichemal gesehen und bisweilen auch mit ihr wird geredet haben; und er hätte solches im Bette erforschen können, wo nicht an der Stimme, so hätte er es doch sonst erkennen und unterscheiden mögen. Unserer Gewohnheit nach pflegen wir in solchen Sachen viel vorsichtiger zu handeln; denn wir führen die Braut nicht im Finstern, sondern mit Trommeln und Fackeln beide in die Kirche und auch zum Brautbette, nehmen auch Leute dabei zu Zeugen, daß keinerlei Betrug daselbst geschehen kann. Daß aber ein solcher Gebrauch bei den Leuten nicht gewesen sei, kann man aus dieser Historie abnehmen; denn sonst wäre Jacob so schändlich nicht betrogen worden.

105. Darnach ist es zu der Zeit gebräuchlich gewesen, daß die Männer und Weiber unterschiedliche Wohnung und in den Häusern jedes sein eigen Gemach inne gehabt; sie haben nicht also untereinander gewohnt, wie bei uns geschieht. Die Knaben haben bei den Vätern gewohnt und die Jungfrauen bei den Müttern besonders allein; in der Schlafkammer aber, da Jacob seine Braut bei sich gehabt, hat er zwar frei- und ohne Scheu mit ihr reden können, und er hätte es auch am Anrühren wohl abnehmen können, daß es Rahel nicht war.

106. Man kann aber daraus schließen, daß zu der Zeit sehr große Scham bei den Leuten muß gewesen sein; oder daß sie der Patriarch Jacob sehr schamhaftig und züchtig gewesen ist; oder aber daß Lea wird vom Vater also unterrichtet gewesen sein, daß sie gar stille geschwiegen, oder je leise reden sollte: und muß aber doch Jacob sehr schamhaftig gewesen sein, daß er mit ihr nicht geredet und sie auch nicht sonderlich berührt hat, sondern hat sich nur allein mit solcher Liebe zu ihr gehalten, wie einem Ehemann geziemet, vor großer Freude, daß er nun endlich mit seiner liebsten Braut der Liebe pflegen möchte.

107. Diese große Scham und treffliche Keuschheit rühmt die Schrift, daß er in das Brautbett ohne allen Argwohn, ohne Geilheit, ohne unzüchtige Begierde, sondern mit einfältiger und ehelicher Begierde gestiegen ist. Denn daß Lyra sagt, daß er drei Nächte sich ihrer solle enthalten haben, das lasse ich fahren. Man soll vielmehr die große Einfalt betrachten, daß er gar an keinem Wahrzeichen gespüret hat, daß es eine fremde Jungfrau wäre, bei welcher er gelegen hat. Denn man muß die Väter nicht also ansehen, wie die Papisten von ihnen urtheilen, daß sie gar verstürzt und unsinnig gewesen seien, Unzucht zu treiben, wie sie, die Papisten, selbst sind; sondern an Jacob ist eine sehr einfältige Keuschheit gewesen und eine sehr reine keusche Liebe zu seiner lieben Braut, und er hat keine Ursache gehabt, solchen Argwohn zu schöpfen, daß ihm sollte eine Andere anstatt seiner Braut beigelegt sein worden. Denn beide Eltern und die Braut selbst auch sind allesamt sehr wohl bekannt, und man hat allezeit an ihnen allen äußerliche große Treue und rechte Frömmigkeit gespürt; und wird dennoch betrogen.

108. Ich denke aber und entschuldige Jacob also, wie ich droben von Lot gesagt habe, da er bei seinen Töchtern geschlafen hat. Denn es kann wohl geschehen, daß bisweilen ein Mensch durch irgend eine starke Einbildung, darauf er in seinem Herzen gar beflissen und befürzt ist, dermaßen verschlungen wird, daß er weder hört noch sieht, und gar nichts fühlt oder empfindet, wenn er auch schon mit offenen Augen steht und mit offenen Ohren hört. Denn ein heftiger und tiefer Gedanke, darauf der Mensch gar beflissen ist, zieht das Herz von allen Sinnen ab, daß der Mensch äußerlich gar nichts sieht oder fühlt.

109. Also geschieht es ganz und gar auch, wenn die heftigen Affecten und Bewegungen den Menschen einnehmen, entweder im Zorn, oder in Traurigkeit, oder auch in der Liebe. Denn solche Affecten und Bewegungen machen, daß der Mensch gar verstürzt wird und gleichsam von Sinnen kommt. Daher kommt es, daß du oftmals im Kriege, wenn man mit den Feinden zusammengetroffen soll, an den Kriegsknechten sehen magst, wie sie sich entsetzen, wie so gar erschrocken und bestürzt sie sind vor großer Furcht. Also war Lot vielmehr von Traurigkeit denn vom Weine trunken, und war in der Traurigkeit gar verschlungen, so er deswegen gehabt hat, daß Sodom so greulich mit höllischem Feuer verbrannt war, und daß er derhalben nichts gefühlt oder empfunden hat; sondern was er gethan, hat er alles unwissend und mit einem fremden Herzen gethan, daß er nicht bei sich selbst gewesen ist, wie wir es pflegen zu nennen.

110. Dieses ist ein wunderlich Ding, ist aber doch nicht ungewöhnlich, und geschieht natürlich an den melancholischen Menschen, wenn sie irgend auf ein Ding ganz und gar beflissen sind und ernstlich darauf denken; und so können oftmals einen solchen melancholischen Menschen alle sein Gut nehmen, wenn er schon dabei wäre, daß er darauf Acht gäbe.

111. Eben auf solche Weise ist auch das Herz Jacobs mit Liebe und Freude eingenommen und darin verschlungen gewesen; und er hat Gott gedankt, daß er die Braut gehabt hat, von welcher er hoffte Erben zu bekommen: er ist trunken gewesen, nicht vom Wein, sondern von Liebe, die er zu Rahel hat, auf die er so lang gewartet und nach der er Verlangen gehabt hatte. Solche Trunkenheit der Liebe ist wahrlich ein großer Schrecken, dadurch er sich entsetzt und gleichsam entzückt worden ist, daß er den Betrug nicht hat merken oder verstehen können. Denn er hat gar keinen Argwohn, daß ihm etwas Böses hätte begegnen mögen, sondern er ist ganz und gar in seine Gedanken vertieft, und ist in der Liebe dermaßen ersoffen, daß er dessen gewiß ist, es könne ihm nunmehr die Braut nicht genommen werden. Darum hat er an nichts weniger denken können, denn daß ihm der alte Schalk die Braut entziehen sollte und ihm Lea an ihrer Statt beilegen.

112. Auf diese Weise habe ich droben auch gesagt von Isaac, da er Jacob für Esau gesegnet hat. Wiewohl derselbe die Stimme Jacobs kannte, und sein Herz ihm zuvor sagte, daß eine List und Betrug darunter verborgen sein müßte, doch, weil er so gar auf den Segen, den er seinem Sohne geben wollte, beflissen und bestürzt war, hat er auf die Dinge, die er gefühlt und gehört hat, keine Achtung gehabt.

113. Solche Exempel kommen oftmals vor in diesem Leben, wenn die Herzen der Menschen von den äußerlichen Sinnen auf die Gedanken zurückfallen, damit sie zuvor behaftet und bekümmert gewesen sind, und meinen, es sei nicht möglich, daß sie könnten betrogen werden, oder sonst irgend einen Schaden nehmen; denn sie glauben ihrem Herzen mehr, als den Augen und Ohren. Darum müssen wir denken, daß Jacob auch dermaßen von Liebe trunken und bestürzt, ja, gleichsam entzückt gewesen ist, daß er an nichts weniger gedacht hat, denn daß er könnte betrogen werden. Diese große Hitze der Liebe hat ihn blind gemacht.

114. Dieses sind also meine Gedanken, die mir von dem Betrug dieses alten Schalkes Laban, wie Jakob damit betrogen worden sei, eingefallen sind. Dabei aber soll man auch zugleich merken das treffliche Exempel der Scham und Zucht an Jakob, welches sonst selten ist. Und man soll die Keuschheit der Patriarchen groß machen und hoch heben über alle Keuschheit und eheliches Leben der Mönche. Denn hier werden die hitzigsten Bewegungen der Liebe und Unzucht überwunden, oder werden sie zum wenigsten etwas inne gehalten an allen Beiden, an Jakob und Rahel, welche, wiewohl sie die rechte Braut war, doch aus dem Brautbette, darauf sie gehofft und das ihr auch gebührete, ausgeschlossen wird, und gleich thun muß, als wäre sie die Braut nicht; sie muß sich selbst trösten, und mit heißen Thränen die elende Liebe und Begierde, so sie zu ihrem Bräutigam gehabt, etwas lindern, so viel als sie kann und mag.
\quellseite{139}

115. An diesem Ort aber ist zu sehen, wie groß die väterliche Gewalt bei diesen Leuten gewesen ist. Denn Lea wird nur mit einem Wort und Wink von ihrem Vater Laban dazu gezwungen, daß sie muß an der Rahel Statt treten, und daß sie hat müssen die Braut werden auch wider die gegebene Treue und Glauben. Denn Jakob hatte vor der Zeit mit Lea kein Wort geredt, daß die Ehe betreffend hätte. Da war zwischen ihnen keine Liebe, kein Trauring, kein Verbindniß vorhergegangen; sondern der Vater nimmt die Tochter hin und legt sie dem Jakob in das Brautbett ohne vorhergehende Bewilligung, ohne allen Vertrag oder Unterredung. Der Vater fragt die Tochter nicht, ob es auch ihr Wille sei, er hört ihre Antwort nicht.

116. Wo dieses bei den Leuten der Gebrauch gewesen ist, so ist es wahrlich ein sehr böser Gebrauch gewesen. Denn das ist gewiß, daß mit Lea und Jakob kein Verlöbniß gehalten und daß keine Zeugen dazu gezogen worden, sondern es wird alles in der Eile und unversehens ausgerichtet durch die Tyrannei und Gewalt des Vaters, der die Tochter hinnimmt und sie zu Jakob führt. Darum ist dieses ein sehr böses Exempel, darin wir sehen, daß Lea eine Braut hat werden müssen aus lauter Tyrannei und greulicher Kühnheit ihres Vaters. Es ist wohl abzunehmen, daß sie zwar leichtlich darein wird gewilligt haben, daß sie möchte eine Mutter werden der heiligen Nachkommenschaft, so Jakob verheissen war, und darum ist sie auch zu entschuldigen: aber die Tyrannei des Vaters soll man durchaus verfluchen und verdammen. Und wiewohl Jakob, der heilige Mann, diese Gewalt leidet, so ist es doch ein verflucht Exempel, dem man mit nichten nachfolgen soll.

117. Lyra aber und die Andern fragen an diesem Ort: Ob auch in der Nacht zwischen Lea und Jakob eine rechte Ehe gewesen sei? Darauf antworte ich also, daß keine rechte Ehe gewesen sei. Ist es denn ein Ehebruch gewesen? Antwort: Mit nichten. Was ist es aber denn gewesen? Antwort: Es ist ein Monstrum, ein Ungeheuer gewesen. Denn wo du die That an sich ansehen willst, so ist Jakob der Lea Ehemann nicht, sondern er schwächt und beschläft die, so ihm nicht ehelich vertraut war: und er hat doch daran keine Schuld, sondern die Schuld ist des treulosen Mannes Laban, der den Bräutigam und die Braut betrügt, und trennt die Ehe damit, daß er dem Jakob eine Andere zufügt und beilegt, denn er ihm vertraut hätte. Derhalben kann Jakob billig kein Ehebrecher genannt und auch für keinen gehalten werden.

118. Und wenn man fragt: Ob Jakob auch damit gesündigt habe, daß er eine Fremde erkannt hat, soll man antworten: Die That an sich sei zwar eine offenbare Hurerei und keine Ehe, er habe aber gleichwohl nicht gesündigt; denn er ist durch eines Andern Untreue betrogen worden, daran er nicht einmal hat denken können. Denn wer wollte befürchten, daß ein Vater seine Tochter, die von ihm geboren ist, solcher Freude, darauf sie so lange gehofft hatte, berauben sollte? Die Tochter hat anders nicht gedacht, denn sie würde die Nacht eine Braut, und sollte nun ihrem allerliebsten Bräutigam zugeführt und beigelegt werden; und Jakob ist auch mit rechter ehelicher Liebe in die Schlafkammer gegangen, hat daran nicht gezweifelt, er würde Rahel als sein liebes Weib herzen. Denn das hat ihm der treulose Bube, Laban, versprochen, und die ganzen sieben Jahre sich nicht anders vernehmen lassen. Aber in einem Augenblick trennt er die Ehe, die er so lange heuchlerisch versprochen hatte.

Das heißen grobe Bubenstücke.

119. Derhalben sollen wir auf diese vorgelegte Frage also antworten: Daß zwischen Jakob und Lea keine Ehe gewesen sei; denn es hat an ihrer beider Willen gemangelt und kein Theil hat eingewilligt. Sie war dem Jakob nicht ehelich vertraut worden, so war Jakob auch der Lea Bräutigam nicht; sondern es ist ein Irrthum geschehen und eine unüberwindliche Unwissenheit, und in einer Sache, die da weltlich ist, davon die Juristen sagen, daß sie den Menschen ganz und gar entschuldigen soll.

120. Denn sie machen einen Unterschied zwischen der groben, unachtsamen Unwissenheit und der unüberwindlichen, unvermeidlichen Unwissenheit. Wo man von der That, von der Person, von den Umständen und allen andern Dingen nichts weiß, mit solcher Unwissenheit sündigt man nicht; sondern so viel das Weltregiment belanget, ist es Unschuld. Aber in der Theologie taugt solches nicht, wiewohl dieselbige Unwissenheit Etliche dahin auch gezogen haben; sondern sie gilt nur außerhalb der Theologie, im weltlichen Regiment, da hat sie Statt: sie soll aber vor GOttes Gericht nicht gezogen werden, sonst würden alle Menschen selig werden. Als zum Exempel: das ist ein unüberwindlicher oder unvermeidlicher Irrthum und Unwissenheit, daß die Juden den Sohn Gottes gekreuziget haben, wie das St. Paulus 1 Cor. 2, 8. bezeuget; denn die Vernunft verstehet es nicht: sie sind aber darum vor GOtt nicht entschuldiget. Darum soll man die unüberwindliche Unwissenheit außerhalb der Theologie setzen und gelten lassen. Im weltlichen Regiment aber ist es eine gänzliche Entschuldigung.

121. Wie diese Vermischung Jakobs und Leas keine Ehe ist und die That an sich Hurerei ist; es wird aber durch die unüberwindliche Unwissenheit entschuldigt: und Lea wird auch entschuldigt, daß man sie für keine Hure halten soll, nämlich, durch die Gewalt des Vaters, der gesagt hat: Ich gebiete dir und zwinge dich, will es kurzum von dir haben als dein Vater; du, Lea, bist die Braut nach meiner Gewalt und Willen. Derhalben geht sie hin in ihrer Einfalt und reinem kindlichen Gehorsam: und wiewohl es nicht glaublich ist, daß sie sich dessen sollte hart geweigert haben, daß sie zu Jakob ist gebracht worden, so wird sie doch durch die väterliche Gewalt und kindlichen Gehorsam entschuldigt.

122. Es ist aber ein sehr schwer und unbillig Ding, daß sie der Vater durch solche Tyrannei zwingt und daß sie sich nicht hat wider ihn setzen dürfen. Darum soll man alle Schuld auf den Geizhals Laban legen, der billiger Nabal heißen sollte, welcher damit umgegangen ist, daß er beide Töchter dem frommen Mann Jakob mit Gewalt hat aufzudringen wollen. An Jakob aber ist dessfalls keine Sünde oder Schuld, sondern ist bei ihm lauter Unwissenheit. Denn er wird unwillkürlich und wider seinen Willen der Liebe seiner freundlichen lieben Braut beraubt, welche er hofft dieselbe Nacht zu herzen.

123. Wie die Römischen Historienschreiber ihre Lucretia auch entschuldigen; wiewohl daselbst keine unüberwindliche Unwissenheit war, sondern der Lucretia ist Gewalt geschehen. Denn es ist zwar ein Ehebruch daselbst begangen, weil zween Personen zusammenkommen, die beide Eheleute sind, und haben auch beide wider die eheliche Pflicht gesündigt: doch ist allein Tarquinius ein Ehebrecher, Lucretia aber hat die Ehe nicht gebrochen; denn sie ist mit Gewalt durch das Schwert gezwungen worden, daß sie mit ihrem Leibe die Schande hat begehen müssen; sie hat von dem Ehebrecher leiden müssen, daß er ihr Gewalt gethan, welcher Gewalt das arme schwache Weiblein sich nicht hat erwehren können.

124. Also erzählt man auch ein Exempel von der heiligen Jungfrau Lucia. Da ihr der Richter hart gedroht hat, daß er sie zu dem Altar der Götzen und mit Gewalt in das gemeine Haus führen wollte, und nun schon etliche Jünglinge versammelt waren, die sie schänden sollten, hat sie geantwortet: Willst du meine Hand mit Gewalt hinreißen, den Abgott Weihrauch zu opfern oder zu räuchern, so bin ich kein Götzendiener, sondern du bist einer; und wo ich in das gemeine Haus mit Gewalt geführt werde, so will ich derhalben keine Hure werden, sondern ich werde eine zweifache Krone der Keuschheit bekommen.

125. Eben auf dieselbe Weise wollen wir auch den Patriarchen Jakob entschuldigen, der aus einfältiger und ehelicher Liebe in die Brautkammer gegangen ist, und denkt, er wolle da seine Rahel finden; aber ohne seinen Willen und daran er nie gedacht hat, vermischt er sich mit Lea, die ihm nicht verlobt oder vertraut war.
\quellseite{140}

126. Da er sich aber darüber hart beklagt, und Laban derhalben beschuldigt hat, wird es noch eine rechte Ehe werden, wenn GOtt damit dienstfertig und übersieht, und diese Vermischung damit bestätigt, daß er ihm Kinder von ihr gibt, als nämlich, Simeon, Levi, Ruben und andere mehr. Da muß es wohlgethan heißen, auf daß alles, was die Heiligkeit thun, ihnen zum Besten dienen möge, sonst ist es weder eine Ehe noch ein Ehebruch, sondern wie gesagt, ein Monstrum, ein Ungeheuer.

127. Und wir sollen dem Exempel mit nichten nachfolgen: man soll die Kinder nicht zwingen, die zu lieben oder zu ehelichen, vor denen sie sich scheuen, oder denen sie nicht gewogen sind; es fällt sonst Gefahr und Mühseligkeit im Ehestande mehr denn genug vor, wenn auch die zusammengefüget werden, die sich unter einander sehr lieb haben, wie das die täglichen Exempel ausweisen.

128. Es wird uns aber auch diese Geduld nicht vorgehalten, daß man derselben solle nachfolgen; denn sie ist allzu hoch und gar zu ritterlich, nämlich, wo Einer leiden und mit Geduld ertragen soll, daß ihm seine Braut genommen werde eben in der Stunde und dem Augenblick, wenn die Hochzeit gehalten wird und das Beilager geschehen soll. Darum mag man diese Geduld mehr rühmen und bewundern, denn ihr nachfolgen.

129. Mose aber beschreibt nun weiter, wann Jakob den Betrug gemerkt und wie ihm sein gutes Muthe gewesen. Denn er wird ohne Zweifel sehr heftig bestürzt worden sein, da er des Morgens, als die Sonne aufgegangen war, gesehen hat, daß ihm eine Fremde zugefügt worden war; und ist nicht allein darüber erschrocken, daß er gesehen hat, daß er betrogen worden und ihm seine allerliebste Braut genommen worden war, sondern der Schrecken im Gewissen ist auch dazu gekommen um der Hurerei oder Blutschande willen. Denn er war ein frommer gottseliger Mann, der sich mehr um die Sünde, denn um die Gewalt, die ihm geschehen war, bekümmert hat. Und er hat nicht alsbald verstehen können, wie doch dieser Wechsel möchte geschehen sein. Er legt aber die Schuld auf Laban und zeigt seine Unschuld an: „Warum hast du mir das gethan?“ spricht er, „habe ich dir nicht um Rahel gedienet? warum hast du mich denn betrogen?“

\begin{quote}
\textbf{V. 26. Laban antwortete: Es ist nicht Sitte in unserm Lande, daß man die jüngste ausgebe vor der ältesten.}
\end{quote}

130. Der Schalk, der treulose Nabal, verspottet ihn noch dazu, da er ihn so sehr und so unbilliger Weise betrübt und ihn an seiner hochzeitlichen Freude gehindert hatte. „Es ist nicht Sitte in unserm Lande“, sagt er. „Warum hast du aber das nicht zuvor gesagt, da Jakob zu dir sagte: ‚Ich will dir sieben Jahr um Rahel dienen‘? Ist denn diese Rede nun jetzt erst aufgekommen? wie Jener in der Komödie sagt. Warum hast du diesen Gebrauch deines Volkes nicht vorher vorhin angeführt? Der alte Schalk kommt jetzt erst mit der Sitte daher, und hat daran nicht genug, daß er seinen frommen Schwiegersohn betrogen hat, er muß ihn auch noch mit Worten verhöhnen und sein spotten, ja, er legt ihm noch durch Betrug und List den Dienst der andern sieben Jahre auch auf nur darum, daß sein Hab und Gut im Hause durch fremde Arbeit gemehret und gebessert werde auch durch Lügen und Betrug. Derhalben wäre er wohl werth gewesen, daß ihn GOtt versucht hätte, da er durch andere Leute Schweiß und schwere Arbeit reich werden will. Aber das gemeine Sprüchwort hat ihm auch gegolten und ist an ihm wahr geworden, da man sagt: Je böser Schalk, je besser Glück.

131. Man sollte aber an diesem Ort auch merken die Art zu reden in der hebräischen Sprache, da im Lateinischen steht: Non facit; auf Deutsch: Man thut nicht also. In der hebräischen Sprache wird oftmals für ein unpersönliches Zeitwort ein Activ oder ein Plural gebraucht. „Sie thun nicht so“, das heißt: „es geschieht nicht also, man thut nicht also.“ Und Luc. 6, 38. steht auch also: Mensuram bonam dabunt in sinum vestrum: „Ein voll und überflüssig Maß wird man in euren Schooß geben.“

\section[Vierter Theil]{Vierter Theil. Wie Laban von Jakob verlangt, daß er Lea behalten und um Rahel ferner dienen soll, und wie Jakob drein willigt}

\subsection{I}

\begin{quote}
\textbf{V. 27. Halte mit dieser die Woche aus; so will ich dir diese auch geben, um den Dienst, den du bei mir noch andere sieben Jahr dienen sollst.}
\end{quote}

132. Mit dem Wort „mit dieser“ meint er Lea. Da steht es. Laban hat gesehen, daß sein Hab und Gut im Hause durch den Dienst Jakobs ist gemehret worden, darum denkt er nun, wie er ihn betrügen und die andern sieben Jahre auch im Dienst behalten möge. Denn er hat seine Rechnung also gemacht: Wo ich ihm meine Tochter Rahel geben werde, so wird er alsbald davon ziehen und mein Haus verlassen, welches durch seine Gegenwart und Treue sehr zugenommen hat; darum werde ich ihn noch etwas länger behalten müssen, es geschehe gleich mit Recht oder mit Unrecht. Zwar nicht darum, daß ich ihm nicht sollte Gutes gönnen; denn ich wollte ihm wohl mehr Töchter geben, wenn ich sie hätte: sondern darum, daß ich ihn länger bei mir behalten und durch seinen getreuen Dienst möge reich werden.

133. Was ist das für ein Schalk! Es hatte Jakob dem Befehl seines Vaters nun genug gethan, daß er ihn hatte heißen ein Weib nehmen aus diesem Geschlecht, und hatte nun sieben Jahre gedienet, wie er darüber mit seinem Schwiegervater einen Vertrag gemacht hat; deshalb hat er gedacht, er wollte nun mit seinem Weibe wiederum heimziehen. Da läßt ihm der Teufel diese Beschwerung vorfallen, daß er ihm mit der Hochzeit nicht allein Ungelegenheiten macht und dieselbe verhindert, sondern ihm auch noch andere sieben Jahr auf den Hals legt.

134. So mußte nun Jakob vierzehn Jahre dienen; denn um Lea hatte er nicht gedienet: das ist ihm wahrlich schwer und verdrießlich gewesen. Er hat nun zwei Weiber und doch gleichwohl nichts Eigenes; er dient mit leeren Händen, welches sehr beschwerlich und verdrießlich ist. Denn der Geizhals Laban, der billig Nabal heißen sollte, reißt Alles an sich, und geht nur damit um, daß er durch dieses frommen Mannes Fleiß und Wohlthaten möge reich werden: er hätte ihm aber nunmehr wiederum auch sollen Gutes thun und die Aussteuer entrichten, aber er spottet sein noch dazu, gibt ihm wider seinen Willen zwei Mägde, und die eine Tochter, welche er nicht lieb hatte. Siehe, was der verfluchte Geiz thut. Darum ist er ein Räuber, ein Dieb und Mörder, der da reich wird durch den Schweiß und Blut des frommen Mannes und seiner Töchter, wie hernach folgen wird. Er hat es alles an sich gerissen. Der versuchte Mannsfrechheit wird also beschrieben und abgemalt, daß nicht Wunder wäre, wenn ihn die ganze Welt in den Abgrund der Hölle hinein versuchte: er ist ein größerer Räuber, denn sonst alle Diebe und Räuber sein mögen; und kann dieser Geiz durch keine Rhetorik genugsam herausgestrichen und groß genug gemacht werden.

\subsection{II}

\begin{quote}
\textbf{V. 28. Jakob that also, und hielt die Woche aus. Da gab ihm Laban Rahel, seine Tochter, zum Weibe.}
\end{quote}

135. Jakob nimmt die Bedingung an, so ihm sein Schwiegervater, der da geizig und ein Schalk war, vorgeschlagen hat, und wird damit also die ganze Sache ausgeglichen. Es wird aber gefragt: Warum Jakob in dies alles also eingewilligt habe? Denn wiewohl Laban den Gebrauch im Lande anführt und vielleicht in selbigem Lande auch eine solche Gewohnheit gewesen sein mag, daß man die ältesten Töchter erst hat ausgegeben: so ist es doch zumal ein loses Argument, und ist Wunder, daß Jakob diesen Gebrauch nicht sollte gewußt haben, und daß er nicht etwa ein Exempel irgend einer Ehe sollte gesehen haben, das entweder mit solcher Gewohnheit übereinstimmt, oder aber derselben zuwider gewesen wäre. Denn innerhalb sieben Jahren kann man ja die Sitten eines Landes leicht lernen. Da er aber diese Sitte anführen hört, willigt er also darein. Er hätte mit höchstem Recht Lea verwerfen mögen und zu Laban also sagen: Die Gewohnheit mag sein wie sie will, du hast mir Rahel vertraut, und ich habe dir nach unserm Vertrage sieben Jahre um sie gedienet; dieselbe will ich behalten und keine Andere.

136. Was soll aber der fromme heilige Mann thun? Er läßt es gehen; er leidet diese Gewalt mit Geduld, auf daß er Rahel noch bekommen möge, welche die einige und rechte Braut war; und er zieht wohl, daß er dieselbige nicht erlangen kann, denn allein unter dieser unbilligen Bedingung, nämlich, daß er Lea behalte, welche er nicht lieb hatte. Denn er hat gedacht: Was soll ich thun? Lea ist durch die Gewalt ihres Vaters bewogen worden, wie ich auch mit List von ihm hintergangen bin; wo ich sie nun verlassen werde, so wird sie stets verlassen sein. Denn ich habe sie beschlafen,
\quellseite{141}
und ihrer besten Aussteuer, nämlich ihrer Jungfrauschaft, beraubt, oder sie wird ja von jedermann für verdächtig gehalten werden, als daß sie ihre Keuschheit verloren habe, wenn ich mich auch schon ihrer enthalten und sie nicht erkannt hätte. Also wird Jakob durch Barmherzigkeit, Liebe und Geduld bewogen, daß er Lea nicht verlassen will. Und ist ein trefflich Exempel einer sonderlichen Barmherzigkeit und Tugend, dadurch er die bei sich behält, von der er doch weiß, daß sie ihm anstatt seiner Braut beigelegt worden ist.

137. Es wird aber nach den Rechten gefragt: Ob es sich auch geziemt habe, und ob nicht solche Barmherzigkeit im Gesetz Moses verdammt werde, der da sagt: Es soll Einer nicht zwei Schwestern zur Ehe nehmen? Die Person wird durch die Barmherzigkeit und Liebe entschuldigt, aber an der That zweifelt man noch; ja, das Gesetz will haben, er soll die Eine verlassen. Darum soll man dieses Exempel nicht dahin ziehen, daß man denselbigen also nachfolgen wollte. Jakob behält Lea und hält sie für sein rechtmäßig Weib, wiewohl er sie gern verlassen hätte; aber es hat ihn das Recht und der Gebrauch desselben Vaterlandes im Wege gelegen, und darnach auch, daß er sie beschlafen hatte. Weil er aber darein willigt, so hebt er an schuldig zu werden, und mit seiner Einwilligung bestätigt nun Jakob die Sünde. Was wollen wir aber nun dazu sagen?

138. Antwort: Es sind etliche Exempel der großen tapfern Helden, die sind hoch und ritterlich; etliche aber sind von gemeinen Sitten. Die Gesetze und guten Sitten muß man einfach halten, und dieselbigen auf keinerlei Weise übertreten lassen, auf daß keine Unordnung werde. Die großen ritterlichen Exempel aber sind, die mit den Gesetzen oder mit dem Rechte nicht übereinstimmen. Denn es kommt oftmals vor, daß irgend ein tapferer Mann, ein großer trefflicher Held, dem Gott einen sonderlichen Muth und Kraft gegeben hat, hindurchbricht und die Regel übertritt: es ist aber keinem Exempel nicht nachzufolgen. Was die gemeinen Sitten, Gesetze und Rechte mitbringen, nach denen soll man sich halten und solchen Exempeln nachfolgen. Aber soviel die hohen ritterlichen Thaten der großen tapfern Helden belanget, denselben ist nicht nachzufolgen: da gilt kein Exempel, es sei denn allenthalben gleich.

Als, wenn du Jakob wirst gleich sein, wenn ein solcher Fall kommt, wenn solche Gelegenheit und Noth vorfällt, alsdann wirst dir auch geziemen zu thun, das Jakob gemeinet hat; bist du aber Jakob nicht allenthalben gleich, so mußt du dich nach den gemeinen Gesetzen und Sitten halten.

139. Man muß nicht mit Eifer hindurchbrechen, und muß kein Exempel machen, dem man alsobald nachfolgen wollte, um irgend eines Falles willen, der sich mit einem tapfern Helden und ritterlichen Manne möchte zugetragen haben. Ein Grammatiker sagt: Die Wörter, darin der letzte Buchstabe ein a ist, sind weiblichen Geschlechts; die aber auf ein um ausgehen, sind sächlichen. Wenn du aber da die Exempel oder Wörter, als da ist, auriga, oder glycerium, willst hinzuziehen, dieselbigen werden sich mit der Regel nicht reimen, wiewohl sie eben dieselbe Endung haben. Denn es sind sonderliche Wörter, die dahin nicht gehören und den Regeln nicht unterworfen sind, sondern die Regeln sind ihnen unterworfen. Ein Poet fragt nach seiner Regel. Dasselbe muß man auch in der Medizin und in Rechten merken. Wenn ein Arzt allezeit stracks nach der Regel fortfahren will, so tödtet er viel Leute. Man muß zwar den Regeln folgen, aber nicht so gar tollkühn und mit Eifer. Denn es kann irgend ein solcher Fall kommen, der der Regel nicht unterworfen sein will. Da muß man auf die geometrische Proportion sehen, das ist, man muß auf die Person und nicht auf die That sehen.

140. Die Mönche sind in ihren Regeln der arithmetischen Proportion gefolgt, das ist, sie haben Einem gethan wie dem Andern: so viel Einer arbeitet, ißt, trinckt und schläft, eben so viel soll der Andere auch arbeiten, essen und schlafen. Aber diese arithmetische Gleichheit gehört auf den Markt, wo man kauft und verkauft, und nicht in das Regiment. Denn also sind viel Mönche, die nicht gleiche Last des Ordens, wie die Regeln mitbringen, haben tragen können, getödtet worden. Einer kann sich genügen lassen, wenn er sieben Stunden schläft; der Andere läßt sich an drei Stunden genügen; ein Anderer aber hat kaum an zehen Stunden genug.

141. Darum hat Augustinus weislich geredet, da er diesen schönen herrlichen Spruch setzt: Non aequaliter omnes, quia non aequaliter valetis omnes, daß ist: Lieben Brüder, ihr müßt nicht Alle gleiche Last tragen, denn ihr seid nicht Alle gleich stark. Und diesen Spruch hat auch der Bischof von Worms gelobt. Denn er reimt sich fein zu der geometrischen Proportion, die nicht ein Ding mit dem andern gleichstellt, wie es auf dem Markte pflegt zuzugehen; sondern vergleicht die Personen mit einander, und nach denselben ordnet und schickt sie die Sache. Die arithmetische Proportion auf dem Markte richtet Alles ohne Ansehen der Personen, das ist, es sei gleich ein Weib, oder ein Mann, oder aber ein Kind, so verkauft sie ein Brod um einen, zwei oder drei Pfennige. Aber mit dem Essen wird die Proportion oder Vergleichung der Personen gehalten. Man findet etwa im Kloster Einen, der sich an Einer Semmel genügen läßt; man findet auch wohl einen Andern, der sich daran nicht kann genügen lassen.

142. Dieses sage ich also zum Exempel. In der Theologie aber und in der Schrift muß man sich davor viel mehr hüten, daß Einer nicht also argumentire oder schließe: Jakob bricht hindurch und thut wider die gemeinen Sitten und das Gesetz Moses, darum geziemt mir dasselbe auch zu thun. Nein, lieber Bruder, es heißt: Nego consequentiam: Das will noch lange nicht folgen. Und ist das die Ursache: weil Jakob nicht eine solche Person ist, die an die Sitten gebunden; sondern er ist ein sonderlicher tapferer Held, der geführt und regieret wird nicht nach den Sitten oder Weise im Gesetz, sondern richtet sich nach einer sonderlichen ritterlichen Weise eines tapfern Helden: darum soll man ihm nicht nachfolgen, sondern man soll ihn dermaßen ansehen, daß man sich seiner verwundere.

143. Du sollst dem Achilles darin nicht nachfolgen, daß er Hector erschlagen hat, oder wenn du dasselbe anfangen wirst, so wird Hector dich erwürgen. Du sollst dich nicht unterstehen, den Himmel oder die Säulen des Hercules zu tragen, weil du kaum so viel Stärke hast, daß du einen Stecken tragen kannst. Das sind die Narren im weltlichen Regiment, die den tapferen Helden nachfolgen wollen, die sie über und außerhalb der Regel einher gehen und sonderliche Thaten ausrichten sehen. Der große Mann, sagen sie, hat also gethan, hat also gedacht, darum will ich dem auch also nachfolgen.

144. Du hast aber einen andern Spiegel, darein du sehen sollst, nämlich, die gemeinen Sitten, dergleichen das gemeine Recht und Gesetz; das sollst du nicht übertreten. Und brich dich nicht heraus, ehe denn dich Gott beruft und dich heißt einen tapfern Helden oder Jakob sein; sonst werden aus solchen Leuten eitel Affen, die beide sich selbst und auch Andere in Schaden führen. Wie jener Mordbrenner von Natur kleinmüthig gewesen, und da er ein Affe geworden und anderer großer Helden tapfern Thaten nachfolgen wollte, hat er damit seine Thorheit an den Tag gegeben und sich selbst verderbet. Denn solche närrische Menschen richten Anderes nichts aus, denn daß sie dem gemeinen Nutzen Schaden thun, und die Welt irre machen und gar verkehren.

145. Derhalben soll man von diesen Exempeln also handeln, daß gleichwohl die gemeinen Sitten unverletzt bleiben, denen man allein nachfolgen soll. Und die Exempel oder Thaten der großen Helden sollen wir ansehen und uns derselben verwundern; aber wir sollen denselben nicht nachfolgen. St. Paulus schlendert die Otter in das Feuer, die ihm an seine Hand gefahren und daran hangen geblieben war, wie in der Apostelgeschichte steht am 28. Capitel V. 5. Thue du es ihm nach.

146. Und man erzählt eine Fabel von einem Arzte, der ein Kindlein über der Taufe gehalten und Gevatter geworden war, und hatte dasselbe die Worte von der Taufe hören sprechen mit dem Gebet und Sprüchen aus der Schrift, darin angezeigt wird, daß wir aus dem Reich des Teufels erlöset und in das Reich des Sohnes Gottes versetzet werden. Da verwunderte er sich der großen Verheißung und Wohlthaten, so uns in der Taufe gegeben und vorgetragen werden, und sagte: Wenn ich wüßte, daß ich auch also getauft wäre und daß dieselbigen Worte über mich auch gesprochen wären, so wollte ich mich vor dem Teufel nimmermehr fürchten. Und da ihn Andere des erinnert und überredet haben, daß er auch eben mit diesen Worten getauft sei, hat er einen solchen großen Glauben gefaßt, daß er ganz und gar bei sich beschlossen, er wollte sich nimmermehr vor dem Teufel fürchten. Endlich aber ist ihm ein Gespenst vorgekommen und der Teufel ist ihm in Gestalt eines Bocks erschienen; davon aber ist er so gar nicht erschrocken, daß er auch das eine Horn genommen und zerbrochen hat, und hat es seinen Gesellen gebracht, von denen er gefangen war. Da nun dasselbe einer von seinen Gesellen gesehen hat, wiewohl er denselben Glauben nicht gehabt, hat er sich doch das große Wunderwerk bewegen lassen, und hebt an, sich auch zu rühmen, daß er getauft sei, und wolle sich vor keinem Gespenst, wie greulich es auch sein möchte, fürchten. Geht derhalben an denselbigen Ort und das Gespenst des Bocks mit den Hörnern kommt wieder. Da er aber nun den Bock bei den Hörnern angreifen wollte, nimmt ihn der Teufel bei den Haaren und erwürgt ihn. Warum ist nun Jener der Gefahr entronnen und Dieser nicht? Antwort: Das ist darum geschehen, daß an diesem große Vermessenheit und kein rechter Glaube gewesen ist; er hat gemeint, er könnte den Teufel mit seinen Kräften und Stärke überwinden. Er hat keinen rechten Glauben gehabt.

\quellseite{142}

147. Daher soll sich ein Jeder selbst prüfen, was für Gaben er habe. Denn wie wir am Leibe, am Verstande und am Gute einander ungleich sind: also sind wir auch mit den geistlichen Gaben einander nicht gleich. Ein Jeder soll in seinem Berufe bei dem Sittengesetz und bei dem gemeinen Rechte bleiben, bis ihn Gott berufen oder zwingen wird, daß er etwas Sonderliches thun soll. Jakob hat dieses nicht aus eigenem Rath oder Willen gethan, sondern er ist durch Gottes Willen dazu hingerissen und gezwungen worden, daß er dieses ohne seinen Willen, und daran er nie gedacht, hat thun müssen. Darum bestätiget Gott diese That hernach, daß er ihn mit Kindern segnet und Lea eine Mutter wird der Patriarchen.

148. Also soll man nun die Väter entschuldigen, aber ihrem Exempel soll man nicht nachfolgen, auf daß nicht eine Unordnung und Zertrennung daraus folge der Gesetze und gemeinen Sitten; sondern wir sollen bei den gemeinen Gesetzen bleiben und dieselben nicht übertreten, es sei denn, daß wir durch einen sonderlichen Beruf davon abgezogen werden, oder daß wir einen solchen starken Muth fühlen, wie etwa die großen tapfern Helden einen sonderlichen Trieb gefühlt haben.

149. Wo aber nun irgend ein Schwärmer aus dem Exempel Jakobs also schließen wollte: Jakob hat zwei Schwestern zur Ehe genommen, darum mag ich dasselbe auch thun: darauf antworte ich also: Nego consequentiam: Das will noch lange nicht folgen. Denn man muß die Personen unterscheiden, deren etliche gemeine Personen sind, die andern aber sind große tapfere Helden, die da sonderliche Bewegungen haben, dadurch sie getrieben werden, sonderliche Thaten zu thun. Denn Gott ist der Navigator und der oberste Regent über uns alle: wenn der irgend Einem etwas wird eingeben, daß wider die gemeine Regel ist, dadurch soll die Regel nicht gebrochen oder aufgehoben werden; sondern man soll immer bei den gemeinen Rechten bleiben, und soll es Gott lassen befohlen sein, daß der Herr heranziehe, welche er wil. Denn er ist es allein, der da recht sagen kann: Sic volo, sic jubeo, sit pro ratione voluntas: Also will ich es haben, also heiße und gebiete ich es, daß und kein Anderes. Er ist es allein, der die Gesetze ändern kann. Wenn er Einen von der Regel ausnimmt, so ist er ausgezogen, sonst nicht.

150. Also hat der Herr, der Münzer, die Historien der Könige und Richter gelesen; hat auf der Kanzel den Bauern eingeschärft die Geschichten Josuas, Simsons, item Davids, und hat ihnen derselben Exempel vorgehalten. Ihr seid Gottes Volk, schrie er, darum sollt ihr dem Exempel der großen Heiligen nachfolgen, als, Simsons und Josuas; sollt die Fürsten todtschlagen, und das weltliche Regiment ändern und in einen andern Stand bringen. Aber die Schlußfolgerung ist falsch und gilt nicht. Denn solche Männer, die so tapfere Helden sind, werden von der Regel ausgenommen, und wir, die der Regel unterworfen sind, können und sollen ihnen auch nicht nachfolgen.

\begin{quote}
\textbf{V. 29. 30. Und (Laban) gab seiner Tochter Rahel seine Magd Bilha zur Magd. Also lag er auch bei mit Rahel, und hatte Rahel lieber, denn Lea; und dienete bei ihm fürder, die andern sieben Jahre.}
\end{quote}

151. Nachdem nun die tapfern und sonderlichen Thaten und wunderbaren Dinge beschrieben sind, wendet sich nun der Heilige Geist zu der Schwachheit, und die Dinge zu beschreiben, so die gemeinen Sitten belanget. Denn wiewohl ein großer Held von der Regel ausgenommen wird, so muß man sich doch immer wieder nach der Regel richten. Wie in der Grammatik, obwohl das Wort poeta nicht weiblichen Geschlechtes ist, wie man nach seiner Endung annehmen sollte, dieweil der letzte Buchstabe ein a ist, so wird es doch gleichwohl declinirt, wie es in der Grammatik heißt, nach der Regel der Wörter, die sich auf ein a enden. Also muß ein großer Held sich auch halten, daß er die gemeinen Sitten nicht zerbreche, ja, daß er sie vielmehr bewahre, vertheidige und dieselben in Schwange helfe erhalten. Wiewohl er für seine Person auf eine sonderliche Weise ausgenommen ist, so soll er sich dennoch wiederum nach der ersten Declination, wie man in der Grammatik redet, und nach der Regel richten. Die Poeten sagen von Hercules, daß er die Ruder zerbrochen habe im Schiffe der Argonauten, und ist doch gleichwohl nach seinem Rathe Alles glücklich fortgegangen: also brechen wohl die großen tapfern Helden durch die Gesetze hindurch, aber sie zerreißen gleichwohl dies gemeine Regiment nicht.

152. So wird nun jetzt dieser gute fromme Mann Jakob beschrieben wie ein Narr; denn die Liebe, so er zu Rahel trägt, bringt ihn dazu, daß er die andern sieben Jahre sich auch verpflichtet, Laban zu dienen. Das würde ein großer tapferer Held nicht gethan haben, sondern der hätte gesagt: Ich habe dir um Rahel gedienet, die will ich haben, daß sie mir soll gegeben werden als meine Braut; und sollte ihm noch wohl gedroht haben zu schlagen, oder würde die Jungfrau, so ihm verlobt war, mit Gewalt weggeführt haben. Welches er denn auch mit allem Recht hätte thun mögen. Aber der große Mann läßt sich herunter, und begibt sich in die äußerste Knechtschaft und dient dem Laban ganze vierzehn Jahr um die eigene Rahel. Denn Lea halben hat er mit ihm keinen Vertrag gemacht. Zuvor war er in den Himmel gestiegen, und hat vor Gott das thun dürfen, das sonst keinem Andern geziemt hätte; jetzt demüthigt er sich wiederum also, daß er sich in Knechtschaft begibt.

153. Wir sollen aber diese Exempel gern lesen, darin uns der Heilige Geist die heilige Schrift hat rühmen und loben wollen, daß wir uns dieselbe sollen lassen befohlen sein, und hat uns an den großen Helden auch fleischliche und schwache Werke zeigen wollen. Denn es ist ein größerer Trost, wo man an den Heiligen solche fleischliche Dinge, denn wo man die hohen und vortrefflichen Thaten an ihnen sieht, denen wir nicht nachfolgen können.

154. Darum will der Heilige Geist nicht haben, daß man es verachten soll, sondern er lobt diese häuslichen und gemeinen Werke dieses Lebens so sehr, daß er auch bisweilen etliche unvollkommene Werke, und da Mängel und Fehler mit unterlaufen, erzählt, daß er damit anzeige, daß er geduldig sei und die Weise der Schwachen dulden könne: daß er die kleinen Kinder nicht verwerfen wolle darum, daß sie etwas Blattern haben, unschlägig und unselig sind, sofern sie nur im Glauben und in der Liebe beharren und beständig bleiben, wie Paulus ihnen sagt: Er hat Geduld mit unserer vielen Schwachheit. Wie auch in der Epistel an die Hebräer am 4. Capitel V. 15. steht: „Wir haben nicht einen Hohenpriester, der nicht könnte Mitleiden haben mit unserer Schwachheit“ etc. Dieses dient nun zum Trost der Kirche, darin nicht allezeit ritterliche Männer und große Helden sind, oder ihrer sind ja nicht viele. Aber dieselbigen lassen sich auch wohl in genere feminino decliniren, wie man in der Grammatik sagt, ob sie wohl generis masculini sind.

155. Also läßt Gott Jakob verstoßen werden, und eine solche Knechtschaft leisten, die schimpflich ist und einem freien Menschen nicht wohl ansteht; und dasselbe geschieht aus fleischlicher Liebe, die er zu der Braut hat, darauf der alte Buhler gar verstürzt ist, und stellt sich so gar leichtfertig, daß wiewohl er jetzt schon seine Rachel mit allem Recht weg hat, begibt sich doch gleichwohl, der Jungfrau zu Gefallen, die andern sieben Jahre in die Knechtschaft begibt. Droben (V. 20.) hat er gesagt, daß er um die sieben Jahre gedäucht haben, als wären es einzelne Tage vor großer Liebe: jetzt, da er nun dessen gewährt ist, daß er sie zum Weibe bekommen hat, welche er so heftig lieb gehabt, begibt er sich noch in die schwere Knechtschaft anderer sieben Jahre. Ist das nicht kindisch, fleischlich und gar närrisch Ding, daß er um der Jungfrau willen, die ihm vertraut und nach allem Recht gegeben war, sich selbst freiwillig in solche Knechtschaft steckt, und seine eigene Hausorge noch immer verschiebet und liegen läßt? Denn das ist ja billig und nöthig, wenn du ein Weib genommen hast, daß du denkest, wie du nun dein eigen Dach und Haus haben und anordnen mögest. Da läßt Jakob der Jungfrau zu Gefallen die Sorge gar liegen, nimmt sich derselben nicht an, und dient mittlerweile dem schändlichen Ungeheuer, dem scheußlichen, geizigen und undankbaren Manne Laban.

\quellseite{143}

156. Solche Leute sind nun die großen Helden, und also werden sie beschrieben uns zum Exempel, auf daß wir wissen, daß sie auch Menschen gewesen sind wie wir, nämlich, bald stark und beständig, bald niedrig und schwach. In den hohen und ritterlichen Werken sind sie unüberwindliche Löwen, die keinen Gleichen haben; in häuslichen Werken sind sie Narren und schwache Leute. Gott steht aber auf beiderlei, und läßt sich den ebensowohl gefallen, der mit geringen Werken umgeht, als der hohe große Werke thut.

\section[Fünfter Theil]{Fünfter Theil. Von Lea und ihrer Fruchtbarkeit; item, von Rahel und ihrer Unfruchtbarkeit}

\begin{quote}
\textbf{V. 31. Da aber der Herr sahe, daß Lea unwerth war; machte er sie fruchtbar, und Rahel unfruchtbar.}
\end{quote}

157. Dies ist sehr fleischlich Ding: der Herr muß in das Spiel sehen. Hat aber Gott Anderes nichts zu schaffen, denn daß er auf diese Niedrigkeit in der Haushaltung sehen müsse? Wo er so viel übrige Zeit hat, so wird er wohl zu schaffen finden, daß er sich wider den Teufel, der ein Gott ist dieser Welt, oder wider die großen Monarchen und Fürsten dieser Welt auflehne. Siehe, was thut er mit dem Mädchen Lea? Antwort: Ein Jeder soll wissen, daß Gott in seinem Beruf für ihn sorge. Denn Gott sorgt für das, so gering, mittelmäßig und groß ist: er ist ein Schöpfer und Regierer über Alles. Wer wollte aber glauben, daß Gott sollte für Lea sorgen? Es sollen es wahrlich alle Menschen gewißlich dafür halten, daß Gott Achtung auf sie hat und für sie sorgt, wie gering, verachtet und niedrig sie auch sind. Denn darum malt uns Mose so fleißig ab und stellt uns vor die Augen, wie es um die Haushaltung Jakobs bestellt gewesen.

158. Jakob hat sein Weib, darnach er so großes Verlangen gehabt, bekommen, und Lea dazu, die er nicht lieb gehabt, nach der er auch nichts gefragt hat. Dasselbe versieht Rahel wohl, und hat die Schlüssel und das Regiment über das ganze Gesinde, ist allezeit bei Jakob. Sie ist das liebe Mäuschen. Auf die Frau aber, die das Regiment hat, sieht das ganze Haus und die Nachbarschaft, Alles richtet sich nach ihr, und wissen, daß Lea verachtet ist und daß auch Jakob ihrer nicht achtet. Denn daß er sie lieb hat, das thut er darum, daß er sich ihrer erbarmt, thut es nicht aus ehelicher Liebe, wie ein Bräutigam seine Braut liebt: es ist nur ein Almosen. Derhalben sitzt die arme Lea mit ihrer Magd in ihrer Hütte, ist elend und betrübt, bringt die Zeit zu mit Spinnen und Weinen. Denn das andere Gesinde und vornehmlich Rahel hat sie gar verachtet, dieweil sie von ihrem Manne ist verachtet worden, der Rahel vorgezogen und dieselbe sehr lieb gehabt. Sie ist nicht schön, ist nicht angenehm, sondern gar veracht und verabscheut, auf Hebräisch snuah; da sitzt das arme Mädchen und man achtet ihrer nicht. Rahel ist hoffärtig, sah sie ungern einmal an. Ich bin die Frau im Hause, hat sie gedacht, Lea ist eine Magd. Dieses sind recht fleischliche Dinge an den heiligen Vätern und Müttern, wie sich in unsern Häusern solche Dinge auch pflegen zu begeben.

159. Wir dürfen aber nicht denken, daß Lea solche Verachtung ohne großen Schmerzen, ohne viele Thränen, Heulen und Weinen habe dulden und ertragen können. Denn das weibliche Geschlecht ist zumal ein schwach Werkzeug, welches von dem Manne heftig will geliebt sein, oder aber will ja zum wenigsten nicht verachtet sein, oder vornehmlich vom Manne, oder vom Gesinde im Hause. Darum hat es Lea sehr wehe gethan, daß sie im Hause so gar verstoßen und verachtet hat sein sollen, und daß Rahel mit dem Gesinde das Regiment allein gehabt und ihr gleichsam getrozet hat. Denn niemand ist, der Achtung auf sie habe und etwas nach ihr frage, auch Jakob selbst nicht. Aber allein der Herr, der muß in das Spiel sehen, muß dem Volk seinen Affect brechen und crucifizieren, dazu auch der Rahel ihren Muth brechen, nämlich, daß er Lea erhebt und vorzieht, die sonst gehaßt und gar verachtet war.

160. Also soll man lernen, daß Gott alles sieht und regiert, aber doch einzig und sonderlich auf dasjenige sieht, welches verachtet und verlassen ist; wie Christus am Kreuze selbst auch ist verachtet gewesen. Denn was die Welt und auch alle heiligen Leute, als Jakob, Rahel, wegwerfen, das lieset er auf und ist ihm eitel Heilthum. Dieses ist ein wunderbarer und großer Trost für die armen, betrübten und elenden Herzen, die sich damit zum Glauben und Hoffnung erwecken sollen; zumal sie hören, daß Gott also abgemalt und beschrieben wird, wie Judith am 9. Capitel V. 13. steht: „Allemal hat dir gefallen der Elenden und Demüthigen Gebet.“

161. Gott ist selbst auch ein sehr großer Held und ist seiner Regel unterworfen; doch lenkt er sich und läßt sich also hernieder, daß er nirgend so scharf sieht, als wenn er auf das, so niedrig und gering ist, schaut. Denn also ist Lea sehr verachtet, und jedermann unangenehm und verhaßt gewesen: wo sie etwa dem Gesinde irgend was geheißen hat, werden sie gesagt haben: Was hast du mir zu gebieten? fragt doch Jakob und Rahel nichts nach dir. Und vielleicht mag auch Rahel mit ihrer Hoffahrt solchen Ungehorsam bei dem Gesinde genähret haben; daß also Lea beide vom Herrn und von der Frau, und dazu vom Hausgesinde ist verachtet worden. Dasselbe hat sie aber in ihrem Herzen sehr betrübt und bekümmert gemacht.

162. Denn das thut einem frommen gottseligen Weibe sehr wehe, wenn sie sieht, daß sie von ihrem Manne soll verachtet werden. Wo sie von Andern beleidigt wird, kann sie es so viel desto leichter ertragen und überwinden, wenn sich nur der Mann freundlich um sie annimmt; denn das ist des Weibes Leben, wenn sie der Mann herzlich lieb hat. Welches ich gleichwohl von frommen ehrlichen Weibern will verstanden haben. Wenn der Mann sauer, seltsam, wunderlich und zornig ist, dasselbe thut einem gottseligen Weibe wehe, denn der Tod selbst. Denn sie will gerne vom Manne erfreuet werden, und fürchtet sich vor Verachtung, als vor der Finsterniß und vor dem äußersten und größten Unglück.

163. Diese Schmerzen und Betrübniß hat nun die arme Lea tragen und leiden müssen, ist darüber betrübt und bekümmert gewesen, daß sie von ihrem Herrn sollte verachtet werden. Höre aber, was Mose sagt: „Der Herr sahe“ etc., sonst sieht es niemand. Gott sieht und segnet die verachtete Lea, die da weint, verhaßt und betrübt war; aber die da hoffärtig war, demüthigt er. Da geht nun ein ander Licht auf der, die da sitzt in Finsterniß; denn Jakob wird mit seiner Liebe und Rahel in ihrer großen Ehre und Herrlichkeit gekreuziget. Der Herr schließet der Rahel den Leib zu und läßt sie unfruchtbar bleiben; welches der Rahel nicht ein geringes Kreuz, ja, der Tod selbst gewesen ist, wie sie sich dessen hernach hart beklagen wird. Er ehrt aber die, so verstoßen und verachtet war, damit, daß er sie fruchtbar und zu einer Mutter macht; denn das heißt den Leib zuschließen oder öffnen, nämlich, Kinder und Früchte des Leibes geben, wie die Erde geöffnet wird, wenn sie blüht.

164. Das heißt die leibliche Wort, und heißt die Verachtung derer, die verachtet und verstoßen sind. Und ihr sollt das fleißig merken, daß Gott genannt wird ein Anseher der Armen, das ist, derer, die verachtet und verworfen sind, allein, daß sie nicht verzagen. Dieses schreibt der heilige Geist den Betrübten zum Trost, nämlich, daß Gott Lea segnet und erhöht, daß sie schwanger wird und Kinder gebieret.

\begin{quote}
\textbf{V. 32. Und Lea ward schwanger, und gebar einen Sohn; den hieß sie Ruben, und sprach: Der Herr hat angesehen mein Elend; nun wird mich mein Mann lieb haben.}
\end{quote}

165. Das arme elende Weib, die da blöde Augen gehabt und die sie auch von wegen der Bekümmerniß mit Weinen verderbet hat, glaubt dennoch, daß sie von Gott angesehen werde. Der Herr, sagt sie, hat gesehen, daß ich verachtet bin. Sie führt kein Verdienst an, sondern nur das, daß Gott die schwach, das ist, die Verachtete angesehen habe. Denn dieses ist der rechte Titel und die eigentliche Definition und Beschreibung unsers Herrn Gottes, nämlich, daß er ein solcher Gott sei, der auf die Verachteten und Elenden sieht. Er sieht auch auf die Hoffärtigen und großen Hansen in der Welt, aber die bedürfen seiner Gnade und Barmherzigkeit nicht, sie verachten sein Ansehen; darum muß er über sie hinsehen, wie der 138. Psalm V. 6. sagt: „Der Herr kennet den Stolzen von ferne.“ Es läßt sich aber ansehen, als sei der Vers im Lobgesange der Jungfrau Maria, da sie Luc. 1, 48. sagt: „Er hat die Niedrigkeit seiner Magd angesehen“, aus diesem Text glossier genommen; denn es sind eben dieselben Worte, die an diesem Ort stehen.

\quellseite{144}

166. Das hebräische Wort ani können wir nicht eigentlich genug verdeutschen. Mose wird im 4. Buch am 12. Capitel V. 3. genannt: vir an, ein geplagter Mensch, das ist, der elendeste Mensch auf Erden; wie St. Paulus redet 1 Cor. 15, 19.: „Hoffen wir allein in diesem Leben auf Christum, so sind wir die elendesten unter allen Menschen.“ Auf Lateinisch heißt es miseria, auf Deutsch Elend. Also singt die Jungfrau Maria: Er hat mein großes Elend angesehen; denn sie war ihrer Eltern beraubt, war eine arme Waise und verlassen, hat andern Leuten dienen müssen; da kommt Gott und sieht eine solche arme elende Magd an.

167. Und Lea sagt an diesem Ort auch: Jakob hat mich nicht angesehen, denn allein, was er aus Gnaden gethan, daß er sich meiner erbarmt hat; Rahel hat mich gar verachtet; darum bin ich arm, elend und verstoßen, auf Hebräisch snuah; aber der Herr hat mich angesehen und hat mir einen Sohn gegeben. Derhalben dankt sie Gott und lobt ihn, und sagt noch dazu: „Nun wird mich mein Mann lieb haben“, nun wird der Haß und die Verachtung in Liebe und Ehre verändert werden; denn ich bin durch Gottes Segen eine Mutter geworden, welches gleichwohl meiner Schwester nicht widerfahren ist. Und deshalb bin ich viel glücklicher denn sie; nicht daß ich solches verdient hätte, sondern weil mich der Herr angesehen hat.

168. Es ist aber allezeit für einen großen Segen Gottes gehalten worden, wo die Weiber fruchtbar sind und Kinder gebären, auch bei den Heiden, und die Weiber sind derhalben in Ehren gehalten worden. Zu Rom ist es den Männern verboten gewesen, Gold zu tragen, den Weibern aber hat man es erlaubt.

169. So ist nun Lea eine große Herrin, weil sie fruchtbar ist und Kinder gebieret vor ihrer Schwester her, welcher jedermann gewünscht und von der man gehofft hatte, daß sie erst gebären würde. Da aber nun Rahel solches gehört, ist sie ohne Zweifel bewogen worden, ihre Schwester zu neiden und hassen, und das ganze Gesinde hat sich auch solcher wunderlichen Veränderung verwundert: aber Gott hat es gethan, der auf das Niedrige und was verachtet ist, sieht.

170. Wiewohl nun dieses alles lauter fleischlich Ding ist, und das zum Hausregiment gehört, und sich ansehen läßt, daß gar nichts Geistliches daran sei, so ist es doch in Wahrheit alles eitel geistlich Ding; wie man aus den Worten der Lea zu sehen ist, die einen sonderlichen großen Trost in sich halten. Denn Gott ist nicht ein solcher Gott, der allein für die Engel sorge und dieselben regiere; sondern er sorgt und regieret auch das Hausregiment, und alle Dinge, so zum Hausregiment gehören, der die Herzen und Urtheile der Menschen also ändern kann, daß die kurz zuvor ganz verachtet und verstoßen gewesen sind, darnach die Allerliebsten und Angenehmsten werden.

171. Und darum ehrt Lea ihren Sohn mit diesem Namen, daß sie ihn Ruben nennt, zum öffentlichen Zeugniß wider die, so sie verachtet und gehaßt hatten. Als wollte sie sagen: Nun seht ihr selber, daß ich von Gott nicht verworfen bin; Gott hat mich nicht verachtet, wie ihr mich verachtet habt. Reu, seht; rau: Sie haben den Sohn gesehen. Siehe, will sie sagen, da habt ihr meinen Sohn, beschauet ihn; denn Gott hat mich angesehen, die ich vor euch snuah, das ist, verachtet gewesen bin: darum heiße ich ihn Ruben, daß ihr aufhöret mich zu hassen und zu verachten; er ist ein Sohn des Gesichtes. Gott siehet, ich sehe, ihr sehet auch, daß ich nicht verstoßen und verachtet bin, wie ihr mich zuvor habt angesehen. Also tröstet der Herr die arme elende Lea, die gar verhaßt war, und sie bekennt es und dankt Gott dafür, und tröst nun wider alle ihre Verächter. Das wird ohne Zweifel Rahel sehr verdrossen haben, welche von Jakob sehr geliebt und allezeit für die rechte Hausmutter gehalten worden, und bleibt doch unfruchtbar; ja, es werden der Lea noch drei Söhne gegeben, wie im Text folgt.

\begin{quote}
\textbf{V. 33. Und (Lea) ward abermal schwanger, und gebar einen Sohn, und sprach: Der Herr hat gehöret, daß ich unwerth bin, und hat mir diesen auch gegeben. Und hieß ihn Simeon.}
\end{quote}

172. Vorhin hat der Herr gesehen, nun hört er auch. Denn ich halte dafür, Rahel werde noch nach ihrem Stolze mit ihr sein umgegangen, unangesehen, daß sie nun einen Sohn geboren und daß Gott selbst sie angesehen hat. Was ist das? wird Rahel und das Gesinde gesagt haben; wer weiß, ob sie auch mehr Söhne gebären wird? Es kann wohl geschehen, daß dieser Sohn bald dahin stirbt, und daß Rahel als die rechte Hausmutter auch Kinder gebären wird.

173. Also hatte Lea die Verachtung noch nicht gar überwunden und war derselben noch nicht gar überhoben, sondern sie verachten noch mit Gewalt und wider allen Witz und Vernunft den Segen und das Ansehen Gottes. Darum wiederholt Lea ihr Gebet zu Gott und sagt also: Lieber Herr Gott, ich habe nun einen Sohn geboren, sie fragen aber nach dieser Geburt nichts; ich bin nochmals nicht angenehm genug um des einigen Sohnes willen: gib mir noch einen, auf daß endlich ich auch für die Hausmutter möge gehalten werden. Jakob hat sie zwar nicht gar verachtet, sondern hat sie geliebt; aber der Haß des Hausgesindes und der Rahel ist noch nicht überwunden. Der erstgeborne Sohn, sagt sie, sollte mir ja Gunst und Ansehen zuwege gebracht haben; aber ich bin noch immer verachtet.

174. Derhalben erhört nun Gott das Gebet und gibt ihr den andern Sohn, nämlich Simeon. Da entfällt nun Rahel der Muth und sie hebt nun an, demüthig zu werden und gleichsam zu sterben, wie man hernach sehen wird, da sie sagt: „Schaffe mir Kinder; wo nicht, so sterbe ich.“ Denn sie ist noch unfruchtbar und Lea geht schon mit ihrem dritten Sohne schwanger.

\begin{quote}
\textbf{V. 34. Abermal ward sie schwanger, und gebar einen Sohn, und sprach: Nun wird sich mein Mann wieder zu mir thun, denn ich habe ihm drei Söhne geboren. Darum hieß sie ihn Levi.}
\end{quote}

175. Lea will also sagen: Nachdem mich nun mein Mann lieb gewonnen, wird er nun auch die Hütte Rahels verlassen; denn er hat bisher bei ihr gewohnt als bei seinem geliebten Weibe, und ist zu mir nicht gegangen, denn nur allein, daß er sich meiner erbarmt hat. Nun wird er sich aber auch zu mir thun und die unfruchtbare Rahel verlassen. Ich weiß zwar wohl, daß er mich jetzt lieber hat denn zuvor, er hat sich aber noch im Bett und am Tisch nicht zu mir gehalten; die Gnade und Ehre hat bisher Rahel gehabt: jetzt aber wird er sich nun zu mir thun und bei mir wohnen. Darum nennt sie den Sohn Levi, das ist, Vereinigung oder Zugesellung. Als wollte sie sagen: Ich will das zuwege bringen, daß er sich nun zu mir gesellen soll im Bette und am Tische. Das hat aber Rahel sehr wehe gethan, und hat gedacht: Nun werde ich müssen snuah, das ist, verachtet sein; mein Mann wird mit mir zürnen und sagen: Du bist wohl meine liebe Braut gewesen, der Herr aber segnet dich nicht, sondern Lea, und will haben, daß ich Lea lieben soll. Da ist es mir im Hause Jakobs wahrlich sehr wunderlich gegangen, ja, es hat sich darin gar umgekehrt.

\begin{quote}
\textbf{V. 35. Zum vierten ward sie schwanger, und gebar einen Sohn, und sprach: Nun will ich dem Herrn danken. Darum hieß sie ihn Juda. Und hörte auf Kinder zu gebären.}
\end{quote}

176. Also wollte sie sagen: Ich hätte mir an einem, an zwei und aber sonderlich an drei Söhnen genügen lassen, darüber aber bekomme ich auch noch den vierten Sohn; nun lobe und danke ich dem Herrn, ja, jetzt werde ich nun auch gelobt, vorgezogen und die liebe Hausmutter werden. Ich bin nun meiner Niedrigkeit und Elende erhöhet, und meine Schwester ihrer Ehre und Herrlichkeit beraubt worden. Derhalben wird nun Rahel begehren zu sterben vor Zorn und Ungeduld, und wird zu ihren Mann sagen: „Schaffe mir Kinder; wo nicht, so sterbe ich“ 20.

\quellseite{145}

177. Bisher aber haben wir nun gehört, wie Jakob nach der Hochzeit, so er mit Lea und Rachel gehalten, vier Söhne gezeugt habe von Lea, die verachtet und verstoßen war. Mit welchem Exempel GOtt anzeigt, daß er den Elenden, und die betrübt und verachtet sind, gnädig und günstig sei, daß sich Alle, die mühselig und beladen sind, nicht zu beklagen haben, vornehmlich wenn sie an GOtt glauben und sich die Ungeduld und den Zorn nicht überwinden lassen. Denn es ist viel besser, betrübt und geplagt sein, wo GOtt lacht und gnädig dabei ist, denn groß Glück haben, und daß GOtt dabei zürnt und ungnädig ist.

178. Denn was ist die Welt mit ihrer Gunst, Freude und aller ihrer Wohllust gegen das gnädige und freundliche Ansehen, damit GOtt die Elenden und Niedrigen ansieht? Über das Fleisch läßt es nicht zu, daß wir solches verstehen oder glauben könnten; ja, es murret zum öftern um des Unglückes willen, so gegenwärtig ist, und untersteht sich, die erschrockenen betrübten Herzen vom Glauben und Trost mit Gewalt abzutreiben, und dahin zu bringen, daß sie wider GOtt murren und mit ihm zürnen sollen. Denn so viele und so mancherlei Widerwärtigkeiten machen, daß wir in unserer Schwachheit sehr verzagt und betrübt werden.

179. Und, wollte GOtt! daß wir lernen und uns gewöhnen könnten an den Kampf wider das Fleisch, den St. Paulus lehrt. Denn wo wir schon nicht alles also ergreifen und thun können, was der Geist haben will, so würden wir doch gleichwohl das verstehen und empfinden, daß uns GOtt nicht will verlassen, und nicht gestatten, daß wir vom Fleische überwunden werden sollten. Ich wollte wahrlich mich gerne mit St. Paulo freuen, rühmen und triumphiren im Tode, im Kreuz und Leiden; doch „der Geist ist willig; aber das Fleisch ist schwach“, Matth. 26, 41.

180. Derhalben sollen wir uns daran genügen lassen und damit trösten, daß GOtt mit uns Geduld hat, und der Heilige Geist unserer Schwachheit hilft, auf daß der Geist wider das widerspenstige Fleisch siegen möge, und nicht auf die eine oder andere Seite weiche, das ist, nicht sicher, aber auch nicht verzagt und ungeduldig werde. Das ist unser Ruhm, nämlich, GOttes gnädige Verzeihung, der uns sonst die Bosheit unseres Fleisches könnte zurechnen; aber um des Geistes willen, der dem Fleische widerstrebet, rechnet er es uns nicht zu, sondern vergibt und verzeiht es uns gnädiglich. Und also hat der Heilige Geist Lea aufgerichtet und getröstet, ihr Elend angesehen und ihr Gebet erhört; darnach, daß er ihr Mut gemacht, daß sie gesagt hat: „Nun will ich dem Herrn danken“ 21. Welches alles Worte und Reden sind des Geistes, der da widerstrebet und wider das Fleisch streitet; wiewohl er noch nicht siegt und triumphirt, so weicht er aber gleichwohl nicht, und willigt auch in die Schwachheit und Bosheit des Fleisches nicht.

181. Lyra unterscheidet die Anfechtungen, so Lea gehabt, christlich genug auf vierfache Weise, je nachdem die vier Personen sind, von denen ihr dieselben widerfahren sind; er habe es gleich von den Juden genommen, oder aber von sich selber nach seinem Verstande erfunden. Sie ist von der Rachel verachtet, von ihrem Manne verstoßen, von den Nachbarn und Gesinde für fremd gehalten worden und endlich von GOtt auch verworfen gewesen. Diese Theilung oder Unterschied gefällt mir wohl, und reimt sich fein mit der Erfahrung derer, die da betrübt und angefochten werden. Denn wo ein Grund zur Traurigkeit vorliegt, so schließt der angefochtene Mensch alsbald Eines aus dem Andern. Denn weil Lea von Rachel verachtet wird, so schließt sie notwendig, es müsse ihr der Mann auch nicht sehr wohl gewogen sein, und deßhalb sei nicht sie, sondern Rachel die rechte Hausmutter. Zum Dritten ist auch noch dazu gekommen, daß sie die Nachbarn und das Gesinde geschmäht haben, von denen sie hat hören müssen, daß sie mit List und Betrug dem Jakob wider seinen Willen wäre beigelegt worden, und werde derhalben von Rachel billig verachtet, welche dem Jakob ordentlich verlobt und ehelich vertraut war. Dieses alles wird der Teufel ihr also eingeschärft und gemehret haben und hat damit die Anfechtung sehr groß gemacht, und ihr dazu eingegeben, daß ihr GOtt selbst auch nicht gewogen wäre.

182. Aber wider diese vier Gründe zur Traurigkeit und Verzweiflung gibt GOtt der Lea vier andere Gründe des Trostes, nämlich, vier Söhne. Zum Ersten, wider die Verachtung der Rachel steht er ihr Elend an. Zum Andern, dieweil sie noch von ihrem Manne verachtet ward, so hört er ihr Gebet. Zum Dritten, auf daß sie nicht für fremd angesehen werde, und als Eine, die gar verworfen und ausgeschlossen sei, so gibt GOtt, daß sich der Mann zu ihr thut und sich zu ihr hält. Zum Vierten, triumphirt und siegt sie wider alle Anfechtung, ist nun nicht mehr so verachtet und verschmäht, sondern lobt GOtt und dankt ihm.

183. Dem Exempel sollen wir auch folgen, und lernen, daß GOtt nach der Anfechtung die Erlösung und den Trost pflegt überschwänglich zu geben, dadurch alle Traurigkeit und Schrecken überwunden wird. Es ist aber dem Fleische schwer und wird ihm sehr sauer, auf den Trost, so GOtt geben will, zu harren. Darum werden uns solche Exempel des Glaubens vorgehalten, daß wir sehen sollen, wie Lea geschrien, geglaubt, geharrt, und alle Gewalt und Plage von ihrer Schwester erlitten habe im Glauben und in Hoffnung, und wie GOtt sie endlich, da sie auf ihn gehofft, angesehen und reichlich getröstet habe.

\chapter{Das dreißigste Capitel}

\section[Erster Theil]{Erster Theil. Von dem Neide Rachels gegen ihre Schwester, und von ihrer Ungeduld, und wie sie Jakob darüber bestraft}

\begin{quote}
\textbf{V. 1. Da Rachel sahe, daß sie dem Jakob nichts gebar, neidete sie ihre Schwester, und sprach zu Jakob: Schaffe mir Kinder; wo nicht, so sterbe ich.}
\end{quote}

1. Dieses Capitel ist eines von den allerschwersten in diesem Buche, beides, so viel die Worte und auch den Verstand und Historie belanget. Derhalben fallen viel Fragen ein, beide was die Grammatik und auch was die Historie betrifft, welche historischen Fragen noch viel schwerer sind denn die Fragen in der Grammatik. Auf daß wir aber nicht dafür angesehen werden, als wollten wir es gar übergehen, so wollen wir so viel thun, als wir können.

2. Und zum Ersten erheben die fleischlichen und Epicurischen Menschen diese Frage, wie sie solche Historie ansehen als ein Exempel der äußerlichen Schande und Unreinigkeit. Denn Jakob nimmt zwei Schwestern und darüber noch ihre zwei Mägde, und wird also ein Ehemann zu vier Weibern: darum richten sie Jakob darnach, wie sie geartet sind, und erdichten, daß er unkeusch soll gewesen sein, der sich Tag und Nacht nur ganz und gar der Liebe zu den Weibern beflissen habe, sintemal er sich nicht habe an einem Weibe genügen lassen, sondern deren vier hat nehmen müssen. Dadurch werden sie sehr geärgert und nehmen daher Ursache zu lästern und dagegen ihren Cölibat und ehelosen Stand zu rühmen und zu loben wider die Geilheit und Unkeuschheit, wofür sie es halten, der sehr heiligen Patriarchen und vornehmlich Jakobs.

3. Aber ihr Urtheil von diesen Dingen gilt so viel, als wenn irgend eine Sau oder Esel von einem trefflichen Harfe- oder Lautenschläger urtheilen wollte. Es ist ein solches Volk, das ganz und gar in Unzucht, Hurerei und Ehebruch ersoffen ist, die sich Tag und Nacht nur von ihrer Unzucht träumen lassen und denken, was sie wohl thun würden, wo ihnen solche Freiheit oder Muthwille gestattet würde, daß sie alle Nächte mit den Weibern wechseln und mit ihnen schätzen möchten nach der Brunft ihres Fleisches, wie sie mit ihren Huren spielen: also machen sie keinen Unterschied zwischen dem heiligen Leben der Eheleute, und ihrer schändlichen Hurerei, die sie treiben.

4. Lyra aber läßt es sich hier sehr sauer werden, Jakob zu vertheidigen. Aber man soll den großen und unzüchtigen Säuen also antworten: Daß Jakobs Keuschheit größer sei, der vier, fünf, oder hundert Weiber hat, denn in ihrem ganzen Cölibate und ehelosen Stande gefunden wird, wenn sie sich schon ihrer Huren enthalten. Denn wir wollen den Fall setzen, daß irgend Einer in Wahrheit ohne Weib lebe, der sich gänzlich enthalte: so ist doch gewiß, daß Jakob hundertmal keuscher sei. Denn der außerhalb der Ehe ohne Weib lebt, brennt Tag und Nacht von Unzucht: wenn er schläft, so hat er seine Befleckung; wenn er aufwacht, so fühlt er die Brunst wieder. Was ist aber nun das für eine Keuschheit, wo man also lebt und brennt mitten im Feuer der Unzucht? Wo er irgend ein schön Weib ansieht, wird er ganz und gar entzündet; und wiewohl er sich bezwingt und des Werkes enthält, so bringt doch solche Brunft ihre Befleckung und Unreinigkeit mit, nicht allein wenn er schläft, sondern auch wenn er wacht, wie Hieron das bezeugt.

\quellseite{146}

5. Derhalben ist solches keine freiwillige, sondern eine erzwungene Keuschheit. Denn er wollte wohl lieber mit Gewalt über die Weiber fallen, als wäre er unsinnig: er kann die Brunst nicht bezwingen oder inne halten. Ich sage von denen, die da keusch sind und der Unzucht widerstreben, wie St. Bernhard und Bonaventura gewesen sind; welche, da sie sonst heilige und fromme Männer gewesen, doch nicht haben können ohne Brunft sein. Ja, es klagt Hieronymus auch über die Träume, darin ihm geträumet hat, wie er zu Rom mit den Mädchen getanzt hätte. Desgleichen zeigen Ambrosius und Augustinus heimlich an, wiewohl sie beide sich der Weiber enthalten haben, daß es ihnen gleichwohl nicht an der Brunst und Befleckung, so sie des Nachts gehabt, gemangelt habe. Wo nun solches den besten und heiligsten Männern widerfahren ist, was geht uns denn der erheuchelte Cölibat und ehelose Leben dazu auch das Urtheil dieser unflätigen Säue an, die da meinen, daß sich die Väter gegen ihre Weiber dermaßen gehalten haben, wie sie sich halten würden, wo ihnen gestattet würde, daß sie bei so vielen Weibern wohnen möchten?

6. Was aber hernach folgt, und auch was droben ist gesagt worden von der Keuschheit Jakobs, damit wird angezeigt, daß er ganz keusch gelebt habe. Denn ist das nicht eine Keuschheit, wo Einer ohne Weib lebt bis in das achtzigste Jahr seines Alters, und daß er sich enthält nicht aus Zwang, sondern freiwillig, und daß er dazu im Geist geneigt ist? Dasselbe halte du auch also von deiner ersten Jugend an bis in das Alter, so wollen wir deinen Cölibat und deine Keuschheit rühmen. Derhalben hat der Cölibat der Papisten gar keinen Ruhm, wenn man denselben mit der Keuschheit der Patriarchen vergleichen wollte.

7. Hernach aber werden wir hören, wie sich Jakob beklagt des mühseligen und verdrießlichen Lebens halben, da er mit Laban etwas hart redet und ihn beschuldiget, und spricht: „Des Tages verschmachtete ich vor Hitze, und des Nachts vor Frost, und durfte nicht schlafen“; da er sein ganzes Leben beschreibt, das er nicht im Müßiggange, ohne Sorge, Beschwerung, Mühe und Arbeit zugebracht, und daheim nicht müßig gesessen bei dem Feuer oder Nachsessen und etwa mit Lea gescherzt habe, sondern wie er Tag und Nacht draußen gewesen, des Viehes gehütet und alle Arbeit bei dem Viehe verrichtet, Frost und Hitze unter dem Himmel hat müssen leiden, und ist gar selten in die Hütte unter das Dach gekommen. Das soll einem den Ekel wohl vertreiben, wenn er schon jung ist.

8. Derhalben, sage ich, hat Jakob mit so vielen Weibern ganz keusch gelebt, und viel keuscher, denn sonst irgend ein Anderer mit Einem Weibe. Denn wo man des Tages Hitze und des Nachts soll Frost leiden, soll dazu derselben sein mit vieler Arbeit in fremden und schweren Diensten unter einem solchen Herrn, der ein Schalk und ein Geizhals ist, das ist wahrlich eine große und unerträgliche Last.

9. Darum entschuldiget Lyra Jakob recht, daß er keusch gelebt und seiner Weiber und Mägd[e] mäßig gebraucht habe. Denn er hat sich nur darum zu ihnen gehalten, daß er begehrete Kinder mit ihnen zu zeugen; welches auch wohl zu sehen ist an dem großen und kümmerlichen Verlangen der Weiber, daß sie immer gewünschet, daß sie Kinder haben möchten. Denn siehe, wie sie sich so kümmerlich nach dem Manne gesehnet, daß sie nur möchten Kinder von ihm haben. Darum ist er selten in den Hütten gewesen, da die Weiber ihre Wohnung gehabt, und hat bei seinen Weibern gewohnt in der keuschsten Liebe, daß er begehrte, Kinder mit ihnen zu zeugen. Und wir sehen, daß er Lea lieb gewonnen um der Kinder willen, die zuvor gehaßt und verachtet war, und mit ihr mehr Gemeinschaft hat als mit der Rahel, nachdem er gesehen und erfahren hat, daß sie fruchtbar ist. Derhalben wollen wir auf die erste Frage bei diesem Texte also geantwortet haben, daß den unflätigen Menschen das Maul verstopft werde, die von dem Ehestande der Patriarchen nach ihrer Unzucht und unreinem Cölibat urteilen wollen.

10. Von der andern Frage ist vorhin gehandelt worden, nämlich, daß Jakob über die zwei Weiber auch noch zwei Mägde genommen hat. Ich kann aber davon nichts Gewisses sagen, sondern es läßt sich ansehen, daß zu der Zeit in selbigem Lande der Gebrauch gewesen sei, daß die unfruchtbaren Weiber ihren Männern die Mägde gegeben und beigelegt haben; wie Laban jeglicher Tochter eine Magd mitgegeben hat, vielleicht zu dem Ende, wo die Tochter nicht gebären würde, daß alsdann die Magd an ihre Statt treten möchte, daß von ihr das Haus erbaut und gemehret würde. Also hat droben, 1 Mos. 16, 2., Sarah dem Abraham eine Magd gegeben, nicht die da fremd war, sondern aus ihrem Hause und Gesinde. Isaak ist dieser Weise nicht nachgefolgt, und Jakob würde dieselbe auch nicht gebraucht haben, wo er nicht von Laban wäre betrogen worden und die Weiber dies nicht so sehr begehrt hätten.

11. Denn das muß man auch betrachten, daß man zu der Zeit die Fruchtbarkeit für einen großen Segen und sonderliche Gabe Gottes gehalten. Wie das 5 Mos. 28, 4. zu sehen ist, da Mose unter die Segen die Fruchtbarkeit des Leibes auch mit zählt; und 2 Mos. 23, 26. spricht er: „Es soll nichts unträchtig noch unfruchtbares sein in deinen Lande“ c.

12. Wir achten es heutiges Tages so groß nicht. An dem Viehe haben wir es gerne und begehren es auch sehr, daß es fruchtbar möge sein; aber unter den Menschen sind ihrer Wenige, die es für einen Segen halten, wo die Weiber fruchtbar sind; ja, man findet ihrer Viele, die den Segen nicht gern haben, und halten es für ein sonderlich Glück, wo die Weiber unfruchtbar sind. Welches wahrlich auch wider die Natur, und viel weniger göttlich und heilig ist. Denn solche Neigung ist von Gott in die menschliche Natur gepflanzet, daß sie wünscht und begehrt, daß sie möge fruchtbar und gemehret werden.

13. Derhalben ist es unmenschlich und gottlos, wo man die Kinder verachtet. Wie neulich Einer sein Weib, da sie oftmals geboren, eine Säu genant hat. Der nichtswürdige und unreine Mensch! Also sind die heiligen Väter mit nichten gesinnet gewesen; denn sie haben es erkannt, daß es ein sonderlicher Segen Gottes wäre, wenn ein Weib fruchtbar war, und dagegen haben sie auch die Unfruchtbarkeit der Weiber für einen Fluch gehalten. Und solche Erkenntniß und Urtheil ist aus Gottes Wort hergeflossen, 1 Mos. 1, 28.: „Seid fruchtbar und mehret euch“ c. Daher haben sie verstanden, daß Kinder eine Gabe Gottes wären.

14. Das Exempel aber oder die That, da die Mägde den Männern sind gegeben worden anstatt der Weiber, welche unfruchtbar waren, soll nicht dahin gezogen werden, daß man demselben nachfolgen sollte. Es soll niemand sagen: Dieses hat Jakob gethan, darum geziemet mir es auch zu thun; wie droben von Münzer ist gesagt worden, der die Bauern ermahnete, sie sollten nach dem Exempel Josuas und Simsons die Fürsten todtschlagen: sondern du sollt denken, daß dir gebühren will, bei dieser Regel zu bleiben: Ein Jeder soll sein eigen Weib haben.

15. Wenn du Josua und Simson willst gleich sein, so schaffe auch, daß alle Umstände dich zur Veränderung des gemeinen Regimentes und die Obrigkeit todt zu schlagen treiben und bewegen, wie dieselbigen großen Helden durch sonderlichen Beruf bewogen worden sind; sonst gilt das Exempel nicht. Denn du hast mehr und größere Exempel, die da zeugen, daß man die Obrigkeit nicht tödten und die gemeinen Regimente nicht ändern soll. Derhalben sollst du der Obrigkeit gehorsam sein wie andere Unterthanen; denn du bist von der Regel nicht ausgenommen. Darum gedenke und bleibe im Gehorsam, werde kein großer Held: wie Münzer und Andere sich dessen unterstanden haben c.

16. Darum wird nun dies nicht gelehrt zum Exempel, sondern daß wir uns der Exempel und des Nachfolgens enthalten: wir sollen uns derselben verwundern, aber nicht nachfolgen. Denn etliche Dinge sind, denen wir nachfolgen, etliche aber, deren wir uns verwundern sollen. Hoffe, glaube und rufe Gott an, wie Lea gethan hat; aber vier Weiber sollst du nicht nehmen, wie Jakob gethan hat. Denn dasselbe gebühret dem Jakob allein, und denen, die Gott von der gemeinen Regel hat ausnehmen wollen. Wir sollen uns üben im Glauben, Geduld und Hoffnung, so uns an den Vätern vorgehalten wird, und sollen uns der hohen Exempel, wo sie als große Helden sonderliche Thaten gethan haben, enthalten.

\quellseite{147}

\subsection{I}

17. Dieses sind die Fragen, die man im Anfang dieses Capitels hat erwägen und davon etwas sagen müssen. Nun wollen wir auch den Text besehen, welcher auch etwas dunkel ist. Im Texte steht: „Sie neidete ihre Schwester“, oder eiferte wider sie. Denn im Hebräischen steht das Wort kana, neiden, eifern; kanno, Eiferer etc.

18. Es wird aber auch gefragt: Ob die heiligen Weiber, die Gott im rechten Glauben angerufen haben, auch neidisch und gehässig gewesen seien? wie das der Text klar anzeigt. Lyra macht zweierlei Neid: der eine ist, damit ich einen Andern neide, wenn es ihm wohl geht; der andere, wenn ich Unmuths werde und zürne um meines Kummers und Elendes willen. Das ist aber eigentlich kein Neid. Wir entschuldigen es also, wie wir sonst pflegen zu thun: Die heiligen Väter und Mütter sind noch unter dem Fleische gefangen gewesen, wie wir auch das Fleisch haben, so wider uns streitet und uns gefangen nimmt in der Sünden Gesetz.

19. Es geht aber mit allen Anfechtungen, vornehmlich was hohe Anfechtungen sind, also zu, daß sie dem Menschen schier alles Empfinden der Göttlichkeit wegnehmen. Also, wer mit Kreuz und Unglück beladen ist, der fühlt keine Geduld: wer des Glaubens und Hoffnung halben angefochten wird, der empfindet keinen Glauben und keine Hoffnung, er läßt sich nicht dünken, daß er glaube oder hoffe, sondern es dünket ihm, als sei er ganz gottlos. Wer geplagt wird mit Verzweiflung und Ungeduld wider Gott, wenn Gott irgend etwas auf eine andere Weise schickt und regieret, denn wir es vorgeschrieben und bei uns bedacht hatten, der fühlt in seinem Herzen sehr großen Zorn und Widerwillen wider Gott, und läßt sich dünken, daß er nicht ein Fünklein Glaubens oder Hoffnung habe.

20. Solche Gedanken pflegen betrübte Menschen zu haben, wenn die Anfechtung mit Gewalt daher geht. Sonst, wenn keine Anfechtung vorhanden ist und ich alsdann vom Glauben und von der Hoffnung bisweilen disputire, so dünkt mich nichts Leichteres zu sein, denn der Glaube und die Hoffnung; aber im Kampf verschwinden bald Alles miteinander. Also geht es mit andern Anfechtungen auch, als mit Ungeduld, Zorn und Unkeuschheit: da sieht und fühlt man die Brunst und Nachgier. Und es ist wahr, solche Affecten und Begierden sind nicht der Glaube, Hoffnung und Liebe, sondern sind das Fleisch selbst. Also hat Rahel die Anfechtung gehabt mit dem Neide, damit sie gewünscht hat, daß es geschehen möchte, daß doch ihre Schwester keine Kinder gebäre, wiewohl sie gleichwohl unter der Anfechtung nicht ist liegen geblieben; sondern es ist ihr widerfahren, daß St. Paulus Röm. 7, 19. sagt: „Das ich nicht will, das thue ich, und das ich will, das thue ich nicht.“

21. Wer mit Verzweiflung angefochten wird, der fühlt das unaussprechliche Seufzen, damit er wünschet, daß er mit solchem Zweifel nicht möchte also geplagt werden; und fühlt doch gleichwohl, daß der Zweifel noch etwas stärker ist und ihn gefangen nimmt in der Sünde und des Unglaubens Gesetz. Aber dawider wird durch den Heiligen Geist im Herzen ein heimlich Seufzen erweckt, das wiederum murrt und widerspricht, dazu auch zürnt mit dem Unglauben und sagt: Pfui dich, du leidiger Unglaube, sprich nicht also: Wer weiß? es ist erlogen etc. Und dasselbe geschieht mit einem kleinen Aechzen, und mit unaussprechlichem Seufzen, wie es St. Paulus Röm. 8, 26. nennt; welches Seufzen das Herz, so angefochten wird, erhält, daß es nicht gar gottlos werde und Gott gar verachte, wie die Epicurer pflegen zu thun, die in jeder geringen Anfechtung den Glauben und Hoffnung gar dahin fahren lassen.

22. Es ist aber dieser Streit dem Fleische gar schwer und verdrießlich, welches lieber einen solchen Glauben haben wollte, der nicht angefochten würde, oder der Stöße nicht also erwarten dürfte; sondern der Löwen und Bären erwürgen könnte, wie David, Simson und Andere gethan haben. Mein Glaube aber ist so gering und schwach, daß er sehr leichtlich von einem Löwen oder Bären verzehret und hingerissen werden möchte. Es ist aber noch das zerstoßene Rohr und das kleine glimmende Docht vorhanden, Jes. 42, 3., und Gott ist auch gegenwärtig, der das Seufzen versteht, der da erkennt, was der Heilige Geist für uns bittet. Darum ist es vor Gott ein groß Feuer, ja, es ist heißer und größer, denn Himmel und Erde sein kann.

23. Darum soll man nun die Heiligen also entschuldigen nach dem Spruch Pauli Röm. 7, 23.: „Ich sehe ein ander Gesetz in meinen Gliedern, das da widerstreitet dem Gesetz in meinem Gemüthe, und nimmt mich gefangen.“ Item, da er v. 25. sagt: „Mit dem Fleisch diene ich dem Gesetz der Sünde.“ Denn also sind die heiligen Mütter Menschen gewesen, die im Fleisch gelebt, aber nicht nach dem Fleisch gewandelt haben, uns zum Exempel, daß wir nicht verzagen und auch dem Fleische nicht weichen sollen. Rahel neidet zwar ihre Schwester und hat einen Eifer wider sie, aber der Geist streitet und kämpft wider den Eifer des Fleisches. Sie wollte gern den Segen gehabt haben, damit die Schwester begnadet ist, und zwingt sich doch selbst, daß sie nicht wünscht, daß sie sollte unfruchtbar werden.

24. Darum aber mag man sehen, daß die sehr heiligen Weiber nicht unfeusch gewesen sind, sondern haben Kinder und den Segen begehrt. Denn das ist die Ursache des Neides an Rahel, welche, wo sie andern Weibern, wie solcher zu unserer Zeit sehr viele anzutreffen, gleich gewesen, würde gesagt haben: Was liegt mir daran? ich gebäre oder gebäre nicht; wenn ich nur Hausmutter bleibe und sonst alles andern Dinges genug habe, daran genüget mir wohl. Aber Rahel hat so großes Verlangen nach Kindern, daß sie auch lieber sterben wollte, denn daß sie sollte unfruchtbar bleiben. Und ich weiß nicht, ob ich in irgend einer Historie auch dergleichen Worte mehr gelesen habe.

\subsection{II}

25. Derhalben ist es ein Exempel eines sehr frommen und keuschen Weibes, welches nur den einigen Eifer und Verlangen hat, auch bis in den Tod, daß sie möchte Kinder haben; wie droben 1 Mos. 16, 2. Sarai sich auch also hat vernehmen lassen, daß sie solch groß Verlangen hätte nach Kindern. Und ist an allen beiden dieser Affect und Begierde sehr zu loben. Wenn ich keine Kinder haben soll, sagt Rahel, so muß ich sterben; ich will lieber nicht leben, denn keine Kinder haben; und wo ja mein Fleisch zur Geburt untauglich und unfruchtbar ist, so soll doch mein Mann mit der Magd Kinder zeugen; wird es ihm aber mit derselben auch versagt werden, so begehre ich, daß mich Gott aus diesem Leben fodern wolle.

26. Und sie hat nicht geringe Ursachen gehabt dieses großen Verlangens. Denn Jakob hat ohne Zweifel ihnen allen beiden gerühmt und verkündigt, wie er die Verheißung habe, daß der gebenedeyete Same sollte von ihm geboren werden. Da sie nun diese Predigt gehört, ist dadurch in ihnen das Verlangen und Begierde entzündet worden, daß sie möchten Kinder bekommen. Und vornehmlich ist das Verlangen an Rahel groß gewesen, so daß sie weniger Hoffnung sie gehabt, desto größer ihr Verlangen gewesen ist, auch also, daß sie sich vornimmt, entweder zu gebären oder zu sterben; wie sie denn auch hernach über der Geburt dahin stirbt. Dies Verlangen und Begierde des frommen gottseeligen Weibes ist gut und heilig; ist aber so gar heftig, daß sie darüber in einen harten Kampf geräth, und mit Neid und Eifer wider ihre Schwester angefochten wird. Und man soll dies für kein schändlich Exempel halten, wie die Papisten diese heiligen Menschen beschuldigen, daß sie sollen unkeusch gewesen sein. Denn sie haben im Ehestande nicht auf die schändliche und elende Lust des Fleisches gesehen, sondern auf den Segen, daß sie Kinder möchten haben um des verheißenen Samens willen.

\begin{quote}
\textbf{V. 2. Jakob aber ward sehr zornig auf Rahel, und sprach: Bin ich doch nicht Gott, der dir deines Leibes Frucht nicht geben will.}
\end{quote}

27. Jakob wird zornig, darum daß Rahel der Kinder halben so gar hart auf ihn dringt, und antwortet ihr darauf also: Die Schuld ist nicht mein, du bist doch mein allerliebstes Weib, bei der ich ja mehr gewohnt habe, als bei deiner Schwester; bin ich doch nicht Gott, der ich dir die Fruchtbarkeit, Kinder zu gebären, geben könnte.

28. Dies ist aber alles sehr gering und unflätig, und wie hoch man es auch loben kann, so ist es doch menschlich, fleischlich und weibisch; und ist ein Wunder, daß sich der Heilige Geist damit bekümmern mag, solch Ding zu beschreiben, nämlich, wie Rahel ihre Schwester Lea geneidet habe, und wie die Eine geboren und aber die Andere unfruchtbar gewesen sei; item, daß Rahel schier gelästert und allzu ungestüm auf den Mann gedrungen habe, daß sie möchte Kinder von ihm haben, also daß Jakob auch nach dem Fleisch sehr zornig auf sie wird und also zwischen denselben Eheleuten ein rechter Zank entsteht. Ist das nicht kindisch Ding?

\quellseite{148}

29. Von dieser Frage aber haben wir droben gehandelt, und man soll sie allezeit bewegen und den Leuten einschärfen, nämlich, warum doch der Heilige Geist, der ja einen sehr reinen Mund hat, mit so großem Fleiß von diesen Dingen reden möge, deren der heiligste Vater, der Pabst, mit seinen keuschen Mönchen und Nonnen ungern einmal gedenken wollten als solcher Dinge, die gar unflätig und fleischlich sind vor ihren Augen. Denn sie wandeln in großen Dingen ihres Cölibats und edelsten Lebens. Diesen Unflat aber, wie die Weiber schwanger geworden, Kinder geboren und wie die Eheleute etwa mit einander gezürnt, achten sie nicht werth, daß man es lesen sollte. Es hätte, sagen sie, der Heilige Geist seiner Heiligkeit nach von himmlischen und sonst von andern höheren Dingen reden können und nicht von solchen geringen und fleischlichen Dingen; er sollte ein Mönch oder Nonne geworden sein; aber nun erzählt er nur, wie es um die Haushaltung gestanden habe und wie es mit Jakobs Ehestande zugegangen sei. Das ärgert uns heilige und englische Männer, die wir über die Wolken herwandeln in Weisheit und Geistlichkeit der Engel.

30. Dieweil aber sie diese geringen Dinge verachten und einen Ekel daran haben, darum ist der Heilige Geist nun den so hochsinnigen und ruhmredigen Heiligen feind und erkennt sie nicht für die Seinen, er läßt sie in ihrer Herrlichkeit, Hoffahrt und Eitelkeit immer hinfahren und fährt hernieder zu seinen Creaturen; dafür sorgt er und ziert sie. Denn er hat den Erdboden geschaffen, er hat Mann und Weib geschaffen, und sie gesegnet, daß sie sollen fruchtbar sein; er hat ihnen die Welt unterworfen, und er ist es, der noch Alles erhält; ernährt und gibt der Mutter Milch, das Kindlein damit zu nähren und zu erhalten. Diese Creatur verachtet er nicht, er verachtet sein Werk nicht; aber doch gleichwohl dem Cölibat und ehelichen Stande nicht zur Schmach. Denn derselbe ist auch eine Gabe Gottes, und wir loben beiderlei Stände, jeden nach seiner Ordnung und Maße.

31. Wir loben den Ehestand nicht also, daß wir den Cölibat schelten oder verwerfen sollten, welches die Papisten nicht thun; denn dieselben loben ihren Cölibat also, daß sie dagegen den Ehestand lästern und schmähen, und die Keuschheit rühmen sie nicht als eine Gabe Gottes, sondern als ihr eigen Werk. Das ist eine rechte Teufelslehre; wie das auch an den Früchten wohl zu sehen ist, die daraus gekommen sind. Denn weil sie sich in der Werke des Fleisches haben rühmen wollen, darum leben sie heutiges Tages so schändlich, und sind unzüchtiger, denn alle Hurenswirth und Hurensbälge sein mögen. Warum aber das? Antwort: Darum, daß sie die Keuschheit nicht gelobt haben als ein Werk Gottes, sondern als ihr eigen Werk, und haben den Ehestand, der ein Werk und Ordnung Gottes ist, verachtet und gelästert.

32. Derhalben hat nun Gott Lust, solche geringe Dinge zu beschreiben, daß er damit anzeige und bezeuge, daß er nicht verschmähe, auch seinen Abscheu habe, oder auch nicht weit sein wolle von der Haushaltung, von einem frommen Ehemanne, und von Weib und Kindern. Warum thut er aber das? Antwort: Darum, daß er es geschaffen hat, derhalben regieret und erhält er es auch als seine Creatur; wiewohl das Fleisch der Sünden halben verderbet ist, und das Fleisch derer, die im Cölibat leben, ist auch nicht rein und dieselben sind auch eben von solchem Ursprung hergekommen. Warum verdammen sie ihn denn und halten von ihrer unreinen Keuschheit mehr als vom heiligen Ehestande? Darum will uns nun der Heilige Geist lehren und bezeugen, da er von diesen geringen, menschlichen und gemeinen Dingen handelt, daß wir wissen sollen, daß er wolle bei uns sein, für uns sorgen, und beweisen, daß er unser Schöpfer und Regierer sei. Das sehen die Papisten nicht, sondern verachten es; darum müssen sie auch billig die Strafe tragen für solche Verachtung.

33. Zum Andern sehen sie auch nicht, wie mit so reinem und mütterlichem Herzen diese Weiber um Kinder bitten; welche Begierde und Worte der Rahel anzeigen, daß ihr Herz weit und fern von der Unzucht gewesen ist. Das sehen sie aber nicht. Warum das? Antwort: Dieweil geschrieben steht: Hinweg mit dem Gottlosen, daß er nicht sehe oder höre, was Gott redet, was er thut oder wirkt. Sie sind nicht werth, daß sie die Herrlichkeit Gottes sehen. Also hören sie uns heutiges Tages, wie wir rufen, das Wort GOttes lehren und singen; sie sehen die Werke Gottes und sehen sie doch nicht. Warum? weil sie es nicht werth sind. Also sehen sie auch nicht, daß uns Rahel vorgehalten wird zum Exempel eines sehr feinen, reinen und mütterlichen Herzens und Keuschheit, welche nur vom Fleisch Kinder begehrt wie Lea auch. Sie hören, daß diese Weiber von nichts Anderem reden, als: Der Herr hat es gegeben, der Herr hat es gesegnet; und hernach, da sie sich zanken: „Der Herr wird Richter sein zwischen mir und dir.“ Alle ihre Worte und Reden sind von Gott, von seinen Wohlthaten, Gaben und Werken, und halten doch gleichwohl die schändlichen Papisten dieses alles für Sünde, und achten diese heiligen Matronen dessen nicht werth, daß ihre Exempel in der Kirche und Gemeinde Gottes den Leuten sollten vorgehalten werden.

34. Was kann doch in der Gemeinde Gottes Besseres und Nützlicheres gelehrt werden, denn das Exempel einer gottseligen Hausmutter, die da betet, seufzt, schreit, Gott dankt, das Haus regieret, thut, was das Amt eines frommen Weibes mit sich bringt, begehrt, daß sie möge Kinder haben, mit großer Keuschheit, Dankbarkeit und Gottseligkeit? Was sollte sie mehr thun? Aber der Pabst, Cardinäle und Bischöfe sollen das nicht sehen; denn sie sind es nicht werth. Der Heilige Geist läßt sie immer in wunderlichen, großen und überhimmlischen Dingen wandeln, und daß sie sich nur ihrer Keuschheit verwundern und dieselbe hoch rühmen, die wohl werth wäre, daß man sie damit irgend in das gemeine Hauß wiese; diese Dinge aber sollen sie keinesweges sehen. Unterdeß aber führt und regieret der Heilige Geist die heiligen Weiber also, daß er damit bezeuge, daß dieselben seine Creaturen sind, die er regieren will nicht allein nach dem Geist, sondern auch nach dem Fleisch; daß sie Gott anrufen, beten, ihm für die Kinder danken und ihrem Mann gehorsam sein sollen etc. Das heißt, diese unflätigen Dinge der Gemeinde zu lesen geben, auf daß Gott seine großen und wunderbaren Werke in der Kirche und Gemeinde beweise, wie er in den fleischlichen und irdischen Dingen himmlische und geistliche Dinge wirken könne. Die Gottlosen aber sehen die fleischlichen Dinge wohl, und was geistlich und göttlich ist, sehen sie nicht.

35. Zum Dritten sehen sie das auch nicht, mit welcherlei Unfall, Jammer und Elend diese armen Weiber gekämpft haben, sondern lassen sich träumen, sie haben in Ueppigkeit, Müßiggang und Wohlust gelebt, wie sie ein Leben führen. Sie denken, Jakob sei ein großer Herr gewesen und habe in großem Gut mit seinen Weibern in eitel Wohlust gelebt; Rahel sei wie eine Königin mit köstlichem Schmuck geziert und mit großem Geld und Gut begabt gewesen. Denn ein solch Leben und solche Weiber begehren sie zu haben und ein großer Theil unter ihnen lebt auch in solcher Wohlust und gutem Gemach. Aber siehe die Historie an, so wirst du finden, daß Jakob ein armer elender Mann gewesen, der auch nicht einen eigenen Keller gehabt hat. Also sind die Weiber auch sehr arm gewesen; wie sich Rahel und Lea hernach beide beklagen 1 Mos. 31, 14. 15.: „Wir haben doch kein Theil noch Erbe mehr in unseres Vaters Hause. Hat er uns doch gehalten als die Fremden; denn er hat uns verkauft, und unsern Lohn verzehret.“ Sind nun die Töchter arm und Bettlerinnen gewesen, so ist es viel mehr glaublich, daß die Mägde werden arm und nothdürftig gewesen sein.

36. Derhalben ist große Armuth da gewesen, und immer viel harte schwere Arbeit dabei, die sie Tag und Nacht getragen und über das noch Hunger, Durst und Frost gelitten haben. Dem folge du auch also nach, der du diese sehr heiligen Menschen beschuldigst, daß sie sollen unkeusch gewesen sein. Nimm vier Weiber, und diene darnach ganze vierzehn Jahre einem solchen Herrn, der da geizig und undankbar ist. Aber niemand ist, der das sehen oder betrachten wolle. Und offt heutiges Tages keiner, der auch nicht mehr denn ein Weib, geschweige denn vier, nehmen wollte mit solcher Bedingung, daß er nur dabei dienen und immer ein Knecht sein sollte und mit seinem Schweiß und Arbeit einen Andern reich machen, und darnach nicht sehen, wie er sich und sein Weib nähren möchte.

37. Darum ist dies ein überaus großer Glaube, daß Jakob so viel Weiber nimmt auf die Hoffnung, daß sich Gott in seiner Armuth helfen und ihn sammt seinem ganzen Hausgesinde ernähren werde. Ein Weib nehmen ist zwar leicht zu thun; aber dasjenige darnach mit den Kindern und Gesinde ernähren, das hat Mühe und Arbeit. Da ist aber nun niemand, der auf diesen Glauben Jakobs wollte Acht geben; ja, Viele sind, die an ihren Weibern das hassen, daß sie fruchtbar sind und viel Kinder gebären, nur um der Ursache willen, daß man die Kinder nähren muß und auferziehen. Denn also pflegen sie gemeiniglich zu reden: Warum sollte ich ein Weib nehmen, dieweil ich arm und ein Bettler bin; ich will lieber allein meine Armuth tragen und mein Elend und Armuth nicht größer machen. Aber diese Schuld legt man nicht billig auf den Ehestand und die Frucht des Leibes. Warum klagst du nicht vielmehr über deinen Unglauben, daß du der Gütigkeit Gottes nicht vertraust, und führst also großer Elend über dich selbst damit, daß du den Segen Gottes strafst? Denn wo du der Gnade Gottes und seinen Verheißungen vertrauen könntest, so würdest du ohne Zweifel wohl ernähret werden; dieweil du aber auf den HErrn deine Hoffnung und Vertrauen nicht setzest, so wirst du nimmermehr Glück haben können.

\quellseite{149}

38. Auf die Weise sehen wir nun, daß in diesen unflätigen Dingen die größten Tugenden hervor leuchten, als nämlich ein vortrefflicher Glaube, gewisse Hoffnung und eine unüberwindliche Geduld gegen Gott und Menschen; welches weder die Papisten noch Andere leichtlich sehen werden, es sei denn, daß sie die Umstände fleißig ansehen werden. Die aber, so diese Dinge gar nicht achten, auch gar nichts glauben, die mögen hinfahren und die unreinsten Hurenjäger bleiben.

39. Wir sollen lernen, daß nicht allein zur Kirche oder Religion und zum weltlichen Regiment, sondern auch zur Haushaltung ein gläubiger Mann und gläubiges Weib gehört. Denn ein ungläubiger Mensch thut nimmermehr etwas recht und daß Glück dabei sein könne. Ich gedenke noch wohl, daß Staupitz mir eine Historie erzählte von einem Prior in einem Kloster. Da derselbe stets klagte, daß die Zinsen und Einkünfte des Klosters zu gering wären, die Mönche damit nach Nothdurft zu unterhalten, hat endlich Staupitz die Register seiner Einnahme und Ausgabe von ihm gefordert, darin er gesehen hat, daß des Klosters Gut jährlich merklich gemehret gewesen; darum hat er ihn vor sich fordern lassen und ihn des Amtes entsetzt, und also zu ihm gesagt: Du bist kein gläubiger Mensch, darum ist es unmöglich, daß du dem Kloster wohl vorstehen könnest.

40. Darum steht alles Regiment und das ganze Leben im Glauben; denselben aber zu erwecken und zu stärken, dazu dient überaus wohl das Exempel Jakobs. Denn derselbe hat auch nicht einen Heller, der sein eigen sei, und muß doch gleichwohl vier Weiber ernähren in dem schweren Dienst, davon er nichts mehr hat, denn kümmerlich seine Nahrung und Kleidung, daran er sich vierzehn Jahre hat genügen lassen. Das wird ihm nicht leichtlich Einer unter uns nachthun, auch nicht mit Einem Weibe, geschweige mit zweien oder mit vieren. Denn niemand ist, der auch nur Eine Woche ohne gewisse Besoldung leben und darben wollte; wie sollte Einer denn ganze vierzehn Jahre nur um das tägliche Brod dienen? So gar sind wir mit diesen Menschen nicht zu vergleichen, sondern wir sind gar ungläubig.

41. Wenn nun der Heilige Geist von schlechten, geringen und verachteten Dingen redet, so mengt er gleichwohl mit ein die köstlichsten Edelsteine sehr großer schöner Tugenden, welche die groben unflätigen Säue, die Papisten, nicht sehen. Denn sie sehen nur auf die fleischlichen Dinge, und anstatt des Honigs saugen sie aus der Rose eitel Gift. Denn in diesen Exempeln der Väter strafen sie, daß sie sollen Unzucht getrieben haben, daß sie viel Weiber gehabt; und wenn Rahel und Lea Kinder begehren, strafen sie solches als geile Beischläfe. Sie denken daran nicht, wie sie Gott angerufen und gepreiset, wie sie so großen Glauben und Geduld gehabt haben.

42. Wie sollten aber solche Leute anders urtheilen, die im Müßiggange und Wollust dieses Lebens gar ersoffen sind, wie die Cardinalen und Dompfaffen, die ohne Mühe und Arbeit, dazu auch ohne alle Anfechtung ihr Leben zubringen? Warum sollen wir die heilige Schrift auf eine andere und bessere Weise lesen. Und eben darum wiederhole ich dieses so oft und schärfe es euch ein, daß wir diese geringen und gar fleischlichen Dinge nicht mit fleischlichen, sondern mit geistlichen Augen ansehen sollen; alsdann werden wir den wunderbaren Rath des Heiligen Geistes in solchen Beschreibungen sehen, viel mehr denn sonst in hohen und sehr geistlichen Dingen.

\section[Zweiter Theil]{Zweiter Theil. Wie Rahel ihre Magd dem Jakob gibt, welche ihm zwei Söhne gebiert; item, wie Lea ihre Magd dem Jakob beilegt, welche ihm auch zwei Söhne gebieret}

\subsection{I}

\begin{quote}
\textbf{V. 3. Sie aber sprach: Siehe, da ist meine Magd Bilha; lege dich zu ihr, daß sie auf meinem Schooß gebäre, und ich doch durch sie erbauet werde.}
\end{quote}

43. „Auf den Schooße gebären“, ist eine sonderliche Art und Weise zu reden in der hebräischen Sprache, damit Rahel zu verstehen gibt, daß sie wolle Mutter werden, und mit diesen Worten nimmt sie das Kind, so von der Magd geboren ist, als ihr Kind an. Denn diese Worte bedeuten das Amt einer Mutter, insondermal die Mutter das Kindlein in ihrem Schooße trägt; wie im Propheten Jesaia am 66. Capitel V. 12. steht: „Auf den Knieen wird man euch freundlich halten“ etc. Die Mutter hat das Kind auf dem Schooße. Darum hat sie gedacht: Vielleicht ich ja von mir keinen Sohn haben kann, so will ich doch Mutter werden des gebenedeyten Samens durch meine Magd.

44. Daraus kann man sehen, wie groß und heißes Verlangen sie gehabt haben nach dem verheißenen Samen; item, wie gottselig sie gewesen, und wie so einen trefflichen Glauben sie gehabt, damit sie sich an die Verheißungen von Christo, so Jakob geschehen waren, gehalten haben. Darum müssen ja gewißlich diese Weiber sehr fromm gewesen sein, sintemal sie eines solchen armen Mannes Weiber haben sein können und in ihres Vaters Hause dienen wie Mägde. Dieses ist eine solche Heiligkeit gewesen, dergleichen man sonst nicht finden kann. Es ist ein großer wunderbarer Glaube, Geduld, Hoffnung und Liebe gewesen. Unsere Bürgerstöchter oder Bauerntöchter sind so zart, daß sie solches nimmermehr thun würden.

\begin{quote}
\textbf{V. 4. Und sie gab ihm also Bilha, ihre Magd, zum Weibe; und Jakob legte sich zu ihr.}
\end{quote}

45. Das würde auch kein Weib thun; denn sie würde lieber ohne Kinder sein, als daß sie sollte die Ehre der Mutterschaft und des Ehebettes der Magd einräumen und zulassen. Und dieses ist ein wunderlicher Affect und Begierde, davon ein Exempel droben, 1 Mos. 16, 2., an Sarah auch beschrieben worden ist, daß sie die eheliche Ehre und Herrlichkeit den Mägden übergeben und eingeräumt haben. Das wird kein Weib nach dem Fleisch thun, sondern es ist ein Werk des Geistes, dadurch sie erweckt und bewogen worden ist, daß sie gedacht hat: Ich will auch eine Mutter sein, wenn ich auch schon alle Ehre meines Ehebettes und der Mutterschaft eine Zeitlang übergeben sollte.

\begin{quote}
\textbf{V. 5. 6. Also ward Bilha schwanger, und gebar Jakob einen Sohn. Da sprach Rahel: Gott hat meine Sache gerichtet, und meine Stimme erhöret, und mir einen Sohn gegeben. Darum hieß sie ihn Dan.}
\end{quote}

46. Dieses sind nicht Worte einer unzüchtigen Hure. Denn sie dankt Gott, daß sie erhöret worden, und zieht Gott zum Richter an und schreibt ihm die Gabe der Fruchtbarkeit zu. Ich habe, sagt sie, geschrieen und gebetet, aber Gott sei gelobt, der meine Sache gerichtet und meine Stimme erhöret hat.

47. Es ist eine sehr schöne Danksagung; sie hängt aber noch am Fleisch, darum sagt sie: „Gott hat meine Sache gerichtet“, als wollte sie sagen: Wiewohl meine Schwester das rechte Weib nicht ist, hat sie doch wider mich stolzirt, wie droben, 1 Mos. 16, 4., Hagar wider Sarah auch gethan hat, darum daß sie fruchtbar gewesen ist denn ich: der HErr hat aber meine Sache auch gerichtet. Aber, meine liebe Rahel, warum freust du dich so sehr? Ist es doch nicht dein Sohn, er ist ja von deinem Fleisch nicht geboren, dein Leib und deine Brüste wissen von diesem Sohne nichts. Aber was liegt daran? Sie ist doch gleichwohl fröhlich bei sich, daß sie einen Sohn von der Magd hat, den sie für ihren Sohn annehmen und halten will. So große Liebe hat sie zu der Frucht des Leibes gehabt.

48. Dan heißt so viel als Gericht. Denn „Gott“, sagt sie, „hat meine Sache gerichtet“ und meine Schmach angesehen; wie sie hernach (V. 23.) sagt: „Gott hat meine Schmach von mir genommen.“ Denn die Unfruchtbarkeit ist für eine Schmach gehalten worden, um des Mordes willen 1 Mos. 1, 22.: „Seid fruchtbar und mehret euch.“ Aber GOtt hat gerichtet, das ist, er hat für mich gestritten und meine Sache ausgeführt; denn ich werde nun auch etwas gelten, wiewohl mir dieser Sohn von der Magd geboren ist. Denn sie hat also gebetet: Lieber HErr GOtt, dieweil du mich ja nicht willst lassen Mutter werden, wiewohl ich die Hausmutter bin und das vornehmste Weib, so gib mir doch von der Magd Kinder, wenn ich ja dessen nicht werth bin; denn warum willst du mich doch also verwerfen, verstoßen, zu Schanden machen und verschmähen? Dieses ist ihr Gebet gewesen, so sie mit vielem Seufzen und Thränen gethan hat; und sie zeigt an, daß dasselbe Gebet sey erhöret worden. Ob nun wohl dieses für fleischlich und unzüchtig Ding angesehen wird, so geschieht es doch alles mit hohem großem Geist, im Glauben und wunderbarer Geduld, darnach mit Lob GOttes und Danksagung dafür, daß sie GOtt von der Schmach der Unfruchtbarkeit erlöset hatte.

\quellseite{150}

\begin{quote}
\textbf{V. 7. 8. Abermal ward Bilha, Rahels Magd, schwanger, und gebar Jakob den andern Sohn. Da sprach Rahel: GOtt hat es gewandt mit mir und meiner Schwester, und ich werde es ihr zuvor thun. Und hieß ihn Naphtali.}
\end{quote}

49. Nun wird erst die Freude Rahels groß, da der andere Sohn geboren worden ist. Aber bei den Hebräern ist eine Frage, was Naphtali heiße? Darum zanken die Grammatikern und sind dermalen der Sache noch nicht eins geworden. Santes Pagminus und Münster vermuthen es also: Ich habe einen heftigen Streit gehabt mit meiner Schwester. Daß es aber ein Kampf GOttes genannt wird, deß zeigen sie diese Ursache an, daß in der hebräischen Sprache alles, was groß und herrlich ist, göttlich genannt wird. Daher Minne eine Stadt GOttes genannt wird, das ist, eine große, herrliche Stadt. Item, Zeder GOttes, das ist, ein hoher und großer Edelbaum. Wiewohl ich mir nun nicht zuschreibe, daß ich der hebräischen Sprache großen Verstand habe, so bin ich doch dem feind, wo ich sehe, daß jemand so tollkühn ist, daß er verfälscht, was einfältig geredet ist. Gleich als ob Minne nicht auch eine Stadt GOttes sei in der Vorhaut, und daß die, so beschnitten sind, allein GOttes Volk wären. Dieses kommt alles her von der großen Hoffart und Ruhmredigkeit der Juden, die nicht leiden können, daß etwas sollte groß oder herrlich genannt werden, denn allein was sie angeht und betrifft. Wer hat aber den Grammatikern die Gewalt gegeben, daß sie „Kampf GOttes“ so deuten sollen, als ob es so viel heiße als ein großer heftiger Kampf? Warum will ich dem einfältigen Verstande folgen, den die Grammatik gibt.

50. Das hebräische Wort bahal heißt, wechseln, verkehren; wie Ps. 18, 26. 27. steht: „Bei den Heiligen bist du heilig, und bei den Verkehrten bist du verkehrt.“ Wo sich Einer von GOtt abwendet, und anders thut oder lehrt, denn GOtt lehrt und thut, mit etwas Besseres und Trefflicheres thun und verkehrt es doch alles: da wird GOtt auch verkehrt und ändert sich, daß er dem nicht folgt, das ein solcher verkehrter Mensch vorhat, sondern läßt ihn also in seinen Gedanken hinfahren. Daher kommt das Wort „Naphtali“, das heißt auf Lateinisch permutatus, wie wir sagen: Es hat sich gewandelt. Also sagt Rahel auch: Gelobet sei GOtt, denn es hat sich eine Veränderung zugetragen zwischen mir und meiner Schwester. Und es scheint, als habe sie die Worte geredet, da Lea aufgehöret hat zu gebären: da der vierte Sohn geboren war, da bekommt Rahel Kinder durch die Magd, nämlich, Dan und Naphtali. Darum hat sie gesagt: Es hat sich nun gewandt: meine Schwester hat aufgehöret zu gebären, ich aber fange an und fahre nun fort auch Kinder zu bekommen; zwar nicht durch meinen Willen, sondern durch GOttes Gnade und Gabe, der mein Gebet erhört hat.

51. Nun fällt aber eine Frage ein: Wie Jakob innerhalb sieben Jahren zwölf Kinder habe zeugen können? Denn die ersten sieben Jahre hat er umsonst gedient ohne Weib und Kind; in den andern sieben Jahren hat er zwei Schwestern genommen, und läßt es sich ansehen, daß Joseph im letzten Jahr der Knechtschaft Jacobs geboren sei, nach den andern Söhnen allen, das ist, im vierzehnten Jahr. Denn also steht im Text die Ordnung der Historie. Wie hat aber dasselbe geschehen können? Dieweil Lea allein in den ersten vier Jahren vier Söhne geboren hat, so haben ja die andern acht Söhne in drei Jahren müssen geboren werden, welches unmöglich ist. Darauf antworte ich also: Mose hat hier die Figur gebraucht, die man nennt anticipatio (Vorwegnehmen), welche bei den Hebreern hysteron proteron ist; wie wir das hernach erklären wollen. Es ist aber eine solche Figur, wenn man das Erste zuletzt setzt und also die Ordnung umkehrt. Und daraus wird man verstehen, was dieses für eine Veränderung sei, davon Rahel diesem Sohn den Namen gibt. Denn daher kommt es, daß dieses so schwer und dunkel ist, weil es scheint, daß diese Patriarchen alle in den andern sieben Jahren geboren seien. Wenn das nun erklärt sein wird, so wird der Name Naphtali auch etwas klarer werden.

52. Damit aber solches geschehen könne, so muß man zu den andern sieben Jahren noch sechs Jahre setzen, in welchen Jahren Jakob und das Vieh gedienet hat. Denn dieses ist die rechte Ordnung der Historie. In den ersten vier Jahren sind die vier Söhne von Lea geboren, als Ruben, Simeon, Levi und Juda; darnach hört Lea auf zu gebären, und der Rahel wird von der Magd Dan geboren, welcher der fünfte Sohn Jacobs ist, und das fünfte Jahr, darin Lea nicht geboren hat. Dieses Jahr aber füllt Bilha, die Magd Rahels, aus, da sie ihren Sohn Dan geboren hat. Und da hat sich Rahel gefreut von wegen derselbigen Fruchtbarkeit, und kommt noch über das ein anderer Sohn dazu, nämlich Naphtali, im sechsten Jahr, in welchem Lea auch nicht gebieret, daher der Sohn Naphtali genannt wird. Also wollte Rahel sagen: Meine Schwester hat in diesen zwei Jahren aufgehört zu gebären, als, im fünften und im sechsten Jahr: nun wendet sich aber meine Sache und wird damit besser, wie zuvor von diesem Wort, was es eigentlich heiße, aus dem 18. Psalm ist gesagt worden, und heißt in den Sprüchen Salomos Cap. 8, 8. eben auch so viel. In den Worten der Weisheit ist „nichts Verkehrtes“, das ist, GOttes Wort ändert sich nicht, es wendet sich auch nicht, sondern ist schlecht und recht. Allhier aber, spricht Rahel, geschiehet eine Veränderung, welches eine Veränderung GOttes ist. Denn da meine Schwester aufgehöret zu gebären, habe ich nun Kinder bekommen, dieweil ich den Vorzug gehabt, und jetzt immer zunehme, wachse und im Flor stehe, da sie nun verdirret. Also dankt das gottseelige Weib GOtt, daß er ihr Gebet erhöret und ihr Kinder gegeben hat zum wenigsten von der Magd Fleisch.

53. Man soll aber das große Verlangen, so diese Weiber nach den Kindern gehabt, nicht für eine schändliche Begierde oder Geilheit halten; sondern es ist an ihnen beiden eine sehr keusche Begierde und Verlangen gewesen nach den Kindern, also daß sie auch den Mägden übergeben haben das Recht der Mutterschaft und des Segens, dadurch die Nachkommen Jacobs das Land Canaan besitzen sollten; welche Verheißung Jakob ihnen wird fleißig eingeprägt haben. Darum hat eine jegliche begehrt des gebenedeieten Samens und der Nachkommen Mutter zu sein.

54. Und bis daher hat Mose die Linie oder Ordnung der Historie gehalten in der Beschreibung der sechs Jahre. Aber dieselbe Ordnung wird nun geändert und wird ein hysteron proteron daraus; welche Figur, wo sie nicht fleißig beim Lesen der heiligen Schrift, ja, auch sonst beim Lesen aller Historien gemerkt wird, oftmals zuwege bringt, daß der Verstand sehr dunkel wird. Denn jetzt sollte nach der Ordnung dieser Historie folgen die Geburt Josephs nach Naphtali: aber damit verziehet Mose bis in den dritten Paragraph, und kommt wieder auf Lea, welche in zwei Jahren, und auch in dem Jahr, da Joseph wird geboren werden, nicht geboren hat. Und das geschieht durch die Figur hysteron proteron, daß ist, da das Hinterste zuvor gesetzt wird, dadurch oftmals in den Historien die schwersten Stellen erklärt werden.

55. Joseph sollte der Erste sein und nach Naphtali gesetzt werden, aber er wird hernach gesetzt: Lea mit ihrer Magd Silpa sollten hernach stehen, die werden vorher gesetzt. Denn in den Historien sind oftmals mancherlei Dinge auf Eine Zeit, aber doch nicht an Einem Ort geschehen, etliche zu Rom, etliche in Griechenland: da werden solche Dinge nicht nach gewisser Ordnung beschrieben, sondern man muß eine Historie zum Ende führen, darnach muß man dieselbige Zeit wiederholen, und die andere
\quellseite{151}
Historie, wie dieselbe erzählt wird, auch ausführen. Denn es ist wohl einerlei Zeit, aber die Dinge sind mancherlei, die auf eine Zeit zugleich geschehen sind.

56. Derhalben soll man es also in eine Tafel setzen: Die ersten vier Jahre sind Leas, das fünfte und sechste Jahr sind der Magd Rahels, im siebenten Jahr ist Joseph geboren; denn unten (Cap. 41. V. 46.) sagt er zu Pharao, daß er dreißig Jahr alt sei, welches mit der Zeitrechnung sein übereinstimmt. Im siebenten Jahr nach der Hochzeit, so Jakob mit Rahel gehalten, ist Joseph geboren worden, welches ganz offenbar ist. Aber dessen schweigt Mose still, und ehe er von der Geburt Josephs redet, läßt er vorher gehen die andern zwei Söhne von der Lea und von ihrer Magd, welche in den sechs Jahren geboren sind, die zum Dienst Jakobs gethan worden, da er um das Vieh gedienet hat; denn hernach sagt der Text, Jakob habe wiederum in sein Land reisen wollen, da Joseph geboren war; er habe aber mit Laban aufs neue einen Vertrag gemacht noch auf sechs Jahre. So wird es nun klarlich aus den Umständen bewiesen, daß Joseph vor diesen letzten sechs Jahren geboren sei, nämlich im siebenten Jahre, und daß er sechs Jahre lang ernähret und aufgezogen worden sei in Mesopotamien. Benjamin ist sieben Jahre nach Joseph geboren im Lande Canaan, da die Mutter über der Geburt gestorben ist. Also ist es klar und offenbar, daß Jakob seinem Schwiegervater Laban zwanzig Jahre gedient hat, vierzehn Jahre um die Töchter und sechs Jahre um das Vieh. Diese Jahre muß man in dieser Rechnung alle zusammen setzen.

57. Dieweil aber Lea drei Jahre mit der Geburt hat eingehalten, hat sie dem Exempel der Rahel nachfolgen wollen, und also gedacht: Siehe, meine Schwester hat einen Sohn geboren, nämlich Joseph, und wird derhalben von ihrem Manne geliebt als das vornehmste Weib, von welcher er die rechten Erben hofft zu bekommen, und hält jetzt seinen Sohn Joseph für den Erstgebornen: der ist das Licht und der Fürst, dieweil er der Königin und der rechten Hausmutter Sohn ist. Darum hat Lea geeifert, und ist bekümmert gewesen, wie sie ihres Mannes Gemüth ändern und von ihrer Schwester abwenden möchte, und hat darum diesen Rath erdacht, daß sie ihm die Magd geben wollte. Mit derselben hat Jakob den achten Sohn gezeuget, nämlich Gad. Im neunten Jahre wird der neunte Sohn, Asser, geboren, da ist Ruben, der erstgeborne Sohn von Lea, ein Knabe nun acht oder neun Jahre, auf das Feld gegangen, und hat da ohngefähr Dudaim gefunden und bringt sie heim; und da Rahel dieselbigen empfangen hat, macht sie mit Lea einen Bund, daß sie die Nacht bei Jakob schlafen sollte. Denn ein solcher Zank und Zorn ist im Hause unter den Weibern entstanden um einer ehrlichen Ursache willen. Darum wird Lea im zehnten Jahre schwanger und gebieret den Isaschar, im elften Jahre aber den Sebulon. Im zwölften Jahre wird Dina geboren. Also hat Jakob innerhalb zwölf Jahren elf Söhne gezeuget und eine Tochter, nämlich Dina.

58. Aber das bewirt den Leser, daß nach der Geburt der elf Patriarchen die Geburt Josephs zuletzt beschrieben wird, so er doch sechs Jahre zuvor geboren ist, ehe denn Mose solches erzählt. Aber das ist nun klar aus der Regel, davon ich gesagt habe; welche wahrlich viel gebraucht wird in den Büchern der Könige und der Propheten, ja, auch in den heidnischen Historien und derselben, daß eine Geschichte ganz zu Ende gebracht werde, während andere davon unterschieden Geschichten sich zu gleicher Zeit zugetragen haben.

\subsection[Eine Tafel aus dem Chronikon D. Luthers, welche die Ordnung der …]{Eine Tafel aus dem Chronikon D. Luthers, welche die Ordnung der Historie erklärt:}

\begin{center}
\begin{tabular}{lcl}
Lea hat geboren: & & Rahel hat geboren: \\
Ruben. & 1. & 0. \\
Simeon. & 2. & 0. \\
Levi. & 3. & 0. \\
Juda. & 4. & 0. \\
0. & 5. & Dan. \\
0. & 6. & Naphtali. \\
& & Von Bilha. \\
0. & 7. & Joseph. \\
Von Silpa: Gad. & 8. & 0. \\
Asser. & 9. & 0. \\
Isaschar. & 10. & 0. \\
Sebulon. & 11. & 0. \\
Dina. & 12. & 0. \\
0. & 13. & 0. \\
0. & & Benjamin. \\
\end{tabular}
\end{center}

59. Es haben aber diese Leute wunderlich hausgehalten. Denn wir sehen, daß die Weiber Macht gehabt, den Männern ihre Mägde zu geben und dieselben den Männern wieder zu nehmen. Denn da Rahel ihren erstgebornen Sohn Joseph gebieret, hat sie die Magd nicht mehr wollen lassen zum Manne gehen; und Lea auch nicht, da Gad und Asser von der Magd geboren waren. Daran wohl zu sehen ist, wie keusch sie gewesen sind. Denn Jakob hat seines Gefallens, wenn es ihm möchte gelüstet haben, nicht dürfen bei den Mägden liegen, es sei denn mit seiner Weiber Willen geschehen und daß sie es ihm haben zugelassen; welches die Ehebrecher und Hurer nicht pflegen zu thun. Es ist die Lust des Fleisches an Jakob getödtet gewesen, und hat ihr eigen Maß und Ziel gehabt. Er ist nicht heraus gefahren, dahin er selbst gewollt und da ihn seine Lust hingezogen hat. Diese zwei Stücke, nämlich die Figur, anticipatio genannt, und daß hier gerühmt wird die Keuschheit dieser heiligen Menschen, sind fleißig und wohl zu merken. Denn sonst wird niemand diese Historie leichtlich verstehen.

60. Darnach soll man auch erkennen, daß Gott auch bei diesem geringen Stande gegenwärtig sein wolle, welcher sowohl als die andern voller Anfechtung, Jammer und Trübsal ist. Denn Gott regiert, schützt und erhält die Heiligen auch in geringen Dingen, damit sie umgehen, auf daß die Leute, so in geringem Stande sind, nicht verzagen, und die, so in hohem Stande sind, nicht hoffärtig und stolz werden.

61. Denn Gott regiert die Welt also, daß er, was hoch ist, stürzt, und was niedrig und gering ist, erhebt. Er macht zunichte die Weisheit der Weisen und verwirft den Verstand der Verständigen. Diesen Wechsel hält unser Herr Gott. Die Kirche ist in der Welt der Knechtschaft unterworfen; ja, sie ist, wie Paulus sagt, dem Tode ergeben; und gleichwohl, wenn die Welt ohne Kirche wäre, so würden Könige und Fürsten mit aller ihrer Macht und Gewalt müssen zu Boden gehen. Denn sie, die doch jedermann muß unterworfen sein, herrscht und trägt den ganzen Erdboden. Denn die Gottseligen erhalten die Welt und sind doch jedermanns Knechte; wie in einem griechischen Verse steht, den wir also verdeutscht haben:

\begin{quote}
Der Herr muß selber sein der Knecht, Will er im Hause finden recht, Die Frau muß selber sein die Magd, Will sie im Hause schaffen Rath etc.
\end{quote}

62. Also auch, in der Stadt ist allein die Obrigkeit jedermanns Knecht und Unterthan: die Andern, die da angesehen werden, als ob sie Knechte sind, haben gut Gemach, genießen des Friedens und der Ruhe in der Stadt, im Lande und in allen Herrschaften; der aber das Regiment hat, ist servus servorum, ein Knecht aller Knechte.

63. Denn also regiert Gott den Erdboden, daß ihm dienen muß alles, was hoch, mittelmäßig und niedrig ist; die Höchsten sind die Niedrigsten und die Niedrigsten sind die Höchsten. Derhalben hat Herzog Friedrich, Churfürst von Sachsen, weislich und herrlich zu D. Staupitz gesagt, daß seines Trachtens, da er in seinem Herzen die ganze Welt und alle Stände der Menschen hätte übersehen und betrachtet, die Bauern am allerbesten schienen daran zu sein, die sonst in diesem gemeinen Leben, so dem weltlichen Regiment unterworfen ist, die Allergeringsten sind: und zeigt deß diese Ursache an, daß sie allein Frieden und gute Ruhe genießen. Sie werden nicht geplagt und gemartert mit Sorgen und Gefahr von wegen des gemeinen Regimentes. Im Sommer warten sie des Ackerbaues, im Winter mögen sie bei dem Feuer oder Kachelosen sitzen, und von demjenigen leben, das sie durch Gottes Segen erworben haben. Ob sie nun wohl nicht so herrlich leben mit Essen und Trinken, wie Könige und Fürsten, so genießen sie doch der allerbesten Güter, als Friede und Ruhe, und wohnen innerhalb ihres Zaunes viel sicherer und glückseliger, als die Könige und Fürsten in ihren Schlössern oder Festungen wohnen.

64. Dieses ist eine wunderbare Regierung oder Wechsel, so unser Herr Gott hält. Denn was man für das Geringste und Elendeste ansieht, das ist das Allerhöchste und Glückseligste. Es ist in der Welt kein elender Ding denn die Kirche. Denn daher kommen die Klagen der Heiligen und der Kirche, daß sie zu Gott schreien und sagen: Warum wendest du dein Angesicht von mir? Warum vergissest du meiner? Ich predige, bekenne meinen Glauben, thue und leide alles, was mir Gott gebietet; aber es ist kein Mensch so geplagt und verlassen, als ich bin. Da antwortet aber Gott und spricht: Ich will dein nicht vergessen; wie ganz lieblich geschrieben steht im Propheten Jesaia am 49. Capitel, V. 14. 15.: „Zion aber spricht: der Herr hat mich verlassen, der Herr hat mein vergessen. Kann auch ein Weib ihres Kindleins vergessen, daß sie sich nicht erbarme über den Sohn ihres Leibes? Und ob sie desselbigen vergäße, so will ich doch dein nicht vergessen“ 2c.

\quellseite{152}

65. Man sieht es aber äußerlich nicht, daß GOtt so große Sorge für uns habe und sich unser so fleißig annehme. Derhalben muß man lernen und sich gewöhnen an solchen Wechsel und Veränderung der Dinge und Menschen in dieser Welt. Die Heiden haben solches etlichermaßen aus der Erfahrung gelernt, daß der Knecht im Hause Herr ist und der Herr Knecht sein muß; ja, die Mönche klagten dasselbe von ihren Knechten auch. Denn sie hatten etliche Brüder, die sie Conversen hießen, denen die Küche und sonst die Hausarbeit befohlen war; aber dieselbigen herrschten auch über die Andern, daß auch ein Sprüchwort bei ihnen geworden ist, daß sie sagten: Conversus Dominus, aus den Worten, die im Evangelium Luc. 22, 61. stehen: Conversus Dominus, respexit etc.: „Der HErr wandte sich um, und sahe ihn an“ 2c. Damit wollten sie so viel sagen: Der conversus ist unser Herr, der über uns herrscht. Also geht es mit unserer Haushaltung auch zu: da regieren die Mägde die Frauen; die Kinder regieren das ganze Haus und Gesinde, welches ihnen dient.

66. Mit solchen Exempeln lehrt uns GOtt, daß er gegenwärtig bei uns sei und für uns sorge, und daß er unser nicht vergessen wolle, wir leben gleich in welchem Stande wir leben mögen. Die Kirche ist die rechte Königin in der ganzen Welt, aber man sieht an ihr nichts weniger; denn die Welt herrscht und regieret überall: wo aber die Kirche mit ihrem Gebet und Lehre die Welt nicht erhielte, so würde in Einem Augenblick Alles verderben und zu Trümmern gehen. Dasselbe sieht aber und glaubt die Welt nicht; ja, der Türke und der Franzose meinen, sie seien es, so die Welt regieren und erhalten, bis sie dermal einst dahin fallen und zu Schanden werden in ihrem Rathe und Vornehmen; alsdann werden sie erst erfahren und inne werden, daß sie nichts sind: aber es wird alsdann zu lange geharrt sein.

\subsection{II}

\begin{quote}
\textbf{V. 9—11. Da nun Lea sahe, daß sie aufgehöret hatte zu gebären, nahm sie ihre Magd Silpa, und gab sie Jakob zum Weibe. Also gebar Silpa, Leas Magd, Jakob einen Sohn. Da sprach Lea: Rüstig. Und hieß ihn Gad.}
\end{quote}

67. An diesem Orte sollt nun die Geburt Josephs erzählt worden sein; aber es werden allhier, nach der Figur hysteron proteron, die Söhne der Lea aufgezählt. Man soll aber allen Argwohn ablegen und fahren lassen, als sollten diese heiligen Menschen mit unzeitlicher Vermischung und Hurerei eingegangen sein. Denn Jakob begehret die Magd nicht, steht ihr auch nicht heimlich nach ihrer Ehe, daß er sie derselben berauben möchte, wie die Hurer pflegen zu thun; sondern er hätte sich gerne genügen lassen an der Liebe der einzigen Rahel; er geht aber zu der Magd aus Befehl und aus Gehorsam gegen sein Weib, welche der Magd das Ehebett erlaubt und eingeräumt hatte; denn sie ist nun fast daran verzagt, daß von ihrem Leibe Kinder sollten geboren werden. Denn sie war drei ganze Jahre unfruchtbar und gibt zwar dem Manne die Magd nicht gern zum Weibe; sondern wenn sie Hoffnung gehabt hätte, daß sie hinfort mehr sollte schwanger werden, würde sie das nimmermehr gethan haben. Dasselbe wird man sehen an dem, so hernach folgt: da sie selbst wiederum gebieret, schafft sie die Magd wieder ab, und will nicht haben, daß sie von ihrem Manne mehr Kinder gebären soll. Also thut Rahel auch: da Joseph geboren ist, will sie nicht haben, daß sich die Magd wiederum zu ihrem Manne halten soll.

68. Und haben wahrlich die heiligen Mütter ein sehr heißes und starkes Verlangen gehabt, daß sie nur möchten Kinder haben. Darum, da Lea in drei Jahren nicht gebieret, ist sie betrübt, daß ihr derhalben der Mann sollte fremd oder feind werden: und dieweil keine Hoffnung ist, daß von ihrem Leibe mehr Kinder sollten geboren werden, so gibt sie ihm die Magd zum Weibe. Denn da Rahel nur den einigen Joseph geboren hatte, wird sie davon stolzer, denn Lea, da ihr vier Söhne geboren waren: auf denselbigen Sohn sehen Vater und Mutter allein, dazu auch das ganze Gesinde; dessen verwundern sie sich und loben ihn hoch; denn er war von Rahel geboren, welche das liebe und vornehmste Weib war, von der man auch den rechten Herrn im Hause erwartete und den Erben des Königreichs und Priesterthums. Darum ist Lea bekümmert und denkt: Ach, ich armes elendes Weib, nun bin ich wiederum gedemüthigt, meine Kinder werden verachtet und der Mann ist mir auch abgeneigter geworden. Wie sie hernach (V. 15.) Rahel beschuldigt und schilt, da sie sagt: „Du hast mir meinen Mann genommen“ 2c. Und läßt es sich wahrlich ansehen, daß dieses ein fleischlicher Zorn und Eifer gewesen ist; GOtt hat aber denselbigen gnädiglich geduldet, und gar unangesehen solches Zornes und Eifers segnet er beide Lea und Rahel.

69. Derhalben hat nun Lea bei sich selbst also gedacht: Ich werde wahrlich den Mann wiederum in meine Hütte ziehen müssen; ich sehe wohl, daß es dem Vater um den Joseph allein zu thun ist; an dem hat er seine Lust und Freude, weil er von meiner Schwester geborn ist, die mir allezeit vorgezogen worden: ich will ihm aber Silpa zum Weibe geben, daß ich ihn also von Rahel abziehen und mit mir wiederum versöhnen möge. Darum mag man an diesen heiligen Weibern sehen, nicht daß sie unkeusch gewesen sind, sondern daß sie großes Verlangen und Liebe zu Kindern gehabt haben.

70. Wenn die Manichäer und Papisten über solchen Text in der heiligen Schrift kommen, so stopfen sie sowohl Ohren als Augen zu; und nach ihrem Urtheil sinken diese fleischlichen Dinge also, daß sie dieselbigen kaum lesen oder wissen mögen. Augustinus in seinen Bekenntnissen schreibt, daß die Manichäer die heiligen Patriarchen ganz giftig geschmäht und gelästert haben, darum daß sie viel Weiber und Kinder gehabt haben. Da ich, sagt Augustinus, noch ein Manichäer war, spottete ich der Patriarchen und belachte sie als die geilsten Menschen; aber du Manichäer spottetest mein wiederum auch, da ich deinen abenteuerlichen Lügen nachfolgte, daß wenn ich einen Apfel oder Birne vom Baume gebrochen hatte, ich glaubte, daß der Baum darum weinte 2c.

71. Also geht es mit den Papisten auch zu: indem sie an den Eltern hoch verachten, daß sie dazu Liebe haben, daß sie mögen mit einander Kinder zeugen, so wälzen sie sich derweile in allerlei schändlicher Unzucht; ja, es haben sich auch nicht weniger Gerson, Bonaventura, Hugo, Origenes und Hieronymus an diesen Exempeln geärgert. Darum gehen sie daran vorüber und mögen sie nicht mit Fleiß ansehen, ziehen sie nur auf die Allegorien oder heimliche Deutung.

72. Sie sehen nicht, daß im Text klarlich gesagt wird, wie die heiligen Weiber zu GOtt schreien, beten, ihm danken und vertrauen. Item, daß im Text steht: Der HErr hat gehört, GOtt hat angesehen, GOtt hat es gewendet. GOtt, der da angerufen wird, der das Gebet erhört, der gelobt und gebenedeyet ist, der ist mitten in den Historien. Daran sollten sie abnehmen und daraus urtheilen lernen, daß diese Historien rein, züchtig und keusch sind, und daß darin nicht böse Lust und Unzucht gefunden wird, wie sie sich träumen lassen, und damit Hurenjäger pflegen umzugehen. Es ist wohl wahrlich ein wunderlich Ding, daß niemand von den Vätern dieses jemals gesehen hat. Und das noch mehr ist, im Text wird gesagt: Lea hat Jakob ihre Magd zum Weibe gegeben; derhalben sind es ja keine Huren gewesen, sondern eheliche Weiber, so Jakob sind gegeben worden, Kinder mit ihnen zu zeugen.

73. Ich habe aber oftmals gesagt, daß ich einen großen Abscheu habe vor den Allegorien, und daß ich es verdamme, daß man sich derselben befleißigt. Denn die Exempel und Zusprach der Väter schrecken mich und machen mich scheu, welche mit ihren Allegorien die Lehre verdunkeln, dadurch die Liebe, Geduld und Hoffnung zu GOtt sollte erbauen werden, wenn sie uns mit ihrer Speculation und eigenen Gedanken der Allegorien von der Lehre und rechtem Verstande der Worte abführen.

74. Und desselben befleißigen sich vornehmlich Hieronymus und Origenes, ja, Augustinus wäre auch bald dahin gebracht worden, wo ihn nicht der Streit und die Disputation, so er mit den Ketzern hat müssen halten, davon abgezogen hätte. Dieweil aber ich dieselbigen Männer für die höchsten Theologen und Lehrer der heiligen Schrift gehalten, bin ich anfänglich ihnen auch also nachgefolgt. Wenn ich in der Bibel pflegte zu lesen, folgte ich dem Verstande nicht, den die Worte und Buchstaben mit sich bringen, sondern ihren Exempeln nach machte ich Alles zu eitel Allegorien.

75. Derhalben will ich die Theologie Studirenden ermahnt haben, daß sie sich vor solcher Auslegung der heiligen Schrift hüten wollen. Denn die Allegorie oder heimliche Deutung ist schädlich, wenn sie sich mit der Historie nicht reimt; vornehmlich aber, wo man sie anstatt der Historie gebraucht, aus welcher man die Gemeinde besser lehren und unterrichten kann von der wunderbaren Regierung Gottes, wie er alle Stände in diesem Leben regiert, als, das Haus- und Weltregiment und dazu auch die Kirche.

\quellseite{153}

76. Dieweil aber dieselbigen Ausleger in den Historien solcher Dinge nicht achten, müssen sie von Gott wegen Alles auf die Allegorien und fremden Verstand ziehen. Wie in diesem Exempel: weil sie nicht sehen den Rath und die Regierung Gottes, so unter der geringen Gestalt der Haushaltung und des Ehestandes bedeckt ist, so fahren sie zu und dichten einen fremden Verstand von dem beschäulichen und wirklichen Leben; denn dieses sind ihre Sätze, dahin sie das Uebrige alles pflegen zu ziehen. Rahel, sagen sie, soll das beschäuliche Leben bedeuten, Lea aber das wirkliche. Ja, das ist die treffliche große Weisheit, damit sie die Lehre, so der Kirche und Gemeine Gottes nöthig ist, als, vom Ehestande, von der Haushaltung und was dazu gehört, zugedeckt und gar verdunkelt haben.

77. Sie sagen, Jakob habe Rahel lieb, darum daß sie schöne Augen hat. Dieser eignen sie das beschäuliche Leben zu; aber der Lea das wirkliche, welche von Jakob nicht geliebt, sondern verachtet wird. Das haben sie darum gethan, daß sie sich über alle Stände erheben möchten. Denn das wirkliche Leben gehört den Leuten, die mit häuslichen und weltlichen Werken umgehen, die auf dem Felde und im Hause arbeiten: das ist ein sehr verdrießlich Leben, und das keine schöne Gestalt hat, wie Lea blinde Augen gehabt. Das beschäuliche Leben aber findet man in den Klöstern und in den Stiftkirchen der Pfaffen, die keine Mühe haben des Haus- und Weltregimentes halben, die da müßig sitzen und ihre Speculation und hohe Gedanken von Gott haben: sie beten, fasten, haben ihre Gesichte, Offenbarungen und Erleuchtungen. Und ist endlich die Unsinnigkeit und Teufelsbetrug mit der Allegorie von dem beschäulichen Leben so weit gekommen, daß man keinen für einen rechten Mönch gehalten hat, es wäre denn, daß er sonderliche Offenbarung hätte. Und vornehmlich ist das bei den Minoriten gemein gewesen, die haben sich von nichts Anderem träumen lassen, denn allein von den Gesprächen, so Christus, die Engel, die verstorbenen Heiligen und die Jungfrau Maria mit den Seelen gehalten haben. Und man hat noch ein Buch von den Offenbarungen der Brigitte, darin verfaßt ist ein Gespräch Christi, so er mit den Seelen soll gehalten haben. Es ist aber lauter Teufelsbetrug, damit ich doch selbst beinahe wäre gefangen worden, da ich noch ein Mönch war, wo mich nicht D. Staupitz davon gezogen hätte, der mich zum öffentlichen Lehren der Theologie gebracht hat, da ich nach seinem Rath und Befehl war Doctor der Theologie geworden.

78. Derhalben ist nun der Mönche Leben die schöne und liebe Rahel genannt worden, aber das Haus- und Weltregiment hat man mit der häßlichen Lea verglichen. Ja, in den Klöstern hat man auch noch einen Unterschied gemacht zwischen den Priestern und Conversen, und hat diesen das wirkliche, jenen aber das beschäuliche Leben zugeeignet. Aber auf diese Weise werden wir nicht der Lea oder Rahel, sondern den Türken und Tartaren gleich werden. Denn unter den Türken sind auch viel Religiöse, die sich dessen allein befleißigen, daß sie Mahomets Alkoran nach der Allegorie deuten und auslegen, auf daß man von ihnen so viel desto mehr halten möge. Denn die Allegorie ist wie irgend eine schöne Hure, die den Menschen also liebkost, daß man Lust und Liebe zu ihr haben muß. Und vornemlich was müßige Leute sind, die keine Anfechtung haben, dieselbigen meinen, sie seien mitten im Paradiese und sitzen Gott im Schooße, wenn sie irgend mit solcher Speculation und hohen Gedanken umgehen.

79. Und zwar sind erstlich die Allegorien von den närrischen und müßigen Mönchen hergekommen, und sind zuletzt so gar weit ausgebreitet worden, daß auch Etliche aus dem Buche Ovids, Metamorphosen genannt, eitel Allegorien gemacht haben. Aus einem Lorbeerbaum haben sie die Jungfrau Maria gemacht und Apollo sollte Christus sein. Wiewohl nun solches ungereimt Ding ist, so geschieht es doch gleichwohl, wenn man es den Jünglingen, die noch unerfahren sind, aber doch zum Studiren und zu guten Künsten Lust und Liebe haben, vorhält, daß sie sich solches anfänglich sehr wohl gefallen lassen und sich ganz und gar auf solche heimliche Deutung begeben.

80. Darum bin ich den Allegorien feind. Wo aber ja derselbigen jemand gebrauchen wollte, der sehe zu, daß er es mit rechtem Verstand thue. Denn zu allererst soll man den Verstand suchen, so die Historie mitbringt, welcher uns recht und gewiß lehrt und unterrichtet, der da streitet, vertheidigt, gewinnt und bauet. Wenn der recht und rein ist, so mag man darnach Allegorien und heimliche Deutung suchen. Aber nicht wie die Mönche gethan haben und auch nicht von dem beschäulichen Leben; sondern daß sich die Allegorie mit der Historie reimen möge, und daß darin begriffen werde das Heiligtum vom heiligen Kreuz, das ist, die rechte Lehre vom Kreuz, von Glauben, Hoffnung, Liebe und Geduld. Sie soll keinem Mönch in sich begreifen, der da sitzt und sich träumen läßt, er rede mit Christo, und sich seiner Offenbarung rühmt, und wie er die Stimme der Engel oder der Jungfrau Maria gehört habe; denn er sitzt müßig und ist ohne Kreuz und Anfechtung. Ein solch beschäulich Leben ist verflucht und verdammt.

81. Derhalben soll die Allegorie mit der Historie übereinstimmen, und darauf soll als auf den rechten Grund Edelsteine, Gold und Silber gebaut werden: sonst ist die Allegorie schädlich und bringt nur Irrthum, oder wo sie auch am besten ist, so sind es doch lauter Stoppeln und eitel Narrenwerk. Und man soll ja fleißig unterscheiden den Grund oder die Lehre an sich selbst, und das Gebäu, so darauf gebaut wird. Die Allegorie soll der Grund nicht sein, sondern soll auswendig der Rand sein; es sei gleich Gold, Edelsteine, oder auch Stoppeln, wenn es nur auf dem rechten Grunde steht; sonst ist es lauter Teufelsbetrug, wo sich die Allegorien mit dem Grunde nicht reimen.

82. Denn wenn Gott wird nicht einem Jeglichen sonderliche Offenbarung geben, er wird seine neuen sieben Gebote geben; sondern er hat uns an dieses Gebot gebunden, das man vom Himmel herab gehört hat, da Gott der Vater von seinem Sohn Christo sagt: „Den sollt ihr hören.“ Item Luc. 16, 29.: „Sie haben Mosen und die Propheten.“ Er will haben, daß wir sein Wort hören und demselbigen glauben sollen; daselbst hast du zu speculieren übrig genug, da bist du abgezogen nicht allein von äußerlichen Geschäften, sondern auch von den Gedanken deines eigenen Verstandes und Vernunft, daß du da den rechten Sabbath hälst; und daß dein Wahn und Verstand aufhöre und gar nicht gelte.

83. Das wirkliche Leben aber ist die Liebe, oder der Glaube, der durch die Liebe thätig ist; desgleichen, daß man Geduld habe, wo das Kreuz vorhanden ist. Dieses ist die Uebung mit allerlei Leben, als, mit dem beschäulichen und wirklichen Leben, nicht allein in den Klöstern bei den Mönchen und Conversen, sondern auch im Haus- und Weltregimente bei allen Menschen, die im Glauben leben, und mit äußerlichen Werken umgehen, wie gar keinen Schein dieselben haben und wie unschätzig sie auch sein mögen. Aber die menschliche Natur will sich an der gemeinen Lehre und Wort des Evangeliums nicht genügen lassen, will auch mit den Sacramenten nicht zufrieden sein, welche durch den Sohn Gottes und die Apostel gegeben und durch ihre Nachfolger ausgebreitet worden sind; sondern will immer etwas Sonderliches haben. Da kommt denn der Teufel und redet mit dir. Wo du ihm nun nicht wirst antworten können: Also sagt Gottes Wort, du redest, das dem Worte zuwider ist, hebe dich weg Satan: so wirst du bald verführt werden.

84. Dieses rede ich und lehre es mit sonderlichem Eifer, und daß ich über die ungereimten Allegorien unwillig und zornig bin, darum daß ich gesehen habe, wie die Mönche aber und Papisten sich neben den Manichäern dieser Zeit mit ihrer Phantasie und Narrenwerke so schändlich beschämet haben. Die Mönche aber und Papisten sind den Manichäern gleich geworden, und haben die Kirche erfüllt mit eitel Brigittischen und anderen dergleichen Büchern, damit sie das Haus- und Weltregiment verachtet und gestraft haben, wiewohl sie solches Regimentes niemals haben entbehren können.

85. Wenn du aber diese Historie etwas fleißiger ansehen wirst, so wirst du darin schöne Exempel sehr großer Keuschheit im Ehestande sehen, dergleichen Keuschheit nie ist gefunden worden bei denen, die den Stand der Jungfrauschaft und Keuschheit so hoch gerühmt haben, etliche sonderliche und wunderliche Exempel aufgenommen, als, Johannes des Täufers, der Jungfrau Maria und dergleichen Anderer mehr. Die Andern behalten immer noch die Brunft und Lust zur Unzucht. Darum kann gar keine Keuschheit mit den Exempeln dieser Weiber verglichen werden. Denn dieselben werden mit sehr schönen Worten beschrieben und gezieret, nämlich, wie sie beten, seufzen, schreien, glauben, und im Ehestande auf GOtt hoffen, daß er ihnen werde Kinder geben. Die Unkeuschheit aber denkt nicht an GOtt, ruft ihn nicht an, dankt ihm auch nicht; sondern ist blind, toll und thöricht. Diese Weiber aber beten, weinen und sagen: O GOtt! gib uns doch Kinder \&c. Das sind nicht Worte eines unzüchtigen Weibes, sondern einer sehr keuschen und ehrlichen Matrone.

\quellseite{154}

86. Darum sind ihre Exempel sehr nützlich und nöthig in der Kirche, weil dadurch bezeuget wird, daß GOtt bei dem Stande des Hausregiments sein, und denselben regieren und schützen wolle; vor welchem dein ehelicher Stand, dein Mönch- und beschäulich Leben vor GOtt stinket, ja, auch in deinem eigenen Gewissen. Und ist billig und recht, daß die Mönche um solcher Verachtung willen des Ehestandes dermaßen gestraft werden. Denn dieweil sie diese und dergleichen Exempel verlassen und die Werke GOttes vor ihnen sinken, so werden sie wiederum auch verworfen und verachtet von GOtt, bis sie in scheußliche Unzucht und Unkeuschheit gefallen sind.

87. Was aber nun das Wort gad seiner rechten Ableitung nach eigentlich heißen soll, darüber zanken die Rabbinen der Juden und sind in der Auslegung nicht eins. Ich halte aber dafür, daß es recht verdolmetscht werde, daß es so viel heiße, als gegürtet, gewaffnet, gerüstet zum Kriege: daher kommt gedud, Soldat, Kriegsknecht. Hieronymus verdolmetscht dasselbe Wort im Lateinischen: Excursores in exercitu, das ist, Kriegsknechte, die zum Streit gerüstet sind. Wir haben es aber mit gutem Bedacht und Wissen also verdolmetscht: „rüstig“ oder gerüstet; wiewohl ich weiß, daß es bisweilen so viel heißt, als Zerschneidung, bisweilen auch wohl Glück. Wie es Münster und die Andern verdolmetscht haben. Und dasselbe geben wir ihnen gern zu; denn es hat es Augustin auf dieselbe Weise auch verdolmetscht: und wir strafen ihre Meinung nicht, folgen aber gleichwohl derselben nicht nach, wie wir es auch Jes. 20. und 65. anders gemacht haben denn sie: da in der gemeinen und der andern lateinischen Dolmetschung also stehet: Qui ponitis fortunae mensam (Ihr richtet dem Glück einen Tisch zu), das haben wir im Lateinischen gegeben: Qui ponitis Marti mensam: Ihr richtet dem Mars einen Tisch \&c. Denn die Juden haben auch den Abgott Mars geehrt nach dem Exempel der Griechen. Wir sind bei der Wurzel des Worts geblieben, welche bedeutet, gerüstet sein. Und diese unsere Auslegung bestätiget Mose hernach im 49. Cap. V. 19., da im Segen der Kinder Gad gesagt wird: „Gad gerüstet, wird das Heer führen.“ Dieweil wir nun ein solch ansehnlich Zeugniß haben von Mose, wollen wir uns daran wider das Zanken der jüdischen Rabbinen genügen lassen. Gad soll vorher gehen, soll Vorgang sein; darum haben wir es verdolmetscht „rüstig“, es will angehen. Wer den Andern aber folgen will, der mag es gar wohl thun. Rahel hatte gesagt, da Naphtali geboren war: „Ich werde es meiner Schwester zuvor thun“, ich habe ihr angekommen. Und Lea sagt hier nun auch: Ich werde es nun meiner Schwester wieder zuvor thun; ich habe jetzt den Ueberwind. Sie will sagen: Es geht rüstig an, daß ich auch einen Knecht habe, der da gewinnen und mein Geschlecht ausbreiten wird.

88. Also haben die heiligen Weiber unter einander gestritten und geeifert, daß sie beide vor großem Eifer ihrer mütterlichen Herrlichkeit sich verziehen und dieselbige eine jegliche ihrer Magd übergeben haben; denn das war eine große Ehre und Herrlichkeit, daß ein Weib von Jakob hat mögen schwanger werden und Kinder gebären, welcher die Verheißung hatte von dem gebenedieten Samen. Welche Verheißung Jakob ohne Zweifel ihnen fleißig und stets wird eingeprägt haben: Also hat mir Abraham, Isaak und GOtt selbst gesagt, daß von meinen Erben der Heiland der ganzen Welt kommen wird. Denn ein solcher heiliger gottseliger Ehemann hat ja seinen Weibern die Dinge nicht verschwiegen, so zum Hausregiment gehören und auch was die Erben und Nachkommen belanget hat.

89. Aus solcher Predigt sind die heiligen Weiber im Glauben brünstig geworden, daß es ihnen nicht um die fleischliche Lust, sondern um die Kinder ist zu thun gewesen, um der Herrlichkeit willen der künftigen Erben und Nachkommen. Denn wo sie nur allein das Beischlafen gesucht hätten, so würden sie sich um die Kinder nicht also bekümmert haben, hätten auch dem Manne ihre Mägde nicht zu Weibern gegeben. Rahel würde gesagt haben: Was liegt mir daran? wenn ich schon nicht gebäre, so halte ich mich doch gleichwohl zu meinem Mann; daran habe ich genug. Es wird aber klar angezeigt, wie sie die Wollust und das Beischlafen verachtet haben, und sind von diesem Verlangen gar entfernt gewesen, daß sie möchten Mütter des verheißenen Samens werden; darüber haben sie untereinander gezankt, auch bis an den Neid. Darauf geben aber die Papisten keine Achtung; denn der Gottlose soll die Herrlichkeit GOttes nicht sehen.

\begin{quote}
\textbf{V. 12. 13. Darum gebar Bilha, Leas Magd, Jakob den andern Sohn. Da sprach Lea: Wohl mir, denn mich werden selig preisen die Töchter. Und hieß ihn Asser.}
\end{quote}

90. Bilha, die Magd Leas, gebieret ihrem andern Sohn, um desswillen sie sich abermal sehr freut; denn sie haben nur begehrt, Kinder zu haben. Darum soll man in dieser Historie nicht Exempel fleischlicher Unzucht, sondern der Keuschheit suchen. Lea begehrt, daß ihr Mann an ihr ein Wohlgefallen haben möge von wegen der Fruchtbarkeit; welches ja eine löbliche Tugend ist an einem Weibe, das da begehrt bei ihrem Manne zu wohnen und keinem Andern nachzufolgen, ja, die darum bekümmert ist, daß sie ihm allein gefallen möge, bei ihm allein möge Gunst haben, und vornehmlich einem solchen großen Manne, dem der Heiland der Welt ist verheißen worden. Darum ist solches eine rechte Eheliebe gewesen, voller Gottseligkeit, Keuschheit, Gehorsam, und nicht der fleischlichen Unzucht, wie es die Römische Sau und viele von den Vätern und Mönchen ausgelegt haben.

91. Asser heißt, selig. Nun werde ich es gut haben, sagt sie, und alle Weiber und Töchter werden mich selig preisen; ich werde dem Jakob eine liebe Hausmutter sein, also auch, daß man sagen wird, Jakob habe Lea lieb und achte der Andern als der Unfruchtbaren nicht mehr.

\section[Dritter Theil]{Dritter Theil. Wie Lea von neuem schwanger wird und dem Jakob Kinder gebieret; und wie endlich auch Rahel schwanger wird und gebieret}

\subsection{I}

\begin{quote}
\textbf{V. 14—16. Ruben ging aus zur Zeit der Weizenernte, und fand Dudaim auf dem Felde, und brachte sie heim seiner Mutter Lea. Da sprach Rahel zu Lea: Gib mir der Dudaim deines Sohns ein Theil. Sie antwortete: Hast du nicht genug, daß du mir meinen Mann genommen hast; und willst auch die Dudaim meines Sohns nehmen? Rahel sprach: Wohlau, laß ihn diese Nacht bei dir schlafen um die Dudaim deines Sohns. Da nun Jakob des Abends vom Felde kam, ging ihm Lea hinaus entgegen, und sprach: Bei mir sollst du liegen; denn ich habe dich erkauft um die Dudaim meines Sohns. Und er schlief die Nacht bei ihr.}
\end{quote}

92. Ruben ist ein Knabe von acht oder neun Jahren gewesen, da dieses geschehen ist; er ist kein Ehemann gewesen und hat nicht älter sein können denn neun Jahre. Derselbe ist nun hinaus gegangen zu seinem Vater, da er die Schafe Labans geweidet oder sonst auf dem Felde gearbeitet hat, wie die Kinder pflegen zu ihren Vätern zu gehen; vielleicht mag ihn auch die Mutter wohl dahin gesandt haben. Dieweil es nun zur Zeit der Weizenernte war, so hat er ohngefähr Dudaim gefunden, und der Knabe hat selbst nicht verstanden, was er gefunden und seiner Mutter gebracht hatte. Es scheint aber, daß es irgend eine schöne Frucht gewesen sei, da er an der Farbe Gefallen gehabt hat. Welcherlei Frucht oder Blume es gewesen sei, kann man nicht wissen. In der gemeinen lateinischen Uebersetzung steht, er habe gefunden mandragorae, und so haben es die andern Dolmetscher auch gegeben. Das soll, wie man meint, eine solche Frucht sein, die da wächse vom Wasser, so die Diebe, wenn sie an den Galgen gehängt sind, lassen; wie die Schreiber, da sie diese mandragorae beschreiben, viel ander ungereimtes Ding vorbringen. Aber es sind lauter Fabeln und Lügen, und mögen vielleicht von den Juden hergekommen sein. Es ist aber der Knabe nicht zu dem Galgen gegangen, sondern auf das Feld.

93. Es muß übrigens gewißlich eine solche Frucht gewesen sein, welche um die Weizenernte pflegt reif zu werden, die da etwas langsamer kommt; wie in unsern Landen die Früchte sind, so man nennt, Nachtschatten, Judenkirschen, Brombeeren, Himbeeren. Vom Nachtschatten (solanus) verstehen es auch der Andern viele, dessen Frucht eine sehr schöne Farbe hat und dazu auch einen guten Geschmack, und man gebraucht es viel in der Arznei. Oder es müssen Brombeeren gewesen sein, die auch einen süßen lieblichen Geschmack haben. Dieses will ich aber den Grammatikern und Aerzten befohlen sein; denn ein jeglich Land hat seine sonderlichen und eigenen Früchte und Kräuter. Dem Knaben hat die Farbe wohl gefallen, und da er heim gekommen ist, haben die Weiber über die Dudaim einen Streit gehabt, und Rahel hat so großes Verlangen nach den Dudaim, daß sie Lea den Mann hergibt und sich ihres Rechtes verzeiht, nur daß sie von dem Knaben die Dudaim empfangen möge.

\quellseite{155}

94. Dieses ist aber über die Maßen lächerlich und kindisch Ding, daß man auch nichts Schlechteres oder Geringeres sagen oder schreiben könnte. Warum ist es denn nun gleichwohl beschrieben? Antwort: Man soll allezeit vor Augen haben, daß ich zum öftern pflege zu sagen, nämlich, daß der heilige Geist der Meister sei dieses Buchs, der hat selbst Lust dazu, daß er also spielen und scherzen möge, solche geringe kindische Dinge und die nichts werth sind, zu beschreiben, und dieselbigen hält er uns vor, daß man es in der Kirche lehren soll, gleich als zu großer Besserung. Sind aber dazumal nicht größere und herrlichere Dinge geschehen unter den Königen in Egypten, zu Babel, in Persien und Palästina etc., welche eher lesenswerth wären als dieses geringe Ding? Wozu dient es doch, daß man solch Ding erzählt, wie ein Knabe Dudaim dahin bringe, und zwei Schwestern mit einander einen Streit halten, wie sie mit dem Manne und Bette tauschen wollen?

95. Es sind ohne Zweifel zu der Zeit mehr und größere Dinge geschehen, denn die hier beschrieben werden. Aber diese geringen Dinge sind voll Trostes und guter Lehre, mehr als jene. Und man soll sie nicht oberflächlich übergehen, vielweniger aber soll man sie verachten, als wären es unnützliche Dinge und die nichts werth wären, welches die Papisten thun. Laß es doch geringe Dinge sein, und laß dagegen ihre Dinge groß sein und einen herrlichen Schein haben; laß sie Könige, Bischöfe und Cardinäle sein etc.: der heilige Geist aber und Gott, der Schöpfer, läßt sich also hernieder, und hat Lust, daß er mit seinen Heiligen scherzen und Kinderspiel treiben möge in solchen geringen Dingen; nicht mit Auferweckung der Toten oder anderen großen Wunderwerken, sondern da Rahel und Lea mit einander streiten über den Dudaim; und will haben, daß man solch Ding der Kirche und Gemeinde Gottes soll vorhalten.

96. Was sind aber nun die Historien der Heiden, so Virgili, Homer, Livius und Andere beschrieben haben, wie herrlich sie auch mit Worten geschmückt sind? Es sind Historien der Griechen, von Alexander und von Hannibal; es mangelt ihnen aber an der Herrlichkeit, an der Krone, welches Gottes Wort und Verheißung ist. Ja, an der Krone mangelt es ihnen; darum ist es mehr arm Ding damit, denn daß es sollten rechtschaffene Historien sein. Denn was ist eine Historie ohne Gottes Wort? Wenn der Herr sagt: Dieses ist mein Wille, meine Ehre, dieses gefällt mir wohl, dazu habe ich Lust, hier will ich wohnen; wie kindisch und unsäglich nun dieses Ding ist, so hat es doch alles einen unermeßlichen und unendlichen Nachdruck, nämlich, Gottes Wort. Die Historien aber des großen Alexander, Julius Cäsar etc. sind wohl trefflich und sehr herrlich: es mangelt ihnen aber an dem rechten Schmuck, darum sind sie wie Spreu vom Winde zerstreut und haben das rechte Gewicht oder Nachdruck nicht, es ist nur eine Herrlichkeit des Bauches und Fleisches.

97. Darum sollen wir diese Historien groß achten und uns damit trösten, wiewohl wir arme, elende, verworfene und verachtete Leute sind; wie wir in Wahrheit sind. Denn es ist auf Erden nichts elender und verachteter, denn ein rechter Christenmensch ist, der dem Worte Gottes glaubt. Die Andern alle sind Fürsten, große Herren, sind wohl daran, richten alle ihre Dinge aus nach ihrem eigenen Gefallen; allein das Volk, so das Wort hat und an Gott glaubt, ist ein Kind des Kreuzes und der Plage. Und wo die Welt nicht da ist, daß sie dasselbe Volk kreuzigen und plagen möge, so ist doch der Teufel mit seinen Engeln da. Aber dagegen ist das unser Trost, daß die Christen mit Weib, Kindern und Gesinde bei einander wohnen, und ohne allen Schein haushalten; welches denn äußerlichen Ansehen nach lauter Kinderspiel ist; gleichwohl sind dasselbe die rechten und göttlichen Historien, die auch mit der Ehre und Herrlichkeit des göttlichen Wortes geschmückt sind.

98. Wir sollen aber das Wort annehmen, und uns fest daran halten, und sollen bedenken, daß wir getauft, absolvirt und mit Gottes Wort gelehret sind. Wir sollen Gott danken und fröhlich sein; es falle gleich der Himmel und die ganze Welt über einen Haufen und treffe uns dazu auf den Kopf, so sind wir doch dessen gewiß, daß unser Kinderspiel und armes Elend geschmückt und gekrönt ist mit göttlicher Ehre und Herrlichkeit, nämlich, mit Gottes Wort. Und um solcher Ursache willen werden diese Historien heilig genannt, und sind gemein allen Menschen, die Gottes Wort haben, in welchen Gott wirkt, und dasselbe mit Wohlgefallen, mit seiner Gnade und Barmherzigkeit. Welchergestalt er nicht umgeht mit Alexander, dem Großen, Scipio, Cicero und Hannibal etc. Darum sollten wir uns freuen, und uns an solchem göttlichen Wohlgefallen genügen lassen, dankbar und fröhlich sein. Zudem sollen wir auch mit Geduld leiden, was für Mühseligkeit und Beschwerung vorfallen mag. Und wie schlecht, gering, weibisch und voller Jammers unsere Werke sind, so sollen wir doch diesen Titel dabei schreiben: Gottes Wort, dadurch alles, was wir thun, herrlich und ewig bleiben wird; so dagegen die weltlichen Historien ewig elend und erbärmlich sind.

99. Also werden wir verstehen, daß uns dieses vom Heiligen Geist nicht vergeblich ist vorgehalten worden, daß wir es lesen, lehren und glauben sollen. Der Papst mit seinem Anhange verachtet solches mit großer Hoffahrt, und läßt sich dieweil träumen, daß er jetzt schon seinen Stuhl im Himmel bei St. Petro erhöht habe, und wie Tezel, der den Ablaß hin und wieder austrieft, rühmete, er wollte mit St. Petro nicht wechseln noch wegen der vielen unzähligen Seelen, die er durch den Ablaß selig gemacht hätte. Aber das ist ein eitler und teuflischer Ruhm. Denn Gott sieht nicht auf das Niedrige auf Erden, und sagt also: Ich habe dieses Weiblein geschaffen und habe sie dem Manne zugefüget, die zeugen nun Kinder mit einander, und dulden den Jammer und das Elend, so der Ehestand mitbringt: darum gefallen sie mir darin wohl, daß sie sich im Ehestande zusammenhalten nach meiner Ordnung. Und die Eheleute sind auch bei den Heiden viel angenehmer denn alle Philosophen, die da wunderlich Ding speculiren vom Cölibat und ehelosen Leben. Wie viel mehr sollen sie bei uns angenehm sein, die wir getauft, gelehrt, gekrönt und geschmückt sind mit der Herrlichkeit des Wortes Gottes, und die wir die Engel zu Hütern haben?

100. Darum sollen wir uns rühmen und freuen dieser kindischen und häuslichen Werke halben, sintemal sich der Heilige Geist demüthigt, dieselben zu beschreiben, auf daß er damit anzeige, daß er bei uns sei, die wir glauben und das Wort haben, und daß auch das Kinderspiel von den Dudaim und närrischem Geschwätz, so die Weiber darüber haben, Gott wohlgefalle; allein daß wir nicht verzagen, sondern im Glauben beständig bleiben.

101. So laß den Teufel hergehen und machen, was er nicht lassen kann, wir sind in einem geistlichen Beruf, wir haben das Wort, die Taufe, Absolution, und sind dessen gewiß, daß Gott mit uns und für uns ist. Ich wollte aber lieber mit Gott in der Hölle sein, als ohne Gott im Himmel. „Denn ob ich schon wanderte im finstern Thal“, spricht David im 23. Psalm V. 4., „fürchte ich kein Unglück; denn du bist bei mir, dein Stecken und Stab trösten mich.“ „Bettete ich mir in die Hölle, siehe, so bist du auch da“, wie im 139. Psalm V. 8. steht. Wenn ich dein Wort habe, so frage ich nach dem höllischen Feuer nicht. Wiederum aber gelüstet mich auch nicht im Himmel zu sein, wenn du nicht bei mir wärest. Denn wo Gott ist, da ist auch das Reich Gottes: wo das Wort ist, da ist das Paradies und Alles. Darum sollen wir Gott für das Wort danken; denn wir sind jetzt schon im Paradiese und unter Gottes gnädigem Wohlgefallen, der sich Alles gefallen läßt, wir essen gleich Dudaim, oder Kraut, oder Milch; denn wir sind in seinem Wort.

102. Diese Lehre soll man aus diesem Text nehmen, welche wahrlich trefflich und der Kirche nöthig ist. Nun muß man auch etwas sagen von dem Zank. Denn man muß diese Historien nicht also ansehen, wie sie die Papisten für solchen Scherz halten, wie die unzüchtigen Weiber pflegen zu treiben. Diese heiligen Weiber begehren Anderes nichts denn die natürlichen Früchte ihres Leibes. Denn das Weib ist von Natur dazu geschaffen, daß sie Kinder gebären soll; darum hat sie Brüste, sie hat Arme, Kinder zu nähren, aufzuziehen und zu tragen. Das ist des Schöpfers Absicht gewesen, darauf er gesehen hat, da er den Menschen geschaffen, daß die Weiber gebären und die Männer Kinder zeugen sollten, ausgenommen die, welche GOtt selbst ausgenommen hat, von denen gesagt wird Matth. 19. Sonst ist die Creatur also von GOtt geschaffen, daß sie soll fruchtbar sein. Derhalben ist an diesen Weibern zu sehen, wie sie mit einander streiten nicht aus Unzucht, sondern aus Liebe und Verlangen, daß sie nur Kinder möchten haben; denn sie wissen wohl, daß sie dazu geschaffen sind. Darum begehren sie, daß ihr Leib, Brüste, Hände und alle Glieder GOtt dienen mögen.

\quellseite{156}

103. Und wiewohl sie schon vorhin Kinder gehabt haben, so haben sie doch gedacht, daß dieselbigen sterblich wären und daß sie derselben wiederum möchten beraubet werden. Darum wenn ein Weib auch hundert Kinder gehabt hätte, würde sie doch derselben noch immer mehr begehrt haben. Denn sie haben gesehen auf die Erben und Nachkommen, so den Vätern verheißen waren, auf welche die Predigt gegangen war, die sie von ihrem Manne gehört hatten. Also muß man diese heiligen Weiber rühmen, und den schändlichen Papisten das Maul stopfen, die nach ihrer eigenen Unzucht von der Keuschheit und dem Ehestande der Patriarchen urtheilen wollen; wie auch Augustinus durch den Irrthum der Manichäer dahin gebracht war, daß er es auch dafür gehalten hat, daß der Patriarchen Leben ein unkeusch Leben gewesen wäre. Alhier aber hören wir, wie Lea gebetet, geschrieen und geseufzt habe um Kinder, daß sie auch von GOtt erhört worden sei; welches ja einem unzüchtigen Weibe nicht gebühret, sondern es gebühret vielmehr einer frommen, züchtigen und heiligen Matrone, die der Unzucht feind ist, und nur begehrt, daß sie Kinder gebären möge.

\begin{quote}
\textbf{V. 17. 18. Und GOtt erhörete Lea, und sie ward schwanger, und gebar Jakob den fünften Sohn, und sprach: GOtt hat mir gelohnet, daß ich meine Magd meinem Manne gegeben habe. Und hieß ihn Isaschar.}
\end{quote}

104. Der fünfte Sohn von Lea ist im zehnten Jahre geboren worden. Lea aber dankt GOtt. Und wie sie zuvor gebetet und gemeint hat, also lobt sie nun jetzt den HErrn mit fröhlichem Herzen. Die Grammatici bei den Hebräern disputiren von dem Namen Isaschar, wie man denselben lesen solle. Es ist aber daran nichts gelegen, du liesest gleich Isaschar oder Issachar.

\begin{quote}
\textbf{V. 19. 20. Abermal ward Lea schwanger, und gebar Jakob den sechsten Sohn; und sprach: GOtt hat mich wohl berathen; nun wird mein Mann wieder bei mir wohnen, denn ich habe ihm sechs Söhne geboren. Und hieß ihn Sebulon.}
\end{quote}

105. Dieses sind alles Worte und Werke der sehr heiligen Matronen, nämlich, Beten, Weinen, GOtt danken. Es sind nicht Worte oder Werke schändlicher unzüchtiger Weiber, sondern ehrlicher Matronen. Und was können wir vor GOtt mehr oder Größeres thun, denn diese zweierlei Werke, Beten und Danken? Das Erste ist, daß man das Wort höre, welches uns von GOtt gegeben wird, dazu wir nichts thun, sondern nehmen es nur an, wie es uns vorgetragen wird. Das Andere ist, wo man das Wort gehört und angenommen hat, daß man alsdann beten und GOtt um Hülfe anrufen soll, und wo man sie erlangt hat, GOtt dafür danken und loben. Wir beten aber nicht allein mit dem Munde oder äußerlichen Worten, sondern auch mit herzlichem Seufzen, mit allen Kräften und Gliedern, welches ein Gebet ohne Aufhören ist.

106. Das hebräische Wort zebadani steht sonst an keinem Ort mehr in der heiligen Schrift; wir haben es dem Exempel der Andern nach im Lateinischen verdolmetscht: dotavit: GOtt hat mich berathen. Ob es aber also recht verdolmetschet ist, kann ich nicht eigentlich wissen; denn man kann es mit keinem andern Exempel erklären. Sie hat also sagen wollen: GOtt hat mir etwas Gutes gegeben. Du verstehst nun gleich dadurch die Aussteuer, so ihr der Vater möchte gegeben haben, oder sonst irgend ein Geschenk. Wir sehen hier aber Anderes nichts, denn sehr große Göttlichkeit und wie man des weiblichen Geschlechtes und des Ehestandes recht gebrauchen soll; wie es denn diese Leute gebraucht haben. Die Jungfrauen haben zwar solche Ehre und Ruhm nicht, weder im Alten noch im Neuen Testament.

\begin{quote}
\textbf{V. 21. Darnach gebar sie eine Tochter, die hieß sie Dina.}
\end{quote}

107. Bisher hat Mose erzählt die Kinder der Lea, nämlich, sechs Söhne und eine Tochter, desgleichen zwei Söhne von der Magd; nun folgt, wie Joseph von Rahel geboren sei.

\subsection{II}

\begin{quote}
\textbf{V. 22—24. Der HErr gedachte aber an Rahel, und erhörte sie, und machte sie fruchtbar. Da ward sie schwanger, und gebar einen Sohn, und sprach: GOtt hat meine Schmach von mir genommen; und hieß ihn Joseph, und sprach: Der HErr wolle mir noch einen Sohn dazu geben.}
\end{quote}

108. Droben ist gesagt worden, daß man auf die Ordnung der Historie müsse Achtung geben, und daß Mose die Figur gebraucht habe, die genannt wird Wiederholen (recapitulatio) oder Vorwegnehmen (anticipatio); denn Joseph ist im siebenten Jahre geboren worden, darum sollte seine Geburt lange zuvor beschrieben worden sein. Dieweil aber Mose damit verzogen hat bis hier zuletzt, so kommt er nun wieder auf das, so erst geschehen, welches in der Beschreibung das letzte geworden ist. Denn das ist klar am Tage, daß Joseph im siebenten Jahre des Dienstes seines Vaters geboren ist, sintemal Jakob nach der Geburt Josephs noch sechs Jahre gedient hat. Darum sind die vier Söhne und die Tochter Dina nach ihm in denselbigen sechs Jahren geboren, wiewohl sie in der Historie zuerst aufgezählt worden sind. Diese Figur, Vorwegnehmen genannt, macht oftmals die Erzählung etwas dunkel; darum muß man darauf fleißig Achtung geben, wenn man die Historie lieset, daß der Leser, so darauf keine Acht hat, nicht betrogen oder verhindert werde.

109. Es hält uns aber der heilige Geist abermals vor nicht allein eine Entschuldigung, sondern auch ein großes Lob und Ruhm dieser so heiligen Weiber, die nicht von wegen der Unkeuschheit brünstig sind, sondern großes Verlangen und Begierde haben nach Kindern. Und also wird uns damit auch gerühmt der Ehestand Jakobs, wie derselbe so keusch gewesen, daß man darin schier gar nichts Fleischliches gespürt hat; wiewohl es unmöglich ist, daß im Ehestande gar nichts Fleischliches sein sollte. Jakob aber wird von GOtt also gezieret, daß sein Ehestand anzusehen ist, als ob es ein rechter englischer Ehestand sei. Denn, Lieber, siehe doch, was daselbst für schöne Uebungen sind der Gottseligkeit und des allerhöchsten Gottesdienstes. Rahel betet, seufzet, weint bis in das sechste Jahr; dieselbige Zeit über ist sie wohl geplagt und versucht worden mit großer Geduld. Denn es ist eine große Last und ein sehr schweres Kreuz gewesen, daß sie die Schmach der Unfruchtbarkeit so lange Zeit hat tragen sollen. Wie hat sie unterdes so viel Weiber gesehen, die Kinder geboren haben! Ja, sie hat gesehen, daß ihre Schwester Lea auch mit der Ehe gezeuget gewesen, daß sie eine Mutter war und Kinder gehabt hat, nämlich vier Söhne. Ich halte aber dafür, daß sie von solcher Traurigkeit also matt geworden, daß sie lediglich nicht allein der Unkeuschheit, sondern auch sonst aller Freude im Herzen vergessen und sich derselben begeben hat: sie hat keine Lust weder zu lachen, noch zu scherzen gehabt, sondern Tag und Nacht zugebracht mit stetigem Seufzen, mit Heulen, Weinen, Wehklagen und mit Beten.

110. Derhalben, sage ich, wird die Ehe Jakobs und Rahels als ganz züchtig und keusch beschrieben: und wenn etwa Rahel hat bei dem Manne wohnen und sich zu ihm halten wollen, hat sie das nur zu dem Ende gethan, daß sie möchte Mutter werden und das Haus Jakobs mehren, welches die Verheißung hatte. Und ist ein Exempel einer sehr heiligen Matrone, die im Glauben und Geduld GOtt um Gnade angerufen, auf den Trost gewartet, dazu sehr schwere Anfechtung und Schmerzen erlitten hat.

111. Darum hat Mose ein sehr nachdrücklich Wort gebraucht, da er sagt: „Der HErr gedachte an Rahel“ 2c., als wollte er sagen: Sie hatte bei sich selbst schier verzagt, und in ihrem Herzen also gedacht und beschlossen, GOtt würde an sie nimmermehr denken, ja, er hätte ihrer ewiglich vergessen. Ich werde keine Mutter werden, hat sie gedacht, sondern ich bin die Elendeste unter allen Weibern: ich sollte die rechte Hausmutter sein im Hause, aber GOtt hat meiner vergessen. Also ist sie nun in die Hölle geführt worden. Da sich es ansehen läßt, daß keine Hoffnung der Hülfe mehr vorhanden sei, ergreift sie fast aus Verzweiflung die Magd und gibt dieselbige dem Manne zum Weibe, das sie sonst nicht würde gethan haben, wo sie nicht alle Hoffnung hätte fahren lassen. Sie verzagt aber also, daß sie gleichwohl noch ein Fünklein von Glauben behalten hat: sie behält noch in der Verzweiflung das Seufzen, welches St. Paulus Röm. 8, 26. „unaussprechliche Seufzen“ nennt. Aber dasselbige ist so gar tief begraben und überschüttet mit der Unmöglichkeit und andern widerwärtigen Bewegungen, so sie in ihrem Herzen gefühlt hat, daß sie solch tief Seufzen kaum empfunden hat.

\quellseite{157}

112. Wie Hanna, die Mutter Samuels, auch daran verzagt hatte, daß sie würde Kinder kriegen, und konnte auch ihr Seufzen nicht fühlen, und das große Verlangen, so sie inwendig in ihrem Herzen gehabt, daß sie möchte Kinder haben, 1 Sam. 1, 10. Aber Gott, der die Herzen erforscht, versteht das unaussprechliche Seufzen wohl, welches sonst mit keinerlei Sinnen begriffen und mit keinen Worten ausgedrückt werden kann.

113. Also erzählt Augustinus von seiner Mutter Monica, die ganze neun Jahre darüber geheult und geweinet hat, daß ihr Sohn sollte verloren werden, darum daß er in die Secte und Irrthum der Manichäer gerathen war. Sie hat aber nur von Gott, daß ihr Sohn möchte bekehrt und ein Christ werden, und wollte ihm um der Ursache willen ein Weib, die eine Christin war, ehelich vertrauen, auf daß er durch Weiberwohnung und Gemeinschaft mit einem solchen gottseeligen Weibe und durch ihre Ermahnung wiederum auf den rechten Weg möchte gebracht werden. Und sie hätte sich daran genügen lassen, wenn er nur von derselben Ketzerei der Manichäer zur reinen gesunden Lehre wäre bekehrt worden. Es ließ sich aber ansehen, wie Augustinus sagt, als hätte ihr der Herr solches abgeschlagen und zu geben verweigert: er hatte aber doch dasjenige, daran es gelegen war, und ihr tiefes Verlangen gehört, nämlich, das unaussprechliche Seufzen.

114. Wie aber das Gebet und Seufzen unaussprechlich ist, also ist auch die Erhörung und Freude unermeßlich und unaussprechlich; wie St. Paulus sagt Eph. 3, 20.: „Dem aber, der überschwenglich thun kann über alles, das wir bitten oder verstehen“ etc. Er gibt das nicht, was seine Heiligen oben im Herzen und mit dem Schaume der Worte bitten und begehren, sondern er ist ein allmächtiger Geber und der da sehr reich ist, der da gibt nach dem tiefen Seufzen. Und darum läßt er das Gebet immer heftiger wachsen und gemehret werden, und erhört es nicht so bald. Denn wo er zum ersten Geschrei oder Bitte antworten würde, so würde das Gebet nicht gemehret, sondern würde kalt werden. Darum verzieht er mit der Hülfe, und dann macht das Gebet täglich und wird stark und groß, desgleichen wird auch das Seufzen des Herzens tiefer und brünstiger, bis es gleichsam zur Verzweiflung kommt: da wird dann das Gebet sehr brünstig und stark, wenn es sich läßt ansehen, daß jetzt das Seufzen schier aufhören will.

115. Wenn er uns nun so bald erhörete, so würde das Gebet so stark, so lebendig und brünstig nicht, sondern es wäre nur obenhin ein äußerlich Seufzen, welches noch beten, seufzen und begehren lernt, dermalen noch kein Meister des Gebetes ist: wenn es aber zur Verzweiflung gekommen ist, und das Herz, so angesehen wird, also denkt: Ach, es wird nichts daraus, es ist alles verloren! und doch gleichwohl noch ein Fünklein und das glimmende Docht vorhanden, alsdann gedenke, und sei stark und halte fest. Denn das ist der Kampf, den die Heiligen führen müssen, die da meinen, daß der Strick jetzt alsbald brechen werde, und doch noch das herzliche Seufzen behalten. Darum ist alsdann das Gebet vollkommen und stark.

116. Also wird hier von Rahel auch gesagt: „Der Herr gedachte an Rahel“ etc. Als wollte Mose sagen: Rahel hat Anderes nichts gesucht, denn daß alle ihr Gebet und Weinen vergeblich und umsonst gewesen wäre, und daß sie Gott in seinem Herzen gar ausgetilgt hätte und an sie nicht mehr denken würde. Aber, meine liebe Rahel, du sollst es dafür nicht halten, denn du fehlest weit: dieses dein Fühlen ist fleischlich, es ist noch nicht der Geist, der da seufzt, und das unaussprechliche Seufzen, sondern es ist das Fleisch, das solches fühlt und schwach ist. Gott hatte dein noch niemals vergessen, sondern alsbald im Anfang, da du angefangen hast um Kinder zu bitten, hat er alle Worte deines Seufzens gehört und aufgezeichnet; das Gebet war aber damals noch nicht brünstig und stark genug, darum hat es müssen wachsen und groß werden, und darum hat Gott auch die Erhörung verzogen bis in das andere, dritte, vierte und fünfte Jahr; da hat er dir zuerst durch die Magd einen Trost gegeben. Das Seufzen hat aber noch nicht aufgehört: es hat sich noch lassen ansehen, als hätte sich Gott von dir gewendet. Darum hat nun Gott an dich gedacht, sagt Mose, wiewohl er dein keinen Augenblick vergessen hat. Aber dieses ist endlich die rechte Zeit gewesen, daß Gott dein Seufzen gehört hat, da du dich hast dünken lassen, daß dein Seufzen gar begraben, zugescharrt und vergessen wäre.

117. Auf dieselbige Weise sollen wir auch lernen beten und der Hülfe von Gott warten, wenn Widerwärtigkeit vorhanden ist und der Glaube anfängt zu sinken. Denn wir haben die Verheißung des Evangeliums; wir haben die Taufe und die Absolution, dadurch wir unterrichtet und gestärkt sind; wir haben das Gebot, darin uns befohlen wird, daß wir beten sollen; wir haben den Geist der Gnade und des Gebetes: wenn wir aber erst angefangen haben zu beten, so ist unser Herz in Angst, und klagt, daß es mit dem Gebete nichts ausrichte.

118. Darum muß man lernen, wenn du mit Beten nichts ausrichtest, daß du noch dazu forschest, das ist, suchest; und wenn dich dünkt, daß solch Suchen auch noch vergeblich sei, und Gott sich immer mehr und mehr verbirgt, so fange an, und klopfe, und höre nicht auf, bis daß du die Thüre auflaufest, darin er verschlossen ist. Denn daran ist kein Zweifel, daß unser Gebet nicht sollte alsbald erhört werden, wenn die erste Sylbe ausgeredet ist; wie der Engel zu Daniel sagt, Dan. 9, 23.: „Da du anfingest zu beten, ging dieser Befehl aus, und ich komme darum, daß ich dirs anzeige“ etc.

119. Daß uns aber Gott nicht alsobald gibt, was wir bitten, geschieht darum, weil er will gesucht sein, und will haben, daß man immer anhalten und ihn also heraus klopfen soll; wie das Gleichniß lehrt Luc. 18, 2. ff. von dem ungerechten Richter. Alsbald kommt er und errettet seine Auserwählten, und gibt reichlicher und mehr, denn wir gebeten, gesucht und geklopft haben. Darum verzieht er aber mit der Erhörung, daß das Gebet gemehret und das Seufzen stärker werde, welches uns noch sehr geringe zu sein dünkt, dieweil wir seufzen: es ist aber in Wahrheit sehr hitzig, wie es Sanct Paulus ein Geschrei oder Rufen nennt. Denn wir sprechen nicht nur die Worte aus, indem wir durch die Zunge und Lippen einen Laut hervorbringen; auch lassen wir nicht mit lauter Stimme unsere Gebete erschallen: sondern wir schreien geradezu. Es ist kein Laut oder Stimme des Mundes, sondern ein Rufen im Herzen und ein unaussprechlich Seufzen: unter den linken Zitzen steckt es, wenn das Herz klopft und seufzt, und schier vor großer Angst gar matt wird; und alsdann ist das Gebet recht vollkommen und kräftig.

120. Davon soll man oft sagen und soll es immer wiederholen, auf daß wir nicht alle Hoffnung und Vertrauen von unserm Gebet dahinfallen lassen. Welches Gebet, ob es wohl anfänglich kalt ist und die Hülfe nicht alsbald erlangt, so sollen wir doch wissen, daß die Hülfe darum verzogen wird, auf daß das Gebet desto vollkommener und stärker werde. Denn das Gebet hat eine wunderliche Kraft und Allmächtigkeit. Als, da Rahel ist angesehen worden, als wäre sie von Gott gar verachtet und verlassen, und sie doch noch ein glimmend Docht geblieben, da hat Gott, als er angerufen, gesucht und hervor geklopft worden, müssen aufmachen. Eben aber auf dieselbige Weise hilft er auch allen, die ihn anrufen, und hilft ihnen so reichlich und mildiglich, daß sie bekennen müssen, daß sie solches nie hätten hoffen dürfen. Wie Monica, St. Augustins Mutter, hat sagen mögen: „Ich habe darum nicht gebeten, habe deß auch nicht hoffen dürfen, daß mein Sohn jemals ein solcher großer Lehrer in der Kirche werden sollte“ etc. Darum sollen wir nicht matt werden oder verzagen, sondern im Gebet und mit Bitten und Suchen fortfahren und beharren, bis die Hoffnung und Errettung, darauf wir gewartet haben, sich sehen läßt.

121. Rahel aber zeigt selbst genugsam an, was ihr Kreuz und Anfechtung gewesen sei; denn sie sagt: „Gott hat meine Schmach von mir genommen.“ Im lateinischen Text steht: Collegit opprobrium meum: Er hat es zusammen gerafft. Sie hat fünf oder sechß Jahr nach einander die Schmach und Schande getragen, darum daß sie unfruchtbar war, und sie selbst und andere Leute auch haben sich lassen dünken, sie wäre das elendeste Weib unter allen Weibern. Denn sie hat gesehen, daß ihre Schwester fruchtbar war; sie hat gesehen, daß ihre und ihrer Schwester Magd, dazu alle anderen Weiber auch geboren haben, und daß sie derhalben gelobt, geliebt und in Ehren gehalten würden: ich aber allein, hat sie gedacht, werde für ein verworfen, verachtet und verfluchtes Weib gehalten; der Herr aber hat diese Schmach und Schande alle gesehen, die ich vielfältig und schwer getragen habe, und dieselben hat er zusammen gerafft und von mir genommen; nun bin ich ein rechtes Eheweib und Hausmutter geworden.

\quellseite{158}

122. Also ist ihre Begierde und Verlangen erfüllt und ihr Seufzen erhört worden, da es zuvor durch langen Verzug so vieler Jahre bewährt und genähret ist. Es erhebt auch aus diesem Text, dieweil Rahel sagt, der Herr habe ihre Schmach von ihr genommen, daß zu derselben Zeit gemeiniglich von jedermann dafür gehalten worden ist, daß an einem Weibe die Unfruchtbarkeit eine Schande wäre, vornehmlich in der Linie oder Geschlecht, daraus man die Erben, so verheißen waren, erwartet hat. Im Neuen Testament ist es für keine Schmach gehalten worden; denn da wird die Jungfrauschaft sehr gelobt, vornehmlich im Kirchenamt. Aber bei den Vätern im Alten Testament ist es ein sehr schwer Kreuz und eine große Schmach gewesen. Und solch Urtheil ist hergekommen aus Gottes Wort, 1 Mos. 1, 27. 28.: „Gott schuf den Menschen, ein Männlein und Fräulein, und segnete sie, und sprach zu ihnen: Seid fruchtbar und mehret euch, und füllet die Erde“. Diese Worte hat Adam bei den Nachkommen ausgebreitet und ihnen dasselbe fleißig eingeschärft: darum ist es also immer von Einem zum Andern auf ihre Kinder und Kindesfinder gekommen, daß sie es dafür gehalten haben, daß ein jegliches unfruchtbares Weib von Gott verworfen und verflucht wäre. Das ist das Allerschwerste gewesen, daß derselbe Text zugleich auch den Fluch Gottes in sich begreift. Es ist zwar die Unfruchtbarkeit an sich selbst eine schwere Last und ein großer Jammer; sie wird aber noch viel schwerer, wenn das dazu kommt, daß ein unfruchtbar Weib verflucht sei und Gott nicht gefallen könne, oder wenn sie ihm schon gefalle, daß sie doch gleichwohl unter den andern Weibern im Volke Gottes verachtet und verschmäht hat sein müssen.

123. Derhalben hat Elisabeth, die Mutter Johannes des Täufers, fast mit denselbigen Worten Gott auch gedacht, da sie sagt Luc. 1, 25.: „Also hat mir der Herr gethan in den Tagen, da er mich angesehen hat, daß er meine Schmach unter den Menschen von mir nähme“; und der Engel spricht zu Maria Luc. 1, 36.: „Elisabeth, die im Geschrei ist, daß sie unfruchtbar sei.“ Als wollte er sagen: Die derhalben ein böses Geschrei hat und sehr geschmäht wird, daß sie nicht fruchtbar ist. Derhalben ist zu der Zeit die Unfruchtbarkeit nicht allein eine sonderliche Plage und Kreuz gewesen unter den Leuten, sondern hat auch Gottes Fluch auf sich liegen gehabt.

124. Heutiges Tages findet man ihrer Viele, die Gutes genug haben und von großem Geschlechte sind, die darüber ungeduldig werden und für untauglich nehmen, wenn sie dieses Segens beraubt sind, und haben sehr großes Verlangen nach Kindern, dieweil sie Gutes genug haben, daß sie die Kinder wohl nähren könnten: es ist ihnen aber dieses Kreuz darum noch etwas erträglicher, dieweil im Neuen Testament das nicht dazu kommt, daß sie derhalben sollten von Gott verworfen und verflucht sein.

125. Der andere gemeine Haufe des Volkes ist auch gottloser, als die Heiden gewesen sind. Denn der größte Theil unter den Eheleuten begehrt keine Kinder zu haben; ja, daß noch mehr ist, sie sind den Kindern feind, und meinen, es sei besser ohne Kinder leben, darum daß sie arm sind, und nicht haben, damit sie ihr Gesinde unterhalten möchten; vornehmlich aber die, so zum Müßiggange geneigt und faule Schelme sind, die vor der Mühe und Arbeit des Ehestandes fliehen. Nun ist aber gleichwohl der Ehestand zu dem Ende nicht eingesetzt, daß man darin immer sich ergötzen und müßig sein sollte, sondern man soll darin Kinder zeugen und auferziehen, und das Gesinde ernähren; welches wahrlich eine unermäßliche Last ist, voll großer Sorge, Mühe und Arbeit. Darum bist du von Gott geschaffen, daß du ein Ehemann oder ein Eheweib sein sollt, und daß du lernest, solche Mühseligkeit dulden und ertragen. Die zu den Kindern nicht Lust haben, sind Säue, Böcke und grobe Klötze, sind nicht werth, daß sie Männer oder Weiber heißen sollen, weil sie den Segen Gottes, der den Ehestand geschaffen und eingesetzt hat, verachten.

\section[Vierter Theil]{Vierter Theil. Wie Jakob Laban um seinen Lohn bittet, und um Erlaubniß anhält fortzuziehen; wie ihn Laban sucht bei sich zu behalten; und wie sich Jakob dabei verhält}

\subsection{I}

\begin{quote}
\textbf{V. 25. 26. Da nun Rahel den Joseph geboren hatte, sprach Jakob zu Laban: Laß mich ziehen und reisen an meinen Ort, und in mein Land. Gib mir meine Weiber, und meine Kinder, darumb ich dir gedienet habe, daß ich ziehe; denn du weißest meinen Dienst, wie ich dir gedienet habe.}
\end{quote}

126. Siehe doch, wie züchtig und demüthig der Patriarch um seine Entlassung bittet, und fordert gar nichts mehr, denn seine Weiber und seine Kinder, und zeigt doch daneben an, daß er dem Laban treulich und wohl gedienet habe, und daß er wohl werth wäre, daß er ihm etwas sollte zum Lohne geben. Denn darum sagt er: „Du weißest meinen Dienst, wie ich dir gedienet habe.“ Als wollte er sagen: Es ist billig, daß du denken sollst, daß ich dir nicht allein um Weib und Kind gedienet habe, sondern sollst auch an die Aussteuer, oder an den Lohn, oder aber an irgend ein Geschenk denken, damit du mich, ja, vielmehr deine Töchter und Kindeskinder zieren möchtest für meinen treuen und willigen Dienst. Dasselbige hat ihm Jakob mit diesen Worten sein züchtig und heimlich zu verstehen geben wollen.

127. Warum thut er aber das? Antwort: Darum daß Laban ein Unflat und Geizhals war; er war ein Hund etc., er war gleichsam ein Ungeheuer derselbigen seligen Zeit. Denn er sollte billiger wegen seiner großen Bosheit in dieser letzten bösen Zeit gelebt haben, da Alles durch den Geiz und sonst alle andere Sünde und Schande verderbet ist. Derhalben ist dies nun fleißig zu merken um der That Jakobs willen, so hernach folgt, die wir der Ursache halben entschuldigen werden. Denn der gute fromme Mann hat ganze vierzehn Jahre gedient in der äußersten Armuth, dergleichen Dienst, so mit solcher Armuth verknüpft ist, heutiges Tages niemand dulden würde oder ertragen könnte. Denn er hat die ersten sieben Jahre um Rahel gedient; er ist aber seiner Hoffnung ganz unbilliger Weise beraubt worden und hat von Laban gar keinen Lohn genommen. Denn Laban hat also gedacht: Er ist fremd und elend, derhalben will ich ihn mir dienen lassen, so lange als ich kann; denn er dient mir treulich, und mehret und bessert mir mein Hab und Gut. Darum mögen wir denken, wie viel Jammer und Elend der gute fromme Jakob müsse gelitten haben.

128. Denn Laban (der billig Nabal heißen sollte, daß man das Wort umkehrt) hat die Treue und Fleiß des frommen Mannes nicht erkannt, und gedenkt, ihn jetzt von sich ziehen zu lassen ohne Lohn. Denn zuvor hat er zwar kein eigen Gut oder gar wenig eigenen Gutes erworben; er hat nichts gehabt, denn nur Nahrung und Kleidung; und damit ist es noch schmal und kümmerlich genug zugegangen. Der fromme Mann hat einem rechten Hunde gedient. Und hernach werden die Töchter selbst auch klagen über den Geiz und die Räuberei ihres Vaters, dadurch Rahel ist bewogen und verursacht worden, daß sie dem Vater seine Götzen gestohlen hat; denn er hat sie nicht als seine Töchter, sondern wie Mägde gehalten, deren Arbeit, und was sie mit Wollwerk und sonst anderer Arbeit mehr verdient, er zu seinem Genuß und Nutzen gewendet hat. Da er aber nun mit ihm um den Lohn einen Vertrag machen will, versucht er ihn abermal zu betrügen, auf daß er ja seinem Schwiegersohn und Töchtern für ihren treuen Dienst nichts geben dürfe, dadurch er doch war reich geworden.

129. Diese Unbilligkeit und Gewalt hat Jakob selbst und auch seine Weiber sehr betrübt, und sie werden zum öftern heimlich über die Armuth und Noth geklagt haben, dazu sie durch den Betrug und Räuberei ihres Vaters gebracht worden. Denn dieses zeigen die Umstände, wo man dieselben recht erwägt, klärlich genug an. Und es ist an Laban eine verfluchte Bosheit, die gar wider alle Sinne und Natur des Menschen ist. Denn sonst pflegen ja die Väter ihre Töchter etwas lieber zu haben, weil das weibliche Geschlecht auch schwächer ist. Aber dieser unfreundliche und greuliche Nabal beraubt auch noch darüber seine Töchter und plagt sie mit hartem schwerem Dienst; da aber vom Lohne geredet und gehandelt wird, gedenkt er abermals, wie er Jakob aufs neue hintergehen und betrügen möge. Jakob thut alle seine Dinge treulich und aufrichtig: aber der Betrüger geht nur mit Räuberei und Ränken um; wie sich Jakob hernach darüber beklagt, da er sagt: „Er hat nun zehenmal meinen Lohn verändert“ etc.

\quellseite{159}

130. Darum kann man die Lieblosigkeit, Grausamkeit und den Geiz Labans nicht groß genug machen; und man soll sie darum an den Tag bringen und auf jede Weise groß machen, weil dadurch die Ursache angezeigt wird der That Jakobs, so hernach wird erzählt werden, da er Laban wiederum auch betrügt; auf daß man solches nicht für unrecht und schelmisch halte, wie es sich ansehen läßt, wo man es nur obenhin ansieht. Es ist aber ein nöthiger Betrug gewesen um der unbilligen Räuberei willen seines Schwiegervaters.

131. Jetzt aber bittet er um Erlaubniß, daß er wieder möge an seinen Ort ziehen. Ich habe dir, sagt er, vierzehn Jahre gedienet: die ersten sieben Jahre umsonst; und da du mich mit List überwunden hast, habe ich dir noch andere sieben Jahre dazu gedienet, und dieselbigen sind nun um; so ist es nun Zeit, daß ich auch mich selbst, mein Gesinde und mein Haus versorgen muß. Du sollst aber bedenken, wie treulich ich und meine Weiber dir gedienet haben, das du doch jetzt gar nicht achtest; denn du hast mir gar keinen Lohn dafür gegeben, daß ich dir die vierzehn Jahre gedienet habe. Darum ist es billig und recht, daß ich meinen Abschied nehme, und auf mein eigen Ding auch nun sehen und dafür sorgen möge. „Gib mir meine Weiber und meine Kinder.“ Als wollte er sagen: Wo du mir ja nichts geben wirst an Geld oder sonst meinem gebührlichen Lohn, so wird es doch vermalen kommen, daß Gott darüber richten, und dich um solcher Undankbarkeit willen strafen wird.

\subsection{II}

\begin{quote}
\textbf{V. 27. Laban sprach zu ihm: Laß mich Gnade vor deinen Augen finden. Ich spüre, daß mich der Herr segnet um deinetwillen.}
\end{quote}

132. Du heiliger Pabst, du frommer Bischoff von Mainz, kannst du auch von Gott reden? Laban führt den Namen und den Segen Gottes an. Aber wenn du die Wahrheit sagst und sprichst, daß du durch Gottes Segen um Jakobs willen an deiner Habe hast zugenommen und reich geworden bist, warum dankst du ihm nicht dafür, sondern gedenkest ihn noch zu bearbeiten, zu übervortheilen, auf daß er dir noch immer länger dienen möge? Ja, also pflegen die Heuchler und alle Geizhälse und Gottlosen zu thun: sie führen den Namen des Herrn, aber unmöglich. Darum braucht Mose ein wunderlich Wort, auf Hebräisch nichashti, welches den Namen hat von einer Schlange, auf welche Bedeutung er auch hat sehen wollen, wie es sich ansehen läßt. Denn es heißt auch Wahrsagung und Zauberei; wie 4 Mos. 24, 1. steht: „Er ging nicht hin nach den Zauberern.“

133. Es ist aber ein Gebrauch gewesen, daß sie haben mit den Schlangen gehandhabet und Zauberei getrieben; darum hat Mose wollen anzeigen, daß Laban ein Wahrsager gewesen sei, denn er sagt: Auguratus sum, das ist: Ich habe mich es wohl lassen dünken; es ahnt mich wohl; mein Herz hat es mir gesagt; ich habe es wie an einem sonderlichen Wahrzeichen abgenommen, daß mich Gott um deinetwillen gesegnet hat. Aber er will Jakob die Ehre nicht geben, will ihm auch dafür nicht danken, und gibt auch Gott die Ehre nicht. Also hat mich gedünkt, sagt er, daß mich der Herr durch dich gesegnet habe. Er will für einen solchen Mann angesehen sein, der dessen wohl werth sei, daß er von Gott gesegnet und reich gemacht werde, wiewohl er, ehe Jakob zu ihm gekommen ist, fast arm war.

134. Das hat Mose mit dem wunderlichen und nachdrücklichen Worte anzeigen wollen. Also sind aber die Leute gesinnet, die da falsch, zweizüngig und betrüglich sind, die andere reden, denn sie es im Herzen haben; wie die Gestalt ist der Zunge einer untreuen listigen Schlange. Er sagt nicht sein einfältig: Ich erfahre und verstehe es wohl, ich halte es auch gewißlich dafür, daß ich um deinetwillen gesegnet und reich geworden bin. Denn wo er das erkannt hätte, so würde er den Beutel haben müssen aufthun, und seinem Schwiegersohn, seinen Töchtern und seinen Kindeskindern eine Summe Geldes geben. Er hält es aber nicht gewiß dafür, sondern zweifelt noch daran und disputirt davon, wo Jakob von ihm wolle Lohn fordern, daß er Ursache haben möge, sich mit irgend einem Schein zu weigern. Darum sagt er: Ich lasse es mich wohl bedünken, es kann wohl geschehen sein, daß mich der Herr um deinetwillen gesegnet hat. Das ist der Künste eine, so die Fürsten wissen, nämlich, daß sie solche Worte gebrauchen, die keinen gewissen Verstand haben, die auf beiderlei Seiten nach ihrem Gefallen, und wie sie es gelüstet, gebeugt und gedeutet können werden.

135. Hernach aber werden wir hören, daß Laban auch ein Götzendiener gewesen ist. Derselben ist er vielleicht mit Wahrsagen umgegangen oder hat die Götzen um Rath gefragt, und hat es bei sich selbst dafür gehalten, daß er solches Glück davon habe, weil er den Götzen gedienet hatte. Davon will ich aber nicht disputiren. Das ist aber gewiß, daß er den Namen des Herrn unnützlich führt. Er trachtet darnach, äußerlich sich nach dem Geist und der Art Jakobs zu richten: er sieht, daß derselbe göttlich und fromm ist, darum braucht er eben dieselbigen Worte, welche Jakob auch pflegte zu gebrauchen.

\begin{quote}
\textbf{V. 28. (Und Laban sprach:) Stimme den Lohn, den ich dir geben soll.}
\end{quote}

136. Das hebräische Wort nakab heißt, durchbohren; wie im Propheten Haggai steht am 1. Capitel V. 6.: „Welcher Geld verdienet, der legt es in einen, nakub, löcherichten Beutel“ etc. Es heißt auch ernennen, bestimmen: in welcher Bedeutung es gesetzt wird 3 Mos. 24, 16.: „Welcher des Herrn Namen nennet, der soll des Todes sterben“, das ist, wer in Gottes Namen flucht etc. Also haben wir es an diesem Orte auch verdolmetscht: Stimme oder nenne mir einen Lohn. Es hätte Mose wohl ein ander Wort gebrauchen können; er will aber damit zu verstehen geben, daß Laban mit einem sonderlichen und geistlichen Worte habe verhüten wollen, daß er nicht zu viel Lohn fordern sollte, als ob er damit Gott erzürnen würde, wenn er einen unbilligen Lohn fordern wollte. Denn alle Worte, so Laban redet, kommen aus einem solchen Herzen, das vom Geize ganz verderbet ist, und gehen nur dahin, daß er Jakob übervortheilen und berauben möge. Er setzt aber noch dazu: So will ich dir es geben. Das ist nicht wahr, er gedenkt ihm nicht einmal einen Heller zu geben; wie er hernach einen neuen Betrug wird vornehmen, auf daß er den Bund und Vertrag, was er ihm hat als Lohn geben sollen, nicht halten müsse. Da werden wir sehen, daß er nicht allein geizig und gottlos, sondern der Geiz selbst gewesen ist.

\subsection{III}

\begin{quote}
\textbf{V. 29. 30. Er aber sprach zu ihm: Du weißt, wie ich dir gedienet habe, und was du für Vieh hattest unter mir. Du hattest wenig, ehe ich herkam; nun aber ist's ausgebreitet in die Menge, und der Herr hat dich gesegnet durch meinen Fuß.}
\end{quote}

137. Jakob wiederholt seine vorigen Worte mit einer harten Strafe des Geizes, des gottlosen Wesens und des Wahrsagens, damit Laban umging. Als wollte er sagen: Du darfst dich nicht dünken lassen oder daran zweifeln, ob du diesen Segen mir oder dir selbst zuschreiben; es beweist sich zwar an sich selber genugsam, daß dich der Herr um meinetwillen gesegnet hat; denn du bist zuvor auch mit Wahrsagen umgegangen; wo du aber damit etwas erlangt hättest, so wäre solches auch wohl gekommen, da ich nicht bei dir gewesen bin. Aber es beweist sich an sich selber, daß dein Hab und Gut alsdann erst ist gemehret worden, da ich mich deines Hauses angenommen und meinen Dienst darin verwaltet habe.

138. Darum widerleget Jakob hiermit die gottlose und ungeschickte Antwort Labans von seinem Wahrsagen oder Götzdünken, und rühmt seine Arbeit, so er in seinen Arm gethan hat. Du weißt, sagt er, wie fleißig und treulich ich dir gedienet habe, welches auch das Werk an sich genugsam ausweist, wenn du schon daselbst leugnen wolltest und ich desselben schweigen. Denn was ist es für Vieh gewesen, so mir erstlich befohlen worden ist? Vielleicht mögen es etwa tausend Schafe gewesen sein, die Rahel, deine Tochter, allein weidete, also daß du weder Knecht noch Hirten dazu bedurft hast; so gar wenig war derselben. Was darfst du denn auguriren (wahrsagen), du geiziger Wanst? Denn also hätte Jakob diesen Unflat und Geizhals billiger und schärfer strafen mögen. Es kann zwar die Heerde der Schafe nicht groß sein, die ein junges Mädchen weidet, aber dasselbe Wenige ist nun gemehret oder aufgebrochen; denn so viel bedeutet das hebräische Wort. Als wollte Jakob sagen: Da ich dein Vieh unter mir gehabt, und dasselbe gemehret und bewahrt habe, ist es gleichsam mit Gewalt gemehret und viel geworden, daß es auch vor großer Menge aus den Ställen herausgebrochen ist; derhalben hat sich Gott gesegnet um deinetwillen, da ich bin zu dir gekommen, und nicht, da du deines Wahrsagens gewartet und den Söhnen gedient hast.

\quellseite{160}

139. Diese Widerlegung und Strafe ist scharf genug, damit er diesen sehr geizigen Menschen strafet, der sich beschwert, seinem getreuen und wohlverdienten Knechte seinen verdienten Lohn zu geben, der ihm nicht Ein Jahr, sondern vierzehn ganze Jahre gedient hatte; in welchen Jahren er nichts ausgegeben, sondern Alles an sich gerafft und gespart, dazu seinen Schwiegersohn und seine Töchter beraubt, und sie gar nichts für ihre sehr schwere Arbeit hat behalten lassen. Es wird aber bald hernach ein Exempel göttlichen Gerichtes folgen wider solche greuliche Räuberei. Denn „wer da kärglich säet, der wird auch kärglich ernten“, 2 Cor. 9, 6.

140. Und ist ein solch Exempel, das sich sehr wohl reimet auf unser Leben, so wir jetzt führen. Denn wir sehen, wie die Fürsten, der Adel, Bauern und Bürger, alle miteinander wider den armen Jakob handeln, wie Laban gethan hat. Es wird ihnen aber allen eben also gehen, wie es hier auch gegangen ist. Denn wie heftig sie an sich reißen, scharren und zusammenraffen; so segnet doch GOtt den Jakob, der da beraubt wird, und dem, der Alles geraubt und genommen hat, nimmt er wiederum Alles, wie im 39. Psalm V. 7. steht: „Sie sammeln, und wissen nicht, wer es kriegen wird.“

141. Und ist wahrlich ein sehr großer Jammer, daß die Welt sich dessen nicht will überreden lassen, sondern sie will GOttes Gerichte fühlen und erfahren. Wenn die Leute nur nach großem Gut trachten, sammeln und an sich reißen mit Recht und Unrecht, wie sie nur können und mögen, alsdann halten sie sich für selige Leute, und meinen, sie haben nun das Hinweg, danach sie mit so großer Begierde gestanden haben. Aber der Psalm antwortet ihnen darauf und sagt: „Sie wissen nicht, wer es kriegen wird.“

142. Also sind die Städte und Fürsten, die groß Geld zusammen bringen und gar unsinnig sind mit Wuchern und Geizen. Wem sammeln sie es aber? Antwort: Bruder Zeit, dem gottlosen Kriegsvolk, welches bald hernach Alles hinreißen und unter die Füße treten wird; wie uns dasselbe neulich auch schier also begegnet wäre, im Jahr 42. um die Ostern, in dem Lärm, so sich zwischen den Herzogen von Sachsen erhoben hatte bei der Stadt Murzen. Wo es GOtt nicht gewehret und abgewendet hätte, wie bald wäre Bruder Zeit gekommen und hätte in Einem Tage Alles zerrissen!

143. Pabst Clemens hat den Tempel zu Toledo geplündert, und hat daraus genommen eine überaus große Summe Geldes, die fast dreier Königreiche werth gewesen ist. Das hat er unter dem Schein gethan, daß er dasselbe Geld der St. Petri-Kirche einverleiben wollte, denn er ist ein köstlicher Meister darin gewesen, wie man rauben und an sich scharren soll, und ließ sich ansehen, er könne mit der ganzen Welt nicht gesättigt werden. Aber er hat es auch Bruder Zeit gesammelt, den gottlosen Kriegsgurgeln, die es darnach alles zerrissen und geplündert haben.

144. Eben so wird es den Andern allen auch gehen, die jetzt mit Wucher und unrechtem Gut reich werden und lassen es mit ihrem Scharren und Geizen kein Maß oder Ende haben. Und wo sie etwa ermahnt werden, daß sie den Armen auch etwas lassen sollen, oder daß sie doch mit GOttes Segen etwas sammeln, so hören sie es nicht, sondern denken also, wie der Poet sagt: Quaerenda pecunia primum, virtus post nummos etc.: Ich muß erst nach Geld und Gut trachten, fromm werden kommt wohl hernach; wenn ich irgend eine gute Summe Geldes zusammen gebracht habe, darnach will ich derselben gebrauchen, und will sie dennoch wohl behalten, ich raffe sie gleich zusammen und sammle sie mit Recht oder mit Unrecht. Das wird aber mit nichten geschehen; denn also sagt der Psalm: „Sie sammeln, und wissen nicht, wer es kriegen wird.“

145. Wie viel besser und seliger wäre aber das, wo man GOtt fürchtete und mit göttlichem Segen reich würde? wie der Spruch lautet im 37. Psalm V. 16.: „Das Wenige, das ein Gerechter hat, ist besser, denn das große Gut vieler Gottlosen.“ Aber da sind die Leute gar taub. Denn sie müssen also sammeln und zusammenscharren, daß sie einen pariz haben, wie es auf Hebräisch lautet, der es wieder hinreiße. Darum mögen wir uns nach dem gemeinen Spruch richten: Mitte vadere, sicut vadit etc.: Laß es gehen, wie es geht, die Welt will sich nicht sagen lassen, sie will es erfahren. Dieweil Laban Alles sammeln und an sich reißen will und seinen Schwiegersohn und seine Töchter ihres verdienten Lohnes beraubt, so verliert er Alles; wie man sagt: Wer zu viel haben will, dem wird nichts.

146. Ich habe selbst, der ich nun schon ein Sechzigjähriger bin, viel Exempel solcher Labansbrüder gesehen, welcher großes Gut bald in einem Hui zerstreut und umgekommen ist. Ich könnte auch wohl Exempel namhaft machen, die sich in der kleinen Stadt Mansfeld, welches meine Vaterstadt ist, zugetragen haben. Es wird auch in andern Städten, als zu Leipzig, zu Freiberg, und sonst an vielen andern Orten, da man Bergwerke hat, von merkwürdigen Exempeln solcher gesagt, die auch seit unserem Gedenken reichlich an sich gerissen und gesammelt haben; aber jetzt wäre nicht ein Heller mehr übrig von all ihrem großen Geld und Gut. Ja, das heißt, in einen löcherichten Beutel sammeln, wie der Prophet Haggai sagt am 1. Capitel, V. 6. Warum geschieht aber das? Ei, sagt der HErr, darum, daß ich darein geblasen habe. Ja, wenn der darein bläset, so muß das Gut zerstäuben, wenn auch die ganze Welt voll Goldes und Silbers wäre.

147. Wiewohl wir nun damit nichts ausrichten bei den Heuchlern und gottlosen Leuten, so muß man es doch gleichwohl dem Volk einschärfen, ob vielleicht Etliche wären, die dadurch möchten bewogen werden. Laban hat darnach nichts gefragt; Jakob aber hat nichtsdestoweniger seine Weiber und Gesinde solches gelehrt, ja, er hat Laban selbst auch des Zornes GOttes erinnert wider das gottlose Wesen und wider den Geiz. Sehet, wird er zu seinem Gesinde gesagt haben, wie mein Schwiegervater so überaus geizig ist, der weder GOtt fürchtet noch auch an unsern so treuen Dienst denkt; denn er gibt uns nicht allein keinen Lohn, sondern beraubt und plündert seine eigenen Töchter und Kindeskinder. Was soll ich thun? wenn ich schon den Lohn nennen und bestimmen werde, so wird er doch eine andere List finden, damit er uns übervortheilen und betrügen möge. Denn solche Leute haben nur ihre Lust daran, daß sie Andere berauben, übervortheilen und betrügen mögen; sie erdenken nichts denn eitel List und gehen nur immer mit Betrug um. Und wenn sie Einem einen Heller abschinden können, so meinen sie, sie haben einen Gulden genommen. Es ist eine verfluchte Art Menschen, denen beide GOtt und Menschen feind sind.

148. Das letzte Stück im Text verholmetschen sie im Lateinischen alle: ad ingressum sive introitum meum, das ist: Da ich erst in dein Haus gekommen bin. Aber das Wesen an sich weist es aus, daß Laban nicht alsbald reich geworden sei, da Jakob zu ihm in sein Haus gekommen ist. Darum wird es besser also verdolmetscht: ad pedem vel ad cursum meum; damit zugleich angezeigt wird der große Fleiß und die schwere Arbeit, so Jakob gehabt hat. Als wollte er sagen: Ich habe Tag und Nacht gelaufen, habe nicht geritten noch auf feinem Wagen gefahren, sondern ich bin gar mühe und matt geworden, da ich auf das Feld, auf die Weide zu dem Vieh und heim in das Haus gelaufen, auf daß ich dir recht und getreulich dienen möchte. Des Tages verschmachtete ich vor Hitze und des Nachts vor Kälte. Und um solches meines Laufens und meiner Arbeit willen hat dich der HErr gesegnet.

\subsection{Und nun, wann soll ich auch mein Haus versorgen?}

149. Dies ist eine sehr harte Strafe; denn Jakob will also sagen: Du bist mein Schwiegervater, derhalben solltest du ja von wegen des natürlichen Rechtes mich alsobald, da meine vier Söhne geboren waren, ermahnt haben, daß ich mein Haus auch versorgte; und wenn du ein göttlich oder ehrlich Gemüth hättest, so solltest du billig also mit mir reden: Mein lieber Jakob, ich sehe, daß jetzt unser Geschlecht gemehret wird, dieweil mir vier Kindeskinder geboren sind; ich will dir irgend ein tausend Schafe geben, davon du für dich und dein Gesinde eigen Gut bekommen mögest. Denn wer wollte doch seiner Tochter, und seinem Schwiegersohne, der da ehrlich und fromm ist, und der vornehmlich nicht seinen eigenen Nutzen stiftet, sondern dem Schwiegervater treulich dient, nicht helfen? Das thut aber Laban nicht, darum ist Jakobs Klage über seine Armuth höchst gerecht, da er sagt: „Wann soll ich auch mein Haus versorgen?“ Ich habe zwei Weiber und auch so viel Mägde, dazu sieben Kinder: was nun diese alle erarbeiten und erwerben können, dieses genießest du, und den Nutzen, so wir davon haben sollten, nimmst du hinweg.

\quellseite{161}
Denkst du nicht daran, daß ich nach göttlichen und natürlichen Rechten auch einmal mein Haus werde versorgen müssen? Er wird hernach zu Laban sagen: „Wo nicht GOtt auf meiner Seite gewesen wäre, du hättest mich leer lassen ziehen“ 2c. Ist denn nun das nicht unbillig und tyrannisch, seinen Töchtern, Kindeskindern und seinem Schwiegersohne das Brod mit Gewalt aus dem Munde und Rachen reißen?

150. Denn das ist der Nachdruck dieser Worte, und wiewohl Jakob dieselben ihm fein züchtig und demüthig vorgehalten hat, begreifen sie doch eine harte Strafe in sich. Ich bin, will er sagen, ein Mann zu vier Weibern, bin zu so viel Kindern Vater; für dieselbigen aber arbeite ich nicht, bin ihnen nicht nütze, sondern ich nähre dich und dein Haus: du verschlingst und reißest hinweg, was wir mit unserer Arbeit erworben haben. Wenn auch nur gar ein wenig Vernunft oder gemeiner Witz an dir wäre, oder wenn auch irgend eine Scham oder Freundlichkeit an dir wäre, so würdest du uns nicht also nackend und bloß umher gehen lassen, daß wir darben und so gar arm sein müssen. Das heißt, das heilige Kreuz und das Gold von der Heiligen Füßen stehlen; es heißt, den armen hungrigen Lazarus vollends verschmachten lassen, daß er Hungers sterben muß.

151. Darum ist es ein unglaublich Exempel eines unerhörten Geizes an Laban, welches wir darum groß machen sollen, daß wir lernen die Welt mit ihrem Ungeheuer, dem Geiz, verdammen und verfluchen, und daß wir verkündigen die künftige Strafe, so darauf folgen wird, und wie die Güter, so mit Unrecht erworben sind, verstieben und verschwinden werden. Wie das die Sprüchwörter, so bei allen Völkern gebräuchlich, genugsam bezeugen. Man sagt in einem lateinischen Sprüchwort: De male quaesitis non gaudet tertius haeres, das ist: Uebel genommen Gut kommt dem dritten Erben nicht zu gut 2c. Und ihr jungen Gesellen werdet es dermaßen sehen, daß die Kinder der Geizhälse darben und betteln werden, und daß dagegen auch die Verheißung des Heiligen Geistes erfüllt werden muß im 37. Psalm V. 25.: „Ich habe noch nie gesehen den Gerechten verlassen, oder seinen Samen nach Brod gehen.“ Wiederum heißt es auch also: Ich habe noch niemals einen Gottlosen gesehen, der erhalten worden wäre; wie im selbigen Psalm V. 35. 36. steht: „Ich habe gesehen einen Gottlosen, der war trotzig, und breitete sich aus, und grünte wie ein Lorbeerbaum. Da man vorüber ging, siehe, da war er dahin; ich fragte nach ihm, da war er nirgend funden.“ Und die Erfahrung stimmt mit der heiligen Schrift fein überein. Also wird Laban auch verflucht werden als ein Räuber und Mörder seiner Töchter und Kindeskinder, und seine Güter werden zerstreuet werden.

152. Das ist aber eine sehr schöne Rede, und eine rechtmäßige Ursache, warum Jakob um seine Entlassung bittet: Denn mir ist befohlen, will er sagen, daß ich mein Haus auch versorge und meinen Hausgenossen wohl vorstehen soll, nach göttlichem und natürlichem Recht. Denn „so jemand die Seinen, sonderlich seine Hausgenossen, nicht versorget, spricht St. Paulus 1 Tim. 5, 8., der hat den Glauben verleugnet, und ist ärger denn ein Heide.“ Laban aber hört gar nichts, läßt sich auch die Ehrbarkeit und die Noth seines Schwiegersohnes nicht bewegen; wie man in einem lateinischen Sprüchwort sagt: Venter caret auribus: Der Bauch hat keine Ohren.

\section[Fünfter Theil]{Fünfter Theil. Wie Laban Jakob überredet zu bleiben; wie Jakob einen gewissen Lohn bestimmt und durch Kunst seinen Lohn sehr groß macht}

\subsection{I}

\begin{quote}
\textbf{V. 31. Er (Laban) aber sprach: Was soll ich dir denn geben?}
\end{quote}

153. Er fragt ihn furchtsam genug; als wollte er sagen: Wo du irgend hundert Gulden bedarfst, so bitte ich, du wollest nur vier von mir fordern. Sollte er ihm aber nicht alsobald mit solchen Worten begegnet sein: Ich erkenne deine Treue und deinen Fleiß wohl und bedarf nicht viel Worte? ich will dir von selber gutwillig alle Freundschaft angeboten haben, und in allem, das meinem Verstande nach zu deinem Nutzen und Frommen gereicht, will ich dir als meinem Schwiegersohne und Sohne gerne dienen. Siehe aber, was der Geiz thut; er fragt noch lange: „Was soll ich dir denn geben?“ siehe, daß du nicht zu viel von mir forderst. Ei du schändlicher Hund! Jakob ist dein Schwiegersohn, es sind deine Töchter und Kindeskinder, und du fragst noch lange, wie du ihren Nutzen schaffen sollt? Mose beschreibt solchen trefflichen \unsicher{Eulcio} herzlich, viel besser und eigentlicher, als irgend ein Apelles ihn mit seinen Farben hätte abmalen können.

155. Jakob aber erkennt Labans betrügliches und schalckhaftiges Herze wohl, darum denkt er: Wenn ich schon Geld fordern würde und du mir auch versprechen würdest solches zu geben, so wirst du mir doch nichts geben; du wirst irgend eine List oder Betrug finden, damit du den Vertrag, so ich mit dir gemacht hätte, mit List umstoßen mögest; du wirst jetzt irgend einen Schaden an deinem Gute, bald aber andere Beschwerung vorwenden, damit du mich immer mögest hinhalten; denn die nicht gerne bezahlen, können unzählig viel Entschuldigung finden. Darum bietet er ihm eine andere Bedingung an, wie im Text folgt.

\begin{quote}
Jakob sprach: Du sollst mir nichts überall geben; sondern so du mir thun willst, daß ich sage, so will ich wiederum weiden und hüten deiner Schafe.
\end{quote}

156. Ich will nicht, daß dir an deinem Hab oder Gut etwas soll abgehen; du magst zwar das alles wohl behalten, was ich dir mit meiner Arbeit und Fleiß erworben habe, und ich will bei dir bleiben und will eben mit solcher Treue deine Heerde weiden, wie bisher geschehen ist, und will dafür keinen gewissen Lohn fordern, sondern will des Segens des HErrn warten, und will für gut nehmen, was der mir geben wird.

Dieses ist dem Geizhalze Laban eine sehr angenehme Rede gewesen, da er hört, daß er jetzt kein baar Geld hat ausgeben müssen, und daß er den Knecht, der ihm so treu war, noch etwas länger gebrauchen könnte. Derhalben folgt nun der Bund zwischen Laban und Jakob.

\begin{quote}
\textbf{V. 32. 33. Ich will heute durch alle deine Heerde gehen, und aussondern alle fleckichte und bunte Schafe, und alle schwarze Schafe unter den Lämmern, und die bunten und fleckichten Ziegen. Was nun bunt und fleckicht fallen wird, das soll mein Lohn sein. So wird mir meine Gerechtigkeit zeugen heute oder morgen, wenn es kommt, daß ich meinen Lohn von dir nehmen soll, also daß, was nicht fleckicht oder bunt, oder nicht schwarz sein wird unter den Lämmern und Ziegen, daß sei ein Diebstahl bei mir.}
\end{quote}

157. Labans Hab und Gut sind Schafe, Ziegen, Lämmer, Böcke und ander Vieh gewesen. Dieweil aber der Text etwas dunkel ist, muß man zu allererst einen gewissen Verstand herausnehmen. Denn Hieronymus klagt, daß dieser Text von den siebenzig Dolmetschern ganz unordentlich verdolmetschet sei, und sagt, er habe keinen gesehen, weder unter den hebräischen noch lateinischen Auslegern, der ihn deutlich genug erklärt hätte. Etwa ebendaselbe klagt Burgensis auch. Mit Lyra bin ich zu guter Maßen wohl zufrieden, wie er es ausgelegt hat, wiewohl er noch in etlichen Stücken zweifelt, nämlich, daß es die Grammatik dunkel macht und da die Worte nicht gar heraus geredet werden, und schier sechs Verse oder Zeilen etwas schwer zu verstehen sind.

158. Aber dieses ist die Summa davon: Wir wollen, sagt Jakob, einen Bund machen; ich fordere nichts von dir; sondern was mir durch GOttes Segen zufallen wird, damit will ich zufrieden sein. So laß uns nun durch deine Heerden gehen, da heraus nehmen alles, was fleckicht und bunt ist, also daß auf einem Theil die Schafe bleiben, die nur einerley Farbe haben und nicht bunt sind, auf die andere Seite aber sollen die bunten gestellt werden. Es sind aber viererlei Heerden, nämlich, Widder, Schafe, Ziegen und Böcke, bisweilen wird aber synecdotisch nur Eine genannt, bisweilen mehr, sie werden nicht allezeit alle mit einander aufgezählt; daher kommt es, daß der Text schwer ist zu verstehen. Da aber die Heerden von einander getheilet sind, nimmt Jakob die nicht bunt sind, die aber fleckicht und bunt sind, werden Laban zugeeignet; und Jakob bietet seinem Schwiegervater diese Bedingung an, daß alles, was bunt und fleckicht von den Schafen gezeuget wird, die nur schwarz oder weiß sind, Jakob zufallen solle anstatt seines Lohnes: die andern Lämmer aber von einerlei Farbe, sie seien gleich weiß oder schwarz, sollen Laban, seinem Schwiegervater, zufallen.

\quellseite{162}

159. Diese Bedingung nimmt Laban gern an, und hat gar keinen Argwohn, daß irgend ein Betrug oder List darunter möchte verborgen sein, sintemal er in so vielen Jahren Jakob treu und aufrichtig erfunden hatte. Derhalben steht ihm der Hals offen nach einem neuen Gewinn, und lacht und verwundert sich stillschweigend der Einfalt Jakobs, der die kleine Anzahl Schafe erwählet hatte, und hat ihm die Lämmer von Einer Farbe gelassen, deren ohne Zweifel viel mehr sein würden. Denn das ist natürlich, daß weiße Schafe weiße Lämmer und schwarze Schafe auch schwarze Lämmer zeugen. Das kommt aber gar selten vor, daß die Schafe, die einerlei Farbe haben, bunte Lämmer zeugen, es sei denn, daß die Menschen irgend künstlich und mit sonderlichem Fleiß der Natur helfen, wie hernach wird gesagt werden. Darum hat sich Laban dessen gefreut, daß ihm an seinem Gut und Habe nicht allein nichts würde abgehen, sondern daß ihm solches auch gemehret werden würde; denn die Natur steht auf seiner Seite.

160. Aber eben damit gibt er abermals an den Tag seinen verfluchten Geiz; denn er hätte lieber seinen Schwiegersohn ermahnen sollen, daß er sich selbst nicht wollte übervortheilen, denn daß er solche unbillige Bedingung sollte angenommen haben. Denn er hat für sich die Schafe, so nur weiß oder schwarz waren, ausgelesen und gewählet, von welcher natürlicher Weise nicht bunte Lämmer, sondern den alten Schafen gleiche haben müssen gezeuget werden. Derhalben sollte er gesagt haben: Was willst du thun, mein lieber Schwiegersohn? du willst dich selbst um deinen verdienten Lohn bringen, wenn du die bunten Lämmer anstatt deines Lohnes für dich nehmen willst: es ist ein unbilliger Vertrag und wider die Natur.

161. Aber er hat an der keines gedacht. Ja, das noch mehr ist, er freut sich, daß er Ursache hat, den armen Narren abermals zu betrügen und sein zu spotten, und hoffet, wo er jemals zuvor an seinem Gut habe zugenommen, so wolle er jetzt erst sehr zunehmen und reich werden, und dasselbige auf die Weise, so Jakob selbst hat vorgeschlagen. Ein solch greulich Ding ist es um den Geiz, daß er die Menschen von aller Göttlichkeit und Freundlichkeit abwendet, und macht, daß sie gar greulich und mit teuflischer Bosheit behaftet werden. Denn also wird dieser Laban hier beschrieben und abgemalet, den der Geiz so gar wild, greulich und unfreundlich gemacht hat, daß er seinem Schwiegersohn und Töchtern auch nicht eine Klaue von aller seiner Heerde mit gutem Willen geben wollte. Dieses ist die Meinung dieses Bundes zwischen Laban und Jakob, darin zugleich der greuliche Geiz Labans hervorgehoben und groß gemacht werden soll.

162. Nun wollen wir auch den Text besehen, welcher, wie ich gesagt habe, etwas dunkel ist, weil die Worte so kurz abgebrochen sind. Zum Ersten sagt er, daß er alles Vieh, das da fleckicht ist, von dem andern absondern wolle, wenn es auch nur Einen Flecken hätte. Als, wenn irgend ein schwarzes Schaf an der Stirn oder sonst am Leibe Einen weißen Flecken hätte, oder ein weiß Schaf nur Einen schwarzen Flecken hätte, wo es auch am Leibe sein möchte. Und läßt sich ansehen, als wolle er mit demselben Namen im allgemeinen oder synedochisch begriffen haben alles Vieh, das nicht allein Einen Flecken, sondern auch viel kleiner Flecken hätte; denn das heißt bei den Hebräern, Vieh, das nakod ist. Darnach nennt er auf Hebräisch die Schafthalu, welches wir eigentlich verdeutscht haben, bunte Schafe. Denn also werden diese zweierlei Namen unterschieden: diese haben größere, breitere Flecken von vielerlei Farben, weiß oder schwarz; jene aber haben kleinere Flecken. Zum Dritten sagt er, daß er alle schwarzen Schafe unter den Lämmern aussondern solle; damit meint er die schwarzen Lämmer. Es ist aber ein Unterschied unter den zwei hebräischen Worten, seh und keseb; jenes wird insgemein von allen Schafen gesagt, dieses aber heißt ein jährig Läm(m)lein, oder das noch jünger ist: wenn es über das Jahr kommt, so heißt es nicht mehr keseb: wir heißen es ein Osterlämmlein. Zum Vierten, sondert er auch unter der Heerde der Ziegen die fleckichten und bunten aus, wie mit den Schafen; aber aus den Lämmern nimmt er nur die schwarzen.

163. Darauf folgt aber im Text: „Das soll mein Lohn sein.“ Hier ist eine Ellipse, und man muß hier ergänzen: Alles, was davon gezeuget wird, das denen gleich ist, und was bunt sein wird von den weißen Schafen. Denn mich dünkt, daß nur die Heerde der weißen Schafe dem Jakob geblieben sind; wiewohl die andern haben wollen, daß auch die schwarzen Schafe dazu gethan seien. Und ist dieses die Meinung, daß Jakob sagen will: Alles, was die weißen Schafe von Lämmern tragen werden, das bunt und fleckicht ist, das soll mein sein; was sie aber von weißen Lämmern tragen werden, soll dein sein.

164. Derhalben ist es eine sehr unbillige Bedingung, sintemal es wider die Natur ist, daß weiße Ziegen oder Schafe schwarze oder bunte Lämmer oder Ziegen tragen sollen. Jedoch hat Jakob diese Bedingung dem Geizhalse also angeboten. „So wird mir meine Gerechtigkeit“, sagt er, „zeugen“, das ist, ich werde gerecht sein. Denn mit diesem Vertrag liegt es gar am Glück und wie es fallen wird. Wo es vielleicht also kommen würde, daß mir die weißen Schafe bunte Lämmer tragen würden, daß du den Argwohn nicht schöpfen sollst, als wäre ich damit hinterlistig oder betrüglich umgegangen, wiewohl ich dir nie ein Unrecht gethan habe. Was aber nicht fleckicht, sondern einfarbig und weiß bei mir wird gefunden werden, davon sollst du sagen, es sei dir gestohlen, und sollst mich frei des Diebstahls zu beschuldigen haben.

165. Dieses ist zumal eine harte Verpflichtung, daß sich der fromme heilige Mann gegen den bösen Buben dermaßen verpflichtet hat. Und es scheint daraus, daß Laban Jakob zum öftern wird geschmäht haben, wenn irgend seiner Begierde oder Geiz nicht genug geschehen ist, als sollte ihm etwas gestohlen sein, oder aber, als sollte das Gesinde Jakobs irgend Schafe oder Ziegen von der Heerde geschlachtet und gefressen haben. Oder er hat ja zum wenigsten mit solcher Strafe und Schmähen Jakob bedrohen und bei ihm verhüten wollen, daß er nichts zu seinem eigenen Nutzen verwenden sollte, wiewohl er an seiner Treue und Fleiß nicht hat zweifeln können.

\begin{quote}
\textbf{V. 34. Da sprach Laban: Siehe du, es sei, wie du gesagt hast.}
\end{quote}

166. Es läßt sich ansehen, als habe Laban noch gezweifelt und Jakob nicht getrauet. Ach, sagt er, daß doch jetzt einmal geschehen möchte, wie du gesagt hast. Wohlan, wir wollen es versuchen, bis ich sehe, ob diese Schafzucht auch also gerathen wolle, wie ich hoffe, daß mir es zuträglich sein werde. Wiewohl er sieht, daß Alles seinem Geiz entspricht, nimmt er es doch noch nicht für gewiß an, sondern mit einer Bedingung, auf daß er den Vertrag ändern und aufheben möge, wenn oder wie oft er selbst wolle; wie wir hernach (Cap. 31. V. 41.) hören werden, daß Jakob zu ihm sagt: „Du hast mir meinen Lohn zehnmal verändert.“ Darum wird die Abwechselung des Geizes Labans je mehr und mehr gemehret; denn er sperret sich noch, da ihm sein Schwiegersohn auch einen sehr billigen Vertrag anbietet und der für ihn selbst fast unbillig ist. Er bewilligt ihn nicht förmlich; nicht darum, daß er es dafür halten sollte, daß solcher Vertrag ihm nicht würde nützlich und wohl anzunehmen sein, sondern daß er die Macht behalten will, denselbigen zu widerrufen und sich zu weigern ihn zu halten, es gerathe gleich wie es wolle. Siehst du wohl, was der Geiz thut?

167. Derhalben hat der Poet herrlich und wohl gesagt: Auri sacra fames, quid non mortalia pectora cogis? Du verfluchter Geiz, wozu kannst du die Herzen der Menschen nicht bringen? Denn die Menschen werden durch den Geiz also verändert, daß sie nichts Menschliches an sich behalten, sondern werden gar zu Klößen und Gößen, daß sie gar keine menschliche Empfindung behalten. Und wann kann sein daher ziehen, daß im 115. Psalm V. 4–8. steht: „Die Gößen der Heiden sind Silber und Gold. Die solche machen, sind gleich also, und alle, die auf sie hoffen.“ Die Macherer, Räuber und Geizhälse sind nicht Menschen; „sie haben Augen, und sehen nicht, sie haben Ohren, und hören nicht.“ Und Sirach sagt: Es ist nichts Böserees, Schändlicheres denn ein geiziger Mensch. Denn die Geizhälse sind gar unfreundliche, greuliche und tyrannische Menschen; wenn sie alle Menschen mit einem Heller bei dem Leben erhalten könnten, gäben sie doch den Heller nicht. Darum sind es Mörder und Todtschläger; denn sie rauben und verschlingen, was andere Leute genießen und von dem sie sich erhalten sollen.

\quellseite{163}

168. Wo nun auch ein klein Fünklein menschlicher Freundlichkeit noch an Laban gewesen wäre, so sollte er doch durch die Tugend und die Großmuth seines Schwiegersohnes zur Freundlichkeit und Wohlthaten bewogen worden sein, und ja billig also gedacht haben: Ich sehe, daß mein Schwiegersohn fromm und ehrlich ist, und daß er es ja auf meiner Seite und mit meinen Sachen so gern gut sehen wollte, daß er mir einen solchen Vertrag anbietet, daß er mir nicht allein leidlich, sondern auch sehr nützlich und zuträglich sein möchte; aber da sei Gott für, daß ich den annehmen sollte. Er hat meine Töchter zu Weibern, mit denen er ehrlich und freundlich lebt, und hat mir vierzehn Jahr treulich gedient: wahrlich, ich muß greulicher sein, denn irgend ein unvernünftig Thier sein kann, wenn ich mit seinem Schaden meinen Nutzen und Frommen stiften wollte. Aber er denkt an der keines, ja, läßt sich auch davon gar nichts träumen. Denn der Geiz löscht alle Freundlichkeit, Scham und auch die Natur des Menschen selbst aus, und macht aus den Menschen eitel silberne und goldene Götzen etc.

169. Deshalb gefällt mir sehr die Fabel der Poeten von dem Könige Midas, von dem sie schreiben, daß alles, was er angerührt habe, Gold geworden sei, wie er gewünscht hätte, daß es geschehen möchte. Ja, solche Gesellen sind eben auch die Geizhälse, nämlich, wie silberne und goldene Götzen, die gar keine menschliche Empfindung haben. Und es ist an einem andern Ort von einem heidnischen Schriftsteller herrlich und fein geredet, daß ein geiziger Mensch nimmermehr etwas Gutes thue, denn allein wenn er stirbt. Denn die Geizhälse sind eine unnütze, ja eine schädliche Erdlast, und wo man in Städten und Landen gut Regiment angerichtet hätte, sollte man sie immer wegjagen. Denn wo sie sind, da ist ein gemein Verderben der Gemeinschaft, so die Menschen unter einander haben; und sie sind eine Pest der menschlichen Gesellschaft.

\begin{quote}
\textbf{V. 35. 36. Und sonderte des Tages die sprenglichten und bunten Böcke, und alle fleckichte und bunte Ziegen, wo nur was Weißes daran war, und alles, was schwarz war unter den Lämmern; und thats unter die Hand seiner Kinder; und machte Raum dreier Tagereisen weit zwischen ihm und Jakob. Also weidete Jakob die übrige Heerde Labans.}
\end{quote}

170. Mose gebraucht hier nun andere Namen denn zuvor. Droben (V. 32. ff.) hat er gesagt von den schwarzen Lämmern, item, von den fleckichten Schafen und Ziegen; hier nennt er nun die Böcke: auf daß er damit anzeige den großen Geiz an Laban, der nicht zu erledigen gewesen ist, welcher auch das Vieh ausgesondert hat, so im Vertrage nicht ausdrücklich ist genannt worden. Denn wiewohl wir droben gesagt haben, daß solches durch die Figur synekdoche geschehe, so läßt es sich doch ansehen, daß er nicht vergeblich die vorigen Namen jetzt mit andern Namen geändert habe.

171. Das hebräische Wort akudim bedeutet die großen Böcke, so vor der Heerde hergehen; wie tejaschim die kleinen Böcke oder Ziegenböcklein bedeutet, die man füttert und aufzieht, daß man sie essen will und nicht zur Zucht. Was eigentlich der Ursprung sei des Wortes akudim, weiß ich nicht. Das Wort akad heißt binden; darum deuten sie es, daß es Böcke soll heißen, die an den schwarzen Füßen einen weißen Ring umher haben, oder dagegen einen schwarzen Ring an weißen Füßen, daß es sich ansehen läßt, als ob die Füße damit gebunden seien; oder, welches mir am besten gefällt, er will damit verstanden haben, die über den Rücken einen langen Striemen haben, wie fast alle Böcke pflegen zu haben; daß sich solches auch dazu reime, den Geiz Labans damit arg zu machen, der auch auf das allergenaueste die Böcke ausgesondert hat, deren Flecken man kaum hat merken können.

172. Denn darum sagt er: „Wo nur was Weißes daran war, und alles, was schwarz war unter den Lämmern.“ Also wollte er sagen: Nachdem er schon die bunten Schafe ausgesondert hatte, hat er noch den Fleiß auch darauf gewandt, daß er an einem jeglichen Schafe oder Bock, welche einfarbig waren, auch die einzelnen Haare an den Füßen oder am Barte genau besehen hat, ob etliche fleckicht oder bunt wären. Eine solche genaue und fleißige Sonderung hat er gebraucht wider seine Töchter und seinen Schwiegersohn, auf daß er damit beweisen möchte, wie so ein großer Geizhals und karger Filz in seinem Herzen stecke, dem billig Gott und alle Menschen feind sind. Und bisher haben wir nun von dem Bunde oder Vertrage gesagt zwischen Laban und Jakob.

\subsection{II}

\begin{quote}
\textbf{V. 37—39. Jakob aber nahm Stäbe von grünen Pappelbäumen, Haseln und Kastanien; und schälte weiße Streifen daran, daß an den Stäben das Weiße bloß ward; und legte die Stäbe, die er geschält hatte, in die Tränkrinnen vor die Heerden, die da kommen mußten zu trinken, daß sie empfingen, wenn sie zu trinken kämen. Also empfingen die Heerden über den Stäben, und brachten sprenglichte, fleckichte und bunte.}
\end{quote}

173. Wir haben gehört, wie Laban so große Begierde gehabt und sich geängstiget hat nach einem neuen Raube und Gewinn. Was geschieht aber? Gott zeigt Jakob eine sonderliche Kunst, damit er der Natur helfen und sie ändern kann. Denn er nimmt Stäbe von dreierlei Bäumen und schält dieselbigen, nicht daß er die Schale davon gar abgezogen hätte, sondern daß er sie bunt machte, daß sie weiß und schwarz gewesen sind, also daß man an einer Seite das Weiße gesehen hat, da die Rinde abgeschält war, an der andern Seite aber das Schwarze, da die Rinde geblieben war. Dieselbigen Stäbe legt er in die Tränkrinnen, daß die Schafe solche Stäbe vor ihren Augen liegen haben, und wenn sie die ansehen, empfangen mögen. Es ist eine künstliche Philosophie gewesen oder eine Magie, damit er zuwege gebracht hat, daß die Schafe in der Brunst, da sie die bunten Stäbe angesehen, bunte Lämmer gebracht haben. Und also sind von der weißen oder schwarzen Heerde bunte und fleckichte Lämmer gekommen.

174. Das Wort „sie empfingen“ lautet im Hebräischen, jechemu, welches herkommt von jacham; die Schafe haben sich bei den Stäben gewärmt; das ist so viel gesagt: Sie haben mit einander gelaufen, wie es die Schäferhirten pflegen zu nennen. Und steht dieses Wort auch im 51. Psalm V. 7.: „Meine Mutter, jechemati, hat mich in Sünden empfangen“, das ist, ich bin empfangen worden in der Brunst und schändlichen Lust des Fleisches. Dadurch wird verstanden die Brunst der Kinderzucht, welche vor dem Falle Adams rein und zur Kinderzucht geschaffen und nötig gewesen ist; jetzt aber ist sie vergiftet und verderbt durch die Erbsünde: es ist keine unschädliche Brunst, wie sie anfänglich gewesen ist, sondern ist durch die Geilheit und böse Lust verderbet.

175. Also beträgt nun Jakob durch Kunst und Geschicklichkeit oder durch die natürliche Magie die Kunst oder vielmehr Bosheit Labans, so er an Jakob geübt hatte; welche Kunst die Väter entweder durch langen Gebrauch oder aber durch Unterricht der Vorfahren gelernt haben. Denn Jakob wird ohne Zweifel darin von irgend einem Patriarchen unterwiesen worden sein, oder aber Gott wird es ihm durch den Heiligen Geist eingegeben haben. Und es ist eine gewisse Kunst und die sich mit der Lehre der Aerzte fein reimt, die da sagen, daß in der Empfängniß aller Thiere, nicht allein der unvernünftigen Thiere, sondern auch der Menschen, der Frucht eine sonderliche Gestalt oder Flecken eingebildet soll werden beide aus den Gedanken und auch aus mancherlei Gegenständen, so dem Herzen oder auch den Augen vorkommen, nicht allein in der Brunst der Empfängniß, sondern auch nach der Schwangerschaft.

176. Hieronymus und auch die Physiker erzählten ein Exempel von einer Königin, die ein Kindlein geboren hat, welches eine Gestalt und Angesicht gehabt hat wie ein Mohr; und dasselbe soll geschehen sein, sagen sie, weil sich ihr ein Mohr stark eingeprägt hat, der auf eine Tafel gemalt gewesen ist bei den Betten. Desgleichen sagen sie auch noch von einem andern Weibe, die des Ehebruchs beschuldigt gewesen, nämlich darum, daß sie häßlich war von Gestalt und hätte doch ein schön Kindlein geboren, welches Vater und Mutter und dem ganzen Geschlechte ungleich gewesen; und sie wäre verdammt worden, wo sie nicht Hippokrates errettet hätte, da er eine Vermahnung that, daß man sie fragen sollte: Ob sie auch in der Schlafkammer irgend eine gemalte Tafel gehabt hätte, die sie mit Lust hätte angesehen? und da man die Tafel gefunden hat, ist sie von den Richtern unschuldig erkannt worden.

177. Also sehen wir bisweilen, daß die Kindlein am Angesicht, an den Augen, Backen, oder am Nacken bunte Flecken haben, Blut, oder sonst von anderer Farbe, nämlich, wenn die schwangeren Weiber irgend unversehens etwas Ungewöhnliches angesehen oder davor erschrocken sind und mit den Händen an solche Glieder gegriffen haben. Wir haben hier zu Wittenberg einen Bürger gesehen, der ein Angesicht gehabt wie ein todter Mensch; der sagte, daß seiner Mutter, da sie schwanger war, plötzlich ein menschlicher Leichnam vor die Augen gekommen wäre, davor sie also erschrocken, daß das Angesicht der Frucht in ihrem Leibe davon die Gestalt eines todten Menschen bekommen hätte.

\quellseite{164}

178. Daselbste pflegt auch also mit Fleiß zu geschehen, wenn das Vieh und die unvernünftigen Thiere mit einander laufen. Wie Hieronymus sagt, daß in Spanien die schönen artigen Rosse vor die Mutterpferde gestellt werden, wenn sie laufen, auf daß davon junge Füllen gezeuget werden, die solchen schönen Rossen mögen gleich sein.

179. Darum ist mit den schwangern Weibern nicht zu scherzen, sondern man soll fleißig Acht auf sie haben um der Frucht willen. Denn es ist viel unzählige Gefahr der unzeitigen Geburt und sonst mancherlei scheußlicher Mißgeburt halben. Derhalben soll ein Ehemann alsdann am allermeisten mit Vernunft bei seinem Weibe wohnen, wie St. Petrus 1. Epist. 3, 7. sagt. Ich erinnere mich noch wohl, da ich ein kleiner Knabe war, daß zu Eisenach eine schöne und züchtige Matrone eine Matte geboren hat; welches davon war hergekommen, weil einer von den Nachbarn einer Matte eine Schelle angehängt hat, daß davor die andern sollten fliehen. Dieselbige Matte mit der Schelle ist der schwangeren Frau begegnet, und weil sie um die Sache nichts gewußt, ist sie davon, daß ihr die Matte so plötzlich vorgekommen ist und sie derselben ansichtig geworden, auch also erschrocken, daß die Frucht in ihrem Leibe davon die Gestalt desselbigen Thierleins bekommen hat. Solche Exempel sind sehr gemein, wenn die schwangern Weiber oftmals plötzlich bewegt und erschreckt werden, auch mit Gefahr ihres Leibes und Lebens.

180. Darum soll man sich hüten, daß sie nicht etwa heftig bewegt werden, es sei gleich äußerlich am Leibe, oder innerlich am Gemüth. Denn die auf die schwangern Weiber keine Achtung haben und der zarten Frucht nicht verschonen, das sind Todtschläger und Kindermörder. Wie man etliche Männer findet, die so gar greulich und tyrannisch sind, daß sie sich auch nicht scheuen, die armen schwangern Weiber zu schlagen. Ei, das sind tapfere Helden, die haben wahrlich einen großen Muth gegen das schwache weibliche Geschlecht, aber sonst sind es wohl verzagte Männer. Denn wir haben neulich gehört, daß ein Fürst, der sonst viel andere Sünde und Schande getrieben hat, über seine Gemahlin das Schwert soll gezückt haben, da sie krank gewesen und auf dem Bette gelegen hat. Ei, wie ist mir das so ein tapferer Held und gewaltiger Kriegsmann! ja, das gebühret einem tapferen Helden gar nicht, sondern ist eine große Sünde und Schande. Denn wo du ein rechtschaffener Mann bist, so wirst du deines Gleichen wohl finden, mit dem du dich schlagen kannst. Die großen tapfern Helden sind stark wider andere Stärke und schwach gegen die Schwachen. Was soll das sein, daß du wider ein Kind und schwanger Weib einen Kampf anfängst? Dieses arme weibliche Geschlecht hat sonst ohne das Gefahr übrig genug auf sich liegen, auch bei frommen Ehemännern, die sich wohl wissen zu halten, nämlich, von bösen Nachbarn, vom Teufel, von mancherlei Gespenstern und Bildern der unvernünftigen Thiere. Nun ist es große Sünde, daß durch deine Grausamkeit solche Gefahr sollte gemehret werden.

181. Ja, es haben auch die Heiden an ihren großen Helden diese Tugend gerühmt, daß sie sich gegen ihre Weiber friedlich, still und freundlich gehalten haben. Denn dasselbige bezeuget Homer, da er die großen tapfern Helden, als, den Achilles, Hektor und Andere mehr, beschreibt. Da Hektor seinen Harnisch und Rüstung angehabt, führet er sein Söhnlein. Also sind sie bei den Weibern auch Weiber gewesen, also auch, daß man nichts Weiberisches gesehen hat, als diese großen Helden bei ihren Eheweibern gewesen sind; aber im Kriege, da es Streitens gilt, hat sich Achilles gegen Hektor anders gehalten, denn da er mit seiner Briseis scherzt.

182. Darum ist man denen billig feind, die wider das arme, wehrlose, schwache, weibliche Geschlecht stark und kriegerisch sind und immer balgen wollen. Wir Männer sind nicht dazu geboren, daß wir das schwache weibliche Geschlecht verletzen, sondern schützen und verteidigen sollen; denn das Weib hat einen Leib, der dazu geschaffen ist, daß sie schwanger sein, und Kinder zeugen und nähren soll, und ist vielerlei Gefahr unterworfen; darum soll man mit Vernunft und freundlich mit ihr umgehen. Und bisher ist ein Theil dieses Textes ausgelegt, und Eine Kunst, so Jakob geübt hat, beschrieben worden.

\begin{quote}
\textbf{V. 40—43. Da schied Jakob die Lämmer, und that die abgesonderte Heerde zu den fleckichten und schwarzen in der Heerde Labans; und machte ihm eine eigene Heerde, die that er nicht zu der Heerde Labans. Wenn aber der Lauf der Frühlinger Heerde war, legte er diese Stäbe in die Rinnen vor die Augen der Heerde, daß sie über den Stäben empfingen; aber in der Spätlinger Lauf legte er sie nicht hinein. Also wurden die Spätlinge des Laban, aber die Frühlinge des Jakob. Daher ward der Mann über die Maße reich, daß er viel Schafe, Mägde und Knechte, Kameele und Esel hatte.}
\end{quote}

183. Die andere Kunst, so Jakob gebraucht hat, ist von dem Laufe der Heerde, so im Frühlinge und im Herbste geschieht. Denn Jakob hat seiner Kunst gebraucht im ersten Laufe, das ist, im Frühlinge, wenn die Schafe am stärksten sind; da hat er die Stäbe in die Rinnen gelegt, auf daß die stärksten Lämmer bunt würden, die ihm zugehörten: im Herbste aber und in der Spätlinger Laufe sind die Schafe schwächer, darum daß sie die Hitze der Sonne nicht haben, wie im Frühlinge. Darum hat er dazumal keine Stäbe gelegt, auf daß natürlicher Weise einfarbige Lämmer gezeuget würden, die Laban zugehörten. Es ist aber der Lauf im Frühlinge darum besser, weil derselbige geschieht, da die Sonne noch aufstehet, wenn die Hitze wieder kommt, und die Kraft aller Kräuter und Thiere gemehret wird. Darum hat Jakob verschafft, daß er die besten und stärksten Schafe hätte, nämlich die Frühlinge; wie man die Lämmer, so im Herbste fallen, Spätlinge und Wintershaafe nennt. Es ist übrigens eine fast künstliche Geschicklichkeit, ja, es ist eine Listigkeit und schier eine Schalkheit gewesen; und Jakob hat sich doch daselbst auch nach Billigkeit und solches Maß gehalten, daß er seinem geizigen Schwiegervater keinen Argwohn erregte, daß er also durch Kunst betrogen würde. hat er nichts geändert, damit Laban auch einen Theil behalten möchte und daß er ja nicht gar beraubet würde.

184. Darum ist er fein übervortheilt worden, und hat selbst nicht verstehen können, wie es kommen möchte, daß die Frühlingslämmer bunt gezeuget würden und die Spätlinge der Natur folgeten. Er hat gedacht, es wäre ohngefähr oder aber durch göttlichen Segen also geschehen; darum hat er den Vertrag, wie hernach folgen wird, wohl zehnmal geändert, und wo er jemals zuvor gezeuget hat, daß er geizig wäre, so hat er eben damit ein trefflich augenscheinlich Exempel von sich gegeben seines verfluchten Geizes. Denn dieweil er gesehen, daß die Frühlingslämmer besser und stärker wären, hat er den Bund aufgehoben und hat für sich die bunten gewählt. Da nun Jakob da der Natur mit seiner Kunst nicht geholfen hat, haben die Schafe einfarbige Lämmer getragen und sind abermals Jakob die besten zugefallen. Darum ist Laban abermals betrogen worden, und er hat auch zum dritten, vierten-, ja zum zehntenmal den Vertrag und den Lohn Jakobs geändert. O des schändlichen Geizes, den billig jedermann ansehen sollte!

185. Es schelten und plagen mich aber unsere Rabbinen, daß ich diesen Text im Deutschen nicht recht soll verdolmetscht haben. Denn sie setzen noch die dritte Kunst oder Geschicklichkeit dazu aus dem, daß im Text gesagt wird, daß Jakob die Lämmer geschieden und sie zu Hause unter die Heerde Labans gethan habe etc. Dasselbige legen sie also aus, daß er die fleckichten Lämmer in einen Haufen getrieben und dieselbigen vor der Heerde der Schafe Labans habe lassen hergehen, auf daß die Schafe, so nachfolgeten, durch das Anschauen der bunten, so vorher gegangen sind, andere bunte Lämmer bringen sollten.

186. Aber dies ist von den sehr schändlichen und geizigen Rabbinen hergekommen, die von den heiligen Patriarchen urtheilen, wie sie selbst gesinnt sind, als hätte er sich nicht genügen lassen an der Kunst, so er geübt hatte, sondern hätte auch noch die bunten Schafe vor die einfarbige Heerde gebracht, daß sie die sollten anschauen, auf daß die Lämmer nicht allein von den bunten Stäben, sondern auch von dem Anschauen der bunten Heerde bunt werden möchten; so doch Mose nur von den Stäben rebet, und gar das Gegentheil will verstanden haben, nämlich, daß Jakob die bunte Heerde abgeschieden und an einen andern Ort gestellt habe. Denn er redet von den Lämmern, die er ausgesondert hat; wie sie heutiges Tages die Lämmer auch besonders pflegen zu stellen, wenn sie die alten Schafe in die Weide treiben. Und die zwei Heerden, Labans und Jakobs, sind dreier Tagreisen weit von einander geschieden gewesen, darum hat ja beiderlei Heerde, Jakobs und Labans, an Einem Orte nicht zugleich sein können.

\quellseite{165}

187. Durch solche Kunst aber und Geschicklichkeit, schreibt Mose, sei Jakob sehr reich geworden, oder, wie im Hebräischen steht: Der Mann ist hervorgebrochen, das ist, er ist über die Maßen groß und reich geworden. Und das dünkt mich zwar der rechte eigentliche Verstand dieses Textes zu sein, der sehr dunkel ist; welcher Verstand übereinkommt mit dem, was alle katholischen Lehrer hiervon halten, und ist genommen aus der Vergleichung deß, so hernach folgen wird, welches diesen Verstand auch noch mehr erklären und bestätigen wird; zudem findet sich dieser Verstand auch, wenn man alle Umstände fleißig betrachten und erwägen wird.

188. Es wird aber gefragt: Ob diese That Jakobs auch möge entschuldigt werden, daß er seinen Schwiegervater mit offenbarem Betrug hintergeht? Denn es hat etlichermaßen den Anschein des Geizes, oder noch wohl mehr des Diebstahls oder eines Raubes. Denn warum hält er nicht die Ordnung und Lauf der Natur zu beiderlei Zeit, im Frühlinge und auch im Herbste? Antwort: Aus dem, so droben ist gesagt worden und das hernach folgen wird, kann man mancherlei und etliche Entschuldigung nehmen. Zum Ersten wird er nach menschlichem Recht entschuldigt, welches denen, die da geizigen und ungerechten Herren dienen, erlaubt und zuläßt, wo ihnen die Herren keinen Lohn geben, sondern sie nur berauben und plündern und das Ihre an sich reißen, daß sie, die Knechte, wiederum auch dasjenige nehmen und rauben mögen, was ihnen zum Lohn gebühret; aber doch gebührlicher Weise, daß solches nicht geschehe zu des Herrn Schaden.

189. Also haben die Kinder Israel die Egypter beraubt und solchen Raub als ihren Lohn genommen, daß sie den Egyptern unbilliger Weise, mit großer Beschwerung hatten dienen müssen, dafür ihnen die Egypter noch nicht gelohnet hatten, 2 Mos. 12, 35. 36. Eben auf solche Weise hatte Jakob auch ganze vierzehn Jahre gedienet, und zwar mit vielem Jammer, Beschwerung und Mühsäligkeit geplagt worden, und ist doch gleichwohl dabei seines gebührlichen Lohnes beraubt worden. Darum hat ihm das von Recht wegen gebühret, was er geraubt und genommen hat auch ohne seines Herrn Willen und Wissen. Dies ist Eine Antwort.

190. Zum Andern, wiewohl es eine List und Betrug ist, so hat es doch Jakob aus göttlichem Befehl gethan; denn es ist ihm ein Engel erschienen, und hat ihm diese Magie, die natürlich und ziemlich ist, gezeigt. Darum hat sie Jakob entweder vom Engel oder aber von den heiligen Vätern gelernt, welche große Erfahrung und Verstand der Dinge gehabt haben. Wenn aber nun GOtt die heiligen und gläubigen Männer etwas thun heißt, das ist ohne allen Zweifel heilig und wohl gethan.

191. Zum Dritten, so wird er auch hernach 1 Mos. 31, 42. sagen (denn dies Capitel kann ohne das nachfolgende nicht verstanden werden): „Gott hat mein Elend und Mühe angesehen“; da werden wir hören, was für Angst und Noth ihn dazu getrieben habe, daß er diesen Betrug hat vorgenommen. „Was irgend gestohlen war“, wird er sagen, „habe ich bezahlen müssen. Wo nicht der Gott meines Vaters, der Gott Abrahams, und die Furcht Isaaks auf meiner Seite gewesen wäre; du hättest mich leer lassen ziehen.“ Der gute, fromme und getreue Jakob hat gar keine Hoffnung gehabt, daß er irgend ein wenig eigenes Gut hätte ersparen oder sammeln mögen von wegen der unerträglichen Raublust seines Schwiegervaters. Das ist wahrlich nicht ein geringer Jammer gewesen, dadurch er ist versucht worden zu diesem Raube, vornehmlich, da der göttliche Befehl dazu gekommen ist, darin ihm geboten worden, daß er dies thun sollte. Denn es hat Laban deßenthalben mehr Nutzen gehabt, wird der Engel gesagt haben, denn du ihm nehmen kannst; darum kannst du mit Recht der Kunst und List gebrauchen, auf daß du doch etwas von ihm bringen mögest; nicht das geraubt sei, sondern das dir von GOtt selbst erlaubt und gegeben ist. Dergleichen ist Jakob nach menschlichem und göttlichem Recht, und dazu auch wegen der äußersten Noth entschuldigt.

192. Diesem Exempel soll aber nicht ein Jeder nachfolgen, es sei denn in gleichem Falle. Denn sonst sehen die Schlafsaugen nur auf die That und lassen die Umstände fahren, und wollen es freventlich zum Exempel auf andere Räuberei ziehen. Du sollst aber diesem Exempel mit nichten nachfolgen, es sei denn, daß du Jakob allenthalben gleich seiest, und daß dich alle Umstände in gleichem Fall auch dazu treiben werden. Denn er hat seinem Schwiegervater sein Gut und Habe gemehret und hat dem gierigen und räuberischen Laban so überaus beschwerlich gedienet, der ihm auch des Brodes nicht satt gegeben hat: er hat Hunger, Durst, Hitze, Frost gelitten Tag und Nacht; hat keine Erstattung oder Lohn davon gehabt; ja, es hat Laban über das noch gedacht, ihn des Lohns auch zu berauben, daß sie mit einander waren eins geworden, sintemal er die Macht behält, den Vertrag zu ändern, so oft er selber will. Diese Umstände soll man fleißig ansehen und erwägen, alsdann wird man Jakob des Geizes nicht beschuldigen können und wird niemand leichtlich diesem Exempel nachfolgen wollen. Denn es ist ein solch Exempel und That, die einem großen Helden gebühret, wie deren viel mehr an den Patriarchen gefunden werden. Wir haben aber von solchen sonderlichen Thaten der großen Helden droben geredet. Bisweilen richten sich die ritterlichen großen Helden nach den gemeinen Regeln, bisweilen auch nicht. Jakob aber kündigt hier wider die Regel nicht; denn es kommt ihm das natürliche und weltliche Recht zu Hülfe; und Christus Matth. 10, 10. spricht: „Ein Arbeiter ist seiner Speise“ und seines Lohnes „werth.“

\chapter{Das einunddreißigste Capitel}

\section[Erster Theil]{Erster Theil. Wie Jakob sich unschuldig nach Hause zu ziehen, und von GOtt befehlet bekommt, dies sein Vorhaben zu vollbringen}

\subsection{I}

\begin{quote}
\textbf{V. 1. Und es kamen vor ihn die Reden der Kinder Laban, daß sie sprachen: Jakob hat alle unsers Vaters Gut zu sich gebracht, und von unsers Vaters Gut hat er solchen Reichthum zuwege gebracht.}
\end{quote}

2. Bisher hat Mose erzählt, was in den sechs Jahren geschehen ist, nachdem Joseph geboren war. Denn von derselbigen Zeit an hat Jakob darauf gedacht, wie er sein eigen Gesinde regieren und ernähren möchte, daß er die vorigen vierzehn Jahre nicht hat versorgen können. Nachdem er aber so viel Kinder Vater geworden ist, hat er schier allzulangsam auf seine eigene Haushaltung gedacht. Wie er aber die vierzehn Jahre über mit großer Geduld unter solchem geizigen Herrn viel harte Beschwerungen getragen hat, also hat er nun in diesen sechs Jahren durch wunderbare Kunst und Segen des Herrn ein nicht geringes Gut erworben. Und das geht nun auch ohne Neid nicht zu. Denn er hat sein Gut vor seinem Schwiegervater und dessen Söhnen nicht können heimlich halten. Deshalb, da die sechs Jahre um sind, schickt sich’s, daß er wieder in sein Vaterland ziehen will, wiewohl er damit nicht eilt; und vielleicht würde er auch noch länger haben bleiben wollen, wo nicht der göttliche Beruf dazu gekommen wäre.

2. Aber zu dem großen Jammer und Elende, so er in Labans Hause so viele Jahre lang gelitten hat, kommt nun auch noch das Unglück, daß ihn die Kinder Labans hassen und neiden, und ihn öffentlich schelten dürfen, als habe er alle ihres Vaters Gut geraubt und an sich gebracht. Denn also sagt Mose: „Es kamen vor ihn die Reden der Kinder Labans, daß sie sprachen“ etc. Er hat auch wohl gesehen, daß sich das Angesicht Labans geändert und daß er ganz sauer und unfreundlich gesehen hat. Welches alles er als Zeichen eines solchen gemüthes verstanden hat, daß ihm entfremdet und auch sein müßte.

\quellseite{166}

3. So höre nun die große Lästerung der Kinder Labans, die Jakob dürfen aufrücken und vorwerfen das Gut, welches doch sein Lohn war für den langen schweren Dienst, so er Laban geleistet hatte: so wirst du sagen müssen, daß die ganze gottlose Welt über einem Haufen im Geiz ersoffen sei. Daher sagt Sirach, daß auf Erden nichts Böseres sei denn ein geiziger Mensch. Und ein heidnischer Schreiber, Alimus Publianus, hat auch herrlich und wohl geredet: Avarus nihil recte facit, nisi cum moritur: Ein geiziger Mensch thut nimmermehr etwas Gutes, denn nur wenn er stirbt. Denn also stimmt die ganze Natur und aller Völker einfältige Rede überein, daß sie die Geizhälse verflucht und verdammt haben.

4. Denn es ist ihnen die Abgötterei und alles gottlose Wesen in ihr Herz gepflanzt, nicht allein gegen fremde Leute und die ihnen gar nicht verwandt sind, sondern auch gegen ihre eigenen Hausgenossen und angeborenen Freunde; wie man hier an Laban sieht, der sich gegen seine Töchter, Kindeskinder und gegen seinen Schwiegersohn über die Maßen übel hält. Was könnte doch Böseres und Greulicheres erdacht werden, daß des Verfluchens mehr werth wäre, denn ein solcher Mensch ist, der seinen Töchtern, die von seinem Leibe geboren sind, oder seinem Weibe nicht ein Stück Brodes gönnt? was sollte der wohl einem Unbekannten thun, oder Einem, der ihm gar nicht verwandt wäre? Darum kann ein Geiziger nichts Besseres und nichts Nützliches thun, denn wenn er stirbt. Denn im Leben ist er weder Gott noch andern Menschen, ja, sich selbst auch gar nichts nütze; er kann nicht anders, denn sündigen wider Gott, wider Menschen und auch wider sich selbst. Denn er thut auch seinem eigenen Leibe nimmer etwas zu gute.

5. Derselben gebrauchen nun die Kinder Labans solche Worte, die sehr groß und gewichtig sind, ihren Neid damit an den Tag zu geben. „Jakob“, sagen sie, „hat alle unsers Vaters Gut zu sich gebracht.“ Lüge dich nicht zu Tode! Es ist eine falsche Anklage, die sie über alle Maßen groß machen. O wir armen Kinder, sagen sie, wir sind nun gar verdorben, wovon wollen wir nun hinfort leben und uns erhalten? Dieser hat alle unsers Vaters Gut an sich gebracht. Es ist nun an dem, daß er uns gar auslaugen, und Alles gar hinnehmen und verzehren wird. Auf solche Weise malt Mose Laban ab mit seinen Kindern, nämlich, daß die Kinder dem Vater fast gleich gewesen sind. Er hat Alles weg, sagen sie. Und Jakob hat doch vierzehn Jahre gedient, und hat gar keinen Lohn empfangen, ausgenommen Brod und Kleidung, und damit ist es noch sehr kümmerlich zugegangen; das Andere ist alles Laban zugefallen, ja, Laban würde nimmermehr zu solchem großen Gut gekommen sein, wo Jakob nicht wäre zu ihm in sein Haus gekommen, wie Jakob droben (Cap. 30.) gesagt hat. Darnach hat er ihm nie etwas genommen, sondern hat mit sonderlichem Fleiß die große Heerde der einfarbigen abgesondert von seiner Heerde der bunten Lämmer, welche viel kleiner gewesen ist, und er hat dieselbigen kaum mit Geschicklichkeit und aus göttlichem Eingeben zusammen gebracht. Und ist nichts desto weniger unterdeß die andere Heerde Labans, die sonst groß genug war, auch noch gemehret worden. Jakob hat nur die Frucht seiner Arbeit ausgesondert, und davon hat er dies Gut zusammen gespart, welches zwar gering war, wenn man es mit dem großen Gut Labans vergleichen wollte.

6. So mögen wir nun aus diesen Worten der Kinder Labans abnehmen, was die Geizhälse für ein Gemüth und für eine Art haben. Denn sie sehen darauf nicht, was sie in der Hand und ihrer Gewalt haben; sie sehen nicht auf die große Heerde, auf ihre großen reichen Einkünfte; sondern daß Jakob von ihrer Heerde auch ein wenig Schäflein bekommen hat; welche Heerde er doch zuvor so viel Jahre her mit so großem Fleiß gemehret, und dagegen nur einen Theil der Heerde zum Lohn für seine Arbeit gefordert hat. Das sehen sie an, und daselbst mit großem Schmerz und Neid. Wie sollte ihm denn Laban noch etwas mehr dazu geben, so er ihm doch das mißgönnt und ungern gibt, das dem Knechte nach allem Recht, und vermöge des Vertrages, so sie treulich mit einander aufgerichtet hatten, gebühret hat? Ja, ein solch schändlich, verächtet Ding ist es und Schmach verderben sollen: er kann nicht satt werden, wenn er auch ohne alles Hinderniß alle sein Einkommen und Schätze frei hat, es sei denn, daß er sehe, daß die Andern gar nichts haben.

7. Es ist aber der Neid eine sehr gerechte Strafe auf den Geiz. Denn unterdeß, da der Geizhals begehrt, daß andere Leute darben sollen, braucht er seiner eigenen Güter auch nicht; wie der Spruch des Hieronymus das bezeuget, da er sagt in der Vorrede vor der Bibel: Avaro tam deest, quod habet, quam quod non habet, das ist: Einem geizigen Menschen mangelt es ebensowohl an dem, das er hat, als an dem, das er nicht hat. Also daß irgend ein armer Mann von seinem täglichen Brod besser lebt, denn der Geizhals von all seinem Einkommen. Ich habe selbst auf der Universität zu Erfurt einen Magister gesehen, der sehr viel Gutes hatte, aber der ein so überaus karger Filz war, daß er auch in einem ganzen Jahr nicht einen Bissen Wein zur Gesundheit seines Leibes getrunken hätte. Was ist nun Elenderes an einem Menschen, denn daß Einer in solchem großen Reichthum noch dürsten, hungern und darben soll, und darnach noch mit solcher bösen Begierde geplagt werden, daß er wünscht, daß die Andern alle arm sollten sein, und er sich immer darüber bekümmern und seufzen muß, wenn er sieht, daß ein Anderer auch nur irgend einen Gulden hat, der sein eigen ist. Da siehe nun, ob nicht der Geiz der scheußlichste Greuel unter allen Greueln ist?

8. Und dieser Geiz und Neid der Kinder Labans übertrifft alle anderen Exempel. Denn Jakob, der gute fromme Mann, hat nun mit großer Arbeit, Jammer und Elend so lang gedient und mit seinem sauren Schweiß sein Weib und Kind ernährt, und hat nicht sein eigen Gut, sondern das Haus und die Heerde Labans gemehret und reich gemacht; und muß doch noch leiden und anhören, daß er also gelästert und fälschlich beschuldigt wird, daß sie von ihm sagen: „Er hat alle unsers Vaters Gut zu sich gebracht.“ Es ist wahrlich dies Exempel und der Name Labans werth, daß man ein Sprüchwort daraus mache, wie aus dem Namen des Euclio ein Sprüchwort geworden ist, den großen unnützen Geiz damit anzuzeigen. Denn wie viel ist dessen wohl gewesen, daß Jakob in diesen sechs Jahren mit seiner Kunst nun erworben hat? Er wäre wahrlich wohl werth gewesen, daß man ihm noch viel mehr zugelegt hätte. Was werden wir aber meinen, daß Laban seinem Schwiegersohne wollte gegeben haben für solchen getreuen und langen Dienst, wo nicht dieser Vertrag wäre zwischen ihnen aufgerichtet worden? Er würde ihm eben das gegeben haben, das Jakob hernach sagen wird: „Du hättest mich leer und arm ziehen lassen“; oder er hätte stets für uns Knecht bleiben müssen, von deß Arbeit und Elend Laban hätte mögen Nutzen und Gewinn haben wie zuvor.

9. Denn wir sehen wohl, wie ihnen der Hals offen steht nach diesem Gut Jakobs und wie feindseliger Weise sie dasselbe groß machen. Als wollten sie sagen: Was würde Jakob haben, wenn er ohne unseres Vaters Gut wäre? Er wäre ein armer Bettler und Fremdling. Warum sagen sie aber nicht lieber also: Was würde unser Vater gehabt haben, wo Jakob nicht wäre hieher zu uns gekommen? Denn was Labans Gut gewesen sei, hat Jakob droben (Cap. 30, 30.) angezeigt, da er sagte: „Du hättest wenig, ehe ich herkam.“ Seine Tochter Rahel hat das Vieh und die kleine Heerde gehütet, er hat seinen Hirten oder irgend einen Knecht gebraucht, darum daß die Heerde so gering war. Aber das kehren sie um und fragen: Was hat Jakob mitgebracht, da er zu uns gekommen ist? Er ist leer, nackend und bloß herein gekommen, so mag er nun auch wiederum also hinweg ziehen. Ja, wo soll er denn der Belohnung warten für seinen treuen Dienst, den er bei seinem Hause nun zwanzig Jahre lang geleistet hat? Ei, er hat gar nichts verdient, werden sie gesagt haben; denn er ist ein leibeigener Knecht: wir aber sind frei geboren, und sind Kinder im Hause. Derhalben lasset uns der Güter, so Jakob erworben hat, gebrauchen und ihn hinaus stoßen, arm und nothdürftig, oder aber er mag uns umsonst dienen für und für. Denn also haben sie gemeint, daß sie diesen Dienst von Rechts wegen von ihm fordern möchten und er ihnen denselben auch also schuldig wäre zu leisten, und man könnte mit keinem Recht irgend ein Lohn denken. Ja, sie haben gemeint, es wollte einem Knechte nicht gebühren, daß er so ein groß herrlich Gut haben sollte.

10. Denn darum gebrauchen sie einen \unsicher{son-} herrlichen Namen, auf Hebräisch kabad, welches eigentlich heißt, schwer sein. Aber metaphorisch bedeutet es Ehre. Also steht droben Cap. 13, 1. 2.: „Abraham zog herauf aus Egypten, und war sehr reich von Vieh, Silber und Gold“ 2c. Also sagt man: Das Land wird beschweret mit Theurung, Pestilenz 2c., item: Das Herz Pharaos ist schwer geworden. Und wird dadurch bedeutet, daß die Sachen da seien, und daß es nicht leer und eitel Ding sei. Daher wird es übertragen, daß es auch Ehre heißt; denn Geld und Gut und wo man Alles vollauf hat, das bringt Ehre und große Pracht. Derhalben reden nun die Kinder Labans neidisch und sagen: Er hat aus unsern Gütern alle diese Schwere zuwege gebracht, das ist, diese großen herrlichen Güter. Sie wollen nicht sagen: Er hat eine Frucht, ein Erbe, Habe oder eigen Gut bekommen; sondern sie sagen neidisch und lästerlich, kabad, wie es im Hebräischen lautet. Als wollten sie sagen: Er ist hieher zu uns gekommen wie ein armer Bettler; er hat aber nun in unserm Hause sehr groß Gut gesammelt.

\quellseite{167}

11. Dieses soll man darum etwas fleißiger lesen (denn es läßt sich ansehen, als habe die heilige Schrift in diesem Exempel eine rechte eigentliche Beschreibung des Geizes stellen wollen), nämlich, daß der Geiz eine Welt und Inbegriff aller Schande und Laster, und ein Geizhals ein solcher Mensch sei, der weder Gott, noch sich selber etwas Gutes gönnt; denn was ihm zu seinem Nutzen nicht dient, dasselbe kann er feindseliger Weise groß machen, und ist ihm leid, daß er desselben nicht genießen kann. Derhalben, das man von den undankbaren Menschen sagt: Omnia convicia dixeris, si ingratum dixeris: Wenn du Einen undankbar schiltst, so hast du ihn damit gescholten, daß er aller Sünden schuldig sei: das mag man viel billiger von einem geizigen Menschen sagen. Denn die Undankbarkeit ist eine solche Sünde, da man versäumt, dasjenige zu thun, das man zu thun schuldig ist (Unterlassungssünde), der Geiz aber ist eine Sünde, da man thut, was man nicht thun sollte (Begehungssünde). Es ist eine solche Sünde, damit man greuliche Tyrannei übt. Derhalben übertrifft diese Sünde die andern weit, darum daß sie der ganzen Gemeinde und auch den Leuten insonderheit sehr großen Schaden thut.

\begin{quote}
\textbf{V. 2. Und Jakob sahe an das Angesicht Labans; und siehe, es war nicht gegen ihn, wie gestern und ehegestern.}
\end{quote}

12. Die erste Ursache, warum Jakob aus dem Hause Labans gezogen, ist diese, daß die Kinder Labans über ihn geklagt und ihn fälschlich beschuldigt haben. Die andere Ursache ist, daß er sieht, daß das Angesicht Labans gegen ihn geändert war. Denn der Vater hatte eben solchen Neid und Geiz an sich, wie die Kinder hatten: er kann das Gut seines Schwiegersohnes ohne großen Neid und Schmerz nicht ansehen. Denn er hat begehrt, daß Jakob nur immer sollte arm und ein Knecht sein, daß er seines Jammers und Elendes genießen möchte. Er sieht nicht, daß er schon zuvor durch die Wohlthat Jakobs gewachsen und reich geworden war, und daß er nun anstatt der seinen Tochter vieler Hirten zu einer solchen großen Heerde bedurft hat: aber daß sein Schwiegersohn und seine Töchter ein oder zwei Schäflein, Lämmer oder Ziegen haben, das sieht er wohl und thut ihm solches wehe; er ist darüber nicht fröhlich, wie er zuvor war, da er allein gewachsen und zugenommen hat, und da Jakob für sein eigen Gut gar nicht gesorget hat.

13. Da aber der Herr unwillig gewesen und gezürnt hat, da haben auch die Kinder und das andere Gesinde sich mit ihrem Angesicht nach ihm gerichtet und Jakob auch sauer angesehen. Und dieselbigen haben auch auf ihre Schwestern Lea und Rahel neidisch und feindselig gelauert, wo sie auch mit dem kleinen Finger irgend etwas würden angreifen, das zu Labans Gut gehörte. Sie haben oft falsche Anklagen und Schmähworte erdichtet, und dieselbigen auf den armen Jakob, der bei dem Hausvater verhaßt war, gelegt; derhalben ist ihm Alles in dem Hause zuwider, verdrießlich und gefährlich gewesen. Und dieser Neid und Haß hat nicht nur einen oder zwei Monat gewährt, sondern es ist der Wahrheit gemäß, daß er bald im andern oder dritten Jahre angefangen habe, sobald Jakob angefangen hat zu wachsen und am Gute zuzunehmen.

14. Was kann aber nun einem Menschen Betrübteres und Elenderes widerfahren, denn wo er mit solchem Gesinde umgehen soll, da der Hausvater, die Kinder, Knechte und Mägde sauer sehen, toben und immer zürnen? Und Jakob hat doch gleichwohl die sechs Jahre ausgehalten.

Endlich aber hat er gedacht: Es will nun Zeit sein, daß ich dieses Haus räume, da die sechs Jahre um sind, darin ich mit sehr großem Haß und viel bittern falschen Anklagen beschweret gewesen; wiewohl er nicht eilt, er wäre auch noch damals nicht weggezogen; aber es ist ihm durch göttlichen Befehl geboten worden, daß er wieder in sein Vaterland ziehen sollte.

\subsection{II}

\begin{quote}
\textbf{V. 3. Und der Herr sprach zu Jakob: Zieh wieder in deiner Väter Land, und zu deiner Freundschaft; ich will mit dir sein.}
\end{quote}

15. Aus diesem Texte mag man sehen, wie viel Mühseligkeit und Trübsal Jakob erlitten habe, sintemal daß die Noth erfordert, daß ihn der Herr ansehen und von dieser Angst und Jammer abfordern muß. Denn Gott redet nicht vergeblich, und trägt sein Wort den sichern Herzen und die ohne Anfechtung sind, nicht vor, sondern denen, die es bedürfen, die Mangel haben, die mühselig und betrübt sind, die im schweren Kampfe liegen wider das Fleisch, und wider alles, das man sieht, das ist, wider alles, so nach menschlichem Sinne und Vernunft wider den Glauben streitet, Hebr. 11., auf daß die betrübten Herzen sich an den einigen Trost des göttlichen Wortes halten und damit auch erhalten werden mögen. „Die Weisheit wird nicht funden im Lande der Lebendigen“, spricht Hiob im 28. Cap. V. 13. Denn es ist ein Wort des Creuzes, und da man gar verzagt und Alles verloren gibt. Darum gehört es nur für die, so geplagt, getödtet und verachtet worden, die verzagt und gar beschweret sind, wie Jakob in seinem Elende gewesen ist bei seinem ungerechten und geizigen Schwiegervater, der ihm nichts geben wollte, und versucht hat, ihm Alles zu nehmen, und ihn gar zu schinden und zu berauben, daß Kinder ihn auch heftig vexiret und geplaget haben, und alle damit umgegangen sind, daß sie ihn für und für in der Knechtschaft möchten behalten, daß er ihnen müsse verpflichtet sein und daß sie durch seinen Schweiß, Mühe und Arbeit nur reich werden möchten.

16. In solcher Angst redet der Herr mit ihm, da er nichts Anderes mehr weiß, daran er sich halten möge, denn nur den Glauben und das Gebet; und da er von allen Menschen gehindert und gehaßt wird, nimmt sich der Herr seiner an, und stärkt ihn in seinem Vornehmen, daß er hat wegziehen wollen, und gebietet ihm, daß er das Gesinde und das Haus verlassen solle, welches ihm und seinen Weibern so gar feind war, und daß er sein eigen Haus versorgen solle. Mit welchen Worten GOtt anzeigt und bezeugt, daß ihm das ganze Leben und alle Werke des Patriarchen Jakob wohlgefallen.

17. Derhalben soll man diesen Text fleißig merken. Denn die heilige Schrift predigt nicht allein von großen hohen Wunderwerken, sondern bezeuget auch, daß GOtt für die geringen Dinge und Sachen, das Hausregiment belangend, auch Sorge trage; welches gleichsam eines der größten Wunderwerke ist und ein Zeichen einer sonderlichen Gnade und Gunst GOttes. Und ist dieser Trost droben und sonst auch in der ganzen heiligen Schrift oftmals wiederholet worden, und soll noch oft getrieben werden; denn wer da glaubt wie Jakob, der soll dessen gewiß sein, daß GOtt für alle seine Gedanken und für alle seine Fußtapfen Sorge trage nichts anders, denn als hätte GOtt sonst nichts mehr zu schaffen, denn daß er ihn schütze und regiere. Denn Mose predigt hier nichts vom Königreich Canaan oder Babel, auch nichts von der ganzen Welt: er erzählt nur, wie Jakob geplagt und angefochten worden sei, und wie er sich gerüstet habe, wiederum in sein Vaterland zu ziehen. Darum dürfen wir uns nicht ängsten und bekümmern, auf welcherlei Weise die Königreiche und Herrschaften der Welt regieret werden; sondern wir sollen GOtt dafür danken, daß wir wissen, daß er so genau für uns sorgt, daß er auch alles weiß, was wir thun, leiden und denken.

18. Dessen müssen aber alle Menschen insonderheit gewiß sein, daß sie GOttes Volk oder Glieder der Kirche sind. Dieser Glaube ist vor allen Dingen nöthig, welcher diesen Synagogismus erst stark angreifen muß, daß er also schließen kann: Alles Volk GOttes ist gesegnet, heilig, GOtt angenehm und gefällig, also daß es aus der Hand GOttes nicht kann gerissen werden; wir aber sind GOttes Volk: darum folgt gewiß, daß auch GOtt für uns sorgt. Der Obersatz ist ganz gewiß und wahrhaftig. Denn es ist auch der Tod und das Blut der Heiligen vor GOtt werth gehalten, und alles, was sie thun und leiden, gefällt GOtt wohl. Ja, daß noch mehr ist, es sind auch ihre Fehler und Sünden zugedeckt und vergeben, wie der 32. Psalm V. 1. zeuget.

\quellseite{168}

19. Aber von dem Untersatz ist noch die Frage. Denn der Türke, der Papst und die Juden rühmen sich dessen auch, daß sie Gottes Volk, und ihm wohlgefällig und angenehm seien. Derselben muß man darauf Acht geben, daß wir deßelbigen Untersatzes mögen gewiß sein, nämlich, daß wir Gottes Volk sind, daß wir unter die Heerde und Volk Gottes gehören, unter den einigen Hirten Christum. Wenn du das fest und gewiß schließen kannst, so bist du schon selig. Darum sei getrost, fürchte dich oder erschrick nicht, wenn auch die ganze Welt brechen und über einen Haufen fallen sollte; sintemal du dessen gewiß bist, daß du unter dem Hirten bist, der ein Herr ist Himmels und der Erde.

20. Es ist aber ganz schwer, auch den Schäflein Christi und den Gottseligen selbst, daß sie diesem Satze Beifall geben. Denn man muß stets im Kampfe liegen wider den Zweifel und Unglauben: so ein groß und schwer Ding ist es um den Glauben. Aber der Türke, die Juden und Papisten fühlen diesen Kampf nicht, sie streiten auch nicht, sondern wollen deß gar gewiß sein und rühmen sich mit aller Sicherheit, daß sie ohne allen Zweifel Gottes Volk seien. Darum bedürfen die Gottseligen wohl, daß sie stets getröstet werden mögen, daß sie diesem Untersatze glauben können, nämlich, daß sie gewiß Gottes Volk sind. Denn also wird dir der Obersatz nicht fehlen. Denn der Heilige Geist bezeuget mit diesem und vielen andern Exempeln mehr, daß Gottes Volk auch in den geringsten und schlechtesten Dingen und Werken Gott wohl gefalle. Denn da wird er Alles durch dich wirken, er wird durch dich die Kühe melken und die allergeringsten Knechtswerke durch dich ausrichten, und wird ihm alles, was du thust, gleich angenehm und gefällig sein, beide das Größte und das Kleinste.

21. Mit welchen Argumenten soll man aber den Untersatz beweisen? Wie kann ich gewiß wissen, daß ich unter Gottes Volk gehöre? Vor der Geburt Christi hatten die Väter die Verheißung, daß Wort und die Stimme Gottes; sie hatten dazu auch das Zeichen der Beschneidung. Wir haben jetzt die Schlüssel, die Taufe, das Abendmahl des Herrn und die Verheißungen des Evangeliums: daraus sollst du Argumente und Beweis nehmen, dadurch du gewiß werden mögest, daß du ein Christ und getauft seiest und daß du in einem heiligen und gottseligen Beruf lebest: siehe, daß du das Wort nicht lästerst, sondern dasselbe liebest, und Gott in irgend einem ehrlichen Amte dienest. Wenn das deinem Herzen also eingeprägt ist und darin fest steckt, wiewohl du wider den Zweifel streiten mußt, sollst du doch daraus abnehmen und schließen, daß du der seiest, deß Werke alle Gott wohlgefallen, und daß er dafür auch Sorge trage. Sollst derhalben also denken: Ich will meinem Beruf folgen, ich will Knechtswerke, Herrenwerke und Kinderwerke thun mit Jakob; denn ich weiß, daß dieses alles vor Gott, unserm Herrn, ehrlich und werth gehalten wird.

22. Dieses muß man immer wiederholen um des stetigen Kampfes willen, den man halten muß wider den Zweifel. Denn unsern Sinnen und Vernunft tritt oft das Gegentheil entgegen, anders, als das Wort und der göttliche Trost lautet. Das Fleisch fühlt viel ein Anderes, also daß sich die Gottseligen dünken lassen, daß Alles gar widerwärtig zugehe, sintemal die Gottlosen groß Glück haben, sie aber betrübt und elend sind. Darum gehört der Glaube auf die Dinge, so man nicht sieht, bis sich die Wahrheit der Verheißung am Ende selbst ausweist. Und um der Ursache willen redet der Herr abermals mit Jakob, nicht allein daß er ihn trösten, sondern auch stärken und ihm helfen wolle in dem Kampfe, darin er ohne Zweifel hart gelegen hat.

23. Man soll aber fleißig darauf merken, was der Herr mit Jakob redet: „Zeuch“, spricht er, „wieder in deiner Väter Land.“ Als sollte er sagen: Du hast lange genug und viel gelitten bei dem geizigen Buben, der dich sehr beraubt und geschunden hat. Derhalben ziehe nun hin und werde in deinem Vaterlande ein Hausvater. Dieses ist die Stimme Gottes, damit er Jakob tröstet und aufrichtet. Da er aber den Trost empfangen hat, bekommt er Muth, und darf nun das thun, das er sonst nimmermehr hätte wagen dürfen. Er ist nun so kühn, daß er heimlich, ohne Wissen seines Schwiegervaters und desselben Kinder, getrost aus dem Hause hat dürfen ziehen mit Weibern und Kindern, und allem, was er hat.

24. Aber dieweil ihm das Wort solches gebietet zu thun, darum fürchtet er sich nicht, sondern zieht fort im Namen des Herrn. Denn dieses sind Gottes Worte, der Jakob von der Angst erlöset, und ihm verheißt, daß er ihn hinfort schützen wolle. Gehe hin, sagt der Herr, und nimm alles, was du hast, und ziehe in dein Vaterland. Ja, es wird mir aber Laban nachfolgen, hätte Jakob sagen mögen, und mich wiederum in die Knechtschaft ziehen; denn er ist stark und gewaltig, hat viel Freunde und Verwandte, ich werde nicht heimlich oder unverletzt davon kommen können. Wie es hernach geschehen ist, da ihm Laban nachgeeilet als einem, der da flüchtig geworden war, und als einem Diebe, und ihm droht, er wolle ihn erwürgen. Darum spricht nun der Herr: „Ich will mit dir sein“; wenn schon alle Teufel mit Laban wären wider dich, so will ich doch mit dir sein; laß ihn nur herkommen.

25. Es ist wahrlich ein groß Ding, daß Jakob diesem Wort hat glauben können. Er ist aber ein wohlgeübter Mensch gewesen im Glauben und ist durch viel Anfechtung geprüft worden, und hat doch gleichwohl auch seine Schwachheit gehabt. Wie die andern Väter durch so viel Anfechtung und Kampf bisweilen sind müde geworden: sie haben aber immer wiederum Muth gefaßt und sind stärker geworden, und hat diese Veränderung stets mit ihnen gewähret, daß sie jetzt angefochten und dann wiederum getröstet worden sind.

\section[Zweiter Theil]{Zweiter Theil. Wie Jakob von seinem Vorhaben mit seinen Weibern redet, und wie sich die Weiber dabei verhalten}

\subsection{I}

\begin{quote}
\textbf{V. 4—9. Da sandte Jakob hin, und ließ rufen Rahel und Lea aufs Feld bei seine Heerde, und sprach zu ihnen: Ich sehe eures Vaters Angesicht, daß es nicht gegen mich ist, wie gestern und ehegestern; aber der Gott meines Vaters ist mit mir gewesen. Und ihr wisset, daß ich aus allen meinen Kräften eurem Vater gedienet habe. Und er hat mich getäuschet, und nun zehnmal meinen Lohn verändert; aber Gott hat ihm nicht gestattet, daß er mir Schaden thäte. Wenn er sprach: Die bunten sollen dein Lohn sein; so trug die ganze Heerde bunte. Wenn er aber sprach: Die sprenglichten sollen dein Lohn sein; so trug die ganze Heerde sprenglichte. Also hat Gott die Güter eures Vaters ihm entwandt, und mir gegeben.}
\end{quote}

26. Mich wundert, was die Ursache gewesen sei, daß Jakob mit seinen Weibern nicht im Hause und in den Hütten gerathschlagt; denn er war ja der Schwiegersohn und der Töchter Mann. Derhalben hat er im Hause bisweilen mit dem Gesinde und seinen Weibern reden können und ihnen seine Absicht anzeigen. Vielleicht hat ihm Laban und seine Kinder das Haus verboten, daß er nicht in die Hütten seiner Weiber hat gehen dürfen; darum hat er sie müssen rufen lassen, daß sie zu ihm auf das Feld hinaus gekommen sind, da er seine und seines Schwiegervaters Heerden geweidet hat. Und Laban wird ohne Zweifel argwöhnisch gewesen sein; darum hat er nicht wohl leiden können, daß Jakob mit seinen Weibern redet, auf daß er sie nicht aus dem Haus wegführe. Es sei nun gewesen, was es wolle, so läßt es sich ansehen, daß Mose damit habe anzeigen wollen, daß Laban ein böses argwöhnisches Herz gehabt habe.

27. Wie ist aber das so eine große Bosheit gewesen, daß er so bald den Vertrag aufhebt, da er sieht, daß Jakob am Gut etwas zunimmt? Was können wir wohl denken, daß er die vorigen vierzehn Jahre geraubt und an sich gerissen habe, dieweil er jetzt so gar häßlich und mit Gewalt den Vertrag und Bund verändert? Daselbstige erzählt der Heilige Geist mit Fleiß also, die Bosheit und den Geiz Labans damit groß zu machen. Denn er hat den Vertrag, so sie mit einander gemacht, verändert, wann und wie oft ihm das gelüstet hat. Wenn er gesehen, daß die Heerden bunte Lämmer getragen haben, hat er, wenn die Lämmer waren jung geworden, gesagt: Ich will haben, daß hinfort die bunten mir zufallen sollen. Wiederum aber, wenn er gesehen hat, daß einfarbige sind gezeugt worden, da Jakob keine Stäbe gelegt hatte, hat er die weißen oder schwarzen Schafe gewählet. Also hat er Jakobs Kunst getäuscht und hat ihn alle Jahre betrogen; ja, noch wohl alle Jahre, dieweil er in sechs Jahren Jakob seinen Lohn zehnmal verändert hat. Im ersten Jahre hat er den Vertrag gehalten; aber in den andern fünf Jahren hat er ihn zehnmal verändert.

\quellseite{169}

28. Das ist sehr unbillig gewesen. Darum hat Jakob gedacht: Was sollst du mit dem Menschen machen? Hier werde ich keinen Nutzen schaffen können; derhalben will ich davon ziehen und diesen höchst gottlosen Menschen fahren lassen. Und dies ist ein Exempel eines trefflichen Glaubens und großer Geduld an Jakob, der dieses hat leiden und ertragen können. Es wäre ihm gar nichts nütze gewesen, weder der Vertrag mit seinem Schwiegervater, noch auch die Kunst, die er gebraucht hatte, da Laban so oft und nach seinem Gefallen den Vertrag verändert hat, wo nicht der göttliche Segen wäre dazu gekommen; wie er sagt: „Gott hat ihm nicht gestattet, daß er mir Schaden thäte“, Gott hat mich errettet und meine Heerde gemehret; ja, „er hat die Güter eures Vaters ihm entwandt, und mir gegeben.“

29. Und das Wort: „Er hat es ihm entwandt“, hat einen sonderlichen Nachdruck. Er sagt nicht einfach: Gott hat es ihm genommen; sondern sagt: „entwandt“. Wie im 142. Psalm V. 7. steht: „Errette mich von meinen Verfolgern.“ Also sagt Jakob auch: Gott hat mich errettet und mir geholfen. Und der Heilige Geist will damit zu verstehen geben, daß diese Güter mit allem Recht Jakob und seinen Weibern gehört haben, die ihnen durch Gottes Segen gegeben und mit ihrer Arbeit erworben sind. Denn er hatte ihm vierzehn Jahre gedient und war in diesen sechs Jahren zehnmal betrogen worden. Darum sagt die heilige Schrift, daß solches mit Gewalt geraubt und genommen sei, daß Jakob und seinem Gesinde gebühret hat. Gott hat mit Macht dazu thun müssen: meine Kunst und List, so mir von Gott gezeigt und bewilligt worden ist, würde nichts geholfen haben, wo Gott nicht mit sonderlicher Kraft und Gewalt das Gut eures Vaters errettet und los gemacht hätte, wie man einen Gefangenen los macht; welches Gut er uns zuvor mit Gewalt aus dem Rachen genommen hat. Dieses alles, das er gehabt hat, ist unser gewesen; und das wir jetzt haben, dasselbe habe ich mit meiner Kunst oder Geschicklichkeit nicht erworben, sondern wir haben es durch göttliche Kraft und Gewalt bekommen. Denn Gott hat es mit gewaltiger Hand diesem Mörder und Buben aus der Faust genommen. Das hat Mose mit dem Worte jazzel anzeigen wollen, nämlich, daß der Herr Jakob damit gesegnet, daß er Laban seine Güter entwandt und ihm dieselben gegeben habe. Denn also ist die Hand des Herrn mit seinen Heiligen, daß er sie schützen und erretten möge von dem Betrug und Bosheit der Gottlosen.

\begin{quote}
\textbf{V. 10–13. Denn wenn die Zeit des Laufs kam, hub ich meine Augen auf, und sahe im Traum, und siehe, die Böcke sprangen auf die sprenglichte, fleckichte und bunte Heerde. Und der Engel Gottes sprach zu mir im Traume: Jakob! Und ich antwortete: Hie bin ich. Er aber sprach: Heb auf deine Augen, und siehe, die Böcke springen auf die sprenglichte, fleckichte und bunte Heerde; denn ich habe alles gesehen, was dir Laban thut. Ich bin der Gott zu Bethel, da du den Stein gesalbet hast, und mir daselbst ein Gelübde gethan. Nun mach dich auf, und zeuch aus diesem Lande, und zeuch wieder in das Land deiner Freundschaft.}
\end{quote}

30. Diese Worte sind synedochisch geredet. Und dieweil Mose so kurz redet, läßt es sich ansehen, daß sie also lauten, als hätte Jakob im Traume nichts gesehen, denn nur die einfarbigen Böcke springen. Derhalben muß man den Text nicht collectivisch, sondern distributiv verstehen, das heißt: bei dem einen Laufe sah er die Böcke auf die Einfarbigen springen, bei dem andern Lauf auf die Mehrfarbigen. Und man muß den Theil für das Ganze nehmen. Denn er will alle Böcke damit verstanden haben unter Schafen und Ziegen: wenn er sagte, du sollst diesesmal die akudim, das ist, die Fleckichten haben, so trugen die Schafe akudim, fleckichte Lämmer, und also auch von den andern Arten. Denn es ist dem gleich, das er droben (V. 8.) gesagt hat: „Wenn er sprach: die bunten sollen dein Lohn sein; so trug die ganze Heerde bunte“; wenn er mir aber die schwarzen oder weißen Lämmer gab, so trug die Heerde auch schwarze oder weiße. Also will er damit anzeigen, daß Gott mit ihm gewesen sei, ihm geholfen und den Betrug und List Labans wunderbarer Weise gehindert habe.

31. Es läßt sich aber ansehen, daß Mose hiermit einen andern Beweis beibringen wolle, den Geiz an Laban zu beweisen, nämlich, daß der Geizhals Nabal, gleich als ob er ein besonderes Vorrecht gehabt hätte, und nicht von Vertrags wegen alle einfarbigen Lämmer und Ziegen an sich gerissen habe, in welchem Lauf sie auch geboren waren, und darnach von den bunten Lämmern Jakob nur einerlei Art gelassen habe. Wenn in diesem Lauf viel bunte gefallen sind, die er zuvor seinem Schwiegersohne gegeben hatte, so hat er in dem andern Lauf die bunten für sich behalten, wenn er gesehen hat, daß die Heerden viele derselben getragen hatten. Also hat er immer einen Wechsel und Aenderung gemacht, und in einem jeglichen Lauf, wenn die Heerden getragen hatten, hat er ihm nicht mehr denn einerlei Art von den bunten gegeben, die andern alle hat er behalten. Darum hat es einer sonderlichen Kraft Gottes bedurft, den groben Geiz Labans etwas zu hindern, welcher so groß gewesen, daß er auch alle tückische Räuberei übertritt.

32. Also hat mich aber Gott getröstet, spricht Jakob, daß er seinen Engel gesandt hat, der mich in meinem schweren Kampf und da ich der großen Bosheit willen meines Schwiegervaters bekümmert war, stärken und zu mir also sagen sollte: Laban betrügt dich und thut dir unrecht, dieweil er dir nichts denn nur die sprenglichten Lämmer läßt; aber der Herr hat dich gesegnet, daß alle Lämmer in der ganzen Heerde sprenglicht geworden sind. Item, im andern Jahr, hat er gesagt, sollten mir die bunten Lämmer zufallen; da hat mich abermals der Engel getröstet und gesagt: Die ganze Heerde wird jetzt bunte Lämmer tragen, auf daß die Eine Art, welche Laban dir allein bleiben lässet, an Zahl und an der Menge die andern übertreffen soll. Ja, daß noch mehr ist, es soll die Heerde eitel solche Lämmer tragen, wie er dir von der Einen Art der bunten gegeben hat.

33. Und wo Gott nicht auf solche Weise Jakob geholfen und seine Armut erstattet hätte, so würde ihm Laban gar nichts gelassen haben. Denn es ist ja eine große Tyrannei und wider alle Billigkeit, daß er ihm nur Eine Art von der ganzen Heerde erlaubt hat zu nehmen wider den Vertrag, den sie miteinander gemacht hatten. Derhalben wird hier abermal Laban abgemalt als ein greulicher Verfluchter Räuber und Tyrann wider seine Töchter, Kindeskinder und seinen Schwiegersohn.

34. Das hebräische Wort beruddim haben wir droben noch nicht gehabt; es kommt aber her vom Wort barad, das heißt, hageln. Aus welcher Ableitung leichtlich kann verstanden werden, was die Bedeutung dieses Wortes sei. Denn wie das Wort nekuddim heißt, das einen oder zwei Flecken hat; also bedeutet beruddim die Schafe, die mit vielen kleinen Flecken besprenket sind, gleichsam mit Flecken behagelt, daß ich also rede, auf Deutsch, geschäckt. Wie Virgil die Ziegenböcklein beschreibt, die da kleine weiße Flecken haben, da er sagt: Capreoli sparsis etiam nunc pellibus albi; das heißt: Die Ziegenböcklein sind auch jetzt bei besprenkelten Fellen noch weiß. Dasselbe Wort steht auch im Propheten Sacharja am 6. Capitel V. 2. 3. im Gesicht von den vier Wagen; da waren am ersten Wagen rothe Rosse, Füchse; am andern Wagen schwarze; am dritten Wagen weiße Rosse, auf Hebräisch lebamim; am vierten Wagen geschächte starke Rosse, auf Hebräisch beruddim. Im 1. Cap. V. 8. steht ein anderer Name, auf Hebräisch serukkim, vom Wort sorek, welches der allerbeste Weinreben heißt, ein brauner Weinreben. Serukkim sind die Rosse, die nicht roth oder schwarz sind, sondern fast tanienbraun, auf Lateinisch spadices, wie die Trauben sind an den Reben. Beruddim aber heißt auf Deutsch, apfelgrau, Schimmel; denn an ihren Lenden sieht man solche Flecken, wie Aepfel oder Spiegel gestaltet sind.

\subsection{II}

\begin{quote}
\textbf{V. 14–16. Da antwortete Rahel und Lea, und sprachen zu ihm: Wir haben doch kein Theil noch Erbe mehr in unsers Vaters Hause. Hat er uns doch gehalten als die Fremden; denn er hat uns verkauft, und unsern Lohn verzehret: darum hat Gott unserm Vater entwandt seinen Reichthum zu uns und unsern Kindern. Alles nun, was Gott dir gesagt hat, das thue.}
\end{quote}

35. Nach der Klage Jakobs folgt auch nun, was Rahel und Lea gesagt haben. Denn es ist ja gottselig und recht, daß eine Mutter sich ihrer Kinder annehme und dieselbigen versorge. Denn nach Gottes Gebot ist dies einer Hausmutter oder Eheweibes Amt, daß sie dem Mann helfe die Kinder ernähren und aufziehen, daß Haus zu regieren und versorgen. Es klagen aber diese zwey Weiber ganz erbärmlich über die Grausamkeit ihres Vaters, welche Weiber jetzt schon Mütter sind von zwölf Kindern, und werden doch gleichwohl von ihrem einzigen und ungerechten Vater alles dessen beraubt, was ihnen vonnöthen gewesen ist, eine solche große Familie zu ernähren und zu unterhalten. Vielleicht wird er ihnen auch das Brod und die Milch ungern und kümmerlich gegeben haben, dieweil er alles, was sie mit ihrer Arbeit erworben hatten, an sich reißt, und sie haben doch keine gelegene Zeit gehabt, sich über ihren Jammer und Elend zu beklagen. Denn sie haben ohne Wissen und Willen des Vaters gar nichts thun oder sagen dürfen.

\quellseite{170}

36. Darum sagen sie nun: „Wir haben doch kein Theil noch Erbe mehr in unseres Vaters Hause.“ Als wollten sie sagen: Was sollen wir hier machen? Wir haben eine große Familie, zwey Mägde und zwölf Kinder; unser Vater aber sorgt dafür gar nichts, wie diese möchten versorget werden. Er gibt uns auch nicht einen Heller, er geht mit uns um, als wären wir Fremde, die wir nichts mehr an den Gütern und Erbe unsers Vaters haben.

37. Dieses scheint aber der Ehrerbietung, so die Töchter dem Vater billig erzeigen sollten, gar zuwider zu sein. Denn es ist eine allzuharte und bittere Rede. Darauf muß man also antworten: Lieber, wem gebührt das Gut und Erbe im Hause? Ist es nicht wahr, daß es den Töchtern gebührt? Und wo dieselbigen ja nicht das ganze Erbe haben sollten, so müßten sie doch die Hoffnung haben, daß sie dann irgend einen Theil bekommen möchten. Nun hält sich aber der Vater, sagen diese Weiber, gegen uns also, daß wir von ihm auch gar nichts zu hoffen und zu erwarten haben: je länger wir hier bleiben, je härter und greulicher wir umgetrieben und geplagt werden mit solcher Dienstbarkeit, die für und für währet; welche derhalben auch so viel unerträglicher ist, weil er uns gar keinen Lohn gibt, sondern über das auch das Geld und Gut sich zueignet und verzehrt, das wir mit unserer Arbeit erworben haben. Und er beraubt uns nicht allein des Theiles und Erbes, so uns von Rechts wegen gebührt, sondern er hat uns auch verkauft, gleich als wären wir leibeigen, hat uns gar keine Aussteuer mitgegeben; sondern er fordert und nimmt noch uns, das wir gar erspart haben;

38. Dieses sind wahrlich erbärmliche Klagen, und läßt sich doch gleichwohl ansehen, als ob die Worte etwas härter seien, denn es Töchtern wohl geziemen wollte. Sie scheinen ihre Eltern nicht zu Ehren zu halten, dieweil sie von ihrem Vater solch gottlos greulich Ding reden, und das schier also lautet, als sollte er ein Kindermörder sein; denn es ist ja im vierten Gebot von Gott geboten, daß wir unsere Eltern ehren und sie mit Worten und Werken zieren sollen. Was könnte aber Greulicheres gesagt werden, denn dieses, daß diese Weiber klagen: Der Vater hält uns nicht wie seine Töchter; er hat die Kinder, die von seinem Fleisch und Leibe geboren sind, nicht lieb, will sie auch nicht kennen; er hat keine natürliche Neigung, wie die Eltern Neigung und natürliche Liebe zu ihren Kindern haben; welche Neigung allen Menschen von Natur eingepflanzt ist. Davon oftmals Exempel vorkommen in diesem Leben und in den Historien aller Völker: ja, es wird im ersten Buch der Könige am 3. Capitel V. 16. ff. gesagt von einer Hure, die sich mit einem andern Weibe vor dem Könige Salomo über ein Kind zankte: „Ihr mütterlich Herz entbrannte über ihren Sohn“; denn sie hat lieber die Frucht ihres Leibes entbehren wollen, denn daß sie sollte ihr Kind haben tödten sehen. Ja, solche Neigung ist sehr natürlich am Menschen; und das noch mehr ist, sie ist auch den unvernünftigen Thieren eingepflanzt; welche Thiere, je wilder und greulicher sie sind, je mehr sie wegen solcher Neigung und Liebe zu ihren Jungen toben und wüthen. Die Wölfe, Bären, Löwen und Schweine streiten heftig und ganz wüthend für ihre Jungen. So überaus kräftig ist die Natur.

39. Ist es aber nun nicht ein greulich Ding, daß der Geiz nicht ein unvernünftig Thier, sondern einen Menschen, der eine sehr köstliche Natur hat, also gar ändern soll, daß er aller natürlichen Neigung und Liebe gegen seine Kinder ganz und gar vergessen soll? Und dieses scheußliche Ungeheuer, der Geiz am Laban, hätte nicht heftiger können gegeißelt werden, denn da die guten frommen Matronen, seine eigenen Töchter, klagen, daß sie von diesem Vater geboren seien und er sie doch nicht wie seine Töchter gehalten habe. Denn damit wird angezeigt, daß der Geiz die natürliche Neigung und Liebe der Eltern gegen ihre Kinder ganz überwinde und vernichte. Und dasselbe beweisen sie damit klärlich, da sie sagen: „Er hat uns verkauft und unsern Lohn verzehret.“

40. Das ist wahrlich wider die Natur. Denn es ist ja natürlich, daß ein Vater seine Töchter, als die ein schwach weiblich Geschlecht sind, mit Nahrung und Kleidung versorge; und nicht allein damit, sondern auch mit der Aufsteuer, mit dem Erbe und anderer Nothdurft. Ja, es folgen auch die Vöglein und andere unvernünftige Thiere der Natur darin fein nach, machen Nester, schüben und nähren ihre Jungen. Das thut Laban nicht, sondern er hält seine Töchter wie fremde Mägd, daß sie klagen müssen, es haben es die Mägde besser denn sie. Denn es gebührt sich nach allen Rechten, beide göttlichen und menschlichen, daß man einer Magd für ihre Arbeit auch ihren Lohn und Speise gebe.

41. Uns, sagen sie, gibt man nichts, sondern wir werden gehalten nicht wie Mägd, sondern als ob wir gar fremd wären, gleich als ob er unser Vater nicht wäre und wir auch seine Töchter nicht wären. Da möchte fürwahr das vierte Gebot zu Ende sein. Denn was hat Laban väterliches an sich, das Ehren werth sei? Er ist eine Bestie, ein unvernünftig Thier und ein scheußlich Ungeheuer. Wie könnetst du aber den als einen Vater ehren, der kein Vater sein will, auch dich nicht würdiget, dich für seinen Sohn zu halten, sondern alle Vernunft und natürliche väterliche Neigung abgelegt hat und gar zu einem Tyrannen wird?

42. Diese Tyrannen aber haben die guten frommen Matronen ganze zwanzig Jahre lang also ertragen und leiden müssen. Ich hätte nimmermehr gemeint, daß ein Mensch seiner Natur nach in solche greuliche Unsinnigkeit sollte geraten können, wo es nicht die heilige Schrift allhier bezeugete. „Er hat uns verkauft“, sagen sie, das ist, er hat uns dir gegeben für deine Arbeit und alle Güter, so du ihm erworben hast, und uns ist von Gottes Gnade wohl geschehen, daß wir einen solchen Käufer gefunden haben, der uns versorgt, nährt und schützt. Aber verflucht sei, der uns verkauft hat. Lieber Jakob, du hast um unserwillen dienen müssen, und alles, was du in vierzehn Jahren mit deiner sehr sauren Arbeit erworben hast, das hast du müssen für uns geben. Wenn wir nicht einen solchen gotteslästerlichen und boshaftigen Käufer angetroffen hätten, so hätten wir müssen Hungers sterben, oder aber es hätte ihn des Dienstes und Zornmuth verdrossen und wäre davon gelaufen, hätte uns und seinen Schwiegervater sitzen lassen.

43. Dieses ist wahrlich eine sehr harte und unerhörte Klage, dergleichen Exempel du kaum finden wirst. Denn es ist ein unglaublich Ding, daß ein Mensch also seine Natur ändern und nicht nur zum unvernünftigen Thiere, sondern gar zum scheußlichen Ungeheuer und Naturwunder werden sollte. Das ist eben, was man gemeiniglich pflegt zu sagen: Wenn ein Engel ein Teufel wird, so wird er gar böse; und daß der Teufel eben darum, daß er von solcher edlen und englischen Natur abgefallen, so gar verderbt und böse geworden ist. Also thut auch eine züchtige Jungfrau, wenn sie einmal ihre jungfräuliche Keuschheit verscherzt hat, so scheut sie sich vor nichts mehr, sondern wird gar unverschämt, allerlei Schande zu treiben. Also, wenn aus einem Christen ein Abtrünniger wird, so wird kein größerer Feind der Kirche und der wahren Religion sein, denn er ist.

44. Auf dieselbige Weise ist es mit Laban auch gegangen. Da er sich dem Geiz und Mammonsdienst ergeben hat, da hat er das gemeine Gefühl verloren, welches nicht allein die Menschen, sondern auch alle unvernünftigen Thiere an sich haben, und das in die Natur gepflanzet ist. Aus einem Vater ist er ein Tyrann und Mörder seiner eigenen Töchter geworden. Denn sie sagen: Er hat uns gar nichts gegeben, nicht ein Härlein, nicht einen Bissen Brodes, und den Lohn, den er dir versprochen hat, hat er zehenmal verändert. Alles, was wir haben, was wir sind mit unserem Leib und Leben, das mögen wir alles ganz und gar unserem Käufer danken, der uns mit seiner Arbeit nährt und erhält. Ja, das noch mehr und greulicher ist, „er hat unser Gut verzehret.“ Unser Mann hat für uns seinen sehr schweren Dienst, den er ihm vierzehn Jahre geleistet hat, gegeben; aber daran hat er sich nicht genügen lassen, er hat über das auch die Frucht und den Lohn für unsere Arbeit einnehmen und verzehren müssen. Was uns gebühret hat für unsere und unseres Mannes Arbeit, dessen genießet er und wendet es zu seinem Nutzen.

\quellseite{171}

45. Nun magst du fragen: Ob diese Töchter Labans auch damit gesündigt haben, daß sie ihrem Vater also abgetreten und so gar viel Böses von ihm sagen. Ich will es zwar nicht loben, sondern will gern zugeben, daß an ihnen eine Schwachheit und Ungeduld sei, ja, daß es eine Sünde der Heiligen sei, aber die doch gleichwohl vergeben und geduldet wird. Denn was könnte man doch Gutes sagen von einem Hunde? Ich würde möglich selbst auch sagen: Wie hündisch hält mich mein Vater, wie wölfisch! etc.; er behandelt mich nicht anders denn als einen Hund. Solcher hündischen Wüterei kann man nichts Menschliches nachsagen.

46. Danach haben sie gleichwohl bei fremden Leuten dieses nicht ausgebreitet, sondern haben es nur ihrem Manne heimlich angezeigt; vor demselbigen haben sie zwar ohne Sünde sich beklagen mögen der Beschwerung halben, so sie gehabt haben. Es ist wahrlich zu erbarmen, daß unser Vater alle freundliche und väterliche Liebe gegen uns so gar abgelegt hat, sagen sie, daß er nicht allein nach uns nichts fragt, sondern auch das hinweg nimmt und verzehrt, was wir mit unserer Arbeit erworben haben. Diese Laster wird niemand recht entschuldigen, loben oder vertheidigen können, sondern er wird sie viel billiger strafen und schelten müssen, auf daß Andere dieses schändliche Exempel erkennen und fliehen lernen.

47. Wiewohl wir nun bekennen müssen, daß diese Erzählung, da die Töchter Labans über die Laster ihres Vaters so heftig klagen, eine Schwachheit sei des Fleisches und des schwachen weiblichen Geschlechtes: so ist doch gleichwohl, wenn man die Wahrheit sagen soll, Laban ein rechter Hund, und ist ein Exempel, das jedermann verfluchen und fliehen soll.

48. Wiederum aber trösten sich nun diese Töchter Labans in ihrer Traurigkeit, Jammer und Elende mit Gottes Gnade und Güte, dem sie das zuschreiben, daß er ihnen alle Habe und Gut, was sie jetzt haben, gegeben habe. Wenn nicht der Herr, sagen sie, mit uns gewesen wäre und uns geholfen hätte, so wären auf Erden keine elenderen Weiber denn wir. Diese Wohltat Gottes erkennen und rühmen sie mit dankbarem Herzen, nämlich, daß sie Gott in solchem großen Elende und Armuth nicht verlassen, sondern alles reichlich erstattet habe, was ihnen der Vater geraubt hatte. Gelobet sei Gott, sagen sie, der unsern Seufzen gehöret und unsere Thränen gezählet hat. Alle Güter, die Gott unserm Vater entwandt hat, sind jetzt unser und unserer Kinder.

49. Darum sagen sie nun ihre Meinung auf die Klage, so Jacob gethan hat über die Gewalt seines Schwiegervaters, und was er vorgenommen hatte mit der Reise: „Alles, was dir Gott gesagt hat, das thue.“ Als wollten sie sagen: Du darfst dich davor nicht fürchten, als ob du diese Güter mit Unrecht hättest und mit dir nähmest. Es ist eben, wie du gesagt hast: „Gott hat die Güter eures Vaters ihm entwandt, und mir gegeben.“ Also wiederholen sie eben dasselbe Wort, das Jacob droben (V. 9.) gebraucht hat. Als wollten sie sagen: Diese Güter, die dir Gott durch ein sonderlich Wunderwerk gegeben hat, sind gefangen gewesen unter der Hand und Gewalt des Mammons und des Geizes unsers Vaters; er ist sie uns auch göttlichem und menschlichem Rechte schuldig gewesen: aber durch seine Bosheit ist es also geschehen, daß bisher weder du noch wir dieselbigen gehabt haben, wiewohl du sie mit deiner Arbeit diese zwanzig Jahre über erworben hattest, bis endlich Gott gekommen ist und diese Güter losgemacht oder sie dem geizigen unrechten Herrn entwandt und dieselbigen uns gegeben hat.

50. Es reimet sich aber dieses Exempel sehr fein mit der That, so die Israeliten begangen, da sie die Egypter beraubt haben. Denn das hat ihnen der Herr geboten, wie 2 Mos. 11, 2. 3. steht: „So sage nun vor dem Volk, daß ein jeglicher von seinem Nächsten, und eine jegliche von ihrer Nächsten silberne und güldene Gefäße fordere; denn der Herr wird dem Volk Gnade geben vor den Egyptern“, daß sie ihnen die Gefäße leihen werden. Und also haben sie die Egypter beraubt, nämlich, da der Herr die Güter, die gleichsam gefangen waren, das Silber und Gold aus den Egyptern losgemacht und von den Egyptern entwandt hat. Sie haben sie beraubt nicht aus eigener Lust und Begierde, sondern nach göttlichem Willen und Befehle, der den Egyptern die Güter entwandt und sie derselben beraubt hat, da er durch Mose geboten, daß sie solche Güter von den Egyptern fordern und damit wegziehen sollten. Denn sie sind euch euren Lohn schuldig, sagt er, für eure Arbeit, die ihr so viele Jahre lang gethan habt. Ihr habt ihnen umsonst gedienet und sie haben euch dagegen nichts wieder gegeben: darum will ich es euch nun vergelten, und was sie euch genommen und vorbehalten haben, das entwende ich ihnen und gebe es euch wieder als euren verdienten Lohn.

51. Und dieses ist ein Abbild aller der Wohltaten, die Gott seiner Kirche und Gemeinde erzeiget, daß er sie nähre und erhalte. Der arme Lazarus liegt vor des reichen Mannes Thüre und leidet Hunger. Also wird die arme Kirche in der Welt auch geplagt und betrübt, und muß darben. Die frommen gottseligen Pfarrherren und Kirchendiener müssen mit Weib und Kindern Hunger leiden, oder wenn sie etwas haben, so ist der Neid und die Bosheit der Menschen so groß, daß sie ihnen dasselbige begehren zu nehmen und zu entwenden mit Recht und Unrecht, wie sie nur können und mögen. Und oftmals, wenn sie das irgend mit einem Scheine thun können, berauben sie auch die Kirchendiener, und nehmen die geistlichen Güter, so zu den Pfarren gehören, weg, daß man mit Wahrheit sagen mag: Alles, was die Kirche hat, das hat sie wider des Teufels und der Welt Willen, und daß ihr der Teufel und die Welt darin widerstreben, und sie mit Gewalt, mit List und Betrug hindern, wo sie nur können und mögen. Was sie hat, ist wie ein Raub, den Gott der Welt mit Gewalt nimmt. Denn die Bürger und Bauern sind nicht werth, daß sie sollen milde Hände haben, das Kirchen- und Schulamt zu nähren und zu erhalten; sondern alle Güter, die sie dazu geben, sind wie ein Raub, den Gott den Wölfen aus dem Rachen und den Löwen aus den Klauen gerissen und genommen hat.

52. Wir sehen, wie gar übel und schmählich die Obrigkeit und ihre Diener, als, Steuerbeamten und Amtleute, mit den Pfarrherren umgehen: wenn sie sich ihrer Noth und Armuth halben beklagen, alsbald werfen sie ihnen vor und sagen: Die Pfaffen sind geizig. Und wo Einige an Fürstenhöfen und in Städten neue Künste entdecken und finden können, die Pfarren zu berauben und die armen Priester oder Kirchendiener zu schinden, die werden gelobt deswegen, daß sie so sparsam sind und Alles wohl können zu Rathe halten. Und solche Gesellen können zu großen Ehren nur darum, daß sie die Kirchendiener schelten und drücken können.

53. Wenn irgend die Reichen ermahnt werden, daß sie armen Schülern zu ihrem Studium helfen wollen, hilft Gott, wie hat es damit so große Mühe und Arbeit! Da sollte Einer wohl bedürfen, daß er so beredt wäre, wie Perikles gewesen ist, daß du nur ein wenig erlangen mögest, welches doch noch ungern genug, und mit Widerwillen und großem Stolz gegeben wird. Darum mögen wir billig sagen, daß wir vom Raube leben, damit Gott die Welt wider ihren Willen und da sie ihm hart widerstrebt, beraubt. Denn die Leute sind nicht werth, daß sie dazu helfen und steuern sollen, die arme, elende, betrübte Kirche damit zu nähren und zu erhalten.

54. Wir erfahren zwar jetzt zur Zeit, daß etliche Fürsten mild und wohlthätig sind, die gern geben, was sie nur geben können; wie der Durchlauchtigste, Hochgeborne Fürst, Herzog Johann Friedrich, Churfürst von Sachsen, thut. Und sind doch gleichwohl in diesem und anderen Landen viel Raubvögel, nämlich, die Amtleute und Steuereinnehmer, die so neidisch und böse sind, daß du kümmerlich aus ihren Händen mit großer Mühe und Arbeit erlangen kannst, was der Fürst mit seiner milden Hand gegeben hat, die Kirchendiener damit zu unterhalten. Also leben wir einfach vom Raube, nicht nach unserer Lust und Begierde, sondern nach göttlicher Gnade und Mildigkeit, der solchen Raub wunderbarer Weise aus den Händen der Raubvögel reißt, auf daß die Kirchen und Schulen mögen erhalten werden, und daß die Diener derselben nicht gar verschmachten und Hungers sterben. Das ist nun ein göttlicher Raub, wie an diesem Ort diese Matronen von ihren Gütern sagen.

55. Wer aber vom Kirchen- oder Predigtamt will reich werden, der darf die Laster und Sünden nicht strafen oder schelten, weder des gemeinen Volkes, noch auch der großen Herren, Steuereinnehmer und Amtleute; sondern er muß ihrem Geize heucheln und schmeicheln können, und muß sagen, was man gern hört. Aber was fromme reine Lehrer sind, die der Kirche oder Schule mit gutem und christlichem Eifer dienen, die werden gehaßt und verachtet von jedermann, und werden nimmermehr reich werden. Denn die Labansbrüder geben weder von ihrem, noch von fremdem Gut etwas.

56. Es kommt oft vor, daß irgend ein Abimelech den Abraham oder Isaak ernährt, wie wir droben (Cap. 20, 14. ff.) gehört haben, aber die Andern alle sind Feinde der Kirche. Es kommt auch oft vor, daß ein Land irgend einen oder den andern frommen und gottesfürchtigen Steuerbeamten hat; die andern aber sind Labaniten, die sich dünken lassen, daß sie Pfarrherren und Kirchendiener nicht werth sind, daß sie auch eine Stunde oder zwei leben oder einen Bissen Brods essen sollten. Und sie verhehlen auch ihre Urtheile nicht, sondern sie sagen öffentlich: Man könne in Dörfern und Städten keiner Säu- oder Schafhirten, oder auch seines Büttels, Stadtknechtes entrathen; aber eines Pfarrherrn und Schulmeisters bedürfe man gar nicht. Derhalben werden jetziger Zeit viel fromme Männer von Raubvögeln hart geplagt. Denn sie können oftmals ihre Besoldung nicht ohne große Mühe und Arbeit von denen erlangen, die dazu bestellt sind, daß sie die Kirchengüter einnehmen, dieweil ihnen dieselbigen auch das ungern geben, daß sie doch nicht dazu gegeben haben.

\quellseite{172}

57. Darum lasset uns nun lernen, daß wir solche Gewalt, Unrecht und Neid mit Geduld leiden und ertragen, und lasset uns die wunderbare Regierung Gottes erkennen, der uns speiset und ernähret, wie er Jakob ernährte, da er Laban gedient hat. Bisweilen gibt er auch wohl irgend einen frommen gottesfürchtigen Steuerbeamten oder Amtmann, der etwas Gutes thut, da uns sonst die andern alle neben und widerstreben.

58. Also erzählt nun Mose, wie diese zwei Weiber sich haben gefallen lassen, die Meinung Jakobs und daß Jakob seinem geizigen Schwiegervater die Güter entwandt hatte. Solche That, sage ich, haben sie gelobt, und haben ihren Mann ermahnt, daß er nur sollte fortfahren: „Alles, was dir Gott gesagt hat“, sprechen sie, „das thue.“ Daselbige ist auch aus sonderlicher Güte Gottes geschehen, Jakob zum Trost, nämlich, daß er die Einwilligung seiner lieben Weiber gehabt hat, die sonst keine andere Hülfe oder Trost mehr hatten denn nur ihren Mann. Und ihr Vater war also mit ihnen umgegangen, daß sie nichts lieber gewünscht und begehrt haben, denn daß sie nur aus des Vaters Hause ziehen möchten. Ihrer aller Wille und Gefühle stimmen überein wider den grausamen tyrannischen Vater, und freuen sich, daß sie Gelegenheit bekommen haben, daß sie mögen davon ziehen. Mein lieber Jakob, haben sie gesagt, wir haben schon vorlängst deinen elenden schweren Dienst gesehen, und haben begehrt, daß du uns viel eher solltest weggeführt haben; denn wir sind dessen gewiß, daß Gott mit dir ist, und was du gesagt hast, ist Gottes Wort und Befehl. Derhalben, so bald du immer kannst, siehe, daß du dem Befehl Gottes nachfolgest.

59. Es ist aber, wie ich zuvor auch gesagt habe, zumal ein gefährlich Ding, daß er aus seines Schwiegervaters Hause zieht und seine Weiber, Kinder und Heerde mitnimmt wider den Willen eines solchen Tyrannen. Und Jakob hätte gewiß solches sich nimmermehr unterstanden, wo er nicht das Wort gehabt hätte; wiewohl er vielleicht zuvor auch mit den Weibern wird gerathschlagt haben, wie sie doch aus solchem Gefängniß und Kerker erlöst werden möchten. Denn dieweil er gesehen hat, daß Laban den Vertrag so oftmals geändert hat, ist ihm eingefallen, daß er sich gefürchtet hat, daß er vielleicht bald hernach einen Anlaß oder Vorwand finden würde, daß er ihn alles seines Gutes berauben möchte; wie Jakob hernach anzeigen wird, da er sagt: „Du hättest mich leer lassen ziehen.“

60. Denn wie die neidischen und geizigen Menschen zu thun pflegen, also wird ohne Zweifel Laban oftmals bei sich selbst gedacht haben: Siehe, ich sehe wohl, daß dieser reich wird und an seinen Gütern zunimmt wider meinen Willen; und das geht heimlicher und wunderbarer Weise zu: ich werde etwas erdichten müssen, damit ich ihn in den vorigen Dienst bringen und diese Heerde mir zueignen möge. Und dieweil Laban solch groß Verlangen und Begierde gehabt, Jakob seiner Güter zu berauben, untersteht sich Jakob nicht ohne große Gefahr, aus diesem Gefängniß wegzuziehen, vornehmlich aber mit so großer Familie und so viel Hinderniß, daß er mit sich nimmt. Er wagt es aber im Namen des Herrn und seine Weiber halten es auch mit ihm um des Wortes willen; denn der Herr hat gesagt: Ich habe wohl gesehen, wie dir Laban Gewalt und Unrecht gethan hat; ich habe aber alle deine Thränen in einen Sack gesammlet: darum ziehe immer hinweg, ich habe dir es befohlen und will dich auch schützen.

61. An diesen Trost hat er sich sicher halten können und sich darauf auch gewißlich verlassen, daß sonst nach menschlichem Rath oder Gutdünken an sich selbst schwer und gefährlich war. Denn da ihm Laban nacheilte, hätte er ihm nicht allein alle Güter nehmen, sondern ihn auch noch dazu strafen können, darum daß er ihm die Töchter heimlich weggeführet hätte. Aber der Befehl Gottes bricht durch das alles hindurch, und macht, daß ihm seine Weiber mit gutem Willen folgen trotz aller Hindernisse. Hernach aber wird Laban toben und ihm sieben Tage nachjagen, aber Alles vergeblich. Jakob nimmt mit sich den Raub, den ihm Gott beschert hatte, darum ist der Abzug und die ganze Reise gut und heilig.

\section[Dritter Theil]{Dritter Theil. Wie sich Jakob auf die Reise macht und davon zieht; und wie Rahel Labans Götzen, Jakob aber dessen Herz stiehlt}

\begin{quote}
\textbf{V. 17, 18. Also machte sich Jakob auf, und lud seine Kinder und Weiber auf Kameele. Und führete weg alle sein Vieh, und alle seine Habe, die er in Mesopotamien erworben hatte, daß er käme zu Isaak, seinem Vater, ins Land Canaan.}
\end{quote}

62. Also bestätiget der Text abermal, daß die Güter Jakobs ein rechter göttlicher Raub und ein verdieneter Lohn für die Arbeit, so Jakob gethan hatte, gewesen sind. Denn er nennt es keinen Raub oder Diebstahl, sondern er heißt es sein Gut oder Habe, Kleider, Gold, Silber, Heerde. Dieses ist alles dein, haben die Weiber gesagt; es ist ein solcher Raub, der dir von Gott also beschert und gegeben ist, den du dem geizigen Höllenhunde (Cerberus) aus dem Rachen nimmst, als wäre es gestohlen und geraubt, wiewohl es dein eigen Hab und Gut ist.

\subsection{I}

63. Es setzt aber der Text hinzu: „Daß er käme zu Isaak, seinem Vater.“ Als wollte Mose sagen: Er hat nicht wollen fliehen als Einer, der ihm seine Töchter wollte entführen, oder sonst als ein Dieb oder Mörder; sondern er hat eine solche Reise vorgenommen, dazu er guten Fug und Recht gehabt hat, nämlich, daß er zu seinem Vater ziehen wollte, da er sicherer und friedlicher wohnen möchte, und Allen, die ihn vor Gericht beschuldigen und des Diebstahls halben verklagen wollten, es wäre gleich Laban oder Andere, wer sie sein möchten, Antwort und Bescheid von seiner Reise und auch von seinen Gütern geben könnte; sintemal es sich ansehen ließ, daß es ungerecht wäre, also heimlich wegzuziehen und wider Labans Willen alle sein Vieh und Güter wegzuführen.

64. Also können wir auch mit gutem Gewissen des Raubes dieser Welt gebrauchen, wiewohl wir die Speise den Raubvögeln, der Welt und Teufel, aus den Klauen nehmen müssen. Denn sie sind uns feind, und sagen, sie bedürfen keiner Pfarrherren oder Kirchendiener; Gott aber sagt dagegen also: Ich bedarf der Kirche und ihr seid meine Diener, derhalben leidet solch Unrecht und Gewalt um meinetwillen, ich will verschaffen, daß ich unter euch etlichen Raub austheile; ich will euch irgend einen frommen Fürsten geben, der euch aufnehmen und beherbergen soll; ich will euch den Erdboden oder die Welt nicht geben, wie ich dem Türken, desgleichen auch den Epicurern und Wucherern gegeben habe, denen ich die Güter dieser Welt vorwerfen und sie mästen muß zum ewigen höllischen Feuer; euch aber will ich ein klein Theil geben, das soll wie ein Raub sein. Und wenn euch das große Mühe und Arbeit kosten wird, so bedenkt, daß ich es der Welt mit Gewalt genommen habe, wie die Güter Jakobs und den Raub der Egypter.

\begin{quote}
\textbf{V. 19. Laban aber war gegangen seine Heerde zu scheeren.}
\end{quote}

65. Dieses soll man fleißig merken, daß Jakob Gott nicht versucht, sondern Gelegenheit sucht, daß er die Anstöße vermeide, denen er sich entziehen kann. Denn niemand soll sich Gefahr auf den Hals laden, sondern man soll sich vor Gefahr hüten. Und wiewohl wir Gottes Wort haben, darauf wir uns gewiß verlassen können und sollen, so legt uns doch dasselbe an sich selber viel größere Gefahr auf. Darum soll man keine andere unnöthige Gefahr auf sich laden, und wir sollen auch nicht faul, sicher und müßig sein, sondern sollen natürlicher Klugheit, Geschicklichkeit, dazu auch guten Rathes und menschlicher Hülfe gebrauchen, daß es nicht scheine, als wollten wir Gott versuchen.

66. Jakob hätte mögen sagen: Ich habe Gottes Wort und Befehl, darum will ich aus diesem Hause ziehen, und will nichts fragen nach dem Toben und grimmigen Zorn Labans, sondern will ihn lassen zusehen und toben, und sein Haus also verlassen und seine Töchter mit allen meinen Gütern wegführen: das wäre Vermessenheit, und unnöthige Kühnheit und Versuchung Gottes. Wenn du nun das Wort hast, so thust du recht, wo du demselben gehorchst und folgst. Du sollst dich aber also in die Sache schicken, daß du der Dinge gebrauchst, die vorhanden sind bei und neben dem Worte. Denn also hat Jakob ein Gebot, aus diesem Lande zu ziehen und wiederum in sein Vaterland zu reisen. Er hat auch eine Verheißung, daß ihn Gott schützen wolle, daß er sich nicht fürchten solle, wie große Gefahr und mancherlei Beschwerung, davon wir hernach hören werden, ihm auch begegnen möchten. Jedoch wartet er auf gute gelegene Zeit, da Laban nicht zu Hause ist, nämlich, da er drei Tage reisen weg gereiset war, seine Heerde zu scheeren. Und hat bei sich selber der Sache also nachgedacht: Er wird mir in diesen drei Tagen nicht nachjagen können, wird es auch nicht so leichtlich erfahren, daß ich weggezogen sei; unterdeß will ich so weit reisen, daß er uns am ersten oder anderen Tage nicht wird erreichen und Störung bereiten können.

\quellseite{173}

67. Also hat Jakob mit sonderlicher Klugheit und Ueberlegung auf eine sehr feine Gelegenheit gewartet. Denn darum hat uns Gott die Vernunft und alle Creaturen gegeben, dazu auch alle zeitlichen Güter, daß sie uns zu unserm Nutzen dienen sollen. Wer da reisen will, der muß Zehrung auf den Weg haben, davon er sich Nahrung und Herberge schaffen möge; er würde aber sehr närrisch thun, wenn er denken wollte, daß er keines Geldes oder Speise bedürfte, und ihm alleinshalben durch göttliche Versehung Alles würde begegnen und zufallen. Darum hat aber Gott alles geschaffen, so zu diesem Leben vonnöthen ist, nicht daß du es von ihm ohne Mittel erwarten sollst, sondern daß du der Dinge, so vorhanden sind, gebrauchen sollst, und dasselbe nach der Ordnung, die er selber dir vorgeschrieben hat. Darum soll man den Gebrauch und Dienst der Creaturen mit nichten verachten, sintemal sie Gott darum geschaffen hat, daß sie uns dienen sollen.

\subsection{II}

\subsection{Und Rahel stiehlt ihres Vaters Götzen}

68. Nun kommt ein zweites Verbrechen: Rahel stiehlt ihrem Vater seine Götzen; auf Hebräisch teraphim. Aber siehe doch, wie verständig das Weib ist. Denn sie hat also gedacht: Alles, was wir haben, hat uns der Vater genommen und vorbehalten, jetzt aber ist es ihm von Gott nun wieder genommen wie ein Raub und ist uns gegeben worden. Derhalben mögen wir alles mitnehmen, was er uns jetzt noch schuldig ist. Mir aber, als die ich ordentlicher Weise Braut geworden bin, sollte er ja billig ein golden oder silbern Kleinod gegeben haben; weil er mir aber das nicht hat gegeben, darum will ich nun hinnehmen, das er mir von selber und nach allem Recht hätte geben müssen.

69. So haben nun erstlich diese heiligen Weiber und Jakob sehr schwer wider das vierte Gebot gesündigt damit, daß sie ihren Vater und Schwiegervater beschuldiget und verklagt haben, als, daß er sich gegen seine Töchter nicht vermäßen gehalten hätte, wie es einem rechtschaffenen Vater geziemt und gebühret zu thun; welches sehr hart und schwer ist. Und läßt sich etlichermaßen ansehen, als ob Mose sich selber nicht eins sei, sintemal er im vierten Gebot sagt, daß man Vater und Mutter ehren und ihre Gebrechen und Schwachheit ertragen soll. Aber an diesem Orte erzählt er nun ein Exempel des Patriarchen Jakob und der Töchter Labans, das dem Gebot stracks zuwider ist. Denn diese Töchter scheinen aller Ehre, die sie ihrem Vater erzeigen sollten, gar vergessen zu haben, dieweil sie sich so gar hart und unfreundlich gegen ihn erzeigen mit der heftigen und harten Klage.

70. Die andere Sünde ist, daß Rahel ihres Vaters Götzen stiehlt. Denn der heilige Geist trägt keine Scheu, das Wort „stehlen“ zu brauchen, sondern sagt klarlich, daß Rahel, die sehr heilige Matrone und eines sehr heiligen Mannes Weib, die dem Worte Gottes glaubt und nach demselbigen ihren Mann ermahnt wegzuziehen, ihres Vaters Götzen (teraphim) gestohlen habe. Er sagt nicht nur, daß sie ihm dieselben genommen oder entwandt habe, sondern sagt frei heraus, sie habe es ihm gestohlen, wider das siebente Gebot, da Gott 2 Mos. 20, 15. sagt: „Du sollst nicht stehlen.“ Soll man aber nun dem Nächsten oder einem Fremden nichts stehlen, wie viel weniger soll Einer seinem Vater etwas stehlen! Und sie stiehlt dazu das, welches der Vater in seinem Hause für das Allerfestlichste und Liebste hält, nämlich, seine Götzen, denen er gedienet hat.

71. Wiewohl aber nun zuvor gesagt ist, daß an allen Heiligen noch viel Schwachheit sei, die ich nicht gern entschuldige, so muß man doch gleichwohl das noch hinzusetzen, daß Rahels Vater aufgehört hatte Vater zu sein, da er allerlei Grausamkeit und Tyrannei geübt hat. Denn er hat seine Töchter gehalten, als wären sie Fremde, und hat seinen Schwiegersohn so oftmals betrogen; ja, er hat darauf gedacht, wie er beide, Schwiegersohn und Töchter, behalten möchte, daß sie nur für und für ihm dienen sollten. Derhalben hat Rahel eine sehr gerechte Ursache gehabt dem weltlichen und natürlichen Rechte nach, die Götzen zu stehlen, wie droben von Jakob auch gesagt worden.

72. Nun kommt aber auch die erste Tafel dazu, die macht einen Unterschied der Gebote Gottes und des Gehorsams, nach dem Spruche Christi Matth. 10, 37.: „Wer Vater oder Mutter mehr liebet, denn mich, der ist mein nicht werth.“ Derhalben sind Jakob und Rahel nicht allein nach der andern Tafel gerecht, sondern auch nach der ersten. Denn wo man Gott in der andern Tafel soll gehorsam sein, alsdann hört der Gehorsam gegen die Eltern auf und man ist in solchem Falle an kein Gebot gebunden. Noch viel mehr aber, wo man dem Rechte der ersten Tafel folgt, das ist, wenn Gott etwas mit einem neuen Worte zu thun gebietet; wie diese heiligen Menschen das Wort Gottes gehabt haben, so er durch den Engel geredet hatte, und darum halten sich beide, Rahel und Lea, da sie sagen: Gott hat unserm Vater seinen Reichthum entwandt zu uns und unsern Kindern; item: „Was Gott dir gesagt hat, das thue“, nämlich, da er dir durch den Engel angezeigt, wie er alles gesehen hätte, was dir Laban gethan hat.

73. Es wird aber Mose im Segen Levi sagen, 5 Mos. 33, 9.: „Wer zu seinem Vater und zu seiner Mutter spricht: Ich sehe ihn nicht“ 2c.; „die halten deine Rede und bewahren deinen Bund.“ Das ist so viel gesagt: Wer Gott dienen will, der soll vor dem vierten Gebot, darin geboten wird die Eltern zu ehren, nicht gehorsam sein. Ei, möchte Einer fragen, was soll das sein? Warum hat denn Mose dieses Gebot gegeben, wenn er am anderen Ort das Gegentheil gebietet? Denn dieses ist ja stracks wider einander, und sind rechte widerwärtige Sprüche, die stracks wider einander lauten. Darauf antworte ich also: Wenn die erste Tafel gebietet, daß man Gott dienen und gehorsam sein soll, alsdann soll die andere Tafel der ersten weichen, und man soll zu Vater und Mutter sagen: „Ich kenne euch nicht, ich muß euch hierin nicht gehorsam sein etc.“

74. Darum schreien sie nun so heftig, daß auf solche Weise der Gehorsam im Welt- und Hausregiment zerrissen und aufgehoben werde, und muß derhalben das Evangelium so gelästert werden, daß man es eine aufrührerische Lehre nennt; wie das die Ueberschrift am Kreuze Christi ausweiset, Joh. 19, 19.: „Jesus von Nazareth, der Juden König“, das ist, ein aufrührerischer König. Und auf dieselbige Weise wird auch die Lehre dieses Königs aufrührerisch genannt wider die andere Tafel, dieweil sie sagt, daß man der Obrigkeit und Eltern nicht solle gehorsam sein; ja, sie gebietet, daß man übel und Böses von ihnen reden soll; wie Rahel von ihrem Vater sagt: Er ist unser Feind und Widerfacher. Wie könnte man aufrührerischer reden denn also? Soll ich denn der Obrigkeit und meinen Eltern nicht gehorsam sein? etc. Was wird da für eine Zertrennung und Unordnung mit allen Dingen werden!

75. Darauf antworte ich also: Man muß zwar den Gehorsam leisten, den Gott in der andern Tafel geboten hat. Aber es gibt verschiedene Stufen des Gehorsams und muß man auf solchen Unterschied fleißig Acht haben; denn wo die erste und andere Tafel wider einander gehen, alsdann ist das die einfältige und rechte Ordnung, da geboten wird, daß die andere Tafel der ersten weichen soll. Denn Gott ist der Schöpfer, das Haupt und der Herr über Vater und Mutter, über Welt- und Hausregiment. Dieses soll alles dem Schöpfer unterworfen sein, und wenn man fragt, welches man unter diesen zweien soll fahren lassen, den Schöpfer, oder die Creatur? so ist das meine Antwort, daß man die Creatur verlassen soll. Denn die erste Tafel geht vorher, und wenn man derselbigen genug gethan und dieselbe gehalten hat, alsdann hat die andere Tafel auch ihre Statt; dann sollst du deinen Eltern gehorsam sein, sollst erdulden und leiden, wo sie dir unrecht thun; aber um meinetwillen, spricht Gott, und nicht wider mich und wider die erste Tafel.

\quellseite{174}

76. So kann man nun sein einfältig und klar auf den Diebstahl der Rahel die Antwort geben: Daß die erste Tafel die andere aufhebe. Es wird in der andern Tafel recht geboten: „Du sollst nicht tödten“; wenn aber die erste Tafel gebietet, daß man tödten soll, so muß das Gebot der andern Tafel weichen. Also wird auch das siebente Gebot aufgehoben, daß man nicht stehlen soll, nämlich, wenn Gott gebietet, daß du hinnehmen sollst, was deines Vaters ist. Also wird im achten Gebot geboten, daß man von dem Nächsten und sonderlich von den Eltern nichts Böses reden und ihnen nicht afterreden soll: aber doch sofern solches nicht wider die erste Tafel sei; denn wo dich dieselbige es heißt, sollst du deine Eltern hassen und strafen, wie Rahel thut.

77. Dieser Unterschied ist sehr nöthig und nützlich. Es kommt aber oftmals vor, daß die Ketzer und Rotten dieses mißbrauchen, und sich rühmen, daß sie auch die erste Tafel haben, und mit solchem Schein die andere aufheben; da muß man des Dinges rechten Verstand haben. Denn also rühmt sich der Papst auch, er habe für sich die erste Tafel, und will daraus also schließen, daß niemand schuldig sei, der Obrigkeit Gehorsam zu leisten, sondern jedermann soll dem Stuhle zu Rom gehorsam sein, welchem auch Kaiser und Könige von Noth wegen unterworfen sein müssen. Siehe aber zu, ob er auch von der ersten Tafel recht lehre und warum er die andere Tafel aufhebe. Denn von den Mönchsorden haben sie auch gesagt, man soll zu Christo laufen, und Vater und Mutter versäumen und verachten; man soll die geistlichen Väter hören. Wie der Spruch des Hieronymus lautet: Wenn dir Vater und Mutter begegnen würden, wo du willst in ein Kloster gehen, und dir die Brüste zeigten, und wollten dich mit ihren Thränen wiederum zurückziehen: so sollst du ihre Thränen verachten und in solchem Falle deine Eltern mit Füßen treten, und sollst nackt und bloß zum Kreuze Christi fliehen.

78. Dieses ist gottlos und teuflisch geredet; denn sie haben die heilsame gesunde Lehre gemißbraucht, die Mönchsorden und Teufelslehren damit zu bestätigen. Es ist die Sache und die Lehre an sich selber recht und wahr: siehe aber zu, daß du ein guter Dialektiker seiest und die erste und andere Tafel recht unterscheiden könnest. Desgleichen siehe auch zu, ob die Mönchsorden wahrhaftig ein göttlich Ding, oder göttliche Sache seien, um deswillen die andere Tafel möge aufgehoben werden. Wo es ein göttlich Ding oder Sache ist, alsdann ist recht gesagt: Du sollst die Brüste deiner Mutter mit Füßen treten, sollst ihre Thränen verachten; wo es aber das nicht ist, so bedenke, daß du schuldig bist, deinen Eltern mit aller Ehrbietung gehorsam zu sein, und daß man mit ihren Gebrechen und Schwachheit Geduld tragen solle, wenn sie nur wider die erste Tafel nicht streiten.

79. Also hat Rahel die erste Tafel auf ihrer Seiten stehen; denn sie ist durch Gottes Wort dazu berufen worden, daß sie aus ihres Vaters Hause ziehen soll. Und der Vater sündigt nicht allein wider die andere Tafel, sondern auch wider die erste; denn er ist ein Götzendiener. Alles, was er bisher gethan, hat er nicht allein aus unmäßigem Geize gethan und daß er anderer Leute Gut begehrt, sondern er hat auch mit Abgötterei gesündigt; welche zwei die allergrößten und schwersten Sünden sind. Denn wiewohl der Geiz an sich selber Abgötterei ist, so ist doch noch der Götzendienst dazu gekommen.

80. Deshalb hat nun Rahel weislich also gedacht und geschlossen (denn sie ist ein Weib gewesen ungefähr von fünfzig Jahren): Ich bin durch die erste und andere Tafel dazu berufen worden, daß ich aus meines Vaters Hause ziehen soll, darum habe ich Gottes Wort, und der Vater sündigt schwerlich wider alle beide Tafeln; denn er ist ein Götzendiener und ein Geizhals; denn er hat mir noch keinen Heller gegeben von alle dem, das er mir nach allem Rechte, beide menschlichem und göttlichem, billig hätte geben sollen. Darum ist er mir meinen Lohn schuldig, und dazu auch noch Geschenk nach dem Rechte der andern Tafel und nicht allein der ersten; denn ich bin ja seine Tochter. Derhalben will ich ihm nun etwas stehlen, und dasselbe soll etwas sein, das ihm an seiner Abgötterei hinderlich bereiten soll. Ich will ihm seine Götzen stehlen; denn er ist allzu karg und so gar ein Filz, daß er nicht leichtlich andere Götzen gießen wird. Und also werde ich nach der ersten Tafel recht thun, wenn ich ihm die Götzen nehme und den falschen Gottesdienst umstoße, und nach der andern Tafel werde ich mich versorgen, damit ich hinweg nehme, das mir nach allem Recht gebühret. Und die heilige Schrift nennt es doch gleichwohl einen Diebstahl; wie Christus nennt Vater und Mutter hassen, da er sagt Luc. 14, 26.: „So jemand zu mir kommt und hasset nicht seinen Vater, Mutter, auch dazu sein eigen Leben“ etc., welches im Gesetze verboten ist; so doch dasselbe kein Haß ist.

81. In der andern Tafel sagt man recht also: Du sollst weder einen Andern, noch dich selbst hassen oder tödten; aber in der ersten Tafel gebietet Gott und sagt: Ich will haben, daß du dein Leben hassen sollst wider das Gebot der andern Tafel, da ich geboten habe, daß du dich lieben sollst. Also ist im fünften Gebot der Todtschlag verboten: aber Gott gebietet der Obrigkeit, daß sie das Schwert führen soll, in der ersten Tafel, da er sagt Röm. 12, 19.: „Die Rache ist mein, ich will vergelten.“ Derhalben tödtet die Obrigkeit nicht aus Kraft der andern, sondern der ersten Tafel; und solches ist kein unrechter Todtschlag, wiewohl es in Wahrheit ein Todtschlag ist. Also wird es an diesem Ort ein Diebstahl oder Raub genannt, so viel Laban und die andere Tafel betrifft: aber in der ersten Tafel ist es ein recht Werk, das sich wohl geziemet zu thun, und ist ein Lohn, der ihr gebühret; und wird doch gleichwohl ein Diebstahl genannt. Denn Laban hält es für einen Diebstahl, und wir können den Namen nicht ändern; denn die heilige Schrift redet also. Eben also ist auch der Raub, so die Israeliten den Egyptern geraubt haben, in der andern Tafel und so viel die Egypter belanget, ein Raub; aber in der ersten Tafel ist es eine Schuld, so die Egypter den Israeliten schuldig gewesen sind. Demnach hat nun Rahel nicht gestohlen, sondern hat heimlich ihren gebührlichen Lohn genommen.

82. Es bedarf übrigens einer fleißigen Unterscheidung und einer scharfen Dialektik, daß wir dieses nicht unter einander mengen, oder diesen Unterschied nicht mißbrauchen, wie der Papst und die Schwärmer thun. Denn dieselbigen maßen sich die erste Tafel zu großer Gefahr der Kirche und auch des Weltregiments an; wiewohl sie doch dieselbige verachten und sie in Wahrheit mit Füßen treten. Aber sie wollen derselben gebrauchen, daß sie dieselbe vorzuwenden haben, und unter solchem Scheine alle Gebote und Rechte, so zum Gehorsam gehören, beide gegen die Obrigkeit und Eltern, aufheben mögen. Der Papst schreibt: Ich bin der rechte Lehrer und Haupt der Kirche, darum soll man mir gehorsam sein; ich stehe in der ersten Tafel; ich habe Macht, alles aufzuheben, was in der andern Tafel ist.

83. Aber da bedarf es eines rechten Verstandes, daß man den Unterschied recht anmelde, und daß man die Gebote und Lehre des Papstes auf die Wage lege und sie gegen den rechten Verstand der ersten Tafel halte. Denn ich lehre auch also, daß unser Fürst dem Kaiser nicht soll gehorsam sein, und absolviere ihn von dem Gehorsam nicht allein gegen den Papst, sondern auch gegen das ganze Römische Reich und den Kaiser: aber mit welcherlei Recht? Denn so muß ich zusehen, daß ich den Gehorsam recht beschreibe, was er sei, und denselbigen zu unterscheiden wisse. Ich absolviere aber den Fürsten nach diesen Rechten: weil der Kaiser und Papst Gebote und Befehle ausgehen lassen wider die erste Tafel. Derhalben ist weder der Fürst, noch das gemeine Volk schuldig, solchen Geboten gehorsam zu sein wider den Gehorsam Gottes. Der Papst absolviert die Leute auch, daß sie dem Kaiser nicht müssen gehorsam sein, und hat oftmals die Kaiser von Reich und Kaiserthum entsetzt. Warum aber das? Ei, sagt er, darum, daß sie mir nicht gehorsam sind; ich sitze in der ersten Tafel. Aber er ist ein Heuchler, der nicht recht versteht und verfälscht die Lehre der ersten Tafel, und maßt sich fälschlich an, daß er nicht hat. Wir aber absolviren die Leute recht, nämlich, weil der Kaiser und Papst Gottes Wort nicht annehmen, sondern gebieten solch Ding, das stracks wider das Wort ist. Darum lehren wir, daß man solchen Gehorsam verfluchen soll.

84. Wir müssen aber lernen, daß wir den rechten Unterschied treffen zwischen dem Wort Gottes, damit Gott unmittelbar gerichtet wird, und zwischen dem Gehorsam gegen die Menschen, als da ist der Gehorsam gegen die Eltern, die Lehrmeister, die Herren, die Obrigkeit. Denn dieses sind Creaturen; die Kirche aber und Gemeinde Gottes hat keinen andern Lehrer denn allein Gott; darum hat sie auch keinen andern Gehorsam. Wenn nun der Schöpfer unmittelbar sagt: Dieses will ich von dir gethan haben, du sollst mir allein gehorsam sein, mich sollst du für deinen Gott halten, du sollst an Jesum Christum glauben, den sollst du hören: da soll man niemand anders hören. Wenn du das gehört hast, so steige auch hernieder zum Gehorsam gegen die Creaturen, und sei deinen Eltern und Obrigkeit gehorsam, welche aufhören und weichen müssen, wo man dem Schöpfer dienen soll. Sie bleiben aber fest stehen, wenn man dem Schöpfer gehorsam hat, was ihm gebühret.

\quellseite{175}

85. Auf solche Weise behält der Raub oder der Diebstahl der Rahel den Namen eines Diebstahls in der andern Tafel, aber in der ersten Tafel ist es ein rechter Lohn, so ihr für ihre Arbeit gebühret; und ist doch gleichwohl einerlei Ding. Die Rache oder Strafe, so die Obrigkeit übt, ist in der andern Tafel ein Todtschlag; aber in der ersten Tafel ist es eine Gerechtigkeit und ein gut Werk. In der andern Tafel ist es ein Ungehorsam, wo man den Eltern nicht gehorsam ist, und ihnen flucht oder übel nachredet: in der ersten Tafel ist es eine Ehre und Gottesdienst. Diesen Unterschied lehrt die Dialektik des heiligen Geistes. Darauf muß man aber nun fleißig Acht geben, daß wir Gottes Wort recht haben und nicht Menschensatzung unter dem Scheine und Namen des Wortes; wie der Papst die ganze Welt betrogen hat und das Regiment dieser Welt seiner Tyrannei unterworfen hat.

86. So ist nun Alles daran gelegen, daß wir Prediger oder Lehrer und Schüler haben, die dieses lernen, und man muß Gott bitten, daß er uns solche Pfarrherren und Seelenhirten geben wolle, wie er bisher gethan hat, die solchen Unterschied fein klärlich lehren und sich darin nach dem Lichte des Wortes richten. Denn was ist doch die Welt, wenn sie dieses himmlische und geistliche Licht nicht hat, das wir aus der Theologie haben. Die Lehre der Fürsten und Aerzte hat auch ihre Statt und Ehre, und ist sehr gut und nöthig, gibt dazu auch viel Geld; wenn sich aber jedermann auf diese Studien werfen würde, wo würde denn dieses himmlische Licht bleiben, wenn man keine Schüler und Zuhörer der Theologie haben würde?

87. Was hilft es, wenn man schon Gold, Silber und Gesundheit des Leibes hat; ja, was nützt uns dieses ganze Leben, wo dieses Licht des Wortes verloren ist? Nun sind aber jetzt fast Wenige, die darüber halten, da doch vonnöthen wäre, daß es Viele wären, die darauf Fleiß wenden möchten, und echt schmerzlich zu, daß man solch Licht kann behalten; denn der Teufel und die Welt sind ihm feind. Wie Jakob ein schöner Stern ist, der im Hause Labans geleuchtet hat mit rechter Erkenntniß und Lehre des Wortes; aber wir sehen wohl, wie Laban mit ihm umgegangen ist. Und liegt doch die ganze Welt in Finsterniß und greulicher Blindheit, wo es ohne dieses Licht ist. Was ist der Papst, was ist der Türke ohne dieses Licht? Laß sie reich, weise und sehr gewaltig sein; weil ihnen aber das Licht des Lebens mangelt, sind es die unglücklichsten Leute.

88. Derhalben sollen wir uns nun der heiligen Schrift befleißigen und uns in der Theologie wohl üben, daß wir Schulen und Pfarren behalten mögen, welches der höchste Gottesdienst ist in der ganzen Welt. Ob es übel zugeht, ob wir stehlen müssen, so mögen wir stehlen im Namen des Herrn. Und wo wir ja mit Rauben unser gebührlichen Lohn, ja, unsere Nahrung und was zu diesem Leben vonnöthen ist, der undankbaren Welt rauben müssen, so ist es doch ein Raub und Errettung solcher Dinge, die uns gebühren. Laßt das Andere den Laban und Nabal fressen, laß ihn ein gutes Jahr haben, wir leben vom Raube; welches doch kein Raub ist in der ersten Tafel, wiewohl die Bürger und Bauern es dafür halten, daß sie uns solches geben als Dieben und Räubern. Aber was sie davon halten, soll uns nicht anfechten. Wir sollen uns viel lieber darnach richten, daß St. Paulus Phil. 2, 15. 16. sagt: „Seid ohne Tadel und lauter, und Gottes Kinder, unsträflich mitten unter dem unschlächtigen und verkehrten Geschlechte, unter welchem ihr scheinet, als Lichter in der Welt; damit, daß ihr haltet ob dem Wort des Lebens.“ Und soll sich ein Jeder in seinem Stand befleißigen, daß er die Lehre und das Wort befördern und fortsetzen helfe um seiner Seligkeit und der Ehre Gottes willen; und soll zum Besten deuten diesen Raub, der gering genug ist, damit uns Gott in diesem Leben unterhalten will. Denn wo das Licht des Wortes verlöschen ist, da wird wiederum Alles voll greulicher Finsterniß, türkischer und papistischer, oder anderer Ketzer falscher Lehre werden.

89. Dieses sei gesagt von dem Diebstahl der Götzen, welche die fromme Rahel geraubt hat als einen Lohn für ihren Dienst oder anstatt ihres gebührlichen Erbes. Denn da war gar keine Hoffnung, daß sie auch irgend nur ein klein Theil von ihres Vaters Gütern mit gutem Willen des Vaters hätte bekommen mögen. Auf daß sie nun nicht gar leer und vom ganzen Erbe abgeschnitten aus ihres Vaters Hause ziehen möge, so nimmt sie mit sich die goldenen Götzen. Denn Laban hatte nichts Menschliches mehr an sich, sondern war gar ein unvernünftig reißend Thier geworden, der seinen Schwiegersohn und seine Töchter ganz unbilliger Weise beraubt hatte. Darum war es nicht unbillig, daß er wiederum beraubt wurde von denen, die zu seinem Gute rechte Erben und Besitzer waren.

90. Rahel hat aber dieses gethan, daß Jakob, ihr Mann, nichts davon gewußt hat; wie wir das hernach sehen werden. Denn er würde es wohl schwerlich haben geschehen lassen, wenn er es gewußt hätte. Sie darf es aber wagen, und verläßt sich auf das Wort, und daß sie es ihrem Vater auch wohl abverdient hatte, dem sie, da sie schon ehelich war, vierzehn Jahre gedienet hatte. Darum es ja billig gewesen wäre, daß sie freigebig hätte sollen beschenkt werden.

91. Weiter müssen wir auch nun etwas sagen von der Frage, die zur Grammatik gehört, welche an diesem Orte angeregt wird. Denn dieses ist der erste Text, darin das hebräische Wort teraphim gesetzt wird; und ist gewiß, daß es Götzen heißt. Denn Laban war ein Götzendiener. Es ist aber ein wunderlich Ding, daß Jakob, der fromme heilige Mann, mit seinen frommen heiligen Weibern und mit seiner ganzen Gemeinde in dem abgöttischen Hause gewohnt hat. Derhalben ist in dem Hause zugleich Gott und der Teufel gewesen, Gottes Kinder und Teufels Kinder. Welches oft pflegt also zu kommen. Und ist jetzt noch die ganze Welt zusammengesetzet aus den Kindern Gottes und des Teufels. Wie in Adams Hause Kain und Abel auch bei einander waren; in Abrahams Hause Ismael und Isaak; in Isaaks Hause Esau und Jakob. Und also wird bis ans Ende der Welt die falsche Kirche bei der rechten Kirche sein, die falschen Brüder bei den frommen gottseligen Brüdern. Wir lehren heutiges Tages das Evangelium lauter und rein, dazu mit aller Treue: wir haben aber noch unter uns Wucherer, Papisten, Ketzer und Schwärmer; denn das Unkraut bleibt stets unter dem reinen Weizen stehen: es kann oder mag in der Welt anders nicht zugehen. Jakob hat zu derselben Zeit die himmlische Lehre und die Verheißung vom Messias rein und lauter gelehrt; dieselbige Lehre haben seine Weiber und das Gesinde angenommen; und ist wohl glaublich, daß auch Etliche von der Abgötterei zum rechten wahren Gottesdienst sind bekehrt worden: aber die Kirche oder Gemeinde Jakobs ist in einem gottlosen abgöttischen Hause gewesen.

92. Aber von dem hebräischen Worte teraphim bringen die Juden viele Lügen vor. Dieselbigen wollen wir mit wenigen Worten berühren, damit es nicht scheine, als hätten wir von ihrem Narrenwerke gar nichts gewußt. Lyra beschreibt die teraphim also, daß es ein Haupt gewesen sei eines erstgebornen Knaben, der getödtet und dem Teufel geopfert worden und mit Salz und Würze oder Specerei besprenget gewesen sei, daß es nicht sollte faul werden; welches unter der Zunge eine goldene Spange soll gehabt haben, darauf der Name des Teufels oder Abgotts geschrieben gewesen, bei welchem Teufel man göttliche Antwort gesucht hat: und eben solche Auslegung gebraucht er auch von Micha im Buch der Richter am 17. Cap. 3.

93. Ich glaube aber der keines. Denn teraphim ist ein gemeines Wort in der heiligen Schrift und man soll die Auslegung desselbigen Wortes aus andern Orten der Schrift nehmen. Als, 1 Sam. 19, 13. steht von Michal, die an Davids Statt ein Bild (teraphim) in das Bett gelegt hat und ein Ziegenfell zu seinen Häupten, daß sie also damit ein Bild machte eines kranken Menschen, der zu Bette läge; da kann ja dieses Wort nicht verstanden werden für eines Kindes Haupt, sondern es ist ein Bild eines Menschen gewesen, welches an Davids Statt in das Bett gelegt worden. Item, im Propheten Hosea Cap. 3. V. 4. steht also: „Die Kinder Israel werden lange Zeit ohne König, ohne Fürsten, ohne Opfer, ohne Altar, ohne Leibrock und ohne Heiligtum (teraphim) bleiben.“ Da will es sich nicht reimen, wenn du das Wort teraphim wolltest verdolmetschen, daß es Götzen heißen sollte, als hätte das Volk GOttes Götzen gehabt.

\quellseite{176}

94. Es sei aber, was es wolle, so dünkt mich nach meinem Verstand und nach der Grammatik, daß teraphim ingemein heiße ein Gleichniß oder Bild des Gottesdienstes, er sei gleich recht oder falsch. Laban hat Götzenbilder gehabt. Hosea aber, wie es scheint, hat wollen alle Stücke oder Figuren des Gottesdienstes zusammen fassen, als, die Hütte, den Tisch, den Vorhang, den Leuchter, die Schaubrode etc., die sie alle haben entbehren müssen, da sie im Elend und in fremden Landen gefangen waren. Wie nun die Heiden ihre Götzen mit diesem Namen genannt haben, also zieht Hosea denselbigen Namen auf alles Geräthe der Hütte. Denn es will sich nicht reimen, daß der Prophet diese Strafe dahin sollte verstanden haben, daß sie keine Götzen gehabt hätten; sondern er will damit verstanden haben das heilige Geräth, wie wir zuvor gesagt haben; also daß er damit ingemein ein Bild des äußerlichen Gottesdienstes verstehe mit Zeichen und Figuren.

95. Das Bild aber, das Michal an Davids Statt in das Bett legte, halte ich, sei irgend ein Bild oder Götze gewesen, so noch irgend an einem Ort übergeblieben ist vom vorigen Aberglauben und Abgötterei. Wie bei uns noch viel Bilder übergeblieben sind, die an etlichen Orten noch bewahrt werden zum Gedächtniß der alten Abgötterei, daß man es irgend schauen und die Kinder mit spielen läßt, und nicht zum Gottesdienst. Solche Götzen haben bei Davids Zeiten in den Häusern hin und wieder in Winkeln verborgen gelegen; da ist Michal in der Eile zugefahren, da nichts Anderes vorhanden gewesen ist, und hat das Bild ergriffen und dasselbe an Davids Statt ins Bett gelegt.

96. Denn zu der Zeit sind ohne Zweifel auch noch viele Spuren geblieben von dem Aberglauben und der Abgötterei aus dem Lande und Religion der Zebusiter, und etliche Geschrey, die man nicht gut hat abschaffen können; wie das die ängstlichen Klagen in vielen Psalmen bezeugen, in welchen David mit großem Schmerz über die Abgötterei betet und streitet. Und der Prophet Micha sagt von der Stadt Lachis, daß sie ein Anfang gewesen sei der Sünde der Töchter Zions, und daß man darin die Uebertretung Israels gefunden habe, also auch, daß man unter dem Volke Gottes die Abgötterei nicht hat gar vertilgen können. Ja, siehe auch die Zeit Moses an in der Wüste. Denn davon sagt der Prophet Amos im 5. Cap. V. 25. 26. und St. Stephanus in den Geschichten der Apostel am 7. Cap. V. 42. 43. also: „Habt ihr vom Hause Israel mir in der Wüste die vierzig Jahr lang Schlachtopfer und Speisopfer geopfert? Ja wohl. Ihr trugtet die Hütten eures Molochs, und die Götzen eurer Bilder, und den Stern eurer Götter“ etc. Das nun solches von dem Volk unter Mose hat können gesagt werden, ist wahrlich ein wunderlich Ding. Dagegen aber spricht Bileam 4 Mos. 23, 23.: „Es ist kein Zauberer in Jakob, und kein Wahrsager in Israel.“ Aber das ist synedochisch geredet, wie das in der Schrift gewöhnlich ist.

97. Denn es ist allezeit der gute Same in der Kirche unter das Unkraut gemengt gewesen. Es ist aber daran genug, daß zur Zeit Davids das Königreich oder Regiment und das Kirchenamt rein, und Beides also verordnet und bestellt gewesen ist, daß darin nichts Gottloses, oder das abgöttisch wäre, gelehrt oder geboten worden ist. Also ist unter Mose in der Wüste die Lehre rein und die Obrigkeit fromm gewesen, und die nichts verordnet hat, denn allein was heilig und recht gewesen, ohne alles gottlose Wesen oder Abgötterei; daß auch derhalben Bileam 4 Mos. 23, 23. mit Recht gesagt hat: „Es ist kein Wahrsager in Israel“; item 4 Mos. 24, 5.: „Wie fein sind deine Hütten, Jakob, und deine Wohnungen, Israel!“ Daran ist aber kein Zweifel, daß noch bei etlichen Spuren der Abgötterei werden geblieben sein.

98. Eben auf dieselbige Weise, wenn ein Pfarrherr in unserer Kirche recht lehrt und wider die Abgötterei streitet, desgleichen auch die Laster und Sünden der Menschen straft, da kann ich mit Wahrheit sagen: Zu Wittenberg ist kein Abgott oder gottlos Wesen. Darum ist die Kirche heilig; denn das Predigtamt ist rein, und die weltliche Obrigkeit gebietet oder schützet auch die Abgötterei nicht. Wie Herzog Georg in seinem Lande pflegte zu gebieten, daß man die papistische Religion halten sollte, und die Pfarrherren waren ihm auch darin gehorsam; da war nun keine rechte oder reine Kirche.

99. Wo aber eine Kirche ist, darin die himmlische Lehre lauter und rein gelehret wird, und die Obrigkeit wider die reine gesunde Lehre nicht streitet, da sind die Hütten Jacobs rein und schön. Der Leib an sich selbst ist gesund und stark; er ist aber nicht ohne Unfall, Eiter, Schwären, Speichel, und sonst ohne Dreck und Mist; und wird doch gleichwohl synedochisch ein gesunder Leib genannt. Eben auf solche Weise bleiben auch teraphim in der Welt, ja auch wohl in der Kirche und im Hause Jakobs, bis zum Ende. Dies sei gesagt von der Bedeutung des Wortes teraphim.

100. Sie disputiren aber auch noch von seiner Ableitung. Denn Etliche sagen, es komme her von dem Wort rapha, mit dem Buchstaben a geschrieben, Etliche aber von rapha, mit dem Buchstaben e geschrieben. Jenes heißt heilen oder gesund machen. Daher kommt resphaim, das ist, große Riesen und Aerzte, als, wenn man sagt, Helfer oder Heiland, das ist, Götter. Dieses aber heißt aufhören, stille halten. Und mich dünkt es richtiger, daß teraphim vom Stille sein herkomme; und wird gezogen oder gedeutet auf den Sabbath oder sonst auf den Gottesdienst, daß sie haben müssen stille sein, GOtt dienen und den Sabbath halten, wie im 46. Psalm V. 11. steht: „Seid stille und erkennet, daß ich GOtt bin“ etc., das ist, gehorcht mir, hört auf mit eurem eigenen Vornehmen, laßt ab, laßt mich machen. Also haben die Götzen ihren Sabbath und Gottesdienst gehabt, von welchem Dienst sie teraphim sind genannt worden, daß man ihnen hat feiern müssen.

\subsection{III}

\begin{quote}
\textbf{V. 20. Also stahl Jakob dem Laban das Herz, damit, daß er ihm nicht ansagte, daß er flöhe.}
\end{quote}

101. Das ist nun eine andere Sünde, und ist ein viel größerer Diebstahl, denn Rahel's war. Denn Jakob stiehlt Laban seine Götzen oder Güter nicht, sondern das Herz. Dieses aber ist eine eigene Art zu reden in der hebräischen Sprache, „Einem das Herz stehlen.“ Und heißt so viel, als wenn Einem seine Hoffnung, so er gefaßt hat, fehlet, und daß man ihm allen seinen Rath und Anschlag gar zu nichte macht, daß er gar keine Hoffnung mehr haben kann, das zu genießen, worauf er in seinem Herzen gedacht hat.

102. Denn Laban hat bei sich selbst also gedacht und geschlossen: Siehe, ich habe da meinen Schwiegersohn mit meinen Töchtern, und mit Kindeskindern und mit Mägdgen, durch deren Fleiß und Arbeit will ich sein reich werden; denn sie sind mir nach allem Recht unterworfen, daß sie mir mein Hab und Gut mehren und bessern mögen. Ich habe auch nun erfahren, daß sie getreu sind; denn Jakob ist ein frommer und aufrichtiger Mann, der wird mir nicht Unrecht thun und wird auch seinen jetzigen Dienst oder Stand nicht leichtlich ändern ohne mein Wissen und Willen; wird es auch ohne Gefahr nicht thun können. Wo er mir aber etwas nehmen oder Schaden thun wird, alsdann werde ich erst Ursache haben, daß ich sie derhalben aufhalten werde, daß sie werden müssen verpflichtet sein, mir für und für zu dienen.

103. Dieses ist die große Hoffnung und sind die Gedanken Labans. Es sind sehr kluge und listige Anschläge; und wie er sich hat dünken lassen, sind sie ganz gewiß und haben ihm nicht fehlen können. Aber in einem Augenblick werden sie alsbald zu nichte und schwinden gar dahin. Das heißt die heilige Schrift „Einem das Herz stehlen“, das ist, Einem seine Hoffnung zu nichte machen, wie gewiß und ungezweifelt er sie gefaßt hat, und dasselbe noch ohne sein Wissen. Denn Laban hat sich vor der Dinge keinem gefürchtet, sondern hat sich selbst gewißlich damit getröstet, daß Jakob bei ihm bleiben würde, und daß er sich nichts weniger unterstehen könnte, denn davon zu ziehen. Und Jakob zieht doch gleichwohl darüber heimlich weg, nicht ohne großen Schaden und Unwillen Labans. Davon wird aber Laban ein lang Register vieler großen Sünden machen, die Jakob damit soll begangen haben; wie wir das hernach hören werden.

\begin{quote}
\textbf{V. 21. Also flohe er, und alles, was sein war, machte sich auf, und fuhr über das Wasser, und richtete sich nach dem Berge Gilead.}
\end{quote}

\quellseite{177}

104. Man soll fleißig merken, daß die heilige Schrift allen Werken Jakobs Namen der Laster und gröbsten Verbrechen gibt. Mose hat zuvor gesagt, Jakob hätte dem Laban zu bringen das Herz gestohlen; hier sagt er nun, er sei geflohen ohne Wissen und Willen seines Schwiegervaters und habe ihm seine Güter genommen, und dazu alle seine Hoffnung, daß Jakob ihm hätte für und für dienen sollen, zunichte gemacht. Diese Stücke alle werden also genannt, daß es Sünden gewesen sind, wie Laban dieselbigen bald hernach mit seinen geschmückten Worten vergrößern wird. Wie es Christus Luc. 14, 26. auch einen Haß nennt, damit Einer seinen Vater, Mutter, oder sich selbst auch haßt; welches ja ein Name einer Sünde, so es doch die höchste Tugend ist.

105. Also ist in Wahrheit diese That kein Diebstahl; aber die heilige Schrift redet, wie die Menschen pflegen zu reden, die gemeiniglich von solchen Thaten also pflegen zu urteilen. In den Ohren Labans ist es ein Diebstahl und eine Flucht; aber nach göttlichem Recht und Gerechtigkeit wird es für einen verdienten Lohn gehalten, den Jakob mit allem Recht gebühret hat, wie der, welcher seinen Vater und Mutter um Christi willen haßt, sie nicht recht haßt, sondern liebt: und wer sich selbst verleugnet und sein Leben verliert, der findet eben damit dasselbe, Matth. 10, 39.

106. Das Wasser aber, davon hier Meldung geschieht, ist der Euphrat, welcher Mesopotamien und das gelobte Land trennt. So weit ist nun Jakob gekommen, daß er jetzt schon über den Euphrat fährt, und hat nun große Hoffnung gefaßt, daß Laban ihn und seine Familie nimmermehr nicht werde einholen können. Aber der Geiz ist ein solch Ungeheuer, das nicht müde wird, ist schnell und hurtig, läßt auch Laban keine Ruhe. Da aber Jakob über das Wasser gefahren ist, zieht er weiter nach dem Berge Gilead.

\section[Vierter Theil]{Vierter Theil. Wie Laban Jakob verfolgt, und wie ihm der Herr erschienen bei dieser Verfolgung}

\subsection{I}

\begin{quote}
\textbf{V. 22–24. Am dritten Tage ward es Laban angesagt, daß Jakob flöhe. Und er nahm seine Brüder zu sich, und jagte ihm nach sieben Tagreisen, und ereilete ihn auf dem Berge Gilead. Aber Gott kam zu Laban, dem Syrer, im Traum des Nachts, und sprach zu ihm: Hüte dich, daß du mit Jakob nicht anders redest, denn freundlich.}
\end{quote}

107. Mose beschreibt an diesem Ort nur obenhin, wie Laban dem Jakob nachgejagt sei; denn er wird hernach davon weitläufiger reden. Denn Laban tobt und wird gar erhitzt von Zorn, und erdenkt in seinem Herzen einen Haufen der größten Sünden und Laster, die Jakob soll begangen haben. Nun sieht man, hat er gedacht, was Jakob für ein Mann sei; nun habe ich ihn in seiner eigenen Sünde gefangen und bestrickt; denn dieses ist nun die allergerechteste Ursache, darum ich ihm mit allen Gütern und mit dem ganzen Gesinde wieder zu mir ziehen, und ihn demnach bestricken möge, daß er mir durch einen Eid zusagen und sich verpflichten muß, daß er mir für und für sein Lebenlang dienen wolle. Denn er ist ein Räuber, ein Dieb, der mir meine Töchter und Vieh entführt, ist ein Kirchenräuber und Vatermörder, und ist unter diesen Sünden eine jegliche werth, daß sie mit eigenem Gefängniß gestraft werde.

108. Also ist Laban vom Zorn gar erhitzt, und dichtet sich selbst die gerechtesten Ursachen, wie er diesen allerbösesten Buben strafen möge. Und er würde wahrlich mit diesem frommen Manne greulich umgegangen sein, wenn er seinem grimmigen Zorn hätte folgen dürfen. Ich habe aber droben gesagt, daß Jakob durch Eingebung Gottes und auf göttlichen Befehl, so er vom Engel empfangen hat, weggezogen sei; sonst ist es an sich selber ein gefährlich Ding, daß er hätte sollen fliehen und keinem Menschen in dem ganzen Geschlecht und Land etwas davon sagen, ja, auch wider seines Schwiegervaters Willen; und ihm dazu seine Töchter, Kindeskinder und die ganze Heerde wegführen, nach welcher Heerde Laban mit seinen Kindern heftig gestanden hat. Wo nun Gott solchem greulichen Vornehmen Labans nicht gewehret hätte, so wäre er mit Jakob auf unbillige Weise umgegangen.

109. Dieses soll man darum so viel fleißiger betrachten, auf daß wir in unserm Herzen bedenken mögen, wie Laban so einen grimmigen Zorn gefaßt, und wie heftig er begehrt habe, dieses Unrecht zu rächen. Denn es können viel Sünden zusammen, von denen wir hernach hören werden, wie sie Laban so sehr vergrößern wird. Vielleicht möchte er seines Lebens geschonet haben, aber nur um seines eigenen Nutzens willen, daß er ihm sein ganzes Lebenlang hätte müssen verpflichtet sein.

110. Dieses ist aber ein sehr trefflich Exempel, daraus wir lernen, wie Gott dem Toben und grimmigen Zorn des Teufels und seiner Glieder wehren und steuern könne, daß sie ihres Gefallens nicht toben und wüthen dürfen. Denn also bricht und hindert Gott den zornigen Muth Labans, daß er wider Jakob auch nicht einmal das Maul hat dürfen aufthun. Und dasselbe wird uns also zum Trost vorgehalten, auf daß wir auf Gottes Gnade und Barmherzigkeit gewisse Hoffnung haben, welcher, wie St. Paulus 1 Cor. 10, 13. sagt, getreu ist, der uns nicht lässet versuchen über unser Vermögen, sondern machet, daß die Versuchung so ein Ende gewinne, daß wir es ertragen können. Gott setzt der Versuchung ein Maß, daß wir nicht über unser Vermögen oder aber wider seinen Willen und Rath versucht werden.

\subsection{II}

111. Der Teufel und Laban hatten gewißlich beschlossen, daß Jakob den Tod verdient hätte; sie würden aber die Strafe etwas gelindert haben, doch unter der Bedingung, wenn er sich mit seinem ganzen Gesinde in ewige Dienstbarkeit begeben hätte, welche Dienstbarkeit schwerer gewesen wäre denn der Tod selbst. Dieweil sie aber darauf denken, wie sie ihn unterdrücken wollen, da kommt Gott und macht ihren Rath und ihr greulich Vornehmen zu nichte. Du Laban, spricht Gott, sollst das mit nichten thun, das du dir vorgenommen hast; ja, du sollst ihm auch kein hart oder unfreundlich Wort geben. Also setzt er seinem Zorn ein Ziel; wie Gott zu Hiob sagt im 38. Capitel v. 10. 11. von den Wellen des Meeres: „Da ich dem Meer den Lauf brach mit meinem Damme, und setzte ihm Riegel und Thür, und sprach: Bis hieher sollst du kommen, und nicht weiter; hie sollen sich legen deine stolzen Wellen.“ Denn wo das Meer vom Winde und Ungestüm getrieben wird, schwillet es auf und droht gleichsam dem Damme, der darumher gezogen ist, mit seinen Wellen, als ob es jetzt ausbrechen und über alle Länder laufen wolle. Aber ich, spricht Gott, habe ihm einen Riegel und Thür gesetzt; darum fürchtet sich der Damm vor dem Toben und den Wellen des Meeres nicht.

112. Also erhebt sich der Teufel wider die Kirche und droht ihr sehr greulich, versucht, die Gottseeligen zu überwältigen und gar zu vertilgen. Aber der dem Meer Riegel und Thür gesetzt hat, der hält auch die stolzen Wellen, daß der Teufel seines Gefallens nicht wüthen kann. Auf solche Weise hindert Gott auch alles Toben und grimmigen Zorn Labans mit einem Zaum und macht ihn so stille wie ein Lämmlein, wie unwillig und zornig er auch war. Denn Gott läßt nicht zu, daß die, so da glauben und sein Wort haben, über ihr Vermögen betrübt oder versucht werden; er ist ein getreuer Hüter derer, die sein Wort halten und bewahren. Jakob hatte Gottes Wort, da er zu ihm sagte: „Zeuch wieder in deiner Väter Land“; item: „Ich habe alles gesehen, was dir Laban thut“ etc. Dem Wort gehorcht er und nimmt Alles mit sich, wiewohl es sich ansehen läßt, daß er daran die schwersten Sünden thue; und Gott bewegt die Herzen seiner Weiber, daß sie ihrem Manne und dem Worte folgen und gehorchen, ja, daß sie heftig bei ihm anhalten, daß er wegziehen wolle.

113. Auf solchen Gehorsam folgt die Verfolgung. Der Teufel versucht, ihn mit seinem ganzen Gesinde wieder zurück zu ziehen und ihn zu unterdrücken. Und ist wohl zu glauben, daß Jakob mit seinen Weibern in großer Angst gewesen sei. Denn es wird ihm ohne Zweifel eingefallen sein, daß er also gedacht hat: Siehe, Laban jagt dir nach und wird dich tödten, und Alles verderben und hinreißen: Ei, wie habe ich so unweislich gehandelt, daß ich geflohen bin und ihm meinen Rath verborgen gehalten habe! Was soll ich ihm für Ursache anzeigen, daß ich so heimlich weggezogen bin? Ich habe mich und meine Familie in diese Gefahr gestellt, und habe daran unweislich gethan etc.

114. Dieses ist eine innerliche Anfechtung gewesen, die ihm sein Herz bekümmert hat. Und also ist er anfänglich von Laban verfolgt worden, inwendig aber haben ihn die schweren Gedanken der Verzweiflung gemartert. Und seine Weiber und Mägde und die ganze Familie haben eben dieselbige Anfechtung auch gefühlt; Die werden gesagt haben: Wie viel besser wäre es gewesen, wir wären in Mesopotamien geblieben, denn daß wir hier im Elend jämmerlich umkommen sollen, oder aber wiederum zurückgeführt werden in den vorigen Dienst, der nun viel schwerer sein wird! Wir hätten zwar mit dem vorigen Dienst mögen zufrieden sein; wo wir dabei ja nicht groß Gut hätten haben können, so hätten wir doch Friede gehabt und wären ohne Gefahr geblieben. Es ist ja zu erbarmen, daß wir in einem Augenblick entweder unser Leben, oder aber diese wenigen Güter und alle Hoffnung der Freiheit verlieren sollen.

\quellseite{178}

115. Solche Gedanken werden ihnen ohne allen Zweifel eingefallen sein. Mose hat zwar im Text nicht ausgedrückt, wie sie im Herzen gesinnet gewesen; er hat es aber damit genugsam angezeigt, da er sagt, Gott sei zu Laban, dem Syrer, im Traume gekommen etc. Unterdeß aber hat sich Jakob an die Verheißung gehalten, so ihm Gott gegeben hatte, nämlich, daß er auf GOttes Befehl weggezogen wäre, und wird damit seine Weiber und Familie aufgerichtet und getröstet haben, daß ihm Gott zu Bethel erschienen war; item, daß er zu ihm gesagt hatte: „Zeuch wieder in deiner Väter Land, und zu deiner Freundschaft“ etc. Haltet fest, lieben Weiber und Kinder, hat er gesagt, Gott wird uns nicht verlassen noch verwerfen.

116. Also hat Jakob gekämpft und die Anfechtung durch den Glauben überwunden. Denn das ist der Kampf des Glaubens: zur Zeit des Friedens und außerthalb der Anfechtung thut er gute Werke; zur Zeit der Anfechtung ruft er, ja, er schreit heftig, daß es alle Engel hören. Denn er wird in die Hölle geführt und muß die Gefahr und Noth des Todes erfahren; da hebt er an zu kämpfen und zu schreien, mit welchem Geschrei beides, Himmel und Erde, bewegt wird.

117. Darum kommt nun GOtt zu Laban im Traume des Nachts, da er in seinem Herzen greuliche Anschläge gehabt, und setzt der Verfolgung Ziel und Maß: da werden die stolzen Wellen und die grimmigen Wassermogen im Herzen Labans stille. Denn GOtt hat im Traum zu ihm also gesagt: Ich sehe, daß du grimmig und hitzig bist vom Zorn, und richtest all dein Vornehmen und Gedanken dahin, daß du Jakob willst Gewalt thun. Du läßt dich dünken, du habest zumal die gerechteste Ursache, ihn zu verfolgen und zu unterdrücken; denn du meinst, es sei auf Erden kein böser Mensch denn dein Jakob. Du hast in deinem Herzen einen großen Haufen Sünde und Lasterung bedacht, die du über ihn willst ausgießen. Aber ich gebiete dir, daß du nicht allein die Hände stille haltest und ihn damit nicht beleidigest, sondern daß du ihn auch nicht mit einem Worte erzürnest.

118. Auf diese Weise verläßt GOtt seine Gläubigen nicht, vornehmlich wenn sie ihn im rechten Glauben anrufen; wie wir auch zum öftern erfahren haben, wie wunderbarlich uns GOtt nach seiner Güte bewahrt hat. Denn wie viel und wie greuliche Anschläge der Könige, Päbste und Cardinalen sind bisher wunderbarer Weise verhindert worden, zwar nicht durch unsere Kräfte oder Rath, sondern durch göttliche Kraft! Denn GOtt hat sie entweder gar erschrecket gemacht, daß sie nichts haben dürfen anfangen; oder wenn sie sich etwas unterstanden haben, hat er ihren Rath und Vornehmen doch zunichte gemacht und gehindert.

119. Es sind neun Reichstage gehalten worden von der Zeit an, da das Evangelium in Deutschland wiederum hervorgekommen ist und angefangen hat zu leuchten, darüber die Widersacher greulich getobt und uns sehr heftig gedroht haben. Ja, sie haben zu Augsburg öffentlich gesagt, sie wollten Gut und Blut zusammen setzen. Und sie hätten uns zwar lebendig zerschlungen, wie in 124. Psalm V. 3 steht, wo uns GOtt nicht bewahrt und ihre Anschläge zunichte gemacht hätte. Wie oft haben wir gesehen, wie der Mordbrenner getobt und allerlei Böses vorgenommen hat, bis er endlich aus seinem Lande ist verjagt worden! Also werden sie hinfort noch immer mehr versuchen, und werden nicht ruhen können, bis sie der Türcke zerstreuet wird.

120. Darum sollen wir an dieses Exempel denken, wie Jakob sei errettet worden von Laban, seinem zornigen und greulichen Feinde, durch das Wort, denn er geglaubt hat. Denn wo das Wort ist, da ist auch der Glaube; wo aber der Glaube ist, da ist auch das Schreien um der Anfechtung willen; das Schreien aber muß gewißlich erhört werden. Und wenn das Schreien erhört ist, so zerbricht es alle Gewalt des Himmels und der Erde, und aller höllischen Pforten. Darum wird das lange Register der vielen Sünden, das Laban zusammengetragen hat, Jakob damit zu beschuldigen, bald in einem Hauf zerstreut. Ja, es wird sich Laban dessen noch wohl gefreut haben, daß er von GOtt nur mit Worten ermahnt worden und nicht irgend mit einer sonderlichen Plage oder Kreuz um seines Neides und Geizes willen gegen seinen Schwiegersohn, deß er sich hat schuldig gewußt.

121. Also ist er mit Gefahr des Todes erschrocken und geschlagen worden, daß er jetzt nicht begehrt, wie er seinen Schwiegersohn möchte in Gefahr bringen, sondern wie er der Strafe und dem Unglück, davon er sich gefürchtet hat, entfliehen möchte, wiewohl er durch die Vermahnung nicht frömmer geworden ist. Er ist wohl durch GOttes Gericht erschrocken, aber er hat bald wiederum gethan, wie seine Art war, wie wir hören werden. Er hat zwar die Hände stille gehalten, daß sie nichts Böses gethan haben; er enthält sich aber des Fluchens und der Schmähworte nicht. Das ist eine Galgen- und Judasreue; wie die Päbste, Bischöfe und ihre Trabanten, Könige und Fürsten, gehalten werden und verhütet in ihren Rathschlägen, die sie wider uns halten: es wird ihnen der Muth genommen, wie der 76. Psalm V. 13. sagt; aber sie bessern sich gleichwohl nicht. Das ist der Heuchler und Ungläubigen Buße: die thun Buße wie Laban oder Esau. Es ist keine ernste oder freiwillige Buße, sie bessern sich nicht, sondern werden nur abgeschreckt, daß sie dasjenige nicht vollbringen, deß sie sich unterstanden hatten.

122. Von den Träumen habe ich an andern Orten gesagt, wie sie zu unterscheiden sind. Die rechten Träume bringen mit sich, daß sie den Träumenden so tief eingeprägt werden und die Herzen bewegen, daß sie gleichsam darüber erschrocken und bestürzt werden, daß solche Träume nicht können verachtet werden; wie die Träume Pharaos und Nebucadnezars gewesen sind. Die andern Träume aber sind nur vergebliche Hoffnung und Gedanken, die man verachten soll.

123. Wir sollen aber dieses Exempel zu unserer Lehre und Nutzen gebrauchen; denn es wird uns um der Ursache willen vorgehalten, daß wir sehen sollen, daß GOtt bei den Seinen sein und ihnen helfen wolle, wie elend und schwach sie auch sind. Laban ist viel gewaltiger und stärker, denn Jakob, welcher in großer Gefahr steckt, und gar wehrlos und arm ist: er hat aber das Wort; darum ist auch GOtt und die Schaar der heiligen Engel bei Jakob, der da glaubt und GOtt im Glauben anruft. Wie nun GOtt diesen Patriarchen aus seines Vaters Hause geführt hat, also führt er ihn auch wiederum hinein durch den Schutz der heiligen Engel. Eben auf dieselbe Weise wird er uns auch, wenn wir glauben und mit starkem festem Glauben dem Worte anhangen, erretten, auch in der Pestilenz, im Tode und im Kriege. Denn es ist unmöglich, daß der Mensch sollte verlassen und nicht geschützt werden, der dem Wort GOttes glaubt.

124. Daßselbe wird auch die nachfolgende Historie von den Engeln ganz herrlich bezeugen. Denn sie haben größere Stärke denn alle Feinde; wie das Exempel von der Belagerung Jerusalems unter Sanherib ausweiset. Und Christus spricht zu Petro Matth. 26, 53.: „Meinest du, daß ich nicht könnte meinen Vater bitten, daß er mir zuschickte mehr denn zwölf Legionen Engel?“ Dieweil wir nun unter dem Schutz GOttes sind, so ist daran kein Zweifel, daß wir auch im Schutz der heiligen Engel sind, die über uns wachen, die im Leben bei denen sind, so in Noth und Gefahr stecken, und die Sterbenden zum Frieden und zur Ruhe bringen. Denn David sagt im 91. Psalm V. 11.: „Er hat seinen Engeln befohlen über dir“ etc.; item im 34. Psalm V. 8.: „Der Engel des HErrn lagert sich um die her, so ihn fürchten, und hilft ihnen aus.“ Und dieses sollte uns billig ermahnen und aufwecken, daß wir das Wort gern hören und lieben, demselbigen aber noch lieber glauben sollen mit dieser Zuversicht, daß wir unter dem Schutz der heiligen Engel sind.

125. Dieses sei gesagt von dem göttlichen Befehl, dadurch Laban ist verboten worden, daß er Jakob weder mit Worten noch mit Werken beleidigen sollte. Denn dieses göttliche Urtheil macht dem Laban das Register, darin er die Sünden zusammengetragen, Jakob damit zu beschuldigen, und alle seine Rechnung und Anschläge gar zunichte. Die Worte lauten aber nach Art der hebräischen Sprache also: Du sollst mit Jakob nicht anfangen freundlich zu reden, und darnach die Rede mit Scheltworten beschließen. Aber Laban ist dem Befehl nicht allenthalben gehorsam gewesen, wiewohl er sich der That hat enthalten müssen. Das Register aber wird gar verworfen. Deine greulichen Anschläge und Gedanken, hat die göttliche Stimme gesagt, sind nichts, denn ein lauterer Traum, Eitelkeit, Irrthum und Lüge; darum sollst du deinem grimmigen Zorn nicht folgen. Er wird sich aber doch gleichwohl nicht gar enthalten können, daß er Jakob nicht sollte schelten und ihn mit Schmähworten hart anfahren; denn es folgt zumal eine harte Klage, die er über Jakob thut.

\quellseite{179}

\begin{quote}
\textbf{V. 25. Und Laban nahte zu Jakob. Jakob aber hatte seine Hütte aufgeschlagen auf dem Berge; und Laban mit seinen Brüdern schlug seine Hütte auch auf, auf dem Berge Gilead.}
\end{quote}

126. Nach der Strafe und Ermahnung, so Laban im Traum gehört hatte, daß er mit Jakob nicht hart oder unfreundlich reden sollte, zieht er doch weiter, wiewohl ihm alle seine Rachepläne zerstört sind. Und dieweil er ihm seinen Schaden thun kann, so stellt er sich gleichwohl, als wollte er ihm gern Schaden thun und daß er zornig auf ihn sei damit, daß er ihn droht, Böses zu thun. Denn er will ja den Namen nicht haben, daß ihm seine Anschläge sollten gefehlet haben, oder daß er da müsse wiederum wegziehen, und hätte nicht ausgerichtet, was er sich auszurichten vorgenommen hatte. Er hielt es dafür, es wäre ihm eine große Schande, daß er sich zuvor also sollte gerühmt und mit Drohen und Morden geschnaubt haben, und daß solches sollte vergeblich sein. Darum gibt er großen Zorn vor, als wollte er mit seinem Vornehmen fortfahren, unangesehen, was ihm Gott im Traum gesagt hatte.

127. Derhalben ist dieses ein fein Bild der falschen Buße, wie die Heuchler haben. Denn solche Reue und Leid haben bisweilen die bösen gottlosen Menschen wohl; nicht daß sie rechte wahre Reue und Leid tragen oder ernstliche Buße thun, welcher Meinung ich ehemals gewesen bin, sondern sie stellen sich nur äußerlich also, als ob sie Buße thun und Leid tragen für ihre Sünde. Darum muß man Davids Worte, 2 Sam. 12, 13., die er zu Nathan redet, da er sagt: „Ich habe gesündiget“, viel anders verstehen, denn Sauls Worte, der eben also zu Samuel sagte, 1 Sam. 15, 24.: „Ich habe gesündiget.“ Es sind zwar einerlei Worte, Eine Stimme, und läßt sich ansehen, als ob es einerlei Reue und Buße sei; aber die Herzen sind sehr ungleich.

128. Die bösen gottlosen Menschen thun also Buße, daß sie sich darüber mehr betrüben, daß ihnen ihre böse Lust und Sünden verboten sind, denn daß sie sich damit bekümmern sollten, wie sie die böse Lust und Sünden tödten möchten. Das ist eine Gesetzbuße, die wir pflegen eine Galgenreue zu nennen. Denn wo sich der Dieb vor dem Galgen und vor der Strafe nicht müßte fürchten, so würde er viel lieber stehlen wollen, denn daß er sich anderer Leute Gut enthalten sollte; darum ist es ihm leid, daß er sich um Furcht der Strafe willen des Stehlens enthalten muß.

129. Also wird Laban auch abgemalt, daß er nicht recht habe Buße gethan, sich auch nicht von Herzen gebessert; sondern es ist ihm leid, daß ihm seine böse Lust und grimmiger Zorn durch göttliche Gewalt ist gebrochen und verhindert worden. Darum ist es nur eine äußerliche Buße zum Schein. Wie Saul auch sagt: „Ich habe gesündiget; aber ehre mich doch jetzt vor den Aeltesten meines Volks“ etc., 1 Sam. 15, 30. Es ist ihm angst und bange dabei, daß er vor den Aeltesten soll verhöhnt oder geschmäht werden, und nicht, daß er Gott erzürnt hat.

130. Aber ein recht bußfertig Herz ist also gesinnet, daß es sich vor nichts Anderem fürchtet, als vor Gottes Zorn und Ungnade, und achtet der Schmach oder Schande vor den Menschen gar nicht, wenn es nur wissen mag, daß ihm Gott wolle gnädig sein; wie David im 51. Psalm V. 6. 19. ausgedrückt, daß er gefühlet hat, daß er der Sünden halben also ist gesinnet gewesen.

131. Es ist aber zu merken, daß die Heuchler bleiben, wie sie sind; und wenn sie ja nicht toben und wüthen mögen, wie sie gerne wollten, so lassen sie sich doch vernehmen, daß sie grimmig böse sind und daß sie auch sonderliche Gewalt haben; und daß sie nicht für geringer geachtet werden, denn die, die sie hassen und verfolgen.

132. Der Berg Gilead ist zu der Zeit nicht also genannt worden, sondern er wird hier vormals mit dem Namen genannt, damit er hernach erst ist bezeichnet worden. Es scheint aber, als habe Mose anzeigen wollen, daß Laban den Berg oben eingenommen habe, Jakob aber habe seine Hütten unten am Berge aufgeschlagen; auf daß jener damit anzeigte, daß er der Höhere und Gewaltigere wäre. Nun folgt aber im Text, wie Laban die Sünde Jakobs mit so vielen Worten vergrößert und nach einander hat wissen zu erzählen.

\section[Fünfter Theil]{Fünfter Theil. Wie sich Laban nach der Erscheinung gegen Jakob verhalten, und wie sich Jakob bei diesem Verhalten Labans gezeigt}

\subsection{I}

\begin{quote}
\textbf{V. 26—30. Da sprach Laban zu Jakob: Was hast du gethan, daß du mein Herz gestohlen hast, und hast meine Töchter entführet, als die durchs Schwert gefangen wären? Warum bist du heimlich geflohen, und hast dich weggestohlen, und hast mir's nicht angesagt, daß ich dich hätte geleitet mit Freuden, mit Gesängen, mit Pauken und Harfen? Und hast mich nicht lassen meine Kinder und Töchter küssen? Nun, du hast thörlich gethan. Und ich hätte, mit Gottes Hülfe, wohl so viel Macht, daß ich euch könnte Uebels thun; aber eures Vaters Gott hat gestern zu mir gesagt: Hüte dich, daß du mit Jakob nicht anders, denn freundlich, redest. Und weil du denn ja wolltest ziehen, und sehnetest dich so fast nach deines Vaters Hause, warum hast du mir meine Götter gestohlen?}
\end{quote}

133. Wenn Laban recht hätte Buße gethan, so hätte er auch geglaubt und dem Worte Gottes Beifall gegeben, welches von Jakob sagt, daß er fromm und unschuldig wäre, und daß er nicht verdient hätte, daß ihm solche Gewalt geschehen. Denn er hört ja, daß er von Gott absolvirt, ledig und los erkannt wird von allen Sünden, um welcher willen er ihn verfolgt. Aber auf daß er sich selbst bei Ehren behalten und seine Gerechtigkeit vertheidigen möge, so redet er mit ihm, auf daß er nicht für einen solchen Mann geachtet werde, der seinem Zorn keinen Nachdruck habe geben können, und auf daß Jakob erschrecken und seine Sünde bekennen soll. Das hat der stolze Heuchler und sehr listige Werkheilige haben wollen.

134. Darum sagt er nun zuerst also: „Was hast du gethan?“ Als wollte er sagen: Alles, was du gethan hast, sind die ärgsten Sünden: du bist der allerhöchsten Sünden schuldig, um welcher willen ich dich billig nach meiner Gerechtigkeit und Recht strafen möchte. Er hat ihm aber nicht Uebels thun dürfen, sondern stellt sich männlich und macht sich eine eitele Gerechtigkeit an. So ist nun vorhin gesagt, daß er geizig gewesen sei; jetzt aber sehen wir, daß er dermaßen beschrieben und abgemalt wird, daß er auch ein hoffärtiger stolzer Heuchler und Werkheiliger gewesen sei: und ist doch gleichwohl im Traume durch Gottes Wort ermahnt worden, daß Jakob fromm und gerecht wäre. Aber er erkennt weder seine Sünde, noch die Tugend und Frömmigkeit Jakobs; ja, er macht noch über das die Sünde Jakobs sehr groß: „Du hast mir“, sagt er, „mein Herz gestohlen.“ Das ist ein sehr greulich Laster.

135. Als erste Sünde erdichtet er das Verbrechen der Majestätsbeleidigung; denn ich, sagt er, war ja der Vater und Schwiegervater, du warst mein Knecht: du hast aber aller Ehre und alles kindlichen Gehorsams vergessen, ja auch aller Liebe, so du billig zu mir tragen solltest. Du hast dich unterstanden, ohne Erbietung wegzuziehen, und mir als deinem Vater, Schwiegervater und Herrn nicht ein Wort davon zu sagen; hast mir über das meine Töchter zu entführen gewagt als ein Dieb, Mörder und Räuber. Aber Laban läugt auf das unverschämteste. Denn ob er ihm wohl das Herz gestohlen hat, so hat er doch dazu gute Ursache gehabt, daß er heimlich geflohen ist; denn er hätte sonst von dem geizigen Herrn keine Erlaubniß erlangen können, heim zu ziehen.

136. Die andere Lüge aber ist, daß er sagt: „Du hast meine Töchter entführet, als die durchs Schwert gefangen wären.“ Denn sind es nicht seine Weiber, die du, Laban, ihm ehelich vertraut hast? Und diese Ehe hat ihn zwar sehr viel gekostet. Heißt denn das mit Gewalt oder durch das Schwert entführen, wenn ein Mann sein Weib mit sich führt, die ihm nach allem Recht folgen soll, und sie hat auch von selbst eingewilligt und den Mann ermahnt wegzuziehen? Da siehst du nun, wie Laban so eine gifftige und teuflische Zunge habe, damit er seine Sünden herrlich macht und schmückt, und was wohl gethan ist, lästert. Denn der schändliche Heuchler dichtet und lügt dieses alles, daß er Jakob und seine Familie damit verhöhnen und zu Schanden machen, aber seine Gerechtigkeit und Ehre schmücken möge. Wann sind aber die Töchter entführt oder durch das Schwert gefangen worden? Wann hat sie Jakob geraubt als ein Dieb oder Mörder? Sind es nicht seine ehelichen Weiber, die er ja nicht geraubt oder entführt hat, sondern die von selber und guten Willens ihrem Manne nachgefolgt sind, wie sie ihm billig haben folgen sollen?
\quellseite{180}

137. Derhalben hätte es sich gebührt, daß Laban also sollte gesagt haben: Einmal du ja mein Haus hast verlassen wollen, so hast du daran recht und wohl gethan, daß du deine Weiber nicht verlassen hast, sondern hast sie mit dir weggeführt. Und wiewohl mir das wehe thut, daß du mir mein Herz gestohlen hast, so liebe ich doch das an dir, daß du meine Töchter hast wollen mit dir fliehen lassen. Das hätte ihm besser angestanden, daß er also geredet hätte, und das hätte einem frommen und ehrlichen Mann gebühret. Aber er ist ein gottloser Heuchler, darum hat er es alles umgekehret, und was an seinem Schwiegersohne zu loben und zu rühmen war, das tadelt und straft er auf das allerfeindseligste.

138. Siehe aber, ob er auch das gehalten habe, das ihm von Gott geboten war, daß er mit Jakob nicht hart oder unfreundlich reden sollte? Er trägt zusammen auserlesene und giftige Schmähworte, und was noch das Allerschändlichste ist, die da falsch und erlogen sind. Denn er weiß ja wohl, daß es seine Weiber sind, und da er die wegführt, hat er ja daran nicht gesündigt, vornehmlich da sie ihm gutwillig gefolgt sind: ja, es sollte sie Laban auch wider ihren Willen dazu gezwungen haben, daß sie ihrem Manne folgten.

139. Darum ist es ein verfluchter Heuchler, der einen solchen trefflichen Mann, der so viel große Tugenden an sich hat, mit solchen giftigen Schmähworten lästern und beschweren darf. Und zwar wiederholt er noch einmal, daß er zuvor gesagt hatte: „Du hast mir mein Herz gestohlen.“ Warum hast du so heimlich ohne mein Wissen fliehen wollen? Denn der Unbilligkeit kann er nicht vergessen, die beißt ihn in seinem Herzen. Denn er betrübt sich der Sünde halben nicht, welches auch keine Sünde an Jakob ist; denn Jakob ist sehr fromm und heilig: aber das verdrießt Laban, daß er es an keinem Zeichen hat merken können, daß Jakob diesen Anschlag gemacht hatte mit seinem Abzuge. Ja, sagt er, hast du mir nicht allein das Herz gestohlen, sondern mich selbst auch mit allem meinem Gut. Denn Jakob hatte ihm sein Hab und Gut gemehret, daß es sehr groß geworden. Daß ihm nun der gute Bissen aus dem Rachen entzogen ist, das thut ihm wehe, und denkt also: Wenn ich das hätte wissen sollen, so wollte ich allerlei versucht haben, daß ich dich hätte behalten mögen; denn jetzt ist nun alle mein Heil, dazu das Gedeihen aller meiner Güter hinweg.

140. Da er aber nun fast allen seinen Zorn und Unmuth ausgeschäumt hat damit, daß er Jakob greulich angeklagt und sich selbst gerechtfertigt, rühmt er nun von seiner Liebe und Gewogenheit beide gegen seinen Schwiegersohn und auch gegen die Töchter, daß er damit die Sache, darum er Jakob beschuldigt, so viel schwerer mache, und gibt große Liebe, Freundlichkeit und herzliche Neigung vor, so er zu seinen Kindern trage; und will also sagen: Ich habe dich und meine Töchter darnach geliebet, und euch gedienet, womit ich euch habe dienen können, daß du an mir nichts vermissen könntest, sondern mich allezeit hast als einen sehr guten Freund gespüret. Darum wollte ich jetzt auch alle guten Freunde und Einwohner dieses Landes berufen und dich ehrlich und herrlich geleitet haben, mit Freuden, mit Singen, mit Pauken etc. Ja wahrlich, das würde er gethan haben. Was würde es aber für Freude und Pauken gewesen sein? Antwort: Daß er ihn hätte in Eisen und Gefängniß geworfen und daß er mit einiger Dienstbarkeit hätte müssen beschweret sein.

141. Wenn Jakob dem Laban hätte angezeigt, daß er mit seinen Weibern davon ziehen wollte, so hätte er ihn nimmermehr ziehen lassen. Denn solche Freundschaft und Dienst hat er seinen Freunden pflegen zu erzeigen, wie zuvor ist gesagt worden: er würde ihn wiederum in ewige Dienstbarkeit und Elend gezogen haben, wo es ihm GOtt nicht gewehrt hätte; wie er solche Tyrannei damit genugsam bewiesen hat, da er Jakob seinen Lohn wohl zehnmal verändert hat, und ihn gezwungen, daß er ihm ganze vierzehn Jahre um seine Töchter hat dienen müssen und sechs Jahre um die Heerde Viehes, ohne gewissen gebührlichen Lohn. Derhalben, wo GOtt seinem Geiz nicht gewehret hätte, so würde er seinen Schwiegersohn aller Güter beraubt haben. Ja, das sind seine Pauken, Harfen und Freude, damit er Jakob hätte geleiten wollen.

142. Er erzählet aber dieses alles nur darum, daß er den frommen und heiligen Mann Jakob damit beschweren und verhasst machen, und seine Tyrannei oder Grausamkeit und Geiz schmücken und bedecken will, und sind lauter läppische Lügen. Ja, je freundlicher die Worte sind, die er redet, je mehr Bitterkeit er in seinem Herzen verborgen hat, und alle Worte sind mit sonderlichem Fleiß dahin gerichtet, die höchsten Tugenden an Jakob damit zu lästern und zu tadeln. Denn wie sollte er den haben wollen ziehen lassen, von dem er begehret hat, daß er ihn auch hätte mögen tausend Jahre in seinem Hause behalten? Und das ist ihm vornehmlich leid gewesen, daß er einen solchen getreuen Knecht sollte verlieren, der ihm so überaus nützlich und geeignet war. Das ist der Diebstahl des Herzens.

143. Er setzt aber noch mehr hinzu und will damit gleich auf einen heimlichen Einwand Antwort geben: „Du hast mich nicht lassen meine Kinder und Töchter küssen.“ Als wollte er sagen: Ob du wohl sagst, daß meine Töchter deine Weiber sind, so bist du doch darum zu strafen, daß du mich aller freundlichen Werke, die ein Vater üben soll, beraubt hast, die ich meinen Kindern und Kindeskindern würde erzeigt haben nach der Liebe und natürlichen Neigung, die ich zu ihnen trage. Denn das ist ja bei allen Völkern also ein gemeiner Gebrauch, wo die Töchter aus ihres Vaters Hause ziehen, daß sie alsdann ihre Eltern zuvor grüßen und Abschied nehmen, und begehren, daß sie der Vater küssen, umarmen und segnen wolle; dasselbe habe ich von Herzen begehrt zu thun. Es sind zwar deine Weiber, die du mit all dem Recht hast wegführen mögen: aber es hätte sich gebühret, daß du dich meine herzliche Begierde solltest haben bewegen lassen, und mir zugegeben haben, daß ich zuvor meine sehr lieben Töchter und Kindeskinder hätte küssen mögen, und solltest sie nicht so greulich und auch heimlich ihrem lieben Vater aus seinen Armen, damit er sie hätte mögen umfangen und küssen, gerissen haben.

144. Ja, kommen daher die bitteren Thränen? O wie muß das so ein eisern Herz sein, daß solcher Thränen haben nicht kostet! Der verzweifelte Schalk sollte Pastor oder Bischof von Mainz sein, der sein greulich und tyrannisch Herz mit solchen göttlichen und schönen Worten bedecken kann. Denn wie wolltest du, Laban, deine Töchter küssen? Freilich nicht anders denn mit Judas Kusse? Oder wenn du sie ja hast küssen und deinen Töchtern väterliche Liebe beweisen wollen: warum hast du sie nicht ernährt, gekleidet, und sonst gethan, was du ihnen zu thun schuldig gewesen, da sie dir noch gedienet haben? Ja, da hat er Alles an sich gerissen, hat sie gehalten, als wären sie Fremde, und hat sie verkauft, als ob sie gefangen wären; und wo sie nicht wären davon geflohen, so hätten sie müssen Hungers sterben. Ja, das war die große Liebe des gottlosen, ich sollte sagen, des greulichen und verfluchten Vaters Laban. Jetzt aber, da er sieht, daß er beschuldigt wird, daß er sich offenbar gegen seine Töchter übel und ganz ungebührlich gehalten habe, und daß sein Schwiegersohn durch seine Grausamkeit dazu gezwungen worden sei, daß er mit seinen Töchtern hat fliehen müssen: da will er nun diese seine Grausamkeit zudecken damit, daß er seinen Schwiegersohn beschuldigt, als ob er ihm seine Töchter wider seinen Willen und ohne Abschied entführt; nicht daß er nach den Töchtern etwas fragt oder sie liebt, sondern daß er nur seinen Geiz damit bedecken möge. Denn das bekümmert ihn noch immer, daß er anfänglich gesagt hat, wie ihm Jakob sein Herz gestohlen hätte, damit er gemeint hat, daß ihm der Nutzen und das Gedeihen seiner Güter, so er vom Dienste Jakobs gehabt, genommen und entzogen worden sei.

145. Denn was würde der greuliche Tyrann eben jetzt in dem Augenblick, da sie von ihm gezogen sind, gethan haben? So er doch zuvor innerhalb zwanzig Jahren ihnen gar keine Freundschaft oder väterliche Liebe, wie es einem Vater wohl geziemt hätte, erzeigt hat? Wenn aber noch irgend ein Tröpflein väterlicher Liebe oder Freundlichkeit an ihm gewesen wäre, so hätte er sich noch anders erzeigen müssen: er hätte alles Zornes und Unbilligkeit vergessen, oder aber hätte ja also sagen
\quellseite{181}
müssen: Wohlan, ich erkenne es nunmehr wohl, daß ich mit euch etwas hart und unfreundlich umgegangen bin und daß ihr gleichwohl von mir gar kein Geschenk oder Lohn empfangen habt. Derhalben will ich euch nun als meinen Kindern dieses von meiner Heerde, oder von dem, das mir am allerliebsten ist, geben, als nämlich, etliche Stücke und goldene Kleinode. Das sollt ihr zum Zeichen haben unserer Liebe und Freundschaft, die für und für unter uns währen soll.

146. Was sollte aber der sehr geizige Heuchler geben? Ja, wir werden hernach hören, wie er alle Dinge so genau erforscht und ausfragt, und wollte auch das gern weggenommen haben, das sie vom Raube, so ihnen Gott gegeben hat, bekommen hatten. Unterdes aber weiß er sich fein zu schmücken, als ob er ganz fromm sei und seine Kinder sehr lieb habe; Jakob aber muß den Namen haben, daß er sich unfreundlich und tyrannisch gegen seinen Schwiegervater gehalten habe. Dieses ist eine unglaubliche und unerhörte Bosheit, Heuchelei und Grausamkeit. Gottlose Leute und die Heiden, die der Natur und Vernunft folgen, scheuen sich vor solchen greulichen Sünden; und es läßt sich ansehen, als ob es nicht wohl möglich sei, daß solche große Grausamkeit sollte in eines Menschen Herz kommen können, dadurch Einer seine eigenen Töchter und Kindeskinder nicht allein gar nicht achtet, sondern sie auch noch über das aller schuldigsten Pflicht und verdienten Lohnes berauben dürfte. Das muß aber eine verfluchte Heuchelei und Verstellung sein, die noch solche große Sünde bedecken kann mit einem schönen Schein und Vorwendung des Kusses, der Umarmung und großer väterlicher Liebe. Denn Laban hat seine Töchter darum küssen und sie umarmen wollen, auf daß er sie damit in die Bande ewiger Dienstbarkeit bringen und behalten möchte.

147. Dies sind nun die Tugenden des trefflichen heiligen Mannes Laban, und dagegen auch die Sünden und Laster Jakobs, die er ihm aufrückt und vorwirft. Denn also pflegt es in der Welt zu gehen: fromme Leute werden für Sünder gehalten, und wiederum, Sünder für Fromme. Es ist in der ganzen Welt kein Sünder, denn allein der Gottes Wort hat und an Jesum Christum glaubt. Die aber das Wort verfolgen und demselben feind sind, die sind gerecht und fromm, und wie Christus sagt Joh. 16, 2., meinen, sie thun Gott einen Dienst daran. Es kommt aber auch oft vor, daß die vom Wort nichts wissen und es nicht verstehen, die Heiligen oder Gläubigen hassen und verfolgen. Aber so fromm ist Laban nicht, daß er aus Unwissenheit sollte gesündigt haben, sondern er verfolgt wissentlich den frommen Mann Jakob. Wie jetzt die Cardinäle, Bischöfe und Andere auch thun, die wissentlich wider die Frommen toben und wüthen, und wollen doch gleichwohl für fromme heilige Leute gehalten werden.

148. Die letzte und greulichste Sünde, die Laban seinem Schwiegersohne vorwirft, ist, daß er ihm seine Götter gestohlen hat. Da siehe, wie über alle Maßen groß und greulich er die Sünde machen kann. Die Stücke, sagt er, die ich bisher erzählt habe, sind schwer genug und einem frommen Manne nicht; ich lasse sie aber fahren und will dir es verzeihen. Denn du magst dich vielleicht das große Verlangen nach deinem Vaterlande oder irgend eine andere Lust oder Begierde haben überwinden lassen, daß du dich also heimlich weggestohlen hast; aber dies ist die allergröbste Sünde, die du begangen hast, daß du mir meine Götter gestohlen hast. Was willst du nun darauf antworten? Ich will es gar zugeben, daß du irgend aus Leichtfertigkeit oder Ungehuld, oder aber aus Vorwitz dazu bist getrieben worden, daß du davon gezogen, wiewohl du an andern Orten, oder auch in deinem Vaterlande nicht wohl besser hättest leben können denn eben bei mir: warum hast du mir aber die Götter, das ist, meinen Gottesdienst gestohlen? Denn das ist ja eine überaus große Sünde. Ist es daran nicht genug gewesen, daß du dich gegen mich schalkhaftig erzeigt hast? du hast über das dich und die Deinen mit der Sünde auch noch beflecken müssen, daß du bist ein Kirchenräuber geworden.

149. Da hat er nun aber einen andern und viel edlern und viel heiligeren Schein gefunden, den er vorwenden kann, nämlich, die Religion und den Gottesdienst; wiewohl er sich keiner Gottseligkeit befliessen hat, sondern es hat ihm nur der Geiz und Mammonsdienst sein Herz eingenommen. Solche Gesellen sind unsere Widersacher auch. D. Eck und der Bischof von Mainz, die fragen nach der Ehre Christi und der Wahrheit des Evangeliums gar nichts, sondern trachten nur nach den großen Präbenden und geistlichen Lehen, oder nach anderer großer Würde und Herrlichkeit. Also bekümmert sich Laban um den Gottesdienst auch nichts, sondern es ist ihm um das Gold und Silber zu thun, davon die Götzen gemacht waren.

150. Jedoch stellt er sich gleichwohl, als ob es ihm großer Ernst wäre mit der Gottseligkeit und Gottesdienst. Als wollte er sagen: Ich frage zwar nach dem Golde und Silber nicht; denn ich könnte wohl leichtlich andere Götzen wieder machen. Aber darum ist es mir leid, daß du meine Andacht, meinen Gottesdienst und Gebet verletzt hast. Darum magst du nun sehen, wie schwerlich du damit gesündigt hast, daß du Gott dem Herrn seinen Dienst und Religion, und dazu alle Werke der Gottseligkeit in meinem ganzen Geschlechte und Hause gestohlen hast. Dieses ist ein Kirchenraub und Sünde wider die erste Tafel. So viel an dir liegt, sagt er, beraubst du mich und mein Haus der Religion, und öffnest damit der Gottlosigkeit und dem Teufel Thür und Thor; denn der Teufel begehrt doch, daß wir nur sichere Verächter Gottes und des Gottesdienstes sein mögen.

151. Also deutet Laban diesen Kirchenraub, nicht daß es allein ein Diebstahl sei der Götzen, sondern daß er damit den Gottesdienst verletzt und ganz zerstört habe. Denn es ist nie ein Volk so närrisch gewesen, daß es Holz, Steine, Gold oder Silber, daraus die Bilder gemacht waren, sollte angebetet haben; sondern sie haben die erste Tafel auch dazu gezogen, und haben gedichtet und sich träumen lassen, daß Gott droben im Himmel diesen Gottesdienst ansehe, und das Gebet erhöre, das vor diesem Bilde geschehen sei. Wie Jerobeam, da er die zwei goldenen Kälber aufrichtete, auch eine solche Predigt dabei gethan hat: „Dieses sind deine Götter, die dich aus Egypten geführt haben“, 1 Kön. 12, 28. 30. Er hat wohl gewußt, wie das Volk auch, daß die Kälber, die also gemacht waren, der Gott nicht war, der die Israeliten aus Egypten geführt hatte, 2 Mos. 32, 4., sondern sie haben geglaubt, daß der Gottesdienst, so bei solchen Götzen geschehen ist, Gott mußte wohlgefällig und angenehm sein.

152. Und der König hatte daran ein Gefallen, daß er dieses Bild also erdichtet und gemacht, darum daß die Ochsen oder Kälber das vornehmste Opfer waren. Er hat zwar kein Bild haben können, das sich besser zum Gottesdienst geschickt hätte, denn eben dieses Bild: darum sagt er nicht, daß der Ochse oder das Kalb Gott wäre, sondern er bindet den Gottesdienst an die Kälberbilder, dabei man Gott anrufen, ihm dienen und opfern sollte. Derhalben ist daselbst sinngeschichtlich das Gefäß für den Inhalt zu nehmen, nämlich, daß man Gott in und bei den Kälbern hat dienen sollen. Wie Gott verheissen hatte, daß er das Gebet derer, die ihn anrufen würden in der Hütte bei der Lade und Gnadenstuhle, auch erhören wollte. Also haben alle Völker ihrer Väter Gottesdienst behalten und denselben geschmückt, oder aber vielmehr verfälscht mit ihrer Abgötterei. Zu Ephesus war ein Bild der Abgöttin Diana, davon sie ein Gedicht gemacht haben, daß es göttliche Gewalt habe, daß es Gebet höre und annehme, Apg. Gesch. 19, 35.

153. Also beschuldigt Laban den Jakob, daß er nicht Gold oder Silber gestohlen habe, sondern daß er ihn und seine Gemeinde des ganzen Gottesdienstes oder Religion beraubt hätte, und, so viel an ihm war, hätte er ihn und sein Volk auch des Himmelreiches beraubt. Denn was ist das Andere, wenn man Einen der Religion und des Gottesdienstes beraubt, denn daß man ihn damit Gottes selbst, des ewigen Lebens und Seligkeit beraubt und ihn stracks vor den Teufel und in die Hölle hinein wirft? Diese Sünde ist so greulich und böse, daß sie nimmermehr entschuldigt noch gebüßt werden kann mit irgend welchen Worten oder Opfern.

154. Laban aber macht dies so überaus greulich und groß, nicht daß er von der Religion oder Gottesdienst so viel hielte, als ob damit der Seligkeit der Leute oder der Gnade Gottes ein Abbruch geschehen wäre; sondern nur daß er Jakob damit beschweren möge, daß er verhaßt werde und daß ihm seine Verwandten und Brüder dermaßen feind werden sollen, daß er mit ihnen nimmermehr könne ausgesöhnt werden; welche vielleicht das wohl gewußt haben, daß Laban mit Jakob etwas hart und unfreundlich umgegangen war, und daß Laban an solchem getreuen Knecht viel gelegen gewesen, sein Hab und Gut zu mehren. Derhalben, auf daß er die Verdienste Jakobs und seinen eigenen Geiz und böse Begierde zudecken und verhehlen möchte, so faßt er fein zusammen alle Tugenden oder guten Werke in der ersten und andern Tafel und schmäht sich selbst fein damit, daß nichts heiliger hat mögen angesehen oder gehalten werden nach allen Geboten Gottes denn eben der allerheiligste Nabal. Dagegen aber wendet er alle Sünden auf Jakob, legt sie wie in einem Haufen auf ihn; so er doch in Wahrheit nichts Anderes sucht oder meint, denn das Gold und Silber.

\quellseite{182}

155. Wie heutiges Tages die Papisten den Namen der Kirche hoch zu rühmen wissen, und daß es ihnen darum zu thun sei, daß sie gern die wahre Religion der lieben Väter erhalten wollen, daß man bei Christen und dem rechten unverfälschten Glauben bleiben möchte, zudem, daß man sich mit gebührlicher Ehre und rechtem Gehorsam gegen die Gesetze, gegen die Kirche und Obrigkeit erzeigen möchte (das alles, sage ich, wissen sie fein vorzuwenden), und sie dienen doch niemand, fragen auch nach nichts, denn daß sie nur den Bauch versorgen.

156. Darum ist dieses eine treffliche Beschreibung eines solchen Heuchlers, der sich selbst rühmt als den allergrößesten Heiligen, und beschwert den frommen heiligen Mann Jakob mit allerlei Schmach, mit großen Sünden und Betrug, als der da werth sei, daß er nicht mit menschlicher Strafe, sondern mit höllischer Marter und Pein sollte gestraft werden. Er nennt ihn einen Dieb und Mörder, der ihm seine Töchter entführt und sein Haus geplündert habe, nennt ihn untreu, undankbar, einen Kirchenräuber, und der die Religion und Gottesdienst zerstört habe. Er hätte nicht deutlicher können gescholten und gelästert werden.

157. Und daraus ist wohl zu sehen, was Laban von Jakob gehalten habe; wie gar übel und schändlich er mit ihm umgegangen sei; ja, daß er seiner wie eines armen leibeigenen Knechtes, oder vielmehr wie eines Esels oder sonst eines armen unvernünftigen Thieres gemißbraucht hat, nur zu seinem Genuß und daß sein Hab und Gut durch ihn möchte gemehret werden. Und jetzt schmähet er ihn auf das allerschändlichste, aber Alles mit Lügen und Unwahrheit, für seinen höchsten Fleiß, Mühe und Arbeit. Heißt das Gott gehorsam sein, der ihm verboten hat, er sollte mit Jakob nicht unfreundlich reden? Ja, dieses ist freilich das Allerunfreundlichste, nämlich, daß er ihn aller Sünden und Laster beschuldigt, die man nur erdenken oder sagen kann.

158. Dieses wird aber uns zur Lehre und Exempel geschrieben. Denn dieses ist also für und für eine Regel der Gottlosen, daß sie Gott nur dem äußerlichen Schein nach gehorsam sind, sie bekennen ihn mit Worten, aber mit Werken verleugnen sie ihn, Tit. 1, 16., und verachten seine Gebote. Laban enthält sich wohl öffentlich der Gewalt und legt die Hände nicht an ihn: aber wenn er es hätte thun mögen, er würde sein nicht verschont haben. Denn er ist ganz erhitzt von Begierde, daß er ihm gerne wollte Schaden thun. Es werden aber die Gottlosen etwas gedemüthigt und gezwungen durch Furcht der Strafe: darum thun sie wohl Buße, aber nur zum Schein und mit erdichteter und falscher Demuth. Wie Laban Alles mit lauter Heuchelei thut, und heimlich hat er mit den Zähnen geknirscht, daß er sich seiner Lust und Begierde nach nicht hat dürfen rächen.

159. Die andere Regel der Gottlosen ist diese: Alle Sünden, daran sie in ihren Herzen denken und die sie begehren zu vollbringen, messen sie den Heiligen bei; und dagegen alles, was die Heiligen thun und was ihre rechten Tugenden sind, schreiben sie sich selbst zu. Dieses ist eine gemeine Regel aller Heuchler und Gottlosen. Laban ist selbst ein Götzendiener, ein Dieb, ein Mörder, ein Menschendieb, ein Geizhals, und ungeheuer seiner Zeit; und läßt sich doch dünken, er sei der aller ehrlichste und frömmste Mann, und beschuldigt Jakob, der ganz unschuldig, einfältig und fromm ist, und dazu wohl sehr schönen Tugenden und großer Früchte, die er hervorbringt und beweiset der Kirche, dem Welt- und Hausregiment zu Nutzen und Besserung. Aber da ist niemand, der da sündigt und unrecht thut denn Jakob allein; und niemand ist heilig und fromm denn dieser Nabal.

160. Dieses ist sehr verdrießlich und unbillig. Es ist aber darum geschrieben, „auf daß wir durch Geduld und Trost der Schrift Hoffnung haben“, Röm. 15, 4. Es sündigt in dieser Welt niemand denn nur der eingeborne Sohn Gottes; dagegen aber ist niemand gerecht und fromm denn nur der Teufel: alles, was er sagt oder thut, ist recht und wohl gethan. Darum schreiben sich die Ketzer und Papisten allezeit Gerechtigkeit, Heiligkeit und Weisheit zu, als die allein den rechten katholischen Glauben schützen und vertheidigen. Die rechte Kirche aber muß diese höchste Schande fragen, daß sie ketzerisch, irrig, ärgerlich und aufrührerisch sei, und, wie St. Paulus sagt 1 Cor. 4, 13.: „Wir sind stets als ein Fluch der Welt, und ein Fegopfer aller Leute.“

161. Dieweil es nun Gott also gefällig ist und anders nicht sein kann, denn daß ein solch Volk sein muß, das um Gottes willen ein Fluch, „ein Spott der Leute und eine Verachtung des Volks“ sei, Ps. 22, 7.: so soll es uns nicht verdrießen, daß wir in der Welt für solche Leute gehalten werden. Denn wir haben einen sehr reichen Trost deß, der da gesagt hat Matth. 5, 11. 12.: „Selig seid ihr, wenn euch die Menschen um meinetwillen schmähen und verfolgen und reden allerlei Uebels wider euch, so sie daran lügen. Seid fröhlich und getrost, es wird euch im Himmel wohl belohnet werden.“ Es ist um eine kurze Zeit zu thun: dies Leben ist kurz und elend, aber die Freude und Seligkeit, darauf wir hoffen und harren, ist ewig. Dieweil es nun Gott also wohlgefällt, daß wir die kleine Heerde sein sollen, die gering und verachtet ist, so lasset uns mit geduldigen und fröhlichen Herzen solchen Jammer und Elend in diesem Leben tragen und leiden; wie Jakob die Schmach, Schande und Verachtung von den Labaniten, die Stolze und Heuchler waren, gelitten hat.

162. Wir mögen immerhin für Kirchenräuber und Zerstörer der Religion oder auch des gemeinen Regiments, ja, für die allergrößten Sünder auf der Erde gehalten werden: wir sollen aber diese Schmach leiden, ohne uns zu beklagen; und wir merken sie auch überminden um des Sohnes Gottes willen, der unser Vorgänger und Herr ist, welcher selbst auch solche Schmach gelitten und sie in allen Heiligen überwunden hat. Lasset uns nur aus der Zahl der Jakobiten und nicht der Labaniten sein. Und lasset uns das nicht anfechten, daß der Teufel bei den gottlosen Heuchlern für einen Gott gehalten und geehrt wird, als der heilig und gerecht sein muß, wie die ganze Welt urtheilt.

163. Nun ist aber noch eine Frage da von dem hebräischen Wort el, da Laban sagt: „Und ich hätte, mit Gottes Hülfe, wohl in viel Macht, daß ich euch könnte Uebels thun.“ Ich kann aber hier nichts Gewisses sagen, die Bedeutung dieses Wortes zu erklären; darum will ich es den hebräischen Rabbinen und Grammatikern befehlen. Wiewohl ich der Meinung bin, daß sie selbst die hebräische Grammatik nicht recht vollkommen verstehen, sonderlich bei diesem Worte; so gar seltsam mengen sie die Deutung unter einander und ist damit gar ein verworren Ding. Und sie tragen fast an zwanzig Worte zusammen, damit sie das Eine Wort erklären wollen; unter welchen Worten keines dem andern gleich ist.

164. Daran ist aber kein Zweifel, daß da die hebräische Sprache noch eine lebende Sprache war, viele Worte in der täglichen Umgangssprache ganz anders gelautet haben, denn sie jetzt lauten, da wir sie nur aus den Büchern herausklauben müssen. Wie ein Unterschied ist, wo man nach der Grammatik redet und wo man recht Lateinisch redet. Darum müssen wir nicht sowohl auf die grammatischen Regeln, als auf den Sprachgebrauch Rücksicht nehmen. Ich besorge aber, daß der Juden Grammatik nicht vollkommen sei, darum fehlen die Rabbinen oftmals, sonderlich in etlichen Worten.

165. Das Wort el heißt bisweilen Gott, und an diesem Ort heißt es auf Lateinisch: secundum Deum manus mea: „Mit Gottes Hülfe hätte ich wohl so viel Macht.“ Die Grammatiker und die stracks den Regeln folgen, sagen, es soll auf Lateinisch heißen: secundum facultatem: Es ist im Vermögen. Und das will ich zwar leichtlich zugeben; denn der Text im Propheten Micha am 2. Cap. V. 1. reimt sich auch damit, da er sagt: Fac sicut malum, quia Deus in manibus eorum: „Sie haben die Macht“, sie sind gewaltig, und wie wir es verdeutscht haben: „Sie sind die Herren.“ Denn, also steht auch 5 Mos. 28, 32.: Non erit, leol, fortitudo in manu tua: „Es wird keine Stärke in deinen Händen sein.“ Derhalben muß man es verstehen, daß es eine sonderliche Art zu reden sei.

166. Wie in der lateinischen Sprache viel Worte dahin gerathen sind, daß sie eine andere Bedeutung haben, denn die Regeln der Grammatik mitbringen; als da sind amabo, obsecro, quaeso (ach lieber), age (wohlan), apage sis, Edepol, Hercle, bei welchen Worten du in einen großen Irrthum gerathen würdest, wenn du sie stracks nach der Grammatik deuten wolltest. Denn das Wort amabo ist durch den gemeinen Gebrauch ein Ausruf geworden. Solche Veränderungen kommen in allen Sprachen vor. Die Deutschen sagen: „Gelt, ich will dir es bezahlen: da wird das Wort ‚Geld‘ als Adverb gebraucht. Item: Gott, du sollst es lassen; da ist das Wort ‚Gott‘ ein Ausruf geworden.

\quellseite{183}

167. Derhalben darf ich mir kein sonderliches Urtheil in der hebräischen Grammatik anmaßen, allein daß ich sehe, wie es sich die Grammatik noch sauer werden lassen, und haben es aber dermalen noch nicht herausgebracht, darum daß die Art und Weise zu reden nicht einerlei ist, der die Hebräer in ihrer gemeinen Sprache täglich gebraucht und wie es die Grammatiker ausgelegt haben. Die der Sprache unerfahren sind, zerpflücken die Bedeutungen der Wörter und richten so lauter Wirrwarr an; wie sie mit dem Wort schebeth auch thun, 1 Mos. 49, 10., davon sie über alle Maßen vielerlei Meinung gehabt haben. Denn Etliche deuten es, daß es Stab heißt, die Andern Scepter, die Andern legen es noch anders aus. Darum ist die hebräische Sprache zum größten Theil verloren gegangen und ist noch nicht vollkommen wiederum aufgefunden; und wenn wir das Neue Testament nicht hätten, so würden unsere Rabbinen nimmermehr den rechten Verstand aus dem Alten glauben können. Das Neue Testament hat viel dazu geholfen, daß die hebräische Sprache und das Alte Testament wiederum aufgerichtet und erklärt worden ist.

168. Ich will übrigens darüber nicht streiten, wo jemand das Wort el also verstehen wollte, daß es sollte Macht oder GOtt heißen. Das ist gewiß, daß durch die Figur, syncope genannt, el von dem Wort ajal herkommt, welches so viel heißt als Stärke; und von dem Abstractum, wie man es in der Grammatik nennt, machen sie das Denominativum oder Eigenschaftswort el, das heißt, stark. Wie wir es Jes. 9, 6. verdeutscht haben: Deus fortis, Kraft. Wo es allein steht, ohne das lamed, so heißt es Kraft, auf Lateinisch fortitudo vel Deus. Daß also GOtt von der Stärke el genannt wird.

169. Daran ist kein Zweifel, daß dieses Wort ganz mit Recht von GOtt gebraucht wird, wie man im 22. Psalm V. 2. sehen mag: Eli, eli etc.: „Mein GOtt, mein GOtt, warum hast du mich verlassen?“ Daher kommt auch elohim. Dieses etwa sagen die Hebräer von diesem Text. Wo sie nun den rechten Verstand getroffen haben, so ist es gut: wo nicht, so ist es das gemißte Zeichen, daß sie noch dermalen der hebräischen Sprache keinen rechten und gewissen Verstand haben. Aus dem Alten Testament können sie diese Auslegung nicht bestätigen; aber aus dem Neuen Testament wird es bewiesen, daß dieses Wort el so viel heißt als GOtt.

170. Es wird aber an diesem Ort das lamed vor das Wort el gesetzt, le el, das lautet im Lateinischen also: Adest, vel in promptu est manus mea; auf Deutsch: Da ist meine Hand oder Macht. Ob es nun hier GOtt heiße, oder aber allein Stärke, daß es ein Abstractum wäre, wie man es in der Grammatik nennt, und daß aus dem Namen GOttes ein Ausruf geworden sei, und daß es so viel heiße als Macht oder Vermögen, darüber will ich nicht hart streiten, sondern will dasselbe unentschieden lassen. Wir haben es im Deutschen etwa also wiedergegeben: „Und ich hätte, mit GOttes Hülfe, wohl so viel Macht, daß ich euch könnte Uebels thun.“ Dieses sei nun gesagt von der Grammatik. Er sagt, wie es auf Lateinisch lautet: Adest mea facultas; vel a Deo manus mea: Ich habe von GOtt so viel Macht etc.

171. Aber alles, was der Heuchler redet oder was er hält von dem Namen GOttes, ist alles miteinander erdichtet und erlogen Ding; denn die Heuchler haben keinen GOtt. Ich hätte wohl so viel Macht, sagt er, daß ich mich an dir rächen wollte, darum daß du mir meine Töchter entführt und meine Götter gestohlen hast. Aber deines Vaters GOtt hat mir es verboten, und mich hart bedroht, ich sollte dir keine Gewalt thun.

172. Was ist es aber für eine Sünde, die ich gethan habe, möchte Jakob sagen, die du an mir rächen wolltest? Antwort: Ei, daß du nicht hast ewig mein leibeigener Knecht und mein Esel sein und die Knechtschaft nicht länger ertragen wollen, daß du Tag und Nacht schwitzen mußtest, daß du immer laufen und arbeiten, und dazu von meinen Söhnen und von mir noch gehöhnt und geschmäht werden solltest: das ist die so große Sünde, die du gethan hast. Du solltest mich dieses Dienstes nicht beraubt haben, darin ich dich und meine Töchter gern immer behalten hätte. Das ist wahrlich eine große greuliche Sünde, daß sich Einer nicht in die verächtlichste Knechtschaft des grausamsten Tyrannen begeben will.

173. Aber das ist die rechte Eigenschaft der Heuchler, daß sie andere Leute pflegen zu beschuldigen, darum daß sie ihnen im Wege liegen, und sie hindern, daß sie ihr Mütlein nach ihrem Gefallen nicht kühlen oder ihre Lust nicht büßen mögen. Dieser haben wir nun gehabt die Beschreibung des Nabal, welcher ein rechtes Bild ist des Papstes und eines Heuchlers. Ein rechter Papst und falscher hypocrita ist er gewesen. Nun wollen wir einen frommen Mann auch hören.

\subsection{II}

\begin{quote}
\textbf{V. 31–33. Jakob antwortete und sprach zu Laban: Ich fürchtete mich, und dachte, du würdest deine Töchter von mir reißen; bei welchem aber du deine Götter findest, der sterbe hie vor unsern Brüdern. Suche das Deine bei mir, und nimm hin. Jakob wußte aber nicht, daß sie Rahel gestohlen hatte. Da ging Laban in die Hütte Jakobs, und Leas, und der beiden Mägde, und fand nichts.}
\end{quote}

174. Dieses ist eine kurze Entschuldigung auf den ganzen Haufen von Lästerungen, so Laban wider den frommen heiligen Mann geredet und ausgespieen hatte. Was sollte ich thun? sagt er, „ich fürchtete mich“. Er bekennt fein demüthig, und entschuldigt sich nicht ausführlicher des Raubes halben; er frißt aber alle Schmach und Unbilligkeit in sich und schweigt davon stille, und sagt nur, er habe sich gefürchtet. Wenn er etwas heftig gewesen wäre, wie viele Leute darüber ungeduldig sind, so sie sollen Unrecht leiden, so hätte er alle diese Scheltworte wiederum auf ihn legen und ihn damit anfahren mögen. Er würde gesagt haben: Höre auf, und fluche einem solchen Manne nicht, der sich um dich wohl verdient gemacht hat. Denn du bist selber, wie dich GOtt und alle Menschen dafür halten, ein Kirchenräuber, ein Mörder, ein Tyrann, und ein Ausbund von Geiz, der du deine Töchter verkauft und sie wie die allerverächtlichsten leibeigenen Mägde und unvernünftigen Thiere gehalten hast. Ja, das noch mehr ist, du hast ihnen für ihren sauren Schweiß und Arbeit nicht gedankt ihren gebührlichen Lohn zu geben, sondern hast ihnen denselbigen mit eigenem Gefängniß erstattet und belohnen wollen. Dieses wäre ein unendlicher Stoff zum Streit und Zank gewesen.

175. Derhalben werden wir hier ermahnt, daß wir auch die Schmach, item, die Verachtung und den Stolz der Widerfacher mit Geduld leiden und ertragen sollen; man soll aber doch gleichwohl die Wahrheit nicht verschweigen. Und der Muthwille, den die Widerfacher an uns mit falschen und lästernden geübt haben, wird an ihnen ungestraft nicht bleiben, wiewohl wir ihnen nicht wieder fluchen. Denn es ist Einer, der es an ihnen suchen und richten wird.

176. Darum sagt nun Jakob also: Was ich in dieser Sache gethan habe, das habe ich aus Furcht und Schrecken gethan. Also beschuldiget er zwar Laban nicht, und gibt doch so viel zu verstehen, daß er billig zu strafen und zu verklagen sei, dieweil sich Jakob hat fürchten müssen, daß ihm sein Schwiegervater die Töchter mit Gewalt nehmen möchte, die er ihm zu Weibern gegeben hatte. Das ist wahrlich sehr greulich. Derhalben beschuldiget er ihn sein heimlich und mit einer schönen Redewendung der größten Sünden, also daß er ihn viel heftiger verklagt damit, daß er sich entschuldiget, denn er sonst mit offenbarem Schelten hätte thun können.

177. Denn also pflegen die Entschuldigungen der Gottlosen die allerschwersten Anklagen der Gottlosen zu sein. Wie jetzt der Papst und sein Anhang uns auch beschuldiget, als sollten wir Ketzer, Gotteslästerer, Ungehorsame und Aufrührer sein. Wenn ich aber antworten werde: „Ich bin kein Ketzer, bin auch kein Aufrührer, wie Christus den Juden klärlich und mit ausdrücklichen Worten antwortet, Joh. 8, 49.: ‚Ich habe keinen Teufel‘ (welches wahrlich geschehen muß; denn wir sollen nicht Beifall geben und dazu nicht stille schweigen, wo die Lehre oder göttliche Wahrheit geschmäht und gelästert wird): da zeihe ich den Papst mit einer kurzen Widerlegung aller der Sünde und Schmach, damit er uns beschuldiget hat.

178. Derhalben ist dieses eine treffliche Entschuldigung; denn er beruft sich heimlich auf das Gewissen des Heuchlers selbst, und beschuldiget ihn aus allen seinen Thaten, daß er ein greulicher Tyrann sei. Dein eigen Gewissen, sagt Jakob, gibt davon Zeugniß, daß du mit uns sehr unglücklich umgegangen bist, daß ich auch derhalben gute Ursache gehabt habe, davon zu denken, wie ich davon fliehen möchte. Denn ich dachte bei mir selber also: Wenn Laban erfahren wird, daß ich fliehen will, so wird er mir meine Weiber, und alles nehmen, was ich habe, und mich in das Gefängniß werfen. Gesetzt aber daß einem frommen Mann, daß er mit einem getreuen Menschen, und der da heilig, fromm, fleißig und wacker ist, im Hause alles zu versorgen, was darin auszurichten ist, dermaßen soll umgehen, daß er sich vor der Gewalt und Unbilligkeit seines Schwiegervaters hat fürchten müssen? nämlich, daß derselbe der nahen Freundschaft, damit sie einander verwandt sind, und auch der Pflicht und natürlichen Neigung, so er zu seinem Schwiegersohn und Töchtern billig haben sollte, vergessen, und die Weiber ihrem Manne mit Gewalt nehmen möchte? Ist das an ihm zu loben oder zu entschuldigen? Ja, es ist wohl eine große Schande und ist wider alle Vernunft und Ehrbarkeit. Denn die Töchter, die man ehlich einem Manne vertrauet und gegeben hat, sind nicht mehr unter der Gewalt ihres Vaters: der Vater hat nicht Macht, die Töchter wieder zu fordern oder sie dem Mann zu nehmen. Darum ist dieses eine greuliche Tyrannei gewesen, daß der gute fromme Mann in der Gefahr gewesen ist, daß er wiederum seine lieben Weiber und Kinder verlieren möchte.

\quellseite{184}

179. Und es läßt sich ansehen, als habe ihm der Tyrann etlichemal damit gedrohet. Wenn ich wüßte, wird er gesagt haben, daß du damit umgingest, daß du wolltest davon ziehen, oder daß du meine Sache nicht getreulich ausrichtetest, so wollte ich dir meine Töchter alsbald wieder nehmen. Dasselbe hat zwar Mose mit Worten nicht ausgedrückt; aber es ist doch sehr glaublich, und läßt sich ansehen, als habe der Heilige Geist anzeigen wollen, daß Jakob nicht vergeblich angeführt habe, seine Flucht oder Abzug zu entschuldigen, daß er sich gefürchtet habe.

180. Ei, du heiliger Mann Laban, wie sollen wir nun dich recht loben und preisen? Ist es denn recht, daß ein Vater oder Schwiegervater also gesinnet sein soll gegen seinen Schwiegersohn und Töchter, die ihrem Mann folgen und die unter dem Hausgesinde ihres Vaters für Leibeigene sind gehalten worden, durch welcher Fleiß und Arbeit sein Gut und Habe sehr zugenommen hatte, daß du diejenigen denkst mit Gewalt von ihrem Mann zu reißen? Wahrlich, dieser Heuchler ist das greulichste Ungeheuer, das da leben mag, welches mit seiner Grausamkeit alle unvernünftigen Thiere und alle Tyrannen weit übertrifft; ja, er ist grausamer, denn die Grausamkeit selbst sein kann.

181. Wo nun irgend auch ein klein Fünklein Freundlichkeit oder Gütigkeit an ihm gewesen wäre, so hätte er sich doch bewegen und erweichen lassen sollen die große Frömmigkeit, Gehorsam, Treue und Geduld des frommen getreuen Jakob. Aber er ist durch das alles nicht im geringsten bewegt worden. Denn darum jagt er seinem Schwiegersohn nach, daß er ihn wiederum in das Gefängniß der Dienstbarkeit bringen, oder aber ihm die Töchter mit Gewalt nehmen möge. Davor hat sich nun Jakob gefürchtet, und will so viel sagen: Du bist selber Schuld, daß wir haben fliehen müssen, daß du uns so jämmerlich unterdrückt, und gedroht hast, wenn ich wegziehen würde, daß du mir alsdann Alles nehmen wolltest; darum habe ich mich gefürchtet.

182. Zum Andern, daß du mich nun beschuldigst als einen Kirchenräuber, da will ich also thun, sagt er: „Bei welchem du deine Götter findest, der sterbe hie vor unsern Brüdern.“ Da beschreibt nun der Heilige Geist, wie es menschlicher Weise um die Heiligen steht. Denn wie hoch und stark sie im Glauben und Geist sind, so irren und fehlen sie doch nicht allein in vielen Dingen, sondern fallen auch nach dem Fleisch. Jakob denkt, er habe gar keine Schuld, könne nicht beargwohnt werden und es habe keine Gefahr mit ihm, und in solchem Theil sei Laban nun genug geschehen; darum erbietet er sich, daß der solle getödtet werden, bei welchem er die Götzen finden werde. Aber auf die Weise wird sein liebes Weib von ihrem eigenen Manne dem Tode übergeben. Denn wo dem unser HErr GOtt nicht zuvorgekommen wäre und gehindert, so hätte Jakob unbedachtsam seine liebe Rahel auf die Fleischbank geliefert. Das ist zu weit gegangen. Jakob gibt allzu viel nach; aber er thut es unwissentlich. Solche Sünde haben die Gläubigen, die bisweilen nach dem Fleisch aus Unwissenheit oder durch Irrthum fallen; wie Rahel auch hernach weiblich lügen wird, ihr Leben zu erretten.

183. Mit dieser Antwort sollte übrigens der heilige Mann Laban billig zufrieden gewesen sein und gesagt haben: Was soll ich viel suchen oder forschen? Ich höre, wie hoch ihr eure Unschuld betheuert, und nehme solches auch an. Es sind ja meine Töchter und Kindeskinder, gegen welche ich mich zuvor etwas hart und unfreundlich gehalten habe; ich will euch nun hinfort nicht mehr beschweren oder betrüben, will euch auch keines Diebstahls zeihen. Das hätte einem Schwiegervater und frommen Mann wohl angestanden; aber er kann nicht ruhen, denn er ist voll Teufel. Es hätte sich sonst Jeder an dieser Entschuldigung wohl genügen lassen, wenn er hätte gehört, daß Einer sagte: Ich habe mit deinen Götzen nichts zu schaffen. Nun setzt aber Jakob noch etwas mehr dazu, auf daß er ja nicht möge beargwohnet werden, als ginge er mit List um: Ich gebe es zu, sagt er, daß du suchest und Alles untersuchest. Die Macht würde ihm ein Anderer nicht leichtlich gegeben haben. Ich hätte gesagt: Ich habe deine Götzen nicht, dabei hätte ich es bleiben lassen, und hätte ihm die Macht nicht gegeben, daß er unter all meinem Gut hätte suchen mögen; oder aber, ich hätte selbst allein gesucht. Darum ist es sehr unbillig gewesen, daß diesem scheußlichen Ungeheuer so großer Muthwille hat sollen gestattet werden. Denn war es daran nicht genug, daß er hört, daß Jakob seine Götzen nicht hatte, und daß er sich so hoch erboten hat, man solle den tödten, der sie gestohlen hätte?

184. Aber siehe doch, wie der heillose Mann sich gar nicht schämt; denn er fährt fort und sucht in allen Hütten. Darum kann ich es nicht lassen, ich werde heftig darüber erreget, und ist mir leid, daß ich dieses alles mit Worten nicht erreichen kann und nicht genugsam mit seiner rechten Farbe herausstreichen. Ich würde ihm nichts mehr zugegeben haben, wenn er nicht hätte wollen mit meinem Bekenntniß zufrieden sein, daß so aufrichtig und beständig war. Jakob gibt ihm eigentlich mehr nach, denn er billig sollte gethan haben.

185. Darum geht er nun erst in Jakobs Hütte; und wäre daran genug gewesen, wenn er daselbst nichts gefunden hätte. Er geht aber in der Lea Hütte auch und über das noch in der beiden Mägde Hütte. Zuletzt war noch eine da, nämlich der Rahel, die muß er auch beschmuttern. Und da ist Rahel ihres Leibes und Lebens halben in große Gefahr gekommen.

Darum ist nun der Heilige Geist dabei, und findet bald einen Rath gegen die Erlaubniß Jakobs, der da gesagt hatte: „Bei welchem du deine Götter findest, der sterbe“ etc. Es war zu viel gegeben, es wird aber eine Abhülfe folgen; denn wo Jakob schlummert, da wacht GOtt und die lieben Engel.

186. Und das ist GOttes Werk und seine Kunst, daß er das andere und bessere, was Jakob aus Irrthum verderbet hatte. Er kann böse Sachen gut machen, wenn wir es verderbt und verwahrlost haben. Wie in so großer Schwachheit, so an den Menschen ist, die Heiligen nicht ohne viel und große Sünde, darein sie fallen, sein können. Ich habe wahrlich oftmals viele Dinge unweislich und thörlich ausgerichtet, darüber ich hernach heftig erschrocken bin, und ich konnte nicht sehen, wie ich aus solchen Sachen, die durch meine Thorheit verwirrt und verderbt waren, wiederum möchte frei werden und daraus entkommen: aber der HErr hat eine solche Weise und Weg getroffen, daß es gebessert worden ist, was ich versehen und verderbt hatte.

187. Also ist Rahel eine große und gewisse Gefahr zu errettet gewesen, weil ihr Mann so unweislich zugegeben hatte, daß Laban seine Götzen suchen möchte, wo er wollte. Denn so grausam und tyrannisch wie dieser sehr schelmische Heuchler war, würde er entweder wider die Töchter gewüthet, oder darauf gedrungen haben, daß der Schwiegersohn sich ihm hätte müssen zur Knechtschaft begeben, sie damit zu erretten und am Leben zu erhalten. Aber GOtt rettet und führt seine Heiligen also, daß sie wohl etwa irren und fehlen mögen, aber daß es gleichwohl mit ihnen ein gut Ende nehmen muß, oder aber muß ja ohne großen Schaden abgehen. Denn es müssen den Auserwählten, und die da glauben, alle Dinge zum Besten dienen, auch ihre Fehler und Sünden. Das ist gewiß. Denn GOtt pflegt aus Nichts Alles zu machen; darum kann er auch aus dem, das böse ist, Gutes hervorbringen. Wie Augustinus vornehmlich an solchen Gedanken Lust hat, daß er sagt: So groß ist die Gütigkeit GOttes, daß er nichts Böses würde geschehen lassen, wenn er es nicht wiederum gut machen könnte. Jakob hatte zwar sehr unweislich geirrt, und so viel an ihm war, hatte er sein liebes Weib in die Gefahr des Todes oder ja zum wenigsten in große Schmach und Schande gebracht; es wird ihm aber von GOtt ein Rath und Anschlag angezeigt, der so listig ist, daß beide Laban und der Teufel selbst damit betrogen werden.

\quellseite{185}

\section[Sechster Theil]{Sechster Theil. Wie Laban seiner Söhne wegen Jacobs Hütten durchsucht und nach dieser Durchsuchung ernstlich von Jakob gestrafet wird}

\subsection{I}

\begin{quote}
\textbf{V. 33–35. Und ging aus der Hütte Leas in die Hütte Rahels. Da nahm Rahel die Götzen, und legte sie unter die Streu der Kameele, und setzte sich darauf. Laban aber betastete die ganze Hütte, und fand nichts. Da sprach sie zu ihrem Vater: Mein Herr, zürne nicht, denn ich kann nicht aufstehen gegen dir; denn es gehet mir nach der Frauen Weise. Also fand er die Götzen nicht, wie fast er suchte.}
\end{quote}

188. Hier greift der Satan Rahel an die Gurgel. Denn Laban geht in ihre Hütte und durchsucht darin alles Einzelne auf das allergenaueste. Nun hat sie aber die Götzen in den Kleidern oder Aermeln, oder auch im Strohe nicht verborgen können, da war nichts mehr vorhanden, denn nur ihr eigener Leib, damit sie die Götzen bedeckt hat. Darum legt sie nun die teraphim unter die Streu der Kameele, nicht daß sie daselbst sicher genug hätten liegen können; denn Laban hätte die Streu leichtlich untersuchen können: sondern sie hat sie über das mit dem garstigsten Theile des Leibes bedeckt; denn sie hat sich auf die Streu gesetzt. Wiewohl das auch nicht ohne Gefahr ist. Denn Laban hätte mögen argwöhnen (wie der Geiz und Heuchelei pflegt argwöhnisch zu sein), es wäre irgend ein Betrug darunter verborgen; er möchte gesagt haben: Das muß nicht ohne Ursache geschehen sein, daß ich dich hier sitzen finde an einem ungewöhnlichen Orte. Wie, wenn meine Götzen hier verborgen lägen? Und wo er sie hätte heißen aufstehen, daß er in der Streu hätte suchen mögen, so wäre der Diebstahl an den Tag gekommen und entdeckt worden.

189. Sie erdenkt aber eine sehr feine Lüge: sie gibt vor, daß es ihr jetzt gehe nach der Frauen Weise. Nun begibt es sich aber also und thut um der Gesundheit willen noth, daß zu solcher Zeit ein Weib nicht auf der Erde, sondern viel lieber auf dem Bette oder auf dem Strohe sitze, wenn sie ja kein Bett haben kann. Also hat Rahel in der äußersten Noth und Gefahr bald Hülfe und Rath gefunden; wie das die Erfahrung ausweist, daß die Weiber sehr schnell guten Rath erfinden. Derhalben hilft nun die Natur oder Verstand und Geschicklichkeit, so die Weiber pflegen zu haben, der Rahel auch etwas. Wiewohl aber solche Geschicklichkeit an den Weibern gefunden wird, so ist doch das weibliche Geschlecht von GOtt nicht verordnet zum Regiment weder in der Kirche noch sonst in weltlichen Aemtern, dazu sonderlicher großer Verstand und guter Rath gehört; sie sind aber dazu berufen, daß sie das Haus versorgen und darin fleißig zusehen sollen; denn je länger sie von großen und hohen Dingen rathschlagen, je mehr sie die Sachen verwirren und verhindern. Aber der erste Antrieb ihrer Natur pflegt in plötzlicher Gefahr sehr gut und glücklich zu sein.

190. Also wird nun Rahels Vater durch solche weibliche Lüge und durch redliche Ursache getäuschet und unter solchem Schein betrogen. Denn er darf an den Orten nicht suchen, welche die Natur verborgen hat, und kann auch weiter nichts dagegen aufbringen, sondern wird durch natürliche Scham abgeschreckt. Aber GOtt hat dasselb mitgewirkt; denn der hat ihm seine Sinne verkehret und ihn mit Blindheit geschlagen, daß er nicht genauer gesucht hat. Er glaubt, es sei wahr, was ihm seine Tochter sagt, nämlich, daß sie um ihrer Krankheit willen habe wollen an einem besondern Ort allein sitzen. Darnach war es nicht wohl glaublich, daß sie sollte auf den Götzen sitzen. Denn er hatte gemeint, daß solches ihrer Krankheit halben nicht wohl geschehen könnte. Er hat aber auch nichts weniger glauben können, denn daß seine Tochter, die in seinem Hause auferzogen war, die Götter und den Gottesdienst ihres Vaters stehlen, und dieselbigen so gar verächtlich halten sollte, daß sie dieselben unter die Streu der Kameele legen und darauf sitzen würde. Meine Tochter, hat er gedacht, würde ja meinen Gottesdienst, darin sie von Kindheit aufgezogen ist, nicht so gering achten. Also täuscht er sich selbst ganz fein. Die Tochter hat die Lüge dazu gethan, welches nicht eine Scherzlüge oder schädliche, sondern eine sehr dienstliche und nöthige Lüge war, auf daß sie damit dem Vater Ursache gäbe, daß er sich selbst betrügen mußte.

191. Das ist nun GOttes Werk und Regiment bei gefährlichen Lagen und Thaten der Heiligen. Denn wo er selbst die Widersacher nicht mit Geiz und Blindheit schlüge und ihnen ihre Sinne verwirrte, daß sie daran nicht denken, daran sie am allermeisten denken sollten: so würden oftmals auch die allerlistigsten Anschläge, die man hätte erdenken mögen, gar nichts helfen.

192. Das ist aber gleichwohl ein wunderlich Ding, daß Rahel also stark und beständig gewesen ist, daß sie die Götzen, die sie zuvor für sehr heilig gehalten, hat verachten können, unangesehen der Gefahr, so ihr von ihrem zornigen Vater darüber begegnet ist. Denn wir gehen ja ehrlicher mit unserm Kirchenbildern und Altären um, denn Rahel mit den Götzen ihres Vaters umgegangen ist, die sie in die Streu verbirgt und unter ihren Leib legt. Und es ist eine wunderliche Rede, damit sie sich vor ihrem Vater entschuldigt, daß sie vor ihm nicht könne aufstehen. Sie redet ihren Vater an, und nennt ihn doch nicht Vater, sondern Herr: „Mein Herr, zürne nicht“ 2c. Das ist wahrlich eine schwere Sünde, daß sie ihren Vater nicht fürchtet und ehrt, wie sie billig thun sollte nach dem vierten Gebot. Und wenn die Weiber ihre Krankheit haben, sind sie so schwach nicht, daß sie nicht sollten aufstehen oder gehen können.

193. Mose erzählt dieses nicht vergeblich, und läßt sich ansehen, als ob es schier also laute, daß sie ihres Vaters gespottet habe als eines Götzendieners, der so genau und unsinnig seine Götzen suchte. Aber ich kann davon nichts Gewisses sagen. Sie hat es auch wohl darum thun können, daß der Schein so viel desto größer wäre, daß sie vermassen sich gebärdet hat, als wäre sie nicht allein gar bereit und geneigt gegen ihrem Vater aufzustehen als eine gehorsame Tochter, sondern daß sie sich auch als eine Magd gegen ihren Herrn mit Furcht und ängstlicher Sorge erzeigt, warum daß sie ihrer Krankheit halben verhindert werde, daß sie ihrem Vater seine gebührende Ehre nicht erzeigen könne. Denn die Weiber pflegen in der heiligen Schrift ihre Männer Herren zu nennen, die Kinder aber nicht also. „Mein Vater“, sagt Elisa zu Elia 2 Kön. 2, 12. Was es nun aber für eine Ursache gewesen sei, daß sie ihn einen Herrn nennt, so ist dieser Name doch dahin gerichtet gewesen, daß die Lüge damit geschmückt und beschönigt werden möchte. Da Laban aber vergeblich also gedroht und in allen Hütten seine Götzen gesucht hat, wird er gar unsinnig und böse: darum widerlegt ihm Jakob nun seine Scheltworte etwas schärfer und mit mehr Worten; denn er wird nun wiederum seine Wohlthaten und die Undankbarkeit Labans herzählen.

\subsection{II}

\begin{quote}
\textbf{V. 36. 37. Und Jakob ward zornig, und schalt Laban, und sprach zu ihm: Was hab ich mißgehandelt oder gesündiget, daß du so auf mich erhitzt bist? Du hast alle meinen Hausrath betastet. Was hast du deines Hausraths funden? Lege das dar vor meinen und deinen Brüdern, daß sie zwischen uns beiden richten.}
\end{quote}

194. Endlich wird Jakob auch zornig, und verläßt sich auf die göttliche Verwahrung, so zu Laban geschehen war, und daß GOtt Laban selbst auch gestraft hätte; denn es thut ihm die große Unbilligkeit wehe, und verläßt sich gleichwohl mehr auf GOttes Urtheil denn auf seine eigene Unschuld. Daraus faßt er Muth und Trost, und will also sagen: GOtt hat mein Elend und Trübsal gesehen, und hat dich um deiner Tyrannei willen gestraft, die du an mir geübt hast; denn wo ich vor GOtt irgend einer Sünde wäre schuldig gewesen, so hätte ich die Strafe müssen leiden und du nicht. Nun aber bin ich meiner Unschuld gewiß, und weiß, daß die Arbeit, so ich mit meinem Dienst gethan habe, GOtt wohl gefallen und daß ihm deine Undankbarkeit und alle dein Thun übel gefallen hat. Das ist Jakobs Trost und sein Ruhm, daß er GOttes Urtheil auf seiner Seite hat, der ihn für gerecht hält und seinen Widersacher für ungerecht. Dazu kommt aber auch noch das sehr gewisse Zeugniß der Erfahrung und des Gewissens Labans. Denn er beschreibt das ganze Leben, beide sein selbst und auch seines Schwiegervaters, und bewähret es, daß daraus Laban auch überzeugt und überwiesen werde, daß er ein ungerechter Mann und ein greulicher Tyrann sei.

\quellseite{186}

195. Darum antwortete er nun dem Narren nach seiner Narrheit; aber doch also, daß er gleichwohl auch hält, das Salomo noch dazu sagt Spr. 26, V. 4. 5., nämlich, „daß du ihm nicht auch gleich werdest; sondern, daß sich der Narr nicht dünken lasse, daß er weise sei.“ Denn Salomo heißt uns, daß wir die Sünder und Narren unterrichten, strafen und schelten sollen zur Zeit und zur Unzeit, daß sie sich nicht dünken lassen, daß sie weise seien, das ist, daß sie in der Meinung nicht bleiben, als hätten sie daran recht gethan, daß sie gesündigt haben; sondern man soll ihnen die Strafe Gottes vorhalten und ihnen ihr Leben und böse Thaten aufdecken, daß sie fühlen und verstehen mögen, daß sie gesündigt haben, und daß sie Buße thun. Wer aber den Narren also straft, daß er ihm gleich wird, als, wenn ich ihn auch wieder belügen und ihn lästern wollte: so würde ich ihm gleich sein. Man soll ihn aber also strafen, daß er mit seiner Thorheit zu Schanden werde; und das ist eine göttliche Strafe. Dein eigen Gewissen und dein ganzes voriges Leben straft dich; derhalben siehe zu, daß du deine Sünde erkennest etc.

196. Eine solche Strafe gebraucht Jakob wider Laban und sagt also: Was habe ich gethan? womit habe ich übertreten? was ist meine Sünde? Du hast über die Maßen viel Scheltworte und Beschuldigungen über mich ausgesprochen; du hast mich ausgezupfen für einen Dieb, einen Viehdieb, einen Kirchenräuber, und daß ich dir deine Religion und Gottesdienst soll zerstört haben: der Sache beweise mir nur eines und überführe mich desselben, so will ich meine gebührliche Strafe dafür leiden; wo du aber das nicht thun kannst, so gibst du jetzt deine Leichtfertigkeit und deinen Muthwillen mit Augen an den Tag. Denn mit deinem Stillschweigen bekennst du, daß du gelogen hast; ja, du wirst durch Gottes und deines eigenen Gewissens Zeugniß überwunden. Du kannst nichts vorbringen, darin ich gesündiget hätte, oder daß ich dessen etwas versäumt haben möchte, das ich hätte thun sollen. Alle dein Schelten ist eine lautere Lästerung, es ist eitel erlogen Ding und nur eine Unsinnigkeit gewesen. Denn du wirst durch dein eigen Urtheil und Gewissen zu Schanden gemacht und verdammt: ich aber habe Gott zum Zeugen meiner Unschuld und dazu dein eigen Gewissen. Also soll man den Narren mit Kolben laufen.

197. Wo sind nun die großen geleistenden Worte und die erdichteten Laster: Du hast mir mein Herz gestohlen, du bist flüchtig geworden, wie der Dieb, V. 26., solche Worte mit großem Zorn und Widerwillen ausgestoßen? Sie sind alle verschwunden wie der Rauch. Das hat aber der närrische Laban nicht verstehen wollen, da ihn der Herr im Traum ermahnt, daß er mit Jakob nicht hart oder unfreundlich reden sollte. Darum setzt nun Jakob hier die Glosse hinzu und legt ihm die Worte des Herrn aus, und sagt also: Du hast mich unbillig und thörlich solcher großen Sünde und Laster beschuldigt, und die Stücke alle, um welcher willen du mich verklagst, gehören billiger auf dich; denn das hat Gott gesagt und dein eigen Gewissen überzeugt dich auch davon.

198. Also hat der Narr Laban oder Nabal müssen gezäumt und gehebelt werden, daß er sehen konnte, daß er geirrt und gesündiget hätte. Und das heißt, dem Narren antworten nach seiner Narrheit, daß er sich nicht weise dünken lasse, wie Salomo lehrt in den Sprüchen am 26. Capitel V. 5. Denn man soll zu den teuflischen Scheltworten, so wider die wahre Religion und wider GOtt selbst erdichtet und geredet werden, nicht stille schweigen, sondern man soll sie strafen und widerlegen. Und solche Widerlegung ist keine Sünde: es ist kein Zorn, der unbillig oder sträflich wäre; sondern ist ein heiliger und gerechter Zorn und ein christlicher Eifer, der auf das gottlose Wesen und auf die Sünde zornig ist.

199. Solche Strafe ist aber dazu nütze, daß die Leute im Irrthum und in ihren Sünden nicht verharren, und also des Teufels leibeigen bleiben, der sie dazu getrieben und besessen hat; sondern man soll sie strafen, daß sie mögen bekehrt und von der Gewalt und Tyrannei des Teufels erlöst werden. Wie wir den Papst strafen als den Antichrist und Verführer der ganzen Welt; ja, wir strafen und verdammen allerlei gottlos Wesen und Abgötterei. Wir zürnen darüber, nicht aus eigener Rachgier, sondern aus rechtem Eifer, und daß wir in unserem Gewissen erhitzt sind, die Ehre Gottes zu retten und zu erhalten.

200. Jakob aber bringt hart auf seinen Schwiegervater und stopft ihm nun das Maul mit Fragen, daß er zu ihm sagt: Was habe ich dich gesündiget? Was habe ich verschuldet? Welches ist denn die Sünde, die ich soll begangen haben? Hier will ich denen, die der hebräischen Sprache Verstand haben, empfehlen, daß sie den Unterschied der zwei Wörter, pescha und cheth, untersuchen; denn die Rabbinen bringen dreizehn Worte zusammen, welche alle Sünde und Uebertretung heißen, die nicht hieher gehören. Aber diese zwei Wörter, die Mose gebraucht, pflegt man also zu unterscheiden: das Wort pescha heißt die Schuld, wenn man nicht thut, was man zu thun schuldig ist (culpa omissionis); cheth aber heißt die Sünde, wenn man thut, was man nicht thun sollte (peccatum commissionis): und diese Worte begreifen in sich die Erbsünde. Pescha, wie gesagt, heißt, wenn man nicht thut, was man thun sollte (omittere facienda); cheth aber, wenn man thut, das man nicht thun sollte (facere omittenda).

201. Also steht 2 Kön. 3, 5.: Praevaricatus est Moab in Israel, das ist: „Die Moabiter fielen ab von Israel.“ Denn die Moabiter waren dem Volke Israel unterworfen, daß sie ihm haben Tribut geben müssen; aber sie sind von ihnen abgefallen, das ist, sie haben ihnen nicht mehr gehorsam sein und Tribut geben wollen. Item, also stehet im Propheten Jesaia am 1. Capitel V. 2.: Praevaricati sunt in me: „Sie sind von mir abgefallen“, das ist, sie sollten mich fürchten und ehren, sollten mir als ihrem Vater gehorsam sein; aber den Gehorsam, den sie mir schuldig sind, lassen sie fahren. Ich habe sie zu meinen Kindern angenommen, daß sie mir dienen sollten; aber sie fallen von mir ab und erwählen sich fremde Götter. Was sie nicht thun sollten, das thun sie, und was sie schuldig wären zu thun, das lassen sie liegen. Item, also steht auch Jer. 2, 13.: Duo mala fecit populus meus etc.: „Mein Volk thut eine zwiefache Sünde: mich, die lebendige Quelle, verlassen sie; und machen ihnen Brunnen“ etc. Mich sollten sie ehren und mir dienen als dem lebendigen Brunnen; aber mich lassen sie fahren: darnach suchen sie andere Brunnen; das sollten sie nicht thun.

202. Also sagt Jakob hier auch: Es sind zwei Stücke, darum du mich hättest strafen mögen, nämlich, wo ich etwas gethan hätte, das ich billig nicht sollte gethan haben; oder aber, wo ich etwas versäumt hätte, das mir zu thun gebühret hätte: weise mir nur eins von beiden nach. Du sagst, ich sei dir ungehorsam geworden, ich habe deine Götzen gestohlen, habe dir deine Töchter entführt; welches alles offenbar Lügen sind, und du mußt selbst mit Worten und Werken meine Unschuld beweisen und bezeugen. Dazu kommt das göttliche Zeugniß auch und dein eigen Gewissen: ja, du bist selbst beiderlei Sünden schuldig; denn du hast damit gesündigt, daß du nicht gethan hast, was du mir und deinen Töchtern zu Ehren und Erstattung billig solltest gethan haben; du hast ihnen keine Aussteuer und kein Theil vom Erbe gegeben. Darnach, hast du daran auch gesündigt, daß du uns nachgejagt und gedacht hast Uebels zu thun. Dieses sind große schwere Sünden, die du wahrlich erkennen und um deren Vergebung du bitten solltest.

203. Also ist nun die große Ehre Labans in die höchste Schande und Unehre verändert worden. Und Jakob erzählt nun seine Sünden insonderheit, eine nach der andern. Du bist auf mich erhitzt, sagt er, du hast mir nicht nur nachgejagt oder mich verfolgt, sondern bist von großem Zorn auf mich gar erhitzt und wüthend. Denn so viel heißt an andern Orten das hebräische Wort dalak; wie im 7. Psalm V. 14. steht: Sagittas suas ardentes effecit ignitas etc.: „Seine Pfeile hat er zugerichtet zu verbrennen“, er will Feuer schießen, eine Stadt anzünden. Also geißelt er die große Schuld Labans nicht allein damit, daß er sie ausführlich schildert, sondern auch damit, daß er sie mit starken Namen belegt. Du hast mich, sagt er, verfolgt mit deinen angebornen und andern Freundesfeinden, bist hitzig und grimmig gewesen, hättest gern mit Schwert und feurigen Pfeilen mich angegriffen, wie Einer, der von dem Teufel besessen ist.

204. Er schilt ihn wahrlich hart genug, als wollte er sagen: Wiewohl du an mir der Sünde keine gefunden hast, daß ich entweder nicht gethan hätte, was ich zu thun schuldig gewesen, oder aber das begangen hätte, was ich nicht sollte begangen haben: so solltest du mich wahrlich nicht so grimmig und so hitzig verfolgt haben, sondern solltest ja billig an unsere Verwandtschaft gedacht haben, daß ich dein Schwiegersohn und deine Töchter meine Weiber sind. Aber solcher Verwandtschaft und Liebe unangesehen, hast du mit so großer Unsinnigkeit wider uns getobt, daß du gar vergessen hast, daß du deiner Töchter Vater bist. Wenn aber nun ich dein Schwiegersohn nicht wäre, so bin ich doch deiner Schwester Sohn; denn du bist mein Vetter und nicht allein mein Schwiegervater. Derhalben könnten wir einander nicht näher verwandt sein, ich wäre denn dein Sohn, von deinem Weibe geboren. Und wo Jakob gestorben wäre, so wäre Laban Vater gewesen nach dem Gesetz. Also sind wir nun einander der Blutsfreundschaft sowohl, als auch der Schwägerschaft halben verwandt; aber du begehrest uns nicht allein mit keinem Geschenk, sondern unterstehst dich auch über das, uns noch zu berauben der Güter, die uns von Gott verliehen sind; welches ja wider alle Liebe und Freundschaft ist.

\quellseite{187}

205. Darum sagt er recht: Du zürnest mit mir nicht menschlicher Weise, sondern tobst ganz teuflischer Weise und steckst voll höllischen Feuers, bist gar erhitzt wider die, so unschuldig und dir ganz nahe verwandt sind. Wenn ich einen frommen und ehrlichen Schwiegersohn und dabei einen ungerathenen Sohn hätte, so wollte ich wahrlich dem ungerathenen Sohn alle mein Hab und Gut entbinden und wollte es meinem Schwiegersohn geben. Aber Laban wendet es um und verkehrt die Ordnung aller Werke der Liebe und der ganzen Natur: er beraubt seine ehrlichen Töchter und den sehr frommen Mann, seinen Schwiegersohn. Das ist nun das Erste gewesen.

206. Nun sagt aber Jakob weiter, und macht es noch immer größer, was Laban gethan hat: „Du hast alle meinen Hausrath“, das ist, all mein Hab und Gut, „bestätet“: dasselbe gebühret aber nun einem weisen oder aufrichtigen Mann auch nicht, vornehmlich mit den Gütern seines Sohnes oder Töchter also umzugehen. Denn es ist eine Schande, Alles insonderheit dermaßen zu begehren, anzutasten, und so gar genau zu untersuchen, als wären es die ärgsten Feinde: ja, es würde sich niemand gegen seinen Feind auf solche Art recht verhalten. Was hast du aber nun damit ausgerichtet, sagt Jakob, oder was hast du nun gefunden mit solchem fleißigen Nachsuchen? Allein, daß du damit angezeigt hast, daß du ohne alle Ursache also unsinnig und hitzig bist vom Zorn. Und wenn du ja nach meinem Urtheil nichts sagst, auch nicht gewiß schließen kannst, ob du oder ich unrecht habe: wohlan, so wollen wir die Sache an etliche Schiedsrichter gelangen lassen und mögen unsere Brüder zwischen uns beiden richten. Er bietet ihm das Recht an. Aber Laban wird durch sein eigen Urtheil überwunden und zu Schanden gemacht, daß er auch das Maul nicht darf aufthun.

207. Ja, also soll man mit den Gottlosen und Heuchlern umgehen: der Gerechte aber wird ledig und los erkannt, und der Gottlose wird an seine Statt dahin gegeben; denn das Unglück, damit er einen Andern bedroht hat, es ihm zuzufügen, kommt ihm auf seinen eigenen Kopf. Wie im 7. Psalm V. 17. steht: „Sein Unglück wird auf seinen Kopf kommen, und sein Frevel auf seine Scheitel fallen.“

208. Aber der Brüder halben ist eine Frage: Was Jakob damit meine, daß er sagt: „Zeuge das dar vor meinen und deinen Brüdern.“ Woher hat Jakob in dieser Pilgrimschaft Brüder gehabt, Sintemal er ja zwanzig Jahr lang allein gedienet hat? Und hernach (V. 46.) steht auch: „Er sprach zu seinen Brüdern: Leset Steine auf“ etc. Dieses ist meine Meinung. Denn wo wir keinen gewissen klaren Text haben, müssen wir rathen, und müssen, so gut wir können, Gedanken und Meinungen aus dem Text herauszuholen. Es läßt sich ansehen, daß ganz Mesopotamien fast eitel Weide gewesen sei, und die Leute haben sich darauf begeben, wie sie auch noch thun, daß sie des Viehes gehütet haben auf den Weiden: derhalben sind in allen Dörfern und Städten Hirten gewesen. Und ist wohl glaublich, daß Jakob unter seinem Gesinde viel Blutsfreunde gehabt hat aus dem Hause Labans, die ihm sowohl gedient haben wie seinem Schwiegervater. Und die Dienste sind getheilt gewesen, Etliche haben Jakob, Etliche aber Laban gedient. Die aber nun mit Jakob gezogen sind, werden ohne Zweifel fromm und gottselig gewesen sein. Wie gesagt worden vom Gesinde Abrahams, von Hagar und von Knecht Abrahams, die sich ihm angeschlossen haben um der Religion willen, die bei ihm rein war.

209. So versteht er nun damit die Brüder, die aus der Freundschaft Nahors und Tharahs gewesen sind, die sind mit Jakob wiederum in sein Vaterland gezogen, und er hat ihrer bedurft; daß sie ihm seine Heerde gehütet haben; denn seine beiden Weiber haben zu thun gehabt, daß sie die Kinder auferzogen und versorgt haben, die noch jung und schwach waren. Er aber allein hat so eine große Heerde nicht weiden können, er hat müssen Knechte und Mägde dazu haben. Dieselbigen nennt er nun an diesem Ort Brüder. Wohlan, sagt er, laß uns beiderseits etliche von unsern Verwandten und Blutsfreunden wählen, die diesen Hader entscheiden mögen. Dieses ist der Billigkeit sehr gemäß, daß ihm Jakob eine solche Bedingung anbietet. Und also wird nun Laban überwunden durch das göttliche Urtheil, durch sein eigen Gewissen, durch sein voriges Leben, und durch das gemeine Zeugniß aller Menschen und dazu seiner eigenen Brüder oder Verwandten; also daß er nun gar nichts finden kann, das er darauf hätte antworten mögen, da er dermaßen überwunden und zu Schanden gemacht ist.

\begin{quote}
\textbf{V. 38–41. Diese zwanzig Jahr bin ich bei dir gewesen, deine Schafe und Ziegen sind nicht unfruchtbar gewesen; die Widder deiner Heerde hab ich nie gegessen; was die Thiere zerrissen, brachte ich dir nicht, ich mußte es bezahlen; du fordertest es von meiner Hand, es wäre mir des Tages oder des Nachts gestohlen. Des Tages verschmachtete ich vor Hitze, und des Nachts vor Frost, und kam kein Schlaf in meine Augen. Also habe ich diese zwanzig Jahr in deinem Hause gedienet, vierzehn um deine Töchter, und sechs um deine Heerde; und hast mir meinen Lohn zehnmal verändert.}
\end{quote}

210. Jakob schilt seinen Schwiegervater nicht allein, sondern wirft ihm auch seine greuliche Undankbarkeit vor, indem er alle näheren Umstände hervorhebt, und wie er bei ihm so einen schweren Dienst gehabt habe. Denn er erzählet seine Arbeit und Wohlthaten und streicht dieselbigen heraus, auf daß dadurch die Bosheit und Undankbarkeit Labans gesteigert werde. Denn Mose hält uns aber ein vollkommenes Exempel und Abbild eines frommen und getreuen Knechtes vor, dergleichen man an andern Orten in keinerlei Historien findet. Zu unsern Zeiten aber, hilf Gott! wie ist es da so gar gemein, daß fast alle Hausväter jämmerlich klagen über die Untreue und Bosheit der Knechte.

Ich habe dir, sagt Jakob, mit so großem Fleiß gedienet, daß deine Schafe und Ziegen nie unfruchtbar gewesen sind.

211. Können aber das die Hirten zuwege bringen, daß das Vieh nicht unfruchtbar werde? Antwort: Es ist wahrlich sehr viel gelegen an einem getreuen und fleißigen Hirten. Denn man weiß wohl, wie boshaftig unsere Schäfhirten jetzt sind, die da können zuwege bringen entweder durch böse Künste, oder aber mit ihrem Unfleiße und Nachlässigkeit, daß die Schafe ihrer Herren, der Hausväter, sterben, unfruchtbar und mager werden: dagegen aber, daß die ihr eigen sind, viel Milch und Käse geben, dazu auch viel Wolle und viel Lämmer tragen müssen. Darum ist viel gelegen an der Frömmigkeit und Treue eines Hirten, daß er nicht versäume die Zeit, wenn die Schafe laufen, und ihrer sonst auch wohl warte; denn wo der Schafe wohl gewartet wird, geschieht es natürlich, daß sie auch pflegen fruchtbar zu sein.

212. Darum sagt Jakob: „Deine Schafe sind nicht unfruchtbar gewesen.“ Das habe ich mit meiner Treue und Fleiß zuwege gebracht, daß die Heerden viel Lämmer getragen haben; in zwanzig Jahren sind sie nie unfruchtbar gewesen. Denn in denselbigen Landen wie an etlichen Orten Italiens pflegen die Schafe im Jahre zweimal zu tragen. Dazu gehört fleißige Wartung, daß die Heerden nicht versäumt werden, wenn die Zeit ist, daß sie laufen sollen. Darum, sage ich, ist dieses ein Exempel eines getreuen Knechtes, welches gar kein gleiches hat. Wie sollte aber Laban darüber nicht reich geworden sein, wenn er auch selbst schon schläfrig und müßig gewesen ist, dieweil der treue Knecht nichts gethan hat, denn daß er seinem Herrn sein Gut mehrte und nur immer seinen Nutzen stiftete?

213. Zum Andern, sagt er: „Die Widder deiner Heerde hab ich nie gegessen.“ Denn das Recht haben die Hirten, daß sie bisweilen ein Schaf schlachten oder zum wenigsten die Milch oder Käse davon essen, und sich mit der Schafshaut oder Wolle bekleiden und bedecken. Jakob aber hat dessen nicht gebrauchen wollen, sondern hat es viel lieber alles zu seinem Schaden seinem Herrn lassen zu Nutze kommen; darum hat auch Rahel (V. 15.) nicht vergeblich gesagt, daß der Vater ihren Lohn verzehret habe.

\quellseite{188}

214. Zum Dritten, sagt Jakob: „Was die Thiere zerrissen, brachte ich dir nicht.“ Das hat sich nicht geziemet zu thun, wiewohl er nach allen Rechten das in seine Rechnung nicht sollte genommen haben, was von Wölfen oder andern Thieren zerrissen war. Und ich habe zwar, will er sagen, die Heerde so fleißig nicht bewahren können, daß nicht bisweilen der Wolf das eine oder das andere Schaf hätte nehmen sollen; darum sollte ich billig den Schaden nicht entgelten. Dieweil du mir aber hast vorgeworfen, daß solches durch mein Versäumnis geschehen sei, daß ich den Wolf in den Schafstall hätte fallen lassen, darum habe ich den Schaden müssen erstatten und bezahlen.

215. Also hat Laban keinen Schaden gelitten, weder daß die Heerden im Laufen sollten versäumt, noch auch daß ihm das Fleisch sollte aufgeessen worden sein, und dazu, daß ihm auch die Thiere nichts zerrissen haben. Ich habe allein die Schuld tragen müssen, sagt Jakob, ich mußte es gehen haben, ich habe bezahlen und erstatten müssen alles, was verloren war. Das ist wahrlich ein sehr schwerer Dienst und große Tyrannei gewesen, daß Einer auch das bezahlen muß, das mit dem höchsten Fleiß nicht kann oder mag bewahret werden. Kann man es doch oftmals im Hause nicht verhüten, daß nicht irgend ein Hund etwas nehme oder fresse, oder sonst Schaden thue. Das ist aber Wunder, woher Jakob das genommen habe, damit er es bezahlt und erstattet hat, was vom Hause verloren war. Wir können zwar dasselbe aus der Historie nicht abnehmen, errathen aber mag man es; denn es wird ihm entweder Isaak, sein Vater, Geld gesandt haben, oder Laban mag ihm irgend einen kleinen Lohn gegeben haben, welchen er fast ganz darauf gewendet hat, daß ja sein Herr, der Geizhalß, keinen Schaden nehmen möchte.

216. Zum Vierten, sagt Jakob: „Es wäre mir des Tages oder des Nachts gestohlen.“ Denn also lautet es nach dem Hebräischen. Wiewohl die Rabbinen sagen, der Buchstabe י sei an diesem Ort überflüssig. Ich halte es aber nicht allenthalben mit den Rabbinen. Dieses ist aber der Verstand dieser Worte: Alles, was die Leute des Tages und des Nachts gestohlen haben, das mußte durch meine Schuld und Versäumnis gestohlen sein; ich mußte die Zahl wiederum ergänzen und erfüllen, es wäre gleich von den Thieren zerrissen oder gestohlen, was verloren war. Dir ist daran nichts abgegangen, auch nicht einmal eine Hand voll Wolle, ich habe es alles ersetzen müssen, was nur für Schaden geschehen ist. Also ist nun Laban reich geworden gar durch eines Andern Sorge, große Mühe und Schweiß, ja, zum Schaden und Nachtheil dieses getreuen Hirten. Das ist aber verflucht Gut, das mit anderer Leute saurem Schweiß, und dazu mit Tyrannei, damit Andere beschweret werden, erworben wird.

217. So gebraucht nun Jakob zumal schöne Worte und predigt ihm das Gesetz, daß er ihn zur Erkenntniß seiner Sünde bringen möchte. Aber der Bauchler will allein gerecht sein, und läßt sich dünken, er könne auf keinerlei Weise sündigen. Darum ist es unmöglich, daß man solche Leute zur Erkenntniß der Sünde und zu rechter Buße bringen sollte.

218. Zum Fünften, sagt er: „Des Tages verschmachtete ich vor Hitze und des Nachts vor Frost“ etc. 2c. Das ist ihm die ersten vierzehn Jahre noch erträglicher gewesen, als die letzten sechs Jahre, da er sich selbst auch, und seinen Weibern und Kindern hat dienen müssen, sein Gut zu mehren und sein Haus zu versorgen. Derhalben ist es ihm sehr schwer gewesen, daß er zugleich Hitze und Frost hat leiden sollen, und hat dabei keinen Schlaf noch Ruhe haben können. Denn es bedarf ja der Leib, daß er seine Ruhe haben und schlafen möge, auf daß er nicht gar matt und entkräftet werde, wenn er allzu ermüdet wird. Warum hat er aber so viel gelitten? Antwort: Darum, daß die Diebe und Wölfe der Hitze oder des Frostes nicht achten und weder nach Tag noch Nacht fragen. Darum, sagt er, habe ich müssen bei der Heerde wachen Nacht und Tag, ich durfte mich der Hitze oder des allzugroßen Frostes halben nicht entschuldigen: ich mußte fort; und das geschah alles zu deinem Nutzen und Gewinn, da ich harte Arbeit thun und schwitzen mußte, auf daß nur deine Heerde wachsen und zunehmen möchte, wenn du schliefest.

219. Wo willst du aber heutiges Tages ein Exempel finden, das dieser wunderbaren Treue, so an Jakob gewesen, gleich wäre? Denn es klagen beide, die vom Adel und auch die Bauern, die mit Vieh umgehen, alle über die große Untreue der Schäfhirten und sagen also, daß die Schafe der Herren gründig und unfruchtbar werden und sonst andern Schaden leiden, dazu von den Thieren zerrissen werden; aber der Schäfhirten Schafe nicht also, darum daß sie jener Schafe weniger Acht haben und zu ihrer Herren Schaden ihren eigenen Nutzen und Frommen stiften. Und sie pflegen gemeiniglich, es sei nun eine Fabel oder eine Historie, zu erzählen, daß Einer vom Adel einen Schäfhirten hat gefänglich einziehen lassen um seiner Bosheit und Diebstahls willen; da man aber denselbigen jetzt hat sollen an den Galgen hängen, soll ihn der Edelmann gefragt haben: Ob er nicht irgend einen andern Schäfhirten wüßte, der fromm und getreu wäre, den er möchte über seine Schafe setzen? Da soll der Dieb gesagt haben: Er wisse gar keinen, und sie wären alle eben so fromm, wie er, ja, etliche wären noch wohl etwas ärger. Da, sagen sie, solle der Edelmann durch solche Antwort sein bewogen worden, daß er diesem Schäfhirten das Leben geschenkt und ihn noch behalten hat, weil er ja keinen besseren hat finden können.

220. Darum erklären die Sitten, so jetziger Zeit die Knechte haben, die Tugend und Treue dieses frommen Hirten Jakob sehr fein. Denn es ist gar ein sonderlich Exempel, und ist wohl zu sehen und abzunehmen, wie groß heutiges Tages die Verschwendung sei mit der Haushaltung, da die Leute so überaus untreu sind. Man muß aber mit solchem Schaden Geduld tragen und muß stille schweigen, bis uns Gott aus diesem Leben in ein besseres fordern und bringen, und an den bösen untreuen Leuten die Sünde strafen wird. Denn solche Treue, wie Jakob einem ungerechten und untreuen Herrn bewiesen hat, wird niemand einem Andern erzeigen können, es sei denn, daß er recht an Gott glaube und ganz und gar mit dem Glauben am Wort hange.

221. Jakob sagt noch mehr: „Also habe ich diese zwanzig Jahr in deinem Hause gedienet, vierzehn um deine Töchter, und sechs um deine Heerde; und hast mir meinen Lohn zehnmal verändert.“ Laban hat sich daran nicht genügen lassen, daß er in den vierzehn Jahren durch die Arbeit und Treue Jakobs gar reich geworden; er hat ihn auch noch über das seines gebührlichen Lohnes berauben wollen. Darum fühlt ihn nun Jakob so sehr, als er immer kann, daß er ihn doch dahin bringen möchte, daß er seine Sünde erkennen möchte. Aber Laban hat ein solch Herz, das härter ist denn ein Diamant oder ein Amboß.

222. Und daran mag man sehen, wie große Geduld und auch großen Glauben die Patriarchen gehabt haben. Es haben zwar die Märtyrer viel und greuliche Marter und Pein erlitten: aber es ist alles nicht zu vergleichen mit der sehr schweren Arbeit und großen Gefahr der Patriarchen. Denn das ist eine große Geduld, ganze zwanzig Jahre solche große Tyrannei und Gewalt leiden. Wie es aber mit der Veränderung des Lohnes sei zugegangen, ist droben gesagt worden.

\begin{quote}
\textbf{V. 42. Wo nicht der Gott meines Vaters, der Gott Abrahams, und die Furcht Isaaks auf meiner Seite gewesen wäre; du hättest mich leer lassen ziehen. Aber Gott hat mein Elend und meine Angelegenheiten, und hat dich gestern gestraft.}
\end{quote}

223. „Gott hat mein Elend angesehen“, sagt Jakob, „und hat mich gestraft.“ Im Hebräischen steht es also, daß kein Accusativ dabei ist, und lautet auf Deutsch also: Gott hat es gesehen und darum gescholten. Denn ein solcher greulicher Tyrann ist Laban gewesen, daß er Jakob nicht allein würde haben leer und nackt ziehen lassen, sondern er würde ihn auch mit Gewalt wiederum in den schweren Dienst gebracht haben, wo ihn Gott nicht unterdeß bedroht und gehindert hätte; wie zuvor gesagt ist. Dieses ist übrigens ein sehr trefflicher Text, desgleichen wir noch nicht gehabt haben, um des Namens willen „die Furcht Isaaks“. Das Wort pachad, bei den Hebräern, heißt, sich fürchten oder entsetzen. Wie im 14. Psalm V. 5. steht: „Daselbst fürchten sie sich“ etc., die Heuchler fürchten sich, da keine Furcht ist; oder wie am andern Ort, Spr. 28, V. 1. steht: „Der Gottlose fleucht, und niemand jagt ihn.“ Also wird es gemeiniglich ausgelegt. Aber man muß es vom Gottesdienst verstehen.

224. Was mag aber nun der Heilige Geist damit meinen, daß er die „Furcht Isaaks“ nennt? Daran ist kein Zweifel, er will damit Gott verstanden haben; denn er setzt noch dazu: „Der Gott meines Vaters, der Gott Abrahams, und die Furcht Isaaks.“ Das Wort „Furcht“ wird mit dem Wort „Gott“ verbunden, derhalben muß es so viel heißen, als Gott. Darum bedeutet die Furcht an diesem Ort einfach GOtt. Wo aber dieses Wort etwas Anderes sollte heißen denn GOtt, so würde es Jakob nicht mit dem Wort „GOtt“ verbunden haben. Denn nun darf die Ehre, oder den Gottesdienst keinem Andern zuwenden. Es sind aber unzählige Fragen und Auslegungen, warum er es die Furcht oder Ehre Isaaks heiße, und warum er diesen Namen nicht auch dem Abraham, sondern nur dem Isaak gegeben habe. Ich will aber darin dem einfältigen und gemeinen Verstande folgen, daß er es also genannt habe um des trefflichen und sonderlichen Gottesdienstes willen seines Vaters Isaak, als der mit sonderlicher Gottesfeligkeit GOtt gefürchtet und geehrt hat. Dieses ist der gemeine Verstand und ich lasse mir denselbigen wohlgefallen.

\quellseite{189}

225. Ueber das aber muß man auch diese Worte merken, welche Mose also gesetzt hat, daß sie einen sonderlichen Nachdruck haben, wie ich an andern Orten oftmals auch davon Erinnerung gethan habe. Wie im 2. Capitel dieses ersten Buchs Mose V. 21. steht: „GOtt, der Herr, baute ein Weib aus der Ribbe, die er von dem Menschen nahm.“ Da hätte er wohl ein ander und vielleicht ein besser Wort gebrauchen können, wie es scheint. Er hat aber mit diesem wollern Ausdruck anzeigen wollen, daß ein Weib ein göttlich Gebau ist, nicht allein um des Geheimnisses willen der Kirche, sondern auch was die Haushaltung belanget. Denn die ganze Schrift braucht diese Art zu reden von den Weibern, wie 2 Mos. 1, 21. steht von dem Segen, damit GOtt die Wehmütter gesegnet und belohnt hat: „GOtt baute ihnen Häuser“ etc. Und Rahel hat droben, Cap. 30, V. 3., zu Jakob auch also gesagt: „Lege dich zu ihr, daß ich durch sie erbauet werde.“ Das ist so viel gesagt: daß ich Kinder und Erben kriegen möge im Hause oder in unserem Geschlechte. Denn dazu ist vornehmlich das Weib geschaffen, daß durch sie das Haus soll erbaut und daß Kinder und Erben geboren werden.

226. Dergleichen sagt Jakob an diesem Ort, daß Laban sey gestraft worden, und dasselbe ist auch sehr nachdrücklich geredet; wie kurz zuvor, V. 16., Rahel auch gesagt hat, wie GOtt ihrem Vater seinen Reichthum entwandt habe etc. Denn der Teufel mit den bösen Engeln und auch mit bösen Menschen ist nur darauf beflissen, daß er alle guten Werke verhindere und zerstöre. Wo nun noch irgend was Gutes geschieht, dasselbe pflegt also zu geschehen, daß GOtt den Teufel und böse Menschen straft, wider sie streitet und ihnen das Gute gleichsam aus dem Rachen reißen muß. Diese und dergleichen Art zu reden, die so nachdrücklich und mit besonderem Fleiß gesetzt werden, soll man in der Schrift wohl merken.

227. Derhalben halte ich es dafür, daß an diesem Ort der Verstand nach dem Buchstaben der beste sey, nämlich, daß GOtt die „Furcht Isaaks“ genannt werde um des trefflichen Gottesdienstes willen, so Isaak getrieben hat. Die Andern legen es also aus, daß sich Isaak gefürchtet habe, da er hat sollen geopfert werden. Aber der Verstand reimt sich gar nicht und ist gottlos. Denn Isaak hat sich nicht gefürchtet, ist auch nicht erschrocken, sondern ist gutwillig seinem Vater und GOtt gehorsam gewesen. Denn man muß den Heiligen nicht zuschreiben, daß sie sich sollten fürchten, sondern es ist ihnen ein freiwilliger Gehorsam und Tapferkeit. Denn Furcht heißt die Sünde, die da herrscht, welche Sünde Isaak nicht gehabt hat, sondern der Geist ist in ihm gehorsam gewesen; und wiewohl sich das Fleisch dawider gelegt, so hat doch der Geist den Sieg behalten und geherrscht, der ihm das Fleisch unterworfen hat.

228. Aber außerdem glaube ich, daß auch Christus hier angedeutet wird. Denn die Väter haben diesen Segen ausdrücklich vom leiblichen unterschieden: „In Isaak soll dir der Same genannt sein“, spricht St. Paulus Röm. 9, 7., welcher auch auf das erste Wort sonderlichen Nachdruck legt. Denn der Segen Ismaels mit den zwölf Fürsten, die von ihm kommen sollten, war leiblich. Aber „in Isaak soll dir der Same genannt sein“, welches sonst kein anderer Same ist denn allein Christus, der ist der rechte Same. Denn die Väter haben über den leiblichen Segen, den Ismael und Esau gehabt haben, vornehmlich auf den geistlichen Segen gesehen, welcher Christus selbst ist. Darum verstehe ich hier die heilige Dreifaltigkeit, und desgleichen an andern Orten auch, wo ich nur aus den Worten des Alten Testamentes dieses Geheimniß herausnehmen kann.

229. Also heißt hier die „Furcht“ der Sohn Gottes, der da Mensch geworden ist, auf welchen die Patriarchen und Propheten in den Verheißungen gesehen und über den sie das Volk belehrt und berichtet haben. Wie er anfänglich ist verheißen worden in diesem ersten Buch Mose Cap. 3, V. 15.: „Der Same des Weibes soll der Schlange den Kopf zertreten.“ Item Cap. 12, 3. und Cap. 26, 4.: „Durch deinen Samen sollen alle Völker gesegnet werden.“ Diesen Samen haben die Väter gefürchtet und ihm gedient, welcher gleichsam der rechte Edelstein und der Kern in den leiblichen Verheißungen ist, welche wie Schalen oder Spreu waren, so ihnen gegeben gewesen sind, ihren Glauben damit zu stärken, daß sie Erben um Christi willen haben würden.

230. Darum will nun der Heilige Geist anzeigen, daß Isaak mit großem Verlangen auf diesen Samen gewartet und ihm gedient habe. Und Jakob rühmt sich hier des künftigen Herrn Christi, dem sein Vater gedient hat. So sollen wir nun fleißig daran denken, daß solche vornehme Stellen und Sprüche auf Christum müssen gezogen werden: wie wir deren droben etliche in den Verheißungen gehabt haben. Und hernach im 49. Capitel V. 10. wird vom künftigen Schilo gesagt werden, der ein Sohn ist vom Leibe seiner Mutter geboren. Und lässet uns ja der Auslegung der jüdischen Rabbinen nicht folgen, die Christum den Herrn greulich verdunkeln. Denn die Väter haben auf diesen Edelstein, so in die leiblichen Segen gefaßt gewesen, fleißig Acht gehabt und denselbigen in Ehren gehalten; sonst würde Jakob hier nicht gesagt haben, die „Furcht Isaaks“. Denn das ist der Sohn Gottes, Christus Jesus, der Mensch geworden ist, in der andern Person. Und er sagt nicht darum also, daß ein anderer GOtt sei denn der GOtt Abrahams, sondern das Geheimniß der Dreifaltigkeit damit anzuzeigen.

231. Auf die Weise werden uns diese nachdrücklichen Worte sehr süß und angenehm, welche Mose hin und wieder in seinen Schriften mit untermenget wie Edelsteine, sofern sie nur auf Christum gezogen werden; denn die andern Auslegungen sind kalt und faul, wie die ist vom Opfer Isaaks.

232. Und dieses ist auch der rechte und schöne Schmuck dieser Historien, darin der Heilige Geist solche geringe und verächtliche Werke der heiligen Patriarchen beschreibt, daß sich die Heuchler und Papisten dünken lassen, daß sie gar verachtet und ungöttlich sind: sie begreifen aber in sich die höchste Weisheit und rechten Glauben. Darum alles, was die Heiligen thun, wie gering und unschätzig es auch scheint, ist es doch groß und herrlich. Denn sie thun Alles im Glauben und im Wort, das ist, in Unschuld und Heiligkeit. Denn also spricht Christus Luc. 8, 50.: „Glaube nur“, so wirst du selig werden, wirst rein und heilig werden, und was du thust, wird alles rein, heilig und nützlich sein. Wie an einem Schafe nichts Unnützes ist; denn es hat auch der Mist und die Beine ihren Nutzen, und nicht allein die Milch oder die Wolle: also, wo etwa die Heiligen aus Ungeduld noch murren, das ist der Mist; es thut aber den Heiligen keinen Schaden, sondern muß ihnen Alles zum Besten dienen; denn sie leben, thun und leiden im Glauben an das Wort. „Dein Wort“, spricht Christus Joh. 17, 17., „ist die Wahrheit.“ Wer nun das Wort ergreift und demselben glaubt, der ist auch wahrhaftig heilig und gerecht. Dagegen sind aber alle Widerfacher des Worts, wie heilig sie auch dem äußerlichen Scheine nach sind, doch verworffen und verdammt.

233. Dieses soll man fleißig merken uns und der ganzen Kirche zum Trost und Lehre, auf daß wir wissen, daß unser Leben GOtt angenehm und gefällig ist auch in den geringen häuslichen Werken. Denn davon weiß die Vernunft und die Welt nicht, welche sich nur auf wunderliche Gesichte, Erleuchtungen und seltsame Speculation legt, dazu auch mit harten und scheußlichen Werken, als Fasten und Wachen etc., umgeht. Daher sind die Mönche und bei den Heiden die Philosophen gekommen, die es in solchen seltsamen Werken Andern haben wollen vorzuthun, und haben also damit machen wollen, daß sie das gemeine Volk sollte hoch ansehen und in großen Ehren halten; aber die Weisheit dieser Welt ist vor GOtt eine Thorheit.

234. So lebt nun Jakob ohne allen Schein sonderlicher Werke in Einfalt und in der Knechtschaft, und scheint, als sei er vor GOtt verachtet, wie er in der Welt gar arm und elend ist. Denn die Welt sieht nicht, was der Glaube sei und was er vermöge: so gar ein verborgen und unbekannt Ding ist er vor der Vernunft, nach dem Spruch Christi Joh. 14, 17: „Die Welt kann den Heiligen Geist nicht empfahen“, steht aber hört ihn nicht, verstehet auch von ihm mit keinerlei Sinne etwas. Wir sollen uns aber das angelegen sein lassen, daß die Gottseeligen gelehrt und zum Glauben an das Wort erweckt werden mögen, auf daß sie sich immer finden lassen in der Liebe und Uebung des Wortes. Denn alsdann sind alle ihre Werke, damit sie umgehen beide im Welt- und Hausregiment, sehr gut: du gehest, du schläfest oder wachest, so gefällt es Gott alles wohl; denn das Wort, der Glaube und der Geist, der in dir wohnt, sind gut.

\quellseite{190}

235. Endlich gehört das auch zu unserem Trost, daß Jakob sagt, daß der GOtt seines Vaters auf seiner Seite gewesen sei etc., und habe sein Elend und Arbeit angesehen. Was ist aber seine Arbeit Anderes, denn geringe, schwache und dazu unflätige Werke, die bisher an ihm als einem armen Schäferlein beschrieben worden sind? Von diesen Werken hat der Heilige Geist haben wollen, daß sie der Kirche und den Gläubigen zum Exempel und Stärkung sollen vorgehalten werden, auf daß sie gewiß wissen möchten, daß GOtt alle Werke der Heiligen und Gottseeligen ansehe, ja, daß er auch alle ihre Haare gezählt habe, wie etlichemal ist gesagt worden. So große Bekümmerniß und Sorge hat er für die, so das Wort hören und glauben.

236. Dieser Trost ist aber vornehmlich der Kirche und Gemeinde GOttes sehr nöthig; denn der Teufel mit der ganzen Welt und unserem eigenen Fleisch plagt die jämmerlich, die das Wort haben. Darum darf man auf keinen Trost oder Schutz hoffen von der Welt und den Fürsten dieser Welt, und wir dürfen bei uns selbst auch keinen fleischlichen Rath suchen; denn das ist alles vergeblich. Die aber gleichwohl solchen Rath suchen oder darauf hoffen, die müssen erfahren, daß im 116. Psalm V. 11. steht: „Alle Menschen sind Lügner“; item im 146. Psalm V. 3.: „Verlasset euch nicht auf Fürsten, sie sind Menschen, die können ja nicht helfen.“ Aber die ganze Welt geht dahin, als wäre sie blind, und verdirbt in ihrem eitelen Rath und Vornehmen; denn sie will nicht leiden, daß sie der Teufel, die Gottlosen und ihr eigen Fleisch verfolgen soll.

237. Darum muß GOtt die Seinen ansehen und sie versorgen, die sonst gar verlassen sind und weder Hülfe noch Rath wissen, sondern nur auf sein Wort sehen; denn sie tragen das Kreuz und haben Jammer und Trübsal; wie Jakob nicht allein vom Teufel geplagt wird, sondern auch von seinem Schwiegervater und Blutsverwandten selbst. Darum lasset uns dieses lernen und merken, daß GOtt alle unsere Werke und Gedanken zählt, wie gering und verachtet sie auch sein mögen, wenn wir nur an ihn glauben. Denn wo wir auf sein Wort Acht haben, so hat er wiederum auf uns Acht, und sorget für uns in aller Noth und Widerwärtigkeit, die uns die Welt und der Satan in den Weg legen.

\section[Siebenter Theil]{Siebenter Theil. Labans Antwort auf Jakobs Bestrafung, und wie er mit Jakob einen Bund macht}

\subsection{I}

\begin{quote}
\textbf{V. 43. Laban antwortete, und sprach zu Jakob: Die Töchter sind meine Töchter, und die Kinder sind meine Kinder, und die Heerden sind meine Heerden, und alles, was du siehest, ist mein. Was kann ich meinen Töchtern heute oder ihren Kindern thun, die sie geboren haben?}
\end{quote}

238. Bisher hat uns Mose ein Exempel und Bild eines heiligen Mannes und eines frommen und getreuen Knechtes vorgehalten, nämlich Jakobs, was er gethan und wie viel er gelitten habe, wie er sich nach der Liebe beflissen habe, seinen Widersacher zu gewinnen, und zur Erkenntniß seiner Sünde oder zur Buße zu bringen, und also gethan hat, was ein heiliger und frommer Mann billig thun soll; was er aber damit ausgerichtet habe, wollen wir nun weiter hören.

239. Denn Mose beschreibt uns ein Exempel, das dem gar entgegen ist, nämlich, eines sehr bösen Heuchlers. Wiewohl er ihn deutlich genug abgemalet hat, wie er sich gegen seinen Schwiegersohn, Töchter und Kindeskinder gehalten hat; so lasset uns doch nun sehen, nachdem er ist ermahnt und gestraft worden von GOtt und auch von Jakob, wie viel er sich daraus in seinem Leben gebessert habe. Er sagt zum Ersten: „Die Töchter sind meine Töchter, und die Kinder sind meine Kinder, was kann ich ihnen nun thun?“ Ei du heiliger Mann, St. Laban! Da hört man doch kein Wort von der Buße oder Bekenntniß seiner Sünde. Er sagt nicht: Mein lieber Jakob, ich erkenne und bekenne meine Sünde; denn ich bin zu grausam und zu wenig väterlich mit dir umgegangen, und habe mich auch gegen dich und meine Töchter nicht gehalten, wie mir gebühret hätte. Solche Worte wirst du nimmermehr einem Heuchler können abdringen, oder gar selten. Denn ein Heuchler ist ein solch Ungeheuer, das schlechterdings nicht sündigen kann, und ist auch nicht zu bekehren. Ein Heuchler läßt sich nicht dünken, daß er sündigen oder irren könne, wenn er schon seiner Sünde und Irrthums halben dermaßen ermahnt wird, daß er es mit Händen greifen kann.

240. Denn ein solch Bild wird uns in etlichen andern auch vorgehalten, als, in Pharao und Saul; oder aber, wenn sie ja bisweilen mit Worten ihre Sünde bekennen, so thun sie doch dasselbe gar selten, oder thun es heuchlerischer und verstellter Weise. Es muß der allein, der da heilig und gerecht ist, ein Sünder sein. Ja, das noch mehr ist, die Ungerechten und Gottlosen wollen den Namen haben, daß sie nicht sündigen oder Uebels thun, vornehmlich wenn sie Macht und Gewalt haben. Als, Heinz Mordbrenner, der Bischof von Mainz, Doctor Eckel und Grickel thun nicht Unrecht, die sind vor den Menschen gerecht und heilig.

241. Dieses ist nun eine große Anfechtung und Unbilligkeit, die das Fleisch sehr ärgert und plagt. Aber diese Exempel werden uns zum Trost und Stärkung geschrieben, auf daß wir wissen, daß es nicht anders sein könne, wir müssen bei solchen Heuchlern leben, und wir mit so viel größerem Muth ertragen und dulden mögen, wenn wir solche verkehrte Urtheile der fleischlichen Menschen hören müssen. Jakob muß allein ein Sünder sein; Laban aber ist fromm und gerecht. Also müssen jetzt wir Wittenberger auch allein Sünder sein, wiewohl wir von unserer Lehre Rechenschaft geben jedermann, der das von uns fordert. Wir beweisen und vertheidigen unsere Lehre, wir widerlegen den Widersachern ihre Lehre und überwinden sie nach ihrem eigenen Urtheil und Gewissen. Wir richten aber damit gleichwohl nichts aus. Sie sind fromm und gerecht; wir müssen gottlos und ungerecht sein. Wer nun das nicht leiden und ertragen kann, der kann auch kein Christ sein. Sie müssen Recht haben, Christus und Christen haben unrecht.

242. Nun wollen wir die Worte Labans besehen, welche mir nicht fälschlich, sondern aufrichtig auslegen wollen, nicht nach der Meinung seines Herzens, nach welcher Alles mit großem Stolz und höhnisch geredet zu sein scheint. Denn daß er sagt: „Die Töchter sind meine Töchter“, das ist wohl wahr; aber er hatte sie jetzt Jakob gegeben, daß sie seine Weiber waren. Desgleichen war auch die Heerde sein verdienter Lohn, die er mit großer Arbeit und Schweiß erworben hatte. Darum lügt Laban und ich möchte seine Worte besser also deuten; ich will aber solchem harten Verstande nicht folgen. Wir wollen zwar die Worte also verstehen, als habe er sie aus väterlichem Herzen geredet. Als wollte er sagen: Mein lieber Sohn, du hast ja alles von mir, was du hast; ich bin der Vater, warum sollte ich denn den Meinen Schaden thun? Ich habe dir meine Töchter gegeben, dazu Schafe und Ochsen, du hast Alles von mir bekommen; der Herr segne dich, daß du wachset und groß werdest, und deines Gutes gottseeliglich und glückselig gebrauchen mögest, daß dir es wohl gehe. Denn also pflegt ein aufrichtiger frommer Vater zu seinen Töchtern und Schwiegersohn und Kindeskindern zu reden. Eben solche Worte gebraucht dieser Heuchler auch. Es ist aber alles falsch und eitel erdichtet Ding. Und das wollen wir mit zwei Argumenten beweisen.

243. Zum Ersten, erkennt er noch seine Sünde und Raub nicht, so ihm Jakob vorgeworfen hat. Denn der hat ihn des Diebstahles und sonst vieler schwerer Sünden beschuldigt und daß er sich gegen ihn und seine Töchter greulich und unfreundlich gehalten hätte; aber er erkennt oder bekennt der Dinge keines, sondern will gerecht und fromm sein. Derhalben, obgleich solche Heuchler die Worte frommer aufrichtiger Leute gebrauchen, wie Laban die Worte eines frommen getreuen Vaters auch sein nachreden kann, so reden sie doch nicht von Herzen, was sie reden, sie sind nicht aufrichtig.

244. Zum Andern, hätte auch dem Vater gebühret, daß er seinen Töchtern und Kindeskindern irgend ein Geschenk sollte gegeben haben. Meine liebe Rahel, sollte er gesagt haben, siehe,
\quellseite{191}
da hast du zehen Gulden etc. zum Ehrentheil. Du Ruben nimm hin, daß will ich dir schenken, daß du dabei deines Großvaters gedenken mögest. Aber er bekennt seine Sünde nicht, auch gibt er ihnen gar kein Geschenk zum Zeugniß seiner väterlichen Liebe und Gewogenheit. Es ist nur alles darauf gerichtet, daß der gute fromme Mann soll hintergangen und betrogen werden, auf daß die Weiber und Brüder ihr Urtheil von Jakob ändern, aller Unbilligkeit vergessen und die Freundlichkeit des Schwiegervaters rühmen sollen; aber die Früchte weisen es fein aus, was es für ein Baum ist.

245. Denn Gott will vor allen Dingen haben, daß er bekennen soll, was er ihm genommen hatte, und daß er es auch wiedergegeben sollte; darnach, daß er auch ein Zeichen rechter Gewogenheit und ungefärbter Liebe von sich geben sollte. Aber der geschieht keines; der fromme Heuchler beharrt in seiner Gerechtigkeit und gibt doch kein Zeichen der Gerechtigkeit von sich, weder mit Ja, noch mit Nein. Also ist es gar vergeblich und umsonst, daß ihn Jakob gestraft hat. Das ist nun das Erste, daß wir es aufrichtig auslegen und den harten und feindseligen Verstand fahren lassen. Nun fährt er aber fort, und wird auch geistlich und heilig. Er läßt sich daran nicht genügen, daß er seine Gerechtigkeit vertheidigt wider Jakob, der ungerecht sein muß, als hätte Laban nicht wider ihn gesündigt: jetzt wird man ihn nun in den Himmel heben müssen unter die Chöre der heiligen Engel.

\begin{quote}
\textbf{V. 44. So komm nun, und laß uns einen Bund machen, ich und du, der ein Zeugniß sei zwischen mir und dir.}
\end{quote}

246. Der Heuchler wiederholt und führt oftmals dieses Stück ein, daß er sagt: Es sei ein Zeugniß zwischen mir und dir, wie die ganze Erzählung ausweisen wird. Denn also pflegen die Heuchler zu thun: die weil sie gar keine Früchte tragen, haben sie doch viel Blätter, wie der Feigenbaum, den Christus ersucht. Er wiederholt das zum öftern bis zum größten Ueberfluß, daß er sagt: „Gott sei Zeuge zwischen mir und dir.“ Item: Der Berg, der Steinhaufe und das Malzeichen sei Zeuge zwischen uns. Aber es ist kein rechter Grund oder Wahrheit dahinter. An diese Kunst der Heuchelei müssen sich die gewöhnlichen, die in der Welt leben wollen.

247. Und jetzt hebt er nun an, den frommen Mann verhaßt zu machen und ihn mit Argwohn zu beschweren. Denn er stellt sich, als ob er Gott fürchte, und daß er Jakob auch als einen, der nicht sehr gottesfürchtig ist, mit seiner Vermahnung zur Furcht Gottes reizen wolle. Er will einen Bund machen vor Gott, will einen Stein aufrichten, und macht mit ihm einen sehr genauen Vertrag hinsichtlich der Grenzen, daß Jakob über dieselbigen nicht herfahre und ihm irgend einen Schaden zufüge. Wozu dient es aber, daß er eines solchen Mannes halben so gar besorgt ist, deß Treue und Fleiß er ganze zwanzig Jahre sehr wohl erkannt und erfahren hatte? Was thut es noth, daß er es nöthigen, daß er von seinem getreuen Schwiegersohn will, daß er sich gegen ihn mit einem Eid verpflichten soll und daß er ein so groß Gespänge treibt mit dem Bunde? Er hat ihm zuvor seine Töchter ehelich vertraut und hatte ihm dazu auch noch sein Hab und Gut vertraut zu versorgen: warum fürchtet er sich denn nun vor ihm, da er im fremden Lande wohnt?

248. Es ist Anderes nichts, denn nur ein äußerlicher Schein und unnütz Gepränge, damit er versucht, den frommen heiligen Mann zu verunglimpfen und in Verdacht zu bringen, als ob ihm Gefahr von Jakob drohe. Er maßt sich große Unschuld und Gerechtigkeit an, und stellt sich, als müsse er sich vor den bösen Schalk Jakob sehr fürchten, als ob er ihm gedächte Gewalt und Unrecht zu thun. Wenn er ein frommer getreuer Mann wäre, so würde er vielmehr also gesagt haben: Siehe hin, mein lieber Jakob, der Herr segne dich und die Deinen: wo ich dir in deinen Sachen dienen kann, das will ich herzlich gern thun, und wo irgend Gefahr oder Beschwerungen vorfallen werden, so zeige mir es an, so sollst du erfahren, daß ich dich väterlich und herzlich lieb habe. Ich will dich, und meine lieben Töchter und Kindeskinder nimmermehr verlassen. Aber er verspricht ihnen nichts; ja er beschweret noch über das Jakob damit, daß er ihn verhaßt macht, da er sich stellt, als ob er sich vor ihm fürchten müsse. Denn er geht damit um, daß er den Schein und Namen haben will, als sei er heilig und gottesfürchtig; Jakob aber soll von jedermann für verdächtig gehalten werden, als wollte er Laban gern Schaden zufügen. Also, dieweil er ihn mit Recht nicht überwinden und strafen kann um irgend einer Sünde willen, versucht er gleichwohl, ihm fälschlich eine Schelle anzuhängen, daß er mit solchem Argwohn hinziehen soll, als der ihm wolle gern Unrecht thun, und ein solcher Mann sei, der nach Gott nichts frage.

249. Es fällt mir aber das auch noch ein, daß ich denke, Laban werde etlichermaßen bestürzt gewesen sein, wie die Gottlosen pflegen, die sich schuldig wissen, und auch vor einem rauschenden Blatte zu erschrecken pflegen. Wie Salomo Spr. 28, 1. sagt: „Der Gottlose fleucht, und niemand jagt ihn.“ Und die Poeten schreiben von Orestes, wie er sei unsinnig geworden. Und Juvenal hat die Menschen sehr wohl beschrieben und abgemalt, welche von ihrem eigenen Gewissen gemartert werden, da er sagt: Cur tamen hos tu evasisse putes, quos diri conscia facti mens habet attonitos? etc., das ist: Du sollst nicht meinen, daß die ungestraft davon gekommen seien, die ihr eigen Gewissen beschuldigt und erschrocken macht. Hi sunt, sagt er weiter, qui trepidant et ad omnia fulgura pallent: Das sind die Gesellen, die also erschrocken sind, und sich gar entfärben und bleich werden auf einen jeden Blitz. Denn das ist die Strafe, so auf die Sünde folgt.

250. Derhalben hat das auch wohl geschehen können, daß (wiewohl Laban wußte, daß ihm seine Gefahr von seinem Schwiegersohne drohe, den er ja als getreu und fromm durch Erfahrung kannte) er sich doch gleichwohl vor den Andern gefürchtet hat, als nämlich vor den Eltern, oder Verwandten und Freunden Jakobs. Aber das Römische ist dieses gewesen, daß er sich mit solchem Schein und Heuchelei hat schmücken und Jakob beschweren und verdächtig machen wollen.

251. Mit solchen Leuten müssen nun die Gottseeligen in dieser Welt umgehen, die sie beide leiblich und geistlich betrügen, und dasselbe noch unter dem Schein der Heiligkeit, Gerechtigkeit, Gottseeligkeit und Geistlichkeit; und dann erst pflegen sie zu triumphiren und fröhlich zu sein, wenn sie uns übertäuben können mit dem falschen Wahne, als ob sie die rechte Religion meinten und Gott recht dienen, wir aber wider dieselbige sündigen. Darum ist nun Laban allein gerecht und heilig, in der ersten und andern Tafel. Dieses ist das rechte Conterfei und Bild der Welt.

\subsection{II}

\begin{quote}
\textbf{V. 45. Da nahm Jakob einen Stein, und richtete ihn auf zu einem Mal.}
\end{quote}

252. Jakob ist dazu willig, einen Bund mit ihm zu machen, nur um der Ursache willen, daß er also füglich und bald von dem sehr schmeichelnden Heuchler kommen möge. Denn es ist allen Gottseeligen sehr beschwerlich, wenn sie mit Heuchlern sollen umgehen; und kann zwar auch ein Jeder leichter mit einem offenbaren Feinde umgehen, denn mit einem solchen Heuchler, der sich äußerlich stellt, als ob er Freund sei und gern Wohlfahrt erzeigen wollte, und der doch innerlich im Herzen Feindschaft und Ungunst trägt. Es ist nichts Schwereres oder Verderblicheres, denn wo man mit einem solchen Menschen umgehen soll, bei ihm wohnen, mit ihm essen, der zugleich auf allerlei Gelegenheit steht, wie er nur Einen lästern und betrügen möge, und doch dessen nicht öffentlich überführt werden kann, daß er mit Gewalt oder Betrug umgeht, sondern dessen Heuchelei und heimliche Tücke, so er hinter sich hat, man dulden und ertragen muß; ja, ein solcher Heuchler schwört noch theuer und hoch, daß er der beste Freund, und der wahren Gottseligkeit ganz geneigt und gewogen sei; oder wenn du ihn ja Lügen strafen und ihm sein verborgenes und tückisches Herz vorwerfen wirst, so kann er sich doch also stellen und Alles leugnen, daß er unüberwunden bleibt. Darum geht denn Jakob den Vertrag einfach ein, und eilt, daß der Bund aufgerichtet werde, und er selbst nimmt mit seiner Hand einen Stein und richtet ihn auf zu einem Male und zum Gedächtniß des Bundes, und denkt also: Wollte Gott, ich könnte dieses Buben nur bald los werden! Ja, er ermahnt seine Brüder, daß sie wollen Steine auflesen und einen Haufen machen. Dasselbe müssen wir auch thun, wir müssen die Gemeinschaft der Heuchler fliehen, auf welcherlei Weise es nur geschehen kann oder mag.

\begin{quote}
\textbf{V. 46. 47. Und sprach zu seinen Brüdern: Lest Steine auf. Und sie nahmen Steine, und machten einen Haufen, und aßen auf denselben Haufen. Und Laban hieß ihn Jegar Sahadutha; Jakob aber hieß ihn Gilead.}
\end{quote}

\quellseite{192}

253. Die Hirten und seine Verwandten nennt Jakob Brüder nach Gewohnheit der Schrift, nämlich, die von Nahor hergekommen waren. Und sie hätten zwar beiderseits wohl mögen zufrieden sein, da sie den Bund gemacht hatten; aber der allerheiligste Laban will noch mehr vergewissert sein. Es ist nicht genug, daß sie einen Stein zum Malzeichen und Gedächtnis des Bundes haben aufgerichtet, sondern er ist der Erste, der dem Steinhaufen einen Namen gibt auf syrische oder chaldäische Sprache, nur darum, daß er seinem Schwiegersohne das Gewissen treffen möge. Denn er will also sagen: Siehe, wir haben nun einen Bund gemacht, wir haben Steine aufgelesen und einen Haufen gemacht, dazu auch demselbigen Haufen einen Namen gegeben; nun will dir es gebühren, daß du auch an diesen Bund gedenkest. Nun siehe, wie große wunderliche Geistlichkeit er vorgeben kann, wie dieser Bund solle und müsse gehalten werden: so er doch droben, V. 41., den Lohn wohl zehnmal verändert hat, darum er mit Jakob war eins geworden; ja, wo er nicht wäre vom Gott gestraft worden, so hätte er mit ihm gar keinen Bund gemacht, sondern er würde Jakob wiederum mit Gewalt in die vorige Knechtschaft gebracht haben.

254. Das hebräische Wort gal heißt Haufe. Ed heißt Zeuge. Jakob sieht auf das achte Gebot, daß man kein falsch Zeugnis reden soll. Das syrische Wort aber heißt so viel, als das hebräische, ein Hügelchen von Steinen oder Erde.

255. Er nennt es aber nicht vergeblich einen Haufen des Zeugnisses; denn er will so viel sagen: Gedenke an diesen Haufen, der da aufgerichtet ist; laß dir denselbigen ein ewiges Gedächtnis sein in deinen Ohren und Herzen, daß du dieses Bundes nicht vergessest; denn wo du den nicht halten wirst, so wird dieser Haufen wider dich stehen zum herrlichen Zeugnisse im Himmel und auf Erden. Es wäre zwar meinetshalben nicht vonnöthen gewesen, diesen Haufen zu machen und aufzurichten, ich bin fromm und gerecht: aber wir haben von dir Gefahr zu fürchten. Also gibt er nur lauter Geistlichkeit und Gerechtigkeit vor, und will heimlich dem Jakob den Argwohn und die Schuld anhängen, als ob er ein ungetreuer Mann sei, der nicht Glauben halte.

256. Er fährt noch immer fort, Jakob das Denkmal einzuschärfen, und bestellt Alles sehr fleißig, bestimmt ihm nicht allein die Stätte, sondern auch die Zeit, hernach wird er ihm auch die Person anzeigen. Er ist ein sehr guter Notarius, er kann das Instrument wohl verwahren.

257. Mose aber hat dieses Stück hinzugesetzt: „Daher heißt man ihn Gilead.“ Denn es ist in der heiligen Schrift ein sehr berühmter Name, als, im Propheten Jeremia und an andern Orten mehr. Und wird auch metaphorisch auf die Stadt Jerusalem gezogen, die darum ein Hause des Zeugnisses genannt wird, weil sie eine heilige Stadt und eine Stadt Gottes ist, darin man den Gottesdienst, das Wort und die Verheißung Gottes gehabt hat. Es ist ein Hügel, der erhöht ist, und ein göttlicher Haufen, darin viel Zeugen Gottes gewohnt haben, die Gott haben reden hören und von ihm gezeuget mit ihrem Bekenntnisse, mit Predigten, Dankfesten. Auf diese Weise haben die Propheten dieses Wort als ein Gattungwort auf die Stadt Jerusalem und die ganze Kirche gezogen, da das Volk ist zusammen gekommen und da man Gottes Wort und Zeugniß gehört hat.

258. Es läßt sich aber Laban daran noch nicht genügen, daß er dieses also wiederholt hat; er kann den Bund, den sie gemacht hatten, nicht genugsam befestigen, wie die Heuchler und Ketzer pflegen; denn es ist nur ein unnütz Geschwätz. Wie das an den Sacramentsschwärmern auch wohl zu sehen ist, die sehr reich sind von Worten und ihr Ding nicht groß genug machen und schmücken können; unsere Lehre aber machen sie gar gering und grausam. Darum setzt er noch mehr dazu und sagt:

\begin{quote}
\textbf{V. 49. Und sei eine Warte, denn er sprach: Der Herr sehe darein zwischen mir und dir, wenn wir von einander kommen.}
\end{quote}

259. Mizpah ist auch eine berühmte Stätte in der heiligen Schrift (denn hernach ist es der Eigenname einer Stadt geworden), nicht weit von Gilead, wie im Buch Josua zu sehen ist. Als Gattungswort aber bezeichnet es eine Warte, darauf man stehen kann und die Feinde kommen sehen, desgleichen wenn ein Feuer aufkommt und wo sonst etwas Anderes zu sehen ist. Darum ist es nicht genug an dem Haufen des Zeugnisses, sondern dieweil er einmal von mir gewagt hat wegzuziehen, hat Laban gedacht, und mir meine Töchter zu entführen, so möchte er vielleicht auch dermaßen noch wiederkommen und mir das Leben nehmen, darum wird man einer Warte bedürfen. Ja, das noch mehr ist, er zeigt auch den Namen und die Gegenwart Gottes an. „Der Herr“, sagt er, „sehe darein zwischen mir und dir.“ Ich nehme nicht allein zum Zeugniß den Haufen hier auf Erden und unter den Menschen, sondern ich berufe mich auch auf Gott den Allmächtigen, der unser Zeuge sein und darein sehen soll. Hilf Gott, wie ist an dem Manne so große Geistlichkeit und Heiligkeit! wie ist er so beständig und aufrichtig im Gottesdienst und in der Anrufung Gottes!

260. Ja, also unterstehen sich diejenigen, die sich Gottes und der Göttlichkeit zu rühmen, die nicht Buße thun, sondern in ihren Sünden gar verhärtet und verstockt sind. Denn dieser Heuchler ist gar unbußfertig und ganz sicher wider den frommen heiligen Mann Jakob, und überschüttet ihn mit himmlischen und irdischen Zeugnissen, damit er ihn fein Gewissen beschweren will. Was soll aber nun Jakob dazu sagen oder thun? Laban ist allzu heilig, und darf Gott zum Zeugen rufen und daß er darein sehen soll. „Wenn wir von einander kommen“, sagt er, das ist, wenn du in dein Vaterland ziehen wirst und ich wiederum heim in mein Haus ziehen werde, und du etwas anfangen würdest ohne mein Wissen, wenn ich deine heimliche List, damit du mir möchte nachstellen, nicht zuvor sehen kann; darum will ich anstatt eines gewissen Schutzes und Festung diesen Haufen zwischen mir und dir hier bleiben lassen, darnach auch die Warte zwischen uns und Gott. Also ruft er beide Gott und Menschen zu Zeugen.

261. Diese Gedächtnisse des Bundes sind wahrlich gewiß genug. Wer wollte nun sagen, daß Laban nicht sollte ein heiliger Mann sein, dieweil er so getrost Gott anrufen und sich auf seine Hülfe und Schutz verlassen darf? Aber er thut solches nur darum, daß er Jakob damit schrecken und beschweren möge; denn beschwört er so theuer und hoch, als er immer kann, von dem er doch zuvor erfahren hatte, daß er sehr fromm und getreu war. Aber also pflegt sich die Heuchelei zu schmücken, und alle andern Menschen neben sich zu verdammen und zu schänden. Darnach sind auch alle Heuchler Lügner, darum hört Laban noch nicht auf, sondern setzt noch eine Beschwörung dazu, damit er ihn bezeuget.

\begin{quote}
\textbf{V. 50. Wo du meine Töchter beleidigest, oder andere Weiber dazu nimmst über meine Töchter. Es ist hie kein Mensch mit uns; siehe aber, Gott ist der Zeuge zwischen mir und dir.}
\end{quote}

262. Siehe doch, welch ein kunstreicher Poet mir das ist, allerlei Ursachen zu erdichten. Daß du meine Töchter nicht beleidigest, sagt er; das hat er nur vorgewandt und ist gar ein eitel Gedicht. Denn die Ursache, die er droben, V. 44., angeführt hat, ist der Wahrheit ähnlicher, nämlich, daß ihm entweder Jakob selbst, oder aber seine Verwandten möchten Schaden zufügen. Und er beschweret nun abermals den sehr frommen Mann und zieht ihn heimlich des Ehebruches, als hätte er ihn zuvor auch desselben in Verdacht gehabt, und daß er jetzt nun befürchten müsse, daß er seine Töchter würde sitzen lassen und davon ziehen.

263. Ach! was sollte der sehr fromme heilige Mann thun, der seine Weiber und Kinder so lieb gehabt, daß er um beider Weiber willen ganze vierzehn Jahre gedient und so lange Zeit solchen schweren Dienst ertragen hat; der auch die andere, die er doch nicht liebte, noch zum Weibe begehrte, genommen, und der seinen Schwiegervater viel Dienste erzeigt, mehr als einem Schwiegersohn oder Knecht wohl gebühret hätte? Wer wollte nun einen solchen frommen Mann beschuldigen oder in Verdacht haben, daß er mit seinen Weibern unfreundlich umgehen sollte?

264. Nun gibt aber Laban vor, wie er so väterlich gesinnet sei und rechte Neigung und Liebe zu seinen Kindern trage. Ach! sagt er, wie wird es nun meinen Töchtern gehen, dieweil sie so fern von mir hinweggeführt werden!

\quellseite{193}

Siehe du zu, daß du an diesen Bund gedenkest und sie weder mit Worten noch mit Werken beleidigest. Also beschweret er noch immer den frommen Jakob mit giftigen Lästerungen, und erdichtet und lügt doch alles, was er sagt. Und du hörst nicht ein Wort von ihm, damit er anzeigt, daß er seine Sünde erkennen, oder seinem Schwiegersohn und Kindeskindern für die Gewalt und Unrecht, so er ihnen gethan hatte, wollte genugthun. Er stellt sich, als ob er väterlicher Weise für seine Töchter Sorge trage. Aber er zeigt sich gegen sie gar nicht freundlicher oder wohlthätiger, denn er zuvor gethan hat.

265. Und er setzt das auch noch dazu, daß er sagt: „Es ist hie kein Mensch mit uns“, wir Zwei stehen allein bei diesem Haufen. Man mag aber wohl fragen: Ob sie allein mit einander geredet haben, daß keine Schiedsrichter dabei gewesen seien, oder aber, ob sie mit einander geredet haben, daß alle ihre Verwandten dabei gewesen und die Fremden allein ausgeschlossen sind? Man mag es auf beiderlei Weise verstehen. Wiewohl ich mir den Verstand am besten gefallen lasse, daß er die Fremden ausgeschlossen habe. Wiewohl wir nun allein sind, sagt er, und hier keine Leute sind, vor denen du dich dürftest fürchten, so ist doch Gott bei uns, der drein sieht und Zeuge ist zwischen mir und dir. Es geht alles dahin, daß der fromme unschuldige Mann, Jakob, soll beschweret werden, wo nicht mit öffentlicher Schmach und Schande, die er ihm gern wollte anhängen, so doch zum wenigsten, daß er beargwohnet und beneidet werden soll. Das ist der Heuchler Art; wie der Poet sagt: Crimine ab uno disce omnes etc.: An Einer Sünde magst du sie alle erkennen.

\begin{quote}
\textbf{V. 51. 52. Und Laban sprach weiter zu Jakob: Siehe, daß ist der Haufe, und das ist das Mal, das ich aufgerichtet habe zwischen mir und dir. Derselbe Haufen sei Zeuge, und das Mal sei auch Zeuge, wo ich hierher fahre zu dir, oder du herüber fährest zu mir über diesen Haufen und Mal, zu beschädigen.}
\end{quote}

266. Er wiederholt eben dasselbe, das er zuvor auch gesagt hat, und macht es groß und schmückt es mit seinen Worten. Aber endlich gibt er die rechte Ursache an den Tag, warum er den Bund so gar genau und sorgfältig gemacht habe, nachdem er mit Worten so greulich wider Jakob genommen hat. Er gedenkt nun seiner Töchter nicht, sondern zeigt an, daß er sich fürchte, daß Jakob irgend was Böses wider ihn vornehmen möchte, oder aber seine Verwandten, die dermaßen solche Gewalt und Unbilligkeit möchten rachen wollen. „Wo ich herüber fahre zu dir“, als wollte er sagen: Ich werde zwar zu dir nicht herüber fahren; wenn du auch schon in der äußersten Noth eines Trunkes kalten Wassers würdest bedürfen, so werde ich doch nicht bei dir sein, daß ich dir denselben gebe; wo du aber herüber fahren wirst, mich zu beschädigen, so wird dich dieser Haufe und dieses Mal strafen. Denn dieweil er sich schuldig weiß, so fürchtet er sich, daß es doch allenhalben sicher ist; er vertrauet niemand, und weiß sich gleichwohl fein zu schmücken mit dem Namen und Titel der Religion oder Gottesdienstes und des Bundes.

267. Aber in Wahrheit ist solch Schrecken und Furcht die größte Strafe der Heuchler und der Gottlosen. Denn ein böses Gewissen kann nicht ruhen oder stille sein. Es ist ein Hündlein, das heißt auf Deutsch Teufel; wenn derselbe schon im Leben ruht, so kommt er im Tode und bellt. So hat nun Jakob ein sehr gut und gewiß Zeugniß seines Gewissens gehabt, wie auch Laban selbst heimlich dafür gehalten hat. Laban aber weiß sich vieler Unbilligkeit und Untreue schuldig, die er seinem Schwiegersohn für seine sehr großen Wohlthaten erzeigt hat, damit er ihm sein Hab und Gut gemehret hatte. Denn Jakob hat sich an einen geringen Lohn genügen lassen, hat damit vorlieb genommen, daß ihn Laban hatte leer ziehen lassen. Darum denkt nun Laban: Wie, wenn er dieser Unbilligkeit gedenken würde, mir dermaßenst noch mit seinen Freunden zu schaffen machen möchte, und etwa noch fordern wollen, was ihm bei mir vom Erbtheil, oder von seinem verdienten Lohn rückständig ist? Darum will er sich verwahren, auf daß er nicht herüber fahre und irgend einen Lärm anrichte, oder Schaden thue. Wenn er aber wiederum in die vorige Knechtschaft und Elend kommen würde, damit würde er eher zufrieden sein. Ich will aber zu dir nicht herüber fahren, sagt Laban: ich werde nicht kommen, dich reich zu machen, oder dir Gutes zu thun, wie du zu mir gekommen bist.

\begin{quote}
\textbf{V. 53. Der Gott Abrahams, und der Gott Nahors, und der Gott ihrer Väter sei Richter zwischen uns.}
\end{quote}

268. Wie kann der Tropf das so heiß machen! Mit solchen harten Worten, die wie Donner-Schläge sind, versündigt er, den frommen heiligen Mann zu erschrecken und ihn als einen ungerechten und gottlosen Buben zu verleumden, dagegen aber seine Heiligkeit, die gar engelisch ist, zu rühmen. Er führt aber den Namen des Herrn unnützlich, und mit dem vielen Geschwätze verräth er sich selbst, daß er im Herzen ein Heuchler und loser Mißgönner sei; wie in den Sprüchen Salomos im 14. Cap. V. 23. steht: „Wo man mit Worten umgeht, da ist Mangel“, und wie man pfleget zu sagen: Viel Worte, nichts dahinter etc. Laban aber versteht den „Gott ihrer Väter“, den Tharah und Abraham gedienet hat, da er noch abgöttisch war zu Ur in Chaldäa; dieselbige Abgötterei ihrer Väter führt er hier an, davon im Buche Josua am 24. Cap. V. 2. gesagt wird: „Eure Väter wohneten vorzeiten jenseit des Wassers, Tharah, Abrahams und Nahors Vater, und dieneten andern Göttern“ etc.

269. So haben wir nun bisher ein Conterfei und Bild gehabt eines Heuchlers und Götzendieners, der alles, damit er nur umgeht, erdichtet und den Namen Gottes unnützlich führt. Wie wir aber an dem Heuchler nur eitel vergebliche Worte gehört haben, also macht dagegen Jakob seine Rede kurz und ist voller Gottseligkeit und der besten Früchte. Darum richtet er sein Ding alles fein fürzlich aus; denn er hat Laban zuvor hart genug gestraft. Dieweil er aber nun in seiner Sicherheit, Hoffahrt und gottlosem Wesen beharrt, läßt er ihn fahren; denn der Heuchler redet doch nichts Anderes, als nur eitel Lästerungen wider Gott und Schmachworte wider den frommen und heiligen Mann Jakob.

\begin{quote}
\textbf{V. 54. 55. Und Jakob schwur ihm bei der Furcht seines Vaters Isaak. Und Jakob opferte auf dem Berge, und lud seine Brüder zum Essen. Und da sie gegessen hatten, blieben sie auf dem Berge über Nacht. Des Morgens aber stund Laban frühe auf, küssete seine Kinder und Töchter, und segnete sie; und zog hin, und kam wieder an seinen Ort.}
\end{quote}

270. Jakob schwört nicht bei dem ‚Gott Nahors‘ oder Tharahs, sondern schwört bei einem näheren Gott, das ist, bei dem Gott, welchem sein Vater Isaak fürchtet und dem er dient, nämlich, bei Christo, nach den Befehlen und Verheißungen Gottes. Denn die Heuchler dienen auch Gott; aber mit Geboten und Satzungen der Menschen, wie der Prophet Jesaia am 29. Cap. V. 13. und Christus Matth. 15, 8. sagen: „Dieses Volk nahet sich zu mir mit seinem Munde, aber ihr Herz ist ferne von mir.“ Niemand naht sich mehr zu Gott mit dem Munde, niemand führt den Namen Gottes mehr und öftere, denn eben die Heuchler thun.

271. Aber Gott spricht, ihr Gottesdienst sei abgöttisch und falsch. „Ich höre zwar meinen Namen wohl, spricht er; ich werde von ihnen genannt ein Schöpfer Himmels und der Erden; aber Alles wider das andere Gebot. Denn sie nennen mich nicht also von Herzen, ja, das noch mehr ist, sie verlästern den Gottesdienst mit Menschengeboten. Sie fordern mich zum Rechte, und wollen mit Gott rechten, wie im Propheten Jesaia am 58. Cap. V. 2. stehet. Aber ich will nicht haben, daß man mir dienen soll mit solchem Dienst und Lehre, so die Menschen erwählt haben; sondern in der Furcht vor mir soll man mir dienen, das ist, daß du mein Wort mit Glauben annimmest, wie Jesaia am letzten Capitel V. 2. steht: „Ich sehe an den Elenden, und den zerbrochenes Geistes ist, und der sich fürchtet vor meinem Wort“; solche Leute sollen mir dienen. Item Jesaia am 8. Cap. V. 13.: „Heiliget den Herrn Zebaoth. Den laßt eure Furcht und Schrecken sein, so wird er eine Heiligung sein.“

272. Wenn nun Christus unsere Furcht sein soll, das ist, wenn wir ihm dienen sollen, so müssen wir das Wort und die Verheißung Gottes haben. Und eine solche Furcht oder Schrecken ist der rechte Gottesdienst, der in uns zuwege bringt und wirkt, daß wir alle anderen Schrecken und Furcht verachten, wie Jesaia Cap. 8, 12. 13. sagt: „Laßt euch nicht grauen; sondern heiliget den Herrn Zebaoth“, der wird euch wohl schützen; bedenket, daß man den allein fürchten, ehren und dienen soll, nämlich, der mit euch redet; wenn ihr das thut, so bin ich, der Herr, eure Furcht, denn ihr habt ein gewisses und starkes Wort; wenn ihr euch daran haltet, so werdet ihr im Gottesdienst nicht irren.

\quellseite{194}

273. Auf diese Weise hat Jesaia diesen Text verstanden, da er GOtt unseren Schrecken oder unsere Furcht nennt um des Gottesdienstes willen, der GOtt nach seinen Geboten und Verheißungen geleistet wird, nicht allein mit dem Munde, wie die Heuchler GOtt fürchten und ihm dienen ohne Wort mit Menschgeboten.

274. In solchem Glauben und Bekenntniß opfert nun Jakob, welches der Heuchler nicht thut; denn er läßt sich dünken, daß er frommer und heiliger sei, denn daß er des Opfers oder der Anrufung bedürfen sollte; denn er will mit seinen Sünden noch recht gethan haben. Jakob aber dankt GOtt, opfert und richtet ein Mahl an, das zum Opfer erfordert wurde: er ruft seinen Freunden, daß sie zusammenkommen, und daselbst wird er ohne Zweifel eine Predigt gethan haben. Laban achtet des Opfers nicht, wiewohl er vielleicht auch mitgegessen hat; aber Jakob lobt GOtt und die Furcht seines Vaters, das ist, Christum, daß er dieses Buben ist los geworden. Da nun das Mahl gehalten ist, zieht Laban davon. Nun sind wir des losen Tropfes auch los.

275. Ehe er aber wegzieht, küßt der Heuchler noch seine Kinder und Töchter; zwar nicht aus väterlicher Neigung oder Herzen, sondern daß er sich nur äußerlich also stellt, als sei er so leutselig und freundlich, wie es sich gebühret nach der Weise und Gewohnheit aller Völker. Denn er läßt sie ziehen ohne alle Geschenke. Er segnet sie wohl mit Worten, aber inwendig im Herzen trägt er solchen Neid und bittere Feindschaft, daß er ihnen auch das nicht gönnt, was sie durch GOttes Segen ohne seine Gaben erlangt hatten. Dieses sind nun zwei entgegengesetzte Bilder, nämlich, eines sehr greulichen Heuchlers, der ohne Aufhören sündigt, und ist ein Exempel aller Sünde, der Abgötterei, Heuchelei, des Geizes und endlicher Unbußfertigkeit; darnach eines sehr frommen und heiligen Mannes, der mit solchem Ungeheuer hat kämpfen müssen.

\chapter{Das zweiunddreißigste Capitel}

\section[Erster Theil]{Erster Theil. Wie Jakob wieder in sein Land zieht, und seinem Bruder Esau Boten schickt und ihn grüßen läßt}

\begin{quote}
\textbf{V. 1. 2. Jakob aber zog seinen Weg; und es begegneten ihm die Engel GOttes. Und da er sie sahe, sprach er: Es sind GOttes Heere; und hieß dieselbige Stätte Mahanaim.}
\end{quote}

1. Nachdem Laban hingezogen ist, und der heilige Patriarch Jakob von der schweren und langwierigen Mühsäligkeit und Dienst seines Schwiegervaters erlöset worden, wird er froh, daß er zuletzt, da die Trübsal ein Ende hat, wiederum Friede und Trost bekommen hat, und zieht nun immer seinen Weg fort, auf daß er wieder zu seinem allerliebsten Vater Isaak kommen möge, der schon lange durchs Alter kraftlos geworden war.

2. Mose aber sagt, daß ihm, da er hingezogen sei, die Engel GOttes begegnet seien, und da er sie gesehen, habe er mit großer Freude erkannt, daß es die Heere und Heerscharen GOttes wären, und hat dieselbe Stätte geheißen machanaim; denn er sagete: Diese Heerschar ist GOttes Heerschar; wiewohl wir auf Deutsch das hebräische Wort mit dem Singular nicht wiedergeben können; denn es ist ein Wort, das nur im Plural gebraucht wird, oder vielmehr im Dual, der auch im Griechischen vielfach vorkommt. Machanaim heißen zwei Heere, und sind solche Worte, damit der Patriarch zu verstehen gibt, daß er fröhlich sei und triumphire, ja, daß er nun getröstet und ohne Sorge sei, darum daß ihm GOtt Frieden verliehen hat. Als wollte er sagen: Jetzt erscheinen mir nun die Engel, der Himmel lacht mich an, das Ungewitter und der Winter ist vorüber, jetzt scheint nun ein sehr schönes helles Licht.

3. Zu solcher stillen Freude, dieweil er aller Mühe und Beschwerung ist los geworden, kommt nun das auch noch, daß ihm die Engel und Heerscharen GOttes begegnen. Es hat aber Mose vorhin auch angezeigt, daß den Patriarchen zum öftern die Engel erschienen sind; wie er droben von Abraham 1 Mos. 18, 1., von Hagar Cap. 16, 7., von der Leiter Jakobs Cap. 28, 12. ff., da die Engel auf und nieder fliegen, erzählt hat; item, Cap. 31, 3., da Jakob vom Engel ermahnt worden, daß er wiederum in sein Vaterland ziehen sollte. Darum haben sie oftmals Engel gesehen, und hat auch den frommen und heiligen Vätern vornehmlich noth gethan, daß sie dieselbigen haben sehen mögen. Aber Jakob sieht hier nicht einen Engel, sondern ganze Heerscharen der Engel; wie der Evangelist Lucas sagt im 2. Cap. 13., da er die Geburt Christi beschreibt: „Und alsobald war bei dem Engel die Menge der himmlischen Heerscharen.“

4. So werden nun die Engel Heerscharen genannt, item, Kriegsgefächte, Wächter, Führer und Regenten bei den Creaturen GOttes. Denn das ist hienieden auf Erden ihr Amt. Droben ist ihr Amt, daß sie Gloria in excelsis singen: „Ehre sei GOtt in der Höhe“ 2c.; hiednieden gegen uns und die Creaturen ist ihr Amt, daß sie wachen, regieren und streiten nicht allein für die Gottseligen, sondern auch für die Gottlosen. Wie im Propheten Daniel zu sehen ist, da der Engel sagt, daß er gekommen sei, wider den Fürsten in Persien zu streiten, Dan. 10, 13. Nun streitet aber ein guter Engel wider den andern guten Engel nicht; darum ist der Fürst in Persien ein böser Engel gewesen, und hat unter die Zahl gehört der Engel, von welchen Christus Joh. 12, 31. sagt: „Der Fürst dieser Welt wird ausgestoßen werden.“ Der Satan, der ein Gott ist dieser Welt, hat große Heerschaaren der Teufel, und auch unter den bösen Engeln ist eine Art Monarchie.

5. Denn das weiset sich an sich selbst aus und die Schrift zeiget es auch klärlich an, sintemal wir sehen, daß Alles in der Welt greulich untereinander gemenget, verwirrt und unordentlich vermischt wird; zudem läuft es allenthalben wider einander mit vielen erschrecklichen Sünden und Schanden. Das kommt daher, weil allenthalben die bösen Engel regieren, in den Höfen des Papstes, des Kaisers, der Könige und Fürsten, ja, auch wohl in den Häusern der gemeinen Unterthanen. Das sind und werden auch genannt Werke des Teufels, die wir sehen und erfahren; die Welt aber sieht nicht, daß Todtschlag und andere unzählige Sünden und Schanden mehr geschehen durch des Teufels Gewalt und Anreizung. Das Werk sehen sie alle wohl: aber die Ursache, daher es kommt, sehen sie nicht.

6. Darum ist es gewiß, daß die Fürsten unter den guten Engeln wider die Fürsten der bösen Engel streiten. Denn das bringt die Erfahrung also mit, wie auch die Heiden etlichermaßen einen Schatten von dieser Erkenntniß behalten haben, da sie gedichtet haben, daß ihres Hausgötter oder gute und böse Geister sein sollen. GOtt ist ein Schöpfer und Regierer aller Dinge, der auch die Herrschaften der Gottlosen, als, der Babylonier, Assyrer und Perser, durch die Engel verwaltet; wie Daniel das also bezeuget, und man eben solches an unsern Fürstenhöfen und Regimenten sieht. Daher kommt es, daß zum öftern auch die besten Sachen an den Fürstenhöfen auf mancherlei Weise verhindert, verwirrt und verzogen werden, und aber doch endlich einen glückseligen Ausgang gewinnen durch Anleitung und Rath der guten Engel. Und das ist unseres Herrn GOttes wunderbarer Rath, davon man nicht disputiren soll, warum er die Welt auf diese oder andere Weise regiere.

7. Das ist aber eine große und wunderbare Weisheit, daß Jakob die Engel, so ihm begegnet sind, erkennen und nennen kann GOttes Heere oder unsers HErrn GOttes Kriegsvölk. Hat denn GOtt hier auf Erden Heere und Kriegsvölk? Jakob nennt wahrlich alle Engel also. Und diesen Verstand hat er von Laban, seinem Schwiegervater, nicht gehabt, der ein gottloser Götzendiener war und solche hohe und himmlische Dinge nicht gelehrt hatte; sondern solche Erkenntniß ist von den Vätern immer von Einem auf den Andern gekommen. Denn Abraham hat Noah gesehen länger denn fünfzig Jahre, Sem aber hat auch noch nach Abraham gelebet; und denselben haben Isaak und Jakob gesehen und gehört; der ist ein sehr großer Lehrer, ein Prophet und Priester des Höchsten gewesen, der diese Lehre von Enos, der Adams Enkel gewesen ist, genommen hat.

8. Dies sind die Lehrer und Zuhörer gewesen, die diese Lehre bewährt und dieselbige bei den Nachkommen ausgebreitet haben, und ist auch zugleich mit dazu gekommen die Erfahrung und Trübsal. Wer aber nun solche Führer und Lehrer hat, wird leichtlich viel lernen können. Und Jakob hat es nicht allein von den Vätern gelernt, sondern die Erfahrung hat es ihn auch gelehrt: er hat die Engel sehen auf und nieder steigen; item, in Mesopotamien bei seinem Schwiegervater hat er im Traume einen Engel gesehen und gehört, und an diesem Ort sieht er nun die Heere und Heeresscharen der Engel. Derhalben kommen zusammen beide die That oder das Werk an sich selbst und die Worte, oder die Erkenntniß und das Werk und die Erfahrung.

\quellseite{195}

9. Sonst ist diese Lehre und Weisheit etwas höher, denn daß man sie mit menschlicher Vernunft sollte fassen oder begreifen können. Denn die Vernunft weiß davon nichts, daß die Engel Acht haben und Sorge tragen für die Herrschaften und Regimente, für die Haushaltung, für Menschen und Vieh, und endlich auch für alle Creaturen. Denn also beschreibt sie die Epistel an die Hebräer am 1. Cap. V. 14.: „Sind sie nicht allzumal dienstbare Geister, ausgesandt zum Dienst, um der willen, die erben sollen die Seligkeit?“ Sie sind keine Götter oder Göttinnen, sondern sind Diener, die die ganze Welt regieren, und dasselbe um derer willen, so die ewige Seligkeit ererben sollen. Denn was in diesem Leben geschieht, das geschieht alles um der gottseeligen Menschen willen, und derer, die da sollen selig werden. Um derselben willen scheint die Sonne, es werden die Herrschaften und Regimente bewahret und erhalten; das Land wird fruchtbar, und werden auch um ihretwillen die Ehen und Heirathen gestiftet. Kurz, alles, was im Himmel und auf Erden ist, das ist zu dem Ende verordnet und gerichtet, daß die Gerechten gesammelt, und die Zahl derer, die da sollen selig werden, erfüllt werden möge. Dies ist nun eine rechte himmlische Lehre, und menschliche Vernunft und Weisheit lehrt davon nichts, nämlich, daß in diesem Leben durch den Dienst der heiligen Engel die Herrschaften, dazu das Welt- und Hausregiment, und endlich alles, was diese Welt hat, regieret wird.

10. Nun erhebt sich aber hier eine Frage: Wenn dies also wahrhaftig und gewiß ist, wie es denn eigentlich wahr ist, so folgt daraus, daß die Meinung des Erasmus und aller Epicurer muß höchst falsch sein. Denn dieselben werden daraus, daß alle Dinge so gar ungleich und verkehrt sind in diesem Leben, und daß sie sehen, daß die bösen Menschen mehr Glück haben als die Frommen und Gottseeligen, bewogen, daß sie dahin fahren und diesen Schutz und Dienst der Engel, ja, die Vorsehung Gottes selbst gar verleugnen. Und die gemeine Erfahrung und auch die Klage der Heiligen bezeuget das auch: Je schalkhaftiger Einer ist, je besser Glück er hat. Also sagt auch Cicero: Optime cogitata pessime cadunt: Was man am allerbesten bedacht und betrachtet hat, geräth am allerübelsten. Also findet man, daß Julius Cäsar, Demosthenes und Andere mehr, die sehr weise, gelehrte, tugendreiche und verständige Männer gewesen sind, geklagt haben, daß ihre Rathschläge keinen glückseligen Fortgang und Ende gehabt. Ja, auch ihr Ende war im Verhältniß zu ihren hohen Verdiensten um den Staat sehr traurig. Also werden die Gottseeligen auch unterdrückt, die Kirche wird verfolgt und hart geplagt, die Gottlosen thun sich hervor, stehen in großem Flor, sind reich und gewaltig. Was ist denn dies für ein Regiment, das die heiligen Engel verwalten? Wer sollte das nicht sehen, daß alle Dinge ohne Gottes Rath und Regierung, und ohne Gefahr geschehen und untereinander geworfen werden? Ja, daß dies vielmehr des Teufels Regiment sei? Wie ihn Christus Joh. 16, 11. einen Fürsten dieser Welt, und St. Paulus 2 Cor. 4, 4. einen Gott dieser Welt nennt. Wie kann aber nun das Regiment der Engel bestehen?

11. Antwort: Dies ist, wie ich zuvor gesagt habe, eine wunderbare und unbegreifliche Weisheit Gottes, welche menschliche Vernunft nicht verstehen kann. Wiewohl wenn wir die Augen wollen aufthun, so können wir wohl sehen und urtheilen, daß der Teufel gebunden ist und gefangen liegt, als wäre er in Eisen und Ketten gespannt, also daß er auch nicht ein Haar auf unserem Haupte kann anrühren, es sei denn, daß es Gott haben wolle und ihm es zulasse. Wo uns aber irgend ein Unglück widerfährt, das sollen wir mit nichten dahin deuten, als sollte es daher kommen, daß die Engel in ihrem Amte nicht sollten fleißig genug sein; sondern wir sollten es lieber auf die Anfechtung ziehen, damit in diesem Leben ein jeder gottseelige Mensch und die ganze Kirche geplagt wird. Denn der Teufel hat so große Gewalt nicht, wie es äußerlich scheint. Denn wenn er vollkommene Gewalt hätte zu tödten und zu wüthen, wie es ihm gelüsten möchte, so würdest du nicht eine Stunde leben, du würdest nicht ein Stück Vieh, nicht die Frucht auf dem Felde, auch nicht das Getreide in der Scheune, und endlich der Dinge keines, so zu diesem Leben dienen, unverletzt und unverdorben behalten.

12. Nun betrachte aber in deinem Herzen die ganze Ordnung der Natur und dieses ganzen Lebens, und beschau allerlei Art der Menschen, des Viehes, der Vögel und Fische: so wirst du finden, daß mehr Gutes, als Böses ist; du wirst auch sehen, daß der geringste Theil der Gewalt des Teufels unterworfen sei. Denn er muß ja die Fische im Wasser, die Vögel in der Luft, die Menschen und das Vieh in Dörfern und Städten bleiben lassen, was er nicht thun würde, wo es ohne den Schutz der heiligen Engel wäre. Ja, sprichst du, er richtet gleichwohl bisweilen großen Lärm an: er macht, daß ein Reich wider das andere, ein Fürst wider den andern ist, er verwirret untereinander die Lande, Städte und ganze Geschlechter. Antwort: Es ist wahr, er richtet solchen Lärm an, und verwirret gerne, wo er nur etwas verwirren kann; er kann aber gleichwohl das nicht ausrichten, daß er am allermeisten begehrt, nämlich, daß er Alles möchte umkehren und das Unterste nach oben wenden, und Alles über einen Haufen werfen. So stark sind die Wälle, die Festungen und Zäune, so die heiligen Engel um uns und um alle Dinge gemacht haben.

13. Wenn man nun das Böse gegen das Gute, so in diesem Leben geschieht, sollte halten, so wird des Guten viel mehr sein als des Bösen, ausgenommen was zur Anfechtung, oder zur Strafe und Züchtigung der Sünden gehört, wo etwa die Sünden überhand nehmen, und die Königreiche, Lande, Städte und Fürsten bald gegeben werden, daß sie gestraft und verwüstet sollen werden; wie Gott Sodom und die umliegenden Städte verderbet hat: da hören dann die Wälle und Heerlager der Engel auf. Im Uebrigen so richtet er im weltlichen Regimente und sonst in allen Ständen dieses Lebens mancherlei Tumult und Lärm an: er treibt die Leute zum Ehebruche, zum Diebstahl, zum Todtschlag, und daß sie falsch schwören; er verführt oder verderbet sie aber gleichwohl nicht alle. Denn er ist gezwungen, daß er auch unter den bösen Menschen, die nicht tödten kann, die er wohl gerne tödten wollte. Denn wo er Alles in seiner Hand und Gewalt hätte, so würde er auch die Gottlosen umbringen, wiewohl er sie vorher in seinen Stricken gefangen hält. Darum duldet Gott auch die bösen Menschen und Sünder, seine unermäßliche Güte und Geduld damit zu beweisen; aber bis auf die Zeit, die zur Strafe bestimmt ist. Wenn nun ihre Sünden erfüllt sind, dann zieht er die Hand ab.

14. Auf solche Weise mag man antworten auf die Frage: Wenn die Engel Gottes Herrschaften und dienstbare Geister sind, warum denn gleichwohl noch so viel Böses geschehe, das Gott mißfällt? Den Gottlosen geht es wohl, es wird ihnen das Leben, Ehre und Herrlichkeit gegeben, sie haben groß Geld und Gut; dies wird wahrlich den bösen Menschen alles gegeben durch die guten Engel. Ich antworte aber darauf also: Daß man von dem göttlichen Rath nicht disputiren soll, nämlich, warum er den Bösen auch etwas Gutes gebe, und seine Gaben ausstreue ebensowohl unter die Guten, als unter die Bösen, unter die Dankbaren und Undankbaren; wie Christus spricht Matth. 5, 45. Denn das geschieht darum, daß er damit anzeige, daß er nicht nur menschliche Güte an sich habe, die ihr Ziel und Maß hat, sondern vielmehr, daß seine Güte unermäßlich, unendlich und unbegreiflich sei.

15. Darum sollen wir ihm sein Regiment lassen und seine sehr große Gnade und Barmherzigkeit rühmen, sintemal es offenbar ist, daß mehr Gutes als Böses geschieht auch an den bösen Menschen, und die Gott lästern, welche auch ihre leiblichen Güter haben. Denn wo Gott auch nur an Einem Tage die Welt durch die Engel nicht regierete, so würde bald in einem Hauf das ganze menschliche Geschlecht gar vergehen, der Teufel würde Alles hinreißen und wegtreiben, er würde Alles mit Theurung, Pestilenz, Krieg und Feuer verderben. Und solchen Schaden würden nicht allein die Bösen, sondern auch die Frommen tragen müssen.

\quellseite{196}

Daß wir aber nur solcher großen Gefahr frei und sicher sein könnten, da sollen wir es gewißlich dafür halten, daß solches durch den Schutz der himmlischen Heerschaaren geschehe.

16. Es lassen es zwar die Engel bisweilen zu, daß etwas Böses geschehe; aber solches geschieht darum, daß wir versucht und unser Glaube auch geprüft und geübt werde; ja, daß wir lernen GOtt erkennen in seinem wunderbaren Rath und Werken und daß wir ihm für die wunderbare Regierung danken sollen. Wie uns dieses Exempel des Patriarchen Jakob lehrt; denn derselbe zieht immer fort, ist fröhlich und getrost, darum daß ihm die Sonne aufgegangen ist und die himmlischen Heerschaaren ihm begegnet sind. Es wird aber bald hernach die Sonne wiederum untergehen. Jetzt regiert er und triumphirt, ist voll Glaubens und geistlicher Freude; aber bald wird er in den größten Jammer, Elend und Bestürzung geworfen werden.

17. Eben auf dieselbe Weise führt und regiert GOtt beide, die Gottseligen und Gottlosen. Wenn es sich jetzt ansehen läßt, daß Alles über einen Haufen fallen wolle, so soll man doch nicht gar alle Hoffnung und Vertrauen fallen lassen; sondern du sollst immer harren, dulden und die Sache abwarten. GOtt lebet noch, die Engel haben noch das Regiment und beweisen Schutz. Und wiewohl es sich ansehen läßt, daß sie ihres Amtes, und der Hut, so ihnen über uns befohlen ist, vergessen haben: so hören sie doch in Wahrheit damit nicht auf, sondern es ist eine Versuchung und Prüfung, damit GOtt nicht allein die Gottseligen, sondern auch die Gottlosen übt; wiewohl es zwar die Gottlosen nicht verstehen. Wie es zugeht, wenn plötzlich große Ungewitter der Winde und Platzregen kommen, und die heftigen Donnerschläge dem Himmel und der Erde drohen, daß Alles solle verbrannt, oder daß der ganze Erdboden fallen und untergehen werde: wenn aber das Ungewitter gestillt ist, so grünt das Land, und zeigt an, daß das Ungewitter, wie schrecklich es auch gewesen ist, zur Fruchtbarkeit sehr nütze gewesen sei.

18. Das ist nun die Lehre, die an diesem Ort uns vorgehalten wird, nämlich, daß die Engel dienstbare Geister und Diener der Creaturen sind, die da streiten für die Wohlfahrt und das Heil der Welt und der gottseligen Menschen. Und dasselbe ist hienieden auf Erden ihr Amt. Droben aber im Himmel ist ihr Amt, daß sie singen: Gloria in excelsis Deo: „Ehre sei GOtt in der Höhe“, und: Te Deum laudamus etc.: „HErr GOtt, dich loben wir“; item, wie Christus spricht Matth. 18, 10.: „Die Engel im Himmel sehen allezeit das Angesicht meines Vaters“. Und dasselbe thun sie zu GOttes Ehre, zur eigenen und aller Gläubigen Freude.

19. Die wunderbare Regierung aber dieses Lebens verstehen sie sehr wohl, nämlich, wie sich das Glück der Gottlosen mit dem Unglück der Gottseligen reime; welches wir in diesem Leben, da wir noch das Fleisch am Halse tragen, nicht sehen oder verstehen können; aber am Ende der Welt und nach diesem Leben werden wir sehen die allerschönste Einhelligkeit und wie diese Regierung so fein zusammenstimmt.

20. Wie ein grober bäurischer Mensch den wunderbaren Zusammenklang so vieler Stimmen auf der Orgel oder Harfe auch nicht versteht, darum daß er von der ganzen Harmonie nichts weiß: also meinen wir auch, daß Alles ohngefähre geschehe. Wir lassen uns dünken, daß GOtt schlafe, und der Teufel wache und regiere. Und kommt auch zugleich das Urtheil der Vernunft dazu, welches Urtheil schließt, daß weder GOtt noch Menschen die Welt regieren, sondern daß alles, was mit uns Menschen hier auf Erden geschieht, gar ohngefähr geschehe. Das ist die Weisheit des Fleisches; nicht die himmlische und unvergleichliche Erkenntniß der Regierung GOttes, welche anzeigt: daß der Dinge in der Welt viel mehr sind, die erhalten, denn die verdorben werden; desgleichen, daß auch die Güteigkeit GOttes größer sei, als die Bosheit des Teufels sein kann. Aber die Vernunft und das Fleisch zweifelt und wankt, auch noch in den Heiligen, wenn es solche große Unordnung und Ungleichheit des Geschickes sieht. Das kommt daher, weil wir noch in demselben Licht nicht sind, daß sich die Engel freuen, und darin sie sehen, wie Sünde und Gerechtigkeit, Tod und Leben, Finsterniß und Licht mit einander übereinstimmen.

21. So lasset uns nun lernen, daß die Regierung der Engel gewiß ist, und daß die unermessliche Güte und Barmherzigkeit GOttes die Bosheit des Teufels weit übertrifft; und daß GOtt diese Dinge auf Erden alle wunderbarer Weise erhalte wider allen grimmigen Zorn, Wüthen und Anfall des Teufels. Denn also sehen wir, daß es mit allen Stücken in diesem Leben zugeht. Es verliert bisweilen ein Hausvater irgend ein Schaaf, ein Huhn, oder Geld; er behält aber noch das Haus und seine Aecker, und den größten und besten Theil seiner Habe.

22. Also ist das dieser Stadt ein großer Schmuck und Ehre, daß die hohe Schule hier ist, und so viel feiner gelehrter Leute, die mit der himmlischen Lehre und guten Künsten umgehen; aber unter denen sind etliche Buben, die nicht erkennen, daß sie in dieser Versammlung, als da die Musen ihre Wohnung haben, ja, da der rechte Tempel GOttes und der heiligen Engel ist, sein nüchtern und züchtig leben sollen; sondern richten immer Lärm und Tumult an, als ob sie unter den groben und tyrannischen Leuten wohneten, die Cyclopen oder Centauren sind genannt worden, und schänden diese herrliche Versammlung mit bösen Sitten und vielem Aergerniß. Wie sich aber diese Ungleichheit zusammenreime, das sieht GOtt und die heiligen Engel wohl. Vielleicht, wo die Zeit der Strafe noch nicht da ist, und daß die Schule zerstört werden soll, so ist es eine Versuchung, die bald aufhören wird, und GOtt wird darnach andere Früchte und den Segen geben von denen, die da gottselig und züchtig leben und ihren Lehrern gehorsam sind.

23. Bisher haben wir gesagt von den heiligen Engeln, wie dieselbigen uns behüten und schützen. Solche Erkenntniß, und daß man sich darauf verlassen möge, ist vornehmlich den Gottseligen nütze und nöthig. Wie Jakob diese himmlischen Heerschaaren mit großer Freude angesehen hat, und ist jetzt fröhlich und triumphirt, darum daß er nun alle Mühsaligkeit und alle Teufel unter seinem ungerechten Schwiegervater Laban überwunden hat. Er ist, sage ich, fröhlich und ohne Sorge im Herrn, und hat die heiligen Engel, die sind seine Mitgenossen und Zeugen dieser Freude.

24. Aber man muß fleißig darauf merken, daß das im Leben der Heiligen gebräuchlich ist, daß es sich mit ihnen immer ändert. Denn bisweilen geht es ihnen also, daß sie fröhlich und getrost sind, bisweilen aber trägt es sich mit ihnen zu, daß sie Traurigkeit und Beschwerung haben. Wie denn solche Veränderung auch in der ganzen Natur bei allen Dingen sich befindet, als, die Veränderung des Tages und der Nacht; item, des Lichtes und der Finsterniß, der Berge und Thäler, wie im 104. Psalm V. 8. steht: „Die Berge gehen hoch hervor, und die Berge treten sich herunter, zum Ort, den du ihnen gegründet hast.“ Bisweilen hören wir etwas, das uns erfreut, wo nicht in der Welt, so ist es doch im Himmel durch GOttes Wort: bisweilen kommt auch die Anfechtung und Traurigkeit des Geistes wieder. An solche Veränderung müssen sich die Gottseligen gewöhnen; denn es wird bald eine wunderbare Veränderung folgen.

25. Das Wort aber machanaim verstehe ich also, daß es nicht ein gemein Wort, sondern der Eigenname eines Orts oder Städte in der heiligen Schrift sei. Und man soll diesen Namen darum merken, dieweil David im 34. Psalm V. 8. den Spruch daher genommen hat, da er sagt: „Der Engel des HErrn lagert sich um die her, so ihn fürchten, und hilft ihnen aus“, welches des vornehmsten und merklichsten Trostes einer ist in den Psalmen, den wir auf uns ziehen sollen, daß wir sein gebrauchen mögen.

26. Ja, sprichst du, sehe ich es doch nicht, daß ich den Dienst der Engel um mich habe, und das noch mehr ist, ich empfinde wohl das Gegentheil, und lasse mich dünken, daß ich in des Teufels Gewalt dahin gegeben und gleichsam in die Hölle geführt sei. Antwort: Dafür sollst du es mit nichten halten. Denn wo du dem Teufel überlegen wärest, so würde er dich auch nicht eine Stunde leben lassen, sondern er würde dich in alle Schande und Laster stürzen; ja, er würde dir noch so viel Zeit und Raum nicht lassen, daß er dich in Laster und Sünde führete, sondern er würde dich alsbald tödten. Das du aber noch lebst, sollst du dem Schutz und Schirm der heiligen Engel zuschreiben. Oder, wenn du aus diesem Leben fahren und ja dem grimmigen Zorn des Satans unterliegen mußt, wenn es GOtt also verhängt, so ist doch das eine große Gnade und Wohltat, daß du einen solchen GOtt hast, der dich tröstet und stärkt durch sein Wort.

27. Daß du nun dem Satan also in seine Hände und Gewalt gegeben bist, wie du dir es vor deinen Augen und in deinem Herzen eingebildet hast, daßselbige ist eine Versuchung, die nur eine Stunde und Augenblick währet, nicht zum Verderben oder Verdammniß, sondern zur Prüfung, zur Seligkeit, zu innerlichem Segen. Wie Christus spricht Joh. 12, 24: „Es sei denn, daß das Weizenkorn in die Erde falle, und ersterbe, so bleibets alleine; wo es aber erstirbt, so bringts viel Früchte.“ Denn also ist Christus auch in die Hände der Mörder gegeben worden; aber nur eine Stunde lang und zur Seligkeit. Derhalben wenn du fühlst, daß du vom Teufel geplagt und angefochten wirst, so bete fleißig und danke Gott, daß du nicht verderbet, sondern nur versucht und bewähret bist; wie im 3. Capitel der Klagelieder Jeremias V. 22. steht: „Die Güte des Herrn ist, daß wir nicht gar aus sind; seine Barmherzigkeit hat noch kein Ende“ etc. Item St. Paulus sagt 2 Cor. 4, 9: „Wir leiden Verfolgung, aber wir werden nicht verlassen. Wir werden untergedrückt; aber wir kommen nicht um“ etc.

\quellseite{197}

28. Daselbige lehrt uns auch das Exempel Hiobs, dem erstlich der Teufel alle sein Gut mit den Kindern hinweggenommen, und läßt ihm das lästige, verdrießliche, schmähsüchtige Weib; zudem plagt ihn der Teufel auch noch an seinem Leibe mit sehr bösen Schwären. Und spricht doch gleichwohl der Teufel zu Gott: „Hast du doch ihn, sein Haus und alles, was er hat, rings umher verwahret“ etc., Hiob 1, 10. Da klagt der Teufel, daß die Gottseeligen und dazu die Gottlosen rings umher verwahret sind, und droht ihm hart mit heimlichem Zorn, daß er ihn am Leib und Gut verderben wolle. Als wollte er sagen: Wenn Hiob in meiner Gewalt wäre, so wollte ich ihm eine Macht an ihm beweisen, und wie er dich loben sollte: ich wollte bald zuwege bringen, daß er dir öffentlich unter die Augen fluchen sollte. Es verhängt ihm aber der Herr, daß er ihn am Leibe plagen und sein Gut nehmen mag, das Leben darf er ihm aber nicht nehmen.

29. Solche Exempel lehren nun, daß alle Bosheit und Plage, so der Teufel an uns übt, nur eine Züchtigung und Strafe sei, dadurch wir erweckt und munter gemacht werden, daß wir nicht schnarchen und faul werden, und also unter Andern, die sicher und faul sind, auch untergedrückt und verderbet werden. Es wird aber hiermit auch zugleich bestätigt und bleibt fest und gewiß stehen diese Lehre, daß die Engel Gottes Heerscharen sind, unangesehen, daß gleichwohl in diesem Leben noch viel Böses geschieht um der bösen Engel willen.

\subsection{II}

\begin{quote}
\textbf{3—5. Jakob aber schickte Boten vor ihm her, zu seinem Bruder Esau, ins Land Seir, in der Gegend Edom. Und befahl ihnen, und sprach: Also saget meinem Herrn Esau: Dein Knecht Jakob läßt dir sagen: Ich bin bis daher bei Laban lange außen gewesen, und habe Rinder und Esel, Schafe, Knechte und Mägde; und habe ausgesandt, dir, meinem Herrn, anzusagen, daß ich Gnade vor deinen Augen fände.}
\end{quote}

30. Maleach ist eben dasselbe Wort, damit er droben die Engel genannt hat. Denn es ist ein Name des Amtes und nicht der Natur. Denn nach dem Amt sind sie Boten, sie sind Kriegsknechte, die auf Posten und Schildwache stehen für die ganze Welt; aber nach ihrer Natur sind es Geister.

31. Jakob aber ist noch fröhlich, und zieht immer fort mit Freuden, fürchtet sich nicht, erschrickt auch nicht; denn er zieht aus Mesopotamien auf Gottes Befehl, der ihn im Traum durch den Engel gestärkt hat, und der Laban gestraft und inne gehalten hat. Zudem hat er gesehen die Heerscharen Gottes. Darum hat er nun mit seinem Bruder Esau gütlich und freundlich handeln wollen, und ist wohl glaublich, daß er ernstlich und von Herzen Verlangen gehabt hat, seinen Bruder zu sehen; denn er denkt, er werde nun in diesen zwanzig Jahren wiederum sein Freund geworden sein. Darum schickt er Boten vor sich her und läßt ihn grüßen, daß er daraus vernehmen möge, daß er ihn als seinen Bruder lieben und in Ehren halten wolle. Hernach aber wird er ihm Geschenke schicken, da er hört, daß Esau ihm mit vierhundert Mann entgegen ziehe. Dasselbe wird etwas schwer sein und Jakob hätte sich dessen nicht versehen.

32. Das lose Geschwätz der Juden, dem Lyra an diesem Orte, wie es scheint, folgt, ist eitel unnütz Ding; denn sie sagen, Jakob sei erschrocken, darum daß er zu lange bei Laban, dem Götzendiener, gewohnt hatte und von ihm etlichermaßen wäre befleckt worden; item, daß Jakob sollte gedacht haben, daß die Verheißungen eine solche Bedingung hätten, daß darin ein schuldiger Gehorsam gegen Gott gefordert würde. Man soll aber von den Vätern mit nichten also halten. Wiewohl der Unterschied recht ist, daß zweierlei Verheißungen sind, nämlich, etliche, die man heißt bedingte, die eine sonderliche Bedingung haben, und etliche, die einfach ohne Bedingung sind, als, die Verheißung des Gesetzes hat eine Bedingung, aber die Verheißung der Gnade ist einfach ohne Bedingung. Wenn dem Arbeiter ein Lohn versprochen wird, so ist da eine solche Bedingung, so die Arbeit und Pflicht fordert, wie der Vertrag mit sich bringt; wo nun die Arbeit und schuldige Pflicht nicht folgt, so wird auch der Lohn nicht ausgegeben und entrichtet.

33. Aber solche Verheißungen sollen wir fahren lassen, wenn man mit Gott im Gewissen handeln soll; denn wir werden damit bald zu Schanden gemacht werden, sintemal wir auch nicht einen Augenblick unsere Pflicht thun. Darum ist die Glosse Lyra's falsch und sehr unrecht, die er von den Juden genommen hat, die da lügen, daß es Jakob dafür nicht habe halten können, daß die Verheißung der Gnade einfach und ohne Bedingung wäre, und daß sein Gewissen erschrocken gewesen sei, darum daß er sich mit der Abgötterei seines Schwiegervaters bekleidet hätte. Über dies wird nicht recht gesagt von dem heiligen Patriarchen Jakob; denn er wird ohne allen Zweifel vor der Abgötterei geflohen sein, und wird sie verabscheuet haben, so viel es immer ist möglich gewesen.

34. Und die Heiligen sind zwar nicht ohne Fleisch, aber sie zwingen die fleischlichen Begierden, daß sie nicht herausfahren und das Werk begehen; wie St. Paulus Gal. 5, 24. sagt: „Welche Christo angehören, die kreuzigen ihr Fleisch sammt den Lüsten und Begierden“; denn sie werden getödtet durch den Geist, der über sie herrscht. Es ist wohl Lust und Begierde des Fleisches da, die nicht aufhört zu murren; sie fährt aber nicht heraus in das Werk. Die Heiligen streiten nicht nach dem Fleische.

35. Dies ist aber eine Verheißung der Gnade, wenn Gott also sagt: Du hast nichts gethan, hast nichts verdient; aber dies will ich dir thun und schenken allein aus lauter Gnade und Barmherzigkeit. Solche Verheißungen sind ohne alles Verdienst. Und denselbigen sind die Verheißungen der Patriarchen Abraham, Isaak und Jakob gleich gewesen, wie sie droben, Cap. 27. V. 28. ff., sind erzählt worden: Die Kinder deiner Mutter werden sich vor dir neigen; item: Ich habe dich mit Wein und Oel gestärket. Dazu kommt keine Bedingung: Wenn du dies thun wirst, so sollst du gesegnet werden; sondern: Du hast diese Verheißung und Segen umsonst. Mose ist zwar voll solcher Verheißungen, wie im Gesetze die Verheißungen sind, die eine sonderliche Bedingung haben; aber die Patriarchen haben einfache Verheißungen, die ohne Verdienst sind. Darum ist Jakob um solcher Verheißung willen getrost und fröhlich, er freut sich des sehr großen reichen Trostes halben und daß er den stärksten Schutz hat von den heiligen Engeln. Wie denn David auch also gesinnt gewesen ist, da er im 30. Psalm V. 7. sagt: „Ich werde nimmermehr darnieder liegen.“ Gott lacht, die Engel springen, frohlocken und sind fröhlich, der Geist freuet sich im Herrn. Darum fürchtet sich Jakob nicht, sondern glaubt, daß in so langer Zeit sich der Zorn seines Bruders etwas gestillt habe, vornehmlich, da ihm inzwischen an seiner Gewalt und Reichthum nichts abgegangen war und er daheim im Vaterland geblieben, daß ob er außerhalb Landes im Elende einen armseligen Dienst gehabt und denselbigen ausgewartet, und dazu noch jetzt in der Irre ziehe und keine gewisse Wohnung habe. Denn Esau hatte alles Gut und das ganze Land seines Vaters inne. Darum ist Jakob ohne Sorge, und ist gewiß, daß sein Bruder mit ihm werde versöhnt sein. Und darum läßt er ihn durch die Boten lieblich und freundlich grüßen, ist unerschrocken und hat gar keinen bösen Argwohn auf ihn.

36. Die Form und Weise des Grußes ist nach Art und Gewohnheit der hebräischen Sprache gestellt, da sich Einer des Andern Knecht nennt, der höher oder größer ist als er, damit seine Ehrerbietung anzuzeigen; wie auch die Italiener sehr höflich grüßen. Es sind aber im Gruße Jakobs nicht Worte eines, der sich fürchtet, sondern sind Worte der Ehrerbietung und Freundlichkeit.

37. Und wiewohl Esau daran gedacht hat, daß ihm der Segen genommen war, so wird ihm doch zugleich auch eingefallen sein, daß er gedacht hat: Was ist es denn, das er genommen hat? Ich gebrauche ja aller Güter, und alles, was in meines Vaters Hause ist; er aber ist ein Fremdling und in Elend; er hat um zwei Weiber einen gottlosen und abgöttischen Manne dienen müssen; ich bin in diesem Lande ein Herr, ich habe Nachbarn, Verwandte und sonst andere Freunde, die sehr gewaltig sind. Jakob ist ein armer elender Mensch; er mag den Segen immerhin haben, dieweil er ihm Anderes nichts gibt oder nützt, denn daß er im Elende Armuth leiden muß und einen schweren verachteten Dienst verwalten; ich bin reich und gewaltig, mir mangelt gar nichts.

\quellseite{198}

38. Also pflegt sich das Fleisch zu freuen, wo der Nächste Schaden und Elend leidet. Wohlan, was hat mir denn mein Bruder genommen? Es hat mir es wahrlich GOtt herrlich erstattet, so er mir etwas genommen hat. Ja, es scheint, als ob die Ordnung gar umgekehrt sei, also daß meines Bruders nichts geachtet wird und ihm die göttliche Gnade entzogen worden ist; nach mir aber fragt GOtt viel, macht mich immer größer und herrlicher, und mehret mir dazu mein Gut. Und eben dasselbige hat Jakob auch mögen einfallen, daß er daran auch gedacht hat, als der sehr wohl gewußt hat, wie stolz, ehrgeizig und sicher sein Bruder war; und darum läßt er ihn auch mit solcher Ehrerbietung und ganz demüthig grüßen.

\section[Zweiter Theil]{Zweiter Theil. Was die Boten dem Jakob für eine Botschaft gebracht, und wie sich Jakob dabei verhalten}

\begin{quote}
\textbf{V. 6, 7, 8. Die Boten kommen wieder zu Jakob und sprachen: Wir kamen zu deinem Bruder Esau; und er zeucht dir auch entgegen mit vier hundert Mann. Da fürchtete sich Jakob sehr, und ihm ward bange; und theilete das Volk, das bei ihm war, und die Schafe, und die Kinder, und die Kameele in zwei Heere, und sprach: So Esau kommt auf das eine Heer, und schlägt es, so wird das übrige entrinnen.}
\end{quote}

39. Bisher hat Mose erzählt, wie der Patriarch Jakob, da er wieder aus Mesopotamien gezogen, sehr fröhlich und freudig gewesen und gleichsam triumphirt hat, daß er von der so schweren Anfechtung, so er bei Laban gehabt, erlöset ward, da ihm auch zugleich die Heere GOttes begegnet waren, ihm seine Fröhlichkeit zu mehren, und daß er möchte ohne Sorge sein und ein gewisses Vertrauen haben, daß ihn GOtt schützen würde. Jetzt kommt aber nun vor die schöne klare Sonne eine finstere Wolke und Ungewitter, und Jakob fällt gleichsam vom Himmel wiederum in die Hölle.

40. Mose macht es aber fast lang, und braucht viel Worte, da er diese Anfechtung beschreibt. Denn es ist ein merkwürdiges Exempel, dergleichen uns in der heiligen Schrift viele vorgehalten werden, nämlich, wie GOtt seine Heiligen pflege zu versuchen und ihren Glauben mit mancherlei Anfechtung und Trost zu prüfen. Wie im 94. Psalm V. 19. steht: „Ich hatte viel Bekümmernisse in meinem Herzen“ etc. Also ist dies Leben auf Erden voller Jammer und Trübsal, und ist darum das künftige Leben verheissen, auf daß wir durch Geduld und Trost der Schrift Hoffnung haben, Röm. 15, 4. Unterdeß werden diese zwei Dinge untereinander gemengt, nämlich, die Geduld oder der Trost und die Trübsal. Denn wenn der Trost vorhanden ist, so ist allenthalben Glück und Freude. Wiederum aber, wenn die Finsterniß und Trübsal herein fallen, so regiert der Teufel und Christus wird gekreuziget.

41. Darum werden wir durch solche Exempel unterrichtet, daß wir nicht verzagen, sondern uns des Aeußersten versehen und auf gewisse und reiche Erlösung warten und hoffen sollen, das Unglück werde gewiß ein Ende nehmen, wiewohl man im Unglück und Trübsal selbst kein Ende sieht, noch daß GOtt daselbst gegenwärtig sein sollte; sondern die Gegenwart GOttes ist vor dem Fleische verborgen, daß es dieselbige nicht verstehen kann.

42. Darum muß man das Wort lernen und sich darin üben, und wir sollen wissen, daß dieser Wechsel oder Veränderung im Leben der Heiligen und aller Gläubigen, die GOtt gefallen wollen, immer für und für also währet. Wie St. Paulus 2 Tim. 3, 12. sagt: „Alle, die gottselig leben wollen in Christo Jesu, müssen Verfolgung leiden.“ Und in den Geschichten der Apostel am 14. Cap. V. 22.: „Wir müssen durch viel Trübsal in das Reich GOttes gehen.“ „Aber GOtt ist getreu, der uns nicht lässet versuchen über unser Vermögen“, 1 Cor. 10, 13, wiewohl sich vor der Vernunft das Gegentheil sehen läßt, und dünket uns alle Anfechtung allzu viel, groß und unleidlich zu sein. Wie St. Paulus 2 Cor. 1, 8. 9. von sich selber sagt: „Wir waren über die Maßen beschweret und über Macht, also, daß wir uns auch des Lebens erwegen, und bei uns beschlossen hatten, wir müßten sterben.“ Du darfst dich aber davor nicht fürchten, daß dich GOtt über die Maßen beschweren sollte; denn der Spruch steht fest: „GOtt ist getreu“ etc., 1 Cor. 10, 13.

43. Dieses lehrt oftmals in diesen Historien wieder, darum muß man es auch oft wiederholen, und ist rechter Ernst und ein hohes Ding: es sind nicht eitel vergebliche Worte, sondern es wird hier gehandelt vom Tode oder vom ewigen Leben. Und ist niemand, der diese Weisheit also vollkommenlich erschöpft hätte, daß er keiner Uebungen oder Lehre mehr bedürfte. Denn man muß allezeit lernen, und in der Lehre wachsen und zunehmen; und diese Weisheit muß man eigentlich aus GOttes Wort lernen, sie kann mit den Augen oder fleischlichem Verstand nicht begriffen werden.

44. Nachdem wir nun den großen Triumph und die Herrlichkeit Jakobs gesehen haben nebst dem Heer GOttes und himmlischen Heerscharen, so wollen wir nun mit ihm auch in die Hölle fahren und seine Traurigkeit und Schrecken sehen. Die Boten kommen wieder und verkündigen Jakob, was sie gesehen haben. Wir bringen, sagen sie, eine traurige Nachricht: „Dein Bruder zeucht dir entgegen mit vier hundert Mann.“

45. Da wird Jakob sehr erschreckt, wiewohl keine Gefahr vorhanden ist; wie wir hernach hören werden, wie er aus dieser Angst und Noth entronnen sei. Es wäre zwar die Gefahr groß gewesen, wenn GOtt verhängt hätte, daß die vierhundert Mann Jakob mit seiner Heerde und Gefährten, die um ihn waren, hätten mit Gewalt angegriffen; und dies allein ist Jakob in dem Augenblick also vor Augen gestanden und hat ihm in seinem Herzen gelegen, daß ihm die bewaffneten Männer könnten Schaden thun, und daß sie daher kämen, ihn mit Weib und Kindern zu erwürgen. Aber GOtt gedenkt viel ein Anderes, und wird mit Jakob ein ander Ende nehmen, welches er zuvor nicht sieht.

46. Das Fleisch kann es nicht lassen, es muß zittern und erschrecken. Denn Esau war sich wohl bewußt, daß zwischen ihm und seinem Bruder ein Unwille war; darum hat er daran nicht gezweifelt, es würde ihm gewißlich der Untergang bevorstehen. Denn er wird gedacht haben: Was mag es bedeuten, daß mein Bruder mit so viel gerüsteten Männern daherziehet? Warum läßt er mich nicht irgend mit einem Wort grüßen oder vollkommen heißen durch die Boten? Also deutet er es zum Aergern, wie das Fleisch immerzu zu thun pflegt, das bald alle Hoffnung fallen läßt, und wird erschrocken und kann nicht anders gesinnet sein. Denn St. Paulus sagt von sich selber Röm. 7, 23.: „Ich sehe ein ander Gesetz in meinen Gliedern, das da widerstreitet dem Gesetz in meinem Gemüthe.“ Das Gesetz im Gemüthe sagt also: Hoffe und vertraue, zweifele nicht, fürchte dich nicht, GOtt lebt und regieret, die Engel wachen über dir. Ich höre aber ein ander Gesetz in meinen Gliedern, das da sagt: Es ist alles gar verloren, die Sache ist verdorben, es ist aus mit dir. Das ist das Gesetz in den Gliedern; das Fleisch ist ungläubig und will immer zagen, und kann sich selber nicht verheißen oder glauben, daß es sollte erlöset werden. Eben dasselbige ist demnach dem Patriarchen Jakob auch begegnet: der zittert und wird erschrocken, da er hört, daß sein Bruder mit so großer Menge Volks daherziehet, und wider ihn und sein armes, elendes, wehrloses Häuflein seiner Weiber, Hirten und kleiner Kindlein streiten werde, und daß er nicht allein schnauben solle mit Streiten, sondern auch mit Morden und Würgen; wie er hernach sagt: „Daß er nicht komme und schlage mich, die Mütter sammt den Kindern.“

47. Wie reimt es sich aber, daß ein solcher heiliger Patriarch also zittern soll, daß er von solchem großen Triumph und Fröhlichkeit, nach welcher er sich so hoch gefreut hat, in eine solche tiefe Verzweiflung fallen soll? Ist das recht, daß man also wanken und zweifeln soll? Ja, wären wir da gewesen, wir hätten es wohl ärger gemacht. Dies sind Exempel, die uns lehren und unterrichten; wie Petrus, da ihn der Herr hieß auf dem Wasser zu ihm kommen, Matth. 14, 28. ff., eilt er bald zu Christo, triumphirt und ist sehr kühn als ein Herr des Meeres und der Wellen, ja, der Teufel und der ganzen Welt; aber bald, da er einen starken Wind gesehen hat, fängt er an zu zweifeln. Darum straft ihn auch Christus und spricht: „O du Kleingläubiger, warum zweifeltest du?“ Ich habe dich heißen zu mir kommen und du kommst daher; dieweil du aber empfindest, daß der Wind dem Wort und der Erfahrung zuwider ist, da wanket dein Glaube und hebst an zu sinken.

\quellseite{199}

48. Auf dieselbe Weise läßt sich der Mensch dünken, wenn er ohne Anfechtung ist, daß er sein gesund und stark genug sei, er darf den Teufel verachten und mit Füßen treten, und kann andere Menschen auch stärken; wie im 30. Psalm V. 7. steht: „Ich aber sprach, da mir’s wohl ging: Ich werde nimmermehr darnieder liegen“, ich gebe um alle Teufel und um den Tod nicht ein Klipplein. Wenn aber ein Ungewitter und ein starker Wind daher geht, oder daß irgend ein ungewöhnlich Gespenst sich zeigt, da schreien wir: Ach weh, es ist alles aus, der Tod will mich fressen! und, wie im 30. Psalm V. 8. folgt: „Da du dein Antlitz verbargest, erschrak ich.“

49. Also ist es Mose auch gegangen bei dem Rothen Meer. Die Kinder Israel waren mit großen Freuden ausgezogen. Egypten war gedemüthiget und gestraft; sie hatten den Raub mitgenommen und davon gebracht; sie riefen alle: „Der HErr, unser GOtt, ist mit uns“ etc., da ist es eitel köstlich Ding. Da sie aber sehen, daß ihnen Pharao mit Heereskraft nachjagt, und daß sie zu beiden Seiten von Bergen umgeben waren, da sagen sie: „Waren nicht Gräber in Egypten?“ 2 Mos. 14, 11. Nach solchem großen Triumph und fröhlichen folgt der Schrecken des Todes und der Hölle. Ja, so stehe Mose selbst auch an, zu dem der HErr 2 Mos. 14, 15. sagt: „Was schreist du zu mir?“ Es wollte ihm das Herz auch entfallen und in tausend Stücke zerspringen, dieweil er das Volk Israel aus Egypten geführt hatte, und da sie noch kaum herausgekommen sind, sieht er, daß sie greulich auf die Fleischbank müssen geopfert werden. Da hat er gedacht: Nun bin ich doch eine Ursache, daß so viele Leute müssen erwürget werden. Was habe ich nun gethan? Ich habe den egyptischen König erzürnt und ihn an das Volk gehetzt. Darum schreit er so heftig und sein Herz ist erschrocken. Aber GOtt antwortet ihm auf sein Schreien, und spricht: „Was schreist du zu mir?“ Ich will den Tyrannen wohl in Schranken halten und zwingen. Ja, wir sehen aber das Meer vor uns, zu beiden Seiten sind wir von Bergen umringt, wie können wir entrinnen? „Heb deinen Stab auf“, spricht der HErr 2 Mos. 14, 16., und recke deine Hand über das Meer, und theile es von einander, daß die Kinder Israel hinein gehen, mitten hindurch auf dem Trockenen.“ Ja, das kann unser HErr GOtt.

50. Und dieses soll die Kirche und Gemeinde GOttes lehren und stärken in aller Noth und Widerwärtigkeit. Denn darum wird es uns vorgehalten, daß wir alsdann glauben und GOtt vertrauen sollen, wenn es sich ansehen läßt, daß Alles unmöglich und verloren, und daß weder Rath noch Hülfe mehr vorhanden sei. Wie im 107. Psalm V. 27. steht: „Sie wußten keinen Rath mehr.“ Da sollst du lernen also schließen: Wiewohl Alles mit großer Schwachheit zugeht, und es jetzt scheint, als ob Alles verloren sei, so ist doch GOtt nichts unmöglich. Denn also sagt der Engel zu der Jungfrau Maria Luc. 1, 37.: „Bei GOtt ist kein Ding unmöglich.“ Denn er hat einmal aus Nichts Alles geschaffen: die Macht und Gewalt hat er noch, und auf dieselbige Weise erhält und regiert er noch alle Dinge. Denn das bei uns Nichts ist, das ist bei GOtt Alles; was uns unmöglich, das ist ihm sehr leicht. Er wird auch am jüngsten Tage mit eben solcher Kraft die Todten auferwecken, die vom Anfang der Welt in dem Staub der Erde gelegen haben.

51. Dies sei nun gesagt von dem Schrecken Jakobs in der sehr großen Noth und Gefahr, welcher Schrecken ein Anzeichen ist, daß auch in den Heiligen das Fleisch noch kräftig und lebendig ist. Denn um solcher Ursache willen ist dieser sehr heilige Mann nach solchen großen Verheißungen und starkem Trost schier in Verzweiflung gefallen.

52. Nun sehe auch, wie es um das ganze Gesinde und die Kirche oder Gemeinde, so Jakob in seinem Haus gehabt, bestellt sei. Denn da sie sehen, daß das Haupt und der Hausvater selbst wankt und zittert, was sollten denn sie nicht thun? Der Bischof, der Lehrer und Tröster der Kirche und Gemeinde wankt und fürchtet sich: was sollen denn die für Hoffnung haben, die an solchem ihrem Lehrer und Hauptmanne hangen?

Derhalben wird unter den Weibern Jakobs und seinen Hirten oder Verwandten ein jämmerlich Trauern gewesen sein. Denn er hat auf dieser Reise eine große Menge Volkes bei sich gehabt: um der großen Heerde willen, dieselbige zu treiben, hat er bei hundert oder zweihundert Hirten haben müssen; dieselbigen sind auch erschrocken und ungeduldig geworden, und werden vielleicht gesagt haben: Was soll es sein, daß er uns aus Mesopotamien geführt hat? Wollte GOtt, wir wären noch bei Laban, da würden wir sicher und wohl leben, da dürften wir uns vor diesen gerüsteten Männern nicht fürchten. Solche Klagen pflegen ein erschrockenes und betrübtes Herz mehr zu plagen und zu martern, denn die Noth und Gefahr an sich. Wiewohl sich nun Jakob seiner Person halben nicht gefürchtet hatte, so ist ihm doch sein Herz weich und verzagt geworden von solchem Geschrei, Heulen und Weinen der Weiber, und der Schrecken des Fleisches ist dadurch immer gemehret worden, wie man hernach sehen wird.

53. Und das ist es, das Mose sagt, daß sich Jakob nicht allein sehr gefürchtet habe, sondern daß ihm auch angst und bange geworden sei; denn er beschreibt damit einen sehr schwachen Glauben. Das hebräische Wort zur heißt auf Lateinisch angustia, Angst, und Mose will sagen: Es ist ihm bange gewesen, sein Herz hat keinen Rath gewußt, und man hat keinerlei Weise oder Wege sehen können, wie er hätte entrinnen mögen, oder wie ihm hätte mögen geholfen werden; welches die Schrift an andern Orten auf Lateinisch pfleget latitudo zu nennen, als, im 4. Psalm V. 2.: In tribulatione dilatasti mihi: „Der du mich tröstest in Angst.“ Denn die Schrift setzt die hebräische Worte gegen einander, latitudo und angustia, das ist, den Schrecken des Herzens und den Trost. Wie Jakob, da er in der höchsten Angst war, erschrocken gewesen ist, darum daß er keinen Rath mehr gewußt hat, wie der 107. Psalm V. 27. sagt. Darum ist dieses eine ganz jämmerliche Geschichte, da dieser Siegesmann und Überwinder also abgemalt wird, daß er gleichsam in die Hölle gefallen ist, und so viel herrlicher Verheißungen und Trostes, die vorhergegangen sind, gar vergißt.

54. Aber dieses wird uns darum vorbehalten, daß wir unser Vertrauen nicht auf uns selbst stellen, wie St. Paulus sagt 2 Cor. 1, 9. Denn Paulus hatte auch mancherlei Gefahr und Noth erfahren, war auch auf mancherlei Weise daraus errettet worden; bisweilen aber erzeigt er sich, daß er so einen großen starken Muth habe und sich vor nichts fürchte, daß es sich mit ihm ansehen läßt, daß er auch alle Gefahr und Angst ohne alles Fürchten verachten könne; wie er sagt 2 Cor. 6, 11. 12.: „O ihr Corinther, unser Mund hat sich zu euch aufgethan; unser Herz ist getrost. Unserthalben dürft ihr euch nicht ängsten. Daß ihr euch aber ängstet, das thut ihr aus herzlicher Meinung“ etc. Da ist sein Herz voll Freude, und will haben, daß sie sich alle mit ihm freuen und getrost sein sollen. Alles, was im Himmel und auf Erden ist, lacht ihn an. Wiederum aber sagt er 2 Cor. 1, 8. 9.: „Wir waren über die Maaße beschweret, und über Macht, also, daß wir uns auch des Lebens erwegten;“ item: „Wir hatten bei uns beschlossen, wir müßten sterben. Das geschah aber“, sagt er, „darum daß wir unser Vertrauen nicht auf uns selbst setzten, sondern auf GOtt, der die Todten auferwecket.“ Was ist das, lieber Paulus? Warum bist du nicht fröhlich und getrost? Warum machst du die Andern nicht auch fröhlich? Soll denn Paulus, ein solcher großer Apostel, dermaßen gedemüthiget werden, daß er lieber sterben wollte denn leben? der eben voll Heiligen Geistes war, mit dem läßt es sich nun ansehen, daß er gar ohne Geist sei.

55. Ja, wir haben das erste und vornehmste Bild desselbigen Kampfes vorgestellt an Christo, der im Garten Blut geschwitzet hat. Und ist zwar um das Leben aller Heiligen also bestellt, daß die Exempel uns darum vorbehalten werden, daß wir daraus gelehret werden und lernen sollen, daß wir unser Vertrauen nicht auf uns selbst stellen, sondern auf GOtt, denn das ist die Kunst der Gottseligen, daß sie sich selbst mißtrauen und GOtt vertrauen, und einem solchen GOtt, der die Todten auferweckt, und „rufet dem, das nicht ist, daß es sei“, Röm. 4, 17. Denn das sind GOttes Werke. Wenn du nun willst gottselig sein, und GOtt für deinen Vater erkennen und ehren, so bedenke, daß du mit den Teufeln, mit der Welt, mit dem Fleisch, mit der Sünde und Gesetz wirst streiten müssen. Und unter diesen allen ist die Welt der geringste Feind; wiewohl zu der Verfolgung der Welt allezeit auch der grimmige Zorn und die Wuth des Teufels kommt, der der Welt hilft und die Gottseligen schreckt. Denn wir leiden auch nicht die allergeringste Verfolgung, es sei denn der Teufel auch mit dabei.

\quellseite{200}

56. Derhalben sollen wir bedenken, daß wir uns das vornehmlich lassen angelegen sein und dasselbige auch lernen, daß wir in aller Noth und Anfechtung Gott vertrauen, deß eigen Werk ist, daß er die Todten auferweckt, daß er die, so sehr erschrocken sind, tröstet und zufriedenstellt, daß er aus höchst elenden Menschen höchst selige Menschen, und aus denen, die schier verzweifeln wollen, fröhliche Leute macht. Darum gebraucht Mose nicht vergeblich so viel Worte und beschreibt nicht umsonst so genau diese Anfechtung. Denn er will haben, daß wir dieses Exempel fleißig in unser Herz einprägen sollen.

57. Was thut aber Jakob, da sein Glaube also angefochten und schwach wird, da das Fleisch herrscht in seinem Gesetz, das da zweifelt, verzagt und lästert? Antwort: Er thut erstlich, was an ihm ist, wie gemeiniglich die Papisten pflegen zu reden, aber in einem andern Verstande. Denn sie verstehen diese Art zu reden nicht recht, wissen sie auch nicht recht zu gebrauchen. Denn wir haben droben gesagt, wiewohl wir die Verheißung und Gebote Gottes haben, so soll man doch gleichwohl Gott nicht versuchen, das ist, man soll die Mittel nicht versäumen oder verachten, sondern man soll diejenigen gebrauchen, die man gebrauchen kann; sintemal Gott die Vernunft, und auch den Rath und Hülfe der Vernunft nicht darum gegeben hat, daß du sie verachten sollst. Dieses thun diejenigen, die vermessen sind, oder die da verzweifeln, wenn sie also sagen: Ich thue gleich, was ich wolle, so kann ich doch das nicht verhindern, das nothwendig und durch ein Schicksal geschehen soll und muß. Das sind Worte derer, die da verzweifeln. Die aber vermessen sind, pflegen solche Worte zu führen: Wenn ich schon nicht esse, so werde ich doch wohl leben, wenn ich leben soll; Gott hat mir das Leben verheißen, darum ist daran nichts gelegen, ich esse oder esse nicht. Wir sollen uns aber wohl vorsehen, wenn wir göttliche Verheißung haben, daß wir Gott nicht versuchen durch Vermessenheit und daß wir auch mit Verzweiflung nicht sündigen. Wenn du eine Leiter hast, ist es nicht nöthig, daß du dich woltest zum Fenster heraus stürzen; darfst auch nicht mitten durch die Elbe gehen, dieweil du eine Brücke hast; sondern ein Jeder soll das thun, was ihn die Vernunft lehret oder weiset, und soll das Andere Gott befehlen, der wird wohl geben, daß es einen guten Ausgang nehme.

58. Desselben ist dies ein sehr schön Exempel. Denn Jakob unterläßt der Dinge keines, das er thun kann. Er theilt das Volk, das bei ihm war, dazu sendet er seinem Bruder Geschenke; er zieht selbst vorher, und thut alles, was ihm zu thun möglich ist. Ein Anderer, der da verzagt wäre, würde gesagt haben: Du willst gar nichts thun; wo es also versehen ist, daß du sollst erschlagen werden, so wirst du dich davor nicht hüten können. Wie die Andern aus der Verheißung auch einen solchen Schluß ziehen, daß sie sagen: Wenn ich versehen bin, so werde ich nicht können verloren werden, ich thue gleich, was ich wolle. Dies sind teuflische Worte, davor soll man sich hüten. Das ist wohl wahr: was zuvor versehen ist, wird also geschehen; man muß aber das hinzusetzen, daß dir solches unbekannt sei. Wie du nicht wissen kannst, ob du morgen sterben oder leben werdest; und Gott hat nicht haben wollen, daß du es wissen solltest.

59. Derhalben ist es ja ein närrisch Ding, daß du darnach forschen willst, daß Gott aus sonderlichem Rath vor dir verborgen hat. Dieweil du aber nicht weißt, wie lange du leben werdest, darum sollst du der Mittel gebrauchen, die zum Leben nöthig sind. Wo es also versehen ist, daß du nach Ausgang eines Monats sterben sollst, so sollst du gleichwohl Gott nicht versuchen; sintemal du dessen noch nicht gewiß bist, sondern du sollst die Mittel gebrauchen, die dazu vonnöthen sind, daß man dies Leben erhalten möge.

60. Derhalben soll man die Thorheit oder vielmehr die Gottlosigkeit derjenigen strafen, die Alles auf die Versehung ziehen wollen, welche in sich begreift den heimlichen Rath und Regierung Gottes, die uns unbekannt ist. Wenn du Gott wärst, alsdann möchte es zu leiden sein, daß du dich darnach richtetest. Dieweil aber Gott haben will, daß du davon nicht wissen sollst, darum denke darauf, daß du in deinem Berufe bleibest, und nicht weiter gehest, denn dich das Wort weiset, und gebrauche die Mittel und den Rath, so Gott verordnet hat. Ich kann zwar nicht zuvor sehen, was für Frucht aus meiner Lehre folgen wird, welche Menschen bekehrt und welche nicht bekehrt werden sollen. Wenn ich nun sagen wollte: Die da sollen bekehrt werden, die werden wohl ohne mich bekehrt werden; wo aber Etliche nicht sollen bekehrt werden, was ist es denn vonnöthen, daß ich vergeblich arbeite? Wenn ich also sagen wollte, das wäre gar ein ungereimt und gottlos Ding. Denn wer bist du, daß du nach solchem Ding forschen willst? Thue, was dir in deinem Amt zu thun gebühret: wie es aber hinaus gehen und gerathen werde, das sollst du Gott befehlen. Es gefährhet dir nicht, daß du sagst: Wenn dies geschehen soll, so wird es geschehen. Gott hat in seinem Wort durch den heiligen Paulus befohlen: Gehe hin und richte dein Amt aus, dazu du berufen bist: „Predige das Wort, halt an, es sei zu rechter Zeit, oder zur Unzeit; strafe, dräue, ermahne mit aller Geduld und Lehre“, 2 Tim. 4, 2. Diesem Befehl soll man von Nöthen wegen gehorchen und den unnützen Vorwitz fahren lassen.

61. Dieses soll man fleißig lehren, die gottlose und türkische Weisheit, dazu heutiges Tages viele Lust haben, damit zu widerlegen. Denn die Türken sind also gesinnet, wenn sie in den Krieg ziehen. Wo ich sterben soll, sagen sie, so werde ich sterben; wo aber nicht, so wird mir kein Geschoß schaden können. Darum greifen sie den Feind an, als wären sie unsinnig, mit großer Gefahr, wenn sie sich auch vor solcher Gefahr wohl hüten könnten. Aber Gott will haben, daß wir der Gelegenheit und Mittel, die uns gewiesen sind, die Gefahr zu meiden, gebrauchen sollen.

62. Wie Jakob mit höchstem Fleiß alles thut, was er nur für Rath hat finden können, wiewohl ihm damit, was er erdenkt und thut, nicht geholfen ist, wird auch dadurch nicht errettet. Denn wo Gott das Herz Esaus nicht geändert hätte, so würde Jakob mit seinem Fleiß nichts ausgerichtet haben. Doch versucht er in solcher großen Angst und Noth alles, was er kann und mag. Er hört nicht auf, für seine Familie fleißig zu sorgen, wie er sie schützen möchte; er ist mit Händen und Füßen nicht stille; sondern was zur Hand liegt, das thut er, auf daß er Gott nicht versuche. Er theilt das Volk in zwei Heere, desgleichen auch die Heerde, als, Schafe, Rinder und Kameele. Das hebräische Wort machanoth heißt hier im Femininum „Haufe, Schaar“.

63. Was ist aber das für ein Rath, daß er sagt: „So Esau kommt auf das eine Heer, und schlägt es, so wird das übrige entrinnen“? Das ist nur ein menschlicher Rath und den die Vernunft und nicht der heilige Geist erfunden hat. Denn der heilige Geist hat dazumal etwas Anderes zu thun gehabt, nämlich, daß das Herz Esaus möchte gelindert und erweicht werden. Doch hat gleichwohl der heilige Geist dieses auch regiert, also daß Jakob Gott nicht versucht hat. Denn er hat gedacht: Wenn mein Bruder an die Heerden kommen wird, so wird er sie alle mit den Hirten und mit der ganzen Familie im ersten Sturm und grimmigen Zorn erwürgen; derhalben will ich sie in zwei Heere theilen, daß so er das eine Heer schlagen wird, er sich an solcher Rache genügen lasse und des andern Heeres verschone, und aufhöre zu toben und zu wüthen.

64. Denn es kann wohl geschehen, daß ein menschlich Herz gesättigt wird, wenn von einem Theil so Viele erschlagen sind. Wie oftmals auch die grausamsten Sieger nach dem Siege dem Kriegsvolk geboten haben, nicht mehr Grausamkeit mit Würgen und Morden zu üben. Da die Schlacht bei Canna geschehen war, und da Hannibal gesehen hatte, daß die Römer geschlagen waren, hat er dem Kriegsvolk geboten, daß sie sollten aufhören zu würgen. Dieses ist nun menschlich Ding, daß Jakob in seinem Herzen also schließt, daß der Feind, wie grausam er auch sei, damit gleichwohl gesättigt werden könnte, wenn er sieht, daß der halbe Theil Volks niedergelegt und geschlagen ist. Also hütet er sich davor, daß er Gott nicht versucht, und setzt doch den halben Theil der Niederlage und dem Raube aus; den andern Theil aber hofft er zu erhalten.

65. Daß er sich aber solch greulich Ding von dem grimmigen Zorn seines Bruders einbildet, das kommt vom Fleisch und vom Teufel her. Gott läßt Jakob von Esau einen solchen Argwohn schöpfen; er selbst aber gedenkt unterdeß viel ein Anderes. So ist nun dies der erste und ein schöner Rath, so viel menschliche Weisheit
\quellseite{201}
belanget und wie die Vernunft davon urtheilt, damit Jakob versucht, den Zorn und Grimm seines Bruders zu lindern.

66. Das Andere ist, daß er sich zum Gebet wendet, und damit hütet er sich auch, daß er GOtt nicht versuche. Denn wenn wohl der Glaube schwach ist, so scheint er doch daselbst hervor, da er schier von Angst und Schrecken gar unterdrückt ist, da sich Jakob fürchtet, er möchte mit dem ganzen Volk, so bei ihm war, und auch mit der Heerde erschlagen werden. Und das ist das rechte Werk, so in der Noth und Verfolgung vornehmlich vonnöthen ist. Er nimmt sich nichts Wundersames vor und daß er mit seinem Schwert über sein Vermögen etwas thun wollte, und er versucht auch GOtt nicht, daß er sollte verwegen und allen Rath der Vernunft verwerfen: sondern er fängt an, sich mit dem Glauben aufzurichten und zum Gebet zu halten, und liegt doch noch immer im Kampf.

67. Wir haben aber nun oftmals gesagt, daß der Glaube alsdann am stärksten sei, wenn er am schwächsten ist. So gar wunderlich sind die Werke GOttes. Jesaia Cap. 60, 22. spricht: „Aus dem Kleinsten sollen tausend werden, und aus dem Geringsten ein mächtig Volk.“ Item, Paulus spricht 2 Cor. 12, 9. 10.: „Ich will mich am allerliebsten rühmen meiner Schwachheit; denn wenn ich schwach bin, so bin ich stark.“ Also ist Jakob hier auch sehr schwach in seinem Schrecken und Angst, ist über die Maaßen sehr betrübt und kleinmüthig, und ist doch sein Glaube nie stärker gewesen. Denn der Glaube, der da wider den Unglauben streitet, der setzt auf das allererste.

68. Das selbige Seufzen aber versteht niemand, weder Jakob, noch sonst jemand anders. Und das ist das unaussprechliche Seufzen, davon St. Paulus sagt Röm. 8, 26.: „Der Geist selbst vertritt uns aufs beste mit unaussprechlichem Seufzen.“ Es ist keine Stimme Eines, der da fröhlich ist und triumphirt, sondern ist in Wahrheit ein solch Seufzen, da einem betrübten Herzen dünkt, daß es nur gar kleinlaut sei und schwerlich Odem holen könne. Der Prophet Jesaia nennt es Cap. 42, 3. ein glimmendes Docht und ein zerstoßen Rohr. Dieweil aber nichts schwächer ist denn solch Seufzen (denn es ist gleichsam der letzte Odem), so ist es doch gleichwohl ein unaussprechlich Seufzen.

69. Dieses ist stracks wider einander, oder, opposita in adjecto, und ist doch gleichwohl wahr, und ist kein Mensch, der es mit Worten erreichen oder ausdrücken könnte, wie kräftig dieses Seufzen sei. Denn meine Kraft spricht GOttes Wort 2 Cor. 12, 9., „ist in den Schwachen mächtig.“ Das Fleisch widerspricht solchem wohl und sagt, es verstehe es nicht, ja, es empfinde und erfahre gar das Gegentheil. Es ist aber gleichwohl gewiß, wenn die Kraft schwach ist, so ist sie vollkommen, und wiederum, die vollkommene Kraft ist lauter Schwachheit. Wer hat doch jemals also geredet? Ja, hüt dich vor dem: Ich kann nicht. Wenn dich der Teufel dahin bringt, daß du nicht kannst, so hast du schon verloren. Siehe aber zu, daß du zu solcher Schwachheit die Verzweiflung nicht kommen lässest; denn es ist ein unaussprechlich Seufzen. Also redet die Schrift. Wenn die Verzweiflung nicht dazu kommt, sondern der ist ein Seufzen, so ist es die vollkommenste Kraft in der allerschwächsten Schwachheit.

70. Also spricht der Herr zu Mose im andern Buch am 14. Cap. V. 15.: „Was schreiest du zu mir?“ Nun schriee aber Mose gar nicht, sondern er war stumm vor Furcht; man hat daselbst keinen Laut irgend einer Stimme gehöret; er war, als wäre er todt. Der Heilige Geist aber, der dasselbe Seufzen verstehet, der sagt, daß es ein groß Geschrei sei. Wie schrieest er aber denn? Antwort: Abba Pater: Ach du lieber GOtt! Da ist die Verzweiflung noch nicht vollständig, sondern es ist noch ein klein Vertrauen da, gar ein kleiner Odem, ein kleines Seufzen, das zwar in deinen Augen sehr klein ist, aber in den Ohren Gottes ist es sehr groß, und ist ein Geschrei über alles Geschrei, welches Himmel und Erde erfüllet, daß es GOtt nicht unterlassen kann, er muß solch Geschrei erhören, und darauf antworten: „Was schreiest du?“ Also betet nun Jakob auch mit unaussprechlichem Seufzen, wie im Text folgt.

\begin{quote}
\textbf{V. 9—12. Weiter sprach Jakob: GOtt meines Vaters Abraham, und GOtt meines Vaters Isaak, der du zu mir gesagt hast: Kehre wieder in dein Land, und zu deiner Freundschaft, ich will dir wohlthun; ich bin zu geringe aller Barmherzigkeit und aller Treue, die du an deinem Knechte gethan hast; denn ich hatte nicht mehr, weder diesen Stab, da ich über diesen Jordan ging, und nun bin ich zwei Heere worden. Errette mich nun der Hand meines Bruders, von der Hand Esau; denn ich fürchte mich vor ihm, daß er nicht komme, und schlage mich, die Mutter samt den Kindern. Du hast gesagt: Ich will dir wohlthun, und deinen Samen machen wie den Sand am Meer, den man nicht zählen kann vor der Menge.}
\end{quote}

71. Der schwache und fast verzweifelnde Glaube Jakobs verzweifelt doch gleichwohl noch nicht gar: das Docht glimmt, verlischt aber noch nicht; das Rohr ist zerbrochen, aber nicht verworfen; wie St. Paulus sagt 2 Cor. 4, 8. 9.: „Wir haben allenthalben Trübsal; aber wir ängsten uns nicht. Uns ist bange; aber wir verzagen nicht. Wir leiden Verfolgung; aber wir werden nicht verlassen. Wir werden untergedrückt; aber wir kommen nicht um“ etc. Derselben erholet sich Jakob ein wenig, aber doch mit großem Schrecken.

72. Es ist aber ein trefflich Exempel eines herrlichen Gebets, daß alle Eigenschaften hat, die zum guten Gebet gehören. O wer es thun könnte! Es sind uns im Neuen Testament die Lehre und Verheißungen vom Gebet sehr reichlich vorgehalten; wer kann es aber also halten? Das hat Jakob selbst auch nicht thun können. Er ist aber also erschrocken, daß er erstlich Alles bestellt und verordnet, ehe denn er an das Gebet kommt, so doch das Gebet sollte vorhergehen: Erbarm dich mein, o HErr GOtt etc. Aber das Fleisch, das uns gefangen nimmt in der Sünde Gesetz, macht, daß wir die Ordnung umkehren. Darum ist es ein Gebrechen an ihm, daß er mit dem Gebet länger gewartet hat, denn er billig sollte gethan haben. Es ist aber eine Schwachheit des Fleisches, die ihn erstlich auf den Rath der Vernunft gebracht hat.

73. Er hat aber das Gebet also gestellt, daß er sagt: „GOtt meines Vaters Abraham“ etc. Im Hebräischen heißt es „Götter“. Es ist ein brünstiges Gebet nach dem Fünklein des Glaubens, der so heftig gestritten hat. Es ist eben das Geschrei Moses am Rothen Meer. Erstlich aber ergreift Jakob mit dem Glauben den GOtt Abrahams und Isaaks, und nennt ihn darnach auch seinen HErrn. Also erinnert er GOtt der Verheißungen, so Abraham und Isaak geschehen waren, und vornehmlich die GOtt ihm selbst auch gethan hatte, welches das Gebet sehr stark und brünstig macht.

74. Er sagt aber im Gebet weiter also: „Du hast zu mir gesagt: Kehre wieder in dein Land“ u. s. w. Als wollte er sagen: Du hast mich in diese Noth, in diese Anfechtung und Gefahr meines Lebens gebracht. Das ist aber ein großer Trost, wenn Einer also sagen kann: HErr GOtt, ich bin ja nicht aus eigenem Frevel oder Tollkühnheit in diese Noth gerathen, dazu auch durch keines Weisen oder Narren Rath; sondern du hast es gesagt, du hast mich geheißen: darum habe ich daran recht gethan, daß ich Laban verlassen habe. Und dies ist nun deine Sache; deine Verheißung und Treue wird jetzt angefochten, du wirst deine Treue und Glauben und nicht meinen Glauben allerdings retten müssen.

75. Dieses ist der schwache Glaube und das Seufzen, davon gesagt worden ist. Das Seufzen beweget Himmel und Erde und ist ein sehr angenehmes Gebet. Du hast es gesagt; ich gehe in deinem Gehorsam; ich soll in mein Vaterland ziehen auf deinen Befehl, und wie du zu mir gesagt hast, lieber HErr. Siehe aber, wie so viel große Hindernisse vorfallen. Ich bin in solche Angst und Beschwerung gekommen, daraus ich mich mit meinen Kräften und Rath nicht zu erretten weiß; darum bedarf ich deiner Hülfe. Die Andern werden vielleicht nicht gebetet, sondern ihm vorgeworfen haben: Wären wir bei Laban geblieben, so hätten wir mögen sicher sein. Aber Jakob fragt darnach nichts und ist dies ein sehr starkes Gebet.

76. Was hat er nun aber mehr gesagt? „Ich will dir wohlthun.“ Als wollte er sagen: Das ist ja dein Wort, daß du mich nicht verderben und mir keinen Schaden zufügen, sondern mir wohlthun wollest in meinem Vaterlande, darin du mich hast ziehen heißen. Und daraus sieht man, daß er nicht allein vom Teufel ist angefochten und geplagt worden um der Nachricht wollen, daß sein Bruder ihm entgegen zöge, sondern daß ihn seine Weiber, Kinder, und das ganze Gesinde auch werden geplagt haben. Die werden gerufen haben: Ach, lieber Vater, wo nun hinaus? lieber Vater, lieber Mann, lieber Herr, was hast du gethan? warum hast du uns in solche große Gefahr gebracht? Die Klagen haben ihm die gewaltigen Worte abgedrungen, daß er sagt: Lieber HErr GOtt, höre doch, wie diese Leute zittern und mir mein Herz auch martern, wiewohl ich weiß, daß du mir Hülfe und Trost verheißen hast.

\quellseite{202}

77. Also scheint der Glaube hervor und läßt sich hören, wiewohl er schwach ist. Denn er wird sie ohne Zweifel also angeredet haben: Wie seid ihr doch so gar erschrocken? wollt ihr denn gar verzagen? wir müssen wahrlich die Hoffnung nicht fallen lassen, daß wir nicht errettet werden sollten. Ich will zwar nicht verzagen, wenn ihr schon verzagen würdet. Auf solche Weise hat er sich selbst aufgerichtet, und hat dem Weinen und der Ungeduld der Seinen damit gewehret. GOtt hat mich nicht heißen wieder in mein Land ziehen, daß er mir damit wollte Schaden thun. Es muß noch gut werden: GOtt wird uns helfen, lasset uns nur der Hülfe gebrauchen, die vorhanden ist, und zu GOtt schreien.

78. Ein solches Gebet aber ist GOtt sehr angenehm, welches geschieht in der äußersten Noth und wenn die Gefahr am größten, da fast Alles verloren ist. Das ist das unaussprechliche und sehr gewaltige Seufzen, damit sich die Gottseligen selbst aufwecken wider die Verzweiflung, daß sie Muth fassen und sagen: Nun muß es dennoch nicht so sein, wir werden dennoch nicht verderben: „Ich werde nicht sterben, sondern leben“ etc., wie im 118. Psalm V. 17. steht. Ich habe ja die Verheißung, der HErr hat gesagt, er wolle mir wohl thun; weinet und heulet, und zittert nicht also; GOtt hat uns heißen aus Mesopotamien und wieder in mein Vaterland ziehen.

79. Dieses ist der Gottseligen Kampf, darin sie ihren Glauben gewaltig erwecken mit dem Andenken und Vertrauen auf die Verheißung, und daß sie von GOtt Befehl haben. Ich muß und soll predigen, der Teufel aber wehret; wohlan, so muß gepredigt sein und sollte die Welt zerreißen. Diese sind die, so Gewalt thun und das Himmelreich zu sich reißen, und Matth. 11, 12. Fleischliche Menschen lesen solch Ding schläfrig und verstehen davon gar nichts; denn es sind grobe Leute und haben solche Anfechtung nicht erfahren. Sie wissen nicht, wie dem zu Muthe ist, der in der äußersten Noth liegt. Da bläset der Teufel zu und gibt dem ein, der angefochten und betrübet ist: Es ist verloren, was willst du viel schreien? es ist doch mit dir gar verloren. Aber der Geist spricht wiederum: Ei, es ist nicht verloren, hart dagegen; ich weiß, daß GOtt mit mir ein Anderes im Sinne hat und mir auch viel Anderes verheißen.

80. Das ist eine große Gewalt und Kraft des Geistes in der Schwachheit, und ist GOtt ein sehr angenehmes Opfer, wie im 51. Psalm V. 19. steht, denn es ist ein Opfer der Tödtung: „Die Opfer, die GOtt gefallen, sind ein geängsteter Geist, ein geängstet und zerschlagen Herz wirst du, GOtt, nicht verachten.“ Ja freilich wirst du es nicht verachten; denn es ist ein sehr angenehmes Opfer, das da wohl riecht, welcher Geruch Himmel und Erde füllet; wie im Hohen Liede Salomos am 3. Capitel V. 6. steht: „Wer ist die, die herauf gehet aus der Wüste, wie ein gerader Rauch?“ Es ist zwar ein kleines Räuchlein, aber es riecht hinauf in den Himmel, und macht unserm Herrn GOtt die Nase so wohl, daß er sagt: Höre auf, ich kann nicht mehr hören. Es ist, sage ich, ein klein gering Seufzen; aber es riecht beide mit Geruch und Geschrei vor GOtt hinauf in den Himmel, daß GOtt dadurch erweckt wird zu helfen.

81. So lasset uns nun lernen stark sein und einen unverzagten Muth haben, wie viel und großes Uebel und Gefahr uns auch begegnen und wie oft auch Verzweiflung im Herzen sich regen mag. Und lasset uns halten, daß der heidnische Poet sagt: Tu ne cede malis, sed contra audentior ito, das ist: Weiche vor dem Unglück nicht, sondern gehe im Gegentheil um so kühner voran. Denn wer nicht weicht und die Hände nicht gar abläßt, der ist selig, und ist ein trefflicher Krieger, der das allerköstlichste und beste Opfer geopfert hat. Jakob ist nie ein heiligerer Priester gewesen denn hier an diesem Ort, wiewohl er zuvor auch durch viel Anfechtung ist getödtet worden. Darum gebraucht Mose so viel Worte, da er diese Historie beschreibt, und wird hernach den Kampf, da er mit dem Engel gerungen hat, erzählen, und werden größere Dinge geschehen, denn Jakob hat bitten oder hoffen dürfen. Auf solche Weise will GOtt haben, daß unser Glaube soll geübt und erweckt werden, auf daß wir von Tage zu Tage immer wachsen und stärker werden.

82. Bisher haben wir nun die Verheißung gehabt, die zum Gebet gehört, wie im Neuen Testament gelehret wird, und dazu auch das Gebot. Denn man muß nicht also beten, wie die Mönche ein unnützes Gemurmel und viel Worte zu ihrem Gebet gemacht haben, die an die Verheißung, oder das Gebot, oder auch an die Noth, die uns zum Gebet treiben soll, gar nicht gedacht haben. Das heißt nicht beten; wie ich selbst vor dieser Zeit auch gebetet habe. Da ich ein Mönch war, rief ich zwar GOtt an in der Noth; ich wußte aber von der Verheißung und vom Gebot nichts: wir haben nur einfach die Worte gemurmelt. Ein recht Gebet aber soll aus einem gläubigen Herzen herkommen, und das sich beides die Noth und das Gebot GOttes vorhält, dadurch das Herz erweckt wird, im Glauben zu beten, da alle Worte insonderheit erwogen werden; nicht wie die Mönche oder Nonnen ihr Gebet pflegen zu murmeln, daß das Herz ferne davon ist, nichts davon weiß oder versteht.

83. Und ein recht Gebet fragt auch nichts nach vielen Worten, sondern macht nur viel Seufzens, darauf seine Worte folgen, denn die nur sehr kleinlaut sind; wie an diesem Gebet zu sehen ist. Denn Jakob hat nicht allein mit diesen Worten gebetet, die Mose erzählt, sondern er hat die ganze Nacht und den ganzen Tag geseufzet. Es ist ein lang Gebet gewesen, so viel das Seufzen belanget, wiewohl der Worte sehr wenig sind. Darnach gehört auch zum Gebet Dankbarkeit und daß man der vorigen Wohlthaten darin gedenket soll. Wenn man die Exempel zusammen lieset, so erwecken sie den Glauben wunderbar; und das ist GOtt auch sehr angenehm.

84. So sagt nun Jakob: „Ich bin zu geringe aller Barmherzigkeit und aller Treue, die du an deinem Knechte gethan hast“; das ist: Du hast mir die größte Gnade und Treue erzeigt. Dieses ist die Dankbarkeit oder das Dankopfer; wie das erste ein Opfer der Tödtung gewesen ist, mit welchem Seufzen er getödtet wird und an seinen und aller Creaturen Kräften verzaget. Unter andere Wohlthaten aber zählt er diese und sagt: „Denn ich hatte nicht mehr, weder diesen Stab, da ich über diesen Jordan ging, und nun bin ich zwei Heere worden.“ So groß und reich bin ich bei Laban geworden, da ich anfänglich nicht mehr gehabt habe denn einen Stab, da ich aus meines Vaters Hause zog, das ist, da war ich gar allein. Ich hatte ja wohl Zehrung und ein wenig Geld (denn er war so gar arm nicht, daß ihm seine Eltern nicht sollten Zehrung auf den Weg gegeben haben), aber ich war allein und hatte keinen Gefährten. Es ist eine Art zu reden in der hebräischen Sprache, welche Mose auch gebraucht 2 Mos. 21, 3.: „Ist der Knecht ohne Weib kommen, so soll er auch ohne Weib ausgehen.“ Es bedeutet nur Einsamkeit und nicht Armuth an Geld.

85. Derhalben erkenne ich, spricht er, daß dieses mein Werk nicht sei, dieweil ich mit meinem Stabe allein über diesen Jordan gegangen bin, daß ich solch groß Gut erlangt habe; sondern daß es von deiner Barmherzigkeit und Treue herkommt. Ich habe Laban vierzehn Jahre umsonst gedienet, da war ich arm und nothdürftig, und hatte kaum das Brod zu essen; und du, mein lieber GOtt, hast mich innerhalb sechs Jahren so reich gemacht, daß ich in dieser kurzen Zeit zwei Heere geworden bin. Das ist ja ein offenbar und greulicher Segen, o GOtt! Denn es ist ein groß Ding, daß ein Hirte von Milch, Wolle, Käse und Butter ohne Geiz dermaßen machen und zunehmen soll. Und Jakob hatte vier Weiber, elf Kinder und einen großen Haufen Knechte.

86. Dieses Dankopfer hat er nun mit schönen Worten herausgestrichen und groß gemacht, welches sonst hier kürzlich beschrieben wird, wie er es hernach auch kurz gemacht hat. Denn das Gebet hat Tag und Nacht gewähret und wird an diesem Ort nur einmal erzählt. Jakob aber hat nicht einmal gebetet oder diese Worte geredet, sondern er hat ohne Aufhören gebetet.

87. Man weiß aber wohl, was es für eine Art zu reden ist, wo diese zwei Worte stehen, „Barmherzigkeit und Treue.“ Denn sie werden in der heiligen Schrift allezeit zusammen gesetzt, wie im Propheten Micha am 7. Cap. V. 18. und sonst an andern Orten in den Psalmen. „Barmherzigkeit“ heißt die Wohlthat selbst oder die Wohlthätigkeit, wie Matth. 9, 13. Christus spricht aus dem Propheten Hosea 6, 6.: „Ich habe Wohlgefallen an Barmherzigkeit, und nicht am Opfer.“ Das hebräische Wort chasid heißt beiderlei, nämlich, der seinen Nächsten liebt und ihm Gutes thut, und dem Gutes geschieht, dem GOtt große Barmherzigkeit und Wohlthat erzeigt. Darum sagt nun Jakob also: Lieber HErr GOtt! deiner Barmherzigkeit und Wohlthat, so du mir erzeiget hast, ist keine Zahl, und kommt nun dieses auch dazu, daß du mich armen Mann innerhalb sechs Jahren so reich gemacht hast, daß ich nun zwei Heere geworden bin.

\quellseite{203}

88. „Treue“ ist die Verheißung; wie St. Paulus redet Röm. 15, 8. 9.: „Ich sage aber, daß JEsus Christus sei ein Diener gewesen der Beschneidung um der Wahrheit willen GOttes, zu bestätigen die Verheißung, den Vätern geschehen: daß die Heiden aber GOtt loben um der Barmherzigkeit willen.“ Christus war den Juden verheißen, darum ist ihnen solche Verheißung treulich gehalten worden. Den Heiden ist er nicht verheißen gewesen, sondern er ist ihnen aus Gnaden gegeben worden. Unser Väter keiner hat die Verheißung gehabt. GOtt hat wohl verheißen, daß er aller Menschen GOtt sein wolle, welche Erkenntniß den Menschen in ihr Herz gepflanzet ist; wie St. Paulus Röm. 1, 19. ff. anzeigt. Und die Werke und Gottesdienste aller Völker bezeugen das auch, daß „GOtt“ kein Anderes nichts sei, denn den Menschen Gutes thun. Denn darum ruft Einer den Jupiter, der Andere den Mars an etc., aus keiner andern Meinung, denn allein darum, daß sie wollen Hülfe haben. Also verstehen alle Menschen von Natur und halten es dafür, daß GOtt eine solche göttliche Gewalt sei, die da wohlthätig sei, von der man alles Gute bitten und erwarten soll. GOtt ist, der da verheißt, und ist wahrhaftig, daß ist, er verheißt allen Menschen im Gesetz der Natur, das da sagt: Du sollst GOtt anrufen und ihm dienen, 5 Mos. 6, 13.

89. Wiewohl sie nun in der Person Gottes irren um der Abgötterei willen, so ist doch gleichwohl der Dienst da, der dem rechten GOtt gebühret, das ist, die Anrufung, und daß sie alles Gute und Hülfe von ihm erwarten. Diesem Volk aber sind die Verheißungen von GOtt gegeben worden, da ist die Wahrheit und rechte Erkenntniß, wer der rechte GOtt sei, ohne Abgötterei und Irrthum ausgedrückt. Das ist aber eine große und unermeßliche Wohlthat Gottes. Darum rühmt Jakob die Barmherzigkeit und Treue GOttes und will also sagen: Du hast mir verheißen und diese Treue auch gehalten, vornehmlich in diesen sechs Jahren; desgleichen hast du mir auch verheißen, daß du mir wohlthun wollest, wenn ich wieder in mein Vaterland ziehen werde. Um solcher Barmherzigkeit und Treue willen, die du mir erzeiget hast und noch weiter erzeigen wirst, bitte ich, du wollest mich erretten.

90. Diese Dankbarkeit begreifst ferner auch die dritte Tugend des Gebets in sich, welche auch zur Tödtung gehört, nämlich, daß er ohne sein Verdienst erhöret werde. Denn man muß nicht also beten: Lieber HErr GOtt, siehe mich an; denn ich bin ein heiliger Mönch, ich bin eine keusche Jungfrau, ich bin ein trefflicher Bischof. Wie der Phariseer Luc. 18, 12. sagt: „Ich faste zweimal in der Woche, und gebe den Zehnten“ etc., oder gebe so viel Almosen. Solchen Greuel soll man weit vom Gebet wegthun und ausfegen, sonst wird das Gebet verderbt, und böserlei Opfer, nämlich, der Tödtung und der Dankbarkeit, wird mit solchem Stand und Unflat gar zu Schanden gemacht.

91. Darum soll man mit nichten also sagen, sondern wir sollen dem Exempel Jakobs nachfolgen, der da sagt: „Ich bin zu gering aller deiner Barmherzigkeit“, als wollte er sagen: Ich bin auch nicht Einer Wohlthat oder Treue werth, die du mir jemals erzeigt hast, oder noch jetzt erzeigst und hinfort erzeigen wirst ohne mein Verdienst; denn es ist unmöglich, daß ich etwas sollte verdienen können: darum verlasse ich mich auf meine Würdigkeit nicht, sondern auf deine Verheißung und Barmherzigkeit.

92. Das ist ein rechtschaffen Herz und wahrhaftiges Gebet; denn dies ist in Wahrheit unsers HErrn GOttes Urtheil von uns, welches St. Paulus ausgedrückt hat 1 Cor. 4, 7.: „Was hast du, das du nicht empfangen hast? So du es aber empfangen hast, was rühmest du dich denn?“ Darum soll man also beten: Lieber HErr GOtt, was du mir bisher gegeben hast, das ist deine Verheißung und Barmherzigkeit, an mir ist gar keine Würdigkeit. Derhalben, wie du bisher gethan und mir, wiewohl ich dessen nicht werth bin, Alles gegeben hast: also wollest du mich Unwürdigen hinfort auch erhöhen und mir gnädiglich helfen. Sonst wenn du anders betest, wird dich GOtt strafen und sagen: Was stolzirest du, der du Asche und Staub bist? du bist doch Erde und wirst wieder zur Erde werden, 1 Mos. 3, 19. Habe ich dich nicht aus Dreck geschaffen? Wer hat dir dein Leben erhalten? Alles, was du hast, ist einfach mein.

93. So sind nun dies dreierlei Eigenschaften eines guten Gebets, die da machen, daß es GOtt sehr angenehm und der süßeste Geruch vor ihm ist, und ein solch Gebet muß erhört werden. Die erste Eigenschaft aber ist, daß du die Verheißung ergreifest. Darnach, daß du in deiner Angst getödtet werdest. Zum dritten, daß du GOtt dankest, und erkennest, daß du auch nicht Einer Wohlthat werth seiest; sondern daß du bittest und hoffest, daß dir möge geholfen werden allein aus Gnade und Barmherzigkeit. Es sind zwar diese Tugenden zu loben und auch von GOtt geboten worden, als nämlich, Keuschheit, Nüchternheit, und daß man freigebig ist, den Armen zu helfen. Man soll aber darauf sich nicht verlassen. Unser Vertrauen soll nicht stehen auf dem Gesetz und seinen Werken, wiewohl die Werke sollen da sein; sondern das Vertrauen soll stehen auf GOttes Barmherzigkeit und Treue. Aldann ist das Gebet und das Seufzen ein golden Opfer.

94. Nun folgt die Bitte. Was bittest du denn nun, Jakob? Du erkennst dich meiner Barmherzigkeit unwürdig; du erkennst, daß ich der sei, der dir die Verheißung gethan hat; du bekennst, daß du Alles von mir aus lauter Gnade empfangen: was bittest du denn nun? was mangelt dir? Das Kind muß einen Namen haben, soll man es taufen. Gleichniß als ob GOtt nicht wisse, was wir bedürfen. Er weiß es wahrlich wohl. Denn also spricht Christus Matth. 6, 8.: „Euer Vater weiß, was ihr bedürfet, ehe denn ihr ihn bittet.“ Warum will er denn haben, daß ich seufzen, schreien und anklopfen soll, und dazu mich auch seiner Hülfe unwürdig erkennen?

95. Antwort: Das will er darum haben, daß du auch bedenken sollst, was dir mangelt und fehlt; sonst sind wir sicher und schnarchen, und lassen uns träumen, daß wir die seligsten Leute seien, so wir doch die elendesten sind. Das wissen wir aber nicht, lernen es auch sonst nicht, denn allein zur Zeit der Anfechtung. Jakob hatte von GOttes Gnade das Leben, er hatte zwei Heere; aber in dem Augenblick verliert er beides und wird gar zunichte gemacht. Daselbige hat er zuvor nicht gewußt, da er also rühmte und triumphirte; jetzt wird er aber nun zur Schule geführt, und wird gelehrt, was und wie viel er vermöge, auf daß wir lernen unser Vertrauen nicht auf uns selbst stellen, sondern auf GOtt, daß wir uns der Gaben also freuen, daß wir uns derselben gleichwohl nicht rühmen. Denn dasselbige ist sehr schwer also zu treffen und der verderbten Natur willen, die durch die Erbsünde dermaßen verderbet ist, daß wir leichtlich in Verzweiflung fallen können, wenn wir angefochten und betrübt werden. Wenn es allenthalben glücklich und wohl zugeht, so sind wir sicher und stolz. Darum werden die Heiligen bisweilen niedergeschlagen und gedemüthigt, und bisweilen dann auch wiederum aufgerichtet, auf daß sie lernen halten, daß man pflegt zu sagen: Medio tutissimus ibis: Wenn man fein in der Mitte bleibt, das ist das Allerbeste; das geschieht nicht, denn allein durch Anfechtung. Darum sollen wir nicht vermessen sein, sondern unser Vertrauen allein auf GOttes Barmherzigkeit und Treue stellen. Denn dieselbige steht fest und währet ewig, vergeht nicht, und unsere Hoffnung und Vertrauen, so wir darauf setzen, kann nicht fehlen.

96. Was mangelt aber nun Jakob? Antwort: Es mangelt ihm Alles. Mein Bruder, will er sagen, wird mich, meine Weiber und Heere erwürgen; davor fürchtet er sich. Wo ist aber nun die Verheißung GOttes, die er mit dem Glauben ergriffen hat? Weißt du nicht, daß GOtt deinem Bruder sein Herz ändern kann? GOtt hat Jakob Barmherzigkeit und Treue verheissen und erzeigt, doch zittert und seufzt das Fleisch und flieht vor der Tödtung: das geschieht von wegen der Erbsünde. Er sieht nur die Gefahr an. Ich habe wohl zwei Heere, will er sagen, aber ich sehe auch wohl, daß es bald geschehen kann, daß ich Alles verliere und daß ich auch selbst der Gefahr des Todes kaum entrinnen kann. Denn er sagt also: „Ich fürchte mich, daß er komme, und schlage mich, die Mütter sammt den Kindern.“ Es wird der Glaube abermals angefochten und wird schwach, verzweifelt aber doch gleichwohl nicht.

\quellseite{204}

97. Denn also pflegt es zu gehen, wenn die Anfechtungen vorhanden sind: der Glaube kämpft mit dem Fleische, und die Gedanken laufen zusammen, deren etliche verzagen und etliche gute Hoffnung haben: es sind widerwärtige Winde des Glaubens und Unglaubens, der Hoffnung und der Furcht. Er möchte mich, die Mütter sammt den Kindern schlagen, sagt er. Daraus sieht man, was seine Angst und Furcht gewesen ist. Er hat bald in einem Hui Alles verloren: da ihm die Nachricht kommt, daß ihm sein Bruder entgegen ziehe, da fällt alle sein Gut und Hoffnung dahin; doch bleibt er noch an der Verheißung hangen: „HErr, du hast gesagt: Ich will dir wohlthun, und deinen Samen machen“ etc. Das hält er gegen den Zweifel, und da er sich läßt dünken, daß er gar vernichten müsse, dieweil er sagt: „Daß er nicht komme“ etc. Denn das sind Worte des Fleisches gewesen, das da streitet, und eines sehr schwachen Glaubens. Er ergreift aber das Wort: Deine Treue, sagt er, wolle mich erhalten, aufrichten und trösten, dadurch du verheißen hast, daß du mir wollest wohlthun, mich und meinen Samen segnen. Das wird aber nicht geschehen, du wirst mich und meinen Samen nicht segnen, wenn mich Esau aufreiben und vertilgen wird.

98. Nun hätte aber GOtt den Samen auch wohl vermehren und den Segen ausbreiten können, wenn schon Jakob wäre vertilgt worden; wie die Verheißung des Samens und der Erben an Isaac auch gegangen hat, und es wird doch von GOtt geboten, daß er sollte geopfert werden, und Abraham sollte gleichwohl den Glauben behalten, daß GOtt auch aus der Asche Isaaks Samen hervorbringen könnte, wie droben gesagt ist. Das ist wahrlich eine sehr schwere Anfechtung gewesen. Jakob aber hätte schier in dieser Furcht des vergessen, daß GOtt gesagt hatte: Er wollte sein GOtt sein, und wie er droben, V. 9., das Gebet angefangen hat: „HErr, der du zu mir gesagt hast“ etc.

99. Er hat aber nicht allein also gebetet, sondern wird auch etliche Predigten und Ermahnungen gethan haben an seine Familie, auf dieselbige Meinung. Er wird oftmals zu Rahel gesagt haben: Meine liebe Rahel, sei getrost, hoffe auf GOtt und sein Wort. Ja, wird sie wiederum gesagt haben, er wird uns aber erwürgen? Da hat Jakob geantwortet: Er wird es nicht thun. Ich habe den HErrn gebeten, daß er uns schützen und erretten wolle. Bedenke, daß GOtt mit mir geredet und daß er mir solche verheißen hat. Mit solcher Vermahnung, Predigten und Gebet hat er sich und seine Familie aufgerichtet und getröstet.

100. Und dieses sind sehr treffliche Stellen und auch sehr merkwürdige Exempel in der Historie der Väter, die uns lehren, daß wir vor allen Dingen GOttes Wort ergreifen und uns daran fest halten sollen, daß wir das ja nicht verlieren oder wegwerfen, wenn wir auch schon das Andere alles verloren haben, wenn uns schon Esau die Schafe todtschlägt. Wie greulich uns auch die Feinde drohen, so sollen wir uns doch den Glauben an die Verheißung nicht nehmen lassen, sondern sollen am Worte hangen bleiben und darin leben, wie die heiligen Väter, Abraham, Isaak und Jakob, darin gelebt haben. Denn ob du wohl schwach bist, sollst du doch deswegen das Wort nicht wegwerfen, der Sache wird wohl gerathen; wirf du aber das Wort weg, oder wirst dir es lassen von den Augen und aus dem Herzen wegnehmen, so wird der Teufel triumphiren und wird dich zur Verzweiflung treiben.

101. Mit dieser Anfechtung und Gebet ist nun Jakob die Nacht über zu Mahanaim geblieben, welches eine Stadt ist, über dem Bach Jabok gelegen, dahin David geflohen ist, da ihn Absalom aus dem Lande verjagt hatte. Es ist eine berühmte Stadt in der heiligen Schrift, nicht weit vom Jordan gelegen; denn der Jabok fließt in den Jordan. Gilead, Mizpa, Mahanaim sind Städte, die allernächst dabei liegen.

\section[Dritter Theil]{Dritter Theil. Wie Jakob auf mancherlei Rath bedacht ist, Esau zu versöhnen, und ihm endlich Geschenke entgegen schickt; und wie er seine Familie über das Wasser setzt}

\begin{quote}
\textbf{V. 13—15. Und er blieb die Nacht da, und nahm von dem, das er vorhanden hatte, Geschenk seinem Bruder Esau: zwei hundert Ziegen, zwanzig Böcke, zwei hundert Schafe, zwanzig Widder, und dreißig säugende Kameele mit ihren Füllen, vierzig Kühe, und zehn Farren, zwanzig Eselinnen mit zehn Füllen.}
\end{quote}

102. Ich habe droben gesagt und man muß es oft sagen, daß man GOtt nicht versuchen, sondern daß man Rath, Mittel und Hülfe gebrauchen soll, die vorhanden sind, auf daß wir den Türken nicht gleich werden. Denn es ist eine sehr böse Versuchung, daß Einer Hände und Füße gar will stille halten und sich in Gefahr nicht fahren begeben, daß du dich vor der Gefahr nicht hüten oder derselbige meiden willst, wenn du wohl könntest, und willst darnach GOtt die Schuld geben, als ob es sein Wille gewesen sei. Darum soll man die Gottseligen daran erinnern, daß wir Menschen nicht wissen können, was die Sachen, so da vorfallen, für einen Ausgang gewinnen mögen und was zukünftig sei. GOtt weiß allein, was zukünftig ist und was geschehen soll. Dieweil wir aber das nicht wissen, so sollen wir gleichwohl thun, was wir können. Und daher gehören die Sprüche, daß der Mensch thun soll, was an ihm ist; nicht in Sachen, die GOttes Gnade betreffen, sondern in äußerlichen Dingen, und wo Noth und Gefahr vorhanden ist, da soll man der Vernunft folgen, auf daß wir GOtt nicht versuchen, 1 Cor. 10, 9. Und Augustinus sagt von diesem Texte also: Wir sollen durch diese Exempel ermahnt werden, daß wiewohl wir an GOtt glauben, wir doch gleichwohl thun sollen, was zu thun ist; auf daß, wenn wir solches unterlassen, es nicht scheine, als ob wir GOtt versuchten.

\subsection{I}

103. Denn also hat Jakob sich mehrerer Auswege bedient. Erstlich, da er an die Grenze des Landes kommt, da Esau gewohnt hat, schickt er Boten vor sich her, die sich erkundigen sollen, wie sein Bruder möge gesinnet sein, und die Freundschaft und Versöhnung bei ihm suchen sollen. Aber der Ausweg ist mißglückt. Denn sein Bruder zieht ihm entgegen mit vierhundert Mann, ist viel gewaltiger und Jakob viel feinder, wie es scheint, denn zuvor. Er hat aber doch daran recht gethan, daß er die Boten ausgeschickt hat. Und wo er dieses Mittel nicht gebraucht hätte, so hätte er damit gesündigt.

104. Der andere Ausweg ist, daß er das Volk theilt und desgleichen die Heerde auch in zwei Heere (denn er verzagt nun am Frieden und an der Freundschaft seines Bruders), auf daß er nur den Einen Theil retten möchte. Und dasjenige hat ihm auch nicht geholfen; denn GOtt bestätigt das Herz seines Bruders. Er hätte aber keinen bessern Ausweg oder Weise finden können, Etliche aus seinem Volk von der Tyrannei seines Bruders zu erretten. Denn Esau war allzu gewaltig, und wird ohne Zweifel ein Bundesgenosse gewesen sein der Fürsten zu Seir und der andern umliegenden Völker. Er hat Kinder, Rinderfinder und Schwiegertöhne gehabt; denn sie sind beide zu der Zeit, da sie zusammen gekommen, hundert Jahre alt gewesen, oder je zum wenigsten acht und neunzig Jahre. Esau hat sechzig Jahre ein Weib gehabt; in solcher langen Zeit ist seine Familie sehr groß geworden, und er hat am Gute und Gewalt zugenommen von wegen seiner Schwiegerväter und guten Freunde, die ihm verwandt gewesen sind. Jakob hat im vierundachtzigsten Jahre seines Alters zwei Weiber genommen, darum führt er mit sich Kinder von zwölf oder vierzehn Jahren; Joseph war sechs Jahre. Was ist nun dieses arme elende Häuflein gegen ein so großes Heer, oder je gegen die Kinder und Kindeskinder Esaus, die Jakob gleicher Gestalt gehabt und ihm feind gewesen sind? Darum hat Jakob daran weislich gethan, daß er sein Volk und Heerde getheilt hat; wiewohl das auch noch umsonst gewesen wäre, wenn Esau seinem grimmigen Zorn hätte folgen können.

105. Auf solche Weise sollen wir bedenken, daß wir auch, wenn Noth und Gefahr vorhanden ist, versuchen und thun sollen, was die Vernunft und menschliche Geschicklichkeit uns weiset oder zeigt; wiewohl bisweilen auch der beste Rath, so wir haben erdenken mögen, zunichte wird. Und das sollen wir auch um der Ursache willen thun, auf daß wir hernach unsere Nachlässigkeit und Thorheit nicht zu spät an uns selbst strafen dürfen, wie oftmals solche Klagen gehört werden, daß Mancher sagt: Wenn ich dieses oder jenes erst versucht hätte, so wäre dieses Unglück nimmermehr gekommen. Solche Leute rathschlagen allzu langsam, wenn ein Ding schon geschehen ist, und beschweren damit auch ihre Gewissen, daß sie in Zeiten seine Hülfe oder Rath gebraucht haben. Derhalben soll man rathschlagen und mit allem Fleiß thun alles, was man nur aus menschlichem Rath erdenken kann, auf daß du in deinem Herzen zufrieden sein und sagen mögest, du habest dich zuvor umgesehen, bedacht, und alles gethan, was da zu thun möglich gewesen.

\quellseite{205}

106. Der dritte Ausweg ist, daß er ihm Geschenke schickt. Wiewohl aber nun dasselbige auch nicht wohl geraten ist, so ist es doch gleichwohl sehr fein bedacht gewesen. Und Mose beschreibt es alles ganz fleißig, wie er in dieser ganzen Erzählung etwas mehr Worte gebraucht, als er sonst zu thun pflegt: denn es ist eine heftige und starke Anfechtung und eine sehr schwere Tödung. Nachdem ihm nun die zweierlei Auswege mißfällig waren und der zornige Muth Esaus nicht hat versöhnet werden können durch die Boten, so Jacob zu seinem Bruder geschickt hatte, so sendet er ihm nun ein Geschenk zu schicken und dasselbige mit Worten und äußerlichem Gepränge herrlich zu schmücken, auf daß Esau, der zornig war, dadurch etwas erweicht werden möchte, daß er gelinder würde. Und dasselbige erzählt Mose mit sonderlichem Fleiß. Denn er schickt ihm nicht allerlei Vieh; wie heutiges Tages gepflegt zu geschehen, wenn irgend einem Fürsten hundert Ochsen, oder zwei hundert Schafe geschenkt werden \&c., und das ist ein sehr köstlich Geschenk: sondern er nimmt aus einer jeglichen Art eine gewisse Anzahl und schickt ihm also ein Geschenk von mancherlei Heerde Viehes.

107. Es haben aber die Geschenke an sich selbst die Art und Natur, daß Andere damit können versöhnet werden; wie auch der heidnische Poet sagt: Munera, crede mihi, placant hominesque deosque, das ist: Mit Geschenken kann man beide Götter und Menschen versöhnen. Derhalben hat kein besserer Ausweg oder Weise können erfunden werden als das Anbieten der Geschenke. Denn diese zwei Stücke, nämlich, daß man sich persönlich bemühte, oder einen Andern abschickte, wo er zornig ist, und daneben auch Geschenke gibt, sind sehr kräftig, und dienen wohl dazu, daß damit die Herzen, so erbittert sind, versöhnt werden mögen. So schickt er ihm nun erstlich etliche Ziegen und Böcke, und also nach einander eine gewisse Anzahl von allerlei Heerde insonderheit; wie im Text erzählt wird.

108. Es sagt aber Mose: „Er nahm von dem, das er vorhanden hatte“; das hat Hieronymus gegeben: Er schied. Und ist dieses der Verstand, daß er nicht das Allerbeste von der Heerde ausgelesen, sondern ist durch die ganze Heerde gegangen, darin alles Vieh wohlgehalten und gemästet war, und daraus hat er es gelesen, wie ihm ein jedes vorgekommen ist.

109. Das hebräische Wort minchah heißt Geschenk. Droben, Cap. 4. V. 3. 4., wird es gebraucht von den zwei Brüdern, Abel und Kain. Es heißt aber ein Opfer und Gabe, so geopfert wird. Es heißt nicht allein ein Opfer, so Gott geopfert ist, sondern auch eine Gabe oder ein Geschenk, so den Menschen gegeben wird. Die zwei hebräischen Worte izim und parim reimen sich mit den deutschen Namen, Geiß und Farre. Wir haben es von den Juden genommen. Darum ist es nun ein herrlich Geschenk, und ist sein ordentlich unterschieden nach einer jeglichen Art des Viehes, so zur Haushaltung gehört.

110. Siehe, wie fleißig er Alles verordnet, auf daß er seines Bruders Herz, der sehr zornig war, desto leichter erweichen möge, wenn er solchen seinen Anblick vor sich sehen wird. Er sondert die Heerden aus und gibt einer jeglichen Heerde einen sonderlichen Hirten, und ordnet es also, daß Raum sei zwischen einer Heerde nach der andern, auf daß Esau, wenn er ihnen entgegen zieht, die Heerde nicht ineinander gemenget sehe, sondern daß ihm erstlich die Ziegen, darnach die Böcke, zum Dritten die Schafe, und so weiter, daß zwischen jeder Heerde ein Raum gelassen war, die Heerden begegnen mögen. Das hebräische Wort revah heißt nicht eigentlich ein Raum, wie sie es verdolmetschen, sondern Ruhe. Wie 2 Mos. 8, 15. steht von der andern Plage, da Pharao und die Egypter mit den Fröschen geplagt worden sind: „Da Pharao sahe, daß er Luft krieget hatte.“ Ich halte aber dafür, daß Mose dieses Wort darum habe gebrauchen wollen, weil dieses Vieh sich gern in einen Haufen versammelt, also daß sie die Köpfe einander zukehrten, daß sie kaum Luft haben könnten. Wie im 49. Psalm V. 15. steht: „Sie liegen in der Hölle wie Schafe“ \&c. Derhalben will er haben, daß sein Geschenk soll gezieeret werden nicht allein damit, daß es fein unterschieden ist, sondern daß auch der Heerden möge verschont werden, auf daß sie nicht verletzt werden oder umkommen, wenn sie so gar gedrängt in einem Haufen bei einander stehen. Er setzt aber nun auch noch seine geschmückte Worte dazu.

\begin{quote}
\textbf{V. 17–20. Und gebot dem ersten, und sprach: Wenn dir mein Bruder Esau begegnet und dich fraget: Wem gehörst du an? und wo willst du hin? und wes ist’s, das du vor dir treibest? sollst du sagen: Es gehöret deinem Knechte Jakob zu, der sendet Geschenk seinem Herrn Esau, und zeucht hinter uns hernach. Also gebot er auch dem andern, und dem dritten, und allen, die den Heerden nachgingen, und sprach: Wie ich euch gesagt habe, so sagt zu Esau, wenn ihr ihm begegnet; und saget ja auch: Siehe, dein Knecht Jakob ist hinter uns. Denn er gedachte, ich will ihn versöhnen mit dem Geschenk, das vor mir hergehet; darnach will ich ihn sehen, vielleicht wird er mich annehmen.}
\end{quote}

111. Dieses sind zwar kurze Worte, aber sie sind sehr künstlich gesetzt. Er befiehlt ihnen, daß sie das Geschenk mit Worten zieren und ihn auch ehrlich anreden sollen. Ihr sollt nicht sagen, spricht er, diese Heerde, in fünf Haufen getheilt, gehört deinem Bruder Jakob. Also sollt ihr mit nichten sagen, auf daß ihm nicht einfalle zu gedenken, wie ihm der Segen ist entzogen worden; sondern ihr sollt sagen: Diese Heerde gehört deinem Knechte, deß Herr du bist. Dieses ist eine große Demuth, daß sich Jakob also herunter wirft; dadurch sollte wahrlich das erzürnte Herz beweget und erweicht werden.

112. Und diese Worte befiehlt er nicht allein Einem Hirten, sondern ihnen allen, daß sie ja sein bescheiden und demüthig reden, und Jakob einen Knecht, Esau aber einen Herrn nennen. Und also war auch ihrer beider Rüstung. Denn Esau zog daher mit einem großen Haufen wie ein Herr, mit Roß und Wagen; Jakob aber kommt zu Fuß wie ein Schafhirte, und sendet seinem Bruder kein königlich Geschenk, sondern ein Hirten geschenk. Damit hat er gehofft, seinen Bruder erweichen zu können, daß er bedenken sollte: Was willst du mit dem machen, der mehrmals und gar elend und arm ist, der ein solches Häuflein führt, das zum Kriege gar nicht dient, als, Kindlein, Weiber, Mägde und Knechte? deren muß man allerdings verschonen, dieweil sie mir zu Füße fallen. Dieses ist wahrlich eine große Demuth, daß sich ein solcher großer Patriarch dermaßen herunter wirft, der so große Verheißungen und dazu auch Geschichte der Engel gehabt hat. Aber die Tödung demüthiget ihn also, und erinnert ihn auch zugleich dessen, daß er in der Noth und Gefahr, so gegenwärtig vorhanden ist, Rath schaffe, wie er nur kann und mag.

113. Er hat fleißig verhütet, daß Esau auch nicht mit Einem Wort mehr möchte erzürnt werden; darum heißt er sie, daß sie sich enthalten sollen, daß sie Esau nicht Jakobs Bruder oder Freund nennen, und daß sie ihm der Namen keinen geben, woraus man irgend Hoffahrt abnehmen möchte. Nun war ja Jakob mit nichten geringer als sein Bruder Esau; denn er war vom Vater gesegnet worden, und in dem Segen wird ausdrücklich gesagt: „Deiner Mutter Kinder müssen dir zu Füße fallen“, 1 Mos. 27, 29. Hier sieht man wahrlich gar das Gegentheil von dem, das ihm verheißen worden ist. Ja, er wird noch mehr gedemüthigt werden. Ich habe aber zuvor gesagt, daß dieses die Zeit der Tödung sei, und daß sich Jakob darum so herunter wirft, daß er seinen Sachen Rath schaffen möge in der Noth und Gefahr, so jetzt gegenwärtig vorhanden war, auf daß es nicht scheinen möchte, als hätte er sich leichtfertig in Schaden und Verderben gebracht.

114. Er setzt noch mehr dazu und sagt: „Er selbst zeucht hinter uns hernach“: er flieht nicht, sondern folgt uns nach und versieht sich Gutes zu dir, und hofft, Gnade und Frieden bei dir zu finden. Dieses sind auch sehr freundliche und liebliche Worte, die ein Herz zufrieden stellen und erweichen könnten, wie hart es auch möchte verletzt sein. Und er ermahnt sie fleißig, daß sie ja dermaßen mit seinem Bruder reden wollen, wie er ihnen die Worte vorgesagt hatte. Hütet euch, daß ihr nicht sagt, daß sein Bruder euch nachfolge, oder daß ich bei Laban reich geworden sei und daß ich ein großes Gesinde habe. Seid nicht stolz und hoffärtig, und macht mich nicht verhaßt, sondern demüthigt euch und sagt: „Dein Knecht ist hinter uns.“

\quellseite{206}

115. Denn „ich will ihn versöhnen“, spricht er, „mit dem Geschenk, und darnach will ich ihn sehen, vielleicht wird er mich annehmen.“ Das ist Jakobs Plan: „Ich will ihn mit dem Geschenk versöhnen“, daraus will ich leichtlich abnehmen, ob er etwas besänftiget sei. Wo er ja dem ersten Hirten nicht antworten wird, so wird er ja dem andern, dem dritten, oder denen, die hernach folgen, antworten; und wenn das Geschenk überreicht ist, so will ich selbst auch kommen und mich vor ihm demüthigen: „Vielleicht wird er mich annehmen“, das ist, er wird mich zu Gnaden nehmen, er wird sich mit mir versöhnen. Denn dieses ist in der Schrift eine sonderliche Art zu reden: suscipere faciem, das heißt, Einen zu Gnaden annehmen. Facies peccatorum sumere, heißt, wenn man der Sünder Sünde lobt und sie nicht straft.

116. So sehen wir nun, wie der Mensch thun soll, das an ihm ist. Denn also sendet Jakob seinem Bruder Geschenke und ordnet solches Geschenk auf das allerkünstlichste, und endlich ziert er es auch mit Worten, so viel er immer kann. Das heißt thun, was an ihm ist, und Gott nicht versuchen. Wenn er das nicht gethan hätte und ein Todtschlag geschehen wäre, so würden diese Gedanken ihn stets in seinem Gewissen gemartert haben: Warum hast du dich vor deinem Bruder nicht gedemüthigt und ihn mit gebührlicher Ehrerbietung angeredet? Wenn er gesehen hätte, daß du ihm wärest zu Füße gefallen, so würde er deiner ohne Zweifel verschont haben: dieweil du aber nachlässig gewesen bist, Rath und Hülfe zu finden, so geschieht dir jetzt eben recht, daß du dieses leiden mußt.

\subsection{II}

\begin{quote}
\textbf{V. 21—24. Also ging das Geschenk vor ihm her, aber er blieb dieselbe Nacht beim Heer, und stund auf in der Nacht, und nahm seine zwei Weiber, und die zwo Mägde, und seine eilf Kinder, und zog an die Furt Jabbok, nahm sie, und führte sie über das Wasser, daß hinüber kam, was er hatte; und blieb allein.}
\end{quote}

117. Da Jakob die Boten mit dem Geschenk voraus gesandt hatte, folgt er nicht alsbald nach, sondern bleibt dieselbe Nacht im Lager, daß eine Tagereise zwischen Jakob sei, und dem Geschenke, so vorausgeht, und daß er aus solchem Thun, so zwischen ihnen war, seines Bruders Gemüth erfahren möge. Er ist aber noch nicht zufrieden, sondern wird noch mehr getödtet. Denn in der Nacht steht er auf und setzt mit Weib und Kindern über die Furt Jabbok. Und läßt es sich angehen, als wolle der Text anzeigen, daß Jakob heimlich, da andere Leute geschlafen haben, zu Mitternacht habe wollen über den Bach fahren, und es scheint, als habe er sich gefürchtet, daß wenn sein Bruder des Morgens kommen würde, er ihm den Zug verhindern möchte. Darum hat er ihm den Zug verhindern, seine Weiber und Kinder, und alles, was er hatte, überzuführen, und dasselbe hat er innerhalb zween Stunden kaum ausrichten können.

118. Mit Nachdruck aber sagt Mose: „Er führete sie über das Wasser“; und lautet dies, als hätte er seine Familie erstlich an den Bach geführt und wäre er allein hinüber gefahren, und hätte versucht, an welchem Ort man am leichtesten überfahren könnte, und hätte darnach erst die Weiber und Kinder übergeführt. Da er sie aber übergeführt, hat er sich nun der ganzen Familie an einen besondern Ort begeben und hat daselbst mit großer Bekümmerniß gebetet; denn er ist noch voll Schreckens und Schwachheit. Darum wählt er einen solchen Ort, der da einsam ist, darum daß ein solcher einsamer Ort denen, die da beten, sehr wohl gelegen ist, auch am Tage. Denn wenn das Gebet ernst und brünstig ist, so haben wir nicht gern, daß man die Worte, die wir vor Gott etwa ungeschicklicher Weise heraus schütten, hören soll.

119. Darum ist er nun des Nachts aufgestanden, und hat eine Stunde oder zwei zu thun gehabt, daß er seine Weiber, und was er sonst mehr gehabt, übergeführt hat: und so viel Stunden hat auch der Kampf gemacht, den er geführt hat. Es schickt sich aber die Zeit bei der Nacht vornehmlich wohl zu Anfechtungen; wie droben im 15. Capitel, §. 12. gesagt worden, da Mose sagt: „Da nun die Sonne untergegangen war, fiel ein tiefer Schlaf auf Abram; und siehe, Schrecken und große Finsterniß überfiel ihn“ etc. Die Finsterniß an sich macht die Herzen der Menschen und auch der unvernünftigen Thiere betrübt und verzagt; wiederum aber wenn der Tag und das Licht anbricht, das erfreut alle Thiere. Darum sagt man im deutschen Sprüchwort recht: Die Nacht ist niemandes Freund. Es ist eine Zeit, die da wohl dient zur Anfechtung und zum Gebet, dazu die Anfechtung und Trübsal den Menschen treibt. Darum ist nun Jakob brünstiger geworden zum Gebet, aber doch in großer Schwachheit.

120. Wir sehen aber, wie Mose, ja, der Heilige Geist selbst, mit so großem Fleiß auch die allerkleinsten Werke und Leiden der Patriarchen beschreibt, an welchen man keine wunderbaren Werke, und die einen großen Schein haben, sieht; wie die Mönche und Werkheiligen sich solcher gleisenden Werke haben zu rühmen wissen. Aber daher ist dieses Leiden der Heiligen sehr trefflich und golden, darum daß es diese Verheißung hat, daß vor dem Herrn nicht allein ihr Tod oder Blut werth gehalten ist, sondern daß auch die Haare auf ihrem Haupte werth und gezählet sind. Darum hält es der Heilige Geist nicht für unnütz, daß er so lange bei diesen häuslichen und Hirtenwerken und Leiden verweilt; denn in demselbigen wird der Glaube recht geprüft, und wird ein Opfer, das Gott sehr angenehm ist.

121. Darum sollen uns diese Exempel der Patriarchen immer vor unsern Augen liegen, auf daß wir lernen, daß auch unsere Werke, die im Glauben geschehen, desgleichen unsere Trübsal, so wir im Glauben leiden, Gott alle angenehm und gefällig sind, und wie ein stetiges Opfer. Dieses ist ein großer Trost der Gläubigen. Das Leben der Gottseeligen scheint ein unnützes Leben zu sein und gar keine Früchte oder Herrlichkeit zu haben. Aber dieses ist eine treffliche Ehre und Herrlichkeit, daß wir wissen, daß Gott auch unsere Thränen, und ein jegliches Tröpflein, so aus unsern Augen fällt, zählt, und daß es alles vor Gottes Augen angeschrieben und gleichsam in ein golden Gefäßlein gesammelt sei.

122. Es ist aber kein Mönchsfasten von sieben oder mehr Tagen, ist kein Wachen oder eheloses Leben, deren keines aus dem Glauben herkommt. Und darum verfolgt auch der Teufel solche Mönchswerke nicht, sondern hegt sie, daß sie dabei ohne Gefahr sein mögen: Gott aber gefallen sie sehr übel; denn sie werden alle vorgenommen und geschehen auch alle wider das Wort. Aber solche Werkheiligen sind hin zu zweifältig arme elende Leute; denn dieweil sie in des Teufels Stricken gefangen liegen, werden sie viel heftiger geplagt denn alle Gottseeligen; denn was ist elender und erbärmlicher, denn dem Teufel gehorsam und unterworfen sein?

123. Es ist aber auch zu verwundern, daß der Patriarch so eine gehorsame Familie und Gesinde gehabt, das mit so großer Geduld so viel Noth und Gefahr ertragen hat, und vornehmlich diese Boten, die sich mit Geduld dazu gebrauchen lassen, daß sie dem zornigen Bruder entgegen gesandt werden. Dieselbigen hätten wieder nach Mesopotamien ziehen mögen, und daselbst sicher und ohne Gefahr sein: sie tragen aber mit ihrem Herrn die gemeine Gefahr, der sich dessen auf das höchste befliessen hat, daß er sie im Glauben und Erkenntniß Gottes erhalten möchte, der ja selbst die Verheißung gegeben hatte, darauf sie zugleich mit ihrem Herrn gewartet haben. Derhalben ist Jakob ein trefflicher und sehr nützlicher Hoherpriester, der so viel Menschen aus Mesopotamien errettet und sie durch reine Lehre und rechten Gottesdienst selig macht. Und seine Hausgenossen sind rechte Märtyrer gewesen, die mit solchen großen Beschwerungen um ihres Herrn willen gekämpft haben.

124. Du wirst jetziger Zeit in diesem Stand der Menschen, der fast verderbet ist, Wenige finden, die solcher Geduld und Beständigkeit nachfolgen. Wo aber das Wort gelehret wird, da ist der Glaube, und wo der Glaube ist, da sind auch solche Nachbarn, Knechte, Weiber und Kinder. Darum müssen ja zu unserer Zeit derseelbigen auch etliche sein. Denn wie der Regen nicht um eines Halmen, sondern um vieler Aecker willen fällt, daß sie fruchtbar werden und viel reiche Früchte bringen sollen: also wird auch das Wort nicht darum gegeben, daß nur Einer oder Zwei dadurch sollen bekehrt werden, sondern um vieler Menschen willen in allen Ständen. Und ihr sollt auch nicht denken, wenn ich und D. Pommer nicht werden da sein, daß eben dann auch alle gottseeligen und frommen Männer gestorben sein werden: unser HErr GOtt hat mehr fromme Leute. Es sind aber doch die Gottseeligen in der Welt eitel Wunderwerke und Gaben GOttes; denn der größte Theil ist gottlos.

\quellseite{207}

\section[Vierter Theil]{Vierter Theil. Von dem Kampfe Jakobs}

\begin{quote}
\textbf{V. 24. Da rang ein Mann mit ihm, bis die Morgenröthe anbrach.}
\end{quote}

125. Es wird von jedermann dafür gehalten, daß dieser Text einer der dunkelsten im Alten Testamente sei. Und das ist kein Wunder; denn es wird darin gehandelt von der hohen schweren Anfechtung, da der Patriarch Jakob streiten muß nicht mit Fleisch und Blut, auch nicht mit dem Teufel, sondern wider GOtt selbst. Das ist aber ein greulicher Kampf, wenn GOtt selbst kämpft und feindlicher Weise wider den Menschen streitet, als wollte er ihm das Leben nehmen. Wer in solchem Kampfe bestehen und gewinnen will, der muß wahrlich ein heiliger Mann und ein rechter Christ sein. Derhalben ist nun diese Historie dunkel, dieweil die Sache, davon gehandelt wird, so groß ist. Und darum daß der Text so dunkel ist, sind alle Ausleger vorüber gegangen. Derhalben möchten wir ihn auch überhüpfen; wir wollen aber doch davon sagen, so viel wir können.

126. Lyra hat diesen Kampf, so gut er konnte, dargestellt; er hat aber nicht Alles recht erklären können. Augustinus dichtet eine Allegorie oder heimliche Deutung; aber man soll in der heiligen Schrift nicht bloße Allegorien suchen. Denn wo sie nicht eine Historie und gewisse Sache, davon gehandelt wird, zum Fundament haben, so sind es Anderes, nichts denn Aesopische Fabeln. Darnach wenn wir schon die Historie haben, so ist es doch nicht einem Jeden gegeben, Allegorien daraus zu machen. Origenes hat kein gut Glück damit, Hieronymus noch weniger; denn sie haben keinen vollkommenen Verstand der Historien, ohne welchen Verstand niemand mit den Allegorien recht umgehen wird. Darum muß man vor allen Dingen den historischen Sinn herausziehen, der lehrt, tröstet und stärkt; darnach kommt die Allegorie, die ziert und erklärt solchen Verstand wie ein Zeuge; die Historie aber ist wie der Meister, oder Hauptstück und Fundament der ganzen Sache.

127. So wollen wir nun versuchen, ob wir den rechten Verstand dieses Textes und die Lehre desselbigen herausfinden können. Und wo wir ja in denselbigen Verstand nicht vollkommen erreichen werden, so wollen wir doch nicht weit fehlen. Zum Ersten aber muß man die Hindernisse so mancherlei Meinungen beiseite schaffen.

128. Der Prophet Hosea Cap. 12. V. 4. 5. führt diesen Text an, da er sagt: „Er hat von allen Kräften mit GOtt gekämpfet. Er kämpfte mit dem Engel und siegete, denn er weinete und bat ihn; daselbst hat er ihn ja zu Bethel funden“ etc. Daraus schließen alle Ausleger, daß dieser Mann, der mit Jakob ringt, ein Engel gewesen sei.

129. Wir halten aber diese Regel für und für, nämlich, so oft die heilige Schrift der Geschichte oder Erzählungen gedenkt (wie droben Cap. 18. V. 8. ff. dieses Buches, da drei Männer vor Abraham gestanden sind) und der Name „Engel“ nicht ausdrücklich gesetzt wird, da deuten wir es nicht als Engel. Denn es wird klärlich im Text gesagt: „Und der Herr erschien ihm“ etc., und nicht ein Engel. Hernach aber Cap. 19. 1., da zwei Engel gen Sodom kamen, da geben wir zu, daß es Engel gewesen seien.

130. Wie Cap. 28. V. 12. 13., da Jakob die Engel hat sehen auf und nieder steigen, und der Herr steht oben auf der Leiter, da verstehen wir den Herrn nicht als einen Engel, wie die, so auf und nieder steigen, ausdrücklich Engel genannt werden; sondern wir verstehen es von dem Sohne Gottes, der da soll ein Mensch werden, nämlich, von unserm Herrn Jesu Christo, der wahrhaftiger GOtt und wahrhaftiger Mensch ist, und wir sagen per communicationem idiomatum (vermöge der Mittheilung der Eigenschaften): Der Mensch ist droben und über alle Creatur und GOtt ist der Unterste. Welches das Geheimniß ist, „welches auch die Engel gelüstet zu schauen“, 1 Pet. 1, 12. Denn um der Einigkeit willen der Person in Christo sehen sie GOtt hienieden und dagegen den Menschen droben. Und darum sagen wir auch von Christo: Der Mensch sitzt zur Rechten GOttes des Vaters; item, GOtt ist niedergefahren zur Hölle und aufgefahren gen Himmel. Dieses nennt man communicationem idiomatum (die Mittheilung der Eigenschaften), davon wir droben weiter gesagt haben.

131. Wo nun der Name „Engel“ nicht ausdrücklich gebraucht wird, da verstehen wir es auch nicht von den Engeln. Und an diesem Orte wird klärlich gesagt: „Du hast mit GOtt gekämpfet und bist obgelegen.“ Item, hernach V. 30. wird folgen, daß Jakob sagt: „Ich habe GOtt von Angesicht gesehen.“

132. Den Text aber in Propheten Hosea legen wir also aus, daß die Worte, so der Prophet erzählt, solche Worte sind, so er Andern nachredet: nicht daß er seine Lehre und Meinung von diesem Exempel damit angezeigt haben wolle, sondern daß er die falschen Propheten schelten und strafen will, die solche Worte führten und mit dem Exempel des Patriarchen Jakob ihre Abgötterei gegen den Tempel zu Jerusalem und die göttliche Ordnung des Gottesdienstes vertheidigten. Denn sie werden gesagt haben: Es ist nicht vonnöthen, daß wir zu Jerusalem opfern, wie es die Priester fordern und haben wollen um ihres Geizes willen, und verdammen unsere Opfer, als die unrecht und gottlos sein sollen, darum daß wir das Volk von ihrem Gottesdienste abführen und ihnen damit an ihren Einkünften viel abgeht. Also haben sie die rechten Lehrer und gottseeligen Priester, welche über dem wahren Gottesdienst hielten, verleumdet, und das Zeugniß dieser Historie auf ihre Seite gezogen und auf sich gedeutet, sintemal GOtt an demselbigen Ort zu Bethel oder Pniel mit Jakob geredet, ihm dazu die Leiter und Engel gezeigt hätte, und daß er mit GOtt gekämpft, und der Engel geweint und ihn gebeten hätte; denn dieses Letzte haben sie selbst dazu gedichtet.

133. Solche Predigten haben die falschen Propheten gethan, und dieselben zu widerlegen haben sich die Propheten sehr bemüht, vornehmlich Hosea; denn alle seine Predigten gehen dahin, daß er damit die Abgötterei, so im Königreich Israel war, strafete, darum daß sie sich selbst eigenen Gottesdienst und sonderliche Stätte erwähleten, da vor Zeiten GOtt mit den Vätern geredet hatte, da die Engel erschienen waren, oder da sonst um anderer Ursachen willen die Stätten berühmt waren. Nun hatte aber GOtt durch Mose hart verboten, daß sich niemand unterstehen sollte, neue Gottesdienste zu erdichten. Als, da er sagt 5 Mos. 12, 4. 5. 6.: „Ihr sollt dem HErrn, eurem GOtt, nicht also thun; sondern an dem Orte, den der HErr, euer GOtt, erwählet wird, aus allen euren Stämmen, daß er seinen Namen daselbst lässet wohnen, sollt ihr forschen, und dahin kommen; und eure Brandopfer und eure anderen Opfer“ etc. dahin bringen.“ Und dazumal war der Tempel zu Jerusalem zum Gottesdienste verordnet; daselbst war die Lade des Bundes, und GOtt hatte verheißen, er wolle von dem Gnadenstuhl aus reden, auf daß er das Volk von den mancherlei Götzendiensten und zu dem einigen rechten Gottesdienste bringen möchte, dazu er sie alle versammeln wollte, wo nicht leiblich und mit äußerlichen Werken, so doch mit dem Herzen und mit rechter Anrufung, wo Etliche ihre Wohnung von demselbigen Ort etwas weit hätten.

134. Wider diesen Verstand aber stritten die falschen Propheten und sagten: GOtt ist allenthalben, darum kann man ihm auch an allen Orten dienen und ihn anbeten, beide zu Bethel und auch sonst auf einem jeglichen Berge. Derhalben sahen sie nicht auf GOttes Gebot. Denn wo GOtt eine sonderliche Weise und auch einen gewissen Ort bestimmt zu seinem Gottesdienste, so muß man nicht fragen: An welchem Ort ich nur GOtt diene, so wird es ihm wohlgefallen, wenn ich daselbst allein gottseelig und andächtig thun werde; oder: An welchem Ort mich es gelüsten wird, da will ich opfern. Wie der Prophet Jesaia im 57. Cap. V. 5. solche Unsinnigkeit hart straft: „Ihr laufet in der Brunst zu den Götzen unter alle grüne Bäume“ etc.

135. Und auf solche Weise pflegen heutiges Tages Türken und Juden auch zu reden, nämlich, daß sie außerhalb der Einigkeit des Glaubens und der christlichen Kirche GOtt wohl dienen können. Mahomet sagt, ein Jeder könne in seiner Religion wohl selig werden, wenn er nur bete, Almosen gebe, und andere gute Werke thue, so sei es nicht vonnöthen, daß er ein Christ sei oder daß er sich in der Einigkeit Christi und seiner Kirche finden lasse. Also ist es im Papstthum eben auch zugegangen, da alle Winkel voll Capellen, Klöster und allerlei Abgötterei gewesen sind.

136. Darum schreit nun Hosea dawider und spricht also: Dieses Exempel Jakobs bestätigt mit nichten eure Abgötterei: er hat zwar mit dem Engel gekämpft, ihr sollet aber daran denken, daß der Herr, der Gott der Heerschaaren, durch einen Propheten das Volk Israel aus Egypten geführt hatte. Wo ist aber nun der Gott zu finden und anzutreffen? Antwort: Da er das Gedächtnis seines Namens aufgerichtet hatte. Wo ist denn dieses Gedächtnis? Antwort: Da ist es, da auch die Lade des Bundes ist.

\quellseite{208}

137. Das ist wohl wahr, daß Gott nicht gebunden ist, weder an Jerusalem, noch an irgend einen andern Ort; und daß er auch an andern Orten helfen könne, das wird niemand leugnen: versuche es aber, und sehe, was du kriegen wirst. Wo du aber nach deinem eigenen Willen und Gutdünken Gottesdienste errichten wirst, so wirst du das ohne Gefahr und Gottes Zorn nicht thun. Gott könnte wohl durch seine Allmacht das menschliche Geschlecht selig machen ohne Christum, ohne die Taufe, ohne das Wort des Evangeliums; er hätte inwendig die Herzen durch den Heiligen Geist erleuchten und die Sünde vergeben können ohne das Predigtamt und ohne Kirchendiener: er hat es aber nicht thun wollen. Und Gott hat mit großem Ernst verboten alle selbsterdachten Gottesdienste.

138. Wenn nun die Heuchler sagen: Alles, was aus guter Meinung geschieht, gefällt Gott wohl; so soll man solche selbstverwandte Andachten verdammen, und soll die Leute ermahnen, daß sie dahin sehen, da sich Gott geoffenbart hat. Man soll nicht sagen: Paulus hat zu Rom gepredigt, darum ist daselbst Christus; St. Jakob liegt zu Compostel in Spanien begraben, derhalben soll man an demselbigen Orte Gott dienen mit Anrufung St. Jakobs. Also soll man mit nichten sagen. Es sind aber, sagen sie, daselbst viele Wunderwerke geschehen? Antwort: Gott hat einen Greuel und verdammt alle solche wilde ausschweifende Gedanken der Menschen außerhalb der einigen Offenbarung, so im Wort und Sacramenten geschehen ist, in welche er uns hat versammeln und einschließen wollen. Denn darum hat Christus seine Jünger ausgesandt mit diesem Befehle Matth. 28, 19. 20.: „Gehet hin und lehret alle Völker und taufet sie im Namen des Vaters und des Sohnes und des Heiligen Geistes. Und lehret sie halten alles, was ich euch befohlen habe.“ „Wer da glaubet und getauft wird, der wird selig werden; wer aber nicht glaubet, der wird verdammt werden“, Marc. 16, 16. Er will, daß wir uns zum Wort und zur Taufe halten, als einem gewissen und unfehlbaren Wahrzeichen, daß er uns helfen und selig machen wolle. Wie er im Volk Israel verheißen hatte, daß er sie bei dem Gnadenstuhle erhören wollte.

139. Demnach wenn du willst von deinen Sünden absolvirt sein, so gehe zum Pfarrherrn, oder zu deinem Bruder und Nächsten, wenn du ja den Pfarrherrn nicht haben oder hören kannst: der hat Befehl, daß er dich absolvieren und trösten soll. Du sollst dir keine sonderliche Absolution dichten. Wenn du willst das Abendmahl des Herrn gebrauchen, so gehe in die gemeine Versammlung der Kirche und gebrauche es daselbst, und erdichte dir keine sonderliche Verwaltung der Sacramente. Denn Gott will nicht haben, daß wir umher in der Irre gehen sollen mit eigenen und selbstverwählten Werken oder Gedanken. Und darum hat er uns in die Grenzen seines Worts beschlossen und versammelt, auf daß wir uns nicht mögen und wiegen lassen von allerlei Wind der Lehre, Eph. 4, 14.

140. Denn also ist es uns im Papstthum gegangen, da wir die Taufe, das Abendmahl des Herrn, Vergebung der Sünde verachtet haben und sind hiebei zu St. Jakob, gen Rom und Jerusalem wallfahrten gegangen, als ob Gott nicht in allen Kirchen und gemeinen Versammlungen wäre, die sein Wort und Sacrament haben. Endlich sind über die Maßen viele Secten und Orden gewesen, da ein jeder Orden seine eigene und sonderliche Weise gehabt hat, anders als die gemeine Weise und Ordnung Gottes ist. Und solche unzählige Formen und Weisen der Hurerei, wie die Schrift davon redet, kann man nicht verhüten, es sei denn, daß wir unsere Werke fallen lassen und verwerfen. Darum will Gott haben, daß man sie alle wegthun soll, und hält uns sein Wort vor, das da sagt: An diesem Orte sollst du anbeten, mir dienen und opfern. Siehe, da hast du im Wort, im Abendmahl, in der Taufe Vergebung der Sünden, daran wirst du dich müssen halten und genügen lassen, so du willst selig werden.

141. Eben auf solche Weise haben die Propheten den Juden mit höchstem Fleiß eingepräget und vorgehalten den Ort, welchen Gott selbst erwählt hatte wider die selbst erfundenen Wallfahrten zu den verschiedenen Götzen und Opfern. Aber ihre Zuhörer sind übel mit ihnen umgegangen, haben sie verjagt und erwürgt. Wie heutigen Tages der Papst auch verdammt, verbannt und erwürgt diejenigen, so die schändlichen Lüste und Hurerei der großen Abgötterei strafen. Aber nichtsdestoweniger nehmen gleichwohl die, so recht gottselig sind, das Licht des Worts mit dankbarem Herzen an, und begeben sich wie die Küchlein unter die Flügel unserer Henne, des Herrn Jesu Christi, des Sohnes Gottes, und da finden sie unter derselben Flügeln Ruhe und Schutz wider alles Anlaufen und hitzige Anschläge des Satans, Eph. 6, 11. Dies sei gesagt von dem Text des Propheten Hosea, den man nicht also verstehen soll, als ob der Prophet aus seiner Meinung also geredet habe, sondern daß er es den Andern also nachgeredet hat.

142. Man soll auch der Juden Glosse verwerfen, die da dichten, daß der Mann, so mit Jakob gerungen hat, ein guter Engel gewesen sei, der auf Esau gewartet und ihn bewahrt hat; derselbe habe sich wider Jakob gelegt und mit ihm gekämpft, wie er denn Amtes halben schuldig gewesen für Esau zu streiten. Da aber der Engel überwunden worden, soll er gemeint und gebeten haben, daß er ihn verlassen wollte. Aber dieses lose Gedicht hat viel ungereimtes Dinges in sich. Sie lassen sich träumen, der Engel habe nicht gewußt, daß der Segen Jakob gehört hat. Item, daß Gott Jakob dermaßen gestärkt hätte, daß er den Engel überwinden möchte, und dagegen daß der Engel sei geschwächt worden und derhalben gemeint habe. Das ist eine närrische und frivole Einbildung, nämlich, daß man dichtet, als sollte der Engel von dem verheißenden Segen nichts gewußt haben, oder daß er so böse sollte gewesen sein, daß er der göttlichen Ordnung hätte widerstreben wollen. Lyra hat es etwas besser getroffen: der macht aus diesem Engel einen Mann, aber in der Person Gottes, und daß er habe fliehen wollen, aber er sei von Jakob aufgehalten worden, bis er ihn gesegnet hat.

143. Dieses ist aber unsere Meinung, daß dieser Kämpfer der Herr der Ehren sei, nämlich unser Herr Gott selbst oder der Sohn Gottes, der da sollte Mensch werden, der den Vätern erschienen ist und mit ihnen geredet hat. Denn Gott geht mit seinem auserwählten Patriarchen Jakob nach seiner großen Güte ganz freundlich um und versucht ihn, als spiele er mit ihm auf das allerfreundlichste. Aber dieses Spiel ist ihm ein unermeßlicher Schmerz und sehr große Angst in seinem Herzen, und ist doch gleichwohl in Wahrheit ein Spiel, wie das Ende solches sein ausweiset, wenn er nach Pniel kommen wird. Denn da wird es offenbar werden, daß es nur lauter Zeichen einer sehr vertrauten Liebe gewesen sind. Und er spielt mit ihm also, darum daß er seinen Glauben damit prüfen und stärken will. Wie irgend ein frommer Vater seinem Kindlein einen Apfel nimmt, dazu das Kind Lust gehabt, nicht daß es behalten von dem Vater wegnehmen und ihm sein werden soll, sondern daß es dadurch mehr gereizt werde, den Vater zu herzen, zu bitten und zu flehen: Ach mein Vater! gib mir wieder, was du mir genommen hast. Da hat der Vater Lust zu solcher Versuchung, und das Kindlein, wenn es den Apfel wieder empfängt, gewinnt es den Vater immer desto lieber, wenn es sieht, daß der Vater zu solcher Liebe und Kinderspiel Lust hat.

144. Dieses Spiel ist in der Haushaltung sehr gemein, aber in Sachen und im Kampfe der Heiligen ist es sehr schwer und verdrießlich. Denn Jakob weiß nicht, wer der ist, der mit ihm kämpft; er weiß nicht, daß es Gott ist. Denn er fragt ihn hernach, wie er heiße. Da er aber den Segen empfängt, spricht er: „Ich habe Gott von Angesicht gesehen.“ Da ist eine neue Freude und Leben geworden aus dem Tode selbst.

145. So dünkt mich nun dies die Lehre zu sein aus dieser Historie, wenn ich sie nur nach ihrer Würde recht erklären könnte, daß Gott mit seinen Heiligen, nach dem Exempel Jakobs, bisweilen pfleget zu spielen, so viel ihm belanget, gar mit einem Kinderspiele; uns aber, die er auf solche Weise versucht, scheint es viel anders. Es ist aber eine sehr gute und heilsame Uebung, dazu die allervollkommenste Unterweisung, die das aller seligste Ende gewinnet, nämlich, daß man erkennen möge, welches da sei der gute, der wohlgefällige und der vollkommene Wille Gottes. Denn Fleisch kann zwar anders nicht denken, denn daß es ein böser, verdrießlicher und betrübter Wille sei: aber alsdann lacht Gott freundlich, wenn wir weinen, und er „hat Wohlgefallen an denen, die ihn fürchten und auf seine Güte hoffen“, Ps. 147, 11.

\quellseite{209}

146. Die Anfechtung aber der Verzweiflung, die da pfleget nicht unterzulassen, macht den Schmerz und Schrecken des Fleisches immer größer, nämlich, wenn ein solch betrübt Herz klagt, daß es von Gott verlassen und verworfen sei. Das ist die letzte und auch die allerschwerste Anfechtung des Unglaubens und der Verzweiflung, damit die allergrößten Heiligen pflegen versucht zu werden. Und wer da bestehen und beharren kann, der kommt zur vollkommenen Erkenntniß des Willens Gottes, daß er mit Jakob sagen kann: „Ich habe Gott von Angesicht gesehen“ etc. Ich glaube nicht, daß es unser Herr Gott so gut mit mir meinte; aber ehe er dahin kommen, wird es uns sauer. Darum ist nun die Lehre in dieser Historie offenbar und klar, nämlich, von den Anfechtungen der größten Heiligen, die mit großer Süßigkeit schmecken, wie freundlich der Herr sei, Ps. 34, 9. Ob nun wohl jedermann diesen schweren Kampf nicht fassen oder verstehen kann, so sind doch solche Leute deswegen nicht zu verwerfen.

147. Von den Erscheinungen aber ist etlichemal insgemein gesagt worden. Und wir sollen wissen, daß man nicht allen Erscheinungen glauben soll, es sei denn, daß sie dem Glauben ähnlich seien. Ich will das Wort Gottes behalten und damit zufrieden sein, damit will ich sterben, damit will ich leben. Es ist süße und Trost genug in der Verheißung Gottes nicht allein wider den Teufel, das Fleisch und die Welt, sondern auch wider diese hohe Anfechtung. Denn wenn Gott einen Engel senden würde, der da sagte: Du sollst dieser Verheißung nicht glauben, so wollte ich ihn von mir stoßen und sagen: „Heb dich weg von mir, Satan“, Matth. 4, 10. Oder wenn mir Gott selbst erscheinen würde in seiner Majestät, und also sagte: Du bist meiner Gnade nicht werth, ich will meinen Rath ändern und will dir die Verheißung nicht halten: da sollte ich nicht weichen, sondern sollte wider Gott selbst auf das allerhärteste streiten. Wie Hiob Cap. 13, 15. sagt: „Wenn er mich schon tödten würde, so will ich doch auf ihn hoffen.“ Wenn er mich auch in die Tiefe der Hölle würde werfen und mich mitten unter die Teufel setzen, so glaube ich doch, daß ich werde selig werden. Denn ich bin getauft, bin absolviret; ich habe empfangen das Pfand meiner Seligkeit, nämlich, den wahren Leib und Blut des Herrn Christi Abendmahl. Derselben will ich sonst nichts sehen oder hören, sondern will in dem Glauben leben und sterben, es sage mir gleich Gott, oder ein Engel, oder aber der Teufel, das dem zuwider ist.

148. Dieses, halte ich, sei die rechte Erklärung dieses Textes. Denn also wird es Jakob am Ende selbst auslegen, da er v. 30. sagen wird: „Und meine Seele ist genesen.“ So ist nun dies Bild des Streites und des Kampfes in den Heiligen sehr voller Trost. Wie an einem andern Orte ein Exempel erzählt wird von einer Nonne, Mechthilde genannt. Dieselbe ward vom Teufel geplagt, daß sie vom Glauben gar nichts wußte oder fühlte. Das war eine Anfechtung des Unglaubens, welches ein sehr harter Schmertz und Stein des Gewissens ist. Denn die Herzen werden vom Schrecken und Zweifel verzehrt. Und allein die Erfahrung zeigt an, was es für ein Schmerz sei, mit Worten kann man es nicht erklären. Jedoch war dieselbige Anfechtung diesem Kampf, so Jakob gehalten hat, noch nicht gleich; denn Gott hat noch nicht wider sie gestritten, wie hier gesehen ist, sondern der Teufel, der die Menschen treiben kann zum Mißglauben, Verzweiflung und zur Lästerung wider den Heiligen Geist. Wie ich selbst etliche gesehen habe, und sonderlich Weiber, die geradezu darüber klagten, daß sie verdammt und verworfen wären. Denn sie wurden geplagt vom Satan mit dem Geist der Lästerung.

149. So hat sich nun die Nonne wider diese Anfechtung mit nichts Anderem gewehret, denn mit diesen Worten, daß sie sagte: Christiana sum: „Ich bin eine Christin, das ist, ich bin getauft mit dem Blute des Sohnes Gottes, ich bin gespeiset mit dem Leib und Blut Christi, daran halte ich mich fest, mit dem Troste bin ich zufrieden, wenn schon Gott selbst anders sagen würde.“ Wie Abraham geboten ward, daß er den Sohn, in dem ihm die Verheißung gegeben war, tödten sollte, und doch gleichwohl auf das festeste glaubte, daß Gott seine Verheißung halten würde.

150. Dieses muß man fleißig treiben um derer willen, die etwa noch Pfarrherren werden sollen. Denn es werden allezeit Etliche sein, die diese Anfechtung haben werden, daß dieselbigen durch der Pfarrherren Wort aufgerichtet und gestärkt werden mögen. Als nämlich auf diese Weise, daß sie zu ihnen sagen: Sei getrost, mein Sohn, glaube, daß du getauft bist, daß du mit dem Abendmahl des Herrn geweidet und gespeiset, dazu auch absolviret bist, mit Auflegung nicht meiner, sondern Gottes Hände, der zu dir gesagt hat: Ich vergebe dir deine Sünde und verheiß dir das ewige Leben. Wenn sie das mit festem Glauben fassen und annehmen, so werden die Anfechtungen und Lästerungen des Geistes verschwinden.

151. Denn Abraham hat ohne Zweifel auch von ganzem Herzen groß Schrecken und Zagen gefühlt, da ihm geboten war, daß er seinen Sohn tödten sollte. Unterdeß aber hat er diesen Glauben behalten, dadurch er gewiß geschlossen hat, wenn schon sein Sohn zu Asche würde, so würde er doch gleichwohl ein Vater vieler Erben und Nachkommen werden laut der göttlichen Verheißung. Wie, ist denn Gott ein Lügner geworden, wird er denn seine Verheißung nicht halten? Solche Gedanken sind ihm eingefallen. Und er gehorcht doch gleichwohl dem Befehl und denkt also: Wie der Herr diesen meinen Sohn vom Leibe einer veralteten und unfruchtbaren Mutter und aus meinen Lenden hat kommen lassen, also wird er ihn auch wiederum erwecken können, wenn er schon verbrannt und zu Asche geworden ist. Auf solche Weise, sage ich, hat Abraham auch mit Gott gekämpft, welches der allergrößten und schwersten Kämpfe einer ist. Aber Gott zeigt darin an seinen vollkommensten und höchsten Willen.

152. Wiewohl nun dieser Kampf nicht verstanden oder ertragen werden kann, denn allein von den Heiligen: so muß man doch diese Lehre und Trost haben, uns damit zu stärken, daß wir vom Teufel nicht verschlungen werden; wiewohl Gott getreu ist, der uns nicht läßt versuchen über unser Vermögen, 1 Cor. 10, 13. Denn das lehrt uns dieses Exempel Jacobs, der zu diesem Kampf sehr schwach war, und wird doch gleichwohl nicht überwunden. Es hält sich aber Gott gegen ihn also, daß er es nicht erkennen kann, daß Gott der Kämpfer sei; er meinet, es sei ein Engelchen. Aber es ist Gott, der sich vernehmen läßt, daß er sein Widersacher sei, gleich als wollte er ihn tödten, der Verheißung und des Segens berauben und denselbigen seinem Bruder Esau geben. Und niemand kann mit Worten erreichen, was er werde für Gedanken gehabt haben. Aber solche Gedanken werden ihm ohne Zweifel eingefallen sein: Was bin ich denn für ein armer elender Mensch? Bin ich denn nur dazu geschaffen, daß ich immer Unglück haben soll? Muß ich denn nur immer ein Unglück über das andere haben, und damit also geplagt werden, daß ich nimmer zur Ruhe kommen kann? Ist doch kein elenderer Mensch auf Erden, denn ich bin. Ich sehe, daß mein Bruder Esau herrscht, triumphirt, zunimmt und groß wird mit großer Herrlichkeit, mit großem Gut, mit Kindern, Kindeskindern und mit großem Einkommen. Wie, wenn unser Herr Gott wäre andern Raths geworden, mich verworfen, meinen Bruder aber zu Gnaden genommen hätte?

153. Dieses sind Jacobs Gedanken gewesen; es sind aber doch allein Gedanken geblieben. Denn derselbigen kann sich die Natur und der schwache Glaube nicht enthalten, wie sie auch andere Affecten und Bewegungen der Ungeduld, des Zornes und böser Lust nicht leichtlich ablegen kann. Aber man lasse es nur Gedanken bleiben, daß es keine gewissen Sprüche werden, die endlich bestätiget werden durch unser Urtheil und Gewissen. Ich kann mich dessen nicht erwehren, daß mein Herz nicht sollte mit wunderlichen Gedanken und Anfechtungen bekümmert und geplaget werden.

154. Derhalben muß man dem Rath jenes Eremiten oder Einsiedlers folgen, vor dem sich ein Jüngling beklagte, daß ihm oftmals böse Gedanken einfielen von Unzucht und von andern Sünden. Da hat ihm der alte Mann geantwortet: Du kannst den Vögeln nicht verbieten, daß sie über den Kopf fliegen; aber laß sie nur immerhin fliegen und laß sie nicht nisten in den Haaren deines Haupts. Laß es Gedanken sein und Gedanken bleiben: aber laß keine feste Ueberzeugung daraus werden. Denn dieses ist die Art derer, die da verzweifeln, als Sauls, Judas und anderer, die aus den Gedanken Schlüsse machen: sie sagen: „Meine Sünde ist größer, denn daß sie mir vergeben werden möge“, 1. Mos. 4, 13.; „Ich habe unschuldig Blut verrathen“, Matth. 27, 4. Dieselbigen verwerfen das Wort, den Glauben und das Gebet. Das heißt aus der Anfechtung solche Urtheile machen, wie im Gericht gesprochen werden. Jakob aber läßt den Glauben nicht fallen, wiewohl er sehr schwere Anfechtung und viel unzählige Gedanken gehabt hat, welche diejenigen nicht verstehen, die solche Dinge nicht selbst erfahren haben.

\quellseite{210}

155. So ist nun der vornehmste Verstand und Meinung dieser Historie das Exempel der vollkommenen Heiligen, und der hohen Anfechtungen, nicht wider Fleisch und Blut, wider den Teufel, oder auch irgend wider einen guten Engel, sondern wider GOtt selbst, der da erscheint in feindseliger Gestalt. Denn wiewohl Jakob nicht weiß, wer dieser Mann ist, so fühlt er doch das wohl, daß er von GOtt verlassen sei, oder daß GOtt ihm zuwider sei und mit ihm zürne.

156. Also behalten wir den Verstand, so die Grammatik oder Historie gibt, und wo man denselbigen hat, so finden sich die Allegorien fein von sich selbst. Wir bleiben einfach bei den Worten im Text, nämlich, daß es eine Gestalt eines Mannes und kein Engel gewesen sei. Denn wie es Jakob erzählt hat, also hat Mose die Historie auch beschrieben, nämlich, daß Jakob sagt: Da hat mich plötzlich ein Mann angegriffen und mit zugesezt. Das ist also eine Gestalt oder Schein gewesen. Er disputirt aber davon nicht, wer der Mann gewesen sei; denn er weiß davon nichts. Aber später, da er ihm hinten nach sieht, da versteht er es, und sagt: „Ich habe GOtt von Angesicht gesehen.“ Darum ist dies kein anderer Angriff gewesen, denn als hätte ihn sonst irgend ein großer starker Mann angegriffen; wer aber derselbe sein möchte, wußte er nicht.

157. Und Jakob war auch von gehöriger Stärke und vollkommenen Kräften, dergleichen in einem gesunden starken Leibe zu sein pflegen. Dieweil er nun ohne Schwert oder sonst andere leibliche Wehre kämpft, hat er ihn hart widerstanden, und er hatte noch etlichermaßen Glauben an die Verheißung. Wiewohl er erschrocken war und zitterte, so waren doch gleichwohl seine Gedanken noch nicht feste und gewisse Ueberzeugung geworden. Sein Glaube ist zwar hart angefochten und versucht worden, er hat aber fest gehalten an diesem Schluß: Ich habe die Verheißung. Darum tröste er sich, und versagt unserm Herrn GOtt den Kampf, so er ihm anbietet, nicht, sondern wagt es, denkt: Ermüret er mich, so ermüret er mich.

158. Was aber die Angst seines Herzens gewesen sei mitten im Kampf, das wird nicht beschrieben; wir können aber wohl abnehmen, daß die Anfechtung ein Kampf des Glaubens gewesen sei wider den Mißglauben. Denn der Mann hat ihm sein Herz stürmen und überwinden wollen unter dem Namen Gottes, daß er alle Hoffnung der Barmherzigkeit und Gnade Gottes sollte fallen lassen. Wie der Satan Christum auch angreift Matth. 4, 9., da er zu ihm sagt: „Das alles will ich dir geben, so du niederfällst und mich anbetest.“ Da verstellt sich der Teufel in die Gestalt und Majestät GOttes. Und ist die höchste Anfechtung, die man nur fühlen kann, mit Worten aber kann man sie nicht ausreden. Und ist es kein Wunder, daß der Mensch in solchem schweren Kampf mit seiner Vernunft nichts schließen oder rathen kann. Denn da kämpft der Geist selbst, und der Heilige Geist ist dabei, der hilft unserer Schwachheit; und wo der nicht dabei wäre, so würde die Anfechtung anders hinaus gehen, denn der Mensch bei sich beschlossen hat.

159. Es ist aber wohl glaublich, daß der Mann weiter zu Jakob also mag gesagt haben: Jakob, du wirst sterben müssen; denn du bist der nicht, dem GOtt die Verheißung gegeben hat, oder dem er sie auch halten wollte etc. Das sind feurige Pfeile gewesen. Und dieser heilige Patriarch war ja eigentlich der, welchem die Verheißung und der Segen gehörte. Er wird aber hier angefochten auf das äußerste Fünklein, bis auf das allerkleinste Docht, Jes. 42, 3., es glimmt aber noch und ist noch nicht erloschen, und ist noch kein zerbrochen Rohr.

160. Dieses sind aber höhere Dinge an den Heiligen, als wir verstehen können; wiewohl man es gleichwohl lesen und hören soll. Mose hat eben die Worte so eigentlich nicht ausgedrückt, die sie mit einander gewechselt haben, dieweil der Kampf gewähret; es will sich aber doch nicht reimen, daß sie sollten gar stille geschwiegen haben, sondern der Mann wird ohne Zweifel mit erschrecklicher Stimme sich haben hören lassen: Jakob, du mußt herhalten. Darauf wird Jakob gesagt haben: Nein, das wollt GOtt nicht; ich werde nicht also umkommen. Ja und Nein ist da auf das allerschärfste und heftigste auf einander gegangen. Das kann man mit dem Munde und mit Worten nicht genugsam ausdrücken, vornehmlich, wenn GOtt selbst selbst sagt: Du wirst sterben müssen; und der Geist dawider schreiet: „Ich werde nicht sterben, sondern leben“ etc. Ps. 118. V. 17., und V. 13. „Man stößet mich, daß ich fallen soll; aber der HErr hilft mir“ etc. Ich mag getrieben, gestoßen und angegriffen werden, daß ich fallen soll, gleichwohl werde ich nicht sterben. Also haben sie mit einander gekämpft mit den Armen und auch zugleich mit Worten, wie es pflegt zu gehen, wo Zwei mit einander kämpfen. Unterdeß aber hat auch der Glaube gekämpft mit Beten und Schreien; wie Mose am Rothen Meer geschrieen hat.

161. Dieses ist der Höhepunkt des Kampfes gewesen, darin sich der Glaube mehr bemüht hat denn die Arme, daß er immer angehalten, und dieses wiederholt und gesagt hat: Nein, nein, ich werde nicht also sterben müssen, GOtt hat mich es geheißen, mich dazu berufen und gesandt, daß ich wieder in mein Vaterland ziehen soll; ich will dir nicht glauben und auch nicht beifallen. Und wenn mich GOtt schon tödten wird, so mag er mich immerhin tödten, so will ich aber doch gleichwohl leben. Auf solche Weise werden sie andere und aber andere Worte gebraucht haben die zwei Stunden; und das werden ohne Zweifel scharfe und heftige Reden gewesen sein. Jakob hat nicht gewußt, wer der Mann wäre, und die Vernunft hat darin nicht schließen oder rathen können. Denn in solchem Kampf sind alle Sinne verwirrt und die Vernunft ist auch nicht bei sich selbst. Wie im 107. Psalm V. 27. steht: „Sie wußten keinen Rath mehr.“ Jakob hat nicht gewußt, wie ihm geschehen sei, allein, daß der Glaube sagen hört: Der Segen gebührt nicht dir, sondern deinem Bruder Esau, du bist nichts etc.

162. Darum lebt er allein durch das Wort, und sein Glaube ist zugleich sehr schwach und auch sehr stark. Denn wenn das Docht glimmt, so ist es nahe daran, daß es angezündet werde. Und er ist in sehr großer Angst gewesen; denn er wird hernach V. 30. sagen: „Meine Seele ist genesen“, als wollte er sagen: Meine Seele ist in höchster Angst gewesen; jetzt aber habe ich nun ein wenig Luft bekommen. Denn GOtt versucht nicht der Meinung, daß er den Menschen damit tödten wolle, sondern daß er ihn erquicken, bestätigen und stärken will. Dieses kann man aus dem Text selbst abnehmen. Und dieses ist nun der Kampf Jakobs kürzlich also beschrieben. Denselben wollen wir benennen, so darin erfahren sind, befehlen, ihn weiter und mit mehr Worten zu erklären.

163. Das hebräische Wort, welches die siebenzig Dolmetscher gegeben haben, kämpfen oder ringen, kommt her von abak, das heißt, Staub, der erregt ist, daß er in der Luft sich von einander theilt. Es heißt nicht der Staub, der da stille liegt, sondern der erregt wird von denen, die da arbeiten oder kämpfen. Wie im Propheten Jesaia am 5. Cap. V. 24. steht: „Ihre Sprossen fahren auf, wie Staub.“ So heißt nun dieses Wort so viel, als mit Kämpfen den Staub erregen, daß er auffahre. Wiewohl ich es dafür halte, daß es der Worte eines sei, die man nicht allzumal versteht. Denn die hebräische Sprache ist noch nicht vollkommen wiederum aufgerichtet, und die Juden verfälschen viel Dinges mit ihren „aequivocationes“, das ist, mit den Worten, die mehr denn einerlei bedeuten sollen, und sonst mit ihrer Auslegung, vornehmlich in den dunkeln Texten vom Messias.

\begin{quote}
\textbf{V. 25. Und da er sahe, daß er ihn nicht überwältigte, rührte er das Gelenk seiner Hüfte an; und das Gelenk seiner Hüfte ward über dem Ringen mit ihm verrenkt.}
\end{quote}

164. Jakob ist stärker als der Mann, der mit ihm gerungen hat; nicht, daß er in Wahrheit stärker sei denn der HErr, sondern es wird der Vergleich mit dem Manne festgehalten. Denn der HErr hat nicht größere Kraft geübt, als sonst ein Mann üben kann oder pflegt zu üben: er will nur Eines Mannes Kraft und Vermögen versuchen, dieweil er als der HErr eigentlich stärker ist. Und ich wollte wahrlich nicht gern mit einem gläubigen Manne, der göttliche Verheißung hat, ringen, ob ich gleich ein starker Riese wäre; denn da arbeitet die Natur über ihr Vermögen.

165. Derhalben beweiset GOtt seine Gewalt nicht, sondern nur solche Kunst und Verstand, soviel einem starken Manne geziemt, und bringt zugleich mit sich eine erschreckliche Anfechtung des Geistes; dieselbe stärkt Jakob erst recht, daß er ringt über sein Vermögen.

\quellseite{211}
Denn kein Mensch, wenn er gesund und bei Sinnen ist, kann das thun, das die Natur thut, wenn sie Noth leidet und beschweret ist. Wenn die Natur geschwächt ist und arbeitet bis in den Tod hinein, da kann sie alle Beschwerden überwinden, welche sie sonst nicht zu ertragen vermöchte. Darum pflegen sie zu sagen, wenn das Kämpfen geschehen ist und die harten Stöße vorüber sind, wenn sie wieder zu sich selbst kommen: Ei, wie matt und kraftlos bin ich! Denn es folgt darauf, daß der Mensch schwach und matt wird in allen Kräften, dieweil die Natur über ihr Vermögen gearbeitet hat.

166. Man erzählt sich wohl eine Geschichte oder eine Fabel, nämlich, daß auf einem steinernen Altare fünf Finger sollen abgedrückt sein, daß man das Zeichen davon sehen kann; dasselbe soll Einer gethan haben, der begehrt hat, daß man ihm so eine große Kostie im Abendmahl geben wollte, wie sie der Priester empfinge; und da ihm solches ist versagt worden, soll er entweder erschrocken oder aber erzürnt worden sein und mit der rechten Hand auf den Altar geschlagen haben, und davon soll das Zeichen seiner Hand im Stein abgedrückt für und für auf dem Altar geblieben sein. Denn ein Mensch, der von schweren Anfechtungen, oder aber von großem Schrecken bestürzt ist, thut das wohl, daß ihm sonst außerhalb der Anfechtung zu thun unmöglich ist.

167. So hat nun dieser Mann, welcher Christus, unser Herr, gewesen ist, mit Jakob gerungen eben auf solche Weise, wie ein Mann mit dem andern ringt. Da aber der Herr gesehen hat, daß Jakob stärker war, vornehmlich im Glauben, welchem Gott nicht gern widerstrebt, gebraucht er eine sonderliche Kunst, wie die Kämpfer pflegen, und rührt das Gelenk der Hüfte Jakobs.

168. Das hebräische Wort kaph heißt auf Lateinisch vola, das ist, eine Hand, die wie ein Löffel gekrümmt ist. Daher nennt es Mose kapoth, auf Lateinisch coclearia, auf deutsch Koten (Gelenkknöchel); und um der Vergleichung willen heißt es das Gelenk oder das Bein, das ausgehöhlt ist, das die Griechen ischion heißen, darin das Oberbein eingefügt ist und darin es sich bewegt. Es sind aber etliche Muskeln darum her, daß die zweierlei Knochen zusammengehalten werden und nicht von einander springen. Und durch daselbstige Gelenk regt sich der Fuß. So will er nun sagen, daß ihm der Mann die Hüfte verrenkt habe, daß ihm das Bein aus dem Gelenk gekommen ist. Wo aber das Gelenk also verrenkt ist, so werden die Muskeln auch aus ihrer Lage gerückt und zusammen gezogen, und dieweil so viel Nerven zusammenkommen, so entstehen sehr heftige Schmerzen, und folgt daraus, daß daselbstige Glied lahm wird.

169. Da nun der Mann Jakob nicht überwinden konnte, hat er ihm die Hand an die Hüfte geschlagen, und davon ist die Hüfte mit ihren Muskeln verrenkt worden. Er hat ihm einen harten Druck gegeben und das Bein wollen lahm machen. Das sind der Kämpfer Kunstgriffe. Und es hat mir ein glaubwürdiger Mann zu Eisenach ein ähnliches Exempel erzählt, da ich in Natmos war, nämlich, daß er einen starken Riesen gekannt hätte mit einer rauhen Brust und langem geschorenen Haar, der immer lange Zappen an den Haaren hangen gehabt hat und immer mit bloßem Haupt gegangen ist, der auch Eisen hat pflegen krumm zu machen und zu brechen. Und er sagte, daß zu der Zeit dergleichen viele in der Türkei sollen gefunden worden sein. Es war aber zu derselben Zeit am Hof bei Kaiser Friedrich ein trefflicher Jude und berühmter Kämpfer, der in Kämpfen alle jungen Edelleute am Hof überwunden und erlegt hatte mit eben solcher Kunst, wie in dieser Historie beschrieben wird. Und vielleicht mag er dieselbige Kunst hieraus gelernt haben, daß er das Gelenk der Hüfte angerührt, auf daß er seinen Widersacher mit dem Schmerz, der an dem Orte pflegt sehr heftig zu sein, matt machen und überwinden möchte. Da nun derselbige Riese oder Kämpfer an des Kaisers Hof gekommen ist, und ihn die Hofdiener ermunterten, daß er mit dem Juden ringen sollte; wiewohl er sich vernehmen ließ, daß er mit solchen Leuten nicht wollte zu thun haben, so hat er sich doch endlich durch ihr Bitten überreden lassen und hat es mit dem Juden gewagt. Da hat der Jude gebeten, man wollte ihm zu gut halten, wo er vielleicht den fremden Kämpfer, und der daselbst ein Gast war, etwas hart angreifen würde. Der Riese aber hat ihn dagegen auch ermuntert, daß er alle seine Kunst und Stärke seines Leibes gebrauchen, und sein gar nicht schonen sollte, er wollte beweisen, was er könne, und wie stark und gewandt er sei. Alsbald aber, da sie mit einander zusammentreffen, hat der Riese den Juden ergriffen und ihm das Gelenk im Nacken zerbrochen, und ihn also gebödtet. Dieses sage ich zum Exempel, diese Historie damit zu erklären, daß wir wissen, daß die Kämpfer dergleichen Kunstgriffe viele haben, als nämlich, da sie die Hüfte, das Gelenk und andere Glieder am Leibe anrühren, damit oftmals die, so nicht sonderlich stark sind, Andere, die viel stärker sind als sie, überwinden können.

170. Im Uebrigen hat Jakob in solchem großen Schrecken nicht gefühlt, daß ihm das Gelenk ist angerührt worden. Denn an einem angefochtenen Menschen werden Sinne und Vernunft verwirrt und bestürzt, daß er irgend aufmerken und auch gar nichts hören oder sehen kann. Und das pflegt vornehmlich also zu kommen, wenn auch geistliche Anfechtungen dazutreten, wie solche Anfechtungen Jakob zu der Zeit gefühlt hat. Darum hört er nicht auf mit Ringen, da ihm auch schon das Gelenk seiner Hüfte verrenkt ist.

\section[Fünfter Theil]{Fünfter Theil. Was auf Jakobs Kampf gefolgt}

\begin{quote}
\textbf{V. 26. Und er sprach: Laß mich gehen, denn die Morgenröthe bricht an.}
\end{quote}

171. Wenn dieser Mann, der mit Jakob ringt, Christus ist, was geht ihn denn an, daß die Morgenröthe anbricht? Antwort: Er hat die Person eines Mannes angenommen, dieselbige behält er auch, bis er sich offenbart. Darum ist ihm in solcher Person die Morgenröthe auch angebrochen und muß nach der Weise aller Völker zu seiner gewöhnlichen Arbeit greifen.

\begin{quote}
\textbf{Aber er antwortete: Ich lasse dich nicht, du segnest mich denn.}
\end{quote}

172. Ist doch das Gelenk deiner Hüfte verrenkt, warum willst du ihn nicht lassen? Du bist nun lahm geworden, was willst du nun machen? Ich empfinde keine Schwachheit, spricht Jakob. Lieber, wer stärkt dich denn also? Antwort: Der Glaube, die Verheißung; ja, eben die Schwachheit des Glaubens.

173. Und auf solche Weise wird GOtt überwunden, nämlich, wenn der Glaube nicht aufhört, nicht müde wird, sondern immer anhält und GOtt drängt. Wie an dem cananäischen Weiblein zu sehen ist, mit welcher Christus einen Kampf hält, da er zu ihr sagt: Du bist ein Hund, das Brod der Kinder gebühret dir nicht; da weicht sie gleichwohl nicht, sondern legt sich wider ihn und sagt: „Essen doch die Hündlein von den Brosamlein, die von ihrer Herren Tische fallen“, Matth. 15, 27. Darum gewinnt sie auch und hört ein schön herrlich Lob, das ihr der Herr nachsagt V. 28.: „O Weib, dein Glaube ist groß.“ Solche Exempel lehren uns, daß der Glaube nicht weichen und daß er nicht aufhören soll mit Anhalten, wenn er auch schon den Zorn GOttes fühlt, und nicht allein den Zorn und die Sünde. Und das ist die Kraft und Stärke des Geistes.

174. Dieweil nun Christus mit Jakob noch ringt, seine Ohnmacht verbirgt, und haben will, daß er ihn lassen soll, so antwortet Jakob: „Ich lasse dich nicht, du segnest mich denn.“ Ei, warum das, lieber Jakob? Antwort: Darum daß du gesagt hast, ich sei verflucht und vermaledeyet; derhalben meine Seele betrübt gewesen ist: darum so widerrufe daselbstige Urtheil und segne mich. Du mußt mir einen Widerruf thun, oder ich werde dich nicht lassen. Ich habe dich an der Stärke des Leibes überwunden, ich will auch die Worte deines Mundes überwinden. Denn meine Seele hat sich etwas heftiger bemüht (von welcher du gesagt hast, daß sie verloren und verdammt wäre), denn mein Leib oder meine Arme. Darum will ich dich nicht lassen, es sei denn, daß du mir einen Widerruf thust, und gebest mir deß ein Zeugniß, daß ich vor GOtt gesegnet sei.

\begin{quote}
\textbf{V. 27. 28. Er sprach: Wie heißest du? Er antwortete: Jakob. Er sprach: Du sollst nicht mehr Jakob heißen, sondern Israel. Denn du hast mit GOtt und mit Menschen gekämpft, und bist obgelegen.}
\end{quote}

175. Wie heißest du? Bist du nicht Jakob? Woher hast du denn solche große Kraft und so große Stärke? Denn die Anfechtung ist jetzt etwas gelinder geworden, und Christus hebt an die Larve abzulegen und tröstliche und lebendig machende Worte zu reden. Er erklärt sich nun, wie er im Evangelium auch thut gegen das canaanäische Weiblein, da er Matth. 15, 28. sagt: „O Weib, dein Glaube ist groß.“ Ihr Glaube war sehr hart angefochten, da er sie eine Hündin nannte. Sie sagt aber dagegen also: Dieweil du mich einen Hund nennst, so gib mir auch die Brosamlein von der Herren Tische, die den Hunden gebühren. Das ist wahrlich ein sehr schöner und herrlicher Glaube gewesen, dazu auch ein trefflich Exempel, das da anzeigt, welches die Weise und Kunst sei, mit GOtt zu kämpfen. Denn wir sollen nicht alsbald bei dem ersten Streiche den Muth und alle Hoffnung fallen lassen, sondern man muß anhalten, beten, suchen und klopfen; und wenn er schon will davon gehen, so höre doch du nicht auf, sondern folge ihm immer fleißig nach; wie das canaanäische Weiblein that, vor welcher sich Christus nicht hat verbergen können, sondern sie kam in das Haus, darin er war, sagt Marcus Cap. 7. V. 25., und fiel nieder zu seinen Füßen.

\quellseite{212}

176. Und wenn er sich auch im Hause in die Kammer verbirgt, und nicht haben will, daß jemand zu ihm kommen soll, so weiche du doch nicht, sondern folge immer nach: will er dich nicht hören, so klopfe an die Kammerthüre, mache ein Geschrei. Denn das ist das höchste Opfer, daß man mit dem Gebete nicht aufhöre, und daß man immer anhalte und suche, bis wir ihn überwinden und obliegen. Und er hat sich uns jetzt schon ergeben, daß wir des Sieges können gewiß sein. Denn er hat sich an die Verheißung gebunden und hat seine Treue und Glauben mit einem Eide bestätigt, da er Joh. 16, 23. sagt: „Wahrlich, wahrlich, ich sage euch, so ihr den Vater etwas bitten werdet in meinem Namen, so wird ers euch geben.“ Item, Marc. 16, 16.: „Wer da glaubet und getauft wird, der wird selig werden.“ Diese Verheißungen werden dir nimmermehr fehlen, es wäre denn, daß du nicht folgen und suchen wolltest. Da verläßt du durch deine eigene Schuld mit deinem Schnarchen und Faulenzen die gewissesten Verheißungen und Christum, den HErrn, selbst, dieweil du dich weigerst, auf den Platz zu treten und mit GOtt einen Kampf aufzunehmen, da der Nutzen dieser Verheißung geschaut und geübt wird.

177. Wie dieser Mann den Jakob übt und versucht, bis sich die rechte Stärke und Kraft des Glaubens hervorthut. Und darum verändert er ihm den Namen und spricht: „Du sollst nicht mehr Jakob heißen.“ Du bist zuvor ein Untertreter genannt worden um deines Bruders willen, sie haben dir aber noch den rechten Namen nicht gegeben: „du sollst Israel heißen“; denn das hast du verdient mit deiner Stärke und unüberwindlichem starkem Muth, damit du GOtt und Menschen überwunden hast.

178. So erklärt sich nun dieser Kämpfer, da er sagt, daß Jakob mit GOtt und Menschen gekämpft habe, und gibt ihm von dem Kampfe einen neuen Namen. Denn das hebräische Wort el heißt GOtt. Und bei allen Vätern ist gefragt worden: Was dieses Wort „Israel“ heißen sollte? Hieronymus beklagt sich bescheiden, wie er durch das große Ansehen und Gelehrsamkeit etlicher sehr beredter Männer eingeschüchtert werde, die dieses Wort auslegen, daß es auf Lateinisch heiße: Vir videns Deum: Ein Mann, der GOtt sieht. Er setzt aber doch dazu, daß er die Deutung des Engels oder GOttes lieber annehmen wolle, denn der Andern, wenn sie auch die allerberedtesten Männer sein mögen.

179. Das hebräische Wort sarah heißt herrschen, obliegen, ein Oberster oder Fürst sein; wie Jes. 9, 6. steht: sar schalom, Friedefürst. Wenn aber der Buchstabe j davor gesetzt wird, daß man sagt jeschar el, so vergeht das erste a, und wird alsdann ein Eigenname, israel; wie auch in andern Worten geschieht, wo der Buchstabe j davor gesetzt wird, als, Josaphat, Jesus, Jesaia. Jisrael heißt Kämpfer, Fürst, Herr oder Herrin GOttes. Sie haben es sich auch alle sauer werden lassen, dieses Wort auszulegen; denn es hat ihnen gedäucht, ein ungereimtes Ding zu sein, daß wir sollten Herren und Ueberwinder GOttes genannt werden. Und es ist wahrlich auch wahr, wenn wir davon urtheilen wollen nach der Philosophie oder Vernunft; aber im Geiste und nach der Theologie sagt man recht und wohl von GOtt, daß er von uns überwunden werde.

180. Jedoch wird er nicht also überwunden, daß er uns müßte unterthänig sein, sondern sein Gericht oder Zorn und Grimm, und alles, was uns hindert und im Wege liegt, wird von uns überwunden mit Bitten, Suchen und Klopfen, daß er aus einem zornigen Richter, wie er vorher ist angesehen und gehalten worden, ein sehr lieber und freundlicher Vater wird und sagt: „O Weib, dein Glaube ist groß“; item: „Dein Glaube hat dir geholfen“; „Wie du glaubest, so geschehe dir.“ Ei, wie thust du mir so wehe mit deinem Geschrei! Das ist nun voller Trost, daß uns GOtt also übt und ermahnt, daß wir kämpfen sollen, und läßt sich vernehmen, daß ihm ein sehr angenehmes Opfer sei, daß er von uns möge überwunden werden.

181. Dieses darf man aber nicht nach der Vernunft beurtheilen, die sich auch untersteht, GOtt zu überwinden mit eigenen Kräften und Werken; sondern er wird also überwunden, wenn er sich uns zuerst ergeben und gefangen hat im Worte, in der Verheißung, in der Taufe. Da ist es denn übrig, daß du dasjenige überwindest, das dir diesen GOtt nehmen will, nämlich, durch die Wahrheit der Verheißungen und des Glaubens. Oder wenn er sich stellt, als sei er dir ungewogen, und zürne mit dir, und daß er dich nicht hören, dir auch nicht helfen wolle, so sollst du sagen: Lieber Herr GOtt, du hast dieses ja in deinem Worte verheißen; derhalben wirst du ja deine Verheißung nicht ändern, ich bin getauft, bin absolviret. Wenn du auf diese Weise hart wirst anhalten und immer also auf ihn dringen, so wird er überwunden werden und sagen: „Dir geschehe, wie du willst“; denn du hast die Verheißung und den Segen, ich muß mich dir ergeben. Denn wo man beständig bleibt, und immer sucht und fleißig bittet, das ist ihm ein süßes angenehmes Opfer.

182. Dieses ist sehr fein und lieblich in der Theologie, daß die Christen überwinden GOttes genannt werden, auf Hebräisch sarah oder sarah GOttes. Denn sarah heißt sonst auch eine Hausmutter, oder eine solche Herrin, die im Hause regiert. So sind nun die Gläubigen Fürsten und Herren GOttes; denn sie kämpfen mit GOtt und siegen. Nach der andern Deutung fragen wir aber nichts. Denn es gibt keinen Sinn, wenn du dieses Wort also auslegen willst: der Herr mit GOtt; sondern es wird hier ausdrücklich also gesagt: „Du hast mit GOtt und mit Menschen gekämpft.“ Der Kämpfer nennt sich nicht allein GOtt, sondern auch einen Menschen, als wollte er sagen: Dieweil du mich überwunden hast, so hast du GOtt und Menschen überwunden. Er setzt es beides zusammen: Du hast zuvor einen Menschen, nämlich, deinen Bruder, mit Untertreten überwunden, daher du den Namen Jakob bekommen hast; jetzt aber bist du zugleich GOtt und Menschen obgelegen.

\begin{quote}
\textbf{V. 29—31. Und Jakob fragte ihn und sprach: Sage doch, wie heißest du? Er aber sprach: Warum fragest du, wie ich heiße? Und er segnete ihn daselbst. Und Jakob hieß die Stätte Pniel; denn ich habe GOtt von Angesicht gesehen, und meine Seele ist genesen. Und als er vor Pniel überkam, ging ihm die Sonne auf.}
\end{quote}

183. Der lateinische Dolmetscher setzt im Text das Stücklein hinzu: quod est mirabile. Es steht aber im Hebräischen nicht, darum haben wir es auch ausgelassen. Dieweil Jakob aber so sehr erschrocken und bestürzt war, versteht er die Worte noch nicht. Denn er hat die Anfechtung des Mannes noch nicht ganz überwunden, welche er bisher erlitten, und die so schwer war, daß er nicht anders gefühlt hat, denn daß er müßte verloren sein und sterben. Denn solche Erschwingungen geschehen eben auf dieselbige Weise, wie dem zu Muth ist, der da angefochten wird, oder wie das arme schwache Fleisch denkt und fühlt: Denn es muß beides geschehen und gefühlt werden, nämlich, die Tödung des Fleisches und die Lebendigmachung des Geistes. Wie oftmals in der heiligen Schrift gesagt wird, Ps. 118, 21.: „Ich danke dir, daß du mich demüthigest, und hilfst mir“; item 1 Sam. 2, 6.: „Der HErr tödtet, und macht lebendig, führet in die Hölle, und wieder heraus.“ Das sind unseres HErrn GOttes Werke.

184. Wenn nun Jakob Schwachheit des Glaubens fühlt und nicht fest genug an der Verheißung hängt, sondern nach dem Fleisch noch etlichermaßen wankt und an dem Willen GOttes zweifelt, wie das Fleisch zu thun pflegt: alsdann erscheint ihm das Angesicht GOttes sehr unfreundlich und sauer, kämpft und stellt sich, als ob es ihn tödten wolle. Eben dasselbe widerfährt auch Allen, die angefochten werden, wenn sie sich der Gütigkeit GOttes, die sie vor Augen sehen, nicht freuen, sondern erschrecken und fürchten sich vor dem zornigen Gott; da dünket sie, daß ihnen Alles den gewissen Tod drohe, da ist ihnen Alles schwarz, dunkel, traurig, unglückselig und kläglich im Himmel und auf Erden, es scheint ihnen weder Sonne noch Mond; denn das Herz, so im Fleisch unterdrückt wird, stirbt dahin, und wenn dasselbe todt ist, so fällt Alles dahin. Also, da Jakob noch ein Fünklein Geistes hat, das in ihm lebte, kämpft und streitet er in großer Schwachheit; welche Schwachheit doch sehr stark ist, wie droben gesagt ist.

\quellseite{213}

185. Dieses sind heimliche und wunderbare Dinge, und die allein denen bekannt sind, so die Verheißungen haben, darin sie also versucht und gedemüthigt werden; und in solcher Demüthigung werden sie doch überwinder, auch über Gott selbst. Denn ich, sagt dieser Mann, bin Gott und Mensch, den du obgelegen hast. Aber dieweil der Kampf währt, versteht es Jakob nicht, und sieht auch seine Stärke nicht wegen der Schwachheit; sondern er hebt an zu fragen: Wer bist du? Sage mir deinen Namen, weil du sagst, daß ich Gott und Menschen obgelegen habe. Denn da hat der Kämpfer erklärt, daß er Gott und Mensch sei. Aber Jakob versteht es noch nicht, und liegt noch gar in Finsterniß, da ihm die Sonne dermalen noch nicht aufgegangen ist. Er hat noch immer die Gedanken, daß der Widersacher vorhabe, daß er ihn tödten wolle; und doch gleichwohl, dieweil er hört, daß er obgelegen habe, sagt er: Ich wollte gleichwohl gerne wissen, wer du bist, dieweil du da dermassen mit mir gerungen hast, daß du mir das Leben wolltest nehmen, und bekennst doch nun gleichwohl, daß du überwunden seiest.

186. Der Kämpfer aber verhehlet es noch vor ihm, wer er sei, und läßt ihn ungewiß und im Zweifel stehen; und wie Jakob gesegnet ist, also wird ihm auch das Bild dieses Mannes gezeigt. Jakob versteht es nicht, darum antwortet ihm der Mann auf die Frage nicht, was den Namen belanget. Als wollte er sagen: Ich habe dir es schon zuvor gesagt, daß du mit Gott und Menschen gekämpfet und obgelegen bist; daraus solltest du ja schließen und abnehmen können, wer ich sei. Den Namen will er ihm nicht sagen, sondern fährt noch immer fort ihn zu versuchen, zu üben und zu züchtigen. Wenn aber Jakob bei sich selbst wäre und die Wolken den Schein der Sonne nicht verdunkelt hätten, so würde er es desto leichter verstanden haben, daß er ein Nebenmensch Gottes und der Menschen genannt worden war. Die Worte versteht er zwar wohl, aber die Sache an sich versteht er nicht; denn er weiß nicht, daß dies der HErr sei. Er denkt: Wer bist du denn, indem ich Gott und Menschen überwunden habe? Es folgt aber endlich die Erscheinung des HErrn, da er ihn segnet. Denn da legt er die äußerliche Larve ab, und offenbart sich, daß dieser Mann, der mit ihm gerungen hat, Gott und Mensch sei, der sich in künftiger Zeit in unserm Fleisch offenbaren würde und den die Juden kreuzigen würden.

187. Was es aber für ein Segen gewesen sei, hat Mose nicht beschrieben. Es ist aber ohne Zweifel der Segen gewesen, damit die Väter gesegnet worden sind, nämlich: „Durch deinen Samen werden alle Völker gesegnet werden“ etc. Als wollte er sagen: Lieber Jakob, ich bin der, der dich gesegnet hat und segnen will. Das nun Jakob zuvor, da dieser noch ein Mann war, nicht hat sehen können, sondern disputirt, und fragt aus Zweifel und dieweil er es noch nicht gewiß war, dasselbe versteht er nun aus dem Wort und Segen und sieht es mit großer Freude. Bist du der? denkt er, ei du himmlischer Vater und HErr! ich meinte, es wäre irgend ein Gespenst, oder sonst ein Mann. So bist du nun derselbe Segner, der mich und meinen Vater Isaak, dazu auch meinen Großvater Abraham gesegnet hat?

188. Dieses ist ein sehr fröhliches Ende und Beschluß dieses Kampfes. Denn jetzt kommt nun Jakob aus der Hölle in den Himmel, aus dem Tode in das Leben. Und ist wahrlich ein sehr harter und schwerer Kampf gewesen, den er bisher geführt hat. Darum dankt er nun Gott und bekennt seine Angst und Noth, und will nicht haben, daß diese Stätte ihren alten Namen länger behalten soll, sondern er nennt sie Pniel. Als wollte er sagen: Diese Stätte soll nicht ein Kampf oder Gesicht der Hölle, sondern ein Gesicht Gottes genannt werden. Und darum sagt er: „Ich habe Gott von Angesicht gesehen“; jetzt sehe ich klarlich, daß der Kämpfer, der mich versucht hat, Gott selbst gewesen ist. Warum bin ich also bestürzt? Warum habe ich also gezittert? Ich habe nicht gewußt, daß dieses der HErr, mein Gott, war.

189. Und nun wollen wir ohne allen Widerspruch sagen, daß dieser Mann nicht ein Engel, sondern unser HErr Jesus Christus gewesen, der ewiger Gott ist und ein Mensch hat werden sollen, den die Juden kreuzigen würden. Der ist mit den heiligen Vätern sehr vertraut gewesen, ist ihnen oftmals erschienen, hat oft mit ihnen geredet. Darum hat er sich den Vätern in solcher Gestalt gezeigt, daß er damit anzeigen möchte, daß er dermaleinst im Fleisch und menschlicher Gestalt bei uns auf Erden wohnen würde.

190. Dieses ist die rechte Auslegung, die wir nicht erfunden oder von Andern genommen haben, sondern dieser Mann hat sich selbst also geoffenbart, da er Jakob den Namen gibt, daß er Israel heißen soll. Und Jakob sagt an diesem Orte selbst auch: „Ich habe Gott von Angesicht gesehen.“ Eben derselbe unser HErr Jesus Christus hat Jakob also versucht, nicht daß er ihn verderben, sondern stärken und Kraft geben wollte, und daß er in diesem Kampfe die Kraft der Verheißung etwas besser lernen möchte; ja, er hat Jakob diese Kraft und Stärke gegeben, daß er siegen und diese Geschichte des HErrn mit Freuden rühmen möchte.

191. In der Noth und im Kampfe hat er nicht also geredet; wie Andere auch unter dem Kreuze niemals solche fröhliche Worte pflegen zu reden. Es müssen aber doch die Gottseligen und Geistlichen endlich im Kampf überwinden, und dahin kommen, daß sie Gott auch von Angesicht sehen. Wie zuvor von dem cananäischen Weiblein gesagt ist, die Christum den HErrn gesehen und gehört hatte; und da er in ein Haus ging, da er wollte verborgen sein, folgte sie ihm nach, und hielt immer an, bis sie Christo sein Herz abgenommen hat, welches sie, wie es schien, allzu hartnäckig abwies.

192. Und wenn ein solcher Kampf überwunden ist, so können wir uns mit dem Patriarchen Jakob auch freuen, und rühmen, daß wir Pniel, das ist, das Angesicht des HErrn gesehen haben. Denn also geschah auch oftmalen der sehr schweren Kampf, den er erfahren hatte, und auch seine Freude. Und will so viel sagen: Lieber HErr Gott, wie bin ich in so großer Angst und Beschwerung gewesen, dazu in großer Bestürzung! aber Gott sei Lob und Dank, ich bin nun heraus gekommen und genesen, meine Seele ist erlöset und errettet aus aller dieser Angst, nun danke ich dem HErrn, meinem Gott.

\begin{quote}
\textbf{V. 31. 32. Und er hinkte an seiner Hüfte. Daher essen die Kinder Israel keine Spannader auf dem Gelenk der Hüfte, bis auf den heutigen Tag; darum, daß die Spannader an dem Gelenk der Hüfte Jakobs gerühret ward.}
\end{quote}

193. Da der Kampf ein Ende hat, fühlt der Patriarch, daß ihn die Hüfte schmerzet. Er bekommt aber auch Trost, die himmlische Sonne oder das Angesicht Gottes, und das Licht der Sonne, so täglich und natürlich ist: Himmel und Erde ist ihm eitel Licht. Er hat hierinnen an diesem Orte von der Spannader, davon er droben, V. 25., nichts gesagt hat; denn er hat nur des Gelenkes gedacht. Hier setzt er nun die Spannader dazu; denn die Gelenke haben sehr starke Adern um sich. Wir haben es auf deutsch gegeben, die Spannader. Die Juden essen davon nicht, die Deutschen aber scheuen sich davor gar nicht, ja, es sind Viele, die eine Delicatesse daraus machen.

194. Was das hebräische Wort nascheh eigentlich heiße, weiß ich nicht genau. Der Name Spannader ist daher aufgenommen, der heißt auf Lateinisch, obliviscens, vergessend. Die Juden wissen selbst auch nicht eigentlich, was dieses Wort heiße. Etliche legen es aus, daß es auf Lateinisch heißen soll, nervus oblivionis; als wenn man sagen wollte: das seine Stätte vergessen hat um des Verrenkens willen. Aber es ist jüdisch, das ist, ungeschickt und faul Ding. Denn die Juden haben harte und scharfe Deutungen. Die Andern haben es auf Lateinisch verdolmetscht, nervus contractionis, das wir eine Spannader nennen, und dieselbige Deutung ist besser. Denn man pfleget den Fuß zusammen und wieder auseinander zu ziehen und also zu regen. Ich halte dafür, es komme von dem Worte nasah, das heißt, aufheben, daß es heiße, die Ader, so über der Hüfte erhaben ist; ich sage es aber nicht als gewiß. Hieronymus hat es in Lateinisch nicht recht gegeben: Nervus emarcuit in femore: Die Ader ist ihm in der Hüfte verdorret. Denn Mose will sagen, sie ist verrenkt; und nach solchem Verrenkten pfleget ein sehr großer Schmerz zu folgen. Das ist, das Mose sagt, daß die Spannader an der Hüfte Jakobs sei gerühret worden.

195. Das ist die Historie und der einfältige Verstand dieses Textes, welche Historie mir erklärt, wie wir am besten gekonnt haben; und ich hoffe, sie werde Andern auch klar genug sein. Nun sollte man auch etwas von der Allegorie oder heimlichen Deutung sagen. Ich habe aber oftmals gesagt, daß die Allegorien gefährlich sind, wo man die heilige Schrift auslegt. Denn Einer folgt Diesem, der Andere folget einem Anderen: Einer legt diesen, der Andere einen andern Grund; und bisweilen hauen sie Gold, bisweilen Stoppeln darauf. Und St. Augustinus sagt fein: Figura nihil probat, das ist: Mit Figuren oder Gleichnissen kann man nichts beweisen.

\quellseite{214}

196. Warum wenn man die Lehre unseres Glaubens beweisen und bestätigen soll, da gelten die Allegorien oder andere Gleichnisse nicht; sondern man muß solchen Beweis und Zeugniß haben, die man aus den Quellen der heiligen Schrift genommen hat. Denn ein Jeder kann seines Gefallens Figuren und Allegorien entweder auslegen, oder aber selbst dichten. Darum ist nichts Gewisses daran, damit man bestehen könnte.

197. Jedoch muß man sie aufführen und gebrauchen wie schöne Lichter und Zierrathe, damit die Lehre oder der rechte Verstand nach dem Buchstaben geziert und erklärt wird. Wie St. Paulus Gal. 4, 22. ff. das Exempel von Abraham, Hagar und Sarah anführt, daß er damit die Lehre von der Gerechtigkeit des Glaubens und von den zwei Testamenten zieren und erklären möge. Und wenn sich die Allegorie mit der Lehre reimt, so ist es ein sehr schöner Schmuck, damit die Lehre gleichsam gewürzt wird.

198. Ehe wir aber von der Allegorie handeln, wollen wir noch etwas länger bei dem Exempel bleiben, das uns in dieser Historie vorgehalten wird, dasselbe weiter zu erklären. Denn es muß allezeit der buchstäbliche Sinn vorher gehen, als die Grundlage der Allegorien. Und dies ist ein trefflich und sehr heilsam Exempel, das man der Kirche und Gemeinde Gottes mit sonderlichem Fleiß vorhalten soll. Denn wir werden darin ermahnt, daß wir uns in unserem Leben auch also schicken und an dem Bilde des Kampfes Jakobs die Gemeinde Gottes erkennen lernen. Denn Gott verbirgt die Kirche oder Gemeinde und auch unsere Seligkeit unter einer dunklen und erschrecklichen Decke; daran wir uns gewöhnen müssen, auf daß wir nicht verzagen oder mißtrauen auch in der größten Noth und Widerwärtigkeit, so uns der Teufel, die Welt, oder auch wohl Gott selbst widerfahren läßt. Denn also sagt St. Paulus Col. 3, 3.: „Euer Leben ist verborgen mit Christo in Gott.“

199. Dieses ist die Summa von diesem Exempel. Wir leben nicht im Fleische, oder wenn wir schon darin wohnen, so leben wir doch nicht nach dem Fleische, und wir dienen der Weisheit und der Gerechtigkeit des Fleisches nicht, und vielweniger dem Geiz, dem Zorn oder dem Wucher; das sind unsere Waffen nicht: sondern wir wohnen im Fleische und leben im Glauben des Sohnes Gottes. Denn ich weiß von keinem anderen Leben, denn daß ich glaube an Gott, den allmächtigen Vater, und an Jesum Christum, seinen Sohn, unsern Herrn, der für mich gelitten hat etc.

200. Ja, möchte Einer sagen, dieses steht man aber nicht. Antwort: Es ist ja verborgen in Gott, aber doch also, daß der Teufel uns dieses Leben, wie sehr es auch verdunkelt ist, nicht nehmen oder hinreißen kann. Denn wir werden nicht gegeißelt, nicht gekreuziget im Fleisch, freuen uns auch im Fleisch nicht; sondern wir werden getödtet, sterben, leiden; wir werden geschändet und gelästert; wir werden verspottet, verspeiet, und für die allerärgsten Menschen gehalten, welche sie als Ketzer und Missethäter mit Schwerdt und Feuer verfolgen und umbringen, wenn sie nur können. Ist nicht das eine sehr schwere Decke dieses Lebens? Aber darunter liegt verborgen ein sehr schöner und köstlicher Edelstein, der also genannt wird: Ich glaube an den Sohn Gottes.

201. Denn wie Christus, unser Haupt, verborgen ist in Gott, also ist die Kirche und Gemeinde Gottes bedeckt und liegt auch verborgen unter solcher Decke, und sie muß in der Welt leiden, daß sie aufrührerisch, irrige, ketzerisch, und ein Haufe der bösesten Menschen genannt wird, die je auf Erden gelebt haben. Diese Schmach, sage ich, muß sie tragen. Das sind ihre schönen Edelsteine, die sie trägt auf Erden, ihr Kleinod und goldene Ketten, damit sie Gott in diesem Leben zieret. Denn also wird es in diesem Exempel abgemalt, wie Gott uns in sonderheit und auch insgemein die ganze Kirche versuche nach dem Exempel Jakobs, nämlich, daß er mit ihr kämpft, und sich als einen Widersacher und Feind erzeigt, der sie verlassen, verwerfen, ja, gar verderben wollte.

202. Denn wo du recht ansehen wirst, wie es um unsere Kirche steht und gelegen ist, so scheinen wir gar nichts zu haben, denn nur das Wort und die Sacramente; und wir haben viel unzählige Widersacher, als, Fürsten, Edelleute, Bürger, Hausgenossen, Schüler, und endlich unser eigen Fleisch, damit wir uns schleppen. Denn das ist unser Ruhm, daß das Wort geprediget und verspottet werden auch von unsern eigenen Hausgenossen. Das ist unsere Decke, darum wir uns dünken lassen, daß uns Gott mit nichten wolle für die Seinen erkennen und halten. Denn es reimt sich wohl nichts weniger zu der Kirche denn eine solche Gestalt.

203. Ich habe selbst oftmals gedacht, da ich solche große Verachtung, Ueberdrus und Haß des Wortes bei den Leuten gesehen habe, nachdem das Licht des Evangeliums wiederum ist hervor gekommen: Lieber, warum habe ich doch angefangen zu lehren und zu predigen? oder, warum fahre ich doch damit fort? dieweil die Leute immer mehr und mehr wider uns toben, als wären sie unsinnig, und alle Tage ärger werden.

204. Aber dieses muß man leiden und also schließen: Das ist der rechte Gott, der verborgen ist, das ist seine Art und Eigenschaft, er ist in Wahrheit verborgen und dennoch nicht verborgen. Denn das Fleisch liegt uns im Wege und hindert uns, daß wir ihn nicht ansehen können; dasselbe murret, ist betrübt, ist zornig und ungeduldig, und schreit: Ich bin der allerelendeste Mensch, bin der allerverachtetste und verworfenste. Und ist wahrlich auch wahr, die Pfarrherren und Kirchendiener mit andern gottseeligen Menschen haben keinen Schutz von den Fürsten oder Königen, ja, sie haben auch von den Bürgern und Bauern den Schutz nicht, so ihnen von Rechts wegen gebühret. Und das noch mehr ist, was sie auch jetzt schon haben, wird ihnen oftmals genommen, und müssen also allerlei Räuberei und Gewalt unterworfen sein. Darum scheint es, als habe uns Gott gar verlassen und verworfen; denn er ist vor uns verborgen und wir werden mit ihm zugleich verborgen: aber im Glauben, im Worte und in den Sacramenten offenbart er sich und läßt sich sehen.

205. Die Vernunft, Weisheit, Gerechtigkeit des Fleisches und das Licht der Sonne hält Gott für einen finstern und dunklen Ort; aber da kommt das Wort dazu, das leuchtet wie ein Flämmlein mitten in Finsterniß und theilt den Glanz umher aus durch die Lehre und Sacramente; dieselben heißt uns Gott ergreifen. Wenn wir sie annehmen, alsdann ist uns Gott im Geiste nicht mehr verborgen, sondern nur im Fleisch.

206. Deshalb sollen wir nun auch lernen, wenn wir mit Unglück und Jammer überfallen und mit Finsterniß bedecket werden (die uns unausbleiblich macht, daß wir sollten die Kirche sein, oder auch, daß wir Gott gefallen sollten), daß wir alsdann das Wort ergreifen, und sinken und fallen lassen, was da fällt, und daß wir uns nicht anfechten lassen, wenn schon Andere abfallen und die Lehre verleugnen; sondern denken, daß wir an einem dunkeln Orte, 2 Petr. 1, 19., sind, da uns das Licht des Wortes vorleuchtet: „Wer da glaubet und getauft wird, der wird selig werden“, Marc. 16, 16. Das ist das einige Licht, welches die Sonne nicht, menschliche Vernunft nicht sieht, sondern es leuchtet im Herzen. Ohne das Wort sollen wir nichts wissen und nichts sehen; denn wo das nur leuchtet, hat es keine Gefahr, und es wird eine Stunde kommen, darin wir werden hervorbrechen und mit Freuden sagen: „Ich habe den Herrn von Angesicht gesehen, und meine Seele ist genesen.“

207. So hat nun Jakob der Kirche und Gemeinde Gottes ein sehr nützliches und schönes Exempel gegeben, nämlich, eines solchen Glaubens, der in Schwachheit kämpft, auf daß wir nicht denken, wie die Mönche gedacht haben, daß die Väter und Propheten Steine oder Klötze gewesen seien, an welchen gar keine Schwachheit sollte gewesen sein. Denn sonst müßten wir verzagen, die weil wir in unserm Fleisch noch sehr große Schwachheit fühlen; denn das Fleisch murret, und will nichts leiden oder sich unter die Füße treten lassen. Darum klagen wir und schreien, daß man unbilliger Weise mit uns umgehe. Deshalb lasset uns die heiligen Patriarchen ansehen und uns mit ihrem Exempel trösten, sintemal sie auch nicht stets fest gestanden und stark im Glauben gewesen sind.

208. Jakob war mit sehr reichen Tröstungen überschüttet, nicht allein mit denen, welche er selbst empfange, sondern auch mit denen, welche seinem Vater und Großvater, Abraham und Isaak, gegeben waren: und kämpft doch mit großer Schwachheit. Darum sollst du also denken: Ich bin nicht allein, der da angefochten wird mit dem Zorn Gottes, mit der Vorsehung und Mißglauben; ich bin es nicht allein: alle Heiligen, so viel ihrer sind, die jemals geglaubt haben oder jetzt noch glauben an den Sohn Gottes, die erfahren solchen Kampf der Anfechtung auch, damit sie oder die ganze Kirche versucht wird. Denn was ist diese ganze Versammlung, so die Kirche genannt wird? Es ist ein kleines Häuflein sehr armer elender Menschen, die vor der Welt die verachtetsten und verzagtesten Leute sind. Was ist das Häuflein gegen die ganze Welt zu rechnen? Was ist es gegen das türkische Reich und Frankreich? ja, auch gegen unsere Widersacher, die Papisten?

\quellseite{215}

209. Wenn du nun fragst, wo die Kirche sei? so siehst man sie nirgends. Man muß aber auf die äußerliche Gestalt nicht sehen, sondern auf das Wort, auf die Taufe, und da muß man die Kirche suchen, da die heiligen Sacramente recht und vollkommen gereicht werden, da Zuhörer, Lehrer und Bekenner des Wortes sind. Wenn man es noch nicht sehen kann, so bedenke, daß unsere Güter verborgene Güter sind, und daß man es in diesem Leben nicht sehen kann, wie groß sie sind. Der Mann ringt noch, es ist aber ohne Gefahr, ja dieser Kampf gewinnt einen sehr fröhlichen Ausgang.

210. Denn also sollen wir aus dieser Historie ein Sprüchwort machen: Wenn man denkt, unser Herr Gott habe Einen verworfen, so soll man dafür halten, unser Herr Gott habe ihn in den Armen und herze ihn. Also fühlt und denkt Jakob nicht anders, denn daß er werde verloren sein; aber da er sich umsieht, so ist er unserem Herrn Gott in den Armen. Dasselbige lehrt uns auch das Exempel Hiobs, der auch gedemüthigt und sehr betrübt worden ist. Denn also führt Gott seine Heiligen wunderlich, Ps. 4, 4. Wenn wir meinen, es sei mit uns aus und Alles verloren, alsdann herzt und küßt uns Gott als seine herzlieben Kindlein. Das ist es, das St. Paulus sagt: „Wenn ich schwach bin, so bin ich stark“; wenn ich sterbe, so lebe ich, 2 Cor. 12, 10. 2 Tim. 2, 11.

211. Wir verstehen es aber nicht, und ist das die Ursache, weil uns das Fleisch im Wege liegt, welches nicht leiden kann, daß es getödtet werde, und hindert den Geist, daß er die unermessliche Liebe und den guten geneigten Willen Gottes gegen uns nicht sehen kann, bis er aus diesem Kampfe heraus kommt und das Hinderniß des Fleisches abgeschafft werde. Aber dasselbe wird hin und wieder in der ganzen heiligen Schrift reichlich gelehrt. Denn also steht in der Offenbarung Johannis am 3. Cap. V. 19.: „Welche ich lieb habe, die strafe und züchtige ich.“ Und in den Sprüchen Salomos am 3. Cap. V. 11. 12.: „Mein Kind, verwirf die Zucht des Herrn nicht; und sei nicht ungeduldig über seiner Strafe. Denn welchen der Herr liebet, den straft er, und hat Wohlgefallen an ihm, wie ein Vater am Sohn.“ Diese und dergleichen Sprüche mehr sind sehr trefflich und merkenswerth.

212. Heißt aber das ein Wohlgefallen haben am Sohne, wenn man ihn mit Ruthen streicht und schlägt? Also lehrt wahrlich die Schrift und die Erfahrung bezeugt solches auch. Denn was fromme und getreue Väter sind, dieselbigen züchtigen ihre Kinder mit allem Ernste; sie wollen nicht leiden, daß sie ungehorsam sollen werden und das Ihre übel durchbringen etc. Und das thun sie aus keinem bösen Herzen, sondern meinen es also, daß die Kinder dadurch unterwiesen, und was Arges an ihrem Leben ist, gebessert werden möge, auf daß sie endlich fromme Männer werden mögen. Der Knabe aber oder das Kind, das gestrichen und gezüchtigt wird, verstehet den guten geneigten Willen und die herzliche Neigung des Vaters nicht, sondern meint, der Vater sei ein Tyrann und habe alle Sorge und Liebe zu ihm gar fallen lassen. Also thut auch ein frommer getreuer Schulmeister. Der haßt an dem Schüler, daß er faul und schläfrig ist; darum fleißigt er sich, daß er den Knaben kindlicher Weise züchtige, nicht daß er ihm feind wäre, oder Lust hätte, die Knaben zu schlagen oder zu bläuen (wie ich, da ich ein Knabe war, vor Zeiten in den Schulen deren Viele gesehen habe), sondern er hat den Knaben sehr lieb, und solche Liebe bewegt ihn dazu, daß er ihn züchtigt.

213. Dieses ist also die Wahrheit an sich, und die Erfahrung, die da lehrt, daß die Züchtigung noth thut. Und der himmlische Vater hält eben dieselbige Weise auch, die Gottseiligen zu züchtigen. Wer nun diese Züchtigung ertragen und beharren kann, der wird endlich erfahren, daß in der Trübsal an die Hebräer am 12. Cap. V. 11. steht: „Alle Züchtigung aber, wenn sie da ist, dünkt sie uns nicht Freude, sondern Traurigkeit sein; aber darnach wird sie geben eine friedsame Frucht der Gerechtigkeit denen, die dadurch geübet sind.“

214. Eben dasselbe sehen wir auch in dieser Historie geschehen: da der Fluch weicht und die Versuchung vorüber ist, folgt darauf der Segen und Freude, daß Jakob sagt: „Ich habe Gott von Angesicht gesehen; und meine Seele ist genesen.“ Denn also deutet ein gottseeliger Mensch, wenn er der Anfechtung entronnen ist: Ei, wie habe ich doch so einen gnädigen Gott! Ich war schon verzagt und ging in die Hölle hinab, nun sehe ich aber, daß mir dieser Kampf zum Leben nützlich gewesen ist. Ich hätte immermehr gemeint, daß Gott so nahe bei mir wäre. Darum ist es sehr lieblich, daß man auf solche Weise die Zucht der Züchtigung erfahren mag. Ach du himmlischer Vater, bist du mir so nahe und ich wußte es nicht? wie ist mir jemand so wohl! Laß nun Esau kommen und alle Teufel, so fürchte ich mich nicht; denn ich habe den Herrn, meinen Gott. Zuvor habe ich ihm auf den Rücken gesehen, in der Gestalt und Person eines Mannes, in welcher Person ich mich dünken ließ, daß er mir den Tod drohte, und mein Herz war in großer Angst, daß er mich irgend in die Hölle würde stoßen: jetzt aber sehe ich sein Angesicht.

215. Denn die äußerliche Gestalt pfleget sich eben so sehen zu lassen, wie ein gottseeliger Mensch innerlich in seinem Herzen denkt oder gesinnet ist. Wie den Jüngern widerfahren ist: da Christus auf dem Meere ging, da erschracken sie und waren in Angst, dieweil sie den Herrn nicht kannten, und darum dünkte sie auch Christus irgend ein Gespenst oder der Teufel zu sein. Da sie aber hören, daß er sagt: Ich bin es, da wird die Gestalt verändert und ihre Augen werden geöffnet, daß sie Christum erkennen und alle Furcht fahren lassen; das heißt, den Herrn sehen von Angesicht zu Angesicht, nämlich, wiederum aus der Hölle geführt und in den Himmel gesetzt werden, daß das Herz nun also schließe: Nun habe ich nicht einen solchen Gott, der mich versuchet oder mit mir zürnt, sondern der mich segnet und mir gnädig ist. Zu solcher Erkenntniß Gottes aber gehört Uebung und Erfahrung; welche Erkenntniß zwar immer fleißig soll gelehrt und getrieben werden. Es muß aber Erfahrung dabei sein und daß man sich in diesem Leben daran gewöhne, sonst kann man es nicht verstehen; sintemal die Dinge, so uns dünken Zeichen zu sein des göttlichen Zornes und daß er uns sollte feind sein, die gewissen Anzeichen der größten Liebe und geneigten Willens Gottes gegen uns sind.

216. Ich habe oftmals gesehen, daß die besten Menschen gewiß sind geplagt worden mit Schrecken, Trübsal und sehr schweren Verfolgungen, so auch, daß sie schier in ihrem Herzen gar verzagt worden sind. Dieses muß man aber lernen, auf daß wir solche Leute trösten können, und die Versuchungen also deuten, daß es gleichsam eine sonderliche Weise sei, da Gott mit uns pfleget zu kämpfen in der Gestalt eines Verderbers, und daß wir sie ermahnen, daß sie fest an der Verheißung halten, oder an dem Zeichen und Fürbild des Wortes, mit solcher Hoffnung, daß gewißlich die Erlösung folgen werde. Denn Gott „führet in die Hölle, und wieder heraus“, 1 Sam. 2, 6. Jetzt siehst du ihm auf den Rücken, und es läßt sich ansehen, als ob dir Gott feind sei und sich von dir gewendet habe; aber bald hernach wirst du ihn auch unter Augen und sein Angesicht sehen. Ja, das heißt diejenigen lieben, die er züchtigt, und diese Liebe muß man durch Erfahrung lernen.

217. Und man muß vor der Züchtigung nicht fliehen, und sich nicht weigern, dieselbe zu dulden. Wie man von einem Bauern sagt, welcher, da er von seinem Pfarrherrn diesen Trost hörte, daß die Trübsal und Jammer, damit uns Gott straft, Zeichen wären der Liebe, geantwortet hat: Ach, wie gerne wollte ich, daß er Andere dermalen liebte und mich nicht!

218. Das ist eine närrische und gottlose Antwort gewesen. Denn man soll nicht also denken oder reden. Man soll die Werke Gottes in uns nicht also deuten und verstehen, sondern man soll wissen, daß es eine sehr heilsame Tödtung sei, dadurch wir gezüchtigt werden zum Leben und Seligkeit und nicht zum Verderben. Wie St. Paulus Röm. 12, 2. zeugt, da er sagt: „Daß ihr prüfen möget“, nicht allein, daß ihr mit Worten gelehret werdet, sondern daß ihr es auch durch Erfahrung lernen möget, welches da sei der gute, der wohlgefällige, und der vollkommene Gotteswille.“ Denn das ist Gottes Wille, unsere Tödtung und Heiligung. Wir können aber nicht geheiligt werden, es sei denn, daß das Fleisch und der sündliche Leib getödtet werde, welcher sündliche Leib in diesem Leben mit ganzer Gewalt sich auf allerlei Sünde gibt, als, Ehebruch, Unzucht, Diebstahl etc.

\quellseite{216}

219. Derhalben richtet, züchtiget und straft uns Gott, bis wir lernen, welches da sei der gute, der wohlgefällige, und der vollkommene Gotteswille, auf daß wir mit David singen, wie im 119. Psalm V. 71. steht: „Es ist mir lieb, daß du mich gedemüthiget hast, daß ich deine Rechte lerne.“ Ich wollte es gerne überhoben sein, mein Fleisch scheut sich vor der Anfechtung; ich weiß aber, daß dieses der gute, wohlgefällige Wille Gottes ist. Item, wie er im 118. Psalm V. 18. sagt: „Der Herr züchtiget mich wohl; aber er gibt mich dem Tode nicht“; V. 13.: „Man stößet mich, daß ich fallen soll; aber der Herr hilft mir“, daß meine Seele das Angesicht Gottes sehen und genießen möchte; V. 21.: „Ich danke dir, daß du mich demüthigest, und hilfst mir.“

220. Bisher haben wir nun gesagt von dem Stempel des heiligen Patriarchen, welches die Heiligen, so im Glauben schwach sind, sehr tröstet. Wie sie aber Trost empfangen haben und die Versuchung ihnen zur Seligkeit gerathen ist, also sollen wir uns dessen auch befleißigen, daß wir ihnen im Kampfe und in der Erlösung auch gleich werden. Denn so werden wir lernen, was da gesagt sei, das im 34. Psalm V. 9. steht: „Schmecket und sehet, wie freundlich der Herr ist. Wohl dem, der auf ihn trauet.“ Also verstehen wir aber und empfinden, wie freundlich und lieblich der Herr sei, wenn wir aus der Finsterniß und Schatten des Todes kommen; denn so ist das Licht lieb und sehr angenehm. So küßt denn das Kindlein die Ruthe. O du liebe Ruthe, wie machst du so fromme Kindlein! Dieses sei gesagt von dem Stempel und von der Historie.

221. Zu den Allegorien habe ich nicht sonderliche Lust; ich will aber doch ein wenig davon sagen. Augustinus setzt diesen Grund seiner Allegorie, daß Jakob der Leib sei seiner Nachkommen; und so weit hat er recht. Denn es muß eine Person sein, welche die Figur ist einer ganzen Versammlung. Das er aber dazu sagt von dem Kampfe, so Jakob mit dem Manne gehalten, und vom Siege, nach welchem Jakob gesegnet wird, das lautet sehr hart und ist allzuweit hergeholt. Denn er meint, es solle dadurch verstanden werden die Kreuzigung Christi, daß Jakob dem Engel oder Gott obgelegen hat; welche Kreuzigung durch die bösen Nachkommen geschehen sollte: dieselbigen haben Christum überwogen, das ist, sie haben ihn gekreuzigt, und der gekreuzigte Christus hat doch gleichwohl Jakob gesegnet, nämlich, den besseren Theil der Nachkommen.

222. Aber dieses ist eine allzuharte und erzwungene Allegorie. Er macht zweierlei Nachkommen, nämlich, gute und böse: die bösen haben obgelegen und gesiegt, die Guten sind gesegnet worden, da Christus den Tod und die bösen Ueberwinder überwunden hatte. Diesen Theil der Allegorie lassen wir uns nicht gefallen, und läßt sich ansehen, daß die Väter, so dieser Allegorie nachgefolgt sind, betrogen sind durch die falsche Uebersetzung mit dem lateinischen Worte latitudo (denn also haben sie das hebräische Wort kaph übersetzt, welches Wort, wie droben gesagt ist, das Gelenk der Hüfte heißt), und aus dem Worte latitudo haben sie zweierlei Nachkommen erdichtet. Hieronymus und Origenes haben dergleichen Allegorien viel erdichtet, aber es ist ihnen damit nicht wohl gerathen. Derhalben bedarf es wohl eines sehr guten Verstandes, die Allegorien zu unterscheiden, welche bei den Vätern gefunden oder von Andern erdichtet worden.

223. Mich dünkt aber, daß es sich nicht übel reime, wenn wir aus Jakob eine Person oder Leib machen der guten Nachkommen und nicht der bösen, wie die Schrift einen Unterschied macht. „Denn es sind nicht alle Israeliter, die von Israel sind. Nicht sind das Gottes Kinder, die nach dem Fleisch Kinder sind; sondern die Kinder der Verheißung“ etc., Röm. 9, 6. 8., sonst hätten auch die bösen Nachkommen den Segen Gottes empfangen; sondern Jakob ist eine Figur und Vorbild der rechten Nachkommen Israels, das ist, aller Heiligen, und unser auch, die wir an Christum glauben. Darum ist Israel in dieser Anfechtung mit Jakob, wenn entweder die Glieder insonderheit angefochten werden, oder aber insgemein die ganze Kirche.

In solchem Kampfe überwinden wir Gott, wie zuvor gesagt ist; denn wir haben die Verheißung. Wenn wir daran mit dem Glauben hangen, so werden wir auch in der Schwachheit stärker, also daß wir unsern himmlischen Vater und den Sohn Gottes überwinden nach dem Exempel des cananäischen Weibleins, zu welcher Christus Matth. 15, 28. sagt: „O Weib, dein Glaube ist groß“; item, wie er sagt Marc. 9, 23.: „Alle Dinge sind möglich dem, der da glaubet.“

224. Diese Allegorie ist nützlich und gut, welche die Gewissen unterrichtet und stärkt. Dieselbige soll man in die Uebung bringen und soll sie allzeit vor Augen haben, daß wir es gewißlich dafür halten, daß ein gläubiger Mensch Gott mit seinem Glauben und Gebet überwindet. Denn Gott hat verheißen, daß er die Gläubigen schützen, ihnen helfen und alles Gute geben wolle. Darum will und kann er sich nicht leugnen. Wo er aber in einer andern Gestalt erscheint, oder daß du es anders empfindest, und daß dich dünkt, daß er dir widerstrebe, so sollst du in deinem Herzen davor nicht erschrecken, und sollst auch nicht weichen, sondern im Glauben wider ihn kämpfen, daß du mögest obliegen und Israel werden. Wie aber? Zwar nicht durch die Kräfte oder Waffen deines Fleisches und deiner Natur, sondern durch das Vertrauen auf die Sache, die du mit Gott hast, nämlich, daß er dir verheißen und geschworen hat, er wolle dein Gott sein. Mit solchem Vertrauen wirst du obliegen, welches nicht von Natur, sondern von der Verheißung herkommt. Wenn er dir nun begegnet wie ein Kämpfer und dich verderben will, oder seinen Namen und Verheißung verbergen, so sei stark und halte fest am Worte, ob du wohl große Schwachheit fühlst: so wirst du obliegen.

225. Darnach wirst du in diesem Kampfe auch das empfinden, daß die Spannader oder das Gelenk der Hüfte verrenkt wird und verdorrt. Diese Allegorie ist recht, und reimt sich mit der Art und Weise zu reden, so die Schrift hat, die oftmals durch die Hüfte die Nachkommen versteht. Darum hat Augustin die Spannader der Hüfte nicht übel auf die Nachkommen gezogen. Denn also redet Mose im andern Buch am 1. Cap. V. 5.: „Und aller Seelen, die aus den Lenden Jakobs kommen waren, der waren siebenzig.“

226. Die Spannader aber, so verderbt und lahm geworden ist, lehrt die, so da kämpfen und gläubig sind, daß die fleischliche Geburt zu diesem Siege gar untauglich ist; denn die Hüfte, welche die fleischliche Geburt bedeutet, ist verrenkt und verdorret. Daher wissen die Christen, daß sie nicht Gottes Volk oder Kinder Gottes werden durch die Geburt oder Gerechtigkeit des Fleisches, auch nicht, daß es ihnen von Fleische angeerbt wäre, wie Joh. 1, 13. geschrieben stehet: „Welche nicht von dem Geblüt, noch von dem Willen eines Mannes geboren sind“; es ist alles getödtet. Darum rühmen sich die Heuchler und Juden solcher Dinge vergeblich, und sind doch die Juden noch jetzt gar unsinnig der fleischlichen Ehre und Gerechtigkeit des Gesetzes halben. Denn dieser Mann hat mit der Hand die Hüfte Jakobs angerührt, dadurch ist die Spannader getödtet worden. Die Kinder Gottes aber werden von Gott geboren. Also sollen nun die Gottseligen und die Nachkommen Jakobs lernen, daß sie sich des Fleisches nicht rühmen, und auch das Leben in den Wohlthaten und Ruhme des Fleisches nicht suchen, sondern im Segen Gottes; der macht uns zu Kindern Gottes nach der Verheißung 1 Mos. 22, 18.: „Durch deinen Samen sollen alle Völker gesegnet werden“; item Joh. 1, 13.: „Die von Gott geboren sind“ etc. Dieses ist ein fein Stempel, Allegorien zu suchen; denn es lehrt und tröstet zugleich.

227. Und das reimt sich auch wohl dazu, daß die Kinder Israel keine Spannader essen, daß ist, sie glauben nicht. Denn „essen“ heißt an anderen Orten glauben, als Joh. 6, 53.: „Werdet ihr nicht essen das Fleisch des Menschensohns“ etc., „so habt ihr kein Leben in euch.“ Wir essen nicht, glauben auch nicht, daß die Spannader über die Hüfte erhöhet ist. Wir predigen die Weisheit, Ehre, Gerechtigkeit und Gewalt des Fleisches nicht, sondern enthalten uns von allem Vertrauen und Rühmen fleischlicher Gerechtigkeit. Diese Weisheit verstehen und haben allein die guten Nachkommen Jakobs. Denn denselben ist allein die Hüfte verrenkt, die Andern alle ehren und essen die Spannader der Hüfte, die sehr zäh und hart ist; denn sie ist ihnen noch nicht getödtet. Allein die Christen enthalten sich der verrenkten Hüfte und Spannader. Denn also lehrt das Neue Testament Joh. 1, 13.: „Welche nicht von dem Geblüt, noch von dem Willen eines Mannes, noch dem Willen des Fleisches, sondern von Gott geboren sind.“ Diese Allegorie ist besser und näher, die sich auch mit der Lehre des Evangeliums reimt. Denn das Volk Gottes kämpft mit Gott, und siegt und empfängt nach dem Siege den Segen, und verwirft die Gerechtigkeit und Weisheit des Fleisches, ja, es kreuzigt das Fleisch sammt seinen Lüsten und Begierden, Gal. 5, 24. Wer gottseelig ist und die Gabe hat, Allegorien zu dichten, der wird leichtlich andere finden können.

\quellseite{217}

\chapter{Das dreiunddreißigste Capitel}

\section[Erster Theil]{Erster Theil. Wie Esau dem Jakob mit einem versöhnten und geänderten Herzen entgegen kommt; und wie sich Jakob hierbei verhält}

\subsection{I}

\begin{quote}
\textbf{V. 1. Jakob hub seine Augen auf, und sahe seinen Bruder Esau kommen mit vier hundert Mann.}
\end{quote}

1. Was in diesem Capitel erzählt wird, ist an sich klar und leicht; denn es sind solche Dinge, die dieses Leben, bürgerliche und äußerliche Sitten belangen. Darum wollen wir es kürzlich überlaufen, dieweil weder die Worte noch die Historie an sich dunkel oder schwer ist. Denn Mose beschreibt zum Exempel und Stärkung unseres Glaubens, wie Gott das Seufzen des Glaubens, der da kämpft, ob er schwach ist, erhöre; welches wohl ein unaussprechlich Seufzen ist, aber doch nicht ohne große Frucht.

2. So gehört nun Jakob in die Zahl derer, von welchen Christus sagt Marc. 9, 23.: „Alle Dinge sind möglich dem, der da glaubet.“ Denn er hat durch den Glauben Gott und Menschen überwunden, also daß ihm weder Gott noch Mensch wollte oder konnte Schaden zufügen. Gott hat ihn gesegnet; sein Bruder Esau aber ist dermaßen geändert worden, daß er ihm nicht allein keinen Schaden thun will, sondern will ihm auch helfen, ihn lieben und wohlthun; sein Zorn ist verändert worden in eine freundliche Brüderschaft.

3. Ja, das ist die Hand des Höchsten, dadurch der, so zornig ist, verhindert wird, daß er nicht thun kann, was er sich vorgenommen hatte. Denn Gott hat aller Menschen Herzen in seiner Hand, daß sie nicht über das Ziel, so ihnen gesetzt ist, schreiten können. Wie im Buche Hiob am 38. Cap. V. 10. 11. steht, daß Gott dem Meer mit seinem Damm Riegel und Thür gesetzt, und gesagt habe: „Bis hieher sollst du kommen, und nicht weiter; hie sollen sich legen deine stolzen Wellen.“ Wie das Meer anstehen, wenn es von Wellen und Sturmwinden ungestüm ist, lassen sich dünken, daß es den Damm, der darumher gezogen ist, zerbrechen und allenthalben überfließen werde. Und werden doch gleichwohl solche Wassermogen und Wellen bezwungen und inne gehalten innerhalb ihres Riegels und ihrer Thür. Also toben die Herzen der Menschen auch greulich, wenn sie zornig sind; aber Gott hat ihnen ein Ziel ihres Grimmes und Zornes gesetzt, darüber sie nicht schreiten dürfen. Dasselbe sieht man an diesem Exempel Esaus auch. Und man hat noch andere mehr hin und wieder in den Historien der Heiden, die da anzeigen, daß die Weisheit und Gewalt der Menschen niemals über das Ziel, so ihr von Gott gesteckt gewesen, habe kommen können.

4. Hannibal war fast an die Mauern der Stadt Rom gekommen, darüber beide die Stadt und auch ganz Italien dermaßen erschrocken und bestürzt worden, daß er gewisse Hoffnung und Macht gehabt, beides zu erobern, wenn er nur einen solchen Muth gehabt hätte, wie er zuvor im Streit gezeigt hatte. Aber er zog wieder zurück, nämlich, da Gott sein grimmiges Vornehmen gebrochen und gehindert hat. Darum schrieen seine Freunde und sprachen: Vincere scis, Hannibal, victoria nescis uti: O Hannibal, siegen kannst du wohl, aber du weißt des Sieges nicht zu gebrauchen. Sie haben aber nicht gewußt, daß darunter eine geheime und göttliche Kraft verborgen war, die ihn bezwungen hat. Also ist der Pabst bis an sein Ziel gekommen; desgleichen werden die Türken auch nicht weiter kommen, denn ihnen das Ziel von Gott gesteckt ist. Ja, es wird des Teufels Grimm und Toben auch so lange währen, als es Gott gefallen wird, daß es ihm zu seiner Ehre und unserer Tödung und Seligkeit dienen kann.

5. Eben auf solche Weise redet der Herr auch vom Könige der Assyrer, Jes. 37, 24.: „Du sprichst: Ich bin durch die Menge meiner Wagen herauf gezogen auf die Höhe der Berge, an den Seiten Libanons; und habe seine hohen Cedern abgehauen“ etc. Aber der Herr antwortet ihm V. 29. also: „Weil du denn wider mich tobest, und dein Stolz herauf vor meine Ohren kommen ist, will ich dir einen Ring an die Nase legen, und ein Gebiß in dein Maul, und will dich des Weges wieder heimführen, deß du kommen bist.“ Dem hat es weder am Willen, noch an der Macht gemangelt, aber die Macht und Gewalt Gottes hat ihm gewehret. Also spricht der Herr Jes. 10, 15.: „Mag sich auch eine Axt rühmen wider den, so damit hauet? oder eine Säge trotzen wider den, so sie zeucht?“ Du darfst dich nicht über die Maßen erheben um deines Glückes und Gewalt willen. Denn es ist eben so viel, als wenn die Säge oder Axt rühmen wollten, daß sie das Holz haueten oder von einander sägeten. Denn wenn ich die Säge nicht führte und das Holz hauete, so würde die Axt oder Säge nimmermehr etwas ausrichten.

6. Also war dieser Sieg Jakobs von Gott auch verheißen umsonst und nur aus lauter Gnade, nämlich, daß er seinen Bruder überwinden würde. Darum wird Esau nicht überwunden durch menschliche Kraft, Geschicklichkeit, Rath, böse Tücke, oder daß sich Jakob höflich gegen ihn erzeigt hätte, als wäre er sein Freund, sondern allein durch Gottes Gnade und Güte; denn sein Wille wird geändert.

7. Und das ist der allerherrlichste Sieg unter allen, wo die Menschen dahin gebracht werden, daß ihr Herz und Wille geändert wird. Und ist nie ein besserer Sieg gewesen, als wenn aus unwilligen und widerspenstigen Leuten gutwillige Leute geworden sind. Das ist aber allein ein Werk göttlicher Gewalt und Majestät, nämlich, ein Herz und Muth, der zornig und grimmig ist, dermaßen ändern, daß er freundlich und gütig wird. Denn sonst ist das menschliche Herz so gar unbändig, daß es durch keinerlei Gewalt geändert werden kann: nicht durch des Teufels Macht, oder des Todes, oder sonst irgend eines Unglücks, wie stark es auch sein möchte, es sei denn, daß es mit Gewalt gezwungen werde. Denn des Menschen Wille ist unüberwindlich und verhärtet, und kann durch gar keine Weisheit, auch durch keinerlei Kraft erweicht werden, wo er nicht mit Güte überwunden wird. Das kann aber allein Gott thun, wie an diesem Exempel zu sehen ist.

8. Denn Jakob ist billig Israel genannt worden, weil er Gott überwunden, und empfunden hat, daß er mit ihm lacht und fröhlich war, den er zuvor zum Feinde gehabt und der ihm nach dem Leben gestanden hatte. Zudem hat er auch den Zorn seines Bruders Esau, und aller, die mit ihm gekommen sind, überwunden mit nichts Anderem, als mit dem Glauben und Gebete, damit Gott selbst und die Menschen überwunden werden. Daher kommt es, daß Salomo sagt in den Sprüchen am 16. Cap. V. 7.: „Wenn jemandes Wege dem Herrn wohlgefallen, so macht er auch seine Feinde mit ihm zufrieden.“

9. Und dasselbe haben wir bisher auch erfahren. Denn deren sind von Gottes Gnaden Viele mit uns zufrieden gemacht worden, die zuvor dem Evangelium heftig feind waren und dasselbe verfolgten; und das ist durch keine Gewalt geschehen, sondern mit sehr angenehmen und schönem Siege, dadurch der Wille des Menschen von selber geändert wird. So große Kraft hat das Gebet: damit hat Jakob beide Gott und seinen Bruder Esau überwunden; er hat Himmel und Erde gewonnen. Sein Bruder Esau gedenkt jetzt nicht einmal ihn zu verletzen, und kann es auch nicht thun. Denn alle sein Sinn und Wille ist geändert worden; wie wir hören werden.

10. Darum ist nun Jakob bei sich selbst fröhlich und guter Dinge über den Sieg, den er von Gott selbst hat und nicht vom Engel; wie er zuvor gesagt hat: „Ich habe den Herrn von Angesicht gesehen“ \&c. Und diese Versöhnung seines Bruders zeigt vielerlei, nämlich, daß dieses ein wahrhaftig Gesicht und nicht irgend ein vergeblich Gespenst gewesen ist, und darnach, daß auch das Gebet und das unaussprechliche Seufzen mehr erlangt hat, als Jakob hat bitten oder hoffen dürfen. Denn er hat sich selbst so viel Trostes weder bei Gott, noch bei seinem Bruder jemals verheißen. Er hat nur darum, daß sein Bruder abziehen möchte und ihn mit seinem Gesinde unverletzt lassen. Er hat zwar die große Freundschaft nicht erwarten dürfen, daß ihm sein Bruder Esau entgegen laufen, ihn herzen und um den Hals fallen, und ihn dazu auch noch küssen und weinen würde.

\quellseite{218}

11. Deshalb soll man glauben und hoffen, und im Gebet beharren und anklopfen, der Herr wird gewißlich kommen und nicht verziehen. Denn es ist unmöglich, daß ein solch Gebet, damit man immer anhält und darin man beharrt, nicht sollte erhört werden. Dieweil wir aber nicht glauben, darum erfahren wir auch diese Beharrung nicht, empfinden dazu die große Güte und Hülfe Gottes auch nicht. Darum lasset uns unsere Herzen aufwecken zum Glauben und zum Gebet; und wir sollen wissen, daß Gott dazu Lust hat, wo man also beharrt, und daß er die Beharrung geboten hat, da er Matth. 7, 7. sagt: „Suchet, so werdet ihr finden; klopfet an, so wird euch aufgethan.“ Und das Gebet wird nimmer reichlicher erhört, denn in solchem Kampf und Seufzen des Glaubens, der also mit Gott kämpft.

12. Es hat die Kirche und Gemeinde Gottes bisher Linderung vieles großen Unglücks erlangt, so vorhanden gewesen vom Türken und von andern Widerfachern; darum ist unser Gebet viel anders und reichlicher erhört worden, denn wir gebeten oder selbst verstanden haben. Wie St. Paulus sagt Eph. 3, 20. und Röm. 8, 26. 27.: „Wir wissen nicht, was wir beten sollen, sondern der Geist selbst vertritt uns aufs beste mit unaussprechlichem Seufzen. Der aber die Herzen forschet, der weiß, was des Geistes Sinn sei.“ Denn wir bitten allezeit weniger, als wir bitten sollen; und wir hoffen darauf nicht, daß uns Gott aus gutem geneigtem Willen gibt. Darum bitten wir nicht, wie es sich gebühret, und verstehen es auch nicht; denn es sind größere Dinge, als wir in unserem Herzen begreifen können; denn wir nur auf kleine und geringe Dinge denken. Der Herr ist groß und hoch, darum will er haben, daß man auch von ihm was Großes bitten soll, und will solches auch gern geben, auf daß dadurch seine Allmacht bewiesen werde. Dieses ist das erste Stück vom Siege Jakobs.

\subsection{II}

13. Zum Andern, auf daß sich Jakob davor hüten möge, daß er Gott nicht versuche, und daß er nicht stolz werde von der großen Freude und von dem herrlichen Siege, und also seines Bruders Herz, das noch kaum erst gestillet und versöhnt war, erbittern möchte: so stolzirt oder trotzt er gegen seinen Bruder nicht, wie die Narren pflegen zu thun. Er sagt nicht also: Ich habe die Verheißung vom Siege, ich frage nach deinem Zorn und Drohen nichts; denn von solcher Verachtung und Stolz würde die Wunde, so erst mit schwerer Mühe und Arbeit geheilet war, wiederum aufgebrochen worden sein: sondern er gebraucht aller Freundlichkeit, wie er nur kann und mag, seines Bruders Freundschaft zu mehren und zu erhalten; er richtet eine andere sehr schöne Procession und Gesänge an, damit er ihm entgegen ziehet mit seinem Gesinde. Das heißt thun, was an Einem ist (facere, quod in se est), auf daß wir den Segen Gottes nicht verlieren.

14. Wenn wir durch Gottes Gnade und Güte reichen Trost und Freude haben, so sollen wir nicht trotzen, oder uns gehen über die Maßen erheben, sondern sollen des Segens Gottes gebrauchen mit Furcht, mit Demuth und Danksagung; wie der 2. Psalm V. 11. lehrt: „Dienet dem Herrn mit Furcht und freuet euch mit Zittern“, auf daß wir nicht in Sicherheit und Hoffahrt fallen und Gott versuchen.

15. So ehrt nun Jakob seinen Bruder, und gebraucht die Mittel, die er gebrauchen kann, seine Freundschaft zu behalten; und geht ihm damit an dem Siege gar nichts ab, wiewohl er sich gegen seinen Bruder demüthig und freundlich erzeigt. Er theilt aber sein Gesinde in die Ordnung, in welcher er, da das Geschenk vorausgeschickt war, seinem Bruder begegnet; und dasselbe thut er mit allem Fleiß, auf daß er das Geschenk zieren möge, Esau damit zu erweichen und seinen Zorn zu lindern, wiewohl er des Sieges jetzt schon gewiß ist.

\begin{quote}
\textbf{V. 1—3. Und theilete seine Kinder zu Lea, und zu Rahel, und zu beiden Mägden; und stellete die Mägde mit ihren Kindern vorne an, und Lea mit ihren Kindern hernach, und Rahel mit Joseph zuletzt; und er ging vor ihnen her, und neigete sich siebenmal auf die Erde, bis er zu seinem Bruder kam.}
\end{quote}

16. Hier theilt er sein Volk mit einem ganz andern Herzen und auf eine ganz andere Weise, als er droben gethan hat. Denn da war er noch ungewiß, was sein Bruder für einen Willen und Gemüth gegen ihn gehabt. Darum war er erschrocken und verzagt, und rathschlagte also bei sich selbst, daß er sein Volk in zwei Haufen theilte, und dachte: Wo Esau das eine Heer schlagen wird, so wird das andere entrinnen. Er blieb allein dahinten am Ufer stehen und kämpfte mit der Verzweiflung. Da war das Fleisch schwach und wanckend und war um der Schwachheit willen des Glaubens zur Verzweiflung geneiget. Jetzt aber tritt er nun zu seinem Bruder mit fröhlichem Herzen, er theilt sein Hausgesinde in eine Ordnung nicht mit verzagtem, sondern mit fröhlichem und sicherem Herzen; aber doch also, daß er Gott nicht versuche, noch auch seinen Bruder aufs neue erzürne, den er mit dem schwachen und kämpfenden Glauben überwunden hatte.

17. Darum hat er diese Theilung des Volkes nur dazu angerichtet, daß er damit ein Gepränge machen wolle und daß es ein Anzeichen der Freundschaft sein sollte. Und bleibt nun nicht mehr dahinten, sondern geht vor den Andern allen her, und zieht seinem Bruder entgegen. Denn die Anfechtung ist nun überwunden und die Sonne des Trostes und der Freude geht ihm auf. Und ist ein sehr schönes Gepränge. Denn die freundlichen lieben Weiber, die lieben Kindlein und Mägde folgen dem Hausvater nach, und da sie vor seinen Bruder gekommen sind, neigt er sich siebenmal auf die Erde, bis er zu ihm kommt.

18. Das hebräische Wort für „anbeten“, welches in der heiligen Schrift sehr gemein ist, kommt von dem Wort schachach her, das heißt, sich neigen; wie im 38. Psalm V. 7. steht: „Ich gehe krumm und sehr gebückt: den ganzen Tag gehe ich traurig.“ Von solcher Gebärde des Rückens hat die Anbetung, das heißt, Ehrerbietung, so denen erzeigt wird, die größer und höher sind denn wir, oder auch die man GOTT erzeiget, ihren Namen; und geschieht entweder, daß man steht und sich neigt, oder aber daß man die Kniee beugt. Denn diese Gebärde ist ein Anzeichen der Ehrerbietung und der Demuth, so denen erzeigt wird, die wir dafür erkennen, daß sie Ehren werth sind. Wir pflegen unserer Gewohnheit nach das Haupt zu entblößen und den Hut abzuziehen, welches die Juden nicht thun, wenn sie auch schon im Tempel oder in der Kirche sind. So nennen wir nun „Anbetung“ alle äußerliche Ehrerbietung, so denen erzeigt wird, die größer sind als wir, oder aber auch wohl unsers Gleichen sind, es geschehe gleich, daß man ein Knie beuge, wie bei den Fürsten, oder aber daß man beide Kniee beuge, wie im Gebet geschieht. Auf solche Weise neigt sich nun Jakob siebenmal vor seinem Bruder als dem Aelteren und Herrn, des Wohlthat und Freundlichkeit er erkennt, daß er mit ihm versöhnt ist, und sich nun recht brüderlich und freundlich vernehmen läßt.

19. Mithin fällt demnach eine Frage ein: Wie sich dieses reime mit dem Segen, so droben 1 Mos. 25, 23. ist verheißen worden, daß der Größere dem Kleinern dienen sollte? und Cap. 27, 29., daß Jakob ein Herr sein würde über seine Brüder? Denn es scheint, als ob mit dieser Demuth, da sich Jakob vor seinem Bruder also neigt, der Segen umgekehrt und gar geändert werde. Und dieses Argument kann auf mancherlei Weise groß gemacht und bekräftigt werden. Denn man liest an keinem Orte, daß sich Esau vor Jakob geneigt und ihn für seinen Obersten oder Herrn erkannt habe. Das Gegentheil geschieht aber an diesem Ort von Jakob, der sich vor Esau neigt. Ueber das hatte Esau eine viel größere Gewalt. Jakob führte seine Kindlein mit sich, unter welchen Ruben der Aelteste war, von zwölf oder aufs höchste von dreizehn Jahren, die andern waren jünger.

20. Dagegen aber Esau hat Kindeskinder, und die vielleicht auch noch weiter Kinder gehabt haben und unter denen Fürsten und Herren gewesen sind, die er gezeugt gehabt. Denn er hat im vierzigsten Jahr seines Alters Weiber genommen. Von der Zeit an hat er viel Kinder gezeugt, bis Jakob ist sieben und siebenzig Jahre alt geworden, in welchem Jahr er erstlich den Segen empfangen hat. Ja, Esau hat schon viel Kinder gehabt, da er noch in Jakobs Vaters Hause war, ehe er ihn den Segen genommen hat. Darum ist er ein Herr und Fürst gewesen, der den Heiden im Lande Canaan durch Schwägerschaft verwandt und mit den Ismaëliten und andern Fürsten im Bunde gewesen ist. Solchen großen Fürsten begegnet Jakob mit seinen Kindern und einem armen, wehrlosen Haufen, da er aus der Dienstbarkeit und Elend kommt, und ist viel geringer, als Esau, der daher zieht mit vierhundert gerüsteter Mannschaft.

\quellseite{219}

21. Hernach aber wird Jakob wiederum in Egypten ein Fremdling sein müssen. Unterdeß hat Edom viel treffliche Fürsten und Könige bekommen und ist seine Gewalt täglich größer geworden; ja, da sie aus Egypten ziehen, ist den Kindern Israel verboten worden, das Land Edom zu berühren, und der Herr sagt, er wolle ihnen davon auch nicht einen Fuß breit geben. Ist nun das das Herrschen, so Jakob treiben soll? Davon ist der Erste, der die Edomiter bezwungen hat, 2 Sam. 8, 14.; aber vor seiner Zeit waren sie sechshundert und vier und zwanzig Jahr frei gewesen, und sind doch gleichwohl unter der Gewalt des Volkes Israel nicht länger geblieben als hundert und zwanzig Jahre. Denn unter dem Könige Joram haben sie das Joch wiederum abgelegt und einen eigenen König gemacht, und haben den Königen Israel keinen Tribut mehr gegeben, 2 Kön. 8, 20. Derhalben läßt es sich ansehen, als sei der Segen aufgehoben und umgekehrt worden.

22. Etliche antworten auf diese Frage also: Daß dieser Segen auf Christum gehe, in dem er endlich erfüllt worden sei. Man muß aber auf diese Historie antworten. Denn diese Unterthänigkeit Jakobs, und daß er sich also gedemüthigt hat, welches die ganze Zeit seines Elendes und Pilgerschaft gewährt hat, ist nur eine Versuchung, wie gesagt ist. Wie wir Christen die allerelendesten Menschen und jedermann unterworfen sind: es ist aber keine Verwerfung, sondern nur ein Kampf; denn in Wahrheit sind und bleiben wir auch Herren Himmels und der Erde. Die Versuchung ist keine Veränderung, oder daß dadurch die verheißenen Güter sollten umgestoßen werden. Wie das Gefängniß des Volkes Israel unter dem Könige der Assyrer auch eine Versuchung gewesen ist, aber nicht eine Veränderung des Volkes. Also in den Schrecken der Sünde und des Todes vertausche ich den Tod nicht mit dem Leben, Christum mit dem Teufel; wiewohl es in der Anfechtung nicht anders scheint. Also ist Abraham auch versucht worden 1 Mos. 22, 2.: da ihm von Gott geboten ward, daß er seinen Sohn Isaak opfern sollte, verliert er damit seinen Sohn nicht.

23. Und also sind ingemein alle Anfechtungen; wenn man sie erst ansieht, scheinen sie sehr traurig, erschrecklich und die den Verheißungen gar ungleich sind: aber der Ausgang reimet sich doch endlich mit Gottes Worte, welches fest und unbeweglich bleibt. Wie es pflegt zu gehen, wenn sich plötzlich ein Ungewitter erhebt, da läßt es sich auch ansehen, daß der greuliche Donner und Blitz Himmel und Erde, Menschen und Vieh bedrohe, daß Alles über einen Haufen fallen und untergehen solle; aber es bringt so gar keine Hindernisse oder Veränderung, daß es auch das Land feuchtet und fruchtbar macht.

24. So viel aber die Historie an sich belanget, ist zwar Israel niemals unter der Gewalt Edoms gewesen: Esau hat nie über das Volk Israel herrschen können; ja, zur Zeit Davids bis auf den König Joram ist Edom den Kindern Israel unterworfen gewesen. Derhalben ist der Segen fest und gewiß stehen geblieben und das Regiment hat bei den Nachkommen Jakobs gewähret bis auf Herodes und die Zukunft des rechten wahrhaftigen Israels, nämlich, des Messias. Wie die Prophezeiung lautet 1 Mos. 49, 10., da Jakob sagt: „Es wird das Scepter von Juda nicht entwendet werden, bis daß der Held komme.“ Da nun die Zukunft desselben vorhanden war, so hat auch die fleischliche Herrschaft Israels aufgehöret und da hat erst Herodes das Regiment eingenommen. Der Segen ist übrigens ununterdrückt geblieben bis auf die Zukunft des Heilandes, der mit Jakob gerungen hat. Da ist ein neu und besser Reich angefangen und angerichtet worden, darin der rechte Israel regiert über Himmel und Erde, über Engel und Teufel. Denn im Namen Jesu beugen sich alle derer Kniee, die im Himmel und auf Erden, und unter der Erde sind, Phil. 2, 10. Da ist der Segen Jakobs recht erfüllt worden, und ist doch gleichwohl unter der Figur vor der Zukunft Christi auch ununterdrückt geblieben, also daß Israel den Edomitern nie ist unterworfen worden; vielmehr es von ihnen oftmals angefochten und geplagt worden ist.

25. Daß sich aber an diesem Ort Jakob also demüthigt und vor seinem Bruder neigt und ihn einen Herrn nennt, soll man nicht also verstehen, als ob ihm damit an seinem Segen und Ehre etwas abgebrochen wäre; sondern, wie man gemeiniglich pflegt zu sagen, es sind Ehrenworte, die binden nicht. Und diese Ehrbietung oder Gehorsam bringt nicht mit sich, daß ihm damit eben das Erbe und die Herrschaft wäre übergeben worden. Wiewohl Jakob sich einen Knecht, Esau aber einen Herrn nennt: so hat doch Esau darum das Regiment noch nicht und hat dasselbe auch niemals bekommen.

26. Darum kommt diese Ehrbietung aus sonderlicher Demuth, auf daß es nicht scheine, als ob Jakob wegen dieses Friedens und Glückes sich überhoben und damit stolziert hätte; und kommt nicht davon her, als sollte dadurch der Segen ihm entwendet sein, und ist nur ein Ehrengedicht. Darnach pflegt das auch in diesem Leben zu geschehen, daß die, so höher und größer sind als Andere, nicht allein ihres Gleichen, sondern auch wohl die geringer sind als sie, freundlich und ehrlich anreden und ihnen ihren Dienst anbieten.

27. Denn also pflegt die Liebe auch denen, die geringer sind als wir, gleichwohl allerlei Dienst und Ehrerbietung zu erzeigen, ihnen auch zu helfen und zu fördern. Also thut ein Pfarrherr oder Kirchendiener: wenn er irgend Einen sieht, der da angefochten und betrübt ist mit Schrecken des Gewissens und der in solcher Noth Trost von ihm begehrt, denselbigen betrübten Menschen bietet er auch allerlei geringen und knechtlichen Dienst an, auf daß er das arme, schwache und verzagte Herz aufrichten und stärken möge. Aber auf solche Weise übergibt er ihm gleichwohl das Kirchenregiment nicht, sondern er behält selbst das Amt, so ihm befohlen ist.

28. Das sind Werke der Gottseligkeit und Freundlichkeit, damit man oftmals großen Unwillen verhüten kann. Und sind solche Werke und Dienstbarkeit von Gott nicht verboten; ja, uns ist geboten, daß wir uns Andern unterwerfen sollen, die uns gleich sind, oder auch die geringer sind als wir, dem Nächsten zu Nutz und Besserung. Und mit solcher Demuth geht uns an unserer Ehre gar nichts ab; wiewohl wir uns etwas lenken und neigen einem Andern zu Nutz und Wohlfahrt, ja, auch wohl zu Gottes Ehre. So bleibt nun Jakob der Segen ununterdrückt und fest bestehen. Und dieses alles, dadurch er scheint seinen Bruder unterworfen zu werden, soll nicht eine Veränderung, sondern nur eine Versuchung genannt werden.

\section[Zweiter Theil]{Zweiter Theil. Von der Ausführung Esaus und Jakobs; von ihrem Gespräch, und von dem Geschenke, so Jakob dem Esau gegeben}

\subsection{I}

\begin{quote}
\textbf{V. 4. Esau aber lief ihm entgegen, und herzte ihn, und fiel ihm um den Hals, und küssete ihn; und sie weineten.}
\end{quote}

29. Mose gebraucht volle und nachdrückliche Worte, die Versöhnung der zwei Brüder zu beschreiben, also auch daß daran niemand zweifeln darf, daß Esau nicht Alles von Herzen thun und reden sollte, und daß er nicht sollte in Wahrheit geändert sein; denn er läuft zu seinem Bruder, herzt ihn und fällt ihm um den Hals, küsset ihn, und können sich Beide vor Freude des Weinens nicht enthalten. Denn das brüderliche Herz thut sich auf, und ihrer beider Herz entbrennt in rechter Liebe und Freundschaft, daß sie alles Zorns und Unbilligkeit vergessen. Es wird auch der vierhundert Mann nicht mehr gedacht und man sieht kein Zeichen mehr des Zorns oder Widerwillens, sondern es ist alles voll rechter und brüderlicher Liebe und Freundschaft. Denn Mose gebraucht solche Worte, die da anzeigen, daß dieses alles aufrichtig und getreulich zugegangen und geredet worden ist.

30. Denn wiewohl auch die Heuchler sich stellen können, als haben sie Einen lieb, so zeigen es doch allhier die Gebärden und alle Worte insbesondere an, daß an ihnen beiden ganz brennende Liebe gewesen sei, und daß Esau in seinem Herzen recht besänftigt und versöhnt gewesen ist. Denn er denkt also: Was sollt du mit diesem Manne machen? Es ist ja mein rechter leiblicher Bruder. Wo es nun Gott also gefallen hat, daß er hat den Segen vom Vater empfangen sollen und mir denselben entziehen, so will ich es ihm nicht mißgönnen, will ihn darum nicht neiden oder mit ihm zürnen, er mag ihn immerhin haben und für sich behalten.

\quellseite{220}

31. Die Andern aber werden gedacht haben: Ach, der arme elende Jakob, der vor zwanzig Jahren rühmte, daß er den Segen empfangen und Esau entzogen hätte, der muß sich nun jetzt vor unserm Fürsten und Herrn Esau auf das allerdemüthigste neigen. Wo ist nun der große Ruhm und die herrliche Hoffnung Jakobs? Wir haben den rechten Segen; denn das Regiment ist bei Esau laut der Verheißung 1 Mos. 27, 39. 40. „Deines Schwertes wirst du dich nähren. Und es wird geschehen, daß du auch ein Herr, und das Joch von deinem Halse reißen wirst. Du wirst eine fette Wohnung haben auf Erden, und vom Thau des Himmels von oben her.“ Die fette Wohnung auf Erden, desgleichen Gold und Silber haben wir, werden sie gesagt haben; Esau nähret sich seines Schwertes, ist ein trefflicher Kriegsmann, der sich seinen Bruder unterworfen hat, und wird andere Völker auch unter sein Joch bringen; er hat das Joch seines Bruders vom Halse gerissen.

32. Und ist wohl glaublich, daß Esau auf dieselbige Meinung gelehrt hat von der Verheißung, die ihm geschehen war; er wird sie ohne Zweifel groß gemacht haben, wie fleißig und ernst er immer gekonnt hat. „Denn die Kinder dieser Welt sind klüger, denn die Kinder des Lichtes“, Luc. 16, 8., und sie können die Sprüche über die Maßen wohl drehen und wenden, wo deren etliche sind, die sich dünken lassen, daß sie ihnen auf ihrer Seite dienen; ja, sie können aus einem Zweiglein einen ganzen Wald, aus einem Blümlein eine ganze grüne Wiese und aus einem Worte die ganze Bibel machen.

33. Wir, die wir die rechten und viel reicheren Verheißungen und Segen haben, verwundern und freuen uns derselben so hoch nicht; wir bringen nicht so fest darauf; wir breiten die großen herrlichen Dinge nicht so aus, wie die Gottlosen ihre Güter und Gaben müssen groß und herrlich zu machen. Denn siehe die päpstlichen an, wie die so heftig dringen auf den einigen Spruch und denselbigen müssen groß zu machen: „Du bist Petrus, und auf diesen Fels will ich bauen meine Gemeine“ etc., Matth. 16, 18., den Primat des Papstes damit zu verteidigen: da hat man ganze Bücher von Büchern, darin dieser einzige Spruch ausgelegt wird. Aber dagegen, wie kalt und nachlässig gehen wir mit der Lehre des Evangeliums um? Wie ganz schläfrig lesen oder hören wir, daß Christus für uns gelitten hat und für unsere Sünde gestorben ist, 1 Petr. 3, 18.; item, daß geschrieben steht Gal. 3, 11.: „Der Gerechte wird seines Glaubens leben.“ Es sind die Lehrer der andern Künste etwas fleißiger, die Lehre in ihrer Profession zu üben, und ist auch nützlich und nöthig in diesem Leben, daß die guten Künste getrieben und die, so einen guten Verstand haben, darin fleißig geübt werden. Also finden wir an den Krämern und Kaufleuten, wie sie so fleißig und genau rechnen können. Wir aber, sage ich, sind sehr faul und schläfrig, unsere Güter und Gaben zu preisen und die seligen groß zu machen, welche Faulheit und Nachlässigkeit von unserem Fleische herkommt. Darnach, so handeln wir auch von verborgenen und himmlischen Dingen, die man nicht sieht oder fühlt: sie aber handeln von gegenwärtigen, sichtbaren und greiflichen Dingen; darum sind sie so viel wacher und fleißiger darauf.

34. Auf solche Weise hat nun Esau die Verheißung ganz fleißig überlegt und erwogen. Denn die Erfahrung hat sich auf die Verheißung gereimt. Mir ist das Schwert verheißen, wird er gesagt haben, darum bin ich billig ein Kriegsmann und ein Oberster, mein Bruder aber ist ein Hirt, ist arm und elend; das fette und fruchtbare Land gebührt mir, und ich bin keinem Menschen unterworfen und habe jetzt das Joch meines Bruders auch vom Halse gerissen. Solcher Ruhm und Trotz kommt vom fleischlichen Segen; denn derselbe steht jedermann vor Augen und hat in der Welt einen großen Schein und Gespränge. Jakob hat zwar auch einen leiblichen Segen gehabt, aber der hatte solchen großen Schein nicht. Denn Gott hat Könige und Fürsten gehabt vor seinem Bruder Israel, und der geistliche Segen ist allezeit bedeckt unter einer sehr geringen und elenden Gestalt.

35. Und dieser Ruhm Esaus ist sehr nützlich gewesen zur Versöhnung. Denn er hat gedacht: Ich bin der Herr, und dieweil ich mich vor meinem Bruder nichts zu besorgen habe, warum sollte ich denn unfreundlich mit ihm umgehen? Ich habe nicht allein die Worte des Segens, wie Jakob hat, sondern ich habe auch die That zugleich mit den Worten. Denn ich habe viel Geld und Gut, große Gewalt, viele Güter, Frieden, ich bin geboren vom Fleisch und Geblüte der heiligen Väter. Darum hat mich Gott nicht verworfen, wiewohl mir Jakob den Segen weggenommen und entzogen hat.

36. Und es ist wohl glaublich, daß Esau noch endlich selig geworden sei. Darum hat aber Gott diesen Unterschied gemacht zwischen Jakob und Esau, item, zwischen Isaak und Ismael, auf daß er damit anzeigen möchte, daß sein Reich schlechthin geistlich ist. Und ist daran kein Zweifel, es werden ihrer Viele vom Geschlechte und Nachkommen Esaus selig geworden sein. Denn der Edomiter haben sich viele zum Volk Israel geschlagen und sind beschnitten worden; es sind ihrer jährlich viele hinauf gegangen gen Jerusalem auf die Feste, so man daselbst gehalten hat, die auch daselbst angebetet haben.

37. Darum hat Esau gehofft, daß er der Gnade Gottes und des geistlichen Segens auch wolle theilhaftig werden, und hat auch empfunden und erfahren, daß er leiblich ist gesegnet und gemehret worden: darnach hat er seines Bruders Demuth gesehen. Denn alles dazu gedient hat, daß sein Herz dadurch erweicht und gelindert worden ist, auf daß er wahrhaftig und von Herzen mit seinem Bruder versöhnt würde, und also bei sich selbst geschlossen hat: Warum soll ich meinen Bruder erwürgen? Warum soll ich an seinen lieben Kindern und Weibern Tyrannei üben? Behüte mich Gott, daß ich ja kein Brudermörder werde. Und geschieht Esau daran eben so lieb, daß Jakob Gut und Glück bekommen hat, als ob es ihm selber widerfahren wäre.

38. Zuletzt ist auch ein ernstes und brünstiges Gebet dazu gekommen, damit Jakob Linderung erlangt hat. Der Reichthum selbst und die leiblichen Güter, die er hat, hätten das erbitterte Herz nicht zufrieden stellen können, wiewohl sie bisweilen den Zorn etwas pflegen zu lindern; das Gebet aber hat größere Kraft gehabt. Denn wo wir ernstlich und beständig beten, so muß gewißlich aus dem Feinde ein Freund werden. Sofern wir nur zu Gott schreien, und hoffen nicht auf unsere Würdigkeit, sondern wie das glimmende Docht allein auf Gottes Gnade und Barmherzigkeit, alsdann erhört uns Gott gewißlich. Das ist nun eine rechte und ungefärbte Versöhnung dieser Brüder, wie Mose gewisse Zeichen erzählt des guten geneigten Willens und brüderlicher Liebe an ihnen beiden. Nun ist die Auslegung dieses Capitels leicht.

\subsection{II}

\begin{quote}
\textbf{V. 5. Und (Esau) hub seine Augen auf, und sahe die Weiber mit den Kindern, und sprach: Wer sind diese bei dir? Er antwortete: Es sind Kinder, die Gott deinem Knechte bescheret hat.}
\end{quote}

39. Dieses ist ein sehr schön und freundlich Gespräch dieser zwei Brüder und ist voller Göttlichkeit. Jakob bleibt noch immer in der Demuth und daß er seinen Bruder mit Freundlichkeit möge einnehmen. Denn Gott will haben, daß man auch die weltliche Herrlichkeit oder, wie sie es gemeiniglich genannt haben, das sanctum saeculare, das ist, das weltliche Heiligthum, in seiner Ordnung und gebührlicher Weise ehren soll. Es hat seine Ehre und seine Hoheit. Die nun darin die Andern übertreffen, sollen mit nichten verachtet werden, sondern man soll ihnen ihre gebührliche Ehre geben nach Gottes Ordnung.

40. Das hebräische Wort chanan und chen heißt eigentlich, Einen begnaden. Also sagt hier Jakob: Damit mich Gott begnadet hat. Also erkennt er zugleich Gottes Gnade und Wohlthat und dankt ihm dafür; seine Sache aber schmückt er mit Worten, wie er nur kann und mag, seines Bruders Freundschaft damit zu erhalten. Und lehrt uns mit seinem Exempel, daß wir unsern Obersten auch gebührliche Ehrerbietung erzeigen sollen. Wir sollen gegen sie nicht hart, rauh, oder halsstarrig und ungehorsam sein, sondern uns befleißigen, daß wir mit Demuth und Ehrerbietung ihre Gunst erlangen mögen. Das gehört in die Sittenlehre.

41. Warum fragt aber Esau, weß die Kinder sind, als ob er das nicht sollte wissen? Sind doch die Boten vorausgesandt gewesen, die ihm das Geschenk gebracht haben, von welchem er ohne Zweifel erforscht hat, was Jakob für
\quellseite{221}
Gefährten habe, und von andern Dingen mehr? Antwort: Er fragt darum nach den Kindern, daß er also Gelegenheit haben möge, mit seinem Bruder freundlich und brüderlich zu handeln und sich mit ihm zu unterreden. Denn da muß aller Argwohn beigelegt werden, als sollte noch Betrug oder Falschheit darunter verborgen sein. Er hat nicht fälschlich wie ein Heuchler gehandelt, und ich halte dafür, er werde endlich noch selig geworden sein, wie ich von Ismael auch halte.

42. So ist nun der Wille Esaus wahrhaftig überwunden und geändert worden. Wo aber der Wille geändert ist, so weicht das Andere von selber und folgt Alles sein hernach, da steht Jakob das Haus und alles Gut Esaus offen; denn der Wille ist die Königin, die über die Güter herrscht. Wenn ich weiß, daß mich Einer lieb hat, so zweifle ich daran nicht, es muß mir alles bereit und willig sein, was mir von einem freundlichen und gutwilligen Herzen herkommen kann. Dieweil nun das Herz Esaus mit Jakob versöhnt ist, indem auch Tag und Zeit, wie man zu sagen pflegt, die Bekümmerniß und den Zorn lindert, hernach die Demüthigung des Bruders: so darf man nichts in dieser Versöhnung für falsch oder heuchlerisch ansehen, sondern daß alle Gebärden und Worte aufrichtig und getreu gewesen, die aus rechter Liebe und gutem geneigten Willen hergekommen und nicht gefärbt gewesen sind. Denn darum ist er seinem Bruder entgegen gelaufen, hat ihn umarmt, geküßt und geweint, und sucht jetzt, auf welcherlei Weise er nur kann, Gelegenheit, sein brüderliches Herz und Liebe gegen Jakob zu beweisen. Das ist der sehr herrliche und rechte Israelitische Sieg, da durch das Gebet und den Kampf des Glaubens Gottes und der Menschen Wille überwunden worden ist.

\begin{quote}
\textbf{V. 6—8. Und die Mägde traten herzu mit ihren Kindern, und neigten sich vor ihm. Lea trat auch herzu mit ihren Kindern, und neigten sich vor ihm. Danach trat Joseph und Rahel herzu, und neigten sich auch vor ihm. Und er sprach: Was willst du mit alle dem Heer, dem ich begegnet bin?}
\end{quote}

43. Dieses Gepränge richtet Jakob auch an, seinen Bruder damit zu versöhnen. Wiewohl er zuvor durch göttliche Kraft versöhnt war, natürlich, durch den Sieg und Kampf des Gebetes und des Glaubens, so hat er doch auch das äußerliche Anzeichen der Freundschaft nicht wollen unterlassen, auf daß es nicht scheine, als wollte er Gott versuchen, oder den Frieden und die angefangene Einigkeit verwirren und mit Hoffahrt seines Bruders Herz aufs neue betrüben. Darum neigt sich nun Jakob mit seinen Weibern, Mägden und Kindern vor Esau, so hoch er der Sohn und Erbe ist des bessern Segens. Denn er hatte den geistlichen Segen um des verheißenen Christi willen, welcher an dem leiblichen hing; Esau aber hatte die Verheißung nicht von Christo, daß derselbe von seinem Fleisch sollte geboren werden. Und doch gleichwohl demüthigt sich der, der größer ist, und hält sich also, als ob er der Kleinere wäre.

44. Auf solche Weise sollen wir auch aller göttlichen Ordnung unterthänig sein um Gottes willen; denn es ist Gottes Creatur. Darum sollen wir die Obrigkeit ehren, mit welcherlei Ehrerbietung wir können; aber doch also, daß wir um ihretwillen das Wort Christi und die Verheißung der Gnade nicht verwerfen oder verleugnen, auf daß wir den geistlichen Segen nicht verlieren. Denn man soll der Obrigkeit alle Ehrerbietung, dazu auch allerlei Dienst erzeigen mit gutem Gewissen und fröhlichem Herzen, wenn sie bleiben bei dem Ziel, so ihnen gesetzt ist, daß ist, daß der göttliche Gehorsam und das Bekenntniß des Wortes unverletzt bleibe. Wo sie aber Gott wollen gleich sein, soll man sie strafen, und soll ihnen dagegen vorhalten den Spruch der Apostel in deren Geschichten am 5. Cap. V. 29.: „Man muß Gott mehr gehorchen, denn den Menschen.“

45. Also demüthigt sich nun der Größere, der von Gott mit dem geistlichen Segen begnadet ist, gegen seinen Bruder, der vor Gott der Kleinere ist. Wie wir den Königen, Fürsten, Edelleuten, Bürgermeistern, Rathsherren ihre gebührliche Ehre erzeigen, heißen sie, wie Christus Luc. 22, 25. spricht, „Gnädige“, Durchlauchtige, Edle, Weise, „Herren“ etc. Denn es ist eine weltliche Ordnung und Herrlichkeit, von Gott geordnet und eingesetzt, die Gott von uns will geehrt haben. Wenn sie aber ihren Stuhl über die Sterne Gottes erhöhen und dem Allerhöchsten gleich sein wollen, wie Jes. 14, 13. 14. steht, soll man ihnen mit nichten gehorsam sein. Aber in ihrer Ordnung, Stätte und Maß erzeigen wir ihnen gebührliche Ehrbietung und haben sie herzlich lieb.

46. Wir schmeicheln ihnen nicht, wenn wir sie Gnädigste Herren und Durchlauchtige heißen, sondern wir ehren den Stand und die Person von Herzen, die von Gott zu diesem Amte verordnet ist. Auf solche Weise heuchelt Jakob dem Esau als dem Obersten auch nicht, sondern nennt ihn einen Herrn, sich aber einen Knecht aus rechter Demuth und Ehrerbietung gegen seinen Bruder. Und ist ein schön Exempel einer sonderlichen Höflichkeit, so in dieser Sprache und Volke sehr gemein ist; welche Höflichkeit auch die gebraucht haben, die einander gleich gewesen sind, wenn sie mit einander geredet haben. Es ist bei uns nicht gebräuchlich, daß sich Einer des Andern Knecht nennen sollte. Dieses ist die einfältige und gemeine Form, die wir gebrauchen, daß wir sagen: Mein Freund, mein Herr etc.

47. Esau aber fragt abermals, was Jakob mit alle dem Heere wolle? auf daß er also Ursache habe, freundlich mit seinem Bruder zu reden, und Worte mit ihm zu wechseln aus Liebe und Freundschaft, so Einer zum Andern gehabt hat. Und es läßt sich ansehen, als habe er dermalen das Geschenk, so ihm Jakob gesandt, noch nicht empfangen. Was aber die Ursache gewesen sei, kann man nicht wissen; ob er noch erhitzt gewesen ist vom Hasse, und die Boten, so zu ihm gesandt waren, nicht habe hören wollen, ehe er gesehen hat, daß sich der Bruder selbst demüthigt; oder aber, da er nun etwas gelinder geworden, und die Heere, so ihm sein Bruder unter die Augen schickt, gesehen hat, aus sonderlicher Güte und Barmherzigkeit sei bewogen worden und derhalben das Geschenk nicht habe nehmen wollen: das weiß ich nicht gewiß. Jedoch läßt es sich ansehen, daß er mehr aus Freundschaft denn aus Haß sich geweigert hat, das Geschenk anzunehmen. Was sollte ich von ihm solches annehmen? wird er gesagt haben. Warum schickt er mir solch Geschenk, so ich doch von Gottes Gnaden viel reicher bin, denn er ist? Und darum wird er die Boten nicht gehört haben, die um der Ursache willen zu ihm gesandt waren, nämlich, daß sie dermalen noch nicht verstanden haben, wie er gegen seinen Bruder gesinnet gewesen.

48. Also will ich es auslegen, auf daß dieses gemäß bestehen bleibe, daß dieser Sieg solch Wunderwerk gewirkt habe, daß das Herz Esaus hat müssen geändert und versöhnt werden. Darum kann es wohl geschehen sein, daß er, da der Zorn noch hitzig gewesen, mit den vierhundert Mann fortgezogen sei, seinen Bruder zu schrecken oder zu ermürgen. Da er aber gesehen hat, daß er sich demüthigt und daß er ihm Geschenke sendet, da ist die Hand Gottes dazu gekommen und hat ihm sein Herz geändert. Warum sollte ich meinen Bruder Gewalt thun, oder was sollte ich dem frommen und getreuen Manne nehmen, der elend und arm ist? Der HErr hat mich reichlicher gesegnet als meinen Bruder; es wäre billig, daß ich mit meinem Reichthum ihm in seiner Armuth sollte zu Hülfe kommen: warum sollte ich denn Geschenke von ihm nehmen?

\begin{quote}
\textbf{Er antwortete: Daß ich Gnade fände vor meinem Herrn.}
\end{quote}

49. Jakob beharrt noch immer dabei, daß er seines Bruders Freundschaft haben und behalten möge. Und ihre Herzen werden nun mit einander gar versöhnt, daß aller Haß und Widerwille ausgelöscht wird. Das wirkt Gott innerlich durch seinen Geist, auswendig aber durch Mittel, nämlich, durch das Gepränge der Ehrerbietung und daß er ihm Geschenke hat lassen anbieten; wiewohl es sich ansehen läßt, als habe Esau das Geschenk nicht angenommen vor großer herzlicher Liebe, damit er seinen Bruder, da er nun mit ihm versöhnt war, umfaßt.

\begin{quote}
\textbf{V. 9. Esau sprach: Ich habe genug, mein Bruder, behalt, was du hast.}
\end{quote}

50. In diesen Worten ist nichts, das gefärbt oder falsch sei, sondern es ist alles als aufrichtig und von Herzen geredet zu verstehen. Denn die Schrift dichtet nichts; sondern wenn sie sagt, daß Jakob beide Gott und Menschen obgelegen habe, will sie haben, daß wir es dafür halten sollen, daß Esau gänzlich versöhnt und geändert sei. Denn Jakob hat beide, nämlich, Gott und Menschen, als einen Widersacher und Feind gefühlt und erfahren; aber der Kämpfer sagt, er hätte Gottes und des Menschen Herz überwunden und geändert. Darum verstehen wir es also, daß alle Worte und Werke Esaus ohne Falsch und Heuchelei sind.

51. Dieses treibe ich darum etwas fleißiger, euch der Juden Größe damit zu erinnern, die da pflegen die Werke der Heiden und Unbeschnittenen zu verkleinern und zu verachten, und ihre Gebrechen groß zu machen, sich selber aber wissen sie hoch zu rühmen und über die Maßen sehr zu erheben. Darum halten sie es dafür, daß sich Esau nur äußerlich also freundlich gestellt habe und nicht von Herzen; und daß dem Jakob alle solche Worte, damit er ihn so liebkost, verdächtig gewesen sind und daß er seinem Bruder nicht habe Glauben geben wollen.

\quellseite{222}

52. Aber ich halte dafür, daß Esau wahrhaftig und von Herzen geändert worden sei, wiewohl er die gerechteste Ursache gehabt, daß er Jakob feind war und mit ihm zürnte. Denn ihm als dem Erstgebornen gebührte der Segen. Aber er war ein großer, feiner, tapferer Mann, der ohne Zweifel unterrichtet gewesen ist aus der Lehre und Predigt Isaaks und der andern Väter, bei welchen er erzogen worden ist, und gelernt hat, seine Lust und böse Begierde im Zaum zu halten und zu zwingen. Darnach ist auch dieses Gepränge dazu gekommen, daß darauf angerichtet war, daß er dadurch sollte versöhnet werden; item, der Kampf und das Gebet Isaaks, und endlich die göttliche Regierung: welches alles das Herz Esaus dazu getrieben hat, daß er solcher Ungerechtigkeit vergessen hat. Darum hat er dieses alles aus gutem reinem Herzen gethan und geredet; wie Jakob bald hernach hoch rühmen wird, da er sagt: „Ich sahe dein Angesicht, als sähe ich Gottes Angesicht.“ Denn dieweil es ein göttlich Werk ist, so kann es ja nicht falsch oder gefärbt sein, sintemal Gott nichts thut, denn was recht und wahrhaftig ist.

53. Wo man nun den Wahn bei Seite legt, als sollte Esau hinterlistig gehandelt und sich nicht von Herzen so freundlich erzeigt haben, so rühmen wir Esau billig, daß er sich ehrlich und freundlich geweigert habe, das Geschenk anzunehmen. Ich habe, sagt er, alles Dinges genug, darum will ich dein Geschenk nicht annehmen; „behalte, was du hast“; ich will es deinen Kindern und Gesinde lassen, daß sie dessen genießen mögen. Dieses ist alles aus einem freundlichen und brüderlichen Herzen hergenommen.

\begin{quote}
\textbf{V. 10. 11. Jakob antwortete: Ach nicht, hab ich Gnade funden vor dir, so nimm mein Geschenk von meiner Hand; denn ich sehe dein Angesicht, als sähe ich Gottes Angesicht; und laß dir’s wohlgefallen von mir. Nimm doch den Segen von mir an, den ich dir zugebracht habe; denn Gott hat mir’s beschert, und ich habe alles genug. Also nötigte er ihn, daß er’s nahm.}
\end{quote}

54. Jakob hält hart an, daß sein Bruder das Geschenk, so er ihm zugebracht hat, annehmen wolle. Denn ich, spricht er, gebe dir es von dem Segen und Gute, so mir Gott bescheret hat; und du sollst dich meine Armut und Nothdurft nicht bewegen lassen, daß du derhalben dich weigern wollest, diese Verehrung von mir anzunehmen. Denn ich habe von Gottes Gnade Alles genug, und biete dir solches nicht an, als sollte mich die Noth oder Furcht dazu bewegen, sondern ich thue es aus gutem geneigtem Willen und rechter herzlicher Liebe, die ich zu dir trage. Darum bitte ich dich, du wollest es nicht verschmähen, sondern dir gefallen lassen. Denn ich schenke dir es aus keiner andern Meinung, denn daß du daraus verstehen sollst, daß ich dich brüderlich lieb habe, und wo du etwa in deinem Herzen erzürnet wärest, daß du mir verzeihen wollest. Das Geschenk ist zwar an sich gering und sollte deiner Herrlichkeit nach wohl etwas größer sein: aber du mußt mein Herz ansehen und wie ich gegen dich gesinnet sei. Das hebräische Wort minchah ist droben, Cap. 4, V. 3. 4., gebraucht worden, da vom Opfer Kains und Abels geredet wird. Es heißt eigentlich ein Geschenk.

55. Es setzt aber Jakob eine Ursache dazu, damit er versucht, seinen Bruder zu überreden, daß er das Geschenk, so ihm zugebracht wird, annehmen wolle, und will sagen: „Ich habe an dem allein genug, und ist mir eben so lieb, als sonst alles Gut sein kann, daß ich Gnade vor dir gefunden habe: Denn ich sahe dein Angesicht, als sähe ich Gottes Angesicht.“ Dieses ist eine wunderbare und sehr liebliche Rhetorik. Und man soll nicht denken, daß Jakob lüge oder daß er etwas erdichtet hätte, sondern man soll es dafür halten, daß er ernstlich und von Herzen rede. Nicht, daß ich dafür halte, daß man ihm in seinen Sünden heucheln oder liebkosen sollte; sondern daß man die Tugend und Gebrechen recht lerne unterscheiden.

56. Denn es ist nun oftmals gesagt worden, daß wir die Heiligen nicht allenthalben entschuldigen. Und wir haben droben gehört, daß Jakob seine Schwachheit genug bewiesen hat, da der Unglaube und das Fleisch dem Geiste widerstrebt, welche Schwachheit wir nicht entschuldigen. Denn ein heiliger oder gläubiger Mensch, der da weiß, daß ihm seine Sünden vergeben sind und daß er ein Erbe ist des ewigen Lebens, soll nicht zittern, erschrecken oder verzagen, sondern mit festem Glauben auf die Verheißung kämpfen und den Zweifel überwinden.

57. Dieweil wir aber noch im Fleische sind, das unrein und von Sünden befleckt ist, darum werden wir verhindert, daß wir nicht mit solcher Inbrunst vertrauen, hoffen und beten können, wie wir thun sollten, sondern es ist noch viel Schwachheit und Gebrechen an uns. Und Gott duldet doch mit Langmuth das glimmende Docht und den schwachen Glauben, sofern er nur nicht wider das Wort streitet, oder das Bekenntniß und den Glauben des Wortes nicht verwirft. Die Schwachheit thut uns keinen Schaden, sondern die Kraft Gottes ist in den Schwachen mächtig, 2 Cor. 12, 9. Aber man muß sich vor Heuchelei hüten, nämlich, daß wir weder die Gebrechen und Sünden der Heiligen entschuldigen, noch auch übel auslegen, daß sie manchmal und recht geredet haben.

58. Es fragt Augustin an diesem Ort: Ob Jakob ein Götzendiener gewesen sei und ob er seinem Bruder geschmeichelt habe, da er sagt, er habe sein Angesicht gesehen, als sähe er Gottes Angesicht? und läßt es sich sehr sauer werden, daß er Jakob entschuldigt. Es ist aber eine sonderliche Art zu reden in der Schrift, welche man recht verstehen soll. Denn droben, Cap. 32, V. 28., hat der Kämpfer gesagt: Du hast Gott und Menschen überwunden, das ist, es haben dir Zwei den Rücken zugewendet, nämlich, Gott und dein Bruder, welches ein Anzeichen ist des Zorns oder eines abgeneigten Herzens. Und diese Art und Weise zu reden ist in allen Sprachen gemein. Wenn ich einem feind bin, so wende ich das Angesicht und die Augen von ihm und kehre ihm den Rücken zu, das ist ein Zeichen des Zorns und Ungnade. Auf diese Weise, wenn von Gott gesagt wird, daß er das Angesicht von einem wendet und ihm den Rücken zukehrt, so wird damit angezeigt, daß man die Gnade oder die Erkenntniß der Gnade, der Gunst und der Barmherzigkeit Gottes, dazu der Freude und Danksagung verloren habe, und daß anstatt dieses allen komme Schrecken, Traurigkeit und daß man an Gottes Willen zweifelt; daß ein betrübt Gewissen sagt: Ich weiß nicht, wie ich mit Gott daran bin.

59. Darum heißt nun das Angesicht nicht das Anschauen des göttlichen Wesens, sondern ist eben das, so im 67. Psalm V. 2. steht: „Er lasse uns sein Antlitz leuchten, und sei uns gnädig“, das ist, er rede freundlich mit uns; item im 4. Psalm V. 7. 8.: „Herr, erhebe über uns das Licht deines Antlitzes. Du erfreuest mein Herz.“ Als, wenn ich zu dem, der angefochten wird, sage: „Sei getrost, mein Sohn, dir sind deine Sünden vergeben“ Matth. 9, 2.; item: „O Weib, dein Glaube ist groß“ etc., Matth. 15, 28. Alsbald leuchtet uns das Antlitz Gottes, wenn er uns so freundlich anredet, und zeigt mit Worten und Werken an, daß er uns von Herzen günstig sei. Da, sage ich, zeigt er uns sein Antlitz, und dasselbe geschieht in diesem Leben nur im Wort und in den Sacramenten.

60. Dagegen aber, wenn er sagt: Du hast gehuret, hast gestohlen, du sollst des Todes sterben; „die Hurer, Ehebrecher“ etc. „werden das Reich Gottes nicht ererben“, 1 Cor. 6, 9. 10.; item 1 Thess. 4, 6.: „Niemand vervortheile seinen Bruder im Handel; denn der Herr ist der Rächer über das alles“; item Matth. 25, 46.: „Diese werden in die ewige Pein gehen“: vor solchen Sprüchen und Worten allein erschrickt und zittert das Herz. Darum ist da das Antlitz Gottes von uns gewendet.

61. Wie nun Jakob droben, Cap. 32, V. 30., vom Angesicht Gottes sagt, wie es sich zu ihm gewendet habe: „Meine Seele ist genesen“, darum daß ihn der Herr freundlich angeredet und ihm einen sehr reichen Trost in der Verheißung und im Segen gegeben hat: also geht ihm auch bei seinem Bruder die Sonne, das Licht und Leben auf aus dem Wort und Zeugniß seines guten geneigten Willens. Und ist dieses die Meinung und Verstand seiner Worte, daß er sagen will: Mein lieber Bruder, ich habe gesehen, daß sich dein Angesicht zu mir gewendet hat, nicht anders, denn wie ich zuvor Gottes Angesicht gesehen habe, da es sich auch zu mir gewendet hat. Und da ich jetzt nun dein Angesicht ansehe, däucht mich, daß ich wahrlich GOttes Angesicht sehe. Dieses ist keine Schmeichelei, sondern es ist ernstlich und von Herzen geredet; denn er will also sagen: Ich habe mich deines Angesichtes dermaßen gefreut, als hätte ich gesehen, daß mir GOtt begegnet wäre.

\quellseite{223}

62. Man muß aber das freundliche Angesicht GOttes erkennen in seinen Verheißungen, in den Sacramenten, desgleichen auch in äußerlichen Wohltaten und Gaben, als, an einem gnädigen Fürsten, und wenn der Nächste, item, Vater und Mutter, mit uns versöhnet sind. Wenn ich sehe, daß meiner Eltern Angesicht freundlich ist, so sehe ich damit zugleich auch, daß mir GOttes Angesicht freundlich ist und er mich anlacht; wie Jakob zuvor gesagt hat, er habe den HErrn von Angesicht zu Angesicht gesehen. Dasselbe Angesicht GOttes sieht er jetzt am Angesicht seines Bruders Esau auch; denn er sieht das Wohlgefallen des göttlichen Willens in der Freundlichkeit und Gunst seines Bruders. Auf dieselbe Weise scheint in allen Creaturen das Angesicht GOttes; denn es sind GOttes Wercke und Zeugnisse des göttlichen Willens und Gegenwart. Damit erzeigt er sich so freundlich gegen uns als mit dem äußerlichen Anschauen, wie er innerlich im Herzen durch das Wort und Verheißung uns sein freundliches und gnädiges Angesicht zeigt.

63. Das ist die rechte Erklärung dieser Art und Weise zu reden, damit es niemand also deute, daß es eine Schmeichelei oder Abgötterei gewesen sei. Und Jakob setzt noch eine Erklärung hinzu, damit er die Liebe und Freundschaft an seinem Bruder rühmt, da er sagt: „Und du hast mir wohl gewollt.“ Das hebräische Wort razon heißt guter Wille oder Wohlgefallen; wie im 51. Psalm V. 20. steht: „Thue wohl an Zion nach deiner Gnade“, oder nach deinem guten Willen; item Matth. 17, 5.: „Dies ist mein lieber Sohn, an dem ich Wohlgefallen habe“; item Luc. 2, 14.: „Den Menschen ein Wohlgefallen.“ Und Jakob will an diesem Ort so viel sagen: Ich sehe, daß du Gefallen an mir hast und nimmst mich mit Gnaden an; du umarmst mich, küßest mich und weinst mit mir; zudem weigerst du dich, meine Geschenke anzunehmen, bist geneigter zu geben denn zu nehmen. Also erzeigt du gegen mich deinen guten Willen und Wohlgefallen. So nimm nun dieses kleine Geschenk zum Zeichen einer Gegenliebe, daß du anzeige, daß ich deinen guten Willen erkenne. Er nennt aber das Geschenk einen Segen, da er sagt: „Nimm den Segen von mir an, den ich dir zugebracht habe; denn GOtt hat mirs bescheret“, denn GOtt hat mir Alles genug gegeben.

64. Dieses sind wahrlich Worte eines Herzens, das über die Maßen fröhlich ist und große Freude hat. Er freut sich von ganzem Herzen, daß er der brüderlichen Liebe genießen mag. Darum ist dieses ein sehr schönes klares Wetter, so auf das Ungewitter und Finsterniß des Kampfes gefolgt ist, den er zuvor mit seinem Bruder und mit dem Manne gehalten hat. Ich sehe, spricht er, daß GOtt und mein Bruder aus großer Güte mit mir versöhnet sind. Darum wird große Freude damit angezeigt, da er droben, Cap. 32, V. 30., gesagt hat: „Und meine Seele ist genesen.“ Also wollte er sagen: Nun dünkt mich, daß ich wahrhaftig wiederum lebendig werde aus dem Tode und der Hölle, dieweil ich das freundliche und liebliche Angesicht GOttes gesehen habe. Und es leuchtet in diesen Worten eine große wunderbare Demuth. Denn er will also sagen: Ich habe zwar kein herrlich und groß Geschenk gebracht, deß man sich hoch verwundern möchte; sondern es ist GOttes Segen: und ich bitte, du wollest es also annehmen, daß es ein Zeichen und Gedächtniß sei meiner Dankbarkeit und des Segens GOttes, von deß Gabe und Gnade ich dir dieses alles, was es nur ist, anbiete. Darum nimm in GOttes Namen, was mir GOtt bescheret hat, nicht anders, als wenn es dir von GOtt gegeben wäre, daß es bei dir auch ein Segen sei und unermesslich gemehret werden möge, und daß diese Freundschaft und brüderliche Liebe, so Einer zum Andern hat, zwischen uns für und für währen möge.

65. Zuvor aber hat Esau gesagt: Ich habe genug, auf Hebräisch rab; Jakob sagt aber: Ich habe Alles, auf Hebräisch kol: der macht seinen Segen groß. Ich habe nicht allein viel, sondern ich habe Alles. Das sind Worte Eines, der sich freut und vor Freuden springt und beweisen, daß ihm GOttes und seines Bruders Angesicht gezeigt war und zwar in Gnade, Barmherzigkeit und gutem geneigtem Willen. Wenn ich nur GOttes Gnade und deine Gunst habe, will er sagen, so habe ich genug und übrig genug. Und es sollte mir gar kein Schade sein, wenn ich dir auch schon Alles gäbe: mir geht daran nichts ab, mir wird auch nimmermehr nichts mangeln, dieweil du mit mir versöhnt bist. Ich bin ein reicher Herr, weil ich GOtt und dich zum Freunde habe.

66. Dieses aber alles zeigt an, wie der Glaube in Jakob wieder lebendig und erweckt worden sei, und jetzt wiederum herrscht und triumphirt. Denn Jakob hält es jetzt dafür, daß sein Segen mit allen Gütern der ganzen Welt nicht zu vergleichen sei. Zuvor hatte er nichts, da er sagt: „Ich fürchte mich vor meinem Bruder, daß er nicht komme, und schlage mich, die Mütter sammt den Kindern“ und dem ganzen Gesinde. Da, sage ich, hat er gar nichts. Denn es war alles eitel verloren und verzagt Ding. Hier aber wird ihm wiederum Himmel und Erde erfüllt mit allen Gütern und Reichthum der ganzen Welt. Denn der Unglaube beraubt uns der Gnade und Barmherzigkeit GOttes, des Bruders und aller Güter.

67. Mose setzt aber noch dazu, wie Jakob seinen Bruder genöthigt habe, das Geschenk zu nehmen. Denn er wird also gesagt haben: Verachte den Segen GOttes und mein Geschenk nicht, auf daß du nicht angesehen werdest, als ob du nicht von Herzen mit mir versöhnet wärest. Denn wiewohl ich daran nicht zweifele, so wird doch unserer beider Familie denken, es sei noch ungewis Ding mit deiner Gnade und Gunst. Darum bitte ich dich um der Liebe und brüderlichen Freundschaft willen, daß du dieses kleine Geschenk nicht verschmähen wollest.

\section[Dritter Theil]{Dritter Theil. Von dem Anerbieten, so Esau dem Jakob macht; wie Jakob solches ausschlägt und gen Sucoth zieht}

\subsection{I}

\begin{quote}
\textbf{V. 12. 13. Und er sprach: Laß uns fortziehen und reisen, ich will mit dir ziehen. Er aber sprach zu ihm: Mein Herr, du erkennest, daß ich zarte Kinder bei mir habe, dazu Vieh und säugende Kühe; wenn sie einen Tag übertrieben würden, würde mir die ganze Heerde sterben.}
\end{quote}

68. Esau will dankbar sein und seinem Bruder auch wiederum Wohlthat erzeigen. Denn er erbietet sich, daß er mit ihm ziehen und ihn geleiten wolle bis an den Jordan. Und dieses geschiehet nicht hinterlistig; wiewohl es glaublich ist, daß bei der Familie beider Theile noch allerlei Argwohn mag gewesen sein.

69. Aber Jakob ist aus der Verheißung dessen gewiß, daß sein Bruder mit ihm versöhnt war, und versteht wohl, daß dieses alles aus einem freundlichen und getreuen Herzen herkommt, und bittet doch seinen Bruder, daß er sich dermaßen nicht bemüthigen und ihm diesen Dienst nicht leisten wolle. Als wollte er sagen: Es ist nicht nöthig, daß du dich also bemüthigen sollst, und es würde dir allzuschwer fallen; denn mein Herr weiß, daß ich zarte Kinder bei mir habe, dazu \unsicher{trächtiges} Vieh und Schafe, saugende Lämmer und Kälber, so ohne Gefahr nicht folgen können, wenn wir so eilends fortziehen. Ruben war dreißig Jahre und nicht viel darüber; Levi war zwölf Jahre; Juda, zehen Jahre; und so immer fort bis auf Dina, die auf dieser Reise noch nicht das vierte Jahr erreicht hatte. Derhalben führt er mit sich einen Haufen Kindlein mit großer Mühe; die Kinder aber, die so zart sind, können dem Haufen, so Esau gehabt von reisigem und gerüstetem Volk, nicht folgen. Also, lieber Bruder, sagt Jakob, werden wir dich hindern und dir beschwerlich sein.

70. Daß sich nun Jakob also weigert, den Dienst, so ihm sein Bruder angeboten, anzunehmen, ist ehrlich und wohl gethan. Denn es ist viel beschwerlicher, einen Haufen wehrloser und zarter Kinder zu führen denn etliche Heere reisigen Volks; sintemal sich vielerlei zutragen kann, dadurch man verhindert wird, daß man nicht so fortziehen kann und daß man bisweilen etliche Tage oder Stunden muß still liegen. Es ist unendliche Sorge dabei, wo man mit seiner Familie reisen soll; und vornehmlich ist es dazumal beschwerlich gewesen, da man keine Wagen gehabt hat, wie jetzt bei uns, sondern sie haben ihr Gespäck auf die Kameele und Esel gelegt und die Kinder auch darauf gesetzt.

71. Darum ist es nicht erdichtet oder erlogen, daß Jakob sagt: Wir werden dir beschwerlich sein und dich verhindern, werden dich auf der Reise aufhalten, wir wollen dich mit diesem verdrießlichen Haufen nicht beschweren; sondern er bittet ihn ernstlich, daß er diese Mühe nicht wolle auf sich laden, auf daß Esau keine Entschuldigung bringen soll, daß er etwas langsam ziehen wolle oder etwas Anderes dergleichen. Und will damit heimlich zu verstehen geben, daß solches dem gerüsteten Volke, so Esau bei sich gehabt, werde verdrießlich sein. Ja, er erzählt auch die Gefahr, nämlich, da er sagt: „Wenn sie einen Tag übertrieben würden, so würde mir die ganze Heerde sterben.“ Das hebräische Wort heißt schlagen, wie wir pflegen an eine Thüre zu klopfen. Denn das Vieh wird getrieben mit einer Ruthe oder Stecken. Wenn ich es nun also mit Gewalt treiben würde, sagt er, daß die Heerden so schnell wie das reisige Volk gehen müßten, so würde mir die ganze Heerde und Kinder in einem Tage sterben. Auf Deutsch heißt es, übertreiben. Darum schickt es sich besser, daß du mit deinem Volke vorherziehest, so will ich allemal nachfolgen.

\quellseite{224}

\begin{quote}
\textbf{V. 14—16. Mein Herr ziehe vor seinem Knechte hin. Ich will sänftiglich hin nachtreiben, darnach das Vieh und die Kinder gehen können, bis daß ich komme zu meinem Herrn in Seir. Esau sprach: So will ich doch bei dir lassen etliche vom Volk, das mit mir ist. Er antwortete: Was ists vonnöthen? Laß mich nur Gnade vor meinem Herrn finden. Also zog des Tages Esau wieder um seines Weges gen Seir.}
\end{quote}

72. Das hebräische Wort nahal heißt eigentlich geleiten, wie im 31. Psalm V. 4. steht: „Um deines Namens willen wollest du mich leiten und führen“; und im 23. Psalm V. 2.: „Der Herr führet mich zum frischen Wasser.“ Und an diesem Ort sagt Jakob: Ich will sie sänftiglich hin nachtreiben lassen, wie die Schafe auf die Weide oder zum Wasser geführt werden, daß das Vieh und Kinder nicht verderben. Denn das ist eines Hirten Amt. Und also sagt David von Absalom 2 Sam. 18, 5.: „Fahret mir säuberlich mit Absalom.“

73. Das er aber weiter sagt: „Darnach das Vieh und die Kinder gehen können“, das lautet nach dem Hebräischen auf Lateinisch also: juxta pedem operis et parvulorum: das ist: nach dem Fuß des Werkes und der Kleinen. Und er nennt all sein Hab und Gut auf Lateinisch opus; auf Hebräisch, melachah, auf Deutsch, Werk. Und eben dasselbe Wort steht auch im dritten Gebot. Daher kommt das hebräische Wort melecheth, die Kraft des Himmels, die da auf Erden wirket, das Land fruchtbar zu machen, daß es reichlich Frucht trage. Und Christus redet auch auf die Weise, Luc. 21, 26.: „Der Himmel Kräfte werden sich bewegen.“ Hier aber nennt Jakob sein ganzes Gesinde, und alle sein Hab und Gut ein Werk auf Lateinisch opus, ausgenommen die Kinder und Weiber. Wie im dritten Gebot gesagt wird, 2 Mos. 20, 9. 10.: „Sechs Tage sollst du arbeiten, und alle deine Werke thun“, das ist, du sollst darin ausrichten, was du zu thun hast; „aber am siebenten Tage ist der Sabbath“, da sollst du ruhen und das Wort des Herrn, deines Gottes, hören.

74. Es wird hier ferner von Andern eine Frage erregt: Ob Jakob gelogen habe, da er sagt, er werde zu Esau nach Seir kommen, und man findet es doch nirgend geschrieben, daß er dahin gekommen sei. Ich habe aber zum öftern gesagt, daß ich die Heiligen nicht aller Sünden entschuldigen könne. Denn wir wissen, daß die allergößten und heiligsten Männer oftmals greulich gefallen und befleckt worden sind, nicht allein durch gemeinen Irrthum und Schwachheit, sondern auch mit großen Sünden wider den Glauben, Hoffnung, Liebe und Geduld, nämlich, mit Mißglauben, Zweifel, Zorn und Ungeduld wider Gott. Da wollen wir sie zu Gesellen haben und Exempeln, deren wir uns trösten mögen. Und wir sollen nicht denken, daß sie Stöcke, Steine oder Klötze gewesen seien, sondern uns gleich. Elia erlangte von Gott, daß in drei Jahren kein Regen gefallen ist; doch hat er hernach gefühlt, daß sein Fleisch zittert, und begehrt, daß er nur bald sterben möchte, 1 Kön. 19, 4. Aber wo sie recht gethan und wohl geredet haben, und was nicht zu tadeln ist, das soll man mit nichten hinderlistig oder fälschlich deuten. Darum sollen wir Jakob auch nicht im Verdacht haben, als ob er fälschlich oder betrüglich mit seinem Bruder gehandelt hätte, sondern es ist an ihm ein sehr getreu Herz, und das sehr fröhlich ist um der Versöhnung willen, so zwischen ihm und seinem Bruder aufgerichtet war. Er ist nun fröhlich und befreit sich seiner Gefahr mehr, darum hat er keine Ursache, seinem Bruder zu schmeicheln oder zu lügen.

75. Warum ist er aber nicht gen Seir gezogen, so er es doch seinem Bruder verspricht? Antwort: Er verspricht ihm das nicht, daß er auf dieser Reise straks gen Seir ziehen wollte, sondern er heißt Esau vor sich hinziehen, bis daß er ihm folgen und auch gen Seir kommen werde. Er ist aber erst heimgezogen zu seinem alten Vater und hat denselben besucht, ehe er zu seinem Bruder gen Seir gekommen ist. Es reimt sich aber nicht übel, daß er etlichemal werde zu seinem Bruder gen Seir oder sonst irgend an einen andern Ort gezogen sein. Und ist daran kein Zweifel, sie sind beide zusammengekommen, da ihr Vater ist begraben worden.

76. Esau aber, auf daß er abermals seine Freundschaft und guten Willen beweisen möge, sagt, daß er doch etliche von seinem Volke bei Jakob wolle bleiben lassen, die ihn beschützen mögen. Aber Jakob weigert sich, solches anzunehmen, darum daß ihnen der lange Verzug auf der Reise um so vieler Hindernisse willen würde verdrießlich sein. Ich habe daran, sagt er, gar genug, daß ich nun weiß, daß du mein Freund bist. Also ziehen sie nun beide hin, da sie von einander Abschied genommen haben, Esau gen Seir, Jakob aber gen Suchoth etc.

\subsection{II}

\begin{quote}
\textbf{V. 17. Und Jakob zog gen Suchoth, und bauete ihm ein Haus, und machte seinem Vieh Hütten; daher heißt die Stätte Suchoth.}
\end{quote}

77. Die Stadt Suchoth ist damals also genannt worden von der That. Denn Jakob hat daselbst eine Hütte gemacht für sein Vieh und ein Haus gebaut für seine Familie. Und der Text lautet, als ob Jakob an dem Orte eine Zeitlang geblieben wäre, und daß es nicht allein ein Abwegelein oder Herberge auf eine Nacht oder Tag gewesen sei, sondern er gedenkt des Bauens, so Jakob gethan hat. Darum gibt er zu verstehen, daß er daselbst eine Zeitlang gewohnt habe. Lyra sagt, nach der Meinung der Hebräer, er habe daselbst anderthalb Jahre gewohnt; aus welcher Meinung die vorige Frage gekommen ist: Ob Jakob gelogen habe, sintemal er versprochen hat gen Seir zu kommen und doch gleichwohl anderthalb Jahre zu Suchoth geblieben ist?

78. Es ist aber ungewiß, wie lange er daselbst geblieben, und ob er zu seinem Bruder gekommen sei oder nicht. Und hat doch wohl geschehen können, ja, es ist also anständig und bei den Menschen gebräuchlich, daß sie ihre Liebe bezeugen mit Dienst und Freundschaft, so sie einander erzeigen, die mit einander versöhnt sind. Auf die Weise ist es wohl glaublich, daß diese Brüder auch oftmals sind zusammengekommen und Freundschaft mit einander gehalten, und dieselbige auch unter ihren Kindeskindern und Schwägern aufgerichtet haben. Das, dünkt mich, folgt besser daraus, daß der Text sagt: Jakob habe zu Suchoth gewohnt, daselbst eine Hütte gemacht und ein Haus gebaut, denn daß er sollte gelogen haben und zu seinem Bruder nicht gekommen sein.

79. Daß ich aber dies noch nebenbei sage: Ich bin etwan der Meinung gewesen, daß Benjamin in diesem Jahre, da Jakob wieder aus Mesopotamien gekommen, geboren wäre: aber ich habe gefunden, daß die Meinung falsch ist und daß er erst bei acht Jahren darnach geboren ist. Denn das ist ein nützlicher und nöthiger Fleiß, daß man in den Historien der heiligen Schrift forsche und an die Ordnung der Zeit denke, wie Eines nach dem Andern ergangen und geschehen sei. St. Paulus verbietet die Geschlechtsregister der Personen und nicht die Rechnung der Zeit. Denn solche Rechnung ist sehr nützlich, die heilige Schrift recht zu verstehen. Die Geschlechtsregister aber sind wegen der unendlichen Unordnung ungewiß und schier unmöglich bei solcher großen Veränderung, da sie zum andern- und zum drittenmal geheirathet haben, desgleichen auch bei dem Unterschied der Geschlechter, die auf mancherlei Weise untereinander gemenget sind. Es kommt auch oftmals vor, wie Lyra anzeigt, daß ein Mensch zwei oder drei Namen hat, damit er genannt wird, oder daß ihrer Viele einen Namen haben, damit sie genannt werden. Darum ist es mit den Geschlechtsregistern ein verworren Ding und ist verboten, und wer darnach mit Vorwitz forschen wollte, der wird damit vergebliche Arbeit thun. Die Rechnung aber der Zeit ist nöthig, auch sonderlich um der Ursache willen, daß die Juden der Zukunft Christi halben damit überweisen werden mögen, davon man sehr klare Zeugnisse haben kann aus der Ordnung der Zeit. Wiewohl man zwar dieselbigen auch nicht so gar genau fassen kann; denn: jedoch ist wenig daran gelegen, wenn man es auch so gar genau nicht haben kann.

\quellseite{225}

Hernach aber wollen wir von der Zeitrechnung Jakobs mehr sagen.

80. Suchoth, Pnuel und Mahanaim sind Städte, über dem Jordan gelegen, im Stamme Gad. Dasselbe Land haben hernach die Araber eingenommen bis an den Jordan, darum hat es den Namen Arabien bekommen. Im Buch der Richter am 8. Capitel V. 15–17. wird erzählt, daß Suchoth und Pnuel von Gideon sind zerstöret worden darum, daß sie dem Volke, das müde war, nicht haben Brod geben wollen.

81. So hat nun Gott dem Patriarchen Jakob endlich nach der Knechtschaft, Elend und so viel Mühe und Anfechtung Frieden und Ruhe gegeben, daß er mit seiner Familie und Vieh sicher und friedsam hat wohnen und für sich selbst seiner Haushaltung warten können, wie einem Hausvater gebühret. Derhalben baut er nun ein Haus, lehrt und regiert.

82. Und der Heilige Geist schämt sich nicht, dieses alles so fleißig zu beschreiben. Zum ersten darum, daß er damit anzeigen möge den Wechsel oder Veränderung in dem Leben der Heiligen, da Trübsal und Trost also mit einander pflegen zu wechseln und Eins auf das Andere zu folgen. Wie es Gott in der Natur verordnet hat, daß Tag und Nacht, Winter und Sommer auch also mit einander wechseln. Und Gott versucht seine Heiligen darum aus großer Güte und Barmherzigkeit also, auf daß sie im Frieden und Ruhe nicht sicher und faul werden, in der Noth und Widerwärtigkeit aber nicht verzagen sollen. Denn das Kreuz ist nöthig, das Fleisch zu demüthigen, daß es dem Geiste nicht widerstrebe noch über ihn herrsche. Die Trübsal aber, die kein Ende oder Abwechslung hätte, würde den Geist und das Herz gar verzagt machen und zur Verzweiflung treiben. Darum mäßiget es Gott also und wechselt Trübsal und Trost mit einander aus großer Güte, nach der Regel Pauli 1 Cor. 10, 13.: „Gott ist getreu, der uns nicht lässet versuchen über unser Vermögen.“ Er läßt zwar die Versuchung oft zum höchsten und zum unaussprechlichen Seufzen kommen; er macht aber doch, „daß die Versuchung so ein Ende gewinne“, daß wir nicht gar versinken, sondern „daß wir es können ertragen“; wie St. Paulus sagt 2 Cor. 4, 9.: „Wir werden untergedrückt, aber kommen nicht um“ etc.

83. Dieses ist die Lehre der ganzen heiligen Schrift, so uns an diesem Ort in dem Exempel Jakob vorgehalten wird, nämlich, daß es also immer mit einander wechselt, Trübsal und Trost, um der Ursache willen, daß der sündliche Leib getödtet werde, auf daß er sich nicht überhebe und hoffärtig werde; darnach, daß der Geist von Traurigkeit nicht verschlungen und von Schrecken nicht gar ersticket werde. Denn Gott will haben, daß wir uns vor beiderlei hüten sollen, nämlich, daß wir nach dem Fleisch nicht stolz und hoffärtig werden und nach dem Geiste nicht gar verzagen, sondern daß wir auf der Mittelstraße bleiben zwischen der Traurigkeit und Freude, zwischen der Ehre und Schande. Denn also hat der Patriarch Jakob viel Anfechtungen gehabt; aber nach denselbigen allen hat er wiederum Trost empfangen, da er das versöhnte Angesicht Gottes und seines Bruders gesehen hat. Denn da ist seine Seele genesen, die Macht ist vergangen und die Sonne scheint ihm nun wieder, da er an Pniel vorüber ist. Bald hernach wird eine andere Anfechtung folgen.

84. Die andere Ursache, darum uns der Heilige Geist dieses vorhält, ist, daß er damit anzeigen will, daß ihm alle Werke der Heiligen oder Gläubigen, wie gering und kindisch sie auch sein mögen, angenehm und gefällig sind wie gute Früchte vor Gott. Es werden nicht allein die hohen und theologischen Tugenden der Heiligen beschrieben, als, da sie mit dem Tode, mit den Sünden gekämpft und andere Anfechtungen gehabt, dazu dieselbigen auch mit großer Gefahr überwunden haben; sondern auch diese schlechten, häuslichen und geringen Werke: auf daß wir lernen unser Leben dermaßen anrichten, daß wir dessen mögen gewiß sein, daß wir Gott in allem unserem Dienste und Werken wohlgefallen. Ich bete nicht allezeit, ich rede nicht allezeit von Gesetz des Herrn, und kämpfe auch nicht stets mit der Sünde, Tod und Teufel; sondern ich ziehe meine Kleider an, ich schlafe, spiele mit meinen Kindlein, ich esse und trinke etc. Aber dieses alles, wo es im Glauben geschieht, wird von Gott dafür gehalten, daß es recht und wohl gethan und eitel gute Werke vor ihm sind.

85. Dieses soll man den Leuten fleißig einprägen wider die greuliche Heuchelei der Mönche, welche Heuchelei mit ihrer abenteuerlichen Gestalt der Menschen Augen und Sinne bezaubert, daß die Leute gar eingenommen werden und sich derselben Heuchelei hoch verehren, die täglichen und häuslichen Werke aber verachten und davor einen Ekel haben, als die da weltlich und von Gott verachtet seien: so doch die Mönche und Heuchler in ihrem Leben viel ungeschicktes Dinges und ungereimter Werke haben. Aber man soll wissen, daß auch ein jedes vergebliche Wort der Gottlosen unserm Herrn Gott mißfällt; ja, es ist auch ihr Opfer und der allerheiligste Gottesdienst vor Gott eitel Sünde und Greuel. Den Gläubigen aber dient Alles zum Besten, Röm. 8, 28, auch ihre Schwachheit, Irrthum und Sünde. Davon weiß die Vernunft und Weisheit des Fleisches nichts; darum urtheilt sie auch nicht recht und kennt auch die Werke der Heuchler und Heiligen nicht recht unterscheiden.

86. Jakob aber zieht von Suchoth gen Salem und kauft daselbst ein Stück Ackers, und richtet einen Altar auf, den Namen Gottes daselbst anzurufen und die himmlische Lehre zu predigen. Dasselbe Stück halte ich, daß es billig zu dem folgenden Capitel soll gesetzt werden und daß man damit das vierunddreißigste Capitel anfangen soll.

\chapter{Das vierunddreißigste Capitel}

\section[Erster Theil]{Erster Theil. Von Jakobs Aufenthalt und Verrichtungen in Sichem, und wie ihm seine Tochter Dina geschmäht worden}

\begin{quote}
\textbf{Cap. 33. V. 18. Darnach zog Jakob gen Salem, zu der Stadt des Sichem, die im Lande Canaan liegt (nachdem er aus Mesopotamien kommen war), und machte sein Lager vor der Stadt.}
\end{quote}

1. Die erste Frage in diesem Capitel gehört in die Grammatik, nämlich von der Bedeutung dieser Worte, „Salem“ und „Sichem“. Und die hebräischen Grammatiker und die andern Ausleger alle sind nicht eins mit einander, was die besondere Bedeutung dieser Namen sei. Denn ich habe droben auch etlichemal gesagt, daß der Worte viel sind, deren Bedeutung die Hebräer auch selbst nicht wissen. So gar ist der Gebrauch und Verstand dieser Sprache gefallen, daß er niemals recht vollkommen wiederum kann aufgerichtet werden. Und sind nicht allein die Worte, sondern auch die Art zu reden und wie die Worte sollen zu einander gesetzt werden, vielfältig und auf mancherlei Weise verdunkelt; daher es kommt, daß wir die Bedeutung und Sinn und Zweck vieler Redensarten und Sprüche nicht wissen können. Oder wenn ja der Verstand dieser Sprache auf einigerlei Weise wieder könnte aufgerichtet werden, so muß solches durch die Schriften geschehen, die aus dem Neuen Testament den Verstand der heiligen Schrift haben.

2. Die Juden verdunkeln den rechten natürlichen Verstand mit ihren Aequivocationen, indem sie die Worte hin und her ziehen, daß sie mancherlei Bedeutung haben sollen; und dasselbe thun sie mit Fleiß, auf daß sie damit allerlei Fragen und Irrthum erregen mögen.

3. Aber ein fleißiger und guter Grammatiker soll Fleiß anwenden, daß er die mancherlei Bedeutungen auf Eine ziehen möge, die da gewiß sei; denn es muß ja Ein Wort Eine rechte und eigentliche Bedeutung haben. Darum soll er einen jeglichen Ursprung der Worte anzeigen, dazu ihre Eigenschaft, darnach auch die Figuren, Metaphern, Anspielungen und Redensarten. Das thun der Juden Rabbinen nicht. Denn wo die Dinge fallen, daselbst fallen auch die Worte und ihre Bedeutung. Wie niemand die Rede der Juristen, der Ärzte und Theologen verstehen könnte, wenn man die Dinge, davon sie reden, verloren hätte. Und deß haben wir ein Exempel an der Theologie der Scholastiker. Denn jetzt ist unter unseren Zuhörern niemand, der den Scotus, Thomas oder dergleichen Lehrer verstehen könne. Denn die Dinge, davon sie reden, und der Gebrauch ihrer Worte ist gar veraltet und nunmehr ungewöhnlich. Daselbige würde mit andern Wissenschaften auch also geschehen, wenn man sie nicht immer gebrauchte und übte. Denn wo die Dinge, davon geredet wird, vergehen, da wird auch der Verstand der Worte immer mehr und mehr verdunkelt.

\quellseite{226}

4. So wird nun gefragt: Ob das Wort „Salem“ an diesem Ort zu anderen Worten vermaßen gesetzt sei, daß es ein Gattungswort oder aber ein Eigenname sein soll? Wo man es als ein Gattungswort versteht, heißt es Friede oder friedsam, stillkommen, bisweilen auch Vergeltung. Wenn man es nun in solcher Bedeutung adjectivisch verstehen will, so wird dieses der Verstand sein: Jakob ist unverletzt, mit gutem Frieden und sein stille, das ist, ohne Anfechtung, dazu frisch und gesund über den Jordan gezogen und in die Stadt Sichem gekommen. Diese Auslegung schickt sich nicht übel. Denn er hat nun die Anfechtung und den Kampf überwunden, welchen Kampf er mit dem Manne, dem Kämpfer, geführt hat: er hat nun Trost empfangen und steht Alles wohl und gut mit ihm. Darum ist er in seinem Herzen wohl zufrieden, baut ein Haus und richtet eine Hütte auf.

5. Die Andern ziehen es auf die Historie von der Hüfte Jakobs, die ihm verrenkt war, und legen es also aus: Jakob ist unverletzt, das ist, frisch und gesund gekommen, also daß er nicht mehr gehinkt hat. Und diese Auslegung strafe ich auch nicht. Denn es hat wohl geschehen können, daß ihm die Hüfte ist wiederum zurechtgebracht und gesund geworden durch die Arzneikunst; welche Kunst ohne Zweifel die Väter sehr wohl gekannt haben. Denn sie haben nicht allein die Kunst gehabt, sondern auch sehr viel der besten Arzneien: sie haben Balsam und dergleichen viel andere Dinge mehr gehabt, die wir jetzt nicht haben. Darum disputire ich von der Sache an sich nicht, und glaube gern, daß Jakob die Hüfte wiederum geheilt und gesund geworden sei. Das ist aber gleichwohl die besondere Bedeutung des Wortes nicht. Und dieweil die Hebräer andere Worte haben, damit dieser Verstand eigentlich hätte können angedrückt werden, so würden sie ohne Zweifel auch ein ander Wort, das sich noch etwas besser geschickt hätte, gebraucht haben. Denn es scheint, daß dieses eigentlicher sei, wenn du es also geben würdest: Jakob ist fein stille und friedsam gekommen.

6. Die dritte Meinung, daß es der Name einer Stadt sei, kann man nicht leichtlich beweisen oder annehmen. Es ist Jerusalem auch Salem genannt worden, wie wir droben in 24. Capitel gehört haben, und man kann nicht gewiß sagen, daß Sichem bei den Alten Salem sollte genannt worden sein. Wie an einem andern Ort von der Stadt Luz gesagt wird, die hernach Bethel ist genannt worden. Denn man findet nicht alle Namen der Städte bei den Juden. Ich will aber darüber mit niemand streiten; denn es ist unserem Glauben oder unserer Religion ohne Gefahr, man verstehe dieses Wort als ein Gattungswort oder aber als einen Eigennamen.

7. Aber der Name der Stadt Sichem ist bei allen Nachkommen geblieben und ist sehr berühmt, vornehmlich in den Historien der Könige. Sichem nennt es ein närrisch Volk, das zu Sichem wohnt. Wie die Deutschen von den Schwaben und Bayern zu sagen pflegen. Es ist ein stolz hoffärtig Volk gewesen, das sich nur auf Wohllust begeben hat. Und sie wollen sagen, daß diese Stadt nahe bei den zwei Bergen Garizim und Ebal gelegen gewesen. Die Andern sagen, daß sie bei Jericho gelegen gewesen sei.

8. Das weiß man aber wohl, daß Silo, da die Hütte gewesen ist, sehr nahe bei Sichem gelegen gewesen, und ist sehr glaublich, daß Silo eben derselbe Ort gewesen ist, da Jakob gewohnt und da Christus mit dem Samaritanischen Weibe geredet hat, Joh. 4, 7. ff., und ist eben daselbst die Hütte gewesen unter Josua und Eleasar eine lange Zeit bis auf Samuel. Und also haben die Nachkommen auch daher Ursache genommen zu mancherlei Abgötterei, dazu sie durch das Ansehen und Exempel der Väter bewogen worden sind, denen sie haben nachfolgen, und an demselben Ort wider Gottes Wort opfern und Gott dienen wollen. Wie hernach Jerobeam Sichem wieder gebaut und dahin den Gottesdienst gewendet hat vom Tempel zu Jerusalem, 1 Kön. 12, 25. ff.

9. Wiewohl ich aber, da so mancherlei Meinungen sind, nichts Gewisses zu antworten oder zu sagen weiß, weil die Rabbinen der Juden selbst der Sache nicht eins sind: so gefällt mir doch diese Auslegung am besten, daß Jakob von so mancherlei Plage sei erlöst worden, und voll Glaubens und Trostes, dazu auch gesund und stark in die Stadt gekommen sei, die da Sichem genannt wird von den Fürsten, oder aber von dem, der sie gebaut hat; wie viele andere Städte mehr den Namen bekommen haben von denen, die sie gebaut oder gestiftet haben. Und diese Stätte ist zu allen Zeiten unter dem Volk sehr berühmt gewesen, wie gesagt ist.

\subsection{I}

\begin{quote}
\textbf{Cap. 33. V. 19. Und kaufte ein Stück Ackers von den Kindern Hemor, des Vaters Sichem, um hundert Groschen; daselbst richtete er seine Hütte auf.}
\end{quote}

10. Es hat bisher der Patriarch Jakob in diesem Lande nichts Eigenes gehabt, aber nun kauft er endlich ein Stück Ackers; wiewohl die Verheißung gewiß da steht, daß ihm Gott das ganze Land geben würde, so kauft er doch gleichwohl dieses Stück; wie droben, Cap. 23, 17., sein Großvater Abraham auch einen Acker gekauft hat gegen Mamre über und eine zwiefache Höhle, der Sarah zum Begräbniß. Darum sagt Stephanus recht in den Geschichten der Apostel am 7. Cap. V. 5.: „Er gab ihm kein Erbtheil drinnen, auch nicht eines Fußbreit“, nämlich, bei ihrem Leben; sondern sie haben es beides gekauft, was sie darin gehabt haben, Abraham zum Begräbniß, und Jakob zur Wohnung, da er nun zur Ruhe gekommen und den großen schweren Anfechtungen entronnen war.

11. Er gibt aber hundert Groschen für das Stück Ackers; daraus wohl zu sehen ist, daß es nicht sehr groß wird gewesen sein. Die Andern sagen, er habe hundert Lämmer dafür gegeben. Du magst es aber nun verstehen, daß er Geld oder Lämmer dafür gegeben habe, so hat doch der Acker nicht können groß sein, denn er für so ein geringes Geld gekauft hat, oder sie werden wohlfeiler gewesen sein, denn sie jetzt bei uns sind. So hat er nun eine eigene Wohnung bei der Stadt Sichem, die in der heiligen Schrift sehr gerühmt wird um dieses Erbes willen, so Jakob daselbst gehabt, und um Silos willen, das nahe dabei gelegen gewesen ist, da die Hütte war und da das Volk zusammen gekommen ist.

\begin{quote}
\textbf{Cap. 33. V. 20. Und richtete daselbst einen Altar zu, und rief an den Namen des starken Gottes Israel.}
\end{quote}

12. Jakob gedenkt an diesem Ort zu bleiben, darum richtet er daselbst einen Altar zu. Er ist aber noch nicht wieder zu seinem Vater Isaak gekommen; sondern hernach, halte ich, werde er beide, seinen Vater und auch seinen Bruder, oftmals besucht haben, da er sein Gesinde und Vieh an diesem Ort gelassen hat. Und Hebron ist nicht sehr weit von Sichem gelegen gewesen, nämlich, bei zehen Meilen; Hebron ist fünf Meilen von Jerusalem gelegen gewesen, Sichem aber drei oder vier Meilen. Darum hat er seinen alten Vater und Bruder leichtlich besuchen können; wiewohl die Schrift das ausdrücklich nicht sagt, so folgt es doch sehr wohl, sintemal sich die Brüder mit einander versöhnt haben. Und ich halte dafür, daß Jakob zum wenigsten acht Jahr an diesem Ort werde gewohnt haben.

13. Mose aber beschreibt auch, wozu Jakob den Altar aufgerichtet habe, da er sagt: „Und rief an den Namen des starken Gottes Israel.“ Das lautet im Hebräischen: Und rief an den Gott der Götter Israels. Und dieses ist ein trefflicher Text; denn es wird uns fast an allen Orten im Alten Testament unter dem Namen Gottes Christi geoffenbart. Und Mose begreift an diesem Ort zugleich die drei Personen, wie er sonst an andern Orten auch oftmals zu thun pflegt. Der Mann, der mit Jakob droben gerungen hat, ist Christus Jesus, Gottes Sohn, unser Herr, selbst gewesen, der also hat wollen geoffenbart werden, daß er eine unterschiedene Person wäre vom Vater und Heiligen Geist und doch gleichwohl ein einiger GOtt wäre. Also deute ich hier das erste Wort el, daß es der Sohn Gottes sei, der da hat sollen Mensch werden. Aber das Wort elohe deute ich auf den Vater und Heiligen Geist. Und diese drei Personen sind der Eine GOtt Israel.

14. Die Juden aber lesen und der Text lautet auch also, daß er den Altar selbst also genannt habe, wie sie den Text im Propheten Jeremia am 23. Cap. V. 6. und am 33. Cap. V. 15. 16. auch auslegen, da also steht: „Siehe, es kommt die Zeit, spricht der Herr, daß ich dem David ein gerecht Gewächs erwecken will; und soll ein König sein, der wohl regieren wird, und Recht und Gerechtigkeit auf Erden anrichten. Zu derselbigen Zeit soll Juda geholfen werden, und Israel sicher wohnen. Und dies wird sein Name sein, daß man ihn nennen wird: „Herr, der unsere Gerechtigkeit ist.“ Da steht das Tetragrammaton, wie sie es nennen, Jehovah, das heißt der rechte wahre Gott, der uns gerecht macht. Aber dieses treffliche Zeugniß von der Gottheit des Sohnes Gottes widerlegen sie mit solchen Sprüchen, daß dieser Altar genannt wird deus deorum, der Gott der Götter, der starke Gott Israel; item, Jerusalem, wie Jeremia am 33. Cap. V. 16. steht: Zu derselbigen Zeit soll Juda geholfen werden, und Jerusalem sicher wohnen. Und man wird ihn nennen: der Herr, der unsere Gerechtigkeit ist“, item 2 Mos. 17, 15.: „Und Mose bauete einen Altar, und hieß ihn: Der Herr Nissi“, das heißt, meine Erhöhung, meine Stärke; item 1 Mos. 22, 2., der Berg Morija ist die Ehre Gottes genannt worden.

\quellseite{227}

15. Mit solchen Sprüchen versuchen die Juden, die gewaltigen Zeugnisse von der wahren Gottheit Christi zu verkleinern und zu widerlegen. Aber Jerusalem oder der Altar ist nicht der Herr, ist auch unsere Erhöhung oder Gerechtigkeit nicht: und die Juden bringen doch mit diesem Argument hart auf uns und wollen sich dasselbe nicht nehmen lassen. Wir können es aber leichtlich auslegen aus der Lehre, die wir im Neuen Testament haben; wie an andern Orten gelehret wird.

16. Da nun Jakob Trost empfangen hat, und da er genesen, auch frisch und gesund und erlöset worden ist von dem Schrecken des Todes, und sich die Schmerzen der Hüfte, die ihm verrenkt war, gelegt hatten: richtete er einen Altar auf, daß er Gott danken möge dafür, daß er wiederum war gesund worden und genesen. Denn darum hat er den Altar zugerichtet, daß daselbst das Wort sollte gelehret werden, daß man daselbst beten, Gott danken und sonst andern Gottesdienst verrichten sollte.

17. Denn wo die heilige Schrift sagt, daß Altäre aufgerichtet seien, ist es eben so viel gesagt, als wenn sie sagte, daß man Schulen oder Kirchen aufgerichtet hätte, da man den Gottesdienst verrichten sollte und was dazu gehört. Die Altäre sind nicht aufgerichtet worden, Schauspiele oder unnütze Gepränge daselbst zu treiben, auch nicht zum Abgötter, sondern das Wort daselbst zu predigen. Also wird nun Jakob beschrieben, wie er in seinem Herzen fröhlich und getrost gewesen sei nach der Erlösung; und vergißt doch in solchem Glück Gottes nicht, der ihn erlöset hat, sondern richtet einen Altar zu, lehrt seine Kinder, regiert sein Gesinde und Gemeinde, und ermahnt sie, Gott anzurufen und ihm zu danken. Dieser Friede und Ruhe aber währt ihm nicht lange; denn jetzt neigt sich der Tag zum Abend und die finstere Nacht geht daher.

\subsection{II}

\begin{quote}
\textbf{Cap. 34. V. 1. 2. Dina aber, Leas Tochter, die sie Jakob geboren hatte, ging heraus, die Töchter des Landes zu sehen. Da die sahe Sichem, Hemors Sohn, des Heviters, der des Landes Herr war, nahm er sie, und beflecket sie, und schwächte sie.}
\end{quote}

18. Dieses ist nun eine andere und viel schwerere Anfechtung, denn die droben, Cap. 32. V. 6. ff., gewesen ist, da er erschrocken war, darum daß ihm sein Bruder mit vierhundert Mann entgegen zog; und ist diese Anfechtung voll großer Gefahr, Traurigkeit und Jammer. Denn er sitzt jetzt daheim, der gute fromme Patriarch, hat seine Weiber, Kinder und die Gemeinde in seinem Hause um sich, ist wohl zufrieden, da er nach dem schweren Kampf neulich Trost empfangen hat; und sitzt zwar nicht müßig, sondern richtet das Amt aus eines Hausvaters und Lehrers mit großem Fleiß, nämlich, mit Regieren, Lehren und Beten: und wird alsbald wiederum in eine sehr betrübte Anfechtung geführt, das er nie erwartet hatte. Derhalben ist dieses Exempel uns zur Lehre geschrieben, auf daß wir lernen Geduld haben in Trübsal und Widerwärtigkeit. Denn wo ein solcher großer Mann so greulich umgetrieben und angefochten wird, so muß uns nicht Wunder oder unbillig zu sein dünken, wenn uns auch dergleichen dermaßen widerfahren wird.

19. Es wird aber hier erstlich gefragt nach dem Alter der Dina, darin sie geschwächt ist. Denn es ist nicht wohl glaublich, daß sie schon mannbar gewesen sei, da Jakob erst gen Sichem gekommen war, sintemal er sie nicht länger, als drei oder vier Jahr auf das höchste gehabt hat, sondern sie hat eilf oder zwölf Jahre müssen alt sein, da sie ist beschlafen worden. Und folgt daraus, daß Jakob zu Suchoth oder Sichem zum wenigsten acht Jahre wird gewohnt haben, ehe Benjamin ist geboren worden, deß Geburt im nächstfolgenden Capitel erst beschrieben werden wird. Das ich nun in der ersten Ausgabe meines Chronikons angezeigt hatte, daß Jakob straks zu seinem Vater Isaak gezogen sei, im ersten Jahr, da er wieder gekommen ist, dasselbe habe ich darnach geändert. Und ich halte, daß er fast bei zehen Jahren in diesem Lande gewesen sei, da er aus Mesopotamien gekommen war, und daß er acht Jahre zu Sichem gewohnt habe, da er einen eigenen Acker gehabt, da er sein Vieh und Gesinde hat bleiben lassen; und ist er bisweilen zu seinem alten Vater und zu seinem Bruder gegangen, sie zu besuchen.

20. Daß wir nun diese Jahrzahl also machen, dazu zwingt uns Dina, die wir zum wenigsten müssen zwölf Jahre lassen alt gewesen sein, daß sie mannbar hat sein können. Wiewohl dasselbe auch fast selten und schier unmöglich ist, vornehmlich zu unserer Zeit, daß ein Mägdlein von zwölf Jahren schwanger werden sollte, darum daß es mit der Geburt viel Mühe und Arbeit hat und gemeiniglich die Gefahr Leibes und Lebens dabei ist: derhalben auch mehr Stärke und ein vollkommenes Alter dazu gehört. Jedoch mag es wohl sein, daß zu der Zeit die Natur und Kräfte der Menschen etwas stärker gewesen sind, darum daß sie bessere Lust gehabt haben, dazu auch bessere und gelegene Städte, da sie gewohnt haben, und da sie sich mit Essen und Trinken mäßig haben wissen zu halten. Und die bürgerlichen Rechte halten, daß Knaben von vierzehn Jahren, Mägdlein aber von zwölf Jahren zum Ehestande sollen tauglich und alt genug sein. Aber in dem Alter verstehen sie kaum, daß sie leben oder daß sie Mägdlein sind; darum können wir uns heutiges Tages nach diesem Zeitmaß nicht richten.

21. Und wenn schon Dina dazumal zum Ehestande ist alt genug gehalten worden, so wird sie doch im Text auf Hebräisch Jaldah genannt, das heißt, Mägdlein. Denn es ist kein Gebrauch gewesen, daß man die Mägdlein so bald sollte zur Ehe gegeben haben. Denn also haben wir droben, Cap. 29., von Rebekka gehört, daß sie fast bei dreißig Jahren gewesen ist, da sie den Isaak ehelich vertraut worden. Darnach ist die Natur der Menschen nicht allein in diesem letzten Alter der Welt etwas schwächer und kraftloser geworden, sondern ist auch über das sehr geschwächt worden deswegen, daß jetzt die Leute nicht so mäßig leben und daß auch jetzt so gesunde Luft nicht ist; zudem sind die Städte, da wir wohnen, auch nicht so gelegen und gesund. Daher kommt es, daß wir auch vor der Zeit alt werden, und daß wir am Leibe mit mancherlei Krankheit beschweret und geplagt werden.

22. Es disputiren aber ihrer Viele davon: Ob dem Alter der Menschen etwas abgegangen sei und ob die Kräfte und Jahre der Menschen an ihrem Leibe abgenommen haben? Und Etliche halten es dafür, daß daran nichts Sonderliches abgegangen sei, sintemal heutiges Tages noch viel Leute gefunden werden, die an ihr hundertstes Jahr gekommen sind, oder die das neunzigste Jahr erreichen; wie ich solche Leute selbst gesehen habe. Dieweil aber Mose im 90. Psalm V. 10. sagt: „Unser Leben währet siebzig Jahr, und wenn’s hoch kommt, so sind’s achtzig Jahr“; und daß zu derselbigen Zeit die, so viele Jahre alt geworden, geachtet sind, daß sie an die äußerste Grenze des Alters gekommen sind: so scheint es, daß die Jahre und Kräfte der Menschen an ihrem Leibe jetziger Zeit gar nicht abgenommen haben, da man deren Viele sieht, die siebzig und achtzig Jahre alt werden.

23. Die Andern ziehen den Spruch an aus dem 102. Psalm V. 26. 27.: „Herr, du hast vorhin die Erde gegründet, und die Himmel sind deiner Hände Werk. Sie werden vergehen, aber du bleibest. Sie werden alle veralten, wie ein Gewand; sie werden verwandelt, wie ein Kleid, wenn du sie verwandeln wirst.“ Da zeugt ja die heilige Schrift, daß Himmel und Erde veralten wie ein Kleid; darum ist es gewiß, daß die Kräfte Himmels und der Erde, und aller Creaturen, die im Himmel und Erde sind, abnehmen, dazu auch die Thiere, Menschen, Zweigelein, die da gepflanzet werden, und Kräuter, an welchen allen so viel Stärke und Kraft nicht mehr wie vor Zeiten gewesen ist. Adam hat neunhundert und dreißig Jahre gelebt und alle Väter haben insgemein vor der Sintfluth länger gelebt; nach der Sintfluth aber hat das Alter der Menschen abgenommen, und alle Creaturen, dazu Himmel und Erde, haben auch abgenommen. Die Sonne hat solche Kraft nicht, die Erde ist so fruchtbar nicht mehr wie zuvor. Denn von der Sintfluth ist die ganze Natur allerdings sehr besteckt und verderbt worden. Ja, daß noch mehr ist, die Sintfluth hat auch das Paradies verwüstet und zerstört. Die vier Wasser, davon Meldung geschieht in diesem ersten Buch Mose am 2. Cap. V. 10. ff., sind nach der Sintfluth nicht mehr aus dem Paradies geflossen, das Land hat eine andere Gestalt gewonnen, man hat es auf eine andere Weise bauen müssen nach solcher greulichen Zerstörung und Verwüstung, und das Alter und die Kräfte aller Creaturen nehmen immer allgemach mehr und mehr ab.

\quellseite{228}

24. Darum bekommt eine Mannsperson erst im vierundzwanzigsten Jahr solche Stärke am Leibe, daß sie zum Ehestande und sonst andere Dinge auszurichten tauglich ist. Ein Mägdlein wird erst im siebzehnten oder achtzehnten Jahr tauglich, die schwere Mühe und Arbeit der Geburt zu ertragen. Zur Zeit Davids ist das gemeine Alter eines Menschen siebenzig Jahre gewesen, und solche Leute werden jetziger Zeit auch gefunden, aber doch sehr wenig, und ist jetzt fast ein Wunder, wo die Menschen siebenzig oder achtzig Jahre alt werden. Der gemeine Lauf ist, daß sie fünfzig oder sechzig Jahr alt werden. Die darüber noch länger leben, die überschreiten das gemeine Ziel. Mose aber redet in seinem Gebet im 90. Psalm V. 10. von dem gewöhnlichen Ziel; denn das Leben über siebenzig oder achtzig Jahr ist gar untauglich, daß der Mensch darin etwas ausrichten sollte, item, zur Kinderzucht und auch zum Hausregiment. Und David ist auch gestorben, da er siebenzig Jahr alt gewesen. Heutiges Tages sind fast die Sechzigjährigen zu allem Regiment untauglich. So gar sehr haben die Kräfte am Leibe der Menschen, und sonst an allen Dingen abgenommen: und kommt über das auch noch dazu, daß jetzt die Leute so gar unmäßig leben, daß sie so unmäßig fressen und saufen, damit wir unsere Leiber wie mit großen Sturmböcken selbst verderben.

25. Nun wollen wir auf die Historie kommen, darin ein sehr betrübtes Unglück des Patriarchen Jakob begriffen wird, welches nicht allein ihm beschwerlich und betrübt, sondern auch dem ganzen Geschlechte, Kindern und Gesinde schändlich gewesen ist. Denn es ist die größte Schmach und Schande, daß ein Mägdlein, so noch fast ein Kind ist, soll beschlafen werden. Wie sie denn Mose nicht bethulah oder almah nennt, sondern jaldah auf Hebräisch, das heißt ein kleines Kind; und Sichem ist auch noch ein Knabe gewesen.

26. So sagt nun Mose, Dina sei heraus gegangen, die Töchter des Landes zu sehen, das ist, die Weiber des Landes. Es will damit anzeigen, daß sie sicher gewesen ist, und sich keiner Gewalt, und viel weniger, daß sie sollte beschlafen werden, besorgt hat, sintemal sie noch nicht mannbar war; und darüber ist sie zu Falle gekommen. Darum wird die Schmach davon größer, daß Sichem ein solch Mägdlein genommen und beschlafen hat, die so jung und noch damals zum Ehestande nicht tauglich war. Und dieses ist auch die Ursache gewesen, daß ihre Brüder darüber dermaßen sind bewegt und erzürnt worden. Sie werden sich aber, wie wir hören werden, allzugrausam in ihrer Rache an Sichem erzeigen.

27. Der Vater wird ohne Zweifel darüber den größten und sehr erbärmlichen Schmerz gehabt haben. Denn er wohnte im Lande Canaam acht Jahre, nachdem er aus Mesopotamien gekommen war, und in denselbigen Jahren baut er ein Haus und richtet eine Hütte auf für sein Vieh; und ist nicht mehr ein Fremdling oder Gast, sondern wohnt im Lande, richtet daselbst eine Kirche oder Gemeinde an, richtet dem Herrn einen Altar zu, predigt von der Verheißung und den Geboten Gottes, daß dieselben unter allen Nachkommen ausgebreitet werden mögen, und hat Frieden und Ruhe zum wenigsten nach gemeiner Weise, wie in diesem Leben bei dem Welt- und Hausregimente pflegt Fried zu sein, da gleichwohl viel Beschwerung und Schaden mit unterläuft beide im Welt- und Hausregimente. Er thut auch häusliche Werke dabei, mehret seine Heerde, sammelt Gut, schafft seinem Gesinde ihre Nahrung, duldet die Beschwerungen, die da pflegen mit unterzulaufen, als, Feindschaft der Nachbarn, Diebstahl, Räuberei und sonst heimliche Nachstellungen. Dieses ist das gemeine Unglück, so täglich vorfällt, welches nicht gebessert oder verhütet werden kann, sondern man muß solches zugleich mit andern Leuten tragen, und Gott danken, daß unterdeß noch gemeiner Friede und Ruhe bleibt. Es ist übrigens eine solche Ruhe, die nicht ohne Schweiß des Angesichts, ohne Dornen und Disteln sein kann. Und solchen Frieden und Ruhe hat zwar Jakob diese acht Jahre über gehabt.

28. Jetzt aber kommt nun eine ungewöhnliche und greuliche Anfechtung, die man nicht zählen kann unter die täglichen Dornen und Disteln. Denn es mag es ein Jeder bei sich selbst abnehmen, wie das einem Vater so ein großer Schmerz und ein schweres Kreuz sei, wenn er eine einzige Tochter hat, und soll an derselbigen die höchste Schmach sehen und daß sie ihm dermaßen geschändet wird. Es ist wohl keine geistliche Anfechtung des Glaubens, der Hoffnung und Geduld, wie die droben gewesen sind: aber unter den Anfechtungen, so sich im Hausregimente begeben können, ist es doch die aller greulichste und unleidlichste, nämlich, wo einem Vater sein einzig Töchterlein, die noch nicht mannbar ist, mit Gewalt soll geschwächt werden, und daß solches geschehen soll an einem Orte, da es still und guter Friede ist, nicht von den Feinden, sondern von einem Nachbar, von einem Freund, ja, von einem, der ein Fürst und Beschützer und Vater des Vaterlandes ist, an welches Herberge und Treue er bisher gutes Genügen gehabt und damit wohl zufrieden gewesen war. Von dem muß er nun wider alles Erwarten solche große Schmach leiden. Darum ist dieses eine sehr böse und unbillige Anfechtung.

29. Es werden aber die größten und heiligsten Männer dermaßen beschrieben und abgemalt, daß sie nicht allein dem Schweiß ihres Angesichts, Disteln und Dornen, so täglich im Hausregiment vorfallen, müssen unterworfen sein, sondern daß sie auch auf das härteste und übelste von Feinden und gottlosen Menschen müssen geplagt werden.

30. Es ist wahrlich ganz unbillig und kläglich, daß ein solcher großer Patriarch dermaßen soll geschmäht werden, der so viel treffliche Verheißungen gehabt, daß er in Gottes Schutz sein würde, daß ihn die Engel schützen und bewahren sollten. Wo sind nun die Heerschaaren der Engel und Machanaim? Wo ist nun der große herrliche Sieg, damit er Gott und Menschen überwunden hat? Wo ist nun die herrliche Frohlockung, da zu ihm gesagt worden ist: „Du hast mit Gott und mit Menschen gekämpft, und bist obgelegen“? Wer bewahrt ihn hier? Wer wacht über ihn? Es sehen Gott und die Engel durch die Finger, sie verstellen sich; Gott weiß nichts darum, und hält sich nicht anders, denn als ob er es nicht wisse und auch nicht sähe, daß Jakobs Tochter hingriffen und beschlafen wird. Denn er läßt solches geschehen, da die Engel stille sind und ihr Amt nicht ausrichten. Es ist fürwahr ein groß und erbärmlich Unglück, das billig nicht frommen gotteligen Menschen, sondern Gottes Feinden, Heiden und gottlosen Leuten widerfahren sollte.

31. Darum muß man lernen Geduld haben, wenn etwa uns auch solche Betrübniß und ungewöhnlich Unglück widerfahren würde, damit es sich ansehen läßt, daß man es auf keinerlei Weise und mit keiner Geduld ertragen könnte; wie zwar diese Anfechtung gewesen ist, die alle menschliche Geduld übertrifft; wie das hernach die greuliche That der Brüder Dinas bezeugen wird, damit sie die Sünde, die an ihrer Schwester begangen war, gerächt haben.

32. Denn also müssen die allerheiligsten Männer solchen greulichen Unglück unterworfen sein, und läßt sich die Vernunft dünken, es werde das Reich Gottes dermaßen verwaltet, daß die Gottseeligen und Gottlosen gleich geachtet werden, ja, daß die Gottlosen viel besser Glück haben denn die Gottseeligen.

33. Warum läßt es aber Gott geschehen, daß dieser so heilige Patriarch mit solchem Kreuz hat müssen beschweret werden, als ob er nicht heilig, und Gott angenehm und gefällig wäre? Antwort: Das ist um unserwillen geschehen, auf daß wir lernen Geduld haben und uns trösten können in unserer Widerwärtigkeit, und den Finger auf den Mund legen, wenn uns solches Unglück auch begegnen würde. Denn wir sind nicht besser denn diese großen Männer, darum sollen wir auch kein sonderlich Glück begehren, sondern sollen uns gewöhnen an solche Uebung und Prüfung des Glaubens, des Trostes und der Geduld; sintemal es Gott also wohl gefallen hat, daß er über den Schweiß des Angesichts und die gewöhnlichen geistlichen Trübsale auch solch ungewöhnlich und unerhörtes scheußlich Unglück über uns verhängt.

34. Es wird aber ohne Zweifel dazumal dem frommen Patriarchen eingefallen sein, daß er gedacht hat: Ach ich armer elender Mann, was habe ich doch gesündigt? Ich bin vielleicht sicher gewesen, habe nicht ernstlich und brünstig gebetet, oder unter meinem Gesinde und Gemeinde ist etwas begangen worden ohne mein Wissen. O lieber Herr Gott, erbarme dich mein, und verzeihe mir die verborgenen Fehle, Pf. 19, 13. Mit solchen Gedanken werden die Gottseeligen geängstet, daß sie erschrecken und zittern, als ob sie um irgend einer sonderlichen Sünde willen in solchen großen Jammer und Unglück geworfen werden. Und diese Bekümmerniß und Schmerz, so Jacob gehabt, ist davon größer geworden, weil er Dina als seine einzige Tochter und die noch nicht mannbar war, so heftig lieb gehabt, zudem auch, daß er ein Fremdling war und diese Gewalt nicht hat rächen können, und hat vielleicht gedacht, er müßte sich noch größeres Unglück besorgen.

\quellseite{229}

35. Es wird aber auch davon disputirt: Was das für ein Ausgehen gewesen sei, da der Text sagt, Dina sei herausgegangen, die Töchter des Landes zu besehen? Denn die Schrift erklärt es nicht, und Einer sagt dieses, der Andere sagt ein Anderes dazu. Die Juden sagen, sie sei aus Vorwitz herausgegangen, die Sitten und Herrlichkeit des Volkes oder der Stadt zu beschauen. Und ist wohl glaublich, daß etliche Spiele, Tänze oder Hochzeiten daselbst gehalten worden sind. Da hat Dina sehen wollen die Töchter des Landes, wie dieselbigen geziert und geschmückt gewesen, und wie schön sie wären. Und dasselbige zeigt der Text auch an, wie es scheint, nämlich, daß sie vorwitzig gewesen sei, sintemal sie ohne Erlaubniß ihres Vaters und ihrer Mutter und dazu auch allein ohne Gefährten ist hinausgegangen. Sie ist also sicher und trotzig, weil es noch ein junges Mägdlein war, und sich vor solcher Gefahr nicht gefürchtet hat, daß sie ihre Keuschheit verlieren möchte. Darum scheint es, als habe sie aus Vorwitz gesündigt, darum daß sie zu den Töchtern und ihrer Gesellschaft herausgegangen ist und hat ihre Eltern nicht zuvor um Rath gefragt. Wie die Mägdlein von Natur dazu Lust haben, daß sie mit andern Jungfrauen, die ihnen gleich und ihre Nachbarn sind, Gemeinschaft haben mögen. Ihre Brüder sind unterdeß nicht daheim gewesen, sondern auf dem Felde; Lea ist mit ihrem Töchterlein Dina dahin, die geht ohne Vorwissen der Mutter heraus.

36. Darum sollen die Mägdlein dieses Exempel wohl merken und man soll ihnen dasselbe fleißig einschärfen, daß sie nicht wild seien, daß sie nicht Freundschaften suchen und gern unter der Thür stehen, sondern daß sie lernen im Hause bleiben und nirgendhin gehen ohne Erlaubniß ihrer Eltern und ohne Gefährten. Denn der Teufel stellt der Keuschheit dieses weiblichen Geschlechtes heimlich nach, welche Geschlecht von Natur und Art schwach, leichtsinnig und närrisch ist, und derhalben auch den listigen Nachstellungen des Teufels sonderlich unterworfen ist.

37. Dina hat wollen mit den andern Mägdlein, ihren Nachbarinnen, Bekanntschaft machen und hat wollen ihren Sitten und Schmuck lernen nachfolgen. Denn ihr Vater wohnte außerhalb der Stadt; sie aber waren hinausgegangen zum Gespränge auf irgend eine Hochzeit: da ist Dina auch hinzugelaufen, solch Gespränge anzuschauen, und ist also von Sichem gesehen worden, der in demselben Gespränge oder Tanze auch gewesen ist mit den Töchtern des Landes. Da er sie nun gesehen hat, raubt und demüthigt er sie. Denn also reden die Hebräer von der Schwächung. Daselbst ist nun sehr unbillig und übel gethan. Da er sie gesehen hat, hätte er sie lieb haben und zum Weib begehren mögen; aber er folgt seiner bösen Lust und Unzucht: da er sieht, daß sie vor den Töchtern des Landes sehr schön von Gestalt war, wird er entzündet von Liebe und nimmt sie mit Gewalt hin, bittet oder begehrt weder der Eltern, noch des Mägdleins Willen, sondern schwächt sie mit Gewalt.

38. Da pflegt Mord und Jammer zu folgen; wie das alle Historien bezeugen, daß es niemals ohne Mord und Blutvergießen sei abgegangen, wo irgend Jungfrauen weggeführt und geschändet sind worden. Und man hat dessen an dieser That auch ein Exempel. Denn Gott übersieht solchen Muthwillen und Unsinnigkeit der Unzucht nicht und läßt es nicht ungerächt, wenn jemand den hitzigen Begierden und der Lustseuche nachhängt. Denn er will haben, daß wir dem Fleische widerstreben sollen, das wider uns streitet, auf daß es uns nicht in allerlei Laster werfe, so man ihm den Zaum läßt.

39. Mose sagt hier deutlich: Er raubte sie; „er nahm sie.“ Das ist zu hart, eines Andern Tochter nehmen, der in derselben Stadt oder dabei ein Einkömmling oder Gast war, der unter dem Schatten und Schutz des Herrn Hemor in guten Frieden und Ruhe gelebt hat: daß nun demjenigen solche Schmach und Gewalt widerfahren soll, ist wahrlich ganz unbillig. Der Ungehorsam aber und Vorwitz des Mägdleins wird hart genug gestraft; denn sie wird mit Gewalt überfallen und geschwächt.

\section[Zweiter Theil]{Zweiter Theil. Wie Sichem freundlich mit Dina redet; wie er seinen Vater bittet, ihm dieselbe zur Ehe zu nehmen; wie Hemor deswegen mit Jakob und seinen Söhnen redet; und wie sich Jakobs Söhne dabei verhalten und was sie von Hemor und Sichem verlangen}

\subsection{I}

\begin{quote}
\textbf{V. 3. Und sein Herz hing an ihr, und hatte die Dirne lieb, und redete freundlich mit ihr.}
\end{quote}

40. Sichem sieht, daß er gesündigt hat; denn er fühlt und versteht, daß das Mägdlein wider ihren Willen solche grobe Schmach leidet; er sieht, daß sie jämmerlich heult und klagt: jedoch dieweil er sie so heftig lieb hat, hängt sein Herz an ihr und redet er freundlich mit ihr, auf daß er sie dadurch reizen möge, daß sie ihn wiederum auch lieb gewinne.

41. Der junge leckere Sichem wird muthwillig und übel erzogen sein, ohne Ruthe und Zucht. Solche Jünglinge verderben dadurch, wenn die Eltern zu weich sind und ihnen zu viel übersehen, und wachsen auf in allerlei Sünden und bösen Begierden, denen sie ohne alle Furcht und Scheu nachhängen. Darum sagt Salomo in den Sprüchen am 29. Capitel V. 15.: „Ruthe und Strafe gibt Weisheit; aber ein Knabe, ihm selbst gelassen, schändet seine Mutter.“ Wenn nun dem Kinde seinen Willen läßt: Es ist gut, was unser Töchterlein thut: das ist ein gewisses Hürlein. Danach aber kommt es, daß der Vater die Schande und die Mutter die Schmach tragen muß. Das ist dann die Strafe, daß sie ihr Kind nicht wohl erzogen haben. Gott will haben, daß die Jugend regiert und mit guter Zucht gezwungen soll werden; denn dasselbe Alter ist sehr schwach und unerfahren, und denkt nur auf närrische, tückische und häßliche Dinge. Darum kann es sich selbst nicht regieren, kann auch nicht sehen, was ihm nützlich und gut sei. Darum hat aber Gott die Eltern und Schulmeister verordnet, die auf die Jugend sehen und sie in ihrem Leben und Sitten regieren sollen, daß sie thun, was ihnen in ihrem Amte gebührt, und das mit nichten überschreiten.

42. Darum hängt nun Sichem an dem Mägdlein mit jugendlicher Liebe. Es ist ihm aber noch nicht leid, was er gethan hat, sondern er meint, es sei ihm frei erlaubt, seines Gefallens zu sündigen, als dem Herrn im Lande, und habe ihn niemand darum zu strafen; wie man solche Sprüche hat, daß sie sagen: Sanctitas, pietas, fides privata bona sunt; reges quo libet eant; das ist: Heiligkeit, Gottseligkeit, Treue und Glaube sind sonderliche Güter oder Tugenden für die Unterthanen, die Könige und Fürsten mögen gehen, wohin sie nur wollen. Item: Es ist häurisch, daß ein Fürst sich nach den Gesetzen der Gottseligkeit und Gerechtigkeit richten soll. Herren sollen einen Vorzug haben, sie gehörten nicht unter die Gesetze: darum mögen sie thun, was sie nur gelüstet. Ja, Blitz und Donner! Solches geziemt noch wohl der Obrigkeit weniger als den Unterthanen, und ihr Muthwille und Freiheit soll mit größerem Ernste gestraft werden; denn sie sind diejenigen, die über Zucht und guter Ordnung halten und Anderen ein Exempel sein sollen aller Tugenden. Aber da schreien sie: Ich bin ein Fürst und Herr im Lande, warum sollte ich mir selbst nicht etwas nachsehen können? Ja, ich will dir wohl ein Anderes sagen, nämlich: Je größer du bist, je mehr du dich demüthigen sollst, daß du mit Tugend und guten Exempeln Andern vorgehen mögest, nach dem gemeinen Spruche: Regis ad exemplum totus componitur orbis; das ist: Nach dem Exempel des Königs richtet sich die ganze Welt. Die bösen Exempel der großen Herren und Fürsten sind die allerschädlichsten und erstrecken sich sehr weit; denn sie verderben damit viele Leute, und bringen sie dahin, daß sie auch muthwillig sündigen. Die Exempel der gemeinen Unterthanen thun so großen Schaden nicht und können leichter gebessert werden: die Sünden der Fürsten sind gar teuflisch, darum werden sie auch greulich gestraft; wie Sichem greulich ist gestraft worden, da er meinte, daß er als ein Fürst oder Herr alles thun möge, was ihn nur gelüstet. Es kostet ihm seinen Hals; er und seine Eltern, und dazu die ganze Stadt kommen jämmerlich um. Dieses ist nun eine Lehre von guten Sitten.

\quellseite{230}

43. Es steht aber im Text weiter: Er hat dem Mägdlein zum Herzen geredet, das ist: Er hat dem Mägdlein seine freundliche, süße und tröstliche Worte gegeben, die ein betrübtes Herz pflegen zu erfreuen. Es sind Buhlworte gewesen, da hat er große Versprechungen und etliche Geschenke lassen mit unterlaufen, auf daß er damit das Mägdlein, so betrübt und geschmächt war, trösten möchte. Aber das ist alles vergeblich; denn das Mägdlein bleibt betrübt und traurig.

\subsection{II}

\begin{quote}
\textbf{V. 4. 5. Und Sichem sprach zu seinem Vater Hemor: Nimm mir das Mägdlein zum Weibe. Und Jakob erfuhr, daß seine Tochter Dina geschändet war; und seine Söhne waren mit dem Vieh auf dem Felde, und Jakob schwieg, bis daß sie kamen.}
\end{quote}

44. Das hebräische Wort Jaldah heißt Töchterlein, Mägdlein; es heißt Kindlein, die noch jung sind. So spricht nun Sichem zu seinem Vater: „Nimm mir das kleine Mägdlein zum Weibe“, das noch fast ein Kind ist und mit den Kindern umgeht. Nun bittet er zu spät seinen Vater, daß er ihm dieses Mägdlein verschaffen wolle, daß er sie möge zum Weibe nehmen. Warum hast du nicht zuvor des Mägdleins Willen, dazu auch ihrer Eltern und deines Vaters Willen haben wollen? Es heißt: Sero sapiunt Phryges: Wenn der Schade geschehen ist, so wird man erst witzig.

45. Jedoch ist es ihm noch nicht leid, sondern er macht die Sünde noch größer, die sonst durch die Buße verschwindet und getilgt wird. Denn die Buße nimmt die Sünde weg; dagegen aber macht die Unbußfertigkeit die Sünde groß, breitet sie aus, daß sie immer schwerer wird. Er spricht nicht: Ich habe gesündigt; ich erkenne, daß ich daran sehr übel gethan habe, daß ich des sehr frommen Mannes Tochter geschwächt habe. Nein, man hört da von keinem Bekenntniß der Sünde, von keiner Buße; ja, er denkt, er habe recht gethan, oder dieses sei eine geringe Sünde, ein Mägdlein entführen und beschläfen; er verwickelt also seinen Vater Hemor auch mit in dieselbe Sünde, der seinen Sohn lieb hat und sein zart mit ihm umgeht, als der Herr und künftige Erbe sein soll; wie hernach, V. 19., gesagt wird von Sichem: „Er war herzlich gehalten über alle in seines Vaters Hause.“ Darum straft oder züchtigt er ihn nicht, sondern macht sich fremder Sünde theilhaftig, und willigt auch mit darein. Das ist das Aergerste.

46. Endlich aber kommt diese böse That auch vor den alten Vater Jakob und vor seine Söhne Simeon, Levi und die andern. Denn dieses Geschrei hat nicht lange können heimlich bleiben, darum daß diese Personen von beiden Seiten berühmt waren. Da wird nun ohne Zweifel das Haus Jakobs mit Wehklagen und Traurigkeit erfüllt worden sein und werden die Eltern und das ganze Gesinde jämmerlich geweint haben. Die beiden Weiber sind noch am Leben, und werden hoch betrübt, da sie diese Nachricht gehört haben, und vornehmlich hat es der Lea sehr wehe gethan. Die Knechte, Mägde und Hirten sind alle sehr ungeduldig darüber geworden. Die Andern haben es mit Heulen und Weinen bezeugt, daß es ihnen über die Maßen wehe gethan hat.

47. Jakob selbst ist vor großer Traurigkeit erschrocken, daß er still schweigt. Denn er findet nirgends Rath oder Trost, und er kann diese Gewalt, so ihm an seiner Tochter geschehen war von dem großmächtigen Herrn, nicht rächen, denn auch daran nicht, daß er es an ihm rächen wollte. Derhalben thut er nichts dazu: er redet mit dem Thäter nicht darum und läßt sich auch die That nicht gefallen, sondern schweigt still dazu und leidet es mit Geduld, und hofft Rath und Hülfe vom Herrn zu bekommen. Denn was sollte er thun, dieweil er ein Fremdling und Gast im Lande war, dazu auch an Volk und Gewalt viel geringer war denn der Thäter? Und Hemor wird ihn ohne Zweifel verachtet und viel geringer gehalten haben denn sich selbst. Und Sichem hat gedacht: Was ist denn das so ein groß Ding, daß ich ihm sein Töchterlein geschwächt habe? Will ich sie doch wohl zum Weibe nehmen. den sehr zornig, daß er eine Narrheit an Israel begangen, und Jakobs Tochter beschlafen hatte; denn so sollt es nicht sein.

48. Das Geschrei ist vor die Brüder der Dina gekommen gleich in derselben Stunde, da die That begangen und da es im Hause Jakobs auch ruchbar geworden ist. Es wird irgend der Knechte einer zu den Söhnen gelaufen sein und gesagt haben: Ach! wie hat sich in eures Vaters Hause so ein groß Unglück zugetragen: Sichem hat ihm seine Tochter Dina genommen und sie beschlafen etc. Darum sagt der Text: Da sie gehöret haben, was geschehen war, haben sie die Heerde verlassen und sind angekommen, und sind sehr betrübt und zornig gewesen. Und hier steht im Hebräischen das Wort, welches Mose droben im 6. Cap. 3. 6. auch gebraucht hat, da im Lateinischen steht: Et tactus dolore intrinsecus: „Und es bekümmerte ihn in seinem Herzen“; item im 3. Cap. 3. 16., da Gott zum Weibe spricht: In dolore paries filios: „Du sollst mit Schmerzen Kinder gebären“; denn er zeigt damit an, daß dem Herzen angst und bange ist, daß es keinen Rath und Gedanken weiß. Es heißt, bekümmert sein, angst und bange sein, daß du in deinem Herzen an nichts denken kannst. Auf solche Weise sind die Söhne Jakobs auch so gar bestürzt gewesen, daß sie nichts haben finden können, was sie denken, rathen oder thun sollten.

49. Zu solcher Bestürzung und Angst ist auch der Zorn gekommen, der so viel hitziger war, je weniger sie Rath oder Hülfe gesehen haben. Denn das war ihnen zum höchsten beschwerlich, daß ein unbeschnittener Mann ihrer Schwester diese Gewalt gethan hatte, der die Sünde nicht erkannt und gar kein Zeichen der Buße von sich gegeben hat. Der große Herr, Sichem, behält ihre Schwester daheim in seinem Hause, da er sie gefangen hält, und sie können dieselbe mit Gewalt von ihm nicht fordern; und dieweil sie keinen Richter haben, so wissen sie niemand, vor dem sie sich dieser Gewalt halben, so ihnen an ihrer Schwester geschehen war, beklagen mögen. Hemor aber kommt und will sich mit Jakob und seinen Söhnen vergleichen; aber dieser Schanden will sich nicht heilen lassen, und die Brüder werden sich auch nicht lindern oder versöhnen lassen, bis sie diese Gewalt an Sichem gerächt haben.

50. Operari stultitiam, eine Narrheit begehen, ist bei den Hebräern ein Sprüchwort, und sie verstehen dadurch eine große Schande, Schmach und Hohn, davon man singt und sagt. Gleich als wollte Einer sagen: Ei, wie ist diesem Manne in seinem Geschlechte so ein großer Hohn und Schande widerfahren! Darum wird im Texte dazu gesetzt: „Denn so sollt es nicht sein.“ Die Juden fehlen dieser Auslegung; denn sie verstehen die Bibel nicht. Das ist aber der rechte Verstand dieser Worte, wie man im Lateinischen sagt: Non decebat fieri hoc, das ist: Also sollte man es nicht machen, daß ist zu viel und hart. Ich halte aber dafür, daß die Hurerei bei Leibesstrafe sei verboten gewesen nicht allein im Hause Jakobs, sondern auch im ganzen Lande.

51. Darum wird mit diesen Worten angezeigt, daß neben dem Hausgesinde Jakobs auch die andern Nachbarn und Einwohner desselben Landes bestürzt und sehr zornig geworden seien, daß einem solchen heiligen Manne, der ihr Gast war, so ein großer Hohn widerfahren war. Also hat die Unbilligkeit der Sünde andere Leute bewogen, Mitleiden zu haben mit der Familie Jakobs; und sind auch die Fremden dadurch ihre Freunde geworden, so die Gerechtigkeit, Billigkeit und Zucht lieb gehabt haben. Wie wir droben, Cap. 26. V. 11., vom Könige Abimelech gehört haben, wie heftig er dem Ehebruche feind gewesen, und ein sonderlich Gebot hat lassen ausgehen, darin er verboten, daß niemand Rebekka berühren sollte. Denn wir sollen nicht denken, daß die Leute so gar grob und wild gewesen seien, daß sie gar keine Zucht, ehrliche Sitten und gute Gesetze hätten leiden können. Und es ist unmöglich, daß äußerlich in der Welt sollte Friede erhalten werden ohne Gerechtigkeit und mittelmäßige Zucht. Darum mißfällt diese That den Einwohnern desselben Landes, und ihrer aller Herz und Wille neigt sich zum Hause Jakobs, Mitleiden mit ihm zu haben. Die frommen sind darüber Sichem feind geworden und haben diese Sünde an ihm verabscheut. Denn sie halten alle dafür, man hätte es nicht also machen sollen, daß ist, sie halten es dafür, daß er damit wider die guten Sitten und Gesetze gesündigt habe.

\subsection{III}

\begin{quote}
\textbf{V. 8–10. Da redete Hemor mit ihnen, und sprach: Meines Sohns Sichem Herz lehnet sich
\quellseite{231}
nach eurer Tochter; lieber, gebet sie ihm zum Weibe. Befreundet euch mit uns; gebet uns eure Töchter, und nehmet ihr unsere Töchter, und wohnet bei uns. Das Land soll euch offen sein; wohnet und merket, und gewinnet drinnen.}
\end{quote}

52. Hemor redet sie an ohne alle Vorrede, und vergißt, daß er sie zuvor mit freundlichen Worten sollte einnehmen, daß sie ihm möchten gewogen werden, und erkennt die Schuld nicht, bekennt auch die Sünde nicht und bittet vielmehr um Verzeihung. Er ist allzu stolz und hoffärtig und darum wird auch die Sünde so viel schwerer. Und die Söhne Jakobs, so durch diese große Schmach und Schande heftig erzürnt waren, werden nicht versöhnt, sondern immer mehr erbittert. Denn er behält Dina im Hause, und sagt nur: „Gebet sie meinem Sohne zum Weibe“; so er doch zu allererst für seinen Sohn um Verzeihung sollte gebeten haben: Lieben Freunde und Brüder, verzeiht uns diese schwere Sünde, damit euch mein Sohn erzürnt hat. Ich will meinen Sohn Sichern züchtigen, will ihn ins Gefängniß werfen und in die Eisen schlagen, auf daß er lerne mit unseren Gästen hinfort ehrlich und friedsam umgehen. Der Dinge thut Hemor keines; sondern es scheint, als wolle er gleichsam die That noch vertheidigen, und macht sich also der ganzen Sünde theilhaftig und führt damit eine greuliche Strafe über sich.

53. Darum soll man lernen, daß die Eltern ihren Kindern nicht allzu weich sein, sondern sie zwingen sollen, wie sie von ihren Vätern und Voreltern in Zucht und Zwang sind gehalten worden. Denn dieser Hemor vergißt ganz und gar aller Ehrbarkeit und seiner schuldigen Pflicht, und als wäre die Sache wohl ausgerichtet, bittet er um das Mägdlein, daß sie dasselbe seinem Sohne wollen zum Weibe geben, und meint, damit solle die große Sünde, daß er es geraubt und beschlafen hatte, bezahlt sein, wenn das beschlafene Mägdlein seinem Sohne würde zum Weibe gegeben werden. „Befreundet euch mit uns“, spricht er. Im Hebräischen stehet hitchathenu, das lautet, als wollte er auf Lateinisch sagen: affinescite, das ist, werdet unsere Schwäger, lasset eure Söhne bei uns Bräutigame und eure Töchter Bräute werden. Das Wort chathan heißt eigentlich Bräutigam oder Schwager: daher kommt das Hithpaël, wie es in der hebräischen Grammatik genannt wird, als wenn man auf Lateinisch wollte sagen: sponsificate, befreundet und verschwägert euch mit uns; wir wollen uns also mit einander befreunden und Ein Volk werden. Wo bleibt aber die Erkenntniß der Sünde? Er denkt: Mein Sohn hat nichts Böses gethan. Ist das Sünde, wenn man eine Jungfrau zum Weibe begehrt? Ja, das ist vielmehr eine schöne löbliche Tugend.

54. Und er sagt weiter: „Und wohnet bei uns. Das Land soll euch offen sein; werbet und gewinnet drinnen“ 2c. Wo dieses sein Ernst ist, daß er sagt, so gibt er ihnen Versprechungen genug. Und ich will gerne glauben, daß es sein Ernst gewesen sei, daß er also geredet hat; denn er verstehet wohl, daß sein Sohn eine große Sünde begangen hat. Aber aus sonderlicher Hoffart will er es sich nicht lassen sein, will sich nicht demüthigen.

55. Also pflegen die stolzen Heuchler zu thun: sie können eher an allem ihrem Gut Schaden leiden, denn daß sie ihre Sünde sollten erkennen. Denn sie meinen, es könne ihnen keine größere Schande widerfahren, denn wo sie ihren Irrthum und Fehler sollten bekennen. Nun ist es aber in Wahrheit ehrlich, die Sünde erkennen; wo man sie aber vertheidigen, verleugnen und schmücken will, das ist ein erschrecklich und gottlos Ding. Dieser Hemor bietet diesen Leuten alle Gewalt im Lande, die Verheirathung seines Sohnes mit der Dina und alle Güter an; unterdes aber schweiget er von der Sünde stille und verhehlet dieselbige.

\begin{quote}
\textbf{V. 11. 12. Und Sichern sprach zu ihrem Vater und Brüdern: Lasset mich Gnade bei euch finden; was ihr mir saget, das will ich geben. Fordert nur getrost von mir Morgengabe und Geschenk, ich wills geben, wie ihr heißet; gebet mir nur die Dirne zum Weibe.}
\end{quote}

56. Sichern, der Sohn Hemor, erbietet sich auch, daß er ihnen die Morgengabe und andere Geschenke geben wolle, nur daß sie ihm die Tochter zum Weibe geben. Die Sünde aber bekennt er nicht, gedenkt derselben auch gar nicht. Wie Sirach sagt im 17. Cap. V. 26.: „Die Todten, als die nicht mehr sind, können nicht loben“ 2c. Nun hätten sie aber zu allererst sich mit Gott sollen versöhnen und darnach dann mit dem Vater und seinen Söhnen; sonst können die Herzen nicht fest vereinigt werden, weder in Bündnissen noch im Ehestande, es sei denn, daß man zuvor mit Gott versöhnet sei durch demüthige Erkenntniß der Sünde, und mit den Menschen auch versöhnet sei also, daß man sie um Verzeihung gebeten. Denn wo man die Sünde noch vertheidiget, oder ja zum wenigsten verhehlet, dadurch werden die Herzen nicht vereinigt, sondern immer mehr von einander gerissen und erbittert, und vor Gott wird eine Sünde auch immer schwerer durch Unbußfertigkeit, ob sie gleich sonst sehr leicht und gering wäre.

57. Derhalben ist es Hemor und seinem Sohne Sichern kein rechter Ernst, sich mit diesen Leuten zu versöhnen, sondern ist nur eine lautere Heuchelei und Verstellung; und dadurch wird Sichems Sünde immer größer. Denn Unbußfertigkeit ist der Tod und Teufel selbst, die auch eine geringe Sünde sehr groß und schwer macht. Dagegen aber werden große greuliche Sünden vergeben und weggenommen, wenn rechte wahrhaftige Buße dazu kommt. Wie jedermann Mitleiden hat, wo irgend ein Mord von Ohngefähr geschehen ist, und gern solche Sünde vergibt dem, der sie ohngefähr und nicht gern begangen hat, wenn er sie bekennt und sich solche Sünde läßt leid sein, wenn er sich beklagt, wie er durch des Teufels List betrogen und überwunden worden sei. Wo man aber die That noch entschuldigen und vertheidigen will, wie die Sentauern und Edel-leute, die da meinen, sie mögen nur frei ihres Gefallens toben und wüthen mit den armen Bauern, Gewalt an ihnen üben und sie nur todtschlagen: solcher grausamen Tyrannen sind beide Gott und Menschen feind, und sie finden nirgends Raum, da ihnen solche Tyrannei verziehen werden möchte. Also fährt hier der stolze und gottlose Mann Hemor mit seinem Sohne auch fort, der ein Verächter ist Gottes und der Menschen. Darum ist es mit den großen Worten und Versprechungen eitel erdichtet Ding, die er vielleicht nicht würde gehalten haben, wenn ihm die Jungfrau wäre zum Weibe gegeben worden.

\begin{quote}
\textbf{V. 13–17. Da antworteten Jakobs Söhne dem Sichern und seinem Vater Hemor betrüglich, darum, daß ihre Schwester Dina geschändet war; und sprachen zu ihnen: Wir können das nicht thun, daß wir unsere Schwester einem unbeschnittenen Manne geben; denn das wäre uns eine Schande. Doch dann wollen wir euch zu Willen sein, so ihr uns gleich werdet, und alles, was männlich unter euch ist, beschnitten werde. Dann wollen wir unsere Töchter euch geben, und eure Töchter uns nehmen, und bei euch wohnen, und Ein Volk sein. Wo ihr aber nicht willigen wollet euch zu beschneiden, so wollen wir unsere Tochter nehmen und davon ziehen.}
\end{quote}

58. Bisher haben wir gehört, wie sich Sichern und sein Vater Hemor fälschlich und nicht von Herzen gedemüthigt haben; welche Demuth die Liebe, so er zu der Jungfrau gehabt, zuwege gebracht hat, daß er versprochen, was sie nur fordern würden, dafür zu geben, daß sie ihm die Jungfrau zum Weibe vertrauen sollten. Aber die Liebe ist es, die solches verspricht, und ist eine Unmäßigkeit, dadurch er die Sünde nicht erkennt, und richten also der Sohn und Vater nichts aus, damit die Sünde möchte gebüßt oder Gnade erlangt werden. Derhalben wird die Sünde, so an sich greulich ist, durch die Unbußfertigkeit immer größer, da er dieselbe nicht bekennt und dafür nicht genug thut; denn er behält nichts desto weniger die Jungfrau, die er mit Gewalt wider Recht und Billigkeit genommen und beschlafen hatte, und die eines solchen großen Mannes Tochter ist, der ein Patriarch, ein Prophet und Priester Gottes war. Diese Umstände machen die Sünde sehr groß, und wenn die Sache vor das weltliche Gericht gebracht würde, alsdann sollte man erst recht sehen, wie greulich diese That wäre. Wir wollen aber nach der Schrift davon handeln, wollen nicht viel verprorene Fragen davon machen, damit heutiges Tages die Juristen alle Rechte verwirren, daß man sich nicht recht darnach richten kann. Wo eine Jungfrau mit Gewalt weggeführt wird, das ist an sich in allen Rechten, beide göttlichen und weltlichen, eine solche Sünde, damit man eine Leibesstrafe verdient; wie Troja darum zerstört worden, daß Helena weggeführt war. Und zu allen Zeiten ist diese Sünde greulich geachtet worden.

\subsection{IV}

59. Es fällt aber nun hier eine wichtigere und schwerere Frage ein, nämlich: Ob die Söhne Jakobs daran recht gethan haben, daß sie so betrüglich geantwortet haben, ohne Vorwissen ihres Vaters Jakob, der dazu schweigt. Von Sichem und Hemor ist kein Zweifel, daß sie nicht sollten ungerecht sein, dieweil sie ihre Sünde nicht erkennen, nicht Buße thun, auch nicht sagen, daß sie gesündiget haben, zudem auch für die Sünde keine Genugthuung leisten. Derhalben wird nun eine Sünde gegen die andere gehalten, Gewalt gegen Gewalt, Schmach gegen Schmach, dazu auch eine Ursache gegen die andere, und kein Theil thut, was billig und recht ist.

\quellseite{232}

60. Den Mord der Sichemiten betreffend, kommt die Meinung der Juristen und der heiligen Schrift überein: „Die Rache ist mein; ich will vergelten“, spricht Gottes Wort Hebr. 10, 30. Item: Was recht ist, das sollst du auch auf die rechte Weise ausführen. Denn niemand soll sich die Gewalt anmaßen, daß er selbst Richter sein oder sich rächen wollte, wo ihm von Andern Unrecht geschehen ist. Die Deutschen sagen: Wer wieder schlägt, hat Unrecht. Es haben sich aber doch ihrer Viele sehr bemüht, daß sie diese Frage haben erklären wollen. Und Lyra und Burgensis streiten hart mit einander, welcher Theil möge gerecht sein. Hernach aber wird Jakob selbst diese That seiner Söhne strafen, da er sagt: „Ihr habt mir Unglück zugerichtet, daß ich stinke vor den Einwohnern dieses Landes, den Cananitern und Pherefitern“ etc. Da läßt es sich ansehen, daß er den Mord strafe, den seine Söhne begangen hatten. Und im 49. Capitel, da Simeon und Levi hat segnen wollen, vergibt er dieser That nicht, und entschuldigt sie auch nicht, sondern sagt V. 5. 6. also: „Die Brüder Simeon und Levi; ihre Schwerter sind mörderische Waffen. Meine Seele komme nicht in ihren Rath; denn in ihrem Zorn haben sie den Mann erwürget, und in ihrem Muthwillen haben sie den Ochsen verderbet.“ Da hat er diesen Mord, der zu Sichem begangen war, gemeint, wie es alle Lehrer auslegen, und darum verflucht und verdammt er sie, ehe er stirbt, so sie doch eine sehr gerechte Ursache haben zu solchem Morde.

61. Wiederum aber, da er auf den Segen Josephs kommt, spricht er 1 Mos. 48, 22.: „Ich habe dir ein Stück Landes gegeben außer deinen Brüdern, das ich mit meinem Schwert und Bogen aus der Hand der Amoriter genommen habe.“ Er sagt, er habe es mit dem Schwerte erworben, und gibt es Joseph als einen Segen, als ein wohl erworben Gut; wie dessen auch Joh. 4, 5. gedacht wird, da also steht: „Bei dem Dorflein, das Jakob seinem Sohne Joseph gab.“ Jetzt aber an diesem Ort sagt Mose, daß die Söhne Jakobs betrüglich geredet haben, und hernach empfangen sie um dieser That willen den Fluch, und Joseph empfängt gleichwohl den Besitz desselben Landes, als das mit gutem Rechte erworben sein soll. Wie wollen wir dieses mit einander vergleichen, daß es sich zusammenreimen möge?

62. Wir wollen die Antwort aus dem Text und aus den Worten Mosis nehmen, so viel es möglich ist, diese Frage zu erklären. Mose sagt: Sie haben betrüglich geantwortet auf die Bitte Hemor, die doch billig genug war, wie es schien. Das ist so viel gesagt: Sie haben sich verstellt, sie haben gelogen mit Worten und auch mit der That; und ist solches nicht eine Scherz- oder Rothlüge, sondern eine schädliche Lüge gewesen, die den Sichemiten zu Schaden und Verderben geraten ist. Das ist wahrlich eine Todsünde und kann nicht anders verstanden werden. Darum sagt Mose nicht vergeblich: „Sie haben betrüglich geantwortet“, das ist, sie haben das nicht von Herzen gemeint, das sie mit Worten geredet haben, sondern sie haben das Gegenteil gemeint. Sie haben sich gestellt, als wollten sie sich zu ihnen halten und sie unwerth und unwerth lassen, und hatten doch im Sinne, sie zu erwürgen; wie das die That an sich wird ausweisen.

63. Und der Text sagt weiter: Und sie redeten darum, daß ihre Schwester Dina geschändet war, das ist, sie haben darüber disputirt und gezankt, haben gesagt: Du hast die Tochter eines Priesters und heiligen Propheten entführt und beschlafen, welchen Gott in diesem Lande zu einem Lehrer und Kirchenvater gesetzt hat. Sie haben diese Gewalt, so an ihrer Schwester begangen war, groß gemacht nach allen Umständen, wie sie nur gekonnt haben. Denn es ist eine andere und viel eine schwerere Sünde, eines Königes Tochter zu schwächen, denn irgend eines Bauern Tochter. Und Simeon und Levi waren Söhne eines Propheten, so die Verheißung hatten desselben Landes; darum hatten sie gewisse Hoffnung gefaßt aus den väterlichen Predigten, daß sie würden Herren und Priester werden in dem verheißenen Lande. Denn Jakob wird ohne Zweifel mit großem Fleiß diese Verheißung gelehrt und sie bei den Nachkommen ausgebreitet haben, und hat zugleich auch mit den Verheißungen die Gebote Gottes gelehrt, wie die Leute ihr Leben und Sitten gottselig anrichten sollten. Solche Lehre und Hoffnung haben die Söhne aus des Vaters Lehre und Predigt geschöpft und haben daraus also geschlossen: Wir sind Herren in diesem Lande; ob wir wohl jetzt Fremdlinge sind, so gebührt doch uns nach göttlichem Recht die Verheißung des Landes. Sollen wir denn darin von euch solche große Gewalt und Unrecht leiden, davon jedermann zu sagen weiß, die ihr doch jetzt ein fremd Land inne habt? Dieses werden sie mit großem Trotz und Stolz aus der Verheißung des Vaters gerühmt haben. Wir sind es, werden sie gesagt haben. Wie wir hernach hören werden, wie ganz trotzig sie ihrem Vater antworten, da er sie straft, daß sie diesen Mord begangen hatten: „Sollten sie denn mit unserer Schwester, als mit einer Hure, handeln?“ Und solche Hoffart ist stets bei dem jüdischen Volke gewesen, daß sie sagen: Wir sind Gottes Volk; wir sind das königliche Priesterthum etc.

64. So haben nun diese Brüder mit großem Stolze disputirt von ihrer Schwester, die beschlafen und geschändet war, die eines solchen großen Mannes Tochter war, die des Geschlechtes, der Schönheit und auch der Gottseligkeit halben alle anderen Jungfrauen in selbigem Lande übertroffen hat. Denn so groß und schwer haben sie diese Sünde gemacht, die vielfältig unterschieden ist von einer solchen gemeinen That dadurch, daß sie eine sonderliche Person betrifft und ein Unterschied ist derselbigen und anderer Personen. Denn diese Jungfrau ist keines gemeinen Bürgers oder Bauern Tochter, sondern des höchsten Propheten im ganzen Lande; und ist nicht eine einzelne Person insonderheit geschmäht, sondern das öffentliche Amt ist dadurch verachtet worden. Wer nun einen solchen großen Patriarchen seine Tochter entführen darf und ihm dieselbe mit Gewalt zu Schanden machen, und sie dazu wider seinen Willen noch behalten, was sollte der Andern nicht thun dürfen?

Und was hat man nun Anderes von ihnen zu erwarten, denn daß sie bald hernach uns auch das Leben und all unser Hab und Gut mit Gewalt nehmen werden? Dieweil wir die Verheißung haben, daß wir in diesem Lande sollen Herren sein, wohlan, so wollen wir wieder ein Thälten thun. Warum greifen sie aber die Sichemiten so hinterlistig an und nicht öffentlich? Denn wiewohl es scheint, als ob sie mögen entschuldigt werden der Ursache halben, daß man den Feind wohl betrügen mag, nach dem gemeinen Spruche, da man sagt: Dolus, an virtus, quis in hoste requirat? Wer will das an dem Feinde suchen, ob er Hinterlist oder aber Stärke gebraucht habe? so ist es aber doch nicht recht, daß sie in dem freundlichen Gespräche und Unterredung, so sie mit den Sichemiten gehalten, solche Hinterlist gebrauchen, da sie sich damals gegen sie noch nicht erklärt hatten, daß sie ihre Feinde wären. Darum halte ich sie nicht für unschuldig; wie sie auch der Vater selbst verdammt wird. Und ist doch gleichwohl groß Wunder, daß Jakob das Land behält. Denn er hat dasselbe eingenommen, da die Bürger erschlagen waren, und hat daselbst seine Weide gehabt für die Heerde, da er bei seinem Vater Isaak zu Hebron gewohnt hat. Und er wird darnach selbst sagen, daß er es nicht um sein Geld gekauft, sondern mit seinem Schwert und Bogen erworben habe. Dasselbe gibt aber nun den Juristen große Ursache zu vielen Fragen und daß sie weitläufig davon disputiren.

65. Mose aber gibt selbst zu verstehen, daß die Söhne Jakobs nicht recht gethan haben, daß sie mit Hemor und Sichem betrüglich geredet und sie betrogen haben, daß den Sichemiten zu Schaden und Verderben geraten ist. Denn wiewohl es scheint, daß die Sünde der Strafe fast gleich sei, so ist doch ein Unterschied zwischen den Rechten und der Vollstreckung, wie man das Recht ausführen soll, item, zwischen der Sünde, und der Strafe, so einem um der Sünde willen auferlegt wird. Denn es steht geschrieben: „Die Rache ist mein, ich will vergelten“, 5. Mos. 32, 35. und Röm. 12, 19.; item: Was Recht ist, das sollst du auf rechte Weise ausführen. Es ist nicht unrecht, daß Sichem gestraft wird; wie auch die Trojaner nicht mit Unrecht verheert und zerstöret worden sind um der Entführung der Helena willen; aber das weiß man noch nicht, ob sie auch in rechter und gebührlicher Ordnung ihr Recht ausführen. Denn es mangelt am Richter, und sie vergelten einfach Böses mit Bösem und rächen eine Gewalt mit der andern, wie man sagt, Stück um Stück.

\quellseite{233}

66. Erst aber reden sie aufrichtig und deutlich genug, da sie sagen: „Wir können das nicht thun, daß wir unsere Schwester einem unbeschnittenen Mann geben.“ Daraus wohl zu sehen ist, daß sie ihre Töchter nur denen haben pflegen ehelich zu vertrauen, die sich beschneiden ließen. Aber mit dieser Einfalt stellen sie ihnen Netze, und gebrauchen also dieses Vormandes, auf daß sie Gelegenheit haben mögen, sie anzugreifen und zu überfallen. Darum sagen sie weiter: „Doch dann wollen wir euch zu Willen sein, so ihr uns gleich werdet“ etc. Das ist der Anfang der Hinterlist, damit sie sündigen, und das ist es auch, daß der Vater an ihnen straft und verdammt. Das wäre recht und ehrlich gewesen, wenn sie also gesagt hätten: Gebt uns unsere Schwester wieder heraus, oder wir wollen sie euch mit Gewalt nehmen; denn wir dürfen sie keinem unbeschnittenen Mann geben. Dieses ist aber die Hinterlist, da sie sagen: „Wo alles, was männlich unter euch ist, beschnitten wird, so wollen wir unsere Töchter euch geben“ etc.

67. Und ist es ein Wunder, daß die Söhne Jakobs so keck gewesen sind, daß sie solches haben fordern dürfen, oder daß sie solche greuliche Rache in ihrem Herzen haben mögen vornehmen. Denn, siehe ihr Alter an: Jakob ist hundert und acht Jahre alt; Ruben ist bei neunzehn Jahren alt und der ist der älteste; Simeon ist ein Jahr jünger, und so immer Einer nach dem Andern: es sind junge Lecker. Die jungen Lecker nehmen sich vor, eine solche große Gewalt zu üben wider eine ganze Stadt, die doch des Krieges und anderer Dinge gar unerfahren gewesen sind. Das ist mir eine wunderliche Kühnheit. Wiewohl, zu derselbigen Zeit ist die Natur der Menschen noch stärker gewesen, da ein Jüngling von vierzehn Jahren zum Ehestande ist tauglich gewesen. Zu unseren Zeiten ist noch ein Knabe von achtzehn Jahren nirgend zu tauglich, ich will sagen schweigen, daß er die Mühe und Arbeit im Kriege sollte ertragen können.

68. Warum sind sie denn so stolz? Antwort: Die Gewalt, so an ihrer Schwester war geübt worden, ist sehr groß gewesen; und dazu ist gekommen, daß sie sich darauf hoch verlassen haben, daß sie wußten, sie hätten einen solchen Vater, dem die Verheißung gegeben war, daß er das Land erben sollte. Und sie haben diese greuliche That nicht allein unternommen, sondern sie haben die Knechte und Hirten, und vielleicht auch wohl etliche von den Nachbarn, die sie dazu gefordert haben, zu Hülfe gehabt. Und ist wohl glaublich, daß Jakob unter seinem Hausgesinde hundert Mann wird gehabt haben; denn zu so viel Heerden muß man viel Hirten haben: dieselbigen sind ihre Freunde und Gesellen gewesen. Nun aber können hundert gerüstete und beherzte Männer (wie von diesen gesagt wird, daß sie beherzt in die Stadt sind gegangen und die Leute plötzlich überfallen haben) leichtlich andere sechshundert Mann schlagen, vornehmlich, wenn Gott um der Sünden willen solchem Spiele zuführet, welche Sünde er auf die Weise strafen will.

69. Da sie sich aber nun dürften lassen, daß sie nicht stark genug seien, die Sichemiter zu unterdrücken, so gebrauchen sie Hinterlist, und fordern von den Sichemitem, daß sie alles, was männlich unter ihnen ist, beschneiden sollen, daß sie also geschwächt werden, auf daß sie nicht gerüstet sein können, ihnen Widerstand zu thun. Das heißt hinterlistig und betrüglich gehandelt. Denn wiewohl der Schein, den sie vorwenden, dem Ansehen nach ehrlich genug ist, so ist er doch eitel und unnöthig. Denn sie haben doch hernachmals ihre Töchter unbeschnittenen Männern geben müssen.

\section[Dritter Theil]{Dritter Theil. Wie Hemor und Sichem in das Verlangen der Söhne Jakob willigen und ihr Volk auch dazu überreden; wie sie allgemein jämmerlich darüber erwürget werden und ihre Stadt geplündert wird; wie Jakob deshalb mit seinen Söhnen redet und was er für Antwort erhält}

\subsection{I}

\begin{quote}
\textbf{V. 18. 19. Die Rede gefiel Hemor und seinem Sohn wohl. Und der Jüngling verzog nicht solches zu thun; denn er hatte Lust zu der Tochter Jakobs. Und er war herrlich gehalten über alle in seines Vaters Hause.}
\end{quote}

70. Dieweil der Jüngling so große Lust und Liebe hat zu der Jungfrau, so nimmt er die Bedingung mit großer Begierde an, und wird also mit List hintergangen und betrogen, daß er die Beschneidung annimmt; ja, der Vater und Sohn verwickeln das ganze Volk auch mit in solche Sünde und Gefahr. Denn wo damals Sichem und Hemor weise gewesen wären, so hätten sie sich geweigert, solch hartes und schweres Ding anzunehmen, auch um der Ursache willen, daß es des gemeinen Wesens Nutzen nicht sein würde, daß die Religion und Sitten ihrer Vorfahren und ihres Vaterlandes dermaßen sollten geändert werden. Sie sollten gesagt haben: Wir wollen euch lieber eure Schwester herausgeben, die Sünde bekennen und um Verzeihung bitten. Also hätten sie sich und das ganze Volk nicht in solche große Gefahr gebracht.

71. Aber also pflegt es zu gehen, wenn irgendwo eine Stadt oder Land um der Sünde willen soll gestraft werden, daß alsdann zuerst beide die Regenten und Unterthanen zu Narren werden, daß sie insgemein sündigen und darüber auch sämtlich gestraft werden; wie das die Historien bei allen Völkern bezeugen. Also, da das Königreich Israel und Juda hat sollen zerstöret werden durch den König zu Babel, wurden zu Jerusalem und an andern Orten mehr die Propheten erwürgt, es ward unschuldig Blut vergossen, und das ganze Land voll Geizes, Mordes, Sünde und allerlei Mißbrauch. Item, da die letzte Zerstörung Jerusalems und des ganzen jüdischen Volkes vorhanden war, da mußte erst der Sohn Gottes gekreuzigt werden. Also auch zu unsern Zeiten, da die türkischen Heere Deutschland verderben sollen, muß zuvor Alles mit Wucher, mit Untreue, Bosheit, Verrätherei und Verfolgung erfüllt werden, daß also der Amoriter Sünden erfüllt werden, wie die Schrift, 1 Mos. 15, 16., sagt. In Summa: Wir müssen es also machen, daß unsere Sünden reif werden.

72. Also in diesem Exempel auch: den Bürgern drohte Strafe wegen mancherlei Schandthaten; darum wird das ganze Volk auch mit dem Herrn in gleiche Thorheit, Sünde und greuliche Strafe verwickelt. Wie der heidnische Poet sagt: Quicquid delirant reges, plectuntur Achivi: Daß die Fürsten Narren sind, das muß das Land tragen. Also ist das Volk Israel auch geschlagen worden um Davids Sünde willen, da er das Volk zählt. Und die Thorheit der Fürsten bringt auch zuwege, daß wir müssen gestraft werden, wie wir vorlängst verdient haben, daß die Fürsten toll und thöricht sind und den Türken in das Land kommen lassen. Und daß beide Theile gestraft werden, Unterthanen und Obrigkeit, so müssen sie irgend mit einer endlichen Thorheit die Sünde häufen und derselben immer mehr machen; wie in dieser Historie Sichem der Erste ist, der den närrischen Rath annimmt und anfängt; denn er selbst läßt sich beschneiden, und überredet das ganze Volk, daß sie dasselbe thun.

\begin{quote}
\textbf{V. 20—22. Da kamen sie nun, Hemor und sein Sohn Sichem, unter der Stadt Thor, und redeten mit den Bürgern der Stadt, und sprachen: Diese Leute sind friedsam bei uns, und wollen im Lande wohnen und werben, so ist nun das Land weit genug für sie; wir wollen uns ihre Töchter zu Weibern nehmen, und ihnen unsere Töchter geben. Aber dann wollen sie uns zu Willen sein, daß sie bei uns wohnen, und Ein Volk mit uns werden, wo wir alles, was männlich unter uns ist, beschneiden, gleichwie sie beschnitten sind.}
\end{quote}

73. Ei, wie ist mir das so ein schöner Redner! Er sagt aber davon nichts, daß er die Jungfrau wollte wieder herausgeben, daß er sich mit dem Vater versöhnen und ihm seine Schmerzen lindern wollte: er beharrt noch immer in der Verachtung Gottes und der Menschen; darum ist auch der Zorn über sie gekommen. Es wird aber das närrische Volk auch zugleich in den Irrthum geführt, daß sie mit einwilligen und den Raub und die Schändung Dinas billigen. Das bringt ein toller und närrischer Fürst zuwege um der Sünden willen des Volks, die mit der Thorheit der Könige gestraft werden; wie das alle Historien und alle Schriften der Poeten ausweisen. Denn dieses ist also der Lauf der Welt vom Anfang bis zum Ende.

74. Also heutiges Tages auch, da so viel Sünden unter dem Volk im Schwange gehen und ungestraft bleiben, als, Sicherheit und Verachtung Gottes und der Menschen: so wird es endlich kommen, daß durch irgend eines Fürsten Thorheit Strafe und Unglück über uns kommen wird, darunter wir verderben müssen. Dieweil es das Volk so arg macht, wird ein Fürst eine Thorheit thun, wird ein Unglück anrichten, darüber werden wir alle mitgehen müssen. Und dieweil wir ja nicht hören wollen, uns auch durch Gottes Wort und Drohung nicht bessern wollen, darum wird endlich die Strafe kommen müssen.

\quellseite{234}

75. Es ist eine greuliche Sünde gewesen, daß Hemor und Sichem die Sünde verschwiegen und nicht erkennen, ja, den frommen heiligen Mann Jakob mit großer Sicherheit verachten. Was Jakob? Was Jakobs Tochter? Und zwar, wer vor solcher großen Sünde nicht erschrickt, und sich weder Gottes Gericht noch das Urtheil der Menschen bewegen läßt, der sollte wahrlich wohl dem Vater auch nach dem Leben stehen, wie er zuvor der Tochter heimlich nachgegangen und sie um ihre Ehre gebracht hat.

76. Auf solche Sicherheit folgt Blindheit und Thorheit, daß sie der heimlichen List und Gefahr, so vorhanden ist, nicht achten, noch dieselbe verstehen. Sichem ist gar unsinnig vor großer Lust und Liebe, die er zu der Jungfrau hat; diesem sieht der Vater und das ganze Volk durch die Finger, und ist niemand, der daran dächte, wie die Sünde möchte gebüßt werden. Sie sind toll und thöricht, es thut keiner die Augen oder Ohren auf, daß doch Einer gesagt hätte: Ei, lasset uns gemach thun, lasset uns ihnen die Jungfrau wiedergeben, lasset uns den Zorn des Vaters und der Brüder lindern. Wie Etliche bei den Trojanern waren, die den Rath haben, man sollte Helena wieder herausgeben und des Vaterlandes Wohlfahrt verschonen, welches ja mit nichten um eines Menschen Lust willen in Gefahr sollte geführt werden. Die Andern aber haben solches widerrathen und haben ihrer Lust und Unsinnigkeit nachgegeben; darum ist das Verderben und die Zerstörung darauf gefolgt. Wie in diesem Exempel das Volk auch ihrem Herrn folgt, und darum müssen sie auch mit einander die Strafe tragen und über einen Haufen umkommen; und geschieht ihnen eben recht.

77. Das ist zu viel. Und Lyra und Burgensis disputiren darüber: Ob auch Sichem und Hemor hinterlistig gehandelt haben? Denn es scheint also, daß sie betrüglich handeln, oder daß sie ja zum wenigsten ihre Bürger betrügen. Denn das war nicht mit im Vertrage, daß alles, was die Söhne Jakobs hatten, sie auch mit ihnen haben sollten, daß es ihnen mit einander gemein wäre. Wo nun dieses ihr Ernst gewesen ist, daß sie es gesagt haben, so ist es ein offenbarer Betrug und Bosheit, damit sie nicht allein die Sünde nicht erkennen, daß sie die Tochter Jakobs mit Gewalt genommen und geschändet hatten, sondern gedenken auch noch, Jakob zu überfallen und ihm sein Hab und Gut zu nehmen. Also wird wahrhaftig die Sünde gemehret und immer größer werden.

78. Wo sie aber nur also mit solchen zierlichen Worten Scherz treiben und nicht von Herzen reden, so haben sie das Volk betrogen wollen, auf daß sie dasselbe so viel leichter überreden möchten mit der Hoffnung, daß sie würden das Vieh und die Güter Jakobs bekommen. Jedoch scheint es, daß dieses darauf gerichtet sei, daß sie an den Nachbarn Tyrannei üben und sie unterdrücken wollen. Denn wozu dient es, daß sie sagen: „Ihr Vieh, und alles, was sie haben, wird unser sein“, wenn sie allein das gesucht haben, daß sie mit ihnen wollten Gemeinschaft halten und Ein Volk mit einander sein, und nicht gedacht hätten, ihre Güter einzunehmen? Darum handeln sie betrüglich, entweder daß sie die Söhne Jakobs betrügen, oder aber ihren Bürgern Lügen vortragen.

79. Und vielleicht ist dieses die Rechnung und Argwohn der Söhne Jakobs gewesen, daß sie gedacht haben: Wo sie die Tochter, die noch nicht mannbar war, haben rauben und schänden dürfen, was werden sie uns oder dem Vater nicht thun? Denn also schließt die Vernunft: Wer das Ernste thut, der wird das Kleine auch wohl thun. Es ist ja ein geringer Ding, Einen seine Güter nehmen und ihn dienstbar machen, denn irgend eine Jungfrau mit Gewalt wegführen. Und scheint es, daß sie auf solche Meinung droben etwas hart mit ihnen geredet und disputirt haben. Aber dieses alles verschweigen die Esel, Hemor und Sichem (denn Hemor heißt ein Esel), und predigen dem Volk nur und versprechen ihnen, daß sie Herren werden sollen nicht allein dieser Jungfrau, sondern auch ihrer Brüder und aller ihrer Güter.

80. Also wird nun die Sünde immer größer und zieht bald eine andere Sünde nach sich, wenn sie nicht durch die Buße getilgt wird. Also geht es mit einem Hurer oder Ehebrecher: wenn er in der Brunft der bösen Lust und Unzucht bleibt, und daß ihm der begangene, vor dem er sich besorgt, daß er ihn hindern oder verrathen möchte, so wird er die Hände nicht wohl stille halten können, sondern wird sich unterstehen ihn zu erwürgen, auf daß er seine Lust frei büßen möge. Und man hat einen herrlichen Spruch Gregors, der lautet also: Ubicunque serpens caput inseruerit, eodem et caudam et totum corpus facile inseret; das ist: Wenn der Teufel den Kopf hat in das Loch gebracht, so wischt die Schlange bald hernach mit dem Schwanze.

\begin{quote}
\textbf{V. 24. Und sie gehorchten dem Hemor, und Sichem, seinem Sohn, alle, die zu seiner Stadt Thor aus und eingingen, und beschnitten alles, was männlich war, das zu seiner Stadt aus und einging.}
\end{quote}

81. Das Volk muß hier die Narrheit des Königs tragen, wie man im Sprüchwort sagt. Denn sie sollten lieber also gesagt haben: Das Unrecht sei zu schrecklich; man sollte lieber die Jungfrau ihrem Vater wieder geben und die Religion und alte Gewohnheit der Väter behalten; man könnte wohl eine andere bessere Weise vornehmen, diese Sünde zu büßen. Wir wollen ihre Güter nicht begehren; denn es wird nimmermehr rechte Einigkeit unter uns sein, wenn wir heimlich nach ihren Gütern stehen wollen. Aber sie sind alle taub, stumm und blind, und fallen dem Teufel in sein Netz und dem Tode in seine Stricke; denn Alle, die zur Stadt Thor aus und eingingen, lassen sich beschneiden, das ist, Alle, die in der Stadt wohnen. Denn es ist bei den Hebräern eine sonderliche Art und Weise zu reden, wenn sie sagen, zur Stadt Thor aus und eingehen, welches so viel gesagt ist, als in der Stadt ein Bürger sein, der auf seinen Acker aus und eingeht.

82. Und also gehorchen die armen Bürger dem närrischen Herrn, der das ganze Volk mit in seine Unbußfertigkeit verwickelt, daß sie an dieser Sünde alle schuldig werden. Denn der Thäter und der in die That willigt, werden zugleich gestraft; wie St. Paulus sagt Röm. 1,32.: „Die Gottes Gerechtigkeit wissen (daß, die solches thun, des Todes würdig sind), thun sie es nicht allein, sondern haben auch Gefallen an denen, die es thun“: Stehler und Hehler, Thäter und Verräther gehören unter Eine Strafe. Die Könige treiben Thorheit und sündigen wider Gott und Menschen, und bekümmern sich um die Sünde nicht, sind gar sicher, als hätten sie wohl gethan, meinen nicht, daß sie betrogen seien, sondern daß sie vielmehr die Fremdlinge und Gäste überwältigt haben; darum folgt nun eine greuliche Strafe.

\subsection{II}

\begin{quote}
\textbf{V. 25. 26. Und am dritten Tage, da sie es schmerzte, nahmen die zween Söhne Jakobs, Simeon und Levi, der Dina Brüder, ein jeglicher sein Schwert, und gingen in die Stadt thürstig, und erwürgten alles, was männlich war. Und erwürgeten auch Hemor, und seinen Sohn Sichem mit der Schärfe des Schwerts, und nahmen ihre Schwester Dina aus dem Hause Sichem, und gingen davon.}
\end{quote}

83. Es ist ein groß Wunder, woher diese Jünglinge solchen großen Muth genommen haben, oder wo sie mögen gelernt haben, Schwert und Bogen zu gebrauchen wider den Feind, so sie doch Schafhirten und Bauern gewesen sind. Von dem dritten Tage machen die Juden viel unnütz Geschwätz; aber es lehrt es die Erfahrung in allerlei Krankheiten, daß am dritten Tage der Schmerz am allerheftigsten ist. Also ist es auch mit den Weibern, die fühlen nach der Geburt und nach der Arbeit, so die Natur in der Geburt verrichtet hat, am dritten Tage die größte Schwachheit und Schmerzen.

84. Also geht es auch zu, wenn die Natur des Menschen arbeitet und sich bemüht, daß der Stein dem Menschen abgehen soll, so fühlt man die Schwachheit über der Arbeit nicht; wenn aber die Natur den Stein abgetrieben hat, so sind alle Glieder und der ganze Leib matt und müde. Ob nun dasselbe in der Beschneidung auch also zugehe, darum mag man die Juden fragen. Dieweil aber die Schrift dasselbe auch bezeugt, so dürfen wir daran nicht zweifeln, die Erfahrung werde es auch also mitbringen.

\quellseite{235}

85. Diese Gelegenheit nehmen nun Simeon und Levi wahr: da die Sichemiten zu Bette gelegen, und ihrer keiner zum Streit gerüstet gewesen, sondern allesamt der Pflege und Ruhe bedurft haben, da greifen sie die armen, wehrlosen und schwachen Leute an. Und die zwei Jünglinge, Simeon und Levi, sind nicht allein gewesen, sondern haben die Hirten, die Knechte und auch etliche Nachbarn bei sich gehabt. Und Mose sagt, sie seien muthig und trotzig in die Stadt gegangen, nicht allein um solcher Ursache und Gelegenheit willen, die wehrlosen Leute zu überfallen und zu erwürgen, sondern auch dieweil sie sich auf ihre Ehre und Herrlichkeit verlassen haben, daß sie eines Patriarchen Söhne waren und Herren über die ganze Welt. Und sie haben nicht allein gestritten, daß sie ihre Schwester wieder heraus haben wollten, sondern sie haben diese Leute als Feinde der Verheißung und des Volkes Gottes angegriffen.

86. Und haben alles, was männlich war, erwürgt. Es ist nicht allein Hemor, der Herr, oder sein Sohn erwürgt worden, sondern alle Mannspersonen, die kleinen unmündigen Kindlein, die noch in der Wiege gelegen und geschrieen haben, sind auch mit erwürgt worden. Das ist allzu greulich und tyrannisch. Aber also geht es unter den Menschen zu, wenn ein Feind gegen den andern tobt und wüthet. Wie wir an des Türken Tyrannei auch sehen, der beide Mann und Weib, die er gefangen hat, jämmerlich erwürgt oder wegführt. Und ist ein heimlich Gericht und Zorn Gottes, daß er verhängt, daß diese Bürger zu Sichem alle erwürgt werden. Vielleicht haben die Söhne Jakobs solches von ihren Vorfahren gelernt, die alle Bürger in Städten, welche sie erobert haben, dermaßen pflegten zu erwürgen. Und 4 Mos. 31, 17. 18. wird Josua geboten, daß er alle Männer und Weiber der Midianiter erwürgen und allein der Jungfrauen verschonen solle. Desselgleichen wird dem Saul auch geboten, daß er das ganze Volk der Amalekiter, Großes und Kleines, Jung und Alt, erwürgen und ausrotten solle. Und man sagt, daß solches noch heutiges Tages der Perser Gewohnheit sei.

87. Also werden nun die armen Bürger um der Thorheit und Sünde willen ihres Königs jämmerlich gestraft: Fürsten haben genarrt, populus (das Volk) muß leiden. Es hatte aber gleichwohl das Volk diese Strafe verdient, darum daß sie sich fremder Sünde hatten theilhaftig gemacht, und hatten mit ihrer eigenen Sicherheit und Bosheit diese Strafe und Verderben über sich geführt.

88. Wie wir heutiges Tages unsere Bürger lehren und strafen, daß sie doch Gottes Wort mit dankbarem Herzen hören und annehmen wollten, und aufhören, mit ihrem Betrug, Räuberei und Wucher den Zorn Gottes zu erregen. Aber es ist weder Maß noch Aufhören mit dem Geiz, Räuberei und Betrug, so in Kaufen und Verkaufen, und im ganzen Leben der Menschen durch und durch getrieben wird. Darum richten wir mit unserem Lehren und Strafen nichts aus, sondern die Sünden werden immer gehäuft und nehmen überhand mit Gewalt ohne alle Strafe. Wenn nun irgend ein Fürst kommt und besetzt die Stadt mit Kriegsvolk, und dasselbe Kriegsvolk greulich wider diejenigen tobt und wüthet, die es billig schützen und vertheidigen sollte: alsdann schreien wir, daß uns Unrecht geschehe, daß wir vermaßen unterdrückt werden sollen, und klagen heftig über die Tyrannei der Fürsten. Aber das sollte man längst zuvor bedacht haben, und man sollte Gott den Herrn gehört haben, der solches den Leuten hat zuvor verkündigen lassen und hat sie mit seinem Wort zur Buße ermahnt. Also sind zu derselbigen Zeit die Sünden der Sichemiten auch erfüllet gewesen, darum ist die Strafe über sie gekommen durch die Thorheit ihrer Fürsten, darein sie auch mit sind gemengt worden, daß sie ihrer Sünden, des Raubes und Schändung der Dirne, theilhaftig wurden, und daß sie also auch zugleich mit ihnen haben müssen umkommen.

89. Und dieweil wir jetzt in unserem Muchern, Stehlen und sonst allerlei anderen Sünden kein Maß haben, wir hören Gottes Wort nicht, fragen auch nach dem Zeugniß unsers eigenen Gewissens nicht, daß unserer bösen Lust und Begierde widerstrebt, so werden wir auch nicht ungestraft bleiben, sondern werden die Strafe tragen müssen nach Gottes rechtem Gericht, und werden alsdann zu lange gewartet haben, und vergeblich klagen, daß wir unbilliger Weise beschweret und unterdrückt werden.

90. Simeon aber und Levi sündigen auf der anderen Seite auch. Denn wiewohl die Sichemiten solche Strafe billig leiden und diese Brüder die allergerechteste Ursache haben, und ihnen auch billig wehe thut, daß Dina, ihre Schwester, dermaßen weggeführt und geschändet war: so thun sie doch gleichwohl unrecht, darum daß ihr Vater es nicht hat haben wollen und sich auch diesen Mord nicht hat gefallen lassen, sondern sie thun dieses aus eigenem frevelhaften Vornehmen, ohne Befehl ihres Vaters.

\begin{quote}
\textbf{V. 27–29. Da kamen die Söhne Jakobs über die Erschlagenen, und plünderten die Stadt, darum, daß sie hatten ihre Schwester geschändet. Und nahmen ihre Schafe, Rinder, Esel, und was in der Stadt und auf dem Felde war. Und alle ihre Habe, alle Kinder und Weiber nahmen sie gefangen, und plünderten alles, was in den Häusern war.}
\end{quote}

91. Nachdem die Bürger erwürgt worden sind, wird nun die Stadt auch geplündert, und werden Weiber und Kinder gefangen genommen und alles Gut geplündert: und dieses geschieht alles ohne Vorwissen und wider den Willen des Vaters. Denn wo er es gewußt hätte, so würde er es ihnen ohne Zweifel verboten haben. Ich halte aber dafür, daß die andern Söhne Jakobs auch mit bei der Plünderung gewesen sind, vornehmlich die von der Lea geboren waren, Joseph und Juda ausgenommen; wiewohl Etliche unter ihnen noch kleine Knaben gewesen sind. Aber die ältesten haben das Gesinde und die Hirten in der Nähe umher zu sich genommen und die Stadt geplündert. Simeon und Levi sind die Rädelsführer gewesen; darum nimmt sie der Vater in seinem letzten Segen 1 Mos. 49, 5.: vasa iniquitatis, das ist, er sagt: Sie haben unrecht gehandelt mit ihren mörderischen Waffen 2c., und verdammt sie als Mörder, und droht ihnen auch, er wolle sie zerstreuen in Israel, welches auch hernach geschehen ist. Levi ist des Erbes beraubt worden, dafür hat man ihm das Priesterthum gegeben, den Bettelstab; denn die Leviten haben nicht von ihrem eigenen Einkommen gelebt, sondern von Opfer und Zehnten. Simeon ist auch mit stetiger Armuth geplagt worden. Also ist diese Sünde hart genug an ihnen gestraft worden bei allen Nachkommen.

92. Die Juden machen diese That so gering, als sie immer können, und wollen gar kein Maß zu halten, ihr Volk zu rühmen, sondern sind damit sehr gelinde; sich selbst schreiben sie das Lob der Gerechtigkeit zu, den Heiden aber nehmen sie allen Ruhm, und schmeicheln sich selbst auch so sehr, daß sie den Text aus dem 49. Cap. V. 5., da Jakob sagt: Simeon und Levi haben unrecht gehandelt mit ihren mörderischen Waffen 2c., von Joseph auslegen, daß er diese zwei Brüder verflucht habe, darum daß sie ihn nach Egypten verkauft hatten. Das thun sie darum, auf daß sie damit die Schmach und Schande ihres Volkes zudecken und daselbst als heilig und gerecht rühmen mögen. Ja, Burgensis zeigt auch an, daß die Juden ein solch Gesetz gemacht haben, daß es erlaubt, ja, auch recht und Gott gefällig sei, einen Heiden umzubringen und zu erwürgen, auf welcherlei Weise oder Gestalt solches geschehen möge von einem Juden. So gar voll Schmeigerei und Haffes sind sie wider uns, als wären sie unsinnig. Josephus, der sonst ein trefflicher Mann gewesen ist, lügt doch auch oftmals schändlich wider die heilige Schrift in den Historien, nämlich darum, daß er seinen Volke so wohl gewogen gewesen ist, als daß er aus Saul so einen großen Herrn und heiligen Mann macht.

93. Es scheint aber doch, als ob die Söhne Jakobs gleichwohl noch mögen einigermaßen entschuldigt werden, als hätten sie recht gethan, dieweil Gott stille dazu schweiget, und mit diesem Exempel die Städte der Heiden erschreckt hat, auf daß sie Jakob nicht verfolgten. Nein, Jakob wird nur sagen, daß sie es zugerichtet haben, daß er vor dem Einkommen des Landes fürchten müsse; aber hernach, da er hat sterben sollen, versucht er sie: darum zeigt er damit offenbar an, daß sie unrecht gethan haben.

94. Und da die That geschehen war, folgt erst die Disputation, ob sie auch recht gethan haben? Wie man im deutschen Sprüchwort zu sagen pflegt: Urlaub kam hernach, und wie bei etlichen Völkern ein Gebrauch gewesen, daß sie zuvor den Dieb an den Galgen gehängt haben und darnach erst um das Recht gefragt. Nem, die Fürsten sagen, daß vieles in der That geschieht, so nicht nach Recht geschieht; und wenn das sollte aufgehoben werden, so wäre das Amt der Fürsten vergeblich. Denn vor der That bedarf man ihres Amtes oder ihrer Kunst nicht, aber nach der That fragt man erst: Ob, das geschehen ist, recht gethan sei oder unrecht? Da ist der Fürsten Amt, daß sie Rath geben, damit Einer sein Recht vertheidigen könne. Darum wird hier gefragt: Ob die Söhne Jakobs mögen entschuldigt werden, welche ja nicht nach den Rechten, sondern mit Hinterlist, mit Lügen und Gewalt gehandelt haben? Wir halten es einfach dafür, daß sie gesündiget haben; denn sie haben über diese That nicht Gericht halten lassen, welches Gericht bei den nächsten Städten umher hätte mögen gesucht werden. Wie sie Jakob auch verdammt, und Gott hat allen ihren Nachkommen solche Strafe aufgelegt, die schwer genug gewesen ist.

\quellseite{236}

95. Aus dieser Frage aber wächst noch eine andere und viel schwerere Frage, nämlich: Warum Jakob seinen Söhnen nicht geboten habe, alles, was sie geplündert hatten, zurück zu geben, vornehmlich was man hat zurück geben können, als, Weiber, Jungfrauen, Knechte, Vieh, und anderes Hab und Gut der Stadt? Und das noch mehr ist, Jakob hat hernach dieses Stück Land der Sichemiter Joseph gegeben; und von Joseph wird gesagt, daß er in dem Lande begraben sei, auf daß er dasselbe Erbe noch möchte einnehmen, da er schon gestorben war. Wo sie nun unrecht gehandelt haben, warum geben sie denn nicht wieder heraus, daß sie mit Gewalt genommen haben? Hierauf mögen wir recht antworten: Daß die Heiligen viel thun mögen, darum sie nicht zu strafen sind, das gleichwohl Anderen nicht zu thun geziemt. Danach, so wird hier das göttliche Urtheil vermischt mit der Ungerechtigkeit der Menschen. Das Erbe und die Güter der Sichemiter liegen ledig, haben keinen Herrn, da diejenigen um der Sünde willen erwürget sind, die sie besessen haben; darum gibt sie Gott billig denen, die zuvor Unrecht und Gewalt gelitten haben. Und die Ursache haben die Söhne Jakobs, um welcher willen sie durch Gottes Verhängniß, der ihnen damit zusieht, die Güter der Sichemiter einnehmen, und da sie dieselben eingenommen haben, behalten. Wie die Kinder Israel von selbst die Egypter nicht haben berauben mögen, sondern die Egypter sind derhalben beraubt worden, weil sie die Strafe verdient hatten, und dem Volke Israel noch ihren Lohn schuldig waren, daß sie ihnen so viele Jahre gedient hatten. Endlich ist auch das Gebot Gottes dazu gekommen, der die Egypter strafen und den Hebräern ihren Dienst vergelten wollte.

96. Auf solche Weise waren auch die Söhne Jakobs und ihr Vater sehr gehöhnt worden, und hatten groß Unrecht und Gewalt gelitten, da ihnen die Jungfrau mit Gewalt beschlafen und entführt war; welches die Erfüllung gewesen ist der Sünden der Amoriter, davon droben im 15. Cap. V. 16. ist gesagt worden: „Die Missethat der Amoriter ist noch nicht alle“ etc. Was soll aber nun Jakob für solche große Schmach haben? Darauf antwortet Gott und spricht zu ihm also: Nimm diesen Raub zur Erstattung für die Schmach, die du und deine Kinder getragen haben. Also wird Jakob auf göttlichen Befehl ein Herr über diese Güter, und wird hier nicht disputirt vom Rechte, damit er sie besitzt.

97. Denn also hält sich Gott gegen uns, wie wir uns pflegen gegen ihn zu halten. Wir thun zuvor, was uns wohlgefällt, und nach der That fragen wir erst um das Recht, und begehren, daß er es sich wolle gefallen lassen, was wir gethan haben. Also straft er uns, ehe denn er nach den Rechten fragt, wie der 18. Psalm V. 27. sagt: „Bei den Verkehrten bist du verkehrt.“ Danach aber läßt er uns schreien und klagen über Gewalt und Unrecht, so wir von denen leiden, die uns nach dem gerechten Urtheil Gottes plagen. Warum werden wir erwürgt? Warum werden uns unsere Güter genommen? Antwort: Die Macht hat Gott wohl, daß er es thun mag und die Ordnung also umkehren; denn wir lassen auch die That vor dem Rechte hergehen.

98. Wiewohl nun die Söhne Jakobs unrecht thun, so straft doch Gott die Sichemiter nicht mit Unrecht, die er zuvor ermahnt hatte, daß sie wider das Recht und seine Gebote nicht thun sollten, daß sie nicht Hurerei treiben, und nicht eines andern Tochter mit Gewalt nehmen und behalten sollten. Dieweil sie aber solcher Ermahnung nicht gefolgt haben, darum werden sie gestraft. Und Jakob, der Gottes Gericht und Rache wohl versteht, so über die Widersacher, die ihn beleidigt hatten, geht, behält von Gottes wegen den Raub, wiewohl nicht ohne Seufzen und Weinen; wie das seine Worte ausweisen, darin er seine Söhne beschuldigt und hart mit ihnen redet.

\subsection{III}

\begin{quote}
\textbf{V. 30. Und Jakob sprach zu Simeon und Levi: Ihr habt mir Unglück zugerichtet, daß ich stinke vor den Einwohnern dieses Landes, den Cananitern und Pherizitern; und ich bin ein geringer Haufe. Wenn sie sich nun versammeln über mich, so werden sie mich schlagen. Also werde ich vertilget sammt meinem Hause.}
\end{quote}

99. Aus diesen Worten mag man sehen, wie sich Jakob die That seiner Söhne habe gefallen lassen. Die Juden versuchen, diese Worte zu bemänteln, und erdichten Glossen darauf, daß es sollen Worte sein, daß sich Jakob gefürchtet und sie nicht gestraft habe. Und darum streiten sie auch hart darüber, daß Jakob seine Söhne nicht soll verdammt haben. Aber hernach, 1 Mos. 49, 5., wird er sie Gefäße der Gewalt und des Unrechts nennen, da er sagt: Die Brüder Simeon und Levi, unrecht haben sie gehandelt mit ihren mörderischen Waffen etc., als sie den Sichemitem Gewalt und Unrecht gethan haben, und wird sie und alle Nachkommen verfluchen. Derhalben ist es gewiß, daß er es nicht aus Furcht geredet hat, da er sagt: „Ihr habt mir Unglück zugerichtet, daß ich stinke vor den Einwohnern dieses Landes“ etc., sondern daß ihm diese Gewalt und Tyrannei mißfallen habe; denn darum nennt er sie Mörder, Räuber, und die Gewalt geübt haben, und macht diese That sehr groß, daß es eine solche Thorheit sei, die zu großer Gefahr und Schmach gerathen werde. Es ist, sagt er, eine zweifache Thorheit, die erstlich wider die Sichemiter unbilliger Weise begangen, und zum Andern, mir selbst auch sehr gefährlich ist.

100. Und dieses ist nun eine neue Anfechtung. Zum Ersten, ist ihm, da er acht Jahre in dem Lande gewohnt hatte, die Tochter geschändet worden, davon er einen sehr großen Schmerz gehabt hat. Da er nun diesen Schmerz und Bekümmerniß etlichermaßen überwunden hat, und nun denkt, wie er dieses Uebel heilen möchte, da richten ihm seine Söhne einen andern und viel größern Lärm an, und so viel an ihnen ist, sind sie daran Schuld, daß ihre Eltern sammt dem ganzen Gesinde verderben müssen. Denn also wird der gute fromme Patriarch geklagt haben: „Ich habe vorhin meine Tochter verloren, nun habt ihr es zweywege gebracht, daß ich auch in Gefahr meines Lebens gekommen bin und daß ihr alle mit mir werdet vertilgt werden.“

101. Darum ist es eine sehr scharfe Strafe und Beschuldigung, und sind nicht solche Worte, die er aus Furcht oder Schrecken geredet haben möchte. Ihr habt, will er sagen, nicht allein Andern Gewalt, sondern auch mir und meinen Weibern, dazu dem ganzen Hausgesinde, und unserem Hab und Gut Schaden gethan; ja, es wird diese Gewalt euch selbst auch zu großem Nachtheile gerathen. Was soll ich nun thun? Ich habe da meine lieben Weiber, von denen ihr geboren seid, nämlich Rahel und Lea; ich habe mein Hausgesinde und die Gemeinde in meinem Hause: wenn die Städte, so allernächst umher liegen, uns mit gesammten Kräften angreifen würden, als die wir diesen Lärm und greulichen Mord angerichtet haben, wie wollen wir solcher großen Menge widerstehen, die wir ein geringer Haufe sind? Ihr solltet euch demnach nicht gerächt haben; welches uns und unsern Gütern zu großem Schaden und Nachtheil geraten wird: sie werden mich tödten, spricht er, und ich sammt meinem ganzen Hause werde vertilget werden.

102. Derhalben hält er die Sünde seiner Söhne nicht gering, sondern macht sie groß und sehr greulich, und da er nun hat sterben sollen, rückt er ihnen dieselbige noch auf und hat derselben nicht vergessen können, und nennt sie Gefäße der Gewalt und des Unrechts, das ist, die unrecht gehandelt haben mit ihren mörderischen Waffen. Denn es hat gewißlich eine große Gefahr daraus entstehen können. Daß aber nun Gott solche Gefahr aus sonderlicher Güte vom Hause Jakobs abgewendet hat, ist eine sonderliche und gar eine göttliche Wohlthat. Also können oftmals die Bürger in einer Stadt einen Lärm anrichten, der der ganzen Stadt zu großem Schaden und Nachtheil gerathen kann; wenn aber die Gefahr abgewendet wird, auch mit einem Gewinne und Nutzen, da ist ein anderer Meister, der ein Helfer ist in der Noth zur rechten Zeit; wie sich Gott allhier gegen Jakob auch erzeigt. Es kommt oft vor, daß Einer in Gefahr geführt wird durch eines Andern Thorheit und Frevel, wie Jakob durch die Thorheit seiner Söhne in Gefahr seines Leibes und seines ganzen Hauses kommt: aber es ist Gottes Werk und seine Wohlthat, der aus dem, das böse ist, etwas Gutes macht und das Böse zum Guten wendet, und macht aus einem bösen Anfange ein fröhlich Ende.

\quellseite{237}

103. Dieses geschieht nicht aus Jakobs, Levis und Simeons Rath, sondern allein aus Gottes Güte; denn welchen Gott lieb hat, der wird auch in der höchsten Gefahr ein gutes glückseliges Ende finden. Und um Eines frommen Mannes willen wird oftmals einem ganzen Lande Rath geschafft; wie die Historie der drei Könige Judas, Israels und Edoms 2 Kön. 3, 9. ff. zeugt, die durch den Weg in der Wüste Edom ihr Heer geführt haben wider der Moabiter König, und da sie sieben Tage reisen zogen, hatte das Heer und das Vieh, das unter ihnen war, kein Wasser, und kamen deshalb in die äußerste Gefahr und Noth. Und der König Israel verzagte schon, klagte und sprach, V. 10.: „O wehe. Der Herr hat diese drei Könige geladen, daß er sie in der Moabiter Hände gäbe.“ Und sie wären auch alle jämmerlich dahin gestorben vor Durst; aber Josaphat war ein guter Bürge da, wie Elisa zum Könige Israels sprach V. 13. 14.: „Was hast du mit mir zu schaffen? Gehe hin zu den Propheten deines Vaters“ etc. „So wahr der Herr Zebaoth lebet, vor dem ich stehe, wenn ich nicht Josaphat, den König Juda, ansähe, ich wollte dich nicht ansehen noch achten.“ Da ist die Gegenwart des Königes Josaphat drei Heeren heilsam, nämlich, daß um seinetwillen Gott Hülfe und Sieg und so viel Wassers gegeben hat, daß nicht allein Menschen und Vieh sind erquickt worden, sondern daß sie auch herrlich wider den Feind gesieget haben. Um weß willen ist es aber geschehen? Ist es um des Königes Israels oder Edoms willen geschehen? Mit nichten, sondern weil der fromme gottlose König Josaphat durch den Frevel der zwei Könige sich dahin hat bringen lassen, daß er mit ihnen wider den Feind gezogen ist, da hat Gott um seinetwillen Hülfe und Sieg gegeben. Also ist Egypten erhalten worden um Josephs willen; Syrien um Naemans willen. Auf dieselbe Weise ist hier auch Alles verloren gewesen, daß Jakob an sich selbst und seiner ganzen Familie gar verzagt. Er hat aber die Verheißung, darum folgt nothwendig die Hülfe, und er muß mit seiner Familie den Völkern und Städten, so daselbst umher lagen, mitten aus den Zähnen gerissen werden.

104. Man hat sich aber auf diese Hülfe und Verheißung nicht also verlassen sollen, daß man dabei Gott versuchen wollte; wie sich die Söhne Jakobs dessen mit großem Stolze gerühmt haben, daß sie einen Patriarchen zum Vater hätten, der die Verheißung hätte, und meinten, daß sie derhalben Alles thun dürften, und die Gewalt, so ihnen an ihrer Schwester widerfahren war, ihres Gefallens rächen möchten. Denn das heißt durch Frevel und Vermessenheit Gott versuchen. Und wiewohl Jakob der göttlichen Hülfe gewiß war, so würde er ihnen doch diese That verboten haben. Daß aber die Gefahr, so von den umliegenden Völkern Jakob und den Seinen hätte widerfahren können, abgewendet wird, das geschieht gewiß nicht um der Söhne, sondern nur um des Vaters willen.

105. Und dieses ist ein trefflich Exempel, daß uns erinnert, daß zuweilen fromme heilige Männer in große Noth und Gefahr kommen nicht durch ihre eigene, sondern anderer Leute Schuld. Und zwar die Andern, die auch mit in solcher Gefahr sind, sehen gar nicht, wie ihnen möchte geholfen, oder wie sie könnten errettet werden, sondern es ist, ihres Erachtens, Alles verloren. Dieweil aber etliche fromme gottlose Menschen, oder auch nur Ein frommer Mensch mit im selbigen Schiffe ist, so muß das Schiff unverletzt und ganz an das Ufer und Land kommen, wiewohl es durch groß Ungewitter hin und her ist getrieben worden; und wenn auch tausend Teufel dawider strebten oder im Schiffe rumorten. Also steht in der Apostel Geschichte am 27. Cap. V. 37., daß um des einigen Pauli willen auch das Schiff, darin er gefahren ist mit zweihundert und sechs und siebzig Seelen, ist erhalten worden.

106. Auf dieselbige Weise währet auch die Welt noch immer um der Kirche willen, sonst würde in einem Augenblick Himmel und Erde im Feuer verzehret werden; denn die Welt ist nicht eines Weizenkorns werth, sintemal sie allenthalben voll Gotteslästerung und eitel gottlosen Wesens ist. Dieweil aber die Kirche mitten unter den Gottlosen ist, so läßt es Gott geschehen und der Kirche willen, daß auch die Gottlosen der gemeinen Güter und Gaben dieses Lebens mit genießen, und alles, was die Welt hat, das hat sie um der Kirche willen. Wie der Engel zu Paulo sagt in den Geschichten der Apostel am 27. Cap. V. 24.: „Gott hat dir geschenkt alle, die mit dir schiffen“, welches wahrlich abgöttische und gottlose Leute waren.

107. Also, daß das Land fruchtbar ist und daß jährlich allerlei Früchte aus der Erde wachsen, solches verdient wahrlich der Papst oder der Türke nicht; denn es sind gottlose Leute; sondern es geschieht darum, weil Gott in der Welt eine kleine Heerde hat, die ihm lieb ist, um derselbigen willen gibt er reichlich alle Wohlthaten und Güter dieses Lebens. Wenn das Häuflein hinweg wäre, so würde die Welt gar zunichte gemacht werden, und würde nie Sodom, und gleich wie Gomorra werden, wie Jes. 1, 9. steht.

108. Darum lasset uns aus diesem Exempel lernen, daß Gott um der Heiligen oder Gläubigen willen auch den Gottlosen hilft, und daß er oftmals den Frevel der narrischsten Menschen, damit sie Lärm und Gefahr erwecken, durch seine Güte bessert und zum Besten wendet; denn also ehret er seine Heiligen. Wie ganz Syrien um des einigen Naeman, Egypten um Josephs und Israel um Josaphats willen erhalten werden. Und werden doch gleichwohl die Gottlosen ihre Strafe leiden müssen; wie die Söhne Jakobs, Simeon und Levi, auch sind gestraft worden, da der Vater sie und alle ihre Nachkommen verflucht.

109. Darnach zeigen diese Worte auch an, daß Jakob, der fromme heilige Mann, abermals kämpft mit einer sehr großen Anfechtung des Geistes und des Glaubens, welcher schier erloschen ist: das Docht glimmt, das Rohr ist zerstoßen und fast zerbrochen. Denn die Worte, die er redet, sind nicht Worte des Glaubens, sondern nur des murrenden Fleisches und des kämpfenden und sehr schwachen Glaubens. Es sind ihm die großen herrlichen Verheißungen fast entfallen, da der Herr zu ihm gesagt hat: „Ich will dich segnen“ etc., und „Du sollst ausgebreitet werden.“ Die Sonne und Sterne des Wortes und der Verheißungen sind verdunkelt und bedeckt mit schwarzen Wolken solches Zitterns, Zweifels, und daß er schier gar verzagt ist. Denn er sagt: Ach, sammelt mein Haus, werde erwürgt und gar vertilget werden; da ist er voll Furcht und Mißglaubens, und ist doch gleichwohl solches keine Verzweiflung.

110. Aber dieses ist die Natur und Art des Glaubens: wenn die Sonne scheint und der Himmel fein klar ist, so ruhen wir in dem Troste, der uns vorgehalten wird, und singen: Te Deum laudamus etc. (Herr Gott! dich loben wir etc.), und wie im 34. Psalm V. 2. steht: „Ich will den Herrn loben allezeit“ etc. Wiederum aber: „Da du dein Antlitz verbargest, erschrak ich“, sagt der 30. Psalm V. 8. 9. 10. „Ich will, Herr, rufen zu dir; dem Herrn will ich flehen. Was ist nütze an meinem Blute, wenn ich todt bin?“ Die Worte lassen sich ansehen, als ob es Worte seien derer, die verzweifeln. Wie David im selbigen Psalm V. 4. sagt: „Herr, du hast meine Seele aus der Hölle geführet; du hast mich lebendig behalten, da ich in die Hölle fuhren“, als wollte er sagen: Es war schier mit mir aus.

111. Auf diese Weise kämpft der Glaube mit der Schwachheit, wenn der Teufel und das Fleisch herrschen, und scheint vor dem, der also kämpft, daß sein Glaube mehr da sei. Wie Jakob hier klagt: Ich hinke und bin ein Greuel geworden vor allen Nachbarn umher: ich bin in die äußerste Gefahr gekommen, und fühle Anderes nichts, denn daß ich gar verderben und zu Boden gehen muß. Wo ist da der Glaube und daß er an die Verheißungen gedacht hätte? Warum denkt er nicht in seinem Herzen lieber also: Du hast mit Gott gekämpft und bist obgelegen, wie viel mehr wirst du den umliegenden Völkern obliegen? Diesen ist ihm alles entfallen, die Anfechtung verschlingt die herrlichen Verheißungen und die schönen Siege des Glaubens. Und er verzweifelt doch gleichwohl noch nicht, wiewohl er Einen, der da verzweifelt, nicht gar ungleich ist. Er mißtrauet Gott noch nicht, und ist nicht ungläubig, wiewohl es so scheint; sondern es ist eine Anfechtung. Wie St. Paulus von sich selbst sagt 2 Cor. 1, 8. 9.: „Wir waren über die Maßen beschweret, daß wir uns auch des Lebens erwegten, und bei uns beschlossen hatten, wir müßten sterben“ etc., „auf daß wir unser Vertrauen nicht auf uns selbst stellten“ etc. Die Heiligen müssen also versucht werden, müssen in die Hölle und in die Tiefe fahren und von dannen wieder in den Himmel geführt werden.

112. Denn Jakob behält noch immer das glimmende Docht in seinem Herzen. Da ihm ist angesagt worden, daß Simeon und Levi die ganze Stadt geschlagen und geplündert hätten, hat er sich dessen nicht enthalten können, er hat müssen heftig darüber bewegt und bekümmert werden, darum daß so viel Leute erwürgt worden waren. Derhalben ist es ohne Zweifel eine sehr große Schwachheit gewesen; da ist Klagen, Heulen und Schreien gewesen zu Gott, und wenn er in seinem Kämmerlein allein gewesen, hat er mit Thränen also gebetet: Lieber Herr Gott, hilf mir, der du mir die Verheißung hast gegeben; denn ich bin jetzt in großer Noth und Gefahr. Das Gebet ist erhört worden nach der Verheißung, denn er hat den Geist des Glaubens, oder der Gnade und des Gebets gehabt. Darum geht ihm nun wiederum der Morgenstern und die Morgenröthe auf, sein Glaube wird aufgerichtet und überwindet alle Gefahr. Denn die Sünde seiner Söhne wird zugedeckt, und der gute fromme Vater wird vor den Nachbarn beschützt, und ist niemand, der sich dürfte wider ihn legen, wie wir im folgenden Capitel hören werden. Den Schutz hatten die zwei Brüder, Simeon und Levi, nicht verdient. Aber dem alten Vater Jakob werden alle Sünden seiner Söhne geschenkt mit dem Raube der Sichemiter, doch also, daß sie gleichwohl beide gestraft werden: die Sichemiter werden gar vertilgt, die Söhne Jakobs aber sind vom Vater verflucht worden.

\quellseite{238}

\subsection{IV}

\begin{quote}
\textbf{V. 31. Sie antworteten aber: Sollten sie denn mit unserer Schwester, als mit einer Hure, handeln?}
\end{quote}

113. Höre doch, wie die Söhne dem Vater antworten. Die hoffärtigen Tropfen erkennen ihre Sünde nicht, und lassen sich auch den Mord und unrechte Gewalt, die sie geübt hatten, nicht leid sein, sondern vertheidigen sie noch dazu. Als wollten sie sagen: Wir haben daran recht gethan, daß wir die Sichemiter erwürgt haben; denn greuliche Sünden sollen billig auch greulich gestraft werden. Und darum machen sie es auch sehr greulich und groß, daß Sichem Dina entführt hatte, so er sie doch nicht zur Hure, sondern zu seinem Eheweib zu haben begehrte; und sie trösten auch den frommen alten Mann nicht einmal, lassen sich auch nicht vernehmen, daß es ihnen leid wäre, daß sie den Mord begangen hatten, lindern solche That gar nicht damit, daß sie doch die Sünde erkannt und die Schuld dem Vater abgelegen hätten. Sie thun keine Buße und begehren nicht, daß ihnen möchte verziehen werden, daß sie mit Mord und Räuberei die Nachbarn hart erzürnt und den Vater in seinem Herzen sehr betrübt hatten.

114. Der gute alte Mann hat solche stolze Antwort und daß sie ihre Sünde noch vertheidigen wollen, ertragen müssen; welches ihm sein Herz, das schon des vorigen Schreckens halben betrübt und matt war, noch mehr geplagt und betrübt hat. Jakob muß die Sünde seiner Söhne tragen und um seiner willen verzieht es sich mit der Strafe, und die stolzen muthwilligen Söhne werden geduldet bis zu seiner Zeit.

115. Es leuchtet aber in diesem Exempel eine wunderbare große Güte und Gnade Gottes, der den Gottlosen hilft und sie errettet, auch dann, wenn Alles scheint gar verloren zu sein; wo nur das Docht glimmend bleibt, und daß man sich nicht zur Bosheit und Lästerung kehrt, sondern ächzt, seufzt und schreit. Solches Geschrei erfüllt Himmel und Erde. Alle Sachen dieses unsers Patriarchen haben so gar übel gestanden, daß es schien, als könnten sie nicht wiederum zurecht gebracht werden: aber Gott verläßt den nicht, dem er zum öftern Hülfe und Errettung verheißen hat. Es ist zwar erbärmlich gewesen, daß er in solche Gefahr und Angst hat gerathen sollen, und in der Noth und Gefahr von seinen Söhnen gleichsam verspottet und verachtet werden, und daß immer ein Schmerz zum andern gekommen ist: aber er erträgt dieses alles mit wunderbarer Geduld und Gott errettet ihn also. Unser Herr Gott greift darein, so wird es alles gut; wie nun folgt.

\chapter{Das fünfunddreißigste Capitel}

\section[Erster Theil]{Erster Theil. Wie Gott mit Jakob redet, und ihm Befehl gibt, einen Altar aufzurichten; wie Jakob darauf an seine Gemeine eine Predigt hält}

\subsection{I}

\begin{quote}
\textbf{V. 1. Und Gott sprach zu Jakob: Mach dich auf, und zieh gen Bethel, und wohne daselbst, und mache daselbst einen Altar dem Gott, der dir erschienen, da du flohest vor deinem Bruder Esau.}
\end{quote}

1. Bisher hat der Patriarch Jakob großen Schmerz gehabt und ist in großer Angst gewesen, da immer eine Anfechtung auf die andere gekommen ist. Es ist ihm seine einzige Tochter geschwächt worden; die Sichemiter sind erwürgt worden; und aus demselben Morde ist eine große Gefahr entstanden, dadurch er und sein ganzes Haus vertilgt und umgekommen wäre, wo ihn Gott in solcher Gefahr nicht beschützt und errettet hätte. Denn es hat ihm da kein Mensch rathen und helfen können, und war auch niemand, der den armen, elenden, betrübten Vater getröstet hätte. Seine Söhne antworten ihm auf die Bestrafung mit großem Stolz und rechtfertigen noch ihre sehr böse Sache. Also ist er nun von Sorgen und greulichen Schmerzen bekümmert, daß er gleichsam zwischen Himmel und Erde hängt; und kommt die Schwachheit des Glaubens auch dazu, damit er so hart kämpft, daß er aller Verheißungen vergißt, die wahrlich sehr herrlich waren, vergißt auch dazu der vorigen Erlösung, darin er die Hand und gnädige Hülfe Gottes erfahren hatte.

2. Diese Exempel aber der Väter werden uns darum vorgehalten, daß man sie in der Gemeinde Gottes lesen und lernen soll, auf daß die Gottlosen sehen mögen, daß sie wunderbarlich versucht worden sind durch Schwachheit und durch Stärke, daß sie bald gesiegt, und auch bald verstoßen worden sind durch Glück und Unglück, daß sie endlich durch mancherlei schwere Anfechtung Märtyrer würden, ein süßer Geruch, und „gerader Rauch“, wie im Hohen Liede Salomos steht am 3. Cap. V. 6., und daß also das Fleischliche an ihnen vollkommenlich getödtet werden möchte. Darnach lehren uns diese Exempel auch, daß unsere Leiden kindisch und gering sind gegen solche schwere Trübsale der Väter; darum sollen wir daraus Trost nehmen und unsere Herzen bereiten zur Geduld. Denn wir haben mit solcher großen Bekümmerung, Jammer und Noth noch nicht gekämpft, wie sie dieser eine Patriarch in seinem Leben getragen hat.

3. Jetzt aber kommt nun Gott und hilft seinem Patriarchen in der äußersten Noth und da es mit ihm gleich alles verloren war; wie er denn genannt wird und auch in Wahrheit ist ein rechter Helfer in der Noth, wie der 9. Psalm V. 10. 11. sagt: „Der Herr ist des Armen Schutz, ein Schutz in der Noth“ etc. Denn er verlässet seine Heiligen nicht, die auf ihn hoffen, wiewohl es scheint, als ob sie gar verlassen und verstoßen seien; denn das heißt nicht verlassen, nämlich, wenn sie fast verlassen sind. Aber Gott macht aus unserm Nein ein Ja. Also, wenn die Gottlosen sagen: Es ist keine Hülfe bei Gott, ich bin gar verloren: da antwortet Gott also: Du bist nicht verloren, sollst auch nicht verderben, wie du urtheilst; sondern ich will dir Mund und Weisheit geben auch in dem äußersten und höchsten Unglück, daß du nicht sollst verlassen werden. Es scheint wohl, als ob es eine Verlassung sei, aber in Wahrheit ist sie es nicht. Darum läßt sich Gott abermals mit seinem Worte hören und tröstet den betrübten und bekümmerten alten Mann.

4. Und wir haben nun oft gesagt, daß man in den Legenden der heiligen Väter vornehmlich darauf Acht geben soll, wenn Gott mit ihnen redet. Darum übertreffen die Historien in der heiligen Bibel aller andern Völker Historien. Denn sie sind heilig und nützlich um der Ursache willen, daß Gott daselbst redet. Und man soll sie nicht nur so flüchtig überlaufen, wie der stummen Heiligen Legenden, da GOtt nicht redet; sondern man soll den Werth und Schatz, der diese Historien theuer macht, erkennen und betrachten, nämlich GOttes Wort, das er mit den heiligen Vätern redet. Ja, das ist es, daß diese Legenden gelten macht. Und wenn man die Exempel der Eremiten oder Einsiedler hiermit vergleichen wollte, werden solche Exempel gar verdunkelt. Dieselbigen Einsiedler mögen zwar wohl nach den Geboten GOttes und im Glauben gelebt haben und werden derhalben auch billig für große Männer gehalten: sie sind aber diesen Vätern mit nichten gleich, welchen GOtt in der Noth begegnet und eine neue Sonne aufgehen läßt. Die andern Legenden haben ein groß Gepränge wegen der großen und ungewöhnlichen Wunderwerke; aber die Historien der Väter muß man nicht nach solchen Werken beurtheilen, sondern vielmehr nach der Heiligkeit, welche GOttes Wort ist.

\quellseite{239}

5. Da nun GOtt sieht, daß Jakob verlassen ist, und von seinen Söhnen verachtet und verspottet wird, und daß sie sich darüber nicht sehr bekümmern, daß sie ihren Vater in sehr große Gefahr gebracht und sehr betrübt hatten, da kommt er nun zu rechter Zeit und tröstet ihn. Denn das Seufzen hat Himmel und Erde erfüllt. Lieber Jakob, will er sagen, dein Gebet und Thränen zwingen mich, daß ich dir helfen muß. Wovon dich dünkt, daß es dein Verderben werde sein, soll dir gar nichts schaden. Ich bin der HErr, dein GOtt; durch mein Verhängniß sind die Sicherheiten erzwungen worden, und ist dieses ganze Spiel dazu angerichtet, auf daß ich dich prüfen, dich versuchen, und dich auserwählt und bewährt machen möchte. Aber sei du nur getrost; ich will allen grimmigen Zorn, Toben und Wüthen des Volkes wehren und inne halten, und will deine Traurigkeit und Seulen in Freude verändern, und du sollst ein Herr zu Sichem sein; denn also habe ich es verordnet. Wie Rehabeam, da er die zehen Stämme wieder an sich bringen wollte, wird ihm solches zu thun verboten. Denn es kam Gottes Wort zu Semaja, dem Mann GOttes, und sprach, 2 Chron. 11, 3. 4.: „Sage Rehabeam, dem Sohn Salomos: So spricht der HErr: Ihr sollt nicht hinauf ziehen, noch wider eure Brüder streiten. Ein jeglicher gehe wieder heim, denn das ist von mir geschehen“ etc., nämlich, daß ich damit Salomo, den Vater, und seinen Sohn, Rehabeam, strafen möchte. Also läßt GOtt bisweilen den Türken und Pabst wüthen; und wiederum, wenn er will, hält er sie beide zurück. Denn der dem Meer sein Ziel gesetzt hat, der hat auch dem Zorn und Wüthen des Teufels und der Welt ihr Ziel gesetzt. Wie Christus im Garten sagt, Luc. 22, 51.: „Lasset sie doch so ferne machen“ etc.

6. Dieses läßt aber GOtt alles darum geschehen, daß er uns damit lehren will, daß wir ihn anrufen, zu ihm schreien und seufzen sollen. Denn dieses ist der „gerade Rauch, so herauf geht aus der Wüste“, wie im Hohen Liede Salomos steht am 3. Cap. V. 6. Auf solche Weise wird hier das sehr große Ungewitter gestillt, und Jakob empfängt einen sehr angenehmen Trost und bekommt Frieden in seinem Herzen, da GOtt zu ihm sagt: Die Bürger der umliegenden Städte sollen dich nicht überfallen; verachte nur solch Ungewitter und den grimmigen Zorn der Völker; gehe hin, predige, lies und übe den Gottesdienst, bete; laß mich sorgen, wie dieses Ungewitter, und die Wellen, so daher gehen, mögen gestillt werden.

7. Bethel aber ist nicht weit von Sichem gelegen, und wir haben droben vom Namen disputirt. Jerusalem hat zween Namen, und Viele meinen, daß Bethel Jerusalem, oder daß es der Berg Morija sei, darauf der Tempel gebaut gewesen ist; und behält noch heutiges Tages den alten Namen bei den Türken. David hat denselbigen Hügel von Arauna, dem Jebusiter, gekauft und hat daselbst einen Altar gebaut, 2 Sam. 24, 24. 25. Denn da hat Abraham den Isaak geopfert, und hat den Namen von der Ehrerbietung und vom Gottesdienst, weil auf demselbigen Berge Adam, Eva, Kain und Abel geopfert haben; und hernach haben ihn die nachfolgenden Patriarchen für und für heilig gehalten und ist bei ihnen zur Uebung und Gebrauch des Gottesdienstes verordnet gewesen. Darnach aber haben ihn die Jebusiter eingenommen; aber David hat ihn wieder erobert und angerichtet.

8. Daß Morija aber Bethel nicht ist, kann man daher wissen: Bethel ist eine Stätte gewesen nahe bei Ai und Sichem gelegen, ist nicht im Stamm Juda gewesen, auch nicht im Stamm Benjamin. Und daher sind die Götzendiener gekommen, mit denen Hosea und die andern Propheten für und für gestritten haben. Denn sie haben sich dessen gerühmt, daß GOtt mit den Vätern an dieser Stätte geredet hätte; darum sollte auch die Stätte besser sein und billig in größeren Ehren gehalten werden denn der Tempel zu Jerusalem. Lyra ist der Meinung, daß Bethel Morija sei. Das däucht mich nicht wahr zu sein, oder wo es ein und dieselbe Stätte ist, so muß hernach der Name verändert worden sein. Nach den Geographen ist Bethel Morija nicht, ist auch Jerusalem nicht, oder ist mit Anspielung auf die Bedeutung des Namens so genannt worden, nämlich darum, daß es GOttes Haus war.

9. Der HErr aber gebietet dem Patriarchen Jakob, er solle zu Bethel einen Altar bauen und erinnert ihn der vorigen Trübsal und Noth und tröstet ihn mit den Exempeln der vorigen Anfechtungen und Errettungen. Als wollte er sagen: Gedenke an die Verheißungen, so dir vor der Zeit geschehen sind, und wisse dich zu erinnern, was dir zuvor begegnet ist, da du vor deinem Bruder Esau flohest, und wie viel du gelitten hast, und wie gnädig ich dich erlöset und geschützt habe. Von der Zeit aber, da Jakob vor seinem Bruder Esau geflohen war, bis an diese Zeit, sind fast bei dreißig Jahr verlaufen. Denn zwanzig Jahr ist er in Mesopotamien gewesen, acht Jahr hat er zu Sichem gewohnt, zwei Jahr ist er auf dem Wege gewesen, und ist fast die ganze Zeit über mit der Feindschaft seines Bruders geplagt gewesen. Darum gedenke, spricht der HErr zu Jakob, daß du zuvor auch in Bethel eine Verheißung gehabt hast, daß ich dich schützen wollte, und wie väterlich ich dich aus so vielen und großen Anfechtungen errettet und meine Verheißung getreulich gehalten habe. Da Jakob die Worte hört, kommt er wieder zu sich selbst und faßt Muth, und wird ihm die Wunde, so ihm von den Söhnen geschlagen war, alsobald gestillt. Denn GOttes Wort macht lebendig und tröstet die, so bekümmert sind.

10. Es wird aber im Text nicht ausgedrückt, wie GOtt mit Jakob geredet habe, ob es im Traum geschehen sei, oder aber durch ein deutliches Gesicht, oder durch irgend einen Propheten. Denn diese dreierlei Weisen pflegt GOtt zu gebrauchen. Nun ist aber zu der Zeit gar kein Prophet oder Patriarch vorhanden gewesen; denn sie waren alle gestorben. Es wäre denn, daß vielleicht Isaak, der Vater, zu Jakob, seinem Sohne, einen Boten gesandt, oder ihm Briefe geschickt hätte, darin er ihm befohlen hätte, umzukehren und gen Bethel zu ziehen. Auf welcherlei Weise es aber nun geschehen ist, so ist es daran genug, daß wir gehört haben, daß GOtt mit ihm geredet hat. Er rede nun gleich, auf welche Weise er will, so ist es doch gewiß, daß solches GOttes Wort ist.

11. Allein, daß die menschliche Natur so gar verkehrt und verderbt ist, daß wir, leider! nicht glauben, daß wir GOttes Wort hören, wenn er etwa durch einen Menschen mit uns redet. Denn wir bewegen das Wort nach dem Ansehen und Größe deß, der da redet: den Menschen, der da redet, hören wir als einen bloßen Menschen, und denken, es sei das Wort eines Menschen. Und darum verachten wir es auch und werden dessen überdrüssig, so wir doch GOtt dafür danken sollten, daß er sein göttlich Wort dem Menschen oder Diener in den Mund legt, welcher uns allenthalben gleich ist, der mit uns reden, und mit dem Wort trösten und aufrichten kann.

12. Wir erkennen aber die Größe dieser Wohlthat nicht, danken auch GOtt nicht dafür. Darum sind in dem Stück die Legenden der heiligen Väter auch gar vortrefflich, nicht allein um des Wortes, sondern auch um des Glaubens willen der Patriarchen, damit sie sich fest und gewißlich an das Wort gehalten haben. Wir können es ja, leider! nicht so gewiß dafür halten, wenn wir einen Diener des Evangeliums predigen und absolvieren hören, daß wir dann von der göttlichen Majestät selbst gelehrt und absolvirt sind, und Trost empfangen haben. Wir sagen wohl, es sei ein Wort, das wir in der Kirche hören, aber des Pfarrherrn Wort. Nun ist es aber nicht des Pfarrherrn, noch St. Petri, noch auch eines andern Dieners Wort, sondern der göttlichen Majestät selbst.

13. Derhalben sind nun unsere Historien kälter um unsers Glaubens willen, welcher kälter ist. An den Vätern aber ist der Glaube viel brünstiger gewesen; wie in ihrem ganzen Leben ein großer wunderbarer Glaube leuchtet. Darum werden auch die Historien der Väter billig heilig genannt und sind auch heilig. Wenn wir nun glauben könnten, daß GOtt mit uns redet durch unsere Eltern, durch die Pfarrherren und Diener des Worts, so würden wir empfinden und erfahren, daß unsere Herzen mit großer Freude würden entzündet werden. Denn wir würden uns also rühmen und sagen: Ich bin getauft. Von wem, vom Pfarrherrn? Antwort: Mit nichten, sondern vom Heiligen Geist; ich bin absolviret vom Heiligen Geist und von GOtt selbst. Warum soll ich mich denn fürchten? Vor wem soll ich erschrecken? Warum soll ich mich meiner Sünden halben bekümmern? Darum sind wir, mit den Vätern verglichen, kaum kleine Kinder und Säuglinge. Denn sie haben bei den Heiden, so allernächst umher gewohnt haben, viel Aergerniss der Abgötterei gesehen, und haben doch das Wort GOttes so fest angenommen und sich daran gehalten. Wir haben vor unsern Augen viel unzählige Exempel des Glaubens aus unsern und auch aus vorigen Zeiten und glauben doch gleichwohl nicht.

\quellseite{240}

14. So soll man nun zu allererst fleißig auf die Worte GOttes merken, der da redet; darnach auf den Glauben derer, die geglaubt haben: alsdann wirst du die rechten Wunderwerke sehen in den Historien der Heiligen, die durch den Glauben an das Wort allerlei Unglück der Welt und der Teufel überwunden haben. Und solchen Sieg haben sie erlangt nicht durch ihre Weisheit; denn sie haben keinen Rath gewußt, wie der 107. Psalm V. 27. sagt: sondern um des Worts willen, so zu ihnen gekommen ist und sie getröstet hat, und durch den Glauben, damit sie an dem Wort gehangen haben. Denn „unser Glaube ist der Sieg, der die Welt überwunden hat“, 1 Joh. 5, 4. GOtt macht uns zu Siegern und Ueberwindern durch sein Wort, damit er uns aufrichtet und stärkt. Darum ist das ein großer herrlicher Trost, den GOtt allhier dem Jakob gibt; denn es ist eben so viel, als ob er sagen wollte: Ich habe Tod und Leben, und alles Ungestüm, Ungewitter und Trübsal in meiner Hand: glaube du nur, so wird dir geholfen werden.

15. Er gibt ihm aber auch ein Gebot, einen Altar aufzurichten, das ist, den Gottesdienst anzurichten. Dabei muß man diese Regel merken: Ohne den Heiligen Geist soll man keinerlei Art oder Weise vornehmen, GOtt zu dienen. Wie Jakob nicht aus seinem eigenen Willen erdichtet, daß an diesem Ort das Haus GOttes sei; sondern er hat es aus der Offenbarung gelernt, da er droben, 1 Mos. 28, 11. 12., am selbigen Ort geschlafen und die Engel gesehen hat auf und nieder fahren; da sagt er V. 16.: „Gewißlich ist der HErr an diesem Ort“, und hier wohnt gewißlich GOtt. Darum soll man nicht aus eigenem Vorwitz Altäre oder sonderliche Gottesdienste aufrichten, sondern man soll auf GOttes Gebot sehen.

16. Jerobeam hat in nachfolgenden Zeiten zu Bethel einen Kalb aufgerichtet, auf daß das Volk daselbst opfern und anbeten sollte, 1 Kön. 12, 28. ff. Und damit er das Gebot GOttes, den Tempel zu Jerusalem zu bauen, das Salomo hatte, verdunkeln möchte, so hat er dieses Gebot, so dem Patriarchen Jakob gegeben gewesen, einen Altar und GOttesdienst zu Bethel aufzurichten, hoch gerühmt. Denn dasselbe, sagt er, wäre ja viel älter, und derhalben müßte es auch köstlicher und herrlicher sein denn das neue Gebot, den Tempel zu Jerusalem zu bauen.

17. Wie vor dieser Zeit, da man GOtt in seinem Wort und Sacramenten hat suchen und erkennen sollen, hat der Pabst beides, Wort und Sacramente, verworfen, und hat das Heiligthum und Gedächtniß der Heiligen errichtet, damit er das arme Volk betrogen hat, daß sie glauben sollten, daß an solche Dinge der Gottesdienst gebunden wäre: und sind darüber die Leute in solche Thorheit gerathen, daß sie an etlichen Orten Josephs Hosen, die Milch der Jungfrau Maria und des St. Franciscus Unterkleid als Heiligthum gerühmt und ausgerufen haben. Dazu sind die erdichteten und wunderlichen Werke der Heiligen gekommen, als des Franciscus, Dominicus und Anderer mehr, welcher Heiligkeit sie nach den Werken und äußerlichem Leben beurtheilt haben. Denn also hat man die Leute betrogen und narren müssen, die sich von der Wahrheit abgewendet hatten. Wenn du aber die heilige Schrift fleißig lesen wirst, so wirst du sehen, daß GOtt ohne das Wort nichts gefällt.

\subsection{II}

\begin{quote}
\textbf{V. 2. Da sprach Jakob zu seinem Hause, und zu allen, die mit ihm waren: Thut von euch die fremden Götter, so unter euch sind, und reiniget euch, und ändert eure Kleider.}
\end{quote}

18. Jakob hält sich an den Trost, den er im Wort GOttes gehört hat: derhalben hört das Ungewitter auf und das Meer wird ganz stille, der Winter ist vergangen, die schönen Blumen lassen sich sehen, die Nachtigall singt. Darum ist er nun fröhlich und dankt GOtt: Gelobet sei GOtt der HErr, die Traurigkeit und Angst hören auf; denn das Licht der Sonne ist aufgegangen und vertreibt die Finsterniß der Trübsal. Und wiederholet auch nun die Worte selbst, die er GOtt hat reden hören, und erzählt die Anfechtung, die ihm in ohngefähr dreißig Jahren widerfahren sind. Er ist aber noch nicht an allem Jammer und Elend vorüber; denn es ist noch das greuliche Ungewitter und Unglück vorhanden mit seinem Sohn Joseph. Doch dankt er GOtt dafür, daß er das vorige Unglück überwunden hat, und rüstet sich gen Bethel zu ziehen, da ihn der HErr zuerst getröstet hatte.

19. So sagt nun Mose, wie Jakob zu seinem Hause, und Allen, die mit ihm waren, geredet habe. Da macht er einen Unterschied zwischen den Hausgenossen und den Fremdlingen. Die Hausgenossen waren seine Weiber, Kinder, Knechte und Mägde, unter den Andern aber, die mit ihm waren, versteht er die, welcher Herz GOtt gerühret hatte, daß sie sich zum Hause Jakobs geschlagen, entweder daß sie der Hoffnung gewesen, sich mit ihnen zu verheirathen, oder aber daß sie vom Raube und Beute der Sichemiter herzu gekommen waren. Denn ich habe oftmals gesagt, daß es sehr glaublich sei, daß da die Patriarchen gelehrt und gepredigt haben, Viele von den Heiden zu ihnen gekommen sind, die gesehen haben, daß die Patriarchen fromme, gottselige und heilige Männer gewesen und daß GOtt mit ihnen war, und haben darum ihre Lehre gehört und angenommen.

20. Denn dieweil sie von GOtt als Gesandte und Prediger des Worts in die Welt gesandt gewesen sind, sollen wir nicht denken, daß ihr Dienst und Amt ohne Frucht abgegangen sei. Und haben sich nicht allein die zu ihnen gehalten, so vom Geblüt der Patriarchen waren, sondern auch Fremde; wie droben, Cap. 14. V. 13., gewesen sind die Bundesgenossen Abrahams, Escol und Aner, die ohne Zweifel alle das Wort werden gehört haben, deßgleichen auch Abimelech. Und hernach hat Joseph in Egypten, Daniel zu Babel, Jona zu Ninive die Lehre vom rechten wahren GOtt gelehrt. Derhalben hat nun GOtt die Kirche in der Welt versammelt nicht allein aus dem einen Geschlecht der Patriarchen, sondern aus allerlei Völkern, zu welchen das Wort gekommen ist.

21. Auf solche Weise haben sich Viele von den Cananitern zu Joseph geschlagen und haben seinem Wort geglaubt, welche GOtt errettet hat von dem Unglück, so vorhanden gewesen ist, und von der Abgötterei. Wie sich die Hure Rahab im Buch Josua am 2. Cap. V. 1. ff., da sie gesehen, daß die ganze Stadt in Gefahr und Noth war, zu den Kundschaftern geschlagen, und dieselbigen gebeten hat, daß sie sie und ihre Freundschaft schützen und erhalten wollten. Denn es werden allezeit Etlicher Herzen gerühret, zu welchen die, so gottselig sind, in der Welt kommen. Und solche Leute, halte ich, sind die gewesen, die Jakob mit sich geführt hat. Denn das Wort wird ohne Frucht nicht gelehrt, sondern sammelt einen Haufen in der Welt nicht allein von denen, die das mündliche Wort hören und die in der Kirche und Gemeinde der Patriarchen sind, sondern bringt auch die herzu, die fremd sind, nämlich, die dazu versehen sind.

22. Danach ist auch dazu gekommen, daß sie sich mit den Cananitern befreundet hätten um des Heirathens willen der zwölf Patriarchen, welche heidnische Weiber zur Ehe genommen haben. Juda hat die Cananitin Thamar genommen, Joseph die Tochter eines egyptischen Priesters. Also ist das Geblüt der Heiden und der Israeliten zusammen gekommen, daß auch die Moabitin Ruth eine Mutter Christi wurde. Denn GOtt ist nicht allein der Juden GOtt, sondern auch der Heiden, Röm. 3, 29., wiewohl er durch die Juden, als durch seine Diener, seine Gnade und sein Wort in die Welt ausgebreitet hat. Endlich haben auch die Patriarchen Gemeinschaft gehabt mit den Heiden. Daher haben sie oftmals Gelegenheit genommen, sich mit einander zu unterreden von GOtt und von Religionssachen. Denn sie sind nicht stumme Klöße und Steine gewesen; sondern sie haben mit den Leuten einen solchen Wandel geführt, daß sie die, so fremd gewesen sind, mit sonderlicher Freundlichkeit zu sich gezogen haben, daß sie ihre Mitgenossen in der Lehre und Religion werden möchten.

\quellseite{241}

23. Es ist aber zumal eine feine Predigt, die der Patriarch Jakob seinem Hause und Gemeinde thut, die versammelt war von seinen Hausgenossen und auch von Fremden; in welcher Predigt er sie ermahnt, sich zu bessern, und von den Sünden abzustehen, die noch in derselbigen Kirche und Gemeinde gewesen sind. Und die erste Sünde, die er straft, ist wider das erste Gebot, nämlich Abgötterei; da will er haben, daß sie sich darin bessern und solche Sünde abschaffen sollen. „Thut von euch“, spricht er, „die fremden Götter, so unter euch sind.“

24. Denn das ist das Hauptstück und der Ursprung, davon alle anderen Sünden herkommen. Und alle Reformation oder Besserung, so vorgenommen werden mag, ist vergeblich, wo nicht erst die Lehre gereinigt wird. Denn siehe an die Thorheit des Papstes und aller nachfolgenden Concilien, welche zu allererst etliche äußerliche Ceremonien verordnen, als, daß sie den Priestern gebieten, lange Kleider zu tragen, ihre sieben Zeiten und Messe fleißig zu lesen, verbieten ihnen, zu spielen und Hurerei zu treiben. Dasselbe heißen sie eine Reformation der Kirche. Und wo man noch dermaleinst ein Concil halten wird, so wird man solch Ding darin abhandeln und verordnen. Denn die Bischöfe und Cardinäle sind grobe ungeschickte Leute, die keine Gedanken haben auf das Wort und auf die Lehre, verstehen es auch nicht und fragen gar nichts darnach.

25. Darum „muß die Axt den Bäumen an die Wurzel gelegt werden“, Luc. 3, 9., und man muß den Kopf der Schlange angreifen. Denn wie hart nun sie auf den Leib schlägt, kann sie doch solches alles ohne Gefahr ertragen; wenn du ihr aber auch nur mit einer kleinen Ruthe den Kopf triffst, so stirbt sie bald dahin. Nun ist aber die vornehmste Gewalt des Satans dahin gerichtet, daß er wider das Wort und Lehre streitet, dieselbige zu vertilgen; welche Lehre im ersten Gebot begriffen ist. Die greift er sehr hart an. Derhalben soll man sich vor allen Dingen befleißigen, daß wir die rechte und gewisse Lehre von Gott haben mögen. Da mag man dann eine rechte Reformation und Kirchenordnung anrichten.

26. Wenn sich der König Ferdinand rüstet wider den Türken zu ziehen, so läßt er gebieten, daß man fasten, item, Processionen und Heiligenfahrt halten soll. Wozu dient aber solch Ding anders, denn daß es ein Gespött des Teufels ist? Das sollte man aber lieber vornehmen, daß der rechte Glaube und die reine Lehre vom rechten Gottesdienste möchte gelehrt werden. Aber das fahren sie um, vertheidigen die offenbaren Irrthümer in der Lehre und bestätigen noch immer die Abgötterei: darnach wollen sie dann mit äußerlichen Ceremonien und Bittfahrten den Zorn Gottes stillen. Also wird Gott zwiefältig erzürnt und verspottet.

27. Derhalben so lasset uns an dieses Gebot Jakobs denken, der da sagt: „Thut von euch die fremden Götter“, das ist, „schafft die Aergernisse der Lehre und des gottlosen Wesens ab, auf daß das Volk lernen möge Gott recht fürchten und ihm vertrauen; alsdann so werden die langen Kleider, Platten und andere dergleichen äußerliche Dinge mehr auch ihre Statt finden.“

28. Die zwei Worte aber „fremde Götter“ stehen im Hebräischen nicht in demselben Numerus. Das Wort „Götter“ steht in der Mehrzahl, das andere Wort aber „fremd“ steht in der Einzahl und Genitiv, und lautet also: „Die Götter des Fremden“ oder „einer fremden Sache.“ Das hebräische Wort aber, nekar, hat eine doppelte Bedeutung, deren eine das Gegentheil der andern ausdrückt; denn es bezeichnet beides: unbekannt oder fremd, und: bekannt. Wie hernach in Einem Verse, Cap. 42. V. 8., beiderlei Bedeutung steht, Joseph erkannte seine Brüder, und er entfremdete sich seinen Brüdern, das ist, er machte sich unkenntlich. Es ist dies die Redefigur antiphrasis, die sich in allen Sprachen findet. Also nennen die Lateiner bellum, Krieg, als das gar kein gut Ding ist; und lucus heißt ein Wald, darum daß er gar finster ist und nicht leuchtet. Die Deutschen haben solcher Worte viele, als, wenn sie sagen: Halt, ich will dir rathen, will dir helfen; das ist auch durch die Figur antiphrasis geredet, und ist dieses die Meinung, daß ich einen Andern strafen oder verderben will. Item, wenn man sagt: Der Henker kann dem Diebe wohl rathen etc.

29. Eben also verhält es sich auch mit der Bedeutung dieses Wortes: die Götter eines unbekannten oder bekannten Dinges. Den rechten Verstand aber muß man nehmen von der Art und Weise zu reden. Und an diesem Orte muß das hebräische Wort nekar nothwendigerweise unbekannt und fremd bedeuten. Denn wie das lateinische Wort bellum seinen Namen hat vom Schönen, und der Krieg, bellum, den Unversöhnlichen sehr schön, hellum, und angenehm dünkt zu sein, wie denn auch die Rache süß ist (aber die Erfahrung bringt es mit, daß wir es fühlen und empfinden, daß es ein sehr betrübt und elend Ding ist um den Krieg; denn es wird bellum genannt durch die Figur antiphrasis, es heißt honig gesäet und süß aufgesäet): also gelüstet auch anfänglich die Leute der neue Gottesdienst sehr wohl, und läßt sich ansehen, als ob es die rechte lautere Wahrheit sei; es scheint ein lieblich süß Getränk zu sein. Denn also pflegen die Rottengeister viel zu versprechen und zu rühmen von ihrer Lehre, daß nichts Gewisseres, nichts Besseres oder Köstlicheres sei; aber am Ende erfahren wir es erst, daß es nur Finsterniß, Irrthum und Unverstand ist.

30. Wenn aber der Mensch in seinem Herzen einmal des rechten, wahren und einigen Gottes gefehlet hat, so muß eine große unzählige Menge und Unordnung mit vielen Göttern daraus folgen. Wer einmal Gott verlassen wird, da ist es unmöglich, daß derselbige bei einer gewissen Gottheit, das ist, bei Einer Religion oder einem gewissen Gottesdienste bleiben sollte. Wie aus dem Unglauben, oder wo man den Glauben verloren hat, unzählig viel Sünden folgen müssen. Dasselbe mag man in allen Historien sehen. Da Israel von Gott ist abgefallen, haben sie dem Baal, Damos, Ashtaroth und andern Göttern gedienet. Da wir im Papstthum von Christo, dem Sohne Gottes, unserm Heilande, sind abgefallen, da ist Dominicus, Franciscus, Vincentius, Christophorus, Clara und andere vermeinte Heilige mehr an Christi Statt gekommen, und dazu viele andere unzählige neue Gottesdienste, welche täglich die Pfaffen und Mönche erdichtet haben, ohne Ende und ohne Maß. Denn die Lust und der Vorwitz der Menschen hat auf keinerlei Weise können gestillt werden.

31. Derhalben, wo man von dem einigen Gott abweicht, das ist eben so viel, als daß man in der Menge vieler unbekannten Götter gar ersäuft wird. Und um solcher Ursache willen hat Jakob die fremden Götter genannt Götter eines fremden, unbekannten und thörichten Gottesdienstes, der doch sehr wohl bekannt und heilig schien. Darum sollen wir unsere Herzen und allen Fleiß auf den einigen Gott wenden, welcher ist Gott Vater, Sohn und Heiliger Geist, und sollen bleiben bei dem einigen Mittler Christo Jesu. Dieses ist der erste Theil der Reformation, so Jakob seinem Hause vorgehalten hat.

32. Die Juden verstehen hier unter den fremden Göttern die Bilder der Götzen, deren sie für und für viele gehabt; und sie haben so viele Götter gehabt, als Teufel in der Hölle sind. Wie wir im Papstthum auch gesehen haben. Jakob aber versteht nicht allein die goldenen und silbernen Bilder, sondern das denen anhängt, oder die Neigung, so die Menschen dazu haben: welche Neigung wie das Haupt ist im silbernen oder goldenen Bilde, nämlich, das Anbeten. Denn sonst ist an dem Bilde keine Sünde, sondern ich mag daselbsten gebrauchen; wie droben Cap. 31. V. 34., die Götzen Labans, ihres Vaters, gebraucht hat, die sie unter die Streu der Kameele legte und sich darauf setzte. Aber das ist eine verdammliche Sünde, nämlich, das Vertrauen des Herzens, so auf das Bild gewendet wird, wiewohl das Bild selbst zu solcher Sünde Ursache gibt; aber die Lust zu solchem Bilde und der Götze ist erst im Herzen.

33. Wie das Volk vor dieser Zeit und jetzt noch etliche Könige zu den stummen Götzen laufen und denselbigen dienen, als, zur Maria zu Regensburg, oder da sonst anderer Heiligen Bilder sind. Das heißt in Wahrheit, Holz und Steine anbeten. Denn der Glaube, der dem einigen Gott gebühret als dem Schöpfer, wird da auf das silberne Bild gewendet und nicht auf den Schöpfer. Was könnte oder möchte aber närrischer genannt werden oder geschehen? Denn warum betest du nicht auf dieselbige Weise auch die Thäler an und das Geld, das du im Beutel hast? Darum sieht Jakob nicht allein auf die äußerlichen Bilder, sondern will haben, daß vor allen Dingen Gott recht soll geehrt werden und daß die Lehre rein sei; und will also die Reformation vom ersten Gebote anfangen.

\quellseite{242}

34. Der andere Theil, da er weiter sagt: „Und reiniget euch“, begreift in sich die Sünde wider die andere Tafel, und will haben, daß man sich darin auch bessern soll. Denn erstlich, sagt er, die Götter, von denen ihr meint, daß sie wohl bekannt sind, sind sehr unbekannt; darum thut sie von euch. Darnach, so thut auch Buße und steht ab von den Sünden, so wider den Nächsten geschehen. Es sollen unter euch keine Diebe, Hurere oder Ehebrecher sein, auf daß der Leib und das Herz rechtschaffen, keusch, rein und unschuldig sei von solchen Sünden, damit der Nächste verletzet wird. Das ist die äußerliche Gerechtigkeit.

35. Das dritte Stück, von der Veränderung der Kleider, ist gar ein äußerlich Ding und belanget die Ceremonien: Leget die unsauberen Kleider ab, zieret euch auch mit einem äußerlichen Schmuck. Denn solche Ceremonien, wiewohl sie niemand gerecht machen, sind sie doch nöthig. Denn es soll in äußerlichen Ceremonien, Gebärden und Sitten auch Ehrbietung gehalten werden, daß man fein ehrlich und züchtig zusammenkomme an dem Orte, da das Wort gelehret wird, da man betet und Gott anruft, und da sonst der andere Gottesdienst geschiehet. Vor allen Dingen soll zum Ersten das Herz dermaßen geschickt sein, daß es rechten Glauben habe zu Gott und mit ihm versöhnet sei, darnach Liebe habe zum Nächsten, welche vollkommen und ohne alles Aergerniß oder Anstoß sein soll. Zum Dritten, wo man zusammenkommt, soll man äußerliche Zucht halten, auf daß Alles ehrlich und ordentlich zugehe, 1 Cor. 14, 40. Die aber solches verachten und versäumen, zeigen damit an, daß sie nichts glauben und daß sie Gott und seine Kirche verachten. Denn wir sollen uns ja in der Kirche, da man zum Gottesdienste zusammenkommt, nicht also halten, wie in der Schenke; sondern dazu gehört Ernst und rechter Anstand.

36. Das ist nun eine schöne Reformation, da man erstlich die Lehre reformirt und reinigt; darnach werden auch die Sünden ausgefegt, als Wucher, Diebstahl, Räuberei und Unzucht. Zuletzt geschiehet auch eine Ermahnung, daß sich das Volk und die Priester im Tempel fein ehrlich kleiden sollen. Unsere Reformatoren kehren die Ordnung um; denn sie fangen ihre Reformation an von Schuhen und Kleidern, und was das Vornehmste und Beste ist, das lassen sie liegen. Man sollte es viel lieber also machen, daß das Herz innerlich durch den Glauben rein werden möchte, äußerlich aber der Leib und die Glieder durch die Liebe und äußerlichen Ceremonien auch in Kleidung und Sitten.

37. Endlich soll man auch in dem Gebote, so Jakob von Gott gegeben wird, die Art und Weise zu reden merken. Denn wir sollen uns des befleißigen, wenn wir die Bibel lesen, daß wir unsere Lehre von den drei Personen in der Gottheit aus dem hebräischen Texte vertheidigen und beweisen mögen. Denn also lauten die Worte aus dem Hebräischen: Die Götter sprach: Mache Gott einen Altar. Diese Art zu reden ist fleißig zu merken, wenn Gott von Gott redet. Denn wiewohl die Juden diese Art zu reden verachten und vernichten, und sagen, es sind hier ein Wechsel der Personen statt, der in der heiligen Schrift üblich ist; wovon zwar wir auch selbst bekennen, daß es wahr sei: so halten wir uns doch nicht mit den jüdischen Tifteleien auf. Denn wir können sie nicht bekehren, wie wir andere halsstarrige Menschen auch nicht bekehren werden. Und dieweil solche Art zu reden mit dem Neuen Testamente übereinstimmt, so behalten wir an diesem Orte die Worte, wie sie an sich lauten.

38. Wo nun Gott oder die Götter von Gott reden, da werden die drei Personen ausgedrückt und einiger Gott. Unter dem Gott aber, dem Jakob einen Altar machen soll, verstehen wir den Sohn Gottes, welcher der einige und wahre Gott ist mit dem Vater und dem Heiligen Geist; wie wir droben von dem Manne gesagt haben, der mit Jakob gerungen hat. Dieses unterweiset und stärkt uns sehr herrlich. Denn da Mose dieses geschrieben hat, ist er nicht trunken gewesen, und der Heilige Geist, der durch ihn geredet hat, auch nicht. Und es rede gleich ein Engel oder Mose mit uns, so ist es Gottes Wort. Darum will nun der Herr so viel sagen: Wenn du den Gott ehren wirst, der dir zu Bethel erschienen ist, so wirst du den Vater und den Sohn und den Heiligen Geist ehren. Denn wo gesagt wird: Elohim, Götter, sprach, da wird darunter verstanden, daß in der Gottheit mehr denn Eine Person sei.

\begin{quote}
\textbf{V. 3. Und lasset uns auf sein, und gen Bethel ziehen, daß ich daselbst einen Altar mache dem Gott, der mich erhöret hat zur Zeit meiner Trübsal, und ist mit mir gewesen auf dem Wege, den ich gezogen bin.}
\end{quote}

39. Es ist fleißig zu merken, daß in der Historie so vieler Patriarchen nicht beschrieben oder gerühmt wird, daß sie irgend ein Thier geschlachtet hätten zum Opfer, so doch die heilige Schrift zum öftern sagt von Abraham, Isaak und Jakob, wie sie Altäre gebaut haben. Von Noah, Adam, Kain und Abel lesen wir, daß sie Opfer gebracht, und zwar Brandopfer geopfert haben; welches ohne Zweifel die Väter gethan haben nach der Weise, die Mose hernach beschrieben hat, also daß sie Schafe, Ochsen und Ziegen geschlachtet und gekocht, und ein Theil davon verbrannt haben dem Herrn zum Brandopfer, ein Theil aber haben sie davon gegessen. Denn wozu könnten sie sonst die Altäre gebraucht haben? Wo die Schrift derselben gedenkt, da ist kein Zweifel, es werden Brandopfer gewesen sein, und Opfer, da sie Kuchen, Semmelmehl, Weihrauch und Anderes mehr geopfert haben, welche Opfer 3 Mos. 1. 2. 3. ff. beschrieben werden. Dieses ist aber sonderlich wohl zu merken, daß diese Patriarchen auch Altäre aufgerichtet haben, und Mose gedenkt doch gleichwohl gar keiner Opfer. Er läßt es aber dem Leser befohlen sein, daß er daraus schließen mag, daß daselbst Opfer, Brandopfer und Hebopfer gewesen, wo Altäre aufgerichtet worden sind.

40. Wiederum aber, wenn Mose vom Opfer und Brandopfer der Leviten handelt, so beschreibt er solche Opfer der Thiere, die sie geschlachtet haben, so fleißig, daß er auch die Haare und Haut, und den Fuß, so von den Thieren kommt, nicht übergeht oder verschweiget; da wirst du nichts finden, denn nur von eitel Räuchwerk und Schlachten. Er sagt nicht, was Gott geredet habe, was die gelitten oder gethan, die geopfert haben; dagegen in den Historien von den Vätern gar nichts gesagt wird vom Schlachten der Ochsen oder Schafe, welchen Theil sie behalten und welchen sie verbrannt haben: sondern das vornehmste Stück wird dazu gesetzt, nämlich, die Predigt, so die Väter bei den Altären gethan haben. Mose kocht, brät, siedet. In diesen Legenden aber wird nichts gesetzt denn nur von dem Altar, und wird weder Weihrauchs, Weins, noch sonst eines Opfers gedacht.

41. So will nun Jakob einen Altar bauen. Wozu denn? Antwort: Er will ihn Gott bauen, „der mich erhöret hat“, sagt er, „zur Zeit meiner Trübsal“ etc. Ja, das sind die rechten Altäre, gegen welche die Opfer der Ochsen und Böcke nur Schatten und äußerliche Zeichen sind. Aber Predigen, das Wort hören, Gott danken und beten, das sind die rechten Opfer; dieselbigen beschreibt Mose an den Vätern: die äußerlichen Opfer aber rührt er kaum mit einem Wort an. Ich will daselbst rühmen, spricht Jakob, wie Gott mit mir gewesen sei auf dem Wege, den ich gezogen bin, wie er mich erhöret habe. Derhalben ist Opfern in Wahrheit Anderes nichts, denn Gott danken und loben, daß er unser Gott ist, der uns nach seiner Verheißung erhören wolle, und bei uns sein in aller Trübsal und uns aus dem Tode erretten. Denn das heißt, daß er unser Gott ist, und daß man ihm danke, daß er uns bisher versorget und bewahret hat, und hinfort auch versorgen will; wiewohl er uns läßt versucht werden schier über unser Vermögen.

42. Eben dergestalt und zu solchem Ende werden die Kirchen recht gebaut. Wenn aber diese Opfer nicht da sind, so mögen es billiger Kuh- oder Sauställe genannt werden denn Altäre oder Kirchen. Denn dazu ist es verordnet, daß man in der Kirche zusammenkommt, daß wir darin die Leute von Gott also lehren und berichten sollen, daß er unser Gott sei, der uns läßt versucht werden uns zu gute. Das ist die Predigt vom Glauben, von unserer Geduld, von Gottes Gnade, und wie er die Seinen führt und regiert. Wo diese Lehre nicht gehört wird, da sollst du es dafür nicht halten, daß eine Kirche oder Altar sei. Wie im Abthum die Kirchen nur die Theater des Teufels gewesen sind, der die Seelen verspottet und jämmerlich erwürget hat.

43. Das Schlachten aber und die Opfer der Leviten hat Gott verordnet auf eine bestimmte Zeit zum äußerlichen Zeichen, auf daß das große Volk durch einen gewissen Schein zum rechten Gottesdienst und Erkenntniß Gottes möchte gezogen werden. Wie er uns die Taufe, die Schlüssel, Absolution und das Abendmahl gegeben hat nicht um des äußerlichen Werkes willen an sich, wie sich die Papisten träumen lassen: sondern daß wir dabei der Wohlthaten Christi gedenken sollen, der da sagt Luc. 22, 19.: „Das thut zu meinem Gedächtniß“, daß unter euch die himmlische Lehre gehöret werde; daß ihr mich anrufet, mir danket; daß ihr auf mich hoffet, und geduldig seiet, das Kreuz zu tragen, bis ich kommen werde und euch von allem Uebel erlöse. Dieses treiben und lehren wir fleißig, öffentlich in der Gemeinde und sonderlich bei den Kranken. Wir schlachten keine Ochsen oder Schafe, sondern loben Gott.

\quellseite{243}

44. Denn also sagt hier Jakob, daß er nicht darum einen Altar aufrichten wolle, daß er darauf opfern wolle; wiewohl dasselbe auch damit begriffen wird: sondern daß er Gott preisen und loben möge, der ihn erhöret hat zur Zeit der Trübsal, nämlich, da er vor seinem Bruder Esau hat fliehen müssen. Er erkennt, daß ihn Gott erhöret habe; darum hat er immer für und für angehalten mit Beten, Suchen, Klopfen, daß er Gott angerufen hat. Esau und der Satan waren seine Meister, die ihn getrieben und gezwungen haben, zu Gott zu schreien und ihn um Hülfe anzurufen. Denn Gott will haben, daß wir ihn anrufen sollen, und will uns erhören und unser Gott sein.

45. Also haben wir gehört, daß Jakob oftmals in großer Beschwerung und Angst gewesen ist und daß er in seinem Schrecken etlichermaßen Gott verloren hat. Wie er ein sehr elender Mensch war, da er nach Mesopotamien hat fliehen müssen. Darnach ist er bei Laban unbilliger Weise mit unendlichem Unglück betrübt worden, damit er hat freiten müssen. Endlich, da er wieder aus Mesopotamien gezogen, hat er mit dem Sohne Gottes kämpfen müssen. Derhalben ist sein ganzes Leben elend und armselig gewesen; und das meint er, da er sagt: „zur Zeit meiner Trübsal“. Denn er hat fast bei dreißig Jahren immer ein Unglück über das andere tragen müssen.

46. Scheint es aber nun auch da, als ob Gott gnädig und barmherzig sei? Ja, er ist wahrlich ein solcher Gott. Denn Jakob setzt hier im Text dazu, daß er das Gebet erhöre. Wie im 66. Psalm David von Gott rühmt, da er V. 16. 17. 19. 20. sagt: „Kommt her, höret zu alle, die ihr Gott fürchtet; ich will erzählen, was er an meiner Seele gethan hat. Zu ihm rief ich mit meinem Munde, und priesete ihn mit meiner Zunge. Darum erhöret mich Gott, und merkt auf mein Flehen. Gelobt sei Gott, der mein Gebet nicht verwirft, noch seine Güte von mir wendet.“ Er ist ganz geneigt, uns zu erhören, wenn wir nur auch geneigt wären, zu ihm zu schreien. Mit der Erhörung ist es gewiß; denn er selbst hat gesagt, er selbst hat uns befohlen, daß wir bitten sollen, da er sagt Joh. 16, 23. 24.: „Bittet, so werdet ihr nehmen. Wahrlich, wahrlich, ich sage euch, so ihr den Vater etwas bitten werdet in meinem Namen, so wird ers euch geben.“ Item Joh. 15, 7. spricht Christus: „So ihr in mir bleibet, und meine Worte in euch bleiben, werdet ihr bitten, was ihr wollt, und es wird euch widerfahren.“ Derhalben sind wir billig zu strafen, daß wir nicht bitten, nicht schreien, nicht anhalten. Das Exempel Jakobs aber soll uns dazu locken und ermahnen, der da Gott rühmt und preist nicht seiner Werke oder Opfer halben, welche gleichwohl auch geschehen sind, daß er sich damit dankbar erzeigt, daß ihn Gott errettet hatte; sondern er rühmt vornehmlich, daß er einen solchen Gott habe, der ihm seine Verheißung gegeben, ihn erhört, geschützt und erlöst habe aus allem Unglück, darin es bisweilen schien, als ob es gar verloren wäre.

47. Das sollen wir fleißig lernen und merken. Und das Exempel unserer ersten Eltern lehrt uns eben dasselbe auch. Denn da ihnen Gott nach dem Falle gerufen hat, da er sagt 1 Mos. 3, 9.: „Adam, wo bist du?“ waren sie in Angst des ewigen Todes. Darnach aber, da Gott V. 15. spricht: „Der Same des Weibes wird der Schlange den Kopf zertreten“, und anzeigt, daß ihr Geschlecht bleiben solle und daß er das menschliche Geschlecht erhalten wolle; zudem daß auch ein solcher Same solle geboren werden, der der Schlange den Kopf zertreten soll: da werden sie wiederum aufgerichtet und erfreut, und erschrecken nicht, da sie die Worte hören V. 19.: „Im Schweiß deines Angesichts sollt du dein Brod essen“ etc.; item, V. 16.: „Du sollst mit Schmerzen Kinder gebären“ etc. Denn Eva wird gesagt haben: „Ich will die Schmerzen der Geburt und des Todes gerne tragen, wenn nur das schwere Gericht und Zorn Gottes mag aufhören; welchen Zorn wir auf uns geladen haben, da wir in den Apfel gebissen und Gottes Gebot übertreten haben: wenn aufstatt der ewigen Strafe eine leibliche kommt, so will ich sie mit aller Geduld tragen.

48. Also ist unser ganzes Leben wahrlich ein recht elend Leben, und ein rechter saurer Apfel, darein wir beißen müssen. Aber es ist doch gleichwohl noch das rechte Leben dabei. Der andere Apfel der Eva war sehr süß, und war sehr lustig und lieblich dem Scheine nach anzusehen; er brachte aber mit sich den ewigen Tod. Darum ist es besser, Jammer und Elend leiden mit Hoffnung ewiger Erlösung, denn vor solchem Elende und Trübsal fliehen und in die ewige Verdammniß hinein rennen.

49. Wir sehen, daß die Propheten und Patriarchen auch eben dasselbe gelitten haben und daß sie sich des Leidens rühmen; wie das die Psalmen und viele Sprüche der Apostel bezeugen, als, Röm. 5, 11. und an andern Orten mehr. Darum sollen wir nun in unserem Herzen darauf auch denken, daß wir uns vor keinem Unglück also fürchten, daß wir derhalben den Gottesdienst und das Gebet wollten liegen lassen. Denn dieweil ich vor dem ewigen Tode und Verdammniß sicher und davon erlöst bin, so lasse nur immer hergehen Türken und Tartaren, Papst und allerlei Unglück, ich will davor nicht verzagen, dieweil ich die Erstattung und den sehr seligen Wechsel habe, daß ich weiß, daß der Tod in das Leben, die Verdammniß in ewige Seligkeit und die Sünde in Gerechtigkeit verändert ist; denn ich weiß, daß den Gottfeligen Alles muß zum Besten dienen.

50. Es ist aber bemerkenswerth, daß wir so gar undankbar sind, daß wir nicht erkennen, noch Gott dafür danken, daß der Sohn Gottes den ewigen Tod in eine zeitliche Strafe verändert hat. Denn wir wollen gerne in diesem Leben immer unsere Lust haben und in Freuden leben, und erschrecken und fliehen vor einem jeglichen Kreuz, wie gering es auch sein mag. Nun thut aber dem Fleische vonnöthen, daß es getödtet werde, und solchen Gehorsam sollten wir im Kreuze mit Geduld leisten, auf daß wir Gott dankbar wären, dieweil wir wissen, daß Christus für unsere Sünde gestorben ist. Und er sagt selbst, daß sein Joch sanft und seine Last leicht sei, Matth. 11, 30. Und endlich, dieweil Jakob mit solcher großen Freude rühmt, daß ihn Gott aus leiblicher Trübsal errettet habe; wie viel mehr wird uns gebühren, daß wir Gott loben und ihm danken für die ewigen und zeitlichen Wohlthaten!

51. Er setzt aber zu der Erhörung das auch, daß er weiter sagt: Der mit mir gewesen ist „auf dem Wege, den ich gezogen bin“. Ja, er ist in Wahrheit mit ihm gewesen und hat solches mit dem Werke und Erlösung bezeugt; aber sehr heimlich und daß man es nicht gesehen hat. Wie Jeremia sagt in seinen Klageliedern am 3. Cap. V. 22.: „Die Güte des Herrn ist noch nicht aus.“ Denn soviel das äußerliche Ansehen und Schein belanget, war der Teufel, Laban und die Verzweiflung mit Jakob. Er sagt aber: Diese ganze dreißig Jahre über ist der Herr mit mir gewesen.

52. Also hat es sich nun eigentlich mit dem Opfer verhalten, welches bei den Vätern gebräuchlich gewesen ist, und hat auch zu der Zeit also gemehret, da sie das Gespränge mit dem Levitischen Opfer getrieben haben. Denn Mose hat nicht vornehmlich auf den Tempel und Altar gesehen, auch nicht auf die Schafe und Ochsen, die David und Salomo geopfert haben. Denn das sind nur die Sporen und Kniffe oder die Schalen vom rechten Gottesdienst, die Gott eine Zeitlang fordert und deren er gebraucht als äußerlicher Zeichen um unsers Willen, daß er damit unsere Sinne aus dem Schlaf erwecke.

53. Der rechte Gottesdienst aber ist im Neuen Testamente ausgedrückt, da der Herr Christus sagt Luc. 22, 19.: „Das thut zu meinem Gedächtniß“; item, Matth. 28, 19.: „Taufet sie in meinem Namen“. Das ist aber erschrecklich, daß der Papst aus der Messe ein Opfer und äußerlich Werk macht, und den rechten Kern verwirft, welcher ist, daß man Gott danken, ihn anrufen, hoffen und im Kreuz und Trübsal ihn auch rühmen soll. Die Kindle haben erhalten, nämlich, das Werk unserer Hände und das Opfer; so wir doch unser Fleisch und unser Blut hier nicht essen oder trinken.

\quellseite{244}

54. Der rechte Gottesdienst aber wird im 50. Psalm V. 14. beschrieben und lautet also: „Opfere Gott Dank und bezahle dem Höchsten deine Gelübde.“ Dazu sollen die äußerlichen Ceremonien dienen und sollen dich erwecken, Gott zu loben, der dich erhört, erlöst, regiert und dir hilft, auf daß du ihm vertraust, seine Hand duldest, daß du lernest getödtet werden und sterben, und in dem allen geduldig seyn, wie es gottseligen Menschen ziemet. Das ist das Opfer unseres Leibes, und das lebendig ist, davon St. Paulus sagt Röm. 12, 1.

\section[Zweiter Theil]{Zweiter Theil. Wie Jakobs Gemeinde die fremden Götter von sich thut, welche Jakob vergräbt; wie Jakob darauf sicher auszieht; wie er gen Bethel kommt und einen Altar aufrichtet}

\subsection{I}

\begin{quote}
\textbf{V. 4. Da gaben sie ihm alle fremde Götter, die unter ihren Händen waren, und ihre Ohrringe; und er vergrub sie unter eine Eiche, die neben Sichem stund.}
\end{quote}

55. Da siehst du, daß der Patriarch eine gehorsame Kirche oder Gemeinde gehabt und die Gottes Wort geglaubt hat. Denn sie thun alsbald die fremden Götter von sich und nehmen die Reformation ihres Lehrers an. Darum ist die ganze Gemeinde allenthalben wohl reformirt und fein angerichtet gewesen. Und Mose hat in seinen Büchern hernach viel aus der Weise und Ceremonien der Väter zusammen gelesen, wie Einer, der aus allen Büchern hin und wieder etwas lieset und dasselbe zusammen trägt. Unter andern Bildern aber sind auch die teraphim, Götzen, gewesen, die Rachel droben, 1 Mos. 31, 19., ihrem Vater gestohlen hat. Und die Götzen oder Bilder haben sie Götter genannt; nicht daß sie so thöricht oder unsinnig gewesen sind, daß sie Holz, Gold oder Silber für Gott gehalten oder angebetet hätten, sondern sie die silbernen oder goldenen Bilder an den Gott gebunden, der das Gebet erhört und den Gottesdienst ansehen sollte, der an dem Ort aus menschlichem Vorwitze errichtet war.

56. Wie denn das Papstthum solcher silbernen, goldenen, hölzernen und steinernen Götzen voll gewesen ist. Denn wiewohl wir gewußt haben, daß solche Götzen alle mit Menschenhänden geschnitzt und gegossen waren, so sind wir doch der Meinung vor den Bildern niedergefallen und haben sie angebetet, daß Gott auf dieses oder jenes Bild sehen würde. Und wir ließen uns dünken, daß St. Barbara, Anna und St. Christophel, Jeder auf seine Bilder sähe und unser Gebet erhörete. Dieses ist die Finsterniß in Egypten gewesen, und ganz und gar dieselbe Unsinnigkeit und Raserei, welche vor Zeiten bei den Heiden auch gewesen ist, da sie so viel unzählige Götter erdichtet haben; denn also haben sie Juno, Bacchus, Ceres, Priapus und Andere mehr geehrt.

57. Ich halte aber nicht dafür, daß in einem jeglichen Hause goldene Bilder gewesen sind, sondern hölzerne, dergleichen unsere Leute vor dieser Zeit in ihren Häusern gehabt haben. Es wird irgend ein hölzerner Götze gewesen sein, geschmückt mit silbernen und goldenen Kleinoden. Wie man von Dionysius lieset, der das Bild Jupiters geplündert, da er den goldenen und silbernen Schmuck davon genommen hat; wie man bei uns gebraucht goldene Borden, Spangen, Krallen, Paternoster etc. Und ich halte dafür, daß die Ohrringe auch ein solcher Schmuck gewesen sind, also daß man dieselben auf die Bilder ziehe, und nicht auf die, so die Bilder geehrt haben. Wiewohl Lyra sich dahin neigt, daß Mose den Schmuck oder den Hut und priesterliche Kleider verstehe, die sie gebraucht haben, wenn sie haben beten oder opfern wollen, entweder die Priester oder auch andere Leute. Denn solch Ding ist in allerlei Gottesdiensten gebraucht worden, beide die da falsch und recht gewesen sind. Jedoch däucht mich dieser Verstand wahrhafter zu sein, nämlich, daß die Ohrringe der Götzen Schmuck gewesen seien.

58. Das hebräische Wort nesem heißt ein kleines Mündlein, oder ein goldener Halbzirkel, wie der Mondo, wenn er sich gekrümmt hat, damit sie die Stirne von einer Seite zur andern geschmückt haben wie mit einer Krone. Wir nennen es jetziger Zeit ein Perlein oder eine goldene Borde und Haarband.

59. Aaron hat unter anderem priesterlichen Schmucke auch einen Ohrring, nesem, gehabt. Danach haben sie die Könige auch gebraucht und haben daraus die Kronen gemacht. Derhalben ist es zu der Zeit gebräuchlich gewesen, daß sie die Götzen oder Bilder an den Köpfen mit solchen Mündlein geschmückt haben um göttlicher Ehre willen, daß sie eine Gestalt und Schein hätten der Heiligkeit; wiewohl das auch sein mag, nach Lyras Meinung, daß die Priester solchen Schmuck auch noch mögen gebraucht haben.

60. Jakob aber nimmt und vergräbt dies alles auf einen Haufen. Und ist wahrlich eine treffliche Frömmigkeit und großer Gehorsam an diesen Leuten gewesen, die so willig und gern von sich gegeben haben, daß sie doch so theuer und werth gehalten haben. Der größere Theil ist geraubt worden, da sie die Sichemiter geplündert, welche Götzendiener gewesen sind. Er hat aber seine Hausgenossen und Fremde, die mit ihm waren, überredet mit der herrlichen Predigt, die er that, daß sie solche Dinge von sich sollten thun, und daß er also mit seiner Gemeinde, wenn dieselbe gereinigt wäre von der Abgötterei, gen Bethel ziehen möchte.

61. Warum vergräbt er aber die Götzen unter die Eiche, die golden und silbern, oder hölzern und mit Golde geschmückt gewesen sind, und warum hat er sie nicht lieber eingeschmolzen und das Geld etwa den Armen ausgetheilt, oder sonst zu etwas Anderem gebraucht? Denn wozu dient das Gold und Silber, wenn man es in die Erde vergraben hat? Antwort: Das hat er um des Mißbrauchs willen gethan; denn er wußte, daß Gott demselben sehr feind ist: und er selbst hat ganz und gar verachten wollen das Gut, welches übel gewonnen war, dieweil er durch Gottes Segen sonst genug hatte, das er unter die Armen hat können austheilen.

62. Mose gebietet 5 Mos. 7, 5. alle Bilder der heidnischen Götter zu verbrennen, und will nicht haben, daß das Volk Israel des Silbers oder Goldes, das daran ist, begehren soll oder es in ihr Haus bringen, auf daß sie nicht auch Gott ein Abscheu werden, wie diese Bilder waren. Daraus ist klarlich zu sehen, daß etliche Bilder hölzern gewesen sind, die sie mit Spangen, Hütlein und Mänteln, von Gold und Silber gemacht, geschmückt haben; denn das Gold verbrennt man nicht. Die Könige aber mögen vielleicht goldene und silberne Götzen gehabt haben. Und Mose hat solch Gebot von den Vätern, daß man die Bilder der Götzen einfach von sich thun und verbrennen solle als ein gottlos Ding, und daran Gott einen Greuel hat; wiewohl man das Gold oder Silber irgend wozu gebrauchen könnte. Es gar rein hat Gott zu der Zeit die erste Tafel haben wollen, daß er auch das Gold und Silber, davon die Götzen gemacht waren, die Heiligen nicht hat wollen gebrauchen lassen; wiewohl es an sich selber gute Creaturen Gottes sind. Das hat er aber darum nicht haben wollen, dieweil solche Creaturen mit abgöttischem falschem Gottesdienst befleckt gewesen sind, mit welchem falschen Gottesdienst Gott heftig erzürnt wird. Auf die Weise vergräbt Jakob die fremden Götter in die Erde, auf daß er damit bezeuge, daß er von der Sünde rein wäre, die da heißt Abgötterei. Und daselbst ist auch im Gesetze Moses also gehalten worden denn groben und unverständigen Volke zum Schrecken, auf daß es lernete, solchen Greuel zu fliehen und zu verfluchen.

63. Wir im Neuen Testament wissen, daß ein Götze nichts ist in der Welt, wie St. Paulus sagt 1 Cor. 8, 4. Darum haben die heiligen Märtyrer und Bischöfe die Götzen eingeschmolzen und zu dem rechten gottseligen Gebrauch verwendet, die Kirchendiener, oder aber sonst andere gottselige Menschen und Arme damit zu erhalten. Und daran haben sie recht gethan. Denn solche Gesetze von äußerlichen Ceremonien haben nun aufgehört, welche Gesetze das grobe unverständige Volk gehabt, dem solcher äußerlicher Zwang vonnöthen war, dieweil es den Geist nicht gehabt, sondern nur mit dem Gesetz ist regiert und mit äußerlicher Disciplin gezwungen worden, daß es nicht in Abgötterei fallen möchte. Denn also pflegen wir die Kinder auch äußerlich zu schrecken, daß sie sich der Sünden, oder was sonst ungebührlich ist, enthalten, nämlich, wenn wir ihnen gemalte Drachen oder sonst mißgestaltete scheußliche Ungeheuer zeigen, und mit denselbigen die Dinge, so schändlich sind, vergleichen, daß sie erschrecken lernen, wenn sie solcher Dinge auch nur gedenken, hören oder dieselben ansehen. Eben eine solche äußerliche Kinderzucht hat dieses Volk auch gehabt. Denn Mose hat nicht haben wollen, daß sie die Götzen oder Bilder begehren, oder daheim in ihr Haus tragen sollten, auf daß sie lernten, daß Gott solchen Götzen sehr feind wäre und daß er daran einen Greuel hätte.

\quellseite{245}

64. Die aber nun erwachsen und verständig sind, bedürfen dessen nicht; darum halten wir daselbstige Gesetz nicht und bedürfen desselben auch nicht. Denn was haben wir jetzt in unsern Kirchen, das nicht mit der allerschändlichsten Abgötterei der Päpste erworben sei? Und ist nirgend ein Fürst, Stadt oder Gemeinde, die ihre Kirchen- oder Schuldiener mit eigener oder neuer Besoldung ernähre oder unterhalte. Wiewohl von den Kirchengeräten auch etliche ehrlich erworben und zur Kirche gegeben sind, die guten Künste und Schulen zu erhalten. Was aber zu Messen, Vigilien und zum Heiligendienste gegeben ist, das ist Anderes nichts, denn Purgatorium und sehr schändliche Abgötterei. Solch Silber und Gold ist alles nesem, wie das hebräische Wort lautet, das ist, es sind Ohrringe, Spänglein der Götzen, so dazu gegeben sind, daß der Aberglaube und das gottlose Wesen oder Abgötterei möchte erhalten werden. Wir gebrauchen aber darinn der evangelischen Freiheit, daß wir solch Gold nicht verwerfen oder in die Erde vergraben, sondern wir gehen damit um, daß es zum göttlichen und rechten Gebrauche gewendet werde, und was bisher den abgöttischen Meßpfaffen gegeben worden ist, hinfort unter die frommen gottseeligen Lehrer und Diener des Evangeliums ausgetheilt werden möge. Und wollte GOtt, daß wir solches erlangen und erhalten könnten bei den Raubvögeln, die mit Gewalt und Unrecht die Kirchengüter an sich reißen, und den armen Pfarrherren und Schuldienern kaum so viel lassen, daß sie sich des Hungers erwehren können.

65. Im Churfürstenthum Sachsen werden davon die Pfarrern und Schulen erhalten; und haben wir erlaubt, daß was daselbst übrig ist, entweder der Fürst behalten mag, oder aber daß er es unter die armen Schüler soll austheilen. Wenn man nun fragt: Warum wir nicht auch die übel gewonnenen Güter der babylonischen Hure wegwerfen? so antworten wir darauf also: daß sie jetzt christlich und wohl angelegt werden, und daß wir es dahin gebracht haben, daß das Geld, so zuvor dem Teufel gedienet hat, hinfort nun Gott dienen muß. In Gesetz Mosis ist vormals gewesen, daß das Volk mit härterer Zucht hat müssen in Schranken gehalten werden, um der Gefahr willen, darein sie hätten gerathen mögen wegen der Heiden, so allernächst um sie her gewohnt haben, da damals solch groß Licht des Wortes nicht war, wie jetzt ist, und da sie diese Zeit der Gnade nicht gehabt, die von GOttes Gnade uns heutiges Tages erschienen ist. Darum binden uns die Gesetze der äußerlichen Ceremonien nicht.

\subsection{II}

\begin{quote}
\textbf{V. 5. Und sie zogen aus. Und es kam die Furcht Gottes über die Städte, die um sie her lagen, daß sie den Söhnen Jakobs nicht nachjagten.}
\end{quote}

66. Dieses ist eine glückselige und sichere Reise gewesen unter einem solchen Vorgänger und Hauptmanne, den Jakob hatte. Und es ist ein treffliches und merkwenrthes Exempel, welches anzeigt, daß die Gläubigen und die an GOttes Wort hangen, Schutz und Schirm haben von dem Schöpfer aller Dinge, wider die Pforten der Hölle, Matth. 16, 18. Denn der Prophet Elisa hat recht gesagt 2 Kön. 6, 16.: „Derer ist mehr, die bei uns sind, denn derer, die bei den Widersachern sind“; wenn wir es nur gewißlich dafür halten und fest genug glauben könnten. Denn solche Schwachheit des Glaubens und der Mißglaube macht die Herzen der Gottseligen sehr matt und verzagt. Es sind aber doch allezeit Etliche, die diesen Trost gleichwohl etlichermaßen, wie sie können und mögen, behalten. Jakob selbst ist in der Gefahr und Noth etwas schwächer gewesen im Glauben, wie wir etlichemal gehört haben, denn sonst jemand von seinen Hausgenossen hat sein können; und das unaussprechliche Seufzen in seinem Herzen hat gleichwohl Himmel und Erde erfüllt, und GOtt erweckt, ihm Trost und Hilfe zu erzeigen.

67. Dieses dient nun uns zur Lehre und zum Troste, daß wir es gewißlich dafür halten, daß wir nicht können verlassen werden, wir glauben gleich fest oder nicht fest. Denn die Schwachen werden auch geschützt; wie St. Paulus lehrt Röm. 14, 1., daß man den Schwachen im Glauben solle aufnehmen; wiewohl dieselbe Schwachheit in den Neulingen oder jungen Christen ist, die noch erst angefangen haben, die christliche Lehre zu fassen. Aber die Schwachheit der großen Männer, als, Jakobs und der andern Heiligen, ist nicht im Worte oder in der Lehre, sondern in der Versuchung, wenn es scheint, als ob sie den Glauben schier verloren haben. Das ist ein Glaube, der da kämpft, und ist keine Schwachheit der jungen Christen, die erst angefangen haben; und dieselbigen nimmt doch GOtt auch auf und schützt sie, wie dieses Exempel und viele andere unzählige mehr anzeigen und bezeugen.

68. Also hält aber GOtt das Leben der Heiligen in Maß und Schranken, auf daß wir nicht stolz werden, und daß wir unser Herz trauen nicht auf uns selbst stellen, sondern auf GOtt, der die Todten auferwecket, wie St. Paulus sagt 2 Cor. 1, 9. Denn das ist des Schäfers Werck, daß er die Armen reich macht, die Todten auferweckt, die Betrübten aufrichtet und tröstet, und die Hoffärtigen stürzt.

69. So hat er nun Jakob erstlich mit dem Worte gestärkt, da er ihn hat heißen gen Bethel ziehen, und das Zeichen des Schutzes und der göttlichen Hülfe dazu gethan hat, nämlich, daß Gebot, einen Altar aufzurichten. Denn dieweil er einen Altar hat bauen und opfern sollen, so hat er ja leben müssen. Als wollte der Herr sagen: Die umliegenden Völker sollen dich nicht vertilgen, sondern du sollst gen Bethel ziehen und daselbst opfern. Zum Dritten kommt auch die That dazu und die Erfahrung. Denn der Herr macht, daß sich die Nachbarn fürchten, die zuvor grimmig und zornig waren um des Mordes willen, den die Fremdlinge in der Nachbarschaft und vor ihren Augen angerichtet hatten. Diese Nachbarn, sage ich, hat der Herr so furchtsam gemacht, daß sie den Söhnen Jakobs und dem alten Vater mit seinen armen und wehrlosen Häuslein nicht haben nachjagen können. Denn der Herr ist der Vorgänger und oberste Hauptmann dieses Heeres, das Jakob führt, und dasselbe stärkt er mit seinem Worte und sagt zu ihnen also: Seid getrost und habt guten Muth; ich will für euch streiten ohne Wehre und ohne äußerliche Gewalt; ich will nur ihre Herzen erschrocken machen, daß sie euch nicht dürfen angreifen oder nachjagen.

70. Lieber, wer sollte mit einem solchen Fürsten oder Heerführer streiten dürfen, der erstlich den großen starken Muth überwindet und niederlegt? Wenn aber derselbe niedergelegt ist, so müssen alle Kräfte und Glieder am ganzen Leibe weich und matt werden. Denn die des Krieges erfahren sind, pflegen darauf Acht zu geben und daraus zu erkennen, was es für einen Ausgang mit der Schlacht gewinnen werde: wenn von beiden Theilen die Schlachtordnung wider einander gestellt und angerichtet ist, da können sie bald sehen, welcher Haufe gewinnen werde; denn am selbigen Haufen, sagen sie, spürt man sonderliche Fröhlichkeit, beide an Menschen und auch an Rossen. Dagegen aber bei dem Haufen, so überwunden werden soll, sagen sie, scheine Alles matt und gleichsam todt. Das sagen die Kriegsleute, die des Krieges wohl erfahren sind; und solches reimt sich auch mit den Sprüchen in der heiligen Schrift. Denn der Sieg kommt vom Himmel, wie im 144. Psalm V. 10. steht: „Der du den Königen Sieg gibst“; item im 33. Psalm V. 17.: „Rosse helfen nicht, und ihre große Stärke errettet nicht.“ Wenn GOtt will Sieg geben, so gibt er auch dabei einen fröhlichen und starken Muth. Wenn es GOtt gibt, so ist es gegeben. Wenn er aber dem Feinde den Muth nimmt oder bricht, so fällt ihm auch alsbald die Wehre aus den Händen, wie im 76. Psalm V. 6. 7. steht: „Von deinem Schelten, GOtt Jakobs, sinkt in Schlaf beide Roß und Wagen. Die Stolzen müssen beraubt werden und entschlafen, und alle Krieger müssen die Hände lassen sinken.“

71. Auf solche Weise tröstet GOtt Jakob mit Worten und Werken in seiner höchsten Anfechtung und Noth. Denn da er mit seinem Volk fortgezogen ist, sind die Nachbarn also erschrocken, daß niemand unter ihnen das Maul hat dürfen aufthun, ich will geschweigen, daß Einer sollte das Schwert haben zücken dürfen. Denn der Herr hat da die Macht der Bogen, Schild, Schwert und Streit zerbrochen, daß sie nicht haben fertig und muthig sein können zu streiten, da sie doch sonst sehr gewaltig und streitbar waren. Und solches ist nicht geschehen um Simeons und Levis willen, als sollten es die mit ihrer That oder Hinterlist zuwege gebracht haben; sondern aus Kraft der Verheißung und durch den Glauben des Patriarchen und der Seinen, wiewohl derselbe schwach war.

72. Darum ermahnt uns nun der Heilige Geist mit diesem Exempel, daß wir den Artikel von der Schöpfung recht und wohl lernen sollen, nämlich, daß Alles in GOttes Hand und Gewalt stehe, und daß wir uns gewöhnen und erwecken, unsern Schöpfer zu vertrauen, welches Vertrauen in uns noch recht schwach und gering ist. Denn wenn wir es fest dafür hielten, daß GOtt der Schöpfer ist, so würden wir wahrlich auch gemächlich glauben müssen, daß er in seinen Händen und Gewalt hat Himmel und Erde, und alles, was darin ist. Ja, daß noch mehr ist, wenn wir sehen würden, daß die Welt mit allen Elementen brechen und über einen Haufen fallen wollte, und uns schon auf dem Halse läge, so würden wir doch sagen: Du wirst nicht fallen, wenn du schon fällst, es sei denn, daß es GOtt haben wolle. Und wenn sie uns auf dem Kopf liegen würde, so würden wir sagen: Du wirst mir keinen Schaden thun und wirst mich auch nicht unterdrücken; oder wenn es ja GOtt also wohlgefällt, daß ich mit solcher deiner Last überfallen und unterdrückt werden soll, so geschehe im Namen GOttes, was dem Herrn wohlgefällt. „Meine Zeit stehet in deinen Händen“, Ps. 31, 16. Wo es ihm aber anders gefällt, so will ich dir, Himmel und Erde, mit aller Gewalt des Türken und Papstes und andern grimmigen Zorn der ganzen Welt Trotz bieten.

\quellseite{246}

73. Denn also war es mit dem Hause Jakobs bestellt, daß es gewißlich und gänzlich hätte müssen zerstöret werden. Denn es hat den Canaaniten an der Stärke und an guten Willen, dazu auch am Rath, an Wehre oder Kriegsrüstung und an Häusern nicht gemangelt. Was hat ihnen aber im Wege gelegen? Antwort: Der Herr sagt: Ich will es nicht haben. „Bis hieher sollst du kommen, und nicht weiter; hie sollen sich legen deine stolzen Wellen“, wie Hiob sagt im 38. Cap. V. 11., und wie im 65. Psalm V. 8. 9. steht: „Der du stillest das Brausen des Meeres, das Brausen seiner Wellen, und das Toben der Völker, daß sich entsetzen, die an denselben Enden wohnen, vor deinen Zeichen.“ Wir leben mitten im Meer; und wie das Gewässer der Sintfluth fünfzehn Ellen hoch über die allerhöchsten Berge gegangen ist, daß sie sind bedeckt worden, also ist jetzt noch das große Meer drei Ellen höher als die Erde. Warum überfüllt es uns aber nicht? Antwort: Darum, daß GOtt dem Meer sein Ziel gesetzt hat. Es sind zwar seine stolzen Wellen aufgeschwollen und stoßen sehr hart an das Ufer, als wollten sie jetzt bald ausbrechen und über das Land laufen: aber solche große Gewalt der stolzen Wellen wird von GOtt gezwungen und inne gehalten.

74. Eben dasselbe mag man auch sehen in allem grimmigen Zorn der Welt, so wider die Kirche tobt und wüthet. Der Türke ist wie ein ungestümes Meer und das von Wellen aufgeschlagen ist, und wenn wir denselben mit unsern Kräften zwingen und zurückhalten sollten, so wäre es schon vorlängst mit uns aus. Denn wir Deutschen liegen und schandern, sind schläfrig und trunken, und haben keine solche Fürsten oder Feldherren, die Weisheit, guten Rath und Muths haben solche große wichtige Sachen recht auszuführen könnten. Und wenn der Türke fortfahren würde, so hätte er schon vorlängst ganz Deutschland eingenommen. Denn es ist niemand, der uns schützen könnte, weder Kaiser, noch Könige, noch Fürsten. Allein GOtt streitet für uns; und der Türke wird Deutschland nicht überfallen, denn allein sofern es GOtt beschlossen hat und haben will. Sonst hätte er uns schon vorlängst ohne Mühe und Arbeit unterdrückt, dieweil wir so faul und schläfrig sind; denn es sind beide, die Fürsten und der Adel, gar verrückt und verdorben in fleischlicher Üppigkeit, Wohlust und Geiz. Daß wir aber beschützt und erhalten werden, das geschieht durch die Kraft und Güte GOttes, und denselben Schutz erlangt die Kirche mit dem Gebet und Glauben. Also wird der Kaiser im Niederlande nicht weiter kommen können, wird nicht weiter toben und wüthen, denn es von GOtt versehen ist. Dieses wird uns zum Trost vorgehalten und unsern Glauben damit zu erwecken, auf daß wir GOtt so viel brünstiger anrufen und ihm vertrauen sollen.

75. Und man sieht die Vorsehung und Regierung GOttes leuchten auch in den Historien der Heiden. Es hätte Hannibal Rom erobern können ohne alle Mühe und Beschwerung, da er die allerwehrhaftigsten Hauptleute und Heere der Römer geschlagen hatte: aber er ist von GOtt zurückgehalten worden. Und dieweil die Andern dasselbe nicht verstanden haben, so schrieen sie und sagten: Vincere scis, Hannibal, victoria nescis uti: Gewinnen kannst du wohl, Hannibal, aber du weißt des Sieges nicht zu gebrauchen. Aber er sollte es nicht höher bringen.

76. Der Kaiser Carl hat den König von Frankreich gefangen, da er in einer großen Schlacht ihn vor Pavia überwunden hatte, im Jahr 1525. Im andern Jahr hernach hat er Rom mit Gewalt erobert und das Kriegsvolk hat die Stadt ausgeplündert. Also hat er zwei der allergelugsten Monarchen in seiner Gewalt gehabt. Und hat an guter Gelegenheit und an Heereskraft auch nicht gemangelt, wider den Türken einen herrlichen Krieg zu führen, da man im Jahr 1532 auserlesen Kriegsvolk versammelt hatte aus allen Orten des Römischen Reichs: aber er hat solcher Gelegenheit nicht gebraucht und wird hinfort in künftigen Zeiten solche gute Gelegenheit vergeblich suchen. Denn GOtt ist es, „der den Fürsten den Muth nimmt, und schrecklich ist unter den Königen auf Erden“, wie im 76. Psalm V. 13. steht. Derselbe bricht den großen streitbaren Kriegsfürsten, wie Pyrrhus, Hannibal und Andere mehr gewesen sind, den Muth mit einem Wort oder Wink.

77. Das ist die rechte Erkenntniß und der rechte Glaube von der Schöpfung. Welche Lehre fleißig zu betrachten und zu üben ist, auf daß wir uns damit erwecken, GOtt anzurufen und ihm zu vertrauen. Denn wo man sich auf die Festungen und Wälle und Büchsen verlassen will, das ist ein gottlos und vergeblich Ding. Und wenn GOtt dermaleinst verhängen würde, daß der Türke Deutschland überfallen sollte, so werden uns die Wälle, die mit so großen Kosten und Arbeit gebaut sind, nicht schützen: und mich würde wahrlich alsdann nicht gelüsten, in dieser Stadt zu wohnen, ich wollte eher hinausfliehen. Wenn wir aber unsere Kniee beugen und zu unserm Schöpfer schreien, der kann feurige Mauern um uns machen, wie der 125. Psalm V. 2. sagt: „Um Jerusalem her sind Berge; und der Herr ist um sein Volk her, von nun an bis in Ewigkeit“; item Ps. 34. V. 8.: „Der Engel des Herrn lagert sich um die her, so ihn fürchten, und hilft ihnen aus.“ Ich bin den großen Wällen und Festungen feind; denn es ist Anderes nichts, denn daß man das Geld damit verschwendet, und ein Ruhm sehr großer Thorheit. Wir sollen uns viel lieber das lassen angelegen sein, daß wir es fest und gewißlich dafür halten, daß wir in der Hand und Gewalt unsers Schöpfers sind: und nicht wir allein, sondern auch unsere Feinde und die Teufel mit allen Pforten der Hölle. Denn wo der Teufel freie Macht hätte zu toben und zu wüthen, so würde unser keiner einen Augenblick lebendig bleiben.

78. Darum lasset uns lernen unser Vertrauen auf GOtt stellen, der unser Schöpfer und Schutz ist, und nicht auf die Festungen, Büchsen, Weisheit oder Gewalt; es ist lauter Thorheit. Der Türke und der Teufel brechen ohne alle Mühe und Arbeit hindurch, zerreißen und verderben das alles. Denn GOtt will nicht haben, daß sich die Seinen auf etwas Anderes verlassen sollen denn auf ihn allein, darum daß sie den Namen führen, daß sie solch Vertrauen haben auf GOtt, den Schöpfer, durch seinen Sohn, durch welchen er uns zu Gnaden aufgenommen und einen Bund mit uns gemacht hat; welcher Bund ist, daß wir ein solch Vertrauen haben, daß unser Leben allein an GOtt hängt wider alle listige Anschläge und Gewalt des Teufels und der Welt. Und wenn GOtt haben will, daß ich soll vertilgt werden, so ist es nicht vonnöthen, daß er dazu Kriegsleute schickt; wo aber nicht, Trotz allen Türken, Tod und Teufeln in der Hölle!

79. Derhalben sollen nun die, so göttliche Verheißung haben, gedenken, daß sie sich darauf verlassen, daß GOtt für uns sorgen wolle, daß er unser Schutz, unser Hirt und Vater sein wolle um seines Sohnes willen; wie Christus selbst sagt, Joh. 16, 27.: „Der Vater hat euch lieb, darum, daß ihr mich liebet“ etc.; z. B. 33.: „In der Welt habt ihr Angst; aber seid getrost, ich habe die Welt überwunden“; Cap. 15, 18.: „So euch die Welt hasset, so wisset, daß sie mich vor euch gehasset hat“, und ist doch auch euer Leben dieses schändlichen Kreuzes in meiner Hand und Gewalt. Warum erschrecken und fürchten wir uns denn also vor dem Drohen der Gewalt unsrer Feinde? Wir sollen uns lieber freuen im Herrn, der uns berufen hat durch seinen heiligen Beruf, und hat uns gelehret, daß wir ihm vertrauen, und in ihn überwinden und siegen sollen.

80. Dieses ist die Lehre aus diesem Text, darin Mose mit sonderlichem Fleiß angezeigt hat, daß Furcht über die Völker gekommen sei, die so gewaltig waren; unter welchen ohne Zweifel Viele zornig gewesen sind, und gesagt haben: Sollen wir denn diese große Schmach unserer Nachbarn ungestraft und ungerächt hingehen lassen? Soll dieser Haufe der losen Leute und Bettler also frei in unserer Nachbarschaft toben und wühlen, und ungestraft bleiben? Aber der Heerführer Jakobs ist der Herr der Heerscharen und des ganzen himmlischen Heeres; wie des Propheten Elisa Knabe 2 Kön. 6, 17. gesehen hat, daß der Berg voll feuriger Rosse und Wagen war um Elisa her. Und der Spruch im 91. Psalm V. 11. ist sehr lieblich, da David sagt: „Er hat seinen Engeln befohlen über dir, daß sie dich behüten auf allen deinen Wegen.“ Und wir wissen, daß dieses wahr ist, und wird so oft in der heiligen Schrift wiederholt, und glauben es doch nicht. So gar ein großer und erbärmlicher Jammer ist es um den Unglauben. Die Sache an sich ist ganz gewiß, die Worte sind wahrhaftig: und das Herz glaubt doch gleichwohl nicht.

\quellseite{247}

\subsection{III}

\begin{quote}
\textbf{V. 6. 7. Also kam Jakob gen Lus im Lande Canaan, die da Bethel heißt, sammt alle dem Volk, das mit ihm war, und baute daselbst einen Altar, und hieß die Stätte El Bethel darum, daß ihm daselbst Gott offenbaret war, da er flohe vor seinem Bruder.}
\end{quote}

81. Jakob ist frisch und gesund in Lus angekommen, sammt alle dem Volke, das mit ihm war. „Denn der Schutz Gottes und die Heere der Engel umher beschirmen sie, daß sie auch nicht eine Klaue verlieren. Die Nachbarn sind sehr hart erzürnt, der Raub von Weibern und Kindern ist weggeübt: aber die starken und streitbaren Helden dürfen nicht einmal das Maul aufthun. Warum aber das? Antwort: Dieweil es Gott dem HErrn also wohlgefällt.

82. Da sie aber fortgezogen sind, werden ohne Zweifel unterwegs die Mitwisser der Sichemiten, die erwürgt waren, mit Weinen und Heulen ihr Elend und Schmerzen und daß sie nun gar einsam wären, angezeiget haben. Da hat Jakob ihre Traurigkeit und Wehklagen mit Worten etwas gelindert, wie er nur gekonnt hat, und es sind deren Etliche hernach mit den Patriarchen verheirathet worden. Jakob hat seine Canaanitin zum Weibe genommen; aber seine Söhne und Hausgenossen werden deren Etliche zu Weibern genommen haben. Simeon aber, Levi, Juda und Ruben sind dazumal schon Ehemänner gewesen. Und mit solcher Hoffnung hat Jakob die armen Mitwitwen freundlich getröstet, nämlich, daß er sie seinen Hausgenossen und Anderen ehelich vertrauen wollte; wie sich denn die Könige Israel, so von Jakob hergekommen sind, gleichergestalt auch gnädig und freundlich erzeigt haben gegen die Fremden. Denn Salma hat hernach Rahab und Boas Ruth zum Weibe genommen, Matth. 1, 5. Und sind also diese Mitwitwen auch zum Erbe und Gemeinschaft der Kirche gekommen. Welches ja eine sehr schöne Erstattung gewesen ist des Unfalles, daß sie ihrer vorigen Männer beraubt waren und sonst viel leibliches Unglück erlitten hatten.

83. Diese Freundlichkeit ist an Jakob und an den Vätern billig zu rühmen, daß sie mit dem weiblichen Geschlechte so ehrlich und freundlich umgegangen sind, auf daß sie in ihrem Elende und Gefängniß Trost haben möchten. Denn Jakob wird sie also getröstet haben: Lieben Töchter, leidet diesen Schaden mit Geduld und seid zufrieden; denn es hat Gott also wohlgefallen. Und was ihr leiblich für Schaden gelitten habt, das soll euch fein erstattet werden mit andern, geistlichen und viel größern Wohlthaten. Denn Jakob ist nicht ein solcher grober und unmenschlicher Feind gewesen, wie die Türken und Spanier sind; sondern er hat ihnen allen Rathen und anzeigen wollen, wie ihnen möchte geholfen werden. Er ist voll Glaubens, Geduld und Freundlichkeit gewesen; und eben dieselbigen Tugenden haben auch geleuchtet an seinen Weibern, Lea und Rahel, welche die Mitwitwen freundlich angenommen, und sich beflissen haben, wie sie ihren Schmerz mit allerlei freundlichem Dienst und Gottseligkeit lindern und sie mit ihnen zur Gemeinschaft aller Wohlthaten und der ewigen Seligkeit bringen möchten.

84. Lus ist eine Stadt, die von Alters her also ist genannt worden, wie wir im 28. Cap. V. 19. gesehen haben. Jakob ist der Erste, der den Namen verändert, und nennt sie Bethel. Droben, Cap. 28. V. 17., hat er geweissagt, daß daselbst gewißlich Gott wohnte; jetzt aber wird nun dieser Name von Gott bestätigt. Wie der Engel Jakob auch Israel genannt hat; aber derselbe Name ist auf ihm nicht geblieben: darum wird er von Gott wiederholt, auf daß er für und für auf ihm bleiben und beständig sein möge. In den Büchern der Richter und Josuas steht die Historie, wie die Kinder Israel die Bürger und Einwohner der Stadt Lus vertrieben haben; darnach hat sie den Namen Bethel behalten.

85. Von der Frage aber: Ob Lus und Jerusalem eine Stadt gewesen sei, ist droben Cap. 28. geredet worden. Alle Historien der Könige zeigen an, daß es verschiedene Städte gewesen sind. Unter Abraham und Josua sind Bethel, Ai und Hai Städte gewesen, so nahe bei einander gelegen haben; Lus aber ist Bethel und wird in der ganzen Schrift durch und durch nahe bei Ai gesetzt. Von Jericho ist sie zwölf Meilen gelegen gewesen, wie Hieronymus anzeiget, und lag an dem Weg nach Jericho hinab. Dahin ist aber nun Jakob frisch und gesund gekommen, und beladen mit Gold, Silber und mit dem Raube der Sichemiten, und hat das Erbe desselben Landes behalten. Denn es folgt hernach, daß Joseph zu seinen Brüdern ist gesandt worden, welche die Heerde ihres Vaters geweidet haben zu Sichem.

86. Hier soll man wiederholen, was droben gesagt ist von dem Bau des Altars. Denn Mose übergeht andere Ceremonien und Opfer, die er sonst im dritten Buche weitläuftig beschreibt, und erzählt nur dies einige Stück, welches ist, daß man Gott lobt, die Schwachen tröstet, die Irrigen stärkt, die Unwissenden lehrt und nur einfach das Wort treibt. Denn daselbst wird die Sünde und der Tod zerstört und Gerechtigkeit und ewiges Leben wieder aufgerichtet; der Teufel wird unter die Füße getreten und Gott wird erhöht. Dieser Gottesdienst ist Gott sehr angenehm. Darnach kommt die Anrufung und das Gebet der Kirche auch dazu, damit wir bitten, daß wir in dieser Lehre und Glauben gestärkt und erhalten, und dazu auch von Gott in aller Trübsal geschützt werden mögen. Diese zwei Stücke sind der rechte, wahre und höchste Gottesdienst, nämlich, die Lehre und das Gebet. Nach denselben folgt dann weiter die Tödtung des Fleisches und daß der alte Sauerteig in unserm Fleische ausgefegt werden möge. Und Mose sieht vornehmlich auf das Wort und auf der Seelen Seligkeit. Es reimt sich aber damit nicht übel, daß nicht allein die Hausgenossen Jakobs, sondern auch die Nachbarn dahingekommen seien, da der Altar gebaut war. Denn Gott ist nicht allein der Juden, sondern auch der Heiden Gott.

87. Der Name der Stätte el bethel ist ein neuer und sonderlicher Name. Denn das Wort el wird vorhergesetzt, welches auch im Namen bethel mitgesetzt ist. Wir haben aber nicht allenthalben die hebräischen Eigennamen behalten, sondern haben in unserer deutschen Uebersetzung bisweilen Gattungsworte daraus gemacht: wie an diesem Orte der lateinische Dolmetscher dieses im Lateinischen gegeben hat, domus Dei, Gottes Haus. Diesen Namen aber halten wir fest; denn er ist merkwürdig, weil el dazugesetzt ist. Da die Stätte zuvor Bethel genannt worden, nennt er sie „Gott Bethel“, und ist in der hebräischen Sprache sehr gemein, daß die Eigennamen zu den Namen Gottes gesetzt werden in einem Worte; als, Jesaia, Gabriel, Daniel, Hesekiel, Jeremiah, auf Lateinisch, exaltatio Dei, die Erhöhung Gottes. In der deutschen Sprache haben wir deren wenig. In der griechischen Sprache aber sind solcher Namen viele, als, Theodosius, Theodorus etc.

88. Nun kann aber diese Stätte nicht Gott sein, und wird doch also genannt el bethel. Darum disputiren wir mit den Juden von solchem Namen, um des Textes willen im Propheten Jeremia am 23. Cap. V. 6. und Cap. 33. V. 16., da also steht: „Und man wird ihn nennen: Der HErr, der unsere Gerechtigkeit ist.“ Daraus wir bewähren und erhalten wollen, daß Davids Gemähch oder der Messias, von welchem, wie sie alle zugeben, dieser Text soll verstanden werden, wahrer Gott sein müsse, wiewohl Jeremia sagt, man werde ihn HErr nennen, und es wird daselbst das nomen tetragrammaton, wie es nennen von den vier Buchstaben, nämlich, Jehovah, gesetzt, welches allein Gott zugehöret. Darum halten wir fest darüber, daß Christus wahrer Gott sei, der uns gerecht macht, dieweil er den Namen hat, welcher allein Gott gebühret.

89. Dieser Text ist fest und stark genug auf unserer Seite. Wir können aber der unbeständigen Lästerung der Juden so genau nicht begegnen, daß sie nicht immer ihre Ausflucht suchen, wie sie uns entgehen mögen. Denn sie werfen uns dergleichen Zeugnisse vor, darin der Name Gottes auch andern Dingen zugeeignet wird: wie diese Stätte eh bethel genannt wird, da el ein Eigenname Gottes ist. Und droben im 33. Cap. V. 20. hat Jacob einen Altar gebaut, und hat denselben genannt, den starken Gott Israel. Nun folgt es aber, sagen sie, daraus nicht, daß man also schließen will: Das Ding hat den Namen Gottes, derhalben ist es Gott. Denn es ist ein Anderes, den Namen Gottes haben, und ein Anderes, Gott sein. Mose nennt 2 Mos. 17, 15. den Altar: Dominus vexillum meum etc.: Der Herr mein Panier, Triumph, Sieg und Erhöhung; so doch der Altar der Herr nicht ist, der Sieg und Triumph auch nicht. Solche Tempel legen sie uns vor und streiten damit wider den Text Jeremia.

\quellseite{248}

90. Unra läßt es sich sehr sauer werden und antwortet darauf also: Christus wird unsere Gerechtigkeit und unser Herr genannt in seiner Person; andere Dinge aber oder Stätten sind keine Person, können auch diesen Namen nicht haben, sondern bedeuten nur irgend eine göttliche Wirkung, welche durch solchen Namen eingeführt wird. Die Juden aber lassen sich das nicht bewegen, daß sie derhalben nicht sagen sollten, daß Christus ein bloßer Mensch sei. Und ich habe zwar bei mir noch nicht beschlossen, wie man doch diese schlüpfrigen Male, die Juden, dererfassen und halten möchte. Wir haben zwar daran genug, daß dieses Zeugniß oder dieser Spruch mit dem Neuen Testament übereinstimmt, darin offenbar gesagt wird, daß Christus unsere Gerechtigkeit und auch wahrer Gott sei, bis wir sie aus andern Orten der Schrift überzeugen mögen. Es ist aber unmöglich, daß man den Secten oder Rotten alle Wege versperren sollte, daß sie uns nicht entwischen könnten; wie Christus selbst die Pharisäer und Sadducäer nicht hat überwinden können, da sie auch schon überwunden gewesen sind.

91. Wir wollen aber doch also antworten, daß man diesen Namen von Christo rühmen und ausbreiten soll, wenn er gleich mit dem Eigennamen nicht Jehovah genannt würde, wie die Juden lästern. Denn sein Name ist gewesen, Jesus von Nazareth, ein Sohn Matthias. Wir Christen aber legen die Worte Jeremias also aus: „Man wird ihn nennen“, das ist, man wird von Christo rühmen, man wird ihn preisen und ausbreiten, daß er der Herr sei, der uns gerecht macht. Die Juden bleiben nur an den Silben hangen, nämlich, daß er nicht also mit dem Namen sei genannt worden. Es hat aber mit diesem Namen eine andere Bewandtnis, als mit anderen Namen, wie Bethel und dem Altar, daß der unser Gott sei. Denn das ist der gemeine Name nicht, damit der Altar genannt wird, sondern es ist nur eine äußerliche Benennung und historische Erinnerung. Denn das ist niemals vom Altar gesagt oder gerühmt worden, daß er sollte unser Herr sein. Also sind auch die andern Namen, welche der Prophet Jesaja Christo zueignet, solche Namen, die von Christo gesagt werden. Als, da er im 9. Cap. V. 6. sagt: „Und er heißt Wunderbar, Rath, Ewig-Vater“ etc. Wenn das die Juden nicht annehmen, so mögen sie immer hinfahren mit ihrem Gespött, daß sie treiben. Dieses ist dennoch gleichwohl gewiß, daß ein Unterschied sei zwischen einem wirklichen Eigennamen, wie er den Christen in der Taufe gegeben wird, und zwischen einem auszeichnenden Beinamen.

92. Warum richtet aber Jacob den Altar auf? Antwort: Darum, daß sich Gott daselbst offenbart. Er will GOtt danken, daß er ihm das Wort gegeben, und will den GOtt, der sich ihm offenbart hat, rühmen, will sein gedenken, und ihn bekennen. Das heißt el bethel: Hie wohnet Gott; es sei gleich ein hölzern oder steinern Haus. Wenn man daselbst das Wort bekannt und predigt, oder daß Gott da gelobt und angerufen wird, alsdann ist es ein Tempel oder Kirche, wenn es auch schon unter dem Himmel oder unter einem Baume ist, als, unter der Eiche, darunter Jacob die fremden Götter vergraben hat. Und dieser Weise haben darnach die Götzendiener und gottlosen Pfaffen auch nachgefolgt. Denn sie haben die allerlustigsten Stätten auf den Bergen, Thälern, oder bei den Wassern und Bäumen erwählt, dahin sie das Volk gezogen haben, unter dem Himmel zu opfern nach der Weise der Väter.

93. Im Hebräischen aber wird ausdrücklich gesagt: Quoniam ibi revelati Dii: Darum, daß ihm daselbst die Götter offenbart waren. Und solche Stellen in der Schrift soll man fleißig merken, wenn die Schrift von Gott redet als von vielen Göttern. Wie droben im 1. Cap. V. 26. steht: „Laßt uns Menschen machen, ein Bild, das uns gleich sei“; item, im 3. Cap. V. 22.: „Siehe, Adam ist worden als unser einer“; item, im 11. Cap. V. 7.: „Wohlauf, lasset uns hernieder fahren, und ihre Sprache daselbst verwirren“; item, im 18. Cap. V. 2. 3.: Abraham sahe drei Männer gegen ihm, und sprach zu ihnen: „Herr, hab ich Gnade funden vor deinen Augen, so gehe nicht vor deinem Knecht über“ 2c. Das sind sehr wichtige und merkenswerthige Stellen. Denn sie zeigen an, daß die Väter eben dieselbe Erkenntniß und denselben Glauben von der Gottheit gehabt haben, den wir heutiges Tages auch lehren. Und ist zumal ein ungeschickt Ding, wenn man uns vorwirft, daß im Alten Testament nirgend die Lehre von der Dreifaltigkeit verkündigt sei, so es doch aus so vielen Sprüchen offenbar ist, daß die Gottseligsten damals auch geglaubt haben, daß in der Gottheit mehr denn Eine Person sei. Daß aber dieses Geheimniß eben nicht in allen Buchstaben ausgedrückt worden ist, sollen wir wissen, daß solches dem Neuen Testament hat müssen vorbehalten werden, darin die klarsten Offenbarungen sind.

94. Es ist aber bei den Vätern nicht ungebräuchlich gewesen, daß sie gelehrt haben, daß mehr Personen wären denn Eine in dem einigen göttlichen Wesen. Wiewohl der mehrere Theil diese Lehre nicht geglaubt hat; wie noch heutiges Tages die Türken und Juden uns darin nicht beifallen. Ja, daß noch mehr ist, unter denen, die unsere Predigt hören, sind auch Viele, die solche Lehre nicht achten noch glauben. Denn es widerfährt ihnen eben dasselbe, das Jesaia von seinen Zuhörern im 6. Cap. V. 10. sagt, daß ihr Herz verstockt sei, und ihre Ohren dicke, und ihre Augen verblendet, daß sie nicht sehen mit ihren Augen 2c. Denn die Menschen, so durch ihre Affecten und Begierden, oder aber durch eigenen Wahn bethört sind, sehen mit ihren Augen und hören mit ihren Ohren nicht, sie leben mit andern lebendigen Menschen, wie Klötze und Stöcke.

95. Es hören jetzt viel Leute unsere Predigt, lesen unsere Bücher, die wir schreiben, aber mit den Herzen sind sie unterdeß in der Küche oder sonst an andern Orten. Daher kommt es, daß wenn sie schon stehen und unsere Predigt hören, verstehen sie doch davon gar nichts; welches nicht ein natürlicher, sondern freiwilliger Unverstand ist, dadurch Einer sein Herz abwendet von den Dingen, so gegenwärtig vorhanden sind, und speculirt, oder denkt und begehrt derweil andere hohe Dinge. Solche Herzen nennt der Prophet Jesaia dicke oder verstockte Herzen, die gleichsam mit Schlamm überzogen sind, und Augen, die beschweret sind, wie die Augen der Jünger waren im Garten, die darauf keine Acht geben, was Christus redet. Und also ist der mehrere Theil der Menschen beschaffen, beide unter den Juden und auch unter den Christen: derer sind fast wenige, so auf die Lehre der Väter Acht geben oder ihrem Glauben nachfolgen.

96. Wir sollen aber bedenken, daß an diesem Ort klarlich gesagt wird, daß Jacob die Götter offenbar sind; denn im Hebräischen steht es in der Mehrzahl. Und auch das Zeitwort ist bedeutsam: Dii detecti, die aufgedeckten Götter; denn dasselbe Wort heißt offenbar oder aufgedeckt werden, und auswendig sich sehen lassen, wie Einer gestaltet sei; aufdecken, das zugedeckt gewesen ist. Also steht 3 Mos. 18, 7.: „Du sollst deines Vaters Scham nicht blößen“ 2c. Also ist Gott dem Jacob geoffenbart und aufgedeckt, ist ihm erschienen, wie er ist, Vater, Sohn und Heiliger Geist, in dem Sohn, der da sollte Mensch werden. Wie Jacob auf der Leiter drei Personen gesehen hat, und daß Gott würde Mensch werden von seinem Samen, der alle Geschlechter auf Erden segnen würde. Also wird gesagt, daß sich Götter geoffenbart haben, und ist doch Ein Gott. Denn wozu ist es vonnöthen gewesen, daß Mose dieses alles geschrieben hat, so er noch wohl andere Worte hätte gebrauchen können, wenn er nicht hätte wollen beschreiben, daß dieses große Geheimniß den Vätern geoffenbart gewesen sei?

97. Es hat aber Jacob diese Offenbarung in vielen schönen Predigten gerühmt. Denn darum ist der Altar aufgerichtet worden zum Gedächtniß solcher trefflichen Gabe, nämlich, daß ihm die Dreifaltigkeit offenbar war und die Menschenwerdung Christi, und daß ihm geistliche und leibliche Verheißungen zugleich gegeben waren. Und da werden auch Etliche gewesen sein, die solches schläfrig und mit dicken Ohren gehört haben. Der Patriarch aber hat ein groß Licht und Trost davon empfangen, und dieweil er der Gefahr entronnen ist, darein ihn seine Söhne geführt hatten, so lebt er nun ohne Sorge und in gutem Frieden und hält sich an die Verheißungen und Offenbarungen Gottes. Es wird aber bald ein anderer Unfall folgen. So gar ist nichts im Leben der Heiligen, das lange währe, und ist Alles voll mancherlei Streites; wiewohl die Erlösung auch wieder kommt.

\quellseite{249}

\section[Dritter Theil]{Dritter Theil. Von dem Tode der Deborah; wie Gott dem Jakob von neuem erscheint und Jakob den Ort der Erscheinung heiligt}

\subsection{I}

\begin{quote}
\textbf{V. 8. Da starb Deborah, der Rebekka Amme, und ward begraben unter Bethel, unter der Eiche; und ward genannt die Klageiche.}
\end{quote}

98. Diese Anfechtung ist etwas leichter und erträglicher, denn die, so hernach folgt, die er vom Tode der Rahel gehabt hat. Wozu dient es aber, daß man dieses alten Weibleins Tod hat schreiben und wissen müssen? Antwort: Es ist darum geschehen, daß wir wissen sollten, daß das Leben und alle Werke der Heiligen unseren Werken gleich sind. Sie sind wohl Leute gewesen voll Glaubens: sie haben sich aber bisweilen in der Anfechtung sehr schwach erzeigt, und uns zum Troste, daß wir in unserer Schwachheit nicht verzagen, dieweil sie eben mit demselben Fleische bekleidet gewesen sind, damit wir uns heutiges Tages auch noch schleppen. Sie haben wohl trefflichere Gaben gehabt denn wir: aber in der Gemeinschaft der Gnade und Hülfe, desgleichen auch des Schutzes Gottes sind wir einander gleich. Es wird aber Debbora, die Amme, hier eingeführt und beschrieben als ein heiliges Weib und die dem Patriarchen Jakob sehr lieb gewesen ist, darum daß sie vielleicht Alters und sonderlicher Gottseligkeit halben sehr geschickt gewesen, ihn zu stärken und zu trösten in seinen allerschwersten Anfechtungen. Darum trägt Jakob Leid über ihren Tod, und nennt die Eiche, darunter sie begraben wird, die Klageiche.

99. Und das ist auch eine Ursache gewesen, warum Deborahs Tod ist erzählt worden, nämlich, auf daß die heilige Schrift anzeigen möchte, daß die Patriarchen nicht Klöße und Stöcke, oder gar unfreundliche Leute gewesen sind, die keine natürliche Neigung gehabt hätten, wie die Mönche gewesen sind. Denn dieselben haben das für eine große Ehre gehalten, daß sie ihre Eltern, Brüder und Schwestern haben verachten können und haben des Spruches Christi dazu schändlich gemißbraucht, da er Luc. 14, 26. sagt: „So jemand zu mir kommt und hasset nicht seinen Vater, Mutter, der kann nicht mein Jünger sein.“ Daraus haben sie gedichtet eine sonderliche und evangelische Heiligkeit, wenn sie ihre Eltern hassen und verachten würden. Was ist aber das Anderes gewesen, denn sein eigen Fleisch und das ganze menschliche Geschlecht hassen? Wie an einem Ort erzählt wird von einem, der nicht hat ertragen können, daß er einen Menschen hätte ansehen sollen.

100. Dieselben Mönche haben also gelehrt: Wer Lust und Liebe hat mit den Menschen umzugehen, der wird mit den heiligen Engeln keine Gemeinschaft haben können. Denn das einsame Leben haben sie ein englisch Leben genannt; aber es ist in Wahrheit ein recht teuflisch Leben gewesen. Denn sie haben den Menschen die Neigung genommen, so Gott in die menschliche Natur geschaffen und gepflanzet hat. Gott will nicht haben, daß du aus eigenem Rath und Vornehmen, wenn dich dein Beruf nicht dazu treibt, Vater und Mutter verachten oder verlassen sollst (wie Hieronymus die ermahnt, so in die Klöster gehen wollen); denn solches ist ein selbsterwähltes Werk. Wenn dich aber irgend ein Tyrann zwingen wollte, die reine gesunde Lehre zu verleugnen, oder wenn dich deine Eltern dazu zwingen würden: alsdann heißt dich Gott dein Leben und Gut in Gefahr setzen und daß du lieber deine Eltern in solchem Falle verlassen sollst denn das Bekenntniß des Wortes. Da hat der Spruch Christi Luc. 14, 26. keine Statt, wenn man entweder Gott verleugnen oder aber seinen Vater verlassen soll; alsdann soll man zu Vater und Mutter sagen: „Ich kenne euch nicht“, 5 Mos. 33, 9., ich habe keinen Vater oder Mutter. Und dasselbe soll man thun nicht aus eigener Wahl, sondern wenn dich dein göttlicher Beruf dazu treibt.

101. Dieses sage ich um der Historien, oder, wie sie es genannt haben, um der Legenden willen, darin das Leben der Väter, so vor uns gelebt haben, beschrieben ist, darin solcher Unsinnigkeit und Ketzerei viel gefunden wird. Die heilige Schrift aber beschreibt die Heiligen also, daß es solche Leute sind, die bereit sind allerlei Gefahr zu leiden und alle ihr Gut zu verlassen. Wo aber solches die Noth nicht fordert, so sind sie voll Neigung und Liebe, die Gott in die Natur gepflanzet hat, als, daß sie ihr Gemahl, Kinder und Blutsfreunde lieben, und Mitleiden haben mit denen, die in Unfall gerathen und betrübt sind etc. Denn also hat Mose den Patriarchen Jakob abgemalt, daß er seine Weiber und Kinder sehr lieb gehabt, so doch nichts Fleischlicheres sein kann, denn solch Ding ist. Er schreibt, daß er geweint und Leid getragen habe um des Todes willen dieses heiligen Weibes, der Debora. Aber in solcher Neigung und Liebe ist sein Herz doch gleichwohl dem göttlichen Willen unterworfen gewesen und bereit zum Gehorsam in seinem Beruf.

102. Also werden andere Menschen auch heftig bewegt und bekümmert, wenn ihnen ihre Freunde und Verwandten, und noch viel mehr, wenn ihre Gemahlinnen oder Kinder mit dem Tode abgehen. Und solches ist keine Sünde, sondern sind gute natürliche Neigungen: wenn der Vater weint um seines Sohnes willen, der gestorben ist, wenn er Leid trägt, daß er sein Weib verloren hat, da ist er nicht zu strafen. Ja, das Laster, so dieser Neigung zuwider ist, nämlich, die Lieblosigkeit, wird von St. Paulo verdammt, welcher haben will, Röm. 12, 15., daß wir mit den Weinenden weinen, und mit den Fröhlichen uns freuen sollen, das ist, wir sollen uns anderer Leute Unglück oder Wohlfahrt bewegen lassen.

103. Wenn ich höre, daß der Türke greulich wüthet und der armen Christen Blut vergießt, daß er die kleinen Kinder an die Zaunstecken oder Pfähle steckt und mit den Weibern greuliche Schande treibt, so gebührt mir wahrlich, daß ich seufzen soll. Und ein menschlich Herz kann es nicht lassen, es muß heftig bewegt werden, wenn es solch greulich Ding hört. Hier möchte mir aber irgend ein Stoiker vorwerfen: Ei, du bist zumal ein zarter Märtyrer, weinst und heulst wie ein Weib: da gebührt einem Manne, der ein männlich Herz hat, daß er solch Ding verachte. Was geht es uns an, daß die Jungfrauen und auch die Weiber weggeführt, und daß sie geschändet und erwürgt werden? Darauf antworte ich also, daß solches vielmehr eine Unsinnigkeit und Raserei, und keine Stärke ist. Denn was kann gesagt oder gedacht werden, das aller menschlichen Freundlichkeit mehr zuwider, denn daß man sich anderer Leute Jammer und Unglück und das so schwer und greulich ist, nicht bewegen läßt? Es ist nicht gottselig oder christlich, daß man dazu lache, wenn andere Leute in Noth und Gefahr sind, als, wenn wir sehen, daß Stücke in Irrthum gerathen sind; oder daß der Kaiser und der Türke greulich toben wider die armen Christen; item, daß der Papst so viel tausend Seelen in die ewige Verdammniß stürzt: da gebührt mir, daß ich ein Mensch sei, und daß ich es dafür halte, daß ich auch menschliche Gefühle und Empfindungen haben soll, wie jener in der Komödie sagt. Viel mehr aber gebührt es mir, ein Christ zu sein, und mir solch Unglück anderer Leute, und daß so Viele an Leib und Seel verderben sollen, zu Herzen gehen zu lassen. Und Christus lehrt selbst solche Freundlichkeit, da er Luc. 6, 36. sagt: „Seid barmherzig, wie auch euer Vater barmherzig ist.“ Und St. Paulus sagt Röm. 12, 13.: „Nehmet euch der Heiligen Nothdurft an.“

104. Darum ist das unfreundliche Volk der Mönche billig zu verfluchen, welches gar ohne alle Gottseligkeit und Gemeinschaft des menschlichen Geschlechtes dahin lebt. Denn wo alle Menschen ihnen nachfolgen wollten, und in die Winkel kriechen, daß sie keine Menschen sehen möchten, wer wollte denn bald hernach der Kirche vorstehen, dazu das Welt- und Hausregiment verwalten? Gott hat uns in diese Welt berufen und gesetzt, und nicht außerhalb der Welt. Und Johannes der Täufer und Christus der Herr selbst werden aus der Wüste berufen zu lehren und zu predigen in der Welt. Dahin gehört auch nun die Beschreibung der heiligen Väter, darin angezeigt wird, daß es freundliche Leute gewesen sind, welche die Gemeinschaft der Menschen lieb gehabt und ganz freundliche Neigung zu den Ihren und auch zu Fremden getragen haben. Denn es ist ein gottlos und unmenschlich Ding, wo man mit andern Leuten nicht Mitleiden hat, und nicht mit ihnen weint, wenn sie elend und betrübt sind; oder wo man nicht verlangen hat nach denen, die aus diesem Leben ihren Abschied genommen haben, dieweil sie uns im Leben sehr lieb gewesen sind.

\quellseite{250}

105. Wir sehen, wie der Tod dieser frommen Matrone Debora dem Patriarchen Jakob so wehe gethan und wie er sich darüber betrübt hat; welche Debora, wie ich dafür halte, ihn bisweilen in seiner höchsten Angst und Noth mit Worten und Trost wird erquickt haben. Denn die Weiber haben ihre Gaben, und GOtt hat oftmals des weiblichen Geschlechts gebraucht, große Wunderwerke durch sie auszurichten, wie Agatha, Anastasia, Agnes, Lucia mit ihrem Glauben und Beständigkeit im Bekenntniß und in der Marter auch großer Männer Legenden verdunkelt haben. Der Heilige Geist hat ihre Herzen angezündet und gestärkt, daß sie der Marter und sehr greulichen Pein gelacht haben.

106. Also halte ich dafür, daß diese Debora eine weise und gottselige Matrone gewesen sei, die Jakob gedient und gerathen hat, das Gesinde im Hause zu regieren; und in großer Gefahr und Beschwerung hat er sich oftmals von ihr zureden und trösten lassen. Denn die Weiber, so die Gottseligkeit lieb haben, pflegen auch sonderliche Gnade zu haben, Andere zu trösten und ihnen ihre Schmerzen zu lindern; und das Gespräch der Weiber bewegt schier den Menschen mehr als das der Männer. Darum hat GOtt im Paradies dem Adam die Eva zugegeben, daß sie sollte um ihn sein, nicht allein die Hausarbeit auszurichten, sondern auch zum Trost in gemeinem Jammer und Unfall.

107. Debora aber ist sehr alt gewesen und das Gesinde hat sie gleichsam für die Großmutter im Hause gehalten. Denn sie ist der Rebekka Amme gewesen und ist derhalben auch selbst dem Jakob und den Kindeskindern sehr lieb gewesen; denn dieselben Kinder hat sie erzogen und ihrer ganz freundlich gewartet. Und die Großmütter pflegen größere Neigung zu den Kindeskindern zu haben denn die Mütter selbst, darum daß in solcher Neigung die Liebe immer herabsteigt und in solchem Niedersteigen gegen die gewöhnliche Regel immer brünstiger wird. Ihr Alter kann man leichtlich aus der Zeitrechnung abnehmen, dieweil sie der Rebekka Amme gewesen ist. Denn in diesem Jahr ist Jakob zum wenigsten hundert und sechs Jahr alt gewesen: seine Mutter, Rebekka, ist im dreißigsten Jahre ihres Alters dem Isaak zum Weibe gegeben worden, da er ein Vierzigjähriger gewesen ist. Derselbe hat im sechzigsten Jahre seines Alters Esau und Jakob gezeuget und im fünfzigsten Jahre der Rebekka. Wenn man dazu hundert und sechs Jahre thut von ihrer Geburt an bis auf diese Zeit, so kommen einhundert und sechs und fünfzig Jahre der Rebekka: darüber ist die Amme noch bei vierzig Jahren älter gewesen; darum ist sie damals ein hundert und sechs und achtzig oder ein hundert und neunzig Jahre alt gewesen. Nun ist sie aber um des Alters und anderer Gaben willen von der ganzen Gemeinde und allem Gesinde in Ehren gehalten worden. Denn sie hatte eine lange Zeit die sehr großen Patriarchen gesehen und gehört, Sem mehr denn hundert Jahre, mit welchem Jakob fünfzig und Rebekka ein hundert Jahre gelebt hat, die Amme aber ein hundert und dreißig Jahre. Eine solche Matrone hatte einen großen Schatz von Erkenntniß und Erfahrung. Darum haben sie dieselbe billig für ihre Großmutter gehalten.

108. Und ist wohl glaublich, daß auch Rebekka zu der Zeit oder kurz zuvor werde gestorben sein, da sie vielleicht von Bekümmerniß des Herzens und von Sorgen matt geworden ist; wiewohl ihres Todes nicht gedacht wird. Es scheint aber, daß da Jakobs Dienstjahre verlaufen gewesen, Debora gesandt worden sei, Jakob wieder aus Mesopotamien zu rufen; wie ihm Rebekka droben im 27. Cap. V. 45. versprochen hat, da sie sagt: „Ich will darnach schicken, und dich von dannen holen lassen“ etc. Denn sie war aus Syrien gekommen und der Rebekka von ihrem Bruder und Mutter zur Gefährtin mitgegeben worden, da sie dem Isaak sollte ehelich vertraut werden, wie droben Cap. 24, 59. steht. Derhalben ist sie einige Zeit bei Jakob geblieben, daß sie ihn und sein Gesinde trösten und erquicken möchte.

109. Da sie aber nun wieder heimziehen, ist sie unterwegs gestorben: und das Begräbniß unter der Eiche, so die Klageiche genannt worden ist, zeigt an, daß Jakob sammt dem ganzen Gesinde darüber groß Leid getragen haben und daß sie ihren Schmerz mit Thränen bezeuget haben. Welches der Heilige Geist darum erzählt, daß wir wissen sollen, daß es freundliche Leute gewesen sind; und daß es ein ungöttlich Ding sei, wo man sich anderer Leute Jammer und Elend nicht zu Herzen gehen läßt, darüber nicht Leid trägt und weint; denn man findet solche Neigung auch an den unvernünftigen Thieren. Derhalben nennt er die Stätte des Begräbnisses die Klageiche, zum Zeichen, daß sie die fromme Amme lieb und werth gehalten und Verlangen nach ihr gehabt haben; und daß er damit anzeige, daß er sammt dem ganzen Hausgesinde die sehr liebe Amme beweint und über ihr Absterben Leid getragen habe. Sie wollten sie wohl lieber wieder mit heimgeführt haben zum alten Vater Isaak; aber GOtt hat es anders haben wollen.

110. So lasset uns nun lernen, daß wir uns mit den Fröhlichen freuen und mit den Weinenden weinen sollen, daß wir nicht Stöcke und Klötze werden; wie man der Exempel viel hat in dem Buche, „Leben der Väter“ genannt, welche außerhalb des Berufs sich in die Wüste begeben haben von der Gemeinschaft der Menschen, und da sich mit den Dingen nicht bekümmert haben, damit sonst andere Menschen umgehen.

\subsection{II}

\begin{quote}
\textbf{V. 9. 10. Und GOtt erschien Jakob abermal, nachdem er aus Mesopotamien kommen war, und segnete ihn, und sprach zu ihm: Du heißest Jakob; aber du sollst nicht mehr Jakob heißen, sondern Israel sollst du heißen. Und also hieß man ihn Israel.}
\end{quote}

111. Es ist nun oft gesagt worden, daß GOtt die Heiligen zum öftern aufrichte und tröste, und sie über ihr Vermögen nicht versucht werden läßt. Darum wiederholt und bestätigt er jetzt nun die Verheißung, so Jakob droben, Cap. 28, V. 13. ff., ist gegeben worden, da er ihm erschienen ist im Traume auf der Leiter. Und das thut er darum, daß er ihm seinen Schmerz und Traurigkeit damit lindern möge, welche Traurigkeit er gehabt hat von dem Tode der Debora. Denn das Leben der Heiligen ist Anderes nichts, denn daß sie niederfahren zur Hölle und wieder herauf fahren; finstern Licht und Finsterniß, Anfechtung und Trost immer auf einander folgen. Es ist aber der Patriarch noch an allem Unglück nicht vorüber, sondern es ist noch viel größerer Jammer und Elend vorhanden, das er wird leiden müssen. Darum wird er nun von GOtt gestärkt, auf daß er dieses gegenwärtige Unglück und auch den künftigen Unfall mit so viel größerem Muthe ertragen möge.

112. Hier wird ferner mit klaren Worten gesagt, daß GOtt Jakob erschienen sei; im Anfange dieses Capitels aber nur, daß er gesagt habe: „Mach dich auf, und zeuch gen Bethel.“ Und es scheint, als sei zwischen diesen zwei ein Unterschied. Denn GOtt hat durch Isaak, Jakobs Vater, oder durch die Amme Debora reden können, von welcher Debora die heilige Schrift nicht vergeblich rühmt, daß der Patriarch Jakob sie beweint habe; und ist sehr wohl glaublich, daß sie eine sehr weise Matrone und Prophetin gewesen sei, voll Heiliges Geistes, die Jakob über viele Dinge Vorhalt gethan haben wird. Was aber die Heiligen reden, das soll man dermaßen verstehen, als habe es GOtt selbst geredet. Als, wenn wir das Evangelium predigen, wenn wir taufen, oder Diener zum Predigtamte berufen und ordiniren, so predigen, taufen und ordiniren wir nicht, sondern GOtt redet durch uns. Darum wird es auch GOttes Wort, GOttes Sacrament und GOttes Amt genannt; und man sagt recht also: GOtt redet, GOtt tauft, wenn er dasselbe durch die Diener thut; sintemal das alles GOtt zugeeignet wird, was die heiligen Menschen gesagt haben. Als zum Exempel, das sind Marci Worte, da er droben Cap. 2. V. 23. sagt: „Das ist doch Bein von meinen Beinen“; item V. 24.: „Darum wird ein Mann seinen Vater und seine Mutter verlassen, und an seinem Weibe hangen“: und Christus spricht gleichwohl Matth. 19, 5., daß GOtt solches geredet habe. Derhalben, wenn wir GOttes Wort reden, soll man das nicht als eines Menschen Wort annehmen.

113. Und es ist wohl glaublich, daß Jakob auf diese Weise gen Bethel zu ziehen befohlen worden, ohne, ohne Erscheinen, allein durch einen prophetischen Geist, entweder durch Isaak, seinen Vater, oder durch Debora, welche voll Heiliges Geistes ihm wird gerathen haben, daß er von Sichem fliehen solle. An diesem Ort aber sagt die heilige Schrift also: „GOtt erschien Jakob“, nämlich, Christus der Herr selbst in der Gestalt eines Engels. Das ist nun eine offenbare Erscheinung gewesen, und viel eine andere, als wenn er durch einen Menschen redet.

114. Darum, da nun Debora gestorben ist, die Jakob sehr wohl hat ermahnen und trösten können, da das Rüstzeug der Prophezeiung
\quellseite{251}
oder Weissagung verlöschen ist; da tritt GOtt an die Stelle, und verläßt Jakob nicht; sondern da Deborah verschwunden ist, erscheint er selbst, auf daß eine neue Erscheinung komme zu denen, die er droben gehabt hat, sonderlich auf der Leiter, dieweil er sieht, daß Jakob Leid trägt und weint, darum daß er die fromme Prophetin verloren hatte, und redet sehr freundlich mit ihm.

115. Der Segen aber soll verstanden werden vom geistlichen Segen, nämlich, vom künftigen Samen. Denn der leibliche Segen ist nur um des geistlichen willen gegeben worden. Wie wir in diesem ganzen Leben zu dem Ende leben, daß wir getauft werden, daß wir glauben, daß wir geheiliget und selig werden: das ist Gottes Ordnung in den Eltern, Lehrern und Kirchendienern. Denn Gott hat den Dienst aller Creaturen, auch der Engel, dazu geschaffen, daß sein Reich komme, sein Name geheiliget, und wir selig werden und das Erbe des ewigen Lebens erlangen sollen.

116. Derhalben wird vornehmlich der geistliche Segen darin begriffen, und Gott will haben, daß er darauf warten und sich daran halten solle. Als wollte er sagen: Sei getrost; denn ob sie schon alle sterben würden, so lebe ich doch noch, und wiederhole jetzt meine Verheißung dir, der du traurig bist und weinst, auf daß du es gewiß dafür haltest, daß von deinem Fleische Erben und der gesegnete Same kommen werde; wie ich Abraham und Isaak verheißen habe. Und auf daß Jakob ein gewiß Pfand und Zeugniß haben möge des gnädigen göttlichen Willens, verändert er ihm den Namen und nennt ihn Israel. Welche Veränderung des Namens anzeigt, daß Gott nicht von geringen, gemeinen, sondern von sehr großen und sonderlichen Dingen mit ihm handele. Das ist ein gemeiner Name, den er in der leiblichen Geburt durch seine Eltern bekommen hat; aber das ist sehr trefflich und voller Trost, daß er ihm nun einen neuen Namen gibt.

117. Wie wir in diesem Leben geboren werden Söhne und Töchter unserer leiblichen Eltern, und wie St. Paulus Eph. 2, 3. sagt: „Kinder des Zorns“; da wir mit uns bringen den Namen der sündlichen und verderbten Natur wegen der Erbsünde; wenn wir aber getauft werden, alsdann erlangen wir einen neuen Namen, und wird nun der, so getauft ist, nicht mehr ein Sohn Johannis oder Petri etc. genannt (wiewohl wir unter den Menschen in der Gemeinschaft dieses Lebens solcher Namen bedürfen), sondern vor Gott wird der Name geändert und aufgehoben, und wird der Mensch mit einem neuen Namen wegen der Taufe und des Glaubens ein Christ genannt. Der Name erinnert uns der Taufe, welche Taufe in täglichen Trübsalen und Anfechtungen geübt wird und ihre Frucht und Wirkung hervorbringen muß, auf daß wir wachsen und ein neuer vollkommener Mensch werden, und daß also der christliche Name vollendet werde, bis unser Name und alter Adam gar aufgehoben und abgeschafft werde.

118. Darum ist hier zu merken, daß Gott selbst dem Patriarchen Jakob diesen Namen gibt, mit welchem neuen Namen ein Unterschied gesetzt wird zwischen dem Fleisch und dem Geist. Denn Israel ist ein göttlicher Name, und Gott hat eine andere Weise Namen zu geben, denn es sonst bei den Eltern, Verwandten, Nachbarn, Begleichen in eines Jeden Vaterlande und nach eines Jeden Stande gebräuchlich ist; sondern der Name, damit uns Gott nennt und erkennt, ist ein sonderlicher Name. Wie er zu Mose spricht 2 Mos. 33, 12.: „Ich kenne dich mit Namen.“ Und dahin hat St. Paulus gesehen, da er Röm. 9, 6. sagt: „Es sind nicht alle Israeliter, die von Israel sind.“ Es sind nicht alle Israeliter, die vom Fleische herkommen und den Namen Israel führen, sondern die rechten wahren Israeliter sind geistlich und Gott hat ihnen einen geistlichen Namen gegeben; wie das ausführlich und herrlich ausgelegt wird Röm. 9.

119. Gott will aber so viel sagen: Ich habe dir bisher deinen Namen noch nicht geoffenbart, damit du bei mir genannt wirst, sondern ich habe es geschehen lassen, daß du mit deinem väterlichen Namen und nach dem Fleische bist genannt worden: jetzt aber gebe ich dir nun einen Namen nach dem, wie ich mit dir gehandelt habe, wie ich dich wunderbarlich geführt habe, nicht als einen fleischlichen Sohn, sondern wie ich dich geführt, regiert, betrübt, gesegnet und geheiliget habe durch vielerlei Plage, daß ich aus dir einen neuen Menschen und eine neue Creatur machen, und dir nun jetzt einen neuen Namen nicht des Fleischs, sondern des Geistes geben möchte.

120. Was die Ableitung des Wortes, ist droben erklärt worden. Das hebräische Wort sarah heißt streiten, oder vielmehr überwinden. Denselben Worte und Verstande ist der Name Jakob nicht sehr ungleich; denn Jakob heißt auch ein Untertreter, das ist, ein Ueberwinder. Hier aber sagt er, daß er Gott überwinde: Du wirst ein Ueberwinder Gottes sein. Und das ist ein gemeiner Name aller Christen; denn wir sind nicht allein Ueberwinder des Teufels, der Sünde, des Todes, der Menschen und dieses Lebens, sondern auch Gottes; denn er hat verheißen, daß er uns helfen wolle, ja, er hat sich uns gar ergeben. Wenn es ohne das wäre, so wären wir keine Ueberwinder Gottes; dieweil er aber unser Gott ist durch die Verheißung, und zu uns gesagt hat: „Ich will dein Gott sein“, sei du nur getrost etc.: daher kommt es, daß wir Gottes Ueberwinder werden.

121. Denn wenn er uns seine Verheißung gegeben hat, so stellt er sich darnach, als sei er uns gar feind, läßt uns versucht werden und verbirgt sein Angesicht gar vor uns, als hätte er aller seiner Verheißung vergessen. Da muß man denn bitten, suchen und anklopfen; und wenn wir uns dünken lassen, daß unser nichts geachtet werde und daß wir gar verlassen müssen, so muß man doch beständig bleiben, beharren und immer anhalten im Glauben mit Gebet und Geduld. Auf solche Weise überwinden wir Gott, wenn wir die Verheißung, oder Gott, der uns die Verheißung gegeben hat, nicht verlassen; und wenn wir also beharren im Gebet und Glauben, so wird Gott uns aus einem verborgenen Gott ein erscheinender Gott, der uns tröstet, und thut, was wir haben wollen, wie in 145. Psalm V. 19. geschrieben steht: „Der HErr thut, was die Gottesfürchtigen begehren“, item Joh. 15, 7.: „So ihr in mir bleibet, und meine Worte in euch bleiben, werdet ihr bitten, was ihr wollt, und es wird euch widerfahren.“ Das sind rechte Israelitische Stücke und damit Gott überwunden wird. Es gehört aber dazu, daß Christus dabei sagt: „So ihr in mir bleibet, und meine Worte in euch bleiben“, item Joh. 14, 23.: „Wer mich liebet, der wird mein Wort halten“; und wir werden zu ihm kommen, und Wohnung bei ihm machen.“ Denn so ist GOtt überwunden; und alsdann setzt er uns und ruft uns aus als Israeliter, wenn er uns um des Glaubens willen an seine Verheißung erhört und erhält.

122. Solche Leute sind aber die Israeliten nach dem Fleische nicht. Denn sie würden das nicht thun können, das zum Israel gehört, der darum ein Ueberwinder ist, dieweil er im Worte bleibt, und im Bitten, Suchen, Anklopfen und Anhalten beharret. Solche Beständigkeit und Sieg gibt Gott Jakob, und spricht also zu ihm: Du bist lange genug und sehr viel versucht und wohl geübt worden; ich habe mein Angesicht oftmals vor dir verborgen. Dieweil du aber dich so fest an die Verheißung gehalten hast, habe ich dir müssen weichen, dich erhören und dir helfen, und rufe nun jetzt aus und erkläre, daß du ein rechter Israel und Ueberwinder Gottes bist. Nicht daß er nicht lange zuvor auch ein solcher Israel gewesen sei; sondern er offenbart jetzt diese Ehre Jakobs ihm zum Troste. Denn droben, 1 Mos. 32, 28., ist prophetischer Weise zu ihm gesagt worden: „Du sollt nicht mehr Jakob heißen“ etc., aber hier wird solches nun erfüllt und geoffenbart, da ihm Gott selbst erscheint. Du bist, spricht er, bisher ein solcher Israel gewesen; du sollt aber nun auch also heißen.

123. Hier erhebt sich aber nun eine Frage: Warum der HErr sage, daß er hinfort nicht mehr soll Jakob heißen, so er doch von Mose, wie in diesem Capitel, und von allen Propheten ohne Unterschied Jakob und Israel genannt wird? Lyra disputirt ohne Ursache und vergeblich von dem wirklichen und beschaulichen Leben und sagt: Er werde Jakob genannt um des wirklichen Lebens willen, daß er die Affecten und Lust des Fleisches untertrete; Israel aber werde er genannt um des beschaulichen Lebens willen, darum daß er jetzt Gott sehe. Freilich führt er diejenige Ableitung des Wortes nicht an, welche Eusebius und die Grammatiker haben, nämlich, daß sie sagen, „Israel“ heiße so viel, als ein Mann, der Gott sieht; sondern er legt es aus, als heiße „Israel“ so viel, als auf Gott gerichtet, das ist, daß er mit Gott umgehe mit seinen lustigen Gedanken und Speculation und das künftige Leben beschauet, wofür er Gregorius anführt.

124. Aber das lasse ich fahren; denn ich halte, daß die Väter selbst nicht verstanden haben, was sie von zweierlei Leben gelehrt haben. Denn die Mönche haben würkliches Leben genannt, das darin besteht, daß der Mensch die bösen Lüste und Begierden im Zaum halte und zwinge, und äußerlich das Leben mit guten Sitten recht führe durch Fasten und Kasteiung des Fleisches, durch unsaubere Kleidung und dergleichen andere Dinge mehr; welches der Rottengeist Münzer „entrogen“ nennt, das ist, daß man sich die groben Sünden betreffend bessern sollte. Die Heiden haben es eine äußerliche Zucht genannt, die in diesem Leben nützlich und nötig ist, vornehmlich der Jugend, die Unterweisung und Lehre bedarf, daß sie nicht gar grob und unhöflich, und zu allen Werken dieses gemeinen Lebens ungeschickt werde.
\quellseite{252}

125. Aber nach der heiligen Schrift muß man das würkliche und beschäuliche Leben auf eine andere Weise beschreiben. Denn dasselbe ist vornehmlich jetzt ein solch Leben, da Einem die Augen übergehen. Es sind nicht so lustige Gedanken, Speculation und Trost, wie der Mönche Speculation gemein ist; sondern wer ein rechter Israel sein und GOtt überwinden will, der muß mit solcher Uebung umgehen, davon wir zuvor gesagt haben, nämlich, mit Bitten, Suchen, Anklopfen und daß er an dem Worte und Verheißung fest halte. Das ist das rechte beschäuliche Leben der Gottseeligen, da die Vernunft und unsere Einbildung oder Gedanken nichts vermögen, da Sinn und Verstand des Menschen ertödtet wird mit allen Kräften, und der Mensch allein vom Worte GOttes lebt, wie 5 Mos. 8, 3. geschrieben steht: „Er demüthigte dich, und ließ dich hungern, und speisete dich mit Manna, das du und deine Väter nie erkannt hatten; auf daß er dir kund thäte, daß der Mensch nicht lebe vom Brod allein, sondern von allem, das aus dem Munde des HErrn gehet.“

126. Das andere Speculieren, das die Mönche treiben, sind kalte, todte und gefährliche Gedanken. Denn sie kommen her aus eigenem Willen und Vernunft, ohne Wort und Verheißung. Darum soll man sich davor hüten, und wer gottseeliglich und recht von GOtt speculieren und Gedanken haben will, der soll es nicht thun ohne das Wort, sondern er soll den Sohn GOttes, der Mensch geworden ist, ergreifen, soll an der Krippe und Windelein anfangen, darein der Sohn gewickelt ist, bis daß er kommt an die Himmelfahrt. Also wird er es im Glauben empfangen, und wird im Worte und der Verheißung, so uns gegeben ist, bleiben. Ich kenne Etliche, die sich ihrer Speculationen rühmten, daß sie daran so große Lust gehabt hätten, daß sie sich dünken ließen, wie sie etlichermaßen entzückt und in den Himmel erhoben wären. Solche Speculationen habe ich nicht lieb, will mir deren auch nimmermehr wünschen; denn sie sind ohne Wort. Und die denselbigen folgen, sind Leute, die vom Gesetz und Evangelium nichts wissen.

127. Die rechte Speculation aber, die einen Israeliten macht, ist eine Tödtung aller Kräfte, die in unsern Sinnen und Vernunft sind, und daß man allein im Glauben und Hoffnung auf die Verheißung hangen bleibe. Wenn der Mensch in den letzten Zügen liegt und mit dem Tode ringt, da steht man, wozu solche Speculationen nütze sind; denn da fühlen wir bei uns gar keinen Rath oder Hülfe. Darum ist solche Speculation nöthig, das ist, daß man nur das lautere Wort ergreife und sich unsere Sinne und Fühlen nicht irren lasse, nach dem Spruche Christi Joh. 8, 51.: „So jemand mein Wort wird halten, der wird den Tod nicht sehen ewiglich.“

128. Aber das Gesetz oder die Gebote GOttes setzen und ordnen das würkliche Leben: nicht Satzungen der Mönche, sondern die Werke, so im äußerlichen Leben zu unserem Berufe gehören. Da soll ein Jeder betrachten, was sein Stand sei, den er führt, und soll darin seines Berufes warten, und sich fleißig üben in den Werken, die in den Zehen Geboten GOtt befohlen hat. Darnach, wenn die Anfechtungen kommen, wirst du ein Israelit werden, wenn du am Glauben auf die Verheißung fest hältst und dich darin übst. Wie du in den Zehen Geboten ein Jakob gewesen bist, also wirst du ein rechter Israel werden in den Verheißungen, die dich zum Gebet treiben werden, daß du GOtt anrufest. Deß haben wir ein trefflich Exempel an dem cananäischen Weiblein Matth. 15, 25—27., welches sehr fein geübt wird in dem beschäulichen Leben, und hält also an, daß sie sich auch mit seinen harten Worten zurücktreiben und abweisen läßt; sie klopft so lange an die Thüre, gleichsam mit Gewalt, bis Christus weichen muß und sie erhören, und ihren Glauben und Beständigkeit, daß sie so fein beharrt hatte, vor allem Volke rühmt und ausruft.

129. Auf solche Weise lehrt man recht vom würklichen und beschäulichen Leben, nämlich, wenn man es unterscheidet nach den Zehen Geboten und nach der Verheißung, welches die Mönche nicht verstanden haben. Denn die Zehen Gebote zeigen einem jeglichen Menschen die Werke seines Berufes; die Anfechtungen aber lehren uns nicht die Zehen Gebote verstehen, sondern daß wir GOtt mit dem Glauben ergreifen und uns seines Verdienstes oder Verheißung des Gesetzes rühmen, wie der Pharisäer thut Luc. 18, 11., und wie sie sagen Jes. 58, 3.: „Warum fasten wir, und du siehest es nicht an?“ Denn solcher Ruhm ist von GOtt verdammt. Außerhalb der Verheißungen mag ich dessen gedenken, was ich gethan oder gelitten habe in der Verfolgung und Trübsal: aber nicht nach dem Exempel der Mönche, die in den Werken ein Verdienst des ewigen Lebens erdichten.

130. Und dieses hat mich gut gedäucht aus den Vätern zu wiederholen, auf daß wir nicht dafür angesehen werden, als hätten wir ihre Schriften gar verworfen. Denn man soll solche Schriften nicht verachten, sondern mit rechtem Verstand lesen. Es scheint, daß Cyra und Gregorius nicht verstanden haben, was das würkliche oder beschäuliche Leben sei; denn sie setzen beiderlei Leben in die Ceremonien und Satzungen der Mönche: aber die Wirkung soll in den Zehen Geboten bleiben, die Contemplation oder Beschäulichkeit in der Verheißung. Jakob wird in dem Gesetz ein Untertreter und in den Verheißungen ein rechter Israel.

131. Aber nun will ich wieder auf die Frage kommen: Warum Mose oder die nachfolgenden Historien nicht allenthalben den Namen Israel behalten haben? Und antworte darauf also: GOtt redet von dem vornehmsten Namen, nämlich, von den Verheißungen, darin Jakob geübt war und noch weiter soll geübt werden. In derselben Verheißungen willen soll er vornehmlich und eigentlich ein Israel sein und auch also genannt werden. Er ist aber Jakob genannt worden nach gemeiner Gewohnheit der Menschen und um der Werke willen. Und das ist nicht der vornehmste Name: wie die Werke, welche die Gottseeligen thun nach dem Gesetz, auch die vornehmsten Werke nicht sind, wiewohl sie zu diesem Leben gehören, sondern sind Jakobs Werke. GOtt will beiderlei haben, daß wir glauben und auch gute Werke thun sollen, und daß wir also Israeliten und Jakobiten seien. Aber vom Glauben und Verheißung werden wir vornehmlich Israeliten genannt, und von den Werken, das vor der Obrigkeit und den Eltern gehorsam sind, daß wir ein ehrlich und züchtig Leben führen etc., werden wir Jakobiten genannt. Und dieweil solcher Gehorsam nicht vollkommen ist, macht er nicht gerecht, ist aber gleichwohl nöthig als eine Frucht des Glaubens: der Glaube aber ergreift die unbefleckte und unsträfliche Gerechtigkeit des Sohnes GOttes, welche rechte Israeliten macht. Dieses sei gesagt vom geistlichen Segen, welcher Jakob den Namen gibt vom Glauben, damit er diesen Segen ergreift. Was nun hernach folgt, gehört zum leiblichen Segen.

\begin{quote}
\textbf{V. 11. 12. Und GOtt sprach zu ihm: Ich bin der allmächtige GOtt, sei fruchtbar und mehret dich; Völker und Völkerhaufen sollen von dir kommen, und Könige sollen aus deinen Lenden kommen; und das Land, das ich Abraham und Isaac gegeben habe, will ich dir geben, und wills deinem Samen nach dir geben.}
\end{quote}

132. Die vorigen Verheißungen werden hier Jakob zum Trost wiederholt, und reimen sich fast alle Stücke insonderheit und ein jegliches Wort mit der vorigen Historie, droben, Cap. 17., da Abraham die Beschneidung gegeben wird, da sowohl ihm als auch der Sarah der Name verändert worden ist; und braucht daselbst GOtt auch denselben Namen schaddai, den er hier gebraucht hat; und die Auslegung desselben Namens magst du auch droben suchen. Warum er aber sage, daß Könige aus seinen Lenden kommen sollen, und doch hernach, da das Volk begehrt, daß ihnen möchte ein König gegeben werden, es darum von GOtt hart ist gestraft worden: darüber zu disputiren, ist hier nicht der Ort.

133. Er setzt aber dazu: „Und will deinem Samen nach dir geben“, damit niemand an der Verheißung zweifeln dürfe, als ob sie nicht gewiß wäre, dieweil man weiß, daß Abraham und Jakob in demselben Lande Fremdlinge gewesen sind, und auch nicht eines Fußes breit darin gehabt haben, ausgenommen das Stück Landes zu Sichem, das Jakob mit seinem Schwert und Bogen eingenommen hat, wie er selbst hernach sagen wird Cap. 48. V. 22. Darum legt er die Verheißung aus, daß er dasselbe Land dem Samen und Nachkommen Jakobs geben wolle. Denn was dem Sohne gegeben wird, dessen mäßt sich der Vater auch mit Recht an. Unterdeß aber hält er sich an die Verheißung vom Namen Israel, die viel besser ist. Denn GOtt war sein Land und Erbe. Und also ist nun Jakobs Schmerz, so er gehabt hat über Dina, da ihm dieselbe beschlafen war, und daß ihm Debora gestorben, gelindert worden. Aber es sind noch schwerere Stücke dahinten.
\quellseite{253}

\begin{quote}
\textbf{V. 13. Also fuhr GOtt auf von ihm, von dem Ort, da er mit ihm geredet hatte.}
\end{quote}

134. Die heiligen Patriarchen sind fleißig gewesen und haben für und für mit Lehren und Beten angehalten, und zu solchen Uebungen haben sie täglich die Frühstunden angewandt. Vornehmlich aber ist das Volk am Sabbath zusammen gekommen zur Predigt, und daß sie in der Gemeinde gemeinschaftlich gebetet und GOtt angerufen haben. Wenn sie aber die Gemeinde haben gehen lassen, hat Jakob allein gebetet und die Augen aufgehoben gen Himmel, hat sich mit dem Wort beschäftigt und davon geredet: da ist ihm GOtt erschienen, und hat nicht mit ihm geredet durch seinen Vater Isaak oder durch Debora, sondern ist gegenwärtig bei ihm gewesen, wie droben Cap. 32. V. 30., da er sagte: „Ich habe GOtt von Angesicht gesehen.“ Es ist ein solch Gesicht gewesen, wie zuvor, da der Mann mit ihm gerungen hat, allein daß er es verstanden hat, daß es GOtt selber wäre, der da redete, welches er im vorigen Kampf nicht so bald gesehen hat. Und daß es eine göttliche Offenbarung gewesen sei, bezeuget Mose selbst, da er sagt: „GOtt fuhr auf von ihm.“

\begin{quote}
\textbf{V. 14. 15. Jakob aber richtete ein steinern Mal auf an dem Ort, da er mit ihm geredet hatte, und goß Trankopfer drauf, und goß ihm mit Oele. Und Jakob hieß den Ort, da GOtt mit ihm geredet hatte, Bethel.}
\end{quote}

135. Es ist zu merken, daß dieser Ort sonderlich gerühmt wird vor andern Orten; denn jetzt wird er zum drittenmal Bethel genannt. Sollte aber das nicht falsche Propheten und Abgötterei machen? Jerobeam hat wahrlich dessen sehr gemißbraucht, da dieser Ort so gerühmt worden, und hat den ganzen Gottesdienst und Religion dahin gewendet: darum Hosea und die andern Propheten, die zu seiner Zeit gelebt haben, so großen Streit und schwere Arbeit gehabt haben, die Abgötterei zu Bethel zu strafen, zu verdammen und abzuschaffen. Und sie haben diesem Ort den herrlichen Namen genommen, und haben ihn Betheven geheißen, das ist, ein Haus des Schmerzes und der Trübsal; und richten doch damit wenig aus.

136. Denn die falschen Propheten warfen den frommen gottseligen Priestern und Propheten vor und sagten: Was ist es, daß ihr viel rühmt vom Tempel zu Jerusalem? Lest das erste Buch Mose, da werdet ihr finden, daß GOtt Jakob an diesem Ort erschienen ist und derselbe darum Bethel und El Bethel ist genannt worden. Was könnte man Gewisseres oder Stärkeres vorbringen? Und scheint es, als könne man dieses Argument nicht auflösen. Was ist es, sagten sie, daß ihr uns Mose oder das neue Gebot Salomos vorhaltet? Meint ihr, daß sie den Gottesdienst von einem solchen heiligen Ort haben können oder wollen gen Jerusalem wenden? Denn an diesem Ort haben die Patriarchen, Abraham, Isaak und Jakob geopfert, und der König Jerobeam ist um solcher Gottseligkeit willen zu loben, sagten sie, der ein frommer König gewesen ist, und das Volk an denselben Ort geführt hat, welchen die Patriarchen gerühmt haben, an welchen GOtt selbst gekommen ist, den er sich hat gefallen lassen und da er mit den Vätern geredet hatte.

137. Das ist dazumal die Gelegenheit und Ursache gewesen, die sie vorgewendet haben, die frommen gottseligen Lehrer und Propheten zu verfolgen und zu erwürgen; wie Amazia den Amos heftig schilt: „Du Seher“ spricht er, „gehe weg, und fleuch ins Land Juda, und iß Brod daselbst, und weissage daselbst. Und weissage nicht mehr zu Bethel; denn es ist des Königs Stift und des Königreichs Haus“ 2c., Amos 7, 12. 13., das ist, es ist ein Haus GOttes, das der König Jerobeam gestiftet und aufgerichtet hat.

138. Die rechte Antwort aber auf das, so die falschen Propheten vorgenommen haben, welches einen großen Schein gehabt hat, ist diese gewesen: GOtt hat mit den Vätern geredet zu seiner Zeit und an seinem Ort; aber darnach hat er durch Mose klärlich gesagt 2 Mos. 20, 24.: „In welchem Ort ich meines Namens Gedächtniß stiften werde, da will ich zu dir kommen, und dich segnen“; nicht an einem jeden Ort, Berge, grünen Aue oder Wasser, sondern du sollst auf den Ort sehen, welchen der HErr erwählen wird. Und darum hat er ihnen die Hütte, Altar und das Levitische Priesterthum gegeben, auf daß er die Herzen der Menschen daran binden möchte. Und dasselbe haben die Propheten dem Volk auch fleißig eingeschärft; aber es ist vergeblich und umsonst gewesen um des sehr großen Scheines willen, welchen die Götzendiener hatten vom Gottesdienste der Väter und daß dieser Ort so sehr berühmt war. Ja, sie haben daraus noch weiter geschlossen und gesagt: Wo Bethel ein heiliger Ort ist, so sind eben um solcher Ursache willen Gilgal, Mizpa, Silo und Thabor auch heilig, da die Väter große herrliche Dinge ausgerichtet haben. Daher ist eine unzählig große Menge Götzen gekommen, wie es bei uns mit den Wallfahrten ist zugegangen. Denn sie wollten auf das Gebot Moses nicht sehen, der die Hütte und andere Ceremonien auf GOttes Befehl geordnet hatte, und auch darum die Hütte auf Hebräisch moed genannt, 2 Mos. 29, 10., das ist, gewiß beschlossen und bestimmt, auf daß er die irrenden Herzen dahin führen möchte, daß sie auf einerlei gewisse Weise GOtt dienen sollten.

139. Da aber nun Jakob das steinerne Mal hat aufgerichtet, hält er damit dieselbige Weise, die er droben, Cap. 28. V. 18., auch gebraucht hat; denn er begießt ihn mit Oel: und hat vielleicht auch denselben Stein genommen. Das Gießen aber geschieht mit dem, das da fließt, als vornehmlich mit Wein und Oel. Und ist eine solche Weise gewesen, die sie gebraucht haben zu weihen und zu salben, welche darnach Mose auch angenommen hat; denn derselbe salbte die Priester und Kleider 2c. Wie Jakob den Stein begießt und salbt, auf daß es ein solcher Ort sei, der dazu geweiht und GOtt zugeeignet ist, daß daselbst die Gemeinde zusammen kommen sollte, GOttes Wort zu hören, GOtt anzurufen und GOtt zu danken. Das Geheimniß oder die heimliche Deutung übergehe ich. Dieses ist der leibliche Segen und wie die Kirche oder Gemeinde soll angerichtet werden. Und ich glaube gerne, daß der Patriarch Jakob an diesem Ort mehr gethan habe, denn Mose geschrieben hat. Denn droben, Cap. 28. V. 20. ff., hat er ein Gelübde gethan, daß er den Zehnten von seinen Gütern geben wollte, wo er gesund und unverletzt würde wieder kommen: darum halte ich dafür, er werde den Kindern Eber den Zehnten gegeben haben, die Kirche damit zu erhalten.

\section[Vierter Theil]{Vierter Theil. Wie Rahel über der Geburt stirbt und Jakob ihr ein Grabmal setzt}

\begin{quote}
\textbf{V. 16. 17. Und sie zogen von Bethel. Und da noch ein Feldwegs war von Ephrath, da gebar Rahel. Und es kam sie hart an über der Geburt. Da es ihr aber so sauer ward in der Geburt, sprach die Wehmutter zu ihr: Fürchte dich nicht, denn diesen Sohn wirst du auch haben.}
\end{quote}

140. Das hebräische Wort kibrath ist der Worte eines, die weder die Juden noch wir verstehen. Und die ganze hebräische Sprache würde noch viel dunkler sein, wenn wir Christen dieselbe mit der lateinischen und deutschen Uebersetzung der Bibel nicht erklärt hätten. Denn wo die Juden irgend an einem Worte zweifeln, so geben sie denselben vielerlei Bedeutungen und verdunkeln es mit ihren Glossen viel mehr. Darum streiten die Grammatiker über dieses Wort auch und sind der Sachen noch nicht eins. Hieronymus macht zwar Worte genug davon, wiewohl er es übel auslegt. Er sagt, daß es die siebenzig Dolmetscher gegeben haben, hippodromus, das ist ein solcher Ort, da man mit Rossen in die Wette rennt; so doch die Ableitung des Wortes solches nicht mit sich bringt. Und ich kann nicht verstehen, was sie damit gemeint haben, denn vielleicht das allein, daß es so weit gewesen sei von dem Orte, da die Rahel geboren hat, gen Bethel, wie ein Roßlauf weit ist.

141. Unter den Rabbinen bringt Einer dieses, der Andere ein Anderes herbei. Einer macht ein Wort daraus, das von mehreren Worten zusammengesetzt sei; der Andere will haben, es sei ein einfaches Wort. Die da haben wollen, daß es von andern Worten zusammengesetzt sei, sagen, es komme her von k, welches die Aehnlichkeit bedeutet, und vom Wort barath, quasi barath terrae, id est, spatium terrae ad vendendum in Ephratha, etwa ein Feldweges von Ephratha. Wie weit oder viel aber desselbigen sei, ist damit noch nicht ausgedrückt. Die Andern sagen, es komme her vom Wort birjah, welch Wort 2 Sam. 13, 8. gesetzt wird von Thamar und Ammon, und Speise oder Suppe heißt; als wenn er auf Lateinisch sagen wollte: spatium terrae, unius jentaculi etc., das ist, so viel Weges als Einer gehen kann bis zum Morgenbrod, das ist ungefähr bei einer Viertelmeile Weges. Diese Ableitung des Wortes gefällt mir sehr wohl, wo man anders irgend eine Ableitung annehmen soll. Die Andern deuten es insgemein als eine Strecke Weges; sie sagen aber nicht, wie weit es sei. Hieronymus sagt, es komme her vom Wort bahur, das heißt auf Lateinisch electus, auserlesen, und will haben, daß es eine auserlesene Zeit im Jahr heißen soll, nämlich den Lenz oder die Zeit des Mai, wann die lustigste Zeit ist im Jahr, formosissimus annus, wie der Poet sagt, und wann die Menschen über Feld gehen, und Blumen brechen auf dem Lande, so allernächst umher gelegen ist. Aber das reimt sich nicht zu dem Maß oder Weite, die Mose hat anzeigen wollen. Darum will ich denen folgen, die es also deuten, daß es so viel Weges heißen soll gen Bethlem, so weit man gehen kann bis zum Morgenbrod: ein klein Weglein, daß Einer möchte zu der Suppe lustig werden.
\quellseite{254}

142. Nun läßt uns die Anfechtung ansehen, die auf den Frieden und Sicherheit des Gewissens folgt, welchen Jakob von der göttlichen Erscheinung und Verheißung bekommen hat. Wir können aber diesen Text nicht vollkommenlich erklären, bis wir an das 48. Capitel kommen werden, da Jakob anzeigt, daß er zu dieser Zeit die Worte der Verheißung nicht verstanden habe, wiewohl sie voll Trostes gewesen sind. Die Ursachen dieser Dunkelheit sind die folgenden Anfechtungen. Denn es wird ihn der Segen verheißen und daß er Samen haben soll, nämlich, da GOtt spricht: „Sei fruchtbar und mehre dich“; und geschieht doch dem äußerlichen Ansehen nach das Gegentheil deß, so ihm verheißen wird. Warum wollen wir verziehen bis in das 48. Capitel mit der Auslegung. Denn da wird Jakob erst durch die Erfahrung lernen, wie die Verheißung hat sollen verstanden werden. Jetzt ist er damit zufrieden, daß GOtt zu ihm herab kommt und ihm Trost und Hülfe verheißt; wie er es aber verstehen soll, wird er hernach erfahren.

143. Nach der Erscheinung rüstet er sich, mit friedlichem Herzen wiederum zu seinem Vater zu ziehen, und war schon nahe bei Hebron gekommen. Denn Bethel ist zwischen Jerusalem und Hebron gelegen, welche auf sechs deutsche Meilen von einander liegen. Jakob ist von Hebron drei oder vier Meilen weit, und ist wohl glaublich, daß er vor der Heerbe und dem ganzen Gesinde, nachdem er solche hinter sich gelassen, sei hergezogen und geeilt sei zu seinem Vater, da er nur etliche wenige Gefährten mit sich genommen hat. Da sie nun von Mitternacht und von Bethel nach Hebron ziehen, kommt unversehens dieser Unfall, welcher der Verheißung stracks zuwider ist, da er sicher auf der Schwelle in seines Vaters Hause war. Er brachte Freude und Trost mit sich von der Verheißung, und war fröhlich, daß er sein sehr liebes Weib, das schwanger war, zu seinem Vater bringen sollte und ihn damit erfreuen, wenn er sie sehen würde. Auf demselben Wege ist es der Rahel angekommen, daß sie hat gebären sollen, und da sie von den Schmerzen in der Geburt überwunden ward, ist ihr die Seele ausgegangen.

144. Ist aber nun das der Trost, ist das ein Segen oder Vermehrung, daß ihm sein vornehmstes und allerliebstes Weib genommen wird, die wohl erzogen war und sehr gute Sitten gehabt hat, dazu auch sehr wohl erfahren gewesen ist, das Hausregiment zu verwalten? Denn sie war daheim bei ihrem Vater Laban an Hausarbeit und des Viehes zu warten gewöhnt. Zudem so hatte sie sich mit ihrem Mann ihrer sauren Arbeit genährt und fast kümmerlich gelebt. Darum wird nun Jakob seines Weibes als einer treuen Gesellin beraubt, von der er gehofft hat, daß er sie in seines Vaters Haus einführen sollte, dem alten Vater Isaak auf seinen Schooß, vornehmlich da sie schwanger war, auf daß er den Vater in seinem Alter mit dem lieben Kindlein erfreuen möchte. Sie selbst hat auch von Herzen gewünscht, den Schwiegervater zu sehen. Aber da fällt ihm bald alle Hoffnung und Freude dahin. Und werden also Jakob, Isaak und Rahel der süßesten Wünsche und Hoffnungen beraubt, so sie gehabt haben.

145. Heißt aber das fruchtbar sein und von GOtt gesegnet werden? Ja, das ist ein sehr böser Fluch. Aber unser HErr GOtt hat eine hebräische Sprache, das Hinterste zu vörderst, das Leiden zuvor und die Herrlichkeit hernach. Das Kreuz und Leiden muß vorher gehen, wie St. Petrus sagt 1 Petr. 1, 11., daß der Geist Christi in den Propheten zuvor bezeuget hat die Leiden, die in Christo sind, und die Herrlichkeit darnach. Und zu Mose spricht der HErr: Zuletzt werdet ihr meinen Rath verstehen. Wie Jakob hernach in Egypten das Letzte auch versteht, da Joseph wiederum lebendig wird, von dem er meint, daß er verloren und gestorben wäre.

146. Unterdeß aber muß er alles Trostes beraubt sein und das ganze Haus wird mit Wehklagen und Traurigkeit erfüllt, und nachdem die Sonne untergegangen, ist die allerbetrübteste Finsterniß eingefallen, da ihm sein allerliebstes Weib stirbt, welches die ersten sechzig Jahre unfruchtbar gewesen und zehn Jahre nach Joseph geboren hatte; welches Weib der Mann noch immer lieber genommen hätte, darum daß er täglich gewartet und Verlangen gehabt hat, daß sie ihm mehr Kinder gebären würde, dieweil sie schwanger war. Darnach ist sie dem ganzen Gesinde eine sehr liebe Herrin gewesen, eine Hoffnung und Trost im Hause. Dieselbe aber wird aus diesem Leben abgefordert unversehens und daß sie nicht gehofft hatten, und stirbt dahin an einem ungelegenen und unglückseligen Orte. Denn wo sie vor der Zeit in Syrien gestorben wäre, so wäre es doch etwas leidlicher gewesen. Jetzt wird der Schmerz um so viel größer, dieweil sie gleichsam auf der Schwelle im Hause des Vaters dem Jakob dahin stirbt.

147. Ich will jetzt schweigen der Traurigkeit und Wehklagen, so er darüber gehabt hat, daß ihm das fromme gottselige Weib genommen war; welche Traurigkeit an sich selbst sehr schwer ist. Denn in diesem Leben ist nichts besser oder angenehmer, denn wo Mann und Weib freundlich und in Einigkeit mit einander leben. Wie Sirach sagt im 25. Cap. V. 1. 2.: „Drei schöne Dinge sind, die beide GOtt und den Menschen wohl gefallen: wenn Brüder eins sind und die Nachbarn sich lieb haben, und Mann und Weib sich mit einander wohl begehen.“

148. Man hat nicht sehr großen Schmerz darüber, wo Einem ein eigeninniges, ungeschlachtet, böses Weib stirbt: aber unter diesen Eheleuten ist sehr große brünstige Liebe und ganz freundliche Neigung gewesen, so Eines zum Andern getragen hat. Sie ist züchtig, heilig, in ihrem Amte fleißig und dazu den Nachbarn und dem Gesinde lieb gewesen, hat ihren Mann lieb gehabt, ist ihm in allen Dingen gehorsam gewesen. Darum ist des Patriarchen Jakob Schmerz und Traurigkeit überhalb groß gewesen und er hat nun alle Hoffnung und seine größte Freude verloren, und also gedacht: Was ist es, daß GOtt sagt, daß mein Geschlecht gemehret und groß werden soll, so er mir doch die Mutter nimmt, durch welche ich fruchtbar und gemehret werden sollte?

149. Er hat die Verheißung nicht verstanden, und hat auch diese zwei Dinge, die eigentlich stracks wider einander sind, nicht zusammenreimen können, nämlich, daß ihm GOtt einen Samen und Erben verheißt, und doch gleichwohl die Hausmutter wegnimmt, von welcher er und das ganze Haus mit Isaak, dem Großvater, Kinder und Erben erwartet haben. Und Jakob hat vornehmlich um des Trostes willen, der ihm neulich gegeben war, sich seiner Gefahr besorgen können mit seinem Weibe; ja, sie haben alle einträchtig Rahel gerühmt, wie sie so selig wäre, darum daß sie vieler Kinder Mutter werden würde, sinnemal GOtt gesagt hatte, daß Jakob würde gemehret werden. Da hat niemand gemeint, daß es anders geschehen würde denn durch Rahel. Und Jakob wird oftmals zu ihm gesagt haben: Meine liebe Rahel, laß uns GOtt danken, der mir jetzt abermal aufs neue Erben verheißen hat. Nun erst wirst du fruchtbar werden. Und um solcher Ursache willen, wo er sie zuvor jemals lieb gehabt, so hat er sie jetzt erst recht lieb gewonnen. Aber diese schöne liebliche Hoffnung, so die Eheleute und das Gesinde gehabt haben, hat ihnen gefehlet. Wie es der Eva auch ergangen ist, da Kain geboren war; denn sie sprach also 1 Mos. 4, 1.: „Ich habe den Mann, den Herrn.“ Als wollte sie sagen: Der Herr ist selbst gekommen, der der Schlange den Kopf zertreten wird. Aber es hat ihr nicht allein ihre Hoffnung gefehlet, sondern sie hat auch von dem Sohn ein anderes großes Unglück erwarten müssen.
\quellseite{255}

150. Also spricht Gott hier zu Jakob: Du wirst fruchtbar und gemehret werden, ich will dir aber das Weib nehmen. Diese Dinge sind stracks wider einander. Rahel war die Mutter des erstgebornen Sohnes Joseph und war die rechte Hausmutter, und sie haben gehofft, daß jetzt bald Benjamin und darnach noch andere Kinder mehr von ihr kommen sollten. So gebieret sie nun auf dem Wege, da sie zu Ephrath, dem Großvater, ziehen wollten. Aber es ist ihr in der Geburt sauer geworden und ist darüber in Gefahr Leibes und Lebens gekommen, und kann das Kindlein dem Großvater nicht in den Schooß legen, wie sie doch heftig begehrt hat; sondern sie stirbt dahin von großem Schmerz der Geburt und wird zugleich mit ihr die Freude und sehr schöne Hoffnung auch gar verlöscht. Die Wehmutter sieht, daß sie mühselig und erschrocken ist aus Furcht des Todes, darum tröstet sie sie und spricht: „Du wirst deinen Sohn auch haben.“ Aber der Trost ist umsonst. Denn Rahel fühlt, daß sie kraftlos und matt wird, und daß sie der Tod allegemach überfallen hat. Darum heißt sie den Sohn Benoni, nämlich, vom Schmerz.

151. Dieses ist das Ende, so Rahel genommen hat, und ist wahrlich ein erbärmlich und betrübt Ende beide dem Manne selbst und auch sonst dem anderen Gesinde. Sie entschläft in der Geburt, in ihrem Beruf und Stande, darein sie von Gott gesetzt war, und erliegt unter dem Kreuz, welches Gott den Weibern auferlegt hat, da er sagt: „Ich will dir viel Schmerzen schaffen, wenn du schwanger bist“ etc. Sie ist von der Gefahr und dem Jammer nicht frei, den die gottlosen Weiber tragen müssen. Und dasselbige ist auch zu erbarmen, daß ein solch heilig Weib, das eine Ehre und schöner Schmuck desselben Hauses und der ganzen Gemeinde gewesen ist, eben auch dieselbige Beschwerung tragen soll, welche den Gottlosen zu widerfahren pflegt.

152. Dieses Exempel soll man fleißig betrachten, und das Wehklagen und den Schmerz des frommen Patriarchen wohl bedenken, welchen Schmerz alle Umstände in dieser Historie groß und schwer machen, und ist des Jammers noch kein Ende. Denn da er die Mutter verloren hat, verliert er alsbald und fast im selben Jahre auch Joseph, den erstgebornen Sohn von dieser Mutter. Zum Dritten, wird ihm Bilha, der Rahel Magd, beschlafen von Ruben, seinem erstgebornen Sohne von andern Weibe. Heißt das segnen? Ja, es hat Jakob denken mögen: Es hat nicht Gott, sondern der Teufel mit mir geredet; wie denn irgend ein ungelaubiger Mensch hätte leichtlich herausfahren mögen und solche Worte reden. Jakob ist zu seinem Vater nicht gekommen, ehe ihm diese zwei großen Unfälle begegnet waren, als, der Tod der Rahel, und die Blutschande, so Ruben mit Bilha begangen hat, welche Bilha an der Rahel Statt war gekommen. Das ist wahrlich allzuviel gewesen, daß er mit drei so greulichen Unfällen in so kurzer Zeit nach einander hat sollen überfallen werden, von denen einer allein einen ungeduldigen und fleischlichen Menschen tödten könnte.

153. Wir aber sind zarte Leute, etwas zu leiden (delicati martyres). Denn wir verstehen dieses große Unglück nicht und wir würden es viel weniger ertragen können. Das Leiden Christi ist höher, denn alle Leiden sein können, und es beweget uns auch so sehr nicht, dieweil er GOtt ist. Aber die Väter sind uns gleich gewesen und haben alle menschliche Affecten und Neigungen gehabt. Als, Jakob hat seine Rahel lieb um der trefflichen Tugenden willen, so sie gehabt hat. Sie, Rahel selbst, hat gute Hoffnung, daß viel Kinder von ihr sollen geboren werden; aber gleich in derselben Hoffnung und in der Geburt stirbt sie dahin.

154. Darum sollen nun andere Leute wissen, mit ihrem Schmerz so viel desto mehr Maß zu halten, wenn etliche ihrer Gemahlin oder Kinder beraubt werden, dieweil wir sehen, daß die Väter eben dasselbe auch gelitten haben. Denn wir sind viel geringer denn sie, und haben die Verheißung nicht vom Königreiche, vom Samen oder Nachkommen und vom Priesterthum, wie sie gehabt haben; welches dazumal Jakob alles verloren hat.

155. Solche Anfechtungen aber werden immer härter und heftiger, wenn der Schrecken im Gewissen auch noch dazu kommt, wenn der Satan Oel in das Feuer gießt. Denn wo der Teufel sieht, daß die Herzen angefochten und erschrocken sind, so schießt er sie bald auch mit seinen feurigen und vergifteten Pfeilen, daß sie denken sollen, GOtt zürne mit ihnen und sei ihnen feind geworden.

156. Darum ist es ein großer Trost in aller Anfechtungen, daß wir es dafür halten können, daß GOtt bei uns und uns auch günstig sei. Aber das menschliche Herz kann schwerlich solchen Trost annehmen, wenn unser Herr GOtt Einen also herzt, daß ihm die Seele ausgehen will. Aber das ist wahrlich eine solche Liebe, wie er an seinem eigenen Sohne erzeigt hat; wie Christus sagt Joh. 15, 9.: „Gleichwie mich mein Vater liebet, also liebe ich euch“; nämlich nach häuslicher Art zu reden, davon wir droben gesagt haben, das ist, daß das Leiden vor der Herrlichkeit hergehe; welche Art zu reden die Welt nicht versteht. Denn GOtt füllet den Gottlosen, die er nicht lieb hat, den Bauch mit leiblichen Gütern, Ps. 17, 14.; er setzt aber dazu und sagt Luc. 6, 24. 25.: „Wehe euch, die ihr hie lachet und voll seid“ etc.

157. Auf solche Weise werden uns die Trübsale der heiligen Väter zum Exempel vorgehalten, auf daß wir ihre Uebung und trefflichen Glauben in solchem Jammer und Unglück sehen, und daß wir lernen, daß wir das Angesicht Gottes, damit er uns anlacht, segnet und Gnade verheißt, nicht verlieren, wenn er uns das wegnimmt, was uns wohl thut und lieb ist, und das er mit seiner Verheißung bestätigt hat. Ja, der Glaube soll also schließen: Wiewohl mir dieses Unglück sehr schwer ist, so glaube ich doch, daß GOtt wahrhaftig sei. Aber in solchen schweren Fällen, die so unversehens kommen, wird der Glaube sehr hart gedrückt und läßt sich kaum sehen, wie das glimmende Docht. Darum muß er mit solchen Exempeln auferweckt werden.

\begin{quote}
\textbf{V. 18. Da ihr aber die Seele ausging, daß sie sterben mußte, hieß sie ihn Benoni; aber sein Vater hieß ihn Benjamin.}
\end{quote}

158. Mose überläuft es alles kürzlich. Es sind aber ohne Zweifel viele Worte gefallen und ist auch Seufzen und Weinen mit untergelaufen, ehe das Kind Benjamin ist genannt worden. Rahel selbst hat GOtt ihre Seele befohlen, und in den letzten Zügen und da es ihr über der Geburt so hart gegangen ist, hat sie sich an den Trost gehalten, den sie von Jakob gehört hatte. Denn sie ist voll Glaubens gewesen, und ist zuvor hart genug angefochten worden mit gemeinem Unglück und Jammer ihres Mannes im Hause Jakobs, das sie mit ihm gelitten und getragen hat. Wie Sarah und Abraham auch gemeine Noth und Gefahr mit einander getragen haben. Denn fromme gottselige Weiber lassen sich zu Herzen gehen die Anfechtung und Beschwerung ihrer Männer, und haben mit ihnen Mitleid, weinen mit ihnen. Solche natürliche Neigung ist an Rahel auch sehr trefflich gewesen. Darum hat ihr Tod den Mann, das Gesinde und auch die Nachbarn sehr betrübt, denen ihre Tugenden und Gottseligkeit bekannt gewesen ist.

159. Wir sollen aber den Trost lernen, damit wir uns aufrecht erhalten mögen, wenn uns unser Gemahl, Kinder und gute Freunde mit dem Tode abgehen. Denn wiewohl beide unser und auch unserer Verwandten Tod ein sehr betrübt Ding ist, so ist doch das gewiß, daß wir ewig leben werden. Es ist zwar besser gewesen, daß Rahel also in rechter Anrufung und Glauben an GOtt gestorben ist, als wenn sie in diesem Leben noch allerlei Freude sollte gehabt haben, und dazu auch die Freude, so GOtt verheissen hätte. Denn nach dem Tode werden wir viel seliger und über mehr Güter Herren sein, denn wir in diesem Leben verlassen haben. Denn das sind zeitliche und vergängliche Güter, dafür wir ewige Güter bekommen.

160. Aber das Fleisch kann sich dessen nicht enthalten, es muß heulen und Leid tragen; wie dieses Exempel anzeigt, da sie diesem Sohne einen Namen geben vom Schmerz. Und ich halte dafür, er werde eine Zeitlang also genannt worden sein, bis der Schmerz etwas gelindert worden, und Jakob ihn verändert und gesagt hat: „Ich will in meinem Hause nicht mehr ein Trauergedächtniß haben, das mich für und für erinnern soll der letzten Worte und Schmerzen meines Weibes, das ich verloren habe, und mir also meinen Schmerz auch immer erneuern; er hieß ihn aber Benjamin, das ist, den Sohn seiner Rechten.

161. Sie disputiren aber von der Ursache dieses Namens. Ich halte dafür, daß es nach dem Hausregiment also könne verstanden werden: Lea ist das Weib der Linken und Rahel das Weib der Rechten gewesen, das ist, die erste und vornehmste. So ist nun Benjamin der Sohn der Rechten, als wollte er sagen, seines allerliebsten Weibes. Und Jakob hat haben wollen, daß bei diesem Namen vielmehr der Liebe, so er zu seinem allerliebsten Weibe getragen hat, denn des Schmerzens und Wehklagens sollte gedacht werden. Also verändert er den traurigen Namen in einen fröhlichen Namen.
\quellseite{256}

162. Die andere Ursache dieses Namens gehört zum Welt- oder Kirchenregiment. Als wollte er sagen: Dieser Sohn, der von meinem allerliebsten Weibe geboren ist, wird das verheißene Erbe empfangen und das dem Sohn des vornehmsten Weibes gebührt. Er wird der Erbe sein des Königreichs und des Priesterthums. Denn Jakob ist der Meinung, daß die Verheißung und der Segen Joseph gebühre, wie er den Segen vor seinem Bruder Esau erlangt hatte. Und wenn Joseph sterben würde, meint er, daß Benjamin an Josephs Statt kommen werde. Er versteht noch nicht, daß Juda der Vater Christi sein werde, bis er es in Egypten erst lernen wird, da es ihm der Heilige Geist eingeben wird. Er wird aber in dieser seiner Meinung so viel mehr bestärkt, weil er keinen Sohn mehr zu bekommen hofft. Darum hat er ihn sonderlich lieb und daher kommt auch der Haß seiner Brüder gegen Joseph. Und da derselbe Haß und Neid ausgebrochen und an den Tag gekommen ist und Joseph von seinen Brüdern ist verkauft worden, hat sein Vater Jakob über die Maßen großen Schmerz darüber gehabt.

\begin{quote}
\textbf{V. 19. Also starb Rahel, und ward begraben an dem Wege gen Ephrath, die nun heißt Bethlehem.}
\end{quote}

163. Rahel ist gestorben in rechter Anrufung und Glauben an Gott, und ist aufgenommen worden in den Himmel, in den Schooß Abrahams, ihres Vaters. Sie ist aber begraben worden nicht weit von Bethlehem, fast ein Viertel Weges einer deutschen Meile: da ist sie gestorben und begraben auf dem Felde, so allerseits dabei gelegen, oder auf dem Wege; gleich als wenn irgend eines Hirten Weib auf dem Felde oder unter dem Vieh dahin stürbe: da ist kein Hauß oder Herberge gewesen, ausgenommen die Hütte Jakobs. Das ist der rechte Weg der Heiligen zum Himmel.

164. Ephratha und Bethlehem sind zweierlei Worte, sie haben aber schier einerlei Bedeutung. Ephratha heißt fruchtbar. Eben so viel heißt auch Bethlehem, ein Brodhaus, da eine gute Schnabelweide ist.

\begin{quote}
\textbf{V. 20. Und Jakob richtete ein Mal auf über ihrem Grabe; dasselbe ist das Grabmal Rahels bis auf diesen Tag.}
\end{quote}

165. Die Väter haben die Gräber herrlich geziert und in Ehren gehalten: sie haben die Todten nicht hingeworfen wie das Aas der unvernünftigen Thiere, sondern sie haben ihnen Grabmale aufgerichtet zu eigenem und unsterblichem Gedächtniß, auf daß es Zeugnisse wären der künftigen Auferstehung, die sie geglaubt und darauf sie gehofft haben.

166. Darum soll man die Ceremonien und das Leichengepränge behalten, desgleichen auch, daß man weine und Mitleid habe mit denen, so ihre lieben Freunde verloren haben. Nicht, daß man für die Todten bitten soll, wie Lyra haben will, oder daß wir uns vor dem Tode fürchten sollten; sondern daß wir vielmehr im Tode lernen den Glauben üben, der da kämpft mit dem Schrecken des Todes. Und sollen also denken, daß wir zwar sterben und verächtlicher Weise begraben werden, aber daß wir mit großer Herrlichkeit wiederum auferstehen werden. Wie Jakob diesen Trost behalten hat, als wollte er zu seinem Weibe sagen: Meine liebe Rahel, die Hoffnung deiner Freude in diesem Leben hat dir gefehlet: aber das ist ein seliger Fehler. Denn die leiblichen und vergänglichen Güter hast du mit den ewigen vertauscht.

167. Zu dem Ende und um solcher Ursache willen sind die Male bei den Begräbnissen aufgerichtet worden, auf daß wir lernen, nach dem Exempel der Väter dieses gegenwärtige Leben gering schätzen und verachten um des andern bessern Lebens willen, so noch dahinten ist. Darnach, daß wir uns auch an die Verheißung halten und mit Geduld den Willen Gottes ertragen lernen. Denn was kann der Teufel oder der Türke mehr thun, wenn sie uns auch schon das Leben nehmen? „Denn so wir glauben, daß Jesus gestorben und auferstanden ist; also wird Gott auch, die da entschlafen sind durch Jesum, mit ihm führen“, wie St. Paulus 1 Thess. 4, 14. sagt.

\begin{quote}
\textbf{V. 21. Und Israel zog aus, und richtete eine Hütte auf jenseit des Thurms Eder.}
\end{quote}

168. Der lateinische Dolmetscher hat aus dem Wort „Eder“ einen Gattungsnamen gemacht, aber es ist ein Eigenname eines Schlosses oder Thurms, welcher auf dem bethlehemischen Felde eine Warte gewesen ist. Denn es ist gewiß, daß auf dem bethlehemischen Felde allezeit Weide gewesen ist. Hieronymus schreibt, daß hernach Eder zu einem Tempel oder Kirche verändert worden sei, die zu seiner Zeit auf Lateinisch Angelus ad Pastores genannt worden, darum daß daselbst die Engel den Hirten erschienen sind in der Nacht, da Christus geboren ist. Denn es sind daselbst allezeit viele Hirten und Heerden gewesen, dieweil gute Weide da gewesen ist für die Schafe. Darum hat dieser Thurm davon den Namen gehabt, daß die Heerde daselbst ist geweidet worden. Das Schloß oder der Thurm ist jenseit Bethlehem gelegen gewesen, nach Mittag hin. Lyra sagt eben auch dasselbe und folgt dem Hieronymus nach.

169. Die Andern sagen, das Wort „Eder“ habe verschiedene Bedeutung, und deuten es auf die Stadt Jerusalem, darum daß im Propheten Micha am 4. Cap. v. 8. also steht: „Und du Thurm Eder, eine Feste der Tochter Zion“ etc. Welcher Irrthum daher gekommen, daß der Text im Lateinischen dunkel und übel verdolmetschet ist, und daß sie auf die Grammatik nicht Acht gegeben haben. Denn der Prophet Micha redet vom Thurme der Heerde als einem Gattungsnamen durch die Figur, metonymia genannt, wie das aus den nachfolgenden Worten zu sehen ist: „Du Thurm der Heerde Ophel“, das ist, die Feste oder Burg der Stadt. Darum ist es ein Eigenname einer Stätte auf dem Felde, nahe bei Bethlehem, da heutiges Tages ein Tempel oder Kirche sein soll, die auf Lateinisch genannt wird Angelus ad Pastores.

\section[Fünfter Theil]{Fünfter Theil. Von der bösen That Rubens; von den zwölf Söhnen Jakobs; und von Isaakobs Tode und Begräbniß}

\subsection{I}

\begin{quote}
\textbf{V. 22. Und es begab sich, da Israel im Lande wohnete, ging Ruben hin, und schlief bei Bilha, seines Vaters Kebsweibe; und das kam vor Israel.}
\end{quote}

170. Nach dem Tode Deboras und Rahels folgt nun ein anderer und greulicher Unfall des heiligen Patriarchen: und wie greulich derselbe ist, kann ich mit Worten nicht erreichen oder ausreden. Ich will es den Rednern befehlen groß zu machen; denn es ist eine greuliche Sünde. Der erstgeborne Sohn Ruben, der zur Königsherrschaft und Priesterthum verordnet war, verderbt sich selbst mit allen seinen Nachkommen. Es ist eine greuliche That, aber dem Thäter ist sie am schädlichsten.

171. Wir wollen es aber doch ein wenig lindern und wollen sagen, daß er aus Thorheit oder böser Lust sündigt habe. Und ich halte dafür, Bilha werde schön gewesen sein und nicht sehr beständig; denn sie hat sich den närrischen Jüngling bald überwinden lassen. Wo sie ihn aber etwas hart gestraft und ihm gedroht hätte, sie wollte es seinem Vater sagen, so hätte sie ihn mögen abweisen und zurücktreiben. Aber es lautet gar sehr übel. Darum ist Bilha sowohl zu strafen als Ruben.

172. Mose nennt sie ein Kebsweib. Sie ist aber doch Jakobs rechtes Eheweib gewesen und eine Mutter der zwei Söhne Dan und Naphtali; welches zwei große Stämme gewesen sind, davon Simson und Barak hergekommen sind. Ihr Bruder Ruben aber hat sie dermaßen geschmäht und gehöhnt, daß sie die Schande für und für haben tragen müssen. Debora und Rahel waren erst gestorben und begraben und sind ihnen beiden die Beine gleichsam noch warm gewesen, da diese Sünde begangen wird. Denn die zwei Weiber waren die Zuchtmeisterinnen im Hause gewesen. Da die gestorben sind, ist Bilha, der Rahel Magd, da geblieben, und wird alsbald eine Veränderung und Unordnung der Zucht, so in der Hütte Rahels war gehalten worden; welche Unordnung sie leichtlich hätte verhüten können, und würde nicht gelitten haben, daß die Magd so geil geworden wäre.

173. Der Teufel hat die große Schande angerichtet, den frommen heiligen Mann Jakob damit zu plagen, welcher die Verheißung gehabt, daß von ihm Völker und Könige sollten geboren werden. Was hat aber dieser Segen für einen Anfang? Es stirbt ihm sein Weib, auf welcher alle Hoffnung der Erbschaft oder Kinder gestanden hat. Und jetzt wird nun die Hütte seines vornehmsten Weibes schändlich verunreinigt. Denn Jakob hat sich hinfort zu dem Kebsweibe nicht mehr gehalten. Und es wäre nicht Wunder gewesen, wenn er von dem großen Schmerz gar wäre unterdrückt worden. Ich würde wahrlich solches nicht ertragen oder überwinden können. Denn es ist allzu hart oder schwer und gar unleidlich.
\quellseite{257}

174. Und die Wunde und Schmerz, so er gehabt hat vom Tode Rahels und Deboras, war dazumal noch neu. Da macht ihm sein erstgeborner Sohn Ruben eine neue Wunde und Schmerz, da ihm die Thränen noch die Backen herabgeflossen sind. Bilha aber hat willig diese Schande begangen, ist dazu nicht hart genöthigt oder gezwungen worden. Und ist dies keine Hurerei oder Ehebruch, oder eine bloße Unzucht, sondern ist eine überaus schändliche Blutschande, daß der Bube bei seines Vaters Weibe schläft, da der Vater noch am Leben gewesen ist.

175. Dieses wird uns aber zum Trost geschrieben, auf daß wir wissen, daß unser Leiden und Trübsal nicht das alleräußerste oder schwerste sei, und daß es gar nicht zu vergleichen mit dem Jammer und Trübsal der Patriarchen. Darnach sollen wir auch wissen, daß die Sünden vergeben werden; aber doch nicht ohne Strafe. Denn dieser Stamm ist hernach schwach und sehr verachtet und im ganzen Volk Israel der geringste gewesen; und ist die Erstgeburt, das ist, das Königreich und Priesterthum von ihm genommen worden, also daß es dem Ruben viel besser gewesen wäre, auch siebenmal zu sterben. Denn diese Historie wird für und für bis an der Welt Ende gelesen und ausgerufen: in allen Kirchen und Schulen hat man diese große Schande hören erzählen, der sich Kindeskind und alle Nachkommen haben schämen müssen. Und man findet keine Historie, darin gemeldet werde, daß Gott in diesem Stamme etwas Treffliches ausgerichtet habe. So bleibt nun Ruben am Leben, und der Stamm, so von ihm hergekommen ist, auch; aber sie haben in ewiger Schmach leben müssen. Und diese Strafe ist wahrlich schwer genug, daß alle Nachkommen haben hören müssen, daß man von ihres Vaters Schande auf allen Kanzeln gesagt hat.

176. Wo sind aber nun die Söhne Jakobs, die starken Helden, Unrecht und Sünde zu rächen und zu strafen? Warum erwürgen sie nun ihren Bruder nicht, der eine greuliche Blutschande begangen hat? Sicherlich ist erwürgt worden, darum daß Dina, Jakobs Tochter, geschlafen hatte: das war aber noch eine solche Sünde, die mit einer ehrlichen Heirath hätte mögen beschönigt, zugedeckt und geheilt werden. Es war ja keine Blutschande oder irgend ein verbotener Grad. Zuhier hätte man wohl einer sonderlichen Strafe bedurft, diesen Buben zu strafen um der Blutschande willen, die er begangen hatte. Und Dan und Naphtali hätten genug gute Ursache gehabt. Denn man hat ihnen immer vorwerfen können, daß ihre Mutter Blutschande getrieben hätte; wiewohl sie aus der Blutschande nicht geboren waren, so hat aber doch ihre Mutter hernach die Blutschande begangen. Und die Schele hat ihnen Ruben angehängt. Darum ist es ein Anfang und Ursprung gewesen eines sehr großen Hasses und heimlicher Feindschaft, die für und für zwischen ihnen genähret hat.

177. Die Juden versuchen, diese Sünde an Ruben zu entschuldigen; wie sie die Thränen über die Maßen sehr pflegen zu loben, wie vor Zeiten die Griechen, und heutiges Tages die Italiener alle anderen Völker gegen sich verachten, und sie barbar nennen, das ist, grobe ungeschickte Leute. Also thun die Juden auch: wo an ihrem Volk auch offenbare Laster gefunden werden, so versuchen sie doch, dieselben gering zu machen, zuzudecken und zu schmücken. Darum sagen sie, Ruben habe keine Blutschande begangen, sondern er habe seinem Vater Gewalt und Unrecht gethan. Denn da Rahel gestorben, soll Jakob der Bilha Bett in der Lea Schlafkammer gesetzt haben; und darum sei Ruben zornig geworden und habe das Bett Bilhas weggenommen, und habe nicht bei ihr geschlafen etc. Also verfälschen sie die heilige Schrift mit ihrem sehr ungeschickten lügenhaftigen Geschwätz. Aber im Hebräischen sagt der Text offenbar: Er hat bei Bilha geschlafen: wie des Potiphar Weib zu Joseph spricht, wie hernach steht im 39. Cap. V. 7.: „Schlafe bei mir“; item, wie hernach im 49. Cap. V. 4. steht: „Er fuhr leichtfertig dahin, wie Wasser. Du sollst nicht der Oberste sein; denn du bist auf deines Vaters Lager gestiegen, daselbst hast du mein Bette besudelt“ etc.; item, 1 Chron. 6, 1.: „Ruben war der erste Sohn, aber damit, daß er seines Vaters Bette verunreinigte, ward seine Erstgeburt gegeben den Kindern Josephs, des Sohns Israel; und er ward nicht gerechnet zur Erstgeburt.“

178. Darum kann solche große Sünde nicht zugedeckt oder beschönigt werden; und einen wie großen Schmerz der fromme Patriarch gehabt hat von dieser Schande seines Sohnes und Weibes, können wir leichter denken, denn mit Worten ausreden. Und ist wahrlich wohl zu verwundern, daß solche greuliche Fälle auch den heiligen Leuten begegnen sollen, da der Heilige Geist wohnt, redet und sie regieret, da die rechte Kirche ist, die da rein und heilig ist: und vermag doch der Teufel daselbst auch so viel auszurichten, daß er solchen greulichen Lärm anrichtet. Es war aber dazumal im Hause Jakobs kein ander Weib, die da Rahels, so gestorben war, Statt hat kommen mögen, von welcher man die Erben, von Gott verheissen, hat erwarten mögen, denn diese Bilha allein; und die wird nun beschlafen und befleckt von seinem Sohne. Was könnte doch einem Menschen in diesem Leben Traurigeres und Greulicheres widerfahren?

179. So ist er nun bisher Jakob genannt worden: hier wird er erst Israel genannt in der Historie. Da Mose sagt: „Und Israel zog aus, und richtete eine Hütte auf“ etc.; item: „Da Israel im Lande wohnete“ etc. Diesen Namen prägt uns die Schrift mit sonderlichem Fleiß ein, die schändliche Blutschande damit groß und schwer zu machen. Denn sie will damit zu verstehen geben, daß Ruben diese Sünde begangen habe, da sein Vater herrlich verklärt war, da ihm der Name verändert, und die Verheißung gegeben war, daß er fruchtbar sein und gemehret werden, und daß er Könige zeugen sollte. Nach solcher großen Herrlichkeit geht Ruben hin und beschläft Bilha, und verunreinigt die vornehmste Hütte der Rahel. Darum wird Jakob an diesem Ort nicht vergeblich Israel genannt. Denn es predigest den Heiligen Geist, daß der Vater, den Gott also erhöht hat, so schändlich von seinem Sohne soll beflecket werden an seinem Weibe, welches an der verstorbenen Rahel Statt gekommen war.

180. Es ist wahrlich Wunder, daß er so lange hat leben können. Er hat aber im Wort der Verheißung gelebt, sonst wäre es unmöglich gewesen, so viel und so schwer Unglück und Jammer zu dulden und zu überwinden; jetzt lernt er, was wachsen und gemehret werden heiße, nämlich, daß es so viel heiße, daß er gar kein Kind mehr zeugen könne; denn Rahel ist gestorben, von welcher er viele Kinder zu bekommen gehofft hatte, Bilha aber ist befleckt und besudelt worden. Darum hat Jakob hernach allein ohne Weib gelebt, und hat sich zu Bilha, die befleckt war, nicht mehr gehalten, noch auch sonst zu einem andern Weibe. Darum setzt der Heilige Geist bald hernach die Kinder Jakobs nach einander gleichsam in einem Register, daß er damit zu verstehen gebe, daß ihrer nicht mehr gewesen sind denn zwölf; dabei soll es bleiben. Also hat David auch um der Blutschande willen, so Absalom begangen hat, ohne Weib gelebt bis an das Ende seines Lebens. Und daß die Leute so häufig der Weiber enthalten nicht so sehr Alters halben, als wegen ihres Schmerzes und daß ihm solche schändliche That seines Sohnes und Weibes so überaus wehe gethan hat.

\subsection{II}

\begin{quote}
\textbf{V. 22–26. Es hatte aber Jakob zwölf Söhne. Die Söhne Leas waren diese: Ruben, der erstgeborne Sohn Jakobs, Simeon, Levi, Juda, Isaschar und Sebulon. Die Söhne Rahels waren: Joseph und Benjamin. Die Söhne Bilhas, Rahels Magd: Dan und Naphtali. Die Söhne Silpas, Leas Magd: Gad und Asser. Das sind die zwölf Söhne Jakobs, die ihm geboren sind in Mesopotamien.}
\end{quote}

181. Im Hebräischen steht: Die Söhne Jakobs, der geboren ist. Da wird die Einzahl für die Mehrzahl gesetzt. Er ist zuvor Israel genannt worden wegen des greulichen Lasters; nun wird er wiederum Jakob genannt. Und Mose will so viel sagen: Also endet sich die Kinderzeugung des Patriarchen Jakob, des unvergleichlichen Mannes, der die Verheißung hat, daß er fruchtbar und gemehret werden soll und daß aus seinen Lenden Könige kommen sollen. Wie wird aber nun die heilige Schrift erfüllt? Da Benjamin geboren und die Mutter über der Geburt gestorben ist, hat Jakob aufgeführt Kinder zu zeugen und hat nur zwölf Söhne: wie wird er denn gemehret? Hier siehst du, wie Jakob aus solchem Jammer und Trübsal den rechten Verstand der Verheißung genommen hat. Denn da Rahel gestorben und ihre Hütte verunreinigt war, hat er verstanden, daß die schon geboren wären, von welchen die Könige herkommen würden, nämlich, Joseph und Juda, unter welchen dieser der Vater der Könige in Juda ist, der Andere aber der Vater der Könige in Israel. Dieselben sind aus den Lenden Jakobs gekommen, und die von ihnen geboren sind, gelten auch als Nachkommen Jakobs. Also ist Jakob durch seine Trübsal gelehrt worden, daß er sich hiefür des Heirathens enthalten sollte; darum hat er sich zu Lea auch nicht gehalten. So haben wir nun bisher an Jakob ein trefflich Exempel des allergrößten Unglücks und Trübsal gehabt und einer greulichen Sünde an Ruben.
\quellseite{258}

182. Dieses sind aber nun die zwölf Patriarchen, welcher Heiligkeit in der Schrift gerühmt wird, die Väter gewesen sind des heiligen Volkes Israel. Siehe aber, in wie große scheußliche Laster und Sünde sie der Teufel eingewickelt hat. Denn er macht aus Ruben einen solchen Sohn, wie man nicht leichtlich in den Historien der Heiden einen finden wird. Denn es haben auch die Heiden einen Greuel vor dieser Sünde; wie St. Paulus 1 Cor. 5, 1. sagt: „Es gehet ein gemein Geschrei, daß Hurerei unter euch ist, und eine solche Hurerei, da auch die Heiden nicht von zu sagen wissen, daß einer seines Vaters Weib habe.“ Daselbst hat sich aber zugetragen in der Gemeinde des sehr heiligen Patriarchen Jakob; darnach ist Mord und greuliche Lüge gefolgt. Denn die andern Söhne alle, ausgenommen Joseph und Benjamin, vornehmlich die Söhne von Lea, sind Mutter- und Vatermörder. Denn so viel an ihnen ist, tödten sie ihren alten Vater Jakob mit ihren greulichen Sünden, dadurch er mehr bestürzt und betrübt worden, denn sonst aller andern Trübsale halben. Sie erwürgen ihren Bruder Joseph auch, indem dem Vater seinen Rock, als ob er wahrhaftig erwürget wäre, und trösten ihn nicht allein gar nicht, sondern spotten noch des frommen alten Vaters dazu.

183. Dieses wird uns aber darum vorgehalten, auf daß dadurch die unendliche Langmuth und unschätzbare Gnade und Barmherzigkeit Gottes gerühmet und gepreiset werde, damit niemand seiner großen und vielen Sünden halben erschrecke und verzage, oder den Glauben und Vertrauen zu Gott fallen lasse. Denn diese Patriarchen sind auch schändliche böse Buben und große Sünder gewesen. Und es hat ihnen am Wort nicht gefehlet; dazu am Lehrmeister, an der Zucht oder Disciplin und guten Exempeln ihres Vaters hat es auch nicht gemangelt. Sie haben an Deborah eine Zuchtmeisterin gehabt; sie haben gesehen, wie ihnen Gott so mancherlei Hülfe und Wunderwerke erzeigt hat; sie haben gehört, daß den Vätern so viel Verheißungen geschehen waren. Was hat aber Ruben und die Andern gelernt, da sie so viel schöne Predigten gehört haben? Es sind böse Buben gewesen, und sie sind auch hart genug gestraft worden und haben es dürfen büßen müssen.

184. Derhalben gehört dies alles zu unserer Lehre und Strafe, daß wir uns vor dem Zorn und Strafe Gottes fürchten lernen. Darnach dient es auch zum Trost, wo Etliche sind, die böse ungehorsame Kinder haben, daß sie solchen Unfall und Trübsal mit Geduld tragen, sintemal wir ja nicht besser sind denn die Patriarchen. Denn ein Jeder unter ihnen ist mit greulichen Sünden befleckt gewesen: Asser, Dan, die Benjaminiten (wie das die zwei letzten Capitel im Buch der Richter bezeugen) und noch viel mehr die Ephraimiten. Denn der Teufel setzt der Kirche und Gemeinde Gottes viel härter zu denn den Fremden.

\subsection{III}

\begin{quote}
\textbf{V. 27—29. Und Jakob kam zu seinem Vater Isaak gen Mamre in die Hauptstadt, die da heißt Hebron, da Abraham und Isaak Fremdlinge innen gewesen sind. Und Isaak ward hundert und achtzig Jahr alt, und nahm ab, und starb, und ward versammelt zu seinem Volk, alt und des Lebens satt. Und seine Söhne Esau und Jakob begruben ihn.}
\end{quote}

185. Hier wird abermals der Name Jakob gesetzt. Aber droben, V. 22., hat er den Namen Israel der Sünde Rubens entgegen gesetzt: und scheint es, als sei der Patriarch etwas langsam fortgezogen mit seinem Vieh, und was er sonst vom Gesinde und Geräthe mehr gehabt; oder aber daß er selbst entweder zu Fuß, oder auf einem Kameele vorhergezogen sei, und seinen Vater etlichemal besucht habe, ehe er das ganze Gesinde zu ihm gebracht hat.

186. Mamre ist der Eigenname eines Mannes, eines der drei Brüder, die mit Abraham im Bunde gewesen sind, wie droben im 14. Cap. V. 13. steht. Von diesem Mamre, dem Erstgebornen, hat die Stadt den Namen, darum daß entweder er, oder seine Nachkommen zu Hebron regiert haben.

187. Isaak aber, der jetzt alt und blind war, hat Jakob und seine Kindeskinder nicht gesehen; er hat nur gehört, wie er ihm erzählt hat die Anfechtung, Jammer und Trübsal, die er bisher nach einander gelitten hatte, und daß Deborah und Rahel gestorben wären etc. Da hat ihn solche große Trübsal gejammert, und hat seinen Sohn getröstet und gesagt, daß solches eitel häßlich und teuflisch Unglück sei; aber Gott sei mächtiger denn der Teufel, der wolle nicht zulassen, daß wir über unser Vermögen versucht werden, 1 Cor. 10, 13. Denn der in uns ist, der ist größer, denn der in der Welt ist etc. Und sein Vater selbst hat in seinem ganzen Leben immer mit dergleichen Beschwerung und Trübsal auch streiten müssen, hat derhalben mit sonderlichem Ernst und mit Furcht seinen Sohn trösten können.

188. Da er aber nun daheim etliche Jahre mit seinem Vater zusammen in Ruhe gelebt, stirbt der Vater Isaak und ist nun Jakob allein noch am Leben: die andern Väter alle von der Sintfluth bis auf diese Zeit sind gestorben. Derhalben liegt die ganze Last des Testaments allein auf Jakob, der noch lebt, bei so großer Bosheit seiner Söhne, und im Gebet beharrt mit Lehren, Ermahnen und Drohen wider den Teufel, der sehr getobt und gewüthet hat. Er hat in seinem hohen Alter große Mühen tragen müssen, und hat nun an Isaak oder an einem andern der Patriarchen keinen Gesellen mehr.

189. So stirbt nun Isaak in stillem und geruhigem Alter, da er des Lebens satt war: nicht wie die Menschen, so immer in Wohlust leben, sondern ist dieses Lebens überdrüssig, und von großer Mühe und vieler Trübsal gar entkräftet gewesen. „Ich bin, wird er gesagt haben, des Lebens müde, ich mag es nicht mehr; darum hat er gebeten, daß ihm Gott erlösen und ihn endlich in den Hafen und besser Leben führen wolle.“

190. Diese Art aber und Weise zu reden: „Er ist zu seinem Volk versammelt worden“, zeigt an und bezeugt die künftige Auferstehung der Todten, sintemal ein Volk ist, zu dem wir versammelt werden. Denn wenn wir sterben, so verschwinden wir nicht in die Luft. Darum sagt der Heilige Geist nicht, er ist verschwunden, da er nicht mehr gelebt hat, sondern „er ist versammelt“, nicht zerstreut, hin- und her geworfen, oder geplagt worden, wie in diesem elenden und armseligen betrübten Leben, sondern er ist erlöst von allem Uebel, ist zu seinem Volk versammelt worden; wie die andern Väter, die im Frieden schlafen, welche Gott in seinen Schooß versammelt hat, da sie fein stille liegen und ruhen. Das ist der rechte Verstand dieser Worte, welche sehr nachdrücklich gesetzt werden. Es ist ein Volk der Todten, unter welchem sind Adam, Seth, Abraham, Debora u. s. w. Das Volk ist versammelt im Schooß und Armen Gottes, da sind sie in guter stiller Ruhe und werden zu seiner Zeit wiederum aufstehen.

191. Darum zeigt nun die heilige Schrift mit dieser Figur und Art zu reden an, daß die Väter gestorben sind, nicht wie die Heiden sterben, sondern daß sie versammelt und bewahret sind in der Hand Gottes. Wie der Prophet Jesaia im 57. Cap. V. 2. sagt: Die Gerechten „kommen zum Frieden, und ruhen in ihren Kammern“, denn sie haben richtig vor sich gewandelt. Und wird eine Stunde kommen, da sie wieder erscheinen und aus den Gräbern gehen werden, schöner als die Sonne und Sterne, nämlich, „die richtig vor sich gewandelt haben“. Ruben ist zu der Hoffnung der künftigen Auferstehung nicht gekommen, er sei denn durch Buße gereinigt worden: wie Paulus ein Verfolger und Feind Christi gewesen ist, der viel unschuldiges Blut vergossen hat; er hat aber dieselbige Sünde mit dem Blute des Sohnes Gottes abgewischt.

192. Darum soll man solche Stellen in der Schrift fleißig merken, mit den Worten, die der Heilige Geist gebraucht, daß Gott die Heiligen nicht verwirft oder zerstreuet, sondern sie also versammelt, daß auch nicht ein Bein oder Haar von ihnen umkommen muß.

193. Es kommen aber die Brüder, Esau und Jakob, zusammen zum Begräbniß, und Esau wird im Texte vor Jakob gesetzt. Und ist daran kein Zweifel, er werde auch zuvor seinen Vater besucht haben, und daß ihn der Vater auch oftmals liebreich gestraft und vermahnt habe, daß er den Haß und die Nachgier ablegen sollte. Darum kommt er zum Begräbniß, daß er damit seinen Gehorsam und Ehrerbietung gegen den Vater beweisen möge, und ist ein gewiß Zeichen, daß er mit seinem Bruder ausgesöhnt gewesen, und daß er sich zu der rechten Kirche und Gemeinde gehalten hat, auf daß er der geistlichen Verheißung auch theilhaftig werden möchte, wo nicht aus der Verheißung, doch gleichwohl aus Gnaden; wie wir Heiden zu Gnaden angenommen werden, nicht daß uns solches versprochen wäre, sondern aus lauter Gnade und Barmherzigkeit.
\quellseite{259}

\chapter{Das sechsunddreißigste Capitel}

\section[Erster Theil]{Erster Theil. Von Esau und seinem Geschlechte überhaupt, und von Esaus Weibern und seinen Kindern insonderheit}

1. In diesem Capitel muß man die Ordnung der Historie ansehen; denn Mose hat die Figur gebraucht, die hysteron proteron genannt wird. Denn was nach diesem Capitel folgt bis auf das 40. Capitel, das sollte vorher gesetzt sein worden, vor die Worte: „Und Isaak ward hundert und achtzig Jahr alt.“ Denn alles, was in den nächstfolgenden drei Capiteln beschrieben wird, ist vor dem Tode Isaaks geschehen, welcher geschehen hat, daß Dina beschlafen worden, daß die Sichemiten erwürgt und die Söhne Jakobs damit einen blutigen Sieg erlangt haben; item, daß Debora und Rahel gestorben und über das auch Joseph nach Egypten ist verkauft worden, und zu dem auch die Blutschande, so Juda mit Thamar begangen hat.

2. Dieses alles kann man aus der Zeitrechnung abnehmen. Denn sonst scheint es, daß Isaak alsbald gestorben sei, da Jakob zu ihm gekommen ist; welches falsch ist. Denn er hat zum wenigsten noch dreißig Jahre gelebt, da Jakob wieder aus Mesopotamien gen Hebron gekommen ist. Darnach ist es gewiß, daß Isaak ein hundert und achtzig Jahre alt gewesen, da er gestorben ist. Esau aber und Jakob sind geboren worden im sechzigsten Jahre des Alters Isaak. Wenn dieselbigen von den hundert und achtzig Jahren abgezogen werden, so bleiben noch hundert und zwanzig Jahre; das ist das Alter Esaus und Jakobs, da Isaak stirbt. Joseph ist dreißig Jahre alt gewesen, und Isaak ist gestorben dreizehn Jahre, nachdem Joseph nach Egypten verkauft gewesen, im letzten Jahre vor der Erlösung Josephs aus dem Gefängnisse.

3. Es ist aber zumal ein erbärmlich Leben gewesen dieses alten Mannes, der seines Gesichtes und Lichtes beraubt gewesen von der Zeit an, da er Jakob gesegnet hat, von den sieben und siebzigsten Jahre seines Alters. Denn dazumal waren ihm seine Augen schon dunkel geworden, wie wir gehört haben. Darum sollen wir merken, daß die drei nachfolgenden Capitel zwischen diese Worte sollen gesetzt werden, da Mose sagt: „Jakob kam zu seinem Vater Isaak“, und: „Isaak ward hundert und achtzig Jahr alt.“ Denn die Zeitrechnung erinnert uns vieler Dinge mit Nutzen und erklärt die Historien, gibt dazu oftmals Ursache zu vielen schönen Gedanken. Denn die frommen Patriarchen und Propheten sind rechte heilige Märtyrer gewesen, die viel Jammer und Elend erlitten haben, dazu wir mit nichten stark genug sind, solches zu leiden und zu ertragen. Dieses habe ich zuerst von der Historologia, das heißt, davon sagen müssen, daß das Vorherste zu hinderst gesetzt ist. Denn dieweil Mose im vorigen Capitel gesagt hat, daß Esau und Jakob zusammen gekommen sind, ihren Vater zu begraben, hat er gedacht, es wollte sich an diesem Orte sehr wohl reimen, daß er von Esau rede, denn er in diesem Capitel seinen endlichen Abschied gibt und dessen er hernach nicht mehr gedenken wird.

\subsection{I}

\begin{quote}
\textbf{V. 1. Dies ist das Geschlecht Esau, der da heißt Edom.}
\end{quote}

4. Der Name Edom ist stets in diesem Lande geblieben. Die Lateiner und Griechen nennen es Idumaea, wie Virgilius Psalmen aus Idumäa nennt. Die geistlichen und pestilenzialischen Juden aber nennen in allen ihren Auslegungen das Römische Reich und die Christenheit Edom, so doch sie wahrhaftig das vermaledeite und verdammte Volk sind. In ihren Gebeten beten sie wider Edom und bitten, daß ihr Messias uns wolle erwürgen und umbringen, und haben über die Maßen närrische Meinungen vom Ursprung des Römischen Reiches. Mose aber hat Esau nicht vergehen wollen, dieweil er Abrahams Kindeskind ist, und Esau und Jakob sind einander der Geburt und des Geblütes halben gleich. Esau ist geboren von dem sehr edlen Blute Sarahs und Rebekkas, Abrahams und Isaaks. Darum hat ihn Mose nicht wollen verwerfen, sondern erzählt nach einander seine Kinder und Kindeskinder, und gedenkt des leiblichen Segens, davon droben 1 Mos. 27, 40. ist gesagt worden: „Deines Schwertes wirst du dich nähren; und es wird geschehen, daß du deines Bruders Joch von deinem Halse reißen wirst“ etc.

5. Im Anfange dieses ersten Buches Mose ist auch von dem Geschlechte Kains gesagt worden. Wiewohl dieser verworfen, verbannt und verflucht war auf Erden, hat ihm doch Gott Raum gegeben zur Buße, daß er und seine Nachkommen die Wohltaten der geistlichen Verheißung erlangen möchten. Also wird an Esau auch gepriesen das Exempel göttlicher Geduld, auch an denen, die verworfen sind. Denn daran ist kein Zweifel, Esau und Ismael werden die Lehre Abrahams, Isaaks und Jakobs mit sich genommen haben, desgleichen auch die Ordnung und Weise der Opfer und vornehmlich der Beschneidung. Darum sind viel vortreffliche und herrliche Männer vom Geblüte Esaus geboren. Aber darin ist der Unterschied, daß sie die Verheißung des Messias nicht gehabt haben. Die Söhne und der Heiland der Welt ist ihnen nicht verheißen worden; sie sind aber darum nicht verworfen, weder sie, noch ihre Nachkommen, wie andere Heiden. Viele Mohren, Ammoniter und Edomiter haben sich zum rechten Bekenntniß und Religion des Gottes Israel begeben nach dem Beruf Gottes. Viele von den Midianiten und Babyloniern sind selig geworden durch Gottes Gnade und Barmherzigkeit, so ihnen zufällig widerfahren ist, wiewohl Christus den Heiden nicht verheißen ist wie dem Samen Abrahams. Es ist ihnen auch Gottes Wort nicht vertrauet worden; denn diesen Ruhm müssen wir dem Samen Abrahams geben, davon der 147. Psalm V. 20. sagt: „So thut er keinen Heiden“ etc. Daher hat das Volk Israel Patriarchen, Propheten, Gottes Wort; und St. Paulus sagt Röm. 15, 8. 9.: „Ich sage aber, daß Jesus Christus sei ein Diener gewesen der Beschneidung um der Wahrheit willen Gottes, zu bestätigen die Verheißung, den Vätern geschehen. Daß die Heiden aber Gott loben um der Barmherzigkeit willen, wie geschrieben stehet: Darum will ich dich loben unter den Heiden, und deinem Namen singen“ etc. Aber vom Samen der Heiden ist Christus nicht gekommen.

6. Darum muß man einen Unterschied machen zwischen der Verheißung und der Wahrheit, und zwischen der Gnade und Barmherzigkeit, so zufällig kommt. Wir haben Barmherzigkeit ohne Verheißung, die Juden haben sie mit der Verheißung; „denn das Heil kommt von den Juden“, Joh. 4, 22., der Ueberwinder des Todes und des Teufels kommt vom Stamme Juda. An dieser Verheißung hat es Esau und Ismael gemangelt: von der Barmherzigkeit aber sind sie nicht ausgeschlossen; denn die Historien bezeugen das Gegentheil. Wie auch die nicht alle, so diese Verheißung gehabt, sich zu der Gemeinde Gottes gehalten haben; denn die Kinder nach dem Fleische sind nicht Gottes Kinder, sondern die Kinder der Verheißung. Ja, in dem Volke sind viel grobe fleischliche Laster und greuliche Sünden geschehen, dergleichen auch in den Historien der Heiden nicht gelesen werden, als, daß David von seinem eigenen Sohne aus dem Lande verjagt wird, und das ganze Haus mit vielen Morden befleckt und unglückseliger ist, denn des Cantalus oder der Pelopiden im Griechenlande.

7. Darum folgt das nicht, daß man sagen wollte: Wir haben die Verheißung, darum sind wir alle Erben derselben. Wie das auch nicht folgt: Wir haben die Verheißung nicht, darum sind wir verworfen. Gott hat aus allen Völkern etliche angenommen, auf daß wir uns nicht rühmen sollen; sondern, wer sich rühmet, der rühme sich des Herrn“, 1 Cor. 1, 31., der sowohl, der die Verheißung hat, als der sie nicht hat, „auf daß Gott allein gerecht sei, und gerecht mache den, der da ist des Glaubens an Jesu“. Röm. 3, 26.
\quellseite{260}

8. So beschreibt nun Mose das Geschlecht Esaus nicht allein um der Ehre Abrahams willen, Isaaks und Jakobs, welcher Kindeskind und Sohn er ist, sondern vielmehr um der Ursache willen, daß wir wissen sollen, daß er des Segens theilhaftig geworden ist, den Jakob gehabt hat. Denn er hat allezeit einen freien Zutritt gehabt, sich zu der rechten Kirche und Gemeinde Gottes zu begeben, wenn er sich nur demüthigt und mit Ernst gesagt hat: Gott wolle sich mein nur erbarmen, und mich lassen theilhaftig werden der Brosamen, die da fallen vom Tisch der Gemeinde, welche im Hause Gottes ist. Solche Demuth und Gehorsam gefällt Gott wohl, wiewohl du die Verheißung nicht hast.

9. Dagegen aber wenn die Kinder der Verheißung stolz werden, verlieren sie den Segen. Denn Gott will haben, daß St. Petrus sagt 1 Petr. 5, 6.: „So demüthiget euch nun unter die gewaltige Hand Gottes, daß er euch erhöhe zu seiner Zeit“ etc. Und um der Ursache willen demüthigt und versucht er seine Heiligen also, daß sie nicht stolz und hoffärtig werden. Also halte ich dafür, daß unter den Nachkommen Edoms viel treffliche Männer gewesen und die auch voll Heiliges Geistes gewesen sind. Wiewohl es auch an solchen nicht gefehlt, die dem Volk Israel sehr und bitter feind gewesen sind. Wie heutiges Tages die Papisten uns auch spitzfeind sind, haben einen Greuel an uns, halten uns für verfluchte und verdammte Leute: jedoch sind ihrer auch viele, die uns gewogen sind und uns lieb haben, und unsere Lehre annehmen. Und dieses, glaube ich, sei die Ursache gewesen, warum Mose das Geschlecht Esaus auch mit unter die Zahl der Heiligen gesetzt hat, und hat beide ihn und etliche seiner Nachkommen mit der Gemeinde der Heiligen zusammengefüget und verbunden; denn sie sind theilhaftig geworden nicht der Verheißung, sondern der Barmherzigkeit, welche die Verheißung gibt.

10. In der lateinischen Uebersetzung ist dieses Capitel sehr dunkel und verworren. Mose aber beschreibt erstlich die Nachkommen Esaus, so bei seinem Leben gewesen sind; darnach die Schwägerschaft oder das Geschlecht der Weiber. Zum Dritten, die Fürsten und Könige. Unter denselben, zeigt er an, sind etliche Fromme gewesen, die dadurch, daß sie das Wort lange Zeit gehört haben, bekehrt worden sind. Denn Seir ist nicht weit von Hebron gelegen gewesen, darum haben sie leichtlich zu ihrem Vater Isaak und zu seiner Gemeinde gehen können. Und wenn das Wort irgend an einem Ort gegenwärtig wird, so breitet es seinen Glanz aus und kommt hin und wieder zu den Nachbarn. Darum gibt Mose zu verstehen, daß nicht Alle verdammt seien, die von Edom hergekommen sind, wiewohl er von der Verheißung verworfen gewesen und die Erstgeburt verloren hatte.

11. Und ist wahrlich ein wunderbares Ding, daß fast in allen Historien der heiligen Schrift die Erstgebornen unglücklich sind. In Egypten sind alle Erstgebornen erwürgt worden. Es scheint aber, daß dieses die Ursache sei, dieweil das Gesetz eine Bedingung bei sich hat, und die das Gesetz haben, werden gemeiniglich ärger; denn sie werden stolz um des Gesetzes willen, als ob es darum gegeben sei, daß es der Erstgeburt zur Hoffart helfen sollte, so es doch darum gegeben ist, daß man sich in der Erstgeburt demüthig halten soll. Darum hat Gott den Stolz und Hoffart der Erstgebornen unterdrückt, wenn sie solcher Ehre und Herrlichkeit mißbraucht haben. Denn alles, was geschieht und gesagt wird, geht auf den einigen erstgebornen Sohn Gottes, Jesum Christum, dem gebühret allein alle Ehre, Kraft und Gewalt. Nun wollen wir auf das Geschlecht Esaus kommen.

\subsection{II}

\begin{quote}
\textbf{V. 2. 3. Esau nahm Weiber von den Töchtern Kanaans, Ada, die Tochter Elon, des Hethiters; und Aholibama, die Tochter Ana, die Nichte Zibeons, des Heviters; und Basmath, Ismaels Tochter, Nebajoths Schwester.}
\end{quote}

12. Eben was hier gesagt wird, ist droben im 26. Cap. V. 34. auch gesagt worden, aber mit andern und veränderten Worten. Denn das erste Weib Esaus wird daselbst Judith genannt und ihr Vater Beeri, ein Hethiter, wird hier Elon genannt. Denn die Veränderung der Namen ist in allen Sprachen sehr gewöhnlich, darum sind von Andern etliche nützliche Regeln gegeben worden, dadurch viel dunkle Texte ausgelegt und erklärt werden. Die erste Regel ist: Alle Menschen haben gemeiniglich zwei Namen, als, Martinus Luther, Philippus Melanchthon. Die andere Regel ist: Die Väter und Mütter, Söhne und Töchter sind solche von Natur oder nach dem Gesetz. Die dritte Regel: Es kommt oft vor, daß Namen gegeben und verändert werden zufälliger Weise, als, Jakob wird Israel genannt von der Geschichte, die sich mit ihm zugetragen hat. Die Andern haben von ihren Aeltern Namen bekommen, als, Pastor, Pfarrherr; Judex, Richter; Consul, Bürgermeister. Es kommt auch oft vor, daß nach dem Alter Veränderungen geschehen: die Jungfrauen pflegen wir anders zu nennen in ihrer Kindheit, denn wenn sie erwachsen sind. Bisweilen bekommen die Leute Namen von der Kleidung, von der Farbe und Gestalt, bisweilen auch von ihrer Gesellschaft. Als, Eusebius ist genannt worden Eusebius Pamphili um seines Gesellen willen in der Marter.

13. Also sollen wir denken, daß eben dasselbe auch in diesem Volk oftmals geschehen sei. Judith wird hier Ada genannt, und kann wohl sein, daß sie vom natürlichen Vater zu einem Stief- oder Pflegevater gekommen ist und daß daher der Name verändert worden sei. Ja, siehe unsere Exempel an im Papstthum, wie da so eine große Thorheit gewesen ist, die Namen derjenigen zu verändern, die dem Papst einen Eid geschworen hatten, seine Satzungen zu halten. Ich bin in der Taufe Martinus genannt worden, darnach im Kloster Augustinus. Was könnte doch Schändlicheres oder Ungöttlicheres geschehen, als daß man den Taufnamen verwirft und fahren läßt um der Kappe willen, die er angezogen hat? Also schämen sich die Päbste, daß sie mit ihren rechten Namen sollen genannt werden, denn sie in der Taufe erlangt haben; denn sie sind zwar dessen auch nicht werth, daß sie den Taufnamen behalten sollen, darum erhalten sie andere Namen, die sie sich selbst geben. Julius der Zweite ist zuvor Julianus genannt worden. Mit solcher Veränderung der Namen haben sie angezeigt, daß sie von Christo und der Taufe abgefallen und abtrünnig geworden sind, und dasselbe ist in allen Klöstern sehr gemein gewesen.

14. Also ist das erste Weib Esaus Judith genannt worden, das ist Judaea, Jüdin, auf Lateinisch Laudatrix, Lobfängerin. Wie Lea ihren Sohn Juda nennt darum, daß sie die Wolthat Gottes erkannt hat und sich dankbar hat erzeigen wollen. Es ist ein schöner Name gewesen. Es scheint aber, daß Judith unartig geworden sei, und daß daher die Veränderung des Namens gekommen sei, aber nicht zum Besseren. Denn Ada heißt, Halsband, Schmuck. Sie hat Lust gehabt zu schöner Kleidung und weiblichem Schmuck; da sie eine Hausfrau geworden ist, hat sie wollen Ada sein: die geputzte Frau, Sammetfrau, Kettenfrau etc. Also hat sie ihren schönen Namen verloren, daß sie nicht mehr Judith, sondern Ada ist genannt worden, als die für und für hat vor dem Spiegel sitzen und sich immer putzen und schmücken wollen.

15. Das andere Weib Esaus ist Ahalibama. Das hebräische Wort ohel heißt Hütte; wie Mose kurz zuvor gesagt hatte, daß Jakob eine Hütte aufgerichtet habe jenseit des Thurms Eder. Im Propheten Hesekiel am 23. Cap. V. 4. steht ohlah, das heißt, ihre Hütte, nämlich des Weibes; denn Gott redet daselbst mit großem Zorn wider die Hure, die eine eigene Kirche oder Gemeinde hat und das Gesetz und Priestertum verworfen hat. Mose aber zeigt ferner an, daß dieses Weib die rechte Frau, domina, und Königin gewesen sei. Droben, Cap. 26, V. 34., ist sie Basemath genannt worden; als wenn du sagen wolltest auf Lateinisch, aromatica, suave olens, ein wohlriechendes Weib von köstlicher Würze oder Specerei; denn sie hat einen guten Geruch und einen guten Namen gehabt; ein fein Jungfräulein, ein Würzkränzlein, ein Nelkenblümlein. Jetzt wird sie Ahalibama genannt, die Schloßfrau, a tabernaculo excelso, die von der Höhenburg, gnädige Frau.

16. Das dritte Weib heißt Basmath, ist droben, Cap. 28, V. 9., Mahalaθ genannt worden. Es scheint, als habe Mose die Namen der Frauen der drei Weiber Esaus mit Fleiß also verändert. Dies ist Ismaels Tochter, vom Geblüt Abrahams, des Edlen sie gewesen ist, und ist mit Jakob und Esau Geschwisterkind im dritten Grade. Das Wort an sich selber heißt so viel als, krank, schwach; und Mose gibt damit zu verstehen, daß es eine Dirne gewesen
\quellseite{261}
sei, die schwach gewesen. Darum haben sie gesagt, daß sie nicht lange leben würde, und haben ihr davon den Namen gegeben und sie Schwächlein geheißen, aegrotula. Dies Weib allein verändert ihren Namen zum Bessern; denn jetzt wird sie Basmath genannt, das ist eine wohlriechende balsamische Rose und Blume. Wie die zwei vorigen Weiber untüchtig geworden sind und ihre guten Namen verloren haben: also hat sich dieses dritte Weib gebessert und ihr Leben göttlich geführt, aber nicht ohne Kreuz und Jammer.

17. Droben im 26. Cap. V. 35. ist gesagt worden, daß die beiden Weiber Esaus gegen Isaak und Rebekka morath, das ist, sehr bitter gewesen sind, und die frommen heiligen Menschen und das ganze Gesinde wohl geplagt haben. Davon hat Judith leichtlich können verändert und eine Ada aus ihr werden, da das Gesinde wird gesagt haben: Sie ist Judith nicht, ihre Sitten und Leben reimen sich mit dem Namen nicht; sie ist Ada, eine Tanzfrau. Item, die Andere hat einen bösen Geruch und ist tyrannisch. Die Dritte ist später in das Haus genommen, und ist mehr gelobt und geliebt worden von jedermann. Nun folgen die Kinder dieser Weiber Esaus.

\begin{quote}
\textbf{V. 4. 5. Und Ada gebar dem Esau Eliphas, aber Basemath gebar Reguel. Ahalibama gebar Jehus, Jalam und Korah. Das sind Esaus Kinder, die ihm geboren sind im Lande Canaan.}
\end{quote}

18. Ada ist die Mutter des berühmten Mannes Eliphas, des erstgebornen Sohnes Esaus, welcher ein Urenkel Abrahams und ein Enkel Isaaks und ein Sohn Esaus ist. Viele meinen, er sei einer von den drei Freunden Hiobs. Wenn aber das wahr ist, so muß Hiob gelebt haben lange zuvor, ehe die Kinder Israel aus Egypten gezogen sind. Denn die Zeitrechnung zeigt an, daß Eliphas siebenzig oder achtzig Jahre alt gewesen ist, da Isaak gestorben ist. Da Esau seinem Bruder Jakob entgegen gezogen, da ist Eliphas zum wenigsten sechzig oder fünf und fünfzig Jahre alt gewesen. Denn Esau hat im vierzigsten Jahre seines Alters ein Weib genommen: von der Zeit an bis auf den Tod Isaaks sind achtzig Jahre. Aber die Meinung mit Hiob lasse ich unentschieden; denn man hat davon in der Schrift kein Zeugniß. Es reimt sich nicht übel, daß es irgend einer von den Urenkeln oder Enkeln gewesen sei, der des Vaters Namen wird behalten haben. Aber gebieret nur Einen Sohn, Basmath gebieret auch Einen. Von diesen zwei Weibern hat Esau zwei Söhne. Aber die Frau, domina, und Königin Ahalibama gebieret ihm drei Söhne: die hat es den Andern zuvor gethan, dieselbige eine Mutter gewesen ist dreier Söhne. Und ist wohl glaublich, daß sie vornehmlich im Hause und unter dem Gesinde Isaaks morath, das ist, sehr bitter gewesen sei.

19. Diese Kinder alle werden im Lande Canaan geboren, im Hause seines Vaters Isaak. Denn da Esau vierzig Jahre alt gewesen, hat er seine ersten zwei Weiber genommen, die Hethitinnen waren. Derhalben sind diese fünf Kinder geboren, da er noch bei seinem Vater war. Jakob ist nach Mesopotamien gezogen im sieben und fünfzigsten Jahre, darum sind sieben und dreißig Jahre zwischen der Heirath Esaus, bis ihm der Segen ist genommen worden. Denn so lange hat Isaak damit verzogen, und Eliphas und die andern Kinder Esaus sind auch so viel Jahre bei dem Großvater gewesen und haben ihn predigen hören. Denn haben sie nicht gesehen, sondern Eber und Debora. Also haben auch die Weiber Esaus den alten Isaak und die Rebekka gehört, und sind an ihr hundertstes Jahr gekommen, wo sie noch gelebt haben, da ihr Schwiegervater gestorben ist; und ist wohl glaublich, daß sie in so langer Zeit des Hasses und der Bitterkeit vergessen haben. Und ich halte dafür, daß Mose dieses darum schreibe, daß er damit will anzeigen, daß endlich ihre Herzen gelindert worden, da die Brüder mit einander versöhnt sind, und daß Eliphas ein guter frommer Mann gewesen sei, sintemal man gar nichts Böses von ihm findet an irgend einem Ort in der Schrift.

\begin{quote}
\textbf{V. 6–8. Und Esau nahm seine Weiber, Söhne und Töchter, und alle Seelen seines Hauses, seine Habe, und alles Vieh, mit allen Gütern, so er im Lande Canaan erworben hatte, und zog in ein Land von seinem Bruder Jakob. Denn ihre Habe war zu groß, daß sie nicht konnten bei einander wohnen; und das Land, darin sie Fremdlinge waren, mochte sie nicht ertragen, vor der Menge ihres Viehes. Also wohnte Esau auf dem Gebirge Seir. Und Esau ist der Edom.}
\end{quote}

20. Da hat sich die Uneinigkeit zwischen den Brüdern und der Haß Esaus gegen Jakob und seine Eltern erhoben, nachdem er den Segen verloren hat. Darum ist er in solchem Zorn gewichen, ist dazu auch bewogen worden durch die Menge des Viehes, da er gesehen hat, daß deshalb das Land sie nicht hat ertragen können, und ist also gen Seir gezogen, woher er Ahalibama zum Weibe genommen hatte. Derselbigen folgt er nach in ihr Vaterland. Ich halte aber dafür, daß er gen Seir gezogen sei zwei oder drei Jahre, nachdem Jakob nach Mesopotamien gezogen war; denn so lange ist er bei dem Vater geblieben. Und das Haus Isaaks ist dazumal in großer Aufregung gewesen, da Esau damit umgegangen ist, wie er seinen Bruder erwürgen möchte, und hat die Eltern damit sehr betrübt, daß er seinen Bruder also haßte. Und er hat mit Jakob das väterliche Erbe getheilt, wiewohl Jakob nicht gegenwärtig gewesen; die Eltern aber haben Jakob zwei Theile zugeeignet wegen des Rechtes der Erstgeburt.

21. Darnach zeigt Mose an vom leiblichen Segen, so Esau widerfahren ist, nämlich, da er der Menge des Viehes gewohnt. Denn er ist vom Geblüt eines sehr großen Patriarchen, und hat nichts verloren, ausgenommen die Verheißung vom Samen. Darum ist es gewiß, daß er sehr gemehret und reich gewesen sei. Ja, ich habe keinen Zweifel, daß Viele vom Geschlecht Esaus zu dem Gott Abrahams und Isaaks bekehrt worden sind, wie die Historie der Könige anzeigt. Er hat auch die Beschneidung und die Lehre der Väter von der Verheißung behalten.

\begin{quote}
\textbf{V. 9. 10. Dies ist das Geschlecht Esau, von dem die Edomiter herkommen, auf dem Gebirge Seir. Und so heißen die Kinder Esau: Eliphas, der Sohn Ada, Esaus Weibes; Reguel, der Sohn Basmath, Esaus Weibes.}
\end{quote}

22. Den Anfang dieses Capitels wiederholt Mose, welches er sonst auch zu thun pflegt; wenn er etwas dazu setzen will, so wiederholt er, was vorher ist gesagt worden. Also, da er im andern Buch am 6. Cap. V. 16. sein Geschlecht erzählt, setzt er hernach, daß er anfänglich auch gesagt hätte: „Dies sind die Namen der Kinder Levi in ihren Geschlechtern“ etc., auf daß es nicht scheine, als ob er aus einem andern Stamme geboren wäre denn von Levi. Dasselbe thut er hier auch, dieweil er zu dem Vorigen noch etwas mehr setzen und dasselbe weitläufiger erklären will.

23. Fünf Söhne sind Esau im Lande Canaan geboren worden, unter welchen Reguel noch ein kleines unmündiges Kindlein gewesen zu sein scheint, und daß er erst geboren worden, da Jakob nach Mesopotamien gezogen war. Sein erstgeborner Sohn aber Eliphas ist zu der Zeit sechs und dreißig Jahre alt gewesen, daß er nun alt genug war zum Ehestande. So lange haben aber die Hethitischen Weiber im Hause Isaaks und Rebekkas regiert und von ihren Gütern gelebt. Und Mose erzählt dies darum, daß er uns dadurch erinnern möge des Kreuzes und der Mühseligkeit, so die frommen heiligen Eltern getragen haben wegen der Hoffart und Herrschaft der Weiber ihres Sohnes, davon er droben, Cap. 26, V. 35., gesagt hat; wiewohl es glaublich ist, daß sie endlich mit der Zeit werden etwas gelinder worden sein.

24. Es pflegen aber sonst die Schwiegermütter von Natur ihrer Söhne Weiber zu hassen. Rebekka soll man aber dessen billig nicht beschuldigen; denn sie hat ihren Sohn Esau lieb gehabt und um seinetwillen seine Weiber und Kinder, als ihre Enkel. Und sie hat es nicht verdient, daß sie von den Hethitischen Weibern, deren jegliche eine stolze Frau, domina, gewesen ist, hat sollen also geplagt werden, welche Weiber sich nach den Sitten und der Gewohnheit ihres Vaterlandes gehalten haben; denn die Hethiter sind ein sehr gewaltiges Volk gewesen und das allerärgste unter den Cananitern. Im Propheten Hesekiel am 16. Cap. V. 3. steht: „Deine Mutter ist aus den Hethitern“ etc. Sie haben ihren Namen gehabt von Schreckens; denn sie sind Tyrannen und jedermann verhaßt gewesen: darum haben sie die heilige Mutter Rebekka bei vierzig Jahren sehr geplagt. Basmath ist von den zwei verachtet und geschmäht gewesen, und ist wie eine Magd gehalten worden, die ihnen hat auf dem Fuße nachfolgen müssen; und wird doch vor Ahalibama gesetzt, auf daß damit Mose anzeige, daß er und der heilige Geist den Hethitinnen feind sei. Sonst ist der Name eines Landes, deß Vater Esau genannt wird, und nicht irgend einer Person. Von Esau, der Edom genannt wird, hat dieses Land Seir den Namen bekommen.

\quellseite{262}

\begin{quote}
\textbf{V. 11–13. Eliphas Söhne aber waren diese: Theman, Omar, Zepho, Gaetham und Kenas. Und Thimna war ein Kebsweib Eliphas, Esaus Sohns, die gebar ihm Amalek. Das sind die Kinder von Ada, Esaus Weib. Die Kinder aber Reguel sind diese: Nahath, Serah, Samma, Missa. Das sind die Kinder von Basemath, Esaus Weib.}
\end{quote}

25. Da Jacob wiedergekommen ist, ist Eliphas sechs und fünfzig Jahre alt gewesen. Und diese Söhne sind ihm zumal schon geboren gewesen, und ich halte dafür, daß dieselben große Herren und Scharhansen gewesen sind, mit welchen Esau ausgezogen ist seinem Bruder Jacob entgegen und ihn zu schrecken: und wiewohl sie mit ihm versöhnt sind, so haben sie doch in dem Zuge ihre große Herrlichkeit und Gewalt sehen lassen. Theman ist in der heiligen Schrift berühmt, und ich halte dafür, daß das Land Theman von ihm den Namen bekommen habe, welches gedacht wird im Propheten Jeremia am 49. Cap. V. 7.: „Ist denn keine Weisheit mehr zu Theman?“ Das ist, in Idumäa; denn es ist ein Theil davon.

26. Thimna ist die Schwester gewesen eines von den Fürsten zu Seir, der Horiter. Und Ada und Thimna sind beide Hethitinnen, darum haben sie sich so viel leichter vertragen können. Und Amalek, der Sohn Thimna, ist von Ada, dem ersten Weibe Esaus, zum Sohne angenommen worden. Um der Ursache willen wird Amalek unter die Söhne Adas mitgezählt. Von Amalek aber ist das in den Historien der heiligen Schrift berühmte Volk hergekommen, davon 4 Mos. 24, 20. gesagt wird: „Amalek, die ersten unter den Heiden; aber zuletzt wirst du gar umkommen.“ Es hat eine sehr vornehme Herkunft. Esau, Isaak, Abraham sind seine Ahnen. Das ist ein trefflicher Adel gewesen, so viel das väterliche Blut belanget; und derhalben sind die Amalekiter auch sehr stolz und hoffärtig gewesen. Von der Mutter aber sind sie Hethiter gewesen.

\begin{quote}
\textbf{V. 14. Die Kinder aber von Ahalibama, Esaus Weib, der Tochter des Ana, der Neffe Zibeons, sind diese, die sie dem Esau gebar: Jehus, Jaalam und Korah.}
\end{quote}

27. Mose erzählt die Kinder und Enkel Esaus, auf daß er damit anzeige, was Esau für ein Kriegsvolk gehabt habe, da er Jacob entgegen gezogen ist.

Hier kommt aber eine Frage: Warum Ahalibama eine Tochter Ana und eine Tochter Zibeons genannt werde? Denn wie kann sie zweier Männer Tochter sein? Die Juden treiben ein wunderlich Spiel damit und haben zwei schändliche Meinungen, nämlich erstlich, daß Zibeon ein gottloser Mann und ein Verächter des Ehestandes gewesen sei, und darum habe er bei seines Vaters Seir Weibe geschlafen und den Ana gezeuget, also, daß Ana Zibeons natürlicher Sohn sei und für seinen Bruder sei gehalten worden, in Rücksicht auf Seir. Die andere Meinung der Juden ist, daß Zibeon seines Sohnes Ana Weib soll geschlafen haben, und soll also mit Beiden, der Mutter und des Sohnes Weibe, Blutschande getrieben haben. Darum halten sie Ahalibama für ein sehr blutschänderisches Weib, welcher Vater durch Blutschande von seinem Vater soll geboren sein.

28. Ich bin aber den Juden und den Auslegungen ihrer Rabbinen von ganzem Herzen feind. Denn dieses ist ihre Weise, ja, es ist also eine offenbare Unsinnigkeit an ihnen, daß sie, woher sie nur können und mögen, \unsicher{der ihren Loch und der heiligen Schau zusammenlesen}. Es sind die verfluchtesten Leute, und werden vom Teufel gefangen gehalten und besessen: und das Gemälde oder Gleichniß des Herrn Christ vom Splitter und Balken, Luc. 6, 41., reimt sich wohl auf sie; item, daß er auch sagt Matth. 23, 24. von denen, so die Mücken seigen und Kameele verschlucken. Denn daß der eine Pharisäer von sich selber sagt Luc. 18, 11.: „Ich bin nicht, wie die andern Leute“, dasselbe rühmen ganz unverschämt alle Juden von sich auch und sagen: Wir sind Gottes Volk, so von Gott mit großen Gaben geziert ist. Und einer unter ihren Rabbinen hat gesagt: Es sei Wunder, wo man unter allen heiligen Einen weisen Mann finden sollte. Im Volke Israel aber sollen eitel vortreffliche weise Leute sein.

29. Ja, es ist ein sehr schändlich und schädlich Volk, und ist's auch dazu damals schon gewesen bei der Väter Zeiten unter Jacob und David. Denn, lieber, siehe doch, was sie für herrliche Thaten begangen haben. Ruben sitzt auf der Bilha, seines Vaters Weibes, Lager gestiegen; Juda treibt mit seines Sohnes Weibe Blutschande. Eben solche Schellen hängen die Heuchler dem Ana, der Ahalibama Vater, auch an. Warum sehen sie aber nicht erst den Balken in ihren Augen? Denn was ist schändlicher, als der Aufruhr Absaloms wider seinen Vater David, dadurch er ihn aus dem Lande jagt, und so viel an ihm ist, ihn erwürgt, und schändet ihm seine Kebsweiber? Laß nun die Rabbinen hingehen und von der Sünden Weisheit rühmen. Danach ist es allezeit ein aufrührerisch Volk gewesen wider die Propheten und Könige, und haben wider den Herrn gemurrt, die Propheten aber auf das greulichste erwürgt. Darum bin ich den Auslegungen der Rabbinen sehr feind, darin sie die heilige Schrift mit ihren Lügen häßlich verfälschen.

30. Aus diesem Text, dieweil hier Ahalibama eine Tochter Anas und eine Tochter Zibeons genannt wird, wollen sie schließen, daß sie zwei Väter gehabt habe und daß Ana seine Mutter soll beschlafen haben etc. Dieses ist eine rechte jüdische, das ist, eine eselhafte Folgerung, welche auch von Lyra und von Pagninus verworfen wird. Daß nun solche Laster unter den Heiden gewesen sind, daran zweifelt niemand; aber von Ana hat man keinen klaren Text noch Zeugniß der Schrift, daraus erwiesen und überwunden werden möchte, wie man von Ruben und Juda klare Historien hat, item, auch von den andern Brüdern, die Joseph verkauft und ihren alten Vater durch Traurigkeit und Herzeleid erwürget haben. Denn also beweint er seinen Sohn, den er verloren hatte, und sprach 1 Mos. 37, 35.: „Ich werde mit Leide hinunter fahren in die Grube, zu meinem Sohne“, das ist, ich werde vor großem Herzeleid sterben; item 1 Mos. 42, 38.: „Ihr werdet meine grauen Haare mit Herzeleid in die Grube bringen.“

31. Derhalben sind die Voreltern und Patriarchen der Juden Vatermörder. Denn also wird allezeit aus dem Allerbesten das Aergerste und Böseste. Aus den Engeln werden Teufel, aus dem Volk Gottes werden solche Leute, die den Sohn Gottes kreuzigen, aus den Propheten werden falsche Propheten, aus unsern Zuhörern und Schülern werden Schwärmer, Rottengeister und Ketzer. Und ist nicht vergeblich geredet, daß Christus zu den Juden sagt Joh. 8, 44.: „Ihr seid von dem Vater, dem Teufel.“ Daher weiß man wohl, daß es ein unsinnig Volk ist und daß sie gar dahin gegeben sind in verkehrten Sinn; wie sie auch zu Davids und der Propheten Zeiten überaus böse Leute gewesen sind. Darum pflege ich ihre Glossen zu widerlegen und zu verdammen, wie ich nur kann und mag.

32. Es sollte uns auch wohl nicht schwer sein, vielerlei Meinungen zu erdichten, wenn es sich geziemen wollte. Wir möchten sagen, daß Ana auch bei seiner Mutter geschlafen hätte, aus welcher Blutschande Ahalibama geboren sei, oder dergleichen viel andere Meinungen. Aber in der heiligen Schrift soll man nichts erdichten, wo man nicht offenbare Zeugnisse des Wortes hat. Wir werden hernach hören, daß Ana ein sehr böser Bube gewesen ist, und ein Urheber einer schändlichen Vermischung, nämlich, der Esel und Pferde. Wo er nun keine Scheu getragen hat vor Gottes und auch der Natur Ordnung und Angesicht, sondern hat dürfen mit einander laufen lassen solche Thiere, die gar ungleicher Art sind, welches wider die Natur ist, und dazu auch wider Gottes Ordnung in der Schöpfung, davon die heilige Schrift sagt in diesem ersten Buche Mose am 1. Cap. V. 25.: „Und Gott machte die Thiere auf Erden, ein jegliches nach seiner Art, und das Vieh nach seiner Art, und allerlei Gewürm auf Erden nach seiner Art“: so hat das auch wohl geschehen können, daß er bei seiner Mutter geschlafen hat.

33. Und ich sage zwar nicht, daß die Heiden von solchen Lastern und Sünden frei gewesen sind, dieweil auch St. Paulus Röm. 1, 26. 27. von ihrer unziemlichen Vermischung und Blutschande sagt. Aber dieweil diese Art zu reden, die Mose gebraucht von Ahalibama, auch von ehrlichen Töchtern gebraucht wird, wie hernach, 1 Mos. 36, 39., von der Mehethabeel, die eine Tochter Matreds und eine Tochter Mesahabs gewesen ist: darum will ich die Worte nicht häßlich oder falschlich auf einen andern und fremden Verstand ziehen, sondern will sie ehrlich deuten. Denn es ist in der heiligen Schrift gemein, daß die Töchter ihrer Väter und Großväter Töchter genannt werden. Mose sagt 2 Mos. 6, 20.: „Amram nahm seine Muhme Jochebed zum Weibe“, welche eine Mutter Moses, und eine Tochter Levis ist, seines Großvaters, das ist, seines Kindes Tochter; dieselbe, da sie eine Base war, hat er zum Weibe genommen. Darum haben wir hier keine Ursache, Schande und Laster zu dichten, dieweil uns die Schrift oder die Historie dazu nicht zwingt; wie zuvor in den Exempeln der Patriarchen des jüdischen Volkes gesagt ist.

\quellseite{263}

34. Darnach dient dieses auch zur Ehre Esaus, daß wir es dafür halten, daß er mit nichten solche Schande gelitten habe, und daß er Ahalibama nicht würde zum Weibe genommen haben, wenn er gewußt hätte, daß sie von der Blutsbande geboren wäre. Darum halte ich dafür, daß die Heiditinnen beide ehrliche Weiber gewesen seien, wiewohl sie hernach angefangen haben, stolz und hoffärtig zu werden und sich ihres Geschlechtes zu überheben. Und vielleicht haben sie auch Götzen mitgebracht in Isaaks Haus, und haben die Andern angereizt, dieselbigen zu ehren, dadurch die frommen heiligen Eltern sehr erzürnt worden, und in der ersten Tafel durch Abgötterei, in der andern aber durch Ungehorsam und Hoffahrt sehr wohl geplagt worden sind. So ist aber doch ihr Ursprung und Abkunft ehrlich gewesen, und Ahalibama ist eine Tochter gewesen Anas als ihres natürlichen Vaters und Zibeons Tochter als ihres Großvaters.

\section[Zweiter Theil]{Zweiter Theil. Von der Herrlichkeit und Herrschaft Esaus und seiner Nachkommen; von ihrer Schwägerschaft; und von ihren Königen und Fürsten}

\subsection{I}

\begin{quote}
\textbf{V. 15–19. Das sind die Fürsten unter den Kindern Esau: Die Kinder Eliphas, des ersten Sohns Esaus, waren diese: Der Fürst Teman, der Fürst Omar, der Fürst Zepho, der Fürst Kenaz, der Fürst Korah, der Fürst Gaetham, der Fürst Amalek. Das sind die Fürsten von Eliphas, im Lande Edom, und sind Kinder von der Ada. Und das sind die Kinder Reguel, Esaus Sohns: Der Fürst Nahath, der Fürst Serah, der Fürst Samma, der Fürst Missa. Das sind die Fürsten von Reguel im Lande der Edomiter, und sind Kinder von Basemath, Esaus Weib. Das sind die Fürsten der Ahalibama, Esaus Weibes: Der Fürst Jehus, der Fürst Jaelam, der Fürst Korah. Das sind die Fürsten von Ahalibama, der Tochter des Ana, Esaus Weib. Das sind Esaus Kinder, und ihre Fürsten. Er ist der Edom.}
\end{quote}

35. Mose hat erst das Geschlecht und Familien Esaus genannt, jetzt beschreibt er nun ihre Würde und Macht oder Herrschaft. Droben im 14. Cap. 3. 6. haben wir gehört, daß die Horiter zu Seir gewohnt haben, dieselbigen haben die Kinder Esaus vertrieben und verjagt. Darum ist das ganze Land Seir der Edomiter Erbe geworden; welches aus 5 Mos. 2, 4. 5. zu sehen ist, da also steht: „Und gebeut dem Volk, und sprich: Ihr werdet durch die Grenze eurer Brüder, der Kinder Esau, ziehen, die da wohnen zu Seir; und sie werden sich vor euch fürchten“ etc.; item: „Das Gebirge Seir hab ich den Kindern Esau zu besitzen gegeben.“

36. Das ist der leibliche Segen gewesen, welchen Gott auch darum so viel lieber gegeben hat, weil Eliphas und die Andern die Beschreibung und die göttliche Lehre behalten haben, welche sie von den Vätern empfangen hatten. Denn Esau hat seinen Kinderkindern schöne herrliche Predigten gehalten; ja, Esau hat des Priesteramtes auch gepflegt anstatt seines Vaters in der Abwesenheit Jakobs. Darum soll man sie nicht alle für Verworfene halten. Denn wiewohl sie die Verheißung vom künftigen Samen verloren haben, so haben sie doch gleichwohl zufällige Barmherzigkeit erlangt.

37. So erzählt nun Mose die Fürsten, welche die Horiter vertrieben haben; denn sie haben Weiber genommen von den Horitern. Darnach aber ist unter die Schwäger und Kindeskinder Uneinigkeit gekommen und die Kinder Esaus haben gesiegt wider die Einwohner im Lande Seir; sie sind nicht durch ordentliche Nachfolge oder Erbschaft zum Regiment gekommen. Und aus diesen zwölf Fürsten haben sie darnach einen König erwählt und haben mit vereinigten Kräften die Horiter aus dem Lande und von ihrem Gut getrieben. Der Herr hat ihnen den Sieg gegeben, wie droben im Segen gesagt worden ist: „Deines Schwerts wirst du dich nähren.“ Ich halte dafür, sie haben das Land Seir im Anfang inne gehabt: da Jakob nach Egypten gezogen ist, da haben sie angefangen zu herrschen, und haben gedacht, Jakob werde des Segens vergessen und ihn fahren lassen. Laß ihn immerhin ziehen, werden sie gesagt haben, immer weg mit ihm, wir sind Fürsten und Herren geworden in diesem Lande: nun wird man sehen, welches Geschlecht gesegnet sei. Also geht allezeit der leibliche Segen vor dem geistlichen her.

38. Es sind aber von Eliphas, dem erstgebornen Sohn Esaus, sechs Fürsten gekommen; er selbst aber ist kein Fürst gewesen. Und Amalek ist auch mit in dieser Zahl. Warum aber Korah mitgezählt werde, kann ich eigentlich nicht wissen. Denn es werden allein die Kindeskinder und nicht die Söhne Esaus unter diesem Titel der Fürsten erzählt. Vielleicht ist es also geschehen, daß ihnen andere Namen sind gegeben worden, oder aber, daß sie von Andern sind als Kinder angenommen worden, welches in der heiligen Schrift sehr gemein ist; wie Joseph Ephraim angenommen hat zum rechten Bruder.

39. Bei den Hebräern aber heißt aluph Fürst, von eleph, das heißt, tausend. Darum führt der Evangelist Matthäus Cap. 2, 6. den Text aus dem Propheten Micha am 5. Cap. V. 1. also an: „Und du Bethlehem im jüdischen Lande bist mit nichten die kleinste unter den Fürsten Juda“; während der Prophet sagt: „unter den Tausenden in Juda“; das sind gleiche Worte. Wenn Einer etliche Dörfer und tausend Mann unter sich in seiner Gewalt gehabt hat, der ist ein Fürst gewesen.

\begin{quote}
\textbf{V. 20–30. Die Kinder aber von Seir, dem Horiter, der im Lande wohnte, sind diese: Lothan, Sobal, Zibeon, Ana, Dison, Ezer und Disan. Das sind die Fürsten der Horiter, Kinder des Seir, im Lande Edom. Aber des Lothans Kinder waren diese: Hori und Hemam; und Lothans Schwester hieß Thimna. Die Kinder von Sobal waren diese: Alwan, Manahath, Ebal, Sephu und Onam. Die Kinder von Zibeon waren: Aja und Ana. Das ist der Ana, der in der Wüste Maulthiere erfand, da er seines Vaters Zibeons Vieh hütete. Die Kinder aber Ana waren: Dison und Ahalibama, das ist die Tochter Ana. Die Kinder Dison waren: Hemdan, Esban, Jethran und Cheran. Die Kinder Ezer waren: Bilhan, Sawan und Akan. Die Kinder Disan waren: Uz und Aran. Dies sind die Fürsten der Horiter: Der Fürst Lothan, der Fürst Sobal, der Fürst Zibeon, der Fürst Ana, der Fürst Dison, der Fürst Ezer, der Fürst Disan. Das sind die Fürsten der Horiter, die regiert haben im Lande Seir.}
\end{quote}

40. Das hebräische Wort seir heißt Feldteufel. Im Propheten Jesaia am 13. Cap. V. 21. haben wir es verdolmetscht Feldgeister, die da pflegen in den Wäldern und auf dem Felde zu erscheinen und ganz rauh und haaricht sind. Die Heiden haben es auf \unsicher{Katzenich} lares und genii genannt, Hausgötter; es sind aber Teufel, wiewohl sie sich verstellen, als wollten sie den Menschen wohlthun. Und dieweil Gomor rauh geboren war, ist er Seir genannt worden. Darum bleibt das Wort Seir entweder von den gegenwärtigen oder aber von den vorigen Einwohnern des Landes. Esau aber hatte ein Weib genommen von den Töchtern Seir oder der Horiter.

41. Derhalben wird der Schwägerschaft auch gedacht, Isaak und Abraham zu Ehren, deren Enkel und Sohn er ist. Denn ob er wohl des Segens beraubt war, so ist doch gleichwohl die natürliche Neigung an den Eltern geblieben, so sie zu ihrem Sohne gehabt haben, daß sie ihn nicht gar verwerfen mochten. Wie Kain nicht dermaßen verstoßen worden von seinen Eltern, darum daß er Abel, seinen Bruder, ermordet hatte, daß sie ihn nicht gleichwohl sollten lieb gehabt und für ihn als ihren Sohn gesorget haben. Denn die natürlichen Neigungen werden nicht verhindert oder aufgehoben, denn allein in den Dingen, die wider Gottes Gebote sind. Darum ehrt Gott auch die Verwandtschaft Esaus von Seiten seines Weibes um Isaaks willen. Denn wozu dient es, daß des Seir und der Horiter gedacht wird, welche unterdrückt und vertrieben sind von den neuen Fürsten, die an ihre Statt gekommen sind? Aber sie werden aufgezählt und beschrieben um weltlicher Ehre und Wohlstandes willen, propter sanctum seculare, wie sie es nennen, oder dem weltlichen Regiment zu Ehren.

42. Die andere Ursache ist, daß man ein Zeugniß haben möchte des Ursprungs oder der Abkunft der Edomiter, auf daß sie nicht für unehelich und Bastarde gehalten würden. Und dieses ist gleichsam der vierte Abschnitt in diesem Capitel, darin die Schwägerschaft Esaus aufgezählt wird.

43. Die Horiter haben ihren Namen daher, daß sie weiß sind. Denn das Wort chur heißt weiß, und das da scheint und groß ist. Die großen Herren haben in den Morgenländern weiße Seide gebraucht, wie heutiges Tages die Türken weiße seidene Hüte gebrauchen. Der Römer Tracht war Purpur, und die Toga. Die Morgenländer haben weiße seidene Kleider gehabt. Also wird im Evangelium von dem reichen Schlemmer gesagt: „Er kleidete sich mit Purpur und köstlicher Leinwand“, Luc. 16, 19. Daher sind die Römer von der Kleidung also genannt worden; wie die Römer torquati (mit Halsketten Versehene) sind genannt worden.

\quellseite{264}

44. Eliphas hat Thimna, die Schwester des Fürsten Lothan, zum Weibe genommen und hat mit derselben Amalek gezeuget. Diese Leute sind nicht gar zu verdammen; denn sie haben sich bessern können und Buße thun aus den Predigten Esau’s. Von Ana ist zuvor gesagt worden, und sind auch der Juden Lügen widerlegt worden, dazu sie Lust haben, daß sie die Edomiter und Heiden damit schänden und schmähen mögen. Es werden aber solche Lügen von Andern auch verworfen.

45. Denn es ist ein Anderes, wo man ein offenbares Zeugniß der Historie hat, wie man der keines hat von Ana, daß er mit seines Sohnes Weibe Blutschande sollte getrieben haben. Und ist auch ein Anderes mit irgend einer heimlichen That, die man nicht sollte an den Tag bringen oder offenbaren; wie man zu sagen pflegt: De occultis non judicat ecclesia: Die Kirche urtheilt nicht von den Dingen, die da heimlich sind; sondern solche Thaten sollen in der Beichte heimlich bleiben und heimlich gehalten werden. Und um der Ursache willen soll man in der Kirche die heimliche Beichte behalten. Denn die betrübten Gewissen, die da beschweret sind von Schrecken der Sünde, offenbaren in solcher Beichte ihre Beschwerungen, und empfangen Trost, den sie in der gemeinen Predigt also nicht ergreifen können. Denselben betrübten Menschen, das wollen wir haben, soll die Beichte eine Zuflucht sein, nämlich, welcher Gewissen der Teufel in seinen Stricken gefangen hält, und die er bezaubert und so gar matt macht, daß sie sich nicht können losmachen und herauswickeln, und fühlen oder sehen Anderes nichts, denn daß sie verloren und verdammt sein müssen. Denn in diesem Leben ist kein größerer Jammer und Elend, denn die Schmerzen und Angst eines solchen Herzens, das verlassen ist und keinen Rath oder Trost weiß.

46. Solche betrübte Herzen sollen nun einen Zutritt haben zu der Beichte, daß sie bei den Kirchendienern Rath suchen und finden mögen. Oder wenn ja die Sache so gar grob ist, daß man sich schämt, dieselbe vor den Kirchendienern zu sagen, sollen sie es sonst einem andern Christen und gottseligen Menschen, wer er auch sei, offenbaren, von dem sie wissen, daß er treu und fromm ist. Vor demselben sollen sie sich beklagen der Dinge halben, die ihnen ihr Gewissen drücken, beschweren und plagen, und sollen bei demselben Rath suchen und also sagen: O rathe, lieber Bruder, meinem erschrockenen und betrübten Gewissen, ich werde mich erhängen, ertränken, oder sonst ein Leid thun und in meinen Sünden dahin sterben etc. Wo da ein solcher betrübter Mensch Gottes Wort wird hören von Vergebung der Sünden, entweder vom Diener des Wortes, oder aber sonst irgend von einem andern Christen, wer er auch sein mag: so wird er aufgerichtet werden und Trost empfangen, dadurch das arme Herz, so der Teufel mit seinen Pfeilen verwundet hat, wiederum geheilt werden.

47. Wir suchen oder fragen darnach nicht, welche oder welcherlei Sünde man erzählen soll, sondern darin steht die vornehmste Kraft und Nutzen der Beichte, daß den armen Herzen, die da seufzen und mit der Verzweiflung kämpfen, guter heilsamer Rath aus Gottes Wort angezeigt werde, daß sie unter schweren Last der Sünden und dem Schrecken des Teufels, das er unendlich größer und größer zu machen pflegt, nicht gar liegen bleiben: vornehmlich aber in sonderlichen und ungewöhnlichen Fällen, wie sich deren bisweilen etliche zutragen pflegen, wenn der Teufel die menschliche Natur treibt und entstellt durch greuliche Sünde. Wie vor Zeiten zu Erfurt sich ein greulicher Fall zugetragen hat, als ich daselbst Beichtvater war, welchen Fall meiner Mitbrüder einer in der heimlichen Beichte gehört hat. Denselben Fall will ich aber jetzt hier erzählen um derjenigen willen, die dermaleinst Pfarrherren oder Lehrer in der Kirche werden möchten.

48. Zu Erfurt hat eine Magd von ehrlicher Freundschaft, und die züchtig und fromm gewesen ist, einer reichen Wittfrau gedient. Derselben Wittfrau Sohn, ein Jüngling, hat sich die Magd gefallen lassen, dieweil sie schön war, und hat sie lieb gewonnen, hat auch begehrt bei ihr zu schlafen. Die Magd aber hat sich der bösen That geweigert (wie einer frommen züchtigen Magd gebührt) und hat den unsinnigen Jüngling etlichemal abgewiesen und zurückgetrieben. Endlich, da er täglich immer mehr und mehr bei ihr angehalten und bei seinem bösen unziemlichen Vornehmen beharrt, hat die Magd, durch die Schändlichkeit der Sache bewogen, damit sie ihr gut Gerücht, Zucht und Frömmigkeit retten und sich solcher großen Schande erwehren möchte, seiner Mutter den ganzen Handel erzählt und sie ermahnt, daß sie ihrem Sohne, der ihr heimlich nachstelle, wie er sie möchte um ihre Ehre bringen, solches untersagen und verbieten wolle. Da nun die Mutter die Sache erwogen hat, macht sie es mit der Magd also, und heißt sie in die Liebe ihres Sohnes willigen, und daß sie ihm eine gewisse Stunde bestimmen soll, daß sie zusammen kommen und die Sache ihres Bedünkens sicher ausrichten möchten. Denn sie, die Mutter, wollte eben auf solche bestimmte Zeit sich in der Magd Bette legen, und durch solche Gelegenheit ihren Sohn zurückhalten, von solchem bösen Vornehmen abzustehen. Den Rath läßt sich die Magd gefallen und macht mit der Mutter einen solchen Bund. Da kommt nun der Sohn in der Nacht auf die bestimmte Stunde, und begehrt bei der Magd zu schlafen, wie sie mit einander übereingekommen waren. Die Mutter aber, die zuvor bei sich beschlossen hatte, mit solcher List die Brunst des Sohnes zu verhindern und zu verhüten, ist durch ihre eigene böse Lust und des Teufels Eingeben überwunden worden, daß sie den Sohn hat bei sich schlafen lassen. Aus solchem unziemlichen Beischlafen und Blutschande wird ein Mägdlein geboren. Dasselbe, das es heimlich weggethan und von Andern ist auferzogen worden, hat endlich die Mutter wiederum zu sich genommen, dazu sie durch mütterliche Neigung und Liebe ist bewogen worden.

49. Derselbe Sohn aber, da er nun alle diese Dinge nichts weiß, gewinnt er dieses Mägdlein auch lieb, und begehrt, man wolle sie ihm zum Weibe geben. Da ist die Mutter in großer Angst gewesen, und hat solches dem Sohne widerrathen und seiner unziemlichen Lust und Begierde widerstanden, wie sie nur gekonnt und vermocht hatte. Er aber fährt zu und verlobt sich mit der Jungfrau, und nimmt sie auch also zum Weibe wider der Mutter Willen; denn zu der Zeit waren die heimlichen Verlöbnisse gemein und sind nicht zerrissen worden. Da aber die Ehe vollzogen und die Hochzeit gehalten war, weiß die Mutter nicht, was sie thun soll, und ist in Verzweiflung gerathen, daß sie sich selbst hat umbringen wollen. Denn sie ist der Blutschande halben bekümmert, so von denselbigen Eheleuten in jeder Nacht geschehen, und konnte gleichwohl die Ehe nicht scheiden. Da sie aber dieses allein gemuth und den großen Schmerz und Angst des Herzens nicht länger allein hat ertragen können, geht sie zu ihrem Beichtvater, erzählt ihm den Handel und begehrt Rath und Trost von ihm. Der weiß nicht, wie er daran ist, und wird gar erschrocken über solchen greulichen Handel, daß er nicht weiß, was er darin rathen soll. Derhalben wird derselbe Fall vor die Juristen gebracht; die zweifeln selbst auch, was darin möchte recht sein, als in einem ungewöhnlichen Fall, und halten es dafür, daß man es göttlicher Güte und Barmherzigkeit befehlen müsse. Daran ist es aber nicht genug gewesen und das arme Gewissen war damit seines Zweifels und Zagens noch nicht los. Endlich ist die Sache vor die Theologen gekommen, dieselben haben darin ein sehr gut Urtheil gesprochen. Erstlich aber haben sie gefragt: Ob auch der Sohn wüßte, daß sein Weib seiner Mutter Tochter wäre, von seinem Geblüt empfangen? Darauf hat die Mutter gesagt, daß er davon nichts wüßte. Es weiß es allein Gott und ich, sprach sie, und die Jungfrau, meines Sohnes Weib, weiß auch nicht, daß sie ihres Mannes Schwester ist. Darnach haben sie weiter gefragt: Ob es auch eine gute Ehe wäre und wie sie sich unter einander vertragen könnten? Darauf sagt die Mutter, daß sie mit einander in gutem Frieden und Einigkeit lebten. Darum haben sie also beschlossen: Man muß dem betrübten Gewissen der Mutter mit der Absolution rathen; dem Sohn aber die Ehe auch zulassen, dieweil dieselbe auch vorlängst durch Verlöbniß und gewöhnlichen Kirchgang oder hochzeitlichen Ehrentag und das Beischlafen bestätiget war. Denn die Mutter wäre eine einzelne Person, welche diese Sache allein nicht beweisen oder bezeugen könnte; zudem so würde daraus groß Unglück folgen, wenn die Eheleute sollten von einander geschieden werden.

\quellseite{265}

50. Diese Sentenz und Urtheil der Theologen, unter denen die Sache geheim verhandelt ist, ist sehr zu loben. Denn das Gesetz muß herrschen und Statt haben, wo man öffentlich vor Gericht handelt, und nicht im Gewissen. Und die Pforte der Absolution, daß ich es also nenne, soll vor den armen Sündern immer mehr zugeschlossen werden. Er falle so tief, als er immer wolle. Und ich habe die Historie um der jungen Theologen willen erzählt, die entweder jetzt das Predigtamt verwalten, oder aber noch dermaleinst zu dem Amt kommen möchten, daß sie vorsichtiglich mit den erschrockenen Gewissen handeln; daß sie denen, so ihre Sünde bekennen, keine Stricke anlegen, oder die bekümmerten und betrübten Herzen nicht mehr oder weiter bekümmern. In heimlichen Sachen ist es sehr schwer zu urtheilen; darum gehören gelehrte und verständige Pfarrherren dazu, welche die armen Gewissen nicht beschweren oder verwirren, sondern sie erretten, aufrichten und wiederum heilen, welche Gewissen der Teufel bethört und in seinen Stricken gefangen hat.

51. Man sagt auch noch von einem andern Falle, der sich mit einer Ehebrecherin soll zugetragen haben, die von einem Ehebrecher war schwanger geworden und einen Bastard geboren hat, und darnach denselben ihrem Ehemann und den andern ihren Kindern dermaßen beigebracht und angegeben hat, als wäre er auch ehelich geboren. Da sie aber in ihrem Herzen bekümmert worden um des Ehebruchs willen, den sie begangen hatte, und daß die ehelichen Kinder des Erbtheils unbilliger Weise sollten beraubt werden, welches Theil dem Bastardo zufallen würde, ist gefragt worden: Wie ihr doch möchte zu rathen sein? Ob sie ihrem Mann und Kindern die Sünde bekennen sollte, und sie bitten, daß sie den Bastard auch wollten lassen Miterben sein zu den Gütern? denn dasselbe haben Etliche gerathen. Dieweil es aber schien, als ob solches nicht ohne Gefahr geschehen würde, so haben die Andern anders gerathen. Endlich aber hat der Heilige Geist dem Weibe einen sehr guten Rath eingegeben. Denn da sie schwer krank gewesen und zu Bette gelegen hat, und keine Hoffnung des Lebens mehr da gewesen ist, hat sie ihre lieben Kinder vor sich zusammenkommen lassen, und zu ihnen also gesagt: Meine lieben Kinder, ich bitte, ihr wollet mir verzeihen. Denn ich bekenne, daß ich aus menschlicher Schwachheit in die große schwere Sünde des Ehebruchs gefallen bin, davon ein unehelicher Sohn geboren, welcher unter euch ist; wer aber der sei, will ich nicht sagen. Ich bitte aber, ihr wollet ihn für ehelich halten und ihn von dem Erbe nicht ausschließen. Da sind sie dazu alle von Herzen willig und geneigt gewesen, und ist dazu einem jeden Kinde sehr lieb gewesen, daß der Bastard nicht mit Namen war genannt worden. Denn niemand unter ihnen hat diese Schande tragen wollen und haben sich darnach die Brüder wohl vertragen können. Die Mutter aber, da sie den Trost von Vergebung der Sünde gehört, hat sie sich in den letzten Zügen und Todesnöthen daran gehalten und ist durch rechten Glauben und Erkenntniß der Gnade und Barmherzigkeit Gottes in Christo selig geworden.

52. Darum soll man, was heimlich ist, nicht so leichtlich nachsagen, wie die Juden mit dem Ana thun, und dasselbe noch fälschlich und mit Unwahrheit. Darnach muß man in solchen Fällen den betrübten Gewissen helfen, daß sie erlöset und dem Teufel aus dem Rachen gerissen werden, daß nicht das theure Blut Christi an ihnen verloren sei.

53. Es wird aber dem Ana zugeschrieben, daß er die Maulthierpferde soll erfunden haben. Denn da er in der Wüste die Heerde seines Vaters geweidet, hat er die Esel mit den Mutterpferden laufen lassen, davon sind die Maulthierpferde gekommen. Plinius im 8. Buch und 44. Capitel beschreibt weitläufig, und mit vielen Worten, wie mancherlei Weise solche Vermischungen bei den Heiden geschehen seien, welches wider die Natur und Schöpfung der Dinge ist. Darum ist es eine schändliche und ungebührliche Erfindung. Und er sagt noch weiter, daß die dritte Art, so von verschiedener Art der Thiere hergekommen, unfruchtbar sei. Denn Gott hat nicht haben wollen, daß solche Art sollte weiter ausgebreitet werden.

54. Wiewohl ich an dem hebräischen Wort jemim zweifle und von seiner Bedeutung nichts Gewisses sagen kann; denn die Juden sind wegen des Hasses, den sie gegen die Heiden tragen, sehr kühn, Laster zu erdichten, sie damit zu beflecken und zu schmähen. Und die üblichen Punkte sind erst später in der hebräischen Sprache erfunden worden, darauf wir uns nicht allenthalben sicher verlassen können. So mag es nun wohl sein, daß er die Seen und Teiche erfunden habe, deren im Evangelium Lucä oftmals gedacht wird. Und Hieronymus ist auch der Meinung und der Text selbst ist nicht dawider; wiewohl dazu gesetzt wird, daß er die Esel gehütet habe. Ich will aber doch nichts Gewisses davon sagen, wiewohl ich in der deutschen Uebersetzung den Rabbinen gefolgt habe. Es scheint aber, als ob es sich sehr wohl reimen wollte, darum daß in der Wüste wenig Wasser ist, daß entweder diesem oder aber einem andern Stamme zugeschrieben möchte werden, daß sie die Seen erfunden hätten. Und die hebräischen Buchstaben und Puncte, wenn sie recht geschrieben werden, reimen sich auch fein dazu; denn jamim heißt Seen: und ist eine sehr große Gabe und Nutzen mit den Seen, vornehmlich an Orten, da es sehr trocken ist, und in der Wüste. Wo er aber die Maulthierpferde wollte verstanden haben, so sollte auch dabei der Mutterpferde gedacht worden sein, welches Mose nicht gethan hat.

\subsection{II}

\begin{quote}
\textbf{V. 31—39. Die Könige aber, die im Lande Edom regieret haben, ehe denn die Kinder Israel Könige hatten, sind diese: Bela war König in Edom, ein Sohn Beors; und seine Stadt hieß Dinhaba. Und da Bela starb, ward König an seine Statt Jobab, ein Sohn Serah von Bazra. Da Jobab starb, ward an seine Statt König Husam, aus der Themaniter Lande. Da Husam starb, ward König an seine Statt Hadad, ein Sohn Bedad, der die Midianiter schlug auf der Moabiter Felde; und seine Stadt hieß Awit. Da Hadad starb, regierete Samla von Masrek. Da Samla starb, ward Saul König, von Rehobot am Wasser. Da Saul starb, ward an seine Statt König Baal Hanan, der Sohn Achbor. Da Baal Hanan, Achbors Sohn, starb, ward an seine Statt König Hadar, und seine Stadt hieß Pagu; und sein Weib hieß Mehetabeel, eine Tochter Matred, die Mechahs Tochter war.}
\end{quote}

55. Bisher hat Mose beschrieben die Fürsten vom Geschlechte Esaus oder Edom; darnach die Fürsten Seir und der Horiter, die Esau und seinen Kindern durch Schwägerschaft verwandt gewesen sind. Und diese tapferen Helden haben die Kinder Edom vertilgt und ihr Land eingenommen, welches darnach Edomä oder Idumäa ist genannt worden. Heutiges Tages hat es den Namen verloren, und von der Zeit des Hieronymus her ist es Gebalena genannt worden. Denn wie die Königreiche verändert werden, also ändern sich auch die Namen. Jetzt nennt er nun auch die Könige Edoms. Denn wir haben zuvor gesagt, daß Gott in ihnen den Patriarchen Isaak und Rebekka geehret habe; wie er droben Abraham und Sarah auch also ehrt, daß er um ihretwillen Ismael segnet und ihm viel Erben verheißt. Denn also spricht er zu Abraham im 17. Cap. v. 20.: Dazu um Ismael habe ich dich auch erhöret. Siehe, ich habe ihn gesegnet, und will ihn fruchtbar machen und mehren fast sehr. Zwölf Fürsten wird er zeugen, und will ihn zum großen Volk machen. Diese Gaben hat Gott den Ismaeliten gegeben um Abrahams willen. Und ohne Zweifel haben auch die Ismaeliten die Beschneidung und übrigen Ceremonien gehalten, vielleicht auch wohl die Lehre und den Glauben von dem verheißenen Samen, und sind also zu der Kirche und Gemeinde Abrahams gekommen; wie dagegen Viele vom Hause und Geschlechte Abrahams abgefallen und Heiden geworden sind.

56. Auf solche Weise hat Gott den Berg Seir oder das Land der Horiter den Edomitern gegeben, von welchen sie vertilgt worden sind, und es scheint, daß sie im Glauben und bei der Religion der Väter, Abrahams und Isaaks, beständig geblieben sind. Da sich aber die Horiter nicht haben wollen bessern, und die abgöttischen falschen Gottesdienste ändern und abschaffen, sind sie vertilgt worden. So wird nun allhier von Mose beschrieben der leibliche Segen Esaus, davon er auch zum geistlichen Segen gekommen ist. Wiewohl man den Herrn Christum vom Samen Esaus nicht hat erwarten dürfen, so ist es aber doch ihm oder seinen Kindern nicht versagt worden, daß sie nicht auch der gemeinen Wohlthaten der Verheißung mit dem Volke Israel haben gebrauchen mögen. Wie wir in den Büchern der Könige sehen, daß die Edomiter bisweilen Freunde und Unterthanen gewesen sind der Könige Juda, als, zu den Zeiten Davids, der Idumäa unter sich gebracht und bezwungen hat, bisweilen aber auch wiederum abgefallen sind, als, unter dem Könige Joram.

\quellseite{266}

57. Darum wird dieses alles also beschrieben um der Väter, Abrahams, Isaaks und Jakobs willen, und denselbigen zu Ehren, und daß niemand meinen soll, daß den Edomitern der Weg verschlossen sei zur Gnade und zum Segen des verheißenen Samens.

58. Es wird aber gefragt: Ob diese Fürsten und Könige vor oder nach Mose gewesen sind? Wo sie nach Mose gewesen sind, so hat er dieses ja nicht schreiben können, sondern diesen Zusatz hat ein Anderer gemacht, wie das letzte Stück ist im fünften Buch Mose. Denn er hat ja von sich selber nicht gesagt 5 Mos. 34, 10.: „Und es stund hinfort kein Prophet in Israel auf, wie Mose, den der HErr erkannt hätte von Angesicht.“ Item, andere Ding mehr, so daselbst vom Grabe Moses erzählt werden. Es wäre denn, daß du sagen wolltest, daß er solches durch einen prophetischen Geist zuvor gesehen und geweissagt hätte. Dieweil aber der Heilige Geist diese Könige und Fürsten beschreibt und rühmt, sollen wir ihr Gedächtnis auch behalten.

59. Die Fürsten aber, so droben, V. 15. ff., aufgezählt worden, sind vor den Königen gewesen, und hat nicht Einer nach dem Andern regiert, sondern sie haben zugleich regiert, also daß ein Jeder über seine Städte, Gesinde und Unterthanen regiert hat an seinem bestimmten Ort, und daß das Land schier in zwölf oder mehr Fürstenthümer ist getheilt gewesen. Unter denselbigen aber sind viel herrliche Männer gewesen, als Isaaks und Rebekkas Enkel, wiewohl ihre Mütter von dem gottlosen Geschlechte der Horiter hergekommen waren.

60. Die Könige aber, so an diesem Orte aufgezählt werden, sind an der Fürsten Statt gekommen. Denn sie haben wohl gesehen, daß vonnöthen wäre, daß man nur Einen obersten Regenten hätte, auf daß das ganze Land unter Einem Haupte in Einer Gemeinde zusammengehalten würde, auf daß sie sich so viel besser schützen könnten wider ihre Nachbarn, die Philister, Araber, Egypter, die ihre Feinde waren. Dieses ist nun die Monarchie und Regiment der Könige Edom.

61. Dieses Verzeichnis aber, da die Könige und Fürsten Esaus also nach einander aufgezählt worden, erinnert uns, daß die geistlichen Verheißungen sehr langsam erfüllt werden. Denn also sagt St. Paulus 1 Cor. 15, 46.: „Der geistliche Leib ist nicht der erste, sondern der natürliche, darnach der geistliche.“ GOtt hat das Geschlecht Esaus mit vielen großen leiblichen Gaben gezieret. Er hat ihm Kinder und Kindeskinder gegeben und dieselben alle Könige und Fürsten. Jacob hat den geistlichen Segen und Verheißung, nämlich, da zu ihm gesagt wird: „Sei ein Herr über deine Brüder“ etc., davon wir droben, Cap. 27., gehört haben: er ist aber ein armer elender Schafhirte, muß mit seinen Kindern umherziehen in der Irre, Fremde und Elend. Wie ist denn nun GOtt wahrhaftig, dieweil er die Verheißung, so dem Jacob geschehen war, nicht erfüllet, sondern segnet und mehrt Esau und macht ihn sehr groß, dem doch der Segen genommen ist? Ja, da haben die leiblichen Nachkommen und Erben Esaus gerühmt, daß, da sie Fürsten im Lande wären, so wäre auch die Erstgeburt bei ihnen, Jacob aber sei gar verworfen. Davon nimmt das Fleisch, daß es sich mit nichten reime mit göttlicher Wahrheit, und wird sehr geärgert durch die Hoffart derjenigen, so der leiblichen Gaben halben groß und reich sind.

62. Man muß aber darauf antworten aus dem 2. Psalm V. 4.: „Der im Himmel wohnet, lachet ihrer, und der HErr spottet ihrer.“ Und wie der heidnische Poet gesagt hat: Tolluntur in altum, ut lapsu graviore ruant, das ist: Sie werden hoch erhoben, auf daß sie so viel härter fallen mögen. Welches sich sehr fein reimt mit dem Spruche im 73. Psalm V. 18.: „Aber du setzest sie aufs Schlüpfrige, und stürzest sie zu Boden“ etc. Ja, das ist der Gottlosen Glück: sie erheben sich und werden stolz, wenn es ihnen wohl gehet, wegen der Herrlichkeit des Fleisches; wie sich die Edomitern gerühmt haben und gesagt: Wir haben das Königreich, wir haben das Land Seir und der Horiter inne; Jacob ist ein Bettler und Landstreicher, der immer geplagt wird mit allerlei Jammer und Unglück, mit Frost und Hitze; er ist allenthalben ein armer unglückseliger Mann, er hat ungehorsame Kinder und die böse Buben sind; Ruben ist auf seiner Mutter Lager gestiegen, Dina ist beschlafen worden von Sichem; derhalben hat er ja mit nichten den Segen.

63. Das läßt aber GOtt darum also geschehen, auf daß der Glaube geübt und geprüft werde; und daß wir lernen allein am Worte hangen und, was sichtbar und greiflich ist, nicht ansehen, sondern uns daran halten, daß man nicht sieht, und daß wir wider alles Gefühl und Urtheil der Vernunft den Trost behalten, davon St. Paulus sagt 2 Cor. 4, 9.: „Wir leiden Verfolgung; aber wir werden nicht verlassen. Wir werden untergedrückt; aber wir kommen nicht um“ etc. Denn was man äußerlich mit Augen sieht, das fehlt und verschwindet; was aber verheißen ist und das man nicht sieht, das ist gewiß und beständig. Man muß dessen aber im Glauben warten, und den Verzug und Aufschub, ehe die Verheißung erfüllet wird, mit Geduld ertragen. Denn zu seiner Zeit wird man es wohl sehen.

64. Was aber nun die Könige und Fürsten der Edomiter für Ehre und Herrlichkeit gehabt haben, so ist doch gleichwohl die Herrlichkeit der Verheißung des Samens viel größer und trefflicher, desgleichen die Herrlichkeit der Propheten, der wunderbaren Erlösung aus Egypten und Errettung auch von andern Feinden, item, der herrlichen schönen Siege, so das Volk Israel gehabt hat, also daß auch der Segen der Edomiter mit diesen Gaben nicht einmal zu vergleichen ist. Jedoch ist Jacob noch wie ein fauler und verdorrter Klotz und in die Erde vergraben; zu seiner Zeit aber wird er grünen. Was bald blüht und lieblich anzusehen ist, das ist nicht beständig und währet nicht lange. Und hier hat der Vers des Poeten auch Statt, da er sagt: O formose puer, nimium ne crede colori: Hüte dich vor dem, was hoch hergeht. Aber diejenigen siehe an und verwundere dich derer, so die Verheißung haben. Denn die göttliche Majestät pflegt die Hoffärtigen also zu erheben, daß sie sich ansehen lassen, als die in großem Flor stehen und immer größer werden. Aber das ist dahin gesetzt, daß zuvor gesagt ist aus dem 73. Psalm V. 18.: „Du setzest sie aufs Schlüpfrige, und stürzest sie zu Boden.“ Also täuscht GOtt den Teufel und die Welt und spottet ihrer.

65. Bei welchen nun die Verheißung und der Glaube ist, die müssen Frost und Hunger leiden, sind verachtet und verschmäht. Wie Einer soll gesagt haben: Es frage nichts Sonderliches nach der Lehre der Religion, verstehe sie auch nicht; aber das müsse er sehr wohl, daß die GOtt und den Kaiser dienen, die elendesten Leute sind. Die Gottlosen aber rühmen sich der Herrschaft und des Regiments. Und das thun sie zwar billig; denn GOtt füllt ihnen den Bauch mit Gütern; er gibt ihnen Königreiche und Güter dieser Welt, wie im 17. Psalm V. 14. steht: „Die da Kinder die Fülle haben, und lassen ihr Uebriges ihren Jungen.“ Er gibt wohl den Heiligen oder Gläubigen auch leibliche Güter, aber langsam und durch viel Trübsal, auf daß ihr Glaube geübt werde, und daß sie die Gaben Gottes verstehen lernen und derselben recht und wohl gebrauchen.

66. Darum läßt uns bei dem Haufen bleiben, der das Wort hat, wiewohl er verachtet und verschmäht ist. Die Gottlosen, die Hoffärtigen und Geizhälse wachsen fein daher und werden groß in dieser Welt; es wächst daher, daß jedermann denkt: Die werden es thun, also daß die Gottlosigkeit anfangen zu grünen und ungeduldig zu werden, wie der 73. Psalm V. 13. sagt: „Soll das umsonst sein, daß mein Herz unschuldig lebt, und ich meine Hände in Unschuld wasche?“ Aber er antwortet darauf und sagt V. 18. 19.: „Ich sehe, daß sie auf das Schlüpfrige gesetzt sind; sie werden plötzlich zunichte; sie gehen unter und nehmen ein Ende mit Schrecken.“ Darum dient dieses uns zur Lehre, auf daß wir uns nicht ärgern an den gemeinen Ackerleuten, so für und für in der Welt währen.

67. Die Könige aber, halte ich, sind zu Moses Zeiten gewesen. Denn die Fürsten sind dem Esau geboren worden vor dem Tode Isaaks, und haben im Lande Seir regiert, da Jacob noch in Mesopotamien gewohnt hat. Denn da er wieder aus Mesopotamien gekommen ist, sind sie ihm mit ihrem Vater gerüstet entgegen gezogen, und hat ihr Regiment bei zwanzig oder dreißig Jahren gewähret bis auf den Tod Isaaks. Da hat es sich angefangen, daß das Fleisch und der leibliche Segen im Fleische gewesen, die Könige aber haben länger denn zweihundert Jahre regiert. Danach sind andere Fürsten gekommen. Darum mag man sehen, wie die Herrlichkeit des Fleisches so groß geworden ist bei denen, so die Verheißung des Samens nicht gehabt haben.

\quellseite{267}

68. Inzwischen ist Jacob elend im fremden Lande, bei dem noch die Blume des Segens war; aber man sieht an ihm nur eitel Dürre und Unfruchtbarkeit. Und da solches die Edomiter gesehen, haben sie ihm noch dazu getrotzt und sein gleichsam gespottet und gesagt: Ist das nun der Segen, welchen Jacob mit List unserm Vater Esau genommen hat? Ei, wie ist Gottes Gericht so gerecht, der solche Gewalt nicht hat ungerächt gelassen! Siehe, wir haben Fürsten und Könige, die Israeliten sind arme Knechte und gefangene Leute, sind allerlei Unglück, Jammer und Noth unterworfen. Und dasselbe haben die Priester und Fürsten auch also gerühmt; wie heutiges Tages die türkischen und päpstlichen Pfaffen von dem großen Glück auf ihrem Theile auch zu rühmen wissen. Es sind der Gottseligen wenig gewesen, die sich haben wider solch groß Aergerniß zu rüsten gewußt, welche auch die Verheißung gesehen, und also geschlossen haben: Gott könne ja seiner Verheißung nicht vergessen, und es werde dermaleinst aus dem dürren Roße ein schöner Zweig kommen, der da endlich grünen und Frucht bringen werde.

69. Ich halte aber dafür, daß diese Könige weltlicher Weise fromme Männer und Etliche unter ihnen auch heilig gewesen sind, sintemal sich der Heilige Geist dazu herabläßt, sie aufzuzählen und zu rühmen. Ihre Historien und Geschichte, desgleichen auch die Beschreibung ihres Regimentes und ihrer Städte sind uns unbekannt; wiewohl daran kein Zweifel ist, daß bei den Juden noch etliche Andenken an sie gewesen sind. Die Kinder Esau werden die herrlichen Thaten und wie es um das ganze Regiment gestanden hat, beschrieben haben.

70. Der erste König ist Behor. Da der gestorben ist, hat Jobab von Bozra an seiner Statt regiert. Diese Nachfolge, daß also Einer nach dem Andern regiert hat, ist droben mit den Fürsten nicht gehalten worden. Bazra ist eine berühmte Stadt in Idumäa, und wird derselben gedacht im Jesaia Cap. 63, 1. und Jeremia; wiewohl Jesaia desselben Worts gebraucht, durch die Figur catachresis genannt, und dadurch Jerusalem verstehet. Es ist auch noch ein ander Bazra, im steinichten Arabien gelegen. Augustinus und Ambrosius halten dafür, daß dieser Jobab vom Geschlechte Esau sei, und meinen, es sei Hiob: und solcher Meinung hat man viel Anzeichen, davon man es abnehmen kann. Denn es reimet sich daher die Buchstaben und Namen der Freunde, von welchen Eliphas, der Themaniter, wie es sich ansehen läßt, hergekommen ist von Theman im Lande Edom; item, Baldad und Sophar. Wo nun das wahr ist, so reimt es sich fein mit der Historie. Denn Hiob ist ein sehr weiser und verständiger Mann gewesen im Gebete und Worte des Herrn, und hat die Lehre und Beschreibung seines Vaters Esau angenommen.

71. Lyra disputirt darüber, weil unter den Buchstaben, damit sich die Namen anfangen, ein Unterschied ist. Denn Jobab schreibt man mit dem Buchstaben j, Hiob aber mit dem h. Darauf antworte ich aber also, daß in der heiligen Schrift nicht ungewöhnlich sei, daß der Buchstabe a gesetzt werde und daß man das j auslasse, und umgekehrt auch also, beide im Anfange und am Ende, und auch mitten im Worte. Darum ist es ein schwacher Beweis. Bei unserm Adel ist dieser Name auch geblieben, Job, das ist, Hiob, Jobst etc.

72. Denn sei aber, wie ihm wolle, die Hebräer sagen, daß Hiob hergekommen sei von Nahor, Abrahams Bruder, vom Geschlechte Bethuel und Laban in Syrien. Dieselbe Herkunft, meinen sie, soll auch Bileam haben, und ist sehr glaublich und reimt sich auch wohl mit der Schrift. Darum, wenn man Vermuthungen anstellen darf, scheint es, daß die Stadt Haran eine Kirche und treffliche Schule gewesen sei in Mesopotamien, da Hiob und seines Gleichen sind zusammen gekommen, item, Bileam. Denn das Wort ist in demselbigen Geschlechte ausgebreitet gewesen.

73. Derhalben wie Abraham und Isaak ihre Kirche bestellt und regiert haben, also hat Nahor auch gethan, und die Lehre ist von diesen Patriarchen und Regenten der Kirche auf ihre Erben gekommen. Aus diesen Kirchen und Schulen sind die Propheten, Lehrer und Patriarchen gekommen, unter denen Hiob auch gewesen ist. Ich kann aber nicht gewiß schließen, ob er von Esau oder von Nahor hergekommen sei: das weiß man gewiß, daß er zu demjenigen Geschlechte Sems gehört, welches gleichsam die Heidnische Christin ist. Denn bei Sem ist die Lehre geblieben, und von ihm ist sie nach Idumäa, Mesopotamien, Arabien und Syrien aufgenommen, darnach durch Joseph auch nach Egypten.

74. Also hat Gott vom Anfange die Kirche und Gemeinde gesammelt und erhalten. Wie im Papstthum die Taufe, die Schlüssel, das Abendmahl und andere reine Artikel des Glaubens geblieben sind, da unterdeß der Papst andere Artikel gemacht hat, als, von den Reliquien und Anrufung der Heiligen, vom Fegefeuer; damit er die gesunde Lehre verfälscht und vermischt hat. Also war Jerusalem voller Altäre und Götzen; und Gott erhält doch gleichwohl allezeit in der Welt unter den Heuchlern und Götzendienern die Kirche.

75. Darum wollen wir die Disputation von der Zukunft Hiobs unentschieden lassen. Die Freunde, deren in seiner Historie gedacht wird, sind ohne Zweifel weise und gelehrte Männer gewesen. Bileam ist auch ein guter Prophet gewesen, aber bis auf das Gold: da das kommt, läßt er sich damit bestechen und verführen. Um des Themaniters willen will aber ich nicht mehr der Meinung, daß Hiob von Esau hergekommen. Denn Arabien liegt näher bei Juda und ist daselbst besser bekannt; wiewohl Mesopotamien auch nahe dabei gelegen ist. Die Araber aber strecken sich aus fast über das ganze Morgenland. Und Hiob klagt Cap. 1. V. 17., daß er seinen Schaden von Saba, das ist, von den Chaldäern empfangen habe. Wir können aber davon nichts Gewisses sagen. Denn es ist ein sehr verworren Ding und eine große Unordnung mit den Eigennamen in der hebräischen Sprache um der Veränderung willen der Buchstaben, auch derer, die literae substantiales genannt werden, der Consonanten, die eigentlich das Wort ausmachen, wie Hadad und Hadar solche Eigennamen sind. Item Mose, 2 Mos. 2, 22., nennt seinen Sohn Gerson, als wenn man auf Lateinisch sagen wollte advena ibi, Hier Gast; und wird doch die ganze Schrift durch und durch Gerson genannt.

76. Hadad hat die Midianiter geschlagen: der Sieg ist ihm gegeben worden von Gott, Isaak zu Ehren, wiewohl Midian auch Abrahams Sohn gewesen ist. Ich halte aber dafür, daß jener fromm, und dieser stolz und hoffärtig gewesen sei, und daß er neben der Lehre Abrahams auch abgöttische Gottesdienste mitgebracht habe. Darum ist er auf dem Felde Moab überwunden und geschlagen worden, da sie mit einander zusammengetroffen sind, indem sie durch einander Land aus ihren Landen gezogen waren, daß sie mit einander kämpfen wollten. Dieses sind Historien, so um die Zeit geschehen sind, da sie hinab nach Egypten gezogen sind.

77. Saul von Rehoboth, das ist, vom breiten Wasser, am Wasser, nämlich, vom Wasser, so nahe dabei gelegen hat, Nahar oder Phrath; wie wir sagen, an der Elbe oder am Rhein. Es ist aber eine Vereinigung und Bundesgenossenschaft gewesen unter den Blutsfreunden, den Idumäern; und sie haben sich auch befreundet mit den Einwohnern in Mesopotamien.

78. Wenn man Baal Hanan umkehrt, so heißt es Hannibal, auf Lateinisch, vir gratiosus, ein holdseliger Mann, Johannes. Unter den acht Königen aber sind nur drei, die Städte haben, nämlich, Bela, Habab und Hadar. Dieselben scheinen gut Haushalter gewesen zu sein, die ihr Land und Städte in gutem Frieden immer größer gemacht und geziert haben, haben Städte gebaut und gute Ordnung mit Gesetzen, Gericht und Gerechtigkeit aufgerichtet.

79. Mehetabeel ist ein vornehm und trefflich Weib gewesen, dieweil Mose dieselbe vor den Andern allein rühmet. Das Wort heißt, Eine, die vor Gott Gutes gethan hat oder von Gottes wegen wohlthätig gewesen ist, als wenn man sagen wollte, Theodosia. Es wird auch ihr Geschlecht genannt, daß sie aus einem berühmten Geschlechte und von ehrlichen Eltern geboren sei. Und ist gleichsam eine solche Beschreibung, wie die der Thalibama gewesen ist, da wir gesagt haben, daß man die giftigen Augen der Juden verwerfen soll. Denn sie kann wohl genannt werden eine Tochter ihres Vaters, und auch eine Tochter ihres Großvaters, welcher an des Vaters Statt getreten ist und sie erzogen hat, da der Vater gestorben gewesen; welches jetzt noch also in fürstlichen Rechten verordnet ist und gehalten wird.

80. Ich halte ferner dafür, daß dieser Hadar der letzte König der Edomiter gewesen sei, und daß Jephthah im Buche der Richter im 11. Cap. V. 13. von demselben rede, daß, da Israel aus Egypten gezogen, der König der Edomiter gewesen sei, daß sie möchten durch sein Land ziehen, haben aber von ihm solches nicht erlangen können. Darum halte ich dafür, daß Gott solche Hoffahrt an ihm gestraft habe und daß die Monarchie und das königliche Regiment Edoms wiederum in etliche Fürstenthümer sei getheilt worden. Zwei Jahre, nachdem sie aus Egypten gezogen, ist das Volk Israel auf dem Einzug nach dem gelobten Lande begriffen gewesen; da sie aber wider den Edomern gesündigt, haben sie wiederum müssen zurück ziehen und um das Land Edom her ziehen. Diese Reise haben sie vollbracht in acht und dreißig Jahren, so lange sind sie um das Gebirge Edom hin- und hergezogen. Also ist, meines Erachtens, dieser Hadar der letzte König gewesen. Da der gestorben ist, sind die Fürsten ihm nachgefolgt, welche zuletzt aufgezählt werden.

\quellseite{268}

\begin{quote}
\textbf{V. 40—43. Also heißen die Fürsten von Edom, in ihren Geschlechtern, Vettern und Namen: Der Fürst Thimna, der Fürst Alva, der Fürst Jetheth, der Fürst Oholibama, der Fürst Ela, der Fürst Pinon, der Fürst Kenas, der Fürst Teman, der Fürst Mibzar, der Fürst Magdiel, der Fürst Iram. Das sind die Fürsten in Edom, wie sie gewohnt haben in ihrem Erblande. Und Esau ist der Vater der Edomiter.}
\end{quote}

81. Nach den Königen, und da die Monarchie und das königliche Regiment um der Sünde willen des Königs Hadar verstört war, der den Kindern Israel nicht gestatten wollte, durch Idumäa zu ziehen, folgen nun diese Fürsten, jeder an seinem Orte, und in seinem Geschlechte und Namen, wie Mose sagt. Und er zeigt damit an, daß nicht eine solche Nachfolge gewesen sei, da Einer nach dem Andern regiert habe, wie unter den Königen gewesen ist.

82. Es reimt sich aber die Ordnung der Historie fein. Denn das Königreich Edom hat zweihundert Jahre gewährt, in welchen Jahren Jakob mit seinen Nachkommen in Egypten gewesen ist. Nachdem sie aber wieder aus Egypten gekommen sind, ist Idumäa wiederum in Fürstenthümer getheilt und durch Fürsten regieret worden bis auf Davids Zeit, der Idumäer bezwungen und unter sich gebracht hat. Wiewohl aber nun die andern Nachkommen vom königlichen Stamme gewesen, so werden sie doch gleichwohl nicht mitgezählet. Hernach, unter Joram, wird des Königes Edoms wiederum gedacht werden, da die Idumäer von den Königen Juda abgefallen.

83. Darum ist das Geschlecht Esaus ganz passend zu der Historie der zwei Patriarchen, Isaaks und Jakobs, gesetzt. Denn darin sehen wir die große Ehre und Herrlichkeit des Fleisches und der gottlosen Menschen in diesem Leben, welche Herrlichkeit, Reichthum und Glück vollauf haben. Sie sind Herren und Regenten in der Welt, und verachten entweder die Gottseligen, oder drücken sie, wie die Edomiter Jakob mit seinen Nachkommen mit großem Stolz verachtet haben. Dagegen aber Jakob, der die Verheißung hat und ein großer hoher Patriarch auch in der ganzen Welt ist, dazu auch ein Regent der rechten Kirche und Gemeinde GOttes, wird von wegen vieles Unglücks in seinem Hause und sonst auch vieler Jammer halb jämmerlich gedrückt und geplagt, also daß es scheint, als sei er des geistlichen Segens und der Gnade GOttes gar beraubt. Aber zuletzt, am Ende folgt, das im Lobgesange der Jungfrau Maria steht Luc. 1, 52.: „Er stößet die Gewaltigen vom Stuhl, und erhebt die Niedrigen.“ Das ist der Beschluß und das Ende mit den Gottlosen und Gottseligen; wie Christus sagt Luc. 6, 24. 25.: „Wehe euch Reichen, denn ihr habt euren Trost dahin. Wehe euch, die ihr voll seid; denn euch wird hungern“ etc. Dagegen aber Matth. 5, 4.: „Selig sind, die da Leid tragen, denn sie sollen getröstet werden.“

84. Darum sollen wir uns befleißigen, daß wir uns lieber zu den Armen und die da Leid tragen, so die Verheißung haben, halten, denn zu den Gottlosen. Denn wiewohl die Gottlosen eine Zeitlang groß Glück haben, so werden sie doch endlich gestürzt und vertilgt: die Gottseligen aber, die eine Zeitlang gedrückt werden, haben gewisse Zuflucht, Hülfe und Trost bei GOtt, der bei den Betrübten ist und sie erhält, ihren Jammer lindert, und die kleine Heerde, so wiederum aus ihrem Jammer und Elend herkommt, endlich mit ewiger Herrlichkeit und mit ewigem Leben begaben wird. Und wir sollen nicht kleinmüthig werden, auch nicht verzagen, ob wir wohl große Mühe und Arbeit, dazu auch viel Jammer und Unglück tragen müssen; sondern sollen es gewißlich dafür halten, daß gegenwärtig bei uns sei der Sohn GOttes, unser Herr Jesus Christus, der alle die, so in dieser Welt mühselig und beschweret sind, behütet, leitet und regieret, wie er verheißen hat Matth. 28, 20.: „Ich bin bei euch alle Tage bis an der Welt Ende.“ Welchem sei Lob und Ehre mit dem einigen Vater und Heiligen Geiste in Ewigkeit, Amen.

\chapter{Das siebenunddreißigste Capitel}

\section[Erster Theil]{Erster Theil. Von Joseph und dessen Zucht und Frömmigkeit, wie ihn sein Vater besonders geliebt, und ihm einen bunten Rock machen läßt}

\begin{quote}
\textbf{V. 1. Jakob aber wohnete im Lande, da sein Vater ein Fremdling immer gewesen war, nämlich im Lande Canaan.}
\end{quote}

1. Dieses ist nun die letzte Historie, von dem letzten Patriarchen Joseph, damit Mose dies sein erstes Buch, die Genesis, beschliessen, als mit einem sehr schönen und fröhlichen Ausgange oder Ende aller Dinge, so Jakob gehandelt, und vornehmlich seines ganzen Lebens und aller seiner Trübsale. Denn nachdem sein Haus und Gesinde ganz verwirrt, und voll großer Müheseligkeit, Angst und Sorgen gewesen, damit der sehr heilige Patriarch, wie wir gehört, hat kämpfen müssen (nicht sowohl von wegen der Feindschaft seines Bruders Esau und der Unbilligkeit seines Schwiegervaters Laban, welches er alles mit großem Muthe, ganz unverzagt, mit festem unüberwindlichen Glauben, dazu auch mit großer wunderbarer Geduld ertragen und überwunden hat; sondern vielmehr von wegen der Trübsale und Bekümmerniss, so sich in seinem eigenen Hause zugetragen und ihm daselbst begegnet, als, daß ihm seine Tochter geschmäht ward, daß ihm Debora, Rebekkas Amme, dahin starb und dazu auch sein liebes Weib, Rahel, daß sein Sohn Ruben solche große unsägliche Blutschande begangen, indem er seines Vaters Ehebette berührt und befleckt); und da er in solchen großen und hohen Beschwerungen diese einige Hoffnung und Trost in seinem höchsten Alter noch übrig gehabt, nämlich an seinem lieben Sohn Joseph, der der Erstgeborene war von seinem verstorbenen Weibe Rahel, welcher mit seinem gottseligen Wandel und heiligen Leben seines alten Vaters Herz, so für und für beschweret und bekümmert gewesen, gleichwohl noch etwas erfreut, getröstet und aufgerichtet hat: da wird ihm plötzlich und unversehens dieser auch hinweg gerissen; daß also der arme, elende, alte Vater, nachdem er erstlich sein liebes Weib verloren, nun auch dazu dieses seines geliebten Sohnes muß beraubt sein.

2. Denn Joseph ist verkauft und nach Egypten geführt worden, daß er dem Vater also aus den Augen genommen. Und ist solches geschehen nicht durch fremder Leute, oder aber derer, so Feinde möchten gewesen sein, sondern seiner eigenen Brüder Verbitterung, Haß und Bosheit. Und da dieselbigen solche ihre große Sünde betrüglicher Weise verhehlten, und ihrem alten Vater das blutige Kleid zeigten, als sollte Joseph von einem wilden Thiere gefressen sein, hat er, der Vater, nicht anders gemeint oder gewußt, denn demselbigen wäre also und sein Sohn Joseph wäre so erbärmlich und betrübt, denn demselbigen wäre also und sein Sohn Joseph wäre so erbärmlich und betrübt, denn dem Vater war, daß er nicht anders hätte hoffen können, als daß er verloren wäre.

3. Also wird der alte betagte Jakob in diese Angst und Bekümmerniss, ja, in die Hölle selbst geworfen, und daselbstige durch die Bosheit und Sünde seiner eigenen Kinder. Daher ist sein Glaube ohne allen Zweifel also hart und schwerlich angefochten und zerrüttet worden, daß er nicht ohne großen schweren Kampf an der Verheißung, so ihm GOtt gethan, wird haben halten können. Denn sein Herz ist nicht allein mit väterlicher Neigung gegen seinen lieben Sohn, sondern auch mit Zweifel und Unglauben bekümmert und angefochten gewesen, nämlich, ob auch GOtt seine Verheißung halten wollte, oder aber, ob er ihn um etlicher sonderlicher Sünden willen gänzlich verworfen hätte; sintemal er nun des Sohnes beraubt sein mußte, von welchem er gehofft, daß er der rechte Erbe sein würde des Segens, so ihm von GOtt verheißen war, und dazu aller seiner Güter.

4. So ist nun diese Historie und letztes Elend dieses heiligen Mannes Jakob die allerbeschwerlichste und sehr erschrecklich gewesen, und hat bald im Anfang einen sehr schweren und heftigen Kampf. Denn sie ist voll schweres und erbärmliches Leides oder Wehklagens, dadurch der heilige Altvater, so zuvor von so viel Kümmerniß, Leiden und Trübsal fast entkräftet, beinahe ganz getödtet ist; und hat sich noch über dies alles auch die schwere Sünde seines Sohnes Juda zugetragen, welcher mit Thamar, seiner Schwiegertochter, Blutschande begangen: welches denn auch dem alten Vater keine Besserung und Schmerzen sehr gemehrt haben wird.

\quellseite{269}

5. Wie aber dem nun sei, hat er gleichwohl noch immer mit wunderbarer Beständigkeit gute Hoffnung und festes Vertrauen behalten, und hat sich selbst aufgerichtet und getröstet mit so viel Exempeln der vielfältigen und augenscheinlichen Hülfe und Erlösung, so er sein ganzes Leben hindurch empfunden und erfahren hatte, bis ihm endlich wiederum ein neues Licht und Trost erschienen und begegnet ist.

6. Denn es ist gleichwohl endlich ein sehr wunderbarer und lieblicher Ausgang gefolgt, daß sein großes Elend und Schmerz, so sich unversehens zugetragen hatte, in sehr große und unerhoffte Freude ist verkehret worden. Denn Joseph, welches Tod der alte Vater betrauert hatte, wird wiederum aus der Finsterniß an das Licht, aus dem Tode zum Leben herausgeführt, und ist von Gott wunderbarlich erhalten, und dazu vom Könige in Egypten zu großer Ehre und Herrlichkeit erhoben und gesetzt worden, nachdem Jakob selbst auch aus vielen und großen Beschwerungen erledigt und erlöset war.

7. Und dies ist nun das Ende des ersten Buchs Mose, und ist das sechste und letzte Buch nach der Abtheilung dieser unserer Erklärung, wie wir dieselbige droben vorgenommen. Gleichwie aber in der ganzen Schrift nichts Schöneres ist denn die ganze Genesis, also ist unter den andern Historien der andern Patriarchen dies Exempel sehr trefflich und merkenswerth, und ist eigentlich damit also beschaffen, daß ich es weder mit Worten noch Gedanken genugsam fassen oder verstehen kann. Derhalben wollte ich wünschen, es wäre Andern vorbehalten auszulegen und zu erklären, die weit gelehrter und beredter sind, denn ich bin.

8. Dieweil aber gleichwohl die Ordnung der angefangenen Erklärung erfordert und mit sich bringet, daß wir zum wenigsten etwas sagen, wollen wir die Historie, welche wahrlich sehr wunderbar und hoch ist, überlaufen, dieweil wir ja dieselbe nach ihrer Würde nicht genugsam handeln und erklären können. Und muß aber nun zuerst aus dem, was droben gesagt, die Ordnung dieser Erzählung wiederholt werden. Denn am Ende des 35. Capitels hat Mose erzählt, wie der Patriarch Isaac, da er nun alt und lebenssatt geworden, gestorben und zu seinem Volk versammelt sei, und wie ihn seine zwei Söhne, Esau und Jakob, im Lande Canaan begraben haben. Von da an geht er zur Beschreibung des Geschlechts Esaus über, und nimmt daher Ursache, dieweil Esau bei seines Vaters Begräbniß gewesen, und ohne Zweifel mit ihm auch alle seine Söhne und Schwäger, welche dazumal mit Jakob versöhnt waren, da sich auch die zwei Brüder wiederum mit einander vereinigt. Derhalben, was nun folgt bis auf das 40. Capitel, ist alles geschehen, da der Patriarch Isaac noch am Leben gewesen. Darum haben wir auch Erinnerung gethan, daß diese Geschichten, nach rechter gebührlicher Ordnung der Historie, müssen dem 35. Capitel mit einverleibt werden, nämlich kurz vor dem letzten Stück, da V. 28. gesagt wird: „Und Isaac ward hundert und achtzig Jahr alt“ etc.

9. Und diese Erinnerung ist sehr nöthig, damit sich der Leser desto leichter in diese Beschreibung, so etwas verworren ist, finden und selbige verstehen könne. Denn die Andern alle, so dies Buch ausgelegt, und vornehmlich Lyra und Augustin, haben sich sehr gemartert mit dieser Stelle, daß auch Augustin, da er sich also darüber bemüht, daher Ursache genommen, schier an der heiligen Schrift und derselbigen Glaubwürdigkeit etwas zu zweifeln.

10. Eben also pflegt es den Antilogisten auch zu gehen, so die Sprüche der Schrift, welche lauten, als ob sie wider einander seien, zusammen lesen, und welche die Gewißheit der heiligen Schrift mit Fleiß zu schwächen und umzustoßen versuchen: wo sie an solche Stellen kommen, bleiben sie darin hangen, sind also verwirrt, daß sie weder hinter sich noch vor sich kommen können, und schreien alsdann heftig, die Predigten und Geschichten in der Schrift hängen nicht fein an einander, es sei alles verworren und eitel ungewiß Ding. Und dies ist auch die Ursache, daß solche Antilogisten zum rechten wahren Licht und Verstand göttlicher Dinge nimmermehr kommen mögen.

11. Es ist auch eine Regel bei den Kabbalisten, In lege non esse prius aut posterius, das ist, daß im Gesetz nichts stehe, das nach rechter Ordnung der Zeit, darin es geschehen, beschrieben sei. Aber diese Regel ist nicht überall und insgemein wahr, soll und muß auch mitnichten zu dem, so allhier erzählt und beschrieben wird, gezogen werden. Denn Isaac hat gewißlich noch gelebt, da Joseph nach Egypten verkauft worden, und noch zwölf Jahre darnach, da er wiederum aus dem Gefängniß erlöst, und zur Herrschaft und Regiment über ganz Egypten erhoben und gesetzt worden ist.

12. Also, die folgende Historie, da Thamar ist geschwächt worden, muß auch auf dieselbige Zeit gezogen werden. Denn die Sünde ist von Juda begangen, da sein Großvater Isaac auch noch am Leben gewesen. Wir wollen aber gleichwohl hernach die Rechnung der Jahre und Personen auch nach ihrer Ordnung erzählen. Denn so jemand die Zahl der Jahre von der Zeit, da Jakob endlich geworden, bis er nach Egypten hinab gezogen ist, rechnen wird, der wird finden, daß innerhalb derselben sechs und vierzig Jahre viele große und wunderbare Dinge geschehen sind. Und wiewohl die Historie denen, so nicht sehr erfahren, und die sich befleißigen, daß sie allzugenau die widerwärtigen Sprüche in der Schrift glauben, däucht ganz verworren zu sein, werden sich doch christliche Leser, so dieser Ordnung erinnert, leichtlich darein finden und die rechte Ordnung auch recht verstehen.

13. Nun geht aber Mose weiter in der Historie des Patriarchen Jakob, und schreibt, wie er seine Wohnung gehabt habe im Lande, darin sein Vater ein Fremdling gewesen. Denn er war zu Isaac gekommen, da er zu Hebron im Lande Canaan wohnte, und daselbst hat er auch sein Geschlecht und Hausgesinde bestellt, daß sie da auch mit ihm gewohnet. Das Land aber wird darum fremd genannt, weil, wiewohl es Abraham, Isaac und Jakob verheissen gewesen, sie gleichwohl darin noch nichts Eigenes oder Gewisses gehabt, ausgenommen des Rechtes über das Begräbniß, so Abraham gekauft, wie wir droben, Cap. 23, V. 17., gehört. Derhalben ist Jakob auch, wie seine Väter, als ein Fremdling von einem Ort zum andern gezogen, wie die Hirten nach Gelegenheit des Viehes und der Weide umher von einem Ort zum andern zu ziehen pflegen, Jakob aber hat gleichwohl in solcher Pilgerschaft seinen alten Vater Isaac nicht verlassen, da er seine Tochter Dina zu ihm gebracht, die geschwächt war, und seinen liebsten Sohn Joseph, welchen Isaac noch in seinem Alter ohne allen Zweifel mit großer herzlicher Lust und Freude gesehen und beherzt haben wird.

\subsection{I}

\begin{quote}
\textbf{V. 2. Und das sind die Geschlechter Jakob: Joseph war siebzehn Jahr alt, da er ein Hirte des Viehes ward mit seinen Brüdern; und der Knabe war bei den Kindern Bilha und Silpa, seines Vaters Weibern, und brachte vor ihren Vater, wo ein bös Geschrei wider sie war.}
\end{quote}

14. Es hat die heilige Schrift diesen Gebrauch, daß sie pflegt die „Geschlechter“ zu beschreiben, wenn sie auch nur einen einzigen Sohn anführt; wie an diesem Ort allein Joseph genannt wird, und die andern Söhne Jakobs ausgelassen werden, gleich als gehörten sie alle nicht auch mit in die Zahl und als wäre dieser Joseph allein der Sohn, dieweil er vor den Andern in diesem Geschlechtsregister aufgezählt wird. Und ist zwar dieser Joseph auch allein her, so billig und mit allem Recht den Andern allen mag vorgezogen werden. Denn er ist ja von dem vornehmsten Weibe Jakobs geboren und haben die Eltern seiner lange großes Verlangen nach ihm gehabt; zudem gebührt ihm auch vermöge der Ehe zwischen seinem Vater und Mutter die ganze Würde und Herrlichkeit der Erstgeburt, wiewohl er dieselbige damals noch nicht inne gehabt, wie wir hören werden.

15. Welches ist denn nun das Geschlecht Jakobs? Er hat einen Sohn, der heißt Joseph. Wie ist es aber um denselbigen bestellt, welches Schicksal ist ihm widerfahren? Höre, was Mose davon sagt: Joseph war siebenzehn Jahr alt, da er ein Hirte des Viehes ward mit seinen Brüdern. Also war es um Joseph bestellt, daß er ein Jüngling von siebenzehn Jahren war, von solchem Alter und Stärke, daß er wohl hätte mögen ein Ehemann sein, wie seine andern Brüder auch waren. Denn Juda, Ruben und die Andern hatten zumal schon Weib und Kinder. Aber dieser hat noch allein gelebt in seiner Jugend, ohne Weib und ohne eigen Hausgesinde, und er allein ist dem alten Vater lieb gewesen, daß er an ihm Lust und Freude gehabt, darum daß er geboren war von seinem allerliebsten Weibe, so verstorben war; und sein ganzes Leben, auch alles, dessen er sich beflissen, mit allen seinen Sitten, welche rein und unschuldig gewesen, haben ihn dem Vater wunderbarer Weise lieb und angenehm gemacht.

\quellseite{270}

16. Zu wem hat sich aber nun dieser Joseph gehalten? Er hat nicht Gemeinschaft gehabt mit den vornehmsten Söhnen, so sein Vater mit Lea gezeuget, nämlich mit Juda, Simeon, Ruben und Levi; denn dieselbigen haben Joseph, den Sohn der Rahel, aus Hoffahrt gar verachtet: zudem steht auch der gute fromme Jüngling, daß die Kinder der Mägde ganz verachtet und verpotet gewesen, und daß sie von den vier Söhnen der Lea für die allergeringsten leibeigenen Knechte sind gehalten worden. – So hat er sich derhalben vornehmlich zu solchen Verachteten gesellt. Denn die Andern sind sehr stolz und trotzig gewesen; wiewohl Ruben etwas gedemüthigt war wegen der Sünde und Blutschande, daß er seines Vaters Bett befleckt; und ich halte dafür, Juda wird gleichwohl auch etwas demüthiger gewesen sein denn die Andern. Aber Levi und Simeon sind sehr stolze Krieger und Todtschläger gewesen, welche auch Sichem und Hemor, seinen Vater, mit allen, so in der Stadt waren, mit der Schärfe des Schwerts erwürgt: die sind sehr grobe, hoffärtige, stolze Bauern und Filze gewesen.

17. Und ist wahrlich ein wunderbares Ding, daß sie so lange Zeit der heiligen Väter Zuhörer gewesen sind, Isaaks, Jakobs, auch Rebekkas und Deboras, ihrer Amme, und davon gleichwohl mit nichten frommer geworden sind und sich nicht gebessert haben. Denn also hat sie Mose hier beschrieben und abgemalt, daß die Sünde und Schande, so sie begangen, auch den schändlichsten Thaten und Exempeln der Heiden fast gleich geachtet werden mögen.

18. Derhalben dürfen wir uns nicht verwundern, wenn wir sehen, daß in dieser letzten bösen Zeit derer, leider! so viele sind, so von der rechten wahren Gottseligkeit abfallen, die so viel schöne und herrliche Predigten so gar ohne alle Frucht und Besserung hören; mögen uns aber daran genügen lassen, und damit zufrieden sein, daß doch etliche Wenige gefunden werden, so sich noch mit Lust zum Wort halten und daraus etwas gebessert werden. Denn wir sehen, daß es auch eben also dem sehr heiligen Patriarchen Jakob begegnet ist, in welches Hause und Geschlecht die, so unter seinen Söhnen die vornehmsten gewesen, ihm, dem alten Vater, den allergrößten Jammer, Trübsal und Schmerz gemacht haben.

19. An dem Jünglinge Joseph aber hat sich eine sonderliche Zucht und Frömmigkeit sehen lassen, daß er solche ungereimte bäurische Hoffahrt und Verachtung sein selbst und der andern Brüder, so von den Mägden geboren, so gar geduldig hat ertragen können, so es doch ihm, als dem erstgebornen Sohne, mehr geziemt hätte denn Simeon und Levi, daß er sich etwas überhoben und hervorgethan hätte.

20. Derhalben wollen wir diese Historie also handeln, daß wir den Stolz und Hoffahrt dieser zwei Söhne Jakobs gehässiger Weise groß machen, damit Gottes Gnade und Barmherzigkeit daraus auch desto mehr gesehen und erkannt werde, daß er aus der größten Bosheit gleichwohl das Beste machen und hervorbringen kann. Wir wollen aber mit nichten ihre Schande und begangene Ubelthat zudecken, sondern dieselbige so groß machen, als wir immer können; denn es ja in allewege unbillig und zu erbarmen ist, daß der gute fromme Vater so gar unbilliger Weise hat sollen behandelt, und von seinen eigenen Söhnen verpotet und betrübt werden; denn sie haben den frommen Jakob, ihren Vater, ja zu herzlich hoch betrübt, und haben mit dem armen alten Manne gespottet, daß es zu erbarmen ist.

21. Und Joseph hat eben dieselbige Ursache auch gehabt, daß er sich nicht hat wollen zu ihnen gesellen, sondern ist mit den andern verachteten Söhnen der Mägde vornehmlich umgegangen, und hat dieselbigen in Ehren gehalten; wiewohl er von der Mutter, so etwas edler gewesen, geboren war: er läßt die Mörder zu Sichem zufrieden. Daraus denn eigentlich die vortreffliche gute Art und der sonderliche Verstand dieses Jünglings hervorgeht; welches auch Mose an ihm allhier vornehmlich hat rühmen wollen. Denn es ist gewißlich eine gute Art an ihm, und ein feiner, sanftmüthiger, demüthiger Verstand gewesen. Dazu denn auch weiter dies gekommen ist, daß er gottesfürchtig und züchtig war, auch seine eheliche und stille Sitten gehabt, damit er nicht allein seines alten Vaters, sondern auch der Nachbarn Herz wird eingenommen haben, daß sie ihm sonderlich sind gewogen worden. Es ist derhalben kein Wunder, daß ihn sein Vater herzlich lieb gehabt, nicht allein deswegen, daß er sein Sohn war, sondern auch um der guten Art und schönen Tugend willen, damit er begabet gewesen.

22. So ist nun Joseph wie ein schöner heller Stern und wie der Morgenstern in dem Hause Jakobs, mit welchem die großen Bauern mit nichten zu vergleichen gewesen. Derhalben läßt er sie auch fahren und hält sich zu den verachteten Brüdern, hütet mit denselbigen des Viehes, wie das seines Vaters Wille und Wohlgefallen gewesen; so er doch der vornehmste Sohn war und dazu König über ganz Egypten hat werden sollen: der große König in Egypten muß in seiner Jugend der Schafe hüten.

23. Weiter wird aber Joseph auch daher gelobt, daß er vor seinem Vater gebracht und ihm angezeigt habe, wo seine Brüder etwas, das nicht recht war, gethan, und derhalben berüchtigt worden sind. Denn er hat die Gerechtigkeit und Ehrbarkeit geliebt, ist voll Liebe und Gehorsam gegen seinen Vater gewesen, also daß er vor ihm nichts hat verschweigen können, wo seine Brüder etwas begangen, dadurch Andere verletzt worden, oder daraus seines Vaters Geschlecht ein böses Geschrei bekommen oder geschmäht werden möchte. Derwegen wo er gehört, daß von seinen Brüdern etwas Böses oder Schändliches, das sie begangen haben sollten, ist gesagt worden, hat er es alsbald vor Jakob, seinen Vater, gebracht, als vor einen Richter, und dem es gebührte, solches in allewege zu verhüten oder zuvorzukommen, damit ja das Haus und die Gemeinde, welche mit so großen und herrlichen Verheißungen geziert und begnadet gewesen, mit solchen bösen Exempeln oder Thaten nicht entstellt oder befleckt würde, und daß nicht von den Fremden oder Ungläubigen der Name Gottes gelästert, und damit auch zugleich der Gottesdienst, so Jakob und seine Hausgenossen geübt und dessen sie sich gerühmt, geschmähet würde. Diese Pflege des rechten Gottesdienstes wollte Joseph ganz unverrückt und unverletzt, dazu auch mit keinerlei Aergerniß befleckt haben.

24. Und hat ihn ohne allen Zweifel sehr betrübt, daß er den Muthwillen und Frevel seiner Brüder hat sehen müssen, nämlich, daß sie alle sammt dem ganzen Hause Jakobs deshalb bei den Heiden so sehr haben müssen geschmäht werden. Derwegen er auch etlichemal dem Vater über den Ungehorsam und unbillige Thaten seiner Brüder geklagt hat, damit sie sich selbst unter einander oder aber die Nachbarn betrübt und beleidigt haben; wird gesagt haben: Siehe, lieber Vater, Ruben, Simeon und Levi, meine Brüder, sind zumal unruhige und händelsüchtige Menschen, richten oftmals Zärm, Unlust und Hader an, ärgern damit also die Völker, so allernächst um uns her wohnen, welche denn auch ganz fleißig und genau auf unser Leben und was wir für Sitten haben, lauern. Nun sollte ja billig dieses Aergerniß aus unserm Hause und Gemeinde abgeschafft werden, daß wir deshalb bei denen, so draußen und unsers Glaubens nicht sind, nicht gelästert werden möchten.

25. Wie zu unsern Zeiten wir auch mit Papisten umgeben sind, welche sehr scharfe Augen haben, sehen sehr genau auf uns, haben auch solche scharfe Nasen, damit sie gar leichtlich riechen können; denn sie lauern feindseliger Weise auf uns, ob sie etwas, das schändlich wäre, könnten von uns riechen oder inne werden, ja, ob sie unsere Worte oder Werke, so auch an sich selbst ehrbar und gut sind, lästern oder tadeln möchten. Derhalben auch der heilige Paulus fleißig ermahnt und sagt Eph. 5. 15.: „Sehet zu, lieben Brüder, wie ihr vorsichtiglich wandelt“, und daß ihr ja keine Ursache gebet dem Widersacher zu lästern. Item Col. 4, 5.: „Wandelt weislich gegen die, die draußen sind.“ Denn die Widersacher gehen fleißig damit um, daß sie auch etwa ein kleines und gar leichtes Geschrei, so von uns auskommen möchte, fassen mögen. Was aber Gutes an uns ist, das sehen sie nicht, achten dessen auch nicht, ja, verachten es wohl gänzlich, und haben nur Lust an dem, was an sich nicht gut ist, und wenn sie nur etwa ein böses Geschrei von uns hören mögen.

26. Derhalben hat der Patriarch Jakob sonderlich darauf Acht haben müssen als ein
\quellseite{271}
Bischof oder Lehrer der Gemeinde, damit in seinem ganzen Geschlecht ein gut Gerücht und ein guter Name bleiben möchte. Denn es geziemt sich ja mit nichten, daß ein Pfarrherr oder Seelenhirt ein unehrlich Leben führen und sein Gesinde oder Hausgenossen schändlich und übel regieren wollte. Und ist viel besser, Verachtung und bittern Haß der Widersacher tragen oder dulden, denn in Schande fallen und in des Teufels Strick, wie der heilige Paulus 1 Tim. 3, 7. sagt.

27. Gleiche Sorge hat Joseph, dieser Jüngling, auch gehabt und solchen Fleiß angewendet, nämlich, wie er allerlei Aergerniß, dadurch das Haus Jakobs, die Gemeinde und auch die Verheißung Gottes verachtet oder gelästert würde, möchte helfen verhüten. Und ist daran kein Zweifel, er wird deshalb von seinen Brüdern sehr gehaßt worden sein, darum daß sie des Vaters Ermahnung und Strafe nicht haben ertragen können, sind darüber ungeduldig geworden. Derhalben sie sich denn auch heftig erzürnt, wenn sie erfahren haben, daß Joseph ihre Uebelthat geoffenbart hatte; werden gesagt haben: Daß den Verräther alles Unglück treffe, der alsobald vor den Vater bringt, davon er hört, daß wir es sollen geredet oder gethan haben! Denn sie unterstanden sich, großen Muthwillen zu üben und allerlei böse Stücke zu treiben ohne Wissen ihres Vaters, haben auch wohl der andern frommen Leute in ihrem Geschlecht gute und freundliche Ermahnung oder Strafe eben also gedeutet und für Verrätherei gehalten; ja, sind noch wohl dazu aus bitterem Haß gegen dieselbigen auch immer ärger geworden, haben täglich ihre Bosheit und Sünde gehäuft, und sich darum gar nicht bekümmert, ob sie damit sich bei den Heiden umher guten oder bösen Namen machen möchten.

28. Es haben Lyra und Burgensis an diesem Ort fleißig und sehr genau dem nachgeforscht, was doch vornehmlich die große Sünde und Uebelthat möge gewesen sein, um derentwillen Joseph seine Brüder vor dem Vater verklagt habe. Denn also steht in der alten Uebersetzung. Es ist aber gleichwohl eine unnöthige Frage; denn aus dem hebräischen Text kann man es nehmen, was der rechte einfältige Verstand sei. Das hebräische Wort dibbath an sich selbst heißt Geschrei, Gerücht oder Meinung, so unter dem Volk ist ausgebreitet worden. Als, 4 Mos. 13, 33., da Mose Kundschafter in das Land Canaan gesandt und da dieselbigen alle wieder kommen, machen sie, daß ein böses Geschrei vom Lande ausgegangen ist, das ist, sie haben dem Lande ein böses Gerücht gemacht; denn sie sagten: „Das Land frißet seine Einwohner“ etc. Eben dasselbige Wort hat an diesem Ort Mose auch gebraucht, da er sagt, wie Joseph vor ihren Vater Jakob gebracht habe, „wo ein bös Geschrei wider sie wäre.“

29. Er hat sich aber gehalten zu den vier Söhnen der Mägd, Gad, Naphthali, Asser und Dan, dieweil dieselbigen etwas demüthiger gewesen, denn die andern vier, welche sich über die Maßen sehr erhoben und stolzirt haben wegen der Erstgeburt, daß sie von der Mutter Lea geboren waren, und machten sich selbst also beide des Lebens und böser Sitten halben ein böses Geschrei, scheuten sich auch vor ihrem Vater gar nicht, sondern meinten, es wäre ihnen frei, nach ihrem eigenen Muthwillen zu thun Alles, was sie nur gelüstet, dieweil sie eines namhaftigen und hochberühmten Patriarchen Kinder wären und so herrliche große Verheißung hätten.

30. Joseph aber, dieweil er Zucht, Friede und Ehrbarkeit geliebt, war betrübt wegen der Bosheit dieser seiner Brüder, und hat oftmals vor den Vater gebracht, wo sie mißgehandelt und gesündigt, hat ihn ermahnt und immer bei ihm angehalten, daß er doch im Hause und der ganzen Gemeinde ein etwas ernstlicheres Einsehen hätte, damit Zucht und Ehrbarkeit, wie sich das gebührt, möchte gehalten werden.

31. Und das ist die einige Ursache gewesen, daß er also heftig von ihnen ist gehaßt worden, nämlich, daß Joseph immer so ein fleißig Aufsehen gehabt, damit es ja in seines Vaters Hause recht zugehen möchte, und daß er begehrt, daß vornehmlich die Ehre bei jedermann ihren guten Namen unverletzt behalten möchte. Er ist fromm, und hat derhalben auch die Frömmigkeit, und was dazu gehört, gern gefördert; Aergerniß aber wollte er auch aufgehoben und abgeschafft haben. Darüber ist es ihm angegangen, wie es den Frommen zu gehen pflegt, und hat damit seiner Brüder Haß auf sich geladen und sich von ihnen darüber verfolgen lassen müssen. Nun wird aber weiter im Texte auch noch eine andere Ursache gesetzt.

\subsection{II}

\begin{quote}
\textbf{V. 3. Israel aber hatte Joseph lieber, denn alle seine Kinder, darum, daß er ihm im Alter gezeuget hatte; und machte ihm einen bunten Rock.}
\end{quote}

32. Dies hat den großen stolzen Junckern allererst schwer und unleidlich zu sein gedünkt; denn sie haben nicht wohl können ertragen, daß Jakob den Joseph lieber gehabt hat denn die andern Kinder, dieweil er ja noch fast ein kleiner Knabe und der Jüngste unter ihnen war. Sie werden gesagt haben: Wir sind ja die rechten Herren, als die wir geboren sind von der Mutter Lea, so unsers Vaters Jakob erstes Weib gewesen, ehe denn Joseph ist empfangen worden. Ich, Ruben, bin bald im ersten Jahre nach der Hochzeit geboren, Simeon im andern, Levi im dritten Jahre und so immer fort; Joseph aber ist am allerletzten geboren. Derhalben sind wir die, an denen das ganze Wesen, oder die vornehmste Macht und auch die Herrlichkeit dieses Hauses hängt. Rahel aber mit ihren Kindern ist wie eine Zugabe oder etwas Zufälliges. Derhalben gebührt uns die Ehre, die Liebe und Gunst des Vaters, und auch dieses ganzen Hauses und Geschlechtes.

33. Um solcher Ursachen willen hat sich bei ihnen erhoben und ist wie ein Feuer angegangen der bittere, feindselige und greuliche Haß gegen ihren Bruder Joseph; denn sie sind nunmehr nicht sehr wohl bei Sinnen, sondern voller Bosheit und vom Teufel gar besessen gewesen.

34. Es gibt es zwar die Vernunft selbst und der gemeine Witz oder Verstand aller Menschen bringt es je also mit, daß Joseph von seinem Vater Jakob recht und billig, und auch natürlicher Weise vor den Andern ist geliebt worden, darum daß Rahel Jakobs vornehmstes Weib gewesen, und Joseph von ihr fast in seines Vaters Alter geboren, da die Eltern mit Verlangen ganze sieben Jahre auf ihn gehofft und gewartet, zudem, dieweil er so guter Art und mit allerlei Tugenden geziert war. Einen solchen Sohn würde ich selbst auch lieber haben denn die andern, der sich so fromm, züchtig und gehorsam erzöge in allen Dingen, ob er schon nicht der erstgeborne Sohn wäre.

35. Derhalben, sage ich, war diese Liebe, so Jakob zu seinem Sohne Joseph getragen, auch von Natur ganz billig und recht, welche doch die Brüder Josephs, so von boshaftigem bitterm Haß ganz entbrannt waren, mit nichten haben ertragen können; denn sie waren so gesinnt, daß sie selbst von Joseph wollten angebetet und verehret sein; wollten von ihm hören, daß er sie also grüßen und anreden sollte: Ihr seid Herren, die vornehmsten Juncker und Regenten in des Vaters Hause; alles, was euch nur gelüstet und wohlgefällt, soll mir auch gefallen.

36. Solche Frömmigkeit und Demuth haben sie an ihrem Bruder gesucht und von ihm haben wollen. Dieweil aber Joseph sich immer vor dem Vater gescheuet und ihm auch gebührliche Ehre erzeigt, endlich auch ein fleißig Aufsehen gehabt, daß es in des Vaters Hause allenthalben züchtig und auch göttlich zugehen möchte: um der Ursache willen haben sie also erbärmlicher Weise den frommen alten Vater geplagt und gemartert, und sind zugefahren, haben auch den frommen Knaben Joseph betrüglicher Weise und mit Gewalt aus dem Hause in das Elend gestoßen. Denn sie haben mit nichten leiden können, daß man ihnen hat zugeredet, sie im Zaume gehalten und ihren Muthwillen nicht gestatten wollen; ja, sie haben ihrer Bosheit halben mehr wollen gelobt denn gestraft sein.

37. Eben solche Bosheit und verkehrtes Wesen geht auch jetzt mit Gewalt an durch alle Stände, wird in der Kirche, im weltlichen und Hausregimente also geübt. Denn jedermänniglich wüthet und tobt wider den, so ihm zuredet und ihn seines Amtes erinnert. Man zürnt heftig mit denen, so die Sünde und Laster, auch die, so gar gemein und grob sind, strafen. Man soll niemand einreden, sondern einen Jeden thun lassen, was er will; und ihren Frevel und Muthwillen zu beschönigen, dürfen sie gar herrlich vorwenden, Einer seinen hohen Stand, der Andere seinen Reichthum, oder ein Anderer sein Vaterland, oder auch seine Eltern, daß man um solcher vermeinten Ursachen willen ihnen müsse durch die Finger sehen, und sie ihrer Sünden und Schanden halben mit nichten strafen oder schelten solle. Ich, sagen sie, bin mit im Regimente, derhalben sollte man ja auch billig des Standes verschonen und denselbigen nicht also schmähen. Ich bin Bürger zu Wittenberg, zu Nürnberg, derhalben ich mir nicht unbillig mehr Gerechtigkeit und Freiheit zeugne, denn ein Fremder oder Einkömmling. Und wiederum, die da ehrbar und bescheiden sind und in den Schranken ihres Amts und der Gesetze bleiben, dieselbigen werden fast von Jedem gehaßt; wie es denn gleichfalls dem frommen Joseph auch in seinem Geschlecht gegangen ist.

\quellseite{272}

38. Darnach ist diese Ahnsucht, so die Brüder Josephs gegen ihn getragen, dadurch auch noch größer geworden, daß ihn der Vater vor den Andern sonderlich geliebt. Denn wiewohl Jakob die andern Söhne auch geliebt, als die von ihm geboren waren, ist doch gleich aus guten rechtmäßigen Ursachen die väterliche Zuneigung zu Joseph viel kräftiger gewesen; und dieweil sie das nicht haben ertragen können, sollte man ihnen billig feind sein.

39. Und ich verwundere mich, daß man sie noch so lange in der sehr heiligen Gemeinde Jakob hat dulden können. Ich hätte sie zwar alle umgebracht oder aus dem Hause gestoßen. Denn, lieber Gott! was ist doch das für große Unsinnigkeit und Bosheit, mit einem solchen Vater zürnen, der seinen Sohn, welcher fromm und gottselig, dazu auch von seinem liebsten Weibe geboren ist, liebt? Derhalben haben wir auch gesagt, daß man diese Bosheit müsse groß und schwer machen, so viel es immer möglich; denn sie machen es zu grob.

40. Uebrigens erhebt sich allhier eine Disputation unter den Grammatikern: Warum doch Mose die Ursache der Liebe, so Jakob zu Joseph getragen, also ausdrücklich gesetzt? nämlich, darum daß Joseph gezeuget und geboren worden sei, da Jakob, sein Vater, nunmehr alt gewesen. Denn man kann aus vielen Argumenten das Gegenteil beweisen und darthun. Und zudem zeigt auch die Historie selbst das Gegenteil an. Burgensis hat angezeigt, daß im chaldäischen Text also stehe: Dilexit eum, quia sit in filiis sapientiae senum, das ist: Joseph war ein solcher Sohn, und war an ihm so große Weisheit und Verstand, daß er es auch wohl einem jeglichen Alten in der Weisheit hätte zuvor thun mögen; und möchte der Spruch Esaiae am 65. Cap. V. 20. wohl hieher gezogen werden: „Die Knaben von hundert Jahren sollen sterben“, das ist, sie werden alle über ihr Alter weise werden; die Knaben werden gleich so weise werden, als wären sie hundert Jahre alt.

41. Darüber streitet Lyra, und will haben, daß Joseph darum ein Sohn des Alters genannt werde, daß ihn sein Vater im Alter gezeuget habe. Ich aber kann nicht verstehen, wie sich dieser Verstand mit dem reime, so hernach auch von Benjamin gesagt wird, welcher beinahe dreizehn Jahre nach Joseph geboren ist: und wird auf gleiche Weise von Juda, seinem Bruder, vor Joseph ein Sohn des Alters genannt.

42. Wiewohl ich nun solche Disputation den Hebraisten zu befehlen und heimzustellen pflege, so scheint doch die Ursache, der Lyra nachfolgt, zu schwach, sintemal Benjamin dreizehn Jahre jünger ist; Joseph aber ist der liebste Sohn, nach Dan und Naphtali, Gad aber und Asser sind entweder nach Joseph oder aber mit ihm auf Eine Zeit geboren. Zuletzt aber sind geboren Gad, Asser, Isaschar, Sebulon und Dina, und sind dieselbigen alle jünger denn Joseph. Wie will aber nun gleichwohl dieser Text bestehen, daß Joseph soll geboren sein, da sein Vater Jakob alt gewesen? Ja, die Ordnung der Historie zeigt über das noch an, daß eben dasselbe auch also von allen Söhnen Jakobs möge gesagt werden, sintemal Jakob da allererst Kinder zu zeugen angefangen hat, da er vier und siebzig Jahre alt gewesen. Derhalben will ich diese Frage den Hebraisten befehlen, daß sie dieselbe auflösen; denn ich weiß, daß die hebräische Sprache mancherlei Art und Weise hat zu reden, so uns jetzo fast unbekannt, und deren eigenthümliche Bedeutung wir auch schwerlich wiederum werden vollkommen herzustellen können.

43. Die Meinung des Burgensis ist allegorisch und sehr schön, daß Joseph ein Sohn des Alters gewesen, das ist, er sei weise und verständig gewesen. Ich wollte aber das wohl gern sagen, und die Worte und die ganze Historie lauten auch fast dahin, als sollte dies der rechte Verstand sein, daß der Patriarch Jakob gemeint, daß er nicht, denn allein dieser zwei Söhne Vater sei, nämlich, Josephs und Benjamins, so von Rahel geboren waren: wie er denn hernach, 1 Mos. 42, 36., selbst sagen wird: „Ihr beraubet mich meiner Kinder. Joseph ist nicht mehr vorhanden, Simeon ist nicht mehr vorhanden, Benjamin wollt ihr hinnehmen; es gehet alles über mich“ etc. Denn also muß man die Worte, so im Text stehen, lesen, bewegen und betrachten, daß sie auf die väterlichen Zuneigungen gezogen werden, daß der alte Vater Jakob gedacht: Rahel war mir von Gott verheissen und gegeben, darum ist sie auch die rechte Mutter, und diese zwei Söhne, so von ihr geboren, werde ich auch für meine rechten und liebsten Kinder halten: nun hat sie aber in meinem Alter den Joseph und Benjamin geboren, darum habe ich dieselbigen auch sonderlich lieb. Das ist eigentlich der Verstand, wie es Jakob gemeint hat, daß man ihm also auch hätte die Wahl gegeben unter den Kindern der Rahel und der Lea, welcher er am allerliebsten hätte können oder wollen enthalten, würde er ohne allen Zweifel diese zwei Söhne seines Alters den Andern allen vorgezogen haben.

44. Auf diese Weise habe ich diesen Text ausgelegt, nämlich von der Zuneigung und Liebe des väterlichen Herzens, daß Jakob nach so vieler und mancherlei Angst, Kreuz und Widerwärtigkeit sich nur zu diesem Sohne gehalten, welcher geboren war von der rechten und vornehmsten Mutter, die auch in der Geburt umgekommen und zu seinem großen Schmerz gestorben war. Derhalben er auch an demselbigen Sohn seinen einigen Trost, Lust und Freude gehabt. Denn er allein hat den alten Vater in so vieler Angst und Sorgen trösten können, da ihn die andern Söhne sehr bekümmert und gemartert, ja, dazu ihm noch wohl täglich neue Bekümmerniß gemacht und dieselbige gleichsam gehäuft haben.

45. Es ist aber die Kunst, geneigter guter Wille und Liebe des Vaters daher noch immer größer geworden, daß er an dem Jünglinge sonderliche Gottseligkeit und allerlei Tugend gesehen hat. Denn es ist auch Joseph in Wahrheit weise, züchtig, fromm, gottselig, dazu mit hohen vortrefflichen Gaben sonderlichen Verstandes gezieret und begabt gewesen. Darum ist er auch Jakobs, seines Vaters, Herz, Freude und Wonne. Denn also pflegen die Eltern zu thun, wo Einer auch von zwanzig Kindern wäre, und unter allen nur Einen hätte, der fromm und gehorsam wäre, da, sage ich, würde der Vater nicht anders gegen das fromme Kind geneigt sein, denn als ob er sonst keines mehr hätte: er würde die gottlosen und ungehorsamen enterben, oder aber zum wenigsten sich aller Sorge und Bemühung um sie entschlagen.

\subsection{III}

46. Nun ist aber noch ein anderer grammatischer Anstoß, nämlich, von dem bunten Kleid. Was aber das für ein Kleid gewesen sei, bekenne ich frei und öffentlich nicht zu wissen. Wir haben es in unserer deutschen Uebersetzung also gegeben: „einen bunten Rock“; wie es denn sonst auch gemeiniglich ist verdolmetscht worden. Lyra sagt, daß es vestis serica, das ist, ein seidener Rock gewesen. Burgensis geht an diesem Ort stillschweigend vorüber, sagt davon nichts. Die Andern halten dafür, es sei ein Kleid gewesen, das gewirkt sei von mehr denn einerlei Garn, das ist, von mancherlei Farben, da der Zettel roth und der Einschuß himmelblau ist; ein Schiller, da sich zwei Farben unter einander gemenget zeigen. Mich däucht aber solche Meinung nicht wahrscheinlich; wiewohl ich auch nichts weiß dawider zu reden.

47. Das hebräische Wort, passim, heißt, ein leinener Rock oder von Leinwand gemacht, welches ein griechischen Text gegeben ist, πολύπορον, das ist, der viele Fäden hat. Ob es aber ein bunt oder zertheilt Kleid gewesen sei, kann aus dem Text eigentlich nicht verstanden werden, und dies Wort steht sonst an keinem Ort mehr in der Schrift, denn allein im 2. Buch Samuelis am 13. Capitel von Thamar, so von ihrem Bruder Ammon geschwächt worden. Denn also steht im Text am selbigen Ort V. 18. 19.: Thamar warf Asche auf ihr Haupt und zerriß den bunten Rock; da steht, sage ich, auch im Hebräischen das Wort passim, und sagt der Text daselbst auch, daß solche Röcke (im Hebräischen meilīm) des Königes Töchter, weil sie Jungfrauen waren, getragen haben. Ich wollte es wohl gerne verdeutschen: einen schönen weißen Kittel, ein weiß Kleid, wie solche Kleidung in den Morgenländern gebräuchlich gewesen; denn in denselben Ländern sind solche weiße Kleider, von feiner Leinwand oder weißer Seide gemacht, die schönste und ehrbarste Tracht gewesen: wie auch wiederum die Völker, so gegen Abend oder Niedergang der Sonne wohnen, Purpurkleider zu tragen Lust gehabt.

48. Derhalben wolle man mir erlauben, daß ich diesem Dinge möge nachdenken, hier etwas errathen und dies also auslegen: daß Jakob seinem allerliebsten Sohn Joseph, als dem künftigen Erben und Priester in seinem Hause, einen solchen Rock gemacht, damit er ihn vor seinen andern Brüdern zierte und also von denselben unterschied. Denn es ist ein priesterlich Kleid gewesen, wie Hanna dem Samuel auch ein solch Kleid gebracht, 1 Sam. 2, 19., einen kleinen Rock (im Hebräischen meil), so dem gleich gewesen, damit Thamar bekleidet war.

\quellseite{273}

49. Ein solcher meil, wie das hebräische Wort lautet, wollte ich gerne sagen, wäre Josephs Kleid auch gewesen. Denn also pflegt man auch das Kleid der Engel zu malen, nämlich schneeweiß, und das von vielem Leinen Garn oder Fäden gemacht oder gewirkt ist; und ich will glauben, daß auch unsers HErrn Christi Rock ebenso gewesen, davon die Evangelisten sagen, er sei ungenähet gewesen, von oben an gewirket durch und durch, Joh. 19, 23., welchen seine Mutter Maria wird gewirket haben: ein gewirkter gestrickter Rock. Und ist mir in der ganzen heiligen Schrift ein fast gemein Wort, wo die Engel sind erschienen; desgleichen auch in der Verklärung Christi, da Matth. 17, 2. Marc. 9, 2. also steht: „Und sein Angesicht leuchtete wie die Sonne, und seine Kleider wurden helle und sehr weiß, wie der Schnee.“ Desgleichen auch das Kleid des Engels, so bei dem Grabe Christi erschienen, davon auch der Evangelist, Matth. 28, 3., sagt, daß es weiß gewesen sei wie der Schnee. Und Salomo sagt auch in seinem Prediger am 9. Cap. V. 8.: „Laß deine Kleider immer weiß sein.“ Mose aber hat die Obrigkeit auch mit einem hebräischen Wort genannt horim, das ist, weiß bekleidet, nämlich daher, daß sie weiße Kleider getragen. Wie man die vom Adel nennen möchte nobilitas torquata daher, daß sie goldene Ketten am Halse tragen, damit sie geziert sind.

50. Gleicher Weise und Gestalt scheint Jakob den Joseph auch geziert zu haben mit einem sonderlichen Kleid, das weiß und priesterlich gewesen, daß er damit anzeigte, wie er zu dem Sohn sonderliche Liebe trage, und daß damit auch die künftige Herrlichkeit, so Joseph haben würde wegen der Priesterschaft, bedeutet werden möchte. Wer aber nun lieber die vorige Deutung behalten will, daß es ein bunt Kleid gewesen sei, mit dem will ich auch nicht hart streiten.

51. Es scheint aber doch gleichwohl, daß der Thamar Kleid weiß gewesen. Denn in derselben Historie allein steht dies hebräische Wort passim: sie wird einen schönen weißen Kittel getragen haben. Darnach bewegt mich auch die Beschreibung des priesterlichen Kleides 2 Mos. 28, 39., da Mose den Priestern zuerst ein Unterkleid zeuget nach dem leinenen Hemde, כתונת, Kutte. Der leinene Rock aber war eng, ein Unterrock, welcher alle Gliedmaßen am ganzen Leibe bis auf die Schenkel bedeckte, der nicht weite, sondern enge Ärmel gehabt, ohne Falten, daß die Priester so viel fertiger und geschickter sein möchten zum Opfer und das Vieh zu schlachten, wie unsere Fleischer.

52. Zum Dritten thut Mose noch dazu das priesterliche Kleid oder des Hohenpriesters Rock. Derhalben hat der Hohepriester auch ein solch Kleid gehabt, so auf Hebräisch meil heißt, welches ganz von gelber Seide war: der schöne recht gelbe seidene Rock. Und ist allein des Hohenpriesters Rock gelb oder goldfarben gewesen, denn auch das Kleid Samuels gleich war, so ihm die Hanna, seine Mutter, gebracht; wiewohl das eine andere Farbe gehabt: ein Chorhemdlein. Denn das auf Hebräisch meil geheißen, ist weiß oder silberfarben gewesen, dessen auch Andere, so nicht Priester waren, haben gebrauchen mögen; aber mit dem Kleid von gelber Seide oder Goldfarbe sind allein die Hohenpriester geziert worden.

53. Also liest man von Hiob, daß er sein Kleid zerrissen habe, Hiob 1, 20.; item, wie David dem Saul in der Höhle einen Zipfel vom Rock geschnitten, 1 Sam. 24, 5.; an welchen Stellen auch dasselbe hebräische Wort meil steht, daraus wir schließen können, daß es ein königlicher und großer Herren Schmuck gewesen sei. Denn in den Morgenländern haben die Könige und große Herren nicht goldene Ketten oder Purpur, sondern ein weiß seiden Gewand getragen. Die Juden haben auch nicht seidene Kleider, sondern Pelze getragen, haben darin dem Exempel Adams und Evas gefolgt. Wie wir zu unsern Zeiten die Kinder auch erst mit Pelzen bekleiden; hiernachst pflegen sich auch die Erwachsen derselben zu bedienen. Aber die Art der Kleidung, so jetzund neulich in Gebrauch gekommen, welche also schändlich zerhackt oder zerschnitten wird, ist zwar nicht werth, daß man sie soll Kleider heißen, sondern sind eitel seltsame, unnatürliche, erschreckliche Ungeheuer dieser letzten unseligen bösen Zeit.

54. Zum Letzten gehört hieher auch, daß Mose verboten hat, daß das Volk nicht sollte Kleidung tragen, so von mancherlei Garn, oder von Wolle und leinenem Garn gemacht; wie er auch nicht haben will, daß man den Samen und die Thiere vermischen soll. Derhalben ist dies einfach mein Gedanke und einfältige Meinung, daß es ein Rock gewesen sei von einerlei Art Garn gemacht, dazu weiß und ungenähet, gewirkt oder gestrickt.

55. Dies sei gesagt von Josephs Rock, damit Jakob hat anzeigen wollen die Liebe, so er zu diesem Sohn getragen, als den er mit solcher Hoffnung auferzogen, daß er sollte Priester werden in seinem Hause. Darum er ihn denn auch mit sonderlichem Schmuck von den andern Häusern seiner Brüder abgesondert und unterschieden hat, daß sie dabei wissen möchten, daß dieser Joseph der oberste Regent und Erbe in seinem Hause werden sollte.

56. Daher ist nun entstanden der große greuliche Haß in den Herzen derer, welche der Teufel selbst getrieben und angehetzt hat, daß sie ihrem leiblichen Bruder mit so erschrecklichem Neide und grimmigem Zorn ohne alle Ursache zugesetzt, unangesehen, daß sie deshalb zum öfteren beide von ihrem Vater und auch andern gottseligen Leuten, so auch mit in der Gemeinde Jakobs gewesen, werden ermahnt und davon abzustehen erinnert worden sein. Ja, sie haben auch darin gar nicht betrachtet die vielfältigen Wohlthaten, damit sie ihr allerliebster Vater, als seine Kinder, täglich überschüttet, und daß sie bei ihm beide mit Weib und Kind, sammt allem ihrem Hausgesinde in großer Anzahl, ihren Unterhalt gehabt, indem er sie so gar väterlich genährt: für diese Wohlthaten alle gedenken sie nämlich dem Vater damit zu danken und ihm dieselbigen zu bezahlen, wenn sie ihren Bruder Joseph zum Hause hinaus stoßen, oder aber ihn sammt dem alten Vater tödten werden.

57. Es wird aber dies Exempel uns zum Trost also vorgehalten, daß wir nicht wehmüthig werden oder verzagen, wenn wir dergleichen auch erfahren müssen; sondern daß ein Jeder unter uns sein Amt fleißig ausrichte, jeglicher nach seinem Berufe, mit Lehren, Strafen, Drohen beides in der Kirche, und auch in weltlichen und Hausregimente. Es wird unsere Arbeit nicht vergeblich sein im HErrn. So aber Etliche sich weigern werden, unserer christlichen Vermahnung zu folgen und zu gehorchen, dieselben müssen endlich zu ihrer großen Gefahr und Schaden lernen klug sein.

\begin{quote}
\textbf{V. 4. Da nun seine Brüder sahen, daß ihn ihr Vater lieber hatte, denn alle seine Brüder: waren sie ihm feind, und konnten ihm kein freundlich Wort zusprechen.}
\end{quote}

58. Ei, welch eine sehr schöne und brüderliche Liebe ist das, so in der Kirche, welche Jakob in seinem Hause gehabt, gespüret worden, da doch ohne Zweifel so eine feine Zucht und Ordnung von den Patriarchen wird gehalten worden sein, und da zugleich mit Jakob auch der ältere Patriarch Isaak noch am Leben gewesen! Denn die Brüder Josephs haben ihn, der doch ganz unschuldig war, so sehr gehaßt, daß sie ihm auch nicht ein freundlich Wort haben können zusprechen; ja, sie haben ihn angeschnautzt, haben alle seine Worte und Werke wissen zu tadeln. Und wie er dem Vater der Allerliebste gewesen, also haben wiederum seine Brüder niemand väterlicher gehalten als ihn. Sie haben ihn sehr heftig gehaßt und dasselbe aus lauter Bosheit; vornehmlich aber Simon und Levi. Denn ich halte gleichwohl, daß sie alle ursprünglich und anfänglich nicht gleich gesinnet gewesen. Denn Ruben war etwas gedemüthigt darum, daß er sich schuldig gewußt wegen der begangenen Sünde, und die Andern, halte ich, seien auch sanftmüthiger und nicht so gar rauh gewesen. Derhalben soll dieser Haß der Brüder Josephs also verstanden werden, wie das Murren der Apostel, da sie darüber unwillig geworden, daß Maria Magdalena den HErrn gesalbet, Matth. 26, 8.; denn die andern Jünger waren auch dazu von Juda überredet und angereizt worden. Und daß ich dies auch beiläufig hinzufüge, daß es Hieronymus dafür hält, Judas Ischarioth komme vom Stamm Simeon her.

59. Wir wollen aber damit diese Sünde nicht klein machen oder einen von diesen Brüdern entschuldigen, sondern haben gesagt, daß ihre Bosheit auf alle Weise soll groß gemacht werden, nicht allein aus Haß wider die Sünde und den Teufel, sondern auch uns selbst zu Trost; sintemal Gott in der Gemeinde und in Geschlechte der heiligsten Patriarchen hat solche giftige böse Leute ertragen können, denen er doch hernach ihre Sünde gnädiglich vergeben. Denn wiewohl sie endlich noch Buße gethan, ist es doch mit nichten zu verschweigen, daß sie zu der Zeit mit solcher Bosheit und so bitterem Haß gegen den unschuldigsten Menschen gleichsam entbrannt gewesen.

\quellseite{274}

60. Nun läßt es aber gleichwohl Gott geschehen und steht gleichsam durch die Finger bei solchem unbilligen und boshaftigen Haß; ja, er gibt ihnen damit große Ursache, ihren Bruder zu meiden, jedoch zu ihrem eigenen Verderben. Denn sie thun wie die Juden, die auch Christum, unsern Heiland, aus großem Haß verfolgten; suchten immer Ursache, in allen seinen Worten und Werken, ihn zu lästern und Uebels von ihm zu reden, so er doch für und für mit großer Güte und Freundlichkeit, auch mit mancherlei Wunderwerken und Wohlthaten das ganze Volk überschüttet und gleichsam überhäuft hat: noch muß Alles unrecht sein. Denn sie haben sich der keines gefallen lassen, bis Gott ihnen endlich verhängt, diesen bittern Haß ihres Gefallens ganz auszuschütten und damit ihr Muthlein zu fühlen, indem er seinen eingebornen Sohn ihnen in die Hände gegeben, daß sie ihn ans Kreuz geschlagen, und also damit die unbillige Grausamkeit und bittern Haß, so sie gegen Christum, den Herrn, gehabt, erfüllten und sich damit also sättigen möchten.

61. Aber, o wie unseelige und elende Leute sind das, welche Gott also sich selbst übergibt und deren Grimm und bösen Lüsten er nicht widersteht! Wiewohl wir zwar alle, leider! also von Natur gesinnet sind, daß wir nicht wohl leiden können, daß wir im Zaum gehalten werden; aber wehe denen, welchen unser Herr Gott dermaßen zusteht, daß sie ihres Gefallens Sünde treiben. Denn es pflegt doch endlich zu folgen, wie man gemeiniglich sagt: Qui, quae vult, dicit, quae non vult, audiet, das ist: Wer alles sagt, was ihm gelüstet, muß zuletzt wiederum auch hören, das er nicht gerne hört. Also auch: Wer alles thut, was ihm gelüstet, muß dermaleinst auch dafür leiden, daß er wohl lieber hätte wollen überhoben sein. Denn wir wollen uns selbst nicht gern im Zaum halten, und wollen auch unsere Lust und böse Begierde nicht einschränken, wie es Andere uns vorschreiben, unterjagen und rathen.

62. Wo uns denn nun Gott auch etwas übersieht, fallen wir darüber in solche Thorheit, Jammer und Noth, daß wir, wie die Juden, mit greulichen Läster- und erschrecklichen Strafen dermaßen umgeben werden, daß wir nicht leichtlich wiederum können herauskommen, und daß wir endlich allzuspat an uns selbst Abscheu haben werden und uns selbst schämen müssen. Also ging es mit den Sodomitern auch: die sonnte man keinerlei Weise enthalten oder zwingen, daß sie nicht also ihres Gefallens tobten wider den frommen und heiligen Mann Lot; aber es währete nicht lange, daß sie greulich gestraft, und plötzlich umgebracht und getödtet wurden. Es mußte gehen, wie sie wollten; des Morgens kam das höllische Feuer vom Himmel, und that, daß sie nicht wollten.

63. Also auch vor der Sintfluth wollte niemand Noah hören, sondern nahmen zu Weibern, welche sie nur wollten: derhalben kam die Sintfluth, und that auch, das sie nicht gerne wollten. Also will man es haben. Und werden eben dasselbe unsere Leute, denen wir jetzund auch predigen und die wir täglich zur Buße ermahnen, gewißlich auch also fühlen und erfahren, dieweil sie ja alle Zucht und Ehrbarkeit wollen ganz und gar aufgehoben haben, wollen sich mit nichten einreden lassen.

64. So ist nun dies gewißlich ein heiliges, gesegnetes und das allerfürtrefflichste Leben, wenn Gott nicht durch die Finger sieht, wo wir sündigen, und das zu thun vorhaben, so sich nicht ziemt; sondern wenn er uns alsbald straft und züchtigt mit seiner Ruthe, mit Trübsal und sonst allerlei Kreuz, oder dasselbe ausrichtet durch Menschen, so uns mit fleißigem Vermahnen wiederum auf den rechten Weg bringen, damit unsere bösen thörichten Lüste also bezwungen werden, daß wir ihnen hinfort nicht mehr folgen oder nachhängen.

65. Wiederum aber ist das ein sehr böses Zeichen, wo einem bösen Buben und der immer Böses zu thun im Sinne hat, Alles glücklich, und wie er es selbst begehrt, hinausgeht. Denn also sagt auch Jeremia im 12. Cap. V. 2. von den falschen Propheten: „Du lässest sie viel von dir rühmen, und züchtigest sie nicht“ das ist, alles, was sie vorhaben, was sie lehren und thun, wird dermaßen von Jedermann gelobt, wird mit solcher Begierde angenommen, als ob es von Gott selbst also beschlossen wäre. „Du züchtigest sie nicht“, sagt der Prophet,“ das ist, was sie Unbilliges und Ungöttliches nach ihren bösen Lüsten vorhaben, dem widerstehest du nicht, hältst sie nicht im Zaum. Aber was ist der endliche Ausgang? Das folgt daselbst im Propheten V. 3.: „Aber du lässest sie frei gehen, wie Schafe, daß sie geschlachtet werden; und sparst sie, daß sie gemürget werden.“

66. Auf solche Weise haben die Söhne Jakobs auch gemüthet, haben ihrem Haß und Neid immer nachgehangen: aber sie haben endlich wiederum auch müssen erfahren, daß sie nicht gern gehabt.

67. Derhalben sage ich, daß Gott nichts so angenehm, und so lieb und wohlgefällig sein kann, als daß man mit Geduld ertrage, wo man vom Herrn gezüchtigt und väterlich gestraft wird; wenn wir aber das nicht ertragen wollen, daß wir im Zaum gehalten werden, werden wir doch endlich in viel greuliche Strafe geraten, darin wir dann auch müssen stecken bleiben. Denn es heißt: Entlauffst du mir, so entlauffst du, aber dem Henker oder dem Teufel nicht. Es kommen zwar bisweilen etliche wiederum auf den rechten Weg und werden bekehrt, aber solches geschieht nicht ohne große greuliche Strafe, damit sie zur Buße gezogen werden.

68. So scheint es nun, daß Simeon, welcher ein Vater des Juda Ischarioth gewesen, und Levi die Väterführer und vornehmsten Stifter dieser schweren Sünde gewesen sein, so die Söhne Jakobs an Joseph, ihrem Bruder, begangen, und daß sie die andern auch, so von den Mägdlein geboren, mit sich an den Reihen geführt, die sonst, wo sie nicht also angeführt, unimmermehr sich des hätten unterstehen dürfen. Und ist dieser böse Rath und Vornehmen auch mit nichten gehindert worden, sondern ihnen glücklich, wie sie nur selbst begehrt, angegangen, dieweil Gott verhängt und zugelassen, daß sie ihren Bruder gefangen und verkauft haben. Ja, Joseph selbst auch hat dies Feuer helfen abblasen und setzt nun die dritte Ursache des Hasses seiner Brüder auch noch hinzu.

\section[Zweiter Theil]{Zweiter Theil. Von den Träumen Josephs}

\subsection{I}

\begin{quote}
\textbf{V. 5–7. Dazu hatte Joseph einmal einen Traum, und sagte seinen Brüdern davon; da wurden sie ihm noch feinder. Denn er sprach zu ihnen: „Höret, lieber, was mir doch geträumet hat. Mir däuchte, wir banden Garben auf dem Felde, und meine Garbe richtete sich auf, und stund; und eure Garben umher neigeten sich gegen meiner Garbe.}
\end{quote}

69. Dies ist nun die rechte Hauptursache der Feindschaft gewesen, und scheint es, als habe Joseph mit vorbedachtem Muth seinen Brüdern diese Ursache also entdeckt und geoffenbart, daß er dieselbige ihrem Stolz und Neid stracks entgegensetzte. Denn Simeon und Levi haben sich selbst vornehmlich das Regiment und die Herrschaft angemaßt, darum daß sie noch damals unschuldig waren, da Ruben schon etwas gedemüthigt war; darum auch ein Jeder unter ihnen wird gedacht haben: Ich bin Herr über meine Brüder, derhalben will ich mir auch die Gerechtigkeit und Herrlichkeit der Erstgeburt anmaßen, und sollen meine Brüder sich vor mir neigen und sich vor mir bücken meiner Mutter Söhne; wie wir droben, Cap. 27. V. 29., in dem Segen Jacobs gehört, und hernach auch von Juda, Cap. 49. V. 8., wird gesagt werden: „Juda, du bist, dich werden deine Brüder loben. Deine Hand wird deinen Feinden auf dem Halse sein; vor dir werden deines Vaters Kinder sich neigen.“

70. Solche Herrschaft, sage ich, hat sich Simeon angemaßt, hat ein stark Vertrauen und gute Hoffnung, er werde sie auch gewißlich bekommen. Darum auch Joseph diesen Traum solcher tollen Vermessenheit und Stolz Simeons entgegengesetzt. Als wollte er sagen: Ich kann nicht leiden, daß ihr euch länger also herumführet und schleppet mit ungewißer und vergeblicher Hoffnung der künftigen Herrschaft, als ob dieselbige auf euch kommen würde; denn ich, so von euch gar verachtet und verstoßen bin, soll endlich euer Herr werden: nicht daß ich mir selbst solche Herrschaft über euch zuschreibe und mir dieselbe anmaße, sondern ich erzähle euch meinen Traum, den ich gehabt.

\quellseite{275}

So gar freundlich redet er mit seinen Brüdern, als ob er ein halber Narr, oder gleichsam ein Kind wäre: „Höret, lieber, was mir doch geträumet hat“ etc. „Mir däuchte, wir banden Garben“ etc. So einen wunderlichen Traum, daß sich niemand hätte versehen können, hat der Knabe erzählt, und dasselbe ganz kindlich, aus keinem Neid oder böser Lust, sondern aus lauter Einfalt und Unschuld. Aber GOtt hat ihre Herzen beschweret, und die Ursache des Hasses also gemehret, daß sie je länger je mehr erbittert und entbrannt worden sind; aber zu ihrem eigenen großen Schaden. Da geht nun das Feuer recht an.

\begin{quote}
\textbf{V. 8. Da sprachen seine Brüder zu ihm: Solltest du unser König werden, und über uns herrschen? Und wurden ihm noch feinder um seines Traums und seiner Rede willen.}
\end{quote}

71. Daß Joseph seinen Traum also erzählt, hat diese großen Juncker sehr verdrossen. Er sollte damit stille geschwiegen, und seinen Brüdern hofiret und geheuchelt haben, und gesagt: Simeon und Levi, die sollen Herren sein und die Herrschaft haben; daß er damit also ihre guten und heiligen Begierden bestätiget hätte. Aber das ist unseres HErrn GOttes gewaltige Hand, der dem Simeon und Levi ein Specklein auf die Falle legt. Denn dieweil sie dem unschuldigen Joseph also heimlich nachstellen, wird ihnen GOtt auch wiederum feind, und verhängt, daß sie hart anlaufen, dazu sie denn ihr eigener Wille getrieben, und daß sie ihres eigenen Bruders Mörder geworden sind.

\begin{quote}
\textbf{V. 9. Und er hatte noch einen andern Traum, den erzählte er seinen Brüdern, und sprach: Siehe, ich habe noch einen Traum gehabt; mich däuchte, die Sonne, und der Mond, und elf Sterne neigten sich vor mir.}
\end{quote}

72. Das Wort „anbeten“ bedeutet eigentlich die Geberde des Anbetens oder der Ehrerbietung und daß man sich an einen andern Ort beweget. Denn „anbeten“ ist und heißt, sich auf die Erde begeben oder Einem zu Fuße fallen, und mit solcher äußerlichen Geberde einem Andern Ehre erzeigen. Wir haben den Gebrauch, daß wir die Kniee beugen, es sei gleich vor den Fürsten oder auch vor GOtt.

73. Wir müssen aber nun auch an diesem Orte etwas insgemein von den Träumen sagen. Und dieweil ich sehr wohl weiß, daß diese Materie sehr reich und weitläufig ist, will ich mir derhalben kürzlich davon helfen. Denn ich bin nicht geschickt Träume zu haben oder auch dieselbigen zu deuten, begehre auch solche Geschicklichkeit oder Kunst nicht, und habe mit GOtt meinem HErrn einen Bund gemacht, daß er mir nur keine Gesichte oder Träume und auch dazu keine Engel senden wolle. Denn ich bin mit dieser Gnade wohl zufrieden und lasse mir daran gern genügen, daß ich die heilige Schrift habe, die mich reichlich lehrt und berichtet alles deß, so beide zu diesem und auch zum künftigen Leben zu wissen vonnöthen ist. Dieser heiligen Schrift glaube ich und bin damit zufrieden, bin auch gewiß, daß ich dabei nicht kann betrogen werden; will aber damit gleichwohl Andern an ihren Gaben nichts abbrechen, da vielleicht Einem außerhalb der Schrift GOtt etwas durch Träume, Gesichte, oder aber durch die Engel möchte offenbaren. Ich gebe zu, daß es Gaben sind; aber ich achte oder begehre derselbigen nicht. Denn dazu bewegt mich, daß so viel unzählige Gespenster, Lügen, Verblendung des Gesichts und andere Betrug gewesen, damit die Welt erschrecklicher Weise im Papstthum eine lange Zeit durch den Teufel ist verführt worden. Zudem habe ich auch noch diese Ursache, daß die heilige Schrift allein genugsam ist; und wo ich derselben nicht glauben würde, so würde ich wahrlich auch nicht leichtlich weder Engeln, noch Gesichten, noch einem Traum glauben. Aber wie gesagt, dies ist meine eigene Ursache, will aber Andern kein Ziel oder Maß vorgeschrieben, noch auch jemand damit etwas genommen haben. Denn ich habe übrig genug Offenbarung, dieweil ich weiß, was ich glauben, hoffen, oder auch Andern rathen soll, wie man sich auch dazu schicken soll, daß man dies Leben christlich und ehrbarlich zubringen möge. Und dasselbe will ich nach meinem geringen Vermögen Andern herzlich gern mittheilen, daß sie selbst für sich dies auch also wissen mögen und die Zehen Gebote, Glauben und Vater Unser verstehen lernen.

74. Wir wollen aber doch gleichwohl auch noch etwas sagen, was die Lehre und Meinung sei der heiligen Schrift von den Träumen; denn sie rühmt an etlichen Orten die Träume, aber an andern Orten schilt und verwirft sie sie auch wiederum. Sirach spricht am 34. Cap. V. 7.: „Träume betrügen viel Leute, und fehlt denen, die darauf bauen.“ Item, Salomo in seinem Prediger Cap. 5, V. 6. sagt auch: „Wo viel Träume sind, da ist Eitelkeit und viel Worte; aber fürchte du GOtt.“

75. Wiederum aber zeugen auch die Historien in der heiligen Schrift, daß etliche Patriarchen und Propheten auch Träume gehabt, als, Jacob, Joseph, Daniel, desgleichen Pharao und andere Gottlose auch, welche ihnen doch von GOtt gezeigt sind. Und 4 Mos. 12, 6. 7. 8. spricht GOtt zu Aaron: „Ist jemand unter euch ein Prophet des HErrn, dem will ich mich kund machen in einem Gesicht, oder will mit ihm reden in einem Traume. Aber nicht also mein Knecht Mose, der in meinem ganzen Hause treu ist; mündlich rede ich mit ihm.“ Da bestätigt GOtt die Träume und Gesichte, und ist eine schöne treffliche Stelle, so wohl einer fleißigen und weitläufigen Auslegung bedürfte. Mose, spricht GOtt, hat etwas Größeres gesehen, hat noch andere und größere Offenbarungen ohne die, so ihr von ihm gehört, die ihr seine Jünger seid, und selbst auch Gesichte und Träume habt: ich weiß wohl, was Mose und ich mit einander pflegen zu handeln und worin wir uns unterreden. Mose hat auch das Leiden Christi gesehen.

76. Weiter sind nun aus diesem Orte hergekommen die gemeinen Weisen zu reden in den Propheten von den Träumen und Gesichten. Also hat Joel die Sendung des Heiligen Geistes verkündigt Joel 3, 1.: „Eure Jünglinge sollen Gesichte sehen, und eure Aeltesten sollen Träume haben.“ Derhalben können wir nicht leugnen, daß GOtt solche Offenbarungen vom Anfange in der Kirche der alten Väter und Propheten gebraucht habe. Und ist aber die erste Weise der Offenbarung die allerhöchste und vornehmste, nämlich, die Weissagung oder Eingebung, wie Daniels und der andern Propheten gewesen, so durch Erleuchtung des Heiligen Geistes die Lehre und göttlichen Verheißungen mit hellen klaren Worten und mit deutlichem Verstande dargegeben haben. Darnach sind dann die Gesichte, oder etliche Bilder und äußerliche Erscheinungen. Zum Dritten, die Träume. Also ist im Neuen Testamente auch das Erste die Lehre oder Predigt. Dazu kommen dann die heiligen Sacramente als äußerliche Ceremonien, welche auch wohl möchten Gesichte genannt werden, als, unsere Taufe, das Nachtmahl des HErrn, die Absolution; das sind, sage ich, auch etliche Gesichte, denn es sind äußerliche Ceremonien.

77. Dieweil aber nun die heilige Schrift diese dreierlei Weisen sowohl lobt, als verwirft, däucht mich, daß man derselben Regel und dem Urtheil auch also folgen müsse, wenn man reden will von den Träumen, nach welchem wir uns zu richten pflegen in den Gesichten und Weissagungen, nämlich, daß man darauf sehe, daß sie ja vornehmlich dem Glauben ähnlich sein mögen.

78. Denn es verhält sich mit den Gesichten und Weissagungen auch also, wie mit den Träumen, daß sie bisweilen wahrhaftig, und bisweilen auch wohl falsch und lügenhaftig sind. Denn sie kommen auch nicht allezeit oder allein von GOtt, sondern auch vom Teufel, der gleichsam unsers HErrn GOttes Affe ist; und wie er auch pflegt Propheten zu erwecken, also erregt er und gibt auch wohl Gesichte und Träume, so es ihm GOtt also aus seinem wunderbaren Rath verhänget. Jedoch sieht der Teufel anderswohin, und meint damit viel ein Anderes denn GOtt, wenn er auch schon die Wahrheit redet durch Träume. Denn GOtt gibt sein Wort und die Zeichen dabei zur Seligkeit und Erlösung der Menschen; der Teufel aber sucht gewißlich nichts Anderes, denn eitel Schaden und Verderben der Seelen, dazu auch, daß die göttliche Wahrheit verdunkelt werden möge; darum lügt er, auch wo er schon die Wahrheit redet. Daher kommt es, daß auch in den heidnischen Historien etliche Träume erzählt werden; wie denn die Träume des Brutus, Cäsars, Sullas und Hannibals gewesen sind.

79. Es schreibt Livius, daß Hannibal im Traume einen Jüngling in göttlicher Gestalt gesehen habe, der da sagte: Er wäre von Jupiter nach Italien zu Hannibal gesandt, daß er sein Heerführer sein sollte, dem er stracks nachfolgen und den er nimmer aus den Augen lassen sollte; und da sich Hannibal umgesehen, soll er hinter sich eine wunderbare große Schlange gesehen haben, die sich hin und her bewegt und dadurch eine gewaltige und große Niederlage unter den Bäumen und Sträuchern angerichtet habe. Das ist eigentlich und gewißlich ein weissagender Traum und der auf ihn selbst auch einen Eindruck gemacht. Denn der Teufel hat sehr wohl gewußt, wie es damals in der ganzen Welt gestanden, und ist leicht gewesen aus dem, so vorher gegangen und geschehen war, zu schließen und abzunehmen, wie es auch hernach gehen würde, und solches dem Hannibal im Traum anzuzeigen. Desjenigen Schlages sind auch des Lichtenbergers Weissagung und auch des Arnolds De villa Dei, welche den Geist der Weissagung gehabt, so man sonst den Wahrsagergeist nennt, und trifft bisweilen sein zu, daß es also kommt, wie sie zuvor gesagt, bisweilen auch wohl nicht: Trifft's, so trifft's; fehlts, so fehlts.

\quellseite{276}

80. Weiter, im Traume des Hannibal hält Livius selbst und alle Italiener dafür, daß die Schlange Hannibal soll bedeutet haben, der Italien verwüsten würde. Aber gleichwohl ist diese Deutung der Wahrheit näher, daß der Teufel oder aber Gott durch den Teufel hat damit anzeigen wollen, daß Hannibal einen schönen Jüngling vor sich hätte, das ist, daß er würde gut Glück haben. Als sollte er sagen: Hannibal, du bist wohl bis daher glücklich gewesen und hast große Dinge ausgerichtet; aber dies Glück wird nicht allezeit währen. Denn es wird dir endlich nachziehen und folgen ein erschrecklich Thier, nämlich, das Römische Reich, welches in einander wickeln und verwüsten wird den ganzen Erdboden. Also hat dieser Traum, der vom Satan als einem Lügner hergekommen, gleichwohl die Wahrheit zuvor verkündigt; er aber hat es nicht verstanden, und die Italiener, die sehr ehrgeizig sind, hören nicht gern, daß durch ein solch schändlich und greulich Thier das Römische Reich soll verstanden werden.

81. Wir wollen aber auf die Bilder in den Geschichten Daniels sehen, da die Königreiche dieser Welt nicht den Menschen oder Engeln, sondern großen greulichen Thieren verglichen werden. Derhalben bedeutet diese Schlange mit nichten den Hannibal, sondern gewißlich die zukünftige Zerstörung, so durch das Thier, das ist, das Römische Reich, geschehen sollte. Und hat der Teufel also zu Hannibal wollen sagen: Dein Glück und Gedeihen wird keinen Bestand haben: du wirst das nicht vollenden, was du angefangen hast; du wirst es nicht hinausführen.

82. Und zu solcher Offenbarung ist des Heiligen Geistes nicht vonnöthen gewesen. Denn der Teufel hat wohl gesehen die Macht, den Rath und Willen, auch weß man sich beflissen, und wie es in allen beiden Regimenten, der Römer und auch der Carthager, gestanden; und Gott hat verhängt, daß solches dem Hannibal im Traume geoffenbart wurde, aber doch ohne Deutung.

83. Solcher Träume findet man viele in den Historien der Heiden, so beide gewiß und ungewiß sind. Denn wie Gott durch gute und böse Fürsten oder Regenten die Welt regiert: also braucht er auch zugleich des Dienstes beide guter und böser Engel. Derhalben er auch den bösen gottlosen Leuten bisweilen wahre, bisweilen auch falsche Träume eingibt; und der Teufel hat auch seine Propheten, welcher Weissagungen aber doch gleichwohl ungewiß sind.

84. Die andere Art der Träume ist natürlich, daraus die Aerzte bisweilen abnehmen und errathen etliche gewisse Feuchtigkeiten, Bewegung des Gemüths; desgleichen, wie die Leiber der Menschen beschaffen sind. Dieselbigen Träume sind gemein und jedermann kann sie wohl haben. Und bei solchen Träumen soll man sich halten nach der Lehre Catonis, der da sagt: Somnia ne cures etc., das ist: Du sollst der Träume nicht achten und darauf nichts geben; wiewohl auch derselbigen etliche bisweilen ihre Bedeutung haben. Aber von denen will ich nicht disputiren, sondern will den Leser gewiesen haben, daß er davon im Macrobius lese, welcher davon in die Länge geschrieben.

85. Die dritte Art von Träumen sind nun die prophetischen und wahren Träume; und die werden in der heiligen Schrift gerühmt, und gehören auch zu heilsamer Regierung in der Kirche, oder zu den Offenbarungen, so beide zu diesem und dem zukünftigen Leben vonnöthen, und sind gleichsam Prophezeiungen, oder Weissagungen, welche mit nichten zu verachten sind.

86. Aber da wird nun gefragt: Wie doch dieselbigen zu unterscheiden und davon zu urtheilen sein möge? Soviel mich belanget, bin ich zwar kein Träumer, kann auch nicht wohl errathen, was die Träume bedeuten mögen, wie zuvor gesagt; darum will ich davon auch nicht urtheilen, denn allein, darnach die Träume erfüllt und vollzogen werden, und nach andern Umständen, so die Kirche und der Seelen Seligkeit belangen, nach der Regel und Anweisung des göttlichen Worts. Denn alle Träume, so von Gott eingegeben, werden auch gewißlich erfüllt, und das also, daß sie auch der Welt und menschlichem Geschlechte nütze und heilsam sein müssen: werden derhalben auch billig und recht geschätzt oder beurtheilt ab effectu, das ist, nachdem sie einen Ausgang gewinnen oder erfüllt werden.

87. Und das sage ich nicht allein von den Träumen, so den Heiligen, als, Daniel und Joseph, eingegeben, sondern auch von denen, so den Gottlosen gezeigt werden, als, Pharao, Nebucadnezar und Andern. Denn dieselbigen Träume gewinnen auch ihren Ausgang und haben auch ihre angehängte Auslegung gehabt; welches ja zwei wahrhaftige Zeichen sind, daß die Träume auch gewiß sind, nämlich, daß wo man Gott die Träume eingibt, er auch die Auslegung und Erfüllung zugleich mitgibt.

88. Zum Dritten ist es um dieselbigen Träume auch also beschaffen, daß sie dem Träumer Kraft angethan machen und bekümmern, prägen ihm die Dinge, davon er Träume gehabt, heftig und hart ein, wie Pharaons Traum gewesen. Denn denselben ist ein solch Schrecken und Zagen im Herzen angekommen, daß er gleichsam darüber entzücht worden ist, daß er gefühlt und verstanden hat, es müßte gewißlich eine Ermahnung sein, so von Gott gekommen. Da hat auch Gott nicht verhängen wollen, daß ein gottloser falscher Prophet den Traum, so doch einem Gottlosen gezeigt, auslegen oder deuten sollte, sondern gibt einen rechtschaffenen wahren Propheten, als, Joseph oder Daniel, welcher Bericht geben muß durch Erleuchtung des Heiligen Geistes, was eigentlich die Bedeutung des Traumes sei.

89. Wo aber nun der Heilige Geist selbst das auslegt und deutet, was er eingegeben, alsdann darf man gar nicht zweifeln, daß solches nicht auch seinen gewissen Ausgang gewinnen werde. Dergleichen wirst du aber in den Träumen der Heiden nicht finden, ausgenommen, daß sie gleichwohl davon etwas heftig bewegt und bestürzt werden. Aber dieweil ihre Träume vom Teufel hergekommen, der auch wohl von zukünftigen Dingen kann und pflegt Erinnerung zu thun, darum sind sie auch gemeiniglich ungewiß und falsch, wird auch dabei keine Deutung gezeigt, wie in göttlichen Träumen geschieht, welche die Propheten gewißlich und treulich erzählen und erklären, wie sie das der Heilige Geist gelehrt. Denn also haben Daniel, Jakob und Joseph aus den Träumen, so ihnen gezeigt, geschlossen und verstanden die rechte Lehre und die heilige Schrift, da sie von großen Dingen redet, nicht allein von den Regimenten, davon sie lehren, daß sie von Gott selbst geordnet und erhalten werden, sondern auch vom Reiche Christi und den vornehmsten Artikeln der christlichen Lehre. Darum liegt es an einem guten Ausleger, wie das deutsche Sprüchwort lautet.

90. So ist nun das, sage ich, des Heiligen Geistes Gebrauch, daß er erstlich pflegt die Träume einzugeben, darnach auch die Deutung anzuzeigen, und zum Dritten, dieselbigen dann auch zu erfüllen. Und solches kann auch sein auf die heilige Dreifaltigkeit gezogen werden, nämlich, daß Gott der Vater die Träume gibt, Gott der Sohn dieselbigen auslegt und erkläret, der Heilige Geist aber sie erfüllt.

91. Und dies heiße ich wahre und prophetische Träume, welche nicht fehlen. Aber von denselbigen, sage ich, will ich mich nicht unterstehen zu urtheilen, denn allein, nachdem sie einen Ausgang gewinnen und erfüllt werden. Die andern Träume aber, deren viel sind, verachte ich ganz und gar. Ich habe ja auch bisweilen wohl etliche Träume, so mich etwas bewegen; ich verachte sie aber doch gleichwohl, und habe mit dem Herrn, meinem Gott, diesen Vertrag gemacht, daß ich Mose und den Propheten glauben wolle. Denn zu diesem Leben begehre ich mir keiner Träume und sind mir dieselbigen auch zum künftigen Leben gar nicht vonnöthen.

92. Ich will aber gleichwohl indessen Andern ihre Gaben auch gern lassen; denn ich weiß wohl, daß Gott oftmals Königen und Fürsten mancherlei Offenbarung durch Träume gezeigt hat. Die aber nun solche Träume auslegen wollen, sollen wohl zusehen, was für ein Geist sei, so sie eingegeben, und woher sie gekommen seien.

93. Die Fürsten haben auch beiderlei, nämlich, Träume und Weissagungen, wie es gehen und stehen, und was es für ein Ende nehmen soll mit dem türkischen Reiche; aber ich kann nicht glauben, daß solche vom Heiligen Geist herkommen. Denn es haben weder Türke noch Papst rechte Propheten, sonst müßten sie geoffenbart werden durch das Bekenntniß rechter Lehre und Glaubens. Denn sei aber nun wie ihm wolle, wir achten solcher Träume nicht, achten auch dessen nicht, von wem sie herkommen, dieweil sie keine angehängte Deutung haben, wiewohl sie bisweilen auch ihren Ausgang gewinnen und erfüllt werden. Denn der Teufel ist sehr wohl geübt und gerüstet mit vielen unzähligen Krempeln göttlicher Regierung; desgleichen weiß er auch, was der Menschen Wille sei und was sie sich befleißigen; weiß auch ihren Rath und Anschläge wohl und regiert auch dieselbigen: derhalben kann er auch leichtlich ersehen, was die Rathschläge und Regimente für Glück und für einen Ausgang gewinnen werden. Er sieht, daß was die Leute bedacht und beschlossen, bisher so oder so hinausgegangen und gerathen ist, und daß dermalen noch keinerlei Anzeichen vorhanden, daß GOtt irgend ein Hinderniß wollte eintreten lassen, dadurch der vorgenommene Rathschlag geändert oder gehindert werden möchte. Darum ist es ihm leicht, daß er von zukünftigen Dingen weissage.

\quellseite{277}

94. Also hat der Lichtenberger den bayrischen Krieg zuvor verkündigt. Denn der Teufel hat wohl gesehen, was dazumal des Kaisers Maximilian, des Pfalzgrafen und anderer Fürsten Rathschläge gewesen. Dasselbige thut er auch zu andern Zeiten, und wo sich GOtt nicht aus sonderlichem Rath dagegen setzt, so geschieht auch wohl und wird erfüllt, was also zuvor verkündigt ist. Darum aber kümmert sich der Teufel nicht, ob ein solcher Traum auch den Menschen zu ihren Sachen nütze und heilsam sei oder nicht; ja, wenn gleich Alles zu Scheitern geht, da fragt er nichts nach.

95. Uebrigens aber ist dies ganze Leben, darin wir leben, nur ein eitler Traum und wie eine dunkele Nacht. Denn es leuchtet sonst in der ganzen Welt kein Licht, ausgenommen dies irdische und leibliche Licht. Das Licht der Sonne ist uns gleichsam eine Ursache des Irrthums, daß wir meinen, wir seien im Licht; aber in Wahrheit gehen wir in der Irre wie in einer finstern Nacht. Derhalben, sage ich, ist dies Leben wahrhaftig ein Schlaf und geistlich nichts Anderes denn eine Nacht; wie im Traume Josephs seine Eltern der Sonne und Mond, seine Brüder aber den Sternen verglichen werden, so in der Finsterniß dieses Lebens scheinen. Und dasselbe bezeuget Mose auch, da er Ps. 90, 5. sagt: „Du lässest sie dahin fahren wie einen Strom, und sind wie ein Schlaf“, das ist, unser Leben ist kurz, flieht und geht bald dahin, wie die Elbe oder sonst ein jegliches großes Wasser immer dahin fließt, wie ein Wasser vorüber rauscht. Also, sagt Mose, hast du sie auch gemacht, wie ein fließend oder rauschend Wasser. Also sagt auch Hiob Cap. 14, V. 1, 2.: „Der Mensch vom Weibe geboren lebt kurze Zeit, und ist voll Unruhe, gehet auf wie eine Blume, und fällt ab, fleucht wie ein Schatten, und bleibet nicht“ 2c. Und das noch ärger ist, setzt Mose hinzu: „Die Menschen sind wie ein Schlaf“, das ist, unser Leben hier auf Erden ist gegen das ewige Leben nichts Anderes denn „wie ein Schlaf“.

96. Derhalben hangen wir und stecken in diesem Leben wie mitten in der Nacht und der allerdicksten Finsterniß, wo uns der rechte Morgenstern nicht aufgeht und scheint, davon der 119. Psalm V. 105. sagt: „Herr, dein Wort ist meines Fußes Leuchte, und ein Licht auf meinem Wege.“ Denn wo wir die heilige Schrift nicht haben, daß sie uns vorleuchte und alle unsere Werke regiere, so ist dies ganze Leben und alles Licht der Vernunft, ja, dazu auch alle Weisheit und aller Menschen Rath nichts Anderes, denn eitel Finsterniß und eine lautere Unordnung.

97. Augustus, Alexander und Julius Cäsar stecken wie in einem dicken Nebel, wissen selbst nicht, was sie thun. Also auch Nebucadnezar, Pharao und alle anderen Könige und Monarchen; allein, daß sie sich gleichwohl dünken lassen, daß sie sehr viel ausrichten und daß sie auch überaus weise seien. Aber es sind eitel Träume, wie im Jesaia Cap. 29, V. 7. steht: „Wie ein Nachtgesicht im Traum, so soll sein die Menge aller Heiden, so wider Ariel streiten.“ Also auch die Geschichten aller großen Helden, so von den Poeten und sonst auch in den Historien der Heiden gerühmt werden, als, des Hector und Achilles, sind wie Träume. Denn gleichwie einem Hungrigen träumet, daß er esse, wenn er aber aufwacht, so ist seine Seele noch leer; und wie einem Durstigen träumet, daß er trinke, wenn er aber aufwacht, so ist er matt und durstig: also sollen sein die Menge aller Heiden, die wider den Berg Zion streiten“ Jes. 29, 8. Denn das kommt oft vor, daß Einer im Traum sich dünken läßt, daß er einen Sack voll Geldes finde oder mit einer schönen Jungfrau spiele: aber wenn er wieder erwacht, flieht und findet er, daß er betrogen ist. Eben auf solche Weise ist auch dies ganze Leben nichts Anderes, denn wie eine Nacht und Schlaf.

98. Das ist aber ein seliger Traum, wenn GOtt in solcher Nacht sein Wort und die heilige Schrift gibt, wie er in diesem natürlichen Leben oder im Schlaf eine sonderliche Offenbarung gibt, wie diese ist, so Joseph im Traum gehabt. Dazu aber, wie ich zuvor erwähnt, ist vonnöthen, daß auch die Deutung des heiligen Geistes angehängt werde und die Erfüllung auch mit folge. Denn so viel Verstand habe ich von den Träumen, nämlich, darnach oder sofern der Ausgang folgt; mehr und weiter gar nichts. Derhalben kann ich auch darüber nicht urtheilen, wann den Träumen zu glauben sei, es sei denn, daß sie GOtt selbst deute oder auslege.

99. Darnach ist das auch zu merken, daß der unsaubere Geist auch Träume zeugt. Denn aus vielfältiger Erkenntniß und mancherlei Erfahrung ist er viel gerüsteter und viel besser geschickt, denn wir sind. Wir wissen nicht, was die Fürsten und andere Menschen für Rathschläge haben, womit sie umgehen und was ihr Vorhaben sei; ja, ich kenne kaum meine eigenen Rathschläge. Aber der Teufel ist sehr listig, weiß und versteht sehr genau alle Rathschläge der Könige, der Weltweisen, Juristen und auch der Theologen, das allein ausgenommen, daß er nicht wissen kann, wie mein Glaube und Hoffnung zu GOtt stehe, wie ich mit GOtt daran bin. Denn der Glaube und GOttes Wort ist ihm wie eine Finsterniß und eine Wolke, dazu er mit seinem Licht nicht kommen kann. Wie er denn das selbst anzeigt durch Krystalle und Spiegel, so specula mercurialia genannt werden, darin er irgend eine Stadt oder eine sonderliche Person zeigt, aber die gleichwohl etwas dunkel und wie mit einem finstern Nebel oder Regen bedeckt ist; daher es auch kommt, daß nicht so klarlich kann erkannt oder gesehen werden, was daselbst geschieht. Also spottet er aber der Menschen, und bekennt doch gleichwohl damit auch, daß er dem Menschen nicht könne in das Herz sehen; da verbrennt er sich Maul und Nase.

100. Eben auf dieselbe Meinung sagt auch der Engel zu Maria Luc. 1, 31. 35.: „Der Heilige Geist wird über dich kommen, und die Kraft des Höchsten wird dich überschatten. Du wirst schwanger werden im Leibe, und einen Sohn gebären.“ Aber auf solche Weise, die gar ein verborgenes Ding und solche Ueberschattung sein wird, daß auch weder Engel noch Teufel etwas davon sehen oder verstehen können. Derhalben sage ich, daß der Teufel mit der Schärfe seines Gesichts in das Wort und ein gläubiges Herz mit nichten sehen kann; und um derselben Ursache willen habe ich mehr Lust zum Wort und Glauben denn zu Träumen, und habe auch den Herrn gebeten, er wolle mir keine Träume geben; denn sie fehlen doch gemeiniglich. Das Wort aber ist ein gewisser Schirm, und ist eine solche Dunkelheit, darein die bösen Geister nicht sehen können, wenn sie auch noch so hoch wären.

101. Derhalben, wiewohl die Weissager und Träumer nicht zu verachten sind, wo ihre Weissagung und Träume dem Worte GOttes und Glauben ähnlich sind: so soll man aber gleichwohl darauf vornehmlich bringen, daß man ja zuvor Mose und die Propheten wohl studire und sie auch verstehen lerne, das ist, daß wir vor allen Dingen das Wort rein haben und dasselbe auch recht verstehen mögen. Daraus werden wir dann leichtlich alle Gesichte, alle Träume und Weissagung, ja auch beider gute und böse Engel beurtheilen können.

102. Wir wollen aber nun wieder auf die Träume Josephs kommen. Der Träume, so er gehabt, sind zwei, unter welchen der erste mit dem Landleben und Ackerbau zusammenhängt. Denn es träumt ihm daselbst nur von seiner Garbe und von seiner Brüder Garben. Seine Garbe richtete sich auf, und stand; gegen dieselbe haben sich die andern geneigt. Dieser Traum hat Joseph nicht sonderlich angefochten, darum hat er auch nichts dagegen geredet, sondern schweiget dazu gar stille. Denn er hat sich auch gereimt mit dem, so er selbst auch mit Joseph vorgehabt, nämlich, daß er beschlossen hatte, Joseph sollte der vornehmste Erbe in seinem Hause werden, dieweil er auch von seinem vornehmsten Weibe geboren war, welchem derhalben auch die Erstgeburt gebührte, dazu das Regiment und das Priesterthum, und zwei Theile des Erbes. Derhalben stimmt nochmals dieser Traum Josephs sein mit den Gedanken Jacobs, seines Vaters, und bestätigt sie auch, und ist der Anfang und gleichsam ein Vorläufer auf den rechten Traum. Denn es wird hier der Eltern noch nicht gedacht, daß sich die sollten vor ihm neigen, sondern nur der Brüder wird gedacht. Nun mag es aber von Natur und auch im Gesetze also geordnet und geboten, daß der erstgeborne Sohn über seine Brüder sollte Herr sein; wie das andere frühere Exempel auch bezeugen.

\quellseite{278}

103. Nun hatten aber Simeon und Levi sich schon vorlängst das Recht der Erstgeburt zugeeignet; derhalben sie auch also von Neid und Haß gleichsam entbrannt worden wider Joseph, wiewohl sie dazu kein ander Recht gehabt, denn allein die Zeit, daß sie eher denn Joseph in diese Welt geboren waren von ihrer Mutter Lea, welche ihrem Vater Jakob zu allererst vermählet, damit auch der arme Jakob unbilliger Weise von seinem Schwiegervater Laban betrogen worden. Denn es soll die Erstgeburt nicht nach der Zeit gerechnet oder geschätzt werden. Und nichts desto weniger haben gleichwohl diese Zwei die Herrlichkeit der Erstgeburt sich selbst zugeeignet und sich dieselbige angemaßt. So hat Levi darauf gehofft, daß er dazu wollte der Nächste sein, und Simeon mit dem Tode abgehen würde. Und hatten also das andere Hausgesinde auch auf ihre Seite gezogen; das hat auch müssen auf sie sehen und sich nach ihnen richten. Darum stimmen die andern Brüder auch mit ihnen überein und setzen sich einmüthig wider Joseph; wie sich die Apostel Matth. 26, 8. von dem einigen Juda haben überreden lassen, daß sie wider Magdalena auch gemurrt und ungeduldig über sie geworden sind. Diese Zwei aber, Simeon und Levi, sind sonderlich mehr denn die Andern heftig erzürnt worden.

104. Zudem hat gleichwohl dieser Traum auch ein fast starkes Nachdenken in ihren Herzen gemacht, und kann nicht gezählet oder gerechnet werden unter gemeine Träume, die nichts Sonderliches werth sind. Darnach ist der Vater darüber nicht ungeduldig oder erzürnt worden, sondern hat ihn mit seinem Urtheil vielmehr im Herzen bestätigt und aus väterlicher Gewalt sich denselben gefallen lassen. Und kommt also dazu auch des Vaters Deutung selbst, welche göttlich ist, und ein gewisses Anzeichen, daß der Traum auch also werde müssen erfüllt werden.

105. Der andere Traum aber Josephs ist nicht so gemein oder gering, sondern ist etwas hoch und gleichsam himmlisch. Denn er hat seine Bilder oder Gleichnisse von der Sonne, vom Mond und von den Sternen; sagt von seinen Garben, welche bedeuten die harte Arbeit derer, so mit dem Ackerbau umgehen, und sonst andere schwere und mühselige Werke der Menschen. Und wird nun in demselbigen Traume auch Vater und Mutter zugleich mit den Brüdern begriffen, da im vorigen Traume ihrer nicht ist gedacht worden. Darum wird nun dieses Traumes halben der Vater selbst auch etwas bewegt oder erzürnt und straft ihn auch darüber, wie nun im Texte weiter folgt.

\begin{quote}
\textbf{V. 10. 11. Und da das seinem Vater und seinen Brüdern gesagt ward, strafte ihn sein Vater, und sprach zu ihm: Was ist das für ein Traum, der dir geträumet hat? Soll ich, und deine Mutter, und deine Brüder kommen, und dich anbeten? Und seine Brüder neideten ihn. Aber sein Vater behielt diese Worte.}
\end{quote}

106. Droben, 4. u. 5. ist erzählt worden, wie die Brüder Josephs ihn gehaßt und geneidet haben. Mithin aber werden sie nun gar unsinnig. Denn das kann weder die Natur, noch der Gebrauch oder Weise einigerlei Völker, wer sie auch sein mögen, ja, Gott selbst nicht vertragen oder dulden, daß sich die Eltern vor ihren Kindern neigen sollten. Darum ist dieser Traum gar abscheulich zu hören und ist wider alle Vernunft und gemeinen Verstand aller Menschen. Er macht aber doch einen gewaltigen Eindruck; denn es werden deutlich namhaft gemacht die Sonne, der Mond und elf Sterne, welches genau ist die Anzahl des Vaters, der Mutter und Brüder, und kommt gar keine andere Zahl darin vor.

107. Lyra sagt, daß Jakob diesen Traum verachtet und ihn gar soll verworfen haben, als der offenbar ungereimt sei; und daß er dazu aus väterlicher Gewalt die Thorheit dieses Jünglings gestraft, welchem aus kindischer Leichtfertigkeit geträumt, daß sich beide Eltern und Brüder vor ihm neigen sollten. Aber im Texte wird das Gegentheil angezeigt, nämlich, daß der Vater diesen Traum nicht allein nicht verachtet, sondern diese Worte fleißig behalten habe; welcher Art und Weise zu reden der Evangelist Lucas auch gebraucht, da er Luc. 2, 51. sagt: Maria habe alle Worte, so von den Hirten sind geredet worden, behalten, und sie in ihrem Herzen bewegt.

108. Es hat zwar Jakob seinen Sohn gestraft; denn es läßt sich wahrlich zuerst ansehen, daß es fast ungereimt und ungeschickt Ding sei; aber der Eindruck, den der Traum machte, war so stark, daß er denselben nicht gänzlich hat verachten können, wiewohl er denselben anfänglich widersprochen. Denn er ist der Meinung, daß Joseph der Erstgeborne sei von Rahel, die ihm auch ehrlich vertraut gewesen, welche er auch heftig geliebt, und um deren willen er auch ganze sieben Jahre sehr harte und schwere Arbeit gethan, und so viel Dienstbarkeit getragen, so ihm ganz beschwerlich gewesen. Dieser Sohn aber, auf den er lange Zeit gewartet, war ihm noch endlich spät und im äußersten Alter geboren. Wen sollten nun diese und dergleichen wichtige Ursachen nicht bewegen, es wäre denn, daß ihm von Gott etwas Anderes geoffenbaret würde!

109. Zum Letzten, so hat auch Jakob die Verheißung gehabt von dem künftigen Samen, in dem alle Völker sollten gesegnet werden, 1 Mos. 28, 14. Dieselbige haben alle Patriarchen hoch geachtet, und Jakob hat ohne Zweifel darauf gehofft, daß der Same von Joseph kommen sollte, und hat gedacht, er würde ein solcher Herr werden, vor dem sich alle seine Brüder neigen würden, dazu auch alle Väter und Mütter in seinem ganzen Geschlechte.

110. Derhalben hat er auch die ganze Sache auf das allerfleißigste bei sich selbst erwogen, und gedacht: Wer weiß dennoch, was das somnium (der Traum) bedeuten möchte? Wer weiß, was dies wird für ein Zeichen oder Anbeten sein, es sei jetund gegenwärtig oder aber noch zukünftig? Wo von Joseph der gebenedeyte Same wird geboren werden, da werden nicht allein ich und die Mutter, sondern auch alle unsere Nachkommen sich neigen und anbeten; wie es denn hernach auch in Christo ist erfüllt worden. Und David, der auch Rathgeber ist, hat diesen Sohn, da er noch nicht geboren war, auch angebetet, Ps. 110. Desgleichen haben auch gethan Adam und Eva, und alle lieben Väter vom Anfang der Welt, welche diesen Sohn noch jetund für und für und in alle Ewigkeit anbeten.

111. Solches ist ohne Zweifel dem Patriarchen Jakob auch also eingefallen. Darum ist es ja ein wunderbarer Traum, da Vater und Mutter verstanden werden unter der Sonne und Mond, die Brüder aber und nachfolgenden Fürsten und Väter unter den Sternen. Darnach wird er auch noch etwas mehr und Höheres daraus geschlossen haben, nämlich, daß die vornehmsten und höchsten Creaturen Zeichen von uns sind, und nicht die Sache selbst, die dadurch angedeutet ist. Wir selbst sind das Ding, und sind so viel vortrefflicher, so viel ein jedes Ding mehr und größer ist, denn sein Zeichen, dadurch solch Ding bedeutet wird. So nun Gott Vater und Mutter und Brüder so großer Ehren werth achtet, daß die allerschönsten Creaturen ihre Zeichen sein müssen, so folgt nothwendig daraus, daß die menschliche Natur von Gott auch groß geachtet werden muß. Die Sonne ist ein Zeichen, so den Vater bedeutet. Derselben folgt, daß die Haushaltung und das weltliche Regiment, und viel mehr die Kirche vor Gott in den größten Ehren und Würden müssen gehalten werden. Dies und dergleichen viel mehr hat, sage ich, Jakob aus dem Traume Josephs, seines Sohnes, geschlossen.

112. Derhalben ist dies eine sehr schöne Erklärung gewesen, so Jakob erstlich des künftigen Samens erinnert um der Verheißung willen, so ihm und den andern Vätern gegeben war. Mein Same, wird er gedacht haben, wird den künftigen Christus anbeten: Sonne, Mond und Sterne, das ist, was von Fürsten, Königen, Vätern und Müttern ist in Reiche und einer jeden Regierung, diese alle mit einander werden seinem Reiche unterworfen sein. Darum hat er traun seinen Sohn nicht verworfen als einen falschen Propheten oder sonst einen losen unnützen Träumer, sondern hat sich wissen zu erinnern der Verheißung vom künftigen Samen, und hat bei sich erwogen, wie groß desselbigen Samens Majestät und Herrlichkeit sein würde. Er sei wohl ein Vater des Samens, wie Abraham auch gewesen; jedoch wolle dies die Noth erfordern, daß er sich vor ihm werde neigen und diesen Samen anbeten müssen. Und das ist ein rechter Gedanke, wiewohl er noch der Person halben geirrt, dieweil er meint, der Same solle von Joseph geboren werden. Aber GOTT hat hernachmals die Person geändert, und doch den Namen und die Deutung ungeändert gelassen, und hat diese Ehre dem Reiche Juda zugeeignet; wie das auch Jakob hernach lernen wird.

\quellseite{279}

113. Darnach ist die Auslegung dieses Traums auch nach der Historie sehr trefflich. Denn das ist der größte Schmuck der Polizei oder weltlichen Regierung und auch der Haushaltung, daß GOTT Vater und Mutter, Brüder und Hausgesinde höher ehrt, denn Sonne, Mond und Sterne. Diese Auslegung rühmt den Hausstand und gefällt mir sehr wohl. Und könnte das vierte Gebot nicht mit größerer Ehre oder Herrlichkeit erklärt und geziert werden. Zum Letzten sind auch in der Kirche rechte und hohe Sterne, nämlich, die Diener und Lehrer, welche doch gleichwohl unter Christo sind, daß er der Alleroberste und allein König sei über den ganzen Erdkreis.

114. Bis hierher haben wir, was zu der Historie gehört, erklärt. Nun aber machen uns die Grammatiker auch etwas zu schaffen. Denn es wird hier gefragt: Wie oder welcher Gestalt doch dies möge erfüllt sein, daß im Traume gesagt wird; wie es denn ja der Heilige Geist selbst eingegeben, daß Vater, Mutter und auch die Brüder sich vor Joseph neigen sollen, so doch die Mutter verstorben ist? An diesem Orte haben sie sich alle abgemüht, wie doch Rahel sich vor ihrem Sohn geneiget, welche beinahe schon zwei Jahre oder, meines Trachtens, noch nicht gar zwei Jahre verstorben, ehe denn Joseph etwas geträumt. Denn bis die höchste ganze Zeit über ist er durch wunderbare Zufälle hin- und hergeworfen worden, welche so groß gewesen, daß sie auch nicht wohl mit Worten ausgedrückt werden können.

115. Joseph ist nun siebenzehn Jahre, denn er war geboren sieben Jahre nach der Hochzeit. Da er aber geboren, ist Jakob noch sechs Jahre lang in Mesopotamien geblieben, und hat darnach acht Jahre zu Sichem gewohnt, wie das aus unserem Chronikon zu sehen ist. Nach der Zeit ist er gen Bethel gezogen, und ist Benjamin geboren, da nun Debora allbereit gestorben und bis die Mutter Josephs in der Geburt mit dem Tode abging. Da Joseph nun sechzehen Jahre gewesen, da ist der Schmerz und die Wehmuth noch frisch gewesen wegen des Todes der Rahel. Denn ich halte dafür, daß nicht mehr denn anderthalb Jahre dazwischen gewesen; und hat doch GOTT diesen Mann wiederum getröstet mit der Hoffnung von der Erstgeburt dieses Sohnes, da kurz zuvor Debora, die fromme heilige Matrone, die Vorsteherin und Beratherin des Hauses, verstorben war, dergleichen, da auch um dieselbige Zeit Dina geschwächt, welches ein Mägdlein gewesen von elf oder zwölf Jahren. Denn zu derselben Zeit ist die Natur der Menschen kräftiger und tüchtiger gewesen; wie wir denn auch hernach hören werden, daß Juda in seinem zwölften Jahre ein Weib genommen.

116. Derhalben, nach so vielen und großen Trübsalen, als daß die zu Sichem erschlagen und ausgerottet, Dina geschwächt, Rahel, sein Weib, und dazu auch Debora, die Amme, welche eine Säule und Zierde der ganzen Familie gewesen, verstorben; deßgleichen, nachdem sein Sohn Ruben die Blutschande begangen: da gibt GOTT dem Patriarchen Jakob nun wiederum einen neuen Trost durch diesen Sohn Joseph, den er billig für den Erstgebornen gehalten; denn Benjamin ist dazumal noch ein Kind gewesen von ein oder zwei Jahren. Und wird offenbar dem Joseph in diesen Träumen die Herrschaft verheißen, dazu auch der Same, so alle Völker segnen würde; wie er es selbst durch Erleuchtung des Heiligen Geistes ausgelegt. Derhalben ist nun dazumal Jakob wiederum gleichsam aus dem Tode und der Hölle erweckt und wieder lebendig geworden, ist nun wieder im Himmel. Siehe, wird er gedacht haben, GOTT will haben, daß dieser dein allerliebster Sohn der Erstgeborne sein soll und der rechte Erbe des Segens. Darum ist er nun so fröhlich, und hat sich dessen sehr gefreut und GOTT von ganzem Herzen gedankt. Bald fällt er wieder in den Abgrund der Hölle eines solchen unermeßlichen Falles halben, so sich ganz unversehens zugetragen.

117. Nach dem hohen und großen Trost, darin erklärt worden ist, daß Joseph größer und höher werden sollte, denn Sonne und Mond, die Sterne, und auch der Himmel und Erde; sintemal von ihm der herkommen würde, welchen Sonne und Mond, alle Patriarchen und Könige, ja, auch die Engel selbst anbeten würden; da er denkt, nun werde ich auch einmal Frieden haben: da fällt diese schöne liebliche Hoffnung sammt allen diesen Gedanken jämmerlich und bald in einem Augenblick wiederum dahin, nämlich, da Joseph weggerissen worden durch die Sünde und Bosheit seiner eigenen Brüder. Denn dieselbigen haben diesen Traum auf eine andere Weise ausgelegt und haben daraus also geschlossen: Rahel, so gestorben ist, wird sich ja vor ihrem Sohn nicht neigen können; derhalben ist dieser Traum vergeblich und ein Betrug des Teufels. Sie haben also dieser Ursache halben ihres Bruders Josephs Theilheit mit lauter Stimme verlacht, als der sich eigene Träume von seiner verstorbenen Mutter gedichtet. Und ist also ihr Stolz, Hoffart und Vermessenheit immer größer geworden des Regiments halben, daß solches gewißlich bei ihnen bleiben würde.

118. Denn solche Gedanken kommen gemeiniglich den Leuten also vor, wenn GOTT an ihnen die Sünde strafen will, nämlich, daß sie GOTTes Wort und die rechte gewisse Wahrheit für ein solch Ding halten, das auch abscheulich anzuhören sei. Also soll es gehen.

119. Es sollten diese Brüder, dem Exempel ihres Vaters nach, den Traum Josephs etwas fleißiger erwogen und ihre verbotenen Lüste im Zaum gehalten haben. Aber sie fallen dahin, als wären sie unsinnig, in große greuliche Sünde, daß sie schändliche Mörder werden auch an ihrem eigenen Vater, Bruder und Mutter. Denn so viel an ihnen ist, haben sie in Wahrheit ihren alten Vater getödtet, welcher erst recht fröhlich geworden war und an Joseph, seinem lieben Sohne, sonderliche Lust und Freude gehabt. Weher hätten sie ihm nicht thun können, wenn sie ein Schwert oder Spieß durch ihn gestochen hätten: größere Pein hätten sie ihm nicht bereiten können.

120. Sie werden zwar genugsam vermahnt durch den Traum, wie solcher von Joseph, ihrem Bruder, ist erzählt worden, und hat ihnen sonst an andern wichtigen Argumenten auch nicht gemangelt; als da ist, die Natur, Gesetz oder das Recht, die Sache an sich selbst, dazu die Unschuld des frommen Jünglings, endlich auch die Strafe des Vaters, welcher sie oftmals zur brüderlichen Liebe ermahnt. Durch diese Stücke alle sollten sie ja billig bewogen und erweicht worden sein, wo sie nicht eiserne Herzen gehabt. Aber wenn GOTT steht, daß Pharao ja will verhärtet und verstockt sein, alsdann läßt er dem Teufel den Zaum, und läßt Ursachen vorfallen, daß er immer mehr verhärtet und verstockt werde.

121. Also haben nun diese Brüder Josephs alle Vermahnung ganz frei verachtet, und gedacht: Joseph ist ein Narr, der Vater selbst ist ein schwachsinniger Greis, es sind grobe ungeschickte Leute, welchen es an Rath und guten Gedanken mangelt. Wir aber sind die Regenten, welchen gebühret im Hause und in der Kirche unseres Vaters zu regieren. Derhalben wollen wir dieses losen Träumers gar nicht achten und nichts auf ihn geben.

122. Was wollen wir aber darauf für Antwort geben, daß sich gleichwohl dies offenbar ansehen läßt, als wolle es sich mit nichten reimen, nämlich, daß Rahel nun schon gestorben war? Ezra sagt, der Vater habe erstlich diesen Traum gesagt. Aber ich bin damit nicht zufrieden. Denn wiewohl er sich anfänglich also geberdet, als habe er des Traums nicht sonderlich geachtet, ja, daß er ihn auch gestraft, hat er denselben doch hernach behalten und in seinem Herzen erwogen, daß er ihn immer in frischem Gedächtniß gehabt. Augustin läßt es sich sehr sauer werden, und bringt auch einen theologischen Verstand herzu, nämlich, daß dies alles solle von dem künftigen Christo verstanden und in ihm auch erfüllt werden, welchen Adam, Eva, und alles, was lebendig und todt ist, auch die Hölle selbst, anbeten; wie Paulus Phil. 2, 10. sagt: „Daß in dem Namen Jesu sich beugen sollen alle derer Kniee, die im Himmel, und auf Erden, und unter der Erde sind.“ Dies ist theologisch und nach dem Verstande der heiligen Schrift recht geredet und ist alles auch also wahr.

123. Was wollen wir aber antworten, soweit es sich auf den Haushalt Jakobs bezieht? Wie kann die verstorbene Rahel sich vor jemand neigen? Da folge ich der Antwort, so in der heiligen Schrift gebräuchlich, solche Fragen aufzulösen, nämlich durch die Figur, synecdoche genannt. Denn wiewohl Rahel nicht lebt, daß sie sich für ihre Person neigen konnte, hat sie sich doch in ihren Söhnen, in dem Mann und in dem ganzen Geschlecht vor ihm geneigt; wie wir auch nach Art derselben Figur sagen, daß alle Brüder den Joseph gehaßt, wiewohl Juda, wie ich dafür halte, ihn nicht gehaßt, und die Andern auch nur in den Haß gewilligt, sind aber doch nicht wie Simeon und Levi Rädelsführer gewesen.

\quellseite{280}

124. So ist nun dies der Verstand des Traums: Es werden sich vor mir neigen Sonne und Mond, und auch die Sterne, das ist, das ganze Geschlecht, Eltern, Brüder, alle Vorfahren und Nachkommen miteinander. Hernach aber am Ende werden wir hören, wie Jakob sich vor seinem Sohn Joseph geneigt, da er Cap. 47, V. 29. sagt: „Habe ich Gnade vor dir funden, so lege deine Hand unter meine Hüfte, daß du die Liebe und Treue an mir thust, und begräbest mich nicht in Egypten“ etc. Denn da redet er mit seinem Sohne wie mit einem Regenten, nach Art und Weise seines Königs, so den Andern steht und als ein Unterthan um etwas bittet. Denn wiewohl er solches aus väterlichem Herzen geredet, hat er es ihm doch gleichwohl nicht geboten, sondern ihn als seinen Oberen gebeten. Und in der Epistel an die Hebräer am 11. Cap. V. 21. steht: „Durch den Glauben segnete Jakob beide Söhne Josephs, und neigte sich gegen seines Scepters Spitze.“ Denn wiewohl diese Stelle auf eine andere Weise möchte verdolmetscht werden, läßt es sich doch gleichwohl also ansehen, daß Jakob seinen Sohn Joseph angebetet, das ist, sich vor ihm geneigt habe. Denn also sagt der Text: Da Joseph sei herzugekommen, habe sich Israel geneigt auf dem Bette zu den Häupten, nicht anders, denn als wäre ein großer Fürst zu ihm gekommen. Und solches hat er nicht allein um des Regiments willen gethan, so dazumal Joseph über ganz Egypten gehabt; sondern auch um des gebenedeyten Samens willen, von dem er gehofft, daß er von Joseph herkommen würde, und davon er auch zuvor eine Offenbarung empfangen hatte. Dies ist eine seine einfältige und richtige Antwort, diese Frage aufzulösen, nämlich, durch diese Figur synecdoche, welche in der heiligen Schrift gebräuchlich ist.

125. Der andere Knoten, so allhier dem Augustin zu schaffen gemacht, ist in der Erzählung, da vom Tode Isaaks geredet wird, daß wir droben beim 35. Cap. V. 29. auch gedacht.

Es ist aber die Erklärung auch leicht, nämlich durch die Figur oder Weise zu reden, so da heißt hysteron proteron, welche gemeiniglich im Lateinischen genannt wird anticipatio oder recapitulatio, das ist, wenn man die Rede umkehrt, und dasjenige, so vorher stehen soll, hernach setzt. Welcher Figur die heilige Schrift oftmals pflegt zu gebrauchen; wie denn die, so der Heiden Historien beschrieben, derselben auch viel gebrauchen. Als, wenn Livius die Geschichte eines Monats oder Jahres beschreibt, kann er nicht zugleich alles erzählen, was auch an andern Orten zur selbigen Zeit geschehen. Also hat Mose zuerst das Leben Isaaks zum Ende gebracht, darnach Josephs; und sind doch gleichwohl noch etliche Dinge zugleich geschehen, welche zugleich nicht haben können beschrieben werden. Wie Zwei wohl zugleich mit einander singen, aber doch nicht zugleich reden können. Ja, es singt wohl zugleich mit einander ein ganzer Chor, da sie aber doch nicht alle zugleich mit einander predigen könnten.

126. Derhalben haben Anra und Augustin die Rechnung recht gemacht, daß Joseph verkauft worden sei zwölf Jahre, ehe Isaak gestorben: sie vergessen aber gleichwohl der Figur, so in den Historien gemeiniglich gebraucht wird, nämlich, da etwa das Vorderste hernach gesetzt wird. Denn droben, Cap. 35, V. 29., eben in der Ordnung und Linie der Geschichte Isaaks sein Tod auch mitgesetzt und beschrieben wird, in welcher Erzählung gleichwohl die Dinge, so auch zur selbigen Zeit geschehen, nicht haben mit erzählt werden können. Darum hätte sich Augustin selbst leichtlich heraushelfen können, wo er den Tod Isaaks, so zuvor, ehe denn Joseph verkauft, erzählt worden, durch diese obgenannte Figur, hysteron proteron, verstanden hätte.

127. Denn man darf deshalben nicht das Ansehen und die Glaubwürdigkeit der heiligen Schrift in Zweifel stellen, sondern der Leser muß sich deß wissen zu erinnern, daß die Schrift nach der Art und Weise der Historien-Schreiber ihre Geschichte auch pflege zu erzählen; wie ich davon an andern Orten oftmals Erinnerung gethan.

128. So hat nun Mose bis hieher die zwei Träume Josephs erzählt, in welchen der heilige Geist gleichsam ein Vorspiel gemacht hat zu der vortrefflichen Historie vom künftigen Regiment oder Herrschaft, so Joseph in Egypten gehabt, oder wie es die Schrift an einem andern Ort, Psalm 105, 17., ausgelegt, von dem Manne, so vor Israel her nach Egypten ist gesandt worden, daß er zur Zeit der Theurung beim Leben möchte erhalten werden. Daraus wir denn sollen erkennen lernen die väterliche Sorge, so Gott für uns trägt, der lange zuvor anzeigt, und was zukünftig ist, verkündigt, und dasjenige auch wunderbarlich nach seiner Güte also mäßigen, schicken und regieren kann, daß es gleichwohl ein solch Ende gewinnen muß, daß wir es auch ertragen können; wie es denn hernach Joseph selbst noch also deuten wird, da er Cap. 50, V. 20. sagt: „Ihr gedachtet es böse mit mir zu machen, aber Gott gedachte es gut zu machen.“ Desgleichen am 45. Cap. V. 7.: „Gott hat mich vor euch her gesandt, daß er euch übrig behalte auf Erden, und euer Leben errette durch eine große Errettung.“ So ist es nun gewiß, daß Gott treulich und fleißig für uns sorgt, ehe wir selbst denken können, wie es mit unsern Sachen geraten möge.

129. Ueber das wird auch mit solchen Exempeln der Menschen Thorheit gestraft, die wir alle so närrisch, ja, auch so hochsichtig sind, daß wir oftmals unsern eigenen Nutzen hindern und unser selbst größte Feinde sind, wo nicht der gnädige GOtt nach seiner unermesslichen Güte unsere Rathschläge umkehrte, und dieselbigen anders regierete, denn wir selbst gedacht. Denn also unsinnig sind diese Brüder Josephs, ihren Bruder zu verderben und ganz auszurotten, der doch ihr, dazu auch des ganzen Landes Canaan Heiland werden sollte: und wo, wo er also umgekommen wäre, hätten sie doch müssen Hungers sterben. Aber Gott hat alles das Böse, so sie gedacht, auch wider und vor ihren Gedanken her wunderbarer Weise zum Guten gewandt, und dasselbe ist etlichermaßen durch diese Träume also zuvor bedeutet und vorbereitet worden.

130. Nun wollen wir fortfahren, und zu dem, was weiter folgt, greifen, und wollen erstlich an Jakob ein wunderbares Exempel göttlicher Regierung sehen, und dasselbige betrachten, nämlich, wie GOtt die Seinen in die Höhle hinein und wieder heraus zu führen pflege, wie er sie tröste und betrübe, auf daß wir uns daran gewöhnen, daß wir allein Gottes Wort vertrauen und demselbigen anhangen. Es muß mit den Göttlichen dahin kommen, daß sie nichts fühlen, nichts Anderes verstehen, nichts sehen, sonderlich im Tode, daran sie sich halten mögen, denn allein das bloße Wort. Darum müssen wir nach dem Fleisch sterben lernen und am Wort halten, daß wenn auch die ganze Welt brechen oder über einen Haufen fallen sollte, wir gleichwohl mit festem Vertrauen sagen könnten: „Ich glaube an GOtt den Vater, allmächtigen Schöpfer“ etc. „Und an Jesum Christum, seinen einigen Sohn“ etc. Jakob aber hat jetzt zu Hebron gewohnt, hat einen Sohn von siebenzehn Jahren, auf den er in seinem Alter Hoffnung und Trost gesetzt, und meint, es sei nun der Jammer aus: so geht er erst ärger an, denn er ihn sein Lebenlang gehabt hat.

\section[Dritter Theil]{Dritter Theil. Wie Jakobs Söhne gen Sichem zur Weide ziehen; wie Joseph von seinem Vater zu ihnen gesandt wird; und wie diese Aussendung voller Gefahr gewesen}

\subsection{I}

\begin{quote}
\textbf{V. 12—14. Da nun seine Brüder hingingen, zu weiden das Vieh ihres Vaters in Sichem, sprach Israel zu Joseph: Hüten nicht deine Brüder des Viehes in Sichem? Komm, ich will dich zu ihnen senden. Er aber sprach: Siehe bin ich. Und er sprach: Gehe hin, und siehe, ob es wohl stehe um deine Brüder, und um das Vieh; und sage mir wieder, wie sich’s hält. Und er sandte ihn aus dem Thal Hebron, daß er gen Sichem ginge.}
\end{quote}

131. Aus diesem Text erhellt, daß die Weide bis daher bei Hebron gewesen, da Isaak und Jakob mit ihrem Hausgesinde gewohnt. Und ist kein Zweifel, daß dies gute fromme Leute gewesen sind, so denn Jakob und seinen Söhnen ihre Aecker verpachtet haben, dieselbe nicht Eigenes in dem Lande gehabt; und wird gleichwohl allhier der Garben oder Aecker gedacht, welche sie gebauet. Es scheint aber, als seien die Einwohner des Landes Nachkömmlinge gewesen der drei Brüder Aner, Escol und Mamre, so gottseelig und fromm gewesen, haben Gottes Wort fleißig gehört, und sich auch zu der Kirche oder Versammlung, so Abraham und Isaak gehabt, gehalten.

\quellseite{281}
132. Es scheint auch, daß dazumal die Weide auf den Ackern und Wiesen zu Hebron fast ganz abgeweidet gewesen; welches auch die Ursache gewesen, daß sie gen Sichem gezogen: oder aber sind darum an einen andern Ort gezogen, daß sie die Gegenwart und Gemeinschaft Josephs nicht haben länger ertragen können, wider welchen sie mit so großem Haß und Neid entbrannt, daß sie auch nicht mehr freundlich und friedlich mit ihm haben reden mögen, ja, haben ihn vor Augen nicht wohl sehen mögen. Denn das kann wohl geschehen, daß Einer deswegen, daß er einem Andern feind ist und ihn neidet, von ihm flieht; wie von Kain gesagt wird, 1 Mos. 4, 5.: „Da ergrimmte Kain sehr, und seine Gebärde verstellete sich“; er hat seinen Bruder Abel nicht freundlich und gütig ansehen können, sondern sobald er seiner ansichtig geworden, ist er ergrimmt, und hat mit solchen sauern unfreundlichen Gebärden seines Angesichtes angezeigt, daß sein Herz mit bitterem Zorn entbrannt. Eben solche Ursache hat diese Brüder auch können treiben, daß sie an einem andern Ort Weide gesucht; Joseph aber ist daheim bei dem Vater geblieben, da die Andern also von ihm gewichen und seine Gemeinschaft gemieden.

133. Darnach sehen wir auch, daß sein Geschlecht gleichwohl Sichem noch inne gehabt, sintemal sie dahin auf die Weide gezogen. Denn da die Bürger zu Sichem erschlagen, ist Jakob mit seinen Söhnen an ihre Statt gekommen, und die Mitten der Sichemiter, und so noch Etliche mehr von dem Volk übergeblieben, sind ohne Zweifel mit Jakob daselbst in ihrem Vaterlande auch geblieben. Und ich halte dafür, daß auch die Nachbarn im Lande Canaan mit ihnen versöhnet gewesen, nämlich, durch die Schwägerschaft, so sie mit Juda gemacht, wie hernach folgen wird in 38. Cap. V. 2. Daher sie die Andern auch geliebt, welche züchtig und fromm gewesen. Aber die allervornehmsten Buben, Simeon und Levi, die Mörder zu Sichem, werden sie ohne Zweifel verabscheut haben.

134. Es hat aber denselben Ort darnach Joseph eingenommen mit seinen Nachkömmlingen, an welchem er beinahe von seinen Brüdern ist erschlagen worden. Denn hernach wird ihm dieses Land von seinem Vater gegeben, dasselbe zu besitzen, welches ja Jakob nicht hätte thun können, wo er es nicht hätte als sein eigen Land inne gehabt. Derhalben sind sie von dem Acker, so ihnen vermiethet gewesen, gen Sichem gezogen auf ihre eigene Weide. Wiederum aber, da sie auf denselben Ackern nicht mehr Weide gefunden, haben sie sich wieder nach Hebron gewendet. Denn sie können und wollen nicht länger bei dem schändlichen bösen Verräther bleiben, welcher haben wollte, daß sich die Andern alle vor ihm neigen sollten.

135. Ferner wird von diesem Lande Sichem das Königreich Israel und Ephraim auch das Königreich Josephs genannt und ist der erste König Jerobeam aus diesem Geschlechte geboren gewesen. Denn sie haben auf die Prophezeiung gesehen, darin hernach die Herrschaft dieses Orts dem Joseph verliehen wird; welchen Ort Jakob, wiewohl er durch große Tyrannei seiner Söhne eingenommen, dennoch behalten, da Gott ihm denselben zugemessen, und also die böse That seiner Söhne zum Guten gewendet. Denn solche bürgerliche und zeitliche Güter, davon Gott nichts Sonderliches verordnet oder geboten, sind nur zufällig, welche Gott austheilt, wem er will. Wer hat, der hat. Jakob hat davon keine Verheißung oder Befehl gehabt, daß die zu Sichem also sollten erschlagen werden; aber so Hemor mit seinem Sohn erschlagen war, hat er das Land behalten mit Bewilligung der Völker, so am nächsten umher gewohnt. Also geht es noch ohngefähr und zufälliger Weise Vielen, daß ihnen viel Gutes widerfährt, wo es Gott also verhängt, davon sie doch gleichwohl kein sonderlich Gebot noch Verheißung Gottes gehabt.

\subsection{II}

136. Was nun weiter im Text noch übrig, ist nicht sonderlich schwer zu verstehen, weder was die Historie noch auch was die Grammatik belanget. Mose beschreibt die Sorge des alten Vaters Jakob, das Hausregiment belangend; welcher, dieweil er hochbetagt und altersschwach gewesen, nicht hat können persönlich bei seinem Vieh sein, doch gleichwohl für seine Söhne und das Vieh Sorge getragen. Dies wird aber also vom Heiligen Geist uns vorgehalten, damit wir uns darnach richten, daß wir nicht alle Sorge um dies zeitliche Leben ablegen; wie die Mönche gethan, welche nur von anderer Leute Arbeit und Sorge gelebt. Denn dies ist ein Exempel, so zum Hausregiment gehöret; und wird allhier sehr fein beschrieben die väterliche Sorge, dadurch Jakob bewogen, auf seine Kinder und Vieh Acht zu geben, wie Alles möchte recht verwaltet und ausgerichtet werden. Darum schickt er Joseph zu ihnen, und sagt: „Gehe hin, und siehe, ob es wohl stehe.“ Denn droben, V. 2., ist gesagt, daß Joseph vor seinen Vater gebracht, wo von seinen Brüdern etwa ein böses Geschrei wider sie gewesen. Er wird aber nicht allein dem Vater angezeigt und verkündigt haben, was seine Brüder Uebels gethan, sondern auch um Knechte und Mägd[e] ungerecht gehandelt, zudem, wo am Vieh oder sonst in der Haushaltung etwa Schade geschehen.

137. Und ist wahrlich diese Sorge mit nichten von denen zu verachten, so dazu berufen und bestellt sind, daß sie in der Kirche, im Welt- und auch im Hausregiment Andern sollen vorstehen, daß sie auch zugleich daheim in ihrem Hause Bischöfe und Aufseher seien, und auf ihre Aecker, Vieh und Anderes mehr ein fleißiges Aufsehen haben; wie allhier Jakob gethan.

138. Der Papst hat einen Unterschied gemacht zwischen weltlichen oder fleischlichen Dingen und den geistlichen, und hat gelehrt, daß man sich vor der Haushaltung hüten solle, als welche Gott nicht sonderlich geachtet, ja, gar verworfen habe; hat also haben wollen, daß seine Geistlichen solcher Hausorge ganz und gar sollten überhoben sein. Denn sie haben gesagt: Wir dienen Gott in himmlischen und geistlichen Sachen; derhalben wollen wir die heidnischen und weltlichen Dinge fahren lassen etc. Was sollen wir aber nun gleichwohl thun, sagten die geistlichen Herren? Ja, wir wollen dieweil unsere Siebenzeiten, Metten und Vesper lesen, und dafür der Welt Güter fressen und verzehren, wollen über Könige und Fürsten herrschen und regieren etc.

139. Derhalben sollen nun fromme gottselige Hausväter wissen, daß ihre Werke alle Gott wohlgefallen, sie warten nun des Viehes, bauen den Acker, oder auch wenn sie im Rath oder der Mistgrube arbeiten; die Hausmutter aber melkt die Ziegen und Kühe etc. Diese Werke alle soll man für geistlich oder gebührlich halten, wenn sie von christgläubigen Leuten geschehen. Denn wo sind wir sonst alle hergenommen, denn aus dem Ehestande und dem Hausregiment, welches die Papisten gar verachten, als ob sie von Klötzen und Steinen geboren wären. Wo die Mütter nicht Kinder gebären, sie säugen, wischen und ernähren würden, müßte doch nothwendig das ganze menschliche Geschlecht vergehen.

140. Dieweil aber nun Gott selbst diese Stände und Werke gestiftet und geordnet, soll man mit nichten denken, daß dadurch der Gottesdienst verhindert werde, sondern soll es dafür halten, daß es sehr gute und wohlgefällige Uebungen seien der Gottseligkeit gegen Gott und die Menschen. Denn Gott will haben, daß die Frucht in Mutterleibe getragen, gesäugt und ernährt werden soll, und solche durch den Fleiß und Sorge der Mütter also ernährt, auferzogen werde und immer zunehme. Und darum hat er auch in die Brüste Milch gegeben.

141. So aber doch jemand die Gabe der Keuschheit verliehen und gegeben ist, daß er ohne Hausorge und Beschwerlichkeit leben kann, derselbige mag traun dieser seiner Gaben gebrauchen. Es ist aber ein Anderes, frei sein von etlichen solchen Beschwerungen dieses Standes, und ein Anderes, den Stand oder das Leben an sich selbst verdammen und verwerfen, als ob derselbe ungöttlich und heidnisch wäre.

142. Wenn wir aber diese Sache recht erwägen wollen, werden wir wahrlich finden, daß nichts ungöttlicher oder weltlicher sei, wie die Papisten pflegen zu reden, denn der ganze Stand der päpstlichen Bischöfe; wiewohl sie einfältige ungelehrte Leute mit dem Titel oder Namen, daß sie sich geistliche Personen nennen, betrügen können; ja, sie selbst sind vielmehr zu verdammen wegen des verkehrten und ungöttlichen Urtheils, damit sie Gottes Ordnung verdammen, davon sie halten, daß sie voller Gefährlichkeit und Müheseligkeit sei.

143. Und man soll und muß immer fleißig sehen auf den, so diesen häuslichen Stand gestiftet, nämlich, daß Gott selbst solche häusliche Werke eingesetzt, dieselbigen lobt, belohnt und ziert. Denn also hat allezeit die heilige Schrift in den Historien der Väter diese scheinbar geringen, närrischen, verachteten und weltlichen Werke mit den hohen großen Werken des Glaubens und Kreuzes, und andern geistlichen Nebigkeiten vermengt, daß sie damit anzeige und bezeuge, daß GOtt sowohl für die größten, als auch für die geringsten Dinge Sorge trage. Denn es ist Ein HErr beides der Großen und der Niedrigen, der Könige und Knechte, der Männer und Weiber etc. Wir haben alle denselbigen einigen GOtt, und wir sind auch eins, in Einigkeit und auf einerley Weise diesem einigen GOtt zu dienen; wiewohl unsere Wercke und Beruf ungleich sind.

\quellseite{282}

144. Es soll aber jeder in seinem Berufe thun, was ihm daselbst gebühret. Wie denn auch Jakob ein heiliger und geistlicher Mann ist, redet vom Gesetz des HErrn, betet, steiget und lehrt die Kirche: hat aber gleichwohl inzwischen andere geringe häusliche Wercke, den Acker und Vieh betreffend, nicht verachtet, oder dieselbigen gar liegen lassen. Welches denn uns also zum Exempel vorgehalten wird, daß wir wissen sollen, alle unsere Wercke im Hausregiment gefallen GOtt auch wohl, und daß sie zu diesem Leben nöthig sind, darin Jedem gebühret, nach seinem Vermögen und in seinem Beruf dem einigen GOtt, der unser aller HErr ist, zu dienen.

145. Dies soll oftmals mit Fleiß getrieben und dem zarten Alter gleich zuerst mit eingepräget werden, daß es von diesen Dingen besser möge lernen urtheilen, denn etliche Rechtsverdreher und Papisten pflegen, welche nicht aufhören können, die Orden und Regeln der Mönche und Nonnen den Wercken vorzuziehen, so zur Haushaltung und weltlichem Regiment gehören, wie hoch auch dieselben von GOtt gelobt und geziert werden. Die Jungen sollen wissen, daß GOtt ebenso wohl gefalle eine Frau, so ihr Kindlein säugt, oder auch eine Magd, welche das Haus kehrt, als eine faule Nonne oder Carthäuser-Mönch, der gute faule Tage hat. Dies, sage ich noch, soll man wissen wider den Greuel des Papstes, so vor dieser Zeit in die Kirche eingeführt, da sie alle anderen Stände weltlich und unnütz genannt, und allein ihre Platten und Kappen hoch gerühmt und gepreiset haben. Wir, sagten sie, sind die Geistlichen, wir beten für die andern: ihr Bürger, Bauern, Fürsten, Adel und alle Weltlichen könnt GOtt nicht dienen. Derhalben gebt uns her Gold, Silber, eure Schlösser, Städte und Lande, so wollen wir Tag und Nacht für euch bitten, dieweil ihr den Acker bauet, euer Hausgesinde und Unterthanen regieret. Denn also habe ich selbst auch gelehrt, da ich noch ein Mönch war. Aber GOtt hat mir nach seiner großen Gnade und Barmherzigkeit meine Sünde vergeben.

146. Man soll aber doch oftmals noch bedenken, wie erschrecklicher Weise die ganze Welt bethört und betrogen gewesen durch die schändlichen Papisten, Canonisten und Sophisten, welche die Anrufung und andere Gottesdienste allein an die Klöster gebunden, als ob GOtt nicht auch könnte angerufen werden im gemeinen Leben. Denn sie hören nicht auf, die heillosen Leute, Juristen und Mönche, sich zu bemühen und noch immer allen Fleiß anzuwenden, daß nochmals dieser erschreckliche Greuel in Klöstern wiederum möchte aufgerichtet werden, damit die Wercke, so GOtt selbst geboten, also verdunkelt und vernichtet werden, als da sind: häusliche Wercke, so die Weiber im Hause ausrichten; oder auch so zum weltlichen Regiment gehören; welches ja in Wahrheit GOttes Wercke sind, sintemal ja das Weib, ein Bürger, Fürst, Edelmann GOttes Creaturen sind und Jeglicher unter uns zu einem sonderlichen Wercke oder Amte geschaffen ist, daß die solche Wercke in ihrem Amte fleißig ausrichten, wissen sollen, sie gefallen damit GOtt wohl, sofern sie sich nur vor Sünden hüten und sich derselben enthalten.

147. Dies Licht ist ganz und gar aus der christlichen Kirche weggenommen gewesen; und das ist durch die Päpste geschehen, also daß auch gar keine Erkenntniß, weder des Schöpfers noch der Creaturen, unter den Leuten mehr geblieben ist. Und die Bischöfe und Mönche werden noch versuchen, solch Licht wiederum auszuleschen, wo nicht ihre gottlosen Rathschläge und Vornehmen auf dem Reichstage, so jetzt gehalten wird, von GOtt verhindert werden. Denn das werden sie mit allen Kräften betreiben, daß das Pabstthum und die Klöster mit allen päpstlichen Decreten wiederum möge aufgerichtet werden; wiewohl sie selbst solche Decrete verachten und verspotten. Denn sie fasten und beten nicht, sondern schlemmen und dämmen, jeden Tag und Nacht im Saufe; sie werden aber doch versuchen, den Kaiser zu betrügen, dieweil sie ihre großen Gottesdienste also hoch zu rühmen wissen, daß dieselben wiederum in Klöstern müssen aufgerichtet werden; gleich als ob die unreinen Buben, die Mönche, allein die seien, so GOtt einen wohlgefälligen Dienst leisten können mit ihrem Heulen und Brüllen in der Kirche.

148. Wider solche seltsame ungeheure Greuel des Pabstthums soll man diese und andere dergleichen Stellen wohl merken, nämlich, vom Leben und den Wercken der lieben Patriarchen, welche GOtt bisweilen sehr hoch geführt. Also, da Jakob mit dem Engel gerungen, daß ist der höchsten Wercke eins gewesen; denn da war nicht Raum oder Zeit, die Ziegen zu melken, sondern zu ringen und zu siegen wider die Sünde und Tod, und dasselbe durch den Glauben und Hoffnung. Bisweilen aber führt sie GOtt auch außerhalb solcher hohen Wercke; wie denn derselbe Jakob auch weltlich und häuslich gewesen, hat des Viehes, der Knechte und des Gesindes gewartet. Denn gottelige Menschen sind ja nicht allein zu einerley Werck bestellt, wie die Mönche sich selbst an sonderliche selbsterwählte Wercke gebunden und verpflichtet haben; sondern sie waren bisweilen der geistlichen Uebung, bisweilen haben sie auch mit bürgerlichen oder häuslichen Sachen zu schaffen.

149. Ueber das wird auch an diesem Ort gerühmt der Gehorsam Josephs in dem, daß er alsobald da ist, da ihm sein Vater gerufen. Denn diesen Gehorsam hat der Heilige Geist mit vielen Worten erzählt nicht allein um der folgenden Historien willen, sondern daß er an diesem Jünglinge Joseph solch kindliche Tugend rühmen möge; welches doch vor den Augen der Papisten nur leibliche und närrische Wercke sind, als wäre in ihrem Leben gar keine menschliche Schwachheit. Es sind aber gleichwohl diese Wercke, so von ihnen verachtet werden, vor GOtt die köstlichsten und besten.

\subsection{III}

\begin{quote}
\textbf{V. 15—17. Da fand ihn ein Mann, daß er irre ging auf dem Felde; der fragte ihn und sprach: Wen suchest du? Er antwortete: Ich suche meine Brüder; lieber, sage mir an, wo sie hüten. Der Mann sprach: Sie sind von dannen gezogen; denn ich hörte, daß sie sagten: Lasset uns gen Dothan gehen. Da folgte Joseph seinen Brüdern nach, und fand sie zu Dothan.}
\end{quote}

150. Joseph zieht sich fort und eilt mit sehr geneigtem und willigem Herzen zu seinen Brüdern, und da er eine Zeitlang auf dem Felde irre gegangen, hat ihm ein Mann, dem er auf dem Felde begegnet, den Weg gewiesen; er aber weiß nicht, wie große Gefahr ihm da werde begegnen und wie nahe ihm das Verderben sei. Denn es hat ihn niemand davor gewarnt, und GOtt selbst hat verhängt, daß er seinen Brüdern, so über die Maßen grausam waren, in die Hände gekommen ist. Also achtet GOtt, wie es scheint, ihrer beider nicht viel, nämlich, des Vaters, Jakobs, und Josephs, seines Sohnes, und stellt sich, als sehe er nicht, was Josephs Brüder schon vollänget im Sinne gehabt, daß sie sich auch öffentlich haben vernehmen lassen; und läßt es also zu, daß der Knabe vom Vater dahin gesandt und also in das Verderben hinein gestoßen wird. Wo sind denn nun die heiligen Engel, die Mahanim, von denen droben, Cap. 32, v. 1. 2., gerühmt wird, daß sie für Jakob wider Esau gestritten haben? Es erscheint hier keiner, der den Jakob warnt und sagt: Hüte dich, schicke deinen Sohn Joseph nicht zu deinen andern Söhnen; denn sie werden ihn tödten und alsdann wirst du müssen deines Herzens Lust beraubt sein. Alle Engel und GOtt selbst schweigen da stille. Dies alles ist wahrlich sehr wunderbar und ganz unerhört Ding, welches ich weder mit Gedanken, noch mit Worten erreichen oder ausdrücken kann. Denn GOtt verhängt allhier, daß beide Jakob, der Vater, und Joseph, sein Sohn, in das gegenwärtige Verderben hinein gerathen sind.

151. Uns aber ist dies ein großer herrlicher Trost, daß wir sehen, daß GOtt die heiligsten Patriarchen nicht anders führt und regiert denn uns auch. Denn er hat ja diesen Patriarchen Jakob nicht nur durch eitel Wunderwercke und hohe geistliche Uebungen geführt, welche den gemeinen Verstand frommer gottseliger Leute übertreffen möchten; sondern er wirft ihn auch hin in gemeine Gefährlichkeiten und Unfälle, wie sie einem Jeden von uns begegnen können. Joseph ist ein zarter, frommer unschuldiger Jüngling gewesen von siebenzehn Jahren, hat den Heiligen Geist, GOtt ist ihm gnädig und hat an ihm Gefallen; hat auch in der Gemeinde Jakobs, seines Vaters, ein sehr heilig Leben geführt: und wird doch mit diesem erbärmlichen Fall so gar hart beschwert, den, wiewohl er vielleicht uns nicht däucht ungewöhnlich zu sein, gleichwohl dieser heilige Mann gar nicht verdient hat.

\quellseite{283}

152. Deshalb steht dies Exempel allhier uns zu Trost geschrieben, und erinnert uns, daß wir daran denken, daß wir arme Menschen sind, und daß wir es deshalb dafür halten sollen, daß alles, was menschlich ist, uns auch begegnen könne. Denn so diese Beschwerungen dem heiligen Patriarchen, der des Heiligen Geistes voll war, widerfahren, was wollen wir uns denn verwundern und darüber erzürnet oder ungeduldig sein, wenn wir auch dergleichen erfahren müssen? Warum wollen wir uns dessen nicht vielmehr freuen und Gott dafür danken, wo wir fühlen, daß wir mit gleicher Gefährlichkeit auch also geübt und umgetrieben werden, damit Gott von Anfang auch die allerheiligsten Menschen versucht hat?

153. In dieser großen Noth und Gefahr aber sehen wir, daß Gott und die Engel dazu ganz stille sind, lassen es geschehen, daß Joseph in das betrübteste Elend, ja, in den Tod selbst gestürzt werde. Sie sehen, daß der Vater über den Tod seines lieben Sohnes sehr elend und betrübt werden wird, und haben doch gleichwohl dieser Brüder Vornehmen nicht gehindert.

154. Darum sollen wir uns unter einander ermannen zur Geduld mit dieser Leute Exempel, welche in ihrem Kreuz, so sie getragen, uns auch gleich gewesen. Und sind diese Exempel uns näher und faßlicher, derhalben sie uns auch mehr bewegen, denn das Exempel des Sohnes Gottes selbst, dieweil dasselbe höher ist und gar keine Vergleichung hat; wiewohl er von sich selbst auch sagt Matth. 26, 24: „Des Menschen Sohn gehet zwar dahin, wie von ihm geschrieben stehet“, gleich als ob weder der himmlische Vater, noch die Engel auf ihn Acht hätten, sintemal er den sehr grausamen Juden also übergeben und vorgeworfen wird, daß sie ihn jämmerlich zerreißen, verwundet und gekreuzigt haben.

155. Eben auf dieselbige Weise gehen diese Zwei auch dahin in den Tod hinein, und Gott schweiget dazu auch stille, die Engel schweigen auch, ja, sie haben sich noch wohl darüber gefreut. Denn es ist nicht geschehen, den Joseph oder Jakob damit zu verderben, sondern Vielen zum Heil. Aber dieser Rath Gottes ist noch heimlich und verborgen, wiewohl er sehr gut und nützlich gewesen. Denn durch diese Versuchung und Strafe wird Gott verhüten, daß Jakob mit seinen Kindern nicht Hungers sterbe, wird auch damit verschaffen, daß Egypten bekehrt und zur wahren Erkenntniß Gottes kommen, und daß also viel Leute gewonnen und selig werden mögen.

156. Dies ist wahrlich ein sehr heilsamer Rath, aber der vor ihren Augen ganz verborgen gewesen. Denn also hat Gott gedacht: Ich will zwar sie alle beide tödten; denn ich werde nicht allein ihnen Gutes thun, sondern will auch dem ganzen Egyptenland durch sie Rath schaffen und helfen. Das ist dazumal unseres HErrn GOttes Meinung und sein Wohlgefallen gewesen, wie eine elende und verdrießliche Gestalt dieser Ansehen es auch gehabt, dermaßen von Gott verachtet und verlassen zu sein.

157. Derhalben sollen wir, so oft wir betrübt und angefochten werden, auf dies Exempel sehen, und sollen bedenken, daß wir unserm Fleische, wo es in Anfechtung zornig und ungeduldig ist, männlich widerstehen, und also sagen: Warum mich jetzt GOtt dermaßen verlassen, weiß ich nicht; daran zweifle ich aber gar nicht, der liebe gnädige Vater wird es aus gutem väterlichem Rath thun, daß es mir selbst auch noch endlich werde nütze und gut sein. Wiewohl es mein sündlich Fleisch nicht sieht, sondern murrt und wider den Geist strebet: so muß doch gleichwohl dies Kreuz mit Glauben und Geduld getragen und überwunden werden; denn ich sehe auch an den heiligen Vätern dergleichen wunderbaren Rath GOttes, damit er sie geführt und regieret. Denn also sind alle diese Drei zugleich dahingestorben, der Großvater Isaak, Jakob der Vater und Joseph der Sohn. Aber siehe und merke darauf, wie viel Gutes unser HErr GOtt daraus gewirkt hat; wie herrlich er doch diesen schweren Fall und Bosheit des Teufels behandelt, daß er diesen armen und gleichsam getödteten Hirten, Joseph, zum Könige in Egypten und zum Herrn vieler andern Völker macht, und was noch das Vornehmste ist, daß damit also die rechte Lehre von GOtt durch dasselbe ganze Königreich ausgebreitet wird; das kommt alles aus der kleinen Versuchung. Derhalben, sage ich, sollen wir selbst auch also schließen: Wohlan, GOtt hat mich jetzt auch mit diesem oder anderem Kreuz beschweret; was liegt daran? Ich will das mit Geduld also leiden. Das Fleisch ist wohl dazu schwach, seufzt, heult und wehklagt: aber GOtt sagt: Du weißt nichts hiervon, bist ein Narr; mein Rath und Wille ist, daß ich aus diesem deinem Kreuz ein groß Gut hervorbringen wolle.

158. Also sagt auch Christus im Neuen Testament von seinen Heiligen und Gläubigen Luc. 18, 7. 8.: „Sollte aber GOtt nicht auch retten seine Auserwählten, die zu ihm Tag und Nacht rufen?“ 2c. „Ich sage euch: Er wird sie erretten in einer Kürze.“ O! wenn wir es glauben könnten, und an diesen Worten Christi, des Sohnes GOttes, festhalten, und einen solchen Geist haben, der mit Freuden sagen könnte: Wohlan, es mag mir mein Leib und Leben, dazu mein Gut genommen werden, es fahre Alles dahin; ich bin deß gewiß, daß GOtt selbst dies thut und solches über mich verhängt, und meint es dazu sehr gut mit mir. Wenn wir also sagen könnten, so wären wir rechte gute Christen und Ueberwinder der ganzen Welt. Und das hieße auch alsdann, sich GOtt recht ergeben in seinen Gehorsam und Wohlgefallen; wie St. Paulus ermahnt Röm. 12, 2.: „Verändert euch durch Verneuerung eures Sinnes.“ Das renovamini (verneuert euch) thut aber wehe. Warum aber das? Folgt: „Auf daß ihr prüfen möget, welches da sei der gute, der wohlgefällige, und der vollkommene GOtteswille.“ Da werdet ihr es mit Leiden, Wehklagen, Seufzen lernen, was GOtt mit euch vorhabe und über euch beschlossen habe.

159. Also ist mir daran gar kein Zweifel, daß nicht das herzliche Seufzen in unsern und vieler gottfeligen Leute Herzen wider den Papst, Türken, und die greulichen Lästerungen der Widersacher, so wir mit großen Schmerzen unsers Herzens hören müssen, da wir seufzen und schreien: Komm, lieber HErr! daß solch Seufzen und Schreien, sage ich, nicht sollte so viel zuwege bringen, daß die Erlösung der gottfeligen Kinder sich bald nahen und herzeilen sollte. Was gilt's, der gemitus (Seufzer) soll die Welt zerreißen. GOtt wird nicht länger können Geduld tragen, spricht Christus Luc. 18, 7. 8., sondern wird seine Auserwählten bald retten. Er wird kommen, er wird nicht ausbleiben; und wenn er wird anfangen zu richten, so wird er eine solche Weise finden und treffen, daß die Widersacher und Lästerer gar zu Grunde vertilgt und seine Kirche gleichwohl noch erhalten werde. Also hat er durch die Sintfluth das ganze menschliche Geschlecht vertilgt, da die Leute am allerschärfsten gewesen und gewißlich gemeint haben, daß nun gar keine Strafe oder Unglück vorhanden wäre, darum daß gleichwohl noch etliche fromme, heilige und weise Leute, wie es denn schien, unter ihnen gewesen. Also ist auch das Königreich Babylon umgekommen; item, Sodom, Jerusalem und Rom sind auch also verheeret und zerstöret worden. Wer hat sie vertilgt und verstört? Das hat dies Seufzen und die Thränen in den Herzen der Heiligen gethan. Mose hat am Rothen Meer auch zum HErrn geschrieen; was ist aber da geschehen? Das Meer hat sich von einander getheilt, und da Pharao hinein gegangen, muß er darin ersaufen; das Volk Israel aber wird wunderbarer Weise erlöset, 2 Mos. 14, 15. ff.

160. Derhalben wäre es sehr nütze und gut, wenn wir lernen und uns daran gewöhnen könnten, daß wir uns auch in unserm Jammer und Trübsal mit solchen Exempeln trösteten, und daß wir daraus gewißlich schließen könnten, daß aus unserm Jammer und Unglück GOtt der HErr einen großen Haufen Güter über uns sammeln und kommen lassen wolle, und daß er solches auch wohl thun könne. Lasset uns nicht murren und zürnen wider GOtt, wie die Juden in der Wüste gethan, davon der 106. Psalm V. 13. sagt: „Aber sie vergaßen bald seiner Werke, sie warteten nicht seines Rathes.“ Er hatte vor, er wollte ihnen helfen, aber sie konnten es nicht erwarten, sie wollten es ohne Verzug haben. Darum blieben sie in der Wüste, sind in das Land Canaan nicht gekommen; denn sie verachteten das liebe Land, glaubten des HErrn Wort nicht und murrten in ihren Hütten.

161. Auf diese Weise ist nun der heilige Patriarch Jakob auch in großer Betrübniß und höchster Angst gewesen, und hätte sich nicht können da herauswinden, wo er nicht Trost gehört hätte von seinem alten Vater Isaac und den andern Hausgenossen, welche ihn der göttlichen Verheißung werden erinnert und ihn damit wiederum aufgerichtet haben, bis er seinen Sohn Joseph wieder bekommen hat, da derselbe nun zu großer Ehre und Herrlichkeit erhoben. Denn das ist ja ein sehr hoch und groß Ding, daß Joseph also zur Herrschaft über ganz Egypten erhoben wird; und wo dasselbe GOtt dem Jacob zuvor geoffenbart und gesagt hätte: Ich will verschaffen, daß der Verlust deines Sohnes dir eine Ursache vieles großen Nutzens und Frommens sein soll. So hätte er es gewißlich mit gutem Willen und mit Freuden den Sohn lassen hinfahren. Aber GOtt hat ihm solches verborgen und Alles vor ihm verhehlet, daß er also den alten Menschen erstlich tödten, und darnach den neuen Menschen wiederum erfreuen und trunken machen und mit neuen unzähligen Wohlthaten erfüllen möge.

\quellseite{284}

162. Inzwischen aber hat es sich gebühret, das Herz zu erwecken zum Glauben, zur Hoffnung, auf göttliche Hülfe zu warten, nach der Lehre, so im 27. Psalm V. 14. steht: „Harre des HErrn, sei getrost und unverzagt, und harre des HErrn.“ GOtt weiß wohl, was die Versuchung für ein Ende und Ausgang gewinnen soll, welchen du selbst nicht sehen kannst. Halte fest; ich will es so mit Zucker durchmengen, daß wenn du gleich wirst sterben müssen, sollst du doch in der Auferstehung der Todten wiederum mit einiger Freude getröstet werden.

163. Denn dazu dienen diese Uebungen, daß wir daran GOttes Gnade und Barmherzigkeit, dazu das Geheimniß des Glaubens und rechter Hoffnung verstehen lernen, und daß wir daraus auch etlichermaßen den unversöhnlichen Rath GOttes über uns begreifen lernen; wie im 4. Psalm V. 4. steht: „Der HErr führet seine Heiligen wunderlich“; item Ps. 16, 3.: „Ich habe an den Heiligen all mein Gefallen“; desgleichen im 17. Psalm V. 7.: „Beweise deine wunderliche Güte, du Heiland derer, die dir vertrauen.“ GOttes Gnade und Barmherzigkeit, und seine Verheißung ist gewiß vorhanden; sie wird aber wunderlich erfüllt über alle unsere Gedanken, Willen und Begierde.

164. Derhalben soll man mit dem lieben Hiob Cap. 13, V. 15. sagen: „Wenn er mich schon tödtet, will ich doch auf ihn hoffen“; und mit dem heiligen Paulus 2 Cor. 4, 8. 9., da er als ein geschickter Redner ganz herrlich von dieser Sache redet: „Wir haben allenthalben Trübsal, aber wir ängsten uns nicht. Uns ist bange; aber wir verzagen nicht. Wir leiden Verfolgung; aber wir werden nicht verlassen. Wir werden untergedrückt; aber wir kommen nicht um.“ Unser Leben ist ganz armselig und elend, ist voll großes Unglückes und Trübsal; aber wir werden doch nicht verlassen. Und also denkt und schließet auch ein christlich Herz in der Noth: Ich weiß, daß ich nicht zu Schanden werden soll, meine Hoffnung und mein Glaube und Vertrauen zu GOtt wird mir nicht fehlen. Denn GOtt selbst, seine Verheißung, Sacrament und Absolution können mir nicht fehlen; denn sie sind auch GOttes selbst. Derhalben wenn ich schon sterbe, werde ich doch wieder leben; wenn ich gleich umkomme, werde ich doch erhalten und selig werden.

165. Denn also wird Jacob auch verworfen, wie es scheint, wird aber doch nicht verlassen; er wird angefochten, verzagt aber gleichwohl nicht; und endlich bleibt er noch gleichwohl GOttes Kind und GOtt bleibt in ihm, aber ganz verborgen und wunderbarer Weise. Und zuletzt, da dies Spiel zum Ende kommt, da Jacob sieht, daß sein Sohn noch am Leben und so ein großer Regent geworden, ist es ihm dazu sehr lieb und freut sich nun auch, daß er zuvor also umgetrieben und so viel Unglück erlitten hat; denkt also: Siehe, solche große Freude wäre mir nie widerfahren, wo mein Haus und Gemeinde nach meinem Rath und Willen regiert worden wäre; Joseph wäre ein Hirte geblieben, wie meine andern Söhne: nun aber wird er zu königlicher Herrlichkeit erhoben und werden durch ihn viele andere Menschen auch erhalten, daß ihnen an Leib und Seele geholfen wird.

166. Wenn wir nun betrübt, angefochten und auch also umgetrieben werden, sollen wir unser Herz wider das Fühlen des Unglücks erwecken und also sagen: „Ich werde nicht sterben, sondern leben“, Ps. 118, 17., ob sich schon das Gegentheil sehen läßt. Und wiewohl ich an mir selbst ganz verzagt muß, will ich doch auf den hoffen, der Alles aus Nichts macht; und wenn ich gar zunichte geworden, kann er mir wohl wieder aufhelfen und dasselbe mir und Andern auch zum Allerbesten.

167. Derhalben, je grausamer das Leiden ist, je größer und wunderbarer Ding es in den Heiligen und Gläubigen wirkt; und ist ein gewisses Anzeichen göttlicher Gnade und Güte, wenn die Gläubigen mit Kreuz und Widerwärtigkeit angefochten und beschweret werden. Denn wo sie im Glauben an die göttlichen Verheißungen beharren und beständig bleiben, folgen bald große unglaubliche Güter hernach; wie der Spruch des heiligen Jacobus Cap. 1, V. 12. lautet: „Selig ist der Mann, der die Anfechtung erduldet; denn nachdem er bewähret ist, wird er die Krone des Lebens empfangen, welche GOtt verheißen hat denen, die ihn lieb haben.“ Und Joh. 12, 24. spricht Christus: „Es sei denn, daß das Weizenkorn in die Erde falle, und ersterbe, so bleibts alleine; wo es aber erstirbt, so bringts viel Früchte“: daher daß das Korn verfault und stirbt, wächst bald und kommt hervor ein schöner Halm, der darnach hier dreißigfältige, dort sechzigfältige Frucht bringt.

168. Dies ist eine solche Lehre, so für und für in der heiligen Schrift getrieben wird; und ist auch GOttes Wille, daß wir nach dem Fleische getödtet und nach dem Geiste wiederum lebendig gemacht werden sollen. Denn dasselbe meint St. Paulus auch, da er 2 Cor. 4, 11. sagt: „Denn wir, die wir leben, werden immerdar in den Tod gegeben um JEsu willen, auf daß auch das Leben JEsu offenbar werde an unserm sterblichen Fleische.“ Item, daselbst bald hernach V. 16. 17. 18. spricht er: „Darum werden wir nicht müde; sondern ob unser äußerlicher Mensch verweset, so wird doch der innerliche von Tage zu Tage verneuert. Denn unsre Trübsal, die zeitlich und leicht ist, schaffet eine ewige und über alle Maße wichtige Herrlichkeit, uns, die wir nicht sehen auf das Sichtbare, sondern auf das Unsichtbare“ \&c.

169. Dies ist also die rechte wahrhaftige und einige Kunst der Christen, dazu sie von Christo berufen und erwählet sind; wie er denn zu seinen Jüngern Joh. 15, 16. sagt: „Ihr habt mich nicht erwählet, sondern ich habe euch erwählet, und gesetzt, daß ihr hingehet und Frucht bringet, und eure Frucht bleibe.“ Darum übt er sie auch also und tödtet sie, auf daß ihre Früchte hernach desto reicher und größer sein mögen.

170. Auf solche Weise ist er auch mit diesen sehr heiligen Patriarchen umgegangen, als kennete er sie nicht; und hat doch gleichwohl seine Gnade und Barmherzigkeit über den Großvater Isaac, über Jacob, den Vater, und Joseph, seinen Sohn, so reichlich ausgegossen; denn sie sind ohne allen Zweifel GOtt sehr angenehm und lieb gewesen. Es wird aber solche Liebe und Gnade GOttes also verborgen und heimlich gehalten, daß GOtt auch nie einen seiner Freunde grausamer und härter gestraft hat, auch nicht die Vatermörder selbst alhier, so Joseph verkauft, noch auch die Ismaeliter oder Midianiter, so ihn gekauft haben.

171. Also erweckt und gewöhnt er nun seine Heiligen und Gläubigen, seinen wunderbaren Rath zu bedenken; den andern Haufen aber läßt er wachsen und groß werden, in ihren Lüsten leben, allerlei Büberei und Räuberei in aller Sicherheit zu treiben, läßt ihnen allen ihren Willen, daß sie thun, was sie nur gelüstet, züchtigt sie nicht, hindert oder wehret ihnen nicht, was sie Ungöttliches in höchster Sicherheit vornehmen. Aber, o welch ein betrübtes und unseliges Nachsehen ist das, wo ihnen GOtt also zusieht, daß sie thun, was sie nur wollen!

172. Dagegen sind die viel seliger, welche GOtt in Betrübniß, Elend und den Tod hineinwirft, und denen er doch dabei seine Verheißung gibt und sagt: Du hast meine Verheißung, daran laß dir genügen. Wo ich dich verwerfen oder verstoßen hätte, würde ich dir meine Verheißung nicht geben und würde dich nicht heißen mir vertrauen; wenn ich wollte, daß du ewig verdammt sein solltest, würde ich dich nicht absolviren und von deinen Sünden los sprechen, nicht taufen, und würde dich nicht zur Gemeinschaft meines lieben Sohnes JEsu Christi berufen. Dieweil du aber das heilige Sacrament und Evangelium hast, so gedenke an meine Worte Matth. 5, 4.: „Selig sind, die da Leid tragen; denn sie sollen getröstet werden“; item V. 11. 12.: „Selig seid ihr, wenn euch die Menschen um meinetwillen schmähen und verfolgen, und reden allerlei Uebels wider euch, so sie daran lügen. Seid fröhlich und getrost, es wird euch im Himmel wohl belohnet werden“ \&c.

173. Es hat aber Joseph nicht allein himmlische und ewige Belohnung, sondern auch zeitlichen Lohn empfangen, der auch alle Hoffnung und Gedanken weit übertroffen. Denn er ist durch diese kurze und leichte Anfechtung ein Vater des Königes in Egypten worden, und der auch zugleich ein Bischof und Lehrer vieler Völker gewesen. Dies ist auch eine leibliche und treffliche große Belohnung der Geduld und Beständigkeit in Trübsal und Widerwärtigkeit. Und also pflegt GOTT sein Angesicht zu verbergen und seine Hand vor seinen Heiligen und geliebten Kindern zu entziehen, läßt sie Leid tragen, verkauft und in das Gefängniß geworfen werden und sterben; nicht anders, als ob sie GOttes Feinde wären, auf welche er zwar mehr denn auf diese seine Geliebten Acht gibt; nicht daß er in Wahrheit also gesinnet, wie es äußerlich scheint, sondern daß er gleichsam zu diesem Spiele in den Werken seiner Hände Lust hat. Denn also sagt er Offenb. 3, 19. selbst: „Die ich lieb habe, die züchtige ich auch.“

\quellseite{285}

174. Dies sollen wir fleißig merken und betrachten, daß nachdem wir nun angefangen haben zu glauben und zu hoffen auf das ewige Leben, welches uns durch das Wort und die heilige Taufe schon gegeben ist, wir auch alles Unglück, so uns darnach begegnen wird, mit Geduld ertragen lernen; und daß wir es gewißlich dafür halten, solches alles geschehe nicht ohngefähr oder vergeblich, sondern aus väterlichem Rath und Willen GOttes zu unserm Heil und Seligkeit; welches wir dann allererst verstehen werden, wenn die Anfechtung ein Ende haben wird. Mitten in der Noth, wenn der Kampf mit der Anfechtung am allerheftigsten ist, versteht man es nicht; denn das Fühlen unsers Fleisches reißt uns ab von der Verheißung. Und kommt das dem Menschen, so in Angst und Noth ist, nicht in den Sinn, daß er gedächte: Ich bin dennoch getauft, habe die Verheißung GOttes: sondern da wird das Herz ganz überladen mit Klagen, Seufzen und Weinen nach dem Fleische.

175. Nun muß aber das Fleisch gekreuzigt und getödtet werden; denn es hindert uns, daß wir die Verheißung und die Wahrheit des Wortes GOttes, so in der Anfechtung gespüret wird, wo man daran mit festem Glauben hält, nicht recht verstehen können, und alsdann folgt auch, das im 34. Psalm V. 9. steht: „Schmecket und sehet, wie freundlich der HErr ist.“ GOTT läßt uns also angefochten werden, daß er Raum habe, uns darnach wiederum zu sättigen, zu trösten und trunken zu machen, wenn wir ganz kraftlos und erschöpft sind, bei uns selbst gar keine Hülfe oder Trost zu finden ist. Denn sonst sind wir sehr schläfrig. Und wenn wir also GOttes Wort mit Ueberdruß hören, werden wir darüber ganz faul und sicher. Immer macht einen, spricht man. Der Hunger ist ein guter Koch.

\section[Vierter Theil]{Vierter Theil. Von den blutigen Aufschlägen der Söhne Jakobs über ihren Bruder Joseph}

\begin{quote}
\textbf{V. 18—20. Als sie ihn nun sahen von ferne, ehe denn er nahe bei sie kam; schlugen sie an, daß sie ihn tödteten, und sprachen unter einander: Sehet, der Träumer kommt daher. So kommt nun und lasset uns ihn erwürgen, und in eine Grube werfen, und sagen, ein böses Thier habe ihn gefressen; so wird man sehen, was seine Träume sind.}
\end{quote}

176. GOTT schweigt gar stille zu den ungöttlichen und blutigen Anschlägen, bösem Rath und Vornehmen dieser Brüder Josephs, zerstört oder hindert dieselben nicht, ist blind und taub, wie es scheint, sieht und hört nicht, versteht auch nicht, womit sie umgehen, und hat doch gleichwohl zu Jakob gesagt: Ich will dein Schuh sein.

177. Nun ist es wahrlich ein großer Jammer, daß man in dem Hause und der Gemeinde, welche damals allein in der ganzen weiten Welt die heiligste gewesen, solche schändliche und abenteuerliche Frucht sehen soll, in welcher Gemeinde der hohe große Patriarch Isaac Pfarrherr und Lehrer gewesen, GOttes Wort und Verheißung von Christo sonderlich im Schwange gegangen ist, der Heilige Geist regieret, zudem auch viel schöne Exempel der Gottseligkeit, der Zucht und eines feinen häuslichen Wandels geleuchtet, daß dennoch daraus solche erschreckliche Ungeheuer gekommen. Denn Jakob ist mit seiner ganzen Kirche und Gemeinde dem Teufel und allen seinen Engeln unterworfen, in dem Grade, daß nicht fremde Leute, nicht die Feinde oder Ketzer, sondern seine eigenen Kinder, von seinem Leibe geboren, die er fleißig auferzogen und in GOttes Wort unterrichtet, ihren armen unschuldigen Bruder getödtet haben.

178. Dies ist wahrlich ein wunderbärlich Ding, welches so hoch und schwer ist, daß es niemand, wie beredt er auch sein mag, ausreden, ja, kein Mensch ausdenken kann. Mich dünkt, es sei noch ärger denn Kain und sein Bruder, und daß diese Brüder weit grausamer denn Kain, der Brudermörder, gewesen sind. Denn so viel Brüder haben den einigen Jüngling Joseph, der noch nicht die Jahre seiner rechten Jugend erreicht, so gar unbilliger und unmenschlicher Weise geneidet und gehaßt, und das ohne alle Ursache oder auch ohne alle seine Schuld. Denn, lieber GOtt, was kann doch das, daß er zwei Träume gehabt, daß er seinem Vater sonderlich lieb gewesen, und vor denselben gebracht hat, wo ein böses Geschrei wider seine Brüder war, für eine Sünde sein? Da ist doch nichts an ihm gewesen, denn eitel Unschuld und Gottseligkeit, ist schlecht und recht und ohne alle Falschheit gewesen; welches ja alles rechte Ursachen sind, daß sie selbst ihn billig auch sollten geliebet haben. Was kann da dafür, daß er Träume gehabt hat? etc. Dennoch sind diese Brüder, so in solchem heiligen Hause des Patriarchen Jakob geboren und auferzogen, die er dazu ernstlich in der Zucht gehalten, also bezaubert gewesen, daß sie diesen ihren Bruder, der so fromm und gottselig war, und zugleich mit ihm auch ihren alten Vater getödtet haben. Das sind eigentlich Vatermörder, Mutter- und Brudermörder, von denen Paulus 1 Tim. 1, 9. sagt.

179. Der gute alte Vater Jakob sitzt zu Hebron, ist in seinem Herzen wohl zufrieden, denkt nicht anders, denn daß alle Sachen wohl stehen. Da er nun in solcher Stille und guter Hoffnung sitzt, da kommt ihm sein lieber Sohn Joseph durch seiner Brüder Haß um und er selbst, der Vater, kommt auch also mit um. Was sollen denn wir armen elenden Leute in diesem Leben nicht auch zu erwarten haben, sintemal auch von Hausgenossen, von Kindern und von unserm eignen Fleisch uns mehr Schaden und Unglück widerfahren kann, denn von Nachbarn oder Fremden? Zacharias, Luc. 1, 71., singt in seinem Lobgesange also: „Daß er uns errettete von unsern Feinden“; aber Jakob möchte hier wohl lieber von seinen liebsten Freunden und Kindern Errettung wünschen und bitten. Denn solcher gestalt wird unser Leben beschrieben und abgemalt, daß es der Eitelkeit, dem Teufel und allem Unglück unterworfen sei, so gar auch, daß wir uns dessen auch von unsern lieben Freunden, Kindern und Hausgenossen zu befürchten haben, daß sie uns eine Ursache des Todes oder steinigen Wehklagens sein werden; welches noch beschwerlicher ist denn der Tod selbst.

180. Denn also haben ja die Kinder Jakobs die aller trefflichsten Lehrer der Kirche und Propheten gehabt, von denen sie täglich GOttes Wort, Verheißung, Gesetz und Anderes gehört haben, an denen sie immer gute Sitten und sehr viele Exempel aller Ehrbarkeit gesehen; und haben doch diese giftigen Spinnen aus solchen schönen Blumen und seinen Rosen nichts Anderes denn Gift gesogen. Und da sie sich ja weder durch Vermahnung noch Exempel hätten bessern wollen und von solcher greulichen Sünde abschrecken lassen, sollten sie doch ihres frommen Vaters und heiligen Großvaters damit billig verschont haben. Aber der Gott dieser Welt hat sie verblendet, daß sie weder hören noch sehen können.

181. Dazu ist auch noch gekommen das große schwere Aergerniß, so aus diesem ihrem Vater bei den Nachbarn, den Ismaeliten, Midianitern und Andern entstanden, zu welchen dies Geschrei hat kommen müssen. Denn diesen Leuten ist Joseph von seinen Brüdern verkauft worden, darum werden sie ohne Zweifel gedacht haben: Siehe, dies sind die Kinder des Segens, wir aber sind aus der Kirche oder Gemeine Abrahams gar verstoßen, sind wie die Verworfenen, und die fremd sind von der Verheißung und Kirche GOttes; diese aber sind die schönen Erben dazu, welche diese große Sünde begangen.

182. Dies ist wahrlich ein sehr groß Aergerniß in ihren Herzen gewesen, daraus sie auch nichts Anderes haben schließen können, denn daß der Ruhm vom Segen, Abraham verheißen, müßte vergeblich sein, und daß sie mit den Ihren vielmehr die rechte Kirche wären, als darin ja solche grausame Exempel so schändlicher Sünde nicht gefunden würden. Denn an den Früchten soll man den Baum erkennen; und war auch das Aergerniß, so zuvor daraus genommen, daß Dina geschwächt worden, nicht viel geringer. Derhalben solche Aergernisse auch ohne allen Zweifel Jakob selbst mehr betrübt denn das Unglück an sich selbst; wie zu unsern Zeiten uns auch mehr grämt und wir darüber auch mehr wehklagen, daß in unsern Kirchen auch, leider! so viele Aergernisse erweckt werden, denn über der Verfolgung der Papisten und Tyrannen. Denn die Papisten thun der Kirche so großen Schaden nicht, als die falschen Brüder, so sich auch desselben Glaubens, Lehre, Hoffnung und Anrufung mit uns rühmen, brauchen auch mit uns einerlei Sacramente und sind in dem allen uns fast gleich. Darum werden die Papisten auch bewogen und verursacht aus solchen Aergernissen, daß sie unser spotten und sagen: Es sind Buben in der Haut, sie beißen sich unter einander, werden auch gewißlich also unter einander verzehret werden.

\quellseite{286}

183. So ist nun Jakob ein Exempel der Geduld und Gedanken, so man von dem wunderbaren Rath Gottes über uns haben soll; und hat Gott haben wollen, daß dies hat sollen geschrieben und in der Kirche gelesen und gelehrt werden, daß wir daran lernen mögen, wie er seine Heiligen so wunderlich führe, und dieselben also regiere, daß auch die Heiligen selbst solchen Rath Gottes, dadurch sie geführt werden, nicht verstehen oder begreifen können; und daß um keiner andern Ursache willen, denn daß unser Fleisch, Sinn und Verstand, und dazu auch unsere Weisheit also getödtet werde, und daß wir uns daran gewöhnen, daß wir einfach die Augen zuthun, und den Verheißungen Gottes vertrauen, wenn er sich schon auch also stellt, als achte er unser nicht, ja, wenn sich auch ganz das Gegenteil sehen läßt. Wie die Braut im Hohenliede Salomos Cap. 2, 9. sagt: „Er stehet hinter der Wand“, er verkriecht sich und versteckt sich hinter das Tuch; und wie das Kindlein Jesus, da es zwölf Jahr alt war, sich auch seinen Eltern entzogen, Luc. 2, 43. Eben also verhält er sich auch noch gegen die Gottseligen und Frommen, daß es sich ansehen läßt, als kenne er uns gar nicht, auch da er uns zuvor große Verheißung gethan; und nachdem er mit uns gewißlich seinen Bund aufgerichtet, läßt er uns doch also umkommen, als hätte er seiner Verheißung gar vergessen, und als achte er unser nicht und hätte uns gar verworfen.

184. Dies ist nun der Christen Weisheit, solchem göttlichen Rath stille halten, und im Glauben auf die Verheißung, so er uns gegeben, beharren. Denn es ist ja solche Verheißung gewiß und beständig genug, und der Bund des Herrn besteht auch fest, wie der 121. Psalm V. 4. sagt: „Der Hüter Israel schläft noch schlummert nicht.“

185. Aber darauf antwortet menschliche Vernunft also: Dies alles ist wohl herrlich und fein geredet, ich erfahre und empfinde aber gleichwohl das Gegenteil. Denn er schläft nicht nur, sondern schnarcht auch noch dazu; ja, es ist da gar kein GOTT vorhanden, der für uns sorge oder Acht auf uns gäbe. Also ist nun Jakob der göttlichen Verheißung gewiß, weiß auch, daß der Bund Gottes fest bestehen wird; und geht doch Gott gleichwohl mit ihm also um, daß es scheint, als habe weder er selbst, noch sein Sohn Joseph solche heiligen Engel zu Hütern, die dem grimmigen Zorn seiner Brüder widerstehen möchten: da schweigen sie alle still, und lassen den Teufel auf die heilige Kirche und Gemeinde einstürmen, wider dieselbige toben und wüthen. Wo ist denn nun GOTT?

186. Derhalben werden wir nun dessen oftmals erinnert, daß wir aus solchem Exempel gelehret werden, daß man die göttliche Verheißung mit festem Glauben ergreifen müsse, und daß man an GOtt, wo er etwas verheißt, gar nicht zweifeln solle. Denn wie GOtt nicht lügen kann, also kann er es auch nicht lassen, er muß für uns sorgen, vornehmlich, wo wir an seiner Verheißung fest halten. Denn wo dieselbige fest ergriffen wird, ist es unmöglich, daß wir sollten verlassen werden, darum daß ja GOtt wahrhaftig ist. Derhalben, wo er es verhängt, daß wir angefochten, in die Hölle geführt und getödtet werden, wie wir an diesem Exempel Jakobs sehen, sollen wir allezeit zu seiner Verheißung Zuflucht nehmen, und weit aus den Augen hinwegthun das erschreckliche Aergerniß, damit wir also gekreuzigt und geplagt werden.

187. Jakob und sein Sohn Joseph werden sehr hart versucht und beschweret, welches ja der göttlichen Verheißung ganz und gar zuwider gewesen, und scheint nichts weniger, denn daß GOtt ihrer etwas achten und für sie sorgen sollte. Denn er sendet ihnen ja keinen Engel, ja, auch nicht ein Blatt vom Baume, damit der Teufel möchte aufgehalten und zurück getrieben werden; sondern thut ihm noch wohl Thür und Fenster auf, daß er nur auf das greulichste gegen Vater und Sohn toben und wüthen möge.

188. Es hat wohl Jakob sehr große und reiche Verheißung, dazu auch derselbigen Erfüllung und lebendige Erfahrung, nämlich, in dem Segen, so dadurch ihm zeigt, daß er zwölf Söhne bekommen hat. Aber er muß gleichwohl noch immer mit vielen Beschwerden, so ihn vorkommen, kämpfen, als ob er keinen GOTT oder auch keine Verheißung hätte. Jedoch müssen wir schließen und es noch dafür halten, daß er gleichwohl von GOtt mit nichten verlassen gewesen. Denn es steht die Verheißung noch immer fest, Jes. 49, 15. 16.: „Kann auch ein Weib ihres Kindleins vergessen, daß sie sich nicht erbarme über den Sohn ihres Leibes? Und ob sie desselben vergäße, so will ich doch dein nicht vergessen. Siehe, in die Hände hab ich dich gezeichnet“ etc.

189. Warum begegnet uns aber dennoch so viel erschreckliches Dinges von unserm eigenen Fleisch, unsern Kindern und Geschlechte? Antwort: Dies ist also die Weise göttlicher Regierung und also ist es auch beschaffen um das Leben der Heiligen und Gläubigen in dieser Welt. Derhalben gehört dazu eine solche Weisheit und Lehre, so alle Vernunft und menschlichen Verstand übertrifft, daß ich durch solche Weisheit sagen könne: Ich bin getauft, bin von meinen Sünden absolvirt und ledig gesprochen; ich habe meines Herrn Christi Leib gegessen und sein Blut getrunken; ich habe das gewisseste Wort Gottes, der wird mir ja nicht lügen oder mich betrügen, wenn es sich auch schon ansehen ließe, als wollte Alles miteinander widrig hinausgehen.

190. Denn also hat Jakob auch gedacht: Es scheint zwar, als ob ich gar verlassen sei, dieweil mir mein Sohn genommen, von dem ich sonderlich die Verheißung, welche jetzt in Zweifel gestellt wird, verstanden; doch ist es nur eine Versuchung. Aber wie viele, meint du, sollte man deren wohl finden, die daselbige auch also thun könnten?

191. Darum, sage ich, ist es vonnöthen, daß unser Fleisch getödtet und der Verstand des Fleisches und Vernunft mit aller menschlichen Weisheit gar zunichte gemacht werde. Es muß doch dahin kommen. Alles ist durch das Wort gemacht, muß auch dadurch wiederum gemacht werden; wir sind durch das Wort geschaffen, müssen uns auch wieder zum Wort kehren und weiden. Also haben die Sophisten auch geredet, haben es aber selbst nicht verstanden: Ad principium, a quo processimus, est redeundum, das ist: Wir müssen wiederkommen zu dem Anfang, da wir hergekommen sind. Das ist wohl leichtlich also geredet, daß man davon speculiren könne; aber wo man es mit der That üben soll, und also zunichte werden, sterben und gar zu Grunde gehen, daß es scheint, als sei nun nichts mehr übrig vom Leben oder vom Gefühl im Fleisch denn nur das Wort, das hat Mühe und Arbeit. Wenn ich sterbe, so fahre ich hinunter in die Hölle und komme gar um. Was soll ich alsdann thun? Da ist keine Hülfe mehr übrig denn allein das Wort: „Ich glaube an GOtt“ etc. Daran halte ich mich fest, wie heftig er auch zürnt, mich verläßt, tödtet, ja, in die Hölle hineinführt. Warum halte ich mich aber an das Wort? Darum, daß ich getauft und absolviret bin, habe dazu gebraucht die heilige Communion. Dem Wort glaube ich. GOtt gebe, es breche Himmel und Erde, so muß man doch das Wort, die Verheißung und Sacramente deshalb nicht hinwerfen oder verleugnen, wenn ich auch schon in die Hölle sollte geworfen werden.

192. Dies sind aber nicht solche Dinge, davon man nur vergeblich speculiren möge, sondern sind mitten aus dem Leben und Erfahrung genommen: und man soll sie nicht nur einmal, dieweil wir leben, hören oder bedenken, sondern dieselbigen oftmals wiederholen und in Uebung bringen. Ich bin einmal getauft, bin angenommen in göttliche Verheißung, werde täglich absolviret, und höre, daß mir Vergebung der Sünden verkündigt wird, werde mit GOttes Wort aufgerichtet und getröstet, bin im Reich der Gnade und Seligkeit. Was geschieht aber? Da begegnet mir täglich so viel und so großes Unglück und Wunder, daß wo ich keine andere Hülfe oder Rath hätte, denn meine eigene Vernunft und Fleisch, würde ich mich bald mit einem Stricke oder sonst mit anderer Wehr selbst umbringen. So voll unzähliger List ist der Teufel, und sind alle Winkel voll mancherlei Art des Todes, der Nachstellung, des Schwerts etc. Doch muß man, wie viel auch solcher Wolken und Finsterniß im Wege liegen, gleichwohl in das Licht des Wortes und göttlicher Verheißung hinein brechen, und also sagen: Nun, im Namen Gottes, ich habe gesündigt, das bekenne ich; doch glaube ich, daß ich getauft bin und daß mein Herr Christus zur Rechten Gottes sitzt etc.

\quellseite{287}

193. Dies ist das Seufzen, davon Gott haben will, daß es in den Herzen der Gläubigen soll erweckt werden, damit sie nicht sicher und träge werden, und von Ueberdruß und geistlicher Faulheit umkommen und verderben. Denn also wird der Geist erweckt, der Glaube geschärft; es wächst die Erkenntniß Gottes und wird der neue Mensch immer von Tag zu Tag verneuert, und man lernt darüber, welches da sei der gute vollkommene Wille Gottes.

194. So ist es dem Mann auch gegangen. Denn es wird nun dieser heilige Mann sterben, und ist dies die letzte und äußerste Anfechtung gewesen, welche ihn auch ohne Zweifel wird getödtet haben. Denn er hat sonst sein voriges Leben in höchster Bekümmerniß, in Finsterniß, Unglück und Schaden des Todes zugebracht. Wie denn ein solches Leben, so in stetigem und betrübtem Wehklagen zugebracht wird, kein Leben, sondern der Tod ist, da die Sonne schwarz wird, da weder Schlaf, noch Essen oder Trinken schmeckt, da man zu den Dingen keine Lust hat, die doch sonst gemeiniglich den Menschen zu erfreuen pflegen. Das ist wahrlich der rechte und stetige Tod; darum geht nun hier Jakob zu Grunde. Es hat ihn allhier die letzte Anfechtung getroffen und betreten, und dieselbe hat ihn auch getödtet; denn er hat seinen allerliebsten Sohn verloren, welcher der Erstgeborne war von seinem ersten Weibe. Und wird ihm auch zugleich damit erneuert sein die schwere Betrübniß, so er zuvor darüber gehabt, daß die Sichemiter erschlagen und daß ihm seine Tochter Dina geschwächt worden. Und ist bei ihm der Schmerz auch noch neu gewesen, so er über den kläglichen Tod der Mutter Josephs empfangen, und zuletzt, daß Ruben Bilha geschwächt und damit die greuliche Untugend angerichtet hat. Diese vier schweren Anfechtungen sind einander bald gefolgt, immer eine nach der andern innerhalb drei Jahren, ohne Aufhören. Das ist ja ein armer Mann, in dem Hause, so doch ganz heilig gewesen und da Gottes Wort und der Heilige Geist regiert, gleichwohl noch so große Gewalt des Teufels die Oberhand haben soll, daß nicht leichtlich auch einem schändlichen Heiden größer Herzleid könnte widerfahren sein. Ich wüßte nicht, wie es Einem übler gehen sollte, der ein Heide wäre; und wenn solches dem Julius Cäsar, dem Octavius, oder Alexander dem Großen begegnet, wäre es wahrlich erschrecklich und grausam genug.

195. Dies sind allerdings große erschreckliche Dinge, welche ich auch nach ihren Würden nicht handeln kann: es ist mir zu hoch. Darum sollen wir uns nicht fürchten; es komme, was da wolle, so kann es uns doch so übel nicht gehen als dem heiligen Manne. Nein, sage ich, härter und beschwerlicher kann es uns nicht gehen. Wir sind ja noch nicht zu solchen schweren Leiden und Kreuzigungen gekommen, wie dieses Patriarchen Leiden gewesen. Wenn schon der Türke käme, Alles raubte und verwüstete, und uns dazu auch noch ermordete, würde doch solcher Jammer noch nirgendes mit dem Jammer und Trübsal zu vergleichen sein, so Jakob erlitten. Ja, über das ist auch die große Qual oder Pein, so die lieben Märtyrer erlitten, noch damit nicht zu vergleichen. Denn sie haben gelitten, als die voll Heiligen Geistes und voller Freuden, ohne Wehklagen und Traurigkeit gewesen, haben über das auch den Tyrannen noch trotzen können und ihrer spotten. Ihr Leiden hat einen kleinen Augenblick gewähret; hier aber ist eine stetige Tödtung gewesen, so für und für ganze zwei und zwanzig Jahre gewähret, bis er nach Egypten hinab gezogen.

196. Also ist nun Gottes Rath sehr wunderbarlich, dadurch die Heiligen oder Gläubigen regiert werden, da der Satan greulich tobt und wüthet, und doch der Heilige Geist daselbst auch zugleich mit gegenwärtig ist, der die heilige christliche Kirche bewahrt, lehrt, ermahnt und wider den Teufel erhält. Denn wiewohl allhier beide Väter und Sohn verlassen werden, haben sie doch gleichwohl, wie zuvor gesagt, das Wort noch ganze zwanzig Jahre behalten, in welchem sie alle beide auch mitten im Tode noch gelebt haben, das ist, nur in der Verheißung und im Wort.

197. Das ist aber ein sehr erbärmlich Leben nach dem Fleisch. Denn der arme Vater stirbt in Wahrheit dahin, sieht vor seinen Augen nichts, denn den Tod seines Sohnes und auch sein selbst. Darum sagt er auch darnach V. 35: „Ich werde mit Leide hinunter fahren in die Grube.“ Joseph aber wird über das noch dazu in das Gefängniß geworfen und des Ehebruchs angeklagt, als wäre es nicht Unglück genug gewesen, daß er im Elende und Dienstbarkeit leben, und also von seinen allerliebsten Vater und ganzem Geschlechte weggerissen und von ihnen abgesondert sein mußte. Also ist Tod über Tod, Hölle über Hölle über ihn gefallen; und Joseph lebt gleichwohl noch. Warum oder wie hat er denn gelebt? Ei darum, daß er die Verheißung Gottes und die Predigten seines Vaters gehört hat: Gott hat Abraham, Isaak, meinem Vater, verheißen, er wolle ihres Samens und Kinder Gott sein ewiglich. Diese Verheißungen haben Jakob, den Vater, und dazu seinen Sohn, Joseph, erhalten und aufgerichtet. Denn Joseph ist wie ein guter Zunder gewesen, der leichtlich Feuer fängt, und ist wie von seines Vaters Worten abgewichen: er hat sich nach seiner Lehre fleißig gehalten. Ich, wird er gedacht haben, bin Jakobs Sohn, dem Gott verheißen, er wolle für mich sorgen; das glaube ich fest, wiewohl Alles wider sinnig geht.

198. Aber solch Leben und Exempel der lieben Väter verstehen die elenden Papisten gar nicht. Sie hören wohl, daß Jakob ein Weib gehabt, dazu auch Kühe und ander Vieh, und solches verachten sie gar als eitel fleischlich Ding. Wohlan, lieben Herren, wollet doch dies Exempel etwas genauer besehen und betrachten, nämlich, welch ein armer elender Mensch Jakob nach dem Fleisch gewesen sei, und schauet doch auch daneben an die schöne Lampe, so in dem finstern Ort Joseph und Jakob vor den Augen scheint. Denn daran halten sie sich: daher kommt es, daß sie so lange Zeit, ganze zwanzig Jahre, im Elende und selbst in der Hölle gelebt haben. Wiewohl Joseph zehn Jahre zuvor aus der Hölle erlöset worden, hat aber doch inzwischen der lieblichen Gemeinschaft seines lieben Vaters beraubt sein müssen. Derhalben soll man in den lieben Vätern die Verheißungen Gottes und den Glauben an dieselbigen ansehen, daß sie solchen großen Jammer und Elend haben tragen und dulden können, durch seines Andern Hülfe, denn allein, daß sie dies Wort gehabt: Ich, der Herr, bin dein Gott, dein Schutz, fürchte dich nicht, setze nur dein Vertrauen auf mich, ich bin mit dir, ich will dich wohl schützen und retten.

199. Wie uns auch in der heiligen Taufe, Absolution und Communion gesagt wird: Ich bin der Herr, dein Gott; bekümmere dich nicht, ich will Rath schaffen; wirf nur deine Sorge und Anliegen auf mich; du hast ja an mir einen solchen Gott, der dir verheißen, er wolle väterlich für dich sorgen. Ja, sagt ein armes bekümmertes Herz, ich sehe gleichwohl das Gegentheil. Da antwortet dann Gott wiederum also: Es ist wahr, liebes Kind, du siehst ja das Gegentheil; aber das ist nur eine Anfechtung, die dir dazu nütze und gut ist, daß du daraus lernen und erfahren mögest, wie freundlich der Herr sei. Denn wo diese Anfechtung nicht gekommen wäre, würdest du in deinem Fleisch ganz grob und unverständig bleiben, würdest nimmermehr verstehen können, was da sei, daß ich gesagt habe: Ich bin der Herr, dein Gott. Also mußt du, liebes Kind, gezüchtigt und gelehret werden, daß du durch solche Uebung und Erfahrung so mancherlei Anfechtungen erkennen lernest, daß ich der Herr, dein Gott, sei, wie 5 Mos. 8, 3. geschrieben steht: „Er speisete dich mit Manna, auf daß er dir kund thäte, daß der Mensch nicht lebt vom Brod allein, sondern von allem, das aus dem Munde des Herrn geht.“

200. Es hat ja die Meinung nicht, daß du solltest verloren oder verdammt werden. Denn die Taufe ist ja gewiß genug, die Verheißung und Absolution kann dir auch nicht fehlen. Was nun? Diese Meinung hat es und damit geht Gott um, daß du lernen sollst, wie kräftig das Leben im Wort sei, und daß du bei dir gewißlich also schließest: Wiewohl ich hart umgetrieben und geängstiget werde, daß ich auch darüber gar verderbe, geschieht es doch nur darum, daß ich an meine liebe Taufe und göttliche Verheißung denken soll. Denn ich habe ja einen solchen gnädigen Gott, der für mich sorgt, daran ich gar nicht zweifle. Wiewohl es sich ansehen läßt, als sei mir Alles zuwider, so sind es doch nur Anfechtungen, dadurch mein Glaube geprüft und bewährt wird, ob ich auch fest glaube, daß Gott mein Schutz und Schirm sei.

201. Dies sind die Exempel der lieben Väter, davon ich gesagt habe, daß man in denselben auf die Verheißung und ihren Glauben, damit sie an solchen Verheißungen gehalten, und darnach auch auf das Kreuz, so sie getragen, Acht geben soll; denn das sind die vornehmsten Stücke in den Legenden der Heiligen. Darnach kann GOtt auch wohl leiden, daß sie mit der Haushaltung und andern geringen Dingen umgegangen, welche doch die Welt und Papisten gar verachten. Denn sie sind auch Menschen gewesen wie andere Menschen, und thun solche Wercke, so andere Menschen Wercken auch gleich sind, welche auch wohl von gottlosen Leuten geschehen können; also daß zwischen Jakob und einem jeglichen andern Hirten gar kein sonderlicher Unterschied gewesen. Aber dies ist allein der Unterschied, welchen nicht bald ein jeder fleischliche Mensch sehen kann, nämlich, daß Jakob die Verheißung hat, ist GOttes Kind, ist unter seinem Schutz und Schirm, und kann solchem GOtt vertrauen; darum ist auch alles, was er thut, daß er Kinder zeugt, die Ziegen melkt, GOtt wohlgefällig und angenehm gewesen. Warum das? Ei darum, daß hier die Verheißung regieret: er ist unter dem Himmel göttlicher Verheißung und glaubt derselbigen.

\quellseite{288}

202. Wenn aber nun die Anfechtungen mit Haufen daher gehen, alsdann sieht man, was für wahrhaftige Zeugniß die Heiligen an sich gehabt: wie allhier keine Ziegen gemolken wird, sondern werden beide, Vater und Sohn, getödtet, das ganze Haus wird zerrüttet und die Kirche oder Gemeinde GOttes gar betrübt. Was geschieht nun daselbst? Es erhebt sich ein groß Schreien und Heulen: Ach, wie sind wir doch so gar nichts, sind verlorne Leute! Aber dagegen sagt der Glaube also: Du bist nicht verloren; bedenke, daß du noch göttliche Verheißung hast, welche über dir ausgebreitet, so weit der Himmel ist. GOtt sorgt noch für dich, wiewohl du das weder siehst noch fühlst. Es ist nur eine kleine Wolke überher gezogen, welche GOtt den Herrn gar verschlungen zu haben scheint.

203. Von diesen Dingen sollte man bei den Historien der Väter sagen und dieselben oftmals wiederholen, daß wir lernten tapfer im Glauben stehen, und die Verheißung, so mir von GOtt empfangen, selbst auch groß achten. Du sollst dich nicht gering achten; wenn du getauft bist, GOttes Wort hast, absolvirt bist, berufen bist: so bedenke, daß das Himmelreich über dir ausgebreitet, und nicht allein GOtt selbst, sondern auch alle Engel ihre Augen auf dich gerichtet haben. Wenn nun schon Alles zerstört wird, Himmel und Erde über einen Haufen geworfen werden, und sich auch alle höllischen Pforten regen, der Papst, Kaiser, Türke auf das allergreulichste toben und wüthen; so sage du also: Ich bin getauft; so hast du es gar. Mit solchem Glauben und Vertrauen wirst du siegen. Denn dieweil GOtt für dich sorgt, wird er dich nicht verlassen, und wird dir auch kein Schade, der dich an deiner Seligkeit hindern möge, widerfahren.

204. Das hebräische Wort, nachal, will ich den Hebräern befohlen sein lassen, daß sie erklären und ausforschen, was es eigentlich bedeute. So viel, als ich es verstehen kann, heißt es, wo man etwas mit List ausrichten will, das wir gemeiniglich heißen, practiciren, seltsame Practiken machen; wie zu unsern Zeiten die Sophisten sind, ja, damit jetzt beinahe alle Bauern und Bürger umzugehen wissen. Diese Brüder machten ihre Anschläge mit listigem Rath, gingen betrüglicher Weise und hinterlistig damit um, daß sie Joseph mit List und tückischer Weise tödten möchten, damit der Vater solchen Rath nicht inne würde und daß sie ja die Schuld nicht tragen müßten. Maleachi am 1. Cap. 2. 14. steht geschrieben: „Verflucht sei der Vortheilhafte, der in seiner Heerde ein Männlein hat, und wenn er ein Gelübde thut, opfert er dem Herrn ein untüchtiges.“ Das hebräische Wort, nachal, heißt auf Deutsch listig, tückisch, betrüglich, wie der Bischof von Mainz und dergleichen andere böse alte Schälke mehr, so mit List umgehen, und Alles mit solcher ihrer Trügerei zu umstricken wissen, selbiges auch listiglich und tückischer Weise bedecken können. Wir pflegen auf Deutsch zu sagen: Sie wischen das Maul etc. Wer hat es gethan? Der Niemand, der Schalk. Niemand ist es, so beide im Haus- und Weltregiment Alles verwirrt und zunichte macht.

205. Also thun diese Gesellen hier auch, tödten den sehr frommen Jüngling, und machen es doch also, daß sie gleichwohl für die Unschuldigen, die gar keine Schuld haben, mögen gehalten werden. Ueber das können sie sich auch noch also stellen, als ob ihnen ihres Vaters Leid sehr zu Herzen gehe und ganz wehe thue. Sie haben ihn nicht mit offenbarer Gewalt umbringen können; denn es würde es der Vater erfahren haben, vor welchem sie diese That mit Fleiß haben verbergen wollen. Derhalben schließen sie nun diesen bösen Rath, daß sie ihn umbringen wollen, und gleichwohl nicht für die angesehen oder gehalten werden, so diesen Todtschlag angerichtet haben. Wir wollen, werden sie gedacht haben, also thun, wollen ihn erwürgen, und sagen, es habe ihn ein wildes Thier, ein Löwe, Bär oder Wolf zerrissen. Wir wissen zwar von seinem Tode nichts; wir wollen sagen, wir haben unterwegs das blutige Kleid gefunden; das wollen wir unserm Vater schicken, daß er zusehe, ob es Josephs Kleid sei oder nicht. Wir wissen aber wahrlich nicht, wo oder wann er möge umgekommen sein.

206. Dies ist aber ein sehr schön Geheimniß und Bild des Israelitischen Volks. Denn die Juden sind eigentlich solche nochlim, wie das an König Saul, und den Juden, so vor Pilato gestanden, zu sehen ist, die da sagten Joh. 18, 31.: „Wir dürfen niemand tödten.“ Dies sind, sage ich, die rechten nochlim und der Schalks-Niemand, die den Schaden thun, und doch unterdes nichts Anderes suchen, denn daß sie wegen ihrer Unschuld und Frömmigkeit hoch mögen gelobt werden. Dieses scheint mir die Bedeutung und Nachdruck des Wortes zu sein.

207. Siehe aber zu, ob nicht diese Brüder vom leidigen Teufel besessen gewesen, dieweil sie damit so tückisch umgegangen, wie sie ihren Bruder Joseph umbringen möchten und ihren alten Vater so bößlich betrügen. Wiewohl ich ohne Sünde glauben mag, daß sie nicht alle zugleich so gar böse gewesen, sondern daß die vornehmsten Rädelsführer dieser Sünde Simeon und Levi worden gewesen sein. Denn es scheint, als ob Juda und Ruben nicht gern in diese greuliche Sünde willigen wollen. Wie denn in der Erzählung, da Mose erzählt, wie es alles ergangen, Ruben entschuldigt wird, und daß Juda gerathen, man sollte ihn lieber verkaufen etc. Aber diese Zwei, Simeon und Levi, haben in ihrem Herzen das bei sich also gewißlich beschlossen, sie wollen kurzum ihren Bruder Joseph erwürgen und gar vertilgen; die Andern aber, als, Gad, Asser, Sebulon, Naphtali, Jaschar haben ihre Entschuldigung daher, daß sie noch jung gewesen. Denn dies ist geschehen, da Joseph siebenzehn Jahre alt gewesen, und Juda ist da zwanzig Jahre, Levi ein und zwanzig, Simeon zwei und zwanzig, Ruben aber drei und zwanzig Jahre gewesen; die Andern unter Joseph sind noch etwas jünger gewesen. Diese aber sind zwei große Hansen und Herren gewesen, Simeon und Levi, das sind die Junker am Hofe, ja, die allerschädlichsten Buben im Hause Jakobs. Die Andern haben ihnen nicht widerstanden und werden gleichwohl also durch diese bösen verzweifelten Buben in dieselbe Sünde auch mit verwickelt. Ja, lieber Herr GOtt, es geschieht oft, ein Bube kann wohl eine ganze Stadt irre machen.

208. Ruben hat wohl etwas dagegen gewehrt, daß sie dies unbillige Vornehmen nicht sollten vollziehen, aber seine Brüder haben ihn gar verachtet. Denn er hatte bei ihnen das Ansehen nun verloren um der Blutschande willen, so er begangen; und die Andern sind ihm so trotziger und vermessener gewesen, weil sie sich rühmten, daß sie fromm und heilig wären, und daß derhalben das Regiment im Hause und Kirche ihnen gehörte.

209. Es ist aber nützlich und gut, daß man diese Sünde der Brüder Josephs nur groß genug mache, daß daraus die christliche Kirche gebessert und getröstet werden möge; sintemal dieses trefflichen, großen heiligen Mannes Kinder so eine greuliche böse That in solcher Sicherheit bedacht und vorgenommen, und auch zugleich mit sich ihre anderen Brüder an den Reihen gebracht und mit in dieselbige Sünde gemengt, daß sie sich also mit ihnen haben verbinden müssen, Joseph umzubringen, nur um der Ursache willen, daß er seine zwei Träume, so er gehabt, erzählt hat und daß er dem Vater sonderlich lieb gewesen. Da sind doch bei ihnen alle guten Wercke gänzlich aufgehoben und verlöschen, so beide in der ersten und andern Tafel geboten, welche alle bei ihnen gar nichts gegolten haben. Es ist an ihnen nichts mehr übrig gewesen, denn die alleräußerste Verkehrtheit und eitel Unfreundlichkeit, so auch die schändlichsten Heiden hätten an sich haben mögen. Da wird GOtt nicht gefürchtet, ja, er wird gar verspottet und dazu auch gekreuzigt.

210. Darnach mag man solche Gesellen billig und mit allem Rechte Muttermörder, Vater-, Großvater- und Urgroßvatermörder nennen. Denn ich glaube, daß Kain viel länger hätte leben können, wo er ohne diese Trübsal gewesen wäre. Also auch die frommen ehrlichen Matronen, Rahel, Rebekka und Debora, wo sie dazumal noch am Leben gewesen, hätten sie doch vor großem Herzeleid sterben müssen; und ist ihnen daran sehr wohl geschehen, daß sie vor solchem Jammer ihren Abschied aus diesem Leben genommen. Denn wo sie das Jahr noch gelebt hätten, wäre es ihnen nicht möglich gewesen, so großen Schmerz und Traurigkeit zu ertragen; sie hätten bei ihren heißen Thränen die Haare aus ihrem Haupte gerissen und wären bald vor Leid und Traurigkeit gestorben. Also haben sich die ungehorsamen Kinder mit dieser einigen Sünde, so sie an ihrem Bruder begangen, zugleich in viel andere große schwere Sünden gewickelt und sich derselben damit schuldig gemacht.

\quellseite{289}

211. Derhalben lernet allhier ihr Junggesellen, daß ihr euren Eltern, Schulmeistern und andern Vorgesetzten gehorsam seiet. Denn das ist eine große erschreckliche Sünde und eine Ursache vieles unzähligen Unglücks, wo die Kinder ungehorsam sind; sintemal dieselbige Sünde Eltern, Brüder, Großväter und das ganze Geschlecht tödtet und freuet dazu noch den Sohn GOttes. Hütet euch!

212. Wie Simeon und Levi Buße gethan, geht nicht aus der heiligen Schrift hervor; aber ich will glauben, daß der größere Theil ihrer Buße und Strafe gewesen sei, daß die Knäblein in Egypten also sind ersäuft worden. Und GOtt hat einen jeglichen unter ihnen zu seiner Zeit strafen und plagen können, daß sie wiederum auch haben weinen und heulen, dazu auch ihre Haare ausraufen müssen, und haben vielleicht also wohl schwerere Strafe erlitten, denn wir denken können, welche eben nicht beschrieben sind. Jedoch scheint das wahrscheinlich, daß GOtt eine Strafe, so ihrer Sünde gleich oder ähnlich gewesen, ihnen habe auferlegt, nämlich, daß der König in Egypten das Gebot habe ausgehen lassen, daß alle Knäblein im Volke Israel, wenn sie eben geboren wären, in das Wasser sollten geworfen und ersäuft werden. Als wollte GOtt sagen: Ihr habt euren Bruder Joseph in eine Grube, darin kein Wasser gewesen, geworfen: nun will ich es wiederum dahin lassen kommen, daß eure Kinder in das Wasser sollen geworfen werden und darin also umkommen.

213. Denn es folgt doch gemeiniglich am Ende nach der Sünde auch die Bezahlung; nicht, sage ich, eine papistische Genugthuung, sondern die Strafe GOttes. Denn GOtt verzeiht und vergibt ja wohl die Sünde, aber damit der Sünder nicht darüber einschlafe und gar sicher werde, und sich der Sünden noch wohl dazu rühme, als hätte er es wohl ausgerichtet und wäre Alles wohl gethan, so kommt GOtt mit der eigenen Ruthe, sucht heim und straft die Sünde der Väter an den Kindern in das dritte und vierte Glied.

214. Warum thut er aber das? Nicht thut er es, als ob man mit solcher Strafe für die Sünde könnte genug thun; denn er läßt sich mit unserer Strafe nicht also zufriedenstellen, sieht solche unsere Genugthuung nicht an: sondern er straft darum, daß der Sünder nicht schnarche und gar sicher werde, oder sich seiner begangenen Bosheit und Sünden noch rühmen soll; sondern daß uns die begangene Sünde beißen, tödten und zwingen soll zu erkennen, wie groß und schwer die Sünde sei; daß wir also darüber weinen, seufzen und GOtt um Gnade anrufen, auf daß sie uns ohne Verdienst, aus Gnaden, um Christi Jesu willen möge vergeben werden, dieweil ja solche Gnade von GOtt verheißen ist denen, so über ihre Sünden weinen, dieselben bekennen, sich leid sein lassen und also wahrhaftige Buße thun.

215. Und wenn GOtt die Sünde nicht also strafte, würden wir nicht allein schnarchen und die Sünde gar verachten, sondern würden sie noch wohl dazu täglich häufen, und deren immer mehr und noch größer machen, denn die vorigen gewesen. So liegt wohl die Sünde vor der Thür und schläft, wie zu Kain 1 Mos. 4, 7. gesagt wird; sie bleibt aber nicht sicher noch verborgen, sondern wird dermaleinst plötzlich aufwachen, und wird dich plagen und kreuzigen; aber doch nicht zum Verderben oder Verdammniß, sondern daß du dadurch zur Buße kommen und dich bessern mögest, daß du deine Sünde bekennest, zu GOtt seufzest und schreiest, und ihn um Gnade anrufest. Also werden sie es wohl bezahlt haben. Also, sage ich, werden die Israeliten ohne Zweifel für ihre Sünden Strafen, die schwer genug waren, gelitten haben, wie das in der Historie 2 Mos. 1, 11. ff. angezeigt ist. Denn wie sie ihren Bruder, Vater und Mutter getödtet haben, also haben sie auch wiederum viel unzähliges Unglück und Plagen in Egypten unter der Tyranney Pharaos leiden müssen. Pharao musterte sie recht in Egypten, daß sie gedenken sollten, welch eine große Sünde sie begangen hätten, und daß sie die Sünde nicht für so ein gering Ding hielten.

216. Denn wo mir die Größe der Sünde nicht mit rechtem Fleiß bedenken werden, so wird sie GOtt selbst groß achten. Werden wir uns nicht selber richten, so wird uns der HErr richten. Und es ist doch noch eine große Wohlthat, daß uns GOtt zu unserer eigenen Seligkeit richtet und straft, und nicht zur Verdammniß, wie Judas und Saul gerichtet sind. Weißt du aber sicher sein und also denken: Was liegt daran, wenn ich schon meine Eltern erzürnt habe? da, sage ich dir wahrlich, schläfst deine Sünde und ist noch gar verborgen. Aber sie wird gewißlich aufwachen, und dich dahin treiben, daß du endlich sagen wirst: Ach wehe mir armen Sünder, was habe ich gethan? Warum habe ich GOtt in meinem Vater und Mutter und Schulmeister so verachtet? Warum bin ich ihnen nicht gehorsam gewesen in dem, was sie mich geheißen? Das wird endlich kommen, es kann nicht ausbleiben.

217. Also haben wir oftmals ernstlich ermahnt der heimlichen Verlöbnisse halben, daß man sie beiderseits meiden und sich die Jungfrauen und auch Junggesellen davor hüten sollen. Aber da sieht man Viele, die wie mit einem Sturme, als wären sie vor Liebe gar blind, auch wider ihr eigen Gewissen dahingehen, daß die Junggesellen nicht ordentlich und im Gehorsam gegen ihre Eltern freien, noch auch die Jungfrauen sich also freien lassen; sondern daß ein Theil den andern gleichsam mit Gewalt raubt, es sei den Eltern lieb oder leid. Was thun nun dieselbigen Anderes, denn daß sie sich mit erschrecklicher Todsünde beflecken, welche wohl eine Zeitlang schläft und verborgen liegt? Denn das Mägdlein gefällt dem Junggesellen wohl und ist sein Herz in der Liebe zu dem Mägdlein ganz entbrannt, daß er weder Tag noch Nacht an etwas Anderes denkt und ihm auch sonst von nichts träumen läßt. Aber wenn nur ein Monat oder zwei verlaufen, so folgt endlich hernach ein jämmerlich Seufzen und Wehklagen; wie Abigail zu David sagt. 1 Sam. 25, 29. 31.: „Die Seele meines Herrn wird eingebunden sein im Bündlein der Lebendigen bei dem Herrn, deinem GOtt“ etc. „So wirds dem Herzen meines Herrn nicht ein Stoß, noch Aergerniß sein, daß du nicht Blut vergossen hast ohne Ursach“. Solche Anstöße und Aergernisse pflegen für und für das Gewissen zu betrüben, da man denkt: Ich habe dies Mägdlein wider der Eltern Willen genommen. Da kommt hernach das Rauen im Nacken, das schwarze böse Hünfletin, der Reuling, das beißt dein Lebetage, hört nicht auf, wenn auch die Sünde schon vergeben ist.

218. Bei denen aber wird der Schmerz oder Traurigkeit des Gewissens noch etwas größer und schwerer sein, die nun den Willen GOttes wissen, daß man ohne Bewilligung der Eltern nicht heirathen oder eine Ehe eingehen soll. Denn zuvor, da wir GOttes Gebot nicht gewußt, möchte vielleicht solche Unwissenheit ihre Entschuldigung haben; aber wie viel besser wäre es nun, mit Wissen und Willen der Eltern ehelich werden, daß wenn du bei dem Mägdlein schlafen wirst, du sie auch mit Freuden und gutem Gewissen, ohne allen Anstoß und Aergerniß berühren mögest, daß beide GOtt und die Eltern dir dasselbe wohl gönnen, und darüber lachen und fröhlich sind. Das, sage ich, wäre ja viel besser, denn daß du dich mit ihr also bissest und heimlich vermischest, da dein eigen Herz, so dich dessen beschuldigt, mit steter Traurigkeit gegen dich also murren wird: Hättest ich doch dies Mägdlein wohl bekommen und mich zu ihr halten mögen, daß die Eltern auch ihren Willen dazu gegeben, und darüber mit uns fröhlich und guter Dinge gewesen wären. Warum habe ich denn meinen lieben GOtt und auch die Menschen damit also betrübt und erzürnt? Denn das schwarze Sündlein, der Reuling, so endlich zu kommen pflegt, hört nicht auf zu widerbellen und das Gewissen zu beißen, ob du schon wohl weißt, daß dir die Sünden vergeben sind.

219. Und wir Prediger, so das Kirchenregiment verwalten sollen, haben oftmals mit solchen verwirrten Gewissen, so mit dergleichen Sünden beschweret sind, viel zu schaffen, bis wir sie wiederum aufrichten und trösten. Derhalben mögen wir mit den Juristen und ihren Canones, welche eigentlich diese Verordnung angerichtet und gestiftet haben, billig zürnen. Sie zünden die Häuser an, daß sie brennen, und lassen darnach uns dafür sorgen, wie wir dieselben wiederum löschen mögen. Wir thun aber, was wir vermögen, absolviren und trösten die, so mit solchen Sünden beschweret sind; noch bleibt das schwarze Sündlein nicht aus: der Schmerz kommt doch wieder und wird immer erneuert, so oft in der Haushaltung ein Unglück oder Schaden vorfällt, wenn etwa ein Kind krank wird, dem Weibe es ungerade mit der Geburt geht oder wenn sie dahin stirbt etc. Das ist die Ursache des ganzen Unglücks, nämlich, die Verachtung der Eltern; mit welcher Verachtung das schwarze Sündlein dir bald unter Augen hält und dir dieselbe aufrückt; und wo das nicht vorhanden wäre, könnten wir wohl geduldiger allerlei Beschwerden und Unfall leiden und ertragen. Denn da ist das Herz still und wohl zufrieden, weiß, daß dieser Stand Gott wohl gefällt, spielt oder scherzt in gutem Frieden mit seinem lieben Gemahl, und ist dessen gewiß, daß alle seine Werke Gott wohlgefallen und angenehm sind. Und im Fall, daß auch etwas Widernwärtiges sich begeben würde, wird doch das Herz derhalben nicht leichtlich bewegt oder erschrocken.

\quellseite{290}

220. Dies sage ich also zum Exempel; man soll es aber auch auf alle anderen Sünden ziehen, auf welche ohne Zweifel endlich auch solcher Schmerz des Gewissens folgen wird, auch in den Menschen, so am allerschiersten und auch sehr greulich und unansehnlich sind. Denn es ist eine solche Stimmung, die man schwer von der Person, die also gesinnt, absondern kann. Darum, lieben Kinder, sehet euch vor, es kommt, hütet euch. Ich sei noch jung, ihr wißt noch nicht, welch ein greulich Thier ein böses Gewissen sei. Gewöhnet eure Herzen daran, daß ihr euren Eltern und Schulmeistern gehorsam seiet. Es ist gar ein fein Ding, wie Juvenal sagt, daß Einer keinerlei Sünden halben bleich oder schamroth werden darf. Es ist nichts Besseres und darnach Einer mehr Verlangen haben soll, als daß man ein solch Leben führen möge, das Gott angenehm und wohlgefällig sei, und daran auch fromme Eltern Lust und Freude haben mögen. Ich könnte hier wohl viel Exempel anziehen, so ich beide selbst erfahren und dazu auch an Andern gesehen habe, welch ein groß Unglück ein böses Gewissen sei; aber ich übergehe dieselben um der Kürze willen.

221. Denn dafür soll man es gewiß halten, daß die Sünde gewißlich der Tod und Verdammniß zum Tode sei. Es ist aber zweierlei Verdammniß: eine des Zornes und die andere der Gnade. Wo es eine Verdammniß ist des Zornes, da fühlt man solche große Kraft und Gewalt der Sünde, daß sie ohne alle Hilfe und Trost das Gewissen überfällt und gar in Grund verderbet, wie Juda, Saul und allen andern Verzweifelten widerfahren ist. Von solcher Verdammniß wolle uns ja Gott gnädig erlösen und davor behüten. Die andere Verdammniß ist eine Gnadenverdammniß; davon sagt Paulus 1 Cor. 11, 31.: „So wir uns selber richteten, so würden wir vor dem Herrn nicht gerichtet.“ Da erhält Gott die, so dermaßen gerichtet werden, durch das Wort und Verheißung, durch den Trost des Wortes und Predigamtes.

222. Wenn der Pfarrherr, oder sonst ein jeder Diener des Worts kommt und ein solch erschrocken Gewissen aufrichtet, kommt es doch gemeiniglich also, wiewohl es aufgerichtet und getröstet wird, daß doch immer das Seufzen auch wider den Willen und alles Fühlen solches Trostes wieder kommt. Ich weiß wohl, was mir begegnet, wenn ich bedenke, wie ich mein ganzes übriges Leben geführt. Denn wiewohl ich das weiß, daß mir meine Sünden vergeben sind, kommt doch gleichwohl das pukah, das ist, der Anstoß immer noch wieder: ohne den Anstoß kann ich nicht sein, ich muß schamroth werden, und sagen: Pfui dich, was habe ich gemacht! Ich weiß mich ja nicht schuldig, als hätte ich mit eines Andern Weibe Unzucht getrieben, jemand todt geschlagen, oder dergleichen anderer grober Sünden weiß ich mich nicht schuldig; da kann mich das schwarze Sündlein nicht heißen: aber anderer Sünden halben beißt es mich, daß auch mein Gewissen nun schon zufrieden gestellt, und die Narbe bedeckt und zugeheilt ist.

223. So macht nun bei Etlichen die Sünde auf zur Verdammniß, bei den Andern aber zur Züchtigung und zur Buße, daß sie gezüchtigt und bekehrt werden mögen, daß sie zu Gott rufen und schreien, und endlich selig werden: wie dem König Manasse seine Sünde auch aufgeweckt und offenbar ist, aber ihm zu Gnaden und Seligkeit. Dergleichen auch dem Petrus, Paulus und der Magdalena.

224. Die aber haben es noch etwas besser und sind desto seliger, die da ohne Anstoß und Seufzen wegen schwerer Sünde leben, so sie begangen, als, fromme und gottselige Eheleute, die da gewiß sind, daß sie mit Bewilligung ihrer lieben Eltern ehelich geworden sind. Dieselben, sagen sie, können allerlei Unglück viel leichter und mit gutem Gewissen ertragen und überwinden, denn Andere, so darin die Gewalt ihrer Eltern verachtet haben.

225. Diese heimliche Sünde aber ist jetzt noch still und gar verborgen, dieweil sie damit umgehen, daß sie diese greuliche Sünde vollbringen mögen. Und wird uns dies Exempel also vorgehalten, Beide uns zu Trost und auch die Gottlosen und Sichern damit zu schrecken. Denn diese Brüder Josephs haben zwar Gnade erlangt, da sie endlich seufzten, Leid trugen und büßten, nicht nach päpstlicher, sondern göttlicher Buße, welche wohl eine Verdammniß zum Tode ist, aber doch also, daß der Sünder dadurch erneuet werde, daß er zu Gott schreie, und um die Gnade anrufe, so er in Christo, seinem lieben Sohne, erzeigt. Also sind die Niniviten bekehrt worden und haben gefastet; aber solch Fasten ist nicht ihre rechte Buße und Genugthuung gewesen, sondern das Seufzen, da sie sagten: „Daß wir nicht verderben“, Jon. 3, 9., ist ihnen zur Seligkeit gerathen, daß es ihnen keine Verdammniß zum Tode war.

226. Sie lassen es aber noch bei dem greulichen Anschlage nicht bleiben, sondern haben sich auch dazu mit bittern und unfreundlichen Worten vernehmen lassen: „Sehet“, sagen sie, „der Träumer kommt daher“ etc. Das böse Sagen „der Träumer“ bedeutet gar eine bittere Verachtung. Denn es ist beide ihr Herz und Zunge voll teuflischer Bosheit und Bitterkeit gewesen. Und ist mit diesen nämlich ganz wohl zu vergleichen der Bischof in unserer Nachbarschaft, welcher an seinem Hofe selbst unterhalten und predigen gehört hat den frommen gottseligen Prediger M. Georg, den er darnach meuchmörderischer Weise hat jämmerlich umbringen lassen. Und da derselbe nun erwürget und umgebracht war, hat sich der Bischof gestellt, als wüßte er von dem Todtschlage gar nichts, und hat sich äußerlich vernehmen lassen, als wäre es ihm herzlich leid, daß der fromme Mann also unbilliger Weise umgekommen.

227. Sobald sie nun ihren Bruder gesehen, den sie je sollten geliebt, freundlich empfangen und geschützt haben, vornehmlich dieweil er noch jung und dazu auch ganz unschuldig war, haben sie alsbald angefangen feindlich und giftig von ihm zu reden und gesagt: Ecce iste Baal somniorum; denn so steht im hebräischen Text: Das ist der Herr und Meister der Träume; der schändliche Dichter, der Träumer kommt. Sie geben ihm Schuld, daß er die Träume erdacht oder gedichtet habe.

228. Baal heißt ein Herr, ein Ehemann, oder eigentlich ein Hauswirth: „Jes. 1, 3.: Ein Ochse kennet seinen Herrn“ und im 5. Buche Mose: „Wenn du deines Bruders Esel siehest irre gehen, sollst du ihn wieder zu seinem Herrn führen.“ Daher hat der Baalsdienst, welcher Götzendienst die weiteste Verbreitung fand, seinen Namen. Denn die abgöttischen Juden hatten selbst einen solchen Gottesdienst erwählt, der vortrefflich, und viel köstlicher und besser sein sollte denn die gemeine Weise des Gottesdienstes. Denn sie wollten auch einen solchen Hauswirth haben, der ihr Buhle sein sollte. Baalath heißt eine Buhle. Wir können uns daran nicht genügen lassen, werden sie gesagt haben, daß Gott unser Herr ist, wie er anderer Völker Herr auch ist; sondern wollen haben, daß er unser Baal, das ist, unser Buhle sein soll, wir wollen ja etwas besser sein denn Andere. Darum haben sie auch Fest- oder Feiertage und gewisse Opfer aufgerichtet. Wie bei uns die Mönche auch gethan, welche sich an dem rechten geistlichen Gottesdienst des Glaubens, der Liebe und Hoffnung nicht haben genügen lassen, sondern über das noch sonderliche neue Arten und Weisen der Kleidung, Unterschiede der Speise, der Zeit und Stätte erdacht: das haben sie eine Vollkommenheit geheißen, als wären sie nun nicht mehr gemeine Christen, sondern wie Engel, weit über den gemeinen Haufen der Christen gesetzt.

229. So sehen wir nun hier an diesen Brüdern Josephs, wie ganz sauer und unfreundlich sie die Träume des einfältigen frommen Jünglings ausgelegt oder gedeutet haben; gleich als hätte er dieselben mit List erdichtet, daß er damit also den Andern das Regiment und Priesterthum nehmen und entziehen möchte. Also ist der Haß, so sie unbilliger Weise auf ihn geworfen, bei ihnen entzündet und immer gemehret worden. Denn Simeon und Levi haben das nicht ertragen können; Ruben aber war wegen der Sünde, so er begangen, gar verachtet,
\quellseite{291}
und Juda war unter ihnen der Jüngste. Darum werden sie gesagt haben: Es wird je das Regiment sonst niemand anders gebühren denn uns zwei, und wenn ich, Simeon, gefangen, wird die Herrschaft auf Levi kommen. Dies sind die rechten schalthaftigen Meister, so diese Tragödie gespielt haben. Und werden derhalben auch hernach von dem Patriarchen Jakob mit ihren rechten Farben herausgestrichen und mit gebührlicher Strafe abgemalt werden, Cap. 49, 5. 6. 7.: „Die Brüder Simeon und Levi; ihre Schwerter sind mörderische Waffen. Meine Seele komme nicht in ihren Rath, und meine Ehre sei nicht in ihrer Kirche“ etc. „Ich will sie zertheilen in Jakob, und zerstreuen in Israel.“ Es hat der ganze Stamm Levi nichts Eigenes gehabt, und im Stamm Simeon sind nur arme Schreiber und Schulmeister gewesen; wie jetzt die Juden noch sagen, es sei sonst im ganzen Volk kein Schriftgelehrter mehr denn allein im Stamm Simeon. Derhalben ist diese Strafe der Armuth und Vetzelei über beide Stämme ergangen, wiewohl ihnen die Sünde, da sie Buße gethan, vergeben worden ist.

230. Und da sie also begehrt, obenan zu sitzen und allein das königliche Regiment über die Andern zu haben, sind sie jämmerlich hinunter gestoßen und in immerwährende Armuth geführt worden. Denn der Zank ist über der Erstgeburt oder über der königlichen und priesterlichen Würde gewesen. Das ist die Ursache des großen Zankes und bittern Hasses gewesen, dadurch diese Brüder getrieben worden, daß sie solche unmenschliche Sünde begangen, so den Mord, welchen Kain an seinem Bruder auch geübt, weit übertroffen hat. Denn derselbe Kain hat um gleicher Ursache willen seinen Bruder Abel auch erwürgt: darum daß sich Gott das Opfer Abels gefallen lassen, Kains aber verworfen hat, darum muß er sterben. Also hat dem Joseph geträumt, er sollte der Fürst und Priester im Hause sein, und ist die Deutung der Träume an sich offenbar; darum sind sie so heftig ergrimmt und erzürnt worden, daß sie zu ihm sagten: „Solltest du unser König werden, und über uns herrschen?“

231. Ebenso pflegt es noch unter Brüdern zuzugehen, wo sie das Erbe mit einander theilen sollen. Denn da sind die Brüder selten mit einander eins und können sich Fremde viel leichter mit einander vereinigen denn Brüder.

Joseph hat wahrlich um dieses Vorzugs willen sterben müssen, der ihm doch nach allen Rechten billig gebührte; denn er war der Erstgeborne von der, so Jakobs vornehmste Ehefrau war, und die er auch erstlich begehrt und geliebt hatte. Denn Lea ist ihm nach derselben betrüglicher Weise, gleichsam mit Gewalt, aufgedrungen und gegeben worden, so er sie doch nie begehrt oder geliebt hat. Jedoch, dies alles unangesehen, versuchen diese zwei Buben, solche Herrlichkeit hinzureißen und sich dieselbe anzumaßen; so doch nach des Vaters Willen und auch wegen der Geburt dieselbe Joseph gehört hat. Derhalben sind sie wahrlich ärger denn Kain gewesen, und die Sünden, so sie begangen, sind viel greulicher denn der Mord Kains. Denn sie denken daran nicht, daß sie ihren armen alten Vater betrüben werden, der doch um ihretwillen vom Anfang seines Ehestandes so viel Unglück erlitten. Da muß ja kein Verstand und keine Gedanken in ihnen gewesen sein. Sie wollen schlecht mit dem Kopf hindurch: er muß sterben. Das ist der Teufel. Dasselbe haben darnach ihre Nachkömmlinge, die Juden, eben auch gethan. Denn unangesehen der vielen Wunderwerke, so Christus der HErr gethan, haben sie doch immer das Wort wiederholt und gesagt, Luc. 19, 14.: „Wir wollen nicht, daß dieser über uns herrsche.“ Er will uns lehren, will unser Hoherpriester und König sein, sagt, er sei Gottes Sohn. „Hinweg mit ihm, kreuzige ihn.“ (Gleiche Verstockung ist allezeit bei denen, so göttlicher Wahrheit und Ordnung widerstreben. Denn dasselbe lassen zu unsern Zeiten die Papisten auch hören: Wir wollen uns kurzum von niemand regieren oder lehren lassen“ etc.

232. Die andern Brüder aber Josephs haben gleichwohl diesen teuflischen Anschlag nicht erstlich gestiftet oder erdacht, sind aber doch mit in das Spiel gekommen; wie es denn in der Welt zu gehen pflegt. Und da sie nun mit einander in ihren Herzen beschlossen, ihren Bruder zu erwürgen, haben sie über das auch noch einen sehr listigen Rath beschlossen, nämlich, wie sie solche ihre böse That und den Leib ihres erschlagenen Bruders verbergen möchten. Und, Lieber, siehe doch, wie erschrecklicher Weise sie verblendet und vom Teufel gar besessen gewesen sein müssen. Denn sie haben noch zu dem, daß sie den Mord an ihrem Bruder begangen, auch eine schändliche Lüge gelogen; ja, sie haben keinerlei Sünde, vornehmlich wider die andere und erste Tafel, allein das sechste Gebot vom Ehebruch ausgenommen, unterlassen.

233. Derhalben, wer da ihre Sünden groß machen wollte, der würde dazu allhier weitläufige Materie haben. Denn sie vergreifen sich an allen Zehen Geboten, berauben den Vater aller Nachkommen seines Erbes und Gutes, dazu seiner Ehre, und wo sie es vermöcht, hätten sie ihn auch noch dazu der Verheißung wohl beraubt. Wider das sechste Gebot haben sie zwar nicht gesündigt; aber in allen andern Sünden, so in den Zehen Geboten von Gott verboten, sind sie ganz ersoffen gewesen, haben darin bis über die Ohren gelegen. Desgleichen haben sie das ganze Geschlecht betrübt und in ein erbärmliches Leid und Wehklagen geführt, und vornehmlich ihren Vater und Großvater, welcher den Jammer gesehen, darein seines Sohnes Kinder gerathen, und hat doch nicht gewußt, wer solchen Jammer gestiftet. Die Ismaeliten und Midianiter haben es zwar wohl gewußt, aber zu großem schweren Aergerniß der ganzen Kirche.

234. Darum haben sich diese Brüder fein tückisch wissen zu schmücken, und sagen: „Laßt uns ihn in eine Grube werfen, und sagen, ein böses Thier habe ihn gefressen.“ Da haben sie ja deß gar nicht geachtet, daß er ihr Bruder war. Es hätte wahrlich dem allergrausamsten Feinde einfallen müssen, daß er gedacht hätte, er wäre ja sein Bruder, und es müßte ja diesesfalls billig ein Bruder des andern verschonen und sich an ihm nicht vergreifen. Aber diese Gesellen würgen ihren Bruder dahin und werfen ihn darnach in eine Grube, erdichten gegen den Vater und Andere eine vermeinte Entschuldigung dieser ihrer bösen That, daß sie sagen: „Wir haben ihn traun nicht getödtet, sondern es hat ihn ein wildes Thier gefressen.“

235. Hier ist abermals ein solcher Ort, dabei man gedenken soll, wie doch unsers HErrn Gottes Rath so wunderbarlich sei; welches wir gar nicht, oder nur ein wenig und im Dunkeln sehen und verstehen können. Denn wo es scheint, als ob Gott zürne, so zürnt er doch nicht von Herzen, sondern behält noch immer seine Barmherzigkeit und Treue über uns; wie denn solches die heilige Schrift von ihm allenthalben bezeugt. Solche Verheißungen leuchten immer in dieser Finsterniß der Trübsal.

236. Darnach ist hier auch ein anderer Rath Gottes, daß er es geschehen läßt, daß diese Brüder dem Joseph heimlich nachstellen, und mit großer Weisheit solchen Mord, so sie an ihm begangen, verhehlen; wie Kain auch gedachte, daß es gar verborgen wäre, daß er seinen Bruder Abel erwürgt, und daß davon weder Gott noch Menschen etwas wüßten. Also zweifeln diese Gesellen auch nicht, es könne wohl verhütet werden, daß ihr Vater nichts davon erfahre, was sie mit Joseph angerichtet. Aber in diesen ihren Rathschlägen und listigen Tücken allen, deren sie sich bei sich selbst höchlich erfreut, werden sie gar bethört und verblendet.

237. Wie zu unsern Zeiten Gott auch den Türken und Papst weise sein und mit eitel List umgehen läßt, und ihnen dabei wunderbarlich groß Glück gibt. Aber eben damit macht er sie gar toll und thöricht, wenn sie sich dünken lassen, daß sie die Weisesten seien. Wer ist aber, der solches versteht? Denn wir sind vielmehr schwach, sind Narren und gar verachtet in der Welt. Wir müssen immer Leid tragen und trauern; sie aber sind fröhlich und getrost wegen der Ehre, so sie haben, nämlich, daß sie allein weise und mächtig sind, ja, daß sie das Regiment haben über die ganze Welt. Aber die, so den Heiligen Geist haben und des HErrn Sinn erkennen, urtheilen davon, wie im 73. Psalm V. 18. 19. steht: „Herr, du setzest sie aufs Schlüpfrige, und stürzest sie zu Boden. Wie werden sie so plötzlich zu nichte“ etc. Derhalben irren alle die im Herzen, so sich auf solche ihre Weisheit und List verlassen. Die Bischöfe und Cardinäle sind mit allen ihren heimlichen Praktiken zu eitel Narren geworden; aber niemand ist, der es sehen könne. Wir sehen es zwar wohl, daß sie in allen ihren Rathschlägen und Vornehmen gar toll und unsinnig sind; aber daß werden wir sie noch lange nicht überreden, daß dem also sei.

238. Wenn nun die Gottlosen uns heftig bedrohen mit dem Tode, Galgen, Feuer und Schwert, damit sie uns gedenken zu tödten und gar zu verscharren, sollen wir gewißlich wissen, daß Gott, der da gesagt hat: Ich will dein Schutz sein, ihrer Thorheit lachen und das Gegentheil vornehmen werde. Denn so haben diese Brüder Josephs gedacht: „Kommt, laßt uns ihn ermürgen“; dagegen aber sagt GOtt also: Lasset mir ihn leben, und bewahret ihn, daß er unverletzt bleibe. Sie haben gesagt: Wir wollen ihn in die Grube verscharren; GOtt aber sagt: Ihr sollt mir ihn von den Todten auferwecken. Sie hoffen damit das Lob der Gerechtigkeit und Unschuld zu erlangen; GOtt aber sagt dagegen also: Verklagt euch selbst und stürzet euch in die ewige Verdammniß hinein. Denn also ändert und verkehrt GOtt den Menschen ihre Rathschläge. Und das verstehen allein die Heiligen, die den Geist haben, und ehemals göttliche Hülfe und Rettung aus der Noth empfunden und erfahren haben.

\quellseite{292}

239. Und dies ist auch die Lehre des 2. Psalms V. 2. 3. 4., welcher beider Rathschläge und Vornehmen abmahlt. Die Könige und Fürsten im Lande toben und schreien: „Lasset uns zerreißen ihre Bande, und von uns werfen ihre Seile“; wir wollen ihren Christum nicht ansehen. „Aber der im Himmel wohnet, lachet ihrer, und der HErr spottet ihrer.“ Liebe Herren, sagt er, euer Verwerfen und Zerreißen ist mein Erhöhen, Auferwecken, Krönen und zum Könige setzen. Das ist ihr Gewinn davon. Euer Begraben oder Verscharren ist mein Auferwecken aus der Hölle heraus, und die Entschuldigung eurer Sünde ist eure ewige Verdammniß.

240. Derhalben, daß diese Brüder mit großem Trotze sagen: „So wird man sehen, was seine Träume sind“, dasselbe haben sie darnach zu ihrem großen Schaden gesehen, da sie schon vorlängst solchen bittern Haß gegen ihren Bruder mit dem greulichen Mord erfüllt hatten. Denn GOtt hat ihr Toben und grimmigen Zorn wohl gesehen, wiewohl sie gemeint, sie könnten denselben vor ihm fein verhehlen. Darum laßt uns lernen das Drohen und die greulichen Rathschläge unserer Feinde verachten, und es gewiß dafür halten, daß GOtt im Himmel droben schon ganz das Gegentheil beschlossen habe und daß er schon jetzt ihrer lache; daß er mit uns aber fein freundlich spiele, daß er unsern Glauben und Hoffnung prüfen und bewähren möge; wiewohl es uns nicht sehr angenehm, sondern ein sauer unfreundlich Spiel ist. Welche aber GOtt also belacht, und die selbst jetzt auch lachen, denen wird solches ein saures und betrübtes Lachen werden.

\section[Fünfter Theil]{Fünfter Theil. Wie Joseph in eine Grube geworfen, und wie er endlich von seinen Brüdern verkauft wird}

\subsection{I}

\begin{quote}
\textbf{V. 21. 22. Da das Ruben hörte, wollte er ihn aus ihren Händen erretten, und sprach: Lasset uns ihn nicht tödten. Und weiter sprach Ruben zu ihnen: Vergießet nicht Blut, sondern werfet ihn in die Grube, die in der Wüste ist, und leget die Hand nicht an ihn. Er wollte ihn aber aus ihrer Hand erretten, daß er ihn seinem Vater wieder brächte.}
\end{quote}

241. Mose entschuldigt Ruben, daß er mit den andern Brüdern nicht sich verschworen habe, Joseph umzubringen; sondern daß er sich nach Vermögen ihrem grimmigen Zorn widersetzt, und da er ihn nicht habe gänzlich erretten können, er doch so viel zuwege gebracht, daß er gleichwohl nicht getödtet worden. Denn er hat auch das nicht gesucht oder gemeint, daß Joseph in die Grube sollte geworfen werden, wiewohl es scheint, als habe er solches gerathen; sondern er hat sich in allewege beflissen, daß er seinen Bruder am Leben erhalten möchte. Derhalben hat ihm auch GOtt solches für die That und vollkommene Erlösung, so viel an ihm gewesen, gerechnet, dieweil er dazu guten geneigten und vollkommenen Willen gehabt, und keines gern verhindert hätte, nämlich, daß Joseph nicht getödtet, noch auch in die Grube geworfen würde. Denn auf diese Meinung sagt auch Augustin: Deus coronat intus voluntatem, ubi foris non invenit facultatem, das ist: GOtt krönet inwendig den guten Willen, wo er auswendig das Vermögen nicht findet. Und Christus deutet auch selbst das fünfte Gebot also Matth. 5, 22.: „Wer mit seinem Bruder zürnet, der ist des Gerichts schuldig.“ „Wer seinen Bruder hasset, der ist ein Todtschläger“, 1 Joh. 3, 15. Denn da ist der Wille vollkommen, welcher nichts Anderes begehrt, denn daß der Bruder nur aufgegeben und getödtet werden möge. Denn dann ist ein Werk schon vollendet, wenn der Wille, solches zu vollbringen, vorhanden ist, es sei zum Guten oder zum Bösen.

242. Derhalben ist Ruben an dem Blut seines Bruders unschuldig, wiewohl ihn diese Unschuld von der Blutschande, so er zuvor begangen, nicht befreit; und es scheint, als hätte er derselben Strafe, so er damit verwirkt, mit dieser That entfliehen oder lindern wollen, ob er sich vielleicht also bei dem Vater wiederum einschmeicheln könnte. Aber es hat der Vater solchen Fleiß und Demuth an Ruben nicht geachtet, und wiewohl er ihm die Sünde vergeben, ist doch gleichwohl die Strafe noch auf ihn geblieben; wie man im deutschen Sprüchwort sagt: Alte Schuld rostet nicht.

243. Also haben die andern Brüder endlich auch ihre wohlverdiente Strafe noch leiden müssen für die Sünde, daß sie an ihrem Bruder den Mord begangen; und ist dies die Strafe gewesen, daß ihre Nachkommen und Kindeskinder in Egypten die schwere Dienstbarkeit haben tragen müssen, und daß ihre jungen Kinder, so Knäblein waren, im Wasser haben ersäuft und umgebracht werden müssen; vornehmlich aber sind die zwei Stämme Simeon und Levi vom Vater besonders verflucht worden. Denn die Sünde wird nicht so gar vergeben und nachgelassen, daß nicht die Strafe der begangenen Uebelthat und eines bösen unruhigen Gewissens noch sollte übrig bleiben. Also, wiewohl David seine Sünde auch vergeben war; denn er hörte Nathan 2 Sam. 12, 13. sagen: „Der HErr hat deine Sünde weggenommen“: ist doch gleichwohl auf dieselbe Sünde noch gar eine schwere Strafe von GOtt gefolgt, nämlich, der Aufruhr seines Sohnes Absalon, welcher den Vater des Reichs und aller seiner Güter beraubt und ihm noch dazu alle seine Kebsweiber beschlafen hat.

244. Derhalben sollen wir uns mit allem Fleiß und Wachsamkeit vor Sünden hüten; denn sie bleiben nimmermehr ungestraft, es sei denn, daß die Strafe durch große ernste Buße und Barmherzigkeit GOttes aufgehoben werde; sonst heißt es: Alte Schuld rostet nicht. Und wo große Noth und Gefahr vorhanden, pflegen wir zu bitten: O GOtt, räche nicht alle Schuld! Gedenke nicht, HErr, an unsere Missethat! Also haben Simeon und Levi frei dahin und ohne alle Scheu gesündigt und wird solche ihre Sünde auch eine Zeitlang heimlich gehalten; aber sie schläft vor der Thür, daß sie dereinst nach der Strafe wiederum erweckt werde. Da findet sie sich alsdann gar bald, und beißt so, daß der Sünder anfängt zu schreien: Ach wehe mir armen Manne! Ich habe, leider, die Strafe mit dieser oder jener Sünde wohl verdient! Auweh, das habe ich damit verdient! Derhalben sage ich noch, daß es um die Sünde ein sehr erschrecklich, betrübt und elend Ding ist, darauf gewiß allezeit göttliche Rache folgt, vornehmlich, wo wir uns nicht selbst richten: wie auch wiederum die Wohlthaten oder guten Werke nimmer unbelohnt bleiben.

245. Also hat die Blutschande Rubens damals noch geschlafen, und er begehret jetzt, daß er mit dem Vater derhalben möge versöhnt werden; aber es ist noch nicht Zeit, den Vater zu versöhnen. Darum hat er auch bei ihm nichts erlangen können. Zuvor, da er sich der greulichen Sünde wohl hätte enthalten mögen, hat er es nicht gethan: jetzt, da er sich nun bei dem Vater verdient machen will, will er es von ihm nicht annehmen. Es ist aber vor GOtt gleichwohl ein gut Werck gewesen, daß er sich beflissen, wie er seinen Bruder erretten und also entschuldigt sein möchte, wie Juda auch thut; aber die zwei sehr bösen Buben, Simeon und Levi, sind die rechten Anstifter und Rädelsführer gewesen dieser großen Uebelthat; und derhalben hat Simeon auch am allerschwersten büßen müssen. Denn er wird hernach in das Gefängniß geworfen werden, darum daß er der Vornehmste war, so diese Sünde angefangen. Und ich halte dafür, daß Christus von niemand anders getödtet worden sei, denn von diesen zwei Stämmen, von den obersten Regenten und Schriftgelehrten. Die Hohenpriester sind vom Stamm Levi, aus dem Stamm Simeon aber die Schriftgelehrten gewesen. Derhalben, wie sie ihren Bruder Joseph geplagt und gemartert, also haben hernach ihre Nachkommen Christum auch gekreuzigt, und sind diese Zwey mit allen ihren Nachkommen die Boshaftigsten gewesen.

246. Hieronymus sagt an einem Orte, daß in der Erzählung der zwölf Geschlechter Simeon ausgelassen werde, um des Verräthers Judas willen. Von den Hohenpriestern ist es gewiß, daß sie vom Stamm Levi hergekommen sind; aber Hannas und Caiphas kommen her von Simeon, und dazu auch Judas, ihr Vorgänger, wie ihn Petrus nennt Ap. Gesch. 1, 16. Das ist nun die Entschuldigung Rubens, welcher für Joseph gebeten, daß sie ihn nicht wollten tödten. Nehmt ihm das Leben nicht, hat er gesagt.

\quellseite{293}

Also sind Isaschar, Sebulon, Gad. Alle auch unschuldig; aber da sie durch das Ansehen und Gewalt ihrer ältesten Brüder verletzet sind, haben sie auch mit eingewilligt. Wie zur Zeit Christi die Hohenpriester das gemeine Volk der Juden auch an sich gehänget und zu ihrer Meinung gebracht, daß es gerufen hat, man sollte ihn immerhin kreuzigen. Nun folgt aber der rechte harte Kampf, so sie mit Joseph gehalten.

\begin{quote}
\textbf{V. 23. 24. Als nun Joseph zu seinen Brüdern kam, zogen sie ihm seinen Rock mit dem bunten Rock aus, den er anhatte; und nahmen ihn, und warfen ihn in eine Grube; aber dieselbige Grube war leer, und kein Wasser drinnen.}
\end{quote}

247. Die Brüder folgen dem Rath Rubens, und werfen Joseph in die leere Grube, darin kein Wasser war, welches an diesem Orte von Mose kürzlich beschrieben wird. Hernach aber, da sie aus Egypten gekommen sind, wird es Ruben selbst etwas weitläufiger wiederholen mit angehängter ernster Strafe seiner Brüder, welche also lautet 1 Mos. 42, 22.: „Sagt ichs euch nicht, da ich sprach: Versündiget euch nicht an dem Knaben, und ihr wolltet nicht hören?“ Es wird ohne Zweifel Joseph ein sehr schöner Jüngling von siebzehn Jahren gewesen sein; welches die rechte Blüthe der Jugend ist und die allerlieblichste Zeit des Lebens. Aber plötzlich, da er sich dessen nicht besorgt, haben ihn seine Brüder hingerissen und stracks auf ihn eingestürmt. Du Verräther, werden sie gesagt haben, du Bösewicht, du mußt sterben. Und wiewohl er ihnen zu Füßen gefallen, sich vor ihnen geneigt, mit gefalteten Händen jämmerlich geweint und geheult, ihnen gefleht, sie um GOttes willen und, wie man sagt, vom Himmel zur Erde gebeten, dazu auch um der Liebe und Ehrerbietung willen, so ja Kinder ihren Eltern zu leisten schuldig sind: sieht er doch, daß weder sein flehentlich Bitten, noch bittere Thränen bei ihnen eine Statt haben gefunden; wie sie das hernach selbst bekennen, da sie unter einander sprachen, 1 Mos. 42, 21.: „Das haben wir an unserm Bruder verschuldet, daß wir haben die Angst seiner Seele, da er uns flehete, und wir wollten ihn nicht erhören; darum kommt nun diese Trübsal über uns.“ V. 22.: „Nun wird sein Blut gefordert.“

248. Bedenket doch, welch eine erbärmliche Gestalt das gewesen, da der Knabe seine Hände und Augen gen Himmel aufgehoben und seine Brüder mit bittern Thränen und herzlichem Seufzen gefleht hat, da er gesehen, daß er von ihnen umringt und jetzt alsbald hat sollen erschlagen werden. Denn wie ist das so ein großer Jammer, nicht von Fremder oder Feinde, sondern seiner eigenen Brüder Händen also dahinfällen und erschlagen werden! Aber diese Brüder sind noch härter denn ein Kieselstein oder ein Diamant: sie lassen sich weder die Angst noch das Seufzen oder Heulen des Knaben erweichen, da man doch erfahren, daß oftmals Feinde dadurch sich haben bewegen und erweichen lassen.

249. Derhalben kann man die Grausamkeit dieser zwei Brüder nicht groß genug machen, welche also wider ihr eigen Fleisch und Blut gewüthet und getobt haben. Es ist zwar das Leiden und Kreuz unsers Herrn Christi auch betrübt und grausam genug gewesen; wiewohl wir wissen, daß ihm solch Leiden nach dem guten Willen seines himmlischen Vaters widerfahren und auferlegt gewesen: aber dieser Brüder Grausamkeit ist noch viel größer. Denn es hätte ihnen ja geziemet, daß sie sich die natürliche Zuneigung zu ihren Brüder sollten haben bewegen lassen. Wie denn Ruben gleichwohl versucht hat, diese Zuneigung und Liebe, so bei ihnen gar erloschen, wiederum zu erwecken, da er für seinen Bruder gebeten, daß sie ihn wollten leben lassen und nicht tödten: aber er hat damit bei ihnen nichts ausrichten können. Denn wiewohl sie den Knaben nicht alsbald erwürgt, wie sie bei sich beschlossen, haben sie ihn doch gleichwohl in die Grube geworfen, welches viel betrübter und greulicher ist. Denn sie gedenken ihn also in der Wüste mit Hunger und Durst zu tödten. Wie viel erträglicher wäre es ihm gewesen, daß er sich hätte alsbald tödten lassen, denn in so viele langwierige Pein und Marter geworfen werden? außer dem, daß ihr Vornehmen und Anschläge von GOtt sind verhindert worden. Denn wir sehen sonst, daß die, so zu hart gepeiniget werden, auch begehren, daß sie nur bald getödtet werden, damit sie also je eher je lieber von der Marter, so sie leiden, los werden mögen.

250. Derhalben hat der Satan in diesen Brüdern, so zukünftig noch Hohepriester sollten werden, die äußerste Grausamkeit geübt. Denn Joseph ist eine Figur Christi gewesen, und ist hier an diesem Orte die Niederfahrt zur Hölle bedeutet worden, oder wie Sacharja Cap. 9, 11. sagt, zur „Grube, da kein Wasser innen ist“, daß er durch das Blut seines Bundes die Gefangenen in der Grube ausließe und erlösete.

251. Erstlich aber haben sie ihm den bunten oder ungefärbten Rock ausgezogen; nicht daß sie zu solchem Raube große Lust hätten, sondern ihren Vater, den frommen alten betrübten Mann, damit zu betrügen und zu verspotten. Denn dieser Rock, so zuvor ein Zeichen der Liebe gewesen, welche sein Vater zu ihm getragen, muß jetzt nun auch ein hartes und böses Zeichen werden, den armen Vater damit zu betrüben und zu martern, welcher an dem Rock große Freude gehabt und damit seinen Sohn Joseph mit sonderlichem Fleiß gezieret; in demselbigen hat er nun auch sollen getödtet werden. Denn also wird er hernach sagen: „Es ist meines Sohnes Rock, ein böses Thier hat ihn gefressen“ etc.

252. Da sie ihm nun den Rock ausgezogen, haben sie ihn in die Grube geworfen, lassen ihn da liegen, und denken: Nun ist er dahin. Denn das ist ihr Vorhaben gewesen, daß er in der Grube vor Hunger und Durst, dazu vor Stank und Unflath hat sterben sollen. Und ist wahrlich durch GOttes Vorsehung also geschehen, daß in der Grube kein Wasser gewesen, daß er nicht alsbald mit Wasser bedeckt würde: sonst wäre er im selbigen Augenblick erstickt und gestorben. Sie nehmen aber den Rath Rubens darum an, daß sie ihn etliche Tage martern und ihre gefaßte Bosheit und Haß nur wohl erfüllen mögen.

253. Auf solche Weise regieret der Teufel im Hause des sehr heiligen Patriarchen Jakob unter den Kindern GOttes, wie S. Ioh. sagt. Unterdeß aber ist Benjamin ein kleines Kind gewesen und hat noch in der Wiege gelegen. Denn innerhalb zweier Jahre sind diese Dinge alle geschehen, nämlich, daß Dina geschwächt, Rahel und Debora gestorben, Ruben die Blutschande begangen und Joseph verkauft worden ist. So ist immer ein Unglück über das andere gekommen. Ruben kann wohl entschuldiget werden, wie droben gesagt worden, dazu die Andern auch, als, Naphtali, Isaschar etc. Und ich halte dafür, daß sie auch unschuldig gewesen sind. Aber Hannas, Caiphas, Judas, das sind die rechten Meister und Stifter dieser Sünde, die auch bei sich selbst ihr Frohlocken darüber gehabt, und ihrem Herrn, dem Satan, dafür gedankt, daß ihnen ihre Rathschläge so fein glücklich fortgegangen, sind fröhlich und getrost über solchen sehr bösen Sachen, gedenken dieweil ihres alten Vaters und Josephs, ihres Bruders, gar nicht, welcher unterdeß jämmerlich geseufzt und geweint, ja, im Tode und der Hölle gemartert wird ohne allen Trost oder Hoffnung der Erlösung. Denn da ist niemand, der den betrübten Knaben getröstet und ihn ermahnet, sich auf GOttes Güte, Hülfe und Beistand zu verlassen, der da gewißlich sein Schutz und Schirm sein würde in der allerhöchsten Angst und Noth; sondern er hat müssen im äußersten Jammer der Einsamkeit und der Hölle stecken bleiben. Den Rock aber haben sie mit sich genommen, daß derselbige noch des frommen Vaters Gift und tödtliches Verderben sein soll, welcher ihn zuvor so oft und so gar lieblich erfreut hatte.

254. Davon wird aber hier im Text nichts gesagt, ob sie das Kleid zerschnitten oder zerrissen haben oder nicht, ehe sie es in das Blut getunkt. Denn haben sie es nicht zerrissen, so haben sie wahrlich sehr närrisch gehandelt, und ist ihnen widerfahren, wie es gemeiniglich pflegt zu gehen und wie man im gemeinen Sprüchwort sagt, daß allezeit der Sünde oder gottlosen Wesen die Thorheit anhängt. Und Hilarius sagt: Wo das gottlose Wesen sich selbst so klüglich rathen könnte, wie gar kühn es ist, so wäre es schwer, sich der Wahrheit anzunehmen und dieselbe zu schätzen. Darum sagt man auch im Sprüchwort: Mendacem oportet esse memorem et intelligentem, das ist: Wer lügen will, der muß ein gut Gedächtniß haben und muß verständig sein. Haben sie aber nun den Rock, den sie in das Blut tunken wollten, ganz bleiben lassen, ist es eine große Thorheit gewesen, damit sie sich selbst verrathen haben. Denn es hätte der Vater also schließen können: Siehe, sie sagen, es habe den Joseph ein böses wildes Thier gefressen; wie, daß es ihm denn nicht auch den Rock zerrissen hat? oder wo der Rock wäre zerrissen gewesen, hätte er doch mögen denken, man sollte ja billig gesucht oder geforscht haben, ob das Thier nicht etwas von seinen Gebeinen oder Fleisch hätte lassen übrig bleiben.

255. Also, sage ich, muß das gottlose Wesen und die Lüge sich immer selbst verrathen. Und wo Jakob in seinem Herzen zufrieden und getrost gewesen, wäre ihm dies auch ohne Zweifel eingefallen. Aber der gute Mann ist so erschrocken und gleichsam gar von Sinnen gewesen, als sollte er jetzt dahin sterben; wie die Leute, wo sie plötzlich erfreut oder aber betrübt werden, auch nicht gar bei sich selbst sind, also daß ihnen alle Sinne und Verstand verschwinden, und sie also erschrecken und erfahren, daß sie weder sehen noch denken können. Eben also ist es Jakob auch gegangen, da ihm erstlich die Nachricht gekommen, wie sein Sohn Joseph umgekommen, daß er deshalb mit Wehklagen und großer Traurigkeit so gar überfallen und gleichsam bedeckt worden, als wäre er gar verschlungen. Er ist nicht bei sich selbst gewesen, daß er auch aus Wehmuth wird gesagt haben: O Herr Gott! ist er todt! Denn sonst hätte er mögen fragen, da ihm der Rock ist gebracht worden: Ob denn so gar nichts wäre übrig geblieben von dem Leichnam Josephs, und wie doch das wäre zugegangen, daß noch das Kleid ganz geblieben; sintemal die bösen wilden Thiere den Leib mit dem Kleide zu zerreißen pflegen. Ja, er würde auch gebeten haben, daß er möchte an den Ort geführt werden, da doch etliche Spuren solches greulichen Todes möchten gefunden werden. Aber ein Herz, das erschrocken ist, bedenkt solches nicht, kann auch darauf keine Acht geben.

\quellseite{294}

256. Derhalben ist dies ein trefflich Exempel der großen Thorheit an den Lügnern, welche mit sich selbst immer eins sind, und derhalben leichtlich von denen, so auf ihre Dinge fleißig Acht geben, ergriffen und gestraft werden können. Diese Brüder meinen, sie haben traun allhier einen sehr feinen und klugen Rath erfunden; aber wenn du es recht bedenkst, ist es gar ein sehr närrischer Rath, und der mit viel Gefahr verbunden war. Sie sollten darüber um Leib und Leben gekommen sein, wo nicht der gute alte Vater zu leichtgläubig gewesen wäre.

\subsection{II}

\begin{quote}
\textbf{V. 25. Und setzten sich wieder zu essen. Indes huben sie ihre Augen auf, und sahen einen Haufen Ismaeliter kommen von Gilead, mit ihren Kameelen; die trugen Würze, Balsam und Myrrhen, und zogen hinab in Egypten.}
\end{quote}

257. „Sie setzten sich wieder zu essen“, sagt Mose, als hätten sie es wohl ausgerichtet. Ihr Gewissen ist gar sicher, die Sünde schläft und ist noch verborgen: GOtt aber hat es alles wunderbarlich geschickt und geordnet. Der arme elende Vater sitzt daheim, weiß nicht, wie es mit seinem Sohn steht: so kann Joseph selbst auch nicht sehen, was sein Unglück für einen Ausgang oder Ende gewinnen werde. Allein GOtt sieht alles, und spielt also auf das allerfreundlichste mit dem Vater und auch mit seinem Sohn; denn es ist eigentlich ein solch Spiel, das voll göttlicher Barmherzigkeit und Güte ist; uns aber ist es ein sehr schwer und betrübt Spiel, so es doch in uns eine ewige und über alle Maßen wichtige Herrlichkeit schafft, 2 Cor. 4, 17.

258. Dieweil sie nun da gesessen, läßt GOtt selbst eine Gelegenheit sich bieten, Joseph aus der Grube zu erretten, und ihm gleichsam aus der Hölle und Tod wieder zu rufen. Wiewohl der Jammer auf ihm liegen bleibt, daß er gleichwohl noch gefangen und weltlich oder bürgerlich gestorben ist; sintemal er von seinem allerliebsten Vater geschieden und gerissen ist, wie der Vater selbst auch solches Unglücks halben beinahe dahin stirbt, nämlich, weil er seines Sohnes, den er sonderlich lieb gehabt, muß beraubet sein. Denn wo ihm nur der Sohn gesund und lebendig geblieben wäre, hätte er den Unfall oder Schaden an den andern Söhnen und aller seiner Habe und Wohlfahrt desto leichter ertragen können. Denn er ist ja der Erstgeborne gewesen, an welchem dem Vater alle Hoffnung und Wohlfahrt gelegen war, von dem man hoffte, daß viel Erben und Nachkommen sollten geboren werden, daher beide die Kirche und äußerliches Regiment sehr wachsen und zunehmen würden. Derhalben ist das ganze Geschlecht Jakobs mit erbärmlichen Wehklagen und großer Traurigkeit, so bitterer und härter ist denn der Tod selbst, beschweret und beladen gewesen.

259. Weiter: die Ismaeliten, so nun den Joseph gekauft, sind Josephs und seiner Brüder Blutsfreunde gewesen; denn Ismael und Isaak sind Brüder gewesen. Darum sind diese nun auch Brüder und Vettern im dritten und vierten Gliede; sie sind Vettern und nahe Freunde. So sind die Midianiter auch ihre Vettern gewesen im selbigen Grade, und scheint es, als habe der alte Haß noch unter diesen Blutsverwandten gewähret. Denn sie haben Joseph bald und mit großer Begierde gekauft, und sind Kaufleute gewesen, die mit Würze, Balsam und Myrrhen umgegangen und gehandelt haben.

260. Was aber nun eigentlich diese drei Worte in der hebräischen Sprache heißen oder bedeuten mögen, würden wir nicht verstehen können, wo wir es nicht aus 2 Mos. 30, 23. ff. könnten etlichermaßen schließen. Sie haben Krämerei getrieben mit Würze, und das Land ist sonderlich zu der Zeit von GOtt reichlich gesegnet gewesen, daß es viel Balsam, Myrrhen und Cinnamet gehabt etc., wie das zu sehen ist 2 Mos. 30, 23. ff., da beschrieben wird das köstliche oder heilige Salböl, welches von Balsam, Cinnamet, Kalmus und Casien gemacht worden. Wie aber diese Wörter zu unterscheiden seien, und welche eigentlich das Genus und welche die Species bedeuten, ist mir nicht bekannt.

261. Das Wort nechot haben Etliche Würze, Etliche sonst ein gewünscht köstlich Ding, etliche Andere auch wohl Wachs gedeutet und ausgelegt.

262. Das Wort zeri heißt Harz, ist ein gemein Wort für alles, so aus den Bäumen fließt oder trieft; wie an unsern Kirschbäumen zu sehen. Aber das ist wild Harz, Weihrauch aber und Myrrhen ist das allerköstlichste. Das Pech ist unter allen das gröbste, damit man die Fässer begießt. Unser Agstein oder Bernstein ist auch Harz; denn es trieft auch aus einem Baum und wird hart. Wie vielerlei Art sie gehabt haben, wissen wir nicht.

263. Der Balsam ist aller Harze oder Gummi, welcherlei Art es auch sein mag, König. Die Schrift nennt es ungeschnittenen Balsam, das ist, solcher Balsam, so nicht durch Oeffnung oder Schneiden aus dem Baum geflossen. Wie denn auch die Aerzte und Historien-Schreiber einen Unterschied machen, und sagen von diesem Balsam, daß er wegen der sehr großen Fruchtbarkeit des Baumes von selbst herausfließt.

264. Der Balsamfast, oder der erste Balsam, so von selbst und ohne Oeffnung des Baums herausfließt, ist der allerköstlichste und edelste Balsam. Also wird im Hohenliede Salomos die außerwählte Myrrhe genannt, das ist die, so erst heraustrieft, so nicht herausgedrückt wird, sondern durch die Rinde des Baumes herausgeflossen und nicht mit Gewalt heraus gezwungen ist. Und dasselbige Wort steht auch im Propheten Jeremias Cap. 8, V. 22.: Ist denn keine Salbe in Gilead? Oder ist kein Arzt nicht da? Warum ist denn die Tochter meines Volks nicht geheilet?“ Die Andern haben es gegeben: Tropfen Weihrauch oder Balsam. Der Prophet aber versteht daselbst die Würze oder Arznei, so aus den Bäumen gekommen, damit Jerusalem hätte können geheilet werden, wie im 51. Cap. v. 8. steht, da er der Babylonier troppet und sagt: „Nehmet auch Salben zu ihren Wunden, ob sie vielleicht möchte heil werden.“

265. Das Wort lot hat auch mehr denn einerlei Bedeutung. Denn bisweilen deuten sie es, daß es ein Saitenspiel sein soll, in welchem Verstande es oft in den Psalmen gebraucht wird; bisweilen wird es auch wohl gedeutet, daß es eine Eichel oder aber Kastanie heißen soll etc. Es sei nun gewesen, was es wolle, wir verstehen es von einem edlen Saft, so aus den Bäumen fließt, und lassen uns genügen an gemeiner Erkenntnis oder Verstand dieser Dinge, nämlich, daß diese Leute Würzkrämer gewesen, so nach Egypten köstliche Salben und andere wohlriechende Dinge gebracht, welche man zur Arznei gebraucht; wie unser Agstein oder Bernstein große Kraft hat, wenn die Weiber gebären sollen, und wider den Stein und Schlag. Hat nun solcher Saft bei uns, so doch etwas gröber ist, so große Kraft, wie viel mehr ist die Würze und Saft derselbigen Lande den Menschen zu ihrer Gesundheit nützlich und gut gewesen. Und hernach hat Jakob seinen Söhnen befohlen, daß sie von den Früchten des Landes etwas mitnehmen und dem obersten Regenten in Egypten bringen sollen, 1 Mos. 43, 11., welche Früchte in Egypten entweder gar nicht, oder aber weniger gewachsen, und derhalben desto angenehmer daselbst gewesen sind. Also haben auch die drei Weisen, Matth. 2, 11., dem Kindlein Jesu Gold, Weihrauch und Myrrhen geschenkt. Wir haben in Deutschland der Dinge keines, oder was wir davon haben, ist alles verfälscht und verderbt.

266. Nun folgt auch weiter im Text die Entschuldigung Judas, welcher der Gelegenheit, oder des Mittels, so sich dargeboten, gebraucht und seinen Bruder wahrhaftig vom Tode errettet hat.

\quellseite{295}

\begin{quote}
\textbf{V. 26. 27. Da sprach Juda zu seinen Brüdern: Was hilft’s uns, daß wir unsern Bruder erwürgen, und sein Blut verbergen? Kommt, lasset uns ihn den Ismaeliten verkaufen, daß sich unsere Hände nicht an ihm vergreifen; denn er ist unser Bruder, unser Fleisch und Blut. Und sie gehorchten ihm.}
\end{quote}

267. Juda hat etlichermaßen den grimmigen Zorn seiner Brüder gelindert. Es mag ihn vielleicht etwas gereut haben, daß er zuvor in Josephs Tod auch mit gewilligt hat; wiewohl es zwar eine schlechte Frömmigkeit ist, so er damit vorgewandt. Er bekennt frei, daß Joseph ihr Bruder, und ihr Fleisch und Blut sei; und hat doch mit solcher Rhetorik sie beredet, daß sie etwas gelinder mit ihm sind umgegangen. „Er ist“, sagt er, „unser Bruder, unser Fleisch und Blut“, und ist der Sohn, der unserm Vater sonderlich lieb ist. Welches denn wahrlich sehr große und starke Argumente sind, daß auch schwerlich in der ganzen Welt andere stärkere möchten gefunden werden bei vernünftigen Leuten und so noch gemeinen Witz und Verstand haben. Denn sie begreifen in sich alle natürliche Zuneigung oder Liebe, so menschlicher Natur eingepflanzt.

268. Und sollten derhalben diese Brüder sich billig durch diese Rhetorik und Dialektik haben bewegen und erweichen lassen, daß sie von ihrem bösen Vornehmen abgestanden wären. Aber das hat sie mehr bewogen, daß er sagt: „Lasset uns ihn den Ismaeliten verkaufen, daß sich unsere Hände nicht an ihm vergreifen“ etc. Er hat sie nicht gar stillen oder bereden können, daß sie wären anders gesinnt geworden und von ihrem Vornehmen abgestanden. Sie haben zwar so viel zugegeben, daß sie etwas daran ändern wollen, aber doch nur insofern, daß der leibliche Tod in den bürgerlichen Tod möge verwandelt werden. Denn sie bleiben gleichwohl noch bei diesem Vorhaben, daß der, so frei geboren, weggeführt, zur Dienstbarkeit gebracht und verkauft und also von seinen Eltern gerissen werden soll. Nun möchte oftmals Einer viel leichter den leiblichen Tod denn solchen bürgerlichen Tode und Dienstbarkeit begegnen möchte: und ist ein stetiges Kreuz, ein unendlicher Jammer und Dienstbarkeit, daß endlich ein solcher Mensch von seinen lieben Eltern, Brüdern und ganzem Geschlecht muß abgezogen und weggerissen sein. Mir zwar wäre es viel leidlicher, an Leib und Leben gestraft werden, denn also in ewiger Dienstbarkeit und Elend leben.

269. In diese Angst und Noth sind nun diese heiligen Patriarchen, Jakob und sein Sohn Joseph, geworfen; denselben wird ein solch schwer Kreuz durch die heimlichen Tücke und List dieser Brüder zugerichtet. Es sind schändliche Zimmerleute, die ihrem eigenen Vater und Bruder Joseph dies schwere Kreuz zimmern und schmieden. Das soll niemand annehmen und beschönigen. Das soll niemand leiden, denn solche fromme heilige Eltern.

270. Gott schweigt dazu noch immer still, als sähe er solches nicht, als ob er dieses Raths Erdichtern auch selbst hülfe, und sieht doch, daß sie so Böses vorhaben. Warum leidet er aber das? Warum schlägt er sie nicht mit Blitz und Donner, und tritt solch ungöttliches Vorhaben mit Füßen samt denen, so es gestiftet und angefangen? Oder aber, warum läßt er nicht vielmehr die Mörder, Ehebrecher und Tyrannen also geplagt und gemartert werden, und verschont doch damit solcher heiligen Leute? Antwort: Gott will, daß wir betrachten und lernen sollen, wie groß die Liebe sei, so die Eltern zu ihren Kindern tragen, daß wir daraus auch abnehmen und verstehen lernen, wie groß die Liebe Gottes sei, damit er uns geliebt, daß er auch seinen eingebornen Sohn geliebt, daß für uns wollen lassen leiden und creuzigt werden. Denn Joseph ist ein Bild des Sohnes Gottes.

271. Man muß aber alhier gleichwohl nicht verschweigen die Heuchelei, so noch an Juda gefunden wird. Denn er bringt wohl ein gesundes und göttlich Argument vor, welches aber doch ungeschickt, ja, heuchlerisch und teuflisch ist. „Daß sich“, sagt er, „unsere Hände nicht an ihm vergreifen“ etc. Diese heuchlerische pharisäische Theologie hat von Anfang der Welt gewähret und regiert. Denn also hat Kain auch für den nicht gehalten werden wollen, der seinen Bruder Abel erwürget hatte, da er 1 Mos. 4, 9. sagt: „Soll ich meines Bruders Hüter sein?“ Ei, wer wollte seinem Bruder so thun? Desgleichen sagt auch Saul von David, 1 Sam. 18, 17.: „Meine Hand soll nicht an ihm sein, sondern die Hand der Philister.“ Also steht auch in der Historie des Leidens Christi, Joh. 18, 28., daß die Juden nicht haben wollen in das Richthaus gehen, damit sie nicht unrein würden; oder daß sie ja nicht möchten für die Leute gehalten werden, so einen Menschen wider Recht und Billigkeit tödten wollten.

272. Darum sind diese talmudischen und pharisäischen Meinungen vom Anfange auch in diesem heiligen Hause und Gemeinde Jakobs geblieben. Denn daselbst sind solche Theologen gewesen, die es dafür gehalten, daß der kein Todtschläger wäre, welcher sich mit der Hand des Todtschlages enthalte. Ja, haben doch diese Brüder gleichwohl Joseph mit ihren Händen in die Grube geworfen; warum dichten sie sich denn noch solchen Trost und großen Schein der Heiligkeit? Warum, daß sie an ihren Bruder die Hände nicht gelegt, sondern ihn nur verkauft haben. Ei, welch eine schöne Gerechtigkeit ist mir das, wenn du Einen nicht mit der Hand, aber doch gleichwohl mit deinem Willen und Begierde, mit Rath, Hülfe und Einwilligung tödtest! Lasset uns ihn verkaufen, sagen sie, so werden wir unschuldig sein. O Juda, du bist noch nicht rein. Das ist nur auf gut talmudisch oder jüdisch geredet.

273. Aber also pflegt es in der Welt zuzugehen. Es müssen auch in der heiligen Kirche und Gemeinde noch Heuchler sein. Es ist derhalben nicht Wunder, daß auch zu unsern Zeiten dergleichen noch so viel sind; und hat doch Christus, unser HErr, wie in der Evangelisten Historie zu sehen, keinerlei Leute härter gestraft als eben die Heuchler. „Wehe euch Heuchlern“ etc., sagt er Matth. 23, 13. ff. Mit den armen Sündern aber hat er gegessen und getrunken, freundlich mit ihnen geredet und umgegangen; hat bei ihnen Wunderwerke gethan; aber mit den Heuchlern hat er gar keine Gemeinschaft. „Siehe“, sagt David im 51. Psalm V. 8., „du hast Lust zur Wahrheit“, bist feind der Heuchelei. Auch ist Christus hernach selbst von den Heuchlern getödtet worden.

274. Ja, so sein thun sie Buße, diese heiligen Väter: sie verändern den leiblichen Todtschlag in einen bürgerlichen, und ziehen eine Mosaische Decke vor die Augen mit diesem seinen Schein, daß sie gleichwohl die Hände an ihn nicht gelegt; so sie doch daran viel schwerer sündigen und haben damit auch Josephs Jammer und Trübsal gemehret. Also soll es den Gottlosen gehen, wenn sie sich schmücken wollen, daß sie es liebermal ärger machen. Denn nach ihrem Gutdünken sind sie die, so allein wollen weise sein. Gott verstehet ihre Heuchelei nicht, kann so scharf nicht sehen, ihr Vater Jakob sieht es auch nicht, wird es auch nimmermehr inne werden. Also haben sie sich selbst schmeicheln können; und je mehr sie sich schmücken und ihre grausame Sünde zu bedecken versuchen, je mehr sie ihre eigene Grausamkeit an den Tag geben.

275. Mose aber sagt gar nicht, wie sich Joseph dabei gebärdet und was er in solcher großen Gefahr geredet habe: er redet von ihm nicht anders als von einem Stummen, oder von einem Stein oder Klotz, der nichts redet, nicht gebeten, nicht geschrien, der gar nichts dazu gethan, damit er seine Brüder hätte stillen und versöhnen mögen; sondern dies alles hat Mose dem Leser befehlen wollen zu bedenken, darum daß auch mit Worten nicht Alles genugsam hat ausgedrückt oder beschrieben werden können. Was aber nun die betrübten elenden Gebärden, das erbärmliche Weinen, die traurigen Klageworte, dazu das Seufzen und Schluchzen Josephs, damit er sie flehentlich gebeten, gewesen sein mögen, mag ein Jeder bei sich selbst bedenken und erwägen. Denn er ist von Natur fast zart gewesen, und hat sonderlich eine große Zuneigung und Liebe zu seinem Vater gehabt, von dem er also plötzlich und unversehens muß hingerrissen werden; wird dazu beraubt des Kleides, so ihm sein Vater gegeben; wird in die Grube geworfen; ist wiederum herausgezogen und verkauft, und seinen allerliebsten Eltern weit aus den Augen und Armen gerückt worden, daß er auch gar keine Hoffnung mehr gehabt, daß er wieder zu ihnen kommen möchte. Lieber Gott! welch ein Seufzen und Weinen wird da gewesen sein? Ach mein lieber Vater, wie geht es mir! O wie wirst du schreien und klagen!

276. Diese Stücke sind in der Historie Josephs sehr traurig und sollten billig an den Tag gegeben worden sein. Aber Mose geht da vorüber, und zeigt damit an, daß Gott selbst auch dazu stille geschwiegen. Hernach aber werden solche Stücke von den Brüdern Josephs selbst mit großem Schmerz erzählt werden, da sie sagen, 1 Mos. 42, 21. 22.: „Wir haben die Angst seiner Seele, da er uns flehete, und wir wollten ihn nicht erhören; darum kommt nun diese Trübsal über uns“ 2c. „Nun wird sein Blut gefordert“ 2c. Er hat sich wahrlich jämmerlich gestellt, hat nicht dazu gelacht oder still geschwiegen.

\quellseite{296}

277. Und ist dies wahrlich eine sehr erschreckliche Tragödie gewesen ohne alle Hülfe und Hoffnung der Erlösung. Jakob, der Vater, ist nicht dabei, weiß nichts von den Dingen, so daselbst geschehen, so doch dieser ganze Jammer endlich allein auf ihn kommen wollte. Also wird das Unglück immer größer und hat Joseph auf vielerlei Weise sterben müssen. Siehe, wird er gesagt haben, ich werde dahingerissen, und muß Hebron, dazu meinen armen elenden Vater verlassen, der mein nun auch wird beraubt sein müssen! Denn Joseph wird dazumal schon vorgesehen haben, was Jakob, sein Vater, über ihn für Leid tragen und was für Traurigkeit er haben würde. Also bedeckt und beladet demnach Gott die Heiligen und Gläubigen, und die, so er am allerliebsten hat, mit großem Jammer und Trübsal. Und dies soll man auf Christum ziehen, daß wir bedenken, was das sei, daß Gott für uns seinen Sohn dahingegeben, daß er ist in die Grube geworfen, darin kein Wasser ist, Sach. 9, 11. Es haben die Propheten nicht vergeblich diese Dinge also hoch gerühmet, und sich verwundert, daß der Sohn also in der Egypter Hände dahin ist gegeben worden.

278. Danach muß man auch an diesem Ort eine Frage abhandeln, so zur Grammatica gehört, nämlich: Ob die Ismaeliten eben dieselbigen seien, so auch Midianiter genannt werden? Welche Frage von Augustin und Andern allhier ohne Noth ist abgehandelt worden. Ich wollte zwar lieber, daß andere Fragen möchten erregt werden, so wichtiger und nützlicher sind, nämlich, von Gottes Willen und von seinen wunderbaren Werken, vom Kreuz und Glauben der lieben Väter, und was doch Gott damit meine, daß er die Heiligen und Gläubigen immer mit so großem Jammer also übt, nämlich, uns zu Trost und daß er damit auch das Leiden Christi hat wollen abbilden. Aber doch wollen wir dasselbe auch besehen.

Lyra sagt, nach der Juden Meinung, Joseph soll dreimal verkauft worden sein. Erstlich haben ihn seine Brüder den Ismaeliten verkauft, zum Andern sei er von den Ismaeliten wiederum den Midianitern, und zum Dritten sei er dem Potiphar in Egypten verkauft worden. Und ich lasse es mir gefallen, daß die teuflische Bosheit dermaßen nur groß genug gemacht ist zur Verachtung der Heuchler, und den Neid dieser Brüder damit abscheulich zu machen. Aber ich halte dafür, daß er nur zweimal sei verkauft worden; wie das aus der Historie leichtlich zu verstehen, die da zeugt, daß Potiphar Joseph von den Ismaeliten gekauft habe. Ich will aber damit die andere Meinung doch nicht verworfen haben, daß er sollte dreimal verkauft sein. Denn von den Midianitern wird auch gesagt werden, daß sie ihn dem Potiphar verkauft haben. Aber die Frage von Keturah und Hagar, ob das Eine Frau gewesen, ist droben erklärt worden.

279. Mose hat es unter einander gemenget, daß er bald die Midianiter, bald aber die Ismaeliten nennet. Ich halte dafür, daß es sei Eine Gesellschaft gewesen, so doch mit zwei Namen genannt worden; wie es denn mit Kaufhändeln zuzugehen pflegt. Und bewegt mich dazu sonderlich diese Ursache, daß Mose zuvor gesagt, daß sie nicht Ismaeliter, sondern einen Haufen Ismaeliter gesehen haben. Denn das hebräische Wort, orchah, heißt eigentlich die, so über Land reisen, oder vielmehr einen Haufen Wandersleute oder Gefährten, wie Luc. 2, 44. von Christo gesagt wird: „Sie meineten aber, er wäre unter den Gefährten.“ Und geschieht es oftmals, daß auch Leute aus mancherlei Ständen und Nationen zusammenkommen, und mit einander über Land reisen. Eine solche Gesellschaft ist auch diese der Ismaeliter und Midianiter gewesen, so mit einander hantiert, gekauft und verkauft haben; wie denn solches auch wohl unter mancherlei Völkern, so weit von einander wohnen, zu geschehen pflegt. Es haben nun dieselbigen oder aber Andere den Joseph verkauft, daran ist wenig gelegen; aber das weiß man und ist offenbar, daß er eigentlich ist verkauft worden.

280. Und ist nicht so viel gelegen an dieser grammaticalischen Frage, als an dem, so sonst an diesem Ort soll betrachtet werden, nämlich: Wie doch die Ismaeliten oder Midianiter Joseph von seinen Brüdern mögen gekauft haben? Ob sie auch zuvor von Judas, das ist, von Simeon, dem Verräther, und Levi, das ist, von Caiphas und Hannas erforscht und gefragt haben, wer doch dieser Knabe wäre, den sie wollten verkaufen? oder woher derselbe in ihre Hände gekommen wäre? ob er gestohlen, oder jemand entführt sein möchte? Denn daran ist kein Zweifel, daß sie nach gemeinem Gebrauch der Kaufleute nicht allein nach der Kaufsumme, sondern auch wie die Waare beschaffen und wie groß dieselbe gewesen, fleißig werden geforscht haben. Und wo dem also ist, werden ohne Zweifel Judas, Hannas und Caiphas herausgefahren sein und Joseph, ihren Bruder, heftig und schwer verklagt haben. Und wenn die Kaufleute gefragt: Ist dieser euer Bruder, warum wollt ihr ihn denn verkaufen? werden sie darauf diese Antwort gegeben haben: Wenn dieser nicht ein Uebelthäter wäre, wollten wir ihn dir nicht haben übergeben. Also haben sie ihm allerlei Sünde und Uebelthat aufgeleget; wie die Juden dem Herrn Christo auch gethan vor Caiphas. Auf dieselbige Weise werden diese Gesellen auch geschrieen haben, dieser Joseph sei ein böser verwegener Bube, der sich immer wider seine Brüder gelegt, den Eltern ungehorsam gewesen, auch daheim im Hause Alles verwirret, einander geworfen und verdorben habe, und endlich den frommen heimlich nachgestellt, dem Vater nach dem Regimente und seinen Brüdern auch nach ihrer Nachfolge oder Erbgerechtigkeit gestanden: und wenn er die unterdrückt hätte, alsdann wollte er das Regiment allein haben. Wo dies also die Ismaeliten oder Midianiter angehört, und solche unrechtmäßige und falsche Anklage sich haben gefallen lassen, so haben sie wahrlich auch sich gegen das Fleisch und Blut Jacobs verschworen.

281. Die Brüder aber werden sich ohne Zweifel solcher Käufer höchlich gefreut haben, als die Feinde wären des Hauses und der Gemeinde Jacobs. Denn es sind sonst auch in der Nachbarschaft die Cananiter und Amoriter gewesen; denselbigen hätten sie den Joseph auch wohl verkaufen können: aber sie haben ihn lieber denen wollen aufbürden, die ihn fern von den Eltern wegführten, und ihn unfreundlicher und härter behandeln würden als die Nachbarn. Auf solche Weise wollte ich ja gerne die Werke Satans und die unerhörte Bosheit dieser Brüder nur groß genug machen. Wenn wir ihn den Nachbarn verkaufen würden, haben sie gedacht, wiewohl dieselbigen unserm Geschlecht auch feind sind, und eben so begierig würden gewesen sein, ihn zu kaufen, ist doch diese Gefahr dabei, daß es der Vater bald erfahren und ihn alsdann wiederum heimführen möchte. Derhalben ist dies eine sehr feine Gelegenheit, diesen unruhigen aufrührerischen Menschen nur weit hinweg in ferne Lande zu verkaufen, daß er hinfort nicht mehr Lärm anrichte und auch nicht mehr in des Vaters Hause nach dem Regimente stehe.

282. Dies ist das Bündniß des Pilatus und Caiphas wider Joseph; und ist mir recht, so hat der Heilige Geist durch David aus dieser Stelle den 2. Psalm gemacht, da er v. 2. sagt: „Die Könige im Lande lehnen sich auf, und die Herren rathschlagen mit einander wider den Herrn und seinen Gesalbten.“ Es haben Juden und Heiden zugleich mit einander gerathschlagt, die Midianiter und Ismaeliter, Herodes, Pilatus, Caiphas, Hannas, so zuvor einander spinnfeind waren. Denn wir haben droben, 1 Mos. 21, 14. 25, 1 ff. gehöret, wie Ismael aus Abrahams Hause verstoßen worden, und daß auch die Midianiter sammt den Kindern der Keturah mit ihren Gaben abgefertigt worden sind und am väterlichen Erbe kein Theil gehabt haben: daher ein solcher Haß unter diesen Völkern entstanden, so für und für gewähret, wiewohl sie alle von einerlei Geblüt geboren und hergekommen sind.

283. Also sind die Heiden und alle Völker mit stetigem Haß und Grimm immer erhitzt gewesen wider das herrliche heilige Haus und Geschlecht Abrahams, welches Haus noch immer die göttliche Verheißung behalten; wie das die erbärmliche Klage und das brünstige Gebet bezeuget im 83. Psalm v. 6–9. „Denn sie haben sich mit einander vereinigt, und einen Bund wider dich gemacht, die Hütten der Edomiter und Ismaeliter, der Moabiter und Haggariter, der Gebaliter, Ammoniter und Amalekiter, die Philister sammt denen zu Tyrus; Assur hat sich auch zu ihnen geschlagen, und helfen den Kindern Lot.“ Und allhier haben sich auch auf dieselbige Weise vereinigt und einen Bund mit einander gemacht die Brüder Josephs mit ihrem und ihres Vaters Feinden und haben noch dazu unter einander dieses Bundes halben höchlich frohlockt.

\quellseite{297}

284. Ist es aber an dem, daß zwischen ihnen solche Worte nicht gefallen und daß sie gar keinen Bund mit einander gemacht, so mögen die Ismaeliten und Midianiter etlichermaßen entschuldiget sein. Aber ich halte dafür, daß sie auch mit eingewilligt und sich mit den Brüdern Josephs wider ihn verbunden haben. Denn sonst würden sie ja solche greuliche That dieser Brüder gestraft und sich von ihnen gewendet haben, als die daran ein Mißfallen und Abscheu gehabt. Derhalben ist Joseph, seinen eigenen Blutsfreunden und Vettern verkauft worden, gegen die er auch ohne Zweifel die Hände gefaltet, und die er flehentlich um Gnade und Hülfe angerufen, gebeten und gefleht, sie wollten doch sein verschonen und ihn wiederum zu seinem Vater führen um deswillen, daß sie ihm ja Blutshalben so gar nahe verwandt wären. Ach Lieben, ich bin euer Vetter, helft mich erretten. Ich bin derer Bruder, so mich euch verkauft haben! Aber er findet sie nicht gelinder denn seine Brüder. Also verbinden sich die Ismaeliten und Midianiter mit Israel und tödten den Sohn Gottes, wie das im 2. Psalm beschrieben.

285. Mose hat zwar an diesem Ort das jämmerliche Geschrei und Weinen Josephs nicht erzählt; aber daran ist kein Zweifel, er wird kläglich geheult und gar demüthiglich und flehentlich gebeten haben, daß er möchte errettet, und daß doch seiner Brüder Zorn und Haß, so sie gegen ihn gefaßt, etwas möchte gelindert werden. Ach nicht, wird er gesagt haben, ach wehret doch meinen Brüdern! Da sollten wahrlich diese Ismaeliten und Midianiter seinen Brüdern widerstanden und ihnen gewehret haben, daß sie doch nicht so greulich wider diesen so frommen Jüngling wüthen wollten, sollten solche übel gewonnene Waare nicht gekauft, oder ja zum wenigsten den gekauften Knaben wiederum zu seinem Vater gebracht haben. Sie hätten billig sollen sagen: Lieber, was thut ihr doch jetzt? Simeon und Levi, warum wollt ihr doch also wider euer eigen Fleisch und Blut toben und wüthen? Dies ist ja eine unbillige That, ist wider alle menschliche Freundlichkeit und Liebe, so ihr ja billig vermöge eures Namens und Glaubens an eurem Bruder üben und beweisen solltet.

286. Aber so muß die Welt nicht thun. Denn sie ist nicht werth, daß sie etwas Gutes thun soll, sondern vielmehr, daß sie selbst Andern nur Schaden thue, oder aber, daß wo sie sieht, daß Andere Böses thun, sie alsdann sein durch die Finger sehe, und zu solchen Sünden nur helfe und dieselbigen fortsetze, daß sie ja also vollkommenlich ihrem Herrn und dem Gott dieser Welt dienen möge. So haben nun, sage ich, weder die Brüder noch Gefreunden, die Ismaeliten und Midianiter, Joseph erhören wollen, da er so kläglich geheult, gebeten und gefleht hat. Ja, Gott selbst hört ihn auch nicht.

287. Welches Exempel, wie oft zuvor gesagt, zu unserm Troste dient und gehört, daß wir damit auch in unserer Noth und Trübsal gestärkt und nicht so bald den Muth fallen lassen, wo etwa Schade oder Gefährlichkeit vorfallen. Denn wir haben uns greulicheres Unglücks von den Türken oder andern Feinden, sie seien auch wer sie wollen, nicht zu befürchten; und die zu unsern Zeiten Ueberfall leiden vom Türken, erfahren oder empfinden damit noch nichts, das dem Jammer und Trübsal dieser lieben Väter gleich wäre, und unser Leiden ist noch viel geringer denn derselbigen armen Leute, so täglich den Ueberfall vom Türken erwarten müssen. Denn das ist ja unter allem Unglück der schwersten eins, wo die Kinder von ihren lieben Eltern, oder aber wiederum die Eltern von ihren lieben Kindern gerissen werden. Das ist der Tod selbst, und ist viel härter, denn die allergewissesten Plagen, so die lieben Märtyrer erlitten, welche etwa nur Eine Stunde oder auf das höchste zwei mögen gewähret haben. Gott aber hat es also haben wollen, daß seine Heiligen und seine Kirche solche erschreckliche Trübsal leiden sollte, daß damit ihre Herzen erweckt werden möchten zu bedenken, wie ein groß Ding das gewesen, daß der Sohn Gottes selbst um unsertwillen eine Zeitlang auch hingereissen, daß ich es also nenne, in das Elend gesandt und in die Hölle geworfen ist, damit er uns also vom Elende der sehr schweren Sünde und ewigen Verdammniß erlösen möchte.

288. Derhalben werden allhier von Mose die Midianiter und Ismaeliter nicht vergeblich also mit Namen genannt. Denn er sagt nicht, die Kaufleute, die Amoriter etc., schweigt von den Namen der nächsten Nachbarn gar stille, nennt aber mit Fleiß Verwandte, daß er damit anzeige, wie sie alle in das Verderben ihres Vetters gewilligt, und daß aus diesem Exempel jedermann lernen möchte, wie groß die Bosheit der Welt sei. Denn also haben hernach die Juden auch gethan, wiewohl sie auch unsers Herrn Jesu Christi Gefreundte und Vettern gewesen, haben greulich wider ihn getobt und gemüthet, haben sich keinerlei Bitte oder Erbarmung bewegen lassen. Und das waren nicht Ismaeliten, sondern Israeliten, auf daß erfüllet würde, das durch den Propheten Micha, 7, 6., Matth. 10, 36., gesagt ist: „Des Menschen Feinde werden seine eigenen Hausgenossen sein.“ Denn da Pilatus und die Heiden sich des Herrn Jesu erbarmt und Mitleiden mit ihm gehabt, ihn für unschuldig gehalten, da sehen wir, daß seine Vettern, die Juden, greulicher sind, denn die Schlangen und die grausamsten wilden Thiere.

289. Dies ist nun das rechte Bild der christlichen Kirche zu allen Zeiten, welches Gott also in seinem Sohne Christo und seinen Heiligen vorgebildet hat. Wöllt ihr Christen sein, so schickt euch darein. Derhalben lasset uns lernen, Gott in solcher Trübsal und Widerwärtigkeit gehorsam sein und unsere Augen gen Himmel richten. Denn Gott hat ja nicht Lust an unserm Verderben, Elend und anderm Unglück; sondern dies ist ein starker Trost, so für und für währet, so im 36. Psalm V. 6. steht: „Herr, deine Güte reicht, so weit der Himmel ist, und deine Wahrheit, so weit die Wolken gehen“ etc.; und V. 9.: „Sie werden trunken von den reichen Gütern deines Hauses; und du tränkest sie mit Wohlust, als mit einem Strom.“ Und im 37. Psalm V. 18.: „Der Herr kennet die Tage der Frommen, und ihr Gut wird ewiglich bleiben“ etc.

290. Und das ist dazumal auch unsers Herrn Gottes Meinung gewesen mit Joseph und Jakob. Mein lieber Jakob, mein lieber Joseph, ich sehe wohl, was für Unglück ihr auf dem Halse habt, das euch drückt. Ich schlaffe nicht; will aber gleichwohl jetzund solch Unglück nicht wegnehmen und euch trösten. Denn es ist die Zeit der Erlösung noch nicht da, sondern ihr müßt zuvor erfahren und lernen, was der Teufel und was die Welt sei, was Kinder, was Brüder, ja auch, was der Tod sei, auf daß meine Gnade euch desto süßer und angenehmer sei, und daß ihr sehet, daß ich gleichwohl noch immer für euch gesorgt habe, auch in der Noth und mitten im Tode. Denn es wird euch nicht ein Härlein von eurem Haupte fallen ohne meinen Willen, sofern ihr nur nicht verzagt werdet, sondern also Geduld haben und das Kreuz tragen lernt. Ich will das, so ich euch verheißen, gewißlich und treulich halten. Ich habe verheißen, daß ich euch wolle segnen: das will ich auch halten, will euch nicht lügen. Euer Fleisch wird wohl murren und darüber ungeduldig werden, aber demselben widersteht und herrscht im Glauben und Warten auf die Erlösung.

\begin{quote}
\textbf{V. 28. Und da die Midianiter, die Kaufleute, vorüber reiseten, zogen sie ihn heraus aus der Grube, und verkauften ihn den Ismaeliten um zwanzig Silberlinge; die brachten ihn in Egypten.}
\end{quote}

291. Joseph wird noch um ein geringer Geld verkauft denn Christus; und ich halte dafür, daß dieses Kaufgeld bei zwanzig Thalern gewesen. Denn ich mag nicht allzu genau von diesen Silberlingen disputiren: aber es hat ohne Zweifel aus diesem Text Sacharja seine Prophezeiung genommen von Christo, da er Sach. 11, 12. sagt: „Und sie wogen dar, wie viel ich galt, dreißig Silberlinge.“ Denn die Sache und alle Umstände reimen sich daher sehr fein, und kann auch keine größere Aehnlichkeit sein, als zwischen Christo dem Gekreuzigten und Joseph. Das Verkaufen und Tod Beider kommt sein mit einander überein. Denn wie Jesaia Cap. 53, V. 8. von Christo sagt: „Er ist aus dem Lande der Lebendigen weggerissen“, also ist Joseph auch aus dem Lande und seinem Vater aus den Augen hingekommen worden, nicht anders, als ob er nimmermehr wiederkommen und seinen Vater sehen würde.

\section[Sechster Theil]{Sechster Theil. Wie Ruben sich verhalten, da er Joseph nicht in der Grube findet; wie die Söhne Jakobs den Rock Josephs zu ihrem Vater schickten und wie dem Vater hiebei zu Muthe}

\subsection{I}

\begin{quote}
\textbf{V. 29. 30. Als nun Ruben wieder zur Grube kam, und fand Joseph nicht darinnen; zerriß er sein Kleid, und kam wieder zu seinen Brüdern, und sprach: Der Knabe ist nicht da, wo soll ich hin?}
\end{quote}

\quellseite{298}

292. Das hebräische Wort, jeled, heißt ein kleines Kindlein, wie Jes. 9, 6. steht: „Uns ist ein Kind geboren.“ Dies ist aber nun eine andere Ermahnung, so durch Ruben geschehen, welcher sich wahrlich also erzeigt und vernehmen läßt, als ob es ihm ein Ernst gewesen, und der derhalben auch billig entschuldiget sein mag. Denn er hat sich durchaus nicht gefürchtet vor seiner Brüder Haß oder Unwillen, sondern schützet den Knaben Joseph nach bestem Vermögen, und zürnt, da er sieht, daß er weggeführt. Wo ist der Knabe? sagt er; „er ist nicht da!“

293. Das ist ihm aber also in den Sinn gekommen, daß er gedacht, er möchte vielleicht von Andern weggenommen sein, da die Brüder sein verschont und ihn nicht haben tödten wollen, entweder vom Teufel, oder von Mördern, oder aber von den Nachbarn, den Amoriten oder Cananitern, so zuvor auch den Hausgenossen Jakobs heimlich nachgestellt; oder daß ihn irgend ein Hirte gefunden, und den Sichemien oder andern Heiden, so bei ihnen gewohnt, übergeben hätte, ihn umzubringen. Also hat er gemeint, daß der Knabe umgekommen sei nach der Zeit, da ihn seine Brüder in die Grube geworfen. Denn es scheint, als habe Ruben sich gefürchtet vor unbilliger Gewalt und Haß der Nachbarn; und bekennt doch gleichwohl, er habe Joseph erhalten wollen, daß er möchte am Leben bleiben, verdammt auch die böse That seiner Brüder; ja, er beschuldiget sich selbst auch mit, da er sagt: „Wo soll ich hin?“ und wird hernach dieselbige Strafe noch einmal wiederholen. Nicht, daß er auf diese Weise die Sünde der begangenen Blutschande gebüßt habe; denn er steckt noch in schwerer Strafe, nämlich, daß er damit die Erstgeburt verloren und anderer großen Freiheiten mehr beraubt ist; sondern dieweil er diese Sünde erkennt, und sich hütet, daß er in der Andern Sünde nicht willigen und sich derselben nicht theilhaftig machen will, geht er vornehmlich damit um, daß er bezeuge, daß Joseph unschuldig sei, und daß er ihn also erretten und am Leben erhalten möge. Aber Simeon und Levi haben ihn gar verachtet um der Sünde willen der begangenen Blutschande: darum brechen sie durch und fahren mit ihrem Vornehmen fort.

294. Die Strafe aber, damit Ruben seine Brüder angefahren, ist gleichwohl billig und recht gewesen, da er sagt: Was habt ihr gemacht? Seht, der Knabe ist nun aus unsern Händen und Gewalt gekommen, und wenn wir jemand schon alle wollten, daß er möchte gesund, frei und los sein, könnten wir ihn doch mit keinem Gelde lösen, wie viel es auch sein möchte. Denn wo ihn vielleicht die Cananiter angetroffen, so werden sie ihn ohne Zweifel in ferne Lande wegführen, und entweder umbringen, oder aber doch mit erbärmlicher und ewiger Dienstbarkeit plagen.

295. Derhalben entschuldiget nun Mose Ruben, berichtet aber nicht, was seine Brüder ihm geantwortet. Aber es scheint, daß er doch heimlich so viel anzeige, daß Simeon und Levi ihn werden überredet haben, daß er still schweigen und mit den Brüdern zugleich in den greulichen Mord ihres Bruders Joseph willigen soll. Das alles aber miteinander ist aus der verfluchten Theologie hergekommen, welche darnach die Juden auch immer behalten haben, nämlich, daß sie gemeint, wer mit der Hand nicht todtschlage, der sei auch kein Todtschläger. Mit diesem Argument haben sie ihn überredet, daß er nur sollte still schweigen, auf daß er mit seiner Fürbitte den Lärm nicht um so viel größer machte, wo vielleicht seine Worte und dies Geschrei vor ihren Vater möchte gebracht werden, und haben ihn bedroht, wo er nicht würde mit ihnen diese That heimlich halten, würde er eben solche Gewalt und beständigen Haß seiner Brüder auch noch erfahren müssen. Und es wäre wahrlich kein Wunder gewesen, wo sie sich darüber auch selbst unter einander verwundet oder erschlagen hätten. Welches denn der Teufel ohne Zweifel auch mit allem Fleiß wird gesucht haben; er würde solches auch haben zuwege gebracht, wo es Gott nicht gewehret hätte. Derhalben schweigt nun Ruben still und willigt auch mit ein, da er gesehen, daß die Andern auch schweigen und dies verhehlen wollen, damit er nicht Ursache gebe, daß sein Vater desto mehr verwirrt oder betrübt werde, und seine Brüder unter einander nicht selbst auch darüber in weitern Zank gerathen. Aber auf diese Weise hat er sich nun auch theilhaftig gemacht dieses Mordes, so sie beide an ihrem Vater und Bruder begangen.

296. Wiewohl nun dies von Mose nicht so ausdrücklich beschrieben, wird es doch gleichwohl im Texte nicht undeutlich mit angezeigt. Denn er zerreißt sein Kleid und bekümmert sich heftig um solche Gefahr und Tod seines Bruders. Aber dieweil er sich eines großen Lärms oder Todtschlags und auch vor seinen Brüdern selbst gefürchtet, schweigt er still. Der Herr Scharioth, Simeon und Levi, Caiphas, Hannas werden ihn wohl auf das Maul haben schlagen dürfen, und sagen: Schweig, hüte deinen, oder dich soll alles Unglück holen. Denke darauf, daß vor dem Vater diese unsere That verschwiegen bleibe, damit wir nicht mögen dafür angesehen werden, als hätten wir unsern Bruder getödtet. Denn wo er es erfahren würde, möchte er uns alle miteinander alsbald aus dem Hause und Gemeinde verstoßen und in Bann thun; wie er doch hernach thun wird.

297. Derhalben ist das eine große Weisheit, so wir ihnen von Gott verliehen sehen, daß sie diese That vor ihrem Vater dermaßen verhelet haben. Denn alles Andere in der Historie ist böse und schändlich. Dieser einige Rath ist gut (wo doch etwas Gutes in dieser großen Sünde hat sein können), nämlich, daß sie den alten Vater des bereden mögen, so es sei von einem bösen Thiere zerrissen worden. Darum sind sie zu ihrem Vater getreten, und haben sich traun auch gestellt und gezeigt, als ob sie um seiner Betrübniß und seines Jammers willen selbst auch betrübt wären, und haben bezeugt, wie sie ein groß Verlangen nach ihrem verstorbenen Bruder hätten. Ja, sie werden vorgegeben und gesagt haben: Ach hätten wir unsern lieben Bruder noch! wollte Gott, er möchte noch am Leben sein! Aber es ist mit ihnen alles eitel Falschheit und Heuchelei gewesen; und ist doch gleichwohl diese ihre Heuchelei nicht so gar böse gewesen, als wenn sie dem Vater den ganzen Handel geoffenbart hätten.

298. Mit Ruben aber scheint es doch, als ob er seine Entschuldigung habe, wiewohl sie sehr schadhaft ist. Wie er denn hernach allein mit den Andern selbst etwas hart geredet: „Sagt ichs euch nicht, da ich sprach: Versündiget euch nicht an dem Knaben, und ihr wolltet nicht hören?“ etc., 1 Mos. 42, 22. Und es scheint wahrlich, als habe er ihnen weislich und wohl gerathen. Aber es liegt doch die Heuchelei noch darunter verborgen, davon ich zuvor gesagt, daß sie es dafür gehalten, es sei kein anderer Todtschlag, denn allein der, so mit der Hand vollbracht oder begangen wird. Hätte aber Ruben wollen gänzlich ohne alle Schuld sein, sollte er heimlich zum Vater gelaufen sein, und den erinnert haben, daß er fleißig wolle nachforschen unter seinen Söhnen selbst, welcher doch der sein möchte, so diesen Mord an Joseph begangen hätte. Aber er hat sich allzuviel vor ihnen gefürchtet und sich ihr Drohen gar überwinden lassen, daß er mit den Herren, Simeon und Levi, mit Judas, dem Verräther, und Caiphas endlich auch einwilligt. Denn dieselbigen haben auch über die Andern das Regiment gehabt, haben ihnen geboten und verbothen.

299. Darum haben Juda und Ruben, und die Andern alle zu dieser Uebelthat nicht wohl das Maul aufthun dürfen, und hat also Ruben auch sammt den Andern seine Seele und Gewissen mit diesem greulichen Tode beide seines Vaters und Bruders Joseph beschwert und befleckt. Inzwischen haben sie ihre Sache fein können schmücken, und den Vater betrügen, daß er von diesem Handel ja nichts erfahren möchte.

300. Ich möchte nun aber sehr gern wissen, was doch diese Brüder die ganzen zwanzig Jahre über mögen in ihrem Herzen gedacht haben? Denn so viele Jahre sind es gewesen von der Zeit an, da Joseph verkauft worden ist, bis Joseph hinab nach Egypten gezogen. Es ist wahrlich ein sehr großes Aergerniß gewesen, und ich halte dafür, daß sie die ganze Zeit über gar nichts Gutes haben thun, sich auch mit gutem Gewissen nicht einiges Seufzens mit Anrufung und Beten zu Gott, dem Vater im Himmel, vernehmen lassen können; das ist gewiß. Denn das Blut Josephs, der Tod ihres Vaters und das Wehklagen des ganzen Geschlechts hat ihnen immer im Herzen gelegen. Wo aber das Gewissen mit solcher großen Last beschweret ist, da kann es Gott nicht anrufen oder beten.

301. Ich zwar, wo ich mich einer kleinen Verletzung oder Leides schuldig wüßte, würde nicht können beten. Derhalben muß das Herz von allem Widerwillen, Haß und Neid gegen den Nächsten ganz frei sein, daß es niemand Böses, sondern jedermann Gutes gönne und gerne vergebe. Denn wo etwa unser Amt das also fordert, daß wir die Laster der Leute strafen müssen, die wir gern mit solcher Strafe wollten zur Buße bringen, soll man nicht denken, daß da die Liebe aufgehört. Denn solche Strafen geschehen nicht aus Haß, oder daß man damit jemand Schaden zufügen wollte, sondern allein der Meinung, daß man gern wollte, daß dadurch die Leute möchten gebessert werden.

\quellseite{299}

302. Darum sage ich, wenn man beten und Gott anrufen will, muß das Herz von allem Hasse frei sein; wiewohl wir die Sünde und Laster der Menschen hassen und sie verabscheuen müssen. Also beten wir auch recht für die gottlosen Bischöfe. Denn wir sollen ja nicht also gesinnet sein, daß wir ihnen Böses, nämlich ihr Verderben oder Verdammniß, wünschen wollten; strafen aber doch ihr gottloses Wesen und Sünde nur der Meinung und zu dem Ende, daß sie dadurch möchten gebessert werden.

303. Derhalben möchte ich wirklich gern wissen, was doch diese Brüder die zwei und zwanzig Jahre über mögen gethan haben, darin sie gesehen, daß ihr Vater immer im Sack und in der Asche gesessen, kläglich geweint und geheult hat um des Todes seines lieben Sohnes willen. Ich hätte das große und so erbärmliche Wehklagen so lange nicht können ansehen, hätte mich nicht enthalten können, daß ich nicht sollte gesagt haben: Ach wir armen unseligen Leute, was haben wir doch gemacht! Nun hören aber diese Gesellen ihren Vater täglich lehren und predigen, werden auch für Glieder der Kirche und des Hauses Jakobs gehalten: und ist doch ihrer keiner, der Gott anrufen könne, etwas Gutes zu thun; oder da sie gleich etwas Gutes thun möchten, ist doch solches eitel gottlos und verworfen Ding. Denn sie stecken noch immer tief in schwerer Todsünde, nämlich, dem Morde, so sie an ihrem Vater, Bruder und Großvater begangen hatten. Solch schändlich Ding kann die Heuchelei anheben und treiben; ja, die Heuchelei, sage ich, kann solche erschreckliche Dinge ausrichten und in so großer Bosheit noch verharren.

304. Wo Ruben innerhalb derselbigen Jahre gestorben wäre, hätte er vielleicht die Sünde erkannt, seinem Vater sie gebeichtet und um Gnade und Vergebung gebeten. Aber der andere Haufe würde die Sünde verhehlt haben und würde so hingestorben sein. Denn in so langer Zeit hat ihrer keiner Buße gethan, und sind mit der Zeit immer mehr und mehr, da sie im Gewissen die Sünde gefühlt, verhärtet worden. Wiewohl nun der Schrecken des Gewissens oftmals wiedergekommen, welches sie auch nicht gar haben tilgen oder unterdrücken mögen, also daß sie der Sünde sollten gar haben vergessen können, dieweil ihnen täglich vor ihren Augen und Herzen das Wehklagen und die bittern Thränen ihres alten Vaters gestanden: haben doch diese Söhne, so härter denn ein Kieselstein oder Diamant gewesen, solchen Schrecken gar verachtet; denn sie sind gar besessen und verhärtet gewesen vom Obersten der Teufel, sintemal sie so eine lange Zeit haben ertragen können, daß sie ihr Gewissen selbst solcher Sünde immer beschuldiget hat.

305. Nun ist aber daran kein Zweifel, wie die Schrift an andern Orten bezeuget, es müsse Gott heftig darüber gezürnt, und die verschmacht haben, so solche große greuliche Sünden angefangen und getrieben, nämlich, ihren eigenen Vater und Bruder zu tödten. Es ist derhalben ein hartes und erschreckliches Ding, solches also zu lesen und zu verstehen von den großen vornehmen Männern und Obersten in der Kirche, von den Kindern und Erben der Verheißungen Gottes, welche sie in ihres Vaters Hause von Jugend auf mit großem Fleiß und Ernst von ihrem Großvater Isaak und Jakob ihrem Vater gehört und gelernt hatten; noch geschieht es.

306. Wir sind auch durch die Erbsünde verderbet, daß wir auch, leider! also gar geneigt sind zu sündigen. Da sollen wir aber eins thun, und uns befleißigen, daß wir in solcher großen Schwachheit unserer verderbten Natur gleichwohl noch immer ein brüderlich und kindlich Herz und natürliche Zuneigung oder Liebe behalten mögen, damit wir uns ja nicht selbst Strafen und Verderben aufladen. Denn diese Brüder, da sie sich befleißigen, ihren unschuldigen Bruder zu unterdrücken, dieweil sie ihn so gar heftig gehaßt, sind eben damit ihre eigenen heftigen Feinde und Teufel geworden, so gar sehr, daß sie sich auch selbst mehr Schaden gethan, denn Joseph immer hätte thun können, wo er wäre ihr Herr geworden.

\subsection{II}

\begin{quote}
\textbf{V. 31—33. Da nahmen sie Josephs Rock, und schlachteten einen Ziegenbock, und tunken den Rock ins Blut, und schickten den bunten Rock hin, und ließen ihn ihrem Vater bringen, und sagten: Diesen haben wir funden; siehe, ob es deines Sohns Rock sei, oder nicht. Er kannte ihn aber, und sprach: Es ist meines Sohns Rock; ein böses Thier hat ihn gefressen, ein reißend Thier hat Joseph zerrissen.}
\end{quote}

307. Sie haben noch einen andern Betrug und List versucht, damit sie ihre böse That desto leichter bedecken, und ihren Vater um so viel schwerer und härter betrüben und martern mögen. Denn sie nehmen Josephs Rock, und nachdem sie denselben in das Blut eines Ziegenbocks getunkt, schicken sie denselben dem Vater hin, ihn damit also zu überreden, daß Joseph von einem wilden Thiere sei zerrissen worden.

308. Wie weit ist aber wohl Gott von ihren Augen? Wie gar sicher und ohne alle Furcht richten sie doch dies alles aus? Sie achten es nichts, daß sie doch gedacht hätten, wie der gute alte Vater möge sein gekümmert worden, da ihm unversehens solche traurige Botschaft gekommen. Sie hätten sich doch allein davor mögen hüten, daß sie ihn nicht mit so ungewöhnlichem und grausamem Anschauen des blutigen Kleides erschreckten, da sonst dem armen veralteten Manne nichts Betrübteres und Elenderes hat vorgebracht werden können. Es wäre ja noch etwas gelinder gewesen, wenn sie ihm hätten angezeigt, daß der Knabe unterwegs etwa von den Mördern angegriffen und weggeführt worden wäre; oder aber, daß sie ihn nicht gesehen hätten, auch nicht wüßten, wie er doch möchte umgekommen sein: er möchte vielleicht den Feinden, die nahe umher gewohnt, oder den Ismaeliten oder Midianitern, denen sie ihn doch verkauft hatten, in die Hände gekommen sein. Aber der Dinge keines ist ihnen eingefallen, verschonen auch ihres alten Vaters gar nicht, sondern haben mit solchem betrübten und erschrecklichen Gegenstande die Schmerzen seines Herzens immer gemehret und größer gemacht.

309. Derhalben ist dies eine große und gar teuflische Bosheit, daran der Satan Lust hat, daß er also alle Grausamkeit und Muth über den frommen alten Mann ausschütten möchte. Sie aber werden darüber in greulichen erschrecklichen Sünden gar ersäuft, und bleiben darin stecken wider ihr eigen Gewissen, ohne alle Buße, ohne alle Empfindung der Liebe, und haben darnach die ganzen zwei und zwanzig Jahre über gelebt ohne wahre Anrufung, ohne Gottesfurcht und alle geistlichen Uebungen; derselben haben sie sich nicht gottseelig befleißigen können, es sei denn, daß sie zuvor ihre Sünde erkannt und sich mit ihrem Vater und Brüder versöhnt hätten.

310. Was thut aber nun dieweil unser Herr Gott dazu? Wo sind nun die großen herrlichen Verheißungen, 1 Mos. 15, 1.: „Ich bin dein Schild und dein sehr großer Lohn“; Cap. 35, 11.: „Sei fruchtbar und mehre dich“; Cap. 31, 3.: „Ich will mit dir sein“? Ist es nicht wahr: Alles scheint hier diesen großen Verheißungen ganz zuwider? Denn es scheint, als kenne jetzt unser Herr Gott Jakob und Joseph nicht, achte ihrer auch gar nicht; sondern als habe er sie vielmehr dem Teufel vorgeworfen, sintemal er zuläßt, daß der sehr heilige Vater Jakob von der Wuth seiner grimmigen Söhne so erbärmlich gemartert und geplagt wird. Sollen wir nun sagen, daß Gott da an seine Verheißung gedacht habe? Man kann zwar nichts weniger sagen oder urtheilen, denn also. Aber David hat es gesehen und allererst nach der That recht davon geredet, welches sonst im rechten Kampf der schweren Anfechtung nicht hat können gesehen werden. Denn also sagt er im 105. Psalm V. 17.: „Er sandte einen Mann vor ihnen hin“, der Jakob helfen sollte.

311. Was ist aber das für ein Senden? Was ist das für eine Art zu reden: Einen Helfer oder Heiland nach Egypten senden, der Jakob und seinem ganzen Hause helfen soll? Wie wird er doch gesandt? Er wird in die Grube geworfen, wird verkauft, sein Vater wird auch getödtet. Heißt das einen Heiland senden? Ja, es heißt also, aber nach der Art, wie unser Herr Gott zu reden pflegt. Denn Joseph wird ja zu einem großen König gesetzt. Aber Gott sieht es allein, Jakob und Joseph sehen es nicht; die sind in der größten Angst und tragen nur Leid. Derhalben ist dies eine sonderliche und himmlische Sprache, einen Heiland oder Helfer senden und ihn zum Könige setzen eben damit, daß er in die Grube und in die Höhle hinein geworfen wird. An diese Sprache müssen wir unsere Herzen gewöhnen, daß wir verstehen lernen, was David an einem andern Ort, Ps. 37, 12. 13., sagt: „Der Gottlose dräuet dem Gerechten, und beißet seine Zähne zusammen über ihn. Aber der Herr lachet sein; denn er siehet, daß sein Tag kommt“, darin er verderben und zu Boden gehen soll etc. Das können wir nicht thun, wir können nicht so weit sehen und am Glauben halten, denn allein mit großer Schwachheit; aber Gott hält sehr fest über seinen Verheißungen, dergestalt, daß er derselben nicht allein nicht vergisset, sondern auch auf ihre Erfüllung bedacht ist und der Widersacher lacht. Und über das spricht er auch das Urtheil über sie, wie der 2. Psalm V. 4. sagt: „Der im Himmel wohnet, lachet ihrer, und der Herr spottet ihrer.“

\quellseite{300}

312. Ja, sprichst du, sehe ich es doch nicht? Das ist auch recht: darum werden uns diese Exempel vorgehalten, daß wir in unsern Anfechtungen auch denken sollen, daß unsere Widersacher jetzt auch belacht und schon eine Strafe über sie beschlossen wird; daß wir aber, wiewohl wir betrübt und geplagt sind, doch von Gott geliebet werden, und daß er für uns gewißlich Sorge trage und große Acht auf uns habe, aber gar verborgener Weise, wie Jesaia Cap. 45, 15. sagt: „Fürwahr, du bist ein verborgener Gott, du Gott Israel, der Heiland.“

313. Derhalben sollen wir uns daran genügen lassen, daß wir das Wort und die heiligen Sacramente haben, darin sich Gott uns offenbart, die Frucht aber und das Ende solcher Gnadenzeichen wird zu seiner Zeit auch folgen. Inzwischen sollen wir uns mit diesen Gedanken aufrechthalten und trösten: Ich habe das Gnadenzeichen, und dabei das Wort, daran will ich fest halten, wie heftig auch die Welt und der Satan wider mich toben und wüthen, und mich mit allerlei Jammer und Unglück überschütten; allein, daß wir zusehen und uns bescheiden, daß wir solches alles, was es nur ist, mit gutem Gewissen leiden. Denn daran ist kein Zweifel, wir sind gewißlich bei Gott in sehr großen Gnaden, und er hält uns sehr theuer und werth, die Widersacher aber werden von Gott verlacht und verspottet, und dazu behalten, daß sie endlich noch die schwerste Strafe und Plage leiden sollen. Denn das Lachen unseres Herrn Gottes gibt das höllische Feuer, wie im 2. Psalm V. 5. folgt: „Er wird einst mit ihnen reden in seinem Zorn, und mit seinem Grimm wird er sie schrecken.“ Ja, hüte dich vor solchem Lachen.

314. So sind nun Jakob und Joseph vor Gott Kinder der Gnade. Diese Verkäufer aber sind Kinder des Spottes, des Zorns und Verdammniß. Und wollte Gott, daß wir dies also könnten lernen und nur einigermaßen behalten. Denn das Fleisch liegt uns immer im Wege, und ist doch gewiß, daß solch Leben, darin man immer so viel leiden muß, das allerbeste und köstlichste Leben ist, so auch, daß ihm keine Vergebung der Sünden vornehmen. Denn es ist ohne Sünde. Dies sage ich aber von Frommen und Gottseligen, das ist, von denen, die da leiden im rechten Glauben an Christum Jesum. Denn die sündigen ja damit nicht, sondern tragen solches und leiden um anderer Leute Sünde willen, und ist wohl von Socrates geredet: Es ist besser Unrecht leiden, denn Unrecht thun. Denn der da leidet, sündigt ja damit nicht; der aber Andern Unrecht thut und sie betrübt, der sündigt. Wer sich unschuldig weiß, und merkt, daß er unbilliger Weise betrübt und geplagt wird, der kann sich an seine Unschuld halten, und siehet zugleich auch mit auf die Verheißung der Erlösung. Darum fürchtet er sich nicht, verzagt auch in seinem Herzen nicht; denn er weiß, daß er ja niemand beleidigt, sondern nur anderer Leute Unrecht und Sünde trägt. Wenn die rupes und petra, das ist, der Fels, dasteht, nämlich, ein gut Gewissen, so kann nichts schaden; wenn gleich Caiphas und Ischarioth kommen und zürnen, so haben wir doch gewonnen. Ein gut Gewissen ist wie der härteste Steinfels, darauf sich Fromme und Gottselige in ihrer Trübsal verlassen, und verachten mit großem und hohem Muth das Drohen aller Widersacher; wie geschrieben steht 1 Petr. 3, 13. 14.: „Und wer ist, der euch schaden könnte, so ihr dem Guten nachkommet? Und ob ihr auch leidet um Gerechtigkeit willen, so seid ihr doch selig. Fürchtet euch aber vor ihrem Trotzen nicht, und erschrecket nicht.“

315. Wiederum aber, wenn die betrübt und gestraft werden, so Andere unbillig beleidigt haben, alsdann fürchten sie sich, und erschrecken greulich, und wie der Poet sagt, werden sie bleich und entsetzen sich, wenn sie auch nur einen Blitz sehen. Ach, rufen sie, wehe mir, diese Strafe habe ich mit meiner Bosheit, Muthwillen und Ungehorsam verdient. Was soll ich nun thun? Wohin soll ich mich kehren oder wenden? Das ist denn wahrlich eine erbärmliche Marter des Gewissens, welches sich also selbst anklagt und verdammt.

316. Derhalben geht es Jakob und Joseph sehr wohl, wiewohl sie zwar nicht lachen, oder Glück und Wohlfahrt erkennen können, so lange die schwere Trübsal mit ihnen gewähret. Denn der ist wahrlich recht selig, der in der Anfechtung also gesinnet ist, daß er es bei sich selbst dafür halten kann: Wenn schon das Aergste nach unten gekehrt, und die Welt strafen und über einen Haufen fallen sollte, ja, wie heftig die Welt mit ihrem Fürsten und allen höllischen Pforten tobt und wüthet, was geht es mich an? ich weiß mich hier unschuldig; denn ich trage anderer Leute Kreuz, von denen ich unbilliger Weise geplagt werde. Ich thue niemand leid, sündige nicht. Ueber das habe ich meine Absolution und die heiligen Sacramente. Laß hergehen, ich will mich nicht deswegen lassen noch erschrecken. Wehe aber denen, die also sündigen und mich plagen! Denn es gebühret Christen, daß sie nicht zürnen oder ungeduldig werden, wo ihnen Leid widerfährt. Wiewohl doch das Fleisch nach seiner Art und Weise immer pflegt zu murren; aber der Geist zürnt nicht, denkt viel lieber also: Wehe dir! du hast mir es nicht gethan, du hast dir es gethan; hast nicht mich, sondern dich selbst am allermeisten beleidigt. Also hätten Jakob und Joseph auch sagen mögen: Ihr verkauft, verwerft und tödtet mich nicht, sondern euch selbst. Das Fleisch versteht es wohl viel anders: aber der Geist ist eigentlich also gesinnet.

317. Nun möchte aber jemand fragen: Was soll man denn nun thun? Soll man derhalben dem gottlosen Wesen Zaum und Zügel lassen, und dem Grimm und Bosheit der Menschen gar nicht steuern oder wehren, weder mit Gesetzen noch Strafen; sondern noch außerdem einem bösen Buben dafür, daß er uns Schaden gethan, danken und zu ihm also sagen: Du hast mir gar kein Leid gethan, thue nur immerhin, was dich gelüstet, ich will gern Alles leiden? Auf solche Weise werden die bösen Buben gereizt, ihrer Sünden immer mehr zu machen und dieselben zu häufen, sintemal wir ja lehren, daß man den Schaden, so uns von Andern zugefügt wird, mit Freuden und gutem Gewissen leiden soll.

318. Antwort: Uns ist befohlen das Amt des Gesetzes und Evangeliums nach der Lehre Pauli 2 Tim. 4, 2.: „Strafe, dräue, es sei zu rechter Zeit, oder zur Unzeit.“ Nur getrost, schelten und strafen, auch über die Sünden zürnen und über dieselben unwillig sein, das gehört zum Amt des Gesetzes, und ist auch der Väter, Schulmeister und Obrigkeit Amt. Ja, es gebühret auch wohl gemeinen Privatpersonen um der brüderlichen Strafe und Vermahnung willen, so uns allen insonderheit von Gott geboten, nämlich, daß wir nach allem unserm Vermögen, ein Jeder an seinem Orte, dem Uebel steuern und wehren solle, und verhüten, daß die bösen Buben nicht also frei ungestraft hingehen, toben und wüthen wider uns und andere Leute. Wenn das geschehen ist und will doch noch nicht helfen, so soll ich mich zufrieden geben, und sagen: Du wirst nicht mich, der ich dich gelehrt, sondern dich selbst strafen. Und wenn wir mit diesem allen nichts ausrichten können, haben wir das noch vor uns, daß wir in unserem Herzen zufrieden seien und nur auf die göttliche Strafe warten; denn die bösen muthwilligen Buben werden es erfahren zu ihrem großen Schaden, daß sie nicht Andere, sondern sich selbst am allerhärtesten beleidigt. Der Teufel wird sie wohl genung martern.

319. Wir aber haben diesen Trost: je mehr sie uns bekümmern, plagen und martern, um so viel größer machen sie unsere Ehre und Krone im Himmel. Inzwischen aber müssen sie gemahnt und gestraft werden; nicht daß wir ihnen damit Schaden thun, sondern daß wir solche boshafte Leute auf den rechten Weg bringen mögen, daß sie sich bessern, und nicht in Gottes schwerem Zorn und Ungnade fallen, der den unbußfertigen Menschen ein verzehrend Feuer ist. Auf diese Weise widerstehen wir dem Bösen, nämlich, durch das Amt des Worts und weltlichen Schwerters, und leiden doch das Unglück, dem wir nicht wehren können; welches uns zu großem Nutz und Frommen, ihnen aber zu Schaden und Verderben gereichen wird.

320. Dies ist der Christen Theologie und Weisheit, und wiewohl wir dieselbige noch nicht gar erreicht haben, sollen wir doch täglich
\quellseite{301}
darin geübt werden, und uns gewöhnen, daß wir im geistlichen Kampf und Trübsal, so wir leiden, mit einem beständigen und stillen Herzen sagen können: Du kannst mir keinen Schaden thun, ich bin ein Christ, du schadest mir nichts, du förderst mich; siehe dich vor. Was hat es Joseph geschadet, da er verkauft und in das Elend getrieben worden ist? ja, wozu ist es ihm nicht nütze und gut gewesen? oder wie hätten ihm seine Brüder zu größerer Ehre und Herrlichkeit bringen können? denn eben damit sie versucht haben, ihn zu hindern und zu unterdrücken, haben sie ihn sein zu der Hoheit erhoben, daß er zum großen Herrn geworden; wie ihm denn kurz zuvor geträumt hat.

321. Auf diese Meinung wird der Spruch Gregors auch gerühmt: Die Gottlosen thun uns Gutes, auch wenn sie Böses thun. Und Augustinus sagt von den unmündigen Kindlein, so von Herodes getödtet worden: Es hätte der Feind Herodes mit aller Macht und höchsten Vermögen seines Königreichs diesen Kindern keinen größeren Nutzen zufügen können, denn eben diesen, daß er sie getödtet.

322. Also erniedrigt nun Gott die Seinen, daß er sie erhöht; er tödtet sie, daß er sie wieder lebendig mache; macht sie zu Schanden, daß er sie zu Ehren setze; wirft sie herunter, daß er sie erhebe. Aber es ist eine Kunst über alle Künste und Weisheit über alle Weisheit, die nicht ohne große Arbeit, und von so gar Wenigen pfleget gelernt und verstanden zu werden; sie ist aber doch wahrhaftig und gewiß, wie dies Exempel zeuget. Denn Gott hat wahrhaftig, wie im 105. Psalm V. 21. steht, den Joseph zum König in Egypten, und zum Herrn und Helfer ihrer Völker gesetzt. Wie denn? Daß er verkauft, verworfen und getödtet worden; das sind Gottes Werke, die nicht können verstanden werden, es sei denn, wenn sie erfüllt und vollbracht sind. Indem sie aber geschehen, können sie nicht begriffen oder verstanden werden, denn allein mit dem Glauben. Denn darum muß man sich einfach halten, daß wir sagen: „Ich glaube an Gott den Vater, allmächtigen Schöpfer Himmels und der Erden.“

323. Auf dieselbige Weise, wenn ich aus diesem Leben werde abscheiden, halte ich mich an diesen Trost, daß ich glaube an den Sohn Gottes. Ja, ich werde in die Erde begraben, von den Würmern gefressen, und muß gar im Stanke verfaulen und verzehrt werden; wie Hiob Cap. 17, V. 14. sagt: „Die Verwesung heiße ich meinen Vater, und die Würmer meine Mutter und meine Schwester“; da sehe ich Gottes Rath nicht, daß wenn ich schon sterben und verfaulen muß, ich doch noch einmal lebendig werden soll. Aber Gott hat es verheißen und gesagt: Du wirst wieder leben; denn „ich lebe, und ihr sollt auch leben“, Joh. 14, 19. „Ich bin der Herr, dein Gott.“ Wie werden wir aber leben? Antwort: Im ewigen Leben, und in einem solchen Leibe, der schöner und klarer ist denn die Sonne. Dies sehe oder fühle ich jetzo nicht, ich glaube es aber und ertrage diesen sehr geringen Verzug. Denn das Leben ist schon bereit, und wird inzwischen auch die Krone des Reichs und Herrlichkeit bereitet, welche mir der Herr an jenem Tage, der gerechte Richter, geben wird, nicht mir aber allein, sondern auch allen, die seine Erscheinung lieb haben“, wie Paulus sagt 2 Tim. 4, 8.

324. Dies geschieht aber alles verborgener Weise; darum soll man es mit Geduld leiden und ertragen, daß sich Gott also wunderbärlich verbirgt. Jakob und Joseph haben die Krone des Reichs nicht gesehen, es zeigen sich auch seine Brüder nicht vor ihm, wie hernach in Egypten geschehen wird; aber mit diesem Verkaufen wird das Zukünftige vorbereitet. Und ist bei unserm Herrn Gott eben, als sei es geschehen; es kommt gewißlich.

325. So ist nun dies uns zur Lehre geschrieben, daß wir den Glauben recht verstehen und üben lernen, welcher den Tod und alles Unglück leiden kann, und doch in dem allem auf das Leben und die Erlösung warten, welcher Glaube, sage ich, die Gewalt und Unrecht, so uns von Andern begegnet, ertragen kann; wiewohl er sie darüber Amts halben ernstlich schilt und straft, haßt aber doch niemand, sondern verzeiht gerne, segnet, hofft, wünscht und thut Gutes ohne alle Rachgier. Denn das sehen wir klärlich an dem Exempel Josephs, welcher ja seines Unrechts, so ihm widerfahren, gedenkt, begehrt keiner Rache; sondern thut auch noch alles Gute den allerbösesten Buben, so ihn verkauft und getödtet hatten.

326. Auf diese Weise sollen ihre Herzen gewöhnen zur Geduld, Glauben und Liebe diejenigen, so Gewalt haben, Andere zu züchtigen und zu strafen, es sei mit dem Wort, oder mit dem Schwert, oder aber mit der Ruthe. „Die Geduld“, sagt St. Jacobus Cap. 1, 4., „soll fest bleiben bis ans Ende.“ Denn der geduldig ist, sündigt nicht. „Wer gestorben ist, der ist gerechtfertiget von der Sünde“, Röm. 6, 7. Ein geduldiger Mensch ist im Glauben Christi wahrlich recht heilig, da ist keine Sünde mehr übrig an ihm; denn alles, was er leidet, ist lauter Gerechtigkeit.

327. Diese Dinge sind also in den Historien der lieben Heiligen geschrieben und geschehen; es ist aber alles sehr schwach und unvollkommen zugegangen; wie das an Jakob und Joseph zu sehen. Die haben zwar ihren Söhnen und Brüdern nichts Böses gewünscht oder gethan: es ist aber doch wohl abzunehmen, daß in Beiden ein großer Streit muß gewesen sein wider die Schwachheit des Fleisches. Denn wie oft, meinst du, wird Joseph mit weinenden Augen gen Hebron, da der Vater wohnte, gesehen haben, da er vorüber geführt ist? Wie kläglich wird er gethan haben, da er so plötzlich und unversehens dem Vater aus den Augen und aus seinem Hause weggeführt und in anderer Leute Gewalt genommen ist! Das Fleisch wird sich des Murrens nicht haben enthalten können; aber der Geist, so darüber gestritten und geseufzt, hat doch gewonnen und den Sieg behalten, hat die unbillige Gewalt seiner Brüder ohne alle Nachgier ertragen; da die sehr bösen Buben gar unsinnig gewesen, und immer ihrem Haß, so sie wider ihn gefaßt, nachgehängt und gefolgt sind.

328. Und Mose hat auch dies noch hinzugesetzt, wie sie den elenden alten Mann so jämmerlich und unbillig haben gemartert damit, daß sie fälschlich erdichtet, wie es mit Joseph zugegangen sei. Denn das zeugt die erbärmliche und tödtliche Klage, so Jakob gethan, da er sagt: „Es ist meines Sohns Rock; ein böses Thier hat ihn gefressen“ etc. Dies hat er genommen aus den Worten seiner Söhne, die sie geredet, da sie ihm das blutige Kleid gezeiget. Da setzt er das dazu, solches müsse ja ein Argument und Anzeichen sein, daß Joseph von einem bösen Thiere zerrissen. Nun wär es doch noch viel wahrscheinlicher, daß ihn ein Mensch müßte erwürgt haben; und wo der Vater etwas fleißiger der Sache hätte nachgeforscht durch alle und jede Umstände, würde er den Betrug bald erfahren haben. Denn er hätte mögen sagen: Wie daß der Rock ist ganz geblieben? Und, wo habt ihr ihn denn bekommen? Man sieht ja keine Wahrzeichen von Zähnen oder Klauen irgend eines wilden Thieres? Denn wo ihn ein böses Thier zerrissen hätte, so hätte es ihn ja den Rock nicht zuvor angezogen, würde auch denselben nicht haben lassen ganz bleiben, sondern würde ihn zugleich mit dem Leibe weggerissen haben. Aber daran hat der gute Vater nicht gedacht. Denn er ist gar ersoffen und verschlungen gewesen in großem herzlichen Schmerz und Wehklagen. Darum hat er sich in so großer Traurigkeit und der schwersten Anfechtung leichtlich überreden lassen.

\subsection{III}

\begin{quote}
\textbf{V. 34—36. Und Jakob zerriß seine Kleider, und legte einen Sack um seine Lenden, und trug Leid um seinen Sohn lange Zeit. Und alle seine Söhne und Töchter traten auf, daß sie ihn trösteten; aber er wollte sich nicht trösten lassen, und sprach: Ich werde mit Leide hinunter fahren in die Grube, zu meinem Sohne. Und sein Vater beweinete ihn. Aber die Midianiter verkauften ihn in Egypten dem Potiphar, des Pharao Kämmerer und Hofmeister.}
\end{quote}

329. Nun haben sie es wohl ausgerichtet, die Junker. Da sie ihren Bruder verkauft und getödtet haben, sind sie des auch eine Ursache geworden, daß ihr armer alter Vater für und für Leid getragen und auch zu Tode ist gemartert worden; ja, sie haben ihn zwar noch greulicher geplagt als ihren Bruder. Denn sie machen, daß er ganze zwei und zwanzig Jahre nach einander Leid getragen, und haben ihn die lange Zeit über immer sehen im Sacktuch und großer erbärmlicher Traurigkeit einhergehen, daß er immer seufzt, schluchzt und jämmerlich weint. Denn ich halte dafür, er habe nach der Zeit niemals weder fröhliche Kleidung gebraucht, noch ein fröhlich Angesicht gehabt. Hilf Gott! wie müssen das so gar unbekannte und greuliche Leute gewesen sein, daß sie haben ansehen können, daß ihr alter Vater so gar unbilliger Weise und so gar lange Zeit hat müssen trauern und Leid tragen! Ei! daß die Leute sollen so voller Teufel sein!

\quellseite{302}

330. Ich hätte nicht geglaubt, daß jemals ein Mensch solche große Sünde hätte thun können, daß er wissentlich und vorsätzlich vor seinen Augen ansehen könnte, seinen alten schwachen Vater sterben und in die Grube steigen, dazu er mit seinen Sünden Ursache gegeben, und in dem allen mit ihm kein Mitleiden haben, ja, darüber an sich selbst gar keine menschliche Freundlichkeit, oder natürliche Neigung oder Liebe, so doch allen Thieren von Gott eingepflanzt, gegen Vater und Großvater fühlen oder empfinden sollte. Heißt das Vater und Mutter ehren? Und dazu der alte Vater Israel muß das auch sehen. Denn sein Schmerz und Wehklagen ist auch hinzugekommen, welcher, nachdem Joseph verkauft worden ist, noch zwölf Jahre gelebt, und zuletzt ohne Zweifel vor herzlicher Wehmuth, darum daß er seines Sohnes Sohn verloren, wird gestorben sein. Dieser treffliche Mann, der billig der Welt Preis möchte genannt werden, soll sich noch zu Tode weinen und grämen? Es wäre kein Wunder, wenn sie unser Herr Gott hätte Sodom und Gomorra lassen werden, und wenn Feuer vom Himmel gefallen wäre und sie also nach dem Exempel der Sodomiter vertilgt hätte. Und sie wären auch ohne Zweifel eben auf solche Weise umgekommen und verderben, wo nicht etliche wenige fromme gottselige Leutlein, so dies Geschlecht noch aufgehalten, nämlich, Isaac, der Großvater, und Jacob, der Vater, und Joseph, sein Sohn, gewesen wären. Dies sind zu der Zeit die rechten drei Atlanten gewesen; die drei Männer tragen es, sonst würde Gott mit Schwefel und Pech darein geschlagen haben.

331. Und ist wahrlich nichts wunderbarer, denn daß solche Gesellen so viele Jahre das Wort göttlicher Verheißung und Segens gehört und sich doch daraus dieser Brüder kein einiger gebessert hat, da sie doch oft durch ernstliche Lehre und Gottseligkeit des Vaters und Großvaters sind ermahnt worden, haben sich auch durch solche klägliche Traurigkeit der frommen alten heiligen Männer nicht bewegen lassen, daß sie doch mit ihnen auch Mitleiden gehabt und sich ihrer erbarmt hätten. Ihr seid harte liebe Gesellen; also soll man dem Teufel dienen. Nun haben sie aber damit weder dem Vater noch dem Großvater, noch Joseph, sondern sich selbst am allermeisten Schaden gethan. Simeon und Levi sind dazu getrieben worden durch Hoffnung auf die Erstgeburt und Begierde darnach; Simeon ist Papst gewesen, Levi Cardinal; Juda ist etwas aufrichtiger, hat sich aber doch vor der zwei Brüder Macht und Gewalt gefürchtet; Ruben hat bei ihnen gar kein Ansehen gehabt, sondern ist gar verachtet gewesen wegen des begangenen Schandfleckes, darum hat er müssen stille schweigen um des Scharioth und Caiphas willen. Es wäre kein Wunder, wenn sie unser HErr GOtt hätte Sodom und Gomorra lassen werden, und wenn Feuer vom Himmel gefallen wäre und sie also nach dem Exempel der Sodomiter vertilgt hätte. Und sie wären auch ohne Zweifel eben auf solche Weise umgekommen und verdorben, wo nicht etliche wenige fromme gottselige Leutlein, so dies Geschlecht noch aufgehalten, nämlich, Isaac, der Großvater, und Jakob, der Vater, und Joseph, sein Sohn, gewesen wären. Dies sind zu der Zeit die rechten drei Atlanten gewesen; die drei Männer tragen es, sonst würde GOtt mit Schwefel und Pech darein geschlagen haben.

332. Also muß man ihre Sünde, so viel immer möglich, nur groß und schwer genug machen, uns zum Trost. Denn dieweil sie Vergebung solcher Sünde endlich erlangt haben, sollen wir auch um unserer Sünde willen nicht verzweifeln. Aber doch also, daß du gleichwohl Gott nicht versuchest. Denn da Judas verzweifelt, hat er sich selbst den Strick um den Hals gelegt und ist also in Verzweiflung erschrecklich dahin gestorben. Derhalben siehe zu, daß du dir selbst nicht auch das Unglück und Verderben aufsetzest. Wo aber ein solcher Fall kommt, daß ein Mensch etwa von einem Fehl übereilet wird, wie Paulus Gal. 6, 1. sagt, da thut es solcher Exempel nothwendig, daß du bedenkt, die heiligen Patriarchen sind auch arme Sünder gewesen und sind erschrecklich gefallen. Ich weiß mich zwar solcher greulichen Sünden nicht schuldig, außer daß ich fünfzehn Jahre dem Antichrist gedient, da ich täglich die abscheulichen Opfermessen gehalten. Aber wo mich mein Gewissen schreckt, halte ich mich an solche Exempel und Verheißungen des Evangeliums und tröste mich damit. Es hätte zwar keine Blutschande so schwer sein können, wie solche unbillige böse That dieser Brüder gewesen, damit sie ihren Vater und Bruder, dazu den Großvater durch Traurigkeit und herzliche Schmerzen also gemartert und getödtet haben. Und endlich werden sie ohne Zweifel einen großen greulichen Schrecken ihres Gewissens gefühlt haben, da sie ihren Bruder und den König in Egypten angesehen und vor ihnen erschrocken sind. Und über das ist alles Unglück, so sie sammt ihren Nachkommen unter Pharao erlitten haben, die rechte Strafe gewesen, so ihnen Gott um dieser Sünden willen auferlegt. In Egypten werden sie es wohl gefühlt haben.

333. Bis daher haben wir nun gehört, wie Joseph nach Egypten geführt, und von den Ismaeliten und Midianitern dem Potiphar, dem Küchenmeister oder Hofmeister des Königes, verkauft worden. Das läßt Gott also geschehen, und schweigt stille dazu, schläft, ist taub und ganz harthörig, hat gar kein Mitleiden, und thut, als kenne er auch solchen lieben Sohn nicht. Alle Engel schweigen auch stille, und wird ein solcher großer Schatz nach Egypten hingerrissen und um ein geringes Geld, nämlich, zwanzig Silberlinge, verkauft, der doch so ein großer Patriarch, Prophet und Regent werden sollte. Was macht doch unser Herr Gott mit seinen Auserwählten? Was ist doch das für eine wunderbare Regierung, damit die auserwählten Gotteskinder geführt und regiert werden? Was ist es doch, daß er sie also verlässt und so jämmerlich plagt?

334. Der arme elende Joseph wird von seinen Käufern dahingerrissen, und muß an Hebron vorüberziehen, da ihm doch wird eingefallen sein, daß er gedacht hat: Siehe, da wohnt mein lieber Vater, er weiß nicht, wie es jetzt mit mir geht, und ich kann ihn doch nicht anreden, kann ihn auch nicht anschauen, oder von ihm den letzten Abschied nehmen. Das ist wahrlich ein großer erbärmlicher Jammer gewesen. Ich will des alten Vaters schweigen, welcher, da er den Handel erfahren, mit weinenden Augen schreit und sagt: „Ich werde mit Leide hinunter fahren in die Grube, zu meinem Sohne“ und mich mit ihm verscharren lassen, und dieweil ich diesen Sohn verloren habe, wird mir nun dies Leben nicht mehr lieb oder angenehm sein.

335. Das mögen aber auch Kaufleute heißen, die einen solchen theuren Schatz an Hebron vorüberführen und nach Egypten bringen, welcher darnach des ganzen Königreichs Heil wird, so den Leuten an Leib und Seele geholfen. Denn er wird daselbst eine Kirche und rechte Lehre von Gott anrichten, und wird große Dinge ausrichten, so demselben ganzen Königreiche nütze und gut sein werden. Darum hat er erstlich also müssen aufgekaufet und getödtet werden, ehe der Tag seiner Auferstehung und Herrlichkeit gekommen ist, damit er nicht stolz werde, sondern bedenke, was er zuvor gewesen, und woher er zu solcher großen Herrlichkeit sein erhoben worden.

336. Es wird dies also beschrieben uns zum Troste. Denn in diesem erschrecklichen Kreuze des Vaters und Sohnes ist Gott dieweil taub und stumm gewesen, kennt nicht, weiß auch nicht von den Dingen, so da geschehen sind. Aber der Glaube ist noch da, und Gott spricht ihm innerlich im Herzen zu und sagt: Lieber Joseph, harre doch, glaube nur und verzage nicht, halte fest an der Verheißung, so du von deinem Vater gehört. Also, sage ich, redet Gott mit ihm durch das Wort seines Vaters: Gott hat dem Urgroßvater und Großvater und dazu deinem Vater einen Samen verheißen, nun denke darauf, daß du in derselbigen Verheißung mit festem Glauben beharrest und beständig bleibest. Aber er redet solches mit ihm durch wunderbares Stillschweigen, da er nichts sieht oder hört. Denn Gott ist, als wäre er gar blind und todt, und Joseph lebt und hält sich allein an die gemeine Verheißung: Gott hat Abraham einen Samen verheißen 2c. Ich glaube an Gott, an den meine Väter auch geglaubt haben. Darnach wird er dann herrlich und mit der That mit ihm reden, wenn er ihn zum Könige und Helfer über ganz Egypten setzen wird. Jetztund aber ist Joseph begraben und gestorben, hat seinen Charfreitag und Sabbath; sein Vater stirbt auch dahin: aber sie werden beide durch göttliche Kraft, welche auch die Todten kann lebendig machen, wieder aufstehen. Das Herz eines Gläubigen muß doch leben und wieder auferstehen von den Todten, ob es gleich von der unermeßlichen Last der ganzen Welt unterdrückt da läge.

337. Und werden uns diese Exempel darum vorgehalten, daß wir uns gewöhnen an Geduld in Trübsal, daß wir ja nicht sollen ungeduldig sein und wider GOtt murren, wie viel und groß auch die Trübsal, Angst und Noth sei, damit wir beschweret sind. Wehe thut es; wie es dem jungen Herzen auch ohne Zweifel wird wehe gethan haben. Es kann zwar ein menschlich Herz solches nicht ertragen oder gar überwinden ohne große Pein und Schmerzen. Wie denn Joseph ohne Zweifel auch große Schmerzen gehabt und heftig betrübt gewesen sein wird, da er gefühlt, daß er so unbilliger Weise von seinem Vater gerissen und Fremden vorgeworfen worden ist, und daß er also in ewige Dienstbarkeit dahingegeben, da er nimmermehr etwas Eigenes zu bekommen, oder wiederum auf Freiheit und Erledigung hat hoffen können, sondern mußten servus servorum, Knechtes Knecht sein. Denn die Dienstbarkeit ist an sich selbst schwer und erbärmlich genug, wenn schon andere Beschwerungen nicht dazu kommen. Sie ist aber diesem frommen Jünglinge so viel verdrießlicher gewesen, daß er noch in der Zeit seiner blühenden Jugend seiner Eltern und aller Bequemlichkeiten und Vortheile dieses zeitlichen Lebens hat müssen beraubet sein. Läßt denn solches unser Herr GOtt über seine Kinder gehen; so laßt uns auch mit geduldigem Herzen ertragen, wenn uns etwa traurige und widermärtige Dinge begegnen. Denn das sind nicht Zeichen des Zorns, oder daß uns GOtt verlassen hätte, sondern vielmehr Anzeichen der Gnade und Prüfungen unseres Glaubens.

\quellseite{303}

338. Darum ist ein solch Leben das allerheiligste, darin auch die höchsten Patriarchen gelebt, mit welchen unsere Mönche und papistischen Bischöfe gar nicht zu vergleichen sind. Es ist nichts mit dem Fasten und mit den Wunderwerken, so in den Legenden unserer Heiligen, als, des Franciscus, Ambrosius, Augustinus \&c., erzählt werden. Das heißt gelebt: jenes ist eitel Kinderspiel. Dies sind, sage ich, Exempel, die da lehren und anzeigen, was ein recht christlich Leben sei, welches die rechten wahrhaftigen Uebungen der Gottseligkeit und christlicher Geduld seien. Wohlan, Joseph ist hin; da lassen wir ihn in der Grube ruhen. Er ist nun begraben, wir wollen ihn ruhen lassen in der Grube, im scheol, wie sein Vater sagt, in seiner Schule.

\chapter{Das achtunddreißigste Capitel}

\section[Erster Theil]{Erster Theil. Wie Juda gen Odollam zieht, ein Weib nimmt und drei Söhne zeugt}

\begin{quote}
\textbf{V. 1—5. Es begab sich um dieselbige Zeit, daß Juda hinab zog von seinen Brüdern, und that sich zu einem Mann von Odollam, der hieß Hira. Und Juda sahe daselbst eines Cananiters Mannes Tochter, der hieß Suah; und nahm sie. Und da er sie beschlief, ward sie schwanger, und gebar einen Sohn, den hieß er Ger. Und sie ward abermal schwanger, und gebar einen Sohn, den hieß sie Onan. Sie gebar abermal einen Sohn, den hieß sie Sela; und er war zu Chesib, da sie ihn gebar.}
\end{quote}

1. Ich habe etlichemal Erinnerung gethan, daß man fleißig soll Acht geben auf die Ordnung der Historie in der Erzählung von dem Tode Isaaks, und wie Joseph verkauft worden ist. Denn Mose gebraucht die Figur, so hysteron proteron genannt wird, nach der Weise der andern Geschichtschreiber. Am Ende des 35. Capitels, V. 28. 29., hat er gesagt: Da Isaak hundert und achtzig Jahr alt gewesen, sei er gestorben. Hernach im 37. Cap. V. 35. sagt er: „Und sein Vater“ Isaac „beweinete ihn.“

2. Um der Ursache willen muß man dies Capitel auch mit Fleiß ansehen und merken, daß was darin begriffen, lange zuvor geschehen sei, ehe Isaak gestorben, oder Joseph verkauft worden ist. Denn Isaak hat noch immer gelebt, bis Joseph zum Regiment über Egypten erhoben und gesetzt worden ist, das sind beinahe dreizehn Jahre gewesen. Und Juda ist hinab gezogen zu einem Mann von Odollam, mit Namen Suah, und nimmt desselben Mannes Tochter zum Weibe, zeugt mit ihr drei Söhne; davon zwei erschlagen worden, da Sela, der Großvater, noch gelebt und sein Bruder Joseph auch noch bei dem Vater frisch, gesund und frei gewesen ist; erst sieben oder acht Jahr nach dieser Zeit wird Joseph verkauft. Darum ist alles, was in diesem Capitel steht, dem Tode Isaaks vorangegangen; welches aus der Rechnung der Jahre Jakobs klarlich geschlossen und bewiesen werden kann. Und diese Ordnung sollte man allezeit vor Augen haben, gleich als in einer Tafel vorgeschrieben, damit der Leser nicht verwirret werde, wenn er darauf nicht würde gute Acht geben.

3. Es scheint aber, daß Juda hinab gezogen sei von seinen Brüdern, so dahzumal nicht zu Hebron, sondern zu Sichem gewohnt haben, an welchen Ort Jakob aus Mesopotamien gekommen war und da er eine Hütte aufgerichtet hatte. Und ich halte dafür, daß er nicht länger denn zwei Jahr bei dem Vater geblieben sei, da er noch ein Jüngling gewesen; sondern habe seines Vaters Haus und seine Brüder verlassen, und sei mit seiner eigenen Habe zu den Nachbarn, den Odollamitern, gezogen, vielleicht von ihnen selbst gebeten oder geladen. Denn man weiß nicht eigentlich, was sonst die Ursache möge gewesen sein, daß er dahin gezogen.

4. Es ist aber doch wahrscheinlich, daß zwischen denselbigen Heiden und ihren Gästen, den Juden, Freundschaft und Bekanntschaft muß gewesen sein. Und ist ein gewisses Anzeichen vortrefflicher Höflichkeit und Freundlichkeit, daß sie solche fremde Gäste in ihrem Lande haben dulden können, die bei ihnen den Acker bauen, Hantierung treiben und sonst mit andern gewöhnlichen Werken dieses gemeinen Lebens sind umgegangen. Es scheint, als seien sie dazumal noch nicht so reich und gewaltig gewesen, wie droben auch gesagt ist. Es sind der Amoriter Sünden noch nicht erfüllt, und sind noch fromm und gastfrei gewesen, bis sie darnach großes Geldes und Gutes, dazu unmäßigen Ueberflusses halben gar verderbt worden sind. Derhalben wird damit angezeigt, daß sie freundlich und gutmüthig gewesen seien, die Juden aufzunehmen und dieselben bei sich zu dulden, und haben das gethan aus natürlichem Gesetz ohne alles geschriebene Recht. Solche Leute waren Aner und Mamre, die Amoriter und Heiden gewesen, 1 Mos. 14, 13. Und droben, Cap. 20, 15., haben wir auch gehört, daß Abimelech und Andere mehr die Väter freundlich beherberget und ihnen Hülfe erzeigt haben.

5. Es kann aber auch wohl die Ursache sein dazu gekommen, daß es Juda verdrossen hat, daß seine Brüder so gar stolz und tyrannisch gewesen, vornehmlich Simeon und Levi, welche sich selbst die Herrlichkeit und den Vorzug der Erstgeburt zugeeignet, daß er darüber also von ihnen gewichen, dieweil er für seine Person etwas stiller und demüthiger war, und ihren Stolz und so gar wilde und scheußliche Sitten nicht wohl hat ertragen können.

6. Lyra hat sich sehr darüber bekümmert, warum doch Juda aus seines Vaters Hause gewichen und weggezogen sei? und führt hier der Juden kindische Possen an, damit sie beinahe die ganze heilige Schrift belehrt haben, nämlich, daß die andern Brüder aus Mitleiden mit dem Vater, dieweil sie gesehen, daß er von Schmerz und Leid, so er getragen, fast verzehrt worden, angefangen haben mit Juda zu zürnen, ihn gescholten haben, darum daß er sie dazu verursacht, daß sie Joseph verkauft hatten. Denn wo er ihnen den Rath nicht hätte gegeben, wollten sie ihn traun nicht verkauft haben. Aber mit der Rechnung der Jahre und rechter Ordnung der Historie werden diese Fabeln sein widerlegt. Denn Juda hat dazumal schon ein Weib und drei Söhne gehabt, ehe Joseph nach Egypten ist geführt worden.

7. Sichem aber liegt ziemlich weit von Odollam. Denn es sind Städte, so nicht in einerlei Stamm gelegen waren. Sichem ist gelegen in Ephraim, welche hernach ist Neapolis genannt worden. Odollam aber ist gegen Abend gelegen gewesen, da der Stamm Juda gewohnt hat. Und scheint es, als sei er durch sonderliche Schickung GOttes in dasselbe Land gezogen, welches darnach seine Nachkommen inne gehabt.

8. Er ist aber eingeführet bei einem frommen ehrlichen Mann, der ihm sehr gütlich gethan und freundlich mit ihm ist umgegangen; und bei demselben ist er eine Zeitlang geblieben, meines Erachtens, nicht als wäre er sein Knecht gewesen, sondern als ein fremder Gast, der seines Viehes gehütet bei Odollam, der Stadt Juda, und daselbst hat er eines Cananiters Tochter gesehen und dieselbe zum Weibe genommen.

9. Aber daran möchte vielleicht Einer zweifeln: Ob er auch seinen Vater hätte um Rath gefragt, beide seines Abschieds halben und da er auch hat wollen an fremden Orten ehrlich werden? Und es scheint, daß die Worte im Text beinahe dahin lauten, als habe er solches ohne den Vater um Rath zu fragen gethan. Denn es wird ja der väterlichen Gewalt nicht gedacht, daß er sie auf des Vaters heißen oder Rath zum Weibe genommen, wie hernach gesagt wird, daß Juda seinem Sohne Ger ein Weib gegeben habe. Und kann doch das auch nicht bewiesen werden, daß er ohne Einwilligung seines Vaters eine Cananitin zum Weibe genommen habe; oder im Fall, daß solches möchte bewiesen werden, wird doch daraus nicht folgen, daß man eben nach diesem Erempel daselbst auch thun möge. Denn er ja selbst mit seinem Sohn Her eine andere und bessere Ordnung damit gehalten. Und dieweil er für und für bei seinem Vater Gunst gehabt, wird er ohne Zweifel auch seiner Lehre und Vermahnung mit höchstem Fleiß gehorcht haben. Denn wir werden hernach noch hören, daß der Vater der andern Brüder keinem seinen Sohn Benjamin hat befehlen oder vertrauen wollen denn allein dem Juda, und mit großen Sorgen auf seine Wiederkunft aus Egypten gewartet hat.

\quellseite{304}

10. Und ich halte nicht dafür, daß er also gewichen und weggezogen sei, daß er darnach zum Vater nimmer wieder gekommen wäre; sondern wird ihn ohne Zweifel bisweilen besucht haben, zu sehen, was er machte, und wie es dem Vater und Großvater Isaac ginge. Und wird vielleicht dazumal auch, da Joseph verkauft worden, zu seinen Brüdern auf Besuch gegangen sein, oder von seinem Wirthe, bei dem er sich etliche Jahre aufgehalten, wiederum in seines Vaters Haus hineingekommen sein. Es wäre denn, daß wir es dafür halten wollten, daß er zu Odollam allein die Weide und sein Vieh gehabt, und wäre also ab und zu, von einem Ort zum andern gezogen, entweder die Schafe zu scheren, oder sonst um anderer Ursachen willen.

11. Weiter aber ist auch des Wortes halben eine Frage, so die Grammatik betrifft, nämlich, ob das Wort „Cananiter“ ein Eigenname sei und verstanden werden soll für Einen von demjenigen Volke der Cananiter; oder aber ob es ein Gattungswort sei, daß es so viel heiße als ein Kaufmann. Denn es wird oftmals in der heiligen Schrift als ein solches Gattungswort gebraucht. Als, in den Sprüchen Salomos am 31. Cap. V. 24.: „Sie macht einen Rock und verkauft ihn; einen Gürtel gibt sie dem Cananiter“, das ist, dem Krämer. Derhalben ist es ungewiß, ob dieser ein Krämer gewesen sei, oder nicht. Denn es steht nichts dabei, daraus man solches abnehmen und verstehen möchte.

12. Das Wort „Suah“ ist nicht der Tochter, sondern des Vaters Name; des Weibes Name aber wird nicht ausdrücklich gesetzt, wird auch im Geschlechtsregister Christi nicht aufgezählt. Aber Thamar ist berühmt und ist eine Mutter Christi geworden. Die Söhne Juda haben die Ehre auch nicht müssen haben, wiewohl ihr Geschlecht ziemlich groß gewesen. So ist nun Judas Weib schwanger geworden, und hat ihren ersten Sohn geboren, welchen sein Vater Juda Er hat heißen lassen. Den andern zwei Söhnen, Onan und Sela, hat die Mutter ihre Namen gegeben. Und der Name Onan ist hergekommen von Schmerzen und Schwere der Geburt; wie Rahel, da sie von großen Schmerzen der Geburt überwunden ward, ihren Sohn auch Benoni genannt hat. Dem Dritten hat sie einen besseren Namen gegeben und läßt ihn Sela heißen, darum daß sie denselbigen etwas leichter und glücklicher geboren hat.

13. Das letzte Stück im Texte: „Und er war zu Chesib, da sie ihn gebar“, hat Hieronymus also gegeben: Da der geboren war, hat sie aufgehört zu gebären. Denn das pflegt den Dolmetschern also zu begegnen, daß sie bisweilen ein Wort, das ein Eigenname ist, als ein Gattungswort deuten, darum daß gemeiniglich die Eigennamen von den Gattungswörtern herkommen und in derselbigen Bedeutung auch mögen verstanden werden. Chesib aber ist der Name einer Stadt in Juda, welche sonst Achsib genannt wird, Micha 1, 14., da auf diesen Namen der Stadt gesehen, und dabei angezeigt wird, es solle der König Israel daselbst fehl schlagen; denn also lauten die Worte des Propheten: „Der Stadt Achsib wirds mit den Königen Israel fehlen“ etc. Es ist auch noch ein Achsib gewesen, im Stamme Asser gelegen, gegen Mitternacht; aber diese in Odollam gehört zum Stamme Juda, oder vielmehr zum Stamme Dan, gegen Mittag. Hieronymus folgt dem Aquila, und hat es also gegeben, daß es ein Gattungswort sein soll, welches mir nicht gefällt. Denn solche Freiheit durch sonderliche Figuren oder Weise zu reden, und neue Deutungen, die etwas fremd sind, zu dichten, gefällt mir nicht; wie denn das sehr hart lautet: Sie ist in der Lüge gewesen, sie hat aufgehört zu gebären. Denn also lautet es im Hebräischen, wenn es ein Gattungswort ist: Es ist mit ihrem Gebären Lüge gewesen, oder, die Geburt hat ihr gefehlet.

14. Derhalben verstehe ich es also, daß es ein Name sei einer sonderlichen Stätte, welche in der heiligen Schrift oft erwähnt wird. Denn es sind zwei Städte gewesen, die diesen Namen gehabt, wie Hieronymus auch sagt; und dieselbige Deutung reimt sich mit der Historie. Denn Mose will beweisen, was er gleich im Anfang angezeigt hat, nämlich, daß Juda von seinen Brüdern geschieden sei, sintemal er mit seinem Weibe nicht zu Sichem gewohnt, sondern seine Eltern und Brüder verlassen hat, und hinabgezogen ist zu einem Manne von Odollam, da er ein Gast und Fremdling gewesen, und ist daselbst von einem Orte zum andern gezogen, hat des Viehes gehütet nach Gelegenheit der Weide, wo er dieselbe am besten hat haben können. Denn bald hat er zu Achsib, bald zu Thimnath, bald aber zu Odollam gewohnt: welches alles Städte im Stamme Juda gewesen sind, oder aber vielmehr im Stamme Dan gelegen. Simeon hat hernach zu Thimnath Hochzeit gehalten.

15. So wil nun Mose anzeigen, daß Juda keinen gewissen Sitz gehabt, da er gewohnt, wie denn seine Väter auch nicht gehabt; sondern während seine Brüder zu Sichem bei dem Vater geblieben, ist er umhergezogen, und hat mit Hira von Odollam viel Gemeinschaft gehabt, der sein guter Freund und Schwiegervater war. Und es scheint, daß daselbst noch etliche fromme Cananiter übrig gewesen, die Juda viel Wohlthaten erzeigt haben. Denn wo sie unfreundlich mit ihm umgegangen wären, hätte er mit ihnen keine Schwägerschaft gemacht. Und hernach, da die zu Sichem erschlagen worden, sind sie ohne Zweifel den Söhnen Jakobs zu Hülfe gekommen, die Gewalt, so ihnen widerfahren, an den Sichemitem zu rächen, und Jakob zu schützen, daß er nicht mit seinem Hausgesinde von den Nachbarn erschlagen würde.

16. Wer aber behalten will, wie es Hieronymus im Lateinischen gegeben und verdolmetscht hat, und heimliche Deutungen suchen will, der mag es zwar meinetshalben wohl thun: als sollte hiermit angezeigt sein, daß Judas Weibe die Geburt nunmehr gefehlet und daß sie aufgehört habe zu gebären. Denn es hat Juda darnach in neun oder zwölf Jahren nicht mehr Kinder gezeugt; wie wir hernach sehen werden.

17. Es ist aber fleißig zu merken, daß alles, so bisher erzählt, vor dem Absterben der Rahel und Debbora geschehen sei, ehe Benjamin geboren, Bilha geschändet und Dina geschwächt worden ist. Denn Juda ist von seinen Brüdern weggezogen im andern oder dritten Jahr, nachdem sein Vater Jakob zu Sichem angekommen war und acht oder neun Jahre in guter stiller Ruhe gelebt hatte. Aber im achten oder neunten Jahre oder fast um dieselbe Zeit sind die vielen Trübsale als mit einem Sturm zugleich hereingefallen, als, daß Dina geschwächt, Thabel gestorben etc., und zuletzt Joseph verkauft worden ist. Was aber Mose weiter erzählt von der Heirath Judas, und vom Absterben seiner Söhne und auch seines Weibes, das ist alles zuvor geschehen.

18. Darnach soll man auch ansehen die Rechnung der Jahre Jakobs und Judas. Denn wo du das, das in diesem Capitel erzählt wird, mit den Jahren des Alters Judas vergleichen wirst, wird es beinahe als unglaublich scheinen, und daß in so kurzer Zeit zwischen dem, daß Jakob nach Egypten hinab gezogen und Juda geboren ist, nicht Alles habe geschehen können. Denn da Jakob schon vier und achtzig Jahre alt gewesen, hat er Lea und Rahel zu Weibern genommen, und da er nach Egypten hinab gezogen, ist er hundert und dreißig Jahre alt gewesen. Juda aber ist im vierten Jahre geboren, nachdem Jakob ehelich geworden: so bleiben nun schier zwei und vierzig Jahre zwischen der Geburt Judas und daß Jakob hinab nach Egypten gezogen ist. In der kurzen Zeit ist Juda ein Ehemann geworden und hat auch seine Schwiegertochter geschwächt. Denn da sein Weib genommen, war er, so denkt man, zwanzig Jahr alt.

19. Nun wollen wir auch annehmen, daß die drei Söhne innerhalb zweier Jahre geboren seien, welchen darnach, da sie erwachsen und an ihr zwanzigstes Jahr gekommen sind, der Vater die Töchter der Cananiter gefreiet hat. Bald aber, da die Söhne gestorben, ist Thamar ein Jahr Witwe geblieben. Des andern Jahres aber, so da nächst gefolgt, hat Juda, ihr Schwiegervater, Blutschande mit ihr begangen, und hat also mit seiner Schwiegertochter zwei Söhne gezeugt, nämlich, Perez und Serah. Wenn man nun diese Jahre alle zusammenrechnet, kommen vier und vierzig Jahre, so doch nicht mehr denn zwei und vierzig Jahre sind zwischen der Geburt Judas, und daß Jakob hinab in Egypten gezogen. Und ist doch dies noch das andere Geschlecht. Denn hernach im 46. Capitel werden aufgezählt zwei

\quellseite{305}

Knaben, Judas Enkel, nämlich, Hezron und Hamul, so Perez gezeuget, ehe sie nach Egypten gezogen, welchem ja desgleichen auch zwanzig Jahre müssen zugegeben werden, ehe er zur Ehe und Kinder zu zeugen auferwachsen und Alters halben tauglich geworden ist. Und wenn diese zwanzig Jahre zu den vorigen gethan werden, wird es sechzig Jahre machen. Das ist eine solche Zeit, so die gemeine und wahrhaftige Rechnung der Jahre weit übertrifft.

20. Derhalben schließen wir nun daraus nothwendig, daß zu der Zeit die Natur beide an Männlein und Fräulein viel stärker und vermöglicher gewesen, denn sie jetzund ist, also daß die Eheleute haben im zwölften und dreizehnten Jahre ihres Alters Kinder zeugen können. Denn wir werden sonst auf keine andere Weise die Zeitrechnung mit der Historie ausgleichen können. Derhalben müssen wir es dafür halten, daß Juda nicht mehr denn zwölf Jahr alt gewesen, da er ehelich geworden ist, und daß innerhalb zwei Jahren darnach seine drei Söhne geboren sind; zu welcher Zeit er vollkommenlich vierzehn Jahre alt gewesen. Dazu sind gekommen die zwölfte Jahre seines Sohnes Ger, welchem er in demselben Alter Thamar zum Weibe gegeben, da Juda vier und zwanzig Jahre alt gewesen; in welchem Jahr auch Ger und Onan erschlagen worden sind, oder aber ja des nächsten Jahres darnach, nämlich, im fünf und zwanzigsten Jahr. Derhalben hat Thamar nicht geboren bis an das sechs und zwanzigste Jahr Judas, und hat unterdeß gewartet, bis Selah erwachsen ist, der zwei Jahre jünger gewesen denn sein Bruder Ger, vielleicht von elf oder zwölf Jahren. Wir wollen aber annehmen, daß sie in ihrem Hause ein Jahr lang als Wittwe gesessen, und alsdann will folgen, daß Juda sieben und zwanzig Jahre gewesen, da er seine Schwiegertochter geschwächt. Denn wir müssen der Wittwenschaft Thamar nicht mehr Zeit geben, darum daß zur selbigen Zeit die Weiber sich \unsicher{erküfftet} vor der Unfruchtbarkeit, Verachtung und Fluch, welchem sie sich auf das höchste beflißen zu entfliehen. Darnach weiter sind aus der Blutsschande im acht und zwanzigsten Jahr zugleich Zwillinge geboren worden, von welchen der eine, Perez, im zwölften Jahr auch ein Weib genommen, welches das vierzigste Jahr Judas gewesen; und der hat Hezron und Hamul gezeuget. Denselbigen müssen auch zwei Jahre zugegeben werden; es wäre denn, daß sie vielleicht Zwillinge gewesen. Zuletzt aber, da Juda nun zwei und vierzig Jahre alt gewesen, ist er nach Egypten hinab gezogen. Diese Rechnung der Jahre scheint der Ordnung der Historie am allerältesten zu sein, und wer da will, mag dieselbe also in unserem Chronicon suchen, darin sie noch etwas genauer und fleißiger ausgerechnet ist.

21. Das aber muß man allerdings zugeben, daß die Väter zu der Zeit viel eher Kinder gezeuget haben, da die Natur noch stärker gewesen, denn sie jetzund in dieser letzten Zeit der Welt ist. Also halte ich es wahrlich dafür, daß es ein sonderlich Werk Gottes gewesen sei, daß Juda in demselben Alter tauglich hat sein können, Kinder zu zeugen. Jetzt zur Zeit sind ihrer noch sehr Wenige, die im achtzehnten Jahr ihres Alters reif sind, ehelich zu werden und Kinder zu zeugen. Die Patriarchen aber sind dazu innerhalb weniger Jahre der Stärke halben tauglich geworden, daß sie auch, ehe sie recht mannbar waren, Kinder gezeuget haben.

22. Dieweil aber dies Capitel nichts Anderes in sich begreift, denn allein die Erzählung, wie Juda Kinder gezeuget, und von seinen Brüdern gewichen und weggezogen sei, und über das auch die allerschändlichste Blutschande, so er mit Thamar, seiner Schwiegertochter, begangen, fällt abermals hier eine gemeine Frage ein: Warum doch Gott und der Heilige Geist diese so schändlichen und unsäglichen Dinge habe schreiben und verwahren lassen, daß sie in der Kirche erzählt und gelesen würden? Denn es scheint, als habe er daran nicht recht gethan, daß er andere größere Dinge übergangen, und diese, so nicht werth sind, daß man sie an den Tag bringen und im Gedächtniß behalten sollte, dessen gleichwohl werth geachtet habe. Denn wer will glauben, daß solche Dinge zur Seligkeit oder Besserung der Kirche können gelehrt werden? Er erzählt, wie Juda von seinen Brüdern entwichen, eine Fremde zum Weibe genommen und darnach seine Schwiegertochter durch Blutsschande befleckt habe. Es wäre ja besser gewesen, daß dieses zugedeckt und auf ewige Zeiten vergessen worden wäre.

23. Die Canonisten und Papisten können sich dieser Frage leichtlich entledigen, als die da sagen: Das Alte Testament sei ein Buch vom Beschlafen und Hurerei der Juden. Aber sie geben damit ihren Unverstand an den Tag, und zeigen genugsam, daß sie Säue und grobe Esel sind, indem sie von solchen Dingen reden, davon sie doch keinen Verstand haben und davon sie gar nichts wissen, und schmähen damit auch den Heiligen Geist, als ob er nicht wüßte, was man schreiben oder der Kirche und Gemeinde Gottes vorhalten solle. Hieronymus hat auch auf diese Frage geantwortet. Aber es hat die Ursache nicht allein, so er hier angezogen, daß Christus sein Geschlechtsregister vom Anfang in der heiligen Schrift habe wollen rühmen lassen; und dieweil Thamar seine Mutter gewesen, habe deshalb ihr Name jedermann sollen bekannt sein: sondern über das wird uns hier ein großer Trost vorgehalten, daß die heiligen Väter und solcher großen Patriarchen Kinder, Juda und die Andern, beschrieben werden als solche Menschen, so voller Schwachheit und voll großer Mängel oder Sünde, damit diese arme Natur beladen ist, gewesen, und daß sie Gott wunderbarlich mit dem Heiligen Geiste regiert und geführt hat, doch also, daß er verhängt, daß sich an ihnen ihre eigene angeborne Neigung, das ist, die Sünde und Frucht der Erbsünde, gereget und hervorgebrochen ist. Denn das ist ja eine erschreckliche Sünde Beider, nämlich, Judas und Thamars, welche man doch mit nichten gar entschuldigen kann, wenn sie gleich etwas verkleinert werden mag.

24. Es werden aber solche Exempel erzählt uns zur Lehre und Trost, und auch zu Stärkung unsers Glaubens, daß wir dabei die große Gnade und Barmherzigkeit Gottes betrachten sollen, welcher nicht allein die Frommen und Gerechten, als, Abraham, Isaak und Jakob selig gemacht, sondern auch die Ungerechten, als, Juda, Thamar, Ruben, Simeon, Levi, welche sehr große Sünder gewesen. Daß also niemand seiner Gerechtigkeit oder Weisheit halben vermessen oder stolz sein und wiederum auch niemand seiner Sünden halben verzagen soll.

25. So werden nun in der heiligen Schrift die schönsten Exempel der lieben Väter gerühmt, Abraham, Isaak und Jakob, welche wie Lichter der ganzen Welt, und in der Kirche oder Gemeinde Gottes wie Sterne, die Sonne oder Morgenröthe, vollkommen im Glauben, Hoffnung und Liebe gewesen sind. Es werden aber darin auch zugleich große erschreckliche Sünder beschrieben, die so böse und schändlich sind, als jene frommen und gerecht waren. Ruben thut seinem Vater wohl so eine große Schalkheit als die Gerechtigkeit Abrahams gewesen ist. Also wird uns die Gerechtigkeit der größten Heiligen vorgehalten, und auch die Sünde derer, so auf das äußerste böse gewesen, die doch aus einem Geschlecht und einerlei Geblüt gekommen sind.

26. Dies ist die Predigt von der Buße und vom Glauben oder Vergebung der Sünde, damit ja niemand um seiner Gerechtigkeit willen vermessen oder stolz sei, und daß auch die, so gefallen sind, deshalb nicht verzagen sollen. Denn es werden auch Judas und der andern Väter Gebrechen und Fehler, ja, auch ihre großen ungeheuren Fälle oder Sünden erzählt. Und auf diese Weise sollen die armen Herzen, so sich der Sünden schuldig wissen, aufgerichtet und getröstet werden, nämlich, daß man zu einem armen Sünder sage: Verzage nicht, denn Gott heißt dich ihm vertrauen und seiner Verheißung glauben; und er kann dich auch gerecht, heilig und selig machen, wie er jene armen Sünder selig gemacht. Es ist eben derselbe Gott, der die Stolzen erniedrigt und die Niedrigen erhöht. Er will aber nicht haben, daß wir uns auf unsere Gerechtigkeit verlassen, oder in Sünden verzweifeln sollen, sondern daß wir auf seine Gnade und Barmherzigkeit allein vertrauen.

27. Darum sind diese Exempel der Kirche und Gemeinde Gottes sehr nöthig. Denn wo wollten wir bleiben? Und was könnten wir doch für Hoffnung haben, wo nicht Petrus Christum verleugnet hätte, die Apostel alle sich an ihm geärgert? wo nicht Mose, Aaron und David auch gefallen wären? Derhalben hat Gott mit diesen Exempeln die armen Sünder trösten wollen und also sagen: Bist du gefallen, kehre wieder um, die Thür der Gnaden steht dir offen; du aber, der du dich deiner Sünden schuldig weißt, sei deshalb nicht vermessen oder stolz; sondern verlaß dich und setz dein Vertrauen auf meine Gnade und Barmherzigkeit.

28. Dies ist die rechte Ursache, warum solche Erzählungen, so voller schändlicher

\quellseite{306}

Aergernisse sind, mit untergemengt werden in die Legenden oder Historien der heiligen Patriarchen. Mose hat mit sonderlichem Fleiße zugleich die allerbößesten Früchte mit den besten erzählt. Wie im Neuen Testament Judas auch mit ist unter den Aposteln, zum großen wunderbaren Exempel, der Kirche zum Trost, damit ja die Gottseligen nicht verzagen. Aber nicht, daß daran die Gottslosen ein Exempel haben sollen, dem sie also nachfolgen müßten; sondern daran sich die Bußfertigen halten mögen, daß sie nicht alle Hoffnung auf Gottes Gnade und Barmherzigkeit fallen lassen, als ob sie nicht könnten selig werden.

29. Der Heilige Geist hat diese Stücke nicht wollen schreiben lassen, damit ein Exempel zu stellen, darnach auch also zu leben, sondern zum Trost; nicht den Muthwilligen hindert die Freiheit zu sündigen damit zu bestärken, sondern die Hoffnung und den Glauben damit zu erwecken. Denn er befiehlt und hält uns dadurch vor Vergebung der Sünden, so ja stark genug ist; sintemal er uns damit also tröstet, die, so gefallen sind, in der Kirche wiederum aufzurichten, zu erwecken und zu stärken. Und diese Lehre findet man in den Büchern der Heiden oder Papisten nicht; der Heilige Geist hat sie allein im Bibelbuch gelehrt.

30. Danach wird uns auch vor Augen gestellt, wie es um die allerheiligsten Menschen gestanden, nämlich, daß sie ja nicht steinern oder eisern, und nicht grobe Klöße gewesen sind, ohne Verstand, ohne Fleisch und Blut, sondern sind Menschen und uns allenthalben gleich gewesen: und wo sie darüber etwas Gutes gethan, daß sie solches nicht aus eigenen Kräften unseres Fleisches, sondern durch die Gabe des Heiligen Geistes gethan haben. Derhalben kann es wohl geschehen, daß uns eben dasselbe auch begegne, was ihnen widerfahren ist, da sie noch in demselbigen Fleische gelebt, welches wir auch an uns tragen.

31. Zum Andern hat auch der Heilige Geist gesehen auf den Messias und die Geburt des Sohnes Gottes; welches die vornehmste Ursache ist. Denn dieser Fall hat müssen geschehen auch in derselben Linie, darin der Sohn Gottes sollte geboren werden. Der höchste Patriarch Juda, ein Vater Christi, hat diese unsägliche Blutschande begangen, damit Christus geboren werde aus einem solchen Fleische, das überaus sündig und mit einer sehr schändlichen Sünde befleckt wäre. Denn Juda hat mit einer Hure, so mit Blutschande befudelt, nämlich, mit seiner Schwiegertochter Zwillinge gezeuget, daher nochmals die Linie unsers Herrn und Heilandes Christi geführt worden ist. Da muß Christus in seinem Fleische ein Sünder werden, so schändlich, als er immer werden kann. Das Fleisch Christi kommt aus einem solchen Beilager, so mit greulicher Blutschande befleckt; desgleichen das Fleisch der Jungfrau Maria, seiner Mutter, auch und aller Nachkommen Judas: auf daß damit angezeigt würde der unaussprechliche Rath Gottes von seiner Gnade und Barmherzigkeit, da er das Fleisch oder menschliche Natur hat angenommen von einem unreinen Fleisch, welches greulicher Weise befleckt gewesen.

32. Denn die Scholastiker disputiren davon, ob Christus von sündlichem oder reinem Fleisch geboren sei? Ob Gott vom Anfang der Welt ein rein Tröpflein Fleisches behalten habe, davon Christus habe sollen geboren werden? Darauf antworte ich nun also: Daß Christus wahrhaftig von wahrhaftigem und natürlichem menschlichem Fleisch und Blut geboren sei, welches in Adam der Erbsünde halben verderbt gewesen; aber doch also, daß es wiederum hat können geheilet werden. Wie mir denn, die wir mit sündlichem Fleisch umgehen sind, glauben und hoffen, es werde geschehen, daß am Tage unserer Erlösung das sündliche Fleisch gereiniget, und von aller Schwachheit, vom Tode und Schmach abgesondert werden soll. Denn die Sünde und Tod sind solche Beschwerden, so von uns abgesondert und davon wir noch können erlöset werden. Derhalben, da es an die Jungfrau kommen ist, oder an das Tröpflein des jungfräulichen Blutes, da ist erfüllt worden, was der Engel Luc. 1, 35. gesagt hat: „Der Heilige Geist wird über dich kommen, und die Kraft des Höchsten wird dich überschatten.“ Es ist zwar der Messias nicht geboren worden aus Kraft Fleisches und Blutes, wie Joh. 1, 13. gesagt wird: „Nicht von dem Geblüt, noch von dem Willen eines Mannes“ etc., aber doch gleichwohl aus unserem Fleisch und aus dem verderbten Blut hat er wollen geboren werden. Aber in dem Augenblick der jungfräulichen Empfängniß hat der Heilige Geist das sündliche Fleisch gereinigt und geheiliget, und das Gift des Teufels und Todes, welches die Sünde ist, davon abgewaschen wird.

33. Derhalben so wird dies um Christi willen geschrieben, welchen der Heilige Geist hat wollen in die Sünde so tief stecken, als es immer möglich gewesen. Darum hat er müssen als mit Blutschande befleckt und aus solchem befleckten Fleisch geboren werden.

34. Zum Dritten ist auch das zu merken, daß der Heilige Geist in dieser Historie hat anzeigen wollen, daß Gott also den Samen Abrahams angenommen und erwählt habe, daß er gleichwohl die armen Heiden nach dem Fleisch nicht verwerfen wolle. Die Juden haben nicht Ursache, sich zu rühmen, daß sie allein Abrahams Same seien. Denn nach Jsaat sind sie nicht mehr allein Abrahams Same; denn Rebekka ist ja nicht von Abraham, sondern von Nahor hergekommen. Also, Lea und Rahel sind Bethuels Kinder gewesen; da sind die Samen vermischt worden: der Same der Heiden ist mit Abrahams Samen vermischt. Desgleichen Thamar ist ein canaanisch Weib.

35. Derhalben hat Christus diese canaanische Mutter nicht verschmähet, sondern hat wollen aus dem Samen des verworfenen Volkes geboren werden, daß sich ja die Juden nicht über die Maßen erheben sollten und von ihrem Geblüt rühmen. Denn sie müssen bekennen, daß sie auch Idumäer und Parther sind; und daß auch dem natürlichen Rechte nach die Heiden unsers Herrn Jesu Christi Mutter, Brüder, Vettern, Muhmen und Schwestern sind, zum wenigsten um der Mütter willen. Denn sie sind alle hergekommen von dieser Thamar, welche die Mutter des ganzen Stammes Juda ist, und das Fleisch Perez ist von der Canaanitischen und von Abrahams Samen zusammengesetzt. Derhalben muß Israel von Ruth wegen die Heiden für seine Blutsfreunde, Verwandte und Brüder erkennen und dasselbe um der Mütter willen. Denn die Mutter macht auch Schwestern, Brüder und Ehefrauen, nicht weniger als der Vater. Wiewohl aber nun die Ordnung des Geschlechtsregisters nicht von den Müttern hergeführt wird und die Weiber unter den Nachkommen nicht mit gerechnet werden: doch werden sie auch mit aufgezählt und wird ja gleichwohl dieser Weiber ausdrücklich und sonderlich gedacht, nämlich, der Thamar, Ruth, Rahab, Bathseba. Da kommen die Heiden mit in die Gemeinschaft des Israelitischen Volkes, und daß sie ihre Mitgesellen werden nicht allein der Religion und Glaubens, sondern auch einerlei Fleisches. Denn es werden ja wahrhaftig Mann und Weib ein Fleisch. Darum sind nun die Juden und Heiden auch Ein Fleisch und von einem Fleisch geboren.

36. Auf diese Weise hat Gott den Stolz und die Thorheit des allerhoffärtigsten Volks, der Juden, wollen dämpfen, welche den Heiden immer spinnefeind gewesen. Denn sie werden auch wider ihren Willen zugeben müssen, daß Juden und Heiden auch natürlich und nach dem Fleisch Brüder sind; denn sie werden ja diese Mutter, Thamar, nicht verwerfen können; wiewohl die Vaterschaft und Geschlechtsregister von den Müttern nicht hergeführt wird. Moab sagt von Ruth: Dies ist meine Tochter, und eine Mutter des Messias, darum folgt, daß auch der Messias mein Sohn ist. Dasselbige hat auch Canaan von Thamar sagen können. Und endlich Manasse und Ephraim, die Söhne Josephs, sind die nicht auch von einer egyptischen Mutter geboren? Darum folgt auch, daß der königliche Stamm in Israel egyptisch ist, und Egypten herrscht über das Volk Israel, wo nicht aus der Linie der Väter, doch wegen der Blutsverwandtschaft und der Natur. Zu dem kommt das auch noch, daß die Aerzte sagen, die Mutter gebe dem Kindlein mehr Nahrung als der Vater; denn sie nährt das Kind zehn Monat lang in ihrem Leibe, allein von ihrem Geblüt, und darnach, wenn das Blut in Milch verändert wird, säugt sie das Kindlein damit.

37. Derhalben haben die Juden keine Ursache, sich zu rühmen, sondern sich zu demüthigen und das mütterliche Geblüt zu erkennen. Denn vom Vater sind sie Israeliten; aber von der Mutter sind sie Heiden, Moabiter, Assyrer, Egypter, Cananiter. Und das hat GOtt auch wollen anzeigen, daß der Messias ein Bruder und Bruders Sohn oder Vetter sein würde beider der Juden und Heiden, wo nicht nach dem väterlichen Geschlechtsregister, doch wegen der mütterlichen Natur, also daß sein Unterschied zwischen Juden und Heiden sei, denn allein daß Mose hernach das jüdische Volk mit sonderlichem Gottesdienste und politischem Regiment von den Heiden abgesondert hat.

\quellseite{307}

38. Dies aber ist darum geschrieben, daß jedermann wissen sollte, daß der Messias Juden und Heiden zugleich in ein und dieselbe Kirche versammeln würde; wie sie auch von Natur und Geblütes wegen mit einander verbunden sind. Und darum habe ich es auch droben dafür gehalten, daß das Wort „Canaan“ nütze als ein Gattungswort verstanden werden. Denn Christus hat wollen von den Cananitern geboren werden, als einem heidnischen Volke, und hat auf solche Weise wollen anfangen zwischen Juden und Heiden Frieden zu machen. Denn es muß doch Alles in Christo wiederum versammelt und gleichsam in Eine Ordnung zusammengebracht werden; denn er ist nicht allein der Juden GOtt, sondern auch der Heiden. Wir Heiden sind zwar eigentlich von Japheth hergekommen; aber Ham und Canaan sind dem Sem blutsverwandt gewesen. Wie wir nun etwas fremder sind, kommen wir doch eben desselben Rechtes wegen auch zur Gemeinschaft und Einigkeit Christi, wie die Assyrer und Egypter dazu gekommen sind.

39. Dies habe ich gesagt wider das eitele Rühmen der Juden von Abrahams Samen. Denn das ganze Volk Israel ist von Lea und Rahel hergekommen; welches ja nicht Abrahams Geschlecht gewesen; der Stamm Juda hat seinen Ursprung von Thamar, und endlich David selbst ist auch von Moab hergekommen. Derhalben ist der Juden Stolz billig sehr zu strafen, damit sie die Heiden ganz feindlicher Weise verachten und hassen, so sie doch ihre Blutsfreunde sind um der Thamar willen, ihrer Mutter, und auch um Anderer willen, welche mir zuvor aufgezählt, durch welche der Heiden und Juden Geblüt vermischt worden.

40. Und dies sind nun die Ursachen, warum der Heilige Geist diese Historie nicht hat übergehen wollen: Erstlich denen, so gefallen sind, und den Bußfertigen zum Trost; darnach zum Preis der Heiden. Wiewohl die Mutter in das Geschlechtsregister nicht mit gehören, welche Ehre allein den Vätern gebührt, so haben aber doch gleichwohl beide Vater und Mutter einerlei Fleisch und Blut, und in den Nachkommen hebt sich die Blutsfreundschaft mit Perez und Juda an. Thamar ist nun nicht mehr Schwägerin, sondern Mutter nach der Natur. Und ist die Mutter ebensowohl des Kindes Substanz, Blut und Leib, als der Vater.

41. Dasselbige aber sollte man den schändlichen unflätigen Canonisten wiederum vorgehalten haben, welche dergleichen Historien, so vom Heiligen Geist ausgegangen und beschrieben sind, für schändlich halten und verachten, und halten vielmehr hart über dem gottlosen Gesetze der Päpste von dem unreinen Ehelibat oder ehelosen Stande, damit sie die rechte Keuschheit und das ganze Leben der Eheleute also verfinstert haben, daß auch viele gute fromme Herzen sich vor dem Ehestande als vor einem gottlosen Stande gescheuet haben, und daß die Jugend nicht mehr zur Haushaltung und andern guten Werken ist auferzogen worden, so dazu gehörig und nothwendig sind. Daraus denn erschreckliche Stricke der armen Gewissen und sonst viele unzählige Unordnungen gefolgt sind in Sachen, das Haus- und Weltregiment belangend, da man oftmals das Erbe nicht recht ausgetheilt.

42. Darum ist es des leidigen Satans Wort, wenn sie sagen, daß in der Genesis allein erzählt und geschrieben sei vom Beischlafen der Juden. Denn eben der Satan hat sich zu allen Zeiten unterstanden aus großem, bitterem Haß des menschlichen Geschlechts, diese arme Natur zu entstellen, zu beklecken und zu verwüsten, darum daß er gesehen hat, daß GOtt dieselbige annehmen und mit dem ewigen Leben ehren würde. Das muß dem Teufel wehe thun und verdrießen. Und wir sehen zwar, daß diese arme menschliche Natur durch GOttes Verkündigung dem Spott und Gewalt des Teufels unterworfen ist; aber es hat doch gleichwohl GOtt angefangen sie zu zieren, und wird sie noch im zukünftigen Leben mit größerer Herrlichkeit überschütten und krönen. Es sollte billig in diesem Leben jedermann von der Kinderzucht viel halten und dieselbe auch begehren; wie sie denn sonst auch von allerlei Thieren begehrt wird. Aber unter den Menschen sind ihrer viele, so sie gar verachten; das hat der Papst und der Teufel angerichtet. Es ist zwar kein Bauer, der gern eine Kuh, so unfruchtbar ist, haben oder dulden würde; sollte denn nun die Fruchtbarkeit in der Natur der Menschen nicht viel mehr angenehm und zu begehren sein als in unvernünftigen Thieren? Hiezu kommt auch noch die große Schwachheit und Unordnung an den Menschen, so von der Sünde hereingeführt worden ist. Aber harre doch, GOtt wird dies arme Fleisch wiederum reinigen und erlösen, wie er es schon zuvor allbereit mit der größten Ehre und Herrlichkeit, von wegen der unaussprechlichen Vereinigung der menschlichen Natur mit der göttlichen, überschüttet hat.

43. Dies sage ich darum, damit niemand die sehr heiligen Patriarchen der Schande beschuldige; denn wir sehen, wie mit so großer Begierde auch die heiligen Weiber nach Kindern verlangt, also daß sie auch beinahe unsinnig geworden, daß sie ja gerne wollten Kinder haben, wie droben, Cap. 30, V. 1., von der Rahel gesagt ist; denn das ist allein ihr einiger und sonderlicher Schmuck und Herrlichkeit gewesen. Und Salomo hat unter den drei Dingen, so nicht zu ersättigen sind, auch den Leib eines Weibes mit gezählt, Spr. 30, 16. Denn das ist ihr also in die Natur gepflanzt und sie ist dazu geschaffen, daß sie soll fruchtbar sein und Kinder gebären. Darum kann sie nicht anders, denn solche Mehrung begehren; wie denn GOtt selbst 1 Mos. 1, 28. sagt: „Seid fruchtbar und mehret euch.“ Und ein Weib, die nicht gebieret, begehrt doch ohne Aufhören, daß sie Kinder gebären möge, und wird damit nimmer gesättigt. Sie hört nicht auf. Bei uns aber ist diese Ordnung der Natur verändert worden wegen böser Gewohnheit, und um der päpstlichen Canonen willen, daraus wir gelernt, Haut- und Weltregiment zu verachten. Derhalben soll man diese Exempel fleißig treiben und wiederum heranziehen aus dem Unflathe der Cananiten, daß wir von denselben recht urtheilen und in solchen Erzählungen des Heiligen Geistes Absicht sehen und verstehen können.

\begin{quote}
\textbf{V. 6. 7. Und Juda gab seinem ersten Sohn ein Weib, die hieß Thamar. Aber er war böse vor dem Herrn; darum tödtete ihn der Herr.}
\end{quote}

44. Was an diesem Orte erzählt wird von der Hochzeit Ers, ist gefolgt, nachdem Joseph verkauft war, und nach andern Trübsalen mehr, so kurz hiebevor sind erzählt worden. Droben aber, da die Heirath Judas beschrieben wird, gedenkt Mose der väterlichen Gewalt nicht. Aber ich habe darum keinen Zweifel, er habe sich mit seines Vaters Jakob Willen und Rath in diesen Stand begeben. Alhier ist aber nun ein klarer Text, daß Juda selbst seinem Sohne ein Weib gegeben habe; und ich halte dafür, daß solches geschehen sei, da sein Sohn Er zwölff Jahre gewesen und Isaak auch noch gelebt, und da Joseph schon verkauft und nach Egypten geführt war.

45. Das Weib hat Thamar geheißen: wessen Tochter aber dieselbe gewesen sei, weiß man nicht. Wir haben die Vermuthung, sie werde eine Cananitin gewesen sein; wie denn Judas Weib auch war. Daß aber ihrer also gedacht wird, ist hoch vonnöthen, und wird dies Capitel allein um der Thamar willen geschrieben; denn sie ist eine Mutter unsers Heilandes, des Sohnes GOttes, um dess willen die ganze heilige Schrift gegeben worden ist, daß derselbe möchte bekannt und gerühmt werden. So ist nun, sage ich, der Messias von dieser Thamar hergekommen und geboren worden, wiewohl durch Unzucht und Blutschande. Denselben sollen wir in diesem Buche suchen und erkennen lernen; und wiewohl es ungewiß ist, ob sie eine Cananitin oder Amoritin gewesen (denn der Name ihres Volks, von dem sie hergekommen, wird nicht so klarlich ausgedrückt), hat sie doch ohne Zweifel Canaan angehört.

46. Also wird nun Christus theilhaftig des Blutes der Cananiter, und hat also an sich selbst einen Leib, zusammengezeugt vom Samen Abrahams und Hams oder Canaans, daß er damit bald im Anfang bezeugen möchte, daß er die Heiden nicht verwerfen wollte, die er aufgenommen, und gewürdigt hat, in seine eigene Person anzunehmen. Die Bösen hat er zwar gewaltig gestraft; wie dasselbe der gottlosen Cananiter und Sodomiter Exempel anzeigen, welche er vertrieben oder gar vertilgt hat; wie er denn der Juden als der natürlichen Zweige auch nicht verschont hat, daß er sie nicht auch sollte verstoßen haben. Und der Ruhm des Geblütes, es sei der Väter bei den Juden, oder der Mütter bei den Cananitern, ist ihnen gar nichts nütze gewesen; denn Christus will nicht Kinder des Fleisches haben. Er allein ist ein Kind des Fleisches ohne alle Sünde des Fleisches; von den Andern allen bleibt dieser Spruch fest stehen, Joh. 1, 13.: „Welche nicht von dem Geblüt, noch von dem Willen des Fleisches geboren sind“ etc. Also mögen nun beide Juden und Cananiter sich des Fleisches rühmen, daß die Sünde weggegangen und sie davon gereinigt werden; wie der Sohn Gottes darum menschliche Natur hat an sich genommen, so von aller Sünde ganz rein gewesen, auf daß er damit unser Blut und Fleisch, beide der Heiden und auch der Juden reinigte. Da steht nun die liebe Thamar; und ist Christus der Juden Retter, nicht Schwager, sondern Verwandter. Juda aber, da er seinem Sohne ein Weib gegeben und ein Großvater oder Schwiegervater geworden, ist vier und zwanzig Jahre alt gewesen.

\quellseite{308}

47. Nun setzt Mose noch eins hinzu, nämlich, daß Ger sei böse gewesen vor dem Herrn, darum habe ihn auch der Herr getödtet. Warum hat er aber dieses heiligen Geblütes nicht verschont, als des Sohnes des Patriarchen Juda? Darauf antworte ich also: Die Kinder des Fleisches werden nicht für Samen gerechnet, sondern des Geistes, Röm. 9, 8. Es muß der Heilige Geist dazu kommen, der das Fleisch reinige, sonst werden die fleischlichen Kinder Abrahams sowohl als die Heiden verworfen und getödtet. Nun folgt im Text:

\begin{quote}
\textbf{V. 8. Da sprach Juda zu Onan: Lege dich zu deines Bruders Weib, und nimm sie zur Ehe, daß du deinem Bruder Samen erweckest.}
\end{quote}

48. Das hebräische Wort, jabam, so an diesem Orte gesetzt, steht sonst nirgends mehr, denn 5 Mos. 25, 5., da geboten wird, daß ein Bruder seines verstorbenen Bruders Weib soll zur Ehe nehmen etc. Dies ist aber nicht Moses Gesetz gewesen, sondern der Vater vor dem Gesetze Moses; daraus denn erhellt, daß Mose viel Ceremonien und Satzungen wird gesammelt haben, so zuvor die heiligen Väter aufgerichtet oder gestiftet gehabt. Von den Gesetzen ist dies eines gewesen, so uns bewährt oder erhalten und von den hohen Patriarchen immer von einem auf den andern gekommen. Es ist aber wohlgleich über die Maßen ein schwer und verdrießlich Gesetz gewesen; denn es zwingt dich, daß du auch wider deinen Willen deines Bruders Weib, wo er ohne Erben verstorben, zur Ehe nehmen sollst, wiewohl du weder Lust noch Liebe zu ihr hast; so es doch für unmöglich gehalten wird, die mit fleischlicher und ehelicher Liebe zu lieben, welche du selbst weder ermählt noch begehrt hast, es möchte denn traun aus heftiger Brunft der Unzucht geschehen. Wie aber, wenn sie unfruchtbar wäre, wie die gewesen, davon Matth. 22, 24. ff. Meldung geschiehet? Derselben, sage ich, ist das ein sehr hartes und schweres Gesetz gewesen. Aber in diesen Gesetzen der äußerlichen Ceremonien muß man auch auf Christum sehen, welchen man darin suchen muß.

49. Wenn aber im 5. Buch Mose Cap. 25, 3. 5. geboten wird, daß ein Bruder seinem verstorbenen Bruder Samen erwecken sollte, wird die Polygamie mit eingeschlossen, das ist, daß Einer hat mögen mehr denn Ein Weib haben, und wird dasselbe auch von Mose bestätigt. Und ist dies auch die Ursache gewesen, warum Salomo so viel Weiber gehabt, deren etliche ihm auch gewißlich durch Geblüt oder Schwägerschaft werden verwandt gewesen sein, und so etwa ein armes Weiblein in dem Geschlecht gewesen, dieselbe wird sich zum Könige, als zu ihrem Blutsfreunde oder Schwager, gehalten haben. Da haben sie sich zu dem Retter geschlagen, der wird viele ernährt haben, und sind doch auch ohne Zweifel etliche seine Kebsweiber gewesen neben der Königin um dieses Gesetzes willen: welches wahrlich ein verdrießlich und unleidlich Ding gewesen, nämlich, mit so vielen Ehefrauen oder Kebsweibern beladen oder beschweret sein, zu welchen du doch weder Lust noch Liebe hättest.

50. So heißt nun das Wort jabam ein solch ehelich Bündniß machen, da ein Blutsfreund gezwungen wird vermöge des Gesetzes, die Wittwe, so sein verstorbener Bruder ohne Kinder hinterlassen, zur Ehe zu nehmen; oder wo kein Bruder mehr vorhanden gewesen, sind laut desselben Gesetzes die andern Blutsfreunde schuldig gewesen, solches zu thun. Wer aber dies Gesetz erstlich gestiftet, oder woher es gekommen sei, weiß ich nicht. Juda gebietet seinem Sohne Onan, er solle seines Bruders Weib zur Ehe nehmen; welches er ungern gethan: und es scheint, als sei er auch derhalben mit dem Weiblein desto unfreundlicher umgegangen, oder da er gezwungen worden, daß er sie ja hat nehmen müssen, habe er sich doch geweigert, sie zu beschlafen. Denn er hat dies verdrießliche Gesetz nicht ertragen können.

51. Darum ist auch Thamar, so doch von Natur sehr fruchtbar gewesen, von dieser Zwey, Ger und Onan, nicht schwanger geworden, und sind vielleicht auch noch wohl andere Sünden dazugekommen. Wie denn solches beider Ende anzeigt, daß ihre Bosheit sehr groß gewesen; denn der Herr hat sie beide in einem Jahre getödtet. Wir wollten denn traun sagen, daß der Eine im Anfang des andern Jahres getödtet sei, da er zwölf Jahr alt gewesen, und vom Vater gezwungen ward, Thamar zur Ehe zu nehmen.

52. Es ist wahrlich eine wunderbare Freveligkeit gewesen, so die Weiber damit gehabt, daß sie nach einander haben freien mögen alle Brüder oder Verwandten ihrer vorigen Ehemänner, und wo das Weib unfruchtbar gewesen, ist es wahrlich eine sehr unangenehme und schwere Last gewesen. Und das heißt recht, die böse Lust zwingen, welche gezwungen wird, da sie nicht will, und wenn sie also gezwungen ist, flieht sie doch und wendet sich ab von denen, so ihr ohne ihren Willen aufgedrungen worden. Derhalben halte ich es dafür, daß dies Gesetz nur von den Frommen und Gottfürchtigen wird gehalten worden sein; aber die Bösen und Gottlosen werden es gar verachtet und übertreten haben. Es steht in der Historie Ruth Cap. 4, 9. 1. ff. geschrieben, wie der, so ihr nächster Verwandter war, da er sie wegen der Verwandtschaft zum Weibe nehmen sollte, habe lieber den Acker oder das Stück Feld wollen fahren lassen, und den Schuh ausziehen, ein Barfüßermönch werden, und dem Andern geben (welches die Strafe war derer, so dasselbe Gesetz übertreten), ehe er solche Dienstbarkeit hat tragen wollen. Das würde wahrlich eine erschreckliche Zucht sein wider die bösen Lüste, wenn dies Gesetz uns auch binden sollte. Denn wir erfahren, wie viel Mühe und Arbeit sonst die Liebe der Eheleute, so anfänglich freiwillig und sehr hitzig gewesen, pflegt mit sich zu bringen, da oftmals der Satan die Herzen irre macht und sie zu vielem unbilligem Haß, Hader und Zank reizt.

53. So sagt nun Juda zu Onan: „Lege dich zu deines Bruders Weib, und nimm sie zur Ehe“ (denn so viel bedeutet das hebräische Wort, nämlich, nicht allein bloß ehelich werden, sondern sich zu seiner Blutsfreundin thun und sie zum Weibe nehmen), „daß du deinem Bruder Samen erweckest.“ Das ist die endliche Meinung gewesen dieses Gesetzes, und haben die Brüder oder sonst die Verwandten, so noch am Leben gewesen, diese Sorge tragen müssen, damit des verstorbenen Bruders Name im Lande nicht vertilgt oder ausgelöscht würde. Denn im Volke Israel hat Jeder müssen Kinder nachlassen. Nun folget aber im Texte eine scheußliche und schändliche That des bösen Buben Onan.

\begin{quote}
\textbf{V. 9. 10. Aber da Onan wußte, daß der Same nicht sein eigen sein sollte, wenn er sich zu seines Bruders Weib legte, ließ ers auf die Erde fallen, und verderbete es, auf daß er seinem Bruder nicht Samen gäbe. Da gefiel dem Herrn übel, daß er that, und tödtete ihn auch.}
\end{quote}

54. Das muß ein muthwilliger verzweifelter Bube gewesen sein; denn dies ist ja eine schändliche Sünde, noch viel greulicher, denn eine Blutschande oder Ehebruch. Wir heißen es eine Sünde der Weichlinge, ja, auch wohl eine Sodomitische Sünde. Denn er hat sich zu ihr gelegt, sie beschlafen, und da es zum Samen kommt, hat er denselben auf die Erde geschüttet und verderbet, damit das Weib nicht schwanger würde. Wahrlich, zu der Zeit soll man ja der Ordnung der Natur, so von Gott zur Geburt geschaffen ist, folgen. Darum ist es, sage ich, die allerschändlichste Sünde gewesen, den Samen locken und das Weib damit reizen, und sie doch bald im selben Augenblick desselben berauben. Er ist vom bösen Neid und Abgunst gar entbrannt gewesen, darum hat er nicht leiden wollen, daß man ihn zwingen sollte, diese unleidliche Dienstbarkeit zu tragen. Derhalben ist er billig von Gott getödtet worden. Es ist eine böse That, darum hat ihn Gott gestraft.

55. Und das ist es auch, das ich gesagt habe, daß die sonderlich fromm gewesen sein müssen, so dies Gesetz gehalten. Denn es ist ja eine große schwere Last, einem Andern dienen, Erben zu erwecken, und desselben Nachkommen zu erhalten, fremde Kinder zeugen, dieselben aufziehen und ernähren, ihnen das Erbe lassen, und das alles unter dem Namen des verstorbenen Bruders. Von solcher Liebe weiß die Welt gar nichts. Es ist ein verdrießlich Ding, der Waisen Vormünder zu sein, dieselben regieren und versorgen; welches zu unsern Zeiten nach äußerlichen Rechten zu geschehen pflegt. Hilf Gott! wie viel Klage und wie große Untreue findet man dabei in aller Welt! Denn es ist ganz schwer und ein trefflich Werk der Liebe, daselbst Treue und Fleiß anzuwenden, daß man Andern ihre Güter bewahre und sie dabei erhalten möge. Derhalben begreift dies Gesetz große brünstige Liebe in sich, welche dieser Schalk nicht hat üben oder beweisen wollen; sondern will sich lieber selbst mit so schändlicher scheußlicher Sünde beflecken, denn daß er seinem Bruder Samen erwecken wollte. Nun hat aber Juda auch gesündigt, darum wird bald eine greuliche Strafe folgen.

\quellseite{309}

\section[Zweiter Theil]{Zweiter Theil. Von Thamar, und von der Blutschande Juda's mit Thamar}

\subsection{I}

\begin{quote}
\textbf{V. 11. Da sprach Juda zu seiner Schnur Thamar: Bleibe eine Wittwe in deines Vaters Hause, bis mein Sohn Sela groß wird. Denn er gedachte, vielleicht möcht er auch sterben, wie seine Brüder.}
\end{quote}

56. An diesem Orte disputiren die Schreiber, wer doch der Thamar Vater möge gewesen sein, und die Juden haben nach ihrer Weise Lust zu Possen und Fabeln, und dichten, sie sei des Priesters Senn Tochter gewesen, und wollen es daraus schließen, daß Juda hernach befohlen, sie hervor zu bringen, daß sie um des Ehebruchs willen verbrannt würde, darum daß man der Priester Töchter, wenn sie gehurt haben, hat pflegen zu verbrennen.

57. Aber diese Possen lasse ich fahren; denn sie können ebenso leichtlich verachtet werden, als sie leichtlich sind erdichtet worden. Denn sie ist einmal fünfzig Jahren alt gewesen, Ger aber und Onan sind zwölf Jahre: denselben wird Juda nicht zugelassen haben, daß sie Thamar, so fünfzigjährig gewesen, sollten gefreit haben. Darum ist sie eine Cananitin gewesen wie Rahab, und wie Ruth eine Moabitin von heidnischem Geblüte geboren war.

58. Nun ist aber dies auch ein sehr schwer und scharf Gesetz gewesen, daß wenn die Töchter nicht fruchtbar waren oder ihre Männer starben, sie leicht verstoßen und ihren Eltern wieder zugestellt wurden. Darum war es damals ein schwer Ding, die Töchter ehelich zu berathen. Aber andere Völker haben auch andere Gesetze. Uns bindet dies oder andere dergleichen Gesetze gar nicht. Juda aber läßt sich dünken, er habe seiner Schwiegertochter den dritten Sohn billig versagt; denn sonst würde er ihr denselben nicht also zu geben verweigert haben: und sündigt doch gleichwohl daran, daß er damit so streng verfahren. Und ist dies eine sehr kalte und lose Entscheidung gewesen, daß er sagte: „Vielleicht möchte der dritte Sohn auch sterben“ und ich müßte also ohne Kinder bleiben, sintemal mir mein Weib auch gestorben. Es wird aber aus dem folgenden Text noch offenbar werden, daß er seinem Sohn Sela ein ander Weib gegeben habe, und Thamar also verachtet und sitzen lassen. Denn er hat sie wiederum in ihres Vaters Haus gesandt, und dermaßen verachtet und verlassen. Er hat dies Gesetz nicht wollen halten, darin geboten war, daß ein Bruder seines verstorbenen Bruders Weib zur Ehe sollte nehmen etc. Denn Mose wird V. 14. sagen: Sela sei groß geworden, und Thamar sei ihm doch nicht vertraut worden.

59. Derhalben ist diese Thamar ein armes unseliges Weib gewesen, und scheint, daß sonst auch überall in dem jüdischen Volk das weibliche Geschlecht fast dienstbar und verachtet gewesen sei, und daß sie aus hochbringender Noth sich mit ihrer Arbeit, als, mit Spinnen und Weben, haben nähren müssen. Denn also scheint es mit Thamar auch geschehen zu sein, daß sie sich daheim bei ihrem Vater nach weiblicher Art, als, mit Spinnen und Waschen, ganz schwerlich und kümmerlich erhalten habe.

60. Und droben, Cap. 31., haben wir gehört, daß Laban auch von der Arbeit seiner Töchter Rachel und Lea an seiner Nahrung zugenommen und reich geworden sei. Desgleichen steht auch Jes. 4, 1.: „Sieben Weiber werden zu der Zeit einen Mann ergreifen, und sprechen: Wir wollen uns selbst nähren und kleiden; laß uns nur nach deinem Namen heißen.“ Und das ist auch ein hartes Gesetz gewesen. Es sind so treffliche herrliche Fräulein. Bei den Heiden aber und unserer Gewohnheit nach soll dies Gesetz nicht gehalten, auch nicht zum Exempel angeführt werden, daß man demselben wollte nachfolgen; sondern unsern Rechten und Gesetzen, und jetziger Obrigkeit sollen wir folgen und derselben gehorchen, sofern die weltlichen Ordnungen an sich selbst nicht ungöttlich und dem Wort Gottes nicht zuwider sind. Wir sollen nicht sehen, was die Syrer oder Philister für Statuten und Ordnungen gehabt; denn ein jeglich Land hält sich nach seinen eigenen Rechten und Sitten, wie man im deutschen Sprüchwort sagt: So manches Land, so manche Sitten.

Also ging Thamar hin, und blieb in ihres Vaters Hause.

61. Dies alles ist schier innerhalb dreier Jahre geschehen. Ger und Onan sind fast in Einem Jahre, darin sie waren ehelich geworden, umgekommen. Auf den Dritten hat Thamar gewartet, bis er würde groß werden, und ist derweil eine Wittwe geblieben. Juda aber hat damit durch die Finger gesehen, und ist also vom geschriebenen Gesetz abgetreten und hat seinem Sohn eine Andere getrauet (denn also mögen wir hieraus schließen); daß also Thamar gar keine Hoffnung mehr hat haben können. Und wiewohl sie gut Recht gehabt, den dritten Sohn zu fordern nach dem Gesetz, wird sie doch gedacht haben: Ich bin eine Cananiterin, bin verachtet, verworfen und verdammt. Was soll ich thun?

\begin{quote}
\textbf{V. 12. Da nun viel Tage verlaufen waren, starb des Suah Tochter, Juda Weib. Und nachdem Juda ausgeweinet hatte, ging er hinauf, seine Schafe zu scheren, gen Thimnath, mit seinem Hirten Hira, von Odollam.}
\end{quote}

62. Es ist oben gesagt, daß die Zeit dieser vielen Tage nicht mehr denn Ein Jahr zusammen macht; denn sonst würde die andere Zeit nicht genug sein, die Kinder und Kindeskinder Judas zu zeugen, welcher, da er seine Schwiegertochter geschwächt, sieben und zwanzig Jahr alt gewesen, und im acht und zwanzigsten Jahre wird Thamar gebären; welcher Thamar Söhne weiter auch Kinder zeugen werden in den Jahren, so noch übrig sind, bis auf das zwei und vierzigste Jahr. So sei das nun Ein Jahr dieser vielen Tage, davon im Text gemeldet wird, daß darin Judas Weib gestorben sei, nachdem seine zwei Söhne getödtet waren. Wo du aber diese Tage zählen oder rechnen wirst, nicht von der Zeit an, da Thamar verworfen, sondern da Juda Hochzeit gehalten, so werden dergleichen Jahre daraus; dasselbe Jahr aber wird der Mitwissenschaft Thamars zugeschrieben.

63. Nun ist aber Juda auch ein Wittwer gewesen, wiewohl nicht länger als ein halb Jahr, sintemal sein Weib erst nach der Verstoßung Thamars und nach dem Tode seiner zwei Söhne Ger und Onan gestorben ist. Es ist auch im Gesetz also verordnet gewesen, daß sie etliche Monden haben müssen Leid tragen. Und dies ist eine sehr löbliche Gewohnheit bei allen Völkern, die Verstorbenen, sonderlich wo es Verwandte oder Eltern sind, zu beweinen. Mose gedenkt nicht, daß Juda habe Hochzeit gehalten, so er doch noch vollkommenen und vermöglichen Alters gewesen, nämlich von sieben und zwanzig Jahren. Darum hat er sich des Ehestandes enthalten, aber gleichwohl nicht ohne Weiber gelebt. Denn also hat der Väter Leben müssen beschrieben werden, daß sie uns allenthalben gleich seien nach dem Fleisch, nach der Sünde und Tode, auf daß die unermessliche und unaussprechliche Liebe und Barmherzigkeit Gottes, damit er uns geliebt, gerühmt werden möge.

\subsection{II}

64. So sagt nun Mose: Juda sei fröhlich, und mit aufgerichtetem Haupte hinaufgegangen, seine Schafe zu scheren. Denn dieselbigen Tage sind jährlich als Feiertage gehalten worden, an welchen sie auch etwas herrlicher und fröhlicher gelebt haben mit Essen und Trinken, wie 1 Sam. 25, 2. ff. steht von Nabal, daß er seinen Schäffcheren ein Mahl zugerichtet, dahin auch David seine Jünglinge gesandt etc. Wie bei uns in der Zeit der Ernte auch geschieht. Denn denen, so von der Arbeit und Hitze müde geworden, thut Essen und Trinken noth. Dies sind aber Reden von guten äußerlichen und bürgerlichen Sitten.

65. Das Vieh aber, wie droben, Cap. 37, 13., gesagt, ist nicht zu Sichem bei dem Vater gewesen, wiewohl er denselben bisweilen wird besucht haben, sondern zu Abdollam, dabei allernächst Thimnath gelegen. Achsib aber ist im Stamme Juda gewesen, im Thal des Landes Juda. Denn darum wird auch von Simson gesagt, daß er sei hinab gezogen. Und Kyra legt es also aus, daß es von Jerusalem bis gen Thimnath herab gegangen sei und wiederum vom Meer gen Jerusalem soll es bergauf gegangen sein.

\quellseite{310}

66. Das hebräische Wort, reeh, habe ich im Deutschen verdolmetscht, hirten, um welches Wortes willen die Hebräer heftig zanken und damit sie uns wohl plagen. Denn sie sagen, es soll ein Freund heißen, oder der unser Nächster, aber doch gleichwohl unser Blutsfreund nicht sei. Wie es denn im Evangelio recht gegeben und verdolmetscht ist, Matth. 22, 39.: „Du sollst deinen Nächsten lieben, als dich selbst.“ Hieronymus handelt davon fleißig und sagt, es heiße, Hirte, wenn man es mit dem Pünctlein schreibt, so die Hebräer Cholem nennen, also roeh. Wenn man aber sehe reeh, alsdann soll es ausgelegt werden, daß es einen Freund heiße. Und er hat es selbst im Lateinischen gegeben, pastor seu opilio, das ist, Hirt oder Schäfer. Ich will es den Hebräern befehlen. Es scheint, als sei es ursprünglich hergekommen von dem hebräischen Wort reeh, und es erinnert an diese Bedeutung, wenn es im Anfang dieses Capitels heißt, daß Juda bei einem Mann von Abdollam eingekehrt sei, welcher ein ehrlicher frommer Bürger gewesen und Juda beherbergt habe. Es scheint aber doch, daß Juda etwas höher gewesen sei als er. Denn Juda hat durch ihn der Thamar den Ziegenbock gesandt, wie durch seinen Knecht. Derhalben halte ich dafür, es sei nicht sein Freund oder Nachbar, sondern der etwas geringer gewesen denn er, welcher den Juda in Ehren gehalten. Aber ich will doch dieser Meinungen keine für so gar gewiß halten, will es lieber unentschieden lassen, wie Hieronymus auch thut, und den Hebräisten anheimstellen, darüber zu urteilen. Es schickt sich zwar beiderlei Auslegung wohl; wiewohl ich zu der einen mehr Lust habe, wenn man das Wort reeh also versteht, daß es ein Hirte heiße. Denn also ist das Leben der Patriarchen zum andern umhergezogen sind, wie ihre Väter auch gethan und wie zu unsern Zeiten die Schafmeister und Ochsentreiber thun. Also sind sie auch in der Irre umhergezogen, haben täglich neue Weide gesucht, haben Wiesen und Acker gemiethet, dazu auch andere Stätten, so ihnen zu ihrem Vieh und Heerden gelegen gewesen. Darum sagt David in 39. Psalm V. 13.: „Ich bin beide dein Pilgrim und dein Bürger, wie alle meine Väter.“ Und also sind Jakobs Söhne auch zerstreut gewesen hin und wieder durch mancherlei Orte und Geschlechter um ihrer Heirath und Schwägerischafft willen, so sie gemacht hatten; denn sie haben am Lande nichts Eigenes gehabt. Nun folgt im Text:

\begin{quote}
\textbf{V. 13. 14. Da ward der Thamar angesagt: Siehe, dein Schwäher gehet hinauf gen Thimnath, seine Schafe zu scheren. Da legte sie die Witwenkleider von sich, die sie trug, deckte sich mit einem Mantel, und verhüllte sich, und setzte sich vor die Thür heraus an dem Wege gen Thimnath; denn sie sahe, daß Sela war groß worden, und sie ward ihm nicht zum Weibe gegeben.}
\end{quote}

67. Juda hat droben Unrecht gethan, daß er seinem Sohn seine Schwiegertochter, Thamar, nicht hat wollen zum Weibe geben nach dem Gesetz, wie er selbst hernach bekennen wird. Nun aber kündigt sie desgleichen auch. Das hebräische Wort, zaiph, steht droben auch in der Historie von Rebekka, Cap. 24, V. 65., die aus Mesopotamien gekommen ist, und da sie den Isaak auf dem Felde hat gehen sehen, sagt der Text, daß sie den zaiph, das ist, den Mantel, genommen und sich damit verhüllt habe, Ehrbarkeit und Zucht damit anzuzeigen. Denn das sind dieses Volkes Sitten gewesen, daß ein Weib anders bekleidet gewesen als eine Jungfrau und eine Wittwe auch noch anders, welches ein Anzeichen ist, daß es seine züchtige Leute müssen gewesen sein, die wohl erzogen sind, ehe sie aus der Art gefallen und in allerlei Sünde gefallen sind.

68. Ich halte aber dafür, daß dies Wort bedeuten müsse den großen Schleier, so die Weiber gebraucht, damit sie nicht allein das Haar eingebunden, sondern auch das ganze Haupt bis auf die Achseln hinab haben pflegen zu bedecken: es ist ein langer Schleier oder groß Gehülle um den Kopf gewesen. Wir werden hernach hören, daß auch die Männer ihr eigen Gehülle gehabt, damit sie das ganze Haupt verhüllt. Wie in den Hochschulen die Studentenkappen gewesen, damit sie auch den ganzen Kopf und die Achseln bedeckt haben. Ein solch Gehülle, dafür halte ich es, ist auch bei den Chaldäern gebräuchlich gewesen; wie Hes. 23, 15. erzählt wird, daß die Fürsten der Chaldäer gekommen seien, mit spitzen, bunten Hüten auf ihren Köpfen etc. Es werden große weite Kappen gewesen sein. So haben nun beide Männer und Weiber ihren sonderlichen Schmuck gehabt, damit sie die Häupter bis auf die Achseln bedeckt. Und zu unsern Zeiten ist an etlichen Orten in Deutschland auch gebräuchlich, daß man etliche Gehülle oder Kappen trägt, damit auch beide Hals und Mund bedeckt wird, daß man nur die Augen sehen kann.

69. Dies ist aber nicht ein Schmuck gewesen für die Wittwen, sondern für andere Matronen, so nicht Trauerkleider, sondern fröhliche und Feierkleider haben tragen wollen. So hat nun Thamar ihre Trauerkleider abgelegt und den Mantel (auf Hebräisch zaiph) angezogen, und hat nicht allein das Haupt verhüllt mit einem für Matronen bestimmten und ehrlichen Schleier, sondern hat auch den ganzen Leib fein hübisch und ehrlich geschmückt. Ich höre, daß auch bei den Türken soll gebräuchlich sein, daß die Weiber beide das Haupt und den ganzen Leib also verhüllen. Und es ist eine Erzählung vorhanden Eines, der in der Türkei gefangen gewesen, darin er bezeugt, er habe nie gesehen, daß eine Hausmutter in bloßem Haupte oder Angesicht mit ihrem Ehemann geredet habe. Das griechische Wort heißt theristrum; dasselbe legen sie also aus, daß es ein Sommerkleid sei, so den Schnittern in der Ernte dienlich und bequem ist.

70. Da aber nun Thamar also geschmückt und geziert hervorgetreten ist, Juda zu reizen oder zu locken, sagt Mose, habe sie sich vor die Thür gesetzt (im Lateinischen steht: in ostio oculorum), das ist, auf die Wegscheide, welches eine sonderliche Figur oder Weise zu reden in der hebräischen Sprache ist und im Lateinischen gar nicht gebräuchlich, und lautet, als ob man auf Lateinisch wollte sagen: in janua vel foribus duorum oculorum. Denn es ist der Dual, nämlich, daß man es im Lateinischen also möchte ausdrücken: una via sinistri oculi, altera dextri, auf Deutsch: ein Weg des linken, der andere des rechten Auges. Darum daß auf der Wegscheide die Augen sich auf zwei Seiten, nämlich, zur Rechten und Linken, pflegen zu richten. Die Andern haben es im Lateinischen gegeben: in divisione fontium, bei der Theilung der Duellen. Hieronymus aber hat es im Lateinischen verdolmetscht: in bivio sive in compito, das ist, auf der Wegscheide, da Einer, so über Land reist, fleißig muß zusehen, welchen Weg er gehen wolle.

71. So hat nun Thamar auf der Wegscheide gesessen, aber also, daß sie sich doch am meisten geneigt und gegeben hat auf den Weg, so nach Thimnath gegangen. Dies ist aber ihr Vorhaben, daß sie ihren Schwiegervater hat betrügen wollen. Welches wahrlich vermögen und gefährlich genug gewesen, dazu auch nicht ohne große Sünde. Und ist derwegen Thamar gar ein thörichtes und freveles Weib gewesen, welche GOtt so trotziglich versucht hat.

72. Und es scheint, als habe sie nicht die Unzucht dazu getrieben, daß sie solches hat thun dürfen. Denn die Ursache wird im Texte klärlich ausgedrückt, nämlich, wie sie gesehen habe, daß sie verachtet sei worden, und daß sie den dritten Sohn Judas, den Sela, nicht habe freien mögen, daß ihr doch von Rechts wegen und nach dem Gesetze gebühret hätte. Denn dies Gesetz ist auf das strengste gehalten worden, nämlich, dem verstorbenen Bruder Erben zu erwecken und die Nachkommen zu erhalten, und daß solches hat müssen geschehen durch den Bruder oder sonst der nächsten Freunde einen. Ob aber die Heiden an dasselbe Gesetz auch gebunden gewesen, ist ungewiß. Die Väter haben ehrlich und sehr hart darüber gehalten; denn darum hat Juda so ernstlich seinem Sohn Onan geboten, er sollte sich zur Thamar legen, so seines verstorbenen Bruders hinterlassene Wittwe war, und sie zur Ehe nehmen.

73. Da aber nun Thamar sieht, daß ihr das fehlet, darauf sie gehofft, und daß sie den dritten Sohn nicht hat mögen zur Ehe haben, hat sie diesen listigen Anschlag gemacht, nicht als hätte sie die Hütze der Unzucht dazu gereizt oder getrieben (denn sonst hätte sie sich wohl können an einen Andern hängen), sondern daß sie ihr Recht hat wollen ausführen. Denn sie hat das
\quellseite{311}
in keinen Zweifel gestellt, daß sie die gerechteste Ursache gehabt zu zürnen und den dritten Sohn zum Ehemann zu fordern. Denn das Verlöbniß war nicht allein die praesent, wie sie das pflegen zu nennen, daß sie solches einander selbst gelobet, sondern es war auch die Bewilligung der Eltern dazu gekommen. Denn Gottes Gesetz gebietet, daß Sela dieser Thamar Bräutigam seyn soll vermöge göttliches Rechts, und ist solch Gebot viel stärker, denn wo die Eltern, oder auch die Bewilligung des Bräutigams und der Braut da gewesen wäre; denn sie werden beiderseits als von Gottes wegen dazu gezwungen. Darum hat Thamar eine herrliche Ursache, den Sela zur Ehe zu begehren; denn er war ihr Bräutigam, und war wegen des Gebots Gottes schuldig, daß er sie hätte sollen zur Ehe nehmen.

74. Dies Gesetz, sage ich, ist zwar sehr scharf und schwer, ist aber doch von Andern streng gehalten worden, und nach demselben Gesetz ist Sela der Ehemann, Thamar aber die Ehefrau, daß sie nicht haben mögen geschieden werden, nach der Regel Matth. 19, 6.: „Was nun Gott zusammen gefüget hat, das soll der Mensch nicht scheiden.“ An dasselbige Gebot hat sie sich gehalten und sich dessen in ihrem Herzen auch gerühmt. Dieser Sela, wird sie gesagt haben, ist mein Ehemann, ich bin sein Eheweib und bin die rechte Hausmutter in diesem Hause, ich gehöre in das Haus; warum sollte ich denn leiden, daß er mir sollte genommen oder entzogen werden? Juda wird das Spiel versehen haben, sündigt wider dies Gesetz, und will, das also zusammen gefüget ist, scheiden, gibt meinem Sohn eine Andere, verachtet und verwirft diese.

75. Darum tobt und zürnt Thamar also nicht ohne Ursache. Wiewohl sie aber hürisch handelt, hat sie doch in ihrem Herzen mehr Zorn als Lustseuche gefühlet; denn sie denkt also: Siehe, ich bin nun verachtet und verworfen ohne einige Ursache, und werde diese Schmach für und für tragen müssen, daß ich ja nach dem Gesetz die rechte Ehefrau gewesen, aber doch verstoßen worden sei von meinem Schwiegervater, ohne alle meine Schuld und da ich solches mit keiner Sünde verwirkt habe. Und also werde ich nun beraubet werden der Frucht meines Leibes, und ist dazu keine Hoffnung, daß ich sonst könnte ehelich berathen werden, sondern werde mein ganzes Leben in erbärmlichem Wehklagen zu bringen müssen, daß ich also für und für verlassen und verachtet sein werde. Darum versucht sie es, so hoch sie kann, und versucht, ob sie Juda betrügen möchte. Und ich will glauben, wo sie den Sela selbst auch hätte können betrügen, sie würde das auch mit allem Fleiß gethan haben. Das Gesetz macht ihr Muth, wiewohl sie sich doch billig des Geblüts halben davor sollte gescheut haben; denn es ist eine Blutschande: aber aus dem Zorn und Ungeduld, daß sie solche Schmach und Hohn hat sollen tragen, hat sie sich unterstanden, eine solche scheußliche That zu begehen, welches ihr doch mit nichten geziemet hat.

76. Es werden aber solche schwere Sünden der Heiligen so vielfältig und mit so vielen Worten erzählt uns zum Trost, daß wir wissen sollen, daß die Patriarchen und heiligen Matronen uns gleich gewesen sind, welche bisweilen große hohe Dinge ausgerichtet, die wir nicht verstehen und denen wir nicht nachfolgen können; wie derselben etliche erzählt werden Hebr. 11, 7. ff. Bisweilen aber haben sie auch närrische und lästerliche Dinge, ja die greulichsten Sünden begangen, damit Gott allein gepriesen und gerühmt werde in unsern Werken, beide guten und bösen: in guten, die er selbst in uns gethan; in bösen aber, die er uns aus Gnaden und Barmherzigkeit vergeben hat.

77. Und sind Juda und Thamar große Sünder geworden, nicht daß sie Gott solches geheißen hätte, sondern er hat es also geschehen und darnach schreiben lassen, hält es auch seiner Kirche also vor zum Exempel und zum Trost. Thamar kann sich zwar nicht entschuldigen, hat sich aber doch die Ungeduld und die Schmach, daß sie also einsam und unfruchtbar bleiben sollte, überwinden lassen. Denn sie ist ganz entbrannt gewesen von großem Verlangen und Begierde, Kinder und Erben zu zeugen, und will in dem Hause eine Mutter werden, darin ihr das Gesetz das Mutterrecht zugeeignet und gegeben hat. Sie will es haben, sollte sie es auch vom Schwiegervater selbst kriegen.

78. Solches ist eine weibliche Ungeduld, damit sie sehr sündigt wider die Gesetze, darinn allen Menschen geboten wird, züchtig und schamhaftig zu sein, und kann nicht geleugnet werden, daß diese Sünde nicht sollte eine Blutschande gewesen sein. Denn da sie nicht hat können eine Mutter werden vom Sohne oder nächsten Blutsfreunde, reizt und bewegt sie den Schwiegervater selbst dazu, daß er sie beschlafen hat; welches unerhört und unzüchtig genug ist.

79. Also bringt eine Sünde die andere: die Sünde Judas bringt mit sich die Sünde und schändliche Blutschande der Thamar, und geschieht eben dasselbe, so man zu sagen pflegt: Nulla calamitas sola: Es kommt kein Unglück allein. Denn das arme Weib läßt sich die Verachtung und Schmach überwinden, dieweil sie sieht, daß Sela, wie Mose sagt, war groß geworden, und sie verstoßen wird, daß sie ihm nicht zum Weibe gegeben worden ist. Ich halte aber dafür, Juda werde dies heimlich ausgerichtet haben, und daß Thamar nicht öffentlich vor Gericht verworfen und dem Sela eine Andere gegeben sei, und aber sie gleichwohl solches also vermuthet und gemerkt habe; denn es war kaum ein Jahr vergangen nach dem Tode des Weibes Judas: wie denn der Text nicht von Jahren, sondern nur von Tagen sagt; wiewohl sonst die Schrift etwa Jahre zu verstehen, da sie Tage nennt. Da nun ihr voriger Ehemann, der andere Sohn Judas, gestorben, ist Sela bei elf Jahren gewesen, und Thamar hätte noch ein Jahr auf ihn warten können, daß er zum Ehestande alt genug geworden wäre. Aber sie sieht, daß gar keine Hoffnung vorhanden, daß sie ihn bekommen möchte; darum wird sie nun getrieben und bewogen, diese große Sünde zu begehen wegen der Schmach und daß sie vermöge göttliches Rechts hätte sollen in dem Hause die Hausmutter sein. Da denkt sie: Ich will auch etwas thun, das soll gethan heißen.

\begin{quote}
\textbf{V. 15. Da sie nun Juda sahe, meynete er, es wäre eine Hure, denn sie hatte ihr Angesicht bedeckt.}
\end{quote}

80. An diesem Orte wird gefragt: Wie doch das zugegangen sei, daß Juda Thamar, seine Schwiegertochter, nicht gekannt habe, die er doch ja zum wenigsten an den Augen hätte kennen sollen, welche offen und unverhüllt waren, desgleichen an der Stimme; denn er hat mit ihr geredet, und ist ihm zuvor auch wohl bekannt und zween Jahre seine Hausgenossin gewesen, da sie seiner zwei Söhne, Ger und Onan, Weib war, da er ja täglich ihre Sprache gehört und dieselbe von andern leichtlich hat unterscheiden können? Denn es ist so ein großer Unterschied der Stimmen, daß auch die Vögelein ihre Mutter an der Stimme erkennen. Darum ist es ja ein wunderlich Ding, daß durch das alles Juda nicht aufmerksam geworden ist. Was ist denn nun die Ursache gewesen? Antwort: Wo man auf etwas sonderlich verstürzt ist, da hat man weder Sinne noch Gedanken. Denn wer auf einerlei Ding so gar ernstlich beschlossen ist, der sieht, oder hört gar nicht, was ihm auch sonst vor die Augen oder Ohren kommt; wie man das in der täglichen Erfahrung dieses Lebens sehen mag, und man gemeiniglich also sagt: Pluribus intentus minor est ad singula sensus, das ist: Wenn Einer auf vielerlei Dinge beschlossen ist, kann er eines jeden insonderheit nicht so gar wohl warten. Wenn ich eines Worte höre und fleißig darauf merke, so ist mein Gemüth in meinen Ohren, und kann ich mich nicht so sehr damit bemühen, was mir sonst vor Augen kommt. Also auch, wenn mein Gemüth durch einen der andern äußerlichen Sinne, welcher es auch sein mag, in Anspruch genommen ist, ist seine Wirkung in den andern schwächer. Derhalben ist Juda ganz der Meinung, als sähe dasselbige eine Hure, und dieweil ihm das allein im Sinne und vor Augen ist, hat er weder auf die Stimme noch Augen des Weibes Achtung geben können. Wo aber jemand an solcher Antwort auf diese Frage nicht Genüge hat, der mag es dafür halten, es sei ein Wunderwerk gewesen, oder aber, daß Juda entweder von Gott, oder aber vom Teufel verblendet gewesen, daß er weder gesehen noch gehört hat.

\begin{quote}
\textbf{V. 16—18. Und machte sich zu ihr am Wege, und sprach: Lieber, laß mich bei dir liegen. Denn er wußte nicht, daß sie seine Schnur wäre. Sie antwortete: Was willst du mir geben, daß du bei mir liegest? Er sprach: Ich will dir einen Ziegenbock von der Heerde senden. Sie antwortete: So gib mir ein Pfand, bis daß du mir's sendest. Er sprach: Was willst du für ein Pfand, daß ich dir gebe? Sie antwortete: Deinen Ring, und deine Schnur, und deinen Stab, den du in den Händen hast. Da gab ers ihr, und lag bei ihr; und sie ward von ihm schwanger.}
\end{quote}

81. Nach dieser Erzählung scheint es, daß Juda der Unzucht und Hurerei halben hat mögen beargwohnt werden, als ob er zuvor auch mit dem Laster beladen gewesen. Denn wie ist es möglich, daß Thamar in ihrem Herzen sich selbst diese Verstößung sollte gemacht haben, daß Juda wohl würde zu betrügen sein unter dem Scheine, als ob sie eine Hure wäre, wo sie nicht gemußt, daß er damit pflege umzugehen und sich ungebührlicher Vermischung beflissen hätte?

\quellseite{312}

82. Und warum ist Juda nicht stracks hingegangen zu der Heerde, dieweil er sich in seinem Herzen vorgenommen, allein dies einige Geschäft auszurichten, nämlich, die Schafe zu scheren, sondern hat dasselbige liegen lassen, und ist plötzlich, da er die Hure gesehen, mit böser Lust entzündet worden? Wie kommt er so plötzlich darauf? Wir aber wollen solchen Argwohn nicht haben, wollen lieber also schließen: Thamar habe den Juda für einen guten frommen Mann gehalten, und nur versuchen wollen, ob sie auf die Weise ihr gebührliches Recht erhalten möchte, da sie das sonst zum öftern vergeblich versucht hatte. Denn es scheint, als habe sie mancherlei Anschläge gemacht, und etlichemal angehalten durch den Vater, durch die Schwäger, die Söhne Jakobs, durch ihre Blutsfreunde, ob sie mit gutem Willen von ihrem Schwiegervater etwas möchte erlangen; es ist aber alles vergeblich gewesen. Derhalben greift sie nun die Sache auf eine andere Weise an, nicht daß sie sich im Verdacht der Unzucht gehabt hätte, sondern daß sie möchte ihren Willen schaffen, es geschehe gleich wie es wollte.

83. Wo aber jemand dünkt, daß dieser Argwohn billig auf Juda gelegen, will ich ihn auch nicht entschuldigen. Denn es ist gleichwohl nicht ein geringes Anzeichen der Leichtfertigkeit, so gar bald solche schändliche Brunst fühlen, als ob er dazu von Natur und seiner Lebensart halben sehr geneigt gewesen. Es sei aber nun gewesen was es wolle, so hat Thamar doch gedacht, sie wollte auf das äußerste alles versuchen und ansehen, was sie nur könnte, dieweil ihr sonst die andern Anschläge alle gefehlet. Sie hat aber doch sich selbst nicht vertrösten können, wie dieser Anschlag geraten oder was für einen Ausgang derselbe gewinnen würde, hat sich auch nicht auf göttliche Regierung verlassen können. Denn sie hat GOtt nicht können mit Hülfe anrufen, wie Juda auch nicht gebetet hat; denn der Heilige Geist reizt oder treibt niemand zur Hurerei oder Blutschande: und wo man in Sünden liegt, da kann kein Gebet recht geschehen. Darum schließen wir, daß sie beide Sünder und schuldig gewesen sind; wiewohl ihnen die Sünde hernach ist nachgelassen und vergeben worden.

84. Juda aber bittet, daß er möge zu der Thamar eingehen, wie die Schrift zu reden pflegt, das ist, bei ihr liegen, und wird nicht hinzu gesetzt, an welchem Orte sie die Blutschande begangen, dieweil Mose sagt, wie sie öffentlich am Wege gesessen, daß sie die, so vorüber gegangen sind, gesehen haben. Ich halte nicht dafür, daß sie nach der Weise der Philosophen, so Cyniker geheißen, öffentlich vor den Leuten bei einander gelegen, sondern werden vielleicht etwa in ein Häuslein oder Höhle, oder aber in den nächsten Wald gewichen sein. Und daselbst ist sie schwanger geworden durch die schändliche Blutschande, und ist das Fleisch aus den Lenden Judas gekommen, von welchem Christus sollte geboren werden, und ist also immer, mit Sünden befleckt, fortgepflanzt worden bis auf die Empfängnis Christi.

85. So geht unser HErr GOtt mit unserm Heiland um. Er läßt zu, daß er in solcher schändlichen Blutschande empfangen wird, auf daß er wahrhaftiges Fleisch an sich nehmen möchte, wie unser Fleisch auch in Sünden herausgefüttert, empfangen und ernährt ist. Darnach aber, da die Zeit gekommen, daß unser Fleisch im jungfräulichen Leibe seiner Mutter hat sollen annehmen, ist es durch die Kraft des Heiligen Geistes gereinigt und geheiligt worden, wie die Worte des Engels lauten, Luc. 1, 35., da er zu Maria sagte: „Der Heilige Geist wird über dich kommen, und die Kraft des Höchsten wird dich überschatten“ etc.; und ist doch in Wahrheit ein beflecktes Fleisch gewesen von Juda und Thamar.

86. Derhalben, sage ich, ist dies alles um Christi willen beschrieben worden, auf daß man wüßte, daß er wahrlich von sündlichem Fleisch würde geboren werden, aber doch ohne Sünde. Darum, was David von sich sagt im 51. Psalm V. 7.: „Siehe, ich bin in Sünden empfangen“, dasselbe wird auch recht vom Fleisch Christi gesagt, daß es noch nicht angenommen, auch noch nicht gereinigt war, sondern wie es im Leibe Thamars gewesen: solches hat er aber darnach angenommen, da es gereinigt war, auf daß er die Strafe der Sünden für uns an seinem Leibe tragen möchte.

87. Mose erzählt aber auch weiter, was Juda Thamar für ein Pfand gegeben, bis er ihr den versprochenen Ziegenbock senden würde, nämlich, einen Petschiering, seine Schnur und Stab. Das hebräische Wort pathil heißt ein Faden. Im Plural pethilecha, wie es an diesem Orte gesetzt ist, panniculus tuus contortus ex filis factus, ein geflochten, gewirkt Tuch, und ist eine Kleidung gewesen, so die Männer gebraucht, das Haupt damit zu bedecken. Es ist kein Trauerkleid gewesen. Denn Mose hat zuvor gesagt, wie Juda ausgetrauert und sich über das Absterben seines Weibes habe trösten lassen. Derhalben hat er sich nun wiederum herrlich gekleidet. Und ich halte dafür, es werde ein geflochten leinen Tuch gewesen sein, damit er sein Haupt bedeckt, nach gemeinem Gebrauch derselben Leute. Denn die Hebräer deuten es also, daß es ein doppelt und geflochten Tuch sein soll, wie Luc. 16, 19. vom reichen Manne gesagt wird, daß er sich mit Purpur und köstlicher Leinwand bekleidet: welches eigentlich eine Kleidung des Hauptes und Halses gewesen. Und jetzt brauchen die Türken auch solche Kleidung. Dergleichen steht auch Ap. Gesch. 19, 12. von Paulo geschrieben, daß sie auch von seiner Haut die Schweißtüchlein und Koller über die Kranken gehalten, daß sie gesund worden. Also begehrt Thamar, daß Juda ihr das geflochtene leinene Tuch geben wolle, auf Hebräisch pethilecha, welches zugleich das Haupt und den Hals umgab, wie sie in Griechenland auch beiderlei bedeckt. Wir haben bei uns Hüte oder Barette, deren wir gebrauchen. Dieser Ort wird nicht weit von Thimnath gewesen sein, darum hat er gedacht, er würde sein Pfand bald wieder fordern können, als, den Ring, das Tuch, damit er das Haupt bedeckt, und den Stab, so insonderheit den Männern zugehören. Arrhabo ist ein griechisch Wort, welches hernach die Lateiner auch gebraucht, es ist aber ohne Zweifel von den Hebräern hergekommen, die auf ihre Sprache ein Pfand erabon heißen.

88. Neben dem mag man auch allhier sehen, was zu derselben Zeit das Volk für ein Leben geführt. Es ist sehr guter und schöner Friede im ganzen Lande gewesen, darum haben sie keine Waffe oder Schwert gebraucht, und hat dieser Patriarch weder Tasche, noch Geld oder Waffe getragen, da es sich doch gebührt hätte, daß ein solcher trefflicher Mann, der so ein groß Hausgesinde zu regieren gehabt, mit Zehrung sollte gerüstet gewesen sein, zum wenigsten mit einem oder mit zwei Gulden. Aber er trägt nichts mehr bei sich, denn seinen Ring, den Stab mit dem Kleide. O welch ein feiner seliger Friede ist das! Darum kann er denn Weiße feil Geld geben, sondern verspricht ihr einen Ziegenbock. Und es sind auch zu derselbigen Zeit die Kaufhändel und Hantierungen unserer Kaufmannschaft nicht gleich gewesen, sondern sie haben die Waaren vertauscht, das Vieh um Getreide, Brod und andere Nothdurft gegeben.

89. Thamar aber fragte nach keinem Gelde oder Pfande, sondern will Samen haben. Sie will eine Mutter sein, und macht sich auf, geht von dannen hinweg, wartet nicht auf den Ziegenbock, so ihr versprochen war, behält aber das Pfand; damit sie anzeigen wollen, wie es scheint, daß sie sich an dem Beischlaf noch nicht wollen genügen lassen. Und ich wundere mich, wie sie das hat wissen können, daß sie schwanger war; es ist ein wunderbarlich Weib gewesen, und es scheint, daß sie dieses noch mehr hat versuchen wollen, wo ihr dies gefehlet, darauf sie gehofft hatte. Aber sie hat es auf das alleräußerste versucht und ist ihr wohl geraten. Denn es war bei dem Volk eine große Ehre, wo die Weiber fruchtbar gewesen, und des Segens und Gesetzes Gottes willen. Bei uns wird es nicht groß geachtet, darum können wir auch solche Historien nicht verstehen. Es war bei ihnen nichts Schändlicher, denn das ehelose Leben und die Unfruchtbarkeit, darum daß das weibliche Geschlecht zur Geburt und Kinderzucht geschaffen ist. Wir haben wohl Lust und Liebe zur Fruchtbarkeit an dem Vieh, auf daß damit unser Gut und Nahrung gemehret werden möge, und das Vieh, so unfruchtbar ist und nicht Junge zeuget, wird bald zur Schlachtbank geführt; im Ehestande aber achtet man des nicht.

90. Die Heiden haben auch ein sehr fein Gesetz gemacht (sie mögen es von den Vätern genommen, oder aber darin der Natur gefolgt sein), wie solches aus der Komödie zu sehen ist: vermöge desselben wurden sie gezwungen, daß sie ihre nächsten Blutsfreundinnen haben ehelichen müssen, oder aber ihnen den Brautschatz oder Aussteuer geben.

91. Derhalben ist die Verachtung des Ehestandes zu strafen, sonderlich wo es also zugeht, daß sich die Leute mit unziemlicher Barmherzigkeit beflecken. Kannst du keusch leben, so ist es gut: wo nicht, so ist es viel besser und ehrlicher, ehelich werden denn unziemlicher Brunft nachzuhängen. Und ob du wohl dabei eben nicht groß Gut haben wirst, so wirst du doch also in Gottesfurcht züchtig und keusch leben können. So sind nun solche Sitten und Gesetze nütze und heilsam gewesen in gemeinen Regimenten, welchen es ja sehr nütze ist, daß die Städte voll Kinder und junges Volks sein mögen. Und wenn jetzt solche Ordnung und Gesetz bei uns im gemeinen Leben auch gebräuchlich und bestätigt wäre, wie daran das jüdische Volk gebunden gewesen, so würde auch der Schande und Unzucht weniger sein. Aber der sehr große Muthwille, so jetzt getrieben wird, geht mit Gewalt durch die ganze Welt und entstellt die menschliche Natur mit erschrecklichen Aergernissen. Die Leute, so zur Wohlust und guten Tagen geneigt sind, fürchten sich vor der Arbeit und Müheleidigkeit, so der Ehestand mitbringt; oder wenn sie sich schon in denselben Stand begeben, sehen sie doch vielmehr dahin, daß er ihnen eine Arznei wider die Unzucht sein möge, denn daß sie sollten begehren, Kinder darin zu zeugen. Und das sind die rechten Früchte des ehelosen Standes der Papisten.

\quellseite{313}

92. Die heiligen Väter haben Weiber genommen um Begierde und Liebe willen, Kinder zu zeugen, die sie auch mit großem Fleiß auferzogen, und haben ihre Familie mit großer Arbeit regieret. Das ist ein eigen Ding bei ihnen gewesen. Wir sind Pfaffen geworden, wir haben begehrt ein Pfaffenleben zu führen, haben uns zum Müßiggang, Völlerei und sonst anderer Wohlust dieses Lebens begeben, und haben die Sorge und Hausarbeit bei Seite gelegt. Die Pfafferei hängt uns noch sehr an, und wird dieser Schade und Unrath mit keinerley Hülffe oder Weise können erstattet werden, welchen das ehelose Leben der Papisten diesem heiligen Stande, so Gott selbst eingesetzt, gebracht hat. Denn da ist nicht mehr zu helfen, wenn es so weit kommt, daß man aus öffentlichen Lastern Gewohnheiten macht, nur frei ohne alle Scheu zu sündigen, wie Seneca sagt; wie es auch bei uns zugeht, die wir ja den rechten Preis des Ehestandes aus der heiligen Schrift suchen und verstehen sollten. Denn die Leute, so in Wohlust und Müßiggang verdorben sind, verachten beide Haus- und Weltregiment: darnach spotten sie noch dazu der Eheleute als solcher Leute, die gar in Wohlust und Unzucht ersoffen sein sollen.

93. Es ist aber der Ehestand mit nichten so ein zartes Leben, sondern ist voller unzähliger Sorge und Mühe, ja voll Schweißes und schwerer Arbeit; die Lust büßt sich da wohl. Aber da Mönch oder Pfaffe, der du ein ehe-loses Leben führst, bist recht im Müßiggang und Wohlust ersoffen; wie denn Einer recht und weislich soll gesagt haben: Die Mönche führen nicht ein geistlich, sondern ein wollüstig Leben; sie verachten und verkleinern die andern Stände, auf daß sie von fremden Gütern sich nähren, und das Regiment über die ganze Welt haben und in höchster Herrlichkeit leben mögen. Das ist eigentlich der höchste und der allergarste Müßiggang, nämlich, anderer Leute Arbeit genießen und der ganzen Welt Güter zu Hause scharren.

94. Dies soll man um der Sünde willen der Thamar desto fleißiger betrachten, daß wir nicht meinen, sie habe aus Unkeuschheit gesündigt. Sie hat zwar sehr schwer gesündigt; aber sie hat es gethan aus Ungeduld der Schmach, so auf ihr gelegen, welche ihr zu dieser Sünde Ursache gegeben. Denn sie ist darüber betrübt, daß ihr die höchste Ehre der Weiber soll genommen werden, nämlich, ein Eheweib zu sein und sonderlich dieses Juda Sohn, und daß sie also hat sollen beraubet werden alles Schmuckes, damit dazumal eine Hausmutter geziert gewesen. Das ist ein rechtsmäßiger Zorn gewesen, der auch schier wohl zu entschuldigen ist. Denn wir sehen, wie übel es uns verdrießt, wenn wir auch nur einen geringen Schaden leiden an unserer Ehre, Würde oder Gütern. Was hat aber nun das arme Weib noch übrig, da ihr die Herrlichkeit genommen ist, daß sie nicht kann eine Mutter sein, und des Hauses, der Kinder oder Nachkommen, und dazu des ganzen Hausregiments muß beraubet sein? Derhalben hat sie große Ursache zu zürnen und soll nicht so leichtlich verdammt werden; wiewohl sie doch auch nicht zu entschuldigen ist. Denn ein solch Weib muß für und für gefangen und gleichsam leibeigen sein. Sie war dazu von Gott berufen und verordnet, daß sie sollte desselben Hauses Ehre und Preis sein, wird aber nun von Juda verworfen. Und also haben sie beide schwerlich wider Gottes Gebot gesündigt.

\begin{quote}
\textbf{V. 19. Und sie machte sich auf, und ging hin, und legte den Mantel ab, und zog ihre Witwenkleider an.}
\end{quote}

95. Daraus ist es offenbar, daß Thamar muß ein fromm ehrlich Weib gewesen sein; denn sie wird nicht stolz, hat sich auch Andern nicht gemein gemacht, sucht nicht ungebührliche Wohlust, geht nur damit um, daß sie in dem Hause eine Mutter werden möge, darein sie von Gottes wegen gesetzt war. Darum legte sie nun die Kleidung ab, damit sie sich bekleidet, und zog wiederum ihre Witwenkleider an, sitzt und trägt Leid, und dieweil sie schwanger geworden, wartet sie nun weiter auf Gottes Segen.

96. Es ist ein wunderbarer Fleiß des Heiligen Geistes, diese schändliche unzüchtige Historie zu beschreiben, daß er auch Alles bis auf das Aeußerste so gar genau ausgeführt, daß er sich auch nicht gescheut, von der Geburt der Zwillinge und wie die Nachgeburt einen Riß gemacht, zu reden. Warum hat sich doch der allereinste Mund des Heiligen Geistes also hernieder gelassen zu solchen niedrigen verachteten Dingen, ja, die auch unzüchtig und unflätig, und dazu noch verdammlich sind, als ob solche Dinge dazu etwas sollten nütze sein, daß dadurch die Kirche und Gemeinde Gottes möchte gelehret werden? Was hat die Kirche damit zu thun?

97. Darauf antworte ich, wie zuvor auch, nämlich, daß dies alles um Christi willen also erzählt werde, welcher durch die ganze heilige Schrift als unser Bruder, unser Blutsfreund und Verwandter beschrieben wird, der uns auch noch etwas näher verwandt ist, denn sonst jemand von unsern geliebten Freunden sein mag. Denn er ist nach der Linie Abrahams aller Kinder Israel Bruder und Vetter. Nach der Mutter ist er ein Bruder, Vetter, Geschwisterkind aller Egyptier, Canaaniter, Amoriter; denn er ist geboren von Thamar, Ruth, Rahel und Bathseba, Urias Weib.

98. Mit uns aber, die wir von Japheth herkommen, ist er nicht Geschwisterkind, ist aber doch unser Schwager auch nach dem Fleisch. Ja, wenn wir hinauf gehen bis zu Adam, ist er eben sowohl auch unser Vetter und mit uns Geschwisterkind als der Juden, und mögen uns dieser Verwandtnis nicht weniger rühmen denn sie. Wiewohl solche Verwandtschaft nichts zur Sache thut, ja, sie ist gar verdammt, Joh. 1, 13: und Röm. 9, 8., wo nicht der Glaube dazu kommt und das Wort vom gebenedeiten Samen, dadurch verneuert und wiederum auferweckt wird das Fleisch, so vermesslich, sündlich und verdammlich ist, weswegen allein sich die Juden mit solchem großen Stolze rühmen. Was nützt es aber, ein Bruder oder Mutter sein vom fleischlichen Fleische, das dem ewigen Gerichte Gottes unterworfen ist? Wiewohl das nützlich ist, daß man wisse, daß solcher Ruhm den Juden nicht eigen ist, oder der ihnen insonderheit allein gebühret; sondern daß er den Heiden auch gemein und ihnen auch zugehörig ist, wo nicht nach Abraham, doch aber nach Adam und allen Patriarchen bis auf Noah; denn da wird das ganze menschliche Geschlecht unter drei Brüder getheilt. Uebrigens sind wir alle zugleich in den Lenden Noahs gewesen. Aber es war der gebenedeite Same noch nicht da, welcher den schweren Fall unserer verderbten Natur wiederum aufrichten sollte. Darum, sage ich, dient es nirgends zu, daß man sich der Gemeinschaft der Natur, so noch nicht wiederum auferweckt, rühmen will; und ist dies allein die rechte Ehre, daß Christus der Eva verheißen und gegeben ist.

99. Da wird nun der gebenedeite Same beschrieben, welcher von dem vermaledeyten, verlornen und verdammten Samen und Fleische hergekommen ist, doch also, daß derselbe Same an sich selbst ohne Sünde und unverderbt geblieben. Der Natur nach ist das Fleisch Christi einerlei mit unserm Fleisch: aber in seiner Empfängniß ist der Heilige Geist dazu gekommen, und hat das Fleisch, so er von der Jungfrau genommen, überschattet und gereinigt, daß er mit der göttlichen Natur vereinigt würde. Derhalben ist an Christo ein ganz heiliges, reines und sauberes Fleisch, aber in uns und allen Menschen ist es gar verderbt und unrein, außer allein so viel es in Christo wiederum verneuert wird.

100. Und also steigt der Heilige Geist da hernieder mit seinem allerreinsten Munde und greulichen Blutschande, damit anzuzeigen, daß das Fleisch Christi auch in Juda und Thamar gewesen sei, und daß er auch bei den Leuten sei, die da sündigen; wiewohl die Sünde nicht allein eine große Schande und scheußlicher Unflat, sondern auch verdammlich und vor GOtt eine schwere Schuld ist. Es ist der Heilige Geist nahe, sowohl bei den verderbten als unverdorben, bei den gerechten und ungerechten, den schändlichen und ehrlichen, bei scheußlichen und schönen Menschen; wiewohl er solches nicht thut, ist kein Anfänger oder Stifter der Sünde und des Verderbens, sondern wo er nicht daselbst gegenwärtig und nahe bei den Sündern wäre, wer wollte denn das schändliche und sündliche Fleisch wiederum reinigen und bessern?

\quellseite{314}

101. Derhalben zeugt die Erzählung in diesem 38. Capitel von der Gegenwart des Heiligen Geistes, daß er so nahe sei, daß er auch seine Zunge hergibt und dieselbe zu solchen scheußlichen und schändlichen Dingen braucht. Ja, wo er solche Dinge beschreibt, braucht er dazu viel mehr Worte als sonst in andern Dingen, so auch, daß es scheint, es sei ein Uebriges und nicht vonnöthen, also in die Länge davon zu reden; auf daß wir wissen sollen, daß GOtt immer auf seine Creatur sieht und sonderlich auf den Messias und Christus, welchen er also durch die Sünder und sündliche Geburt herführt.

\begin{quote}
\textbf{V. 20—22. Juda aber sandte den Ziegenbock durch seinen Hirten von Odollam, daß er das Pfand wieder holte von dem Weibe; und er fand sie nicht. Da fragte er die Leute desselbigen Orts, und sprach: Wo ist die Hure, die außen am Wege saß? Sie antworteten: Es ist keine Hure da gewesen. Und er kam wieder zu Juda, und sprach: Ich habe sie nicht funden; dazu sagen die Leute desselben Orts, es sei keine Hure da gewesen.}
\end{quote}

102. Ich habe droben gesagt, daß die Hebräer darüber disputiren, ob man im Hebräischen solle ro oder re lesen; denn Etliche haben es gegeben oder verdolmetscht, Freund, die Andern aber, Hirte. Und das ist auch des Hieronymus Auslegung, und es scheint, daß dieser Text mit ihm überein stimme, sintemal der Mann von Odollam als ein Diener von seinem Herrn gesandt wird. Darum gefällt mir das auch besser als die Auslegung der Juden. Denn wiewohl Juda bei ihm geherbergt, muß man es doch dafür halten, daß er eines großen Ansehens gewesen und sonderliche Gunst bei dem Wirthe gehabt, und dazu auch bei den andern Cananitern, welche ihn in großen Ehren gehalten, und ihm freundlich gereicht und mitgetheilt, was ihm zur Haushaltung vonnöthen gewesen.

103. Vielleicht hat er auch bei ihnen des Lehr- und Priesteramts gepflegt nach der Weise der andern Patriarchen, so an allen Orten den Gottesdienst geübt mit Lehren und Beten; wie man dergleichen viele Exempel in ihren Historien findet. Denn wo der Glaube ist, schweigt er nicht stille, sondern lehret und unterrichtet das Volk, vermahnt die Leute zum Glauben, zum Gebet, zu opfern und GOtt zu loben; wie Röm. 10, 10. geschrieben steht: „So man von Herzen glaubet, so wird man gerecht; und so man mit dem Munde bekennet, so wird man selig.“ Wenn du fest glaubst und hast das Wort im Herzen, so ist es unmöglich, daß du dasselbe nicht auch mit dem Munde bekennen und davon reden solltest. Also kommt aus dem Worte im Herzen auch das Bekenntniß des Mundes, daß du GOtt lobst und ihm dankst für seine Wohlthaten, daß er die Lehre seines Worts geoffenbart, und daß du dieselbige Andere auch lehrest und öffentlich also bekennest; wie Juda sein Hausgesinde, seinen Wirth und Nachbarn gelehrt und sie zur Erkenntniß GOttes gebracht hat, und wie die Apostel in der Welt umhergezogen, die Lehre ausgebreitet und GOtt angerufen haben. Sonst hätte Juda in dem fremden Lande nicht sicher und ohne Gefahr wohnen können, sondern wäre von denen, so nicht fromm gewesen, verstoßen worden, wo er nicht Etliche damit zu Freunden gemacht, daß er sie recht gelehrt und mit dem Worte ihnen gedient hätte. Derselbigen ist Einer sein Wirth auch, er sei nun gleich ein Hirt oder sonst ein guter Freund gewesen, der durch das Wort, so Juda gelehrt, zu dem GOtt Abrahams ist bekehrt worden.

104. Es braucht aber Mose hier abermals derselben Weise zu reden, wie er kurz zuvor auch gethan, welche bei den Hebräern gebräuchlich gewesen, und lautet im Lateinischen: ostio-oculorum super viam, das ist, auf der Wegscheide, da sich dieselbe anfängt und gleichsam geöffnet wird; das heißen sie ostium oculorum, wie man es im Lateinischen geben möchte. Da hat Thamar gesessen, an dem Orte, so nach Thimnath geht. Daß aber dieser Ort nicht weit von Thimnath gelegen, wird dadurch angezeigt, daß während Juda dahin gegangen, die Schäfe zu scheren, Thamar hinweg gegangen ist. Denn so bald Juda dahin gekommen, ist es glaublich, daß er ihr den Ziegenbock gesandt. Dies alles ist in Einem Tage geschehen. Darum muß die Wegscheide nahe bei der Heerde zu Thimnath gewesen sein. Ich halte dafür, es sei fast so weit gewesen wie ein Viertheil einer deutschen Meile, und daß Juda gemeint, sie würde daselbst verharren, bis er den Ziegenbock senden würde.

105. Man möchte auch wohl diesen Argwohn schöpfen, daß sie sich an demselben Ort ein Häuslein oder Hütte gemacht hätte, so zu solchem unzüchtigen Gewerbe gelegen und bequem gewesen, darin sie am Wege gesessen; in dasselbige ist vielleicht Juda gegangen, daß der Hirte draußen auf ihn gewartet. Denn es stimmt sehr wohl hiermit überein, daß sie den Betrug nicht von sich selbst erdacht, sondern daß sie dazu ihre Freunde oder Verwandten etwa unterrichtet und bestellet haben, welche gehört hatten, wie sie verworfen und wiederum in ihres Vaters Haus sollte gesetzt sein. Darum werden sie sich der Thamar und ihres Anfalles erbarmt und dem Juda also heimlich nachgestellt haben; und haben ihn auch mit solchem Betrug und List in den Strick gebracht und gefangen. Denn es ist nicht glaublich, daß sie sollte öffentlich unter dem freien Himmel gesessen haben.

106. Und diese meine Muthmaßung wird durch die Antwort der Nachbarn bestätiget. Denn da der Hirte nach der Hure gefragt, antworten die Leute desselbigen Orts, es sei keine Hure da gewesen, so sie doch ohne Zweifel Thamar werden gesehen und gekannt haben, sondern haben solches noch verhehlet. Und wie es scheint, haben sie also sagen wollen: Wir wollen es noch nicht anzeigen; ihr werdet es aber bald hernach selbst inne werden.

107. Auf diese Weise, halte ich, werden sie Thamar abgerichtet haben, daß sie dieselbe durch solchen Betrug wieder in ihre gebührliche Stätte und Ehre setzen möchten. Und es ist beinahe ein solcher Betrug, als der war, damit Jakob seinem Bruder Esau durch seiner Mutter Eingeben den Segen entzogen hat. Denn Thamar war gar mit ewiger Schmach verstoßen, und hat müssen daran verzagen, daß sie noch würde Kinder zeugen, wie die Worte Judas dahin lauten, da er zu ihr sagte: „Bleibe eine Wittwe in deines Vaters Hause“; daran er schwer gesündigt. Denn er sollte sie daheim in seinem Hause behalten und ihr seinen Sohn Sela zur Ehe gegeben haben nach göttlichem Recht, welches er selbst übertritt, und thut ihr mit großem Schimpf und Lügen Unrecht. Darum scheint es, daß dieser Betrug recht und billig sei. Denn Thamar hat sonst an ihrer Ehre verzagen müssen. Darum sind ihre Freunde und Verwandten bewogen worden, in Ansehung solches ihres Elendes ihr zu rathen und zu helfen, daß sie mit Recht oder Unrecht, wie sie nur könnte oder möchte, ihre mütterliche Ehre wiederum bekommen und sich zeugen möchte. Und dieweil Juda in großem Ansehen und auch zu gewaltig war, als daß sie ihn hätten zwingen können, Thamar wiederum in ihre gebührliche Ehre zu setzen, greifen sie ihn mit List an, und überwinden ihn auch damit.

108. Das nehme ich aus der Antwort, da die Nachbarn sagen: Was fragst du uns viel nach der Hure? Meinst du, daß wir hier solche lose Weiber unterhalten? Diese Worte zeigen an, daß sie um diese That ein Mitwissen gehabt und dieselbe angestiftet haben, sintemal sie ihm mit solchem Stolz geantwortet: Meinst du, daß wir Huren hier haben? Denn sie halten den Betrug heimlich, decken ihn zu, sagen nicht, daß sie ihre Blutsfreundin sei, und wünschen also heimlich, daß sie nur schwanger sein möchte, alsdann wollten sie das Ding wohl offenbaren; wie es hernach dem Juda verkündiget wird, wer die Hure gewesen sei.

109. Das hebräische Wort kadeschah haben wir bis daher in diesem Buche der Genesis noch nicht gehabt. Denn droben hat Mose gesagt, Juda habe gemeint, Thamar wäre eine zonah, auf Hebräisch, welches eigentlich eine Hure heißt, wie Rahab auch gewesen. Aber die Arabinen wollen haben, es sei ein solch Wort, das vielerlei Bedeutung habe, und heiße, Hure und auch Wirthin. Und daher kommt es, daß Etliche disputiren, Rahab sei keine Hure gewesen. Aber aus dem, so wir droben gehört, ist offenbar, was eigentlich die rechte Bedeutung dieses Wortes sei.

\quellseite{315}

110. Das Wort kadeschah aber, welches Mose in der hebräischen Sprache allhier braucht, heißt, heilig, wie Jesaia am 6. Cap. V. 3. steht: „Heilig, Heilig“ etc., und ist ein sehr gemein Wort. Von denselben kommt ein ander hebräisch Wort, mikdasch, welches heißt, das Heiligtum oder der Tempel; und die Geheiligten heißen diejenigen, so zu etwas verordnet oder zubereitet sind; als, wenn man sagt: Der Medien Könige sind geheiligt, das ist, sie sind zum Kriege gerüstet. Also braucht Mose dies Wort oftmals in mancherlei Bedeutung. 5 Mos. 23, 17: „Es soll unter den Kindern Israel nicht sein, weder Hure noch Bube“, das lautet im Hebräischen weder kadesch noch kadeschah, es sollen alle Eheleute, Ehemänner und Eheweiber sein.

111. Es scheint aber, als habe der Hirte etwas züchtiger und ehrlicher reden wollen: darum nennt er sie auf Hebräisch kadeschah, eine Wirthin, nicht wie Juda, eine Hure, wie das hebräische Wort sonst eigentlich heißt. Denn es sind in der hebräischen Sprache viele Wörter wie auch in andern Sprachen, die nicht allein mancherlei Bedeutung haben, sondern auch also gebraucht werden, daß man das Gegenteil verstehen muß; welche Figur antiphrasis heißt. So sind in der lateinischen Sprache, bellum, lucus, sacrum, etc., wie der Poet sagt: Auri sacra fames, heiliger Hunger nach Gold, das ist: Du verfluchter und verdammter Geiz. Eben dasselbige ist in der deutschen Sprache auch sehr gemein; als, wenn man sagt: Ei! das ist ein fromm Kind, ein frommer Gesell; ei! der will gut werden. Denn da wird das Wort „gut“ oder „fromm“ verstanden für schändlich oder böse, durch diese Figur, antiphrasis genannt. Also heißt er an diesem Orte dies Weib eine heilige Hure, als wollte er sagen: Ei! das heilige Kindchen, ein fromm Töchterlein; denn die Nachbarn haben listig geantwortet, wie gesagt, nämlich, sie haben daselbst eine Hure gesehen. So gar fleißig hat der Heilige Geist Alles erzählt, daß es Wunder ist, daß er in solchen Erzählungen so viel Worte verloren hat. Aber er trägt davor keine Scheu, daß er auch von ungöttlichen und schändlichen Dingen handelt, auf daß dieselben auch der Kirche zur Lehre und Trost mögen vorgehalten werden.

112. Juda sprach: Sie habs ihr, sie kann uns doch ja nicht Schande nachsagen; denn ich habe den Bock gesandt, so hast du sie nicht funden.

112. Juda antwortet: „Siehe, ich will nicht Streit haben. „Sie habs ihr“, sagt er, wie sie will, und setzt noch eins hinzu: „Sie kann uns doch ja nicht Schande nachsagen.“ Denn das ist der rechte Verstand des hebräischen Textes. Was aber das gesagt sei, weiß ich nicht. Unser Dolmetscher hat es gegeben: „Sie kann uns ja nicht Lügen strafen, als hätte ich ihr den Ziegenbock nicht gesandt. Die Rabbinen der Juden legen es aus: Daß diese Schande nicht offenbar werde und an Tag komme. Welche Auslegung nicht allein falsch und unrecht, sondern auch lächerlich ist. Denn dies Ding ist ja offenbar, wie sollte es denn heimlich können gehalten werden? Juda sendet ihr den Bock und der Hirte fragt mit deutlichen Worten nach der Hure, gibt damit zu verstehen, daß er den Bock gebracht und damit den Hurenlohn bezahlen wolle. Das heißt ja nicht die That verhehlen oder heimlich halten. Die Nachbarn haben gefragt, was er mit dem Bock meine? da hat er geantwortet: Der Bock sei der Hure gesandt worden; da ist gar nichts Heimliches gewesen, weder an Worten noch an Werken.

\section[Dritter Theil]{Dritter Theil. Von Thamars Schwangerschaft; wie sie darüber soll verbrannt werden; und wie sie endlich zwei Söhne gebieret}

\subsection{I}

\begin{quote}
\textbf{V. 24. Ueber drei Monden ward Juda angesagt: Deine Schnur Thamar hat gehuret; dazu siehe, sie ist von Hurerei schwanger worden.}
\end{quote}

113. Dies werden ohne Zweifel der Verwandten und Freunde Thamars Worte gewesen sein, welche auf sie gesehen und für sie gesorgt haben. Denn sonst würden sie sich darum nicht bekümmert haben, ob sie ehrlich oder unehrlich gelebt. Aber dieweil sie dies dem Juda sobald angesagt, scheint es, daß sie diesen Rath mit der Hurerei werden gestiftet haben. Zuvor haben sie es zwar mit Fleiß verhehlet, da sie noch nicht eigentlich gewußt, daß sie schwanger wäre. Denn wo Thamar nicht wäre schwanger geworden, wußten sie wohl, daß es viel besser wäre, diese That zu verhehlen, denn daß sie offenbar würde. Jetzt aber, da sie nun merken, daß sie schwanger sei, haben sie solches dem Juda mit Freuden angesagt. Da gehet das Gaudium an.

114. „Ueber drei Monden“, sagt Mose; da das Kindlein hat angefangen zu leben, haben sie nicht gewartet, bis sie an ihrem Leibe wäre etwas stärker geworden und das Kindlein auch gewachsen wäre, sondern sie eilen und haben getrotzt, und Juda gleichsam gespottet, daß er von der Thamar also listig hintergangen war, und haben es nun geoffenbart, wer die Hure sei, welche sie zuvor nicht haben kennen wollen, da sie der Hure gesucht. Also wollten sie sagen: Du eilst allzu sehr mit dem Hirten und mit dem Bocke, du kommst zu früh. Höre nun jetzt zu: Thamar, deine Schwiegertochter, hat gehurt. Wer aber die Sünde mit ihr begangen habe, da magst du nachsehen; denn es ist niemand, den wir dieser Sünde können zeihen, ohne allein dich. Da spricht nun Juda ein Urtheil über seinen eigenen Hals; wie David auch gethan, da Nathan zu ihm kam, 2 Sam. 12, 5—13.; derselbe hat aber auch seine Sünde bekannt.

\subsection{II}

\begin{quote}
\textbf{Juda sprach: Bringet sie hervor, daß sie verbrannt werde.}
\end{quote}

115. Da siehst du, daß dies Gesetz auch von den Vätern hergekommen sei, wie das andere auch, daß ein Bruder seinem verstorbenen Bruder hat müssen Samen erwecken. Denn es ist dazumal auch eine Strafe gesetzt gewesen auf die Töchter der Priester, wenn sie gehurt hatten, nämlich, daß man sie mit Feuer verbrennen sollte, auch die schon ehelich gewesen sind. Derhalben wird Thamar auch gefordert, daß solche Strafe über sie ergehen solle. Denn sie ist eines Priesters Weib gewesen, und vielleicht in der andern Linie und Geschlechte der Canaaniter auch eines Priesters Tochter. Indes daß kann man nicht genau wissen.

116. Denn daß sich die Juden träumen lassen, sie solle von Messiasgeschlecht geboren sein, ist allzuweit hergeholt, ist auch ungereimt zu hören. Denn Melchisedech ist dazumal bei fünfzig Jahren todt gewesen; und wenn wir schon wollten sagen, daß sie sollte in dem letzten Alter gezeuget sein, wollte daraus folgen, daß sie müßte fünfzig Jahre gewesen sein, da sie Juda's Söhne zur Ehe genommen, welches unmöglich ist.

117. Ich halte aber dafür, daß nicht allein die Töchter, sondern auch die Eheweiber und sonst alle Weiber in den Geschlechtern der Priester diesem Gesetze und Urtheile unterworfen gewesen seien. Denn damit hat Gott also einen sonderlichen großen Unterschied machen wollen zwischen weltlicher Obrigkeit und dem Amte der Priester, und hat der priesterliche Stand vor dem andern mit größerem Ernste über die Gesetze der Zucht und Ehrbarkeit halten müssen vermöge solcher Gesetze. Darum hat er auch den Uebertretern größere Strafe verordnet. Jakob aber ist ein Priester gewesen mit allen seinen Söhnen, so hin und wieder zerstreuet waren, und hat Jeder an seinem Orte gelehrt, gebetet und geopfert, wie Ismael und Isaak im Hause Abrahams gethan, und war unter ihnen Einer der Erstgeborne, ein Herr und Regent über die Andern alle. Also ist Ruben, ehe denn er vom Vater verworfen, der Erstgeborne gewesen, welchem auch zwei Theile des Erbes gebührten. Die Heiden, zu denen sie gekommen sind, haben sie auch für Priester gehalten, und da sie ihre Lehre gehört, ist es ihnen dazu sehr lieb gewesen und haben sich dessen gefreut, daß solche heilige Leute zu ihnen wären gekommen, die das Volk von der Erkenntniß Gottes, und von der Anrufung oder vom Gebet lehrten und berichteten, und für sie beteten möchten; und deshalb haben sie dieselben mit Ehrerbietung aufgenommen und sind mit ihnen freundlich umgegangen.

118. Derhalben, da Thamar den Kinderkindern Jakobs vermählt gewesen, die da Priester waren, und der Hurerei halben bezüchtigt worden, wird ihr die Strafe alsobald zuerkannt, daß sie solle verbrannt werden um des Schandfleckes willen, damit sie das ganze Geschlecht und Priesterstand beflecket hätte. Vielleicht hat Juda auch diese Ursache zum Befehl genommen, da er gesehen, daß sie also mit gutem Fug hingerissen und getödtet werden möchte, damit sie ihm und seinem Geschlechte hinfort nicht mehr beschwerlich sein möchte: darum verdammt er sie halb zum Tode ohne alle Berathschlagung, ja, auch ohne alles Mitleiden; und hat wahrlich ein betrübtes und greuliches Urtheil über dies Weib gefällt, die doch mit schwerem Leibe gegangen, so man ihrer doch um der Frucht willen sollte verschonet haben. Es sei denn, daß dies der Gebrauch gewesen, daß sie im Gefängniß verwahret und behalten worden wäre bis sie geboren hätte. Ist aber die Frucht zugleich mit der Mutter gestraft worden, so ist es wahrlich ein überaus hartes Gesetz gewesen.

\quellseite{316}

\begin{quote}
\textbf{V. 25. Und da man sie hervor brachte, schickte sie zu ihrem Schwäher, und sprach: Von dem Manne bin ich schwanger, deß dies ist. Und sprach: Kennet du auch, wes dieser Ring, und diese Schnur, und dieser Stab ist?}
\end{quote}

Nun wird Juda im Strick gehalten, darein er durch Betrug und List der Thamar und ihrer Verwandten gekommen war. Nun, wohlan, werden sie zu ihm gesagt haben, du hast wollen diese Hure sehen, wer sie wäre: es ist unsere Blutsfreundin und deine Schwiegertochter gewesen, welche du verworffen. Nimm den Stab, den Ring \&c., erkenne doch, woher die in unsere Hände gekommen sind?

\begin{quote}
\textbf{V. 26. Juda erkannte es, und sprach: Sie ist gerechter, denn ich; denn ich habe sie nicht gegeben meinem Sohn Sela. Doch beschlief er sie nicht mehr.}
\end{quote}

Juda hat sich selbst sehr hart beschuldiget, und Thamar absolvirt und losgesprochen; er entschuldiget oder verhehlt seine Sünde gar nicht. Und es ist wahrlich eine wunderbare Beschuldigung, daß er lieber alle Schuld auf sich legt denn auf Thamar. Derhalben wird nun gefragt, welcher unter ihnen am schwersten gesündigt habe, Juda oder Thamar? Er hat keinen Ehebruch oder Blutschande begangen, sondern bloß Hurerei getrieben; Thamar aber ist an beiderlei Sünden schuldig, darum daß sie eine Braut war und ein Eheweib des dritten Sohnes Judas, nämlich, Sela’s, welcher nach göttlichen Rechten ihr Ehemann ist und der sie des Ehebruchs hat beschuldigen mögen; daß sie aber bei ihrem Schwiegervater gelegen, ist eine sehr schändliche Blutschande.

Derhalben können sie auf keinerlei Weise entschuldiget werden, wiewohl die Sünde Judas etwas geringer ist. Wo aber ein rechtschaffener Redner dazu gekommen wäre, der sich der Thamar angenommen, ihre Sache rechtlich auszuführen, der hätte dies noch etwas heftiger und größer machen können. Und werden vielleicht die Gefreunden der Thamar mit Juda also geredet und ihn beschuldiget haben: Du hast nicht allein Thamar geschmäht, sondern auch unser ganzes Geschlecht; du hast uns der Ehre beraubt, daß Thamar keine Mutter geworden ist, hast uns deines Sohnes beraubt; also hast du uns allesammt unsere Nachkommen verderbt, so viel an dir gewesen. Solche Stücke werden sie ohne Zweifel dem Juda aufgerückt haben, wiewohl sie nicht geschrieben sind, und werden gesagt haben: Du bist ein Priester und Lehrer des göttlichen Worts, regierst andere Leute mit Lehre und Anweisung; aber du entstellst mit solchen bösen Exempeln dein Amt schändlich. Es hätte dir gebühret, daß du fest solltest gehalten haben über den Gesetzen, darin Gerechtigkeit und Ehrbarkeit geboten; aber dieser Gesetze ungeachtet, schmähst du jetzt die, welche du solltest geehrt und genähret haben. Derhalben solltest du auch viel billiger gestraft werden denn Thamar, unsere Verwandte, welche du ihrer höchsten Ehre beraubt hast, nämlich, der mütterlichen und priesterlichen Ehre, und wo diese Ehre verloren, da ist auch Einem dies Leben bitter und ganz zuwider.

122. Darum hat Juda recht gesagt: „Sie ist gerechter, denn ich“, wiewohl die Blutschande eine große Sünde ist. Aber er hat mehr und größere Sünden begangen, welche stracks wider GOttes Gesetz und Recht sind, und möchte wohl des Kirchenraubes, Todtschlages und sonst allerlei Sünden und Unrechts beschuldiget werden. Also ist aber Juda gedemüthigt worden, damit er seine Sünde erkennen und bekennen möchte, nämlich, daß er sie aus dem Hause verstoßen und sie des Priester- und Mutterrechts beraubt hätte, ja, auch zuletzt all Nachkommen, Kinder und Erben.

123. Und auf solche Weise hat dieser Patriarch müssen beschrieben werden, auf daß er wüßte, daß er ein Sünder wäre, damit er andere Sünder vor sich nicht verachtete, als ob er selbst nicht auch fallen könnte, dieweil er priesterliche Gewalt und Ehre gehabt und dazu das Lehramt geführt. Denn es sind alhier viel Sünden, die alle groß und greulich auf einem Haufen zusammen gekommen sind.

124. Dies aber wird uns zum Trost vorgehalten. Die großen Heiligen müssen große Fälle thun, daß GOtt damit anzeige, daß er haben wolle, daß alle Menschen sollen gedemüthigt werden und im Register der armen Sünder auch mit stehen; und wenn sie das erkennen und bekennen werden, sollen sie bei ihm Gnade und Barmherzigkeit finden. (Geschieht etwa ein Fall, wie soll man ihm thun? Haben doch die, so Christum an das Kreuz geschlagen, dies Gebet Christi gehört: „Vater, vergib ihnen, denn sie wissen nicht, was sie thun.“) Man soll sich zwar fleißig vor Sünden hüten; ist aber jemand gefallen, der soll darum nicht verzagen. Denn GOtt hat beiderlei verboten, nämlich, Verzweiflung und auch Vermessenheit, zur Linken und zur Rechten weichen. Auf der rechten Seite soll man nicht stolz oder vermessen sein; auf der Linken aber auch nicht verzagen oder verzweifeln. Auf der geraden Heerstraße soll man bleiben. Ein armer Sünder soll das Vertrauen auf GOttes Barmherzigkeit nicht fallen lassen; ein Gerechter soll nicht stolz sein. „Der Herr hat Gefallen an denen, die ihn fürchten, die auf seine Güte hoffen“, wie der 147. Psalm V. 11. sagt. Er hält die, so vermessen und stolz sind, und liebt die, so noch in der Furcht, Hoffnung und Zuversicht beharren, nicht auf uns selbst, sondern auf GOttes Barmherzigkeit.

125. Also werden nun die Heiligen gedemüthigt und ihre Sündenfälle bekannt gemacht, daß sie uns mit ihren Exempeln unterweisen und die rechte Mittelstraße zwischen Verzweiflung und Hoffart oder Vermessenheit zeigen.

126. Zuletzt, da Juda seine Sünde erkannt, legt er sich selbst eine Strafe auf, daß er damit anzeige, daß er ernstliche Buße gethan; denn er hat sich hinfort Thamars und auch anderer Weiber enthalten. Er wird aber auch ohne Zweifel seine Verwandten um Verzeihung gebeten haben, darum daß er ihre Freundin geschwächt; und sie werden wiederum ihn auch um Verzeihung gebeten haben, darum daß sie ihn also listig betrogen und zu solcher sehr schweren Sünde gereizt hätten.

127. Und also wird hier ein großer Haufen Sünden zusammengetragen, nämlich, Judas, Thamars und ihrer Gefreunden. Aber die Sünde Judas ist größer um des Lehramts und Kirchenregiments willen, welches er bis daher verwaltet. Da aber nun die Sünden alle, damit sie einander verletzt, vergeben, sind sie beiderseits wiederum versöhnt und Thamar auch in ihre vorige Ehre im Hause Juda wieder gesetzt worden; wiewohl sie lieber seinen Sohn Sela geheiratet hätte. Dieweil aber solches nicht mehr hat geschehen können, ist ihr doch dasselbe angenehm gewesen, daß sie wiederum aus ihres Vaters Hause, da sie in Gefahr ewiger Unfruchtbarkeit gestanden, in das Haus geführt wird, darin ihr die mütterliche Ehre und das Hausregiment von GOtt war zugeeignet worden.

128. Was aber das Amt der heiligen Matronen, als Sarahs, Rebeks, Leas, Rahels und Thamars gewesen sei, ist aus den Historien der Patriarchen offenbar. Es sind Hirtenweiber gewesen, darum haben sie ihr Haus regiert, des Viehes gewartet, die Kühe und Ziegen gemolken, die Wolle, Milch, Käse, Butter und dergleichen Anderes gesammlet und zu Rathe gehalten. Dies ist die Hausarbeit gewesen bei den Patriarchen, und solcher Arbeit hat Frau Thamar vorgestanden. Des Sela Weib hat wohl im Hause auch ihre eigene Städte gehabt; aber ich halte dafür, Thamar werde die Vornehmste gewesen sein. Es sind keine herrliche Matronen gewesen, sind nicht in Mißgängen, Unkeuschheit oder guten Tagen verdorben, sondern sind fleißig im Amt der Haushaltung gewesen. Es ist auch mit der Viehzucht ein weitläufig und mühselig Ding und voll schwerer Arbeit gewesen, dazu viel Hausgesinde gehört. Denn also sagt hernach Jakob vor Pharao 1 Mos. 47, 3.: „Deine Knechte sind Viehhirten, wir und unsere Väter.“

129. Darum kann man aus nichts besser abnehmen, was diese Matronen für ein Leben geführt, denn eben daraus, was ihre Ehemänner für Habe und Güter gehabt. Sie haben nicht eigene Häuser, Acker, Holz oder Bergwerke gehabt, sondern Vieh, Milch, Butter, Käse. Es sind eitel Viehmägde gewesen, wie jetzt bei uns auch etliche Matronen mit solcher Hausarbeit und Viehzucht umgehen. Und zu derselben Zeit haben sich auch Könige und Fürsten damit bemüht und sich desselben beflissen; und ist dazumal keine andere Hanterung gewesen, denn daß sie Vieh um andere Waare gegeben und also mit den Gütern getauscht haben. Also gibt Juda einen Ziegenbock zum Hurenlohn, der ein großer Prophet und Priester war und daheim eine große und zahlreiche Familie hat, hat aber doch weder Heller noch Pfennig, den er dem Weibe hätte geben mögen. Daran ist aber kein Zweifel, es werden diese Matronen für uns in ihrem Amt sehr fleißig gewesen sein, sintemal zur Viehzucht sonderlich großer Fleiß, Sorge und Arbeit gehört. Und sie haben sich nicht geschämt, solche geringe unschätzige Arbeit auf sich zu nehmen; denn Gott hat auch Gesellen daran und ist ein Herr und Hüter dieser sehr heiligen Matronen und Erzhausmütter gewesen, da er doch dieweil anderer Fürsten und Könige nicht sonderlich gedacht hat.

\quellseite{317}

\subsection{III}

\begin{quote}
\textbf{V. 27—30. Und da sie gebären sollte, wurden Zwillinge in ihrem Leibe erfunden. Und als sie jetzt gebar, that sich eine Hand heraus. Da nahm die Wehmutter, und band einen rothen Faden darum, und sprach: Der wird der erste herauskommen. Da aber der seine Hand wieder hinein zog, kam sein Bruder heraus; und sie sprach: Warum hast du um deinetwillen solchen Riß? Und man hieß ihn Perez. Danach kam sein Bruder heraus, der den rothen Faden um seine Hand hatte. Und man hieß ihn Serah.}
\end{quote}

130. Siehe doch, wie fleißig der Heilige Geist diese erbärmliche und elende Geburt beschrieben hat. Denn er hat Lust an seinem Geschöpf und Creatur. Wenn wir aber dieselbe ansehen, haben wir daran nicht allein keine sonderliche Freude, sondern sind auch so gar verblendet, daß wir die Werke Gottes nicht sehen. Es ist ein groß wunderbärlich Ding um die Zeugung und Geburt auch an den unvernünftigsten Thieren, als, wenn eine Kuh ein Kalb, das Schaf ein Lämmlein bringt; wie viel mehr ist sich dessen zu verwundern und ehrlich zu halten, daß ein Weib schwanger wird und einen Menschen gebieret! Denn wir werden ja alle aus dem Leibe des Weibes auf die Welt geboren durch die Geburt, das ist, durch den Tod; sintemal die Mütter mit den Kindern in gewisser Gefahr des Todes stehen. Und ist auch bei den Heiden eine sonderliche Hochachtung und Ehrfurcht vor diesem Werk geblieben; wie die Römer um der Geburt willen den Weibern erlaubt haben, goldenen Schmuck zu tragen, daß sie damit geehrt und getröstet werden sollten wegen dieser schweren Last, Mühe und Arbeit, so sie tragen müssen, wenn sie Kinder gebären. Dieweil aber Gott dies Wunderwerk also über das ganze menschliche Geschlecht, dazu auch über alle Thiere auf dem ganzen Erdboden ausdehnt und so gemein macht, darum hält man es, dieweil es täglich geschieht, also verächtlich: aber der Heilige Geist zieret und ehrt sein Geschöpf, und hat Lust, dasselbige anzuschauen und zu rühmen.

131. Es wird aber an diesem Ort nicht allein die Geburt also beschrieben um der Kinder willen, sondern wird auch die große Gefahr der Mutter, da sie gebieret, erzählt. Thamar hat selbst nicht gewußt, daß sie Zwillinge in ihrem Leibe getragen; wie denn die schwangeren Weiber solches nicht wissen können, ehe sie gebären. Da sie aber nun in der Geburt ist, lassen sich die Zwillinge sehen; da wird der Jammer angegangen sein. Die Sache hat sie beide, Juda und Thamar, in große Furcht und Schrecken gebracht. Denn sie sind in der größten Angst und Gefahr gewesen, und werden auch ohne Zweifel sehr hart geplagt worden sein, da sie an ihre Sünde gedacht, deren sie ihr eigen Gewissen beschuldigt hat. Ach, hat Juda gesagt, was haben wir nun gethan! Es wird der Herr die Mutter tödten und ich werde der zwei Söhne auch beraubt werden. Denn allezeit, wenn Trübsal und Gefahr vorfällt, kommt das Gewissen und der Neuling auch dazu, daß dem Menschen dann alleerst angst und bange wird, darum daß er die Sünde begangen hat. Und kann sich dessen das menschliche Herz nicht enthalten, es muß da zittern, erschrecken und sich selbst beschuldigen und sagen: Siehe, ich habe gesündigt; nun muß ich für solche Sünde auch diese Strafe leiden.

132. So sehen wir nun, daß um solcher Gefahr willen diese Geburt sehr beschwerlich gewesen, welche nicht allein die Eltern, sondern auch das ganze Geschlecht in so große Schrecken und Angst gebracht. Sela, der Sohn Judas, dazu der Hirte oder Wirth und die andern Leute zu Odollam haben kümmerlich geseufzt und mit Seufzen und Beten Gott um Hülfe angerufen. Und das ist die rechte Uebung der Priesterschaft gewesen, daß sie mit Ermahnung, Anrufung und Bekenntniß der Sünden umgegangen sind, haben gesagt: Lieber Herr Gott, wir bekennen, daß wir gesündigt haben; erbarme dich unser und vergib uns unsere Sünde. Diese Anzeichen, dabei wir solches abnehmen, sind nicht ungewiß und vergeblich. Denn solche Dinge pflegen oftmals im Leben der Heiligen vorzufallen, und regt sich das Gewissen der Sünde halben nimmer so sehr und ist nimmer so kräftig und stark, als wenn man in Trübsal, Angst und Schaden gerathen ist. Wenn es uns wohl geht, so ruht die Sünde und schläft, und wir leben immer so sicher dahin. Denn ohne das Gesetz ist die Sünde todt, wie St. Paulus von sich selbst sagt Röm. 7, 9.: „Ich lebte etwan ohne Gesetz.“ Aber die Trübsal oder Unglück öffnet Einem die Augen, und weckt das Gewissen recht auf, daß es nicht mehr ruhen kann, sondern immer beißt und martert.

133. Wir müssen denken, daß es allhier in diesem ganzen Geschlechte auch also ergangen sei. Denn da Thamar angefangen hat zu gebären, lassen sich in ihrem Leibe die Zwillinge sehen und Einer von ihnen hat die eine Hand heraus gegeben. Da ist gewißlich große Noth vorhanden gewesen und sie haben den Tod selbst vor Augen gesehen. Denn das ist große Gefahr, wenn in der Geburt erstlich der Fuß oder die Hand des Kindleins hervorkommt, und die Mutter kann das Kind nicht recht ordentlich gebären, es sei denn, daß das Haupt erstlich hervorkomme, daß die Wehmutter dasselbe ergreifen könne. Wo aber Hände oder Füße zuvor herauskommen, ist es sehr ungelegen und gefährlich, so daß man da Gott nicht aus Gnaden wunderbärlich hilft, oftmals die Frucht bei einzelnen Stücken aus der Mutter Leibe gerissen werden muß. Da vergeht zwar den armen Weibern der \unsicher{Kibel} wohl, und da treten heran große tödtliche Schmerzen anstatt der vorigen Lüste, die sie mögen gehabt haben. Da ist der Scherz aus und geht das Sterben an. Wie an diesem Orte drei Seelen zugleich sterben, als, der zwei Zwillinge und der Mutter. Wie soll dem Juda da das Herz im Leibe gezappelt haben?

134. Die Wehmutter hat zwar Fleiß gethan, was ihr Amts halben gebührte, sie hat sich gestellt, als sei ja noch Hoffnung und Stärke vorhanden, daß sie damit der Mutter möchte ein Herz machen. „Der“, sagte sie, „wird der erste herauskommen.“ Denn das sind Worte, damit sie die Mutter getröstet, und angezeigt, daß noch gute Hoffnung vorhanden; wie auch die Wehmutter zu Rahel sagte, da sie dieselbe trösten, aber doch nicht zuvor sagen wollte, wie es gehen würde: „Fürchte dich nicht, denn diesen Sohn wirst du auch haben“, 1 Mos. 35, 17. Aber Rahel ging darüber dahin, ward überwunden von der schweren Arbeit, daß sie also über der Geburt sterben mußte.

135. Es gebührt aber den Wehmüttern, daß sie beherzt seien, und wie Virgili von Aeneas sagt, sich äußerlich stellen, als hätten sie noch gute Hoffnung, und daß sie die armen Mütter in der Arbeit trösten und sagen, daß sie selbst auch gute Hoffnung haben wollen. Denn die Hoffnung macht, daß diejenigen, so Schaden, Kummer und Noth leiden, noch immer wunderbarer Weise in ihren Herzen sich trösten; wie Aeneas daselbst in Virgili weiter sagt: Dabít Deus his quoque finem, das ist: Gott wird dermahleinst auch diesem Unglück ein Ende machen. Solche Klugheit und Stärke soll auch eine Wehmutter an sich haben, daß sie auch in gegenwärtiger offener Noth und Gefahr noch immer in ihrem Herzen gute Hoffnung habe, und daß sie mit tröstlichen Worten der Mutter in der Arbeit auch solche Hoffnung mache.

136. Nun ist aber das Kind noch nicht geboren, und es scheint, als ob die Gefahr immer größer werde, da das Kindlein seine Hand wieder hinein zog. Wiewohl nun die Wehmutter sie tröstet und sie auch getrost sein heißt, ist doch die Hoffnung sehr ungewiß. Darum werden sie da immer mehr geängst und gebetet haben, und gesagt: Herr, nun hilf! Es war zuvor noch ein wenig Hoffnung vorhanden; nun aber ist dieselbe gar dahin gefallen. Darum ist die Sünde auch wiederum lebendig geworden und hat geherrscht, und da sie an die Blutschande gedacht, hat solches beide Judas und auch Thamars Herz sehr bekümmert und geplagt. Und ist zwar diese Buße und Genugthuung einmal schwer genug, dazu auch ein großer Jammer und ganz bittere Reue gewesen. Denn da war die Gefahr vorhanden, daß die Schmerzen der Geburt aufhören würden, ohne welche keine Mutter gebären kann. Und ist die erste Arbeit der Natur, da sie hat gebären sollen, gar vergeblich gewesen, nachdem das Kind die Hand wieder hinein gezogen. Darum hat Juda in großer Angst und mit Thränen geschrieen und gebetet, daß ja die Mutter sammt den Zwillingen nicht möchte umkommen. Und dazumal hat die Anfechtung der Lustseuche und anderer Schande oder Leichtfertigkeit nicht Statt gehabt, dazu sonst die Leute im Müßiggange und guten Tagen, wo man Alles vollauf hat, ihre Lust zu haben pflegen. Da gilt es Leib und Leben, daß Mutter und Kinder bei einander bleiben, wenn die Wehen vorüber sind. Das ganze Haus ist da mit Weinen, Seufzen, Heulen und betrübtem Wehklagen erfüllt worden.

\quellseite{318}

137. Es ist aber in demselbigen Augenblick Gott der Herr dabei gewesen und hat zu solcher Geburt durch ein neu Wunderwerk noch Hülfe erzeigt. Das Erste, da der Text sagt: „Und als sie jetzt gebar, that sich eine Hand heraus“, ist erbärmlich und kläglich: das Andere aber, das also lautet: „Und siehe, da kam sein Bruder heraus“, ist voll Trostes und Freude, so nach dem erbärmlichen Schrecken und Angst gefolgt ist. Denn dieser ist recht und glücklich geboren worden, darum hat sie des Vorigen halben nicht wollen verzagen, welcher nun der Andere geworden ist. Und die Wehmutter hat sich gleichsam darüber entsetzt und mit großer Verwunderung gerufen: „Warum hast du um deinetwillen solchen Riß gerissen?“ Als wollte sie sagen: Du Schalk, daß dich Gott behüte, du hast schier einen bösen Riß gethan.

138. Das Wort pariz heißt bei den Hebräern ein Mörder, Straßenräuber, Straußdieb, der nicht auf der rechten Heerstraße wandelt, sondern im Holze und auf den Beiwegen hin und wieder geht, den Leuten, so über Land ziehen, heimlich nachzustellen, sie zu überfallen und zu beschädigen: oder, die sich sonst weigern, gemeinen Gesetzen und Ordnungen sich zu unterwerfen, sondern gehen mit Gewalt dahin, toben und wüthen nach ihrem tollen wilden Kopf. Derhalben, sagte die Wehmutter, ist dieser Riß um deinetwillen geworden, damit du zuerst herauskämmest. Es ist hier die rechte Ordnung der Natur nicht gehalten worden, sondern du bist durch einen solchen Riß hervorgekommen und geboren worden, und um deinetwillen haben beide Mutter und Zwillinge große Gewalt erlitten; denn du hättest sie beinahe mit solcher Gewalt gar getödtet.

139. Also hat sich die Wehmutter gefreut und fröhlich gemacht um solcher glücklichen Gewalt willen, so beide Mutter und Kindern angenehm gewesen. Denn sie hat einen fröhlichen und lustigen Ausgang genommen, die Wehen, so sie denen Leib und Leben gebracht, derenthalben sie sich besorgt, daß sie umkommen und sterben würden. Es ist nun wohl gerathen, sie denken nun der Angst nicht mehr, sonderlich die so erschrecklich und gefährlich gewesen, wie Christus sagt Joh. 16, 21. Denn es ist über alles Hoffen und ganz unversehens die Mutter mit den Zwillingen durch diesen Riß am Leben erhalten worden und es ist mit der andern Geburt auch glücklich und leicht zugegangen.

140. Bis daher geht nun die Historie von Juda und Thamar, welche ich lieber einfältig habe erklären wollen, denn daraus sonderliche Allegorien und heimliche Deutungen dichten. Denn es ist viel nützlicher und nöthiger, daß man die Historie handele und erkläre allein um der ungeschickten und ungereimten Glossen willen der Juden, welche die Ordnung der Historie gar unter einander werfen und verwirren. Derhalben sollen wir uns wissen zu erinnern, wie anfänglich in dieser Historie und sonst auch an andern Orten mehr von uns ist angezeigt worden, daß Juda ehelich geworden sei, ehe denn Joseph ist verkauft worden: welches Mose droben hat angefangen zu erzählen und das er in nachfolgenden Capiteln vollends ausführen wird. Denn wie dies unter allem andern Unglück und Beschwernissen, so Jakob gehabt, das beschwerlichste ist, also ist es auch das letzte gewesen.

\chapter{Das neununddreißigste Capitel}

\section[Erster Theil]{Erster Theil. Wie Joseph an Potiphar zur Knechtschaft verkauft worden}

\begin{quote}
\textbf{V. 1. Joseph ward hinab in Egypten geführt; und Potiphar, ein egyptischer Mann, des Pharao Kämmerer und Hofmeister, kaufte ihn von den Ismaeliten, die ihn hinab brachten.}
\end{quote}

1. Mose wiederholt im Anfang dieses Capitels etliche der Stücke, davon er droben angefangen zu reden, nämlich, wie Joseph verkauft worden; welches Verkaufen er vor der Erzählung des vorigen Capitels durch die Figur hysteron proteron beschrieben, so es doch erst geschehen, da Juda schon von seinen Brüdern geschieden und ehelich geworden war und da auch seine drei Söhne geboren gewesen. Hernach aber, da Er und Onan gestorben und Thamar von ihrem Schwiegervater geschändet, ist Joseph vier und zwanzig Jahre alt gewesen und schon nach Egypten geführt worden; daß also der erste Theil des vorigen Capitels vorhergegangen, ehe Joseph ist verkauft worden; der andere Theil aber ist nachgefolgt, wie dasselbe aus der Zeitrechnung zu sehen.

2. Daraus aber ist abzunehmen, daß Joseph dreimal verkauft worden: Erstlich den Midianitern, darnach den Ismaeliten und zum Dritten Potiphar. Denn es sind vielleicht die Midianiter nicht in Egypten hinein gezogen, darum haben sie ihn den Ismaeliten, so da hinein gereist, wiederum übergeben. Wir wollen aber erstlich die Worte erklären, welche den Dolmetschern viel zu schaffen geben. Das Wort „Potiphar“ ist schier ein deutsch Wort, pot, fett, pharr, als wenn man sagen wollte, ein fetter dicker Farre. Wir können aber nicht wissen, was die richtige Ableitung dieses Wortes sei; denn es ist ein egyptisches Wort.

3. Das hebräische Wort saris deuten Alle, daß es so viel heiße als Eunuche. Und ist dies Wort fast gemein in der heiligen Schrift, und die Juden verstehen es gemeiniglich, daß es heiße, Verschnittener. Aber sein Ursprung, daher es kommt, ist etwas ehrlicher. Denn es kommt her vom Kammerdienste. Darum haben wir es verdeutscht, Kämmerer, denn beschlossen ist, in der Kammer aufzuwarten. Und ich halte dafür, daß dies die rechte und eigentliche Bedeutung gewesen sei; hernach aber, da die königlichen Geschlechter getheilt worden, hat man die sarisim, das ist, die Eunuchen, in Frauenzimmer der Königinnen gebraucht. Sonst wird der Verschnittenen in der Schrift nicht gedacht, denn nur Jes. 56, 3. und Matth. 19, 12. An anderen Stellen heißt es allezeit Kammerdiener oder sonst andere Diener der Könige. Wie denn in der Historie Esther Cap. 1, V. 14. auch sieben Kämmerer aufgezählt werden, so vor dem Könige Ahasveros gedient, welche dazu bestellt waren, daß sie auf des Königs Leib warteten, und seine nächsten und treuesten Diener waren.

4. Darum muß ich Lyra loben, der sagt, daß das Wort „Eunuche“ auf zweierlei Weise verstanden und gebraucht werde: einmal für die, so verschnitten gewesen, welche auf die Königinnen in ihren Kammern gewartet; zum Andern aber für die Kämmerer der Könige. Und dieselben werden gemeiniglich in der heiligen Schrift durch dies Wort verstanden, als Jes. 39, 7., da zum Könige Hiskia gesagt wird: „Dazu werden sie deine Kinder, so von dir kommen werden und du zeugen wirst, nehmen, und müssen Eunuchen Kämmerer sein im Hofe des Königs zu Babel.“ Da halte ich nicht dafür, daß dies Wort zu verstehen sei für die, so im Frauenzimmer aufwarten; sondern daß damit angezeigt werde, daß sie werden Kämmerer sein müssen im Hof oder im Saal des Königs zu Babel. Denn es will sich nicht reimen, daß Daniel und seine Gesellen verschnitten und bestellt gewesen seien für Frauendienste, sondern sie sind um den König gewesen, wie das aus Daniel zu sehen. Und in den Historien der Könige Joram und Ahab wird auch gesagt: Er hat einen von den Eunuchen gegeben, das ist, von seinen nächsten Secretären oder Kammerdienern.

\quellseite{319}

5. Der Juden Narrenwerk bin ich feind und verdamme dasselbe ganz und gar. Hieronymus erzählt aus ihren Auslegungen, daß Potiphar diesen Knaben, dieweil er hübsch von Gestalt gewesen, soll gekauft haben, seiner schändlich zu mißbrauchen. Derhalben soll ihn GOtt also gestraft haben, daß ihm sein männlich Glied verdorret, damit er des Jünglings nicht mißbrauchen könnte. Hernach aber soll er zum Hohenpriester erwählt worden sein in der Stadt On, und das derselben Tochter die Asnath sei, welche Joseph zum Eheweibe genommen, und daß man die Priester habe pflegen zu verschneiden, wie bei den Heiden den Priestern der Cybele geschehen, welche man Gallien genannt. Dies aber ist alles ganz falsch und eitel erlogen Ding. Denn es ist ein Anderer der Potiphera gewesen, welches Tochter den Joseph zur Ehe genommen; welcher den Namen von Pharao gehabt, als wenn man auf Lateinisch sagen wollte: crassities Pharaonis, die Dicke des Pharao. Die Andern, damit sie darauf antworten mögen, wie er habe können ein Eheweib haben, so er doch ein Eunuche gewesen, machen ein solch Gedicht, daß da er verschnitten worden, er gleichwohl die Sehne noch unverletzt behalten, daß er sich gleichwohl hat können zu seinem Eheweibe halten, wiewohl er unfruchtbar gewesen. Dies sind schändliche und scheußliche Lügen, die sich zwar für jüdische Säue geziemen, welche alles, was schändlich ist, von den Heiden pflegen zu erdichten, und ihr eigenes Lob über die Maßen zu rühmen. Ich habe es aber darum erzählt, auf daß der Leser wissen möchte, was ihr Verstand sei von den Eunuchen, und alsdann so viel lieber diese wahrhaftige Auslegung desselben Worts annähme, nämlich, daß dieser Potiphar einer von den nächsten und allertreuesten Kämmerdienern Pharaos gewesen sei. Und dazu kommt das auch, daß ihn der König so hoch geehrt und zum Obersten über die Fleischer gesetzt hat.

6. Alhier aber werden wir nun abermals mit den Grammatiken müssen zu schaffen haben. Hieronymus hat die hebräischen Worte, sar hattabachim, verdolmetscht: einen obersten Heerführer oder Feldherrn. Aber es ist gewiß, daß das hebräische Wort also lautet und eben so viel heißt, als, würgen und schlachten, was geopfert werden soll. Darum haben sie es auch auf Lateinisch verdolmetscht: principem mactatorum, das wir auf Deutsch gegeben haben, Hofmeister, der an Pharaos Hofe regiert, ja, nicht allein am Hofe, sondern auch im Tempel und was das Opfer belanget, daß er den Metzgern herausgegeben, was man hat schlachten sollen, es wäre gleich zum weltlichen, oder aber zum geistlichen und heiligen Gebrauch. Ein solcher Regent oder Amtmann, so ehrlicher und herrlicher ist gehalten worden als andere Hofdiener, halte ich, sei dieser Potiphar gewesen, ein Marschall oder Haushofmeister. Hernach wird noch gedacht werden, 1 Mos. 40, 2., des Obersten über die Bäcker und Schenken. Das sind sonderliche, unterschiedliche Aemter gewesen im Königreich Egypten, und mögen vielleicht also getheilt gewesen sein, daß Einer auf das Schlachten der Metzger, wozu man das Geschlachtete auch hat brauchen wollen, gesehen; ein Anderer aber die Bäcker und wieder ein Anderer die Schenken regiert hat. Und dasselbe nicht allein zu Hofe, sondern auch in einer jeglichen Stadt und im ganzen Königreiche; wie denn dasselbe sehr wohl bestellt und mit den schönsten Gesetzen geordnet gewesen.

7. Darum halte ich dafür, daß dies nicht allein auf die Hofhaltung müsse gezogen werden, welche meines Erachtens so weitläufig und herrlich nicht wird gewesen sein, wie jetzt unserer Fürsten Hofhaltung ist; denn zu derselben Zeit hat man sich sehr mäßig müssen zu halten: sondern über das ganze Land Egypten haben die Bäcker, Fleischer und Schenken ihre eigenen Regenten und Obersten gehabt, welche vermöge ihres empfangenen Ansehens nach gemeinen Rechten und Ordnungen desselben Königreichs regiert. Und dieselben Regenten haben auch die jährlichen Zinsen und den Tribut, so man dem Könige hat müssen geben, von jeglichem Stande zu Unterhaltung des Königreichs eingefordert. Also werden wir hernach hören, daß Joseph ist zum Regenten über das Korn und Getreide gesetzt worden. Und ein solcher Hofmeister des Königs zu Babel ist auch gewesen der Nebusar Adan, 2 Kön. 25, 8. ff. Das ist aber falsch und erlogen, was die Juden plaudern, daß Joseph diese Aemter der Bäcker, Schenken und Köche oder Metzger erst soll eingesetzt haben. Denn zuvor, ehe Joseph nach Egypten gekommen, ist dies schon alles dermaßen verordnet und bestellt gewesen.

8. Und bis daher haben wir von den Worten gesagt. Nun wollen wir wiederumb auf die Historie kommen, welche wahrlich wohl eines andern Auslegers oder Meisters bedürfte. Und ich zwar wollte auch gern Andern Raum geben und weichen, welche noch etwas vermöglicher an Kräften wären und auch von GOtt größere Gaben haben; aber mir wollen thun, so viel wir können, zu seiner Zeit wird mir GOtt der HErr wohl Ruhe und Feierabend geben und mir dieses Amt erlassen.

\begin{quote}
\textbf{V. 2. Und der HErr war mit Joseph, daß er ein glückseliger Mann ward; und war in seines Herrn, des Egypters, Hause.}
\end{quote}

9. Alhier hebt sich nun an das treffliche und herrliche Exempel, desgleichen man in allen andern Historien nicht beschrieben findet. Joseph ist ein sehr glückseliger Jüngling gewesen, und war auch von seinem Vater wohl unterrichtet und auferzogen. Ei! die Zucht ist wohl gerathen. Es ist nicht vergeblich gewesen, daß Jakob diesen Sohn vor den andern allen sonderlich geliebt hat; denn wir sehen, daß er durch Mitwirkung des Heiligen Geistes die Lehre seines Vaters auf das Allertreulichste und Fleißigste gehört und behalten hat. Darum leuchtet und erzeigt sich auch an ihm ein sehr brünstiger Glaube, da er in seiner ersten und blühenden Jugend nicht allein mit äußerster Trübsal, sondern auch mit höchster Wohlflust angefochten und versucht wird, und hat doch gleichwohl dadurch nicht mögen überwunden und zu Sünden getrieben werden.

10. Dies Exempel sollten sich die jungen Gesellen vorstellen und sich darnach richten. Ei! wie ist das so großer Glaube, Zucht, Keuschheit und Geduld gewesen, daß man es auch mit Worten nicht genugsam ausreden kann. Denn da er verkauft worden, ist er ein Jüngling von siebenzehn Jahren gewesen, und wird also von seinem Vater und Großvater, welche er geehrt und geliebet hat, weggerissen, und in schwere Dienstbarkeit geworfen ohne alle Hoffnung der Erledigung, oder daß er dermaßen wiederumb möchte zu den Seinen kommen. Da ist wahrlich ein fester Glaube, unglaubliche Stärke und Geduld vonnöthen gewesen.

11. Und das Wort, so sein Vater Jakob in des Knaben Herz gepflanzet, ist ganz rein und sehr fruchtbar gewesen, und hat sehr viel reiche Früchte gebracht, wie Luc. 8, 15. geschrieben steht. Denn der Andern Einer, so auch gleiches Alters mit Joseph gewesen, wäre entweder durch Ungeduld oder Schmerzen des Herzens überwunden worden, oder aber würde bald die Gelegenheit, Unzucht zu treiben, angenommen haben. Aber an Joseph siehst du der keines, sondern es leuchten an ihm so viele schöne Tugenden zugleich, als nur bei dieser unserer Natur möglich sind. Und ist, sage ich noch, sonst nirgendes ein Exempel, das mit diesem möchte verglichen werden. Derhalben haben auch die andern Brüder nicht mit solchem Fleiß oder Furcht dem Vater zugehört, wiewohl es einerlei Lehre und Wort gewesen, welches Jakob ohne Zweifel seinen ganzen Hause stets, und mit Fleiß und Treue mir vorgehalten haben.

12. Es ist ein großer Jammer, daß er dem Vater so weit aus den Augen gerissen wird, und wird freilich seines Vaters Haus nicht ohne den größten Schmerz, Wehklagen und Thränen beide sein selbst und auch des Vaters verlassen haben; GOtt schweigt dazu stille, thut, als sehe er diesen Jammer nicht; Joseph seufzt und heult, und wird geachtet, als der gar verworfen sei; da ist nichts. Warum ist er aber in seinem Herzen nicht zaghaft geworden und schreit nicht als ein Verzweifelnder: Ich sehe, daß ich von GOtt und den Menschen verlassen bin, und daß ich ohne mein Verdienst mit solchem unbilligen Unglück beschweret werde, der ich doch meinen Vater in Ehren gehalten und ihm in Allem gehorcht habe 2c.; wie denn die Verzweifelnden solche ungestüme Klage pflegen zu führen. Und daran ist kein Zweifel, es wird sein Herz auch solche harte Stöße und Stürme erlitten und ausgestanden haben. Denn er ist ein Jüngling gewesen von siebenzehn Jahren, und wir dürfen ja nicht denken, daß er ein Kloß oder Stock gewesen sei, daß er dies große Unglück gar nicht sollte gefühlt haben, da er in diese greuliche Anfechtung der Verzweiflung, Zorns und Feindschaft wider GOtt ist geworfen, dazu auch mit Unzucht angefochten worden und endlich die erbärmliche Dienstbarkeit hat tragen müssen, da er fremden und unbekannten Leuten zu dienen, wie der allergeringste und leibeigene Knecht, gezwungen wurde, war dazu von seinen leiblichen Brüdern verrathen und verkauft. Siehe, ob das nicht ein erschrecklicher und unerträglicher Jammer sei, damit unser Unglück noch lange nicht zu vergleichen ist?

\quellseite{320}

13. Also ist er nun wahrhaftig ein Erbe seines Vaters, das ist, der Trübsale so der Vater auch erlitten. Denn wiewohl Juda auch nicht ein klein Kreuz gehabt, ist es doch nichts gegen dies Kreuz zu rechnen. Joseph ist noch zumal jung, es ist zu frühe mit ihm; denn er ist noch ein Jüngling gewesen von siebzehn Jahren, ist aller der Bequemlichkeiten, Vortheile und Hülfe, deren die Jugend bedarf, beraubt gewesen, ist in ein fremdes unbekanntes Land gekommen zu gottlosen abgöttischen Leuten, da er weder Lehre noch Exempel gesehen, daraus er hätte mögen unterrichtet oder gestärkt werden; sondern er hat vor seinen Augen nur lauter Barbarei und die schändlichsten Exempel sehen müssen; inwendig aber wird er auch des Teufels List und seine feurigen Pfeile gefühlt und erfahren haben.

14. Und doch hat noch in seinem Herzen das Wort, so er von seinem Vater gehört, regiert, welches er fest behalten. Mein Vater, wird er gesagt haben, hat mich gelehrt: Es lasse mich Gott, wie lange er wolle, so will ich harren und mit Geduld seiner Hülfe warten: „Harre des Herrn, sei getrost und unverzagt, und harre des Herrn“, sagt der 27. Psalm V. 14. Denn er wird den ganzen Psalter auswendig gewußt haben; sintemal er mit der That und im Werke das alles ausgerichtet und bewiesen hat, was in den Psalmen vom Glauben, Geduld und Hoffnung, und wie man des Herrn harren soll, gelehrt wird. Er harrt und richtet sich immer auf an den göttlichen Verheißungen, so er von seinem Vater gehört, verzagt nicht, murrt auch wider Gott nicht.

15. Und mit solchen großen Anfechtungen hat er müssen allein kämpfen, hat niemand, der solches hätte mit gewußt, mit dem er sich hätte unterreden können und der ihn doch hätte können ein tröstlich Wort zusprechen. Darum ist die Gefahr, darin Joseph gewesen, viel größer und schwerer, als des Patriarchen Lot, welcher, wie St. Petrus 2. Epist. 2, 7. 8. sagt, hat unter den schändlichen Leuten zu Sodom wohnen müssen, hat von ihnen gottlose Worte hören und schändliche Exempel sehen müssen. Aber in solchen großen Bekümmernissen sind nicht weit von ihm gewesen sein Gesinde, sein Weib und Töchter; die haben ihn können helfen der vorigen Worte und Exempel gedenken; und über das hat ihn Abraham, der daselbst noch gegenwärtig gewesen, mündlich auch aufrichten, trösten und stärken können. Wenn man einen solchen Mann an der Hand hat, das ist eine große Zuflucht, Hülfe und Trost. Derhalben ist Lot selig gewesen, wenn man ihn mit diesem frommen heiligen Jüngling will vergleichen, welcher allein mitten unter Löwen, Drachen, Krokodilen und Teufeln ist gewesen. Ja, er übertrifft auch noch den Daniel.

16. Dies sind wahrlich rechte Wunderwerke. Denn da ist der Heilige Geist in Joseph, und das Wort, so in sein Herz gepflanzt, nicht also darin eingedrückt, gewurzelt und verwahrt, daß es ein unbeweglicher Fels wird wider den Teufel, beide zur Rechten und auch zur Linken. Das ist, sage ich, ein Wunderwerk über alle anderen Wunderwerke, so bisher der ganzen Welt vorgehalten wird. Es wird zwar die Historie an sich selbst einfach dahin erzählt; sie begreift aber die größten Dinge in sich.

17. Es ist aber zum andern Unglück dies noch gekommen, daß er in solchem Alter auch den Schaden der Erbsünde wird gefühlt haben, dadurch alle Menschen zur Unzucht geneigt sind. Dieselben Flammen wird er ohne Zweifel auch gefühlt haben vom siebenzehnten Jahre an bis auf das sieben und zwanzigste und wird damit solche lange Zeit gekämpft haben. Vielleicht mag er auch wohl durch viel andere Anreizungen und böse Exempel nicht allein der Frau seines Herrn, sondern auch wohl der Mägd[e] zu unziemlichem Beilager angefochten worden sein, wiewohl solches in der Historie nicht ausgedrückt wird. Denn es muß ja ein Knecht, so knechtliche Arbeit thut, es sei gleich im Hause, oder draußen auf dem Felde, mit Weibern und Mädchen, die entweder Hausgenossen oder aber Nachbarn sind, umgehen; und wir wissen wohl, wie groß der unmögliche Zorn des Satans ist, welcher immer der Jugend heimlich nachstellt, so sehr, wenn schon weder Weib noch Magd sie zur Unzucht reizt, daß es doch der Teufel unsäglicher Weise versucht, wie er sie um ihre Keuschheit bringen möge, oder dieselbige an ihnen verletzen. Das überwindet aber der Knabe Joseph alles miteinander durch ein sonderlich und göttlich Wunderwerk.

18. Nun ist er aber in solcher Knechtschaft fast bei zehen Jahren gewesen, bis daß er das siebzehn und zwanzigste Jahr seines Alters erreicht, und hat mit großem Fleiß und Treue dem ganzen Hausgesinde, Knechten und Mägd[e]n vorgestanden und dieselben regiert beide im Hause und auf dem Felde, hat sich dabei züchtig und schamhaftig beide mit Worten und Werken, ja, auch mit allen Gebärden gehalten. Denn die Gebärden sind gemeiniglich Anzeichen, wie der Mensch im Herzen gesinnet, oder wozu er geneigt ist.

19. Mose aber beschreibt die Knechtschaft Josephs und wie er dieselbe verwaltet, und sagt davon also: „Der Herr war mit Joseph“, der Herr hat es gethan. Er hat das Wort von seinem Vater Jakob aufgenommen, welches er festiglich im Herzen behalten, daraus er gelehrt worden, zu glauben und sich auf Gottes Barmherzigkeit zu verlassen. Mit diesem Glauben hat er Alles überwunden. Nicht, daß sein Fleisch darüber nicht sollte erschreckliche Stürme vieler schweren Anfechtungen gefühlet haben; sintemal es fast weich und zart ist, kann die Anfechtung nicht wohl vertragen, kann sich auch nicht wohl enthalten, daß es nicht sollte herausfahren und wider Gott toben; aber der Heilige Geist hat ihn gestärket. Darum sagt auch Mose: „Der Herr war mit ihm“, denn der Herr hat Lust, mit solchen Leuten umzugehen. Mit Ruben, Simeon und Levi, und den andern Brüdern Josephs ist er nicht gewesen: aber bei diesem ist er, welcher Gott glaubt, und ihn für wahrhaftig, gnädig und barmherzig hält.

20. Daraus auch weiter die Früchte gefolgt sind, ob sie wohl an diesem Orte nicht erzählt werden, nämlich, rechte Anrufung und Glaube, so er von seinem Vater gelernt. Und wiewohl er nicht Raum oder Zeit gehabt, in solcher schweren Dienstbarkeit zu opfern, hat er doch im Felde und sonst an allen andern Orten beten, danken und Gott loben, auch das Wort, so er gehört, in seinem Herzen wiederholen, und sich also in der Gottseligkeit, Glauben und rechtem Bekenntniß immer üben können. Denn das läßt der Heilige Geist nicht, er reizt und erweckt die Gottseligkeit immer zu solchen geistlichen Uebungen. Darum sagt hier der Text, der Herr sei mit Joseph gewesen. Wo aber der ist, da ist auch der Geist der Gnade und des Gebets, und da hört auch die Anrufung und Danksagung nicht auf.

21. Darnach hat er auch in der ganzen Hausregierung das priesterliche Amt nicht liegen lassen, wiewohl er ein geborner Priester und Erstgeborner war. Darum hat er nicht allein für sich selbst und für seinen Herrn gebetet, sondern wird ohne Zweifel auch seine Mitknechte und Mägd[e] gelehrt haben, wie man Gott anrufen und ihm recht dienen solle; hat auch an ihnen die Sünden gestraft, und wird also etliche aus ihnen zu rechter Buße und Erkenntniß Gottes besehrt haben.

22. Derhalben siehe, wie viel gute Werke er gethan habe, wie er sich in so großem Jammer und Trübsal noch immer der Gottseligkeit und Heiligkeit befliessen; und hat ihn Gott in dem Allen bewahrt, daß er nicht in Ungeduld gefallen, hat ihn auch vor des Teufels Nachstellungen gnädiglich behütet, nämlich, zur Knechtschaft vor Unzucht und Geilheit der Frauen und Mägd[e] im Haus, und auch zur Unken vor Ungeduld in der Trübsal. Diese Dinge werden nur mit dunkeln Worten an diesem Orte angezeigt; wir sehen aber, daß Alles daraus zu schließen durch eine gute und offenbar[e] Folge. Denn wo der Herr ist, da folgen auch die rechten Früchte des Worts, als, Bekenntniß, Lehre, Danksagung, und allerlei Opfer und rechter Gottesdienst. Dies sei erstlich gesagt von der Theologie und Priesterthum Josephs, das ist, wie er der Gottseligkeit und dem Bekannte obgelegen und derselben gepflegt habe.

23. Nun wird aber weiter im Texte dies auch noch hinzugefügt, daß Joseph ein glückseliger Mann gewesen sei. Hier gedenkt Mose erstlich insgemein der Glückseligkeit, hernach wird er dies auch insonderheit erklären, so viel die zeitliche Wohlfahrt belanget. Aber das soll man nicht allein auf die Glückseligkeit der Güter seines Herrn ziehen, sondern soll es auch auf Joseph selbst deuten. Denn alles, was er gethan, hat nicht allein dem Herrn, sondern auch seinen Mitknechten und ganzem Hausgesinde, ja, über das auch der Frau selbst wohlgefallen, also daß sie ihn auch lieb gewonnen. Darum,
\quellseite{321}
sage ich, haben sie ihn alle lieb und werth gehalten, die seine Frömmigkeit, Geduld und andere Tugenden mehr an ihm gesehen haben, dadurch sie auch sahen, daß er gesetzt und bewogen worden, ein ehrliches Leben zu führen, Sünde und alles Böse zu meiden. Da haben sie einmüthiglich alle miteinander alles Gute von ihm gesagt, und solchen frommen gott-seeligen Jüngling gelobt, welchem alles, was er nur angefangen, so glücklich und wohl gerathen ist, nicht allein im Hausregimente, was seines Herrn Nutzen und Frommen gewesen, sondern auch sonst in seinem ganzen Leben und allem seinem Thun, also daß er darüber zugleich von GOtt und Menschen ist geliebt worden.

24. Warum setzt aber Mose das hinzu: „Und Joseph war in seines Herrn, des Egyptens, Hause“? Antwort: Warum, daß er nicht schalkhaftig, ungetreu oder treulos gewesen, der darauf gedacht hätte, wie er seinem Herrn möchte entlaufen, oder der da sonst begehrt, jemand Leid oder Schaden zu thun; sondern ist freiwillig in seinem Amte geblieben und hat solche Knechtschaft mit geduldigem Herzen ertragen, und bei sich selbst also gedacht: HErr GOtt, du hast mich von meines Vaters Hause, Brüdern und ganzem Geschlechte weggeführt; ich weiß wohl, daß daheim mein Vater und Großvater seufzen, heulen und weinen, und groß Verlangen nach mir als einem Abwesenden haben. Was soll ich aber thun? Soll ich mich selbst von dieser Dienstbarkeit los machen und davon laufen? Nein traun, das will ich nicht thun. Warum ist er in seines Herrn, des Egyptens, Hause geblieben, das ist, er hat sich selbst GOtt aufgeopfert zu einem Opfer, und also gedacht: Dieweil ich unter fremde Gewalt verkauft bin, will ich da bleiben und mit diesem meinem Glück und Stande zufrieden sein, wiewohl es erbärmlich und verdrießlich ist, und will mich befleißigen, daß ich nicht allein den Knechten, sondern auch meinem Herrn und meinem GOtt gefallen möge, wie St. Paulus Eph. 6, 5. Col. 3, 22. lehrt, daß die Knechte ihren Herren sollen gehorsam sein, nicht mit Dienst vor Augen, sondern als vor GOtt etc.

25. Ein solcher getreuer Knecht ist Joseph gewesen, und wenn jemand solche Knechte haben möchte, die Joseph gleich wären, derselbe würde leichtlich können reich werden. Aber jetzt zu unsern Zeiten findet man derselben wenig. Denn sie sind gemeiniglich alle Diebe und Mörder, stehlen, verderben und veruntreuen alles, und sind doch frei geboren, sind nicht leibeigene oder erkaufte Knechte, wie zu derselbigen Zeit die Knechte gewesen. Und ist kein schändlicher Volk in der Haushaltung, beide in diesen Landen und sonst auch über ganz Deutschland, denn eben die Knechte sind. Und ist herrlich wohl geredt, das man gesagt hat: Ein fauler Dieb ist besser, denn ein fauler Knecht. Denn jener stiehlt langsam und wenig, dieser aber alle Tage, ja versäumt wohl alle Stunden dasjenige, so zu seines Herrn Nutzen und Frommen dienen möchte.

26. Derhalben sind sie nicht werth, daß sie Knechte heißen, sondern sollen vielmehr Diebe und Mörder heißen. Und ist wahrlich Potiphar sehr glücklich gewesen darum, daß er so einen frommen getreuen Knecht an Joseph gehabt, der sein leibeigen gewesen wie eine Kuh oder Ochse, und mit solchem seinem Stande so sein hat können zufrieden sein und sich daran genügen lassen; der sein befohlen Amt so treulich ausgerichtet, und mit nichten darauf gedacht hat, wie er sich rächen wollte, oder aber wiederum zu seinen Eltern laufen; denn „er war“, sagt Mose, „in seines Herrn Hause“; und der billig den Namen eines Knechtes nicht sollte gehabt haben, ist doch ein Knecht aller Knechte geworden. Nun folgt aber weiter im Text von seinem zeitlichen Glück und Wohlfahrt.

\section[Zweiter Theil]{Zweiter Theil. Wie Joseph wegen seines treuen Dienstes von Potiphar erhoben und geehrt wird, und wie Potiphars Haus Josephs wegen gesegnet wird}

\begin{quote}
\textbf{V. 3. 4. Und sein Herr sahe, daß der HErr mit ihm war; denn alles, was er that, da gab der HErr Glück zu durch ihn, also, daß er Gnade fand vor seinem Herrn, und sein Diener ward. Der setzte ihn über sein Haus, und alles, was er hatte, that er unter seine Hände.}
\end{quote}

27. Hier hebt nun Joseph an, sich allemahl hervorzuthun, zu wachsen und groß zu werden, und ist ohne Zweifel ein groß Glück und reicher Segen gewesen, damit GOtt alles, was er gethan, geziert hat. Sein Herr hat gesehen, daß er züchtig, mäßig, schamhaftig, ehrliebend gewesen, und daß er Alles recht und glücklich ausgerichtet und verwaltet hat; darum deutet er, er habe nun einen Diener bekommen nach seinem Willen. Das ist ein Knecht für mich, hat er gedacht. Und lauten schier die Worte also, als habe er ihn frei gelassen darum, daß er ihm so treulich und aufrichtig gedient hatte. Wiewohl ich das nicht als gewiß sagen kann. Der HErr aber hat sein Gebet erhört, und segnete erstlich den Herrn, darnach Joseph auch, also daß er über das ganze Haus, und alles, was Potiphar hat, gesetzt ward.

28. Dies ist eine große wunderbare Wohlthat gewesen, daß der Egypter alle sein Gut und Habe, und alles, was er zu schaffen gehabt auf dem Felde und im Hause, einem leibeigenen Knechte und der fremd war, befohlen und vertraut hat. Nun ist Joseph ein großer Amtmann, darum muß seine Treue und Glaube sehr groß und sehr bekannt gewesen sein, um deßwillen er auch vor seinem Herrn Gnade gefunden hat, welcher sich nun von dem Knecht seiner Gefahr oder Betrug mehr besorgt: hat auf ihn keinen bösen Argwohn, daß er ihm vielleicht etwas entführen und veruntreuen möchte, und darnach alsobald davon laufen, wie sonst Knechte zu thun pflegen, nachdem er ihm alle sein Gut und Habe in Verwahrung gethan, daß er darüber regieren sollte. Er hat sich aber ohne Zweifel vor der Untreue der andern Knechte gefürchtet; aber von Joseph hat er solches nie in keinen Sinn genommen. Denn er thut ihm freiwillig und mit vollem Vertrauen alles unter seine Hände, als ob er selbst über das Gesinde regierte und persönlich dabei wäre. Denn er hatte seine Treue wohl erfahren, daß er aufrichtig war. Darnach hat er gesehen, daß er gethan hat, was einem gott-seeligen Menschen und der GOtt fürchtet, zu thun gebühret. Nicht, daß dieser Egypter, Potiphar, GOtt den HErrn erkannt hätte, oder daß er also zu Joseph gesagt hätte: Ich sehe, daß der HErr mit dir ist; sondern Mose und die Propheten reden nach ihrer Art und Weise. Denn die Gottlosen und groben Leute halten es nicht also, sondern wenn sie einen frommen heiligen Menschen sehen, achten sie denselben, daß er werth sei, daß ihn die Götter lieb haben. Eben dasselbe hat Potiphar von Joseph auch gehalten.

29. Dies ist nun ein sehr großer, herrlicher, und zugleich auch ein sehr schwerer Stand und Regiment, so Joseph geführt. Denn wer könnte das so leichtlich leiden und ertragen, so er gelitten und ertragen hat? Er hat diese schwere erbärmliche Knechtschaft elf Jahre lang getragen, hat dabei gar nichts verdient; es sei denn, daß er darnach, da er ist frei gelassen worden, etwas noch bekommen oder erlangt hätte: dieweil er aber gedient, ist alle seine Arbeit nicht ihm, sondern nur seinem HErrn zum Nutzen und Frommen gerathen. Derselben hat er davon nichts mehr gehabt, denn nur Nahrung und Decke; und damit ist es noch schmal genug zugegangen; welches denn auch an ihm zu loben ist, nämlich, daß er solche äußerste Armuth mit Geduld getragen. Denn er hätte billig in dem Alter sollen zum Ehestande greifen, Kinder zeugen, und sein eigen Hausgesinde unterhalten und nähren. Deß allen aber hat er mangeln müssen und sieht noch dieser schweren Dienstbarkeit kein Ende; und doch hat er es gleichwohl also sein ertragen und sich mit fester Hoffnung und Glauben an göttliche Hülfe gehalten. Denn er denkt bei sich selbst also: Wenn ich schon diese Last etwas länger würde tragen müssen, will ich doch deshalb nicht kleinmüthig werden noch verzagen. Denn ich habe ja von meinem Vater gelernt, wie ich des HErrn harren, auf ihn hoffen und ihn anrufen soll. Diese Lehre will ich also behalten und mich darin üben, so lange ich lebe.

30. Wer nun diese Dinge fleißig bedenken wird, der wird leichtlich sehen, welch eine große Last es mit solcher schweren erbärmlichen Dienstbarkeit gewesen sei. Denn das ist ja schwer genug, ein Gefangener und fremder Gewalt unterworfen sein, und über das noch Armuth und schwere Arbeit tragen, dadurch du gar erschöpft, und der Andere, der ein Tyrann, darüber nur reich, groß und herrlich werde.

31. Wir klagen viel über den Türken; aber mit der Ungeduld werden wir es nicht besser machen, werden damit des Jammers nicht abkommen oder denselben überwinden. Lasset uns aber viel lieber dies Exempel Josephs ansehen und uns dasselbe vorstellen. Denn wiewohl unser Jammer und Dienstbarkeit auch sehr betrübt ist, ist es doch nicht zu vergleichen mit dem Jammer, so Joseph erlitten. Derhalben, wenn wir an die Trübsal denken, so Joseph gehabt, mögen wir uns unserer Ungeduld halben wohl schämen, die wir dem Teufel pflegen zu dienen und Gott vielfältig zu erzürnen. Denn lieber, siehe doch, was diesem Priester und Jüngling, dem Gott und Menschen günstig gewesen, begegnet und widerfahren sei.

\quellseite{322}

\begin{quote}
\textbf{V. 5. 6. Und von der Zeit an, da er ihn über sein Haus und alle seine Güter gesetzt hatte, segnete der Herr des Egypters Haus, um Josephs willen; und war eitel Segen des Herrn in allem, was er hatte, zu Hause, und zu Felde. Darum ließ er es alles unter Josephs Händen, was er hatte; und er hatte nichts davon, denn nur das Brod, das er aß.}
\end{quote}

32. Dieser Text lehrt uns, wie Gott um eines gottseligen und frommen Mannes willen auch unwürdige und gottlose Leute zu segnen pflegt; wie droben Jakob zu seinem Schwiegervater Laban sagt, 1 Mos. 30, 30.: „Du hattest wenig, ehe ich herkam; nun aber ists ausgebreitet in die Menge, und der Herr hat dich gesegnet durch meinen Fuß.“ Wie nun Jakob ein Segen im Hause Labans gewesen, also ist Joseph auch bei seinem Herrn, dem Egypter, gewesen. Das ist unsers Herrn Gottes Weise. Also liebt er die Seinen, die ihn anrufen und auf seine Hülfe und Gaben warten, daß er um Eines Menschen willen einem ganzen Lande Gutes thue und dasselbe segne. Wiewohl Joseph selbst nichts empfangen hat, sondern er bleibt ein Knecht, es sei denn, daß er vielleicht frei gegeben worden, da er über das Haus Potiphars, und über alles, was er hatte, gesetzt worden.

33. Alles aber, was der ganzen Welt Gutes widerfährt, soll man erkennen und mit dankbarem Herzen rühmen, und wissen, daß Gott um seiner Auserwählten willen vielen tausend Menschen Friede gibt und sonst auch andere Wohlthaten erzeigt. Wiewohl er doch auch der Auserwählten nicht verschont; wie er Joseph so viele Jahre hat dienen, und ihn noch endlich zum Lohne für seinen treuen Dienst mit höchster Schmach in Kerker werfen lassen. Also belohnt er die allerheiligsten Menschen, welche sehr nütze sind; ja, die der rechte Segen der Welt sind, die kriegen so schändlichen Lohn.

34. Derhalben ist fleißig zu merken, daß Gott um seiner Kirche willen die Gottlosen segnet. Denn wo die Welt ohne dieselben wäre, und wo nicht das Evangelium gelehrt würde, so wäre sie schon vorlängst vergangen. Also ist zu unserer Zeit Gott auch mit uns; denn wir haben das Wort und die heiligen Sacramente: er redet und wirkt durch uns, erlöst also viele Leute vom Tode und ewiger Verdammniß. Wir sind die Helfer und Heiland der Welt, und alles, was Deutschland Gutes hat, soll es nicht seiner großen Gewalt, nicht seiner Tugend oder Frömmigkeit, auch nicht seiner Weisheit, sondern den verworfenen, verachteten und verfluchten Menschen, welche Christen genannt werden, zuschreiben und denselben dafür danken.

35. Der Papst, Cardinäle und das ganze Geschwärm derselben bösen Buben genießt der Ehren, Gutes, Gewalt und aller Lust allein um des Evangeliums willen, und dasselbe müssen unsere Leute auch bekennen. Und geht doch der Papst mit den Seinen ihres höchsten Vermögens nur allein damit um, daß sie uns dafür recht danken, nämlich, mit Schwert, Feuer, Hunger, und allerlei andern Strafen und Plagen, wie sie solche nur erfinden mögen. Ja, siehe auch unsere eigenen Zuhörer an. Die Bauern und der Adel, Bürger und Obrigkeit plündern und berauben diejenigen, welche sie ja billig nähren, denen sie helfen, die sie unterhalten und denen sie zulegen sollten, daß also die Kirche möchte erhalten werden, daran ja alle ihr Heil und Wohlfahrt gelegen ist. Es muß so sein, da mir nichts aus.

36. Im Papstthum war kein Maß und es hatte kein Ende mit Geben und Zulegen, den ungöttlichen falschen Gottesdienst damit zu erhalten. Daher sind gekommen die Stifte und Klöster, die fürstlichen Bisthümer, die gewaltigen Reiche der Cardinäle und Päpste; da hat der Teufel können zublasen wider Christum und die Leute zu solcher Wildigkeit und ungöttlichen Gaben antreiben. Denn solches hat ihm wohl geschmeckt, und die Leute dienten damit der Abgötterei und Gotteslästerung, so zu der Zeit öffentlich gelehrt worden. Darum haben die Dompfaffen, Bischöfe und Aebte Alles vollauf gehabt, so sie doch ein gemeiner Schaden und Verderben der ganzen Welt gewesen sind. Sie sind dick, stark und fett geworden: wir aber, so das Wort treulich und mit großem Fleiß lauter und rein lehren, und damit nichts Anderes suchen, denn allein die Wohlfahrt und Seligkeit des deutschen Landes, müssen dagegen den bittersten Haß und Neid tragen von allen Ständen, und könnten sie uns den Bissen Brod nehmen, den wir essen, sie thäten es. Ei, warum das? Antwort: Darum, daß es also muß zugehen. Die Welt muß der Teufel sein und soll die Ihren ehren, wie Christus Joh. 15, 19. sagt: „Wäret ihr von der Welt, so hätte die Welt das Ihre lieb.“ Jenen losen Haufen aber der bösen schändlichen Menschen nährt und ehrt sie gern. Denn sie sind vom Teufel und von der Welt. Wenn etwa ein frommer Joseph kommt, demselben gibt sie ungern einen Brosamen trockenen Brodes.

37. Wo aber irgend ein Fürst, der mit seiner Mildigkeit die Kirche nährt und unterhält, desselben Herz ist ohne Zweifel vom Heiligen Geist gerührt worden; die andern fast alle sind nur räuberisch, daß sie Alles an sich reißen. Siehe doch die Pfarrherren auf den Dörfern hin und wieder an, wie ihrer so viel erbärmlicher Weise vor Hunger und Durst schier verschmachten müssen, haben oft nicht, daß sie ihren Kindlein ein Hemd kaufen. Der Adel reißt die Kirchengüter hin, nicht allein der Klöster, sondern auch der Pfarren, deren sie sich doch billig enthalten sollten, daß davon die Pfarren und Kirchendiener genährt und unterhalten werden möchten. Darum wird es bald dazu kommen, daß die Pfarrherren werden müssen Hungers sterben. Es ist kein Ernst, kein Fleiß, kein Herz da zur Gottseligkeit; denn es nimmt sich derselben niemand mit Ernst und von Herzen an.

38. Wir fordern und nehmen auch billig von unsern Zuhörern unsern gebührlichen Lohn, wie Christus Luc. 10, 7. sagt: „Ein Arbeiter ist seines Lohnes werth“, das ist, er kann denselben mit gutem Gewissen annehmen; item 1 Cor. 9, 14.: „Also hat auch der Herr befohlen, daß, die das Evangelium verkündigen, sollen sich auch vom Evangelio nähren“; item 1 Tim. 5, 17.: „Die Aeltesten, die wohl vorstehen, die halte man zweifacher Ehre werth.“ Ja, wahrlich sind sie zweifacher Ehre werth. Aber wo das? Antwort: Bei Gott; bei der Welt aber achtet man sie, als die des Schwerts, des Galgens, der Hölle oder wo noch etwas Betrübteres wäre, werth seien.

39. Was fragen wir denn darnach? Wir, so der undankbaren Welt dienen, haben die Verheißung und Hoffnung des Himmelreichs, und wird die Erstattung und Vergeltung dieses unseres Jammers so groß sein, daß wir auch uns sehr schelten werden, daß wir um solcher Verachtung und Undankbarkeit willen der Welt uns jemals eine Thräne oder Seufzer haben entfallen lassen. Warum, werden wir sagen, haben wir nicht noch etwas Schwereres erlitten? Hätte ich doch nimmermehr geglaubt, daß eine solche große Herrlichkeit im ewigen Leben sein würde; denn sonst wollte ich davor keine Scheu getragen haben, wenn ich auch noch viel mehr hätte leiden sollen.

40. Dies ist der Trost und Hoffnung der Gottseligen, nämlich, daß unsere Seligkeit und Herrlichkeit nicht in diesem Leben gesucht wird, sondern wir sehnen uns nach einem andern und bessern Leben. Alles aber, was wir in dieser Welt thun, das richten wir miteinander nur zu Gottes Ehre, damit viel Leute mögen bekehrt und selig werden. Darum essen wir auch und trinken, thun die Werke, so beide zum Haus- und Weltregiment gehören, desgleichen sammeln wir eine Kirche mit Lehren und Predigen, Alles darum, daß wir nur ein ander Leben suchen mögen. Wie Joseph, welcher seine Hoffnung gehabt auf Freiheit oder Erledigung, oder auch auf das Regiment in Egypten, doch aber auf die ewige Glorie und Seligkeit im Himmelreich gehofft und gewartet hat. Eben darauf wollen wir auch hoffen; wiewohl Gott auch in diesem Leben die Gottseligen mit vielen Gütern überschüttet. Das ist aber ihr vornehmster Lohn nicht, sondern das ewige Leben und ewige Freude ist endlich der Gottseligen rechter Lohn im Himmel, welchen auch Christus und alle Heiligen empfangen haben.

41. Derhalben mag die Welt immer hinfahren mit ihren räuberischen Geizhälsen, da sie hingehört. Lasset aber uns unser befohlen Amt ausrichten; wie wir ja von Gottes Gnaden noch immer lehren und dabei mit höchster Geduld große Trübsal leiden. Denn wir gönnen ja in allewege den gottlosen und undankbaren Leuten noch immer gerne Gutes, damit ihnen möchte gethan und geholfen werden, und sind damit zufrieden, daß wir wissen, daß solcher unser Fleiß und Arbeit Gott wohlgefällt. Was wir nicht können fortbringen, das mögen wir dahinten lassen, wenn wir nur einen günstigen und gnädigen Gott haben, der da sagt: Siehe, liebes Kind, dein Dienst gefällt mir wohl. Ich will dein Lohn sein, ich will dir es reichlich bezahlen, daß du den undankbaren Menschen gedienet hast. Diese Verheißung und Trost ist viel reicher, denn alle Schätze der ganzen Welt sein mögen.

\quellseite{323}

42. Das soll man zum Ersten an dieser Stelle lernen, daß die Welt sich dessen mit nichten rühmen kann, daß sie dieses zeitlichen Lebens werth sei, auch nur einen Augenblick; sondern daß Gott um des Evangeliums, um der Taufe und Vergebung der Sünden willen Alles gibt, und das auch den allerschändlichsten und ärgsten Feinden der Kirche, als Türken und Papisten; er wirft es in die \unsicher{Napuse}.

43. Und das sei nun gesagt von dem Segen, so um des gottseeligen frommen Joseph willen im Hause der Gottlosen gewesen. Denn Potiphar, der Egypter, wird groß und reich, nicht durch seinen, sondern durch Josephs Fleiß und Arbeit, da er dieselbe gute Tage gehabt und unterdessen im Müßiggange geschmachtet hat. Wie auch Laban 1 Mos. 30. die ganze Last seiner Haushaltung dem Jacob allein auferleget hat, also daß er auch das von ihm forderte, so ihm etwa durch Diebstahl weggekommen war. Also hat der Herr dem gottlosen Könige in Syrien auch Glück und Heil gegeben um Naemans willen, 2 Kön. 5, 1. Und hätte auch unser Gebet dem Wort gethan, Deutschland sollte wohl anders stehen, es würde darin gar übel zugehen. Unser Evangelium ist der rechte Naeman und Joseph. Solches erkennen zwar die Widersacher nicht, glauben es auch nicht; lasset aber uns es erkennen und Gott dafür dankbar sein.

44. Darnach ist das auch zu merken, daß der Segen des Herrn gleichwohl denen nicht widerfährt, die da schlafen und faul sind, wie man sagt von Polykrates, der ein Heide gewesen, daß er im Schlaf das Netz gezogen habe, und im 127. Psalm N. 2. geschrieben steht: „Der Herr gibts seinen Freunden schlafend.“ Denn da hat David die Arbeit nicht verboten, sondern will haben, daß das Gewissen soll stille und friedsam sein: es soll nicht eine solche Arbeit sein, so den geistlichen Schlaf des Glaubens und Hoffnung hindern sollte; sondern es soll zugleich auch mit dabei sein ein stilles Herz und ein gut friedsam Gewissen, wie Christus sagt Matth. 6, 25. 34.: „Sorget nicht für euer Leben“ etc., seid guter Dinge: „Es ist genug, daß ein jeglicher Tag seine eigene Plage habe.“ Machet euch nicht selbst Unglück. V. 31.: Schlafet, und saget nicht: „Was werden wir essen, was werden wir trinken?“ Ich will aber doch nicht haben, daß ihr sollt faule Brüder sein und müßig gehen, sondern ihr sollt arbeiten. Denn Gott gibt seinen Segen denen nicht, die nach dem Fleische liegen und schnarchen und faul sind. Er will, daß das Fleisch immer geübt und getödtet werden soll, wie St. Paulus sagt Gal. 5, 24.: „Welche Christo angehören, die kreuzigen ihr Fleisch sammt den Lüsten und Begierden.“

45. Wiewohl aber nun die Arbeit das nicht ausrichtet, daß der Segen gibt und mitbringt, will doch Gott haben, daß wir unser befohlen Amt ausrichten und fleißig arbeiten sollen, damit also das Fleisch immer geübt werden möge, auf daß es nicht schnarche und gar faul werde; wie der heidnische Poet sagt, daß alle Menschen ihrer Art nach zu guten faulen Tagen geneigt seien und daß sie nur der Arbeit mögen überheben sein. Darum sollen wir uns nicht zum Müßiggehen und Faulenzen begeben, sondern sollen die Arbeit mit allem Fleiß und Treue auf uns nehmen, wie dieselbe eines Jeden Beruf mitbringt und fordert: und dabei sollen wir auch noch in aller Mühseligkeit, Arbeit und Verfolgung ein fröhlich Herz haben und also des Segens vom Herrn erwarten. Denn also sind Jakob, Joseph, Rahel und Lea auch nicht faul gewesen. Lea sagt zu Jakob 1 Mos. 31, 15.: „Hat nicht unser Vater Laban unsern Lohn verzehret?“ Als wollte sie sagen: Wir haben müssen die Mühe und Arbeit im Hause tragen, nämlich, des Viehes hüten, Kühe und Ziegen melken etc. Was aber nun mit solchem unserm Fleiß erworben ist, das alles hat Laban, unser Vater, verzehret. Und sind ohne Zweifel Rebecca und Sarah auch so fleißig in ihrer Arbeit gewesen; aber doch mit gutem Frieden und Ruhe ihres Gewissens.

46. Auf diese Weise ist nun Joseph sehr fleißig in der Regierung über das Hausgesinde gewesen; ist der Erste gewesen, so des Morgens aufgestanden, hat das ganze Haus allenthalben besichtigt, hat verordnet, was beide im Hause und auf dem Felde hat verrichtet werden sollen, und Alles insonderheit fleißig ausgetheilt und bestellt, hat so lange nicht gewartet, bis das Gesinde von selbst thäte, was ihm gebührete. Denn das pflegt schier nimmer zu geschehen. Wie das an unsern Leuten wohl zu sehen ist, die so überaus faul und nachlässig sind, daß ich sie an keinen Ort träger gesehen habe, als in diesem Lande und in dieser Stadt. Wenn man nur bis morgen zu essen hat, thun sie gar nichts, oder aber richten doch ihre Arbeit, so ihnen gebühret, untreulich aus; darum sind es auch lose Leute und arme Bettler, und alles, was sie haben, das haben sie vom Gebet der christlichen Kirche, können nimmermehr reich werden, dieweil sie so gar zum Faulenzen, Schlaf, Fressen und Saufen geneigt sind. Es ist ein faul, schändlich, träges Volk etc. Die aber emsig und fleißig sind in der Arbeit, können leichtlich reich werden. Denn solchen Leuten will Gott den Segen geben, wie Salomo sagt Spr. 10, 4.: „Der Fleißigen Hand macht reich.“

47. Und wir sehen, daß es eben auch also an der Fürsten Höfen zugeht, da gemeiniglich Zwei oder Drei die ganze Last im Regimente tragen, als etwa der Fürst selbst, oder etliche von den Schreibern; der andere Haufe liegt und schnarcht Tag und Nacht, frißt und säuft; oder wo sie ja etwas thun, thun sie doch dasselbe nicht um ihres Amtes, sondern nur um Geizes und eitler Ehre willen: darum gehen auch so große Unkosten auf die Hofhaltung.

48. Joseph aber ist ein sehr fleißiger Hausvater gewesen, der selbst fleißig auf alle Arbeit des Gesindes gesehen hat, darob Rechenschaft gefordert: welches Aufsehen auch mehr nützt, als die Handarbeit thun kann; wiewohl auch dazu großer Fleiß gehört und daß man damit immer anhalte. Derhalben ist Joseph nicht allein fromm und züchtig gewesen, und hat auch sehr fleißig für seinen Herrn, für den König und für das ganze Egyptenland zu Gott gebetet, sondern hat auch sehr fleißig gefördert und verwaltet allerlei Hausarbeit, da er auch Alles aufs fleißigste nach der Zahl, nach dem Gewichte und Maß gerichtet und abgerechnet hat. Welches denn auch eigentlich das Amt ist eines guten Hausvaters; wie dasselbe auch die gemeinen Sprüche bezeugen, da man sagt: Das Auge des Herrn, und der Mist, so von seinen Füßen fällt, macht den Acker fett und fruchtbar. Item: Frons occipito prior est, das ist: Die Stirne ist näher als der Hintertheil des Haupts. Denn wo nicht ein rechter Joseph ist, da wird nichts recht ausgerichtet.

49. Darum, sage ich, ist vonnöthen, daß man immer wache und fleißig sei. Wie denn die Alten das Scepter sein abgemalt haben. Oben an der Spitze, da man jetzt eine kleine Krone hinsetzt, hat man ein wachend und offen Auge gegeben, damit anzuzeigen, daß zur Verwaltung der Polizei oder Regiments durchaus vonnöthen sei, auf alle Dinge ein fleißiges Aufsehen zu haben. Wenn ein Fürst nicht wacher ist, so wird nichts daraus. Alexander der Große hätte nimmermehr so große Dinge in so kurzer Zeit können ausrichten, wo er nicht selbst in eigener Person bei allem Dinge mit sonderlichem Fleiß und Sorgen gewesen wäre. Die Regenten müssen traun oben im Scepter ein Auge tragen. Und was die Haushaltung belanget, pflegt man auch von der Gegenwart der Hausmutter also zu sagen: Der Frauen Augen kochen wohl; der Magd Augen nimmermehr.

50. Das ist nun die andere Lehre an diesem Ort, nämlich, daß der Segen Gottes denen, so ihn lieben, und dabei wachsam sind und gerne arbeiten, gegeben werde. Denn wiewohl der Fleiß an sich selbst allein nichts taugt, will ihn doch Gott gleichwohl haben, auf daß also das Fleisch zerknirscht und gekreuziget werde; welches denn fromme gottelige Leute thun um Gottes willen. Darum sind sie auch das Heil und der Segen der ganzen Welt. Wie Joseph allhier beschrieben wird als das Bild eines rechtschaffenen und vollkommenen Menschen, und der mit allerlei Tugenden geziert gewesen; und da solches sein Herr gesehen, hat er ihm auch alles, was er hatte, unter seinen Händen gelassen, wie im Text steht.

51. Das Stücklein aber, so im Text hernach folgt: „Und er hatte nichts davon, denn nur das Brod, das er aß“, ist auf mancherlei Weise durch das verdrießliche Gemäsch der Juden behandelt worden. Etliche legen es also aus: Joseph habe nichts denn trocken Brod gegessen, und sei doch durch sonderlichen Segen Gottes schön von Gestalt und auch ganz gesund dabei gewesen. Die Andern wollen es verstehen von Potiphar, nämlich, daß derselbe solchem getreuen Diener Alles hat lassen befohlen sein, also daß er sich um nichts Anderes bekümmert, denn allein, daß er daheim im Hause gegessen und darnach bald wieder an den Hof gegangen sei. Noch Andere verstehen durch das Brod das Weib, also daß der Herr in seinem Hause nichts gekannt habe denn nur sein Weib; und auf dieselbe Weise legen sie auch aus die Worte in den Sprüchen Salomos am 5. Cap. V. 15.: „Trinke Wasser aus deiner Grube“ etc. Aber die Juden sind werth, daß sie solchen Lügen glauben und also lehren.

\quellseite{324}

52. Ferner, so heißt das hebräische Wort lechem nicht allein Brod, sondern insgemein alle Speise, als, Brod und Wasser, Essen und Trinken; wie Dan. 4, 9. ff. steht, daß der König im Gesichte einen Baum gesehen, davon Alles zu essen hatte etc., das ist, der Baum hatte Aepfel und andere Früchte getragen, davon alle Thiere gegessen haben; und Jer. 11, 19.: „Laßt uns den Baum mit seiner Speise verderben“, das ist, mit den Früchten.

53. Derhalben halte ich es dafür, daß diese Worte auf Joseph müssen gezogen und gedeutet werden, daß derselbe allerlei Mühe und Arbeit im ganzen Hause getragen, und daß durch ihn der Herr den Potiphar gesegnet: wie droben Laban auch zugenommen hat durch seines Schwiegersohns und seiner Töchter Arbeit, und hat ihnen doch, da sie von ihm gezogen, für solchen ihren schweren und treuen Dienst, so sie ihm so viele Jahre geleistet, auch nicht einen Heller oder Pfennig zu Lohne gegeben.

54. Also hat Joseph von seinem unbarmherzigen Herrn auch seinen Lohn bekommen. Denn die Welt begehrt wohl sehr des Segens und der Arbeit frommer heiliger Leute und hat dazu Lust; wie die Drohnen, so nicht arbeiten, gleichwohl des Honigs auch mit genießen, welchen die andern rechten Bienen machen und daran sie sich ergötzen. Wie heutiges Tages die Fürsten und Tyrannen großen Nutzen vom Evangelium haben, welches jetzt gelehrt und immer mehr und mehr ausgebreitet wird. Denn sie werden reich, entweder von vielen großen Geschenken der Bischöfe und Cardinalen, oder aber vom Raube der Kirchengüter, die sie an sich reißen; ja, endlich brauchen sie auch mit höchster Wohllust des Friedens und guter stillen Ruhe, Gutes und allerlei Wohlfahrt; welches alles die christliche Kirche allein mit ihrem Gebet von Gott erlangt. Aber den armen bedrängten Kirchendienern geben sie gar nichts. Wenn man einem armen Pfarrherrn soll ein Hemd geben, so ist niemand dabei. Eine solche grausame Unfreundlichkeit und Geiz ist auch an Laban und diesem Potiphar, dem Egypter, gewesen, welcher hernach die Arbeit seines frommen treuen Dieners mit Banden und Kerker belohnen wird.

55. Dies Glück haben alle frommen gottseligen Kinder in diesem Leben, nämlich, daß sie arbeiten und der undankbaren Welt den Segen erlangen, und dabei mit dem täglichen Brod sich genügen lassen. Inzwischen aber können sie sich an diesem Trost aufrichten, daß sie wissen, ihre rechten Güter sind im Himmel, da auch ihr Wandel ist, Phil. 3, 20.; denn sie sind nicht von dieser Welt, sondern sie sind vom Himmel, von dannen sie auch des Heilandes warten, welcher ihnen den rechten Lohn geben wird; lassen sich dieweil genügen an der \unsicher{Parteten}, daß sie Nahrung und Kleider haben, 1 Tim. 6, 8., welches ihnen gleichwohl unser Herr Gott immer verordnet und gibt, aber doch also, daß bisweilen Viele darüber noch Hunger und Kummer leiden müssen.

56. So ist nun das der Verstand dieser Worte: So lange Joseph bei diesem gottlosen Herrn gewesen, hat er gar nichts Eigenes gehabt oder gehabt, ist nicht räuberisch oder geizig gewesen, hat auch nichts begehrt, denn nur sein tägliches Brod, das er gegessen hat, das Andere hat er alles Gott befohlen. Darum siehst du hier wohl, welch ein greulich Thier die Welt ist, welcher Gerechtigkeit diese ist, daß sie nur anderer frommer Leute Arbeit frißt und verzehrt; eines frommen Knechts nicht anders, denn wie eines Esels oder Ochsen gebraucht, und läßt ihm dagegen kümmerlich seine Nahrung und Kleider zur äußersten Nothdurft. Denn das heißt „tägliches Brod“, Leibesnothdurft, was zu diesem Leben und vonnöthen ist, nicht zum großen Gute, Reiche, Hoffahrt oder Geize, sondern nur zur Nothdurft. Und dasselbe hat ihm noch dieser gottlose Herr kümmerlich und häßlich gegeben. Er hat sich davon nicht tödtlich kleiden oder wohl leben können. Denn wer da arbeitet, dem gibt man die Spreu und nur die übrigen Brocken, die Andern aber leben von seiner Arbeit, ob er wohl selbst schwerlich so viel hat, davon er leben möge.

57. Darum hat Joseph eine große schwere Last auf sich liegen gehabt, da ihm, wie Mose sagt, sein Herr alles, was er gehabt, unter seine Hände gethan hat. Dem Potiphar aber ist das zumal ein fein sanft Leben gewesen, da er von Anderer oder fremder Arbeit, Sorge und Mühselligkeit seine Nahrung gehabt, wie die Drohnen sich auch vom Honige der andern rechten Bienen nähren. Denn er wird also gedacht haben: Siehe, dies Knechtlein habe ich um ein gering Geld gekauft; ich sehe aber, daß er mit einer solchen Geschicklichkeit begabt ist, daß mir selbst nicht vonnöthen ist, auf die Arbeit meines Hausgesindes zu sehen. Denn dieser Knecht richtet von selbst Alles aus und thut solches mit großem Fleiß. Eben dasselbe wird hernach Pharao auch thun.

58. Solcher Drohnen, das ist, solcher faulen Brüder, so nur von anderer Leute Arbeit zu leben begehren, ist die ganze Welt voll in allen Ständen. Wer zur Arbeit geschickt und willig ist, wird fast bei jedermann in großem Werth gehalten, und man begehrt sein vornehmlich darum, daß die Andern ihre Last und Beschwerung, so ihnen Amtshalben obliegen, es sei gleich zu Hofe oder daheim in eines Jeden Hause, von sich abwenden mögen und ihm auf seine Achseln legen, und sie sind gleichwohl darauf fast sehr beflissen, daß sie nur seine Arbeit am meisten genießen mögen, so doch das ein ausdrückliches Gebot Gottes ist: „So jemand nicht will arbeiten, der soll auch nicht essen“, 2 Thess. 3, 10. Ja, es sind ihrer noch viele, die sich stellen, als ob sie von nirgendes was wüßten, oder sie begehren auch gar nichts zu wissen, allein darum, daß sie ja nicht auch zu gemeiner gewöhnlicher Arbeit dieses Lebens gezogen werden. Wie ich noch wohl gedenke eines Mönchs zu Erfurt, da er von Andern gestraft ward, daß er immer vor der Hausarbeit, so im Kloster ausgerichtet war, flöhe, so er doch groß und stark war vom Leibe, daß derselbe darauf also geantwortet: Er wäre zu solcher Arbeit nicht geschickt, müßte auch wohl, wo er einmal die Hände an die Arbeit legen würde, so müßte er es darnach oftmals mehr thun. Denn wenn man merkt, daß Einer ein Ding thun kann, so läßt man ihn dasselbe gern thun.

59. So ist nun Joseph ein Jüngling, der alle Tugenden an sich gehabt, und was an ihm gewesen, ist alles herrlich und löblich. Er läßt sich genügen an dem, was ihm bescheret ist, leidet mit großmüthigem Herzen Armuth, hat dabei Geduld und Glauben zu Gott, und ist milde und gütig gegen seinen tyrannischen Herrn Potiphar in dem, daß er das Haus mit großem Fleiß und Treue regiert, und selbst arm und nothdürftig dabei ist. Derselben ist er der allertheuerste Schatz des ganzen Königreichs Egypten gewesen, wiewohl er von niemand dafür erkannt worden ist, und hat dazu mal in der ganzen Welt keinen seines Gleichen gehabt, ausgenommen seinen Vater Jacob und seinen Großvater Isaak, und war doch derselbe große Schatz vor der Welt unbekannt und gar verborgen. Nun hätte zwar dies herrliche Licht der schönen Tugenden nicht also können hervorscheinen, wo er nicht einen solchen geizigen Herrn gehabt, welcher mit seiner Grausamkeit weckt, daß er sich in dem schweren Dienst, so er gehabt, so geduldig und sanftmüthig in seinem Herzen erzeigt.

60. Diese Tugenden sind fürwahr wunderbarlich, und übertreffen weit die Wunderwerke und Geschichten, so er bei Pharao ausgerichtet, welche hernach werden erzählt werden. Jene scheinen wohl äußerlich vor Augen, als ob sie herrlicher seien: diese aber sind verdeckt und verborgen, sind wie das Wasser Siloah, das stille geht, und nicht wie die starken und vielen Wasser des Stroms der Herrlichkeit und Hoffahrt der Assyrer, Jes. 8, 6. 7. Darum sieht oder verstehst sie auch niemand, so es doch ganz herrliche und sehr große Wunderwerke sind, welche Joseph alle Tage thut. Potiphar sieht sie wohl, aber mit geizigem unbarmherzigen Herzen wie Laban. Da liegt der köstlichste Edelstein in Finsterniß verborgen, ist mit der Knechtschaft wie mit Mist oder anderm Kothe gar überschüttet, daß ihn niemand sehen kann.

61. Die Werke, so hernach folgen, werden etwas mehr ins Auge fallen, da er die Träume deutet und Erinnerung thun wird von der zukünftigen Theurung, und das ganze Königreich mit Getreide versorgen wird, welches die sieben fruchtbaren Jahre über behalten und verwahrt worden ist. Aber in Wahrheit sind diese Werke noch etwas größer. Darum ist es herrlich geredet, das Augustinus sagt, daß die Wunderwerke, so täglich in der Welt geschehen, größer seien, als die von Christo, da er noch auf Erden gelebt, geschehen sind.

\quellseite{325}

62. Das ist ein groß und wunderbares Werk gewesen, daß Christus mit sieben Broden vier tausend Mann gespeist hat: aber er speist noch alle Tage die ganze Welt. Deß würden wir uns billig viel mehr verwundern, wo es nicht bei uns in Verachtung genommen wäre, darum, daß es täglich geschieht und so gar gemein ist. Derhalben sind die Wunderwerke, so man vor Augen sieht, viel geringer, als die, davon man hört, daß sie geschehen sind. Die Werke, so Joseph gethan, da er dem Potiphar gedient hat, sind größer, als die, so er gethan hat, da er das Regiment gehabt; wiewohl dieselben den Leuten mehr zu Herzen gehen und sie auch mehr bewegen, als jene verborgenen Werke thun. Denn hier sieht man an Joseph seine sonderliche große Gestalt oder Herrlichkeit. Er geht schlecht einher, und hat die Last des Hauses regiments auf sich liegen, ist dabei arm und gar verachtet, bis hernach die große Ehre und Herrlichkeit seiner Wunderwerke hervorbricht; welche doch alle miteinander mit nichten zu vergleichen sind mit dem, daß Joseph ein Kind GOttes ist, dazu ein Priester voll Heiliges Geistes, Gottseligkeit und der höchsten Tugenden, um deßwillen das ganze Haus Potiphars so reichlich gesegnet und groß gemacht worden.

63. Darum sagt Augustinus: GOtt habe sich selbst etliche Wunderwerke, so etwas kleiner und seltsamer sind, vorbehalten, daß er dadurch, dieweil sie also seltsam sind, die Welt der größeren erinnern möchte. Es ist jeder Mensch an sich selbst ein groß Wunderwerk, und wenn auch die ganze Welt und die Engel selbst alle ihr Vermögen zusammen thun würden, könnten sie auch nicht ein Härlein am Menschen machen. Also sehen wir, daß der Himmel beweget wird, daß die Sonne leuchtet: aber dadurch werden wir nicht sonderlich bewegt. Wenn aber aus GOttes Wort gelehrt wird, daß diese sehr schönen himmlischen Körper aus Nichts geschaffen sind, alsdann können wir etlichermaßen abnehmen, wie groß das wunderbare Gebäude sein möge, welches man mit Augen nicht begreifen kann, sondern das allein mit dem Gehör, wo es durch GOttes Wort gezeiget wird, erkannt und begriffen werden kann.

64. Ein solch Wunderwerk ist Joseph bei Potiphar in Egypten auch gewesen. Darum ist dies ein herrlich Ding, daß er hier also beschrieben wird. Und ist solches durch sonderlichen Rath des Heiligen Geistes der ganzen Kirche, ja, der ganzen Welt zum Exempel vorgestellt. Denn es steht sonst in der heiligen Schrift kein ander Exempel mehr, daß diesem möchte gleich sein. Das Leben der andern Väter beide vor und nach der Sintfluth ist nicht so gar fleißig und genau beschrieben worden, daß darin zugleich, wie sie sich in der Jugend und jungen Jahren befliessen und gethan haben, erzählt würde, ohne allein, daß gemeiniglich die Zeit ihrer Geburt angezeigt wird. Aber weiß sie sich vornehmlich in der Jugend befliessen haben und wozu sie am meisten geneigt gewesen, dazu, was für Gefahr, Kümmerniß, oder Trübsal sie erlitten, solches wird anderweit nirgends also erzählt als in dieser Historie Josephs allein. Seine Brüder sind zwar auch abgemalt und beschrieben worden als sehr böse Buben und mit unsäglichen großen Sünden befleckt: aber dieser Joseph allein scheint hervor, wie eine schöne Rose unter den Dornen, der mit allerlei Tugenden geziert gewesen, als, mit Geduld und Beständigkeit über der rechten Lehre der Gottseligkeit, wie er dieselbe von den Vätern empfangen, fest zu halten, dazu mit Keuschheit, Zucht und Mäßigkeit. Denn er ist zu der Zeit hingerissen worden, daß er mit seinem Vater keine Gemeinschaft hat haben können, ja, daß sie auch einander nicht gesehen, da er noch der Lehre und Vermahnung seines Vaters am allermeisten bedurft hätte. Wie zu unsern Zeiten der Türke die Kinder hinreißt in dem Alter, da sie leichtlich die Lehre, so sie erstlich gehört, vergessen, und zu türkischer Lehre gereizt und gewöhnt werden mögen.

65. Denn die zarte Jugend, wiewohl es ihr an Zuchtmeistern und Andern, so auf sie warten und sie vermahnen, nicht mangelt, ist doch gleichwohl, wie der Poet sagt, wie Wachs, daß sie sich leichtlich zu Sünden neigen läßt, ist frevelhaft, widerspenstig und ungehorsam, und die Erbsünde fängt alsdann auch an sich in ihnen zu regen und zu wüthen, mit welcher Sünde ihr der Teufel heimlich nachstellt, und die arme Jugend mit den Stricken seiner Schmeicheleien und Wollüste des Lebens zu fangen sucht. Derhalben ist es ein großer Jammer und große Gefahr gewesen, daß dieser Jüngling, so mit rechter Lehre und guten Exempeln der Zucht und Disciplin noch nicht vollkommen befestigt, gleichsam mitten aus der besten und heiligsten Lehre der Unterweisung hingerissen und von seinen Brüdern in die Hände der Gottlosen überantwortet wird, welche ihn in ein unbekannt Land, so voller Abgötterei gewesen, weggeführt, da er Teufelslehre hat müssen sehen und hören, so der Lehre, welche er von seinem Vater gehört, ganz und gar zuwider gewesen. Dergleichen hat er die schändlichsten Sitten und Leben der Heiden ansehen müssen; und hat doch dabei mit großer wunderbarer Beständigkeit die reine Lehre und Gottesdienst noch behalten, wie er es von den Vätern empfangen hat.

66. Darum rühmt der 81. Psalm V. 6. 7. Joseph wunderbarer Weise, da er sagt: Joseph habe in Egypten fremde Sprachen gehört und daß seine Hände der Töpfen sind los geworden. Er hat müssen die unbekannte Sprache der Egypter hören, und nicht allein die, so ihre Muttersprache war, sondern auch fremde Lehre und Gottesdienst, die sich mit der Lehre seines Vaters nicht gereimt hat: bei welchen fremden Gottesdienst er noch gleichwohl mit größtem Glauben und Gottseligkeit die Lehre seines Vaters bewahrt, und weder durch die heiligen noch weltlichen Reden der Egypter verführt worden ist. Da hat wahrlich der Heilige Geist sein müssen, der in ihm den Glauben und rechten Verstand der Gottseligkeit gestärkt und bestätigt hat, damit das zarte schwache Herz nicht weich würde, und sich den vermeinten Gottesdienst und die große Herrlichkeit der gottlosen Lehrer nicht bewegen ließe, von rechter Gottseligkeit abzufallen. Er wird oftmals, wo er entweder den Gebrauch ihres Gottesdienstes oder aber ihr verderbtes heidnisches Leben gesehen hat, gedacht haben an die heiligste Disciplin oder Zucht und Lehre seines Vaters, und wird sich damit bewahrt und beschützt haben wider allerlei Aergerniß. Das muß ja ein fein Kind gewesen sein. Der hat fleißig und wohl studirt; da muß auch der Heilige Geist gewesen sein.

67. Jetzt, zu unserer Zeit, da, leider! die Jugend so gar muthwillig und frech ist, können wir noch schwerlich die guten Gesetze von Zucht und Ehrbarkeit erhalten, wenn man auch schon die besten Zuchtmeister dabei hat. Aber dieser Jüngling von siebzehn Jahren hat mit so großen Fleiß und Glauben über der Lehre seines Vaters gehalten, daß er auch mit keinerlei Exempel oder anderer Lehre von der rechten wahren Gottseligkeit hat können abgeführt werden; wiewohl er nicht ohne großen Schmerz seines Herzens so viel Aergerniß und Beschwerungen außerhalb seiner Kirche oder Gemeine wird ertragen haben.

68. Er hatte auch dies Kreuz mit seinen Sünden nicht verdient; sondern hat gelitten, wie die Christen pflegen zu leiden, die um der Wahrheit und Gerechtigkeit, und nicht um Sünde oder Uebelthat willen verfolgt und geplagt werden. Denn davor, lehrt der heilige Petrus, sollen sich Christen hüten, da er 1. Ep. 4, 15. sagt: „Niemand aber unter euch leide als ein Mörder, oder Dieb, oder Uebelthäter, oder der in ein fremd Amt greifet.“ Da heißt das griechische Wort eigentlich einen Anführer, der sich in eines Andern Amt drängt: das sind die Gesellen, die regieren wollen, da man ihnen nichts befohlen hat.

69. So sind nun die Ursachen, darum die Frommen und Gottseligen in diesem Leben also geübt und geplagt werden, nicht ihre Sünden und Uebelthaten, sondern es ist allein der Wille GOttes, Gerechtigkeit und ihr gutes züchtiges Leben. Also verfolgen uns heutiges Tages der Türke, Pabst, Ketzer und Rotten, nicht daß sie uns mit Wahrheit um Dieberei oder Unzucht, oder anderer Laster willen beschuldigen könnten, sondern um der Lehre und rechten Gebrauchs willen der heiligen Sacramente, so wir lehren. Wir müssen doch nichts desto weniger die Schmach- und Scheltworte von ihnen leiden, daß sie uns als Aufrührer, Rebellen, oder ungehorsame und halsstarrige neidische Leute schelten; so müssen wir heißen. Denn die Sacramentenschwärmer wissen uns auch keine andere Ursach vorzuwerfen, denn daß wir wider die Liebe handeln, und die Einigkeit trennen und verhindern etc. Und wiewohl solches alles wahrhaftig auf ihnen ruht, daß sie dessen billig mögen beschuldigt werden, messen sie es doch nur uns zu. Sie thun es, wir aber müssen es tragen. Also soll es sein.

70. Joseph wird betrübt und geängstiget nicht um Ungehorsams willen, sondern daß er seinem Vater gehorsam, und ein frommer, heiliger, züchtiger Jüngling ist; wie er bald hernach um Zucht und Keuschheit müssen auch in das Gefängniß wird geworfen werden. Und bis daher hat ihn der Teufel auf der Linken angegriffen, welcher ihn mit Armuth und Jammer versucht hat in der schweren und verdrießlichen Knechtschaft, und er ist noch in den obersten Kerker der Hölle; bald aber wird er auch in den untersten geworfen werden. Es muß noch ärger mit ihm werden.

\quellseite{326}

\section[Dritter Theil]{Dritter Theil. Wie Potiphars Weib in böser Lust gegen Joseph entbrannt und denselben zu verleiten sucht, und wie sich Joseph dabei verhalten}

\begin{quote}
\textbf{V. 6. Und Joseph war schön und hübsch von Angesicht.}
\end{quote}

71. Dies Stücklein hat Mose mit hereingesetzt, daß er desto füglich das, so hernach folgt, auch mit anhängen möchte. Er sagt, Joseph sei schön und hübsch von Angesicht gewesen. Denn wiewohl er nicht köstlich gewesen, noch auch sonst sonderliche Wartung gehabt, ist doch sein schön Angesicht und freie Gestalt an sich selbst ein großer Schmuck und Ehre gewesen, vornehmlich an ihm, als einem Jünglinge und der noch in der rechten Blüthe seiner Jugend stand. Und in welcher Herzen der Heilige Geist wohnt, dieselben pflegen auch ein frei, aufgerichtet, fröhlich Angesicht zu haben. Denn GOtt erweckt in ihnen die rechte Freude, welche auch das Angesicht erfreut; wie denn die Traurigkeit hinwiederum das Angesicht auch gar entstellt. Und kommt das auch noch dazu, daß er von sehr frommen Eltern geboren und seine Mutter Rahel selbst auch schön von Gestalt gewesen: und neben dem, daß Joseph schön war von Angesicht, ist er auch zugleich dabei gottseelig und von guten ehrlichen Sitten gewesen. Daher ist seine Frau, Potiphars Weib, in ihrem Herzen gegen ihn mit unzüchtiger Liebe also entzündet, welche sich an ihm so vieler herrlichen Tugenden verwundert, so um der Schönheit willen an ihm höchst mehr hervorgezogen haben. Wie denn der Poet Virgili recht sagt: Gratior est pulchro veniens e corpore virtus. Die Tugend ist um so angenehmer, wenn Schönheit des Leibes dazukommt. Wenn viele seiner Tugenden dazu kommen, alsdann wird Schönheit und Gestalt so viel trefflicher und höher. Wie ein schönes Weib das nicht so hoch ziert, daß sie schön ist von Leibe, als daß sie feine, züchtige, demüthige Sitten hat, nicht bitter oder unfreundlich ist. Solche Tugenden übertreffen weit alle Schönheit des Leibes, und machen, daß die, so ungestaltet sind, gleichwohl ihrerhalben gerühmt und geehrt werden.

72. So ist nun Joseph aller Gaben oder Tugenden halben, beide äußerlich am Leibe und auch innerlich an der Seele, sehr vortrefflich gewesen, ist GOtt und Menschen wohlgefällig und angenehm gewesen. Denn es ist noch ein größerer Schmuck am Menschen, daß er tugendsam ist, als daß er schön von Gestalt ist. Wenn aber die Tugend und Schönheit zusammen kommen, so ist es ein unüberwindlich Band der Liebe, es muß Einen fangen. Solcher Leute aber findet man jetzt zur Zeit wenig, die mit so vielen Tugenden geziert, und bei jedermann so holdselig und angenehm sind. Denn derer sind viel, die von Unzucht und Zorn entbrannt, dazu mit Ungeduld, mit Fressen und Schweglerei ihren schönen Leib auf das schändlichste entstellen und verderben; inwendig aber schleppen sie sich auch mit einem solchen Gewissen, daß mit mancherlei Sünden befleckt ist. Aber Joseph ist auch in der rechten Blüthe seiner Jugend ein trefflich Exempel aller Tugenden, und an dem die Weiber und auch seines Herrn Frau ein lieblich Schauspiel gehabt. Nun kommt aber der Teufel und versucht ihn auch auf der Rechten, nachdem er ihn auf der Linken mit schwerer Arbeit, Knechtschaft und Elend, darin er gelebt, nicht hat überwinden oder zu Falle bringen können. Und dies sollen wir fleißig merken, damit wir den lieben Eltern auch mögen gleichförmig werden, beide mit Wohlthun und auch mit Leiden.

\begin{quote}
\textbf{V. 7—10. Und es begab sich nach dieser Geschichte, daß seines Herrn Weib ihre Augen auf Joseph warf, und sprach: Schlafe bei mir. Er weigerte sich aber, und sprach zu ihr: Siehe, mein Herr nimmt sich nichts an vor mir, was im Hause ist; und alles, was er hat, das hat er unter meine Hände gethan; und hat nichts so groß in dem Hause, das er vor mir behalten habe, ohne dich, indem du sein Weib bist. Wie sollt ich denn nun ein solch groß Uebel thun, und wider GOtt sündigen? Und sie trieb solche Worte gegen Joseph täglich. Aber er gehorchte ihr nicht, daß er nahe bei ihr schlief, noch um sie wäre.}
\end{quote}

73. Es hat der Teufel bis daher gesehen, daß dieser fromme heilige Jüngling mit großer unglaublicher Geduld und starkem Muthe alles Unglück und Beschwerung seines Elendes überwunden hat, darin er von seinen Brüdern und Gesfreunden weggerissen, dazu der größten Herrlichkeit und Hoffnung der Erbgut beraubt sein und unter den groben abgöttischen Leuten wohnen müssen; durch welche Beschwerungen alle doch gleichwohl sein Herz schier ganze zehen Jahre lang nicht hat können ermüdet oder geschwächt werden. Denn wiewohl er vom siebzehnten Jahre an bis in das sieben und zwanzigste, welches er dazumal erreicht, vom Satan mit vielen unzähligen Pfeilen und sonst harten Stößen angegriffen worden, hat er doch allein dieses alles überstanden, ist beständig und unüberwindlich geblieben, so er doch niemand um sich gehabt, der ihn hätte mögen ermahnen, oder aber der mit ihm auch in gleicher Trübsal gewesen; ja, er hat endlich gar keine Exempel vor sich gehabt, sondern dieselben, so er täglich gesehen, sind gar widerwärtig gewesen, dadurch er leichtlich von rechter Lehre der wahren Gottseligkeit hätte mögen abgeführt werden; derhalben muß gewißlich eine große Kraft des Geistes und göttlicher Gaben an ihm gewesen sein. Denn es wäre ihm ja etwas leidlicher gewesen, wo er doch nur einen einigen Mitgesellen gehabt. Aber er allein hat ohne Vermahnung, ohne Opfer und ohne heilige Predigten mit großer Beständigkeit noch immer die wahre Gottseligkeit, und wie man göttlich leben soll, behalten, und ist in dem so gar nicht anderes Sinnes geworden, daß es scheint, daß er immer, je länger je mehr, zugenommen und in solcher Beständigkeit gar befestigt worden sei.

74. Jetzt aber greift ihn nun der Teufel auch auf der rechten Seite an, das ist, mit Schmeichelei oder Anreizungen zur Unzucht und Wollust, welche Anfechtung viel schwerer und gefährlicher ist, vornehmlich bei einem Jünglinge, denn sonst alle die Stürme und harten Stöße sein können, so von der Linken hergekommen. Denn um das achtzehnte Jahr fängt die Erbsünde an sich zu regen und zu wüthen, und finden sich alsdann in den Herzen der Jünglinge erschreckliche Bewegungen und Gedanken, damit sie zur Unzucht oder unziemlicher Vermischung gereizt werden. Und hat dazumal Joseph unter den ärgsten Exempeln der Heiden gelebt, dadurch die Natur, so von sich selbst zur Unzucht geneigt, noch immer mehr hat können gereizt und entzündet werden, vornehmlich an einem solchen seinen gesunden Menschen. Denn also sagt Augustinus an einem Orte: Je gesunder wir am Leibe sind, je größer und unsinniger diese Seuche wird. Sie ist nicht so gar hitzig in einem schwachen Leibe, und der mit Sorge und Bekümmerniß beschweret ist, als in einem starken gesunden Leibe.

75. So hat nun dieser Jüngling ganze zehen Jahre lang gelebt in solcher Anfechtung böser Lust, und hat männlich gestritten, da er nicht einmal, oder von Einem Weibe allein, sondern zum öftern und von vielen zu unsüchtiger Liebe ist gereizt worden. Denn die Mägdchen pflegen dies Unglück und Anfechtungen auch zu fühlen, und wo sie bei jungen Gesellen wohnen, versuchen sie sich an denselbigen auf mancherlei Weise, damit sie ihre Herzen jetzt hieher, dann dorthin bewegen, ob sie dieselbigen zu Liebe reizen möchten, vornehmlich wo sie sehen, daß sie schön von Gestalt und stark von Leibe sind, also daß die jungen Gesellen oftmals größere Mühe haben, solcher Anreizungen sich zu erwehren, als ihren eigenen Lüsten zu widerstehen. Denn also werden diese Dinge in den Schriften der Poeten abgemalt, wie jener Bauer im Virgil sagt:

\begin{quote}
Malo me Galatea petit lasciva puella, Et fugit ad salices, et se cupit ante videri.
\end{quote}

das ist:

\begin{quote}
Galatea, die vorwitzige Metz, Treibet mit mir fast vielen Scherz, Wirft mich mit einem Apfel fein, Will damit von mir gesehen sein, Thut sich darnach verborgen schon Wohl hinter einem Weidenbaum, Als habe sie das nicht gethan, Und will doch mein nicht müßig gahn.
\end{quote}

Dazu kommt dann auch der Teufel mit seiner Kunst und List, damit er die Mägdchen unterrichtet und anreizt, die jungen Gesellen, so ihrer Sachen keine Acht haben, zu betrügen. „Der Teufel kann wohl eine Metze auf Einen abrichten.“

\quellseite{327}

76. Ueber die gemeinen Anreizungen aber der Mädchen, so entweder Hausgenossen mit ihm oder sonst Nachbarinnen gewesen sind, von welchen Joseph ohne Zweifel auf mancherlei Weise versucht worden, seine Keuschheit und Zucht hintan zu setzen, wird er von der Frau im Hause selbst hart angefochten, die es versucht, ihn mit ihren Schmeichelworten zur Unzucht und Hurerei zu reizen. Das ist nun gar aus der Weise. Diese Anfechtung ist unter den andern allen die allerstärkste und gewaltigste um solcher seinen sichern Gelegenheit und Urlaub willen zu sündigen; denn die Frau und die rechte Hausmutter selbst hat Lust, mit ihm zu buhlen. Wahrlich, wer sich eines solchen Weibes entschlagen und enthalten kann, da magst du wissen, daß ein solcher sich in seinem Leben wohl weiß recht zu halten, daß er nirgends zu viel oder zu wenig thue. Denn was Andere in solchen Fällen willig und gern angenommen hätten, oder dem sie mit höchstem Fleiß und Begierde von selbst gefolgt wären, wie denn die Jugend der Unzucht pflegt nachzuhängen und frei unbedachtsam hinzuzugehen, das hat Joseph frei verachtet und abgeschlagen, ungeachtet, daß ihm unverschämter Weise solches zugemuthet und angeboten worden, und die Frau bei ihm täglich fast ungestüm angehalten, und ihn immer genöthigt und getrieben hat.

77. Nun war aber Joseph ein Jüngling hübsch von Gestalt und stark von Leibe, mitten in der Blüthe der Jugend, da er leichtlich hätte mögen erweicht und überwunden werden von einem solchen ansehnlichen und berühmbten Weibe, die ihn so anhaltend und unablässig zur Unzucht gereizt. Und er hätte also Macht gehabt, ihren Willen zu thun, daß ihm solches nicht allein ohne alle Gefahr gewesen, sondern er hätte davon auch noch wohl großen Nutzen haben mögen. Denn es schien, als ob ihm nichts Anderes so wohl hätte dienen mögen zu großem Ehren und höherem Stande zu kommen, als wo er der Frau Gunst hätte und daß ihn dieselbe möchte gewogen sein, die eines so großen Mannes Weib war, der bei dem Könige vor Andern in großem Ansehen und Gnaden gewesen, und zum Obersten über die Köche und Metzger am Hofe gesetzt war. Welche Ehre und Würde die Hebräer sehr groß machen und sagen, er sei in Egypten wie ein Hoherpriester oder Papst gewesen. Derhalben hätte Joseph bei ihm durch mancherlei Gelegenheit von der Frau können gerühmt und befördert worden sein.

78. Joseph aber hat diese bequeme Gelegenheit und so ihm ganz sicher und ohne Gefahr gewesen, auch große Gewalt, Gunst und Wohlust hätte bringen mögen, verachtet, dieselbe gar abgeschlagen und fahren lassen. Das ist wahrlich ein klares Anzeichen der größten Frömmigkeit und Keuschheit; welche Tugenden daher ihren Ursprung gehabt, daß in seinem Herzen Gottes Wort, dazu rechte Ehre und Furcht Gottes gehabt, und daß daselbst der Glaube durch den Heiligen Geist regieret hat. Ich habe auch wohl anderer Jünglinge Exempel gehört, welche keusch und züchtig bei Weibern und Jungfrauen gelebt: aber hier ist die schändliche Liebe, so ihn angereizt und getrieben hat, welche auch leichtlich und in einem jeden Augenblick das junge Herz und das Fleisch, so von selbst von Unzucht entbrannt ist, hätten verkehren und überwinden mögen. Darum ist es eine sehr böse Anfechtung gewesen.

79. Da lerne nun, wer lernen kann, und lasset uns mit Fleiß an dies Exempel denken, und dabei lernen, wie angenehm es Gott dem Herrn sei fest am Glauben halten, ihn fürchten, sich vor seinen Geboten scheuen, und von ihm nicht abweichen, sondern in rechter wahrhaftiger herzlicher Gottseligkeit verharren. Dazu uns denn Christus Luc. 21, 19. ermahnt, da er sagt: „Fasset eure Seelen mit Geduld“; item Matth. 24, 13.: „Wer beharret bis ans Ende, der wird selig.“ Denn der Teufel wird nicht müde, kann nicht ruhen, sondern hält immer an, sucht Ursache und Gelegenheit, die Leute zu Sünden zu treiben. Es ist verwegen nicht genug, nur einen Tag oder ein Jahr sich gegen ihn mit Mühe und Arbeit zu wehren und zu streiten; sondern es ist vonnöthen, daß man dabei Geduld habe und in solchem Streit mit Beständigkeit verharrre, bis er entweder gewinne, oder aber überwunden und ihm der Kopf zertreten werde. Es ist aber einem Jünglinge gar schwer, ganze zehn Jahre mit so großen Feinden, als, dem Fleisch und dem Teufel, zu kämpfen, die mit so viel Geschossen, und so gar bequemer und guter Gelegenheit auf ihn einstürmen. Denn hier ist der Dinge keines unterlassen worden, die dazu gehören, daß Joseph hätte mögen überwunden werden.

80. Es hat ferner dies Weib den Jüngling ja sehr unverschämt angeredet, da sie sagt: „Schlafe bei mir.“ Wiewohl ich es nicht dafür halte, daß sie gleich mit diesen Worten angefangen habe, damit sie ihn angereizt und zur Unzucht genöthigt hat, sondern wird sich beflissen haben, ihn anfänglich mit sonderlicher List und Schmeichelworten einzunehmen, wird sich gegen ihn also erzeigt haben, als wäre sie ihm sonderlich gewogen, hätte ihn nicht anders denn in Ehren lieb, wie das einer frommen Matrone wohl ansteht, und zwar wegen seiner großen Tugenden und Geschicklichkeit. Sie mag ihn auch bisweilen angelacht und freundlich mit ihm geredet haben und gefragt, wie es mit dem Hausregiment beschaffen wäre, wie sich die andern Knechte in ihrem Amt hielten, ob sie desselben auch fleißig warteten? Wird ihm auch bisweilen etliche Geschenk gemacht haben, damit anzuzeigen, daß ihr und ihrem Herrn sonderlich sein treuer Dienst lieb und angenehm wäre. Sie ist nicht bald mit diesen Worten heraus gefahren: „Schlafe bei mir“; sondern hat mit vielem Schmeicheln, mit Gesprächen und mancherlei heimlicher List sein Herz angestochen, und hat sich damit also einen seinen Zutritt bereitet, diese Bitte darauf an ihn zu stellen.

81. Diemeil Joseph aber von Natur sanftmüthig und züchtig war, hat er ihr freundlich geantwortet. Denn es hat dem Knechte nicht geziemet, daß er sich hätte sollen mit stolzen Worten gegen seines Herrn Frau vernehmen lassen. Darum daß sie gesehen, daß er bisweilen auch freundlich dazu gelacht und in aller Demuth mit ihr geredet, hat sie daraus in ihrem Herzen immer größere Hoffnung geschöpft, bis sie es dafür gehalten, er würde nun durch ihre Güte und freundliches Gespräch, dazu auch die Wohlthat, so sie ihm erzeigt, überwunden sein, und daß sie ohne Zweifel seiner nun wollte mächtig werden. Nun ist es Zeit, hat sie gedacht, und fährt heraus: „Schlafe bei mir.“ Denn sie denkt, er sei nun gar eingenommen und überwunden mit so viel Schmeichelworten, Geberden und Winken, deren die Weiber wohl wissen zu gebrauchen. Welches Mose auch hat wollen anzeigen, wie es scheint, da er sagt, daß sie ihre Augen auf Joseph geworfen, das ist, sie hat ihre Augen auf diesen Jüngling gerichtet, hat immer auf ihn gesehen, hat sich seiner schönen Gestalt und Tugend verwundert, hat ihn oftmals gelobt wegen seiner ehrbaren Sitten, so er an sich gehabt, und auch seiner Unschuld halben; hat dazu auch an ihm die vortreffliche große Geschicklichkeit gelobt, und hat also sein Gemach sich selbst damit den Weg bereitet zu der Schande, so er mit ihr hat begehen sollen.

82. Darum hat er die vorige Freundlichkeit nun verwandelt in eine heilige und männliche Beständigkeit; welches ohne Zweifel seine Frau sehr verdrossen haben wird, sie hat aber doch nicht aufgehört mit ihrer unbilligen und schändlichen Bitte immer weiter bei ihm anzuhalten, wiewohl Joseph es versucht hat, ihr solche zu widerrathen und sie von der schändlichen Liebe abzuleiten. Wie er denn eine ernste und heilige Predigt thut, damit er sie von der Brunst ihrer gefaßten Unzucht möge abwenden.

83. Denn es sind nicht Schmeichelworte, wie davon die Grammatici unnütz machen; sondern es ist eine sehr ernste und weise Rede, damit er die schändliche und unzüchtige Brunst der Unzucht an ihr auszlöschen suchte. Aber diesem Unglück kann mit Worten nicht gesteuert werden, und hilft auch dazu sonst gar kein ander Mittel, wie der Poet sagt: Aegrescitque medendo: Je länger man daran heilt, je kränker man wird, und je mehr Arznei man dazu gebraucht, je mehr es sich regt; wie das zu sehen ist aus der Beschreibung Didus im Virgil. Denn alles, was aufgebracht und dagegen gesetzt wird, diese Seuche zu lindern oder aber wegzunehmen, dadurch wird sie vielmehr gemehrt und immer mehr angezündet, nicht anders, als wenn man Oel in das Feuer thut. Es ist gar eine böse unsinnige Brunst, welche man nicht leichtlich bezwingen kann; denn sie ist blind und taub, und ist einfach eine Leidenschaft und Raserei, vom leidigen Teufel eingegeben. Es ist keine züchtige und heilige Liebe, da man ein Weib oder jungen Gesellen zur Ehe begehrt, sondern vielmehr eine böse Seuche, ja, eine Raserei und Unsinnigkeit; wie sie denn auch dermaßen von den Heiden beschrieben ist.

\quellseite{328}

84. So predigt nun dieser fromme heilige Jüngling dem unsinnigen Weibe und hält ihr drei Argumente vor. Das erste ist, daß er sagt: Das sei ferne von mir, daß ich sollte wider GOtt sündigen. O liebes Weib, bedenke doch, welch eine große Sünde das sei, die Furcht GOttes fallen lassen, sich vor ihm nicht scheuen und ihn erzürnen um solcher geringen kurzen Lust willen. Aber da hilft kein GOtt, sie bedenkt nicht, was GOtt verheißet oder was er gedroht hat, denkt auch an keinerlei Belohnung oder Strafe. Das andere Argument ist: Du bist meines Herrn Weib, denselben sehe an wie dich selbst. Fürchtest du dich nicht vor GOttes Zorn, bist auch der Meinung, daß du deiner Seligkeit nicht achten dürfest, so erforsche doch dein eigen Gewissen, und bedenke, was der Name und Herrlichkeit einer Hausmutter erfordere: bedenke, daß du deinem Ehemanne allein zugefügt und gegeben bist, und daß ein solcher trefflicher Mann, dessen Würde dir bisher zu großer Zierde gereicht, billig in allen Ehren und Würden solle gehalten werden. Das ist wahrlich eine keine weise Rede. Das dritte Argument ist noch das allertapferste, nämlich, daß er nun die Undankbarkeit und Untreue anführt und sagt: Mein Herr hat mich in eurem Hause hoch geehrt, hat mich zum Herrn und Regenten gesetzt über das ganze Haus; wie würde aber nun das so eine große Sünde sein, solcher Wohltat vergessen, und meinen Herrn in seiner Hoffnung, so er auf mich gesetzt, betrügen!

85. Diese Argumente sind ja gewiß genug, und ich kann nicht sehen, was doch Kräftigeres möchte angeführt werden, die Unzucht damit zu überwinden. Er hält ihr die Furcht GOttes entgegen, das Amt eines frommen Eheweibes und ihr eigen Gewissen, endlich die Undankbarkeit und die Untreue, damit er sich nicht beladen will. Aber je tapferer, ernstlicher und heiliger er ermahnt wird, und je mehr sie sieht, daß Josephs Herz befestigt ist, Zucht und Keuschheit zu halten, je heftiger sie entbrennt. Derhalben ist dies ein heftiger Kampf gewesen zwischen der Beständigkeit, Keuschheit zu halten, und der Unzucht, so sich auf das allerköstigste dawider gelegt. Es sind daselbst zwei feindliche Geister mit aller Gewalt aneinander gerathen. Der Teufel ist auf Joseph eingestürmt sehr freundlich und mit seinen glatten Worten; aber es ist doch eine erschreckliche Grausamkeit. Denn wo er sieht, daß die frommen gottseeligen Kinder mit keinerlei Pein oder Marter können überwunden werden, so greift er sie zuletzt an, und versucht, ob er sie mit Unzucht und Wollust überwinden möge; wie Hieronymus an einem Orte sagt.

86. Und Cyprian schreibt, daß die Heiden die Märtyrer mit langsamer Strafe haben pflegen zu peinigen, und haben nicht zugelassen, daß sie bald gar dahin stürben, da sie doch sehr heftig begehrt haben zu sterben. Darnach aber, da sie dieselben nicht haben können mit peinlicher Strafe überwinden, sollen sie unzüchtige Weiber, die sehr schön von Gestalt waren, herzu haben bringen lassen, welche mit Küssen und anderer Schmeichelei die böse Lust des Fleisches reizen sollten. Denn sie haben wohl gewußt, daß die Christen die Keuschheit loben, sagen und halten auch, daß GOtt unkeuslicher Vermischung feind sei; derhalben wollten sie, daß auf diese Weise frommer Christen unverletzte Gewissen sollten befleckt und die Lehre auch geschändet werden. Eben auf solche Weise ist Joseph von diesem unzüchtigen Weibe auch angefochten worden, welches mit Gewalt ihm zugesetzt und der Dinge keines unterlassen hat, damit auch die heiligen Märtyrer sind angefochten und hart bestürmet worden.

87. Hieronymus erzählt, daß einer von den Märtyrern, da er gefühlt, daß ihm Hände und Füße gebunden gewesen, und er also mit Gewalt zur Unzucht hingeführt worden ist, sich selbst soll die Zunge abgebissen und dieselbe der Hure in das Angesicht gespien haben.

88. Die rohen gottlosen Leute haben wenig Sorge darum, wie sie die Keuschheit erhalten mögen; darum macht ihnen auch der Teufel auf der Teufel damit nichts zu schaffen; denn er hat sie schon zuvor in seinen Stricken. Du bist bereits mein, denkt er. Er versucht sie durch so gar offenbare Gelegenheit nicht, darnach sie doch groß Verlangen haben; ja, er treibt sie, daß sie auch mit großer Gefahr, mit höchster Beschwerung Leibes und Seele, dazu auch mit Verlust ihrer Gesundheit erlangen, was sie begehren. Denn er weiß, daß sie ihm nun schon unterworfen sind: darum geht er etwas härter mit ihnen um, daß sie also den rechten Lohn ihres gottlosen schändlichen Wesens empfangen mögen, und daß sie ja nicht ohne Krankheiten, ohne Schande und andere Verhinderung Gelegenheit finden mögen, ihre böse Lust zu erfüllen; welche Gelegenheit er sonst von selbst und ohne alle Beschwerung denen entgegen bringt, die gottseelig sind und den Heiligen Geist haben.

89. Das Stücklein im Tertze, da Joseph sagt: Siehe, mein Herr weiß nicht, was im Hause ist, ist etwas dunkel, und der Juden Rhabinen haben es auf mancherlei Weise ausgelegt. Ich halte dafür, dies sei einfältig der rechte Verstand: Mein Herr weiß nicht, was ich weiß, hat kein Mitwissen um die Dinge, so im Hause geschehen. Ich sage ihm nicht von den Dingen, so das Hausregiment belangen, ich regiere das Haus frei aufrichtig, er fordert von mir keine Rechenschaft von allen und jeden Dingen; darum weiß er nicht, was ich weiß, wie hoch es läuft mit seinem eigenen gewonnenen Gute, wie viel Vieh er habe, wie man säe oder ernte; ich weiß solches allein. Das, habe ich gesagt, sei ein Argument, so er von der Undankbarkeit und Untreue genommen. Als wollte er sagen: So große Ehre hat mir mein Herr erwiesen, daß er mich nicht allein zum Herrn über Alles gesetzt hat, obwohl ich sein leibeigener Knecht bin, sondern hat mir auch die Freiheit des ganzen Hausregimentes übergeben, daß er von mir davon nimmer eine sonderliche Rechenschaft fordert; denn er sieht, daß sein Gut und das ganze Hausgesinde wunderbarer Weise zunimmt und gemehret wird, und denkt nicht, wo das herkomme oder wie es möge zugehen. Warum sollte ich denn, sagt er, mit dieser Sünde und Untreue solchen gütigen und freundlichen Herrn erzürnen?

90. Nun folgt aber auch im Tertze, wie er dies Argument herausstreichen, da er sagt: „Und mein Herr hat nichts so groß in dem Hause, das er vor mir verborgen habe“, welches Hieronymus im Lateinischen etwas allzu frei verdolmetscht hat. Nach dem Hebräischen lautet es: Er hat nichts vor mir verborgen, ohne allein dich, weil du sein Weib bist. Wir haben in unserm Deutschen den Accusativum, eum, mit dem Dativ, ei, verdolmetscht, und ist dies der rechte Verstand der Worte: In diesem Hause hat mein Herr nichts so groß und lieb, so köstlich, daß er mir nicht sollte vertrauen und mich darüber regieren lassen. An andern Orten, liegenden Gründen oder Stätten mag er vielleicht noch etwas Anderes lieber haben, denn daß er mir solches befohlen oder vertrauen wollte; aber in diesem Hause ist unter allen Dingen nichts, darüber ich nicht Macht hätte, ohne du allein.

91. Das hebräische Wort, chasach, wird oft gebraucht, als, 1 Mos. 1, 4.: „Da schied GOtt das Licht von der Finsterniß.“ Es heißt aber oftmals so viel als, verbergen oder verhehlen, welches geschieht durch die Figur, metaphora genannt. Denn was in Finsterniß ist, wird verhohlen, daß man es nicht sehen kann, als, 1 Mos. 22, 16., da GOtt zu Abraham spricht: „Dieweil du solches gethan hast, und hast deines einigen Sohnes nicht verschonet“ etc. Da lautet es nach dem Hebräischen auch also: Non tenebrasti filium tuum, was der lateinische Dolmetscher gegeben hat: Non peperceristi, und ist so viel gesagt: Du hast ihn nicht verhehlet oder verborgen, hast ihn mir nicht entzogen, sondern mir zum Opfer gegeben, wie ich dir geboten habe. Also lautet es an diesem Orte auch im Lateinischen: Herus non posuit a me aliquid in tenebris etc., das ist: Mein Herr hat nichts vor mir versteckt oder verhehlet, sondern Alles in meine Gewalt und unter meine Hände gethan. Wir haben es im Lateinischen gegeben: Nihil excepit a me: Er hat nichts so lieb in diesem Hause, das er vor mir verberge. Es ist mir Alles übergeben, dich ausgenommen, du allein bist in Finsterniß verborgen und wirst vor mir heimlich gehalten: ich soll dich nicht sehen oder kennen, indem du sein Eheweib bist. Er leugnet nicht, daß er ihr nicht sollte schuldig sein, Gehorsam und sonst allerlei Dienste zu leisten, dergleichen sie selbst den Kindern bedürfte; aber das allein ausgenommen, sofern du sein Eheweib bist, sagt er, da bist du mir verboten. Die eheliche Pflicht und Recht an deinem Leibe zu haben, ist mir verboten.

\quellseite{329}

92. Auf diese Weise predigt Joseph der Frau, die gar unsinnig war, und führt an den Fleiß, Gehorsam und Treue, so er seinem Herrn schuldig; erinnert damit sie selbst auch ihrer ehelichen Treue. Als wollte er sagen: Es hätte dir gebühret, daß du mich solltest ermahnt und erweckt haben, in meinem Amte fleißig und treu zu sein, und solltest mich zu solcher Untreue des Ehebruchs und Undankbarkeit nicht treiben, welche auch sonst unzählige Laster in sich begreift; wie man sagt: Omnia convitia dixeris, si ingratum dixeris, das ist: Wenn man Einen undankbar schilt, so schilt man ihn damit, als der sonst auch mit allen Lastern behaftet sei. Und Bernardus sagt: Nullum vitium perniciosius est ingratitudine, eo, quod exsiccet fontem bonitatis, das ist: Kein Laster ist schädlicher denn die Undankbarkeit, darum daß sie den Brunnen der Güte oder Wohlthat gar austrocknet und leer macht. Darum sehen wir, daß Joseph ein sehr heiliger und weiser Mann gewesen ist, der allerlei Tugend an sich gehabt; ein solcher Mann, den unser Herr Gott gemacht hat.

93. Nun schließet er aber mit ganzem Ernste und sagt: „Wie sollte ich denn nun ein solch groß Uebel thun, und wider Gott sündigen?“ Als wollte er sagen: Es ist eine große Sünde, meines Herrn Weib schänden, der mir die ganze Hausforge befohlen und auferlegt hat: aber es wäre noch eine viel größere Sünde wider Gott. Das ist auch das allervornehmste und stärkste Argument. Aber in der Welt ist es das leichteste; denn die Welt achtet des Zornes Gottes nicht, die Gottseligkeit aber bewegt das Gericht und die Rache Gottes sehr. Wenn die Welt hört, daß der göttlichen Drohungen und Strafen gedacht wird, kann sie dazu noch wohl mit lauter Stimme lachen, und ist eben der Meinung, wie jener Strauchdieb auch war, welcher dem Wandersmanne seinen Stock raubte, und da derselbe rief, er würde für solche Uebelthat am jüngsten Gerichte Rechenschaft geben müssen, antwortete der Mörder: O wenn es noch so weit dahin ist, bis diese Sünde an mir soll gestraft werden, so gib mir das Hemd auch her. Also geht es, wenn Meister Hans nicht mit dem Strange und Schwert da ist.

94. Aber in den Herzen der Gottseligen hilft es sehr viel dazu, daß sie sich vor Sünden hüten, wenn sie der Strafe Gottes gedenken, welcher ernstlich mit denen zürnt, so da sündigen. Denn also hat Joseph auch sein Herz und Augen auf Gott gerichtet, da er doch sehr große Ursache und Gelegenheit, so ihm ohne Gefahr gewesen wäre, zu sündigen gehabt, dadurch er seiner Frau Gunst und Gnade, dazu ihre Güter hätte bekommen mögen. Er hat aber in seinem Herzen beschlossen, er wolle ihrem Willen und bösen Lüsten nicht folgen, wenn sie ihn auch alles Gut des ganzen Königreichs geben wollte. Aber dies Argument würde bei den Tyrannen und andern gottlosen Leuten nichts gelten; wie denn dadurch dies Weib auch nicht bewogen wird, denn sie hört nicht auf, diesen Jüngling zu treiben, daß er ihren Willen thun soll. Darum sagt nun Mose: „Sie trieb diese Worte gegen Joseph“, die, die, das ist, alle Tage, „täglich“. Denn dies ist eine gebräuchliche Art und Weise zu reden, die insgemein soll verstanden werden. Als, 2 Mos. 8, 14. steht: Fecerunt tumulum, tumulum, das ist, sie haben da einen Haufen, dort einen Haufen gemacht, und viel unzählige Haufen. Und im Propheten Hesekiel steht auch, homo, homo, der Mensch, jener Mensch, das ist, insgemein alle Menschen. Item virgam, virgam (Ruthe, Ruthe).

95. Auf dieselbe Weise hat dies Weib den Jüngling auch alle Tage getrieben und genöthigt mit dieser sehr gewaltigen Anfechtung. Er aber hat ihr nicht gehorchen wollen, hat bei ihr weder schlafen noch wohnen wollen, sondern hat sein Bett an einen besondern Ort gesetzt, daß sie bei der Nacht nicht hat dazu kommen können. Und das ist wahrlich das größte Wunder, welches hoch zu rühmen ist, daß in einem Jünglinge so große Keuschheit gewesen, der doch noch nicht erfahren hat, was es für ein Ding um die Liebe der Weiber wäre. Derhalben hat daselbst der Geist der Unzucht mit dem Geiste der Keuschheit auf das allerheftigste gestritten. Und bei solcher Gelegenheit zu sündigen ist nicht ein geringer Wunder gewesen, die Flammen der Unzucht zu überwinden, als daß die drei Männer, davon Daniel schreibt, aus dem Feuerofen gesund und unverletzt wieder gekommen sind.

96. Bernhard, über das Hohelied Salomos, schilt etliche Geistliche, die bei den Weibern wohnten, und doch vorgegeben haben, daß sie keuschheit hielten. Wie vor dieser Zeit die papistischen Pfaffen ihre Verwandten oder Schwägerinnen zu sich genommen in der Meinung, daß sie mit denselbigen wohl könnten keusch leben. Und da ich ein junger Knabe war, gedachte ich noch wohl, daß die Priester nicht in Verdacht des Ehebruchs oder der Hurerei standen, wiewohl sie bei Weibern wohnten, bis sie darnach ganz unverschämt Blutschande, Ehebruch und Hurerei getrieben, dazu auch andern Männern ihre ehelichen Weiber entführt haben. So sehr hat die Unzucht und Geilheit der Pfaffen seit unserm Gedenken zugenommen. Darum sagt Bernhard zu seinen Geistlichen: Es ist nicht wohl zu glauben, daß ihr solltet keusch leben; denn bei einem Weibe wohnen und sie nicht erkennen oder berühren, ist mehr, als Tote auferwecken.

97. Das redet er aber von der ersten Brunst derer, so die Flammen der Unzucht noch nicht gefühlt, denen ist es unmöglich, daß sie nicht sollten von Liebe entbrannt werden. Einem Ehemann ist es etwas leichter, der nun eine Zeitlang bei dem Weibe gewohnt. Ueber das ist es auch um die Liebe im Ehestande also beschaffen, daß sie selten pflegt beständig zu sein und für und für zu währen, wie Jesus Sirach Cap. 25, 2. unter den Dingen, so beide Gott und Menschen wohlgefallen, dies auch mit aufzählt: „Wenn Mann und Weib sich mit einander wohl begehen.“ Und wo die Deutschen sehen, daß eine Ehe wohl gehalten ist, pflegen sie dasselbe als das allerhöchste Gut zu rühmen. Sie sagen sie, das ist eine gute Ehe. Denn der Satan pflegt gemeiniglich den Samen des Hasses und des Zankes unter sie zu streuen; dadurch denn die Herzen von einander gerissen werden, daß sie hernach auch mit Ehebruch und sonst unziemlicher Unkeuschung sich beflecken und versündigen. Im Anfang, ehe sie ehelich mit einander geworden, sind sie wohl von Liebe gegen einander gar entbrannt; wenn sie aber die Ehe vollzogen und nun ehelich bei einander wohnen, so ist der Teufel bald da, der alsdann auf mancherlei Weise versucht, solche hitzige Liebe auszulöschen. Denn er kann wohl aus einem Munde beide kalt und warm blasen.

98. Und wenn die Jugend von Liebe und Unzucht so unsinnig ist, daß ist keine einfache natürliche Liebe, sondern ist des Teufels Odem, damit er sie angeblasen, wie Hiob am 41. Cap. V. 12. sagt: „Sein Odem ist wie lichte Lohe, und aus seinem Munde gehen Flammen.“ Derhalben bläst er der Hur und dem Hurer in das Blut und erweckt also in ihnen ein sehr heftiges Feuer. Wiederum aber bläst er in die Herzen der Eheleute einen kalten Odem, davon die Hitze der ehelichen Liebe kalt wird und vergeht. Daher sind gekommen diese und dergleichen Worte der Ehebrecher: Uxor est honestatis nomen, mulier vero voluptatis, das ist: „Eheweib“ ist wohl ein feiner ehrlicher Name, aber „Weib“ ist ein Name zur Wollust. In solchen Ehebrechern hört die Brunft des Weibes gegen den Mann auf und wiederum auch des Mannes gegen das Weib; denn der Teufel bläst aus dem Norden eine erschreckliche Kälte und Frost daher. Derhalben wer da will ein Weib nehmen, der soll Gott um Gnade bitten, daß er mit dem Weibe möge in Frieden leben und daß sie sich unter einander lieben, damit ja die Herzen nicht von einander gerissen werden. Und wenn zwischen Eheleuten Einigkeit bleibt, daß der Teufel mit seinen Gesägfen nicht darein kommt, alsdann ist es ein offenbarer Segen und Gabe Gottes.

\quellseite{330}

99. Es redet aber Bernhard von denen, die nicht ehelich mit einander verbunden sind, welche der Satan immer mehr antreibt zu unziemlicher Liebe und Unzucht, die noch nicht geschmeckt haben, was für Fälle oder Honig der Liebe sei, und viel mehr zu erlangen hoffen, als sie darnach finden, wenn sie ihres Wunsches theilhaftig geworden.

100. Cyprian disputirt auch wider die Kirchendiener zu seiner Zeit, die auch pflegten bei nackten Weibern zu schlafen. Und Augustinus erzählt, daß etliche sonderliche Ketzer gewesen, die sich von Abel genannt haben Abelonitten. Denn sie wollten nach dem Exempel Abels, der nicht ehelich gewesen, ein eheliches Leben führen, und wohnten doch bei Weibern, wie eheliche Leute pflegen, allein daß sie keine Kinder zeugten. Allenthalben aber unter in Flecken und Dörfern haben sie fremde Kinder zusammen gesucht, welche sie zu Erben gemacht. Und das, sagt Augustinus, sei den Nachbarn, welche arm waren, sehr angenehm gewesen und habe darüber ihre Secte sehr zugenommen.

101. Aber ich will nun wieder zum Texte kommen. Es muß gewißlich Joseph ein Jüngling gewesen sein von großem trefflichem Geiste, nicht allein Unglück und Widrigkeiten zu leiden, sondern auch die fleischliche Anreizung zur Unzucht zu überwinden, dazu die harten Stöße oder Stürme des Satans, beide zur Linken und zur Rechten, zurück zu treiben. Denn er hat ihn nicht nur einmal angegriffen, sondern hat ihn alle Tage mit der gewaltigsten Schmeichelei und Anreizung zur Unzucht genöthigt. Derhalben ist dies Exempel der Keuschheit an Joseph gar selten und unvergleichlich; wie im Jesua Sirach Cap. 42, V. 14 geschrieben steht: „Es ist sicherer, bei einem bösen Mann zu sein, denn bei einem freundlichen Weibe, die ihn zu Hohn und Spott machet“, als wollte er sagen: Es ist sicherer mit bösen seltsamen Männern umzugehen, als mit einem Weibe, die sich freundlich und holdselig erzeigt, und vornehmlich, wenn sie auch noch dazu schön ist. Denn sie kann die Herzen heftig an sich ziehen und dieselben einnehmen. Damit aber Joseph die schwere Anfechtung überwinden möchte, hat er immer mit großer Beständigkeit auf Gottes Gebot gesehen, und da er sieht, daß sein Rath oder Hülfe zu finden ist, sich der unsinnigen Brunst seiner Frau zu erwehren und derselben zu widerstehen, hat er sich doch mit größtem Fleiße von ihr gesondert, hat alle Gelegenheit gemieden, mit ihr unzuziehen oder zu reden, und hat ihr gewehret, daß sie auf keinerlei Weise hat zu ihm kommen können.

\section[Vierter Theil]{Vierter Theil. Wie das Weib Potiphars Joseph mit Gewalt zur Unzucht nöthigen will, wie Joseph davon läuft, und das Weib darauf ihn schrecklicher Dinge beschuldigt}

\subsection{I}

\begin{quote}
\textbf{V. 11. 12. Es begab sich der Tage einen, daß Joseph in das Haus ging, sein Geschäft zu thun; und war kein Mensch vom Gesinde des Hauses dabei. Und sie erwischte ihn bei seinem Kleid, und sprach: Schlafe bei mir. Aber er ließ das Kleid in ihrer Hand und flohe, und lief zum Hause hinaus.}
\end{quote}

102. Siehe doch, wie frevelhaft und unsinnig das Weib ist, welches aufs neue und gar unversehens diesem Jüngling hart zusetzt, auf daß sie ihn unversehens mit Gewalt und heimlicher List unterdrücken und zwingen möge, bei ihr zu schlafen. Darum ist es eine sehr gewaltige Anfechtung. Es steht aber im Hebräischen: secundum diem hanc, um diesen Tag; wohin nun dasselbe solle gezogen werden, weiß ich nicht genau. Es haben vielleicht die Hebräer aus der Väter Satzung das genommen, daß sie sagen, es sei ein sonderlicher Feiertag gewesen, und die Frau habe sich gestellt, als wäre sie krank, daß sie mit den Andern nicht dürfte unter die Leute ausgehen, und sei den Joseph frei gestanden, daß er bei den Götzenopfern nicht habe sein müssen; und derhalben, da die Andern alle zum Opfer oder Götzendienst gegangen, sei Joseph allein im Hause geblieben. Und ist diese Erzählung ganz wahrscheinlich. Denn es scheint, als habe sie lange zuvor mit sonderlichem Fleiß Ursache und dazu auch einen gelegenen Tag gesucht, damit sie, wenn sonst niemand dabei wäre, den Jüngling entweder mit Drohen oder Schmeichelworten dazu zwingen möchte, ihres Gefallens Unzucht mit ihr zu treiben. Joseph aber wird an nichts Böses gedacht und sich vor nichts gefürchtet haben, denkt, es habe nun gar keine Gefahr mit ihm, dieweil die Frau krank gelegen. Denn sonst würde er mit ihr nicht allein im Hause geblieben sein; wie er sich denn zuvor mit höchstem Fleiß ihrer Gemeinschaft enthalten hat.

103. Aber da die Frau solche gelegene Zeit und Stätte bekommen, hat sie gedacht: Er wird mir nun diesen Tag nicht entlaufen, er habe denn meinen Willen gethan. Darum greift sie ihn an, da er dessen keine Sorge gehabt und sich davor nicht gefürchtet hat; sondern da er seiner gewöhnlichen Geschäfte im Hause hat ausrichten wollen, da ist sie ihm um den Hals gefallen, und hofft ihn gewißlich also zu überwinden. Aber der Heilige Geist ist gleichwohl noch stärker. Denn der Teufel ist nun mit seinem Mittagshöden schon überwunden, darum bläßt er von einem andern Ort eine neue Anfechtung daher. Denn Joseph flieht davon und läßt das Kleid in des Weibes Hand. Daraus denn hervorgeht, daß sie ihn mit ganzer Gewalt wird angegriffen und gehalten haben, und daß sie versucht hat, ihn zu erweichen nicht allein mit Küssen und daß sie ihm ist um den Hals gefallen, sondern hat ihn auch gleich mit Gewalt wollen behalten, und daß sie ihm gedroht, wo er nicht würde ihren Willen thun, so wollte sie mit lautem Geschrei alle Nachbarn herzu rufen, daß sie ihn züchtigen, wie sie nur könnten, als der ihr nach ihrer Ehre gestanden, sie derselben zu berauben.

104. Dies ist wahrlich ein sehr heftiger Kampf gewesen, so Joseph mit dieser schändlichen bösen Hure gehalten. Denn sie hat gesehen, daß alle ihre, auch die gewaltigsten, Argumente, heimliches Nachstellen, Gelegenheit, ihre Begierde, Gedanken und alle ihre Anschläge vergeblich und umsonst gewesen. Darum verzagt sie nun und gibt es gar verloren, ein solch hartes diamantenes Herz zu überwinden, und wird nun ihre Liebe zur Wuth und Tollheit; wie es denn zu gehen pflegt, wenn solche Bälge mit ihren hurenischen Schmeichelworten nichts ausrichten können, werden sie gar unsinnig, daß sie auch diejenigen, mit denen sie der schändlichen Liebe nicht pflegen können, nur gerne gar umbringen wollten. Also lieset man von St. Agnes, daß da sie sich männlich gewehret, und der Unzucht widerstanden hat, so ihr des Richters Sohn zugemuthet, und dazu keinerley Weise können gebracht werden, daß sie sich schändlich mit ihm hätte vermischen wollen: da hat er sie wollen zum Opfer zwingen, und dieweil sie dasselbe auch mit großmüthigem Herzen verachtet, hat er es endlich dahin gebracht, daß sie sollte am Leibe gestraft werden.

105. Denn damit geht der Teufel um, daß er die Furcht und Ehre GOttes, mit welcherlei Künsten er nur kann oder mag, in den Herzen der Menschen auslöschen und sie zu allerlei Sünden und Schanden treiben möge. Und wo er beide zur Rechten und auch zur Linken zurück getrieben wird, daß er solches nicht bald kann zuwege bringen, versucht er doch, mit stetigem Anhalten diejenigen matt und müde zu machen, die er einmal hat angefangen anzufechten und zu bestreiten: oder wo er mit Gewalt nicht kann ausrichten, braucht er sonst Betrug und List dazu, damit er doch endlich einen solchen Kriegsmann, der nun lange Zeit also versucht worden, daß er darüber fast matt geworden, gar zu Boden schlagen möge. Derhalben wendet er sich nun auch wiederum auf die linke Seite, dieweil er es zur Rechten vergeblich versucht hat.

106. Denn Joseph kommt nun in Gefahr Leibes und Lebens darum, daß er des Ehebruchs wird angeklagt werden. Das ist nun ein rechter Heiliger, welcher um Gottesfurcht willen, und deswegen, daß er seinem Herrn Treue und Glauben gehalten, dazu um anderer trefflichen Tugenden willen endlich diesen Lohn davon bringt, daß er in das Gefängniß geworfen wird, und bald hernach hat sollen hingerissen werden, daß er auch am Leibe gestraft würde. Also ist es um das Leben der Christen, und derer, die an GOtt glauben, beschaffen, nämlich, daß sie um der Tugend, Wahrheit und Gottseligkeit willen und nicht um Sünde oder Schande willen leiden müssen.

\subsection{II}

\begin{quote}
\textbf{V. 13—15. Da sie nun sahe, daß er sein Kleid in ihrer Hand ließ, und hinaus entflohe; rief sie dem Gesinde im Hause, und sprach zu ihnen: Sehet, er hat uns den ebräischen Mann herein gebracht, daß er uns zu Schanden mache. Er kam zu mir herein, und wollte bei mir schlafen; ich rief aber mit lauter Stimme. Und da er hörte, daß ich ein Geschrei machte, und rief, da ließ er sein Kleid bei mir, und flohe, und lief hinaus.}
\end{quote}

\quellseite{331}
107. Die Hure sieht, daß der Jüngling also beständig ist, daß er nicht hat können überwunden werden und daß er sein Kleid in ihren Händen gelassen; darum wendet sich das weibliche Herz nun zu andern Kunstgriffen. Denn sie macht bei sich eine solche Rechnung: Vielleicht wird er ein Geschrei machen und sagen: Ich habe ihn mit Gewalt dazu treiben wollen, daß er bei mir schlafen sollte, und daß er darüber entflohen sei und sein Kleid bei mir gelassen habe. Dasselbe Argument will ich ihm verkehren, und will demselben eine andere Gestalt geben, die jener gar soll zuwider sein; will sagen, er habe versucht, mich um meine Ehre zu bringen, er sei aber vom Geschrei erschrocken und also davon geflohen, und habe sein Kleid in meinen Händen gelassen. Dies ist zumal eine schöne Rhetorik, ja, es ist eine rechte Teufelskunst.

108. Nun kommt aber Joseph auch seines Gewissens und guten Gerüchts halben vor der Welt in große Gefahr, und stirbt damit nicht allein der Glaube dahin, dadurch er an Gott den Herrn geglaubt, sondern auch der Glaube, damit er seinem Herrn ganz treulich gedient hat; ja, alle seine vortrefflichen Tugenden verlöschen dadurch auch, und wird nun an derselbigen Statt des Ehebruches, dazu der Heuchelei und allerlei Sünden mehr angeklagt: als hätte er mit sonderlicher List sich mit solcher Einfalt, Glauben und Treue in seinem befohlenen Amte also erzeigt, daß er durch solche Gelegenheit seiner Frau ihre Ehre rauben möchte.

109. Also muß er die Schande dieser Hure tragen, dadurch seine allerbesten Werke schändlich entstellt und mit des Teufels Gift besudelt werden: was er in dem Hause und sonst sein Lebenlang Gutes gethan, das hat er um der Schalkheit willen müssen gethan haben, als hätte er sich damit also den Weg wollen bahnen, diese Sünde zu vollbringen.

110. Eben dasselbe ist Christo auch widerfahren. Denn warum hat er die Todten auferweckt, die Kranken gesund gemacht und sonst andere Wunderwerke mehr gethan? Da antworten Hannas und Caiphas: Er habe solches um nichts Anderes gethan, denn daß er dadurch wollte König werden, und diese Städte, sagen sie, und uns alle in Grund verderben. Hinweg, hinweg mit ihm, kreuzige ihn. Er hat uns mit seiner Heuchelei betrogen und verderben wollen, damit er bisher das arme tolle Volk bethört hat. Ja, es hat auch hernach das Volk selbst also geschrieen: Wir meinten, er hätte uns wollen Gutes thun; aber eben mit solchen Wohlthaten hätte er uns schier gar in Grund verderbet. Denn ein solcher Heuchler ist der Teufel, welcher bisweilen auch Gutes thut, auf daß er damit desto mehr Schaden thun möge. Diese Schellen aber hängt er darnach den Werken der gottseeligen Kinder an, als ob sie auch dasjenige, so sie gottselig und nützlich ausrichten, darum thun, daß sie damit Schaden thun und die ganze Welt verderben wollen.

111. Wie zu unsern Zeiten die Wiedertäufer unsere guten Werke auch auslegen und deuten. Es sind ja wohl fromme gelehrte Leute, sagen sie, aber sie stehen doch den gemeinen Sachen und der christlichen Kirche nicht wohl vor, sondern verwirren und verhindern sie vielmehr. Darum schreien sie auch über uns, wir seien Ketzer und die allerschädlichsten Leute, auf daß sie uns nur unter einem gerechten Scheine verfolgen und unser Blut vergießen mögen, und also den teuflischen Haß, so sie gegen uns gefaßt haben, erfüllen können.

112. Daselbst sehen wir in diesem Exempel auch also geschehen. Das Weib ruft ihren Hausgenossen, und spricht: Sehet, er hat uns den hebräischen Mann hereingebracht, und der hat sein Kleid bei mir gelassen. Das ist des Teufels Dialektik und Rhetorik, und nicht allein ein Wahrscheinlichkeitsbeweis, sondern ein zwingender Beweis. Denn das ist eine zwingende Folgerung: Hat er das Kleid da gelassen, so hat er es gewißlich der Frau Gewalt anthun wollen. So ist sie nun die züchtige, keusche und heilige Lucretia, dieweil sie schreit, seufzt und heult: Joseph aber muß der unverschämteste Ehebrecher sein; denn er ist geflohen und hat sein Kleid fallen lassen. Das heißt der Teufel. Wer könnte solches verhüten, denn allein der Heilige Geist, der solches nicht verbietet oder hindert, denn nur zu seiner Zeit, und läßt es bisweilen geschehen, daß er damit die Seinen also versuchen und prüfen möge?

113. Wie hat sie aber dem Gesinde im Hause so bald rufen können, so doch diese Lucretia und Joseph allein gewesen sind? Ich halte dafür, daß dieweil diese Dinge geschehen sind oder kurz darnach, das Gesinde wiederum heimgekommen sei; oder daß sie eben die Stunde, da Joseph geflohen und von ihr gegangen ist, nicht geschrieen habe, sondern erst hernach, da das Hausgesinde wieder gekommen, da hat sie das also erdichtet, wie er eben in dem Augenblick davon gelaufen, so er doch vor einer Stunde oder noch eher sich davon gemacht und hinaus gegangen war. Und wird gesagt haben: O liebes Gesinde, wie seid ihr da eben zur rechten Zeit gekommen! wie hat mich so sehr nach euch verlangt! denn wäret ihr nicht also zu rechter Zeit gekommen, so wäre ich meiner Ehre und dazu meiner Wohlfahrt beraubt worden. Jetzt ist der schändlichste Ehebrecher davon gelaufen, hat mich und sein Kleid hier bleiben lassen.

114. Ja, das ist die List der Weiber, und ist eigentlich des Teufels heimliches Nachstellen und seine Hinterlist. Sie hat es nun dahin gesetzt, sie wolle ihn zu Schanden machen und um das Leben bringen, darum sucht sie zu Hause, wie sie nur kann und mag, allerlei Sünde und Schande, die sie möge auf ihn legen, wie sonst bis daher und auch noch zu derselben Zeit an Joseph viele Exempel schöner herrlicher Tugenden geleuchtet haben.

115. Diese Dinge werden aber in der Kirche vom Heiligen Geiste also vorgehalten uns zur Lehre, daß wir lernen an Gott glauben und ihn fürchten, darnach aber immer kämpfen, und dem Teufel widerstehen, welcher unsern Glauben und Keuschheit immer anficht; wie er hier Josephs Glauben angefochten hat durch Unzucht. Denn wo durch Ehebruch und andere unziemliche Vermischung Zucht und Keuschheit verletzt und geschwächt werden, da folgt, daß sofort auch der Glaube verletzt wird. Deshalb lasset uns fest stehen, und uns bei dem Wort, bei rechtem Glauben und gutem Gewissen halten, beide zur Rechten und zur Linken, durch gute Gerüchte und böse Gerüchte. Es hat zwar Joseph bis daher einen guten Namen gehabt wegen seines feinen ehrlichigen Lebens, dadurch auch der Frau Herz selbst wunderbarlich entbrannt und bewogen worden zur Unzucht, und ist wahrlich deshalb ein seliger und glücklicher Mensch gewesen und hat ihn alles Hausgesinde lieb und werth gehalten. Nun aber, da er das gute Gerüchte verloren, muß er mit dem bösesten Gerüchte, so ihm auferlegt wird, kämpfen: da geht er auch hindurch. Dies ist nun um Josephs willen nicht geschrieben; denn er nicht gewußt, daß Mose jemals dies hat schreiben sollen: es wird aber uns vorgehalten, daß wir solchen Exempel der Gottseligkeit und Keuschheit auch nachfolgen sollen. Denn zu dem Ende soll man solche Geschichte in der Kirche behalten, dieselbe immer treiben und lernen.

116. Nachdem nun die andern Hausgenossen die Klage und Beschuldigung über Joseph von ihrer Frau gehört, da werden die, so fromm gewesen sind unter den Knechten, sich ohne Zweifel über solchen schweren Fall betrübt, oder aber werden die Klage der Frau für verdächtig gehalten haben, dieweil sie wohl gewußt, wie Joseph so ein fein, keusch, züchtig und unschuldig Leben geführt. Ach, daß ihm die Frau nur nicht unrecht thäte! haben sie gedacht. Denn es ist wohl glaublich, daß er viele von Gott und rechtem Gottesdienst werde recht gelehrt, sie auch damit von der Abgötterei und andern Sünden abgeführt haben. Derselben Etliche werden diesen erbärmlichen Fall beklagt haben, die Andern aber mögen sich auch daran wohl geärgert haben, und gesagt: Hat der Schalk das im Sinne gehabt? Solche Aergernisse werden ohne Zweifel dazu gekommen sein; denn wer kann demselben wehren? Wie wir denn zu unsern Zeiten die giftigsten Lästerungen des Schalks und Faber auch nicht verhüten können, daß sie nicht gen Rom und an den Kaiser selbst auch schreiben, und in ihren Briefen uns und unsere Lehre lästern, als sei es eine aufrührerische und verworrene Lehre. Christus selbst ist auch gekreuzigt worden als ein Mörder. Also geht es diesem trefflichen großen Doctor auch, welcher in diesem Hause Viele mit großen Ehren unterrichtet und zur Gottseligkeit und Keuschheit angehalten hat, der hat nun mit seinem Exempel das ganze Hausgesinde sehr schwerlich ärgern müssen, daß auch darüber beide Lehre und Vermahnung mit allen ihren Früchten vergehet und dahinfallt; wie alle seine Tugenden durch dies böse Geschrei, so die Frau über ihn ausgebracht, bald veraltet und verlöschen sind.

117. Das ist ein groß gewaltig Aergerniß, und ist um so viel ärgerlicher, soviel der, so gefallen ist, auch größer und trefflicher ist. Das Martyrium thut trefflich wehe noch über die vorige Marter, so auch greulich genug gewesen, nämlich, daß er von seinen Eltern weggeführt und unbekannten abgöttischen Leuten zu ewiger Knechtschaft verkauft worden ist. Er hat sich aber dazumal durch seine Tugend und Zucht wieder aufgerafft und hat sich einen herrlichen Namen und Ruhm gemacht bei dem ganzen Hausgesinde. Aber derselbe Name fällt nun auch dahin und Joseph wird also geistlich getödtet, und sterben mit ihm auch die schönen herrlichen Exempel der Tugenden, so an ihm gewesen. Siehe, werden sie gesagt haben, wie haben wir so einen bösen Meister oder Lehrer gehabt, wie sollte uns der die Gottseeligkeit und gute ehrliche Sitten gelehrt haben, dieweil er selbst so ein böser Bube ist? O! das sind rechte Passiones (Leiden), die thun sehr wehe.
\quellseite{332}

118. Ueber solch schwer Leiden klagt Christus selbst auch im 109. Psalm V. 2. 3. 4.: „Sie reden wider mich mit falschen Zungen, und sie reden giftig wider mich allenthalben, und streiten wider mich ohn Ursach. Dafür, daß ich sie liebe, sind sie wider mich.“ etc. Denn also pflegen die Feinde der Wahrheit die Frucht und das Gedeihen des Evangeliums, des Heilands der Seeligkeit bei frommen einfältigen Zuhörern zu verdunkeln, zu Schanden zu machen und zu verderben. Zu solcher Lästerung sind Eck, Cochläus, Pighius und viel Andere mehr rechte gute Meister, die sich selbst zieren mit falschem und geschicktem Lobe, und aber schänden sie, auf daß sie uns bei denen, so unserer Lehre nicht beifallen, mögen verhaßt machen. Derhalben, was das Merckwürdigste und Beste sein möchte, die Herzen guter einfältiger Leuten einzunehmen, dasselbe reißen sie heimlich weg, nämlich, die Gunst und den guten geneigten Willen der Leute, dadurch wir Viele mit dem Worte würden gewinnen und unterrichten können. Wir müssen sinken, auf daß sie schön sein. Aber der Herr wird doch zuletzt alle Heuchelei ausrotten, und die Zunge, die da stolz redet, wie der 12. Psalm V. 4. sagt. Wie D. Eck neulich ein erschrecklich Ende genommen, so er mit Worten und Werken wohl verdient, da er ohne Erkenntniß und Anrufung Gottes von Sinnen gekommen und also jämmerlich dahin gestorben ist. Eines solchen Endes müssen die Andern auch gewarten.

119. Also ist nun zu der Klage, so die Ehebrecherin und Todtschlägerin über Joseph geführt, auch das noch gekommen, daß die Hausgenossen mancherlei unbillige Urtheile haben hören lassen. Wiewohl Joseph den Argwohn von seiner Frau gar nicht hat schöpfen können, daß sie so gar toll und unsinnig werden sollte. Er hat gehofft, er wollte sein Kleid heimlich wieder bekommen und sie würde noch wieder versöhnt werden können; weiß nicht, daß sie gar entbrannt ist von Haß wider ihn, und daß sie nach dem unschuldigsten Blut dürste und Verlangen habe.

120. Wo aber bei den Hausgenossen rechter Verstand gewesen wäre, hätten sie die Klage und Beschuldigung, die das Weib geführt, etwas fleißiger erwogen, und da würden sie befunden haben, daß diese Klage nicht recht wahrscheinlich sei, sondern daß die Argumente vielmehr auf das Weib selbst zu ziehen wären. Denn das war ja nicht wohl glaublich, daß das ganze Hausgesinde aus dem Hause gegangen sein sollte, es wäre denn traun mit vorbedachtem Rath von der Frau selbst an einen andern Ort gesandt, auf daß sie also allein mit dem einzigen Jünglinge möchte im Hause sein: oder wenn es sich ohngefahr hätte zugetragen, oder daß vielmehr Joseph aus Frevel es also hätte angerichtet, daß er allein im Hause geblieben und zu ihr in die Schlafkammer gekommen wäre der Meinung, daß er sie schänden wollte, hätte er doch wohl mit Gewalt solches an ihr vollbringen können, wie hart sie sich auch dawider gelegt, und ihm vergeblich widerstanden oder gewehret hätte; wie die Historie der Römer von der Lucretia sagt. Und 5 Mos. 22, 25—27. steht auch geschrieben von einer Dirne, so auf dem Felde mit Gewalt beschlafen worden, daß man ihr nichts thun soll; denn sie hat, sagt daselbst der Text, „keine Sünde des Todes werth gethan: sondern gleichwie jemand sich wider seinen Nächsten erhöbe, und schlüge seine Seele todt, so ist dies auch. Denn er fand sie auf dem Feld, und die vertraute Dirne schrie, und war niemand, der ihr half.“ Denn wo die böse Lust zur Unzucht noch mit Unsinnigkeit vermischt wird, da bricht sie hindurch und kann leichtlich Alles überwinden.

121. Wo nun Joseph hätte einen Fürsprecher gehabt, der ihn vertreten hätte, der hätte seine Frau Lügen strafen können, und die Argumente, so sie angezogen, umkehren und wider sie selbst brauchen. Denn wo er ihr hätte Gemalt gethan, oder solches jemals in Sinn genommen, daß er es thun wollte, würde er nicht sein Kleid bei ihr gelassen haben, damit der Ehebruch nicht möchte an den Tag kommen, oder aber daß er dessen mit keinerlei Wahrzeichen möchte überwunden werden, daß er hätte Gewalt an sie legen wollen. Darum ist es aus der Erzählung selbst klar und offenbar, daß das Weib schuldig, Joseph aber unschuldig ist. Danach so wäre auch das Weib des Vermögens nicht gewesen, daß sie ihm das Kleid mit Gewalt hätte nehmen und behalten können.

122. Mose aber zeigt an, daß Joseph eine sonderliche Absicht müsse gehabt haben, daß er lieber hat wollen davonlaufen, als daß er mit der Ehebrecherin hätte streiten wollen. Denn sie wird ohne allen Zweifel versucht haben, mit vielerlei Schmeichelworten, Küssen, oder sonst unzüchtigen Geberden sie immer gekonnt hat, ihn bei sich zu behalten und zu erweichen. Nun ist es aber ein schwerer Kampf, daß dermaßen ein Jüngling mit einem schönen Weibe kämpfen soll. Darum ist es eine große Beständigkeit und der Klügste Rath gewesen, daß Joseph gedacht, es wäre besser davon zu laufen, denn daß er einen solchen Streit halten sollte mit seines Herrn geilem und schmeichlerischem Weibe. Denn die Jugend ist an sich selbst zur Unzucht geneigt; und kommt der Teufel dann auch dazu, der macht mit seinem Odem und weibischen Schmeichelworten, daß die Kohlen brennen müssen. Darum ist das der nächste und beste Rath, daß man nur davonlaufe, wie Paulus sagt 1 Cor. 6, 18.: „Fliehet die Hurerei.“ Es ist besser, daß man sich der Gemeinschafft solcher Huren gar enthalte, als daß man lange mit ihnen streiten wollte. Und wenn dies dazumal auch Einem eingefallen wäre, so wäre dasselbe auch sehr nütze und gut gewesen, Joseph damit zu entschuldigen und seines Herrn Weib zu beschuldigen.

123. Sie hat ferner unter ihre Klage auch etliche Lügen gemengt, da sie sagt, wie sie habe gerufen; denn daß sie solches erdichtet, ist jedermann wohl bewußt. Also ist keine Sünde allein, sondern zieht immer eine die andere nach sich. Und das pflegt zu geschehen in allen Sündenfällen. Wer nicht bald wieder aufsteht, Buße thut und sich bessert, derselbe fällt bald dahin in eine andere Anfechtung.

Nach der Sünde der Unzucht und Ehebruchs folgt bald die Lüge, darnach der Todtschlag, und daß man pflegt unschuldig Blut zu vergießen, und endlich folgt dann auch die Verzweiflung. So geht es mit den Kezern auch. Niemand wird ein Ketzer um Eines Artikels willen: sondern wo sie in Einem Artikel des rechten Verstandes fehlen, daß sie sich den Teufel haben bethören lassen, da irren sie auch bald in andern Artikeln mehr. Also kommen in den Ehebrechen viele Sünden zusammen über einen Haufen, als, Lüge, Diebstahl, Todtschlag und dergleichen andere mehr. Derhalben ist die Erzählung dieser Ehebrecherin nicht allein falsch, kraftlos und sich selbst zuwider, sondern ist auch erdichtet, falsch und erlogen Ding.

\begin{quote}
\textbf{V. 16. Und sie legte sein Kleid neben sich, bis sein Herr heim kam.}
\end{quote}

124. Mose hat an diesem Ort ein sehr nachdrücklich Wort gebraucht, welches 2 Mos. 20, 11. auch steht, nach: „Es ruh[e]te der Herr am siebenden Tage.“ Und will hier sagen, sie habe das Kleid mit Fleiß bewahret, damit es niemand anrührete oder beiseite legte. Und wird auch damit ihre Begierde angezeigt, daß sie Lust gehabt zum künftigen Todtschlage, und daß sie in ihrem Herzen nicht habe ruhen können, bis sie Joseph umgebracht hätte, da sie gesehen, daß ihr nunmehr alle Hoffnung abgeschnitten gewesen, ihres Gefallens Unzucht mit ihm zu treiben.

\begin{quote}
\textbf{V. 17. 18. Und sagte zu ihm eben dieselben Worte, und sprach: Der hebräische Knecht, den du uns hereingebracht hast, kam zu mir herein, und wollte mich zu Schanden machen. Da ich aber ein Geschrei machte, und rief; da ließ er sein Kleid bei mir, und flohe hinaus.}
\end{quote}

125. Diese schändliche Hure redet sehr häßlich und schmählich, damit sie sich schmücken und den Jüngling verhaßt machen möge. Darum brauchte sie eine wunderbare und gar teuflische Dialektik und Beredsamkeit, die Sünde, deren sie ihn fälschlich beschuldigt, nur groß und schwer zu machen. „Dieser Ehebrecher, dein Knecht“, sagt sie, von welchem du so viel gehalten, „den du uns hereingebracht hast.“ Als wollte sie sagen: Wir haben sonst auch ägyptische Knechte, deren nie einer eine böse That begangen hat: dieser Fremdling aber, der nicht werth ist, daß er bei uns wohnen soll, der von seinen eigenen Mitbürgern verstoßen und verkauft worden ist, hat eine solche unbillige böse That thun dürfen. Also nimmt sie das von seinem Namen und Vaterlande, auf daß sie damit die Sünde groß machen möge; denn die Israeliten sind bei den Heiden sehr verhaßt gewesen, wie die Christen auch sind.
\quellseite{333}

126. Darnach nimmt sie auch ein Argument von ihrer Person, nämlich also: Wenn irgend ein egyptischer Knecht oder auch dieser Hebräer etwa eine Nachbarin oder Magd geschwächt hätte, könnte ich es ihm leichtlich verzeihen und zu gute halten, und es wäre auch die Schmach so groß nicht: aber der Frau oder Hausmutter selbst etwas Unzüchtiges zumuthen, das ist eine große Sünde und die nicht leichtlich kann verziehen oder gebüßt werden. Darum liegt der Nachdruck in dem Wort „mir“. Als wollte sie sagen: Ist das nicht ein unbillig Ding, daß eine solche vornehme Frau vermessen soll geschmäht werden, die ich doch in diesem ganzen Königreich die keuscheste Lucretia und ein Creuzel guter Zucht allen andern Matronen gewesen bin. Und dies ist, sagt sie zu ihrem Herrn, gar deine eigene Schuld, dieweil du diesem Fremdling so leichtlich geglaubt, und alles vertrauen hast, was im ganzen Hause zu thun ist, dazu auch die Ehre der Hausfrau selbst. Es ist gewiß daraus abzunehmen, daß du diesen hebräischen Knecht in größerem Ehren und Würden gehalten hast, deß Laster und Uebelthat dir vielleicht nicht unbewußt gewesen, als mich selbst. Bin ich denn bei dir so gar verachtet und verschmäht, daß du davor keine Scheu getragen, mich einem solchen losen, schändlichen, unflätigen Knechte zu vertrauen, unangesehen, wie mir es auch an meiner Wohlfahrt und Ehre gehen möchte? Höre aber zu, wie große Arbeit ich gethan habe. So viel Stärke war an mir nicht, daß ich ihn mit Gewalt hätte überwinden und hinausstoßen können; aber da ich ein Geschrei machte, ist er erschrocken, und ließ sein Kleid bei mir liegen und floh hinaus.

127. Dies ist die schwere Klage, so Potiphars Weib über Joseph geführt hat, und da der Mann dieselbe gehört, hat er sich darüber sehr erzürnen müssen. Denn es ist nichts so unleidlich, als daß Einer hören soll, daß seinem Weibe Gewalt geschehen sei und daß ihr jemand nach ihrer Ehre gestanden habe: da wälzt Fleisch und Blut. Darum kann daselbst die Vernunft nicht bedeuten oder bei sich schließen, daß auch der andere Theil möchte gehört oder gefragt werden; welches wahrlich billig geschehen sein sollte. Denn das ist ja verdächtig, daß sie ihm ein fremd Kleid zeigt, und man hätte fragen mögen, wer denn Jünglinge das Kleid ausgezogen? Und wiewohl sie der Sache damit hätte einen Schein geben wollen und sagen, er selbst habe das Kleid von sich geworfen, auf daß er so viel fertiger wäre davon zu laufen: doch wo Potiphar ein jeglich Stück in der Sache etwas fleißiger erwogen hätte, und dazu den Jüngling auch zu sich gefordert und auf der Frau Klage antworten lassen, so hätte der Jüngling gar eine herrliche Vertheidigung oder Entschuldigung entgegen setzen können. Endlich, wo er ihr hätte wollen Gewalt thun, so wäre er nicht davon geflohen. Dem Geschrei aber hätte er leichtlich rathen und wehren können; wie Tertullianus zu Lucretia sagt: Ferrum in manu est etc.: Siehe, ich habe ein Schwert in der Faust; wirf du nun schreien, so sollst du doch von meinen Händen sterben. Aber Potiphar wird überwunden durch die Zuneigung oder Liebe, so ein Ehemann zu seinem Weibe hat, dazu bewegen, daß sie Sache an sich selbst so gar unbillig und böse ist: das, sage ich, hat dem Potiphar im Wege gelegen, daß er nicht fleißiger durch alle Umstände nach der Sache gefragt.

\begin{quote}
\textbf{V. 19. Als sein Herr hörte die Rede seines Weibes, die sie ihm sagte und sprach: Also hat mir dein Knecht gethan; ward er sehr zornig.}
\end{quote}

128. Potiphar hat nicht des Jünglings, sondern seines Weibes Wort allein gehört; darum gewinnt sie die Sache, und bewegt ihres Mannes Herz, daß er zornig und ungeduldig wird, und er scheint billig darüber zornig geworden zu sein. Wo sind nun die großen reichen Verheißungen, daß er seine Heiligen und Gläubigen liebe, bewahre, behüte wie seinen Augapfel, daß er barmherzig sei und wohlthue in tausend Glied denen, die ihn lieben? Denn jetzt wird Joseph gar verlassen und wird aller Hülfe und Trostes beraubt. Erstlich ist er von seinen Eltern hingerrissen und weggeführt worden, daß er in fremder abgöttischer Leute Knechtschaft gekommen ist. Was ist das für eine Liebe? Darnach hat er eine Linderung gefühlt in solcher Knechtschaft, nämlich, da ihm sein Herr alles, was er hatte, unter seine Hände gethan: aber nach derselben Erquickung ist er wiederum viel schwerer und erbärmlicher gefallen, als zuvor je geschehen. Denn da wird er hingeführt, daß er wie ein Ehebrecher hat sollen gestraft werden, und muß also des schändlichsten Todes sterben darum, daß sein gut Gerücht, dadurch er bei seinem Herrn und dem ganzen Hausgesinde Gunst gehabt und das Lob der Gerechtigkeit, Heiligkeit, Keuschheit und Klugheit erlangt hatte, zugleich mit ihm gar verlohren ist.

129. Dies ist wahrlich das größte und schmerzhafteste Kreuz, und zumal ein unbilliger Lohn, so ihm für seinen treuen Dienst, den er ganze dreizehn Jahre treulich geleistet, gegeben wird: da wird ihm nicht allein seine Mühe und Arbeit, sondern auch seine schönen Tugenden mit solcher Strafe und schändlichem Tode bezahlt. Nun diene der Welt, wer diesem will: aber es gehört kein Diener in die Welt, denn der Teufel und seine Mutter. Joseph hat seinem Herrn seine Habe und Gut fleißig verwahrt und das Haus ganz treulich regiert: aber diese schändliche Hure gibt ihm dafür den ärgsten Lohn.

130. Derhalben dürfen fromme Christen auch keinen andern Lohn von der Welt erwarten, sollen auch nicht denken, daß sie um der Welt willen in ihrem Amt treu und fleißig sein wollten. Denn sie ist dessen nicht werth, daß die Christen ihr auch mit einem Fingerlein dienen sollten, sondern daß vielmehr der Teufel mit allen seinen Engeln über sie herrsche; welches er auch fürwahr erschrecklicher Weise thut. Aber wie Paulus sagt Röm. 8, 20.: „Die Creatur ist der Eitelkeit unterworfen, und um deß willen, der sie unterworfen hat“, dienen wir, lehren, vermahnen, leiden, trösten, und thun, was uns von Gott befohlen ist, nach dem mehreren Theil unwürdigen Leuten, da wir mit unserm getreuen Dienste nichts gewinnen, denn Haß, Neid, und daß wir dazu noch wohl in das Elend getrieben werden. Und ist unser ganzes Leben Anderes nichts denn nur eitel Verlust aller Wohlfahrten.

131. Derhalben darfst du nimmermehr hoffen, daß es dahin kommen werde, daß die Welt deine Treue und Fleiß erkenne und dieselben belohne. Denn sie thut ganz das Gegenteil, wie dies Exempel anzeigt. Darum siehe zu, daß du mit deinem Dienst und Leben anders wohin sehest, denn auf der Welt Gunst und Gnade, welche gar bald in grimmigen Zorn kann geändert werden, und denke also: Wenn ich zum Pfarrherrn, oder Schulmeister, oder sonst zum andern Amte berufen und bestellt werde, so will ich mein Amt fleißig ausrichten, nicht daß ich dafür den Lohn wolle erwarten von denen, welchen ich diene; denn darauf habe ich gar keine Hoffnung, daß mir dieselbigen sollten dankbar sein: sondern ich will der Gütigkeit meines himmlischen Vaters nachfolgen, der seine Güter austheilt, und gibt Gold, Silber, Früchte, Friede, Gesundheit auch den undankbaren und bösen Menschen, und will denken an die Vermahnung Christi, der da sagt Matth. 5, 48.: „Darum sollt ihr vollkommen sein, gleichwie euer Vater im Himmel vollkommen ist“, das ist, ihr sollt auch den Bösen, Unwürdigen und Undankbaren dienen, und nicht allein den Würdigen, Dankbaren und Frommen. Es werden sich aber doch noch etliche Wenige finden, die werden es erkennen und dankbar dafür sein; die Andern aber werden dir auch wohl nach dem Leben stehen; wie dies Exempel Josephs lehrt, was man in der Welt erwarten solle auch für die allerhöchsten Wohlthaten, nämlich, Bande und Gefängniß.

132. Also sind jetzt Viele, die im Anfange das Evangelium mit aller Treue gemeint und dasselbe heftig lieb gehabt, welche nunmehr beide der Lehre und auch den Lehrern spinnfeind sind. Darum müssen wir unser Leben auf diese Weise einrichten, daß wir uns auf dergleichen Ausgang und Frucht unserer Arbeit bereiten. Und wo uns etwas Anderes widerfahren wird, darauf wir doch nicht gehofft hatten, das mögen wir alles für eitel Gewinn rechnen. Du wirst die Welt nicht ändern, sondern sollst also sagen: Ich muß hindurch; denn ich habe dies Amt nicht angenommen um der undankbaren Welt, sondern um Gottes willen. Wie der Weinstock reichlich und überflüssig den seinen lieblichen Wein gibt auch den allerbösesten Menschen: also wir auch, dieweil wir Gottes Kinder sind, sollen wir solcher Mildigkeit und Güte unsers himmlischen Vaters und dem Exempel der andern Creaturen auch nachfolgen, und über solcher Undankbarkeit der bösen Welt nicht ungeduldig werden, sintemal uns Gott in dies Leben hat lassen geboren werden, da Jeder an seinen sonderlichen Ort gesetzt und bestellt ist, daß wir der Welt dienen und die Auserwählten versammeln sollen zur ewigen Seligkeit. So nun etliche fromme Leute sind, die unsern Dienst und Amt erkennen, dessen haben wir uns zu freuen; ist es aber der Fall, daß der mehrere Theil uns verflucht, verdammt und noch wohl dazu erwürgt, sollen wir wissen, daß solches unser Lohn sei, und daß wir damit sollen gleichförmig werden dem Sohne Gottes und allen Heiligen, welche auch also für ihre höchsten Wohlthaten die allergrößte Undankbarkeit der Menschen erfahren und empfunden haben.
\quellseite{334}

\section[Fünfter Theil]{Fünfter Theil. Wie Joseph ins Gefängniß geworfen worden}

\begin{quote}
\textbf{V. 20. Da nahm ihn sein Herr, und legte ihn ins Gefängniß, da des Königs Gefangene innen lagen; und er lag allda im Gefängniß.}
\end{quote}

133. Da hat er seinen Lohn. Joseph wird in das Gefängniß geworfen, nicht daß sein daselbst wohl gewartet werden soll und daß er es darin soll gut haben, sondern daß er des Tages erwarten soll, da man über ihn werde Gericht halten, und ihm die Strafe auferlegen, so über die Ehebrecher beschlossen war. Denn sein Herr hat gedacht, er wollte ihn mit dem Gefängnisse und Leibesstrafe also peinigen und endlich gar tödten lassen. Derselben stirbt er gleich jetzt dahin durch ein sehr erbärmlich Unglück, und dazu in der schändlichsten Sache, als hätte er das ganze Haus mit Ehebruch und schändlichster Unzucht befleckt und verunreinigt, so doch zuvor dafür ist gehalten worden, daß er desselben Hauses Säule und Heil gewesen ist.

134. Wo ist denn nun Josephs Gott, der ihn erlösen und erretten möge? Warum verläßt er ihn also? Antwort: Er verläßt ihn nicht, sondern versucht ihn zum Exempel und Trost allen andern Heiligen und Gläubigen, auf daß wir lernen beharren, hoffen und auf Gott vertrauen, es gerathe wie es wolle, in großer Geduld, in Trübsal, in Nöthen, durch Ehre und Schande, durch böse Gerüchte und gute Gerüchte, wie Paulus ermahnt 2 Cor. 6, 4. 8. Denn das ist gewiß, daß gleichwohl die Kirche und die Engel Gottes uns ehren müssen: daran sollen wir uns genügen lassen und sollen nicht kleinmüthig werden. Denn uns der Teufel und die Welt mit ihren Lügen und Lästerungen schelten werden, oder uns gar tödten und verderben, sieht doch Gott mit offenen Augen auf unsere Angst und Noth, lacht uns an und hat Lust an unserer Tugend und Sieg in unserer Trübsal.

135. Derhalben sollen wir gewißlich wissen und es dafür halten, daß Gott mit uns nicht zürne, wenn wir betrübt und geängstigt werden. Denn also hat Joseph auch gedacht: Es ist zwar dies Gefängniß, diese Schande und Schmach meinem Fleische sehr verdrießlich; aber ich will auf die Verheißungen sehen, die ich gehört, und auf die Exempel, die ich selbst gesehen habe, und will mich daran halten und vom Herrn der Hülfe warten. Denn er ist in die alleräußerste Noth und Gefahr gekommen, es kann nun nicht ärger mit ihm werden. Ueber alles vorige Unglück, so er erlitten, muß er nun auch mit Gewalt getödtet werden; darum ist er ein rechter Märtyrer. Denn jetzt spricht seine Frau von nichts Anderem, denn davon, daß er sterben soll: und er wird auch darauf in das Gefängniß geworfen, und weiß nun fürwahr, daß er darin sterben und verderben muß. Denn es war ein Gefängniß, darein man die Uebelthäter, so am Leibe sollen gestraft werden, gelegt hat; wie hernach aus diesem Gefängniß der Amtmann über die Bäcker geführt und an den Galgen wird gehängt werden.

136. Darum hat er wie ein rechter Märtyrer sein Blut vergossen um einer Hure willen; wie Johannes der Täufer um der Herodias willen auch sterben mußte. Jetzt ist er schon aus dem Lande der Lebendigen genommen noch seinen eigenen und Anderer Urtheil und Empfindung. Und da hat er wahrlich vor Gott und den Engeln das köstlichste Opfer geopfert. Wie sollen sich die Engel gefreut haben? Aber sein Vater sitzt dieweil daheim, wird von Leid und bittern Thränen gar verzehrt, während sein Sohn von einer Stunde zur andern im Gefängniß die Leibesstrafe erwartet.

137. Da sollen wir aber unsers Herrn Gottes Güte und Kraft erkennen lernen, der die Todten wiederum auferweckt, und die, so schon zum Tode verdammt sind, wiederum lebendig macht. Und das hat er mit diesem Exempel lehren wollen, was er für ein Gott sei und wie er wolle geehrt sein, nämlich mit Leiden und daß man in Leiden auf ihn harre; wie die Schrift hin und wieder dazu ermahnt, als, im 27. Psalm V. 14.: „Harre des Herrn, sei getrost und unverzagt, und harre des Herrn“; item Hab. 2, 3.: „Die Weissagung wird ja noch erfüllet werden zu seiner Zeit“ etc. „Ob sie aber verzeucht, so harre ihrer; sie wird gewißlich kommen, und nicht verziehen.“ Darnach halte ich dafür, daß auch dazu gekommen sei, daß je zum wenigsten Etliche unter dem Hausgesinde werden Mitleiden mit ihm gehabt und sonst andere Dienste und Willen ihm erzeigt haben, damit sie das schwerste Kreuz, so Joseph getragen, zu lindern versucht.

138. Und wir werden bald hernach hören, daß auch der Kerkermeister ihn lieb gehabt. Daraus denn hervorgeht, daß sie nicht alle der Klage und Beschuldigung, so die Frau über Joseph geführt, Glauben geschenkt haben. Wiewohl er aber nun noch nicht gar erlöst ist, hat es ihm doch an Erquickung und Trost nicht gemangelt, damit er sich etwas hat erholen können. Und nun folgt ein sehr schöner Text, der da zeigt, daß Gott auch auf das Gefängniß, auf die Bande und auf den allerschändlichsten Tod seine Augen gerichtet habe, als auf die schöne helle Sonne. Die Hölle, der Tod und die Verdammniß stehen dem Herrn offen und vor ihm ist nichts dunkel oder todt. Joseph ist zwar vor der Welt verdammt: aber unser Herr Gott urtheilt viel anders als die Menschen; wiewohl weder Joseph noch die andern Heiligen solche Gnade göttlichen Einsehens in der Anfechtung fühlen oder verstehen. Denn wo wir es dafür gewißlich halten könnten, daß Christus der rechte Bischof unserer Seelen wäre, und wären alle mit einander und ein Jeder insonderheit also gesinnet, daß wir sagen könnten: „Ich weiß, daß der Herr der ganzen Welt, Schöpfer aller Dinge und ein Ueberwinder des Teufels der rechte Bischof ist, der auf meine Seele, und Leib und Leben sieht und daß er auch meinen Tod und die ganze Hölle ansieht, so würden wir wahrlich alles Leiden, und dazu alle Trübsal und Noth überwinden können.

139. Mose aber braucht hier eine Repetition, daß er wiederholt und sagt: „Und Joseph lag allda im Gefängniß,“ auf daß er damit anzeige, daß Joseph alle Hülfe, alle Hoffnung und Trost genommen und entzogen worden sei, und daß niemand da gewesen, der sein etwas gedacht hätte. Derhalben kann man an Joseph nichts Anderes sehen als den Tod, und daß er seines Lebens, seines guten Gerüchts und aller Tugenden, welche kurz zuvor das ganze Hausgesinde an ihm gelobt und deren sie sich verwundert, nun verlustig werden muß. Da kommt nun Christus und leuchtet in die Hölle hinein mit gnädigen Augen.

\begin{quote}
\textbf{V. 21–23. Aber der Herr war mit ihm, und neigte seine Huld zu ihm, und ließ ihm Gnade finden vor dem Amtmann über das Gefängniß, daß er ihm unter seine Hand befahl alle Gefangene im Gefängniß, auf daß alles, was da geschah, durch ihn geschehen mußte. Denn der Amtmann über das Gefängniß nahm sich keines Dinges an; denn der Herr war mit Joseph, und was er that, da gab der HErr Glück zu.}
\end{quote}

140. Christus, der ein Bischof ist unserer Seelen und der auch in die Hölle und Tod hinein sehen kann, sieht allein und gibt auch allein Acht auf Joseph, und ist darüber froh, daß ihm ein so schön Opfer gebracht wird. Derselben, wenn es scheint, daß Alles verloren sei, und keine Hülfe oder Trost mehr vorhanden ist, da hebt sich die Hülfe des Herrn an, der da sagt: Siehe, lieber Joseph, ich bin bei dir; laß dir daran genügen, daß ich auf dich sehe. Darum ist dieser Text voller Trostes und Freude, da er sagt: „Der Herr war mit Joseph.“ Denn wiewohl er noch nicht erlöst wird, so denkt doch schon jetzt Gott darauf, wie er ihn wolle aus dem Tode und Hölle erwecken, nachdem sein Blut nun lange genug gen Himmel gerufen, und so viel zuwege gebracht, daß der Herr ein Bischof seiner Seele geworden ist und seine Huld zu ihm neigt.

141. Das hebräische Wort, chesed, heißt Wohlthat, Gabe, Gnade, Barmherzigkeit. Es ist nicht das Wort chanan, sondern chesid, daher kommt chasid, welches Paulus deutet, daß es so viel heiße als, Wohlthat, 1 Tim. 6, 2.: „Die der Wohlthat theilhaftig sind.“ Es heißt nicht allein Gunst oder Gnade, sondern auch die Gabe des Heiligen Geistes; denn die Gnade wird allein nicht empfunden oder gefühlt, sondern die Gabe muß auch dazu kommen. Christus hat ihm den Heiligen Geist im Gefängniß gegeben, den Geist der Wahrheit, der ihn erhalten hat im Tode, in der Schande und Schmach, und hat ihm diesen Sinn eingegeben: Fürchte dich nicht, sei getrost und unverzagt, harre des Herrn.
\quellseite{335}

142. Dies ist ein ganzer und sehr großer Trost, daß Gott ihn ansieht und seine Gabe zu ihm neigt, das ist, bläst ihm ein den Geist der Stärke, des Raths, macht einen festen, starken, lebendigen Heiligen aus ihm. Denn er macht ihn lebendig mitten im Tode, er tröstet ihn in der Verdammniß, daß sein Herz schließen kann: Wie heftig auch mein Herr wider mich tobt und wüthet, so werde ich doch nicht sterben, meine Tugend und gut Gerücht wird nicht vergehen. Dies redet \unsicher{chedes} oder der Geist in sein Herz, daß er keine Gedanken an Hölle oder Tod gehabt, sondern nur an Leben und gute stille Ruhe.

143. Also wird er im Geist lebendig gemacht und auferweckt. Denn die jämmerliche böse Hure hat ihm schon mit ihren Lügen das Leben genommen, und dazu sein Gewissen erschrecklicher Weise irre gemacht, da er gesehen, daß er solcher großen Schande halben beschuldigt und verdammt worden; da ist er beide vor seinen und auch vor anderer Leute Augen und Sinnen schon gestorben, und hat damit auch seinen guten Namen und alles, was fromme Leute Gutes von ihm gehalten, verloren. Nun aber thut sich sein Gewissen wiederum auf und wird erfreut; denn Gott redet ihm in das Herz, das ist, er redet freundlich und tröstlich mit ihm, 2 Cor. 12, 9.: „Laß dir an meiner Gnade genügen“, ich gebe dir meinen Geist, daß du nicht verderben sollst, ich will dir ein fröhlich Herz geben, daß du bestehen könnest. Und dies ist der erste Trost gewesen, da das Gewissen inwendig erfreut, zufrieden gestellt und geheilt worden ist; und ist ihm also das Leben und Heil selbst wieder gegeben worden.

144. Der andere Trost ist äußerlich, da der Text weiter sagt: „Und der Herr ließ ihn Gnade finden“, das ist, hat ihn angenehm und hold gemacht in den Augen des Amtmannes, der über das Gefängniß war; denn er hat gemerkt, daß an diesem Jünglinge eine sonderliche feine gute Art wäre. Denn wo die Gabe des Heiligen Geistes den Menschen gegeben wird, leuchtet sie auch äußerlich, daß man solche Gnade auch am Gesicht, an Worten, an dem Munde und den Gebärden sehen kann. Derhalben, da die Andern, so in selbigem Gefängniß auch gefangen gelegen, gesäuft, geheult und vor Ungeduld getobt haben, ist Joseph allein gutes Muthes und wohl zufrieden gewesen, hat die Andern gelehrt, getröstet, und wiewohl er auch in derselbigen Gefahr mit ihnen gewesen, hat er sich doch nicht anders gehalten, denn als wäre er ganz frei und seiner Befreiung gar gewiß.

145. Ueber solche Beständigkeit und Freudigkeit des Herzens hat sich der Amtmann verwundert, hat gedacht, es müsse etwas Sonderliches mit ihm sein, und hat ihn gefragt, was er doch verwirkt hätte, darum er in die Haft und Gefängniß gekommen wäre? Da wird ihm Joseph die ganze Sache ordentlich nacheinander erzählt haben, und hat seine Unschuld leiblich darthun können. Darum findet er Gnade bei dem Amtmann, daß ihm derselbe die Pein etwas gelindert. Er ist nun der Hure nicht mehr unterworfen, sondern hat einen Herrn bekommen, der zwar Macht hat, ihn am Leibe zu strafen; aber er hat an seinen Werken, Geschick und Gebärden gesehen und gespürt, daß er fälschlich und mit Unrecht bei seinem Herrn müßte angegeben und vor ihm verklagt worden sein, und der Heilige Geist macht, daß er in seinem Reden, Vermahnen und Trösten holdselig gewesen. Darum hat ihm der Amtmann sein ganzes Amt befohlen, wie zuvor Potiphar ihm auch das ganze Hausregiment zu verwalten befohlen hatte.

146. Dieser Amtmann aber ist der oberste Richter gewesen, und hat Macht gehabt, die Mißthäter zu strafen; welchen die alten Deutschen einen „Grimmwalt“ geheißen haben, das ist, eine Macht des Zorns. Und da derselbe gesehen, daß dieser Jüngling in Worten geschickt war, dazu ein fröhlich aufrichtig Angesicht gehabt und von sonderlicher guter Art gewesen ist, hat er gedacht, wie er das Halsgericht etwas länger aufschieben und ihn gar los machen möchte; und hat ihn auch darum über die andern Gefangenen gesetzt. Muß also Joseph ein Herr werden, wenn er schon in die Hölle käme. Denn es hat müssen gehen, wie der Traum gelautet, darin ihm verheißen war, daß sich vor ihm beide seine Eltern und Brüder neigen sollten, 1 Mos. 37, 10. Derhalben muß er über Andere regieren nicht allein in seiner Knechtschaft bei Potiphar, sondern da er auch im Gefängniß gebunden gewesen; sintemal ihm daselbst das Amt befohlen gewesen, die Mißthäter mit dem Gericht und Schwerdt zu strafen, also daß er alles, was dem Richter zu thun gebühret hat, selbst gethan und verrichtet, hat den Gefangenen Essen gegeben, und wo sie mit einander gehaderet, hat er müssen entscheiden, die Unschuldigen freisprechen und die Schuldigen strafen: kurzum, er muß Herr sein.

147. Er hat aber dies Amt mit großen Ehren verwaltet, also daß der Amtmann über das Gefängniß von ihm keine Rechenschaft seiner Verwaltung gefordert. Er hat ihm vollkommenen und freie Macht gegeben über alles, was daselbst zu verwalten gewesen. Darum sagt auch Moise, der Amtmann habe unter seinen Händen nichts behalten, das ist, alles, was der Richter in seinen Händen oder Macht gehabt, das hat er alles Joseph befohlen und denselben zum obersten Richter gesetzt. Er hat, sagt Moise, in seinen Händen nichts behalten, wie droben von Potiphar auch gesagt ist, daß er, er hat sich der Dinge nicht angenommen, so zu seinem Amt gehörten. Denn er hatte einen Statthalter, der sehr fleißig und getreu war. Und ist nun auf diese Weise der heiligste Märtyrer Joseph wiederum aus dem Tode zum Leben gekommen, und wird nach der Auferweckung vom Tode auch zum Himmel fahren und zum Herrn oder Regenten über Egypten gesetzt werden.

148. Und dies ist die Art und Weise, so Gott zu halten pflegt, da er seine Heiligen führt und regieret. Derhalben soll man das Exempel dieses Jünglings mit Fleiß merken; denn es ist voll allerley Tugend und Gottseligkeit. Daran zu sehen ist, wie doch unser Herr Gott diesen sehr lieben und löblichen Edelstein polirt und zugerichtet habe, ja, wie er damit, daß er ihn geschlachtet und getödtet, ein schön, lieblich, angenehm Opfer vollbracht habe. Denn diese Dinge werden eben der Ursache halben der Kirche Gottes und den Gewissen gottselliger Leute vorgehalten, daß sie sollen verstehen lernen, wie es sich verhalte mit unsers Herrn Gottes Rath und Werken, nämlich also, wie er zu Mose 2 Mos. 33, 20. 23. sagt: „Mein Angesicht kannst du nicht sehen, du wirst mir aber hinten nach sehen.“ Und Hagar sagt droben im 16. Cap. V. 13.: „Gewißlich, hie hab ich den Rücken gesehen deß, der mich siehet.“ Da wird nichts Anderes aus.

149. Der Anfang göttlicher Wirkung ist dieser, daß er sein Angesicht von den Seinen abwendet, und daß er scheint, nicht Gott, sondern der Teufel zu sein. Also ist in dieser Historie sein Angesicht auch zu der Hure und den Tyrannen gewendet, dieselben erhält und beschirmt er allein, und achtet Josephs nicht, wie Jeremia auch klagt Cap. 12, 2.: „Du lässest sie viel von dir rühmen“ etc. Darum überheben sie sich und rühmen, Gott sei bei ihnen, sei ihnen günstig und gnädig. Hier wohnt Gott, schreien sie; aber Joseph, Jakob, Abraham haben dies Angesicht nicht auf sich gewendet. Da heißt es: Hier wohnt der Teufel: „Mein Angesicht kannst du nicht sehen“ etc., 2 Mos. 33, 20.

150. Denn also pflegt Gott die Seinen zu führen und zu regieren, wie im 77. Psalm V. 20. geschrieben steht: „Dein Weg war im Meer, und dein Pfad in großen Wassern, und man spürete doch deinen Fuß nicht.“ Und Jes. 30, 20. 21.: „Deine Augen werden deinen Lehrer sehen, und deine Ohren werden hören das Wort hinter dir her sagen.“ Und Christus sagt zu Petro Joh. 13, 7.: „Was ich thue, das weißest du jetzt nicht; du wirst aber hernach erfahren.“ Du willst mir in das Angesicht sehen, du willst, ich soll thun, was dich dünkt nüge und gut zu sein: aber ich will es also machen, daß du dich sollst dünken lassen, als habe dies etwa ein Narr gethan und nicht Gott. Du sollst mir in Stücken sehen, nicht in das Angesicht. Du sollst meine Werke und meinen Rath nicht sehen, damit ich dich nach meinem Wohlgefallen mache und wiederum mache. Es soll dich närrisch dünken, du sollst solche meine Werke nicht anders aufnehmen, denn als ob es der Tod und Teufel selbst wäre.

151. Derhalben lasset uns diese Regel und Ordnung, so Gott in der Regierung seiner Heiligen zu halten pflegt, lernen. Denn ich habe auch oftmals versucht, unserm Herrn Gott gewisse Weise vorzuschreiben, nach der er sich halten sollte in der Regierung seiner Kirche, oder sonst anderer Dinge. Ach Herr! habe ich gesagt, ich wollte gern, daß dies nach der Ordnung möchte geschehen und daß es möchte einen solchen Ausgang gewinnen. Aber Gott that, das dem, so ich gebeten hatte, stracks zuwider war. Da gedachte ich daselbige: Nun ist ja gleichwohl mein Rath nicht wider Gottes Ehre, sondern wird sehr nütze sein, daß dadurch der Name Gottes geheiliget, sein Reich gesammelt und gemehret werde und daß auch die Erkenntniß seines Wortes weiter ausgebreitet werden möchte; kurzum, es ist ja ein sehr schöner Rath und sehr wohl bedacht. Aber Gott hat ohne Zweifel solcher Weisheit gelacht und gesagt: Wohlan, ich weiß wohl, daß du ein verständiger gelehrter Mann bist; aber ich habe den Gebrauch noch nie gehabt, daß mich Petrus, oder D. Martinus, oder wer es auch sein möchte, lehren, unterrichten, regieren oder führen müßte. Ich bin nicht ein solcher Gott, der ich mich also wolle lehren, oder regieren lassen, sondern der da pflegt Andere zu führen, regieren und zu lehren.
\quellseite{336}

152. Nun ist es aber sehr verdrießlich, daß sich meine Weisheit dermaßen nur allein soll lehren und leiten lassen, und daß ich dieselbe soll tödten. Und sind wahrlich Viele darum greulich gefallen, da sie diese Tödung nicht haben ertragen können. Denn also thun die Sacramentenschwärmer. Die können Gott fein lehren, wie doch das zugehen möge, daß der Leib und Blut Christi im Brod und Wein sein könne, so doch Christus gen Himmel gefahren ist? wenn sie also denken: Dieweil ich es nicht begreifen kann, daß im Abendmahl des Herrn der Leib und Blut Christi gegenwärtig sei, darum will ich Gott eine Weise und Maß vorschreiben, dadurch er könnte da sein.

153. Also zwingen sie Gott fein, daß er sich nach ihrer Ordnung und Lehre halten soll, wollen ihm einen sonderlichen Ort im Himmel zuweisen, und wissen doch nicht, was der im Himmel der Himmel sei. Warum thust du nicht viel lieber Augen und Ohren zu und hörst auf das Wort, und wenn du das gehörst, folgst demselben und bist damit zufrieden, wie im 37. Psalm V. 7. gesagt wird: „Sei stille dem Herrn und warte auf ihn“ (er soll nicht auf dich warten, dessen unterstehe dich nur nicht); so geht es recht hinaus, alsdann gibt Gott viel mehr, als wir mit unserm Rath und Weisheit ausrichten können. Wenn die Schwärmer nur Augen und Ohren zuthäten, und bei sich selbst also schlössen: Siehe, ich will das Brod und den Wein nehmen, und glauben, daß ich da den Leib Christi esse und sein Blut trinke wahrhaftig, so könnten sie von ihrem Irrthum leichtlich erlöset werden. Dieweil sie aber die Worte Christi nach mathematischer Weise abmessen und verstehen wollen und viel vom Himmel und Erde disputiren, darum können sie den rechten wahren und eigentlichen Verstand der Worte Christi nimmermehr erreichen, dieweil sie durch das Urtheil der Vernunft, welches blind ist, gar bethört sind.

154. Derhalben thue die Augen zu, so wird dir alles Disputiren vom Himmel und Erde verschwinden, ja, wie sie vor sechstausend Jahren nichts gewesen, also werden sie jetzt auch, wenn du deine Augen zugethan, gleich als wären sie gar zunichte geworden, deinen Glauben und die Worte Christi nicht hindern. Denn das Wort, welches Alles geschaffen hat, soll man gegen die Creatur halten, welche gar nichts ist, wenn man sie mit dem Wort vergleichen wollte. „Er gebeut, so wirds geschaffen“, sagt der 148. Psalm V. 5.

155. Dies ist unseres Herrn Gottes Weise, so er stets für und für gehalten, deren man sich billig zu verwundern hat, nämlich, daß er Alles aus Nichts mache und Alles wiederum zunichte mache. In dieselbige Weise soll man sich gewöhnen und dem Wort anhangen, daß wir die Dinge, so den Glauben verwirren oder verhindern wollen, wie großen Schein dieselben haben mögen oder wie groß sie auch an sich selbst sind, nur aus den Augen hinweg thun. Wo Joseph gesagt hätte: Herr Gott, laß mich zu Hebron bei meinem Vater bleiben, oder hätte wider Gott getobt wegen der unbilligen Anklage und Gefängniß, und Gott hätte ihm seine Begierde erfüllt, und ihn dessen gewähret, daß er nicht wäre in solche Dienstbarkeit gekommen, auch nicht in das Gefängniß geworfen worden wäre, nach der Weisheit seiner Vernunft, welche unserm Herrn Gott diesen Gehorsam ganz schwerlich leisten kann: so wäre er nimmermehr zu solcher großen Herrlichkeit erhoben worden. Dieweil er aber voll heiliges Geistes gewesen, hat er den Willen des Herrn mit Geduld ertragen, bis er endlich ist erhöht worden, nicht allein ihm selbst zum Trost und Ehre, sondern auch dem Egyptenlande zu zeitlichem und geistlichem Heil und Wohlfahrt: welches Landes heilsames Werkzeug er geworden ist durch diese Trübsal und Gefängniß, dieweil er den Stücken gesehen, und darauf geharrt, bis Gott sich offenbarte und sein Heil zeigte viel reicher und herrlicher, denn er selbst jemals hätte bitten oder hoffen dürfen. Darum ist unsere Weisheit wahrlich närrisch, und wir sind nicht allein Narren, sondern auch die Narrheit selbst.

156. Diese Exempel aber werden uns vorgehalten zu unserer Lehre und daß wir dadurch gestärkt werden, daß wir lernen Gott glauben und auf ihn hoffen. Der Glaube aber ist nicht derer Dinge, so man sieht, sondern derer, so man nicht siehet, Hebr. 11, 1. Wie Joseph an den Herrn geglaubt, den er nicht sieht, und hofft, er habe bei ihm Gnade gefunden, die er doch auch nicht sieht, sondern fühlt und empfindet, daß ihm Alles widrersinnig geht; und hat doch gleichwohl endlich das angefangen, was er geglaubt und gehofft hat.

157. Warum hat er bisher wie ein Blinder in der dicksten Finsterniß gehandelt, und hat weder Gott, noch Vater, noch sonst etwas, denn allein den Tod und eitel Verderben vor seinen Augen gesehen; er ist aber bei dem Worte geblieben, daß er von seinem Vater gehört hatte: Ich bin der Herr, dein Gott und deiner Väter Gott. Das Wort ist sein Leben gewesen und daraus wird er hernach zu großer Ehre und Herrlichkeit erhoben werden.

158. Das ist nun ein Exempel der Werke Gottes, welches anzeigt, wie Gott seine Heiligen so wunderbarlich führt, daß sie lernen in der Anfechtung durch Glauben und Hoffnung Geduld zu haben; wiewohl die Hoffnung eigentlich die Geduld selbst ist. Denn unser Leben soll also verborgen sein, und soll heißen: Thue die Augen zu, und halte dich an das Wort, nicht allein wenn Noth und Gefahr vorhanden ist, sondern auch in den höchsten Artikeln der christlichen Lehre.

\chapter{Das vierzigste Capitel}

\section[Erster Theil]{Erster Theil. Von dem obersten Schenken und Bäcker, und von ihrer Gefangenschaft und Träumen überhaupt; besonders von dem Traum des obersten Schenken und dessen Auslegung}

\begin{quote}
\textbf{V. 1—4. Und es begab sich darnach, daß sich der Schenke des Königs in Egypten, und der Bäcker, versündigten an ihrem Herrn, dem Könige in Egypten. Und Pharao ward zornig über seine beiden Kämmerer, über den Amtmann über die Schenken, und über den Amtmann über die Bäcker, und ließ sie setzen in des Hofmeisters Haus ins Gefängniß, da Joseph gefangen lag. Und der Hofmeister setzte Joseph über sie, daß er ihnen dienete; und saßen etliche Tage in Gefängniß.}
\end{quote}

1. In diesem Capitel ist nichts Sonderliches, dabei man lange verweilen müßte. Denn es hält nur in sich, was die Gelegenheit gewesen, dadurch Joseph erhöhet worden ist, und wie ihm unser Herr Gott sein Angesicht geoffenbart hat.

2. Wir haben droben von dem Amtmann über die Metzger gesagt; wiewohl man es so gar eigentlich nicht wissen kann, wie es sich mit der Polizei oder Hausregiment in Egypten verhalten habe. Es ist zwar meines Erachtens dieser Amtmann über alle die gesetzt gewesen, die das Vieh geschlachtet, entweder im ganzen Königreiche, oder aber in der Stadt allein. Ihm ist befohlen gewesen, ein Aufsehen zu haben auf das Schlachten, es sei nun gleich zum Opfer, oder aber das sonst täglich im Hause gebraucht worden: er hat den Metzgern Befehl gegeben, wie sie das Fleisch verkaufen; wie sie das Vieh, so man hat opfern wollen, mästen sollten. Wir heißen es einen Marschall oder Hofmeister. Dies ist Josephs Herr gewesen, wie aus dem, so droben gehört, zu sehen ist, welcher ihn gefaßt und in das Gefängniß geworfen hat.

3. Alhie wird aber nun auch etlicher anderer Amtleut gedacht. Und ich halte dafür, daß dies nicht einfache Schenken oder Bäcker gewesen seien, sondern die Obersten, so über die andern Bäcker und Schenken alle gesetzt gewesen. Und es hätte zwar diese Stadt Wittenberg solche Leute auch wohl nöthig, damit die Bäcker, Weinschenken und Fleischer nicht so greulich uns selbst mehr als das Vieh entleibten und schlachteten. Denn sie richten sich nach keiner rechtmäßigen Ordnung, sondern schlachten, kochen und backen nur nach ihrem eigenen Gefallen. Derhalben muß das Königreich Egypten sehr gute Ordnung gehabt haben. In den kaiserlichen Rechten wird auch gedacht eines Bischofs über das Brod, das ist, der über die Bäcker gesetzt gewesen und Aufsicht über sie gehabt hat. Und in der Polizei der Stadt Nürnberg werden dieselben Aemter mit großen Ehren verwaltet.
\quellseite{337}

4. Ich halte aber dafür, daß dieser Amtmann über die Schenken zuvor des Königs Schenke zu Hofe gewesen sei, und daß der Andere Brod und Kuchen gebacken habe, so der König selbst und auch das Hofgesinde gebraucht; er ist der oberste Schenke und oberste Bäcker gewesen, wie Josephs Herr ein Oberster über die Fleischer, und Bischof oder Aufseher über das Fleisch gewesen ist.

5. Diese Beiden aber haben wider den König gesündigt. Was aber ihre Sünde gewesen sei, zeigt Mose nicht an. Ezra erzählt aus den Juden Commentarien, der Schenke habe damit gesündigt, daß eine Fliege in des Königs Becher gefallen; der Bäcker aber soll Kuchen gebracht haben, darin der König kleine Sandkörner gefunden. Das sind solche Träume, deren die Juden wohl werth sind, welche die beiden verspotten und gegen sich gering achten; rühmen nur ihre eigene Weisheit und Heiligkeit hoch.

6. Es ist aber aus dem Texte wohl abzunehmen, daß es müssen große schwere Sünden gewesen sein, sintemal der Eine seiner Sünde halben ist gehängt worden. Vielleicht werden sie dem Könige Gift unter das Essen und Trinken gemischt haben. Denn es ist ja der König so närrisch oder tyrannisch nicht gewesen, daß er um so geringer Schuld willen wider seine Diener so gar greulich hätte toben können; sondern ich halte dafür, es sei ein weiser, verständiger und frommer König gewesen, vornehmlich dieweil ihn unser Herr Gott also geehrt, daß er einen großen Patriarchen, nämlich Joseph, zu ihm gesandt, welcher große Wunderwerke thut; darum scheint es, daß er ein frommer gütiger Fürst gewesen sei. Und solche Fürsten werden leichtlich verachtet und werden von den Hofdienern für Narren gehalten, daß sie meinen, sie haben Macht zu thun, was sie nur wollen. Darum, halte ich dafür, werden sie dem Könige heimlich nachgestellt haben, ihn am Leben zu beschädigen, und da man solches erfahren und sie darüber ergriffen, sind sie deshalb in das Gefängniß geworfen worden, darin man die Uebelthäter verwahret, so am Leibe gestraft werden sollten.

7. Hieraus aber erhellt, daß Josephs Herr, der Oberste über die Bäcker, auch Kerkermeister gewesen sei, als der oberste Richter, der die Gefangenen in seiner Gewalt gehabt und jetzt nun etwas sanftmüthiger geworden ist. Im vorigen Capitel hat Mose eines andern Amtmanns gedacht über die Gefangenen: alhier aber sehen wir, daß der Herr, so den Joseph gekauft, auch vom Könige zum Kerkermeister gesetzt gewesen. Und daß derselbe Joseph verziehen und ihm das Leben geschenkt habe, ist daraus zu sehen, daß er ihn über die andern Gefangenen gesetzt hat. Und es scheint, als habe ihn der vorige Amtmann mit diesem seinem Herrn ausgeführt, welcher gesagt haben wird: Es ist ein sehr frommer Jüngling, Lieber, schone sein und nimm ihm ja das Leben nicht; denn einen solchen fleißigen getreuen Menschen wird man sehr wohl brauchen können.

8. Diese wunderbaren Werke Gottes soll man erkennen lernen und dieselben hoch rühmen. Denn er kann auch die, so Feinde sind, zum Frieden bekehren; wie Salomo in seinen Sprüchen Cap. 16, V. 7. sagt: „Wenn jemands“, der da gläubig ist, „Wege dem Herrn wohlgefallen, so macht er auch seine Feinde mit ihm zufrieden“; oder wo ja nicht alle, so werden doch etliche bekehrt. Diese Freundschaft aber oder Gunst des Herrn ist daher gekommen, dieweil, wie droben Cap. 39, 2. gesagt ist, der Herr mit Joseph war; wie im Buch der Weisheit Cap. 10, 13. 14. geschrieben steht: „Die Weisheit verließ den verkauften Gerechten nicht; sondern behütete ihn vor der Sünde, fuhr mit ihm hinab in den Kerker; und in den Banden verließ sie ihn nicht, bis daß sie ihm zubrachte das Scepter des Königreichs und Obrigkeit über die, so ihm Gewalt gethan hatten.“ Und gefällt mir der Text im selbigen Buch sehr wohl darum, daß er es der Weisheit zuschreibt, daß Joseph also wunderbarlich ist regiert und behütet worden. Der Herr, der die Weisheit ist, das ist, Gottes Sohn, der Mensch werden sollte, der war mit ihm. Denn Salomo in seinen Sprüchen nennt ohne Zweifel auch die Weisheit den Sohn im göttlichen Wesen, und wir lehren auch aus dem heiligen Paulus, daß Christus habe das Volk Israel aus Egypten geführt.

9. So ist nun die Weisheit, Gottes Sohn, der da hat Mensch werden sollen, mit Joseph im Gefängniß, und erlangt ihm Gnade, bis er erhöht wird. Dies sind Gottes Werke, die uns lehren sollten, daß wir in Anfechtung Geduld haben, und am Glauben fest halten, auf daß wir uns ja das Wort nicht nehmen lassen, es sei gleich innerlich im Herzen, da wir es bedenken, oder aber äußerlich, da man es in die Uebung bringen soll.

\begin{quote}
\textbf{V. 5—8. Und es träumte ihnen beiden, dem Schenken und Bäcker des Königs zu Egypten, in Einer Nacht, einem jeglichen ein eigener Traum; und eines jeglichen Traum hatte seine Bedeutung. Da nun des Morgens Joseph zu ihnen hinein kam, und sahe, daß sie traurig waren; fragte er sie, und sprach: Warum seid ihr heute so traurig? Sie antworteten: Es hat uns geträumet, und haben niemand, der es uns auslege.}
\end{quote}

10. Dies ist der Anfang der künftigen Herrlichkeit und Ehre Josephs. Es hat diesen Beiden in Einer Nacht geträumt und der Traum mit seiner Auslegung hat sie auch beide betroffen. Da aber Joseph gesehen hat, daß sie trauriger gewesen, denn sie zu sein pflegten, hat er diese beiden Eunuchen oder Amtmänner gefragt, was doch die Ursache solcher Traurigkeit wäre. Droben ist gesagt, daß sie nicht verschnitten, sondern des Königs Diener und große Hofjunker gewesen seien, wie deren zu dieser Zeit an des Türken Hof viele sind.

11. Der Ort aber, da die Gefangenen sind verwahret worden, ist im Hause des Obersten über die Fleischer gewesen, das ist, in dem Hause von Josephs Herrn, welches der König seiner Treue befohlen hatte. Denn aus den Historien der Könige und Jeremias ist klarlich zu sehen, daß die Könige Gefängnisse gehabt, welche sie ihren Edelleuten oder Amtmännern verliehen und befohlen. Joseph aber ist über das Gefängniß gesetzt. Darum wird ihm auch die Sünde verziehen gewesen sein, um welcher willen er von seiner Frau, Potiphars Weib, verklagt war, und vielleicht wird sie auch müssen gebüßt und bezahlt haben, da der Herr inne geworden ist, daß sie selbst der Sünde schuldig war.

\begin{quote}
\textbf{V. 8. Joseph sprach: Auslegen gehöret GOtt zu, doch erzählet mir.}
\end{quote}

12. Hieraus ist klar, daß Joseph nicht gefeiert oder unnütz gewesen sei; denn der Heilige Geist ruht nicht, sondern regt ihm sein Herz und treibt ihn dazu, daß er die Andern auch wird gestraft und gelehrt haben, wie diese zwei Eunuchen oder Amtmänner auch. Er hat nicht schweigen können, da er gesehen, daß sie zur Abgötterei so gar geneigt und sich derselben befleißen haben; wie er es denn auch nicht von der Frau hat vertragen können, daß sie ihn zum Ehebruch reizte, sondern hat ihr davon abgerathen mit einer ernsten und tapferen Predigt. Und ist daran kein Zweifel, er wird diesen auch gepredigt haben; wiewohl das Gespräch, so er mit diesen Amtmännern gehalten, mit kurzen Worten beschrieben wird. Sie haben sich beklagt, daß sie niemand gehabt, der ihnen die Träume auslegen könnte, und wo sie der Träume könnten los werden, wollten sie zu den egyptischen Priestern, zu den Wahrsagern und Sterndeutern gehen, daß ihnen dieselben die Träume ausgeleget, und sie dabei erinnerten, wie es hinfort mit ihnen gehen würde. Da hat Joseph geantwortet: Ihr werdet mit euren Wahrsagern und Abgöttern nichts ausrichten, ihr dürft auch die Ausleger nicht um zukünftige Dinge fragen, sie können euch nichts erinnern, es sind Teufelsdiener. „Auslegen gehöret GOtt zu“; bei demselben muß man die Auslegung, pithronim, suchen, und dieselbe von ihm bitten.

13. Auf diese Meinung hat er eine sehr schöne Predigt angerichtet; denn der Heilige Geist ist geschäftig und kräftig in ihm gewesen. Und er führt alsbald alles auf GOtt; wie er auch droben 1 Mos. 39, 9. zur Frau sagte: „Sollt ich wider GOtt sündigen?“ Denn GOtt war in seinem Herzen, in seinem Munde und Werken; darum sind ihm alle seine Rathschläge und Werke wohl gerathen, nach der Lehre des 1. Psalms V. 3. Und er hat das Bekenntniß des Namens Gottes nicht verhehlen können oder verschweigen; denn er ist ein Priester und Prophet. Darum muß er auch so viel Unglück leiden, wie im 116. Psalm V. 10. steht: „Ich glaube, darum rede ich. Ich werde aber sehr geplagt“ etc. Und doch wird er aus all diesem Jammer erlöst. „Auslegen“, sagt er, „gehöret GOtt zu“, denselben rufet an, dem befehlet euer Leben, und bittet, daß er euch einen Ausleger eurer Träume geben wolle.
\quellseite{338}

14. Wohlan aber, wiewohl ich weiß, daß ihr zum Aberglauben und Abgötterei geneigt seyd, erzählet mir doch eure Träume, ob ich vielleicht dieselben auslegen könnte. Denn er hat schon dazumal gefühlt, daß GOtt bei ihm war, wie droben, Cap. 39, 21., gesagt ist: Der HErr hat seine Gnade zu ihm geneigt, hat ihm den Heiligen Geist gegeben, darum hat er angefangen, alles Gute zu verstehen und zu hoffen, und hat wahrlich beiderlei Träume ganz herrlich ausgelegt. Vielleicht werden ihm bisweilen auch andere Träume vorgetragen worden seyn, und er hat zugleich auch an seine eigenen Träume gedacht von den Garben, von der Sonne, Mond und Sternen. Daraus abzunehmen, daß er in solcher Auslegung fertig und wohl versucht wird gewesen seyn; und sagt doch gleichwohl: „Träume auslegen gehöret GOtt zu.“ Welches sonderlich wohl zu merken ist. Denn er schließet nicht, daß es bei menschlicher Weisheit stehe, Träume auslegen. Denn unsere Auslegungen sind ungewiß und können fehlen.

15. Danach eignet sich auch der Teufel die Macht zu, durch Geschichte und Träume zu weissagen. Denn er kennt die Ordnung der Ursachen. Er sieht die Ursachen, so gegenwärtig und vorhanden sind, daraus schließt er, wie es ferner hinausgehen werde. Er sieht, wie die Fürsten einander hassen und unter einander heimliche Feindschaft haben, daß einer dem andern nicht gut ist; er sieht des Kaisers Praktiken; und daraus macht er seine Rechnung und schließet: Es müssen in Deutschland große Empörungen werden. Und er fehlt zwar darin nicht, wo nicht ein guter Engel solchen Unglück wehrt. Also, sage ich, argumentirt und schließet der Teufel aus sichtbaren Ursachen. Aber die göttlichen Träume kommen her aus unsichtbaren Ursachen. Darum muß man derselben Auslegung von GOtt bitten. Deshalb sagt nun Joseph: „Erzählet mirs“, ich bin GOttes Diener; vielleicht wird er mir offenbaren, was die Träume bedeuten.

\begin{quote}
\textbf{V. 9—11. Da erzählete der oberste Schenke seinen Traum Joseph, und sprach zu ihm: Mir hat geträumet, daß ein Weinstock vor mir wäre, der hatte drei Reben, und er grünete, wuchs und blühete, und seine Trauben wurden reif; und ich hatte den Becher Pharao in meiner Hand, und nahm die Beeren, und zerdrückte sie in den Becher, und gab den Becher Pharao in die Hand.}
\end{quote}

16. Dies ist eine Beschreibung eines grünen, blühenden und wachsenden Weinstocks. Der Weinstock bringt hervor die Reben, dieselben grünen darnach und gewinnen Weinblätter. Zum Dritten läßt er die Augen hervor schießen, wenn die kleinen Blätter zugleich mit den Trauben heraus brechen. Das hebräische Wort eskol heißt Traube. „Seine Trauben wurden reif.“ Sachath heißt ausdrücken, wie wir das aus diesem Ort sehen und abnehmen können. Denn dies Wort steht sonst in der Bibel an keinem Orte mehr, welcher Wörter wir droben viel gehabt, deren Bedeutung bei den Juden ungewiß ist.

\begin{quote}
\textbf{V. 12—15. Joseph sprach zu ihm: Das ist seine Deutung. Drei Reben sind drei Tage. Ueber drei Tage wird Pharao dein Haupt erheben, und dich wieder an dein Amt stellen, daß du ihm den Becher in die Hand gebest, nach der vorigen Weise, da du sein Schenke warest. Aber gedenke meiner, wenn dirs wohl gehet, und thu Barmherzigkeit an mir, daß du Pharao erinnerst, daß er mich aus diesem Hause führe. Denn ich bin aus dem Lande der Hebräer heimlich gestohlen; dazu habe ich auch überdise nichts gethan, daß sie mich eingesetzt haben.}
\end{quote}

17. Vor diesen Worten wird ohne Zweifel Gebet und Anrufung hergegangen sein. Denn dies ist nicht so gar bald geschehen, wie es Mose kürzlich beschrieben hat. Nachdem Joseph den Traum hat erzählen hören, ist er an einen Ort besonders gegangen, hat seine Knie gebeuget und gebetet. Denn die Heiligen oder Gläubigen thun nichts ohne den Segen und Gebet, fallen nicht sobald auf die Händel, so sie zu verrichten haben, sondern greifen dieselben mit GOttesfurcht an; und dieweil sie auf den HErrn hoffen, darum rufen sie ihn auch an. Und wiewohl die äußerlichen Gebärden nicht allezeit dazu kommen, beten sie doch im Herzen und mit Seufzen; wie er zuvor das auch bezeuget, da er sagt: Die Auslegungen sind Gaben GOttes; und Träume auslegen kommt nicht her von unserer Weisheit oder Kunst, sondern der Heilige Geist muß dazu kommen, ohne welchen keine Gabe oder keine Auslegung ist. Darum wird er gesagt haben: Ich will GOtt bitten, daß er mir wolle Gnade verleihen, solche Deutung hervor zu bringen, so sich eigentlich zu diesem Traum schicken und reimen will.

18. Das hebräische Wort mispath wird oftmals gebraucht, daß es Weise oder Gewohnheit heißt. Als an diesem Orte, da im Lateinischen steht, secundum judicium suum, das ist, nach der vorigen Weise; und 2 Mos. 21, 9. von dem Hausvater, der seinem Sohne seine Magd vertrauet, juxta judicium, das ist, der soll Tochterrecht an ihr thun. Und im 119. Psalm V. 132.: secundum judicium diligentium nomen tuum, das ist: „Wie du pflegst zu thun denen, die deinen Namen lieben“, und wie du pflegest mit denen umzugehen, die dich lieben; nicht, wie du pflegest die Gottlosen zu strafen.

19. Daß aber Joseph bittet und begehrt, daß der Schenke seiner gedenken wolle, und Pharao erinnern, daß er an ihm wolle Barmherzigkeit thun etc., dasselbe übertreiben die Juden und sagen, er habe daran gesündiget und habe um der Sünde willen zwei Jahre länger in diesem Gefängniß bleiben müssen. Aber er wird deshalb unbillig und nicht mit Recht getadelt und gestraft. Er hat sein Vertrauen nicht auf den obersten Schenken gesetzt; dieweil wir so oft gehört, daß er dies Kreuz mit Geduld getragen, und bereit gewesen, auch den Tod und die äußerste Noth und Gefahr zu leiden: sondern dies ist ein Exempel der Lehre, die Paulus gibt 1 Cor. 7, 21.: „Bist du ein Knecht berufen, sorge dir nicht; doch kannst du frei werden, so brauche deß viel lieber.“ Denn wir sollen uns selbst kein Unglück oder Gefahr aufladen, sondern wenn wir ohngefähr oder nach GOttes Willen betrübt oder geängstigt werden, alsdann sollen wir mit beständigem Herzen alles Unglück leiden, doch also, daß wir gleichwohl keinen guten Rath und ordentliche Hülfe verachten, dadurch wir könnten erlöset werden. Denn das heißt GOtt versuchen, wenn man die Mittel, so GOtt zeigt, dadurch uns möchte geholfen werden, verachten und derselben nicht gebrauchen will.

20. Also hat uns GOtt den Leib gegeben, nicht daß wir ihn tödten sollen mit Hungern und vielem Wachen, sondern wir sollen sein warten mit Essen, Trinken, Kleidung, Schlafen und mit Arznei, wie Paulus Röm. 13, 14. Col. 2. 23. die Heuchler straft, die ihres eigenen Leibes nicht verschonen. Wie denn ich vor Zeiten auch ein solcher schändlicher Todtschläger und Verfolger meines eigenen Leibes gewesen bin. Denn so fastete, betete, wachte und machte mich selbst matt und müde über mein Vermögen; das heißt sich selbst tödten und umbringen.

21. Wir sollen aber zu beiderlei bereit und gerüstet sein, unser Leben zu erretten und zu erhalten und auch den Tod mit Geduld zu leiden nach GOttes gnädigem Willen und Wohlgefallen. Wie es denn GOtt haben will, daß unser Leben seinen über Tod und Leben, wie Paulus sagt Röm. 14, 8.: „Wir leben oder sterben, so sind wir des HErrn“, und nicht bloß des HErrn, sondern wir sind selbst auch HErren. Denn so wir des HErrn sind, wenn wir sterben und leben, wie Christus selbst ist, so herrschen wir auch wahrlich über Leben und Tod. Wir können und sollen uns selbst des keines zueignen oder aufladen, sondern beiderlei soll man mit aller Demuth und Geduld aufnehmen. Also haben wir droben Cap. 20, 2. 11. von Abraham gehört, wie er sein Leben in guter Hut gehabt, da er Sarah geheißen, daß sie sagen sollte, sie wäre seine Schwester, damit er also von der Gewalt des Königes ungefärdet bleiben möchte. Also ist Jacob auch geflohen vor seinem Bruder Esau, da er auf ihn erzürnet war, Cap. 28., und darnach Cap. 32. schickt er ihm Geschenke, ihn damit zu versöhnen, und versucht alles, was er kann und mag, daß er nur vor ihm sicher sein möchte. Derselbe hat Isaak Cap. 26. auch gethan, ist geflohen, hat gegeben und geschenkt, damit er der Nachbarn und des Königs Abimelech Gunst und Freundschaft haben möchte.

22. Derhalben hat Joseph daran nicht gesündiget, daß er von dem Schenken des Königs begehrt, daß er ihm bei dem Könige förderlich sein wollte, damit er auch des Gefängnisses möchte erlediget werden. Denn im Anfang seines Elendes hat er bald gesuffzet und seine Brüder gebeten, daß er ja nicht verkauft werden möchte: dieweil er aber sich dazumal dessen nicht hat erwehren können, daß er von seinen Brüdern nicht in solche betrübte und erbärmliche Dienstbarkeit fremder Leute hineingegeben und verstoßen würde, hat er zwar das Kreuz, so ihm vom Herrn auferlegt, mit Geduld getragen. Nun aber, da er sieht, daß eine andere Gelegenheit sich bietet, daß er möchte frei werden, nämlich, durch diesen obersten Schenken an des Königs Hofe, dieweil er wußte, daß derselbe bei dem Könige würde in Gnaden sein, da denkt er, daß er es versuchen wolle, ob der Herr zu dieser Gelegenheit Gedeihen geben wollte und ihn dadurch aus solchem beschwerlichen Gefängniß erretten. Und wird ohne Zweifel der andere Amtmann über die Bäcker eben dasselbe auch begehrt und darum gebeten haben, wiewohl vergeblich. Wo aber nun Joseph diese Weise nicht hätte vorgenommen, daß er dadurch möchte frei werden, so hätte er Gott versucht, welcher haben will, daß wir der Mittel und Gelegenheit, so er sich bieten läßt, gebrauchen sollen, und wiederum auch gern entbehren, welche er uns versagt; und sollen also unsere Geduld im Kreuze üben, es sei gleich zur Rechten oder zur Linken.
\quellseite{339}

23. Darum ist es mit der Juden Gedicht gar ein ungereimt Ding, es ist falsch und erlogen, und ist der Rath, welchem Joseph gefolgt, zu loben, daß er den Schenken der Wohlthaten erinnert, so er ihm im Gefängniß beide leiblich und geistlich erzeigt, und bittet, er wolle derselben eingedenk sein. Denn Joseph ist ein Prophet gewesen, voll Heiligen Geistes, der nicht allein die Gefangenen gespeiset und, wie er gekonnt, ihrer gewartet hat, sondern hat sie auch getröstet und gelehrt, wie sie sollten getrost sein, gute Hoffnung haben und auf göttliche Hülfe warten. Man soll aber dies Exempel mit Fleiß merken, da Joseph diese Gelegenheit mit sonderlichem Fleiß gesucht, daß er möchte frei werden, damit er Gott nicht versuchen möchte. Denn sonst würde ihn sein Gewissen hernach selbst beschuldigt haben, als ob er dies Mittel und diese Gelegenheit seiner Befreiung verachtet und versäumt hätte.

24. Auf solche Weise, habe ich gesagt, soll man auch des Leibes warten, sein gut Gerücht, dazu Weib und Kind und den ganzen Vorrath im Hause bewahren, und soll alles thun, was sich zu Erhaltung desselben gebühren will, so viel immer möglich. Denn wo ein jeglicher seines Dinges nicht achten wollte, sein Amt beräumen, es wäre gleich ein öffentlich oder privat Amt, möchte ich auch sagen, wenn ich mein befohlen Amt liegen ließe: Was ist es vonnöthen, daß wir viel predigen, lehren und taufen? Nun hat uns ja Gott nicht dazu geschaffen und in diese Welt gesetzt, daß wir im Müßiggang, Wollust und Ueberfluß leben sollen, und vermahnen alle Sorge nur auf Gott legen, als müßten wir gar nichts thun und nur gute faule Brüder sein: sondern wenn ein Jeder sein Amt mit Fleiß ausrichtet, alsdann soll man Gott das Gedeihen und wie es alles hinausgehen und gerathen möge, anheimstellen und befohlen sein lassen. Also thut der Ackermann auch: der befiehlt die Sorge des Ackerbaues Gott dermaßen nicht, daß er selbst der keines thun wollte, was zum Ackerbau gehört, das Land oder den Acker nicht pflügen oder bauen wollte; sondern wenn er das alles gethan und ausgerichtet hat, alsdann bittet er erst, daß Gott Gedeihen und gut Wetter dazu geben wolle. Wo uns nun Gott erhört und der Bitte gewähret, sollen wir ihm dafür danken; wo aber nicht, sollen wir es auch geduldig leiden und unserm Herrn Gott den schuldigen Gehorsam leisten, auf daß wir lernen, welches da sei der gute und wohlgefällige Wille Gottes in uns, Röm. 12, 2.

25. Darnach hat sich Joseph auch entschuldiget wegen der Sünde, darum er in das Gefängniß geworfen war: er führt an und beruft sich auf sein Gewissen, und die, so ihn fälschlich angegeben und verklagt haben, rechtfertigt er nicht, damit es nicht scheine, als ob er sich selbst schuldig wüßte.

26. Denn wir müssen traun unser gut Gerücht nicht also verachten oder in Wind schlagen, daß wir die Sünde oder Laster, so uns fälschlich und mit Unwahrheit zugemessen und aufgerückt werden, eingestehen und dazu stille schweigen, und die Ankläger und Lästerer mit unserm Stillschweigen also stärken; sondern wir sollen sagen: „Ich habe nicht gesündigt. „Ich habe keinen Teufel“, sagt Christus zu den Juden, Joh. 8, 49.; und abermals Joh. 18, 23.: „Hab ich übel geredet, so beweise es, daß es böse sei.“ Und man hat davon etliche Sprüche Augustins, die gar gewichtig sind, und deren man billig gedenken und darnach man sich richten soll. Als, da er sagt: Wer selbst sein gut Gerücht nicht achtet, ist ein grausamer Mensch.

Mensch. Stern, er sagt: Uns thut nothwendig, daß wir vor Gott ein gut Gewissen haben und auch ein gut Gerücht um des Nächsten willen. Daß man unschuldig ist, daran ist es genug vor Gott, aber nicht bei den Menschen. Denn da ist vonnöthen, daß man auch ein gut Gerücht habe. Wenn ein Bischof oder Prediger fälschlich und mit Unwahrheit beschuldiget würde wegen Unfleißes, Diebstahls, Ehebruchs, oder daß er sollte den gemeinen Schatzkasten bestohlen haben, soll er zu solchem Unrecht mitnichten still schweigen und solches damit bekräftigen. Wo er aber seine Unschuld dargethan, und der Widersacher will damit nicht zufrieden sein, ei, so soll er es im Namen des Herrn dulden und solche unrechtmäßige Beschuldigung oder Lästerung leiden, und soll denjenigen, so die Lästerung wider ihn erregt und aufgebracht, darum nicht verfolgen, oder sich an ihm rächen. Denn es ist an dem genug, daß er bewiesen und dargethan hat, daß ihm unrecht daran geschehen, daß er dermassen geschmäht worden. Wie die Exempel Christi und der Apostel auch lehren.

27. So sagt nun Joseph: Ich habe das mit nichts verschuldet, daß ich an diesen Ort gekommen bin. Denn ich bin aus dem Lande der Hebräer weggeführt, und haben mir diesen Jammer und Unglück die zumwege gebracht, welche mich billig sollten errettet und geschützt haben. Ich habe auch in meines Herrn Hause nicht gesündigt, sondern bin von einer sehr bösen Hure fälschlich und mit Unwahrheit des Ehebruchs beschuldiget worden. Derhalben bitte ich, du wollest mich bei dem Könige vertreten und für mich bitten. Und dies hat er etwas weitschichtig und rhetorischer Weise können ausbreiten und groß machen, auf daß er Gelegenheit hätte, seiner Unschuld und Wohlthat zu gedenken. Aber durch dies Mittel wird er noch nicht befreit werden, darum bleibt er immer länger im Glauben, Hoffnung und Geduld, gibt sich zufrieden, murret nicht, zürnt nicht, flucht auch weder Gott noch Menschen.

28. Denn bei den sehr heiligen Patriarchen ist der rechte, gewisse, beständige Glaube und Erkenntniß der himmlischen Lehre, welche wir von Gottes Gnaden auch haben, gewesen, daß sie haben an Christum geglaubt und durch ihn das ewige Leben zu erlangen gehofft. Der Artikel von der Auferstehung ist in ihre Herzen geprägt gewesen. Denn sie verharren und bleiben beständig nicht allein im Tode, sondern auch über und nach dem Tode. Denn der Tod wird Joseph so oft zugebracht, er wird so oft getödtet, daß man mit Wahrheit sagen mag, daß er auch weit über den Tod beständig geblieben sei. Er wird aber in seinem Herzen also gedacht haben: Warum sollte ich mich selbst martern, oder wider Gott toben und wüthen? Ich weiß, daß der Herr mein Gott ist nicht allein in diesem, sondern auch im zukünftigen Leben.

29. Also hat der Glaube von der Auferstehung der Todten auch im Alten Testament geleuchtet, wiewohl nicht so gar reichlich und offenbar als im Neuen Testament. Christus sagt Matth. 22, 32.: „Gott ist nicht ein Gott der Todten, sondern der Lebendigen.“ Darum haben auch alle Väter, so auf Gott gehofft, ohne allen Zweifel die Auferstehung der Todten geglaubt. Denn das kann klarlich dargethan und bewiesen werden aus der Disputation Matth. 22., da er also argumentirt und schließet: Abraham, Isaak und die andern Patriarchen sind im festen Glauben verharrt und beständig bis in den Tod geblieben: darum haben sie die Auferstehung der Todten auch geglaubt.

30. Denn wo nach diesem Leben kein ander und besser Leben ist, was bedürfen wir denn des Glaubens, oder was bedürfen wir Gottes selbst, oder der Hoffnung auf Gott über oder nach dem Tode? Und wenn wir der Auferstehung nicht warten oder hoffen dürfen, so ist auch kein Glaube und kein Gott. Und wiederum auch also: „Ist ein Glaube und Hoffnung auf Gott, so muß ja auch ein Gott sein, welcher aller deren, so auf ihn vertrauen und hoffen, Gott ist, und deßwillen sie auch die Hoffnung der Unsterblichkeit und des ewigen Lebens behalten. Denn dieweil alle Menschen sterben müssen und die Gottlosen gleichwohl ihr Vertrauen auf Gott setzen, bekennen sie damit von Noth wegen, daß nach diesem Leben ein ander und ewiges Leben sein werde.

31. Ein Epicurischer Mensch bedarf keines Glaubens an Gott; denn derselbe führt gar ein viehisch Leben. Darum denkt er, es sei ihm daran nicht viel gelegen, er sterbe oder lebe. Julius Cäsar hat gesagt, daß der schnelle Tod der allerbeste sei. Denn er ist ein Heide und ungläubiger Mensch gewesen, der kein Gottes Wort gehabt. Darum ist ihm auch kein Glaube, Hoffnung, Liebe oder Geduld vonnöthen gewesen. Wenn das Leben aus ist, so ist mit den Epicureern Alles aus.
\quellseite{340}

32. Diese Dinge sollen wir darum bedenken, daß wir dabei sehen, wie dieser Artikel von der zukünftigen Unsterblichkeit in den Heiligen des Alten Testaments leuchte, welche ohne allen Zweifel gewißlich auch die Auferstehung der Todten geglaubt, und vornehmlich ihr Vertrauen auf Christum gesetzt haben, welcher ein Herr, Stifter und Anfänger der Auferstehung ist.

33. Danach haben wir gesagt, daß das Exempel der Bitte, so Joseph an den Schenken gerichtet, auch wohl zu merken sei, dadurch wir ermahnt werden, daß man der Güter oder Gaben, und allerlei Gelegenheit oder guten Mittel, so vorhanden sind, gebrauchen solle, unsern Leib, dazu auch die Polizei oder weltlich Regiment, die Kirche, und unser Weib und Kinder zu nähren, zu erhalten und zu erretten. Es soll ein Jeder an seinem Orte und in seinem befohlenen Amte thun, was ihm da gebühret, und was er kann und mag. Ist es der Fall, daß es ihm am Gebeinen mangeln wird und daß es nicht gerathen will, soll er zum Leben oder zum Tode bereit sein, zur Traurigkeit oder zur Freude, und daß er Gottes und des Bösen ertragen könne, also, daß Gott der Herr unsere Hoffnung und Vertrauen sei nicht allein in diesem, sondern auch im zukünftigen Leben; wie geschrieben steht 1 Tim. 4, 8.: „Die Gottseligkeit ist zu allen Dingen nütze, und hat die Verheißung dieses und des zukünftigen Lebens“, das ist, man soll hoffen und vertrauen beide für dies und auch für das zukünftige Leben, daß der Herr unser Gott ewiglich sei.

\section[Zweiter Theil]{Zweiter Theil. Von dem Traum des obersten Bäckers und dessen Auslegung; item von Joseph}

\subsection{I}

\begin{quote}
\textbf{V. 16—19. Da der oberste Bäcker sahe, daß die Deutung gut war, sprach er zu Joseph: Mir hat auch geträumet, ich trüge drei weiße Körbe auf meinem Haupt; und im obersten Korbe allerlei gebackene Speise dem Pharaon; und die Vögel aßen aus dem Korbe auf meinem Haupt. antwortete und sprach: Das ist seine Deutung. Drei Körbe sind drei Tage. Und nach dreien Tagen wird dir Pharao dein Haupt erheben, und dich an den Galgen henken, und die Vögel werden dein Fleisch von dir essen.}
\end{quote}

34. Dieser Traum ist dem vorigen gleich, die Deutung aber ist ungleich. Denn der Bäcker denkt, er wolle auch nach drei Tagen erlöset werden, wie dem Andern die Befreiung widerfahren worden. Aber es hört hier auf die Aehnlichkeit, nicht des Glaubens, sondern der Dinge. Darum schließt er nicht recht: Dieweil der oberste Schenke nicht gestraft worden, so werde er gleicher Weise auch ungestraft davon kommen.

35. Das hebräische Wort chori heißt eigentlich weiß und heißt auch löchericht. Darum hat es Hieronymus im Lateinischen nicht recht gegeben: canistrum farinae, Brodtkorb. Jedoch ist ein Unterschied zwischen diesem Worte und dem andern hebräischen Worte laban, welches auch „weiß“ bedeutet. Denn dasselbe Wort heißt, was sehr weiß ist und das herrlich leuchtet; wie Dan. 7, 9. vom Alten geschrieben stehet, deß Kleid schneeweiß gewesen, hellweiß; es heißt nicht dunkel oder rothweiß, sondern das sehr schön weiß ist. Und in der Historie der Esther Cap. 8, 15. steht: „Mardochai aber ging aus von dem Könige in königlichen Kleidern, gelb und weiß“ etc. Derhalben werden allenthalben in der Schrift die Edlen, große Herren, Fürsten, Oberherren und Regenten des Volks chorim genannt, das ist, weiß; wie das im Buch Nehemia zu sehen ist. Denn es ist ein gemeiner Gebrauch gewesen, sonderlich der Völker im Morgenlande, daß sie weiße Kleider getragen, wie auch vom reichen Manne Luc. 16, 19. gesagt wird. Und jetziger Zeit haben die Türken auch weiße Kleider, wie vor Zeiten die Römer den Gebrauch gehabt, daß sie Purpur getragen haben, und werden noch jetzund die Araber Loban genannt, vom Worte laban.

36. Weiter aber wird auch hier canistrum, ein Korb, mit diesem Namen genannt darum, daß er viel Löcher hat, dieweil er von vielen Ruthen dünn geflochten ist wie ein Netz, ein schneeweißer Korb.

37. Also lautet aber nun der Traum: Im obersten Korbe sind allerlei gebackene Speisen gewesen, Kuchen, leichte Brode. Dies ist noch alles glücklich und gut genug, und reimt sich noch wohl mit dem Traume des Schenken. Denn droben, V. 10., werden drei Reben gezeigt und hier drei Körbe: dort aber war der Becher Pharaos in der Hand des Schenken, allhier ist der Korb in der Hand des Bäckers, aber nicht der Meinung, daß Pharao daraus sollte gespeiset werden, sondern die Vögel sollten aus dem Korbe essen. Darum hat es ihm bedeutet, daß er sollte an den Galgen gehenkt werden, daran sein Fleisch von den Vögeln würde zerfleischt werden.

38. Es ist aber droben insgemein von den Träumen und deren Auslegungen gesagt worden, von welcher Materie Lyra unordentlich genug gehandelt. Nun wird aber gefragt: Ob und wie den Träumen zu glauben sei? Mose sagt 5 Mos. 13, 3., man solle den Träumen nicht glauben, und Sirach im 34. Cap. V. 1. ff. hat davon auch sehr fleißig disputirt. Wiederum aber werden im Propheten Joel Cap. 3, V. 1. die Träume gelobt, und wir wissen, daß viel Träume vom Anfang in der Kirche beide den Gottseligen und auch den Gottlosen gezeigt worden sind.

39. Cyprian hat einen Traum gehabt von Christo, wie er gesessen sei, und ein Feind ihm Stricke gelegt und damit ihm und seinen Heiligen nachgestellt habe. Und Augustin erzählt einen Traum, so seine Mutter Monica gehabt, welcher, da sie gemeint, ein Jüngling erschienen ist mit fröhlichem Angesicht und in einem glänzenden Kleide, der gesagt hat, warum sie weine. Da hat sie geantwortet, wie sie das Verderben oder Verdammniß ihres Sohnes beweine wegen der Secte der Manichäer, welcher Augustin zugethan gewesen. Da hat sie der Jüngling heißen zufrieden sein und zu ihr gesagt: Weine nicht; denn wo du bist, da ist er auch. Diese Stimme hat Monica angenommen als eine himmlische Antwort, und ist fröhlich zu ihrem Sohne gegangen, daß sie ihm den Traum erzählete, darin ihr die Hoffnung gezeigt, daß er würde zur wahren Gottseligkeit bekehrt werden. Augustin aber hat anfänglich desselben Traums gespottet, und ihn also höhnisch verachtet und gesagt: Du hast es nicht recht verstanden. Denn vielleicht hat der Jüngling also gesagt: Wo er ist, da wirst du auch sein. Aber Monica ist noch fest auf ihrer Meinung bestanden. Und Augustin sagt, er sei durch dasselbe Wort seiner Mutter sehr heftig bewegt worden. Denn es ist ein Traum gewesen, welcher durch GOttes Wirkung eingegeben, geglaubt und ausgelegt worden ist. Darum ist es endlich auch also ergangen, wie es pflegt mit göttlichen Träumen zu gehen.

40. Die andern Träume aber, so von GOtt nicht herkommen, sind ungewiss und falsch. Derhalben gehört ein sonderlicher Verstand dazu, die Träume zu unterscheiden.

41. Ich habe aber nun oftmals gesagt, daß ich vom Anfang meiner Sache allzeit den Herrn gebeten habe, daß er mir weder Träume, noch Gesichte oder Engel sende. Denn es haben mich viel Rottengeister angegriffen, deren Einer der Träume, der Andere der Gesichte und aber ein Anderer sonderlicher Offenbarungen sich gerühmt, damit sie sich unterstanden mich zu lehren. Aber ich habe ihnen geantwortet, daß ich solcher Offenbarung nicht begehre, und wo mir selbst auch etliche vorkommen würden, wollte ich ihnen doch nicht glauben. Und habe mit ernstem Gebet darum gebeten, daß mir GOtt rechten gewissen Verstand der heiligen Schrift geben wolle. Denn wo ich das Wort habe, so weiß ich, daß ich auf rechtem Wege gehe, und daß ich nicht leichtlich kann betrogen werden oder in Irrthum fallen. Und ich will lieber Dawids Verstand haben, denn sonst prophetische Gesichte, welche, meines Erachtens, David selbst auch nicht sehr begehrt hat. Siehe aber, wie so einen gewissen Verstand der Schrift er gehabt habe.

42. Die Schwärmer haben versucht, mich mit ihren Träumen, der Eine hieher, der Andere dorthin zu treiben und zu bewegen; und wo ich einem jeglichen hätte gehorchen wollen, so hätte ich wahrlich dreisig oder vierzigmal meine Sache ändern müssen. Da ich sie aber alle verworfen, da haben sie über mich geschrieen, ich wäre halsstarrig und hätte meinen eigenen Kopf, und ließen mich also zufrieden. Darum gebe ich nicht viel auf Gesichte und Träume; und wiewohl sie das Ansehen haben, als sollten sie etwas bedeuten, verachte ich sie doch, und lasse mir genügen am gewissen Verstande und Gläubwürdigkeit der heiligen Schrift. Wenn ich nun das Wort habe, so bin ich gewiß, daß GOtt und Engel bei mir sind, und nicht stichbarlich, daß sie aber doch gleichwohl ihren Glanz geben und mich auf dem Wege der Wahrheit leiten. Dies ist
\quellseite{341}
eigentlich meine Meinung, welche ich nicht zu ändern weiß.

43. Es sind aber doch nicht allerlei Träume gänzlich zu verachten, wiewohl ich dazu keine Lust habe, sondern man muß auf die rechten Wahrzeichen der rechten wahren Träume merken und Acht geben. Denn auf die Träume, so von Gott eingegeben werden, folgt bald hernach, daß sie dermaßen eingeprägt werden, daß der, dem geträumt hat, solchen Traum mit nichten vergessen kann, oder wo ihm der Traum ohngefähr entfallen oder aus dem Sinn gekommen, kommt er ihm doch bald wiederum vor; wie das in der Historie des Königs Nebucadnezar Dan. 2, 1. ff. zu sehen ist. Denn der \unsicher{Same}, von Gott auserworfen, bringt von Noth wegen seine Frucht. Also ist es in diesem gegenwärtigen Exempel auch, daß den zwei Amtmännern und hernach dem Könige auch selbst ihre Träume so fest eingeprägt worden, daß sie darüber sehr bewegt und bestürzt worden sind, und können sich nicht zufrieden geben, bis sie die Auslegung der Träume hören.

44. Darum lasset uns aus diesem Capitel lernen, was Joseph gesagt: „Die Auslegung gehöret Gott zu“, der gibt Träume, und gibt auch zugleich die Auslegung mit, welche sonst niemand anders erfinden könnte. Es hatte Pharao viel weise und gelehrte Leute in seinem Königreiche, als, Priester und andere Regenten über den Gottesdienst der Egypter, wie das hernach wird zu sehen sein, unter welchen gewesen sind Sternkundige und die des Himmelslaufs Erfahrung gehabt, welche Künste fast ausschließlich bei den Chaldäern sich fanden. Aber dieselben Auslegungen sind sehr ungewiß und zweifelhaft, wiewohl sie manchmal übereinstimmen mit dem Ausgange, daß es also erfolgt, wie die Träume gelautet haben, wie man im Sprüchwort sagt: Saepe jaciendo jacitur Venus: Wo man oftmals wirft, muß man ja endlich die Venus treffen.

45. Julius Cäsar hatte einen Traum gehabt, daß er seine Mutter geschwächt hätte. Denselben Traum hat ihm der Satan gezeigt und eingegeben, welcher unseres Herrn Gottes Affe ist. Derhalben haben die Ausleger diesen Traum gezogen auf den Bürgerkrieg, damit er das Vaterland betrübt und geplagt hat. Item, es hat ihm geträumt, daß er flöge. Und werden sonst im Sueton viele andere Träume von Nero und Claudius erzählt. Dies sind nicht prophetische Auslegungen, sondern gelten eben so viel, als Wahrsagung oder die Weissagung, davon Cicero etliche Bücher geschrieben, als, Tusculanae quaestiones und De divinatione.

46. Der Teufel ist ein gewaltiger und wunderlicher Geist, welcher beide die, so da machen, und auch die so schlafen, betrügen kann. Und geschieht zwar bisweilen, daß seine Träume und Auslegungen ihren Ausgang gewinnen, bisweilen aber auch nicht. Denn sie haben ihren Ursprung aus solchen Ursachen, so ihm vor Augen liegen. Er sieht ganze zwanzig oder dreißig Jahre lang an den Rath und Anschläge an den Fürstenhöfen, darnach wie die Werkzeuge an Leib und Seele geschickt sind; desgleichen sieht er, wie die Fürsten auferzogen und gelehrt werden; sieht, was ihre Art und Sitten sind; daraus kann er viel schließen. Wie Herzog Georg zuvor gesehen hat, daß seine zwei Söhne sterben würden noch bei seinem Leben.

47. Darum kann der Satan leichtlich Träume eingeben und darnach eine Auslegung daraus schließen, wie er sonst die Phantasie erregt, und den Leuten mancherlei Bilder, Raserei und Unzucht einbläst. Darnach hat er seine eigenen Pfaffen und Träumer, wie deren viele gewesen sind schier bei allen Völkern und auch im Volke Israel, welche die Propheten heftig schelten und dasselbe gottlose Wesen des Volks strafen. Dieselben Pfaffen aber und Träumer fehlen auch oftmals in der Auslegung der Träume. Und wo Satan diese Träume der zwei Amtmänner gehört hätte, würde er die Ungleichheit darin nicht so leichtlich haben verstehen können; er hätte denn zuvor gemußt, was Pharao und auch die Amtmänner im Sinne gehabt. Es verhindert es auch oft Gott, daß solche Träume nach ihrer Auslegung nicht geraten, und den Ausgang nicht gewinnen, so zuvor verkündigt worden, nämlich, wenn die Ursachen geändert werden, daraus sie der Satan geschlossen hat. Darum soll man solchen Träumen oder ihrer Auslegung keinen Glauben geben.

48. Der erste und vornehmste Grad aber der Träume ist eine göttliche Art der Weissagung, davon Joel Cap. 3, V. 1. sagt: „Eure ältesten sollen Träume haben“, nämlich, göttliche Träume, welches gewisse Einbildungen sind, von Gott eingegeben; und ist bei denselben auch zugleich gesetzt nicht allein die Auslegung, sondern auch der Glaube derer, welchen sie gezeigt worden, die ihnen geglaubt, und endlich auch der rechte Effect, das ist, die Erfüllung und Ausführung. Der andere Grad der Träume ist teuflisch. Der dritte ist natürlich, davon die Aerzte disputiren und daraus sie abnehmen und schließen, wie es um das Geblüt oder um den ganzen Leib des Menschen beschaffen sei. Darum wird da nicht gefragt, was solche Träume bedeuten nach dem äußerlichen Leben, oder nach der Schrift. Und von diesen zwei Graden der Träume soll man sich nach der Regel halten, so im Sirach steht, und die Satz auch gegeben hat, da er sagt: Somnia ne cures etc.: Du sollst der Träume nicht achten.

\subsection{II}

49. Wir wollen nun wiederum auf Joseph kommen. Bis daher haben wir gehört, wie Gott seine Gnade auf ihn gewendet, und daß er ihm auch Gnade gegeben vor dem Obersten über die Metzger und über das Gefängniß, desgleichen vor den zwei Gefangenen, welche an ihm gesehen, daß er sonderlich Glück und seinen Verstand gehabt in all seinem Thun. Darnach hat er auch innerlich Trost gehabt, dadurch er gefühlt und empfunden, daß er mit dieser großen Gabe der Weissagung durch den Heiligen Geist begabt gewesen.

50. Die Weisheit aber ist mit ihm in die Grube gefahren. Darum hat er gute Hoffnung, und verspricht sich die Befreiung wegen dieser herrlichen Gabe, und wegen der Gnade und Barmherzigkeit, damit, wie er gewußt, ihn der Herr angesehen hatte. Er glaubt auch menschlichen Gedanken nach, daß er an diesem Schenken auch werde Hülfe und Trost haben. Welches zwar eine menschliche Zuversicht ist, die doch mit nichten zu verachten oder zu strafen ist. Denn er hat also geschlossen: Gott hat mich mit prophetischem Geiste erleuchtet, darum wird er mich nicht nicht verlassen, und läßt mir über das diese Gelegenheit sich bieten, daß ich dem obersten Schenken seinen Traum deuten soll, welcher mich bei Pharao wird losbitten können. Solche Rathschläge und Gedanken sind zu loben. Denn wir sollen die Mittel, so vorhanden sind, nicht verachten, sondern sie auch möglich und natürlich sind, auf daß wir Gott nicht versuchen. Ist es aber der Fall, daß wir durch solche Gelegenheit oder Hoffnung betrogen werden, sollen wir damit auch zufrieden sein und darüber nicht ungeduldig werden.

51. Nun ist aber dies Vertrauen, so Joseph gehabt, nicht so gar menschlich. Denn er verlässt sich auf die Gabe, so ihm von Gott verliehen ist, da er denkt: Siehe, Gott hat mich mit Gnaden angesehen und bestellt mir nun diesen obersten Schenken zum Helfer. Darum will ich es versuchen, ob ich bei dem Könige etwas erlangen möchte; und wo mir das geraten wird, wohl und gut, wo aber nicht, so geschehe des Herrn Wille.

52. Und dies Exempel wird uns zum Troste vorgestellt, auf daß wir wissen, wir mögen ohne Sünde Hülfe suchen, und der Mittel gebrauchen, so uns von Gott gegeben oder vorgestellt werden, es sei gleich natürlich, uns zu nähren, zu erretten, unsern Leib zu erhalten, oder aber sonst bürgerlicher und häuslicher Weise, Frieden zu erhalten, und daß wir beide für uns und auch ingemein an unserer Habe und Gute unbeschädigt bleiben mögen. Und sollen die Worte der gottlosen Leute, als, der Türken und Anderer mehr, nicht achten, die da sagen, daß diejenigen Dinge, so einmal von Gott bestimmt und beschlossen sind, nicht können gehindert werden. Es ist, sagen sie, ein Fatum, eine Schickung von Gott, darum wird dies oder jenes nothwendig geschehen müssen. Soll es geschehen, so geht es. Dir aber ist solches nicht geoffenbart. Denn sonst würde der Glaube, Hoffnung, Geduld und die Liebe keine Statt haben, sondern würde daraus in den Herzen der Menschen folgen entweder Verachtung Gottes, oder aber, daß er bei uns gar würde verhaßt werden, daß wir nicht glauben, oder keine Hülfe von ihm erwarten würden. Derhalben ordnet und schickt Gott die zufälligen Dinge, so geschehen sollen, auf daß damit Glaube und Hoffnung ihre Uebung haben mögen.

53. Und wir haben gesagt, daß Joseph unnütziger Weise von den Juden gestraft werde, als hätte er sein Vertrauen auf einen Menschen gesetzt. Denn wo man etwas von einem Fürsten bitten soll, muß man ja an seiner Güte und Gnade nicht verzagen, sonst würde ich nichts dürfen bitten; sondern ich soll getrost sein und mich zu ihm versehen, er werde meine Bitte gnädig erhören. Werde ich nun damit betrogen, ist es mir ohne Sünde. Wo ich erlange, was ich bitte, ei! so danke ich dafür meinem lieben gnädigen Gott.

\quellseite{342}

\begin{quote}
\textbf{V. 20—23. Und es geschah des dritten Tages, da beging Pharao seinen Jahrtag; und er machte eine Mahlzeit allen seinen Knechten, und erhub das Haupt des obersten Schenken, und das Haupt des obersten Bäckers unter seinen Knechten; und setzte den obersten Schenken wieder zu seinem Schenken, daß er den Becher reichte in Pharao Hand, aber den obersten Bäcker ließ er henken, wie ihnen Joseph gedeutet hatte. Aber der oberste Schenke gedachte nicht an Joseph, sondern vergaß sein.}
\end{quote}

54. Hier wird nun abermals der fromme gottselige Joseph angefochten; denn er merkt, daß seine Hoffnung und Zuversicht, so er gehabt, ihm gefehet, nämlich, daß die Zeit und Stunde seiner Befreiung jetzt vorhanden wäre. Er hatte den Schenken gebeten, daß er seiner gedächte wollte, wo es ihm würde wohl gehen. Aber was geschieht? Josephs wird gar vergessen, und Gott läßt zu, daß er über solcher Hoffnung und Trost abermals zwei ganze Jahre angefochten und betrübt wird. Lieber Herr Gott, will es denn nicht schier werden? es währet auch gar zu lange.

55. Man pflegt gemeiniglich zu sagen: Patientia saepius laesa fit furor, das ist: Wo man es Einem zu viel macht, kann er in die Länge nicht geduldig sein, sondern aus Geduld wird Ungeduld. Welches denn von menschlicher oder heidnischer Geduld recht gesagt wird. Denn dieselbige wird verwandelt in Zorn und Ungeduld, wo man allzu viel Unrecht leiden soll. Es soll aber nicht sein; sondern die christliche Geduld soll immer verharren und beständig bleiben, und ist derselbe Spruch des Hieronymus heidnisch und nicht christlich.

56. Uns wird allhier auch eine Lehre vorgehalten von der Undankbarkeit der Welt, was sie pflege für einen Lohn für die allerhöchsten Wohlthaten zu geben. Und ist wahrlich ein erschrecklich Gemälde oder Bild der höchsten Undankbarkeit der Welt. Es hatte Joseph seiner Frau, Potiphars Weibe, gedienet, nicht allein nach ihrem Gefallen, sondern so, daß sie von Lust und Liebe unsinnig wurde. Was gibt sie ihm nun wieder dafür? Den Tod und allerlei Unglück. Wer will denn der Welt dienen? Wer will nun des Regiments und der Kirche warten? Wer will Kinder und Schüler lehren und auferziehen, wenn er keinen andern Lohn soll davon bringen, denn die allerhöchste Undankbarkeit für die höchsten Wohlthaten? Das wird freilich die Welt nicht thun; denn die will traun ihre Wohlthaten nicht verlieren, sondern will dafür Danksagung, Ehre und Lob haben. Und wo ihr solches nicht widerfährt, hört sie bald auf, Andern Gutes zu thun. Darum auch Bernhardus sagt: Undankbarkeit ist ein solch Laster, welches den Brunnen der Güte gar austrocknet; welches zu verstehen ist von menschlicher und nicht von göttlicher Güte, auch nicht von der Güte der Kinder Gottes, welche nicht auszuhoffen ist. Denn sie hören nicht auf, Gutes zu thun. Aus diesem Weinstock wird keine Distel oder Dorn, sondern ein guter Baum bringt stets gute Frucht.

57. So hat nun Joseph diese Bosheit und Untreue auch dulden müssen, daß der Schenke sein vergessen, und er also immer länger in der allerschwersten Dienstbarkeit, und höchster Noth und Armuth gewesen würde. Denn er hat im Gefängniß ja nichts gehabt, denn nur kümmerlich das Brod und ein schlecht grob Kleid, und will der Jammer mit ihm nicht aufhören.

58. Dagegen aber ist die Welt also geartet, daß sie verschlingt und an sich reißt allerlei Wohlthaten, welche sie doch nimmer erkennt. Und gegen Gott hält sie sich auch also, der doch seine Sonne aufgehen läßt über die Guten und über die Bösen, ja, er schüttet den ganzen Erdkreis gar aus, und theilt ihn unter die hoffärtigsten und undankbarsten Menschen, welche Verächter und Lästerer Gottes sind. Aber da ist niemand, der da singe, was im 34. Psalm V. 2. steht: „Ich will den Herrn loben allezeit“, und im 148. Psalm V. 3.: „Lobet ihn, Sonne und Mond“ etc. Ja, Fluchen, Schelten, Lästern und daß man seine Heiligen tödtet, das wird ihm gegeben für seine unzähligen leiblichen und geistlichen Wohlthaten. Ja, endlich ist der Sohn Gottes selbst, da er in die Welt gesandt, getödtet worden, und er ist doch von unendlicher und unermesslicher Güte. Darum ermahnt uns Christus Matth. 5, 48., item R. 44.: „Ihr sollt vollkommen sein, gleichwie euer Vater im Himmel vollkommen ist.“ Thut Gutes auch den Bösen, bittet für sie.

59. Diese Kunst hat Joseph sehr wohl gekonnt, und hat das Evangelium recht studirt und gelernt von den allerbesten Lehrmeistern. Denn er ist verkauft worden im siebzehnten Jahre seines Alters, nachdem er zehn Jahre lang seinen Vater hat predigen hören, dazu auch die rechte Lehre von Rahel, seiner Mutter, und Debora, der Säugamme, gehört, welche alle ihn mit sonderlichem Fleiß unterwiesen und sehr lieb gehabt haben. Denn sie haben gesehen, daß er ein feiner stiller Jüngling gewesen ist und gar von guter Art. Die Säugamme Debora ist an der Großmutter Statt gewesen. Und ist dazumal das Gedächtniß der Patriarchen Noah, Sem und Eber noch frisch und unverloschen gewesen. Und Isaak war auch noch am Leben, derselbe wird ihm die Exempel der Trübsale Jacobs fleißig vorgehalten haben, samt den Geschichten, Worten, Exempeln und Trübsalen der andern Väter. Und Joseph selbst hat auch zum großen Theil das Leben seines Vaters in Syrien gesehen, und darnach die Trübsale, so Jacob, sein Vater, über Dina gehabt, da ihm dieselbe geschwächt worden, auch über den Todtschlag der Sichemiter und der Blutschande Rubens. Da hat er gesehen die Thränen und Seufzer seines Vaters, hat sein und des ganzen Hauses Klagen gehört und bei ihm allein sind die Predigten des Vaters kräftig gewesen. Denn er ist der beste Theil gewesen des guten Ackers; die Andern haben beide den Bruder und Vater getödtet. Allein Joseph ist der köstlichste Edelstein gewesen, der die Worte und Exempel Isaaks, Jacobs und der andern Väter behalten und bewahret hat. Er ist gewesen wie ein Diamant im Glauben zu Gott, in der Hoffnung und Geduld, und sein Herz hat mit keinerlei Ungestüm vieler schweren Anfechtungen, so er erlitten, überwunden werden mögen.

60. Darum wird im Buch der Weisheit Cap. 10, 13. 14. recht gesagt, daß die Weisheit, das ist, Gottes Sohn, den verkauften Joseph nicht verlassen habe, sondern sei mit ihm hinab in die Grube und Kerker gefahren. Ja, denselben hat er zum Lehrmeister gehabt, der hat das Gedeihen gegeben, daß er das Wort, so er gehört, fest in seinem Herzen behalten hat. Er hat auch noch vor Augen gehabt die Strafe durch die Sintfluth, das Feuer, so die Sodomiter verbrannt, und andere Trübsale mehr; darnach auch mancherlei Errettungen der Gottseligen, wie dieselben sein Vater erzählt. Daher in ihm die Furcht Gottes, Glaube, Hoffnung und andere Tugenden immer mehr zunahmen und gestärkt wurden, daß sein Herz durch keinerlei Ungestüm hat können geschwächt werden.

61. Also wird er nun im Kerker und gleichsam in der Hölle gehalten, und wiewohl er die Hoffnung gefaßt, daß er nun bald würde befreit werden, sieht er doch, daß solche Hoffnung gar vergeblich ist. Er hat aber einen prophetischen Geist, der ihn tröstet, daß er sich immer mit der Lehre der Väter aufrecht erhalten, und gedacht: Jetzt wird Gott helfen, er verzieht aber doch ja zu lange. Nun ist es aber gleichwohl sehr verdrießlich, hoffen und immer also darüber betrogen werden, und daß die Hoffnung vergeblich und umsonst ist; das macht den betrübten elenden Leuten ihren Schmerz immer größer. Darum hat ihn die Vermahnung, im 27. Psalm V. 14. steht, noch gethan: „Harre des Herrn, sei getrost und unverzagt, und harre des Herrn“ etc. Inzwischen aber, dieweil er verzieht, muß man immer um Errettung beten und seufzen. Darnach soll man auch Mittel suchen, sollt dich befleißigen und bemühen, daß du aus gegenwärtiger Noth und Trübsal mögest errettet werden. Denn wo es ohne die Hoffnung und den Fleiß wäre, damit ich begehre erhört und erlöst zu werden, bedürfte ich des Anrufens oder Betens gar nicht. Das Gebet aber soll emsig und unablässig sein, soll nimmer aufhören, nicht müde werden, und soll dem Mißtrauen, Unglauben und Verzweiflung immer widerstehen.

62. Denn auch Joseph, ob er wohl im Geist geduldig war, ist er doch ein Mensch gewesen, wie die andern Heiligen alle, so noch im Fleisch bleiben, ob sie gleich nicht nach dem Fleisch leben, wie Paulus Gal. 5, 17. sagt: „Das Fleisch gelüstet wider den Geist.“ Und währet derselbe Streit in allen Heiligen oder Gläubigen durch ihr ganzes Leben, in etlichen mehr, in etlichen weniger. Wiewohl nun Joseph in seiner Trübsal geglaubt, gehofft und Gott gedankt hat, ist doch daran kein Zweifel, er habe auch noch in seinen Gliedern ein Gesetz gefühlet, daß dem Gesetz in seinem Gemüth widerstritten, Röm. 7, 23. Solch Widerstreiten nach dem Fleisch und etliche Regungen der Ungeduld hat er gehabt; aber er hat es doch überwunden, nicht daß er solches auch stets gefühlt hätte, sondern etwa einen Tag oder halben Tag, oftmals aber auch wohl zwei oder drei Tage, nach Gottes Wohlgefallen, der die Seinen läßt angefochten werden, macht aber, daß die Versuchung so ein Ende gewinnen muß, daß sie es können ertragen, 1. Cor. 10, 13. Es muß ja Einer einen Tag oder zween betrübt sein. Das Fleisch muß sein Murmeln haben. Es soll aber doch der Geist über das Fleisch herrschen, wie 1. Mos. 4, 7. gesagt wird: „Laß du der Sünde nicht ihren Willen“, daß du endlich schließest, als ob du Alles müßtest verloren geben und gar verzweifeln.

\quellseite{343}

63. Das Fleisch, so wir noch umher tragen und damit wir uns noch schleppen, macht, daß die Freude des Geistes nicht kann vollkommen sein, kann auch das nicht ausrichten, was sie sollte. Denn es geht in diesem Leben ohne Murren nicht ab. Und doch muß dasselbe Fleisch im Zaum gehalten und gezwungen werden, damit das Murren und die Ungeduld nicht herrsche, wiewohl es den Menschen immer treibt, reizt und zum Zorn zwingt.

64. Nun sind aber auch geistliche Anfechtungen zu dem Unglück gekommen, so äußerlich das Fleisch müde gemacht und sehr hart betroffen. Denn es ist ein armselig Ding, im Kerker also gedemüthigt werden, Elend und Verachtung leiden, versäumt und allein ohne alle Hülfe und Trost gelassen werden; dazu für die höchsten Wohlthaten nichts denn die höchste Undankbarkeit erfahren müssen. Darum ist dies die größte Geduld gewesen, so kein Mensch von selbst haben kann. Nachdem aber der Teufel gesehen, daß Joseph mit der Anfechtung des Murrens und der Ungeduld gestritten, da hat er erst seine feurigen Pfeile auf ihn gerichtet und ihm damit zugesetzt; wie das auch derselbigen Pfeile einer gewesen, da er Matth. 4, 3. zu Christo sagt: „Bist du Gottes Sohn, so sprich, daß diese Steine Brod werden.“ Also wird er auch hier zu Joseph gesagt haben: „Du stehst diese erbärmliche und langwierige Noth; siehst, daß dein Gebet vergeblich und umsonst ist; du schreiest, rufst Gott an und glaubst an ihn: aber es ist alles vergeblich. Jetzt hattest du die Hoffnung, du würdest durch die Deutung des Traums, so der Schenke gehabt, befreiet und erhöht werden können. Nun wohlan, laß dich Gott erretten, hat er Lust zu dir. Diese feurigen Pfeile des Satans haben das Herz Josephs, so allbereit schwach und verwundet war, noch viel härter geplagt.

65. Wenn sich die Gottseeligen im Geist freuen, und singen: „Der Herr ist meine Macht und mein Psalm“ etc., Ps. 118, 14., so ist der Teufel weit von ihnen, und hört bei ihnen das Murren und Ungeduld auf. Aber wenn der Baum niedergestrecket ist, so kommt der Schalk. So lange die Worte des Lobes und der Danksagung klingen und gehört werden, daß sie sagen, wie im 34. Psalm V. 2. steht: „Ich will den Herrn loben allezeit; sein Lob soll immerdar in meinem Munde sein.“ Ich will mich freuen des Herrn, meines Heilandes etc.: so lange verschwinden alle Anfechtungen der Traurigkeit und Mißglaubens, und thut sich der Himmel und Paradies weit auf, die Hölle ist erloschen. Wenn sie aber wiederum verlassen werden, alsdann kommt der Gedanke von der Trübsal, Verlassung und Elend wieder. Wie Joseph bei sich selbst gedacht, daß er von seinen Brüdern verkauft worden, die jetzt daheim gute ruhige Tage hätten und in Wollust lebten, da er dieweil jämmerlich gebunden und im Kerker verschlossen liegen mußte, so er doch solches nie verdient und nichts Böses gethan hätte. Das ist der Teufel; wenn das Wasser trübe wird, so fischt er gern, und kommt mit den Mordstößen, nämlich, mit der Anfechtung des Unglaubens und Verzweiflung.

66. Nun ist aber Joseph in diesem Jammer gar allein gelegen, hat niemand gehabt, mit dem er sich hätte unterreden können; da ist kein Beichtvater, da war niemand, der ihn getröstet und aufgerichtet hätte, da er schier von so vielem schweren Unglück gar unterdrückt gewesen. Wohlan, da liegt das liebe Kind noch zwei Jahre: so lange wird er noch im Kerker gehalten. Wo es nur zwei Monate oder zwei Wochen gewährt hätte, so hätte er es etwas leichter, ja, wohl mit gutem Willen alles ertragen können. Aber der oberste Schenke vergißt sein ganz und gar, und würde sein ganzes Lebenlang an Josephs Bitte nimmer gedacht haben, wo nicht dem Pharao von GOtt der Traum wäre gegeben worden, welches ihm eine Ursache und Anfang seines Heils und Wohlfahrt sein wird.

67. Inzwischen aber, da der Schenke befreit worden, hat Joseph mit großen Sorgen alle Stunden, Tage und Monate an den Fingern abgezählt, und mit herzlichem Seufzen gesagt: Ach! wann will es werden, jetzt ist schon ein Jahr verlaufen und geht nun das andere auch daher. Endlich, da die zwei Jahre nun auch vergangen und doch keine Befreiung erfolgt ist, hat er gedacht, es wäre nun gar mit ihm verloren, ihm würde nun nicht können geholfen werden. Nun wird nichts daraus, hat er gedacht. Das sind alles solche Argumente, so den Menschen zur Verzweiflung führen. Und ist ohne Zweifel die höchste und schwerste Anfechtung gewesen und ist auch die letzte.

68. So sitzt nun dies schöne Licht in Finsterniß, und lernt dies Wort verstehen, 5 Mos. 8, 3.: „Der Herr demüthigte dich, und ließ dich hungern“ etc., „auf daß er dir kund thäte, daß der Mensch nicht lebt vom Brod allein, sondern von allem, das aus dem Munde des Herrn gehet.“ Wo Joseph diese Speise nicht hätte mit sich gehabt, und wo der Heilige Geist und der Sohn GOttes nicht wäre mit ihm in die Grube gefahren, und ihm mit dem Worte, so er von seinem Vater gelernt, erhalten hätte: so würde er die harten Stöße des Satans nicht haben ertragen oder überwinden können. Denn die geistlichen Anfechtungen übertreffen weit alle menschliche Stärke und Kraft; wie denn Viele von solchen harten Stößen und Donnerschlägen des Teufels getödtet werden; wenn er den tödtlichen Spieß und feurige Pfeile in des Menschen Herz schießt, und der Mensch kein GOttes Wort hat, wird er in einem Augenblicke verlöschen. Und dieselben tödtet nicht etwa das Schwert oder ein Mensch, sondern der Teufel durchsticht unmittelbar das Herz und Seele mit seinen feurigen Pfeilen.

69. Darum lernt nun Joseph vom Worte GOttes leben, und nicht allein vom Brode, damit er in solchen großen Trübsalen sein Leben nicht hätte erhalten können. Denn wo er nicht Trost vom Heiligen Geiste gehabt, so wäre er vor Traurigkeit gestorben, sonderlich in der Einsamkeit, da er aller Menschen Gemeinschaft und Hülfe beraubt gewesen, und gar keinen Trost gehabt, denn allein von Tieren, Mördern und andern Uebelthätern mehr, welchen er im Gefängniß gedienet. Deshalb ist es nicht an menschlicher Hülfe und Trost gelegen gewesen, daß er überwunden und besteht wider den Satan, die Welt und sein eigen Fleisch, und lebt auch in seinem andern Leben, denn daß in GOttes Wort ist. Denn er bleibt einfach an der Verheißung hangen, ist bereit, auch noch etwas Größeres zu leiden; wiewohl er nach dem Fleische begehrt erledigt zu werden und auch nach dem Geiste. Jedoch dieweil er sieht, daß es GOtt also haben will, ist er mit seinem Willen zufrieden, und tödtet das Fleisch, so da schreit und murrt, und sagt: Der Wille des Herrn geschehe. Aber es ist ein trefflich schwer fiat (geschehe). „Meinem Vater ist von GOtt verheißen worden, wird er gesagt haben, daß er solle Erben und Nachkommen haben und gesegnet werden; und mein Vater hat mich gelehrt, auf solche Verheißung warten, bis sie erfüllt werde. Nun werde ich nach dem Fleische hingehalten und geschwächt; aber im Worte werde ich wiederum gestärket und muthig gemacht.“

70. Und also ist der sehr heilige und fromme Joseph in diesen zwei Jahren gekreuziget, gestorben und begraben, und niedergefahren zur Hölle. Nun wird der HErr kommen und ihn erretten, zu Ehren setzen und groß machen; wie er ihn berufen, gerecht gemacht und den Heiligen Geist und seinen Sohn gegeben hat, welcher mit ihm in das Gefängniß gestiegen. Nun ist die Marterwoche aus! Denn er wird nun bald lebendig gemacht werden und auferstehen. Und ist wahrlich ein sehr schön Exempel, welches der christlichen Kirche und allen betrübten und angefochtenen Menschen, welcherlei Weise sie auch beschwert sein mögen, vorgehalten werden soll. Wir würden schwerlich so viel leiden. Und auf solche Weise soll man GOtt erkennen lernen, daß er wunderbarlich sei in seinen Heiligen, wie die Worte im 4. Psalm V. 4. lauten: „Erkennet doch, daß der Herr seine Heiligen wunderlich führet; der Herr höret, wenn ich ihn anrufe.“ Das exaudiet (Hören) folgt wohl, aber das clamavero (Anrufen) währet zumal lange. Joseph hat nun dreyzehn Jahre im Elende gelebt, und zwei Jahre im Gefängniß wegen seines treuen und fleißigen Dienstes und wegen seiner Keuschheit.

\quellseite{344}

Aber Gott erhört die Seinen also, daß er mit ihnen wunderbarlich umgeht, wie das alhier an Joseph zu sehen ist. Er selbst weiß zwar die Ursache nicht, warum der Herr die Errettung so lange verzieht, sondern hat dieselbige Gott befohlen und ruft ihn dieweil an.

71. Dies Exempel soll uns nun auch stärken in unsern Anfechtungen, die wir noch mit so großer Trübsal und Beschwerung nicht gekämpft haben. Und dieser Patriarchen Leiden übertreffen auch weit nicht allein der Mönche Leiden, sondern auch die Marter der Heiligen im Neuen Testamente, die Apostel ausgenommen, als, Paulus und andere, welche auch solche schwere Anfechtung und große Ehre gehabt, wie die Patriarchen. Denn wir sehen, daß auch des einigen Joseph Leben voller Marter und Leiden ist, damit er leiblich und geistlich auf das allerschwerste ganze dreizehn Jahre umgetrieben und geplagt worden ist: und hat nichts desto weniger solche Wunderwerke gethan, welche auch den Werken der Heiligen in der Kirche, so aus den Heiden versammelt worden, weit mögen vorgezogen werden.

\chapter{Das einundvierzigste Capitel}

\section[Erster Theil]{Erster Theil. Von den Träumen Pharaos}

\begin{quote}
\textbf{V. 1—7. Und nach zweien Jahren hatte Pharao einen Traum, wie er stünde am Wasser, und sähe aus dem Wasser steigen sieben schöne fette Kühe, und gingen an der Weide im Grase. Nach diesen sahe er andere sieben Kühe aus dem Wasser aufsteigen; die waren häßlich und mager, und traten neben die Kühe an das Ufer am Wasser. Und die häßlichen und magern fraßen die sieben schönen fetten Kühe. Da erwachte Pharao. Und er schlief wieder ein, und ihm träumete abermal, und sahe, daß sieben Aehren wüchsen aus Einem Halm, voll und dicke. Darnach sahe er sieben dünne und versengte Aehren aufgehen. Und die sieben magern Aehren verschlangen die sieben dicken und vollen Aehren.}
\end{quote}

1. Wir haben nun gehört das treffliche, unglaubliche Exempel des Glaubens, der Hoffnung, Liebe, Geduld und aller Tugenden an dem Jünglinge Joseph. Ich glaube aber, daß Isaak, sein Großvater, welchen wir droben, Cap. 35, 29., begraben, wiewohl er, nachdem Joseph verkauft worden, noch am Leben gewesen, um diese Zeit gestorben sei, da er seines Sohnes Sohn, Joseph, noch nicht wieder bekommen hat. Und er wird ohne Zweifel von Bekümmerniß und Traurigkeit des Herzens verzehret worden sein, und ist endlich also im Glauben und Hoffnung auf göttliche Verheißung entschlafen.

2. Dies sind heilige und geistliche Historien, und wenn dieselben fleischliche Menschen, als, die Papisten, lesen, schauen sie nicht in das Herz hinein, und sehen darin nicht die großen Wunderwerke des Glaubens, der Hoffnung und Geduld, daß Joseph solch groß Unglück und so lange Zeit hat mit unverzagtem Herzen, festen beständigen Glauben und Hoffnung leiden und ertragen können.

3. Es ist aber ein sehr groß Exempel, mit welchem andere nicht können verglichen werden um der vielen unzähligen Marter willen, so alle an diesem einigen Menschen zusammen gekommen, die so lange Zeit gewähret und dabei man so gar keine Hülfe oder Trost gespüret. Denn sonst pflegen die Märtyrer, wie die Kindlein, so von Herodes getödtet worden, in einem oder zwei Tagen von ihrer Marter und Pein loszukommen; hier währets ja zu lange. Und du solltest Wenige finden, die so gar armselig und verlassen wären, die sich zum wenigsten des Wurmens, Zorns, oder Ungeduld würden enthalten können. Ja, es hat bisweilen Joseph selbst auch nach dem Fleisch wohl gewurmt, gemeint, und hat ihn das Fleisch dahin getrieben, daß er gesunken und fast kleinmüthig geworden ist. Denn die heiligen oder Gläubigen sind ja keine Klöße oder Steine, sondern sind Menschen, von andern Menschen geboren. Darum sollte man sich nach der Lehre Pauli halten, da er Röm. 6, 14. sagt: „Die Sünde wird nicht herrschen können über euch, sintemal ihr unter der Gnade seyd.“ Und zwar, wo die Gnade geherrscht und regieret hat, haben sie damit das Weinen und den Schmerz des Herzens verschlungen und überwunden.

4. Aus solchen Exempeln aber sind hergeschlossen die schönen trefflichen Ermahnungen in den Psalmen und Propheten. Als, im 27. Psalm V. 14.: „Harre des HErrn, sei getrost und unverzagt“ etc.; item, Hab. 2, 3.: „Die Weissagung wird ja noch erfüllet werden zu seiner Zeit, und wird endlich frei an Tag kommen, und nicht außen bleiben. Ob sie aber verzeucht, so harre ihrer; sie wird gewißlich kommen, und nicht verziehen“; item Jes. 49, 23.: „Du wirst erfahren, daß ich der HErr bin, an welchem nicht zu Schanden werden, so auf mich harren.“ Das Fleisch und menschliche Schwachheit sagt: Es wird ja zumal lang, die Hülfe verziehet sich fünf, zehen oder zwanzig Jahre, und man sieht der Anfechtung kein Ende. Bedenke aber, daß du göttliche Verheißung hast, und daß dazu Gott zum Freunde und Vater, der aus unermesslicher Güte, Gnade und Barmherzigkeit verheißen hat, er wolle für dich väterlich sorgen, du sollst sein Söhnchen sein. Er ist dein Vater, du bist sein lieber Sohn. Ja, es wird aber nichts draus, spricht das Fleisch? Antwort: Harre mit allen Heiligen und Gläubigen, so wird doch endlich das geschehen, was Christus Matth. 24, 13. verheißt: „Wer beharret bis ans Ende, der wird selig.“

5. Also werden wir zu unsern Zeiten auf mancherlei Weise geplagt und betrübt vom Satan und der bösen Welt. Es verfolgen uns Papst, Bischöfe und Cardinäle; die Ketzer oder Rottengeister plagen uns auch, als, Sacrament-schwärmer, Wiedertäufer und andere mehr. Aber was willst du Anderes thun, denn daß du des HErrn harrest? Denn du hast ja die Verheißung, daß Gott dein Vater sei. Da läßt die Hoffnung zum Glauben kommen: zweifele an der Verheißung gar nicht und harre, wenn gleich dasjenige, das verheißen ist, nicht so bald gegeben wird. Denn, wiewohl „die Hoffnung, die sich verzeucht, das Herz ängstet“, wie Salomo Spr. 13, 12. sagt, so muß man doch daran fest halten; dazu die Epistel an die Hebräer im 12. Cap. V. 7. vermahnt: „Gott erbeut sich euch als Kindern.“

6. Darnach sollen wir uns aufrichten, trösten und erwecken mit solchem Exempel, allerlei Beschwerung, Gefährlichkeit und Schmerzen zu leiden; und wo wir schon sollten getödtet werden, ist es doch gewiß, daß wir dadurch wahrhaftig von allem Uebel erlöset und errettet werden. Wie vor Zeiten den lieben Märtyrern auch nichts Böses widerfahren ist, da sie sind getödtet worden, sondern sie haben damit das Ende aller Trübsal erlangt; werden wir aber erhalten, so sollen wir beharren und immer in Hoffnung und Glauben der endlichen Erlösung warten. Denn Gott haßt oder verachtet uns nicht wie ein Feind, ob er sich wohl anders gegen uns erzeigt. Wir sollen aber also schließen und sagen: Ich weiß, daß ich getauft bin, daß ich den Leib des Sohnes Gottes gegessen und sein Blut getrunken habe, daß ich von Gottes wegen absolvirt und von allen meinen Sünden losgesprochen bin; ja, daß mir alle meine Sünden vergeben sind, und daß mir der Sieg wider den Teufel, Tod und Hölle verheißen; was will ich denn nun mehr haben?

7. Also hat Joseph auch gethan: wiewohl es sich mit seiner Befreiung sehr lange verzogen, hat er sich doch an die Lehre und Worte seines Vaters Jakob gehalten und dieselben bewahrt, und hat nicht gezweifelt, es würde die Verheißung, so seinem Vater geschehen, noch erfüllt werden. Und wo Gott dieselbige noch etwas länger verzogen hätte, so hätte er doch darauf geharrt und gehofft; wie Hiob Cap. 13, V. 15. auch von sich selbst sagt: „Wenn er mich schon tödten würde, will ich doch auf ihn hoffen.“ Denn das sind Worte des Glaubens, der Hoffnung und der Liebe gegen Gott und den Nächsten. Wie heftig auch die Welt und ihr Fürst im Papst und Fürsten tobt und wüthet, soll sie mir doch das Vertrauen zu Gott, so ich einmal gefaßt habe, nicht nehmen. So lange sollen es der Satan, viel weniger der Türke, nicht machen. Sie können uns zwar dies Leben, Gut und Habe, Weib und Kind und auch die Gesundheit wohl nehmen; so will ich aber doch für und für auf meinen lieben Gott hoffen.

8. Auf solche Weise erweckt und stärkt die heilige Schrift den Glauben an Gott, daß wir sagen: Herr Gott, du hast mich durch einen Bruder absolviret, getauft und mit deinem Leibe und Blute gespeiset; mache es mit deinem Knecht, wie es dir gefällt; ich will darüber nicht zürnen noch dich lästern, sondern will Alles mit Geduld ertragen. Denn ich will nicht, daß dein Bund, so du mit mir in der heiligen Taufe und Abendmahl des Herrn aufgerichtet, zunichte werden soll. Solche Leute spricht gewiß Christus selbst selig, die auf diese Weise bis an das Ende beharren und beständig bleiben. Von den Andern aber sagt Sirach Cap. 2, 16. also: „Wehe denen, so nicht beharren.“ Denn sie hoffen wohl eine Zeitlang und harren des Herrn, werden aber bald kleinmüthig, sinken dahin, verzagen und verzweifeln gar, wie man deß ein Exempel hat an dem Weibe des Tobias und auch des Hiob, welche vielleicht im Anfang nicht bald gewankt, aber darnach sind sie wegen des langen Verzuges und da das Unglück so gar groß geworden, heraus gefahren und haben greuliche Lästerworte geredet; als, Hiob 2, 9.: „Segne Gott und stirb.“ Als wollte sie sagen: Du hast noch immer Geduld und gute Hoffnung, und wirst doch darüber betrogen; du brauchst gar keinen Rath oder Hülfe, fährst noch immer fort mit Loben und Segen; ei, lobe immerhin und stirb ins Teufels Namen. Aber Hiob antwortete V. 10. und sprach zu ihr: „Du redest, wie die närrischen Weiber reden. Haben wir Gutes empfangen von Gott; und sollten das Böse nicht auch annehmen?“ „Der Herr hats gegeben, der Herr hats genommen; der Name des Herrn sei gelobt!“ Cap. 1, 21.

\quellseite{345}

9. Diese Exempel soll man fleißig merken, und dieselben den Herzen gottseliger Leute vorhalten und einprägen. Denn was fleischliche Menschen sind, die verachten es gar, und sehen nicht auf das Zukünftige der Tugenden, welche geschrieben und der Kirche zum Trost vorgehalten werden, auf daß wir lernen, daß unsere Trübsal ein gewisser Beweis und Pfand sei, daß wir Gottes Kinder sind.

10. Man hat einen Spruch Taulers, wiewohl er darin nicht nach Art der heiligen Schrift redet, sondern gar eine andere und fremde Rede führt. Der Mensch, sagt er, soll wissen, daß er großen Schaden gethan habe, wenn er nicht Gottes Werk sich gefallen läßt, nämlich, wenn ihn Gott will kreuzigen, tödten und den alten Menschen zunicht machen; welches nur durch Kreuz und Leiden geschieht. Denn da erfährst du an dir Gottes Werk, der dich züchtiget, dich hobelt und die groben Aeste abhaut; und der alles, was dich hindern mag an der Erbauung und Besserung zum ewigen Leben, mit der Axt abhaut, mit der Säge und Hobel, wie David sagt im 37. Psalm V. 7.: „Sei stille dem Herrn, und warte auf ihn.“

11. Auf solche Weise haben sich Abraham, Isaac und Jakob, und Joseph unserm Herrn Gott unterworfen, ihm stille gehalten und Hülfe von ihm erwartet. Wie in diesem gegenwärtigen Exempel jetzt Gott auch gegenwärtig ist, und die Zeit und Stunde findet, Joseph zu erretten: welcher selbst keinerlei Wege oder Weise hat sehen können, wie er hätte mögen befreit werden, sondern es einfach Gott allein heimgestellt hat, und dieweil im Finsterniß bleibt, sich damit nicht bekümmert, wie ihm möchte geholfen werden.

12. Das ist ein trefflicher Glaube gewesen. Darum kommt nun Gott und errettet ihn wunderlicher Weise. Welche Weise David im 105. Psalm V. 18–20. auch rühmt, da er sagt: „Sie zwangen seine Füße im Stock, sein Leib mußte in Eisen liegen; bis daß sein Wort kam, und die Rede des Herrn ihn durchläuterte. Da sandte der König hin, und ließ ihn los geben, der Herr über Völker hieß ihn auslassen.“ O welch eine treffliche wunderbare Erlösung ist das! Denn wo Joseph die Wahl im Kerker wäre gegeben worden, ob er wollte, daß sein Vater ihn vertreten, oder aber daß der Schenke der Wohlthat, so er ihm kurz zuvor erzeigt, gedenken sollte, würde er freilich derselben Bedingungen eine, welche es auch gewesen wäre, mit großer Begierde angenommen haben; ja, es hätte auch sein Herr oder Frau, wo sie Buße gethan und sich die begangene Sünde hätte leiden lassen, von seiner Unschuld können Zeugniß geben. Aber Gott hat der Weisen keine gebrauchen wollen, sondern der König selbst und der Herr über Völker hat ihn müssen heißen auslassen, auf daß seine Erlösung recht augenfällig und herrlich wäre; wie er denn auch endlich durch diese Erlösung noch zu königlicher Würde wird erhoben werden.

13. So läßt uns nun lernen dem Herrn vertrauen und auf ihn harren, wie Joseph geglaubt hat; wiewohl er nicht gewußt, auf welche Weise er würde erlöset werden, und hat auch die lange Zeit und Verzug mit Geduld ertragen können. Und lasset uns also schließen: Wiewohl mein Fleisch wird aufsehen zu murren, will ich mich doch inne halten und das Toben des Fleisches zwingen. Solcher Exempel sollen mir uns billig verwundern und dieselben lieb haben. Denn sie übertreffen der Mönche Fasten und harte Kleidung der Carthäuser weit. Denn diese Exempel ändern und zünden die Herzen der gottseligen Menschen durch den Glauben, Hoffnung und Liebe an.

14. Darum folgt nun weiter eine sehr schöne Zubereitung oder Anfang der Erlösung Josephs. Denn der Herr kommt nun nach zwei Jahren der Tage, wie es im Hebräischen lautet, das ist, nach den Tagen zweier Jahre, und gibt dem Könige einen Traum, dessen Auslegung er von Joseph bitten wird, und wird ihm auch die Verwaltung der Dinge, so durch den Traum zuvor verkündigt, gar übergeben.

15. Was nun die Grammatik betrifft, will ich den Lesern befohlen sein lassen, wiewohl die Beschreibung klar genug ist. Und sind sonderlich wohl zu merken die zwei Bedeutungen in den hebräischen Worten rioth und dakkoth, die einander etwas gleich lauten. Da erhellt aber nun und wird offenbar, was Gottes Rath gewesen sei mit Joseph und wie er mit den Menschen pflege umzugehen. Denn Joseph ist verloren und verdammt, gefänglich, gestorben und begraben gewesen; dieweil er aber verharrt und an Gott geglaubt hat, da hat er müssen ein Heiland der Welt werden; welches ihm nie war in Sinn gekommen, daß er es hätte hoffen, wünschen oder begehren dürfen. Denn er wäre gern mit der Wohlthat zufrieden gewesen, wenn er nur aus der Grube des schweren Gefängnisses geführt und seinem Vater wiederum zugestellt worden wäre, und dazu etwa ein Geschenk bekommen hätte, damit er so viel ehrlicher und herrlicher wiederum hätte zu den Seinen kommen mögen. Aber Gott sagt also: Lieber Joseph, du fehlest weit; weißt nicht, daß du einen solchen Herrn hast, der da mächtig ist, und „überschwänglich thun kann über alles, das wir bitten oder verstehen,“ Eph. 3, 20. Denn dermaßen pfleget Gott die Seinen zu erhöhen, und wird auch seine Erhöhung mit diesem Preiß geziert und gerühmt, daß sie unsere Bitte und Verstand weit übertreffe, und daß das menschliche Herz kann das nicht alles fassen oder begreifen, was Gott vom Gebet der Gottseligen gedenkt.

16. Betrachtet unser selbst Exempel, wie gering und schwach wir anfänglich wider den Pabst gewesen sind. Unser Seufzen war damals fast kleinlich, „wie ein gerader Rauch“, davon im Hohenliede Salomos Cap. 3, 6. Meldung geschieht. Aber ein solch groß Licht und Herrlichkeit haben wir nicht dürfen hoffen, wie wir heutiges Tages, Gott Lob, vor Augen sehen, daß die Lehre des Evangeliums die Gewalt des Pabstes unterdrückt und die Tyrannei des Teufels zerbricht, da so viel Menschen aus der Gewalt der Finsterniß sind errettet worden.

17. Also seufzt Joseph und sehnet sich, begehrt, erledigt zu werden und in seines Vaters Hause zu sein, wollte seinen Brüdern gern verzeihen, nur daß er in der rechten Kirche die Lehre von rechter Gottseligkeit hören möchte. Aber Gott antwortete ihm: Ich bin der Herr, dein Gott, welchen du anrufst, der du räucherst und opferst „geraden Rauch“, Hohel. 3, 6. Ich rieche und habe große Lust an dem angezündeten Weihrauch; denn der Geruch deines Gebets, Glaubens und Hoffnung hat nicht allein meine Nase, sondern auch den ganzen Himmel erfüllt. Ich will dich erlösen, daß es soll heißen, über alles, so du verstehen kannst.

18. Es ist aber der Jammer und die Schwachheit unsers Fleisches billig zu beklagen, welche uns nicht glauben läßt das, so wir selbst lehren. Denn wir haben ja das Wort, und wissen, daß es wahrhaftig und gewiß ist: aber im Gebrauch und in der Uebung ist unser Glaube schwach und zittert jämmerlich. Darum ist es eine große Wohlthat Gottes, daß er uns diesen Trost Jes. 42, 3. geschenket hat: „Das glimmende Docht werde ich nicht auslöschen.“ Am jüngsten Tage aber werden wir uns selbst verklagen und strafen um unserer Ungeduld willen, und werden sagen: Wenn ich noch auf Erden im vorigen Leben wäre, wollte ich mich nicht weigern, tausend Tod und Kreuz zu leiden. Aber jetzt ist es Zeit, daß wir unser Herz zur Geduld und Gehorsam erwecken.

19. Es jammert aber den Herrn unsere Schwachheit, und überschüttet uns ohne unser Verdienst mit vielen unzähligen Wohlthaten, daran wir selbst nie gedacht haben. Wie St. Augustins Mutter zeigt, daß ihr von Gott überschwänglich mehr gegeben sei, denn sie habe bitten oder begehren dürfen. Denn sie betet und seufzt für ihren Sohn, daß er nur möchte getauft werden, und da er getauft worden, ist sie gern gestorben. Denn sie war mit derselben Gabe und Bekehrung ihres Sohnes wohl zufrieden, und meinte, es wäre damit ihrer Bitte genug geschehen. Aber nach ihrem Tode ist Augustin in der Kirche ein trefflicher und nach den Aposteln der vornehmste Lehrer geworden, der viel unzählige Seelen aus vielfältigem Irrthume erlöset und dazu viel greuliche Ketzer überwunden hat. Dasselbe hat eigentlich seine Mutter mit ihrem Weinen und Seufzen zuwege gebracht. Ei! unser HErr GOtt ist ein feiner GOtt, wenn wir es nur glauben könnten.

\quellseite{346}

20. Eben dasjenige mag man in dieser gegenwärtigen Historie auch sehen. Es hat der König in Egypten selbst müssen einen Traum durch GOttes Schickung haben nur um der Ursache willen, daß Joseph mit großen Gnaden, Ehre und Herrlichkeit aus dem Kerker möchte geführt und erlöset werden. Und Pharao war zwar ein Heide, hat von dem GOtt der Hebräer gar nichts gewußt; doch läßt sich GOtt dazu herab, daß er auch mit einem heidnischen Könige durch einen Traum geredet hat, nicht um Pharaos, sondern um seines lieben Sohnes Joseph willen, der im Kerker lag. Da muß Pharao diesem Gefangenen dienen, und ganz Egypten muß Zuflucht haben zu diesem Sohne, der aus dem Gefängniß sollte befreit werden, nicht so sehr um Josephs selbst willen, als auch zur Wohlfahrt und Seligkeit des Königreichs Egypten und anderer umliegender Länder.

21. Also bringet nun das Weizenkorn hundertfältige Frucht, wo es in die Erde geworfen wird und erstirbt, Joh. 12, 24. Denn wer könnte alle seine Früchte erzählen? Das ist noch ein Geringes und nur etwas Leibliches, daß er alle Egypten vom Hunger und Theurung errettet, und Pharao selbst auch vom Tode erlöst mit andern Völkern der umliegenden Länder; wiewohl es eine große Wohltat ist, so viel tausend Menschen erhalten, daß sie nicht umkommen; oder wenn ja in solcher Beschwerung der langwierigen Theurung etliche dahin gestorben, daß doch gleichwohl der größere Theil erhalten worden ist. Aber das sind endlich die rechten großen Gaben, welche David im 105. Psalm V. 21. 22. rühmt, da er sagt: „Er setzte ihn zum Herrn über sein Haus, zum Herrscher über alle seine Güter, daß er seine Fürsten unterwiese nach seiner Weise, und seine Aeltesten Weisheit lehrete.“ Das ist ein groß wunderbar Ding, daß Joseph, der in seines Vaters Hause durch seiner Brüder Neid beide Königreich und Priesterthum verloren hat, die Herrschaft und Regiment über ganz Egyptenland bekommen hat, welches er vom Tode und Hunger, beide leiblich und geistlich, errettet. Denn er ist da ein Regent und Lehrer der Verheißungen GOttes geworden.

22. Also hat David in diese Historie etwas tiefer hineingeschauet und erkannt, wie nütze und heilsam Joseph den Königreichen Egypten gewesen sei. Wie viel seine Leute wird er gemacht haben! Er hat die Fürsten und den König selbst gelehrt und sogar den ganzen königlichen Hof zum Glauben bekehrt; ja, er hat den rechten Gottesdienst gezeigt, desgleichen Priester verordnet, welche darnach Andern den Weg zeigen und sie unterweisen möchten. Und kürzlich davon zu reden: er ist der rechte Christus in Egypten gewesen, und noch mehr, wie unser HErr Christus sagt Joh. 14, 12.: „Wahrlich, wahrlich, ich sage euch, wer an mich glaubet, der wird die Werke auch thun, die ich thue, und wird größere denn diese thun.“ Christus hat ein kleines Volk im engen Winkel des einigen Landes Juda bekehret, hat etliche tausend Menschen mit wenigen Brode gespeist: dieser Joseph hat ganz Egypten und die umliegenden Völker und Königreiche gespeist, beides leiblich und geistlich. Das heißt Joh. 15, 16.: „Ich habe euch gesetzt, daß ihr hingehet und Frucht bringet, und eure Frucht bleibe“, auch unter den unbeschnittenen Heiden, die dem Gesetze der Väter nicht unterworfen sind; wie Joseph in Egypten Frucht gebracht.

23. Darum lernet vertrauen, glauben, hoffen und auf GOtt harren, und daß ihr am Wort hanget und mit dem heiligen Paulo Röm. 8, 18. saget: „Dieser Zeit Leiden ist der Herrlichkeit nicht werth, die an uns soll offenbaret werden“, es ist zu geringe; item 2 Cor. 4, 17. 18.: „Unsere Trübsal, die zeitlich und leicht ist, schaffet eine ewige und über alle Maße wichtige Herrlichkeit, uns, die wir nicht sehen auf das Sichtbare, sondern auf das Unsichtbare“ etc.

24. Und Joseph wird ohne Zweifel mit sich selbst gezürnt und gesagt haben, da er frisch und gesund davon gekommen ist: Warum habe ich doch im Gefängniß gemurrt? Ich sehe ja, daß der allergnädigste GOtt mir beigestanden, da ich betrübt gewesen und gemurrt habe. Ich hätte mich lieber sollen freuen, lachen und fröhlich sein in der Noth und Trübsal. Ach lieber HErr GOtt, verzeihe mir meine Schwachheit! Denn diese großen Wohltaten, so du durch mich so vielen Menschen erzeigest, habe ich mit keinem Weinen oder Seufzen, mit keinem Kreuz oder Leiden verdienen können. Du bist ein solcher GOtt, der da mehr gibt, denn ich immer habe hoffen oder begreifen können. Ja, also pflege ich den Meinen zu helfen, sagt GOtt; so sind meine Werke gethan, daß sie auch allen Verstand und menschliche Vernunft übertreffen. Wie denn solches das Exempel Danielis in Babel auch bezeuget. Das ist aber eine große Wohltat, daß GOtt darüber auch noch verschont und verzeiht dem zerstoßenen Rohre und glimmenden Dochte, und ist zufrieden auch mit einem kleinen Fünklein, wenn wir schon murren, wenn wir nur nicht gar vom Glauben abfallen, noch von der Hoffnung abweichen, sondern daß wir mit Hiob Cap. 13, V. 15. sagen: „Wenn er mich schon tödten würde, will ich doch auf ihn hoffen.“

25. Auf diese Weise erweckt nun GOtt dem Jüngling Joseph einen herrlichen berühmten Erlöser, nämlich, den König von Egypten selbst, welchen der HErr würdiget, ihn so freundlich und gnädiglich durch einen Traum heimzusuchen. Solche Träume aber sind wahrhaftig rechte Weissagungen; wiewohl Pharao dieselbigen noch nicht versteht, sondern hält unserm Herrn GOtt nur einfach stille in dem, daß er mit ihm redet.

26. Dieweil aber nun unser HErr GOtt so sehr und so fleißig für uns sorgt und auf uns Acht hat, auch so gar nahe bei uns ist: was ist es denn, daß uns so sehr irre oder verzagt machen könnte? Denn dies Exempel lehrt, daß GOtt nicht allein die, so da wachen und arbeiten, bewahrt und versorgt; sondern auch im Schlafe, da wir gar nichts thun, wenn wir auch nicht wissen, daß wir leben, alsdann hat er auch Acht auf unsere Träume, zählt alle Sorge und Seufzer auch derer, die da schlafen.

27. Ich bin zwar meinem Fleische feind, weil ich weiß, daß dies alles wahr ist, und mir es doch nicht mit Ernst zu Herzen gehen lasse. Denn welchen Menschen wollte doch der große Fleiß und Sorge, so GOtt für uns hat, nicht bewegen, daß er nicht allein die Haare auf unserm Haupte, sondern auch unsere Thränen zähle und in einen Sack sammle, wie der 56. Psalm V. 9. sagt: „Zähle meine Flucht, fasse meine Thränen in deinen Sack. Ohne Zweifel, du zählest sie“; kein Thränchen muß umsonst geschehen, es wird aufgezeichnet mit großen, mächtigen Buchstaben im Himmel.

28. Derhalben seufzen, weinen und leiden wir nicht vergeblich. Es ist alles angeschrieben und verzeichnet in unsers HErrn GOttes Buche. Das heißt ja genau darein sehen. GOtt sieht uns nicht von ferne an, sondern ist nahe bei uns, und zählt meine Angst und Gedanken, Traurigkeit und Schmerzen, auch des Nachts, und will, daß sie sollen vergolten werden viel reichlicher, denn ich es begehrt oder selbst verstanden habe. Daher ist der Spruch Pauli gekommen, welcher uns für und für im Sinne und Herzen liegen sollte, 2 Cor. 4, 17.: Unsere Trübsal, und was wir leiden, ist zeitlich und sehr leicht, dieweil wir wissen, daß die Herrlichkeit und Erlösung viel größer sein wird, denn wir mit einigen Gedanken oder Worten erreichen können; und GOtt hat doch so genaue Acht auf das leichte und kurze zeitliche Leiden der Gottseligen.

29. Wes Schuld ist es denn nun, daß wir so jämmerlich geplagt und geängstiget werden? Nämlich der Schwachheit unsers Glaubens und Hoffnung, und dazu der Bosheit unsers Fleisches. O schlag todt, lieber HErr GOtt! Soll ich einen solchen GOtt haben, der so gar sorgfältig ist, daß er mich auch des Nachts, wenn ich schlafe oder träume, bewahrt? Ist es nicht genug, daß er die, so da wachen, schützt und unsern Gang richtet? Er hat über das auch noch Sorge für unsern Schlaf und Träume? Das heißt wahrlich recht einen gnädigen GOtt haben, der uns bewahrt, und nicht zuläßt, daß ein Haar von unserm Haupte falle, oder daß der Augapfel tröpfele, oder aber doch jeglichen Tropfen, so aus den Augen fällt, zählt. Weinst du, so hat er ein gülden Becken oder Schüssel und fängt die Thränen auf. Wer sind aber die, deren Thränen er so genau auffängt? Antwort: Es sind arme elende Sünder. Wer ist es denn, der sie sammelt? Das ist GOtt, der Schöpfer aller Dinge. Dies sollte man nun oft bedenken, den Glauben, die Hoffnung und Liebe in uns damit zu erwecken, dieweil wir ja das Wort haben, dazu die Exempel und unsere eigene Erfahrung der unermesslichen Güte GOttes.

\quellseite{347}

30. Das Erste aber nun, das GOtt mit Pharao im Traum redet, ist GOttes Wort; und er redet es darum, daß es endlich soll geoffenbart und verstanden werden. Denn es wird ja nicht dazu gegeben, daß es verborgen und dunkel bleiben soll; wie der Herr Jes. 55, 11. sagt: „Mein Wort soll nicht wieder zu mir leer kommen“, sondern seine Wirkung haben, wiewohl zu etlichen Zeiten etwas dunkeler, zu etlichen aber klarer.

31. Wie wir im Papstthume das Wort gehabt, und ohne Zweifel Viele im Glauben an dasselbe selig geworden sind. Und es klingt das Wort auch jetzt in unsern Kirchen von GOttes Gnade, aber viel heller als in den vorigen Zeiten. Wir nehmen aber täglich zu von einer Klarheit zu der andern, als vom Geist des Herrn, 2 Cor. 3, 18. Denn wir verstehen es von Tage zu Tage immer klarer, und kennen doch nichts Anderes, sondern werden in dasselbe Bild verkläret. Wir lernen je länger je mehr daran und doch nichts Anderes. Also, wenn ich im offenen Felde wandele und sehe von ferne einen Menschen daher gehen, da dünkt mich derselbe erst ein Pferd oder Baum zu sein; wenn ich aber näher hinzu komme, fange ich allmählig an, mehr und mehr zu erkennen, daß es ein Mensch sei, bis ich diesen oder jenen Mann, eine bestimmte Person sehe. Also geht es auch mit unserer Erkenntniß: die ist erst noch verworren, darnach aber wird sie immer klarer und bestimmter. Denn anfänglich, wenn uns GOttes Wort vorgetragen wird, wird es nicht alsbald vollkommen verstanden, wird aber doch gefaßt und angenommen, wiewohl in Schwachheit; nicht so, als sollte die Erkenntniß Christi verdunkelt, sondern auf daß dieselbe von Tag zu Tag geübt und gemehret werde, bis sie endlich vollkommen werden möge.

32. Daselbige geschieht auch also in der Veränderung unseres Willens, welche auch täglich zunimmt, wenn wir allemahl mehr und mehr lernen hoffen, vertrauen und geduldig sein. Daher ist es ein und einerlei Glaube, der da anfängt, zunimmt und vollkommen ist. Denn wir werden allezeit in daselbige Bild verklärt; das Wort aber ist darum gegeben, daß es soll geübt werden und daß es dich auch täglich und alle Stunden üben soll, bis du wachest und ein vollkommener Mann werdest. Also verstehst zwar Pharao noch nicht und erkennet auch GOtt nicht, der mit ihm durch einen Traum redet.

\begin{quote}
\textbf{V. 7. Da erwachte Pharao, und merkte, daß es ein Traum war.}
\end{quote}

33. Die Träume, so von GOtt eingegeben werden, werden den Träumenden also eingeprägt, daß sie heftigere und größere Bewegung fühlen und empfinden, denn sonst in gemeinen Träumen anderer Menschen und auch unvernünftiger Thiere. Ein solch Fühlen ist auch an Pharao gewesen, und GOtt hat sein Herz bewegt, daß er auf den Traum Acht gehabt und denselben behalten hat, sonderlich da er ihm zum andernmal vorgekommen ist. Wie Nebucadnezar, da er des Traumes nicht mehr gedacht, hat er doch desselben nicht können vergessen, weil er ihm so tief eingeprägt war. Und solche Träume werden gemeiniglich von GOtt ausgelegt, weil sie auch von GOtt hergekommen sind.

\begin{quote}
\textbf{V. 8. Und da es Morgen ward, war sein Geist bekümmert, und schickte aus, und ließ rufen alle Wahrsager in Egypten, und alle Weisen, und erzählte ihnen seine Träume. Aber da war keiner, der sie dem Pharao deuten konnte.}
\end{quote}

34. Pharao fühlt und empfindet, daß es eine Offenbarung sei, und versteht doch dieselbe nicht. Darum kommt der heilige Geist und legt es ihm aus. Es ist aber diese Betrübniß Pharaos nicht so gar heftig und erschrecklich gewesen, sondern sein Herz ist etwas erschrocken und bezeuget gewesen, daß er gedacht hat, was doch diese Träume bedeuten möchten. Darum schickte er aus und ließ rufen alle Wahrsager und Weisen in Egypten.

35. Harthummim, auf Hebräisch, sind 2 Mos. 7, 11. die Priester, so sich wider Mose gesetzt. Denn es hat eine jegliche Nation ihre Lehrer gehabt: die Magier im Morgenlande, Chaldäer in Babylonien, die Sophöi oder Weisen in Griechenland. Die Sapientes oder Weisen sind bei den Römern berühmt gewesen wie Prälaten oder Priester, wie die Römer neben den Rechtsverständigen auch Priester und Lehrer gehabt, welche die Lehre von der Religion, und auch von andern natürlichen und magischen Künsten öffentlich gelehrt. Und es hat ihnen an Zuhörern nicht gemangelt, so dieselben Künste auch gelernt haben.

36. Oder wo das jemand besser gefällt, mögen wir es dafür halten, daß diese Weisen von Könige dazu angestellt und besoldet gewesen, daß sie die Jugend unterrichten und lehren sollten, und daß sie ihre Schulen gehabt, darin sie die Religion und sonst allerlei Künste, beides gute und böse, gelehrt haben. Denn dies wird gemeiniglich untereinander gemengt, obwohl die Lehre gereinigt und recht angerichtet ist, vornehmlich in großen Landen oder Königreichen. Wie wir zu unsern Zeiten die rechte wahre Religion, die Andern aber, als, die zu Paris und die Lehrer an andern Hochschulen, die papistische und falsche Religion lehren; und führen doch gleichwohl den Titel und Namen, daß sie Theologen, Doctoren und Magister heißen. Also sind vorgzeiten unter den heidnischen Philosophen auch Secten gewesen. Denn der Satan hat allezeit auch seine eigenen Lehrer in allen Arten der Studien.

37. Diese aber sind Vorgänger, Meister und Ausleger der Religion und der Wunder gewesen, so etwas Künstliches zu bedeuten pflegen, welche der König alle hat berufen lassen, und denen er seinen Traum erzählt hat, welchen sie nicht haben deuten können. Aber das ist der rechte Weg, dadurch sie geführt und bereitet werden sollen auf die Schule Josephs, der soll sie alle zur Schule führen, sie unterrichten und die rechte wahre Religion lehren. Und dies ist eine Vorbereitung auf die Frucht, welche Joseph in Egypten hat bringen sollen, und ist der Dank, den er seinen Verräthern und Stockmeistern gegeben hat. Denn er hat die ärgsten und greulichsten Thaten, so sie an ihm begangen, mit den allerbesten Früchten vergolten und bezahlt. Wie unser Herr GOtt seine Sonne scheinen läßt über dankbare und undankbare Menschen.

\begin{quote}
\textbf{V. 9—13. Da redete der oberste Schenke zu Pharao, und sprach: Ich gedenke heute an meine Sünde; da Pharao zornig ward über seine Knechte, und mich mit dem obersten Bäcker ins Gefängniß legte, ins Hofmeisters Hause; da träumete uns beiden in Einer Nacht, einem jeglichen sein Traum, deß Deutung ihn betraf. Da war bei uns ein ebräischer Jüngling, des Hofmeisters Knecht, dem erzählten wirs. Und er deutete uns unsere Träume, einem jeglichen nach seinem Traum. Und wie er uns deutete, so ists ergangen; denn ich bin wieder an mein Amt gesetzt, und jener ist gehenkt.}
\end{quote}

38. Das wird alles gesagt, daß es eine Gelegenheit sein soll zu der Befreiung Josephs. Der Schenke redet mit gebührlicher Ehrerbietung zum Könige und erzählt seine Sünde, dankt dem König für seine Wohlthat und Gnade, dadurch er war erlöst worden. Er hat aber diese Dinge darum so weilläuftig erzählt, daß er also mit Fug des Jünglings Joseph gedenken möchte. Wiewohl er damit lang genug verzogen hat; es sei denn vielleicht, daß er vor derselben Zeit nicht Gelegenheit gehabt, solches zu erzählen, oder, was noch wahrscheinlicher ist, daß der König und das Hofgesinde dießelbige Erzählung als ein unnütz oder vergeblich Ding verachtet, dieweil der König niemals einen Traum gehabt, noch auch eines solchen Auslegers bedurft hat. Droben, Cap. 40, V. 23., hat aber doch Mose gesagt, daß Josephs sei vergessen worden; darum kann dieser Schenke nicht so gar unschuldig gehalten werden; wiewohl es an aller Herren Höfen fast gemein ist, daß sie solche Dinge verachten, bis das Wasser über die Röhre geht. Jetzt aber will es nun die Zeit und hohe Noth erfordern, daß man auf Joseph auch muß Acht geben; da geht es nun an.

\begin{quote}
\textbf{V. 14. Da sandte Pharao hin, und ließ Joseph rufen; und ließen ihn eilend aus dem Loch.}
\end{quote}

39. Man hat durch das ganze Königreich Egypten einen Ausleger gesucht unter allen Weisen und Wahrsagern und wird ohne Zweifel eine große Belohnung darauf gesetzt worden sein. Wie Nebucadnezar und Belsazer in der Historie von Daniel auch thun. Also ist Alles mit großer Herrlichkeit und großer Zubereitung zugegangen. Und dieweil sie alle sich gleichsam in die Wette darum gedrungen, und Einer vor dem Andern den König hat willfahren und bei ihm Dank verdienen wollen, darum eilen sie nun also, Joseph aus dem Gefängniß los zu lassen, und hindert sie nun daran gar nichts, daß er fremd, elend, gefangen und gebunden ist im Kerker. Denn sie merken und verstehen, daß viel daran gelegen sei, daß des Königs Traum möge recht gedeutet werden. Und es ist gewißlich wahr, daß auch die großen Herrschaften und Königreiche ihre Gefahren und Uebelstände haben. Es ist aber gemeiniglich mehr Furcht vor dem Bösen denn Hoffnung auf zukünftiges Gutes. Und diese Träume waren derartig, daß leichtlich hat können gespürt werden, daß die fetten Rinder eine gute glückselige Deutung haben müssen, die mageren aber würden nichts Gutes zu bedeuten haben.

\quellseite{348}

40. Darum sind sie alle etwas betreten und in Sorgen gewesen, haben Hoffnung und doch immer Sorge dabei gehabt, und können es nicht ausdenken, ob die Träume Frieden oder Krieg, Gutes oder Böses bedeuten. Und es hat auch nicht gemangelt an Etlichen, so es haben errathen wollen und die auch mancherlei errathen haben; wie Lyra aus der Juden Commentaren erzählt. Und ein Jeder hat bei dem Könige Dank verdienen wollen und das Beste thun, auf daß er den Ruhm davon bringen möchte. Aber wie der 113. Psalm V. 7. 8. sagt: „Der HErr erhöhet den Armen aus dem Koth, daß er ihn setze neben die Fürsten.“ Er will die Gabe der Auslegung den großen Hansen und Weisen am königlichen Hof nicht geben, noch auch den Priestern und Lehrern in den Kirchen und Schulen; sondern er gibt es dem armen, elenden, verachteten und verdammten Joseph. Und nun folgt auch darauf, daß er plötzlich erhöhet wird.

\section[Zweiter Theil]{Zweiter Theil. Wie Pharao von Joseph die Auslegung seiner Träume verlangt, und wie Joseph sich dabei verhalten}

\begin{quote}
\textbf{V. 14. 15. Und er ließ sich bescheren, und zog andere Kleider an, und kam hinein zu Pharao. Da sprach Pharao zu ihm: Mir hat ein Traum geträumet, und ist niemand, der ihn deuten kann; ich hab aber gehört von dir sagen, wenn du einen Traum hörest, so kannst du ihn deuten.}
\end{quote}

41. Der König hat Joseph sehr gnädig und freundlich angeredet, und erzählt, wie ihn Andere gelobt haben. Andere, sagt er, haben mir gerühmt von deiner trefflichen großen Weisheit und hohen Gaben, daß du nicht allein deine eigenen Träume deuten könnest, sondern auch die, so Andere gehabt oder gesehen haben, welche selbst weder ihre eigenen noch andere Träume verstehen können. Dies ist ein großer trefflicher Ruhm und werden auch große Belohnungen hernach folgen; wiewohl ihm der König jetzt noch nichts verspricht.

\begin{quote}
\textbf{V. 16. Joseph antwortete Pharao, und sprach: Das stehet bei mir nicht; GOtt wird doch Pharao Gutes weissagen.}
\end{quote}

42. Dies ist ein schöner Text, der sonderlich wohl zu merken ist, nämlich, von weltlicher Obrigkeit. Denn es sind nicht Worte Eines, der dem Könige geheuchelt hätte, da er sagt: „GOtt wird doch Pharao Gutes weissagen“, sondern sind Worte Eines, der die Obrigkeit gebührlicher Weise ehrt. Sie sollen auch nicht also verstanden werden, als zweifelte Joseph daran, ob er auch den Traum, so ihm vorgehalten war, verstehen oder deuten könnte. Denn er ist ohne Zweifel ein frommer heiliger Mann gewesen voll Heiliges Geistes und hat nichts nach eigenem Gutdünken und Vornehmen angefangen; wie wir, wenn wir hören von Alexander dem Großen, von Scipio und andern tapfern Helden mehr, welche große herrliche Dinge ausgerichtet ohne GOttes Wort und ohne rechte Anrufung, uns träumen lassen und uns einbilden, wir seien ihnen gleich an Stärke und andern Tugenden, und könnten auch alle ihre herrlichen Thaten ausrichten und ihnen dieselben nachthun. Aber solche ihre Tugenden und Thaten sind nur zufällige und besondere, ihnen eigenthümliche Gaben gewesen, und sind nicht Exempel, so uns nachzufolgen vorgestellt wären.

43. Ein solcher Mann ist Joseph nicht gewesen, der sich auf die Vernunft, oder seines freien Willens Rath und Weisheit verlassen hätte, sondern hat vom Gebet des Herrn geredet Tag und Nacht und mitbrünstigem Glauben GOtt angerufen. Der Glaube aber ist nicht müßig, bleibt nicht ohne Frucht, sondern übt sich immer in Wort und Gebet. Wiederum aber, obwohl gottlose heidnische Menschen, die ohne Wort und ohne Gebet sind, viel herrliche Dinge ausrichten, wie Alexander und andere tapfere Helden, so sind es doch zufällige und vereinzelte Gaben; wie die Sonne ihren Glanz zugleich auf die Guten und Bösen ausbreitet: also sind auch alle zeitlichen Güter.

44. Die Bibel aber oder Schrift lehrt, daß man das weltliche Regiment nicht aus eigenem Gutdünken und Vornehmen verwalten, sondern daß man in allen Dingen göttlich, im Wort und rechter Anrufung einhergehen solle. Darum hat Joseph auch nichts gethan ohne Gebet zu GOtt, vornehmlich da er ist berufen worden, daß er dem Könige seinen Traum deuten sollte; da wird er ohne Zweifel also gebetet haben: HErr GOtt, ich werde jetzt berufen, des Königs Traum zu deuten; gib mir Gnade und den Heiligen Geist, daß ich ihn recht deuten möge. Darum kommt er daher in rechtem Glauben und zweifelt nicht an der Deutung des Traums; wie Daniel Cap. 2, 18. auch mit seinem demüthigen Gebet von GOtt erlangt, daß ihm das Gesicht des Königs Nebucadnezar und zugleich die Auslegung desselben gezeigt und geoffenbart wird.

45. Derhalben sind dies Worte Eines, der sich vor der Obrigkeit demüthigt, da er sagt: „GOtt wird doch Pharao Gutes weissagen.“ Also wollte er sagen: Wenn ich ihn nicht antworten oder den König berichten würde, der ich dieser großen Gaben unwürdig bin, so wird doch GOtt einen andern Ausleger finden.

46. Diese Ehrerbietung Josephs gegen den König ist wohl zu merken und wir sollen dieselbe immer in frischem Gedächtnis behalten. Denn wir sollen es dafür halten, daß weltliche Polizei oder Regiment GOttes Ordnung sei. Wie auch der Ehestand und die Kirche, und alles, was Gutes in diesen Ständen geschieht, göttlich ist, und die Gottlosen haben solches mit ihrem Gebet von GOtt erlangt. Denn also sagt die Weisheit GOttes in den Sprüchen Salomos am 8. Cap. V. 15.: „Durch mich regieren die Könige, und die Rathsherren setzen das Recht.“ Wiederum aber sagt auch Salomo Cap. 28, 2.: „Um des Landes Sünde willen werden viel Aenderungen der Fürstenthümer“; item im selben Capitel V. 15.: „Ein Gottloser, der über ein arm Volk regieret, das ist ein brüllender Löwe und gieriger Bär.“ Und Hiob sagt am 34. Cap. V. 30.: „Er lässet über sie regieren einen Heuchler, das Volk zu drängen.“ Von denen wird nicht gesagt: Durch mich herrschen die Fürsten; sondern wenn unser HErr GOtt einem Lande günstig ist, so regieret er selbst darin durch fromme Obrigkeit; wo er aber zürnt, alsdann machen die Fürsten unleidliche Gebote und Ordnungen, schinden das arme Volk und machen des gottlosen Wesens und Abgötterei nur viel im Lande; wie das am Manasse, Ahab und Andern mehr zu sehen ist.

47. Dieweil nun Joseph weiß, daß die weltliche Polizei ein göttliches Amt und Ordnung ist, demüthigt er sich und schreckt zurück, da er berufen wird, solcher hohen Majestät zu rathen. Denn er ist ein weiser verständiger Mann gewesen, darum hat er gedacht: Was soll ich armer Mensch dem Königreich Egypten und einem solchen gewaltigen König rathen oder helfen können? Es ist ein göttlicher Stand: denselbigen zu verwalten gebühret nicht Kindern, noch gottlosen Leuten. Manchmal geschieht es aber, daß unser HErr GOtt durch eine Gabe von ohngefähr etwa einen tapfern Helden gibt, als, einen Alexander, Scipio oder Augustus \&c. Aber wo man ihrem Exempel nachfolgen wollte, das taugt gar nicht. Wie Absalom seines Vaters Königreich mit Gewalt einnimmt und sich die Gewalt anmaßt, Gericht zu halten und sonst das ganze königliche Amt zu verwalten, 2 Sam. 15, 4.: „O wer setzet mich zum Richter im Lande, daß jedermann zu mir käme, der eine Sache und Gericht hat, daß ich ihm zum Rechten hülfe?“

48. Unverständige ehrgeizige Leute verstehen nicht, welch ein groß Ding es sei, solcher göttlichen Ordnung in den Regimenten zu dienen. Denn entweder träumt ihnen das, so in der Komödie steht, wie jener Narr sagt: Me regem esse oportuit, daß ist: Ich sollte billig ein König sein; oder wo die Könige und Fürsten verstehen, daß sie zum Regimente geboren seien, können sie sich selbst auch bald zuschreiben und anmaßen solche Weisheit und Vermögen, so zum Regimente dienlich und vonnöthen ist. Nun ist es doch wahrlich nicht an eines Menschen Rath, Weisheit oder Verstand gelegen, über andere Leute zu regieren, und ist eine öffentliche Thorheit und Unsinnigkeit dazu an den gottlosen Leuten, die da begehren zu regieren und sich selbst ins Regiment eindrängen, welches sie endlich umkehren und gar verderben, es sei denn, daß unser HErr GOtt ohngefähr sich solcher gemeinen Noth erbarme, und etwa einen Alexander, oder dergleichen andern tapfern Helden gebe. Doch ist es gleichwohl nicht recht, daß man solchen Exempeln nachfolgen, und etwas Gewisses daraus schließen will und sagen: Jene haben ohne göttliches Licht glücklich und mit großem Ruhme das Regiment über Andere verwaltet: darum werde ich solch Glück auch haben; das, sage ich, solt nicht; denn es ist eine große, unendliche Ungleichheit.

\quellseite{349}

49. Die derhalben zum Regimente oder sonst zu andern Aemtern greifen, die zu Hofe zu Rathe gezogen werden, sollen sich ja vor allen Dingen zur Demuth und Gottesfurcht schicken; wie Pharao, der doch ein heidnischer König ist, sich auf seine Klugheit nicht verläßt, sondern sucht mit Ehrerbietung Rath bei Andern, die in weltlichen Sachen, das Regiment belangend, und auch in der Religion Verstand haben. Und Joseph hat Gott noch viel mehr gefürchtet und geehrt, und antwortet deshalb dem Könige auf seine Bitte mit aller Demuth: Ach! was bin ich! Gott kann es wohl durch einen Andern machen.

50. Diese Exempel soll man fleißig merken, und nicht zweifeln, daß die weltliche Obrigkeit ein göttlicher Stand sei, von Gott geordnet zum Nutzen dieses zeitlichen Lebens, und auch der lieben Kirche, welcher die Obrigkeit dient, wenn sie den Frieden liebt und erhält; wiewohl die Kirche für sich selbst ein ander Amt hat, welches zum künftigen Leben gehört. Darnach weil man weiß, daß es Gottes Ordnung ist und nicht in unserer Macht, Rath oder Weisheit steht, sollen wir unsere Kniee beugen, und mit aller Ehrerbietung und Furcht Gott anrufen und sagen: Lieber Herr Gott, ich bin ein Fürst, Regent, Rath, Amtmann, Prediger, Doctor oder Lehrer; nun ist aber das Regiment dein, dein ist das Reich, Gericht und aller Rath; gib mir die Gnade und Macht, daß mein Rathen und meine Thaten in diesem Amte glücklich sein mögen.

51. Man soll nicht also gesinnet sein, wie Viele sich selbst ganz thöricht rühmen, daß sie so viel Verstand, Rath und Kunst bei sich haben, daß sie wohl gehen Regimenten vorstehen und dieselben zu verwalten wüßten. Solche Ueberhebung oder Ruhmredigkeit ist eine Ursache, daß so viel Unordnung und Unruhe im weltlichen Regimente, in Ländern und Städten vorfallen.

52. Unsere Fürsten und Monarchen haben die Gaben nicht, daß sie ihr Geld, ihre Macht, ihre Wehr und Herrschaften auf den auswendigen fremden Feind, den Türken, wenden und für die Kirche streiten; sondern sie fahren zu und wüthen unter einander mit Würgen und Todtschlagen, und verderben also selbst das liebe Vaterland; denn sie beten nicht, sondern sind stolz und vermessen. Sie wissen, daß sie Könige und Herren sind, träumen und denken dabei, es stehe das Regiment in ihrer Macht; darum fallen sie über einander her und schlagen einander selbst todt. So recht, also soll es gehen. Sie verlassen sich auf ihre Weisheit ohne Gottesfurcht: da hat die Stimme der Weisheit bei ihnen keine Statt: „Durch mich regieren die Könige, und die Rathsherren setzen das Recht“, Spr. 8, 15. Diese setzen und ordnen zwar auch, aber eitel unbillige unnütze Dinge, so dem Vaterlande und den Gemeinen hin und wieder nur schädlich sind.

53. Wenn jezund David das Regiment bei uns haben sollte, da würde er beten, und das Reich immer weiter ausbreiten und wider die Feinde schützen. Unsere Könige lassen es geschehen, daß die Türken mit ihrem Kriegsvolke herein fallen, und in Ungarn, Oestreich und Deutschland Schaden thun, und schützen dagegen ihre armen Unterthanen nicht. Denn sie bedenken nicht einmal, daß das Regiment oder Verwaltung ihrer Herrschaft Gottes Werk sei, sondern sei vielmehr ein Werk ihres Verstandes und ihrer eigenen Klugheit. Daher ist erstlich gekommen das gemeine Sprüchwort, daß man sagt: Aut regem, aut fatum nasci oportuit: Es hat entweder ein rechtschaffener König, oder ein Narr müssen geboren werden. Denn es kann nicht anders sein, es sei denn, daß Gott von ungefähr eine sonderliche Gabe in die Welt austheile und gebe etwa einen weisen und glückseligen Fürsten ohne sein Wort.

54. Dies ist die Lehre, so uns im Exempel Josephs vorgehalten wird, der sich demüthigt, und unwürdig achtet, daß er zu dem Amte berufen werden soll, daß er dem Könige seine Träume deuten und eines solchen großen Königes Rath sein soll. Also wollte er sagen: Gott bedarf meines Dienstes nicht, er kann wohl einen Andern finden, der viel weiser und gelehrter ist, denn ich; ich will aber doch sonst thun, was ich vermag. Auf solche Weise soll man zum Kirchenamte, zu weltlicher Herrschaft und sonst zu allerlei andern Aemtern greifen, mit rechter Demuth und mit Gottesfurcht. Denn wo der HErr nicht die Rathschläge regieret, wird der Ausgang nicht glücklich und heilsam sein, wie er Spr. 8, 14. selbst sagt: „Mein ist beide Rath und That“; item: „Ich habe Verstand und Macht“ und Sieg. Wenn ich nicht rathe, so geht man irre; wenn ich es nicht thue, so richtet man nichts aus.

55. Also auch im Ehestande sollst du dir nicht vornehmen, daß du ein solcher Ehemann sein werdest, der dazu von selbst geschickt genug sei, das Haus zu regieren. Du sollst auch nicht vornehmlich sehen auf die unsinnigen Weiber fleischlicher Lust, damit du das Weib freundlich umfangen und allerlei Wollust pflegen mögest ohne allen Verdruß, Beschwerung und Kreuz; sondern sollst fleißig behalten, daß es ein göttlich Amt sei und daß Gott selbst dich und das Weib zusammengefüget habe. Denn Gott hat Mann und Weib geschaffen und sie zusammengefüget. Darum sollst du also beten: Herr Gott, dieser Stand ist deine Ordnung; ich bitte dich, daß du mir ein solch Mägdlein gebest, damit ich friedlich und ehrlich, dazu auch in gegenseitiger Liebe leben möge; und daß mein Wille ihr Wille sei, und wiederum ihr Wille mein Wille auch sei; wöllest mir auch durch deinen Segen Kinder und Erben geben, die ich christlich und wohl auferziehen möge.

56. Aber wie die Fürsten in weltlicher Polizei diese Lehre verachten, also denken auch die Ehemänner im Hausregimente daran nicht. Darum sehen wir auch, daß die Welt voll Klagens und Jammerns ist. Und dasselbe kommt durch Gottes gerechtes Gericht, sintemal die Leute nicht im Worte und Gebet wandeln, sondern nur in eigener Vermessenheit einhergehen, und sich selbst groß Glück dichten und daß ihr Ding einen herrlichen Ausgang haben werde, so es doch damit endlich viel anders geht, denn sie gedacht hatten. Sie glauben nicht, daß Gott Mann und Weib geschaffen habe, sondern denken, es sei eine solche Zusammenfügung, welche ohngefähr geschehen sei, und dafür unser HErr Gott keine Sorge trage und deren er nicht sonderlich achte.

57. Derhalben geschieht es nicht umsonst, daß man sonderliche Ceremonien und Ordnungen gebraucht in der Kirche, wo man die Ehen eingeführt und aufgerichtet; desgleichen wo man die Diener des Wortes ordinirt. Denn wir segnen Bräutigam und Braut, wünschen ihnen Glück, lesen die Worte von der Einsetzung dieses Standes, rufen Gott an, daß er darüber halten, und denselben schützen und bewahren wolle. Den Dienern aber des Wortes legen wir die Hände auf und thun zugleich unser Gebet zu Gott allein darum, daß wir damit bezeugen, daß es Gottes Ordnung sei, beide in diesen und auch in allen andern Aemtern der Kirche, weltlicher Polizei und Hausregimente.

58. Im Papstthume ist diese Lehre ganz unbekannt und gleichsam begraben und verlöschen gewesen. Denn da haben die Menschen erdacht, daß die Kirche nichts Anderes wäre, denn ein Haufen Päpste, Cardinäle, Bischöfe etc. Sie haben das Licht des göttlichen Wortes nicht gehabt, auch das Gebet oder rechte Anrufung nicht, haben auch zu dieser göttlichen Einsetzung solche Demuth und Ehrerbietung nicht gehabt.

59. Derhalben soll man dies Exempel Josephs fleißig merken, welcher nicht ohne Beruf, auch nicht ohne Gebet, sondern mit gebührlicher Ehrerbietung und Betrachtung des Wortes hinzu getreten ist, dem Könige Rath zu geben und ihm einen Traum zu deuten. Darum ist es auch alles wohl gerathen, da Gott seine Hülfe und Segen gegeben hat.

\begin{quote}
\textbf{V. 17—24. Pharao sagte an zu Joseph: Mir träumete, ich stund am Ufer bei dem Wasser; und sahe aus dem Wasser steigen sieben schöne fette Kühe, und gingen an der Weide im Grase. Und nach ihnen sahe ich andere sieben dürre, sehr häßliche und magere Kühe heraus steigen. Ich habe in ganz Egyptenland nicht so häßliche gesehen, und die sieben mageren und häßlichen Kühe fraßen auf die sieben ersten fetten Kühe. Und da sie die hinein gefressen hatten, merkte mans nicht an ihnen, daß sie die gefressen hatten, und waren häßlich, gleichwie vorhin. Da wachte ich auf. Und sahe abermal in meinem Traume sieben Aehren auf Einem Halm wachsen, voll und dicke. Darnach gingen auf sieben dürre Aehren, dünne und versengte. Und die sieben dünnen Aehren verschlangen die sieben dicken Aehren. Und ich habe es den Wahrsagern gesagt, aber die können’s mir nicht deuten.}
\end{quote}

\quellseite{350}

60. Dies ist eine Wiederholung des vorigen Traumes. Was aber die Grammatik belanget, will ich den Hebräern befohlen sein lassen, daß sie dasselbe auslegen oder erklären. Der Ostwind wird allenthalben so beschrieben, daß er dürr und heiß sei und Alles verzehre. Denn wo dieser Wind weht, wird das Erdreich hart und trocken, und er ist schädlicher als die Hitze der Sonne. Denn daher kommt der Rost oder Brand in die Frucht, verbrannt schwarz Korn, darum daß die Halme wegen der großen Dürre nicht wachsen können. Das hebräische Wort zenumoth will ich den Hebräern befehlen; denn es ist eins von den Wörtern in der hebräischen Sprache, welche gar unbekannt sind, und man muß desselbigen Wortes Auslegung aus den Umständen nehmen, dieweil es sonst nirgends mehr steht denn an diesem Orte. Die Andern haben es gegeben klein oder dünn; denselben sind wir auch nachgefolgt.

\begin{quote}
\textbf{V. 25. Joseph antwortete Pharao: Beide Träume Pharao sind einerlei. Denn GOtt verkündiget Pharao, was er vorhat.}
\end{quote}

61. Dieser Text gehört zu dem, das zuvor auch gesagt worden ist. Denn es ist kein zufälliger, willkürlicher Traum, sondern eine Eingebung und Offenbarung, die tiefen Eindruck machte; ja, es ist eine recht eigentliche Weissagung, dieweil Joseph bekennt, daß dieser Traum GOttes Wort und dazu prophetisch sei; wie dasselbe seine Worte ausweisen, da er sagt: „GOtt verkündiget Pharao, was er vorhat“ etc. Und dies ist die rechte Eigenschaft der Weissagung, welche nichts Anderes ist denn ein Wort von den zukünftigen Thaten oder Werken GOttes. Dieweil aber nun der Traum Pharaos also beschaffen ist, wird er billig und recht eine Weissagung GOttes genannt.

62. Und Joseph hat diese Worte nicht hervorgebracht aus eigener Einbildung oder Wahn, und aus der Lehre der Wahrsager und Priester, welche nur ungewisse Meinungen und falscher Wahn sind; sondern er folgt dem einigen und gewissen Verstande des göttlichen Wortes nach. Darum ist auch der Heilige Geist dabei gewesen, welchen er mit Betrachtung des Wortes und mit dem Gebet erlangt, wie Daniel Cap. 2, 18. 19., und er sieht, daß er ordentlich berufen wird zum Regiment im Königreich Egypten, dazu er sich nun auch ohne allen Zweifel willig begeben.

63. Es hat zwar die Vernunft auch etlichermaßen verstehen können, daß die sieben magern Kühe insgemein etwas Böses bedeuten mußten und wiederum durch die fetten Kühe etwas Gutes müßte verkündigt sein; welches denn Pharao selbst und ohne Zweifel die andern Weisen auch leicht haben errathen können. Aber was es nun im Einzelnen sein werde, davon haben sie nichts Gewisses schließen können. Desgleichen möchte Einer die Kühe auf das weibliche Geschlecht ziehen, als ob sieben Fräulein aus dem königlichen Stamme sterben und andere sieben überbleiben würden; oder daß man es auf sieben reiche und sieben arme und elende Geschlechter deutete.

64. Dies hat wohl ein Ansehen der Weisheit und scheint gerade kein ungereimtes Urtheil zu sein, aber es ist noch gar weit vom rechten prophetischen Verstande. Denn das hat niemand gewiß sagen können, daß die sieben Kühe sollten so viel Wohlfälle, gute, oder aber theure und böse Jahre bedeuten. Aber da die Deutung angezeigt worden ist mit ihrem Ausgang und Erfüllung, scheint sie so einfach und klar, daß man meint, es hätte sie ein Jeder also verstehen können. Und du solltest noch Viele finden, die sich wohl sollten träumen lassen und meinen, sie wollten dies auch so leichtlich und ohne alle Mühe verstanden haben. Wie jetzt die Rottengeister und Schwärmer sich selbst zuschreiben, was sie von uns gehört haben, und wollen drum daher hoch gelobt und gerühmt sein, daß sie anderer Leute Arbeit tadeln und verachten können.

65. Daran ist aber kein Zweifel, daß diese Auslegung auch wohl den Allerweisesten und Gelehrtesten im selbigen Königreich wird genug haben können zu schaffen geben; aber sie haben sich alle vergeblich daran versucht und gearbeitet. Denn die Träume der Könige und der Personen, so der Gemeinde vorstehen, sind Weissagungen, und wollen den Heiligen Geist haben, daß er sie auslege; den Andern aber sind sie gar verschlossen und verborgen. Wie die Weissagungen vom Reich Christi, daß er würde die ganze Welt bezwingen, im Jesaia und andern Propheten, wiewohl sie klar sind, doch nicht leichtlich von jedermann verstanden werden. Denn die Juden haben ihre Träume von Schwertern und leiblicher Gewalt, die er dazu brauchen werde. Und sind solche Prophezeigungen versiegelt und verborgen gewesen, bis Christus und die Apostel alle Schlösser aufgelöset und gelehrt haben, daß es ein geistlich Reich sein würde, welches uns arme Menschen vom ewigen Tode erlösen und mit himmlischen und ewigen Gütern reich machen würde. Den Juden aber ist noch Alles verschlossen und verborgen, was die Propheten beide von Christo und seiner Mutter zuvor verkündigt haben, wie Jes. 29, 11. steht: „Aller Propheten Gesichte werden euch sein, wie die Worte eines versiegelten Buchs.“

66. Derhalben hat Joseph einen prophetischen Geist; darum antwortet er auch Pharao fein richtig und sagt: „Beide Träume Pharao sind einerlei“, so es doch zweierlei Gesichte waren, von Kühen und von Aehren. Aber die Vernunft, wo sie ohne den Heiligen Geist gewesen, würde nothwendig geurtheilt haben, daß es nicht einerlei Träume wären; wiewohl es jetzt auch jedermann dünkt einerlei Traum zu sein, nachdem es von GOtt geoffenbart ist.

67. Wie Lyra auch ganz fein etliche Ursachen der Auslegung dichtet. Denn er sagt, daß durch die mageren und fetten Kühe ganz angenehme die theuren und fruchtbaren Jahre bedeutet werden, sintemal das Land durch die Arbeit der Kühe gebauet wird, und daß eben dasselbige auch durch die Aehren füglich bezeichnet werde, dieweil die Aehren zur Zeit der Ernte gesammelt werden, sie sei nun gleich klein oder groß. Also gibt Lyra Ursache dieser Deutung a posteriori, das ist, von dem, so hernach mit der That gefolgt ist, da dieselbe Deutung zuvor von GOtt ist geoffenbart worden. Wenn ich wüßte, daß durch eine jegliche Kuh ein jegliches Jahr, und durch das Fette ihre Fruchtbarkeit und wiederum durch das Magere Theurung sollte verstanden werden, so wollte ich das Andere alles auch leichtlich errathen und daraus schließen können.

68. Aber wenn nicht der prophetische Geist dabei wäre, so würden sie ohne Zweifel alle fehlen in der Deutung der Kühe. Denn Kühe bedeuten in der heiligen Schrift oftmals Völker, Städte, Obrigkeit; als im 68. Psalm V. 31.: „Die Rotte der Ochsen unter ihren Kälbern“, das ist, unter ihrem Volk. Da werden Ochsen genannt die Fürsten und Obrigkeit, Kühe und Kälber die Städte und Länder. Derhalben hätte es sich nicht übel gereimt, daß er sieben reiche Länder oder Städte durch die sieben fetten Kühe verstanden hätte, und wiederum durch die sieben magern Kühe sieben Länder oder Städte, die auch mager und arm wären. Diese Auslegung betrifft das weltliche Regiment. Dazu kann man auch eine theologische und geistliche Deutung thun, nämlich also, daß der Ochse den Lehrer oder Prediger in der Kirche, die Kühe aber das Volk oder Opfer bedeuten. Wer will aber darüber urtheilen, welche die rechte eigentliche und gewisse Meinung oder Verstand sei? so doch der prophetische Geist nicht soll zweideutig sein, und geziemet ihm nicht, daß er ungewiß rede, sondern soll eigentlich, deutlich und verständlich reden, also daß etwas Gewisses daraus möge geschlossen werden.

69. Derhalben würde dies niemand aus seinem eigenen Verstande oder Scharfsinn haben erfinden und begreifen können, daß das Gesicht von den sieben Kühen von der Haushaltung zu verstehen sein sollte. Das ist allein des Heiligen Geistes Werk gewesen, der Josephs Herz erleuchtet hat; wie er das selbst mit Ehrerbietung und Danksagung bekennt, da er sagt: „GOtt verkündiget Pharao, was er vorhat“, und nicht ich. Und droben, Cap. 40, 8., sagt er auch: „Auslegen gehöret GOtt zu.“ Als wollte er sagen: Dieser Traum ist wahrhaftig eine göttliche Weissagung, darin GOtt lehrt und berichtet, was für gute und böse Zeiten bevorstehen, und zeigt damit an, daß er für sich Sorge trage. Dies ist ein göttlich Gesicht und nicht menschlich. Darum ist die Deutung auch gewiß und wahrhaftig, nämlich, die GOtt selbst auch geoffenbart und gezeigt hat.

70. Dies soll aber ferner auch noch zu der Lehre gezogen werden, davon wir droben gehandelt haben, nämlich, wenn GOtt einem Lande günstig ist, so gibt er dem Volk fromme Fürsten und regiert dieselbigen, wie die Worte Spr. 8, 15. lauten: „Durch mich regieren die Könige, und die Fürsten oder Rathsherren setzen das Recht.“ Welches nicht geschehen kann denn allein durch die göttliche Weisheit, welche die Fürsten regiert den Unterthanen zu Nutz und Wohlfahrt, nämlich, wenn sie mit den Königen und Fürsten redet, und nicht allein redet, sondern auch ihre Rede deutet. Solche Regimente sind wahrlich die besten und steht damit sehr wohl. Denn die Fürsten, wie Jes. 32, 8. steht, „werden fürstliche Gedanken haben, und drüber halten.“ Da muß man aber solche Regenten
\quellseite{351}
oder Räthe haben, wie hier Josephs einer gewesen ist, welche von GOtt gegeben worden, der selbst den Fürsten muß Weisheit und Rath eingeben, entweder in eigener Person, oder aber durch einen Naeman, den Syrier, oder einen Nathan, wie David gehabt, welche den Rath GOttes und sein Wort und Weissagung haben. Das ist des Landes Heil und Wohlfahrt und der rechte Segen, da unser HErr GOtt nicht allein Polizei und Regiment einsetzt und verordnet, sondern auch dasselbe mit sonderlicher Güte und Gnade regieret.

71. Darum ist das eine treffliche Predigt, die man mit Fleiß merken und behalten soll, welche Joseph vor dem Könige thut: GOtt, sagt er, redet mit dir, lieber König, und verkündigt dir sehr große Dinge, die er thun wird in deinem Königreich. Ich bin der Mann nicht, der dir solches erklären wird, sondern GOtt selbst offenbart dir es im Traum durch ein Gesicht. Nun will dir gebühren, daß du diese Wohlfahrt erkennest.

72. Derhalben ist unter diesem Volk Weissagung gewesen und Joseph ist Egyptenlandes Heiland gewesen. Wenn es aber daran mangelt, daß die Weissagung nicht im Lande ist, alsdann werden die Regimente verstört, wie Salomo sagt Spr. 29, 18.: „Wenn die Weissagung aus ist, wird das Volk wild und wüst.“ Wenn aber die Weissagung da ist, das ist, GOttes Rath und Wort, und der Fürst denselbigen gehorcht, ein solch Land oder Regiment magst du wahrlich wohl selig preisen und sagen, daß es um dasselbe sehr wohl stehe. Und ein solch Regiment ist dazumal im Königreich Egypten gewesen, welches mit großen Lasten, Sünden oder Schanden nicht befleckt oder verunreinigt war, um derentwillen GOtt die Herrschaften und Regimente zu strafen und durch die Tyrannen zu verderben pflegt; sondern es ist ein solch Regiment in Egypten gewesen, welches GOtt angenehm und wohlgefällig war; sintemal er dasselbe Land gegen die künftige Theurung so gnädig versorgt.

73. Darum sagt Joseph zum Könige: GOtt hat dich lieb, dieweil er dir sein Wort und Willen offenbart. Und also lobt und zieret er auch die Polizei oder das weltliche Regiment. Denn er will mit Demuth und GOttesfurcht darin dienen und rathen, und ermahnt auch den König selbst, daß er bedenken soll, er sei in einem göttlichen Amt, und daß ihm darin vonnöthen sei, daß er fleißig bete und GOtt um Hülfe anrufe. Denn was wäre ihm mit aller Weisen und Priester Rath geholfen gewesen, dieweil sie selbst von der zukünftigen Theurung nichts gewußt? Ja, es hätte der König selbst und das ganze Königreich zugleich verderben müssen. Derhalben, lieber König, bedenke, daß du dich demüthigest, und GOtt erkennen lernest, der das Wort und auch die Auslegung gibt. Denn wo du dieselbige annehmen wirst, so wird dein Königreich glückselig und löblich sein bei GOtt und den Menschen.

74. Darnach ist an diesem Ort das auch zu merken, daß Joseph den Egyptern das Gesetz der Beschneidung oder andere Satzungen der Väter nicht auferlegt hat; denn er ist ein Lehrer in Egypten gewesen und auch der umliegenden Völker, wie der 105. Psalm V. 22. zeugt: sondern er hat ihnen nichts Anderes vorgehalten, denn die Erkenntniß des rechten wahren GOttes, Schöpfers Himmels und der Erden. Er hat ihnen nicht geboten, daß sie sich sollten beschneiden lassen, oder die Satzungen der Väter halten, sondern hat sie einfach gelehrt die rechte Lehre von dem GOtt, welchem er selbst auch gedient. Er ist ein Bischof über ganz Egypten gewesen, hat den König unterwiesen, dazu die Fürsten und Priester, und das ganze Volk. Darum ist es zu der Zeit im selbigen Königreich sehr wohl gestanden nicht allein wegen der Fruchtbarkeit des Landes und anderer leiblicher Güter, sondern auch daher, daß sie das Licht des göttlichen Worts gehabt haben. Und wird ohne Zweifel der König selbst auch mit Ernst über guter Zucht und Ordnung gehalten haben; wie das an dem Exempel des Amtmannes über die Väter zu sehen ist. Darum ehrt und begabt ihn auch GOtt mit seinem Wort und Weissagung von dem, was er vorgehabt, wie droben der Text sagt.

75. Solche Exempel soll man der Juden und Sabbather Unsinnigkeit und Irrthum entgegen setzen, welche neulich in Oestreich aufgestanden und sich vorgenommen haben, die Leute zur Beschneidung zu zwingen, als ob sie nicht könnten selig werden, sie müßten sich denn beschneiden lassen. Es haben droben auch Abimelech und Pharao, da Abraham noch am Leben gewesen und gelehrt hat, die rechte Lehre von GOtt angenommen ohne Beschneidung. Desgleichen der Minister und Babylonier Könige sind auch nicht beschnitten worden; sondern haben sich daran genügen lassen, daß sie das Wort und die Erkenntniß des wahren GOttes gehabt, ohne die äußerlichen Ceremonien des Volks Israel und ohne Beschneidung.

\begin{quote}
\textbf{V. 26—31. Die sieben schönen Kühe sind sieben Jahr, und die sieben guten Aehren sind auch die sieben Jahr. Es ist einerlei Traum. Die sieben mageren und häßlichen Kühe, die nach jenen aufgestiegen sind, das sind sieben Jahr; und die sieben mageren und verbrannten Aehren sind sieben Jahr theure Zeit. Das ist nun, daß ich gesagt habe zu Pharao, daß GOtt Pharao zeiget, was er vorhat. Siehe, sieben reiche Jahr werden kommen in ganz Egyptenland. Und nach den sieben werden sieben Jahr theure Zeit kommen, daß man vergessen wird aller solcher Fülle in Egyptenland; und die theure Zeit wird das Land verzehren, daß man nichts wissen wird von der Fülle im Lande, vor der theuren Zeit, die hernach kommt; denn sie wird fast schwer sein.}
\end{quote}

76. Die Siebenzahl bedeutet siebenmal sieben Jahre, und haben beiderlei Zeichen einerlei Bedeutung, als wenn zwei Zeugen zu einerlei Sache gezogen werden. Es ist in beiden Geschichten keine Ungleichheit. Übrigens wird bei der Deutung der sieben Aehren und sieben schönen fetten Kühe die Fruchtbarkeit nicht ausdrücklich erwähnt, was bemerkenswerth ist. Bei dem Zeichen aber der magern Kühe wird gesetzt die Bedeutung der theuren Zeit, und aus dem Gegensatz können leichtlich die sieben guten fruchtbaren Jahre verstanden werden.

77. Er wiederholt aber, daß GOtt dem Lande Egypten beides theure Zeit und fruchtbare Jahre geben werde. Das, sagt er, steht nahe bevor; darum zeigt dir es GOtt durch ein Gesicht, und mit diesem wunderbaren Werke ehrt und ziert er dich vor andern Königen der Assyrier und Babylonier. Darum danke unserm HErrn GOtt, er meint es gut mit dir.

78. Saba und scheba sind zwei Worte wegen des Unterschiedes der Punkte, und heißen so viel als: sieben, Fülle oder Sättigung, und, Eid. Daher ist droben im 21. Cap. V. 31. der Ort Bersaba also genannt um dreierlei Ursachen willen: Zum Ersten um der Sättigung willen, darnach wegen des Eides, zum Dritten um der Lämmer willen. Es werden übrigens diese Worte mit einerlei Buchstaben geschrieben ohne Punkte; welche Gleichheit der Worte einen Propheten oder Ausleger auch wohl hätte betrügen können, wo nicht der Heilige Geist dabei gewesen wäre. Aber diese Auslegung und Prophezeiung ist gewiß, da der Heilige Geist durch Joseph mit Pharao redet.

\begin{quote}
\textbf{V. 32. Daß aber dem Pharao zum andernmal geträumet hat, bedeutet, daß solches GOtt gewißlich und eilend thun wird.}
\end{quote}

79. Er gibt nun die Ursache an, warum der Traum zweierlei sei und doch nur einerlei Ding bedeute. Denn was GOtt redet, das ist gewiß und wird auch eilends seinen Ausgang haben. Darum sind es gleichsam zwei Zeugen in einer Sache, und will Joseph fürzlich zum Könige also sagen: Es soll gewiß sein und bald kommen, GOtt wird dich nicht in Zweifel oder ungewissen Wahne stecken lassen.

80. Das hebräische Wort, nachon, heißt, gewiß und beständig. Als, im 51. Psalm V. 12.: „Gib mir einen neuen gewissen Geist“; und im 5. Psalm V. 10.: „In ihrem Munde ist nichts Gewisses“, das ist, sie haben keine gewisse Lehre, sondern nur ungewissen Wahn, Zweifel, Aberglauben, ungewisse Meinung: es sind battologi, das ist, unnütze Wäscher, sie lehren aber nichts Gewisses oder Beständiges, darauf sich das Gewissen verlassen möchte.

81. Derhalben ist das ein rechter gewisser Geist, der uns gewiß macht des Willens GOttes, läßt uns nicht zweifeln, sondern richtet das aus, davon St. Paulus Röm. 14, 5. ermahnt, daß ein jeglicher seiner Meinung gewiß sei; item 1 Thess. 1, 5.: „Unser Evangelium ist bei euch gewesen, nicht allein im Wort, sondern beide in der Kraft und in dem Heiligen Geiste.“ Welche Bergversicherung auch in allerlei Lehre vonnöthen ist und vornehmlich in christlicher Lehre. Denn ich soll des gewiß sein, was ich von GOtt halten soll, oder vielmehr, was er von mir halte.

82. Und ist es ein greulicher Irrthum gewesen in der papistischen Lehre, damit sie die Leute so weit gebracht haben, daß sie an Vergebung der Sünde und GOttes Gnade zweifeln. Du sollst, haben sie gesagt, erkennen, daß du ein Sünder bist, und ein solcher Sünder, der seiner Seligkeit nicht kann gewiß sein. Also ist die ganze Welt ersoffen gewesen in solchem Zweifel und andern irrigen Meinungen von GOtt.

\quellseite{352}

83. Darum soll man lernen, daß GOtt kein ungewisser, zweifelhaftiger oder wandelbarer, zweideutiger GOtt sei, wie ein ungewisses Rohr sei; sondern daß er gewiß und zuverlässig sei: der da sagt: Ich taufe dich im Namen des Vaters, und des Sohnes, und des Heiligen Geistes; ich absolviere und spreche dich los von deinen Sünden etc. Da irren GOtt der Vater, der Sohn und der Heilige Geist nicht, werden nicht vom Winde hin und her getrieben, sondern sind wie ein harter Fels und Sela; wie GOtt in den Psalmen oft genannt wird, daß er ganz gewiß sei, auf welchen du dich gewißlich magst verlassen, und sagen: Ich bin heilig und selig, bin ein Kind und Erbe GOttes; denn ich bin ja getauft.

84. Und nun soll die Zweifel der Mönche weit hinten setzen, soll nicht also sagen: Siehe, ich habe gethan, was ich habe thun sollen; ob es aber nun GOtt gefalle oder nicht, weiß ich nicht. Man soll nicht auf das Ungewisse laufen, oder als die in die Luft streichen, wie Paulus sagt 1 Cor. 9, 26. Unser Gang soll gewiß und beständig sein, daß wir mit gewisser Zuversicht sagen: Ich schlafe im Namen des HErrn, und weiß, daß auch mein Schlaf GOtt wohlgefällt. Wenn ich aber wache, und meine gewöhnliche Arbeit thue in meinem Berufe, mit Schreiben, Lesen, Betrachten und mit Beten, zweifle ich daran nicht, daß solche Arbeit GOtt auch angenehm sei; und wenn ich wüßte, daß es ihm mißfällig wäre, wollte ich mich dessen viel lieber enthalten. Ich bin aber deß gewiß, daß ich GOtt wohlgefalle mit alle meinem Thun, nicht um meinetwillen, der ich solches thue, sondern um GOttes willen, der sich meiner erbarmt, mir die Sünde vergibt, mich liebt, führt und mit dem Heiligen Geist regiert.

85. Diese Gewißheit und Plerophorie (Freudigkeit) soll man behalten. Denn sonst ist die Taufe, die Absolution und auch der Brauch des Abendmahls des HErrn unnütz und vergeblich. Wie es denn uns im Papsthume also ergangen ist, wie Paulus sagt 2 Tim. 3, 7.: „Sie lernen immerdar, und können nimmer zur Erkenntniß der Wahrheit kommen.“ Denn es eine unchristliche Blindheit und Irrthum gewesen, der durchaus zu verdammen ist, wenn auch sonst nichts mehr unrecht oder sündlich gewesen wäre in des Papstes Lehre, nämlich, daß sie gelehrt haben, wir sollen immer hin und her im Zweifel gehen, wanken, ungewiß sein und an unserer Seligkeit zweifeln. Denn solche Ungewißheit oder Zweifel nimmt mir meine Taufe und GOttes Gnade. Ich bin vergeblich ein Christ, arbeite und lebe vergeblich.

86. Darum betet der Prophet im 51. Psalm V. 12.: „Gib mir einen neuen gewissen Geist“, das ist, gib mir einen rechten gewissen Glauben, der nicht hinke auf beiden Seiten; wie die Baaliten thaten, die in ihrem Gottesdienste zweideutig und unbeständig waren: sie arbeiteten, opferten, thaten ihrem Leibe wehe, freuzigten ihr Fleisch, und waren doch ungewiß, ob sie auch GOtt gefällig wären.

87. Derhalben lernt vor allen Dingen und gewöhnt euer Herz daran, daß euer ganzes Leben gewiß und beständig sei, nach dem Spruche Hebr. 12, 12.: „Richtet wieder auf die lässigen Hände und die müden Kniee, und thut gewissen Tritt mit euren Füßen“ etc.; und 2 Petr. 1, 10.: „Machet euren Beruf fest“ etc. Bleibet in eurem Berufe, und schließet bei euch gewißlich also: Ich bin getauft, ich habe den Leib und Blut Christi empfangen, bin absolvirt, GOtt wird mich nicht betrügen. Warum sollte ich denn mich selbst betrügen und zweifeln, dieweil ich weiß, daß GOttes Werke gut und gewiß sind? Was ich darnach thun werde, davon kann ich sagen, daß es ein gut Werk sei und ohne Zweifel GOtt wohlgefalle darum, daß ja der Glaube nicht zweifelt. Wiederum aber, alles, „was nicht aus dem Glauben gehet, das ist Sünde“, Röm. 14, 23. Wo du zweifelst, wirst du immer ohne Unterlaß sündigen. Denn du glaubst nicht, daß du GOtt wohl gefällest, und wäre viel besser, daß du nur aufhörtest und mit allen Werken stille hieltest, bis du gewiß würdest, denn daß du also auf das Ungewisse laufen und arbeiten solltest. Denn wir müssen rechte neben thun, wie das hebräische Wort lautet, das ist, gewiß, beständig und fest sein.

88. Das hat nun Joseph gemeint, da er zu Pharao sagt: GOtt hat dir zweierlei Träume gegeben, nicht daß du solltest ungewiß sein, oder denken, daß durch die Hecken ein Anderes bedeutet würde als durch die Rinder; sondern darum hat es dir zum andermal geträumt, daß du solltest wissen, daß es bald seinen Ausgang gewinnen werde, was dir zuvor ist verkündigt worden. Darum hat er keine Ursache zu wanken oder zu zweifeln. Denn dieweil der König GOtt wohlgefällt, eben darum gibt er ihm eine gewisse beständige Auslegung, welche ihm und dem ganzen Königreiche nützlich und heilsam ist. Wie unser HErr GOtt um der Sünde willen des Volkes pfleget gottlose, ungeschickte Könige und Regenten zu geben; aber wenn die Unterthanen gläubig und gottselig sind, gibt er ihnen fromme und heilsame Fürsten. Er hat in der Policei des Volkes Israel immer fromme Propheten erweckt, die das Volk mit dem Worte regiert haben außerhalb der königlichen Gewalt; darum wurden sie auch von den Königen getödtet.

89. Die Regel soll man übrigens fleißig merken, daß es in der heiligen Schrift nicht vergeblich sei, wo ein Ding mit einerlei Worten zwei oder mehrmal geredet oder wiederholt wird; sondern daß es darum geschehe, daß es eine Stärkung unsers Glaubens sein soll. Wie Jesaia schier allezeit einerlei Ding auf zweierlei Weise auszudrücken pflegt; nicht daß er derhalben wollte gerühmt werden, daß er so beredt wäre, oder so viel Worte machen könnte, sondern uns damit in unserm Glauben zu stärken, auf daß wir dessen ja gewiß seien, das uns GOtt verheißt und mit uns redet.

90. Es wird aber ohne Zweifel Joseph aus dieser Weissagung viele herrliche Predigten gefaßt und genommen haben. Denn er hat sie fleißig und mancherlei Weise gehandelt mit den vornehmsten Herren und Priestern desselbigen Königreichs unterredet, und also das Wort und rechten Gottesdienst in Egyptenland gepflanzet hat.

91. Nun gibt er aber auch einen Rath zu der Auslegung. Denn es hat der Heilige Geist Lust dazu, daß er durch ihn mit dem Könige reden mag, und gibt dem Könige mehr, den er hat hoffen oder begehren können. Er hatte allein gebeten den Traum zu deuten; solche Deutung wird ihm reichlich gegeben, und über das gibt ihm auch nun Joseph Rath, wie er des Segens der sieben fruchtbaren Jahre gebrauchen solle, nämlich, der Armuth und theuren Zeit, so in nachfolgenden Jahren kommen würde, damit zu steuern und zu rathen.

\section[Dritter Theil]{Dritter Theil. Von dem Rath, den Joseph Pharao gegeben, und was dieser Rath sowohl, als die Auslegung der Träume bei Pharao und seinen Knechten gefruchtet}

\subsection{I}

\begin{quote}
\textbf{V. 33–36. Nun sehe Pharao nach einem verständigen und weisen Mann, den er über Egyptenland setze; und schaffe, daß Amtleute verordne im Lande, und nehme den Fünften in Egyptenland, in den sieben reichen Jahren; und sammle alle Speise der guten Jahre, die kommen werden, daß sie Getreide aufschütten in Pharao Kornhäuser, zum Vorrath in den Städten, und bewahren es; auf daß man Speise verordnet finde dem Lande in den sieben theuren Jahren, die über Egyptenland kommen werden, daß nicht das Land vor Hunger verderbe.}
\end{quote}

92. Chokmah, das hebräische Wort, heißt Weisheit oder Gericht und Geschicklichkeit, oder ein fleißiges Aufmerken, welches zu der Kunst der Dialektik gehört, daß man darauf fleißig Acht geben, und sich hüten kann vor den Stricken, vor List und Sophisterei, so erdacht und vorgewendet werden, die Weisheit damit zu verkehren und zu unterdrücken. Einen solchen verständigen und weisen Mann, das will nun Joseph haben, soll der König wählen, der vorsichtig sei und scharfen Verstand habe, welcher nicht leichtlich betrogen oder übervortheilt werden könne. Er will haben, daß ersichtlich ein verständiger, scharfsinniger, geschickter Mann über das ganze Land soll gesetzt werden und daß darnach auch andere Amtleute verordnet werden. Dieser Rath würde Pharao oder irgend seiner Nähe einem nimmermehr in Sinn gekommen sein. Und Joseph hat weislich bedacht, daß vonnöthen sein würde, daß man vielen andern Amtleuten, so über das Getreide gesetzt, einen Obersten, Vizhof oder Aufseher zugeben sollte, nach dem sich die andern richten müßten.

93. Das hebräische Wort, pekidim, haben wir im Deutschen verdolmetscht, Amtleute; wie es denn auch schier also im 109. Psalm V. 8. gebraucht wird: „Sein Amt müsse ein anderer empfahen.“ Aber im 111. Psalm V. 7. wird es etwas anders gebraucht, da also stehet: „Alle seine Gebote sind rechtschaffen.“ Das will ich aber den Hebräern befohlen sein lassen.

\quellseite{353}

94. Ich halte aber dafür, daß sie selbst, die Hebräer, die vollkommenste Kenntniß ihrer Sprache noch nicht ganz wiedererlangt haben. Denn daß sie ein und denselben Wort in vielerlei Bedeutung geben, zeigt genugsam an, daß sie zum guten Theil ihre eigene Sprache nicht kennen, was Ausdrücke und Redefiguren anlangt. Denn das kommt oftmals vor, daß je nach der Verbindung, in welcher ein Wort steht, dessen Bedeutung verschieden ist. Wie in der deutschen Sprache das Wort „Recht“ sich sehr weit ausbreitet, und oftmals von seiner rechten natürlichen Bedeutung abgeht. Als, wenn ich sage: Das ist der rechte Schalk; da ist seine Bedeutung gar anders als die vorige. Oder wenn ich sage: Der Acker wäre mir recht. Das verstehen wir zwar leichtlich und stoßen uns nicht daran, daß dies Wort mehr denn einerlei Bedeutung hat: wenn aber in der hebräischen Sprache auch eine solche Veränderung in der Bedeutung der Worte vorkommt, die uns aber nicht so gemein oder gebräuchlich ist, da können wir die wahre Bedeutung der Ausdrücke nicht erreichen oder verstehen.

95. An diesem Orte aber bedeutet das Wort pekidim die Amtleute oder Superintendenten, die auf das Getreide in Egypten haben sehen und dafür Sorge tragen sollen, Kornmeister oder Kornschreiber. Denn droben haben wir gehört, daß solche Bischöfe oder Obersten gesetzt gewesen sind über die Weger, Scheunen und Bäcker. Nun gehört aber vornehmlich großer Fleiß dazu, daß auf das Getreide Acht gegeben und dasselbe wohl versorget werden möge. Dieselbigen Exempel hat nun Joseph in Egypten gesehen und ist ihnen auch nachgefolgt. Denn solche Ordnungen gehören zum politischen und äußerlichen Leben. Und ist wahrlich dies ein göttlicher Rath und Weissagung des Heiligen Geistes, daß er sieht, daß solche Verwaltungen durch die Amtleute, Aufseher und Bischöfe vonnöthen sein wollen, dieweil kein Regiment der Könige oder Fürsten verwaltet werden kann ohne durch solche Visitatoren, Aufseher und Bischöfe, die aufsehen, wacker sind und hüten.

96. Wir haben auch Rentmeister oder Steuerbeamte, Bürgermeister und andere Amtleute, welche auf Befehl des Fürsten auf Zins und Recht sehen und sonst die andern Ämter verwalten. Aber dieselbigen Amtsverwalter sind einander sehr ungleich. Denn etliche sind fromm, etliche böse; sie sind wohl hart und streng genug auf alle Dinge zu sehen, aber nicht dem gemeinen Nutzen und Fürsten zum Besten, sondern daß sie nur damit ihren eigenen Nutzen stiften und bald reich werden mögen; sie thun wie die Wölfe, Füchse, Geier und Weihen.

97. Joseph aber will solche Amtleute oder Aufseher haben, die den fünften Theil vom Lande einnehmen sollen, welches die Hebräer mit einem Worte ausdrücken, chimesch; als wenn man im Lateinischen sagte: quintare: von allem Getreide, sagt er, das in den fruchtbaren Jahren wachsen wird, soll man den fünften Theil behalten auf die künftige theuere Zeit. Wie es aber damit zugegangen sei, daß sie also den fünften Theil eingenommen, das wird von Mose nicht erzählt, und ist auch nicht wahrscheinlich, daß Joseph diesen Rath sollte gegeben haben, daß Pharao dem Volk das Getreide mit Gewalt nehmen sollte; sondern ich halte dafür, es sei den Amtleuten befohlen gewesen, daß sie um des Königs Geld und auf seine Kosten den fünften Theil jährlicher Frucht oder Einkommens im ganzen Königreich kaufen sollten. Und in den fruchtbaren und reichen Jahren haben sich die Unterthanen an vier Theilen können genügen lassen und damit zufrieden sein, also daß der fünfte Theil in die gemeinen Scheunen gesammlet würde. Es hat der Heilige Geist durch Joseph diese Weise also angezeigt, daß in solchen reichen fruchtbaren Jahren zu täglicher Nothdurft an vier Theilen würde genug sein, und daß man den andern übrigen Theil sollte aufheben, um damit dem künftigen Hunger und Armuth der Unterthanen abzuhelfen.

98. Denn dies ist eine politische und nöthige Lehre, so hier den Fürsten gegeben wird, welchen für das Volk zu sorgen befohlen ist und daher die Unterthanen mit Leibesnothdurft, als, Fleisch, Getreide, Wein, versehen sollen, vornehmlich wenn das Getreide theuer und nicht wohl zu bekommen ist; das ist der Könige und Fürsten Amt, desselben hat Joseph hier den Pharao erinnert. Denn sie sollen Väter des Vaterlandes sein und Hirten, keine Bären oder Wölfe. Und dies ist des Heiligen Geistes Ordnung, der den Fürsten die Sorge für das Getreide und andere tägliche Nothdurft der Unterthanen befohlen hat. Und ist hierbei des Exempels des Durchlauchtigsten Herrn, Herzog Friedrichs von Sachsen, wohl zu gedenken, welcher nicht allein Scheunen und gemeine Kornhäuser, sondern auch etliche Gruben im offenen Felde dazu gemacht, und dieselben mit Getreide und die Keller mit Wein hat füllen lassen. Da er aber von Staupitz und den Räthen derhalben gestraft worden, hat er geantwortet: Er thäte solches nicht um Geizes oder Gewinnstes willen, sondern wegen der Faulheit beider der Bürger und Bauern, die gar nicht an künftige Theurung dächten, sondern in den Tag hinein leben, nur schlecht dahin, von der Hand in den Mund, wie man zu sagen pfleget. Er aber ließe darum das Getreide sammeln und verwahren, auf daß das Volk in der Theurung seine Nothdurft haben möchte, sein Leben zu erhalten. Und es ist wahrlich ein sehr kluger Rath gewesen, dazu dem ganzen Lande nützlich und heilsam. Denn er hat mit seiner Vorsicht das verhütet, daß bei seinem Leben und die Zeit seiner Regierung keine drückende Theurung im Lande geworden ist.

99. Das sind Werke, deren sich die Fürsten in ihrem Amt billig befleißigen sollen; und dieselben sind auch sonderlich nöthig zu diesen Zeiten. Denn wir Deutschen sind rechte Biersäuche, proci Penelopes, fruges consumere nati, wie der Poet sagt, das ist: Wir sind gute volle Brüder, schlämmen und dämmen, leben immer im Sause und verschwenden Alles. Wenn darnach eine theure Zeit kommt, so haben wir gar nichts mehr übrig, damit wir uns des Hungers erwehren und Andern in ihrer Noth auch etwas helfen möchten. Derhalben gebühret es den Fürsten, daß sie den Armen wohl vorstehen und vornehmlich der Jugend, damit sie nicht Hungers sterben. Wie Joseph dem Könige Rath gibt, nicht daß er den fünften Theil mit Gewalt hinnehme, sondern denselben um sein Geld aufkaufen, sammeln und verwahren soll um des zukünftigen Misswachses und Theurung willen, daß also Egyptenland und andere umliegende Länder erhalten und genährt werden möchten.

100. Wir, die wir so klug und vorsichtig nicht sind, bedenken auch nicht, wenn das Getreide wohl gerathen und wohlfeile Jahre sind, daß gefährliche Theurung hernach kommen möchte, müssen oftmals erfahren, daß wir plötzlich und unversehens mit Theurung geplagt werden. Denn wir haben nicht solche Scheuern und Kornhäuser; daher kommt es auch, daß bald das Getreide aufsteigt und theuer wird, wenn auch schon das Einkommen oder jährliche Gewächs ziemlich gut ist. Also ist die Vorsicht, welche uns in diesem Exempel der Heilige Geist lehrt, kein Geiz, sondern eine Weisheit und Gottseligkeit, so frommen Fürsten wohl ansteht und ihnen auch ziemt.

101. Das hebräische Wort, bar, heißt, Weizen, und die Amtleute haben befohlen, daß der fünfte Theil Getreides in Pharaos Gewalt, das ist, in seine Scheunen zusammen sollte geführt werden. Vier Theile haben sie dem Volk bleiben lassen zur Speise und Saat. Ohne Zweifel werden auch etliche Bauern und Bürger in ihrer Haushaltung gleiche Vorsicht gebraucht haben. Denn das Einkommen und Gewächs in den sieben guten fruchtbaren Jahren ist unermäßlich groß gewesen, welches in den folgenden sieben Jahren alles verzehrt worden ist; und sind in allen Städten solche Amtleute gewesen, die also den fünften Theil eingenommen haben von den Ackern, so allernächst umher gelegen. Denn das Königreich ist sehr schön, wohl bestellt und angeordnet gewesen. Der oberste Bischof aber war Joseph, welcher durch ganz Egyptenland die Kornhäuser visitirt und besehen, und von den Amtleuten Rechnung gefordert hat.

102. Es wird aber hernach im 47. Capitel V. 24. wiederum gedacht werden des fünften Theils. Wir wollen an diesem Ort bei der Meinung bleiben und den Verstand behalten, daß es kein ordentlicher Zins, sondern nur die sieben Jahre über auf Josephs Rath also verordnet gewesen sei. Hernach aber wird es eine gemeine Ordnung werden und ist diese sieben fruchtbaren Jahre über damit sehr erträglich gewesen.

\quellseite{354}

103. Zu unsern Zeiten sind die Steuern in Deutschland allzuh hoch gestiegen. Bis daher ist es ja etwas leidlich und mäßig gewesen, da man von zwanzig Gulden einen gegeben hat nach den kaiserlichen Rechten, darin von hundert Gulden fünf gefordert werden, das ist nicht der fünfte, sondern der zwanzigste Theil. Die Juden hatten den Zehnten, haben vom zwanzig drei gegeben, beide vom Getreide und sonst von aller ihrer Habe und Gut. Darum sind wir Deutschen bis daher glücklicher gewesen als sie. Danach ist Egypten härter beschweret worden als die Juden; wiewohl es noch erträglich genug gewesen, wenn Egypten mit keiner andern Steuer ist beschweret worden, denn allein, daß sie den fünften Theil gegeben haben. Wir müssen zu unsern Zeiten, da Wucher, Steuern und Geiz so sehr gestiegen sind, nicht allein den fünften, sondern den vierten und schier den halben Theil geben. Die Wucherer zu Leipzig und der Adel saugen uns gar aus, nehmen zwanzig Gulden von hundert, welches der fünfte Theil ist. Die Obrigkeit nimmt ja noch so viel nicht, aber die Privatbuben, so nicht im Amt sind, dieselben nehmen mit ihrem Wucher schier mit Gewalt den fünften Theil. Dazu kommen denn hernach auch noch andere Steuern und gewöhnlicher Zins. Die mir nun zuvor den zwanzigsten Theil gegeben haben, müssen nun gewissermaßen den dritten und halben Theil geben, werden also gar ausgesogen. Die von Adel und die Kaufleute, und andere Tyrannen mehr machen, daß wir schier keinen Heller behalten.

104. Der Juden Zins ist auch über den Zehnten gelaufen; denn sie haben den Priestern den Zehnten gegeben und den Königen den achten oder neunten Theil; darnach ist noch von jeglichem Haupt auch ein sonderlicher bestimmter Zins gewesen, welchen sie haben geben müssen. Endlich sind sie durch Geiz und andere Räuberei und Schinderei auch gar ausgesogen worden; wie denn zu unsern Zeiten bei den Deutschen solcher Räuberei und Schinderei viel aufgekommen ist, und sind wahrlich die Unterthanen an ihren Gütern und auch die Regimenter an ihrem Vorrath jämmerlich beschnitten und geschwächt worden.

105. Nun ist gleichwohl dies noch keine gemeine Ordnung in Egypten gewesen, den fünften Theil zu geben, sondern das Getreide ist also gesammelt worden auf des Königs Kosten. Aber dies ist nun auch ein kluger und nöthiger Rath gewesen, daß er hat haben wollen, daß solche Aufseher sollten verordnet werden. Denn siehe doch unser gemein Stadtregiment an, da mangelt es an solchen Amtleuten über Fleischer, Bäcker und Schenken, da ist niemand, der Ordnung mache, wie theuer man eine jegliche Waare geben soll. Darum fahren die Bauern, Adel und Wucherer zu, berauben und saugen uns nach ihrem Gefallen gar aus.

106. Darnach will Joseph haben, daß diese Sammlung des Getreides dem ganzen Lande die sieben theuren Jahre sein soll zu pik(k)adon, wie es im Hebräischen lautet, das ist, zu einer Beilage oder Schatz desselben. Wiewohl es nicht so eigentlich ausgelegt werden kann, daß es eine Beilage heiße, wie es 3 Mos. 6, 2. 4. gebraucht wird und an andern Orten mehr. Ich wollte es gern also geben, daß es hieße repositum oder sepositum, das ist, bei sich hingelegt, nämlich, daß solch Getreide nicht habe sollen gebraucht werden, weder irgendwem noch kein jemand insonderheit, sondern hat sollen beigelegt und verwahrt werden auf die zukünftige Noth; und also, daß es zugleich mit sich einschließe und verböte das Aufsehen und Amt des Pfaffthums oder der obersten Verwaltung.

107. Hieronymus hat dies ganze Capitel mehr mit weitläuftigen Worten beschrieben, denn daß er es sollte ausgelegt haben. Wir haben dem griechischen Wort ἀποθησαυεῖν gefolgt, und haben es in Deutschen also gemacht: „Auf daß man Speise verordnet finde beim Lande“, das ist, daß es verwahret und verordnet sein möge auf die künftige Theurung. Das Andere will ich den Hebräern und Juristen lassen befohlen sein, nämlich, von dem fünften Theil. Sie werden aber in ihren Büchern das Wuchern und Aufsaugen, so heutiges Tages die Wucherer und Fälscher treiben, nicht finden. Die alten Rechte sagen: Man solle von zwanzig Gulden Einen geben; jetzt aber wollen sie zwei oder drei haben, und werden täglich neue Künste und andere Schindereien erdacht.

\subsection{II}

\begin{quote}
\textbf{V. 37. 38. Die Rede gefiel Pharao und allen seinen Knechten wohl. Und Pharao sprach zu seinen Knechten: Wie könnten wir einen solchen Mann finden, in dem der Geist Gottes sei?}
\end{quote}

108. Das ist auch eine Gabe GOttes, daß die Rede dieses Propheten Joseph dem Könige und allen Fürsten, Räthen, Amtleuten und Unterthanen wohlgefallen hat; denn diese alle nennt er des Königes Knechte. Darum ist es nicht allein eine sonderliche Gnade GOttes, Andere lehren können, sondern daß man auch dem Lehrer zuhöre und folge, dieweil hier der Teufel ausgeschlossen wird, der da mächtig ist mit Lästeren und die Zuhörer zu hindern, und ihre Herzen von dem Wort, das sie gehört haben, abzuführen, vornehmlich wo Gottes Wort und die Weissagung, nicht allein die geistliche, sondern auch in der Kirche und weltlichem Regiment, erst aufgeht. Darum ist es eine sehr hohe Gabe, daß der König mit seinen Räthen dieses göttliche Wort so bald annimmt und sich dasselbe gefallen läßt; wie Salomo Spr. 20, 12. sagt: „Ein hörend Ohr und sehend Auge, die macht beide der HErr.“

109. Es sind göttliche Werke, ein sehend Auge und hörend Ohr; denn der Teufel pflegt gern beiderlei zu hindern, und hat Lust an Augen, die nicht sehen, und Ohren, die nicht hören, sondern taub sind. Der Papst hat auch Ohren, wie er in seinen Bullen zu brüllen pflegt: Es ist vor unsere Ohren gekommen etc. Aber es sind Augen und Ohren der goldenen und silbernen Götzen, die da Ohren haben, und hören nicht etc., Ps. 115, 4. 6.

110. Wenn aber GOtt gnädig ist und uns in Gnaden heimsucht, so gibt er ein sehend Auge und hörend Ohr. Wie im Königreich Egypten Joseph ein sehend Auge ist und Pharao mit seinen Fürsten und Knechten ein hörend Ohr. Denn sonst, wenn das Ohr gern hören wollte, und hat kein Auge, das ihm vorleuchte und richte, das ist ein großer Jammer. Wie wir im Papstthum ganz begierig waren, daß wir auch nur Eines Psalmen Auslegung hätten hören mögen. Einem solchen Manne wollte ich bis an der Welt Ende gefolgt sein. Aber der Papst war blind, demselben folgten wir nach und sind also beide in die Grube gefallen. Wiederum aber, wenn das Auge sieht und das Ohr nicht darauf merkt, das heißt den Geist vergeblich ausschütten; welches man, wie Sirach Cap. 32, 6. sagt, nicht thun soll: Ubi non est auditor, ne effundas sermonem: Wo man nicht zuhöret, da halte auch still mit reden.

111. Derhalben ist dies eine treffliche große Wohlthat, damit GOtt Joseph ehrt, da er dem sehenden Auge auch ein hörend Ohr gibt, nämlich, einen solchen großen König. Diese Gabe hat der Poet Hesiod auch gerühmt, der da sagt, daß dreierlei Leute seien: die ersten sind die allerbesten und trefflichsten, die von selbst weise und verständig sind; die andern sind, welche dem, der sie ermahnt, gern folgen und gehorchen; aber die dritten sind die allerbösesten, nämlich, die weder von selbst wissen oder verstehen, was sie thun sollen, hören auch Andere nicht und gehorchen den Sehenden und Verständigen nicht.

112. Auf diese Weise wird nun Joseph ein sehr berühmter und glückseliger Lehrer im Königreich Egypten, ja, wie es scheint, fast in der ganzen Welt, und werden ihm Ohren gegeben, die da offen stehen zu hören. Denn alle Fürsten oder Amtleute im Lande haben seinen Rath angenommen und abgebilligt. Das heißt recht das Geheißen geben zum Wort und Predigtamt. Und wird Joseph erstlich nach dem Geist erhöht, da er so viel und große Schüler bekommen hat, die seine Lehre gehört haben, nämlich, den König selbst, seine Fürsten und Räthe. Da sind viel schöne Predigten gefallen.

113. Die erste Ehre ist geistlich, so die Kirche angeht, nämlich, da er ein Prophet, ein Lehrer und Seelenhirt wird. Und ist zwar der andern Patriarchen keiner mit so großer Ehre und Herrlichkeit geziert gewesen. Es haben ja wohl Abraham, Isaak und Jakob ihre Kirche und Gemeinde gelehrt und auch etliche von den Heiden bekehrt, als, Cananiter, Tyrier, Philister etc., aber ihre Schule ist nicht zu vergleichen mit dieser Zuhörerschaft, so Joseph gehabt, aus ganz Egypten und andern umliegenden Ländern, wie im 105. Psalm V. 22. geschrieben steht: „Er hat seine Fürsten unterweisen nach seiner Weise, und seine Aeltesten Weisheit gelehrt“, ja, er hat das ganze Königreich Egypten mit dem Könige selbst und seinen pekidim, das ist, Amtleuten, und den Völkern, ja allermächtigst unterlegen, unterwiesen und gelehrt.

114. Das heißt wahrlich Einen recht erhöhen nach der Trübsal, wenn wir nur in den Anfechtungen beharren könnten; wie Joseph mit großmüthigem Herzen beständig geblieben ist auch in der höchsten Noth und Beschwerung. Darum folgt auch nun eine unaussprechliche große Frucht, so Kirche, Polizei und Hausregiment angeht. Da magst du nun gar eine neue Gestalt aller Dinge in der Welt sehen, darum daß Joseph unserm Herrn GOtt stille gehalten und sich von seiner göttlichen Hand hat zubereiten lassen. Nun wird Pharao zu GOtt bekehrt; und zu allererst erkennt er den Heiligen Geist, damit GOtt seine Kirche regieret, da er sagt: „Wie könnten wir einen solchen Mann finden, in dem der Geist GOttes sei?“ Wiewohl es scheint, daß er zuvor auch fromm und göttlich gewesen sei; aber da wird ohne Zweifel Abgötterei mit untergelaufen sein: allhier aber, da er GOttes Wort predigen hört, nimmt er nun an die rechte Erkenntniß des wahren GOttes und läßt die Götzen und Abgötterei fahren.

\quellseite{355}

115. Dies schöne Licht werden ohne Zweifel die andern Fürsten und vornehmen Leute im Lande sammt dem Könige auch erkannt haben. Denn es ist eine große Gnade und Wohlthat, wenn im Lande der Gottesdienst und was sonst mehr zur Religion gehört, wiederum gebessert und aufgerichtet wird; wie ihnen damals der Glaube und das Licht der himmlischen Lehre ist gegeben worden. Denn wie hat Pharao den Geist GOttes erkannt? Es mag zwar wohl sein, daß er ihn zuvor auch etlichermaßen erkannt, wie die Vorfahren davon gelehrt, und dasselbe von Abraham gemeldet und immer weiter mag ausgebreitet gewesen sein: aber die Erkenntniß wird in ihm so groß nicht gewesen sein, daß er hätte können den Heiligen Geist in den Heiligen und Gläubigen erkennen.

116. Derhalben ist ihm solche Erkenntniß allgemach geoffenbart und allmählich in ihm gemehret worden durch das Gespräch, so er mit Joseph gehalten. Darum urtheilt er von Joseph und seiner Lehre bald und sagt: Das redet GOtt; das ist ein Mann, da redet der Geist GOttes. Dies ist das Größte und Vornehmste in dieser Historie; wie man denn vor allen Dingen am ersten das Reich GOttes suchen soll, Matth. 6, 33. Was aber die weltliche Polizei und Hausregiment belanget, davon wollen wir hernach handeln. Hier wird Pharao die Thür geöffnet zum Reich und zur Erkenntniß GOttes und des Heiligen Geistes ohne Beschneidung und ohne alle anderen äußerlichen Ceremonien der Altväter. Denn Joseph hat ihn dazu nicht gezwungen, dieselben auch von ihm nicht gefordert, da ihm der Heilige Geist eingegeben, daß solche Ceremonien nur ihm und seinen Nachkommen auferlegt und befohlen wären: aber der König in Egypten oder andere Völker sollten mit nichten damit beschweret werden. Die Beschreibung ist dem Samen Abrahams nöthig gewesen um des zukünftigen Messias willen, welcher von demselben Samen sollte geboren werden. Aber dieser Pharao wird ein Christ und wird unterwiesen im Wort vom Glauben und rechter Göttlichkeit, wie andere Heiden auch sind selig geworden ohne Beschneidung, und redet der Heilige Geist aus ihm.

\begin{quote}
\textbf{V. 39. Und (Pharao) sprach zu Joseph: Weil dir GOtt solches alles hat kund gethan, ist keiner so verständig und weise, als du.}
\end{quote}

117. Siehe doch, wie dies Ohr so fleißig zuhört und so gehorsam ist, wie es doch mit so großem Fleiß das sehende Auge und dazu die Gaben GOttes an ihm erkennt. Es bekennt GOtt, der, was verborgen ist, offenbart, und der dies Auge sehend und so weise gemacht hat, und das so wohl thun kann. Wiewohl, sagt er, ich ein König bin über ganz Egyptenland, will ich doch dich mit gutem geneigten Willen gern hören und dir gehorchen. Das geht schön zu. Hätte aber ein solcher großer König nicht billiger also sagen mögen: Das Regiment steht bei mir, dazu gebührt mir auch zu rathen und die Gewalt, Klugheit oder Verstand steht alles bei mir: dir will gebühren, daß du uns gehorsam bist und dich nach unserm Mandat und Befehl richtest. Wie es jetziger Zeit bei dem Adel, Grafen und Baronen zugeht, die hören traun einen Andern nicht so bald, sonderlich wenn er fremd und eines geringen knechtischen Standes und Ansehens ist. Ja, unsere Kaufleute, Bürger und Bauern verlassen sich auch darauf: wenn sie berühmt sind, viel Geld und Gut haben, wollen sie deshalb nicht leiden, daß man ihnen einreden, sie strafen oder ermahnen soll. Ich, sagen sie, bin des Richters im Dorfe oder des Bürgers Sohn; soll mich der Pfaffe regieren? ich kann es so wohl lesen als er.

118. Denn also wird hin und wieder das Predigtamt verachtet und mit Füßen getreten, beide von hohen und niedrigen Ständen, die, weil sich die tollen unsinnigen Leute träumen lassen und meinen, es werde diese große himmlische Gabe den Menschen angeboren und erbe durch fleischliche Geburt von Einem auf den Andern, so doch die Gaben in weltlicher Polizei dermaßen von Einem auf den Andern nicht erben können. Denn das folgt noch lange nicht: Dieser ist von einem Vater geboren, der ein edler, berühmter und großmüthiger Mann gewesen, darum wird er ihm an Tugend und Großmüthigkeit auch müssen gleich sein. Nein, die Natur gibt ihm solche Tugenden nicht; wiewohl diese Geburt natürlich ist, nämlich, daß von einem Fürsten oder Grafen auch ein Fürst oder Graf geboren werde. Aber die Gaben und sonderlich die da geistlich sind, kommen nicht also auf die Erben der äußerlichen Güter, darum daß die Erbsünde im Wege liegt und solches verhindert. Und eben daselbst muß man auch also von andern großen äußerlichen Tugenden der tapferen Helden verstehen. Pyrrhus ist seinem Vater Achilleus nirgends gleich; auch ist Astyanax dem Hector nicht gleich. Ja, es pflegt wohl das Gegentheil zu geschehen, wie solches das gemeine Sprichwort bezeuget: Heroum filii noxae: Großer Leute Kinder gerathen selten wohl.

119. Derhalben ist es gar ein vergeblicher Ruhm, daß man sich von großem Geschlechte oder Adel der Eltern viel rühmen will; wie sich die Juden auch ihres Vaters Abraham, Joh. 8, 39., rühmten. Aber Christus antwortete ihnen darauf also: „Wenn ihr Abrahams Kinder wäret, so thätet ihr Abrahams Werke.“ Wo nun die Tugenden und ehrlichen Thaten dem Adel der Vorfahren nicht gleich sind, da verwundern und rühmen sich die Nachkommen vergeblich der großen Titel, Schild und Wappen ihrer Vorfahren.

120. Es ist übrigens gemeiniglich solchen Leuten die Hoffahrt von Natur angeboren, daß sie um der fleischlichen Geburt willen, wenn dieselbige etwa ein wenig herrlich ist, zu stolzieren pflegen; wie das an den Nachkommen Josephs selbst auch zu sehen ist. Denn ich halte dafür, daß im Volk Israel bei Andern niemals größere Hoffahrt und Stolz gewesen sei als bei den Kindern und Kindeskindern Josephs. Denn also haben sie sich gerühmt, sonderlich die Ephraimiten: Unser Vater ist der weiseste, gelehrteste, frömmste Fürst in Egypten gewesen, darum mögen wir ja billig rühmen und stolzieren; wie das in den Historien der Könige und Richter zu sehen ist. Hilf GOtt, wie haben sie mit so großem Stolz wider Gideon und Jephthah gemurret! Und da noch ihr Vater am Leben gewesen, haben sie sich wider ihre Nachbarn, die Philister, gerüstet, welche auch in der Schlacht alle Ephraimiten todgeschlagen, da sie das zehende Auge nicht hören wollten, nämlich, ihren Vater Ephraim. Es ist darnach allein Beriah übrig geblieben, welchem der Name gegeben worden ist von Schmerz, Jammer und Trübsal, 1 Chron. 8. 13. Stolze Esel werden daraus, aus den Söhnen der Helden, die sich der Tugend ihrer Väter rühmen, und doch derselbigen nicht nachzufolgen gedenken, sondern sich träumen lassen, sie seien auch große Helden, weil sie von so tapfern Helden herkommen und geboren sind.

121. GOtt aber theilt die Gaben aus nach seinem Wohlgefallen, und werden von Natur nicht auf Andere geerbt und fortgepflanzt. Denn auf solche Weise hätte Pharao auch wohl zum Schenken sagen können, welcher die Weisheit des jüdischen Jünglings rühmte: Was ist es, daß du des armen geringen Knechtleins Kunst so hoch rühmest, gleich als ob ich selbst nicht weise genug sei, oder als ob meine Fürsten und Räthe nicht auch daselbst ausrichten könnten? Da wäre wahrlich kein hörend Ohr gewesen, und würden Alle Josephs Weisheit verlacht und sein damit gespottet haben; wie unsere Leute zu thun pflegen. Ja, es würde über das der Teufel auch leichtlich etliche Mißgönner und Lästerungen haben finden können, damit er dasjenige, so Joseph gesagt und gerathen hätte, aufziehen mögen. Aber Pharao hat dies arme Knechtlein, der elend und gefangen war, nicht verachtet, sondern mit aller Demuth weicht er des Schaf- oder Hirten Sohne, wie dazumal Josephs Vater war. Einem Hirtenknaben hört er zu, welche wahrlich bei den Egyptiern dazumal gar verächtlich gewesen sind; wie wir hernach hören werden. Und das haben Pharao und seine Fürsten wohl gewußt, und haben sich doch der Reden dieses verachteten Menschen und des Wortes GOttes, so in ihm ist, so höchlich verwundert, und sich das Wort GOttes in ihm herzlich wohlgefallen lassen. Das ist wahrlich eine herrliche und vortreffliche Gabe eines sehenden Auges und eines hörenden Ohres. Da thut unser HErr GOtt, als wüßte er von seinem Edelmanne oder Fürsten, nimmt einen Kuhhirten, und regieret durch ihn Könige und Fürsten, und den besten Theil der ganzen Welt.

\quellseite{356}

122. Da scheint es allerdings, als thue GOtt nicht recht, daß er nicht auch zugleich mit dem Adel und andern herrlichen Ständen Weisheit gibt. Wie zu unsern Zeiten die Fürsten wegen ihrer Herrschaft und Gewalt, so sie haben, sich selbst auch solche Klugheit und Verstand zuschreiben, daß sie von allen Sachen sehr wohl urtheilen können. Auf daß wir aber nun wissen mögen, daß wir in der Erbsünde geboren sind, werden wir von GOtt also regieret, daß die Gabe der Natur nicht allezeit folgt und die Kinder den Eltern auch nicht allezeit in Weisheit und Tugenden gleich sind.

123. Das hat Pharao an Joseph wohl erkannt; darum ist es nicht ein geringeres Wunderwerk, daß Pharao den Joseph gehört und ihm gehorcht, als daß Joseph gelehrt hat. Denn GOtt regieret sein Herz also, daß sich der König daran nicht geärgert, daß er mit seinen Fürsten dieses sehr armen Hirten, der zu eigenem Gefängniß verdammt war, Schüler wird. Das würde jedwund keines Richters Sohn auf dem Dorfe thun. Du würdest solche Demuth und Bescheidenheit jetzt auch in den allergeringsten Ständen unserer Amtleute oder Steuerbeamten nicht leichtlich finden; geschweige denn bei den Fürsten, Grafen- und Edelleuten. Aber „GOtt“, sagt der 113. Psalm V. 5–8., „hat sich hoch gesetzt, und sieht auf das Niedrige, daß er den Armen aus dem Kothe erhöhe, und ihn setze neben die Fürsten.“

124. Also war David auch ein armes Knechtlein seines Vaters, des Kuhhirten, und ist zu höchster Gewalt und Herrlichkeit erhoben worden. Wiewohl seine Nachkommen sich auch mit Stolz und Hoffahrt versündigt und derhalben auch zum öfteren erschrecklicher Weise gestraft worden sind.

125. Also sind die Römer auch gar thöricht und unsinnig gewesen ihres großen Rühmens und Hoffahrt halben. Und der Türke thut jetzt auch also. Aber doch stolziren die Juden vornehmlich sehr und rühmen, daß sie Abrahams Same seien, wie sehr auch demselben Christus widersprochen hat, da er Joh. 8, 39. sagt: Ihr seid wohl Abrahams Kinder: aber ihr thut ja Abrahams Werke nicht. So werden auch Josephs Nachkommen gesagt haben: Ich bin Josephs Sohn, gesegnet von Jakob, meinem Großvater, darum muß ich der Heiden Land einnehmen. Und dazu haben sie die rechten schalkhaftigen Meister und Zuchführer gehabt, nämlich, die falschen Propheten. Aber, wie droben gesagt ist, sie sind darüber oftmals gar hart gestraft worden, haben es theuer genug bezahlen müssen; sie wurden darob rein todtgeschlagen.

126. Denn das ist der rechte Lohn, so den Hoffärtigen gebühret, und wird solches eben auch unserm Adel dermahleinst widerfahren, dieweil sie alle Andern, so von geringerem Geschlechte und Stande sind denn sie, mit dem größten Stolz verachten und gleichsam mit Füßen treten: welches ein gewiß Zeichen ist ihres Verderbens. GOtt wird diese Hoffahrt in die Länge nicht dulden können, darum werden sie zu Grunde gehen.

127. Es muß zwar ein Unterschied sein der Person und Stände in diesem Leben; wie ein Jude recht unterschieden wird von einem Heiden und was derselben Stände mehr sind. Aber wer um einiger Hoheit oder Vorzugs willen stolziret, er sei gleich Papst oder sonst der allergewaltigste Monarch, derselbe ist vor GOtt schon gestorben. Denn es steht Joh. 1, 13. geschrieben: „Welche nicht von dem Geblüte, noch von dem Willen eines Mannes geboren sind“ etc. Es ist zwar unser Ursprung von dem Geblüte, und die unterschiedenen Stände der Könige, Fürsten, Doctoren etc. sind auch nöthig, welche alle sonst nirgends herkommen denn von dem Geblüte. Aber wie herrlich und adelig dieser Ursprung immer sein mag, soll man darum doch nicht hoffärtig sein; sintemal das allein unser einiger Ruhm und Seligkeit ist, von GOtt geboren sein und nicht von dem Geblüte.

128. Darum ist dieser König sehr zu loben, daß er von sich selbst so wenig gehalten und so demüthig ist, und den armen Hirten gar nicht verachtet. Denn unser Herr GOtt will gefürchtet sein, welcher seine Gaben austheilt ohne alles Ansehen der Person; wiewohl die Welt doch nicht sein kann ohne Unterschied der Personen. Denn es müssen pekidim, wie es im Hebräischen lautet, das ist, Fürsten, Amtleute und andere Aufseher sein. Die Gaben aber folgen demselben Unterschied nicht nach. GOtt sieht die Person nicht an, sondern sieht auf das Niedrige im Himmel und Erden, das ist, beide in der Kirche und auch in weltlicher Polizei. Er kann in seinem Reiche keinen Stolz oder Hoffahrt leiden; und wenn ich anfange stolz zu werden, so kommt GOtt, und erweckt bald einen Andern, der geringer ist denn ich, und wirft mich gar weg. Er stößt einen Doctor gar bald in die Hölle. Also müssen wir auf Erden pekidim, das ist, Oberste und Regenten haben; aber denen, die da hoffärtig sind, widersteht GOtt und gibt den Demüthigen Gnade.

129. Auf diese Weise ist nun der König zum Glauben und Gottseligkeit bekehrt worden, und hat durch sonderliche Gnade GOttes Joseph vor allen seinen andern Weisen und Räthen gelobt und gerühmt. Und wird nun erstlich mit kurzen Worten beschrieben seine Theologie und Bäßthum im Geiste, damit Joseph den König, dazu seine Fürsten und ganz Egyptenland unterwiesen hat. Er wird auch dasselbe Amt, so er geführt, in einer Predigt nicht allein verrichtet haben, sondern hat oftmals gelehrt, und dazu Priester verordnet, die auch dasselbe gelehrt haben.

130. Nun folgt hernach die leibliche Erhöhung. Denn dieweil Pharao mit göttlichem Lichte begabt gewesen, hat er ihn gesetzt zum Bischoff und obersten Regenten in seinem Königreiche auch in zeitlichen Dingen. Denn also sagt er zu ihm: Du besehlst mir, ich soll mich nach einem verständigen und weisen Mann umsehen, aber in Wahrheit, du bist der rechte nachon; denn also lautet das hebräische Wort. Ich muß dich anders schmidten, du bist wahrlich ein rechter Prophet und der Heilige Geist redet durch deinen Mund; darum ist es billig, daß ich dir gehorche, und dich allen meinen Fürsten und Hofdienern vorziehe und dich über das ganze Königreich Egypten setze.

131. Also wird Joseph nun wahrhaftig wiederum in seinen vorigen Stand gesetzt, nachdem er nun die Herrlichkeit des Königreichs und Priesterthums oder der Erstgeburt bekommen hat, deren er kurz zuvor in seines Vaters Hause durch seine Brüder war beraubt worden. Denn GOtt hat den Neid und die Abgunst seiner Brüder ganz weislich gebraucht, also daß er ihn dadurch zum Priesterthum und Regiment erheben möchte, und ihm das ganze Egyptenland dermaßen unterwerfen, daß er dazu keinerlei äußerlichen Gewalt hat brauchen dürfen. Darum wird nun auch beschrieben, wie ihn das weltliche Regiment zu verwalten befohlen worden.

\section[Vierter Theil]{Vierter Theil. Wie Joseph von Pharao zu großen Ehren erhoben wird}

\begin{quote}
\textbf{V. 40. Du sollst über mein Haus sein, und deinem Wort soll alle mein Volk gehorsam sein; alleine des königlichen Stuhls will ich höher sein, denn du.}
\end{quote}

132. Das ist nun eine sehr große Herrlichkeit, welche in sich begreift die Herrschaft oder Regiment über ganz Egyptenland. Und Pharao setzt hier Joseph über alle andern Fürsten und Amtleute seines ganzen Königreichs, ihn, den König selbst, allein ausgenommen.

133. Es zanken aber die Grammatiker mit einander über dem hebräischen Wort jischak, und ich weiß zwar seinen Ursprung nicht, und haben denselbigen auch die Hebräer, so die Wörter, welche mehr als Eine Bedeutung haben, zu erklären pflegen, noch nicht ausforschen können, sind auch nicht tüchtig oder geschickt dazu, daß sie in andern Worten oder Stellen rechten gewissen und eigentlichen Verstand der Art und Weise zu reden sollten fassen und anzeigen können; aber die Worte und rechten Verstand derselbigen zu verdunkeln und zu verwirren, dazu sind sie sehr geschickt. Es soll aber ein rechter Lehrer nicht zweideutig reden, das ist, mit solchen Worten umgehen, die mehr denn einerlei Bedeutung haben; denn er wird sonst nichts Rechtes lehren: sondern man soll alle Ungewißheit der Wörter von den Prädicanten, wie die Dialektiker in den Schulen reden, hinweg thun, das heißt, ein Lehrer oder Redner soll sich derselben enthalten; denn er muß solche Worte brauchen, so gewiß und deutlich sind. Er muß sich befleißigen, daß er eigentlich und verständlich rede. Darum lehren die hebräischen Ausleger sehr übel, sintemal alle ihr Studiren und Fleiß nur darin besteht, daß sie sich mit den Wörtern, die mehr denn einerlei Bedeutung haben, bekümmern, und sind nur auf Sophisterei geschickt, werden unnütze Wäscher daraus und seine rechten Dialektiker oder Lehrer, sondern Sophisten, Betrüger und Verführer.

\quellseite{357}

134. Wir wollen aber doch ihrer Bosheit zu dem Ende gebrauchen, daß wir der Sprache Nutzen und die rechte eigentliche Bedeutung dieses Wortes finden mögen. Wo aber eine sonderliche Figur ist in einem solchen Worte, das mancherlei Bedeutung hat, das ist kein Betrug, sondern ist mehr eine Zierde, also daß gleichwohl die rechte eigentliche Bedeutung des Wortes bestehen bleibt.

135. Derhalben ist das hebräische Wort naschak, welches küssen heißt, ein abgeleitetes Zeitwort, wie es die Grammatici nennen; und sie sagen, es habe seinen Ursprung von dem Stammwort neschek, auf Lateinisch, arma vel instrumenta bellica, das ist, allerlei Wehre oder Kriegsrüstung. Daher, sagen sie, soll naschak so viel heißen als, sich rüsten oder die Wehre zur Hand nehmen; daß dies hier der Verstand sei, als wollte Pharao zu Joseph sagen: Alles Volk soll sich nach deinem Wort oder Befehl rüsten und zur Wehre schicken.

136. Die Andern sagen, dies Wort komme von schakak; aber was die rechte Bedeutung des Worts sei, darüber sind sie ungewiß, und deuten es doch also, daß es so viel heißen soll als, eilends einhergehen oder, mit einem Laut und Geräusch sich bewegen, auf Deutsch, rauschen, wie man sagt: Du hast mich überrauscht, item, laß rauschen etc. Damit wird angezeigt ein leises Murmeln und Rauschen mit den Füßen oder aber mit Wehre und Harnisch; item, das Toben von Männern, wie eines rauschenden Wassers, das da schnell und mit einem Laut oder Brausen daherfließt. In welcher Bedeutung es Joel 2, 5. gebraucht wird, und Jes. 13, 4.; und daher wollte ich es auch ableiten, wenn es vergönnt wäre. Denn ich gebe den Punkten und mancherlei Bedeutungen, so die Juden bei den Worten haben, keinen Glauben. Und unser deutsches Wort „rauschen“ reimt sich sehr wohl dazu; denn es begreift beiderlei Regiment in sich, als, daheim im Frieden, und sonst draußen, wenn man Krieg führen soll, und weiß man sich beides zur Zeit des Krieges und Friedens befleißigen soll, als wollte Pharao zu Joseph sagen: Du sollst der oberste Regent sein im Kriege und auch im Frieden, du sollst beide die Polizei und auch das Kirchenregiment verwalten. Du sollst Priester und Lehrer verordnen, darnach auch in den Städten Regenten setzen, wie das Recht und Polizei dieses Königreichs ausweist; und wo es die Noth erfordern wird, sollst du auch Krieg führen und den Feinden wehren. Auf diese Weise erhebt der König die Wunderwerke Gottes und ehrt Joseph hoch, wie er denn hernach sagen wird: „Ohne deinen Willen soll niemand seine Hand regen“, das ist, alles Volk soll sich regen und bewegen, thun und handeln, wie du sagst und haben willst. Es soll bei dir das Mauschen stehen.

137. Daraus folgt, daß Joseph auch ein Oberster im Kriege gewesen sei. Denn wo man ein Land oder Königreich regieren soll, da bedarf man nicht allein Besetzung mit Proviant, sondern man bedarf dazu auch Kriegsrüstung, wenn etwa die Gefahr kommt, daß die Feinde herein brechen und das Land verderben wollen. Und Salomo sagt auch, Spr. 21, 31.: „Rosse werden auch zum Streittage bereitet.“ Zur Zeit des Friedens muß man die Dinge zubereiten, so vonnöthen sein werden, wo man den Feinden wehren und widerstehen soll. Also läßt es sich nun ansehen, daß beiderlei Sorge dem Joseph sei befohlen gewesen, daß er ein Vater des Vaterlandes sein soll, das Volk zu nähren und zu regieren; darnach sollte er auch ihr Alexander der Große und Achilles sein, wenn Krieg entstehen würde.

138. Darum ist Pharao ein sehr weiser König gewesen, sintemal er gesehen, daß der Heilige Geist in Joseph war und daß seine Rathschläge und alle sein Thun glücklich hinausgehen würde. Er denkt, der Mann wird es thun. Und wer wollte wider einen solchen Mann streiten, da der Heilige Geist Heerführer und Rathgeber ist? Ich wollte wahrlich lieber unter Joseph, wo derselbe sollte Feldherr sein, mich als ein Kriegsmann gebrauchen lassen, als unter Alexander dem Großen. Denn wiewohl man bisweilen auch Glücke finden möchte, die von Natur unmännlich und treibar sind, sind doch derselben fast wenige. Aber es sind hin und wieder über die Maßen viel Thrasone, das ist, die sich hoch rühmen, pochen und schmarotzen können: sobald aber dieselbigen den Feind sehen, nehmen sie die Flucht. Das sind die Störcahnen, Marterhansen etc. Wo aber unser Herr Gott einen solchen Heerführer und Kriegsmann gibt, wie Joseph oder Mose gewesen, der kann mit wenigen Kriegsvolk, welches doch fürchtsam wie die Kirche ist, einen großen Haufen Löwen, das ist, der großen Eisenfresser, überwinden und niederlegen. Denn es steht Spr. 28, 1. geschrieben: „Der Gerechte ist getrost, wie ein junger Löwe“ etc.

139. Also hat ohne Zweifel Joseph einen sehr großen Muth gehabt; und dazu ist auch gekommen die große Erfahrung und Uebung in vielen Dingen, da er fein klug und gefänglich mit so unverzagtem Herzen ertragen hat. Er fährt dem Teufel und Tode in den Rachen hinein und fürchtet sich doch nicht, verzagt auch nicht; denn der Heilige Geist ist bei ihm, welcher ihm den Glauben gibt. Unter demselbigen Heerführer hätte er drei gewaltige Haufen niederlegen können.

140. Also wird Mose geboten, daß er das Volk Israel sollte aus Egypten führen auch wider des Königs Willen; welches weder Scipio, noch Hannibal, noch Alexander der Große genung hätte. Denn es hat nicht geschehen können ohne große Einspörung und Schrecken des Königreichs; welches auch derselbigen Annerung halben sehr hart gestraft worden. Zu solcher herrlichen That gehört wahrlich ein unverzagt Herz und starker beständiger Glaube aus dem Heiligen Geist, welches alles Mose hat haben müssen.

141. Und wenn etwa ein Kaiser oder oberster Feldherr solche Kriegsleute haben möchte, die im Glauben und Hoffnung göttlicher Hilfe also beständig wären, der sollte wahrlich mit einem geringen Haufen die ganze Welt bekriegen und überwinden. Denn solche Kriegsleute haben nicht allein das an sich, daß sie den Tod und auch andere Gefährlichkeit verachten können; sondern haben dazu auch noch ein solch Herz, das des Sieges und Regierun[g] göttlicher Gewalt gewiß ist. Ein solch Herz lehrt ihre Hände streiten und ihre Fäuste kriegen. Diese Stücke sind an Alexander nicht, welcher, ob er wohl einen großen Muth gehabt, thut er doch Alles zufällig, ohngefähr, ohne Glauben und rechte Anrufung. Abraham oder Joseph und andere fromme gottselige Heerführer greifen es an und treffen mit den Feinden zusammen in starkem Glauben und ungezweifelter Hoffnung. Darum brechen sie auch hindurch, siegen und behalten das Feld.

142. Und lieber, siehe doch an das kleine Häuflein des jüdischen Volks mit der Stadt Jerusalem, da die Heiden rings umher gelegen, wie wenig Landes sie gehabt und wie mit großer Menge und Gewalt der Feinde sie umgeben gewesen, welchen alles es weder an großer Macht noch an Weisheit, noch auch an dem bittersten Haß, dies Volk gar zu vertilgen, gemangelt hat. Und während das Land Canaan nicht größer gewesen denn dreißig Meilen in die Breite und fünfzig in die Länge, hat es sich doch gleichwohl gegen so viel herrliche und gewaltige Königreiche gewehret und gehalten. Eine Hand voll Bettler sitzen da zum Trotz so viel Königen und schlagen sich mit ihnen; denn der Geist Gottes hat sie geschützt.

143. Daher sind gekommen so viel herrliche Siege der Richter und anderer Helden, welche gesiegt haben, und unverletzt geblieben und davon gekommen sind. Wie das zu sehen ist an Gideon, Jephthah, Danid etc., die niemal[s] von Feinden verwundet worden; die Andern aber, so erschlagen worden, sind wegen ihres Unglaubens und verzagten Herzens umgekommen.

144. Ein solcher Held ist ohne Zweifel Joseph auch gewesen; wiewohl ich nicht weiß, ob er auch Krieg geführt habe. Aber wo es einmal die Noth erfordert, daß man wider die Mohren und Araber hat streiten müssen (wie es denn wohl glaublich ist, daß etliche Feinde oder Widerspenstige im Reich gewesen sind, die man hat halten und zwingen müssen), da halte ich wahrlich dafür, daß er nicht geringer gewesen sei, denn Josua oder David. Mit solchen Feldherren wollte ich gern in Schlachtordnung stehen, wenn schon der Türke mit allen seinen Heerschaaren uns bekriegen würde. Denn sie haben einen großen Muth gehabt, den Tod und alle andere Gefährlichkeit zu verachten, wegen ihres starken unüberwindlichen Glaubens, und endlich um der Hoffnung des gewissen Sieges willen, so sie gehabt haben.

145. Derhalben hat Pharao seinem Königreich damit sehr weislich und sehr wohl gerathen und vorgegangen, da er denselben diesen Mann zum Regenten gesetzt, an welchem er gesehen, daß er den Geist Gottes gehabt. Und wie er ihn zum Bischof oder Aufseher gesetzt hat über das Getreide und Kornhäuser: also hat er auch ganz Egypten durch ihn zum Krieg gerüstet, nachdem er das Predigamt bestellt, und angezeigt, wie man GOtt recht dienen sollte.

\quellseite{358}

146. Darum geht diese Weise zu reden, populus movebitur, das ist, das Volk wird sich aufmachen und nach werden, insgemein auf alle Wercke des Volcks, es sei gleich, da sie haben Predigt hören, Versammlungen halten, oder auch in Krieg ziehen sollen. Denn diese Dinge alle können nicht geschehen, daß man sich darüber nicht bewegen sollte und ein Geräusch machen. Und hat das Volck und der König selbst auch an solchem obersten Hauptmanne oder Regenten guten Gefallen gehabt: sie sind gern in das Feld gezogen, sind wohl zufrieden, sicher und wacker gewesen, wenn sie gegen die Feinde haben streiten sollen. Denn sie haben an ihm sonderliche Frömmigkeit gesehen, daß er mit allen Sachen recht umgegangen ist, und in allem, was er gethan, Glück und Heil gehabt.

147. Solche Gaben sind sehr selten, und ihrer sind wenig, denen sie von GOtt gegeben werden, welchen Mose 2 Mos. 15, 3. diesen Titel und Namen gibt: „Der HErr ist der rechte Kriegsmann, HErr ist sein Name.“ Wenn der in das Feld kommt, so hüte sich Roß und Mann. Wiederum aber steht auch von ihm geschrieben im 76. Psalm V. 13.: „Der den Fürsten den Muth nimmt, und schrecklich ist unter den Königen auf Erden“; alsdann ist allerdings kein Glück im Kriege. Denn da hat GOtt die Pfeile des Bogens, Schild, Schwerdt und Streit zerbrochen, daß die Stolzen beraubet werden und entschlafen sind, V. 4. 6. Sie lassen Harnisch und Spieß fallen.

148. Und die des Krieges erfahren sind, sagen selbst auch, daß man es an etlichen sonderlichen Zeichen sehen und merken könne vor der Schlacht und ehe denn man mit dem Feinde zusammentrifft, welcher Haufen siegen und das Feld behalten werde. Denn an dem einen Haufen, der geschlagen werden soll, spürt man, sagen sie, daß derselbe ganz still und furchtsam ist, beide Mann und Roß, und daß da gar kein Geschrei, weder der Menschen noch Rosse, gehört werde. Aber bei dem andern Haufen sind sie sehr lustig und wacker zur Schlacht, und haben ein solch Herz, welches ihnen etlichermaßen zuvor sagt und verheißt, daß sie siegen sollen. Denn der HErr ist der rechte Kriegsmann, der die Hände lehrt streiten: welcher auch den Fürsten den Muth gibt und nimmt, wann er will.

149. Also haben wir droben, 1 Mos. 14, 14. ff., gehört, daß Abraham mit einem geringen Haufen, nämlich mit dreihundert und achtzehn Knechten, in seinem Hause geboren, die zuvor an den Krieg nicht gewöhnt und die er damals zum erstenmal dazu gerüstet, hat dürfen vier gewaltigen Königen nachjagen und dieselben angreifen; und da er sie in die Flucht geschlagen und niedergelegt, hat er doch alle sein Hausgesinde und dazu seinen Bruder Lot mit alle seiner Habe und großer Ausbeute unverletzt davon gebracht. Wollen wir es nun dafür halten, daß solches ohne große und heldenmüthige Stärke der Seele geschehen sei? Wahrlich, er wird sie nicht angegriffen haben mit zweifelhaftigem Herzen, das gezappelt und gesagt hätte: Wer weiß, was der HErr thun wird? Vielleicht werden wir siegen und das Feld behalten. Ja, er wird vielmehr als ein Mann voll Geistes mit starkem Muthe und gleichsam mit einem Glaubensstürme fortgezogen sein, und hat diese gewisse Ueberzeugung durch den Heiligen Geist mit in die Schlachtordnung gebracht, daß er gewißlich siegen und seinen Bruder erretten würde. Wenn GOtt solche Kriegsleute uns auch geben würde, da wollten wir uns auch vor der großen Macht und Gewalt der Türken gar nicht fürchten oder erschrecken. Die Stötzmesser und Marterhansen thun es nicht; wie wir das zu unserm großen Schaden wohl erfahren haben.

150. Das ist nun Josephs Ehre, welcher dreizehn Jahre lang ist gedemüthigt gewesen bis in den Tod, da er zwei Jahre oder noch etwas länger im Gefängniß und in der Hölle alle Stunden gewartet, daß er hingerichtet und getödtet werden sollte. Aber er hat all diesen Jammer und so viel Trübsal mit großem unverzagtem Muthe überwunden. Denn des Todes und der Hölle spotten und sie überwinden gehört nicht einem schwachen verzagten Herzen zu, sondern das muß ein unüberwindlich unerschrocken Löwenherz sein. Darum sagt der König zu Joseph: „Deinem Wort soll alle mein Volk gehorsam sein“, es sei gleich, wenn man in Krieg ziehen, oder sonst bei Friedenszeiten zusammen kommen soll in der Kirche zum Gottesdienste, oder auf dem Rathhause, da es die Polizei belanget. In Summa: Es soll bei dir stehen und du sollst Macht haben, zu jeder Zeit allerlei Leute im Lande, wes Standes sie seien, zu versammeln, und ihnen Befehl zu thun, wie sie sich sollen halten. Darin sollen sie dir Gehorsam leisten.

151. Also verstehe ich das hebräische Wort ohne alle Nebenbedeutung, welche ich fahren lasse; und ist diese Auslegung auch gewisser und besser, als der Juden, die solchen eigentlichen und gewissen Verstand nicht beibringen können. Denn ich pflege vornehmlich das zu betrachten, ob sich auch die Auslegung eines Spruches mit dem Glauben reime und demselben gemäß sei. Da wird die Zweideutigkeit gehoben, und wird da der rechte einfältige natürliche Verstand der Worte und Sätze ans Licht gebracht, welchen die Juden nimmermehr erreichen.

152. Er setzt aber das dazu: „Deinem Wort soll alle mein Volk gehorsam sein; alleine des königlichen Stuhles will ich höher sein, denn du.“ Darum halte ich nicht dafür, daß ich darin irre, daß ich die vorigen Worte ausgelegt habe von beiderlei Regiment, im Frieden und in Kriegsläuften, dieweil ja beides zu königlicher Gewalt gehört, welche der König dem Joseph gibt. Derhalben ist er der oberste Regent in Egypten beide in Kriegsläuften, und auch wenn Friede ist. Und ist wohl zu glauben, daß er etwa Krieg geführt und auch herrliche Siege gehabt habe.

153. Wie auch von Mose gesagt wird, daß er Krieg geführt habe wider den König des Mohrenlandes. Wiewohl desselbigen Krieges in der heiligen Schrift nicht gedacht wird, so ist es doch sehr wahrscheinlich, daß er die zehn Jahre lang, nämlich, vom dreißigsten Jahr, da er in das Land Midian geflohen, nachdem er den egyptischen Mann erschlagen, bis auf das vierzigste Krieg geführt und daß er bei den Midianitern große herrliche Thaten verrichtet habe. Denn er war ein Hofdiener, da ihn des Königs Tochter zum Sohn angenommen hatte, und bald von Jugend auf in Kriegshändeln geübt. Und ich halte dafür, daß ihn zumal der König auch selbst oftmals werde gebraucht haben. Darum glaube ich leichtlich, daß er werde ein Oberster oder Hauptmann im Kriege wider die Mohren gewesen sein, wie die Juden erzählen. Denn er ist ein Mann gewesen voll Glaubens und Geistes, dazu eines unverzagten Herzens. Aber da es kund ward, daß er den Egypter erschlagen hatte, und er nun von einem Andern hören mußte: „Willst du mich auch erwürgen, wie du den Egypter erwürget hast?“ 2 Mos. 2, 14.: da hat er Egypten verlassen und ist in das Land Midian geflohen. Denn er war nicht gerüstet mit Waffen, war auch nicht mächtig und wollte den HErrn nicht versuchen; nachdem er aber wieder gekommen, da hat er das Volk aus Egypten herausgeführet.

154. Daraus können wir abnehmen, daß er mit sonderlicher Geschicklichkeit, Rath und Stärke wird begabt gewesen sein. Denn die Leute, so den Muth verloren, dürfen gar nichts wagen; die aber einen hohen Geist haben, und recht männlich oder ritterlich sind, die richten auch viel große herrliche Dinge aus. Es wird 5 Mos. 33, 5. von Mose gesagt: „Er verwaltete das Amt eines Königs, und hielt zusammen die Häupter des Volkes, sammt den Stämmen Israel.“ Denn es ist ein groß, hochtrefflich Werck gewesen, ein solch groß Volk ausführen und dasselbe so lange Zeit regieren: dazu gehört ein großer starker Muth und der Geist GOttes, und auch eine herzliche Zuversicht zu GOtt. Auf diese Weise sind die heiligen Historien in der heiligen Schrift wohl kürzlich beschrieben, aber sie begreifen große herrliche Dinge in sich, die alle sehr wichtig sind.

155. Du siehst aber in dieser gegenwärtigen Historie, wie herrlich und hoch GOtt diejenigen ehrt und erhebt, so auf ihn harren und hoffen, und seine Hand und väterliche Ruthe dulden und ertragen können. Derhalben soll man dies Exempel fleißig einprägen und vorhalten allen Menschen, in weltlicher Polizei und auch in der Kirche, auf daß sie lernen des HErrn harren und in der Anfechtung beständig bleiben. Halt fest! halt fest! „So du glauben würdest“, spricht Christus zu Martha Joh. 11, 40., „solltest du die Herrlichkeit GOttes sehen.“ Es muß der Glaube vorhergehen, und hernach das auch folgen, daß man des HErrn harre, wie Joseph im Stock und Gefängniß des HErrn geharrt; da hat sein Leib in Eisen liegen und sein Herz sich demüthigen müssen, und hat doch immer des HErrn geharrt und ist beständig geblieben.

156. Ei! wie ist es nun so recht geredet, das St. Paulus Eph. 3, 20. sagt: „GOtt kann überschwänglich thun über alles, das wir bitten oder verstehen.“ Unser Gebet ist allzu klein und gering; und Joseph hat sich nicht unterstanden das zu bitten, das er doch endlich erlangt hat, sondern hat in seinem Herzen ein glimmend Docht und zerstoßen Rohr; ja, er hatte das unaussprechliche Seufzen, welches gewesen ist wie ein gerader Rauch, so gen Himmel hinaufgeht. O! ein schön thuribulum und Rauchfaß ist sein Herz gewesen, und Gott hat sehr großen Gefallen an solchem Geruch unsers herzlichen Seufzens und Demuth, daß Einen dünkt, er werde verstinken. Aber es schadet nichts. Halt fest: er hält noch; verzweifle nur nicht. Darum soll man daran fest halten, dazu uns der 27. Psalm V. 14. ermahnt: „Harre des Herrn, sei getrost und unverzagt.“ Denn der Herr wird das glimmende Docht nicht auslöschen, sondern wird es anzünden, und wird das zerstoßene Rohr nicht zerbrechen, sondern dasselbige stärken.

\quellseite{359}

157. Denn es ist Gottes Art und Natur, daß er eben nicht gibt nach unserm schwachen Bitten. Wie Joseph damals begehrt, daß er möchte errettet und aus dem Gefängniß geführt werden, und wiederum zu seinem Vater kommen: Größeres wird er nicht begehrt und gebeten haben. Da läßt ihn unser Herr Gott im Himmel eine Zeitlang Knauwerck und guten Geruch anzünden; aber er antwortet ihm doch mit den Worten Pauli: Du weißt nicht, was du bittest. Ich pflege überschwenglich mehr zu thun, denn ihr bitten oder verstehen könnt. Darum muß ich dich noch etwas länger hinhalten, du mußt noch baß halten. Ich begehre noch mehr des gerechten Taues, Hohel. 3, 6., der da aufsteigt gen Himmel. Hernach aber hat Joseph empfangen, was er zuvor nicht verstanden, auch nie hoffen oder bitten dürfen.

158. Dies sind wahrlich hohe und wunderbare Exempel, welche wir für uns selbst lernen und auch Andern vorhalten sollen, auf daß wir die göttliche Weisheit, Güte, Barmherzigkeit und große Macht erkennen mögen, welche gewißlich bei uns sind. Bonitas, das ist, seine Güte ist gewißlich da, seine Weisheit ist gewißlich da und seine Macht auch. Aber Gott handelt damit nicht nach unserm Willen. Wenn er uns regieret nach seiner Güte und Weisheit, so verstehen wir es nicht; denn wir meinen, Gott kenne uns nicht, wolle auch unser nicht achten, oder darauf denken, wie unsere Anfechtung möchte ein Ende gewinnen. Also urteilt unsere Vernunft. Aber das heißt nicht nach christlicher Lehre oder im Geiste Gott erkennen. Ich kann auch nicht rühmen von großer Vollkommenheit, und bin meinem Fleische feind, welches stets dem Geiste widerstrebt. Darum soll man diese Dinge oftmals bedenken und sich darin üben, auf daß wir prüfen mögen, welches der gute, wohlgefällige und vollkommene Gotteswille sei, Röm. 12, 2.

159. Denn was das sei, daß Gott weise, gütig, mächtig und barmherzig ist, davon versteht niemand etwas, denn allein, sofern er davon speculiren und sich Gedanken machen kann; wie man sich auch etwa von der Gütigkeit oder Weisheit eines Fürsten etlichermaßen eine Vorstellung machen kann, welches doch nur eine lautere Speculation ist und nichts mehr. Aber wenn es zur Praxis, das ist, an die Uebung kommt, daß man es erfahren und mit der That empfinden soll, als, da Gott Joseph vorgreift von seinen Eltern, Großvater, Vater und dem ganzen Geschlecht, und er in fremden Landen, als sollte er die Ehe gebrochen haben, in das Gefängniß geworfen wird, da er dann lange Zeit in großer Furcht gelegen und täglich des Todes hat erwarten müssen: würde auch wohl Einer dasselbe deuten, daß es ein guter gnädiger Wille Gottes wäre?

160. Darum sollen wir nun wissen, daß sich Gott also pflegt zu verbergen unter der Gestalt des ärgsten Teufels, auf daß wir lernen, daß man Gottes Güte, Gnade und Macht nicht mit Speculiren begreifen oder verstehen kann, sondern man muß es mit der That erfahren. Halte du nur fest und harre des Herrn, laß dir genügen an seinem Worte; wie Joseph das Wort des Glaubens hat, er weiß, daß er Jakobs Sohn ist, welcher ja die Verheißung von Gott empfangen von dem zukünftigen Samen und dem Segen. Und da es sich alles widersinnig ansehen ließ mit dieser Verheißung, die seinem Vater gegeben war, hat er doch ohne Zweifel geseufzt und sich beklagt, daß es ihm so gar übel ginge, daß es zu erbarmen wäre. Aber der Herr hat ihm den Trost eingegeben: Harre des Herrn, sei unverzagt, halte fest! Aber da wird ihm wiederum in den Sinn gekommen sein, daß er geseufzt und begehrt hat: Ach wie gerne wollte ich wiederum zu meinem Vater ziehen! Da hört er, daß ihm Gott antwortet in seinem Herzen: Halte fest! Oder daß ich doch nur aus dem Gefängniß kommen möchte! Halte fest! Ja, wie wenn ich im Gefängniß mit Schanden verderben müßte? Halte fest!

161. Dieser Wechsel von Trost und Bekümmerniß ist ihm also immer in seinem Herzen gewesen, bis sich endlich Joseph unversehen und mit großem Verwundern gesagt hat: Ach, habe ich doch nie hoffen können, daß ich dermaßen würde errettet werden, habe auch die große Macht und Gütigkeit Gottes nicht verstehen können, die er an mir geübt und bewiesen hat! Da springt sein Herz in lauter Freuden, freut sich von ganzem Herzen und ist ihm nun wohl dabei, daß er in solcher Schmach, im Tode und Gefängniß gewesen ist. Und ist also Joseph auf beiden Seiten ein trefflich Exempel, indem er geseufzt, und indem er sich wiederum gefreut, da die Anfechtung ihr Ende genommen.

162. Derhalben lernet nun, daß Gott auf solche Weise barmherzig und gütig ist, daß er überschwenglich mehr geben könne und wolle, denn ich verstehen oder bitten kann, Eph. 3, 20. Seine Gnade und Barmherzigkeit ist viel größer, denn daß ich dieselbige Größe mit Gedanken sollte erreichen können.

163. Dabei muß man aber das zugleich auch merken, wie er seine Heiligen versuche und so wunderbar führe, ja ihrer gleichsam spotte. Daher sehen wir, wie die Propheten geschrieben, als, Jesaia, Jeremia und die andern auch: „Herr, warum schläfst du?“ Ps. 44, 24. Und Jer. 12, 2.: „Du lässest sie viel von dir rühmen, und züchtigest sie nicht.“ Und wenn sie schon also seufzen, klagen und schreien, beharren sie doch gleichwohl in der Anfechtung, halten immer fest am Glauben und Hoffnung. Wir sehen zwar sonst nichts, denn allein das Räuchwerk, so sie Gott in ihrer Angst und Noth opfern: aber endlich werden sie doch viel herrlicher errettet, denn sie sich selbst hätten verheißen oder denken können.

164. Das ist nun das Exempel Josephs, der ein Regent über ganz Egyptenland worden, und zum Hohenpriester gesetzt wird, daß er den Kirchen und Pfarrherren mit rechter Lehre vorstehen soll: darnach auch zum obersten Regenten bei Zeiten des Friedens und auch in Kriegsleuften, der da Hauptleute im Krieg bestellen, Gericht, Recht und Obrigkeit in Städten verordnen soll. Und wo er nicht also versucht oder angefochten wäre, hätte er dies alles nicht mit so großem Lob und so großer Geschicklichkeit verwalten können. Tentatio, die Versuchung, macht einen Mann, wie Sirach sagt: „Wer nicht versucht oder angefochten worden ist, was kann der wissen?“ Denn die Sicherheit macht eitel Schwarmgeister, wie Carlstadt und Andere gewesen sind; die unsern Herrn Gott nach ihrer Speculation abmessen und in ihre Gedanken einschließen. Darum sind es unnütze Wäschler, sind gar unerfahren beide in geistlichen und auch in leiblichen Sachen; aber in der Anfechtung zittern sie, sind furchtsam und verzagen endlich gar. Derhalben lerne, wenn dermaleinst die Anfechtung dahergeht und dich trifft, daß du einen guten fröhlichen Muth habest und unerschrocken seist: oder wo du das ja nicht kannst, so halte doch fest, siehe, daß du gleichwohl beharrest und beständig bleibest.

165. Ja, sagst du, wie soll ich denn beharren und beständig bleiben in so großem Jammer und Trübsal? Halte am Worte, und schließe bei dir selbst also und sage: Ich bin getauft, ich glaube an Jesum Christum, den Sohn Gottes. Es gehe darnach der Teufel und Tod, so thuts nichts. Es wird dir die Auserwählung, der Tod und Hölle nicht schaden können; denn der Herr erhält dich mit seiner Hand, so lange du sein Wort hast. Also wirst du erkennen, welches der gute wohlgefällige vollkommene Wille Gottes sei. Aber diese Erkenntniß wirst du nicht erlangen ohne die Erneuerung. Denn das geht vorher, Röm. 12, 2.: „Verändert euch durch Erneuerung eures Sinnes“; und dann allererst wirst du Gottes Güte erkennen. Wie werden wir aber erneuert? Ei, wenn wir den alten Menschen tödten und ablegen; welches dem Fleische wohl sehr wehe thut, wir werden aber von Tage zu Tage immer neuerer; wie vor der Erlösung Josephs die schwerste Tödtung im Gefängnisse und Eisen auch vorhergegangen.

166. Derhalben sollen wir aus seinen Exempel, welches voll aller Tugend ist, die göttliche Weisheit und Güte erkennen lernen, und die Weise der Regierung, wie Gott mit uns pflegt umzugehen, und daß man immer harren soll ohne Aufhören; und daß keine Zeit der Angst und Trübsal so groß oder langwierig sein könne, die uns sollte weich und kleinmüthig machen, oder uns das Wort des Glaubens nehmen, welches eine Kraft GOttes ist zur Seligkeit Allen, die daran glauben, Röm. 1, 16. Darum soll man fest halten.

\quellseite{360}

167. Es ist aber eben dasselbe droben auch gesagt worden, da wir an den Vätern große treffliche Exempel des Glaubens, der Hoffnung und Geduld gehabt haben. Aber in dieser Historie, nachdem Joseph erhöht worden ist, werden wir nun gelehrt, daß GOtt nicht allein helfen könne, sondern auch ganz willig, geneigt und bereit sei, die Seinen zu erretten und zu erhöhen, nur daß er sich eine Zeitlang verberge. Und wer ein Christ sein will, soll sichs merken, daß er diese Dinge fleißig erwäge, bedenke und gar auswendig lerne.

168. Wir haben aber die schönen herrlichen Tugenden Josephs noch nicht alle zu Ende gebracht. Denn deren wir bis daher gedacht, sind noch mehr verborgen und nicht so augenfällig. Jetzt aber werden nun die folgen, so viel höher und herrlicher sind, welche ich mit keinen Worten erreichen oder genugsam ausreden kann. Nachdem er so viel Trübsal, Widrigkeit und Unglück überwunden, wird er nun auch geziert und gleichsam überladen mit unermeßlicher Herrlichkeit, Ehre und Glück. Da gehört nun erst ein großer unüberwindlicher Muth dazu: ja, da gehört erst ein Mann dazu. Du möchtest wohl Etliche finden, die in der Anfechtung und Widrigkeit leicht beharren und Joseph mit Geduld und Glauben etlichermaßen nachfolgen können: aber in diesen Tugenden, als, in solcher Bescheidenheit, Demuth, GOttesfurcht, und daß er sich bei großem Glück so niedrig hat halten können, da hat er nie seines Gleichen gehabt.

169. Denn in so großer Ehre und Glück, da ihn jetzt die ganze Welt gleichsam anlacht, ihn ehrt und so hoch vor Augen hält, sich selbst so fein inne halten, zwingen, nicht stolzieren, sich gar nicht erheben, sondern immer an unsern Herrn GOtt, und was er auch von sich selbst sei, denken, daß ihn GOtt allein aus Gnaden zu solchem großen Glücke gebracht und erhoben habe: das steht in keines Menschen Verstande oder Kräften. Und der heidnische Poet Virgil hat dasselbe auch wohl gesehen, da er sagt: Nescia mens hominum fati sortisque futurae, et servare modum rebus sublata secundis, das ist: Das menschliche Herz kann nicht zuvor wissen, was für Glück oder Unglück ihm begegnen werde, kann auch nicht Maß halten, wenn es Glücks halben erhoben wird, daß es nicht deshalb bald stolz werden sollte. Der Poet schließt einfach, daß kein Mensch sich im Glück mäßigen und inne halten könne, daß er sich deshalb nicht etwas überheben sollte. Alexander der Große, welcher wahrlich einen großen Muth gehabt, ist doch ein rechter armer Knecht seines Muths gewesen, hat sich nicht können bezwingen. Siehe Achilles an, welcher der stärkste und männlichste unter allen Griechen gewesen, wie gar kindisch sich derselbe gehalten, sein eigen Herz zu überwinden, da er um der Priesterin willen den ganzen Haufen seines Kriegspolks hat ziehen lassen.

170. Derhalben ist es in unserer Natur nicht, daß sich ein Mensch selbst im Glück und guten Tagen überwinden könnte. Themistokles und viele andere große Leute haben es auch nicht gekonnt. Ich weiß nur den einigen Scipio in allen Historien der Alten, der diesen Ruhm gehabt; die Andern alle sind des Glückes halben aufgeblasen und stolz geworden.

171. Joseph aber wird zum Regenten gesetzt über ein schön Land, darin es allenthalben sehr wohl gestanden; und ist seine Gewalt und Majestät groß gewesen, sintemal er die Verwaltung des Regiments gehabt, in Kriegsläuften und auch im Frieden, er hat Gesetze und Ordnungen gemacht, hat Strafen verordnet über die, so gesündigt hatten, und hat dieselbigen auch mit großem Ernste vollzogen. Das ist eine treffliche Ehre und sehr große Gewalt gewesen, und hat doch dabei solches Maß gehalten, ist so demüthig, wie er vor Zeiten zu Hebron bei seinem Vater auch gewesen, und weiß sich deß noch fein zu erinnern, daß er eines Schafhirten Sohn ist, gedenkt noch immer seines Elendes, seiner schweren Dienstbarkeit und Gefängnisses, gedenkt an GOttes Wort, welches lehrt, daß wir in Sünden empfangen und geboren sind. Darum behält er auch immer GOttesfurcht und die rechte Demuth; wie er auch zuvor im Gefängnisse und schwerer Dienstbarkeit gehabt hat. Er ist ja wohl etwas fröhlicher und dankt GOtt für die Erlösung, aber alles mit Ehrerbietung und GOttesfurcht. Wir werden hernach sehen, wie er sich so freundlich gegen seine Brüder gehalten, mit ihren geredet und geschätzet hat; denn der große Fürst hat sich dessen nicht geschämt, daß er der Schafhirten oder Kuhhirten Bruder ist.

172. Siehe aber unsern Stolz und Hoffahrt an: wer bei uns mehr Geld und Gut hat, wer von höherem Stande oder Geschlechte ist denn Andere, der wollte gern den Kopf gar in den Himmel heben. Also überheben sich die Fürsten über alle Maßen wegen ihrer großen Gewalt, scheuen sich gar nicht, weder vor GOtt noch vor Menschen. Also auch, die etwas Verstand vor Andern haben, gelehrt sind, Theologen, Juristen, Poeten, die sich mit ihren Schriften einen Namen gemacht, lassen sich dünken, sie seien so groß, daß sie auch weit über den Himmel reichen können. Warum sehen wir aber das nicht an, wo wir hergekommen sind und wie es sonst um uns bestellt sei? Denn sind wir nicht alle von Weibern geboren und mit Weinen und Schreien auf diese Welt gekommen? wie im Buche der Weisheit Cap. 7, V. 3. 4. 5. 6. steht: „Ich habe auch, da ich geboren war, Odem geholt aus der gemeinen Luft; und bin auch gefallen aufs Erdreich, das uns alle gleich trägt; und Weinen ist auch, gleichwie der andern, meine erste Stimme gewesen; und bin in den Windeln auferzogen mit Sorgen. Denn es hat kein König einen andern Anfang seiner Geburt; sondern sie haben alle einerlei Eingang in das Leben, und gleichen Ausgang.“

173. Derhalben sollten wir immer an unsere eigene Niedrigkeit, Elend und Jammer denken, und auf unsere Geburt sehen, daß wir von Natur alle Kinder der Vermaledeyung, der Sünde, des Todes und Teufels, und der ewigen Verdammniß schuldig sind. Warum sind wir denn noch so stolz? Wenn Einer allen Jammer dieses Lebens auf Erden, und den Stolz, so die Menschen gegen einander üben, mit einander vergleichen würde, derselbige würde nothwendig die Unsinnigkeit und Thorheit der hoffährtigen Menschen verlachen müssen, die man billiger in Ketten und Eisen spannen sollte, denn sonst die andern armen Leute, die nicht wohl bei Sinnen sind. Denn dieselbigen hoffährtigen heillosen Leute sind auch, als wären sie gemeinen Witzes und Vernunft beraubt, daß sie nicht sehen können, woher sie ihren Ursprung haben, wer sie sind und was sie endlich für ein Ende haben werden. Sie bedenken nicht, daß sie auch zu allerlei Jammer geboren sind, und daß sie aufwachsen werden sterben müssen, und daß sie fast alle Augenblicke in Gefahr des Todes stehen um der Sünde willen, und auch von wegen Satans, der da erschrecklicher Weise wider das menschliche Geschlecht tobt und wüthet.

174. Darum ist es offenbar, daß wir mit dem Teufelsgifte der Hoffahrt also vergiftet sind, daß wir vielmal unsinniger sind, als sonst die armen Leute, so ihrer Sinne beraubt sind. Denn dieselbigen sind kaum die Schatten dieser wahnsinnigen und rasenden Leute. Wenn du aber einen stolzen ruhmredigen Thrasio, der in einem köstlichen Kleide und mit einer goldenen Kette einherginge, der von seiner Gewalt, seiner Schönheit, und seinem großen Gelde und Gut zu rühmen wüßte, ansehen wolltest, würdest du denken, daß solches alles eitel Larven seien, damit der allerbösartigste, das ist, der unsinnigste Mensch bekleidet ist, nicht anders, denn als wenn du einen unsinnigen Menschen sähest, dem Hände und Füße gebunden wären. Denn wo der stolze Tropf bei Sinnen wäre, würde er bei sich selbst bedenken, woher er doch gekommen wäre, wie es doch um ihn stünde und was es zuletzt für ein Ende mit ihm nehmen werde. Denn das kommt oftmals vor, daß wer heute gesund und stark ist, in Wohlust und großer Herrlichkeit lebt, morgen etwa plötzlich von Krankheit befallen werden kann, daß er bald den Tod zu begehren gezwungen wird.

175. Und das kommt auch noch dazu, daß wir uns mit solcher Hoffahrt GOttes Zorn auf den Hals laden und damit seine Gnade und Gunst verlieren. Darum ist niemand ärmer und unseliger denn eben ein hoffährtiger Mensch. Denn er kann GOtt nicht anrufen oder ihn vertrauen; denn er weiß nicht, daß er der Sünde und dem Tod unterworfen ist; nennt auch GOtt nicht seinen Vater, sondern schließt sich selbst von solcher Gemeinschaft aus und selbstwillig beschließt er sich wiederum in des Teufels Dienstbarkeit und Gehorsam; ja über das, was das Allerschwerlichste ist, erkennt er diesen seinen Jammer nicht, fühlt oder versteht ihn auch nicht, sondern tritt einher aufgeblasen und stolz, gleich als wäre er ein GOtt, den man anbeten müßte. Wahrlich, solche Leute führen sich selbst mit ihrer Unsinnigkeit und Thorheit in großes Verderben. Darum sollte uns der Blindheit solcher armen Leute billig jammern, und sollten beweinen,
\quellseite{361}
daß sie so gar übel daran sind, und ihre Sache deshalb so übel steht.

176. Was soll man aber dazu thun, dieweil aller Menschen Natur zur Hoffahrt geneigt ist? Und pflegt sich gemeiniglich solche Hoffahrt in uns zu mehren, wenn wir etwa im Regimente Ehre und Gewalt haben, oder sonst in andern Ständen zu sonderlicher Herrlichkeit oder Hoheit gekommen sind. Soll man aber darum das Amt der Regierung fliehen und sich weigern dasselbe anzunehmen, wenn Einer dazu ordentlich würde berufen werden? Antwort: Mitnichten. Denn Joseph hat große Herrlichkeit und sonst allerlei Ehre, so das Regiment mitbringt, im Ueberfluß gehabt: er hat müssen auf des Königs Wagen fahren, es ist ihm auch königliche Ehre erzeigt worden, endlich hat ihm auch Pharao das königliche Siegel gegeben. Und dies Leben ist allerdings also beschaffen, daß es Fürsten, Könige, Edelleute, Obrigkeiten, gelehrte Leute in der heiligen Schrift, in guten ehrlichen Künsten, in der Philosophie haben muß, und kann derselben nicht entrathen; wie Egyptenland einen Fürsten hat haben müssen, wie Joseph war.

177. Aber du sollst dich befleißigen, daß Josephs Exempel und seine Sitten oder Leben dir immer vor Augen seien und du dasselbe im Herzen habest, und an die Furcht Gottes denkest, und daß du ihn vor Augen habest und auch dein selbst nicht vergessest: sondern bedenke bei dir selbst, wer du seiest, was du gewesen und was du noch endlich wiederum werden mußt. Denn es hat nie Einer so große Herrlichkeit oder Ehre gehabt, welchen gleichwohl nicht endlich die Würmer verzehrt und gefressen hätten. „Was erhebet sich denn“, sagt Sirach am 10. Cap. V. 9., „die arme Erde und Asche“, du Kind der Vermaledeyung und des Zorns? Ja, sagst du, ich bin gleichwohl von edlem und großem Stamm geboren, bin ein Doctor der Rechte, bin ein großer Philosoph. Es ist recht und gut, aber dieser Gaben aller soll man nun gebrauchen und nicht daraus Genuß ziehen; wie Augustin einen Unterschied macht, und wie der heilige Paulus Ap. Gesch. 13, 36. von David sagt: „Da er zu seiner Zeit gedienet hatte dem Willen Gottes“ \&c.

178. Also ist Joseph nicht hoffärtig gewesen von wegen der königlichen Pracht und Ehre, und daß er eine solche hohe Person, Fürst und Regent über Egypten war; sondern er diente dem Willen Gottes. Welches ist aber der Wille Gottes? Das ist er, daß zu Erhaltung dieses Lebens vonnöthen sei, daß man Obrigkeit, Doctoren, gelehrte Leute und Fürsten habe, auf daß man damit dem Muthwillen des gemeinen Pöbels wehren könne, die Armen zu schützen und die bösen Buben zu strafen. Denn wo es ohne diese Stände wäre, würde kein Friede und Ruhe sein können. Darum ist es Gottes Wille, daß wir durch solchen Unterschied der Stände im weltlichen Regiment ein stilles und ruhiges Leben führen, einander nicht beleidigen, nicht tödten oder berauben.

179. Die Obrigkeit ist gleichsam eine Ruthe Gottes, wie sie Paulus Röm. 13, 4. „Gottes Dienerin“ nennt, „eine Rächerin zur Strafe über den, der Böses thut.“ In diesem Stande soll man aber also leben, daß du immer denkest an Gottes Furcht und dich vor ihm demüthigest, und bedenfest, daß du ein Kind Adams, der Sünde, des Todes und Verdammniß seiest, sowohl als andere Menschen auch, wie Mose in seinem 90. Psalm lehret; doch also, daß du Gott dienest, der dich zum Regimente berufen, dadurch Friede in Städten und Landen erhalten, und die Diebe, Mörder und Todtschläger gebändigt und gestraft werden mögen; desgleichen, daß man auch Kriege führen soll wider die Feinde, welche den weltlichen Frieden und Ruhe verwirren und hindern. Und dies alles soll geschehen mit Demuth und herzlichem Anrufen oder Gebet zu Gott.

180. Und soll ein frommer Regent also sagen: Herr Gott, ich bin eines armen Kuhhirten Sohn, und ein Kind der Sünde; dieweil du es aber also haben willst, daß die Welt soll regiert, Friede erhalten, und die unordentlichen, unruhigen, bösen Buben gestraft werden, und beruffest nun mich zu demselbigen Amt, so will ich dir gern folgen. Ich wollte zwar lieber ohne Amt und eine Privatperson sein: aber dieweil nach deinem Willen und Befehl gehorchen soll, so will ich mein befohlen Amt ausrichten in deiner Furcht und mit aller Demuth. Da bleibt die Demuth unter der schönen goldenen Kette, Krone, Ringen etc. Denn ich bin ja mir selbst kein Fürst, sondern dem Willen Gottes und den Menschen zu Dienste in ihrer Noth. Ich bedarf für mich keines Königreichs, keiner Kanzel, keiner Versammlung in der Kirche und keines Predigtamts. Denn ich kann wohl daheim bei mir selbst dasselbe lesen, was ich öffentlich in der Gemeinde lehre oder predige. Ja, warum predigst du denn? Antwort: Darum, daß es Gott also haben will, der mich zum Pfarrherrn oder zum Prediger seines Worts gesetzt hat; und darum, daß das Volk Unterricht und Lehre bedarf. Also wirst du wahrlich und recht in Gottesfurcht und in der Liebe zum Nächsten wandeln.

181. Aber jetzt wirst du der Leute wenige finden, die ihre Werke dermassen allein auf Gott richten. Denn sie sind jetzt alle nur sich selbst Fürsten und richten Alles nur auf ihre Lust und Eigennutz. Darum sind es rechte Harpyien, Raubvögel, ja, es sind Teufel, die da Fürsten und Väter des Vaterlandes sein sollten, dieweil sie sich so hoch erheben und aufgeblasen sind wegen ihrer Würde, oder Herrlichkeit, und um ihres hohen Stammes, und auch anderer innerlichen Gaben willen, welche zu weltlicher Regierung gehören und vonnöthen sind. Ich, sagen sie, bin von königlichem oder fürstlichem Stamm, bin von Adel, ich habe auch einen Kopf, ich bin ein großer Doctor, bin ein Jurist etc. Ja, solche Leute sind nur sich selbst gelehrt, sind sich allein Könige und Edelleute. Aber in ihre Titel, die sie so hoch rühmen, sollte man also schreiben: O wie bist du so über alle Maßen ein toller unsinniger Mensch! Denn sie sind durch ihre Ruhmredigkeit also verblendet, daß sie auch das nicht wissen, daß doch billig alle Menschen sehr wohl wissen sollten, nämlich, den Jammer, darin sie empfangen und geboren sind, darin sie auch noch leben, und schuldig sind des ewigen Todes und Zorns Gottes, der auf ihnen liegt; sehen auch nicht, wie schwerlich sie daran sündigen, daß sie ihrer hohen Person und Amtes also schändlich mißbrauchen, den armen Unterthanen oder Zuhörern zum Schaden und Verderben.

182. Darum ist kein elenderer, unseligerer Mensch auf Erden, denn ein hoffärtiger Doctor, König oder Fürst. Denn über den gemeinen Jammer, den wir alle unterworfen sind, verdammt sie auch die Person oder Amt, und macht sie schuldig vieles unzähligen Unglücks, damit sie doch dem Willen Gottes gehorchen und dienen sollten. Deshalb schließe ich also, daß die Welt nichts Anderes sei, als ein Kerker voller unsinniger Leute, die da blind sind, und weder den gemeinen noch ihren eigenen Jammer und Elend erkennen.

183. Wie ist es aber so eine große Tugend an großen vortrefflichen Männern, wenn sie sich fein demüthigen und der Hoffahrt erwehren können! Daß etwa ein Knecht demüthig ist, das ist nicht so auffällig, ist auch ein gering Ding: aber wenn ein König, ein Fürst, Edelmann, Doctor, oder sonst auch ein großer reicher Mann demüthig ist, das ist die schönste Sache, da leuchtet solche Demuth, wie Sonne und Mond scheinen. Wenn ein Regent, oder auch sonst Einer, der höheren Standes ist als Andere, bei sich selbst also denkt: Ach, wer bin ich denn? Warum wollte ich mich Andern vorziehen, denn allein, daß ich des Regiments halben eine schwere Last auf mir liegen habe? Es ist ja die Ehre und Herrlichkeit wohl größer als bei Andern, aber derhalben werde ich auch so viel größere Rechenschaft geben müssen von dem Amt, das ich verwaltet habe.

184. Aber wie viel sind deren wohl, die also gesinnt sind? Es ist wahrlich gar eine seltene Gabe, solche herzliche Demuth, damit man doch bei Gott und Menschen Segen und Ruhm erlangen kann; so dagegen den Hoffärtigen, sie seien gleich Tyrannen oder Doctoren, beide Gott und Menschen feind sind. Und läßt es sich ansehen, daß sie sich selbst freiwillig solche Feindschaft machen und dieselbige immer mehren; denn sie sind auf das alleräußerste toll und unsinnig.

185. Derhalben hast du an Joseph ein trefflich und merkwürdig Exempel rechter Demuth, nachdem er so hoch erhoben war. Denn wiewohl er äußerlich zu großer Würde, Ehre und Herrlichkeit gekommen, und solche Ehre immer ist gemehret worden, läßt es sich doch eigentlich ansehen, daß zugleich mit solcher Würde und Herrlichkeit an ihm auch die Demuth größer geworden sei. Darum hat er sich fleißig nach dem gehalten, dazu Sirach uns im 3. Cap. V. 20. ermahnt, da er sagt: „Je höher du bist, je mehr dich demüthige.“

186. Solche Demuth ist aber in der Welt an denen nicht, die hoch erhoben und hohe große Aemter haben; sondern an Abraham, Joseph, David, Josaphat, da findet man sie. Und läßt es sich zwar ansehen, daß gar kein Unterschied sei zwischen ihrem Amt und dem Regiment der Tyrannen; denn sie führen auch Krieg, strafen die bösen Buben und sind Werkzeuge des Zorns Gottes. Und sind doch gleichwohl im Herzen viel freundlicher und demüthiger, denn sonst jemand unter dem gemeinen Pöbel sein kann, die auch mit den Leuten weinen; wie sich Joseph ganz freundlich erzeigte gegen seine Brüder, und weinte, da er sie sieht, verzeiht ihnen dazu leiblich und gern, da sie sich an ihm so hart versündigt.

\quellseite{362}

187. Derhalben ist dies Exempel wie ein großes gewaltiges Geschöpf, so auf die erschreckliche Hoffahrt dieser Welt gerichtet ist, und macht alle Fürsten, Könige und Doctoren zu Schanden. Denn diese rechtschaffene Tugend an Joseph übertrifft die andern alle weit, damit er zuvor die Noth und Angst des Gefängnisses überwunden hat. Allhier leuchtet die Furcht Gottes, vor dem er sich scheuet, und den er zugleich auch geehrt hat, dazu sieht man auch an ihm Verachtung sein selbst und die Liebe gegen den Nächsten. Denn er weiß wohl, daß er sich noch schleppt mit dem Fleische, welches voller Sünde ist, dazu der Vernaledeiung und Zorne Gottes unterworfen, wie andere Menschen auch. Danach hat er die aller schwerste Last des Regiments über ganz Egyptenland beide geistlich und weltlich auf sich liegen gehabt, hat wohl gewußt, daß er davon würde müssen Rechenschaft geben. Dies alles hat er sehr wohl verstanden, und hat zwar dazu viel Ursachen gehabt, daß er sich also gedemüthigt; darauf wir nicht Acht geben, und wir bedenken nicht, daß ein Jeder auch werde Rechenschaft geben müssen seines Amts und wie er dasselbe verwaltet. Denn, wo wir daran gedächten, würden wir uns viel demüthiger halten und erzeigen, als wir sonst thun.

188. Es sind zwar die äußerlichen Hüllen nöthig, nämlich, daß man Richter, Obrigkeit, Lehrmeister, Doctoren und Fürsten haben muß: aber dieser Aemter soll man nur gebrauchen, dieselben recht zu verwalten, und nicht daraus Genuß ziehen. Denn du bist ja der Mann nicht, der eben von Andern müßte angebetet werden; sondern das ist Gottes Wille, daß dieses Leben also regiert und erhalten werden soll, auf daß die Werke des Teufels zerstört, und daß Friede und Zucht und Ehrbarkeit erhalten werde. Dies soll das Ende sein, dahin alles Regiment soll gerichtet werden. Denn Gott sieht gar keine Person an, wenn er uns die Seligkeit gibt. Warum schafft er denn nun und verordnet diese Stände und Stufen der menschlichen Gesellschaft also? Antwort: Er thut es darum, daß er dich unter der Last der demüthigen, und zur Erkenntniß deiner Schwachheit, Jammer und Elendes führen will. Aber desselbigen Rathes Gottes mißbrauchst du, und kehrst ihn gleichsam um und brauchst ihn gar zum Gegentheil, nämlich, zur Hoffahrt und Stolz: so es doch Gottes Wille ist, daß du unter solchen äußerlichen Hüllen der Ordnung Gottes und der Menschen zu ihrer Nothdurft dienen sollt; welcher Nothdurft der Menschen Gott damit hat wollen zu Hülfe kommen, und darum befohlen hat, daß man die Obrigkeit in Ehren halten soll; wie Röm. 13, 7. geschrieben steht: „Gebet Furcht, dem die Furcht gebühret, Ehre, dem die Ehre gebühret.“

189. Darum sollst du solche Aemter zu Gottes Ehre und gemeinem Nutz brauchen, und also denken: Herr Gott, ich wollte lieber ohne Amt etwa in der Stille verborgen sein, ohne Scepter, ohne Krone, Gold, Silber und ohne alle solche große Herrlichkeit und Ehre. Aber du hast mich ja zu diesem Stande berufen, hast haben wollen, daß ich über Andere regieren soll: darum will ich nun deinen Willen gehorchen, und will dem gemeinen Nutzen und nicht meiner eigenen Begierde oder Wollust dienen. Und denke ja immer fleißig daran, wie es eigentlich um alle Menschen stehe, wie gebrechlich und elend sie sind, und in welcher Gefährlichkeit du selbst, wie eine große Last du zu tragen habest an dem Amt, so dir zu verwalten befohlen ist, auf daß du dich demüthigest und erweckt werdest, Gott um Hülfe und Beistand anzurufen.

190. Wie dies treffliche Exempel der Demuth an Joseph in so großer Herrlichkeit und Hoheit lehrt, welches weit herrlicher ist, denn das Exempel seiner Geduld gewesen, so er im Kreuz und Leiden gehabt. Denn es ist viel schwerer gewesen, daß er nach der Errettung am Wort fest gehalten, als daß er solches im Unglück und Anfechtung gethan; welche er mit Geduld ertragen und die er sich nicht hat lassen überwinden, ist auch derhalben von Gott nicht abgefallen, und hat nicht mit Ungeduld, Zorn oder Murren wider Gott gesündigt. Aber dazu gehört noch ein viel größerer Muth, da er es sich selbst überwunden hat. Denn das ist die höchste und schönste Tugend, wie Salomo Spr. 16, 32. sagt: „Ein Geduldiger ist besser denn ein Starker, und der seines Muths Herr ist, denn der Städte gewinnet.“

191. Es ist in diesem Leben die Ehre, so man im Kriege erlangt, wenn man sich darin ehrlich und wohl hält, die allerhöchste; wie an Achilles, Alexander dem Großen, und Scipio die herrlichsten Sachen gewesen sind: aber Salomo hat denselbigen starken Helden den Joseph, David und Mose vorgezogen und ist auch der keiner mit Joseph zu vergleichen. Denn er ist seines Muths ein Herr durch die Demuth.

192. Darum soll man nun erstlich diese Exempel fleißig bedenken, darnach sich auch wissen zu erinnern, woher wir gekommen sind, und welchem Jammer und Elend wir alle insgemein unterworfen sind. Daselbst, sage ich, soll uns immer im Herzen liegen, woher wir gekommen sind, was wir gewesen und was wir wiederum werden müssen, und endlich welch ein schwer mühselig Amt ein Jeder in seinem Stand auf sich liegen habe. Also werden wir uns allemal an Demuth gewöhnen, und daß wir uns fein inne halten und nicht hoffärtig werden.

193. Aber du wirst fast an allen Menschen das finden, das Virgili sagt: Nescia mens hominum etc., das ist, daß sich die Leute der Hoffahrt nicht wohl erwehren können, wenn sie zu großen Aemtern, Glück und Ehre kommen. Und man muß doch dies Exempel den Leuten immer vorhalten, daß Joseph allhier mit der größten Würde und Herrlichkeit geehrt, und mit Gold, Edelsteinen und größter Gewalt geziert wird; welche Stücke alle die Menschen zum Ehrgeiz, Stolz und Hoffahrt zu reizen pflegen. Aber Joseph hat diese großen Ungeheuer, dadurch sonst die Hoffahrt der Menschen angezündet wird, bezwungen können, welches weder Achilles noch Alexander jemals gethan haben.

194. Denn bedenke doch, wie er in so großen Ehren gehalten worden ist! Aller Augen und Herzen im Lande haben auf ihn sehen müssen, die ihn auch höchlich geehrt und viel von ihm gehalten haben. Er ist in Purpurkleidern und mit goldenen Ketten daher gegangen, ist auf des Königs Wagen geführt worden: und ist doch dabei also gesinnet gewesen: „Herr, dir will ich dienen und mich meines Nächsten in seiner Noth annehmen. Ich bin ja wohl ein Fürst, aber wenn ich es recht sagen soll, bin ich mehr ein Knecht aller Knechte.“

195. Und also soll ein jeglicher Fürst, Pfarrherr oder Kirchenbiener, Oberherr oder Regent schließen und denken. Was ist doch ein Schulmeister Anderes, denn ein Diener des Schülers? Dieser ist Herr, der Andere ist Knecht. Also will es Gott haben, daß wir in diesem Leben einander dienen und Gemeinschaft mit einander halten sollen. Wenn ich der Obrigkeit Ehre erzeige und Gehorsam leiste, so bin ich nicht eines Schafhirten Sohn, Fürsten oder Edelmann gehorsam, sondern dem Willen Gottes und gemeiner Nothdurft der Menschen.

196. Und dieweil die Ehrerbietung und Gehorsam gegen die Oberherrn unter dem gemeinen Pöbel ganz nöthig ist, müssen in allewege solche äußerliche Hüllen, Personen oder Aemter verordnet werden: und dieselbigen muß man äußerlich auch zieren und ehren mit Purpur, köstlicher Leinwand, königlichem Wagen, mit großen Titeln und dergleichen anderer Ehre mehr; nicht daß der, dem solche Ehre erzeigt wird, davon soll stolz und hoffärtig werden, als ob die Ehre sein eigen sei und ihm vornehmlich gebühre, sondern man soll an ihm Gottes Werk ansehen, und daß das Amt, so er verwaltet, auch Gottes ist.

197. Denn wo ich die Obrigkeit nicht werde ehren, so kann ich nicht leben, daß ich unverletzt und sicher bleiben möge vor den Mördern; ja, ich kann wohl vor einem bösen Nachbar nicht Frieden haben oder sicher sein. Darum ehre ich den Fürsten oder den Pfarrherrn um meines Nutzens willen, auf daß ich ein stilles, geruhiges und ehrliches Leben führen möge, daß ich die wahre Gottseligkeit und gute Künste lernen könne. Also muß man dies alles auf Gottes Willen ziehen und auf die gemeine Noth der Menschen. Die nun das bedenken, und sich fleißigen denselben also nachzufolgen, das werden gute löbliche Regenten und Lehrer werden beide im weltlichen Regimente und auch in der Kirche, wie Joseph einer gewesen.

\quellseite{363}

\section[Fünfter Theil]{Fünfter Theil. Von den übrigen Gnadenbezeugungen und Ehrentiteln, so Joseph von Pharao empfangen}

\begin{quote}
\textbf{V. 41—43. Und weiter sprach Pharao zu Joseph: Siehe, ich habe dich über ganz Egyptenland gesetzt. Und that seinen Ring von seiner Hand, und gab ihn Joseph an seine Hand, und kleidete ihn mit weißer Seide, und hing ihm eine güldene Kette an seinen Hals. Und ließ ihn auf seinem andern Wagen fahren, und ließ vor ihm her ausrufen: Der ist des Landes Vater. Und setzte ihn über ganz Egyptenland.}
\end{quote}

198. Pharao hat beiderlei Amt zusammengefaßt und begriffen, wie es Joseph hat verwalten sollen zur Zeit des Friedens und darnach auch in Kriegsläuften. Daheim oder im Lande, wenn es Friede gewesen, war er Baumeister und Bischof oder Aufseher über die Bäcker, Schenken, Fleischer und über das Getreide. Im Kriege hat er das Kriegsvolk geordnet oder Schlachtordnung gemacht, hat angezeigt, wie hin und wieder die Grenzen müßten bestellt und verwahrt werden, um die Macht und Gewalt der Feinde abzutreiben. Das war die königliche Gewalt, dazu ist ihm vonnöthen gewesen, daß er sonderliche Geschicklichkeit und Verstand hat haben müssen.

199. Ja, es hat auch der König seinen Ring Joseph über seine Finger gezogen, daß er Kanzler oder Secretär des Königreichs Egypten wäre. Derhalben sehen wir, daß der König in seinem Dinge höher gewesen ist als Joseph, allein des königlichen Stuhls halben, wie er droben gesagt hat. Was sonst das Andere belanget, da hat er die ganze Last des Landes und Regiments auf Joseph gelegt.

200. Darum ist er ein großer vortrefflicher Mann gewesen und ein rechter Heiland in Egypten. Und in andern zukünftigen Leben werden wir erst sehen, wie dieser große vortreffliche Mann, der so vielen Leuten nützlich und mit großem Lobe und Ehren vorgestanden, im Himmel mit großem Lichte und Herrlichkeit scheinen wird. Denn man muß ihn seinen Vätern Abraham, Isaak und Jacob weit vorziehen. Er hat allein das Amt der Regierung verwaltet mit Gottesfurcht, Demuth, und dazu mit großem Fleiße und Weisheit, dazu sonst viele Leute gehören, dieselbe Last zu tragen. Derhalben sollen wir nicht meinen, daß er viel werde müßig oder faul gewesen sein, sondern daß er in allen seinen Geschäften allezeit wacker und fleißig gewesen.

201. Das ist nun wahrlich ein wunderbarlich groß Lob, daß er zu solcher großen Herrlichkeit in Egyptenland ist erhoben worden, welche auch des Königes Gewalt gleich gewesen; und ist seine Ehre noch viel herrlicher gewesen als Davids oder der andern Könige Juda. Wahrlich, solche Kinder, wie Joseph und Daniel gewesen, die sind eigentlich ihrem Vater Abraham gleich, und können sich sein billiger rühmen, denn der andere gemeine Haufe der Juden.

202. Es hat ihm aber der König auch noch dazu eine güldene Kette und weiß seidene Kleid gegeben. Und es ist nun oftmals von der Ungleichheit der Kleidung, so in den Morgenländern, und bei den Völkern, die gegen Niederfang der Sonne wohnen, gebräuchlich gewesen, gesagt worden. Denn jene haben weiße Kleider und Gold getragen; diese aber Purpur und Gold. Und man sagt, daß sie jetzt noch im Egyptenlande solche Kleidung gebrauchen.

203. Darnach hat ihn Pharao auf dem andern Wagen nächst ihm fahren lassen und derselbige ist auch mit königlicher Herrlichkeit und Pomp geziert gewesen. Also ist Joseph ganz und gar mit Ehre und Gewalt des Königreichs Egypten geschmückt worden; wie von ihm geschrieben steht im Buche der Weisheit am 10. Cap. V. 13: „Bis daß sie ihm zubrachte das Scepter des Königreichs und Obrigkeit über die, so ihm Gewalt gethan hatten“ etc. Und wird noch sein Herz unter so vieler Anreizung der Hoffahrt und des Uebermuths dazu nicht bewogen, daß er sich darin über die Maßen erheben sollte. Das ist aber eine große Stärke eines großmüthigen und heldenhaften Gemüthes.

204. Endlich ist das auch noch dazu gekommen, daß er hat das Volk vor ihm lassen ausrufen: Der ist des Landes Vater. Und wird ohne Zweifel solches ein feierlicher Aufzug gewesen sein, da Joseph also ist eingeweihet worden vor den Fürsten und Amtleuten des ganzen Egyptenlandes und alles Volks. Alle haben sie gesehen, daß der König ihm den Ring, das Scepter, die goldene Kette und das seidene Kleid gegeben, ihn dazu auf den königlichen Wagen sitzen heißen.

205. Hier aber lassen sich die Juden allzu lieblich träumen, als hätten sie auch gleiche Ehre und Herrlichkeit mit Joseph. Denn solcher Ehre begehren sie alle sehr; aber es ist vergeblich und umsonst. Denn sie begehren nur die äußerliche Hülle und Person des Regiments, nämlich dazu, daß sie der Heiden Blut vergießen möchten und dieselbigen in ihre Dienstbarkeit bringen: sie fragen nach dem Regimente an sich selbst nicht, achten auch nicht der Furcht Gottes und der Demuth, so in Joseph war; das heißt siebenmal toll und unsinnig sein, wie ich droben gesagt habe.

206. Es haben aber auch die Hebräer mit dem Worte abrech sich und uns ein Kreuz und Marter zugerichtet. Denn sie legen es auf mancherlei Weise aus. Der größte Theil hält es dafür, daß es Ein Wort sei, die Andern machen zwei Worte daraus, ab, das heißt, Vater, und rech so viel als, zart. Vater, wegen der Weisheit, und rech, um der Jugend willen. Der Rabbihaften Grammatik bin ich feind. Sie sagen auch, daß in der syrischen und chaldäischen Sprache das Wort rech eben so viel heiße, als bei den Lateinern das Wort rex, König. Aber wenn man also möchte fabuliren, so wollte ich noch leichter das Wort rech auf das deutsche Wort „reich“ ziehen. Darum will ich denen folgen, die es für Ein Wort nehmen, und es auslegen, daß es heiße, ein sehr zarter Vater des Vaterlandes.

207. Die Andern leiten es her vom Worte barach, das heißt segnen, oder die Kniee beugen. Dieweil es aber scheint, daß der Buchstabe a dawider sei, darum daß er der Buchstabe bemanthi, wie sie dieselben in der hebräischen Grammatik heißen, einer ist, so sagen sie, daß a anstatt des Buchstaben h gesetzt sei; welcher Wechsel oder Veränderung der Buchstaben bei den Hebräern gebräuchlich ist. Dieselben legen es nun also aus, daß Pharao jedermann geboten habe, daß sie vor dem Fürsten Joseph die Kniee beugen sollten; welche Ehre auch noch jetzt allen Fürsten erzeigt wird an allen Fürstenhöfen. Und der Papst hat derselbigen Ehre auch begehrt, so heftig, daß die vor ihm die Kniee nicht beugen, mit einem Stecken geschlagen werden, wenn der heiligste Vater auf seinem Stuhle über die Gassen getragen wird, welchem es traun nicht gezimt, mit Füßen auf der Erde zu gehen. Denn also haben sie aus allen Historien allerlei Ceremonien und königlich Gepränge genommen oder gezogen, und haben es dem verdammten Menschen zugeneigt.

208. Ich verstehe dies Wort also, daß es ein Ausruf gewesen sei, daß ihm das Volk zugerufen habe, wenn sie ihn mit lauter Stimme einen Vater des Vaterlandes geheißen und gerühmt haben. Wie jetzt zu unsern Zeiten, wo man dem Kaiser oder Fürsten huldigt, das Volk antwortet und sagt: „Gnädiger Herr“, damit zu bezeugen und anzuzeigen, daß sie ihn für ihren Herrn erkennen und halten wollen, ihm treu und hold zu sein. Also werden die Fürsten und Amtleute sammt allen Unterthanen in Egypten, da Joseph ist eingeweihet worden, auch gerufen haben: abrech, das ist, gnädiger Herr etc.

\begin{quote}
\textbf{V. 44. Und Pharao sprach zu Joseph: Ich bin Pharao; ohne deinen Willen soll niemand seine Hand oder seinen Fuß regen in ganz Egyptenland.}
\end{quote}

209. Dies ist durch eine Figur geredet, die da heißt periphrasis, als wenn man ein Ding mit viel Worten ausdrücken oder beschreiben muß, da man es mit Einem Wort nicht eigentlich ausdrücken kann. Denn äußerlich nach gemeiner Art des Hausregiments kann das nicht verboten werden, daß nicht beide Hände und Füße sollten geregt werden, wenn man die Hausarbeit ausrichten soll. Darum muß dies der Meinung nicht verstanden werden, wie äußerlich die Worte lauten. Denn Füße aufheben oder regen, heißt in die Höhe steigen; und er will allhier so viel sagen und zu verstehen geben: Es soll niemand im Lande ohne deinen Willen und Ordnung Richter, Fürst oder Obrigkeit sein, sie sollen von dir das Recht empfangen, darnach sie regieren sollen. Alle, so im Regiment sind, daß sie Andern vorstehen. Du sollst Obrigkeiten und Superintendenten oder Aufseher bestellen und setzen hin und wieder über ganz Egyptenland, und soll sich niemand ein Amt oder Regiment anmaßen, dasselbe zu verwalten, ohne deinen Willen. Also muß man diese Figur oder Weise zu reden auf das bürgerliche Leben beziehen, und sie begreift in sich die größte Gewalt und Herrlichkeit, welche Joseph ist gegeben worden.

\quellseite{364}

\begin{quote}
\textbf{V. 45. Und (Pharao) nannte ihn den heimlichen Rath.}
\end{quote}

210. Nun ist das ein ander Kreuz, damit sich die Hebräisten martern; denn sie wissen nicht, ob dies ein hebräisch oder egyptisch Wort sei. Hieronymus disputirt davon, daß es auf Hebräisch heiße, einen, der erfunden oder errathen könne, was heimlich, dunkel oder verborgen ist. Aber ich sehe keine rechte Ursache dieser Auslegung, vornehmlich in dem Wort paneach. Denn der andere Name scheint eher diesen Sinn zu ergeben. Es sind aber viel Wörter, die vielen Sprachen gemein sind. Wie die Hebräer, die Lateiner, Griechen und Deutschen das Wort „Sack“ behalten. Also heißt nun auf Deutsch „Nein; fenestra, Fenster. Denn die Sprachen sind nicht so gar von einander gerissen, daß nicht bisweilen etliche einander gleich seien und sich nicht mit einander reimen sollten. Darum halte ich dafür, daß dieser Name bei vielen Völkern gebräuchlich gewesen sei. Ich kann aber davon nicht urtheilen, zu welcher Sprache er eigentlich gehöre, zur egyptischen oder zur hebräischen, und die Grammatiker heben hierin auch nichts Gewisses. Aber Hieronymus ist der Meinung, daß es egyptische Worte seien, welche so viel heißen, als bei den Lateinern salvator mundi, Heiland der Welt, und zwar, weil Joseph die Welt diesmal von dem bevorstehenden Untergang, so aus der Theurung entstehen können, errettet. Ich bekenne aber mit Lyra einfach, daß ich es nicht weiß. Wiewohl daran kein Zweifel ist, daß es ein solch Wort sei, welches eine sonderliche Würde, oder ein hohes Amt, oder aber ein Verdienst und Belohnung wegen der Deutung der Träume und sonderlicher Weisheit, so Joseph gehabt, bedeutet.

211. So hat nun Mose bis daher beschrieben, wie Joseph sei erhöhet worden über ganz Egyptenland, so daß ihm nicht allein die Bürger, Bauern, und andere gemeine Unterthanen, so Privatpersonen und ohne Amt gewesen, sondern auch die Fürsten und Amtleute sind unterworfen gewesen. Er selbst aber hat niemand zu gehorchen oder unterthänig zu sein brauchen, denn nur dem Könige. Und ist wahrlich ein wunderbarlich Ding, wie er doch allein solche große Dinge habe ausrichten können. Er wird ohne Zweifel etliche verständige geschickte Männer erwählet haben, die guten Rath haben geben und ihm helfen können. Denn es kommt oftmals vor, daß die geringsten Herrschaften vieler Andern Hülfe und Rath bedürfen; wie in den Dörfern die Richter die Schöppen pflegen zu sich zu nehmen. Wie viel mehr ist vieler Leute Hülfe vonnöthen gewesen zur Verwaltung des ganzen Regiments in so einem großen Königreich! Also hat Mose, da er im Volk Israel das Amt eines Königs verwaltet, auch andere Fürsten und Räthe zu sich genommen, item Hauptleute über tausend, und sonst Hauptleute über zehn und dergleichen Andere mehr. Denn es ist Mühe und Arbeit genug, wenn Einer allein nur über zehn Andere regieren und ihnen vorstehen soll.

212. Es läßt sich aber mit diesem Text ansehen, als sollte derselbe etlichermaßen den Ehrgeiz und die Hoffart beschönigen und sie verteidigen, darum daß Joseph so leichtlich und willig ohne alle Weigerung den Ring und goldene Ketten angenommen, und dazu das schwerste Amt, beide die Obrigkeiten und Unterthanen des ganzen Königreichs Egypten zu regieren, auch damit auf sich genommen, also daß niemand im ganzen Lande ohne seinen Willen und Befehl etwas hat thun mögen: das, sage ich, scheint viel zu hoffärtig und ehrgeizig. Denn warum hat er solches nicht auch mit gleicher Demuth abgeschlagen und sich desselben geweigert, wie er droben zum Könige sagt, da er ihm seinen Traum sollte deuten: „Das stehet bei mir nicht, Gott wird doch Pharao Gutes weisen“?

213. Darauf antworte ich also: Es ist zuvor gesagt, daß man in allen Regimenten vornehmlich auf Gottes Ordnung und Willen sehen soll, und wird an diesem Ort uns eben dieselbige Lehre auch also vorgehalten und befohlen, uns darnach zu richten, nämlich, daß GOtt aus sonderlichem Rath hat haben wollen, daß in diesem Leben solche äußerliche Hüllen oder Personen, das ist, unterschiedliche Stände und Aemter unter den Menschen sein sollen, welche bürgerlicher Weise mit einander Gemeinschaft haben.

214. Denn es ist vonnöthen, daß die Obrigkeit, Könige und Fürsten mit sonderlichem Scheine, mit äußerlichem Gepränge und Schmuck umgeben und gezieret seien. Also, wenn dem Richter die Stadtdiener nachfolgen, das hat wohl einen Schein der Hoheit und Ehrgeizes; aber Gott will haben, daß solcher Unterschied der Stände sein solle. Also trägt der König Scepter und Krone, der Adel hat auch seinen Schmuck und Wappen vor Andern. Denn solche Stände sind durch Gottes Willen und Ordnung eingesetzt. Wie wir in den hohen Schulen und derselben Ursache willen auch die Grade behalten, daß man etliche zu Magistren und Doctorn macht; wiewohl in Wahrheit solches an sich schlecht und gering Ding ist: aber man muß damit den Leuten in der Noth dienen, welche solchen Unterschied der Personen fordert und haben will. Denn sonst würden die Unterthanen die Obrigkeit verachten und mit Füßen treten, wenn sie nicht des Standes und der Würde oder Herrlichkeit halben etwas höher wären.

215. Darum sagt Sirach recht: Welcher unter den Brüdern den Andern vorsteht und ihr Oberster ist, soll auch die größte Ehre haben. Denn unser Herr GOtt will haben, daß die Regenten sollen in Ehren gehalten werden, wie Paulus lehrt Röm. 13, 7. Er will die Obrigkeit unversehrt haben. Denn der Teufel, der ein Fürst dieser Welt ist, untersteht sich die Königreiche, Regiment und Haushaltungen zu verwirren und gar umzukehren, oder zu zerstören. Und wir würden nicht können unverletzet bleiben und sicher sein vor Dieben und Mördern, wenn man ohne die Regenten wäre. Es würden alle Städte und Geschlechter mit Hurerei, Ehebruch, Dieberei und Mord erfüllt werden. Denn es kann so noch schwerlich Zucht erhalten werden, da man doch Obrigkeit hat, daß nicht bisweilen böse unruhige Menschen unsern Leib und Gut angreifen, uns daran zu beschädigen. Derhalben müssen solche böse Buben mit Feuer, Schwert, Rad, Gefängniß und allerlei andern Strafen gezwungen und gestraft werden, durch welche Strafen der Teufel in den ungehorsamen und unbändigen Leuten zurückgetrieben wird. Denn sonst würden bald alle Laster mit Gewalt überhand nehmen, daß die Leute ungestraft würden thun wollen alles, was sie nur gelüstet.

216. Derhalben müssen die Personen und Aemter immer bleiben, davon St. Paulus sagt Röm. 13, 3. 4.: „Willst du dich nicht fürchten vor der Obrigkeit, so thue Gutes; so wirst du Lob von derselbigen haben; denn sie ist Gottes Dienerin dir zu gut. Thust du aber Böses, so fürchte dich; denn sie trägt das Schwerdt nicht umsonst.“ Es sind äußerliche Hüllen oder Aemter des Zorns wider die Bösen und des Trostes den Frommen, und denen, so unterdrückt oder beschweret werden, zu gut.

217. Das hat Joseph sehr wohl verstanden. Denn er ist ein Mann voll heiligen Geistes, Darum hat er das Regiment über ganz Egyptenland in Gottesfurcht und Demuth angenommen, und hat in demselben Amt viel mehr Arbeit gehabt, als zuvor in der Knechtschaft und Gefängniß; wie die Heiden auch geklagt haben. Denn davon hat man noch den Spruch des Bias, welcher einer der sieben Weisen der Griechen gewesen: Magistratus ostendit virum, das ist: Ob Einem ein obrigkeitlich Amt befohlen wird, so weiset es sich mit ihm aus, was er für ein Mann ist. Und Cicero sagt auch von sich: O! me falso, quod nunquam fui, sapientem habitum: O! wie bin ich so unwürdig für einen weisen Mann gehalten worden, daß ich doch nie gewesen. Und Demosthenes soll auch gesagt haben: Wo ihm zwo Wege würden vorgestellt werden, da einer zum Tode und der andere zum Regimente führete, so wollte er den Weg gehen, welcher zum Tode führete.

218. Also hat sich Joseph in diesem Regiment gegen GOtt ehrerbietig und demüthig gehalten. Denn er hat wohl gewußt, daß ihm diese Last um Gottes willen auferlegt war. Er wird auch oft gewünscht haben, er wollte lieber im Gefängniß, oder aber daheim bei seines Vaters Heerde oder Vieh sein und Ziegenmilch essen, als in solchem hohen Stande, dazu er gekommen war, wo er wohl und glücklich regieret hat. So ist er ohne Zweifel also gesinnet gewesen, daß er gedacht und gesagt hat: Ich wollte lieber, daß dies Amt einem Andern befohlen würde und daß ich für mich selbst allein in der Stille leben möchte. Denn wer bei seinem Amt sorgfältig ist, wie Paulus ermahnt Röm. 12, 8., also daß er mit Ernst die bösen Buben straft, und über dem Recht, guter Ordnung, Zucht und Gerechtigkeit hält, der wird manniglich die bösen unruhigen Menschen, wie auch den Teufel reizen und auf sich laden. Darum ist es ein männlich Werk und eine große Tugend, daß man sich daran nicht bald ärgere und kleinmüthig werde, wenn es kommt, daß sich der Teufel mit ganzer Gewalt wider uns legt.

\quellseite{365}

219. Ja, es ist auch im Hausregiment, da man das Gesinde, Weib und Kind regieren soll, vonnöthen, daß Einer ein männlich Herz habe und unverzagt sei. Darum haben alle Weltweisen geklagt über die Beschwerung, so im Regiment ist; und daher pflegen auch die Tyrannen zu kommen. Wenn sie sehen, daß ihre Rathschläge und ihr Thun, das alles sehr fein verordnet, keinen Fortgang oder Glück haben, oder aber daß ihnen Andere Widerstand thun, so werden sie gar toll und unsinnig, und werden aus frommen Fürsten Tyrannen, die darnach mit Gewalt und anderer Leute Schaden, von denen sie meinen, daß sie ihnen im Wege stehen, hindurch zu brechen und damit ihre Gewalt also zu erhalten suchen. Denn es sind nicht tapfere Helden, die sich selbst zwingen könnten, sondern hängen und folgen ihrem Zorn und Begierden nach. Also ist das Regiment Josephs nicht gewesen.

220. Ich gedenke noch wohl, daß D. Staupitz von seiner Regierung sagte, er hätte in den ersten drei Jahren nach der Strenge regieren wollen, es habe ihm aber also nicht geraten wollen. Im andern Jahr habe er regiert nach den Gesetzen, nach Rath der Vorfahren und Aeltesten, dasselbe habe ihm auch nicht wollen geraten. Da habe er im dritten Jahr angefangen zu regieren nach Gottes Willen, habe dabei Gott angerufen und gebetet; und gleichwohl sei dabei auch nicht allezeit Glück gewesen. Da sagt er nun zuletzt: Ich that, was ich thun konnte, da ich an alle diesen Rath und Weise verzagt hatte. Denn wiewohl Gott haben will, daß man ihn soll anrufen, erhört er uns doch eben nicht nach unsern Wünschen und Gedanken, in welchen wir ihm sehr klugen und schönen Rath vorzuschreiben pflegen.

221. Darum sollen wir thun, was wir können, und wenn Hindernisse eintreten, sollen wir uns zur Geduld schicken, und uns an das gemeine Sprüchwort halten: Mitte vadere, sicut vadit, quoniam vult vadere, sicut vadit: Laß gehen, wie es geht; denn es will doch so gehen, wie es geht. Und wie man im deutschen Sprüchwort sagt: Wer den schweren Stein nicht heben kann, der lasse ihn liegen. Denn es ist nie ein Regiment ohne Klage und größte Beschwerung gewesen. Und Joseph wird ohne Zweifel oftmals gewünscht haben, daß er nur den goldenen Wagen möchte fahren lassen und daheim in der Stille unter den Schafhirten sein Wesen haben. Denn der Teufel haßt und zerstört gern alle Regimente und Obrigkeit, welche auch die Werke des Fürsten dieser Welt verbietet, und Todtschlag, unrechte Gewalt und Diebstahl etc. straft. Dieweil legt sich der Teufel mit ganzer Macht: Du sollst mir nicht wehren, sagt er; denn das ist mein Amt, daß ich Alles unter einander menge und über einen Haufen werfe mit Haß, Neid, Hurerei und Ehebruch. Dieweil du aber mir in solchem Amt widerstehest, sollst du wissen, daß ich mich für und für auch wider dich setzen wolle. Da werden zwar die Regenten, welche kleinmüthig sind, bald verzagt, schlafen auf beiden Ohren, werden darüber sicher, faul und nachlässig, lassen alle Dinge gehen, wie es geht, lassen alle Sorge liegen, so sie im Regiment billig haben sollten. Aber dieselbigen sündigen damit schwerlich wider die Lehre, welche Salomo in seinem Prediger gegeben hat.

222. Sie sagen vom Herzog Friedrich, dem Churfürsten von Sachsen, daß er auch endlich wegen der großen Last und vielen Sorge müde geworden, und gesagt habe, wie er oftmals allerlei Stände der Menschen betrachtet, und gefragt, welches doch das allerbeste und seligste Leben sein möchte im menschlichen Geschlechte, darnach auch zuvor alle Gelehrten und Philosophen mit großem Fleiß geforscht haben. Er soll aber davon eben also geurtheilt haben, wie Virgili in seinem Buche Georgica, da er sagt: O fortunatos nimium, sua si bona norint, agricola[s] etc.: Die wir wären doch die Landleute überglücklich, wenn sie ihre eigenen Güter kennten. Denn er sagte, daß fromme Fürsten und Oberherren zumal ein armseliges Leben und Stand führten: aber die mit dem Ackerbau umgehen, dieselben hätten das glücklichste Leben und gute Tage, darum daß sie sein sicher und in der Stille unter dem Schutze ihrer Herren mit Weib und Kind lebten.

223. Denn die Obrigkeit ist allenthalben und ist doch nirgends, wenn sie ihr Amt recht und treulich ausrichten will. Wer aber nur sich selbst zu Ruhe regieret, der wird nimmermehr ehrlich oder glücklich regieren. Denn es ist nicht genug, eine goldene Kette oder Purpurkleid tragen und sich auf Wagen fahren lassen, oder auf großen Rossen daher reiten, daß du nur deiner eigenen Lust oder Begierden folgen und nachhängen mögest; sondern es gehört dazu, daß du deinen Unterthanen und dem gemeinen Nutzen dienen sollst.

224. Wenn aber Andere sehen, daß die Last größer ist, das Regiment zu verwalten, als die Ehre, so man davon hat, fliehen sie vor solcher Mühe, wiewohl sie wissen, daß sie zu demselben Stande geboren und von Gott dazu berufen sind; wie es die Edelleute und die Fürsten sind. Denn warum, sagen sie, soll ich so große Arbeit vergeblich auf mich nehmen, da ich nicht weiß, wie es gerathen möchte? Ich will lieber einsam leben und geruhige Tage haben, will eine Kappe anziehen etc. Aber siehe, was daraus folget und was die Mönche mit solcher Flucht ausgerichtet; und ist ein wahres Wort, so man gemeiniglich zu sagen pflegt: Desperatio facit monachum: Die Verzweiflung bringt die Leute dazu, daß sie Mönche werden. Denn sie sind nicht allein um der Sünde willen verzweifelt, sondern auch der Haushaltung und weltlichen Regiments halben, daß sie sich damit nicht haben bemühen wollen.

225. Also sehen wir, daß deren viele sind, die im Ehestande nur auf fleischliche Wollust sehen, lassen sich träumen und warten auch darauf, daß in demselbigen Stande allein eine Lust sein soll. Wenn es ihnen aber nicht also gelingt, solche Lust zu haben, sondern täglich allerlei Trübsal vorfällt, da fangen sie an, über den Ehestand zu klagen, denselbigen zu verfluchen und sagen, daß sie jämmerlich betrogen seien, sie hätten traun nicht gehofft, daß es ihnen also gehen würde.

226. Wider solche Aergernisse, so in allen Ständen vorzufallen pflegen, lehrt die Schrift, daß man keineswegs davon fliehen soll, so man zum Regimente berufen wird, wie Sanct Paulus sagt 1 Tim. 3, 1.: „So jemand ein Bischofsamt begehret, der begehret ein köstlich Werk“, das ist, welches müße und Gott angenehm ist. Dir mag es vielleicht beschwerlich, unlustig und nicht gut sein; aber es gefällt doch Gott wohl und ist der Kirche und ganzen Gemeinde in Landen und Städten sehr nütze. Nicht wie die Dompfaffen und Mönche nach Bischümern und Präbenden stehen, oder Pfarren an sich zu bringen suchen, welche des Teufels Rath sind, und in dem Stande nicht selig werden können, darin sie außerhalb Gottes Ordnung leben, nur ihrer Wollust mit Saufen und Fressen und anderer Uppigkeit mehr warten.

227. Darum soll man ein solch müßig, zart, Epicurisch Leben nicht führen, da man sich allerlei Regiment entzieht. Denn Gott sagt 1 Mos. 1, 28. also: „Füllet die Erde, und macht sie euch unterthan, und herrschet über Fische im Meer“ etc. Auf daß ich aber deine böse Lust und Ehrgeiz etwas im Zaume halten möge, soll dir der Acker Dornen und Disteln tragen, und „im Schweiß deines Angesichts sollst du dein Brod essen, bis daß du wieder zu Erden werdest, davon du genommen bist“ etc., Cap. 3, V. 18. 19.

228. Ja, wer kann aber die scharfen Dornen, Traurigkeit, allerlei Mühe und Arbeit also ertragen? da will ich lieber fliehen, sagen die Leute. Nein, du sollst davor nicht fliehen, sollst dich aber bald von Kindheit auf gewöhnen nicht an Wollust, sondern dem gemeinen Nutzen in diesem Leben zu dienen. Denn darum gehts zu Schulen, lernst gute Künste und übst dich darin, daß du dermahleinst der Gemeinde dienen sollst, es sei gleich in der Kirche oder aber in weltlicher Polizei: und dazu gehört ein gelehrter, geduldiger und starker Mann, der Muth habe. An solche Tugenden mußt du dein Herz gewöhnen und dazu erwecken. Willst du aber ein Mönchsleben erwählen und dasselbige lieber haben wollen, da wird dich der Satan gewißlich aus dem Paradiese kopfüber in die Hölle stoßen. Denn die Mönche sind nichts Anderes denn inutilia pondera terrae, eine unnütze Erblast, welche der Wittwen Häuser und der armen Waisen Gut und Habe fressen und verzehren; das ist ja kein göttlich, sondern ein teuflich Leben.

229. Derhalben sollst du bedenken, du seiest dazu geboren und von Gott berufen, daß du deinem Nächsten dienen sollst. Und ist zwar unter andern Strafen der Erbsünde nicht die geringste, ein Regent sein; es ist fast die andere Strafe nächst dem Tode. Ja, Etliche unter den Heiden, als Demosthenes, haben den Tod nicht für so gar bitter gehalten als die Verwaltung des Regiments. Und Augustinus, welcher ein sehr weiser und glückseliger Kaiser gewesen ist, hat auch darauf gedacht, wie er von dem Heidentum oder Kaisertum möchte frei werden und sich davon zurückziehen; die Andern aber, als Nero und Caligula, sind von teuflischer Unsinnigkeit gar zu Weibern geworden, und haben ihrer bösen Lust und Begierden den Zaum frei gelassen, und alles, was sie nur gelüstet, frei dahin gethan.

\quellseite{366}

230. Wir sind zwar von Natur und wegen der Erbsünde begierig, daß wir regieren mögen, und haben dieses Gift eingesoffen durch Eingebung des Satans im Paradeise, da er 1 Mos. 3, 5. sagt: „Ihr werdet sein wie Götter“, wenn ihr von dieser Frucht essen werdet. Da hat Gott wiederum auch gesagt: Wohlan, ich will dir deine Lüste büßen; dieweil du ja begehrt, Gott gleich zu sein, und daß du Andern mögst vorstehen und über sie herrschen, so will ich dir zu regieren genug geben.

231. Darum ist es wahrlich nicht eine leichte oder geringe, sondern eine harte und greuliche Strafe wider die Erbsünde, ein Regent sein und Andern vorstehen, auf daß wir lernen, wie groß die Bosheit der Menschen sei, die da regiert werden, und auch Satans. Denn dasselbe versteht oder glaubt niemand ohne allein die, so im Regimente sind, welchem der Teufel mit ganzer Macht widersteht. Und der Fürst Michael streitet nicht allein in der Kirche, sondern auch in weltlicher Polizei mit den bösen Engeln. Denn Räuberei, Diebstahl, Betrug, Todtschlag sind Werke des Teufels, welche zwar die Obrigkeit aus dem Herzen der Menschen nicht wegnehmen kann, jedoch straft sie mit dem Schwert und andern Strafen die äußerlichen Laster, wie sie am allerbesten kann und mag. Denn sie trägt darum aus göttlichem Befehl und von Gottes wegen das Schwerdt.

232. Darum sollen die, so im Regimente sind, sich immer wissen ihres Amtes zu erinnern und sich darin männlich halten. Denn darum hat sie Gott vor Andern geschmückt und geehrt mit sonderlicher Würde und Herrlichkeit, mit Purpur, Gold und anderem Schmuck, den Teufel und die bösen Buben damit zu schrecken. Denn die Erbsünde mühet und ist unsinnig, und der Teufel treibt die Menschen zu allerlei Sünde und Schande.

233. Ja, das ist es nun, daß Gott 1 Mos. 3, 22. sagt: „Siehe, Adam ist worden, als unser einer.“

234. Darum hat mich der Text in der Apostelgeschichte am 13. Cap. V. 36. oftmals getröstet, da von David gesagt wird, er habe dem Willen Gottes gedient, so hoch sein Regiment sehr erbärmlich gewesen, voller Aufruhr und allerlei Verwirrung beides im politischen und Hausregimente, also daß da kein großer Unterschied war zwischen seinem und dem Regimente des Herodes. Wie oft mögen wir meinen, daß er geseufzt und gesagt habe: O! wenn ich jetzt möchte daheim meines Vaters Schafe hüten, wie gern wollte ich dieses königliche Scepter und Krone einem Andern übergeben!

235. Das ist aber doch unser Trost, daß wir wissen, wir haben diese Strafen im Paradeise verdient, da wir uns der Gottheit haben gelüsten lassen. Derhalben heißt uns nun Gott, daß wir auch Götter und Regenten sein sollen; und wendet doch solche Strafe in Gnade und Barmherzigkeit, und nicht zum Zorne oder Ungnade, ja, er hält es für einen Gehorsam und angenehmen Dienst. Du, sagt er, sollst mir in dem Amte dienen, ich will haben, daß du also dem gemeinen Nutzen vorstehen sollst. Da wird dann das große Unglück und der gemeine Jammer etlichermaßen gelindert. Ja, es wird alles herrlich gezieret, sintemal wir des Willens Gottes können gewiß sein, daß er seine Ordnung im Amte der Obrigkeit für den angenehmsten Gehorsam annehmen wolle. Darum weiß ich nun auch also, daß ich dem Willen Gottes diene, wenn ich nach meinem Vermögen mein Amt ausrichte. Und wenn etwa mein Fleiß und Rath verhindert wird, so wird damit meine böse Lust im Zaume gehalten: wo mir es aber glückt und von Statten geht, so gereicht solches alles zu Gottes Ehren.

236. Dies sage ich um der fleischlichen Träume willen der Juden, welche nur die Lust und den äußerlichen Schein der Regenten hoch rühmen können, begehren auch, daß sie möchten auf goldenen Wagen fahren, und hören, daß das Volk vor ihnen her sie möchte ausrufen und ihnen frohlocken: aber die schwere Last, Mühe und Arbeit, so dazu gehört, bedenken sie nicht zu tragen und wollen sie auch nicht tragen. Nun muß aber wahrlich doch ein Regent ein rechtschaffener beherzter Mann sein, der mit großem starken Muthe die Beschwerung und Last des gemeinen Regiments tragen könne. Er muß saure Suppen essen, und in einen sauren Apfel beißen können, wie wir im Paradeise in einen süßen gebissen haben. Gott hat die Strafe, welches schwer genug ist, im Paradeise selbst verordnet, und wendet sie doch zu seinem Dienste, und will haben, daß es eine Arznei oder Hülfe sein soll wider die Sünde, daß damit also die Hoffahrt und der Ehrgeiz der Menschen getödtet werden möchte, nach welcher wir haben wollen Götter sein. Darum, sage ich, soll man darin keine Lust suchen, da man dem Willen Gottes dienen soll, sondern die Töd[t]ung unser selbst, göttliche Ordnung und daß damit dem Nächsten geholfen werde.

237. Und ich zweifle nicht, es werde die Ehre und Herrlichkeit den Joseph viel mehr geplagt haben, denn das Gefängniß und sonst sein Jammer und Elend. Darum ist er uns zum Exempel vorgestellet, daß wir ja nicht kleinmüthig und verzagt werden, wenn wir ein Regiment zu verwalten haben; und daß wir auch unserer eigenen Lust nicht nachhängen und vor dem Regimente fliehen: sondern wenn wir dazu berufen sind, daß wir alsdann unser Amt mit Beständigkeit ausrichten, es gehe uns gleich wohl oder übel dabei.

238. Und sollen die jungen Gesellen bald von Jugend auf sich in ihrem Herzen das also vornehmen und bei sich beschließen: Ich will fleißig studiren in der heiligen Schrift und andern guten Künsten, auf daß ich in meinem Amte Gott dienen möge. Denn ich bin auch der armen Sünder einer im Paradeise, darum muß ich auf meinem Acker auch Dornen und Disteln haben. Inzwischen aber will ich mit dieser Ehre zufrieden sein, darum daß ich in diesem Amte und Gehorsam der Gnade und des gnädigen Willens Gottes gegen mich gewiß bin.

\section[Sechster Theil]{Sechster Theil. Von Josephs Verheirathung, und wie er umzieht das Getreide zu sammeln}

\subsection{I}

239. Joseph war nun frei gegeben, und war sein selbst mächtig, und war dem ganzen Egyptenlande vorgesetzt zum Regenten, und hat doch ohne seines Königs Rath und Willen sich nicht in den Ehestand begeben wollen, auf daß er damit seinen Gehorsam beweisen möchte. Denn er hat sich selbst die höchste und königliche Würde nicht zuschreiben, oder sich dieselbe anmaßen wollen. Er mißbraucht auch seine Gewalt nicht, das Königreich unruhig zu machen, oder Pharao durch Aufruhr zu vertreiben, was bisweilen von unruhigen Leuten unter den Gottlosen zu geschehen pflegt, da es ihm doch an Gelegenheit nicht gefehlet hat, Tumulte oder Unruhe anzurichten. Denn er hatte das ganze Egyptenland in seiner Gewalt, und wird ohne Zweifel bei Vielen ein großes Ansehen und große Gunst gehabt haben, und haben ihm vielleicht Etliche das Scepter und die königliche Krone angeboten. Aber er bleibt fein in Gottesfurcht und gibt dem König seine gebührliche Ehre, scheut sich vor ihm mit Ehrerbietung; und dieweil er auch ohne des Königs Willen nicht thut, was ihm doch wohl wäre frei gewesen zu thun, wie sollte er denn ohne des Königs Befehl etwas Größeres anfangen haben?

240. Darnach wird hier nicht allein das angezeigt, wie sich Joseph mit höchstem Willen beschließet, dem Könige gehorsam zu sein; sondern es ist auch ein treffliches Anzeichen der Liebe und Gunst, so Pharao zu Joseph getragen, dieweil er ihn auf diese Weise sonderlich hat ehren wollen, nämlich, da er ihm eine ehrliche, tugendsame Jungfrau von gutem Geschlecht zur Ehe gibt. Denn sie ist eine Tochter eines großen vortrefflichen Mannes und Hohenpriesters unter andern Priestern gewesen, welche bei den Egyptern in großen Würden und Ehren gehalten worden sind; wie hernach wird zu sehen sein im 47. Cap. V. 22., da der Priester Land befreit wird, daß sie davon den fünften Theil nicht haben zur Schatzung geben dürfen.

\quellseite{367}

241. Darum hat er aus sonderlichem Rath und geneigtem Willen des Königs Pharao diese Jungfrau zum Weibe genommen, welche von hohem Geschlecht geboren war. Denn es hat sich der König dessen beflissen, daß er ihn nur möchte hoch ehren, wie er immer gekonnt hat, damit sein Ansehen und Gewalt viel größer und herrlicher wäre, auch um desjenigen seines Schwiegervaters willen, welcher im Königreich Egypten ein hoherpriester war.

242. Es haben aber fast alle, so dies Buch ausgelegt, an diesem Ort disputirt, ob dieser Potiphar derselbige Mann sei, welcher Josephs Herr gewesen und ihn von den Ismaeliten gekauft hätte? Die Juden und auch Hieronymus, welcher ihnen nachfolgt, bringen allhier große Lügen vor. Denn nachdem derselbige Oberste über die Fleischer und Josephs Herr droben, Cap. 39., auf Hebräisch genannt wird saris, das ist, ein Eunuch oder Verschnittener, fragen sie, wie es möglich sei, daß derselbe habe können eine Tochter zeugen. Denn die heidnischen Lasterer und hoffärtigen Leute befreien sich nur dessen, daß sie den Heiden übel nachreden, und daß sie auch was Gutes an ihnen ist und da sie recht gethan haben, doch lästern mögen. Denn sie wollen allein GOttes Volk sein, und gönnen den armen Heiden auch dasjenige nicht, was sie jemals Ehrliches oder Löbliches ausgerichtet und gethan haben. Darum widerstehe ich ihrem Geist gern, dieweil ich weiß, daß sie aus hoffärtigen stolzen Herzen den Heiden immer Uebeles nachreden. Also haben sie aber nun gedichtet, wie Potiphar den Jüngling Joseph gekauft habe, und dieser sehr schön von Gestalt gewesen, seiner habe wollen zur Unzucht mißbrauchen. Aber GOtt habe solcher Sünde gewehret und dieselbe verhütet, nämlich damit, daß er dem Potiphar sein männlich Glied verdorren lassen, daß er nicht habe können Kinder zeugen. Dies haben sie also leichtlich erdichtet; aber sie bringen gar keine Ursache oder Zeugniß, ihr Gedicht damit zu beweisen.

243. Augustin hat es sich auch lassen sauer werden, diese Frage aufzulösen, und beschließt doch endlich also: Wo es Einer dafür halten wollte, daß zwei Potiphar gewesen, derselbige kündigt damit nicht wider die heilige Schrift, und wenn man es dafür halten will, kann solches ohne Gefahr des Glaubens geschehen.

244. Ich glaube entschieden, daß ihrer zwei Eines Namens gewesen seien, und daß dies ein gemeiner und gebräuchlicher Name auch bei vielen Andern gewesen sei; wie unsere Leute auch gern pflegen von Königen und Fürsten Namen zu nehmen und dieselbigen zu gebrauchen, als, Maximilian, Ferdinand, Carl, aber noch viel mehr von denen, so heilig gewesen, als, Heinrich, Kunigunde, Conrad. Und in Sachsen ist von Lothar oder Luther, wie ihn Cäsar nennt, ein berühmter Name und groß Geschlecht geboren, nämlich, die Lüder. Also haben in Egypten ihrer Viele den Namen Pharao mit einem Zusatz gebraucht. Phera ist eigentlich Pharao, und Poti: was aber die rechte Etymologie desselben Wortes sei, will ich den Hebraisten auszuforschen befehlen.

245. Darum halte ich dafür, dieser Priester Potiphar sei ein anderer Mann denn der Potiphar droben, also daß Joseph seines Herrn Tochter nicht gefreiet habe. Und es scheint, daß er sich damit weislich vorgesehen und davor gehütet habe, dieweil das Exempel der Mutter, welche eine Ehebrecherin gewesen, am Tage war. Derhalben halte ich dafür, der fromme, ehrliche und züchtige Jüngling werde sich mit Fleiß und gutem Bedenken derselben Heirath enthalten haben mit der Tochter, so eine unehrliche Mutter gehabt.

246. Denn also hat jener Mönch gerathen: Wer um ein Mädchen freien wollte, der dürfte nicht auf den Vater sehen, sondern vielmehr auf die Mutter. Denn es folgt gemeiniglich die Tochter dem Leben ihrer Mutter nach. Und die Deutschen sagen auch in einem Sprüchwort: Das Bier schmeckt nach dem Faß. Ist die Mutter eine Hure, so ist die Tochter nicht fromm. Jedoch kann es wohl geschehen, daß eine unehrliche Mutter oder gottloser Vater einen frommen ehrlichen Sohn oder Tochter zeugen, und umgekehrt. Dies ist nun meine Meinung, daß zwei Männer gewesen seien Eines Namens: jener vorige ist ein Amtmann gewesen, der über ganz Egyptenland auf das Fleisch gesehen hat; dieser aber ist ein Bischof gewesen über die Kirche und der auf den Gottesdienst gesehen hat, daß Josephs Tochter auf Rath des Königs gefreiet hat. Ich will es aber unentschieden gelassen haben, ob diese Meinung wahr sei oder nicht; denn dabei steht der Glaube und christliche Religion nicht in Gefahr.

247. Man sollte viel lieber darnach gefragt haben, warum doch dieser heilige Mann Joseph ein abgöttisches Weib genommen habe, die eines heidnischen und gottlosen Priesters Tochter gewesen? Denn das ist wohl glaublich, daß alle Priester in Egypten, wie auch bei den andern heidnischen Völkern, gottlos gewesen sind, und die von den Jerobeam, so den Vätern gegeben waren, gar nichts gewußt haben. Warum hat er nicht lieber auf den Rath seines Vaters eine Andere aus seinem Volk und Geschlecht genommen?

248. Hierauf aber antworte ich also: Daran ist kein Zweifel: diese Priester alle, wiewohl sie von der Lehre und Verheißung nichts gewußt, werden doch weltlicher Weise ehrliche und rechtschaffene Männer gewesen sein. Und es ist keine Sünde gewesen, eines Heiden Tochter zur Ehe nehmen. Denn Abraham hat Sarah auch zum Weibe genommen zu Ur in Chaldäa, welches eine räuberische und gottlose Stadt der Babylonier war. Ja, er selbst, Abraham, ist auch anfänglich ein gottloser Mann gewesen, wie im Buche Josua am 24. Cap. V. 2. geschrieben steht: „Eure Väter wohneten vor Zeiten jenseit des Wassers, und dienten andern Göttern.“ Das würde zwar mir nicht geziemen von großen Männern zu sagen: aber Josua trägt keine Scheu, ihnen ihr voriges gottloses Wesen aufzurücken. Und ist eine große Wohlthat gewesen, daß GOtt Abraham von Ur in Chaldäa berufen und ihm die Verheißungen von Christo geoffenbart hat.

249. Also ist Potiphar in des gottlosen Betuel Hause gewesen; Lea und Rahel sind von Laban geboren, welcher ein abgöttischer Mann war. Aber Jacobs Söhne haben nicht Weiber genommen aus ihrem Geschlechte, sondern aus den Geschlechtern der Canaaniter. Abraham hat Hagar zu sich genommen, die eine egyptische Magd war; Boas hat Ruth genommen, eine Moabitin; Salma, der Fürst in Juda, hat die Hure Rahab genommen; und sind hernach in Egypten solcher Heirathen mehr gestiftet worden. Wiewohl nun Asnath auch gottlos und noch nicht unterrichtet gewesen ist in der rechten Erkenntniß der wahren Gottseligkeit, hat sie doch Joseph ohne Sünde zum Weibe genommen und sie darnach den rechten Gottesdienst gelehrt, und wie man den wahren GOtt anrufen sollte. Nun kann es aber wohl sein, daß viele von den egyptischen Priestern, da sie Joseph predigen gehört und sich so viel schöner Tugenden an ihm verwundert haben, zur rechten Gottseligkeit bekehrt worden und von ihrer Abgötterei abgetreten sind; wie David im 105. Psalm V. 22. sagt, daß Joseph ein Bischof aller Bischöfe und Lehrer über alle anderen Lehrer in Egypten gewesen sei.

250. Und daß auch der König gottseelig gewesen, oder ja etlichermaßen ein Licht davon gehabt, kann aus seinem Worten verstanden werden, dieweil er droben, V. 38., gesagt, daß Joseph den Geist GOttes habe. Daraus denn erhellt, daß er nicht viel Götter wird gehabt oder geehrt haben. Denn es sind Worte eines solchen Mannes, der GOtt lauter allein dient. Und sind vielleicht in Egypten noch etliche Anzeichen und ein kleines Ueberbleibsel geblieben von der Lehre der Väter, so Abraham daselbst gelehrt; wiewohl sie allgemach etwas verdunkelt worden. Wie es denn damit zu gehen pflegt, wo die Lehre der Religion ausgebreitet wird, daß da der Teufel für und für sein Unkraut mit untermengt. Zu unserer Zeit hat die Lehre des Evangeliums, da sie wiederum gereiniget und aufgerichtet, viel Leute gewonnen und herzugebracht, so zuvor mit des Antichrists, des Papstes, Tyrannei beschweret gewesen; aber es sind auch zugleich von uns ausgegangen Wiedertäufer, Sacramentsschwärmer und andere Rottengeister, welche von der heiligen Dreifaltigkeit, und Menschenwerdung Christi solche Lehre hervorgebracht, so offenbar falsch und ungöttlich ist. Denn sie sind nicht von uns gewesen, wiewohl sie eine Zeitlang bei uns waren, haben aber die reine Lehre nicht gemeinet oder gesucht, sondern nur ihre Ehre und einen großen Namen.

251. Also haben die Egypter erstlich die reine Lehre und Verheißungen von Abraham gehört; aber dieselben sind bei ihnen hernach verdunkelt worden durch die falschen Lehrer, bis endlich Joseph solche Lehre wiederum erneuert hat.

252. Und daraus schließe ich nun, daß Joseph wird ein heiliges und gottseeliges Weib zur Ehe gehabt haben, sie sei nun gleich vor der Heirath oder aber hernach von Joseph unterrichtet gewesen. Und sie ist eine Mutter geworden in Israel und eine rechte Altmutter, wie Rahel, Lea und Thamar auch gewesen.

\quellseite{368}

Denn sie hat Manasse und Ephraim geboren, welche zweier sehr mächtiger Stämme Väter gewesen sind.

253. Und damit soll man dem Stolze der Juden wehren, auf daß sie wissen, sie thun daran unrecht, daß sie die Heiden so gar verachten, sonderlich die Nachkommen Hams. Denn sind nicht Ephraim und Manasse von einer heidnischen Mutter geboren worden? Zudem mögen sie auch die Hure Rahab ansehen und die abgöttische Ruth, item Thamar, die Cananitin, welche mit der Blutschande ihres Schwiegervaters befudelt gewesen. Warum bist du denn so stolz, du heiliger Jude, so du doch aus heidnischem Samen geboren bist, dazu ernähret und gesäuget worden mit der Milch eines Weibes, die eine Fremde gewesen ist?

254. Also wird Asnath auch mitgezählt unter die heiligen Altväter und Matronen des Hauses Israel. Denn sie ist eine Mutter zweier sehr berühmter und königlicher Stämme. Nachdem Ruben die Erstgeburt verloren, sind Ephraim und Manasse an seine Statt getreten; wie dasselbe 1 Chron. 6, 1. 2. erzählt wird: „Die Kinder Ruben, des ersten Sohns Israel (denn er war der erste Sohn, aber damit, daß er seines Vaters Bette verunreinigte, ward seine Erstgeburt gegeben den Kindern Joseph, des Sohns Israel; und er ward nicht gerechnet zur Erstgeburt).“ So groß hat Gott die heidnischen Weiber geachtet, auf daß er damit bezeugete, daß er die Heiden nicht verwerfen hätte, als die da sollten vermaledeyet sein; wie sie denn die Juden, welches ein verstockt und verblendet Volk ist, gar verfluchen. Denn da sie uns das heidnische Geblüt aufrücken, sehen sie nicht, daß sie eben damit ihren eigenen Ursprung und Geburt verunreinigen und vermardemmen.

255. Es hat aber unser Herr Christus vom Geblüte mancherlei Völker wollen geboren werden. Denn er hat Rahab, Ruth und Thamar zu Müttern. Und dieweil er derselbigen nicht schämt, und in seinem Geburtsregister auch egyptische, cananitische und moabitische Weiber mitgezählt werden: warum haben denn die heiligen Juden so einen großen Abscheu vor unserm Geblüt, als vor etwas Beflecktem und Besudeltem?

256. Eben dasselbe mag man auch von David und den andern Königen sagen, so auch von denselbigen Müttern geboren gewesen. „Denn bei Gott ist kein Ansehen der Person“, Ap. Gesch. 10, 34. Col. 3, 25. Und dies ist uns zum Trost geschrieben, auf daß wir wissen, daß Gott nicht allein der Juden, sondern auch der Heiden Gott ist, Röm. 3, 29. Denn wiewohl die Heiden nicht beschnitten gewesen, sind sie doch Gott eben so lieb und angenehm als die Juden, und ist in dem allerberühmtesten und größten Geschlecht Judas eine Vermischung des Geblütes Thamars geschehen. Zuletzt, obgleich den Heiden keine Verheißung geschehen ist, so werden sie doch auch zugleich darin mitbegriffen. Denn also lautet ja die Verheißung: „In deinem Samen sollen alle Völker auf Erden gesegnet werden.“

\begin{quote}
\textbf{V. 45. Also zog Joseph aus, das Land Egypten zu besehen.}
\end{quote}

257. Das erste Jahr ist er vor Pharao gestanden, und hat bestellt, was zum Hausregiment, Zucht und guter Ordnung am Hofe des Königes gehörete, ehe er sich des politischen oder Weltregiments angenommen hat. Derhalben ziehet er noch nicht aus, das Getreide zu sammeln, sondern nur das Land zu besehen, sich zu erkundigen, wie es um die Religion, Gerichtshändel und Obrigkeit beschaffen, und was hin und wieder in den Städten und bei allen Ständen für ein Leben und Sitten wären; und ist doch gleichwohl an des Königes Hofe geblieben. Und ist dies, meines Erachtens, im ersten Jahre geschehen, ehe die fruchtbaren reichen Jahre angefangen haben.

\begin{quote}
\textbf{V. 46. Und er war dreißig Jahr alt, da er vor Pharao stund, dem Könige in Egypten.}
\end{quote}

258. Zu der Zeit, da Joseph die Träume gedeutet und ein Weib genommen hat, ist er dreißig Jahre alt gewesen: welches Mose darum erzählt, auf daß man daraus schließen könne, wie lange er im Gefängniß und Knechtschaft gewesen sei. Im siebenzehnten Jahre seines Alters ist er verkauft worden, im dreißigsten Jahre aber wird er zu königlicher Herrlichkeit erhoben. Derhalben hat er seinem Herrn zehen Jahre treulich und wohl gedienet, und hat ihm sein Haus gebessert, sein Gut und Habe gemehret, daß er reich geworden ist. Er ist aber drei oder vier Jahre im Gefängnisse gewesen, da er des Ehebruchs von seines Herrn Weibe, der Ehebrecherin, war beschuldigt worden.

259. Das ist nun uns zur Lehre geschrieben, auf daß wir den großen Jammer und Elend, so Joseph erlitten, verstehen lernen, davon der 105. Psalm V. 18. sagt: „Sie zwangen seine Füße im Stock, sein Leib mußte in Eisen liegen.“ Du magst das Eisen nun verstehen, daß er am Leibe mit Ketten gebunden gewesen, oder aber in dem Verstande, wie Simeon sagt zu Maria Luc. 2, 35.: „Und es wird ein Schwert durch deine Seele dringen“; so reimt es sich nicht übel, daß er im ersten Sturme und Grimme seines Herrn in den Stock werde gelegt worden sein, hernach aber hat er Gnade gefunden bei dem Kerkermeister, der hat ihm die Fessel los gemacht und ihm die Gefangenen befohlen zu verwahren.

260. Dies Exempel wird uns vorgehalten, auf daß wir uns an Geduld gewöhnen, dazu auch an Hoffnung und Beständigkeit im Kreuz und Unglück, wiewohl der lange Bezug, ehe uns geholfen wird, unsere Herzen sehr plagt und martert. Denn Gott sieht das Ende unserer Trübsal, wie er zu den Kindern Israel Jer. 29, 10. 11. sagt: „Wenn zu Babel siebenzīg Jahre aus sind, so will ich euch besuchen und will mein gnädiges Wort über euch erwecken, daß ich euch wieder an diesen Ort bringe. Denn ich weiß wohl, was ich für Gedanken über euch habe, spricht der Herr, nämlich Gedanken des Friedens und nicht des Leidens, daß ich euch gebe das Ende, deß ihr wartet.“ Es soll, sagt er, das Gefängniß zu Babel nicht ewig währen, es ist um siebenzig Jahre zu thun; es thut aber wehe, Geduld tragen und des Endes erwarten.

261. Derhalben wird nun dies alles geschrieben, daß wir es lesen und betrachten sollen, und daß wir lernen die Hand Gottes tragen und dulden, der uns heimsucht und väterlich stäupt. Denn die Zeit der Strafe oder Trübsal ist nicht eine unendliche Zeit, wie es scheint, sondern Gott hat gewisse Stunden und Augenblicke bestimmt. Darum ist nur vonnöthen, daß wir beharren und beständig bleiben, wie Christus Matth. 24, 13. sagt: „Wer beharret bis ans Ende, der wird selig.“

262. Siehe an das Exempel aller andern Patriarchen, wie mit großer Beständigkeit sie in ihrem Jammer und Trübsal beharreten; so wirft du an ihnen auch die rechte Frucht der Geduld und Beständigkeit sehen. Wie dies göttliche und hohe Werk der Ehre und Erhöhung Josephs anzeigt, welcher nimmer zu solcher großen Ehre gekommen wäre, wo er nicht also geplagt, getödtet und gedrückt worden wäre. „Es sey denn“, spricht Christus Joh. 12, 24., „daß das Weizenkorn in die Erde falle, und ersterbe, so bleibts alleine; wo es aber erstirbt, so bringts viel Früchte.“

263. Dies wird uns fleißig und immer vom Heiligen Geiste also vorgestellt, auf daß wir wissen sollen, daß wir Gläubigen wohl versucht, aber doch nicht verderbet und zu nichte gemacht werden, damit wir ihrem Glauben auch können nachfolgen, und uns daran halten und aufrichten. Es sagt Kirill von Aleaß, daß er seine Gesellen auch aufgerichtet, getröstet und gestärkt habe mit Exempeln der vorigen Errettung und gesagt: O lieben Gesellen, wir wissen ja wohl, was für Unglück wir zuvor erlitten; o! wir haben ehmals viel mehr erlitten; Gott wird dieses Jammers und Unglücks auch ein Ende machen. Item: Superanda omnis fortuna ferendo est, das ist: Man muß alles Unglück mit Geduld überwinden, und man muß nicht aufhören noch nachlassen. Haben nun die Heiden mit so großem starkem Muthe allerlei Unglück und Beschwerung überwunden: warum sollen es denn wir Christen damit an uns fehlen lassen?

\subsection{II}

\begin{quote}
\textbf{V. 46—49. Und (Joseph) fuhr aus von Pharao, und zog durch ganz Egyptenland. Und das Land that also die sieben reichen Jahr: und sammelten alle Speise der sieben Jahre, so im Lande Egypten waren; und thaten sie in die Städte. Was für Speise auf dem Felde einer jeglichen Stadt umher wuchs, das thaten sie hinein. Also schüttete Joseph das Getreide auf, über die Maße viel, wie Sand am Meer, also, daß er aufhörte zu zählen; denn man konnte es nicht zählen.}
\end{quote}

264. Dies ist nun ein anderer Auszug, denn der droben gewesen. Denn jetzt greift er nun zum Regiment im Lande, und bleibt nun nicht nur vor Pharao, sondern hat das Hofregiment einem Andern befohlen, und er kommt selten wieder zu Hofe. Denn er muß nun wiederum durch ganz Egyptenland ziehen, welches er zuvor nur obenhin beschäftigt; und jetzt befiehlt er von einem Flecken zum andern und von Städten zu Städten die Scheunen oder Kornhäuser und andere Dinge mehr, und zählt das jährliche Einkommen im ganzen Lande. Es geht aber alles sein glückselig von Statten, und die Unterthanen im ganzen Lande haben sich ihm selbst willig unterworfen, haben ihm gehuldigt und sind ihm mit höchster Willigkeit gehorsam. Also ist daselbige Königreich sehr schön geordnet und allenthalben wohl bestellt gewesen, und es scheint wahrlich, daß dazumal eine goldene Zeit in Egypten gewesen sei.

\quellseite{369}

265. Es hat aber Joseph eine große Last getragen, welche er nicht hätte tragen können, wo ihm Pharao mit seiner Macht nicht beigestanden und ihm geholfen hätte. Denn damit hat Pharao die Herzen der Unterthanen gewonnen, daß sie Joseph auch günstig und gehorsam geworden sind. Und wo diese zwei Stücke zusammen kommen, da folgt letztlich das auch hernach, daß man die Obrigkeit auch ehrt und sich vor ihr scheut, und daß auch bei allen Rathschlägen, und bei allem, was man handelt oder thut, Glück und Heil ist. Denn um der Strafe willen wird die Obrigkeit nicht gefürchtet; sondern sie wird geehrt und vor Augen gehalten um ihrer Tugend und Weisheit willen. Wenn aber ein Scherge, auf Hebräisch maabir, regiert, mit welchem Namen Seleucus Philopator genannt wird Dan. 11, 20, da ist kein williger Gehorsam. Denn es wird da die Person verachtet, welche nichts thut, das der Obrigkeit geziemt zu thun. Darum sollen an einem Regenten diese Stücke bei einander sein, nämlich, Ansehen und Gewalt und Ehrbietung, wiewohl das ein gemein Gebot des heiligen Paulus ist, Röm. 13, 7., da er sagt: „Gebet Ehre, dem die Ehre gebührt, Furcht, dem die Furcht gebührt.“

266. Es werden aber die Leute erweckt und bewogen zur Ehrerbietung durch die Tugend, Weisheit, Gerechtigkeit und Mäßigkeit derer, so im Regimente sind. Denn darum hat Joseph bei Pharao so ein groß Ansehen, dergleichen auch bei andern Regenten und Unterthanen im Lande, daß wenn er nur jemand angesehen und gewinkt, oder aber nur einen Finger geregt hat, alsbald bei jedermann der Gehorsam gefolgt und Alles glücklich von Statten gegangen ist. Er hat sie mit Gewalt oder Strafe nicht gezwungen; welches die Tyrannen pflegen zu thun, deren Regiment nur mit Gewalt erzwungen wird, und die damit doch gleichwohl nicht zuwege bringen, daß ihnen die Unterthanen von Herzen gehorsam sind, ja, sie machen damit vielmehr, daß sie ihnen nicht können günstig sein.

267. Das ist aber eine große Wohltat, wenn sich das Volk verwundert über der Fürsten Klugheit, Gottseligkeit und Glück, und Lust hat an ihrem Anblick und Anschauen, an ihrem Regiment und Lobe, also daß sie es dafür halten, sie müssen ihren Herren aus sonderlicher Gunst und Liebe gehorsam sein. Wenn aber die Fürsten um ihrer Sünde und Schande willen bei den Unterthanen sowohl, als sonst hin und wieder bei jedermann ein bös Geschrei haben, daß man sagt: Unser Fürst ist ein Bösewicht, ein Mordbrenner, das ist ein sehr böses, unglückseliges Regiment.

268. Darum fährt nun Joseph fort, das Amt eines frommen Regenten auszurichten, da er das Getreide einsammelt, und in allen Städten sonderliche Scheunen oder Kornhäuser dazu bestellt. Denn er ist nicht nachlässig oder faul gewesen. Und wiewohl er ein sehr groß Ansehen und große Gewalt gehabt wegen seiner Weisheit und sonst allerlei Tugenden, und auch wegen des heiligen Geistes, welchen Pharao ihm gespüret hat, gehört doch noch dazu, daß man fleißig und wacker sei.

269. Denn die Leute richten ihr Amt immer so fleißig aus, daß ihnen daneben nicht sollte vorkommen sein, daß der Herr oder Regent auch selbst mit darzu sehe. Welches die Alten damit angezeigt haben, da sie gedichtet, daß Jupiter in seinem Scepter ein Menschenauge führe, und ist ein sehr schön Zeichen. Denn also pflegt man im gemeinen Sprüchwort zu sagen: Oculus domini impinguat equum: Des Herrn Auge macht das Pferd fett; item: Frons occipito prior est: Die Stirne ist uns näher als das Hinterhaupt. Das ist so viel gesagt: Es ist in allen Regimenten nützlich und auch nötig, daß man visitire und ein fleißiges Aufsehen habe auf alle Dinge, auf daß es allenthalben recht zugehe. Darum hat auch Joseph im ganzen Lande Alles fleißig besehen, er hat seine Augen auf das Land gehabt und nicht außerhalb des Landes, welches faule und vornwitzige Leute zu thun pflegen.

270. Also gehören auch zum Kirchenamt getreue Bischöfe und Aufseher, welche über der reinen Lehre und rechtem Brauch der heiligen Sacramente mit Ernst halten, die falschen ungöttlichen Lehren von rechter Lehre zu unterscheiden wissen. Und müssen solche Bischöfe so viel fleißiger wachen, wie viel begieriger auch die reißenden Wölfe nach den armen Schäflein trachten, dieselben umzubringen. Wie wir sehen, daß die Ketzer und der Papst sehr wachsam sind, schlafen nicht sicher, schnarchen auch nicht, sondern warten sehr gespannt auf allerlei Gelegenheit, dadurch sie die Herzen der Gottseligen betrügen und verführen möchten.

271. Aber ich will nun wiederum zum Text kommen; denn da gefällt mir gar nicht, daß unser Dolmetscher im Lateinischen sich gar zu viel Freiheit genommen, und von selbst, seines Gefallens, diese Stelle ausgelegt durch eine Paraphrase, das ist, da er ein Wort durch viele andere Worte ausdrückt, und doch selbst nicht verstanden hat, was er geredet. Denn also lautet unser lateinischer Text: Venitque fertilitas septem annorum, et in manipulos redactae segetes congregatae sunt in horrea Aegypti, das ist: Und es kam eine Fruchtbarkeit von sieben Jahren, und die Ernte wurde in Garben gebunden und in die Kornhäuser Egyptens gebracht. Es ist eine unzeitige Rhetorik, die er allhier gebraucht hat. Denn man soll die Worte im Text behalten, und dieselbigen auslegen. Darum soll man es aus dem hebräischen Text im Lateinischen also geben: Et fecit terra in septem annis fertilitatis ad acervos, et congregavit, scilicet Joseph, wie wir es auch verdeutscht haben: „Und das Land that also die sieben reichen Jahr; und sammelten alle Speise der sieben Jahre.“ Denn Hieronymus und Lyra verstehen es auch also, daß es heißen soll: segregare in manipulos, das ist, in die Garben sammeln.

272. Aber da es gedroschen ist, hat er große mächtige Haufen zusammen getragen in sieben Jahren. Nachdem das Land getragen, hat er alle Speise den Leuten zu Nutz in die Scheunen oder Kornhäuser gesammelt. Es ist ein sehr verständiger Mann gewesen, daß er nicht alles Einkommens der sieben Jahre auf Einen Haufen gebracht, sondern hat in einer jeglichen Stadt die Speise des Ackers hinlegen lassen, welcher in selbigem Bezirk oder Gericht gelegen gewesen, das ist, Joseph hat es in einer jeden Stadt also verordnet, daß er darin insonderheit das Getreide gesammelt, so auf den umliegenden Aeckern gewachsen war, daß man es nicht collectiv, sondern distributiv verstehen soll, nämlich also, daß eine jegliche Stadt ihre eigenen Kornhäuser gehabt.

273. Auf diese Weise wird die Ordnung gelobt, so Joseph gehalten, das Getreide zu sammeln. Welche Ordnung auch sonst in andern Dingen aufs höchste gerühmt wird. Denn was man ordentlich ausrichtet, dasselbe geräth immer besser und geht glücklicher von Statten beide im Welt- und auch im Hausregiment, als daß man ohne gewisse Ordnung, also ohngefähr hin thut. Man verschwendet Alles mehr, wenn die Regimente nicht recht und fein ordentlich bestellt sind, wie an unseres Fürsten Hofe zu sehen ist. Denn wo dieser denselbigen reformiren würde und in eine rechte Ordnung bringen, da würde man so große unzählige Kosten nicht bedürfen; denn gute Ordnung erhält die Nahrung, aber Unordnung zerstreuet und bringt Alles durch. Also sehen wir, daß die, so gute Hausväter sind, und ihre Haushaltung fein ordentlich gefaßt haben, darnach sie Alles regieren, leiblich reich werden.

274. Und dies sage ich wider Hieronymus, welcher damit, daß er den Text so weitläufig gegeben, diese Ordnung verdunkelt und verkehrt hat. Denn es wird ohne Zweifel alles Ding sehr wohl bestellt und geordnet gewesen sein, nämlich, mit der Ernte, mit Dreschen und wie man es hat sollen zu Haufen sammeln, ja, auch wie man mit Pferden und Wagen, und sonst mit den allgeringsten Dingen hat umgehen sollen. Er läßt einer jeden Stadt die Einkünfte von ihren Aeckern, nimmt nur den fünften Theil davon ohne allen Schaden der Leute, und dasselbe sammelt er und legt es hin, schüttet es auf die Kornhäuser, auf daß also der Vorrath der Speise gegen die künftige Theurung immer zunehmen und größer werden möge. Denn darum setzt er im Text hinzu: „Also schüttete Joseph das Getreide auf, über die Maße viel, wie Sand am Meere, also, daß er aufhörete zu zählen“; welches nicht von einzelnen Körnern zu verstehen ist, sondern von den Scheffeln des Getreides, so gesammelt worden ist. Und hat doch dasselbe viele Getreide sehr schön geordnet und jegliches an seinen Ort geschickt. Dies ist das erste Stück der Weisheit Josephs, so er in seiner Regierung gebraucht.

\quellseite{370}

\section[Siebenter Theil]{Siebenter Theil. Von den zwei ersten Söhnen, so Joseph geboren, und von der Theurung, so nach den sieben reichen Jahren entstanden}

\subsection{I}

\begin{quote}
\textbf{V. 50. Und Joseph wurden zween Söhne geboren, ehe denn die theure Zeit kam, welche ihm gebar Asnath, Potiphera, des Priesters zu On, Tochter.}
\end{quote}

275. Wir haben gesagt von den sieben reichen fruchtbaren Jahren, mit welcher Geschicklichkeit und Klugheit Joseph den Weizen und anderes Getreide in die Scheunen, so in Egypten verordnet, gesammelt habe. Was aber nun erzählt wird von den Söhnen Josephs, scheint gar unnütz und vergeblich. Und reimet sich auch die Worte nicht sehr wohl. Denn also lautet es nach dem Hebräischen: Joseph natus est duo filii: Joseph ist zween Söhne geboren etc. Ich will aber dasselbe den Grammatikern befehlen auszurichten und zu erklären, warum doch so ein trefflicher großer Redner, wie Mose gewesen, so unbequem geredet habe.

276. Dies aber gehört zu dem, davon auch droben ist gehandelt worden, wider die tolle unsinnige Hoffart und Rühmen der Juden, welche alle Heiden gar verachten, unangesehen ihres eigenen mütterlichen Geblüts, welches egyptisch, chaldäisch, ammonitisch und moabitisch gewesen. Eben als wenn ich die Juden verdammen wollte, und sie mich noch mit der heiligen Schrift davon auch überzeugen würden, daß ich von einer jüdischen Mutter geboren wäre, da möchte ich billig wegen meines Frevels beschuldigt und angeklagt werden. Aber sie sind unverchämte Leute, lassen sich mit keinerlei Zeugniß oder Weise bewegen, ihren Irrthum zu erkennen und sich zu bessern.

277. Derhalben ist dies ein trefflicher Text, der da zeuget, daß zween sehr berühmte und namhafte Stämme Israels von einer egyptischen Mutter geboren sind, den Juden zur Schande und den Heiden zur Ehre. Und wir werden im zukünftigen Gericht und Leben hören, wie David und die andern Patriarchen diese heidnischen Weiber gerufen werden, nämlich Rahab, Ruth, Thamar und Asnath, die Mütter, von welchen sie geboren sind nach dem Fleisch.

Wiewohl aber daran kein Zweifel ist, Joseph werde beide Söhne haben beschneiden lassen, so ist er doch so weise wohl gewesen, daß er die andern Egypter zur Beschneidung nicht gezwungen hat. Denn das liest man an keinem Orte, daß der König selbst, oder sonst jemand anders im Lande wäre beschnitten worden. Und wir werden hernach im 43. Cap. V. 32. hören, daß es sich bei den Egyptern nicht geziemet, ja ein Greuel bei ihnen gewesen sei, mit den hebräischen Schafhirten zu essen. Und wiewohl den Söhnen Jakobs erlaubt worden, daß sie haben mögen zu Hof gehen, daß sie daselbst essen möchten, um des großen Ansehens, Ehre und Herrlichkeit willen, so ihr Bruder Joseph daselbst gehabt; haben doch Joseph und seine Brüder besonders an einem eigenen Ort allein gegessen.

279. So hat nun Joseph das Königreich Egypten regiert, welches der Juden Ceremonien und Sitten verachtete und einen Greuel davor hatte, soviel die Beschneidung belanget, und hat doch sonst die andern Stücke der rechten Lehre von Joseph gehört. Es hat ihn GOtt in dieser Regierung wunderlich geführt, da seine Religion eines Theils angenommen und auch eines Theils verworfen worden, nämlich, daß der Glaube und Erkenntniß GOttes angenommen, die Beschneidung aber verworfen worden ist.

\begin{quote}
\textbf{V. 51. Und (Joseph) hieß den ersten Manasse; denn GOtt, sprach er, hat mich lassen vergessen alles meines Unglücks, und alle meines Vaters Hauses.}
\end{quote}

280. Dies ist eine wunderbare Rede, daß Joseph sagt, er habe nicht allein alles Unglücks, sondern auch seines Vaters Hauses vergessen. Denn Manasse wird zum wenigsten zwei oder drei Jahre alt gewesen sein, ehe die theure Zeit angefangen, da die Zeit der reichen Jahre halb verlaufen gewesen. Warum hat er aber inzwischen nicht etliche Gesandte gen Hebron gesandt, die nach seinem Vater und Brüdern hätten fragen mögen? Ja, er zeiget allhier an, daß er gleichsam mit Fleiß und auf GOttes Befehl seines Vaters Hauses sammt allem Jammer und Unglück vergessen habe. Denn so viel heißt bei den Hebräern das Wort amal; wie im Buch Hiob am 5. Cap. V. 7. auch steht: „Der Mensch wird zu Unglück geboren, wie die Vögel schweben empor zu fliegen.“ Also nennt Joseph die ganze Historie der dreizehn Jahre amal, das ist, Unglück, Trübsal und Jammer.

281. Aber es scheint doch, als habe dieser sehr weise und heilige Mann, der ja ohne Zweifel seinen Vater und Brüder geliebet, allhier solche Worte geredet, die sich mit der Gottseeligkeit nicht reimen. Denn also darf niemand außer ein Gottloser reden, daß er seines Vaters Hauses vergessen habe. Nun ist aber Mempis, da der königliche Sitz in Egypten gewesen, vom gelobten Lande durch der Philister Land und durch die Wüste, das steinichte Arabien genannt, nach der Rechnung der Geographen, nicht weiter denn vierzig oder fünfzig Meilen gelegen. Hebron, da Jakob, Josephs Vater, gewohnt, ist schier an der Grenze Egyptens gelegen gewesen. Warum hat denn Joseph, der so ein gewaltiger großmächtiger Herr war und Roß und Wagen in seiner Gewalt gehabt, nicht Kundschafter ausgesandt gen Hebron zu seinem Vater? Ist das die Ehrerbietung, damit Kinder ihre Eltern ehren sollen, und daß du dich darüber noch rühmen darfst, du habest deines Vaters Hauses vergessen? Es wäre noch etwas leidlicher gewesen, daß er gesagt: Er hätte alles seines Kummers und Unglücks vergessen, und hätte doch das nicht hinzugefügt von seines Vaters Hause, daß er desselben auch vergessen.

282. Ich weiß zwar nicht, was ich dazu sagen soll, denn allein, daß es scheint, daß aus diesem Text die Worte im 45. Psalm V. 11. genommen sind, da also steht: „Höre, Tochter, schaue drauf, und neige deine Ohren, vergiß deines Volkes und deines Vaters Haus.“ Es ist aber zumal ein sehr hartes Wort, darin er sich rühmt, er habe seines Vaters Hauses vergessen. Und ich sehe hier keine Ursache, weshalb er sich nicht sorgfältig um sein Vaterhaus bekümmern haben sollte. Ob er vielleicht gedacht, es müßten beide Vater und Brüder vor großem Herzeleid gestorben sein; oder aber, ob er sich aller Hoffnung begeben, sein väterlich Erbe zu erlangen. Denn also lauten die Worte: „Ich habe alle meines Vaters Hauses vergessen“, das ist: Ich werde nun kein Erbe sein zu meines Vaters Gut, wie die Träume, so mir vorzeiten eingegeben worden, zu bedeuten scheinen, daß ich würde ein Herr werden über Brüder und Eltern. Da wird er gedacht haben, solches wäre Phantasie und nur menschliche Träume oder Trügerei gewesen, die nicht einen solchen Ausgang haben würden, wie er wohl gehofft hätte. Und kann deshalb dies also verstanden werden, daß es ein neues Kreuz oder Tödtung sei, damit GOtt Joseph getödtet, und dahin treibt, daß er es für gewiß halten solle, es werde derselbe Traum vergeblich und umsonst sein. Da wird er gedacht haben: Ich will in Egypten bleiben, da wird dermaßen, wenn ich nun gestorben bin, GOtt meinem Nachkommen wohl mit einem andern Erbe versehen. Also hat er voll Glauben auf die Ehre und Herrlichkeit der Erbgeburt, so er in seines Vaters Hause gehabt, Verzicht geleistet, und verspricht sich selbst und seinen Kindern einen andern Segen Gottes auf eine viel andere Weise, denn er zuvor gehofft hatte.

283. Jedoch bleibt noch die Frage, warum Joseph nicht Gesandte oder Kundschafter zu seinem Vater ausgesandt habe? Denn wiewohl es eine große Augen an ihm ist, daß er alle Hoffnung, so er auf sein väterliches Erbe gehabt, gar hat fallen lassen, und sich daran genügen läßt, was ihm sonst in Egypten anderswoher zu Theil werden möchte: so ist es doch noch wunderlich, daß er nicht etliche gen Hebron gesandt, die seinen Vater und Brüder besucht und von ihrem Leben und Zustand wieder Nachricht gebracht hätten. Es wird entweder GOtt solche Gedanken verhindert haben, oder es haben ihn seine vielen Geschäfte und Händel, so ihm Amts halben obgelegen, aufgehalten; und doch sind dieselbigen Geschäfte so groß nicht gewesen, daß er derhalben nicht hätte etliche Diener gen Hebron abschicken können. Oder es mag die Ursache gewesen sein, daß er niemand dahin gesandt, dieweil die Egypter die Juden und Schafhirten verachtet und einen Abscheu vor ihnen gehabt. Darum hat er auch keine Lust gehabt, sie zu sich zu fordern. Oder aber er wird durch den prophetischen Geist gesehen haben, daß bald seine Brüder sammt dem Vater zu ihm nach Egypten kommen würden. Oder wo er ja deshalb sollte beschuldigt werden, daß er daran gesündiget hätte, daß er niemand zum Vater geschickt, wollen wir ihn unter diejenigen mit
\quellseite{371}
zählen, so damit sündigen, daß sie nicht thun, was sie billig thun sollten. Es sei nun gleich, was für eine Ursache es wolle, so ist es wahrlich eine wunderliche Rede, darin er anzeigt, er habe alle seines Vaters Häuser vergessen.

\begin{quote}
\textbf{V. 52. Den andern hieß er Ephraim; denn Gott, sprach er, hat mich lassen wachsen in dem Lande meines Elends.}
\end{quote}

284. Er behält noch immer im frischen Gedächtnis das Elend und Jammer, so er erlitten, und will so viel sagen: Siehe, ich bin elend und arm gewesen, und Gott, der die niedrigen, elenden, betrübten, verlornen und verdammten Menschen anzusehen pflegt, der hat sich meiner auch erbarmt, und mir einen Sohn gegeben, mit Namen Manasse, daß ich alles Unglücks vergessen sollte. Darnach hat er mir auch noch einen andern Sohn gegeben, Ephraim, daß ich wachsen und groß werden sollte im Lande meines Leidens und Elendes.

285. Denn so viel heißt das hebräische Wort onji, und wird eben in solchem Verstande auch gebraucht Jes. 53. Ich bin, sagt Joseph, im Kerker gefangen gelegen, da ich alle Hoffnung meines Lebens und Befreiung weggeworfen und verloren hatte, auch also, daß ich schier ein todter Stamm oder Klotz hätte mögen genannt werden, der nimmermehr grünen oder Frucht tragen würde. Aber wie Jesaia von der Ruthe, so vom Stamme Isai ausgehen wird, redet, also sagt Joseph von sich auch, daß von einem unfruchtbaren und erstorbenen Stamme gleichwohl noch Kinder und Erben herkommen und geboren werden, und daselbst noch von einem egyptischen Weibe; das ist ein wunderbarerlicher Gott.

286. Dies sind aber Exempel, die dazu gehören, daß der Glaube dadurch aufgerichtet und gestärkt werde. Denn wo wir auch also in Anfechtung beharren könnten, so würden wir wahrlich denselbigen Namen und die Art oder Natur Gottes auch erkennen, nämlich, daß er nicht allein machen kann, daß wir unsers Elendes und Unglücks vergessen, sondern daß diejenigen, so gar verdorben und zu nichts geworden sind, auch wiederum wachsen und gemehret werden müssen.

287. Und das ist ein treffliches Lob und große Ehre Gottes gewesen, damit Joseph die wunderbaren Werke Gottes gerühmt, nämlich, daß er aus Nichts Alles machen kann; welche Werke wir auch erkennen und uns darin üben sollen. Denn wir müssen gedemüthigt, getödtet, und wie das Weizenkorn in die Erde geworfen werden und sterben, und darnach mit vielfältiger Frucht wiederum auferstehen. Denn er ist ein Gott Ephraim, das ist, des Wachsens und der Mehrung; aber im Lande onji, das onji, das ist, der Tödung und des Elendes. Und dies ist eine gewisse Regel, die für und für also währet mit den göttlichen Werken, nämlich, aus Nichts Alles machen.

288. Denn also werden wir begraben und in die Erde verscharrt, daß wir darin verfaulen und von den Würmern verzehrt werden: wir werden gesäet in Schwachheit, verweslich, in Unehre und Verwesung, so verächtlicher und schändlicher Gestalt, daß schier zu wünschen, wir wären nie geboren, denn daß wir in solcher Unehre von Würmern gefressen und verzehrt werden sollen und verwesen. Also werden wir gar zu nichts gemacht. Aber Gott sagt: Ich will dich aus dem Nichts und onji, aus der Verwesung, Würmern, Staub und Erde wiederum erwecken, und will aus dir nicht allein einen Manasse, sondern auch einen Ephraim machen, viel schöner und heller, denn die Sonne ist. Und das wird endlich also gewiß geschehen; denn wir haben die gewisse Verheißung. Wiewohl es uns in diesem Leben nicht allenthalben nach unserm Willen geht und wie wir es gerne haben wollten, so werden wir doch im zukünftigen Leben dies auch rühmen und sagen: Gott hat mich lassen vergessen meines Vaters Hauses, und alle meines Unglücks.

289. Es sollte uns aber dies Vergessen billig in diesem Leben auch allgemach ankommen. Denn wiewohl wir uns jetzo unds nicht enthalten können, daß wir an die Würmer und Verwesung nicht denken sollten, dieweil wir sie vor Augen haben, wird es doch dermaßen dazu kommen, daß Gott alle Thränen von unsern Augen wird abwischen, wie in der Offenbarung St. Johannis am 7. Cap. V. 17. steht. Derhalben muß der Glaube hier anfangen, der Thränen und Schande zu vergessen, welche er nicht sieht. Wiewohl die Augen das Verwesen sehen und die Ohren das Klagen und Seufzen hören, die Nase den Stank vom erstorbenen Leibe riecht: so muß doch der Glaube sagen: Ich weiß davon nichts, ich sehe dieses nicht; ja, ich sehe eine Mehrung und eine Klarheit, welche die Sonne und Sterne übertrifft.

290. Derhalben werden uns solche Exempel dazu vorgehalten, daß wir daraus lernen sollen, daß Gott ein Schöpfer aller Dinge sei, der die Todten lebendig macht und ein Verklärer der Würmer und der schändlichsten Verwesung ist. Und er will haben, daß wir in diesem Leben dasselbe also im Glauben erkennen und rühmen sollen, hernach aber werden wir es im zukünftigen Leben in der Wahrheit und mit der That auch erfahren und empfinden.

291. Joseph ist ein solcher großer Fürst in Egypten nicht geworden und Jakob hat den Samen eines so großen Volks auch nicht eher gesehen, sie haben denn zuvor beide müssen getödtet werden. Und das hat Joseph aus eigener Erfahrung gelernt. Darum sagt er: Es muß dahin kommen, daß Einer vergesse alles, was er auf Erden zu hoffen hat. Ich war der Erstgeborne und Erbe von der vornehmsten Mutter; das ist alles verloren, und ist gar keine Hoffnung mehr, das Erbe oder Erstgeburt in meines Vaters Hause zu erlangen.

292. Darum nennt er seinen Sohn Manasse Gott zu Ehren, darum daß ihn Gott hätte zu nichts gemacht und es dahin gebracht, daß er alle seines Vaters Hauses vergessen. Dieweil er ihm aber wiederum eine Mehrung gönnt, da er ihm Kinder gegeben hat, so nennt er nun den andern Sohn Ephraim, welchem hernach Jakob, sein Großvater, einen schönen trefflichen Segen geben wird. Und es ist auch wahr, daß dieser Stamm der gewaltigste gewesen; wie dasselbe in den Büchern Josua, der Richter und Könige zu sehen ist. Jerobeam ist von dem Stamme Benjamin gewesen, und wird auch daher in den Propheten das Reich Israel also genannt, unter welchen Propheten schier die aller trefflichsten aus diesem Stamme geboren sind, als, Elia, Elisa, Hosea und andere mehr.

293. Wir sollen aber hier lernen, daß die Tödung pflegt vor der Erhöhung her zu gehen. Denn wir müssen zuvor Manasses, das ist, zu nichts gemacht werden; wiewohl das Fleisch solche Tödung nicht ohne große Schmerzen dulden und ertragen kann. Wie denn die Natur überall die Verwesung und ihr Verderben flieht; und wird auch nicht ein Baum oder Strauch umgehauen, daß er nicht einen Schall von sich gebe. Und Christus ruft selbst auch am Kreuze: „Mein Gott, mein Gott, warum hast du mich verlassen?“ Matth. 27, 46.

294. Darnach sollen wir uns bescheiden, daß wir der herzlichen Demuth, so an Joseph gewesen, auch nachfolgen. Denn er hat sich der großen Ehre, der Herrlichkeit und des Glücks gar nicht überhoben und damit nicht stolziret, sondern ist sein demüthig geblieben.

295. Es ist übrigens seine Regierung in dem wunderbar gewesen, dieweil er die Egypter hat lassen in der Vorhaut bleiben, und seine zwei Söhne, Ephraim und Manasse, beschneiden lassen und seines Vaters Gesetz gehalten. Denn also ist er jedermann allerlei geworden, daß er ja Etliche möchte selig machen, wie Paulus von sich selbst auch sagt, 1 Cor. 9, 22. Denn er hat seines Vaters Hause dermaßen nicht vergessen, daß er damit auch zugleich der Verheißung, von Gott gegeben, sollte vergessen haben; sondern dieweil er den Heiligen Geist, und dazu auch große Weisheit und Verstand in geistlichen Sachen gehabt, hat er die Verheißung von Christo und seinem zukünftigen Reiche mit höchstem Fleiße behalten. Darum hat er im Gefängniß den Tod gelitten und hat sich doch den Glauben und Hoffnung der Auferstehung nicht nehmen lassen. Was aber sonst für andere Stücke im Gesetze in seines Vaters Hause sind gehalten worden, und was die Erstgeburt und leiblichen Verheißungen betrifft, dieselben hat er vergessen und nur die Verheißung des ewigen Lebens behalten.

296. Joseph ist ein Mann gewesen voll Weisheit und geistlichen Verstandes, und hat den Unterschied wissen zu halten, davon Sanct Paulus sagt 1 Cor. 9, 20.: Den Heiden bin ich geworden, als ein Heide, den Juden, als ein Jude etc. Ich zwar hätte so einen großen Geist oder Glauben nicht gehabt, würde auch Egyptenland nicht also haben regieren können, daß ich ihre Sitten oder Weise und weltliche Rechte, und vielleicht auch wohl einen Theil ihrer Religion, welcher doch nicht gänzlich der wahren Erkenntniß Gottes zuwider gewesen, hätte liegen lassen und hintan gesetzt, da er für sich und seine Kinder eine sonderliche eigene Religion gehabt. Ich hätte gesagt: Ich kann die unbeschnittenen Heiden nicht regieren. Wie die Apostel Petro vorwerfen, Apost. Gesch. 11, 3., und sagen: „Du bist eingegangen zu den Männern, die Vorhaut haben, und hast mit ihnen geessen“, oder wie Etliche lehren, Apost. Gesch. 15, 1.: „Wo ihr euch nicht beschneiden lasset nach der Weise Mosis, so könnet ihr nicht selig werden.“

\quellseite{372}

297. Er hat aber einen apostolischen Geist gehabt; wie Petrus, da er durch ein Gesicht vermahnt war, sagt: „Ihr wisset, wie es ein ungewohnnet Ding ist einem jüdischen Mann, sich zu thun oder zu kommen zu einem Fremdling; aber Gott hat mir gezeiget, keinen Menschen gemein oder unrein zu heißen“ etc., Ap. Gesch. 10, 28. Joseph ist nicht allein zu den Unbeschnittenen gegangen, sondern hat sie auch regiert; welches der Apostel seiner würde gethan haben, ehe ihnen der Heilige Geist gegeben worden ist. Und Joseph hätte es auch nicht gewagt, wo nicht ein prophetischer und apostolischer Geist in ihm gewesen wäre.

298. So hat nun bis daher Mose die Historie von den zwei Söhnen Josephs mit eingeführt. Denn sie haben von Jakob sollen zu Kindern angenommen werden und sollten dazu Fürsten zweier Stämme werden. Darum werden sie an diesem Orte erwähnt, auf daß also die Verheißung, so Jakob gegeben war, erfüllt werden möchte: „Du sollst wachsen und gemehret werden“ etc.

\subsection{II}

\begin{quote}
\textbf{V. 53—56. Da nun die sieben reichen Jahr um waren im Lande Egypten; da fingen an die sieben theuren Jahr zu kommen, da Joseph von gesagt hatte. Und es ward eine Theurung in allen Landen, aber in ganz Egyptenland war Brod. Da nun das ganze Egyptenland auch Hunger litt, führte das Volk zu Pharao um Brod. Aber Pharao sprach zu allen Egyptern: Gehet hin zu Joseph, was euch der saget, das thut. Als nun im ganzen Lande Theurung war; that Joseph allenthalben Kornhäuser auf, und verkaufte den Egyptern. Denn die Theurung ward je länger je größer im Lande.}
\end{quote}

299. Wir haben droben gesagt, daß das hebräische Wort lechem nicht allein Brod heiße, sondern insgemein allerlei Speise; welches denn auch an diesem Orte klärlich zu sehen ist. Es hat aber in der lateinischen Uebersetzung Einer, der sich dünken lassen, daß er auch etwas wüßte, für das Wort panis, Brod oder Speise, das Wort fames, das ist, Theurung, gesetzt. Denn daran ist kein Zweifel, Hieronymus wird es im Lateinischen verdolmetscht haben, esca, das ist, Speise. Es war damals im Egyptenlande noch keine Theurung, sondern war noch ein reicher Vorrath da, den Mangel an der Speise damit zu erstatten. Hernach aber folgt im Texte: „Da nun das Egyptenland auch Hunger litt“ etc.

300. Es ist aber an diesem Orte zu merken, daß Joseph die sieben reichen Jahre über nicht alles Getreide im ganzen Egyptenlande gesammelt habe. Denn es haben auch Privatpersonen, so seinem Wort geglaubt, daheim ein Jeder in seinem Hause, ohne die gemeinen Kornhäuser des Königes, so im Lande hin und wieder aufgerichtet waren, ihre eigenen Scheunen oder Speisekammern gebaut. Und daß gesagt wird, Joseph habe das Getreide in Egyptenlande gesammelt, ist synedochisch geredet, da man das Ganze für einen Theil versteht, wie wir droben vom fünften Theil des Getreides gehört haben. Denn was den gemeinen Unterthanen von dem fünften Theil übergeblieben ist, dasselbe haben sie auch in den sieben reichen Jahren zum Vorrath behalten.

301. Derhalben werden sie zum wenigsten Ein Jahr lang Speise gehabt haben, damit sie ihr Leben haben erhalten können, da die gemeinen Kornhäuser noch nicht aufgethan gewesen, und sind auch mit der Theurung nicht ganze sieben Jahre, sondern auf das höchste fünf oder sechs beschweret worden. Denn daß ein jeder Bürger oder Bauer für sich auch etwas gesammelt, hat in einem oder zwei Jahren sehr viel geholfen. Es hat übrigens die Theurung die Lande, so nächst umher gelegen, bald gedrückt, dieweil sie sich gegen die künftige theure Zeit nicht versorget hatten: aber den Egypten hat die Erinnerung Josephs von der künftigen Theurung sehr viel geholfen. Darum haben sie die Theurung noch so bald nicht gefühlet, da schon die umliegenden Länder Mangel gehabt. Denn in Egypten war noch ein Bethleem, das ist, ein Brodhaus oder Vorrath von Speise. Im dritten und vierten Jahr aber, da hat endlich die Theurung überhand genommen; wie Joseph hernach sagen wird im 45. Cap. V. 6.: „Es sind noch fünf Jahr, daß kein Pflügen noch Ernten sein wird“, und dieselbige Theurung wird immer größer geworden sein, und in den letzten Jahren wird sie ohne Zweifel am allerbeschwerlichsten gewesen sein. Darum hat das Volk aus großer Noth bitten müssen, daß man ihm wollte Getreide lassen aus des Königes Kornhause.

302. Und da hat man öffentlich sehen können, wozu Josephs Rath nütze gewesen, und wozu es ihnen gedienet, daß der König selbst auch fromm und gottelig war, welcher mit aller Treue und höchstem Fleiß dafür sorgte, daß seinen Unterthanen durch diesen Mann möchte geholfen und gerathen werden, denn er das ganze Regiment auferlegt und befohlen hatte; wie er sie auch hat heißen zu Joseph gehen, und thun, was ihnen derselbe sagen würde.

303. Das ist aber eine große schwere Last gewesen, so Joseph die ganze Zeit über getragen, und hat ihm schier größere Mühe und Arbeit gegeben, das Getreide auszutheilen in den sieben Jahren der Theurung, denn zuvor, da er es gesammelt. Denn die Welt ist zu allen Zeiten gleich böse, und in solcher großen beschwerlichen Theurung, da so viel arme Leute gewesen, die Mangel gelitten, hat es an Dieben und Räubern auch nicht gefehlt; wiewohl er nicht allein in der Zeit der Theurung, sondern auch, da die reichen Jahre noch gemähret, das Volk immer fleißig ermahnt und in vielen vortrefflichen Predigten bei ihnen angehalten, daß sie doch die Wohltaten Gottes erkennen wollten, sich vor ihm demüthigen und Buße thun wegen der Theurung, so vorhanden war. Der Pöbel wird aber solches ohne Zweifel verachtet haben, und inzwischen nichts desto weniger im Saufe gelebt, geschlemmt und geprasset haben; ja, es würden auch wohl die Gottsäligen nichts vom Getreide gesammlet haben, wo sie nicht seinem Worte geglaubt hätten.

304. Er wird aber bei und neben der Lehre auch das Gebet fleißig getrieben haben, daß er für die Leute Gott angerufen, daß ihnen möchte geholfen werden. Und hernach, da die Noth mit der Theurung herangekommen, hat er sie erinnert, daß die gegenwärtige Strafe nun übereinstimme mit dem, was er ihnen zuvor gesagt hätte, auf daß sie desto mehr wollten darauf denken, wie sie sich möchten ihr Leben und Wandel bessern, und die rechte Lehre von der Erkenntniß Gottes annehmen und sich dieselbige gefallen lassen. Und er hat nicht allein selbst gelehrt, sondern durch die Amtleute oder andere Diener das Volk auch unterrichten lassen, nicht allein in Egyptenland, sondern auch in umliegenden Ländern allernächst bei Egypten. Denn er ist ein frommer, gottseliger und fleißiger Mann gewesen beide das weltliche und Kirchregiment zu verwalten, und hat solches beides mit Beten und Predigen ausgerichtet, daß Alles möchte gottselig und nach GOttes Wort und Willen zugehen; wiewohl Viele das alles werden verachtet und ihr Gespött damit getrieben haben.

305. Da aber nun die Theurung über ganz Egyptenland gegangen ist, und die Kornhäuser, so ein Jeder für sich gehabt, leer geworden, that Joseph alle Kornhäuser auf (wiewohl im hebräischen Text das Wort „Kornhäuser“ ausgelassen ist; ich weiß nicht, wie dasselbe mag genommen sein). Es ist aber also zu verstehen, daß er hin und wieder in einer jeglichen Stadt die Kornhäuser aufgethan und das Getreide daraus verkauft habe. Denn das Königreich Egypten ist sehr wohl bestellt und Alles darin wohl geordnet gewesen durch Josephs Weißheit, wie droben ist gesagt worden. Es ist bei dem Kaufen des Getreides keine Unordnung gewesen, daß alle Städte hätten müssen Getreide holen nur in Einer Stadt, sondern an allen Orten sind Kornhäuser gewesen, und Joseph hat in einer jeglichen Stadt Eines lassen aufthun, und die andern dieweil zugelassen, bis in nachfolgenden Jahren die Nothdurft wiederum erfordert, daß aus denselbigen den Leuten zu ihrer Speise auch Getreide ist gereicht worden. Denn er hat denen im Lande erstlich geholfen, damit sie nicht Hungers stürben, da inzwischen viel Andere im Lande Canaan und den nächst umliegenden Ländern darüber gestorben sind, ehe sie daran gedacht, daß sie von Joseph Getreide kaufen wollten.

\begin{quote}
\textbf{V. 57. Und alle Lande kamen in Egypten zu kaufen bei Joseph; denn die Theurung war groß in allen Landen.}
\end{quote}

306. Hernach, 1 Mos. 45, 7., wird Joseph zu seinen Brüdern sagen, wie ihn Gott nach Egypten gesandt, daß er sie und ihre Nachkommen beim Leben erhalten möchte. Wiewohl nun die Gottlosen und Sünder umgekommen, sind doch die Frommen, sein Vater Jakob und dazu seine Brüder, erhalten worden.

\quellseite{373}

307. Das hebräische Wort, scheber, will ich den Grammatikern befehlen; denn wir haben im Mose dasselbe Wort bis daher nicht gehabt. Das Wort bar hat vielerley Bedeutung. Es heißt, Sohn, Ausvermähler, rein oder unbefleckt, und heißt auch wohl Getreide, wie an diesem Ort. Warum aber das Wort scheber also ausgelegt wird, daß es so viel heiße als verkaufen, weiß ich nicht. Es heißt sonst, zerbrechen, oder etwas, das zerknirscht ist. Ich halte dafür, es sei ein solch Wort oder Weise zu reden, wie in der deutschen Sprache ist, da man sagt: Er hat weder zu beißen noch zu brochen, vom Brechen. Also sind auch diese gekommen zu brechen, das ist, zu kaufen, daß insonderheit die Materie verstanden werde für dasjenige, so man daraus macht. Wie man das Wort „Eisen“ auch pflegt für „Schwert“ zu gebrauchen. Aber das will ich den Hebraisten befohlen sein lassen.

\chapter{Das zweiundvierzigste Capitel}

\section[Erster Theil]{Erster Theil. Wie Jakobs Söhne nach Egypten reisen}

\begin{quote}
\textbf{V. 1. 2. Da aber Jakob sahe, daß Getreide in Egypten feil war, sprach er zu seinen Söhnen: Was sehet ihr euch lange um? Siehe, ich höre, es sei in Egypten Getreide feil; ziehet hinab, und kaufet uns Getreide, daß wir leben und nicht sterben.}
\end{quote}

1. Alles, was hier geschieht, geht dahin, daß es eine Gelegenheit sein soll, dadurch Joseph zu Ehren gesetzt werde, und daß damit die Wunderwerke Gottes angezeiget werden. Und ist die Historie an sich leicht genug, außer daß man etliche Stücke, so zur Grammatik und auch zur Lehre der heiligen Schrift gehören, handeln und erklären muß. Mose hat am Ende des vorigen Capitels gesagt, wie die Theurung in allen Landen sei groß gewesen, und daß darum auch aus allen Landen Völker herzugekommen seien, in Egypten Speise zu kaufen. Inzwischen aber hat Jakob mit den Seinen auch Hungers halben Noth gelitten. Darum redet er seine Söhne an, wie im Text steht, und ermahnt sie, daß sie auch hinab ziehen, Getreide zu kaufen.

2. Es kommt aber hier abermals vor das hebräische Wort scheber, welches wir droben mit „brechen“ verdolmetschet haben, und lauten die Worte an diesem Ort nach dem hebräischen Text also: Er sahe, daß in Egypten ein Brechen war. Wir aber, die der hebräischen Sprache unerfahren sind, müssen uns an diese sonderlichen Weisen zu reden gewöhnen, und an die neuen Worte, daß wir ihre Bedeutung und Eigenthümlichkeit recht verstehen lernen.

3. Die Juden machen ein Geschwätz und sagen, es heiße darum ein Brechen, daß es den Hunger breche oder stille. Ich bin aber solchem losen Geschwätz und der Leichtfertigkeit der jüdischen Rabbinen ganz feind, welche für und für versuchen, uns ihre Possen mit Gewalt aufzudringen, daß wir dieselbigen für die Grammatik und rechte Lehre der heiligen Schrift halten sollen, unangesehen, was sonst die heilige Schrift für sonderliche Art und Weise zu reden hat, daraus man die rechte Auslegung dieses Worts und anderer mehr nehmen kann. Denn also reden die Propheten, als, Jesaia Cap. 58, V. 7.: „Brich dem Hungrigen dein Brod.“ Und Jeremia in seinen Klageliedern Cap. 4, 4.: „Die jungen Kinder heischen Brod, und ist niemand, der es ihnen breche“, das ist, austheile. Denn die Juden haben Brod gehabt wie Kuchen, und haben keiner Messer bedurft, dasselbe Brod damit zu schneiden. Und ist eben solches auch in der Historie des Evangeliums Luc. 24, 30. 35. zu sehen, da die Jünger Christum erkennen an dem Brodbrechen, item, da von ihm Luc. 22, 19. gesagt wird, er habe das Brod genommen und gebrochen. Darnach ist es auch auf die Seele gezogen worden, als im 14. Psalm V. 3.: Es ist in ihren Wegen eitel Zerknirschen und Unglück. Item, ein „zerbrochener Geist“, das ist, ein betrübter, elender Geist, der hin und her in viel Gedanken und Rathscläge getheilt ist, da man nicht weiß, wo aus.

4. Derhalben soll man diese Art und Weise zu reden merken und behalten, nämlich, daß bei den Hebräern „Brod brechen“ eben so viel ist, als austheilen; oder aber, daß man hin und wieder Brod oder Speise kauft und verkauft, und ist von einem Gleichniß hergenommen worden. In den Propheten aber pflegt es öfter vorzukommen, denn sonst an andern Stellen der Schrift. Salmon in seinen Sprüchen am 11. Cap. V. 26. hat auch diese Art zu reden gebraucht, da er sagt: „Der Segen kommt über den, so Korn verkauft“, das ist, der es austheilt. Denn die Hebräer brauchen nicht allein eigentliche, sondern auch solche Worte, so durch sonderliche Figuren geredet werden; wie dasselbe in andern Sprachen gebräuchlich ist. Ja, es scheint, daß sie mehr als Andere Lust haben, solche figürliche Rede zu gebrauchen.

5. Danach ist das auch zu merken, daß die Hebräer das Hülfzeitwort „ich bin“ nicht haben, und an desselben Statt das Fürwort „ich“ gebrauchen. Als: „Ich der Herr, dein Gott.“ Da wird das „Ich“ gesetzt anstatt des Zeitworts: „Ich bin.“ Item, das Wort „selbst“ wird etwa gesetzt anstatt des Zeitworts: „Er ist.“ Als, der Herr dein Gott, der ist es. Darum wird auch das hebräische Wort, jesch, an dieser Stelle nicht als Hülfzeitwort gebraucht, sondern es heißt so viel, als: Es ist vorhanden, oder, es ist reichlich da, Getreide ist übrig genug da. Daher kommt das Wort tuschijah Spr. 8, 14., da die Weisheit sagt: „Mein ist beide Rath und That“, ich gebe Rath und auch tuschijah, das ist, die That, das Glück und die Ausführung in einer jeden Sache, daß es wohl geraten und vollzogen werden möge. Als, Gott gibt uns im mündlichen Worte Rath, und berichtet uns, wie wir uns halten sollen: wenn wir nun dasdselbige Wort festhalten, so folgt auch die tuschijah, daß uns dasjenige, so er verheißt, auch gewißlich gegeben wird und daß es nicht fehlt. Also auch an diesem Orte sind es nicht bloße Worte oder ein ungewisser Geschrei von dem Brechen oder Austheilen des Getreides, sondern es ist die Wahrheit und die That auch dabei, und ist in Egypten Getreide und Proviant vollauf.

6. Darum sagt nun Jakob zu seinen Söhnen: Warum seht ihr einander noch lange an? Denn das hebräische Wort ist ein Zeitwort der vierten Conjugation im Hithpael und lautet im Lateinischen also: Quare facitis vos videre? das ist, was sehet ihr euch lang um? Und Mose will damit anzeigen, daß Jakobs Söhne mit einander gerathschlagt und disputirt haben, was zu thun sei, dieweil die Theurung Ueberhand genommen. Du Ruben, du Juda, werden sie gesagt haben, was giebt du für Rath, daß wir Hülfe in dieser Theurung finden, und daß wir unser eigenes und auch alles unsers Hausgesindes Leben erretten mögen? Da sie aber also gerathschlagt und einander angesehen haben, hat der Vater gesagt: Ach, wie sehet ihr und seht Einer den Andern an? Dieweil ihr also vergeblich rathschlagt, müssen wir Hungers sterben. Siehe, ich habe gehört, daß in der Wahrheit Getreide genug auszutheilen und zu verkaufen vorhanden ist. Ziehet hinab und kauft uns Getreide. Gehet, daß ihr derselbigen Austheilung auch mit genießen möget.

7. Es zeigt aber Mose an, daß diese elf Patriarchen alle mit ihrem Vater Jakob zu Hebron gewohnt haben, da sie noch ehrlich und Hausväter gewesen sind, Kinder und Kinderkinder, dazu Knechte und Mägde gehabt haben. Derhalben ist es ein Wunder, wie doch so viele Geschlechter, darin so viele Personen gewesen, in einem fremden Lande alle zugleich haben wohnen können? Ich halte aber dafür, daß zumal Lea und die andern alten Matronen im Geschlechte Jakobs gestorben gewesen, und daß andere junge Weiber an ihre Statt genommen waren, nämlich, seiner elf Söhne Ehefrauen, desgleichen auch andere Knechte und Mägde. Benjamin, der unter ihnen der jüngste gewesen, hat zehn Söhne gehabt, und wir werden hernach hören, daß Jakob mit siebzig Seelen hinab nach Egypten gezogen sey, und da sind nicht mitgezählt worden die Weiber, Knechte, Mägde und andere Leute mehr aus den Heiden, die sich zu diesem Geschlechte gehalten um der Lehre willen von der rechten Göttlichkeit und wegen des rechten Gottesdienstes. Denn es werden allein die Seelen gezählt und gerechnet, so aus den Lenden Jakobs gekommen waren, als da sind seine Söhne und Enkel; die Enkelinnen aber werden nicht mitgezählt, und der andere Haufe von Knechten und Mägden wird auch nicht klein gewesen sein.

\quellseite{374}

8. Warum haben sie nun alle bei einander gewohnt zu Hebron als Gäste und Fremdlinge unter den cananitischen Völkern. Und dasselbe ist heutiges Tages bei uns sehr wunderlich und schier unglaublich, daß Jakob mit so vielen Söhnen und seiner Söhne Söhnen im fremden Lande unter den Heiden hat wohnen können, und daselbst Aecker und Wiesen miet(h)en, Gärten haben und Vieh halten. Und ich kann mich auch nicht genugsam verwundern, wie sich doch diese elf Brüder und das große Gesinde mit den Leuten derselbigen Orts haben vertragen können. Es wird nicht erzählt, wie die Heiden mit ihnen umgegangen sind; aber doch scheint so viel angezeigt zu werden, daß sie sich haben leiden müssen und Geduld tragen. Und vielleicht werden die guten Matronen, als, Lea, Silpa und Bilha, vor Herzeleid und Wehmuth gestorben sein. Denn sie haben viel Unrecht nicht allein von den Nachbarn, sondern auch von ihren eigenen Hausgenossen leiden müssen, dadurch sie, wie es scheint, gar niedergeschlagen, gebeugt und verzehrt worden sind; sintemal ihrer in dem Zuge, da Jakob nach Egypten hinab gezogen ist, gar nicht gedacht wird.

9. Warum sagt aber Jakob zu seinen Söhnen: Kauft uns Getreide, daß wir leben und nicht sterben? Warum vertraut er Gott nicht, und warum fürchtet er sich vor dem Tode, dieweil er doch eine göttliche Verheißung hat, daß ihn Gott erhalten und schützen wolle? Dazu hat er auch die Erfahrung, daß er es bisher also empfunden und erfahren hat, daß ihm Gott geholfen und ihn regiert hat im fremden Lande Canaan mit so vielen Söhnen und Enkeln, und mit so großem Hausgesinde. Warum erschrickt er denn hier, und glaubt nicht den Verheißungen Gottes, die er ohne Zweifel im frischen Gedächtniß gehabt, und die er auch in der Kirche, oder Gemeinde, so er in seinem Hause gehabt, immer fleißig gelehrt hat? Denn er wird zu seinem Völkchen gesagt haben: Werdet nicht kleinmüthig in der Anfechtung oder Trübsal; der HErr, unser Gott, hat uns das Land und dazu auch Speise nach unserer Nothdurft verheißen. Wo ist denn nun, lieber Jakob, dein Glaube? Wo ist die Verheißung?

10. Antwort: Wir haben den Befehl, daß wir glauben sollen und auf Gottes Güte trauen: aber wir sollen auch Gott nicht versuchen. Denn wir können nicht leben nach der Regel, die er hat, sondern nach Gelegenheit und wie es die Zeit gibt; jedoch, daß der Glaube und Hoffnung unverletzt bleibe. Darum sagt Jakob nicht: Bleibet da und wartet; der HErr ist mächtig genug, Brod vom Himmel herab zu senden; vielleicht wird er in diesem Lande Getreide regnen lassen und uns also ernähren. Nein, die Meinung hat es mit der Verheißung auch nicht. Es ist zwar daran kein Zweifel, daß dich Gott nicht sollte können oder wollen ernähren und schützen; du sollst aber die Gelegenheit und Mittel, dadurch dir Schutz und Hülfe widerfahren kann, nicht fliehen oder fahren lassen, welcher du ohne Sünde wohl gebrauchen kannst; denn sonst würdest du Gott versuchen.

11. Darum soll man nicht fragen, ob uns Gott nach seinem unabwendbaren Rath und Willen helfen und erhalten wolle; sondern dafür sollen wir es halten und sagen: Ich glaube, daß mich Gott erhalten wird; aber sein Rath ist mir unbekannt, dazu auch die Mittel, dadurch er seine Verheißung erfüllen wird. Derselben soll man die Mittel, so vorhanden sind, gebrauchen: ein Jeder soll seine Nahrung suchen mit Arbeit und Fleiß in seinem Beruf, und soll sich Nothdurft schaffen, das Leben damit zu erhalten, als, Milch, Käse, Wolle etc. Man soll den Acker bauen, die Frucht oder Getreide einsammeln; man soll nicht müßig oder faul sein, als ob du der Hausorge nicht bedürftest, oder als ob man das Gesinde nicht dürfte arbeiten heißen, dieweil du Gottes Verheißung hast, daß er dich ernähren wolle; denn Gott will nicht haben, daß du faul und träge sein sollst, sondern er hat 1 Mos. 3, 19. gesagt: „Im Schweiß deines Angesichtes sollst du dein Brodt essen, bis daß du wieder zu Erden werdest, davon du genommen bist.“ Item V. 18.: „Dornen und Disteln soll der Acker dir tragen“ etc. Ich sage dir zu, sprichst du, die Nahrung und Erhaltung deines Lebens; aber ich will, du sollst der Creaturen und Mittel gebrauchen, so vorhanden sind, und sollst thun, so viel du vermagst, daß du mich nicht versuchest. Wo du aber ohne alle Hülfe und Rath sein wirst, alsdann allererst will ich dich auch durch ein sonderlich Wunderwerk versorgen und erhalten. Ist aber irgend ein gut Mittel vorhanden, so denke, daß dir dasselbe von mir gezeigt und gegeben sei, darum daß du sein gebrauchen sollst, damit du mögest erhalten werden.

12. Auf diese Weise hat Jakob auch eine gewisse Verheißung, die ihm nicht fehlen kann; er thut aber doch, so viel an ihm ist. Er hört, daß in Egypten Getreide verkauft wird; deshalb hütet er sich davor, daß er Gott nicht versuche, wenn er würde dahin sitzen bleiben und sonst anderswoher auf Nahrung warten, sondern er sendet seine Söhne nach Egypten, daselbst Getreide zu kaufen. Was wollt ihr euch lange bedenken? sagt er zu ihnen. Es wird nicht Getreide oben herab aus den Wolken regnen. Geht hin, glaubt an Gott, vertraut ihm, und thut dabei auch, was an euch ist: gebraucht der Creaturen, die in diesem natürlichen Leben von Gott erschaffen und uns gegeben sind, daß wir derselben in unserer Noth gebrauchen sollen.

\begin{quote}
\textbf{V. 3. 4. Also zogen hinab zehn Brüder Josephs, daß sie in Egypten Getreide kauften. Aber Benjamin, Josephs Bruder, ließ Jakob nicht mit seinen Brüdern ziehen; denn er sprach: Es möchte ihm ein Unfall begegnen.}
\end{quote}

13. Hier kommt Jakob mit seinen Söhnen und Enkeln in Noth und Gefahr, und scheint es, als ob er Gott nicht vertraue. Denn er behält Benjamin, den elften Sohn, bei sich und will denselbigen mit den andern zehn Söhnen nicht in Gefahr dahin geben; welches ja der Verheißung und dem Glauben zuwider ist. Und nichts desto weniger glaubt Jakob der Verheißung.

14. Wie wollen wir denn dasselbe nun mit einander ausgleichen, daß es sich zusammen reimen möge? Der Glaube zwar zweifelt nicht; aber der gläubige Jakob zweifelt etwas. Wir haben droben eben dasselbe vom Patriarchen Abraham auch gehört, da er zu Gerar um der Sarah willen bekümmert war, daß ihr etwa ein Unfall begegnen möchte, und Isaac ist in gleicher Gestalt und Weise angefochten worden; wiewohl Beide klare gewisse Verheißung gehabt. Scheint es aber nun nicht, als hätten sie unsers Herrn Gottes gespottet, dieweil sie sich nicht ohne alle Furcht und Zittern an dieselbige Verheißung gehalten? Und Jakob hat oftmals gezweifelt und ist erschrocken, wenn er Esau, seinen Bruder, gesehen hat: da hat er sich nicht anders gehalten, denn als hätte ihm Gott nie etwas verheißen. Und sind die heiligen Patriarchen also uns allenthalben gleich geworden. Denn es ist Jakob an diesem Ort so angst geworden, als ob er gar keinen Glauben hätte.

15. Ich halte aber dafür, daß die Ursache gewesen sei, daß sich der Patriarch Jakob also gefürchtet und daß ihm so angst gewesen, weil er gewußt, daß die Verheißung müßte verstanden und gezogen werden auf den Sohn, so von seinem vornehmsten Weibe, Rahel, geboren war. Denn da er den erstgebornen Sohn Joseph verloren, hat sich nicht geziemen wollen, daß der andere nun auch in Gefahr sollte gesetzt werden, damit er nicht der Nachkommen und des verheißenen Samens gar beraubt werden möchte.

16. Warum bedenkt er aber dies nicht vielmehr bei sich selbst, daß ihm Gott verheißen hatte, ewiglich einen Samen zu geben, und daß er derselben ihm seinen Sohn auch mitten im Tode und in der Hölle erhalten werde? Es ist wahrlich eine große Schwachheit, daß er Benjamin dabei behält unangesehen der göttlichen Verheißung, und die zehen Söhne, so ihm Bilha, Lea und Silpa geboren, in Gefährlichkeit des Todes dahin gegeben hat. Denn er hat für sich selbst keine gewisse Hoffnung haben können, daß sie gesund und unverletzt wiederum zu ihm kommen würden. Soll man denn nun solche menschliche und fleischliche Rathschläge, solchen Zweifel und Schwachheit nicht strafen?

17. Mich dünkt zwar nicht, dieselbige zu entschuldigen; denn sie wird uns darum vorgehalten, daß wir auch an den allergrößten Vätern Exempel sehen sollen beide des Glaubens und auch der Schwachheit, und dazu den Streit des Fleisches mit dem Geist. Denn das ist wahrlich ein Exempel der größten Schwachheit an dem Vater Jakob und seinen Söhnen, die zu einem unbekannten Volk gegangen sind, welche, obschon sie das Wort der Verheißung gehört, doch noch zittern, hoffen und sorgen, glauben und zweifeln noch immer dabei. Denn also pflegen wir alle stark und männlich zu sein zu der Zeit, wenn der Geist über das Fleisch siegt und triumphiert; aber wenn der Unglaube unsere Herzen anficht, da sinken wir bald dahin.

\quellseite{375}

18. Es ist gar ein menschlicher und närrischer Rath Jakobs, daß er Benjamin bei sich behält, damit er nicht umkomme: als ob er nicht daheim auch wohl sterben könnte? oder als ob Benjamins Wohlfahrt und Leben in Jakobs, seines Vaters, Gewalt stünde? Ei! daß sich der große Patriarch, und der so große Verheißung gehabt, so gar schwach erzeigen soll. Was meinst du, wie dir oder mir werde zu Muthe sein, wenn die Gefahr und Anfechtung daher geht?

19. Darum ist dies alles uns zur Lehre und Trost geschrieben, auf daß wir lernen, Gott beständiglich vertrauen, und daß wir gleichwohl auch Sorge tragen für unsere Nachkommen, und für das Amt, so uns von Gott befohlen. Denn man muß nicht also denken: Ich bin ein Pfarrherr oder Kirchendiener, und weiß, daß Gott für die Kirche sorgt; derhalben will ich nichts dazu thun, will auch für mein Amt gar keine Sorge tragen. Wie ein Hausvater auch nicht träge oder faul sein darf; sondern es soll jedermann wacker sein und in seinem Beruf fleißig arbeiten. Denn Gott verbietet die Arbeit nicht, und will nicht haben, daß wir sollen müßig gehen, wiewohl er mit seiner Gegenwart und nach seinem Willen Alles regiert. Aber das Fleisch soll und muß geübt und versucht werden mit Gefährlichkeit, Anfechtung und Schrecken, und mit dem ungewissen Ausgang unserer Arbeit. Jedoch muß die Verheißung und der Glaube unverletzt bleiben, welcher vorher gehen muß, und muß auch das Gebet und die Arbeit, so eines Jeden Beruf erfordert, immer damit geübt und getrieben werden. Und wiewohl solche Arbeit nichts zur Sache thut, will es doch Gott also haben, daß wir im Schweiß unsers Angesichts unser Brod essen sollen, 1 Mos. 3, 19.

20. Ason, das hebräische Wort, heißt Tod oder sonst ein Schade, als, 2 Mos. 21, 22. steht: „Wenn sich Männer hadern, und verletzen ein schwanger Weib, daß ihr die Frucht abgehet, und ihr doch kein Schade widerfährt“ etc. Da steht auch im Hebräischen das Wort ason, das ist, ein tödtlich Unglück oder Herzleid. Sonst wird es in der heiligen Schrift nicht viel gebraucht.

21. Jakob fürchtet sich, daß Benjamin umkommen möchte, daß er vielleicht von einem wilden Thier zerrissen, oder von Feinden möchte gefangen werden; wie er meint, daß Joseph auch also müsse umgekommen sein. Diese Gefahr hat er bei sich selbst also bedacht, und daß er sich also gefürchtet, ist nicht gar vergeblich geschehen. Denn er war umgeben von unbekannten Nachbarn und von Feinden im fremden Lande. Und es scheint, daß er sich davor auch gefürchtet habe, daß er die Verheißung zugleich mit dem Sohne, von dem er meint, er würde der Erstgeborne werden, verlieren möchte. Denn er wußte damals noch nicht, daß Joseph noch am Leben war und daß dieselbige Herrlichkeit der Erstgeburt ihm gebühren sollte; welches er hernach in Egypten lernen und erfahren wird.

22. So haben nun die lieben heiligen Väter auf solche Weise reiche Verheißungen gehabt, und ob sie schon gläubig gewesen, haben sie doch noch Schwachheit des Geistes und Fleisches gefühlt, damit sie noch immer zu thun gehabt, aber doch also, daß den Auserwählten und Gläubigen Alles hat müssen zu gute kommen, und daß ihr Glaube geübt und versucht würde, wachsen und zunehmen möchte. Jakob fürchtet sich vor den wilden Thieren oder daß sein Sohn möchte gefangen werden und umkommen; denn er hat die Verheißung in großen Ehren gehalten, nämlich, in dem Sohn von dem vornehmsten Weibe, auch so, daß er die Gesundheit dieses Sohnes der Wohlfahrt und dem Leben aller andern zehn Söhne vorgezogen hat.

\begin{quote}
\textbf{V. 5. Also kamen die Kinder Israel, Getreide zu kaufen, sammt andern, die mit ihnen zogen; denn es war im Lande Canaan auch theuer.}
\end{quote}

23. Sie sind nach Egypten nicht gezogen ohne Gefährten, sondern haben sich zu andern frommen und ehrlichen Nachbarn aufgeschlagen, welche sie für treu erkannt haben. Es sei nun gleich solches geschehen auf die Ermahnung ihres Vaters, oder aber aus ihrem eigenen Bedenken, damit sie fromme getreue Gefährten haben möchten, dieselbe Reise zu vollbringen, dieweil der Weg durch der Cananiter und Midianiter Land gefährlich war.

24. Nun ist aber dieser Rath und Hülfe auch menschlich, wird aber vom Heiligen Geiste darum geschrieben, daß wir lernen sollen, den Verheißungen Gottes zu glauben und nicht zweifeln. Wiewohl wir im Glauben schwach sind, sollen wir doch dasselbe immer täglich wiederholen, und uns selbst fleißig einschärfen, daß Gott in seinen Verheißungen wahrhaftig sei. Hernach aber, wenn etliche Gelegenheiten und äußerliche Mittel sich darbieten, sollen wir wissen, daß man derselbigen auch gebrauchen soll und daß menschlicher Rath nicht zu verachten sei. Denn darum sind die Kinder Israel diesen Weg nicht allein gezogen, sondern haben sich vereinigt mit den Hebronitern und Cananitern, auf daß sie solcher Gefährten haben möchten sicher sein. Und es ist wahrscheinlich, daß Viele denselbigen verwandt gewesen sein wegen der Heirath, so sie mit den Töchtern der Cananiter eingegangen waren. Denn solche Hülfe ist ihnen vonnöthen gewesen, und daß sie sich mit ehrlichen Leuten vereinigt und befreundet haben wider Gewalt und Unrecht, so sie als Gäste und Fremdlinge im Lande von bösen Buben haben erwarten müssen. Darum haben sie in Noth und Gefährlichkeit auch menschlichen Rath und Hülfe gebraucht.

\begin{quote}
\textbf{V. 6. Aber Joseph war der Regent im Lande, und verkaufte Getreide allem Volk im Lande.}
\end{quote}

25. Das hebräische Wort, schalith, ist auch ein neu Wort, welches wir bisher noch nicht gehabt. Es kommt vom Zeitwort schalath, das heißt, herrschen, und wird hernach oftmals in derselbigen Bedeutung in der heiligen Schrift gebraucht werden; als, im Prediger Salomo an 8. Cap. V. 9.: „Ein Mensch herrschet zu Zeiten über den andern, zu seinem Unglück.“ Daher kommt das Wort schalith, Herr; und die Saracenen nennen ihren Kaiser Sultan; und 1 Mos. 49, 10. steht: „Es wird das Scepter von Juda nicht entwendet werden, bis daß der Held (Sultan) komme.“ Daher kommt auch das Wort Schild, im Lateinischen scutum oder clypeus, welches Wort die Deutschen behalten haben. Derhalben ist Joseph der schalith, das ist, der Sultan über ganz Egyptenland, und ist der maschbir, wie es im Hebräischen lautet, der da Getreide verkauft. Von welcher Art und Weise zu reden, und von dem Fürwort, so anstatt des Zeitworts „ich bin“ gebraucht wird, droben geredet worden ist.

26. Das ist aber ein groß Lob der Weisheit und sonderlichen Fleißes an Joseph, welcher in einem so großen Königreich Alles allein regiert. Pharao ist auch ein Sultan, ist aber kein maschbir, kümmert sich darum nicht, wie das Getreide oder Speise ausgeteilt und verkauft werde; die Sorge und Zucht trägt Joseph allein. Es ist aber Wunder, wie er doch bei so großer Menge der Käufer Jeden habe können ansehen und kennen lernen. Denn unter den andern Haufen derer, so daselbst angekommen sind, hat er selbst seine Brüder erkannt. Daraus abzunehmen, daß er auf die Fremden und dazu auf Jeden insonderheit fleißig Acht gegeben haben müsse; dazu ihm denn sehr großer Fleiß, und daß er hat wacker sein müssen, vonnöthen gewesen.

27. Ich halte aber dafür, er werde darin diese Weise und Ordnung gehalten haben, daß er ein Kornhaus nach dem andern aufgethan, auf daß er so viel leichter auf einen Jeden könnte Acht geben, wenn die Käufer an einen bestimmten Ort haben kommen müssen. Und wiewohl er viel Amtleute und Diener gehabt, die das Getreide auf seinen Befehl verkauft und desselben immer haben warten müssen, so haben doch die Käufer alle vor ihn selbst kommen müssen, daß er ihnen erlaubte, dasselb Getreide zu kaufen. Es ist eine sehr schöne Ordnung gewesen und ein nöthiger Fleiß, daß nicht diejenigen, so des Königreichs Feinde gewesen, mit den andern Fremden und Nachbarn auch hereinschleichen und etwa einen Lärm oder Unfrieden im Lande anrichten möchten.

28. Dasselbe Exempel soll die Obrigkeit an einem jeden Ort mit Fleiß merken, und demselben auch also nachfolgen, daß sie bedenken, es gebühre ihnen zu wachen, und daß sie nicht allein auf ihre Bürger und Unterthanen, sondern auch auf die Ausländischen und Fremdlinge gute Acht haben, damit dieselbigen nicht also leichtfertig und ohngefähr hin und wieder irre gehen zu Gefahr dem gemeinen Wesen.

29. Denn also ist das Königreich Egypten bestellt und geordnet gewesen, es sei nun solche Ordnung von Joseph aufgerichtet, oder aber von ihm erst angefangen worden. Und wo sie Joseph erst gestiftet und gemacht hat, so sehen wir offenbar, daß der heilige Geist auch mit bei der leiblichen Regierung gewesen. Im Anfang seiner Regierung hat er in den sieben reichen Jahren die Kirchen bestellt, darnach auch die weltliche Polizei und Hausregiment. Jetzt, da nun das Getreide verkauft werden soll, sieht er, daß vonnöthen sein will, auch auf die fremden Leute, so ab- und zugegangen sind, Acht zu geben. Und mag vielleicht wohl ein böses ungezüchtigtes Volk gewesen sein, das eines solchen Aufsehens bedurft hat; darum hat er dafür auch gesorgt.

\quellseite{376}

\section[Zweiter Theil]{Zweiter Theil. Von dem Verhalten der Söhne Jakobs gegen ihren Bruder Joseph, und von dem Verhalten Josephs gegen seine Brüder}

\begin{quote}
\textbf{V. 6. 7. Da nun seine Brüder zu ihm kamen, fielen sie vor ihm nieder zur Erde auf ihr Antlitz. Und er sahe sie an, und kannte sie, und stellte sich fremd gegen sie, und redete hart mit ihnen, und sprach zu ihnen: Woher kommt ihr? Sie sprachen: Aus dem Lande Canaan, Speise zu kaufen.}
\end{quote}

30. Nun folgt endlich der rechte Ausgang und Erfüllung des Traumes, welchen Joseph vorher von seinen Brüdern gehabt. Und Mose eilt, daß er diese Erfüllung der Weissagung und Träume Josephs beschreiben möge. Es kennen aber die Brüder Josephs ihn noch nicht, wissen auch nicht, vor wem sie niederfallen zur Erde auf ihr Antlitz. Derhalben ist dies nur der Anfang und eine unvollkommene Erfüllung, da sie vor ihrem Bruder, den sie noch nicht kennen, auf die Erde niederfallen. Hernach aber werden sie auch vor ihm niederfallen, da sie ihn erkannt haben, nach Aufhebung ihres Vaters, auf daß ja dieser Ausgang mit den Träumen Josephs vollkommen übereinkomme.

31. Dies soll man aber nicht so obenhin ansehen oder lesen, sondern dieses angenehme Exempel und Schauspiel, oder vielmehr dieser schöne Spiegel dieser göttlichen Regierung soll fleißig gemerkt werden. Und ist hierin eine klare Allegorie oder heimliche Deutung und Figur Christi, welche sich der Leser auch zugleich vorstellen soll. Aber die vornehmste Ursache, darum der Heilige Geist diese Historie hat schreiben wollen, ist gewesen, daß er damit hat zeigen wollen ein wunderbarlich Schauspiel, und ein gewiß glaubwürdig Exempel der Regierung, wie Gott seine Heiligen pflege zu führen. Welches eigentlich mit dem Spruch des 4. Psalms V. 4. übereinstimmt, darin also steht: „Erkennet doch, daß der Herr seine Heiligen wunderlich führet.“

32. Denn hier sehen wir, daß Joseph mit seinen Brüdern wunderlich umgeht, und spielet mit ihnen ein wunderlich Spiel; aber doch ein solch Spiel, das sie sehr demüthigt und wohl übt, und große ernste Sachen mit sich bringt. Denn auf solche Weise werden sie zur Erkenntniß ihrer Sünden gebracht und in Gefahr des Todes und der Hölle geworfen.

33. Und eben auf solche Weise verhält sich auch Gott gegen seine Heiligen und Gläubigen in Anfechtungen, wie sich Joseph gegen seine Brüder gehalten, welcher nicht verfahren mit ihnen gespielt, daß er begehrte sich an ihnen zu rächen, oder daß er ihnen feind wäre: darum, sage ich, hat er sich gegen sie so hart nicht erzeigt. Denn er meint endlich und sein Herz im Leibe bewegt sich gegen sie. Aber er stellt sich äußerlich als ein Tyrann, der sie gar verderben und umbringen wolle um der schändlichsten Sünde willen der Verrätherei, damit sie die höchste Ungnade, ja auch den Galgen und allerlei Strafe verdient hätten. Dies thut er, sage ich, nicht aus untreuem oder feindseligem, sondern aus brüderlichem und freundlichstem Herzen, und aus keiner andern Ursache, denn daß er damit ihre Buße erforschen und prüfen will, und sie also zur Erkenntniß ihrer Sünde, und zu der Gnade und Barmherzigkeit Gottes treiben möge.

34. Eben ein solch Spiel treibt Gott in der Welt auch nach seiner göttlichen Güte, wie die Weisheit sagt in den Sprüchen Salomos am 8. Cap. V. 31.: „Ich spielte auf seinem Erdboden, und meine Lust ist bei den Menschenkindern.“ Also redet die Weisheit des Vaters, der Sohn Gottes, da er ist Mensch geworden, und zeigt damit an, daß er mit den Menschen spiele und seine Lust habe, und nur freundlich und lieblich mit ihnen umgehe, und daß er an solchem Spiele die größte Lust habe.

35. Uns aber ist es ein sehr betrübter elender Tod. Denn heißt das Spielen, den Menschen gar zu nichte machen, in den Tod werfen, mit viel unzähligen Jammer und Plagen betrüben? Es ist ein Spiel der Katze mit der Maus, welches der Maus Tod ist. Darum ist dies geschrieben, daß wir sollen verstehen lernen den Rath göttlicher Weisheit und wie wunderbarlich Gott seine Heiligen führt: welches so gar betrübt und bitter scheint, daß die Engel, die solches sehen, dazu die Teufel und Welt meinen, daß wir gar verschlungen und verdorben seien; ja, wir selbst auch lassen und dünken, wir seien gar verlassen, verachtet und von Gott verworfen.

36. Daher sind eigentlich diese Klagworte gekommen, als, im 31. Psalm V. 23.: „Ich bin von deinen Augen verstoßen“; und im 44. Psalm V. 24.: „Erwecke dich, Herr, warum schläfst du? Wache auf, und verstoße uns nicht so gar“ etc. Solch Schreien zu Gott wird Jakob ohne Zweifel die ganze Zeit seines Lebens oftmals wiederholt haben in so viel schwerem Jammer und Noth. Und dies ist das stetige Seufzen und Klagen aller Patriarchen vom Anfange gewesen. Denn was könnte doch Elenderes und Erbärmlicheres gesagt oder gedacht werden, denn die ersten Väter gewesen sind, als, Adam, Noah und Abraham?

37. Und man sieht doch am Ende desselben Spieles, daß Gott mit ihnen ganz freundlich gespielt hat, und ist nicht anders mit ihnen umgegangen, denn als mit seinen allerliebsten Kindern. Denn also lehrt die Epistel an die Hebräer am 12. Cap. V. 5. 6. 7.: „Mein Sohn, achte nicht gering die Züchtigung des Herrn, und verzage nicht, wenn du von ihm gestraft wirst. Denn welchen der Herr lieb hat, den züchtiget er; er stäupt aber einen jeglichen Sohn, den er aufnimmt.“ So ihr die Züchtigung erduldet, so erbeut sich euch Gott, als Kindern“ etc.

38. Aber wie es im Hausregimente zugeht, wo man die Kinder züchtigt, ist es ihnen gar nicht angenehm; also läßt man sich dünken, daß die Züchtigung des Herrn sehr hart, bitter und schwer sei. Aber, „selig ist der Mann“, sagt St. Jacobus Cap. 1, V. 12., „der die Anfechtung erduldet; denn nachdem er bewähret ist, wird er die Krone des Lebens empfangen, welche Gott verheißen hat denen, die ihn lieb haben.“ Darum sollen wir es gewißlich dafür halten, daß unser Unglück oder Trübsal, Seufzen und Klagen, und auch der Tod selbst nichts Anderes ist, als ein sehr lustig und schön Spiel göttlicher Güte mit uns.

39. Wie viel sind aber derer, die dies verstehen oder glauben? Wenn wir das könnten in unser Herz bringen, so wären wir wahrhaftig selig; alsdann wären wir auch bereit, mit fröhlichem Herzen alles Unglück zu leiden. Wenn ich aber daran denke, daß ich ein armer Sünder bin und um der Sünde willen gestraft werde, da urtheile ich denn viel anders. Denn da empfinde ich nicht, daß Gott mein Vater sei, daß er gütig und barmherzig sei, sondern daß er der Teufel selbst sei.

40. Darum sollst du wissen, daß Gott allmächtig ist, und ihm derhalben ein solch ernstes Spiel wohl gezieme und dasselbe sich mit solcher großen Majestät auch fein reime. Und man muß wahrlich lernen, und sich an solche seine Lust und Spiel gewöhnen, wie der 4. Psalm V. 4. saget: Ihr müßt lernen. Die Natur oder Vernunft lehrt es nicht; sonst wäre es nicht vonnöthen, daß ihr es von mir müßtet lernen. Der Psalm sagt: Höret, lasset euch sagen, wenn ihr mit mancherlei Beschwerung, mit allerlei Jammer und Noth angefochten werdet, so bedenkt, daß Gott mit euch spielet; welch Spiel euch zwar wunderlich, Gott aber lustig ist. Denn wo er nicht ein väterlich Herz zu euch hätte, würde er nicht also mit euch spielen.

41. Derhalben ist es ein Anzeichen der großen und unaussprechlichen Barmherzigkeit Gottes gegen dich, daß du unter die gezählet bist, zu denen Gott Lust hat, bei welchen seine Lust ist. Darum gibt er dir seine Verheißung, sein Wort und Sacramente als die gewissesten Zeichen und Zeugnisse seiner Gnade, nämlich, daß er dich zu seinem Kinde habe angenommen: und fordert von dir nichts mehr, denn daß du sein Spiel ertragen und dulden wollest, welches ihm lustig, dir aber heilsam ist.

42. Ja, sagst du, ich verstehe aber dies nicht? Antwort: Das ist recht. Denn der Heilige Geist nennt diesen Herrn nicht vergeblich einen Wunderbaren. Derhalben wird in dieser Historie den Gottlosen ein großer Trost und treffliche Lehre vorgehalten, wenn sie sehen, daß Joseph mit seinen Brüdern etwas hart und unfreundlich umgeht, nicht, daß er begehrte, sich an ihnen zu rächen, oder vorhätte, sie zu verletzen oder zu verderben, sondern vielmehr, daß er ihnen an Leib und Seele helfen und sie erhalten wolle. Denn sie waren fett, dick und starr, haßten die Züchtigung oder Strafe, und hatten ihres Vaters Rede verachtet, thaten keine Buße, sondern meinten, ihre Sünde, welche so lange Zeit geschlafen hatte, wäre nun ganz und gar vergessen. Aber es heißt: Alte Schuld rostet nicht. Es läßt sich wahrlich die Sünde nicht zudecken oder austilgen; sondern wenn du sie schon verleugnest oder verbirgst, so schläft oder ruht sie doch vor der Thür, auf daß sie bald hernach geoffenbart werde, wie 1 Mos. 4, 7. steht: Bist du nicht fromm, so bleibt die Sünde nicht sicher noch verborgen.

\quellseite{377}

43. Darum geht er nun erstlich damit um, daß er sie zur Erkenntniß ihrer Sünden bringen möge. Darum versucht und erforscht er ihre Herzen wunderbärlich, auf daß er möge erfahren, wie sie doch gesinnet sind. Und dieweil er sieht, daß sie der begangenen Sünde halben noch nicht Buße thun, versucht er sie so lange, bis er sie unterrichte und ihrer Sünden erinnere, und sie also zur Buße treibe.

44. Man soll aber hier nicht disputiren, wie die Papisten zu thun pflegen, von der Genugthuung, als ob Gott für die Sünde eine Genugthuung von uns forderte, die wir ja gar nicht leisten können. Denn es hat Joseph dieselbe von seinen Brüdern auch dermassen nicht gefordert, da er sie doch, wo er hätte streng mit ihnen verfahren wollen, mit höchstem Rechte hätte tödten mögen um deswillen, daß sie ihm so große Gewalt und Unrecht gethan. Aber das thut er nicht. Denn er weiß, daß im menschlichen Geschlechte keine andere Genugthuung ist für die Sünde, denn allein die da weltlich oder bürgerlich ist, welche die Obrigkeit fordert von denen, so übertreten und übel gehandelt haben. Aber diese Genugthuung hat mit der Theologie und Lehre der heiligen Schrift gar nichts zu thun. Denn darin ist nur das einige Opfer, nämlich, unser Herr Christus, mit welchem dem Zorne Gottes genug gethan worden. Denn unsere Sünden sind viel größer, denn daß sie könnten durch einige Genugthuung bezahlt oder ausgetilgt werden.

45. Darum sollte man das Wort „Genugthuung“ aus der Kirche hinweg thun, oder aber desselbigen mit größerer Vorsicht gebrauchen, als von Cyprian und Andern geschehen ist, zu großem Schaden der Kirche. Es ist die Vergebung der Sünden viel ein höher Ding, denn unsere Genugthuung sein kann. Denn wo dieselbige schon geleistet und gehalten wird, bleibt doch die Sünde noch da.

46. Derselben versucht nun Joseph seine Brüder so viel härter, daß er sie damit zur Buße und Erkenntniß der Sünden treiben möge und an ihnen die Sicherheit tödte, und sie zur Verzweiflung treibe, da er ihnen mit dem Tode und sonst allerlei Strafen droht: nicht, daß sie auf diese Art genugthun für die Sünde, wie die Papisten ihr unnützes Gewäsche treiben; sondern daß er in ihnen Schmerz und Reue erwecken möge wegen ihrer begangenen Sünde, auf daß sie also selbst Hülfe suchen, wie sie von Sünden mögen gereiniget werden.

47. Also betrübt uns Gott auch mit mancherlei Unglück, nicht, daß er uns damit ernstlich strafen wolle; wiewohl es in Wahrheit eine Strafe ist, er hat aber daran keine Lust. Was meint er aber denn mit so viel Jammer, Plagen und Krankheiten? \&c. Antwort: Er geht nur damit um, daß du zur Erkenntniß deiner Sünden mögest gebracht werden. Er weiß wohl, du kannst dafür nicht genug thun, er straft uns auch nicht nach unserm Verdienste. Denn wir verdienen nichts Anderes, denn den Tod und die Hölle. Aber die Sünde, so in unserer Natur steckt und uns anhängt, ist vor unsern Augen verborgen; dieselbige bringt er an das Licht, wie im Gebet Mosis steht, im 90. Psalm V. 8.: „Unsere Missethat stellest du vor dich, unsere unerkannte Sünde in das Licht vor deinem Angesichte“, als wollte er sagen: Du siehst unsern Jammer und Unglück, welches wir selbst nicht sehen. Wir sind vor deinem Angesichte unrein und aufs schändlichste befleckt; es ist aber solche Mißgestalt und Schande unsrer verderbten Natur uns selbst unbekannt. Derhalben, auf daß sie möge geoffenbart und gereiniget werden, so braucht Gott dazu solche gewaltige und bittere Arznei. Und soll er das Uebel ausfegen, so muß er einen Strohwisch und scharfen Sand nehmen und scheuern, daß das Blut kommt.

48. Hierher gehören nun die Strafen und das Unglück, so auf dem menschlichen Geschlechte liegen, als Pestilenz, Krieg und theure Zeit, auf daß damit die Sünde an uns gestraft und geoffenbart werde, daß man es greifen möge, bis wir es verstehen lernen, was wir vor Gottes Augen sind. Ja, wir müssen auch etwa einen sehr schändlichen großen Fall thun, daß wir unsern Jammer und Schwachheit erkennen lernen. Also ist David erschrecklicher Weise gefallen wegen seiner Sicherheit und Hoffahrt, auf daß er damit auch unterrichtet würde, was Sünde sei, und das von Herzen sagen könnte Ps. 51, V. 6.: „An dir allein hab ich gesündigt“; als wollte er sagen: Ich sehe nun wohl, daß wir alle in deinem Gesetze verklagt und verdammt werden; aber da ist kein Mensch, der das verstehe. Ich aber bekenne nun, daß ich nichts mehr denn nur ein armer Sünder bin, der nichts denn eitel Zorn verdient und des ewigen Todes werth ist; kurzum daß ich nichts Anderes bin, denn nur eitel Sünde: „auf daß du Recht behaltest in deinen Worten, und rein bleibest, wenn du gerichtet wirst“ von denen, die fromm und gerecht sein wollen.

49. So lasset uns nun lernen an diesem Tröste fest halten, wiewohl die Schuld unser Gewissen beißt, der Stachel des Gesetzes und Todes uns treibt; daß wir wissen, solches geschehe nicht zu unserm Verderben, sondern vielmehr zu unserer Lehre und Unterrichtung, und daß wir endlich von Sünden gereiniget werden mögen, daß wir zur Erkenntniß unser selbst und unsrer verderbten Natur kommen mögen. Darum sollen wir die Hand Gottes dulden, der uns also fegt, das ist, der mit uns lacht und spielet. Denn die sich weigern, die Hand Gottes zu dulden, dieselben disputiren von der Genugthuung, und wenn dieselbige darnach nicht genugsam ist, alsdann müssen sie nothwendig verzweifeln. Es fegt zwar Gott die Heuchler und Gottlosen auch wohl; aber ehe die Reinigung gar vollzogen wird, zerbrechen sie, wie ein Glas zerbricht.

50. Also haben vor Zeiten die Sophisten gelehret, man müsse sich so lange bemühen und es sich sauer werden lassen, bis die Reue über die Sünde genugsam wäre. Aber was hilft es, über die Sünde Reue und Leid tragen und nicht das kleinste Fünklein von Gottes Gnade und Barmherzigkeit verstehen? Denn also werden Saul und Judas auch versucht, geplagt und gefegt: aber sie haben nur Reue gehabt der begangenen Sünden halben; darum werden sie auch bald kleinmüthig und verzagt; denn da herrscht allein das Gesetz mit seinem Amte. Der Teufel hat auch seine Reue und zittert, wie Jacobus Cap. 2, V. 19. sagt; aber das ist keine Buße zur Seligkeit. Derhalben muß das auch dazu kommen, das der Herr sagt Hes. 18, 23.: „Ich habe nicht Gefallen am Tode des Gottlosen, sondern vielmehr, daß er sich bekehre von seinem Wesen, und lebe.“ In aller Noth und Schrecken muß das kleine Fünklein göttlicher Barmherzigkeit in unserm Herzen haften bleiben, dadurch das Herz, so über die Sünden Reue hat, wiederum aufgerichtet und getröstet werde; wie sich damit Manasse auch aufrichtet, da er sagt in seinem Gebet V. 8.: „Du hast die Buße nicht gesetzt den Gerechten, Abraham, Isaak und Jakob“, sondern den armen Sündern \&c.

51. Das soll man in diesem gegenwärtigen Exempel fleißig und wohl betrachten, auf daß wir die Ursachen unsrer Trübsal verstehen können. Denn wiewohl wir mit der That und äußerlichen Werken nicht sündigen und etlichermaßen äußerliche Zucht haben, ist doch das Fleisch, damit wir uns schleppen, unrein, und wir verstehen nicht, wie groß unsere Sünde und Jammer ist, welcher hinweg gethan werden muß. Es sind ja wohl alle Sünden vergeben und zugedeckt, aber sie sind noch nicht gar ausgefegt; und bleibt noch immer an uns hängen und kleben nicht allein die grobe Hefe der Unzucht, Hoffahrt, Hasses, Zorns und anderer bösen Begierden, sondern auch die innerlichen bösen Stücke und heimlichen Befleckungen, als da ist, Zweifel an Gott, Unglaube, Ungeduld, Murren wider Gott. Welches Unglück sich dann allererst offenbart und hervor thut, wenn das Gewissen durch das Gesetz und Schrecken der Sünde angegriffen und geplagt wird. Wir achten deßselbigen nicht, fegen auch nicht um der schändlichen Befleckung willen: Gott aber sieht es. Darum versucht er auch, unsere unreine Natur zu fegen, und denkt also: Du bist ja nun wohl erleuchtet und getauft; stinkst aber noch, und dein Fleisch ist voll großer Gebrechen, Mängel und Fehler; darum muß ich es reinigen und fegen. Denn was unrein und befleckt ist, wird nicht in das Himmelreich gehen, Offenb. 21, 27. Da hebt sich denn unsers HErrn GOttes Spiel an, nicht daß wir verderben oder verloren werden sollen, sondern daß er uns also versuchen, prüfen und zur Erkenntniß unser Schande bringen möge; doch also, daß wir nicht gar verzagen, sondern vielmehr zu ihm schreien und ihn um Gnade und Barmherzigkeit anrufen; dazu, daß wir auch lernen, daß er ein solcher GOtt sei, der seine wunderliche Güte also beweisen will, wie der 17. Psalm V. 7. sagt; denn dasselbe hat er gemeint mit seinem Spiele.

\quellseite{378}

52. Dies haben die Sophisten, wie droben gesagt, nicht verstanden. Denn sie haben also gelehrt: Wenn du vollkommene Reue hast, so giebst dir GOtt seine Gnade ein de congruo, das heißt, wenn auch nicht als verdienten Lohn, doch nach Billigkeit; und wollen also haben, daß der Mensch allein mit der Strafe GOttes Gnade verdienen könne. Das ist aber unrecht; sondern es muß der Glaube und rechte Anrufung dazu kommen, und daß man GOttes Güte und Verheißung empfinde; daraus denn kommt das Schreien und Seufzen, wenn das Gewissen geängstet und angefochten wird, daß sich damit auch zugleich der Geschmack des ewigen Lebens anfange.

53. Aber was soll ich rufen oder schreien, möchte Einer sagen. Ich fühle wohl, daß ich getödtet werde. Darauf antwortet GOtt und spricht Hes. 18, 23.: Ich will nicht haben den Tod des Sünders etc. Ja, was denn, du lieber GOtt? Ich will haben, daß du zu mir schreien sollst, auf daß ich dich wunderlich führen möge, und daß du deine Unreinigkeit verstehest, welche ich auf diese Weise wegnehmen will. Und ist zwar dieselbe Weise mir ein Spiel, dir aber ist es der Tod, und ist dir doch nütze zu wissen und zu erkennen, daß dir solche Reinigung noth thut. Das „erkennet“ im 4. Psalm V. 4. muß also gelernt sein. Es verhält sich aber mit denen viel anders, mit welchen ich dermaßen nicht spiele. Denn dieselben werden auch hin und her getrieben, werden versucht und geplagt; aber sie schreien nicht zu mir, denken auch nicht, daß ich mit ihnen spiele.

54. Derhalben werden wir dies Schauspiel an Josephs Exempel fein sehen. Denn wie GOtt gegen seine Heiligen oder Gläubigen gesinnet ist, also ist Joseph gegen seine Brüder auch gesinnet, voller Barmherzigkeit, herzlicher Liebe und freundlicher Neigung. Er sucht mit dem reinsten Herzen nichts Anderes, denn ihre Wohlfahrt und Seligkeit an Leib und Seele. Darum schreckt er sie zuerst, daß er sie also erinnern möge der greulichen Sünde, damit sie ihren alten Vater vorlängst herzlich betrübt, ihren Bruder aber in Gefahr Leibes und Lebens gebracht hatten. Und vielleicht mag er sich besorgt haben, Benjamin, sein Bruder, möchte durch ihre Bosheit auch umgekommen sein, oder daß der Vater vorlängst vor großem Herzeleid gestorben wäre.

55. Darum spielt er mit ihnen ein solch göttlich und scharf Spiel, welches doch sehr heilsam und zu GOttes Ehre ist. Denn sie werden damit zur Erkenntniß ihrer Sünde geführt, aber doch ohne Verzweiflung, und kommen also zu rechter Gottseligkeit und Heiligkeit. Denn sie sind nach solcher Reinigung besser und weicher geworden, da sie zuvor ihre Sünde verbergen, und dieselbe gar verschweigen und vergessen wollten, und also in ewigem Jammer und Verdammniß liegen bleiben. Derhalben ist vonnöthen gewesen solcher Arznei, die stark hat sein müssen, wie dieselbige in den Propheten beschrieben wird, als, Hosea Cap. 6, V. 5.: „Darum haue ich sie durch die Propheten“, spricht der HErr, „und tödte sie durch meines Mundes Rede“ etc.; und Hes. 24, 12. nennt GOtt die Stadt Jerusalem einen solchen Topf, darin das Angebrannte klebt und nicht abgehen will, wo es nicht zuvor wiederum verschmolzen wird etc. Da wird dies Spiel unseres HErrn GOttes beschrieben, dadurch wir geplagt und getödtet werden, bis wir wieder zur Erden werden, davon wir genommen sind, und begraben werden.

56. Warum handelt aber GOtt so ungnädig mit uns? Darauf antwortet er selbst im Propheten Hesekiel und sagt, daß der Topf oder das Angebrannte vom Topfe sonst nicht abgehen wolle, es ist zu hart angebrannt. Darum sagt er, daß er ihn im Feuer schmelzen wolle, auf daß er möge ausgekocht und gereinigt werden.

57. Also sagt er auch von der Zukunft Christi, Mal. 3, 3.: „Er wird sitzen und schmelzen, und das Silber reinigen; er wird die Kinder Levi reinigen und läutern, wie Silber und Gold.“ Darum setzet der Prophet das noch hinzu V. 2.: „Wer wird den Tag seiner Zukunft erleiden mögen?“ etc. So spielt nun GOtt mit uns ein väterlich Spiel, wenn er Peinlichkeit, theure Zeit, Krankheit, Traurigkeit des Herzens, und etwa einen schweren Fall mit Kindern oder sonst allerlei Unglück in diesem Leben über uns kommen läßt, auf daß er uns also schmelze und läutere. Wer kann aber dasselbe ertragen? Auf einen solchen Messias hatten die Juden nicht gehofft, sondern sie wollten einen haben, der sie sollte zu Königen und Herren der ganzen Welt machen. Nein, sagt Maleachi, damit werden sie fehlen; denn er wird scharf, und wie es scheinen wird, greulich sein. Und wie greulich diese Läuterung sein mag, spielet doch unser Heiland mit uns, nicht daß wir dadurch sollen verderbet, sondern gebessert werden.

58. Also ist nun Joseph gegen seine Brüder sehr freundlich und sanftmüthig. Denn er wendet sich etlichemal um und weint bitterlich; und zeigt mit solchem seinem Weinen an, daß er ganz geneigt sei, ihnen zu verzeihen. Und geht doch gleichwohl äußerlich so ungnädig mit ihnen um, tobt und zürnt mit ihnen, als wären sie Kundschafter, wirft Simeon in das Gefängniß, und schreckt sie, wie er nur immer kann und mag. Darum ist dies ein wunderbarer Exempel; denn er kann gar zu großer Liebe und Barmherzigkeit sich des Weinens nicht enthalten. Aber wenn du die Worte und äußerliche Geberde ansehen willst, so erzeiget er sich gegen sie viel stärker und härter, denn ein Kieselstein oder ein Diamant sein kann.

59. Eben auf solche Weise führt GOtt seine Heiligen auch wunderbarlich. Ja, die Schulmeister und Eltern züchtigen die Kinder auch mit Ruthen und Schlägen, auf daß etwas Gutes aus ihnen werden möge. Wie das deutsche Sprüchwort lautet: Je lieber Kind, je schärfer Ruthe. Und doch ist dies ein Spiel der göttlichen Majestät, auch frommer Eltern und Lehrmeister, welche ohne Zweifel ihre Kinder damit nicht wollen getödtet haben. Denn das läßt ihnen die herzliche Neigung, so sie zu ihnen tragen, wie dieselbe den Eltern von GOtt eingepflanzt ist, nicht zu; sondern es fordert es die väterliche Liebe und Barmherzigkeit also, nämlich, daß sie mit Schlägen und Sträufen bei den Kindern müssen anhalten, auf daß sie dadurch mögen fromm werden; wiewohl solches ohne Schmerz und ohne Zorn nicht geschehen kann. Sieh, sie lachen stellt, als ob er seinen Brüdern gar fremd und feind sei, und sträfst sie auch sehr hart, droht ihnen mit dem Tode und allerlei Strafen; und doch ist sein Herz ganz freundlich und lieblich, und denkt nichts weniger, denn daß er sie strafen, oder sich an ihnen rächen wollte. Ja, er geht damit um, daß er sie erhöhen und zu Ehren setzen, und auf allerlei Weise zieren möge.

60. Darum sollen wir bedenken, daß dies ein recht Zeichen, Figur oder Vorbild aus der heiligen Schrift sei, nämlich, wie GOtt die Gottseligen pflege zu führen und zu regieren, auf daß wir nicht bald kleinmüthig werden und verzagen, wenn wir fühlen, daß wir von GOtt auch gezüchtigt, oder wohl gar verstossen, und verräumt und verachtet werden, wie David klagt im 31. Psalm V. 23.: „Ich sprach in meinem Zagen: Ich bin von deinen Augen verstossen“ etc. Denn in Wahrheit spielt GOtt mit den Gottlosen nicht also; sondern dieselbigen „werden alt bei guten Tagen“, wie Hiob sagt Cap. 21, V. 13., „und erschrecken kaum einen Augenblick vor der Hölle.“ Er hat traun mit Sodom und Gomorra nicht also gespielt, da guter Friede, Hoffahrt und Alles vollauf war; Lot aber und Abraham wurden hart geplagt, die müssen sich leiden.

61. Also werden wir endlich GOtt erkennen und seinen guten gnädigen Willen verstehen lernen. Denn es hat mit der Strafe so GOtt über die Seinen kommen läßt, die Meinung gar nicht, wie die Papisten mit ihrer thörichten und närrischen Buße vorgeben, als ob sich GOtt mit unserm Genugthun wollte versöhnen lassen. Sie wissen nicht, was Sünde sei; wissen nicht, daß GOtt mit den Heiligen und Gläubigen also spielt, nicht, daß sie sollen verderben oder verloren werden, sondern daß an ihm erkannt werde, was er von sich selbst rühmt und sagt Jer. 3, 12.: „Ich bin barmherzig, und will nicht ewiglich zürnen“; und darnach das Andere auch, so er noch hinzu thut und spricht V. 13.: „Alleine erkenne deine Missethat, daß du wider den HErrn gesündiget hast“ etc. Desgleichen in den Klageliedern Jeremias Cap. 3, V. 31—33.: „Der HErr verstößt nicht ewiglich; sondern er betrübet wohl, und erbarmet sich wieder nach seiner großen Güte; denn er nicht von Herzen die Menschen plaget und betrübet.“ Und Hiob am 10. Cap. V. 8. 9.: „Deine Hände haben mich gearbeitet, und gemacht alles, was ich um und um bin; und versehest mich so gar. Gedenke doch, daß du mich aus Leimen gemacht hast, und wirst mich wieder zu Erden machen“ etc.; V. 12. 13.: „Leben und Wohlthat hast du an mir gethan, und dein Aufsehen bewahret meinen Odem. Und wiewohl du solches in deinem Herzen verborgen; so weiß ich doch, daß du des gedenkest“ etc. Als wollte er sagen: Das weißt du; wie stellst du dich denn dazu? Du wirst mich ja für deine Creatur halten, willst mich ja nicht verwerfen, als wäre ich nicht deine Creatur und du der Schöpfer; wiewohl du dich also stellst, als wüßtest du davon nichts, und erzeigest dich anders als ein Schöpfer.

\quellseite{379}

62. Diese Weisheit aber und dies Spiel Gottes geht die an, so da seufzen und schreien, und die ein zerbrochenes und zerschlagenes Herz haben, und die ein glimmendes Docht sind. Diese sind es, die es fühlen, denselben muß man diesen Trost vorhalten, auf daß sie nicht verzagen, sondern stark und männlich seien, wie der 27. Psalm V. 14. ermahnt: „Sei getrost und unverzagt, und harre des Herrn.“ Denn das ist ein gewiß Zeichen des Lebens, wo man Gottes Zorn fühlt. Der Papst, Bischöfe und Cardinäle fühlen ihn nicht, die Fürsten und alle Gottlosen, Ehebrecher, Hurer, Geizhälse etc. auch nicht; darum leben sie auch nicht. Denen aber, die da fühlen, daß Gott zornig ist, soll man sagen: Lieben Kinder, erkennet und lernet doch, daß Gott die Menschenkinder nicht von Herzen also demüthiget und verstößt. Davon hast du hier ein sehr lieblich Exempel. Joseph geht mit seinen Brüdern sehr hart um, wirft sie in das Gefängniß, und darnach, da sie Benjamin auch zu ihm bringen, erzeigt er sich noch viel greulicher. Du sollst aber wissen, daß es sein Ernst nicht ist, wie er sich äußerlich stellt: es ist seine Meinung gar nicht; also daß man wohl sagen möchte: Er redet oder thut dies nicht von Herzen. Ja, wer es glauben könnte!

63. Derhalben so lerne und gewöhne dich daran, daß du in der Anfechtung beharren könnest, alsdann wirst du erfahren, daß Gott in Wahrheit ganz freundlich ist und dich sehr lieb hat. Ja, sagst du, warum spielet er denn so gar unfreundlich? Liest er denn, wenn er sich also stellet? Mit nichten, wie man Joseph auch keiner Lügen zeihen kann. Augustin und Andere mehr disputiren davon. Aber frage du einen Vater, der sein Kind züchtigt, ob er’s lüge? Er schlägt seinen Sohn wohl, aber nicht der Meinung, daß er ihn damit wolle verderben und umbringen, sondern auf daß er ihn fromm mache; und lügt doch nicht. Er hat ihn recht herzlich lieb, wiewohl er ihn ganz hart straft. Denn Thorheit steckt dem Knaben im Herzen; aber die Ruthe der Zucht wird sie ferne von ihm treiben“, Spr. 22. 15. Also hängt uns die Erbsünde an, so wir im Paradiese bekommen, der Sauerteig des Teufels, damit unsere Natur vergiftet ist; das ist die Thorheit des Herzens und unseres eigenen Willens, nämlich, die Schlauigkeit, Sicherheit und Verachtung Gottes, daß wir uns träumen lassen, wir seien fromm und rein von allem Unflat. Da thut die Ruthe der Zucht vonnöthen, wie Salomo sagt, und daß der Knabe gezüchtigt werde, auf daß er seine Sünde und Irrthum erkenne, darum daß er seinen Vater nicht fürchtet, oder ihn auch nicht wiederum liebt, welcher doch den Sohn herzlich lieb hat; nicht, daß damit für die Sünde genuggethan werde, wie die Papisten davon unnütz waschen können, sondern daß der Knabe fromm und die Thorheit abgethan werde.

64. Denn es wird mit uns geboren eine erschreckliche Finsterniß und Unwissenheit, und eine gänzliche Abwendung von Gott. Da sagt denn der Herr: Ich will dir es aus Gnaden verzeihen ohne alle dein Verdienst, nicht um deiner Reue oder Genugthuung willen. Denn es ist um die Sünde nicht also gethan, daß sie von uns selbst könnte gebüßt, oder mit unsrer Reue und Genugthuung bezahlt werden; sondern der Sohn Gottes selbst hat sie büßen und bezahlen müssen.

65. Das will ich aber thun, spricht Gott: wenn ich dir die Sünde vergeben werde, will ich dich zum Erben und Kind Gottes machen in meinem Reiche, daß ich damit meine Liebe an dir beweise, doch also, daß ich dir erstlich deinen Unflat abwasche. Ich muß dich zuvor wischen und waschen; wie die Mutter das Kindlein nicht in die Wiege legt, sie habe es denn zuvor gewischt und gewaschen, und fragt diesweil nichts darnach, daß das Kind jämmerlich heult und weint, daß sie es derhalben nicht waschen und reinigen sollte.

66. Also sind wir zum Reich Gottes berufen worden, haben Vergebung der Sünden, sind Kinder und Erben Gottes: aber es hängt uns noch an, daß wir träge, faul und schläfrig sind an Leib und Seele. Diese böse Sucht gedenkt unser Herr Gott wegzunehmen und abzuthun in unserm ganzen Leben.

67. Aber diese Lehre ist vorzeiten im Papstthum bei den Sophisten und ihren Zuhörern oder Schülern gar unbekannt gewesen. Und da ich anfänglich etliche kurze Sätze im Druck ausgehen ließ und lehrte, wie das ganze Leben der christgläubigen Menschen eitel Buße wäre, sind darüber die Papisten schier toll und thöricht geworden. Denn sie haben gelehrt, daß die Erbsünde vergeben und gar weggenommen sei in der Taufe, und sei davon gar nichts mehr übrig an uns, denn nur ein kleiner Zunder, das ist, eine natürliche Schwachheit in den Menschen. Darum haben sie nichts gedacht und auch gar nichts gelehrt von der Sünde, wie dieselbe an uns müsse geläutert und ausgetilget werden; die wirklichen Sünden aber, als Diebstahl, Hurerei und dergleichen, sagten sie, sollte man büßen mit guten Werken, mit eigenem Genugthun, mit Wachen, Beten, Klostergelübden etc. Das ist ihre Lehre gewesen von der Buße, daß wir mit unserm Verdienst und Tugenden für die Sünden sollten genugthun; so doch eine viel größere Bezahlung und Genugthung dazu gehört, nämlich, der Sohn Gottes selbst. Aber dieweil sie nicht gewußt haben, was Sünde sei, haben sie dieselbe also gering und klein geachtet, und haben sich träumen lassen, es sei nur eine Schwachheit oder Gebrechen und Trägheit unserer Natur. Wie aber die Sünde müßte ausgetilget werden, darum haben sie sich wenig bekümmert. Also sind die Leute durch solche erbärmliche Finsterniß und Verfälschung der gesunden Lehre greulicher Weise verführt worden, nur wenige ausgenommen, die noch sind erhalten worden, welchen dies Licht von Gott gesandt worden, daß der Mensch allein durch den Glauben müßte gerecht werden.

68. Darum lehren wir dies fleißig, nämlich, wie der alte Sauerteig der Sünden an uns müsse ausgetilget werden durch unser ganzes Leben; wir lehren aber nichts von der Würdigkeit unserer eigenen Kräfte und Genugthuung, als sollte man dadurch die Sünde tilgen und bezahlen können, wie die Papisten fälschlich gelehrt, wenn zuvor die Erbsünde durch das Blut Christi und die Taufe weggenommen wäre. Und haben sie dahin gezogen den Spruch des Hieronymus, da er sagt, daß die gefallen sind, wie nach dem Schiffbruch, das andere Brett der Buße, das ist, der Reue und Genugthuung ergreifen und sich daran halten müßten.

69. Dies ist die Ungewißheit der Sünde und die Unreinigkeit der Mönchslehre im Papstthum, damit sie die Leute überredet haben, daß an ihnen keine Sünde mehr übrig wäre, ausgenommen ein kleiner Zunder; welches wahrlich ein sehr böser und scheußlicher Irrthum, ja, eine lautere Lüge ist. Denn wir sind zu Gnaden angenommen durch die Taufe, nicht allein zur Vergebung der Sünden, sondern daß auch dieselben an uns noch immer sollen ausgetilget werden. Die Vergebung kommt aus lauter Gnade, und widerfährt uns allein um des Sohnes Gottes willen, ohne alle unsere Würdigkeit, Verdienst und Reue.

70. Nach dieser Vergebung folgt Strafe, Angst, Verfolgung, Tödtung; welche alle dahin gehen, daß die Sünde abgethan werde, also daß sie nicht allein vergeben und verziehen sei durch Gottes Gnade und Barmherzigkeit, sondern daß sie auch noch immer durch die Gabe des Heiligen Geistes ausgetilget werde, und daß du lernest verstehen, wie groß die Bosheit der Menschen und wie die menschliche Natur so gar verderbt sei. Denn die große Finsterniß und Blindheit, so uns angeboren ist, läßt uns weder unsern eigenen Jammer, noch die unermessliche Gnade und Barmherzigkeit Gottes erkennen, und hat Leib und Seele ganz starr und gefühllos gemacht.

71. Damit nun diese Stumpfheit und Gefühllosigkeit weggenommen und abgethan werden möge, dazu sind heftige und peinliche Strafen vonnöthen, welche uns die Schlafseuche und Faulheit vertreiben können. Es gehört eine scharfe Geißelruth dazu, daß sie durch Mitwirkung des Heiligen Geistes in uns erwecken die Erkenntniß und Bewunderung der unendlichen großen Güte Gottes, und der großen Liebe, damit er uns geliebet. Denn Judas und Saul haben diese Läuterung nicht ertragen können, sondern sie verzweifeln und rufen überall, daß sie ewig von Gott verworfen seien, darum daß sie also geplagt, genäret und getödtet werden.

\quellseite{380}

72. Darum soll man den Betrübten, die also geplagt werden, dies fleißig einprägen, daß GOtt mit ihnen nicht zürne, und daß sie die Verlassung dahin deuten, daß sie nichts Anderes sei, denn das gewisse Zeichen, daß sie GOtt zu Gnaden angenommen habe; sintemal er einen jeglichen Sohn, den er aufnimmt, stäupt, Hebr. 12, 6., damit er in dem Stumpfsinn und in der Blindheit der Erbsünde nicht verderbe und verloren werde; welche Erbsünde die Leute aus eigener Vernunft weder sehen noch verstehen können.

73. Auf diese Weise thut nun Joseph Anderes nichts, denn daß er mit diesem harten Werk und Rede die Sünde seiner Brüder offenbare. Denn sie ziehen hinab nach Egypten mit andern Gefährten, Getreide zu kaufen, sind sicher, und achten der greulichen Sünde, deren sie sich doch schuldig gewußt, gar nicht, als hätten sie nie gesündigt wider ihren alten Vater oder unschuldigen Bruder. Sie denken, Joseph werde vorlängst gestorben sein, der Vater aber wisse von allen den Dingen gar nichts. Was geht es uns an? werden sie gesagt haben. Die groben Gesellen thun keine Buße. Diese harten Kiefersteine und Diamanten müssen zerbrochen und zerknirscht werden, und ihre Augen müssen geöffnet werden, auf daß sie die Grausamkeit ihrer Sünde sehen können. Und da Joseph dasselbe ausgerichtet, hat er sich bald mit Worten und Gebärden freundlicher erzeigt und hat ehrlich mit ihnen gehandelt.

74. Dies ist aber auch wohl eine Ursache gewesen, daß er ihre Herzen also geforscht und versucht hat, daß er sich besorget, sie möchten sich vielleicht noch mehr an ihrem Vater und jüngsten Bruder versündigt haben. Denn da er von ihnen nach Egypten ist verkauft worden, und gesehen hat, daß sie sich dasselbe noch nicht haben lassen leid sein, hat er Sorge, sie möchten solcher Sicherheit und Unbußfertigkeit halben in andere größere Sünde gefallen sein, und daß sie auch dem andern Sohne, so von Rahel geboren war, welchem die Erstgeburt gebührte, zugesetzt und Gewalt gethan hätten deßwegen, daß sie gern wollten Herren sein und das Regiment allein haben. Anderer Dinge halben hat er wenig Sorge getragen. Er fragt nur nach dem Vater und Bruder. Denn dieweil er weiß, daß sie durch seinen Tod nach der Erstgeburt gestanden und derselben begehret haben, hat er die Sorge getragen, sie möchten vielleicht aus Bosheit und Ehrgeiz bewogen worden sein, daß sie beide Vater und Bruder umgebracht hätten, und daß also das ganze Haus seines Vaters sehr verwirrt und zerrüttet wäre. Um der Ursache willen scheint er droben, 1 Mos. 41, 51., gesagt zu haben, er habe seines Vaters Hauses vergessen, und nicht einmal gedacht, wiederum in sein Vaterland zu kommen.

75. Derhalben hat diese Sorge und Grausamkeit der Sünde den Joseph bewogen, daß er seiner Brüder Herz so viel genauer erforscht und untersucht hat, also, daß er ihnen damit nicht allein die vorigen Sünden, so sie an ihm begangen hatten, sondern auch ihre bösen Gedanken hat zu Gemüth führen und sie derselben erinnern wollen. Und ist wahrlich diese Inquisition und Erforschung sehr unwillkommen und scharf genug gewesen, und ist doch aus dem gelassensten Herzen hergekommen. Ich zwar wäre noch etwas härter mit ihnen umgegangen.

76. Aber diese Bitterkeit oder Unfreundlichkeit, so er ihnen äußerlich bewiesen, geht dahin nicht, als hätte er sich an ihnen rächen wollen der Gewalt halben, so sie an ihm geübt; sondern dient nur dazu, daß sie dadurch möchten zur heilsamen Buße kommen und demüthigt werden, auf daß sie auch sagten, wie im 32. Psalm V. 5. steht: „Ich sprach: Ich will dem HErrn meine Uebertretung bekennen. Da vergabest du mir die Missethat meiner Sünde“ etc.

77. Darum haben wir gesagt, daß dies zu der Lehre von der Buße soll gezogen werden, welche darin steht, daß man herzliche Reue habe über die Sünde und dieselbige ernstlich erkenne, und dazu auch die göttlichen Verheißungen in Christo mit dem Glauben ergreife und annehme. Solcher Stücke ist keines in der Buße, so Judas gehabt, gewesen. Denn er hat sich die Sünde nicht recht herzlich leid sein lassen, ist mehr der Strafe als der Sünde feind gewesen: es ist ihm leid, daß er nun hat müssen verstoßen sein, und ergreift doch die Gnade und Barmherzigkeit GOttes in Christo nicht.

78. Darnach soll man alle Strafen, Plagen, Jammer und Unglück auch dahin ziehen und deuten. Denn damit wird dies zuwege gebracht, daß die Menschen also zu ihrer eigenen und auch zu GOttes Erkenntniß geführt werden. Denn die Strafen und Unglück, so wir leiden, streiten wider das Fleisch und die Sünde, welche nicht durch solche Genugthuung bezahlt, sondern ausgefegt werden soll, und erinnern uns, daß wir GOttes Güte und Gnade erkennen, welcher uns nicht verwirft oder verderbt, wenn er uns stäupt, sondern uns dadurch fromm machen will. Denn es ist kein Zorn, sondern nur eine Züchtigung; keine Enterbung, sondern eine Läuterung. Wenn du nun dies also in deinem Herzen fleißig bedenken und dir selbst einprägen wirst, alsdann wirst du leichtlich eine jegliche Züchtigung des HErrn dulden und ertragen; ja, du wirst noch wohl dazu wünschen und begehren, daß GOtt nur weiblich haue, brenne und den übrigen Sauerteig der Erbsünde ausfegte. Da werden die rechten Kinder daraus, die nicht abfallen, sondern in der Anfechtung und Trübsal beharren und beständig bleiben.

79. Dies soll man fleißig treiben und den Leuten immer vorhalten um der groben Esel, der Papisten, willen, so die große Würdigkeit ihrer Reue und Beichte zu rühmen pflegen, als könnten sie damit für die Sünde de congruo et condigno genugthun, das ist, so, daß GOtt ihnen nach Recht und nach Billigkeit gnädig und hold sein müßte. Aber dies ist eine eigentliche Weisheit der Christen, nämlich, daß sie wissen, sie sind in Sünden geboren, und daß die Sünde noch immer im Fleische stecke und demselben anhange bis in den Tod, und daß wir davon nicht vollkommen können errettet und gereinigt werden, denn allein durch Tod, Würmer und letztes Feuer. Inzwischen aber hat GOtt noch immer Geduld mit uns nach seiner unverdienten Gnade und Barmherzigkeit, verwirft oder verderbt uns nicht, sondern läutert und wäscht uns täglich mit seiner Hand.

80. Und dies ist die rechte Lehre in dieser ganzen Historie, welche wahrlich höher ist, denn ich verstehen oder begreifen kann, geschweige, daß ich sie nach ihrer Würde handeln und auslegen sollte. Joseph erzeigt sich gegen seine Brüder wunderlicher Weise hart und unfreundlich; und doch mit solcher Freundlichkeit, daß er auch weicher ist, denn Moses sein kann, wiewohl er sich dermaßen erzeigt, als ob er härter sei denn Eisen; und thut dasselbe nur darum, daß er ihnen damit an Leib und Seele helfen möge.

81. Denn man soll diese That nicht für ein Exempel annehmen, derselben also nachzufolgen, als ob man strafen und sich selbst rächen möchte an denen, so uns beleidigt und Unrecht gethan haben. Denn es ist hier keine böse Lust an Joseph, sich selbst zu rächen, sondern es ist das freundlichste Spiel mit seinen Brüdern, wie er hermach den silbernen Becher in Benjamin’s, des Jüngsten, Sack hat legen heißen. Ist aber das nicht das freundlichste und lieblichste Spiel? Aber die Stunde, darin es gespielt wird, scheint es nicht also, sondern da ist es der Tod selbst. Wenn aber die Anfechtung ein Ende hat, da sehen sie nichts denn die allergrößte Gunst und Liebe. Ei, sagen sie, wie hat es Joseph, unser Bruder, so freundlich mit uns gemeint! Eben dasselbe empfindet mir auch also, daß wir nach der Erlösung von GOtt dermaßen auch urtheilen und halten müssen, welcher uns darum läßt angefochten und geplagt werden, daß wir prüfen mögen, welches da sei der gute und wohlgefällige Wille GOttes.

82. Die Disputation von der Lüge ist an diesem Orte unnöthig; denn Etliche fragen, ob Joseph lüge, da er sich stellt, als kenne er seine Brüder nicht, oder als sei er ihnen gar fremd. Aber er lügt nicht, wie auch ein Vater nicht lügt, wenn er seinen Sohn züchtigt, wie ich droben gesagt. Denn man möchte gleich nach derselbigen Weise auch fragen, ob GOtt lüge, wenn er uns mit allerlei Unglück, Jammer und Noth straft, wie Hiob am 10. Cap. V. 12. 13. sagt: „Leben und Wohlthat hast du an mir gethan; wiewohl du solches in deinem Herzen verbirgest“; welches wahrlich auch heißt, GOtt Lügen strafen. Dasselbige möchte man auch wohl sagen von einem Vater, der da zürnt und seinen Sohn züchtigt.

83. Darauf antworte ich also: Wir sind Lügner. Aber strafen und züchtigen heißt nicht lügen, sondern anzeigen und offenbaren, wer und was für Leute wir sind. Denn wir verbergen unsere Sünde, wollen der Sünden, so wir begangen haben, nicht schuldig sein; ja, wir lassen uns bedünken, wir haben GOtt oder unsere Eltern niemals erzürnt. Wenn wir aber unsere Sünde bekenneten und sagten:

\quellseite{381}

Ach! lieber Vater, ich habe gesündigt, verzeihe mir es, so würden wir nicht gerichtet oder gestraft vom Herrn, 1 Cor. 11, 31. Dieweil ich aber solches verhehle und lüge, als hätte ich niemals gesündigt; ja, ich habe noch wohl Lust zur Sünde und daß der Vater nicht weiß von meinen Sünden; wie die Brüder Josephs deshalb auch mit einander frohlockten: darum macht uns Gott zu Schanden, und bringt solche Heuchelei an das Licht, daß sie offenbar werde.

84. Derhalben lügt Joseph nicht, sondern geht nur damit um, daß er seine Brüder von ihrer Sünde reinigen möge. Gott lügt auch nicht, sondern ich lüge selbst, dieweil ich sicher bin und schnarche, nicht anders, denn als wäre gar nichts Böses an mir. Er spielt aber mit uns und sagt: Dieweil du selbst an deiner \unsicher{Seelenleib} so guten Gefallen trägst, liebstest dich selbst, und lässest dir träumen, du seiest ganz rein von allen Sünden: so will ich dir es offenbaren und anzeigen, was und wer du vor mir seiest, und will dir die Larve der Sicherheit und Heuchelei ausziehen.

85. So ist nun das eine Lüge, wenn man die Schuld oder Sünde nicht will erkennen, und ist bei der Lüge auch noch diese Thorheit, daß man die Strafe, Züchtigung und Arznei nicht verstehen will. Es fährt aber der Arzt nichts desto weniger fort mit Reinigen, Brennen und Schneiden, wiewohl dasselbe nicht ohne Schmerz geschieht. Aber dieweil du den Rath Gottes nicht verstehst, und lässest noch über das dazu in dem, daß du deine Sünde nicht erkennst: darum ist da einer solchen Arznei vonnöthen, deine Thorheit damit an den Tag zu bringen und zu offenbaren.

\section[Dritter Theil]{Dritter Theil. Wie Jakobs Söhne Joseph nicht kennen, wie Joseph sehr scharf mit ihnen redet und wie sie sich hierbei verhalten}

\begin{quote}
\textbf{V. 8. Aber wiewohl er (Joseph) sie (seine Brüder) kannte, kannten sie ihn doch nicht.}
\end{quote}

86. Dies hat Mose angezeigt, als ein sonderlich Wunderwerk. Denn es waren schon zwei und zwanzig Jahre vergangen, daß die Brüder Josephs ihn nicht gesehen hatten, und er kennt doch einen jeglichen unter ihnen; keiner unter ihnen hat ihn gekannt. Es hat ihnen aber im Wege gelegen und sie gehindert nicht allein das Alter und Gestalt des Leibes, welche in so langer Zeit war geändert worden, sondern auch die königliche Majestät und Herrlichkeit an Joseph, damit sie ihn geziert gesehen haben.

87. Darnach ist das Uebel der Blindheit auch dazu gekommen, damit, wie wir droben gehört haben, die Sodomiter auch sind geschlagen gewesen. Darum denken sie nichts weniger, denn daß dieser Fürst ihr Bruder Joseph sein sollte, welchen sie vor Zeiten verkauft hatten. Ja, sie haben sich selbst in ihrem Herzen beredet, es müßte Joseph vorlängst gestorben sein.

88. Zum Dritten redet er mit ihnen durch einen Dolmetscher, daß sie nicht vielleicht die Stimme kennen. Darum ist es kein Wunder, daß sie ihn nicht gekannt haben. Denn die Melancholiker pflegen dermaßen oft mit ihren Gedanken spazieren zu gehen, oder haben etwa sonst ernste Gedanken, daß sie weder sehen noch hören, daß sie doch schon zuvor gesehen und gehört haben. Derhalben war es nicht vonnöthen, daß er sich so fremd stellte, da er noch unter dem großen herrlichen Schein und Gewalt verborgen.

89. Aber daß er sich also gestellt, geht dahin, daß er damit hat anzeigen wollen, wie er gesinnet wäre. Denn er hat ihnen das freundliche und brüderliche Herz dermaßen noch nicht offenbaren wollen, darum hat er sich mit Worten und Gebärden gegen seine Brüder härter erzeigt denn sonst gegen andere Fremdlinge. Denn er ist über das Getreide gesetzt, dasselbe zu verkaufen, und hat anstatt des Königes einen jeglichen Käufer insonderheit gefragt, wer er wäre, was sein Thun wäre, und von wannen er käme, aus dem Lande Canaan oder aber aus \unsicher{Sundmaa}. Und dies ist wahrlich eine sehr schöne Ordnung gewesen, deren in unsern Landen und Städten um des gemeinen Nutzens willen auch wohl vonnöthen wäre, nämlich, daß man nicht so leichtlich jedermann einnähme in die Städte, wo man sich nicht zuvor um ihn wohl erkundigt hätte. Denn daß solches verachtet und nicht gehalten wird, ist eine Ursache vieles Unglücks. In Egypten hat Jeder vor Joseph kommen und Rechenschaft von seiner Hantierung geben müssen. Und auf solche Weise hat er auch seine Brüder gefragt: „Woher kommt ihr?“

\begin{quote}
\textbf{V. 9—11. Und Joseph gedachte an die Träume, die ihm von ihnen geträumet hatten; und sprach zu ihnen: Ihr seid Kundschafter und seid kommen zu sehen, wo das Land offen ist. Sie antworteten ihm: Nein, mein Herr; deine Knechte sind kommen, Speise zu kaufen. Wir sind alle Eines Mannes Söhne, wir sind redlich; und deine Knechte sind Kundschafter gewesen.}
\end{quote}

90. Das hebräische Wort, meraglim, kommt her vom Worte regel, welches Fußknecht heißt. Derhalben, daß wir Kundschafter verdeutscht haben, lautet nach dem Hebräischen eigentlich so viel als ein Fußknecht, oder Einer, der zu Fuße geht. Aber die hebräische Sprache braucht dies Wort für Kundschafter und Verräther, welches wir doch aus der Ableitung oder Grammatik nicht haben abnehmen können, wo nicht die Umstände und Exempel solches lehreten. Darum können wir den Nachdruck und Verstand vieler Worte in derselben Sprache nicht wissen. Denn das Wortverständniß ist den Hebräern zum großen Theil abhanden gekommen. Ich wollte es haben verdolmetscht, Läufer: aber sie nennen die Kundschafter also, die umher laufen, die Festungen zu besehen und wo das Land offen ist.

91. Denn das ist der rechte Verstand des Wortes ervath, welches Mose oftmals gebraucht hat, als, 3 Mos. 18, 6. 7. und an andern Orten mehr, und heißt eigentlich, nackt oder bloß. Dieweil aber von einem Menschen nicht gesagt wird, daß er nackt oder bloß sei, wo nicht seine Scham aufgedeckt wird, und wo die Feigenblätter oder die Schürze abgelegt werden, darum braucht es Mose für die Scham. Also bedecken wir die Augen nicht, denn sie sind nicht bloß; wie wir die andern ehrlicheren Theile am Leibe auch nicht bedecken. Die Ehe-Sünde aber macht uns nackt, wenn die Glieder, so zur Geburt verordnet, entblößt werden. Derhalben nennt er es hier durch ein Gleichniß die Blöße des Landes, wo das Land offen steht, das ist, solche Orte, die nicht sonderlich befestigt und ungebaut sind. Also droht Gott in den Propheten, er wolle die Scham der Heiden entblößen, das ist so viel gesagt, er wolle ihr Land verwüsten und verderben, daß es häßlich, ungebaut, schändlich und gar wüste stehen solle.

92. Darnach ist auch hieraus zu sehen, was zu derselbigen Zeit für Fleiß auf die Kundschafter gewendet worden ist und wie gute Acht sie auf dieselben gehabt. Darum scheint es auch, daß Joseph nicht allein ein Regent gewesen sei zu der Zeit des Friedens, sondern auch ein Obrister im Kriege, der da Macht und Befehl gehabt, die Unterthanen zu schützen und zu nähren, und hat eine sehr schwere Last auf sich liegen gehabt. Derselben hat er auch große Weisheit, Geschicklichkeit und Gottes Gnade vonnöthen gehabt zu der Regierung des großen weiten Königreichs, das in so großem Flor gestanden.

93. Das gehört ferner auch zu jener Härtigkeit und Unfreundlichkeit, damit Joseph seine Brüder versucht, daß er sie beschuldigt, daß sie Kundschafter oder Verräther seien, als ob sie nicht nach Egypten gekommen wären, Getreide zu kaufen, sondern das Land auszukundschaften. Und er hätte zwar andere unbekannte Leute dermaßen recht und billig fragen und versuchen mögen; hier aber hat er gar keine Ursache zu solcher Erforsschung. Denn er wußte ja wohl, wer sie waren und von wem sie gekommen waren. Er wußte wohl, daß sie der theuren Zeit halben genöthigt worden, Speise zu suchen, und daß sie nichts Böses vorhatten. Wie, hat er denn daran gesündigt oder gelogen, da er seine Brüder einer solchen Sünde zeiht, deren sie nicht schuldig gewesen, und hat die Unschuldigen also geschmäht?

94. Wir haben droben auch geantwortet, daß man gleicher Weise fragen möchte: Ob Gott daran auch sündige, wenn er zu uns sagt: Du bist des ewigen Todes schuldig und bist verdammt? Daran lügt er wahrlich nicht, sondern ist die offenbarste Wahrheit: aber mein Herz ist lügenhaftig und heuchlerisch, welches meine Eitelkeit und Thorheit verhehlt und zudeckt.

95. Derhalben kann oder soll auch dies an Joseph nicht gestraft werden, daß er sich gegen seine Brüder so rauh und unfreundlich gestellt. Denn es ist ein sehr freundlich und lieblich Stellen, und er suchte damit Anderes nichts, denn seiner Brüder Wohlfahrt und Seligkeit an Leib und Seele: es ist ihnen keine schädliche, sondern eine heilsame und dienstliche Lüge.

96. Darnach ist die andere Ursache, daß er so hart und scharf mit ihnen geredet, welche Mose selbst hinzugesetzt, nämlich, dieweil er an die Träume gedachte, die ihm von ihnen geträumt hatten. Denn dieselben will er recht und vollkommenlich erfüllt haben: er will von Herzen von ihnen geehrt sein, und daß sie sich vor ihm neigen sollen nicht allein als vor einem Fürsten in Egyptenland, sondern auch als vor ihrem Bruder Joseph; wie sie ihm denn auch diese Ehre hernach erzeigen werden, nach Absterben ihres Vaters. Denn er wußte wohl, daß ihm die Träume von Gott eingegeben waren und daß es prophetische Träume wären, darum hat er begehrt, daß sie auch wahrhaftig und nicht heuchlerischer Weise haben sollen erfüllt werden.

\quellseite{382}

97. Zuletzt hat Joseph eine rechtmäßige Ursache seines Argwohns gehabt, wie dasselbe aus der Historie erhellt. Denn er hatte elf Brüder daheim gelassen und kommen jetzt nur zehen zu ihm. Er sieht, daß der jüngste Bruder Benjamin nicht da ist, welchen er am meisten zu sehen begehrt hat. Dies ist ein natürlicher Argwohn gewesen, ohne alle Sünde. Darum hat er die Wahrheit erfahren wollen, nämlich, ob auch Benjamin und Jakob noch am Leben wären, oder aber, ob sie beide durch dieser seiner Brüder Bosheit umgekommen wären. Da ist es rechter Ernst gewesen, damit er umgeht; und möchte einem jeglichen wohl eingefallen sein, es sei zu besorgen, daß diese so bösen Buben noch größeres Unheil angerichtet hätten. Denn es ist keine Sünde, auf böse schalkhafte Leute den bößesten Argwohn zu haben. Wie droben, 1 Mos. 20, 2. 11., Abraham Sarah geheißen, daß sie sagen sollte, sie wäre nicht sein Weib; welches doch eine sehr schändliche Lüge war: aber er gibt selbst Antwort darauf und entschuldigt diese Lüge, da er sagt: „Vielleicht ist keine Gottesfurcht an diesen Orten.“ Denn er hatte zu Babylon und Haran wohl erfahren, daß allenthalben in allen Menschen die Erbsünde regierete: warum sollte denn Joseph von Todtschlägern eine gute Meinung haben?

98. Nicht, daß man leichtlich und bald von Andern urtheilen und Böses reden solle: aber einen Argwohn haben auf die, so offenbar böse sind, das ist keine Sünde, sondern ist höchst billig und recht. Und der Evangelist Johannes sagt von Christo selbst also, Joh. 2, 23. ff.: „Viele glaubten an seinen Namen, da sie die Zeichen sahen, die er that. Aber Jesus vertraute sich ihnen nicht; denn er kannte sie alle; und wußte wohl, was im Menschen war.“

99. Derhalben sage ich noch, daß es gut und recht sei, daß man von unbußfertigen verstockten Menschen und die offenbar böse sind, nichts Gutes halte und sich zu ihnen keines Guten versehe. Wo ich ihnen aber glaube und darüber Schaden nehme, so ist es meine eigene Schuld, daß ich in Gefahr komme, dieweil ich weiß, daß man den gottlosen bösen Menschen eben so wenig glauben soll, als dem Teufel selbst, ob sie wohl etwa solche Dinge versprechen und vorgeben, die wahrhaftig und gut scheinen. Denn also straft Christus im Evangelium, Marc. 1, 24. 25., die Teufel, so ihn bekennen, daß er Gottes Sohn sei, und gebietet ihnen, daß sie stillschweigen sollen. Eben also soll man auch mit denen umgehen, die offenbar böse sind, wo sie nicht Buße thun, und sollen uns vielmehr Böses denn Gutes zu ihnen versehen, auf daß sie uns mit keinerlei Schein, unter dem Vorwand der Gottseligkeit und Ehrbarkeit, betrügen können.

100. Derhalben handelt Joseph mit seinen Brüdern sehr weislich, da er sie wegen ihrer Bosheit in Verdacht hat, und erforscht, wie sie sich inzwischen gegen den alten Vater und Benjamin, seinen Bruder, gehalten haben. Und in solcher Erforschung, welches in dieser Historie das Vornehmste ist, bedenkt er dieweil, wie ihnen möge zu ihrer Besserung und Seligkeit geholfen werden. Denn darum hat er sich ihnen desto langsamer zu erkennen gegeben, auf daß er Ursache und Raum haben möge, ihnen also ihre Sünden zu zeigen und dieselben anzuzeigen, damit sie doch nicht so gar sicher seien, und daß für so gering Ding achten, den Bruder umbringen und tödten: sondern es hat diese Sünde groß und schwer müssen gemacht werden, auf daß sie selbst auch erkennen möchten, wie groß dieselbige wäre.

101. Darum ist es nicht vergeblich geschehen, daß sich Joseph gegen die Gäste und Fremdlinge, ja, gegen seine Brüder so hart und greulich erzeiglich erzeigt hat. Denn er hat mit den bößesten Leuten zu thun, welche Vater-, Bruder- und Muttermörder waren. Und ist unmöglich, daß Augustin und Lyra hier sich so viel mit der Lüge zu schaffen machen. Denn dies ist billig und recht gewesen, daß sich Joseph also gestellt hat, damit er sie ihrer Sünde möchte erinnern und sie also vom ewigen Tode und Verdammniß erretten.

102. Sie antworten ihm aber fein bescheiden und demüthiglich und sagen: Wir sind keine Kundschafter, sondern deine Knechte sind gekommen, Speise zu kaufen. Da hat Mose abermals das hebräische Wort, scheber, gebraucht, welches so viel heißt als, brechen. Es müssen aber mehr Worte zwischen ihnen gewechselt worden sein, welche hier nicht beschrieben werden. Denn es scheint, daß sie nicht vergeblich erzählt haben, daß sie alle Eines Mannes Söhne seien. Vielleicht hat Joseph daraus auch Ursache seines Argwohns genommen, daß ihrer so viel zugleich gekommen sind? Was ist das, wird er gesagt haben, daß ihr so zahlreich auf mich hin einströmt? Es scheint, als habet ihr euch zusammen verbunden und geschworen, mir oder dem ganzen Königreiche etwas ein Unglück zuzufügen. Es wäre genug gewesen, daß Einer oder Zwei besonders hereingekommen wären. Ich weiß nicht, was dieser Haufe bedeuten mag, oder was er vorhat. Ihr kommt daher unter andern Käufern, aus vielen Dieben und den ärgsten Buben auf einen Haufen versammelt: ihr werdet etwas im Sinne haben etc. Dies hat er also erdichtet, sie damit zu schrecken, und spielt doch mit ihnen aus freundlichem und brüderlichem Herzen. Sie haben wiederum gesagt: Wir sind zwar kein Haufen Mörder oder böser Buben, die mit Frevel daher gekommen wären, andern Leuten Schaden zu thun. „Wir sind alle Eines Mannes Söhne, wir sind redlich.“ Das hebräische Wort heißt kenim, das ist, vollkommen, fromm und aufrichtig.

103. Es ist aber dies Wort in der ganzen Schrift sehr gewöhnlich, und wir haben desselben gedacht in unserer deutschen Vorrede über den Psalter, da wir die Ursache unserer Dolmetschung angezeigt haben. Es steht im 5. Psalm V. 10.: „In ihrem Munde ist nichts Gewisses“; und im 51. Psalm V. 12.: „Gib mir einen neuen gewissen Geist“; item 1 Kön. 2, 46.: „Und das Königreich ward befestiget durch Salomos Hand“, das ist, beständig, gewiß und rechtschaffen, das nicht eitel leer oder hohl ist wie ein Rohr. Wir haben es an diesem Orte verdeutscht: fromm, redlich. Wir sind, sagen sie, kenim, das ist, rechtschaffen, fromm, aufrichtig und redlich, wir reden nicht anders, denn wir es im Herzen haben, werden nicht hin und her getrieben wie ein unbeständiges und loses Rohr, sondern sind in Wahrheit also gesinnet, wie wir uns mit Worten und äußerlichen Geberden vernehmen lassen. Es ist hie kein falscher Schein und keine Heuchelei darunter verborgen, sondern es ist in uns ein rechtschaffener, gewisser und beständiger Geist, da man sich gewißlich darauf verlassen darf. Es steckt nicht hohl Ding in uns; du darfst dich nicht fürchten, sondern magst dich auf uns verlassen. Dies ist eine sehr feine Antwort und Entschuldigung.

\begin{quote}
\textbf{V. 12. 13. Er sprach zu ihnen: Nein, sondern ihr seid kommen zu besehen, wo das Land offen ist. Sie antworteten ihm: Wir deine Knechte sind zwölf Brüder, Eines Mannes Söhne im Lande Canaan, und der jüngste ist noch bei unserm Vater; aber der eine ist nicht mehr vorhanden.}
\end{quote}

104. Joseph dringt weiter auf sie ein, und macht die Sünde, damit er sie beschuldigt, immer härter, dieweil sie ihre Sünde noch nicht bekennen. Ihr seid, sagt er, nicht darum gekommen, daß ihr Speise kaufen wollt, sondern vielmehr „zu besehen, wo das Land offen ist“; dazu habt ihr diesen Haufen Diebe und böser Buben an euch gehängt. Dies sind aber auch nur halbe Worte und ist hier etwas ausgelassen. Denn er wird sie ohne Zweifel noch etwas mehr gefragt haben. Lieber, wird er gesagt haben, wie kommt es, daß ihr eben in einer solchen Anzahl zu mir kommt? Warum bringt ihr nicht mehr oder weniger Leute mit euch? Da geben sie diese Antwort: Es sind unser zehen, aber in Wahrheit sind wir zwölf Brüder, Eines Mannes Söhne, und sind keine Landstreicher, sondern unser Vater wohnt gewöhnlich im Lande Canaan. Dies hat er also durch etliche Fragen heraus gebracht, und er hat Ursache gehabt, nach dieser Zahl zu forschen; denn Benjamin ist damals drei und zwanzig Jahre alt, und Vater von zehn Kindern gewesen, daß das Alter nicht im Wege war, daß er nicht auch mit den andern Brüdern nach Egypten gezogen wäre. Darum sagt nun Joseph: Ihr seid lose betrügliche Leute, die anders reden und doch anders denken. Ich höre, daß eurer Brüder sollen zwölf sein; warum sind aber ihrer so viel nicht zugegen?

\quellseite{383}

105. Darauf antworten sie also: Siehe, „der Jüngste ist noch bei unserm Vater.“ Also erklären sie sich, daß sie in dieser Sache rechtschaffen, redlich und aufrichtig, und nicht lose oder Heuchler seien. Es ist aber noch nicht Alles heraus, er muß noch das herausdreschen. Darum setzet sie nun auch hinzu: „Der eine ist nicht mehr vorhanden.“ Da sehen wir, daß sie an Josephs Leben und daß er noch sollte gesund sein, gar verzagt haben. Denn sie denken nicht, daß er noch am Leben sein sollte; darum kann ihnen das nicht einfallen, daß dieser Fürst ihr Bruder Joseph sein sollte, da ihnen vornehmlich so viel Umstände im Wege gelegen, als, der große herrliche Schein seiner Majestät, und daß ihre Herzen so gar mit den Gedanken bestimmet gewesen, daß Joseph müßte gewißlich todt sein; welches allein wohl hat eine Ursache sein können, daß sie gar unbedachtsam gewesen.

\begin{quote}
\textbf{V. 14–17. Joseph sprach zu ihnen: Das ist's, das ich euch gesagt habe; Kundschafter seid ihr. Darum will ich euch prüfen; bei dem Leben Pharao, ihr sollt nicht von dannen kommen, es komme denn her euer jüngster Bruder. Sendet einen unter euch hin, der euren Bruder hole; ihr aber sollt gefangen sein. Also will ich prüfen eure Rede, ob ihr mit Wahrheit umgehet oder nicht. Denn wo nicht; so seid ihr bei dem Leben Pharao Kundschafter. Und ließ sie beisammen verwahren drei Tage lang.}
\end{quote}

106. Joseph bleibt noch immer bei seiner Härte. „Das ists“, sagt er, „das ich euch gesagt habe: Kundschafter seid ihr.“ Darauf bleibe ich, Verräther seid ihr. Jetzt habt ihr gesagt, ihr seid fromm und redlich, und seid zwölf Brüder, Eines Mannes Söhne, und sind doch nur zehen hieher gekommen: ihr seyd wahrlich mit Lügen um. Er hört nicht auf mit ihnen zu spielen, bis er das Bekenntniß der Sünden herausbringt.

107. Was sollen sie aber nun hier thun, dieweil sie nicht gefährdet werden und er ihren Worten keinen Glauben schenkt? Sie sagen die Wahrheit, aber Joseph glaubt ihnen doch nicht; denn sie hatten zuvor ihrem Vater selbst auch Lügen gesagt und ihn betrogen, da sie ihm verkündigt, ein wildes Thier hätte ihren Bruder Joseph zerrissen, 1 Mos. 37, 32. 33.

108. Ja, also muß der grobe Unverstand unsrer Natur und die Erbsünde herausgebrochen werden, welche nicht versteht, daß sie unrein und verdammt sei, sondern vielmehr denkt, daß sie in ihren Werken allenthalben gerecht und rein sei. Darum muß sie an das Licht gezogen werden, daß solches offenbar und bekannt werde. Derhalben hält Joseph immer an, da er sagt: „Bei dem Leben Pharao, ihr sollt nicht von dannen kommen, es komme denn her euer jüngster Bruder.“ Nach dem, der todt sein soll, frage ich nicht; aber dieweil der Kleinste noch daheim ist und ihr ihn nicht mitgebracht habt, verrathet ihr euch selbst, gleich einer Spinnmaus, mit eurer Stimme. Denn warum habt ihr denselben allein daheim gelassen? Das muß nicht wahr sein.

109. Hier widerspricht aber Joseph der erkannten Wahrheit, denn er weiß ja, daß sie die Wahrheit reden, allein, daß er den Argwohn hat und sich noch immer besorget, sie möchten etwa den jüngsten Bruder auch getödtet haben. Dieweil ihr euch nun rühmt, daß ihr aufrichtig, fromm und redlich seid, woblan, damit ich meinem Amt und Pflicht genugthun möge, so will ich es versuchen, ob ihr Kundschafter seid oder nicht. Wenn ich euren Bruder sehen werde daher kommen, alsdann will ich euch Glauben schenken: wo er aber nicht erscheinen wird, so wird man es gewiß wissen können, daß ihr müßt lose Leute und Lügner sein. Also hat er sie in ihren eigenen Worten und Rathschlägen gefangen und bestrickt.

110. Und daraus kann man ein Gleichniß und Exempel göttlicher Werke und Regierung nehmen. Denn dies war mit gutem, wohlbedachtem Rath und Willen des Vaters geschehen, daß Benjamin mußte daheim bleiben, auf daß ihm kein Unfall begegnete. Und Jakob wird das hernach selbst also wiederholen, da er V. 38. zu ihnen sagt: „Mein Sohn soll nicht mit euch hinab ziehen, denn sein Bruder ist todt, und er ist allein überblieben; wenn ihm ein Unfall auf dem Wege begegnete, darauf ihr reiset, würdet ihr meine grauen Haare mit Herzeleid in die Grube bringen.“ Aber durch diesen Rath, der so klug und fein richtig war, wie es schien, kommen diese Brüder in die größte Gefahr und Beschwerung, welche jetzt nur das wünschen und begehren, daß doch ihr Bruder auch möchte da sein, und wenn derselbige da gegenwärtig wäre, so würden alle Dinge leicht und ohne alle Gefahr mit ihnen sein. Aber durch die Vorsicht des Vaters sind sie in Gefahr Leibes und Lebens geführet worden.

111. Also soll es gehen, daß wir anlaufen, wenn wir am allerklügsten sind. Und auf solche Weise pflegen oftmals die Dinge, so wir am allerbesten berathschlagt, am allerübelsten zu gerathen. Wie Jakob aus großer Weisheit und höchst rechtmäßiger Ursache den jüngsten Sohn daheim behalten, denkt, er werde also gemächlich sicher, frisch und gesund bleiben. Aber Gott hat ihm aus diesem seinem eigenen Rath seines Sohnes Verderben und schweren Unfall zugerichtet.

112. Darum gilt kein Rath wider den HErrn, wenn er uns strafen will; es hilft kein Rath, keine Vernunft, keine Weisheit wider Gott, wenn er uns plagen will. Es war dieser Rath Jakobs wahrlich sehr gut und weislich, aber es ist doch kein Glück dabei: es hilft nicht. Denn der HErr pflegt auch die allerweisesten Rathschläge zu zerbrechen, und dagegen solche aufzurichten und zu bauen, die sehr schwach und verächtlich sind, und an welchen jedermann verzagen muß und die am allerverächtlichsten sind, und will von uns haben, daß er im 50. Psalm V. 15. und andern Oertern mehr geboten hat, da er sagt: „Rufe mich an in der Noth; so will ich dich erretten“; item im 37. Psalm V. 5.: „Befiehl dem HErrn deine Wege, und hoffe auf ihn; er wirds wohl machen“; desgleichen im 55. Psalm V. 23.: „Wirf dein Anliegen auf den HErrn. Der wird dich versorgen.“ Denn wenn Alles verloren und verdorben ist, und, wie der 107. Psalm V. 27. sagt, wenn Menschen gar keinen Rath wissen, wenn alles verderbet ist, daß man schreit: Wir müssen verderben und zu Grunde gehen: so ist Gott da und sagt: Ich will nicht haben, daß schläge zerbrechen und zunichte machen, und dagegen andere, die sehr verächtlich sind, erwecken und damit den Leuten helfen.“

113. Also auch in diesem gegenwärtigen Exempel ist Alles so gar übel gestanden, daß sie weder Rath noch Hülfe gewußt, meinten, sie wären ganz verdammt, und würden jetzt alsbald sterben und verderben: aber bald wird im 45. Cap. V. 3. folgen: „Ich bin Joseph.“ Und da sie das einige Wort hören, werden sie wiederum lebendig, und werden die höchst armseligen und verlornen Rathschläge aufgerichtet, dadurch den Sachen gerathen wird. Zuvor aber waren sie sicher und fürchteten sich vor keiner Gefahr. Aber Gott findet Schuld an ihnen und Ursache wider sie, auch da sie ihr Ding am allerbesten berathschlagt und wohl bestellt hatten.

114. Derhalben soll man Gott vertrauen, wenn wir auch schon mit dem allerbesten und klügsten Rath gerüstet sind, und nicht verzagen, wenn es uns an Weisheit und menschlicher Hülfe mangelt; sondern sollen also schließen und sagen: „Ich glaube an Gott den Vater, den Allmächtigen, der die allerschönsten Rathschläge zunichte machen und in die Hölle verändern, und wiederum aus solchem Stande, der am allerverächtlichsten ist, den Himmel machen kann. Denn er ist ein allmächtiger Schöpfer. Ich kann ihm nicht zu tief fallen, er kann mich herausheben; ich kann ihm nicht zu hoch sitzen, er kann mich stürzen.“

115. Darum sollen wir nicht vermessen oder stolz sein unserer Weisheit halben, auch um unsrer Sünde willen nicht verzagen. Denn er ist über, und außer, und unter allen Dingen. Vor ihm ist Alles gar Nichts, und Nichts ist Alles. Denn er macht die Todten lebendig, und rufet dem, das nicht ist, daß es sei“, Röm. 4, 17.

116. Dasselbe haben wahrlich Jakob und seine Söhne auch also erfahren. Denn sie meinten nicht, daß sie anlaufen und etwa in Gefahr kommen könnten, sondern da sie jetzt meinten, sie haben Joseph auf seine Fragen gar wohl geantwortet, damit er werde zufrieden sein, da kehrt er ihre Antwort um und schließt daraus, daß sie lose Leute und Lügner sein müssen. Sehr beschnitten, sagt er, daß eurer Brüder zwölf und Eines Mannes Söhne seien: wohlan, wo bleibt denn der, von welchem ihr sagt, daß er noch am Leben sei? Aus eurem Munde richte ich euch, daß ihr leichtfertige Leute und Lügner sein müßt. Dasselbe haben sie nicht zuvorsehen oder behüten können; ja, es ist ihnen solches niemals in den Sinn gekommen, nämlich, daß der Fürst, so über das Getreide gesetzt war, nach ihrem Bruder Benjamin, welchen sie daheim gelassen hatten, fragen würde.

\quellseite{384}

117. Auf diese Weise spielet GOtt mit uns auch, daß wir erkennen und glauben sollen, daß er ein allmächtiger Schöpfer und dazu auch barmherzig sei, der da mächtig ist, die, so gedrückt oder beschwert sind, zu erheben, und die Hoffärtigen zu stürzen, wie gut und gerecht die Sache sein mag. Darum sollen wir uns demüthigen vor dem HErrn und sagen: Du bist allein der Höchste; darum bist du auch höher, denn meine Gerechtigkeit, Weisheit und Rath ist, ja, denn die ganze Welt sein kann.

118. Ich habe auch in meinem Leben nicht wenig Gefahr gelitten bei der Lehre, die ich bekenne und Andern vortrage: ich habe aber etlichemale erfahren, daß solche Gefahr einen sehr glückseligen Ausgang genommen hat, und daß es damit wohl gerathen ist, da doch die Sachen, wie es mir schien, dermaßen gestanden, als wäre es gar verloren. Wiederum aber ist es oftmals gekommen, daß die Sachen, so zuvor am allerbesten bedacht und berathschlagt gewesen, hernach am übelsten gerathen sind.

119. Darum sollen wir uns GOtt unterwerfen, auch da wir uns dünken lassen, daß wir mit Gerechtigkeit und Weisheit wohl gerüstet und bewaffnet sind; denn er ist über Alles. Und wir sollen Fleiß anwenden, daß wir ja die rechte Mittelstraße gehen, auf daß wir uns nicht versündigen, weder mit Vermessenheit, noch mit Verzweiflung.

120. Also lieben sich die Söhne Jakobs selbst sehr wohlgefallen, daß sie ihres Vaters und Bruders gedacht hatten, und hofften nicht anders, denn Joseph würde damit wohl zufrieden sein, daß sie ihm solches erzählt hätten: aber er hat es umgekehrt, und ist das Umkehren so hart gewesen bis an den Tod. Denn er läßt die Worte fahren und greift zu Schlägen, geht mit ihnen stracks zum Gefängniß zu, da er ihnen aufrückt, wie sie sich mit einander verbunden und dies alles erdichtet und gelogen hätten nur, daß sie damit der Sache einen Schein geben möchten, und den Betrug, damit sie umgangen, also bedeckten. Sie meinen, sie haben es wohl ausgerichtet: da kommen sie erst in ein recht Bad, darin sie wohl schwitzen. Bei dem Leben Pharaos, sagt er, entweder bringt euern Bruder, der noch vorhanden ist, her zu mir, oder ihr sollt alle sterben. Also wird ihre Sünde allgemein an das Licht gebracht, und sie kommen zu ihrer selbst Erkenntniß.

121. Augustin fragt an diesem Orte: Ob es sich auch geziemet habe, daß Joseph bei dem Leben Pharaos habe schwören mögen, vornehmlich dieweil er lügt und diesen Eid nicht gehalten hat. Denn er hat ja die neun Brüder ziehen lassen und nur Einen behalten. Aber es thut dieser Frage nicht noth. Denn erstlich mag man wohl schwören, nicht allein bei GOtt, sondern auch bei dem Himmel, bei der Seele, bei dem Haupte; wie diese Formeln erzählt werden von Christo, Matth. 5, 34—36. Und dieweil er solche Eide verbietet, zeigt er damit an, daß sie recht seien, wir schwören gleich bei unserm Haupte oder bei unsrer Seelen Seligkeit. Also sagen die Hebräer: Vivit Dominus. So wahr GOtt lebt; so wahr deine Seele lebt; bei dem Leben des Königs. Solche Eide, sage ich, sind an sich selbst nicht falsch oder unrecht; sondern Christus straft die Leichtfertigkeit bei dem Schwören, und die Eitelkeit derer, die da schwören. Also sagt Paulus 1 Cor. 15, 31.: „Bei unserm Ruhme, den ich habe in Christo Jesu, unserm HErrn, ich sterbe täglich“ \&c. Da schwört St. Paulus bei der Kirche oder Gemeine; wie es eine Gewohnheit gewesen, zu schwören bei Jerusalem, der Stadt eines großen Königes; und im 63. Psalm V. 12. steht: „Wer bei ihm schwöret, wird gerühmet werden.“

122. Derhalben mag man wohl schwören, sonderlich vor der Obrigkeit und in Gerichtshändeln, der Wahrheit Zeugniß zu geben. Denn Schwören ist nichts Anderes, denn daß man GOtt anruft; wie Paulus sagt, 2 Cor. 1, 23.: „Ich rufe GOtt an zum Zeugen auf meine Seele“ \&c.

123. Joseph aber kündigt damit nicht, daß er bei dem Leben Pharaos schwört, welcher größer war denn er; wie in der Epistel an die Hebräer 6, 16. gesagt wird: „Die Menschen schwören wohl bei einem Größern, denn sie sind; und der Eid macht ein Ende alles Haders, dabei es fest bleibt unter ihnen.“

124. Warum hält er aber nicht, was er geschworen hat? Antwort: Es ist genug, daß er seiner Brüder Einen bei sich behält. Denn da er denselben behalten, ist es eben so viel, als wären sie alle da geblieben. Laßt die Andern hinziehen, hat er gesagt, und ihren Bruder herbringen; dieser aber soll hier bleiben zum Zeichen des Eides, den ich gethan habe, bis ich erfahre, ob Benjamin noch am Leben sei. Wo er aber umgekommen ist, so ist es gewiß, daß ihr Kundschafter seid; da will ich euch derselbigen Sünde halben auch strafen.

125. Da werden sie ohne Zweifel ihn mit Weinen und Heulen um Gnade gebeten haben, daß er ihnen Barmherzigkeit beweisen wolle. Ach! werden sie gesagt haben, gnädiger Fürst, du wollest ja nicht alle binden lassen! Laß doch geschehen, daß zum wenigsten unsrer etliche diese Sache an unsern Vater bringen mögen! Aber er läßt sich noch so bald nicht erweichen, läßt sie zuvor beisammen verwahren drei Tage lang, welches hart genug ist. Aber er thut es nicht darum, daß er sich an ihnen rächen wollte wegen der Gewalt, so sie ihm gethan hatten; sondern daß er fleißig erforschen will, wie es seinem Vater und leiblichen Bruder gehen möge, und thut daselbste mit großer Aufrichtigkeit und Treue; darnach, daß er sie ihrer begangenen Sünde erinnern will, und sie zu ihrer selbst Erkenntniß bringen und ihre Sünde auslegen möge.

126. Wie, meinen wir nun, daß ihnen die drei Tage über möge zu Muthe gewesen sein? Denn das ist ja ein großer Jammer und schwere Anfechtung, daß sie sehen, wie sie von ihrem Vater, von Weib und Kind, und ganzem Hausgesinde gerissen seien. Und ist nun eben dieselbige Strafe auf ihnen gelegen, damit sie zuvor ihren Bruder geplagt hatten; ja, sie sind noch etwas übler daran, dieweil sie Eheleute gewesen und des Vaters halben bekümmert sind. Denn sie wissen, daß ihm diese Botschaft den größten Schmerz bringen wird, wenn er erfahren wird, daß seine Söhne bei den Egyptern im Gefängniß gehalten werden. Das ist es, das die Schrift 1 Sam. 2, 6. sagt: „Der HErr tödtet, und machet lebendig, führet in die Hölle, und wieder heraus.“ Und geschiehet doch dies alles aus freundlicher und höchster Liebe. Aber dies Spiel ist ihnen der Tod und Hölle gewesen, und sind ohne Zweifel die ganzen drei Tage über auf das höchste betrübt und bekümmert gewesen, haben Sorge getragen für ihr Weib und Kind, und für den alten Vater. Das ist die rechte Marterwoche und Charfreitag in höchster Traurigkeit und Zagen gewesen.

127. Es ist zwar noch ein klein Küstlein guter Hoffnung übrig, nämlich, daß er heißt den elften Bruder herbringen. Wenn sie denselbigen gebracht, haben sie gute Hoffnung, sie würden unversehrt davon kommen; und wissen doch, daß ihn der Vater nicht leichtlich wird wollen ziehen lassen, es sei denn, daß Einer unter ihnen gen Hebron gesandt werde. Da werden sie weiter mit einander disputirt haben, welcher unter ihnen am allerbesten geschickt sein würde, diese Gesandschaft recht auszurichten, und der so viel Gunst bei dem Vater habe, daß er solches von ihm verlangen könne. Und werden also der Dinge viel mit untergelaufen sein, welche Mose verschwiegen hat, nämlich, wie sie geweint, geseufzt und flehentlich gebeten haben, erstlich den Kerkermeister und darnach auch Andere mehr, daß sie doch wollten für sie bitten bei dem Fürsten, sein Herz erweichen, und ihm anzeigen, wie ihr Vater seinen allerliebsten Sohn nimmermehr würde von sich lassen und der Gefahr willen, darein, wie er hören würde, die andern Neun gerathen wären: ja, er würde wohl argwohnen, sie müßten schon alle umgekommen sein und daß jetzt dem Jüngsten auch heimlich nachgestellt würde.

128. Also wird das Fleisch recht getödtet und gereinigt von der Sünde. Denn sie sind dazumal recht in der Hölle und im Fegefeuer gewesen. Es ist ferner solch ihr Klagen Joseph angezeigt worden, welcher auch das Urtheil, so er über sie gesprochen, bald etwas gelindert; denn er hält dafür, daß es sich gebühren wolle, daß er die Gesundheit und Wohlfahrt seines Vaters und des andern Hausgesindes auch bedenken müsse, und gebietet derhalben, daß nur Einer gefangen gehalten werde, bis er gewißlich erfahren werde, wie es um den jüngsten Bruder stehe und ob sie redliche und aufrichtige Leute seien.

\section[Vierter Theil]{Vierter Theil. Josephs Verlangen an seine Brüder; was solches gewirkt; und wie Joseph endlich zu Mitleid bewegt worden}

\begin{quote}
\textbf{V. 18—20. Am dritten Tage aber sprach er zu ihnen: Wollt ihr leben, so thut also; denn ich fürchte GOtt. Seid ihr redlich, so lasset eurer Brüder einen gebunden liegen in eurem Gefängniß; ihr aber ziehet hin, und bringet heim, was ihr gekauft habt für den Hunger. Und bringet euren jüngsten Bruder zu mir; so will ich euren Worten glauben, daß ihr nicht sterben müsset. Und sie thaten also.}
\end{quote}

\quellseite{385}

129. In hebräischen Texte kommt hier abermals die gewöhnliche Art und Weise zu reden vor, nämlich, von dem Brechen der Theurung. Und die Juden sagen, daß diese Weise zu reden ihrer Ableitung nach daher komme, daß es den Hunger, die theure Zeit, breche; geben also damit ihre Unwissenheit an den Tag, und sind Leute, die nicht allein der Lehre der heiligen Schrift, sondern auch des Verstandes, so die Grammatik gibt, unerfahren sind. Darum heißt nun „Brod brechen“ Anderes nichts denn Speise, und soll nicht activisch, sondern passivisch verstanden werden, das ist, nicht, daß da bricht, sondern das gebrochen wird. Wie von den Jüngern Christi gesagt wird Luc. 24, 35., daß sie den HErrn erkannt haben an dem, da er das Brod brach, das ist, da das Brod gebrochen ward. Desgleichen 1 Joh. 6, 12. heißt Christus die übrigen Broden aufheben. Also wird an diesem Orte das Brodbrechen gesetzt und verstanden für die theure Zeit, wie er droben etlichemal die reichen Jahre passivisch genannt hat fractio saturitatis et fertilitatis, das ist, ein Brechen der Sättigung und Fruchtbarkeit.

130. Daß man nun darauf also merkte und Acht gebe, ist sehr nöthig in der heiligen Schrift, nämlich, wann man die Worte activisch, und wann man sie passivisch verstehen soll, wie es in der Grammatik genannt wird. Denn es wird die ganze Rede geändert je nach der activen oder passiven Bedeutung. Also, Röm. 3, 23.: „Sie sind allzumal Sünder, und mangeln des Ruhms, den sie an GOtt haben sollten.“ Wenn das activisch gesetzt wird, so hat das Wort eine andere Bedeutung, und hat auch eine andere, wenn man es soll passivisch verstehen. Darum ist an beiderlei Bedeutung viel gelegen, daß man den rechten Verstand daraus nehmen könne. Also mögen diese Reden, wenn man sagt: „die Gerechtigkeit GOttes“, item, „das Werk GOttes“, „die Kraft GOttes“ auf beiderlei Weise gedeutet werden. Nach der lateinischen und griechischen Grammatik wird durch die Gerechtigkeit GOttes verstanden die Gerechtigkeit, damit er gerecht ist. Aber diese Auslegung ist etwas dunkel und macht den ganzen Text in der Schrift auch dunkel und unverständlich. Wenn man es aber passivisch versteht, dann haben wir daran Theil und Trost. Denn es heißt die Gerechtigkeit GOttes, dadurch ich aus GOttes Gnade und Barmherzigkeit gerecht werde. Also, wenn du wollest die Rede, da man sagt „der Glaube GOttes“, nach der lateinischen Sprache auslegen, daß es so viel heißen sollte, als der Glaube, damit GOtt selbst glaubt oder seine Verheißung hält, so verdunkelt solche Auslegung den Verstand Pauli. Wenn ich es aber verstehe für die Gabe GOttes in mir, dadurch ich an GOtt glaube, da lautet es viel anders und besser. Also ist das GOttes Werk, nicht das er leidet von Andern, sondern das er in mir wirkt.

131. Vor Zeiten, wenn man lesen und beten sollte, das im 31. Psalm V. 2. steht: „Errette mich durch deine Gerechtigkeit“, erschrak ich ganz und gar davor, und war derselben Worte von ganzem Herzen feind und dachte: Errette mich ja nicht durch deine Gerechtigkeit, damit du gerecht bist (activisch). Warum soll man es verstehen für die Gerechtigkeit, dadurch ich gerecht gemacht werde (passivisch).

132. Auf diese Weise pflegt in der hebräischen Sprache die active und passive Bedeutung den Text sehr zu verdunkeln. Aber wer zwischen beiderlei Reden den rechten Unterschied weiß, der soll wissen, daß er ein Verstand der Schrift nicht wenig zugenommen habe. Darum soll man den Spruch Pauli Röm. 3, 23., da er sagt: Alle mangeln des Ruhms GOttes, auch passivisch verstehen, daß es so viel heiße, als wenn ich sage: „Ich habe nichts, daß ich mich vor GOtt rühmen könne. Wir können uns unsers HErrn GOttes nicht rühmen. Dies bringt der heiligen Schrift großes Licht. Dagegen aber verdunkeln der Juden Rabbinen Alles mit ihren Glossen, wenn sie die Worte activisch deuten, welche passivisch zu verstehen sind.

133. Laurentius Valla sagt recht, daß die griechische Sprache diesen seinen Unterschied nicht habe im Gebrauch der Fürwörter: mei und meus, tui und tuus. Denn es ist ein Anderes, wenn ich sage: amor mei, daß ich geliebt werde, und ein Anderes: amor meus, die Liebe, damit ich Andere liebe; item, odium tui et odium tuum, etc., das ist: Haß gegen dich und dein Haß. Nun mangelt es aber den Hebräern an diesen Worten noch viel mehr. Allein die lateinische Sprache sagt fein säuberlich: Hoc facio amore tui, hoc est, quo tu a me diligaris; dies thue ich aus Liebe, die ich zu dir trage, nicht aus deiner Liebe, damit du auch liebst. So ist das auch, wenn ich auf Lateinisch sage: Gratia tui ignosco tibi: Ich will dir es verzeihen, darum daß ich dir günstig bin, nicht, weil du mir günstig bist. Daraus erhellt nun leichtlich, welch ein großer Unterschied da sei zwischen der activen und passiven Bedeutung. Und dazu ist großer Fleiß vonnöthen, daß man dies also in der hebräischen Sprache merkte und darauf wohl Acht gebe. Denn dieselbige Sprach hat diese Fürwörter nicht. Und Paulus ist deshalb in seinen Episteln etwas dunkel und nicht wohl zu verstehen. Denn er hat der Reden viel nach Art der hebräischen und griechischen Sprache. Die deutsche Sprache kann sich auch gewohnter ausdrücken, weil sie das Wörtlein „wollen“ hat; als, wenn man sagt: um deinetwillen.

134. Dies sage ich nicht allein um dieser gegenwärtigen Stelle willen, sondern daß man auf diese Ungleichheit der Art und Weise zu reden an vielen andern Orten auch fleißig Acht gebe. Denn es hat mir vor Zeiten viel zu schaffen gemacht, wenn ich auf diese Sätze gekommen bin, als: die Gerechtigkeit GOttes, die Werke GOttes; item: dein Werk, deine Gerechtigkeit. Denn wo diese Reden activisch verstanden werden, sind sie der Tod; aber passivisch sind sie Leben und Seligkeit.

135. Wenn ich pflegte zu beten: Errette mich, HErr, durch deine Gerechtigkeit; o nein, dachte ich dann, lieber HErr GOtt, nicht durch deine Gerechtigkeit, wenn man auf lateinische Weise oder philosophisch reden will; sondern durch deine Gnade und Barmherzigkeit. Und haben allerdings hier alle Väter, Augustinus, Ambrosius und andere mehr dahin gesehen und sich daran als an ein Aergerniß gestoßen, als im 25. Psalm V. 10.: „Die Wege des HErrn sind eitel Güte und Wahrheit.“ Da haben sie die Worte: „die Wege des HErrn“, activisch gedeutet und gesagt, daß GOtt durch seine Wahrheit den Teufel und alle Gottlosen und verdammte Menschen richte; die Heiligen aber durch Barmherzigkeit. Damit ist aber der rechte Verstand des Psalms ganz verdunkelt worden, und sie haben davon nichts verstanden, dieweil sie das Wort „des HErrn“, wie dasselbe hier gesetzt wird, nicht recht verstanden haben. Denn der Psalm redet eigentlich von den Wegen und Werken des HErrn, die er in uns wirkt; von dem Glauben des HErrn, dadurch er uns selig macht. Also, daß dies der Verstand sei derselben Worte: Alle Wege des HErrn, das ist, was unser HErr GOtt mit uns macht, ist alles Gnade und Barmherzigkeit.

\subsection{I}

136. Derhalben sollen wir bedenken, daß an diesem Orte die Worte, „Brod brechen“, nicht vermögen verstanden werden, daß es das Brod oder die Theurung breche; sondern daß es heiße, die Austheilung des Brodes für oder wider die Theurung. Nehmet hin, sagt Joseph, und bringet heim, was ihr gekauft habt für die Theurung jetziger Zeit. Denn jetzt erzeiget er sich freundlicher, auf daß er dem alten Vater und seinem Hausgesinde Rath schaffen und helfen möge, daß sie inzwischen nicht Hungers sterben. Darum läßt er nun die Andern mit Gnaden hinziehen, behält nur den Einen gebunden in Verwahrung.

137. Er hat aber den Argwohn der Sünde, welche er ihnen vorgeworfen und deren er sie geziehen, noch nicht gänzlich fallen lassen, sondern wendet vor, er fürchte und scheue GOtt, dadurch er auch sei bewogen worden, seinen gefaßten Zorn gegen sie also zu mäßigen. Und spielet nun also mit seinen Brüdern, wie oft gesagt ist, ganz freundlich; welches doch ein wunderlicher Spiel, und dem äußerlichen Ansehen nach gar bitter und unfreundlich ist. Sein Herz ist voll Liebe und brüderlichen geneigten Willens, daß man nichts Freundlicheres oder Lieblicheres finden möchte, denn eben er gewesen ist, dieweil er bald hernach, da er sich von seinen Brüdern gemeldet, bitterlich geweint hat; und hat doch die natürliche Neigung jetzt verbergen können, daß er ihnen droht, sie zu tödten.

138. Darum soll man solche Exempel der Gnade GOttes gegen uns fleißig merken und bedenken. Joseph hat seine Brüder herzlich lieb, will ihnen nicht übel, und geht doch äußerlich so greulich mit ihnen um und läßt sie

\quellseite{386}

schier in die Gefahr des Todes kommen. Sie hatten sich zuvor viel einiger gehalten gegen Joseph; denn sie gedachten ihn umzubringen, wie droben im 37. Cap. V. 18. gesagt worden. Aber Joseph geht hier nun dahin um, daß er sie zum Leben und Seligkeit bringen möge, wiewohl er sich äußerlich anders stellt.

139. Also verhält es sich eigentlich auch um das Reich Christi und ein solcher König ist er; denn er hat immer etwas Gutes mit uns vor, wer es nur verstehen könnte. Und man soll ihm nicht glauben, wenn er zürnt, als meinte er es böse mit uns; denn in Wahrheit zürnt Christus, das ist, Gott, der Mensch geworden ist, nicht: aber außerhalb Christo ist kein Disputiren oder irgend eine Hoffnung von Gott; wie der Juden, Türken und Papisten Religion oder Glaube gewißlich ist.

140. Aber das Reich Christi ist ein Reich der Gnade und Barmherzigkeit. Wiewohl es aber nun scheint, als ob er uns feind sei, uns verachte und uns verwerfen wolle, wenn er Theurung, Pestilenz und Tod, und allerlei Plage über uns kommen lässet: so läßt er uns doch zugleich das zu und heißt uns auch, daß wir ihm vertrauen sollen. Und wir sollen daselbste auch wissen und glauben, dazu seine Güte und Gnade bekennen bis in den Tod. Auf diese Weise läßt er zu, daß Johannes der Täufer um der ärgsten Hurе willen enthauptet wird; item, daß seine liebe Mutter auch sehr betrübt, und die lieben Märtyrer mit jämmerlicher und dazu sonderlich auserlesener greulicher Pein und Marter geplagt worden sind. Scheint es aber nicht, daß er zürne? Nein, er zürnt wahrlich nicht, und du sollst dich dessen nicht überreden lassen; denn es ist kein rechter, sondern ein verstellter Zorn.

141. Also sagt Augustinus von sich selbst und von seiner Mutter: Wenn ich in der Schule gezüchtiget ward, und kam heim und klagte es meiner Mutter, lachte sie desselben und spottete meiner gleichsam dazu, daß ich gestäupt war; und sie wollte mir doch nicht übel. Wie kann aber Vater und Mutter dazu lachen, wenn der Sohn geschlagen wird? Es pflegt wahrlich also zu gehen. Die Eltern züchtigen ihre Kinder, die Schulmeister ihre Schüler, die Obrigkeit straft die, so es verschuldet haben, und ist eine Dienerin des Zorns und Schwerts, daß sie die Diebe hängen und die Mörder soll köpfen lassen; und hat doch zu solchen Strafen nicht Lust. Denn sie wollte lieber, daß jedermann thäte, was ihm in seinem Berufe gebühret, und gar keine Strafe verdiente.

142. Und dies geht zwar im Reiche des Zorns und des Todes also zu; darum hat es eine andere Meinung mit dem Haus- und Schulregimente. Denn die Eltern oder Schulmeister wollen weder meinen Traun das nicht, daß der, so mit Ruthen gezüchtiget wird und Streiche leidet, denkt und glaubt: sie suchen sein Verderben nicht, oder daß er damit für die Sünde sollte genugthun; mit welchem abgeschmackten Irrthum die Papisten diese Lehre besudelt und verfälscht haben: sondern sie thun eben, das Joseph hier auch gethan hat, welcher mit seinen Brüdern hart umgeht und sie wohl versucht, auf daß sie mögen unterrichtet und erinnert werden ihrer begangenen Sünde. Denn ohne Erkenntniß der Sünde ist kein Heil und kann ohne dieselbige niemand selig werden.

143. Also hat Gott das Gesetz gegeben, auf daß er dadurch die Sünde offenbarе und Zorn anrichte. Derhalben ist es eine Züchtigung und Ruthe, oder, wie Paulus redet Gal. 3, 24., „unser Zuchtmeister“, ohne welchen unser Fleisch die Sünde nicht verstehen kann. Wie ein Schüler, an welchem der Lehrmeister die Ruthe spart, nimmermehr in der Lehre zunehmen kann, sondern wird ein Klotz und grober Tölpel bleiben; denn er kennt oder versteht seine eigene Thorheit und Unglück nicht, wenn er ohne die Züchtigung und Strafe des Lehrmeisters ist. Derhalben ist die Erkenntniß der Sünden in allen Dingen nöthig, wiewohl es daran allein nicht genug ist. Aber Gott pflegt seine Gnade nicht zu geben, will uns auch nicht helfen, es sei denn, daß die Thorheit, so dem Knaben im Herzen steckt, und die Bosheit, so auch in des Mannes Herzen ist, geoffenbart werde, auf daß sie der Gnade fähig und theilhaftig werden mögen.

144. Dies ist nun Maß und Ziel aller Strafen und Plagen; wie die Eltern ihre Kinder züchtigen, auf daß sie fromm werden und sich bessern sollen, aber nicht, daß sie dieselben damit verderben oder enterben wollten. Denn solches wäre der Liebe und natürlichen Neigung, so von Gott ihren Gemüthern eingeprägt, gar zuwider. Und St. Paulus ermahnt Eph. 6, 4. und Col. 3, 21. fleißig, daß die Eltern ihre Kinder nicht sollen zum Zorn reizen, sie nicht erbitteren, auf daß sie nicht scheu werden. Und es soll solche Züchtigung sein, daß noch die Liebe und kindliche Zuversicht zu den Eltern und Schulmeistern da bleibe, und daß die Kinder verstehen mögen, daß sie die Strafe nicht leiden zu ihrem Verderben, sondern vielmehr darum, daß sie in ihrem Amt thun, was ihnen gebühret, und daß sie dem Vater mit ihrer Bosheit nicht Ursache geben, sie zu enterben; welches denen pflegt zu widerfahren, die sich weder mit Worten noch mit Streichen wollen zwingen lassen, daß sie fromm werden, und darnach um Dieberei, Mords- oder dergleichen anderer Uebelthaten willen dem Richter und Henker in die Hände kommen, daß sie von ihnen müssen gestraft werden, weil sie zuvor dem Vater nicht haben gehorchen wollen.

145. Auf solche Weise soll man das Exempel Josephs auf das Reich Christi ziehen, welcher auch straft, nicht daß er damit die Leute verderben oder enterben wolle, sondern daß er das Erbe, so er mit seinem Blut erworben hat, erhalten möge. Denn die Thorheit steckt in unserm Herzen, wir sind in Sünden empfangen und geboren, und bleiben darin stecken unser ganzes Leben und können des greulichen Schadens nicht gar los werden. Es ist ein teuflischer Unfleiß oder Unachtsamkeit in uns, wir sind kalt zum Gebet, und in andern Uebungen oder Werken der Gottseligkeit und rechten Gottesdienstes sind wir ganz lau, sind nicht so brünstig zur Liebe Christi: es ist uns damit nicht so ernst, wie mit dem zeitlichen Gut und was zu diesem gegenwärtigen Leben mehr gehört. Solche Faulheit und daß wir so gar schläfrig sind, kommt von unserm verfluchten Fleisch und vom Teufel her. Daher schreit St. Paulus also Röm. 7, 24. 19.: „Ich elender Mensch, wer wird mich erlösen von dem Leibe dieses Todes? Denn das Gute, das ich will, das thue ich nicht, sondern das Böse, das ich nicht will, das thue ich.“

146. Diese Tollheit und Thorheit der Erbsünde ist nicht allein an den Kindern, sondern auch an den Erwachsenen zu strafen. Warum schickt Gott Plage durch die Obrigkeit, Theurung, Pestilenz, Krieg, Türken, Papst und Teufel etc. Warum thut er aber das? Es geschieht es darum, daß wir zu ihm rufen sollen, wenn wir geschlagen werden. Ja, sagst du, hört er uns doch nicht, oder lacht unsers Schreiens und Weinens noch wohl dazu; wie die Eltern Augustinus auch lachten, wenn er klagte, daß er wäre gezüchtigt worden.

147. Also, wenn der Teufel durch einen Diocletian oder durch den Türken wider die Christen tobt und ihnen hart zusetzt, so seufzen sie und schreien: Ach lieber Herr Gott, höre doch unser Schreien, stehe doch an unsern Jammer und Elend! Warum läßt du doch unserer Trübsal und Elendes? Alsdann antwortet Gott: Liebes Kind, ich thue es darum, daß du sollst geläutert werden. Ich läutere dich aber darum, daß du sollst selig werden und zu dem verheißenen Erbe kommen, ich kann dich sonst nicht selig machen. Ich muß dich also tödten und begraben, daß du verfaulest und verwesest. Denn das Gift ist über deinen Leib ausgegossen, welches ich nicht kann wegnehmen, denn allein durch Diocletian, durch die Ketzer, die Arianer, welche ich dazu gebrauche als meine Ruthen.

148. Darum sollen wir bedenken, daß dies Stempel auf unser ganzes Leben soll gezogen werden; denn es lehrt uns Vieles auf eine leibliche Weise und zu unserm großen Nutzen. Christus ist unser sehr freundlicher lieber Joseph, welcher für uns gestorben ist und sein Blut vergossen hat zur Vergebung unsrer Sünden; darum kann er uns nicht abgeneigt sein. Und der Vater ist uns auch nicht feind; denn er sendet uns seinen Sohn. Vielweniger aber ist uns der Heilige Geist feind; denn derselbe lehrt uns dies alles und tröstet uns damit, und verhehlt doch und verbirgt oftmals diese Liebe mit schrecklichen Anzeichen seines Zorns.

149. Also züchtigt Gott diese Väter sehr hart nicht allein um der Erbsünde, sondern auch der wirklichen Sünde willen. Denn sie sind Vater-, Mutter- und Brudermörder gewesen. Und wo Gott nicht zuthat sonderbarer Weise erhalten hätte, so wäre er ohn allen Zweifel vor großem Herzeleid und Traurigkeit gestorben, nämlich, daß er immer ein Unglück über das andere gehabt. Er hat in seinem Jahr sehen müssen, daß Bilha geschwächt worden, Rahel und Debora, die Säugamme, gestorben, und endlich, daß Joseph auch umgekommen ist, welches Simeon und Levi gestiftet hatten. Zu solchen großen Knoten gehört eine starke Arznei, daß sie unten und oben auspeien. Derselben auf daß Joseph diesem großen Unglück, welches nicht zeitlich, sondern höllisch ist, nämlich, dem Gift der Erbsünde, rathen und helfen möge, damit es abgeschafft werde, hat er dazu ein heftiges und starkes Reinigungsmittel gebraucht durch Gefährlichkeit und Furcht des Todes und der Hölle: sonst wird nichts daraus. Denn sie kennen noch keine Buße, darum daß ihr Herz dick und stark geworden. Sie sind dick und faul geworden, daß sie sich wegen der begangenen Sünde nicht bekümmern. Nun macht aber die Sünde, wo sie verachtet wird, die Leute nicht gut oder fromm, die Gnade hat da keine Statt, es sei denn, daß wir unsere Sünde bekennen und uns demüthigen, ja, gar zu nichts werden, daß wir fühlen, wir sind verloren und verdammt. Sonst hält man es für keine Sünde.

\quellseite{387}

150. Es ist wahrlich ein wahrhaftiger und sehr guter Spruch, daß man sagt: Die Sünde ohne Gesetz ist keine Sünde; wie St. Paulus sagt: Die Sünde war in der Welt, aber sie ward nicht erkannt. Denn wo kein Gesetz ist, da ist auch keine Sünde, sondern eitel Gerechtigkeit. Wenn aber das Gesetz oder Gebot kommt, da wird die Sünde, welche lange geschlafen hat, wiederum auferweckt, daß alsdann der Mensch schrei(et) und sagt: Ach, hätte ich doch nimmermehr gemeint, daß ich so gar böse und verdammt wäre; ich habe bis daher mir selbst allezeit geheuchelt und mich für gerecht und fromm gehalten.

151. Also hat Joseph den rechten, wahrhaftigen und vollkommenen Nutzen des Gesetzes an seinen Brüdern geübt. Denn das heißt, in die Hölle führen, tödten und zu Schaden machen, ihnen nicht zum Verderben, sondern vielmehr zu ihrem Leben und Seligkeit. Denn das Gesetz ist nicht dazu gegeben, daß es tödten soll; wiewohl es in Wahrheit tödtet, wenn es Zorn anrichtet und uns die Sünde offenbart. Es tödtet aber nicht also, daß man im Tode bleiben und verloren werden soll, wie Judas und Saul sind verloren worden; sondern nur, daß es sein Amt an uns ausrichte. Dazu denn auch die Verheißung kommen soll, damit die Verzweiflung nicht darauf folge. Wie diese Brüder endlich gar kleinmüthig und schier verzagen werden. Und wo sie nicht diese Stimme gehört hätten, daß er zu ihnen sagt: „Ich bin Joseph“, so wäre es ihnen eine Versuchung zum Tode gewesen.

152. Darum ist das, allein Reue haben über die Sünde, keine rechte Buße; sondern es muß der Glaube auch dabei sein, welcher die Verheißung von Vergebung der Sünden durch Christum ergreift und sich daran hält, daß die bußfertigen Sünder nicht verloren werden. Diese Brüder sind aber noch nicht zu vollkommener rechtschaffener Erkenntniß der Sünden gekommen; darum hat sie Joseph noch etwas länger versucht: sie müssen noch besser durch die Rolle.

\subsection{II}

\begin{quote}
\textbf{V. 21. Sie aber sprachen untereinander: Das haben wir an unserm Bruder verschuldet, daß wir sahen die Angst seiner Seele, da er uns flehete, und wir wollten ihn nicht erhören; darum kommt nun diese Trübsal über uns.}
\end{quote}

153. Nun will es angehen. Hier offenbaren sie sich nun allgemach und fangen an ihre Sünde zu bekennen. Und Joseph hätte mit diesem ihrem Bekenntniß wohl mögen zufrieden sein, dieweil er weiß, daß sie die Wahrheit reden; er verstellt sich aber noch länger, als wüßte er von nichts. Denn sie würden solches nicht geredet haben, wo sie gewußt hätten, daß jener Fürst in Egypten Joseph wäre. Dieweil er aber die hebräische Sprache nicht braucht, sondern mit ihnen durch einen Dolmetscher redet, da denken sie, es sei nicht noth, daß sie selbst ihre eigene Sünde ihm offenbaren, sondern suchen gleichsam einen Deckel darüber, und wollen traun noch fromm, gerecht und ehrliche Leute sein, nicht zwar vor Gott, vor welchem sie die Sünde bekennen, sondern vor dem Fürsten des Landes, welchen sie nicht kennen.

154. Darnach weiß nun Joseph noch nicht gewiß, ob Benjamin und sein Vater Jakob noch am Leben sind oder nicht; denn wir werden hernach klarlich sehen, daß er durch diese ganze Historie immer noch gezweifelt hat an ihrer beider Leben. Darum kann er noch nicht zufrieden sein, dieweil er hört, daß nur sein allein gedacht wird und nicht der Andern auch. Es muß noch besser heraus, und er will haben, sie sollen noch weiter einen andern Bericht thun, der etwas klarer und gewisser sei, von Jedem insbesondere, derenthalben er sich besorgt, daß ihnen in seiner Abwesenheit Unheil möchte widerfahren sein. Dazu hat er auch an ihnen gesehen, daß sie noch nicht genugsame Reue gehabt über ihre Sünde, wiewohl sie sagen: „Wir haben dies verschuldet“ 2c.

155. Das hebräische Wort ascham heißt: Wir sind schuldig, und ist ein Wort, das mehrere Bedeutungen hat, welches Mose und Jesaja oftmals gebraucht haben. Ascham heißt ein Schuldopfer; also, Lev. 5, 3, 10.: „Wenn er sein Leben zum Schuldopfer gegeben hat, so wird er Samen haben, und in die Länge leben.“ 2c. An welcher Stelle dies Wort, wie auch in Mose, so viel heißt als ein Opfer für die Missethat oder Sünde. Auf Lateinisch heißt es reatus, Schuld.

156. Derhalben sagen hier die Brüder Josephs: Wir sind schuldig und sträflich, wir sind des Todes und Verdammniß schuldig. Denn wir haben solche Sünde wider uns, welche eigentlich eine Missethat, oder besser eine Schuld mag genannt werden. Sie sagen: Wir stecken noch in der Schuld, und haben noch nicht bezahlt, die Sache ist noch nicht geschlichtet und vertragen, unser Herr Gott hat noch einen Zuspruch zu uns, wir haben uns verschuldet.

157. Also sind wir auch noch täglich Schuldner; und wiewohl wir gerecht oder fromm sind, bleiben doch noch immer Schulden bei uns; wie Christus lehrt Matth. 18, 35., da er sagt: Also wird euch mein himmlischer Vater auch thun, so ihr nicht vergebet von eurem Herzen, ein jeglicher seinem Bruder seine Fehle. „Wir sind wohl gerecht, und für Kinder des Reichs erklärt; aber die Erbsünde bleibt noch in uns und streitet wider den Geist. Unser Herr Gott hat noch ein Kerbholz, ein Register wider uns, wir sind noch nicht gar rein von allem Fehl und Unflat. Dies ist die Schuld, nämlich, daß wir unserm Herrn Gott noch schuldig sind. Wenn du nun dazu thun wirst die Hoffart, Unbankbarkeit, Unfreundlichkeit, so an uns gespüret wird, gegen unsre Brüder, da wird die Sünde ganz und gar wieder lebendig und die Schuld ist dir noch nicht vergeben. Alsdann aber wird sie dir vergeben, wenn du deinem Nächsten auch vergibst; so soll die Schuld nicht schaden.

158. Der Knecht ist zehen tausend Pfund schuldig; aber der Herr erläßt ihm die Schuld, daß sie ihm nun nicht schadet. Er ist nicht schuldig, sie zu bezahlen oder dafür genug zu thun; und Schuld ist nun keine Schuld mehr. Woher aber? Ei dadurch, daß sie ihm erlassen ist. Aber wo der Knecht seinem Mitknecht die hundert Groschen nicht auch erlassen will, so fordert der Hausvater wiederum die ganze Schuld von ihm, Matth. 18, 23. ff. Also sind wir nicht ohne Schuld, wenn Gott mit uns nach dem Recht handeln will, und nicht nach Gnade und Barmherzigkeit. Wenn wir aber bitten, daß er uns die Schuld erlassen wolle, und glauben, daß er es gerne thun wolle, und verzeihen dann auch unserm Bruder: so werden uns alle unsere Sünden vergeben, wiewohl es Schulden sind. Wenn ich aber stolziren und pochen will, so steht die Schuld wieder da.

159. Dies ist der rechte Verstand des hebräischen Wortes ascham, welches Christus, die Weisheit des Vaters, welches der allerbeste Grammatiker und Hebräer ist, im Vater Unser verdolmetscht hat, Schuld. Also sagen diese Brüder auch: „Wir haben dies verschuldet“; wir sind einer großen schändlichen Sünde schuldig, die wir an unserm Bruder Joseph begangen haben; dieselbige wird nun aufgedeckt und kommt an den Tag, und wir müssen nun dafür Strafe leiden, wie wir es verdient haben: da findet sich unsere Sünde und geschiehet uns eben recht. Das ist wahrlich ein demüthiges Bekenntniß: sie bekennen, daß sie Schuldiger seien, erkennen auch die Grausamkeit, dadurch sie zuvor das erbärmliche Flehen und Bitten ihres Bruders verachtet haben.

160. Es hat aber Mose droben, Cap. 37, 23. ff., verschwiegen, daß er nun an diesem Orte erzählt, nämlich, daß Joseph jenemal seinen Brüdern zu Füßen gefallen und gefleht habe. Aber Simeon und Levi sind so gar halsstarrig gewesen, daß sie unangesehen seines Fußfalles und flehentlichen Bittens mit ihrem Bruder auf das allergreulichste sind umgegangen, und das um keiner andern Ursache willen, denn allein, daß ihn der Vater vor den andern sonderlich lieb gehabt, und daß er solche Träume gehabt, welche seine zukünftige Herrlichkeit bedeutet haben, dazu er sollte erhoben werden. So gar hart und unversöhnlich sind sie gegen ihn gewesen, daß sie ihn nicht haben hören wollen, da er ihnen zu Füßen gefallen, geflehet, bitterlich geweint und in höchster Angst und Noth gewesen ist.

\quellseite{388}

161. Solche große Sünde nun, nämlich, daß sie ihren Bruder verkauften und getödtet haben, da er sie in der äußersten Noth und Gefahr Leibes und Lebens angesehen, müssen sie nun erkennen. Und wahrlich, wenn du dies wirst genauer mit einander vergleichen, ist die Strafe, so Joseph ihnen auferlegt, noch viel geringer, als ihre Sünde gewesen ist, da sie ihren Bruder verkauft haben. Denn ist das nicht wider alle Freundlichkeit und Gottseeligkeit, daß man sich so gar nicht erweichen oder bewegen läßt, einiges Mitleid zu haben mit einem unschuldigen Menschen, der mit gefalteten Händen dich anflehet, und nicht fremde Leute, sondern seine eigenen leiblichen Brüder um Barmherzigkeit bittet? Wahrlich, es ist eine erschreckliche That gewesen, damit sie wider ihren alten Vater und Joseph, ihren Bruder, überaus schwerlich gesündiget haben, deren Grausamkeit sie noch nicht gänzlich verstehen und erkennen.

\subsection{III}

\begin{quote}
\textbf{V. 22—24. Ruben antwortete ihnen, und sprach: Sagt ichs euch nicht, da ich sprach: Verkündiget euch nicht an dem Knaben, und ihr wollet nicht hören? Nun wird sein Blut gefordert. Sie wußten aber nicht, daß es Joseph verstund; denn er redete mit ihnen durch einen Dolmetscher. Und er wandte sich von ihnen, und weinete. Da er nun sich wieder zu ihnen wandte, redete er mit ihnen.}
\end{quote}

162. Es scheint Ruben etwas frömmer gewesen zu sein, denn die andern, und daß er so schwerlich nicht gesündiget habe, wie Simeon und Levi. Denn wir haben droben, Cap. 37, V. 29. 30., gehört, daß er sehr erschrocken und betrübt worden und dazu sein Kleid zerrissen habe, da er wieder zu der Grube gekommen und Joseph nicht gefunden hat; wiewohl er um der Blutschande willen, so er begangen, beide des Vaters Gunst und auch die Erstgeburt verloren hatte. Darum sagt er: Was will es nun werden? „Nun wird sein Blut gefordert.“ Aber das Wort ist Joseph zu Herzen gegangen und hat ihn meinen gemacht, und daß er nun Mitleiden mit ihnen gehabt; da ist er weich geworden und geht ihm das Herz über. Simeon und Levi bekennen die Sünde noch nicht; denn sie meinen, Joseph sei schon vorlängst gestorben: darum lassen sie sich wenig anfechten, daß seiner gedacht wird. Ruben aber sammt den Andern erkennen die Rache und Strafe GOttes. Er sieht, daß die Strafe auf ihm liegt, und ist jämmerlich darüber erschrocken. Derhalben kann sich Joseph nun nicht länger halten, daß er nicht sollte durch brüderliche Neigung gar bewogen werden, da er das Bekenntniß und Bagen gehört, und daß er sich nicht sollte von ihnen wenden und weinen: er kann es nicht über das Herz bringen. Wie unser einer sich in solchem Falle auch nicht des Weinens würde enthalten können.

163. Aber siehst du nun wohl, wie Christo, unserm Herrn, zu Muthe ist, wenn er die Seinen straft, welch ein hitziger Feuerofen großer Liebe da sei, nicht allein geistlicher, sondern auch leiblicher Liebe? Wie Joseph durch die Worte Rubens also bewegt wird, daß er solche Bewegung seines Herzens nicht hat können überwinden. Er kann sein Herz nicht mehr halten. Er denkt: Sie halten mich für todt, und ich lebe in großer Ehre und Herrlichkeit.

164. „Sie wußten aber nicht“, wie der Text sagt, „daß Joseph solches hörete“, das ist, daß er es verstünde; „denn er redete mit ihnen durch einen Dolmetscher“, welcher auf Hebräisch meliz heißet. Lez heißt, Spötter, wie in Psalm 1, V. 1. steht. Wenn aber der Buchstabe m davor gesetzt wird, so deuten es die Hebräer, daß es Dolmetscher heiße; sie wissen aber doch die rechte Ursache nicht, daher diese Auslegung kommt. Im 1. Psalm kann das Wort nicht für Dolmetscher verstanden werden. Denn an den andern Orten heißt es auch Spötter oder Heuchler, der anders redet und anders denkt, oder der ein Ding anders auslegt, denn die Worte der heiligen Schrift lauten. Ob aber dasselbe Wort Bedeutung oder Ursprung sei, darum daß ein Dolmetscher fremde Sprachen gebraucht, kann ich nicht gewiß schließen. An diesem Orte wird es richtig also gegeben, daß es Dolmetscher heiße. Denn Joseph hat sich gestellt, als könnte er die hebräische Sprache nicht, darum hat er einen Dolmetscher gebraucht; welches ihm wahrlich wohl noch gethan, dieweil so mancherlei und so viel fremde Leute daselbst angekommen sind, Getreide zu kaufen.

165. Aber der Juden Rabbinen machen hier ein unnütz Geschwätz, wie sie zu thun pflegen, nämlich, daß Manasse, Josephs Sohn, dieser Dolmetscher gewesen sei, so er doch zu der Zeit auf das höchste nur acht Jahre alt gewesen ist. Wie hat man denn solchen jungen Knaben in so wichtigen Sachen brauchen können? Es hat müssen ein ernster ansehnlicher Mann sein, und der eine vielfältige Erkenntniß in den Sachen und Sprachen gehabt, welcher dazu mal bei dem Fürsten in Egypten dies Amt gehabt.

166. Dies ist aber vornehmlich wohl zu merken, daß Joseph hier also beschrieben wird, als der voll herzlicher Liebe, natürlicher Neigung und brüderlicher Freundschaft gewesen sei. Denn ob er seine Brüder etwas hart angeredet und hart mit ihnen umgegangen ist, so ist doch sein Herz entbrannt gewesen von natürlicher Neigung und geistlicher Liebe. Denn der Glaube und Heilige Geist verderbt oder zerbricht die Natur nicht; sondern da sie verderbet und zerbrochen ist, heilt er sie und richtet sie wiederum auf. Derhalben bleiben die allernatürlichsten Affecten und Neigungen in den Eltern, Brüdern und Eheleibern, welche durch die Gnade nicht aufgehoben, sondern erweckt werden.

167. Die Mönche haben vor Zeiten aus den Menschen Klötze und Steine machen wollen, die weder fühlen noch etwas verstehen könnten, und der Teufel hat haben wollen, daß dies als ein sonderlich Lob von seinen Heiligen sollte gerühmt werden, nämlich, daß sie gar durch keinerlei Affecten beweget würden, die rechten tollen Phantasten und siebenfältigen Stoiker. Die Gnade aber und der Heilige Geist nimmt der menschlichen Natur ihre Bewegungen nicht so gar weg, daß Jakob, der Vater, den Tod seines Sohnes nicht beweinen sollte. Denn das wäre der Natur gar zuwider, die von GOtt also geschaffen ist, daß sie solche Neigungen hat.

168. Wir haben aber vor etlichen Jahren gesehen, daß die Rottengeister sich auch unterstanden, eine solche Apathie in der Kirche einzuführen, daß die Christen sich aller natürlichen Affecten und Neigungen, als, Traurigkeit über das Absterben ihrer lieben Freunde, gar enthalten müßten: wie Münzer sich beflissen, dieselbe in all seinem Leben und Sitten zu erzeigen als eine sonderliche Heiligkeit. Denn da ihm einstmals angesagt worden, daß ihm ein Sohn geboren wäre, ist er vor dem Altar gestanden, als wäre er stumm und taub, hat sich dessen gar nicht gefreut, GOtt dafür nicht gedankt und gar nichts geantwortet, daß er damit anzeigte, er wäre gar ein Stock und Klotz, und hat sich darnach gerühmt, wie seine Natur gar geändert und getödtet wäre. Das ist in Wahrheit gut schwärmerisch gewesen, und viel ärger, denn der Stoiker Philosophen Phantasie hat sein können. Denn GOtt will, daß die Natur soll erhalten und nicht ausgelöscht werden, und heißt, daß man sie bessern soll, auf daß sie immer reiner werden möge, und daß die natürlichen Neigungen in den Göttlichen gemäßiget sein mögen, als in den Leiden, welche sie mit Gottesfurcht und Vertrauen nicht regieren, sondern lassen sich solche Affecten leichtfertig und ohne gewisse Regel des göttlichen Worts führen, wohin sie nur wollen.

169. Derhalben siehst du, daß in diesem Exempel Josephs beschrieben und abgemalt ist die Heiligkeit der allerheiligsten Väter, welche eben dieselbigen natürlichen Neigungen gehabt wie andere Menschen, wiewohl weniger und mäßiger; desgleichen haben sie auch ihre Leidenschaften und Schwachheiten gleich andern Menschen gehabt.

170. Der Evangelist Marcus schreibt am 3. Cap. V. 5., Christus habe die Phariseer umher angesehen mit Zorn, und sei betrübt gewesen über die Blindheit ihres Herzens, darum daß sie davon disputirt haben, ob man auch am Sabbath einen Kranken heilen möchte? etc. Wie kann aber Christus zürnen oder betrübt werden, der doch ohne Sünde ist? Und warum schreibt ihm der Evangelist das zu, daß er betrübt, traurig und zornig geworden sei? Antwort: Der Zorn, wenn rechtes Maaß dabei ist und daß er im rechten ordentlichen Gebrauche bleibt, ist ein natürlicher Affect und Bewegung, so der Natur des Menschen von GOtt eingepflanzt ist; wie die eheliche Liebe und Lust, so Mann und Weib zu einander haben, auch vom Anfang mit andern natürlichen Affecten und Neigungen geschaffen ist. Dieweil aber dieselben Affecten durch die Erbsünde verderbet sind, muß man sich befleißigen, daß sie mögen gebessert werden.

171. Darum ist das eine offenbare Unsinnigkeit des Teufels, daß man solche Heiligkeit rühmen will, da die Menschen gar zu Klötzen und Steinen werden, und aufhören Menschen zu sein. Sie sagen, daß die nicht heilig seien, so sich zum Zorn bewegen lassen; so doch die Stoiker selbst den Menschen vom Zorn und andern Affecten nicht gar haben befreien und ihm dieselben nehmen mögen. Denn sie stellen sich etwa, als zürnten sie nicht, oder aber fliehen die Ursache, so sie zum Zorn bewegen möchte, und die Gemeinschaft anderer Leute, welche sie erzürnen möchten. Auf die Weise wollte ich mich auch des Zorns enthalten, wenn mich niemand dazu reizt. Daher sind aber gekommen so viel Haufen toller und thörichter Heiligen, so in die Wüste gelaufen und aus der Versammlung und Gemeinschaft der Leute weggeflogen sind, auf daß sie für solche Leute angesehen würden, die gar frei wären von allen menschlichen Affecten und menschlicher Empfindung, wie ihre Werke und Reden in den „Leben der Väter“ beschrieben werden.

\quellseite{389}

172. Es wird aber an einem Ort erzählt das Exempel eines Mönchs, welches mir sonderlich gefällt. Denn da er sich seine Eigenthümlichkeit und Ungeduld hat überwinden lassen, und aus dem Kloster gar in die Wüsten gegangen ist, da er gehofft, sein still und friedsam, ohne allen Zorn und Sünde zu leben, soll es sich ohngefähr zugetragen haben, daß er aus der Wüste zu einem Brunn gegangen, in einem Krüge Wasser zu holen, und den Krug unversehens, da er zu viel Wasser hineingegossen, umgestoßen hat: da soll er zornig geworden sein, den Krug genommen, und auf die Erde geworfen und zerbrochen haben. Da er aber nun wiederum zu sich selbst gekommen, hat er es erkannt, daß es nicht der Leute Schuld wäre, mit welchen er umgegangen war, daß er oft zum Zorn bewogen worden, sondern seiner verderbten Natur und sündlichen Fleisches, welches immer von selbst zum Zorn geneigt wäre. Darum ist er wiederum in das Kloster gekommen und hat seine Affecten und Begierden lernen mäßigen, und sich befleißen, mit Andern Geduld zu haben.

173. Also will sich mit uns auch gebühren, daß wir bei den Leuten bleiben und mit ihnen Gemeinschaft und Gesellschaft halten, daß wir lernen die Sünde der Ungeduld, des Zorns und dergleichen anderer Affecten und böser Lust zwingen, brechen und denselben widerstehen; und also lernen, des Zorns und Lust, so Mann und Weib zu einander haben, recht brauchen, nämlich, daß die Hoffahrt, Ehrgeiz, Haß, Unzucht \&c. ausgefegt werden; welches ohne große Mühe, Arbeit und Schmerz nicht geschehen kann.

174. Darum steht man an den Heiligen nicht allein, daß die menschliche Natur mit ihren Affecten und Neigungen also geschaffen ist, welche der Heilige Geist nicht auslöscht; sondern man sieht auch Schwachheit und Verderben an ihnen, damit sie immer streiten, als die stets in Bereitschaft stehen, und sich darin bemühen, wie sie die gebrechlichen Neigungen an sich selbst tödten mögen.

175. Also haben Jakob und Mose auch viel erlitten nicht ohne große Schmerzen, da auch bisweilen Zorn, Toben, Murren und Ungeduld, ja auch wohl Gotteslästerung mit untergelaufen sind: sie haben sich aber in Schranken gehalten und solche sündliche Bewegungen bezwungen.

176. Es hat Joseph ohne Zweifel auch natürliche Unreinigkeit gefühlt, hat aber derselben den Zaum nicht gelassen, sondern den Zorn und Ungeduld, und dazu andere böse Lust gebrochen, hat dem Heiligen Geiste gehorcht und Gehorsam geleistet, welcher ihm das eingegeben, daß er nicht murren oder Böses mit Bösem vergelten sollte. Wiewohl er nun die gerechteste Ursache gehabt zu zürnen, und die Unbilligkeit, so ihn von seinen Brüdern begegnet, an ihnen zu rächen, führt er doch sein Recht so streng nicht hinaus, sondern verändert dasselbe und macht das freundliche Spiel daraus; und da er sieht, daß ihnen Angst geworden und daß sie über ihre Sünden Leid tragen, weint er mit ihnen.

177. Sie erkennen aber doch ihre Sünde noch nicht recht, wiewohl sie dieselbe unter einander heimlich bekennen. Denn sie würden ohne Zweifel dieselbe verhehlt und verschwiegen haben, wo sie gewußt hätten, daß Einer da wäre, der auf ihre Worte merkte und dieselben verstünde. Denn die Sünde hat diese Art und Natur an sich, daß sie die Leute taub und stumm macht, und ist niemand, der gern haben wollte, daß seine Sünde geoffenbart und daß er deshalb von Andern sollte gescholten werden.

178. Also verhehlen diese Brüder Josephs ihre Sünde auch und schmücken sich vor dem Fürsten des Landes. Da ist die Sünde noch stumm und taub; und sie hoffen, es werde die Sünde also für und für können verborgen bleiben. Dieweil sie aber mit dem Vater noch nicht versöhnt waren, ihm auch ihre böse That noch nicht geoffenbart hatten, hat die Noth erfordert, daß solches auch möchte an den Tag kommen, und daß sie getrieben würden, solche große Sünde auch noch zu erkennen und bekennen.

179. Joseph aber weint, da er doch die gerechteste Ursache gehabt, sich zu rächen, und sogar zugleich auch mit für Benjamin, seinen Bruder, indem er durch brüderliche Neigung bewogen worden. Denn er bekümmert sich um ihn, wie es ihm gehen möge, daß er nicht etwa umgekommen sei; und will deshalb ihren Worten noch keinen Glauben schenken, ist auch nicht schuldig ihnen zu glauben, bis er nun das Gegentheil erfahre, nämlich, daß sein Bruder noch lebe. Wie droben, Cap. 20, 2., Abraham auch um gleicher Ursache willen vom Sarah gesagt hat, sie wäre nicht sein Weib, dieweil er sich besorgt, es möchte an demselben Orte bei den Leuten keine Furcht Gottes sein.

180. Also wußte Joseph auch, daß seine Brüder Vater- und Brudermörder waren; darum schenkt er ihnen nicht alsbald Glauben, bis er seinen Bruder Benjamin sehen wird. Inzwischen versucht er sie wunderbarerlich, auf daß er von ihnen die Wahrheit erforschen möge, und daß er sie der verhaltenen Sünde, so nun fast bei zwei und zwanzig Jahren verborgen gewesen, aufs neue erinnern möge und ihnen dieselbe zu Gemüthe führen. Denn die lange Zeit über hatten sie gesehen, daß der höchst elende Vater von großem schweren Herzeleid und Wehklagen fast verzehrt war, deswegen, daß sie ihm vorgesagt und gelogen hatten, daß sein allerliebster Sohn Joseph umgekommen sei.

181. Derhalben sind Ursachen und Gelegenheiten genug gewesen, daß Joseph seine Brüder nach ihrem Verdienste gestraft hätte. Er geht aber noch viel gelinder mit ihnen um, denn sie mit ihm waren umgegangen. Denn er denkt nur darauf, wie er ihnen helfen und ihren Schaden heilen möge.

182. In dem Gespräche aber, das Mose nur obenhin beschrieben hat, zweifle ich nicht, Joseph werde von allen Dingen insonderheit gefragt haben, als, nach dem Vater, Brüder, nach ihrer beider Namen und was sie für ein Leben und Wesen geführt: welches alles hier nicht erzählt wird, aber aus den Umständen und aus der Historie an sich selbst wohl verstanden werden kann. Es ist auch wohl glaublich, daß er mich nach dem Alter, Verstand, Sitten und Leben eines Jeden insonderheit, die daselbst vor ihm gewesen sind, gefragt haben, auf daß er desto leichter ausforschen möchte, was er wissen wollte, nämlich, von ihrer begangenen Sünde, von seinem Bruder und alten Vater.

\section[Fünfter Theil]{Fünfter Theil. Wie Joseph Simeon bindet; wie er die andern Brüder mit ihrem Geld und mit Getreyd wieder fortschickt; wie Einer von ihnen in der Herberge seinen Sack aufthut, sein Geld wieder findet; und wie ihnen allen dabey zu Muthe}

\subsection{I}

\begin{quote}
\textbf{V. 24. Er nahm aus ihnen Simeon, und band ihn vor ihren Augen.}
\end{quote}

183. Dieser Text zeigt an, daß Simeon der vornehmste Bruderörder gewesen sei und der allerbößeste unter den Brüdern. Denn nachdem Ruben die Erstgeburt verloren hat wegen der Blutschande, so er mit seiner Stiefmutter begangen, oder vielmehr mit seiner Mutter, darum daß sie ja mit seinem Vater Ein Fleisch war, hat er damit auch zugleich alle Ehre und Herrlichkeit abgelegt und sich derselben entäußern müssen: und ist allerdings auch hernach nicht allein er, sondern alle seine Nachkommen und ganzes Geschlecht sehr schlecht und verachtet gewesen und hat nichts Merkwürdiges ausgerichtet, wie das in den Büchern der Könige und sonst hin und wieder in der Heiligen Schrift zu sehen ist.

184. Es wird aber Simeon bei sich selbst also gedacht haben: Siehe, Ruben liegt nun gar darnieder und wird für nichts gehalten um der Sünde willen, die er mit seiner Stiefmutter begangen hat; so wird nun die Erstgeburt auf mich kommen. Da wird der Esel stolz, geht auf das Eis und bricht ein Bein. Er hat sich hervorgethan, ist muthwillig und stolz geworden, und hat dies greuliche Spiel vornehmlich gestiftet und angerichtet, und das Feuer ganz und gar angezündet. Denn die Worte, so droben, Cap. 37, 19., erzählt werden: „Siehe, der Träumer kommt“ \&c., schreiben die Hebräer Simeon zu, daß er dieselben soll geredet haben; denn sie haben nicht alle zugleich diese Worte reden können. Darum reimt es sich wohl, daß Simeon, da er Joseph hat kommen sehen, nicht allein in solche Worte ausgebrochen sei, sondern auch die Andern angereizt habe, die große Uebelthat an ihrem Bruder zu bekennen. Also ist er viel grausamer und unfreundlicher gegen Joseph gewesen denn die Andern alle; und wo es ohne ihn gewesen wäre, hätten vielleicht die Andern ihres Bruders verschont, da er sie so flehentlich gebeten, seiner zu verschonen. Aber dieser Simeon hat um der Erstgeburt willen dem Joseph so viel härter zugesetzt, und darum, daß ihn der Vater sehr lieb gehabt; darnach auch, daß er demjenigen, so Simeon zu erlangen begehrte, im Wege gelegen. Darum hat er gedacht, es wäre der allerbeste Rath, daß nur Joseph umgebracht würde.

\quellseite{390}

185. Dies ist sehr wahrscheinlich und kann leichtlich aus der Historie verstanden werden. Denn wie vom Anfang der Welt die Erstgeburt eine Ursache großen bitteren Hasses und Zanks gewesen, also hat sich Simeon derselbigen auch gelüsten lassen. Und wir streiten jetziger Zeit wider den Papst um keiner andern Ursache willen, denn eben auch um die Erstgeburt, das ist, um der Kirche willen. Denn die Papisten schreiben diese Herrlichkeit sich zu, sie wollen die Kirche und Gott die Nächsten sein. Dagegen aber halten wir sie für die, so von Gott verstoßen und verdammt sind, und maßen uns an die Ehre und Namen, daß wir Gottes Volk und seine auserwählten Kinder heißen. Denn dasselbe ist die Erstgeburt, nämlich, die Kirche sein, und die Leute regieren und lehren zur Seligkeit: unter Gottes Volk gerechnet und in die Zahl Abels und anderer frommen Gottes Diener mitgezählt werden. Da legt sich dann etwa ein Ketzer oder Schwärmer wider uns und rühmt sich, er sei ein Doctor und Lehrer der Kirche, verdammt und verachtet die Andern alle, und tritt sie gleichsam mit Füßen. Daher kommt der Streit und Kampf, welchen die rechte wahre Kirche zu allen Zeiten mit ihren Feinden und Widersachern gehabt.

186. So wird nun Simeon gebunden, als der das Haupt und Rädelsführer dieses Unglücks gewesen ist, die Andern aber läßt er hinziehen; nicht daß Joseph Lust gehabt, sich zu rächen oder sein Mäntlein an ihm zu kühlen, sondern es ist ihm nur darum zu thun, daß er ihn zur Erkenntniß der Sünden bringen möge. Ruben und Juda mit den Andern sind traurig gewesen wegen der begangenen Sünde; aber Simeon ist halsstarrig und stolz, und gedenkt noch, die Hoffnung der Erstgeburt nicht fallen zu lassen; denn er ist desselben gewiß, daß Ruben verworfen und Joseph todt sei. Und dieweil Joseph dasselbe gemerkt, läßt er Simeon in die Bande werfen, auf daß ihm also der Stolz und Vermessenheit durchs Gefängniß gebrochen werde, und daß er nicht denken soll, er werde der Erstgeborne sein, sondern daß er ewig werde gebunden und gefangen sein müssen.

187. Danach hat Levi sich selbst mit dieser Herrlichkeit trösten und sich dieselbe verheißen mögen; es wird aber ihnen beiden fehlen, wie man das hernach aus des Vaters Fluch sehen wird. Juda, der vierte Sohn, hat das Regiment bekommen; die Herrlichkeit aber der Erstgeburt, so Ruben verloren hatte, ist Joseph gegeben worden, und sind ihm zwei Theile gefallen, als dem Könige und Priester, durch seine zwei Söhne, Ephraim und Manasse, von welchen zwei Stämme hergekommen sind.

\begin{quote}
\textbf{V. 25. Und Joseph that Befehl, daß man ihre Säcke mit Getreide füllte, und ihr Geld wiedergäbe, einem jeglichen in seinen Sack, dazu auch Nahrung auf den Weg; und man that ihnen also.}
\end{quote}

188. Das Wort „Sac“ ist ein gemein Wort in vier Sprachen, wie wir droben angezeigt. Hier hat Mose das hebräische Wort kele gebraucht, daß es auch Sac heiße: und das ist ein gemein Wort für allerlei Instrumente oder Werkzeuge, als, vas belli, Kriegsausrüstung. Und hernach im 49. Cap. V. 5. wird von Simeon und Levi gesagt: „Ihre Schwerter sind mörderische Waffen“ 2c. Darnach wird noch ein ander Wort folgen, amtachath, welches umschreibend auch Sac bedeutet, und kommt von mathach, welches heißt, ausbreiten. Als, wenn bei uns die Elbe eine Biegung macht, dasselbe nennen wir in unserer Sprache einen Sack, item, eine wässerige Wolke; oder die Segel an den Schiffen, wenn der Wind sie aufschwellet, haben auch die Gestalt eines Sackes. Also, ein Sack, wenn er gefüllt wird, breitet er sich aus, und hat auch daher den Namen, daß er sich also ausbreitet und aufbläst.

\subsection{II}

189. Warum thut aber Joseph Befehl, daß man ihnen das Geld soll wieder geben? Antwort: Dazu ist er nicht allein bewogen worden durch die Cherbietung, so er dem Vater schuldig gewesen, welchen er damit hat beschenken wollen, sondern die natürliche Neigung, so er in seinem Herzen gegen zu Vater und Brüdern, welche er heftig lieb gehabt, getragen, hat ihn auch dazu verursacht. Diese freundliche Liebe, so billig und recht gewesen, ist die Ursache dieser Beschenkung; denn es ist in Wahrheit eine rechte natürliche Neigung, so er zu ihnen gehabt. Und fährt doch gleichwohl immer fort und stellt sich hart und grausam gegen sie. Denn er will haben, daß sie in der Furcht und Schrecken bleiben sollen, daß sie nicht meinen, daß ihnen solch Geld aus einiger Freundschaft und Gnade gegeben sei; sondern daß sie sich fürchten vor dem Argwohn und daß er gegen sie eine Klage vornehmen werde, als hätten sie ihm das Geld gestohlen und entführt. Nicht daß solches Josephs Ernst sei, sondern er stellt sich nur damit also, auf daß sie gedemüthigt werden, und daß die Erbgründe auch geböthet werde, welche gestärkt und gemehrt worden ist durch die wirkliche Sünde, nämlich, daß sie ihren Bruder erwürgt, welches nun bei zwei und zwanzig Jahren verborgen und veraltet war. Darum stellt er sich nun, als wollte er sie mit dieser falschen Anklage und List dazu bringen, daß sie ihm ewig würden müssen unterworfen und dienstbar sein, wie sie ihn zur Dienstbarkeit verkauft hatten. Dies ist die Ursache gewesen, daß er sich also gestellt; dadurch sie wahrlich sehr furchtsam geworden, daß sie sich besorgt haben, sie würden ohne großen Schaden und Gefahr von dannen nicht kommen.

190. Er hat ihnen aber zu dem Getreide und Gelde auch noch die Zehrung und Futter auf die Reise geben lassen, auf daß sie zu ihrem Vater Jakob kommen mögen und die Säcke so lange zulassen. Das hebräische Wort, zedah, heißt hier, Zugabe, und kommt vom Wort zud, jagen; wie 1 Mos. 10, 8. 9. von Nimrod gesagt wird: „Nimrod fing an ein gewaltiger Herr zu sein auf Erden, und war ein gewaltiger Jäger“ 2c. Mose hat das Wort an diesem Ort gebraucht für Speise oder Futter, wie im 132. Psalm V. 15. steht: „Ich will ihre Speise legen, und ihren Armen Brods genug geben.“

Hieronymus hat es am selbigen Orte mit „Jagd“ gegeben. In der gemeinen lateinischen Dolmetschung steht „Wittwe“; aber ich halte dafür, daß solches der Schreiber Schuld sein muß, welche für das griechische Wort κρέας gesetzt haben κήραρ.

191. Es ist aber ein wunderlicher Ursprung und Bedeutung dieses Worts, dieweil es herkommt vom Wort zu und doch Speise heißt; wie wir nun etlichemal vom Wort scheber gesagt haben, welches brechen heißt. Und ich sage es oftmals, daß die Hebräer selbst alle Worte in ihrer Sprache nicht recht verstehen. Denn das ist wider die Art der Grammatik, wenn du wolltest „brechen“ auslegen für Speise, Nahrung oder Getreide.

192. Wir haben diese Ursache angezeigt, daß die Juden die Gewohnheit gehabt, daß sie das Brod gebrochen und nicht mit Messern geschnitten haben. Und diese Ableitung desselben Worts ist besser, denn der Hebräer, die da sagen, es habe daher den Namen, daß es den Hunger breche. Denn die Juden haben Brod gehabt wie Kuchen und Fladen, die wir auch pflegen zu brechen. Also sind die Schaubrode sechs Kuchen gewesen, auf einander gelegt. Das bedenken die Juden nicht und schickt sich doch dieselbige Auslegung sehr wohl.

193. Also wissen sie auch nicht, warum das Wort zeda, Jagd, Speise heißen soll. Die Jagd aber betrifft nicht allein fromme Männer, sondern auch wohl Mörder, ja, unvernünftige Thiere, als, Wölfe, Löwen, Füchse. Und wiewohl ich es nicht für gewiß sagen kann, halte ich doch dafür, daß dies der Ursprung des Worts sei, darum daß die Jagd zu seinem andern Zweck geschieht, denn daß man Speise von den Thieren haben möge. Denn darum jagt man die wilden Thiere, auf daß die Speise und Nahrung hinreichend möge. Eben dasselbe suchen auch die Geier und Weihen, und was solcher Thiere mehr sind, die vom Raube leben.

194. Und daraus ist abermals zu sehen, daß die hebräische Sprache noch nicht so genau erforscht ist, sondern daß die hebräischen Grammatiker selbst viele Weisen zu reden, Figuren und Ableitungen der Worte nicht wissen. Das ist aber nun der erste Theil in dieser Historie, da die Söhne Jakobs nach Egypten kommen,
\quellseite{391}
Getreide zu kaufen, und ihren Bruder Joseph noch nicht gekannt haben. Nun folgt der andere Theil.

\subsection{III}

\begin{quote}
\textbf{V. 26—28. Und sie luden ihre Waare auf ihre Esel, und zogen von dannen. Da aber einer seinen Sack aufthat, daß er seinem Esel Futter gäbe in der Herberge, ward er gewahr seines Geldes, das oben im Sacke lag. Und sprach zu seinen Brüdern: Mein Geld ist mir wieder worden, siehe, in meinem Sack ist es.}
\end{quote}

195. Das Wort malon heißt bei den Hebräern eine solche Stätte oder Ort, da man eine Zeitlang verharrt oder die Nacht über bleibt, das ist, eine Herberge oder Wirthshaus. Sie disputiren aber an diesem Orte davon, welcher der gewesen sei, der den Sack aufgethan. Nun ist daran nichts gelegen, man sage, daß es Levi oder Juda, oder aber der Andern Einer gethan habe. Und es erhellt daraus auch klarlich, daß die Beschwerung mit allen Dingen und die Theurung muß groß gewesen sein, dieweil die Esel und Menschen einerlei Futter mit einander gehabt; welches ohne Zweifel nicht Heu oder Spreu wird gewesen sein, sondern Getreide, damit sie in Egypten die Säcke gefüllt, darein auch das Geld gelegt gewesen.

196. Das hebräische Wort, mispo, ist dem Wort bar, welches wir droben gehabt, in der Bedeutung gar zuwider. Sie sagen, es heiße, Futter, nicht von Gras oder eigentlichem Getreide, sondern das gekörnt ist, als, Hafer, Wicken oder dergleichen.

197. Es wird auch gefragt, warum hernach, Cap. 43, V. 21, die Söhne Jakobs vor Joseph sagen: „Da wir in die Herberge kamen, und unsere Säcke aufthaten“: so doch alhier nur von Einem gesagt wird, der seinen Sack aufgethan. An beiden Orten aber wird gesagt, daß sie alle erschrocken seien. Ezra gibt die Antwort, daß in der heiligen Schrift etliche Dinge erzählt werden durch die Figur, die da heißt anticipatio oder recapitulatio (Vorwegnehmen oder Wiederholen), wenn etwas, das zuvor gesagt ist, wiederholt wird. Welche Regel nöthig ist, und in der Theologie oder heiligen Schrift, ja auch in allen Historien fleißig gemerkt werden soll. Sie ist aber genommen aus dem andern Buch Augustins De doctrina Christiana, da er sagt, er habe sie mit etlichen andern Regeln mehr von einem Keßer und Donatisten, Ticonius genannt, empfangen. Denn also sagt er daselbst: In der Schrift werden etliche Dinge so berichtet, als wenn die Erzählung die genaue Zeitfolge beobachtete, während doch die Erzählung zu früheren Ereignissen, die erst ausgelassen waren, unvermerkt zurückkehrt. Als, 1 Mos. 11, 32. steht geschrieben: „Und Tharah war zwei hundert und fünf Jahr alt, und starb in Haran.“ Und im folgenden 12. Cap. V. 4. wird gesagt: „Abraham war fünf und siebenzig Jahr alt, da er aus Haran zog.“ Da erzeugt die angegebene Zahl der Jahre Tharahs große Verwirrung und Irrung, wenn man nicht recht auf die Ordnung der Erzählung merkt. Derselben ist an dem Orte diese Regel von großer Bedeutung. Denn so Einer den Ausgang Abrahams verstehen wollte nach der Ordnung der Erzählung im Text, wie es sich ansehen läßt, daß die Worte des Stephanus, Ap. Gesch. 7, 4., fast also lauten: so wollte folgen, daß Abraham aus Haran gezogen wäre im fünf und siebenzigsten Jahr nach dem Tode Tharahs, welches unmöglich ist.

198. Darum ist das die rechte Ordnung: Da Mose die Erzählung von Tharah hat wollen zum Ende bringen, sagt er also: Tharah war zwei hundert und fünf Jahr alt und starb. Er erzählt aber der Dinge keines, so vor seinem Tode geschehen waren. Derhalben kommt er darnach wiederum darauf, und erzählt das auch, was in den zwei hundert und fünf Jahren geschehen ist, ehe Tharah gestorben, nämlich, daß Abraham, sein Sohn, da er fünf und siebenzig Jahr alt gewesen, aus Haran gezogen sei, da sein Vater Tharah noch gelebt. Sonst müßte er sechzig Jahre später ausgezogen sein, als er wirklich ausgezogen ist.

199. Darum lauten die Worte nach dieser Figur, recapitulatio (Wiederaufnahme) genannt, als ob er nach seines Vaters Tode ausgezogen sei, so er doch beinahe sechzig Jahre zuvor ausgezogen ist. Stephanus aber, wie es scheint, hat der gemeinen Erzählung derer gefolgt, so diese Figur nicht kennen, hat nicht allzu genau disputiren wollen von der Ordnung, wie die Dinge nach einander ergangen wären, sondern ist damit zufrieden gewesen, daß er nur die Worte nach ihrer Ordnung gesetzt, und ist bei der gemeinen Meinung geblieben.

200. Wir haben droben, Cap. 38, V. 1., noch ein ander Exempel gehabt, da beschrieben wird, wie Juda von seinen Brüdern weggezogen sei, nachdem Mose erzählt, wie Joseph war verkauft worden; so doch Juda bei neun Jahren, ehe sie Joseph verkauft, von ihnen gezogen war. Derhalben erzählt Mose durch diese Figur, recapitulatio, daß Juda sei von seinen Brüdern hinabgezogen, auf daß er also wiederum darauf kommen möge, was zuvor geschehen war, ehe Joseph verkauft worden. Denn Mose hat vorgehabt, die Historie Josephs ohne Unterbrechung bis zu Ende fortzusetzen, und unterdes das auszulassen oder aufzuschieben, was zuvor oder hernach geschehen wäre.

201. Eben auf dieselbige Weise steht auch im 35. Cap. V. 28. 29. von Isaac geschrieben, daß er gestorben sei, da Jakob aus Mesopotamien zu Hebron angekommen ist; und ist doch noch am Leben, und sieht, daß Ruben die Blutschande begangen und Joseph verkauft worden ist, und hat noch gelebt bis auf das Jahr, das zunächst vor der Erhöhung Josephs hergegangen.

202. Darum ist diese Regel in allen Historien wohl zu merken und darauf gute Acht zu geben. Denn es ist eine andere Ordnung der Worte oder Reden und eine andere der Zeit oder Geschichte. Und die Geschichtschreiber, wenn sie sich vorgenommen haben, eine That zu beschreiben, pflegen sie immer in Ordnung und in Zusammenhange die Erzählung fortzuführen. Wie hier Mose auch erzählt, daß die Brüder Josephs in ihren Säcken das Geld gefunden, da sie wiederum aus Egypten gekommen; und sagt doch nach der Figur anticipatio (Vorwegnehmen) zunächst nur von Einem, daß er das Geld gefunden habe, bis er es später dann auch von den Andern berichtet. Und es ist durchaus wahrscheinlich, daß, da in Eines Sack das Geld gefunden worden ist, die Andern darüber auch erschrocken sind, und daß ein Jeder seinen Sack werde aufgethan und darin auch Geld gesucht haben.

203. Es reimt sich auch nicht übel, daß sie mehr Esel gehabt denn zehn. Denn es hat Ein Esel nicht mehr denn zwei Säcke tragen können; welches Gewicht oder Last nicht genug gewesen wäre auch nur einen Monat lang für solch großes Gesinde, in welchem so viel Söhne, Töchter, Weiber, Knechte, Mägde und Vieh gewesen. Derhalben glaube ich, daß ein Jeder drei oder vier Esel mit Getreide mehr beladen haben, damit sie nicht so bald wiederkommen müßten. Denn der Weg ist lang gewesen von Canaan heraus bis gen Helipolis oder Memphis. Und ich halte dafür, es sei bei fünfzig deutscher Meilen von einander gewesen; welchen Weg sie nicht so bald haben ziehen können, dieweil die Esel langsam fortgehen. Darum haben sie so viel gekauft, daran sie etliche Monate haben können genug haben, ihr Gesinde und Vieh damit zu speisen. Wenn man das nun so annimmt, so wäre es nicht so gar ungereimt, daß man es dafür hielte, daß durch Schickung Gottes ein Jeder seinen Sack aufgethan, und doch Einer nur das Geld gefunden, die Andern aber alle hätten nur ohngefähr und zufälliger Weise einen oder mehr Säcke durchsucht, darin doch nichts mehr gewesen denn Getreide. Denn sie haben unterwegs, da sie gezogen sind, so viel Meile nicht haben können, daß sie die Säcke alle so genau besichtigt hätten.

204. Und beweiset mich vornehmlich dazu dies Argument, daß sie zwar hier erschrecken; hernach aber, V. 35., wird Mose abermals sagen, daß sie sammt ihrem Vater erschrocken sind. Denn sie würden nicht zum andernmal erschrocken sein, wenn sie zuvor alle in der Herberge gesehen hätten, daß ihnen das Geld wiedergegeben worden.

\subsection{IV}

\begin{quote}
\textbf{V. 28. Da entfiel ihnen ihr Herz, und erschraken unter einander, und sprachen: Warum hat uns Gott das gethan?}
\end{quote}

205. Dies ist eine neue Art und Weise zu reden: Ihr Herz ging aus, welche man auf Deutsch wiedergeben kann: Das Herz ist ihnen entfallen. Er zeigt aber damit an, daß ihnen alle ihr Trost, Weisheit und Stärke entfallen sei. Denn sie sind nicht allein darüber erschrocken, daß sie das Geld im Sacke gesehen haben, sondern daß ihnen ihre Sünde, so sie begangen, vor Augen gekommen ist. Denn sie haben bis daher versucht, dieselbige zu bedecken und zu verhehlen, und haben bei sich selbst einen närrischen Trost in ihren Herzen gefaßt, nämlich, die Sünde würde ewig verschwigen und vergessen bleiben, soweit nicht Joseph durch das Spiel, das er mit ihnen getrieben, diese Hoffnung schon zweifelhaft gemacht hatte.

\quellseite{392}

206. Es ist ihnen aber doch noch nicht recht leid und sie erkennen ihre Missethat noch nicht. Derhalben wird ihnen nun ihr Herz oder der Trost, so sie bei sich selbst gehabt, als ob die Sünde gar sollte vergessen sein, genommen. Wiewohl sie die Sünde noch nicht bekennen, sondern in ihrem heimlichen Murmeln und Stilleschweigen sind sie noch der Meinung, daß sie werden davonkommen, doch nicht ohne Furcht und Bekümmerniß. Darum sagt Mose, sie seien erschrocken und irre geworden unter einander; Einer hat den Andern angesehen und gedacht: Das geht nicht recht zu. Aber sie geben die Schuld unserm Herrn Gott, welcher der Richter ist. Wie? sagen sie, hat Gott dies also gemacht? Und solches ist hergekommen aus Schrecken und Angst des Gewissens, welches auf die Sünde mit großem Schmerz und herzlicher Pein zu folgen pflegt.

207. So eine böse Bestie und böser Teufel ist das Gewissen. Denn alle Schreiber, beide, so die heilige Schrift und auch heidnische Historien beschreiben, haben dies Ungeheuer, dies greuliche Thier erschrecklich abgemalt, wie das an Orestes und andern Uebelthätern zu sehen ist. Und die Poeten haben darum die greulichen Personen in den Tragödien erdichtet, die Erinnyen und Furien, das ist, höllische Teufel, welche alle Uebelthat rächen. Alle reden von demselben Unglück und Herzeleid, das da heißt: Mens sibi male conscia, das ist, ein Herz, das sich selbst der Sünden schuldig weiß. Die heilige Schrift aber sagt kürzlich also: „Die Gottlosen haben keinen Frieden“ Jes. 48, 22. und 57, 20. 21. Es ist nichts Unfriedlicheres und das mehr Unruhe bringe, als ein erschrockenes Herz, welches, wenn es nur einen Blitz sieht, ja, auch ein Blatt rauschen hört, bleich wird.

208. Also haben die Söhne Jakobs auch ein bös Gewissen, darum zittern sie so greulich und fürchten sich vor göttlicher Rache. Sie fürchten, er werde mit Donner und Blitz dreinschlagen, und bekennen doch die Sünde noch nicht. Das ist aber ein wunderlich Ding: das Gewissen fühlt Gottes Gericht und seinen schweren Zorn; und schämt sich doch, Buße zu thun und die Sünde zu bekennen, wenn es auch in ewiger Verdammniß und Pein sollte stecken bleiben. Und das geschieht wegen der Erbsünde, wie das an Adam zu sehen ist. Denn da derselbe schon überwunden war und vor Gott floh, darum um daß er sich des Falles schuldig wußte, war doch jene erste Sünde in ihm noch stumm, und konnte nicht leiden, daß ihm die Ohren geöffnet würden, wiewohl die Stimme Gottes in seinen Ohren erschollen war: „Adam, wo bist du?“ 1 Mos. 3, 9.

209. Erasmus schreibt an einem Orte von einer lacedaemonischen Jungfrau, welche, da sie durch Hurerei schwanger geworden, das Kindlein, da es geboren gewesen, lieber hat tödten wollen, denn ihre Sünde bekennen; welche That er also deutet, daß solches eine große Stärke- und Mannheit an ihr gewesen sei, so es doch vielmehr eine Schwachheit ist, dadurch sie lieber sich in Gefahr des ewigen Todes und Verdammniß hat begeben wollen, denn daß sie hätte wollen ihre Sünde bekennen und Buße thun. So groß ist die Tollheit, Wuth und Unsinnigkeit der Erbsünde, welche sich nicht will lassen zu Schanden machen, oder zulassen, daß ihre Schande gesehen werde. Und es ist ganz wider die Art und Natur der Sünde, bekennen, daß man gesündigt habe.

210. Also sagen sie, daß St. Martin dem Teufel soll geantwortet haben, da er sein Gespött hat, und an ihm gestraft, daß er den bösesten und falschhaftigsten Menschen ihre Sünde so leichtlich vergeben hat: Wenn du auch könntest von Herzen sagen: Es ist mir leid, daß ich gesündigt habe, so wollte ich dich auch absolvieren und von Sünden ledig sprechen. Denn der Teufel würde auch selig sein, wie wir alle sind, wenn er aus reinem Herzen sagen könnte: Ach GOtt, erbarme dich mein und sei mir gnädig!

211. Die Absolution und Vergebung der Sünde ist gewiß; allein, daß du von Herzen erkennest und bekennest, daß du gesündigt hast. Aber da kommt der Teufel nicht hin. Adam und Eva haben das nicht thun können. Denn also sagte Adam 1 Mos. 3, 10.: „Ich hörete deine Stimme im Garten, und fürchtete mich, denn ich bin nacket.“ Er sagt nicht: Darum, daß ich dein Gebot übertreten habe, bis ihn GOtt selbst dahin bringt und ihn weiter fragt V. 11.: „Wer hat dirs gesagt, daß du nackend bist? hast du nicht gegessen von dem Baum, davon ich dir gebot, du solltest nicht davon essen?“ Noch will er nicht bekennen, und erkennt oder bekennt die Sünde, so ihm vorgeworfen wird, noch nicht, sondern gibt die Schuld dem Weibe, V. 12.: „Das Weib“, sagt er, „das du mir zugesellet hast, gab mir von dem Baum, und ich aß.“ Siehe, da er schon erschrocken und furchtsam gewesen, und gezittert hat, da er fühlte, daß er gesündigt hat, bekennt er doch die Sünde noch nicht.

212. Davon haben wir an Saul auch ein gleich Exempel, der da sagt 1 Sam. 15, 30.: „Ich habe gesündiget; aber ehre mich doch jetzt vor den Aeltesten meines Volks und vor Israel“, als wollte er sagen: Hilf die Sünde zudecken. Dasselbe taugt aber auch vor Menschen gar nichts. Denn Sünde thun und dieselbe verleugnen macht die Sünde noch einmal, ja, siebenmal so groß, als sie sonst an sich ist. Wenn man Menschen, so nicht gar böse sind, erzürnet hat, können sie denken, so sie bitten und die Sünde bekennen, leichtlich verzeihen, so ist die Sünde alsobald todt; wenn man aber die Sünde noch will zudecken und schmücken, so werden die Herzen dadurch immer mehr zum Zorn gereizt. Dies ist die Art und Natur des Gewissens; wiewohl es aber die Sünde fühlt, und das Herz erschreckt und geschüttelt wird bis zur Verdammniß, bekennt es doch die Sünde noch nicht.

213. Derhalben soll ein Jeder lernen im Namen des Herrn, wenn er gesündigt hat, daß er es bekenne; alsdann wird die Sünde nicht allein vergeben sein, sondern alle Menschen werden Mitleiden mit ihm haben und sich seinen Fall jammern lassen. Aber da ist niemand, der das thun will, denn nur ingemein dahin. Insbesondere aber die Schuld zu bekennen ist dieser verderbten Natur Art oder Eigenschaft nicht, wiewohl sie bekennt, vor GOtt schuldig zu sein. Denn solches zeiget das Exempel dieser Brüder fein an: Joseph spielt mit ihnen und wendet ihre Herzen auf mancherlei Weise hin und her; aber sie denken nicht, daß dies Joseph sein sollte; denn sie meinen, er sei vorlängst gestorben. Sie fühlen zwar wohl, daß GOtt ihrer Sünde gedenkt; aber das Bekenntniß ist noch nicht rein. Darum bringt Joseph also auf sie und hält immer an, daß er es gleichsam mit Gewalt heraus bringe und das Bekenntniß der Sünden von ihnen erzwingen möge, dadurch sie zur Vergebung der Sünden kommen mögen.

214. Wie ist aber das so ein großer Jammer und Thorheit an uns, daß wir also fliehen vor der Thür, dadurch wir zur Seligkeit eingehen sollten. Wenn wir aber weise und ohne diesen Schaden und Verderben der Erbsünde wären, und das Gewissen in uns nicht so verstört und ungehorsam wäre: alsdann wären die Sünden an den Göttlichen wahrhaftig nichts Anderes, denn des Heiligen Geistes Besserungen, Unterweisungen und Wirkungen zum Guten, wie Paulus sagt Röm. 8, 28.: „Denen, die Gott lieben, dienen alle Dinge zum Besten“, auch die Sünden. Warum aber das? Ei darum, daß wo wir in Sünden gefallen, werden wir dadurch gedemüthigt und unterwiesen, werden gewarnet, daß wir uns besser lernen vor Sünden hüten, und werden auch dadurch verursacht, GOtt zu danken, wenn wir unsere Sünde bekennen und Buße thun. So gar vielfältige Frucht ist in der Sünde, wenn Einer derselbigen geistlich brauchen könnte. Es ist wohl der Tod und Verdammniß im rechten Gebrauch des Gesetzes, aber man muß den geistlichen Gebrauch oder Nutzen auch dazu nehmen. Das ist so viel gesagt: Man muß die Sünde erkennen, das Herz öffnen, bekennen, Hülfe suchen; alsdann wird uns die Sünde nützlich sein. Und je mehr du gedemüthigt wirst, je eher du auch Gnade und Vergebung erlangen wirst.

215. David schwieg auch eine Zeitlang stille; wie er von sich selbst sagt Ps. 32, 3.: „Ich wollte es verschweigen.“ Ich hatte den frommen Uria todt geschlagen, sein Weib geschändet; ich schmierte und wollte nicht hinan; die Sünde bedeckte und verhüllte ich, konnte sie fein schmücken. Denn die Erbsünde will ihre Ungestalt und Schande nicht gern schauen lassen, und der Teufel kann auch nicht leiden, daß seine Schande und Gestank an Tag komme und geoffenbart werde; sondern will jederzeit für Unschuld und Gerechtigkeit gehalten haben. Das kann unser Herr GOtt nicht leiden. Darum setzt David im 32. Psalm V. 3. das dazu: „Denn da ichs wollte verschweigen, verschmachteten meine Gebeine“ etc., ich wäre schier zu Grunde gegangen. Denn die Pein und Erinnyen des bösen Gewissens wird nicht mit Stillschweigen, Verhehlen oder Verdecken bezwungen. Derhalben, da Nathan zu David kommt, und sagt 2 Sam. 12, 7.: „Du bist der Mann“, der solches gethan hat, da wird er erschreckt und getödtet; und sagt doch durch den geistlichen Gebrauch der Sünde: „Ich habe gesündiget.“ Da der Prophet das Wort hört, daß David keine Scheu trägt, sich selbst ungestraft zu machen, zu schänden und zu verdammen, da hat er bald zu ihm gesagt Ps. 13.: „Der Herr hat deine Sünde hinweggenommen“ und dir dieselbe vergeben. Wo er sich aber hätte entschuldigen und schmücken wollen, wie Saul gethan, und seine Sünde zudecken, und für seinen Sünder gehalten werden, sondern gerecht und unschuldig sein wollen, so wäre er ewiglich verloren und verdammt worden. Darum ist er ja ein frommer GOtt, wenn wir die Erbsünde recht brauchen können; und wenn Einer gefallen wäre, daß er sich alsdann bald zum geistlichen Gebrauch der Sünde wendete, so ist er darum ein reicher GOtt mit seiner Gnade und Vergebung.

\quellseite{393}

216. Darum sollen wir lernen, daß wir des Unglücks und der Sünden uns zu Gutem gebrauchen können; so werden wir erfahren, wie gnädig und freundlich der Herr ist. Bist du gefallen, und kannst von Herzen sagen: Ich habe gesündiget; so kann er auch sagen: Ich habe dir es vergeben. Dagegen aber, wo du die Sünde wirst vertheidigen und verhehlen, so wird er sagen, er könne dir es auch nicht vergeben.

217. Darum ist es besser, man folge den Worten Davids, und sage mit ihm Ps. 51, 6.: „An dir allein hab ich gesündiget, auf daß du Recht behaltest in deinen Worten“ etc. Denn das ist auch vor der Welt löblich und man kann damit leichtlich Gnade erlangen. Denn solches fordert auch GOtt von uns. Laß mich, sagt er, Recht behalten in meinen Worten, und rein bleiben, wenn ich gerichtet werde; alsdann will ich dir verzeihen; wo nicht, so will ich doch Recht behalten und du sollst verloren werden. Wir haben GOttes Zorn, den Tod und ewige Verdammniß verdient; aber GOtt ist barmherzig, langmüthig und vergibt gern die Missethat, wenn wir nur unsere Unreinigkeit und Unwürdigkeit erkennen und bekennen könnten, wenn wir nur könnten das Maul aufthun.

218. Und es ist wahrlich ein großer Jammer, daß unser Herz durch die Erbsünde so gar ausgeschüttet und alles Vertrauen auf unsere Würdigkeit und Verdienst hinweggenommen ist. Denn das Herz wird bleich und zittert, wenn es nur einen Blitz sieht; und ist doch gleichwohl noch so stolz und will sich vor GOtt nicht demüthigen. Derhalben danke GOtt, wenn du dieses Licht hast, daß GOtt deinen Sünde vergeben wolle, welche sie erkennen. Denn er fordert keine Rache oder Genugthuung dafür, wenn du sie ernstlich und von Herzen bekannt hast, denn allein so viel, als vonnöthen ist Andern zum Exempel und Besserung; wie der 32. Psalm V. 5. sagt: „Ich sprach: Ich will dem Herrn meine Uebertretung bekennen. Da vergabst du mir die Missethat meiner Sünde.“

219. Aber die Söhne Jakobs fliehen noch immer vor diesem Bekenntnisse. Es ist der Tod Josephs in ihren Herzen und nimmt ihnen allen Trost. Und sie fühlen, daß was ihnen für Unglück widerfährt, von dem zornigen GOtt, der ein Rächer aller Uebelthaten ist, herkomme. Denn das Gesetz ist inwendig im Herzen, welches die Leute erschreckt und GOttes Gesetz ist. Darum ist es gewiß, daß alles Schrecken und Zittern des Gewissens durch Mißwirkung GOttes geschieht.

220. Derhalben kannst du das Gesetz nicht abschütteln, sondern es erschüttert dir dein Herz. Denn es ist GOttes ewiges und unwandelbares Gericht, und du wirst sein Anklagen und Stürmen nicht leichtlich ertragen können, wie Judas und Andere mehr, so dadurch erschrocken, verloren worden sind, da sie nicht haben überwinden können die greuliche Pein im Herzen, und die höllische Furie des bösen Gewissens, die das Herz immer treibt und jagt durch GOttes Gesetz, welches allmächtig ist, wie er selbst auch, der es gegeben und den Menschen in ihr Herz eingeprägt hat. Und du wirst von denselben Schrecken nimmermehr erlöset werden, denn allein durch die Hülfe, so GOtt selbst gezeigt hat im 32. Psalm V. 5., da er sagt: Bekenne deine Uebertretung wider dich selbst, so will ich dir es vergeben. Sage du erst, Ps. 51, 6.: „An dir allein hab ich gesündiget, auf daß du Recht behaltest in deinen Worten“; so will ich tilgen alle deine Missethat.

221. Dies ist das Heil- und Friede aller Heiligen vom Anfang der Welt her, Adams, Abrahams, Davids, in welchen allen die Erbsünde gewesen, so sie erschreckt und verdammt hat. Aber durch diese Hülfe, welche in der Erkenntniß der Sünden besteht, sind sie errettet und erlöset worden von den Schrecken des Gesetzes.

222. Auf solche Weise fangen diese Brüder auch an, zu fühlen den Stachel des Gewissens und Zorns GOttes, welcher sie straft. Darum sagen sie: „Warum hat uns GOtt das gethan?“ Was hat GOtt im Sinne? Das sind Worte eines erschütterten Herzens, welchem all sein Trost und Zuversicht genommen und die Sicherheit gar niedergelegt ist; wie denn die Worte derer, so wegen der begangenen Sünde erschrocken und betrübt sind, zu lauten pflegen: Ach GOtt, was habe ich gethan? Also sehen diese Brüder auch, daß dies, so ihnen da widerfahren, von Joseph nicht müsse hergekommen sein; denn sie wußten, daß ihm alles Geld bezahlt war, und daß sie ihm nichts gestohlen oder entführt hatten. Darum denken sie, es müsse etwa von GOtt Einer heimlich dazu bestellt sein, ein Mensch oder Geist, der sie vor dem Fürsten in Egypten des Diebstahls beschuldigen und anklagen sollte. Darum werden sie gedacht haben: Was wollen wir thun? Jetzt werden Etliche daher kommen, vom Fürsten gesandt, die uns angreifen und in das Gefängniß ziehen sollen. Also sehen sie, daß sie weder in Egypten, noch im Lande Canaan, da sie gewohnt haben, werden sicher genug sein können. Und dieser Schrecken wird gesteigert nicht sowohl durch die Gefahr des Verdachts, sondern durch die Erinnerung an ihre vorige Sünde und andern Uebelthaten, die sie aus Furcht erdichtet oder wirklich begangen hatten und daran sie damals gedacht.

223. Denn das soll man fleißig merken, daß ein böses Gewissen, welches schuldig ist, wenn es gerichtet und gestraft wird, nicht allein der rechten Sünden halben, die es weiß, daß es sie wahrhaftig begangen hat, Pein oder Marter und des Teufels Beißen fühlt, sondern es dichtet sich auch etwa aus Einer Sünde viel andere unzählige Sünden mehr. Desgleichen pflegt es auch unzähligen Argwohn bei sich zu fassen, damit der Schrecken immer gemehret wird, und ist doch dieses alles lauter erdichtetes Ding.

224. Das geschieht aber darum, dieweil die Schrift sagt 3 Mos. 26, 36.: „Ein rausche[n]d Blatt soll sie jagen.“ Ja, ein dürres Blatt soll das böse Gewissen also erschrecken, daß es meine, es sei eitel Donner und Blitz. Das ist wahrlich ein großer Jammer, daß man vor einem Blatte, wenn man es nur leichtlich rauschen hört, nicht anders zittern und erschrecken soll, als vor Donner und Blitz. Ich habe wahrlich oftmals mit diesem scheußlichen Thiere nicht ohne großen Kampf gestritten, und 5 Mos. 28, 65—67. wird noch mehr dazu gesetzt, nämlich: „Der Herr wird dir ein bebendes Herz geben, und verschmachtete Augen, und verzweifelte Seele, daß dein Leben wird vor dir schweben. Nacht und Tag wirst du dich fürchten, und deines Lebens nicht sicher sein. Des Morgens wirst du sagen: Ach, daß ich den Abend erleben möchte! Des Abends wirst du sagen: Ach, daß ich den Morgen erleben möchte! vor Furcht deines Herzens, die dich schrecken wird.“

225. Auf solche Weise beschreibt der Heilige Geist die erschreckliche Pein und Grimmen eines bösen Gewissens. Daher es kommt, daß unser ganzes Leben immer wie an einem Faden hängt, steht im Zweifel und ist ganz ungewiß, ist alle Augenblicke voller Angst und Schrecken, fürchtet sich vor Abend und Morgen, ja, endlich auch vor allen Stunden.

226. Also sind diese Brüder Josephs in großer Angst, nicht allein der Sünde halben, deren sie sich schuldig gewußt, wiewohl sie dieselbe mit Fleiß zugedeckt und verhehlt haben; sondern sie fürchten sich auch vor allen ihren eigenen Worten und Werken und dichten sich in allen Werken und Gedanken eine neue Schuld. Und in so großem Schrecken und Zagen sind sie doch gleichwohl noch stumm und bekennen ihre Sünde nicht.

227. Eben also that Adam im Paradiese auch: da er beinahe gar erstorben war aus Furcht vor der Stimme und Zorne GOttes, will er doch das Maul noch nicht aufthun und kann nicht sagen: Ich habe gesündiget. Und eben solche Thorheit und verkehrtes Herz folgt in allen Adamskindern auf die Erbsünde, wie bald im Anfange an unserm ersten Vater solche Thorheit und Halsstarrigkeit gespüret worden, bis die geistliche Furcht der Sünden dazu kommt durch das Wort Christi Marc. 7, 34.: „Hephatha, thue dich auf!“ Inzwischen bleibt das Band der Zunge, damit der Teufel den Mund verstopft, daß wiewohl der Mensch im Gewissen sehr erschreckt wird, er doch gleichwohl keine Hülfe sucht. Denn der Teufel läßt den Sünder nicht von Herzen sagen: Herr GOtt, an dir habe ich gesündiget, vergib es mit, es ist mir herzlich leid, was ich gethan habe!

\quellseite{394}

228. Derhalben soll man diese Beschreibung der Sünden fleißig merken, welche diese Brüder dermaßen erschreckt, daß sie anfingen zu schreien: „Warum hat uns GOtt das gethan?“ als ob sie nicht wüßten, warum es GOtt gethan hätte. Lieben Herren, euer Herz weiß wohl, daß ihr euren Vater, Bruder und Mutter getödtet habt. Denn wo Rahel noch am Leben gewesen, wäre sie doch ohne allen Zweifel vor großem Herzeleid gestorben. Den Vater aber hat GOtt wunderbarlich erhalten, daß er nicht selbst auch von Schmerzen und Leid gestorben ist. Davon solltet ihr wohl nichts wissen? Ja, wir schämen uns, daß wir es sollen bekennen, und die Sünde, so wir begangen, dem Vater offenbaren; denn diese verkehrte Schamhaftigkeit hat ihnen im Wege gelegen. Darum, so traget nun diese Kreuz, daß ihr seufzt und schreiet: „Warum hat uns GOtt das gethan?“ Denn meinet ihr, daß euch dies also werde hingehen, daß ihr deshalb solltet ungestraft bleiben, oder daß GOtt eure Sünde entschuldigen oder loben werde? Das wird er nimmermehr thun. Fraget ihr darnach, warum euch GOtt dies thue? Wie, daß ihr nicht viel lieber sagt: Warum haben wir das GOtt gethan? Oder, warum bekennt ihr doch eure Sünde nicht? Dieweil es aber euch noch nicht gelüstet, dies Bekenntniß zu thun, wohlan, so fürchtet euch noch etwas länger, zittert und erschreckt alle Augenblicke, wenn ihr auch nur ein dürres Blatt rauschen hört, und saget des Morgens: Ach, daß ich den Abend erleben möchte! und des Abends saget: Ach, daß ich den Morgen erleben möchte! etc., 5 Mos. 28, 67.

229. Auf diese Weise lehrt uns das Exempel der heiligen Väter, welch ein erschrecklich Ding es um ein böses Gewissen sei. Und ist ein sehr trefflicher Text, welcher werth ist, daß er deutlich gehandelt und ausgelegt werden möchte, dazu ich mich allerdings nicht für fähig halte. Dieweil sie aber noch nicht bekennen, daß sie gesündigt haben, auch GOtt nicht um Gnade, noch den Vater oder Joseph um Verzeihung bitten, so hört auch die Plage und Anfechtung noch nicht auf, bis ihnen das Bekenntniß und die Erkenntniß der Sünden mit Gewalt abgedrungen werde. Joseph muß noch besser auf den Sack dreschen.

\section[Sechster Theil]{Sechster Theil. Wie die Söhne Jakobs nach Hause kommen, und ihrem Vater erzählen, wie es ihnen gegangen; wie sie ihr Getreide ausschütten und ihr Geld wiederfinden}

\begin{quote}
\textbf{V. 29—34. Da sie nun heim kamen zu ihrem Vater Jakob ins Land Canaan; sagten sie ihm alles, was ihnen begegnet war, und sprachen: Der Mann, der im Lande Herr ist, redete hart mit uns, und hielt uns für Kundschafter des Landes. Und da wir ihm antworteten: Wir sind redlich, und nie Kundschafter gewesen; sondern zwölf Brüder, unsers Vaters Söhne; einer ist nicht mehr vorhanden, und der jüngste ist noch bei unserm Vater im Lande Canaan; sprach der Herr im Lande zu uns: Daran will ich merken, ob ihr redlich seid; einen eurer Brüder lasset bei mir, und nehmet die Nothdurft für euer Haus, und ziehet hin; und bringet euren jüngsten Bruder zu mir, so merkte ich, daß ihr nicht Kundschafter, sondern redlich seid; so will ich euch auch euren Bruder geben, und möget im Lande werben.}
\end{quote}

230. Dies gehört auch noch zu der vorigen Beschreibung, darin wir gesagt, daß Mose ganz kein abgemalt habe die Art und Natur der Erbsünde, daß sie ja nicht will Sünde sein. Und hat zwar der Teufel es bei den Papisten dahin gebracht, nämlich, daß der Schaden der Erbsünde ganz ausgetilgt ist, und nicht mehr für verdammlich gehalten wird. Nun bringt aber gemächlich und nothwendig die Unwissenheit der Sünden auch die Unwissenheit GOttes, Christi, des Heiligen Geistes und aller Dinge mit sich, das ist, daß man GOtt, Christum und den Heiligen Geist nicht kennt, und von keinem Dinge etwas weiß, wo man die Sünde nicht zuvor kennen oder wissen will. Denn niemand darf denken, daß er werde ein Theologe werden, oder daß er die heilige Schrift recht lesen oder hören wolle, der den Schaden der Erbsünde so klein macht und geringe achtet, oder aber denselben noch nicht recht versteht. Ja, kein Mensch ist, der die Kraft der Erbsünde genugsam bedenken oder verstehen könne. Denn wir sollen nicht denken, daß es so ein gering oder leicht Ding sei, darum GOtt seinen Sohn gesandt hat, daß er ein Opfer, uns zu erlösen, geworden ist.

231. Die Papisten lehren, es sei die Erbsünde in der Taufe gar weggenommen, und bleibe nur ein kleiner Zunder und Schwachheit des Fleisches, oder eine Neigung zu sündigen. Derhalben verstehen sie auch nicht, was Gnade oder Vergebung der Sünde, auch nicht, was der Heilige Geist und GOtt sei.

232. In diesem Capitel aber ist eine Beschreibung dieses greulichen Schadens, oder Erbsünde, deren man billig immer gedenken soll, nämlich, daß es ein solch Gift ist, so tief in unserer Seele und Leibe steckt, also daß sie auch für Gerechtigkeit will gehalten werden; wie David anzeigt, da er sagt im 51. Psalm V. 6.: „Auf daß du Recht behaltest in deinen Worten.“ Denn warum redet er also? Ist denn GOtt nicht rein und gerecht? Antwort: Er ist ja gerecht; aber bei uns wird er dafür nicht gehalten, daß er also sei: wir sind vielmehr gerecht, GOtt ist aber ungerecht. So gar fein kann sich die Sünde schmücken und ihre Unreinigkeit zudecken. Wie das an diesen Brüdern zu sehen ist, welche traun nicht erschrocken oder verdammt sein wollen. Wiewohl sie sagen: Wir leiden dies billig, und: Warum hat uns GOtt das gethan? bedeuten sie doch immer ihre Sünde und verhehlen dieselbe; und alhier auch vor ihrem Vater wollen sie mit nichten für die gehalten werden, als die da gesündigt hätten.

233. Nun kann aber GOtt nicht rein oder gerecht sein, wo wir nicht unrein und ungerecht sind, das ist, GOtt wird nicht als gerecht erkannt, wenn wir nicht unsere Ungerechtigkeit bekennen. Nicht, daß ich jetzt erst müßte ungerecht werden; denn ich bin vorhin unrein: sondern das ist die Frage: Ob ich auch im Geiste wolle unrein werden, das ist, meine Sünde erkennen und bekennen? Denn wie anfänglich die Erbsünde hat wollen zugedeckt und geschmückt sein; also folgen auch die wirklichen Sünden hernach, wie das Kind seinem Vater und Mutter nachfolgt, und halten sich streng nach dem Spruche, da man sagt: Si fecisti, nega: Hast du es gethan, leugne. Denn wer vor den Menschen seine Sünde leugnet oder verhehlt, der ist schon gerechtfertigt. Aber es ist in Wahrheit der Teufel selbst, der sich also in einen Engel des Lichts verstellt. Es will alles rein sein, und was GOtt thut, das soll unrein sein. Zuletzt muß aber doch dieser Spruch noch fest stehen bleiben und bekräftigt werden, nämlich, daß wir sagen: Herr GOtt, du bist gerecht, du bist rein; ich aber bin ungerecht und unrein; und müssen alle mit David Psalm 51, 6. sagen: „Auf daß du Recht behaltest in deinen Worten.“

234. Aber davon hat St. Paulus Röm. 3, 8. herrlich gehandelt. Und daraus erhellt auch, wie ungereimt und erschrecklich diese Lehre in den Ohren der fleischlichen Menschen sei. Denn da sagten Etliche: Die Apostel lehreten, man sollte nur frei sündigen und Uebels thun, auf daß Gutes daraus komme. Wenn unsere Ungerechtigkeit GOttes Gerechtigkeit preiset, so will ich getrost drauf sündigen. Dies ist aber nicht eine philosophische, sondern eine geistliche und göttliche Lehre. Denn du wirst alsdann ungerecht werden, wenn du fühlest, erkennest und bekennest wirst, daß du wahrhaftig ein Sünder bist, und sagen wirst: Herr, ich habe gesündiget, vergib es mir!

235. Diese übernatürliche Kunst wollen wir haben. Denn siehe doch diese Brüder an, wie höchlich sie ihren alten Vater betrübt und beleidigt haben, und doch dahin nicht gebracht werden können, daß sie zusammen ihre Sünde bekenneten. Sie ermahnen einander nicht, daß sie zum Vater treten und anzeigen wollten, daß sie Brüdermörder seien; mit welchem Bekenntniß beide GOtt und ihr Vater versöhnt worden wären.

236. Darum ist es vonnöthen, daß du mit Wissen und Wollen sagest: Ich bin ein armer Sünder und ungerecht. Warum aber das? Ei, „auf daß du Recht behaltest in deinen Worten“, Psalm 51, 6. Wirf nur die Augen nieder auf die Erde, und sage: GOtt sei mir gnädig; denn ich bin in Sünden geboren, und meine ganze Natur ist verderbt, ist voller Blindheit, böser Lust, Unwissenheit ihrer selbst und auch der Sünden.

237. Und wo ihr in der Theologie und heiligen Schrift nützlich studiren wollt, daß ihr in der Schrift, welche verschlossen und versiegelt ist, nicht anlaufet, so lernt ja vor allen Dingen die Sünde recht verstehen. Und die Epistel an die Römer soll euch sein wie eine Thür und Schlüssel zu der heiligen Schrift; sonst werdet ihr nimmermehr zum rechten Verstande und Erkenntniß der Schrift kommen. Alsdann aber werdet ihr verstehen, welch ein böses Ding die Sünde sei, um welcher willen Gott seinen Sohn gesandt hat, daß wir davon durch das Blut und unaussprechliche Opfer des Sohnes Gottes gereinigt würden. Denn darum muß ja der Schade selbst auch unaussprechlich groß sein.

\quellseite{395}

238. Wie aber diese Brüder halsstarrig sind, und ihre Sünde noch immer zu bedecken fortfahren: also geht diese erschreckliche Plage und Jammer durch diese ganze Welt, daß die Sünde sich ja nicht will lassen aufdecken, will das nicht sein, was sie doch in Wahrheit ist; denn sie will rein und schön sein, wider ihre Art und Natur. Es will schön sein. Daselbst sehen wir auch in Hausregimente und weltlicher Polizei also geschehen. Denn alles, was daselbst gesündigt wird, wird nur auf den „Niemand“ gelegt: der muß Alles gethan haben. Denn Niemand kann leichtlich ertragen, daß er gestraft und getadelt wird um der Sünden willen. Sie sind sonst alle fromm, es will niemand unrecht gethan haben. Also läßt sich in den wirklichen Sünden die böse verderbte Natur auch sehen. Wer hat denn nun gesündigt? Wer wird bezahlen müssen? Allein der Sohn Gottes ist Sünder, sonst niemand. Er allein trägt die Sünde, und sagt: Sei mir gnädig, denn an dir habe ich gesündigt. Item im 40. Psalm V. 13. sagt er: „Meiner Sünden ist mehr denn Haare auf meinem Haupt.“ Der erkennt und bekennt die Sünde frey heraus, und ist ihm damit Ernst nicht zwar für seine Person, die da heilig und gerecht ist, sondern als ein Patron und Fürsprecher vor dem Richter, und bezahlt auch für uns.

239. Darum lernet ja dies Stück von der Erbsünde mit Fleiß, und disputirt davon nicht, warum Gott dieselbe habe zugelassen, darnach die vorwitzigen Köpfe zu fragen pflegen: sondern fraget viel lieber darnach, wie wir von diesem Schaden errettet und erlöset werden mögen, und daß wir wissen, daß Gott mit uns redet, auf daß er uns erwecken möge, denselbigen Schaden der Erbsünde zu erkennen. Denn wo der erkannt wird, so sagt er: „Sei getrost, mein Sohn, deine Sünden sind dir vergeben“, Matth. 9, 2.; denn ich habe meinen lieben Sohn für dich gegeben, zu einem unschuldigen Lämmlein vom Anfange. Wirst du deine Sünden und Missethat erkennen und bekennen, alsdann wird mein Sohn die Versöhnung für deine Sünde sein. Er wird deine Heiligung, Erlösung, Gerechtigkeit und Weisheit sein, 1 Cor. 1, 30.

240. Da ist alsdann die heilige Schrift geöffnet, und du wirst sehen, daß sie nichts Anderes lehrt, denn daß Gott uns seinen Sohn vorhält, der gekreuzigt, auferstanden und für unsere Sünde dahin gegeben ist. Solches wird dir aber nicht zu Herzen gehen wollen, und wirst auch nicht verstehen, was Christus sei, es sei denn, daß dich die heilige Schrift zuvor erinnere und lehre, was du vor Gott seiest. Dasselbe wirst du aber lernen aus 1 Mos. 3, 19. Denn da wird gelehrt, daß dem ganzen menschlichen Geschlechte die Strafe des Todes auferlegt sei um der Erbsünde willen.

241. Wie groß aber die Bosheit derselbigen Sünde sei und wie sie so gar verheeret sei, das zeigen die Früchte an. Denn wie ist doch bald in der Jugend bei den Kindern so ein großes Widerstreben und Ungehorsam! Welches Müthen und Brennen der Lüste, des Hasses, des Geizes, des Neides in der Jugend und im ganzen Leben! Wir werden ja wohl geboren ohne wirkliche Sünden; aber hernach bricht mit der Zeit mit Gewalt eine unendliche Menge der Sünden hervor. Lieber, woher kommen doch solche greuliche Ungeheuer?

242. Die Philosophen Sokrates und Cicero wissen davon nichts; aber die Schrift sagt: Wir werden geboren als Kinder des Zorns, und die dem Verderben der Erbsünde gar unterworfen sind; und je weiser, gelehrter und klüger wir sind, je mehr Hoffahrt und andere Sünden sich an uns finden. Denn die Natur ist gar verderbt. Die Papisten aber weisen uns vor, daß sie in der Taufe wiederum soll geheilt sein und daß auch darin die Sünden vergeben seien. Das ist wohl recht geredet; aber es ist die Sünde noch nicht gar hinweg genommen, wie Augustin sagt. Dieweil sie aber das nicht wissen, können sie nimmermehr zum rechten Lichte oder Verstande der Schrift, des Reichs Christi und der ganzen Theologie kommen.

243. Es sind ja wohl die Wunden des Mannes, Luc. 10, 34., der halb todt war, verbunden und ist Oel und Wein darein gegossen; es hat die Gabe des Heiligen Geistes angefangen: und sind doch die Wunden noch tödtlich. Er ist angenommen, daß er vom Arzte geheilt werden soll: aber er ist noch nicht gar gesund und wiederhergestellt worden. Wenn du nun sagen wolltest, daß keine Wunde und auch keine Gefahr da wäre, lieber, so versuche doch, ob der, so halb todt ist, gehen und stehen, arbeiten, und sonst thun könne, was ein gesunder Mensch zu thun pflegt. Er wird ja getragen von dem Thiere, darauf er ist gesetzt worden, arbeitet nicht, geht oder steht auch nicht. Also sind wir durch die Taufe auch aufgenommen worden auf unsers Herrn Gottes Thier, das ist, auf das alltheuerste Opfer, so für uns geschehen ist, oder auf die Menschheit Christi. Dadurch werden wir getragen. Und sind zwar einmal angenommen, werden aber noch immer von Tage zu Tage verbunden und geheilet.

244. Wir wollen aber nun wiederum auf das Exempel dieser Väter kommen. Es sind rechte verzweifelte Buben. Sie haben ihren Bruder sammt den Eltern getödtet, soviel an ihnen gewesen; und verharren doch noch immer in Unbußfertigkeit, auch so gar gröblich, daß da sie schon überwunden und zu Schanden gemacht sind, sie doch noch nicht frei bekennen wollen, daß sie gesündigt haben. Wenn sie bekannt hätten, daß sie Joseph verkauft, so wären sie viel gnädiger und freundlicher empfangen worden; aber die Sünde ist bei ihnen noch heilig, gerecht und stolz, und ist da niemand unter ihnen, der gesündigt hätte. Darum können sie sich so wunderlicher Weise schmücken, und schwächen die Erzählung, wie man mit ihnen in Egypten umgegangen sei, und werden davon ohne Zweifel unterwegs geratschlagt haben, wie sie dasselbe dem Vater aufs allgelindeste erzählen und des Gefängnisses nicht gedenken möchten, auf daß er nicht verursacht würde, genauer zu fragen, was die Ursache sei, daß man mit den fremden Leuten so greulich umgegangen wäre: und können doch nicht gar verschweigen, daß Simeon ausgeblieben und nicht da wäre, da sie haben wollen, daß er Benjamin mit ihnen soll hinabziehen lassen.

245. Also erzählen sie nun, was ihnen begegnet war, und sagen: „Der Mann, der im Lande Herr ist, redete hart mit uns“ etc. Wir sind hart und schmerzlich angefochten worden; wir heilige und unschuldige Männer haben wunderbare Bitterkeit und Unfreundlichkeit erfahren müssen. Warum aber das? Zwar nicht, daß wir es verschuldet hätten, sondern du, Vater, hast vielleicht etwas mit einer Sünde dies Kreuz verdient. „Und er hielt uns“, sagen sie weiter, „für Kundschafter des Landes.“

246. Sie waren drei Tage im Gefängnisse verwahrt worden und wollen doch jetzt keines Gefängnisses gedenken. Denn es würde sonst der Vater gefragt haben: Warum sie denn in das Gefängnis geworfen wären? Darum haben sie aus sonderlichem Bedenken alles, was ihnen begegnet, auf das allergelindeste erzählt, damit der Vater durch keine ungewöhnliche That oder Handlung nicht erzürnt und bewogen würde, sie um die Ursache solcher ungewöhnlichen Handlung zu fragen. Also thut immer die Sünde, wie ihre Art ist: sie will ja rein, züchtig, schön und heilig sein; dagegen aber beschuldigt sie uns und verdammt Gott, daß er ungerecht sei und unbillig handle.

247. Wir wollen es aber noch etwas gelinder deuten, nämlich, daß sie des Gefängnisses darum geschwiegen haben, auf daß sie damit den Vater nicht betrübten oder erschreckten. Da aber schwanken sie, ungewiß darüber, wie sie ihm das beibringen sollten, was sie wünschten, nämlich, den Benjamin mit sich zu nehmen. Um daher diese Bitte einzuleiten, versuchen sie mit großem Eifer ihrem Vater einzureden, daß ihnen dieses alles als solchen, die es nicht verdient und die unschuldig wären, begegnet sei; auf daß ja die Erbsünde nicht etwa anlaufe.

248. Warum nennt aber Mose Joseph adone, das ist, einen Herrn des Landes, und redet in der Mehrzahl als von vielen? etc. Denn mit solchen Sprüchen suchen die Juden die Mehrheit der Personen in der Gottheit, welche wir daraus zu beweisen pflegen, daß Gott oft in der Mehrzahl elohim genannt wird, lächerlich zu machen. Welcher Art und Weise in der Mehrzahl von sich zu reden allein die deutschen Fürsten und der Papst folgen. Darum sollen wir wissen, daß die hebräische Sprache oft ein Hauptwort in der Mehrzahl mit einem Zeitwort in der Einzahl verbindet; wie die Griechen zu Wörtern sächlichen Geschlechts in der Mehrzahl oft das Verbum „ist“ (in der Einzahl) setzen. Also an diesem Orte: Der Herr oder die Herren des Landes hat geredet. Droben, Cap. 39, V. 16., haben wir eben dieselbige Weise zu reden auch gehabt, da die Ehebrecherin, Potiphars Weib, das Kleid neben sich legte, bis auf die Zeit, da „seine Herren“ kommen würde, so doch sollte da nur von dem einigen Potiphar, Josephs Herrn, reden. An demselben Orte ist auch ein Fürwort in der Einzahl zu einem Hauptwort in der Mehrzahl gesetzt. Es ist auch zwar nicht eine seine Weise und Art zu reden, und die sich so gar wohl schickte; aber die hebräische Sprache hat ihre eigene Art und Weise, davon man so genau nicht kann Ursache geben; wie in andern Sprachen auch nicht können Ursachen angezeigt werden von allerlei Weise zu reden. Es ist genug, daß wir wissen, daß dies derselbigen Sprache Art und Weise sei. Wie die Deutschen den sächlichen Artikel ohne Grund zu einem Hauptwort weiblichen Geschlechts setzen, wenn sie sagen: Das ist die Jungfrau; item: Das ist die Stadt. Denn das sind Idiotismen und Eigentümlichkeiten in dieser Sprache.

\quellseite{396}

249. Sie fahren aber nun fort ihre Antwort zu erzählen, wie sie sich vor dem Manne entschuldigten, nämlich, daß sie redlich und aufrichtig wären, daß sie nichts Böses vorhätten. Und dies alles gehört noch immer dazu, daß sie sich damit geschmückt haben, auf daß der Vater nicht dächte oder den Argwohn faßte, daß sie solche Strafe des Gefängnisses oder daß sie so unfreundlich empfangen worden, selbst verursacht oder verschuldet hätten, und daß er deshalb vornehmen möchte, weiter darnach zu forschen. Darum sagen sie, dies sei ihnen begegnet, ohne ihre Schuld, und wie sie ihre Unschuld auch öffentlich vor dem Herrn des Landes angezeigt haben. Sie erzählen aber nur die Stücke, welche ihnen gedient haben, die Sünde zu verkleinern; ja, damit sie auch einen solchen Namen erlangen möchten, daß sie für heilig und gerecht gehalten würden.

250. Dieweil aber Simeon nicht gegenwärtig war, den sie Joseph gefänglich bei sich behalten, erzählen sie solches auch fürzlich, warum sie denselben in Egypten hinter sich gelassen. Dieweil er uns, sagen sie, für Kundschafter hielt und mit unserer Entschuldigung nicht zufrieden war, befahl er, daß Einer unter den Brüdern so lange bei ihm bleiben sollte, bis wir unsern jüngsten Bruder zu ihm bringen würden. Das lassen sie aber weg, daß sie drei Tage im Gefängniß gelegen hatten, wie auch dieses, daß der Herr des Landes bei dem Leben Pharaos geschworen, er wollte sie nicht ziehen lassen. Und setzen das auch nicht dazu, daß Simeon auf Befehl des Herrn im Lande gebunden und in Gewahrsam genommen war; denn das alles wäre etwas hart und gefährlich gewesen. Darum haben sie dasselbe nicht also ausdrücken wollen, sondern bedecken und lindern es, wie sie nur können und mögen.

251. Derhalben siehst du wohl, wie sehr listig die Erbsünde ist, wie sie ihre Mißgestalt so fein künstlich schmücken kann. Und solche seine Redner sind alle Menschen, ihre Sünde zu verkleinern und zu entschuldigen. Daher kommt es, daß man vor Gericht und in der Polizei so viel Zeugen, Richter, und was sonst des Dinges im Rechte mehr ist, bedarf, damit die Leute überzeugt und getrieben werden, ihre Sünde zu bekennen. Und zeigen solche Dinge nichts Anderes an, denn die große unüberwindliche Halsstarrigkeit der Sünde. Denn wo dieselbe nicht also zugedeckt und geschmückt würde, bedürfte man keine Gerichte, Gesetze oder Rechte, keine Stockmeister, oder Strafe: aber dieweil die Sünde heilig und gerecht sein will, so muß solche Ehre und Unschuld durch das Schwert, Galgen und Feuer etc. aufgedeckt und geoffenbart werden.

252. Auf diese Weise gebrauchen die Söhne Jakobs ein solches Wort, welches sehr gelinde ist, als hätte der Herr des Landes aus freundlichem und gutem Herzen darum gebeten, daß Einer von den Brüdern möchte bei ihm bleiben, welcher auch kein friedsam und sicher ohne alle Furcht und Gefahr daselbst erhalten würde. Denn es sind Heuchler und Werkheilige, die nicht wohl leiden können, daß die Gerechtigkeit und Heiligkeit der Sünde zu Schaden gemacht werde. Es hat daselbst Gottes Gerechtigkeit und Wahrheit gar keine Statt; die Sünde ist allein gerecht und heilig. Sie sagen trauen nicht, daß Simeon gefangen und in den Kerker gelegt sei, sondern daß der Herr des Landes sehr freundlich mit ihm umgehe und sein wohl warten lasse, welcher ihn als einen Bürger bei sich behalten, damit sie wider die Egypter nichts Böses vornehmen möchten: er könnte aber leichtlich los werden, wenn sie nur den Benjamin zu ihm brächten.

253. Es hat aber Mose an diesem Orte den Grammatikern abermals zu schaffen gemacht, daß sie sich plagen mit dem hebräischen Wort raabon; nämlich: Warum er doch das Getreide also genannt habe, so es sonst ein andermal Theurung heiße; und lautet nach dem Hebräischen also: Nehmet die Theurung, nicht: Speise, eurer Häuser. Was ist doch das für eine Grammatik, daß er Getreide „Theurung“ heißt? Es lassen es sich die Hebräer und sonderlich die Rabbinen sehr sauer werden, und ist daraus auch wohl zu sehen, daß sie die Erkenntniß beides der Dinge und auch der Grammatik nicht haben. Und kann die hebräische Grammatik niemand wiederum aufrichten, denn allein die Christen, welche dasjenige, davon die Schrift vornehmlich handelt, verstehen, das ist, die Christum, den Sohn Gottes, erkennen; denn wo derselbe recht erkannt ist, da wird das Andere alles sein klar und durchsichtig. Man kann aber dies leichtlich auslegen, wenn man auf die Art und Weise zu reden, dazu auch auf die Figuren Acht gibt. Denn es hat keine Sprache mehr Figuren denn die hebräische.

254. Nun ist es aber eine gemeine und gewöhnliche Figur nicht allein bei den Hebräern, sondern auch bei den Lateinern und Deutschen, die da heißt antiphrasis, als, wenn wir einen bösen verzweifelten Buben einen frommen Mann heißen. Und wird sonderlich im Hausregiment diese Figur viel gebraucht gegen Knechte und Kinder; als, wenn wir sagen: Lieber Junker, liebes Frömmichen etc., wenn wir ironischer Weise etwas befehlen oder zugeben und doch damit das Gegenteil meinen. Und was ist Ironie oder Sarkasmus Anderes denn antiphrasis, allein, daß diese in einem einzelnen Wort, jene aber in der ganzen Rede liegt.

255. Wie aber nun Mose droben das Brod „Brochen“ genannt, darum daß sie keine Messer gebraucht haben, sondern das Brod, das sie haben essen wollen, gebrochen worden ist; also sagt er an diesem Orte auch: Bringet raabon, das ist, Theurung, oder eigentlich das Gegenteil: Ueberfluß oder Fülle gegen den Hunger. Wie die Aerzte von Pestilenzen reden, nicht weil sie die Pestilenz bringen, sondern weil sie dieselbe vertreiben.

256. Also ist chesed ein sehr schön Wort; als, Matth. 9, 13. und Hos. 6, 6.: „Ich habe Wohlgefallen an Barmherzigkeit, und nicht am Opfer“; item im 4. Psalm V. 4.: „Der Herr führet seine Heiligen (chasid) wunderlich“; desgleichen auch Jes. 40, 6.: „Alles Fleisch ist Heu, und alle seine Güte ist wie eine Blume auf dem Felde.“ Da lautet es und heißt eigentlich so viel als, Ehre, Güte, Weisheit und was nur im Fleisch Gutes sein kann. Aber eben dasselbe Wort wird auch von Mose gebraucht in 3. Buch am 17. und auch im 20. Capitel, da es heißt, ein großes Laster oder Blutschande, und heißt doch eigentlich Barmherzigkeit, Wohltat oder sonst ein gut Werk; als, 1. Tim. 6, 2.: „Dieweil sie der Wohltat theilhaftig sind“, das ist, der göttlichen Gnade und Wohltat. Derhalben, da Mose sagt, daß die, so eine Blutschande (chesed) begangen, damit den Tod verschuldet haben, will er damit zu verstehen geben, daß sie kein gut Werk gethan, sondern ein großes Laster begangen haben, dafür sie den Tod haben leiden müssen.

257. Also wird Christus die Sünde genannt, er doch das Opfer für die Sünde ist, damit die Sünde bezahlt wird. Denn das ist der Gebrauch dieser Sprache, bei welchem es ganz und gar steht, wie er reden will und was die Kraft und Weise der Worte oder der Reden sein soll: Quem penes arbitrium est et vis et norma loquendi, wie der Poet sagt. Und man kann für solche Figuren keine Ursache anzeigen, welche der gemeine Gebrauch also eingeführt und mitgebracht hat; sondern die Regeln müssen demselben Gebrauch weichen, und man muß da aus der Grammatik Figuren suchen, damit man solche Reden, die etwa der Gebrauch gibt und mitbringt, erklären oder aber entschuldigen möge. Und muß in allen Sprachen zugelassen werden, daß sie diese Freiheit in solchen Reden haben mögen, so dem gemeinen Gebrauch etwas zuwider sind. Gleichergestalt müssen wir uns auch aus dieser Marter helfen.

258. Das aber nun weiter im Text folgt, da Joseph sagt: „Und ihr möget im Lande werben“, diese Worte haben auch eine sonderliche Kraft. Denn das Wort sachar heißt eigentlich so viel als, umhergehen, kaufen und hantieren. Und Mose hat droben, Cap. 23, 16., dasselbe Wort auch gebraucht, da er sagt von Abraham, daß er den Acker zum Begräbniß gekauft, und vierhundert Schekel Silbers dafür ausgegeben hat, welches im Kauf gäng und gebe gewesen; das lautet nach dem Hebräischen: Silbers, das durch die Kaufleute in Umlauf gesetzt wird. Denn daher haben die Krämer bei den Hebräern den Namen, daß sie in den Landen pflegen umherzuziehen und von einem Ort zum andern laufen wie die Fliegen. Darum sagt Joseph, er wolle ihnen Macht geben, im Lande umherzuziehen und zu hantieren oder zu werben, wofern sie ihren jüngsten Bruder zu ihm bringen. Denn es geht alles auf Benjamin, daß derselbe nach Egypten möchte gebracht werden.

\quellseite{397}

\begin{quote}
\textbf{V. 35. Und da sie die Säcke ausschütteten, fand ein jeglicher sein Bündlein Geld in seinem Sacke. Und da sie sahen, daß es Bündlein ihres Geldes waren, erschraken sie sammt ihrem Vater.}
\end{quote}

259. Das hebräische Wort, rik, heißt, ledig machen. Daher kommt Racha, Matth. 5, 22., welches ein loser, nichtiger und verachteter Mensch heißt. Hier thun sie nun alle Säcke auf. Denn wir haben droben gesagt, daß sie derselben viele gehabt; aber sie haben sie in der Herberge nicht alle aufgethan. Jetzt aber hat nun ein Jeder seinen Sack ausgeschüttet, und das Getreide wird in einer gemeinen Steinkammer aufbewahrt, darüber der Vater Macht gehabt, der es wiederum herauslangen soll, zu ihrer aller täglichen Nothdurft.

260. Und da ein Jeder sein Geld in seinem Sack gefunden, haben sie gedacht, dasselbe müßte also von Joseph sonderlich bestellt und angerichtet sein, oder es müßte etwa durch einen Engel oder Teufel durch göttliche Schickung ihnen das Geld wiedergegeben worden sein. Derhalben sind sie heftig erschreckt und geängstiget worden, darum daß sie des Diebstahls halben beargwohnet werden sollten; verheben aber noch gleichwohl vor ihren Vater, daß sie Joseph verkauft hatten, und meinen auch nicht anders, denn daß Joseph gestorben sei. Und wiederum zweifelt Joseph selbst auch daran, ob sein Vater und Benjamin auch noch frisch und gesund, oder am Leben seien. Und ist das die Ursache, warum er nicht aufhört, sie immer mehr und mehr zu schrecken und zu ängstigen; welches er alles dahin gemeint, daß er sie dadurch zur Buße und Bekenntniß bringen möchte.

261. Es hat sich aber der alte Vater auch zugleich mit ihnen gefürchtet. Der muß auch mit tragen. Denn also muß es zugehen, wenn die Gottlosen gestraft werden, daß alsdann die Gottseligen auch einen Theil der Strafe tragen müssen. Darum sind sie nun alle erschrocken, beide die Brüder und auch der Vater. Denn Simeon war gefangen, welchen sie gehabt hatten ohne alle Mühe wiederum los zu machen, wenn sie nur Benjamin bringen würden. Da sie aber nun ihre Säcke aufthun, werden sie wiederum bestürzt, und fürchten, man habe ihnen heimlich nachgestellt; und wo sie zu dem Herrn des Landes würden wiederkommen, sollten sie wohl alle wegen des Diebstahls um die Hälse kommen. Dies alles gehört nur dahin, daß sie recht sollen gedemüthigt werden. Denn Joseph will ihnen traun damit keinen Schaden, sondern will ihnen vielmehr Gutes thun und sie ehren, und wenn sie nun wohl gedemüthigt sind, will er sie zur höchsten Herrlichkeit bringen.

\section[Siebenter Theil]{Siebenter Theil. Von Jakobs Verhalten, da er hört, daß Simeon zurückgeblieben, und Benjamin mit nach Egypten soll gesandt werden}

\begin{quote}
\textbf{V. 36. Da sprach Jakob, ihr Vater, zu ihnen: Ihr beraubet mich meiner Kinder; Joseph ist nicht mehr vorhanden, Simeon ist nicht mehr vorhanden, Benjamin wollt ihr hinnehmen; es gehet alles über mich.}
\end{quote}

262. Wir haben vorhin gehört, wie die Söhne Jakobs so fein künstlich geredet haben, und allerdings gemeint, sie würden auf die Weise also davon kommen. Denn sie haben gesagt, daß ohne ihr Verdienst und unverschuldeter Weise der Herr des Landes so hart mit ihnen umgegangen sei, und ihr Bruder sei nicht gefangen, sondern habe nur als ein Bürge da bleiben müssen, ohne alle Gefahr; haben aber von dem Gefängniß gar still geschwiegen, auf daß sie den Vater damit nicht erschreckten, und daß er daraus nicht Ursache nehmen möchte, etwas fleißiger nach allen Dingen zu fragen.

263. Es scheint aber, als ob die Worte seiner Söhne bei ihm nicht viel gegolten haben und daß er ihnen wenig Glauben geschenkt. Denn er sagt: Ihr bringt mich um alle meine Kinder. Und dies ist ja keine erdichtete Klage, sondern er redet aus väterlichem Herzen und Sorge im potentiellen Indicativ (das heißt, er spricht dasjenige als gewiß aus, was sie möglicherweise beabsichtigten), als wollte er sagen: Ihr wollt jetzt Benjamin hinnehmen. Ihr sagt, es habe der Herr des Landes Simeon da behalten aus guter gnädiger Absicht; und da ich jenesmal Joseph zu euch auf das Feld schickte, brachtet ihr mir die Nachricht, wie ihn ein wildes Thier zerrissen hätte. Was könnte ich denn nun anders denken, denn daß ihr Buben damit umgeht, daß ihr mich der Kinder berauben wollt, die ich am liebsten gehabt?

264. Dies sind alles väterliche Gedanken und Argwohn, welche mich selbst auch würden bekümmeret haben. Und also versucht Gott den Patriarchen Jakob durch seine eigenen Söhne, daß er ihnen gar nicht glauben oder trauen darf, sondern muß denken und Argwohn haben, es sei alles eine abgekartete Sache, wie man zu sagen pflegt, daß er nur seiner allerliebsten Kinder soll beraubt werden.

265. Hieronymus hat den rechten historischen Verstand dieser Stelle gar verdunkelt und verfälscht damit, daß er den Text allzu frei übersetzt hat, nämlich also: Simeon liegt gefangen und ihr wollt nun Benjamin auch hinnehmen. Denn der hebräische Text lautet nicht also, und die andern Brüder haben dasselbe mit Fleiß verkehrt, daß Simeon gebunden und gefangen gewesen, auf daß sie dem Vater seine Schmerzen damit nicht größer machten. Darum weiß Jakob vom Gefängniß nichts; er würde sonst Benjamin nicht haben ziehen lassen. Ja, er hatte schon beschlossen, er wollte lieber Simeon in Egypten bleiben lassen, denn daß er wollte zugeben, daß Benjamin mit den andern Brüdern sollte hinabziehen. Wenn sie das gesagt hätten: Wir sind drei Tage gefangen gewesen und Simeon wird jetzt noch gefänglich gehalten, und wird nicht können los werden, es sei denn, daß Benjamin dahin gehe; so würden sie bei dem Vater gar nichts erlangt haben. Darum sagen sie, es sei also aus gutem Bedenken geschehen, daß Simeon in Egypten bleibe als ein Bürge, auf daß sie möchten wiederkommen, und vor dem Herrn des Landes beweisen und darthun, daß sie keine Kundschafter wären.

266. Mit dieser Rhetorik und seiner geschmückten Rede gedenken sie den Vater zu überreden, daß er seinen Sohn Benjamin mit hinsenden soll. Aber Jakob begegnet ihnen wiederum mit einer kurzen Dialectik und sagt also: Ich bin nun so oft gewarnet worden zu meinem Schaden, daß ich darüber auch bin weise geworden, und kann derhalben euch nicht so leichtlich Glauben schenken. Simeon und Joseph sind nicht mehr vorhanden. Er hält dafür, daß sie beide werden umgekommen sein, und läßt sich nicht überreden, daß Simeon als ein Bürge daselbst geblieben sein sollte; sondern sagt: So oft ich euch ausschicke, es sei auf das Feld oder sonst an unbekannte Orte, so kommt ihr wieder, und laßt allezeit Einen aus, welchen ihr entweder verloren oder getödtet habt. Ihr Buben, wird er gesagt haben, ihr werdet selbst einander erwürgen.

267. Also wird der fromme alte Mann hart genug versucht von seinen eigenen Söhnen, und er wird nicht ohne großen Schmerz seines Herzens diese Klage also geführt haben, daß seine beiden Söhne nicht mehr vorhanden wären. Sollte nun Benjamin hinab nach Egypten ziehen, sagt er, so würde er durch eure Bosheit auch umkommen. Und dies ist auch wahrlich ein erschrecklich Exempel gottlosen Wesens gegen ihren Vater. Und dieweil Gott diese bösen Buben noch hat selig gemacht, so wird er uns auch selig machen. Sie haben aber doch noch endlich schwere Strafe genug leiden müssen wegen ihrer Gottlosigkeit. Ruben aber führt nun die angefangene Rhetorik vollends hinaus; wie im Text folgt:

\begin{quote}
\textbf{V. 37. Ruben antwortete seinem Vater und sprach: Wenn ich dir ihn nicht wieder bringe, so erwürge meine zween Söhne; gib ihn nur in meine Hand, ich will ihn dir wieder bringen.}
\end{quote}

268. Ruben, als der Erstgeborne, will Dank verdienen, und erbietet sich dem Vater, er wolle ihm seine zween Söhne zu Bürgen setzen. Aber es ist sehr ungereimt Ding und ist nicht weißlich von ihm gethan. Denn was kann närrischer sein, denn dem Großvater deine Enkel zu Bürgen setzen, daß er dieselben strafen soll, wo du ihm nicht würdest Glauben halten? Oder, woher weißt du, daß du gewißlich wirst wiederkommen? Wie, wenn du mit Simeon und Benjamin auch würdest gefänglich eingezogen werden? Also hat Ruben allen Verstand und gemeinen Witz oder Vernunft verloren; denn es muß viel ein anderer Zustand oder Bürgschaft sein. Derhalben verachtet ihn Jakob dessfalls auch mit seinen zwei Söhnen; denn er hat Benjamin lieber, als Ruben mit allen seinen Kindern. Ja, Benjamin ist allein noch übrig von Rahel, auf welchen der Vater alle seine Hoffnung und Trost gesetzt hatte der Erstgeburt halben; darum will er denn lieber Benjamin behalten, als Ruben und Simeon. Und scheint es, als habe er denselben auch für verloren und umgekommen gehalten.

\quellseite{398}

\begin{quote}
\textbf{V. 38. Er sprach: Mein Sohn soll nicht mit euch hinab ziehen, denn sein Bruder ist todt, und er ist allein überblieben; wenn ihm ein Unfall auf dem Wege begegnete, da ihr auf reiset, würdet ihr meine grauen Haare mit Herzeleid in die Grube bringen.}
\end{quote}

269. Jakob hat fest in seinem Herzen beschlossen, er wolle seinen Sohn Benjamin nicht ziehen lassen, und sagt also: Welche unter euch gehen wollen, die mögen immerhin gehen; diesem aber will ich nicht gestatten, daß er von mir solle gerissen werden. Der andern Söhne achtet er nicht, befehlt sie dem Glücke, es mag ihnen gehen, wie GOtt will; schweiget auch davon gar stille, wie Simeon möchte wiederum gelöset werden. Wiewohl sie ihm, dem Vater, gesagt, er sollte Simeonß halben nur gute Hoffnung haben, denn es habe mit ihm keine Gefahr. Er glaubt ihnen aber nicht, sondern meint, Simeon sei auch umgekommen.

270. Um Benjamin aber ist es ihm nunmehr allein zu thun, um den trägt er große Sorge, zieht ihn allein den Andern allen vor; und hat dazu auch wichtige und nothwendige Ursachen, nämlich, weil er der Rahel Sohn war und die Ehre der Erstgeburt ihm gehörte. Von den Andern aber allen hat er nichts Gutes denken können. Ich sehe, sagt er, daß eure Rathschläge, und was ihr vornehmt, alle unglücklich sind und ein böses Ende gewinnen; darum werde ich Benjamin nicht so leichtlich in solche Angst und Gefahr stecken. Ich will aber euch gleichwohl nicht hindern, daß ihr nicht sollt alle mit einander oder aber Etliche unter euch hinabziehen.

271. Diese Furcht und Argwohn Jakobs ist menschlich und natürlich gewesen. Denn die in der Anfechtung stecken, denken in ihrem Herzen immer an das, so am allerbösesten und betrübtesten ist: wie diejenigen, so gestraft werden und elend sind, auch nichts Gutes von GOtt denken; denn es kommen ihnen in ihrem Sinne nur solche Dinge vor, dadurch ihr Elend und Unglück nur größer gemacht wird. Ja, über das alles lassen sie sich dünken, daß alle Umstände und was sich nur zufällig möchte zutragen, unglücklich und höchst gefährlich sei.

272. Und das ist des Teufels Werk, daß er vornehmlich pflegt zu üben in einem armen Gewissen, welches vor GOtt schuldig ist; dasselbe greift er mit giftigen Bissen und seinen feurigen Pfeilen ganz erschrecklicher Weise an. Wie auch wiederum, wenn Alles mit uns wohl steht und uns nach unserm Willen geht, niemand daran denkt, wie nahe uns da der Tod, der Teufel und allerlei Gefährlichkeit sein möge; wie denn der alte Mann in der Komödie weislich sagt: Quamobrem omnes, cum secundae res sunt, maxime meditari secum oportet, quo pacto adversam aerumnam ferant, pericula, damna, exilia, daß ist: Wir sollen billig alle, wenn es uns wohl geht, an allermeisten daran denken, wie wir doch das Unglück, Gefährlichkeit und Schaden, der uns begegnen möchte, ertragen wollen. Welcher über Feld reiset und wiederum heim kommt, soll allezeit denken, entweder daß der Sohn etwa gesündigt, oder das Weib gestorben, oder aber die Tochter krank geworden, weil dies gemeine Dinge sind und jedermann leichtlich begegnen mögen; und daß er sich in seinem Herzen nicht dünken lasse, daß dergleichen Dinge neu sind. Alles aber, was da anders geschieht, denn wir zuvor gehofft, dasselbige mag man für einen Gewinn rechnen.

273. Dies fällt uns aber so nicht ein, während uns kein Unglück begegnet, ja, wir pflegen da auf beiden Ohren zu schlafen, als ob es gar keine Noth hätte; wie wir in der Anfechtung und wo die Strafe vorhanden ist, zittern und zagen. Denn also bringt es die menschliche Natur mit. Darum sind wir auf beiden Seiten arme Leute, es gehe uns wohl oder übel. Wir dichten uns selbst oftmals aus Einer Sünde, um deren willen wir gestraft werden, viel andere unzählige Sünden mehr. Und wiederum, wenn wir sicher sind und Frieden haben, wiewohl wir da auch mit vielen Sünden beladen sind, denken wir doch nicht an Eine, fürchten uns auch vor keiner Gefahr.

274. Diese Schwachheit sollen wir an uns erkennen, nämlich, dadurch unser Gewissen, wenn es sicher und fröhlich ist, viel unzähliger Freude sammelt: wo es aber traurig und betrübt ist, kann es nichts denn eitel Schaden, Jammer und Unglück sammeln. Denn das zeigen sowohl die Exempel der heiligen Schrift an, als auch eines Jeden eigene Erfahrung. Wenn mein Herz mit Traurigkeit oder Zorn beladen ist, scheinen viele Gedanken ganz wahrhaftig; wenn ich aber darnach wiederum zu mir selbst komme, so kann ich den Irrthum leichtlich verstehen.

275. Es werden aber die Söhne Jakobs daheim etwas länger aufgehalten, denn sie wohl gehofft hätten. Denn sie meinten, sie wollten bald wiederum nach Egypten ziehen, auf daß sie ihre Zusage hielten und ihren Bruder los machen könnten. Und sie hatten zwar rechtmäßige Ursache genug, bald wieder nach Egypten zu kommen; aber es steht ihnen der Vater im Wege, daß sie solches nicht vollziehen können, das sie vorhatten. Denn er will lieber, daß Simeon in Gefährlichkeit gerathe denn Benjamin, und bringt davon ein fein dialectisch Argument vor, welches wider alle ihre Rhetorik, so sie gebraucht hatten, schließt, und sagt also: Denn wo Benjamin auf dem Wege ein Unfall begegnen würde, so würde ihr mich wahrlich tödten.

276. Das hebräische Wort, ason, heißt, Schade oder Verletzung, sie sei gleich tödtlich oder nicht. Lieben Söhne, sagt er, bedenkt doch, welch ein Unglück das sein würde, wenn ich durch eure Schuld von Schmerzen und Leid sterben müßte, welches also geschehen könnte, wenn ihr Benjamin würdet wegführen, und er etwa in Gefahr kommen sollte, wie wir um Joseph und Simeon auch gekommen sind: da würde ich auch gar eines sehr schweren und unwürdigen Todes sterben müssen. Ja, er führt über das auch seine grauen Haare an; welches beweglich genug ist. „Ihr würdet“, sagt er, „meine grauen Haare mit Herzeleid in die Grube bringen.“ Ich will zwar gern sterben, und dieweil ich nun fast Alters und vieler Sorgen halben verzehrt und matt geworden, verlanget mich sehr nach dem Tode. Das ist aber ja erbärmlich, daß meine Söhne daran sollen Schuld haben, daß ich alter Mann nun noch von Traurigkeit und Leid sterben soll, und der ich wegen so vielen Unglücks lange Zeit wohl versucht und fast verzehrt bin, mein Leben noch in so großen Sorgen, Wehklagen und Schmerzen beschließen soll.

277. Dies sind wahrlich schwere Argumente, welche er wider ihre Rhetorik setzt, dadurch sie gedichtet hatten, daß der Herr des Landes mit Simeon, ihrem Bruder, ganz gnädig und freundlich umgegangen, und ihm auch noch günstig und gnädig wäre. Aber Jakob antwortet ihnen darauf und sagt: Es sei, was es wolle; Simeon ist nicht ohne Ursache unter meinen Söhnen allein aufgehalten worden, so doch den andern nichts dergleichen begegnet ist. Ich sehe wohl, daß ihr euch mit einander darüber geeinigt und diese Listigkeit also unter euch selbst erdichtet habt, daß ihr Benjamin möget hinwegs führen und mich also berauben.

278. Solche Argumente sind sehr trefflich, und sind hergenommen von dem Nothwendigen, Unmöglichen, Ehrbaren und Gewissen. Denn das ist ja eine große Sünde, daß Einer seinen alten Vater in solch Leid und Traurigkeit stecken soll, daß er davon ganz verzehrt werde und dahin sterbe.

279. Zum Letzten disputiren auch an diesem Orte die Hebräisten von dem Worte scheolah, welches sie von dem Grabe verstehen und „Grube“ übersetzen. Sie spotten aber unser, daß wir es deuten, daß es die Hölle heißen soll. Wir fragen aber nach den ungelehrten Esels gar nichts, welche in einen verkehrten Sinn dahin gegeben und schon in die äußerste Finsterniß gestoßen sind, also daß sie auch ihre eigene Mutterprache nicht verstehen können.

280. Was muß aber das Wort keber etwas Anderes heißen und scheolah auch was Anderes bedeuten. Und wie das lateinische Wort sepulchrum, Grab, auf Hebräisch keber, ein gemein Wort ist bei den Gottlosen und auch den Gottlosen: also verhält es sich auch mit dem Worte scheolah. Dies ist also gewiß und ist genommen aus gemeinem Gebrauche der heiligen Schrift, allein, daß dies der Unterschied ist: das Wort keber, Grab, bedeutet eigentlich die Erde, so ausgegraben wird, und darein der verstorbene Leib geleget und begraben wird. Denn es sterben Viele, die eben nicht in die Erde begraben werden, als, die von wilden Thieren zerrissen, oder die gehenkt oder verbrannt werden, dieselben können kein Begräbniß haben, und wird doch von Allen gesagt, daß sie aus diesem Leben in die Grube hinab fahren.

\quellseite{399}

281. Derhalben unterscheidet man diese Worte also recht, daß „Grab“ ein Ort sei, dahin der Leib und Gebeine der Verstorbenen in die Erde begraben werden, da eine gewisse bestimmte Zeit, Person und Stätte ist. Scheolah ist ein gemeiner Ort oder Behältniß, nicht allein der Leiber, sondern auch der Seelen, da alle Todten versammelt werden. Augustinus nennt es in seinem Enchiridion ad Laurentium: abdita receptacula animarum, das ist, heimliche Orte, da die Seelen aufbehalten werden. Derhalben ist dies chaos, diese Kluft, uns zwar unbekannt, darin kein Unterschied der Stätte, Zeit oder Person ist, dahin die auch fahren, die keine Gräber haben. Es ist eine gemeine Stätte, wo es anders eine Stätte heißen soll, nicht für den Leib, sondern für die Seele.

282. In dem apostolischen Symbolum sagen wir, daß Christus gelitten, gestorben, begraben und niedergefahren sei zur Hölle; da diese zwei Stücke nach Art der hebräischen Sprache unterschieden werden, nämlich das Begräbniß und die Niederfahrt zur Hölle. Und im 16. Psalm V. 10, spricht Christus: „Du wirst meine Seele nicht in der Hölle lassen.“

283. Derhalben ist scheolah eigentlich, wenn man es recht beschreiben will, der Ort oder Behältniß der Seelen. Was aber das für ein Ort oder Stätte sei, davon haben die Alten viel unzählige Fragen gemacht und Disputationen gehalten. Wir sollen uns an diesem Unterschiede genügen lassen: daß wie der Leib in der Erde erhalten wird, also die Seele, wenn sie vom Leibe geschieden, auch an ihren bestimmten Ort, scheolah, kommt. Und sind nicht unterschiedene scheolah oder Behältniß der Gottseligen und Gottlosen. Die Andern disputiren zwar etwas scharf davon, und sind mancherlei Träummente auf beiden Seiten, welche mir wenig zu schaffen geben.

284. In dem Gleichniß vom reichen Manne und Lazarus Luc. 16, 22. 23. wird gesagt, daß der reiche Mann in die Hölle begraben worden sei, welches nicht ein Grab des Leibes, sondern der Seele gewesen ist; und wir müssen doch anders denken von den Gottseligen als von den Gottlosen, so beide mit einander in eben demselben scheolah sind. Wie hier Jacob sagt und droben, Cap. 37, V. 35., auch: „Ich werde mit Leide hinunter fahren in die Grube zu meinem Sohne“: item: „Ihr werdet meine grauen Haare mit Herzeleid in die Grube bringen.“ Und vom Tode der Gottseligen ist ein ungefähr und sehr klarer Text in dem Propheten Jesaja am 57. Cap. V. 2.: „Und die Gerechten kommen zum Frieden, und ruhen in ihren Kammern“; item 1 Mos. 15, 15.: „Und du sollst fahren zu deinen Vätern mit Frieden, und in gutem Alter begraben werden“; und im 25. Cap. V. 8.: „Abraham ward zu seinen Volk gesammelt.“ Dies ist das scheolah, das ist, der Ort, dahin die Gottseligen versammelt werden, welchen Jesaja „Frieden“ nennt, da er im 57. Cap. V. 2. sagt: „Die Gerechten kommen zum Frieden, und ruhen in ihren Kammern“; denn sie haben richtig vor sich gewandelt. Da redet er ja von den Todten. Also wird gesagt zum König Josia, 2 Chron. 34, 28.: „Siehe, ich will dich sammeln zu deinen Vätern, daß du in dein Grab mit Frieden gesammelt werdest, daß deine Augen nicht sehen alle das Unglück, das ich über diesen Ort und die Einwohner bringen will.“

285. Dies alles ist aus der heiligen Schrift offenbar und gewiß, daß die Heiligen oder Gläubigen nach dem Leibe in ihre Grube (geber) gehen; wie Abraham auf dem Acker Ephron gegen Mamre ist begraben worden, 1 Mos. 25, 9. Nach der Seele aber gehen sie in ihre Kammern, in das scheolah, da sie mit ihren Vätern versammelt werden. Und es haben die Gottseligen auch wohl ihre Gruben, darein sie fahren; aber daß sie darin ruhen sollen. Und wiewohl sie mit Traurigkeit dahin fahren, kommen sie doch daselbst zur Ruhe und Frieden. Und ist zwar dies Niederfahren der Gottseligen nichts Anderes, denn daß dies Leben in einen andern Stand verändert wird, da man nicht mehr unter der Sonne und auf Erden lebt; sondern dieweil sie dieses Leben verlassen, gehen sie nun in ihre Schlafkammer und Ruhebettlein, darin sie schlafen und ruhen, bis Leib und Seele wiederum im zukünftigen und ewigen Leben zusammenkommen.

286. Darum steht in der Schrift gar nichts vom Fegefeuer; sondern sie sagt, daß die Heiligen und Gerechten in ihr scheolah gehen, da sie des allangenemsten und süßesten Friedens und Ruhe genießen. Der ist aber gerecht, welcher, wiewohl er ein armer Sünder ist, doch an Jesum Christum, den Sohn Gottes, glaubt, und die Absolution empfängt, damit er durch das Predigtamt losgesprochen wird von seinen Sünden, getauft ist und das Sacrament des Altars gebraucht hat; und ist also wahrhaftig von Sünden erledigt und auch gerecht vor Gott. Desselben ist aber kein Fegefeuer bereit, sondern das Grab oder Seele, das scheolah der Gottseligen, Friede und Ruhe in seinem Schlafbette.

287. Vor der Zeit der Menschwerdung Christi ist es genannt worden „der Schooß Abrahams“. Und Christus braucht selbst auch diesen Namen, welchen er ohne Zweifel aus den Reden der Gottseligen und der lieben Väter genommen hat, da er sagt Luc. 16, 22.: Lazarus sei in Abrahams Schooß geführt worden; aber des reichen Mannes scheolah waren Feuerflammen. Und ist dieser Zustand und die Beschaffenheit der Gottlosen Gott allein bekannt, und kann von uns nicht begriffen oder verstanden werden; aber der Leib des reichen Mannes ist in das Grab (geber) gelegt worden. Dies ist also meine Meinung, wie ich diese Worte oder Namen verstehe.

288. Es ist aber noch eine andere Frage von der Seele Christi, welche, mit der Gottheit vereinigt, niedergefahren ist zu der Hölle. Und wie es mit derselbigen Niederfahrt zugegangen sei, wissen wir allerdings nicht; denn dasselbe ist an Christo etwas Sonderliches gewesen, davon er im 16. Psalm V. 10. sagt: „Du wirst meine Seele nicht in der Hölle lassen, und auch nicht zugeben, daß dein Heiliger verwese.“ Denn die Hölle hat seine Seele nicht behalten, schlafend oder ruhend, und das Grab hat seinen Leib auch nicht behalten, daß er darin verfaulte, wie Martha sagt von ihrem Bruder Lazarus Joh. 11, 39.: „Herr, er stinkt schon; denn er ist vier Tage gelegen.“ Darum wird er genannt der Erstgeborne von den Todten, Col. 1, 15.

289. Die aber von Christo auferweckt sind, haben nicht anders zeugen können, denn daß sie geschlafen und in einem Ruhebette gelegen haben, und in Abrahams Schooß sind erhalten worden, wie sie jetzt im Schooß Christi bewahrt und erhalten werden.

290. Und so viel können wir zwar wissen von den Orten, da Leib und Seele bewahrt und aufbehalten werden, nämlich, daß Leiber, das ist, das Grab, die Stätte heiße für den Leib; scheolah aber für die Seele, die da schläft, entweder der Gottseligen, oder aber der Gottlosen; aber doch mit einem Unterschied. Und was nun solches für eine Ruhe sei, wissen wir nicht.

291. Ich denke aber oftmals daran, daß wir solch Ding nicht erreichen oder verstehen können, sintemal wir nicht einmal begreifen können, wie in diesem Leben der Mensch leben könne, wenn er schläft. Denn wo wir nicht durch stetigen und täglichen Gebrauch also an den Schlaf gewöhnt wären, so würden wir davor erschrecken, wie vor dem Tode selbst, sintemal der Schlaf nichts Anderes ist, denn gleichsam eine Ausfahrt aus diesem Leben, daß du selbst nicht weißt, wo oder wer du seiest. Dieweil ich lebe, schlafe ich nicht, sondern merke auf das, was um und neben mir ist, und gehe einher wacker und aufrecht: wenn ich aber schlafe, weiß ich von mir selbst nichts, weiß auch die Stätte nicht, wo ich sein möge, und scheint es, als sei ich in ein ander Leben gesetzt worden.

292. Weil aber der Schlaf nun so ein gewöhnlich Ding ist, erschrickt niemand davor, wie wir zu erschrecken pflegen, wenn wir etwa eine ungewöhnliche erschreckliche Krankheit sehen, oder davon sagen hören, wie die Krankheit, Epilepsie genannt, die fallende Seuche ist. Wenn damit die Menschen angegriffen werden, werden sie bald aller ihrer Sinne beraubt, und können sich gar nicht regen oder bewegen, wie man sich sonst in diesem Leben zu bewegen pflegt, ja, sie scheinen gleichsam in ein ander Leben gefahren zu sein, auch so, daß sie gar nichts fühlen, wenn sie schon vom Feuer gebrannt oder in das Wasser getaucht würden, sehen und hören auch nichts: darum sind die, so mit derselbigen Krankheit beladen, lebendig todt.

293. Auf diese Weise schlafen auch die Kindlein im Mutterleibe für und für. Und solche Wunderwerke kann niemand weder mit Gedanken noch mit Worten erreichen. Und wenn wir nicht wüßten, daß der Schlaf nützlich und gesund, und uns ohne allen Schaden wäre, so würden wir uns davor nicht anders scheuen denn als vor dem Tode. Nun ist aber diese Ruhe der verstorbenen Menschen noch viel wunderlicher, welche die Schrift pflegt einen Schlaf zu nennen.

\quellseite{400}

294. Dieweil wir aber nun, da wir noch in dieser Welt leben, bisweilen nichts weniger fühlen oder verstehen können, denn daß wir leben sollten, und Gott doch denen, die da schlafen, und gar nichts thun oder wirken, ihr Leben erhält: warum wollten wir ihm denn nicht solche große Weisheit und Macht zuschreiben, daß er die Seele auch im Tode erhalten könnte? Denn dieweil dieselbige noch im Leibe wohnt, wird ihrer auf mancherlei Weise im Schlaf gespottet, da ihr im Traum etwa lose oder eitele Gestalten vieler Dinge vorkommen, welche nicht wahrhaftig sind; doch lebt sie, unansehen, daß sie also mit dem Schlaf beschweret ist, daß sie ihre Wirkungen und Handlungen nicht üben kann; der Leib aber liegt wie ein Klotz ohne alle Empfindung. Darum wird niemand leichtlich sagen können, was der Schlaf, was Wachen, oder was die Seele sei etc.

295. Bis daher haben wir nun vom Grabe der Heiligen und auch der Gottlosen geredet, und ist allerdings schwer, etwas Gewisses davon zu sagen; uns ist auch daran nicht so groß gelegen, daß wir es so genau wissen müßten, wie es sich um die Gottlosen verhalte, dieweil wir den Worten Christi glauben müssen, der von dem verstorbenen Mägdlein Matth. 9, 24. sagt: „Das Mägdlein ist nicht todt“ etc. In solcher Erkenntniß und Glauben haben die Gottseeligen genug.

296. Die Theologen, so Scholastiker genannt werden, machen vier Sphären, das ist, vier Orte der Hölle. Den ersten nennen sie limbus patrum, darein haben sie die gesetzt, so vor der Geburt, Leiden und Auferstehung Christi verstorben sind, die da gewartet haben auf die Erlösung Christi und Niederfahrt zur Hölle aus diesem Leben. Denn der Himmel war noch nicht offen. Aber sie verstehen selbst nicht, was sie sagen oder halten. Wir heißen es Abrahams Schooß, wie das Evangelium uns lehrt, daß sie limbus patrum nennen, darein alle Heiligen und Gottseeligen gesetzt sind vor der Zukunft Christi. Und ich will des Namens halben nicht streiten, dieweil man wahrhaftig weiß, daß der Schooß Abrahams gewesen, welches nun der Schooß Christi ist.

297. Denn dieweil Mose und Elia mit Christo auf dem Berge reden, Matth. 17, 3., erhellt daraus, daß er dieselben Heiligen alle in seiner Hand und Macht habe und daß er ihr Gott ist. Sie schlafen aber und ruhen, wie man von allen Gottseeligen weiß, also, daß sie Christus auferwecken kann, wenn es ihm wird gefällig sein. Darum ist es nicht vonnöthen, für sie zu bitten, und ist viel weniger vonnöthen, daß man so viel Geld unnützlich ausgibt, die Seelen aus dem Fegefeuer zu erlösen; wie im Psalterium geschrieben ist.

298. Haben aber nun Augustin oder andere Väter eine andere Meinung, so sollen wir wissen, daß man ihre Schriften mit rechtem Verstande lesen müsse und daß man keine Artikel des Glaubens machen soll aus den Sprüchen oder Meinungen der Väter; wie in voriger Zeit durch solche Thorheit greuliche Finsterniß und Irrthum in die Kirche geführt worden ist.

299. Gregor ist allein der Anfang und Stifter gewesen der vielen Lügen vom Fegefeuer und der Opfermessen für die Todten. Er schreibt an einem Orte, daß in seinem Kloster ein Verwalter gewesen sei, der ohngefähr und aus Unachtsamkeit unter seinen Büchern oder Briefen drei Gulden hat liegen lassen. Da Gregorius dieselben nach dessen Tode gefunden, ist er über die Maßen ungeduldig geworden über den verstorbenen Verwalter und schreit laut, er wäre um der Sünde des Diebstahls willen verwahrloset und würde dafür ewige Pein leiden müssen: so es doch glaublich war, daß derselbige Mönch in andern Dingen mehr Gelegenheit gehabt zu stehlen, und doch nie etwas entführt oder veruntreut hat; sondern dies Geld habe also heimlich unter den Briefen oder Büchern gelegen, daß er selbst davon nichts gewußt. Und sollen doch durch dasselbige Exempel und Geschrei Gregors, wie er gesagt, die Andern also bestürzt und erschreckt worden sein, daß hinfort ihrer keiner auch nicht einen Heller hat behalten wollen. Endlich aber hat Gregor befohlen, man sollte dreißig Messen für ihn lesen, und da dieselbigen sind gehalten worden, schreibt er, sei ihm derselbe Verwalter erschienen, und habe ihm höflich gedankt für solche große Wohlthat, dadurch er von der Pein und Fluch Gottes erlöset wäre.

300. Das sind eigentlich die kräftigen Irrthümer gewesen, davon Paulus schreibt, 2 Thess. 2, 11. Und ist von diesem Stempel Gregors in die ganze Welt gekommen die unendliche Menge der Greuel und der Opfermessen für die Todten. Das Gespenst aber, so ihm erschienen, ist der schalkhaftige, leidige Teufel selbst gewesen, der aus der Hölle gekommen, das ganze menschliche Geschlecht zu betrügen und zu verführen. Denn die erschreckliche und greuele Lüge des Fegefeuers und der Messen für die Todten ist so groß, daß sie niemand genugsam verstehen oder ausreden kann.

301. Es ist eine lange Fabel und lang sind die marktschreierischen Umtriebe, damit der Papst und Cardinäle groß unzählig Geld und Gut zuwege gebracht und erlangt haben. Denn es sind derselbigen greulichen Ungeheuer, als Kirchenraubes und sonst anderer Schindereien, unzählig viel gewesen. Ja, siehe alle Kirchen und Klöster an, welche mit keinen andern Kosten gebaut sind, denn mit solchem Gelde, so dazu gegeben worden, daß man Opfermessen damit gekauft hat für die Todten, dazu auch die ganze Welt ihr Geld und Gut mit Haufen hingegeben, nämlich, die Seelen damit aus dem Fegefeuer zu erlösen.

302. Und ist daran kein Zweifel, daß Christus vornehmlich und sonderlich wird geredet haben von der Profanation und Mißbrauch des Abendmahls des Herrn, damit sie aus der Messe ein Opfer gemacht für die Lebendigen und Todten, da er sagt Matth. 24, 15.: „Wenn ihr nun sehen werdet den Greuel der Verwüstung, davon gesagt ist durch den Propheten Daniel, daß er stehet an der heiligen Stätte (wer das lieset, der merke drauf)“ etc. Denn dies ist ein recht Götzenhaus und teuflisch Ungeheuer gewesen, hat aber ihnen sehr wohl gedient zu ihrem Gewinn, Geld und Gut, dazu auch Gewalt über alle Königreiche der Welt damit zu erlangen.

303. Die Lehre aber des Evangeliums, so da zeugt, daß die Gottseeligen schlafen und ruhen in Christo, widerlegt und verwirft all dies gottlose und abgöttische Wesen vom Opfer des Leibes und Blutes Christi, das sie in ihren Messen für die Todten gethan haben. Darum fürchten sich die Gottseeligen, so im Herrn entschlafen sind, vor keinerley Pein des Fegefeuers, daraus sie die Lebendigen mit Gold oder Silber erlösen könnten. Denn es steht geschrieben Jes. 57, 2.: „Die richtig vor sich gewandelt haben, ruhen in ihren Kammern.“ Item, Offenb. 14, 13. steht: „Selig sind die Todten, die in dem Herrn sterben, von nun an. Ja, der Geist spricht, daß sie ruhen von ihrer Arbeit.“ Welchen Spruch sie selbst auch oftmals pflegen anzuführen und haben doch denselben fälschlich auf die Erlösung der Seelen aus dem Fegefeuer gezogen. Derselben ist im Papste und seinen Aposteln eine teuflische Bosheit gewesen, dazu auch eine erbärmliche Blindheit der thörichten Menschen, die gar keine Sinne gehabt, und allzu leichtfertig gewesen sind, zu glauben und zu solchen schändlichen unglaublichen Mißbräuchen Geld zu geben.

304. Wir sollen aber unsere Herzen mit Sprüchen der heiligen Schrift wider solche wunderliche Meinungen und falschen Wahn befestigen, nämlich, daß die Seelen der Heiligen nach diesem Leben nicht in der Pein liegen, sondern ruhen von ihrer Arbeit; wie Christus sagt Joh. 8, 51.: „So jemand mein Wort wird halten, der wird den Tod nicht sehen ewiglich“; item Joh. 11, 25. 26.: „Ich bin die Auferstehung und das Leben. Wer an mich glaubet, wird leben, ob er gleich stürbe; und wer da lebet und glaubet an mich, der wird nimmermehr sterben.“ Er selbst, der Herr Himmels und der Erde, ist vom Vater in diese Welt gesandt worden, daß er dasselbe hat lehren sollen; welches man also fassen und sich daran festhalten soll. Werden wir aber damit betrogen, o! wie wollen wir es dann so gut haben, wenn wir durch unsern lieben Heiland und da wir an den Sohn Gottes glauben, betrogen werden!

305. Von diesen Dingen muß man Erinnerung thun um der jungen Knaben willen, so jetzt daher wachsen, und nicht wissen, wie große Finsterniß, Irrthum und Jammer in der papistischen Religion gewesen ist, welche sie noch mit Gewalt gebeten zu vertheidigen und zu erhalten; und vornehmlich halten sie über die Massen fest an dem Fegefeuer und des großen Gewinnes ihrer Opfermesse um des großen Gewinnes willen, so sie davon über ganz Europa gehabt.

306. Den dritten Ort eignen die Scholastiker den Verdammten zu, welcher uns gar nichts angeht, dieweil Christus allein seine Heiligen erlöset, daß sie den Tod nicht schmecken; die Andern aber wollen und können wir nicht erlösen, sondern befehlen sie dem Gerichte Gottes. Es ist aber kein scheolah der Gottseeligen,
\quellseite{401}
sondern die Grube der Gottlosen. Und wie es um dieselbige beschaffen sei, davon kann man auch nichts gewiß sagen; und diesen Ort nennen sie eigentlich Gehenna. Hieronymus sagt, daß der Herr Christus dies Wort selbst erfunden habe.

307. Der Ursprung aber desselben Orts ist, daß es mag hergekommen sein von dem Thal oder der schönen lustigen Wiese vor Jerusalem gelegen, da ein Wald gewesen gegen Aufgang der Sonne, welches gedacht wird Jer. 7, 31., da er genannt wird das Thal Hinnom. Denn da haben die Juden oft und gern geopfert, und verachteten also das Priestertum und Tempel, so in der Stadt war und Gottes Wort und Verheißung hatte, nämlich, daß Gott daselbst das Opfer und Gebet dieses Volks annehmen wollte. Aber sie hatten mehr Lust zu dem schönen lustigen Thal, das mit vielen Bäumen besetzt gewesen. Welches also eine stetige Unsinnigkeit ist der Leute zu allen Zeiten und bei allen Völkern, nämlich, daß sie außerhalb des Tempels und wider göttliche Ordnung immer gern den schönen lustigen Wäldern, Brunnen und Wasserflüssen nachfolgen, gerade als ob das Opfer und Gebet da Gott angenehmer sein sollte, als an dem Orte, welchen er selbst in seinem Worte zum Gottesdienst bestimmt hat. Da haben sie aber Altar und Capellen gebaut, wie das hin und wieder in den Predigten der Propheten zu sehen ist.

308. Und vor dieser Zeit ist es auch also zugegangen, daß die Leute die rechten Pfarren und das Predigtamt verachtet und in die Wälder und nach den Heiligthümern der verstorbenen Heiligen gelaufen sind. Und die Papisten verfolgen jetzund auch noch die rechte Lehre des göttlichen Worts, und laufen häuffig vom Evangelium und aus ihrem häuslichen und ordentlichen Beruf zu Sanct Jakob, zu unserer lieben Frauen und andern dergleichen Orten mehr.

309. Also hat sich niemand an seiner Taufe genügen lassen, oder an der Kirche, da er das Wort Christi und die heiligen Sacramente gehabt, da er gewesen, getauft und absolvirt ist, da er Gott selbst gehabt, gehört und empfunden hat, daß er mit ihm in dem Worte redete: sondern an allen Orten baute man Kirchen und Altäre, ließ das Wort und die Sacramente fahren; eben wie die Juden thaten, so daß Thal Hinnom auch denn Tempel zu Jerusalem vorgezogen haben, und dahin der König, die Königin, die Fürsten, Priester, Propheten und das Volk mit großer Anzahl gelaufen sind.

310. So groß ist die Gewalt dieses Teufels, welchem dasselbe also von Gott verhänget wird, daß er des Tempels, des Gottesdienstes und Gottes Volks also spotten möge. Und ist wahrlich nicht Wunder, daß viel Leute dadurch geärgert werden, wenn sie diese Unordnung sehen, daß sie entweder sagen, es sei gar kein Gott, oder aber schreien, Gott wisse davon nichts, achte auch nicht, was die Menschen auf Erden handeln, darum daß er dem Teufel so lange zusieht und ihm seinen Willen läßt, und den großen Greuel und Lästerung seines heiligen Namens und des rechten Gottesdienstes nicht straft.

311. Denn siehe doch an das unverschämte gottlose Wesen der Päpste, wie mit großer Kühnheit und Frevel sie noch ihre Messe, Wallfahrten, Möncherei und andere dergleichen Greuel mehr vertheidigen. Mit welchem Greueln allen unser Herr Gott erschrecklicher Weise erzürnt wird, daran kein Zweifel ist; und läßt doch gleichwohl den Teufel so frech und unsinnig sein teuflisch Spiel treiben, Gott selbst zum Hohn und Spott.

312. Ein solches Thal ist es nun gewesen, welches die Juden gehabt haben anstatt der rechten Kirche, und haben es genannt Thopheth oder Kinder Hinnom, und bei demselben ist der Delberg gewesen und das Thal Rephaim, welches sehr fruchtbar war, wie in Thüringen ein Ort ist, der genannt wird die goldene Aue. Sie haben aber daselbst eine Kirche gebaut, welches in der Wahrheit ein rechtes Hurenhaus gewesen, so dazu sehr wohl gedient, daß damit die armen Seelen betrogen und verführt worden. Und sie sind auf denselben erdichteten Gottesdienst so gar befliessen, toll und thöricht gewesen, daß sie daselbst nicht allein unvernünftige Thiere, als, Ziegen, Schafe, Kinder geopfert haben, wie im Tempel zu Jerusalem geschehen und dazumal verachtet ward; sondern sie haben daselbst, als an einem Orte, so etwas heiliger wäre, Menschen geopfert, ja sie haben dazu ihre Söhne und Töchter geschlachtet.

313. Es sagt auch Jona und die Andern, daß in der Kirche, so daselbst gebaut, der Abgott Moloch gewesen sei, welchen sie als einen Gott oder ihren König, oder aber einen gestrengen Richter geehrt haben, den sie also müßten versöhnen mit dem Opfer ihrer Söhne und Töchter; wie Ahas alle seine Söhne geopfert oder durchs das Feuer gezogen hat, daß er auch nicht hat lassen einen Erben zum Königreiche überbleiben. Es ist aber ein Abgott gewesen von Erz und inwendig hohl, nach dem Bilde eines Menschen gemacht, und hinten mit Feuer gefüllt, bis er gar geglüht und gleichsam gebrannt hat. In desselben Götzen Schoß ward das Kind gelegt, welches dem Abgott sollte geopfert werden, und die Priester haben mit Klappern, Schellen und Pauken dazu geläutet, daß die Eltern das Geschrei des sterbenden Kindleins nicht hören sollten, und haben dazu ein sonderliches Geschrei gehabt, damit sie den Eltern zugejauchzt, Glück gewünscht und gedankt haben, daß Gott das Kindlein ohne Schmerz und Pein angenommen hätte. Hilf Gott, wie ist das so eine erschreckliche unmenschliche Tyrannei und Grausamkeit gewesen! Und ist doch gewiß, daß diese Opfer bei ihnen sehr gemein gewesen; wie der 106. Psalm V. 37. 38. sagt: „Und sie opferten ihre Söhne und ihre Töchter dem Teufel und vergossen unschuldig Blut, das Blut ihrer Söhne und Töchter, die sie opferten den Götzen Canaanis“ 2c.

314. Darum droht Gott, er wolle dieselbe Stätte gar zu Grunde zerstören; wie er sagt Jer. 7, 31—33.: „Die Kinder Israel bauen die Altäre Thopheth im Thal Ben Hinnom, daß sie ihre Söhne und Töchter verbrennen; welches ich nie geboten noch in Sinn genommen habe. Darum siehe, es kommt die Zeit, spricht der Herr, daß man es nicht mehr heißen soll Thopheth und das Thal Ben Hinnom, sondern Würgethal; und man wird in Thopheth müssen begraben, weil sonst kein Raum mehr sein wird. Und die Leichname dieses Volks sollen den Vögeln des Himmels und den Thieren auf Erden zur Speise werden, davon sie niemand scheuchen wird“ 2c. Es soll heißen ein Todten- und Würgethal, daß dieselbe heilige Stätte ein Begräbniß werde der Erschlagenen, um der Abgötterei willen.

315. Das Wort „Thopheth“ aber kommt her von toph, welches Trommel heißt oder Pauke, als ob man sagen wollte: Ein Tempel der Trommeln, so geschlagen worden, wenn die Kinder getödtet sind; auf Deutsch, eine Trommelschule oder Pfeiferschule.

316. Also malt Jesaia im 30. Capitel ab nach poetischer Art die Pein und Marter, so über den König der Assyrer hat in der Hölle kommen sollen. Er macht ihm einen Singetanz in der Hölle. „Es wird“, sagt Jesaia Cap. 30, V. 32. 33., „die Ruthe ganz durchbringen, und wohl treffen, wenn sie der Herr über ihn führen wird, mit Pauken und Harfen, und allerlei halben wider sie streiten. Denn die Grube ist von gestern her zugerichtet“ 2c. Welches beinahe Alle, die den Propheten ausgelegt, von der Hölle, Gehenna, verstehen.

317. Und mit diesem Namen hat Christus auch die ewige Verdammniß anzeigen wollen, welcher Name vom leiblichen Todtschlage hergenommen ist. Denn wie daselbst die Gehenna oder Schlachtung der Kinder gewesen, also hat Gott hernach ihr Blut gerächt mit einer sehr großen Schlacht oder Niederlage der Juden, die auf derselbigen Stätte von den Chaldäern erschlagen sind.

318. Ich halte aber dafür, daß die Kinder, so in demselbigen Opfer umgekommen, selig geworden sind; sintemal der 106. Psalm V. 38. sagt: Sie haben unschuldig Blut vergossen, nämlich, das Blut ihrer Söhne. Und wie die Kindlein von Herodes getödtet sind, also sind diese auch geschlachtet und verbrannt worden. Aber die tollen Eltern, die solche Opfer geopfert, haben auch dafür endlich greuliche Strafen empfangen.

319. Das Wort chinnam heißt frei, umsonst, von chanan oder chen; als, 2 Mos. 21, 2. steht von einem hebräischen Knechte, der soll im siebenten Jahre frei ledig ausgehen 2c.; und Jes. 52, 3.: „Ihr seid umsonst verkauft, ihr sollt auch ohne Geld gelöset werden“; item in 96. Psalm V. 5.: „Sie hassen mich ohne Ursache.“ Daher kommt chen oder chamnah, das heißt, Gnade und Barmherzigkeit. Daher kommt der Name Johannes. Ich halte aber, daß dieses Thal darum also genannt worden sei, daß es ein gemeiner Ort war, dahin jedermann frei hat mögen spazieren gehen; wie bei uns die Wälder sind, welche sich niemand insonderheit zugeeignet ausgenommen der Landesfürst, und stehen sonst allen seinen Unterthanen offen, daß sie daselbst aus- und eingehen mögen.

320. Josua aber im 15. Cap. V. 8., da er die Grenze dem Stamme Juda zugetheilet, führt die Grenzlinie herauf durch das Thal des Sohns Hinnom, bei Jerusalem vorbei; doch so, daß dieselbe ausgeschlossen, und gegen Mitternacht liegen bleibt; denn sie ist nicht im Stamme Juda, sondern Benjamin gelegen. Die Andern sagen, Hinnom sei der Name eines Mannes. Dies will ich aber den Grammatikern lassen befohlen sein.

\quellseite{402}

321. Ich wollte lieber das Wort den also auslegen, daß es ein Sohn sollte heißen nicht eines Mannes, sondern nach Gewohnheit der heiligen Schrift, welche nach ihrer sonderlichen Art und Weise zu reden, „ein Sohn des Todes“ pflegt zu sagen, item, „ein Sohn des Köchers“, damit sie einen Pfeil versteht; item, „ein Sohn des Lebens“. Desgleichen im Prediger Salomo Cap. 12, V. 4.: „Töchter des Gesangs“, damit sie die Ohren will verstanden haben. Und in 80. Psalm V. 15. 16.: „Suche heim diesen Weinstock, und halt ihn in Bau, den deine Rechte gepflanzet hat, und den du dir festiglich erwählet hast.“ Da redet er nicht von einem Menschen, sondern von der Stätte. Und Jes. 5, 1.: Vinea facta est dilecto meo in cornu filio olei: „Mein Lieber hat einen Weinberg an einem fetten Orte“, der hoch gelegen ist. Denn filius olei, ein Sohn des Oels, hat für uns keinen Sinn; wir pflegen auch der Figur oder Weise zu reden nicht zu gebrauchen: darum haben wir es recht ausgelegt, daß es heiße, ein fetter Ort, eine Schmalzgrube, darin Alles wohl wächst. Eben so viel heißt auch das Wort „Gethsemane“, Matth. 26, 36., der Ort, da Christus verrathen und gegriffen worden ist, welches Hieronymus gegeben hat, ein fettes Thal, und schreibt, daß daselbst sehr feine lustige Gärten gewesen seien; und Josephus gedenkt desselben auch.

322. Ich halte aber dafür, daß Christus an demjenigen lustigen Orte die Schmerzen der Hölle erlitten habe; darum ist er ihm feind, und hat ihn genannt Gehenna, die Hölle. Denn da er im Garten gebetet, ist er recht in der Gehenna und in der Hölle gewesen. Vielleicht hat auch da gestanden der Baum des Erkenntnisses Gutes und Böses. Daherhalben hat unser Heiland Christus in dieses Thal kommen müssen und blutigen Schweiß schwitzen; welcher Schweiß genau anzeigt, daß er den Tod geschmeckt habe, welcher die Hölle ist. Darum nennt er es nicht ein fettes Thal, sondern kehrt es um, wie Gehenna, und nennt es Gehenna, Hölle, oder Würgethal.

323. Also ist Christus, unser Herr und Erlöser, für uns recht in der Hölle selbst gewesen; denn er hat wahrhaftig den Tod und die Hölle gefühlt an seinem Leibe. Was er aber gethan und gefühlt habe, da er vom Leibe abgeschieden, wissen wir zwar nicht, aber im Leben und am Leibe hat er wahrhaftig die Hölle geschmeckt.

324. Und ich halte dafür, daß sein Leiden am Kreuze nicht so groß gewesen sei als das im Garten; denn es hätte es die Natur nicht ertragen können. Und das ist wahrlich ein großes Ding gewesen, daß er den harten schweren Kampf hat ertragen können, der ihm aus seinem Leibe blutigen Schweiß getrieben; wie Lucas Cap. 22, V. 44. sagt: „Es ward aber sein Schweiß wie Blutstropfen, die fielen auf die Erde“; das war nicht allein Wasser, sondern lauter Blut, welches von seinem Leibe auf die Erde geflossen.

325. Was aber das für eine Anfechtung gewesen sei, weiß niemand, und man soll bitten, daß wir mit solchem Schrecken und Angst nimmermehr kämpfen dürfen. Denn wir würden es nicht ertragen können; ja, es würde unser Leib und Seele alsbald mit einander darüber erlöschen.

326. Darum sollen wir wissen, daß Christus dazumal die Pein der Hölle hat müssen leiden, aber doch im Fleische; denn was die Seele in der Hölle gethan, davon disputiren Viele auf mancherlei Weise. Ob er die Hölle beraubt, und die Seinen erlöst habe, die vor seiner Zukunft in den Schooß Abrahams versammelt gewesen, ist nicht vonnöthen, daß man darnach allzugenau fragen oder forschen wollte. Dies wissen wir gewiß, daß alle Heiligen oder Gläubigen ewig erlöst sind von der Gewalt der Hölle. Und dieselbe Gnade hat uns Christus im Evangelium geoffenbart.

327. Inzwischen aber bleibt das Grab, scheolah, für die verstorbenen gottlosen Menschen. Gehenna aber ist nicht ein Ort der Gottseligen; sondern die Gottlosen haben daselbst ihr Behältniß, da der Ort, Topheth genannt, das ist, die Trommel- oder Pfeifenstriche in ein ewiges Heulen und Zahnklappen verändert werden wird; wie Christus sagt Matth. 8, 12. Die Hölle der Gottlosen ist kein scheolah, wird nicht Friede, Ruhe oder eine Schlafkammer genannt; wir wissen auch nicht, ob sie schlafen oder wie sie schlafen nach dem Tode. Dies habe ich an diesem Orte sagen wollen von dem Unterschiede zwischen den beiden Worten Grab und Gehenna, welches man die Hölle nennt; damit nicht vonnöthen wäre, auf ein andermaal davon sonderlich und weitläuftiger zu handeln oder zu reden.

328. Wir wollen nun wiederum auf den Patriarchen Jakob kommen, in welches Exempel uns eine sehr schöne Lehre vorgehalten wird, nämlich, daß die heiligen Männer auch Menschen gewesen sind und nicht Steine, Klöße, oder Stöcke. Denn Jakob vergißt hier all der herrlichsten Verheißungen und fällt aus menschlicher Schwachheit dahin, und wird aus Israel, der zuvor mit Gott gekämpfet und obgelegen war, 1 Mos. 32, 28., nun wiederum ein sehr verachteter geringer Mensch, disputirt und ist bekümmert, wie es seinen Söhnen gehen möge, ob sie lebendig oder todt seien, und ziert sich so gar kindisch und weibisch, als irgend ein Mensch thun könnte, nicht anders, denn als ob du den Heiden Cicero sähest, wie derselbe seine Tochter Tullia beweint, oder David, da er über seine Söhne Ammon und Absalom hat Leid getragen, an welchen du wahrhaftige und lautere Menschen siehst.

329. Also wird zwar Jakob hier auch beschrieben, daß er ein solcher Mensch auch gewesen, und wird dasselbe auch oftmals gesagt von den allerhöchsten und heiligsten Männern, welche doch eigene Verheißungen gehabt, so insonderheit sie allein betroffen haben. Wie wir die Taufe, das Abendmahl des Herrn und die Verheißung des Evangeliums haben, so uns gegeben ist; daran sollen wir uns festhalten und damit trösten, also daß wir nimmermehr kleinmüthig werden oder verzagen, es gehe uns gleich wohl oder übel. Denn warum wollte ich mich fürchten vor der Pestilenz, vor dem Türken, vor anderer Gefährlichkeit und vor dem Tode meiner Kinder? Warum bin ich so närrisch, so ich doch einen gnädigen Gott habe, der mich schützt, der mir günstig ist und mich tröstet?

330. Jakob beklagt sich aber hier gar kindisch, da er sagt: „Ihr beraubet mich meiner Kinder“. Er könnte nicht fleischlicher geredet haben, daß der Verheißung, so ihm gegeben war, mehr zuwider sein möchte. Denn in derselben Verheißung sagt ja Gott klarlich: Du sollt wachsen und gemehret werden; ich will dein Gott sein. Von dem allen aber weiß er jetzt gar nichts, ja, er weiß von sich selber nichts, gedenkt auch nicht, daß unsers Herrn Gottes Verheißung gewiß und wahrhaftig sein müsse. Er hat es gar aus dem Sinne geschlagen und wird zu einem lautern Kinde.

331. Das ist es, das wir oft gesagt haben und daß die Schrift bezeugt, daß wir noch im Fleische sind, wiewohl wir gerecht sind, den Heiligen Geist und Vergebung der Sünden haben, sintemal wir Christum bekennen, daß er unser Heiland sei; welches unmöglich wäre, wenn nicht der Heilige Geist in uns wäre, welcher Christum verkläret.

332. Derhalben, wer Christum mit Ernst und gutem Gewissen rühmt und ihn preist, der hat einen gnädigen Gott; wie die Schrift und Worte der Apostel und des Sohnes Gottes selbst zeugen, welche uns ja nicht betrügen oder fehlen können. Warum klagen wir denn so heftig und sind so betrübt? Warum klagt Jakob, welcher auch Verheißung hat, und vergißt doch derselben so gar, ist ganz fleischlich? Die Sünde streitet nicht allein in seinen Gliedern, sondern führt einen gewaltigen Krieg und nimmt ihn gefangen. Ist es denn nicht unbillig, daß wir also geübt und geschwächt werden sollen von unserm eigenen Fleisch? Und kommt über das der Teufel auch noch dazu, welcher Ursache nimmt vom Fleisch und den Unglauben in uns immer mehren hilft, daß wir Gott nicht vertrauen können, ja, er macht, daß wir Gottes vergessen.

333. Derhalben ist dies alles uns zum Troste geschrieben, daß Jakob wider seine Verheißung, so er von Gott empfangen, so schwerlich gesündigt. Denn er sollte billig also gesagt haben: Was kümmert es mich, mein Sohn sterbe oder lebe? es ist um einen Sohn zu thun. Ich habe einen Gott, der mir eine Verheißung gethan hat, der mich tröstet, schützt, der gnädig und wohlthätig ist, der mir noch viel mehr und größere Wohlthaten erzeigen und dieselbigen häufen kann. Warum klage ich denn also? Warum bin ich so betrübt? Aber wider dies alles sagt er noch also: „Ich werde mit Leide hinunter fahren in die Grube.“ Ei, daß der große treffliche Mann so gar weibisch und findisch klagen soll und so leichtlich aller Verheißung vergibt! Denn daß seine Söhne so schwerlich sündigen, ist kein Wunder; denn sie sind unbußfertig und gar sicher, wollen nichts wissen von dem großen schweren Verbrechen, das sie nun fast zwei und zwanzig Jahre auf sich liegen gehabt, welches sie noch immer verhehlen und nicht bekennen wollen: es sind rechte böse Buben. Aber daß der Vater also zittert und zagt, steht ihm ja nicht wohl an.

\quellseite{403}

334. Darum wird uns in diesem Exempel die Lehre des Heiligen Geistes vorgehalten, welche in allen Historien der lieben Väter sehr gerühmt, und an andern Orten auch gelehret und angezeigt wird, nämlich, daß die natürlichen Affecten und Neigungen, so dies Klagen und Trauern zuwege bringen, noch immer in den Heiligen oder Gläubigen bleiben, und daß Gott nicht ein solcher Gott sei, der die Natur zerbrechen, sondern dieselbe heilen wolle.

Der Heilige Geist macht aus den Menschen keine Klötze oder tolle Phantasten, wenn er ihnen den Glauben gibt; sondern erhält und mehrt alles, was Gutes in der Natur ist, als solche väterliche und kindliche Affecten und Neigungen sind, welche Eltern und Kinder zu einander haben; denn es sind ja Gottes Creaturen. Sie werden wohl in vielen Menschen geschwächt, und werden vom Teufel hinweggenommen und aufgehoben; aber Gott läßt sie bleiben, und hebt sie nicht auf, wie die Mönche und Rottengeister thun. Denn Joseph ist ja auch ein starker und heiliger Mann, der einen großen Muth gehabt; doch weinet er auch, wie wir droben, V. 24., gehört haben.

335. Wiewohl aber nun diese Schwachheit an Jakob Sünde ist, so ist es ihm doch nütze und gut zur Demuth, zur Furcht Gottes, zum Glauben und rechter Anrufung oder Gebet, und zu einer größern Freude in seinem Herzen; wie wir das in nachfolgenden Capiteln sehen werden.

\chapter{Das dreiundvierzigste Capitel}

\section[Erster Theil]{Erster Theil. Wie Jakob beschließt, daß seine Söhne wieder sollen nach Egypten reisen, und dieselben verlangen, daß Jakob Benjamin mitschicken solle; und wie sich Jakob dagegen verhält}

\subsection{I}

\begin{quote}
\textbf{V. 1—5. Die Theurung aber drückte das Land. Und da es verzehret war, was sie für Getreide aus Egypten gebracht hatten, sprach ihr Vater zu ihnen: Ziehet wieder hin, und kaufet uns ein wenig Speise. Da antwortete ihm Juda und sprach: Der Mann band uns das hart ein, und sprach: Ihr sollt mein Angesicht nicht sehen, es sei denn euer Bruder mit euch.}
\end{quote}

1. Nachdem nun Jakob seinen Söhnen auf ihre Bitte geantwortet, so viel Benjamin betroffen, daß er denselben mit ihnen nach Egypten senden sollte, und ihre rhetorischen geschmückten Ursachen, so sie angezogen, mit den wichtigsten Argumenten widerlegt hat: bietet sich nun eine andere bessere Gelegenheit, den Vater zu erbitten, und ihm den Sohn gleichsam abzudringen, von welchen er ganz und gar in seinem Herzen beschlossen, ihn bei sich zu behalten. Denn die Theurung hat bisher das ganze Land so hart gedrückt, daß sein ganzes Hausgesinde hätte müssen Hungers sterben, wo sie nicht wiederum nach Egypten gezogen wären und von daher Getreide geholt. Und ich zweifele nicht, es werden viel andere Leute im Lande Canaan und andern Orten mehr Hungers gestorben sein.

2. Derhalben zwingen und überwinden sie auch mit diesen einigen Argumente ihren Vater, der doch ganz hart auf seinen Meinungen gestanden. Und ist nämlich dies ihr Argument, daß der Vater selbst mit den kleinen Kindlein und allem Hausgesinde werde Hungers sterben müssen, und Benjamin werde auch nicht lebendig und gesund bleiben können, ob er schon bei dem Vater zu Hause bleibe, es sei denn, daß er ihn jetziger Zeit dem ungewissen Glücke und dem Willen des Fürsten in Egypten befehlen werde. Denn, lieber Vater, werden sie gesagt haben, siehe doch, was du thust, auf daß du uns alle und dich selbst auch nicht verderbest, dieweil du diesen einigen Sohn nicht gern willst in Gefahr setzen, welche doch nicht groß sein kann; oder es könnte auch wohl gar keine Gefahr daraus werden, wenn wir nach Egypten kommen. Derhalben, da der Vater durch diese äußerste Noth überwunden ist, wird er nun anders gesinnt werden, und seinen Sohn Benjamin, den er sonderlich lieb gehabt, von sich lassen, daß er mit ihnen hinab nach Egypten ziehen möge.

3. Warum läßt aber Gott dies also geschehen? Hatte denn nicht Jakob die Verheißung Gottes, daß er zu ihm sagte: Ich will dein Schutz sein, du sollst wachsen und groß werden; item: „Du sollst ausgebreitet werden gegen den Abend, Morgen, Mitternacht und Mittag“ etc.? 1 Mos. 28, 14. 15. Und war seinen Vorfahren, Abraham und Isaak, desgleichen auch verheißen worden, daß ihr Same ewiglich bleiben sollte. Aber Jakob wird nun jetzt diesen Verheißungen allen zuwider in die größte Angst und Noth geführt, daß er auch mit seinem ganzen Hausgesinde in die gewisse Gefahr Leibes und Lebens gekommen ist. Wo ist denn nun sein Glaube auf die Verheißungen Gottes? Warum thut er nicht Wunderwerke und erlangt von Gott mit seinem Gebet, daß seine Kinder und Hausgesinde in der theuren Zeit erhalten werden mögen? Es läßt sich wahrlich ansehen, daß der Ausgang und der Handel an sich selbst auch sich nicht mit den so großen herrlichen Verheißungen reimen wolle, damit ihn Gott gleichsam überschüttet hatte. Es hätte wohl in so großer Angst und äußerster Noth ein Wunderwerk geschehen sollen; aber er ist in gleicher Gefahr und wird auch mit derselbigen Theurung gedrückt, damit die andern Heiden und abgöttische Menschen gedrückt sind.

4. Es ist aber nun oft gesagt worden, daß die Patriarchen wohl ein sehr heilig Leben geführt haben, daß sie voll Glaubens und unsträflich gewesen sind und auch viele herrliche Thaten gethan haben, dazu mit großem Muthe und festem Glauben auf die Verheißung, so sie gehabt, die schwersten Gefährlichkeiten ausgestanden und überwunden haben: aber dies ist an ihnen das allergrößte Wunderwerk gewesen, daß sie im Glauben göttlicher Verheißungen gelebt haben, da es sich auch mit denselbigen hat ansehen lassen, als ob sie gar verschwunden und ganz vergeblich wären; wie Jakob alhier und seine Söhne mit ihren Weibern und ganzen Hausgesinde in Gefahr Leibes und Lebens gekommen, nicht anders, denn als wären sie von Gott gar verlassen und verstoßen.

5. Darum sehen wir in den Historien dieser heiligen Patriarchen nicht, daß sie auf sonderliche bestimmte Tage fastet, oder sich des Hausregimentes oder weltlicher Polizei enthalten, ihr Fleisch auf sonderliche Weise kasteit hätten; welche Werke andere Leute an Franciscus und Bernhardus rühmen und sich derselben an ihnen verwundern. Und wir verdammen oder verwerfen sie auch nicht, sondern disputiren vom Glauben und Anfechtungen, ob auch Bernhardus, Antonius, Dionysius, Dominicus, wenn sie in großer Noth und Gefahr gewesen, solchen großen Glauben in ihren Klöstern gehabt, wie Jakob in seinem Hause? Denn er leidet die äußerste Noth und Gefahr nicht anders, als die andern gottlosen und abgöttischen Menschen dieselbe erlitten. Und ist zwar in solchem Fall kein Unterschied zwischen dem Hause Jakobs und den Geschlechtern oder Hausgesinde der Canaaniter, seinen Nachbarn.

6. Was thut er aber nun? Antwortet: Er thut, was alle Heiligen zu thun pflegen; welcher Exempel uns lehren, wie ich oftmals an andern Orten gesagt, daß man von unserm Herrn Gott nicht urtheilen soll nach göttlicher Verheißung, und die Mittel, so verordnet sind, dies Leben zu erhalten und Gefährlichkeit damit zu verhüten, nicht verachten soll. Wie Viele ganz freventlich und leichtfertig disputiren die fatalia necessitate, das ist, von solcher zwingender Nothwendigkeit, die mit nichten kann oder mag verhütet werden, und sagen: Wenn mich Gott erhalten will, so werde ich in der Theurung oder Pestilenz wohl lebendig bleiben, wenn ich schon weder Speise noch Arznei brauche. Soll ich aber sterben, wird mir das alles gar nicht helfen.

\quellseite{404}

7. Dies sind gottlose Gedanken und von Gott verboten; denn er hat uns in dem Falle seinen heimlichen Rath nicht geoffenbart, nämlich, wie oder wann er dir helfen wolle, und man soll ihm auch keine Stätte, Zeit oder Person bestimmen. Denn er wird es nicht leiden, daß man ihn an die Vernunft und menschliche Weisheit also binden will; sondern er will haben, du sollst deines Berufs warten, und hat darum Himmel und Erde geschaffen, und alles, was darin ist, daß er dich ernähren und erhalten will. Derhalben sollst du wissen, daß du solcher Mittel gebrauchen sollst, und nach seiner verborgenen Weise zu regieren nicht abzulegen forschen.

8. Denn dies Leben steht in seiner Versehung und deinem Glauben, sintemal er verheißen hat, daß er wolle dein Gott sein, dich erhalten und selig machen. Darum darfst du nicht sagen: Ich thue, was ich wolle, so werde ich doch Gottes Willen nicht hindern können, daß nicht dasjenige geschehen sollte, so er beschlossen hat; denn das sind des Teufels Worte und sind verdammt in Ewigkeit. Thue aber lieber also und wandele in einfältigem Glauben, und wenn du in Gefährlichkeit oder Noth gerathen bist, so gebrauche allerlei ordentliche Mittel, die du haben kannst, damit du Gott nicht versuchest. Findest du aber in denselben Mitteln keine Hülfe, die Gott geschaffen hat, die Gefahr oder Krankheit damit zu vertreiben, alsdann wirf dein Anliegen auf den Herrn, befehl ihm dein Leben, daß er dich nach seiner Weisheit und Güte regieren wolle.

9. Die aber, so die Mittel verachten und sich selbst leichtfertig in Gefahr geben, sündigen schwerlich wider die Verheißung und versuchen Gott; wie der 78. Psalm V. 18. 19. von den Juden sagt, die Gott auch eine gewisse Weise vorgeschrieben haben, wie er ihnen hat helfen sollen: „Und versuchten Gott in ihrem Herzen, daß sie Speise forderten für ihre Seelen, und redeten wider Gott, und sprachen: Ja, Gott sollte wohl können einen Tisch bereiten in der Wüste.“ Aber Gott wird nicht leiden, daß man ihn gefangen nehmen will, und zwingen, daß er es eben nach deiner Weisheit machen müsse, und seine Vorwissenheit und Vorsehung geht dich nichts an. Damit du aber Gott nicht versuchest, so thue lieber also: Höre, was er dir zusagt, und siehe, daß du demselben mit festem Glauben angehest und dich darauf verlässet. Denn das lehrt dieses Exempel des Patriarchen Jakob. Denn wiewohl derselbe in Gefährlichkeit des Todes ist, da alle Speise verzehrt gewesen, welche seine Söhne aus Egypten gebracht hatten, verachtet er doch das Mittel nicht, so er zuvor auch gebraucht, und heißt seine Söhne wieder hinab nach Egypten ziehen und Getreide kaufen, auf daß er nicht sich selbst und sein ganzes Hausgesinde durch Hunger tödte und umbringe.

10. Also hält er sich an die Verheißung und ruft Gott im Glauben an, und versucht ihn nicht, hält es gewißlich dafür, er werde nicht verlassen werden und die Verheißungen werden auch nicht vergeblich sein; und verachtet doch das Mittel nicht, so sich ihm darbietet; wie wir droben, Cap. 12, 10., gehört haben, daß Abraham nach Egypten gezogen ist, da ihn die Theurung gedrückt hat.

11. Dies soll man fleißig merken wider die Rottengeister, die da meinen, daß es große Weisheit und Stärke sei, wenn man einhergehalten und ohne Gottes Wort durch Frevel alles wagen und ansehen darf, das groß, prächtig und ungewöhnlich Ding zu sein scheint. Wie man sagt von zwei Brüdern, die mit einander im heidnischen Lande reiset, und da sie nun gehungert hat, und die Einwohner desselben Landes ihnen Brod angeboten und gereicht haben, soll es der Eine nicht haben essen wollen, darum daß es gottlose Leute und Heiden waren, so ihnen das Brod reichten, und wo er nicht vom Himmel herab gespeiset würde, war seine Meinung, wollte er lieber sterben. Derhalben ist er auch Hungers gestorben und in die Hölle begraben worden. Der Andere aber war zwar etwas flüger, welcher die Speise, so ihm die Heiden darreichten, angenommen und gegessen hat und also am Leben erhalten worden ist.

12. Also wäre es hier auch eine große Thorheit gewesen, sich des Getreides oder Speise zu enthalten, so aus Egypten her gebracht war, darum daß die Leute daselbst gottlos und ungläubig waren. Denn ist nicht das Brod der Egypter sowohl eine gute Creatur Gottes, wie andere Speise und Nahrung auch ist? Also gebrauchen wir eben desselben Marcks, Erdreichs, Wassers, Luft und Sonne, da wir wissen, daß dessen auch die Gottlosen brauchen. Unser Herr Gott wird dir keinen Engel vom Himmel senden, der dich auf eine sonderliche Weise speise oder nähre, dieweil sonst hier auf Erden Speise genug reichlich vorhanden ist, davon du leben kannst.

13. Auf diese Weise antwortet Christus dem Satan auch Matth. 4, 7.: „Es stehet geschrieben: Du sollst Gott, deinen Herrn, nicht versuchen.“ Als wollte er sagen: Ich bedarf der Hülfe der Engel dazu nicht, daß sie mich vom Tempel oben herabführen, dieweil eine gute Treppe da vorhanden ist, daran ich wohl hinabsteigen kann.

14. Hiermit soll man denen begegnen, so von der Vorsehung Gottes und Nothwendigkeit göttlicher Vorwissenheit disputiren, so doch dasselbe sie gar nichts angeht. Denn das ist für sie nicht nothwendig, daß bei Gott nothwendig ist, sondern sie sollten die Verheißungen und Gebote Gottes annehmen, und der Creaturen gebrauchen, die dazu geschaffen sind, dies Leben damit zu erhalten. Und wo ihnen damit geholfen wird, ist es gut; wo aber nicht, sollen sie wissen, daß man doch gleichwohl Gott gehorsam sein solle, welcher, so es vonnöthen sein wird, auch wohl Wunderzeichen thun wird und ihren Glauben und Hoffnung reichlich belohnen.

15. Alhier kommt aber nun abermals in der Antwort Judas vor das hebräische Wort jesch, „ist es nun“: daher weiter das Wort tuschija kommt, Spr. 8, 14.: „Mein ist beide Rath und That“; daß ein Ding geschehen möge, dazu gebe ich nicht allein Rath, sondern auch die That. Und Spr. 3, 28.: „Sprich nicht zu deinem Freunde: Gehe hin, und komm wieder, morgen will ich dir geben; so du es doch wohl hast“, das ist, so du hast, daß du ihm geben kannst. Es ist hier nicht das Hülfszeitwort „Ich bin“, sondern es heißt, da sein, oder, vorhanden sein. „Ist nun“, sagt Juda, „daß du unsern Bruder“ Benjamin „mit uns sendest“, das ist, wenn du das bei dir beschlossen hast, so wollen wir hinab ziehen.

16. Und daraus erhellt nun, daß Juda bei dem Vater mehr Ansehen gehabt denn die andern Söhne. Denn Joseph war nun nicht mehr enthalten. Darum halte ich dafür, daß Juda dem Vater lieber gewesen sei denn die andern Söhne, dieweil seine Worte mehr bei ihm gelten, denn des Ruben oder der Andern, welche auch in derselbigen Sache mit dem Vater gehandelt und darum bei ihm angehalten haben.

17. Danach sehen wir, wie wunderbarlich unser Herr Gott seine Heiligen führt. Denn Jakob war seines allerliebsten Sohnes beraubt worden und alle seine Weiber waren ihm gestorben. Denn es ist deren keine nach Egypten mit hinab gezogen; sondern er hat sein Haus regiert allein mit seiner Söhne Weibern. Ich halte aber dafür, daß Dina an der Hausmutter Statt gewesen und im Hausregiment den vier Weibern Jakobs zunächst nachgefolgt sei. Und diese Weiber werden ohne allen Zweifel auch sehr züchtige und ehrliche Matronen gewesen sein, die das Haus Jakobs und was darein gehöret, mit großer Geschicklichkeit, dazu mit Fleiß und Treue regiert, genährt und gebessert haben, sind nicht faul und nachlässig gewesen; wie denn auch zu der Viehzucht großer Fleiß und Sorge gehört.

18. Juda aber bringt das stärkste Argument vor, das, was die Gelehrten in den Schulen nennen, a necessario et impossibili, das von der Nothwendigkeit und der Unmöglichkeit hergenommen ist, welches die andern alle übertrifft, die ab utili, honesto, decoro, das ist, von dem Nutzen, von dem, was sich ziemt und schickt, hergenommen sind. Es legt sich aber gleichwohl der Vater noch dawider. Also wird es alles durch menschliches Rathen und Thun ausgerichtet; wie auch bei uns.

19. Und man mag billig fragen, warum doch der Heilige Geist so gar geringe, kindische, knechtische, weibische, weltliche und fleischliche Dinge erzählen möge von solchen sehr heiligen Menschen, die so schöne klare Verheißungen haben; welche geringe Dinge ihnen sonst auch mit einem jeglichen gottlosen Menschen gemein sind? Warum schreibt er nicht andere tapfere und höhere Dinge von ihnen? Denn was ist das Großes, daß sie sich in solchem Unflat häuslicher Sorge und Arbeit abmühen müssen? Antwort: „Singe mit dem Gottlosen, daß er die Herrlichkeit des Herrn nicht sehe“, Jes. 26, 10. Gott verbirgt seine Heiligen unter solchen äußerlichen Hüllen und fleischlichen Dingen, daß nichts so verächtlich scheint wie sie.

\quellseite{405}

20. Was ist aber denn nun für ein Unterschied zwischen David und Escipio oder Julius Cäsar? Wenn du äußerlich die Hand, Schwert und Blut ansiehst, ist da gar kein Unterschied. Denn David hat nicht Lust zu selbst erwählten und sonderlichen Gottesdiensten, Religionen und Andachten der Mönche, sondern leidet auch gemeinen Jammer und Noth mit andern Menschen, so bei dem Haus- und Weltregiment zu sein pflegt. Dies ist aber der Unterschied. David lebt in der Verheißung und Geboten Gottes: Julius Cäsar hat weder Gott noch Teufel, ja, er ist des Teufels Leibeigener; denn er wird durch Gottes Gebote nicht regiert, glaubt auch den Verheißungen Gottes nicht, sondern richtet Alles zufällig und ohngefähr aus, wie es ihm vorkommt.

21. Von diesem Unterschied wissen die Mönche nichts, verstehen auch die Heiligkeit der Patriarchen nicht; denn sie ärgern sich an der äußerlichen Gestalt der gemeinen häuslichen Werke, daß der Hausvater seine Kinder, Töchter, seiner Söhne Weiber, Mägd(e) und Vieh regiert und sie versorgt, daß er Geld zählt, kauft, verkauft, und lebt also ingemein hin unter den Heiden. Da wird gar keine Heiligkeit gesehen; daß sie aber im Glauben und Gehorsam Gottes wandeln, darauf geben die Papisten keine Acht.

22. Nun sollte wahrlich das unsere erste und vornehmste Sorge sein, daß wir auf die Verheißungen und Gebote Gottes sehen, darnach unser ganzes Leben eingerichtet werden sollte; denn das ist die rechte Heiligkeit gottseliger Menschen. Dazu kommt dann auch der Gebrauch der heiligen Sacramente. Und wo derselbe rechtschaffen und unverfälscht sein wird, so werden sich vielerlei christliche Uebungen und Gottesdienste finden, damit wir unsern Glauben und Liebe vor Gott und Menschen prüfen und beweisen können. Und da wird es auch nicht fehlen, der Teufel wird uns nachstellen und zusetzen, mit welchen wir alsdann Tag und Nacht kämpfen müssen.

23. Aber dies alles verachten die Papisten ganz hochmüthig. Was Zehen Gebote! Was Glaube! sagen sie: das ist alles verlängt gelehrt worden, ist ein gemein Ding; ich will etwas erwählen, das besser und höher sein soll; will mich von dem Hausregiment, weltlicher Polizei und der Kirche abthun, will etwa in eine Wüste oder Kloster gehen etc.

24. Derhalben soll man solche Historien der lieben heiligen Väter mit allerley Rhetorik und Dialektik, wie man nur kann und mag, vertheidigen und hoch heben, auf daß wir erkennen lernen, daß der Stand GOtt der liebste und angenehmste sei, welcher durch göttliche Verheißung und Schutz regiert wird. Ist es aber der Fall, daß dir es nicht in allem, was du vorhast, nach deinem Willen gehen würde, sollst du darum nicht verzagen; oder so du auch etwa aus Schwachheit fallen würdest, kehre wieder um und vertheidige deine Sünde nur nicht, so wird es nicht Noth haben; denn es steht dir der Weg noch offen, daß du zu der Verheißung wieder kommen kannst.

25. Auf diese Weise übertreffen die Legenden der heiligen Väter die närrischen Lügen und Fabeln weit, darüber wir uns im Papstthume hoch verwundert haben. Denn was geht uns Antonius an, der sich wider sein Gewissen in die Wüste begeben hat, so er doch gewußt, daß man außerhalb Gottes Wort und Befehl weder im Leben noch in der Lehre etwas annehmen soll? Wir wissen, daß dreierlei Stände in diesem Leben sind, nämlich, das Hausregiment, die weltliche Polizei und das Kirchenregiment. Wenn man nun an denselben vorübergehen und sie verachten will, und jedermann allein seinen selbst erwählten Werken folgen wollte und sich derselben befleißigen, wer will dann ein Seelsorger oder Pfarrherr sein? Wer will taufen, absondern, und diejenigen, so ihrer Sünden halben beschweret sind, trösten? Wer will das weltliche Regiment verwalten, oder sonst die gemeine Gesellschaft der Menschen schützen? Wer will die Kinder auferziehen, oder den Acker bauen? Und sind doch diese Stände und Werke, so von Gott eingesetzt, befohlen und bestätigt, im Papstthume gar verachtet und verworfen gewesen, und hat der Teufel die seltsamen Wunderwerke der Mönche mit erschrecklicher Unsinnigkeit den Leuten mit Gewalt aufgedrungen, daß sie sich darüber hoch haben verwundern und viel davon halten müssen.

26. Darum soll man von diesen Dingen oft reden und dieselben den Leuten immer einschärfen, auf daß wir lernen, daß GOtt mit diesen Ständen nicht erzürnt wird, sondern vielmehr dazu Lust und Gefallen hat; allein, daß man damit bleibe bei seinen Geboten und darüber nicht schreite.

27. Vornehmlich aber soll man dies den Junggesellen und unerfahrnen Leuten vorhalten, die noch nicht wissen, was Sünde, was der Teufel und Fleisch sei; und soll ihnen die Exempel der Väter, der Propheten und Apostel zeigen, welche das Hausregiment, die Polizei und das Kirchenamt nicht verachtet, auch keinen Abscheu vor den Werken gehabt haben, so zum Ackerbau gehören, die knechtisch und gering sind; und haben doch gleichwohl im Glauben und Geboten Gottes leben können. Sie haben sich aber in diesem gemeinen Leben also gehalten, daß es vor fleischlichen Augen schien, als ob zwischen Abraham, Isaac und andern Vätern, so GOttes Freunde gewesen, und heidnischen Fürsten, die von Gott nichts gewußt, ja, ihn gehaßt haben, kein Unterschied wäre. Denn das Fleisch sieht die Verheißung und Liebe der Heiligen nicht, wie auch die Papisten und verkehrten Fürsten dieselbige nicht sehen, so ihr Leben nach selbst erwählter Regel der Heiligkeit und Gerechtigkeit richten.

\subsection{II}

28. Auf diese Weise haben wir gehört, wie die heiligen Väter und dieser Jacob reden und klagen, wie sonst andere Bauern und arme, schlechte, verachtete Bürger thun; wie Cap. 42, 36. und Cap. 43, 6. steht, da Jacob sagt: „Warum habt ihr so übel an mir gethan? Es gehet alles über mich.“ Welche Worte schier also lauten, als ob er der Verheißung Gottes nicht mehr glaubte. Ein solcher großer und heiliger Mann, der so viele reiche Verheißungen gehabt, soll also kläglich über seinen Sohn weinen, der doch noch am Leben ist, welchen er ja noch nicht verloren hat? Warum sagt er nicht viel lieber also: Wohlauf, lasset Benjamin mit euch hinab ziehen im Namen des Herrn, welchem es also wohl gefallen hat, von dem ich billig viel mehr halte denn sonst von meiner ganzen Familie? Aber alle seine Worte und Gedanken sind fleischlich und weltlich.

29. Der Patriarch Jacob hat noch immer den Argwohn, als hätten etwa seine Söhne Joseph, und auch Simeon, so in Egypten geblieben war, erschlagen, und steht in Sorge, daß sie nun mit einem neuen Betrug umgehen, auf daß Benjamin auch weggebracht und aufgerieben werde. Dies sind ganz weltliche und fleischliche Affecten und Bewegungen, welche allhier einem solchen Manne zugeschrieben werden, der viel reiche göttliche Verheißungen gehabt und doch noch etwas schwach gewesen; der sich mit der Sünde schlägt, welche immer mit ihm gestritten, ihm widerstrebt und ihn gleichsam gefangen genommen hat, doch also, daß er dabei noch heilig bleibt. Denn er fällt ja von GOtt nicht ab, lästert ihn nicht, sondern beweint nur seinen Jammer und Elend.

30. Welches uns also zum Trost vorgehalten wird, daß niemand so schwach sind und unsers Jammers und Elendes halben sehr heftig bekümmern, wir doch unsern Schmerz also mäßigen lernen, daß wir nicht gar verzagen. Denn der Glaube wird also geübt und angefochten mit den schärfsten Pfeilen und großen Sturmböcken, er muß der Sünde und anderen Unglücks halben, dazu wegen der Schwachheit des Fleisches sich also bemühen und kämpfen, und gewinnt und siegt doch endlich noch.

31. Diesen Kampf und Sieg des Glaubens sehen die Papisten und Mönche nicht, die in ihrer eigenen Andacht gar ersoffen und verblendet sind. Wie ich noch wohl gedenke, daß da ich ein Mönch war, ich in meinem Herzen nicht konnte zufrieden sein, wenn ich nicht alle Tage eine Messe gelesen hätte. Was GOtt befohlen hat, das ließen wir liegen; aber was Menschen erfunden und gelehrt hatten, dem sind wir nachgelaufen, als wären wir unsinnig. Wie Alhus alle seine Söhne geschlachtet und getödtet hat wider GOttes Willen, welcher öffentlich angezeigt wird Jer. 7, 31., da er sagt: Ich habe daran in meinem Herzen nie gedacht; ich habe solches weder gesagt, noch geboten.

\quellseite{406}

32. Denn das ist die große Schwachheit oder vielmehr Unsinnigkeit menschlicher Natur; davon auch der heidnische Poet fein gesagt hat in dem schönen Verse: Quod licet, hoc fugio, quod non licet, acrius urit, das ist: Was mir wohl geziemet zu thun, davor fliehe ich; was mir aber nicht geziemet, das liegt mir immer heftig im Sinne, daß ich es thun möchte.

33. Darum verstehen die Papisten solche herrliche Historien nicht. Denn sie fragen weder nach den Verheißungen, noch nach den Geboten GOttes; wie können sie denn wissen, welches gute Werke sind, oder nicht? Denn dieselben liegen verborgen in dem Leben der Patriarchen, und sind bedeckt und gleichsam verdunkelt mit fleischlichen Werken und Worten, damit sie in diesem Leben in ihrer Haushaltung umgegangen sind. Wie wir allhier sehen, daß Jakob alles denkt und redet, was sonst auch einem fleischlichen Menschen in den Sinn kommen möchte.

34. Darum sollen wir nun wissen, daß die Historien der heiligen Schrift etwas fleißiger müssen angesehen und erwogen werden, als die weltlichen Historien oder Geschichten der Heiden. Denn wo diese Historien der Väter von den lateinischen oder griechischen Schriftstellern beschrieben wären, so würden sie gehalten werden, als die nicht werth wären, daß man sie an das Licht bringe und ihrer gedenken sollte, vornehmlich in der Kirche und unter GOttes Volk. Nun aber muß man darauf merken, daß dies Buch einen andern Meister hat, nämlich, den Heiligen Geist, welcher den Geschichten der Heiligen eine große unwandelbare Kraft gibt und ihre Herzen durch Glauben und Verheißung regiert. Und weil ich derselbe nicht scheut oder es für sein unwürdig achte, solche geringe Dinge, und die zu nichts nütze zu sein scheinen, zu beschreiben, so soll uns auch nicht verdrießen, diese Historien mit großem Fleiß zu überlesen und zu erwägen.

35. Es ist zwar ein gering verächtlich Ding, daß ein Vater, der so große Verheißung gehabt, doch so erbärmlich und schier knechtisch weinen und heulen soll. Warum herrscht in seinem Herzen nicht ein solcher seiner fröhlicher Glaube, der also sagen kann: Was liegt daran, wenn schon mein Sohn in Egypten in Gefahr gekommen? hat er doch dasselbige auch GOtt, der ihn schützen und ihm helfen kann. Solcher Worte hört sie keines von ihm; alles, was er redet und was er thut, ist ganz schwach und weibisch.

36. Dies soll man aber in der Kirche unter GOttes Volk darum lehren, auf daß wir daraus lernen und unterrichtet werden, daß man GOtt, der uns seine Verheißung gegeben hat, glauben soll, und daß er immer für uns sorge, auch wenn wir schwach, elend und betrübt sind. Denn wir gefallen GOtt nicht allein, wenn wir stark im Glauben sind und große Wunderwerke thun, sondern auch wenn wir schwach und andern Leuten gleich sind. Wie Jakob nicht verworfen wird, wiewohl er schier den starken Glauben verloren hat, darin er sich GOttes so herrlich hat rühmen können; wie solcher Glaube im 31. Psalm V. 7. 8. auch gerühmt wird, da der Prophet sagt: „Ich hoffe aber auf den HErrn. Ich freue mich und bin fröhlich über deiner Güte“; und Ps. 32, 11.: „Freuet den HErrn, und freuet euch, ihr Frommen“; item Ps. 34, 2.: „Ich will den HErrn loben allezeit“ etc. Ja, das sagt man dann wohl, wenn es uns wohl geht und kein Unglück vorhanden ist; aber in der Zeit der Noth erschrecken unsere Herzen gar bald und wird ihnen angst und bange, also daß Jakob jetzt gar von GOtt verlassen und gar keine Hoffnung und Trost mehr zu haben scheint.

37. Dies wird uns nun also vorgehalten zum Trost und Exempel, auf daß wir lernen, daß GOtt bei uns sei, wiewohl wenn Unglück und Noth vorhanden, als wenn es uns wohl geht. Und das soll man den Leuten ja mit Fleiß einprägen, so oft wir sehen, daß die Heiligen also erschrocken, bemüht, geängstigt und so gar schwach sind. Denn wo du da den äußerlichen Schein ansehen wolltest, so ist es wahr, sie werden eben also gestraft und geplagt, leiden und sterben auch also, wie sonst andere Heiden und fleischliche Menschen thun, außer daß der Glaube noch immer an ihnen leuchtet, wie schwach er auch ist, und in Wahrheit nichts Anderes ist, denn ein zerbrochen Rohr und glimmendes Docht, wie Jesaia Cap. 42, 3. sagt, welches aber doch gleichwohl noch nicht erlöschen ist. Und gehört eigentlich solch Schrecken und Gefährlichkeit dazu, daß wir damit in unsern Anfechtungen und Noth aufgerichtet und gestärkt werden sollen. Denn sonst geht mich die unschätzige und fleischliche Historie, wie sie äußerlich zu sein scheint, nichts an.

38. Der Papst und die Rottengeister dichten und machen solche Heiligen, die gar stoisch, das ist, trotzig und unfreundlich sind und keine Schmerzen noch Freude empfinden sollen. Wie Münzer seine große Stärke zu rühmen wußte, daß er sich des Schmerzes enthalten und die Creaturen überwinden könnte; sagte, daß ihm die Dinge, so gegenwärtig vorhanden wären, welcherlei sie auch seien, gar nichts zu schaffen gäben, und daß er sich dieselben nicht anfechten oder bewegen ließe.

39. Aber GOtt haßt und verdammt diese Apathie, das ist, solche unmenschliche oder unnatürliche Gefühllosigkeit, und hat die Akademiker (die alles natürliche Gefühl zu ertödten suchten) auch nicht lieb, sondern will solche Leute haben, die christlich und freundlich geneigt sein sollen gegen Eltern, Kinder etc., wie GOttes Wort zeuget Hof. 6, 6.: „Ich will Barmherzigkeit haben. Da ist geboten die Liebe gegen den Nächsten, welche wahrlich nicht erfüllt werden kann ohne äußerliche Empfindung und herzliche natürliche Affecten oder Neigungen.

40. Derhalben laß die Rottengeister mit ihrer Apathie immer zum Teufel fahren. Die heiligen Patriarchen lassen sich es zu Herzen gehen und bekümmern sich, wenn es ihnen selbst und andern Leuten übel geht, und alle ihre Werke und Wundzeichen sind lauter Leiden, Jammer und Angst; welches sie mit Glauben, und bei dem Docht, der noch immer glimmt, leiden. Wiewohl solches vor den Menschen gering und verächtlich Ding scheint; aber vor GOttes Augen sind es eitel köstliche und sehr gute Opfer.

41. Also thut auch ein Hausvater: der leidet Armuth und sonst mancherlei Unglück im Glauben; ist derselbe nicht so gar stark, so ist er doch schwach, dieweil er ein Christ, und bereit ist, daß er ein viel Schwereres leiden wollte. Aber einen solchen Mann halten die Papisten nicht für heilig, denn sie sind Klötze und Stöcke, ohne allen Verstand, ohne Sorge, ohne Liebe, die mit andern Leuten gar kein Mitleiden haben, und welche allein dem Bauch und ihrem Fraß dienen.

42. Uns aber unterrichten diese Exempel gar herrlich und stärken unsern Glauben, damit wir glauben, daß GOtt bei seinen Heiligen sei und für sie sorge, auch wenn es scheint, als ob er ihrer nicht achte und nichts von ihnen wisse, oder sie gar nicht kenne; wie im 4. Psalm V. 4. steht: „Erkennet doch, daß der HErr seine Heiligen wunderbarlich führet.“ Er führt sie nicht auf solche Weise, wie die Vernunft und papistischen Heiligen vorschreiben; sondern er gibt uns seine Verheißung, und will haben, daß wir uns daran halten, und uns im HErrn rühmen und freuen.

43. Es ist aber eine andere Zeit der Freude und eine andere der Traurigkeit. Wenn glückselige Zeit ist, so sind sie fröhlich, und singen: „Ich will den HErrn loben allezeit, sein Lob soll immerdar in meinem Munde sein“ etc., Ps. 34, 2., und alsdann freuet sich auch die Vernunft und das Fleisch und rühmt die Barmherzigkeit GOttes und seine Güte. Aber wenn es übel zugeht und Unglück vorhanden ist, da wird ihr Gesang mit Weinen und Klagen vermengt und geht halb zu, heulen halb und singen halb. Es lacht sich übel, wenn man weint. Jedoch muß man es gleichwohl gewiß dafür halten, daß GOtt die Seinen nicht verlassen wolle, mit welchem Unglück sie auch beschweret sind.

44. Gottlose und weltliche Menschen spotten und verachten diese Lehre; denn sie sehen das glimmende Docht des Glaubens nicht, und kennen auch den Heiligen Geist nicht, welcher in den Herzen der Gottlosen kräftig ist, und dasselbe glimmende Docht erhält und erweckt. Darum habe ich oft gesagt, man soll auf den vornehmsten Meister dieser Historien sehen, nämlich, auf den Heiligen Geist, welcher ja keine leichtfertige Zunge hat oder vergebliche Dinge redet, sondern uns die wichtigsten, nützlichsten und heiligsten Dinge lehrt.

\begin{quote}
\textbf{V. 7. Sie antworteten: Der Mann forschete so genau nach uns und unserer Freundschaft, und sprach: Lebet euer Vater noch? Habt ihr auch noch einen Bruder? Da sagten wir ihm, wie er uns fragte. Wie konnten wir so eben wissen, daß er sagen würde: Bringet euren Bruder mit hernieder?}
\end{quote}

45. Allhier erzählen sie nun, wie Joseph nach allen Dingen so genau geforscht habe, also daß er sie auch alle nach einander angeredet, was ein jeder für Eltern, Blutsfreunde und Brüder hätte. Und das hat er nicht allein darum gethan, daß er aus sonderlicher Liebe und Gunst mit ihnen also spielte und scherzte, sondern auch nach seinem Bruder und Vater fragen möchte. Darum haben sie droben, Cap. 42, V. 13., geantwortet: Wir sind nicht ein Haufen loser Buben oder von mancherlei Eltern geboren, sondern wir sind alle Brüder und Eines Mannes Söhne. Denn er hat dessen wollen gewiß sein, ob sie auch noch beide lebendig und gesund wären. Darum haben sie nicht von selbst, sondern da sie mit so viel Fragen sind genöthigt worden, diesen Bericht gethan und gesagt: Der eine Bruder sei gestorben, der andere aber lebe noch und sei daheim bei dem Vater.

\quellseite{407}

46. Daraus kann man nun nehmen die Auslegung des vorigen Capitels von dem Spiel, so Joseph mit seinen Brüdern getrieben, damit er sie in die größte Angst gebracht hat, auf daß er von ihnen erforschen möchte, ob auch der Vater mit seinem Bruder Benjamin noch am Leben wäre. Und da er das aus ihrer Antwort verstanden, hat er sie doch von aller Furcht und Gefahr, darein er sie geführt, noch nicht gar befreien wollen, sondern hat Simeon gefänglich einziehen und verwahren lassen, darum daß er ihren Worten nicht so leichtlich hat können Glauben schenken. Denn er hat sich noch immer besorget, sie möchten entweder beide oder aber ihrer Einer umgekommen sein.

47. Da sie aber nun mit großem Ernst alles, was sie gesagt, bekräftigt hatten, und vornehmlich von ihrem Bruder, den sie daheim gelassen hatten, und da sie nun vorgaben, daß sie ja ganz treu und aufrichtig wären, da sind sie allererst recht fein angelaufen und haben sich selbst in den Strick gesagt.

48. Wie sie an diesem Orte selbst bekennen, da sie sagen: „Wie könnten wir so eben wissen, daß er sagen würde: Bringet euren Bruder mit hernieder?“ Also pflegt Gott unserer Weisheit und unseres Rathes zu spotten und dieselben zu Schanden zu machen, wenn wir uns dünken lassen, daß wir sehr weise seien, und uns selbst wohl rathen können, daß uns dasjenige schädlich sein muß, davon wir gehofft großen Nutzen zu haben.

\section[Zweiter Theil]{Zweiter Theil. Wie Juda den Jakob zu bewegen sucht, Benjamin mitzuschicken; wie Jakob endlich einwilligt und seinen Söhnen befiehlt, Geschenke mitzunehmen; item von Jakobs Gebet und Klage, so er hiebei thut}

\subsection{I}

\begin{quote}
\textbf{V. 8—10. Da sprach Juda zu Israel, seinem Vater: Laß den Knaben mit mir ziehen, daß wir uns aufmachen und reisen, und leben, und nicht sterben, beide wir, und du, und unsere Kindlein. Ich will Bürge für ihn sein, von meinen Händen sollst du ihn fordern. Wenn ich dir ihn nicht wiederbringe, und vor deine Augen stelle; so will ich mein Leben lang die Schuld tragen. Denn wo wir nicht hätten verzogen, wären wir schon wohl zweimal wiederkommen.}
\end{quote}

49. Juda verheidigt seine Sache gewaltig, und widerlegt, was der Vater vorgewendet, und zwingt ihn mit einem unwiderleglichen Beweise. Denn Juda hat bei ihm das größte Ansehen gehabt, dieweil er den Andern etwas ungezogen gewesen wegen mancherlei Laster, so sie begangen hatten. Aber von Juda hat er etwas mehr gehalten, und wird ihn auch hernach loben und den andern Brüdern vorziehen, darum daß aus seinem Stamme Christus hat kommen sollen. Und wiewohl er die Blutschande mit Thamar begangen, hat er es doch unwissend gethan. Sonst halte ich dafür, er werde ein frommer ehrlicher Mann gewesen sein und ein vortrefflich Exempel der Hauszucht und guter Sitten unter der Familie Jakobs.

50. Und hat ihn der Vater nicht allein um der Haushaltung willen lieb gehabt, sondern auch deswegen, daß sonderlicher Verstand und Geschicklichkeit an ihm gewesen in allem, was er gethan hat. Denn dieweil er der rechte und wahrhaftige Erbe der Verheißung war, wird er ohne Zweifel Gottes Wort lauter und rein gelehrt haben, hat die Laster und Sünden seiner Brüder und des Hausgesindes gestraft, und alles sein fleißiglich verordnet und bestellt. Darum hat Jakob viel von ihm gehalten, und Juda hat sich wiederum auch zum Vater Liebe und Freundschaft und gutes geneigten Willens versehen; bittet ihn darauf, und erlangt auch von ihm, daß sonst im ganzen Hause niemand von ihm hätte bitten oder begehren dürfen. Denn er sagt nicht: Laß Benjamin mit uns hinabziehen; sondern er allein tritt herzu und wird Bürge für ihn. Das Ansehen und die Macht würde er bei dem Vater nicht gehabt haben, wo dieser nicht gesehen hätte, daß er vor den Andern tugendsam und gottselig gewesen.

51. Darum sagt er zum Vater: So du ja der Meinung bist, daß du Benjamin nicht leichtlich einem Jeden unter uns vertrauen willst, so befehl ihn doch mir auf meinen Glauben, denn du das Regiment der Kirche und des Hauses befohlen hast. „Ich will Bürge für ihn sein, von meinen Händen sollst du ihn fordern.“ Dies ist eine sehr schwere und harte Bürgschaft, aber doch nöthig; denn es betrifft das Leben beide des Vaters und auch seiner Söhne. Was sollen wir sonst Anderes thun? sagt er. Wo wir nicht werden wiederum nach Egypten ziehen, so werden wir alle sterben müssen, wie viel andere Leute von andern Nachbarn schon Hungers gestorben sind. Und du darfst ja nunmehr an meiner Treu und Glauben nicht zweifeln, dieweil du mich vorlängst aufrichtig und treu erkannt hast, darum daß mir das Predigtamt und das Hausregiment an deiner Statt befohlen ist.

52. So sucht Juda mit allen möglichen Gründen den Vater zu überzeugen, wie sonst etwa ein anderer Mensch den Andern zu überreden pflegt, dies oder jenes zu bewilligen. Und scheint hier nichts Sonderliches oder Heiliges an ihnen; sondern es reden Vater und Sohn mit einander, wie sonst Zwei aus dem gemeinen Haufen anderer Menschen pflegen mit einander Worte zu wechseln, und Einer diese, der Andere jene Ursache anzuzeigen in Geschäften, davon sie zu handeln haben.

53. Es geht nicht anders zu, denn wie es pflegt mit andern gemeinen und gewöhnlichen Händeln in diesem bürgerlichen Leben zu gehen. Und ist doch noch ein großer Unterschied. Denn diese Beiden werden regiert vom Heiligen Geiste und glauben an den rechten wahren Gott, wiewohl mit großer wunderbarer Schwachheit, so, daß es scheint, als haben sie alle Verheißung, so sie von Gott gehört und empfangen hatten, gar verloren und fahren lassen.

54. Derselben soll man die lieben heiligen Väter auf diese Weise entschuldigen, daß wir gleichwohl nicht leugnen, daß sie auch Menschen gewesen sind, die mit Fleisch und Blut gekämpft haben. Das ist aber ein großer Trost, daß solche Fehler alle verschlungen werden durch den unendlichen Reichthum göttlicher Verheißung. Wie ist doch das so eine große Unwissenheit und Irrthum, daß Jakob, der Vater, Juda und seine andern Söhne alle selbst nicht wissen, was sie thun, wissen davon nichts, daß der Fürst in Egypten ihr Bruder Joseph ist! Sie haben Gottes Verheißungen, verstehen aber und wissen selbst nicht, wie nach denselben ihr Leben regiert werde.

55. Denn also ist es mit dem Leben der Heiligen und Gläubigen beschaffen, und alle ihr Jammer und Unglück, damit sie beschweret werden, ist nichts Anderes, denn ein fein lustig Spiel, das Gott mit uns spielt. Wie ein Vater mit seinen kleinen Kindern umgeht, welchen er bisweilen etwas zu thun befiehlt, daß er doch weiß, daß es ihnen zu schwer ist und sie zu denselben nicht stark genug sind. Er heißt seine Söhne, die Kanne herbringen, heißt sie, ihm die Schoßen oder Stiefel ausziehen. Und wenn sich das Kindlein befleißigt, solches auszurichten, unangesehen, daß es dazu viel zu schwach ist, legt der Vater die Hand auch zugleich mit an, lobt das Söhnlein, wie es so fleißig und stark sei; auf daß also allgemach der Knabe zum Gehorsam gewöhnt und angehalten werde, und den Vater wiederum auch soll lieb gewinnen.

56. Solche Spiele sind dies auch, so die göttliche Majestät mit den gottseligen und gläubigen Menschen treibt, welche allein Gottes Kinder sind und ihn als einen Vater anrufen, und an seiner väterlichen Liebe und geneigtem Willen, so er ihnen verheißen hat, gar nicht zweifeln; und dieselben erfahren auch endlich in der Wahrheit, daß er ihnen die Hand reicht und väterliche Hülfe erzeigt. Aber alles, was sie thun, wie gering und schwach es ist, so wird es doch von Gott gelobt und gefällt ihm wohl; wie sich ein leiblicher Vater der kleinen Kindlein Arbeit und Fleiß, daß sie gerne etwas thun wollten, das ihnen doch zu schwer ist, gleichwohl gefallen läßt und sie deshalb lobt.

57. Also wissen Jakob und Juda von ihrer Historie oder Geschichte selbst nichts, denken auch gar nicht, daß ihnen Joseph sollte günstig und gewogen sein, und daß er sie erhöhen, reich machen und zu großer Ehre erheben wolle.

\quellseite{408}

Denn es ist noch nicht Zeit, den wunderbaren Rath GOttes und seine großen Wercke zu erkennen. Inzwischen schlafen sie gleichsam und schnarchen, geben auf die Verheißung GOttes gar keine Acht. Sie lassen sich dünken, daß sie ernstlich betrübt und gestraft werden, und daß es lauter Ernst sei, damit sie umgehen; und geht doch alles nicht anders zu, was sie reden und thun, denn als ob es im Schlafe geschähe. Hernach aber, da Joseph sagen wird, Cap. 45, V. 4.: „Ich bin euer Bruder“, da werden sie aufwachen, und sehen, daß alles, was sie für Ernst gehalten, nur Träume gewesen sind.

58. Also bringen wir dies Leben auch hin wie im Schlafe. Es ist aber an uns noch nicht zu sehen, daß wir schlafen, sind auch vom Schlafe noch nicht erlöst. Es wird aber einmal kommen, daß wir aufgeweckt und unsere Augen geöffnet werden, wenn wir allein dem Wort GOttes glauben könnten. Wenn wir aber ja ein solch groß Vertrauen in unserm Herzen nicht haben können, wie es die Wahrheit und Stärke der Verheißung GOttes fordern und haben will, so lasset uns doch in Schwachheit mit den lieben Vätern glauben, und daß wir ja wider GOtt nicht murren, auch nicht müde werden mit Anrufen, Beten und Schreien zu GOtt. Wie denn Jacob ohne Zweifel auch fleißig wird gebetet haben. Denn wo das Anrufen bleibt und man mit dem Gebet noch immer anhält, da muß auch die Verheißung und Glaube noch sein. Und wiewohl es ein glimmend Docht ist, tobt es doch oder murrt wider GOtt nicht, sondern seufzt; und bei demselben unaussprechlichen Seufzen ist gewißlich der Heilige Geist, der läßt dich nicht murren oder ungeduldig werden.

59. Derhalben, wo du ja so fest nicht glauben kannst, und GOtt nicht so loben und beten, wie du etwa gethan hast, da dir es allenthalben ist wohl gegangen: so glaube doch und bete, wie es pflegt in der Noth zu geschehen und so viel du immer kannst; wie geschrieben steht im 50. Psalm V. 15.: „Rufe mich an in der Noth, so will ich dich“ erhören und „erretten.“

60. Dies soll man der Kirche und Gemeine GOttes vorhalten zu Stärkung unseres Glaubens, daß wir von GOtt nicht abfallen, wenn wir angefochten und etwas hart versucht werden. Denn endlich werden wir sehen, daß unser Leben, wie ärmlich und elend es ist, nicht anders denn wie ein Schlaf gewesen ist; und da wir uns haben lassen dünken, daß wir in der Hölle wären, werden wir bald, wenn wir aufwachen, empfinden und erfahren, daß wir in den Himmel entzükt seien, und werden uns solcher plötzlichen Erhebung verwundern, davon wir uns zuvor nicht hätten dürfen träumen lassen.

61. Ja, es wird auch im 90. Psalm V. 5. 6. vom Leben der Gottlosen gesagt: „Du lässest sie dahin fahren wie einen Strom, und sind wie ein Schlaf“ etc. Der Papst und alle Gottlose, so in dieser Welt ihres großen Guts und Gewalt halben fast herrlich sind, sind gleich, als schliefen sie in solchem Leben und ihrer großen Herrlichkeit, und lassen sich nur also träumen, daß sie die allerseligsten Leute seien; hernach aber, wenn die Zeit des Gerichts kommt und wenn sie aufwachen, so werden sie sehen, daß alle ihr Gut nur ein lauterer Traum gewesen ist, und werden sagen: Was ist es nun? Wir meineten, wir wären reich, so liegen wir im Abgrund der Hölle.

62. Das hebräische Wort arav heißt am andern Orte, lustig und fröhlich sein; hier aber heißt es, für Einen angeloben oder Bürge werden. Nun sagt Juda noch dazu: „So will ich mein Lebenlang die Schuld tragen.“ Denn also lautet es nach dem hebräischen Texte; ist aber nicht zu verstehen von wirklicher Sünde, welche ja nach aller Theologen Meinung vergeht und nicht durch dies Leben hindurch bleibt; sondern es bedeutet eine Schuld, welche für und für währet. Es will damit nicht verstanden haben das Stündlein, darin die Sünde geschieht, oder das Werk der Sünde, so gegenwärtig ist und bald wieder hinweg geht, damit die Sünde vollbracht wird; sondern das meint er, daß er alle sein Lebenlang wolle die Schuld tragen, als hätte er die Sünde in einem jeglichen Augenblick begangen. So viel bedeutet das hebräische Wort chatha. Also steht im Buche Hiob Cap. 41, V. 16. vom Leviathan: Er verachtet und verschlingt alle Fische im Meer; denn sie haben Schuld vor ihm: er verzehrt sie wie die Schuldigen. Item, 1 Kön. 1, 21. sagt Bathseba: „Ich und mein Sohn Salomo werden müssen Sünder sein“, das ist, wir müssen angesehen werden, als die vor dir gesündigt haben; wir werden schuldig sein müssen.

63. Juda verspricht aber seinem Vater beiderlei: „Ich will“, sagt er, „Bürge für ihn sein“, und für deinen Sohn, unsern Bruder, „die Schuld tragen“. Also wird endlich der Vater überwunden. Und wir sehen hier nichts Ungewöhnliches oder das sonst auch nicht gebräuchlich wäre, sondern es geschieht nach menschlicher Weise und mit großer Schwachheit.

\subsection{II}

\begin{quote}
\textbf{V. 11–14. Da sprach Israel, ihr Vater, zu ihnen: Muß es denn ja also sein, so thuts, und nehmet von des Landes besten Früchten in eure Säcke, und bringet dem Manne Geschenke hinab, ein wenig Balsam, und Honig, und Würze, und Myrrhen, und Datteln, und Mandeln. Nehmet auch ander Geld mit euch; und das Geld, das euch oben in euren Säcken wieder worden ist, bringet auch wieder mit euch. Vielleicht ist ein Irrthum da geschehen. Dazu nehmet euren Bruder, machet euch auf, und kommet wieder zu dem Manne. Aber der allmächtige GOtt gebe euch Barmherzigkeit vor dem Manne, daß er euch lasse euren andern Bruder, und Benjamin. Ich aber muß sein, wie einer, der seiner Kinder gar beraubet ist.}
\end{quote}

64. Endlich überwindet und sieget der Glaube und übergibt er dem Juda nun seinen Sohn Benjamin, verzagt nicht, lästert auch nicht; sondern schließet und sagt also: „Ich habe mich bis daher euch widersetzt, und habe alles versucht und thun wollen, ehe ich meinen allerliebsten Sohn nach Egypten senden wollen: ich wollte GOtt nicht versuchen. Ja, ich habe mancherlei Mittel gesucht, daß ich die Gefahr, darein er hätte gerathen mögen, verhüten möchte. Jetzt aber, dieweil uns nun die höchste Noth dringt, will ich ihn göttlicher Versehung und seiner Güte befehlen.

65. Dies ist ein trefflich Exempel, dem wir in allem unserm Thun und ganzem Leben nachfolgen sollen. Denn wenn wir alles gethan haben, was uns möglich war in unserer Angst und Noth, wie sich Jacob allhier mit großem Fleiß dem Rath und Willen seiner Söhne widersezt hat, und wir gleichwohl in denselben Mitteln allen keine Hülfe finden, alsdann mögen wir sagen: Wohlauf, ich habe gethan, was ich gekonnt habe, ich habe GOtt nicht versucht; nun ist es an dem, daß ich mich an die Verheißung halte, die er mir gethan hat, und Alles nur seinem Willen und Wohlgefallen befehlen sein lasse. Denn also wirst du eine rechte gute Entschließung haben, daß die Noth dich gezwungen habe, und GOtt wird dich wunderbarer Weise erretten, oder aber wird dir ein Besseres geben, denn du zuvor hast sehen können. Allein, siehe zu, daß du nicht verzagest.

66. Denn also ist Jacob überwunden worden, nicht so sehr durch seiner Söhne Rede, als um Furcht willen des Todes, daß er nun seinen Sohn Benjamin läßt von sich gerissen werden; und fürchtet sich doch noch immer. Das Docht glimmt, das Rohr ist zerstoßen. Er sollte sich wohl nicht fürchten oder erschrecken; aber dieweil Joseph und Simeon nicht mehr vorhanden waren, da fürchtet er sich jetzt noch allenthalben. Und wiewohl er GOttes Verheißung hat, weiß er doch nicht, was es mit diesem Sohne noch werden will, welchen er ja gern vor den andern lebendig und gesund erhalten wollte. Dieweil er aber sieht, daß er ihn nicht kann bei sich behalten, widerstrebt er nun seinen Söhnen nicht länger, sondern thut, wie sie ihn ermahnen und was sie in solchem Falle von ihm begehren; läßt Benjamin ziehen, und gibt ihnen erstlich Befehl, was sie Joseph für Geschenke bringen sollen; darnach wünscht er ihnen Glück, und bittet GOtt, daß es ihnen möge wohl gehen.

67. Minchah, auf hebräisch, heißt eigentlich, Geschenk, wie wir GOtt Opfer oder Geschenke geben; als, 1 Mos. 4, 4.: „Der HErr sahe gnädiglich an Abel und sein Opfer“; nicht, daß GOtt begehre mit unsern Gaben gesättigt zu werden; denn er bedarf derselben nicht, wie der 50. Psalm V. 9. zeuget: „Ich will nicht von deinem Hause Farren nehmen, noch Böcke aus deinen Ställen.“ Gleich als wenn jemand einem Fürsten oder großen Monarchen einen Heller oder Groschen geben wollte. Daß man ihm aber ein demüthiges und dankbares Herz erzeigt, gefällt ihm wohl, damit Einer anzeigt, daß er den Fürsten als seinen Herrn gern in Ehren halten und ihm gebührliche Ehrerbietung erweisen wolle.

68. Auf diese Weise schickt Jacob Joseph auch Geschenke, daß er daraus vernehmen soll, daß sie ihm auch gebührliche Ehre und Achtung erweisen wollen. Es war fast gering, daß ihm Jacob Honig und Mandeln schenkt. Und dieweil man gemeiniglich sagt, man könne mit Geschenken beide GOtt und Menschen versöhnen und zu Freunden machen, soll dasselbe gezogen werden auf die Demuth und Dankbarkeit des Herzens, damit die Gemüther der Oberherren versöhnt werden.

\quellseite{409}

69. Wie es Gott wohlgefällt, daß ich ihn von ganzem Herzen ehre und ihm diene; und ist ihm Abels Opfer darum auch lieb und angenehm gewesen, daß er es mit Furcht und rechtem Glauben geopfert und mit derselben Gabe sein dankbares Herz hat beweisen wollen. Denn wo das Herz geopfert wird, das ist Gott und dem Könige eine sehr angenehme Gabe. Es ist aber das Herz nicht allezeit so fein aufrichtig und rein in denen, so Andern Geschenke geben.

70. Derhalben ist es bei den Alten gebräuchlich gewesen, daß man die Fürsten auch mit geringen schlechten Dingen hat versöhnen können, und sie haben nicht angesehen, wie groß oder herrlich das Geschenk wäre, sondern nur die Ehrerbietung, Demuth und Dankbarkeit des Herzens. Denn wer Einem das Herz gegeben hat, der wird ihm auch den Beutel geben. Auf diese Weise gibt nun Jakob seinen Söhnen Befehl, daß sie dem Fürsten in Egypten auch sollen Geschenke bringen, ob er vielleicht erzürnt wäre, daß er sich alsdann desto leichter versöhnen lassen wollte und ihm den Sohn wiederum los geben. Und also sind die heiligen Patriarchen durch den Glauben an die Verheißung Gottes erhalten worden, und haben doch gleichwohl alles gethan, was menschlicher Weise hat geschehen können und sollen, daß sie Gott nicht versuchten.

71. Nun sollte man aber auch die Worte auslegen und erklären, so an diesem Orte vorkommen. Ich will sie aber den Hebräern befehlen, daß sie dieselbigen etwas fleißiger und genauer erwägen; und sind auch deren etliche droben, im 37. Capitel, ausgelegt worden. Es hat mit der hebräischen Sprache und Grammatik noch Mühe und Arbeit. Das Wort bothnim wird auf mancherlei Weise ausgelegt. Ich lasse mir deren Meinung gefallen, die da sagen, es sei terebinthina, ein sonderlich Gummi oder Balsam; wiewohl das ein allgemeiner Name ist für allerlei Saft, so aus den Fichten oder Tannenbäumen fließt. Der Balsam (terebinthina resina genannt) wird sehr gelobt, darum daß er einen großen Nutzen hat in der Medicin. Und solche Geschenke pflegen sehr angenehm zu sein, welche die Natur an etlichen gewissen Orten hervorbringt, wie in Arabien der Weihrauch, item, Balsam, Myrrhen, und in Preußen der Bernstein \&c. Wiewohl daran kein Zweifel ist, es wird in Egypten desselben allen auch viel gewesen sein; und ist fast gering, daß er Honig und Mandeln dazu gethan hat. Aber, wie zuvor gesagt, solches ist ein Anzeichen eines Herzens, das Gott fürchtet, und seinen Fürsten auch ehrt, von welchem er nicht gewußt hat, daß es sein Sohn Joseph war.

72. Und also ehrte nun der Vater den Sohn und neigte sich vor ihm, wie der Traum lautet, und thut alles, was nur dazu gehört, auf daß er ihn versöhnen möge; befiehlt auch seinen Söhnen, daß sie das Geld, so sie in den Säcken gefunden, wiedergehen und bezahlen sollen, damit es ja nicht scheine, als hätte er etwas veräumet, oder aber als wollte er Gott versuchen. Aber dies ist das allertrefflichste Geschenk, daß er seinen Sohn Benjamin mitsendet.

\subsection{III}

73. Bis daher haben wir nun gesagt von dem Befehle, so Jakob seinen Söhnen gegeben. Der andere Theil ist nun das Gebet, so er gethan hat. Hier wird nun abermals das Leben dieses sehr heiligen Patriarchen dermaßen beschrieben, daß es sehr gering gewesen und kein Ansehen gehabt habe. Alles, was er thut, ist gering Ding vor unsern Augen; aber vor Gott ist es sehr köstlich und herrlich. Da findet man keine seltsamen Wunderwerke an ihm, er ist mit keinem sonderlichen Fasten oder Wachen umgegangen, da siehst du keine Stricke oder Mönchskappen; welches alles die Mönche also selbst aus eigenem Gutdünken und Satzung etlicher loser ungelehrter Leute ohne und außerhalb Gottes Wort erfunden haben. Aber das ist ein recht seliges und heiliges Leben, welches gerichtet wird nach der Regel des Berufs, und der Gebote und Verheißung Gottes, daran sich Jakob gehalten, und solchen großen schweren Kampf ausgestanden hat, welchen die Papisten gar verachten; er aber leistet damit Gott den angenehmsten Dienst und der in seinem Worte geboten ist. Denn nachdem er gethan, was da hat können und sollen geschehen, die Gefahr zu verhüten, setzt er nun das Gebet auch hinzu; damit er anzeigt, daß er die Verheißung nicht habe fallen lassen oder verworfen, wiewohl er in einer Angst und Noth jämmerlich seufzt, heult und klagt, so, daß er des Schutzes und der Hülfe, so ihm Gott verheißen, gar vergessen zu haben scheint. Denn es scheint hier noch ein glimmendes Docht, so noch nicht verlöschen, und ein zerstoßenes Rohr, so aber noch nicht zerbrochen, dieweil er schreit und betet und doch nicht ohne groß herzlich Seufzen.

74. Es ist zwar ein sehr kurz, aber doch heftig und brünstig Gebet des Glaubens gewesen. Denn es ist hergekommen von einem zerstoßenen Rohr und glimmenden Dochte; und also hat er zugleich Worte des Glaubens und Unglaubens herausgeschüttet. Denn aus dem Glauben ruft er Gott an; aber der Unglaube macht, daß er noch in Angst ist und beinahe an Benjamin verzagt.

75. Warum fürchtet er sich aber und zittert also, als wäre die Sache gar verloren? Warum freut er sich nicht vielmehr, hüpft und springt um der Verheißung willen, so er gehabt? Ei! darum, daß wir derselben, wenn es uns wohl geht, leichtlich gedenken können; aber wenn Gott sein Antlitz verbirgt, ach! so werden wir bald kleinmüthig und betrübt. Jedoch, so ist es genug, daß das Docht nicht gar verlischt, und daß wir noch rufen, wie Paulus sagt Röm. 8, 15.: „Abba, lieber Vater.“ Wiewohl wir das Rufen selbst nicht hören; wie Mose von seinem Rufen bei dem Rothen Meer auch selbst nichts gewußt. „Aber, der die Herzen forschet“ und prüfet, „der weiß, was des Geistes Sinn“, Regierde und Verlangen „ist“, Röm. 8, 27. Derselbe versteht wohl, wie groß dies Rufen ist; wir können es wegen unserer Schwachheit nicht wissen; wie der Herr zu Mose sagt, 2 Mos. 14, 15.: „Was schreist du zu mir?“, so er doch den Mund nicht hat können aufthun, da er in der äußersten Noth und in der größten Angst und Gefahr gewesen. Es hat ihm sein Herz im Leibe wunderlicher Weise geschlagen, und ist ihm anders nicht zu Muth gewesen, hat auch vor seinen Augen Anderes nichts gesehen als den gewissen Tod, er habe gleich auf Pharao gesehen, welcher ihnen nachgeeilet, oder aber seine Augen auf das Rothe Meer gewendet. Derhalben ist er um solcher großen Gefahr willen gar erschrocken und verstummt gewesen, und mag vielleicht bei sich selbst gedacht haben: Wie, wenn wir alle habe fallen lassen oder verworfen, 75. (Fortsetzung.) ersäuft oder von den Egyptern erschlagen würden? Da sagt Gott zu ihm: „Was schreiest du zu mir?“ Er schreiet zwar nicht; aber der die Herzen erforschet, der versteht es wohl und erhört auch das Schreien.

76. Auf solche Weise betet Jakob auch und hält sich an die Verheißung, aber doch sehr schwächlich; denn er heult und klagt noch gar erbärmlich. Dieweil aber noch ein Fünklein des Glaubens da ist und das unaussprechliche Seufzen im Herzen, so verzaget er nicht, sondern betet. Das Gebet aber wird er mit so wenig Worten und Gedanken nicht abgemacht haben, wie es Mose kürzlich beschrieben und verfaßt hat; denn er hat Gott angerufen, gedankt, gelobt und mit den einfältigsten Worten für Simeon und Benjamin gebetet.

77. Wir haben aber droben etlichemal diesen Namen Gottes, auf Hebräisch, el schaddai gehabt, welcher genannt wird von der Brust oder vom Nähern; wie die Griechen ihre Diana genannt haben πολυμάστορας, welche sie zu Ephesus und andern Orten mehr geehrt haben, wie Ap. Gesch. 19, 27. zu sehen ist, darum daß sie die Menschen nährte und erhielte. Sie haben aber den Namen von den Juden hergenommen, und haben ihre Götzen eben auf die Weise genannt, wie die Juden den wahren Gott auch genannt haben.

78. Also sagt Jakob el schaddai, das ist, Gott, der allmächtig ist, der alles gibt und thut, welcher Alles nährt, regiert, sättigt, versorgt und erhält, daß es in seinem Stande oder Wesen bleibe: der Gott wolle auch Gnade und Barmherzigkeit verleihen vor dem Manne, der in Egypten ein Herr ist. Durch das Wort schaddai haben sie verstanden eine solche Majestät, die Alles nährt; und rechnen heißt, kommt von racham, das heißt, erbarmen. Denn das ganze menschliche und vornehmlich das weibliche Geschlecht ist geschaffen zu nähren, aufzuziehen und Andere zu erhalten, beide an Vieh und auch an Menschen; wie der 110. Psalm V. 3. sagt: „Deine Kinder werden dir geboren, wie der Thau aus der Morgenröthe“; da stehet im Hebräischen auch das Wort merechem. Derhalben hat das Weib ihren Namen von dem Nähramt, welches die Griechen haben geben wollen mit dem Wort πολυμάστορας, wie sie die Diana genannt. Und Jakob rühmt hier Gott auch also, daß er ein Register sei und der die Menschen erhalte und versorge wie ein mütterlicher Leib; wie Jes. 46, 3. 4. steht: „Die ihr von mir im Leibe getragen werdet, und mir in der Mutter lieget. Ja, ich will euch tragen bis ins Alter und bis ihr grau werdet. Ich will es thun, ich will heben und tragen, und erretten.“ Da sagt Gott, daß er uns in seinem Leibe trage, wie ein Weib das Kindlein zu tragen pflegt; und das nicht allein bis in das Alter, sondern für und für.

\quellseite{410}

79. Also malt sich Gott uns ab in der Gestalt eines Weibes und einer Mutter. Und dasselbe bezeugen auch seine Wohlthaten, so er den Menschen erzeigt; wie Paulus sagt Ap. Gesch. 14, 17.: „Und zwar hat Gott sich selbst nicht unbezeuget gelassen, hat uns viel Gutes gethan, und vom Himmel Regen und fruchtbare Zeiten gegeben, unsere Herzen erfüllet mit Speise und Freuden.“

80. Es ist aber dies Gebet der angenehmste Balsam vor Gott gewesen und der süßeste Geruch. Derhalben sollen wir uns auch gewöhnen, ob wir wohl im Glauben schwach sind, daß wir doch beten; da werden wir alsdann erfahren, daß Paulus sagt 2 Cor. 4, 8. 9.: „Wir haben allenthalben Trübsal; aber wir ängsten uns nicht. Uns ist bange; aber wir verzagen nicht. Wir leiden Verfolgung; aber wir werden nicht verlassen. Wir werden untergedrückt; aber wir kommen nicht um“ etc. „Leidet jemand unter euch, der bete; ist jemand gutes Muths, der singe Psalmen“, Jac. 5, 13. Nur daß wir an Gottes Hülfe und Gnade nicht verzagen, der uns also versucht und straft zu unserm Heil und Seligkeit.

81. In der gemeinen lateinischen Uebersetzung stehen die Worte: in vinculis (in Fesseln); aber im hebräischen Text findet man die nicht. Wie wir droben Erinnerung gethan, daß diese Brüder mit wohlbedachtem Rath sich dessen enthalten, daß sie vor ihrem Vater des Gefängnisses nicht gedacht, auf daß sie ihn damit nicht noch mehr irre machten und erschreckten. Gott hat ohne Zweifel das Gebet Jacobs erhört, und ist ein feiner lustiger Irrthum, daß der Vater nicht weiß, daß sein Sohn Joseph noch am Leben und ein Herr in Egypten ist. Denn also spielet Gott mit seinen Heiligen ein fein lustig Spiel, aber also, daß sie das Gegentheil verstehen. Jakob sendet seine Söhne und Benjamin hinab nach Egypten mit einer Instruktion und mit Geschick und betet daneben fleißig, weiß nicht, was unsers Herrn Gottes Rath oder Wille sei.

82. Die Greuel der Sündfluth und Mönche sind diesem straks zuwider. Denn sie haben von ihrem Stande und Leben, daß sie führen, solche Verheißung nicht; darum haben sie auch solchen Glauben und Anfechtung nicht gehabt. Wenn aber die nicht da sind, so lacht und spottet der Teufel aller ihrer Werke.

\subsection{IV}

83. Das letzte Stück aber ist zumal beweglich, daß Jakob saget: „Ich aber muß sein, wie einer, der seiner Kinder gar beraubt ist.“ Das wird er mit Thränen geredet haben, und sind Worte Eines, der da jämmerlich heult und gleichsam verzagt ist. Und dies ist wahrlich eine wunderbare Historie. Ein solcher heiliger Mann, der allenthalben wohl versucht und bewährt ist im Glauben göttlicher Verheißung, soll zuletzt noch so jämmerlich rufen, schreien und klagen: schacholti! schacholti! Ach! ich bin meiner Kinder beraubt. Ich armer Mann, was soll ich thun? Da mir mein allerliebstes Weib, und Joseph, mein Sohn, den ich auch sonderlich lieb gehabt, gestorben sind, soll ich nun dazu auch noch Simeons und Benjamins beraubt werden!

84. Warum erhört doch Gott dies Geschrei nicht, so er mit Thränen und Angst gethan? Denn er wiederholt die Klage zweimal, wie er seiner Kinder müsse beraubt werden, und sagt: Ich habe keinen Sinn, ich armer Mann; was bin ich nun in meinem Elende? Er verzagt und verzagt doch nicht. Diese Affecten und Bewegungen der Heiligen versteht niemand, es sei denn, daß er es durch die Erfahrung lerne; denn da ist die allergrößte Kraft bei der äußersten Schwachheit gewesen. „Der allmächtige Gott“, sagt er, „gebe euch Barmherzigkeit vor dem Manne“; und bald sagt er wiederum: „Ich aber muß meiner Kinder gar beraubt sein.“ Er betet und verzagt zugleich; schreit und klagt nicht anders als etwa ein ungläubiger Mann und kann doch nicht beten ohne Glauben. Derhalben sind dies treffliche Exempel der Gottseligen, welche uns unterweisen, daß wir es dafür halten sollen, wiewohl alles verloren Ding ist, daß man doch an Gott nicht verzagen soll; und wenn schon die Thränen die Backen herab fließen, soll gleichwohl das Herz sagen: Gott der Allmächtige helfe! mit mir ist es verloren. Doch sofern, daß wir nicht in Worte der Verzweiflung und der Lästerung ausbrechen.

85. Wie hier die Verzweiflung nicht weit gewesen; aber es ist noch ein kleines Seufzen vorhanden, dasselbe ist so stark, daß es beide Ohren Gottes erfüllen kann; ja, es erfüllt Himmel und Erde. Und das ist es, das Salomo sagt im Hohen Liede Cap. 3, 6.: „Wer ist die, die herauf gehet aus der Wüste, wie ein gerader Rauch, wie ein wohlriechender Trochus? Es ist ein sehr dünner Rauch und ein wohlriechendes Kügelchen; man kann nichts Subtiles und Schwächeres erdenken. Wer zwingt aber denselben dünnen Rauch, daß er aufsteigen kann? Das thut die feurige Kraft des Glaubens und des Gebets, welche den ganzen Himmel erfüllt. Jakob hat zwar wenig Worte gemacht: „Ich werde meiner Kinder gar beraubt“, sagt er; aber es ist vor Gott ein sehr groß Seufzen gewesen; vor ihm aber, Gott selbst, und vor uns ganz gering. Solche Gebete sind sehr trefflich, und man soll sie fleißig merken und Acht darauf geben.

86. Diese Klage aber, daß sich Jakob beklagt, er werde seiner Kinder beraubt, ist fast das Ende aller Trübsale, Jammeres und Unglücks, so er gehabt. Nun wird etwas Fröhlicheres folgen, nachdem Mose seine Trübsale nun alle beschrieben hat. Aber dies Gebet wird er die ganze Zeit über oftmals wiederholt haben, dieweil seine Söhne aus gewesen: Morgens und Abends wird er diese Klage immer wiederholt und gesagt haben: Ich armer Mann! ich armer Mann! beide mit lauter Stimme und Seufzen. Und ist ein solch Gebet gewesen, wie Luc. 18, 1. beschrieben wird, da Christus sagt, daß man allezeit beten soll und nicht laß werden. Also hat Jakob auch ohne Aufhören gebetet, bis seine Söhne wiedergekommen sind, und ihm verkündigt haben, daß sein Gebet erhört wäre.

87. Und also wird er noch, da seine Trübsal nun am Ende gewesen, dahin gebracht, daß er dies Leid hat tragen müssen, daß er darüber geschrieen und geseufzt hat: Ach, ich bin meiner Kinder beraubt! ich bin leider beraubt! etc. Er muß seinen Sohn auch aufopfern. Wie Abraham Isaak geschlachtet hat, also opfert Jakob auch beide Söhne, Joseph und Benjamin, damit, da er gewußt, daß alle Hoffnung und Glaube der Verheißung auf ihnen gestanden. Das ist wahrlich ein sehr schwer Kreuz gewesen; darum wird er auch also erhöht.

88. Und diese Exempel sind uns darum vorgehalten, daß wir denselben auch sollen lernen nachfolgen. Und wiewohl wir vor allen Dingen das Leiden Christi stets sollen vor Augen haben, bewegen uns doch die geringen und schwachen Exempel der Patriarchen noch etwas mehr denn das hohe Exempel des Sohnes Gottes. Am Christo ist größere Demuth gewesen; er hat sich tiefer heruntergelassen und erniedrigt, denn sonst alle Patriarchen gethan haben: aber es geht ja in unser Herz nicht; denn es steht entgegen, daß wir wissen, daß Christus wahrer Gott ist. Die Patriarchen aber bewegen unser Herz und Gefühl noch stärker, da sie nur Menschen und uns allenthalben gleich gewesen sind, wenn sie in der äußersten Noth und Beängstigung gesteckt haben. Das Exempel Christi ist wunderbar, dieweil es zugleich Sacrament und Gabe Gottes ist, dadurch Christus uns von Gott gemacht ist zur Erlösung und Seligkeit.

\section[Dritter Theil]{Dritter Theil. Wie Josephs Brüder in Egypten ankommen, und Joseph befohlen, sie in sein Haus zu führen, und wie ihnen hiebei zu Muthe; wie ihnen der Haushalter Josephs Muth zuspricht und Simeon zu ihnen herausführt; item, wie er ihnen die Füße wäscht}

\subsection{I}

\begin{quote}
\textbf{V. 15—17. Da nahmen sie diese Geschenke, und das Geld zwiefältig mit sich, und Benjamin; machten sich auf, zogen in Egypten, und traten vor Joseph. Da sahe sie Joseph mit Benjamin, und sprach zu seinem Haushalter: Führe diese Männer zu Hause, und schlachte, und richte zu; denn sie sollen zu Mittag mit mir essen. Und der Mann that, wie ihm Joseph gesagt hatte, und führete die Männer in Josephs Haus.}
\end{quote}

89. Ich halte dafür, Joseph habe seinen Bruder Benjamin von ferne gesehen, da sie noch nicht nahe vor ihm gekommen waren. Er
\quellseite{411}
ist aber dazumal an dem öffentlichen Orte gewesen, dahin die Fremden zu kommen pflegten, wenn sie wollten Getreide kaufen; darum heißt er keinen Haushalter, er soll sie in sein Haus führen. Und jetzt bewegt ihn ihr Anblick noch nicht; hernach aber wird er weinen.

2. 18. Sie fürchteten sich aber, daß sie in Josephs Haus geführt würden, und sprachen: Wir sind herein geführt um des Geldes willen, das wir in unsern Säcken vorhin wieder funden haben; daß es auf uns bringe, und fälle ein Urtheil über uns, damit er uns nehme zu eigenen Knechten, sammt unsern Eseln.

90. Die Ursache, daß sie sich fürchten, ist, daß er sie nicht in die gemeine Herberge, da die andern Fremden eingekehrt waren, sondern in sein Haus besonders hat führen heißen. Da ist ihnen eingefallen, daß sie des Allerhöchsten und Schädlichsten zu ihrem Bruder versündigt haben; so er doch gar nichts Böses gedacht, sondern ihnen nur alles Gute gönnte. Also malt Mose ab ein böses und unruhiges Gewissen, welches auch vor einem rauschenden Blatte zittert und erschrickt. Denn wer Gott nicht fürchten will mit guter heiliger Furcht in Gottes Namen, der mag ihn fürchten mit falscher Furcht in des Teufels Namen; ja, er mag sich fürchten vor seinen eigenen Augäpfeln. Denn das ist die rechte Strafe der Sicherheit. Dies ist aber eine rechte eigentliche Beschreibung eines bösen Gewissens, welches sich fürchtet, davor es sich nicht fürchten sollte, und da es sich fürchten soll, fürchtet es sich nicht, wie im 14. Psalm V. 5. steht.

91. Es steht ein schöner herrlicher Spruch Augustins in seinen „Bekenntnissen“, daran man billig immer denken sollte: Jussisti Domine, et factum est ita, ut poena sui sit omnis animus inordinatus, daß ist: Herr, du hast befohlen und ist auch also geschehen, daß ein unordentlich und unruhig Herz seine eigene Strafe sein soll. Wer Gott nicht fürchtet, der mag sich doch fürchten, daß sein eigen Haus, oder seines Nächsten Wand oder Säule auf ihn falle. Warum aber das? Ei, Gott hat es befohlen, also soll es sein. Wer sich vor der Wahrheit nicht fürchtet, der mag sich vor der Lüge fürchten, und mag für den rechten wahren Gott etwa ein Bild oder Götzen, oder aber sonst ein Gespenst ehren und sich davor scheuen.

92. Wie wir vor Zeiten im Pasththume mehr erschrocken sind, wenn wir an Christum, unsern Heiland, gedacht, denn wo des Teufels ist gedacht worden. Denn wir hofften, wir könnten demselben wohl entfliehen, oder aber auf irgend eine Weise von ihm erlöset werden. Denn wir hatten uns dessen also bereden lassen, daß wir meinten, Christus würde uns richten und verdammen. Und darum riefen wir die Jungfrau Maria und andere Heilige an. Denn eine falsche Furcht scheut und fürchtet sich auch vor einer Fliege, und vor Gott selbst, wenn er uns Gutes thut und uns schüßt; ja, er fürchtet sich vor der Speise und Trank, damit er genährt wird. Derhalben treibt endlich ein böses Gewissen den Menschen zur Verzweiflung, daß er mit sich selbst streiten, und wie eine Geißel oder Stockmeister sein muß. Wie die Poeten den Orestes abgemalt haben, daß er von den höllischen Furien getrieben worden, und dergleichen mehr.

93. Also hat nun Joseph befohlen, daß seine Brüder sollten in sein Haus geführt werden, nicht daß er ihnen wollte Schaden zufügen, sondern daß er damit anzeigen möchte, daß ihm angenehm und lieb sei, daß sie seinen Bruder Benjamin mitgebracht hatten. Und er hält sie nun nicht mehr für Kundschafter, sondern für redlich, fromm, aufrichtig, beständig und wahrhaftig. Sie aber, dem allen zuwider, erschrecken greulich, und denken an das Geld, so in den Säcken verborgen war: dasselbe hat sie also gemartert, daß sie in ihrem Herzen nicht anders denken, als daß sie zugleich mit Benjamin werden sterben müssen. So groß ist der Schrecken dieses rauschenden Blattes darum, daß sie nicht in die gemeine Herberge, sondern in Josephs Haus geführt worden.

94. Warum, werden sie gesagt haben, hat man uns nicht auch mit andern Fremden in die Herberge ziehen lassen? Wir sind wahrlich nun verdorben. Er wird ohne Zweifel in seinem Register der Einnahme gefunden haben, daß unser Geld gefehlt hat; darum wird er es von uns erforschen wollen und wird uns dasselbe Geld mit Gewalt abbringen. Da kannst du keine Verheißung sehen, daran sie sich gehalten, sondern nur eitel Schrecken und Plage des Gewissens; so doch Joseph auf nichts weniger denkt, denn daß er sie peinlich berühren, martern oder strafen lassen wollte, wie ihnen davon träumt. Sie sehen an Joseph das brüderliche Herz und guten geneigten Willen nicht, sondern das einige Stück liegt ihnen immer im Herzen; dasselbe bewegt sie sehr, sie überlegen es hin und her, nämlich, daß er von ihnen das Geld wiederfordern und sie des Diebstahls beschuldigen werde. Dazu kommt auch das Seufzen und Klagen, daß sie den Sohn Benjamin vom Vater genommen und sie beide in Gefahr des Lebens geführt. Nun werden wir wahrlich, haben sie gesagt, wie unserm Vater das Herz zugvorgesagt, seine grauen Haare in die Grube bringen. Wir sind um des Geldes willen herein geführt, und ist daran kein Zweifel, er wird uns strafen.

95. Das ist nun das rechte Kreuz und Marter eines unruhigen Gewissens. Denn warum haben sie zuvor Gott nicht gefürchtet mit rechter Furcht? Jetzt haben sie nur die Gedanken, es werde ihnen gar nichts Gutes begegnen, und werden also jämmerlich geplagt, so doch niemand ist, der sie verfolge oder beleidige; sondern, wie Spr. 28, 1. steht: „Der Gottlose fleucht, und niemand jagt ihn.“ Denn das ist die Art und Natur eines schuldigen Gewissens, nämlich, fliehen und erschrecken, auch wenn es allenthalben wohl steht und eitel Glück vorhanden ist, und solch Glück selbst in Gefahr und den Tod verkehren.

96. So ein großer Jammer ist es um ein böses Gewissen, das sich schuldig weiß. Und das ist auch die rechte Frucht der Sünde, die den Menschen gar zunichte macht mit allem seinem Rath, Sinn, Verstand und allen Kräften Leibes und der Seelen.

97. Sie machen aber die Ursachen ihres heftigen Schreckens groß, da sie sagen: „Er wird es auf uns bringen, und ein Urtheil über uns fällen“ etc. Es ist eben so viel geredet, als: Er wird uns plötzlich überfallen, wie die Tyrannen pflegen zu thun; wie 1 Mof. 6, 4. von dem Nephilim steht, das ist, von den Tyrannen, daß es gewaltige und berühmt Leute auf Erden gewesen seien, die großen Hansen, die die Leute drückten. Sie werden nicht darum groß und stark genannt, daß sie groß wären von Leibe, wie es die Sünden auslegen; sondern ihrer Werke und Amts halben, darum daß sie die Armen überfallen und unterdrücken; wie denn solche Gewalt im 10. Psalm beschrieben wird.

98. Derhalben sagen sie nun: Er wird es auf uns bringen, uns zu unterdrücken. Was wird er mit uns machen? Wenn es wohl geräth, so werden wir seine leibeigenen Knechte sein müssen. Nun war aber dazumal nichts weniger in aller Welt zu besorgen, denn solcher Gedanke, welcher ganz fremd, eitel, vergeblich und närrisch gewesen; und hat sie doch derselbe so heftig geplagt.

99. Aber warum ergreifen sie das Wort nicht, daßelbe hätte sie können aufrichten, trösten und stärken, daß sie vor dem Tode und anderer Gefahr sich nicht hätten fürchten dürfen? So gar heftig haben sie sich entsetzt, daß sie an die Verheißung und Predigten ihres Vaters nicht haben denken können, und hatten dieselben auch nun lange versäumt und derselben nicht sonderlich geachtet. Dieweil sie aber nun das Wort verachten, so mögen sie die Züge haben und sich vor sich selbst fürchten: sie mögen sich dichten und träumen lassen von großer Tyrannei und allerlei Strafen, so ihnen begegnen möchten.

100. Denn das sind die Grinnnen, die höllischen Teufel, und die große Pein und Marter eines unruhigen bösen Gewissens wegen der begangenen Sünde, damit es sich selbst vergeblich und ohne alle Ursache beschwert und plagt. Wie ein Spruch eines Cardinals in aller Munde ist, da er gesagt hat: Conscientia est mala bestia, quae facit hominem stare contra se ipsum, das ist: Das Gewissen ist ein böses wildes Thier, welches macht, daß der Mensch mit sich selbst uneins wird. Und die Italiener haben ein gemein Sprichwort, daß sie sagen: Oportet conscientiam occidere, et dicere, quod sit nihil, das ist: Man muß das Gewissen tödten, und sagen, es sei nichts.

101. Ja, wie lange meinst du aber, daß dies Unglück gestillet werden könne? Nicht länger, denn bis das Gesetz kommt: wenn dasselbe anfängt mich anzuklagen und zu schrecken, so werde ich getödtet und verfolge mich selbst; ich streite wider mich und martere mich mit eitel vergeblichen und nichtigen Gedanken. Also geht es zu: über das, daß wir jämmerlich sündigen wider GOTT, plagen wir uns selbst auch, stehen und streiten wider uns selbst, bis ein Bruder zu uns kommt, der uns tröstet und mit Gottes Wort aufrichtet, und sagt: Ei, warum bist du so thöricht? Warum läßt du dir also träumen? Du irrst und fehlest weit: GOtt zürnt mit dir nicht; er hat deine Sünde weggenommen etc. Denn solche Hülfe und Trost eines Bruders thut einem erschrockenen Herzen noth, daß sich selbst also frißt und martert.

\quellseite{412}

\begin{quote}
\textbf{V. 19—22. Warum traten sie zu Josephs Haushalter, und redeten mit ihm vor der Hausthür, und sprachen: Mein Herr, wir sind vorhin herab gezogen, Speise zu kaufen. Und da wir in die Herberge kamen, und unsere Säcke aufthaten, siehe, da war eines jeglichen Geld oben in seinem Sack mit völligem Gewicht; darum haben wir’s wieder mit uns gebracht. Haben auch ander Geld mit uns herab gebracht, Speise zu kaufen; wir wissen aber nicht, wer uns unser Geld in unsere Säcke gesteckt hat.}
\end{quote}

102. Sie gehen nicht in Josephs Haus, sondern bleiben vor der Hausthür stehen, vielleicht in der Absicht und Hoffnung, daß sie könnten davonlaufen. Und daraus ist wohl zu sehen, wie erschrocken und bekümmert sie gewesen sind. Es hat aber der unschuldige Benjamin zugleich auch mit ihnen erschrecken müssen; und das Geld, welches doch ein Zeichen der größten Liebe und brüderlicher Freundschaft war, hat ihnen diesen großen Schmerz und Traurigkeit gemacht. Denn ein böses Gewissen legt auch das Allerbeste zum Uebelsten aus, und was seiner Art und Natur halben ganz süß und lieblich ist, das verwandelt es in eitel Wermuth und bittere Galle.

103. Ja, die ganze Creatur, so den Menschen zu Nutz und Lust geschaffen ist, scheint ihnen gar zuwider zu sein und die Pestilenz und den Tod zu drohen; wie im Buch der Richter Cap. 5, 20 steht: „Vom Himmel ward wider sie gestritten, die Sterne in ihren Läuften stritten wider Sisera“ etc. Also sind die Kriegsleute erschrocken, daß sie sich vor ihren eigenen Rossen und Wehre, ja, auch vor den Sternen, Himmel und Sonne gefürchtet haben.

104. Eben also ist ein böses Gewissen auch gegen GOtt gesinnt. Denn wenn er uns anlacht, und günstig und gnädig ist, fürchtet es sich doch vor seinem Zorn, und hält dafür, daß er ihm zuwider sei und mit ihm zürne. Derhalben ist das ein trefflich Wort eines Bruders in unserm Kloster gewesen, der da zu denen, so betrübt und erschrocken waren, zu sagen pflegte: GOtt zürnt nicht mit dir, zürne du nicht mit dir selbst. ihm; denn du zürnst viel mehr mit ihm denn er mit dir. Und die Herzen, so aus Furcht erschrocken sind, fühlen in Wahrheit nicht anders denn also. Denn ein unrein Herz kann GOtt nicht sehen; wie Christus sagt Matth. 5, 8.: „Selig sind, die reines Herzens sind; denn sie werden GOtt schauen.“ Darum sind das die allerseligsten Leute, die kein rein Herz haben; denn sie schauen GOtt nicht, sondern meinen, sie sehen oder hören den Teufel, wenn sie gleich an GOtt denken, von ihm reden oder hören.

105. Und kann einem bösen Gewissen zur Zeit der Anfechtung nichts so Gutes und Lustiges vorkommen oder angeboten werden, daß es nicht umkehre, und sich selbst eitel Furcht, Pein und Schrecken daraus mache, und kann nichts anders, denn nur wider uns selbst zeugen und mit uns streiten; und wenn es endlich darnieder liegt und gar verzagt ist, vermag es die ganze Welt nicht wiederum aufzurichten. Denn es ist der Seele Tod; und die Gewissen aufrichten und erwecken ist nichts Anderes als Todte auferwecken. Derhalben macht die Kirche durch das Predigtamt oder das mündliche Wort, das sie hat und braucht, viel mehr Leute lebendig, als Christus selbst ist in seinem Amt gethan; wie er sagt Joh. 14, 12.: „Wer an mich glaubet, der wird die Werke auch thun, die ich thue, und wird größere denn diese thun.“ Darum geschehen täglich viel größere Wunderwerke, die wir nicht sehen und darauf die Welt gar keine Acht gibt.

106. Das Gewissen ist ein viel größer Ding denn Himmel und Erde, welches durch die Sünde getödtet und durch das Wort Christi wiederum lebendig gemacht wird. Derhalben sollen wir uns alle insonderheit befleißigen, daß wir von diesem sehr bösen und höllischen wilden Thier errettet werden mögen. Denn wo dasselbe Thier nicht wäre, so hätte die Hölle kein Feuer, oder auch gar keine Pein. Dies wilde Thier aber zündet an und stärkt den Tod und die Hölle und rüstet die ganze Creatur wider uns. Denn da zürnt Alles, und ist traurig, störrig, scheußlich, und ist wider uns; nicht der Creatur halben, welche gut ist, und uns weder droht noch schadet; sondern es ist unsere eigene Schuld, die wir also erschrocken sind und vor den Creaturen fliehen.

107. Wiederum aber ein fröhlich und sicher Gewissen bricht hindurch, siegt und verachtet hie Tod und Teufel; wie solches sein herrlich beschrieben wird im 112. Psalm V. 7. 8.: „Wenn eine Plage kommen will, so fürchtet sich der Fromme nicht; sein Herz hoffet unverzagt auf den HErrn. Sein Herz ist getrost und fürchtet sich nicht“ etc. Es zürne oder lache der Teufel und die Welt, so fragt es nichts darnach: will er nicht lachen, so zürne er. Ein Herz aber, das sich schuldig weiß, fürchtet sich auch vor einem guten Gerücht. Denn man pflegt zu sagen: Conscientia mille testes: Das Gewissen ist wie tausend Zeugen, ja, wie tausend Heerlager, so wider uns sind. Wie aber solch Schrecken des Gewissens überwunden werden möge, davon lehrt man an andern Orten in der Lehre des Evangeliums.

\subsection{II}

\begin{quote}
\textbf{V. 23. Er aber sprach: Gehabt euch wohl, fürchtet euch nicht. Euer GOtt und eures Vaters GOtt hat euch einen Schatz gegeben in eure Säcke. Euer Geld ist mir worden. Und er führete Simeon zu ihnen heraus.}
\end{quote}

108. Dieser Haushalter hat wohl gewußt, was Josephs Plan und Vorhaben mit seinen Brüdern gewesen ist. Denn Joseph hätte dies nicht ausrichten können, daß er sich also gestellt und anders gebärdet, denn ihm zu Muth war, wo er nicht einen treuen Diener gehabt; welchen er vornehmlich auch gebraucht haben wird, das Getreide zu verkaufen. Und vielleicht hat derselbe auch das Geld in die Säcke gelegt; sonst hätte er ihnen diese Antwort nicht geben können. Darum sagt er nun: „Euer Geld ist mir worden“, wiewohl er es doch nicht empfangen hat; sondern er sagt es darum, daß er sie damit stärken und erschrecken will. Und es scheint, als habe ihn sein Herr unterrichtet in der Lehre der Gottseligkeit; denn er sagt: „Euer GOtt und eures Vaters GOtt hat euch einen Schatz gegeben“, das ist, der GOtt, den ihr und euer Vater ehrt. Ihr dürft euch nicht bekümmern eures Geldes halben: ich habe es empfangen, „es ist mir worden“. Und er sagt recht daran. Aber davon schwindet er stille, daß er es ihnen wieder gegeben und in ihre Säcke gelegt habe. Also folgt dieser getreue Knecht den Sitten und der Gottseligkeit seines Herrn fein nach, wiewohl er ein Egypter war. Er sagte nicht: Ich habe es, oder: Ich habe es nicht in die Säcke gelegt; sondern sagt: Es kam zu mir, ich kriegte es etc.

109. Es hört aber die Furcht noch nicht auf; ja, es wird der Argwohn noch immer größer, und überkommt sie nun eine neue Bekümmerniß, nämlich, wie doch dieser Haushalter ihr Geld bekommen haben möge, oder wer daselbliche in die Säcke gelegt habe.

110. Das hebräische Wort, mathmon, Schatz, hat seinen Namen vom Verbergen; darum daß er vergraben und verborgen wird in der Erde. Darum ist es keine neue Weise, daß man das Geld vergräbt; welches jetzt zur Zeit die Bauern in ihrer Unsinnigkeit und auf die seltsamsten Weisen thun; sondern ist bei den Alten auch gebräuchlich gewesen. Und ich halte dafür, daß viel Geld, welches also vergraben wird, verloren werde. Und es geht mit dem Gelde eben also zu, wie mit dem Menschen; denn es wird auch in die Erde begraben, im Meer versenkt, wenn Schiffbruch geschehen, eben wie das menschliche Geschlecht umzukommen pflegt. Hiob gedenkt auch desselben Gebrauches, da er im 3. Cap. V. 21. sagt: „Die des Todes warten, und kommt nicht, und grüben ihn wohl aus dem Verborgenen“; item Jer. 41, 8.: „Lieber, tödte uns nicht; wir haben Schätze im Acker liegen von Weizen, Gerste, Oel und Honig“ etc.

111. Darnach führt der Haushalter Simeon aus dem Gefängniß heraus, auf daß er ihre Herzen desto mehr aufrichten und stärken möge. Ich zweifle auch nicht, es wird ein leiblich Gefängniß gewesen sein, und nicht ein solcher Ort, darein man die Missethäter und die den Tod verdient, gesetzt und gebunden hat. Das ist aber ein großes Anzeichen gewesen sonderlicher Gnade und Freundlichkeit, daß ihr Bruder alsbald aus dem Gefängnisse vor ihre Augen gebracht wird; und sie können doch ihre Sorge, Bekümmerniß, und alle Furcht noch nicht gar ablegen und fahren lassen.

112. Also ist dies ein sehr schönes Exempel eines bösen Gewissens, nämlich, wie dasselbe gegen GOtt gesinnet sei; und wiederum auch, wie groß die Gunst und Gnade GOttes sei, so er gegen das erschrockene Gewissen trägt. Denn wenn er uns versucht, läßt er mancherlei Hindernisse vorfallen, auf daß wir ihm nicht stracks unter Augen sehen und zu dem Gnadenstuhle kommen können; gleich als wenn Einer zur Kurzweil mit einem Würmlein spielt, und ihm, wo es auf der Erde kriecht, etwa ein Stücklein oder Sträuchlein vorwirft, daß es nicht könne fortkriechen, dahin es gerne wollte, sondern muß sich mancherlei Weise hin und her wenden und ziehen, und an allen Orten versuchen, wie es doch endlich davon kommen möchte. Aber dies Spiel göttlicher Gnade und seines Wohlgefallens verstehen wir vom Anfange nicht, und die Wohlthaten und Gnade selbst, so uns begegnet und vor Augen gestellt wird, deuten wir zu unserm Schrecken und Verderben.

\quellseite{413}

113. Und widerfährt uns eben dasselbe, daß wir vorzeiten, da ich ein kleiner Knabe war, und meinen Gesellen, mit denen ich die Parteten gesammelt, davon wir uns bei unserm Studium erhalten möchten, auch begegnet ist. Denn da wir zu der Zeit, da in der Kirche das Fest der Geburt Christi gehalten wird, auf den Dörfern von einem Hause zum andern umhergegangen und in vier Stimmen die gewöhnlichen Psalmen vom Kindlein Jesu, geboren zu Bethlehem, zu singen pflegten, geschah es ohngefähr, daß wir vor eines Bauern Hof, so an einem Orte allein und am Ende des Dorfes gelegen war, kamen, und da uns der Bauer singen hörte, kam er heraus, und fragte mit groben bäurischen Worten, wo wir wären, und sagte: Wo seid ihr Buben? und brachte zugleich etliche Würste mit, die er uns geben wollte. Wir aber erschracken vor den Worten sehr, daß wir alle von einander wegliefen; wiewohl wir keine rechte Ursache wußten, warum wir hätten erschrecken mögen, und der Bauer uns die Würste mit gutem geneigtem Willen darreichte und geben wollte, außer, daß vielleicht unsere Herzen furchtsam gewesen vom täglichen Drohen und Tyrannei, so zu der Zeit die Schulmeister mit den armen Schülern zu üben pflegten, und so viel leichter von solchem plötzlichen Schrecken scheu geworden sind. Endlich aber, da wir auf der Flucht waren, rief uns der Bauer wieder, und legten die Furcht ab und liefen herzu, und empfingen von ihm die Parteten, so er uns reichte.

114. Gleicher Weise pflegen wir auch zu zittern und zu fliehen, wenn unser Gewissen schuldig und erschrocken ist: da fürchten wir uns auch vor einer Brotrinde, und vor denen, die unsere Freunde sind und uns alles Gute gönnen.

115. Derhalben sollen wir uns mit allem Fleiß vor Sünden hüten, die das Gewissen verwunden, und sollen Gott lernen vertrauen, und mit feinen, stillen, friedsamen, gläubigen Herzen das annehmen, was uns von Gott zugeschickt wird, es sei gut oder böse; wie der 112. Psalm V. 7. sagt: „Wenn eine Plage kommen will, so fürchtet sich der Fromme nicht; sein Herz hoffet unverzagt auf den HErrn“; item Spr. 28, 1.: „Der Gerechte ist getrost, wie ein junger Löwe.“

116. Denn ein böses Gewissen kann auf keine Weise zufrieden sein oder sich zu Gott wenden, es sei denn gar erstorben; sonst flieht es immer vor Gott, vor welchem es doch weder fliehen noch erschrecken sollte. Sintemal Gott ein solches Gut ist, zu dem alle Verzagten, und die sonst von allen Creaturen verlassen, trostlos und ohne Hülfe sind, allein ihre Zuflucht haben sollen. Aber ein armes Herz, daß sich schuldig weiß, würde sich eher in tausend Gestalten verändern, und viel eher durch Steine, Feuer und Erzberge, ja, endlich zu dem Teufel selbst laufen, ehe es sich zu Gott wendete und stracks zu ihm liefe.

117. Ist das nicht ein großer erbärmlicher Jammer, daß wir uns vor dem fürchten und scheuen, von welchem die heilige Schrift so oft sagt, daß er unsere Zuflucht für und für sei? ja, wenn Vater und Mutter uns verlassen und verstoßen würden, so würde er doch die Verlassenen aufnehmen; wie der 27. Psalm V. 10. sagt: „Mein Vater und meine Mutter verlassen mich; aber der HErr nimmt mich auf.“ Derhalben, wenn auch alle Creaturen nichts Anderes sagten, denn daß du verloren, vermaledeit und verdammt wärest, so sollst du doch mit dem 73. Psalm V. 25, 26. sagen: „Wenn ich nur dich habe, so frage ich nichts nach Himmel und Erden. Wenn mir gleich Leib und Seele verschmachtet, so bist du doch, GOtt, allezeit meines Herzens Trost und mein Theil.“ Was frage ich darnach, wenn ich nur einen gnädigen Gott habe!

118. Lieber, siehe doch, was die Heiden thun, die den rechten wahren Gott nicht kennen, und laufen doch in allen ihren Nöthen zu dem, das sie für Gott halten. Sie rufen Mars, Diana und Jupiter an etc., mit welchem Anrufen und Götzendienst sie ja anzeigen, daß Gott ein solcher Gott sein müsse, zu dem man in Gefahr, Angst und Noth fliehen soll; wiewohl sie den Namen und die Ehre Gottes den Götzen und Creaturen fälschlich zuschreiben; wie Paulus sagt Röm. 1, 21. ff.

119. Derhalben ist diese Erkenntniß in aller Menschen Herz gepflanzt, daß die Betrübten, so in Nöthen sind, ihre Zuflucht zu GOtt haben sollen. Das ist der Obersatz im Syllogismus, wie die Dialektiker in den Schulen zu reden pflegen, und kann in keinen Zweifel gestellt werden, denn nur von den Sadducäern und Epicuräern. Aber der Untersatz ist falsch, wenn man sagen will: Dieser Götze ist GOtt. Da läßt sich unsere Natur sehen, daß sie verderbt ist. Daher kommt es, daß wir des rechten Gottes fehlen und denselben wegen solcher verderbten Natur nicht recht antreffen können. So wissen nun alle Menschen, daß Gott unsere Zuflucht ist, und rufen ihn um Hülfe und Trost an. Aber zu dem Gott, der sich in seinem Worte namhaft macht und geoffenbart hat, will niemand mit rechtem Vertrauen treten. Das Gewissen ist erschrocken und flieht zurück; denn es ist unrein und kennt Gott nicht recht; darum schauet es GOtt auch nicht. Aber „selig sind, die reines Herzens sind; denn sie werden GOtt schauen“, Matth. 5, 8. Die springen nicht zurück, sondern stehen fest im Glauben und rechter Anrufung. Wie denn dieselbige Stärke und festes Vertrauen im 112. Psalm V. 8. gerühmt wird: „Das Herz des Frommen ist getrost, und fürchtet sich nicht“ etc., und in Wahrheit ist Gott bei solchen frommen Herzen und hilft ihnen, wenn sie in Nöthen sind.

\subsection{III}

\begin{quote}
\textbf{V. 24. 25. Und führte sie in Josephs Haus, daß sie Wasser, daß sie ihre Füße wuschen, und gab ihren Eseln Futter. Sie aber bereiteten das Geschenk zu, bis daß Joseph kam auf den Mittag; denn sie hatten gehöret, daß sie daselbst das Brod essen sollten.}
\end{quote}

120. Es ist ein alter Gebrauch gewesen, daß man den Gästen die Füße gewaschen hat, und wird auch in der evangelischen Historie gerühmt, und ist von Christo selbst auch also gehalten worden, Joh. 13, 5. ff., da er den Jüngern die Füße gewaschen. Und Paulus sagt 1 Tim. 5, 10. von der Wittwe, „so sie Kinder aufgezogen hat, so sie gastfrei gewesen ist, so sie der Heiligen Füße gewaschen hat.“ Es ist im jüdischen Volk eine sehr schöne Gewohnheit gewesen; zu unsern Zeiten aber ist es gar abgekommen, und ist allein die Ceremonie noch im Gebrauch in des Papsts Kirche geblieben, welche die Könige, Bischöfe und Aebte am grünen Donnerstage herrlich zu halten pflegen, daß sie eine Schürze umgürten und zwölf Armen die Füße waschen, und denselben darnach zu Tische dienen. Denn damit haben sie also dem Exempel Christi nachfolgen wollen, auf daß sie sich desselben Werks halben rühmen und damit stolzieren möchten. Unterdeß aber tödten sie die rechten Christen, vertheidigen ihre falsche Lehre, verfolgen das Wort, vergießen unschuldig Blut, und schmücken sich derweil mit dieser Ceremonie und bedecken damit ihre grausame Tyrannei.

121. Vorzeiten wuschen die Mönche in den Klöstern den fremden Brüdern die Füße, gaben ihnen dazu neue Messer und Schuhe, desgleichen führten sie dieselben in besondere Herberge, thaten ihnen daselbst gütlich und warteten ihrer wohl. Das war recht dem alten löblichen Gebrauch nachgefolgt, welchen man jetzt ohne große Beschwerung, Mühe und Arbeit nicht würde wiederum aufrichten können. Etliche Fürsten und die Papisten waschen das ganze Jahr über den Armen die Füße nicht; wenn sie aber dieselben auf besondere bestimmte Zeit waschen, so nimmt man Leute dazu, welche zuvor gewaschen sind. Das ist eitel lauter Heuchelei, die sich mit dem Stempel und Willen Christi gar nicht reimt. Fromme gottselige Pfarrherren und arme Schüler, und sonst andere Fremdlinge, so gute Zeugnisse und Empfehlungen von ehrlichen Leuten bringen, deren achtet man nicht; ausgenommen, daß noch etliche wenige gemeine Leute sind, die solche Werke verstehen, und wissen, daß sie Gott wohl gefallen, dieselben thun Gutes, und erzeigen Wohlthaten, wem sie sollen und wo sie können.

122. Man soll aber gleichwohl nicht so gar leichtlich jedermann aufnehmen und beherbergen, es sei denn, daß er auch von Andern gute Zeugnisse habe. Wie die zu Corinth den Apollo nicht ohne Briefe aufnehmen, auf daß sie wüßten, wen sie aufnehmen und wem sie Gutes thäten. Denn die Bosheit der Landstreicher und anderer Betrüger wird dermaleinst fromme gottselige Leute zwingen, daß sie mit ihrer Mildigkeit und Wohlthaten gegen die Armen müssen Maß halten, dieweil so viel Heuchler umherlaufen, die sich gar ans Betteln gewöhnt haben, welche unter dem Namen und Scheine, als wären sie recht elend, den rechten Armen die Wohlthaten und Almosen vor dem Maul hinweg nehmen. Denn diejenigen, so des Waschens und anderer Leute milder Almosen recht wohl bedürfen, werden oft ihrer eigenen Schamhaftigkeit und Bescheidenheit halben abgeschreckt, daß sie sich scheuen, die Gemeinde und sonst gottselige Leute um Hülfe und Steuer zu bitten; sie schweigen stille, und fressen ihre Armuth, damit sie beschweret sind, in sich; laufen nicht umher im Lande, sondern leiden ihre Armuth, Jammer und Elend mit Geduld; wiewohl sie oftmals das Brod nicht zu essen haben.

\quellseite{414}

123. Aber irgend ein Betrüger kann sich leichtlich eindringen, und nimmt das hinweg, was man den rechten Armen geben sollte, und verspielt und verschlemmt es noch wohl dazu, und spottet darnach unser, daß wir ihm so leichtlich geglaubt und gegeben haben, was er begehrt hat. Und in solchen Dingen thut die Obrigkeit auch nicht, was ihr Amtshalben zu thun gebühret; denn sie sollte solche Mörder und Diebe aus ihrem Lande und Städten treiben. Und solchen unbekannten Bettlern sollte man nichts geben, die ohne Zeugnüß frommer gottseliger Männer daher gelaufen kommen; wiewohl derselbigen auch täglich über die Maßen viel werden, so von andern ehrlichen Leuten Brief und Zeugnüß mitbringen. Und was sie nur für Schaden gelitten, kommen sie und begehren unser Almosen; so wir doch selbst und unsere Mitbürger oder Nachbarn auch nicht so viel Geld und Gut übrig haben, daß wir allen Armen, so zu uns kommen, ihre Armuth und Mangel erstatten und sie dessen überheben könnten. Und sind wahrlich bei uns auch nicht wenig arme Bürger, die man ja vornehmlich vor andern ernähren und unterhalten sollte.

124. Aber es ist verloren und umsonst, wenn wir schon lange solche Vermahnung thun; dieweil jetziger Zeit alle Zucht aufgehoben und die Leute so gar muthwillig und frech sind. Ich halte zwar nicht dafür, daß Alle in Egypten solche Mildigkeit und Freundlichkeit insgemein geübt haben, wie die Juden und Joseph gethan; und man kann auch nicht wissen, ob es also ein gemeiner Gebrauch über ganz Egyptenland gewesen sei, daß sie den fremden Gästen die Füße gewaschen haben; oder, es ist kein Zweifel, sie werden damit einen Unterschied gehalten haben. Denn Joseph hat fleißig auf die Personen Acht gegeben, hat sie alle insonderheit gefragt, wer sie wären und von wann sie hergekommen.

125. Ist es nun also ein gemeiner Gebrauch im Lande gewesen, so hat man denselben Dienst leichtlich leisten können; aber doch nur denen, davon man sich zuvor erkundigt, wer sie wären; wie wir droben, 1 Mos. 24, 32., gehört, daß Rebekka Abrahams Knecht in ihres Vaters Haus geführt, und ihm Wasser gegeben, seine Füße zu waschen, und dazu seinen Kameelen Stroh und Futter; denn sie hat gewußt, daß er vom Geschlecht Abrahams dahin gesandt war etc. Und zu der Zeit sind der Landstreicher und losen Buben so viel nicht gewesen, wie derselben jetzt über die ganze Welt viel sind.

126. Es gebührete aber den Fürsten und Staaten, daß sie im Land und Städten ein fleißiges Aufsehen hätten auf solche Betrügerei und Räuberei der vielen Bettler, und sollten solche faule Schelme, die nicht arbeiten und doch immer fressen wollen, ernstlich strafen. Denn dieselben berauben uns, und nehmen den Armen, die bei uns wohnen, das Brod, so ihnen gottselige Leute geben würden, vor dem Maul hinweg. Denn fromme gottselige Leute wissen wohl, daß GOtt haben will, daß wir in der Gemeinde GOttes mild und wohlthätig gegen arme Leute sein sollen; darum geben sie auch den Armen Herberge, Essen und Trinken, Kleider und Schuhe. Das Fußwaschen aber, so bei den Papisten gebräuchlich, ist eine lautere Heuchelei, damit sie weder GOttes Ehre noch des Nächsten Nutzen suchen, sondern rühmen nur dasselbe selbst erwählte Werk.

\section[Vierter Theil]{Vierter Theil. Auf welche Art Joseph seine Brüder, besonders Benjamin, empfängt}

\begin{quote}
\textbf{V. 26—28. Da nun Joseph zum Hause einging, brachten sie ihm zu Hause das Geschenk in ihren Händen, und fielen vor ihm nieder zur Erden. Er aber grüßte sie freundlich, und sprach: Gehet es eurem Vater, dem Alten, wohl, von dem ihr mir saget? Lebet er noch? Sie antworten: Es gehet deinem Knechte, unserm Vater, wohl, und lebet noch. Und neigten sich, und fielen vor ihm nieder.}
\end{quote}

127. Das Wort „anbeten“ heißt, die Kniee beugen, oder, einem Andern zu Füße fallen. Mose sagt, sie haben sich zur Erden geneigt, das ist, sie haben vor Joseph einen Fußfall gethan, daß sie mit den Knieen die Erde berührt haben. Damit ist nun der Traum erfüllt worden, welchen er droben, 1 Mos. 37, 6. 7., seinen Brüdern erzählt hat. Joseph aber hat sie freundlich gegrüßt; interrogat ipsis pacem, wünscht ihnen Frieden, wie es nach dem Hebräischen lautet. Also steht im 122. Psalm V. 6.: Rogate, quae ad pacem sunt, Jerusalem, das ist: „Wünschet Jerusalem Glück“ und Heil. Wie wir mit unserm Gruß auch Andern pflegen Glück zu wünschen, wenn wir sagen: GOtt gebe euch Gnade! GOtt gebe euch einen guten Morgen! guten Abend! GOtt sei euch gnädig! segne euch! gebe euch Glück und Heil! Denn das alles heißt bei den Hebräern das Wort pax, Friede.

128. Also verharrt Joseph noch immer dabei, daß er seinen guten geneigten Willen und sonderliche große Liebe gegen seine Brüder beweiset, auf daß sie es bei sich selbst für gewiß halten, daß er ihnen günstig, gnädig und wohl gewogen sei; und hält sich doch also, wie einem Regenten gebühret, welcher immer über gute Ordnung und Zucht halten muß. Und steht, auch der Obrigkeit vornehmlich wohl an, daß sie daneben auch barmherzig und gnädig sei.

129. Darauf fragt nun Joseph: Ob es auch ihrem alten Vater noch wohl gehe? Denn also lauten seine Worte: „Gehet es eurem Vater, dem Alten, noch wohl? Lebet er noch?“ Darauf antworten sie: „Es gehet ihm wohl“ etc. Und daraus kann man abnehmen, was die rechte Bedeutung dieses Wortes pax, Friede, sei. Wie droben im 37. Cap. V. 14., da Jacob seinen Sohn Joseph zu seinen Brüdern sendet, er zu ihm sagt: „Gehe hin, und siehe, ob es wohl stehe um deine Brüder und um das Vieh“, das ist, ob es ihnen wohl gehe, wie wir auf unser Deutsch sagen. Es hatte zwar Joseph zuvor auch gehört, daß der Vater noch am Leben und gesund wäre; damit er aber dessen desto gewisser würde, fragt er abermals darnach. Und sie neigen sich nun zum andernmal vor ihm oder bückten sich; wie wir zu unsern Zeiten den Hut abziehen oder die Kniee zu beugen pflegen, Andern damit Ehre und gebührliche Achtung zu erzeigen.

\begin{quote}
\textbf{V. 29. Und er hub seine Augen auf, und sahe seinen Bruder Benjamin, seiner Mutter Sohn, und sprach: Ist das euer jüngster Bruder, da ihr mir von saget? Und sprach weiter: GOtt sei dir gnädig, mein Sohn.}
\end{quote}

130. Droben haben wir gehört, wie Joseph seinen Bruder Benjamin von ferne gesehen habe. Hier sieht er ihn nun etwas näher an in seinem Hause, und Benjamin wird genannt seiner Mutter Sohn, darum daß er allein Josephs leiblicher Bruder war; wie die andern Brüder mit ihm von Einem Vater geboren waren. Hier heben sich nun die rechten Affekten und natürlichen Bewegungen an, und da finden sich auch nun die Thränen, da er seines Bruders recht ansichtig wird und da auch des Vaters gedacht wird. „GOtt“, sagt er, „sei dir“ günstig und „gnädig, mein Sohn“ etc.

131. Diese Worte werden ihm sehr zu Herzen gegangen sein. Denn er wird nicht allein dadurch beweget, daß er seinen Bruder sieht, sondern auch, da er der Mutter, des Vaters und alles deß, was daheim geschehen war, gedenkt. Da ist ihm eingefallen, daß seine freundliche liebe Mutter in der Geburt um dieses ihres Sohnes willen kümmerlich gestorben, im Jahre vorher, ehe er nach Egypten war verkauft worden. Da wird er zugleich auch bedacht haben, wie viel und wie großen Jammer seine Eltern und er selbst auch hatten ausstehen müssen ihr ganzes Leben lang, vornehmlich aber die zwei Jahre über, darin die Mutter gestorben, und daß er im Ausgang desselben Jahres, darin Benjamin geboren, den Ismaeliten war verkauft worden.

132. Und ist daran kein Zweifel, dies alles werde ihm also eingefallen und vor Augen gekommen sein, da er seinen Bruder gesehen; und sein Herz ist in ihm entbrannt von zarter Liebe und herzlicher Neigung gegen Benjamin und seine Eltern. Darum sagt er: „GOtt sei dir gnädig, mein Sohn.“ Als wollte er sagen: Ach, wie hat dich deine liebe Mutter mit so großem Schmerz geboren! Wie viel große Beschwerden, Sorge und Angst haben doch mein Vater und ich tragen müssen, bis ich dazu gekommen hin, daß ich euch wieder gesehen und mit euch Gemeinschaft haben möge! Ach Gott, wie sauer ist es deiner Mutter geworden, und mir und deinem Vater, ehe ich dich wieder sehe! Da wird nun sein Herz eitel Wasser werden.

\quellseite{415}

133. Lieber, siehe doch, wie kräftig die Natur sei, welche, je besser und rechtschaffener sie ist, desto trefflichere und brünstigere Neigungen hat. Und die Gnade und der Heilige Geist hebt dieselbigen nicht auf, verderbt sie auch nicht, wie den Mönchen davon geträumt hat; sondern wo die Natur verderbt ist, heilt sie der Heilige Geist und richtet sie fein auf, bringt sie wiederum in ihren rechten Stand. Denn dasselbe sieht man hier an diesem Exempel Josephs wohl, welchen die natürlichen Neigungen so gar überwinden, daß beide sein Herz und Augen eitel Wasser werden und nicht anders denn nur weinen können. Er ist gar nicht bei sich selbst, und er verstummt gleichsam dazu auch vor großem herzlichem Mitleid, so er mit seinem Bruder und Eltern gehabt.

134. Darum sagt Mose nicht vergeblich, daß er geeit habe: er kann die Ehre nicht ausarten, die sie ihm anthun. Nachdem er die Worte, darin er Benjamin gesegnet und ihm Glück gewünscht, ausgeredet hat, eilt er alsbald in seine Kammer. Sie wissen nicht, was das Eilen bedeuten möge und daß er sein Angesicht so plötzlich von ihnen wendet; ja, sie deuten es heimlich für ein Zeichen der Ungnade und Feindschaft. Aus den Worten urtheilen sie, er müsse ein freundlicher gnädiger Herr sein; was aber das bedeuten möge, daß er so plötzlich aufhört zu reden, das können sie nicht verstehen. Er ist aber um keiner andern Ursache willen in seine Kammer gegangen, denn daß er hat müssen ausweinen.

135. Und dies sind wahrlich treffliche Historien der heiligen Väter, die da zeigen, daß die Natur nicht geändert, sondern von Gott erhalten werde. Wie das tolle unsinnige Volk der Mönche gelehrt hat, daß man die Eltern, Weib und Welt verlassen solle und in ein Kloster gehen. Was will aber alsdann die Gemeinschaft sein, so ja billig die Menschen mit einander halten sollen? oder was will doch dann das ganze Leben der Menschen sein? Das ist wahrlich keine rechte Heiligkeit oder Gottesdienst gewesen, wie sie gesagt haben, sondern eine teuflische Unsinnigkeit. Wenn du aber durch Noth und Gottes Gebot gezwungen wirst, deine Eltern oder Weib etc. zu verlassen, als, da Jakob seine Söhne, Joseph und Benjamin, hat müssen verlassen; oder, da Abraham geboten wird, seinen Sohn Isaak zu schlachten; da unser Herr Gott in der ersten Tafel verbietet, man soll ihm weder Vater noch Mutter vorziehen, sondern ihn über Alles fürchten und lieben: alsdann soll man dem Worte Gottes gehorsam sein, daß 5 Mos. 33, 9. steht: „Wer zu seinem Vater und zu seiner Mutter spricht: Ich sehe ihn nicht; und zu seinem Bruder: Ich kenne ihn nicht; und zu seinem Sohn: Ich weiß nicht; die halten deine Rede, und bewahren deinen Bund“ etc. Aber in der andern Tafel, je besser, vortrefflicher und heiliger die Natur ist, je mehr hält sie die Eltern in Ehren und liebt Weib, Kinder und Brüder.

136. Aber die Mönche verachten dies alles als eitel fleischlich Ding und das von Gott verdammt sei, und treten es gleichsam mit Füßen. Sie sind in der ersten Tafel abgöttisch und Gotteslästerer, ehren den Teufel und dienen ihm unter dem Namen Gottes. Und in der andern Tafel lassen sie die Werke liegen, so Gott geboten hat, und folgen ihrem eigenen Gutdünken und Menschensatzungen. Wie Ahas auch den Gottesdienst der ersten Tafel verachtet hat und liegen lassen, und dient Gott mit dem Opfer seiner Kinder. Denn damit gedachte er ein köstlich Werk zu thun, daß er seine Söhne nach dem Exempel des höchsten Patriarchen Abraham opferte. Aber wer hat es ihm geheißen oder befohlen, solches nachzuthun?

137. Derhalben will Gott in der andern Tafel haben, daß die Kinder ihre Eltern herzlich lieb haben und wiederum auch die Eltern ihre Kinder. Und das ist traun kein weltlich oder fleischlich Leben, sondern das allergeis(t)lichste um des vierten Gebots willen: „Du sollst deinen Vater und deine Mutter ehren“; welches das nächste Gebot nach der ersten Tafel ist. Desgleichen gebietet es auch den Unterthanen, daß sie der Obrigkeit sollen gehorsam sein, und daß die Obrigkeit wiederum ihre Unterthanen auch also regieren soll, daß sie dabei fortälttig sei, wie Paulus sagt Röm. 12, 8. Darum soll man vor den Aeltern und Werken des vierten Gebots nicht fliehen, sondern sie mit höchstem Fleiß ausrichten und verwalten. Aber die Welt thut, wie ihre Art ist, und verkehrt dies alles, als wäre sie gar unsinnig.

138. Darum soll man solche Stellen in der Schrift fleißig merken, darin die natürlichen Neigungen und Affecten in den Heiligen gerühmt werden, dadurch sie dermaßen sind bewegt worden, daß sie auch zarter gewesen zu sein scheinen, denn sonst andere Weltmenschen und Heiden sind. Denn sie sind etma so gar betrübt und bekümmert, haben ihre Kinder so herzlich lieb, daß sie sich damit auch schwächer erzeigen, denn irgend ein Heide hätte thun mögen. Ja, man sieht auch an ihnen, daß sie auf beiden Seiten, das ist, im Glück und Unglück, sehr heftig und hitzig sind. Daher soll man nun ziehen die Lehre Pauli Röm. 12, 15.: „Freuet euch mit den Fröhlichen, und weinet mit den Weinenden“ etc.

139. Also wird man widerlegen können die Apathie, das ist, die unnatürliche oder unmenschliche Gefühllosigkeit der Mönche und Einsiedler, um welcher willen vorzeiten ihr einsam Leben, das sie geführt, vom gemeinen Manne ist gerühmt worden; und sie haben viel Leute mit sich in den Irrthum und das Verderben geführt. Und hat sich zu unsern Zeiten Münzer unterstanden, dieselbige unmenschliche Gefühllosigkeit wiederum in die Kirche einzuführen und in Schwang zu bringen.

140. In dem Buche, das sie nennen „Leben der Väter“, werden viel wunderliche scheußliche Exempel der Mönche erzählt, welche beides ihre Unsinnigkeit und großen Unverstand der vorigen Zeit genugsam anzeigen. Als, daß Exempel der zwei Brüder, zu welchen ihre Mutter gekommen ist, daß sie ihre Söhne gerne sehen wollte, dieweil sie dieselbigen in langer Zeit nicht gesehen hatte; sie aber haben die Thüre zugethan und sich verborgen. Und da die Mutter sie gebeten, daß sie doch möchte zu ihnen kommen, haben sie sie heißen weggehen und gesagt: Sie sollte dieweil Geduld tragen, daß sie also ihrer Söhne müßte beraubt sein, bis sie dieselben im ewigen Leben sehen würde. Wenn ich an dem Orte ein Herr oder Regent gewesen wäre, wollte ich dieselben Klosterzellen mit Feuer haben verbrennen lassen. Warum haben die heiligen Leute nicht daran gedacht, daß sie ihrer Mutter Früchte gesogen, und daß sie Menschen wären, die von dieser Mutter in diese Welt geboren waren? Aber solcher Werke können sich die Mönche sehr verwundern und dieselben hoch rühmen, und sagen, daß solche Phantasten starke und beständige Männer gewesen, die ihre Mutter haben verachten können, Gottes Gebot unangesehen, welches nicht gebietet, die Mutter zu verstoßen oder zu verachten, sondern daß man die Eltern ehren und ihnen dienen soll.

141. Darum hat ein anderer Einsiedler seine rechte Strafe für seine Thorheit empfangen, zu dem der Teufel gesagt hatte, daß ein Geist zu ihm kommen würde in der Gestalt seines Vaters, da er vernommen, daß der Vater seinen Sohn hat wollen besuchen. Da nun der Vater gekommen, ist der Mönch hervor getreten und meinte, es wäre ein Gespenst, und ergreift eine Axt und schlägt damit den Vater an den Kopf, und erwürgt ihn also. Hernach aber hat er gesehen, daß es kein Geist, sondern sein eigener Vater gewesen war. Mit solchem teuflischen Betrug haben allerdings diejenigen müssen umgetrieben und verführt werden, die ihre eigenen Wege und Fündlein erwählt und denselben nachgefolgt sind. Warum bleibst du aber nicht viel lieber bei deinen Eltern, Weibe und Kindern, wie dir Gott befohlen hat, bis etwa ein Dieblein kommt, und mit Gewalt beide Vater und Kinder hinweg nimmt? etc. Da wirst du denn recht zu dem höchsten Gottesdienste der ersten Tafel geführt werden, daß du Christum bekennest, und das Bekenntniß seines Wortes und reiner gesunder Lehre vorziehest allem demjenigen, was dir auch in diesem Leben billig am liebsten ist.

142. Uebrigens werden die natürlichen Neigungen durch den Heiligen Geist, durch das Wort und Glauben nicht ausgelöscht, sondern gemehrt. Wie daselbst an diesen Vätern zu sehen ist, die solcher heftigen Affecten haben ganz entbrannt sind. Wie alhier von Joseph gesagt wird, auf Hebräisch, nichtmehrn, das ist: „Sein Herz entbrannte ihm gegen seinem Bruder.“

143. Von demselben Worte aber nennen die Juden unsere Mönche keramim, und die Nonnen kemaroth, auf ihre Sprache, als die da schwarze Kleider tragen. Dasselbe thun sie aber nach ihrer Meinung und eigenem Willen; denn sie haben keine Grammatik und haben viel weniger den rechten Verstand der Worte in der heiligen Schrift. Und wir müssen die rechte Bedeutung der Worte aus dem Texte und Exempeln der Schrift nehmen. Als nämlich, wie man diese Art und Weise zu reden verstehen soll, das kann man nehmen und lernen aus 1 Kön. 3, 26. von der Mutter, welcher Kindlein noch am Leben war; da sagt der Text: Incaluerant viscera ejus ob filium suum, das ist: „Ihr mütterlich Herz entbrannte über ihren Sohn“ etc.

\quellseite{416}

144. Derhalben ist diese Ableitung desselben Worts besser und wahrhaftiger, warum die Mönche kerarim genannt werden, nämlich, wegen ihrer sonderlichen Religion oder Gottesdienstes, und ihres strengen Lebens, so sie geführet. Denn sie wenden vor, daß sie viel hitziger sind und größeren Eifer haben zum Gottesdienste, als sonst andere gemeine Leute. Und darum entweichen sie auch, und begeben sich in die Klöster, entziehen sich also von der Gemeinschaft anderer Leute, wollen keine Eheleute, keine Regenten oder Kirchendiener sein, und sondern sich also ab von den heiligen Ständen, so Gott selbst eingesetzt hat.

145. Eben dasselbe Wort steht auch 2 Kön. 23, 5., da der König Josia die Camarim hat abgethan; welches eine sonderliche Secte der Priester oder Geistlichen gewesen zu sein scheint, die besser wollten sein denn die Andern, dazu in dem Bestreben nach der Göttlichkeit etwas brünstiger und andächtiger; wie das Weib 1 Kön. 3, 26. in ihrem Herzen erbebete und entbrannte über ihrem Kinde etc. Dieselben Camarim hat Josia aus dem Königreiche vertrieben, da er die rechte göttliche Lehre wiederum aufgerichtet, und sonst auch viel andere Abgötterei und Götzen zerstört und abgeschafft hatte. Darum halte ich nicht dafür, daß die Mönche kerarim mögen genannt werden um der Kleidung willen, sondern wegen ihrer sonderlichen Andacht, welche sie vorgenommen haben.

146. Eben solche Götzendiener sind auch gewesen, die sie auf hebräisch miplezet genannt, und die Baalsophaffen, die sich auch vor den Andern sonderliche Religion und Heiligkeit angemaßt und derselben gerühmt haben, daß sie Gott nahe, und ihm dermaßen verwandt wären, daß sie ihn ihren Bräutigam oder Ehemann nenneten. Denn so viel heißt das Wort baal. Und mit demselben Namen haben sie sich sehr erhoben über allen andern Gottesdienst, gleich als ob sonst das andere gemeine Volk unserm Herrn Gott nicht recht brünstig oder andächtig genug dienete. Aber sich selbst nannten sie unsers Herrn Gottes Bräutigame und Bräute. Und derhalben erwählten sie auch sehr strenge und harte Regeln und Satzungen, die sie gehalten haben neben dem Opfer, von Mose geboten; gänzlich wie zu unsern Zeiten die Mönche gethan haben und vornehmlich die Carthäuser, welche man billig und recht Baaliten nennen möchte.

147. Denn sie haben das Wort, die Taufe und Sacramente verachtet, darum daß dieselben so herrlich nicht scheinen und vor den Augen der Menschen so hoch nicht daher prangen. Was Liebe, Hoffnung und Glauben haben sie gesagt, welche die Apostel gelehrt haben; wir tragen Platten und haben uns mit hänfenen Strängen unürtelt, essen kein Fleisch, enthalten uns der Wollust im Ehestande, item, der weltlichen Polizei und sonst, daß wir mit andern Menschen keine Gemeinschaft haben. Wir sind rechte Heilige und Baaliten: die gemeinen Leute gehen uns nichts an, welche sonst keinen sonderlichen Verstand haben, ausgenommen, was die Werke dieses gemeinen Lebens belangen mag.

148. Also ist die Möncherei nichts Anderes, denn nur eine Verachtung der göttlichen und himmlischen Stände, und der Werke, welche Christus selbst eingesetzt und bestätigt hat, als, der heiligen Sacramente, der Zehen Gebote, des Psalters, des Katechismus und dazu der ganzen heiligen Schrift. Sie aber haben das arme Volk jämmerlich betrogen mit lauter Lügen, Träumen und falschem Scheine der Heiligkeit; und haben also gemacht, daß beide die Taufe und das Blut des Sohnes Gottes vor ihren äußerlichen Larven gar verachtet worden ist. Als ob das geringe Herrlichkeit sei, daß Einer getauft und mit dem Blute Christi gereinigt ist, daß Einer ein Christ und Gliedmaß Christi, ja, ein Mitgeselle der Propheten, Apostel und Engel, dazu auch aller Heiligen ist. Das ist ja wahrlich eine große Ehre, welche die Mönche verfinstert und untergedrücket haben mit ihren mönchischen, camaritischen und baalitischen selbst erwählten Gottesdiensten.

149. Darum soll man mit höchstem Fleiße wider sie streiten, und wir sollen die Ehre und Herrlichkeit der heiligen Taufe, des Herrn Christi, der Zehen Gebote, des Evangeliums hoch heben und rühmen. Diese Stücke soll man lernen und üben in diesem Leben, du lebest gleich ohne Ehe oder im Ehestande, du seiest ein Regent oder gemeiner Unterthan, daß du nur darauf sehest, daß du das Evangelium, die Taufe und die Artikel des Glaubens lernest und dieselbigen in großen Ehren haltest und immer in deinem Herzen daran denkest. Wiederum aber sollst du die scheußlichen Mönchswerke, als, Kappen und Stricke, nur mit Füßen treten. Denn damit sind die Leute jämmerlich verführet und verderbt worden von dem Menschen der Sünde und Kinde des Verderbens.

\begin{quote}
\textbf{V. 31. Und da er sein Angesicht gewaschen hatte, ging er heraus, und hielt sich fest, und sprach: Leget Brod auf.}
\end{quote}

150. Das hebräische Wort aphak heißt, sich enthalten. Es ist aber nicht ein solch Enthalten, da man sich mit Essen, Trinken und Kleidung mäßig und eingezogen hält, sondern besteht innerlich im Herzen, in Barmherzigkeit und Liebe, und ist eigentlich nichts Anderes, denn daß man sich äußerlich stellt, als ob man gar beständig und stark sei, so doch innerlich das ganze Herz voll eitel Mitleid gegen die, so betrübt sind, fließt; wie von Joseph gesagt worden ist. Und er hält sich doch fest, daß seine Brüder nicht inne geworden, daß er geweint hätte. Es heißt eigentlich, festhalten, und sich mit Gewalt anders stellen, denn Sinnen zu Sinne ist. Als die ein furchtsam und erschrocken Gewissen haben, und wollen doch für stark und unerschrocken gehalten werden. Wie die Mönche sich stellen, als flöhen sie vor der Welt. Aber solch Enthalten ist eine rechte stöckische Heuchelei, und ist kein Stoiker, der in Wahrheit das thue, wie er sich mit äußerlichen Gebärden stellt.

151. Also können sich alle Heuchler fein inne halten, dichten sich äußerlich eine große Stärke, und stellen sich, als ob sie sehr hart seien; aber es ist lauter Heuchelei mit ihnen. Joseph aber, da er von innerlichen Affecten ganz entbrannt war, also daß er sich schwerlich hat des Weinens enthalten können (denn er hatte schon Thränen vergossen), hält sich doch ganz fest und läßt ein Mahl zurichten, stellt sich dieweil auch, als ob er stark wäre. Durch welche Stärke er die natürliche Neigung gegen seinen Bruder etwas gemäßigt und bedeckt hat; aber sein Herz läuft niemals von Thränen und feuriger Liebe ganz über.

152. „Brod auflegen“ ist eine hebräische Art zu reden, heißt so viel als, ein Mahl zurichten. Denn das Wort „Brod“ wird syneddochisch verstanden für Speise und Trank.

\section[Fünfter Theil]{Fünfter Theil. Von dem Mahle, das Joseph seinen Brüdern gegeben}

153. An dieser Stelle wird beschrieben die vaterländische Sitte der Egypter, die einen solchen Abscheu vor den Juden gehabt, Gemeinschaft mit ihnen zu haben, daß auch keiner mit einem Juden hat essen mögen. „Denn es ist“, wie Mose sagt, „vor ihnen ein Greuel.“ Denn sie haben nicht allein aus Hoffahrt und Stolz, sondern auch aus Aberglauben alle anderen Völker, geschweige die Juden, verachtet. Und dasselbe bezeugt nicht allein die heilige Schrift, sondern auch die Schriften der Heiden; wie denn im Herodot und andern Schriftstellern ihre Hoffahrt und abergläubische Religion auch klärlich angeführt und beschrieben wird. Und wird im 87. Psalm Z. 4. und Jes. 30, 7. herab genannt, das ist, stolz und hoffärtig: „Ich will predigen lassen Rahab, und Babel“ etc. Und hernach 1 Mos. 46, 33. 34. unterrichtet Joseph seinen Vater und Brüder, daß sie zum Könige sollen sagen, sie seien solche Leute, die mit Vieh umgehen, auf daß sie im Lande Gosen wohnen mögen. „Denn was Viehhirten sind“, sagt er, „das ist den Egyptern ein Greuel.“

154. Darum ist es ein stolz hoffärtig Volk gewesen vor allen andern Nationen in der ganzen Welt, um des großen gewaltigen Königreichs willen, darin es allenthalben so wohl gestanden und das sehr reich gewesen ist, dazu ein fruchtbar Land gehabt hat, darüber alljährlich der Nil geflossen und dasselbe also fruchtbar gemacht hat. Daher muß nothwendig in Egypten Speise und Getreide sehr wohlfeil gewesen sein. Und ist auch Arabien nahe dabei gelegen, und dazu das Rothe Meer, auf welchem allerlei Waare nach Egypten gebracht wurde. Wie jetzt Frankreich das hoffährigste Königreich ist, darum daß es so ein fruchtbar Land sein soll. Derhalben stolziren solche Königreiche auch vor GOtt, dieweil sie solchen unermeßlichen Segen haben; wie die Herrlichkeit und Ehre Egyptens von den Propheten auch gerühmt wird, daß es dem Paradiese GOttes gleich gewesen sei.

\quellseite{417}

155. Die andere und vornehmste Ursache aber, daß die Egypter so hoffärtig gewesen, war ihre Religion. Denn sie meinten, daß sie damit die Andern alle überträfen, und wollten derwegen für reiner und heiliger gehalten sein denn andere Völker. Denn droben ist gesagt worden, wie groß ihr Aberglaube gewesen, daran um auch der König viel Wahrsager, Weise, Magier und Sterndeuter unterhalten.

156. Ich halte aber dafür, daß die Juden bei den Egyptern so verhaßt gewesen seien um der Beschneidung willen. Denn eine jegliche Religion hat ihre sonderlichen Wahrzeichen und Ceremonien, so der Religion der andern Völker nicht gleich sind, und davor die Andern einen Greuel haben. Wie die Türken uns Christen verachten um der Taufe willen, daß sie derhalben die Christen auch nicht werth achten, daß sie mit ihnen sollten über Tische sitzen, kaum daß sie selbige werth achten, ihre Diener zu sein. Also sagte droben, Cap. 39, 17., Potiphars Weib auch mit großer Verachtung: „Der hebräische Knecht, den du uns herein gebracht hast“ etc. Als wollte sie sagen: Der heillose Jude wollte mir Gewalt anthun; es wäre nicht so gar unbillig und schändlich gewesen, wenn etwa ein egyptischer Edelmann sich dessen unterstanden und nicht der beschnittene Mann von dem schändlichsten Geschlecht und Volk der Juden.

157. Sie haben aber nicht allein vor den Juden einen Greuel gehabt, sondern die Viehhirten, so ihre eigenen Landsleute und Mitbürger gewesen, haben sie auch verachtet. Und wie man sagt, daß zu Constantinopel geringe Leute, die unbekannt und arm sind, mit Knütteln vom Schlosse des türkischen Kaisers getrieben werden, damit sie ihm nicht vor seine Augen gehen: also haben die Egypter ihre Viehhirten auch verachtet, als arme verworfene Leute, die gar eines geringen Standes und Vermögens wären, beide ihrer Polizei, Ordnung und auch ihrer Religion nach; also ob diejenigen, so hoch über den Andern sitzen, ihres hohen Standes halben auch heiliger sein sollten denn Andere.

158. Also hielten die Juden dafür, daß der reiche Mann, Luc. 16, 19., der sich mit Purpur und köstlicher Leinwand kleidete, viel heiliger und GOtt angenehmer wäre, denn Lazarus, der vor des Reichen Thüre lag voller Schwären. Und im 144. Psalm V. 15. wird auch solcher Wahn oder Meinung der GOttlosen erzählet, da sie sagen: „Wohl dem Volk, dem es also gehet.“ Also halten uns die Türken auch für verdammte Leute um der Trübsal, Unglück und Jammers willen, so wir leiden müssen; versprechen aber sich selbst die ewige Seligkeit darum, daß sie in diesem Leben ihres vielen Gutes und großer Gewalt halben in solchem Flor sind und so wohl stehen. Das ist der Egypter Philosophie und der Türken Religion, welche die christliche Lehre sein widerlegt, wie anderstwo gelehrt wird.

159. So beschreibt nun Mose den Gebrauch, so bei den Egyptern gewesen, wenn sie das Mahl gehalten und gegessen haben, nämlich, daß die Amtleute und Räthe, deren Arbeit und Rath Joseph in Verwaltung des gemeinen Wesens sich bediente, besonders an einem eigenen Orte gesessen haben. Lyra aber und die Andern zeigen als Ursache, warum die Egypter den Juden so gar feind gewesen, dies an, daß die Juden manches Vieh zu schlachten pflegten, welches die Egypter angebetet haben. Und dieselbe Ursache ist wahrhaftig. Denn hernach, 2 Mos. 8, 26., sagt Mose: „Das taugt nicht, daß wir also thun; denn wir würden der Egypter Greuel opfern unserm GOtt, dem HErrn; siehe, wenn wir denn der Egypter Greuel vor ihren Augen opferten, würden sie uns nicht steinigen?“ Denn die Egypter ehrten den Abgott Apis, welches auf ihre Sprache „Ochs“ heißt, und man durfte bei ihnen keinen Ochsen schlachten, der im Pfluge gezogen hat. Ich halte aber dafür, daß hier noch eine andere Ursache sein müsse, sintemal die Brüder Josephs von seinem Tische essen, und ja daselbst weder Schafe noch Ochsen geschlachtet hatten, dieweil sie Gäste und Fremdlinge gewesen.

160. Es erhebt sich aber hier noch eine andere Frage, die sich zu dieser Erzählung etwas schickt, nämlich: Warum doch Joseph, der fromme heilige Mann und Prophet, in den Gebrauch dieser sehr gottlosen und hoffärtigen Leute gewilligt und demselben gefolgt sei? dieweil droben etlichemal gesagt worden aus dem 105. Psalm V. 22., daß er sie Weisheit gelehret und die Fürsten unterwiesen habe nach seiner Weise. Warum hat er denn diesen unheiligen Gebrauch nicht verdammt? Ja, er hat auch zugelassen, wiewohl er ein Regent war über das ganze Königreich, daß sie ihn selbst auch ausgesondert und von sich ausgeschlossen, und sie allein und besonders gegessen haben. Das ist nun wahrlich ein wunderlich Ding. Sie wußten, daß er ein Hebräer war, vor welchem Volke die Egypter einen Greuel haben, wie er im Gesängniß von sich selbst gesagt 1 Mos. 40, 15.: „Ich bin aus dem Lande der Ebräer heimlich gestohlen“ etc. Und der Schenke sagt auch zum Könige 1 Mos. 41, 12.: „Da war bei uns ein ebräischer Jüngling“ etc. Und ist doch zum Fürsten und Herrn über ganz Egyptenland gesetzt worden; so es doch jedermann bewußt, daß er von demselben verachteten abscheulichen Volke hergekommen und geboren war. Und sind bei solcher Gemeinschaft, so Joseph und die Egypter miteinander gehabt, noch immer solche Gewohnheiten und Gebräuche beider Völker mit seinem Wissen und Willen also unverändert geblieben.

161. Auf diese Frage antworte ich also: Die Historien in der heiligen Schrift zeigen, daß nie ein Volk so gar rein und ohne allen Mangel gewesen sei. Mose hat die Kinder Israel in einer großen Anzahl aus Egypten geführt; und Bileam sagt von ihnen 4 Mos. 23, 23.: „Es ist kein Wahrsager oder Abgott in Israel“; und steht doch im Propheten Amos Cap. 5, V. 26. also: „Ihr truget die Hütten eures Molochs, und die Götzen eurer Bilder, und den Stern eurer Götter, welche ihr euch selbst gemacht hattet.“ Und Stephanus bestätigt solches auch, da er Ap. Gesch. 7, 42. sagt: „GOtt wandte sich, und gab sie dahin, daß sie dieneten des Himmels Heer; wie denn geschrieben stehet in dem Buch des Propheten.“ Darum hat Mose so ein rein Volk nicht gehabt, wie es schien, sondern es sind ja Etliche zum wenigsten in ihrem Unglauben und gottlosen Wesen geblieben. Wie er selbst darüber klagt 5 Mos. 29, 19.: „Und ob er schon höre die Worte dieses Fluchs, dennoch sich segne in seinem Herzen und spreche: Es gehet mir wohl, weil ich will wandeln, wie es mein Herz dünket, auf daß die Trunkene mit der Durstigen dahinfahre“ etc. Da zeiget er an, daß etliche der Juden gottlos und abgöttisch gewesen. Aber, das 4 Mos. 23, 23. steht: „Es ist kein Wahrsager in Israel“, das ist \unsicher{sinnetöchig} geredet.

162. Also wird im 105. Psalm V. 6. ff. gerühmt ihre Gottseligkeit und ihr Glaube mit allen Wohlthaten, so ihnen GOtt aus Gnaden erzeiget hatte. Wiederum, bald hernach, im 106. Psalm V. 21. ff., wird ihr Unglaube gestraft, da der Prophet sagt: „Sie vergaßen GOttes, ihres Heilandes, der so große Dinge in Egypten gethan hatte“ etc. Derhalben dürfen wir nicht denken, daß im Volke Israel Alle fromm gewesen, wiewohl Alle durch das Rothe Meer gegangen sind; sondern um des Theils willen, so fromm und gottselig gewesen, genießt der andere Haufe die Wohlthaten GOttes auch mit, welche er seiner Kirche zu geben pflegt.

163. Und dasselbe geschieht also zu allen Zeiten, und zu unserer Zeit besteht die Welt und werden die Regimente erhalten um einiger anderer Ursache willen, denn daß sich GOtt eine Kirche sammelt mitten unter dem bösen unartigen Volke. Denn daß man noch in Landen und Städten Regimente und Polizeien hat, darf man den Fürsten und Königen nicht zuschreiben oder danken, welche gemeiniglich schier gottlose und die ärgsten Menschen sind; sondern es wird alles erhalten um des Wortes willen und der Taufe, und wegen des heiligen Sakraments, der noch in der Kirche übrig ist. Denn wo die Gottseligen nicht wären und die noch sollen selig werden, würde die Welt nicht bestehen bleiben; und wenn nun die letzten Heiligen leben werden, so wird bald der jüngste Tag kommen. Denn GOtt fragt nicht nach der weltlichen Polizei oder Hausregiment, denn nur allein um der Kirche willen.

164. Es sind übrigens unter dem heiligen und gläubigen Volk allzeit auch GOttlose und Heuchler mit untermenget; wie das in der Historie Davids, der ein sehr heiliger König gewesen, zu sehen ist. Denn siehe doch seine erbärmlichen Klagen an, darin er nicht allein die Tyrannen im politischen Regimente, sondern vielmehr die gottlosen falschen Lehrer beschuldigt. Derhalben, dieweil nun weder Mose, noch David, dazu andere Könige und Propheten, Elia und Elisa etc., nicht alles, was unrein gewesen, haben können abthun, und alles aufrichten und wieder zurechtbringen, was zu ihrer Zeit gebrechlich und sträflich gewesen: was ist es denn groß Wunder, wenn wir zu dieser Zeit auch noch unter gottlosen und bösen abgöttischen Menschen leben müssen?

\quellseite{418}

165. Auf diese Weise wird Joseph auch gebessert haben, was er hat bessern können, und denen er hat können helfen, daß sie möchten selig werden, die hat er unterrichtet und erhalten: aber ein großer Theil ist im Irrthume geblieben und bei dem Aberglauben ihrer Voreltern. Darum muß man das alles synedochisch verstehen, was im 105. Psalm V. 22. steht: „Er unterweisete die Fürsten und Aeltesten nach seiner Weise“ etc., nämlich, daß sie damit nicht alle gemeint sind, sondern nur Etliche. Daran ist es auch genug.

166. Nun gibt ihm Pharao nichts desto weniger vollkommene Gewalt im politischen Regimente und braucht ihn als einen getreuen Diener, wiewohl er in einem sehr schlechten, geringen und verachteten Volke geboren war; läßt ihm zu, daß er seines Gesetzes und Ceremonien gebrauchen mag, daß er sich der Dinge enthalte, welcher sich auch seine Väter enthalten haben; und ändert der König das ganze Regiment im Königreiche nicht, sondern läßt das Wort und die Lehre Josephs ihren Lauf haben, gestattet ihm, daß er seine Religion mag ausbreiten, so viel immer möglich. Was aber die weltliche Polizei belanget, darin hat er ihm das ganze Regiment zu verwalten befohlen. Ein so großes Königreich hat nicht können so gar vollkommenlich gebessert werden, daß darin nicht etliche Gottlose geblieben wären; wie auch ohne Zweifel in der Polizei viele böse Buben gewesen sind.

167. Also mögen wir Gott danken, daß wir auch Viele durch das Wort zu uns bringen; wir müssen aber noch immer bei andern abgöttischen Leuten leben, es sei gleich unter dem Türken, oder unter Kaiser Karl. Es wird wohl für und für also bleiben, was der Prophet Amos dem Volke Israel auch aufrückt, da er Cap. 5, V. 26. sagt: „Ihr habt euren Abgott Remphan gedienet“ etc.

168. Was nun Joseph nicht hat können ändern oder bessern, daßelbe hat er dem Gerichte Gottes befohlen, und hat sich gehalten nach dem gemeinen Sprüchworte: Laß es gehen, wie es geht etc. Wie er diesen Gebrauch nicht hat können abschaffen, daß die Egypter an den Juden nicht sollten Greuel haben. Und es hat ihm auch nicht gebühret, das Volk bewegen zu erregen und Aufruhr anzurichten. Darum müssen wir auch die Bosheit und Unsinnigkeit der Gottlosen noch leiden, wie Elia und die andern Propheten die Baaliten, Moloch und die Samaritan in der größten Menge haben dulden müssen. Denn die sieben tausend, so zur Zeit Elias übrig geblieben, waren nichts, wenn man sie gegen das ganze Volk gerechnet hat. Uns gebühret zu predigen und zu lehren: Gott aber gibt das Gedeihen; eine Zeit mehr, die andere weniger. Wir ermahnen ja wohl, wir schelten, überzeugen und strafen Andere zum Exempel; wenn wir aber sehen, daß wir damit nichts ausrichten können, sollen wir sagen: Herr Gott, ich will dir es alles anheimstellen und befohlen sein lassen, du wirst es alles wohl machen.

169. Derhalben, da es den Egyptern also wohlgefallen, daß sie haben besonders allein essen wollen, hat ihnen Joseph denselben Gebrauch ihres Vaterlandes und Satzung ihrer Vorältern gelassen; die Andern aber, so das Wort gehört und vor ihm keinen Greuel gehabt, dieselbigen hat er freundlich aufgenommen. Und dasselbe pflegt auch also zu allen Zeiten und bei allen Völkern zu geschehen.

170. Es ist aber ein solcher Gebrauch gewesen mit diesem Mahle, so Joseph hat zurücklassen, wie \unsicher{Abschthume} bei den Mönchen war. Sie haben an langen Tischen gesessen, daß sie einander haben ansehen können; wie im Texte folgen wird: „Man setzte sie gegen ihm.“ Die Amtleute aber haben Joseph zur Rechten und zur Linken gesessen. Und es scheint, daß sie mit dem Anrichten, da sie das Essen aufgetragen, diese Ordnung gehalten haben, daß man die Speisen erstlich vor den Fürsten getragen hat, daß sie nach seinem Willen und Wohlgefallen unter die Andern und die Amtleute oder Räthe selbst auch ausgetheilt würden.

171. Dies ist ein löblicher und nützlicher Gebrauch gewesen bei den Völkern im Morgenlande, daß man einem Jeden insonderheit seinen bescheidenen Theil von Speise und Trank gegeben hat. Das ist eine große Mäßigkeit gewesen, damit man viel gespart, und die sich mit der Deutschen Weise nicht reimen würde. Ein Jeder hat sein bescheiden Theil Fleisch und Brod gehabt. Sie haben nicht vierhundert oder fünfhundert Schüsseln aufgetragen; welches heutiges Tages unsere Fürsten mit großem Ueberfluß zu thun pflegen, damit über die Maßen viel verschwendet und unnütz verthan wird.

172. Es scheint, als habe Joseph diese Ordnung also erstlich erfunden und aufgerichtet; und wird derselbigen auch gedacht im Buch Esther, Cap. 1, V. 8. Ja, die Römer und Griechen haben selbst auch einem Jeden sein bescheiden Theil gegeben nach Gelegenheit der Gesundheit, der Stärke, oder darnach ein Jeder groß oder klein vom Leibe gewesen. Das ist wahrlich eine sehr löbliche Ordnung im Hausregimente gewesen, daß nicht ein Jeder nach seinem Gefallen Alles auf einmal verzehren, auffressen und verschlemmen möchte; wie unsere Leute zu thun pflegen. Und sonderlich die Fürsten zu unsern Zeiten sind gar unsinnig mit Schlemmen, sind darauf geschickt, daß sie nur ohne Maaß und Aufhören Alles verbringen und verprassen mögen. Da ist gar keine Hoffnung der Besserung, sondern man muß das wilde unordentliche Wesen in diesem letzten Alter der tollen Welt also dulden und gehen lassen. Wenn wir etwas mäßiger wären, und doch nur die Kosten ersparen könnten, die jährlich auf das unmäßige Getränk, Bier und Wein, gewendet werden, so würden wir viel mehr Geld und Gut haben, beide insgemein und insonderheit.

173. Es haben übrigens Joseph und die Egypter Jakobs nach der Weise seiner Religion gegessen; allein, daß sich die Egypter viel größere Heiligkeit und Reinigkeit angemaßt haben. Joseph hat bisweilen etwa Kalbfleisch oder sonst ander Fleisch essen mögen, welches den Egyptern nicht geziemte um ihres Abgotts Apis willen. Darum haben sie Joseph seine Weise halten lassen, sofern sie nur nicht gezwungen würden, dieselbe auch mitzuhalten. Er wird ohne allen Zweifel treulich und mit höchstem Fleiße gelehrt haben, so viel es immer hat sein mögen; sie aber haben gethan, was ihnen wohlgefallen hat.

\begin{quote}
\textbf{V. 33. Und man setzte sie gegen ihm, den Erstgebornen nach seiner Erstgeburt, und den Jüngsten nach seiner Jugend. Deß verwunderten sie sich unter einander.}
\end{quote}

174. Das ist etwas Sonderliches, daß er sie hat sich gegenüber setzen lassen, einen Jeden nach seinem Alter. Die Hebräer sagen und ist auch sehr wahrscheinlich: Es habe sich Joseph gestellt, als ob er ein Reisewager wäre, und seine Brüder haben gesehen, daß er in großen Ehren gehalten worden ist und ein trefflich Ansehen gehabt, und daß er Egyptenland mit sonderlicher Weisheit und Rath regieret hat, dieweil er droben, Cap. 42, 18., gesagt: „Ich fürchte Gott.“ Darum haben sie ihn für einen gottesfürchtigen Mann gehalten, und haben leichtlich geglaubt, daß er die Gabe gehabt, daß er habe zukünftige Dinge verkündigen können; vornehmlich, da sie gesehen, daß er sonderliche Kunst und Ceremonien gebraucht, welche auch die Propheten gebraucht haben, als, Ceremonia hat eine hölzerne Kette getragen, Jesaia ist barfuß gegangen. Also Joseph auch: dieweil er in Wahrheit ein Prophet gewesen, hat er auch sonderliche Ceremonien gebraucht, daß ihn seine Brüder für einen Weissager oder egyptischen Propheten hielten.

175. Die Hebräer sagen: Da er sie habe nach einander heißen sitzen, wie sie sitzen sollten, soll er an den Becher geschlagen, und mit demselben Klange ein Zeichen gegeben haben, daß der Erstgeborne sollte obenan sitzen, und so immerfort, Einer nach dem Andern. Und es scheint, daß daher gehöre, das im nachfolgenden Capitel V. 5. gesagt wird: „Der Becher, den ihr gestohlen habt, ist eben der, daraus mein Herr trinket, und damit er weissaget.“ Darum ist es sehr wahrscheinlich, daß er sonderliche Geberden müsse gebraucht haben, auf daß er ihnen damit einen Schein vormachte, als ob er dessen von Gott Bericht empfangen, wie sie nach einander sollten zu Tische sitzen.

176. Sie hätten aber schier verstehen können, daß er sich vor ihnen nur also gestellt, dieweil er so gar genau eines Jeden Alter gewußt, als ob er auch mit in ihrem Geschlechte und bei eines Jeden Geburt gewesen wäre. Aber Joseph hat solches damit verhöhlet, daß er an den Becher geschlagen und daß er sich sonderlich tapfer und ernstlich gehalten, und sich gestellt, als ob er zukünftige Dinge wüßte; auf daß sie nicht merken sollten, daß er vor ihnen diesen Schein und Betrug angerichtet hätte.
\quellseite{419}

177. Was aber für eine Weissagung in dem Becher gewesen, zeigt Mose nicht an. Sie haben sich also desto leichter betrügen lassen mit dieser ungemeinlichen Weise, dieweil sie schon vorlängst an ihres Bruders Leben gezweifelt; und haben nicht können denken, daß dies sollte Joseph sein, sondern es müßte sonst ein großer Herr und heiliger Prophet sein, welcher der egyptischen Religion sehr wohl erfahren wäre. Auf solche Weise spielt Joseph mit seinen Brüdern ganz freundlich, beweist damit, daß er gegen sie ein freundliches und gütiges Herz gehabt, verdunkelt aber und verdeckt gleichwohl noch immer solche Freundlichkeit und guten geneigten Willen, daß sie denselben so gar eigentlich nicht haben verstehen können.

\begin{quote}
\textbf{V. 34. Und man trug ihnen Essen vor von seinem Tisch; aber dem Benjamin ward fünfmal mehr, denn den andern. Und sie tranken, und wurden trunken mit ihm.}
\end{quote}

178. Es sind die Schüsseln erst auf Josephs Tisch gesetzt worden, von denselben hat er einem Jeden sein Theil Essen geben lassen. Benjamin aber hat fünfmal so viel empfangen. Und daraus hätten sie auch auf die Vermuthung kommen mögen, daß sich Josephs mehr äußerlich also gestellt hätte, nämlich, dieweil Benjamin vor den Andern solche Ehre erzeigt wird. Aber sie haben vielleicht gedacht, daß solches der Könige und Herren Gebrauch sei, daß sie hoch heben, wen sie wollen, und wiederum auch verstoßen, welche sie gelüstet.

179. Dies ist aber ein sehr lustig Schauspiel und Exempel, welches wohl geistlicher Lehre ist, wenn man es auf die göttliche Steigerung zieht, damit er seine Heiligen wunderlich führt; auf daß wir nicht bald verzagen, wenn wir angefochten und gestraft werden, und auch im Glück und wenn es uns wohl gehet, nicht vermessen oder stolz werden. Denn Gott versucht uns auf beiden Seiten, mit Glück und Unglück, auf daß wir also geübt werden und den Herrn allezeit loben; er gehe gleich freundlich mit uns um, oder aber lasse uns geplagt werden, so sollen wir doch in unserm Herzen zufrieden sein und ihm danken; denn er spielet eben so mit uns, wie Joseph mit seinen Brüdern gethan.

180. Da sie aber nun das Mahl gehalten, sagt Mose, sind sie trunken geworden. Welches nicht von dem übermäßigen Saufen der Deutschen oder Griechen zu verstehen ist, sondern ist geredet nach der Art und Weise, die auch an andern Orten in der heiligen Schrift gebräuchlich ist; als, im 65. Psalm V. 10.: „Du suchest das Land heim, und wässerst es“, inebrias eam, machst es gleichsam trunken etc. Wenn es regnet, dann macht Gott das Land trunken, wässert es und macht es fruchtbar. Und Spr. 11, 25.: „Wer trunken macht, der wird auch trunken werden“, das ist, wer die Armen mildiglich speiset, daß sie satt werden, der soll auch wieder reichlich gesättigt werden. Item in 23. Psalm V. 5.: Calix meus inebrians, id est, exuberans; auf Deutsch: „Du schenkest mir voll ein“ etc.

181. So heißt nun „trunken werden“ nach der Art der Schrift nichts Anderes, denn sich satt essen und trinken, satt und fröhlich werden. Es heißt nicht eine solche Trunkenheit, wie der Griechen und Deutschen Trunkenheit ist, die in ihrem Wohlleben und täglichen Schlemmen nicht allein den Bauch füllen, sondern wenn er voll ist, so übergeben sie sich bald, und speien wieder heraus, was sie gefressen und gesoffen haben. Diese Brüder Josephs aber sind fröhlich gewesen, haben einander ermahnt, zu essen und zu trinken, und haben sich also erfreut über die Gnade und Wohlthat, so ihnen derselbige Herr erzeigt hatte. Aber dies Glück wird mit ihnen nicht lange währen, sondern wird bald wiederum in eine plötzliche und sehr schwere Traurigkeit verändert werden. Wie denn solche Veränderung der göttlichen Werke allezeit im Leben der Menschen pflegt zu sein, da bald Traurigkeit, bald wiederum Freude vorfällt; auf daß wir also immer im Glauben, Hoffnung und Liebe geübt und versucht werden mögen.

\chapter{Das vierundvierzigste Capitel}

\section[Erster Theil]{Erster Theil. Wie Joseph seine Brüder wieder fort schickt, wie er sie aber bald wieder zurückholen läßt und wie Josephs Brüdern hierbei zu Muthe}

\begin{quote}
\textbf{V. 1. 2. Und Joseph befahl seinem Haushalter, und sprach: Fülle den Männern ihre Säcke mit Speise, so viel sie führen mögen, und lege jeglichem sein Geld oben in seinen Sack; und meinen silbernen Becher lege oben in des Jüngsten Sack, mit dem Geld für das Getreide. Der that, wie ihm Joseph hatte gesagt.}
\end{quote}

1. Ich habe nun oftmals ermahnt, und man soll es auch den Leuten immer vorhalten und einschärfen, daß der Heilige Geist von den so großen Patriarchen so scherzhafte und geringe Dinge schreiben läßt, so er doch wichtige und große heilige Dinge erzählen möchte; wie er denn derselben zwar auch bisweilen etliche mit einführt und sie unter die Historien der heiligen Väter mengt. Ein unverständiger fleischlicher Leser, der da meint, daß diese Dinge gar nichts werth seien, ärgert sich daran leichtlich und verwundert sich, daß solches in der Kirche und Gemeinde Gottes gelesen wird, und daß der Heilige Geist die Zeit und Arbeit darüber also mag zubringen, solche geringe nichtige Dinge zu erzählen. Denn warum hält er uns nicht viel lieber große seltsame Wunderwerke vor, von den Mönchen Fasten, von stoischer und Spartanischer Härigkeit und Unfreundlichkeit der Menschen, die gar hart gewesen wie die Eisen; wie die Carthäuser für solche harte Menschen wollen gehalten werden: als ob in diesen lächerlichen und nichtigen Dingen sonderliche große Lehre sein könnte. Desgleichen disputiren sie auch davon: Ob dies Spiel, so Joseph mit seinen Brüdern getrieben, Gott auch könne wohlgefallen, und aus weß Eingeben oder welchem Geiste er das möge gethan haben?

2. Darauf antwortete ich also: Daß Joseph dies darum gethan und vom Heiligen Geiste darum auch sei beschrieben worden, daß wir daraus lernen, wie man vor Gott recht leben solle. Denn das Leben gilt bei unserm Herrn Gott nicht, ist ihm auch gar nicht angenehm, welches wir selbst nach unserer Weisheit, Geschicklichkeit und eigenem Gutdünken erwählen. Zudem, so haßt auch Gott und verabscheut die harte Apathie, das ist, die unnatürliche Unfreundlichkeit der stoischen Köpfe, so den Menschen aller Sinne und natürlicher Affekten oder Bewegungen beraubt, und also in wunderbaren Dingen wandelt, die ihr zu hoch sind. Und Gott sagt eben dasselbe zu dir, daß er auch zu den Juden gesagt hat, Jes. 1, 11. 12.: „Was soll mir die Menge eurer Opfer? spricht der Herr. Ich bin satt der Brandopfer von den Widdern“ etc. „Wer fordert solches von euren Händen?“ Als wollte er sagen: Wer hat euch doch das alles heißen thun? Und Jes. 58, 2. 3.: „Sie fordern mich zum Recht, und wollen mit ihrem Gott rechten. Warum fasten wir, und du siehest es nicht an? Warum thun wir unserm Leibe wehe, und du willst es nicht wissen?“ Denn sie wollten etwas Großes und Vortreffliches thun, das Gott wohlgefallen sollte. Wie Jener, der sich selbst gedachte zu tödten, indem er sein Wasser über Macht hielt. Aber sie haben von Gott hören müssen: Ich habe euch solches nicht befohlen etc.

3. Also da neulich in Spanien die Königin krank gelegen, hat man etliche Personen zu ihr gebracht, die sich selbst mit eisernen Scorpionen gegeißelt, auf daß sie mit ihrem Blute Gott versöhnen möchten, und von ihm damit erlangen, daß er der Königin dadurch ihr Leben wollte verlängern. Dieselben hätte man auch sollen vorhalten das Wort Gottes in Jesaia: „Wer fordert solches von euren Händen?“

4. Darum lasset uns lernen, daß Gott an den selbst erwählten Werken keine Lust oder Gefallen hat, und daß man ihm auch seine Weise vorschreiben soll, nach unserm Gefallen oder Gutdünken ihm zu dienen; sondern wir sollen uns von ihm lehren und führen lassen: unser Herr Gott will dein Schüler nicht sein; er soll es auch nicht sein; er soll doctor und gubernator (Lehrer und Regierer) sein. Aber menschliche Natur ist so gar verkehrt, daß wir uns achten und liegen lassen, was GOtt gebietet; aber was er nicht befohlen hat, das ermahnen wir, dasselbe gefällt uns wohl, das thun wir auch gern.

\quellseite{420}

5. Darauf pflegen sie aber also zu sagen: Wir thun ja das alles aus gutem Eifer und guter Meinung. Denn daher ist allerlei Aberglaube und Ketzerei gekommen, daher sind die Mönche, die falschen Köpfe und Baaliten gekommen, 1 Kön. 18, 28., die ihren Leib mit Pfriemen gestochen, daß Blut heraus gelaufen ist, und hofften, sie wollten darum erhöret werden.

6. Derhalben sollen wir wissen, daß Christen in solchem Stande wandeln und leben müssen, welcher von GOtt geboten und bestätigt ist, auf daß, du wachest oder schläfest, du essest oder trinkest, du gleichwohl sagen mögest: Dies thue ich im Namen des Herrn, nicht aus eigenem Willen oder selbstgewählter guter Meinung. Alsdann wird endlich, was du thun oder leiden wirst, unserm Herrn GOtt wohlgefallen, und du wirst sehen, daß er mit dir das ganze Leben hindurch auf das allerfreundlichste spielen wird.

7. Also ehret Joseph seinen König, und was zum Regimente im selbigen Königreiche gehört und ihm befohlen gewesen, das richtet er treulich aus. Er hat hin und wieder im Lande und Städten das politische Regiment geziert und mit Gesetzen, guten Sitten und mit rechter Lehre in der Religion, und hat wohl gewußt, daß solche Werke GOtt angenehm sind. Ja, er hat sich gehalten, wie Salomo ermahnt in seinem Prediger Cap. 9, V. 8. 9., da er sagt: „Laß deine Kleider immer weiß sein“ etc. „Brauche des Lebens mit deinem Weibe“ etc. Also thut du alles, was du thust mit Worten und mit Werken, in dem Namen des Herrn JEsu, wie Paulus sagt Col. 3, 17. Wie Joseph mit gutem friedsamen Herzen seine gewöhnliche Arbeit ausrichtet und sich gar nichts wider GOttes Befehl vornimmt. Aber die Papisten mit ihrer Apathie und daß sie sich wunderlicher Weise nicht bekümmern oder Leid tragen, unter Vorwendung sonderlicher Heiligkeit, verstehen diese guten Werke Josephs nicht, welcher im Glauben und Geist der Furcht gegen GOtt und seinen König wandelt und die Gewalt und Herrlichkeit, dazu er gekommen war, zu GOttes Ehre und anderer Leute Wohlfahrt und Seligkeit braucht.

8. Also sagt Samuel zu Saul, 1 Sam. 10, 6. 7.: „Wenn dir nun diese Zeichen kommen, so thue, was dir unter Händen kommt; denn GOtt ist mit dir.“ GOtt wird mit dir sein und Glück geben, „da wirst du ein anderer Mann werden“, wenn du zu den Propheten kommen wirst. Derhalben, wenn dein Herz und Geist rechtschaffen und dein Glaube auch rein ist, was du darnach thun wirst, wie gering und verachtet dasselbe auch sein mag, so ist es vor GOtt alles köstlich und gut. Denn dieselben Werke geschehen nach dem Willen und Regierung des Glaubens und des Heiligen Geistes.

9. Aber hier wirst du fragen: Ich bin Joseph nicht, oder ein solcher großer Patriarch. Schaue aber kein Exempel an und siehe auf deinen Beruf. Bist du ein Schulmeister, ein Schüler, ein Regent, ein gemeiner Unterthan, so beständige dich, daß du dein befohlen Amt treulich ausrichtest, und der Zeit und Person brauchst, damit dich GOtt geziert hat. Und wo du in deinem Amte wirst treu und fleißig sein, so wirst du mehr denn genug zu thun finden, und werden dir mancherlei Beschwerung und Mühseligkeit vorfallen, und es wird also nicht vonnöthen sein, daß du dir freiwillige Marter aufladest; wie Joseph seinen Leib nicht selbst gepeinigt hat, wie die Baaliten pflegten zu thun, sondern hat die Plage und Marter mit Geduld getragen, so ihm von GOtt ist auferlegt worden.

10. Das pflegen die Mönche und Heuchler nicht zu thun; denn niemand ist ungeduldiger, als sie sind, wenn sie von GOtt gestraft werden. Aber wenn sie leiden, was sie selbst erwählt haben, als, Geißeln, Fasten und sonst allerlei Kasteiung des Leibes, da geben sie große Andacht und Demuth vor. Das ist aber die rechte Demuth gewesen, wenn die heiligen Väter, Abraham, Isaak und Jakob, oftmals mit großer schwerer Mühseligkeit, Angst und Noth beschweret gewesen, und doch mit großer Stärke und festem Glauben sehr große Geduld darin gehabt haben. Das lassen die Mönche.

11. Darum verstehen die Papisten die Dinge gar nicht, welche der Heilige Geist von den Patriarchen beschreibt, sondern sie sehen dieselben an, als ob es kindische, verächtliche und fleischliche Dinge seien. Denn dies alles ist wider ihre selbsterwählte Gerechtigkeit und eigene Andacht, und lehrt uns auf GOttes Wegen wandeln und unsers Berufs fleißig warten.

12. GOtt verachtet und verdammt alle Werke der Heuchelei, denn sie sind zu nichts Anderem nütze, denn daß sie nur Hinderniß bringen, dadurch die rechte Heerstraße und der Weg GOttes verschlossen wird. Und die Leute werden dadurch verursacht, daß sie allezeit sich zu andern Stegen und Beiwegen kehren. Darum heißt Paulus, daß man sich davor hüten soll, da er Röm. 16, 17. sagt: „Ich ermahne aber euch, lieben Brüder, daß ihr aufsehet auf die, so da Zertrennung und Aergerniß anrichten, neben der Lehre, die ihr gelernt habt, und weichet von denselbigen“ etc.

13. Bist du ein Ehemann oder wartest sonst in einem andern göttlichen Stande deines Amtes fleißig, so wird dir GOtt immer zu thun und zu leiden genug schaffen. Aber daran mangelt es uns, daß wir ja derselben Dinge nichts achten, darauf wir vornehmlich Acht geben sollten; sondern bekümmern uns mehr mit andern Dingen, die uns nicht befohlen sind. Darum bleiben wir auch so unerfahren in geistlichen Sachen, lernen immerdar und kommen doch nimmer zur Erkenntniß der Wahrheit, 2 Tim. 3, 7.

14. Auf die Weise kommt es dann, daß die Dinge, so in den Historien der heiligen Schrift uns vorgehalten werden und uns schlecht und gering dünken, eine Wichtigkeit bekommen, ein ander Ansehen gewinnen, daß wir davon mehr halten denn zuvor; und sehen nun, daß es solche und so große Werke sind, daß sie kein Baalit oder Mönch jemals hätte ertragen können. Wie ich zuvor gesagt habe von der Ungeduld und Hoffart derer, so nach eigener Wahl einhergehen und nur mit selbsterwähltem Gottesdienste umgehen. So ist kein unleidlicher und giftiger Ding auf Erden, denn ein Barfüßermönch. Denn sie haben Lust zu ihren eigenen Wegen, und wollen darin geehrt und gerühmt sein und daß man ihre Träume für die göttliche Wahrheit halten soll, und können keine Strafe leiden, ja, wollen ganz und gar ungestraft sein. Wir sehen aber, daß die Patriarchen im Glauben, Hoffnung und Geduld wohl geübt und versucht gewesen sind nach GOttes Wort und Willen, und daß sie beide Glück und Unglück, wie es GOtt mit ihnen geschickt, nicht anders denn in GOttesfurcht gebraucht haben.

15. Dergleichen denkest, daß dies Spiel, so Joseph mit seinen Brüdern gehalten, GOtt ein sehr angenehm Schauspiel gewesen sei; wie er auch an der Geduld und Strafe seiner Brüder große Lust gehabt, wiewohl sie solches nicht verstanden. Und hernach wird man sehen, wie angenehm dies alles GOtt gewesen sei.

16. Auf dieselbige Weise sollst du es auch dafür halten, daß dein ganzes Leben ein solch Spiel GOttes sei, und daß ihm alles wohlgefalle, was du thust und leidest, sofern es nur im rechten Glauben geschieht; ja, daß auch der Tod vor dem Herrn ein theuer und köstlich Ding sei. Denn also sehen wir, daß GOtt an dem Leben und allen Werken der Patriarchen Lust gehabt. Und wo diese nicht beschrieben wären, würden wir alle begehren zu wissen, was wohl Abraham, Isaak und Jakob gethan haben? Jetzt aber, dieweil schier alle ihre Worte und Werke beschrieben sind, achten wir derselben nun nichts, darum daß sie so schlecht und gering sind, oder verachten sie wohl gar als eitel fleischliche Dinge. Wir lassen uns dünken, wir wollten viel größere Werke thun, wenn wir zu solcher großen Ehre und Herrlichkeit kommen möchten: aber es ist niemand, der auch das allergeringste Werk, so Joseph gethan, verstehen oder demselben nachfolgen könnte.

17. So spielt nun Joseph noch mit seinen Brüdern, auf daß er ihr Herz erforsche, und erfahren möge, wie sie gegen Benjamin gesinnet seien, ob sie denselben auch von Herzen lieb haben, sintemal sie ihn ja nicht verkauft noch getödtet hatten. Ueber das so hört er auch von ihnen, daß sein Vater noch am Leben sei. Darum begehrt er auch zu wissen, wie sie gegen ihn gesinnet und was sie von ihm halten. Und Joseph kündigt daran nicht; denn er thut Alles im Glauben und seine Werke gefallen auch GOtt darum wohl. Wenn aber ein fleischlicher Mensch oder Baalit dasselbe auch thäte, oder wenn er schon Tödte auferweckte, so wäre es doch eine Todsünde. Denn ihr Opfer ist vor GOtt ein Greuel und ihr Gebet muß Sünde sein.

\begin{quote}
\textbf{V. 3—6. Des Morgens, da es licht ward, ließen sie die Männer ziehen mit ihren Eseln. Da}
\end{quote}
\quellseite{421}
sie aber zur Stadt hinaus waren, und nicht ferne kommen, sprach Joseph zu seinem Haushalter: Auf, und jage den Männern nach, und wenn du sie erreichest, so sprich zu ihnen: Warum habt ihr Gutes mit Bösem vergolten? Ist’s nicht das, da mein Herr aus trinket, und damit er weissaget? Ihr habt übel gethan. Und als er sie ergriff, redete er mit ihnen solche Worte.

18. Hieraus ist zu sehen, wie droben gesagt, daß der Haushalter um seines Herrn Absichten gewußt habe, dieweil Joseph gesehen, daß er so verschwiegen und so getreu gewesen. Das ist aber eine sonderliche List, daß er des Morgens diese Brüder hat von sich gelassen satt und fröhlich um der großen Freundlichkeit willen, so ihnen Joseph erzeigt. Und der Ausgang wird sein übereinkommen mit dieser List, welcher ich zwar nicht würde nachfolgen können; oder, wo jetzt Einer unter uns dasselbe auch anfangen würde, derselbe würde eher ungeschickt zu sein als zu spielen scheinen. Denn es ist nicht gleiche Absicht und auch nicht derselbe Geist oder Herz solcher brünstigen Liebe. Denn wir haben zuvor gehört, daß er ganz in Thränen aufgelöst sei, und daß er nichts Böses vorgehabt, sondern nur in höchster und freundlichster Liebe mit ihnen also gespielt. Darum können wir ihn nicht für einen ungeschickten oder fleischlichen, sondern für einen sehr weisen und sehr heiligen Mann halten.

19. Diese Brüder aber, dieweil ihnen so gütlich geschehen und sie so fröhlich erfreuet worden, haben sich selbst mit wunderlicher Rhetorik und seinen geschmückten Worten geschmeichelt, und sich miteinander gefreuet, daß ihnen ihre Rathschläge, ihres Trachtens, so fein gerathen wären, daß sie so einen gnädigen Herrn hätten, der ihnen so wohl gewogen wäre, daß sie alle mit einander möchten heim zu ihrem Vater ziehen. Und hätten ohne Zweifel dasselbe auch also daheim rühmen wollen, nämlich, wie sie so freundlich von dem Manne, der im Lande Herr war, empfangen wären, und wie er sie hätte zu sich geladen, und ihnen lassen gütlich thun nicht in der gemeinen Herberge, sondern in des Fürsten Hofe selbst; und wie vornehmlich Benjamin so fröhlich wäre geehrt worden, da man ihm hätte fünfmal mehr geben lassen, als den Andern zu ihrem Theile geworden.

20. Denn also pflegen wir uns im Glück zu freuen; da sind wir ohne alle Furcht und Traurigkeit. Wie diese Brüder nicht allein mit fröhlichem Herzen, sondern auch mit hohem Geiste rühmen und triumphiren; und werden jetzt mit den Gedanken umgegangen sein, wie sie ihren Vater beschuldigen wollten, darum daß er sich so gar schwierig erzeigt hatte, Benjamin mit ihnen ziehen zu lassen. Nun mag er es ja selbst wohl sehen, werden sie gesagt haben, wie mit großer Ehre und Herrlichkeit wir jetzt wiederum heim kommen.

21. Aber es ist ihnen eben dasselbe widerfahren, das Salomo sagt, Spr. 16, 18.: „Wer zu Grund gehen soll, der wird zuvor stolz; und stolzer Muth kommt vor dem Fall.“ Und also hat auch ein heidnischer Poet gesagt: Tolluntur in altum, ut lapsu graviore ruant: Sie werden hoch gehoben, auf daß sie darnach so viel härter fallen können. Item, also steht im 73. Psalm V. 18.: „Du setzest sie auf das Schlüpfrige, und stürzest sie zu Boden.“

22. Da ist keine Gottesfurcht, sondern sie meinen, sie haben dies Glück erlangt durch ihr Verdienst, Tugenden und Beständigkeit, und darum, daß sie dem Herrn Treue und Glauben gehalten, da sie Benjamin mitgebracht; und denken nun nicht anders, denn daß ihnen, als den Frommen und Gerechten, solche Belohnung billig von Gott gegeben werde. Menschliche Natur weiß nicht, wie sie sich nur zieren soll, wenn es ihr wohlgeht, und weiß nicht, daß gern ein Wetter nachfolgt, wenn die Sonne heiß scheint. Dieser Ruhm aber von ihrer Gerechtigkeit und Weisheit verschwindet bald wieder in einem Augenblick.

23. Darum sollen wir lernen Gott fürchten, es gehe uns wohl oder übel, mit solcher Beständigkeit und Gleichmäßigkeit des Herzens, daß wir nicht allein wenn gute Zeit und Glück vorhanden ist, sondern auch wenn böse Zeit ist und daß es nicht wohl zugehet, in unserm Herzen Gott, unserm Herrn, ein Te Deum laudamus (HErr Gott, dich loben wir) singen. Ist Glück vorhanden und daß dir es wohl geht, so magst du dich gebrauchen und Gott dafür danken: ist aber Unglück da und daß dir es übel geht, so rufe Gott an, daß er das Unglück und Jammer wolle wegnehmen, oder aber dasselbe etwas lindern. Ja, wie dir es nur gehen kann, so siehe zu, daß Gott allein dein einiger und höchsten Trost und Zuflucht sei. Denn die Zuversicht und der Ruhm auf das Verdienst eigener Weisheit und Gerechtigkeit muß getödtet und abgethan werden; wie diese Brüder solche Zuversicht gehabt, daß sie durch ihre Tugend oder Frömmigkeit zu solcher großen Ehre und Herrlichkeit gekommen wären.

\begin{quote}
\textbf{V. 7—9. Sie antworteten ihm: Warum redet mein Herr solche Worte? Es sei ferne von deinen Knechten, ein solches zu thun. Siehe, das Geld, das wir funden oben in unsern Säcken, haben wir wiedergebracht zu dir aus dem Lande Canaan. Und wie sollten wir denn aus deines Herrn Hause gestohlen haben Silber oder Gold? Bei welchem er funden wird unter deinen Knechten, der sei des Todes; dazu wollen auch wir meines Herrn Knechte sein.}
\end{quote}

24. Es verlassen sich diese Brüder auf ihre Gerechtigkeit und Verdienst, das Verbrechen, so ihnen vorgeworfen, damit abzulehnen, und wollen dabei ganz und gar entschuldigt sein. Denn so trotzig ist menschliche Gerechtigkeit und Sicherheit, daß sie denkt, Gott sei so närrisch und schwach, daß er solche heilige Leute nicht werde beschuldigen können um irgend einer Sünde willen. Aber sie sollten das Gebet Davids studirt und gelernt haben, da er sagt Ps. 19, 13.: „Wer kann merken, wie oft er fehlet? Verzeihe mir die verborgenen Fehler!“ item, in Ps. 143, 2.: „Gehe nicht ins Gericht mit deinem Knechte“; und das Hiob sagt im 9. Cap. V. 2. 3.: „Ich weiß fast wohl, daß also ist, daß ein Mensch nicht rechtfertig bestehen mag gegen Gott. Hat er Lust mit ihm zu hadern, so kann er ihm auf tausend nicht eins antworten.“ Er soll wohl tausend Sünden auf mich bringen, davon ich kein Wort wüßte.

25. Und solche Leute sind wir alle von unsrer Empfängniß und Geburt, nämlich, verderbt durch die Erbsünde, daß wir tausenderlei Weise verklagt werden mögen, daß wir vor Gott schuldig sind, auch alsdann, wenn wir uns dünken lassen, daß wir sehr fromm und gerecht seien. Und wenn wir wollen hoffärtig sein und nicht in der Furcht Gottes und im Glauben bleiben, so legt er uns vor ihm eine große Last Sünden auf, davon wir nichts gewußt haben; wie David sagt Ps. 90, 8. Und da ist dann kein anderer Schutz oder Schirm, dadurch wir vor Gottes schwerem Zorn und Ungnade geschützt werden mögen, denn allein, daß wir sagen: Erbarm dich mein, o HErr Gott; item, wie im 130. Psalm V. 3. steht: „So du willst, HErr, Sünde zurechnen; HErr, wer wird bestehen?“

26. Aber diese Brüder sind nicht allein sicher der heimlichen Sünde halben, sondern dürfen auch noch ihre Unschuld anziehen in einer sehr bösen Sache. Es hat keine Noth, werden sie gesagt haben, wir wissen uns unschuldig. Ja, sie fallen auch das schwerste Urtheil über sich selbst, da sie sagen: „Bei welchem der Becher funden wird unter deinen Knechten, der sei des Todes“ etc. Dies ist nun die Frucht eigener Gerechtigkeit, welche die Menschen so sicher und halsstarrig macht, vornehmlich aber an den Mönchen und Baaliten. Denn dieselbigen haben nicht einfache Thorheit, sondern sind gar unsinnig. Durch welche Unsinnigkeit sie sich selbst solche große Reinigkeit und Unschuld zuschreiben, daß sie trauen nicht in Demuth auf der Erde kriechen oder sich vor Gott bücken wollen, wie die armen Sünder zu thun pflegen, sondern dürsten ihr Haupt hoch bis an den Himmel hinan aufrichten. Darum werden sie auch mit Donner und Blitz vom Himmel bis in die unterste Hölle hinunter gestoßen.

\begin{quote}
\textbf{V. 10—12. Er sprach: Ja, es sei, wie ihr geredet habt. Bei welchem er funden wird, der sei mein Knecht; ihr aber sollt ledig sein. Und sie eilten, und legte ein jeglicher seinen Sack ab auf die Erde, und ein jeglicher that seinen Sack auf. Und er suchte, und hub am Größesten an, bis auf den Jüngsten; da fand sich der Becher in Benjamin’s Sack.}
\end{quote}

27. Josephs Haushalter richtet seine Sache richtig aus und spottet ihrer fein. Er hat aber das Urtheil, so sie über sich selbst gefällt, welches sehr scharf war, etwas gelindert. Er soll nicht sterben, sagt er; denn das wäre unbillig, daß er sollte getödtet werden um dieses Diebstahls willen, oder daß ihr alle solltet Knechte werden um fremder Sünde willen: sondern bei welchem ich den Becher finde, der soll mein Knecht sein.

28. Da haben sie aber eilends die Säcke abgelegt und aufgethan; denn sie wollen den Haushalter zu Schaden machen, daß er sie fälschlich und mit Unwahrheit des Diebstahls gezogen hätte. Dieweil er aber von seinem Herrn unterrichtet war, hat er seine Untersuchung recht umständlich und wichtig gemacht. Denn er fängt mit dem Aeltesten an, da er wußte, daß der Becher in Benjamins Sack gelegt war, auf daß er ihnen damit Muth mache, und den Wahn, so sie geschöpft hatten, als ob sie gerecht wären und mit ihrer Sache bestehen würden, bei ihnen mehren möchte, wenn sie sehen würden, daß der Aeltesten keiner schuldig wäre. Denn zu Benjamin versehen sie sich gar nichts Böses. Wo ihn nicht Simeon, oder der Andern einer, so zuvor in Egypten gewesen, genommen hat, werden sie gesagt haben, so wird es mit uns gar keine Gefahr haben und alsdann wollen wir diesen Haushalter frei Lügen strafen.

\quellseite{422}

29. Ueber das haben sie auch gesehen, daß eines jeden Geld in seinen Sack gelegt war, und haben dasselbe für einen sonderlichen Segen gedeutet; wiewohl Mose dessen nicht gedenkt. Sie werden es aber ohne Zweifel auf das frühere Exempel gezogen haben, dieweil der Haushalter selbst gesagt hatte: Ihr GOtt, und ihres Vaters GOtt hätte ihnen denselben Schatz gegeben. Derhalben haben sie gedacht, es könnte oder möchte nicht sein, daß er den Becher sollte bei ihnen finden, und er hätte sehr unweislich und unglücklich gethan, daß er sie der Dieberei beschuldigt, da er doch gesehen, daß sie GOtt nun so oftmals mit unermesslicher Gnade gesegnet und reicher gemacht hätte. Also ziehen sie Alles auf ihre Ehre. Wie die Welt mit großer Begierde menschliche und göttliche Wohlthaten annimmt und dieselben ihrem eigenen Verdienste, Gerechtigkeit und Frömmigkeit zuschreibt; wie diese Brüder Josephs sich selbst auch der großen Wohlthaten und Gnade des Fürsten in Egypten ganz würdig achten.

30. Aber sie werden so aufgeblasen und stolz von dem falschen Wahn ihres eigenen Verdienstes und Unschuld, auf daß sie darnach so viel schändlicher haben fallen müssen, darum daß sie so aufgeblasen gewesen. Welches den größten Monarchen, als Alexander dem Großen, Julius Cäsar und Andern mehr, die sich auf ihre Gewalt und Herrlichkeit verlassen, auch also begegnet ist, daß sie plötzlich haben fallen müssen.

31. Wiewohl es hier eine andere Meinung hat mit diesen Brüdern Josephs. Denn GOtt gedenkt nicht Böses wider sie, geht auch damit nicht um, daß er sie verderben wolle, sondern will sie nur also demüthigen und tödten, auf daß die Ehre, Erlösung und Freude, so hernach folgen und ihnen begegnen wird, so viel größer und angenehmer sei. Derhalben hört der Haushofmeister Josephs noch nicht auf, eines Jeden Sack zu besichtigen und Nachsuchung zu thun, da er schon das Geld gefunden, sondern fährt immer fort, den Becher zu suchen, und findet denselben in Benjamins Sack, welcher unter ihnen allen der Unschuldigste war.

32. Derhalben fällt nun bald in einem Augenblick alle ihre Hoffnung, Sicherheit, Freude und Triumph jämmerlich dahin und vergeht ihnen gar, da sie jetzt auf das allerficherste waren und an der Frömmigkeit Benjamins gar nicht zweifeln konnten, und da die Andern alle unschuldig waren. Da ist ihnen plötzlich der Muth gar entfallen, und wird ihnen all ihr Trost und Zuversicht, so sie geschöpft hatten, genommen. Da ist alle Freude aus und so gar hinweg, daß nichts schändlicher konnte hinweg sein. Das thut unser HErr GOtt.

33. Jetzt wollten sie nun lieber, daß der Becher in Rubens oder Simeons Sack gefunden wäre, denn in Benjamins, welcher wider des Vaters Willen, so diese Gefahr zuvorgesehen, nach Egypten geführt war. Und was können wir denken, wie Juda zumal muß zu Muthe gewesen sein, welcher für seinen Bruder Benjamin vor Bürge geworden? Sie verstehen aber dies Spiel nicht, thun auch nichts dazu, sondern werden nur also umgetrieben und wunderbarer Weise geführt. Denn es ist uns viel nützlicher, daß wir von GOtt getrieben und geführt werden, denn daß wir von uns selbst etwas thun, verstehen, vorsehen und ordnen nach unserm eigenen Rath oder Gutdünken. Unser Leiden ist das allerheiligste Leben. Also sind nun diese Brüder bald in einem Augenblick aus dem Himmel in die Hölle, aus dem Leben in den Tod jählings gestoßen und getrieben worden.

\begin{quote}
\textbf{V. 13. 14. Da zerrissen sie ihre Kleider; und lud ein jeglicher auf seinen Esel, und zogen wieder in die Stadt. Und Juda ging mit seinen Brüdern in Josephs Haus, denn er war noch daselbst; und sie fielen vor ihm nieder auf die Erde.}
\end{quote}

34. Nun triumphiren sie nicht mehr, sondern zerreißen ihre Kleider, und ist ihnen über die Maßen angst und bange geworden; kommen also wiederum in Josephs Haus, da er wohl bedächtig geblieben war, eben an dem Ort, da sie gegessen hatten, und stellt sich nun auch, als sei ihm der Becher gestohlen worden, welches doch unmöglich war. Darum spielt er jetzt etwas gröber mit ihnen; hernach aber wird er sich seinen Brüdern offenbaren und zu erkennen geben. Sie aber waren so furchtsam und erschrocken, daß sie gar nichts haben verstehen können.

35. Juda war für Benjamin Bürge geworden, darum ist er zuerst vor Joseph auf die Erde niedergefallen; und die andern mit ihm sind auch gar voll Verzweiflung, welche kurz zuvor von Freuden gleichsam trunken gewesen waren, und sich fein sanft träumen ließen, daß sie im Paradies und im Himmel wären, da sich der Fürst mit solcher großen Gnade gegen sie erzeigt hatte. Denn sie fühlen und empfinden nun, daß eben derselbe jetzt ihr Feind ist, und ein solcher Tyrann, der sie könne und wolle in ewige Knechtschaft bringen. Also sind sie, ihres Trachtens, gar verdammt gewesen. Und das heißt, die Hoffärtigen recht demüthigen und die Gewaltigen von ihren Stühlen stoßen.

36. Denn wie mögen wir denken, daß sie werden gesinnet gewesen sein, da sie wieder sind zu Joseph gekommen? Wie werden sie mit so großem Schmerz und Traurigkeit Simeon und Levi, ihre Vorgänger und Rädelführer, beschuldigt und zur Rede gestellt haben, die dies Unglück alles gestiftet und dies Feuer angezündet hatten, da sie ihren Bruder Joseph verkauft haben! Denn damit werden sie die Zeit auf dem Wege zugebracht haben, daß sie einander solches also aufgerückt und sie zur Rede gestellt haben. Da werden Simeon und Levi sehr greulich erschrocken sein und jämmerlich geheult haben, und haben die That nicht können läugnen. Ach! werden die Andern gesagt haben, warum haben wir doch eurem Rath gefolgt? Was wollen wir nun thun? Sollen wir nun wiederum zu unserm Vater ziehen und Benjamin hier bleiben lassen? Wie, wenn er uns noch alle gefänglich würde lassen einziehen? Denn sie haben in ihrem Herzen nicht anders gewußt, denn daß sie alle würden sterben müssen. Und darum haben sie auch ihren Zorn auf einmal gar herausgeschüttet auf Simeon und Levi. Sie werden sich weidlich auf dem Wege untereinander gescholten haben. Da wird ihnen eingefallen sein die jämmerliche und letzte Klage Jakobs, daß sie ihn aller seiner Kinder berauben würden. Da sind ihre Herzen übergegangen von großem Leid. Da wird sie der Heulung angekommen sein. Und haben nun Schmerzen gehabt über die Sünde, so sie begangen hatten; welches doch noch die rechte Reue nicht ist; denn sie sehen jetzt nur ihren eigenen Schaden und Gefahr an.

37. Und ich halte zwar dafür, daß Simeon und Levi daselbst die Alleredelsten, Verachtetsten und gar verzagt gewesen seien, darum daß die andern Brüder sie hart angefahren und zur Rede gestellt, und ihnen das oft werden vorgeworfen und aufgerückt haben. Ihr tödtet unsern Vater und wir allesamt mit euch. Wer will nun unsern Bruder Benjamin dem Vater wiederum zustellen? Oder, welcher unter uns allen wird hinfort dem Vater dürfen vor seine Augen kommen? Ach GOtt! ach GOtt! werden sie gesagt haben, was habt ihr angerichtet? Auf diese Weise wird sie nun Joseph, der fromme heilige Patriarch, zu rechter herzlicher Reue und vollkommenem Buße und Demuth bringen; nicht daß er sie damit wolle verderben, sondern vielmehr, daß er sie bessere, und ihre Sünde strafe, auf daß sie dieselbe recht erkennen und Buße thun, sich auch hinfort vor dergleichen Sünden fleißig hüten.

\section[Zweiter Theil]{Zweiter Theil. Wie Joseph mit seinen Brüdern hart redet und aus ihm seine Brüder antworten. Item, von der Rede Judas}

\subsection{I}

\begin{quote}
\textbf{V. 15. Joseph aber sprach zu ihnen: Wie habt ihr das thun dürfen? wisset ihr nicht, daß ein solcher Mann, wie ich bin, errathen könnte?}
\end{quote}

38. Er sagt nicht, daß er ein Weissager sei, sondern redet, was ihm zu seiner Sache, die er vorhatte, diente. Denn Egypten ist allezeit zu mancherlei Aberglauben geneigt gewesen, und hat damit die rechte wahre Religion mannigfaltig verwirrt und gehindert. Und Joseph hat nicht können alle Irrthümer ausrotten; wie auch die andern Propheten und Apostel nicht gekonnt haben. Und wird nimmermehr ein Prediger das alles ändern oder abschaffen, was in der Kirche irrig oder gebrechlich ist. Wir sammeln Etliche aus der Welt durch das Predigtamt; das Netz wird in das Meer geworfen und sammelt zugleich Böse und Gute. Darum ist daran kein Zweifel, es wird das Regiment, so Joseph geführt, sehr gut und löblich gewesen sein, damit er Vieles gebessert und wiederum zurecht gebracht hat; aber doch also, daß er dem Satan sein Theil auch hat lassen müssen; sintemal nie ein Prophet die alle, welche er gelehrt, hat bekehren können. Und ist wohl zu glauben, daß damals viel Zauberer in Egypten hin und wieder umhergegangen sind.

\quellseite{423}

39. Derhalben spielt nun Joseph also, und stellt sich, als sollten sie ihn billig für einen solchen weisen Mann gehalten haben, der errathen könnte. Als wollte er sagen: Ihr sollet billig gedacht haben, daß in Egypten viel Weissager, Wahrsager und Zauberer wären, und daß ja vornehmlich der Fürst selbst in solchen Künsten müsse erfahren sein etc. Also stellt er sich, als wisse er um ihren heimlichen Rath, und durch solch Stellen wird ihr Schrecken also gemehret, daß sie schier vor Schmerzen sterben möchten.

\begin{quote}
\textbf{V. 16. Juda spricht: Was sollen wir sagen meinem Herrn, oder wie sollen wir reden, und was können wir uns rechtfertigen? GOtt hat die Wissenschaft deiner Knechte funden. Siehe da, wir und der, bei dem der Becher funden ist, sind meines Herrn Knechte.}
\end{quote}

40. Hitzadek ist ein Zeitwort der vierten Conjugation; wie das auch, so in 34. Psalm V. 3. steht: Im Herrn wird meine Seele gerühmet, daß ist: Meine Seele soll sich rühmen des Herrn. Wir können, sagen sie, nichts vorwenden, uns zu rechtfertigen; wir sind eigentlich schuldig und bekennen unsere Sünde. Und dies ist nun der andere Theil dieses Spiels. Joseph hat sie jetzt gebracht bis auf den letzten Kampf, daß sie nun sagen: Wir sind Diebe und böse Buben; wie können wir uns doch reinigen oder mit irgend einem Schein entschuldigen? Wenn wir sagen wollten, der Diebstahl sei in deinem Hause geschehen, so werden wir damit entweder dich selbst oder aber deine Hausgenossen beschuldigen. Und sie sind doch unschuldig, und wissen, daß ihnen Unrecht geschieht; aber sie dürfen und können über ihrer Unschuld nicht hart halten und sich selbst damit schützen; können auch nichts Anderes anwenden, denn daß sie sagen: Wir haben gesündigt, GOtt hat unsere Missethat gefunden. Dürfen also weder GOtt noch Joseph beschuldigen.

41. Und dies ist nun endlich das Bekenntniß der Sünde, welches ihnen Joseph gleichsam mit Gewalt abgedrungen hat. Denn sie ergeben sich nun gar zur Knechtschaft, und bekennen mehr, als sie gethan haben, auf daß sie mit solcher Demuth sein Herz etlichermaßen erweichen mögen, welches Gnade sie zuvor auch erfahren hatten.

\begin{quote}
\textbf{V. 17. Er aber sprach: Das sei ferne von mir, solches zu thun. Der Mann, bei dem der Becher funden ist, soll mein Knecht sein; ihr aber ziehet hinauf mit Frieden zu eurem Vater.}
\end{quote}

42. Joseph erzeigt sich sehr hart gegen seine Brüder und gibt sonderliche Gerechtigkeit vor. „Das sei ferne von mir“, sagt er, daß alle sollten meine Knechte sein; ich begehre nichts weniger denn das. Denn ich beschuldige euch ja nicht alle, sondern begehre nur, daß der mein Knecht sein soll, bei dem der Becher gefunden ist.

43. Derhalben ist dies ein sehr fein Spiel, und ein schön vortrefflich Gedicht dieses Poeten, desgleichen ich nicht hätte können dichten; denn ich habe solchen großen Verstand oder Geist nicht. Daraus erhellt nun, daß Joseph ein sehr vortrefflicher Mann und herrlicher Theologe gewesen ist; wiewohl er die Bücher der Propheten noch nicht gehabt, und als ein Jüngling von siebenzehn Jahren seinem Vater aus dem Schooße gerissen worden. Doch hat er sehr treulich und mit höchstem Fleiß die rechte göttliche Lehre von seinem Vater gelernt; und hat nicht allein die Worte seines Vaters von dem verheißnen Samen behalten, sondern hat auch seinen Geist empfangen, und also, daß er auch täglich durch vielerlei Anfechtung, Trübsal und Leiden immer zugenommen hat.

44. Dieweil nun also die gute Art und der heilige Geist zusammen gekommen, hat daraus ein großer trefflicher Poet werden müssen. Und ist daran kein Zweifel, er wird dergleichen Spiel oder seine Gedichte viel mehr, und sonst auch andere kluge Werke die ganze Zeit seiner Regierung verrichtet haben, welche eben in diesem Buch nicht geschrieben sind. Aber an diesem Einen Exempel ist wohl abzunehmen, was er für ein Prophet gewesen sei und welchen hohen Verstand und großen Geist er müsse gehabt haben. Zu welchen Gaben auch noch gekommen ist der Gebrauch und die tägliche Uebung vieler Dinge, so er in seinem vielfältigen Jammer und Plagen erfahren hat.

45. Darum spielt er nun dies Spiel fein freundlich und führt seine Brüder in Verzweiflung, Tod und Hölle hinein; und da Alles verloren war und sie nun gar verzagt, kommt das Mittel in das Spiel, wie es in den Komödien zu gehen pfleget, und treibt alle Gefährlichkeit weg, macht Alles wieder gut. Joseph ist ein sehr weiser Mann gewesen, und wohl erfahren in Sachen, die den Glauben und die Seligkeit belangen; und hat doch die Bücher nicht gehabt, welche jetzt die Kirche hat, und hat Mose nicht ganz gelesen, auch Jesaia nicht: etliche Historien der Väter ausgenommen bis auf Jakob, sonst hat er nichts gehabt. Darum ist dies eine große treffliche Gabe der Weissagung gewesen.

46. Darnach ist auch zu merken, daß man dies Spiel nicht allein auf die göttliche Regierung ziehen soll, daraus zu lernen, wie er seine Heiligen zu führen pflege; sondern daß es auch zu unserm Trost diene, nämlich, wenn die Sachen so gar böse und verdorben sind, daß man keine Hoffnung mehr sieht, wie denselben geholfen oder gerathen werden möchte, daß man alsdann wissen soll, es sei damit nun auf das Höchste im Spiel gekommen, daß es sich bald ändern und zum fröhlichen Ende kommen werde. Denn also ist es um die göttlichen Gedichte beschaffen, wie Paulus fein sagt Eph. 2, 10.: „Wir sind unseres Herrn Gottes Werk“ (Gedicht). Er ist der Poet, wir sind die Verse und Gedichte, die er macht. Derhalben ist daran kein Zweifel: unser Werk und alle unser Thun ist vor GOttes Augen angenehm um sonderlicher Kraft und Gnade des Glaubens willen.

47. Auf diese Weise soll man die Historien der heiligen Väter handeln, welche die Mönche und die der geistlichen Dinge nicht erfahren sind, für gar vergeblich und unnütz halten. Denn sie verwundern sich, daß Joseph mit dem Spiel also hat umgehen mögen und die gute Zeit so übel angelegt hat, welcher lieber die Zeit soll gefasset und gebetet haben. Was ist es, sagen sie, daß so ein weiser Mann, und der voll heiligen Geistes ist, die Zeit so vergeblich damit zubringt, daß er seiner Brüder also spottet und täuscht, daß er den unschuldigen Benjamin und die Andern auch so jämmerlich plagen und martern mag, und sie mit solchen Possen und Scherz gar irre macht und bekümmert?

48. Aber das ist es, daß ich oft gesagt habe: Wer an GOtt glaubt und ihn von Herzen dient, und sich davor hütet, daß er ihn weder mit Worten noch mit Werken erzürne, desselben Werke alle, ja, auch sein Schlafen und Wachen, dazu auch sein Spielen, geschehen im Namen des Herrn und sind GOtt angenehm. Aber irgend ein Anderer, der ohne Geist und Glauben dahin lebt und GOtt nicht fürchtet, derselbe wird aufgeblasen und stolz wegen seiner Gerechtigkeit und Heiligkeit und erwählt große Werke, und hat einen Abscheu vor diesen Spielen und Werken der Heiligen, darum daß sie seines Erachtens allzu gemein, bürgerlich, fleischlich und weltlich sind. Er fliehet auch vor dem Hausregiment, vor der Polizei und Kirchenämtern. Er lebt ohne Ehe, fastet und tödtet seinen Leib. Und endlich versucht er, wie er nur vor allen Creaturen fliehen möge. Wie man sagt von einem Waldbruder, der die Sonne nicht hat anschauen wollen, darum daß sie ein leiblich Licht wäre; aber zuletzt hat er im Gesicht gesehen, daß die Hölle voll Mönche gewesen ist. Vor solcher tollen Heiligen Aberglauben hat GOtt ohne allen Zweifel einen Greuel und verdammt denselben.

49. Derhalben, die da wollen gottselig leben, sollen GOtt fürchten und ihm vertrauen, und darnach ihres Berufs warten; alsdann werden sie genug zu thun finden. Sie sollen dem Herrn ihren Weg befehlen, Abends und Morgens, sollen im Namen des Herrn schlafen und wiederum aufstehen, und thun, was ihnen vor die Hand kommt, es sei in welchem Stande es wolle. Wie Samuel sagt zu Saul, 1 Sam. 10, 7.: „Thue, was vor deiner Hand kommt; denn GOtt ist mit dir.“ So muß dann Alles wohl gethan sein, und alle deine Werke, auch die scherzhaften, werden vor GOtt und seinen Engeln eitel seine lustige Schauspiele sein. Wie denn dies Spiel GOtt auch wohlgefallen hat. Denn sonst würde der Heilige Geist dasselbe nicht also gerühmt haben.

\quellseite{424}

50. Denn wir sind ja nicht geschaffen zum einsamen Leben, daß wir sollen abgesondert sein von der Gemeinschaft anderer Menschen; sondern daß wir gemeine und gewöhnliche Werke dieses Lebens thun sollen mit solchem Fleiß und Lust, wie mir am allerbesten können; wie Salomo sagt in seinem Prediger Cap. 9, 7.: „So gehe hin und iß dein Brod mit Freuden, trink deinen Wein mit gutem Muth; denn dein Werk gefällt GOtt“ etc. Und darum hat auch GOtt einem Jeden sein Amt und Stand gegeben und was er darin für Werke thun soll, auf daß uns nicht sollte frei sein, mit selbsterwählten Werken nach unserer eigenen Andacht und wie es uns gelüstet, umzugehen.

51. Lasset uns aber nun sehen, was Joseph geantwortet habe? „Das sei ferne von mir“, sagt er, „solches zu thun. Der Mann, bei dem der Becher funden ist, soll mein Knecht sein.“ Da trifft er ihnen das Leben. Er nennt es nun keinen Diebstahl, sondern sagt: „Bei dem der Becher funden ist“ etc.; denn es hat nur das Ansehen eines Diebstahls. Die Andern aber hat er absolvirt, los und ledig gesprochen, daß sie mit Frieden zu ihrem Vater ziehen möchten. Dasselbe ist zwar gnädig und freundlich genug geredet; es ist aber doch gar ein unleidliches Wort im Herzen dieser Brüder und vornehmlich Judas gewesen, welcher für Benjamin war Bürge geworden. Wenn der Becher bei Simeon oder sonst bei einem Andern, wer es auch hätte sein mögen, wäre gefunden worden, so hätte man diesem Unglück etwas leichter rathen können: aber hier kommen sie nun in Gefahr Leibes und Lebens, da geht das Wasser über die Körbe; es ist Alles aus und wird sich ein jämmerlich Heulen und Wehklagen erheben.

52. Darum hat nun jetzt Joseph dies Spiel zu dem Ende gebracht, welches er im Anfang vorgenommen. Denn dies ist der rechte Kampf und ist nun mit dieser Historie auf das Höchste gekommen. Jetzt wird er sie nun dahin bringen, daß sie gar verzagen werden, also daß kein Rath oder Hülfe mehr ist, sondern sie müssen selbst verderben und ihren Vater nun gar tödten. Hier ist es Zeit, daß der Strick breche.

\subsection{II}

\begin{quote}
\textbf{V. 18. Da trat Juda zu ihm, und sprach: Mein Herr, laß deinen Knecht ein Wort reden vor deinen Ohren, mein Herr; und dein Zorn ergrimme nicht über deinen Knecht; denn du bist wie Pharao.}
\end{quote}

53. Dies ist die Rede Judas, welche sehr heftig und brünstig gewesen; dieselbe wird er nicht gethan haben ohne viel Thränen und Seufzen. Und ich würde zwar nicht also reden können, kann auch solchen Affect und herzliche Bewegung nicht genugsam auslegen oder erklären, damit er diese Worte vor Joseph ausgeschüttet hat. Denn siehe, was ihn für Ursachen dazu getrieben haben. Er hatte dahin seinen alten Vater sitzen lassen; aber nicht der Hoffnung und daß er dessen hätte sollen gewärtig sein, daß sie alle, oder ja zum wenigsten der Jüngste, welchen er heftig lieb hatte, in Egypten sollten Knechte werden. Nachdem er aber sieht, daß Benjamin wider ihrer aller Meinung in solche jämmerliche Knechtschaft gerathen und gestoßen werden und nimmermehr wieder zu seinem Vater kommen soll, da wird er so traurig und wehmüthig, daß in seinem Herzen nichts denn eitel Jammer und Verzweiflung gefunden worden; und hat ihn solcher Jammer und Herzeleid gezwungen, daß er, dieweil er geredet hat, viel bittere Thränen wird vergossen haben. Darum hat er diese Worte und ein jegliches insonderheit aus rechtem Grunde seines Herzens geredet.

54. Er tritt aber von selber hinzu, ungefährdet, ja, es hätte ihn Joseph schon abgewiesen; denn er hatte zu ihnen gesagt: „Ziehet ihr hinauf mit Frieden zu eurem Vater.“ Aber die äußerste Noth und daß er in seinem Herzen so heftig bewegt worden ist, hat ihn dazu getrieben, daß er nicht gesehen oder gedacht hat, was ihm wohl oder übel anstünde. Zuvor war er furchtsam, demüthigte sich, fiel vor Joseph nieder auf die Erde: jetzt aber hat er es nicht einmal die Kniee gebeugt, sondern tritt stracks zu ihm; denn der große Unfall und Jammer hat ihn gar überwunden.

55. Ach, daß wir doch auch mit solcher großen Inbrunst GOtt anrufen könnten, wie dieser Juda vor seinem Bruder betet! Ich wollte viel darum geben, daß ich vor unserm Herrn GOtt so wohl beten könnte, als hier Juda vor Joseph betet. Denn dies ist ein recht vollkommenen Exempel des Gebets, ja, es ist das rechte herzliche Verlangen oder Begierde, das im Gebet sein soll. Dieser Juda ist kühn und zuversichtlich, und gedenkt nicht der Gewalt oder Macht, so Joseph gehabt, noch seiner eigenen Niedrigkeit. Das Gebet und die große Noth hat ihn kühn gemacht, daß er vor Joseph ohne alle Furcht ganz unerschrocken getreten ist.

56. Wenn wir auch so vollkommen beten könnten, sonderlich im Namen unsers Heilandes und Mittlers Jesu Christi, so wäre es unmöglich, daß ein solch Gebet nicht sollte erhört werden. Wie sich Joseph jetzt nun nicht länger wird enthalten können, er muß weinen und sich nun recht erklären. Denn wir werden zu unsern Zeiten auch viel mit Jammer und Trübsal geplagt, vornehmlich im Predigt- oder Kirchenamt, da immer eine Anfechtung auf die andere folget, und wenn eine aus ist, folgen bald sieben andere, die weit schwerer sind.

57. Nun soll man aber vor allen Dingen gewiß wissen, daß wir das liebe Wort haben; denn das ist der rechte Grund, darauf wir uns verlassen, daß wir Zuhörer sind und daß GOtt mit uns redet. Daran soll niemand zweifeln. Denn wer das nicht weiß oder zweifelt, der wird auch wohl mit Beten ein lang Geplapper machen und mit dem Maul viel murmeln, wie die Heuchler pflegen; aber recht beten kann er nicht.

58. Wo aber dieser Grund gelegt ist, welcher GOttes Wort ist, da ist dann das Gebet unsere äußerste Hülfe, ja, nicht Hülfe, sondern unsere Kraft und Sieg in aller Anfechtung und Noth; wie GOttes Wort selbst lautet, da er sagt Ps. 50, 15.: „Rufe mich an in der Noth, so will ich dich erretten“; und Jes. 65, 24.: „Und soll geschehen, ehe sie rufen, will ich antworten; wenn sie noch reden, will ich hören.“ Und der Engel Gabriel sprach zu Daniel, Dan. 9, 23.: „Da du anfingest zu beten, ging dieser Befehl aus.“

59. Denn wo das Herz ernstlich und brünstig betet, ist es unmöglich, daß solch Gebet von GOtt nicht sollte erhört werden, sobald es nur angefangen wird. Dasselbe kann ich auch wahrhaftig sagen und bezeugen. Ich mag zwar von meinem Gebete nicht groß rühmen; aber daran ist kein Zweifel, daß das Wort des Engels Gabriel, Dan. 9, 23., wahrhaftig sei; wiewohl wir nicht allezeit mit solcher herzlichen Begierde und Inbrunst beten können. Das unaussprechliche Seufzen ist nicht allezeit da; denn es ist auch nicht in unserer Gewalt; sondern wenn der Geist in uns schreit: Abba, lieber Vater, da schweigt wohl der Mund ganz stille und kann dies Seufzen nicht verstehen oder ausreden. Aber wenn das Rufen des Geistes erst herausbricht, wie Jesaia sagt, ist nichts gewisser, denn daß solch heftig Gebet schon erhört ist. Und wenn sich auch der Türke, Papst, Kaiser und alle Pforten der Höllen dawider legten, sollten sie doch nichts ausrichten, wenn die Kirche dies unaussprechliche Seufzen zu Gebet könnte bringen. Denn wir streiten wider den Teufel und die höllischen Pforten. Was ist aber Fleisch und Blut an den Feinden der Kirche? Sintemal wir den Teufel mit dem Gebet tödten können, warum sollten wir nicht auch den Türken und Papst vertreiben können? Denn also sagt Jesaia Cap. 31, 3. von Egypten: „Egypten ist ein Mensch, und nicht GOtt, und ihre Rosse sind Fleisch und nicht Geist.“

60. Eben dasselbe soll man von unsern Widersachern und Verfolgern auch halten; wie wir denn ja bisher ganze vier und zwanzig Jahre mit unserm Gebete die feindselige Gewalt des Kaisers, Türken und Papstes etc. aufgehalten haben. Aber die Kirche ingemein oder die Heiligen und Gläubigen insonderheit sind nicht allezeit so entbrannt mit Inbrunst eines solchen Gebets. Wie die Theologen ein gemein Sprüchwort haben, daß sie sagen: Spiritus Sanctus non semper tangit corda prophetarum, das ist: Der Heilige Geist rührt die Herzen der Propheten nicht allezeit. Die Erleuchtungen der Propheten währen nicht immer für und für, ohne Aufhören. Wie Jesaia nicht immer und stets auf einander Offenbarungen von hohen großen Dingen gehabt, sondern nur auf sonderliche Zeit.

61. Dasselbe zeigt auch an das Exempel des Propheten Elisa, da er von der Sunamitin sagt 2 Kön. 4, 27.: „Laß sie, denn ihre Seele ist betrübt, und der Herr hat mir's verborgen, und nicht angesagt.“ Da bekennt er, daß GOtt nicht allezeit die Herzen der Propheten rühre. Es ist auch wohl bisweilen der Geist gekommen, wenn sie auf der Harfe oder Psalter gespielt und etliche Psalmen und geistliche Lieder gesungen haben.

62. Und dies sage ich als zum Exempel, nämlich, daß das Gebet nicht allezeit so heftig ist, wie es zu sein pflegt, wo der Heilige Geist das Herz sonderlich regt. Und ist doch dasselbe eine sehr gute Hülfe und Rath, das Gebet zu erwecken, auf daß es brünstig und heftig werden möge, nämlich, daß wir gerne lesen, singen und GOttes Wort hören, und daß wir ja den Psalter fleißig lesen, oder auch in die gemeine Versammlung der Kirche gehen: da wird denn das Herz allemal wach werden und der Geist in uns angezündet werden.

\quellseite{425}

63. Denn das ist die gemeine Regel und Gebrauch, daß der Heilige Geist die Herzen nicht allezeit rührt mit dem unaussprechlichen Seufzen; wenn er sie aber damit anrührt, alsdann ist es schon beschlossen und ist gewiß, daß das Gebet durch die Wolken gedrungen, und dasjenige, darum man bittet, im Himmel und auf Erden erlangt hat, oder Himmel und Erde müßte eher brechen.

64. Darum ist das ein trefflich Gebet, daß man sich erstlich das Wort und darnach die hochdringende Noth vorstellt und dieselben im Herzen erwägt. Wenn uns der Teufel mit Gewalt treibt und die Welt uns dazu plagt, je greulicher wir alsdann bedrängt werden, desto brünstiger beten wir. Denn das Herz des Gläubigen wird heftig bewegt und wird darin das Geschrei und unaussprechliche Seufzen erweckt, und das Gebet wird alsdann auch bald, wenn es erst angefangen ist, erhört; wie der Engel sagt zu Daniel Cap. 9, 23.: „Da du anfingest zu beten, ging dieser Befehl aus.“

65. Also wird an diesem Orte beschrieben ein Exempel eines heftigen Gebetes, so Juda gethan, welcher der großen Majestät und Gewalt Josephs gar vergißt; und da ihn die Noth gedrungen, tritt er stracks hinzu und betet ohne alle Ehrerbietung, die er Joseph hätte erweisen sollen.

66. Also geht es mit uns auch zu. Wenn Glück vorhanden ist, und daß es uns wohl geht, denken wir wohl an die hohe Majestät und Ehre GOttes, welchen wir anrufen, wissen uns auch zu erinnern unserer eigenen Unwürdigkeit und fürchten uns vor seiner Gewalt; darum ist alsdann das Gebet nicht so stark und brünstig. Aber wenn wir in der Noth und Anfechtung sind, da haben wir nicht sonderliche Acht auf die hohe Majestät, sondern sagen stracks: Hilf, lieber GOtt! nun hilf GOtt! laß dich das erbarmen im Himmel etc. Da machen wir keine lange Vorrede. Also macht die äußerste Noth kühne und zuversichtliche Leute, dieselben brechen hindurch und siegen auch endlich.

67. Wenn wir nun auf solche Weise beten könnten, wollten wir leichtlich der ganzen Welt Drohen verachten. Wie wir denn hie daher aller unserer Feinde Gewalt und Stürme ertragen haben; und wiewohl wir im Gebete kalt gewesen, ist doch die Kraft des Gebets immer vor GOtt.

68. Es ist zwar der grimmige Zorn Behemoths und aller Teufel so groß über uns, daß es niemand glauben kann; doch können sie uns nicht unterdrücken. Wo ihnen aber GOtt einen Augenblick den Zaum lassen würde, so würden sie bald alle unsere Weisheit, unsere Kraft und Wehre verschlingen. Denn unsere eigene Hülfe oder Wehre und Rath ist wie eine geringe Gestalt, oder äußerliche Larve, damit GOtt der Widersacher spottet, daß sie meinen, es sei zumal eine große Stärke, und fürchten sich vor derselben Larve. Warum aber das? Darum, daß dies unsers Herrn GOttes Titel ist, der den Fürsten den Muth und ihre Stärke nimmt. Psalm 76, 13. Wer richtet aber dasselbe aus? Antwort: Der da glaubt, der das Wort hat, der da betet und GOtt anruft, das ist, der Geist der Gnade und des Gebets selbst, Sach. 12, 10.

69. Wie Juda hier voll Inbrunst hinzu tritt, ohne alle Furcht; denn es treibt ihn die äußerste Noth. Darum denkt er: Siehe, mein Vater Jakob wartet daheim mit großem Verlangen darauf, daß seine Söhne bald wiederkommen werden. Er hat vorhin Joseph verloren, jetzt steht nun Benjamin auch in Gefahr Leibes und Lebens. Ich, als ein Bürge, werde von meinem Vater hoch getödtet werden; er aber, der Vater selbst, wird vor großem Schmerz und Leid auch sterben müssen. Also ist Juda der Allerarmseligste und Betrübteste, und dazu aller Hoffnung und Hülfe gar beraubt gewesen.

70. Darum sagt er nun auf Hebräisch, bi adoni: Mein Herr. Das Wörtlein bi ist ein Ausruf, heißt an andern Orten so viel als, in mir; hier aber ist es ein Wort, damit man fleht und bittet. Wie denn dasselbe Wort auch die Hure vor dem Könige Salomo braucht, 1 Kön. 3, 26., da sie sagt: „Ach, mein Herr.“ Juda ist erschrocken und gar betrübt, darum schreit er so heftig und wiederholt in die Länge die ganze Historie. Und Joseph hört dasselbe an nicht ohne bittere Thränen und Bewegung seines Herzens; wiewohl er solche Bewegung noch immer verhehlt. Bis daher aber wäre die Vorrede, welche zwar sein demüthig gewesen und die doch solcher großen Majestät Josephs nicht geziemet hat.

\begin{quote}
\textbf{V. 19—24. Mein Herr fragte seine Knechte, und sprach: Habt ihr auch einen Vater, oder Bruder? Da antworteten wir: Wir haben einen Vater, der ist alt, und einen jungen Knaben, in seinem Alter geboren; und sein Bruder ist todt, und er ist allein übriggeblieben von seiner Mutter, und sein Vater hat ihn lieb. Da sprachst du zu deinen Knechten: Bringet ihn herab zu mir; ich will ihm Gnade erzeigen. Wir aber antworteten meinem Herrn: Der Knabe kann nicht von seinem Vater kommen; wo er von ihm käme, würde er sterben. Da sprachst du zu deinen Knechten: Wo euer jüngster Bruder nicht mit euch herkommt, sollt ihr mein Angesicht nicht mehr sehen. Da zogen wir hinauf zu deinem Knechte, meinem Vater, und sagten ihm an meines Herrn Rede.}
\end{quote}

71. Benjamin wird genannt ein Sohn des Alters, darum daß ihn Jakob in seinem Alter gezeuget hatte. Denn wiewohl er die Andern auch gezeuget, da er schon alt an Jahren gewesen, hat er sich doch der zwei, so von dem vornehmsten Weibe, Rahel, geboren waren, sonderlich angenommen. Warum wird aber Benjamin ein junger Knabe genannt, so er doch da mit den andern Brüdern von Joseph gestanden und schon zu zehn Kindern Vater gewesen? Das geschieht darum, daß er der Jüngste ist unter ihnen allen und ganz zärtlich vom Vater geliebet worden ist. Dazu ist er zwölf oder sechzehn Jahre später geboren, nachdem die Andern alle auf die Welt geboren gewesen. Da Joseph verkauft worden, hat er noch in der Wiege gelegen, nachdem seine Mutter, Rahel, in der Geburt jämmerlich gestorben war. Darum, weil er in solchem Jammer und Trübsal geboren ist, wird er billig ein junger Knabe genannt.

72. Er gedenkt aber auch Josephs, Benjamins Bruder, welcher doch da gegenwärtig, ihnen aber noch unbekannt ist. Und freilich hat er dies nicht anhören können ohne große Bewegung und Affect seines Herzens, vornehmlich da die Rede Judas voll Seufzens und Thränen gewesen, dadurch er gehindert worden, daß er nicht so wohl hat reden können. Doch aber hält sich Joseph noch immer standhaft.

73. Juda erwähnt auch die verstorbene Mutter, nicht anders, denn als ob sie noch am Leben wäre, auf daß er damit das Herz Josephs desto mehr bewegen möge. Denn also lauten seine Worte von Benjamin: „Und er ist allein übriggeblieben von seiner Mutter.“ Warum fragen etliche Andere: Warum doch Juda dies möge gesagt haben? Über diese Auslegung ist nach der Grammatik sein schlicht und leicht, wenn man den Dativ in einen Ablativ verändert: Relictus est de matre sua: „Er ist von seiner Mutter überblieben“, wie unser Deutsch lautet. Derhalben sind diese Worte von einem Knaben sehr klüglich geredet, nämlich, er ist der Kleinste, der Jüngste, ist von seiner Mutter allein überblieben; sie aber ist jämmerlich dahin gestorben. Und ist dazu sein Bruder auch umgekommen, er ist allein noch übrig, ist allein des Vaters Lust. Dies, sage ich, ist eine bewegliche und heftige Rede, voll höchster Affecten und Neigungen.

74. Und Juda hat mit dieser seiner Rede nicht allein wegen seines Bruders, sondern auch wegen der großen Liebe und Sorge des Vaters gegen denselben Joseph sehr bewegt. „Wo der Knabe“, sagt er, „seinen Vater soll verlassen, so wird der Vater sterben.“ Wir aber haben dir Gehorsam geleistet. Darum bitte ich, du wollest dich unser erbarmen, dieweil wir mehr gethan haben, denn wir hätten thun sollen.

75. Ja, daß noch mehr ist, er führt auch Josephs Befehl, sprechend an: Du hast uns geheißen, daß wir ihn sollten zu dir bringen, auf daß du ihn sehen möchtest und Gnade erzeigen. Dies aber wird kein Anzeichen sein der Gnade oder Gunst, nämlich, dem Vater seinen Sohn, welchen er sonderlich lieb hat und den er mit großem Wehklagen sehr ungern hat wegziehen lassen, mit Gewalt nehmen und ihn desselbigen gar berauben. Und wir hätten zwar dasselbe von unserm Vater nimmermehr bitten oder begehren dürfen, wo du uns solches nicht mit solchem strengen Gebot eingebunden hättest, daß wir uns dazu gegen dich haben verpflichten müssen. Denn also sprachst du zu deinen Knechten: „Wo euer jüngster Bruder nicht mit euch herkommt, sollt ihr mein Angesicht nicht mehr sehen.“ Darum haben wir uns fleißig nach deinen Worten und Befehl gehalten und uns auf deine Zusage und Befehl verlassen; jetzt aber finden wir nun, daß du so ungnädig mit uns und unserm alten Vater umgehst; daß wir uns je nicht versehen hätten, auch nicht verdient haben, daß du uns nun alle gar vertilgen und todt haben willst.

\quellseite{426}

\begin{quote}
\textbf{V. 25—29. Da sprach unser Vater: Ziehet wieder hin, und kaufet uns ein wenig Speise. Wir aber sprachen: Wir können nicht hinab ziehen, es sei denn unser jüngster Bruder mit uns, so wollen wir hinab ziehen; denn wir können des Mannes Angesicht nicht sehen, wo unser jüngster Bruder nicht mit uns ist. Da sprach dein Knecht, mein Vater, zu uns: Ihr wisset, daß mir mein Weib zween Söhne geboren hat; einer ging hinaus von mir, und man sagte, er ist zerrissen; und hab ihn nicht gesehen bisher. Werdet ihr diesen auch von mir nehmen, und ihm ein Unfall widerfähret, so werdet ihr meine grauen Haare mit Jammer hinunter in die Grube bringen.}
\end{quote}

76. Da gedenkt er abermals des Vaters, und erzählt, was der ihnen vorgeworfen; desgleichen seine Sorge und Bekümmerniß, dazu seinen Jammer und alle seine Worte, auf daß er damit Joseph erweichen und bewegen möge. Er hat nichts ausgelassen oder vergessen, das ihm hat dienen mögen, Joseph zu bewegen, und sein Herz einzunehmen und zu bezwingen. „Werdet ihr diesen auch von mir nehmen,“ sagt unser Vater, „und ihm ein Unfall widerfähret, so werdet ihr meine grauen Haare mit Jammer hinunter in die Grube bringen.“ Hiernach aber geht die Erzählung, darinn Juda den Befehl Josephs und seine Bitte, dazu auch des Vaters Antwort erzählt hat. Nun folgt weiter im Text der Beschluß dieser Rede, welcher sehr beweglich und heftig ist.

\begin{quote}
\textbf{V. 30—34. Nun, so ich heim käme zu deinem Knechte, meinem Vater, und der Knabe wäre nicht mit uns, weil seine Seele an dieses Seele hänget, so wird's geschehen, wenn er sähet, daß der Knabe nicht da ist, daß er stirbt; so würden wir, deine Knechte, die grauen Haare deines Knechts, unseres Vaters, mit Herzeleid in die Grube bringen. Denn ich, dein Knecht, bin Bürge worden für den Knaben gegen meinem Vater, und sprach: Bringe ich ihn nicht wieder, so will ich mein Lebenlang die Schuld tragen. Darum laß deinen Knecht hie bleiben an des Knaben Statt, zum Knechte meines Herrn, und den Knaben mit seinen Brüdern hinauf ziehen. Denn wie soll ich hinauf ziehen zu meinem Vater, wenn der Knabe nicht mit mir ist? Ich würde den Jammer sehen müssen, der meinem Vater begegnen würde.}
\end{quote}

77. Dies hat er nun endlich mit schweren Seufzen und heißen Thränen ausgeredet: Wo ich den Knaben nicht wieder mitbringen werde, so wird gewißlich unser Vater vor herzlicher Traurigkeit sterben, und also werden wir seine grauen Haare hinunter in die Grube bringen. Was meinst du denn, daß du werdest für Nutzen davon haben können, wenn wir unsern alten Vater also tödten werden? Das wollest du doch ansehen: Was hilft es dich, wenn wir unsern alten Vater um das Leben bringen? Wir sollen ja in diesem Leben nichts lieber haben und daran uns mehr sollte gelegen sein, als die Wohlfahrt und Gesundheit unsers Vaters; derhalben will ich dir dieselbe auch lassen befohlen sein, und bitte, du wollest mit deiner Ungnade solchen erbärmlichen und unbilligen Unfall und Tod nicht verursachen. Denn Juda ist angst und bange, und bemüht sich sehr, dieweil er für den Knaben gegen den Vater Bürge geworden war. Dies alles weiß Joseph wohl, oder wird ja zum wenigsten ihren Worten Glauben geschenkt haben. Wir werden, sagt Juda, sonst nicht wieder dürfen vor unserm Vater kommen, es sei denn, daß wir den Knaben auch mitbringen. Ich will ewig dein Knecht sein für meinen Bruder; laß mich solches nur von dir erlangen, daß ich dein Knecht sein möge; denn wie könnte ich hinauf ziehen zu meinem Vater?

78. Dies alles wird er mit Weinen und Seufzen geredet haben, wird seine Hände ausgestreckt haben, wie die zu thun pflegen, so einen Andern zu Fuße fallen. Denn dies ist der letzte Theil dieser Rede, darin die allerhöchste Bewegung seines Herzens und auch seiner Brüder gewesen ist. Es sind sehr tapfere und wichtige Worte, die voll herzlicher Affecten und Bewegungen sind, welche niemands, wenn er auch schon bei sich selbst in seinem Herzen wohl zufrieden ist, ohne Thränen lesen kann. Cicero. sagt an einem Orte: Meine Thränen hindern mich, daß ich nicht reden kann etc. Aber das sind nicht wahrhaftige Thränen, so aus den allermächtigsten Ursachen und Affecten hergekommen oder entstanden sind, als, wenn die Worte und das Weinen untereinander gemengt und also dadurch die Rede verhindert wird.

79. Und wollte Gott, daß ich mit gleicher Inbrunst Gott auch also anrufen könnte, so oft, als ich gern wollte. Denn wenn ich einmal auf die Weise gebetet habe, hat mich gedäucht, daß mir sichtlich oder begreiflich diese Antwort gegeben würde: Es geschehe, was du bittest; es soll da sein. Und kein Mensch würde ein solch Gebet verachten können, welches also von Herzensgrunde heraus geht; vielweniger wird es bei Gott können vergeblich oder umsonst sein. Denn solches bezeugt die heilige Schrift, da sie sagt Sir. 35, 21.: „Das Gebet der Elenden dringet durch die Wolken“; item, daß der Engel sagt zu Daniel Cap. 9, 23.: „Da du anfingst zu beten, ging dieser Befehl aus.“ Flugs muß Gabriel da sein, als Daniel betet. Lauf hin, sagt unser Herr Gott, und antworte etc. Ach! daß wir also beten und glauben könnten, so wollten wir alles erlangen, was uns an Leib und Seele vonnöthen sein möchte. Oder, wenn uns schon Gott etwas versagen würde, so würde er uns doch ein Anderes, das viel besser und größer wäre, geben. Ja, es würde sich mit der Erhörung nicht verziehen, wenn ein solch heftig Gebet vor Gottes Ohren kommen würde.

80. Wie das Exempel Daniels zeuget, der da begehrete, daß Jerusalem und der Tempel daselbst wiederum möchte aufgerichtet werden. Das war sein Vorsatz, darum war es ihm zu thun: er sah die erbärmliche Verwüstung, und wurde darüber bewegt, ließ sich solches zu Herzen gehen. Darum hat er auch ohne Zweifel mit schwerem Seufzen und bittern Thränen gebetet. Flugs kommt der Engel und sagt: Du hast nicht allein das erlangt; denn das ist ein Geringes vor GOttes Augen: ich will dir noch mehr sagen: „Siebenzig Wochen sind bestimmt über dein Volk, und über deine heilige Stadt; so wird dem Uebertreten gewehret, und die Sünde zugesiegelt, und die Missethat versöhnet, und die ewige Gerechtigkeit gebracht“ etc., Dan. 9, 24. Gibt ihm die schöne herrliche Weissagung; ja, die schöne helle Offenbarung thut der Engel noch dazu und beschreibt damit die gewisse Zeit der Zukunft Christi. Das ists nun, das ich gesagt habe, daß ein solch heftig Gebet, wie dies ist, sobald es nur angefangen ist, alles erlangt, was es haben will, und noch viel Anderes mehr und Größeres dazu.

81. Derhalben soll uns dies ein Exempel und Anzeichen sein, daß uns auch begegnen und reizen sollte zu beten mit vollkommenem Vertrauen und starkem Muth, ohne allen Zweifel. Es steht zwar nicht allezeit, wie ich gesagt habe, in unserer Macht, mit solchem großem Geist zu beten; wir sollen aber doch Fleiß anwenden, vornehmlich in der Zeit der Noth, daß unser Gebet ganz brünstig sein möge. Die Herzen empfinden nicht allezeit ein solch Vertrauen, dadurch sie also im Gebet entbrannt sein sollten; denn es liegt uns im Wege und hindert uns, daß unser Fleisch und Blut so gar nachlässig und faul ist; item, die Erbsünde verdunkelt unsere Sinne und Verstand, macht denselben schwach, kraftlos und träge zum rechten herzlichen Gebet; und unsere Geschäfte hindern uns auch, wir werden hin und her gezogen von mancherlei Sorgen.

82. Aber wenn die äußerste Noth daher bringt, da kommt der Heilige Geist auch, der wirft es unserm Herrn Gott vor die Füße. Da tritt auch zugleich mit herzu der Geist des Gebets und schreit in uns: Herr Gott, du hast uns heißen bitten, beten, anklopfen und suchen etc. So werdet ihr es wahrlich erfahren, daß solch Gebet nicht vergeblich oder umsonst sei, sondern groß wunderbärlich Ding ausrichten kann, welches die Welt und wir auch nicht sehen können.

83. Auf solche Weise weinen nun und beten diese Brüder Josephs vor ihm mit größter Angst, Seufzen und Schreien ihres Herzens. Da stehen, sage ich, die lieben Brüder in den Seufzen ihres Herzens, winden die Hände, heulen und weinen, zerreißen die Kleider etc. Und sie weinen nicht allein, sondern Joseph vergießet auch viel Thränen, wiewohl er es verhehlet und den tiefen Schmerz seines Herzens noch immer in sich frißt. Er kann sich aber nun nicht länger enthalten, nachdem er sieht, daß an Juda und seinen andern Brüdern ein rechtschaffen Herz und kindhaftiges Gebet, dazu eine ungefärbte Liebe gegen Benjamin, und daß dabei auch nur rechte Buße ist. Denn dies ist nicht allein ein Exempel des Gebets, sondern auch einer wahrhaftigen und sehr großen Buße; wie ich im Anfange erinnert.

\quellseite{427}

84. Und alhier hätten wir nun Ursache, von rechter wahrhaftiger und falscher Buße zu reden; aber wir wollen dasselbe sparen bis an einen andern Ort, und ist auch albereit diese Lehre in andern Büchern reichlich gehandelt worden. Wir wollen jetzt in der Historie fortfahren und weiter betrachten den großen Schmerz und Seufzen dieser Brüder. Denn also werden sie gedacht haben: Ach, was will doch zuletzt hieraus werden! Wenn schon Juda hier bleibt und Benjamin mit hinauf zum Vater zieht, wird doch der Vater noch nicht können zufrieden sein. Und wir sind nun damit noch nicht frei von aller Sorge und Bekümmerniß; sondern die vornehmste Ursache ist noch dahinten, darum wir billig betrübt und traurig sein mögen, welche auch den Vater und uns selbst mit tödten wird. Denn wie sollte Juda hier bleiben können, welcher ja der vornehmste Haushalter in unsers Vaters Hause ist? Simeon, Ruben und Levi gelten nichts; Juda aber ist das Haupt in unserm ganzen Geschlechte. Wo sie es hinbringen, da will es nicht fort: es ist eitel Jammer und Herzeleid.

85. Es ist, werden sie gesagt haben, alles voll großen Jammers und Wehklagens mit uns; denn wenn wir schon vom Fürsten erlangen, daß wir Benjamin hinauf ziehen mögen und Benjamin wieder zu unserm Vater bringen, auf daß er nicht vor großem Herzeleid dahin sterbe: wie aber, wenn er da nun erst vor Leid sterben wird, wo er den Haushalter, und der das ganze Gesinde regiert hat, verlieren soll? Da würde der Vater zu uns sagen: Ihr wollt mich des Sohnes auch berauben, und verursachet mich also, daß ich ja für und für den schwersten Jammer und Herzeleid haben muß. Ihr bringt Benjamin wieder und läßt Juda aus: ihr bringt mich in allen Jammer. Also, dahin sie sich nur wenden, kommt ihnen nichts denn eitel Jammer und Noth vor beide innerlich in ihren Herzen und auch äußerlich vor ihren Augen; so gar auch, daß sie nun von vielem Weinen und Wehklagen matt und müde geworden sind, und nun nichts mehr bitten dürfen. Ich wollte Gott, daß ihnen noch möchte vergönnt und gestattet werden, daß Juda möchte Knecht sein für Benjamin! Sie wissen aber noch nicht, ob sie dasselbe vom Herrn erlangen können und ob er auch Benjamin entlassen wolle? Derhalben wird hier beschrieben der allerhöchste Schmerz und herzliche Reue über die begangene Sünde, damit auch Joseph nun bald überwunden wird. Er ist nun gefangen.

\chapter{Das fünfundvierzigste Capitel}

\section[Erster Theil]{Erster Theil. Joseph gibt sich seinen Brüdern zu erkennen}

\begin{quote}
\textbf{V. 1. 2. Da konnte sich Joseph nicht länger enthalten vor allen, die um ihn her stunden, und er rief: Lasset jedermann von mir hinaus gehen. Und stund kein Mensch bei ihm, da sich Joseph mit seinen Brüdern bekannte. Und er weinete laut, daß es die Egypter und das Gesinde Pharao höreten.}
\end{quote}

1. Ich zweifle nicht, es wird Joseph bei der ganzen Erzählung und Bitte Judas gezittert haben und ist ihm davon sein Herz weich geworden. Denn ein jegliches Wort, so Juda geredet, hat ihm sein Herz gerühret wegen der großen Treue und Liebe, so er zu seinen Brüdern getragen. Es sind eitel Herzensstöße gewesen. Denn er ist ja nicht grausam oder tyrannisch, sondern ganz fromm und freundlich gewesen, voll heiligen Geistes und wunderbarer herzlicher Neigung zu seinen Brüdern.

2. Und Juda hat beinahe alle die Stücke erzählt, die dazu gedienet, daß Josephs Affecten und natürliche Neigungen dadurch möchten erweckt werden, so daß ihm schwerer geworden ist, das Ende dieser Rede zu erwarten, als denn Juda selbst und seinen Brüdern, dieselbe angehört worden. Denn ein jegliches Wort, so Juda geredet, hat ihm sein Herz gerühret wegen der großen Treue und Liebe, so er zu seinen Brüdern getragen. Es sind eitel Herzensstöße gewesen. Denn er ist ja nicht grausam oder tyrannisch, sondern ganz fromm und freundlich gewesen, voll heiligen Geistes und wunderbarer herzlicher Neigung zu seinen Brüdern. gefangene Rede zu beschließen. Darum sagt er zu den Egyptern, so umher gestanden sind: „Lasset jedermann von mir hinausgehen.“ Er kann nicht länger halten. Sein Herz ist ihm gar weich geworden von dem lieblichen Feuer und der angenehmen Feuchtigkeit, so durch die Thränen und Seufzer übergeflossen war. Er heißt sie alle weggehen, und will nicht leiden, daß jemand anders bei ihm sein soll denn nur diese fremden Gäste, vielleicht darum, daß er nicht gern hat haben wollen, daß die Egypter sein Weinen und Seufzen sehen sollten. Denn der Ausbruch seiner Affecten und freundlichen Neigung zu seinen Brüdern und Vater war so stark, daß er nicht hat leiden können, daß fremde Leute haben sollen dabei sein und daßelbe ansehen. Er hebt derhalben nun mit einem kurzen Worte an:

\begin{quote}
\textbf{V. 3. Und sprach zu seinen Brüdern: Ich bin Joseph.}
\end{quote}

3. Da reißt er heraus. Wäre es doch nicht Wunder gewesen, wenn sie alle miteinander, da sie das Wort gehört, beide vor Schrecken und vor Freude gestorben wären. Er kann aber nicht mehr herausreden, obschon die Augen nicht dabei sind, sondern schüttet sich ganz heraus, kann nichts mehr thun, denn nur weinen, seufzen und heulen, und daßelbe mit lauter Stimme, daß es auch die haben hören können, so hinaus gegangen waren, und das ganze Gesinde Pharaos.

4. Derhalben hat er sich bisher mit großem starken Muthe enthalten und gezwungen, während Juda redete und die Egypter, so umher gestanden, hinausgingen. Da geht es nun heraus, wie ein Wolkenbruch, und verwundert sich das ganze Haus desselben Seelens und Schreiens, und fragen: Was doch dem Herrn widerfahren sein möge, darum daß er also weinte? Was doch das Weinen bedeuten möge? Ob ihm vielleicht die fremden Gäste Gewalt gethan hätten?

5. Das heißt recht, die Herzen rühren und bewegen und beide mit Worten und zugleich auch mit bittern Thränen beweisen die großen Bewegungen des Herzens. Ich kann dies Ding mit Worten nicht ausreden oder erklären, ja, kann es auch mit Gedanken nicht erreichen oder genugsam verstehen.

6. Das ist aber ein sehr freundlich Wort, und das zu dem Kampfe, den sie ausgestanden, sehr wohl dient, daß er saget: „Ich bin Joseph.“ Das dient sehr wohl zur Sache. Und er hat auch daselbst der egyptischen Sprache nicht gebraucht, wie er zuvor durch einen Dolmetscher mit ihnen geredet; sondern braucht jetzt die Muttersprache, da doch seine Brüder nicht geglaubt hatten, daß ihm dieselbe sollte bekannt sein, weil er einen Dolmetscher gebraucht hatte. Jetzt aber, da die Mittler und Unterhändler nicht dabei gewesen, redet er mit ihnen auf Hebräisch, sagt nicht: Ich bin Abrech, Zophnath, Paneach, wie ihn ohne Zweifel seine Brüder zuvor werden genannt haben wegen der größten Gewalt und Herrlichkeit, so er in Egypten gehabt; welche Ehre denn auch der Name, so ihm der König gegeben, bedeutet hat; denn „er nannte ihn den heimlichen Rath“, wie droben im 41. Cap. V. 45. steht: sondern er sagt nun zu ihnen: „Ich bin Joseph.“

7. Dies ist nun abermals ein harter Kampf und der Höhepunkt in dieser Historie, da sich Joseph seinen Brüdern so plötzlich zu erkennen gibt, welchen sie zuvor für einen greulichen und erschrecklichen Tyrannen gehalten haben, darum daß er sich auch so gar fremd und so sehr erschrecklich gegen sie erzeigte. Und wiewohl er mit sehr vielen Anzeichen zu der Hoffnung Anlaß gegeben, daß unter solcher feindlichen Gestalt noch Freundlichkeit und Gnade verborgen wäre, haben sie doch solches nicht verstehen können. Hier hat er aber nun keinen langen Umschweif noch auch einen Dolmetscher gebraucht, sondern fährt stracks heraus mit diesen Worten: „Ich bin Joseph.“

8. Darum ist dies ein sehr schön Exempel, nämlich, wie sich Gott pflege gegen uns zu halten. Denn wenn er die Gottlosen straft, und sich stellt, als sei er unser Gott und Vater nicht; oder aber sich dermaßen erzeigt, daß er vielmehr ein Tyrann und strenger Richter sei, der uns nur wohl peinigen und gar verderben wolle: so sagt er doch endlich zu seiner Zeit und wenn seine gelegene Stunde kommt: „Ich bin der Herr, dein Gott.“ Bisher bin ich dich nicht anders mit dir umgegangen, denn als hätte ich dich gar vergessen und in die Hölle stoßen wollen; aber das Spiel pflege ich also mit meinen Heiligen zu treiben. Denn wo ich dir nicht wäre von Herzen günstig gewesen, so wollte ich dermaßen mit dir nicht gespielt haben.

\quellseite{428}

9. Dies wird uns also vorgemalt in der Regierung der Heiligen zu unserm Troste, auf daß wir lernen sollen die Hand Gottes dulden, welche uns unterweist und prüft, daß wir uns selbst erkennen und demüthigen lernen, und daß also in uns möge getödtet werden das erschreckliche Unglück, welches die Erbsünde genannt wird. Denn es ist nicht darauf abgesehen, daß wir sollten verdammt oder verworfen werden; wiewohl unser Kreuz und Strafe, so wir tragen, dem Verderben und Tode sehr ähnlich ist: sondern es muß die Sünde, so uns von Natur anhängt, ausgestoßen werden, auf daß wir lernen, was das gesagt sei, daß der Herr sagt 1 Sam. 2, 6. 7.: „Ich tödte, und mache lebendig, führe in die Hölle, und wieder heraus; ich mache arm, und mache reich.“

10. Das sind Gottes Werke, davon Jesaia, Paulus und Christus selbst im Evangelium so herrlich disputiren, Matth. 10, 38.: „Wer nicht sein Kreuz auf sich nimmt, und folget mir nach, der ist mein nicht werth“; item Joh. 12, 25.: „Wer sein Leben lieb hat, der wirds verlieren“ etc. Aber da fragt die Vernunft: Warum denn Gott die unschuldigen Christen also plage? Antwort: Das thut er nicht vergeblich oder umsonst, sondern daß er damit der Sünden weniger machen und dieselben abthun will. Darnach, dieweil wir noch immer mancherlei Fehler und Schuld auf uns liegen haben, können wir uns mit nichten der Unschuld rühmen. Denn siehe nur diejenigen an, welche er so sicher und ungestraft nach ihrer eigenen Lust und Begierde hingehen läßt, läßt ihnen den Zaum und daß sie gar keine Anfechtung haben. Lieber, was thun doch dieselbigen? Ei, sie fallen dahin von einer Sünde in die andere, in Unzucht, Todtschlag, Ehebruch, Haß und andere erschreckliche scheußliche Sünde und Schande. Wie daselbige heutiges Tages am Stuhle zu Rom und allen Curtisanen? päpstlichen Haufen zu sehen ist. Denn was sind die Stifte der Dompfaffen und Cardinalen Anderes als eitel Hurenhäuser? Und Rom ist ein solcher Ort, da alle Schälke und Buben, dazu alle Laster überhand genommen haben. „Denn sie sind nicht in Unglück wie andere Leute, und werden nicht, wie andere Menschen, geplagt“, wie der 73. Psalm V. 5. sagt. Es wird die Sünde an ihnen nicht also aufgedeckt und sie haben dabei keine Züchtigung; wo aber die Tödung und Reue aufhört, da ist kein Rath oder Hülfe zur Seligkeit mehr vorhanden.

11. Darum müssen wir Christen lernen, daß die heilige Schrift allerlei diese Züchtigung lehrt, dadurch wir gedemüthigt werden. Wo wir aber gerecht, und unter diejenigen gezählt sind, wie wir denn ja gewiß sind, die da begehren das Sündenreich zu erlangen: so sollen wir wissen, daß wir in den Tod Christi getauft sind, Röm. 6, 3. Desselbigen müssen wir, als unserm Hauptmanne, nachfolgen, daß der Leib der Sünden aufhöre. Und unser Leben ist vor Gottes Augen nichts Anderes, denn daß wir nur immer zerstört, getödtet und in die Hölle geführt werden müssen. Nicht darum, daß wir dadurch verderbet werden sollen, sondern uns dadurch zu helfen, zu erretten und zu läutern. Aber solche Tödung muß nicht allein im Worte, sondern auch im Werke bestehen; wie St. Paulus 1 Cor. 4, 20. sagt: „Das Reich Gottes stehet nicht in Worten.“ Es hängt uns die Sünde noch immer an, sintemal wir darin empfangen und geboren sind; unser Fleisch ist verderbet, welches der Teufel zu allerlei Sünden und Laster reizt: was sollte denn nun Gott sonst mit uns thun, denn daß er uns immer läutere? wiewohl solches ohne Schmerz und Pein nicht geschehen kann. Es muß ja wehe thun und muß gefühlt werden.

12. Auf solche Weise ist Joseph auch wunderbarlich versucht und wohl geübt worden, bis er gedemüthigt, gezüchtigt und endlich auch wiederum erhöht würde. Denn also ist er zu solcher großen Erkenntniß und Verstand der Gottseeligkeit, und zu solcher großen Weisheit gekommen, daß er auch Andern rathen und vorstehen, ja, auch die ganze Welt regieren könnte. Das ganze Königreich Egypten war nicht groß genug für seine Weisheit; wiewohl er darin sehr vielen Leuten nütze gewesen, und nicht allein das gemeine Volk, sondern auch den König selbst und die Priester unterwiesen hat in der wahren Seligkeit. Woher hat er aber das gehabt, daß er solches vermocht hat zu thun? Antwort: Daher, daß er zuvor getödtet und in die Hölle geführt worden, war verachtet und in verzweifelter Lage; und hat doch mit großem starkem Muthe noch immer in der Auferstehung rechten Glauben und Hoffnung behalten.

13. Eben auf solche Weise versucht und demüthigt er seine Brüder auch, auf daß er sie erhöhen möge. Derselben ist dies ein rechtschaffener und trefflich Exempel der göttlichen Regierung, „welcher auch seines eigenen Sohnes nicht hat verschonet, sondern hat ihn für uns alle dahin gegeben“, Röm. 8, 32., daß er für unsere Sünde genugthun sollte. Wir aber werden nicht darum gedemüthigt, daß wir damit für die Sünde genugthun sollen, sondern daß wir nur geläutert und gefeget werden. Er hat für sich selbst keiner Reinigung oder Läuterung bedurft, dieweil er ganz unschuldig gewesen; sondern mit seinem Leiden hat er verdient, daß unsere Demüthigung, Tödung, Verwerfung und Verdammung Gott angenehm ist.

14. Das soll man nun lernen, wie es in der heiligen Schrift gelehret wird. Denn die Papisten und Türken wissen von diesen Dingen gar nichts, meinen, sie seien selig, vollkommen und heilig und daß sie mit ihrem Gehorsam und vermeinten Gottesdienste Gott wohlgefallen ohne solche Läuterung. Darum fragen sie nach Gott nichts, fürchten ihn auch nicht. Und ich rede jetzt nicht vom Papste Julius oder Clemens, als scheußlichen Ungeheuern unter den Päpsten; sondern lässet uns die Exempel der guten frommen Heiligen ansehen, als Augustini, Bernhards und viel Andere mehr gewesen sind, den Glauben an ihnen ausgenommen. Dieselben denken, sie seien fromm und gesund, lehren ja wohl Geduld, aber es geht sehr kalt zu. Und wenn du fragst, warum sie fasten, warum sie ihren Leib kasteien, warum sie sich geißeln lassen, bringen sie die Ursache hervor, und sagen: Auf daß sie damit verdienen mögen, daß sie über andere gemeine Christen im Himmelreiche erhöhet werden mögen.

15. Das ist aber nicht die rechte endliche Ursache des Kreuzes, so den Christen auferlegt wird; sondern vielmehr, daß sie dadurch am Fleische getödtet werden und daß die Sünde aufhöre, und der Erbsünde, so uns von Natur anhängt, auch etwas gesteuert werde. Und je mehr du gefegt oder geläutert wirst, je mehr du im zukünftigen Leben gesegnet werden wirst. Denn es wird ohne Zweifel Ehre und Herrlichkeit nach dem Jammer und Klage folgen, so wir in dieser Welt tragen müssen. Aber aller dieser Trübsale endliche Ursache ist die Läuterung, welche ganz nöthig und nützlich ist, auf daß wir nicht schandbar, träge und faul werden von der Schlafsucht unsers Fleisches. Denn wo wir Friede und Ruhe haben, so beten wir nicht, wir betrachten das Wort nicht, sondern gehen kalt mit der Schrift um, und mit allem, was Gott betrifft; oder aber gerathen doch endlich gar in schändliche und schädliche Sicherheit.

16. Darum müssen wir geplagt und gedemüthigt werden. Geschieht es nicht durch Blutvergießen und Gefängniß, wie der heiligen Märtyrer Leiden gewesen, so muß es doch durch geistliche Anfechtung, durch Traurigkeit, Schmerz und Angst, so wir im Herzen haben, geschehen. Denn sonst würden wir in Sünden verderben, dieweil unser Fleisch verderbt, vergiftet und gar ausfätzig ist, und eines Arztes bedarf, der dem faulen Fleische vorkomme durch Kreuz, Marter, Traurigkeit, Schande, Schmach; denn das ist die rechte Arznei, damit Gott die Sünde auflegt.

17. Solcher Exempel ist die heilige Schrift voll, nämlich, daß die Heiligen erhalten müssen. Denn außerhalb der Anfechtung vergehen sie der geistlichen Uebungen, beten, glauben und loben Gott desto unfleißiger. Wiederum aber in der Noth und Trübsal, da schreien, seufzen und klagen, und werden also im Worte geübt und unterwiesen; wie Jesaia Cap. 28, 19. sagt: „Alleine die Anfechtung lehret aufs Wort merken.“

18. Zuletzt aber, wenn die Anfechtung ein Ende hat, da läßt sich der allgnädigste Vater und Heiland unserer Seelen sehen, und sagt: „Ich bin der Herr, dein Gott; ich habe dich gestäupt und gestraft.“ Wie Joseph hier auch mit lauter Stimme sagt: „Ich bin Joseph“; ich habe euch sehr umgetrieben und auf mancherlei Weise versucht; nicht aus feindseligem Gemüthe oder daß ich Lust hätte euch zu verderben, sondern vielmehr, daß ich in euch habe die Buße und die Läuterung der Sünden erwecken wollen, auf daß ihr gerecht und fromm werden möchtet. Darum fürchtet euch nicht; denn ich bin der Mann nicht, für den ihr mich haltet und wie ihr euch eingebildet habt, nämlich, ein erschrecklicher und greulicher egyptischer Tyrann. Nein, so fehlet ihr weit.

19. Eben auf dieselbe Weise verhält es sich auch mit den Gedanken, Murren und Argwohn der Gottseligen, so sie etwa bei sich selbst von GOtt schöpfen. Denn das ist alles eitel falsch und erdichtet Ding. Denn wir können nicht verstehen, was eigentlich GOtt mit seinem Rath meint. „Was ich thue, das weißest du jetzt nicht“, sagt Christus zu Petro Joh. 13, 7., „du wirst es aber hernach erfahren.“ „Ich werde dir jetzt die Füße waschen; was aber dasselbe Waschen bedeutet oder was ich damit meine, verstehst du jetzt nicht.“

\quellseite{429}

20. Also verstehen wir auch GOttes Rath und Willen nicht zur Zeit des Leidens. Darum zittern diejenigen so erschrecklich, welche betrübt und gestraft werden, und kämpfen mit dem Unglauben, Zorn und Ungeduld wider GOtt, und dazu auch mit Verzweiflung. Wenn sie aber mit GOttes Hülfe am Glauben und Hoffnung festhalten, auf daß sie nicht gar verzagen, so hat es keine Noth. Denn GOtt rechnet ihnen das Klagen und Murren für keine Sünde. Darnach aber kommt er, und sagt: Ich bin der Herr, dein Heiland.

21. Das wußte ich aber nicht, sagt die Vernunft, und ließ mich nicht anders bedünken, denn daß der Teufel also wider mich tobt und gewüthet hätte. Denn unsere Natur kann nicht mit Geduld ertragen, daß sie verderbt und zerstört werden soll; und verstehet auch nicht, daß die Versuchung nur eine Läuterung ist: darum erschrickt sie davor und zürnt oder tobt darüber. Derhalben gehört von Noth wegen dazu der Glaube an das Wort, auf daß wir die Hand des HErrn dulden und ertragen können, und wissen, woran wir uns halten und worin wir beruhen sollen.

22. Die gestraft und geängstigt werden, fühlen nichts Anderes und klagen auch, sie müssen verderben und gar zu Boden gehen, und die Natur kann anders nicht denken, wenn sie getödtet wird; sie kann nicht anders sagen denn: Ich bin verloren. Aber, o mein seliges tröstliches Wort und Stimme der Christen, wenn sie können am Worte hangen und sagen: Ich bin getauft; ich glaube an GOtt den Vater; ich glaube an Jesum Christum. Wenn das nur im Herzen fest und gewiß bleibt, alsdann wird Alles gut werden.

23. Das Fleisch wird zwar geängstiget und geplagt; es thut wehe; die Vernunft verzagt; unser Wille murrt; und endlich, alle Sinne sind gar kleinlaut und bestürzt. Wie David von sich selbst sagt im 31. Psalm V. 23.: „Ich sprach in meinem zagen: Ich bin von deinen Augen verstoßen.“ Solche Worte aber und Gedanken sind zu strafen, und man muß der verderbten Natur und Vernunft widerstehen, welche nichts Anderes verstehen oder reden kann. Was soll ich aber dann thun, möchtet du sagen, wenn ich jetzt schon in die Thür der Hölle und in Verzweiflung gesetzt bin? Soll ich sagen, wie im Psalm steht: Ich bin verloren, ich bin gar verstoßen? etc. Mit nichten, sondern du sollst bei dir also schließen und sagen: Ich bin getauft, bin berufen durch das Wort; ich glaube an den Sohn GOttes, der für mich gelitten. „Mein Vater und meine Mutter verlassen mich; aber der HErr nimmt mich auf“, Ps. 27, 10.

24. Ja, sagst du, ich sehe aber den gnädigen und geneigten HErrn nicht. Das ist recht also. Höre aber, daß er ein solcher Heiland und Helfer ist: Wiewohl er die Barmherzigkeit und Hülfe verhehlt und verbirgt, und sich also erzeigt, als müsse er von dir gar nichts; so sitzt er doch zur Rechten GOttes des Vaters, regiert dich, sorgt und ist um dich bekümmert. Christus klagt selbst auch im 38. Psalm V. 12.: „Meine Lieben und Freunde stehen gegen mir, und scheuen meine Plage, und meine Nächsten treten ferne.“

25. Das ist den Gottseligen sehr beschwerlich, wenn dies alles nicht zum Schmerz, sondern im Ernst geschieht. Und dieweil es Wirklichkeit und nicht nur eitler Schein ist, darum thut es dem Fleische sehr wehe. Als, wenn ein Knabe mit der Ruthe gezüchtigt wird, darum daß er nicht fleißig studirt hat, so fühlt er wahrlich den Schmerz wohl, schreit und heult; und es kann gleichwohl nicht anders zugehen, wo man Zucht und Gehorsam halten muß. Viel weniger wird bei den Erwachsenen, wo sich die Erkenntniß zeigt, die Sünde ohne Schmerz ausgeleget, daß sie nicht schreien sollten, sie seien von GOtt verstoßen und gar verloren. Aber man soll mit nichten also schließen. Wie das Kind um der Züchtigung des Vaters oder Zuchtmeisters willen auch nicht denkt, daß ihm der Vater feind sei. Und was sind wir vor GOtt Anderes denn Kinder?

26. Darum sagt der 118. Psalm V. 18. ganz herrlich: „Der HErr züchtiget mich wohl, aber er giebt mich dem Tode nicht.“ Und Salomo Spr. 22, 15.: „Thorheit steckt dem Knaben im Herzen; aber die Ruthe der Zucht wird sie ferne von ihm treiben.“ Eben dasselbe soll man auch halten von göttlicher Züchtigung, vornehmlich weil er uns das Wort gegeben hat, welchem wir anhangen sollen, nämlich, daß es kein Schmerz, und Traurigkeit oder Buße zur Furcht und Verzweiflung sei; sondern wie hart er auch zürnt, daß wir doch glauben, er sei unser Heiland und lieber Vater. Wie Hiob sagt: „Wenn er mich tödten würde, so will ich doch auf ihn hoffen“, Hiob 13, 15.

27. Denn das ist der Unterschied zwischen den Kindern GOttes und den Ungläubigen, welche sich wider die Strafe auflehnen und die Züchtigung des HErrn nicht tragen wollen. Den Gottseligen aber thut die Züchtigung noth um des verderbten Fleisches willen, welches voll Sünde und Sicherheit, Verachtung GOttes, Zorn und Unzucht ist. Daher es auch kommt, daß sie nicht beten, nicht glauben, nicht hoffen und GOtt nicht achten.

28. Diese Gebrechen und schweren Seuchen und Krankheiten hangen den Heiligen und Gläubigen noch immer an. Darum, auf daß GOtt dieselben an ihnen heilen möge, gebraucht er dazu solcher Strafe, und daß er sie tödtet und in die Hölle führt. Welche Arznei und Werke GOttes in der heiligen Schrift und Exempeln der Gottseligen hin und wieder vorgebildet sind; sie sind aber von der Welt und Vernunft nie verstanden worden; denn sie haben das Buch nicht, daraus sie diese Kunst lernen könnten. Und es ist auch kein ander Buch, welches dies lehrt, denn nur das Buch, darin der Heilige Geist mit uns auf diese Weise redet: „Ich bin ein GOtt, der da tödtet, und lebendig macht, der in die Hölle führt, und wieder heraus, der da arm macht, und reich macht“, 1 Sam. 2, 6. 7. Und solches thue ich nicht besonders oder Eines ohne das Andere, als, daß ich allein sagen sollte: Ich tödte; nein, das wäre der Teufel: sondern ich bin ein solcher GOtt, der da tödtet, und wiederum auch lebendig macht. Ich führe wohl in die Hölle, aber doch also, daß ich wieder herausführe.

29. Dies soll man oftmals wiederholen und fleißig treiben. Denn es ist der Christen Weisheit und eigene Lehre, welche einen solchen starken Muth haben sollen, daß ihnen nichts Böses widerfahren möge, daß sie sich desselben nicht rühmen sollten, und darüber triumphiren und fröhlich sein, darum daß sie das Wort, dazu die Absolution, die Taufe, Vergebung der Sünden und Verheißung des ewigen Lebens haben, und was für Unfall sie leiden sollen, es sei im Tode oder in der Hölle, keinerlei Unglück auszunehmen, sie werden gleich innerlich oder äußerlich bedrängt oder geängstigt; daß sie doch sagen: Ich bin fröhlich, triumphire und stolziere; denn ich bin ein Herr des Todes, der Sünden, der Höllen, der Welt, dazu über alles Unglück; daß bei ihm kein Unterschied zwischen Leben und Tod, Reichthum und Armuth, Schande und Ehre sei; und daß also ein Christ ein Herr werde, mächtig im Streit, welcher über den Schrecken des Todes, der Hölle und alle Widersacher gesetzt ist. Und dasselbe zwar nicht von uns selbst; denn die Natur kann solches nicht ausrichten; sondern weil unser Heiland, unser HErr und Siegesmann, der Sohn GOttes, zu der Rechten des Vaters sitzt, der uns errettet, schützt und regieret. Derhalben, so mag uns nun der Tod erschrecken, der Satan mag uns in die Hölle hinreißen: was liegt daran? Mein HErr und Erlöser lebt noch.

30. Diese Kunst übertrifft alle Vernunft und menschliche Weisheit; und die Heiligen, als, David und Paulus, haben dieselbige auch nicht vollkommenlich erreichen oder verstehen können. Paulus, der doch ein hoher Apostel gewesen, bekennt von sich selbst 2 Cor. 12, 7—9.: „Mir ist gegeben ein Pfahl ins Fleisch, nämlich des Satans Engel, der mich mit Fäusten schlage, dafür ich dreimal dem HErrn geflehet habe, daß er von mir wiche.“ Er wäre gerne sein Los gewesen. „Und er hat zu mir gesagt“, nämlich, der HErr zur Rechten GOttes: „Laß dir an meiner Gnade genügen.“ Wofür hältst du das, lieber Paulus? daß ich dich liebe? Meine Kraft kann nicht mächtig sein denn nur in eurer Schwachheit. Du mußt schwach sein, du mußt leiden, seufzen, jämmerlich und schwach sein, dir selbst zu gut, auf daß du mit Leiden und Streiten endlich gesiegest und ein großer Apostel werdest. Wo du nicht schwach sein wirst, so hat meine Kraft an dir nichts zu thun. Wenn ich dein Christus sein soll und du wiederum mein Apostel, so wirst du deine Schwachheit mit meiner Kraft, deine Thorheit mit meiner Weisheit, mein Leben mit deinem Tode zusammen reinen müssen.

\quellseite{430}

31. Dies ist der Christen Lehre, welche weder die Philosophen noch Juristen lehren; sie verstehen auch der Ursachen keine, weder wo sie herkommt, noch wohin sie gerichtet werden soll, noch auch wodurch man dieselbe zuwege bringen möge. Die Philosophie lehrt ja wohl etwas von Geduld; sie fordert aber auch zugleich dabei, daß man unschuldig sein soll. Darum soll nun das fleißig merken, daß wir wohl als die Unschuldigen leiden, aber die gleichwohl dieser Ruthe bedürfen. Wie der HErr im Jeremia wider Edom sagt, Jer. 49, 12.: „Siehe, so ich es nicht verschuldet hatte, den Kelch zu trinken“, wie Daniel und etliche Andere, „müssen trinken“, und du solltest ungestraft bleiben? Du sollst nicht ungestraft bleiben, sondern du mußt auch trinken.“ Item, also sagt er auch wider Jerusalem, Jer. 2, 35.: „Siehe, ich will mit dir rechten, daß du sprichst: Ich habe nicht gesündiget.“ Als wollte er sagen: Darum will ich dich strafen, daß du deiner Sünden nicht vergessen sollst, darin du empfangen und geboren bist. Denn sonst würdest du sagen: Warum strafst du mich oder warum muß ich also gepeiniget werden, warum legt du diese Strafe nicht vielmehr Moab und Edom auf? Ich bin ja dein Volk, ich opfere dir, lobe und rühme deinen Namen. Das thue ich, spricht der HErr, nur darum, auf daß du dich deiner Unschuld vor mir nicht rühmen sollst. Ich will dich aber nicht in meinem Grimm züchtigen, sondern mit Maßen im Gericht, wie ein Vater sein Kind liebt und züchtigt. Du sollst darum nicht verloren sein.

32. Ja, möchtest du sagen, was habe ich denn gesündiget? was habe ich verschuldet? Antwort: Frage den Apfel im Paradiese. Moab und Babylon haben zwar keine Ursache, dich zu strafen und zu verfolgen; du bist aber vor mir nicht unschuldig um der Sünde willen deiner ersten Eltern im Paradiese. Darum sollst du mir die Ehre geben und deine Unreinigkeit mit Demuth erkennen, und sagen: HErr GOtt, du thust recht, daß du uns züchtigst; wie Daniel sagt Cap. 9, 7.: „Du, HErr, bist gerecht, wir aber müssen uns schämen.“

33. GOtt weiß und erkennt, was in unsern Herzen verborgen ist und wie verderbt unsere Natur ist; darum straft er uns gnädig und väterlich. Und zuletzt, wenn die Strafe aufhört, offenbart er sich uns und sagt: Ich bin dein GOtt, der ich dich liebe und freundlich umfasse; ich habe es herzlich und freundlich gemeint. Da werden denn die Seelen wiederum allgemach aufgerichtet und können sich alsdann mit der Gnade und Freundlichkeit GOttes trösten. Und das ist der Gottseeligen Trost in diesem Leben. O wie wird aber ihre Freude an jenem Tage so groß sein, wenn der Sohn GOttes erscheinen, und sagen wird: Siehe, ich bin dein Heiland und Erlöser, welchen du beschuldiget hast, als sollte ich nach meiner Kirche und Gemeinde nichts fragen, und keine Acht auf dieselbe haben, als ob ich kein GOtt wäre und auch für euch nicht sorgte. Siehe, hier bin ich. Darum aber habe ich dich gezüchtiget, auf daß deine Sünde geläutert und ausgefegt würde, und daß du mich für deinen GOtt und Heiland erkennen solltest.

34. Dies Spiel und wie wunderlich GOtt seine Heiligen regieret, hat der Heilige Geist in diesem Exempel Josephs abmalen wollen. Und lieber, gedeutet doch, welch ein lieblich angenehmes Wort in den Ohren der Gottseeligen das am jüngsten Tage sein wird: „Ich bin Joseph“ 2c., ja, was jetzt das für eine Freude ist, wenn GOtt auch ein klein Tröpflein Trostes durch sein Wort denen gibt, so betrübt und angefochten werden. Wie wird doch Einer so herzlich froh, wenn die Anfechtung vorüber ist! Derhalben wird das viel eine größere, ja, eine unaussprechliche Freude sein, da der GOtt alles Trostes sich offenbaren und sich über uns im ewigen Leben gar ausgütten wird.

35. Es haben die Söhne Jakobs eine kleine Linderung ihres Jammers und Schmerzes, so sie zuvor gehabt, gefühlt, nämlich, da Joseph sie an seinen Tisch gezogen, und gnädig mit ihnen umgegangen ist und sie köstlich tractirt, ihrem jüngsten Bruder fünfmal so viel gegeben als ihrer Einem, und endlich einem Jeden insonderheit sein Geld wiedergeben lassen.

36. Und eben auf solche Weise spielt GOtt mit uns auch, wenn er uns zeitlichen Trost sehen läßt und verleiht. Aber was will es werden an jenem Tage, da nicht ein klein Tröpflein, sondern der ganze Himmel wird aufgethan werden, und die ganze Gottheit sich wird offenbaren und sehen lassen? Da werden wir fröhlich sein und uns dieses Wortes freuen: „Ich bin Joseph“ 2c. Unser lieber HErr GOtt wolle uns stärken und Geduld verleihen, daß wir dieser großen Hoffnung, so uns vorgestellet ist, mit solchem Glauben und Beständigkeit warten mögen, wie wir billig thun sollten. Denn diese Freude wird hin und wieder in den Psalmen auch gerühmt, und werden auch etliche Vermahnungen dazu gethan, die Herzen der Gottseeligen damit zu stärken. Als, im 31. Psalm V. 20.: „Wie groß ist deine Güte, die du verborgen hast denen, die dich fürchten, und erzeigest denen, die vor den Leuten auf dich trauen“; und im 27. Psalm V. 14.: „Harre des HErrn, sey getrost und unverzagt, und harre des HErrn.“

37. Unterdeß aber bleibt in diesem Leben der Wechsel mit Trost und Trübsal; denn sonst müßten die Gottseeligen wegen steter Anfechtung gar verderben, wenn sie dabei nimmermehr ein wenig Trostes hätten. Wie Joseph jetzt seine Brüder zu sich geladen und sie freundlich empfangen hat; darnach hat er sie bald wiederum etwas erschreckt. Also führt GOtt uns auch wunderbahrlich, und erhält uns, daß wir in der Noth und Trübsal nicht verzagen, da wir also von unsern Sünden geläutert werden; denn solche Läuterung wird im zukünftigen Leben aufhören. Nun folgt im Text:

\begin{quote}
\textbf{Lebet mein Vater noch?}
\end{quote}

38. Er fragt nach dem Vater, nicht, als hätte er daran gezweifelt, daß er noch am Leben war, sondern daß er damit beweise, daß er sie von Herzen lieb habe und aus solcher Liebe jetzt auch mit ihnen rede. Es ist aber auch zugleich ein solch Wort, das eine Verwunderung in sich begreift. Als wollte er sagen: Ich habe schwerlich glauben können, daß mein Vater noch sollte am Leben sein; habe allezeit gedacht, er hätte keineswegs so lange leben können. Denn ich bedenke noch wohl, wie heftig lieb er meine Mutter gehabt, und da dieselbe in der Geburt jämmerlich dahin gestorben, habe ich gesehen, wie erbärmlich er über ihren Tod Leid getragen, und habe dazu angesehen, daß er sonst auch im Haus viel Jammer und Unfall gehabt, dadurch er heftig ist betrübt worden. Denn seine Natur ist ganz freundlich, und pflegt die Seinen herzlich lieb zu haben, nicht allein aus natürlicher Neigung, sondern auch aus Liebe des Heiligen Geistes. Derhalben ist mir oftmals eingefallen, da ich weggekommen war, er würde

39. Denn die Schwachheit der Natur hängt auch den Heiligen noch immer mit an, und die Traurigkeit und Bekümmerniß ist so gar stark und kräftig, daß sie viel Leute tödtet. Es sterben in diesem Leben ihrer Viele, auch etliche gottseelige Leute, vor Bekümmerniß ihres Herzens, welches an sich selbst der rechte Tod, und der viel schärfer ist, denn ein Schwerdt sein kann; wie Sirach zeuget im 30. Cap. V. 25., da er sagt: „Traurigkeit tödtet viel Leute, und dienet doch niemand zu“; und Paulus 2 Cor. 7, 10.: „Die Traurigkeit der Welt wirket den Tod.“ Und so je zarter die Natur an den Menschen ist, je eher sie dahinsterben.

40. Es sind aber droben die sehr schweren Trübsale und Bekümmernisse, so Jakob gehabt, beschrieben worden, als, daß ihm Dina und Bilha sind geschändet worden, daß ihm Rachel und Debora, die Säugamme, so kümmerlich dahin gestorben. Nun hat aber der treffliche Mann solcher Trübsale und vielen Unfalls halbsein sehr müssen bekümmert werden, wieviel er seiner natürlichen Affecten und Neigung halben muthwillig und mitleidig, und dazu auch wohl Heiligen Geistes gewesen ist. Denn solche Leute werden sehr bewegt zur Barmherzigkeit, und haben mit Andern, so betrübt und bekümmert sind, groß herzlich Mitleiden. Solche Affecten aber oder Bewegungen zerrütten die Natur auch in den Heiligen.

41. Darum verwundert sich Joseph dessen billig, daß sein Vater die zwei und zwanzig Jahre über, da er nicht war daheim gewesen, in solchem großen Unfall und so vielem Jammer hat können lebendig bleiben. Daran zweifle ich nicht, wird er gesagt haben, mein Vater wird in seinem Herzen so gar bestürzt und erschrocken gewesen sein, da er mich verloren hat, daß er sich auch nicht leichtlich hat trösten lassen. Darum habe ich gedacht, er müsse gewißlich gestorben sein, wo nicht Hungers halben, so müsse es doch aus Traurigkeit und Leid geschehen sein. Joseph wird gesagt haben: Ach GOtt! ich habe allezeit gedacht, mein Vater wird sich zu Tode gegrämt haben. Denn wiewohl die Gottseeligen an GOtt glauben, ist doch die Natur so gar schwach, daß sie die großen Schmerzen und Traurigkeit nicht
\quellseite{431}
ertragen können; und Christus sagt selbst auch Matth. 26, 38.: „Meine Seele ist betrübt bis an den Tod.“

\section[Zweiter Theil]{Zweiter Theil. Wie den Brüdern Josephs zu Muthe ist, da Joseph sich ihnen zu erkennen gibt; und wie Joseph mit ihnen umgeht}

\subsection{I}

\begin{quote}
\textbf{V. 3. Und seine Brüder konnten ihm nicht antworten, so erschraken sie vor seinem Angesicht.}
\end{quote}

42. Daß sich Joseph seinen Brüdern so plötzlich offenbart, daß ihrer keiner nimmermehr gemeint oder sich versehen hätte, das hat in ihrem Herzen eine solche große Freude erweckt, daß sie vor Furcht und Freude zugleich erschreckt worden und schier gestorben sind. Es ist die Seele etwas Vortrefflicheres als der Leib, und der Leib kann solche Bewegungen und herzliche Freude, wenn dieselbe im höchsten Grad ist, nicht ertragen. Was aber die Seele sei in der Speculation oder, wie die Philosophen reden, metaphysisch, übernatürlich, können wir nicht wissen, sondern aus den Affecten oder Bewegungen und Früchten können wir abnehmen, daß sie ein sehr köstlich, edel und wunderbarlich Wesen sein muß, daß so stark und kräftig ist, daß der Leib desselben Gewalt oder Sturm nicht ertragen kann. Alles, wenn etwa im Herzen ein großer unmäßiger Schmerz entsteht und in dem Menschen herrscht, so muß nothwendig der Leib dahinsterben. Und dasselbe geschieht auch also bei großer, übermäßiger Freude; wie sehr viele Exempel in den Historien der Heiden bezeugen.

43. Solche Bewegungen sind auch in Joseph und seinen Brüdern gewesen; und er wird ohne Zweifel dem Tode nahe gewesen sein um der heftigen Bewegung willen seines Herzens; also daß er auch gar nichts mehr hat reden können, denn allein, daß ihm noch die Thränen über die Backen gelaufen sind.

44. Solche große Kraft hat die Seele, daß wo sie allein im Menschen herrscht, so muß der Leib mit allen seinen Sinnen, Riechen, Schmecken, Tasten etc. gar erstarren. Es ist ein wunderlich Creatürchen. Es hat zwar Aristoteles nach der Speculation leichtlich beschreiben können, was die Seele sei, da er sagt:

Anima est actus primus corporis organici physici etc. (Vierter: Die Seele ist das Princip des Lebens, die erste Form, Entelchie, des Körpers.) Aber er kann es doch nicht recht ausdrücken, was sie sei, denn allein obiectiv, das ist, wenn er sie mit etwas Anderem vergleicht. Aber recht verstehen oder empfinden, was die Seele sei, so viel ihre Affecten und die großen heftigen Bewegungen belangen, gehört nicht zu solcher Erkenntniß der Speculation oder da man sie mit etwas Anderem vergleicht, sondern will practicirt, geübt und versucht sein. Alles, wenn eine große unmäßige Freude den Leib einnimmt und schier tödtet, oder wenn Einem die Thränen vor großer Freude aus den Augen fließen; davon die Münche und andächtigen Väter viel schreiben, daß nichts süßer oder lieblicher sei denn solche Freudenthränen; wiederum aber, daß auch kein betrübter und bitterer Ding sei, denn das Weinen derer, die da Leid tragen. Darum ist, sage ich, die Seele eine wunderbare Creatur und uns unbekannt.

45. Dies erinnere ich darum, daß wir diese Offenbarung Josephs nicht nur oberflächlich ansehen oder lesen sollen, sondern sollen die große Bewegung der Herzen betrachten, dadurch beide Joseph und seine Brüder so gar erschrocken sind. Ich wüßte nicht, wie ich mich dazu hätte stellen sollen, da er sagt: „Ich bin Joseph.“ Und sie wissen es auch nicht. Was will denn werden, wenn unser Herr und Heiland Jesus Christus kommen wird, welchen uns in diesem Leben auch auf mancherlei Weise versucht, und uns sehr geplagt, gekäuft und getödtet werden läßt. O welch eine große Freude wird werden, wenn er sich so plötzlich und unversehens ändern wird, da wir zuvor empfunden und uns dünken lassen, daß er der greulichsten Tyrann wäre, welcher uns alle miteinander in einem Hui wollte umbringen! Da wird er sagen: „Ich bin Joseph“, euer Heiland.

46. Ich kann diese Offenbarung Josephs nach ihrer Würde nicht genugsam auslegen; sondern man soll sie oft lesen und bedenken, auf daß wir lernen nach seinem und seiner Brüder Exempel fleißig beten, hoffen und des göttlichen Trostes harren, und daß wir inzwischen mit starken beständigen Herzen die Hand Gottes dulden, der mit uns spielt und uns also versucht um der Schlafsucht willen, damit wir behafet, und um der Kälte willen der Erbsünde, davon die Herzen der Gottlosen faul und träge, dazu auch schwach und matt werden zum Glauben, Hoffnung und zum Gebet, und andern geistlichen Uebungen. Denn wo dies göttliche Spiel nicht ist, so schnarchen wir und sind gar kalt. Deshalb schickt GOtt gleichsam mit diesem Stachel und treibt den groben faulen Esel fort, unser sündlich Fleisch, welches uns mit seiner Last sehr beschweret.

47. Nun müssen wir aber auch etwas sagen von der Gottseligkeit Josephs, welche wahrlich wunderbar, und wohl werth ist, daß man dieselbe wohl erwäge. Denn der große Fürst und Heiland in Egypten und umliegenden Ländern hätte zwar wohl mit großem Stolz seine Brüder gar verachten, und seine Herkunft und Verwandtschaft vor ihnen verheben können. Wie oftmals Exempel von solchen gefühllosen, lieblosen Leuten vorkommen, oder die solche Klöße und so gar steinhart sind, daß sie ihren Eltern oder Verwandten gar keine Ehrerbietung oder Freundlichkeit erzeigen können oder wollen. Denn wiewohl die Natur also von GOtt geschaffen ist, daß sie freundliche Neigung hat, und noch am Menschen ein Brünnlein guter Affecten oder Bewegungen bleibt (davon die Scholastiker den Spruch im 4. Psalm V. 7. verstanden haben: „Das Licht deines Antlitzes ist über uns erhoben“; als, daß der Vater von Natur sein Kind lieb hat, und wiederum das Kind den Vater auch; desgleichen sind allen Menschen solche natürliche Neigungen eingepflanzt zwischen Brüdern und Schwestern): doch werden solche Neigungen durch Wollüstigkeit der Menschen, durch Zorn, Hoffahrt, unbillige Feindschaft und Neid bisweilen verwirrt und verderbt.

48. Joseph aber war des Heiligen Geistes voll und hat die natürliche Neigung nicht ausgelöscht, sondern dieselbe gemehret. Denn der Heilige Geist wird nicht darum gegeben, daß er die Affecten und Neigungen, so in die Natur gepflanzt, wegnehmen und aufheben soll; sondern daß er die erstorbene Natur wieder lebendig machen, die schwache und schlafende Natur auferwecken soll, und wo sie irrt und aus dem rechten Wege getreten ist, bessere und wiederum zurechtbringe. Darum ist der Heilige Geist und die Gnade Gottes eine Arznei der Natur, dadurch dasjenige alles, so in dem angebornen Verderben durch die Sünde in uns geschwächt und verdorben ist, wiederum aufgerichtet und zurechtgebracht wird.

49. So hat nun Joseph deshalb seine Brüder oder Vater nicht verachtet, weil sie Schafhirten, unbekannte Leute und fast von den Allergeringsten im Volk waren, vor welchen die Egypter einen Greuel gehabt. Denn er hat sich müssen fürchten, daß ihm solches bei den Räthen und Amtleuten im selben Königreiche Schaden thun würde, und daß seine Ehre und Herrlichkeit, und das große Ansehen, so er gehabt, dadurch geschwächt werden möchte. Denn wir haben gehört, daß die Räthe mit ihm nicht gegessen haben, und scheinen sie jetzt zweierlei Ursachen haben ihn billig verachten und hassen zu dürfen: weil er ein hebräischer Mann war und weil er seinen Ursprung von dem geringsten Geschlecht der Viehhirten gehabt. Darum ist das eine schöne herrliche Tugend an ihm, daß er mehr gehalten hat vom vierten Gebot und natürlicher Neigung, denn von aller seiner großen Ehre und Herrlichkeit, ja, auch von seiner Gewalt, welche nicht gering, sondern fast die höchste und gleich des Königes Gewalt im Egypten gewesen ist. Er schämt sich dran keines Vaters und Brüder nicht, darum daß sie Schafhirten waren. Er hat wohl bisher ein trauriges und scharfes Spiel mit ihnen gespielt, doch aus großer Liebe und rechter natürlicher Neigung, so er zu ihnen getragen, da er sich besorgt, sie möchten vielleicht seinen alten Vater und jüngsten Bruder auch etwa getödtet haben. Da er aber das alles nun erfahren, daß dieselben noch am Leben sind, da thut sich nun sein Herz recht auf, und bricht hervor die große herzliche Liebe und höchste Freude gegen seine Brüder, so, daß das ganze Haus Pharaos ihn vor gar zu großer Freude hat weinen und heulen hören.

50. Also sehen wir in dieser ganzen Historie Josephs nicht denn eitel große unzählige Tugenden an ihm, und dasselbe in allen Lagen seines Lebens, im Gefängniß, in Schmach und Elend und da er gar verlassen gewesen, und darnach auch in seiner großen Ehre, Hoheit und Gewalt. Denn also ist er auf und nieder, vor sich und hinter sich getrieben worden; und ist doch allezeit beständig geblieben, und hat GOtt vor Augen gehabt und auf seine Hülfe mit sehr großem Glauben und Geduld geharrt.

\quellseite{432}

51. Es wird uns aber hier nicht allein ein Exempel vorgehalten aller Tugenden, so Joseph an sich gehabt, sondern es wird uns auch sein vor Augen gestellt eine rechte Beschreibung GOttes, auf daß wir wissen mögen, was GOtt sei. Die Philosophen disputiren viel und forschen nach GOtt mit ihrem Speculieren, und lernen ihn einigermaßen kennen; wie Plato die göttliche Regierung ansieht und anerkennt. Aber es ist alles lauter Speculation ohne Anwendung; es ist noch lange die Erkenntniß nicht, so Joseph gehabt, nämlich, daß GOtt für uns sorge, daß er die, so betrübt und elend sind, erhöre und ihnen helfe. Das hat Plato nicht schließen können; der bleibt immer bei seinem Speculiren und übernatürlichen Gedanken, wie eine Kuh ein neues Thor ansieht.

52. Das ist aber die rechte Erkenntniß GOttes, dadurch wir den Willen GOttes kennen, und daß er allein weise sei, und helfen könne und sich unser auch gern erbarmen wolle, also, da Joseph das für gewiß hält, und sagt: GOtt wird mich nicht verlassen, wenn ich auch sterben sollte; item, daß er also kann schließen: GOtt sieht dich darum nicht an, oder sorgt nicht darum für dich, daß du so groß und mächtig bist, wenn du auch schon über die ganze Welt ein Herr wärest. Auf der linken Seite verzagt er nicht, wo ihn das Unglück drückt; zur rechten aber erhebt er sich auch nicht, wenn es ihm wohl geht und gut Glück vorhanden ist.

53. So ist nun dies die rechte Erkenntniß GOttes, daß man wisse, daß seine Natur und Wille der ist, welchen er in seinem Worte geoffenbart hat, da er verheißt, daß er wolle mein Herr und mein GOtt sein, und heißt mich, denselben seinen Willen im Glauben fassen. Denn da ist der rechte gewisse Grund gelegt, damit die Herzen zufrieden sind. Es folgt aber auch die praxis, das ist, die Uebung und Erfahrung hernach, da es ins Werk gesetzt wird, da Joseph nun mit Elend, Gefängniß, Gefahr und falscher Anklage geübt und versucht worden ist. Da hat er also gedacht und bei sich geschlossen: Ich sehe, daß dies GOttes Wille ist, der ohne Zweifel etwas Größeres und Höheres damit bedenkt und vorhat, als ich jetzt verstehen kann; ich bin dessen gewiß, daß er mich nicht verlassen oder verwerfen wird; sein Wille ist unbegreiflich, aber doch sehr gut; und ich zweifle gar nicht, er will und kann mich aus diesem Unglück auch wohl erretten.

54. Also hält er sich aufrecht und besteht fest in der äußersten Noth und Gefahr seines Leibes, seines guten Namens, dazu aller seiner Güter und Wohlfahrt, so er in dieser Welt gehabt. Denn nachdem er ins Elend und Knechtschaft verstoßen worden, hat er dazu auch das Lob und guten Namen der Keuschheit und Unschuld, derwegen er bei jedermann ist gerühmt und lieb und werth gehalten worden, verloren, so daß er vor den andern Hausgenossen und vor der ganzen Welt hat Sünden müssen und gar verachtet sein. Und über das, daß er ein solcher verachteter elender Mann gewesen, vor welchem die Egypter einen Greuel gehabt, wird er auch noch um des Ehebruchs willen, damit ihn Potiphars Weib beschuldigt, werth geachtet, daß ihn billig jedermann sollte feind sein und daß er mit allerlei Plage und Marter billig zu strafen sei. Doch beharrt er noch immer in dieser Hoffnung und Glauben, daß er schließet und sagt: GOtt wird mich nicht verlassen; was er aber mit mir vorhat, weiß ich allerdings nicht. Ich werde in der Welt hin und her geworfen und umgetrieben, und wohl versucht; ich weiß aber, daß er auf ein Besseres und Schöneres denkt, denn ich jetzt wissen oder verstehen kann. Ich weiß, daß er mir nach seinem gnädigen Willen wohl gewogen ist; darum befehle und ergebe ich mich dem Herrn, meinem GOtt, der allein weise, allmächtig und gnädig ist.

55. Derhalben, wie er glaubt, also geschieht es auch. Denn er hat nimmermehr denken oder hoffen können auf die große Ehre und Herrlichkeit, dazu er gekommen ist. Dieweil er aber die Hand GOttes hat ertragen können bis in den Tod, darum wird er auch zeitlich erhöht und mit leiblicher Ehre geziert. Und das heißt GOtt recht erkennen, nämlich, daß man es also im Werke und mit der That erfahre, daß er allein weise sei, daß sein Wille der beste und angenehmste, unserer aber närrisch und sehr böse sei. Wie St. Paulus sagt Röm. 12, 2.: „Daß ihr prüfen möget, welches da sei der gute, der wohlgefällige und der vollkommene Gotteswille.“ Darum ist Joseph ein trefflicher großer Doctor und Lehrer der Weisheit und des Willens GOttes.

56. Wiederum aber, da er nun erhoben wird zur höchsten Würde des Erdkreises, auf welchem das Königreich Egypten das blühendste Reich gewesen ist, demüthigt er sich nun also, daß er die große Ehre gering achtet, und denkt, daß sei noch das rechte Gut oder Erbe nicht, so ihm von GOtt verliehen war. Darum wird er nicht stolz, schämt sich auch seiner Freundschaft und Geschlechtes, davon er geboren war, gar nicht, und bleibt in vierten Gebote sein beständig; wie er zuvor auch mit höchstem Fleiß und auf das Allerelrlichste das sechste und siebente Gebot gehalten, nämlich, da er seinem Herrn treulich und aufrichtig gedient und ihm sein Gut und Habe gemehret, und seine Frau bei Ehren behalten. Sind das nicht große treffliche Tugenden? Er denkt nicht also: Wenn ich meinen Vater und Brüder erkennen werde, so werden mich die Egypter bald von diesem hohen Amte, dazu ich gekommen bin, stoßen und mich derselben Herrlichkeit entsetzen; sondern er zieht seinen Vater und Brüder vor, hält von ihnen mehr, und läßt sich dieselben lieber sein, denn seine große Herrlichkeit, wie geringe und schlechte Leute sie auch gewesen sind, vor welchen das ganze Königreich Egypten einen Greuel gehabt; und er bekennt gern, daß er eines Hirten Sohn sei, und erzeigt seinem Vater nach GOttes Gebot große Ehrerbietung.

57. Das ist wahrlich der größte Muth, der auf ein viel Höheres hofft, denn das Königreich Egypten war, oder die ganze Welt sein kann; denn er hat den Glauben und die Hoffnung von der Auferstehung der Todten.

58. Dies alles lehrt die Historie, und die heftigsten Affecten und Bewegungen Josephs. Er hat sich im Gefängniß vor dem Tode nicht gefürchtet: da er oben nun erlöst worden, hat er auf das Königreich auch nicht gehofft, achtet auch der Gewalt und Macht dieser Welt nicht; sondern das Höchste und Größeste in der ganzen Welt achtet er in seinem Herzen gering. Darum hofft und wartet er auf ein ander und besser Leben. Denn das ist je ein groß Ding, daß er das Königreich Egypten hat verachten können um seines Vaters willen, der ein Schäfhirte war, und um seiner Brüder willen, so ihn verkauft, und den Vater mit ihm, so viel an ihnen gewesen, getödtet hatten. Aber also ist das Herz geschickt, wenn der Heilige Geist dazu kommt, der erzeigt den Eltern darnach Ehrerbietung von ganzem Herzen.

59. Joseph hätte sie wohl auf andere Weise versorgen können, wenn er ihnen etliche tausend Gulden gegeben hätte, auf daß seine Ehre und großer Name nicht gefährdet werden möchte. Aber er hat sich so tief und innig gedemüthigt und seine herzliche Neigung gegen seinen Vater und Brüder bezeugt, daß er darüber den ganzen Hof Pharaos mit großem Heulen und Schreien erfüllt hat. Das heißt recht seine Eltern geehrt und zwar mit großer Liebe und Ehrerbietung. Und solches ist sonderlich bei einem so großen Mann eine große Zierde.

60. Darum sehen wir, daß er mit den höchsten Tugenden, welche niemand mit Worten genugsam ausreden kann, begabt gewesen, und dasselbe auf beiden Seiten, nämlich im Glück und Unglück: welches denn uns zum Exempel und Lehre vorgehalten wird.

61. Darnach wird uns auch in diesem Bilde GOtt sein herrlich abgemalt, nämlich, was seine Art und Natur sei, auf daß wir ihn lernen erkennen und ihm vertrauen. Denn seine Art und Natur ist allezeit einerlei und wird uns unsertwillen auch nimmermehr geändert werden. Das ist aber seine Natur, daß er langmüthig ist, daß er die Seinen versucht und wunderlicher Weise übt und umtreibt, daß er sie allgemach zum Tode führt. Wiederum aber, wenn es an die äußerste Gefahr Leibes und Lebens gekommen ist, daß er alsdann plötzlich und eilends denen, so dermalen abgeschnitten werden, zu Hülfe kommt und sie errettet. Darum soll man auf GOtt hoffen und sein harren; wie die Exempel und Vermahnungen in den Psalmen und Propheten uns desselben erinnern.

62. Ueber das ist er auch ein solcher GOtt, der nicht allein aus Nichts Alles macht, sondern auch aus Allem Nichts machen kann. Als, anfänglich hat er Joseph zu Nichts gemacht, und da es nun schien, als ob er gar verdorben und als ob es mit ihm ganz verloren wär, da macht er wiederum Alles aus ihm, das ist, er macht ihn zum allergrößten Mann in der Welt. Derselbe große Mann aber äußert sich nun wiederum auch aller seiner Ehre und Herrlichkeit, steigt von seinem königlichen Stuhl herab, und nimmt sich der allergeringsten Leute an und hält sich zu denselbigen, nämlich seinem Vater und Brüdern, und sagt: „Ich bin Joseph.“ Hier wird er nun abermals zunichte gemacht, und ist wiederum eben derselbe Joseph, so zuvor in die Grube geworfen und verkauft worden war, und die schwere Knechtschaft hat tragen sollen. Er wird wieder eines Hirten Sohn, spricht: Ich bin euer Bruder, eines Hirten Sohn, gleich wie ihr.

\quellseite{433}

63. Derhalben sehen wir hier, was ein recht göttlich und männlich oder ritterlich Herz ist, nämlich, das immer rechten Glauben und Vertrauen hat zu GOtt auch in höchster Armuth und allerlei schweren Anfechtungen; wiederum aber, im Glück sich auch nicht erhebt oder stolz wird, sondern GOtt fürchtet. Solche Leute können wahrlich die Welt wohl regieren, die mit solchem Herzen Alles verwalten, daß sie nicht zufrieden sind mit den Gütern dieses Lebens, sondern hoffen und warten auf ein ander Reich, und behalten doch noch immer die natürlichen Neigungen und Affecten, so in die Natur gepflanzet sind. Denn der Heilige Geist löscht dieselben nicht aus, sondern richtet sie wunderbarerweise wiederum auf, zündet sie an und erhält sie auch. Welches fleißig zu merken ist um der teuflischen Lehrer willen im Papstthum, so gelehrt haben die Eltern verachten, als sollte dasselbe eine treffliche Gottseligkeit sein; und haben die Klöster, darin sie gelegen, hoch gerühmt, vor welchen man die Eltern wohl möchte verachten und gleichsam unter die Füße treten.

64. Wie man einen erschrecklichen Spruch des Hieronymus hat, welcher also lautet: Wenn schon dein Vater vor deiner Thüre liegt und jämmerlich weint, und die Mutter ihren Leib und die Brüste zeigt, damit sie dich gesäuget, nur immer dieselben mit Füßen getreten und stracks zu Christo gelaufen. Behüte GOtt, welch ein hartes Wort ist das! Und wiewohl ich glaube, daß er selig geworden, wollte ich doch nicht gern eines solches Spruchs Urheber sein.

65. Und sie sind nicht zu Christo gelaufen, wie sie meinten, sondern sind von Christo in das Kloster gelaufen. Denn, sind die Eltern, Weib und Kind nicht auch heilig und Christen? St. Petrus hat Weib und Kind gehabt; aber da er gefordert und hingerissen worden ist zur Marter und Leiden, da hat er sie recht verlassen. Oder, ist es denn so ein köstlich Ding, in das Kloster gehen, was soll man von den Patriarchen und andern Heiligen halten? Und sie haben doch vor Zeiten also gelehrt, und der Papst hat solch gottlos Wesen bestättigt mit seinen Decretalen, darin er die Heiligkeit und höchsten Tugenden der Christen in das Klosterleben und Orden der Mönche gesetzt hat. Das ist aber keine Gottseligkeit, sondern ein rechter Greuel und kräftiger Irrthum; wie St. Paulus zuvor gesagt 2 Thess. 2, 11.

66. Wenn du nun Christi willen deine Eltern verlässest, so thust du recht und wohl daran; aber wegen deines Aberglaubens, oder aus freiem Willen und selbst erwählter Andacht sie also verachten und gleichsam mit Füßen treten, das ist weder christlich noch menschlich, sondern teuflisch. Ich habe oft gesehen, daß die Mütter ihren Söhnen, wenn sie in das Kloster gegangen, weinend nachgefolgt und häßlich und schmerzlich bekümmert gewesen sind: dieweil aber ich auch mit demselben Irrthume behaftet war, pflegte ich solche weibische Wehmuth zu strafen, und hielt es dafür, daß die Mütter daran nicht recht thäten, da sie sich darüber betrübten, kunten ihre Söhne heilig und selig wären. Denn also waren wir gelehrt und verführt worden, daß dies der höchste Gottesdienst und Gerechtigkeit sei, wenn man die Gewalt der Eltern, so sie über ihre Kinder haben, verachten könnte.

67. Die Mönche in Egypten haben diese Irrthümer erstlich erdacht und aufgebracht; der Papst aber hat sich dieselben wohlgefallen lassen und mit sonderlichen Gesetzen bestättigt. Und thut eben dasselbige damit auch, daß er die heimlichen Verlöbnisse vertheidigt, auf daß er ja wider GOtt und das vierte Gebot streite, und daß er in Wahrheit derjenige sei, wie er von Paulo 2 Thess. 2, 3, 4. genannt wird, nämlich, „der Mensch der Sünden, und das Kind des Verderbens, der da ist ein Widerwärtiger, und sich überhebet über alles, das GOtt und Gottesdienst heißet.“

68. Darum sollen wir nun wissen, daß GOtt die natürliche Neigung in dem Menschen nicht verwirft oder aufhebt, welche er der Natur in der Schöpfung eingepflanzet hat, sondern daß er dieselbige erweckt und erhält. Wenn es aber die Noth also erfordert, alsdann sollst du mit Joseph sagen: „GOtt hat mich lassen vergessen alle meines Vaters Hause“, 1 Mos. 41, 51. Oder, da Abraham geboten wird, mit einem sonderlichen Befehle GOttes, daß er seinen Sohn Isaak schlachten soll, alsdann soll er sagen: „Isaak, ich kenne dich jetzt nicht;“ oder umgekehrt auch also, wie 5 Mos. 33, 9. geboten wird. Wenn dich aber GOtt nicht beruft außerhalb oder über das vierte Gebot, da sollst du wissen, daß du ihm nothwendig gehorsam sein mußt.

69. Derhalben ist der Papst und die Mönche ein schändlich Ungeheuer, ja ein wüster Greuel, welcher Greuel steht an der heiligen Stätte, wie Christus daselbst Matth. 24, 15. aus dem Propheten Daniel Cap. 9, 27. angezogen hat. Denn es kann niemand glauben, welch ein scheußlicher und greulicher Irrthum das gewesen sei, der doch für den höchsten Gottesdienst ist gehalten worden, darum daß die Vernunft leichtlich kann gefangen und betrogen werden mit dem äußerlichen Scheine der Heuchelwerke.

70. Darum soll man solche Exempel fleißig ansehen, auf daß die Herzen wider diese Unsinnigkeit des Satans gerüstet werden. Wir sehen, wie kräftig die natürlichen Affecten und Bewegungen an Weibern, an Joseph und seinen Brüdern, sind, welche der Heilige Geist mit sonderlichem Fleiß hat rühmen wollen. Denn Joseph hat sie endlich mit diesem Worte aufrichten und trösten wollen, da er sagt: „Ich bin Joseph.“ Ich bin nicht der Abrech, der große Herr über Egypten, der bisher mit euch durch einen Dolmetscher geredet hat; jetzt brauche ich nun meiner angebornen Muttersprache. Aber eben mit denselbigen Worte hat er sie getödtet.

71. Nun setzt er noch einen andern Trost hinzu, da er sie fragt: „Lebet mein Vater noch?“ Dasselbe ist nicht allein so zu verstehen, als ob es nur eine Frage wäre; sondern wie ich gesagt habe, es ist ein solch Wort, damit angezeigt wird, daß er sich verwundert habe, daß er noch leben sollte. Lieber, sagt er, lebt er noch? Als wollte er sagen: „Ich habe zwar meines Vaters Hauses vergessen, darum habe ich so fleißig nach allem Dinge geforscht; aber euer Vater ist mein Vater auch, ihr seid meine Brüder.“ Da erstaunen sie gar und sterben schier dahin; darum sagt Mose: „Sie konnten ihm nicht antworten, so erschraken sie vor seinem Angesicht.“

72. Und das ist kein Wunder gewesen, dieweil sie dazu sehr viele und wichtige Ursachen gehabt, daß sie sich dermaßen entsetzt haben, nämlich, daß sie solchen großen Fürsten zum Widersacher gehabt, der mit ihnen gespielt bis an den Tod, und nun so plötzlich herausfährt und sich ihnen in seiner Muttersprachen offenbart. Denn sie hatten daran nicht gedacht, daß sie in dem Lande je die hebräische Sprache hören würden. Derhalben, da sie so heftig bewegt und in ihrem Herzen so gar erschreckt worden sind, macht derselbe Schrecken so viel, daß sie schier zu Steinen und Klötzen werden. Denn das Innerste ihres Herzens wird in ihnen gar bewegt vor großer Freude und Furcht, also daß sie nicht antworten können, und wären ihr Lebenlang da gestanden, da sie dasselbe Wort gehört: „Ich bin Joseph.“ Sie können das gar nicht fassen oder verstehen; darum muß er noch so viel Worte brauchen, auf daß er sie erwecken möge.

73. Sie sind noch erschrocken und dürfen ihr Maul nicht aufthun, und können auch noch nicht glauben, daß dieser ihr Bruder sei, darum daß ihnen die große Majestät seiner Person im Wege liegt. Und es will sich doch gleichwohl nicht geziemen, ihn zu fragen, oder aber ihm zu widersprechen. Derhalben stellt er sich selbst ihnen vor die Augen und gibt ihnen eine genaue Beschreibung seiner Person. Ich bin kein neuer oder erdichteter Joseph, sagt er, sondern bin eben der Joseph, deß Vater Jakob ist, welches Söhne ihr auch seid; der Joseph bin ich, der euer Bruder ist. Dies hören sie alle; aber der große Schrecken hindert sie, daß sie ihr Herz dahin nicht richten können, dieser Sache nachzudenken, oder ihm darauf Antwort zu geben.

74. Er thut aber noch mehr dazu, und sagt: „den ihr verkauft habt.“ Ist aber das nicht eine solche Beschreibung, so klar genug ist? Ich halte zwar dafür, sie werden es bald verstanden haben, daß dieser nun wahrhaftig ihr Bruder wäre; denn sie hatten ja keinen Andern verkauft denn Joseph. Darum sollte ja bei ihnen daran gar kein Zweifel mehr gewesen sein. Zudem, so hätten sie solches aus den Dingen, so bisher zwischen ihnen von beiden Seiten geredet und gehandelt waren, abnehmen können, als, daß das Geld in die Säcke gelegt, daß der Becher in Benjamins Sack geworfen war. Doch haben sie ihm noch nicht Beifall geben können; denn der Schrecken war gar zu groß und hat sich allzu tief und hart in ihr Herz gesetzt. Ei, lieber GOtt, soll das Joseph sein! denken sie stets.

\quellseite{434}

75. Es hat aber Mose droben, V. 2., gesagt, daß die Egypter, so hinaus gegangen waren, sein Schreien und Heulen gehört haben. Jetzt aber redet er sie nun an nicht mit so gar lauter Stimme und heißt sie auch darum etwas näher herzu treten. Wie Lyra angezeigt hat, hat er solches darum gethan, daß die Egypter nicht hören sollen, daß er dies hinzu gesetzt hat, wie sie ihn verkauft hätten. Joseph ist ein Mann voll Heiligen Geistes gewesen; derhalben, dieweil es die Noth erfordert, daß er sich ihnen geoffenbart, daß er sie hat müssen erinnern, wie sie ihn verkauft hätten, auf daß sie seinen Worten desto leichter möchten Beifall geben, hat er gleichwohl nicht haben wollen, daß bei den Egyptern dieselbige Sünde ruchbar werden sollte.

76. Das ist die rechte Brüderschaft, nämlich, seiner Brüder Sünde verzeihen und zudecken. Er will nicht, daß sie darüber sollten in ein böses Gerücht kommen, sondern begehrt, daß ihre Ehre und guter Name bei den Egyptern unverletzt bleiben möge, auf daß sie nicht verdammt oder gescholten werden als Brüder- und Vatermörder. Darum ist dies ein trefflich Exempel rechter wahrhaftiger Liebe und Vergebung der Sünden.

77. Er kann nun so ihnen deutlich beweisen, daß er ihr Bruder sei, daß er zugleich ihrer Sünde gedenkt: darum redet er auf das sanfteste und freundlichste mit ihnen, hat die Sünde nicht groß gemacht nach ihrer Würde. Er sagt nicht: Ich bin der, welchen ihr boshaftig in die Grube geworfen, denn ihr heimlich nach Leib und Leben gestanden, die ihr euren Vater so bitterlich betrübt habt; sondern sagt nur fein und stille: „den ihr verkauft habt.“ Er läßt alles fahren, was etwas hart gelautet. Als wollte er sagen: Ihr wißt gar wohl, wie ich hieher gekommen bin: ich habe mich von selbst von eurer Gemeinschaft nicht gesondert, oder meinen alten Vater verlassen; sondern bin hieher geführt worden, da ihr mich verkauft hattet.

78. So ist der Heilige Geist in den Vätern, die ihren Brüdern oder Nächsten übersehen, ihnen vergeben und ihre Sünde zudecken können. Er hat nothwendig des Verkaufens gedenken müssen, da er sich ihnen geoffenbart, auf daß sie nicht meinten, daß dies alles erdichten wäre. Darum, auf daß sie ein gewisses Zeichen oder ein solches Anzeichen hätten, dem sie nicht widersprechen könnten, so hat er dies heimlich zu den andern Worten hinzugefügt, welches bei den Egyptern gar verborgen und unbekannt sein soll. Darum heißt es hier zu treten, dieweil sie zuvor etwas weiter von ihm gewichen, und sich nicht unterstanden hatten, zu ihm zu nahen. Das ist wahrlich rechte wahrhaftige Vergebung der Sünden und eine reine ungefärbte Liebe.

\begin{quote}
\textbf{V. 5. Und nun bekümmert euch nicht, und denket nicht, daß ich darum zürne, daß ihr mich hieher verkauft habt.}
\end{quote}

79. Hier hat er sie nun absolvirt und von ihrer Sünde los gesprochen, und hat sie nicht allein absolvirt, sondern verkündigt ihnen auch sehr reiche Gnade und daß er ihnen von Herzen gewogen sei. Ich bin der, den ihr verkauft habt, sagt er; eure Sünde ist offenbar und jedermann wohl bekannt. Wiewohl sie nun dasselbe nicht bekannt hatten, so kann es doch nicht geleugnet werden, es ist an sich offenbar.

80. Siehe aber, wie gnädig und freundlich er mit denen redet, die ihn doch verrathen und verkauft hatten. Da spürt man an seinen Worten gar keine Nachjagd, keinen bösen Willen, oder daß er vorhätte sie zu beleidigen; sondern nur Barmherzigkeit und Güte, so von Grund seines Herzens herkommt. Es ist wahr, sagt er, ich bin von euch verkauft worden; ihr wollet euch aber darüber nicht betrüben. Und steht hier an diesem Ort eben das Wort, welches Christus auch braucht Joh. 16, 21., da er sagt: „Ein Weib, wenn sie gebiert, so hat sie Traurigkeit.“ 1 Mos. 3, 16.: „Ich will dir viel Schmerzen schaffen, wenn du schwanger wirst“ etc., das ist, viel Kummer und Jammer etc. Dieses Wort heißt ezeb, und wird hernach auch auf Götzendienst oder Abgötterei übertragen. Alles, Ps. 16, 4., da in unserer lateinischen Uebersetzung steht: Multiplicatae sunt infirmitates eorum etc., da sollte lieber stehen idola, wie wir es verdeutscht haben: „Jene, die einem andern nacheilen“ etc. Denn Abgötterei und Aberglaube, so die Gottlosen treiben, ist Jammer und Noth, damit sie geängstigt und gemartert werden, wie ein Weib, das da gebären soll. Item im 127. Psalm V. 2.: „Es ist umsonst, daß ihr frühe aufstehet, und esset euer Brod mit Sorgen“, das ist, euer Brod wird euch zu sauer: ihr sollt euch nicht also ängstigen.

81. Also sagt Joseph hier auch: Ich will nicht, daß sich euer beladen und beschwert Gewissen ängstigen, fressen und martern soll; sondern ich sage und will haben, daß ihr mit meiner Gunst und gutem geneigtem Willen zufrieden und fröhlich darüber sein sollt. Denn ich will euch diese Sünde gern vergeben; ihr dürft euch darum nicht bekümmern. Dies ist eine sehr schöne Absolution. „Bekümmert euch nicht“, sagt er, das ist, denkt nicht, daß ich darum zürne, seid guter Dinge, seid nicht unmuthig; ihr habt keine Ursache, traurig und betrübt zu sein. Wie ein armes Gewissen, das sich schuldig weiß, billig ängstlich ist und zittert. Denn daher kommt aller Schrecken und alle Traurigkeit nach dem gewöhnlichen Spruch der Theologen: Si est dolor et tristitia, nemo novit, nisi mala conscientia, das ist: Wo Schmerzen und Traurigkeit ist, weiß davon niemand etwas denn ein lauter böses Gewissen. Und wiederum, ist Fröhlichkeit auch nirgendwo als in einem guten Gewissen. Wie Salomo sagt Spr. 14, 10.: „Wenn das Herz traurig ist, so hilft keine äußerliche Freude“; und Spr. 15, 15.: „Ein guter Muth ist ein täglich Wohlleben“ etc. Was sonst andere Freude ist, die kann nicht vollkommen sein; sie mag den Menschen wohl äußerlich etwas kitzeln, aber sie kommt inwendig in das Herz nicht.

82. Es ist auch kein schwerer Leiden denn Traurigkeit des Herzens; denn das ist der Tod und die Hölle selbst. Darum muß man denselben zuvorkommen, es beschwere gleich dich selbst oder über einen andern. Da schließe die Hölle zu und den Himmel auf, wer da kann, auf daß ein solch schwach betrübtes Herz nicht gar dahinsterbe, wenn es die Sünde fühlt und darüber solche Schmerzen hat.

83. Wie Joseph alhier auf das freundlichste und lieblichtste mit seinen Brüdern handelt, auf daß er ihnen Muth mache; denn er ist ein weiser Mann gewesen, der sich auf solchen Kampf des Gewissens wohl verstanden hat: darum hat er in der Zeit die rechte Arznei gebraucht, dadurch dieser Schmerz gelindert werden möchte. Bekümmert euch nicht, sagt er, und denkt nicht, daß ich darum zürne; es ist alles verziehen und vergeben, dazu auch gar vergessen; ich habe euch brüderlich und von Grund meines Herzens lieb.

84. Das heißt nun recht, von Herzen seinem Bruder die Schuld vergeben, vornehmlich wenn er in solcher Angst und so gar bekümmert ist. Damit denn Joseph bis daher umgegangen ist, daß er seine Brüder demüthigen möchte, auf daß sie an ihrer Wohlfahrt und Wiederkunft zu ihrem Vater bei sich selbst gar verzagen, und nichts Anderes vor ihren Augen haben, denn nur gewisse ewige Knechtschaft und Wehklagen. Da sie aber dermassen betrübt und erschreckt worden, hat er ihnen nun die Sünde mit solcher Gnade und Freundlichkeit vergeben, daß er selbst auch mit den Weinenden weint.

85. Derhalben ist dies ein trefflicher Text, den man billig mit sonderlichem Fleiß merken soll. Denn es kann es niemand sagen oder glauben, wie es so ein schwer Ding sei, wo man ein verzagt Herz und Gewissen, das in Traurigkeit gar ersoffen ist, aufrichten und wiederum aufmuntern soll. Die ohne Anfechtung sind, können leichtlich fröhlich sein; aber das ist Mühe und Arbeit, wenn man den, so mit geistlicher Traurigkeit und Betrübniß beschwert daniederliegt, wiederum dahin bringen soll, daß er zu rechter Erkenntniß und Glauben kommen und sich der Gnade Gottes und seiner Güte trösten könne.

86. Wir sehen, wie Christus nach seiner Auferstehung mit so großer Langmuth und Freundlichkeit mit seinen Jüngern umgegangen ist, wie freundlich er sie angeredet, ihnen seine Hände und Seiten zeigt und sich von ihnen betasten läßt; er läßt sich hören und sehen, ist und unterredet sich mit ihnen auf das allerfreundlichste. Doch haben sie sich nicht bald wiederum aufraffen und zurechtkommen können. Und da der Engel den Weibern sagt, Marc. 16, 7., sie sollen den Jüngern verkündigen, daß Jesus auferstanden wäre, heißt er sie dasselbe ja vornehmlich Petro anzeigen, darum daß derselbe in der größten Angst war; der konnte das eben wohl decliniren. Er hatte dreimal den Herrn verleugnet, hatte sich selbst verflucht; die Andern waren davon geflohen. Darum wird von Noth wegen und zu rechter Zeit vom Engel das hinzugesetzt, daß sie es ja sonderlich Petro verkündigen sollen. Ach wie herzlich redet Sacharia Cap. 4, 6. ff. mit Zorobabel und den andern übrigen Juden, so aus dem babylonischen Gefängniß wiederum in ihr Land geführt waren, ehe er sie versammeln, stärken und gleichsam von den Todten erwecken kann. Denn es ist viel schwerer, ein betrübtes Gewissen trösten, denn Todte auferwecken.

\quellseite{435}

87. Also thut St. Paulus auch. Nachdem er in der ersten Epistel an die Corinther Cap. 5, 1. ff. den Mann, der seine Stiefmutter gefreiet und damit der christlichen Gemeinde groß Aergerniß gegeben, hart gestraft hatte, zeigt er in der andern Epistel Cap. 2, 5. ff. wiederum seine große Sorge an, so er für denselben Mann getragen, daß er ihm vielleicht nicht zu viel mit der heftigen Strafe gethan haben möchte; wiewohl er wider sein Amt nicht gethan hatte, und auch die vorige Strafe damit nicht aufgehoben hat, als hätte er damit unrecht gethan, sondern heißt die Corinther, ihn wiederum aufrichten und trösten, auf daß er nicht in der Traurigkeit gar ersaufen und verderben möge.

88. Den harten unempfindlichen Köpfen aber kann man nimmermehr genug sagen, daß sie ihre Sünde bereuen und sich demüthigen sollen; denn sie fürchten sich vor Gottes Zorn nicht, wenn schon Gott mit seinem Blitz zu ihnen hereinschlägt. Wie denn der Pabst und die Cardinäle nach seinem Drohen fragen. Wo aber Einer mit Blitz vom Himmel getroffen würde, der würde auch fühlen, daß das Gewissen, so ohne allen Trost ist, sich heftig bemüht und sehr bekümmert ist, und daß ihm nicht leichtlich zu rathen noch zu helfen ist. Denselbigen muß man nun mit Trost zu Hülfe kommen, so man aus dem Wort Gottes nimmt, als, Matth. 18, 14.: „Es ist nicht der Wille Gottes, daß jemand von diesen Kleinen verloren werde“; und wie Christus Petro befiehlt, Luc. 22, 32.: „Wenn du dermaleinst dich bekehrest, so stärke deine Brüder.“ Gott will haben, daß du sollst gedemüthigt werden, und nicht verderben oder verloren gehen. Er zürnt mit dir nicht; er hat dich auch mit diesem Blitz in seinem Grimm nicht getroffen, sondern mit dem Wort.

89. Dies sind Exempel des Heiligen Geistes, so in der heiligen Schrift erzählt werden uns zur Lehre und Trost. Denn es ist viel schwerer, einen aufrichten als tödten, wie die tägliche Erfahrung lehrt. Wir schlagen zwar oft mit dem Hammer göttlichen Worts auf den Fels, Cardinalen, Bischöfe und auf die Wucherer, gleich als auf einen Ambos, aber es ist alles vergeblich. Diejenigen aber, deren Herzen etwa mit dem Wort gerühret werden, damit man zu binden pflegt, erfahren, daß es wahr sei, was Gott durch den Propheten Jeremia Cap. 23, 29. sagt: „Ist mein Wort nicht wie ein Feuer, spricht der Herr, und wie ein Hammer, der Felsen zerschmeißt?“ Da ist kein Fels zu hart, er muß brechen. Paulus war ein sehr harter Fels; aber da ihn das Licht vom Himmel umleuchtet, und die göttliche Stimme gehört wird: „Ich bin Jesus, den du verfolgest“ etc., Ap. Gesch. 9, 5. ff., da fällt er bald darnieder, als wäre er todt.

90. Darum tröstet nun Joseph seine Brüder abermals, da er ihre Anfechtung und Angst sieht; wiewohl sie ihre Sünde noch nicht bekannt hatten. Ihr wollet euch, sagt er, nicht betrüben; wollet auch nicht denken, daß ich hiermit etwas geredet oder gethan habe, euch damit zu verderben, sondern vielmehr, daß ich euch damit zur Buße bringen möchte. Jetzt aber ist nun der Tod weggetrieben, die Hölle geschlossen, die Sünde vergeben und der Himmel aufgethan. Dies ist der rechte Pabst und Hohepriester und allhier hat die Absolution ihre rechte Statt.

91. Das aber wir hier im Text aus dem Hebräischen im Lateinischen gesetzt: Ne sit ira in oculis vestris: Bedeutet nicht, daß ich darum zürne: dasselbe haben die Andern nicht recht gegeben, da sie sagen: Non irascantur oculi vestri (Laßt eure Augen nicht zürnen). Denn Joseph redet von seinem und nicht von seiner Brüder Zorn. Und wie sollten sie zürnen, so sie sich doch jämmerlich fürchten und erschrecken vor Joseph, daß sie auch ihr Maul nicht haben aufthun dürfen. Und eben darum thut ihnen auch diese Stärkung noth, daß Joseph ihnen sein brüderliches und geneigtes Herz anzeige und offenbare. Ich dürft nicht denken, will er sagen, daß ich des Verkaufens euch zur Schmach gedacht habe, sondern vielmehr darum, daß ihr mich gemächlich kennen und mir vertrauen möchtet. Ich bin nicht erhitzt oder zornig auf euch, bin euch nicht abgeneigt, wie ich mich bisher angestellt habe; darum seid nicht traurig.

92. Jetzt aber tröstet und absolvirt er sie nicht allein, sondern hat auch noch dazu eine lange Predigt vom Rathe Gottes gehalten, dessen Weisheit und große Güte er nun auch rühmt. Denn das ist in der Christen Schule die vornehmste und höchste Weisheit, wie ich droben gesagt habe, nämlich, daß man lernen möge, was Gott sei, nicht nach der Speculation, sondern nach der Praxis, daß man es mit der That und im Werke erfahren möge. Derhalben wird er nun seine Brüder entschuldigen und die ganze Sache auf Gottes Rath und Weisheit legen.

\section[Dritter Theil]{Dritter Theil. Wie Joseph seinen Brüdern predigt von dem Rathe Gottes, der durch seine Verkauffung hinausgeführt ist}

\begin{quote}
\textbf{V. 5. Denn um eures Lebens willen hat mich Gott vor euch her gesandt.}
\end{quote}

93. Siehe doch, welch ein wunderbarer Doctor und trefflicher Theologe und Ausleger dies ist beider der Worte und auch der Werke Gottes. Ihr, sagt er, habt mich verrathen und verkauft, das ist ja gewiß: nun hört aber und lernt Gottes Rath in dem auch erkennen, was ihr gethan und nicht gethan habt. Ihr habt mich verderbet und getödtet, welches ja nicht geleugnet werden kann: aber dies Verderben und Tödteln, so ihr angerichtet, was ist doch dasselbe vor Gottes Augen gewesen? Denn es ist ein Anderes um Gottes Rath und ein Anderes um euren Rath in allerlei Werke. Gott hat euren sehr bösen Rath und Willen zum Leben gebraucht; nicht meinem Leben allein, welches ein geringes Ding wäre, sondern insgemein sagt er: Um des Lebens willen hat mich Gott hieher gesandt. Ihr habt mich nicht gesandt, sondern der Herr hat mich durch euch gesandt.

94. Das heißt wahrlich recht mit neuen Zungen reden. Denn was ist doch das für eine Art und Weise zu reden? Es ist wahrlich eine ungewöhnliche Art und fremde Weise zu predigen, Einen zum Heil und Leben senden, nicht allein des Gesandten, sondern auch der ganzen Welt. Eine solche Zunge ist nirgendes zu finden. Denn die Menschen würden also sagen: Die zornigen und neidischen Brüder verrathen, verkaufen und tödten ihren armen unschuldigen Bruder. Aber der Heilige Geist sagt: Dies Verkaufen soll das Heil und Leben Egyptenlands und der ganzen Welt sein. Dies ist eine neue unerhörte Rede und eine unglaubliche Antwort, die niemand hätte hoffen können. Es ist aber eine göttliche Art und Weise zu reden.

95. Derhalben ist dieser Text sonderlich wohl vom Rathe Gottes zu merken, welcher im Propheten Jesaia am 55. Cap. V. 8. 9. auch gerühmt wird: „Meine Gedanken sind nicht eure Gedanken, und eure Wege sind nicht meine Wege, spricht der Herr; sondern, so viel der Himmel höher ist, denn die Erde, so sind auch meine Wege höher, denn eure Wege, und meine Gedanken, denn eure Gedanken“; und Joh. 3, 31.: „Wer von der Erde ist, der ist von der Erde, und redet von der Erde; der vom Himmel kommt, der ist über alle.“ Das ist eine irdische Zunge, die da sagt: Joseph, der verkauft ist, der ist auch verloren und ist gar aus mit ihm. Aber eben dasselbe heißt auf himmlische Sprache eine Sendung zum Heil und Leben. Das ist unser Herr Gott.

96. Darum laßt uns lernen, daß wir unsern Gedanken nicht folgen, und daß wir auch nach unserm Rathe die Fälle oder Werke, und wie dieselben einen Ausgang gewinnen, nicht abmessen, richten und verstehen sollen.

97. Es wird zwar diese Sache vor den Menschen wohl menschlicher Weise ausgerichtet und die menschliche Zunge nennt diese Sendung ein Verderben und Untergang Josephs. Aber es ist zweierlei Vaterland und sind unterschiedene Orte des Himmels und der Erden. Darum sind die Werke und Weise zu leben auch nicht einerlei. Christus wird vor der Welt getödtet, wird verdammt und fährt nieder zur Hölle; aber das ist vor Gott Heil und Seligkeit der ganzen Welt vom Anfange bis zum Ende.

\quellseite{436}

98. Darum, in allen Werken und Gedanken, vornehmlich im Leiden, wenn die Christen geplagt werden, wo du gläubig bist, sollst du von deinem Leben und Werken nicht urtheilen; du wirst sonst irren. Du bist fromm, bist ein Narr, wirst versucht, bist gleichsam gefangen und kannst nicht recht reden oder urtheilen von deinen eigenen Sachen. Es heißt: „Harre des Herrn“, Ps. 27, 14., und ärgere dich nicht, murre nicht und verzage auch nicht; denn du gibst deinen Werken oder Leiden den rechten Namen nicht. Dein Urtheil ist falsch, deine Rede ist irrig, deine Weisheit ist Thorheit. Das ist aber Gottes Wille, daß der alte Mensch verderbet und das Fleisch getödtet werden soll. Wenn es aber also verderbet und getödtet wird, so redet es falsch und urtheilt eitel närrisch Ding.

99. Auf solche Weise läßt Gott den Joseph gefreuzigt und in das Gefängniß geworfen und geschmäht werden; welches ja nichts Anderes ist, denn ausgetilgt, weggeräumt und gar ausgelöscht werden. Denn vor der Welt geht dies auch anders nicht zu: aber derselbe ist vor Gott, wie hier steht, Einen zum Heil senden.

100. Dies ist eine solche Art und Weise zu reden, die Gott gemein und gebräuchlich ist. Mose redet nicht also, die Juristen und Philosophen auch nicht; sondern es ist eine Theologie der Verheißung, und sind Worte Gottes, der die Verheißung thut, welche Worte er unter seinen Engeln im Himmel redet; wir aber verstehen sie nicht, bis sein Rath vollbracht wird. Daher kommt es, daß so oft in den Propheten der Rath des Herrn gerühmt wird. Und im 107. Psalm 3. 11. steht auch: „Sie haben das Gesetz, oder Rath, des Höchsten geschändet.“ Dies soll man nun lehren, auf daß wir daraus berichtet werden, daß wir wissen mögen, was die rechte Theologie und Lehre des Evangeliums sei.

101. Die Lehre der Philosophen und des Gesetzes ist eine andere Lehre. Mose hat nicht so eine fertige Zunge, sondern sie ist ihm etwas langsam, und er lispelt und stammelt mit seiner Zunge; er kann dies Wort nicht recht ausreden. Der aber die Verheißung thut und der Heiland, Jehowah, redet also.

102. Die Natur aber, so närrisch, und das Fleisch, so durch die Sünde verderbet ist, verstehet solches nicht, denn nur, wenn das Leiden ein Ende hat; wie der Herr zu Mose sagt: „Du wirst mir hinten nachsehen“ 2 Mos. 33, 23., da er begehrete, daß ihm Gott sein Angesicht zeigen wollte, das ist: Du wirst erst nach der That meine Gedanken sehen. Wie hier Joseph und Jakob mit seinen Söhnen nicht sein Angesicht, sondern ihm auch hinten nachsehen. Denn da Joseph verkauft ward und sein Vater deshalb Leid trug, da war das Angesicht Gottes verborgen und ließ sich daselbst gar kein Gott sehen; sondern die ganze Welt schien voller Teufel. Nun aber, da sein Rath erfüllt und geoffenbart ist, sehen sie klärlich und erkennen auch nun den guten wohlgefälligen und vollkommenen Willen Gottes.

103. Also sagt droben, 1 Mos. 16, 13., Hagar auch: „Gewißlich hie habe ich den Rücken gesehen deß, der mich siehet.“ Als wollte sie sagen: Da mich Gott ansah, habe ich ihn nicht angesehen; jetzt sehe ich nun den Rücken deß, der mich bisher angesehen, da ich von allem seinem Rath und Werken gar nichts gewußt habe. So sagt auch Christus zu Petro Joh. 13, 7.: „Was ich thue, das weißest du jetzt nicht.“ Du wirst mir zukommen, und mich lehren, was ich thun soll: das ist weit gefehlet, lieber Petrus; denn dir gebühret, daß du meine Hand dulden oder ertragen sollst. Laß mich machen, darnach aber wirst du es wohl erfahren und verstehen, was ich mit diesem Fußwaschen gemeint habe.

104. So ist nun dies die Art und Weise, wie die Heiligen von Gott regieret werden, und ist die rechte Weisheit der Kirche und Gemeine Gottes, nämlich, daß man sich an dem Rath Gottes nicht ärgere, und sich an seinem Angesicht nicht stoße, damit er uns begegnet. Denn er pflegt dasselbe bisweilen einzuwickeln und zu verbergen, daß wir es nicht erkennen oder ansehen können. Er ist ja wohl ein Gott des Lebens, der Ehre, des Heils, der Freude und des Friedens, welches das rechte Angesicht Gottes ist; aber er verbirgt es etwa, und ziehet eine andere Larve an, damit er uns begegnet als ein Gott des Zorns, des Todes und der Hölle. Darum siehe zu, daß du weise und verständig seiest, lernest und hörest, was er doch meine oder vorhabe mit solcher ungewöhnlichen und fremden Gestalt. Denn damit wird das gesucht, daß du sollst gedemüthigt werden, und die Hand des Herrn und sein Angesicht, das er dir offenbart, ertragen und darauf harren lernest.

105. Es hätte sich zwar Jakob nicht unbillig an dem wunderbaren Werk ärgern mögen, daß sein Sohn Joseph also verkauft worden. Denn Gott hatte ja ihm und seinen Brüdern verheißen, daß sie sollten wachsen und gemehret, und ihr Same ewig erhalten werden. Was geschieht aber? Sein Sohn, welchen er sonderlich lieb hat, wird dahingerissen, der von seinem vornehmsten Weibe geboren war und auf dem allein die ganze Hoffnung der Nachkommen gestanden hat. Warum handelt denn Gott wider die Verheißung, so er gegeben hat? Er hat ja verheißen, er wolle Jacobs Gott sein, da er droben, 1 Mos. 35, 11. 12., zu ihm sagt: „Ich bin der allmächtige Gott, sei fruchtbar und mehre dich; Völker und Völkerhaufen sollen von dir kommen, und Könige sollen aus deinen Lenden kommen; und das Land, das ich Abraham und Isaak gegeben habe, will ich dir geben, und will des deinem Samen nach dir geben.“ Heißt aber nun das segnen und mehren, da er nicht allein den Vater verläßt, sondern auch den Sohn in die Knechtschaft und in das Gefängniß wirft und dasselbe unter gottlosen abgöttischen Leuten?

106. Wie könnte aber da ein Mensch in seinem Herzen anders denken oder gesinnet sein, denn daß Gott zürnen müßte und daß er seine Verheißung widerrufen wollte? Zwar unser Fleisch hält es eigentlich dafür und redet auch also davon; aber der Geist überwindet, und schließet also: Was unser Herr Gott mit seinem Rath meint, verstehe ich nicht, begehre es auch nicht zu wissen; sondern ich will die Hand des Herrn dulden, und sagen: Du bist mein Gott, deine Verheißung und dein Wort bleibt ewiglich. Damit will ich mich trösten wider das greuliche Aergerniß.

107. Auf diese Weise wird ohne Zweifel auch Joseph nach dem Fleisch auf das heftigste bewegt worden sein, darum daß er so verkauft worden und die große Gefahr seines Leibes und guten Namens tragen müssen. Er hätte die Verheißungen und Predigten von seinem Vater Jakob gehört, aber in so viel großem und mancherlei Jammer und Elend hat er gar das Gegentheil erfahren. Darum hat sich sein Fleisch entsetzen und klagen müssen, weil er so erbärmlicher Weise verlassen worden ist und die große Last des vielen Unfalls tragen müssen.

108. Das Volk Israel aber that in der Wüste viel anders, von welchem im 78. Psalm 3. 18. 19. gesagt wird: „Sie versuchten Gott in ihrem Herzen, daß sie Speise forderten für ihre Seelen, und redeten wider Gott, und sprachen: Ja, Gott sollte wohl können einen Tisch bereiten in der Wüste!“ Die murrten und zürnten mit Gott, und schrieen, sie müßten Durstes und Hungers sterben; könnten den Rath des Herrn nicht ertragen.

109. Wider solch Murren und Zorn rüstet und stärkt uns die heilige Schrift, und ermahnt uns, daß wir lernen und uns gewöhnen sollen an solche wunderbare Regierung Gottes. Denn dasselbe zeigt sein Name an, nämlich, daß er wunderbar sei und allein Wunderwerke thue im Himmel und auf Erden. Wie wir zu unsern Zeiten auch in großer Gefahr stehen müssen um des grimmigen Zorns willen des Pabsts, der Bischöfe und ihres Anhangs, welche bald hernach wider uns erschrecklich toben und uns hart zusetzen werden. Und ist daran kein Zweifel, wir werden unsers Theils noch mancherlei Unglück erwarten müssen, so uns begegnen wird.

110. Was sollen wir aber thun, wenn nun die Verfolgung daher gehet? Sollen wir dann sagen, daß unsere Lehre und Religion falsch und unrecht sei, darum daß die Könige und Fürsten dieselbe verdammen und greulich verfolgen? Mit nichten, sondern wir sollen also schließen: Der Herr lebet noch; ich bin getauft, ich habe Gottes Wort, wie toll und unsinnig auch die Welt über mich ist und mir allein Gut und Leib und dazu das Leben mit Gewalt nimmt. Und wenn schon die Welt darüber brechen und gar zu Trümmern gehen sollte, dazu Alles mit Blut und Mord erfüllet werden sollte: was kümmert es mich? Aber das Fleisch kann solche Beständigkeit nicht haben. Und wenn wir doch diesen Glauben fest behalten könnten, alsdann würden wir in Wahrheit empfinden, daß Gott, an den wir glauben und deß Wort wir haben, allmächtig sei.

111. Was er damit meint, daß er den Türken und Pabst also wider die Christen toben läßt, weiß ich freilich nicht. Sein Angesicht ist sehr erschrecklich, wenn der Türke, der Pabst und die Tyrannen mit dem Schwert und dem Tode wüthen, die Ketzer aber mit ihren Irrthümern, Ketzereien und Aergernissen auch einreißen: wer will uns da das Angesicht Gottes zeigen, welches wir ja gerne schauen wollten? Aber da wird nichts aus; Du wirst mir hinten nachsehen, sagt der Herr. Hänge dich an mein Wort und glaube demselben, und laß mir befohlen sein, daß ich Alles regiere und erhalte. Laß mich machen, du aber sollst ungezweifelt eher Weib und Kind, und alles, was du hast, in Gefahr setzen, ehe du mein Wort wolltest fahren lassen und verleugnen. Siehst du nicht, Petrus, was ich jetzt thue, so wirst du es doch hernach erfahren.

\quellseite{437}

112. Es ist ein groß wunderbärlich Ding, daß in Egypten die Kinder und Nachkommen Jakobs unter Joseph, der im Lande Herr war, so wohl gestanden, und in so gutem Frieden und stillem geruhigen Leben, dazu auch in solcher Ehre und Herrlichkeit gelebet haben. Aber da Joseph gestorben, ist an seine Statt ein anderer Herr und Tyrann gekommen, der weiß nichts von Joseph; da werden die unmündigen Kinder hart geplagt und getödet, und wird das Volk Israel mit harter Dienstbarkeit beschweret.

113. Warum läßt aber Gott dasselbe geschehen? Wie kommt es, daß er die Egypter der so großen herrlichen Gaben und Wohlthat vergessen läßt, damit sie Joseph überschüttet hatte, und daß dem Volke Israel Böses für Gutes vergolten wird? Halte stille. Es scheint zwar, als ob dieselbe Undankbarkeit eine gar unbillige und unleidliche Sünde sei; aber bald hernach kommt dann Gott wieder durch seine Ehre und Herrlichkeit und durch viele größere und höhere Wunderwerke, und führt das Volk Israel durch Mose aus Egypten heraus mit herrlicher Erlösung, läßt den König mit Roß und Wagen und seinem ganzen Heere im Rothen Meere ersaufen.

114. Also auch nach der Zerstörung, so durch die Babylonier geschehen, ist das Volk Israel nicht wieder zu seiner vorigen Herrlichkeit gekommen, welche sie zu den Zeiten Davids und Salomos gehabt haben, nachdem das Volk mit ihren Propheten und Königen sehr geschwächt und beinahe vertilgt gewesen, und von der Zeit an mancherlei Jammer und Angst erlitten. Aber nach fünfhundert Jahren, da kam Christus, der Welt Heiland.

115. Das ist nun der einige Trost in der höchsten Angst und Noth, daß wir alle unsere Sorge auf den Herrn werfen und uns seinem Willen und Wohlgefallen ergeben, er versuche uns gleich ingemein, oder aber einen Jeden insonderheit. Laß nur gehen; danke Gott, daß du das Wort und seine Verheißung hast. Da merke auf, daß du mit großem starkem Muthe alle Gewalt des Satans, des Todes und der Tyrannen ertragen und überwinden mögest.

116. Wie der heilige Petrus solchen Fleiß und Glauben an das Wort sehr lobt, da er 2. Epist. 1, 19. sagt: „Wir haben ein festes prophetisches Wort, und ihr thut wohl, daß ihr drauf achtet, als auf ein Licht, das da scheinet in einem dunkeln Orte, bis der Tag anbreche, und der Morgenstern aufgehe in euren Herzen.“ Er heißt uns unsere Augen und Verstand nur auf das Wort richten, und auf die Taufe, auf das Abendmahl des Herrn und die Absolution, und daß wir das Andere alles nur für eitel Finsterniß halten. Ich verstehe nicht, achte auch dessen nicht, was die Menschenkinder in dieser Welt ausrichten; denn sie freveln und plagen mich. Derselben erschrecklichen Larve, so das Angesicht Gottes vor mir verbirgt, kann ich nicht entfliehen, kann sie auch nicht abziehen; sondern ich muß in Finsterniß und in der größten Dunkelheit hangen bleiben, bis mir ein neu Licht scheint.

117. Also haben Jakob und Joseph gar eine traurige Gestalt vor ihren Augen und Herzen gehabt. Aber wie mit großer Freude hat Gott sie seinen Rücken sehen lassen? Also ist ihrer Nachkommen Betrübniß, Jammer und Noth in Egypten auch sehr erbärmlich gewesen: aber der Ausgang dieser Versuchung ist sehr herrlich gewesen. Denn es ist gar fein tröstlich geredet, das St. Paulus sagt 1 Cor. 10, 13.: „Gott ist getreu, der euch nicht lässet versuchen über euer Vermögen, sondern machet, daß die Versuchung so ein Ende gewinne, daß ihrs könnet ertragen“, daß man fröhlich wird.

118. Darum sehen wir, wie lieblich und so voller Trostes die Antwort Josephs ist, welcher so gar der unbilligen Mißhandlung vergißt, so seine Brüder an ihm geübt hatten, daß er nur mit kurzen Worten des Verkaufens gedenkt, und jetzt nun seinen Geist gar herauschüttet und dieselbe Sünde mit der allerschönsten Deutung auslegt. Ich zürne euch nicht, sagt er, habe auch daran keine Lust, daß ihr erschrocken, zu Schanden gemacht und geschmäht werdet. Und warum sollte ich mit euch zürnen, dieweil ich weiß, daß dies alles aus göttlichem Rath geschehen ist? Wie er denn hernach, 1 Mos. 50, 19. 20., zu ihnen sagen wird: „Fürchtet euch nicht, denn ich bin unter Gott. Ihr gedachtet es böse mit mir zu machen; aber Gott gedachte es gut zu machen“ etc. Und an diesem Orte sagt er auch: „Um eures Lebens willen hat mich Gott vor euch her gesandt“, auf daß ihr möchtet erhalten werden, das sie wahrlich nicht verdient hatten. Und man sieht an David, daß er sonderliche Freude an dem Worte „senden“ gehabt; da er im 105. Psalm V. 17. sagt: „Er sandte einen Mann vor ihnen hin; Joseph ward zum Knechte verkauft.“ Gott wird der Sender genannt, und dasselbe Andern zum Trost und Heil, da er den frommen Jakob Joseph hat verkauft werden lassen und getödtet werden, und daß der Vater darüber ganz jämmerlich erschrocken und betrübt worden ist. Dies ist weder eine griechische, noch lateinische, noch hebräische Weise zu reden. Die Art und Weise aber zu senden ist wunderbar und unerhört.

119. Derhalben sehen wir, wie gar freundlich und lieblich der heilige Geist ist, der diese erschreckliche That, daß Joseph also von seinem Vater ist weggerissen worden, und das sehr schwere Kreuz, dadurch sie schier beide getödtet worden, eine Sendung zum Leben nennt. Und es ist wahrlich eine rechte Benennung, welche auch Gott sehr angenehm ist. Darum soll sie uns auch wohlgefallen; denn sie ist ein Wort und Lehre des Heils. Denn also hat Joseph lernen reden, da er von den Todten auferstanden war, nachdem er den Rath des Herrn und Gott hintennach gesehen hat, nämlich, was er mit der sehr schweren Anfechtung gemeint hätte, und daß er bei den Seinen sei in ihren höchsten Nöthen, wenn sie gehabt, von ihren Brüdern getödtet, verkauft, ins Gefängniß gestoßen und geschmäht werden und sonst in allerlei Unglück gerathen.

120. Und dies soll man der Kirche und Gemeinde Gottes vorhalten, auf daß die Gottseligen durch solche Exempel erweckt werden mögen zum Glauben und zur Geduld, und glauben lernen an Gott den Vater, allmächtigen Schöpfer Himmels und der Erden etc. Wir können bei Gott nichts verlieren, wenn wir glauben, sondern aller Verlust ist hundertfältiger Gewinn. Jakob hatte seinen Sohn verloren und er selbst war auch verloren, wie er selbst und alle Menschen davon nicht anders haben urtheilen können; aber Gott antwortet, und sagt von Joseph: „Er ist nicht verloren noch verdorben, sondern ich brauche sein als eines Gesandten in Egypten, auf daß durch ihn nicht allein seinem Vater, Brüdern und Hausgenossen, sondern auch demselben ganzen Königreiche und allen umliegenden Ländern gerathen und geholfen werden soll. Und dasselbe also, daß solches nicht allein eine leibliche, sondern auch eine geistliche Hülfe sein soll, nämlich, daß er den König, die Fürsten und das Volk im Lande unterweise im Glauben und Erkenntniß Gottes. Das ist wahrlich eine sehr treffliche und große Gesandtschaft. Denn also deutet nun Joseph selbst das Verkaufen. Als wollte er sagen: Ihr habt mich zwar verkauft; aber Gott hat mich gekauft, hat mich erlöst und sich zum Hirten zugeeignet, daß ich ein Fürst und Heiland der Völker sein sollte, eben in demselben Rath, nach welchem ich verloren und verdorben schien.

121. So lehrt uns nun dies Exempel, daß wir auch dadurch nicht verloren oder verderbet, sondern erhalten werden, wenn wir betrübt und geplagt werden, sofern wir nur recht glauben. Denn der Glaube muß da sein; der Unglaube denkt an die Dinge nicht. Denn dies sind geistliche Exempel des allervollkommensten Glaubens, welche weder Juden noch Papisten verstehen, deren sie auch nicht groß achten. Wir aber sollen dies fleißig betrachten und in unsere Herzen prägen, nämlich, daß dies eine Sendung sei zum Heil und Trost vieler Leute, da Joseph nach Egypten verkauft worden ist.

122. Die heilige Schrift redet eben auf solche Weise auch vom Tode und Begräbniß Christi. Denn also sagt Jesaia im 11. Cap. V. 10.: „Und seine Ruhe wird Ehre sein“, das ist, sein Begräbniß, welches bei allen Menschen für die äußerste Schande und Armseligkeit gehalten worden, ist vor Gott eine Ehre. Denn also sagt er selbst auch im 16. Psalm V. 10. 11.: „Du wirst nicht zugeben, daß dein Heiliger verwese. Du thust mir kund den Weg zum Leben“ etc.

\quellseite{438}

123. Derhalben soll man die Kraft und Güte Gottes unserer Schwachheit, Armuth und endlich auch dem Tode selbst entgegen halten. Wie die Deutschen ein fein Sprichwort haben, da sie sagen: Gott hat noch mehr, denn er je vergab. Damit haben sie anzeigen wollen, wie groß und unerschöpflich die Güte Gottes sei. Denn wenn wir sterben, so werden wir ja nicht verloren, sondern leben, wenn wir seiner Verheißung glauben, und wenn er uns schon das Leben, Gut, Weib und Kind durch den Pabst oder Türken mit Gewalt nehmen läßt. Was geschieht denn? Dies gegenwärtige Leben wird uns genommen, auf daß uns ein anderes und viel besseres wiederum gegeben werde. Denn Gott kann und wird auch ohne Zweifel noch mehr geben, denn er je gegeben hat.

124. Und ist sonderlich jetzt zur Zeit diese Lehre sehr nöthig, da wir sehen, daß der Pabst mit seinem Anhange so greulich tobt und wüthet, und daß viel fromme und heilige Männer mit Weib und Kind in das Elend getrieben, und aller ihrer Habe und Gutes beraubt und dazu greulich getödtet werden. Desgleichen wüthet der Türke auch ganz greulich. Darum sind wir gleichsam mitten im Tode, und die Welt droht, als wollte sie gar über einen Haufen fallen. Aber der Herr lebet noch. Es ist Alles des Schöpfers, der da allmächtiger Vater ist. Darum sollen wir nicht ansehen, wie groß das Unglück sei, und wie mächtig die sind, so uns verfolgen. Es soll uns auch unsere eigene, oder unserer Weiber und Kinder Gefahr nicht so gar erschrecken; wie Joseph nicht auf das Gefängniß, Schmach, Tod, und was ihm sonst mehr für Traurigkeit und erschrecklicher Unfall begegnet, gesehen hat; sondern bleibt nur immer am Worte hangen und glaubt an Gott, den allmächtigen Vater. Wenn nun der Türke oder Pabst uns schon lebendig verschlingen würden, lieber, was liegt daran? Sie verschlingen unsern Tod, Jammer und Elend, und machen einen Wechsel, nämlich, daß wir für dies elende Leben ein ewiges bekommen, so wir anders an Gott, den allmächtigen Vater, glauben.

125. Ja, sagt du, inzwischen wird gleichwohl unsere Religion zerstört? Antwort: Was liegt daran? Laß Gott walten und dafür sorgen, der Herr wird sich das wohl lassen angelegen und befohlen sein, daß seine Kirche erhalten werde und daß noch immer auf Erden ein heiliger Same bleiben möge. Es ward wahrlich Joseph zu den Egyptern, einem großen und gottlosen Volk, gesandt, welches von Gott nichts gewußt, da große Gefahr gewesen, daß etwa das junge zarte Herz mit falscher Lehre und Aberlauffen verchristet und verführt werden möchte: aber der Ausgang und Ende zeigt an, daß eben dieselbige Sendung und Gefahr eine Ursache sehr vieles Guten gewesen sei.

126. Also ist es Gott noch möglich, daß er Etliche unter uns auch erhalte, die der Same sein sollen, dadurch die Kirche wachsen und zunehmen möge, wenn auch der Türke und alle Feinde Gottes gar hinunter und zu Trümmern gegangen sind. Denn der läßt sie um seiner andern Ursache willen also toben und wüthen, denn daß sie selbst dadurch zerstört und verderbet werden sollen. Unterdes wird aber das Licht des Evangeliums immer mehr und mehr ausgebreitet. Lasset uns nur fortfahren mit Glauben, Lehren, Leiden und Sterben; denn die Feinde und Widersacher werden selbst auch umkommen müssen.

127. Ja, sagst du, nichts desto weniger wirst du darüber erwürgt. Was liegt daran? Ich weiß doch wohl, daß ich von diesem Leben abscheiden muß. Darum frage ich darnach nichts, ich werde gleich vom Pabste oder Türken getödtet. Ich glaube an Jesum Christum, denselben bekenne ich auch und rufe ihn an, es zürne oder lache, wer da will.

128. Auf solche Weise haben die heiligsten Väter alle ihr Vertrauen und Hoffnung auf Christum gesetzt, und haben mit dem Glauben den Tod, die Hölle und den Satan mit großem starkem Muthe überwunden, und also die ewige Herrlichkeit erlangt, welche ihnen und uns auch im Himmel beigelegt ist. Wenn Joseph in seines Vaters Hause und bei seinem Hausgesinde geblieben wäre, so hätte er daselbst als eine Privatperson oder gemeiner Mann sein Leben zugebracht und wäre nur wenig Leuten nütze gewesen. Da aber das Weizenkorn in die Erde gefallen, da wird es ein Heil vieler Länder, beide leiblich und geistlich.

129. Und was sollen wir doch meinen, daß es mit den Türken und Pabst noch werden wolle? Denn das Böse, so sie gesündet und vorhaben, kann Gott leichtlich zum Guten wenden. Es steht an einem Ort ein schöner, herrlicher und heiliger Spruch Augustins, welcher also lautet: Tam bonus est Deus, ut nullo modo permittat malum, nisi inde norit elicere bonum, das ist: Gott ist so gut, daß er kein Böses geschehen läßt, wenn er daraus nicht etwas Gutes nehmen kann. Denn also hat er auch aus Josephs Gefängniß, Knechtschaft und Elend dem ganzen Egyptenlande Heil, Leben und große Herrlichkeit gegeben.

130. Also schickt er uns auch nach seinem wunderbaren Rath Gefahr Leibes und Gutes, dazu auch unsers guten Namens zu. Aber der Glaube im Herzen bricht hindurch und steht wiederum auf aus dem Tode in das Leben.

131. Dies soll man merken von der Art und Weise zu reden, damit Joseph den wunderbaren Rath Gottes deutet und erklärt, da er sagt: „Gott hat mich vor euch her gesandt“, euch dadurch am Leben zu erhalten, sonst hättet ihr alle umkommen müssen. Ihr habt aber bis daher unsers Herrn Gottes Rath nicht sehen oder verstehen können. Wie Christus Joh. 13, 7. zu Petro sagt: „Du wirsts aber hernach erfahren.“ Und von den Jüngern wird auch gesagt: „Sie verstunden das Wort nicht; denn Jesus war noch nicht von den Todten auferstanden.“ Und Gott thut solches darum, daß er damit unsern Glauben also üben und prüfen will. Denn wo er uns sein Angesicht von vorne her zeigen würde, so hätte der Glaube keine Statt. Darum können wir nun aus dem, so hernach folgt, (a posteriori) schließen, daß uns so viel Leides nicht widerfahren oder gethan werden kann, daß uns Gott dagegen nicht viel mehr und größere Güter geben und uns damit gleichsam überschütten könne.

132. Darum ist dies nun der Trost der Brüder Josephs, damit er ihnen nicht allein Schuld und Pein vergibt, sondern heißt sie auch getrost und gutes Muthes sein, darum daß er nach Gottes Willen und Befehl ihnen als Verräthern gegeben sei, daß er sie am Leben erhalten und zu Ehren setzen soll, und daß also Josephs Tod in das Leben verändert worden sei. Denn ein solcher allmächtiger Schöpfer ist unser Gott, der aus Nichts Alles macht, der den Tod, die Hölle und alles Unglück in seiner Hand hat, und geht nur damit um, daß den Auserwählten alles muß zum Besten gerathen, was ihnen nur begegnet.

\begin{quote}
\textbf{V. 6. Denn dies sind zwei Jahr, daß es theuer im Lande ist; und sind noch fünf Jahr, daß kein Pflügen noch Ernten sein wird.}
\end{quote}

133. An diesem Ort disputiren sie viel, was die Ursachen dieser Theurung mögen gewesen sein, und nehmen dieselben daher, daß der Nil ausgetreten sei, dadurch Egyptenland gefeuchtet worden, da sonst der Himmel keinen Regen gegeben hat. Denn, wie sie schreiben, daß Babylon fruchtbar sei, wenn die Wasser Euphrat und Tigris austreten: so sagt man auch vom Nil, daß derselbe mitten im Sommer, um die Zeit der Sonnenwende, und wenn die Sonne durch die drei Zeichen des Krebses, des Löwen und der Jungfrau geht, auch austrete; und wenn er weniger als zwölf oder mehr als sechzehn Ellen wächst, soll es Theurung bedeuten, wenn es aber dreizehn oder vierzehn Ellen sind, sagen sie, daß sei eine Ursache der größten Fruchtbarkeit und einer guten wohlfeilen Zeit, darum daß die kleinen Wasser nicht alles Land feucht machen können, die großen aber, indem sie zu langsam ablaufen, machen, daß der Acker nicht gebauet werden kann. Wie in unsern Landen unfruchtbare Jahre bisweilen von großer Dürre kommen, bisweilen auch, wenn es allzu feucht ist. Und wir müssen schier alle Jahre dieser Schaden einen tragen. Jedoch hat es mehr Gefahr auf sich, von diesem als von jenem Schaden leiden.

134. Diese Verkündigung aber der Theurung in Egypten ist prophetisch gewesen, nämlich, daß es denen, die den Acker bauen, nicht gerathen würde, welcherlei Ursache es auch haben möchte; welches er prophetisch also ausdrückt, daß er sagt, daß kein Pflügen oder keine Ernte sein werde; damit er so viel meint, daß kein Pflügen helfen und die Ernte sehr gering sein werde. Denn dieselbe Plage ist nicht allein den Egyptern, sondern auch den andern umliegenden Ländern auferlegt worden; und solches ist geschehen aus sonderlichem Rath und Zorn Gottes.

\begin{quote}
\textbf{V. 7. Aber Gott hat mich vor euch her gesandt, daß er euch übrig behalte auf Erden, und euer Leben errette durch eine große Errettung.}
\end{quote}

\quellseite{439}

135. Joseph wiederholt das Wort: „Gott hat mich gesandt“ \&c., und braucht sehr viele Worte, diesen wunderbaren Rath Gottes groß zu machen, auf daß er sie in ihren Herzen zufrieden stellen, und dessen vergewissern möge, daß ihnen ihre Sünde vergeben, und Gott ihnen gnädig und barmherzig sein wolle. Denn durch Gottes Barmherzigkeit ist ihr Verkauf, das voll Neides und Hasses gewesen, in eine heilsame Sendung verändert worden; und scheint fast in einem jeglichen Worte die größte Göttlichkeit und Ehre, so er Gott erzeigt hat. Denn er kann sich nicht genugsam verwundern, auch nicht hoch genug rühmen, wie groß die Ehre sei, daß Gott aus solcher großen Sünde, da man schwerlich eine erdenken möchte, die größer und greulicher sein könnte, so viel großes Gutes, Leben und Heil hat erwecken und hervorbringen können. Denn das ist ja die Art und Natur der Sünde nicht, daß daraus was Gutes kommen sollte; aber dies ist die Wirkung göttlicher Gnade und Güte, die das Böse zum Guten wenden kann; wie Joseph hernach im letzten Capitel V. 20. sagen wird.

136. Darum ist dies eine Lehre allein für die Christen und Kinder Gottes, die da können über sich gehen lassen und in sich fressen, was ihnen nur für Unglück zukommt, und dieselb mit rechten Vertrauen darauf harren, daß sie für einen Schaden oder Unfall tausenderlei Nutzen wiederum haben werden. Denn wir glauben an den Gott, der da ein Allmächtiger Schöpfer ist, welcher aus Nichts kann Alles machen, und aus dem, das Böse ist, Gutes, und aus dem, das gar verdorben und verloren ist, Heil und Seligkeit wirken. Wie Röm. 4, 17. ihm solches zugeschrieben wird, da St. Paulus sagt: „Er rufet dem, das nicht ist, daß es sei“; und 2 Cor. 4, 6.: „Gott, der da hieß das Licht aus der Finsterniß hervor leuchten.“ Nicht aus der Kohle ein Fünklein, sondern „aus der Finsterniß das Licht“, sagt er; wie er auch aus dem Tode das Leben, aus der Sünde Gerechtigkeit, aus der Dienstbarkeit des Teufels und der Hölle das Himmelreich und Freiheit der Kinder Gottes machen und hervorbringen kann.

137. So bleibt nun Joseph noch immer bei diesem wunderbaren Rath Gottes, und bringt nur auf denselben und rühmt ihn sehr hoch. Zuvor hat er gesagt: Ihr sollt darüber nicht traurig sein, daß ihr mich verkauft habt; denn Gott hat mich vor euch her gesandt um eures Lebens willen insgemein: hier aber wiederholt er es nun und schärft es, und zieht es auf sie, da er sagt: „Gott hat mich vor euch her gesandt, auf daß er euch übrig behalte.“ Als wollte er sagen: Die gemeine Hülfe, welche Gottes Güte und Gnade gewirkt hat nach seinem unaussprechlichen Rath für alle Völker, gehört auch für euch. So gar hat eure Sünde nicht geschadet, wiewohl sie groß und erschrecklich war, daß sie auch allen Völkern umher und vornehmlich euch Hülfe und Trost gebracht hat. Als ob er sagen wollte: Gott hätte euch auf anderem Wege wohl erhalten können; aber dies ist bei ihm also beschlossen und bestimmt gewesen, daß ihr durch mich habt sollen erhalten werden; sonst wäret ihr Hungers und wegen eures bösen Gewissens gestorben, und unser sehr heiliger Vater wäre von Traurigkeit des Herzens aufgerieben. Derhalben sollt ihr an Gottes Gnade und Barmherzigkeit nicht zweifeln, der euch eure Sünde vergeben und verziehen hat.

138. Und sollt wissen, daß ich nicht dazu gesetzt oder bestellt bin, die Sünde an euch zu rächen oder zu strafen; denn das sucht Gott nicht, will es auch nicht haben; sondern dieweil ihr Buße gethan habt, sind euch eure Sünden alle vergeben, auch so, daß ihr nun hinfort deswegen sicher und zufrieden, dazu auch ohne alle Furcht der Strafen sein mögt, und euch nur gewisser Hülfe versehen sollt.

139. Auf diese Weise schärft Joseph den göttlichen unaussprechlichen Rath: und wiederum lindert er die Sünde und macht sie gering, als ob es ein ehrlich und löblich Ding wäre; und dasselbe zwar nicht aus ihrer Kraft und Natur, sondern vielmehr aus Gottes Güte und Barmherzigkeit, dadurch sie also gering gemacht und weggenommen wird. Und Paulus sagt Röm. 5, 20.: „Das Gesetz ist neben einkommen, auf daß die Sünde mächtiger würde. Wo aber die Sünde mächtig worden ist, da ist doch die Gnade viel mächtiger worden“; und Röm. 3, 7. 8.: „So die Wahrheit Gottes durch meine Lüge herrlicher wird, zu seinem Preise“ \&c., „so lasset uns Uebel thun, auf daß Gutes daraus komme.“

140. Aber von denen, die solcher überschwänglichen Gnade Gottes mißbrauchen, sagt Paulus Röm. 3, 8. selbst auch: „Welcher Verdammniß ist ganz recht. Denn Gott hat ja nicht geboten oder verheißen, daß er aus deiner Sünde wunderbare, herrliche und heilsame Werke hervorbringen wolle; oder wenn er sich deren, so Sünde und Schande getrieben haben, erbarmt, auf daß er damit seine große Gnade und Barmherzigkeit beweisen möge, so ist dasselbe ein sonderlicher göttlicher Rath, darauf solches kommt, und geschieht nur den armen Sündern zu Trost, die gefallen sind, auf daß dieselben nicht verzagen, sondern sich zu Gottes Barmherzigkeit wenden sollen, und daß sie ja nicht in andere größere und schwerere Laster, gegenwärtig oder zukünftig, fallen. Denn Gott kennt die große Kraft und Tyrannei der Sünde wohl, welche den Menschen bald tödtet, wenn sie wiederum lebendig wird. Also hat sich Saul selbst umgebracht, Judas hat sich selbst an den Strick gehängt, da sie beide vom Schrecken der Sünde überwunden worden.

141. Darum muß man demselben Unglück zuvorkommen, und die armen verzagten Herzen, die gar darnieder geschlagen sind durch die Kraft des Gesetzes und der Sünden, wiederum aufrichten, wenn sie sagen: „Ich habe leider! gesündigt, das bekenne ich; Gott aber ist gnädig und barmherzig, zu dem will ich fliehen, und bitten, daß er mir es aus Gnaden vergeben wolle. Wenn aber Einer aus Vermessenheit, absichtlich und wissentlich sündigt, dazu Gottes Ordnung und seinen schweren Zorn verachtet, derselbe sündigt tödtlich, und ist nicht weit von der Sünde wider den Heiligen Geist.

142. Derhalben soll sich jedermann wohl vorsehen, und in seinem Herzen fleißig erwägen, was Sirach sagt Cap. 5, 2—7.: „Folge deinem Muthwillen nicht, ob du es gleich vermagst, und thu nicht, was dich gelüstet, und denke nicht: Wer will mirs wehren? Denn der Herr, der oberste Rächer, wirds rächen. Denke nicht: Ich habe wohl mehr gesündiget, und ist mir nichts Böses widerfahren; denn der Herr ist wohl geduldig, aber er wird dich nicht ungestraft lassen. Und sei nicht so sicher, ob deine Sünde noch nicht gestraft ist, daß du darum für und für sündigen wolltest. Denke auch nicht: Gott ist sehr barmherzig, er wird mich nicht strafen, ich sündige, wie viel ich will. Er kann bald also zornig werden, als gnädig er ist, und sein Zorn über die Gottlosen hat kein Aufhören.“ Der Zorn und das Gericht Gottes werden durch das Gesetz angezeigt, daß wir nicht sündigen sollen; wenn wir aber gefallen sind, so haben wir das Wort göttlicher Verheißung, darin geboten wird, daß wir nicht verzagen sollen, sondern bedenken, daß Gottes Gnade und Barmherzigkeit unermesslich groß sei, der seinen Sohn für uns gegeben hat, daß er unsere Genugthuung sein soll, und durch sein Blut unsere Sünden gar verschlungen und getilgt werden sollen.

143. Man soll und muß Gott nicht versuchen, daß man auf seine Barmherzigkeit vermessen sein wollte; sondern man soll im Glauben und rechter herzlicher Demuth und Buße um Vergebung der Sünden bitten und darauf hoffen. Wie Joseph, da er die wunderbare Regierung Gottes rühmt, nicht angezeigt hat das Verdienst, die Sünde oder den Trotz seiner Brüder, sondern nur Gottes Gnade und Barmherzigkeit, der nicht haben will, daß die armen Sünder verloren, sondern daß sie alle selig werden sollen.

144. Und er sagt, er sei darum gesandt worden, daß wenn schon die Andern alle umkommen, er gleichwohl noch Etliche, so überbleiben, in der Kirche und Gemeinde Gottes erhalten möge aus dem Hause und Geschlechte Jacobs, welcher ein Pfarrherr und Bischof gewesen, und die Verheißung, so er von Gott empfangen, bewahret hat. Und zeigt also an, daß Gott dieses kleine Häuflein erhalten wolle, und um desselben willen Regiment, Polizei und Friede geben, und dazu Hausregiment, Regeln, fruchtbare Zeit, wie jämmerlich auch Jakob versucht und ungetrieben worden ist. Denn wo es ohne dieselbe Kirche gewesen wäre, so wäre die Welt tausendmal bald vergangen. Darum, sagt Joseph, wird das Land Canaan, Egypten und andere Länder mehr um euretwillen erhalten, auf daß also noch immer ein Häuflein überbleiben möge von Gottes Volk, durch welches darnach viele Andere auch mögen bekehrt werden.

145. Denn wir haben droben gehört, daß die heiligen Patriarchen zu derselben Zeit wie große Lichter in der Welt gewesen sind, daß sie nicht allein ihnen mit Lehren und daß sie das Wort ausgebreitet, gedient, sondern es haben sich auch viele Andere mehr zu ihnen gehalten, die nicht selig geworden sind, wie Etliche unrecht und ungöttlich sagen wollen, in ihrer heidnischen Religion und Aberglauben, sondern darauf kommen, daß sie Jakob und Joseph gehöret, welche sie Gottes Wort gelehrt haben. Solche Leute heißen in der Schrift, Judengenossen, Proselyten, oder Fremdlinge, die nicht von den Linien und Geschlechtern der Patriarchen, sondern Heiden gewesen; und sind ihnen doch Freundschaft oder Schwägerschaft halben verwandt gewesen, als, Mägde und Knechte, so sich zu der Kirche der Väter haben versammeln lassen, daß sie derselben auch mit einverleibt worden.

\quellseite{440}

146. Also hat Joseph dem ganzen Königreiche Egypten sehr viel genützt, da er den König, die Priester und Fürsten im Lande bekehrt hat. Darum sagt er: „Gott hat mich vor euch her gesandt, daß er euch übrig behalte auf Erden.“ Ihr seid aber die, so übriggeblieben sind, und das kleine Häuflein, um euretwillen bin ich gesandt worden.

147. Wie zu unsern Zeiten Kaiser, Papst und Mahometisten um keiner andern Ursache willen leben und noch in der Welt bleiben, denn allein um der Kirche, und des Wortes und der Sacramente willen. Wenn dieselben aufgehoben und weggenommen sind, so werden sie bald gar vergehen; wie Daniel zeuget Cap. 12, V. 7.: „Wenn die Zerstreuung des heiligen Volks ein Ende hat, soll solches alles geschehen.“ Inzwischen aber bleiben noch die Polizeien und Regimente, es bleibt die Zeugung und Erziehung der Kinder. Denn die Zahl der Auserwählten ist noch nicht erfüllt, deren viele auf dieser Welt noch nicht geboren sind. Wenn aber nun die Zahl erfüllt ist, so soll Papst und Türke wohl sehen, wo sie bleiben.

148. Die Kirche aber allein erhält alle Dinge, und ist wie ein Ort, da man pflanzt oder säet in der Welt, dadurch das Wort und die Sacramente ausgebreitet und erhalten werden. Zu derselben Kirche treten täglich etliche Proselyten, die auch die rechte himmlische Lehre hören und annehmen. Unterdeß toben auch noch wider uns, als wären sie unsinnig, der Papst, Kaiser, Türke, und streiten wider diejenigen, welche sie doch wie starke Säulen erhalten. Und unser Herr GOtt würde sie nicht einer Bohne werth achten; sondern allein um der Kirche willen setzt er ihnen die Krone auf, gibt ihnen Geld und Gut, Gewalt und Ehre mit Haufen. Dafür danken sie darnach der Kirche Gottes also, daß sie die noch übrigen armen Christen jämmerlich verfolgen, und wollten dieselben gern ganz und gar verderben und ausrotten, wenn sie nur könnten.

149. Es setzt aber der heilige Patriarch noch mehr dazu, auf daß er ja das wunderbare Werk Gottes groß mache. Ich bin gesandt, sagt er, daß insgemein alle Völker umher durch mich sollen erhalten werden; aber diese Wohltat geht vornehmlich euch an: es wird damit am allermeisten euer Heil und Wohlfahrt gesucht und um euretwillen müssen es auch die Feinde gut haben. Ich bin hieher geführt worden um eures Lebens willen, auf daß ihr bei dem natürlichen Leben möchtet erhalten werden; welches denn geschieht durch eine wunderbare Errettung. Denn das ist eine große unermessliche Gabe Gottes, nämlich, daß euch Gott errettet von der zukünftigen Theurung, so fünf ganze Jahre währen wird; dadurch ganz Egypten und ihr Canaaner mit andern umliegenden Völkern auch verderben würdet, wo Gott nicht aus unermesslicher Barmherzigkeit diese Errettung durch große Wunder angerichtet hätte, dadurch so viele Völker um euretwillen erhalten werden.

150. Daß aber Joseph diese Errettung so groß macht, gehört dazu, daß der Glaube in ihren Herzen dadurch aufgerichtet und erweckt werden möge. Denn es wird es niemand leichtlich glauben, es sei denn, daß er es erfahren habe, welch ein greulich Unglück die Sünde sei, wenn sie geoffenbart und erkannt wird, die das Gewissen so gar darnieder schlägt, daß es sich nicht leichtlich wiederum kann trösten und aufrichten lassen. Wie einer von den Scholastikern über den Spruch Spr. 15, 15.: „Ein guter Muth ist ein täglich Wohlleben“, angezeigt und gesagt hat, daß keine größere Freude sei denn die Freude des Gewissens, und keine größere Lust denn im Geist oder Herzen. Denn die Lust, damit sich das Fleisch belustigt, ist mit nichten zu vergleichen mit der Lust, so im Herzen ist; ja, dieselbe Lust ist gemeiniglich allezeit mit Wehmuth und Galle vermischt. Aber „ein guter Muth“ und fröhlich Herz „ist ein täglich Wohlleben“, das sehr wohl zugerichtet ist. Aber dagegen ein Herz, das sich begangener Sünden schuldig weiß, dasselbe fühlt auch nur Schmerzen und eitel Traurigkeit; wie das an Juda und Saul zu sehen ist. Denn die Sünde, wo sie erkannt wird, ist nichts Anderes, denn der Höllenrachen, so weit offen steht.

151. Eben dasselbe ist auch an diesen Brüdern zu sehen, die sich nicht wiederum können aufrichten, wiewohl Joseph ihre Sünde lindert und dieselbe gar gering macht, und sie eine Sendung aus Gottes Rath verordnet nennt, und die zu ihren und anderer Leute Heil und Wohlfahrt geschehen sei. Sie stehen da noch bei ihm, sind furchtsam, zittern, sind erschrocken und haben sich gar entsetzt. Und er hat bei ihnen wenig ausgerichtet, ob er sie schon lang getröstet und geheilt hat, ob er wohl ihre Sünde gering gemacht und dieselbe zugedeckt hat, da er oftmals die göttliche Ordnung angeführt, und wie dies Ding so ein schön glückseliges Ende genommen, welches vielen Völkern so sehr nütze und heilsam sei, daß er beides sie selbst stärken und darnach auch den betrübten und erschrockenen Vater durch sie trösten möchte.

152. So viel Mühe und Arbeit hat es, wo man die armen Herzen, die sich vor Gottes schwerem Zorn und seinem Gericht fürchten, aufrichten und stärken soll. Ja, siehe die Apostel an, mit denen Christus vierzig Tage umgeht nach seiner Auferstehung, stärkt und lehrt sie vom Reich Gottes, und waren auch dazu gekommen die Erscheinungen und Predigten der Engel: und haben sich doch nicht können fest daran halten und damit recht zufrieden sein. So gar ein betrübt und jämmerlich Ding ist es um die Sünde und um ein böses Gewissen.

153. Darum sollen wir uns mit allem Fleiß und Furcht Gottes vor Sünden hüten, oder wo wir gesündigt haben, darüber sich das Herz bekümmert, da sollen alle Menschen mit einander rathen und helfen, daß ein solch bekümmert Herz von seiner Bekümmerniß und Schmerz errettet und erweckt, oder ermahnet werden möge zum rechten Glauben auf Gottes Gnade und Barmherzigkeit. Wie Joseph seinen Brüdern hier auch auf das fleißigste predigt, da er sagt: Ach, lieben Brüder, bekümmert euch in eurem Herzen nicht, seid doch getrost, und glaubt mir gewißlich, daß euch eure Sünde vergeben ist, die ja an sich selbst wohl groß ist, aber die Hülfe und Erlösung, beide leiblich und geistlich, ist noch viel größer. Diese trefflichen Predigten hören sie von ihm, und zweifeln doch noch, und verstehen sich nicht, ein Wort zu reden.

\begin{quote}
\textbf{V. 8. Und nun, ihr habt mich nicht hergesandt, sondern GOtt, der hat mich Pharao zum Vater gesetzt, und zum Herrn über all sein Haus, und einen Fürsten in ganz Egyptenland.}
\end{quote}

154. Joseph zieht zum öftern an und prägt es ihnen wohl ein, daß seine Brüder ihn verkauft und wie dasselbe geraten sei etc. Und sagt: Daß ihr mich verkauft habt, das hat mich an diesen Ort nicht gebracht; eure Sünde hat mir keinen Schaden gethan: sondern Gott hat dieselbe angenommen, und dieselbe gebraucht, den Satan damit zu schanden zu machen, der euch verführt hat, auf daß er aus einem so großen Uebel ein großes Gut hervorbringen möchte. Gott hat durch seine unermessliche Güte diese Gelegenheit, daß ihr mich also verkauft, ergriffen, und Pharao, den Egypten und andern Völkern mehr zu allem Guten, Wohlfahrt und Seligkeit geraten lassen.

155. Und also läßt sich Joseph das mit höchstem Fleiß angelegen sein, daß er seiner Brüder Evangelist und Gnadenprediger sei, der Gottes Barmherzigkeit vor ihnen rühmt und groß macht. Denn er sagt: Auf daß ihr möchtet erhalten werden, „hat mich Gott Pharao zum Vater gesetzt“, daß er mich höre und mir gehorche, wie ein Kind seinem Vater oder ein Schüler seinem Schulmeister gehorcht. Das ist aber eine große Ehre, nämlich, einen solchen großen König und so viel Fürsten im Lande haben, die seine Zuhörer sind. Dies aber, sagt er, habe weder ich noch ihr ausgerichtet, sondern Gott hat mich mit dieser Herrlichkeit beehrt. Es sind eitel göttliche Werke. Seid stille und erkennet, denn er thut allein Wunder. Und dies sind auch wahrlich größere Wunderwerke denn die Auferweckung der Todten.

156. Derhalben schreibt er Alles Gott allein zu. Denn was sollte ich thun? sagt er; ich war verkauft und zur Knechtschaft gebracht, war dazu verklagt und verdammt, als hätte ich einen Ehebruch begangen, ward darüber in das Gefängniß gelegt: daraus bin ich aber erlöset und nun zu dieser Ehre und Herrlichkeit erhoben worden, daß ich soll ein Vater, Meister, Lehrer und Regent sein aller Fürsten und Priester und dieses ganzen Königreichs. Ich sehe wohl, daß ihr sehr erschrocken seid und euch dessen verwundert, und dazu meinen Worten schwerlich Glauben schenken könnt, daß ich ein Vater und Fürst über dies ganze Land sein soll.

\quellseite{441}
Ich sage aber die Wahrheit, und dies alles ist am Tage, daß man es ja vor Augen sehen kann.

157. Und ist zwar kein Wunder, daß sie daran noch etwas zweifeln: ich hätte selbst auch das nicht leichtlich glauben können. Darum hat er sie dessen nicht ohne große Mühe und Arbeit bereden können, daß er Joseph, ihr Bruder, wäre; vornehmlich da sie gesehen, daß er ein Fürst über Egypten war, welchen der König und die andern großen Herren im selbigen Königreich in großen Ehren gehalten und nach dessen Rath und Willen sie Alles verwaltet haben. Darum macht er die Barmherzigkeit GOttes so groß und rühmt sie hoch, auf daß er damit bei ihnen so viel schaffen möge, daß sie doch seinen Worten Glauben schenken und es dafür halten, daß alles, was geschehen, nicht nach menschlichem Rath oder Willen geschehe, sondern wie es die göttliche Majestät selbst ordnet und regieret.

158. Und bis daher hat nun gewähret die Predigt und das Wort der Verheißung von der Gnade und Barmherzigkeit GOttes, von Vergebung der Sünden, und wie ihnen GOtt alle Wohlthaten erzeigt, auf daß sie dessen gewiß sein mögen, daß sie von der Theurung errettet und GOtt sie mit vielen Gaben gleichsam überschüttet habe, welches wir das Evangelium heißen. Da aber dieselbe Trostpredigt ein Ende hat, gibt er ihnen nun Befehl, darin er anzeigt, was er von ihnen gethan haben wolle.

159. Wie wir auch, da man das Evangelium lehrt, zu allererst hören, daß uns Vergebung der Sünden verkündigt wird. Denn wir werden gerecht durch den Glauben, aus GOttes Gnade und Barmherzigkeit, so er uns in seinem Sohn Christo erzeigt hat, wenn zuvor unsere Sünde erkannt und unser Gewissen derhalben erschrocken gewesen. Wie diese Brüder auch erschrocken sind und sich fürchten, da sie durch Erkenntniß ihrer Sünden gar darnieder geschlagen gewesen. Darum kommt nun bald das Evangelium hernach, auf daß sie nicht verzagen, wie Saul und Judas gethan; sondern daß sie glauben, da sie nun aufgerichtet, gerecht und getröstet worden sind, dazu Vergebung der Sünden und den Heiligen Geist empfangen haben. Dies ist ein schön trefflich Exempel eines auch sehr trefflichen Lehrers, Buße zu predigen, oder die Leute zur Erkenntniß der Sünden und zum Glauben an GOttes Gnade und Barmherzigkeit in Christo zu ermahnen. Nun folgen Worte, darin der Befehl steht, was sie thun sollen.

\section[Vierter Theil]{Vierter Theil. Wie Joseph seinen Brüdern Befehl gibt, zum Vater zu ziehen und ihm Josephs Herrlichkeit zu verkündigen; wie er seine Brüder stärkt, mit ihnen weint; und wie endlich seine Brüder mit ihm reden}

\subsection{I}

160. Er heißt sie eilends hinauf zum Vater ziehen; denn derselbe hat auch nicht weniger Trostes bedürft als seine Söhne. Ihr, sagt er, habt nun GOtt erkannt und mich auch; darum gebt hin, und übet nun dies Werk auch, daß ihr gehöret habt; und was ich euch in das Ohr gesagt habe, dasselbe gebühret euch nun auf den Dächern zu predigen, Matth. 10, 27.

161. Derhalben verkündigt es meinen Vater, und seinem ganzen Hausgesinde und Nachbarn. Wie Christus seinen Jüngern befiehlt Marc. 16, 15.: „Gehet hin in alle Welt, und prediget das Evangelium“ etc.; so sendet Joseph seine Brüder auch aus, und sagt: Eilet, säumet nicht, sondern ohne allen Verzug breitet aus, was ihr gehöret habt. Denn wenn wir GOtt erkannt haben in seinem Sohne Christo Jesu, nun Vergebung der Sünden und den Heiligen Geist empfangen, der die Herzen fröhlich macht und zufrieden stellt, durch welchen Frieden wir die Sünde und Tod verachten, was ist dann noch mehr zu thun? Da soll alsbald folgen: Gehe nun hin, schweige nicht stille, daß nicht du allein selig werdest, sondern daß der andere Hause auch mit dir möge erhalten werden.

162. Darum befiehlt nun Joseph seinen Brüdern dies allein mit großer Sorgfältigkeit, nämlich, daß sie ja dies Wort des Trostes und der Erlösung seinem Vater vortragen und ihn verkündigen, daß sein Sohn zum Herrn gesetzt sei über ganz Egyptenland. Er hat aber wohl gewußt, daß es seinem Vater auch würde schwer sein, solches zu glauben. Denn er dachte also: Mein Vater hat nun fast zwei und zwanzig Jahr um mich Leid getragen, und ist schier vor Traurigkeit seines Herzens gar entkräftet und erstorben; denn er hat vorlängst an meinem Leben und Wohlergehen verzagt, dieweil er in so langer Zeit nichts Gewisses oder Fröhliches von mir hat erfahren können. Denn das Jagen hat sein Herz so gar eingenommen, daß er nicht wohl glauben oder hoffen kann auf einige Hülfe oder Trost; viel weniger aber, daß er sollte glauben können, daß ich also sollte leben als ein Fürst und Herr über Egypten. Darum prägt er daselbe seinen Brüdern fleißig ein, daß sie ja dem Vater dies alles einzeln erzählen und verkündigen wollen.

163. Und diese Erzählung war ganz schwer und unglaublich; wie auch die gewesen ist, so der Engel bei dem Grabe gethan, Luc. 24, 5., da er zu den Weibern sagt: „Was suchet ihr den Lebendigen bei den Todten?“ Und hernach, da die Weiber dasselbe auch den Aposteln verkündigen, glauben sie es doch nicht. Denn die hohen unbegreiflichen Werke und Rath GOttes können nicht so bald begriffen werden: sie sind so hoch und groß, daß wir sie nicht so bald glauben können.

164. Also geht es zu unsern Zeiten auch zu, da die Kirche unbilliger Weise geplaget und verfolgt wird. Die Fürsten und Papisten vergießen vieler Gottseeligen Blut und geht ihnen solches frei aus. Auswendig“, wie Paulus sagt 2 Cor. 7, 5., „ist Streit“ und ist allenthalben voll der ärgsten Erempel; „inwendig“ ist Schrecken und „Furcht“ des Herzens: welches alles uns dies Licht und Leben sauer macht, daß wir demselben gar feind werden.

165. Was wird aber endlich daraus werden? Antwort: Nach so vielem Jammer und Trübsal wird ohne Zweifel eine unermeßliche Freude und die größte Ehre folgen, darin wir uns freuen und triumphiren werden durch wunderbaren und gewünschten Sieg, wie 1 Pet. 1, 11. gesagt wird vom Leiden Christi, und der Herrlichkeit, so darnach folgen werde.

166. Und wollte GOtt, daß nur der Tag unserer Erlösung und jüngsten Gerichts gar bald kommen möchte; denn nach demselben Tage habe ich viel größer Verlangen, denn sonst nach leiblicher Erlösung. Denn wir kriegen doch einen andern Mahomet und Papst. Es wäre besser, daß Alles auf einmal über den Haufen geworfen und aufgehoben würde mit allem Jammer und Unglück dieses Lebens, denn daß sonst noch eine andere zeitliche Veränderung geschehen sollte. Denn wir wissen, daß uns in alle Ewigkeit eine unaussprechliche Freude und unverwelkliche Krone bereitet und beigelegt ist, die wir jetzt nicht sehen. Wir sehen nur eitel Seufzen, Weinen, Jammer, Unterdrückung und Beschwerung der Armen. Wir sehen dem Teufel in den Hintern; GOttes Angesicht aber sehen wir nicht. Lasset uns aber glauben, hoffen, harren, beten, das Wort fleißig hören und demselben anhangen. Denn das meint GOtt, da er sagt im 55. Psalm V. 23.: „Wirf dein Anliegen auf den HErrn, der wird dich versorgen“; und wie Petrus 1. Epist. 5, 7. auch sagt aus demselben Psalm: „Er sorget für euch, fürchtet euch nicht“; item Paulus Phil. 4, 4—6.: „Freuet euch in dem HErrn allewege, und abermal sage ich: Freuet euch. Eure Lindigkeit lasset kund sein allen Menschen. Der HErr ist nahe. Sorget nichts; sondern in allen Dingen lasset eure Bitte im Gebet und Flehen mit Danksagung vor GOtt kund werden.“ Lasset ihn sorgen. Wir sind viel zu närrisch und zu arm oder gering dazu, als daß wir über die schwere Last der Sorge in unsern eigenen Sachen tragen sollten, auch nur einen Augenblick.

167. Derhalben ist das Allerbeste, daß wir nur GOtt vertrauen; wiewohl wir ohne Thränen, ohne Traurigkeit und Schmerz nicht sein können. Wie sollen wir ihm thun? Es ist um unser Fleisch also gethan, daß es sich des Seufzens, Weinens, Betrübniß und Klagens nicht enthalten kann. Ja, auch die unvernünftigsten Thiere, wenn sie geschlachtet werden, zeigen sie mit ihrem Geschrei an, daß es ihnen wehe thut, und die Creaturen, so kein Leben haben, als Holz etc., wenn sie abgehauen oder zerbrochen werden, geben sie auch einen Laut von sich und krachen. Wir aber sollen uns am Wort und Glauben halten, und nicht zweifeln, es sei schon bei GOtt beschlossen, daß er den Schmerz, Jammer und Kreuz, so nur eine kurze Zeit währet, damit uns der Pabst, Kaiser und Mahomet beschweren, in ewige und die allerhöchste Freude wenden und kehren wolle. Mit solchem Vertrauen und Hoffnung wollen wir unsern Jammer und Bekümmerniß leichtlich etwas lindern: er wird es wohl machen.

\quellseite{442}

168. Es hat aber Jakob nicht leichtlich glauben können, daß diese Botschaft wahr wäre, darum daß er ganze zwei und zwanzig Jahre am Leben seines Sohnes gar gezweifelt hatte. Denn es kommt oft vor, daß wir in nur einem Jahr derer vergessen, um welcher willen wir Leid getragen. Darum ist dieselbige Freude gar unversehens gekommen und war auch nicht wohl glaublich. Derhalben ist auch wohl gewesen, daß Joseph so oft wiederholet hat, daß er gesandt und von GOtt gesetzt wäre zu einem Herrn über Egyptenland.

169. Er heißt sie aber nicht gedenken, wie er sei von ihnen verkauft worden; sonst hätten sie dem Vater erzählen müssen, durch welche Gelegenheit und Ordnung alle Ding in derselben Sache geschehen wären. Denn ein solch verkehrtes Ding ist es um die Erbsünde, daß sie ihre Schande nicht will sehen lassen, sondern schmückt und bedeckt dieselbe auf das allerfleißigste. Darum gibt ihnen Joseph keinen Befehl, daß sie sollen sagen, wie sie ihren Bruder verkauft oder getödtet hätten; sondern er selbst hilft auch die Schande ihrer Sünde zudecken, und heißt sie nicht, daß sie die Sünde bekennen sollen, sondern sollen nur das erzählen, daß er ein Fürst über ganz Egyptenland und ein Vater Pharaos sey. Diese großmächtigen, wunderbaren Werke GOttes, das will er haben, sollen gerühmt werden, nämlich, daß der Sohn eines armen hebräischen Schäfleins zu solcher hohen großen Ehre und Herrlichkeit erhoben worden sei.

170. Dieweil er aber wohl gewußt, daß der Vater denjenigen, so ihm die Söhne erzählen würden, schwerlich würde Glauben schenken, wie man hernach sehen wird: darum hat er etliche Wahrzeichen oder Anzeichen hinzugethan, damit sie machen können, daß er ihren Worten glauben müsse. Ihr sollt, sagt er, im Lande Gosen wohnen. Was aber das für ein Land sei, weiß ich nicht. Etliche sagen, es sei dasselbe Land, das sie Thebais nennen. Das ist aber gewiß, es wird nicht weit von Joseph gelegen gewesen sein; wie der Text sagt: „Du sollst nahe bei mir sein.“ Darum wollen wir von dem Namen dieses Orts nicht disputiren, dieweil daran kein Zweifel ist, daß beide in diesen und andern Landen der Name der Orte oftmals ist verändert worden. Bisher haben wir nun gehört, was er ihnen für Befehl gegeben; nun thut er auch noch eine Bestätigung hinzu.

\begin{quote}
\textbf{V. 12. 13. Siehe, eure Augen sehen, und die Augen meines Bruders Benjamin, daß ich mündlich mit euch rede. Verkündiget meinem Vater alle meine Herrlichkeit in Egypten, und alles, was ihr gesehen habt; eilet, und kommt hernieder mit meinem Vater hieher.}
\end{quote}

171. Er will haben, daß sein Bruder Benjamin sein Zeuge und anstatt eines Bürgen sein soll bei dem Vater, auf daß er desto leichter glauben möge, daß Joseph noch lebe. Er stärkt aber und bestätiget sie mit diesem Zeichen, daß er in seiner Muttersprache mit ihnen redet, und jetzt keinen Dolmetscher braucht, wie er zuvor gethan. Und wir sollen nicht denken, daß er sich aus eitler Ehre also gerühmt habe, da er sie heißt, daß sie dem Vater alle seine Herrlichkeit erzählen sollen; welches ehrgeizige und ruhmredige Leute zu thun pflegen. Denn wir sehen ja, daß die höchsten und sonderlich ritterlichen Tugenden an ihm beschrieben werden, verglichen an Achilles, Scipio oder an Alexander nicht gewesen sind.

172. Da Joseph in der höchsten Ehre war, hat er sich doch gegen seine Brüder sehr demüthig und niedrig gehalten, daß er sie nicht verachtet, sondern auf das allerfreundlichste anredet, und vor großer Freude anfängt bitterlich zu weinen, da er von ihnen gesehen und erkannt worden ist; ja, daß noch mehr ist, er küßt sie und nimmt sie freundlich in den Arm. Darum wirst du kein ander Exempel finden, darin größere Freundlichkeit, auch lieblichere Art und Sitten sind denn an Joseph. Er geräth gar nicht des Todtschlags oder anderer Sünden, so sie an ihm, ihrem Bruder, begangen hatten, sondern dieselbigen alle werden gar verschlungen durch das unermessliche Meer der Sanftmüthigkeit, Freundlichkeit und Liebe, so an Joseph gewesen.

173. Diese Tugenden und vortrefflichen Werke leuchten an ihm und werden beschrieben uns zum Exempel und Besserung; wie Paulus sagt 2 Tim. 3, 16. 17.: „Alle Schrift, von GOtt eingegeben, ist nütze zur Lehre, zur Strafe, zur Besserung, zur Züchtigung in der Gerechtigkeit, daß ein Mensch GOttes sei vollkommen, zu allem guten Werke geschickt.“ Denn das bezeugt der Heilige Geist reichlich genug, der dies hat schreiben lassen, nämlich, daß GOtt an solchen geringen und fleischlichen Dingen Lust habe, die zwar äußerlich schlecht und verächtlich genug anzusehen sind, aber vor GOttes Augen sind sie sehr köstlich, dieweil sie im Geist und Glauben geschehen.

\subsection{II}

\begin{quote}
\textbf{V. 14. 15. Und er fiel seinem Bruder Benjamin um den Hals, und weinete; und Benjamin weinete auch an seinem Halse. Und er küssete alle seine Brüder, und weinete über sie. Darnach redeten seine Brüder mit ihm.}
\end{quote}

174. Siehe doch, welch eine treffliche Freundlichkeit und wunderbare Liebe dies ist, so Joseph zu seinen Brüdern gehabt. Denn er küßt nicht allein Benjamin, welchen er sonderlich lieb gehabt und in dem er auch seine Mutter Rahel und seinen Vater geehrt hat, sondern die Andern alle küßt er auch ohne Unterschied. Und mich wundert, wie er hat mögen gesinnet sein, da er mit dem Kusse und Weinen an Simeon und Levi gekommen ist. Denn dieselbigen waren die rechten Rädelsführer und Vorgänger gewesen der Sünden, so sie an ihm begangen, wie wir droben, Cap. 37., gehört haben; und er hat sie doch auch um den Hals genommen und geküßt wie die Andern.

175. Es ist wahrlich in solchem großen Fürsten, der so weise und heilig war, ein sanftmüthiger und demüthiger Geist gewesen. Denn es ist das nicht groß Wunder, daß er weinend und mit so großem Mitleiden Simeon, seinen Feind und Todtschläger, hat küssen können?

176. Das hat aber nicht unterbleiben können, sein Herz hat müssen sehr beweget werden, da er seinen leiblichen Bruder Benjamin um den Hals gefallen und ihn geküßt hat; denn demselben ist er mehr schuldig gewesen, dieweil er der Rahel Sohn war, als den Andern, so von der Lea geboren waren. Denn unser HErr GOtt läßt diesen Vorzug in dem Fleischlichen bleiben, und dasselbe zwar billig, nämlich, daß die natürliche Neigung unter denen größer sei, so von einerlei Eltern geboren sind, denn sonst gegen die Fremden. Wie Joseph Benjamin mehr geliebt hat denn die Andern, und sonderlich die Feinde, so ihm nach dem Leben gestanden waren. Doch nimmt er nichts Böses gegen sie vor, gedenkt auch ihnen gar nichts Böses zu thun, sondern küßt sie ganz freundlich. Welches alles der Heilige Geist also hat rühmen wollen, auf daß er damit verdunkeln und unterdrücken möchte die scheußlichen selbstvermählten Werke der Mönche. Denn dieweil diese Werke nach GOttes Wort und im Heiligen Geiste geschehen, sind sie höher, denn sonst alle Werke und Tugenden der allergrößten Helden.

177. Darum ist nun Simeon, der Joseph getödtet hätte, ihm noch gleichwohl sehr lieb und angenehm. Der Vater aber wird sich etwas härter gegen sie erzeigen, und wird ihnen die Sünde nicht so leichtlich vergeben und derselben so gar vergessen; denn er wird Simeon und Levi verfluchen.

178. Und da Mose die Kinder Israel gesegnet, gedenkt er Simeons gar nicht; denn Stamm Levi aber, davon er selbst auch geboren war, hat er reichlich gesegnet. Und an allen Orten, da er der zwölf Geschlechter gedacht, wird Simeon nirgends ausgelassen; aber da er seinen Segen soll erzählen, da ist Mose gar stumm. Hieronymus sagt, entweder nach seinem eigenen Gutdünken oder aber nach der Andern Meinung, es habe Mose Simeons darum nicht wollen gedenken, daß von ihm würde geboren werden der Verräther Judas.

179. Und ist auch jetzt noch unter dem losen Haufen der übrigen Juden der Stamm Simeon der allerverächtlichste; wie Lyra dasselbe bezeugt über 1 Mos. 49., da er sagt: Die Hebräer thun noch etwas dazu von der Zerstreuung Simeons und sagen: Es ist nirgends kein armer Schreiber oder Zuchtmeister denn nur von Simeon. Das ist darum gesagt, daß die vom Stamme Simeon umherliefen durch die Orte, so den andern Stämmen durchs Loos zugethan waren, daß sie ihre Nahrung mit Schreiben und Unterricht der Kinder gesucht haben; wie die armen Cleriker bei den Christen zu thun pflegen. Dies sagt Lyra an abgemeldetem Orte.

\quellseite{443}

180. Darnach findet man in der ganzen heiligen Schrift nirgendes, daß sie je etwas Treffliches sollten ausgerichtet haben. Daß also Simeon vor GOtt nicht ist trefflich gehalten worden, sondern trotzt mit unter den andern Zwölfen als ein armer Bettler. Denn Jakob hat ihn verflucht und Mose schweiget gar stille von ihm. Darum ist Joseph viel gelinder und gnädiger denn sein Vater Jakob oder Mose gegen seinen Bruder, sintemal er sie alle küßt; und da er dem Bruder mörder um den Hals gefallen, hat er auch über ihn geweint.

181. Ich halte aber dafür, sie werden nicht alle mit gleichem Gemüthe das Küssen von ihm angenommen haben. Simeon wird viel anders gesinnet gewesen sein denn Juda oder Issachar, darum daß sein Herz verwundet gewesen, dieweil er der Sünde halben ein böses Gewissen gehabt, davon er ohne Zweifel Beißen und Stechen gefühlt, und hat derselben begangenen Sünde nimmer vergessen können. Wie sie denn hernachmals, da ihr Vater gestorben und begraben war, nicht haben können zufrieden sein, vornehmlich Simeon. Denn da sie aus dem Lande Canaan wieder gekommen waren, haben sie ihn gebeten, daß er ja der Sünde, so sie an ihm begangen, nicht mehr gedenken wollte, und führen ihren Vater an, daß derselbe ihnen soll befohlen haben, ihn darum zu bitten. Da zeigen sie mit an, daß die Wunde ihres Gewissens noch nicht gänzlich geheilt gewesen, so er ihnen doch so viel Wahrzeichen gezeigt hatte, daran sie gewißlich nichts Anderes denn große Freundlichkeit, Barmherzigkeit und Wohltat haben spüren können.

182. Judas und Benjamins Herz ist viel friedsamer gewesen. Aber Simeon und die Andern haben sehr gezittert, und gemeint, er habe sich vielleicht um des Vaters willen bisher etwas demüthiger erzeigt: wo aber derselbe sterben würde, alsdann würde er wohl thun, wie er sich zuvor auch erzeigt hatte, und würde die Sünde an ihnen rächen. So gar armselig und argwöhnisch ist ein schuldiges Gewissen, und kommt ihm die Sünde und Uebertretung viel leichter an, als ihm das ankommt, daß es sich wiederum soll heilen lassen. Die Sünde geht süß ein, aber sauer aus.

183. Derhalben siehst du, daß Joseph mit Vergebung der Sünden seinen Vater und auch Mose übertrifft, dieweil er dessen, daß sie ihn verkauft hatten, nicht gedenkt. Der Vater aber wird Ruben, Simeon und Levi verschonen, und Mose schweigt Simeons gar. Allein Joseph redet sie freundlich an und geht auf das lieblichste mit ihnen um. Ach, mein lieber Bruder Simeon, Levi, wird er gesagt haben, ich will euch Alles von ganzem Herzen vergeben und vergehen. Sie waren damals in seiner Gewalt, und er hätte mit gutem Rechte die Schmach, so ihm begegnet war, an ihnen rächen mögen. Aber das hat er nicht thun können wegen des großen Mitleidens und seines freundlichen Herzens, so er zu ihnen gehabt, welches an nichts Anderes gedacht denn nur, daß er ihnen die Sünde vergeben, wohl thun und sie trösten möchte. Denn ihr Gewissen ist schwerlich verwundet gewesen und sie haben immer mit großen heftigen Schmerzen ringen müssen; welches nicht leichtlich gelindert oder aufgehoben werden kann.

184. Darum hat er sich ihrer erbarmt, und denkt bei sich selbst also: Wohlan, ich will ein Theologe, ein Prediger, ein Arzt, ein Knecht aller Knechte werden, auf daß ich die gewinnen möge, so mich haben verderben wollen; und er hat nicht allein mit Worten, sondern auch mit der That ihr Heil und Seligkeit gesucht. Es werden aber gleichwohl noch etliche Strafen der Gerechtigkeit folgen, so GOtt durch Jakob und Mose über sie wird gehen lassen. Denn unser HErr GOtt will es nicht so hingehen lassen.

\subsection{III}

185. Nach dem trefflichen Exempel aber der Sanftmüthigkeit und Liebe an Joseph fangen sie nun an, ihrem Bruder zu antworten. Denn sie haben sich bis daher entsetzt, und sind gar erschrocken gewesen wegen der großen Freude, so ihnen unversehens widerfahren, daß sie darüber verstummt und nicht ein Wort haben reden können. Aber die große Sanftmüthigkeit und Wohlthätigkeit, so sie an Joseph erfahren haben, macht ihnen Muth, daß sie nun desto getroster mit ihm reden.

186. Es gedenkt aber die Schrift der Reden nicht, so sie unter einander gehabt haben; man kann es aber leichtlich abnehmen aus dem, wie ein Jeder gesinnet gewesen, wie die Unterredung möge gelautet haben. Die, so unschuldig gewesen sind, als Dan, Issachar, Benjamin etc., werden vielleicht gesagt haben: Ach mein lieber Bruder Joseph, wir haben uns von ganzem Herzen vor der großen Sünde gescheut; aber du, Simeon, weißt wohl, wie höchlich du unsern Vater betrübt hast. Simeon wird solches erkannt, und wird es ihm mit aller Demuth abgebeten und gesagt haben: Ach, lieber Joseph, ich bin der Bube etc., ich wäre werth, daß du mich rädern und übern liefest. Juda und Ruben haben gesagt: Wir hätten dich wohl gerne wollen erretten und wieder zum Vater bringen; aber das war vergeblich. Und sie haben doch gesündigt, darum daß sie den Rath gegeben, daß man ihren Bruder verkaufen sollte.

187. Sie werden auch oftmals wiederholt haben die große Sorge und schwere Bekümmerniß, dazu auch die bittern Thränen ihres Vaters und der andern Mütter, als Bilha und Silpa etc., welche inzwischen vor Traurigkeit des Herzens gestorben waren, und werden noch von andern Dingen mehr geredet haben, die Haushaltung belangend. Und kürzlich davon zu sagen: Es wird eine demüthige Beichte oder Bekenntniß ihrer aller gewesen sein, und daß sie ihm haben abgebeten, was sie an ihm gesündigt hatten.

188. Und er sieht auch selbst ihren Jammer und Angst wohl, die schändliche Wunde; darum läßt er sich dünken, er müsse sich nur bemühen, daß er die, so erschrocken waren, wiederum aufrichte und heile, wie er droben, V. 5., gesagt hat: „GOtt hat mich vor euch hergesandt zu eures Lebens willen, und daß er euch erhalte auf Erden“ etc. Denn wie „ein guter Muth“, wie Salomo Spr. 15, 15. sagt, „ein täglich Wohlleben ist“, ist auch ein beschämender Gewissen und das sich der Sünden schuldig weiß, wie eine stetige Hölle, welche auch nichts Anderes sein wird denn ein böses Gewissen. Wenn der Teufel kein böses Gewissen hätte, so wäre er im Himmel. Das böse Gewissen aber zündet das höllische Feuer an, und erweckt im Herzen die erschreckliche Pein und höllischen Teufelein, die Grimmen, wie sie die Poeten genannt haben. Von außen her wird sich der Teufel weder vor Steinen, Feuer, oder anderer Pein, welcherlei sie auch sein möge, nicht fürchten; aber inwendig im Herzen wird er recht gemartert und gepeinigt. Der Zorn GOttes ist die rechte Hölle des Teufels und aller Verdammten.

189. Dagegen aber, wie große Freude und Stärke ein friedsam Herz habe, das sieht man an den heiligen Märtyrern wohl, die mitten im Leiden und Tode gelacht und GOtt gelobt haben, welches die Gottlosen nicht thun können wegen ihres bösen schuldigen Gewissens; daher es kommt, daß sie GOtt lästern, mit ihm zürnen und ihm von ganzem Herzen feind sind.

190. Derhalben lasset uns doch lernen auch um Furcht willen der Pein, und um des großen Schadens willen, der nicht zu heilen ist, welches ist ein verwundetes Gewissen, daß wir uns vor Sünden hüten. Denn es ist so ein überaus schwerer Kampf, welchen man mit einem Gewissen halten muß, das sich der Sünden schuldig weiß. Der Fall ist leicht geschehen; aber wieder aufzurichten und wiederum ein Verlangen oder Begierde gewinnen nach dem Troste göttlicher Güte und Barmherzigkeit, das hat große Mühe und Arbeit.

191. Wir haben jetzt etlichemal das Exempel der Jünger Christi angezogen und vorgehalten, welche erschrocken und bestürzt waren, nicht anders, denn als hätten sie selbst mit ihren Händen ihren Herrn erwürgt, darum daß sie von ihm gewichen waren. Was thut aber Christus dazu? Er redet mit ihnen ganze vierzig Tage vom Reiche GOttes, von Vergebung der Sünden, von GOttes Gnade und Barmherzigkeit, item, von seinem Leiden. Ihr habt mich nicht getödtet, sagt er, auch nicht erzürnt; aber dies ist meines himmlischen Vaters Rath und Wille gewesen, daß ich habe sollen den Tod leiden und wiederum auferstehen, auf daß ihr in meinem Namen Vergebung der Sünden gepredigt würde. Doch hat er sie damit nicht alle aufgerichtet; sie haben sich sehr bekümmert, ohne Zweifel darum, daß sie ihn also verlassen hatten. Ach, werden sie gesagt haben, es ist uns leid! und haben sich noch so hin mit den Worten Christi trösten lassen. Es hat aber gleichwohl ihr Herz so gar friedsam und rein nicht sein können, als wenn sie dessen gar nichts begangen hätten.

192. Da Petrus zum drittenmal von Christo gefragt ward Joh. 21, 16.: „Simon Johanna, hast du mich lieb? Weide meine Schafe“, da steht der Evangelist Johannes hinzu, er sei darüber sehr betrübt geworden. Ei, warum das? Antwort: Der Stachel der Sünde und daß er an dieselbe gedacht, das hat ihn erschreckt; wiewohl ihn Christus ganz freundlich anredet und eben damit seinen Glauben hat aufwecken und stärken wollen, macht doch gleichwohl, da er an seine Sünde denkt, daß er sagt: „Herr, du weißest, daß ich dich lieb habe.“ Zuvor hat er gesagt, Luc. 22, 33.: „Ich bin bereit, mit dir ins Gefängniß und in den Tod zu gehen“; und Christus hätte ihm solches wohl mögen aufdrücken.

\quellseite{444}

193. Aber er hält sich eben also gegen Petrus und die andern Jünger, wie sich Joseph hier gegen seine Brüder auch hält. Er versammelt sie, da sie erschrocken waren, tröstet sie, schüttet über sie gar heraus seine herzliche Barmherzigkeit und den Reichthum der Gnade von Vergebung der Sünden; dennoch hat er vierzig Tage an ihnen zu binden und zu heilen. Also hat sich Joseph auf etliche Tage beflissen, daß er die Wunde seiner Brüder heilen möchte durch die freundlichsten Unterredungen; und ist auf das vertraulichste und brüderlich mit ihnen umgegangen, auf daß er durch sie dem Vater seine Bekümmerniß auch wegnehmen möchte, der auch seines Unfalls halben, nicht von wegen der Sünde oder Missethat, so er möchte begangen haben, sondern vor Traurigkeit und sonst vieler Angst und Noth im schweren Kampf gelegen hat.

194. Darum ist das Exempel Josephs ein wahrhaftiges und schön klar Bild der Auferstehung Christi. Und geschehen solche Werke noch täglich in der Kirche. Denn das ist eigentlich der Bischöfe und Pfarrherren oder Prediger Amt, daß sie lehren, aufrichten und trösten sollen, nicht die, so verhärtet und unempfindlich sind, die sich mit keinen Worten wollen trösten lassen, welche man viel lieber dem Stockmeister oder Henker befehlen mag; sondern daß sie den betrübten und bekümmerten Herzen den Balsam der heiligen Schrift reichen und mittheilen sollen und zu ihnen sagen: „Fürchte dich nicht“ etc. „Sei getrost, mein Sohn, deine Sünden sind dir vergeben.“ Matth. 9, 2.

195. Aber wie ichs immer dasselbe angehe, daß man solchen Trost der Schrift an sich bringen möge, das bezeuget beides meine und auch anderer Leute Erfahrung. Ich habe die Bibel nun bei dreißig Jahren mit höchstem Fleiß und Mühe gelesen; doch bin ich dermalen noch nicht also gar geheilt, daß ich mit vollkommenem Vertrauen mich an die heilsame Arznei, so uns Gott selbst gezeiget hat, halten könnte. Ich wollte wohl gern stärker sein und einen größern Glauben haben, wollte gern etwas hoffkräftiger in Christo sein, daß ich mich sein mehr trösten und rühmen könnte, aber ich kann es nicht zuwege bringen.

196. Derhalben wird die Wunde geheilt, da man in der Herberge Oel und Wein drein gießet. Denn die Kirchen sind nichts Anderes, denn solche gemeine Herbergen, darin das Volk geheilt wird, welches seine Sünde und den Tod fühlt, dazu auch den Schrecken und die Plage des armen verwundeten Gewissens. Da müssen wahrlich die Kirchendiener zusehen, daß sie dieselben Wunden treulich und fleißig heilen, daß sie Oel und Wein drein gießen.

197. Denn die Sünde ist die aller greulichste Krankheit, und das Gewissen ist ein sehr zart Ding und wie ein kleines Holzwürmlein, wie der alte Dolmetscher der Bibel von David sagt. Denn wo ein solch Würmlein nur ein wenig an der Haut verletzet wird, so stirbt es bald dahin. Und der gemeine und gute Spruch lehrt auch dasselbe, welcher also lautet: Non patitur ludum fama, fides, oculus, das ist: Mit dem guten Namen, dem Glauben und dem Auge ist nicht gut zu scherzen. Denn wie der Glaube, das ist, das Gewissen, in einem kleinen Augenblick kann verletzt werden: also ist auch das Auge und ein gut Gerücht ein sehr zart Ding, welches leichtlich kann verletzt oder beschädigt werden. Es begibt sich oftmals, daß aus einem Wörtlein, so unbedachtsam geredet oder aber nicht recht verstanden worden, viel unzähliche Lästerungen und Ärgerreden hergekommen. Also ist nun die Versöhnung und Heilung dieser Brüder aufgerichtet worden, wiewohl sie noch nicht so gar zurechtgebracht sind. Und ist ein sehr schön Exempel der Heilung, so in der Kirche durch Diener und Lehrer zu geschehen pfleget.

198. Mose hat aber sehr lieblich abgemalt, wie Joseph und Benjamin einander um den Hals gefallen sind. Denn also lauten die Worte im Texte: „Und er fiel seinem Bruder Benjamin um den Hals, und weinete; und Benjamin weinete auch an seinem Halse.“ Also haben auch Gad, Naphtali, Sebulon und die andern an ihm gehangen, und er hat mit großer Lust eines Jeden Rede von seinem Vater gehört, von den verstorbenen Müttern und von andern Dingen mehr. Die Andern, so älter gewesen, haben nicht so nahe hinzutreten dürfen, wiewohl sie davon auch etwas erzählt haben, wie es um den Vater gestanden, und wie ängstlich und bekümmert er gewesen wäre.

199. Daher ist aber das Gerücht über den ganzen Hof des Königes ausgekommen, wie sich Joseph und seine Brüder erkannt, und wie sie mit einander Gespräche gehalten, und der hebräische Schafhirte und arme Bettler, der kaum ein armes Schäferhüttlein im Lande Canaan gehabt hat, ist über ganz Egyptenland hoch gerühmt worden. Zuvor hat davon niemand etwas gewußt, wer dieser Joseph wäre, oder wer seine Eltern wären, von denen er geboren war.

200. Und ist wahrlich Wunder, daß er so hoch erhoben gewesen und das Regiment in Hand gehabt, daß er nicht Boten ausgesandt hat, die erforscht hätten, wie es um das ganze Geschlecht stehe. Aber er gibt selbst Ursache, warum dasselbe nicht geschehen sei, da er den Namen Manasse, seines Sohnes, deutet und sagt 1 Mos. 41, 51.: „Gott hat mich lassen vergessen alles meines Unglücks, und alle meines Vaters Hauses.“ Denn da zwei und zwanzig Jahre seit der Zeit verlaufen waren, da er verkauft worden, hat er gedacht, der Vater würde inzwischen vor Bekümmerniß und Leid gestorben sein, Benjamin aber würde durch die Bosheit seiner Brüder auch umgekommen sein, auf daß Simeon das Regiment allein haben möchte.

201. Und ist das ganze Spiel, so er mit seinen Brüdern gespielt, dahin gerichtet gewesen, daß er desto genauer forschen möchte, wie es um den Vater stünde, ob er auch noch mit dem ganzen Hausgesinde frisch und gesund sein möchte. Also hat er fast alle Hoffnung fallen lassen, daß seine Eltern und Brüder noch am Leben sein sollten, und hat gedacht, es müßte das ganze Haus seines Vaters jämmerlich zerrüttet und gar verwirrt sein. Denn ihr habt gehört, wie gar fleißig und genau er nach Benjamin gefragt, und wie er versucht, ihnen auch mit Banden und Gefängniß abzudringen, daß sie ihm wahrhaftigen Bericht gethan, wie es eigentlich um Benjamin stehen möchte, bis er ihn selbst gegenwärtig sehen könnte.

\section[Fünfter Theil]{Fünfter Theil. Wie das Geschrei von Josephs Brüdern in Pharaos Haus gekommen und Pharao sammt seinen Knechten sich solches wohlgefallen läßt; item, wie Pharao Befehl gibt, Jakob sammt seiner Familie nach Egypten zu holen}

\subsection{I}

\begin{quote}
\textbf{V. 16. Und da das Geschrei kam in Pharao Haus, daß Josephs Brüder kommen wären, gefiel es Pharao wohl, und allen seinen Knechten.}
\end{quote}

202. Das Geschrei von den Brüdern Josephs ist über den ganzen königlichen Hof gekommen, und sie werden hin und wieder über ganz Egypten und in den Ländern nächst umher mit ihrem Vater Jakob sehr gerühmt, so er doch im gelobten Lande, darin er dazumal gewohnt, gar verachtet und unbekannt gewesen. Aber in dem Lande, darin er ein Fremdling ist, wird er nun groß und steht seine Sache daselbst ganz wohl. Darum ist es ja ein wunderbärlich Ding, daß er an dem Orte, so ihm sonderlich von Gott verheißen und bestimmt war, so gar verachtet und unbekannt ist, und nicht ein Haar breit von den Gütern besitzt, so er vom Lande Canaan hat erwarten sollen. Ja, daß noch mehr ist, er wird daselbst greulich geplagt und gemartert, da Bilha und Dina geschwächt sind und Joseph verkauft und verloren worden, und dazu Rahel, sein allerliebstes Weib, in der Geburt jämmerlich gestorben ist. Da ist gar keine Ehre, sondern die höchste Schande mit großem Schmerz und Traurigkeit.

203. Und Joseph hat zwar zuvor von seinem Geschlecht und Ursprung nicht rühmen können, dieweil er es dafür gehalten, daß sein Vater vorlängst von großem Verlangen, so er nach ihm gehabt, gestorben sein müßte, und hatte darum seinen erstgebornen Sohn Manasse geheißen, und sagt nur, er wäre ein hebräischer Mann, und heimlich weggeführt und also nach Egypten kommen. Da er aber nun jetzt seine Brüder erkannt und wieder gesehen, und das Geschrei von ihnen schon etliche Tage über das ganze Königreich gegangen war, da wird nun Jakob erhöhet und im fremden Lande hin und wieder von jedermann gerühmt; ja, das noch mehr ist, Pharao selbst und alle seine Knechte haben diese Botschaft gern gehört und sich dieselbe sehr wohl gefallen lassen.

\quellseite{445}

204. Die Ursache aber, warum dem König dies so sehr wohlgefallen hat, zeigt Lyra an: Daß wiewohl er Joseph lieb gehabt wegen der großen Wohlthaten, so er dem ganzen Königreich erzeigt, hat er es doch dafür gehalten, daß ihm nicht sehr rühmlich wäre, daß ein erkaufter Knecht, der nämlich aus dem Gefängniß geführt, ein Herr über ein solch großes Königreich sein sollte. Denn wir sehen, wie Cicero sagt, wie die Leute, so ihrer Tugend und Geschicklichkeit halben aufkommen und groß werden, so sie doch sonst von geringen Eltern herkommen, bei etlichen vom Adel so sehr gehaßt und angefeindet werden. Darum, sagt Lyra, sei dem König dazu sehr lieb gewesen, und habe sich dessen gefreut, daß er nun Josephs Geschlecht und Eltern, von denen er hergekommen, habe kennen mögen.

205. Und müssen gewißlich beide der König und die Edelleute in Egypten fromme gottselige Männer gewesen sein, die Joseph so sehr lieb gehabt, und sich an ihm der unglaublichen Weisheit, Heiligkeit, Gottseligkeit und sonst allerlei Tugenden höchlich verwundert haben. Jedoch wird es auch an etlichen Mißgönnern und Widersachern nicht gefehlt haben. Denn der Text sagt nicht, daß diese eben allen Knechten gleich wohlgefallen habe. So ist auch der Hofteufel nicht todt, und an der Könige und Fürsten Höfen regieren die selben Hofsünden viel mehr denn sonst unter den Familien anderer gemeiner Leute. Und wie groß daselbst die Gewalt des Teufels sei, zeigen die Worte an im Propheten Daniel am 10. Cap. V. 20., da der Engel sagt: „Weißest du auch, warum ich zu dir kommen bin? Jetzt will ich wieder hin, und mit dem Fürsten in Persienland streiten; aber wenn ich wegziehe, siehe, so wird der Fürst aus Griechenland kommen.“

206. Also ist ohne Zweifel der Teufel in Egyptenland auch gewesen; aber der Segen des Herrn hat seinen grimmigen Zorn überwunden, und Gott hat diesen Fremdling hoch gehoben, und hat ihm einen König gegeben, der sanftmüthig und freundlich gewesen, und dazu etliche Fürsten und Herren im Lande, die auch ehrliebende und gottselige Männer gewesen sind, die ihn lieb gehabt und in Ehren gehalten haben; wie dasselbe aus dem 105. Psalm V.  21. zu sehen ist.

Die Andern aber, so seine Mißgönner gewesen und ihn angefeindet, haben ihn auch oftmals hart geplagt; wie es Daniel im Königreiche Babylon auch ergangen ist. Denn wiewohl er ein heiliger, frommer und unschuldiger Mann war, ist er doch endlich in die Löwengrube geworfen worden durch falsche unrechtmäßige Anklage und Beschuldigung, damit ihm die Hoffschranzen bei dem Könige angegeben hatten.

207. Das ist also gebräuchlich an allen Fürstenhöfen, daß die Frommen durch Lügen und Trügen der Lästermauler unterdrückt werden. Droben, Cap. 39, 20., ist die schwere Anfechtung, da er in das Gefängniß geworfen worden ist, beschrieben: aber was ihm hernach im Regiment für Unglück oder Widerstand begegnet, wird im Text nicht erzählt.

208. Wiewohl nun das ganze Geschlecht der Viehhirten bei allen Egypten ein Greuel gewesen, ist doch der König so gottseelig und so fromm, daß er sich an solchem geringen verächtlichen Stande nicht geärgert hat. Laßt mir den Schäfer bezähmen, wird er gesagt haben. Er ist ihm lieber gewesen, denn sonst viel andere vom Adel und große Herren in selben Königreiche. Ja, das ist die Hand Gottes, „der den Geringen aufrichtet aus dem Staube, und erhöhet den Armen aus dem Koth, daß er ihn setze neben die Fürsten, neben die Fürsten seines Volks“, wie der 113. Psalm V. 7. 8. sagt.

209. So ist nun daran kein Zweifel, er wird dem König sein ganz Geschlecht erzählt haben, bis zurück auf Tharah und Nachor, und daß er angezeigt und bewiesen habe, wie er aus dem berühmten Geschlechte der allervornehmsten Männer geboren sei. Und solches ist dem Könige sehr angenehm und lieb gewesen. Denn, daß Lyra hinzusetzt, wie die Brüder Josephs so gar hübsch und schön von Angesicht gewesen, daß man an solcher schönen Gestalt leichtlich hat abnehmen können, daß sie berühmte Leute wären, dasselbe däucht mich nicht der Wahrheit gemäß. Denn es sind Schafhirten gewesen, die eine schlechte Gestalt und Kleidung gehabt, wie dieselben Leute zu haben pflegen. Darauf haben der König und seine Edelleute nicht vornehmlich gesehen, sondern haben Joseph lieb gehabt um seiner Tugend und großen Weisheit willen.

\subsection{II}

\begin{quote}
\textbf{V. 17. 18. Und Pharao sprach zu Joseph: Sage deinen Brüdern: Thut ihm also, beladet eure Thiere, ziehet hin; und wenn ihr kommt ins Land Canaan, so nehmet euren Vater, und euer Gesinde, und kommt zu mir; ich will euch Güter geben in Egyptenland, daß ihr essen sollt das Mark im Lande.}
\end{quote}

210. Wiewohl Joseph droben, V. 9., aus eigener Macht seinen Brüdern befohlen hatte, daß sie den Vater aus dem Lande Canaan zu ihm bringen sollten, und ihnen versprochen hat, er wollte ihn mit seinem ganzen Haus und allem Geschlecht ernähren, wird es doch nun hier etwas klarer ausgedrückt, daß er solches nicht ohne Bewilligung und Befehl auch des Königs selbst thun will. Darum ist er nicht hoffärtig gewesen, und hat sich nicht über die Maßen erhoben wegen der großen Herrlichkeit und Herrschaft, so er über ganz Egyptenland gehabt; sondern hält sich fein demüthig und niedrig, ehrt den König und nimmt nichts ohne seinen Willen vor. Wiederum wird auch hier die Freundlichkeit des Königs gerühmt, welcher dem Joseph in demselbigen bald beigestallen und ihm sammt seinem Vater herzlich gern alles Gute thun will.

211. Es ist an Joseph solcher Stolz und Frevel nicht gewesen, wie man zu unsern Zeiten bei den Steuerbeamten oder Rentmeistern findet, die aus eigener Vermessenheit und Bosheit allerlei in der Fürsten und ihrer Herren Sachen zu thun wagen, da oft die Fürsten und Herren selbst nichts darum wissen, oder gar etwas ohne ihren Willen geschieht. Ein solcher Mann ist Joseph nicht gewesen. Denn er ist nicht aufgeblasen oder stolz wegen der großen Ehre, dazu er gekommen war, sondern hat sich fein mäßig und demüthig zu halten wissen, und hat erkannt, daß er schuldig wäre, den König in Ehren zu halten und auf ihn zu sehen.

212. Ein solch Herz ist wahrlich \unsicher{kein selbst noch mächtig}, daß sich im Glück nicht überhebt, sondern sich fein in demselben Maße hält, wie es sich zuvor auch gehalten, ehe es zu solcher Herrschaft und Ehre gekommen ist. Aber das ist eine sehr seltene und wunderbare Tugend. Wie die heidnischen Poeten und Geschichtschreiber auch über die große Schwachheit und Unbeständigkeit der Herzen der Menschen klagen. Virgil sagt:

\begin{quote}
Nescia mens hominum fati sortisque futurae, Et servare modum, rebus sublata secundis;
\end{quote}

das ist: Das menschliche Herz kann nicht zuvor wissen, was ihm für Glück oder Unglück begegnen werde, kann auch nicht Maß halten, wenn es vom Glück erhoben wird.

213. Des Menschen Herz kann es nicht lassen, es muß sich überheben, wenn Glück vorhanden ist; und wiederum im Unglück und Widerwärtigkeit wird es gar kleinmüthig und verzagt; wie daselbst alle Historien auch der allergrößten Männer bezeugen. Sie sagen vom Cajus Cäsar, wenn es ihm nicht nach seinem Willen gegangen sei, soll er immer sauer gesehen und den Kopf haben hängen lassen. Darum ist es ein großes Ding, daß ein Mensch beständig ist, es sei im Glück oder Unglück. Und dazu gehört ein solch Herz, daß ein gewiß vollkommen Vertrauen auf Gott allein hat, und daß Gottes nicht vergisst, wenn ihm Alles glücklich und nach seinem Willen geht, und auch nicht verzagt, wo es fehlt und vernimmt, daß sein Ding so glücklich nicht hinaus gehen will, wie er es wohl gerne wollte. Ja, der Glaube ist auch stärker in Noth und Widerwärtigkeit denn im Glück, nach dem Spruch 2 Cor. 12, 9.: „Meine Kraft ist in den Schwachen mächtig.“

214. Derhalben ist Joseph ein großer und wunderbarer Mann gewesen, der dieses Maß in seinem großen Glück hat halten können, daß er zuvor im Gefängniß und Knechtschaft auch gehalten hatte.

215. Und dieser Pharao wird billig auch unter die frommen gottseligen Könige gerechnet, welchen auch ohne Zweifel Gott um Josephs willen gesegnet haben wird, und weil er sich selbst auch sonderlich der Gottseligkeit befliessen; wie er denn von Joseph beide in Religions- und Polizeisachen unterwiesen und berichtet war. Denn Joseph ist in Egypten nicht müßig gewesen, sondern hat die rechte göttliche Lehre mit höchstem Fleiß ausgebreitet, und dieselbige dem König, den vornehmsten im Lande und dem ganzen Volk vorgetragen. Der Herr hat ihm einen gehorsamen König gegeben, der ihm gern gehorcht, und dazu viel Fürsten, Priester und andere Regenten, so nach derselben rechten Lehre begierig gewesen sind, daß sie dieselbe haben hören und lernen mögen.

\quellseite{446}

216. Und es ist kein Wunder, wenn etliche andere Gottlose und Ungläubige auch darunter gewesen sind. Es ist daran genug, daß das Königreich ingemein das Wort gehabt. Wie bei uns auch viel böse verkehrte Leute sind; dennoch ist die Kirche und Schule der Hauptsache nach gottseelig. Denn es wird ja darin nichts gelehrt, das ungöttlich oder unzüchtig wäre, sondern was göttlich und ehrlich ist, das lehrt man darin. So aber Etliche unter den Häusern sind, die solcher Gaben und Wohlthaten GOttes mißbrauchen, nach derselbigen Bosheit fragen wir nicht, lassen uns solches nicht anfechten, daß wir darum nicht urteilen sollten nach der rechten Substanz oder Wesen. Wie sich am Leibe des Menschen mehr Unflat, Wust und Dreck sammelt und von Essen und Trinken ausgeworfen wird denn Fleisch und Blut, und wird doch vom Leibe gesagt, daß er gesund und rein sei, wenn seine Substanz oder Wesen unversehrt bleibt. Denn wir urteilen nicht nach dem, das schändlich ist, sondern nach der Substanz und rechtem Wesen. Wenn aber die Predigten, Gericht und Obrigkeit, oder die Aemter in der Kirche und Polizei verderbt und darum Mangel gefunden wird, das ist das Allerhöchste und kann daselbst leidlich eine teuflische Unordnung und Fälscherei gefunden werden, wiewohl Alles einen großen Schein der Heiligkeit hat; wie sich die Ketzer ihrer guten Werke höchlich rühmen können.

217. Darum ist das Königreich Egypten zu der Zeit sehr schön und in großem Flor gewesen, nicht allein leiblich, sondern auch geistlich. Nun aber wird es noch etwas größer und herrlicher wegen der Ehre, so der Patriarch und Prophet Joseph gehabt. Und GOtt sammelt nun die ganze Kirche in dasselbe Königreich, und ist diesem Könige die große Ehre und Herrlichkeit widerfahren, daß er ein Wirth und Pfleger geworden ist, der die ganze Kirche beherbergt und unterhalten hat; Sintemal er Jakob mit seinen Söhnen und Töchtern und allem Hausgesinde hat heißen zu sich nach Egypten kommen.

218. Denn Jakob ist nicht allein ein Schafhirt oder Hauswirth gewesen; welches die Papisten allein an ihm betrachten, und können nichts mehr, denn daß sie den Ehestand und Kinderzucht, dazu die Hausforge und dergleichen Dinge mehr, nur verachten. Aber wir sollen ihn deshalb desto mehr in Ehren halten, weil er ein Patriarch, Prophet und Engel GOttes gewesen ist, welcher allenthalben, wo er nur gemein ist, GOttes Wort bei sich gehabt; welches nichts Anderes ist, denn GOtt selbst und seine heiligen Engel bei sich haben. Denn wo die Kirche ist, da ist auch das Amt und Dienst der lieben Engel; wie das der 34. Psalm V. 8. bezeuget: „Der Engel des HErrn lagert sich um die her, so ihn fürchten, und hilft ihnen aus“; item im 125. Psalm V. 2. steht also: „Um Jerusalem her sind Berge; und der HErr ist um sein Volk her, von nun an bis in Ewigkeit.“

219. Darum hat der fromme heilige König Pharao nicht einen geringen verachtet und aufgenommen und beherbergt, da er Jakob und seinen Kindern die Herberge gegeben hat, sondern ist ein solcher Wirth geworden, der die Kirche GOttes und der heiligen Engel beherbergt; welches wahrlich eine große Ehre für Egyptenland gewesen ist. Hernach aber gedenkt dieses Land der empfangenen Wohlthat gar nicht (wie denn die undankbare Welt zu thun pflegt), und hat seinen Gästen nicht gedankt, wie es billig sollte gethan haben.

220. Derhalben muß man an den Vätern nicht allein das so sehr ansehen, daß sie Schafhirten gewesen, die Ziegen gemolken und andere häusliche Arbeit gethan haben, als vielmehr, was an Jakob und den Andern auch das Vornehmste gewesen ist, nämlich, die göttliche Verheißung, die also lautet: „Durch deinen Samen sollen alle Völker gesegnet werden“, 1 Mos. 26, 4. Es soll der Heiland der Welt von dir herkommen. Denn der vornehmste und allermächtigste theuerste Schatz und Schmuck in den Historien der Patriarchen ist das Wort GOttes.

221. An diesem Lichte von GOttes Willen mangelt es den Mönchen ganz und gar; und daß GOtt ihren Orden bestätigen und sich denselben gefallen lassen sollte, das fehlt weit; denn sie können nicht rühmen, daß ihr Stand ein GOttes Wort oder Verheißung gegründet sei. Und wenn du dasselbe fleißig betrachten oder bewegen wirst, so wirst du sehen, daß sie ganz fleischliche Menschen sind und ohne GOttes Wort. Darum ist daran kein Zweifel, sie müssen vom Teufel sein. Ja, sagen sie, ich schlaffe doch bei keinem Eheweibe, habe nichts zu schaffen mit täglichem Sorgen oder häuslicher Arbeit, ich warte nur meines Gebets, Lesens und Medicirens. Es ist recht gethan; aber du kannst gleichwohl dessen nicht gewiß sein, daß dein Stand von GOtt eingesetzt, oder durch das Wort der Verheißung bestätigt sei.

222. Wie viel besser thätet ihr aber daran, lieber Franciscus, Dominicus, und alle Päbste und Cardinalen, wenn ihr die Kühe melket, das Haus kehrt, oder sonst allerlei andere Hausarbeit thätet. Denn ein Hausvater hat auch zeitliche Verheißung, darum daß GOtt den Ehestand im vierten Gebot eingesetzt hat, da er sagt: „Du sollst deinen Vater und deine Mutter ehren“ etc. Darum ist je die Vaterschaft oder das Amt eines Hausvaters und der ganze Ehestand mit allen häuslichen Werken geheiligt, geziert und eingesetzt im vierten Gebot. Und über das alles ist noch eine sonderliche Verheißung zugesetzt, nämlich, daß die Kinder, so ihre Eltern in Ehren halten, lange Leben sollen; das man von den Dompfaffen, Päbst und Cardinalen gar nicht sagen kann.

223. Endlich aber, wenn sie in den letzten Zügen liegen und nun sterben sollen, haben sie gar kein rechtes Licht oder beständigen Trost, sondern müssen verzagen. Wie ich vor dieser Zeit einen Dompfaffen zu Erfurt gekannt habe, der allein für seine Person zwei und zwanzig geistliche Lehen gehabt: da er aber in den letzten Zügen lag, und nicht wußte, wohin er aus diesem Leben fahren würde, hat er mit jämmerlichem Seufzen und Heulen diese Worte geredet, die voll Verzweiflung gewesen: O daß ich die Zeit meines Lebens etwa eines armen Schäfers oder Kuhhirten Knecht gewesen wäre. Denn sie leben also roh dahin, wenn sie gesund sind, ohne GOtt, ohne Licht, ohne Kreuz oder Leiden und sterben dann auch also dahin.

224. Darum sollen wir uns dessen befleißigen, daß wir in unserm Stande auch leibliche Verheißungen von GOtt haben mögen, darauf sich das Herz verlassen könne. Wie St. Paulus sagt 1 Tim. 4, 8.: „Die Gottseligkeit ist zu allen Dingen nütze, und hat die Verheißung dieses und des zukünftigen Lebens“; wie auch die Väter Verheißung des Hausregiments, der Polizei und auch der Kirche oder das ewige Leben belangend gehabt.

225. Darum hat Pharao nicht einen geringen gemeinen Mann aufgenommen, der nur ein Hausvater oder weltlicher Regent wäre, sondern er hat die Kirche und Gemeinde GOttes selbst aufgenommen, welche er erhält, nährt und schützt. Das ist eine trefflich große Ehre und königliche Herrlichkeit, damit er billig den Namen erlangt, daß er ein heiliger gottseeliger König genannt worden ist und den GOtt sonderlich hoch geehrt hat. Denn er ist ja von GOtt zu diesem großen Amte erwählt worden, daß er die Kirche hat erhalten sollen, daß er ihr Stätte, Zeit und Schuh gebe, daß dieselbe Kirche gesammelt, gemehrt und erbaut werden möchte.

226. So ist dies nun ein sehr schöner Text an diesem Orte, daß Pharao so große Lust und Gefallen an dem Geschrei von Josephs Vater und seinen Brüdern gehabt. Und wiewohl die Egypter vor den Hirten einen Greuel gehabt, so hat er doch dies Haus und Geschlecht Jakobs und Josephs groß geachtet. Denn wo ich deines Vaters Haus ansehe, wird er gesagt haben, so sehe ich, daß es sehr schön und von GOtt herrlich geziert und geehrt ist. Du, Joseph, aber bist voll Heiligen Geistes, übertriffst mit Weisheit, Gottseeligkeit und mit Rath alle anderen Fürsten und große Herren im ganzen Königreiche; welche zwar vor den Hirten einen Greuel haben; und ich kann auch diese weltliche Satzung oder Landesgebrauch nicht strafen, bessern oder ändern. Aber wie dem allem sei, so habe ich gleichwohl dazu Lust, daß ich deinen Vater und dein ganzes Geschlecht, wiewohl es Schafhirten sind, sehen, und sie aufnehmen und beherbergen möge.

227. Ach, wollte GOtt, daß wir zu unsern Zeiten solcher Fürsten auch viele hätten! Denn dies ist eine große treffliche Tugend an Pharao gewesen, dergleichen man nicht viel findet, die billig von jedermann gerühmt wird. Und daß er voll Glaubens und Heiligen Geistes gewesen sei, erhält klarlich daraus, daß er so gutwillig und gar umsonst, und daß er dafür nichts begehrt, die Kirche genährt und unterhalten hat. Denn zu dieser Zeit wirst du Viele finden, die sich der Lehre des Evangeliiums auch annehmen und derselben so lange gewogen sind, bis sie die Pfarren und Pfarrgüter, Stifte und Klöster gepackt und beraubt haben. Aber deren sind leider! gar wenige, die um der Gottseeligkeit wollen und von Herzen den Lehrern, Pfarrherren oder Kirchendienern recht günstig wären. Die aber solche sind, folgen billig dem Exempel Pharaos nach, welches wahrlich das Allerchristlichste und Heiligste ist, das sie nur thun können. Denn Pharao hat ja nicht seinen eigenen Nutzen gesucht, steht auch sonst nirgends anders hin, da er die Kirche nährt und erhält, denn nur auf die Gottseligkeit und den rechten Gottesdienst, daß nur derselbe gefördert werden möge. Er nimmt oder reißt nichts an sich, er gibt vielmehr sein eigen Gut dahin und wendet es darauf, daß das Haus Jacobs also erhalten und genährt werden möge.

\quellseite{447}

228. Dies ist wahrlich ein solch Exempel, das billig alle Könige, Fürsten, Obrigkeiten und auch die von Adel sich selbst vorstellen und dem sie mit höchstem Fleiß nachfolgen sollen, nämlich, fromme und gottselige Pfarrherren, Kirchendiener und Schüler zu nähren und zu unterhalten. Aber es geschieht leider zu unsern Zeiten gar das Gegenteil nicht allein von den Päbstlichen, sondern auch von denen, die die Reformation und Besserung in den Kirchen aufrichten, und unter solchem Schein die Kirche der Güter berauben, so billig zur Erhaltung des Predigtamts und der Schulen angelegt werden sollten.

229. Durch das „Haus Jakobs“ aber versteht er nicht Holz und Steine; denn er hat mit seinen Kindern nicht in Häusern, sondern in Hütten gewohnt nach der Weise der Hirten in Arabien, oder der Völker, die Nomaden heißen, welche nicht an einem gewissen Orte bleiben konnten, sondern an mancherlei Orten Weide haben suchen und derselben immer nachfolgen müssen. Derhalben verstehen wir auch durch das „Haus Jakobs“ sein ganz Geschlecht, seine Söhne und Töchter, dazu der Söhne und Töchter Kinder, Blutsfreunde und Schwäger, die alle mit ihm gewesen; und ihrer wird ohne Zweifel eine große Anzahl gewesen sein, die mit Jakob hinab nach Egypten gezogen sind. Hernach wird gesagt: „Alle Seelen des Hauses Jakob, die in Egypten kamen, waren siebzig.“ Aber die Weiber, die Töchter und Enkelinnen, Knechte, Mägde, und Schwäger mit ihren Weibern werden nicht mitgezählt, sondern nur die vornehmsten Hauptpersonen sind gezählt worden.

230. Ueber das verspricht Pharao, er wolle ihnen, „das Mark im Lande“ zu essen geben, das ist, er wolle ihnen den allerfruchtbarsten und besten Theil in Egyptenland geben. Das hebräische Wort chalab heißt Milch, aber cheleb heißt Mark oder Fett; wie die Schrift 1 Mos. 4, 4. von Abel sagt: Er habe dem Herrn Opfer gebracht von den Erstlingen seiner Heerde, und von ihrem Fetten, das ist, von den allerfettesten Schafen: er hat von den Erstlingen seiner Heerde das Beste und Niedlichste ausgelesen. Also ist Pharao auch so mild und wohlthätig, daß er ihnen gern das Allerbeste und auserlesenste im ganzen Lande geben und auf sie wenden will. Kommt her, hat er gesagt, ich will die Kirche, Gott und seine Propheten ehren, und alles, was ich nur Gutes habe, will ich euch reichen und mittheilen. Ihr sollt die allerbeste Weide erwählen, und alles, was zu eurer Viehzucht vonnöthen sein wird. Man soll euch das Beste geben.

231. Darum ist dies ein sehr trefflich Exempel, darauf auch der Prophet Jesaia gesehen hat, da er im 49. Cap. V. 23. sagt: „Die Könige sollen deine Pfleger sein“ etc. Solche Könige sind vor Zeiten Theodosius, Constantinus und viel Andere mehr gewesen. Unsere Fürsten reißen jetzt nur Alles hinweg und berauben die Kirchen. Und wiewohl der Kirchen viele sind, die ein reiches Einkommen haben, sind doch der Herren wenige, die davon so viel herausgeben, daß die Kirchendiener und Pfarrherren damit zur Nothdurft unterhalten werden mögen. Man gibt gar zu wenig davon, das Fett aber und das Allerbeste, so die Klöster und Stifte gehabt, nehmen sie hin; welches wir zwar jetzt, da die Fürsten so große Unkosten haben müssen, einigermaßen dulden mögen. Sie sollten aber doch billig an dies Exempel denken, wenn sie ja das Fett oder das Allerbeste nicht geben wollten, daß sie doch für das Brod nicht gar die Spreu geben.

232. Wir haben ja von Gottes Gnaden einen sehr guten frommen Fürsten, und sind auch Andere mehr, die noch mit Ernst für die Kirchen sorgen, und die Pfarrherren schützen und nähren; aber deren sind fast wenig. Der Geiz und die Hoffahrt, so an den Fürstenhöfen regiert, ist eine Wurzel alles Uebels. Und wenn die Priester auch überflüssig haben, so geben sie sich gar auf den Geiz, wie vor Zeiten auch absehen, und studiren nicht so fleißig, als sie billig thun sollten. Nun wendet sich Pharao zu Joseph.

\begin{quote}
\textbf{V. 19. 20. Und gebeut ihnen also, nehmet zu euch aus Egyptenland Wagen zu euren Kindern und Weibern, und führet euren Vater, und kommt; und sehet euren Hausrath nicht an, denn die Güter des ganzen Landes Egypten sollen euer sein.}
\end{quote}

233. Dies ist eine Wiederholung des vorigen Befehls, darin Pharao dem Joseph seinen gnädigen Ring und Wagen übergeben und ihn zum Herrn und Regenten des ganzen Königreichs gesetzt hat. Wie denn kurz zuvor Joseph selbst diese Gewalt, so ihm gegeben, gerühmt hat, nämlich, daß er zum Vater des Königs gesetzt sei; wiewohl er ohne seinen Willen und Befehl nichts gethan hat. Derhalben will Pharao so viel sagen: Du weißt, daß ich mir Alles gefallen lasse, was du im ganzen Königreiche machst, und daß solches so viel gilt, als hätte ich es selbst gemacht; darum wirst du dies auch befohlen sein lassen. Du hast Befehl und Macht über alle Dinge; derhalben sage, daß sie Wagen nehmen, ob sie vielleicht Armuth oder anderer Beschwerung halben nicht herkommen könnten, oder aber daß sie zu langsam mit den kleinen Kindern und Vieh ziehen müßten. Wie sich droben, Cap. 33, 13. 14., Jakob gegen seinen Bruder Esau entschuldiget, daß er allgemach „Fuß für Fuß“ ziehen müßte wegen seines Hausgesindes und Viehes, welches ihn aufgehalten hat.

234. Eben dasselbe bedenkt hier Pharao auch, und betrachtet von selbst, was sie an ihrer Reise aufhalten oder hindern möchte. Darum wird er gesagt haben: Sie haben kleine Kinder, dazu ihre Weiber, und vielleicht auch Etliche, die schwach und matt sind, so nicht zu Fuße werden gehen können, und haben auch eben keine Wagen in Bereitschaft und werden dieselben auch nicht so leichtlich von den Nachbarn erlangen können. Derhalben nehmt Wagen mit aus Egypten, auf daß ihr diese Reise um so bequemer vollbringen möget. Also ist Mose sehr weitläufig und macht viel Worte, da er die Freigebigkeit des Königs rühmt, auf daß er dieselbe ja fleißig und genau uns zum Exempel abmalen und schmücken möge.

235. Ueber das, sagt er, weiß ich, daß ihr auch mancherlei und viel Hausgeräthe habt (wie denn der Bauern und Schäfer Hausrath zu sein pflegt); dasselbe werdet ihr nicht alle ohne Schaden hieher bringen können: darum „sehet nur euren Hausrath nicht an“ und schonet sein nicht. Nach dem hebräischen lautet es im Lateinischen also: Oculi vestri non parcant vasis etc.: Eure Augen sollen der Gefäße nicht schonen, das ist: Lasset euch nicht gereuen, ob ihr etwas hinter euch lassen solltet, das ihr nicht würdet mitnehmen können; „denn die Güter des ganzen Landes Egypten sollen euer sein.“

236. Das Wort „Gefäß“ heißt nach Art der hebräischen Sprache allerlei Hausgeräthe und Werkzeuge, oder sonst Vorrath im Hause. Hernach 1 Mos. 49, 5. werden Simeon und Levi „Gefäße der Ungerechtigkeit“ genannt. Und im Neuen Testamente Röm. 9, 22. 23. heißen unsere Leiber auch „Gefäße des Zorns“, und „Gefäße der Barmherzigkeit“; item 1 Thess. 4, 4. 5. sagt Paulus: „Ein jeglicher wisse sein Faß zu behalten in Heiligung und Ehren, nicht in der Lustseuche“, das ist, seinen Leib und seines Weibes Leib. Also nennt an diesem Orte Mose auch den ganzen Hausrath „Gefäße“. Und Pharao heißt sie, was sie nicht von dannen mitnehmen oder verkaufen können, das sollen sie nur den Andern daselbst zurücklassen.

237. Die Juden und ihre Rabbinen legen dies Stück also aus, daß ihr Verstand diesem trostwürdig ist; da im lateinischen Texte steht: Ne parcatis vasis vestris, das deuten sie also: Sehet zu, daß ihr ja nichts im Lande Canaan hinter euch lasset. Aber derselbe ist ja zumal jüdisch, und reimt sich nicht mit dem, so hernach folgt. Darum wird es viel besser also ausgelegt: „Sehet nur euren Hausrath nicht an“, daß derselbe euch nicht aufhalten oder hindern möge. Warum aber das? Ei darum, daß „die Güter des ganzen Landes Egypten sollen euer sein.“

238. Derhalben haben sie nun erstlich den Befehl von Joseph, darnach auch vom Könige. Welches Mose fleißig beschrieben hat, auf daß er damit den König rühme, darum daß er gastfrei gewesen und der Kirche große Gnade und Wohlthat erzeigt hat. Denn es ist auch eine große treffliche Wohlthat gegen die Fremdlinge gewesen; so er doch wohl wußte, daß die Leute in seinem Königreich vor denselben einen Greuel gehabt. Aber er dient damit der liebe Kirche um des Worts und Gottesdienstes willen und GOtt zu Ehren. Er ist ein rechter heiliger König.

\quellseite{448}

\section[Sechster Theil]{Sechster Theil. Wie Joseph seine Brüder zubereitet zur Rückreise; wie dieselben bei ihrem Vater Jakob ankommen, ihn verkündigen, daß Joseph noch lebe, und wie dem Jakob hiebei zu Muthe ist}

\subsection{I}

\begin{quote}
\textbf{V. 21. 22. Die Kinder Israel thaten also. Und Joseph gab ihnen Wagen nach dem Befehl Pharao, und Zehrung auf den Weg. Und gab ihnen allen, einem jeglichen, ein Feierkleid; aber Benjamin gab er dreihundert Silberlinge, und fünf Feierkleider.}
\end{quote}

239. Droben, V. 19., in des Königs Befehl ist nur der Wagen bedacht worden, Joseph aber hat ihnen über das auch noch Zehrung auf den Weg gegeben; ohne Zweifel auch auf Befehl des Königs, wiewohl dasselbe nicht so förmlich ausgedrückt war. Das hebräische Wort, daſid, heißt eigentlich, Zehrung, Proviant auf den Weg. So wird 1 Mos. 10, 9. von Nimrod gesagt: „Er war ein gewaltiger Jäger.“ Da steht im Hebräischen das Wort zajid; und 1 Mos. 27, 3. 4. sagt Isaak: „Ich will von deinem Wildpret essen.“ An diesem Ort aber heißt es, Proviant oder Zehrung, von der Jagd; und von dieser Bedeutung kann man eigentlich keine Ursache anzeigen.

240. Denn das wird fast in allen Sprachen also gehalten, daß sie die Worte auf mancherlei Weise gebrauchen, unangesehen, woher sie eigentlich ihren Ursprung haben. Also steht 1 Mos. 3, 21.: Sie machten sich chagoroth, das ist eigentlich: Feld- oder Lagergeräte. Aber dabei kann man sich an der Stelle nichts denken und aus der Etymologie kann man die Bedeutung des Wortes nicht erholen; aber der Gebrauch zeigt, daß es Röcke bedeutet.

241. Eben auf solche Weise nennen die Lateiner, durch die Figur antiphrasis, bellum, Krieg, als das gar nicht gut ist, oder, daß denen, so unerfahren sind, der Krieg däucht gut und köstlich Ding sein. Die jungen Krieger sind gar thöricht, die heißen es bellum. Also sind derselben Worte viel mehr, die darum dies oder jenes bedeuten, daß sie also in den Gebrauch gekommen sind; wiewohl es sich mit ihrer eigentlichen Ableitung nicht reimt. Haben aber nun diese Leute sich von der Jagd genährt, so reimt es sich so gar übel nicht, wenn man das Wort „Jagd“ nimmt für „Proviant“ oder „Zehrung“; wie im 132. Psalm V. 15. steht: „Ich will ihre Speise segnen“; da lautet es auch aus dem Hebräischen also: Ich will seine Jagd segnen; es ist aber so viel gesagt: Ich will das Reich Christi mit Essen, Trinken, Kleidung, und allem, was ihm zu diesem Leben vonnöthen sein wird, versorgen.

242. Nun sieht man aber auch an Joseph sonderlich, wie er sich so gar gottselig gegen seinen Vater und Brüder erzeigt hat, da er ihnen auf Befehl des Königs hinlängliche Zehrung und Proviant mitgibt, und über das auch noch Geld, dazu Futter für ihr Vieh, und alles, was ihnen vonnöthen gewesen auf dieser Reise, in das Land Canaan und aus demselben wieder heraus nach Egypten zu ziehen. Ja, er wird auch noch Jedem seine Kleidung geben und seinem Bruder Benjamin gibt er dreihundert Silberlinge; dem Vater aber schickt er zehn Esel mit Getreide und andern Gut aus Egypten beladen. Darum ist der Text an diesem Orte sonderlich wohl zu merken um des schönen Exempels willen des vortrefflichen Pflegers der Kirche und Gemeinde GOttes.

243. Die Kleider, so der Text nennet, chalipoth, Wechselkleider, haben wir in unserer deutschen Bibel verdolmetscht, „Feierkleider.“ Es sind nicht schlechte, gemeine oder tägliche Kleider gewesen, so man zu gebrauchen pflegt, wenn man tägliche Hausarbeit thut; sondern die man am Sabbath oder Feiertage zum Schmuck braucht, damit man den Leib etwas zierlicher bekleidet denn sonst.

244. Also steht im 102. Psalm V. 26. 27.: „Du hast vorhin die Erde gegründet, und die Himmel sind deiner Hände Werk. Sie werden vergehen, aber du bleibest. Sie werden alle veralten, wie ein Gewand.“ Himmel und Erde werden allgemach verwandelt, sie veralten und kommen dem Verderben immer näher, hernach aber werden sie ein neu Kleid anziehen. Denn GOtt wird nicht allein die Erde, sondern auch den Himmel selbst viel schöner machen, wenn er nun das alte Kleid ablegen wird. Das ist nur sein Werkskleid, hernach werden sie einen Ostrock und Pfingstkleid anziehen. Es hat jetzt das Licht der Sonne so eine große Klarheit und herrlichen Schein nicht, wie derselbe nach diesem Leben sein wird, und die andern Creaturen haben auch so große Kraft nicht, wie dieselbe im Anfang der Schöpfung gewesen. Wie solches die Bauern auch bezeugen, daß die Natur allgemach alt werde, und daß die Bäume beide in Wäldern und auch in den Gärten vor dreyßig oder vierzig Jahren viel eher aufgewachsen sind als jetzt, und daß jetzt die Bäume in sechs oder sieben Jahren kaum so viel zunehmen, als sie in voriger Zeit in nur Einem Jahr gethan.

245. Darum, sage ich, wird die Erde mit einem andern Kleide geschmückt werden und der Himmel und die großen Lichter werden auch klarer werden; wie im Propheten Jesaia Cap. 30, 26. steht: „Des Mondes Schein wird sein wie der Sonnen Schein, und der Sonnen Schein wird siebenmal heller sein denn jetzt.“ Einen alten Rock, der zerreißt, muß man mit einem neuen flicken.

246. Eben auf dieselbe Weise reinigt oder fegt uns GOtt auch von Sünden, und erlöst uns vom Tode; welches das Werk der sechs Tage in dieser Welt ist. Wenn aber dieselben nun verlaufen und wir in seine Ruhe eingegangen sind, alsdann wird unser zerrissen und unflätiges Kleid in das Kleid des ewigen Sabbath's verwandelt werden.

247. Was aber dazumal die Silberlinge, so Joseph seinem Bruder Benjamin gegeben hat, gegolten haben mögen, ist mir nicht bewußt. Zur Zeit Judas, des Verräthers, kann man es etlichermaßen nachrechnen, daß dazumal ein Silberling so viel werth gewesen ist als jetzt ein Joachimsthaler oder rheinischer Gulden. Und Joseph hat Benjamin vor den Andern darum mit diesem Geschenk beehrt, dieweil er sein leiblicher Bruder ist, von einerlei Vater und Mutter geboren, unter welchen die Verwandtschaft und Freundschaft etwas näher und größer zu sein pflegt. Die Deutschen heißen es, ganze Brüder.

\begin{quote}
\textbf{V. 23. Und seinem Vater sandte er dabei zehn Esel, mit Gut aus Egypten beladen, und zehn Eselinnen mit Getreide, und Brod und Speise seinem Vater auf den Weg.}
\end{quote}

248. Das hebräische Wort, chamor, heißt, Esel; daher kommt es, daß es auch der Name eines Maßes ist, darein so viel geht, als ein Esel tragen kann, ein Scheffel oder Malter. Aber das Gemeinum, auf Hebräisch athon, auf Lateinisch asina, Eselin, folgt bei den Hebräern der Ableitung nicht, daher die Lateiner und Deutschen ihre Bezeichnungen genommen haben. Was aber die Güter in Egyptenland gewesen sein, weiß ich nicht gewiß. Er wird ohne Zweifel Getreide und andere Früchte der Erde gesandt haben, damit er seinen Vater ehren, und ihn also überzeugen möchte, daß er noch am Leben wäre. Es wird aber nicht eine geringe Last der besten Güter aus Egypten gewesen sein, so viel als zehn Esel haben tragen können. Joseph wird gedacht haben, wie es denn auch die Wahrheit gewesen: Mein Vater wird sich schwerlich überreden lassen, daß ich noch leben soll, und sonderlich wenn er hören wird, daß ich zu solcher großen Ehre und Herrlichkeit soll gekommen sein. Darum richtet er so ein groß, herrlich und königlich Geschenk an, auf daß sein Vater daraus verstehen könne, daß es nicht eines geringen gemeinen Mannes, der kaum einen Esel gesandt haben würde, Geschenk sei, sondern eines Fürsten oder Herrn in Egypten. Darum sendet er ihm nicht allein Speise und Kleidung, sondern auch die allerbesten und edelsten Früchte Egyptens.

249. Mose nennt dreierlei Art der Speise: Erstlich, lechem auf Hebräisch, welches eigentlich allerlei Speise bedeutet, hier aber für Brod verstanden wird, oder gebackene Semmel. Darnach bar, welches Getreide heißt, Säcke voll Korn, und bedeutet eigentlich auserlesen Getreide, als Weizen oder Roggen. Denn das Eigenschaftswort heißt, auserlesen und rein. Dies hat er ihnen dazu gegeben, daß sie für sich auf dem Wege und daheim bei dem Vater davon haben leben sollen, und was sie übrig behalten würden, das sollten sie im Lande Canaan verkaufen, wenn sie wiederum würden auf sein und nach Egypten ziehen. Das Wort mazon wird genannt vom Schelten, oder daß man Einem etwas darreicht und gibt. Als, im 144. Psalm V. 13.: „Ihre Kammern sind voll, und geben heraus einen Vorrath nach dem andern“; damit er verstehen will allerlei Vorrath im Hause, als, Brod, Fleisch, Getränke, Butter, Käse, Eier und dergleichen mehr, daß man zu täglicher Nothdurft haben muß.

\quellseite{449}

250. Darum bringt er nun hervor allerlei Vorrath für seinen Vater und Brüder; denn sie haben unterschiedliche Familien gehabt, und haben nicht alle in einem Hause wohnen und an Einem Tische mit Jakob essen können. Benjamin war schon von sieben Kindern Vater; also sind die Andern alle auch Hausväter gewesen, da ein Jeder seinen sonderlichen Vorrath mit ihren Weibern, Kindern und Viehe gehabt 2c. Jakob, der ein Wittwer war, hat doch auch insonderheit seine eigenen Knechte und Mägde gehabt.

251. Bis daher haben wir nun gehört vom Geschenk, so Joseph seinem Vater und Brüdern gesandt, und damit er sie beehrt hat. Und hat sich der Heilige Geist dazu gern herabgelassen, daß er dasselbe mit so viel Worten beschrieben hat; wiewohl es fast schlecht und gering Ding zu sein scheint, so viel die Sache an sich selbst belanget. Denn was ist uns oder der ganzen Kirche so groß daran gelegen, daß Joseph seinem Vater Jakob Käse, Butter und allerlei Vorrath gesandt hat? Warum hat dieser große Redner, der Heilige Geist, nicht viel lieber etwas erzählen wollen vom Gebet, Versuchung und andern wunderbaren Werken solcher großen Männer?

252. Aber dies sind die närrischen Werke Gottes, davon wir nun oft gesagt haben, dazu er Lust hat, daß er seine Heiligen also nähren möge, und damit anzeigen und bezeugen, daß ihre Werke, so vor der Welt die verächtlichsten, vor ihm die köstlichsten sind; darum daß sie alle geschehen, wie St. Paulus Col. 3, 17. vermahnt, im Namen unsers Herrn Jesu Christi. Denn darum gefällt eines armen getreuen Knechts tägliche Hausarbeit unserm Herrn Gott viel besser denn aller Mönche scheußliche Gottesdienste und Fasten. Denn es sind solche Werke, die im Geiste und rechtem Glauben geschehen. Und daß die Werke etwas gelten und Gott wohlgefallen sollen, besteht nicht darin, daß sie groß und derselben viel sind, sondern es ist am Glauben und Geiste gelegen.

253. Also auch eine Hausmagd, die das Haus kehrt oder die Kühe melkt, thut Gott einen sehr angenehmen Dienst; wie Salomo sagt im Prediger am 9. Cap. V. 7. 8.: „So gehe hin und iß dein Brod mit Freuden, trink deinen Wein mit gutem Muth; denn dein Werk gefällt Gott. Laß deine Kleider immer weiß sein“ 2c. Wenn das Herz dessen gewiß ist, daß Gott wohlgefällt, was du thust an allen Orten, so sollst du daran nicht zweifeln, es sind die allerbesten und heiligsten Werke. Also, ein Schüler, der seinem Schulmeister zuhört im Glauben, soll wissen, daß er daran ein goldenes Werk thut. Wiederum aber, ein Mönch ist mit all seinem vermeinten Gottesdienste vor Gott ein Greuel; denn er geht einher und wandelt ohne Wort, ohne Gehorsam, so von Gott geboten.

\begin{quote}
\textbf{V. 24. Also ließ er seine Brüder, und sie zogen hin; und sprach zu ihnen: Zancket nicht auf dem Wege.}
\end{quote}

254. Das hebräische Wort ragaz heißt, toben, wüthen und zanken, und steht im 2. und 4. Psalm. Und hier sagt der lateinische Text: Nolite commoveri in via: Machet kein Lärm. Joseph ist ein weiser verständiger Mann gewesen, und hat wohl gewußt, daß an seinen Brüdern Fleisch und Blut war, wie auch an andern Menschen; darum hat er gedacht, es könnte leichtlich in einem Augenblicke Zank und Hader unter sie kommen, dieweil alle wohl verstanden haben, wie groß und greulich die Sünde war, sie begangen hatten, und viele unter ihnen unschuldig waren, als Benjamin und die andern Jüngsten, Issachar, Sebulon und Naphtali. Darum hat er sich besorgt, sie möchten unterwegs mit einander in Zank gerathen und einander aufrücken, was geschehen war, und sich also selbst unter einander beschuldigen und lästern. Denn also pflegen sich die, so unschuldig sind, zu entschuldigen und zu schmücken, und alle Schuld auf diejenigen zu legen, so die Sünde angefangen und angestiftet haben.

255. Darum legt er die Regel der Liebe und Vergebung dazwischen, und sagt also: Wohlan, ich habe euch die Sünde vergeben und bin mit eurer Demuth und Bekenntniß zufrieden, und habe euch nicht allein verziehen und vergeben, sondern habe euch viel Gutes und reiche Wohlthaten erzeigt. Derhalben bitte ich, ihr wollet solchem meinem Exempel nach unter einander auch friedsam sein und euch Zanks und Haders gar enthalten.

256. Dies ist ein heilsames und sehr nöthiges Gebot, damit er allerlei Ursache des Zankes abgeschnitten und aufgehoben hat, und will, daß Alles soll begraben und vergessen sein. Simeon war der vornehmste Stifter dieser Sünde; da hat Joseph gesehen, daß demselben die Andern alle zusetzen würden. Darum ermahnt er sie, daß sie sein schonen wollen, und die schwere betrübte Wunde, so schon etwas bewachsen und geheilt war, nicht wiederum erneuern oder aber einen neuen Zank erregen; vornehmlich jetzt, da nun die ganze Sache dem Vater soll vorgetragen und angezeigt werden, welcher mit großem Fleiß erforschen wird die Ursache des ganzen Handelns, so sich zugetragen hatte, bis er nun dessen gar vergewissert, daß Joseph noch am Leben und zu solcher großen Ehre und Herrlichkeit gekommen wäre; da er gewißlich auch darnach fragen würde, wie doch und durch welche Gelegenheit Joseph nach Egypten möchte gekommen sein, und wer doch die wären, die solches möchten angestiftet haben. Darum hat er gesehen, daß Simeon und Levi vor den Andern würden berüchtigt und beschuldiget werden, und daß daraus auch folgen würde, daß sie darüber billig in des Vaters Ungnade kommen würden. Derhalben hat er dasselbe auch verhüten wollen, auf daß sie mit ihrem Zanken und Hadern sich selbst und dem Vater keinen neuen Unwillen und Jammer anrichten möchten.

257. Dies ist wahrlich ein Exempel der vollkommensten und brünstigsten Liebe, dazu eines sehr reinen Herzens, daß er so gar leichtlich vergibt und aus dem Gemüth tilgt und fahren läßt alle Bitterkeit und Hader, und mit höchstem Fleiß nur damit umgeht, daß er sie unter einander und mit Jakob, ihrem Vater, auch versöhnen möge. Das sollten die Mönche auch ansehen und sich darüber verwundern, welche sonst auf nichts Anderes sehen, denn nur auf die Käse und den Vorrath Jakobs.

\subsection{II}

\begin{quote}
\textbf{V. 25. 26. Also zogen sie hinauf von Egypten, und kamen ins Land Canaan zu ihrem Vater Jakob, und verkündigten ihm, und sprachen: Joseph lebet noch, und ist ein Herr im ganzen Egyptenlande. Aber sein Herz dachte gar viel anders, denn er glaubte ihnen nicht.}
\end{quote}

258. Sie werden ohne Zweifel alle vor großer Fröhlichkeit mit diesen Worten herausgefahren sein: „Joseph lebet!“ „Joseph lebet!“ Vornemlich aber die, so unschuldig waren, werden einträchtig, da ihnen vor großer Freude die Thränen die Backen herabgeflossen, mit lauter Stimme gerufen haben: „Joseph lebet noch, und ist ein Herr im ganzen Egyptenlande.“ Und sie haben auch anders nicht können denken, denn es würde der Vater ihren Worten alsbald Glauben geben. Aber Mose setzt noch hinzu, und sagt: „Sein Herz dachte gar viel anders.“

259. Es disputiren aber die Hebräer von diesen Worten: vajaphag libo, und legen sie auf mancherlei Weise aus. Etliche haben es also gegeben: Sein Herz ist ihm gar weich geworden. Die Andern, als Hieronymus, haben es also gemacht: Es ist ihm zu Muthe gewesen, als ob er aus einem schweren tiefen Schlaf erwacht wäre. Aber nicht dünkt, ich habe die rechte Bedeutung und das Ende, dahin diese Worte gerichtet und gemeint sind, erreicht aus der Betrachtung und Vergleichung mit andern Exempeln der heiligen Schrift. Denn ich meine, daß diese Worte bedeuten die Geberde, so gewöhnlich ist, wenn wir das Haupt bewegen oder mit den Augen winken, damit anzuzeigen, daß uns nicht dünkt glaublich sein, was gesagt worden ist, und daß unser Herz weit davon ist; daß wir uns dünken lassen, daß es nicht wohl glaublich und gar unmöglich sei; den Kopf schütteln, und Hm! sagen, ei, es ist eine Mährenschein, es wird nichts draus. So lautet eigentlich das hebräische Wort, pug, nämlich, anders denken, und das Gegentheil von dem halten, das gesagt wird.

260. Rabbi Salomo kommt sehr nahe an diese Erklärung heran; denn er legt es also aus: Sein Herz ist ihm entgangen, er hielt es viel anders; es ging ihm nicht ein. Ich wollte es wohl gern also geben: Er achtete es nicht; wie im Propheten Habakuk am 1. Cap. V. 4. steht, tapug thorah, das Gesetz muß wancken 2c., als wollte der Prophet sagen: Wo ist Recht in der Welt? Hier ist es nicht. Darüber wir denn zu unsern Zeiten auch billig klagen mögen, daß die Ungerechtigkeit allenthalben in der Welt regieret und überhand genommen hat. Wie man gemeiniglich im deutschen Sprüchwort sagt: Gerechtigkeit ist todt, der Glaube leidet Noth. Dasselbe Sprüchwort ist jetzt ganz wahr geworden, und die Fürsten sagen auch: Judicium non venit ad victoriam. Das Recht kann dazu nicht kommen, daß es gewinnen möge. Das Gesetz oder Recht, so in den Büchern geschrieben steht, lehrt und ermahnt wohl, was man thun soll; aber aus dem Recht wird nichts. Und im 38. Psalm V. 9., da in der gewöhnlichen lateinischen Uebersetzung also steht: Afflictus sum, et humiliatus sum nimis: Ich bin bekümmert und zu sehr gedemüthigt, da wird auch mit demselben Wort angezeigt, daß es anders um ihn gestanden denn zuvor. Darum haben wir es in Deutschen also gegeben: „Es ist mit mir ganz anders“, denn zuvor, „und bin sehr zerstoben.“ Wie diejenigen pflegen zu seufzen und klagen, so mit dem Unglauben und Verzweiflung kämpfen, und ihr Herz nicht aufrichten noch Hoffnung schöpfen können. Item, also steht in den Klageliedern Jeremias 2, 18.: „O du Mauer der Tochter Zion, laß Tag und Nacht Thränen herab fließen, wie ein Bach; höre auch nicht auf, und dein Augapfel lasse nicht ab.“ Item Cap. 3, 49.: „Meine Augen fließen und können nicht ablassen; denn es ist kein Aufhören da.“ An denselben Stellen steht auch das hebräische Wort pugath.

\quellseite{450}

261. Also auch an dieser Stelle, da Jakob die Botschaft hört, daß sein Sohn Joseph noch lebe, schüttelt er den Kopf, und zeigt damit an, daß es ihm nicht dünke glaublich sein, darum daß er ganze zwei und zwanzig Jahre weg gewesen war, und wo er inzwischen nicht wäre umgekommen, so würde er es seinem Vater vorlängst angezeigt haben. Darum setzt Mose mit ausdrücklichen Worten hinzu: „Denn er glaubte ihnen nicht“, das ist, er hielt es für eine Lüge, die sie also müßten erdacht haben. Seine Gedanken waren, es wäre nichts; denn er glaubte es nicht. Wie Habakuk Cap. 1, 4. auch sagt vom Gesetz: „Darum gehet es gar anders, denn recht, und kann keine rechte Sache gewinnen“, als wollte er sagen: Gesetz ist nichts, das Recht ist dahin.

262. Also hat Jakob zu dieser Sache, so ihm seine Söhne vorgetragen, auch sein Herz. Es heißt aber nicht aus dem Schlaf erwachen, oder daß ihm sein Herz weich, oder schwach, oder aber verzagt geworden wäre, wie es die andern auslegen; sondern heißt eigentlich: Es wollte ihm nicht eingehen und er denkt gar das Gegentheil.

\begin{quote}
\textbf{V. 27. 28. Da sagten sie ihm alle Worte Josephs, die er zu ihnen gesagt hatte. Und da er sahe die Wagen, die ihm Joseph gesandt hatte, ihm zu führen; ward der Geist Jakobs, ihres Vaters, lebendig, und Israel sprach: Ich habe genug, daß mein Sohn Joseph noch lebet; ich will hin, und ihn sehen, ehe ich sterbe.}
\end{quote}

263. Da ist aber nun eine lange Erzählung nöthig gewesen, die ganze Historie vom Anfang zu wiederholen, wie sie waren nach Egypten hinab gekommen und wie es ihnen daselbst gegangen. Denn sie werden es dem Vater alles offenbart haben, was er zuvor nicht gewußt: wie Joseph mit ihnen gespielt und sie zweimal gefänglich eingezogen lassen; dazu, wie er mit dem Becher gereiniget; und das Andere alles miteinander. Zudem werden sie von ihrer Sünde auch nicht gar haben still schweigen können. Darum haben sie den Vater gebeten, er wolle ihnen die Sünde vergeben, sintemal ihre Handlung, wie sie Joseph verkauft, den Brüdern allen bekannt, welche eben dasselbe auch ihren Kindern und Hausgesinde werden erzählt haben; also daß sie diese ihre Sünde verschwiegen, auch dieselbe demüthig bekannt, und die große überschwängliche Gnade und Barmherzigkeit Gottes, so ihnen widerfahren war, gerühmt und gepreiset haben.

264. Auf diese Weise hat nun der Vater erfahren, wie es seinem Sohn so seltsam und wunderlich gegangen war; und kommt nun auch ein augenscheinlicher Beweis dazu, damit daßelbe klarlich bewiesen wird, nämlich, die Wagen, so Joseph gesandt hat, ihn darauf zu holen, desgleichen die großen herrlichen Geschenke, welche sonst niemand senden konnte, denn der ein Herr über Egypten wäre.

265. Derhalben ist nun sein Geist wieder lebendig geworden, und vor großer Freude, so er darüber gehabt, da er hört, daß sein Sohn noch lebe und zu solcher Ehre gekommen sei, verzeiht und vergibt er ihnen die Sünde gern, läßt allen Zorn fallen, versöhnt sich mit ihnen und nimmt sie zu Gnaden auf um Josephs willen. Er vergißt aber doch der Sünden nicht gar, so seine Söhne begangen hatten; wie dasselbe hernach wird zu sehen sein, da er denen, so daran vornehmlich schuldig gewesen, gar schwere Strafe verkündigen wird. Denn das ist eine göttliche Regel: Die Sünden sollen und müssen wir einander vergeben, auf daß uns Gott auch vergebe. Aber die Strafe nachzulassen thut nicht allezeit noth.

266. Wenn ich aber Gottes Wort also ergreife, daß ich weiß, daß mir Gott so gnädig ist, daß er mich um der Sünde willen nicht verdammen will, so schadet mir das gar nicht; wenn ich Andern zum Exempel gleichwohl darum gezüchtigt oder gestraft werde; wie der 118. Psalm V. 18. sagt: „Der Herr züchtiget mich wohl, aber er gibt mich dem Tode nicht.“

267. Also ist es Mose ergangen. Wiewohl er ein heiliger Mann war, ist er doch in das gelobte Land nicht gekommen, und hat diese Ruthe Gottes leiden müssen, daß er zu ihm sagt: Du wirst in das gute Land nicht kommen, wiewohl du große Wunder gethan hast; wie er das selbst erzählt 5 Mos. 3, 23—26., da er sagt: „Und ich bat den Herrn zu derselben Zeit, und sprach: HERR, HERR, du hast angefangen zu erzeigen deinem Knechte deine Herrlichkeit, und deine starke Hand. Laß mich gehen, und sehen das gute Land jenseit des Jordans“ 2c. „Aber der Herr sprach zu mir: Laß genug sein, sage mir davon nicht mehr.“ Denn das ist eine große Gnade und Wohlthat, daß Gott die Sünde vergibt, und mir zusagt, er wolle mir gnädig und gütig sein. Ich will dir kein Leid oder Schaden thun, will dich nicht verderben; aber daß du gleichwohl immer eine Erinnerung an deine Sünde mögest vor Augen haben und daß auch Andere durch solch Exempel unterwiesen werden, muß ich dich also züchtigen.

268. Daselbst lehrt auch das Exempel Pauli 2 Cor. 12, 7—9., der da klagt über den Pfahl, so ihm in das Fleisch gegeben war, und über den Engel des Satans, der ihn mit Fäusten schlug: „Dafür“, sagt er, „ich dreimal dem Herrn geflehet habe, daß er von mir wiche. Und er hat zu mir gesagt: Laß dir an meiner Gnade genügen.“ Höre nur auf, mich darum zu bitten, diese Trübsal ist dir gut; denn meine Kraft ist in den Schwachen mächtig: ich bin dir nicht ungnädig. Wer nun dasselbe in seinem Herzen gewißlich glaubt, und es dafür halten kann, daß er einen gnädigen und gütigen Gott hat, derselbe wird leichtlich alle Gewalt des Teufels, und der ganzen Welt Zorn und Grimm verachten.

269. So hat nun Jakob sein väterlich Herz damit bewiesen und erzeigt, daß er seinen Söhnen ihre Sünde vergibt; aber sie müssen es dennoch nicht vergessen. Darum zürnt er hernach mit ihnen und straft sie noch darüber; wiewohl er jetzt ihrer schont, da er nach dem Exempel Josephs gelinder geworden war, welches besondere Freundlichkeit und Gnade sie nun alle erfahren hatten. Und er ist nun durch ihre Beweise, so sie angezeigt haben, überwunden worden, daß er zufrieden ist, und sagt: „Ich habe genug, daß mein Sohn Joseph noch lebet.“ Ich begehre in diesem Leben gar nichts mehr, denn daß ich Joseph sehen möge, um deß willen ich bisher ganze zwei und zwanzig Jahre wohl geplagt und gekreuzigt worden bin.

270. Denn also versucht Gott seine Heiligen wunderbarer Weise. Er schickt ihnen ja wohl Anfechtung und allerlei Gefahr zu; aber doch also, daß er sie wiederum erquicke und aufrichte mit großem reichem Trost. Wie Jakob alle Hoffnung, daß sein Sohn Joseph noch lebe, hatte fallen lassen; und da er kurz zuvor Rahel verloren, und Lea und Bilha auch dahin gestorben waren, hatte er daheim in seinem Hause als ein Wittwer sein Leben in stetigem Wehklagen zugebracht und durch fremde Leute Hülfe die Haushaltung verwaltet. Das ist ihm die zwei und zwanzig Jahre über ja sauer genug geworden.

271. Darum folgt nun nach solcher langwierigen Plage ein großer trefflicher Trost, nämlich, daß Joseph nicht allein noch am Leben, sondern auch ein großer Regent sei. Das kann unser Herr Gott thun. Und dies sind unsers Herrn Gottes Werke, die man zuvor nicht verstehen kann, darum daß er uns bald im Anfange nicht sehen läßt, wie reichlich er uns trösten könne und wolle; sondern, wie er uns trösten kann 2 Mos. 33, 23. 20.: „Du wirst mir hinten nachsehen; aber mein Angesicht kannst du nicht sehen; denn kein Mensch wird leben, der mich siehet.“

\quellseite{451}
\chapter{Das sechsundvierzigste Capitel}

\section[Erster Theil]{Erster Theil. Wie Jakob nach Egypten reist und dem Herrn opfert; und wie ihm der Herr erscheint}

\begin{quote}
\textbf{V. 1. Israel zog hin mit allem, das er hatte. Und da er gen Bersaba kam, opferte er Opfer dem Gott seines Vaters Isaak.}
\end{quote}

1. In diesem Capitel steht sonst nichts, denn das Geschlechtregister Jakobs und wie er nach Egypten hinabgezogen sei; darum scheint es vor den Weisen dieser Welt von geringer Bedeutung zu sein. Denn dieselben achten dies für ein gering Ding, und daß nicht werth sei, daß es der Heilige Geist sollte haben schreiben lassen und daß es auch in der Kirche vor Propheten und Aposteln sollte gelesen werden. Aber die Gottlosen, so die Schrift und Werke des Heiligen Geistes genauer und fleißiger ansehen, dieselben sehen vornehmlich daraus und verwundern sich dessen, daß das ganze Leben der Heiligen geführt wird im Glauben an Gott, der seinen Sohn verheißen hat. Denn also hat Jakob im Glauben an Christum gelebt; darum haben auch seine Werke, wie verachtet und unsaftig sie gewesen, unserm Herrn Gott wohlgefallen, darum daß die Verheißung Alles heiligt und vor Gott Alles köstlich und theuer macht, und kann nichts so gering sein, daß in solchem Stande oder Berufe geschieht, welchen Gott geordnet und eingesetzt hat, es gefällt Gott alles wohl; und wenn es im Glauben und Wort geschieht, so übertrifft es weit aller Päpste und Mönche Wunderwerke.

\subsection{I}

2. Dasselbe ist auch der heiligen Patriarchen Lehre gewesen, die dies alles herrlich und wohl verstanden und dazu fleißig und oft in ihren Kirchen gelehrt haben. Wie im Anfang dieses Capitels das Opfer, so Jakob gethan, beschrieben wird, dabei auch ohne Zweifel eine Predigt und Unterricht oder Lehre wird geschehen sein von den Verheißungen Gottes und auch von seinem ganzen Leben. Wie aber die Mönche und Papisten die häuslichen Werke der Väter verachten: also müssen sie auch gar nichts von ihren geistlichen Uebungen, so sie in der Kirche gehabt. Denn da Mose schreibt, wie Jakob Opfer geopfert habe dem Gott seines Vaters Isaak, denken sie an nichts Anderes, denn wie er Ochsen und Schafe geschlachtet und wie dasselbe Fleisch auf dem Altar sei verbrannt worden; denn sie verachten die heilige Schrift, und sind deshalb über solcher Verachtung selbst auch verworfen worden, daß sie nicht werth sind, daß sie die Schrift lesen und verstehen könnten.

3. Aber aus dem, das droben ist gesagt worden, ist offenbar, was das für Altäre und Opfer gewesen sind, so Adam und Noah, und die andern nachfolgenden Väter gehabt haben, nämlich, daß es damit so beschaffen gewesen wie mit unsern Schulen und Kirchen, daß sie daselbst sind zu Hause gekommen zu lehren, beten und danken: da sie nicht allein Ochsen und ander Vieh geschlachtet und geopfert haben, wie die Worte lauten (und wie sich ein Weltlicher davon möchte träumen lassen, daß sie daselbst wären hinzugegangen wie Köche oder andere Metzger), sondern sie sind Hausväter, Propheten und Priester gewesen, die diesen Gebrauch gehabt, wo sie zusammen gekommen sind, daß sie daselbst Altäre gebaut und aufgerichtet, und darauf Ochsen oder Schafe geopfert, aber solches ist nicht ohne Lehre und Predigt geschehen. Denn sie haben zugleich auch dabei gepredigt, ihre Familie und Nachbarn gelehrt und unterrichtet, welche das Wort von der Verheißung des künftigen Heilandes Christi und des Samens Abrahams gehört. Und sie haben gänzlich eben dieselben Predigten gethan, die wir jetziger Zeit der Kirche und Gemeinde Gottes auch vorhalten, allein, daß jene vom künftigen Christo, der noch kommen sollte, gelehrt haben, wir aber nun von ihm also sagen: Christus ist gekommen, da sie gesagt haben: Er wird kommen.

4. Darum hat Jakob droben im 35. Cap. V. 2. seinem Hausgesinde also gepredigt: „Thut von euch die fremden Götter, so unter euch sind, und reiniget euch, und ändert eure Kleider“ etc., das ist, thut Buße, seid fromm, höret das Wort, glaubet an Gott, rufet ihn an und danket ihm. Denn Jakob war ein hoherpriester voll Heiligen Geistes, der seine Kirche und Gemeinde fleißig gelehrt, und Feiertage gehalten hat für sein Hausgesinde, Nachbarn, und Andere, so auch dazu gekommen sind, und bei demselben Gottesdienste haben sein, Predigt hören und auch mit beten wollen.

5. Da nun Mose von Jakob sagt, daß er dem Gott seines Vaters geopfert habe, soll man dasselbe also verstehen, daß er das Evangelium von dem Sohn Gottes gepredigt, und seine Zuhörer vermahnet habe, daß sie Gott fürchten, an ihn glauben und auf ihn hoffen sollten; dazu, daß sie auch ein heilig und ehrbar Leben führen, beten und Gott danken, und auf ein ewiges Leben hoffen sollten.

6. Welchem Gott hat er aber also geopfert? Antwort: Er hat keinem Götzen oder Abgott geopfert, welchen er selbst gemacht oder erfunden hat, sondern „dem Gott seines Vaters“. Denn dasjenige hat Mose mit Nachdruck hinzugefügt: Er hat dem Gott geopfert, der die Verheißung gegeben hat, welchem Isaak, Abraham, Sem, Noah auch gedient haben. Denn das ist auch der rechte Gott, der den Segen verheißen hat allen Völkern in dem gebenedeiten Samen, welcher der Schlange den Kopf zertreten sollte. Demselben hat er gedienet, von dem hat er gepredigt und ihn angerufen; darum ist das auch das höchste Werk der ersten Tafel gewesen.

7. Warum hat er sich aber das so sehr lassen angelegen sein, daß er also opfert und Gott anruft, so doch keine Noth oder Gefahr vorhanden ist? Denn es haben ihn ja beide Pharao und Joseph nach Egypten fordern lassen mit seinem ganzen Hausgesinde und sind jetzt schon die Wagen da vorhanden, desgleichen ist auch Kleidung und Getreide genug dabei? Es möchte. Denn es war nicht genug, daß er von seinem Sohn Joseph und auch vom Könige in Egypten selbst gefordert war; sondern man hat darnach auch fragen müssen, ob es auch Gott gefallen würde, daß er diese Reise thun sollte? Da wird er wahrlich mit sich selbst eine große Disputation gehalten haben, und vielleicht hat er davon auch mit seinen Söhnen, da er gen Bersaba gekommen, Unterredung gehabt, da er also wird gesagt haben: Sehet zu, lieben Söhne, was wir thun. Es ist uns ja das Land Canaan verheißen worden, daß uns Gott darin segnen wollte: wie aber, wenn mich vielleicht Gott also versuchen wollte, daß ich sollte daselbst herausziehen und in Egypten sterben? Denn ich ziehe dahin ohne Gottes Wort, ja, wider das Wort, aus eigenem Rath und Willen, ohne göttlichen Beruf. Ich sehe wohl, daß ich wider die Verheißung thue, in welcher ja das Land Canaan begriffen ist und nicht Egyptenland.

8. Mit diesen Gedanken, so scheint es, war sein und auch seiner Söhne Herz bekümmert, und er wird deshalb gesagt haben: Wir wollen ein wenig hier verziehen, und ich will Gott anrufen und fragen: Du lieber barmherziger Vater! O Gott, der du in deinem Rath wunderbärlich bist, ich begehre ja meinen Sohn Joseph zu sehen, ehe denn ich von diesem Leben abscheiden werde; ich bin aber darüber noch ungewiß, ob ich daran auch sündige, oder aber ob es dir auch also wohlgefalle: darum regiere und schicke du mein Vornehmen und meine Werke nach deinem Willen; hilf, lieber Herr Gott.

9. Also sieht man, daß er ein rein heiliges Gewissen gehabt und das nur auf Gott und seinen Willen gesehen, in dessen Furcht und Ehre er immer gewandelt hat. Aber er kann gleichwohl bei sich selbst nicht schließen, was er thun und ob er es wagen soll. Denn es ist je ein groß Ding, die Verheißung des Landes also fahren lassen, darin Abraham und Isaak, Jakobs Vater, gewohnt haben; darum hat er gedacht: Siehe, dies Land ist Abraham und mir gegeben worden und ich verlasse nun dasselbe mit meinem ganzen Hausgesinde: wie, wenn ich in Egypten erschlagen würde mit allen meinen Kindern und Nachkommen?

10. Dies sind die Anfechtungen, so die Heiligen gehabt; und wird dies alles uns zum Exempel vorgeschrieben, nämlich, wie sie geopfert und wie der Vater mit seinen Söhnen mit einträchtigem Herzen gebetet haben: Lieber HErr GOtt, du wollest uns regieren, daß wir ja unserm Fleisch und Blut nicht folgen, da wir jetzt sollen in ein fremd Land ziehen! O lieber HErr, daß wir ja nicht verderben, und daß ich mich ja daran nicht versündige und darüber in das Verderben gerathe, darum daß mich nach meinem Sohn so gar sehr verlangt und ich ihn so sehr liebe! Denn daß er die Gedanken gehabt und also geredet habe, kann man aus dem nachfolgenden Texte wohl abnehmen.

\quellseite{452}

11. Beer und scheba sind zwei unterschiedene Worte. Man kann aber aus der Geographie sehen, daß Jakob, der nach Egypten hat ziehen wollen, gen Bethel hat kommen müssen; aber er hat sich allzu sehr nach Mitternacht gewendet, da er sollte nach Morgen gezogen sein. Darum halte ich dafür, er sei den nächsten Weg nach Jerusalem gezogen, durch Etham, dadurch GOtt das Volk Israel nicht hat führen wollen, da es aus Egypten zog, um der Philister willen, sondern er hat sie durch die Wüste geführt.

\subsection{II}

\begin{quote}
\textbf{V. 2—4. Und GOtt sprach zu ihm des Nachts im Gesicht: Jakob, Jakob! Er sprach: Hie bin ich. Und er sprach: Ich bin GOtt, der GOtt deines Vaters; fürchte dich nicht, in Egypten hinab zu ziehen, denn daselbst will ich dich zum großen Volk machen. Ich will mit dir hinab in Egypten ziehen, und will auch dich herauf führen; und Joseph soll seine Hände auf deine Augen legen.}
\end{quote}

12. Hieraus ist nun wohl zu sehen, was die Anfechtung gewesen sei, so Jakob gehabt. Denn die göttlichen Verheißungen und Tröstungen sind nicht absolut, sondern relativ, das heißt, sie setzen immer und ganz gewiß die Anfechtung voraus. Denn GOtt redet von den größten Dingen, und ist kein unnützer Wäscher, der da leichtfertig und ohngefähr allerlei herausredet; sondern wenn er redet, so hält er uns die größten Dinge vor. Dieweil er denn nun zu Jakob sagt: „Fürchte dich nicht“, da du jetzt willst nach Egypten ziehen: so folgt nothwendig daraus, daß Jakob greulich muß erschrocken gewesen sein und sich sehr gefürchtet haben, sonst hätte er ja des Trostes nicht bedurft, daß ihn GOtt heißt getrost sein, und ihm verbietet, er solle sich nicht fürchten. Daran ist kein Zweifel, er wird sich gefürchtet haben, er möchte die Verheißung verlieren, dieweil dieselbe auf das Land Canaan gegründet war, welches er jetzt verlassen und daraus er ziehen wollte, und wird deshalb, da er gar irre darüber geworden, gesagt haben: Ach lieben Söhne, wir werden uns versündigen! Thut lieber also und fordert meinen Sohn Joseph hieher zu uns. Denn mit derselben Anfechtung und Furcht hat er das ganze Haus erschreckt, daß sie nichts Gewisses mehr von ihrer Reise haben schließen können. Darum hat er gesagt: Lasset uns zusammen kommen, lasset uns opfern, beten und GOtt anrufen, auf daß wir uns selbst mit dem Worte von GOttes Gnade und Verheißung stärken und trösten mögen; und lasset uns bitten, daß er unser Vornehmen nach seinem Willen regieren und schicken wolle.

13. Darum erscheint ihm nun GOtt in Gesicht, und sagt zu ihm also: „Fürchte dich nicht“ und verzage nicht, als ob du wider meine Verheißung sündigest; sondern ich heiße dich fortfahren und will deinen Zug in Egypten hiermit bestätigt haben. Ich will dich segnen, daß du deinen Sohn Joseph sehen sollst, und seine Hände sollen deine Augen zudrücken.

14. Daraus kann man aber nun wohl abnehmen, wiewohl GOtt in seinen Verheißungen nicht lügen oder fehlen kann, daß er sich doch das vorbehalt, daß er dieselben erfüllt und zur That macht, und das, so er verheißet hat, also ausrichtet, daß Alles ganz widersinnig und dem Wort zuwider scheint. Deshalb gehört nun dazu ein solch Herz, das da weise sei und das den Sinn und Rath GOttes verstehen könne. Denn GOtt hält seine Verheißung treulich, kann auch nimmermehr geändert werden, und läßt doch das Gegentheil vorfallen, damit er uns versucht.

15. Dies ist unsers HErrn GOttes wunderbares Regiment, davon Paulus mit großer Verwunderung ruft, Röm. 11, 33. 34.: „O welch eine Tiefe des Reichthums, beide der Weisheit und Erkenntniß GOttes! Wie gar unbegreiflich sind seine Gerichte, und unerforschlich seine Wege! Denn wer hat des HErrn Sinn erkannt? oder wer ist sein Rathgeber gewesen?“ 2c. Denn ich bin getauft, bin absolvirt von meinen Sünden aus göttlicher Kraft, Matth. 16, 19.: „Alles, was du auf Erden lösen wirst, soll auch im Himmel los sein.“ Ich habe eine solche Verheißung, die mir nicht fehlen kann. Der Seelenhirte tröstet mich und sagt: Du bist von GOtt angenommen, daß du sollst sein Kind sein, ja, du sollst ein Erbe des ewigen Lebens sein. Nichts desto weniger aber wirft mich vielleicht GOtt in eine schwere Anfechtung, macht mich zweifelhaft, bekümmert und verzagt, da er mir solche Dinge läßt vorfallen, die seiner Verheißung stracks zuwider sind. Wie er in den Propheten zuvor verkündigt: Ich will den König zu Babel wider euch senden. Item, der Satan streitet wider uns mit allerlei Waffen und List. Da fühle und sehe ich schon, daß mich GOtt straft und mich in sehr große schwere Noth und Gefahr kommen läßt; doch erbarmt er sich mein, denn ich habe seine Verheißung. Wer kann aber diese Dinge, die so gar wider einander sind, verstehen, oder dieselben mit einander ausgleichen? Ich soll hoffen auf das ewige Leben, und muß jämmerlich umkommen und dahin sterben: dünkt dich denn, daß GOtt auf solche Weise seine Verheißung erfülle?

16. Darum muß man fleißig merken, daß eine andere Lehre sei von GOtt nach seinen Verheißungen, als sonst, wenn man lehrt von der Gerechtigkeit des Gesetzes. Das Evangelium predigt also: Du sollst sterben und im Sterben gleichwohl leben; du sollst leben, und im Leben doch sterben. Dies ist eine sonderliche Weisheit, so die Vernunft und allen menschlichen Verstand weit übertrifft, welche die Juristen und Papisten nicht verstehen. Also gibt GOtt dem Abraham zugleich mit der Verheißung den Samen Isaak, in dem alle Völker der Erden sollten gesegnet werden; alsbald aber gebietet er ihm: Gehe hin, und schlachte deinen Sohn; er ist eben derselbe GOtt, der ihm den Sohn gegeben hat und ihn denselben auch wiederum nimmt.

17. Das ist nun eine sonderliche Lehre für die Kirche, nämlich, daß man wissen soll, daß GOtt ein allmächtiger Schöpfer ist, der zu diesem Leben schafft, und darnach wiederum zerbricht, was er geschaffen hat, auf daß er es zum andern Leben wiederum lebendig mache. Das ist die Verheißung, da er sagt: Du sollst leben. Aber die Erfahrung, so der Verheißung stracks zuwider ist, lautet also: Du sollst sterben. Heißt denn das die Verheißung halten? Es ist wahrlich sehr schwer, solche Dinge zu lehren und zu glauben, die stracks wider einander sind, und ist niemand, der in demselben Glauben beständig bleiben könne, denn allein der, so des HErrn Sinn versteht; wie Paulus sagt 1 Cor. 2, 16.: „Wer hat des HErrn Sinn erkannt? oder wer will ihn unterweisen? Wir aber haben Christi Sinn.“

18. Und reimt sich daher wohl, daß Jakob in diesem Opfer und Gebet gedacht und sich vorgestellt habe das Exempel seines Vaters Isaak, und daß er also geredet habe: HErr GOtt, wie du meinen Vater Isaak aus der Asche und Brandopfer erlöset hast, also regiere mich und hilf mir auch! Siehe, ich gehe hin in großem Zittern und Angst, dieweil ich keinen gewissen Beruf habe. Denn das ist ein groß schwer Ding, aus eigenem Vermessen wider die Verheißung etwas vorzunehmen zu thun. Und wäre diese Entschuldigung nicht genug gewesen, wenn er gesagt hätte: Er wollte dies also aus gutem Wohlmeinen thun, aus väterlicher Liebe und Eifer zu seinem Sohne; sintemal er dazu keinen Befehl hat von GOtt. Darum erscheint ihm nun GOtt, da er opfert, und sagt zu ihm also: „Fürchte dich nicht! Was fürchtest du dich? Zeuch hinab, ich will dich erlösen und wieder herausführen.“

19. Auf diese Weise sehen wir, daß es um den göttlichen Trost also beschaffen ist, daß wo derselbe gegeben wird, daselbst Traurigkeit sein müsse; darum muß man solchen Trost im Glauben aufnehmen und fassen. Und man kann hieraus auch gewißlich schließen und abnehmen, daß Jakob auch wird betrübt gewesen sein bis in den Tod. Denn GOtt richtet niemand auf, stärkt oder tröstet auch niemand, denn allein die, so betrübt sind, so jetzt sterben sollen, mit denen es gar verloren ist. Denn das Wort des Lebens und Heils gehört denen, die in Angst und Verzweiflung sind, zu welchen recht gesagt wird: Du fürchtest dich und dein Gewissen martert dich, der Teufel mit seinem Stachel und dazu das Fleisch plagt dich auch; sei getrost, verzage nicht, GOtt zürnt nicht mit dir 2c.

20. Denn also ermahnt GOtt an diesem Orte den Jakob nicht allein, daß er hinab nach Egypten ziehen soll, und tröstet ihn dazu, und erlöst ihn von der Angst, so er in seinem Gewissen gehabt; sondern hat ihn auch gleichsam überladen mit einer reichen Verheißung, und gibt ihm dazu eine feine herrliche Glosse auf die vorige Verheißung, und sagt also: Daß ich dir das Land Canaan verheißen habe, mußt du so gar streng und genau nicht verstehen, als ob vonnöthen wäre, daß du eben allezeit darin bleiben müßtest. Ich bin nicht allein ein Gott der Juden, sondern auch der Heiden: ich will mir auch aus dem Egypten eine Kirche sammeln. Wie das Volk hernach ist zerstreut worden in Babylonien und unter die Assyrer, auf daß sie Viele möchten zum Evangelium berufen. Da, sagt nun Gott, will ich viel Leute sammeln durch dich und deinen Samen, und will geben, daß sie sich zu einer Kirche bekehren sollen und sich zu euch halten. Denn die Verheißung ist wohl euch allein gegeben worden, aber ihr sollt euch derselben nicht allein annehmen. Denn das taugt nicht, daß ihr also sagen wollet: Wir sind allein beschnitten, darum sind wir auch allein Gottes Volk. Ich will die Heiden ebensowohl berufen, als die Juden, daß sie auch herzu kommen, und sich zu mir bekehren sollen. Darum sündigst du daran nicht, daß du nach Egypten ziehst; sondern fahre nur fort, daß auch die Egypter und andere Leute mehr selig mögen werden; „denn daselbst will ich dich zum großen Volk machen.“

\quellseite{453}

21. Dies ist aber nun der vorigen Verheißung nicht zuwider, da Gott gesagt hatte: Ich will dir dies Land geben; welches nicht allein vom Lande Canaan kann verstanden werden; denn Gott ist ein Herr der ganzen Welt. Und ist wahrlich dies ein wunderbarer Rath Gottes gewesen, daß durch die Juden, wenn sie also hin und her zerstreut worden, die Heiden auch sollten versammelt und zu Gott bekehrt werden. Denn also hat Daniel der Perser Könige bekehrt, Joseph Faraonen, Christus und die Apostel die ganze Welt; wie St. Paulus sagt Röm. 11, 11.: „Aus ihrem Fall ist den Heiden das Heil widerfahren“; item V. 12.: „Ihr Schade ist der Heiden Reichthum.“ So oft Gott die Juden um ihrer Sünde willen gestraft, ist solche Strafe und Sünde eine Ursache des Heils vieler anderer Völker gewesen.

22. Derhalben, da Jakob meint, daß die Verheißung und sein Beruf nur auf das Land Canaan sollte gehen, hört er vom Herrn, daß er zu ihm sagt: „Ich will dich daselbst zum großen Volk machen“, ich will bei dir sein. Wie Christus seiner Kirche auch verheißt, da er Matth. 28, 20. sagt: „Siehe, ich bin bei euch alle Tage, bis an der Welt Ende“; aber doch nicht ohne Anfechtung. Und Gott sagt weiter zu Jakob: Du darfst dich darüber nicht bekümmern, wie die Verheißung könne erfüllt werden, deweil hier befohlen wird, nach Egypten zu ziehen zu deinem Sohn, nach welchem dich so viel Jahre herzlich verlangt hat; denn „ich will dich wiederum heraufführen“; und soll gleich viel sein, ich führe dich todt oder lebendig heraus. Denn du wirst auch in Egypten leben, wenn du schon stirbst, sintemal ich ein Gott bin der Lebendigen und nicht der Todten. „Und“ endlich „soll Joseph seine Hände auf deine Augen legen.“ Du sollst einen schönen Tod haben, sollst in gutem Frieden eines seinen, sanften, und stillen Todes sterben.

23. Es ist aber ein wunderbärlich Ding, daß er ihn heißt nach Egypten hinab ziehen, da sein Sohn ihm, wenn er sterben würde, die Augen sollte zudrücken, und daß er ihn doch wiederum wollte herauf führen mit seinen Nachkommen, wenn er würde gestorben sein. Darum ist Gott des Patriarchen Jakob Gott, der da lebt, wenn er auch schon jetzt gestorben ist, und da Jakob von diesem Leben abgeschieden, ist er nicht gestorben.

24. Es erhellet aber aus dieser großen und sehr reichen Verheißung, daß Jakob bei dieser Reise mancherlei gedacht und gefühlt hat und schrecklich verwirret war. Darum verheißt ihm Gott, daß er ihn wolle groß machen in Egypten und bei ihm sein, ihm helfen und regieren, und daß er endlich sein sanft sterben solle, also daß Joseph die Hände auf seine Augen legen wird.

25. Denn also pflegen die sterbenden Menschen zu verscheiden, mit gebrochenen und offenen Augen, da pflegen dann die Verwandten oder Freunde, so dabei sind, ihnen die Augen zuzudrücken, welches bei den lieben Vätern der vornehmsten Leute in ihrem Geschlecht Amt gewesen ist, und nicht der alten Weiber, oder sonst anderer Weiber, so dazu gedinget werden, daß sie über die Todten Leid tragen; sondern Väter oder Mutter haben ihren Kindern die Augen zugedrückt um der Hoffnung und Ehre willen der zukünftigen Auferstehung. Und ist ein feiner, löblicher und ehrlicher Gebrauch gewesen und eine feine Ceremonie oder Ordnung, daß entweder die Eltern den Kindern, oder aber die Kinder und vornehmlich der Erstgeborne den Eltern die Augen zugedrückt. Es ist eine sonderliche und die letzte Wohltat, so den Verstorbenen in ihrem letzten Stündlein kann erzeigt werden. Und ist dieser Gebrauch ohne Zweifel von den Vätern hergekommen, und ist eine solche Weise, die man billig in der christlichen Kirche behalten soll um der Hoffnung willen der Auferstehlichkeit, so uns verheißen ist, und der Auferstehung der Todten. Wie in Deutschland an etlichen Orten noch der Gebrauch ist, wo man die Todten zur Erde bestattet, daß da die nächsten Freunde zu allererst Erde in das Grab werfen und mit solchem Zeichen bezeugen die Hoffnung und Erwartung der Auferstehung.

26. Auf diese Weise beschreibt Mose ganz fleißig nicht allein die häuslichen Wercke, so Jakob gethan, sondern auch den Gottesdienst und die geistliche Uebung, so er gehabt, sein Gebet, und wie er Gott gelobt und ihm gedankt hat. Welche Uebungen die Wunderwercke der päpstlichen Heiligen weit übertreffen.

27. Wir aber werden aus diesem Exempel gelehrt, daß man in der christlichen Kirche nichts annehmen oder anfangen soll wider oder ohne Gottes Wort. Derhalben sollen wir unser Leben also anrichten, daß wir dessen mögen gewiß sein, daß wir nach der Regel des göttlichen Worts wandeln, wir machen oder schlafen. Und endlich, daß wir wissen, daß alles, was wir thun, aus Gottes Befehl geschehe, welcher uns dazu berufen und gesegnet habe. Denn das Leben der Menschen auf Erden ist armselig, ist dem Teufel, dem Tode und vielen unzähligen Anfechtungen unterworfen, und ist voll des schwersten Jammers, Mühe und Arbeit. Darum müssen die Gottseligen gerüstet sein mit Gottes Wort wider den Satan, die Sünde, und allerlei Noth und Gefährlichkeit. Und wo sich die Papisten nach dieser Lehre gehalten hätten, so wäre die Welt nicht so voll menschlicher Satzung geworden, damit sie wie von einer Sintfluth überfluthet ist, und die Kirche wäre auch nicht mit so viel greulichem Irrthum besudelt worden.

28. Darum ist es eine große Gabe Gottes, daß man im Lichte des Worts und göttlichem Beruf lebe. Denn das ist die rechte höchste goldene Krone, so da scheint auf den Häuptern der Heiligen, wie man sie gemeiniglich pfleget zu malen, nämlich, das Wort, welches Tag und Nacht den Lauf unsers Lebens richtet, wie David sagt in 119. Psalm V. 105.: „Dein Wort ist meines Fußes Leuchte“, und 2 Petr. 1, 19.: „Ihr thut wohl, daß ihr auf das Wort achtet, als auf ein Licht, das da scheinet in einem dunkeln Ort, bis der Tag anbreche etc.

29. Die Vernunft liegt gar in Finsterniß, darum thut ihr das Licht des göttlichen Worts noth, daß sie führen und regieren muß. Denn also thut Jakob: da er sich fürchtet, diese Reise wider Gottes Befehl zu thun, bleibt er eine Zeitlang zu Bersaba, bis daß er opfere, bete und Gottes Wort höre. Und da er das gehört hat, ist er nun darüber fröhlich; denn er weiß nun, daß diese Reise durch Gottes Segen gesegnet und bestätiget ist.

30. Dies müssen wir fleißig merken, auf daß wir wissen, daß man nichts leichtfertig thun, rathen oder vornehmen soll, wir haben denn Gottes Wort; es sei gleich in Buche geschrieben, oder von man etwas Sonderliches soll vornehmen, daß solches von den Eltern oder Obrigkeit befohlen werden möge, daß wir ja wider den Gehorsam und Ehre Gottes gar nichts thun. Man darf den Teufel nicht zu Gaste bitten.

\section[Zweiter Theil]{Zweiter Theil. Das Register derer, so mit Jakob nach Egypten gezogen}

\begin{quote}
\textbf{V. 5—7. Da machte sich Jakob auf von Bersaba; und die Kinder Israel führten Jakob, ihren Vater, mit ihren Kindern und Weibern, auf den Wagen, die Pharao gesandt hatte, ihn zu führen; und nahmen ihr Vieh und Habe, die sie im Lande Canaan erworben hatten; und kamen also in Egypten, Jakob und all sein Same mit ihm. Seine Kinder, und seine Kindeskinder mit ihm, seine Töchter, und seine Kindstöchter, und all sein Same, die brachte er mit sich in Egypten.}
\end{quote}

31. Wie Jakobs Gebet erhört und Gottes Wort erfüllt worden sei, wird hier beschrieben. Es fallen aber hier auch Fragen ein, so zur Grammatik und Arithmetic gehören. Denn Mose erzählt, wie Jakob nach Egypten gekommen sei mit alle seinem Hausgesinde, Gut und Habe, so er im Lande Canaan erworben hatte. Wie droben, 1 Mos. 12, 5, von Abraham und Lot auch gesagt wird: „Also nahm Abraham sein Weib Sarai, und Lot, seines Bruders Sohn, mit aller ihrer Habe, die sie gewonnen hatten, und Seelen, die sie gezeuget hatten in Haran“ etc. Hier aber nennt Mose insgemein „seine Habe“, nicht allein sein Gut, sondern auch seine Knechte und Mägde, die sie zusammen pflegten zu kaufen und zu verkaufen, wie jetzt in der Türkei geschieht. Und es scheint wahrlich kein geringer Haufe gewesen zu sein. Denn die zwölf Patriarchen haben Kinder, Kindeskinder, Knechte, Mägde und vielleicht auch wohl Anverwandte gehabt.

\quellseite{454}

32. Es wird aber gefragt: Warum Mose sage, daß Jakob Töchter und Kindestöchter mit sich genommen habe, so doch im folgenden Geschlechtsregister nur zwei Töchter aufgezählt werden, nämlich Dina, und Serah, die Tochter Assers, Jakobs Enkelin. Wie hat er denn also geredet in der Mehrzahl von einer einigen Tochter, als ob ihrer viele wären, so doch über diese in der ganzen Zahl ihrer nicht mehr gewesen? Antwort: Dies ist also der heiligen Schrift Gebrauch, daß sie bisweilen synedochisch die Einzahl braucht, wenn sie von Vielen redet, wie 2 Mos. 4, 22. steht: „Israel ist mein erstgeborner Sohn“; item: Ephraim sündigt und betet die Kälber an; welches doch vom ganzen Volke und Stamme Ephraim zu verstehen ist. Also sagt man in der lateinischen Sprache auch: Romanus proelio victor, das ist: Die Römer haben in Streit den Sieg behalten.

33. Das Andere ist aber etwas seltener und geschieht nicht so oft, nämlich, daß man von Vielen redet, wo man nur Einen will verstanden haben; wie hernach, Cap. 46, V. 23., steht: „Die Kinder Dan: Husim“, da nur Einer erwähnt wird, so er doch vorgegeben, daß er von Vielen reden wolle. Und an diesem Orte steht auch: „Seine Kindstöchter“, welches allein von Serah geredet ist. Dies ist der Gebrauch, und eine sonderliche Freiheit, so die heilige Schrift hat, nämlich, daß sie etwa eine ganze Linie, Geschlechtsregister oder Geburt nennt, wenn auch nicht mehr, denn nur Einer aufgezählt wird.

34. Augustin meint, man müsse dadurch verstehen die Weiber der Söhne Jakobs; aber der Verstand reimt sich mit dem Texte nicht, welcher ja ausdrücklich von den Töchtern und seiner Kinder Töchtern sagt. Nun folgen die Namen der Kinder Jakobs; dieselben wollen wir kürzlich überlaufen.

\begin{quote}
\textbf{V. 8–15. Dies sind die Namen der Kinder Israel, die in Egypten kamen: Jakob und seine Söhne. Der erstgeborne Jakobs Sohn, Ruben. Die Kinder Ruben: Hanoch, Pallu, Hezron und Charmi. Die Kinder Simeon: Jemuel, Jamin, Ohad, Jachin, Zohar, und Saul, der Sohn von dem canaanäischen Weibe. Die Kinder Levi: Gerson, Kahath und Merari. Die Kinder Juda: Ger, Onan, Sela, Perez und Serah. Aber Ger und Onan waren gestorben im Lande Canaan. Die Kinder aber Perez: Hezron und Hamul. Die Kinder Isaschar: Thola, Phua, Job und Simron. Die Kinder Sebulon: Sered, Elon und Jahleel. Das sind die Kinder von Lea, die sie Jakob gebar in Mesopotamien, mit seiner Tochter Dina. Die machen allzumal mit Söhnen und Töchtern, drei und dreißig Seelen.}
\end{quote}

35. Es wird nur von Saul hier gesagt, daß er der Sohn eines canaanäischen Weibes gewesen sei, und haben dafür die Rabbinen der Juden mancherlei Ursachen erdichtet. Ich gebe aber ihrem Gedicht keinen Glauben, weiß auch nicht, was man hierzu sagen soll. Etliche sagen, daß dies canaanäische Weib Dina sein soll, so von Sichem, dem Heviter, beschlafen worden ist: dieselbe soll darnach einem canaanäischen Manne sein gegeben worden; endlich aber soll sie Simeon zum Weibe angenommen haben, daß sie seinen Hause hat sollen vorstehen und dasselbe regieren, welches mir nicht wahrscheinlich dünkt.

36. Wenn man aber etwas darauf dichten mag, so wollte ich lieber das dazu sagen, daß dies eine ehrliche schöne Matrone gewesen sei, welche Simeon lieb gehabt und die er auch zum Weibe genommen, es habe gleich sein voriges Weib noch gelebt, oder aber sei todt gewesen, welche bereits zuvor von einem andern Manne diesen Sohn Saul gehabt; also daß derselbe in diesem Stamme fremd und von beiderseits canaanäischen Eltern geboren sei. Und dafür habe ich vornehmlich diese Ursache, daß Saul in dieser Zahl der drei und dreißig Seelen übrig ist; denn wo du sie alle nach wirst zählen, so sind im Texte vier und dreißig Personen.

37. Und ich mag dem Lyra hier auch nicht nachfolgen, der Dina gar ausschließt wider die klaren Worte im Texte. Und über das ist Dina, Jakobs Tochter, eine Israelitin und keine Cananitin gewesen. Wo man aber eine Person soll auslassen und von den andern absondern, so halte ich es dafür, daß dieser Saul vielmehr müsse ausgeschlossen werden. Mich wohl ich nichts Gewisses dazu sagen will; denn es ist im ganzen ersten Buch Mose kein Ort, der von den Auslegern auf so mancherlei Weise hin und her gezogen worden sei, als dieser ist, darüber sie sich alle heftig bemüht und den Kopf zerbrochen haben, nämlich, Lyra, Augustin, Burgensis und Hieronymus.

\begin{quote}
\textbf{V. 16–18. Die Kinder Gad: Ziphion, Haggi, Suni, Ezbon, Eri, Arodi und Areli. Die Kinder Asser: Jemna, Jesua, Jesui, Bria, und Serah, ihre Schwester. Aber die Kinder Bria: Heber und Malchiel. Das sind die Kinder von Silpa, die Laban gab Lea, seiner Tochter, und gebar Jakob diese sechzehn Seelen.}
\end{quote}

38. Serah ist eine Tochter Asser. Es werden in dieser Linie nur zwei Weiber aufgezählt, und eine jegliche wird für sich selbst allein genannt „Töchter“, in der Mehrzahl, als ob man von Vielen redete. Dasselbe geschieht nach dem Gebrauch der Schrift. Wenn aber zu den vorigen vier und dreißig Seelen die sechzehn gethan werden von dem andern Weibe, so bringen sie zusammen fünfzig Seelen. Nun folgen die Söhne und der Söhne Kinder von Jakobs vornehmstem Weibe Rahel.

\begin{quote}
\textbf{V. 19–27. Die Kinder Rahel, Jakobs Weibes, Joseph und Benjamin. Und Joseph wurden geboren in Egyptenland Manasse und Ephraim, die ihm gebar Asnath, die Tochter Potipheras, des Priesters zu On. Die Kinder Benjamin: Bela, Becher, Asbel, Gera, Naaman, Ehi, Ros, Muppim, Huppim und Ard. Das sind die Kinder von Rahel, die Jakob geboren sind; allzusamt vierzehn Seelen. Die Kinder Dan: Husim. Die Kinder Naphtali: Jahzeel, Guni, Jezer und Sillem. Das sind die Kinder Bilha, die Laban seiner Tochter Rahel gab, und gebar Jakob die sieben Seelen. Alle Seelen, die mit Jakob in Egypten kamen, die aus seinen Lenden kommen waren (ausgenommen die Weiber seiner Kinder), sind alle zusammen sechs und sechzig Seelen. Und die Kinder Joseph, die in Egypten geboren sind, waren zwei Seelen; also, daß alle Seelen des Hauses Jakob, die in Egypten kamen, waren siebzig.}
\end{quote}

39. Da siehst du, wie der Heilige Geist der Zahl so große Ehre thut, daß er sie Jakobs Weib nennt, welches er von Lea, Silpa und Bilha nicht gesagt hat. Denn sie ist die rechte Frau im Hause gewesen, welche Jakob auch herzlich lieb gehabt, begehrt und ihm auch vertraut worden ist: aber Laban hat durch einen gottlosen Tausch ihn betrogen, wie wir droben, Cap. 29, 23., gehört haben.

40. Er sagt aber hier nun abermals: „die Kinder Dan“, so er doch nur den einigen Sohn Husim gehabt hat, darum daß in der heiligen Schrift die Geburt oder Geschlechter in der Mehrzahl genannt werden, als reden man von Vielen, da doch nur Einer erwähnt wird.

41. Und bis daher währet das Register der Namen und die Aufzählung der Personen, so nach Egypten hinab gezogen sind. Aber die Rechnung ist so gar verworren, daß sie kein Ausleger hat können zurechtbringen. Jedoch hat es Lyra am nächsten getroffen; allein, daß er Dina ausschließt, welche doch der Text klarlich und ausdrücklich mitzählt. Denn wo du sie alle zusammen rechnest, so ist Eine Person übrig, fünfmal im Texte vier und dreißig gesetzt werden, so doch Mose allein drei und dreißig zählt. Ich halte aber dafür, daß man viel lieber Saul sollte auslassen denn Dina; wie ich zuvor auch gesagt habe.

42. Nun folgt noch eine andere Zahl, die auch nicht weniger verworren ist, da er sagt, daß sechs und sechzig Seelen mit Jakob nach Egypten gekommen sind. Aber damit wir dieselbe Zahl verstehen, erscheint es am bequemsten, daß wir diese Rechnung theilen, also, daß das erste Register die betreffe, so von Jakob geboren sind, darin Joseph mit seinen Söhnen begriffen ist; das andere aber betreffe diejenigen, so nach Egypten gezogen, da Joseph mit seinen Kindern ausgeschlossen wird, deßgleichen auch Ger und Onan, und bleiben noch übrig fünf und sechzig Personen.

43. Von Lea sind hinein gekommen zwei und dreißig Seelen, wenn Saul davon ausbleibt; von der Silpa sechzehn; von Rahel elf (denn Joseph und seine zwei Söhne waren schon in Egypten); von Bilha sieben. Wenn man nun diese alle zusammengezählt werden, machen es sechs und sechzig Personen. Wo wollen wir aber die Andern hernehmen, daß sie siebenzig voll werden? Wenn Joseph und seine zwei Söhne dazu gethan werden, so machen es neun und siebzig. Jakob aber, der Letzte, wird der siebenzigste.

\quellseite{455}

44. Darum muß man den Text disjunctive und copulative verstehen, das ist, daß man Theile und wieder zusammenzähle. Alle Seelen, so aus Jakobs Lenden gekommen waren, ihn selbst ausgenommen, sind sechs und sechzig. Denn Jakob ist ja selbst aus seinen Lenden nicht gekommen. Joseph aber und seine Söhne, wiewohl sie von ihm geboren sind, ziehen doch nicht mit in Egypten; desgleichen da Er und Onan in Canaan gestorben sind, ist Saul an ihre Statt gekommen. Er wird aber in der ersten Rechnung ausgelassen, und werden die zwei mitgezählt, darum daß sie zu den Nachkommen Jakobs auch mit gehören.

45. Die siebenzig Dolmetscher der Bibel haben verdolmetscht fünf und siebenzig Seelen etc., und Stephanus, Ap. Gesch. 7, 14., führt auch daraus dieselbe Zahl an. Das ist aber ein offenbarer Irrthum oder Bosheit, daß jene für die zwei Seelen, so Joseph in Egypten sind geboren worden, sich unterstanden haben, neun zu setzen. Und 5 Mos. 10, 22. werden ausdrücklich nur siebenzig Seelen erwähnt. Augustinus versucht sie zu entschuldigen, und hält es dafür, daß ein sonderlich Geheimniß oder heimliche Weisheit darunter verborgen sei. Ich weiß aber nicht, ob dieser Irrthum aufgehen oder herkommen sei von den Dolmetschern, oder aber aus Unwissenheit der Bücherschreiber, oder daß es sonst etwa möchte verfälscht sein. Der Text sagt klärlich von den Seelen, so nach Egypten gekommen sind. Und obwohl Ephraim im andern wohlfeilen Jahr geboren worden, Jakob aber im andern Jahr der Theurung zu seinem Sohne gezogen ist, so ist es doch unmöglich, daß ein Knabe von acht Jahren sollte ein Weib oder Kindeskinder haben, auf welche er doch gesehen haben.

46. Derhalben scheint es, daß die siebenzig Dolmetscher allzufrüh gewesen sind mit ihrem Dolmetschen; oder aber, es wird etwa ein böser Bube ihr Buch verfälscht haben. Denn man muß zu diesem Irrthume nicht still schweigen. Und sind auch der Seelen noch mehr, darin sie auch gezählet haben entweder aus Unwissenheit oder aber aus Schalkheit. Darum haben sie bei mir ein klein Ansehen und gelten bei mir so gar viel nicht.

47. Wir wollen aber die gemeine Zahl behalten, die mit dem fünften Buch Mose übereinkommt. Und damit wir uns so viel leichter aus diesem verworrenen Irrthume helfen können, habe ich gesagt, daß man zweierlei Rechnung soll machen, die eine der Geburt, daß ich es also nenne, da zwei sterben, Ger und Onan, und werden doch mit gezählt und ziehen gleichwohl nicht nach Egypten; wiewohl sie auch aus den Lenden Jakobs gekommen waren. Die andere Rechnung ist mit den sechs und sechzig und siebenzig Seelen, so Jakob und Joseph mit seinen zwei Söhnen dazu gethan werden.

48. Münzer folgt der Juden Rabbinen nach, und sagt: Man soll noch eine Person dazu thun, welche in Egypten geboren sei, nämlich Jochebed, die Tochter Levi und Mutter Mosis. Aber derselbe ist eine kühne und frevele Erdichtung der Rabbinen, die uns gar nichts angehet.

49. Wozu war aber in diesem Geschlechtsregister solche Genauigkeit vonnöthen? Oder: Warum hat Gott haben wollen, daß wir uns dieser Zahl halben also üben und bemühen sollten? welches gar nichts nütze zu sein scheint; denn was kümmert es mich, es ziehen gleich sechzig oder siebenzig Seelen nach Egypten? Ich habe oft gesagt und soll auch allezeit den Leuten eingeschärft werden, daß Gott hat haben wollen, daß das Israelitische Volk einen gewissen und herrlichen Ursprung haben sollte um des Herrn Christi willen, der aus diesem Volke sollte geboren werden. Darum wird mit solchem großen Fleiß die Ordnung oder Linie der Väter in der ganzen Historie der heiligen Schrift, in den Büchern der Richter, der Könige und der Propheten erzählt, auf daß jedermann wissen soll, wer des Herrn Christi Voreltern gewesen sind, daß er nicht an einem unbekannten und ungeachteten Orte geboren, sondern von Gott selbst gesandt und berufen sei, und in dieser Welt himmlische Lehre und göttliche Wunderwerke gethan habe. Denn er ist ja nicht in einem Geschlechte unehrlicher loser Leute und Landstreicher geboren worden, wie der Papst und Mahomet, sondern aus dem allerberühmtesten Volke.

50. Danach ist dies um unsertwillen geschrieben, auf daß wir lernen sollen, wie Gott für seine Heiligen und Gläubigen so fleißig sorgt, daß er auch nicht Eine Person ausläßt, daß er sie nicht sollte mit höchstem Fleiß und Treue auch wollen mitgezählt haben. Wie Christus sagt von unser aller und eines Jeden Haar insonderheit, Matth. 10, 30.: „Es sind auch eure Haare auf dem Haupt alle gezählet.“

51. Derhalben ist Christus die vornehmste Ursache, daß diese Rechnung so genau beschrieben wird. Die andere aber ist, daß die Gottseligen daraus sollen getröstet werden, welchen Gott damit will Muth machen, daß sie sich vor dem Papst, Türken und Kaiser nicht sollen fürchten, sintemal sie also genau beschrieben sind in Gottes Händen, daß nicht allein kein Finger, kein Ohr, sondern auch kein Haar von ihrem Haupte fallen und umkommen soll wider den Willen ihres himmlischen Vaters. Wie der schöne liebliche Spruch im Propheten Jesaia am 49. Cap. V. 14. 15. 16. lautet: „Zion aber spricht: Der Herr hat mich verlassen, der Herr hat mein vergessen“; da antwortet GOtt: „Nein, ich habe dein nicht vergessen: „Kann auch ein Weib ihres Kindleins vergessen, daß sie sich nicht erbarme über den Sohn ihres Leibes? Das wird sie wahrlich nicht thun, die Natur läßt solches nicht zu. Oder wenn sie nun schon eine Hure oder lose Bübin wäre, und ihres Kindes vergessen würde, so will ich doch dein nicht vergessen. Siehe, in die Hände hab ich dich gezeichnet; deine Mauern sind immerdar vor mir“ etc. Als ob er sagte: Alles, was ich thue im Himmel und auf Erden, ist alles dahin gerichtet, daß dir es dienen soll. Es ist mir allein um dich zu thun. Ich kann und will euer nicht vergessen: ich habe so große Sorge für dich und habe dich so lieb, daß wenn du auch nun in die Erde verscharrt und von den Würmern gefressen bist, soll doch nicht ein Haar umkommen.

52. Und also sagt auch der 34. Psalm V. 21.: „Er bewahret ihn alle seine Gebeine, daß deren nicht eins zerbrochen wird.“ Wiewohl die Welt mit Schwert und Feuer wider die Gottseligen tobt und wüthet, wie Johannes Fuß im Concil zu Costnitz von den Papisten ist verbrannt worden, so hat doch gleichwohl Gott alle seine Gebeine, dazu die aller kleinsten Glieder und Härlein in ein Bündlein zusammen gelesen, auf daß sie von den Feinden nicht zerstreuet würden. Laß sie nur brennen, Haut und Haar dahin gehen, es ist alles in ein Bündlein eingewickelt, daß sie nicht sollen umkommen. Ja, daß noch mehr ist, er sammelt auch die Klauen von dem Vieh, nicht allein die Haare der Gottseligen; wie 2 Mos. 10, 26. steht: „Es soll nicht eine Klaue dahinten bleiben.“ Endlich zählt er auch ihre Thränen, wie der 56. Psalm V. 9. sagt: „GOtt, zähle meine Flucht, fasse meine Thränen in deinen Sack. Ohne Zweifel, du zählest sie.“

53. Ist aber das nicht eine große Ehre der Heiligen und Gläubigen, daß alles, was wir thun und leiden im Namen Christi, wie in eine Tafel geschrieben und festes vor GOttes Augen ist. Denn also wird das Geschlecht Jakobs, dieweil es von GOtt berufen und von ihm auch regiert wird, so fleißig ausgezählet, daß auch jetzt noch sein Gedächtniß währet. Daß man nun diese Lehre immer behalte, und dieselbe in der christlichen Kirche fleißig treibe, solches ist sehr nütze, den Glauben und die Hoffnung des zukünftigen Lebens und ewiger Herrlichkeit damit in uns zu stärken und zu erhalten.

54. Ich will aber jetzt die Frage, welche Lyra und die Andern erregt haben, nämlich, von den siebenzig Seelen, so aus Jakobs Lenden gekommen sind, absichtlich fahren lassen. Desgleichen auch, daß ein Plural, das ist, ein Wort, das Viele bedeutet, gesetzt wird zu einem Singular, der nur Eins bedeutet; als, daß hier im Lateinischen steht, septuaginta anima etc., das heißt: siebenzig Seele. Denn das ist in der heiligen Schrift sehr gemein, und die hebräische Sprache, die nicht so eine strenge Regel hat, daß ein Eigenschaftswort und ein Hauptwort, wie es sich doch im Lateinischen verhält, bei einander stehen müssen in derselben Zahl; wie wir droben davon etliche Exempel gehabt, als: Vixit septuaginta annum, et unum annos: Er hat gelebt siebenzig Jahr und ein Jahre. Und hier steht auch: Omnes anima sexaginta sex: Alle Seele sechs und sechzig. Und daßelbe wird gemeiniglich also gehalten in den Wörtern, damit man zählt, oder Zahlwörtern. Dies gehört zu der Grammatik und ist an sich selbst einfach, daß es jedermann verstehen kann.

\quellseite{456}

55. Das hebräische Wort nephesch heißt allenthalben in der heiligen Schrift eine Seele, und der ganze Mensch wird also genannt, von dem wir sagen, daß er Leib und Seele habe, und nennen ihn entweder einen lebendigen Leib, da eine Seele drinnen ist, oder ein beseeltes Wesen.

56. Wir wollen aber hier nicht disputiren, ob die Seele vom Leibe herkomme, daß ist, vom Samen des Vaters; von welcher Disputation die Sententiarier viel handeln, nämlich, ob die Seele auch herkommen ex traduce, wie sie reden, vom Samen wie der Leib, und quälen sich darüber mancherlei Weise; wiewohl ich nicht verstehen kann, mit was Furcht oder Nutzen sie davon also untereinander zanken. Darum wollen wir diese Frage aufschieben bis an einen andern Ort. Etliche halten dafür, daß sie nicht herkomme ex traduce, vom Samen, sondern wenn der Leib bereitet und geschaffen sei mit allen seinen Gliedern, so werde die Seele demselben Menschen bald und in einem Augenblick eingegossen, und wenn sie also eingegossen werde, so werde sie auch gleich damit geschaffen, also, daß beiderlei auf einmal und zugleich geschehe. Augustin will nichts gewiß schließen, und bekennt frei, daß er es nicht wisse.

\section[Dritter Theil]{Dritter Theil. Wie Jakob den Juda zu Joseph voraus sendet; und wie endlich Jakob und Joseph zusammen kommen}

\subsection{I}

\begin{quote}
\textbf{V. 28. Und er sandte Juda vor ihm hin zu Joseph, daß er ihn anweise zu Gosen; und kamen in das Land Gosen.}
\end{quote}

57. Juda wird von seinem Vater vorausgesandt zu Joseph, welche Ehre billig den ersten drei Söhnen Ruben, Simeon und Levi gebühret hätte. Aber er zieht diesen vor; wie er auch droben, Cap. 44, 18. ff., vor Joseph allein für die Andern das Wort geführt und ihre Sache ausgerichtet hat. Darum ist aus dieser Erzählung wohl zu sehen, daß der Vater viel von ihm gehalten, und Juda ein sonderlich Ansehen bei ihm gehabt haben; denn er sich auch mit sonderlichem Fleiß und Geist seiner angenommen hat; wie das hernach wird zu sehen sein an dem Segen, so der Vater über ihn gesprochen hat, darin er weissagen wird, daß Christus von seinem Stamm kommen werde.

58. Juda ist dazumal ohngefähr zwei und vierzig Jahr alt gewesen und ist Vater von zwei Enkeln aus der Blutschande, und ist zugleich König und Priester in des Patriarchen Jakob Hause wider die rechte gemeine Ordnung der Geburt, und hat die Kirche oder Gemeinde, so in Jakobs Hause versammelt gewesen, gelehrt und regiert. Derhalben sendet Jakob den Sohn, so unter den andern der ehrlichste war, vor sich hin zu Joseph, lebendig, wie im Hebräischen steht, daß ist, zu unterweisen (denn jarah, auf Hebräisch, heißt, lehren, unterrichten, daher kommt das Wort moreh, Lehrer, und torah, Gesetz), daß Juda seinem Bruder berichten soll und ihm anzeigen, daß sein Vater jetzt komme, und ihn ermahnen, daß er gen Gosen kommen wollte; denn daselbst wolle der Vater verziehen und auf seinen Sohn warten. Denn man hat Joseph erinnern müssen der Ankunft seines Vaters, auf daß er gewisse Orte bestimmen und die Herberge bestellen möchte.

59. Nun scheint dies ein geringes Ding und nichts auf sich zu haben, daß Jakob Boten vor sich hin sendet zu Joseph, seinem Sohne. Aber die heiligen Väter sind sehr weise und voll heiligen Geistes gewesen. Und wir haben zuvor auch gehört, daß Jakob greulich erschrocken gewesen ist, da er nach Egypten hat ziehen sollen. Denn diese Leute sind nicht so unbedacht gewesen in ihrem Rath, Handlung und Werken wie wir: sie sind nicht also zugeplumpt, als wir tolle Narren. Denn nachdem Jakob das Land Canaan nun mit all seinem Hausgesinde verlassen und zwar mit großer Lust und Freude, hat er gleichwohl unterwegs zu Bersaba verzogen, da er großer Sorgen und Bekümmerniß halben sehr erschrocken gewesen, daß er ohne Verheißung, ohne göttlichen Befehl und Zeugniß in ein fremd Land ziehen sollte. Als er aber daselbst mit Gottes Wort aufgerichtet und gestärkt worden, ist er doch der Furcht noch nicht gar los geworden, sondern ein neues Fürchten und Bekümmerniß befällt sein Herz. Warum aber das? Ei darum, daß geschrieben steht im 146. Psalm V. 3.: „Verlasset euch nicht auf Fürsten.“ Denn aus dem Trost, so ihn Gott zur Antwort gegeben, hat er nicht können gewiß wissen, was Pharao und der Egypten Wille sein möchte, wiewohl er nunmehr an Gottes Willen nicht gezweifelt, und hatte auch zuvor oftmals mit der That erfahren, wie wunderbärlich Gott seine Verheißungen zu erfüllen pflege.

60. Denn er war von Jugend auf, und von der Zeit an, da er den Segen empfangen, auf mancherlei Weise geübt und versucht worden. Er hatte seines Vaters Haus verlassen aus Furcht vor seinem erzürnten Bruder Esau und war ganze zwanzig Jahre im Elend gewesen, und hatte in solchem Elend dazu noch die schwerste Knechtschaft getragen unter seinem Schwiegervater Laban, der ein gottloser und geiziger Mann war. Dazu ist auch gekommen, daß ihm Bilha, sein Weib, und seine Tochter Dina beschlafen worden sind; item, der Tod Rachels; und endlich, daß ihm der Sohn, den er insonderheit herzlich lieb gehabt, ist verkauft und nach Egypten weggeführt worden.

61. Diese wunderbare Regierung Gottes kann die Vernunft nicht begreifen, kann auch nicht verstehen die Weise, die Zeit, den Weg und Rath Gottes, und muß endlich das Fleisch herausfahren und also sagen: Gott ist ein Lügner. Denn heißt das seine Verheißung halten oder erfüllen, den Sohn, welcher der Erbe der Verheißung ist, wegnehmen? Und wer so oftmals ist getroffen worden, derselbige wird endlich weise durch seine eigene Gefahr, so ihm begegnet, lernt sich hüten und wird leichtlich geärgert; wie man pflegt zu sagen: Wer sich so oft verbrennt, der wird auf die letzte scheu.

62. Wiewohl nun Jakob eine sehr gewisse Antwort von Gott seine Reise nach Egypten betreffend gehabt, in welcher Antwort ihm auch zugleich verheißen worden, daß ihn Gott wolle wiederum herauf führen und zum großen Volk machen (welches ja alles wahrhaftig und gewiß ist, und ihm nicht hat fehlen können), doch zweifelt er noch etwas an Pharao. Darum denkt er: Was mag es mit dem Könige und den Obersten im selben Königreich werden wollen? Wie mancherlei und viel Hindernisse werden noch kommen und deswegen vorfallen, daß die Menschen so mancherlei Willen haben und so sehr unbeständig sind in dem fremden Lande und großen herrlichen Königreich, darin sie nicht leichtlich werden leiden oder ertragen fremde Leute, und daß sie von so geringer Abkunft bei ihnen sollten groß werden, vornehmlich, wenn sie sehen würden, daß dieselben wachsen und bei ihnen regieren sollten. Denn solches bezeugen die Historien von allen Zeiten her und die Erfahrung bringt es auch also mit.

63. Joseph ist zwar bis daher ein Herr und Lehrer gewesen der Fürsten und Priester im Lande und auch des Königes selbst etc., und wird ohne Zweifel große Furcht gehabt haben. Er hat hin und wieder viel Kirchen bestellt und geordnet, hat aber das ganze Königreich nicht können bekehren; denn es ist doch allezeit der größte Theil abgöttisch und des Teufels. Wie daran auch kein Zweifel ist, daß Daniel und seine Gesellen zu Babel, Mardachai, item, Esther nicht auch sollten Viele bekehrt haben. Aber die ganzen Königreiche haben sie gleichwohl nicht allenthalben bessern können. Also hat Jona den König und die Fürsten zu Ninive bekehrt, aber doch nicht alle, und ist der größte Theil hernach wiederum in die vorigen Sünden und Laster gefallen.

64. Darum ist es genug, daß wir ein Theil des Volks und dazu noch wohl das geringste bekehren mögen, daß sie sich zur Besserung schicken und fromm werden. Unterdeß regiert Gott wunderbarer Weise und läßt die Abgötterei bleiben. Er sendet aber seine Fischer, als, die Daniel, Joseph, Esther, Mardachai, welche wie arme Würmlein sind, und sammelt doch bisweilen Könige und andere große Monarchen in seine Kirche. Aber das Reich dieser Welt bleibt des Teufels Reich. Cyrus ist nicht durch Gottes Wort zum Königreich gekommen, sondern als er ein König und großer Monarch geworden ist und ein Herr über Babylon, da hat er erst Daniel und seine Gesellen gehört und ihre Lehre angenommen; wie der Herr sagt im Propheten Jesaia Cap. 45, 4.: Ehe denn du mich erkannt hast, habe ich dir Königreiche und große Städte gegeben.

65. So ist nun der König und die Fürsten in Egypten auch fromm gewesen, scheinheydisch geredet; denn es ist nur der geringere Theil durch Joseph, oder aber von Jakob und seinen Söhnen bekehrt worden.

66. Also haben wir zu unsern Zeiten auch das Wort, so die Propheten und Apostel gehabt, bekehren damit viele Fürsten und Herren. Und auch wohl einen oder den andern König: aber doch bleibt der größere Theil gottlos und tobt immer für sich hin, bis Gott einmal kommen wird und Alles zunichte machen. Denn das Wort, so aus des Herrn Munde gehet, soll nicht leer wieder zu ihm kommen, sondern soll thun, was ihm gefällt, wie Jes. 55, 11. steht. Die Papisten aber, und andere Widersacher und Verächter des Worts sollen doch endlich umkommen. Sie werden darüber zu Grunde gehen.

\quellseite{457}

67. Darum wird Jacob also gedacht haben: Siehe, wir ziehen hinab nach Egypten, da werden wir etliche fromme Fürsten finden, die bisher Joseph haben predigen gehört; wie aber, wenn sie wiederum von seiner Lehre abgefallen wären? Wer weiß, ob er auch noch in Gnaden sei bei dem Könige und den Fürsten? Was will werden nach seinem Tode? Wie, wenn die andern Prälaten, Räthe oder Priester den so frommen König wiederum verführt hätten? Darum will er alle Dinge zuvor ausforschen und schickt denselben Juda vor sich hin, und befiehlt demselben, daß er Joseph aller Dinge berichten wolle, und bei ihm und den Einwohnern Egyptens Alles auf das fleißigste bestellen und zubereiten. Denn wiewohl er die Verheißung hat, kann er doch bei sich noch nicht schließen, daß er derhalben alsbald ohne alles Bedenken so hindurch reisen wolle.

68. Also sollen wir auch thun lernen. Wiewohl wir dessen gewiß sind, daß Gott für uns sorgt, sollen wir doch wissen, daß man gleichwohl dabei der Dinge und Mittel gebrauchen soll, die Gott verordnet hat, auf daß wir Gott nicht versuchen. Wir sollen nicht hinein plumpen, und sagen: Ei, es muß wohl geschehen, dieweil ich die Verheißung habe von Gott; darum will ich mich nur mitten unter die Löwen und Wölfe dringen, sintemal mir Gott verheißen hat, daß er mich erhalten wolle. Ich will nicht essen; denn Gott wird mich sonst wohl ernähren. Ja, Gott will haben, daß wir der Hülfe gebrauchen sollen, so zu Erhaltung und Regierung dieses Lebens geschaffen sind: und wenn wir damit nichts ausrichten, soll man alsdann noch darum nicht verzagen, sondern sich daselbst an Gottes Hülfe fest halten, die er uns verheißen hat.

69. Denn also haben die Väter auch gethan. Da sie gewisse Verheißung gehabt, haben sie gleichwohl die ordentlichen Mittel nicht verachtet, sondern haben ihrer Geschicklichkeit und Verstandes gebraucht, sich vor Gefährlichkeit zu hüten und sonst ihr ganzes Leben zu regieren. Sie haben aber dabei zugleich das Wort auch immer vor Augen gehabt und sich darauf verlassen, wenn ihre Arbeit und was sie gewußt haben, nicht genug gewesen ist.

70. Denn unser HErr Gott regiert uns also, daß er ja nicht haben will, daß wir sollen müßig sein. Er gibt uns Nahrung und Kleider; aber doch also, daß wir pflügen, säen, ernten und dreschen sollen. Er gibt uns über das auch Kinder, welche geboren werden und aufwachsen durch Gottes Segen, und müssen doch von den Eltern genähret, versorget, auferzogen und unterrichtet werden. Wenn wir aber gethan haben, was an uns ist, alsdann sollen wir das Andere Gott befehlen, und unser Anliegen auf den HErrn werfen, der wird es wohl machen, er wird uns wohl versorgen, wie der 55. Psalm V. 23. sagt.

71. Also könnte Gott die Kirche auch wohl regieren durch den Heiligen Geist ohne das äußerliche Predigtamt; er will aber das nicht thun ohne Mittel. Darum sagt er zu Petro, Joh. 21, 16.: „Weide meine Schafe“; gehe hin, predige, taufe und absolviere die Leute. In der weltlichen Polizei sagt er zu der Obrigkeit: Du sollst wachen, schüßen, des Schwerts gebrauchen. Darum nennt Paulus die Apostel „Gottes Mitarbeiter“, 1 Cor. 3, 9. Er arbeitet wohl allein, aber er thut es gleichwohl durch uns.

72. Also hätte er auch wohl Kinder schaffen können ohne Adam; wie Adam anfänglich dazu ganz nichts gethan hat, da er aus dem Erdkloß ist gemacht worden, und Eva auch nicht, da sie aus der Rippe Adams geschaffen ist; aber darnach hat er zu ihnen gesagt, 1 Mos. 1, 28.: „Seid fruchtbar und mehret euch“, als wollte er sagen: Ich will nun hinfort Kinder schaffen durch eure Mitwirkung.

73. Also verhält es sich in allen andern Werken dieses gemeinen Lebens auch. Ich soll den Baum nicht mit der Nase von einander hauen, soll auch keinen Strohalm dazu nehmen; sondern soll eine Art oder Säge nehmen; und darum hat Gott dem Menschen Vernunft, Sinn und Stärke gegeben; derselben sollst du gebrauchen als Mittel und Gaben Gottes.

74. Dies soll man allezeit wiederholen und treiben in der Kirche oder Gemeinde Gottes, daß wir ja auf dem rechten Mittelwege oder Heerstraße bleiben und weder zur Rechten noch zur Linken weichen. Denn man soll ja an der Verheißung nicht zweifeln, und soll auch Gott nicht versuchen, daß man die Mittel, so Gott verordnet, wollte liegen lassen oder verachten.

75. Auf diese Weise hat nun Jacob, da er des Königs Herz nicht völlig hat erkennen oder verstehen können, seinen Sohn Joseph lassen seine Ankunft anzeigen, daß er ihm wolle entgegen kommen in das Land Gosen, darin er ihm mit seinem Hausgesinde hat heißen bleiben. Und wird ohne Zweifel gehört haben, daß Joseph mit den stolzen Egyptern nicht an einem Tisch gegessen hat, von welchen wir gehört, daß sie sehr stolz und trotzig gewesen sind wegen des Königreichs Egypten, welches in so großem Flor allerlei Güter, Gewalt, und sonst vieler unzähligen Gaben Gottes halben gestanden hat. Sie sind abgöttisch und die allerhoffärtigsten Leute gewesen, so die Andern neben sich verachtet und nur von sich selbst viel gehalten. Sie haben die Griechen und Israeliten für grobe Leute gehalten. Wie jetzt das lose Volk der Italiener aus sonderlichem Stolz und angeerbter Bosheit die Deutschen, Spanier und andere Völker auch grobe unverständige Leute nennen.

76. Die aber von den Egypten sind bekehrt worden, sind ohne Zweifel treffliche fromme Leute gewesen; wie es denn pflegt zu geschehen in der Bekehrung derer, die zuvor sonderlich böse gewesen sind.

77. Es ist übrigens ein wunderlich und schwer Regiment gewesen, da ein fremder Mann, vor dem das hoffärtige Volk einen Greuel gehabt, in Egypten regieret hat. Aber seine Tugend und prophetische Weisheit, und daß ihm alle sein Rathen und Thaten so fein glücklich fortgegangen und so wohl gerathen ist, das hat ihm Ansehen und Gunst gebracht. Darum haben sie sein Regiment wohl ertragen können, und haben ihn geehrt um der sonderlichen und vortrefflichen Gaben willen, die sie an ihm gesehen haben; wiewohl sie sich doch nicht desto demüthiger gehalten, sondern bei ihres Landes Gebrauch geblieben sind.

78. Wo der Ort gewesen sei, da hier Gosen genannt wird, weiß ich nicht und die Geographen disputiren auch davon. Es scheint an das Land Canaan gegrenzt zu haben. Hieronymus schreibt, daß dieser Ort zu seiner Zeit Ptolemais genannt worden sei; welches ich den Geographen will lassen befohlen sein. Es ist aber gewiß, daß er nicht weit vom Lande Canaan, bei dem Wasser Nil muß gewesen sein. Denn sie sind noch nicht nach Egypten gekommen, sondern sind vorn geblieben, da man erst hinein zieht an den Ort, da hernach Raemses ist gebaut worden, da man nach Palestina durch das Rothe Meer und durch die Wüste gezogen ist.

79. Darum bleiben sie so lange zu Gosen, bis Juda nach Egypten gekommen ist. Und hat sich also Jacob durch sonderliche Vorsicht davor gehütet, daß er Gott nicht hat wollen versuchen, da er sich gefürchtet, es möchte vielleicht etwas durch Unfleiß oder Nachlässigkeit geschehen, das ihm könnte Schaden bringen; darum will er haben, daß Joseph alle Dinge fleißig bestellen und verordnen soll. Das ist ein gottseliger, guter und heiliger Rath, wie alle Werke der heiligen Väter auch heilig und gut sind. Denn wo sie im Glauben wandeln, so gefällt auch Gott ihre Geschicklichkeit und Weisheit der Vernunft wohl. Wiederum aber haßt Gott und hat einen Greuel vor den Werken der Gottlosen, ob auch dieselben schon äußerlich sehr heilig scheinen.

80. Es werden aber beide Vater und Sohn ohne Zweifel heftig begehrt haben, daß sie haben mögen zusammen kommen. Wie wird ihnen beiden das Herz gepocht haben? Wie denn solche ihre Begierde und Verlangen, dazu auch ihre unermessliche Freude an dem, so im Text hernach folgen wird, zu sehen ist.

\subsection{II}

\begin{quote}
\textbf{V. 29. 30. Da spannete Joseph seinen Wagen an, und zog hinauf seinem Vater Israel entgegen gen Gosen. Und da er ihn sahe, fiel er ihm um seinen Hals, und weinete lange an seinem Halse. Da sprach Israel zu Joseph: Ich will nun gerne sterben, nachdem ich dein Angesicht gesehen habe, daß du noch lebest.}
\end{quote}

81. Es ist wahrscheinlich, daß Jacob seinen Sohn Joseph kaum wird erkannt haben, da er vor ihn gekommen ist; denn sein Angesicht wird innerhalb zwei und zwanzig Jahren etwas verändert worden sein, vornehmlich da er die königliche Kleidung und Schmuck angehabt. Es würde aber ihre große Begierde, so sie zu einander gehabt, damit abgemalt, daß sie einander also umfangen haben: da Joseph seinem Vater um den Hals gefallen und immer an ihn gehangen, ist er darüber ganz freudenvoll geworden, auch also, daß er sich nicht hat abhalten wollen von seinem alten Vater loszureißen, da er ihn also umfangen und geküsset hat. Und der Vater, so bisher sein Leben in der größten Traurigkeit und Schmerzen hatte zugebracht, wird jetzt gleichsam wiederum aus der Hölle und Tode zum Leben auferweckt.

\quellseite{458}

82. Darum ist dies wohl eine leibliche Freude und Fröhlichkeit; aber dieselbige ist auch voller rechtschaffener, herzlicher und freundlicher Neigung gewesen, welche Gott über die Maßen wohl gefallen hat. Und ist niemand, der das mit Worten erlangen und ausreden, oder auch verstehen könnte, wie groß diese herzliche Neigung gewesen sei, es sei denn, daß er es auch einmal erfahren werde.

83. Der Vater sagt: Ich will nun fröhlich und gern sterben, wie es unser Dolmetscher im Lateinischen gegeben hat: Jam laetus moriar, etc. Als wollte er sagen: Ich wollte nun diese Stunde oder Augenblick sterben, wenn es Gott also haben wollte. Es sind sehr bewegliche Worte, daß er sagt: Ich bin bereit, jetzt in den Tod zu gehen; denn mein armes Herz ist bisher ganze zwei und zwanzig Jahre lang mit erbärmlichem Wehklagen und Traurigkeit jämmerlich geplagt worden, da ich an deinem Leben gar verzagt hatte: nun aber, da ich dich wiederum lebendig und gesund sehen mag, däucht mich, daß ich vom Tod in das Leben gekommen und alles Jammers und Beschwerung gar los geworden bin.

84. Er hat sich aber die Ehre und Herrlichkeit Josephs noch nicht lassen anfechten oder bewegen. Denn da er gesehen, daß sein Sohn noch am Leben gewesen, hat er so groß Verlangen zu ihm gehabt, daß er sich um die große Herrlichkeit und königliche Gewalt, so Joseph gehabt, nicht sonderlich angenommen hat, und ist damit allein zufrieden, daß nur sein Sohn noch lebt. Ja, das noch mehr ist, er denkt also: Wenn du auch der allernämeligste Mensch und gar verlassen wärest, so hätte ich doch daran übrig genug, daß ich sehen mag, daß du noch lebt und gesund bist; ich achte der Herrlichkeit und des großen Gutes nicht, so nicht allein in Egypten, sondern auch in der ganzen Welt ist; denn ich will nun die Welt gern verlassen und billig und fröhlich aus diesem Leben abscheiden.

85. Solches werden ohne Zweifel die Worte gewesen sein, so Jakob aus väterlichem Herzen geredet hat, welche Mose nur kürzlich beschreibt. Wie er auch des Andern alles schweigt, was geredet oder geschehen ist, da sie sind zusammen gekommen. Es wird aber wahrlich eine lange Disputation gewesen sein, und werden ein Gespräch mit einander gehalten haben von mancherlei Fällen und vielen Veränderungen. Joseph wird seinem Vater in die Zähne erzählt haben, wie es allenthalben um seine Sachen gestanden und wie er so wunderbarlich Glück gehabt. Und sie sind dieselbe Nacht zu Gosen geblieben und haben mit einander zu Nacht gegessen. Und endlich ist das ganze Haus Jacobs und alle sein Gesinde erfreut worden, und haben mancherlei Unterredungen gehabt von allen Dingen, so zuvor geschehen waren, von den Brüdern und ihren Weibern, vornehmlich aber von dem großen Leide, so der Vater die lange Zeit über getragen. Dies alles hat Mose nur eben angezeiget und hat nicht anders können zugeben. Denn es thut sehr wohl, wenn man des Jammers, so nun vorüber ist, und der vergangenen Mühe und Arbeit gedenket. Wie der Poet sagt: Forsan et haec olim meminisse juvabit, das ist: Es wird uns vielleicht dermaleinst wohl thun, daß wir an diese Noth und Gefahr, so wir erlitten, gedenken mögen. Da werden Fragen gefallen sein von der Mutter Bilha, Silpa, Lea etc. Joseph hat aber des andern oder nächstfolgenden Tages zu seinem Vater und Brüdern gesagt: „Ich will hinauf ziehen“ zum Könige etc.

\begin{quote}
\textbf{V. 31—34. Joseph sprach zu seinen Brüdern, und zu seines Vaters Hause: Ich will hinauf ziehen, und Pharao ansagen, und zu ihm sprechen: Meine Brüder und meines Vaters Haus ist zu mir kommen aus dem Lande Canaan, und sind Viehhirten, denn es sind Leute, die mit Vieh umgehen; ihr klein und groß Vieh, und alles, was sie haben, haben sie mitgebracht. Wenn euch nun Pharao wird rufen, und sagen: Was ist eure Nahrung? so sollt ihr sagen: Deine Knechte sind Leute, die mit Vieh umgehen; von unserer Jugend auf bisher, beide wir und unsere Väter; auf daß ihr wohnen möget im Lande Gosen. Denn was Viehhirten sind, das ist den Egyptern ein Greuel.}
\end{quote}

86. Joseph hat wohl gewußt, wie hoffärtig die Egypter waren und was sie für eine Weise gehabt, nämlich, daß sie die Hebräer pflegten zu verachten, nicht allein deshalb, daß sie fremd waren, sondern auch wegen ihres Standes und Lebens, so sie geführt haben. Dieweil er aber auch wohl weiß und erkennt, daß die höchste Gewalt bei dem König war, braucht er abermals menschliche Geschicklichkeit und Klugheit und beweist die Ehre, so er seinem Obersten zu erzeigen schuldig gewesen; wie er denn zuvor auch nie etwas gethan hat ohne des Königs Willen und Befehl. Und hat zwar solche Klugheit wohl noch gethan um mancherlei Beschwerung und Unlust willen, so den Regenten pflegt in solchen großen Herrschaften vorzufallen.

87. Darum sagt er zu seinen Brüdern und seines Vaters Hause: Ich will hinauf ziehen, und hören, was des Königes Gemüth und Wille sei. Denn ihr kommt in ein sehr groß Königreich, da alles Volk vor euch einen Greuel hat um der Viehzucht willen: darum will vonnöthen sein, daß ich euch Zutritt zum König verschaffe. Ueber das muß man euch dessen auch erinnern, daß ihr euch mit höchstem Fleiß bemühen wollet, daß ihr euch richtet nach der Art und Weise der Egypter, die euch beherbergen werden, daß sie sich ja nicht ärgern an euren groben bäurischen Sitten; sondern ihr müßt euch an Leutseligkeit, Freundlichkeit und Höflichkeit gewöhnen: ihr werdet euch an die Art und Weise dieses hoffärtigen Landes gewöhnen müssen, damit sie nicht bald im Anfang etwas an euch finden, das sie tadeln und strafen mögen.

88. Ihr werdet vor den König und die Fürsten des Königreichs Egypten gestellt werden, da sollt ihr euch dessen befleißigen, daß ihr euch ja fein züchtig und demüthig haltet, und sollt männlich also sagen: „Wir sind Leute, die mit Vieh umgehen.“ An demselben Worte werden sich die Egypter wohl ärgern: es ist aber besser, daß ihr frei öffentlich und aufrichtig bekennet, von wannen ihr hergekommen und was euer Stand oder Amt sei, denn daß ihr etwas Anderes wollet vorwenden oder dies verhehlen. So kommt ihr am besten durch.

89. Ueber das will ich auch dem König sagen, daß ihr eure Herde und Vieh mitgebracht habt. Und desselben will ich euch darum zuvor erinnert haben, auf daß ihr eine gewisse Form haben möget, was ihr für Antwort geben sollt, wenn euch der König wird rufen; denn ihr sollt ihm nichts verhalten, sondern sagen: „Deine Knechte sind Leute, die mit Vieh umgehen, von unserer Jugend auf bisher“ etc. Bekennet gleich zu; und daß ihr daneben auch diesen Ort Gosen möget inne behalten, welcher nicht mitten im Königreich liegt, sondern an der Grenze des Landes Canaan. Denn wo ich oder aber der König selbst um meinetwillen euch an einen bessern Ort, der erblicher sein möchte, setzen wollte, das würde den Fürsten und Unterthanen im Königreich Egypten zuwider sein und sie würden euch darüben neiden. Ich gehe aber damit um, daß ihr möget an einem solchen Ort bleiben, da ihr gute Weide, und euch nicht zu befürchten habt, daß euch die Egypter etwa beleidigen oder neiden möchten. Denn was Viehhirten sind, das ist ihnen ein Greuel; nicht daß sie gar keine hätten, sondern darum, daß sie dieselben verachten. Wie bei uns auch die Kuhhirten und Schäfer für die allergeringsten Leute gehalten werden; oder wie Stockmeister, Metzger und Scharfrichter unserm Gebrauch nach nicht zu ehrlichen Aemtern oder Gesellschaft gezogen werden: also sind unter allen Ständen desselben Königreichs die Viehhirten auch die allerverächtlichsten Leute gewesen.

\quellseite{459}

\chapter{Das siebenundvierzigste Capitel}

\section[Erster Theil]{Erster Theil. Wie Joseph dem Pharao seines Vaters Ankunft meldet; wie er fünf seiner Brüder vor Pharao stellt; und wie Pharao mit ihnen ein Gespräch gehalten}

\subsection{I}

\begin{quote}
\textbf{V. 1. Da kam Joseph, und sagte es Pharao an, und sprach: Mein Vater und meine Brüder, ihr klein und groß Vieh, und alles, was sie haben, sind kommen aus dem Lande Canaan; und siehe, sie sind im Lande Gosen.}
\end{quote}

1. Siehe doch, wie beständiglich Joseph dabei beharret, daß er den König vor Augen und in Ehren gehalten hat. Denn er hat ihm solche Ehre erzeigt, daß er gar nichts gemacht hat vorzunehmen oder anzufangen aus eigenem frevellem Sinn, er habe es denn zuvor alles an den König selbst gelangen lassen, als der gar nichts hat thun wollen ohne seinen Willen.

2. Das ist aber eine sehr treffliche Tugend, und die sehr selten ist an allen Amtleuten, Steuerbeamten oder Rentmeistern, so jetziger Zeit diesem Exempel Josephs nicht nachfolgen, sondern das Ziel ihres Amtes und Befehls oftmals freventlich und unverschämt übertreten. Und wenn Einer eine Hand breit Befehl hat vom Fürsten, so nimmt er eine ganze Elle lang. Daher kommt auch so viel Unordnung in Regimente in Landen und Städten, und muß daraus auch folgen, daß sie dermaßen gar ungeschickt und verderbet werden. Denn die Amtleute sollten Treue und Glauben halten, ihrer Fürsten Befehl auszurichten, und sollten gar nichts anfangen oder vornehmen ohne ihren Willen und Befehl. Wenn sie das thäten, so würde auch hin und wieder in den Regimenten und Herrschaften so viel weniger Unglück, Jammer und Noth sein.

3. Wie ohne Zweifel Egypten ein sehr schön Königreich gewesen ist, das allenthalben im Flor gestanden, nicht allein Reichthums halben, welches das Geringste gewesen, sondern guter Ordnung und Gerechtigkeit haben. Und Joseph ist ein sehr trefflicher Mann gewesen, der dies Königreich sehr fein geordnet, und selbst auch mit schönen Exempeln der Jugend und des Gehorsams den andern Amtleuten und Unterthanen vorgeleuchtet hat, dieweil er ohne den Willen und Befehl des Königs nichts gethan hat.

4. Zu unsern Zeiten ist fast an allen Fürstenhöfen mehr Klagens über die Schinderei und ungerechte Gewalt der Steuerbeamten und anderer Amtleute, denn so die Bürger und Bauern unter einander üben. Man kann aber doch allerlei Ursache zwischen den Beamten und Unterthanen auch an die Fürsten nicht gelangen lassen, die ohne das eine solche Last auf sich liegen haben, welche schwer genug ist, nämlich, daß sie Befehl geben, heißen und verbieten, was man thun und nicht thun soll. Zudem aber müssen sie auch viele Klage hören von der Untreue derer, die sie Kindern zu Amtleuten vorgesetzt haben. Wir haben in unsern Landen und gemeinen Regimenten solche Josephsleute nicht, wie dieser gewesen ist. Wie wir auch jetzt in der Kirche keine Paulusse oder Petrusse haben, oder aber es gibt derselben fast wenig und sind zumal selten.

5. Derhalben, da Mose dies Exempel so sehr rühmt, will er damit anzeigen, daß Joseph nicht allein gelehrt, sondern auch selbst den Andern ein gut Exempel der Lehre und Gehorsams gegen den König gegeben, der, wie er gewesen hat, fromm und gottselig war. Und ist wahrlich dies ein schön Bild eines frommen Fürsten, daran man immer denken sollte, welcher für seine Person den Unterthanen wohl vorgestanden hat und nütze gewesen ist und mit gutem Exempel, dadurch er auf Andere gewirkt hat; also daß die andern Räthe oder Amtleute erweckt und ermahnt worden sind, daß sie in ihrem Berufe auch getreu und fein bescheiden gewesen sind. Denn nachdem sein Vater und Brüder gekommen sind, hat er alle Dinge also geordnet, daß sie ohne des Königs Willen im Lande Gosen nicht wohnen sollten.

Und wiewohl er schon aus sonderlichem Rath beschlossen, daß er sie dahin setzen wollte, will er es doch nicht einnehmen, oder darin etwas thun ohne des Königs Befehl; sondern heißt sie daselbst verziehen, bis ihnen wird angezeigt werden, was des Königs Wille sei.

6. Er mißbraucht weder seiner Gewalt, so ihm vom Könige gegeben war, noch seiner Gunst oder Gnade; sondern tritt fein demüthig und mit gebührlicher Unterthänigkeit zum Könige, und sagt: Mein Vater und meine Brüder sind jetzt aus dem Lande Canaan gekommen; dasselbe will ich dir derhalben angezeigt haben, daß du uns weißest, was du ihnen für einen Ort des Landes eingeben wollest, da sie wohnen möchten. Das ist eine schöne Demuth und Gehorsam, damit Joseph sich selbst und das ganze Haus seines Vaters Pharao unterwirft. Und wenn jetzt zu unsern Zeiten diese Demuth und Gehorsam auch also gehalten würde, sollte wahrlich viel mehr Glück und Seligkeit bei allen Ständen dieses Lebens sein, denn sonst ist.

\subsection{II}

\begin{quote}
\textbf{V. 2. Und er nahm seiner jüngsten Brüder fünf und stellte sie vor Pharao.}
\end{quote}

7. Nach dem hebräischen Text lautet es also: Er hat vom Ende seiner Brüder genommen; und kann von allerlei Enden, entweder dem obersten, oder aber dem untersten verstanden werden, das ist, von den ersten oder von den letzten fünf Brüdern.

8. Ich bin der Meinung, daß Joseph seine jüngsten Brüder vor Pharao gestellt habe. Denn das Wort „Ende“ heißt in der hebräischen Sprache das Aeußerste oder das Letzte in einem jeden Dinge; wie unter diesen Brüdern Dan, Naphtali, Ascher und Sebulon die letzten gewesen sind. Aber Andere mögen ohne alle Gefahr eine andere Meinung haben und derselben folgen.

9. Hier wird aber nun abermals gefragt: Was doch Joseph für eine Ursache möge gehabt haben, daß er dies gethan hat? Der Juden Rabbiinen bringen daher ein unnützes Geschwätz und Augendienung, nämlich, daß er die ausgelesen habe, so unter den Andern allen das geringste Ansehen und auch von Person die Geringsten gewesen sind. Denn wo er die Stärksten und so ein männlich Ansehen gehabt, ihm hätte vorgestellt, möchte vielleicht Pharao dieselbigen zum Kriege gebraucht haben. Aber wenn du ihrer aller Alter betrachten wirst, ist ihrer keiner unter dreißig oder vierzig Jahren alt gewesen, und ist ein geringer Unterschied von zehn Jahren, die Einer mag mehr gehabt haben als der Andere; welche Zeit unter den Brüdern keinen großen und offenbaren Unterschied hat machen können. Allein Benjamin ist unter ihnen der Jüngste; an den Andern hat man schwerlich einen Unterschied erkennen können, wenn man sie äußerlich angesehen und die Personen mit einander verglichen hätte. Denn sie sind beinahe alle innerhalb zehn Jahren nach einander auf diese Welt geboren worden. Und ist sonderlich an denen, so erwachsen sind, am Unterschied dieser Jahre nicht viel gelegen. Benjamin war aber schon jetzt Vater von zehn Kindern.

10. Darum halte ich es dafür, er werde die Letzten genommen, und seinen Bruder Benjamin damit geehrt haben, welchen er sonderlich vor dem Könige hat loben wollen, auf daß er sich in denselben vor den Andern wollte lassen befohlen sein, und hat ihm die andern vier jüngsten Brüder zugegeben, so von Bilha und Lea geboren waren.

11. Daß sich aber die Juden träumen lassen, Joseph solle diesen Plan gehabt haben aus Furcht, oder daß er sich etwas besorgt, ist zumal ein loses Geschwätz. Denn warum wollte er sich dessen besorgen, daß sie der König zum Kriege brauchen sollte, so doch nun oftmals gesagt ist, daß die Egypter vor allen Viehhirten einen Greuel gehabt haben? Derhalben hat es damit keine Gefahr, daß sie sollten zum Kriegswesen gezogen werden, dieweil er gewußt, daß der Adel in Egypten und die Hauptleute oder Obersten im Kriege alle Viehhirten gar verachtet haben. Und der König wird selbst hernach, V. 6., sagen: „So du weißest, daß Leute unter ihnen sind, die tüchtig sind, so setze sie über mein Vieh.“ Er weiß wohl, daß sie bei seinen Unterthanen gar verachtet sind, darum heißt er sie, sein klein und groß Vieh weiden.

12. Und daraus kann man nun sehen, daß der Greuel, so die Egypter vor den Viehhirten gehabt, also verstanden werden muß, daß sie dieselben gegen andere Stände der Menschen wohl verachtet haben, relativ und insonderheit, nämlich, so viel die großen herrlichen Stände, oder die vornehmsten berühmtesten Aemter belanget, daß sie dazu nicht sind gezogen oder gebraucht worden, sondern bei denselben gar verachtet gewesen. Wie jetzt auch noch die vom Adel und Bürger in Städten nicht gern große Gemeinschaft haben mit dem Bauernvolke und die Bauern selbst auch die Schaf- oder Sauhirten verachten und sie für geringer halten denn sich selbst. Und dasselbe geschiehet auch also mit mancherlei Völkern und die nicht aus einerlei Lande sind, als, die Franzosen achten sich besser und höher denn die Deutschen; und vornehmlich die Italiener verachten die andern alle mit einander auf eine hochmüthige Art und haben auch einen Greuel vor ihnen.

\quellseite{460}

13. Es können aber doch und sollen auch etliche unterschiedene Aemter oder Stände sein in diesem äußerlichen und bürgerlichen Leben. Denn es ist ja ein ander Amt, daß ein Hauptmann hat, und ein anderes eines gemeinen Kriegsmanns; ein Bürger hat ein anderes Amt denn ein Schafhirte. Also, wo sich die Egypter nicht mit Hoffahrt verkündigt hätten, haben sie ohne Sünde die Viehhirten geringer achten mögen, denn die Hauptleute im Kriege, oder sonst Andere, so hohen Standes und Adels halben höher und größer gewesen sind denn jene.

14. Und Joseph selbst ist von diesem Greuel frei gewesen, wiewohl sein Ursprung und Geschlecht vor den Egyptern greulich war. Es haben wohl die Egypter mit ihm nicht geessen, aber er hat vor ihnen und in ihrem Beisein oder Gegenwart sein Mahl gehalten; und er war wegen seiner Tugend und vortrefflichen Gottseligkeit zu solcher großen Herrschaft gekommen und erhoben worden. Wie der Kaiser Justinus soll aus einem Kuhhirten ein Hauptmann geworden sein und endlich das Kaiserthum erlangt hat. Und Sforza, Herzog zu Mailand, soll eines Bauern Sohn und ein Loch oder Sudeler (Gartoch) im Kriege gewesen sein; dieweil er aber stark von Leibe und dazu einen großen, freudigen und kriegerischen Muth gehabt, ist er zum Herzogthum gekommen und hat des vorigen Herzogs Tochter zum Weibe genommen. Also ist auch Matthias, der König von Ungarn, aus dem Gefängniß zum Königreich gekommen. Und Salomo sagt auch also im Prediger Cap. 4, V. 14.: „Es kommt einer aus dem Gefängniß zum Königreich.“

15. Denn das sind sonderliche Werke Gottes, davon gesagt wird im 113. Psalm V. 5. 6. 7.: „Wer ist, wie der Herr, unser Gott? der sich so hoch gesetzt hat, und auf das Niedrige siehet im Himmel und auf Erden; der den Geringen aufrichtet aus dem Staube, und erhöhet den Armen aus dem Koth.“ Das hat Gott wohl Macht, und hält doch den Unterschied der Aemter und Stände in diesem Leben, als, der Bürger, Bauern, Fürsten, Edelleute etc. Wiewohl es heutiges Tages auch zugeht, wie der alte Vers sagt: Et genus et formam regina pecunia donat, das ist: Geld, so allenthalben regiert und oben schwebt, kann beides großes Geschlecht und schöne Gaben geben; Geld kann die Leute edel und schön machen etc. Es bleiben doch die gemeinen Unterschiede der menschlichen Gesellschaft, daß ein Jeder sein Amt behalten und dasselbe verwalten soll, welches ihm von Gott zugeordnet ist, er sei ein Fürst oder Bauer. Dies will ich gesagt haben wider der Juden unnützes Geschwätz und Gedicht.

16. Denn Joseph hat dem Könige so viel anzeigen wollen, daß er und die Seinen einen ehrlichen Ursprung und Abkunft hätten; wiewohl der Ort, da sie geboren waren, sehr gering und unbekannt war. Darum ob sie wohl Viehhirten gewesen, wären sie doch fromme ehrliche Männer, die er in seinem Königreich ohne alle Gefahr wohl gebrauchen möchte. Denn diese Sorge müssen die Regenten auch haben, daß sie fleißig darauf sehen, was sie für fremde Gäste aufnehmen oder beherbergen. Darum hat nun Joseph diese Fremden auch vor den König gestellt, auf daß er sehen soll, daß er und sein ganz Königreich sich ihretshalben nichts zu besorgen habe, und das noch mehr ist, daß er beweisen will, daß es solche Leute, die Gott angenehm sind, die der Kirche nütze und dem ganzen Königreich Egypten auch heilsam sein können.

17. Und es scheint, als habe Jakob diesen Rath gegeben, daß Joseph seine ältesten Brüder nicht sollte vor Pharao stellen, als, Ruben, Simeon und Levi, welchen er ein wenig ungezogen war. Wiewohl ich nicht weiß, was er vom Juda gehalten; es sei denn, daß du diese Worte im Text „vom Ende der Brüder“ also deuten wollest, daß sie so viel heißen sollten, daß er Etliche von den Aeltesten und Etliche von den Letzten oder Jüngsten genommen habe. Ich kann aber davon nichts Gewisses sagen, und wie ich droben gesagt, ist hierbei gar keine Gefahr, man denke es auf welche Weise man wolle. Denn wiewohl den Aeltesten unter diesen Brüdern ihre Sünden vergeben waren, so haben sie doch der Strafe noch gewarten müssen; wie wir hernach hören werden.

\subsection{III}

18. Es werden aber ohne Zweifel zwischen dem Könige und den Brüdern Josephs viel mehr Gespräche gehalten worden sein, denn hier im Text beschrieben sind. Der König wird sie gefragt haben um Joseph, wie doch derselbe also von seines Vaters Hause sei weggerissen und nach Egypten geführet worden, und worüber sich dasselbe zugetragen etc. Es wird an diesem Orte nur die vornehmste Hauptfrage erzählt.

\begin{quote}
\textbf{V. 3. 4. Da sprach Pharao zu seinen Brüdern: Was ist eure Nahrung? Sie antworteten: Deine Knechte sind Viehhirten, wir und unsere Väter; und sagten weiter zu Pharao: Wir sind kommen, bei euch zu wohnen im Lande; denn deine Knechte haben nicht Weide für ihr Vieh, so hart drücket die Theurung das Land Canaan; so laß doch nun deine Knechte im Lande Gosen wohnen.}
\end{quote}

19. Pharao fragt nach ihrem Stande und Hantierung, so sie treiben; und sie antworten ihm darauf und sagen: Ein Schafhirte sind deine Knechte (denn also lautet es nach dem Hebräischen; wie wir droben auch andere Exempel gehabt, da nicht allezeit die Worte also bei einander gesetzt sind, wie die Regeln in der Grammatik fordern). „Deine Knechte“, sagen sie, „sind Viehhirten“; und gedenken dabei auch ihrer Väter: „wir und unsere Väter.“ Als wollten sie sagen: Wir sind je nicht solche Leute, die so gar neulich oder plötzlich aufgekommen wären; sondern unser Geschlecht und alle unsere Vorfahren sind Viehhirten gewesen. Darum können wir von uns Anderes nichts sagen und uns für nichts Anderes ausgeben, denn daß wir können das Vieh weiden und der Schafe hüten.

20. Also geben sie aber damit zu verstehen, daß sie nicht nach großer Herrlichkeit oder hohen Ständen der Egypter trachten oder stehen wollen. Derhalben ist der Juden Gedicht ein unnützes vergebliches Geschwätz, als sollten sie aus Furcht also geredet haben, auf daß sie ja nicht mit Gewalt zum Kriege möchten gebraucht werden; sondern sie wollen so viel sagen: Wiewohl unser Bruder der oberste Regent ist in deinem Königreich, begehren wir doch solcher oder dergleichen Herrschaft gar nicht; denn wir haben sonst unsers Bruders halben Ehre genug wegen der Herrlichkeit und großen Ehre, dazu er gekommen ist.

21. Und ich zweifle nicht, Joseph wird ihnen diese Instruction gegeben und sie also unterrichtet haben nicht allein um Demuth, sondern auch um ihres Nutzens willen. Denn wer sich davor hüten kann, daß er sich an keines Fürsten Hof begebe, der ist wahrlich ein seliger Mann, nicht allein zu unsern, sondern auch sonst zu allen Zeiten. Es hat ja wohl einen großen Schein, und die Herrlichkeit des Hoflebens ist wohl schön und nöthig, und die Menschen haben immer mehr Lust zu dem, das schön und herrlich scheint, denn zu dem, das nütze und nöthig ist; denn sie begehren nur groß Geld und Gut und ihren eigenen Nutzen, und dazu ein solch fein sanftes Leben, dabei so große Ehre und Herrlichkeit ist. Aber um die Aemter und Arbeit, dieselben zu verwalten, nimmt sich niemand an und begehrt derselben auch nicht; ja, vielmehr sind sie der Arbeit alle feind und fliehen davor, wie sie nur können und mögen. Und darüber haben allezeit alle Fürsten, Könige und Kaiser geklagt. Wie sie sagen, daß Kaiser Maximilian einmal soll geantwortet haben seinem Adel, da sich derselbe vor ihm beklagt und hart mit ihm darum geredet, daß er in allen Aemtern und Geschäften einen solchen Mann gebrauchte, der einer geringen Abkunft und nur ein Pfaffe gewesen, welcher darnach noch ein Cardinal und Bischof aus Salzburg geworden ist. Denen, sage ich, soll er geantwortet haben: Warum thut ihr es nicht? Es hätte euch gebühret, diese Aemter und Arbeit auf euch zu nehmen und dieselben fleißig zu verwalten. Ich muß einen geschickten und fleißigen Mann um mich haben, durch den ich meine Geschäfte ausrichten könne, ihr aber flieht vor solchen Aemtern; darum habe ich diesen Schreiber erwählen und annehmen müssen, der die Arbeit, so in meinem Hofe vorfällt, auf
\quellseite{461}
sich nehme und mir dieselbe ausrichte. Wenn es der Adel nicht thun will, so muß es der Schreiber oder der Pfaffe thun. Ritter oder Reisige kann ich leichtlich machen; aber solche Männer, die zum Regiment tüchtig sind und die Arbeit und Last des Reichs tragen können, die kann ich nicht machen. Denn es ist wahrlich ein sehr schwer Ding um das Regiment.

22. Darum hat Joseph seine Brüder fleißig und klüglich gewarnet. Er ist vielleicht für seine Person auch etwas verdrossen gewesen der Unlust und schweren Last halben, so ihm bisher obgelegen, und will nicht haben, daß solches seinen Brüdern auch sollte aufgelegt werden, oder daß sie etwa mit Gewalt zu solcher großen Gefahr sollten gezogen werden, so dem gemeinen Regiment obliegt und anhängt.

23. Und wer sich hüten kann, der mag es wahrlich wohl thun. Aber nicht auf die Weise, wie Diocletian und viele andere Könige, und auch der türkische Tyrann Mahomet gethan haben, die sich in die Wüste oder Einöde begeben und gesetzt haben, wie die Mönche oder Einsiedler, auf daß sie also vor der Beschwerung, Sorge und Arbeit des Regiments fliehen und davon befreit sein möchten. Denn die thun nicht recht, so zum Regiment gezogen und ordentlich berufen sind, und die davon abstehen und solch Amt verlassen, wenn sie kleinmüthig und verzagt werden wegen der großen Last und Sorgen, so darauf stehen. Dieselbigen sollen aber viel lieber in ihrem Amt, dazu sie von Gott berufen und bestellt sind, bleiben, und sich erzeigen wie tapfere und seine beständige Männer, die wegen der Gefährlichkeit und Ungestümigkeit, damit oftmals die gemeinen Regimente hart beschweret werden, nicht so bald kleinmüthig oder verzagt werden. Denn es ist vonnöthen, daß man in diesem Leben Herrschaften und Regimente habe, und gehören dazu fleißige und arbeitsame Leute und die geschickt sind. Und es gehöret auch sonst äußerliche Herrlichkeit, Ehre und Gut zu einem jeden solchen Amt, es sei gleich ein weltlich oder Kirchenamt: aber es will nicht faule Leute haben.

24. Unser Landesfürst, Herzog Johannes Friedrich, Churfürst zu Sachsen, arbeitet in so viel großen Geschäften so sehr, daß in Wahrheit von ihm wohl möchte gesagt werden: „Nachbar wird ein beinerner Esel sein“ etc., 1 Mos. 49, 14., nämlich, die unermäßliche Last der vielen Arbeit auf ihm liegen zu haben und zu tragen. Denn die Königreiche und Regimente müssen in allewege verwaltet und regiert werden. Es müssen die Sachen und Haber in Gerichten und auf dem Rathhause erkannt und entschieden werden, welches nicht geschehen kann ohne großen Verdruß, Beschwerung, Mühe und Arbeit. Darum hat der Kaiser Maximilian recht und wohl gesagt: Er wollte viel leichter Ritter oder Knechte machen denn Schreiber. Ich kann einen Ritter machen, sagt er, aber einen Schreiber kann ich nicht machen. Unsere Edelleute und Schreibhansen sind beide daheim im Frieden und auch in Kriegsläuften sein nütze, weder zum Regiment, noch sonst Sachen auszurichten: geben sich entweder auf den Geiz, der nicht zu sättigen ist, oder aber auf Saufen, Schlemmerei und Dämmen; aber vor der Last und Gefahr, so bei dem gemeinen Regiment ist, können sie sich fein hüten, und wenn es etwa im Kriege auch zum Treffen kommt, so sehen sie wohl zu, daß sie die Ersten an der Spitze nicht sind.

25. Und wir haben zwar droben auch gesagt, daß es fast ein schwer Ding sei, das gemeine Regiment zu verwalten; wie dasselbe die Leute, so etwas geübt und erfahren, und nicht ehrgeizig sind, sondern vor hohen Ständen fliehen, wo sie nur können oder mögen, wohl wissen.

26. Aber die thun darum nicht recht, die ihr Amt und Stand, dazu sie ordentlich berufen sind, verlassen. Es ist zwar verdrießlich genug, ein Prediger zu sein, ein frommer Hausvater zu sein, oder sonst irgend ein ander Amt treulich und gottseelig zu verwalten: aber man muß ja gleichwohl vor der Arbeit nicht fliehen, sondern stracks hinzutreten und dieselbe auf sich nehmen, und was für Beschwerung vorfallen mag, muß man mit großem starkem Muth überwinden. Wer das nicht thun will, der bleibe davon. Denn Gott hat seine Aemter und Stände nicht dazu geordnet und eingesetzt, daß die, so dieselben verwalten, müßig dabei sein sollten. Wie St. Paulus sagt Röm. 12, 7.: „Hat jemand ein Amt, so warte er des Amts“, seinem Nächsten damit zu dienen. Die zu den Aemtern berufen und bestellt werden, sollen nicht faule müßige Hummeln sein, die nur von anderer Leute Arbeit leben wollen; sondern sollen thun, was ihnen befohlen ist in ihren Amt, und wissen, daß sie durch göttlichen Beruf und Ordnung in das Amt gesetzt sind, daß sie stark und muthig sein sollen, und nicht bald kleinmüthig werden, wenn Gefährlichkeiten und Ungestümigkeiten vorfallen, damit die Regimente angefochten und zerrüttet werden. Und das sollen sie darum auch thun, auf daß sie mit der That die Leute sein mögen, wie sie genannt werden, nämlich, Gestreng, Edel, Großmächtig, Durchlauchtig etc., daß sie recht Feste und Ehrenfeste seien in ehrlichen Dingen. Denn es muß in allewege ein Fleiß zur Gottseligkeit und Tugend, und zu Gottes Furcht bei allen Regimenten sein.

27. Darum hat Joseph seinen Brüdern nämlich und treulich gerathen, daß sie in ihrem Beruf bleiben, und in dem Amte, darin sie bisher gewesen; daß sie in einem fremden Lande nur nicht nach großer Ehre oder Herrlichkeit stehen, sondern mit ihrem Stande und Amte wollen zufrieden sein, und daß sie sich nicht in Sinn nehmen, daß um der großen hohen Ehre willen, darin sie ihren Bruder sitzen sehen, sie eben auch dem nach müßten nach gleicher Ehre und Herrlichkeit trachten.

28. Und dasselbe wird jetzt bei uns auch nicht gehalten und man sieht jetziger Zeit solche Exempel nicht. Denn jetzt sieht du unter den Hofleuten deren viele, wenn sie etwa zu einem Ehrenamt gekommen sind, daß sie alsdann bald versuchen, ihre Verwandten und nächsten Freunde auch hoch zu heben und reich zu machen. Joseph aber heißt seine Brüder mit ihrem Stande zufrieden sein. Und wiewohl sie die Verheißung haben vom Lande Canaan, müssen sie doch der Zeit nicht zuvorkommen, so Gott vorhin versehen hat; wie hernach in den Büchern der Chronica, 1 Chron. 8, 20. 21., von den Kindern Ephraim erzählet wird, daß sie sind vertilget worden, da sie vor der rechten Zeit das Land Canaan haben einnehmen wollen.

29. Derhalben haben die Söhne Jacobs mit dem Könige dermaßen geredet, daß sie ihm zu erkennen gegeben, daß sie gar einen Abscheu hätten vor allen hohen Aemtern, und daß sie aus hochgedrungener Noth dahin gekommen wären, daß sie hätten und begehrten, daß sie auch am geringsten Ort im Lande bleiben und nur ihres Hirtenamtes warten möchten. Wir sind, sagen sie, nicht gekommen, daß wir hier regieren oder sonst etwas Großes vornehmen wollten, auf daß du ja keinen bösen Argwohn auf uns habest; sondern begehren nur hier eine Zeitlang zu bleiben, dieweil wir in unserm Lande, so wir verlassen, jetzt keine Weide oder Speise haben. Es hat uns der Hunger und Kummer und daß wir aller Nothdurft haben mangeln müssen, da heraus getrieben. Wir sind gar arm, und in das Elend getrieben worden; darum bitten wir, daß uns nur möge vergönnt und gestattet werden, irgend an einem Ort in deinem Lande, wenigstens im Lande Gosen, eine Zeitlang unsere Wohnung zu haben. Dies alles hat Mose nur mit wenig Worten überlaufen, aber es ist ohne Zweifel viel mehr geredet worden; und sie werden auch etwas weitläufiger geantwortet haben auf des Königs oder der andern Herren Fragen.

30. Du siehst aber, daß hiermit abermals Josephs Tugend und Frömmigkeit gerühmt wird, als der alles mit höchster Weisheit und Treue regiert und verwaltet hat. Denn diese schönen ritterlichen Tugenden stehen einem solchen großen Fürsten wohl an. Und dies Exempel ist eigentlich viel größer, denn sonst irgend anderer großer Helden, wer auch dieselben möchten gewesen sein, in allerlei weltlichen Regimenten. Wie hat er doch so weislich und wie mit großem starkem Muth und Tugend alle Dinge ausgerichtet, und wie große Gottseligkeit ist bei dem allen gewesen! Welch eine schöne Erkenntniß der rechten himmlischen Lehre hat an ihm geleuchtet, auch so, daß er ein rechter geschickter politischer Mensch und dazu ein vortrefflicher Theologe hat mögen genannt werden!

31. Aber hier kommt nun eine andere Frage, nämlich: Was doch in Egypten möge für Weide gewesen sein, da im Land Canaan und den nächst umliegenden Orten keine mehr gewesen, und dazu auch in Egypten selbst nicht wegen der unfruchtbaren Jahre? Auf diese Frage antworte ich also: Ich halte nicht dafür, daß ganz Egyptenland so gar verdorret gewesen sei, daß nicht etliche Aecker und Wiesen noch sollten übergeblieben sein, die noch Weide gehabt, dieweil sie überfeuchtet worden vom Ueberfluß des Wassers Nil, wie gering derselbe auch hat sein mögen; wie droben vom Unterschied des Ueberfließens gesagt ist. Wiewohl dasselbe auch wenig hieher zu dienen scheint, und wir nicht gewiß wissen können, wie sie die ganze Zeit über das Vieh mögen gemeidet haben. Aber das Land Gosen wird vielleicht nahe bei dem Nil gelegen haben an einem niedrigen Ort, da es so viel mehr und länger hat können gefeuchtet werden. Und es ist ein groß Ding, daß sie daselbst Weide gehabt, daran es doch sonst beinahe in ganz Egypten gemangelt hat. Wie hernach wird gesagt werden, daß die Einwohner Egyptens noch Vieh gehabt, das sie für Getreide gegeben haben; damit angezeigt wird, wie es scheint, daß dazumal noch Weide müsse übrig gewesen sein für das Vieh, und daß die Unfruchtbarkeit also allein auf den Mangel des Getreides oder auf die Fütterung für das große Vieh gezogen und davon verstanden werden soll.

\quellseite{462}
32. Jetzt aber, zu unsern Zeiten, ist Egyptenland gar verändert worden, wie der 107. Psalm V. 33. 34. sagt von göttlicher Strafe und Veränderung der Lande: „Welchen ihre Bäche vertrocknet, und die Wasserquellen versiegelt werden, daß ein fruchtbar Land nichts trug, um der Bosheit willen derer, die drinnen wohneten“ etc. Denn wo GOtt ein Land straft, nimmt er nicht allein die Menschen hinweg, sondern auch den Saft von der Erde. Wie zu unsern Zeiten Italien, und auch unsern Thüringerlande, welches in ganz Deutschland fast das fruchtbarste Land gewesen, an der Fruchtbarkeit auch etwas abgehet. Denn sie sagen, das Einkommen von sieben Jahren solle jetzt kaum so viel bringen, als vor dieser Zeit drei Jahre eingebracht haben. Dies geschieht um der Bosheit willen der Menschen, vornehmlich aber um des Wuchers und Geizes willen, welchem niemand steuern kann. Denn derselbe verfäußerische Geiz wagt alles, was er nur will, und saugt diese Lande jämmerlich aus.

33. Gosen aber, wie ich droben auch gesagt, scheint der Ort gewesen zu sein, da sie hernach Raemses hingebauet haben, da sie auch ausgezogen sind. Und es ist an den äußersten Grenzen Egyptens, gegen das Land Canaan und Arabien gelegen gewesen, da das Wasser Nil seinen ersten Eingang in das Meer gehabt, gegen das Morgenland; welcher Eingang Pelusium genannt wird. Es scheint, daß es ein sehr fruchtbar und fett Land müsse gewesen sein.

\begin{quote}
\textbf{V. 5. 6. Pharao sprach zu Joseph: Es ist dein Vater, und sind deine Brüder, die sind zu dir kommen; das Land Egypten stehet dir offen, laß sie an den besten Ort des Landes wohnen, laß sie im Lande Gosen wohnen; und so du weißest, daß Leute unter ihnen sind, die tüchtig sind, so setze sie über mein Vieh.}
\end{quote}

34. Der König hat Joseph ganz gnädig und mildiglich geantwortet: „Es sind“, sagt er, „deine Brüder“, mit demselben machst du es halten, wie dir es wohlgefällt. Wo unter ihnen Leute sind, die etwas tüchtig und geschickt sind, so setze sie über mein Vieh. Denn wiewohl meine Unterthanen sie verachten, will ich sie doch nicht verstoßen, dieweil ich dich also hoch erhoben habe, der du doch eben auch daher geboren bist, da sie hergekommen sind. Sie aber, wie sie selbst bitten und begehren, sollen gemeine Leute bleiben und kein Regiment haben. Du sollst deine Gerechtigkeit und Vorzug behalten; denn du bist durch sonderliche Freyheit, so allein auf deine Person geht, aus einem Schäflein ein Fürst geworden. Denn GOttes Gaben haben eine Ausnahme, welche GOtt sich selbst allein vorbehalten hat. Es gehört vor unsern Herrn GOtt. Darum soll ein Jeder in seinem Ort oder Stand bleiben, bis das Glück aus der gemeinen Ordnung herauszieht, welche GOtt will herausgezogen haben.

35. Es will aber der König geschickte und fleißige Hirten haben. Einen seiner Hirten, der mit seinem Amt wohl wüßte umzugehen, und der nicht allein wüßte soll, wie es sich mit seinem Amte halte, sondern der auch wacker sei, und fleißig thun wolle, was ihm befohlen wird: einen solchen Hirten will er haben, daß man ihm erzählen solle.

36. Und hier siehst du, daß auch die Hirten Männer der Tugend, taugliche und tugendhafte Männer genannt werden, das ist, die arbeitsam, fleißig und treu sind: die eines Dinges warten. Wie Spr. 31, 10. f. auch von einem tugendsamen Weibe gesagt wird, das ist, die ihr fräulich Amt ausrichtet, endlich, und eine seine Herrin und Hausmutter ist, die des Hauses fleißig wartet, die Kinder wohl aufzieht; und die fleißig und treulich thut, was einer frommen gottseeligen Matrone zusteht. Also wird Bogs, Ruth 2. 1., ein Mann der Tugend, das ist, ein ehrlicher Mann genannt. Und ein Jeder soll sich in seinem Amt und Stande auch also halten.

37. Solche fromme treue Hirten und Diener begehren heutiges Tages die Könige, Fürsten und Hausväter auch, und beklagen sich fast alle, daß jetzt keiner unter allen Dienern gefunden wird, der fromm und getreu sei. Und mag vielleicht dies auch die Ursache gewesen sein, daß Pharao einen guten, geschickten und getreuen Hirten begehrt hat, darum daß auch zu derselben Zeit unter einem solchen glückseligen Könige und in solchem sehr schönen Königreich doch ganz wenig fromme und getreue Diener sind gefunden worden.

\section[Zweiter Theil]{Zweiter Theil. Wie Joseph seinen Vater vor Pharao stellt, und Pharao mit ihm ein Gespräch hält; item, wie Jakob mit den Seinen versorgt wird}

\begin{quote}
\textbf{V. 7. Joseph brachte auch seinen Vater Jakob hinein, und stellte ihn vor Pharao. Und Jakob segnete den Pharao.}
\end{quote}

38. Wir haben droben gehört, daß Jakob nicht von selbst nach Egypten gezogen ist, sondern daß er sehr furchtsam gewesen, und daß ihn die hohe Noth gedrungen habe dahin zu ziehen, und ist auch nicht gern in Egypten geblieben, hat auch daselbst nicht wollen begraben werden. Denn der heilige Geist hat ein gutes Gedächtniß, und gedenkt noch an die Verheißung, so dem Abraham geschehen ist, 1 Mos. 15, 13. 14., da GOtt zu ihm also sagte: „Das sollst du wissen, daß dein Same wird fremd sein, in einem Lande, das nicht sein ist; und da wird man sie zu dienen zwingen, und plagen vier hundert Jahr. Aber ich will richten das Volk, dem sie dienen müssen.“ Also hat nun der Glaube auch ein trefflich gutes Gedächtniß. Denn mit dem Glauben und Verheißungen hält es sich also, daß Eins ohne das Andere nicht sein kann, und hält sich immer Eins an das Andere. Darum zieht Jakob wohl aus dem Lande Canaan heraus, aber mit Furcht, bis er zu Bersaba Trost empfangen. Denn wo derselbe Trost nicht gekommen wäre, so wäre er nimmermehr fortgezogen.

\subsection{I}

39. Da er aber nun nach Egypten kommt, schickt er sich in die Zeit, und hält sich, wie die gegenwärtige Noth erfordert, denkt also: Siehe, ich komme hier in ein fremd Land, darin ich und mein Same wird geplagt werden. Denn derweilen der Text sagt 1 Mos. 15, 16.: „Sie sollen nach vier Mannesleben wiederum hieher kommen“, hat er darauf diese Rechnung gemacht: Das erste Mannesleben ist Abrahams gewesen, das andere Isaacs, das dritte ist mein, das vierte meines Sohns Joseph. Darum gehe ich jetzt hin, gleichsam in das Fegefeuer und in die Hölle, weiß nicht, wie mir es gehen wird und wann ich werde wiederkommen. Er hat wohl den König geehrt, und doch zweifelt er daran nicht, er werde endlich noch seinen Samen oder aber je seine Nachkommen verfolgen.

40. Was sollte er aber thun? GOtt hat ihn heißen hinzuziehen, und ihm verheißen, daß Joseph ihm seine Augen sollte zudrücken, wenn er gestorben wäre. Damit ist er zufrieden, und denkt nun darauf, wie er nur je eher je lieber wiederum aus Egypten möge geführt werden, wenn er auch schon würde gestorben sein, wo ihm ja solches bei seinem Leben nicht widerfahren mag.

41. Nachdem nun Joseph seine Brüder vor Pharao gestellt, da führt er nun seinen Vater auch hinein. Wie ohne Zweifel dasselbe der König also hat haben wollen und wird es auch befohlen haben; denn er war fromm und gottseelig, der die Gaben GOttes an Joseph wohl verstanden und erkannt hat, nachdem er bekehrt war und die rechte Lehre und himmlische Weisheit von ihm gelernt hatte. Der andere König aber, so nach diesem regieren wird, wird Joseph nicht kennen und wird dazu seine Nachkommen hart plagen. Dieweil er aber jetzt seine Brüder gesehen hat, begehrt er nun seinen alten Vater auch zu sehen und ihn anzureden, der einen solchen Sohn hätte, welcher mit so großen Gaben von GOtt begnadet, so ein großer Prophet und voll Heiligen Geistes war. Denn je verständiger und weiser Einer ist, je mehr man begehrt zu sehen den Vater eines so frommen heiligen Sohnes.

42. Darum folgt nun Jakob seinem Sohne Joseph nach und kommt vor den König; er redet aber nichts, bittet ihn auch um nichts, sondern segnet ihn nur, und ehrt die weltliche und königliche Majestät, davon er weiß, daß es GOttes Ordnung ist. Und wird ohne Zweifel der König viel mehr mit ihm durch seinen Dolmetscher geredet und nach vielen Dingen gefragt haben, nämlich, wer sein Großvater und sein Vater gewesen und was jeder gethan habe; was sie für Jammer, Mühe und Arbeit gehabt und erlitten haben? Welches von Mose nicht erzählt wird, der allein das schreibt, wie Jakob den König gesegnet habe, daß ist, er habe ihm Glück und Heil gewünscht, Gott angerufen und gebeten, daß er den König und das ganze Königreich wolle erhalten, und daß es ihm, dem Könige, allenthalben möchte wohl gehen.

\quellseite{463}
43. Sonst gebühret es sich aber also, daß der Größere den Kleinern segnen soll; wie in der Epistel an die Hebräer steht, Cap. 7, 7.: „Es ist ohn alles Widersprechen also, daß das Geringere von dem Besseren gesegnet wird“; also, daß Melchisedek den Abraham segnet, 1 Mos. 14, 19. Aber solches muß verstanden werden von dem göttlichen Segen. An andern Orten aber heißt segnen so viel als, daß man Einem Glück und allerlei Segen wünscht. Ein ander-mal heißt es auch wohl, danken, loben, als, in den Psalmen, Ps. 103, 1.: „Lobe den HErrn, meine Seele.“ Item Ps. 117, 1.: „Lobet den HErrn, alle Heiden“, das ist, danket ihm, lobet ihn.

44. Derhalben segnet Jakob den König als der Geringere den Größeren, und thut eine Rede vor ihm, darin er seine Stunde rühmt, und dankt ihm für seine Wohlthat, so er von ihm empfangen, nämlich, daß er Joseph, seinen Sohn, also erhöht, ihm aber und seinen Söhnen so gnädiglich aufgenommen und empfangen hätte, und ihnen einen sonderlichen Ort im Lande eingeräumt, da sie wohnen möchten, und zu dem, daß er ihnen auch Getreide und Zehrung in das Land Canaan gesandt hätte etc. Das ist wahrlich eine tapfere und feine Rede gewesen. Denn wir dürfen nicht denken, daß Jakob ein Kind oder ungelehrter Bauer gewesen sei; sondern er ist ein Prophet, Bischof und gelehrt im Geiste gewesen. Darum hat er wohl reden können von seinen Sachen, und hat mit einer herrlichen langen Rede dem Könige für alle Wohlthat, damit er ihn gleichsam überschüttet hatte, gedankt; hat daneben eine Bitte mit angehängt, nämlich, daß der König ihm und allen seinen Nachkommen hinfort auch solche Gnade erzeigen wollte.

\begin{quote}
\textbf{8—10. Pharao aber fragte Jakob: Wie alt bist du? Jakob sprach zu Pharao: Die Zeit meiner Wallfahrt ist hundert und dreißig Jahr; wenig und böse ist die Zeit meines Lebens, und langet nicht an die Zeit meiner Väter in ihrer Wallfahrt. Und Jakob segnete den Pharao, und ging heraus von ihm.}
\end{quote}

45. Dies ist nun eine andere Rede, so der König gethan hat zu Jakob, in welcher er ihn fragt um sein Alter. Es werden aber ohne Zweifel auch andere Reden mit untergelaufen sein; wie es denn zu gehen pflegt in solchen freundlichen Gesprächen. Und hat vielleicht der König seine grauen Haare und Runzeln angesehen, die da angezeigt haben, daß er viel älter sein müßte, denn er zuvor von seinen Söhnen gehört, oder er selbst auch hat rechnen können. Derhalben fragt er ihn, und sagt: Ich sehe, daß du ein sehr alter Mann bist; Lieber, sage mir, wie viel Jahre bist du alt? Darauf antwortete Jakob: Ich bin ja so gar alt nicht, wie mir es am Gesichte angesehen wird. Ich habe hundert und dreißig Jahre auf mir und nicht mehr. Meine Väter haben etwas länger gelebt: Isaak ist an seine hundert und achtzig Jahre gekommen; Abraham ist hundert und fünf und siebenzig Jahre alt geworden. So weit bin ich nicht gekommen, wiewohl ich jetzt ein betrübter und gebrochener alter Mann bin.

46. Denn Pharao, wie es scheint, wird daher Ursache genommen haben, nach seinem Alter zu fragen, da er an ihm die langen grauen Haare und das verwitterte runzlichte Angesicht gesehen hat. Es sind aber zweierlei graue Haare. Etliche kommen vom Alter, die andern kommen vom Angst und Trübsal. Man sagt von den Leuten, daß sie von Sorgen und Bekümmerniß ihres Herzens grau werden; nach deutscher Redensart: „Ich möchte grau werden.“ Und im 6. Psalm V. 8. steht: „Meine Gestalt ist verfallen vor Trauern, und ist alt worden; denn ich allenhalben geängstet werde.“ Als wollte er sagen: Ich bin mit Feinden und allen Plagen also umgeben, daß immer, wenn die eine weg ist, bald drei andere kommen, so viel schwerer sind. Denn kein Unglück kommt allein, wie man zu sagen pflegt.

47. Solche Trübsal aber und Jammer machen, daß der Mensch auch so viel eher graue Haare und Runzeln bekommt und daß er auch wohl desto eher sterben muß. Daher kommen die jämmerlichen Klagen derer, so betrübt und traurig sind, daß sie sagen: Du wirst mich grau machen, wirst mich tödten. Denn schwere Anfechtungen kränken wahrlich und tödten auch manchen Menschen. Darum ermahnt Salomo auch in seinem Prediger am 11. Cap. V. 9. und sagt: „So freue dich, Jüngling, in deiner Jugend, und laß dein Herz guter Dinge sein in deiner Jugend.“ Ein junger Mensch soll fröhlich sein, soll sich selbst nicht fressen oder martern mit Angst und beissenden Sorgen, die ihm seine Kräfte verzehren und den Saft seines Lebens aussaugen. Gehe nicht frühe an, dich zu plagen; wenn das Alter herankommt und allerlei Sachen und Geschäfte vorfallen, so werden auch Sorgen und allerlei Jammer genug folgen. Denn wo Geschäfte sind, wird es sich selbst wohl lernen.

48. Darum sagt nun Jakob: Ich bin ja alt, aber nicht so sehr von Jahren, als von vielem Unglück und Jammer, welches den Menschen viel eher tödtet, denn sonst viel Jahre thun können. Und wenn wir jetzt unsere Zeit wollen ansehen, wie viel sind jetzt wohl Leute, die sechzig Jahre erreichen? Wie viel tausend Menschen kommen um durchs Schwert? Wie viel sterben an der Pestilenz, item, von Traurigkeit und Trübsal vor der rechten Zeit? Und ist jetzt ein groß Alter, wo die Leute fünfzig oder sechzig Jahre alt werden. Von Traurigkeit aber sagt Salomo Spr. 17, 22.: „Ein betrübter Muth vertrocknet das Gebeine.“ Denn von Bekümmerniß des Herzens verschmachtet auch der Leib, und wenn dem Herzen angst und bange ist, das nimmt den Saft des Lebens gar hinweg. Darüber beklagt sich Jakob hier auch, da er sagt: „Wenig und böse ist die Zeit meiner Wallfahrt“; denn ich habe viel schweren Unfall und Jammer erlitten, und mein ganzes Leben ist nichts Anderes gewesen denn eine Wallfahrt. Wie denn dasselbe aus seiner Historie wohl zu sehen ist, daß er ein sehr unglücklicher betrübter Mann gewesen fast bis an sein vier und siebenzigstes Jahr. Da er noch jung war, ist er aus seines Vaters Hause verstoßen worden von seinem Bruder Esau, da er zuvor auch allezeit war verachtet worden; und hat sich deshalb des Ehestandes und Haushaltung enthalten bis in das fünf und vierzigste Jahr. Da er aber den Segen empfangen, hat er noch zwanzig Jahre gedienet in Mesopotamien. Ach! wie hat er daselbst von seinem Schwiegervater Laban so große Unbilligkeit leiden müssen! Darnach, da er von ihm gen Sichem gezogen, haben ihn daselbst noch viel schwerere Unfälle betroffen. Denn da ist ihm seine Tochter Dina beschlafen worden; Simeon und Levi haben den Hemor und die Sichemiter erschlagen; Rahel und Debora sterben dahin; ich will noch des Ringens und Kampfes geschweigen, so er mit dem Engel gehalten. Welches alles solche Stücke sind, daß derselben wohl ein jegliches einem könnte grau machen auch vor dem rechten Alter. Endlich hat Ruben auch seines Vaters Weib geschändet; Joseph ist verkauft worden; und da derselbe also weggekommen ist, weiß der Vater nun nichts Anderes, denn sein allerliebster Sohn müsse gar umgekommen sein. Das ist wahrlich der allergrößte Jammer, so er unter andern gehabt hat. Dies sind nun die bösen Tage, so einem sein Leben verkürzen und die Menschen häßlich, grau, runzlicht und alt machen. Denn da brechen die bösen Tage herein, und erfüllen, was sonst der Jahre halben am Alter mangelt.

49. Daß er aber seine Jahre oder Tage die Zeit seiner Wallfahrt nennt, das sind Worte des Geistes und Glaubens, der auf ein ander Leben weist. Denn er will dies armselige elende Leben, so ganz voll Kreuz und Jammers ist, nicht einmal würdig achten, daß er es wollte ein Leben nennen; sondern nennt es nur „eine Wallfahrt“, einen gar betrübten Wandel, dem man doch dienen muß, darum daß Gott also befohlen hat und sein Wille ist, der darin einem jeglichen Menschen seinen eigenen Ort und Zeit zueignet, auf daß er also erhalten und regieret werden möge, so lange es denselben lieben Gott wird gefällig sein. Es ist übrigens dies Leben so greulicher Weise erbärmlich, beschwerlich und betrübt wegen mancherlei Trübsal und Plage aller Teufel und der ganzen Welt. Und ist kein Zweifel, es wird Jakob auch viel mehr andere Trübsal erlitten haben, die nicht alle hier beschrieben worden ist. Darum ist dies Leben kein recht Leben, sondern ist eine Tödtung und Plage des Lebens. Diese Hoffnung erhält uns noch immer, nämlich, daß wir wissen, daß es nur unsere Wallfahrt ist. Wie die Epistel an die Hebräer beide diesen Text und andere dergleichen Stellen mehr sehr fein auslegt im 11. Cap. V. 8–10.: „Durch den Glauben ward gehorsam Abraham, da er berufen ward, auszuziehen in das Land, das er erben sollte; und ging aus, und wußte nicht, wo er hinkäme. Durch den Glauben ist er ein Fremdling gewesen in dem verheißenen Lande, als in einem fremden, und wohnete in Hütten mit Isaak und Jakob, den Miterben derselbigen Verheißung. Denn er wartete auf eine Stadt, die einen Grund hat, welcher Baumeister und Schöpfer GOtt ist.“

\quellseite{464}
50. Auf diese Hoffnung, daß sie eines bessern Lebens gewartet, haben sie die allerschwerste Wallfahrt getragen. Darum nennt Jakob sein Leben nicht Tage einer Wohnung, der Ruhe und Heile, wenn er auch schon ein Königreich hätte inne gehabt. Wie David von sich selbst auch sagt im 39. Psalm V. 13.: „Ich bin beide, dein Pilgrim und dein Bürger, wie alle meine Väter.“ Denn dies Leben ist denen, so da glauben und göttliche Verheißung haben, nur eine Wallfahrt, darin sie erhalten werden in der Hoffnung eines zukünftigen und bessern Lebens. Und es scheint, daß David aus diesen Worten Jakobs, da er sagt: „Die Zeit meiner Wallfahrt langet nicht an die Zeit meiner Väter in ihrer Wallfahrt“, diese Worte des Psalms genommen habe, da er sagt: „Wie alle meine Väter.“ Darum sind wir nun vor GOtt alle Bürger in der Hoffnung; aber vor der Welt sind wir nur arme Pilgrimme in der That.

51. Diesen müssen die Gottseligen mit Fleiß bedenken und ihr Leben nach dieser Weise anrichten. Denn Wenige wollen dasselbe betrachten. Die Meisten folgen dem Nuzen und Lust dieser Welt nach und suchen denselben; welches sie für das Paradies und für ihr Himmelreich halten und lassen die Hoffnung des andern bessern Lebens immer hinfahren.

52. Danach kann man auch aus diesem Texte die Zeit der Geburt Josephs rechnen. Denn der Heilige Geist hat nicht vergeblich haben wollen, daß die Jahre der Väter sollten gezählet werden; sondern es ist darum geschehen, daß die Gottseligen eine gewisse Ordnung und Rechnung der Zeit hätten. Darum wird so nachdrücklich und so genau die Zahl der Jahre im Texte gesetzt. Derhalben ist Joseph jetzt bei neun und dreißig Jahren, Jakob aber hundert und dreißig. Daraus kann man nun leichtlich rechnen, wann Joseph ist verkauft worden; wie wir das in unserm Chronikon angezeigt haben. Denn es gefällt GOtt wohl, daß man die Zeit in dieser Welt weiß. Und wiewohl wir die Stunde und Augenblicke nicht wissen, zeigt uns doch GOtt durch die Zeichen, so Christus selbst erzählt Luc. 21, V. 25. ff., die Zeit, wann sich das jüngste Gericht nahen wird.

53. Endlich geht Jakob vom Könige wiederum heraus mit Danksagung, wie er ihn anfänglich auch freundlich gegrüßt, und sich gegen ihn erklärt, daß er ihm von Herzen dankbar sei für seine königliche Mildigkeit, so er ihm und den Seinen erzeigt hätte. Er gehet ihn nicht nach der Weise eines Großen, sondern eines Kleinern. Da sind nicht solche Bitten, darin etwas gegeben oder verheissen wird; sondern sind solche Segen, darin man etwas wünscht. Denn wir bitten, daß GOtt den Segen geben wolle, aber wir schließen oder gebieten damit nichts; denn dasselbe gebühret nur allein GOtt. Wie hernach Jakob seine Söhne segnet mit gewissen wirklichen Segen, darin er etwas Gewisses bestimmt, anzeigt und ihnen verkündigt.

\subsection{II}

\begin{quote}
\textbf{V. 11. Aber Joseph schaffte seinem Vater und seinen Brüdern Wohnung, und gab ihnen ein Gut in Egyptenland, am besten Ort des Landes, nämlich im Lande Raemses, wie Pharao geboten hatte.}
\end{quote}

54. Hier beschließet nun Mose die Reise, wie Jakob nach Egypten hinab gezogen sei, und will nun wieder kommen auf die Theurung, so dazumal noch im ganzen Königreiche war. Denn Joseph richtet den Befehl Pharaos aus, und setzt seinen Vater und seine Brüder in ein fruchtbar Land, nämlich, zu Raemses, davon die Geographen sagen, das es in Gosen sei, von wo sie hernach wieder ausgezogen sind, wie das im 2. Buch Mose zu sehen ist, da sie das erste Osterlamm gegessen haben. Da hat Jakob und seine Söhne mit ihrem Viehe gewohnt durch Mildigkeit des Königs und Josephs.

\begin{quote}
\textbf{V. 12. Und er (Joseph) versorgte seinen Vater, und seine Brüder, und das ganze Haus seines Vaters; einen jeglichen, nachdem er Kinder hatte.}
\end{quote}

55. Es ist ein Hebraismus, wenn es heißt: „gemäß dem Munde des Kindes.“ Da müssen wir rathen, was eigentlich damit bedeutet oder verstanden wird. Und erstlich findet hier eine Vertauschung des Numerus statt, da das Wort „des Kindes“ gesetzt wird für „der Kinder.“ Er will aber damit so viel sagen, daß Joseph seinen Vater und Brüder gespeiset habe nach der Weise und Maß, wie man die kleinen Kinder pflegt zu speisen; wie es im Hausregimente zugeht, da die kleinen Kinder der Arbeit ihrer Eltern mit genießen, essen das Brod; aber sie selbst arbeiten nichts, erwerben das Brod nicht, sondern helfen es essen: sie bauen den Acker nicht, thun auch sonst gar keine Arbeit, sondern essen mit, was die Andern mit ihrer Arbeit und Fleiß erworben haben. Und da Paulus 2 Thess. 3, 10. sagt: „Wer nicht arbeitet, soll auch nicht essen“, sollen wir wissen, daß gleichwohl daselbst die kleinen Kinder ausgenommen sind, die auch mit essen, ob sie wohl nicht arbeiten; ja, das noch mehr ist, sie thun ihren Eltern noch wohl allerhand Verdruß dafür, daß sie von ihnen so treulich, und mit so großer Mühe und Arbeit auferzogen werden.

56. So ist nun „nach Art der Kinder speisen“ so viel, als wollte man sagen: Einen ohne seine Arbeit speisen oder versorgen. Denn Jakob und seine Söhne mit ihrem ganzen Hausgesinde haben keine Aecker gehabt, die sie hätten bauen können, sind nur Viehhirten gewesen, haben von Joseph das Getreide oder Speise umsonst gehabt, auf Befehl und aus Mildigkeit des Königs. Also hat er nun seinen Vater und Brüder versorgt und ernährt, wie man die kleinen Kinder pflegt zu nähren, daß sie dagegen keine Arbeit thun dürfen.

57. Nun ist aber Jakob nach Egypten geführt worden, wiewohl er dasselbe Land nicht geliebet und nicht gern hinein gezogen ist; sondern dieweil ihn die Theurung gezwungen und er Verlangen gehabt, seinen Sohn zu sehen. Denn er denkt an die Verheißung, daß die Kinder Israel in demselbigen Lande hart sollen geplagt und umgetrieben werden, und weiß, daß nun eigentlich das vierte Mannesleben vorhanden ist. Er wird ja wohl mit Freuden empfangen und aufgenommen, wie Christus am Palmsonntag zu Jerusalem auch mit Frohlocken empfangen worden ist, aber hernach ist man mit ihm auch umgegangen, wie mit Christo am Charfreitage. Darum begehret er mit den Seinen nicht lange in Egypten zu bleiben, sondern wünscht, daß er so bald als möglich wiederum ins gelobte Land kommen möge.

\section[Dritter Theil]{Dritter Theil. Wie Joseph alles Geld, Vieh und Acker der Egypter zusammen bringt, die Egypter leibeigen macht, und diese Leibeigenschaft bestätigt}

\begin{quote}
\textbf{V. 13. 14. Es war aber kein Brod in allen Landen; denn die Theurung war fast schwer, daß das Land Egypten und Canaan verschmachteten vor der Theurung. Und Joseph brachte alles Geld zusammen, das in Egypten und Canaan funden ward, um das Getreide, das sie kauften; und Joseph that alles Geld in das Haus Pharao.}
\end{quote}

58. Mose kommt nun wieder auf die Beschreibung der Theurung. Es ist aber nun oftmals gesagt worden, daß in diesem ganzen Buche gerühmt wird die wunderbare Regierung GOttes, daß gemeiniglich allwege zur Zeit der größten Väter und Propheten die größten Theurungen gewesen sind; wie das zu sehen ist in der Historie Abrahams, Isaaks und Jakobs; darum daß GOtt die Welt mit Theurung straft, wo das Wort in der Kirche in Schwange geht. Jedoch macht er es mit der Strafe also, daß er gleichwohl die Seinen noch nährt und erhält; wie der Spruch lautet im 37. Psalm V. 19.: „In der Theurung werden die Frommen genug haben.“ Und dasselbe zeigt dies gegenwärtige Exempel auch an. Also ist zur Zeit Elias und Elisae, desgleichen auch Davids große greuliche Theurung gewesen.

59. Darum ist es ein wunderlich Ding, wenn GOtt den Menschen die höchste Gnade erzeigt damit, daß er ihnen sein Wort sendet, da diese Welt keinen theureren Schatz hat, denn derselbe ist, und daß er doch unter solchem hellen Lichte des Wortes noch dabei Theurung, Pestilenz und viele andere unzählige Plagen oder Strafen schickt.

60. Aber damit versucht er also seine Heiligen und Gläubigen und stärkt auch zugleich damit die Undankbarkeit der Welt. Wir werden zu unsern Zeiten reichlich begnadet mit der Predigt des lieben Wortes Gottes; aber wie dasselbe von den undankbaren Zuhörern angenommen wird, das sehen wir auch wohl. Denn es reißt mit Gewalt herein und kommt in die Kirche die Tyrannei der Fürsten, der Wucher, Trügerei, Räuberei, so die Leute unter einander insonderheit üben und auch öffentlich insgemein betrieben wird. Durch solche Schinderei und Räuberei wird ganz Deutschland ausgezogen; und wo die Leute also werden fortfahren, ihren eigenen Fürsten den Zaum zu lassen und solchen Muthwillen unter einander zu üben, so wird ohne Zweifel eine greuliche Veränderung hernach folgen. Es wird Germania viel eine andere Gestalt gewinnen, soll es lange also stehen. Denn es wird dermaleinst dazu kommen, daß die vielen Jahrmärkte und die große Üppigkeit mit dem übermäßigen Bauen, Kleidung, Speise und anderem Gepränge wird niedergelegt werden; aber das wird nicht geschehen ohne großen Schaden und Verderben Deutschlands.

\quellseite{465}
61. Die Welt versteht es nicht, dankt auch Gott nicht für die unermesslichen und unzähligen Wohlthaten, damit sie täglich von Gott überschüttet wird; und die Sünden der Menschen sind, leider! allzuviel und unmäßig groß. Denn man sündigt jetzt nicht aus Unwissenheit, wie unter dem Papstthum geschehen ist: wir wissen und erkennen jetzt zu unserer Zeit die Wahrheit. Darum ist am Verstande der Irrthum nicht, ist auch kein Mangel an dem, daß man es nicht sollte wissen; sondern wider allen Verstand wollen sich die Leute dazu nicht bewegen lassen, daß sie sich möchten zur Besserung schicken und fromm werden, wie treulich sie auch gelehrt und ermahnt werden. Denn wir arbeiten ja täglich daran, daß wir sie mit Lehren und Ermahnen in diesem schönen großen Licht und Reichthum des Evangeliums unterrichten und erwecken mögen, Buße zu thun. Aber da hilft kein Wehren: je mehr man predigt, je ärger man wird.

62. Die Papisten verfolgen ohne Aufhören die erkannte Wahrheit und toben wider die Gottlosen, als wären sie unsinnig, wie zu den Zeiten der Propheten: sie sind entzündet von Wahnwitz und Thorheit. Darum werden sie über ihrem Vornehmen und Begierde noch immer unsinniger, die Christen zu vertilgen. Also will Germania hindurch und Deutschland folgt darin dem Königreich Israel und Juda nach. Denn wiewohl in denselben Königreichen die Propheten mit Gottes Wort klärlich und auf das allersteißigste haben angehalten, richteten sie doch damit gar nichts aus. Es half alles nichts. Wie ihnen der Herr solche greuliche Unsinnigkeit aufgerückt und sie deshalb straft, Jer. 2, 23. 24., da er sagt: „Siehe an, wie du es treibst im Thal, und bedenke, wie du es ausgerichtet hast. Du läufst umher wie eine Kameelin in der Brunst; und wie ein Wild in der Wüste pflegt, wenn es vor großer Brunst lechzet und läuft, daß niemand aufhalten kann.“ Du läufst dahin, wie ein toller Hund. Gott hat bei ihnen dies Wort oftmals wiederholt: Bekehre dich doch und sündige nicht also, höre doch mein Wort. Aber das alles haben sie nicht hören wollen, bis der König aus Assyrien kam, der ihren tollen Lauf aufgehalten und ihnen denselben verboten hat. Also werden jetzt unsere Deutschen auch nicht aufhören zu sündigen, bis dermaleinst ein Feind aus fremden Landen daherkommen, und ihr gottlos Wesen und das Toben und Wüthen hindern und abschaffen wird.

63. Wir aber, die das Wort haben und lieben, sollen beten und seufzen, daß sich Gott unser erbarmen wolle, oder je zum wenigsten den Weizen in seine Scheunen sammeln, und daß er uns ja mit der Spreu nicht in das ewige Feuer werfe und verbrenne. Und Gott hat schon jetzt angefangen zu dreischen und seine Tenne zu fegen. Was es aber mit der Spreu und Stoppeln für ein Ende nehmen wird, dasselbe werden endlich diejenigen wohl fühlen, so sich jetzt nicht wollen ermahnen lassen und sich gar nicht zu bessern gedenken. Wir sind wohl entschuldigt.

64. Also hat zu der Zeit das Königreich Egypten einen sehr trefflichen Propheten gehabt, der alle Dinge sehr fein geordnet, so viel das Priesterthum und weltliche Aemter belangt hat; der das ganze Königreich hat beten gelehrt und, wie sie den rechten wahren Gott anrufen sollten; hat ihnen dazu Gesetze gegeben, wie sie recht und ehrlich leben sollten. Aber sie werden dieselbige Lehre ohne allen Zweifel angenommen und sich dabei also gehalten haben, wie die Welt pflegt. Wie solches das Ende auch fein ausgewiesen hat, da die große und schwere Theurung gewesen ist.

\subsection{I}

65. Dies ist aber eine wunderbare Beschreibung, die ich nicht genugsam verstehen kann. Denn wir sind nicht gegenwärtig dabei gewesen, haben auch die Korn und wie das Königreich Egypten ist bestellt gewesen, nicht vor Augen gesehen. Und ich kann nicht denken, wie doch in Deutschland ein solch Reich und Regiment könnte geordnet oder angerichtet werden. Denn, erstlich, wird alles Geld zusammengebracht; darnach, das Vieh; und zum Dritten, werden auch die Leute von dem Könige leibeigen gemacht. Darum wird es müssen ein schön wohlgeordnet Königreich gewesen sein, dieweil in der Theurung das Geld, Vieh, Aecker und Leute sind verkauft worden. Und ist wahrlich eine große schwere Knechtschaft gewesen; und wie gesagt, ich kann nicht verstehen, wie doch in Deutschland eine solche Weise möchte vorgenommen werden ohne Empörung und Aufruhr des gemeinen Pöbels.

66. Es erhellt daraus, daß Joseph mit sonderlichem Fleiß und Ernst über guter Ordnung, Zucht und Gehorsam müsse gehalten haben, und daß er durch das Amt der Priester und Lehrer das Volk bei rechter Lehre und in Gottesfurcht gehalten. Denn dies ist eine sehr große Schatzung, so dazumal das Volk hat geben müssen.

67. Die Hebräer aber sind mit einander nicht eins über der Auslegung des Worts thelah. Denn Etliche haben es im Lateinischen gegeben „unsinnig werden“; die Andern, „bedrücken“. Wir halten es dafür, daß man es recht deuten möge „verschmachten“, wie wir es im Deutschen auch darnach verdolmetscht haben: „verschmachteten vor der Theurung.“ Und es scheint, daß die Theurung in andern Landen so greulich und so groß nicht gewesen sei, wie in Egypten und Canaan, welche Lande das Wort gehabt, so ihnen Jakob und Joseph vorgetragen haben.

68. Da es nun im ganzen Königreich am Getreide gemangelt, hat Joseph alles Geld zusammengebracht, das in Egyptenland gefunden ward; denn die Unterthanen haben dasselbe zusammengetragen, Getreide dafür zu kaufen und ihr Leben also zu retten. Denn die Theurung, so sieben ganze Jahre währet, nimmt alles Getreide hinweg, wie großen Vorrath man desselben immer haben mag. Diejenigen, so für sich selbst etwas in ihre Scheunen gesammelt hatten, haben vielleicht dazumal dieselben schon leer gemacht und ihren Vorrath gar verzehrt, daß sie von Joseph haben Getreide kaufen müssen. Also ist nach drei oder vier Jahren auch fast das Geld der gemeinen Unterthanen in Egypten und Canaan ausgegeben worden, damit sie sich der Theurung erwehret haben.

69. Und dasselbe Geld hat Joseph zusammengebracht, nicht mit Wucher, Räuberei, oder unbilliger, unrechtmäßiger Schatzung; sondern nur damit, daß er die vorigen Jahre in ganz Egypten das Getreide gesammelt auf Kornhäuser, so hin und wieder im Königreich aufgerichtet gewesen, nicht allein von den Aeckern des Königs, sondern auch der gemeinen Unterthanen, die das Getreide in den wohlfeilen Jahren verkauft, dessen sie für ihr Hausgesinde nicht bedurft. Dazu hat Joseph des Königs Geld gebraucht und ausgegeben, welches er nun wieder eingenommen, da er den Leuten das Getreide verkauft hat.

70. Dies ist wahrlich eine große wunderbare Beschwerung gewesen, dadurch mit der Frucht oder Gemächs des Landes auch das Geld der Unterthanen weggegangen ist. Wir haben zu unsern Zeiten dergleichen nicht gesehen, unsere Vorfahren auch nicht. Und sind doch daselbst, in Egypten, Leute gewesen, die sehr mäßig gelebt haben, die genau und kümmerlich haushgehalten haben: sind keine Schlemmer, Fresser und Säufer gewesen, wie wir Deutschen sind, deren Einer in einem Tage so viel Speise und Tranks verthut, damit sich hundert Egypter zur Nothdurft hätten behelfen können.

71. Denn, lieber, siehe doch, was in diesem unserm Städtlein geschieht, da die Bürger nach gehaltener Rechnung gefunden haben, daß alljährlich mehr denn vier tausend Gulden ausgegeben werden für Gerste. Was ists doch, daß man das Geld so gar unnütz verthun soll? Wir kaufen Tag und Nacht und füllen den Bauch mit Bier. Wenn wir aber Lust hätten zur Sparsamkeit und Mäßigkeit, wie wir zum Ueberfluß haben, so könnten wir alle Jahre zwei- oder dreitausend Gulden ersparen und behalten. Wie viel Wein kaufen aber die vollen Brüder hinweg außer dem Bier? Was geht auf den Luxus in Kleidern und anderen unnützen Dingen, so unsere Kaufleute aus der Fremde hierher bringen in unsere Lande? Ja, wie viel Geld fressen die Franfurter Meßen, da in einer jeglichen, wie gesagt wird, an die dreißigmal hunderttausend Gulden aus Deutschland geführt werden. Ich will der Leipziger Märkte und der andern schweigen. Und es scheint doch, daß dies alles noch nicht zu vergleichen sei mit dem großen Ueberfluß der Kleidung, des Weins, Biers und anderer Dinge, die wir ohne allen Nutzen schändlich und übel durchbringen.

\quellseite{466}
72. Es haben wahrlich die Wucherer eine sehr harte Reformation angerichtet; und wo unser Herr GOtt also fortfahren wird mit uns zu zürnen, werden sie uns so gar ausgehen, daß wir kaum das Wasser behalten werden zu täglicher Nothdurft. Darnach kommen dann auch dazu die Schatzungen, so die weltliche Obrigkeit den Unterthanen auflegt. Darum wird Deutschland jetzt zur Zeit hart beschwert und aufgezogen, welches wahrlich bis daher mehr Freiheit gehabt, denn andere Königreiche und Lande, als Frankreich, Spanien, Italien etc. Aber jetzt bricht die Schatzung und Knechtschaft mit Gewalt herein. Denn die Leute wollen sich um der Sünde willen nicht mehr strafen lassen; und wenn die Sünden also frei im Schwange gehen, so ist es auch unmöglich, daß man die Strafe aufhalten und derselben wehren könne, daß sie nicht überhand nehme. Darum mögen wir die Strafe leiden und den Schaden tragen, dieweil wir von Sünden nicht ablassen und uns nicht bessern wollen.

73. Wie Augustin gar herrlich davon sagt, da er den 85. Psalm auslegt, darin V. 11. also steht: „Gerechtigkeit und Friede küssen sich.“ Darüber sagt Augustin: Sei fromm und übe Gerechtigkeit, so wirst du Friede haben; auf daß sich also Gerechtigkeit und Friede küssen. Denn wo du die Gerechtigkeit nicht lieben wirst, so wirst du auch nicht Frieden haben. Denn diese beiden haben einander lieb und küssen sich, daß wer Gerechtigkeit übt, auch Friede finden soll, welcher die Gerechtigkeit küßt. Es sind zwei gute Freundinnen. Du willst vielleicht gerne die eine haben, befleißigst dich aber der andern gar nicht; denn niemand ist, der nicht gerne wollte Frieden haben, sie wollen aber nicht alle fromm sein und Gerechtigkeit üben.

74. Dieweil es uns nun also gelüstet, daß wir frei und ungestraft sündigen wollen, und also wider GOtt toben und wüthen, so mögen wir dagegen auch die Strafe tragen. Denn Ungerechtigkeit und Unfriede füllen sich auch. Die Sünde und die Strafe über die Sünde hängen auch an einander, wie Paulus sagt Röm. 6, 23.: „Der Tod ist der Sünden Sold“; und: „Der Stachel des Todes ist die Sünde“, 1 Cor. 15, 56. Nach dem Willen unsrer bösen Lust und Begierden muß man auch das Wüthen der Strafe leiden.

75. Wie jetzt zu unserer Zeit unter den Leuten gar keine Zucht, keine Gerechtigkeit und Scham mehr ist. Wir predigen, rufen und schreien, halten immer an, zur rechten Zeit und zur Unzeit, wie Paulus redet 2 Tim. 4, 2. Die Obrigkeit aber schweigt stille dazu und sieht durch die Finger bei den Lastern. Wir sagen was wir wollen, die Obrigkeit schweiget still: es will nirgends fort. Darum muß auch die Strafe darauf folgen. Da werden wir dann rufen, aber GOtt wird uns, wenn wir in der Strafe zu ihm schreien, nicht hören; denn wir haben ihn auch nicht hören wollen in unsern Sünden, da er uns durch seine Diener hat gerufen. Wir wollen die Sünden hinaus führen, so müssen wir die Strafe auch hinaus haben.

76. Derhalben ist es ein großer Jammer gewesen, daß Joseph allemal dem Unterthanen Geld alles zusammen gebracht hat, nicht zwar durch Schatzung, sondern daß er ihnen Getreide verkauft, das sie selbst gesammelt, getheilt und verkauft hatten. Und diese Sammlung des Geldes, so Joseph vorgenommen hat, ist billig und rechtmäßig gewesen; denn er hat des Königes Getreide verkauft, dafür er billig und mit Recht auch das Geld gefordert. Und er hatte ihnen diese Theurung lange zuvor verkündigt, und verordnet, daß ein Jeder so viel sammeln sollte, was er zu täglicher Nothdurft würde haben müssen auf ganze sieben Jahre. Aber Etliche werden dem, was er verkündigt hat, geglaubt, die Andern aber werden es gar verachtet haben. Und ist doch eine große Gnade gewesen, daß sie noch Getreide um Geld bekommen können, welches ihnen Joseph ohne Zweifel das erste, das andere und auch die nächstfolgenden Jahre zu gleichem Preise wird verkauft haben. Dies ist ein ordentlicher und rechtmäßiger Contract und Kauf gewesen, daß er ihnen das Getreide also unter einem seinen ehrlichen Titel verkauft hat, und hat sie nicht alle Monden und Jahre mit neuer Schatzung oder Schinderei übertheuert: er hat es gemacht, daß die Leute haben mögen zukommen. Aber mit der Zeit, dieweil die Theurung so lange gewähret, sind gleichwohl die gemeinen Unterthanen allmählich in ihrer Habe und Gut erschöpft und ausgezogen worden. Und es hat nicht anders können zugehen; denn man hat ja in solcher Noth den armen Unterthanen helfen müssen, daß sie nicht durften Hungers sterben.

\subsection{II}

\begin{quote}
\textbf{V. 15—17. Da nun Geld gebrach im Lande Egypten und Canaan, kamen alle Egypter zu Joseph, und sprachen: Schaffe uns Brod. Warum lässest du uns vor dir sterben, darnach, daß wir ohne Geld sind? Joseph sprach: Schaffet euer Vieh her, so will ich euch um das Vieh geben, weil ihr ohne Geld seid. Da brachten sie Joseph ihr Vieh; und er gab ihnen Brod um ihre Pferde, Schafe, Rinder und Esel. Also ernährete er sie mit Brod das Jahr um alle ihr Vieh.}
\end{quote}

77. Da nun das Geld zusammen gebracht war vom Getreide, so Joseph den Cananitern und Egyptern verkaufte, da hat der König einen herrlichen Schatz gehabt. Ob aber die Leute auch noch etwas übrig behalten, weiß ich nicht; denn ich kann nicht denken, daß die Unterthanen von all ihrem Gelde gänzlich sollten sein entblößt worden und so gar ausgezogen sein, daß sie gar nichts mehr gehabt hätten, davon sie ihren Tagelöhnern, Handwerksleuten, und sonst Andern, die sie täglich zur Arbeit gebraucht, ihren Lohn hätten geben können. Oder man möchte es vielleicht also verstehen, daß es synecdochisch geredet wäre, nämlich, daß er alles Geld hat zusammen gebracht, ausgenommen, was sie zu ihrer täglichen Nothdurft bedurft haben, Kleider, Schuhe und andere Nothdurft zu kaufen. Oder sie werden vielleicht mit Getreide, das sie gekauft hatten, ihren Ackersleuten und Arbeitern gelohnet haben. Darum bringen sie jetzt nun ihr Vieh, dafür ihnen Joseph Brod gegeben und sie damit ein Jahr lang versorgt hat.

78. Es drückt aber der Text nicht aus, welches Jahr dasselbe gewesen sei; sondern nachdem das vorige Jahr das Geld verzehrt, bringen sie nun auch das Vieh herzu: nicht daß der König ihr groß und klein Vieh an einen Orte habe versammeln lassen, sondern er redet vom Preise des Viehes, so sie selbst behalten und geweidet haben. Wie solches in unsern Contracten oder Verträgen auch pflegt zu geschehen, als, wenn man die Einkünfte verkauft von einem liegenden Gute, das jährlich Frucht trägt, und behält doch selbst das Eigenthum daran. Darum haben sie das Jahr ihr Vieh verzehrt oder Geldschaß davon gegeben. Und ich halte dafür, es sei das sechste Jahr gewesen und bis auf eins das letzte, nachdem im fünften Jahre das Geld verzehrt gewesen. Die ersten zwei Jahre haben sie vielleicht von dem Vorrath leben können, so sie in ihre eigenen Scheunen gesammelt hatten; desgleichen das dritte Jahr auch, außer, daß es etwas knapper zugegangen ist; im vierten Jahre aber ist das Kaufen und Verkaufen angegangen.

\subsection{III}

79. Es ist wahrlich eine greuliche Theurung und wunderbarlich Regiment gewesen, so Joseph damit gehalten, daß dadurch alle Habe und Gut, dazu auch die Leute selbst für ihre Person durch das ganze Königreich in des Königs Gewalt gekommen und alle seine Knechte geworden sind; also daß sie gar nichts Eigenes gehabt haben, weder daheim noch draußen im Felde, noch auch am Viehe. Und es scheint doch, daß diese Armuth, da sie alles Dinges haben mangeln müssen, mit ihnen nicht lange gewähret habe. Denn sie werden hernach wiederum so groß und reich, daß sie viel Gold und Silber gehabt haben, da das Volk Israel aus Egypten gezogen und sie desselben beraubt hat. Darum hat diese Theurung nicht allezeit gewähret, außer daß in Egypten der Gebrauch geblieben ist, daß sie den fünften Theil von ihren Gütern haben geben müssen.

\begin{quote}
\textbf{V. 18—20. Da das Jahr um war, kamen sie zu ihm in andern Jahr, und sprachen zu ihm: Wir wollen unserm Herrn nicht verbergen, daß nicht allein das Geld, sondern auch alles Vieh dahin ist zu unserm Herrn; und ist nichts mehr übrig vor unserm Herrn, denn nur unser Leiber und unser Feld. Warum lässest du uns vor dir sterben, und unser Feld? Kaufe uns und unser Land uns Brod, daß wir und unser Land leibeigen seien dem Pharao; gib uns Samen, daß wir leben und nicht sterben, und das Feld nicht verwüste. Also kaufte Joseph dem Pharao das ganze Egypten. Denn die Egypter verkauften ein jeglicher seinen Acker, denn die Theurung war zu stark über sie. Und ward also das Land Pharao eigen.}
\end{quote}

\quellseite{467}
80. Das andere Jahr ist nicht das, so dem ersten am nächsten gefolgt, sondern ist das Jahr, so gefolgt ist nach dem, darin sie dem Könige das Vieh gegeben. Das ist aber eine große, schwere und die alleräußerste Noth gewesen, daß sie klagen, sie haben nichts mehr übrig, denn nur ihren Leib und ihr Feld. Das nimmt Joseph auch hinweg; und doch ist dasselbe keine Räuberei, sondern sie haben es verkauft und ist ein rechtmäßiger Contract gewesen. Sie haben auch den Gebrauch ihres Feldes nicht verloren, sondern haben nur dem Könige den fünften Theil von den Früchten, so ihnen darauf gewachsen sind, geben und entrichten müssen. Wiewohl es also scheint, als sei das Feld geradezu verkauft worden, und daß der König diejenigen, so das Feld gebauet, für seine Knechte gebraucht habe, daß sie ihm die Aecker haben bauen müssen.

81. Dies Regiment ist zwar hart und schwer genug, aber doch rechtmäßig und nicht tyrannisch gewesen, und da niemand unrecht geschehen ist: und ist Gottes Zorn und Strafe über ihre Sünde gewesen, dadurch Gott dies Volk also mit Theurung gestraft, und sie also zu Knechten gemacht hat, welche Knechtschaft sie um der Sünde willen haben tragen müssen. Denn wo wir sündigen, sollen wir wissen, daß man dafür auch muß Strafe leiden. So haben sie nun nichts Eigenes mehr, weder Geld, noch Vieh, noch Aecker, und auch nicht ihre eigenen Leiber, der König hat Alles eingenommen; nicht daß er es ihnen mit Gewalt geraubt oder genommen hätte, sondern hat den Unterthanen dafür Getreide verkauft. Und dies ist bis auf eins das letzte und darnach auch gar das letzte Jahr gewesen.

\begin{quote}
\textbf{V. 20. 21. Und ward also das Land Pharao eigen. Und er theilete das Volk aus in die Städte, von einem Ort Egyptens bis ans andere.}
\end{quote}

82. Dies ist nun etwas Neues und Sonderliches, daß Joseph wie ein Siegel dazu gethan, das Kaufen und Verkaufen, so sie gethan, damit zu bestätigen, nämlich, daß die Unterthanen im Königreich also von einem Ort zum andern haben ziehen müssen, auf daß sie wissen sollten, daß all ihr Hab und Gut des Königes wäre, wegen des Kaufs, so sie mit einander abgeschlossen hatten. Er läßt aber das Volk aus einer Stadt in die andere ziehen zum Zeugniß und öffentlichen Wahrzeichen, damit er die Rechtsschaft des Volks versiegelt, als, wenn jetzt die Bürger von Wittenberg nach Thüringen und die Thüringer nach Sachsen ziehen müßten. Da hat er ihnen Aecker und Vieh zugestellt, so des Königes gewesen, doch also, daß sie den fünften Theil der Einkunft dem Könige, als von seinen liegenden Gründen, haben geben müssen.

83. Das ist wahrlich eine wunderliche und schwere Versiegelung gewesen, dadurch sie erinnert worden sind, daß sie das Eigenthum aller ihrer Güter verloren hatten, nicht durch Josephs oder des Königs Schuld, sondern wegen des Kaufs, so zwischen ihnen geschehen war. Und das ist zwar noch eine große Gnade und sonderliche Wohlthat gewesen, dieweil sie leibeigene Knechte waren, daß er sie doch seines Gefallens nicht getödtet, sondern genährt und erhalten, und dazu noch auf die königlichen Aecker und Güter gesetzt hat. Darum hat es dazumal in diesem Königreich also gestanden, daß es um der Sünden willen die Strafe und Züchtigung Gottes hat tragen und leiden müssen, so ihnen Joseph lange zuvor verkündigt, und sie dazumal auch durch die Priester und Fürsten oder Amtleute im Lande hat unterrichten und ermahnen lassen, daß sie ja diese Knechtschaft geduldig leiden wollten.

\begin{quote}
\textbf{V. 22. Ausgenommen der Priester Feld, das kaufte er nicht; denn es war von Pharao für die Priester verordnet, daß sie sich nähren sollten von dem Benannten, das er ihnen gegeben hatte; darum durften sie ihr Feld nicht verkaufen.}
\end{quote}

84. Dies ist ein trefflicher Text, den man sonderlich wohl merken soll. Denn was jetzt zur Zeit wir Deutschen denken oder was wir für Leute sind, weiß ich nicht. Das ist ja gewiß, daß im natürlichen und göttlichen Recht also verordnet ist, daß die dem Altar dienen, auch vom Altar leben sollen; wie Paulus 1 Cor. 9, 13. 14. sagt: „Wisset ihr nicht, daß, die da opfern, essen vom Opfer? und die des Altars pflegen, genießen des Altars? Also hat auch der Herr befohlen, daß, die das Evangelium verkündigen, sollen sich vom Evangelio nähren.“ Und an diesem Orte ist im Text mit großem Fleiß geschrieben und zum andernmal wiederholt, daß den Priestern in Egypten ihr sonderlicher beschiedener Theil verordnet und benannt gewesen sei, daß sie ihr Feld nicht haben verkaufen dürfen.

85. Darum ist dies ein sehr schöner Text, aber wir folgen ihm nicht nach und lassen denselben in Deutschland nicht zu. Man pflegt es zwar in andern Ländern also zu halten, daß die, so öffentliche Aemter verwalten, als, der Herren Knechte, Stockmeister, Kriegsfrnechte und die Oberherren, selbst auch von öffentlicher Besoldung unterhalten werden. Desgleichen nehmen auch die Fürsten Zoll und Schoß von den Unterthanen. Aber wie es mit den Schulleuten und Dienern des Worts gehalten wird, das sehen wir vor Augen. Wenn nicht noch des Raubes aus Egypten etwas überblieben wäre, so wir dem Papst geraubt haben, so müßten wir wohl alle Hungers sterben. Denn es ist jetzt keine Stadt oder Fürstenthum, die ihre Priester und Schulen nähret oder unterhalten, unangesehen die große schwere Arbeit, dazu den treuen Fleiß und Dienst frommer gottseliger Pfarrherren und Seelsorger. Und wo sie von der Zulage des gemeinen Volks sollten unterhalten werden, so würden sie wahrlich ärmlich und kümmerlich leben müssen. Darum sage ich noch, daß wir von dem Raube Egyptens genähret werden, so unter dem Papstthum gesammelt worden; und wir doch daselbst, so davon noch übrig ist, von der Obrigkeit auch hingerrißen; die Pfarrherren und Schulen werden beraubt, nicht anders, denn als wollten sie uns Hungers sterben lassen. Wird es nun so bestehen, so wird man es wohl sehen.

86. Pharao, der König in Egypten, wird am zukünftigen jüngsten Gericht aufstehen und wird die Fürsten und Obrigkeiten in Deutschland verdammen, darum daß er seine Priester in Ehren gehalten und genährt, sich auch an ihren Gütern nicht vergriffen hat, so er doch dieselben vermöge des Contractes, so da rechtmäßig aufgerichtet war, an sich in seine Gewalt hätte nehmen mögen; sondern er gibt ihnen Freiheit und verordnet ihnen dazu einen benannten Theil Getreide und Speise, davon sie ohne allen ihren Schaden gelebet haben. Welchen wollt ihr mir nun unter allen Königen und deutschen Fürsten geben, der diesem Pharao gleich sei?

87. Unser Landesfürst nährt uns von dem Seinen nicht, sondern vom Raube Egyptens, und würde uns auch von dem Seinen nicht nähren oder halten können. Daß man aber den Priestern und Kirchendienern schuldig sei, sie zu nähren, daran ist kein Zweifel; und wo wir dazu nichts geben werden, so wird uns dieser Text verdammen; davon wir denn große Schande haben werden. Gott will haben, daß, der unterrichtet wird mit dem Wort, allerlei Gutes mittheilen soll dem, der ihn unterrichtet, Gal. 6, 6. Und Paulus setzt daselbst V. 7. noch eins hinzu und sagt: „Irret euch nicht; Gott läßt sich nicht spotten.“ Desgleichen sagt auch Christus zu seinen Aposteln Luc. 10, 7.: „Wo ihr in ein Haus kommet, in demselbigen bleibet, esset und trinket, was sie haben. Denn ein Arbeiter ist seines Lohnes werth.“ Was sollen wir aber thun? Sie haben und hören uns auch nicht gern; und wenn wir auch unsern Zuhörern gleich Besoldung und Geld zugeben wollten, so will man uns doch nicht gern hören oder leiden.

88. Darum werden auch die Sünden in Deutschland reif, und die Strafe wird sich nicht säumen und wird nicht ausbleiben. Denn es werden schon jetzt die Fürsten von den Bauern gar ausgezogen, und kommt täglich mancherlei Schätzung und Schinderei auf, und geschiehet uns also daran eben recht. Was man unserm Herrn Gott nicht gibt, das gebe man dem Teufel; wie man sagt: Quod non tollit Christus, tollit fiscus: Was nicht wegnimmt Christus, nimmt hinweg der Fiscus. Willst du dem Schüler nicht geben, so gib dem Bruder Veit und Scharbhansen, die es mit Haufen verschwenden und gehen, dreysig, fünfzig Gold aus einmal nehmen. Da man Gott zu ehren und zu unserer Seelen Seligkeit einem armen Schüler oder Diener des Worts nicht gern einen Groschen geben will, da wirst du hernach einem gottlosen Kriegsknechte und dem Satan selbst tausend Gulden geben müssen.

89. Derhalben, sage ich, soll man diesen Text fleißig merken, nämlich, daß der so fromme König in dieser höchsten Beschwerung, so dazumal auf dem Lande gelegen, doch das Wort also geehret und den Gottesdienst erhalten hat, da er die Priester von der großen Noth und sehr schweren Last der Theurung und Knechtschaft befreit. Er erhält ihre Kinder und Gefinde, und gibt ihnen nicht allein Freiheit, sondern verordnet auch einem jeglichen sein bescheiden Theil Getreide, damit sie also genährt werden mögen. Es ist wahrlich ein sehr gottseliger König, desgleichen wir zu unsern Zeiten nicht haben. Ja, unsere Fürsten rauben und reißen alles hinweg, was sie in Klöstern und Pfarren finden, und sind ihrer sehr wenig, die etwas dazu geben, damit das Predigtamt erhalten werden möchte, oder daß doch etwa ein einiger Schüler unterhalten werden möchte. Darum wird es dazu kommen, daß sie wiederum auch darben werden. Es wird dahin kommen, du wirst auch zu wenig haben.

\quellseite{468}
\section[Vierter Theil]{Vierter Theil. Wie Joseph den Egyptern ihre Aecker wiedergibt; und ob der Priester Feld dem Pharao eigen geworden}

\subsection{I}

\begin{quote}
\textbf{V. 23—25. Da sprach Joseph zu dem Volk: Siehe, ich habe heute gekauft euch und euer Feld dem Pharao; siehe, da habt ihr Samen, und besäet das Feld. Und von dem Getreide sollt ihr den fünften Pharao geben; vier Theile sollen euer sein, zu säen das Feld, zu eurer Speise, und für euer Haus und Kinder. Sie sprachen: Laß uns nur leben, und Gnade vor dir, unserm Herrn, finden; wir wollen gerne Pharao leibeigen sein.}
\end{quote}

90. Hieraus erhellt nun, daß da sechs Jahre um gewesen, eine Erleichterung gefolgt sei, nämlich, daß man dem Volke die Aecker wiederum eingeräumt hat, doch mit der Bedingung, daß sie davon dem Könige den fünften Theil geben sollten. Welches man zwar in unsern Landen leichtlich sollte ertragen können, vornehmlich, wo die Aecker dem Könige zustünden. Wir thun jetzt unsere Aecker Andern um die Hälfte der Früchte aus, so darauf wachsen; doch also, daß die Aeckerleute für die Kosten stehen und alle Arbeit thun müssen. Das ist noch schwerer, nämlich, da man die Hälfte der Güter fordert, also daß der Ackermann den Acker erblich besitzen und gebrauchen möge. Etliche nehmen noch mehr als den halben und dritten Theil. Wiewohl nun Joseph die Aecker der Unterthanen mit gutem Rechte dem Könige gekauft hat, hat er sie doch den Egyptern wiederum verliehen, daß sie dieselben gebraucht, also daß sie davon dem Könige nicht mehr als den fünften Theil der Einkünfte entrichtet haben.

91. Das ist wahrlich eine große Wohlthat und sonderliche Gnade gewesen, welche auch \unsicher{Inya} erwünscht hat. Denn es waren ja die Aecker des Königs, und ist darauf keine andere Schätzung gelegt worden, ausgenommen, daß sie davon den fünften Theil geben; welches ein Geringes gewesen, vornehmlich von solchem fruchtbaren Lande. Sie haben ohne Zweifel auf die Weise ihre Aecker wiederum eingelöst und haben auch wiederum etwas an ihrer Nahrung zugenommen; sonderlich diejenigen, so etwas mäßiger gewesen und eingezogener gelebt, und ihre Güter nicht also mit Fressen und Saufen verzehrt und verschwendet haben, wie wir Deutschen zu thun pflegen, und deshalb auch mit den Schatzungen kein Maaße zu halten wissen, da man gebietet den vierten und fünften Theil zu geben.

92. Denn vor etlichen Jahren ist es gebräuchlich gewesen, den zwanzigsten Theil zu geben, da man von zwanzig Gulden oder Groschen einen gegeben hat. Die Juden haben den Gebrauch gehabt, den Zehnten zu geben. Die Römer centesima, das ist auch der zehnte Theil gewesen. Bei den Juden hat man den Priestern den Zehnten gegeben, und es ist beinahe so weit gekommen, daß man den Neunten hat geben müssen. Da man aber bei uns den zwanzigsten Theil gegeben hat, dasselbe ist noch leidlich gewesen. Wollte Gott! daß wir dasselbe jetzt auch noch erhalten könnten, daß es dabei bleiben möchte, vornehmlich dieweil die Aecker oder Güter den Leuten eigen sind. Und in Egypthen ist es auch keine unbillige Auflage oder Zins gewesen, daß sie den fünften Theil gegeben haben, vornehmlich in solcher Knechtschaft. Ja, es ist eine sonderliche Klugheit und Gnade an Joseph gewesen, der die Aecker also unter die Unterthanen ausgetheilt und ihnen dieselben wieder zugestellt hat, daß sie selbst vier Theile behalten und den fünften dem Könige geben sollten.

93. Dasselbige wird jetzt unserer Könige oder Fürsten keiner thun. Denn es kommen jetzt schon viele Schatzungen über das auf, daß man den zwanzigsten Theil geben muß, also daß ich schier nicht weiß, ob wir den halben oder den dritten Theil geben müssen, oder aber, ob man es gar hinnehmen will, was wir haben. Die Biersteuer hat lange gestanden und ist des Schindereins kein Ende.

94. Und der Wucher wächst über alle Maaße sehr; wie ich denn neulich von einem Wucherer gehört habe, der von hundert Gulden achtzig genommen hat. Das ist ja beinahe noch schlimmer, als wenn man das Ganze nimmt. Zu Leipzig ist eine gemeine Regel, daß sie von zehen Gulden vier zu fordern pflegen. Also muß Deutschland gar erschöpft werden; denn die Fürsten und Adel, und auch die Bürger werden allgemach ausgesogen durch die viele Schinderei der Wucherer. Danach ist das räuberische, diebische und vielspäßige Gesinde, das nimmt beinahe auch den dritten, wo nicht den andern Theil unserer Güter hinweg. Derhalben würden wir es viel besser haben und reicher sein, wenn in Deutschland dieser Gebrauch gehalten würde, daß man den fünften Theil gäbe, wie dazumal in Egypten geschehen ist.

95. Und wir klagen nicht allein über das Hausgesinde, welches der Leute Schaden und Verderben in der Haushaltung ist, sondern auch über viel andere unzählige und offenbare Dieberei, als da ist, die große Listigkeit und Betrug, so in Hantierung, in Kaufen und Verkaufen, und was sonst die Leute, Bürger, Bauern und Nachbarn mit einander zu thun haben, geübt wird. Dies sind greuliche Strafen, so jetzt Deutschland leiden muß, darum daß wir Gott nicht geben wollen, was wir ihm zu geben schuldig sind, nämlich, was zu Erhaltung des Predigtamts vonnöthen ist; denn dazu will jetzt niemand etwas geben.

96. Es hat aber Joseph bei den Unterthanen desselben Königreichs sich großen Dank verdient. Denn sie antworten ihm, und sagen: Du hast uns die sieben Jahre durch deine Klugheit und Güte am Leben erhalten, daß wir nicht gestorben sind. Also ist seine Gewalt und Ansehen, so er gehabt, über die Maßen sehr bestätigt worden, und haben ihm Alle mit einander einmüthig die Ehre erzeigt und ihn einen Vater des Vaterlandes genannt, darum daß er die große Menge Volks, da sie so gar elende und betrübte Leute gewesen, ernährt und am Leben erhalten hat. Denn daß wir noch leben, haben sie gesagt, mögen wir allein dir danken; du wollest nur also mit uns fortfahren und uns vor dir Gnade finden lassen, wir wollen Pharao gern dienen, wollen ihm mit höchstem Willen den fünften Theil geben. Und sie haben solches auch wohl thun können ohne alle Beschwerung von solchem fruchtbaren Lande; wie davon zuvor gesagt ist. Darum haben sie nun Joseph also geehrt und für einen Heiland erkannt, als der Egyptern erhalten habe; welches Lob und Ehre er um solcher Wohlthat willen erlangt hat durch das ganze Königreich Egypten.

\begin{quote}
\textbf{V. 26. Also machte Joseph ihnen ein Gesetz bis auf diesen Tag über der Egypter Feld, den fünften Pharao zu geben; ausgenommen der Priester Feld, das ward nicht eigen Pharao.}
\end{quote}

97. Mose erzählt hier, daß dies Gesetz gewähret habe bis auf diesen Tag, das ist, bis an die Zeit, da er diese Historie beschrieben hat. Ob es aber hernach sei geändert worden, wissen wir nicht. Es ist wahrscheinlich, daß sie die Aecker wiederum gelöst haben, und auch wiederum an Gütern zugenommen und reich geworden sind.

\subsection{II}

98. Wir haben aber gesagt, daß dazumal die Priester in sehr großen Ehren gehalten und durch des Königs Gebot von der gemeinen Knechtschaft des ganzen Königreichs frei geblieben sind; denn sie haben auch bei den Heiden ein groß Ansehen und Würde gehabt. Die Könige und Fürsten haben ihren Rath und Weisheit gehört und haben sie immer bei sich gehabt, wie dasselbe der Heiden Sittierung ausweisen. Denn sie haben die Unterthanen in der Religion unterrichtet, und sie den Gottesdienst gelehrt, der anfänglich so unrein nicht gewesen, wie er in nachfolgenden Zeiten geworden ist.

99. So möchte nun hier gefragt werden: Ob Joseph auch recht gethan habe, daß er den abgöttischen und gottlosen Priestern diese Freiheit gelassen hat? Denn es scheint, daß im selbigen Königreich noch viel abgöttische Leute gewesen sind unter den Priestern und auch dem gemeinen Volk. Es ist aber nun oftmals gesagt worden, daß David dem Joseph ein schön herrlich Lob gegeben hat, da er von ihm im 105. Psalm V. 20—22. sagt: „Da sandte der König hin, und ließ ihn los geben, der Herr über Völker hieß ihn auslassen. Er setzte ihn zum Herrn über sein Haus, zum Herrscher über alle seine Güter, daß er seine Fürsten unterwies nach seiner Weise, und seine Aeltesten Weisheit lehrete.“

\quellseite{469}
100. Auf diese Weise hat David, der vom Heiligen Geiste gewesen ist, Joseph gerühmt und macht ihn zum Lehrer, Priester und Fürsten über die andern Fürsten. Und dies ist ein sehr herrlich und wichtiges Zeugniß von Joseph, der bei dem Volke Gnade und Gunst gehabt, darum daß er solche große Menge von dem Schaden und Verderben, so dazumal ihnen bevor gestanden, errettet hat. Aber diese Ehre ist noch viel größer, daß er ein Lehrer gewesen ist, der die Fürsten, Aeltesten und Priester unterrichtet hat.

101. Er hat sie aber nicht die Kunst der Astronomie, vom Laufe des Himmels, oder Astrologie gelehrt, wie die Juden träumen; sondern er hat ihnen vorgehalten, wie man den wahren Gott recht erkennen und ihm dienen soll. Und er hatte dazu eine herrliche Veranlassung von dem schönen großen Wunderwerke, das Gott durch ihn sieben theuren Jahre über gewirkt hatte. Denn es ist zu der Lehre auch die lebendige Erfahrung gekommen, und die Leute haben vor ihren Augen gesehen, wie wunderbarlich das ganze Königreich und dazu etliche umliegende Völker waren erhalten worden. Sie haben sich dessen alle verwundert, woher Joseph solche große Weisheit bekommen hätte, daß er so gar genau die Theurung hat zuvor verkündigen, und zudem auch mit so großer Geschicklichkeit bestellen und zubereiten können, was zur Unterhaltung dieses Lebens dem Volke vonnöthen sein würde.

102. Dazu wird Joseph viele schöne Predigten gethan haben, so die Priester von ihm gehört, und die Lehre vom rechten wahren Gott angenommen und dieselbe darnach weiter über die ganze Welt ausgebreitet haben. Denn er hat die wahre Religion und die göttliche Weisheit gelehrt, und hat sie aber doch nicht gezwungen, die jüdischen Ceremonien, Beschneidung oder andere äußerliche Dinge, zu halten; sondern das hat er ihnen fleißig eingeprägt, daß sie an den Gott glauben sollten, der den zukünftigen Samen verheißen hatte, und hat vielleicht noch wohl etliche Ceremonien, die nicht gar ungöttlich gewesen, gehen lassen und damit Geduld gehabt, dieselbigen nicht verdammt, sondern angezeigt, wie man derselben recht gebrauchen und Gott damit dienen sollte.

103. Darum ist Joseph ein trefflicher großer Mann gewesen, vornehmlich in geistlichen Sachen; wiewohl er in politischen Sachen auch nicht geringer gewesen ist; wie David Ps. 105, V. 21. 22. beiderlei an ihm rühmt. Er ist auch nicht faul oder müßig gewesen, sondern wo er nur gekonnt und Gelegenheit gehabt, hat er das Volk unterrichtet und gelehrt, dazu sie ermahnt, die reine Lehre anzunehmen. Und ist daran kein Zweifel, es werden sehr viele Egypter durch ihn zur wahren Erkenntniß Gottes gebracht und selig geworden sein.

104. Er hat sie aber von dem Gott gelehrt, der bald im Anfang, im Paradeise, die Verheißung gegeben und dieselbe auch hernach oft wiederholt, und mit Abraham, Isaak und Jacob geredet hat. Denn so viel hat er allein gewußt, und hat nichts von der Bibel gehabt, denn nur die Historie von Abraham, die im 12. Capitel dieses 1. Buchs Mose angehet. Und dazu hat er auch die vorigen zehen Capitel von der Schöpfung gehabt; und was er davon gewußt und verstanden, hat er daher gehabt, daß es die Väter nach einander gelehrt hatten und also von Einem auf den Andern gekommen war. Denn es ist dazumal die Bibel gar ein klein Buch gewesen, hat aber die allergrößten und wichtigsten Dinge in sich begriffen.

105. Wir dürfen aber gleichwohl nicht denken, daß er ganz Egypten und alle und jede Fürsten und Bürger bekehrt haben wird. Denn desselben haben wir nirgend ein Exempel, daß ein Königreich oder irgend ein Land mit allen Unterthanen also gebessert und so wohl bestellt oder geordnet gewesen sei, daß darin gar keine gottlosen Leute oder Abgötterei übergeblieben sein sollte. Siehe das Königreich Juda und Israel an, auch unter David und den allerrühmsten Königen und Propheten; denn da dieselbigen regiert und gelehrt haben, sind gleichwohl noch immer viele böse und gottlose Leute mit untergelaufen, und ist doch ein heiliges und gottseliges Königreich gewesen um der Kirche und Gemeinde Gottes willen, so zu derselbigen Zeit im Volke Israel gewesen ist.

106. Deshalb, wiewohl viel und große Wunderwerke geschehen, so glauben doch die Leute nicht alle und bleiben immer noch etliche Gottlose und Feindlicher unter ihnen. Die Zuhörer des Propheten Jesaja sind mit ihm gar schmählich und übel umgegangen, wie er ihnen dasselbe aufrückt im 57. Cap. V. 4., da er sagt: „An wem wollt ihr nun eure Lust haben? über wem wollt ihr nun das Maul aufsperren, und die Zunge heraus recken?“ Und David wiederholt oftmals die erbärmliche Klage über die falschen Lehrer und unrechten Gottesdienste. Mose aber sagt auch 5 Mos. 29, 18. also: „Daß nicht ein Mann, oder ein Weib, oder ein Gesinde, oder ein Stamm unter euch sey, deß Herz heute sich von dem Herrn, unserm Gott, gewandt habe, daß es hingehe und diene den Göttern dieser Völker, und werde vielleicht eine Wurzel unter euch, die da Galle und Wermuth trage.“ Und im 32. Cap. V. 17. redet er hart mit dem Volk, und beschuldigt sie und sagt: „Sie haben den Feldteufeln geopfert, und nicht ihrem Gott; den Göttern, die sie nicht kannten“ etc. Stephanus Ap. Gesch. 7, 43. sagt auch: „Und ihr nahmet die Hütte Moloch an, und das Gestirn eures Gottes Remphan, die Bilder, die ihr gemacht hattet, sie anzubeten.“

107. Darum darf man nicht denken, daß nicht etliche böse Menschen unter den Frommen sein sollten; sintemal der Satan unter den Kindern Gottes mit ist, auch im Hause Hiobs. Und ist genug, daß ein Theil, wenn er auch gar gering ist, gebessert und zum rechten wahren Licht oder Erkenntniß Gottes gebracht wird. Denn also hat Joseph den König selbst und etliche Fürsten, und der Priester auch viele unterrichtet, und hat durch dieselbige die Erkenntniß des Gottes immer weiter ausgebreitet, der den Messias oder Christus verheißen hatte; die Andern aber haben es verachtet.

108. Darum wollen wir nun auf diese Frage also antworten, da gefragt wird: Ob Joseph daran recht gethan habe, daß er diese Wohlthat der Freiheit auch auf die undankbaren und gottlosen Priester gewendet hat? nämlich, daß er dasselbe nicht gethan habe um der abgöttischen oder gottlosen, sondern um der frommen willen, die er bekehrt hatte. Wiewohl ich das auch leichtlich glauben will, daß dazumal in Egypten weniger Greuel und scheußlicher falscher Gottesdienste gewesen, denn in nachfolgenden Zeiten aufgekommen sind. Denn was später für seltsamer wunderbarer Aberglaube in der Religion gewesen sei, das kann man aus den griechischen und lateinischen Poeten wohl schließen und verstehen; welcher Aberglaube und seltsame Abgötterei gleichwohl nicht ist angenommen oder gelobt worden, da Joseph die Leute gelehrt und unterrichtet hat.

109. Denn das pflegt also zu gehen, wie solches die Erfahrung von aller Zeit her mitbringt und lehrt: wenn ein gottloses Geschlecht und das Gott erkannt hat, aus ist, daß alsdann bald ein anderes an die Statt kommt, welches viel böser und gottloser ist. Wie im Buch der Richter Cap. 2. V. 10. gesagt wird: „Da auch alle, die zu der Zeit gelebt hatten, zu ihren Vätern versammelt worden, kam nach ihnen ein ander Geschlecht auf, das den Herrn nicht kannte, noch die Werke, die er an Israel gethan hatte.“ Wenn aber die Nachkommen von der rechten Erkenntniß Gottes abfallen, so entstehen sie greuliche, wunderliche und scheußliche Gottesdienste; und wenn diejenigen hinweg sind, die zuvor Gott treulich gedienet hatten, so behalten sie zwar wohl die äußerlichen Werke der Väter, aber sie vermehren dieselbigen immer mit neuem Gottesdienst, und vermischen also den falschen Gottesdienst mit dem rechten wahren Dienste, und fallen also in scheußliche Abgötterei; wie solches an dem Volk Israel wohl zu sehen ist, daß der Astaroth und dem Moloch geopfert hat um keiner andern Ursache willen, denn daß sie Gott etwas andächtiger und geistlicher dienen wollten, denn ihre Väter gethan hatten.

110. Also ist es in der Kirche im Neuen Testament auch zugegangen. Da sind bald nach den Aposteln die Ketzer herein gekommen; item, die Bischöfe, so den Herrn nicht recht erkannt haben, hernach die Mönche, und endlich das ganze Papstthum, dazu der ganze Greuel, so an heiliger Stätte steht. Und dieselben haben alle gerufen: Wir wollen Gott etwas eifriger und mit größerer Heiligkeit dienen, denn vor uns geschehen ist. Denn die Apostel, haben sie gesagt, haben viel gelassen, daß die Kirche auch noch zu ihrer Lehre stehen soll. Die ersten Väter haben wenig Andacht und Ceremonien gehabt, wir müssen traun noch mehr äußerliche Ceremonien und Gottesdienste aufrichten. Und auf solche Weise sind in der Kirche so viel neue Ceremonien aufgekommen, da der Teufel allerlei Verfälschung unter die Wahrheit gemischt hat, und sind die Nachkommen allgemach dahin gefallen und immer ärger geworden.

\quellseite{470}
111. Also haben die Egypter, da Joseph gestorben ist, Kälber geopfert, und andere Opfer der Väter auch beibehalten: die Nachkommen aber haben solche Opfer vermehren wollen, und haben schreckliche Ungeheuerlichkeiten dazu gethan, so hin und wieder in den Historien erzählt werden. Denn wo man Gott verliert, der uns seine Verheißung gegeben hat, und dazu den Glauben auch fahren läßt, so bleibt nichts mehr denn nur ein unseliges Nachfolgen der äußerlichen Werke. Wie die Kinder Israel, da sie die Verheißung und das Wort verloren, in welchem Gott Abraham hat heißen den Isaak opfern, sind sie endlich dahingefallen, daß sie nur zum Schein dieses Werk äußerlich nachgeahmt und Menschen oder auch ihre eigenen Kinder geopfert haben. Das hat, sagten sie, Abraham gethan, darum wollen wir dasselbe auch thun. Und in der päbstischen Kirche sehen wir Anderes nichts, denn daß sie nun auch äußerlich die Werke nachthun, so die Heiligen zuvor gethan haben: die haben, sagen sie, gefastet, haben sich auf diese oder andere Zeit enthalten; denselbigen müssen wir wahrlich darin nachfolgen, das will sich also gebühren.

112. Man soll aber vielmehr und vor allen Dingen das halten, das in der Epistel an die Hebräer steht am 13. Cap. V. 7.: „Gedenket an eure Lehrer, die euch das Wort Gottes gesagt haben, welcher Ende schauet an, und folget ihrem Glauben nach.“ Denn an den Heiligen soll man vornehmlich auf ihren Glauben Acht haben. Man soll sehen, was sie geglaubt und worauf sie gehofft, und was sie für ein Wort gehabt haben: da sollt du ihnen nachfolgen und dann hernach opfern. Aber dasselbe verachten sie, als ob sie solches gar nichts anginge, und werden rechte Affen der Heiligen, die nur dem äußerlichen Schein nachfolgen ohne Wort und Glauben. Also auch da die Egypter das Wort verloren, haben sie zwar die äußerlichen Ceremonien, wie sie dieselben von Joseph empfangen, behalten, hernach aber haben sie noch mehr dazu gethan und ihren Abgott Apis und die vierfüßigen kriechenden Thiere auch angebetet etc.

113. Ja, das sind die rechten Früchte, so aus den Ceremonien und Menschen Satzungen kommen, dazu die Leute Lust haben, und lassen das Wort und den Glauben liegen. Was noch unsere Ketzer und Schwärmer thun werden, da wird man Wunder ansehen. Denn die Wiederläufer und Sacramentsschwärmer verachten jetzt auch das Wort, lassen die Lehre des Glaubens liegen; und können doch unterdeß den größten Schein der Heiligkeit und der Ehrbarkeit vorgeben. Und ich höre, daß bei den Schwärmern eine sehr strenge Ordnung und Zucht gehalten werde, was das Aeußerliche belanget; denn sie spielen nicht, sie laufen und schlemmen nicht also, erheben sich nicht überflüssiger Kleidung, oder unmäßigen Wohlleben mit Fressen und Saufen. Das ist ihre Religion, deren sie sich höchlich wissen zu überheben und zu rühmen, als die da weit frömmer denn wir wären. Ja, wo ist aber das Wort? Der Leib Christi, sagen sie, ist im Himmel, das Brod aber ist Brod, Wein ist Wein: sie glauben nichts, sind vom Wort und Glauben abgefallen.

114. Also werden unsere Nachkommen, wenn sie die reine Lehre des Worts verloren haben, auch seltsame und scheußliche wunderliche Meinungen erdichten, wie die Egypter gethan; und also wird Gott die Undankbarkeit und Verachtung seines Worts mit erschrecklichen Plagen strafen.

115. Wir lehren jetzt also, daß Gott nicht allein der Juden, Abrahams, Isaaks und Jacobs Gott sei, wiewohl er allein Abraham und seinem Samen das Wort und die Beschneidung gegeben hat. Denn St. Paulus sagt Röm. 3, 2.: „Ihnen ist vertrauet, was Gott geredet hat.“ Und im 147. Psalm V. 19. steht auch also: „Er zeiget Jakob sein Wort“ etc. Es hat wohl Gott das Predigtamt in diesem Samen und Volk Israel verordnet: aber desselbigen haben sich nicht allein die Juden, sondern auch die Heiden zu erfreuen und mögen desselben auch mit genießen, wie wir droben von Abimelech und von den Cananitern gesagt haben. Und wir haben denselben auch noch mehr Exempel, als, an Hiob, Bileam, Jethro, dem Priester zu Midian und Moses Schwiegervater, welchen das Wort und Predigtamt nicht ist vertraut gewesen, und sind doch dessen auch theilhaftig geworden, da sie dasselbige Wort gehört haben.

116. Also sagen wir auch, daß viel Heiden selig geworden sind, auch aus der Linie oder Geschlechte Kains. Denn es sind viele aus zufälliger Gnade und Barmherzigkeit bekehrt worden, da es sich mit ihnen also zugetragen, daß sie sich auch zu der Kirche oder Gemeine gehalten, die Gottes Wort gehabt. Wiewohl ihnen die Verheißung von Christo nicht gegeben war, so ist doch die Frucht solcher Verheißung auf die Heiden gekommen, da sie die Predigten und Lehre der Väter gehört haben.

117. Eben dasselbe soll man von den Egyptern auch halten, welche die Verheißung und den Gebrauch der Beschneidung nicht gehabt. Denn da Jakob oder Joseph bei ihnen das Wort gelehrt, haben sie dieselbigen gehört, und haben sich auch mit ihnen verheirathet und Schwägerschaft mit ihnen gemacht, auf daß aus Israeliten und Egyptern Eine Kirche werden möchte. Das hat Gott gethan, ehe Christus ist geboren worden; so wird er ja viel mehr dieselbige Weise und Ordnung halten in Sammlung seiner Kirche, nachdem er in das Fleisch gekommen und Mensch geworden, und da nun das Evangelium unter alle Völker ausgebreitet und gepredigt wird.

118. Aber das glauben die Juden nicht, und unsere Schwärmer haben diese Lehre auch gemißbraucht und sind darüber in einen schändlichen Irrthum gefallen. Wie Zwingli neulich geschrieben hat, daß Numa Pompilius, Hektor, Scipio und Herkules sich auch mit Petrus und Paulus und den andern Heiligen im Paradies der ewigen Seligkeit freuen werden. Welches nichts Anderes ist, denn daß sie damit öffentlich bekennen, daß kein Glaube und kein Christenthum sei. Denn wo Scipio und Numa Pompilius, so abgöttische Leute gewesen, selig geworden sind, warum hat denn Christus leiden und sterben müssen? oder was ist es vonnöthen, daß sich die Christen müssen taufen lassen, oder daß man von Christo viel predigt und die Leute auf ihn allein weiset? So gar greulich fallen die Schwärmer dahin, wenn sie das Wort fahren lassen und verlieren, und wissen nichts vom Glauben, sondern halten und lehren eben dasselbe, so im Pabstthum auch gelehrt worden ist: wenn der Mensch thut, was an ihm ist, so wird er dadurch selig.

119. Also ist dies ein sehr schädlicher Irrthum, den wir keineswegs leiden oder vertragen können. Und ich höre doch, daß Zwingli meine Auslegung über dies erste Buch Mose anziehe und sich darauf berufe, da ich gesagt habe, daß etliche vom Geschlecht Kain selig geworden sind; wie ich denn dasselbe jetzt auch noch also lehre. Ich sage aber nicht, daß sie selig geworden sind als Kainiter oder Egypter, sondern als die da eingeleibt und mit der Kirche und Gemeine der Gottseeligen vereinigt gewesen sind.

120. Denn wir haben droben gehört, so oft Mose erzählt, daß Abraham und die andern Väter Altäre aufgerichtet haben, daß sie daselbst ihre Familie gelehrt haben, und daß auch andere Leute zu denselbigen gekommen sind, die Predigt angehört, und das Wort, so die Väter gelehrt, angenommen, und daß sie auch mit den Gottseeligen den rechten wahren Gott angerufen und ihr Gebet mit jener Gebet vereinigt haben. Ich habe traun das nicht gesagt, daß die Heiden von sich selbst, oder aus sich selbst, oder aber durch ihre Ceremonien sollten selig geworden sein, sondern durch das Wort der Väter: welches denn auch der 105. Psalm V. 22. von Joseph so hoch rühmt, da er sagt: Er hat den König, die Fürsten und Aeltesten in Egypten unterwiesen, und hat sie Weisheit gelehret etc.

121. Was ist aber das für eine Theologie, daß man keinen Unterschied zwischen dem Wort und da man kein Wort hat, zwischen Licht und Finsterniß machen will? Denn wo das Wort ist, da ist auch das Licht, wo aber kein Wort ist, da ist eitel Finsterniß. Soll man aber denn nun diejenigen, so das Licht und Wort haben, zu denen setzen, so in Finsterniß sind? Also geht es gemeinlich zu, nämlich, wenn man Eines Artikels fehlet und darinn irrt, daß alsdann des Irrens keine Zahl oder Maß ist. Denn also sind auch die Egypter in immer schändlichere Abgötterei gefallen, bis sie sich endlich auch nicht geschämt haben, Katzen und Mäuse anzubeten. Wie denn auch der heidnische Poet Juvenal solcher ihrer Unsinnigkeit gespottet und sie verlacht hat, da er sagt: Felices, quibus haec nascuntur in hortis, numina! hoc est: Ei, wie sind das so selige Leute, denen diese Götter in ihren Gärten wachsen! denn sie haben auch Knoblauch und Zwiebeln als
\quellseite{471}
Götter geehrt. Derhalben, die einmal vom rechten Weg abgewichen sind, dieselben werden allgemach vom Satan zu andern und schwereren Fällen getrieben. Und die Zwinglianer werden auch noch viel greulicher und schändlicher fallen, darum daß sie noch immer fortfahren, ihre Irrthümer zu vertheidigen.

122. Darum sage ich mit Zwingli nicht, daß die Römische Kirche, oder Numa Pompilius und dergleichen andere Heiden mehr selig und Erben des Himmelreiches geworden sind; sondern daß etliche fromme Männer und Weiber aus dem Geschlecht und Stamm Kains Gottes Wort und die Lehre der lieben Väter gehört haben und durch solchen Glauben auch zur Gemeinschaft des Himmelreichs mit der Kirche oder Gemeinde der Patriarchen gekommen sind.

123. Denn es hat sich Gott allezeit auch aus den Heiden eine Kirche zu sammeln gepflegt; wie Ruth eine Moabitin gewesen, Rahab eine Cananitin, die doch beide im Geschlechtsregister Christi aufgezählt werden. Und haben sich diese nicht allein zu den Gottseligen gehalten, sondern auch noch viele andere Canaiter mit ihnen. Nicht daß Ruth oder Rahab in ihrem Irrthum oder Abgötterei Vergebung der Sünden erlangt und derselbigen theilhaftig geworden sei; sondern da sie bekehrt worden ist, hat sie das Wort von den Israeliten angenommen. Denn das heißt recht, aus einem Heiden und Ungläubigen gläubig werden. Denn nachdem sie dem Wort, das sie gehört, geglaubt hat, ist sie nicht mehr eine Heidin, sondern ein rechtes Gliedmaß der Kirche gewesen. Also ist Naeman durch des Propheten Elisa Wort und Lehre auch bekehrt worden.

124. Ich weiß, daß vor fünfzehn Jahren Viele der Meinung gewesen sind, daß jeder Mensch in seinem Glauben würde selig werden. Was ist aber das Andere, denn daß man aus allen Feinden Christi eine Kirche machen will? Daraus wird denn bald folgen, daß uns das Wort und der Sohn Gottes, unser lieber Herr Jesus Christus, vergeblich und umsonst gesandt und gegeben sei, und wird auf die Weise nun gar kein Unterschied sein zwischen Türken, Papisten, Juden, und uns, die wir Gottes Wort haben.

125. Ich schließe zwar die Heiden nicht aus, aber ich sage, daß sie sonst auf keine andere Weise können selig werden denn allein durch das Wort Christi. Jethro, der Priester in Midian, wird entweder durch Mose oder Joseph zur Erkenntniß Gottes gekommen sein. Denn Joseph ist nicht müßig oder faul gewesen, sondern hat fleißig und treulich gelehrt, und hat viel Geschlechter der Menschen, mit welchen er gelebt und umgegangen ist, zur Erkenntniß Gottes gebracht. Das Wort kommt nicht von den Heiden, sondern, wie Christus sagt Joh. 4, 22.: „Das Heil kommt von den Juden.“ Es ist aber doch nicht bei den Juden allein geblieben, sondern ist unter die Heiden auch ausgebreitet worden, die haben es von den Juden empfangen.

126. Darum ist es eine große Kühnheit der Zwinglianer und sind unverschämte Leute, daß sie solches lehren und meines Ansehens und Exempels also zu mißbrauchen sich unterstehen. Denn ich habe nie anders gelehrt, verstanden oder gehalten, denn wie ich es nun etlichemal in dieser Historie angezeigt und wiederholt habe aus dem 105. Psalm V. 22., nämlich, daß die Fürsten in Egypten von Joseph unterwiesen und gelehret worden sind. Denn dieweil sie Heiden gewesen, sind sie mit nichten selig geworden durch ihre gottlosen abgöttischen Opfer; sondern haben die rechte Weisheit und die wahre Erkenntniß Gottes aus Josephs Predigten gelernt.

127. Und eben auf dieselbige Weise hat er auch den größten Theil des gemeinen Volks gelehrt, und sein Wort hat eben daselbst rechte Statt gehabt, nämlich, da sein Amt und Gewalt, so ihm befohlen war, durch die wunderbare Regierung die ganze Zeit der Theurung und auch der wohlfeilen Jahre bestätiget worden ist. Da ist nun Jakob mit seinen Söhnen und seiner Familie auch dazu gekommen und ist eine sehr schöne Kirche und Gemeinde gewesen. Der König ist fromm und gläubig, die Fürsten sind auch gottselig, das Volk hält von Gott, was recht von ihm zu halten und zu glauben ist; und bleiben doch ohne Beschneidung und andere jüdische Ceremonien: ja, es ist ihnen auch nachgegeben und zugelassen worden, daß sie etliche egyptische Ceremonien behalten und auch etwa unvernünftige Thiere geschlachtet und geopfert haben.

128. Wiewohl sie hernach die rechte, wahre und teure Lehre verworfen und allein die äußerlichen Werke und Opfer behalten haben, welchen auch die Nachkommen der heiligen Väter nachgeahmt und verfolgt haben, als, da sie das goldene Kalb angebetet und was sonst oftmals dessen mehr geschehen ist. Sie haben zwar die Opfer von den Vätern gehabt, aber sie haben dieselben mißbraucht und sind also hernach in greuliche Abgötterei gefallen. Denn also pflegt der Teufel aus der Kirche Gottes die Hölle zu machen. Und wir Deutschen haben uns deß auch zu befürchten, die wir das so große Licht des Evangeliums zu genießen haben, und täglich ermahnt werden, daß wir Gott dafür sollen dankbar sein und mit höchstem Fleiß an dem Wort halten. Und Gott droht auch, daß er die schwersten Strafen und Finsterniß wolle kommen lassen über die Verächter seines Worts. Aber dahin kann es leichtlich kommen, daß bald hernach Alles wiederum voll Irrthum, Phantasei, Secten und Schwärmerei der Rottengeister werden wird, welche noch jetzt bei unserm Leben öffentlich sich unterstehen, ihrem Muthwillen sehen zu lassen, und hin und wieder Verfälschung der rechten Lehre auszubreiten. Darum lasset uns Gott dankbar sein, und lernen, uns vor ihrer Listigkeit zu hüten.

129. Und bis daher hat nun gewähret die Historie von der Theurung und auch von den wohlfeilen Jahren. Nun wird Mose wieder kommen auf die Kinder Israel und auch auf Jakob, den Vater selbst. Und von diesem Ort an bis an das Ende dieses Buchs werden nun die Dinge beschrieben, so zwischen denen gehandelt worden sind, die dazumal die Kirche gewesen und Gottes Wort gehabt haben.

\section[Fünfter Theil]{Fünfter Theil. Wie Jakob sein herannahendes Ende merkt, Joseph zu sich fordern lässet, eine Bitte an ihn thut; und wie er endlich sich auf dem Bette neigt}

\subsection{I}

\begin{quote}
\textbf{V. 27. 28. Also wohnte Israel in Egypten, im Lande Gosen, und hatten es inne, und wuchsen, und mehreten sich sehr. Und Jakob lebte siebenzehn Jahr in Egyptenland, daß sein ganzes Alter ward hundert und sieben und vierzig Jahr.}
\end{quote}

130. Nun wohnen sie im Lande Raemjes. Da werden sie auch ohne Zweifel egyptische Weiber genommen haben, dieweil sie aus ihrem Geschlecht nicht andere haben können; und die andern Egypter haben auch das Wort gehört und haben diesen fremden Gästen die Werke der Liebe erzeigt. Ich glaube gänzlich, daß es eine schöne wohlgeordnete Kirche wird gewesen sein. Und wird auch allhier die Zeit angezeigt, wie lange Jakob in Egypten gelebt habe, nämlich, siebenzehn Jahr; darin hat er vielen Leuten können nützlich sein. Denn dieweil Joseph die vierzehn Jahre über in den wohlfeilen Jahren und auch in der Theurung so große reiche Frucht geschafft hat, wird viel mehr sein Vater mit den andern Söhnen und seinem ganzen Hausgesinde großen Nutzen geschafft haben, da er die Egypter immer gelehrt und gestraft hat, auf daß sie ja Gott recht möchten erkennen lernen und ihn anrufen. Sein ganzes Alter aber ward hundert und sieben und vierzig Jahr. Und was nun weiter folgen wird bis ans Ende, ist alles im letzten Jahr Jakobs geschehen.

\begin{quote}
\textbf{V. 29. 30. Da nun die Zeit herbei kam, daß Israel sterben sollte, rief er seinen Sohn Joseph, und sprach zu ihm: So hab ich Gnade vor dir funden, so lege deine Hand unter meine Hüfte, daß du die Liebe und Treue an mir thust, und begräbst mich nicht in Egypten; sondern ich will liegen bei meinen Vätern, und du sollst mich aus Egypten führen, und in ihrem Begräbniß begraben.}
\end{quote}

131. Der fromme heilige Patriarch Jakob hat erkannt, daß nun die Zeit vorhanden war, da er von diesem Leben abscheiden sollte; welches er entweder durch eine Weissagung erkannt hat, oder aber da er gefühlt, daß er Alters halben sehr schwach geworden.

132. Und dies ist eine sehr schöne Beschreibung seines Todes. Denn es hat ihn kein plötzlicher oder gewaltsamer Tod betroffen, sondern ist mit seinem, sanftem, stillem und friedsamem Herzen entschlafen; also ist er aus dieser Welt gezogen, darin er ein gar bitter Leben, das voll schweren Jammers gewesen ist, geführt.

133. Siehe aber, wie er mit seinem Sohn redet und wie er seine Majestät begehrt. Er hätte ihm aus väterlicher Gewalt wohl gebieten mögen, aber er erkennt, daß er billig soll geehrt werden darum, daß er ein Fürst und Herr war in Egypten, ohne welches Herrschaft und Gewalt er nicht hat können in das Land Canaan geführt werden. Darum ehrt er die königliche Majestät an seinem Sohn, und hält ihn vor Augen, und gebietet ihm nicht, was er thun soll, sondern bittet ihn und sagt: „Hab ich Gnade vor dir gefunden“ etc. Er unterwirft sich seinem Sohn, wie Einer, der ihm zu Füßen fallen und ihn flehentlich um etwas bitten wollte; denn er weiß, daß er von Gott erhöhet ist über seinen Vater und Brüder.

\quellseite{472}

134. Und also wird nun erfüllt der Traum, da er gesehen hat, daß sich elf Sterne, die Sonne und der Mond vor ihm neigeten. Denn seine Mutter, Rachel, war schon gestorben; der Vater aber neigt sich vor dem Sohn, und bekennt, daß er unter dem Sohne sei. Denn das heißt eigentlich sich neigen, nämlich, daß man nicht allein die Kniee beuge, sondern auch viel von Einem halte, als von seinem Obersten und gebietenden Herrn. Also war Joseph der Gewalt und Regiments halben höher denn sein Vater, wiewohl er des Gehorsaṃs und Ehren halben, so die Kinder ihren Eltern zu erzeigen schuldig sind, unter ihm war. Und wiewohl es scheint, als ob dies auch zu der väterlichen Gewalt gehöre, da er von seinem Sohn den Eid nimmt, so hat er es doch auch von ihm nur bittweise verlangt.

135. Die zwei lateinischen Worte, misericordia und veritas, heißen eigentlich auf Deutsch, Wohltat und Treue; darum, daß er zu ihm sagt: Facito bonitatem, et serva, ist so viel gesagt, als: Halt darüber, daß Treue und Glauben sei.

136. Warum er ihn aber, da er schwört, die Hand hat heißen unter seine Hüfte legen, davon haben die Juden viel unnützes Geschwätz gemacht; und wenn sie unsere Bücher nicht hätten, daraus sie heimlich sehr viel nehmen, so wären es die allerärmsten und ungeschicktesten Leute, die da leben mögen. Derhalben dichten sie, daß mit dieser Geberde die Beschneidung solle bedeutet sein, so Jakob damit habe ehren wollen. Aber die Väter haben etwas viel Größeres damit gemeint, das beide der Ehren und Zeit halben höher und mehr ist denn die Beschneidung. Denn die Verheißung vom Samen ist zuerst dem Abraham gegeben worden, 1 Mos. 15, 6. und Röm. 4, 3., da also stehet: „Abraham hat Gott geglaubet, und das ist ihm zur Gerechtigkeit gerechnet.“ Da war schon die Verheißung und der Glaube von dem Samen, so aus den Lenden Abrahams sollte geboren werden. Die Beschneidung aber ist hernach das Zeichen derselbigen Verheißung gewesen. Darum ehrt Jakob die Hüfte um der Verheißung willen, so die Hüfte gehabt, und um des Glaubens willen an solche Verheißung. Denn von den Hüften Abrahams, Isaaks und Jakobs sollte der Same kommen, welcher alle Völker segnen würde. Denn Gott hat sein Wort, Heiligkeit, Wahrheit und seine Güte an die Hüften Abrahams, Isaaks und Jakobs gebunden. Darum ist es ein recht Heiligtum gewesen; nicht darum, daß es Fleisch und Blut war, sondern um der Verheißung willen, so an diese Hüften gebunden und denselben gleichsam eingeleibet gewesen; und haben im Namen Christi und des Herrn geschworen, von dem gesagt wird 5 Mos. 6, 13.: „Du sollst den Herrn, deinen Gott, fürchten, und ihm dienen und bei seinem Namen schwören.“ Denn dieser Gott oder Gottes Sohn sollte herkommen und Mensch werden aus den Hüften Jakobs.

137. Also steht droben, 1 Mos. 32, 25., auch, daß da Jakob mit dem Engel gerungen hat, ihm das Gelenk seiner Hüfte über dem Ringen verrenkt worden ist. Denn Christus ist vom Fleisch und doch nicht vom Fleisch gekommen. Er ist hergekommen durch die Hüften der Väter von Abraham bis auf die Jungfrau Maria, da hört die fleischliche Fortpflanzung auf. Denn wiewohl Maria vom Stamm und Geblüte der Väter war, hat sie doch keinen Samen von menschlichem Ursprung, sondern ist vom Heiligen Geiste schwanger geworden. Also haben die heiligen Väter Christum erkannt und an ihn als an Gott den Herrn geglaubt; welcher wohl vom Fleisch hat sollen geboren werden, aber doch nicht durch das Fleisch, sondern durch den Heiligen Geist, also daß von ihm nicht konnte gesagt werden: „Was vom Fleisch geboren wird, das ist Fleisch“, Joh. 3, 6.

138. Auf diese Weise ehrt nun Jakob, der Vater, seinen Sohn, Joseph, und nimmt den Eid von ihm in Namen des Sohnes Gottes, unsers Herrn Jesu Christi, und Joseph schwört ihm auf seine Hüften in Glauben und Bekenntnis des zukünftigen Samens. Daraus leichtlich abzunehmen ist, daß die Väter fleißig gelehrt und dem Volke fest und beständig vorgetragen haben, daß Christus würde in das Fleisch kommen, der da wahrer Gott und Mensch wäre. Mose hat wohl dies alles kürzlich erzählt; aber die Väter werden ohne Zweifel weitläufig und ernstlich davon geredet und mit einander gehandelt haben.

139. Joseph aber ist hier eine Figur Christi, vor welchem, wiewohl er Jakobs Sohn ist, sich doch der Vater neigt: wie David Christum einen Herrn nennt, so er doch sein Sohn war, Matth. 22, 45. Auf dieselbe Weise sind auch Abraham, Isaak, Jakob und alle Patriarchen Väter Christi gewesen, dem sie gedienet, den sie geehrt und angerufen haben.

\begin{quote}
\textbf{V. 30. 31. Er (Joseph) sprach: Ich will thun, wie du gesagt hast. Er aber sprach: So schwöre mir. Und er schwur ihm.}
\end{quote}

140. Es ist wahrscheinlich, daß der Patriarch Jakob in der Väter Begräbniß habe liegen wollen um der Hoffnung willen der Auferstehung. Denn diese großen heiligen Männer und Helden der Kirche und Gemeinde Gottes haben an der Unsterblichkeit der Seele gar nicht gezweifelt, darum sind sie des Begräbnisses halben so sorgfältig gewesen. Denn wo der Heilige Geist dasselbe in ihnen nicht fest und gewiß gelehret hätte, so wäre es überflüssig gewesen, daß sie so gar ängstlich wollten dafür gesorgt haben, wie oder wo ihr Leib möchte begraben werden, als welcher am Menschen der alleringste Theil ist. Wiewohl die Heiden die Gräber der Verstorbenen auch ehrlich und herrlich zugerichtet haben. Dasselbe haben sie aber von den Vätern genommen und nur solche vergängliche Ehre der herrlichen Begräbnisse damit gesucht, wiewohl sie keine Hoffnung des ewigen Lebens gehabt. Die Väter aber haben fest und gewißlich geglaubt die zukünftige Auferstehung der Todten, also daß sie auch an der Auferstehung des Leibes nicht gezweifelt haben.

141. Darnach hat Jakob wohl gewußt, daß er in Egypten einen großen Samen haben würde, hat aber doch, wenn er sterben würde, in der Väter Begräbniß liegen wollen, da Abraham, Isaak, Lea etc. auch geruht haben. Und ist wohl glaublich, daß Jakob deren Einer gewesen sei, so mit Christo auferstanden sind. Darum hat er dafür so fleißig gesorgt, daß er ja im Lande Canaan möchte begraben werden, da Christus ist von den Todten auferstanden, und viele von den Vätern und ihrem Geschlechte mit ihm auch auferstanden sind.

\subsection{II}

\subsection{Die neigte sich Israel auf dem Bette zu den Häupten}

142. An diesem Orte sind die Dolmetscher der Sachen nicht eins deßwegen, daß die Punkte im hebräischen Texte nicht gleich sind, welche in der hebräischen Sprache vielfältige Zweifel gebären; denn die Weise, die Punkte zu setzen, ist sehr ungewiß. Und wenn wir das Neue Testament nicht hätten, so wäre uns die hebräische Sprache kein nütze und wir würden daraus nichts Gewisses schließen oder verstehen können. Welches an den Hebräern wohl zu sehen ist, so die heilige Schrift jämmerlich und greulich verfälschen und verdunkeln, darum daß sie des rechten wahren Lichts der Schrift beraubt sind, haben in solcher Dunkelheit nichts Anderes, denn sie folgen können, denn nur die Weise mit den Punkten, welche doch fast zweifelhaftig und ungewiß ist. Den Vätern und Propheten ist die heilige Sprache sehr wohl bekannt gewesen und sie waren derselbigen ganz gewiß; aber die Sprache ist zugleich mit ihrem Reiche umgekommen. Und nehmen heutiges Tages die Juden viel aus dem Neuen Testamente, ohne welches sie auch gar nichts wissen oder verstehen würden.

143. Ich kann zwar nicht leugnen, daß es nicht sollte eine sonderliche Art und Weise haben mit den Punkten in der hebräischen Sprache, aber ich habe gar keine Lust dazu, darum daß solche Weise so gar zweifelhaftig ist; welcher Zweifel die Schrift denn auch ungewiß macht und reißt sie hin und her gar von einander.

144. Münster führt an einem Orte einen jüdischen Rabbi an, der da sagt: Sine supra et infra non potest intelligi Scriptura sancta, das ist: Die heilige Schrift kann nicht verstanden werden ohne die oberen und unteren Punkte. Und dasselbe ist wahr bei den Hebräern. Sie zeigen aber nicht an, wer der sei, der das gelehrt oder geordnet habe, daß man diese Worte also nach den Punkten lesen soll; sie bringen auch keine gewisse Argumente oder Beweise, warum man eben auf diese Weise die Punkte hinzuthun müsse; ja, sie lassen es gar frei und in eines Jeden Gefallen stehen, ob man die Punkte gebrauchen wolle oder nicht. Ich glaube, daß man ein ander Wort etwa auch anders lesen und verstehen müsse, wie hier mittah und mattah. Wer will aber mir gewiß sagen, welche Lesart die beste sei? Wiewohl sie sagen, man solle die Worte lesen nach den Punkten, so über und unter einem jeglichen Buchstaben gesetzt sind.

\quellseite{473}

145. Zu der Zeit des Hieronymus hat man, wie es scheint, noch keine Punkte gebraucht, sondern die ganze Bibel ist ohne Punkte gelesen worden. Die neuen Hebräer aber nehmen ich nicht an, die sich selbst zuschreiben, daß sie der hebräischen Sprache rechten Sinn und Verstand haben, so sie doch nicht Freunde, sondern Feinde der heiligen Schrift sind. Darum spreche ich die Worte oftmals aus den Punkten stracks zuwider, es sei denn, daß die vorige Meinung mit dem Neuen Testamente übereinkomme und sich damit reime. Denn aus den Punkten kann man mehr nicht haben, denn daß man nachrathen mag, ob man mattah oder mitteh lesen soll; und also auch von vielen andern Worten mehr, so mit einerlei Buchstaben geschrieben werden. Darum frage ich nicht viel nach der jüdischen Rabbinen Oberhalb und Unterhalb. Es wäre besser, man ließe die Schrift nach dem Urnerhalb; und das Neue Testament gibt uns desselben rechten innerlichen Verstand, nicht den oberen oder unteren. Dieser Text aber bleibt im Zweifel wegen der mancherlei Punkte, darüber sie nicht können eins werden.

146. Die siebzig Dolmetscher haben es also gegeben: Et adoravit Israel fastigium virgae ejus; und behalten die Präposition al, super, fastigium, und lesen matten, das heißt, Ruthe oder Stamm; welches Wort in Mose oftmals gebraucht wird, und sind über das auch noch viel andere Worte, so eben dasselbe auch bedeuten. Darum verstehen sie es also, daß Jakob, da ihn sein Sohn geschworen, das Haupt, oder den Leib geneigt habe gegen das Scepter Josephs und habe sich also vor seiner Hoheit gebückt.

147. Augustin hat sich auch darüber den Kopf zerbrochen und läßt es sich sehr sauer werden, und bringt einen andern Verstand daher, nämlich, daß Jakob, der nun ein alter krummer Mann gewesen, sich, da er gegangen ist, an einem Stecken gehalten habe, wie die Alten auf drei Beinen zu gehen pflegen, und daß er sich also über seinen Stecken geneigt und GOtt angerufen habe. Solche Uneinigkeit, daß Einer diesen, der Andere aber einen andern Verstand über diese Worte gehabt, kommt daher, daß sie in der Grammatik mit den Punkten also zweifeln.

148. Die siebzig Dolmetscher aber haben von ihrer Auslegung ein Gleichniß aus der Historie Esther, da der König Ahasveros das goldene Scepter in seiner Hand gegen Esther reckt, und sie, die Königin, herzurütt und mit ihrer Hand die Spitze des Scepters anrührt. Esth. 5, 2.

149. Es mag vielleicht ein solcher Gebrauch gewesen sein, wie es jetzt gehalten wird, wenn man Magistrater oder Doctoren macht, daß sie haben auf das ausgereckte Scepter schwören, oder dasselbe küssen müssen, die Demuth und Gehorsam gegen ihren Oberherren damit zu beweisen.

150. Derhalben kann ich der Grammatik halben aus diesem Texte nicht kommen, da man nur nachrathen und den Umständen folgen muß; wie es scheint, daß die siebzig Dolmetscher auf die Gebärden gesehen, so Esther gehabt, welche die Spitze des Scepters angerührt und geküßt hat zum Zeichen der Ehrerbietung gegen den König. Ich behalte aber unsere gemeine Dolmetschung, wie wir es im Deutschen gegeben haben, und verstehe das Wort mittah also, daß es Bett heiße. Und bewegt mich dazu vornehmlich die Ursache, daß Mose bis daher das andere Wort, matten, nicht gebraucht hat. Denn droben, 1 Mos. 30, 37., da Jakob die Stäbe schält, item 1 Mos. 32, 10., da er mit seinem Stabe über den Jordan gegangen, steht das Wort makkel. Darum will ich den Grammatikern ihre vielerlei Bedeutungen und Uneinigkeit überlassen.

151. Ich habe droben gesagt, daß Jakob, der Vater, sich vor seinem Sohne geneiget; denn er redet ihn sehr demüthiglich an, da er sagt: „Hab ich Gnade vor dir funden.“ Welches Worte sind eines Geringeren und der die Majestät dessen, den er größer und würdiger achtet, ehrt. Und wir wollen mit diesem Exempel vergleichen, was von David gesagt wird 1 Kön. 1, 47. 48., da also steht: „Und der König hat angebetet auf dem Lager.“ Auch hat der König also gesagt: Gelobet sei der HErr, der GOtt Israel, der heute hat lassen einen sitzen auf meinem Stuhl, daß meine Augen gesehen haben.“ Daraus kann man sehen, was da gesagt sei, adorare super lectum, auf dem Lager anbeten; denn es ist ein Gebet oder Danksagung zu GOtt für die empfangene Wohlthat. Und denselbigen Dienst oder Anbetung hat David GOtt erzeigt auf seinem Lager, da er jetzt weder hat stehen noch gehen können.

152. Eben auf dieselbige Weise ist es hier auch zugegangen mit Jakob: da er Alters und sonst viel Jammer halben, so er in diesem Leben erlitten, matt gewesen, hat er nicht stehen oder auf einem Stuhl sitzen können, sondern hat im Bette gelegen, als der jetzt hat sterben sollen; wie er denn jetzt nichts Anderes thut oder redet, denn daß er sich nur auf die Reise aus diesem Leben schicke und bereite. Denn es hat ihm nun sein Sohn geschworen, da er die Finger auf seine Hüfte gelegt, das ist, im Glauben der Verheißung Christi, der die Todten auferwecken wird, welcher GOtt und Mensch aus den Hüften Jakobs hergekommen ist; welche Hüften gesegnet sind um der Verheißung willen, so daran gehängt und mit dem verheißenen Samen verbunden ist. Darum beweist er seinen Glauben, Gottseligkeit und Ehrerbietung, mit welchen äußerlichen Geberden er nur kann und mag. Und dieweil er es mit Kniebeugen nicht thun kann, dieweil ihm seine Glieder fast schwach und matt gewesen, hat er doch, wie er am besten gekonnt, das Haupt bewegt und sich gegen das Bett geneigt.

153. Lyra und die Juden sagen: Er habe sich nach dem Lande Canaan gewendet und also angebetet. Darüber will ich nicht sehr streiten. Aber was sie von der Beschreibung plaudern, davon habe ich zuvor gesagt, daß solches ein unnützes und vergebliches Geschwätz sei; denn Jakob und Joseph haben beide auf die Verheißung des Segens in dem zukünftigen Samen gesehen, welcher Segen in den Hüften oder Lenden der Väter war. Darum hat Jakob diese Verheißung ehren wollen, da jetzt sein Glaube ist angefochten worden.

154. Wie wir thun, da wir nicht allein, wenn wir beten, sondern auch, wenn wir taufen, absolvieren und absolvirt werden, auch wenn wir zum heiligen Abendmahl gehen, dazu auch, wenn wir die göttliche Verheißung oder den Text des Evangeliums lesen hören, unsere Kniee beugen, oder je zum wenigsten aufgerichtet stehen sollen zum Zeichen der Anbetung, Ehrerbietung und Dankbarkeit.

155. Derhalben, wenn schon im Abendmahl des HErrn nichts Anderes gereicht würde, denn nur Brod und Wein, wie die Sacramentsschwärmer lästern: so ist doch je die Verheißung und GOttes Wort daselbst, und der Heilige Geist wirkt auch und ist kräftig durch das Wort im Abendmahl. Darum will es sich ja gebühren, daß wir mit aller Ehrerbietung hinzu gehen. Wie viel mehr aber gebührt sich solches zu thun, wenn wir glauben, daß daselbst auch mit dem Worte der wahre Leib und das wahre Blut Christi ist! Denn also haben die Väter GOtt und seine göttliche Verheißung angebetet und dazu die Zeichen der Verheißung auch; denn sie wußten, daß es GOttes Wort und Verheißung war. Sie haben das Heiligtum der verstorbenen Heiligen nicht geehrt, wie im Papstthume geschieht, sie haben die Beine oder Kleider der Todten nicht angebetet; sondern haben die Verheißung und Werke GOttes heilig gehalten, und die wohl werth wären, daß man sie billig ehren und anbeten sollte.

156. Wie hier von einer göttlichen Sache ist gehandelt worden nämlich, von dem Eide unter den Hüften Jakobs, im Glauben Christi, der da sollte Mensch werden und die Auferstehung der Todten und das ewige Leben sein sollte. Da nun desselben ist gedacht worden, bewegt Jakob seinen Leib mit einer solchen Geberde, so man pflegt zu haben, wenn man anbetet. Wie David, da er sieht, daß sein Sohn Salomo nach ihm sollte König werden, neigt er sich auch mit dem Haupte und der Brust, ruft GOtt an und dankt ihm dafür.

157. Also ehrt hier Jakob die Verheißung, und neigt sein Haupt, nicht hinter sich, sondern vor sich auf das Angesicht; zeigt mit eben derselbigen Geberde an und bezeugt damit, daß er an Christum glaube und auf denselben allein sehe, der da kommen sollte und gekreuziget werden, und auch wiederum auferstehen sollte, auf daß er allen Todten das Leben gäbe. In solcher Hoffnung und Glauben faltet er seine Hände, neigt das Haupt, nicht mitten oder am Ende des Bettes, sondern zu den Häupten; daß es gar nichts Anderes sei, denn nur eine solche Geberde der Ehrerbietung und Anbetens in solcher großen wichtigen Sache, nämlich, da er an die göttliche Verheißung denkt, die er dazumal im Worte gehabt und die hernach auch mit der That hat sollen erfüllt werden.

\quellseite{474}

158. Also müssen wir auch thun. Wenn wir Gottes Wort hören, sollen wir dasselbe, wo nicht mit Kniebeugen, doch mit demüthigem Herzen, mit Ehrerbietung, dazu auch mit sonderlicher Gottseligkeit aufnehmen. Also ist es auch gut, daß man das Sacrament des Altars mit Kniebeugen ehre. Denn da ist der wahre Leib und Blut des Herrn; desgleichen ist auch daselbst gegenwärtig der Heilige Geist und die Verheißung über das göttliche Wort, welches man ja mit gebührlicher Ehrerbietung hören soll. Denn Gott wirkt daselbst und der Herr läßt sich da sehen; welches bei Mose etlichemal das Angesicht des Herrn genannt wird, nämlich darum, daß Gott dabei ist und sich mir offenbart. Da gebührt mir wahrlich, vor ihm aufzustehen, oder aber auf die Kniee nieder zu fallen. Und sind eigentlich die Offenbarungen und Erscheinungen Gottes, so wir haben, allen Offenbarungen im Gesetze Moses gleich, ja, sie sind noch größer.

159. Wenn ich zur Taufe gehe, so soll ich es wahrlich dafür halten, daß daselbst nichts Menschliches geschehe; sondern das Wasser in der Taufe ist wie eine Decke oder ein Mittelinstrument, wie auch das Wort, damit Gott bedeckt wird. Dahinter steht unser HErr Gott, und das Wasser und Wort ist Gottes Angesicht, dadurch er mit uns redet und wirkt, mit einem jeden Menschen für seine Person insonderheit. Er tauft mich, absolvirt mich, reicht mir seinen Leib und Blut durch des Dieners Zunge und Hand. Denn Gott wirkt in der Taufe das Heil und Seligkeit der Menschen. Und dies ist die Gegenwart oder Gestalt und Erscheinung Gottes in diesen Mitteln oder Instrumenten. Darum ist es billig und recht, wo wir je den Leib nicht neigen, daß wir uns doch von ganzem Herzen neigen, und Gott ehren, der daselbst mit uns redet.

160. Und zwar in der Taufe soll man allezeit die Augen und das Herz richten auf die schöne helle Offenbarung im Jordan, da die Stimme Gottes des Vaters vom Himmel herab gehört, das Fleisch des Sohnes Gottes und der Heilige Geist in der Gestalt der Taube gesehen wird. Also auch in unserer Taufe. Und im Abendmahl des HErrn ist das Erstere das äußerliche Mittel und Zeichen, damit mir Christus seinen Leib gibt; und in Wahrheit weidet und speist uns Gott also selbst und nicht der Diener. In der Taufe hört man die Stimme der heiligen Dreifaltigkeit Gottes und sollen auch die Worte der Taufe nicht anders gehört oder verstanden werden. Daherben ist diese Ehrerbietung vonnöthen, und wo sie nicht äußerlich mit Kniebeugen geschieht, daß man doch geistliche Ehrerbietung zeige.

161. Zu unsern Zeiten aber, bei so großer Verachtung des Worts und der heiligen Sacramente, sind der Leute wenige, so darauf Acht geben, ob sie sitzen oder aufstehen, wo man das göttliche Wort liest oder predigt. Und im Papstthum, wie auch unter dem Volke Israel, hat man die rechte Erkenntniß und Verstand des Worts und der wahren Gottseligkeit liegen lassen, und haben sich die Leute nur der äußerlichen Ceremonien beflissen, beide da man den Text des Evangeliums gelesen und auch wo man die heiligen Sacramente gereicht und gebraucht hat.

162. Das im lateinischen Texte steht, caput lectuli, kann wohl geistlich oder figürlich Christum bedeuten; aber ich verstehe es einfältig, daß es heiße, der vorderste Theil des Bettes, das ist, oben zu den Häupten.

163. Was wollen wir aber sagen zu dem Texte in der Epistel an die Hebräer am 11. Cap. V. 21., da also steht: „Durch den Glauben segnete Jakob, da er starb, beide Söhne Josephs, und neigte sich gegen seines Scepters Spitze“? Da wird ja nicht gesagt, daß er unmittelbar Gott habe angebetet, sondern es wird daselbst die Spitze des Scepters gesetzt für dasjenige, davor er sich geneiget hat. Denn es ist wohl glaublich, daß Joseph ein Scepter getragen habe, sowie den Ring, den ihm Pharao oben gegeben; wiewohl dasselbe bei Mose nicht so klärlich ausgedrückt wird. Das Scepter aber ist das rechte Zeichen nicht allein der obersten Gewalt, sondern auch der andern geringen Aemter, als, der Richter, so einen Stab zu tragen pflegen. Also hat der Papst in seiner Kirche eingesetzt, daß die Beichtväter auch Stecken in den Händen haben tragen müssen; welchen Gebrauch er aus dem Spruche Christi genommen hat, da er seinen Jüngern geboten, sie sollten nichts bei sich tragen auf dem Wege denn allein einen Stab. Und daher ist auch der Bischofsstab gekommen. Also hat Joseph ein Scepter getragen, es sei gleich hölzern oder golden gewesen, und vor desselbigen Spitze hat sich Jakob geneiget. Was wollen wir aber an diesem Orte zu dem Texte sagen?

164. Augustin treibt sein Spiel mit Allegorien oder heimlichen Deutungen, die doch nichts zur Sache thun. Die Andern sind auch der Sache mit einander nicht eins. Einer hat es gedeutet, daß es soll „Bett“ heißen, der Andere „Stab“, nachdem die Punkte auch nicht einerlei sind. Wir aber folgen dem Verstande nach, welcher der einfältigste ist und der sich mit der Schrift und Exempel Davids am allerbesten reimt. Die Andern folgen dem Exempel des Königs Ahasveros im Buche Esther. Und ich gebe mich dazu für keinen Richter aus, welcher Verstand der beste sei. Das wollte ich aber ganz gern wünschen, daß man nur Eine Dolmetschung haben möchte; denn ich habe zu solchem Texte, der so mancherlei und ungleich ist, gar keine Lust; ja, ich bin aller Aequivocation, das ist, da die Wörter mehr denn einerlei Verstand haben, sehr feind. Und man soll sich befleißigen, so viel als immer möglich ist, daß man aus den Worten des Textes Einen rechten einfältigen Verstand nehmen und fassen möge, der sich fein mit der Grammatik reime; und wenn man denselben haben, so schadet es darnach nichts, ob man auch Figuren und heimliche Deutungen im Texte suchen wollte.

165. Wie St. Paulus 1 Cor. 15, 55., da er sagt: „Tod, wo ist dein Stachel?“ etc., nach seinem Geiste eine Figur macht aus dem Texte des Propheten Hosea. Und Bernhardus ist ein wunderbarer Meister damit, nämlich die Sprüche aus der Schrift also auf eine sonderliche Weise zu brauchen. Denn es kommt oft vor, daß er einen Text in der Schrift, der sonst auf eine gewisse Lehre sollte gezogen werden, etwa auf einen gemeinen Verstand zieht. Auf solche Weise mag man zwar diesen Text auch, durch die Figur, catachresis genannt, anderswohin ziehen, welches auch ein guter Verstand ist; doch daß die Grammatik in ihrem einfältigen Verstande unverletzt bleibe. Und ich halte dafür, daß der Schreiber, so die Epistel an die Hebräer geschrieben hat, dies Wort auch also durch dieselbige Figur catachresis, daß es mehr denn Eine Bedeutung habe, verstanden habe, und daß er gelegen hat, virga (welches wir verdeutscht haben: die Spitze seines Scepters), da wir lesen, lectulum, das ist, auf dem Bette zu den Häupten. Denn ich weiß sonst nicht, wie ich mir könnte daraus helfen und auf die vorgelegte Frage antworten.

166. Ich wollte aber von Herzen wünschen, will auch ermähnt und gebeten haben Alle, so die hebräische Sprache lernen wollen, daß sie ja einen jeden Text auf einerlei einfältigen und gewissen Verstand ziehen. Hernach wollen wir dann auch einem jeden leichtlich zulassen, daß er (wie die Poeten pflegen) ein wenig seinen Willen haben möge und etwa eine freiere Auslegung suche; doch so, daß nur ein einiger, rechter und eigentlicher Verstand bleibe. Nun hat es aber wohl können geschehen, daß auch an diesem Orte die Uneinigkeit des Verstandes dieser Worte aus solchem Mißbrauche gekommen ist, daß sie also gleichsam ihr Spiel mit einander der unteren und oberen Spitze haben getrieben. Ich aber, wie droben ist gesagt worden, habe mehr Lust zu dem innerlichen Verstande der Schrift.

167. Da nun Jakob seines Alters und viel Jammeres halben, so er erlitten hatte, fast matt geworden, hat er nun auch das letzte Ziel und das Ende seines Lebens erreicht, davon der Poet sagt:

\begin{quote}
Optima quaeque dies miseris mortalibus aevi Prima fugit, subeunt morbi tristisque senectus. Et labor, et durae rapit inclementia mortis.
\end{quote}

Das ist: Die beste Zeit dieses Lebens läuft den armen Menschen auf Erden bald mit dem allerersten hinweg; darnach kommen mancherlei Krankheiten und das traurige Alter, dazu viel Mühe und Arbeit; endlich kommt dann der grausame bittere Tod und reißt sie gar hinweg.

168. Denn bis daher ist Jakob sein ganzes Leben durch und durch vieler unzähligen Sorge und Unfalls halben fast umgetrieben und entkräftet worden; welche Sorgen machen, daß der Mensch soviel eher alt werden muß und dahinstirbt. Dieser Ursache halben, wie er sich droben, V. 9., beklagt, hat die Zeit seiner Wallfahrt nicht gelangen können an die Zeit seiner Väter in ihren Wallfahrt. Darum, da er nun jetzt gar zu Bette gelegen, hat er die Ehrerbietung gegen seinen Sohn auf seine andere Weise erzeigen können, denn mit dieser Gebärde, daß er sich auf dem Bette zu den Häupten geneiget hat. Denn die Alten, so nun Alters halben kaum einhergehen können, müssen endlich gar auf dem Bette liegen bleiben, oder aber wenn sie noch gehen können, müssen sie doch einen Stecken zur Hand nehmen und also auf drei Beinen gehen (wie die kleinen Kindlein erstlich auf vier Füßen zu gehen pflegen), da denn allmählich alle Sinne und Kräfte des Leibes und Gemüths abnehmen und zum Grabe eilen. Wie die schöne treffliche Beschreibung des Alters, so im Prediger Salomo Cap. 12, 3. steht, dasselbe bezeuget und ausweist: „Zur Zeit, wenn die Hüter im Hause zittern, und sich krümmen die Starken, und müßig stehen die Müller, daß ihr so wenig worden ist, und finster werden die Gesichter durch die Fenster“ etc.

\quellseite{475}

169. Joseph aber steht bei seinem Vater, und schwört ihm einen Eid, daß er thun wolle, wie der Vater gesagt hatte. Darnach ist er wieder von ihm gegangen, da er noch nicht tödtlich krank gewesen, sondern da er allgemach abgenommen, und da es sich mit ihm zum Ende genaht hat, daß er aus diesem Leben abscheiden sollte. Bald hernach wird er wieder zu ihm kommen, wie im nachfolgenden Capitel wird erzählt werden.

\chapter{Das achtundvierzigste Capitel}

\section[Erster Theil]{Erster Theil. Wie Joseph zu seinem kranken Vater Jakob reist; und wie dieser den Ephraim und Manasse an Kindes Statt annimmt}

\subsection{I}

\begin{quote}
\textbf{V. 1. 2. Darnach ward Joseph gesagt: Siehe, dein Vater ist krank. Und er nahm mit sich seine beiden Söhne, Manasse und Ephraim. Da ward es Jakob angesagt: Siehe, dein Sohn Joseph kommt zu dir. Und Israel machte sich stark, und setzte sich im Bette.}
\end{quote}

1. Mose fährt nun fort, daß er auf die Beschreibung komme, wie Jakob seine Söhne gesegnet habe, und sagt: Nach etlichen Tagen, im selbigen Jahre, habe Einer dem Joseph verkündigt (denn also lautet es im Hebräischen), daß die Krankheit und Schwachheit des Leibes bei Jakob immer zunehme und größer werde; und haben ihm ohne Zweifel die Aerzte auch angezeigt, daß eine tödtliche Verschlimmerung eingetreten sei und daß es jetzt fast am Ende mit ihm sei; auf welches Ende auch Jakob schon vorlängst mit starkem beständigem Muthe gewartet und gehofft hatte. Derhalben befürchtet nun Joseph, es möchte vielleicht der Vater also dahinstürben, ehe er ihn und seine Söhne gesegnet; und hatte auch schon zuvor dieser Sache wegen mit ihm geredet, vornehmlich und am allermeisten um des Stammes Levi willen, daß der von den andern abgetheilt und abgesondert werden möchte. Darum bringt er nun seine beiden Söhne mit sich, daß der Tod ja den Vater nicht ereile, ehe er es nach seiner väterlichen Gewalt mit einem jeglichen Stamme verordne, wie es damit nach seinem Willen, oder aber nach dem Rathschlage, so sie zuvor darüber mit einander gehalten, gehen sollen.

2. Denn diese Ordnung ist ganz hoch vonnöthen gewesen, damit unter den Nachkommen keinerlei Zank oder Uneinigkeit erwachsen möchte. Und ist solches aus sonderlichem Rath des Heiligen Geistes geschehen, daß der Stamm Levi abgesondert und zum Priesterthum verordnet würde, und warum auch keine eigenen Güter oder Aecker haben sollte, wie die andern Stämme gehabt. Auf daß aber die Zahl der zwölf Söhne Israel wieder möchte erfüllt werden, hat Jakob den einigen Joseph in zwei Stämme getheilt, da er seine zwei Söhne, Manasse und Ephraim, an Kindes Statt angenommen, welche an Josephs und Levis Statt gesetzt worden. Und die heilige Schrift braucht darnach ohne Unterschied die zwei Namen, Ephraim und Joseph; wie das im Propheten Hosea zu sehen ist, und auch im 77. Psalm V. 16., da also steht: „Du hast dein Volk erlöset gewaltiglich, die Kinder Jakobs und Josephs.“

3. Hieronymus hat an diesem Orte die siebenzig Dolmetscher billig gestraft, darum daß sie einerlei Wort anders und aber anders verdolmetscht haben. Droben verstehen sie das Wort mattah, daß es „Stab“ heiße; hier wollen sie nun haben, daß mattah „Bett“ heiße, und werden es hernach auch auf dieselbige Weise deuten. Sie haben zwar keine Ursache, daß sie so mancherlei Deutung machen und damit so unbeständig sind, und gefällt mir solche Unbeständigkeit auch gar nicht; denn es scheint, als ob sie sich dessen sonderlich befliessen haben, daß sie so mancherlei Deutung einführen wollen.

4. Da aber nun Jakob so viel Kraft und Stärke nicht gehabt, daß er hätte können aufstehen oder auf dem Stuhle sitzen, macht er sich doch stark, so gut er kann, und sitzt im Bette, weiß sich zu erinnern, daß er sich vorgenommen und bei sich beschlossen, die beiden Söhne seines Sohnes an Kindes Statt anzunehmen.

\subsection{II}

5. Der Segen aber, welchen Joseph mit seinen Söhnen empfangen, ist nicht der sonderliche und vornehmlichste Segen, welcher im folgenden Capitel V. 22. ff. wird beschrieben werden; sondern es ist nur ein allgemeiner Segen, darum daß Jakob die zwei Söhne Josephs zu Kindern angenommen. Und wiewohl zuvor solcher Exempel vielleicht in den Geschlechtern der Patriarchen mögen mehr gewesen sein, ist doch dies die erste Stelle in der heiligen Schrift, da dessen gedacht wird. Und ich zwar disputire und halte fest darüber, daß Sarah, Abrahams Weib, auch des Tharah Tochter gewesen sei, also daß er sie als Tochter angenommen und auferzogen habe. Denn darum nennt sie Abraham seines Vaters Tochter und nicht der Mutter, und nennt sie auch seine Schwester. Derhalben ist dies das erste Exempel von der Adoption oder Annahme an Kindes Statt in der heiligen Schrift.

6. Hierauf gründet sich aber der große Ruhm und die Hoffahrt des Stammes Ephraim, welcher um dieses Textes und des Segens willen, so bald hernach folgen wird, sehr stolz gewesen ist, sich darüber sehr erhoben hat, und sie sind auch um solcher Hoffahrt willen oftmals hart gestraft worden. Nun sind sie gleichwohl nicht ohne Ursache so stolz gewesen; denn es ist hier eine große Herrlichkeit, daß sie also sonderlich von ihrem Großvater erzählet, und als Jakobs Enkel mit ihren Vettern zusammen gefügt und denselben gleich gemacht werden, das Erbe einzunehmen und das Land zu besitzen. Daher ist der Stolz und die große Hoffart Ephraims gekommen; wie sie sich unter Jephtah und Jerobeam königliche Gewalt angemaßt und hart darnach gethan haben; und das David rühmt sie auch hoch und zeigt an, daß dieser Stamm unter den andern allen der mächtigste gewesen sei, da er sagt: „Ephraim ist die Macht meines Haupts“, Ps. 60, 9. Und 1 Chron. 8, 21. wird erzählt, daß die Kinder Ephraim aus eigenem Frevel und Vermessenheit sich das Recht genommen haben, die Verheißung zu erfüllen; denn sie wollten die Vorgänger derer sein, die aus Egypten zogen, und das verheißene Land vor der rechten Zeit einnehmen. Darum sind sie auch von Gott hart gestraft und beinahe gar vertilgt worden von den Männern zu Gath oder Einheimischen im selbigen Lande.

7. Derhalben sehen wir, wie die Nachkommen der Patriarchen allemal unartig und ihrer Vorfahren gar ungleich geworden sind. Denn es ist solche Bescheidenheit und Demuth an ihnen nicht gewesen, wie an Abraham, Isaac und Jakob war, die ihr Gemüth haben zwingen und inne halten können, wenn es ihnen nach ihrem Willen wohl gegangen ist. Aber diese sind aufgeblasen und stolz geworden, wenn sie große Ehre und Glück gehabt.

8. Denn es ist viel schwerer, im Glück nicht stolz werden, denn im Unglück ein friedsam Herz haben, geduldig und beständig dabei bleiben; und das menschliche Herz kann viel leichter allerlei Unglück ertragen, denn groß Glück und wenn es vollauf hat; wie man im deutschen Sprüchwort sagt: Es müssen starke Beine sein, die gute Tage ertragen können. Also sollen wir denken, daß auch die Ephraimiter dieser Adoption halben, daß sie Jakob zu Kindern angenommen und gesegnet hatte, sind stolz geworden; wiewohl sie Gott nicht darum mit so viel Gütern überschüttet hatte, daß sie davon sollten stolz werden.

\begin{quote}
\textbf{V. 3–6. Und (Jakob) sprach zu Joseph: Der allmächtige Gott erschien mir zu Luz, im Lande Canaan, und segnete mich, und sprach zu mir: Siehe, ich will dich wachsen lassen, und mehren, und will dich zum Haufen Volks machen; und will dies Land zu eigen geben deinem Samen nach dir ewiglich. So sollen nun deine zween Söhne, Ephraim und Manasse, die dir geboren sind in Egyptenland, ehe ich herein kommen bin zu dir, mein sein, gleichwie Ruben und Simeon. Welche du aber nach ihnen zeugest, sollen dein sein, und genannt werden, wie ihre Brüder in ihrem Erbtheil.}
\end{quote}

\quellseite{476}

9. Also fängt Jakob seine Rede an, die er sehr kurz abgebrochen, und erzählet, wie ihm der HErr im Lande Canaan erschienen sei, und die Verheißung gegeben habe, daß er seinen Samen mehren wolle; welches man wiederholen muß aus dem, was droben ist gesagt worden. Denn er hat zweierlei Erscheinungen gehabt, zu Luz und zu Bethel. Erstlich, da er den HErrn auf einer Leiter sah, daran die Engel Gottes auf und nieder stiegen, droben im 28. Cap. V. 13. Zum andernmal, da er wieder aus Mesopotamien kam, Cap. 35, 1., da Simeon und Levi mit der Schärfe des Schwerdts die Sichemiter erwürget hatten und der Schande willen, so sie mit ihrer Schwester Dina begangen, Cap. 34, V. 26. ff. Denn da sie denselben Lärm angerichtet, ist Jakob darüber also bestürzt worden, daß er sich befürchtet, wenn vielleicht die Nachbarn solche unbillige Tyrannei rächen wollten, so würde er mit seinem ganzen Hause umkommen und erschlagen werden, und hat deshalben daselbst mit seinen Söhnen hart geredet und sie beschuldigt, da er sagt Cap. 34, 30.: „Ihr habt mir Unglück zugerichtet, daß ich stinke vor den Einwohnern dieses Landes, den Cananitern und Pheresitern.“ Da aber nun Jakob also bestürzt und in solcher Furcht gewesen, redet ihn GOtt an, und heißt ihn hinauf gen Bethel ziehen und daselbst wohnen, und auch einen Altar daselbst machen, Cap. 35, 1. Da er aber das gethan hat, ist er ihm noch zum drittenmal erschienen und hat mit ihm geredet, da er ihm seinen Namen änderte und sprach, Cap. 35, 10.: „Du sollst nicht mehr Jakob heißen, sondern Israel sollst du heißen.“

10. Auf dieselbige Erscheinung hat nun hier Jakob eigentlich und vornehmlich gesehen. Darum soll man diesen Text fleißig merken, dieweil er die Verheißung also wiederholt. Denn daselbst sagt GOtt zu ihm Cap. 35, 11.: „Sey fruchtbar und mehre dich“; hier aber sagt er: „Ich will dich wachsen lassen und mehren.“ Und vielleicht ist das die Ursache, darum er so kurz abzubrechen scheint, wie ein alter kindischer Mann zu reden pflegt. Denn er fängt so hoch an, und hört doch so plötzlich wiederum auf zu reden.

11. Es ist aber nun oftmals gesagt, warum GOtt auf Hebräisch genannt wird el schaddai, und steht derselbe Name auch im 91. Psalm V. 1., da es unser Dolmetscher im Lateinischen also gegeben hat: Sub protectione Dei coeli commorabitur, „Wer unter dem Schatten des Allmächtigen bleibt.“ Und der HErr hat droben auch etlichemal den Namen gebraucht, da er mit Abraham und Jakob geredet; und kommt her von dem Wort schad, das heißt, Brust, es heißt auch sonst wohl, Wüste oder Geist; wie jetzt die Juden den Teufel auf Hebräisch sched nennen. Es ist ein solch Wort, das mancherlei Bedeutung hat, aber an diesem Ort heißt es nur einerlei; als wenn wir sagen wollten: GOtt ist ein solcher GOtt, der da nährt, Deus alumnus, wie die Griechen die Göttin Diana auf ihre Sprache polymastos, auf Lateinisch mammosa genannt haben, die da viele Brüste habe, darum daß sie eine Säugamme oder Mutter Asiens und des ganzen Erdkreises wäre, die allen Thieren ihr Futter gäbe und sie alle nährte. Und der Name gebührt GOtt allein; denn er ist auch allein, der da alles nährt und erhält, was er geschaffen hat. Wir haben es verdeutscht: „Der allmächtige GOtt“ etc.

12. Wiewohl es aber nun scheint, daß Jakob seine Rede sehr kurz abgebrochen, so wird doch der, so die Worte etwas fleißiger ansieht, finden, daß es nicht Worte eines alten kindischen Mannes sind, oder der Alters halben seine Rede nicht zu Ende führt; sondern daß es eigentlich Worte sind eines Menschen, der voll Heiligen Geistes gewesen. Denn aus dieser Wiederholung, da er die Verheißung von Vermehrung seines Samens wiederholt, geht klarlich hervor, daß er aus dem Heiligen Geist verstanden habe, daß es dazu kommen würde, daß Ephraim, Josephs Sohn, in die Zahl der zwölf Stämme der Brüder und seiner Kinder sollte angenommen werden.

13. Und das kann man oftmals in der heiligen Schrift sehen, da die Propheten durch sonderliche Erleuchtung und Gabe des Heiligen Geistes aus Vergleichung der Dinge und Worte viel haben heraus lesen können. Als, da im 51. Psalm V. 9. David sagt: „Entsündige mich, HErr, mit Ysopen“ etc., hat er dieselbe Art und Weise zu reden genommen von der Besprengung des Wassers und Asche, davon 4 Mos. 19, 4. 6. steht; denn daraus schließet er: Es wird eine viel andere Besprengung sein, denn die gewesen ist, davon das Gesetz Mosis sagt. Also haben sie auch bei dem Opfer des Osterlamms abnehmen können und verstehen das rechte wahre Lamm, welches die Sünde der Welt trägt.

14. Und sie haben eine sehr schöne Regel solcher Figuren und Vorbilder in Mosi selbst gehabt, da er sagt 2 Mos. 25, 40.: „Siehe zu, daß du es machest nach ihrem Bilde, das du auf dem Berge gesehen hast“, nämlich, die Hütte und das Andere alles. Daraus haben die Propheten leichtlich verstehen können, daß dasselbe ganze Gebäude ein ander Muster und Bild bedeuten müßte, nämlich, den HErrn Christum. Das äußerliche Bild nach dem Gesetz ist, daß sie Weihrauch, Blut, Opfer und Hütte gehabt, damit sie umgegangen sind. Aber GOtt hat Mose eine andere Bedeutung in das Herz und auch in den Mund gegeben; wiewohl er nichts Anderes geben kann denn nur der Schatten, ein äußerlich Gemälde und Bild der Dinge, die er gesehen hat; wie er von sich selbst sagt, 5 Mos. 18, 15.: „Einen Propheten, wie mich, wird der HErr, dein GOtt, dir erwecken“ etc. Als wollte er sagen: „Ich bin der rechte Prophet noch nicht; ich bin nur ein Schatten und äußerlich Bild: es wird aber ein Anderer kommen, welchen GOtt aus dir erwecken wird“ etc.

15. Auf diese Weise haben die Väter und Propheten durch Anregung und Erleuchtung des Heiligen Geistes aus der Vergleichung der Sachen und Worte gar herrlich und wohl verstanden, was im Gesetz mit äußerlichen Ceremonien geschehen ist. Und eben in solchem Verstande wiederholt Jakob hier auch die Worte der vorigen Erscheinung, und schließet daraus, daß dieser Sohn Ephraim, so von Joseph geboren war, ihm eigentlich zustehe. Also wollte er sagen: Nun erkenne ich erst, warum GOtt zu mir gesagt zu Bethel: „Sei fruchtbar und mehre dich.“ Denn da ihm die Verheißung zu Zug gegeben war, als Jakob zu seinem Vater Isaac zog, ist ihm unterwegs sein Weib Rahel gestorben, von welcher er gehofft Kinder und Erben zu erlangen; daß er nun also nicht mehr hat Kinder zeugen können deswegen, daß ihm sein Weib gestorben und er auch sehr alt war.

16. Verhalben hat er nun also gedacht: Ich sehe, daß wunderbarbarer Weise diese noch soll erfüllt werden, was mir zu Bethel ist verheißen worden. Denn darum hat GOtt Joseph sonderlich erhalten, daß wir noch sollten von ihm Erben geboren werden. Und wie er unter Jakobs Söhnen der größte gewesen: also ist der Stamm Ephraim auch der größte und königliche Stamm gewesen, da er mit Manasse verbunden worden. Das hat Jakob im Geist gesehen, und hat daraus geschlossen, daß die Mehrung, so ihm verheißen war, würde in Ephraim geschehen, des Name herkommt von dem hebräischen Wort parah, das heißt, Frucht bringen; wie es Joseph droben in 41. Cap. V. 52. auslegt, da der Text sagt: „Den andern hieß er Ephraim; denn GOtt, sprach er, hat mich lassen wachsen in dem Lande meines Elends.“ Und um der Ursache willen wird er seinen Bruder Manasse vorgezogen werden, daß er also die Bedeutung seines Namens erfüllen möchte.

17. Darum hängen die Worte in der Rede, so Jakob hier gethan, recht fein an einander. Denn wenn du fragen willst: Was er dazu für Ursache habe, daß er Josephs Söhne zu seinen Kindern annimmt? antwortet er darauf also: Darum, daß der allmächtige GOtt mir zu Luz erschienen, und verheißen hat, daß ich sollte gemehret werden. Also kommt beides die Erscheinung und Verheißung Gottes mit dem Namen des Sohnes Josephs fein überein. Danach legt er sie auch aus. Als wollte er sagen: Du, Joseph, warst ein verlorner und verkauster Sohn, aber ich habe dich nun durch eine reiche und herrliche Erstattung wieder bekommen; denn ich will deine zwei Söhne, so von dir geboren sind, in die Zahl meiner Söhne und der zwölf Stämme annehmen.

18. Darum wird auch 1 Chron. 6, 1. gesagt, daß die Erstgeburt Rubens sei gegeben worden den Kindern Josephs, des Sohns Israel. Es gebührten aber dem Erstgebornen nach dem Gesetz zwei Theile am Erbe, nämlich das Königreich und Priestertum; und Joseph war vom rechten Weibe Jakobs, die er am liebsten gehabt, geboren; obwohl Ruben, der Erstgeborne von Lea, der Zeit nach vor ihm hergegangen. Da er aber seines Vaters Bett verunreinigt, hat er sein Recht und Herrlichkeit verloren. Und im nächstfolgenden Capitel V. 4. wird der Vater zu Ruben sagen: „Du sollst nicht der Oberste sein“; und Mose setzt noch eins dazu, 5 Mos. 33, 6., und sagt von Ruben: „Und sein Pöbel sei gering“ etc. Darum wird die Vermehrung dem Ruben genommen und dem jüngsten Sohne und seinem Bruder Joseph gegeben. Wie denn Jakob daselbst auch nicht auf die andern Brüder, sondern sonderlich auf Joseph und seine Söhne zieht, die da sollten wachsen und groß werden vor den andern Stämmen. Und das ist nun die Ursache, daß Jakob die Söhne Josephs zu Kindern angenommen. Darum sagt er hier zu Joseph: Du und deine zwei Söhne sollen mein sein. Ei, warum das? Darum, sagt er, daß mir Gott erschienen ist zu Lus, und sprach zu mir: Siehe, ich will dich wachsen lassen und mehren etc. Also hängt nun Alles sein an einander, was Jakob hier redet, wenn man die Verheißung betrachtet, nämlich, als die Ursache und die Wirkung.

\quellseite{477}

19. Er nennt aber flärlich seine zwei Söhne, Ruben und Simeon, daß an derselbigen Statt Ephraim und Manasse kommen sollen. Denn es hätte zwar die Erstgeburt ihnen gebühret; aber der Vater hat Ruben um der Blutschande willen, und Simeon deswegen, daß er beim Verkauf Josephs ein Rädelsführer gewesen war, gestraft und sie solcher Herrlichkeit beraubt. Darum sind auch diese zwei Stämme die allergeringsten und verächtlichsten im ganzen Volk gewesen, und sind auch jetzt noch die geringsten in jener Hefe der Juden, wie wir schon etlichemal gesagt haben. Die Stämme aber, so an ihre Statt gekommen, damit die Zahl erfüllt worden ist, werden bisweilen nach Ephraim, bisweilen nach Joseph, und bisweilen auch wohl nach Jakob genannt.

20. Die andern Söhne aber, die Joseph noch zeugen würde, sagt Jakob, sollen mit ihrer Brüder Namen genannt werden; also daß sie gleichsam die Neffen derer würden, die er seinen Söhnen gleich macht, da er sie zu Kindern annimmt; und diese sollten Kinder Israel geschrieben und genannt werden und Brüder der andern Söhne Jakobs, wiewohl sie seine Enkel waren. Derhalben sind Ephraim und Manasse Brüder Judas und ihres Vaters Josephs um der Kindschaft willen, daß sie Jakob an Kindes Statt angenommen; welche Kindschaft ihren Grund und Kraft von Gott und das Zeugniß der Schrift und des Heiligen Geistes hat, und haben darum auch mit den andern Stämmen das Gelobte Land geerbt und besitzen und Ephraim hat dazu auch die königliche Herrlichkeit erlangt.

21. Also ehrt nun Gott den frommen heiligen Mann Joseph, daß seine Söhne nicht geringer oder niedriger geachtet werden, denn Juda und die andern Söhne Jakobs.

22. Dies letzte Stück aber haben Lyra und die Andern, so über dies Buch geschrieben, anders ausgelegt, nämlich, daß die andern Söhne Josephs nicht sollen einen eigenen Stamm machen, welches recht gesagt ist, sondern daß sie sollten mit ihren Erben unter Ephraim und Manasse mitgerechnet werden. Das ist aber falsch und unrecht geredet. Sondern so will Jakob sagen: Diese zwei Söhne Josephs, Ephraim und Manasse, will ich annehmen, daß sie meine und nicht deine Söhne sein sollen. Du bist ja wohl ihr leiblicher Vater von Natur: sie sollen aber nicht geachtet oder gerechnet werden nach solcher natürlichen Kindschaft, sondern wie ich sie zu Kindern habe angenommen; und sollen hinfort nicht deinen Namen behalten, sondern nach mir genannt werden, nämlich Kinder Israel; und sollen also das Erbe im Lande Canaan, so mir verheißen, mit meinen Kindern einnehmen.

23. Das aber Jakob dazu sagt: „Diese sollen genannt werden, wie ihre Brüder“, dieses soll nicht auf Ephraim und Manasse gezogen werden, wie es Lyra auslegt; sondern er will sagen: Die andern Brüder, diese zwei ausgenommen, sollen deinen Namen haben. Wie denn solches das Fürwort genugsam anzeigt, so dazu gesetzt ist: „in ihrem Erbtheil.“ Als wollte er sagen: Die andern werden ihre eigenen Güter haben, so unterschieden sein werden von dem Erbtheil Ephraims und Manasse, welches mein ist im Lande Canaan und das ich unter meine zehen Söhne austheile, zu welchen ich diese zwei, so von dir geboren sind, auch thun will; denn sie sollen in deren Zahl mitgerechnet werden. Die andern werden ihre Güter sonst an andern Orten besitzen, es sei gleich in Egypten oder in Arabien etc. Und sie mögen entweder nach deinem oder aber sonst eines Andern Namen genannt werden, so sollen sie doch nicht Kinder Israel heißen; sie mögen aber wohl Josephiten heißen, wo etliche nach diesen Zwei, so Israeliten sollen genannt werden, entweder auf diese Welt schon geboren sind, oder aber hinfort noch von dir werden geboren werden.

V. 7. Und da ich aus Mesopotamien kam, starb mir Rahel im Lande Canaan, auf dem Wege, da noch ein Feldweg war gen Ephrath; und ich begrub sie daselbst an dem Wege Ephrath, die nun Bethlehem heißt.

24. Dies scheint auch sich nicht fein zu dem, was droben gesagt worden ist, zu schicken. Darum erregen die Andern, so über dies Buch geschrieben, mancherlei Fragen, die ich alle stehen lasse; denn sie zeigen nicht an, was Jakob in seiner Rede möchte ausgelassen haben. Jakob erzählt, daß ihm auf dem Wege, da er von Luz zu seinem Vater Isaak gezogen, Rahel zu Hebron in der Geburt gestorben sei; wie das aus dem, so droben gesagt worden, offenbar ist. Bethlehem aber oder Ephratha ist nicht weit von Hebron und Jerusalem gelegen gewesen, und ich halte dafür, daß diese beiden Städte ohngefähr fünf Meilen von einander gelegen sind. Bethlehem aber ist etwa auf dem halben Wege zwischen den zwei Orten gelegen und liegt zwei Meilen von Jerusalem.

25. Sie disputiren auch von dem Worte kibrat; aber wir verstehen es einfältig, daß es heiße ein Feldwegs, wie wir es auch verdeutscht haben; auf Lateinisch stadium.

26. Was meint aber Jakob mit dieser Erzählung, die an diesen Ort gar nicht zu gehören, und auch nicht mit dem zusammenzuhängen scheint, davon hier doch gehandelt wird? Wir haben zuvor die Ursache angezeigt, warum er die Verheißung habe angezogen, nämlich, daß er darauf hat wollen die Kindschaft Ephraims und Manasses gründen. Dahin gehört nun dies auch. Denn er ist nicht so ein alter kindischer Mann, sondern er will so viel sagen: Da mir Gott verheißen hatte, daß er mich wollte lassen wachsen und mehren und mein Weib mir schon zwei Söhne geboren hatte, hoffte ich mehr Kinder mit ihr zu zeugen, daß also meiner Kinder viele würden und mein Geschlecht laut der Verheißung gemehret würde: aber daßelbige Weib ist mir unversehens, daß ich nie erwartet hätte, dahin gestorben, indem sie durch die Schmerzen der Geburt ist überwunden worden. Darum ist vonnöthen gewesen, daß ich deine Söhne zu Kindern annähme, da ich selbst keine leiblichen Kinder gehabt, und die Verheißung nicht konnte verstanden werden, daß die Vermehrung von mir selbst sollte herkommen, da nämlich mein Weib gestorben und begraben war. Also gehört nun dies auch, was Jakob hier sagt, zu der Ursache, warum er Josephs Söhne zu Kindern angenommen.

27. Der Juden Rabinen bringen unnützes Geschwätz vor, nämlich, daß Jakob sich vor seinem Sohn Joseph sollte entschuldigt haben, warum er seine Mutter bei Bethlehem begraben hätte und nicht in der zwiefachen Höhle. Aber dies ist ein vergeblich loses Geschwätz. Denn er hätte leichtlich den Leib der verstorbenen Rahel etliche Meilen Weges können führen lassen zum Begräbniß der Patriarchen. Derhalben ist es keine Entschuldigung, sondern eine einfache Erzählung. Sie haben auch an diesem Orte nicht recht angezeigt den Text aus dem Propheten Jeremia von dem Heulen und Schreien der Rahel; sondern diese unsere Meinung und Verstand halte ich für ganz gewiß, nämlich, daß Jakob, da er sehr armselig und elend war mancherlei Anfechtungen und Trübsal halben, so er erlitten, sich selbst auf diese Weise aufgerichtet und gestärkt habe: Soll mein Same gemehret werden, will vonnöthen sein, daß ich Andere zu Kindern muß annehmen; denn das Weib, von der ich diese Verheißung verstanden, ist gestorben. Dann der Meinung bin ich nie gewesen, daß die Verheißung sollte Lea angehen; ich habe aber einen Sohn von Rahel, nämlich Joseph, und noch einen, Benjamin, in welches Geburt sie kümmerlich dahin gestorben ist. Dieweil mir aber nun meine Hoffnung gefehlet, sagt Jakob, so will ich nun deine zwei Söhne zu Kindern annehmen. Und das aus zweierlei Ursachen: zum Ersten, dieweil mir die Verheißung ist gegeben; darnach darum, daß mir mein Weib gestorben und die Verheißung nunmehr aus mir selbst nicht kann erfüllt werden.

28. Dies ist eine solche Ursache und Meinung, die die Historie selbst mit sich bringt und damit sie sich auch fein reimt, und ist der erste Theil dieses Capitels. Denn bis daher ist der Grund dieser Kindschaft gelegt worden; welcher steht in der göttlichen Verheißung, welche Jakob durch Erleuchtung des Heiligen Geistes von den Söhnen Josephs versteht.

\quellseite{478}

V. 8—12. Und Israel sahe die Söhne Josephs, und sprach: Wer sind die? Joseph antwortete seinem Vater: Es sind meine Söhne, die mir GOtt hie gegeben hat. Er sprach: Bringe sie her zu mir, daß ich sie segne. Denn die Augen Israel waren dunkel worden vor Alter, und konnte nicht wohl sehen. Und er bewachte sie zu ihm. Er aber küssete sie, und herzte sie, und sprach zu Joseph: Siehe, ich habe dein Angesicht gesehen, daß ich nicht gedacht hätte; und siehe, GOtt hat mich auch deinen Samen sehen lassen. Und Joseph nahm sie von seinem Schooß, und neigte sich zur Erde gegen sein Angesicht.

29. Dies ist der andere Theil dieses Capitels, darin gehandelt wird von der Adoption, dazu Jakob bisher eine Rede gemacht hat. Warum fragt er aber also: „Wer sind die?“ als hätte er seines Sohnes Kinder nicht gekannt, die schon zu ihren Jahren gekommen waren. Denn er war siebenzehn Jahre in Egypten gewesen, und Manasse ist geboren im andern oder dritten wohlfeilen Jahr, Ephraim aber im vierten Jahr, und sind jetzt drei oder vier und zwanzig Jahr gewesen. Mose hat selbst die Ursache hernach gesetzt, da er sagt: „Denn die Augen Israel waren dunkel worden vor Alter.“ Denn sie sind nicht so gar nahe bei ihm gestanden wie Joseph, sondern sind etwas weiter von ihm geblieben aus Demuth und Ehrerbietung. Die andern Brüder aber sind nicht dabei gewesen. Derhalben, da sie Jakob nun gesehen hat weit von sich stehen, und sie am Angesicht nicht hat erkennen können, fragt er: „Wer sind die?“ und heißt, man solle sie beide zu ihm bringen, auf daß er sie doch zum wenigsten möge betasten und anrühren, dieweil er ja mit seinen Augen so übel hat sehen können.

30. Denn das Alter macht die Augen dunkel, und macht, daß die Ohren und alle Sinne auch stumpf werden. Nach dem Hebräischen lautet es also: Oculi gravescebant prae senio, erant gravati: Die Augen waren ihm vor Alter schwer geworden, die nicht wohl noch steif sehen konnten. Jakob war alt und fält zu sehen, zu hören und zu glauben. Also redet Jesaia auch an 6. Cap. 23. 10.: „Laß ihre Ohren dicke sein, und blende ihre Augen.“ Als wollte er sagen: Sie sind hart das Evangelium zu hören, lassen sich die Predigt des Evangeliums nicht bewegen, fragen nichts darnach, sondern verachten’s, nicht anders, denn als wenn sie etwa von ferne einen unbekannten fremden Laut höreten, den sie kaum hören könnten; wie man in der lateinischen Sprache sagt: graviter audire, übel hören können.

31. Da sie aber nun näher herzu gekommen, da hat er sie, als seines Sohnes Söhne, geküßt und geherzt; und dasselbe hat er mit großer Lust und Freude gethan, wiewohl sie schon bei Jahren waren. Dies alles hat dazu gedient, daß er sie zu seinen Kindern hat annehmen wollen. Darum freuet er sich im Heiligen Geist und sagt also: „Ich hätte das nie gedacht“ oder gehofft. Das hebräische Wort heißt pilalthi, auf Lateinisch: Non judicavi, non feci hoc decretum vel judicium etc., das ist: Ich hätte das bei mir nicht beschließen können, hätte es auch in meinem Sinn nicht dürfen nehmen, daß mir es dazu kommen würde, daß ich dein Angesicht hätte sehen sollen. Also fährt er nun heraus mit einem Ausruf, frohlockt und thut eine Danksagung von ganzem Herzen, da er an diese unermessliche Güte GOttes denkt. Als wollte er sagen: Du wärest doch in meinem Herzen gar todt und ich hätte an dir gar verzagt. Sind aber nun die Werke GOttes nicht wunderbarlich in der Regierung der Heiligen und Gläubigen? O wie ist dies so eine große und wunderbare Wohltat GOttes, der dich mir hat wieder gegeben, daß ich nicht allein dein Angesicht sehen mag, darauf ich gar keine Hoffnung mehr gehabt, daß ich dich hätte wieder sehen sollen; sondern er hat mich auch deinen Samen sehen lassen, daß sehr schöne fromme Jünglinge sind.

32. Dies hat er alles gesagt, da er voller Freude und Heiligen Geistes war, und wird daselbst diese Kindschaft bestätigt durch den Heiligen Geist, wie andere Werke der Patriarchen, deren sie sich beflissen beide im weltlichen und Hausregiment; wiewohl sonst keine Zeugen dabei gewesen, auch weder Briefe noch Siegel dabei gemacht worden.

33. Da er sie aber nun geküßt hatte und sie jetzt segnen wollte, ist Joseph herzugetreten; und dieweil er befürchtet, es möchte vielleicht der Vater mit seinen dunkeln Augen die zwei Brüder nicht unterscheiden können und wohl vielleicht den Jüngern dem Aeltesten vorziehen, da nimmt er sie von seinem Schooß und neigt sich zur Erde gegen sein Angesicht. Das hebräische Wort aph haben sie gemeiniglich im Lateinischen verdolmetscht, facies, nares, Angesicht oder Nase; als, im 18. Psalm V. 9. steht, wie sie es aus dem Hebräischen gegeben haben: Ignis accensus in naribus ejus etc. Joseph hat sich gegen das Angesicht oder Nase geneiget, das ist, Jakob hat sein Angesicht zu Joseph und seinen Söhnen gewendet, und da dieselbigen nun zu Kindern und unter der Zahl der Kinder Israel angenommen waren, kommt Joseph mit beiden hervor und wartet nun auf den Segen des Vaters.

34. Das ist aber ein wunderlich Ding, daß die zwei Söhne Josephs, die nun schon so groß und alt waren, noch gleichwohl auf seines Vaters Schooße gesessen haben. Aber ich halte dafür, daß da er in seinem Ruhebettlein gelegen, sie bei ihm gesessen haben, welches Mose nennt auf dem Schooße oder Knien sitzen, nämlich, zur Seite oder zu Füßen. Derhalben heißt sie nun Joseph von der Seite aufstehen, und fällt mit ihnen auf die Kniee vor Jakob, der in seinem Bettlein gesessen. Denn dies ist also die Weise gewesen, daß sie den Segen empfangen haben mit Kniebeugen und mit solchen Gebärden, als die pflegen zu haben, die da beten wollen.

\begin{quote}
\textbf{V. 13. 14. Da nahm sie Joseph beide, Ephraim in seine rechte Hand gegen Israels linke Hand, und Manasse in seine linke Hand gegen Israels rechte Hand; und brachte sie zu ihm. Aber Israel streckte seine rechte Hand aus, und legte sie auf Ephraims, des Jüngsten, Haupt, und that wissend also mit seinen Händen, denn Manasse war der Erstgeborne.}
\end{quote}

35. Der Gebrauch und Weise mit dem Handauflegen ist ein sehr alter Gebrauch und ist von den Vätern auch in das Neue Testament gekommen; wie das bei Paulus zu sehen ist, da er 1 Tim. 5, 22. sagt: „Die Hände lege niemand bald auf.“ Mose hat Josua die Hände aufgelegt. Und ist solches eine Ceremonie der Väter, Propheten und der ganzen Kirche gewesen, die zu allen Zeiten ist gebraucht worden, wenn sie Einem ein Amt oder Regiment befohlen haben.

36. Es ist aber die Erzählung und was im Text gemeldet wird, leicht zu verstehen. Joseph nimmt seinen Sohn Ephraim als den Jüngsten bei seiner rechten Hand, als der Jakob zur linken Hand sollte gestellt werden; Manasse aber bringt er bei der linken Hand zu Jakob, daß er bei der linken Hand zu Jakobs Rechten, als den Erstgebornen. Denn er weiß wohl, daß der zur Rechten steht, etwas mehr und Größeres pflegt zu empfangen, und befürchtet, es möchte eine Verwechselung oder Unordnung werden zwischen dem Segen des ältern und jüngsten Sohnes. Die Erstgeburt hatte eine doppelte Portion und Vorzug, denselbigen wollte er gern auf Manasse gewendet haben und daß dem Ephraim nur Ein Theil gegeben werde. Daselbstige ist also offenbar aus den äußerlichen Geberden, so Joseph geführt.

37. Siehe aber, welch ein großer Geist der Weissagung in Jakob, dem alten schwachen Manne, gewesen sei. Er sieht, daß Manasse zu seiner Rechten gestellt wird und der Andere zu seiner Linken; er versteht auch wohl, daß Joseph die Herrlichkeit der Erstgeburt gern auf Manasse wollte gewendet haben: aber Jakob läßt die ordentliche Meinung nach dem Gesetz und nach gemeiner Ordnung und nach dem Gesetz von GOtt eingegangen fahren, und nimmt eine andere Weise vor und folgt derselbigen.

38. Es wird 5 Mos. 21, 15—17. das Recht der Erstgeburt beschrieben, da Mose also sagt: „Wenn jemand zwei Weiber hat, eine, die er lieb hat, und eine, die er hasset; und sie ihm Kinder gebären, beide, die liebe und die feindselige, daß der Erstgeborne der feindseligen ist; und die Zeit kommt, daß er seinen Kindern das Erbe austheile, so kann er nicht den Sohn der liebsten zum erstgebornen Sohn machen, für den erstgebornen Sohn der feindseligen; sondern er soll den Sohn der feindseligen für den ersten Sohn erkennen, daß er ihm zwiefältig gebe alles, das vorhanden ist; denn derselbe ist seine erste Kraft, und der Erstgeburt Recht ist sein.“ Nach dieser Regel bleibt der erstgeborne Sohn vermöge des Gesetzes und göttlichen Befehls der Erstgeborne. Aber GOtt, der ein Stifter und Herr ist des Gesetzes, der mächtigt oder zerbricht auch wohl dasselbige wieder, und macht Jakob und Isaak zum Erstgebornen und verwirft Esau und Ismael. Eben dasselbe thut er an diesem Orte auch mit den zwei Söhnen Josephs.

\quellseite{479}

39. Denn hier gilt das nicht, das man zu sagen pflegt: Legem patere, quam tu ipse fecisti. Du mußt dich selbst auch nach dem Gesetz richten, das du gemacht hast. Gott ist dem Gesetz nicht unterworfen, und thut oftmals, das dem Gesetz stracks zuwider ist, daß wir seine Werke, seine Weisheit, Rath und wunderbare Gerichte ehren; ja, daß wir auch in aller Demuth vor ihm wandeln sollen.

40. Obwohl David der allerkleinste und verächtlichste war unter den Brüdern, ist er doch zum Königreich Israel erhoben worden. Derhalben weiß Gott wohl ein Gesetz zu geben und aufzurichten nach seiner Weise, und handelt doch bisweilen wunderbarlich ohne und wider die Regel des Gesetzes. Joseph will haben, daß Manasse soll der Erstgeborne sein; aber es gefällt Gott nicht also, der will es auf eine andere Weise machen. Derhalben bessert und ändert nun Jakob die Meinung, so Joseph gefaßt hatte, und das thut er durch sonderliches Eingeben des Heiligen Geistes, und hat ihn dazu bewogen die Verheißung, so ihm Gott gegeben, daß er fruchtbar sein und wachsen sollte etc., welche Verheißung er oftmals und ganz fleißig in seinem Herzen erwogen und bedacht hat.

41. Und wird also Ephraim, der jüngste Sohn, dem der Vater den Namen gegeben vom Wachsen, der Erstgeborne, und das durch Gottes Gnade, die eine Herrin oder Meisterin ist über das Gesetz, und gibt nicht aus schuldiger Pflicht nach dem Gesetz, oder aus Verdienst, sondern umsonst, Röm. 4, 16. Und das Gesetz soll doch in seiner Zucht, Zwang und Ordnung bleiben; aber die Gnade hat immer ihre Ausnahme und ist über das Gesetz.

42. Da aber Joseph gesehen, daß sein Vater mit den Händen einen Wechsel macht und die gemeine Ordnung umkehrt, hat er mögen denken, er muß seines Verstandes und aller Sinne beraubt sein; wie wir von alten Leuten sagen, daß sie kindisch zu werden pflegen, wenn sie an ihren Kräften, an Gedächtniß, und beinahe auch an Vernunft und gemeinem Verstande sehr abgenommen. Denn Joseph hatte aus sonderlicher Absicht seine beiden Söhne also in Ordnung gestellt, daß er damit diese Bemerkung verhüten wollte; aber sein Vater hat dem jüngsten Sohne die Erstgeburt zugetragen, das Joseph nicht gemeint oder erwartet hatte.

\section[Zweiter Theil]{Zweiter Theil. Wie Jakob Joseph und dessen Söhne segnet}

\begin{quote}
\textbf{V. 15. 16. Und er (Jakob) segnete Joseph, und sprach: Gott, vor dem meine Väter, Abraham und Isaak, gewandelt haben, Gott, der mich mein Lebenlang ernähret hat, bis auf diesen Tag, der Engel, der mich erlöset hat von allem Uebel, der segne die Knaben, daß sie nach meinem, und nach meiner Väter, Abraham und Isaaks, Namen genannt werden, daß sie wachsen und viel werden auf Erden.}
\end{quote}

43. In der gemeinen lateinischen Uebersetzung wird ohne Ursache hinzugesetzt: Filii Joseph: Die Söhne Josephs. Denn Jakob, der Vater, hat Joseph selbst gesegnet, wiewohl er den beiden Söhnen die Hände aufgelegt und sie auch gesegnet hat. Aber im Hebräischen steht: „Er segnete Joseph“, nämlich, in seinen Söhnen. Denn er hatte ihn erwählt, daß von ihm zwei Stämme herkommen sollten, die das Erbe und göttlichen Segen empfangen möchten, und er hat ihm das Recht der Erstgeburt im Lande Canaan gegeben.

44. Derhalben sind diese zwei Stämme erwählt und angenommen nicht durch ihr Verdienst oder Würdigkeit, sintemal der Vater gesegnet wird, und haben eigentlich nicht Josephiten, sondern vielmehr Jakobiten müssen genannt werden; und die Propheten brauchen doch bisweilen den Namen Joseph und auch bisweilen den Namen Ephraim; wie sie denn auch den Namen Jakob behalten, wiewohl derselbige von Gott geändert war, da er zu ihm sagte 1 Mos. 35, 10.: „Du sollst nicht mehr Jakob heißen, sondern Israel sollst du heißen.“

45. Dies sind aber Worte des Segens oder der Einweihung: „Gott, vor dem meine Väter, Abraham und Isaak, gewandelt haben“ etc. Das hebräische Wort stehet in Hithpael, als wollte er sagen: Qui se ambulare fecerunt, wie man es auf Lateinisch geben möchte; das ist: Die sich also darnach geschickt haben, daß sie möchten recht vor Gott wandeln, das ist, an den sie mit festem Glauben geglaubt haben und in deß Vertrauen und Verheißungen sie auch gestorben sind, auf die gewisse Hoffnung der zukünftigen Auferstehung; wie droben vom Glauben Abrahams und Isaaks gesagt worden ist.

46. Und setzt noch dazu: „Gott, der mich mein Lebenlang ernähret hat“ etc., deß Schaf ich bin, dafür er gesorgt hat wie der allerfleißigste Hirte, den er erlöset, errettet und geschützt hat in viel großer Trübsal bis an den heutigen Tag. Darum redet Jakob fast also, als ob zwei Götter wären. Zum Dritten sagt er: „Der Engel, der mich erlöset hat“, auf Hebräisch goel, nicht podeh, sondern ein rechter Erretter und Erlöser, und wie ein Verwandter; davon Mose sagt 5 Mos. 19, 6.: „Auf daß nicht, goel, der Bluträcher, den Todtschläger nachjage“ etc., das ist, derjenige, so Macht hat, von Rechts wegen zu rächen und zu strafen denjenigen, so den Todtschlag gethan hat.

47. Also sagt Hiob auch im 19. Cap. V. 25.: „Ich weiß, daß mein Erlöser, goel, lebet.“ Dies ist aber ein ander Wort denn padah, welches Hosea am 13. Cap. V. 14. stehet: „Ich will sie vom Tode erretten.“ Daselbstige Wort heißt so viel als, der da Macht hat, Andere zu erretten. Wie Christus der Herr unser podeh und goel geworden ist; denn er hat uns nicht allein erlöset, sondern er hat uns auch vom Rechte wegen also befreit und an sich gebracht, daß der Teufel und die Hölle auch nach dem strengen Recht ihn haben müssen los lassen, darum daß sie den unschuldigen Sohn Gottes getödtet hatten. Derhalben hat sich das Gesetz an ihm verbrannt, der Tod hat sich beschissen, der Teufel, die Hölle und Sünde haben sich vergriffen; da sind sie alle schuldig geworden an Gott und diesem seinem Sohne Jesu Christo, der nun gut Recht hat zu seinen Feinden. Denn warum hast du, o Gesetz, den Sohn Gottes gekreuzigt? Du Teufel, Tod und Hölle, warum hast du den Unschuldigen erwürget? „Wir haben ein Gesetz“, sagten die Juden Joh. 19, 7., „und nach dem Gesetz soll er sterben; denn er hat sich selbst zu Gottes Sohn gemacht.“ Darum stehet ihn nun an, da er von den Todten erstanden ist und wider euch den Sieg behalten hat. O, ich, sagt er, bin Gottes Sohn, bin eine solche Person, die unüberwindlich ist. Was willst du, Satan, du Gesetz, Tod und Hölle, nun sagen? „Der Tod ist verschlungen in den Sieg“ etc., 1 Cor. 15, 55.

48. Und hiervon hat Christus nach seiner Auferstehung befohlen, daß es in der ganzen Welt sollte verkündigt und jedermann das Evangelium gepredigt werden. Komm her, sagt er zu uns, glaube an mich, laß dich taufen, ich will dir meinen Sieg schenken, du sollst nicht verdammt werden, sondern wenn du auch schon stirbst, sollst du in meinem Namen leben: du sollst nimmermehr sterben. Denn „ich bin die Auferstehung und das Leben. Wer an mich glaubt, der wird leben, ob er gleich stürbe“, Joh. 11, 25.

49. Dieser ist nun unser goel, unser Erlöser, welchen Jakob an diesem Orte einen Engel nennt. Als wollte er sagen: Der Engel oder Erlöser, der mich mit allem Recht hat erlösen und schützen können von allem Uebel des Gewissens, des Gesetzes, der Sünde und Todes, auch vor der Gewalt meines Bruders Esau, und von allem Jammer und Noth, damit ich mein ganzes Leben über habe kämpfen müssen.

50. Daraus ist nun zu sehen, was für Leute die heiligen Patriarchen gewesen sind, von welchen der 4. Psalm V. 4. sagt: Wahrlich, „der Herr führet seine Heiligen wunderlich.“ Denn wenn sie anfangen zu reden vom Glauben und von den Verheißungen Gottes, fahren sie gleichsam hoch über alle Himmel. Wiederum aber, wenn sie schwach werden, wie droben, da Jakob mit dem Engel gerungen, und sonst an andern Orten mehr, werden sie so gar darniedergeschlagen, daß sie nicht anders denn wie arme Menschen erzeigen; daß es scheint, als ob sie beinahe gar verzagen und in die tiefste Hölle hinunter versenkt seien. Wie wir denn solch Fürchten und wunderbares Schrecken und Zagen droben an Jakob gesehen haben, nicht anders, denn als ob er Einer wäre von denen, so gar verworfen und zu ewiger Pein verdammt sind. Wiederum aber, wie hier an diesem Orte, ist er so fröhlich, triumphirt und frohlockt dermaßen, als ob gar nirgend Gefahr oder Furcht mehr vorhanden sei, die ihm begegnen möchte.

51. Also sind nun die heiligen Patriarchen, wenn es sie gelüstet hat, Hausväter gewesen, und haben das politische Regiment verwaltet, haben ihre Weiber, Kinder und Gesinde regiert, haben etwa auch mit ihren Weibern gezankt, und sind so gar mit weltlichen Händeln umgegangen, daß sie auch für die allergeringsten Leute unter denen, so im Haus- und Regierungsstande anzutreffen, sind angesehen worden. Bisweilen fahren sie hoch her über alle Himmel, bisweilen sind sie in der Hölle, bisweilen aber bleiben sie mitten in der Welt stecken. Es sind eitel Wunderleute.

\quellseite{480}

52. Ihr Leben ist wunderbarlich, ihre Rede auch wunderbarlich. Wie Jakob hier Gott einen Engel nennt, eben auf dieselbige Weise, wie er droben, 1 Mos. 32, 30., sagt, da er mit dem Engel gekämpft hatte: „Ich habe Gott von Angesicht gesehen.“ Denn dieser Engel ist derselbe Herr oder der Sohn Gottes, den Jakob gesehen hat, und der von Gott in die Welt sollte gesandt werden, daß er uns verkündigte die Errettung vom Tode, Vergebung der Sünden und das Himmelreich. Und dieser Engel ist unser goel, unser Erlöser oder Rächer, der uns mit allen Recht errettet und befreit von der Gewalt des Teufels, welcher nun dem Gesetz unterworfen ist, darum daß er den Sohn Gottes erwürget hat; und müssen derhalben nun Gesetz, Tod und Teufel gar verstummen, und die Hände herreichen und Christo dem Herrn, als dem rechten Siegesmann, gewonnen geben.

53. Derhalben ist fleißig zu merken, daß Jakob hier von Christo, dem Sohne Gottes, redet, welcher allein der Engel oder Gesandte ist, in der Zeit ein wahrer Mensch geboren von der Jungfrau Maria, nicht der Vater, auch nicht der Heilige Geist. Denn er unterscheidet förmlich die drei Personen, und setzt doch gleichwohl hinzu: „Der seyne die Knaben.“ Das Segnen schreibt er förmlich allein Gott zu. Er sagt nicht: „sie segnen“, als ob er von Vielen redete, hat auch das vorige nicht wiederholt; sondern setzt in dem Einen Werk des Segnens die drei Personen zusammen, Gott den Vater, Gott, der ihn ernährt und sein Hirte ist, und den Engel. Derhalben sind diese drei ein Gott und ein Segen. Der Engel richtet eben dasselbe Werk aus, das auch der Hirte und der Gott seiner Väter ausrichtet.

54. Also haben die Väter diesen Artikel von der heiligen Dreifaltigkeit recht wohl verstanden, den wir heutiges Tages auch also lehren.

Und Jakob geht in demselben rechten vollkommenen Glauben und Wort daher über alle Himmel. Denn er sagt: Ich segne diese Knaben nicht, sondern Gott segnet sie durch meinen Mund, nämlich, der Gott meiner Väter, der Gott, der der rechte Hirte ist und uns ernährt, und der Engel; der, welcher dreifaltig und doch einig ist, der segnet sie. Derhalben hat Jakob wohl verstanden, daß der Sohn Gottes hat sollen Mensch werden, und in die Welt gesandt, gekreuzigt und wieder vom Tod auferweckt werden. Wiewohl er dasselbe nicht so klar und deutlich ausspricht, wie es jetzt gelehrt wird, sagt er es doch nachdrücklich genug; und die Andern, so diesen Artikel geglaubt, haben es auch deutlich also verstanden.

55. Das Wort „segnen“ aber muß weitläufig und als mit Nachdruck geredet verstanden werden, daß seine Kraft auch recht möge gefaßt und verstanden werden; wie wir droben auch davon Erinnerung gethan haben. Denn wenn der Segen von Gott herkommt, muß man es nicht so absolut, sondern relatie verstehen; und wird damit angezeigt, daß diejenigen, so gesegnet werden, zuvor dem Fluch unterworfen sind. Solchen Nachdruck (emphasis sive epithasis) desselben Worts muß man fleißig betrachten. Denn alle Menschen sind in Adam vermaledeyet, sind der Sünde, dem Gesetz, dem Tode und Gewalt des Teufels unterworfen; wie Paulus sagt Gal. 3, 10.: „Verflucht sei jedermann, der nicht bleibet in alle dem, das geschrieben stehet in dem Buch des Gesetzes.“ Das Gesetz beschuldigt und verdammt alle Menschen; denn es straft und überzeugt sie, daß sie Sünder sind, und bringt ihre Sünde und Missethat an das Licht.

56. Darum wenn Gott sagt: Ich will dich segnen, da wird dem Menschen beides leiblicher und geistlicher Segen angeboten. Wie an diesem Ort nicht allein der leibliche Segen soll verstanden werden, dadurch diese Knaben leiblich wachsen und gemehret werden möchten; sondern daß der Engel wolle bei ihnen sein, der Jakob auch von allem Uebel errettet hat.

57. Und dies ist der Segen des Neuen Testaments, davon St. Paulus in seinen Episteln an die Römer und Galater reichlich handelt, da er diese Stücke gegen einander hält: verflucht und gesegnet sein, Sünde und Gnade. Und der Engel segnet zugleich den Samen Abrahams oder die Juden und auch die Heiden; wie der Spruch lautet Röm. 9, 8.: „Nicht sind das Gottes Kinder, die nach dem Fleisch Kinder sind; sondern die Kinder der Verheißung werden für Samen gerechnet.“

58. Derhalben sind dies sehr hohe und wichtige Worte, in welchen Jakob sagt, daß Gott Vater, Sohn und Heiliger Geist diese Knaben segnen; und daß im Evangelium Johannis Cap. 5, 17. Christus auch selbst sagt: „Mein Vater wirket bisher, und ich wirke auch“; item, daselbst sagt er weiter V. 21.: „Wie der Vater die Todten auferweckt und machet sie lebendig: also auch der Sohn machet lebendig, welche er will.“ Das kommt nun mit diesem gegenwärtigen Text sehr fein überein. Lebendigmachen ist allein Gottes Werk, welches nach Gottes Willen und durch seine Gnade geschieht. Wie nun der Vater will gnädig sein und lebendig machen: also will ich auch thun, spricht Christus, das ist einerlei Wille und einerlei Werk und sind gleichwohl drei unterschiedliche Personen. Darum hat Augustin recht gesagt: Opera divinitatis ad extra sunt individua, das ist: Was Gott gegen die Creatur thut, das thun alle drei Personen ohne Unterschied. Also ist einerlei Segen Gottes und seines Sohnes, der hier ein Engel genannt wird.

59. Dies haben die heiligen Väter recht verstanden und gelehrt. Derhalben sieht du daran die große treffliche Erkenntniß Gottes und den Glauben, so sie gehabt; und wenn sie davon reden, sieht man vornehmlich an ihnen keine anderen, häuslichen, weltlichen oder politischen Werke, keine Sünde, kein Jagen; sondern da werden sie hoch geführt über Wolken und Luft. Denn so sehen wir, daß der Patriarch Jakob gleichsam gar trunken und voll ist im Glauben und Erkenntniß Gottes, und daß er sonst nichts Anderes thut, denn nur Worte des Glaubens redet, und immer auf die Verheißung bringt und dieselbige wiederholt, darin ganz und gar lebt, regiert und selig ist; und theilt nun durch göttliche Gewalt denselben verheissenen Segen auch diesen Kindern mit, und verheißt ihnen dazu auch das zeitliche Leben und darin alle Wohlfahrt.

60. Denn jetzt sind sie nun dessen gewiß und gar versichert, daß sie so lange leben sollen, bis sie Kinder zeugen werden. Also wird der Tod von ihnen weggenommen und mit demselben auch der Zorn Gottes, Teufel und Hölle, und wird unter sie das Erbe ausgetheilt durch diesen Segen des Vaters, welches er selbst nicht erlangt hatte, nämlich, der Besitz des Landes Canaan. Denn er hat daran nicht gezweifelt, daß er es würde empfangen und daß er mit seinen Nachkommen ein Herr und König desselben Landes werden sollte. In welchem Segen allerlei Wohlthaten, beide leibliche und geistliche, sind begriffen worden.

61. Wiewohl nun das Fleisch und die Vernunft der Dinge keines gesehen noch begriffen haben (denn was ist es, das er sagt von der Verheißung Gottes, von dem sogenannten Engel und vom Lande Canaan, dieweil sie jetzt im Elend und Knechte sind in Egypten? Das aber alles unangesehen), sagt gleichwohl Jakob voll Glaubens und mit großem starkem Muthe frei heraus: Ich habe die Verheißung, ich weiß und glaube, daß ich das Land Canaan erben und besitzen soll; dessen bin ich in der Verheißung gewiß, als ob ich es allbereit inne hätte und darin wohnte.

62. Wie wir uns der Güter, so uns Gott verheißen hat, gewisser trösten können und sollen, als ob wir dieselben schon in Händen hätten. Ich habe das Himmelreich, die Taufe, das Wort, das Abendmahl des Herrn: die Güter stehen mir gewisser zu und sind mehr mein, denn dies Leben selbst, darin ich hier lebe. Ja, sagst du, fühle ich es doch nicht. Antwort: Man muß lernen mit den lieben Vätern darauf hoffen, daß man nicht sieht, und muß die Augen hinauf gen Himmel richten. Du mußt mit den Vätern lernen auf den Wolken reiten. Das Himmelreich ist schon dein, dir sind deine Sünden vergeben durch das Blut des Lämmleins; siehe nur zu, daß du recht glaubest. Es ist kein Scherz mit Gottes Verheißungen, Gott lügt nicht. Er sagt, du seiest gerecht geworden, du seiest heilig und selig; wie dieselbigen Güter uns ja auf das klarste und prächtigste durch das Wort und Predigtamt vorgetragen und angeboten werden; aber daran mangelt es nur, daß wir nicht fest am Wort und den Verheißungen Gottes halten und mit rechtem Glauben denselben Beifall geben.

63. Darum haben die Väter das auch sonderlich vor uns gehabt, nämlich, daß ihnen dies Leben hat sollen verlängert werden, daß wir nicht können gewiß sein; denn wir haben kein
\quellseite{481}
Wort Gottes davon. Aber sie haben es gewißlich dafür gehalten, daß sie wachsen und das Land besitzen würden, wie ihre Väter Abraham, Isaak und Jakob dasselbe im Glauben und Verheißung gehabt; welcher Besitz ihnen gewisser gewesen ist, denn da sie das Land äußerlich mit der That selbst besessen haben ohne Glauben; wie das ihre Historie bezeuget, da Ephraim nach Assyrien verstoßen wird und Juda gen Babel geführt worden ist. Das ist zwar ein gering Ding, wenn es nur ein thatsächlicher und juridischer Besitz ist; wenn er aber theologisch ist und nach der Schrift gefaßt wird, da bleibt die Herrschaft, wenn wir auch schon darüber hinsterben.

64. So haben wir nun an den Vätern Exempel eines vollkommenen beständigen Glaubens. Jakob bietet ihnen an und bestimmt ihnen nicht allein dies leibliche Leben, sondern verheißt solches auch ihrem Samen und Nachkommen.

65. Dem Glauben aber folgt endlich das Gut auch nach in der wahren That. Wo aber der Glaube und der geistliche Besitz nicht ist, da kann auch der Besitz, so in der That vorhanden ist, nicht lange bestehen.

66. Jakob setzt auch noch ein Stück dazu, nämlich, daß die Knaben nach seiner Väter Namen sollen genannt werden. Welches nur lauter Worte des Glaubens und der Verheißung sind, und von den Juden nicht verstanden werden, als die sich allein vom fleischlichen Samen rühmen. Aber Paulus widerlegt ihnen denselben Ruhm gar herrlich, da er Röm. 9, 8. sagt: „Nicht sind das Gottes Kinder, die nach dem Fleisch Kinder sind; sondern die Kinder der Verheißung werden für Samen gerechnet.“ Denn die Väter haben die leibliche und geistliche Vaterschaft zu einander gesetzt, das ist, sie haben den Glauben der Verheißung geerbt, und nicht allein das äußerliche Gut, so sie haben besitzen sollen.

67. Darum ist der Glaube die Summa und das Hauptstück des ganzen Handels, welchen zwar das Fleisch fast hart ansicht und sehr hindert: aber wir müssen uns mit der Verheißung und mit dem Wort trösten und aufrichten, und sollen uns vor keinerlei Gefahr fürchten, so lange wir den Engel, der unser goel, das ist, Erlöser ist, nicht verleugnen und von ihm nicht abfallen. Denn wir müssen auch bisweilen etwas schwach werden; wie die Stärke und der große Muth an den Vätern nicht für uns für gewähret hat. Jakob ist an diesem Ort gar stark im Glauben an die Verheißung, da ist jetzt keine Schwachheit, kein Zagen an ihm, es ist alles eitel hoch himmlisch Ding, damit er umgeht. Aber droben, da Joseph verkauft ward, da war Jakob der allerelendeste und schwächste Mensch auf Erden.

68. Darum sollen wir Gott bitten, daß er uns nicht in Versuchung führe, sondern uns von dem Uebel erlösen wolle, und das ungläubige, unsinnige, böse Thier, unser Fleisch und Leib des Todes, damit wir uns schleppen müssen, zwingen und inne halten wolle, welches uns immer plagt und tödtet. Denn es widerstrebt dem Glauben und der Verheißung, untersteht sich, daß es uns nur auf dasjenige, so jetzt gegenwärtig und vor Augen ist, ziehen und reißen möge. Es will nur ein Jurist sein und kein Theolog. Die Gerechtigkeit aber führt uns nicht in den Himmel, sondern die Verheißung und der Glaube, welches himmlische Gaben sind und alle Wunderthaten und alle Werke, beide politische und häusliche, übertreffen, vielmehr aber die scheußlichen selbsterdichteten Werke der Mönche, und die ganze Disciplin oder äußerlichen Ceremonien des Gesetzes.

69. Derhalben lasset uns in der Einfalt und Reinigkeit des christlichen Glaubens bleiben, daß wir Abrahams, Isaaks und Jakobs Kinder sind, nicht nach dem Fleisch, sondern vielmehr nach dem Geist. Denn wenn ich der göttlichen Verheißung glaube, so bin ich dessen gewiß, daß mein Leben Gott wohlgefällt, und besser ist, denn alle Mönchsorden sein können: denn es macht einen himmlischen Mann, der den Tod überwindet, einen Erben des ewigen Lebens und der den Teufel mit Füßen tritt; wie der 91. Psalm v. 13. sagt: „Auf den Löwen und Ottern wirst du gehen, und treten auf den jungen Löwen und Drachen“ etc. Solche Stärke und Gewalt ist den Christen besonders eigen.

70. Wie hier Jakob gar ein Engel ist, ja, er ist ein großer trefflicher Prediger der Gottheit, der seinen Nachkommen und Erben die rechte Kraft der Verheißung und göttlichen Segens gegenwärtig vorträgt. Denn er sieht auf seine Nachkommen oder Samen, nicht nach dem Fleisch, sondern vornehmlich nach dem Geist und Glauben; welchen geistlichen Samen die Juden auch kennen und sich dessen rühmen sollten, nämlich, daß sie Kinder wären der Väter und Erben der göttlichen Verheißungen. Aber dieselben Verheißungen treten sie schändlich mit Füßen, verwerfen und lästern sie, und rühmen sich derweil hoch allein des leiblichen Königreichs halben.

71. Das sind aber rechte Kinder Abrahams, die da Erben sind des Engels, der mit Jakob und durch Jakob geredet hat. Und wenn Gott solche Kinder nicht findet, so nennt er das sein Volk, welches zuvor sein Volk nicht war, und aus einem närrischen Volk erzürnt er den fleischlichen Samen Abrahams, 5 Mos. 32, 20. 21., daß Abraham der Heiden Vater sei, so die Verheißung im Glauben, ohne Fleisch und Wirklichkeit, annehmen; da die Juden, dem Stracks zuwider, dem Fleisch ohne die Verheißung nachfolgen.

72. Was nun weiter folgt, gehört zum leiblichen Segen, wie das, so wir bisher gehört, den geistlichen Segen betrifft und angeht. Denn also lauten die Worte im Text: „Daß sie wachsen und viel werden“ etc. Und gleichwohl will Jakob hier auch angezeigt haben, daß der geistliche Segen auch mit darein sollte gewickelt werden, und daß Ephraim nicht allein soll gemehret werden und wachsen in dem Lande oder sonst in einem andern Winkel oder Ort der Welt; sondern daß über dem ganzen Erdkreis der Glaube der Verheißung Abrahams soll ausgebreitet werden.

73. Das hebräische Wort, dagah, heißt wachsen oder gemehret werden, wie Fische. Als wollte er sagen: Piscescant, wie man es im Lateinischen ausdrücken möchte, daß ist, wie die Fische sollen sie wachsen etc. Es kommt her von dag, das heißt, piscis, ein Fisch. Es ist ein verbum nominale, wie man es in der Grammatik zu nennen pflegt; das heißt, ein von einem Hauptwort gebildetes Zeitwort; wie in jeder Sprache aus Hauptwörtern Zeitwörter werden und umgekehrt. Denn es ist kein Thier, das in so großer Anzahl wachse und gemehret werde wie die Fische. Die Vögel haben auch wohl ihre Fruchtbarkeit und Mehrung; denn es brütet oftmals ein Vöglein vier oder fünf Junge aus; aber unter allen Thieren ist keins so fruchtbar als die Fische. Derhalben wird dasselbige Wort metaphorisch auf eine unermeßliche Vermehrung gezogen und davon gebraucht.

74. Nun tritt aber Joseph hinzu und will weiser sein denn sein Vater. Und nach dem Gesetz urtheilt er zwar recht, aber er läuft doch an und fehlet weit, wiewohl er ein großer und trefflicher Mann ist und ein Prophet voll heiligen Geistes.

75. Hier streiten zwei sehr heilige Männer und zanken mit einander, und sind uneins über der ganzen Hauptsache; wiewohl sie in ihren Herzen übereinstimmen, aber die Sachen sind sehr ungleich. Der Vater will Ephraim, den jüngsten Sohn, segnen, das will Joseph nicht haben. Nun sind sie gleichwohl beide sehr geistliche Männer, aber der Sachen und Meinung halben ist unter ihnen ein großer Zanck, und will einer dies, der Andere das haben.

76. Wir haben droben dergleichen Exempel auch gehabt mit Abraham und Sarah 1 Mos. 16, 5. Isaak und Rebekka 1 Mos. 25, 28., da dieselbigen Eheleute, welche einander sehr lieb und werth gehabt, doch wider einander sind. Abraham will den Ismael vorziehen, aber Sarah hält Isaak für den Erstgebornen. Isaak hält auf Esau, Rebekka aber auf Jakob. Und haben hart mit einander gezankt, nicht zeitlicher oder geringer Sachen halben, sondern wegen der Nachkommen oder Erben und über den göttlichen Segen. Und wenn es solche leichtfertige Leute gewesen wären, die sich leichtlich hätten erzürnen lassen, hätten sie eigentliche Ursache zu einer großen Feindschaft gehabt; denn sie haben unter einander über einer sehr großen wichtigen Sache gezankt.

77. Unsere Könige, Fürsten und Bischöfe erregen oftmals um geringer Beleidigungen willen große schädliche Kriege und vergießen unschuldig Blut; und streiten doch nicht mit einander wegen irgend eines Nutzens, oder auch nicht Wohlust oder Ehren halben, sondern fechten gegen einander nur aus lauter Bosheit und greulicher Unsinnigkeit. Hier aber disputiren diese großen trefflichen Männer mit einander von einer großen wichtigen Sache, und haben beiderseits sehr treffliche und wichtige Argumente. Abraham, Isaak und Joseph stehen fest auf Gottes Gesetz, welches ja die höchste Weisheit ist, so auch mit dem natürlichen Rechte übereinstimmt. Denn sie schließen also: Ismael, Esau, Manasse ist ja der Erstgeborne, darum gebühren ihm zwei Theile am Erbe. Und Isaak hält sehr hart über dieser Rechnung: Esau ist mein erstgeborner Sohn, darum will ihm die Herrlichkeit der Erstgeburt gebühren von Natur, nach dem Gesetz und endlich auch nach allem Recht. Also bleibt Joseph auch darauf bestehen, als wäre er gar darüber verhärtet in der Weisheit des Gesetzes und der Natur; dem widersetzt aber Jakob, wie Rebekka und Sarah dasselbige Recht auch angefochten und sich darüber gelegt haben.

\quellseite{482}

78. Nun sollte aber gleichwohl das Weib solchem großen trefflichen Mann, der ein Vater ist der Verheißung, billig weichen und ihm Recht gegeben haben. Aber Sarah fährt frei heraus und sagt zum Mann: „Stoß die Magd hinaus mit ihrem Sohn; denn der Magd Sohn soll nicht erben mit meinem Sohn Isaak“, Gal. 4, 30. Und Mose sagt, das Wort habe Abraham sehr übel gefallen. Also ist Isaak mit seinem Weibe über solcher Sache auch nicht einig. Denn wer wollte etwas reden oder thun wider denjenigen, der das göttliche Recht und das Gesetz der Natur für sich hat? Das ist gewiß ein schwerer Streit; und sind doch Abraham und Isaak dahin gebracht worden, daß sie ihre Unwissenheit bekennen müssen und das Recht des Gesetzes und der Natur widerrufen.

79. Wie sich Joseph seines Vaters Jakob Meinung hier auch muß gefallen lassen und was er wider das Gesetz thut. Er legt sich wohl erstlich wider den Vater, als wollte er sagen: Lieber Vater, du irrst mit dem Segen. Aber Jakob antwortet ihm: Ich irre nicht; ja, ich thue dies alles mit gutem Rath, und mit Wissen und Willen. Ich sündige nicht aus Unwissenheit, ob mir wohl die Augen dunkel sind; sondern ich sehe und erkenne sie alle beide wohl.

80. Warum thust du aber dies? Darum, daß es also geschehen muß. Denn Jakob hat sehr fest ergriffen und behalten diese Verheißung, so ihm zu Bethel ist gegeben worden, da jetzt Rahel hat sterben wollen, und deshalb keine Hoffnung da war, daß er konnte Erben haben. Er war ein Mensch voll Heiligen Geistes und hat aus der Verheißung viel genommen durch Erleuchtung des Heiligen Geistes, und hat gedacht an die vorigen Exempel mit Abraham und Isaak, da Sarah Isaak vorgezogen, der doch der zweigeborne Sohn war; dieser Isaak aber hat auch den kleinern Sohn Jakob dem Größeren, nämlich Esau, vorgezogen. Von diesen Exempeln hat er gesehen, daß sie sich recht gereimt haben mit dem, so ihm der Heilige Geist eingegeben hat. Und sind auch dazu gekommen die Sünden seiner Söhne, damit Ruben, Simeon und Levi, so die ältesten gewesen, die Erstgeburt verloren haben. Darum hat er gedacht: Ich habe nun keinen erstgebornen Sohn, habe auch keine andern Söhne mehr von Rahel, und Gott hat mir durch eine wunderbare Erscheinung etwas Sonderliches gezeigt, da er von mir selbst also gesagt: „Sei fruchtbar und mehre dich“, 1 Mos. 35, 11. Darum setzt er sich das in seinem Herzen gewiß vor, er wolle Ephraim dem Manasse vorziehen.

81. Dies muß man aber nun ziehen auf die zweierlei Lehren, so in der christlichen Theologie und Rechten gelehrt und gehandelt werden, nämlich, auf die Gnade der Verheißung und auf das Gesetz. Die Lehre des Gesetzes muß man behalten; denn sie ist nötig dazu, daß Zucht und Friede möchte erhalten werden. Darum soll man auf das genaueste darüber halten: wie Abraham auf Ismael, Joseph auf Manasse hält etc. Denn man muß um der Verheißung willen der Gnade das Gesetz gleichwohl nicht wegwerfen; sondern es soll und muß gelehrt werden, damit Zucht und die Lehre von guten Werken auch erhalten werde, und daß wir unterrichtet werden, uns selbst zu erkennen und nach der Sünde zu demüthigen. Das ist der rechte und nothwendige Gebrauch des Gesetzes. Denn wir müssen in diesem Leben Obrigkeit und Eltern haben, die mit Belohnung und auch etwa mit Strafe über Zucht und Ehrbarkeit halten, und die das Gesetz halten, und ihr Leben christlich und wohl regieren und nach dem Gesetze richten. Diese sind gerecht, und Gott belohnt auch solche äußerliche Gerechtigkeit. Denn man soll und muß Gott ehren, und ja nicht so in der Irre gehen wie die wilden Cyclopen; welches weder Schüler noch andern Ständen in diesem Leben soll erlaubt sein, daß sie ohne Gesetz, ohne alle Zucht und Ordnung also wollten dahin leben, als ob Gott nicht gefordert hätte, daß man seinen Geboten sollte gehorsam sein. Denn wo die Gebote nicht gehalten werden, so folgen bald mancherlei Strafen darauf, es geschehe gleich durch die Obrigkeit, Schulmeister oder Eltern, oder aber sonst durch gemeine Noth und Unglück. Denn die gemeine Noth oder Schaden ist wie ein Kerker oder Gefängniß des menschlichen Geschlechts, darein sie geführt werden, wenn das Gesetz übertreten wird und die Obrigkeit mit der Strafe auch säumig ist.

82. Aber Gott selbst ist von dem Gesetz frei, soll auch demselben nicht unterworfen sein, daß er nach dem Gesetz thun müßte. Denn er ist ein Herr über das Gesetz, der dispensiren und anders thun kann, denn das Gesetz gebietet. Es ist ein Anderes um das Reich der Gnade und ein Anderes um das Reich des Gesetzes. Das Gesetz verbietet die Sünde, zeigt die Ruthe und verkündigt den Zorn Gottes und schwere Strafe denen, so da sündigen. Das ist eigentlich das rechte Amt des Gesetzes, und gehört dazu, daß die Buben, ungehorsame, muthwillige, sichere Menschen dadurch gehändigt werden sollen.

83. Aber das Reich der Gnade ist ein Reich der Barmherzigkeit, der Verzeihung, Errettung und Erlösung von Sünden und Strafen, so die Leute mit Sünden verdient haben.

84. Und ist bei den Juristen auch eine treffliche feine Lehre von der epikeia, wie es auf Griechisch heißt, das ist, von der Billigkeit, daß das strenge Recht etwas gelindert wird; darnach soll man sich in allen Regimenten, so beide die Polizei und Haushaltung belangen, fleißig richten. Es thut einem Hausvater noth, daß er in seinem Hausregiment die Stücke beide habe, nämlich, das Gesetz und die epikeia, das ist, die Billigkeit dabei.

85. Aristoteles hat auch von dieser epikeia sehr fein gelehrt im fünften Buch der Ethik, welches das allerherrlichste ist. Die Obrigkeit ist dazu gesetzt, daß sie nach dem strengen Rechte regieren und richten soll, darüber sie auch allewege halten und dasselbe handhaben soll. Dieweil aber der Fälle und Händel unendlich und unzählig viele sind, die nicht alle in Schrift und Gesetz können gefaßt werden um mancherlei Umstände willen, und Wenige sind, die da sehen und verstehen können, wo das Gesetz recht weislich zu lindern sei: darum hat Aristoteles eine feine Weise angezeigt, wie man sich darin halten soll. Wie er auch darauf zieht die Definition oder Beschreibung der Tugend, da er sagt: Virtus est habitus electivus, in mediocritate consistens, quoad nos, ratione aliqua, ut sapiens judicat, das ist: Die Tugend ist eine solche Geschicklichkeit, wenn der Mensch aus wohlbedachtem Rath der Vernunft thut, was recht ist, doch also, daß er die Mitte darin hält und nicht zu viel oder zu wenig thut, wie ein verständiger vernünftiger Mann dasjenige weiß zu treffen und davon zu urtheilen in diesem Leben. Also ist die Stärke, fortitudo, die rechte Mitte zwischen Zorn oder Frechheit und Faulheit.

86. Aber dieses Maß oder Linderung ist den Affecten oder Begierden der Menschen unterworfen; darum setzt Aristoteles dazu: Wie ein weiser verständiger Mann davon urtheilen oder richten kann. Denn solch Urtheil konnte in keinen Gesetzen begriffen werden, sondern da muß ein lebendig Gesetz sein, nämlich, die Obrigkeit, die das Herz oder Leben des Gesetzes ist, und fleißig darauf sehen soll, wo und wie das Gesetz müsse gehalten werden, und dasselbe lindern soll, wo etwa ein unmöglicher Fall eintritt.

87. Wie ein Hausvater seinem Knechte eine gewisse Weise und Ordnung vorschreibt, wie er die Hausarbeit thun soll: heißt ihn des Morgens um fünf aufstehen, darnach des Ackerbaues warten; heißt ihn sein bescheiden Theil Speise nehmen und daß er die andere Arbeit auch treulich und fleißig thue. Dies eine Gesetz aber kann nicht alle Fälle, so sich können zutragen, vorsehen, welche die vorgeschriebene Ordnung hindern und zerstören können. Derhalben tritt etwa dem Knechte plötzlich und unversehens ein Hinderniß entgegen, wenn er will zu Felde ziehen und des Ackerbaues warten: daß sein Herr krank geworden ist, oder daß nicht bequemen oder gelegenen Wetter ist, oder daß der Hausvater sonst den Knecht zu anderer Arbeit brauchen muß, davon ihm mehr gelegen ist. Da denkt und schließt ein solcher Hausvater bei sich also: Ich bin das lebendige Gesetz in meinem Hause, und habe in meiner Hand und Gewalt die epikeia oder Billigkeit, und habe Recht und Macht, das Gesetz zu mäßigen und zu lindern. Derhalben soll die Strenge des Gebotes, so ich zuvor meinem Knechte gegeben, aufheben. Da bricht er ein Loch durch das Gesetz mit der epieikeia oder Billigkeit wegen des Falles, der sich plötzlich und unversehens zugetragen hat.

\quellseite{483}

88. Also suchen auch die Juristen solches Maß oder Linderung des strengen Rechtes, nach dem die Umstände oder Fälle mancherlei sind, und deren etliche fordern und haben wollen, entweder daß man das Recht schärfe, oder aber dasselbe lindere, darum daß das Gesetz etwa unmöglich ist.

89. Darum ist der große Zwang mit den Statuten oder Regeln in den Klöstern nicht zu loben, da die Mönche ohne Unterschied oder Billigkeit allen Klosterpersonen verbieten, Fleisch zu essen, zu baden, leinene Kleider zu tragen etc. Denn es pflegen auch die Aerzte nach Gelegenheit und mit Maß etlichen Menschen diese, etlichen aber andere Speisen, Uebungen und Arznei vorzuschreiben. Und wenn sie sehen, daß ein Mönch am Leibe etwas schwach ist und daß er sich derhalben ohne Verletzung und Gefahr seiner Gesundheit des Fleischessens nicht kann enthalten, befehlen sie, daß die Strenge der Mönchsregel soll gelindert werden, damit der schwache Bruder nicht dahinstürbe. Denn es ist bei den Carthäusern damit so eine große Tyrannei geübt worden, daß sie Viele getödtet haben mit ihrem sauren und harten Fasten, die sie wohl hätten können beim Leben erhalten, etwa mit einem Hühnerbrüblein oder Stücklein Fleisch, oder aber auch mit einem reinen oder saubern Kleide, da sie daran die Aerzte, die zu solchen Kranken gefordert worden, vergeblich erinnert und sie dazu ermahnt haben, nämlich, daß man nach Beschaffenheit des Leibes die Ordnung mit Essen und Trinken allezeit ändern müßte. Das ist wahrlich eine rechte Mördergrube gewesen und kein recht Klosterleben, nämlich, einen armen kranken Menschen tödten, den sie doch mit einem Brühlein oder Suppe bei Leben und Gesundheit hätten erhalten können.

90. Ja, es ist endlich der Aberglaube so gar kräftig und groß geworden, daß etliche Phantasten die Läuse oder die Flöhe nicht haben tödten wollen. Und ich habe einen Pfaffen gesehen, der da meinte, daß er damit Gott einen Gefallen thäte, daß er derselben Ungeziefer verschonte und es nicht tödtete. Denn er reinigte seine Kleider nicht, ja, er setzte die Läuse, so von ihm kamen, wiederum in die Kappe, und zeigte diese Ursache an, warum er sich so gar unsauber hielt, nämlich, dieweil er wüßte, daß seine Eltern auch im Grabe von den Würmern gefressen und verzehrt würden. Es hat aber Gott nicht geboten, daß wir unsere Leiber tödten sollen, sondern will haben, daß wir sie ehren sollen; wiewohl sie bedürfen, daß man sie in Zucht und Zwang halte.

91. Derhalben ist das die höchste Thorheit und Unsinnigkeit des Teufels gewesen, der mit diesen scheußlichen und garstigen Werken des menschlichen Geschlechts gleichsam höhnisch gespottet hat; denn es sind viel gute fromme Leute jämmerlich in diesen Stücken gefangen gewesen. Ich bin selbst auch ein solcher Mensch gewesen, daß ich mich mit Fasten und Kasteien, mit harter schwerer Arbeit und Kleidung beinahe selbst getödtet hatte, da ich meinen Leib greulich verderbet und entkräftet habe.

92. Von diesen scheußlichen Mönchswerken weiß jetzt die Jugend nichts, läßt es sich auch nicht hart anfechten, wenn man derselben gedenkt; ja, über das wird jetzt das junge Volk über die Maßen muthwillig und ungehorsam, ganz und gar wie die wilden Thiere. Man muß den Leib allewege kasteien oder zähmen, und ihn mit Arbeit, Studiren und sonst andern gemeinen Werken dieses Lebens beladen; aber tödten soll man ihn nicht. Warte viel lieber deines Berufs, da wird dir Mühe und Arbeit genug vorfallen. Es ist nicht vonnöthen, daß du dir selbst freiwillig Uebel aufladest, oder den Tod auflegest; der wird dir zu seiner Zeit, wenn es Gott haben will, doch wohl kommen.

93. Darum ist St. Augustin billig zu loben (wie ihn denn der Bischof von Dalburg gelobt und gerühmt hat) um der Regel willen, die er geschrieben, welche anzeigt, daß er müsse ein sehr weiser, verständiger Mann gewesen sein. Denn dies sind seine Worte: Non aequaliter omnibus, quia non aequaliter valetis omnes, das ist: Man soll euch nicht alle gleich halten, oder gleiche Last auflegen, denn ihr seyd nicht alle gleich stark. Man muß die geometrische Proportion halten, das ist, den Unterschied der Personen betrachten, und nach demselben einem Jeden seinen Ort zugewiesen und also die Last und Arbeit austheilen.

94. Aber das war zumal ein tyrannisch Wort, das die Mönche pflegten zu sagen: Die soll geschehen, wie mir, es geschieht hier niemand unrecht. Das heißt nach arithmetischer Proportion handeln, da die Personen einander gleich gerechnet werden und die Last auch gleich wird: da einem Kinde so viel auferlegt wird zu tragen, wie Einem, der erwachsen ist, oder einem starken Knechte; oder aber wenn man das wollte den Weibern auferlegen, das den Männern gebührt. Daselbst sind die Personen ungleich und die Last oder Arbeit ist auch ungleich.

95. Ich habe im Kloster einen Bruder gesehen, der konnte fünf Semmeln verschlingen, da ich an Einer genug hatte. Wenn nun der Prior demselben geboten hätte, daß er sich nur an Einer Semmel sollte genügen lassen; oder wo er mir hätte geboten, daß ich viel Semmeln sollte verzehren wie jener: so hätte wahrlich der Bruder Hungers sterben müssen und ich wäre von vielen Fressen auch umgekommen.

96. Derhalben soll man erstlich auf den Unterschied der Person Acht geben; darnach aber soll man die Aemter oder Ordnung einer jeglichen Person auch betrachten. Wenn einem Carthäuser vonnöthen ist, daß er haben muß, so laß ihn haben; kann er nicht Kraut oder Fische essen, so mag er Fleisch essen, oder im Gegentheil, kann er nicht Fleisch essen, so mag er Kraut oder Fische essen. Und davon handelt Aristoteles sehr fein, nämlich, über die geometrische Proportion, und von der epieikeia oder Billigkeit. Denn dies ist ein Theil der Gnade, die bei der Obrigkeit, Hausregimente und Polizei Statt haben muß. Wenn das Weib im Hause schwach ist, da merke ich das Gesetz hinweg, bis sie wieder gesund wird. Das Gesetz soll gehalten werden, aber doch also, daß gleichwohl die Obrigkeit immer die geometrische Proportion in der Hand behalte, daß sie den Unterschied der Personen betrachte, und rechtes Maß und Billigkeit treffe. Denn Aristoteles sagt, wie droben gehört: Virtus est habitus, consistens circa mediocritatem, ut sapiens determinet: Die Tugend ist eine Fertigkeit, die in der Mitte liegt, wie dasselbe ein weiser vernünftiger Mann recht treffen und urtheilen kann.

97. Auf diese Weise sind Joseph und sein Vater auch uneinig. Joseph hat eine sehr gute Sache und hat das Gesetz nach seinem Verstande für sich; denn die Erstgeburt ist recht von Gott selbst befohlen und gehehrt. Adam hatte seinen erstgebornen Sohn Kain lieb und es stund alle seine Hoffnung und Wohlfahrt auf ihm: aber GOttes Segen zieht den Andern hervor und erwählt Abel; und verliert also Adam seinen Erstgebornen. Manasse war des Patriarchen Joseph Erstgeborner, darum sagt er: Ach mein lieber Vater, was thust du? Du kennst meine Söhne nicht, du fehlest weit. Ich habe es mit sonderlichem Fleiße also geordnet, daß Ephraim sollte zur Linken und der Andere zur Rechten stehen, und also das Gesetz von den Erstgebornen gehalten werden möchte und das göttliche Recht nicht geändert würde. Du sündigst wahrlich jetzt. Ich sollte dir ja wohl gehorchen und folgen, wie der Sohn dem Vater; aber der Gehorsam des Gesetzes ist besser und geht vor. Du bist ja ein Mensch, der fallen und irren kann. Derhalben, lieber Vater, nimm deine rechte Hand von Ephraims Haupte und lege sie auf Manasse's Haupt; denn dies hat seine rechte Ursache und ist in GOttes Gesetz gegründet. Das war wohl geredet und war eine gute Ursache. Ich würde zwar auch dasselbige gethan haben und sollte wahrlich auch über der Ordnung und Meinung des Gesetzes halten. Und wenn Joseph hätte anders gethan, so hätte er unrecht gethan; denn es hat ihm also gebühret, daß er das Gesetz hielte und dem Rechte der Erstgeburt folgete. Darum will er das Gnadenrecht nicht haben.

98. Aber Jakob antwortet ihm darauf also: Mein lieber Sohn, ich verstehe wohl, daß du das Recht der Erstgeburt sehr hart vertheidigest nach dem Gesetze, welches du willst gehalten und gehehret haben; und ich begehre selbst auch, daß es fest und unbeweglich bleiben möge. Aber es ist jetzt des Gesetzes Zeit nicht, es hat jetzt hier keine Statt; sondern hier gilt allein der göttliche Segen, welcher keinem Gesetze unterworfen ist, auch unserm Rechte oder Weisheit nicht. Derhalben strafst oder verwirfst du Josephs Meinung nicht, sondern er läßt sie im Mittel stehen und in ihrer Würde bleiben. Er hebt das Gesetz nicht auf, sondern richtet das aus, was die Verheißung in sich hält und mitbringt.

\quellseite{484}

\begin{quote}
\textbf{V. 19. Aber sein Vater weigerte sich, und sprach: Ich weiß wohl, mein Sohn, ich weiß wohl. Dieser soll auch ein Volk werden, und wird groß sein; aber sein jüngster Bruder wird größer, denn er, werden, und sein Same wird ein groß Volk werden.}
\end{quote}

99. Dies ist nun eine Bestätigung des Segens, dazu er gesetzt hat eine gewisse und deutliche Definition oder Beschreibung der Verheißung, und kommt beinahe überein mit der Antwort, so droben, 1 Mos. 25, 23, der Rebecca gegeben worden: „Der Größere wird dem Kleinern dienen“; wiewohl sie nicht ohne große Beschwerung, Mühe und Arbeit ihrem Sohne Jakob die Erstgeburt erlangt, da ihr der Vater Isaak in dem widerstanden und sein Bruder Esau ihn verfolgt hat. Und Ephraim hat zwar auch neun Söhne verloren, die schon alt geworden waren, von achtzig oder neunzig Jahren, wie Ezra davon disputirt, 1 Chron. 8, 21. 22. Denn da sie das Land Canaan vor der rechten Zeit wollten einnehmen, als wäre die Verheißung jetzt schon erfüllt gewesen, sind sie darüber von den Philistern erschlagen worden. Und der Text sagt daselbst noch dazu: „Und ihr Vater Ephraim trug lange Zeit Leide, und seine Brüder kamen ihn zu trösten.“ Da sind die Erstgebornen erschlagen worden, daß sie das Erbe nicht haben können einnehmen, so ihnen vom Vater und im Gesetz Gottes verordnet und bestimmt war; und die er am letzten gezeuget, die haben die Verheißung erlangt; von welchen darnach ist geboren worden Josua, der Sohn Nun.

100. Gott will, daß das Gesetz soll gehalten werden, aber auf ganz eigene, wunderbare Weise, darum daß wir wegen der Gerechtigkeit des Gesetzes stolz und vermessen werden. Und es kann nicht anders kommen, es muß aus solcher Gesetzgerechtigkeit die Vermessenheit und Hoffart folgen, wie das an geringen und gemeinen Exempeln zu sehen. Denn wo willst du mir eine fromme züchtige Matrone oder eine junge schöne Jungfrau geben oder zeigen, die solcher Gaben halben nicht etwas stolz und aufgeblasen wäre? Der Muthwille aber und Hoffart derjenigen, so die Regimente in Landen und Städten verwalten, ist über die Maßen groß. Denn die Gerechtigkeit des Gesetzes ist ein solch hoffärtig und vermessen Ding; darum macht sie auch Gott zu Schanden, und wenn er sieht, daß du über die Maßen stolz wirst wegen der Gesetzgerechtigkeit, so läßt er dieselbige fahren und handelt nach der Gnadengerechtigkeit.

101. Also haben die Juden der Gerechtigkeit des Gesetzes nachgestanden, sind aber dadurch zur Gerechtigkeit nicht gekommen, und die Erstgebornen sind erschlagen worden, damit Gott also die Hoffärtigen und Stolzen demüthigen möchte. Denn er will haben, daß wir die rechte Daniels am 9. Cap. V. 18. von ganzem Herzen und auch mündlich thun sollen und sagen: „Wir liegen vor dir mit unserm Gebet, nicht auf unsere Gerechtigkeit, sondern auf deine große Barmherzigkeit“; desgleichen in der Epistel Pauli an Titus am 3. Cap. V. 4. 5.: „Da aber erschien die Freundlichkeit und Leutseligkeit Gottes, unsers Heilandes, nicht um der Werke willen der Gerechtigkeit, die wir gethan hatten, sondern nach seiner Barmherzigkeit, machte er uns selig“ etc. Die Werke der Gerechtigkeit sind, die wir gethan haben; Gott aber sieht dieselbigen nicht an. Die Gnade herrscht über das Gesetz, welches den Menschen nicht vollkommen, sondern nur stolz macht. Darum wendet Gott seine Barmherzigkeit auf die armen Sünder und zieht die Gnadengerechtigkeit der Gesetzgerechtigkeit vor. Da wird dann die Welt toll und thöricht, zürnt und schreit: Soll man denn nun keine guten Werke thun? Ist das Gesetz nicht gut? Ja, es ist gut, und wir vermahnen immer dazu, daß du es thun und halten sollst; aber doch also, daß du sagest, wie Christus lehrt Luc. 17, 10.: Wenn wir alles gethan haben, sind wir doch unnütze Knechte.

V. 20. Also segnete er sie des Tages, und sprach: Wer in Israel will jemand segnen, der sage: Gott setze dich wie Ephraim und Manasse.

\begin{quote}
\textbf{V. 20. Also segnete er sie des Tages, und sprach: Wer in Israel will jemand segnen, der sage: Gott setze dich wie Ephraim und Manasse.}
\end{quote}

102. Es ist darum nichts gelegen, ob du diese Rede activisch oder passivisch, wie man nach der Grammatik redet, verstehst: Israel wird segnen, oder wird gesegnet werden in dir. Denn der Vater will so viel sagen: Du sollst ein Exempel sein des Segens, daß man sage: Wie Ephraim und Manasse gesegnet sind, also soll Gott dich auch segnen. Also liest man in der Historie Ruth, daß die Weiber und alles Volk Boas und Ruth Glück wünschen und sagen:

\begin{quote}
„Der Herr mache das Weib, das in dein Haus kommt, wie Rahel und Lea, die beide das Haus Israel gebauet haben“, Ruth. 4, 11.
\end{quote}

Da werden Jacobs Weiber zum Exempel des Segens vorgestellt.

103. Da siehst du aber nun, wie große und wunderbare Männer die heiligen Patriarchen gemein sind, wie gewiß und unausgeweifelt sie frei heraus geredet haben von dem zukünftigen Segen. Denn jetzt redet Jakob gleich wie Gott; denn er thut auch dies alles in der Person Gottes: er gibt seinen Söhnen den Segen und das Erbe des gelobten Landes, welches er selbst noch nicht hatte denn nur im Glauben. Darum habe ich gesagt: Was mir in der Verheißung haben, dasselbige ist so gewiß und besteht so fest, daß auch weder Teufel, noch Tod, noch Hölle solches mit Gewalt wegreißen könnten. Was mir aber in Händen haben oder in unserm Kasten oder Beutel liegt, ist so gar gewiß und wohl verwahret nicht: es können es die Diebe stehlen oder sonst Andere rauben; ja, es können es auch wohl die Mäuse und Motten fressen. Aber was uns in göttlichen Verheißungen angeboten wird, dasselbige kann uns niemand nehmen. Wie Gott den Frommen und Gottseligen verheißet hat im 37. Psalm V. 19.: „In der Theurung werden sie genug haben.“ Daher hoffen sie gewißlich ihre Nahrung und Leibes Nothdurft zu erlangen; sie wissen, daß ihnen Küche und Keller auf das beste bestellt ist, wiewohl daheim in ihrem Haus nichts Sonderliches zugerichtet oder zubereitet ist; ja, sie haben oftmals wohl nicht einen Heller oder Scheffel in ihrer Gewalt.

104. Woher haben sie es denn? Antwort: Sie haben es in dem Worte, da Gott sagt: „In der Theurung werden sie genug haben.“ Wer da glaubt, der hat es ohne allen Zweifel, und wiewohl er der ärmste Mensch ist, hat er doch alles, was ihm vonnöthen ist. Wie Jakob hier dessen gar gewiß ist, daß er werde das gelobte Land einnehmen, und gibt seinen Söhnen denselbigen Segen nicht anders, denn als führte er sie jetzt gegenwärtig hinein, da sie es mit der That möchten einnehmen: so er ihnen doch nichts hat gegeben denn nur diese bloßen Worte. Doch hat er das Land Canaan gewisser inne gehabt, denn dazumal, da es seine Nachkommen besessen haben. Denn da sie dasselbe nicht im Glauben, sondern mit der That besessen haben, da wurden sie hoffärtig und stolz wegen der Gerechtigkeit des Gesetzes und Gottesdienstes und darum haben sie es auch nicht behalten können. Der König von Assyrien und Babylon ist gekommen und hat sie heraus getrieben. Es war ihnen die bloße Verheißung gegeben worden ohne Gesetz und ohne Verdienst; da sie aber das Land Canaan eingenommen, ist der Bund des Gesetzes dazu genommen mit der Bedingung 5 Mos. 8, 19. 20.: „Wirst du das Gesetz nicht halten, so sollst du verderbet werden und umkommen, eben wie die Heiden, die der Herr umbringet vor eurem Angesicht.“

105. Sie haben es ohne Verdienst empfangen und ward ihnen ohne alle Bedingung gegeben, aus Gnaden; aber sie haben es nicht also ohne Bedingung besessen. Derhalben, dieweil das Gesetz nicht gehalten wird, kann auch das Recht der Besitzer nicht vertheidigt werden, sondern ist ein löchrichten Beutel. Aber die Gnade ist ein fester gewisser Grund. Hier sollte man nun hoch heben und rühmen die große Herrlichkeit und Reichthum der Gnaden und Wohlthaten Christi, welche uns auch gegeben wird mit dem Segen der Väter, wo nicht leiblich, doch im Geist. Aber weiter unten findet sich dazu bessere Gelegenheit.

\section[Dritter Theil]{Dritter Theil. Wie Jakob, da er Josephs Söhne segnet, den Ephraim dem Manasse vorgezogen}

\begin{quote}
\textbf{V. 20. Und (Jakob) setzte also Ephraim Manasse vor.}
\end{quote}

106. Dies Stück hat Mose dazu gesetzt, daß große und wunderbare Werk damit zu rühmen, dadurch die Erstgeburt von Manasse ist weggenommen worden, dem sie doch nach gemeinem Rechte gebühret hätte, und auf Ephraim, seinen jüngsten Bruder, gelegt oder gewendet ist, dem sie nach keinem Rechte zugehört. Denn das ist fast wie ein Wunder und einem seltsamen unnätürlichen Dinge gleich, nach allen Rechten und Gesetzen; sintemal die Erstgeburt nach dem Gesetze der Natur und nach geschriebenem Landrechte und auch nach göttlichem Gesetze dem Aeltesten ist zugeeignet worden. Und man hat sonderlich unter demselben Volke schmerzlich wollen zulassen, daß solch Gesetz sollte geändert werden, dieweil beide die königliche und auch priesterliche Würde dem erstgebornen Sohn gegeben ward. Und darum hast auch Esau seinen Bruder Jakob so heftig, ward ihm bitter feind, und verfolgte ihn, als der ihn solcher großen Ehre und des Erbes beraubt hatte.

\quellseite{485}

107. Und wird auch noch jetzt dieselbe Ordnung im Römischen Reiche gehalten und auch sonst in allen andern Reichen. Der Dauphin in Frankreich folgt seinem Vater allernächst im Reiche oder Regimente. Der Herzog zu Sachsen, welcher der Erstgeborne oder Aelteste ist, wird Churfürst, und wenn sie das Erbe theilen, nimmt er die Chur und die Güter, so dazu gehören, an sich vor seinen andern Brüdern, und theilt dieselben mit ihnen nicht. Derhalben ist es jetzt auch noch eine große Herrlichkeit mit der Erstgeburt. Und Joseph verwundert sich nicht vergeblich so sehr über diese Veränderung; denn Ephraim wird aus Gnaden der Erstgeborne, der nach dem Gesetze der Jüngste und Letzte war.

108. Der Mann, so die Epistel an die Hebräer geschrieben hat, er sei wer er wolle, Paulus, oder, wie ich es dafür halte, Apollo, zieht diesen Text fein künstlich an, da er sagt Cap. 11, 21.: „Durch den Glauben segnete Jakob, da er starb, beide Söhne Josephs.“ Er gedenkt der Segen nicht, so im 49. Capitel beschrieben werden; aber dieses Segens gedenkt er, da Jakob Ephraim und Manasse gesegnet. Denn dieser Segen ist auch viel wunderbarer, denn die sind, so hernach folgen werden. Denn Ephraim nimmt die Erstgeburt hin ohne alles Verdienst; Manasse aber wird derselben beraubet ohne alle seine Schuld. Warum thut aber Jakob solches wider das Recht des Gesetzes und wider den gemeinen Gebrauch und Gewohnheit? Es ist wahrlich viel genommen: es würde kein König oder Fürst solch Unrecht leiden, mit was für großen vortrefflichen Gaben oder Tugenden sein jüngster Bruder auch möchte begabt oder geziert sein. Antwort: Jakob hat solches gethan im Glauben und durch Eingebung des Heiligen Geistes; denn die Verheißung, daß er wachsen sollte und gemehret werden, reimt sich fein mit dem Namen Ephraim. Darum muß man die rechte Ursache dieser That aus dem Glauben und Verheißung nehmen, und nicht aus dem Gesetz, noch aus den Rechten, noch auch aus der Natur.

109. Nun lasset uns aber die Ehre und Herrlichkeit des Segens der Väter mit der Herrlichkeit vergleichen, so wir im Neuen Testamente haben. Denn daran ist ja kein Zweifel, daß das Reich Christi höher und herrlicher sei, denn das Reich Mose und der Väter; wie der 145. Psalm V. 10–12. sagt: „Es sollen dir danken, HErr, alle deine Werke, und deine Heiligen dich loben, und die Ehre deines Königreichs rühmen, und von deiner Gewalt reden, daß den Menschenkindern deine Gewalt kund werde, und die ehrliche Pracht deines Königreichs.“ Und es ist doch eine große Ehre gewesen, daß die Väter haben segnen können, das ist, beides leibliche und geistliche Güter verheißen. Wiewohl der Segen nicht vornehmlich auf Fleisch und Blut geht, und geht doch auch den Leib an. Und insofern hat er die Bedingung, das Gesetz zu halten, an sich hangen: aber sofern man ihn auf den Geist zieht, hat er die Bedingung nicht, sondern kommt aus Gnaden.

110. Darum wird Ephraim bei Hosea und andern Propheten hart gestraft und beschuldiget um der greulichen Sünden willen Jerobeams, des ersten Königs aus demselben Stamme, „der da gesündiget hat, und Israel hat sündigen gemacht“, wie geschrieben stehet 1 Kön. 14, 15. 16., da er den Kälbern geopfert; welche Abgötterei alle Propheten verabscheuen und strafen, und ist auch eben darum das Volk endlich zerstreuet und nach Assyrien weggeführet worden, da sie die Bedingung, so an das Gesetz geknüpft war, nicht gehalten haben. Denn diese Drohung steht bei den vorigen: „Der HErr wird Israel schlagen, gleichwie das Rohr im Wasser beweget wird, und wird Israel ausreißen von diesem guten Lande, das er ihren Vätern gegeben hat, und wird sie zerstreuen über das Wasser; darum, daß sie ihnen Haine gemacht haben, den HErrn zu erzürnen.“

111. Und dieselbe Drohung wird hernach vom Propheten Hosea oftmals wiederholt. Aber wiewohl er hart mit ihnen umgeht und allerlei Plage und Unglück droht und daß sie gar umkommen werden, mengt er doch bisweilen Trost mit unter aus den geistlichen Verheißungen; wie der schöne tröstliche Text anzeigt Hos. 11, 8. 9.: „Was soll ich aus dir machen, Ephraim? Soll ich dich schützen, Israel? Soll ich nicht billig ein Adama aus dir machen, und dich wie Zeboiim zurichten? Aber mein Herz ist anderes Sinnes, meine Barmherzigkeit ist zu brünstig, daß ich nicht thun will nach meinem grimmigen Zorn, noch mich kehren, Ephraim gar zu verderben; denn ich bin GOtt, und nicht ein Mensch.“

112. Auf solche Weise mengt er die Verheißung des Heiligen Geistes und ewigen Lebens mit der leiblichen Verheißung, und hat ihnen zuvor verkündigt, daß sie nach derselben leiblichen Verheißung werden zerstört und verderbet werden, darum daß sie die Bedingung, das Gesetz zu halten, in Vergessenheit gestellt und versäumt hatten, und von GOtt, ihrem Schöpfer und der ihnen die Verheißung gegeben hatte, abgewichen waren. Aber dieweil eben in demselben leiblichen Segen auch der geistliche und ewige mit begriffen war: Ich bin der HErr, dein GOtt, und der verheißene Christus, darum sagt er: Ich kann Ephraim nicht gar verwerfen oder verderben, sondern ich muß mich sein erbarmen. Ich bin wohl greulich zornig auf die abgöttischen Leute, aber ich bin GOtt und nicht ein Mensch; derhalben, um meiner Verheißung willen will ich Israel noch helfen und erhalten.

113. Und also macht er einen Unterschied zwischen den gottlosen abgöttischen Juden, und denen, so sich GOtt hat lassen übrig bleiben, die Baal und die Kälber nicht angebetet und ihre Hand nicht geküßt haben. Welches man auch in der Historie der Könige sehen mag, als, da zu Elia gesagt wird: „Ich habe mir lassen übrig bleiben sieben tausend in Israel, nämlich alle Kniee, die sich nicht gebeuget haben vor Baal“, 1 Kön. 19, 18.

114. Wiewohl wir nun jetzt zur Zeit keine vollkommene zeitliche Verheißung zum leiblichen Reich haben, mangelt es doch an vergleichlichen Verheißungen nicht, nämlich, daß die Kirche auch in diesem Leben soll erhalten und ernähret werden. Aber die geistliche Verheißung im Neuen Testament ist viel herrlicher; wie Paulus sagt 1 Tim. 4, 8.: „Die Gottseligkeit ist zu allen Dingen nütz, und hat die Verheißung dieses und des zukünftigen Lebens.“ Es ist uns alles verheißen, was wir zu diesem Leben bedürfen, dasselbe damit zu erhalten; wie der 8. Psalm V. 7. sagt: „Alles hast du unter seine Füße gethan“; und Luc. 10, 6. 7. 8. spricht Christus zu seinen Aposteln: „Wo ihr in ein Haus kommet, in demselbigen bleibet, esset und trinket, was sie haben. Denn ein Arbeiter ist seines Lohnes werth. Und wo ihr in eine Stadt kommt, und sie euch aufnehmen; da esset, was euch wird vorgetragen.“ Und 1 Cor. 9, 14.: „Also hat der HErr befohlen, daß, die das Evangelium verkündigen, sollen sich vom Evangelio nähren.“ Solche Verheißungen haben wir; und wiewohl sie gering sind, haben wir doch daran genug zu nothdürftiger Unterhaltung dieses Lebens.

115. Aber die Verheißungen von den geistlichen Gütern sind sehr herrlich, und wir sollen dieselben auch groß machen und hoch erheben, nämlich, daß die Kirche bleiben soll, daß sie Herberge, Essen, Trinken und Schutz haben soll, wiewohl etliche Verfolgung leiden und getödtet werden; denn die alte Kirche ist auch nicht ohne Marter und Kreuz gewesen. Ob wir aber keine Könige sind, auch kein zeitlich Königreich haben wie die Juden, behalten wir doch das geistliche Königreich, darin wir leben in Christo Jesu.

116. Diese Herrlichkeit soll man sehr groß machen und hoch erheben, und dasselbe so viel mehr, so viel schrecklicher sie vom Papst verdunkelt und unterdrückt worden ist. Denn da hat man in so viel Kirchen, Klöstern und hohen Schulen nie ein Wort oder einige Lehre gehört von GOttes Wort, wie jetziger Zeit, da uns das Licht des Evangeliums wieder gegeben worden ist. Wenn wir hören das Wort nennen, verstehen wir dadurch die Verheißung und das Predigtamt. Aber wenn dazumal unter dem Papstthum des „Wortes“ gedacht ward, wandte sich mein Gemüth und Gedanken nur auf den Donat, darin das Verbum also beschrieben wird: Verbum est pars orationis cum tempore et persona etc.: Das Verbum, d. h. Wort, Zeitwort, ist ein Redetheil mit Zeit und Person u. s. w.; denn es ward daselbst gar nichts gelehrt von den Verheißungen GOttes.

117. Und was soll ich viel sagen? Im Anfang, da das Evangelium erst wieder an den Tag kam, habe ich einen Mönch hören sagen, da er den päbstischen Aberglauben abgelegt, und nun die reine Lehre gehört und wohl verstanden hat: Hilf, lieber GOtt, habe ich mein ganzes Lebenlang nie etwas gehört von den Verheißungen! Und dieses war ihm von Herzen lieb, daß er nun dasselbe Wort „Verheißung“ hören und verstehen möchte.

\quellseite{486}

118. Derhalben soll unsere Theologie und Neues Testament vornehmlich diesen Theil der himmlischen Lehre treiben, wiewohl man das Gesetz auch lehren muß; es besteht aber das Reich Gottes darin nicht, sondern im Wort der Verheißung. Man pflegt jetziger Zeit zu sagen: Er hat das Wort des Evangeliums oder das Predigtamt lieb: aber in des Papstes Decreten oder in seinen Canones wirst du von dem Wort nicht eine Sylbe finden. Sie schreien allein davon, wie man die Sünde beichten, bereuen und dafür genug thun müsse; item, wie man dem Papst gehorsam sein und der Mönche Regel halten solle: aber von der Verheißung des Evangeliums schweigen sie gar tödtlich. Derhalben ist das Reich des Papstes eine greuliche Verwüstung der Kirche gewesen und ist auch jetzt noch bei dem Papste und Cardinalen „Verheißung“ ein unerhörtes Wort.

119. Unser Reich aber des Neuen Testaments soll zwar auch die Lehre des Gesetzes treiben, um äußerliche Zucht und Disciplin, auch Gehorsam in weltlicher Polizei und die Ehre der Obrigkeit und Eltern damit zu erhalten: aber es besteht das Reich Gottes darin nicht, sondern im Wort, das ist, in der Verheißung, welche das rechte und wahrhaftige Amt des Neuen Testaments ist. Es hat aber der Papst die Lehre des Glaubens verdunkelt und gar vertilgt, und hat nur lauter heuchlerische und seltsame Wunder menschlicher Satzungen hereingeführt; außer, daß Gott aus unermesslicher Güte in Etlichen noch ein klein Licht erhalten hat und bleiben lassen.

120. Wie ich noch wohl gedachte, daß mein Vater die Mönche und allen priesterlichen und päpstischen Plunder verachtet, von welchen doch allein in den Decretalen und ihren Schreibern gehandelt wird. Derhalben, da ich zu Erfurt erstlich in das Mönchskloster des Augustinerordens gegangen war, hat solches meinen Vater sehr verdrossen; und darnach, da er auf meine erste Messe geladen ward, wie das der Gebrauch war, und als bei einem Mittagsmahl vornehme Männer des Augustinerordens dasselbe Mönchsleben lobten und sagten, es nähme sie Wunder, warum er darüber so unwillig wäre, daß ich mich in denselben Orden begeben hätte: da antwortete er kürzlich also: Ei, lieben Herren, wisset auch, daß geschrieben steht: Du sollst Vater und Mutter ehren. Oder wißt ihr das Gebot Gottes nicht, daß man die Eltern ehren soll?

121. Darum ist das Reich des Papstes billig zu verabscheuen und man soll davor einen Greuel haben, darum daß es den allerbesten und edelsten Theil der christlichen Lehre vertilgt und anstatt der göttlichen Verheißungen den Leuten vorgehalten hat ungereimte schändliche Gesetze der Decretalen, Klostergelübde und Wallfahrten.

122. Wir haben jetzt zur Zeit aus unermesslicher Gnade Gottes die großmächtige Ehre Christi, wie das aus unsern Predigten und dem ganzen Amt zu sehen ist. Siehe an die Taufe, das Abendmahl des Herrn, die Absolution, das Evangelium. Diese große Gaben des Heiligen Geistes sollen von jedermann gerühmt und gepreiset werden, und man soll an denselbigen Gaben Gott erkennen und loben, der solche Macht den Menschen gegeben hat. Denn ich bin, der Patriarch Abraham, Isaak oder Jakob nicht, siehe aber, was ich thue: Ich nehme ein junges Kindlein, und wenn ich dasselbe taufe, errett ich es vom Tode, vom Teufel, von der Sünde und versetze es aus dem Reich der Finsterniß in das Reich des Lichts. Das thut der Pfarrherr oder Kirchendiener und im Fall der Noth ein jeglicher Christ.

123. Derhalben sollen wir uns solcher großen Barmherzigkeit und Güte Gottes billig verwundern und dieselbe rühmen, daß er uns segnet mit geistlichem und ewigem Segen, also daß der Teufel fliehen und weichen muß, wenn er sieht, daß das Kindlein getauft wird im Namen des Vaters und des Sohnes und des Heiligen Geistes. Denn daselbst zerstöre ich nicht des Türken, des Kaisers Augustus, auch nicht des Königes der Perser Reich, sondern das Reich des Fürsten und Gottes dieser Welt.

124. Auf dieselbe Weise, wenn ein armer Sünder kommt, der seine Sünde bekennt und sich darüber beklagt, sage ich auch zu ihm: Ich absolviere dich von deinen Sünden im Namen des Vaters und des Sohnes und des Heiligen Geistes. Denn daselbst ziehe ich ihn aus der tiefen Hölle heraus, der erschrocken, verloren, verzagt und verschlungen war, und in seinen Sünden ewiglich hätte sterben und verderben müssen. Also, im Abendmahl des Herrn esse ich den Leib und trinke das Blut Christi, des Sohnes Gottes, so für mich gegeben und vergossen zur Vergebung der Sünden, das ist, zum ewigen Leben.

125. Darum ist diese Ehre viel größer und herrlicher, denn die Ehre der Väter gewesen ist. Das aber mangelt uns allein, daß wir die Augen, Ohren und Herzen nicht aufthun, und diese herrlichen Gaben nicht nach ihrer Würde und Herrlichkeit hoch genug achten und halten. Das ist ja an sich selbst wahr und ist ein gewisses Ding: Wer da getauft und absolvirt wird durch die Hand des Pfarrherrn und im Fall der Noth eines jeglichen christlichen Bruders, der ist ohne allen Zweifel heilig und selig. Da ist das Reich des Teufels, der Tod und die Sünde zerstört. Lieber, wodurch? Durch den Mund eines armen Menschen, des Priesters, oder sonst eines Bruders. Dies sind wahrlich große und unermessliche Wunder, und uns ist von Gott eine solche Gewalt gegeben, die gar ungleich größer ist, denn der Väter gewesen, welche doch sehr hoch und groß war, weil Christus zu uns gesagt, Matth. 28, 20.: „Ich bin bei euch bis an der Welt Ende“; item Joh. 14, 28.: „Ich gehe hin, und komme wieder zu euch“; desgleichen V. 23.: „Wir werden zu ihm kommen, und Wohnung bei ihm machen.“

126. Die Kirche soll meine Burg, mein Schloß, meine Kammer sein; wie Jes. 31, 9. steht: „Der Herr hat es gesagt, der zu Sion Feuer, und zu Jerusalem einen Herd hat.“ Derhalben, wenn wir taufen, so tauft Christus selbst durch den Mund und Hand des Dieners.

127. Sind das nicht wahrhaftig die großen Thaten Gottes, davon die Apostel redeten mit mancherlei Zungen? Ap. Gesch. 2, 4. 11. Es ist wahrlich ein wunderbarlich Ding, daß ein Kirchendiener oder jeglicher Bruder ein Diener ist des Reichs Gottes und ewigen Lebens, der Vergebung der Sünden und Zerstörung der Hölle, ja, endlich der Oeffnung des Himmels und des Reichs Gottes. Darum sind wir Kinder des Segens, nicht allein passivisch, daß wir selbst gesegnet sind, sondern daß wir auch Andere segnen. Und also ist die Verheißung, so dem Abraham gegeben, erfüllt worden: „Durch deinen Samen sollen alle Völker gesegnet werden“, 1 Mos. 22, 18.

128. Gott verheißt uns zwar nicht großen Reichthum oder leiblichen Ueberfluß, sondern Nahrung und Kleider; und es ist doch eine große herrliche Verheißung, da Christus sagt Matth. 6, 33.: „Trachtet am ersten nach dem Reich Gottes, so wird euch das andere alles zufallen“, das ist: Errettet zum ersten die armen Sünder aus der Gewalt der Sünden, der Hölle, des Todes und des Teufels und setzt sie in mein Reich, alsdann will ich auch den Leib mit Essen und Trinken versorgen. Wie aber oder womit soll ich es thun? Darauf antwortet Christus und sagt Matth. 28, 19.: „Gehet hin, und taufet sie im Namen des Vaters“, item V. 20.: „Lehret sie halten alles, was ich euch befohlen habe“: ihr sollt sie nicht in eurem, sondern in meinem Namen taufen.

129. Lieber, wer kann doch diese großen reichen Wohlthaten und Wunderwerke mit Gedanken erreichen, oder mit Worten ausreden? Warum schreien wir denn nicht mit dem Apostel Paulus 2 Cor. 9, 15.: „Gott sei Dank für seine unaussprechliche Gabe“? Denn ist das nicht ein wunderbarlich Ding, daß ein Mensch, der mir allenthalben gleich ist, mir das ewige Leben gibt, nimmt von mir hinweg den Tod, die Sünde, Verdammniß, die Tyrannei des Teufels, das Gift der Sünde, so noch in unserm Fleisch übrig ist? Wodurch thut er aber das? Antwort: Durch Gottes Befehl: „Gehet hin und taufet sie“ etc., und durch göttliche Verheißung.

130. Ja, sagst du, ich sehe aber das Gut an sich selbst nicht? Antwort: Wenn du glaubst, so hast du es; willst du aber nicht glauben, so wirst du auch nichts haben. Denn also lautet die Verheißung: „Wer da glaubet und getauft wird, der wird selig werden“, Marc. 16, 16. Da ist das Sacrament des Altars, da ist die Taufe, da ist das Wort in der Predigt: so viel aber hast du davon, so viel du glaubst: wenn du glaubst, so hast du Alles und vermagst Alles.

131. Ich wollte aber dasselbige gern sehen, sagst du. Antwort: Du wirst es nicht sehen; sondern wie ich gesagt habe, was wir im Wort haben, ist gewisser und beständiger, denn das wir in der That haben; denn dasselbe kann uns genommen werden und wird uns endlich auch alles durch den Tod genommen werden; darum ist es ungewiß. Wenn ich aber glaube, so ist es unmöglich, daß ich sollte Hungers sterben; und ob ich schon nicht Alles so überflüssig und gar reichlich habe, wie die zeitigen und reichen Wänste haben, wird mir doch gleichwohl daran nichts mangeln, was mir zu diesem leiblichen Leben wird vonnöthen sein; doch sofern zum allerersten der Name GOttes geheiliget werde 2c.

\quellseite{487}

132. Glaubst du aber nicht, so hast du auch nichts; durch den Glauben aber und die Verheißung hast du schon jetzt das Reich GOttes; die Gemeinschaft der Heiligen ist dir gewiß und fest genug, nicht anders, denn als ob du schon jetzt im Himmel wärest. Denn die Kinder Israel haben diese Verheißung auch nicht erlangt ohne Glauben. Jakob gibt ihnen das Land, aber wie gibt er es ihnen? Er gibt ihnen das Gut mit der That nicht und gibt es ihnen doch. Er übergibt es ihnen nicht baar oder sichtbarlich, wie Einem hundert oder tausend Gulden in die Hände gegeben werden, sondern in der Verheißung gibt er es ihnen; wie es auch Abraham, Isaak, Israel und Joseph empfangen haben. Aber wie ich nun oftmals gesagt habe, sie haben es gewisser in der Verheißung gehabt, denn dazumal, da sie mit der That und sichtbarlich darin gewohnt haben mit angehängter Bedingung des Gesetzes, so ihnen zu halten geboten war.

133. Derhalben sollen wir Fleiß anwenden, daß wir ja die Verheißung des Wortes groß machen, und wie wir immer können oder mögen, hoch erheben. Dazu sollen wir auch GOtt bitten, daß er uns den Glauben mehren und dabei erhalten wolle. Denn wo wir nicht glauben, werden wir auch nichts empfangen. Wie die ungläubigen Israeliten nicht in das gelobte Land gekommen sind, sondern sind dar nieder geschlagen und gestorben in der Wüste, Zwei allein ausgenommen, Josua und Caleb.

134. Also wir auch, dieweil wir getauft sind, wissen wir, daß wir GOttes Kinder und Bürger mit den Heiligen im Reiche GOttes sind. Die Sache ist gewiß. Wer hat dir aber diese Vergewisserung und den Trost gegeben? Antwort: Das hat der Diener des Worts oder Priester gethan. Lieber, was lehrt doch denn derselbe arme Mensch, der so verachtet, der gar ein armer Bettler ist und kaum einen Strich bezahlen kann, daß er nicht wohl könnte ärmer sein? wie sollte dir der das Himmelreich geben? Denn also pflegen gottlose, ungläubige Leute zu reden, wie vom Volk Israel gesagt wird im 78. Psalm V. 19.: „Sie redeten wider GOtt, und sprachen: Ja, GOtt sollte wohl können einen Tisch bereiten in der Wüste!“ wie kann er Brod geben, und seinem Volk Fleisch verschaffen? Und die Capernaïten sagten auch Joh. 6, 52.: „Wie kann dieser uns sein Fleisch zu essen geben?“ Aber derselbige Psalm erzählt auch die Strafen der Ungläubigen, da er V. 33. sagt: „Darum ließ er sie dahin sterben, daß sie nichts erlangeten, und mußten ihr Lebenlang geplaget sein.“

135. Dafür soll man es nun gewißlich und fest halten, daß ein Pfarrherr oder Kirchendiener Macht hat, den Himmel aufzuthun denen, die sich wollen taufen, absolvieren, und mit GOttes Wort lehren und berichten lassen, darum daß er Gewalt und Befehl von GOtt hat, und kann und soll auch absolvieren und Sünde vergeben; wie ihnen denn Christus ausdrücklich diese Gewalt hat gegeben Joh. 20, 23., da er sagt: „Welchen ihr die Sünde erlassest, denen sind sie erlassen“ 2c.

136. Dies wird aber im Papstthum also nicht gelehrt, sondern wird verdunkelt; ja, das Kind des Verderbens hat diese großen Gaben und Wohlthaten, soviel an ihm gewesen, gar unterdrückt und vertilgt. Denn sie haben also gelehrt: Es ist nicht genug, daß man getauft sei und an Christum glaube, sondern du sollst mir über das auch meine Füße küssen. Darum ist man dem Antichrist billig feind, der uns die Herrlichkeit der Verheißung, des Glaubens und Himmelreichs, und dazu des ewigen Lebens weggenommen hat, und heißt mich dieweil Gehorsam leisten einer erdichteten Person und Larve des Menschen der Sünden. Er weiset mich schlecht dem Teufel in Hintern. Das ist ja eine erschreckliche Strafe der Undankbarkeit gewesen.

137. Derhalben sollen wir uns befleißigen, daß wir ja GOtt dankbar seien für die große Herrlichkeit und Ehre des Reichs Christi, so durch das Licht und Lehre des Worts von GOtt selbst wiederum gereiniget worden ist. Denn es werden uns ja jetzt nicht solche Gesetze vorgehalten, wie in des Papsts Decretalen stehen, sondern die Lehre von Vergebung der Sünden. Der Papst aber hat nur der armen Kirche seine Menschensatzungen mit Gewalt aufgedrungen, hat weder Gesetz noch Evangelium gelehrt. Ja, wenn er das Gesetz gelehrt hätte, so möchte man ihn etlichermaßen geduldet haben; aber der rechte Verstand der Zehen Gebote war auch gar aus der Kirche hinweg genommen. Denn siehe Petrus Lombardus an und die andern Scholastiker auch, wie gar kalt sie mit den Zehen Geboten umgehen; ja, sie achten derselben gar nicht und gehen daran vorüber. Darnach haben die Canonisten nicht einmal gedacht an die Auslegung der Zehen Gebote. Und doch haben beide die Welt gefüllt mit vielen großen Büchern, darin nichts steht, denn lauter Finsterniß und Irrthum von den Mönchsregeln und Gelübden, von ihren Gürteln, Kappen und Fasten, da haben sie fast die ganze Welt und vornehmlich die einfältigen unverständigen Leute bewegt, sich desselbigen Narrenwerkes hoch zu verwundern und es zu loben.

138. Darum ist es eine große erbärmliche Blindheit gewesen, und ist billig und recht, daß man nur alle Stifte und Klöster gar umkehre und zerstöre; es wäre denn, daß man derselben etliche wollte stehen lassen zum Gedächtniß der Sünden und Greuel, die darin ihre Herberge gehabt und getrieben worden sind; wie vor Zeiten im Alten Testament die Häuser Baal, Moloch und Thopheth gewesen sind.

139. Derhalben, wenn wir gedenken an die große und wunderbare Ehre der Väter, da sie haben glauben können, daß beides das gegenwärtige und zukünftige Reich, oder die leibliche und geistliche Verheißung an einander hangen, da die leibliche wie eine Decke war, darin die Verheißung von Christo auch mit eingewickelt war: sollen wir vielmehr unsere Taufe, Abendmahl, Schlüssel und Predigtamt hoch rühmen, also daß wir uns daran gewöhnen, und das Wort und die Sacramente ehren und lieb haben lernen, auch davon viel halten und sie hoch rühmen, ehrlich davon reden; welches wahrlich eine sehr herrliche Gewalt ist, die den Himmel aufthut, die Hölle zuschließt, Sünde und Tod wegnimmt, und das Leben gibt. Aber dies alles geschieht nur im Glauben. Denn wir hoffen und warten noch darauf, und sehen die Güter noch nicht vor Augen; wie Paulus sagt Röm. 8, 24.: „Die Hoffnung aber, die man siehet, ist nicht Hoffnung; denn, wie kann man des hoffen, das man siehet?“

140. Auf diese Weise muß man die Verheißungen des Neuen Testaments gegen die halten, so die Väter im Alten Testament gehabt. Wiewohl jene auch in diesen mit begriffen waren, wurden sie doch dazumal noch nicht so öffentlich ausgebreitet über die ganze Welt, wie jetzt geschieht im Neuen Testament. Denn Malachias hatte das jüdische Volk die Verheißung allein, jetzt aber hat nun die ganze Welt die Verheißung, das Wort und die großen Thaten GOttes, und wird dazu in der Zeit der Noth einem jeglichen Bruder das Predigtamt befohlen und vertraut.

141. Wenn ich aber lehre oder predige, so lehre nicht ich und tröste die Leute, sondern Christus, der in uns wohnet. Darum glaube nicht mir, sondern Christo, der da tauft, tröstet, die Sacramente reicht durch mich; wie er verheißen hat, da er sagt Matth. 28, 20.: „Siehe, ich bin bei euch, bis an der Welt Ende.“ Denn das ist die Gewalt und Ehre des Königreichs Christi, davon der 145. Psalm V. 11. sagt.

142. Derhalben wollte ich wünschen, daß sich die Gottseeligen von Jugend auf daran gewöhnen wollten, daß sie auf das prächtigste und herrlichste redeten von unserer Erstgeburt, von unserm Reiche und dem lieben Predigtamt, und daß wir ja GOtt dafür dankbar seien, ja, daß wir diese Gaben mit gebührlicher Ehrerbietung des Herzens und Geberden des ganzen Leibes ehren; denn es sind göttliche Dinge, und solche Gaben, die uns das ewige Reich geben.

143. Das habe ich an diesem Ort sagen wollen, dieweil man von den Verheißungen der Väter etwas hat sagen müssen, daß ich damit euch erinnern möchte der großen Güter, so wir jetzt haben, und der großen Ehre des Reiches Christi. Denn wiewohl ihre leiblichen Verheißungen etwas größer, unsere aber geringer sind (denn im Neuen Testament wird kein leiblich Königreich angerichtet, sondern wird vielmehr verboten, damit unsere Herzen nicht beschweret werden mit Fressen und Saufen): so sind wir doch damit zufrieden, daß Christus gesagt hat Matth. 6, 33.: „Trachtet am ersten nach dem Reich GOttes, so soll euch das andere alles zufallen.“ Also wollte er sagen: Es soll zum allerersten mein Name geheiliget werden, mit aller Ehrerbietung, mit geistlicher Freude, mit Worten, mit Herz und Mund, alsdann sollt du auch das tägliche Brod haben.

144. Vornehmlich aber sollen wir uns dessen hoch rühmen und diese unsere Lehre heftig treiben wider den Papst, daßselbe greuliche Reich zu zerstören. Denn wenn ich schon hundert Zungen und so viel Munde hätte, würde ich doch mit Worten nicht ausreden können, welch Ungeheuerlichkeit des Papstes Reich gewesen ist. Und man kann dawider nicht so scharf oder heftig reden, daß nicht dies gottlose Wesen, und der Greuel, so da steht an der heiligen Stätte, dasselbe alles noch übertreffe. Denn ich habe in denselben päpstlichen Stücken erbärmlich gesteckt, und habe mit ganzem Ernst alles gethan und gelitten, das heutiges Tages der größte Theil im Papstthum nicht achtet und auch nicht mehr hält. Derhalben haben wir nun bisher von dem Segen der Väter geredet, dabei ich gesagt habe, daß man auch unsere und des Neuen Testaments Ehre und Herrlichkeit rühmen soll.

\quellseite{488}

145. Darum wer nun kann und mag, der zürne, fluche und vermaledeye immer den Papst; denn er hat dem Reich Christi mehr Schaden gethan denn Mahomet. Der Türke stößet wohl den Leib, nimmt den Christen ihre Güter mit Gewalt hinweg, verderbet und verwüstet Alles: aber der verfluchte Papst bringt viel greulicher auf seinen Alkoran, nämlich, daß Christus soll verleugnet werden. Sie sind zwar beide Feinde der Kirche und sind des Teufels leibeigene Diener; denn sie verwersen beide das Evangelium: aber der Papst will haben, daß man seine Canones und Decretalen soll anbeten, damit das Licht des Evangeliums gar unterdrückt und ausgelöscht werde. Derselben sei er verloren und verdammt in Ewigkeit, und alle Engel und Heiligen sollen diesen Scheusal des Papstthums verfluchen. Denn ich hätte sterben müssen in derselben Mördergrube, wo mich nicht Gott errettet hätte. Denn ich wußte nichts von göttlichen Verheißungen, von Gebrauch der Nutzen der heiligen Sacramente; wie ich droben erzählt habe die Worte eines Mönchs, der sich verwundert, daß so viel Doctorem im Papstthum nie etwas gelehret hätten von der Verheißung Gottes. Ja, wenn ich das Wort und den Namen unsers Heilands Christi Jesu hätte nennen, erschrak ich davor ganz und gar, darum daß ich es dafür hielt, er wäre mir zum Richter und nicht zum Heiland vorgestellt. Ich verwunderte mich vielmehr und hielt in großen Ehren einen Pfaffen, der im langen Rock daher ging oder für die Lebendigen und Todten opferte, denn die Lehre von Christo nebst den Verheißungen und Sacramenten, ja, ich meinte, dieselbige Lehre ginge mich gar nichts an.

\section[Vierter Theil]{Vierter Theil. Wie Jakob seinen Tod verkündigt und die Rückkehr seines Geschlechts nach Canaan; und wie er dem Joseph ein Stück Landes zum voraus vor aus seinen Brüdern schenkt}

\subsection{I}

\begin{quote}
\textbf{V. 21. Und Israel sprach zu Joseph: Siehe, ich sterbe; und Gott wird mit euch sein, und wird euch wiederbringen in das Land eurer Väter.}
\end{quote}

146. Jakob redet vom Lande der Väter nach der Verheißung. Denn diesmal sie in Canaan elend und fremd waren und das Land Canaan verlassen hatten, wie können sie es denn auf andere Weise inne haben, denn nur in der Verheißung, im Glauben und Hoffnung? Was er aber zu Joseph geredet, muß man ynniglich verstehen; denn er war mit beiden Söhnen da. Darum was zu Joseph geredet worden ist, dasselbe geht sie auch mit an und sie sind damit auch gemeinet.

147. Und, Lieber, siehe doch den Glauben und trefflichen Trost an, so Jakob gehabt. Denn er zweifelt daran gar nicht, Gott werde sich ihrer annehmen und sie versorgen, da er sagt: „Gott wird mit euch sein“; wie wir jetzt im Neuen Testament des gewiß sind, daß Christus bei uns sein und bis an der Welt Ende. Ueber das verheißt er ihnen auch, daß sie wieder in das Land ihrer Väter sollen geführt werden. Als wollte er sagen: Ich verlasse euch nun und gehe von euch hinweg aus diesem Leben; ihr sollt aber doch nicht verlassen oder ohne Trost sein. Denn eben derselbe Gott und Erlöser, goel, der bei mir gewesen ist in so viel großer Trübsal dieses Lebens, wird auch bei euch sein, und euch wiederum in das Land führen, so euch verheißen ist.

148. Das sind wahrlich starke Worte. Als, wenn ich sage zu Einem, der jetzt sterben soll: Ich absolviere dich von deinen Sünden, und befehle dich Gott und allen seinen Engeln, die dich schüben und ins ewige Leben führen wollen. Denn dasselbe kann ich also mit aller Zuversicht frei heraus sagen und soll es auch thun um des Befehls Christi willen, und wegen der unbetrüglichen Verheißung Gottes, darauf solche Absolution steht. Darum können wir diejenigen, so den Worten glauben, absolvieren und in das Reich Gottes senden deshalb, daß dieser Segen der Väter uns auch gegeben worden ist.

149. Der Papst aber weiset die armen betrübten Gewissen auf das Fegefeuer, und verheißt ihnen gar nichts Gewisses von ihrer Seligkeit, sondern fordert von ihnen die Reue für die Sünde, Beichte und Genugthuung, ob vielleicht damit Gott möchte erweicht und versöhnt werden. Deine Reue, sagt er, möchte so groß sein, Gott möchte es ansehen. Also läßt er die armen Herzen immer in Zweifel stecken, und bringt sie endlich dahin, daß sie Gott feind werden und gar verzweifeln.

150. Wir aber richten nun die erschrockenen Gewissen also auf, daß wir zu ihnen sagen: Glaube, daß du getauft seiest auf Christum Jesum; ich absolviere dich von deinen Sünden im Namen Christi, der für dich gestorben und auferstanden ist, der da gesagt hat Joh. 14, 19.: „Ich lebe, und ihr sollt auch leben.“ Das ist ein gewisser und beständiger Trost, darauf sich die gottseeligen Herzen allein verlassen können; darum sieht ihn auch der Teufel für und für an und legt sich immer dawider. Wie neulich die großen Esel von Rom in ihren Artikeln, so im Druck ausgegangen sind unter dem Namen und Siegel des Kaisers Carl, das Fegefeuer vertheidigen wollen, mit welchen Artikeln sie ihren groben Verstand und Thorheit an den Tag geben, und beweisen damit öffentlich, daß sie Gott hundertmal mehr verlästert habe, dieweil sie sich unterstanden, das Licht der Verheißung auszulöschen, welches Gott jezuund aus wunderbarer Gnade und Barmherzigkeit in unsern Kirchen wiederum angezündet hat.

151. Darum soll man diese Worte des starken vollkommenen Glaubens fleißig merken, in welchem Jakob gewißlich schließt, Gott werde ihn und seine Nachkommen versorgen auch im Elend und Tod; daß man solche Worte halte gegen den Zweifel, welchen die päpstlichen Sophisten einführen, und denselben Irrthum damit widerlegen. Denn Jakob sagt: Ich sterbe zwar wohl; aber doch nicht wie eine Epicurische Sau, gleich als ob nichts mehr an mir sei, denn allein der Tod und der verstorbene Leib; ja, ich verkündige euch zuvor mit wahrhaftigem beständigem Glauben, daß Gott lebt, und daß er bei euch, die ihr lebet, und bei mir, wenn ich todt bin, auch sein und uns helfen werde; denn ich glaube an ihn.

152. Daraus kann man sehen, daß der Glaube nichts Anderes ist, denn das rechte wahrhaftige Leben in Gott selbst; und aus diesen Worten soll man das ewige Leben und die Auferstehung der Todten erkennen lernen. Ich sterbe wohl dahin, aber Gott lebt. Also reden die Epicurer oder Verzweifler nicht, sondern dies ist ein Wort eines Gläubigen, der mit gewissem Glauben und Hoffnung auf das ewige Leben und die Auferstehung der Todten wartet. Denn er sagt: Es wird der Tod weder mir noch euch schaden; denn Gott wird mit euch und mir sein, ja, ich werde bei ihm sein. So gar gewiß ist er des, daß Gott sie trösten und wieder in das Land ihrer Väter führen werde, als ob sie jetzt alle mit einander hinauf zögen und das Land in der That vor ihren Augen hätten.

153. Solche große Kraft hat der Glaube, daß er, wenn wir gestorben sind, wiederum lebendig macht: und zwar in derselbigen Stunde, darin wir anfangen zu glauben und das Wort zu ergreifen, heben wir auch an im ewigen Leben zu leben. Denn das Wort des Herrn bleibet ewig, und Gott, der mit uns redet, ist ewig und wird auch ewig bei uns sein.

154. Derhalben sind Tod und Leben näher bei einander gesetzt, denn wir selbst sehen können. Denn Jakob stirbt dahin, aber der Tod ist nahe bei dem Leben, ja, er ist eben das rechte Leben. Wiederum, Adam lebte im Paradies und wußte von seinem Tode, wiewohl ihn Gott zuvor gesagt, wie nahe ihm der Tod wäre. Denn also lauten die Worte 1 Mos. 2, 17.: „Welches Tages du davon issest, wirst du des Todes sterben.“ Er konnte es aber nicht glauben, da er noch das Leben und alles Gute fühlte und das Böse noch nicht erfahren hatte. Aber Gott erinnerte ihn an beiderlei. Als wollte er sagen: Du fühlst das Leben wohl, den Tod aber fühlst du nicht; du siehst das Unglück nicht, so dir begegnen möchte: aber siehe dich vor, der Tod ist gewiß nicht weit. Denn das sollst du eben wissen, daß du gleich desselben Tages sterben wirst, an welchem du von dem Baum essen wirst; so gar nahe sind Tod und Leben einander. Aber dessen hat Gott Adam nicht überreden können, und er gedachte, er würde nimmermehr sterben. Da regierte der Glaube, darum fürchtet er sich vor dem Tode gar nicht.

\quellseite{489}

155. Also sind wir auch im Tode. Wir sehen, daß wir sterben und den Geist aufgeben, daß der Leib verfault und von den Würmern verzehrt wird; doch thun wir eben auch das im Tode, was Adam im Leben gethan hat. Adam glaubte nicht, daß der Tod so nahe wäre, sondern war sicher, und seine Sicherheit und Hoffahrt war eine heilige und rühmliche Sicherheit, da er gar nichts Böses und kein Verderben gefühlt hat. Also sind wir auch zerschlagen und gedemüthigt im Tode, und haben doch das Leben eben so nahe, wie Adam den Tod gehabt; denn die Verheißung ist gewiß und wahrhaftig, Marc. 16, 16.: „Wer da glaubet und getauft wird, der wird selig werden“; da ist uns, die wir im Tode sind, das Leben und Seligkeit ganz nahe, wie Adam der Tod nahe war. Aber wer da glaubt, wie hier Jakob sich mit starkem Vertrauen des Lebens versiehet, derselbe achtet den Tod für nichts; denn der Glaube bringt das Leben nahe herbei, und schließt also: Wiewohl ich sterbe, so bin ich doch getauft und glaube an Christum; darum bin ich lebendig und selig.

156. Wie nun Adam hat rühmen können, daß er lebte, und sollte doch gleichwohl diesem Wort auch glauben: „Welches Tages du davon issest, wirst du des Todes sterben“, und sich vor dem Tod fürchten: also sollen wir dagegen den Tod nicht fürchten, sondern auf das Leben hoffen. Denn Gott, der zu Adam gesagt hat: Der Tod ist dir nahe, derselbe Gott sagt auch zu uns: Das Leben, Friede und Seligkeit ist nahe bei euch. Und Gott, der also mit uns redet, kann nicht lügen oder trügen; wie Paulus sagt Tit. 1, 2.: „Gott, der uns das ewige Leben verheißen hat, ist wahrhaftig, und lüget nicht“, welches ein Wort ist des Glaubens, der das Leben schon in sich hat, wiewohl es verborgen ist; denn der Glaube fängt in uns das Leben an und hat das Leben in sich selbst.

157. Und die wir glauben, haben diesen Anfang, daß ob wir auch jetzt den Tod fühlen, wir uns doch vor ihm nicht fürchten, wie die Andern, die vom bösen Gewissen geplagt werden, auch bleich werden und erschrecken, wenn sie nur den Tod nennen hören. Aber die gottseligen und heiligen Märtyrer verachten und verspotten den Tod. Wie St. Agnes, da sie hingerissen ward in den Kerker und gepeinigt werden sollte, gesagt: Es wäre ihr nicht anders zu Muth, als würde sie zum Tanz geführt. Lieber, vorher mag doch das Mägdlein solchen großen Muth gehabt haben, daß sie sich so gar nicht gefürchtet, noch erschrocken ist, sondern ist fröhlich und guter Dinge, als feste man sie zu Tische, da man ihr sehr gütlich thun wollte? Es ist zwar nicht eine Epicurische Verachtung des Todes gewesen, sondern rechte Weisheit und rechter Verstand, daraus sie hat schließen können und es dafür halten, daß ihr das Leben ganz nahe wäre. Darum hat sie des Teufels und Todes gespottet und ihrer gleichsam gelacht; denn der Tod war in ihr verschlungen durch das Leben.

158. Dies ist unsere Theologie, die wir lehren, so der gar zuwider, über welcher die blinden und närrischen Scholastiker und Papisten so gar hart halten, die hiervon gar nichts wissen und den Glauben verachten. Derhalben lassen uns den Patriarchen Jakob hören, der vom Tode wie vom Schlaf redet. Denn wenn du ihn fragen wirst: Lieber Jakob, wie, daß du vor dem Tode nicht erschrickst? so antwortet er: Ich werde zwar sterben und in das Grab gelegt werden: aber Gott lebet, der uns das Land verheißet hat, darein er euch führen wird; mich aber wird er in ein ander Leben, das viel besser ist, setzen; denn er hat es verheißen.

159. Darum ist alle sein Leben in dieser Verheißung verschlungen gewesen: Gott hat verheißen, er wolle mir gnädig sein, so auch, daß von mir der Messias soll geboren werden; derselben wird er auch das Land und den Ort geben, da dieser Messias geboren werde. In dem Glauben und Vertrauen will ich sterben und in Frieden schlafen.

160. Dies wird uns nun zum Exempel, dem wir nachfolgen sollen, vorgehalten, daß wir unsere Verheißungen auch groß achten. Denn vor Zeiten im Papstthum, da ich ein Mönch war, brauchte man dieser Wörter nicht, und wußte davon niemand etwas zu sagen, was „Wort“ oder „Verheißung“ war; und ich danke Gott, daß ich zu dieser Zeit leben mag, da nun „Wort“ und „Verheißung Gottes“ in meinen und aller Gottseligen Ohren klingt. Denn wer das Wort hört, der kann leichtlich die göttlichen Verheißungen verstehen, die im ganzen Papstthum allen Theologen dunkel und unbekannt waren.

161. Es muß aber das Wort, welches Gott verheißt, mit dem Glauben verbunden sein; welchen Glauben sie auch nicht im Papstthume verstanden haben, so wenig wie das Wort und Verheißung. Denn dies war ihre gewöhnliche Lehre: Die Sacramente geben denen, die sie brauchen, Gottes Gnade ohne Wort und ohne Glauben: wenn Einer getauft wird, der bedarf des Glaubens nicht; die Sacramente haben so viel Kraft, daß sie wegen des äußerlichen Werks Gnade geben.

162. Dies ist aber das Hauptstück unserer Lehre, daß kein Sacrament an sich selbst ohne Glauben Gnade wirken könne. Denn der Kämmerer, Ap. Gesch. 8, 36. 37., hörte von Philippus, da er sagte: „Siehe, da ist Wasser, was hindert's, daß ich mich taufen lasse?“ „Glaubest du von ganzem Herzen, so mags wohl sein.“ Philippus will ihn nicht taufen, er glaube denn. Denn er sagt nicht: Die Taufe ist dir nütze und macht dich gerecht, du glaubest oder glaubst nicht; wie sich die Papisten träumen lassen, daß auch den kleinen Kindern die Gnade eingegossen werde aus Kraft des Sacramentes. Das ist höchst falsch. Denn sie werden selig aus Kraft der Verheißung, und empfangen den Heiligen Geist, darum daß Christus Matth. 19, 14. gesagt hat: „Lasset die Kindlein zu mir kommen, denn solcher ist das Himmelreich.“

163. Derhalben lassen uns an solchen Exempeln die rechte Kraft des Glaubens lernen, dadurch wir errettet werden vom Tode und in das ewige Leben gehen. Wie Jakob hier sagt: Ich sterbe; aber Gott lebt, der wird bei euch sein. Ich gehe in das Grab, ihr werdet mich weder hören noch sehen; ich will euch aber einen viel bessern Vater lassen, der euch viel besser und treulicher wird versorgen können denn ich. Denn weder ich, noch mein Sohn Joseph, sondern der Gott, an den wir glauben und deß Verheißung wir haben, wird selbst euer Gott, euer Vater, euer Führer und Erlöser sein.

164. Dies alles hätte Jakob ohne Glauben nicht sagen können. Denn die Verheißung und der Glaube hangen nothwendig an einander, wie St. Paulus Eph. 3, 17. sagt: „Christus wohnet in euren Herzen durch den Glauben“; und im Buche der Weisheit Cap. 1, 4. steht auch geschrieben: „Gott wohnet nicht in einem ungläubigen Menschen.“ Wenn der Türke hörte, daß Einem etwas verheißen würde, könnte er bald verstehen, daß daselbst der Glaube gefordert würde. Denn was ist eine Verheißung, welcher niemand glaubt, denn nur ein unnütz vergeblich Ding? Alsdann ist es aber eine rechte Verheißung, wenn der Mensch daran glaubt und sich fest darauf verläßt, und hält Gott für wahrhaftig.

165. Der historische Glaube hält sich nicht also an das Wort, setzt auch sein Vertrauen nicht darauf, sondern sagt: Ich höre, daß Christus gelitten hat und gestorben ist etc.; aber der rechte Glaube hält es also und sagt: Ich glaube, daß Christus für mich den Tod gelitten habe und gestorben sei etc., daran zweifele ich auch gar nicht, und in dem Glauben bin ich zufrieden und verlasse mich auf das Wort wider den Tod und die Sünde.

166. Dies Licht und die Gnade hat uns Gott zu dieser Zeit geoffenbart; dafür sollen wir ihm dankbar sein. Und wir haben des Glaubens ein herrlich Exempel in diesen Worten Jakobs, die wahrlich große wichtige Dinge in sich haben. Denn Jakob stirbt dahin, nicht als der in seinem Herzen gezweifelt oder geschwankt hätte, sondern er sagt aus dem gewissen und beständigen Glauben frei heraus: „Gott wird mit euch sein.“ Also hat er im Tode das Leben, und da er schon todt ist, lebt er doch gleichwohl, er wird zu Grabe gebracht voll Glaubens; denn er vertraut Gott und seinen Verheißungen. Dieser Glaube verschlingt den Tod. Denn es ist nicht eine Erkenntniß nach der Historie, welche allein sich begreift die Geschichte, streitet aber nicht, besteht auch nicht fest in der Anfechtung wider den Tod und die Hölle; wie denn der Teufel einen solchen Glauben und Erkenntniß auch hat.

167. Aber das ist der rechte Glaube der Verheißung, der Gott die Ehre zuschreibt, daß er wahrhaftig sei; der seinem Worte glaubt und sein Vertrauen darauf setzt. Ein solcher Glaube überwindet den Tod und trotzt gleichsam mit ihm; wie Paulus thut, und sagt 1 Cor. 15, 55.: „Tod, wo ist dein Stachel? Wo ist dein Sieg?“ Denn wenn du mich schon wirst verschlingen, werde ich doch wiederum an das Licht und zum Leben kommen.

168. Solche Stellen der Schrift soll man der Kirche vorhalten, um der Exempel willen der heiligsten Männer, die uns mit dem beständigen Glauben vorgeleuchtet haben, welchem Glauben wir auch sollen lernen nachfolgen, daß wir auf diese Weise uns rühmen mögen und sagen: Ich bin getauft und absolvirt, darauf sterbe ich. Was mir nun weiter für Anfechtung und Bekümmerniß vorkommen wird, soll mich gar nicht irre machen. Denn der da gesagt hat Marc. 16, 16.: „Wer da glaubet und getauft wird, der wird selig werden“; item Matth. 16, 19.: „Alles, was du auf Erden binden wirst, soll auch im Himmel gebunden sein“; item Matth. 26, 26. 28.: „Das ist mein Leib, das ist mein Blut, das für euch vergossen wird zur Vergebung der Sünden“: derselbe kann weder lügen noch trügen; das ist gewißlich wahr.

\quellseite{490}

\subsection{II}

\begin{quote}
\textbf{V. 22. Ich habe dir ein Stück Landes gegeben außer deinen Brüdern, das ich mit meinem Schwert und Bogen aus der Hand der Amoriter genommen habe.}
\end{quote}

169. Das ist aber nun ein sehr wunderbarer Ding, daß dieser sehr heilige Mann, der jetzt im Glauben ganz verschlungen ist und im Herrn sicher lebt, glaubt und auf die Auferstehung der Todten hofft, doch sich noch bekümmert um ein kleines armes Stücklein des heiligen Landes; und wiewohl dasselbe klein gewesen, hat er dessen doch nicht vergessen wollen.

170. Also haben diese heiligen Männer gethan: wiewohl sie des Lebens und Auferstehung der Todten gewiß gewesen sind, verlassen sie doch ihr Amt und Beruf, und dazu auch die zeitlichen Güter nicht; sie sorgen wohl nun nicht mehr für sich selbst, aber sie bestellen doch und verordnen Alles richtig für ihre Nachkommen, daß sie ja nichts hinter sich lassen, daraus ihre Kinder oder Verwandten möchten Ursache nehmen, mit einander zu hadern oder zu zanken. Also thut auch ein frommer gottseliger Hausvater: ehe er von diesem Leben abscheidet, verordnet und bestellt er Alles in seinem Hause, und will haben, daß von seinem Gute und Habe ein Theil seinem Weibe, ein Theil seinem Sohne oder seiner Tochter soll gegeben werden. Denn er denkt also: Diese Güter sind mir von GOtt befohlen und vertraut worden, dieselben will ich nun mit Frieden und im Glauben recht bestellen und austheilen.

171. Also sagt nun Jakob zu Joseph: Außerhalb der Theilung des Landes, so du mit deinen Brüdern ingemein halten wirst, gebe ich dir dies Stück Landes insonderheit, das sollst du mit den Andern nicht theilen. Dir gebühren sonst auch noch zwei Theile für deine zwei Söhne: aber über diesen sollst du diesen Theil sonderlich für dich haben. Also ehrt er seinen Sohn mit einer sonderlichen Gabe in Ansehung seiner Herrlichkeit und Gottseligkeit. Denn er ist ein großer vortrefflicher Mann gewesen, geziert und begabt mit allerlei Tugenden.

172. Wiewohl nun solche äußerliche Dinge und leiblichen Güter im Hause Jakobs nicht anders denn auch sonst an andern Orten vergänglich gewesen: doch, dieweil daselbst die Kirche und das Reich GOttes sein sollte, und aus dem Hause so viel Propheten, Könige und vornehme Männer haben sollen erweckt werden, welche in der Welt große Wundermänner gewesen sind; hat GOtt um der zukünftigen Nachkommen willen diesen Ort so groß geachtet, da er hat wohnen, reden und Christum den HErrn hinsenden wollen. Das ist wahrlich eine sehr treffliche und sehr große Ehre gewesen, die alle zeitlichen Güter weit übertrifft, nämlich, GOtt selbst gegenwärtig im Reiche und Gottesdienste bei sich haben; wie Christus sagt Joh. 14, 23.: „Wer mich liebet, der wird mein Wort halten; und mein Vater wird ihn lieben, und wir werden zu ihm kommen, und Wohnung bei ihm machen.“

173. Dies sind sehr große unermeßliche Dinge, und wenn man dagegen die Herrlichkeit aller Königreiche, Edelsteine und alles Geldes und Gutes der ganzen Welt und was das alles kosten möchte, halten wollte, würde doch solches alles mit einander gar verschwinden. Denn welch ein groß Ding ist es doch, daß der Sohn GOttes sagt: Ich und der Vater wollen zu ihm kommen, und Wohnung bei ihm machen? Derhalben ist das die vornehmste Ehre desselben Landes gewesen, daß an dem leiblichen Orte das ewige Leben, der Sieg über Sünde und Tod war.

174. Es zankten aber die Grammatiker mit einander über das Wort schekem, und die Hebräer machen einen Unterschied zwischen der Schulter oder Achsel und dem Stücke, so man auf der Achsel trägt. Nun heißt zwar das Wort schekem allerlei Achsel; aber aus welchen Ursachen es auch heißen soll das Stück, so man darauf trägt, kann ich nicht verstehen, denn allein, dieweil es den Juden gefallen hat, dasselbe Wort also auszulegen, darum daß ein Stück oder Theil auf den Achseln getragen wird, wie sie davon fabuliren oder aber unnütz reden; oder aber darum, daß die Achsel ein Theil des Leibes ist. Aber es ist eine sehr harte und jüdische ungereimte Metapher oder Gleichniß.

175. Derhalben hasse ich an den jüdischen Rabbinern das sehr und verachte es auch gar, daß sie so ganz geflissen sind auf die mancherlei Bedeutung der Wörter. Und Quintilian lehrt und ermahnt recht, daß man sich vor einem solchen Worte, das mehr denn einerlei Bedeutung hat, hüten soll, wie vor einer Klippe, da man pfleget anzulaufen. Und sollen sich diejenigen, so Andere lehren, dessen vornehmlich befleißigen, daß sie solche mancherlei Deutung der Wörter fliehen sollen. Man hat auch noch einen alten Vers, der in den Schulen sehr gemein gewesen ist, und lautet also: Erroris mater est aequivocatio semper, das ist: Es kommt allezeit Irrthum von den Wörtern, die mehr denn einerlei Bedeutung haben.

176. Es kann aber wohl kommen, daß etliche Wörter von Art und Natur zweideutig sind, welche in den Schulen aequivoca a casu, das heißt, durch Zufall zweideutig, genannt werden. Solche Zweideutigkeit kann man in den Sprachen nicht meiden, oder sich deren gar enthalten. Wie wir in der deutschen Sprache mit einem Worte ausdrücken die beiden Wörter, molere und pingere, und sagen, „malen.“ Da haben wir nur Ein Wort, und man kann daraus keinen Unterschied abnehmen zwischen den beiden lateinischen Worten molitor und pictor. Desgleichen das deutsche Wort „Reiff“ bedeutet auch in der lateinischen Sprache zweierlei, pruina und maturitas. Also auch, das Wort „Nagel“ bedeutet unguis und clavus. Und dasselbe geschieht in andern Sprachen auch also, daß ein Wort nach gemeinem Gebrauch der Sprache von Art und Natur mehr denn einerlei bedeutet und man eigentlich seine rechte Kraft oder Verstand nicht erkennen kann, denn allein aus den Umständen, welche zugleich mit der Construction, so die Grammatik gibt, das ist, darnach sie zu andern Worten gesetzt sind, die rechte Bedeutung anzeigen. Als, wenn Einer wollte zum Maler sagen: Malet mir das Korn, der würde bald den Pinsel nehmen und ein Einhorn malen. Und wiederum auch also, wenn du dem Müller ein Papier oder Brief bringen würdest, und zu ihm sagen, er sollte dir darauf ein Bild malen, würde er nicht verstehen können, was du damit meinen möchtest. Von solchen Wörtern reden wir an diesem Orte nicht.

177. Es ist aber eine andere Aequivocation, da die Worte nicht von Art und Natur mehr denn einerlei bedeuten, sondern wie es damit der Sprache oder Rede halben gebräuchlich ist. Als, wenn man im Lateinischen saget: Homo et homo pictus, ein Mensch und ein gemalter Mensch. Solche Aequivocation mag man auf Lateinisch heißen, artificialis, voluntaria oder arbitraria; das ist, künstlich und freiwillig. Und derselben haben sich die Rabbinen ohne alle Maßen viel beflißen, daß man ihrer Bemühung hierin derhalben auch billig mag feind sein. Wo sie sehen, daß irgend etwas geredet worden ist durch eine Figur, entweder die da heißt metaphora, oder aber synekdoche, dichten sie alsbald auf solche Worte neue Bedeutungen. Denn es sind ungelehrte Leute, die nichts wissen von guten Künsten; darum verstehen sie die Figuren nicht.

178. Es kommt oftmals vor, daß etliche Worte um deswillen, daß sie einander etwas gleich sind, auch auf andere Dinge gezogen werden, welches in den Metaphern also geschieht. Aber in den Worten, so ihrer Art und Natur halben mehr denn einerlei bedeuten, ist kein Gleichniß, als unter den lateinischen Worten, molitor und pingere, welche beide „malen“ heißen, und man kann dafür keine Ursache anzeigen.

179. Also steht in der Historie Abrahams das Wort beer scheba, 1 Mos. 21, 31., das wird auf Lateinisch gedeutet, fons saturitatis, oder juramenti, oder septem, das ist, ein Brunnen des Überflusses, oder des Eides, oder sieben; denn dies alles heißt das Wort scheba. Seine rechte ursprüngliche Bedeutung ist, daß es Eid heißt; als, im 110. Psalm V. 4.: „Der HErr hat geschworen, und wird ihn nicht gereuen.“ Aber warum es dahin gezogen werde, daß es sieben oder ein Überfluß heißen soll, weiß ich nicht, außer daß in der hebräischen Sprache, wie in andern auch, solche Wörter etwa workommen,
\quellseite{491}
die in den Schulen genannt werden, aequivoca a casu, zufällige Zweideutigkeiten. Wie man aber dieselben auslegen soll, das zeigen die Umstände im Texte an.

180. Also ist schein auch ein Wort, das mehr denn einerlei Bedeutung hat. Denn bisweilen heißt es ein Stück oder Theil, bisweilen aber eine Achsel; welche Bedeutung aber sich an einem jeden Orte am besten schicken wolle, das muß man, wie ich gesagt habe, aus den Umständen urtheilen. Als, im Propheten Zephania am 3. Cap. V. 9. steht also: „Alsdann will ich den Völkern anders predigen lassen, mit freundlichen Lippen, daß sie alle sollen des HErrn Namen anrufen, und ihm dienen einträchtiglich“, da steht im Lateinischen: ut serviant ei humero uno, welches die Ausleger deuten: parte una, vel opinione, vel consensu unanimi, das ist, auf Einem Theil, oder einerlei Meinung, oder einmüthiglich; also daß Viele wie auf Einer Achsel Christum tragen und rechte Christopfer (Christusträger) sind.

181. Auf dieselbige Weise sagt Jakob hier auch in einerlei Bedeutung: „Ich habe dir ein Stück Landes gegeben“, und deutet es doch auch auf den Eigennamen „Sichem“. Denn es ist hier dreierlei Bedeutung, wenn der Name dieser Stadt zu den vorigen Bedeutungen kommt, und spielt also Jakob mit dem Worte, daß so mancherlei Bedeutung hat, und sagt: Ich gebe dir dies Stück Landes, und will damit verstanden haben den Ort, Sichem genannt. Mit dem Gattungswort eignet er ihm dies Theil zu, und sagt: Das soll dein Theil sein. Aber mit dem Eigennamen will er die Stadt Sichem gemeint und verstanden haben. Wie 1 Mos. 21, 30. 31., da Abraham sieben Lämmer besonders stellt, und da Abimelech fragt: „Was sollen die sieben Lämmer, die du besonders dargestellt hast?“ antwortet er: „Daß sie mir zum Zeugniß seien, daß ich diesen Brunnen gegraben habe.“ Daher heißt die Stätte Bersaba, nämlich von den sieben Lämmern oder vom Eide, daß sie beide mit einander da geschworen haben und einen Bund gemacht, auf Deutsch, Eidstatt oder Siebenstatt.

182. Nun sind aber noch etliche Fragen da, so zur Grammatik und nicht zur Theologie gehören. Denn Jakob setzt hinzu, und sagt: „Daß ich mit meinem Schwert und Bogen aus der Hand der Amoriter genommen habe.“

Wir haben aber droben gehört im 34. Cap. V. 25. ff., wie Simeon und Levi Sichem und Hemor mit den Einwohnern der Stadt erwürgt und die ganze Stadt geplündert haben. Von demselben Todtschlage und Raube oder Plünderung redet nun Jakob hier also: Ich habe dies Stück Landes mit meinem Schwert und Bogen gewonnen. Wiewohl, wie Stephanus zeuget Ap. Gesch. 7, 5., Gott ihm kein Erbtheil in dem Lande gegeben, auch nicht eines Fußes breit, und seine Söhne auch nicht auf Gottes Geheiß und Befehl, sondern aus eigenem Vermessen und grimmigen Zorn die Sichemiter erschlagen haben. Und Jakob billigt doch jetzt dieselbe That, und sagt, daß er es gethan habe, und bestätiget also mit väterlicher Gewalt solche böse That. Warum lobt er jetzt, daß er droben gestraft und verdammt hat, da er sagte 1 Mos. 34, 30.: „Ihr habt mir Unglück zugerichtet, daß ich stinke vor den Einwohnern dieses Landes, den Cananitern und Pheresitern“? Dem aber allem gar zuwider rühmt er sich nun hier und sagt: „Ich habe dies Stück Landes mit meinem Schwert und Bogen genommen.“

183. Wir haben aber droben auf diese Frage geantwortet, und dieselbe Antwort soll man hierher nehmen und wiederholen. Denn die Söhne Jakobs haben diesen Todtschlag nicht allein gethan, haben es auch nicht thun können ohne Hülfe der Cananiter, die sich zu ihnen gesammlet, dazu sie die unbillige Uebelthat bewogen, welche Sichem begangen hatte, da er eines solchen großen theuern Mannes, der ein Gast und fremd war in ihrem Lande, seine Tochter geschwächt hatte. Der Eifer hat die Cananiter dazu getrieben, daß sie den Söhnen Jakobs nicht allein geholfen, die Mißhandlung, so sie mit ihrer Schwester begangen, zu rächen, sondern daß sie ihnen auch das ganze Land zugeeignet und eingeräumt haben. Denn es werden ohne Zweifel Etliche den Patriarchen Jakob predigen gehört haben, und haben es dafür gehalten, man müßte in allewege die Uebelthäter hart strafen, damit man Zucht und Ehrbarkeit erhalten möchte.

184. Derhalben, wiewohl der Anfang böse gewesen um der frevelen That willen seiner Söhne, ist doch der Besitz des Landes Jakob geblieben, wie es mit dem Schwert gewonnen war, dazu auch die Nachbarn selbst geholfen und mit eingewilliget hatten. Dies ist meine Meinung, und halte, daß man derselben darum billig folgen müsse, dieweil die Historie sich mit diesem gegenwärtigen Texte reimt. Es wird dieser Stätte auch gedacht Joh. 4, 5., da sie Sichar genannt wird; welches daher gekommen ist, daß in der Grammatik mit dem Schreiben geirret worden; denn der Bücherschreiber hat in den griechischen Texte aus dem Buchstaben μ ein ρ gemacht.

185. Die letzte Frage gehört ganz zur Grammatik, nämlich: Warum er doch sage: „aus der Hand der Amoriter“, so doch droben Sichem ein Heviter genannt wird, und die Amoriter und Heviter verschiedene Völker gewesen sind? und die Schrift zeuget, daß die Amoriter beinahe jenseit des Jordans gewohnet; wie das aus dem fünften Buch Mose zu sehen ist von Sihon, dem Könige der Amoriter. Wie wird ihnen denn Sichem zugeeignet, daß doch diesseits des Jordans gelegen gewesen? Da folge ich der gemeinen Glosse, daß solches synedochisch geredet sei, wenn man ein Ganzes für einen Theil, oder einen Theil für das Ganze versteht; wie 5 Mos. 1, 20. auch gesagt wird: „Ihr seid an das Gebirge der Amoriter kommen“, wo gleichfalls auch Amoriter genannt werden, die im Lande Juda wohnten. Derselben ist das Wort Amorheus, Amoriter, entweder ein Gattungswort gewesen, oder diese Völker sind Schwägerschaft oder Bündniß halben einander verwandt gewesen, also daß Sichem seinen Ursprung vom Geschlechte eines Theils der Amoriter und eines Theils der Heviter gehabt, welches unter den Völkern, so nahe bei einander gelegen, oftmals zu geschehen pflegt, als, unter Sachsen und Meißnern. Denn also pflegen sich die Leute zu halten. Darum ist an dieser Frage nicht viel gelegen.

\chapter{Das neunundvierzigste Capitel}

\section[Erster Theil]{Erster Theil. Wie Jakob seine Söhne zubereitet zu dem Fluch und Segen, den er ihnen verkündigen will; und wie er seinem Sohne Ruben den Fluch verkündigt}

\subsection{I}

\begin{quote}
\textbf{V. 1. 2. Und Jakob berief seine Söhne, und sprach: Versammelt euch, daß ich euch verkündige, was euch begegnen wird in künftigen Zeiten. Kommt zu Hauf, und höret zu, ihr Kinder Jakob, höret euren Vater Israel.}
\end{quote}

1. Dies Capitel ist nun etwas dunkel und schwer, und das um der Ursache willen, daß es viel Meister und Ausleger hat, die über den Verstand gar uneinig sind, und von denen der Eine dies, der Andere das hält, auch so, daß wir aus einem bösen krummen Wege einen ebenen und geraden machen müssen, darauf wir gehen. Wir wollen aber thun, was wir können.

2. Jakob erweckt seine Söhne und macht sie aufmerksam, als der von vortrefflichen, merkenswerthen Dingen reden will, nämlich, von der Verheißung und Drohung. Und ist auch in Wahrheit ein groß hoch Ding, so gar gewiß von zukünftigen Dingen etwas zuvor sagen und verkündigen.

3. Wir wissen aber aus der heiligen Schrift und aus lebendiger Erfahrung, daß die göttlichen Verheißungen und Drohungen, wenn sie erfüllt werden und gleichsam im Schwange und vollen Laufe gehen, sich doch gar mißverständig ansehen lassen. Denn also hat das Volk Israel treffliche Verheißungen gehabt; wenn du jedoch die Historie ansehen wirst und wie es ihnen gegenwärtig ergangen, scheint es, daß sie die ärmsten Leute und gar verlassen sind, und daß GOtt in seinen Verheißungen entweder nicht wahrhaftig sei, oder aber derselben selbst nicht viel achte: wie wir auch von Abraham, Isaak und Jakob davon etliche Exempel gehabt.

\quellseite{492}

Denn ihre Historien zeigen nichts Anderes an, denn daß sie ohne GOtt gewesen, oder aber daß GOtt ihrer gar vergessen und für sie nicht mehr gesorgt habe, und ihren Feinden und Widersachern frei zugelassen, daß sie möchten regieren, triumphiren, und ihres Gefallens den gottseeligen unschuldigen Leuten nur weidlich zusetzen, wider sie toben und wüthen.

4. Auf dieselbige Weise ist das Volk Israel von den Nachbarn, ihren Feinden, als, den Philistern, Ammonitern, Moabitern, Idumäern, Egyptern und andern mehr, hart geplagt worden, daß nichts weniger am Tage war, als daß sie sollten göttliche Verheißung haben. Und was kann doch erbärmlicher sein, denn die Historie des Patriarchen Jakob, die voll allerlei Jammer und Trübsal ist? Ist aber das die Verheißung und der Segen, so er gehabt hat? Heißt das einen gnädigen GOtt haben, Leben und ewige Seligkeit haben? Wahrlich, es haben die Gottseligen gar das Gegentheil vor ihren Augen und, wie es scheint, ist solches alles auch in der That dem allen zuwider, was ihnen angeboten und vorgetragen wird in den Verheißungen, von welchen die Vernunft urteilt, daß sie falsch und erlogen sind, und wie man zu sagen pflegt: Daß es den Frommen übel und den Bösen wohl geht.

5. Derhalben gehört zu der Verheißung der Glaube, der sich halte und lebe in dem Wort der Verheißung. Denn wo der Glaube nicht da ist, der sich an die Verheißung hält und in derselben lebt, so scheint es möglich nicht anders, denn daß die Dinge, so den Gottseligen begegnen, der Verheißung stracks zuwider sind, also daß die Leute getrieben werden, daß sie gar verzweifeln müssen. Darum schreit die Menge der Juden, so andern Göttern räucherten, Jer. 44, 16. 17. 18., und sagt zu dem Propheten also: „Nach dem Wort, das du im Namen des HErrn uns sagest, wollen wir dir nicht gehorchen; sondern wir wollen thun und nach alle dem Wort, das aus unserm Munde gehet, und wollen Melecheth des Himmels räuchern, und derselbigen Trankopfer opfern, wie wir und unsere Väter, unsere Könige und Fürsten gethan haben in den Städten Juda und auf den Gassen zu Jerusalem. Da hatten wir auch Brod genug, und ging uns wohl, und sahen kein Unglück. Seit der Zeit aber wir haben abgelassen, Melecheth des Himmels zu räuchern und Trankopfer zu opfern, haben wir allen Mangel gelitten, und sind durch Schwert und Hunger umkommen“ etc.

6. Dies ist die Ursache gewesen, daher im Volke Israel alle Abgötterei gekommen ist, auch in der Wüste, wie das der 78. Psalm V. 8. bezeugt, daß eine abtrünnige und ungehorsame Art gewesen sei, „welchen ihr Herz nicht fest war, und ihr Geist nicht treulich hielt an GOtt“. Denn sie begehrten, daß sie also möchten geführt und regieret werden, daß sie nicht im Glauben der Verheißung leben müßten, sondern wollten es alles gegenwärtig haben, davon sie leben möchten. Dagegen aber wollte GOtt haben, daß sie sich an den Glauben halten sollten, und nicht gaffen oder stehen nach den Dingen, die ihnen gegenwärtig möchten wohl thun und angenehm sein. Darum tobten und zürnten sie also; als wollten sie sagen: GOtt soll uns nicht mit Worten speisen, sondern soll uns jetzt baar übergeben, was wir bedürfen. Also begehrten sie, was dem Willen und der Regierung GOttes nur stracks zuwider war, und wenn sie alsbald gegenwärtig nicht empfingen, was sie baten und begehrten, so suchten sie von Stund an einen andern Gott. Und dasselbe verkehrte Wesen menschlicher Herzen wird in demselben ganzen Psalm beschrieben: „Sie hielten den Bund GOttes nicht, und wollten nicht in seinem Gesetze wandeln, und vergaßen seiner Thaten und seiner Wunder, die er ihnen erzeiget hatte“, Ps. 78, 10. 11.; desgleichen auch im 106. Psalm V. 13.: „Aber sie vergaßen bald seiner Werke, sie warteten nicht seines Raths“, daß sie so viel gesagt, sie wollten nicht glauben. Und eben dasselbe wirft ihnen Stephanus auch vor Ap. Gesch. 7, 43.: „Und ihr nahmet die Hütte Moloch an, und das Gefäß eures Gottes Remphan“ etc.

7. Darum muß man darauf mit Fleiß merken und Acht geben in der heiligen Schrift, darin ganz andere Historien begriffen sind, denn in den Schriften der Heiden, die keine göttliche Verheißung in sich haben. Dagegen aber gibt GOtt seinem Volke wohl reiche Verheißung, aber er verspricht sie auch zugleich und prüft ihren Glauben, und lehrt sie, daß sie mehr vom Worte denn vom Brode leben müssen, wie Mose 5 Mos. 8, 3. sagt: „GOtt demüthigte dich, und ließ dich hungern, und speisete dich mit Man, das du und deine Väter nie erkannt hattest; auf daß er dir kund thäte, daß der Mensch nicht lebt vom Brod allein, sondern von allem, das aus dem Munde des HErrn gehet.“

8. Das soll nun ja fleißig behalten, daß GOtt wohl verspricht, und doch mit dem, so er verheißen hat, immer etwas verzieht, ehe er es gibt, und versucht uns also, daß wir gegenwärtig Mangel haben und darben müssen, daß er uns im Glauben an die Verheißung unterrichte, und derselbe Glaube gestärkt und bestätigt werde, und GOtt vertrauen lerne, nicht allein wenn es wohl zugeht und genug vorhanden ist, sondern auch wenn es übel zugeht und Mangel oder Noth vorhanden ist. Solchen Glauben hat der Papst und alle seine Doctoren oder Lehrer nicht, wissen auch davon gar nichts.

9. Das ist aber das vornehmste Stück in der ganzen heiligen Schrift, daß man GOtt in seinen Verheißungen erkenne und dieselben recht verstehe. Denn er hilft auch mit der That und endlich im Werke, aber er prüft zuvor den Glauben in der Verheißung, daß er uns mangeln läßt an dem, das uns noth thut, auf daß wir ihm vertrauen lernen, und ihn ja nicht versuchen, wie die Juden thaten; denn dieselben folgten alsobald fremden Göttern nach, wenn ihnen GOtt nicht alsobald ihres Gefallens und wie sie ihn vorschrieben, alles gab, was sie nur begehrten; denn sie konnten seines Raths nicht erharren.

\quellseite{493}

10. Wider diese Ungeduld werden so viel Ermahnungen gesetzt in den Propheten und Psalmen, als, Hab. 2, 3.: „Die Weissagung wird ja noch erfüllet werden zu seiner Zeit, und wird endlich frei an Tag kommen, und nicht außen bleiben. Ob sie aber verzeucht, so harre ihr; sie wird gewißlich kommen, und nicht verziehen.“ Und im 27. Psalm V. 14.: „Harre des HErrn, sei getrost und unverzagt, und harre des HErrn.“ Das sahen aber die Juden nicht, begehrten das auch nicht zu lernen, darum erwählten sie andere Gottesdienste, und suchten solche Götter, als Astaroth, Baal, die ihnen bald und gleich gegenwärtig helfen sollten, ohne Glauben der Verheißung, daß sie die Süße fühlen und mit Händen greifen möchten. Wie sie denn in ihrem Stolze ganz vermessen zu Jeremia Cap. 44, 17. sagten: „Wir wollen thun nach alle dem Wort, das aus unserm Munde gehet“ etc.; wir wollen kurzum dein Wort nicht hören. Als wollten sie sagen: Wir wollen einen GOtt haben, der uns alsbald erhören soll, und alles geben, was wir begehren, nicht, der mit der Hülfe so lange verziehe.

11. GOtt verziehet aber nicht der Meinung, daß er die Seinen damit verderben oder verlassen wolle, sondern vielmehr, daß er seine Verheißung darnach auf das reichlichste erfüllen möge; wie St. Paulus von ihm zeuget, da er sagt Eph. 3, 20.: „Der überschwänglich thun kann über alles, das wir bitten oder verstehen.“

12. Darum will GOtt haben, daß wir solche Leute sein sollen, die mit dem Verzuge Geduld tragen können, und gewißlich darauf hoffen und harren, daß er überschwänglich und reichlich geben werde, was er verheißen hat. Eben dasselbe hat Jakob auch gelernt und erlangt, da ihm die Verheißung vom künftigen Samen ist gegeben worden und seinen Nachkommen das Land Canaan, und da noch bei seinem Leben Joseph zu solcher großen Ehre und Herrlichkeit ist erhoben worden.

13. Es ist aber das Fleisch so gar sicher und böse, daß es nicht allein an den Verheißungen GOttes verzagt, sondern auch die Drohungen verachtet. Denn dieselben werden auch verzogen, daß sie nicht so bald erfüllet werden; darum glaubt das Fleisch GOtt nicht, wenn er schon lange droht. Und wenn die Verächter und thörichten Menschen hören, daß noch ein Gericht und Strafe dahinten sei, so wider die Sünde gehen soll, o! sagen sie, das wird sich noch wohl verziehen! hätte ich diesem Geld zu zahlen! Jedoch will GOtt haben, daß man sich vor seinem Drohen fürchten und auf seine Verheißung harren soll, welches nicht anders geschehen kann denn im Glauben.

14. Die Welt aber fragt nach der keinem, ist ihr eben so viel, als pfiffe sie eine Gans an. Denn GOtt ist geduldig, und verziehet beide mit seinen Verheißungen und Drohungen, ehe er sie erfüllet; er eilet aber darum nicht, und will endlich, daß er so lange verzogen hat, erstatten, entweder daß er die Bösen um so viel härter strafen wird, oder den Frommen, Gottseligen so viel größere und reichere Wohlthat erzeigen. Denn er kommt doch endlich und kommt nur wohl.

15. Darum gebietet und lehrt die Schrift beiderlei, und sagt: Fürchtet GOtt und glaubet ihm; und wie der 147. Psalm V. 11. sagt: „Der Herr hat Gefallen an denen, die ihn fürchten, die auf seine Güte hoffen“; item in 2. Psalm V. 11.: „Dienet dem Herrn mit Furcht, und freuet euch mit Zittern.“ Und der Glaube muß glauben und fürchten, das unsichtbar ist, wie Noah die Strafe der Sintfluth fürchtet, die er nicht sieht, und auf eine Erlösung hofft, die er auch nicht vor Augen gesehen hat. Wir glauben an den Herrn Christum, den wir doch nicht sehen, und glauben, daß er am jüngsten Tage kommen werde, und werde uns vom Tode auferwecken und verklären, wie Paulus sagt Phil. 3, 20. 21., die Gottlosen aber, die Gott nicht erkennen, und dem Evangelio nicht gehorsam sind, strafen werden, 2 Thess. 1, 8. Also fürchten sich die Gottseligen vor den Drohungen und setzen ihr Vertrauen auf die Verheißung: die Gottlosen aber fürchten, oder glauben, oder hoffen nichts, fragen auch nach keinem Gott.

16. So hat nun Jakob die Rede, so er seinen Söhnen gethan, also angefangen: „Versammelt euch“, lieben Söhne, und höret mir zu, ich will euch lehren, was in künftigen Zeiten geschehen wird. Ich mache jetzt mein Testament, und wollte gern haben, daß ihr dasselbe immer möchtet in frischem Gedächtnis behalten, daß ihr den Herrn, den rechten wahren Gott, fürchtet. Ich will euch predigen und sagen von göttlichen Drohungen und Verheißungen, daß ihr einen Gott haben möget, den ihr fürchten und auf den ihr hoffen sollt. Das meint Jakob damit, daß er seine Söhne so aufmerksam macht, auf daß er sie behalten möge in der Furcht vor der Drohung und im Glauben an die Verheißung Gottes. Als wollte er sagen: Ich weiß wohl, daß viele von euch werden böse ungehorsame Ruben werden und wird ein loser Bübel unter euch sein, die werden abgöttisch und Verächter Gottes sein, der mit seinem Zorn droht und auch Gnade verheißt: wie Mose auch, desgleichen Samuel und David und viel Andere mehr von ihren Nachkommen geweissagt haben. Jedoch werden unter diesen meinen Nachkommen Etliche sein, die Gott fürchten und ihm vertrauen werden. Um derselben willen wird dies gesagt. Die Andern verachten diesen Gott, der mit seiner Hülfe also verzogen, ja, sie hätten den Gott nicht angesehen. Darum beteten sie aller Heiden Götter an, wie 2 Chron. 36, 14. gesagt wird. Und die Bücher der Richter und Könige zeugen auch, daß sie die Götzen gleichsam gehäuft haben über das ganze Land.

17. Wie wir im Papstthume auch gethan haben, da wir es noch viel gröber gemacht, und Alles nur mit lauter Abgötterei der verstorbenen Heiligen gehäuft und erfüllt haben. Einer hat Erasmus geehrt, daß er ihm Geld und Gut bescheren sollte; der Andere Margaretha, daß sie den Weibern in Kindesnöthen helfen sollte; und die Jungfrau Maria ist von jedermann geehrt worden als eine Mittlerin und Nothhelferin in allen anliegenden Nöthen. Niemand hat Hülfe und Erlösung von Gott hoffen und zu erlangen erwarten wollen. Ein Weib, die gebären wollte, begehrte nur gegenwärtige Hülfe, darum rief sie nicht Gott an, sondern Margaretha. Die Bauern aber haben Christophorus, Georg und viele andere Heilige mehr angerufen.

18. Also verachtet das Fleisch Gott mit seinem Drohen und da er Errettung verheißt. Denn dieweil er die Hülfe also verzieht, und damit etwas lange ausbleibt, darum wird er verachtet. Es will sich niemand daran gewöhnen, daß er seinen Glauben etwas übe; sondern die Leute wollen ohne Glauben leben, und der Güter, so gegenwärtig und vorhanden sind, genießen. Den Bauch wollen sie voll haben, aber die gewisse Verheißung lassen sie fahren; denn wiewohl sie von unsichtbaren Dingen redet, werden doch dieselben endlich gewiß kommen und erfüllt werden.

19. Derhalben erweckt nun Jakob seine Söhne mit höchstem Fleiß, seine Weissagungen, die er ihnen zuvor verkündigen wird, zu erkennen: daß sie aus den Drohungen, so den drei Söhnen, Ruben, Simeon und Levi, vorgehalten werden, den Zorn Gottes wider die Sünde verstehen lernen; aus den Verheißungen aber, so den Andern gegeben worden, Gottes Gnade und Barmherzigkeit erkennen mögen.

20. Das ist die vornehmste Lehre und der Inhalt dieses Capitels, warum Gott drohe und verheißet, und wohin er das wolle gemeint und verstanden haben, nämlich, daß er damit den Glauben an sein Wort üben oder prüfen möge. Und wer mit Gott handeln will, der muß das lernen, „daß er nicht lebe vom Brod allein, sondern von allem, das aus dem Munde des Herrn gehet“, 5 Mos. 8, 3. Wenn es am Brod mangelt, soll man darum keinen fremden Gott anrufen, sondern soll das Herz stärken mit dem Glauben an Gottes Wort, und also sagen: Gott hat mir verheißen, er wolle mein Gott und mein Herr sein: will er mich nun lassen Hungers sterben, so mag er es thun, ich will doch auf ihn hoffen.

21. Denn es muß der Glaube und die Furcht Gottes in den Herzen der Menschen sein. Denn Verheißung und Glaube gehören zusammen, und Drohung und Furcht gehören auch eins zum andern. Verheißung ist nichts, wenn nicht der Glaube dabei ist, und wiederum ist auch kein rechter Glaube ohne Verheißung; wie auch keine Furcht ist, wo keine Drohung ist. Aber Gott verziehet mit beiderlei, daß er uns versuche und prüfe. Und dieweil die Welt diese Versuchung nicht leiden oder ertragen will, verachtet sie also beiderlei, und fürchtet den nicht, der da droht, und glaubt dem nicht, der die Verheißung gibt.

\subsection{II}

\begin{quote}
\textbf{V. 3. Ruben, mein erster Sohn, du bist meine Kraft, und meine erste Macht, der Oberste im Opfer, und der Oberste im Reich.}
\end{quote}

22. Diese Rede ist voll Drohung, wiewohl sie auch etwas Verheißung und Segen in sich hat; aber es ist mehr Fluch darin. Und dies alles wird geredet den Andern zum Exempel und sie damit zu schrecken. Denn wiewohl Ruben die Sünde vergeben war, dieweil er darüber Leid genug getragen und Buße gethan hatte, dazu auch sich mit höchstem Fleiß bemüht, daß er Joseph hätte erretten mögen, und endlich nach der Sünde, so er begangen, den Vater wiederum in großen Ehren gehalten hatte: doch, dieweil die Sünde so groß und greulich gewesen, war es nicht nütze oder gut, daß sie hätte sollen ungestraft bleiben; sondern er hat müssen gestraft werden den andern Nachkommen allen zum Exempel. Denn es hätten sonst die Nachkommen sagen mögen: Hat Ruben das thun mögen und ist ungestraft geblieben, warum sollte es mir denn nicht auch frei sein und zu thun gehen?

23. Nun war es aber eine große und greuliche Sünde, die Blutschande mit seines Vaters Weibe, damit er das ganze Haus und Geschlecht geschmäht hat, dazu seinen Vater mit den andern Brüdern und Söhnen der Bilha. Was könnte doch Unbilligeres und Schändlicheres gesagt werden, denn daß er seines Vaters Weib, dieweil derselbe noch am Leben ist und da seine andern Söhne auch noch gelebt, geschändet hat? Dieweil es nun so gar eine schändliche, böse und unerhörte That war, hat Ruben dafür nicht müssen ungestraft bleiben, und wird darum an diesem Orte eine sehr scharfe und harte Drohung gesetzt.

24. Der Vater rückt ihn auf und wirft ihm vor alle Ehre und Herrlichkeit, so er zuvor gehabt, und hält ihm eine Strafe vor, von der er haben will, daß sie bei den Nachkommen bleiben soll: nicht daß Gott die auch verderben oder verdammen wolle, so dies greuliche und abscheuliche Exempel sehen würden; sondern daß sie Ruben also sollten als zum Exempel des schweren Gerichts und Zorns Gottes haben, so Andern zum Schrecken vorgestellet sein sollte.

25. Das hat also geschehen müssen, vornehmlich in dem Volk unter dem Gesetz. Denn wiewohl unter demselbigen Volk auch Barmherzigkeit und Vergebung der Sünden ist, so ist doch daselbst nicht lauter Barmherzigkeit, auch nicht allein das lautere Reich des Evangeliums und der Gnade; sondern es ist auch daselbst ein Stück des politischen oder weltlichen Regiments, da man auch Exempel der Strafen haben muß. Es muß daselbst Meister Hans das Schwert führen, Galgen und Rad brauchen, Andern zum Schrecken und zur Lehre, wenn auch schon die Sünde vergeben wird. Wie einem Diebe die Sünde auch verziehen wird, er muß aber doch an den Galgen gehenkt werden. Die Sünde wird zwar denen, so am Leibe gestraft werden, von Gott vergeben; der Stockmeister aber oder Scharfrichter vergibt es ihnen nicht, daß er ihnen die Strafe, so im Recht bestimmt, nicht anthun sollte; Meister Hans vergibt es ihnen nicht, sondern thut ihnen ihr Recht; wie St. Paulus sagt Röm. 13, 4.: „Die Obrigkeit trägt das Schwert nicht umsonst, sie ist Gottes Dienerin, eine Rächerin zur Strafe über den, der Böses thut.“ Und doch wird der Dieb oder Todtschläger nicht verdammt und andere Missethäter auch nicht, wenn sie Buße thun und an Jesum Christum glauben. Und nach dem Tode fühlt oder empfindet der Dieb die Schmach nicht, daß er ist an den Galgen gehängt worden; aber die

\quellseite{494}

Nachkommen sollen die Schmach anschauen, daß sie dabei denken: Wirst du stehlen, so wirst du auch also bezahlen müssen.

26. Dies ist eine nötige Lehre, die man durchaus in der Welt behalten muß. Darum sagen sie vom Kaiser Maximilian, daß so oft er bei einem Gericht hergezogen sei, habe er sein Haupt entblößt, und dasselbe mit diesen Worten gegrüßt und gesagt: GOtt grüße dich, du heilige Gerechtigkeit. Denn wo keine Strafe wäre, würden wir mit unserm Predigen und Vergebung der Sünden nichts ausrichten, und der Pöbel würde der Lehre und Predigt von GOttes Gnade und Barmherzigkeit schändlich zu allem Muthwillen und Sünden mißbrauchen.

27. Dieweil nun das Reich dieses Volks auch politisch oder weltlich war, hat eine sonderliche Strafe müssen gesetzt werden auf eine solche schändliche böse That, den Andern zum Schrecken, daß sie wissen sollten, daß auch ein GOtt sei, der da droht und straft, und nicht allein Gnade verheißt. Die Kirche hat ihre Strafe auch, nämlich, den Bann; wiewohl man darnach weniger fragt denn nach weltlicher Strafe. Also ist nun Ruben um seiner Sünde willen den Andern wie zum Spiegel vorgestellt worden, daran sich die muthwilligen ungehorsamen Ruben und die ihre Eltern schmähen und unehren, spiegeln und beschauen sollen, daß man also damit diesem sehr schändlichen bösen Aergerniß begegnen möchte. Darum hat GOtt nicht haben wollen, daß es im weltlichen Regiment ungestraft bleiben sollte. Und wenn etwa die Obrigkeit mit der Strafe verziehet und säumig ist, werden doch die Uebelthäter von GOtt gestraft.

28. Also redet aber nun der Vater Ruben an: Du solltest der Oberste gewesen sein, wirst es aber nun nicht sein, und du wirst auch das nicht haben, dazu du geboren bist: ich will dir das nehmen und dich alles dessen berauben, was dir von Rechts wegen und der Geburt halben gebühret hätte. Das ist ein sehr harter Fluch: der Größeste unter den Brüdern wird der Allerleinste. Die Geburt hat dir diese Herrlichkeit gegeben, daß du solltest mein erstgeborner Sohn, meine Kraft und Macht sein, und dazu der Oberste im Reich: aber die Sünde hat dir dies alles genommen und dich dessen gar beraubt.

29. Das hebräische Wort on heißt bisweilen Schmerz; als, droben im 35. Cap. V. 18. hat Rachel ihren Sohn benoni genannt, das ist, einen Sohn meines Schmerzens, darum, daß es ihr in der Geburt schwer gegangen war. Aber an diesem Orte bedeutet es etwas Anderes, wie auch 5 Mos. 21., da Mose von den erstgebornen Söhnen zweier Weiber ein Gebot gibt, nämlich, der Einen, die der Mann lieb hat, und der Andern, die feindselig ist. „Er soll“, sagt daselbst V. 17. Mose, „den Sohn der Feindseligen für den ersten Sohn erkennen, daß er ihm zwiefältig gebe alles, das vorhanden ist; denn derselbe ist seine erste Kraft.“ Da redet Mose von keiner Strafe oder Schaden, sondern von Kraft oder Macht. Darum ist das der Erstgebornen Ehre und Ruhm, daß sie genannt werden die erste Kraft oder Macht des Vaters.

30. Das deuten nun Andere etwas unzüchtig, welches mir gar nicht gefällt, nämlich darum, daß im ersten Jahr des Ehestandes die Liebe pflegt am allerhitzigsten zu sein, und soll derhalben die erste und vornehmste Kraft darin sein, Kinder zu zeugen. Ich halte aber dafür, daß der erstgeborne Sohn darum die erste Kraft genannt werde, dieweil er der Erste ist. Die andern Söhne sind auch eine Kraft, dadurch das Geschlecht an Kraft und Macht gemehret wird. Denn der Vater ist alsdann stark, wenn er viel Söhne hat, die alle seine Stärke oder Kraft sind, und sein Haus wird stark von wegen der vielen Kinder. Aber der erste Sohn ist die erste Kraft; wenn der geboren ist, so denken die Eltern bald, daß sie etwa ein Haus bauen mögen und eine Nahrung zuwege bringen, und wie sie ihr Haus und alle Erben oder Nachkommen nähren möchten. Wem es aber also gefällt, daß er der andern Deutung oder Meinung lieber folgen will, nämlich von dem Wert der Kinderzeugung und von der hitzigen Liebe der jungen Eheleute, der mag es meinethalben wohl thun.

31. Ich halte aber dafür, daß diese Deutung eigentlicher sei und sich besser hierher reime. Denn Jacob will also sagen: Meine Kraft hat mit dir angefangen; du bist die erste Säule meines Hauses. Da ward ich ein Hausvater: das fügest du an. Siehe aber nun zu, was du verloren hast. Solcher großen Ehre und Herrlichkeit wirst du jetzt beraubt um deiner Sünde willen. Du hast nur den Vorzug der Zeit halben, mein Haus aber wird durch dich weder gebaut noch gebessert; denn du bist von der Herrlichkeit der Erstgeburt abgefallen. Denn du warst nicht allein mein Erstgeborner oder meine Kraft, meine erste Säule, sondern warst auch ein Herr über alle deine Brüder und dir gehörten zwei Theile am Erbe; und das zweierlei Rechts wegen, nämlich, daß du der Oberste warst im Opfer (das hat der lateinische Dolmetscher gegeben: potior in donis: der Vornehmste in den Gaben, nämlich, welche GOtt gebracht werden. Lyra aber versteht es von den Gaben, so die Menschen einander zu geben pflegen) und um des Regiments und Herrschaft willen. Denn das hat beides dem Erstgebornen gebühret; wie dasselbe Recht geboten wird 5 Mos. 21, 17. Darum sind ihm zwei Theile zugefallen von den Gütern des Vaters und die andern Brüder haben ihm gehorchen und unterthänig sein müssen; er aber war größer, heiliger und reicher denn die Andern alle.

32. Derhalben ist dies die rechte Beschreibung der Erstgeburt, nämlich, was dieselbe sei, da er sagt: „Der Oberste im Opfer, und der Oberste im Reich“, das ist, der Erstgeborne hat die Herrlichkeit des Priesterthums, die Verwaltung und Regiment über die Lehre und Wort, und dazu über die Polizei. Das ist wahrlich eine treffliche große Ehre gewesen, und die Juden haben diese Ordnung festhalten müssen, wiewohl GOtt derselbigen nicht allezeit gefolgt hat. Denn er ist frei vom Gesetz und ist ein Herr über alle Gesetze. Welche Exempel doch die Ordnung, von GOtt eingesetzt, nicht aufheben, sondern GOtt hat sie selbst also ausgenommen und hat dasselbe Gesetz niemand übertreten dürfen. Was aber unser Herr GOtt anders macht, das kann ihm niemand wehren.

33. Im 1. Buch der Chronika am 6. Cap. V. 1. wird erzählt, daß die Erstgeburt Rubens den Kindern Josephs geschenket worden; wie wir droben, Cap. 48, V. 5., gehört haben, da Jacob zu Joseph sagt: „Deine zween Söhne, Ephraim und Manasse, die dir geboren sind in Egyptenland, ehe ich herein kommen bin zu dir, sollen mein sein, gleichwie Ruben und Simeon.“ Also ist Joseph ein doppelt Kind, hat so viel empfangen wie sonst der Andern zwei.

34. Diese Weise aber und Ordnung mit der Erstgeburt ist wohl unter dem Gesetz also gesetzt gewesen, aber im Reiche Christi hört das Gesetz auf. Denn unser Herr, der rechte erstgeborne Priester und König, hebt das Gesetz auf. Aber in der Welt, wenn daselbst viele Söhne sind, erbt der jüngste des Vaters Haus nach dem bürgerlichen Rechte.

35. Es scheint aber dies alles, so Jacob geredet, kurz abgebrochen zu sein. Denn er gebraucht solcher Worte, die in der Grammatik abstracta genannt werden, welche derhalben auch etwas dunkel sind. Bei uns würde es deutlicher sein, wenn er es im concretum, wie man es in der Schule nennt, ausgedrückt hätte. Er will aber so viel sagen: Die Ehre und Herrlichkeit hätte dir gebühret; du solltest mein Erstgeborner, meine Kraft und erste Macht sein, der Oberste im Opfer und der Oberste im Reich, das ist, du warst Priester und Fürst.

36. Ich will aber den Hebräisten ihren Zank lassen; denn ich weiß mir nicht zu helfen und mich los zu machen aus der Verwirrung, so die Grammatiker einführen. Sie ziehen das Wort seeth auf mancherlei Weise, Einer daher, der Andere dorthin. Welches Wort ohne Zweifel seinen Ursprung hat von dem Wort nasa, das eine sehr weite Bedeutung hat. Denn es heißt: nehmen, annehmen, tragen, aufheben, betrügen, vergeben; als, im 8. Psalm V. 3.: „Du hast eine, os, Macht zugerichtet“ (das hebräische Wort, os, heißt Kraft, Stärke, Reich), und im Propheten Daniel Cap. 11, 38. maussim und sonst oftmals an andern Orten mehr: und bedeutet eine solche Stärke, welches eine Gewalt ist zu regieren und zu herrschen. Droben im 4. Cap. V. 7. steht: „Wenn du fromm bist, so bist du angenehm.“ Da steht im Hebräischen seeth, das heißt, Vergebung. Item, von David steht geschrieben: „Ich habe ihn erwählet, daß ich habe ihn erwählet, daß ich ihn erwählet, daß ich ihn erwählet, daß ich ihn erwählt, daß ich ihn erwählt, daß ich ihn erwählt

\quellseite{495}

„Gottes Lamm, welches der Welt Sünde trägt.“ Und droben, Cap. 45, 27., ist von Joseph gesagt worden: „Er hat seinem Vater Jakob Wagen gesandt, ihn zu führen.“

37. Derhalben soll man fleißig darauf Acht geben, ob dies Wort gesetzt werde zu dem Wort „Sünde“, oder zu einer leiblichen Sache, als, Augen, Athem, Hand, Stimme etc., als, levare oculos, heißt, die Augen über sich werfen; levare vocem, heißt, laut reden oder schreien; capita levare, die Häupter aufheben, heißt, fröhlich sein; item, levare faciem vel assumere, das Angesicht erheben, heißt so viel, als, ehren, ansehen; wie im 82. Psalm V. 2. steht: „Wie lange wollt ihr unrecht richten, und die Person der Gottlosen vorziehen?“ das ist, ehren? als wollte er sagen: Ihr sollt den Sünder nicht ansehen. Also steht im andern Gebot: Du sollst den Namen Gottes, deines Herrn, nicht unnützlich führen oder annehmen.

38. Darum soll man auf die Umstände und Construction Acht geben, wie die Worte zu einander gesetzt sind. Wie es sich damit in andern Sprachen auch verhält, vornehmlich in der griechischen und deutschen Sprache, da oftmals einerlei Wörter, nachdem sie mancherlei Weise zu andern Wörtern gesetzt werden, auch mancherlei Bedeutung haben. Als, wenn man sagt in der lateinischen Sprache: gerere rem, eine Sache ausrichten; gerere bellum, Krieg führen; gerere magistratum, das Regiment führen; morem gerere, Einem zu Willen sein etc. Derhalben, wenn dies Wort sehn zu dem Wort „Sünde“ gesetzt wird, so heißt es so viel als, verzeihen oder vergeben und die Sünde wegnehmen. Als, Jerusalem wird genannt werden, neseth avon, das ist, Vergebung der Sünde oder Missethat, und Ps. 32, 1. steht: „Wohl dem, dem die Uebertretungen vergeben sind“, das ist, der seiner Sünden überhoben ist, dem sie verziehen und vergeben sind. Diese Regel pflege ich zu gebrauchen, daß ich die Worte nach der Construction und nachdem sie zu einander gesetzt sind. Ich weiß, daß etwas Anderes angezeigt werde, wenn ich sage von einer leiblichen Sache, levare saccum, den Sack aufheben, und daß ein Anderes sei, levare faciem, Einen fröhlich ansehen; wie 1 Kön. 2.: Non potero levare faciem meam ad Joab, fratrem tuum: Ich dürfte nicht vor sein Angesicht kommen.

39. Es soll sich aber jedermann befleißigen, und vornehmlich die, so mit der heiligen Schrift umgehen, daß er einen gewissen einfältigen Verstand der Schrift haben möge, und daß er nicht hin und wieder umberlaufe und irre gehe, wie der Juden Rabbinen und im Papstthume die Scholastiker gethan haben, und wie die Juristen, so im Rechte lesen, auch thun, die sich für und für bemühen mit den Worten, die mehr denn einerlei Bedeutung haben. Darum ist die Theologie der Scholastiker Anderes nichts denn ein bloßer Wahn. Man muß, sagen sie, also lehren und glauben, oder aber auf eine andere Weise. Eben also geht es auch zu in den Rechten und in der Philosophie. Ich erinnere mich, daß ich einst den Accursius gelesen habe, der an einer dunkelen Stelle nicht weiter konnte und nun verschiedene Meinungen aufzählte: Es könne entweder so verstanden werden, oder so, oder anders; er wisse es nicht, wie er selbst freimüthig bekannte. Und gefiel mir an dem Juristen wohl, daß er dasselbige so frei und aufrichtig bekannte. Denn wo man also mancherlei Meinung einführt und anzieht, das ist doch nichts Anderes, denn daß man es nicht weiß oder recht verstehet.

40. Wer aber ein Ding recht verstehet, zweifelt nicht, bringt auch nicht so mancherlei Meinungen herein, sagt nicht, man muß es so oder so verstehen. Darum soll man sich mit höchstem Fleiß vor den Rabbinen hüten; denn sie hindern das Studium der heiligen Schrift, sie befleißigen sich nicht, daß sie die Worte einerlei Weise und recht eigentlich deuten oder auslegen. Und wenn ich solche ungewiße Schrift lesen soll, darin die Worte so mancherlei bedeuten, will ich lieber Virgil und Ovid lesen. Den mancherlei Opinionen aber und Meinungen, so die Rabbinen in der Schrift haben, bin ich feind und verabscheue sie gar.

41. Darum halte ich dafür, daß an diesem Orte der rechte und wahrhaftige eigentliche Verstand sei, als wollte Jakob zu Ruben sagen: Du bist der Erstgeborne, meine Kraft und meine erste Macht, nicht an meiner Person, sondern die Kraft meines Hauses, der Anfang meiner Nachkommen, an der das Priesterthum und Reich anfangen soll. Er handelt oder redet nichts von leiblicher Kraft, daß sich Jakob vier und achtzig Jahre sollte enthalten haben, und hätte nie gefühlt, daß Samen von ihm gekommen wäre, bis er sich zu Lea gelegt und dieselbe mit Ruben ist schwanger geworden. Denn wozu dient es, daß man solche Dinge an diesem Ort anführt oder glaubt? Ein junger Knabe, der stark und gesund ist, kann schwerlich den Samen bei sich behalten. Und Mose hat es selbst auch nicht thun können, der es doch im Gesetz verbietet. Es hat sich auch Ambrosius dessen nicht enthalten können, der in einem Hymnus, so man in der Kirche singt, betet: Ne polluant corpora etc., das ist, daß unsere Leiber nicht beflecket werden, ja, das noch wohl mehr ist, es begegnet solches bisweilen auch denen wohl, die Ehemänner sind. Denn mit der Natur der Menschen ist es also beschaffen, die solche Leiber haben, daß sie zur Kinderzucht tauglich und geschickt sind.

42. Derhalben hat Jakob gesehen auf die Kraft, dadurch sein Geschlecht würde ausgebreitet und gemehret werden. Als wollte er sagen: Meiner Nachkommen wird unzählig viel werden, davon bist du der erste König und das Haupt, welchem die Ehre dieses so großen Geschlechts gebühret, und bist nicht allein der Anfang des Volks, das von mir herkommen wird, sondern du hast auch die Herrlichkeit, seeth, des Aufhebens, davon ich gesagt habe, daß es zum priesterlichen Amte gehöre. Denn wir wissen, daß im Alten Testamente mancherlei Opfer und Schlachten mit dem Viehe gewesen, dadurch Gott ist versöhnet und die Sünden sind vergeben worden, nicht aus Kraft der Opfer, sondern durch den Glauben an das zukünftige Opfer.

43. Darum hat Jakob so viel sagen wollen: Die Menschen sollten mit Gott versöhnet werden durch dein Opfer, das wegen des Rechts der Erstgeburt auf dich gekommen war: „Du bist der Oberste im Opfer, und der Oberste im Reich“, du bist das Haupt und der Anfang meines Geschlechts: die Vermehrung meines Hauses hebt sich von dir an. Ueber das hat dir auch gebühret die priesterliche Würde, daß du solltest ein Prediger sein, der die Leute mit Gott versöhnt. Derhalben siehe nun zu, welch ein groß Gut du verloren hast.

44. Das ist wahrlich eine sehr ernste Predigt und eine greuliche Strafe, damit er Ruben hart gescholten hat. Ich hätte mich darüber zu Tode geheult, wenn mich Jakob also angeredet hätte. Er hält ihm alle Herrlichkeit vor, die er hätte haben sollen, als, daß er seine erste Kraft sein sollte, dazu die Ehre des Priesterthums und des Reichs sollte gehabt haben, die herrlich und trefflich groß war, nicht um des Kaisers Augustus oder Alexander des Großen, oder Scipios, sondern um des Hauses willen des Patriarchen Jakob, der die Verheißung der ewigen Seligkeit und vom künftigen Messias gehabt.

\begin{quote}
\textbf{V. 4. Er führt leichtfertig dahin, wie Wasser. Du sollst nicht der Oberste sein. Denn du bist auf deines Vaters Lager gestiegen, daselbst hast du mein Bette befleckt mit dem Aufsteigen.}
\end{quote}

45. Hier machen uns die Rabbinen wiederum zu schaffen in der Grammatik, und wenn dieselbe gewiß wäre, könnten wir ohne alle Beschwerung den rechten Verstand herausnehmen. Aber sie machen sie selbst dunkel mit ihren Glossen und Punkten; wie denn dieser Text von ihnen mancherlei Weise zerrissen wird. Denn dieweil sie denselben wider den Geist auf einen fremden Verstand ziehen, haben sie darnach von Noth wegen gezwungene und fremde Auslegungen suchen müssen, welches eine recht teuflische Sophisterei in der heiligen Schrift ist.

46. Hieronymus hat es gegeben: Du bist ausgegossen wie Wasser; und hat darin den Rabbinen gefolgt, welches mir nicht gefällt. Denn er zieht oder deutet es wie auf ein Gefäß, daraus man das Wasser, so darin ist, gar bis auf den Grund ausgießt, daß auch nicht ein Tröpflein darin überbleibe; welches mit Del also nicht geschehen kann. Darum will ich meine Meinung sagen: Das hebräische Wort pachas heißt nicht eigentlich, ausgießen, sondern wild und unsanft einhergehen oder geführt werden; wie die Philosophen in den Schulen sagen von dem, das feucht ist: Humidum est male terminabile termino proprio, das ist: Was feucht ist, kann sich selbst nicht an seinem Ort inne halten; und von dem, das trocken ist, sagen sie: Siccum est male terminabile termino alieno, das ist: Was trocken ist, kann sich selbst inne halten und bedarf keines andern begrenzenden Körpers. Das Wasser kann selbst kein Ziel setzen, sondern wo es nur offen ist, fließt es aus, oder wo es nicht gefaßt ist, da hat es keinen freien Ausfluß, daß es hieher und dorthin fließt; das ist die Art und Natur des Wassers und aller flüssigen Körper. Aber wenn du schon einen Kieselstein oder sonst einen andern Stein in ein Gefäß wirfst, so hat er doch seinen eigenen Ort, da er von selbst bleibt.

\quellseite{496}

47. So heißt nun, sage ich, pachas, was wild und unbeständig ist, das dahin fährt. Also, sagt Jakob zu Ruben, bist du auch ein Maulaffe. Du bist groß wegen der Herrlichkeit, so auf dich gekommen war: aber du fährst dahin wie das Wasser, fällst gleichsam mit einem Sturm, als ein Blinder, auf deine böse Lust und Begierde; bist nicht beständig, bleibst ja nicht an deinem eigenen Orte, dahin du bestellt warst; thust wie das Wasser, das auch also unbeständig ist, und leichtlich dahin fährt, wo es nur ausbrechen kann.

48. Daher kommt es, daß im Propheten Jeremia und auch Zephania die Propheten leichtfertig und unbeständig genannt werden, die leichtfertig und unbedacht dahin machen, was ihnen nur in das Maul kommt, und das Volk auch wild und unbeständig machen; wie im Psalmhuhn das arme Volk jämmerlich ist umgeführt worden, daß es hin und her gelaufen ist, da zu Jakob, dort zu Catharinen oder Barbara etc. Da war auch alles leichtfertig, ungewis, unbeständig Ding, was man gelehrt und gethan hat.

49. Die aber das Regiment haben in der Polizei und auch in der Kirche, müssen nicht pachas sein; denn das steht ihnen nicht wohl an: sie sollen sich nicht wägen oder wiegen lassen von allerlei Wind der Lehre, Eph. 4, 14. Sie sollen nicht sein wie das Rohr, das der Wind hin und her wehet, sondern sollen fest an ihrem bestimmten Ort bleiben, wie ein unbeweglicher Fels, der weder vom Sturmwinde noch von Wasserwogen umgetrieben oder weggerissen werden kann. Die Elbe rauscht überhin, der Steinfels aber bleibt stehen, obwohl das schnelle Wasser, so vom Berge herab brauset, sonst alles mit hinweg nimmt und zu Boden wirft, was es nur auf dem Felde findet. Also soll auch ein Lehrer und Regent sein in der Kirche, im Hause und in weltlicher Polizei, der sein eigen Ziel und Maß habe, dabei er wisse zu bleiben und zu beharren, und daß er nicht mit einem Sturm, wie ein Blinder, bald dahin, bald dorthin laufe.

50. Also soll eine jegliche Obrigkeit sein mit sich selbst einig und beständig sein, und soll den Schmeichlern und dem, was die Hoffschranzen anbringen, nicht bald Gehör geben. Wenn ein Oberherr hier einem Schmeichler, dort einem andern folgt und sich an der Nase führen läßt, der wird kein guter Regent sein. Eben also wird auch gesagt Jer. 23, 32. von den falschen Propheten: „Sie verführen mein Volk mit ihren Lügen und losen Täubungen.“ Solche lose leichtfertige Streicher sind die auch, so des Papstes Lügen und lose Täubinge lehren; wie das an den Sophisten zu Löwen zu sehen ist. Und im Propheten Zephania im 28. Cap. V. 10. werden sie auch getroffen, da er sagt: „Gebet hin, gebet her; harre hie, harre da; hie ein wenig, da ein wenig.“ Derhalben ist es vonnöthen, daß wir beständig seien: wir sollen Constantine, Tapfere und Felsen sein, daß wir uns nicht mögen oder wiegen lassen von allerlei Wind der Lehre, Eph. 4, 14.

51. So sagt nun Jakob zu Ruben: Du warst der Erstgeborne; aber du hast die großen, herrlichen, wichtigen Dinge verachtet, bist leichtfertig in der Irre dahingegangen, wie das Wasser dahinfährt, hast dir nicht genügen lassen an deinem Ehestande. Das ist eine Meinung von diesem Text; denn ich kann dem ungewissen zweifelhaften und losen Wahn der jüdischen Rabbinen nicht beifallen. Es ist im ganzen Leben vonnöthen, daß man allezeit etwas Gewisses und Beständiges habe, daran man sich halte, vornehmlich aber ist solches vonnöthen in der Kirchenlehre; darum betet David im 51. Psalm V. 12.: „Gib mir einen neuen gewissen Geist.“ Wer nicht will beständig sein, der wird endlich das Amt verlieren, und das Wort auch, welches gewiß ist. Darum bittet David um einen gewissen Geist, der nicht zweifele oder wanke, nicht hin und her in der Irre gehe; wie Jakob sagt, daß in Ruben ein solcher ungewisser Geist gewesen sei, als wollte er zu ihm sagen: Du hast gehandelt wie ein leichtfertiger Narr; du hast deinen bösen Lüsten gefolgt, du fürchtest Gott nicht, du glaubst ihm nicht; es hat dir nicht gemangelt an der Gabe der Erstgeburt, sondern am Glauben hast es dir gefehlt; denn du hast Gott nicht geglaubt.

52. Darum setzt er nun dazu und sagt: „Du sollst nicht der Oberste sein.“ Da nimmt er ihm Beides, nämlich die Herrlichkeit des Priesterthums und auch des Regiments. Du sollst in meinem Hause nicht Priester oder Fürst genannt werden, welches dir nach der Natur wohl hätte gebühren sollen. Aber du sollst über deine andern Brüder nicht regieren, sollst auch des Priesteramts nicht pflegen, und sollst weder im Regiment noch in allen andern Händeln über die Andern gesetzt werden. Mose, 5 Mos. 33, 6., sagt im Segen Rubens also: „Ruben lebe und sterbe nicht, und sein Pöbel sei gering.“ Da hat er diesen Text ausgelegt: Er soll nicht gemehret werden, soll den andern Stämmen nicht gleich sein, sondern soll geringer sein denn sie.

53. Das ist zumal ein hartes Wort, und er redet in der Mehrzahl als von vielen Betten, macht es sehr hart. Du bist in meine Schlafkammer gestiegen und hast deines Vaters Lager eingenommen. Die Juden führen hier und an andern Orten mehr wunderliche unzüchtige Dinge ein, deren ich ungern gedenke. So viel will aber Jakob sagen: Die Schlafkammer meines Weibes Bilha war meine, ich sollte bei ihr im Bette liegen; du aber hast mich dessen beraubt. Und wird eigentlich damit nicht das Bett, sondern was im Bett geschieht, verstanden. Denn die Juden lügen, da sie sagen: Ruben habe seine Stiefmutter nicht geschwächt, sondern habe des Vaters Bett aus der Mutter Hütte geworfen. Es ist aber in Wahrheit eine schändliche Hurerei und rechte Blutschande gewesen. Davon hat Jakob droben stille geschwiegen; aber da er nun sterben soll, verschweigt er es nicht mehr.

54. Ueber das Wort chiel erregen die Rabbinen einen Zank, und machen aus einem transitiven Zeitwort ein intransitives, wie man es in der Grammatik nennt. Ich verstehe es als transitiv und setze dazu einen Accusativ und sage: „Du hast mein Bette befleckt.“ Wie auch 1 Chron. 6, 1. von Ruben gesagt wird eben mit denselben Worten: „Er hat seines Vaters Bette verunreinigt“ etc. Dasselbst ist es ein transitives Verbum, und muß an diesem Ort auch ein solches sein. Und daß wir aus dem Lateinischen gegeben haben „mit dem Aufsteigen“, ascendendo, das lesen sie aus, daß es heißen soll, ascendit, er ist aufgestiegen, das ist, wie sie es deuten, es ist nicht mehr mein Bett: dasselbe ist ja ein gar zu gezwungener Verstand. Ich aber habe zu solchem Verstande, der so mit Gewalt erzwungen wird, nicht Lust, und wenn ich ja nichts beibringen oder anzeigen kann, das sich zu einem Texte wohl schicke, bekenne ich lieber und sage frei heraus, daß ich es nicht verstehe.

55. Es läßt sich aber mit diesem Texte ansehen, daß es geredet sei, was Jakob hier sagt, durch die Figur eclipsis oder reticentia, so man aposiopesis nennt, nämlich, wenn die Rede kurz abgebrochen wird, oder daß man die Worte nur halb ausredet und etwas verschweigt; als, wenn Einer zornig wird, oder sonst von anderer Affecten wegen. Wie im Virgil steht: Quos ego! Sed praestat motos componere fluctus. Es ist diese Figur sehr häufig in allen Sprachen. Also sagt an diesem Ort Jakob auch zu Ruben: „Du bist auf deines Vaters Lager gestiegen, daselbst hast du mein Bette befleckt.“ Vor diesen Worten sind ohne Zweifel die andern Söhne alle erschrocken, und werden weinend gesagt haben: Ach! mein lieber Vater, wie sind doch das so scharfe Worte, die du redest; Lieber, ändere oder lindere doch das so schwere betrübte Urtheil, so du wider deinen Sohn gesprochen hast. Darum wollen wir diesem Texte also helfen durch die Figur Eclipsis. Denn die Auslegung der Rabbinen gefällt mir nicht, daß sie sagen: Es ist nicht mehr mein Bett.

56. Dies ist übrigens eine erschreckliche und sehr scharfe Predigt, in wenig Worte gefaßt, und Ruben wird ohne Zweifel mit großen Schanden und erschrockenem Herzen daselbst vor allen seinen Brüdern gestanden sein. Sein Vater Jakob hat bisher still geschwiegen; aber jetzt, da er sterben soll, läßt er wahrlich ein trauriges, scharfes und betrübtes Valete hinter sich. Denn dies ist wahrlich ein erschrecklich Urtheil wider Ruben, und wird doch nicht in einer Silbe irgend ein Trost dazu gesetzt, sondern beinahe in allen Buchstaben und Titeln nimmt er ihm das Leben und tödtet ihn doch gleichwohl nicht; wie Mose auch sagt 5 Mos. 33, 6.: „Ruben lebe, und sein Pöbel sei gering“; und er ist auch der allerschwächste Stamm gewesen, der nie etwas Herrliches hat ausgerichtet. Im Buche der Richter am 5. Cap. V. 15. wirft ihm Debora vor, wie er stolz und faul gewesen, da sie sagt: „Ruben hielt hoch von ihm, und sonderte sich von uns.“ Dieses Gebrechen hat Ruben auch gehabt.

57. Derhalben ist dies wahrlich eine schwere greuliche Strafe gewesen, so auf seine Sünde gefolgt. Es ist ihm ja die Vergebung der Sünden nicht versagt worden, wie wir droben gesagt haben, aber die Strafe ist gleichwohl noch geblieben: nicht daß der Vater daran Lust gehabt; denn er sagt ja mit klaren Worten: „Du bist meine Kraft“, als wollte er sagen: „Ich wollte zwar gerne, daß du diese erste Ehre behalten hättest; aber du hast sie verloren durch deine eigene Schuld.

\quellseite{497}

58. Lieber, wie kommt es denn, daß er gleichwohl mit seinem erstgebornen Sohne so gar hart umgeht? Antwort: Jakob ist ein Patriarch, dem große herrliche Verheißungen gegeben waren, daß er würde groß werden, und der viel Nachkommen haben sollte, und Ruben war der Vater eines Stammes. Darum ist es vonnöthen gewesen, daß die Züchtung also vorher ging. Denn Jakob hat sich besorgt, es möchten die Nachkommen, von denen er wußte, daß sie viel und gewaltig sein würden, diese Sünde seines erstgebornen Sohnes auch nachthun. Darum hat er ihn also gestraft den Kindern zum Exempel, sie damit abzuschrecken, daß sie nicht auch dergleichen Sünde begehen möchten. Denn es sind nicht alle Kinder der Verheißung, die nach dem Fleisch Kinder sind, haben auch nicht alle den Heiligen Geist empfangen, sondern der mehrere Theil des Volks ist ja ganz fleischlich geworden und hat sich allein des fleischlichen Vorzuges halben gerühmt.

59. Derhalben war es nützlich und gut, daß man ein solch Gedächtniß einer erschrecklichen Strafe gehabt, welches immer für und für gewähret. Sonst hätten sie mögen sagen: O es ist nichts daran gelegen, wenn ich schon bei meiner Stiefmutter schlafe; denn eben dasselbe hat unser Vater Ruben auch gethan. Also ist es um der Nachkommen willen vonnöthen gewesen, daß Jakob sich mit Worten so hart gegen Ruben hat vernehmen lassen, nämlich, daß sie sich vor solchen Blutschanden hüten, und darauf denken sollen, daß dieser Stamm so gedemüthigt und sehr geschwächt worden sei um der Blutschande willen, so Ruben mit seiner Stiefmutter angerichtet und begangen hatte.

60. Wiewohl es aber nun glaublich ist, daß die andern Brüder werden für Ruben gebeten haben, hat sich doch der Vater nicht wollen erweichen lassen, darum daß die Sünde so unsäglich groß; nicht daß er darum Ruben wollte verloren und verdammt haben, sondern daß er möchte selig werden und die Größe derselben Sünde erkennen lernen sollte, und daß unter demselben Volke das Gedächtniß dieser Strafe immer bleiben möchte, daß man sagte: Diesen greulichen Fluch hat Ruben verdient mit seiner großen Sünde. Also verkündigen wir auch Vergebung der Sünden den Ehebrechern, Dieben und Todtschlägern: aber damit gleichwohl der unbändige Pöbel und das gemeine Volk durch Furcht der Pein und Strafe etwas im Zaum gehalten und gezwungen werde, verbieten wir nicht, daß man sie nicht hängen oder tödten sollte, damit Andere nicht auch in allerlei Sünde und Schande fallen. Absalom hat noch schwerer gesündigt als Ruben; siehe aber, welch ein jämmerlich und betrübt Ende er genommen hat, da er so erbärmlicher Weise in seinen Sünden gestorben ist.

61. Darum sollen die Leute solche Strafe immer im frischen Gedächtniß behalten, daß sie daran für und für eine Lehre haben mögen, dadurch sie gewarnet werden, sich vor solchen Blutschanden und dergleichen andern Sünden mehr zu hüten. Denn Gott ist denselben feind und hat davor einen Greuel; wie denn dieser Stamm deshalb bei allen Nachkommen der allerärmste und verächtlichste gewesen ist. So ist nun kein gut Wort in diesem ganzen Segen, und möchte wohl eine maledictio, das ist, ein Fluch heißen.

\section[Zweiter Theil]{Zweiter Theil. Von dem Fluch, den Jakob dem Simeon und Levi ankündigt}

\begin{quote}
\textbf{V. 5. Die Brüder Simeon und Levi; ihre Schwerter sind mörderische Waffen.}
\end{quote}

62. An diesem Orte sind die Grammatici auch nicht eins mit einander über den Verstand und die Worte. Denn das Wort mecheroth deuten sie, erstlich, daß es Wohnung heiße; darnach, verstehen sie darunter ein Schwert. Wir wollen es also verstehen, daß Jakob habe sagen wollen: Simeon und Levi sind Brüder nicht allein von Natur, sondern auch in der Schalkheit; es ist Einer so fromm als der Andere; es sind Brüder, ja, Brüder, es sind Gefäße oder Werkzeuge der Ungerechtigkeit. Denn wenn du es wolltest verdolmetschen: ihre Wohnungen, wie es die Juden auslegen, das gibt gar keinen Sinn; sondern man gibt es viel besser also: Ihre Schwerter sind mörderische Waffen. Es haben ferner die Griechen das Wort macheara aus dem hebräischen hergenommen, wie auch das Wort framea hebräisch ist, und die Deutschen haben in ihrer Sprache viel hebräische Worte.

63. Derhalben verwirft nun Jakob diese Zwei auch, und setzt nun hier ein ander Exempel oder Strafe wider die Todtschläger, Vater- und Brudermörder; wie er zuvor wider die Blutschande auch eine Strafe gesetzt hat. Denn beiderlei Sünde ist in den Regimenten in Landen und Städten schädlich. Darum verflucht er diese Brüder um der Kinder willen nach dem Fleisch bei den Nachkommen, so hernach leben würden, daß sie ja solche Exempel nicht mit Gewalt dahin zögen, ihre Sünde und Schande damit zu vertheidigen. Jakob hat wohl gewußt, daß sie nicht alle würden Kinder werden nach dem Geist, sondern daß der größere Theil würde fleischlich bleiben; und dieweil unter denselben gleichwohl die Gemeinde hin und wieder mit Gesetzen und Weltordnung hat müssen bestellt und versehen werden, sind diese Exempel der Strafen dazu nöthig gewesen, daß sie sich der Blutschande und unrechten Todtschlags zu enthalten wüßten. Wie hernach zwischen Juda und Ephraim ein steter Krieg gewesen ist, die doch Brüder waren. Die Priester aber aus dem Stamme Levi haben Christum an das Kreuz gebracht, und, wie Etliche zeugen, soll Simeon Juda Ischarioths Vater gewesen sein; und die wollen haben, daß dies die Ursache sei, warum Simeon nicht auch mit in dem Segen Mosis aufgezählt worden ist.

64. Derhalben ist dem Patriarchen Jakob, da er jetzt hat sterben sollen, eingefallen, daß er gedacht hat: Vielleicht werden die Nachkommen Simeons und Levis auch blutdürstig und Vatermörder werden. Denn es folgen gemeiniglich die Kinder der Art und Natur ihrer Väter nach und wird also ihr Geist und Unart unter alle Nachkommen ausgearbeitet. Darum redet er nicht allein von diesen zwei Brüdern, sondern von allen ihren Nachkommen, als wollte er sagen: Sie werden auch Bluthunde werden wie ihre Väter. Wo sie aber durch den Heiligen Geist geändert werden, alsdann werden sie fromm und gottselig sein; wie droben Ruben etliche mal Zeichen rechter wahrhaftiger Buße von sich gegeben hat und Andere mehr. Jakob hat aber doch seine Prophezeiung oder Weissagung nicht geändert, hebt auch die Strafe nicht auf, ja, er macht ihre Bosheit greulich groß. „Ihre Schwerter“, sagt er, „sind mörderische Waffen“, die nur hauen, stechen und die Leute mit Gewalt todtschlagen. Also sieht er, daß auch die Nachkommen dergleichen werden würden, wo sie nicht vom Heiligen Geist gebessert würden; und darum straft er sie so hart, daß sie durch das Exempel dieser zwei Brüder ermahnt und gewarnet würden, und sich solcher grausamen Thaten und Todtschlages zu enthalten müßten und sich damit nicht also versündigten. Die Brüder Simeon und Levi, sagt er, sind verzweifelte böse Buben, die nur Lust haben zum Todtschlagen und unrechter Gewalt, daß ihre Schwerter Waffen sind, Mord und Tyrannei damit zu üben.

65. Und ist auch wahrlich dieser Segen oder vielmehr dieser Fluch endlich erfüllt worden in Caiaphas, Hannas und den andern hohenpriestern, die Christum und die Apostel heftig verfolgt haben. Denn diese zwei Stämme, Simeon und Levi, haben den ganzen Haufen der Juden vor Pilatus dazu bewogen, daß sie bitten sollten, Barabbas los zu lassen und Jesum zu kreuzigen.

66. Dies alles hat der heilige Patriarch so viel Jahrhunderte zuvor gesehen, und hat können schließen, daß aus diesen zwei Stämmen nichts Gutes kommen würde. Es sind ja wohl viel Propheten und heilige Männer von ihnen hergekommen, aber aus dem Heiligen Geist. Aber das väterliche Blut hat immer geschnaubt mit „Dräuen und Morden“, wie von Paulus gesagt wird Apg. Gesch. 9, 1. Und desselben willen war es nützlich, daß man also ein ewiges Gedächtniß dieser Sünden gehabt und was darauf für eine Strafe erfolgt ist.

67. Jakob sieht die Geschichten, so sich in seinem Haus zugetragen haben; er sieht auch die Exempel seiner Vorfahren an, was Abraham um Ismaels willen begegnet und was Isaak um Esau willen widerfahren ist. Derhalben hat er gedacht: Bei meinem Leben, der ich ein Vater bin der göttlichen Verheißungen, darin so reichlich die ewigen Güter angeboten werden; ja, sage ich, dieweil ich noch selbst ha- bei frisch und gesund bin und die Kirche regiere, geschieht solches: was wird wohl nach meinem Tode geschehen bei den Nachkommen?

\quellseite{498}

68. Und auf diese Weise redet Mose auch 5 Mos. 31, 27.: „Denn ich kenne deinen Ungehorsam und Halsstarrigkeit. Siehe, weil ich noch heute mit euch lebe, seid ihr ungehorsam gewesen wider den HErrn; wie viel mehr nach meinem Tode.“ Also hat Jakob auch gesehen, daß die Nachkommen würden zweierlei sein: ein Theil derselben, die da würden selig werden durch den Glauben an die Verheißung vom Heiland der ganzen Welt, der von Jakobs Fleisch herkommen sollte. Die sind wahrhaftig unter der Zahl der Kinder Israel oder Gottes gerechnet worden; denn sie haben zu der Ehre des Fleisches, daß sie Abrahams Samen waren, auch die Verheißung vom Messias gesetzt, und darum werden sie auch vom Propheten Jesaia die Uebrigen von Jakob genannt. Der andere Theil aber würde die fleischlichen Israeliten sein, die sich des Samens Abrahams und anderer Väter rühmen würden. Wie auch noch heutiges Tages derhalben die Juden stolzieren und gegen sich die andern Völker alle verachten, und verheissen sich selbst, daß sie Herren werden wollen über Himmel und Erde. Ja, höllisch Feuer auf den Kopf werden sie kriegen.

69. Darum straft sie Johannes der Täufer sehr hart, Matth. 3, 9. 10.: „Denkt nur nicht, daß ihr bei euch wollt sagen: Wir haben Abraham zum Vater. Ich sage euch: Gott vermag dem Abraham aus diesen Steinen Kinder zu erwecken. Es ist schon die Axt den Bäumen an die Wurzel gelegt. Darum, welcher Baum nicht gute Früchte bringt, wird abgehauen und ins Feuer geworfen.“ Item, daselbst weiter V. 12.: „Und er hat seine Worfschaufel in der Hand; er wird seine Tenne fegen und den Weizen in seine Scheune sammeln; aber die Spreu wird er verbrennen mit ewigem Feuer.“ Die Spreu auf der Tenne waren die Israeliten nach dem Fleisch; dieselben, sagt er, sollen in das ewige Feuer geworfen werden; welches zuletzt das Ende ausweisen wird. Eben dasselbe haben ihnen auch alle Propheten bis zum Ueberdrus eingeschärft; sie haben aber damit so gar nichts ausgerichtet, daß sie vom Volk und den falschen Propheten endlich darüber sind erschlagen worden.

70. Also ist heutiges Tages die Kirche auch versammelt aus Heiligen und Gottlosen. Aber wir können die Papisten deß nicht überreden, daß sie Simeon und Levi sind: sie wollen traun auch Abrahams Kinder sein und schreiben sich den Namen und Titel der Kirche zu.

71. Derhalben hat nun Jakob ein solch greulich Gedächtnis dieser Sünde und der Strafe seiner Söhne aufrichten wollen, daß sie etwas hätten, ihren ehernen Stirnen und eisernen Nacken vorzuhalten. Aber sie haben durch keinerlei Predigt oder Exempel können erweicht oder gebessert werden; ja, sie behalten noch bis auf den heutigen Tag die eiserne Stirne; wie unsere Papisten auch also verstockt und verhärtet sind. Alle Propheten sind darüber erschlagen worden.

72. Darum denkt Jakob also: So noch bei meinem Leben meine Söhne solche böse Thaten haben thun dürfen: Ruben ist auf seines Vaters Lager gestiegen, welches wahrlich eine böse unerhörte Uebelthat ist (denn was kann das für eine Ehre sein, so man dem Vater und Mutter erzzeigt, eine solche unbillige That begehen, seinen Eltern zu greulicher Verachtung und Schmach, da Jakob noch selbst am Leben und gegenwärtig gewesen und solche Sünde und Schande selbst angesehen hat?), was sollten wohl die Nachkommen nicht thun? Wenn er von einem heidnischen oder cananäischen Vater geboren wäre, so hätte er doch durch Anleitung der Natur demselben billig Ehrerbietung erzeigen sollen: aber nun thut er das mir, seinem Vater, der ich göttliche Verheißung habe und nach derselben meine Kirche und Gemeinde lehre und regiere. Wiewohl ich ihm nun die Sünde verziehen und vergeben habe, so muß doch eine Strafe darauf gelegt werden um der Nachkommen willen, so hernach leben werden, deswegen, daß die Sünde so gar scheußlich und groß ist. Denn Ruben hätte seinen Vater nicht schändlicher verachten und mit Fäusten treten können als eben mit dieser Blutschande. Darum mag er nun dagegen diese Strafe tragen zum Exempel und Andere damit zu erschrecken, daß sie lernen Gott gehorsam sein und ihren Eltern; oder aber sie sollen wissen, daß sie dem Scharfrichter und Teufel selbst unterthänig sein werden.

73. Und können diese Exempel den Leuten nicht genugsam eingeschärft werden zu allen Zeiten, sintemal da Jakob, der Patriarch, selbst gelehrt und geprediget hat, seine Söhne so schändlich gezüchtiget haben. Und wir dürfen nicht denken, daß er solches gesehen und zuvor verkündiget habe ohne großen Schmerz. Ihr wißt euch wohl zu erinnern, wird er gesagt haben, wie mit so großem Fleiß und Treue ich euch ermahnt habe, daß ihr solltet an Gott glauben und euere Eltern in Ehren halten; und ihr habt auch sonst viel andere heilsame Predigten gehört von geistlichen Dingen und auch die weltliche Polizei belangend. Aber diese göttlichen Predigten alle unangesehen und gar hintangesetzt habt ihr euch weder des enthalten, daß ihr euren Vater nicht geschmäht, noch auch, daß ihr euren Bruder nicht getödtet hättet. Ich verdamme euch zwar nicht ewiglich, sondern verkündige euch nur die Strafe Gottes, die über euer Haus und Geschlecht gehen wird, damit Andere von gleichem Sünden abgeschreckt werden mögen.

74. Und daraus ist zu sehen, welch ein eitel und vergeblicher Ruhm das sei, daß sich die Juden gerühmt, daß sie nach dem Fleisch von den Vätern hergekommen waren, wo der Heilige Geist und die Gnade nicht dabei gewesen. Denn siehe doch, die Nachkommen Jakobs an, von welchen doch Christus geboren ist, wie wir im Segen Judas hören werden.

75. Und diese Weissagung oder Fluch ist nicht vergeblich gewesen, sondern ist über das ganze Geschlecht Simeons und Levis gegangen, Wenige ausgenommen, als Mose, Aaron, Samuel und etliche Andere, die sich von den Gottlosen haben abgesondert; und in den Geschichten der Apostel Cap. 6. V. 7. steht geschrieben, daß ein großer Haufe Priester sich zum Gehorsam des Glaubens bekehret worden. Sonst klagen die Propheten hin und wieder über die Tyrannei der Priester, die wider die frommen gottseligen Lehrer gewüthet haben und die immer nach ihrem Blut gedürstet hat. Das sind alle Simeons und Levis Brüder und Nachkommen gewesen.

\begin{quote}
\textbf{V. 6. Meine Seele komme nicht in ihren Rath, und meine Ehre sei nicht in ihrer Kirche.}
\end{quote}

76. Jakob redet hier nicht von der Seele einer einzelnen Person, sondern von allen Nachkommen. Denn er nennt seine „Seele“ und seine „Ehre“ das ganze Volk und alle Heiligen, die nicht allein nach dem Fleisch, sondern auch nach dem Geist rechte wahre Israeliten sein würden; wie sie Paulus beschreibt Gal. 6, 16. Er wußte wohl, daß ihm von Gott verheissen war, daß er viel und herrliche Nachkommen haben sollte. Denselbigen Nachkommen befiehlt er, sich vor diesen Buben zu hüten. Als wollte er sagen: Ich ermahne und bitte euch, daß ihr euch ja wollet hüten vor dem Exempel und Rath dieser zwei Brüder: was meines Leibes und Seelen ist, das hüte sich vor ihnen. Und das will ich mit diesem meinem letzten Willen bezeugt haben, daß meine Seele in ihren Rath nicht kommen soll, soll sich auch mit ihrer Kirche und Versammlung nicht vereinigen. Denn also steht im Hebräischen: „meine Ehre“, nicht, die meine Person allein belanget, sondern des ganzen Volks, das von Jakob würde herkommen. Alle, die von mir jemals geboren werden sollen, die wolle Gott ja bewahren und erhalten, daß sie mit solcher Kirche sich nicht vereinigen und ihr Ding nicht bewilligen; sondern mein Volk und Geblüt soll darauf denken, daß es bei der göttlichen Verheißung und Geboten bleibe und vom Glauben nicht abweiche, noch auch von den rechten Werken des Glaubens sich abwende, und sich zu den bösen verlornen Werken kehre, wie diese Tyrannen und Todschläger, ja, die greulichen Vater- und Brudermörder, Simeon und Levi, Werke begangen haben.

77. Wie diese harte Rede und Bitte dazumal nöthig gewesen für das Geblüt und Nachkommen Jakobs: also soll man auch dieselbe zu allen Zeiten und in der ganzen Welt immer wiederholen und fleißig treiben. Und wir sollen Gott bitten, daß er uns ja nicht in solche greuliche Laster fallen lasse. Gott behüte uns, daß wir ja nicht Brudermörder werden! das sind ja viel zu große Stücke. Es sind wahrlich in dem Volk greuliche Fälle und unerhörte Exempel gewesen des Vatermords und Blutschande, dergleichen man schwerlich in den Historien der Heiden findet. Es ist der Teufel. So groß ist die Gewalt und Tyrannei des Teufels, daß je heiliger das Volk ist, je böser und ärger es ist. Je reicher und heller das Licht des Worts scheint, je größer die Verachtung des Worts ist. Eben desselben besorgt sich Jakob auch von seinen Nachkommen und sonderlich von Simeon und Levi, daß nichts Gutes von ihnen kommen werde, sondern er verkündigt das Allerhöchste von ihnen zuvor.

\quellseite{499}

78. Er redet aber nicht von der Gnade, sondern von der Natur. Denn wir sind von Natur Kinder des Zorns, und sind derhalben auch geneigt und beflissen auf allerlei Sünden, auf greulichen Mord, Unzucht, Räuberei etc., wo uns die Gnade nicht hülfe und errettete. Derhalben verflucht nun Jacob solche verderbte Natur, die so gar schändlich und böse ist, dadurch auch diese Gesäße der Sünden blindlings getrieben worden, und bittet, daß GOtt die Nachkommen bei der Verheißung und in seiner Gnade erhalten wolle.

\begin{quote}
\textbf{Denn in ihrem Zorn haben sie den Mann erwürget, und in ihrem Muthwillen haben sie den Ochsen verderbet.}
\end{quote}

79. Ueber diesen Text haben die Ausleger mancherlei Meinungen. Hieronymus und viel Andere mehr legen den ersten Theil aus von den Männern, so zu Sichem sind erschlagen worden wegen der Schändung Dinas, und sagen, daß hier der Singular synedochisch gebraucht werde für den Plural, das ist, durch das Wort „Mann“ sollen verstanden werden „Männer“; den andern Theil aber ziehen sie auf Joseph: welches mir nicht gefällt, und ich halte dafür, man müsse beiderlei allein von Joseph verstehen. Denn daß diese Brüder Joseph verkauft, und so viel an ihnen gewesen, umgebracht haben, ist so eine große Sünde, daß wohl vonnöthen gewesen ist, daß man davon ein Gedächtniß behalten hat zum Exempel und zum Schrecken böser Menschen, die Simeon und Levi darin nachfolgen würden.

80. Und Jacob hat dazumal vor Augen gehabt und gesehen auf alle Nachkommen. Er hat wohl gewußt, daß die Priester grausamlich toben und wüthen würden wider die Propheten und andere gottselige Menschen, wie Simeon und Levi bei seinem Leben ihren Bruder Joseph mit Füßen getreten und ihren alten Vater betrübt haben bis in den Tod. Welches wahrlich eine große Verachtung gewesen ist väterlicher Gewalt und nimmermehr groß genug gemacht werden kann.

81. Nun ist keine andere Ursache solcher großen Feindschaft wider ihren Bruder gewesen, denn allein die Liebe und Gunst, so der Vater zu diesem Sohn getragen. Darum sind sie also auf ihn erbittert worden und haben ihn so heftig gehaßt, und bei sich beschlossen, daß sie weder des Alters ihres Vaters, noch Josephs an seinem Leben verschonen wollten; sondern wollten nur auf das greulichste wider Joseph wüthen, und den Vater mit Leid und Schmerz, so er zwei und zwanzig Jahre hat tragen müssen, entkräften und tödten, unangesehen aller Gebote und Verheißungen GOttes, welche Jacob, als ein sehr getreuer Bischof und Lehrer, seinem ganzen Haus vorgehalten hatte. Das sind nun die Früchte solches gottlosen und fleißigen Unterrichts gewesen. Es wäre nicht Wunder gewesen, daß Donner und Blitz darein geschlagen hätte.

82. Derhalben verstehe ich diese Worte von niemand anders denn nur allein von Joseph. Denn die Männer zu Sichem hatten die Strafe wohl verdient, und die andern Landsleute und Nachbarn hatten das zugegeben und eingewilliget, daß Jakob ihren Ort Landes inne behalten sollte. Joseph aber hat gegen seine Brüder nichts gesündigt, denn daß er einen Traum erzählt hatte, ohne alle Auslegung, oder Begierde über seine Brüder zu herrschen; welches damit angezeigt und bedeutet worden ist, daß sich die Sonne und elf Sterne vor ihm geneiget. Er ist ein sehr frommer unschuldiger Jüngling gewesen, welchen der Vater sehr lieb gehabt, und Simeon und Levi haben auch derhalben so viel schwerer gesündigt; denn sie haben keine Ursache gehabt zu solcher grausamen Tyrannei, so sie an ihm geübt.

83. Darum sind sie auch billig verworfen und mit diesem leiblichen Fluch gestraft worden. Und wo nicht der Heilige Geist durch sonderliche Gaben aus ihrem Nachkommen etliche Propheten erweckt hätte (wie denn viele und dazu treffliche Männer von ihnen hergekommen sind, als, Mose, Aaron, Samuel und Andere mehr), so wäre dieser Stamm gar unbekannt und verachtet geblieben.

84. Es wird aber hier gefragt: Warum diese Zwei, Simeon und Levi, Brüder genannt werden, so doch die zehn Andern auch Brüder gewesen sind? Darauf antwortet Lyra mit etlichen andern Hebräern sehr wohl, nämlich, daß Mose diese beiden in der greulichen Sünde habe wollen zu einander setzen, dieweil sie Gesellen gewesen, da sie ihren Bruder getödtet haben. Denn Ruben hatte widerrathen und gewehret, daß sie Joseph nicht tödten sollten, und hatte dazu einen Rath erfunden, wie sie thun möchten, daß sie ihn nicht erschlügen oder umbrächten, sondern daß sie ihn viel lieber in die Grube würfen, damit er ihn also aus seiner Brüder Händen und Tyrannei erretten möchte. Aber dieweil er bei ihnen dazumal gestunken und kein Ansehen gehabt wegen der Blutschande, die er begangen hatte, darum haben Simeon und Levi ihn darin auch verachtet.

85. Juda hat aber auch wider Joseph nichts vorgenommen, und er rieth auch dazu, daß sie ihn verkaufen sollten, auf daß er damit seinem Bruder das Leben erretten möchte. So sind die andern Brüder, als, Dan, Asser, Naphthali, auch nicht Stifter oder Anfänger dieser grausamen Tyrannei gewesen. Denn Mose hat gesagt, daß Joseph immer bei den Söhnen der Kebsweiber gewesen sei und sich zu ihnen gehalten. Derhalben Simeon und Levi den Ruben mit den andern Söhnen, so von den Mägden geboren waren, verachtet haben, und mußten also bei dem großen Stolz und Hoffart der zwei Brüder dieselben Söhne der Kebsweiber gar verachtet und verstoßen sein. Wo aber Sebulon und Isaschar etwas gethan haben, so haben sie doch dasselbe nicht aus eigenem Vermessen gethan, sondern auf Rath und Eingeben dieser zwei. Denn sie sind dazumal zwölf oder dreizehn Jahre gewesen, da diese Sünde begangen worden ist.

86. Derhalben sind Simeon und Levi die ersten Stifter und Rädelsführer gewesen, daß Joseph ist verkauft worden, welche nach Ruben an Alter und Herrlichkeit vortrefflicher gewesen sind denn die Andern. Darum zürnt auch nun Jakob mit ihnen, da er die Sache erfahren hat; wiewohl er sie nicht entsetzt, sondern schließt, daß die Kinder nach dem Fleisch und ihre Nachkommen gottlos sein werden, die nichts Gutes stiften werden. Wenige ausgenommen, als, Mose, Aaron, Samuel, die nicht der fleischlichen Geburt halben, sondern von GOttes Gnaden fromm und heilig gewesen sind. Denn GOtt macht bisweilen aus dem Besten das Aergste und aus dem Aergsten das Beste. Wie aus Petro, da er den Herrn Christum verleugnet, ein trefflicher unerschrockener Mann geworden ist, der darnach Christum frei bekannt hat. Also ist aus Paulo, der zuvor ein Lästerer und Verfolger Christi war, ein großer Apostel geworden. Sonst, wo es ohne GOtt ist, bleiben sie die Aergsten, wie das an den Leviten Korach und Dathan zu sehen ist, die sich wider Mose gesetzt, der doch von GOtt berufen war und Wunderzeichen that. Das war Simeons und Levis Geschlecht, denen sie gleich waren, da sie Mose und Aaron umbringen wollten. Ja, siehe Miriam und Aaron an, die selbst auch wider Mose murrten.

87. So hat nun Jacob haben wollen, daß dies Exempel da hat stehen sollen, daß es die Nachkommen immerdar vor ihren Augen hätten und im Herzen daran gedächten. Denn was kann unbilliger sein, sagt Jacob, denn daß Simeon und Levi meinem allerliebsten Sohn, von meinem vornehmsten Weibe geboren, dürfen solche Gewalt thun? dieweil ich selbst noch die Kirche regiere und das Evangelium von dem künftigen Samen lehre, dazu mit so großem Fleiß ihnen die Gebote GOttes einschärfe. Und Joseph hat sich auch gegen seine Brüder und sonderlich gegen Simeon und Levi so gar hart und unfreundlich erzeigt; denn er hatte Sorge, sie hätten den Benjamin auch erwürgt.

88. Derhalben so liegt nun die Schuld alle miteinander auf diesen zwei Brüdern, von denen Jacob sagt: „Denn in ihrem Zorn haben sie den Mann erwürget.“ Denn wiewohl sie ihn nicht wahrhaftig erwürget, so hat es ihnen doch am guten Willen nicht gemangelt, und sie haben sich unterstanden, daß sie ihn wollten umbringen, wo ihnen Ruben und Juda nicht widerstanden hätten. Welches denn aus ihren Worten klärlich zu vernehmen ist, da sie sagen, Cap. 37, 20.: „Kommt, und laßt uns ihn erwürgen, und sagen, ein böses Thier habe ihn gefressen.“ Da war nicht allein der Wille, sondern war vielmehr ein grimmiger Zorn und Unsinnigkeit, ja, es war auch die rechte That; und er hätte in derselben Stunde sterben müssen, wo sich nicht die andern Brüder wider ihr unsinniges Vornehmen gelegt und demselben widerstanden hätten. Und Joseph hat auch nichts Anderes sehen oder denken können, denn daß jetzt die Todesgefahr vorhanden wäre. Darum sagt auch Juda Cap. 42, 21.: „Das haben wir an unserm Bruder verschuldet, daß wir sahen die Angst seiner Seele, da er uns flehete, und wir wollten ihn nicht erhören.“ Item daselbst V. 22. sagt Ruben auch: „Nun wird sein Blut gefordert“ etc. Sie reden traun selbst vor Joseph, den sie für den Herrn des Landes hielten, nicht anders, denn als wäre sein Blut vergossen worden. Denn es war da ein vollkommen Vornehmen, welches für das Werk gerechnet und gehalten wird; denn der Wille, Joseph zu tödten, war bei ihnen vollkommen. Darum sagt Jakob: „Sie haben den Mann erwürget.“ Und nennt Joseph deshalb einen Mann, darum daß er so vortrefflich und groß gewesen ist, als wollte er sagen: Sie haben mir einen vortrefflichen Mann erwürget. Ich meine, er sei ein Mann geworden, den Mann, desgleichen die ganze Welt nicht hat.

\quellseite{500}

89. Und er kann wahrhaftig \unsicher{persönlich} reden mit den hohen großen Männern im Alten und Neuen Testament, mit Mose, mit David. Sollten sie denn dafür ungestraft bleiben, daß sie den heiligen Jüngling und so großen Mann erwürget haben? Es ist noch genug, daß sie nicht gar verderbt und in Grund vertilgt werden. Und dasselbe geschieht um meinetwillen, sagt Jakob, Gott verschonet ihrer, daß sie nicht gar ausgerottet werden, wie sie mit ihren Sünden verdient hätten. Aber ich will haben, daß bei den Nachkommen dieses Todtschlags, den sie an ihrem Vater und Bruder begangen, immer für und für soll gedacht werden, auf daß die Kindeskinder und die von ihnen noch sollen geboren werden, lernen mögen, wie sie ihre Eltern ehren und vor Augen haben sollen, und ihre Brüder lieben sollen. Derhalben, ihr meine lieben Söhne und Enkel, die ihr wollt rechte wahrhaftige Israeliten sein, hütet euch, daß ihr ja nicht Simeoniten und Leviten werdet: fliehet immer weg von diesen Bruder- und Vatermördern.

90. Da die andern Brüder diese traurige und scharfe Predigt gehört, werden ihnen ohne Zweifel die Thränen die Backen herab gelaufen sein, werden jämmerlich gesessen haben, und selbst Joseph; und der Vater wird diesen Fluch auch nicht haben aussprechen können ohne Seufzen und Weinen. Ich würde wahrlich dies weder ausreden noch hören können ohne große Bewegung des Herzens und aller Eingweide.

91. Aber die Welt und die Gottlosen lassen sich den Zorn und Fluch Gottes nicht bewegen; und hat doch dieser Fluch viel gegolten bei denen, die geistliche Israeliten gewesen sind. Alles, David, Salomo und Andere mehr sind durch dies Exempel bewogen worden, daß sie für und für bedacht, wie sie ihre Eltern ehren und vor Augen haben sollten. Absalom aber und Andere mehr haben es verachtet. Ich glaube zwar wohl, daß Simeon und Levi sündig geworden, so viel ihre Seele belanget: aber die Schmach und Makel ist gleichwohl auf ihrem ganzen Geschlecht liegen geblieben. Und es ist wahrlich auch ein solch Exempel gewesen, daran man billig für und für hat denken sollen, auf daß dadurch die Nachkommen erweckt und ermahnt würden, sich mit gebührlicher Ehrerbietung gegen ihre Eltern und Brüder zu erzeigen, daß sie nicht auch denselbigen Fluch und Strafe auf sich laden möchten, welche Ruben, Simeon und Levi getragen haben: Siehe, liebes Kind, wie es mit Simeon und Levi ergangen ist. Wie wir jetzt die Jugend zu ermahnen pflegen, daß sie ihren Eltern und Schulmeistern gehorsam seien, daß sie nicht den Diebesherren und Satan zum Meister kriegen, und denselben darnach folgen und gehorchen müssen.

92. Das hebräische Wort, razon, heißt Wille; als, im 51. Psalm V. 20.: „Thue wohl an Zion nach deiner Gnade“, oder gutem Willen, und bedeutet hier so viel, daß sie ihren Bruder haben erwürgen wollen, wiewohl sie dasselbe nicht vollbracht, und daß derselbe Wille für das Werk gerechnet werde. Wie die Juristen den auch verdammen, der gefunden wird, daß er gewißlich des Vorhabens ist, etwas Böses zu vollbringen, nämlich, wenn es bei dem nicht steht, der etwas vorgenommen hat zu thun, daß er es nicht vollbringen kann.

93. Das aber Jakob hier sagt: „Und in ihrem Muthwillen haben sie den Ochsen verderbet“, hat eben denselben Verstand, den auch die vorigen Worte haben, allein, daß mit andern Worten wiederholt wird, was er zuvor auch gesagt hat. Denn es wird oftmals in der heiligen Schrift einerlei Ding auf verschiedene Weise ausgedrückt, damit anzuzeigen, daß es gar gewiß sei. Wie Joseph zweierlei Träume sind gegeben worden, die doch beide einerlei Bedeutung gehabt, als, durch die Garben und durch die Sonne, Mond und Sterne. Und wir haben gehört, daß in den Träumen Pharaos durch die Aehren und Kühe auch einerlei ist bedeutet worden. Darnach brauchen die Propheten auch sehr oft diese Figur, die da heißt tautologia, wenn man einerlei Ding mit andern Worten wiederholt, damit auszudrücken, daß gewiß sei, was sie verkündigt haben, und daß es also kommen werde. Also spricht Jesaia sehr oft einerlei Ding aus mit zweierlei Sprüchen. Also wird an diesem Ort auch gesagt von dem, daß sie Joseph todt geschlagen hätten: „Sie haben den Mann erwürget, und haben den Ochsen verderbet.“ Dies ist eine sonderliche Art und Weise, welche die heilige Schrift für und für braucht, und gehört dazu, daß damit die Zuhörer gestärkt werden sollen, daß sie nicht zweifeln an dem, was gesagt ist.

94. Das Wort schor schor bedeutet mehr denn einerlei. Die diesen Text von den Männern zu Sichem verstehen, deuten es, daß es Mann heiße; in welcher Bedeutung es hernach wird gebraucht werden in dem Segen Josephs. Schir heißt, Gesang; scharah heißt, Herrschaft, Mauer, oder auch wohl Königreich; im dritten Buch Mose wird es oftmals gebraucht; und das Wort schor wird allenthalben verstanden, daß es Ochse heiße.

95. Die neueren oder jetzigen Juden gebrauchen gemeiniglich das Wort schorbor, dessen Lyra auch gedenkt, und sagen, daß es ein Ochse sei, von dem sie fabeln, daß er gemästet werde mit dem Gras auf den Bergen und auch in Thälern, daß er am jüngsten Tage geschlachtet, und die so davon selig werden, von seinem Fleisch essen werden. Wie Mahomet in seinem Alkoran auch dichtet von dem Fisch Leviathan, welcher ein großer Wallfisch ist im Meer, davon im 104. Psalm steht V. 26.: „Da sind Wallfische, die du gemacht hast, daß sie im Meer scherzen“ etc. Von demselben Fisch läßt sich Mahomet auch träumen, daß er in diesem Leben genähret werde und daß er darnach einer schönen Jungfrau dienen werde etc. Aber die Lügentheidung will ich fahren lassen.

96. Das Wort schor heißt in der hebräischen Sprache nicht allein, Mauer, sondern auch, Ochse. Darum soll man an diesem Ort lesen: Sie haben mir meinen Ochsen schier geplaget. Und Mose 5 Mos. 33, 17. sagt von Joseph: „Seine Herrlichkeit ist wie ein erstgeborner Ochse.“ Es wird in der heiligen Schrift ein erstgeborner Ochse sehr gelobt, darum daß er das Haus ernährt und versorgt, und man durfte keinen Ochsen opfern, den man vor den Pflug spannte; denn der war wie ein Hausvater und eine Säule des Hauses. Und die Heiden haben sich derselben Ochsen, damit man gepflügt, auch enthalten; denn sie pflügten nicht allein, sondern droschen auch. Derhalben werden im Gesetz die Werke oder Arbeit des Ochsen hoch gerühmt. Und Salomo sagt auch Spr. 14, 4.: „Wo nicht Ochsen sind, da ist die Krippe rein; aber wo der Ochse geschäftig ist, da ist viel Einkommens“; item, das Gesetz verbietet auch, 5 Mos. 25, 4., daß man dem Ochsen, der da drischet, das Maul verbinde.

97. Darum will nun Jakob so viel sagen: Den Mann, der in der so langen Zeit der Theurung nicht allein ganz Egypten, sondern auch uns alle ernährt hat, die wir sonst, wo es ohne ihn gewesen wäre, hätten müssen Hungers sterben, und den ich sonderlich lieb gehabt, darum daß er in meinem Alter geboren war, und den jedermann um seiner guten Art und großen Gottseligkeit willen lieb und werth gehalten: diesem Ochsen haben sie seine Kraft genommen, sagt Jakob, haben ihn aber nicht gar getödtet, und hätten gern gehindert, daß er der Erstgeborne, oder Vater und König sein sollte und mir in meinem Geschlechte am nächsten nachfolgte. Nun ist er aber jetzt gleichwohl zu solcher großen Herrlichkeit erhoben worden, daß er noch hat ein Vater in Egypten und in meinem ganzen Hause werden müssen. Etc, wie ist ihr Rath und ihr Wille, ihr secretum, so fein fortgegangen!

98. Solche Exempel und Vornehmen der Gottlosen sollen immer gelehrt und der Kirche oder Gemeinde Gottes vorgehalten werden, daß wir wissen, wie ihre Anschläge von Gott gehindert und zunichte gemacht werden, daß sie keinen Fortgang gewinnen, wiewohl es ihnen weder an großer Gewalt noch an Weisheit gemangelt. Wie dasselbe Jakob seinen Söhnen hier sein vor Augen stellt, nämlich, da er sagt: „Die sein habt ihr, Simeon und Levi, gehindert, daß Joseph nicht der Erstgeborne und der liebe Sohn des Vaters wäre! Sehet, er ist jetzt ein Fürst und ernährt uns alle.“

99. Also richten die Tyrannen mit ihren greulichen und blutdürstigen Edicten und Rathschlägen nichts Sonderes aus, denn daß sie nur die Gottseligen, welche sie verfolgen, immer zu größerer Herrlichkeit und endlich zur ewigen Seligkeit erheben; denn sie wollen auch Eimenge
\quellseite{501}
und Leviten sein, aber nicht ohne ihren eigenen großen Schaden. Also hat der Papst den heiligen Märtyrer Johannes Huß auch erwürgt, hat denselbigen frommen und frommen Mann nicht dulden können. Aber, was hat er für Nutzen von dem Mord? Er ist nun todt. Was hat, sage ich, der Papst davon? Es hat, Gott Lob! von der Zeit an das Licht des Evangeliums in den Herzen der Gottseligen viel klarer angefangen zu scheinen denn zuvor, und Christus hat angefangen, ihn mit dem Geist seines Mundes umzubringen, 2 Thess. 2, 8.

100. Auf diese Weise hat Gott haben wollen, daß in der ganzen Welt ein Exempel sein sollte, wie Simeon und Levi gestraft worden sind dafür, daß sie Vater- und Brudermörder waren, auf daß wir ihnen ja nicht nachfolgen und gleich werden, dieweil mir, leider! sehen, daß solche Sünden oftmals in der Welt mit aller Gewalt getrieben werden, wenn die Leute sicher sind, und vom Satan dazu gereizt werden, daß sie sich weder mit Gesetzen, noch mit Exempeln der Strafe wollen zwingen und im Zaum halten lassen.

\begin{quote}
\textbf{V. 7. Verflucht sei ihr Zorn, daß er so heftig ist, und ihr Grimm, daß er so störrig ist. Ich will sie vertheilen in Jakob, und zerstreuen in Israel.}
\end{quote}

101. Es ist unter dem Volk, und in der ganzen Regierung nichts schädlicher, denn grimmiger Zorn, Grausamkeit und Tyrannei, vornehmlich wider das Blut der Brüder und wider unser eigen Fleisch. Darum nennt es Jakob einen verfluchten Zorn, daß er so heftig und stark, das ist, so tyrannisch und grimmig ist.

102. Wo aber nun der Fluch dieser zwei Brüder und die Zerstreuung sie erfüllt worden, kann man nicht leichtlich aus den Historien sehen. Daran ist kein Zweifel, daß der Stamm Levi mit den übrigen Söhnen Israel kein Erbe und keinen gewissen Ort im Lande zu eigen gehabt hat, und dasselbe weiß man auch von Simeons Nachkommen wohl. Wiewohl nun dieser Fluch sehr gering und leicht scheint, so ist es doch wahr: wenn er recht betrachtet und erwogen wird, so ist darin ein großer heftiger Zorn begriffen.

103. Aber der Gnade ist Alles vorbehalten, die aus den ärgsten Sündern die allerbesten Pfarrherren oder Seelenhirten gemacht hat.

Wie das Exempel Pauli bezeuget, der ein großer Verfolger und Lästerer gewesen ist, und doch durch Gottes Gnade der höchste Apostel geworden ist, und mehr gearbeitet hat, als irgendeiner der andern Apostel gethan hat. Also sind aus dem Stamme Levi große Propheten und Regenten hergekommen, Mose, Aaron, Mirjam, Samuel, Korah, sehr treffliche Männer, daß es auch schien, als hätte Gott dieser Drohung und Fluches vergessen. Und du wirst nicht leichtlich in irgend einem andern Stamm solche treffliche Männer finden, den Stamm Juda ausgenommen. Denn wem wollen wir mit Mose, Aaron, Samuel, item, mit Zacharias, Johannes dem Täufer und vielen andern großen Männern vergleichen? welche alle aus diesem Geschlechte gewesen sind. Und 5 Mos. 33, 8. ff. hat Mose in der Verheißung vom Stamme Levi große Dinge geweissagt, wiewohl Jakob hier nicht ein gutes Wort davon redet, unangesehen solcher trefflichen großen Männer, die aus diesem Stamme hergekommen sind.

104. Darum habe ich droben gesagt, daß man diese Regel fleißig merken soll, nämlich, daß Gott wohl zürnt und straft, aber doch also, daß er dabei Gnade und Barmherzigkeit sein läßt um der bußfertigen Sünder willen. Als, da er verkündigen läßt, daß Ninive soll umgelegt und in Grund verderbt werden, und die Niniviten nach gehörter Ermahnung sich zu Gott bekehren, da wird die Drohung geändert und Gott verschont ihrer. Wo sie aber fortgefahren wären und in ihren Sünden verharrt hätten, so wäre ohne Zweifel das Verderben erfolgt, daß sie hätten untergehen müssen.

105. Dieweil aber Jakob nun verstanden und erkannt hat die Art und Natur Simeons und Levis aus dem, daß sie die Männer zu Sichem erwürgt und Joseph verkauft hatten, hat er es dafür gehalten, daß ihre Nachkommen nicht weniger tyrannisch sein würden, wo sie zum Regiment kommen sollten, wenn nicht die Gnade dazu kommen und solches brechen würde. So hat er, sage ich, gedacht: Kommen sie in das Regiment, so werden sie es nicht gut machen. Sie sind von Natur so gar boshaftig, grausam und blutgierig, daß wenn zu hoffen wäre, daß von ihnen etliche Nachkommen sollten geboren werden, dürfte man daran nicht zweifeln, dieselben Nachkommen würden auch blutgierig, grimmig und tyrannisch werden. Darum behüte Gott vor dem Exempel.

106. Aber Gottes Gnade ist Alles möglich. Also erregt Korah einen großen Aufruhr und er selbst kommt darüber um. Aber es folgt das große Wunder hernach, daß da er mit all seinem Hab und Gute verschlungen worden, seine Kinder gleichwohl noch erhalten werden, die darnach ihrer großen Tugenden halben sind gerühmt worden, da sie viel schöne Psalmen gemacht haben, die man leichtlich mit den Psalmen Davids vergleichen mag; und waren doch vom Geschlechte Levi und Korah. Die aber der grausamen Tyrannei ihres Vaters nach dem Blute und verfluchten Fleische, und der Erbsünde gefolgt, die sind umgekommen. Die Andern aber, denen die Gnade und Barmherzigkeit Gottes zuvorgekommen ist, die sind heilig und selig gewesen.

107. Derhalben hält Gott diese Regel für und für. Er hält den Leuten seine Drohung vor und erwähnt doch aus bösen Menschen und Sündern etwas Gutes. Wie er aus dem ganzen menschlichen Geschlechte, so der Sünden halben verloren gewesen, noch Etliche selig macht. Also ist Kain mit seinen Nachkommen verloren worden und sind doch ohne Zweifel noch Viele davon bekehrt worden. Also sind Ismael, die Edomiter und Moabiter Feinde des Volks Gottes gewesen, und sind doch unter ihnen auch Viele zur wahren Erkenntniß Gottes gekommen und selig geworden. Denn Christus ist von einer Cananitin und Moabitin geboren worden. Dies sind Ausnahmen der Gnade und sonderliche Gaben; die Natur aber geht immer vor sich in dem Fluche, wie Reuben und Simeons Fluch gewesen: die Art ist nicht gut. Ich verdamme, sagt Jakob, die böse Art und Natur, um deren willen auch etliche Fromme fallen und leiden werden, wiewohl sie durch Gottes Gnade von den Andern abgesondert werden sollen.

108. Was aber diese Stämme für Schaden und Unglück haben leiden müssen, weiß man nicht, und ist auch das in der heiligen Schrift nicht alles begriffen. Daran ist genug, daß angezeigt wird, daß sie im Lande Canaan kein eigen Erbe oder Gut gehabt haben. Simeon hat ein klein Stück Landes und gleichsam einen Winkel bekommen im Stamme Juda, welcher Stamm einen sehr großen Theil gehabt, und zugelassen hat, daß Simeon an einem Orte des Landes Juda sich aufhalten möchte; und ist zwar dasselbe nicht ein geringer Schade gewesen, auch zeitlich.

109. Und ist auch dieses nicht eine kleine Anfechtung gewesen, daß die Leviten gar keinen Theil oder Erbe mit ihren Brüdern gehabt, sondern wie 5 Mos. 10, 9. gesagt wird: „Der Herr ist ihr Erbe.“ Denn sie haben gegessen von dem Opfer des Volks, und haben in etlichen Vorstädten gewohnt, so hin und wieder zerstreut gewesen sind durch alle Stämme, ohne Reich und Regiment. Sie sind Hirtenknaben und Kuhmelker gewesen, die sich mußten hin und wieder zerstreut in den andern Stämmen nähren. Wiewohl es ehrlich und rühmlich ist, vom Erbe des Herrn und vom Opfer gespeiset werden. Es ist aber zumal ein erbärmlich Leben wegen der Bosheit und des großen Geizes der Menschen. Gnade Gott denen, die sich davon nähren sollen! Im Nehemia und an andern Orten mehr kann man wohl sehen, wie sie die Leviten genährt haben; wie sie uns jetzt auch speisen.

110. Denn die Christen haben auch ein Gebot und Befehl, daß sie die Priester und Diener des Worts erhalten sollen, wie Christus sagt Luc. 10, 7.: „Ein Arbeiter ist seines Lohnes werth“; und 1 Cor. 9, 14.: „Der Herr hat befohlen, daß, die das Evangelium verkündigen, sollen sich vom Evangelio nähren.“ Aber es werden Etliche besser gespeiset denn Andere und man thut ihnen etwas gütlicher. Etliche werden auch zwar schmal genug gehalten: aber sie werden fast alle miteinander also unterhalten, daß sie wohl mehr bedürften. Ueber das, so sind die Pfarren und die Güter, so dazu gehören, nicht erblich. Eben also müssen wir denken, daß es auch bei den Juden zugegangen sei. Es haben beide Bürger und Bauern in den andern Stämmen die Leviten ärmlich genug gehalten, also daß sie diese Weissagung und Fluch Jakobs hart genug betroffen hat, und daß auch denselben Fluch fromme heilige Männer haben erfahren müssen. Mose, Aaron, Samuel, die nicht viel übriges Gutes gehabt, wiewohl sie viel schwerer Arbeit gethan und um Gottes willen viel große Dinge ausgerichtet haben. Und sie selbst rühmen sich dessen vor Gott; wie 1 Sam. 12, 3. steht, da Samuel sagt: „Du weißt, daß ich nicht jemands Ochsen oder Esel genommen habe, auch niemand habe Gewalt und Unrecht gethan“ ac.

\quellseite{502}

111. Also haben auch die Frommen den Fluch tragen müssen: aber dieselben haben solches mit Geduld gelitten und solche Armuth mit Geduld in Glauben überwunden. Die andern Stämme haben zu ihrer Unterhaltung, wie der Prophet Maleachi klagt Cap. 3, V. 8. 9., nur die Hefen von ihren Gütern gegeben; wie es nach dem Hebräischen lautet. Was sie selbst nicht mochten, war den Pfaffen gut.

112. Daher ist auch das greuliche Unglück der Abgötterei im Stamm Levi gekommen, auch bei dem Theil, der heilig war (denn vom Priestertume Jerobeams rede ich nicht), daß sie solche Hantierung erdacht haben, damit sie viel haben erwerben können. Denn dieweil sie um des rechten wahren Gottesdienstes und der reinen gesunden Lehre willen verachtet und verfämt worden, haben sie andere Gottesdienste erdichtet, die einen großen Schein gehabt, und dem Pöbel äußerlich vor Augen gestanden sind, daß sie daran Gefallen getragen, auf daß dadurch die Priester mehr geehrt werden und Geld und Gut damit erlangen möchten. Dergleichen Gottesdienste sind gewesen, so sie dem Baal Peor, Astaroth und Camos angerichtet haben; dazu hat alles Volk mit höchster Bereitwilligkeit und großer Andacht gegeben, und sind also für solche ihre Thorheit recht gestraft worden.

113. Eben also widerfährt es auch zu unsern Zeiten rechten getreuen Pfarrherren, welche den Bauern und Bürgern nicht viel mehr geben denn nur schwarzes Brod, und wie der Prophet Jesaja sagt im 30. Cap. V. 20.: Sie werden gespeiset in Trübsal mit Brod, und in Aengsten mit Wasser etc. Darum wird Gott auch dadurch zum Zorne gereizt, und verhängt, daß sie falsche Gottesdienste, Opfermesse, Fegefeuer, Ablaß und Anrufung der verstorbenen Heiligen erdichten; wie Paulus sagt 2 Thess. 2, 10—12.: „Dafür, daß sie die Liebe zur Wahrheit nicht haben angenommen, daß sie selig würden, darum wird ihnen Gott kräftige Irrthümer senden, daß sie glauben der Lüge. Auf daß gerichtet werden alle, die der Wahrheit nicht glauben, sondern haben Lust an der Ungerechtigkeit.“ Daher ist des Papstes und seiner Bischöfe Herrlichkeit und Gewalt so groß geworden; denn sie sind jetziger Zeit Fürsten und Herren über die ganze Welt. Warum aber? daß sie das Evangelium predigen und der Seelen so fleißig warten? Mit nichten. Ei warum denn? Darum, daß sie Abgötterei mit Ablaß und Opfermessen anrichten. Das gibt viel Geld und Gut, damit haben sie die großen Stifte und Klöster, und herrlichsten Paläste aufgerichtet, das wohl der Türke mit alle seinem Gelde und Gute nicht thun könnte. Sie wollten die rechten Pfarrherren und Seelsorger nicht hören, darum hat ihnen Gott kräftige Irrthümer gesandt, daß sie den Gott maussim ehreten, wie Daniel im 11. Cap. V. 38. das zuvor verkündigt; und den Gott, davon ihre Väter nichts gewußt haben, ehreten mit Gold, Silber, Edelsteinen und Kleinoden. So wollen wir es haben.

114. Und wenn das Licht des Evangeliums verlöschen würde, und die Finsterniß, davor uns Gott gnädig behüten wolle, wiederum überhand nehmen sollte, so würde doch jetzt das gemeine Volk nicht weniger unsinnig sein, und würden die Leute noch mit größerer Begierde und Tollheit viel Geld und Gut geben, die Irrthümer zu bestätigen, denn zuvor geschehen ist. Denn die Natur der Menschen ist abgöttisch und abergläubisch, flieht vor dem rechten wahren Gott, vor rechtem Gottesdienste und Gottesfurcht: aber das Vertrauen auf ihre selbsterdichteten Gottesdienste und Werke heben sie immer sehr hoch. Davon werden die Pfaffen reich und große Herren. Wie im Papstthum des unnützen Gebens und der Beschenkungen keine Zahl oder Maß gewesen ist. Die Barfüßermönche haben mit ihren Eigenleibdingen das Volk ganz greulicher Weise betrogen und ausgesogen mit den Rosenkränzen und zehen Vater Unser; die funden an allen Thüren. Dazu hatte das Volk Lust und nahm solch Ding mit großer Begierde an, und gab gar freigebig und, wie man sagt, mit beiden Händen dazu. Was aber jetzt auf Schulen und Diener des Worts gewendet wird, das muß alles auf das übelste angelegt sein, wie das gemeine Volk davon urtheilet; ja, sie wollten die Kirchenzier auch noch gern behalten des Raubes Egyptens, den sie noch vom Papstthum übrig behalten haben.

115. Also haben etliche heilige Männer, Priester und Leviten im Volk Israel, diese Strafe getragen um ihres Vaters Sünde willen; wiewohl sie dem Vater nicht gleich waren, haben sie doch dessen Fluch tragen müssen. Die Andern aber, die dem Exempel Mosis und Samuels und auch deren Gottseligkeit nicht nachgefolgt sind, die sind viel ärger gewesen, denn Simeon und die andern Nachkommen. Die haben sich gänzlich darauf begeben, daß sie nur des abgöttischen Gottesdienstes gewartet und die rechten Propheten greulich verfolgt, und haben sich der Sünde gar nicht enthalten, welche Jakob hier so heftig verflucht hat.

116. Wie heutiges Tages der Papst zu Rom nur damit nach seinem höchsten Fleisch umgeht, daß er fromme gottselige Leute erwürge, und ist so ein greulicher Brudermörder, daß niemand seine Grausamkeit, weder mit Worten noch mit Gedanken, begreifen kann. Ueber das beraubt er auch die ganze Welt und reißt mit Gewalt alle irdischen Güter zu sich, so groß ist nun der Zorn Gottes über die Undankbarkeit der Menschen. Denn dieweil sie die guten frommen Pfarrherren nicht hören wollen, so mögen sie dagegen die bösen und Mörder hören, die Simeone und Leviten, die von Natur verflucht sind, und derhalben zur Abgötterei und allerlei Sünden dahin gerissen werden.

117. Aber etliche wenige Gottselige und Fromme in demselben Geschlecht, denen Gottes Gnade zuvorgekommen und die sie erhalten hat, die haben gedacht, es wollte ihnen gebühren, daß sie mit demjenigen zufrieden wären, das Gott ihnen verliehen und zugelassen, und haben also gesagt: Wir wollen uns daran genügen lassen, daß der HErr unser Erbtheil ist, der wird es reichlich belohnen, erstatten und vergelten, was uns in diesem Leben mangelt.

118. Wie jeglicher sich fromme gottselige Kirchendiener in so großen Trübsalen und Beschwerungen sich mit diesem Troste auch erhalten, da sie bei sich selbst also schließen: Der HErr der Ehren ist selbst für uns gekreuzigt und gestorben, warum sollten wir denn nicht auch mit Geduld leiden, wenn uns etwas Widerwärtiges begegnet? die wir zwar unsrer Natur nach allerlei Unglücks und dazu auch der ewigen Verdammniß werth sind.

119. Und wiewohl, wo uns nicht Gottes Gnade und Barmherzigkeit hülfe, würden wir auch in dieselben Sünden fallen, welche Jakob, der Vater, seinen Söhnen, Simeon und Levi, aufrückt und um deren willen er sie so hart schilt. Und wo ich ein Mönch geblieben wäre, so würde ich mit andern Feinden des Evangeliums jetziger Zeit Christum, und fromme gottselige Lehrer, so sich wider das Papstthum legen, auch verfolgen helfen. Aber die Gnade hat meine Natur geändert, daß ich mich nicht zu den Feinden gesellete und auch unschuldig Blut vergießen hülfe. Derselben soll man dies fleißig merken von der Strafe dieser zwei Brüder. Denn das ist die Ursache gewesen vieler großen Sünden bei den Bösen, die darnach falsche Gottesdienste und Abgötterei angerichtet und die Verfolgungen heiliger frommer Leute dazu auch verursacht haben.

120. Es hat aber Mose diesem Fluche entgegengesetzt einen schönen herrlichen Segen, 5 Mos. 33, 9., daß du daraus verstehen mögest, daß dieser Stamm nicht gar alles Segens beraubt ist, wiewohl ihm an diesem Orte Jakob Alles genommen hat; aber er will damit die böse Natur und das alte Exempel verdammt haben. Mose aber hat denen gepredigt, die unter den geistlichen Nachkommen sein würden, da er sagt: „Wer zu seinem Vater und zu seiner Mutter spricht: Ich sehe ihn nicht; und zu seinem Bruder: Ich kenne ihn nicht; die halten deine Rede“ etc. Da legt Mose den Fluch aus und mäßigt oder lindert denselben, als wollte er sagen: Levi, du bist von Jakob verflucht worden, wenn du der Natur folgen willst, und dich dessen rühmst, daß du nach dem Fleisch unter die Nachkommen Abrahams, Isaaks, Jakobs und Levis gehörst etc. Derselben will niemanden sein, daß du eine andere Natur anziehest, und sagest: Ich kenne euch nicht etc., und daß du dich auf das väterliche Blut nicht verlaßest und darauf tröstest, so du die Gebote des HErrn recht lehren und halten willst. Wie im 45. Psalm V. 11. gesagt wird, der hieraus genommen zu sein scheint: „Höre, Tochter, schaue drauf, und neige deine Ohren, vergiss deines Volks und deines Vaters Haus“ etc. Darum wirst du also müssen gar ein anderer Mann werden, unangesehen, wie oder wer deine Großväter, Väter, und dein ganzes Fleisch und Natur gewesen sei. Ja, du mußt also gesinnet sein, daß du sagest: „Ich weiß von meinen Vätern und Brüdern nichts etc., ich kenne allein den Herrn selbst und will mich auch keiner allein rühmen; wie Joh. 1, 13. steht: „Welche nicht von dem Geblüt, noch von dem Willen des Fleisches, noch von dem Willen eines Mannes, sondern von Gott geboren sind.“

\quellseite{503}

121. Aber diese Vermahnung betrachten die Leviten und Simeoniten, darum wissen sie von nichts mehr zu rühmen, denn daß sie von dem Geschlecht der Väter geboren sind und deshalb von Gott nicht können verworfen werden. Mose lehrt aber das Gegenteil und spricht: Du sollst also sagen: „Ich kenne das Blut, dazu meinen Vater und Bruder nicht.“ Und Paulus sagt Röm. 9, 8.: „Nicht sind das Gottes Kinder, die nach dem Fleisch Kinder sind, sondern die Kinder der Verheißung werden für Samen gerechnet.“ Und dasselbe wird auch setzt an diesem Orte also geschlossen. Denn Simeon und Levi sind ja Kinder und stammen von dem gebenedeyten Blute her, und sind doch verflucht und zerstreut.

122. Sie rühmten sich aber nichts desto weniger wegen des Geblüts der Väter, und hatten einen Titel an der Verheißung und Gnade. Wie der Pabst auch das einige Stück immer treibt und widerhohlt: Wir sind die Kirche, wir sind die Kirche; darum können wir nicht irren und die Andern sollen uns gehorchen. Wider diesen Ruhm wird hier gesagt: „Verflucht sei ihr Zorn.“ Deshalb, auf daß du von diesem Fluch errettet werden mögest, so höre auf diese Vermahnung Moses, und sage: „Ich kenne euch nicht“ etc. Wir sind nicht die Kirche wegen der ordentlichen Succession, wie sie von den Vorfahren auf uns gekommen ist, sondern wir glauben an den Sohn Gottes.

123. So haben wir nun bis daher zwei Flüche gehört dreier Söhne Jakobs, die sehr hart und traurig sind, welche Mose hernach etwas lindert, da er 5 Mos. 33, 6. sagt: „Nuben lebe, und sterbe nicht“ etc., welches Jakob auch also will verstanden haben: „aber sein Pöbel sei gering“, sagt Mose. Item, Simeon soll zerstreut werden; wird er aber zu seinem Vater und Mutter sagen: „Ich kenne euch nicht“, so wird er in dem Fluch der Natur und des Geblüts nicht bleiben; wo aber nicht, so wird er kein Glück haben, wird unseelig und gottlos sein; wie hernach Saiphas und Hannas gewesen sind. Nun folgt weiter der vierte Segen, welcher sehr herrlich ist.

\section[Dritter Theil]{Dritter Theil. Von dem Segen und Verheißungen, so Juda empfangen}

\begin{quote}
\textbf{V. 8. 9. Juda, du bist's, dich werden deine Brüder loben. Deine Hand wird deinen Feinden auf dem Halse sein; vor dir werden deines Vaters Kinder sich neigen. Juda ist ein junger Löwe. Du bist hoch kommen, mein Sohn, durch große Siege. Er hat niedergekneiet, und sich gelagert wie ein Löwe, und wie eine Löwin; wer will sich wider ihn auflehnen?}
\end{quote}

124. Das ist nun eine schöne herrliche Verheißung; denn er sagt von Juda nichts Böses, sondern nur das Allerbeste und Herrlichste, so doch dieser Stamm auch sehr böse Könige gehabt. Derselben wie er hier von allem Fluche stille schweigt, also hat er droben des Segens gar nicht gedacht, und hat nur die Natur zum Exempel vorgestellet, wie er nun an diesem Ort die Gnade hoch rühmt. Denn es sind in diesem Stamme gottlose Könige gewesen, als Jerobeam, Ahas, Sallum, Manasse und andere mehr; und wird doch hier nichts denn eitel Gutes von Juda gesagt.

125. Die Ursache aber und Ursprung dieses großen reichen Segens ist die einige Blume, die aus diesem Volk Gottes in die Welt kommen sollte, nämlich Christus Jesus, um deß willen alles Gute nicht der Natur, sondern der Gnade hat müssen gesagt werden. Und dies ist eine sonderliche Weisheit oder Offenbarung, daß er es so kärlich in diesem Stamme von dem Samen unterscheidet; zu welchem nicht allein Klugheit der Vernunft und weltlicher Verstand gehört, sondern ein sonderlicher Geist, so in diesem Hirten war; denn er ist nicht ein politischer oder weltlicher Hirte gewesen, sondern seiner Heerde und seiner Kirche, welcher Kirche er solche große und so gewisse Dinge zuvor verkündigt.

126. Es werden aber ohne Zweifel die Söhne Jakobs mit großem Verlangen und mit Weinen und Seufzen auf das Andere, so ihr Vater noch hat sagen wollen, gewartet haben; denn sie waren sehr betrübt und erschrocken, da sie so gar harte und greuliche Flüche gehört. Und es hat auch wahrlich sehr hart gelautet, daß sich der Vater, da er jetzt sterben sollte, so hart und halsstarrig, und wie es scheint, so greulich erzeigt hat, da er die drei Söhne verworfen. Da werden sie geheult und gebeten haben: Ach nicht, lieber Vater etc. Dennoch fährt er fort, verändert aber daran gar nichts, was er einmal gesagt oder geordnet hat um der Nachkommen willen; denn er hat wohl gesehen, daß deren Laster und Sünden würden müssen mit Gefängniß und sonst anderen Strafen gehindert und verhütet werden.

127. Und Juda hat sich zwar auch des Weines nicht enthalten können, und ist ihm einig gefallen, daß er sich wird gefürchtet haben, was nun mit ihm werden würde, oder auch mit den Andern, dieweil die Drei also verworfen werden? Er wird gedacht haben: Hilf Gott! was wird es mit mir werden! jetzt kommt es an dich. Es werden in seinem Herzen sehr traurige Gedanken aufgestiegen sein von seinen Sünden und Lastern, von der Hurerei, so er mit Thamar getrieben, und von dem Rath, so er gegeben hat, daß sie Joseph verkaufen sollten. So gar ein erschrecklich Ding ist es, wenn die Gewissen erschrocken sind und sich entsetzen, vornehmlich, wenn so ein großer Patriarch predigt; und war zwar gar kein anderer Trost vorhanden, denn daß sie noch am Leben bleiben und Nachkommen haben sollten. Ich wäre wahrlich vor Leid und Thränen gestorben, und hätte gesagt: Höre auf, lieber Vater, ich sterbe vor Leid.

128. Es folgt aber bald ein schöner Thau und lieblicher Balsam und erquickt ihre Herzen wieder. Denn Juda wird zeitlich und geistlich gesegnet, und der Vater sagt zu ihm also: „Juda, du bist's“, daß ist, du bist der erlesene Fürst, Vater und Haupt in meinem Volk. Joseph hat auch viel Segen, und Juda hätte wohl verdient gehabt, daß ihn der Vater sollte verflucht haben; aber auf daß GOtt anzeigen und beweisen möchte, daß er seine Segen aus Gnaden und nicht aus Verdienst gebe, darum wendet er diesen Segen auf Juda; auf daß wir wissen sollen: was wir sind, das sind wir allein aus GOttes Güte und Gnade. Wie Augustinus sagt: Wir wissen, daß wir aus GOttes Gnaden sind, was wir sind. Wir wissen, daß wir Auserkornen sind aus GOttes Güte, nicht aus Verdienst. Das hat GOtt also wohl gefallen, daß sich niemand rühmen soll denn allein im HErrn.

129. Es zieht aber Jakob den Juda vor und setzt ihn über seine andern Brüder, und sagt: „Dich werden deine Brüder loben“, das ist, du wirst berühmt, groß und herrlich sein vor den Andern allen, also daß dich auch deine Brüder rühmen werden etc. Und es ist endlich auch also zugegangen, wie Jakob hier zuvor verkündigt hat.

130. Du sollst aber, sagt Jakob weiter, nicht allein bei deinen Brüdern herrlich gehalten werden, sondern du wirst auch sieghaft sein wider die Feinde. Du wirst viel Feinde haben; aber sei getrost und unverzagt, es hat keine Noth.

131. Aber sehr Wenige haben diesen Verheißungen geglaubt, als Obed, David und die diesen gleich gewesen sind; die Andern haben nicht geglaubt. Und es ist doch ein gewisses und festes Wort, darin er ihm die Ehre und Sieg wider die Feinde verheißt. „Deine Hand“, sagt er, „wird deinen Feinden auf dem Halse sein“, das ist, deine Feinde sollen sich nicht überwältigen, sondern werden vor dir fliehen und du wirst ihnen nacheilen. Und „vor dir werden deines Vaters Kinder sich neigen“, als wollte er sagen: Die, so gleichen Ursprung und Gewalt von ihrem Vater haben, werden alle der Same Jakobs genannt werden, aber du sollst der Hauptmann und der Fürst des Volks sein, sie werden dich in Ehren halten, und sich dein verwundern. Dasselbe ist hernach erfüllt worden beides an den Andern, vornehmlich aber an David. Und hat Jakob sehr herrlich abgemalt die große Kraft und Stärke Judas, da er ihn einen jungen Löwen nennt. Denn es werden ohne Zweifel mehr und größere Dinge ausgerichtet worden sein, als in der heiligen Schrift stehen. Wie auch die Flüche wider Kain und Ismael nicht alle beschrieben sind.

\quellseite{504}

132. Diese Verheißungen aber soll man im Geist und Glauben verstehen, auch mag das Zeitliche belanget. Wie das in der Historie Davids zu sehen ist, da es scheint, als habe GOtt sein, und was er ihm verheißen hatte, gar vergessen, dieweil es so schwerlich zugegangen ist, da er zum König erwählt und bestätigt worden ist. Und da er schon erwählt gewesen, hat er im ganzen Königreiche zehn Jahre lang keinen gewissen Ort oder Sitz haben können. Daher sind eigentlich so viel Psalmen gekommen, die alle gemacht sind um der vielen unzähligen Gefahr und Verfolgung willen, so er vom Saul gelitten. Und endlich, da er nun zum Reich kommen ist, fällt er schändlich in große greuliche Sünde, und wird sehr hart geplagt, da er die Strafe für solche Sünde getragen und in der Noth gesteckt hat.

133. Derhalben habe ich im Anfang dieses Capitels gesagt, daß die göttlichen Verheißungen den Glauben erfordern und haben wollen, und daß die beiden Stücke, Verheißung und Glaube, zusammen gehören: es sind correlativa, wie man es in den Schulen heißt. Denn ohne Verheißung können wir nicht glauben, und ohne Glauben ist die Verheißung nichts und außer Kraft gesetzt, wenn wir dieselbige nach der That und nach dem Werk abmessen wollen. Darum ist vonnöthen, daß ein Mensch, der göttliche Verheißung hat, diese Kunst recht und wohl wisse, die Paulus lehrt und beschreibt in dem Exempel Abrahams, Röm. 4, 18.; auf Hoffnung glauben, da nichts zu hoffen ist, und daß man auf ein Ding hoffe und dasselbe glaube, wenn es schon alles wider-sinnisch ist. Denn wo ein Mensch daselbige nicht wohl weiß, wird er nicht bestehen können und wird auch nimmermehr selig werden.

134. Wie wir im Neuen Testament die Verheißung von Vergebung der Sünden und vom ewigen Leben in den Sacramenten und im Evangelium haben. Wenn du aber das Werk an sich ansiehst und dem Urtheil und Verstand des Fleisches und der Augen nachfolgen willst, wirst du urtheilen müssen, daß nichts weniger wahr sein könne, denn das dir von Vergebung der Sünden und dem ewigen Leben im Evangelium verkündigt und verheißen wird. Denn wir sehen vor unsern Augen, daß die Menschen dahinstürben und begraben werden mit großer Schwach, daß sie auch dazu von Würmern gefressen und verzehrt werden, so, daß keines Thieres Aas schändlicher oder ekelhafter ist als der todte Leib des Menschen. Heißt das ein ewiges Leben haben, also jämmerlich vertilgt und in die Erde verscharrt, verzehrt und wieder zu Staub und Asche werden, und dasselbe so gar scheußlich und schändlich, daß, wie man gemeiniglich die verstorbenen Leiber zu malen pflegt, aus dem Gehirn des Menschen Rötheln und aus den Därmen Schlangen sollen geboren werden? Noch soll der Glaube daselbst gewiß schließen, halten und hoffen, es werde eine Auferstehung sein und der Leib werde wiederum lebendig werden.

135. So gar eine große und wunderbare Kunst ist es glauben und hoffen, und wird der Glaube sehr wohl beschrieben in der Epistel an die Hebräer am 11. Cap. V. 1.: „Der Glaube ist eine gewisse Zuversicht deß, daß man hoffet, und nicht zweifeln an dem, das man nicht siehet“; davon auch daselbst viel Exempel erzählt werden. Denn also hat Noah geglaubt, daß er in der Sintfluth würde erhalten werden, da er doch mit seinen Augen sah, daß darin die ganze Welt ersäuft ward und umkommen ist. Und wer diese Kunst nicht weiß, der ist kein Christ, und wird auch kein Erbe der Güter werden, so den Gläubigen im Wort vorgetragen und verheißen werden.

136. Auf dieselbige Weise, wenn man die Historie Davids mit diesen Worten Jakobs vergleicht, scheint sie sehr wunderlich und dieser Verheißung stracks zuwider. Denn wo David nicht voll Glaubens gewesen wäre, hätte er nimmermehr das gethan und gelitten, was in seiner Historie und Geschichte erzählt wird. Siehe, wie er mit so großem Muth gegen Goliath gestritten und denselben plötzlich mit einem Steinwurf niederlegt hat; wie er die prächtigen Worte und das Drohen desselbigen großen Riesen so fest und getrost hat verpotten und verachten können. Sollten wir uns dünken lassen, daß er solches ohne Glauben gethan habe? Ja, also ist David gesinnet gewesen: Ich habe einen günstigen und gnädigen GOtt, derhalben will ich nur getrost fortfahren, und wider diesen stolzen Philister streiten. Er hat zwar GOtt nicht sehen können, wie er gegenwärtig vor Augen gesehen hat den großen Riesen, der dem Zeuge Israel Hohn sprach und ihm einen Kampf anbot; aber David ist so getrost wider den Feind ausgezogen, daß er auch daran nicht gezweifelt, er wollte ihn mit GOttes Hülfe niederlegen und erwürgen; wie die Historie 1 Sam. 17, 36. ff. zeuget.

137. Also wirst du in den andern Historien Davids auch nichts Anderes sehen, denn nur wunderbare Kämpfe und Streite, und sonderlich da beschrieben wird, wie er die Ehe gebrochen und den Uria getödtet hat. Das ist wahrlich eine erschreckliche Anfechtung gewesen, und Alles ist wider sinnisch und wider die Hoffnung zugegangen, da er sich mit den sehr schweren Lastern des Ehebruchs, Mords und Gotteslästerungen verunreinigt, um welcher Sünde willen er sich besorgt, daß er von GOtt nun gar verworfen und verdammt wäre; darum er den 51. Psalm V. 13. betet: „Verwirf mich nicht von deinem Angesichte, und nimm deinen Heiligen Geist nicht von mir.“

138. Derhalben, wenn wir hätten sein Leben und die ganze Historie mit allen seinen Werken sollen gegenwärtig ansehen, so hätten wir nimmermehr glauben können, daß diese Weissagung Jakobs erfüllt wäre. Die Feinde tobten und wütheten allenthalben wider ihn, beide im Lande drinnen und auch von außen her, und war Alles voll Verzweiflung, daß niemand denken konnte, daß er das Königreich hätte erlangen sollen, und wie er dabei bleiben möchte, da er es schon überkommen hatte. Doch riß er hindurch, und überwand die Philister, Ammoniter, Moabiter, Syrer, und wohin er sich nur wendete, da hat er doch immer wider seine Feinde gesiegt.

139. Darum muß man dies geistlich verstehen, also daß der Glaube fest an der Verheißung halte, wie David selbst sagt Ps. 21, 2.: „Herr, der König freuet sich in deiner Kraft“; item im 44. Psalm V. 5—7.: „GOtt, du bist derselbe mein König, der du Jakob Hülfe verheißest. Durch dich wollen wir unsere Feinde zerstoßen, in deinem Namen wollen wir untertreten, die sich wider uns setzen. Denn ich verlasse mich nicht auf meinen Bogen, und mein Schwert kann mir nicht helfen“ etc. Er ist ein König des Glaubens und sein Reich ist auch ein Königreich des Glaubens, und wiewohl er ein zeitlicher König ist, so ist er doch gleichwohl allenthalben mit göttlichen Verheißungen gerüstet und umgeben gewesen.

140. Und dies ist die rechte Weise, wie man die Historien der Heiligen lesen soll, welche allen Schriftstellern weit vorgezogen werden sollen, darin die Geschichten beschrieben sind von Hannibal, Scipio und Alexander dem Großen, und von andern Heiden mehr. Welche Historien zwar bei den Griechen und auch bei den lateinischen Schriftstellern hoch gerühmt werden, aber doch mit diesen Historien gar nicht zu vergleichen sind; denn es mangelt ihnen an dieser geistlichen Ehre und sie haben keine Verheißung, und darum sind sie von ungefähr und gar ungewiß. Aber wiewohl in den Historien der Heiligen geringe und schlechte Dinge, wie es scheint, beschrieben werden, sind sie doch in Wahrheit sehr groß und herrlich an sich selbst. Denn hier wird nicht ein Volk, oder die Philister und Moabiter etc. überwunden, sondern der Fürst dieser Welt, der dies kleine geringe Reich verderben und vertilgen wollte. Derhalben sind es vor GOtt sehr herrliche und köstliche Historien; denn sie sind wichtiger und haben einen größeren Vorzug, denn alle andern Schriften, welche ungemein sind, und in sich solche Geschichten begreifen, die zufällig sind, vornehmlich vor der That, und können davon nicht eher rühmen, wie sie gerathen werden, oder was sie sollen für einen Ausgang gewinnen, denn allein wenn sie schon geschehen sind.

141. Auch das Königreich des Volks Israel wird mit so großer Schwachheit und vieler Verwirrung angefochten, daß es den Anschein hat, als würde es alle Augenblick dahin fallen, vornehmlich aber wenn die Sünde, und die Strafe, so sie dafür haben leiden müssen, herein gefallen sind, als, nach dem Ehebruch Davids mit Bathseba und sonst oftmals auch mehr. Doch bleibt es noch immer beständig in solchen schweren Anfechtungen um der Verheißung willen. Denn „Juda ist ein junger Löwe“, der vom Raub kommt, oder der große Beute von den Feinden mitbringt, nicht allein da er den Moabitern, Ammonitern und Syrern viel Silber und Gold genommen, damit hernach der Tempel gebaut worden ist, sondern es sind noch viel trefflichere Historien gewesen wider den Teufel, der ihm widerstrebt und sich mit ganzer Macht und Gewalt wider Davids Reich gelegt hat.

142. „Er hat niedergekniert“, sagt Jakob, „und sich gelagert“, das ist: Er hat ein wohlgeordnet Reich, darin er auch etlichermaßen Friede und Ruhe gehabt; „wer will sich wider ihn auflehnen?“ Wer will den jungen Löwen aufwecken? Dessen hat sich zwar nie Einer unterstanden, daß es ihm darüber wohl gegangen wäre, es sei gleich Moab, oder Ammon, oder Absalom, oder auch sonst jemand anders gewesen: Allen, die sich wider ihn aufgelehnt haben, ist es übel geraten, sie haben dabei kein Glück gehabt. Denn er ist ein König der Verheißung und des Glaubens, welcher mächtiger ist, denn die ganze Welt, Teufel und Hölle sein kann.

\quellseite{505}

143. Bis hieher aber gehen die Juden in diesem Verstande mit uns und lassen sich die Auslegung leichtlich gefallen um ihrer Ehre willen; aber das nun weiter im Texte folgt: „Es wird das Scepter von Juda nicht entwendet werden“ etc., haben sie mit ihren Glossen also schändlich verfälscht und zerrissen, daß kein anderer Ort in der ganzen heiligen Schrift von den bösen verzweifelten Buben also schändlich verderbt und befudelt ist.

\begin{quote}
\textbf{V. 10. Es wird das Scepter von Juda nicht entwendet werden, noch ein Meister von seinen Füßen, bis daß Silo kommen; und demselben werden die Völker anhangen.}
\end{quote}

144. Es ist die Sache und der Verstand dieser Worte klarer, denn daß er könnte verleugnet oder verhehlet werden; und wir haben eben davon auch gehandelt in unserm Buch wider die Juden. Aber die Juden streiten darüber, als wären sie unsinnig, ebenso, wie sich heutiges Tages der Papst wider das sehr helle und gewisse Licht des Evangeliums legt. Solche Leute sind gar nicht Menschen, sondern Teufel. Denn daß man sich davon abwendet, das verachtet, und dem feind ist, was Einer nicht versteht, oder etwa durch Irrthum oder unwissentlich dafür hält, daß es nicht wahr sei, das ist menschlich: aber die offenbar erkannte und bekannte Wahrheit verfolgen, das ist gar teuflisch. Derhalben sind es nicht menschliche, sondern teuflische Verfälschungen, damit die Juden diesen Text verderbet haben. Wie die Papisten und groben Esel zu Rom, durch ihr eigen Urtheil und Gewissen überwunden, wohl fühlen und verstehen, daß sie eine böse gottlose Sache vertheidigen, und hören doch nicht auf, die Unsrigen zu verfolgen und zu tödten. Da handelt nicht ein Mensch wider den Menschen, sondern der Teufel und die Fürsten der Luft streiten wider uns. Es ist wahrlich ein erschrecklich Ding, wider die erkannte Wahrheit wissentlich und vorsätzlich streiten, also daß du selbst bekennen und sagen mußt: Ich weiß wohl, daß die Leute eine gute Sache haben; doch wollen wir dieselben auch selbst mit Gott und seinem Wort vertilgen helfen.

145. Also können die Juden nicht leugnen, dürfen auch darüber ihr Maul nicht aufthun, daß wir aus diesem Text schließen, daß das Reich Judas und der Stadt Jerusalem nun fünfzehn hundert Jahre darnieder gelegen sei; welches ja die offenbar Wahrheit, und heller ist, denn das Mittagslicht sein kann. Dieses Scepter Judas, des jungen Löwen, wie im Text steht, welches der Löwe bestellt und mit Recht eingenommen hat, wird bestehen und von ihm nicht entwendet werden, bis daß der Held komme. Wenn aber der kommen wird, so wird das Scepter auch aufhören und nicht länger währen.

146. Nun ist es offenbar und jedermann kund, daß diese Prophezeiung Jacobs vor fünfzehn hundert Jahren ist erfüllt worden. Darum verachten wir nun sicherlich und getrost der Juden Gespött, daß sie wider uns treiben. Denn die Sache ist gewiß, klar und hell, und ist dazu auch mit augenscheinlicher Erfahrung genugsam bestätigt. Das Reich oder Scepter ist gefallen, die Juden sind zerstreuet über die ganze Welt, darum ist gewißlich der Messias gekommen, gegen welches Zukunft das Scepter von Juda sollte genommen werden.

147. Dieweil sie aber dies nicht können leugnen, fahren sie doch wider die Worte, so an diesem Ort stehen, heraus, als wären sie gar unsinnig, und versuchen, wie sie nur können und mögen, dem zu entgehen, daß sie uns ja nicht beifallen, sondern daß sie viel lieber diesen Artikel setzen und sagen: Messias wird erst nach fünfzehn hundert Jahren kommen: so doch an diesem Ort das Gegenteil gesagt wird, daß er kommen werde, wenn das Scepter von Juda werde genommen werden; wie das die Erfahrung klarlich ausweiset. Und das Reich selbst ist mit dem Stamm ausdrücklich gemeint. Jerusalem aber, welches der königliche Sitz war, regieret jetzt nicht, hat nun innerhalb fünfzehn hundert Jahren kein Scepter gehabt, und ist gar kein Anzeichen vorhanden, daß sie je wiederum zum Königreich oder Regiment kommen könnten.

148. Aber das sind wahrhaftig des Teufels Werke, wissen, bekennen und verstehen, daß ein Ding wahr ist, und doch dasselbe nicht annehmen wollen oder öffentlich bekennen, ja, demselben vielmehr unverschämt und halsstarrig widersprechen und sagen: „Wir wollen von deinen Wegen nicht wissen“, Hiob 21, 14.; item Luc. 19, 14.: „Wir wollen nicht, daß dieser über uns herrsche“, und wie im 2. Psalm V. 5. folget: „Er wird einst mit ihnen reden in seinem Zorn, und mit seinem Grimm wird er sie schrecken.“ Aber denselben Drohen glauben sie auch noch nicht, bis ihnen der Glaube in die Hand kommt; wie den Juden das auch endlich widerfahren ist, was ihnen Gott zuvor gedroht hatte. Eben dasselbe wird der Papst mit seinen Gesellen endlich auch erfahren. Aber, wie ich gesagt habe, diese Sache ist so gewiß, daß es keines Beweises bedarf. Derhalben wollen wir nun den Text besehen.

149. Mich gelüstet aber nicht, alle Verfälschungen der Juden zu erzählen, welche vom Teufel hergekommen sind. Es ist ein solch Stürmen und Wüthen wider diesen Text, daß es Wunder ist. Der Satan kann das Licht nicht vertragen, und wo er es etwa sieht aufgehen, stürmt er darüber mit ganzer Macht, daß er es nur auslöschen möge, da wirft er mit Steinen und Dreck zu. Wir setzen den Juden entgegen die fünfzehn hundert Jahre der Erfüllung; die dienen uns und reimen sich mit den klaren Worten dieses Textes.

150. Sie haben aber, ohne andere Lügen-theidinge, vornehmlich zwei Argumente, die sie treiben. Denn erstlich, während wir das Wort „Scepter“ oder Ruthe von einem Reiche verstehen, wie im 45. Psalm V. 7. auch steht, da er das gerade Scepter ein Scepter des Reichs nennt: so deuten sie es, daß es Keule heiße, das ist, Strafe, Angst und Verfolgung. Und die demselben Verstande folgen, gefallen sich selbst darin sehr wohl und werden für die Vornehmsten und Besten gehalten. Aber der Text, so kurz vor diesen Worten steht, ist dem Verstande zuwider: Der Löwe ist ein solch Scepter, wider den sich niemand legen darf. Derhalben redet er mit rechten von einem solchen Scepter, das etwas leiden soll, sondern vom Reich des Löwen, das jetzt schon aufgerichtet, fest, und wider alle Pforten der Hölle unwiderwindlich ist, daß sich niemand wider dasselbe legen konnte. Darum sagt er: „Es wird das Scepter von Juda nicht entwendet werden, bis daß der Held komme“; und setzt noch dazu: „noch ein Meister von seinen Füßen.“ Dasselbe Wort, mechoqek, heißt, Gesetzgeber, Regent oder Fürst; man mag es einen Kanzler, Hofrichter oder Schreiber heißen, das sind, die Gesetze machen und Mandate schreiben. Es kommt her von dem Worte chok, welches im 2. Psalm V. 7. steht: „Ich will von einer solchen Weise predigen.“ Darum wird das Reich und die Fürsten sammt ihren Räthen bleiben. Und wenn schon das Wort „Ruthe“ oder „Scepter“ so sollte verstanden werden, daß es Verfolgung heißen müßte (wie denn die Könige von Egypten und Assyrien dieses Volkes Ruthe gewesen sind), so will doch das Stück, so im Text allernächst auf diese Worte folgt, diese Auslegung nicht leiden, nämlich, da er sagt: „noch ein Meister von seinen Füßen“, das ist, Juda wird Räthe behalten, oder wie sie von den Evangelisten genannt werden, die Hohenpriester und Aeltesten des Volks. Welches Wort aus Mose genommen ist, der zwei und siebenzig Männer aus den Aeltesten in Israel verordnet, daß sie ihm das Volk sollten regieren helfen; und die sind darnach für und für im Königreich Israel geblieben, auch unter den Königen, und die Juden haben sie genannt Sanhedrin, den großen Rath.

151. Wiewohl sie nun eine Verdrehung des ersten Wortes erdenken, können sie doch das andere Wort auf keinen andern Verstand ziehen, und bleibt dieser Verstand und Meinung gewiß und fest. Das ordentliche Reich wird von Juda nicht genommen werden, noch ein Meister von seinen Füßen. Da muß sich nun der Leser nebenbei erinnern lassen der sonderlichen Art und Weise zu reden, so die Hebräer haben; welche darnach der heilige Paulus auch gebraucht in der Apostel Geschichte am 22. Cap. V. 3., da er sagt: „Ich bin erzogen in dieser Stadt zu den Füßen Gamaliels, gelehret mit allem Fleiß im väterlichen Gesetz“ etc. Denn die Räthe, die das Volk lehren und regieren, haben ihre Zuhörer zu Füßen sitzen. Derhalben sagt nun Jakob: Es wird in diesem Reich an solchen Meistern und Regenten nicht mangeln.

\quellseite{506}

152. Aber jetziger Zeit ist nun Alles miteinander gar verwüstet. Wo ist jetzt Jerusalem? Wo ist jetzt der große Rath und die Zuhörer? Dieselben haben die Juden nun fünfhundert Jahre nicht gehabt. Derhalben ist auch ohne allen Zweifel der Messias gekommen.

153. Das Wort schiloh ziehen sie auch hin und her und deuten es auf mancherlei Weise. Aber die Meinung des Bernhard Ziegler gefällt mir wohl, der da sagt, daß es herkomme von schiljah; wie einer von den Rabbinen mit diesem Argument die Gottheit Christi anficht und darüber streitet, nämlich, da er also sagt: Dieweil er vom Haupt bis zu den Füßen in seiner Mutter Leibe, schiljah, gewesen, wie hat er denn können Gott sein? Wir erkennen es die Nachgeburt, wie 5 Mos. 28, 56. 57. steht: „Ein Weib unter euch, das zuvor zärtlich und in Lüsten gelebet hat“ etc., „die wird die Aftergeburt verschlingen, oder dieselbe ihren Söhnen und ihrer Tochter ver gönnen“, das ist, die Kinder, die da Früchte sind ihres Leibes. Daselbst wird gesetzt continens pro contento, wie man es in den Schulen heißt, das ist, die Nachgeburt für die unmündigen kleinen Kinder. Darum gefällt mir diese Etymologie, das ist, die Ableitung des Worts, welche die Grammatik mitbringt, sehr wohl.

154. Aber die Andern werfen uns vor und sagen: j und der andere Buchstabe h, welcher an das Ende gesetzt wird, die zwei Buchstaben sollen Zeichen sein, dabei man das weibliche Geschlecht des Wortes erkennt, in welchem Geschlechte das Wort schiloh hier gesetzt wird. Aber das sich nicht nicht an; denn die Lateiner haben viel Wörter männlichen Geschlechts, sonderlich was Eigen- oder Taufnamen sind, die sich doch enden auf die Endung der Wörter, so weiblichen Geschlechts sind, als, Cotta, Cinna etc. Darum gefällt mir es besser, daß man schiloh lese denn schilo.

155. Denn etliche von den Rabbinen lesen das Wort schiloh auch aus für die Stätte der Hütte; welche Hütte, wie wir wissen, von den Zeiten Josuas her bisweilen in Silo, bisweilen anderswo gewesen ist; wie der Herr zu David sagt, da er ihm wollte ein Haus bauen, 2 Sam. 7, 6. 7.: „Habe ich doch in deinem Hause gewohnt, seit dem Tage, da ich die Kinder Israel aus Egypten führete, bis auf diesen Tag; sondern ich habe gewandelt in der Hütte und Wohnung, wo ich mit allen Kindern Israel hinwandelte“ etc. Da gibt er zu verstehen, daß die Hütte vor Davids Zeiten nicht an einem bestimmten Orte oder Stamme gewesen sei. Derhalben ziehen sie dies Stück im Text: „bis daß der Held komme“, auf den König Saul, auf David und Jerobeam, welche nach der Zerstörung Silos gelebt haben. Aber die Sache an sich und die ganze Historie ist wider diese Auslegung, wie sie selbst genugsam lernen, nur daß sie die einfältigen unerfahrenen Leute in diesen Worten und rechtem Verstand derselben gern irre machen wollten.

156. Denn das Wort schiloh heißt ein Sohn des Weibes oder der Aftergeburt, und endet sich, wie die Wörter thun, die weiblichen Geschlechtes sind, wie man in der Grammatik redet; und das um der ersten Verheißung willen, so Adam gegeben ist, 1 Mos. 3, 15.: „Ich will Feindschaft setzen zwischen dir und dem Weibe, und zwischen deinem Samen und ihrem Samen. Derselbe soll dir den Kopf zertreten“; und wie Eva hernach sagt Cap. 4, 1.: „Ich habe den Mann, den Herrn.“ Da wird ja klarlich gesagt, daß der Same des Weibes, das ist, der Sohn der Jungfrau kommen und die Gewalt und das Reich des Teufels zerstören werde.

157. Diese Beschreibung haben die heiligen Väter und Propheten mit Fleiß angesehen, und haben dieselbe noch deutlicher gemacht, wie Jesaia Cap. 7, 14. sagt: „Siehe, eine Jungfrau ist schwanger, und wird einen Sohn gebären“ etc. Dieselbe haben Abraham und Isaak auch geübt und getrieben, dazu ihren Nachkommen immer eingeprägt, daß sie denselben gemein und sehr wohl bekannt werden möchte.

158. Und auf diese Weise redet Jakob hier auch, da er sagt: Dies Löwenreich und dieser Löwe, so von Gott selbst geordnet und in das Reich gesetzt ist, der wird stehen und regiert werden von Fürsten und Aeltesten des Volks und von seinen Vätern, bis daß der Jungfrau Sohn komme, das ist, der Sohn des Leibes, der Nachgeburt, wie man sagen möchte. Wir sind nicht Kinder des Leibes, nicht Söhne der Nachgeburt, sondern sind Mannessame, so aus den Lenden des Vaters in die Nachgeburt gekommen ist. Er aber wächst und wird geboren ohne männlichen Samen; wie der Engel sagt zu Maria Luc. 1, 35.: „Was von dir geboren wird“; er sagt nicht, was von Joseph geboren wird. Dein Sohn wird einer Jungfrau Sohn sein, wird eine Frucht der Nachgeburt sein, wird ein schiloh sein, das ist, nur eine Frucht des Leibes seiner Mutter, wird keinen Vater auf Erden haben.

159. Derhalben, so siehe nun, wie schöne herrliche Prophetiaten die Väter aus diesem einigen Wort genommen, beide vor und nach der Sintfluth, welche alle mit den Propheten den Messias auch also beschreiben, daß er wahrhaftig Davids Blut sei, aber also, daß er allein ein Sohn der Nachgeburt sei; und haben auf das dritte und vierte Capitel dieses 1. Buches Mose gesehen, da Cap. 3, V. 15. gesagt wird: „Der Same des Weibes wird der Schlange den Kopf zertreten“; und Cap. 4, 1.: „Ich habe den Mann, den Herrn.“ Und ist wahrlich eine treffliche Stelle in der Schrift, und sind sehr schöne Argumente, damit unser Glaube bestätigt und gestärkt wird, und wird auch damit angezeigt, was die heiligen Väter gelehrt haben vom Anfang bis auf die Propheten und Christum selbst, ja, auch bis auf unsere Zeit und bis an das Ende.

160. Darum so verwerfen wir und verdammen gar die Unsinnigkeit der sehr bösen Menschen, die sich unterstehen, uns diese Lehre zu verdunkeln. Denn Jakob hat nicht vergeblich ein sonderlich Wort gebraucht, sondern er hat diesen Sohn rühmen wollen, der ihm und den andern Vätern verheißen war. Und wir genießen heutiges Tages auch von Gottes Gnade des hellen Lichts des Evangeliums, dadurch das Wort schiloh erklärt und ausgelegt wird, daß es heiße der Jungfrau und auch Gottes Sohn zugleich. Diese Grammatik gefällt mir wohl und ich beruhe darin mit höchstem Vertrauen.

161. Das letzte Stück in diesem Text ziehen die blutdürstigen Rabbinen auf ihre bösen Lüste und Grausamkeit, da sie nur mit eitelm Mord und Todtschlagen schnauben. Denn sie erklären es also: Unter dem Reich des Messias wollen wir alle Heiden unter uns bringen, wollen alle Könige der ganzen Welt todtschlagen und umbringen, wie der Papst und Türke auch thun. Darum verdolmetschen sie diesen Text ohne alle Ursache und stracks wider die Grammatik also: Die Völker werden sich vor ihm entsetzen und er wird sie schwächen; dafür daß im Hebräischen steht: „Denselben werden die Völker anhangen“, und Zucht haben.

162. Es will aber die Schrift nicht haben, daß die Könige oder Regenten sollen getödtet werden, sondern sie sollen bleiben, und sich unterrichten und züchtigen lassen im Reich Gottes; wie der 2. Psalm V. 10. sagt: „So laßt euch nun weisen, ihr Könige, und laßt euch züchtigen, ihr Richter auf Erden“ etc. Das hebräische Wort Iakäh heißt, hören, gehorchen, und gehorsam sein, und nach dem Chaldäischen ist es verdolmetscht: Die Völker werden ihm gehorsam sein und gehorchen; und sie verfälschen doch gleichwohl den rechten natürlichen Verstand dieser Worte wider dieselbe chaldäische Dolmetschung; wie denn dieselbe unserer Meinung auch beifällt, und das Wort schiloh klarlich auslegt, daß es der Messias sei.

163. Darum ist dies ein schöner goldener Text und der sonderlich wohl zu merken ist, nämlich, daß das Recht Christi nicht ein solch Reich sein wird, wie Davids Reich gewesen ist; welches zwar bis auf Christum gewähret hat. Aber von David wird 1 Chron. 29, 3. gesagt, daß er ein Kriegsmann gewesen sei und viel Bluts vergossen habe, und darum dem Herrn kein Haus oder Tempel habe bauen können. Derhalben hat das Reich, so mit Waffen oder Kriegsrüstung, mit Schwert und äußerlicher Gewalt ist regiert worden, jetzt aufgehört, und das Reich schiloh, so darauf gefolgt, ist nicht ein solch gewaltig oder blutiges Reich, sondern besteht allein darin, daß man das Wort höre, demselben gehorche und glaube, dadurch es regiert oder verwaltet wird; wie im 110. Psalm R. 2. gesagt wird: „Der HErr wird das Scepter deines Reichs senden aus Zion“; item Jes. 2, 3.: „Von Zion wird das Gesetz ausgehen, und des HErrn Wort von Jerusalem“; desgleichen Röm. 1, 16.: „Das Evangelium ist eine Kraft Gottes, die da selig machet alle, die daran glauben“; item im 68. Psalm V. 34.: „Siehe, er wird seinem Donner“, das ist, seiner Predigt, „Kraft geben.“ In dem Worte ist göttliche Kraft und Weisheit; wer demselben glaubt, der wird selig werden.

\quellseite{507}

164. Derhalben ist es ein sehr gewaltig Reich im Worte wider den Tod, die Sünde und Teufel, dazu wider alle ihre Gewalt und Tyrannei, daß da Kraft hat zu helfen, zu erretten und zu schützen zur ewigen Seligkeit.

165. Von den Dingen wissen der Juden Rabbinen nichts, Papisten und Türken auch nichts: uns aber gebühret, dasselbe den Leuten fleißig einzuprägen, und daß wir diesen trefflichen Unterschied wohl merken und behalten zwischen dem Reiche Christi und den andern Reichen und auch Davids Reich. Denn so viel hat Jakob sagen wollen: Das Reich Davids, meines Sohnes, das ohne Schwert und leibliche Waffen nicht wird können regiert werden, wird nicht allzeit währen, sondern es wird auf dasselbe das Reich Schiloh folgen, welches allein mit dem Worte wird regiert werden. Wie Christus Marc. 16, 15. sagt: „Gehet hin in alle Welt, und prediget das Evangelium aller Creatur.“ Denn das Wort ist sehr kräftig, welches helfen und erretten kann aus der Hand des Todes und Teufels, dazu von der Gewalt der Hölle, und kann die Leute versehen in das Reich Gottes. Derhalben sollen die Völker diesem Könige gehorchen, das ist, sie sollen mit dem Worte regieret werden: es wird mit Predigen zugehen, das wird das Wahrzeichen sein, damit das Reich Christi unterschieden wird von den andern Regimenten dieser Welt, welche mit dem Schwerdte und leiblicher Gewalt regieret werden. „Die Waffen unserer Ritterschaft“, spricht Paulus 2 Cor. 10, 4., „sind nicht fleischlich, sondern mächtig vor Gott, zu zerstören die Befestigungen“ etc. Denn so sagt ein Diener des Evangeliums: Ich absolvire dich, ich taufe dich, ich declarire und erkläre dich, daß du ein Kind seiest des Reichs Christi; ich verkündige dir Vergebung deiner Sünde, ich verheiße und biete dir an den Sieg wider den Teufel, Glaube, so wirst du selig werden.

166. Derhalben redet Jakob sehr eigentlich, als der vom Heiligen Geiste erleuchtet ist, und hat von ferne her gar ein scharf Gesicht gehabt, da er sagt: Velo jikkhat ammim: das ist: Die Völker werden ihm anhangen, und ihm gehorchen. Und steht eben dasselbe Wort auch in den Sprüchen Salomonis am 30. Cap. V. 17.: „Ein Auge, das den Vater verspottet, und verachtet der Mutter zu gehorchen, das müssen die Raben am Bach aussticken.“ An welchem Orte in der gemeinen lateinischen Dolmetschung (denn sie ist nicht von Hieronymus) falsch stehet: partum anstatt auditum matris, das ist: Und verachtet die Geburt seiner Mutter, anstatt: „Und verachtet der Mutter zu gehorchen.“ Aber Salomo redet daselbst von dem vierten Gebot, und von der Ehre, so die Kinder ihren Eltern zu erzeigen schuldig sind. Und dienet sehr wohl den Verstand dieses Textes damit zu bestätigen, daß das jikhat so viel heißt als Gehorsam. Denn das Evangelium ist ein Gehör, wie Jesaia redet, da er Cap. 53, 1. sagt: „Wer glaubt unser Predigt?“ Dieser schiloh wird nicht regieren mit Schwert, mit Feuer, oder sonst mit äußerlicher Gewalt, sondern mit Gehör und Lehre des Glaubens; und werden ihm nicht allein die Juden gehorchen, sondern alle Völker der ganzen Welt.

167. Das hebräische Wort, am, ist ein wohlbekannt Wort, welches nach meinem Verstande Volk heißt, aber mit Rücksicht auf GOtt; wie Hosea am 1. Cap. V. 9.: „Ihr seid nicht mein Volk (ammim), so will ich auch nicht der eure sein.“ Ich will das mein Volk heißen, das mein Volk nicht ist etc.; daß man es also zusammen verstehe, Deus populi et populus Dei, das ist, GOtt des Volks und das Volk GOttes. Also sagt Jesaia am 31. Cap. V. 7. 8. 9., daß alle Götzen über die ganze Welt sollen aufgehoben und weggeworfen werden, und daß in aller Welt eine Einigkeit soll werden im Gehorsam und Dienste des einigen GOttes. Denn also sagt daselbst der Prophet: „Zu der Zeit wird ein jeglicher seine silbernen und güldenen Götzen verwerfen, welche euch eure Hände gemacht hatten zur Sünde. Und Assur soll fallen, nicht durch Manns Schwerdt, und soll verzehret werden, nicht durch Menschen Schwerdt; und wird doch vor dem Schwerdt fliehen, und seine junge Mannschaft wird zinsbar werden; und ihr Fels wird vor Furcht wegziehen, und seine Fürsten werden vor dem Panier die Flucht geben, spricht der HErr, der zu Zion Feuer, und zu Jerusalem einen Herd hat.“

168. Also ist das Reich schiloh ein Reich des Worts; denn er beruf[t] und regiert sein Volk allein mit dem Worte, ohne Waffen und äußerliche Gewalt. Die aber das Wort nicht hören wollen, die gehören auch zum Reiche Christi nicht. Darum soll sich das Volk mit dem Worte ziehen lassen und nicht knechtischer Weise mit Rüthen, mit Gefängniß und Schlägen lassen zwingen; wie in weltlichen Regimenten die Leute mit Gewalt zum Gehorsam gezwungen werden. Und also ist es allerdings auch erfüllet worden, wie es hier von dem Patriarchen Jakob und den Propheten zuvor verkündigt ist.

169. Das Reich der Kirche hat kein Schwerdt geführt, sondern hat seine Kraft allein vom Worte, und zwar von dem mächtigen Worte des Heiligen Geistes, damit GOtt die Herzen gezogen hat, daß sie glaubten. Darum wird hier gar merkwürdig dazu gesetzt im Texte: „Demselben werden die Völker anhangen“, und ihm gehorchen. Das ist so viel gesagt: Das Wort wird dies Reich machen und bestellen; die Leute werden allein durch das Wort versammelt werden und zu diesem schiloh kommen.

\begin{quote}
\textbf{V. 11. 12. Er wird sein Füllen an den Weinstock binden, und seiner Eselin Sohn an den edlen Reben. Er wird sein Kleid in Wein waschen, und seinen Mantel in Weinbeerblut. Seine Augen sind röthlicher, denn Wein, und seine Zähne weißer, denn Milch.}
\end{quote}

170. Mose hat hier den Grammatikern viel Kreuz und Marter gemacht, oder aber sie haben die Plage und Marter sich selbst zugerichtet; wie wir von dem Worte schiloh und dem Gehorsam der Völker gehört haben, da auch unzählig viel Verwirrungen sind wegen des ungleichen mancherlei Verstandes, und auch um der Unwissenheit willen der Ausleger, unter welchen keiner die Sache an sich versteht.

171. Die Juden lassen sich immer träumen vom Schwert und sind blutdürstig, durch die ganze heilige Schrift.

Wie der Papst auch nicht gebendet, die Kirche mit dem Worte zu regieren, sondern mit Gewalt, und beides mit leiblicher und geistlicher Tyrannei, nämlich, mit dem Schwert und Bann. Denn GOtt hat sie dahin gegeben in verkehrten Sinn, daß sie nicht wissen, was die Schrift, was die Kirche, oder was Christus sei. Ja, es hat auch Gregorius, der ein Bischof gewesen ist, ehe denn die Tyrannei des Papstthums angefangen, Christum nicht sehr wohl erkannt, hat auch das Wort des Evangeliums nicht sonderlich verstanden: darum haben die Menschensatzungen über das ganze Papstthum mit großer Gewalt überhand genommen, und ist kein Wunder, daß sie aus der Mater und Verwirrung nicht haben können kommen wegen der Dunkelheit und Unwissenheit in der himmlischen Lehre.

172. Es lehrt es aber die That an sich selbst, daß die Kirche allein durch das Wort und nicht durch äußerliche Gewalt oder Schwert versammelt worden ist; wie wir auch noch jetzt niemand mit Gewalt zur Erkenntniß des Evangeliums zwingen. Und doch werden die Menschen durch das Wort gefangen, und versammelt zu einer Lehre und einerlei Glauben, so viel ihrer sind, die da wollen gottselig und Glieder der Kirche Christi sein.

173. Die Rabbinen aber martern an diesen Orte nicht allein den Verstand, sondern auch die Grammatik oder den Text jämmerlich um des Buchstabens willen, davon sie sagen, daß er überflüssig sei. Derhalben so habe ich bei keinem Ausleger irgend eine Hülfe, und muß allein den rechten natürlichen Verstand dieser Worte suchen und herausklauben.

174. Vor allen Dingen aber verwerfe ich die Auslegung Ebras und der andern Ausleger, auch der ältesten, die diesen Text verstehen vom Stamme Juda, als ob Jakob nach den vorigen Worten nichts mehr von Christo geredet hätte. Denn sie gehen das zu, daß die vorigen Worte geredet werden vom Könige; aber was hernach folgt, das ziehen sie alle auf den Stamm Juda. Es dünkt aber mich sehr ungeschickt Ding sein, wenn wir es dafür halten wollen, daß der Heilige Geist die Rede so bald abgebrochen habe, da er erst angefangen hat von schiloh und Christo zu reden, und daß er sein nicht mehr habe gedenken wollen. Die Glossa ordinaria, wie sie genannt wird, Ebra und Burgensis führen unter einander Streit mit Auslegungen, und will Einer dieses, der Andere ein Anderes haben. Darum muß ich meinem Geist folgen, dieweil ich sonst niemand habe, dem ich folgen könne.

\quellseite{508}

175. Lyra legt es also aus: Es wird der Stamm Juda so fruchtbar und reich sein auch vor der Zukunft Christi, daß sie werden der Eselin Sohn an den Weinreben binden, und daß ein Weinstock oder ein Rebe so viel Wein tragen werde, daß damit ein Esel möchte beladen werden. Aber das ist gar ein unnützes Gewächs. Denn es könnte auch in unsern Landen und nicht in dem Gelobten Lande allein ein Weinstock leichtlich so fruchtbar sein, daß er so viel Wein trüge, damit ein Esel möchte beladen werden. Und ich kann nicht verstehen, wie sich das mit einander reimen wolle, einen Esel an den Weinstock binden, daß das so viel sollte heißen als, den Esel mit Wein beladen. Derhalben, auf daß es nicht so gar ungeschickt und ungereimt sei, sagen sie dazu, es solle hyperbolisch geredet sein, das heißt, daß man ein Ding größer macht, denn es an sich ist; wie es in dem andern Stücke auch sein soll, da er sagt: „Er wird sein Kleid in Wein waschen“; wie sich oben von seinem Glück auch durch diese Figur redet, im 29. Cap. V. 6.: „Da ich meine Tritte wusch in Butter, und die Felsen mir Delbäche gossen.“ Aber dies alles dient gar nicht zum rechten eigentlichen Verstande dieses Textes.

176. Derhalben sollen wir nicht zugeben, daß die Sprüche von Christo also zerrissen werden, von welchem Jakob angefangen hat zu sagen, daß er werde ein Herr sein der Heiden, und daß ihm die Völker über die ganze Welt hören und ihm gehorchen werden. Bei demselben Vornehmen, von Christo zu reden, laßt uns bleiben, und keine fremde Figuren, als, hysteron proteron etc., dichten und hier herziehen. Wo es uns aber am Verstande mangeln wird, wollen wir die Meisterschaft dem Heiligen Geist lassen, nur daß wir nicht zugeben, daß der Text also verwirrt und zerrissen werde; denn ich will lieber bekennen, daß ich es nicht verstehe.

177. Ich verstehe ferner diese Worte also, daß zur Zeit Christi, oder wenn der Messias oder Schiloh kommen wird, dies geschehen werde, daß der Eselin Sohn an den Weinstock solle gebunden werden. Das ist aber nicht wahr, daß im Stamm Juda so überflüssig viel Weinstöcke gewesen sind. Und wenn das schon also gewesen wäre, was hat der Esel mit dem Weinberg zu schaffen? Und er wird ja nicht darum an den Weinstock gebunden, daß er Trauben tragen soll, sondern daß er trinken soll und soll voll Weins werden. Derhalben scheinen die Worte diesen Sinn zu haben, daß der Esel mit Trauben und Wein wie sonst mit anderm Futter solle gefüttert und trunken gemacht werden. Der Esel soll einmal gute Tage haben und Wein saufen.

178. Das hebräische Wort sorek heißt ein auserwählter Weinstock, und ist Eines so viel gesagt als das Andere: „Er wird sein Füllen an den Weinstock binden“, und daß er sagt: „und an den edlen Reben.“ Jesaia am 5. Cap. V. 1. steht auch dasselbe Wort, da er sagt: „Mein Vetter hat einen Weinberg an einem fetten Ort“ etc. Es ist bei den Juden eine sonderliche Art edlen Weinstocks gewesen; daher auch das Wort „Sirach“ kommt. Und ich halte dafür, daß dies der rechte Verstand sei: Zur Zeit des Messias wird eine solche große Herrlichkeit, Reichtum und Fruchtbarkeit sein, daß auch die Esel und Füllen der Eselin Trauben und Weinbeeren fressen werden. Die Esel sollen etwas Liebliches essen, sollen sich voll saufen.

179. Und es ist nicht hyperbolisch geredet, das ist, daß man es größer macht, denn es an sich ist; sondern ist ganz eigentlich geredet, nur daß durch eine Figur die Eselin genannt wird, wie droben „ein junger Löwe“, und will so viel sagen: Die Esel werden zu der Zeit nicht Spreu oder Heu essen, sondern das unter den Früchten das Allerbeste und Köstlichste ist, wie der Wein ist, „der Götter und Menschen fröhlich macht“, etc. In Buch der Richter am 9. Cap. V. 13. gesagt wird. Denn der Wein ist unter allen Früchten die allerbeste in der ganzen Welt, der das Herz des Menschen erquickt und erfreut, davon beide göttliche und natürliche Menschen lustig und fröhlich werden. Sie werden alle fröhlich vom Wein. Und in den Sprüchen Salomos steht geschrieben am 31. Cap. V. 6. 7.: „Gebet stark Getränke denen, die umkommen sollen, und den Wein den betrübten Seelen, daß sie trinken und ihres Elendes vergessen, und ihres Unglücks nicht mehr gedenken.“ Denn der Wein pflegt den Menschen Fröhlichkeit, Muth und Gestalt zu geben. Darum wird er nun also den Esel und das Füllen, und den Sohn der Eselin an den Weinstock und an den edlen Reben binden, daß er daselbst einen Stall habe und gefüttert werde mit auserlesenen Weintrauben, und daß der heilige Esel trunken werde von dem allerbesten Saft.

180. Nach der Natur zu reden ist es sehr unmöglich, daß der Esel an den Weinstock soll gebunden werden, und er hätte dasselbe besser vom Bock oder Hirsch sagen mögen, die zum Weinstock Lust haben. Aber dem Esel dient der Weinberg oder die Trauben nicht, sondern Disteln und Spreu. Derhalben verstehe ich es also, daß Jakob durch den heiligen Geist rede von dem wunderbaren Reich Christi und von seinem Reichtum, in welchem auch die Esel und jungen Esel Trauben essen, und Wein trinken und trunken werden sollen. Wie Paulus sagt zu den Ephesern am 5. Cap. V. 18, 19. 20.: „Und saufet euch nicht voll Weins, daraus ein unordentlich Wesen folget; sondern werdet voll Geistes, und redet unter einander von Psalmen und Lobgesängen und geistlichen Liedern, singet und spielet dem Herrn in euren Herzen; und saget Dank allezeit für alles, Gott und dem Vater, in dem Namen unsers Herrn Jesu Christi“, welches auch von den Aposteln gesagt wird Ap. Gesch. 2, 4.: „Sie wurden alle voll des Heiligen Geistes“, und redeten von den großen Thaten Gottes.

181. Also werden die Gottseeligen voll Heiligen Geistes, daß sie sich nicht enthalten können; sie müssen herausfahren und Gott danken, ihn bekennen und preisen, und daß sie lehren und reden das Wort des Evangeliums. Solche Leute sind die Apostel und Märtyrer gewesen, da sie vom Heiligen Geist trunken waren. Denn das ist die alleredelste Art Weinstock, davon die Weisheit sagt in den Sprüchen Salomos am 9. Cap. V. 5.: „Kommt, zehret von meinem Brod, und trinket des Weins, den ich schenke“, da werdet ihr trinken und trunken werden, aber von geistlicher Trunkenheit.

182. Also verstehe ich diesen Text hier, daß wir in Christo und zur Zeit Christi sollen trunken werden von den reichen Gütern seines Hauses, wie der 36. Psalm V. 9. sagt, das ist, wir sollen den Heiligen Geist empfangen aus dem Wort und Gehör, davon wir andere Menschen werden; wie sich ein trunkener Mensch viel anders hält, denn einer, der nüchtern und hungrig ist. Denn jener ist fröhlich, lacht, freut sich, singt und jubiliert oder jauchzt; diesem aber wird angst und bange, er ist traurig und kann nichts denn klagen.

183. Dieser reichen Güter halben freut sich nun Jakob seinethalben und aller Nachkommen halben, und sagt: „Er wird sein Füllen an den Weinstock binden.“ Als wollte er sagen: Ein solch Reich wird das Reich Christi und meines Schilo sein: er wird nicht regieren mit Gesetz oder Schwert, welches das Instrument oder Werkzeug ist des Gesetzes, und die Leute traurig, mager, matt und verzagt macht, sondern er wird sie erfüllen mit seinen Gütern, er wird die Unterthanen seines Reichs trunken machen und erfreuen, und sie werden gar andere Leute werden, und werden gespeiset werden nicht mit Spreu, sondern mit Weintrauben, und dazu mit den allerbesten Trauben, daß sie voll und voll werden. Aber geistlich sollen sie trunken werden, also daß die Menschen nicht anders meinen oder an ihnen sehen, denn wie sie von den Aposteln urteilen, daß sie voll süßes Weines wären, Ap. Gesch. 2, 13.

184. Denn ist es nicht ein großer Trost vor den Leuten, daß Petrus und die andern Apostel aufstehen, lehren und predigen, da sie zuvor noch nicht um Erlaubniß gebeten bei der ordentlichen Obrigkeit, weder der geistlichen noch weltlichen. Aber sie sind voll Heiligen Geistes. Darum brechen sie also frei heraus und fahren hindurch auf allen Gassen der Stadt Jerusalem, und predigen von dem gekreuzigten Christo, deß Namen doch bei dem ganzen Volk verhaßt und gar unleidlich war. Ja, sie fahren zu und schelten die Juden hart wegen der \unsicher{überthätigen} Gewalt und Grausamkeit, so sie an dem Gerechten geübt, da sie denselben gekreuzigt, und scheuen sich die Apostel in dem allen vor niemand, fürchten sich auch vor der großen Gewalt der Priester und Aeltesten gar nicht. Und endlich, da sie nun gestäupt und gehöhnet waren, und ihnen geboten worden, daß sie gar nicht mehr reden sollten im Namen Jesu, da gehen sie fröhlich von des Raths Angesicht, daß sie würdig gewesen wären, um des Namens Jesu willen Schmach zu leiden.

185. Das wird nunmehr ein Gesetzprediger thun, ja, der wird sich vielmehr fürchten vor Caiaphas, vor den Obersten und Hohenpriestern, dazu auch vor der Versammlung des Volks; da doch dagegen die Apostel aller ihrer Gewalt und Tyrannei trotzen und denken: Himmel und Erde ist unser; denn Christus ist ein König Himmels und der Erde. Sie sind trunken von dem Heiligen Geist, darum reden sie ohne alle Scheu und Furcht vor den Hohenpriestern von den großen Thaten Gottes und erfüllen mit ihren Predigten die ganze Stadt.

\quellseite{509}

186. Darnach aber haben sie auch auf dieselbige Weise das ganze Römische Reich angegriffen, und haben die Abgötterei der Römer, der Griechen und Egypter gestraft, und halten sich in dem allen anders nicht, denn als wären sie Herren der ganzen Welt. Sie treten vor die Könige und Fürsten, und verkündigen ihnen, daß sie sterben und verderben werden, wo sie die Lehre von Christo, der gelitten hat und auferstanden ist, nicht hören und annehmen wollen: darüber schlägt man sie gar redlich todt. Sie verlachen und spotten der Majestät und großen Gewalt des Römischen Reichs und aller Völker; und das thun nicht allein Männer, als, die Apostel und ihre Schüler, sondern auch Weiblein, und sind so stolz und trotzig, daß sie sich vor keiner Marter fürchten und entsetzen: wenn man sie gleich verbrennt, ersäuft, fragen sie darnach gar nichts. Scheint es denn damit nicht, als ob sie trunken und toll seien? Was ist es doch, daß die höchst elenden Leute und Fischer, ja, auch Weiblein und Jungfräulein sich wider eine solche große Monarchie legen und dieselbe angreifen, wie zu der Zeit der Römer Monarchie war, darin so große Gewalt und allerlei Reichthum, Geld und Gut vollauf gewesen? Da hat die Vernunft schließen müssen, sie wären gar toll und unsinnig, und die Weltweisen haben auch anders nicht gedacht, denn daß sie toll und thöricht wären und daß sie der Teufel besessen hätte.

187. Sie sagen von Anastasia, welche eine Jungfrau gewesen, da ihr der Richter gedroht, sie zu tödten, habe sie geantwortet: Sie hätte nichts Böses gethan. Und da er ihr vorgeworfen: Sie habe die Götzen verachtet, welche sie zertreten und verbrannt hätte, da soll sie den Richter verlacht und gesagt haben, wie sie ihre Götter höchlich geehrt hätte; denn sie hätte dieselbigen durch das Feuer gereiniget, und von den Mäusen und Spinnenweben, und vom Staube errettet, damit sie umgeben gewesen, etliche aber habe sie gekocht und den Armen gegeben. Ja, sagte sie, gebt mir eure römischen Götter auch her, mit denen will ich spielen, wie ich mit den Göttern meines Vaters gespielt habe. Denn wo es Götter sind, wie ihr sagt, daß es sein sollen, so laßt sie sich schützen und vertheidigen wider ein armes Weib. Das heißt den Esel trunken machen.

188. Das ist eine Trunkenheit des Glaubens und der Verheißungen des Heiligen Geistes, wie Paulus es auslegt, den er über uns ausgegossen hat reichlich zur Zeit des Messias, und hat gepredigt und ausgearbeitet, daß sie alle voll des Heiligen Geistes geworden, und den Tod und Teufel verachteten und alles Unglück überwunden.

189. Dies ist vornehmlich zur Zeit der Märtyrer also geschehen. Es ist aber der ganzen Kirche und heutiges Tages uns auch vonnöthen, daß wir trunken werden vom Heiligen Geiste, und daß wir uns vor dem Papst und aller Tyrannen und Teufel grimmigen Zorn, Toben und Wüthen nicht fürchten. Denn das ist die Stimme und Herrlichkeit der Kirche, daß wir an den Weinstock gebunden sind, daß wir voll Heiligen Geistes werden sollen. Christus hat uns berufen und mit dem allerlieblichsten Bande an seinen Weinberg gebunden, das ist, an seine geistlichen Gaben, damit er uns erfüllt, daß wir uns vor der Dinge keinem fürchten, die uns schaden können.

190. Ach, wollte Gott, daß wir die göttlichen Verheißungen recht verstünden und in unsere Herzen prägten, nämlich, daß uns unsere Sünden vergeben sind, daß der Tod und Teufel überwunden, und die Hölle verlöscht und vertilget ist; wie das alles herrlich gerühmt wird im 91. Psalm. Psalm 13.: „Auf den Löwen und Ottern wirst du gehen, und treten auf die jungen Löwen und Drachen“ etc. Wer das fest glaubt, der trinkt gleichsam diese Verheißungen und wird fein lieblich gespeist von diesen Weinblöcken, und ist ein Herr des Todes.

191. Ich habe auch den Satan zum Feinde, der mir so heftig feind ist und mir oftmals nach meinem Leben gestanden hat; wenn er mir aber etwas mit Gewalt zugesetzt hat, habe ich solches leichtlich verachten können, und habe darnach nie gefragt, wenn er mir auch dies Leben nehmen sollte, sintemal ich weiß, daß alle Haare auf meinem Haupte gezählt sind. Aber wenn er uns mit falscher Lehre angreift und daß damit die Schwärmer die armen Gewissen plagen, da haben wir denn viel mehr zu schaffen und ist der Kampf viel schwerer, denn wo man sich wider äußerliche und öffentliche Gewalt wehren soll. Und wir betrachten doch hier auch alle Menschen, und alle Weisheit und Hoffart oder Trotz der Welt, und reißen hindurch. Denn wir sind trunken, nicht von dem Wein, der so ein wüstes, wildes und unordentliches Leben macht; sondern dieses ist eine heilige Trunkenheit, dabei wir uns gewöhnen zur Zucht, und Fleiß anzuwenden, unsere Sitten also zu regieren und zu schicken, daß es uns zu diesem Leben nütze und heilsam sein möge, und daß wir uns vor Hurerei und aller Schande und Unzucht hüten, so, daß wir uns auch des allergeringsten Fluches enthalten, und eine solche Zunge haben, die nur segnen, heilen und trösten könne. Denn rechte fromme Christen suchen und lästern nicht, schwören auch keinen falschen Eid, wie die Trunkenbolde und unsinnigen Leute pflegen; sondern aus dieser geistlichen Trunkenheit folgt ein fein heilsames Leben, wenn sie sich aller Laster und Sünde enthalten, und allein reden und rühmen die Befehle Gottes.

192. Diese Trunkenheit soll nun zur Zeit des Messias regieren: wenn ihm die Völker anhangen und ihm gehorsam werden, so werden sie trunken und fröhlich werden, und dem Teufel und Tode trotzen, also daß sie die Welt für unsinnige und wahnsinnige Menschen halten wird, und werden doch gute, fromme und nützliche Leute sein in menschlichen Dingen, die in dieser Welt auch ihre Feinde lieben, und ihr Leben also anrichten, daß es nicht unordentlich, wüst oder wild sei. Wie denn wahrlich wir Christen ja niemand Schaden thun, sondern begehren von Herzen jedermann nützlich zu sein und zu dienen, wiewohl sie uns spinnfeind sind. Denn wir sind geistlich trunken, haben Vergebung der Sünden, Erlösung vom Tode und Hölle, und spotten des Teufels und Papstes, der da greulicher wüthet und tobt, denn Diocletian und alle Tyrannen gethan haben. Wiewohl der Papst mit seinen Cardinalen viel toller und ungeschickter ist, denn Diocletian war, welcher gleichwohl ein weiser kluger Fürst gewesen ist: diese aber sind lauter grobe Esel, ohne alle Sinne und Verstand. Darum achte ich des Papstes gar nicht, und je mehr er zürnt, je weniger gebe ich drauf. Ich weiß, daß wir auch die Engel richten werden, 1, Cor. 6, 3., und dazu auch Kaiser, Könige und andere böse Würmer mit allen Fürsten und jetzigen Junkern.

193. So ist nun diese Verheißung des Heiligen Geistes bedeutet worden durch das Gleichniß der Trunkenheit in diesem Text, da man von dem allerbesten Wein trunken wird. Denn wie ein trunkener Mensch kühner und muthiger wird, denn er sonst ist; wie der Poet sagt: Tunc pauper cornua sumit, das ist: Dann ist auch ein Armer stolz und muthig, und stellt sich, als wollte er Wunden hauen, so, daß sich Einer allein wider hundert Andere setzen darf: also macht diese geistliche Trunkenheit, die auch heilig und heilsam ist, den Gottlosen viel mehr Muth, welche durch göttliche Kraft aufgeblasen werden durch diesen schier, der den Esel an den Weinstock und edlen Reben gebunden hat, daß sie den Teufel mit allen seinen Schuppen verpotten und unter die Füße treten sollen. Dies, halte ich, sei der rechte natürliche Verstand dieses Textes.

194. Nun muß man aber das auch auslegen, warum er uns Esel und Eselinnen nennt. Hieronymus hat es nur ein wenig berührt und deshalb kürzlich gedacht, da er sagt: Daß nämlich JEsus Christus das Füllen der Eselin, darauf er gesessen, das ist, das Volk der Heiden, zu dem Weingarten der Apostel, so von den Juden hergekommen sind, gethan, und an den edlen Reben, das ist, den auserwählten Weinstock, die Eselin gebunden, darauf er gesessen, nämlich, die Kirche aus allen Völkern versammelt. Diese Weise gefällt mir nicht übel. Ein Esel ist ein einfältiges und dummes Thier, ist aber doch sehr köstlich, darum daß es Christum, der Welt Heiland, trägt. Also sind wir auch der Esel, die Eselin und das Füllen, darauf Christus durch die ganze Welt reitet. Aber dies Gleichniß wird etwas gemildert, wenn du bedenken wirst, daß der Patriarch gesehen habe den Gebrauch über die Weise, so zu seiner Zeit und bei seinem Volk gewesen ist, welches am allermeisten der Esel gebraucht hat. Ein Pferd ist ein kriegerisch Thier und zum Streit gerüstet; wie es beschrieben wird im Buch Hiob am 39. Cap. 23. 22. 23. „Ein Roß spottet der Furcht, und erschrickt nicht, und fleucht vor dem Schwert nicht, wenn gleich wider es klinget der Köcher, und glänzet beide Spieß und Lanze.“ Der Esel aber taugt gar nicht zum Kriege, sondern nur solche Last zu tragen, die im Hause zu tragen ist. Wie Jakob mit seinem Hausgesinde und kleinen Kindern auf Eseln nach Egypten geführt worden ist. Denn es ist ein knechtisch und häuslich Thier, wie man heutiges Tages Esel auch noch also gebraucht, wie man dieselben zu der Väter Zeit gebraucht hat, nur daß wir auch dabei Pferde gebrauchen. Derhalben sagt Jakob hier nicht, daß das Roß an den Weinstock soll gebunden werden, sondern richtet seine Rede auf seine Zeit, darin sie vornehmlich Esel gebraucht haben, solche Last zu tragen, wie man im Hause zu tragen hat.

\quellseite{510}

195. Es setzt aber der Patriarch zu dem, was er hier zuvor gesagt, noch etwas mehr, und sagt: „Er wird sein Kleid in Wein waschen, und seinen Mantel in Weinbeerblut.“ Das ist nicht weniger ungereimt und seltsam geredet, denn daß er von dem Esel gesagt hat, wie derselbe an den Weinstock sollte gebunden werden und dazu an den allerbesten Reben. Ja, dies scheint dem Fluch ähnlicher als dem Segen. Denn was kann ungeschickter oder ungereimter sein, denn die Kleider im rothen Wein waschen? Auf die Weise würdest du sie vielmehr besudeln denn waschen. Und die Schrift lobt die weißen Kleider; wie kann es denn von solchem äußerlichen Waschen nach den Buchstaben verstanden werden?

196. Aber das da folgt, ist noch am allerungereimtesten, da er sagt: „Seine Augen sind röthlicher, denn Wein, und seine Zähne weißer, denn Milch.“ Denn das ist vielmehr ein Gebrechen oder Krankheit der Augen, welches die Aerzte einen Fluß nennen, als, in den Sprüchen Salomos am 23. Cap. V. 29. 30.: „Wo ist Weh? Wo ist Leid? Wo ist Zank? Wo ist Klagen? Wo sind Wunden ohne Ursach? Wo sind rothe Augen?“ (da steht in unserm lateinischen Text: suffusio) nämlich, wo man beim Wein lieget, und kommt aufzusaufen, was eingeschänkt ist.“ Solch Röthe der Augen ist schändlich und ist ein Gebrechen oder Fehler an den Augen. Die immer toll und voll sind, die bekommen rothe Augen; daher ist das deutsche Sprichwort genommen: Du hast auf die Augen getrunken, trinkst einmal auf die Ohren. Darum will sich dieses auf den Stamm Juda nicht schicken, daß es hyperbolisch geredet sein sollte, und kann auch auf den Reichthum oder Ueberfluß nicht gezogen werden, du wolltest denn darunter einen Fluch und Fehler oder Gebrechen verstehen.

197. Man soll aber den Verstand und die Sache, davon ich gesagt habe, daß sie gewiß sei, fleißig merken, nämlich, daß der Patriarch redet von den Verheißungen von Christo oder Schiloh, dem man in seinem Wort gehorchen und anhangen wird. Wie Paulus auch redet Röm. 1, 5.: „Durch welchen wir haben empfangen Gnade und Apostelamt, unter allen Heiden den Gehorsam des Glaubens aufzurichten“; item Röm. 10, 16.: „Aber sie sind nicht alle dem Evangelio gehorsam.“ Das Reich Christi besteht im Gehorsam des Glaubens, und eben dasselbe nennt Paulus auch die Verheißung des Geistes; wie auch Christus Ap. Gesch. Cap. 1, V. 4. befohlen, „daß sie nicht von Jerusalem wichen, sondern warteten auf die Verheißung des Vaters“; item, auch daselbst V. 8.: „Ihr werdet die Kraft des Heiligen Geistes empfahen, welcher auf euch kommen wird.“ Diese Verheißung steht in allen Propheten, daß der Segen durch den Samen des Weibes kommen soll auf alle Völker; und diesen reichen und unaussprechlichen Segen hat er mit diesen ungereimten und dunklen Worten anzeigen wollen; wie es denn wahrlich ein groß Ding ist, daß nicht leichtlich verstanden oder begriffen werden kann, daß der Heilige Geist durch Christum allen denen gegeben wird, die sein Wort hören und an ihn glauben. Denn wer diesen Schiloh hören wird, welcher der gebenedeyte Saame ist, der wird die Verheißung empfangen, so Adam, Abraham und den andern Vätern geschehen ist. Dies ist also gar gewiß und beständig.

198. Nun lasset uns aber sehen, wie das Waschen der Kleider sich dazu schicken wolle. Es ist in jenen Landen sehr gebräuchlich gewesen, daß sie weiße Kleider und von köstlicher Leinwand getragen haben, vornehmlich auf die hohen Feste; wie von dem reichen Manne Luc. 16, 19. gesagt wird. Und Salomo sagt auch im Prediger Cap. 9, V. 8.: „Laß deine Kleider immer weiß sein“ \&c. Und eben dasselbe halten die Türken heutiges Tages auch also, und die Völcker alle, so in den Morgenländern wohnen. Dergleichen sind die Engel auch in der Auferstehung Christi erschienen in weißen Kleidern, Marc. 16, 5., welches ein Zeichen des Lebens und unsterblicher Freude ist. Die Römer haben Purpur und scharlachene Kleider gebraucht. Wie wird aber nun das weiße Kleid im rothen Wein gewaschen werden? das ist zumal ein ungereimt, ungeschickt Ding; wie auch das Binden des Esels an den Weinstock.

199. Darum muß dieses Stück auch verstanden werden von dem Reichthum des heiligen Geistes, daß wenn derselbe wird gegeben werden, Alles mit Milch und Honig und köstlichem Wein fließen soll, darin sie auch die Kleider waschen werden; und wird keines Wassers vonnöthen sein, denn es wird allenthalben Wein vollauf sein. Derhalben, dieweil es so ungereimt lautet, zwingt uns die Noth, daß wir von dem fleischlichen Verstande abweichen müssen, wie Augustin zu sagen pflegt, oder von dem Waschen der leiblichen Kleider absehen, und uns auf diesen geistlichen oder allegorischen Verstand geben. Denn dieses ist ja nach den Buchstaben im Stamm Juda nie geschehen, viel weniger unter Christo.

200. Dies ist aber der rechte und einige Verstand, nämlich, daß Christus die Seinen unter das Kreuz geworfen hat; wie er sagt Matth. 5, 3.: „Selig sind, die da geistlich arm sind.“ Ist aber das ein leiblicher Segen, sterben, gekreuziget werden, mit dem Schwert und Feuer vertilgt werden? Ja, das möchten wir billiger einen Fluch nennen. Darum ist es eine neue Rede in denjenigen Reichthum des Geistes, nach welchem die gottseligen Menschen mit neuen Zungen reden.

201. Dem, wie ich gesagt habe, wenn wir glauben könnten, welch ein groß Ding sei Vergebung der Sünden, auch derer, die noch vorhanden und noch übrig sind in unserm Fleisch, nämlich, daß GOtt uns dieselben nicht wolle zurechnen oder uns darum verdammen; ja, daß er noch also mit uns handeln wolle, als ob gar keine Sünde mehr bei uns übrig wäre, und uns mit dem ewigen Leben krönen; daß er wolle unser Vater und wir sollen seine lieben Kinder und ewig selig sein; wie Marc. 16, 16. gesagt wird: „Wer da glaubet und getauft wird, der wird selig werden“; dann allerdings würden wir dies verstehen. Denn die da glauben, daß sind heilige Esel und Füllen, an den Weinstock gebunden, und sind trunken vom heiligen Geist.

202. Wie die heiligen Märtyrer, und die heiligen Jungfrauen, Agatha, Lucia und viele andere mehr, an den Weinstock sind gebunden worden, haben den Tod für einen Scherz gehalten, die Sünde und Hölle für nichts. Und dieweil sie der Vergebung der Sünden und des ewigen Lebens, dazu auch des guten gnädigen Willens des himmlischen Vaters allerdings gewiß waren, sind sie auch mitten im Tode fröhlich und unerschrocken gewesen.

203. Also sagen sie von Vincentius, dem vortrefflichen Märtyrer, daß da er erst mit Knitteln geschlagen und darnach mit einer Folter gemartert worden war, und gezwungen ward, mit bloßen Füßen durch glühenden Kohlen zu gehen, er gesagt habe, daß er auf eitel Rosen ginge. Also hat er der grausamen Marter gleichsam gespottet, und die glühenden Kohlen verlacht und das feurige Eisen für ein Spiel und Scherz gehalten. Dieses alles zeigt an, wie groß der Reichthum der Verheißung des Schiloh sei.

204. Wenn wir nun auch sterben sollten, oder wenn irgend eine Pestilenz, Theurung, Krieg und aller grimmiger Zorn des Teufels und der Hölle herein fällt, so sollen wir auch, wenn wir glauben, sagen: Ich wollte dich bescheißen, du schändlicher Teufel \&c. Die Kirche hat wahrlich also geglaubt und erzeigt sich noch jetzt mit solchem großen Muth. Denn das heißt recht trunken werden vom Wein des heiligen Geistes und mit neuen Zungen von den großen Thaten Gottes reden.

205. Und dahin haben die Propheten auch gesehen, als, Amos Cap. 9, 13.: „Und die Berge werden mit süßem Wein triefen, und alle Hügel werden fruchtbar sein.“ Es ist ein köstlich Regiment das Reich Christi. Wer da glaubet, daß er einen gnädigen GOtt hat, ja, daß er GOtt zum Vater hat; daß Christus, der Sohn GOttes, den Tod, die Sünde, Hölle und den Teufel vertilgt habe: sollte derselbe nicht fröhlich sein und sich darüber freuen? Ja, er sollte durch eiserne Berge und durch allerlei Widerwärtigkeit hindurch gehen mit unerschrockenem und unüberwindlichem Gemüthe, und sollte es gewißlich dafür halten, daß Alles von eitel Honig, Milch und köstlichem Wein fließe. Ja, er soll noch dazu mit höchster Fröh-

\quellseite{511}
206. Also sollten wir zwar alle gesinnet sein. Denn dieses sind die Verheißungen in Christo Jesu, wie Paulus davon redet, die nicht den Beutel voll Gold oder Silber geben, sondern das Herz voll Freuden, Friede und Sicherheit machen des ewigen Lebens halben: welches alles wir an den Aposteln, Petro und Paulo, und an den Märtyrern, Vincentius und andern mehr, sehen, die sich das gar nicht ansehen lassen, wenn schon Hannas und Caiphas zürnen; ja, sie verachten ihr Drohen ganz und gar, nicht anders als die Engel im Himmel.

207. Und wir lachen heutiges Tages auch des greulichen Wütens des päpstlichen Bannens, so über uns geht. Denn wir glauben an Schiloh, und sind mit unter denen, die ihm anhangen, ihm gehorchen und an ihn glauben. Darum bin ich dessen gewiß, daß mir nichts schaden kann, sondern Alles muß mir zum Guten und zur Seligkeit dienen.

208. Daß wir aber gleichwohl noch so zittern, so erschrocken und furchtsam sind, das ist der Schwachheit Schuld, so noch übrig in uns ist, und kommt nicht vom Glauben her, sondern es ist ein Ueberbleibsel vom alten Adam, der die Sünde fühlt, und sich vor Gottes Zorn fürchtet; vornehmlich in uns, die wir unter dem Papstthum dieses Gebrechen der verderbten Natur durch Gewohnheit und Lehre gemehret haben, da wir also sind gelehrt worden zu zittern, zu zweifeln und vor Christo zu fliehen. Und mir ist es jetzt auch noch schwer, die Lehre des Papstes abzulegen und wegzuwerfen, nicht allein nach dem alten Menschen, sondern auch wegen der Schwachheit des Glaubens, nach welcher ich jetzt noch Scheu trage Christum anzusehen. Und wir haben wahrlich kaum recht angefangen zu hoffen, Christum als unsern Heiland anzurufen, daß er kommen wolle durch Tod, Theurung oder Krieg und uns erlösen; welches wir unter dem Papstthum mit nichten haben bitten dürfen. Wiewohl ich jetzt noch in großer Schwachheit an diese Erlösung durch die Zukunft Christi denke und sie begehre, um der schändlichen und verdammten Lehre willen des Papstes.

209. Also müssen wir auf dreierlei Weise streiten wider diesen Schaden der Erbsünde. Erstlich, muß man lernen die Lehre, welche der Papst gar verdunkelt und der auch die Verheißungen weggenommen hat, daß wir nicht gewußt haben, was Christus wäre, sondern ganz und gar Heiden, Türken und Juden waren. Nachdem wir aber nun Erstlinge geworden sind der Creaturen Gottes, wie Jacobi am 1. Cap. V. 18. steht, müssen wir uns des päpstlichen Sauerteigs allgemach entwöhnen und abthun, welcher uns annoch anhängt, bis daß er gänzlich ausgefegt wird, und müssen uns gewöhnen an den Reichthum und Glauben der Verheißung. Zum Dritten, muß man dann auch das Murren des Fleisches zwingen und inne halten.

210. Ihr aber seid nun so viel seliger, die ihr zu diesem schönen Lichte der reinen Lehre nur die verderbte Natur daher bringt, nicht böse Gewohnheit und gottlose Lehre. Derselben sollen wir Gott dafür danken, daß wir getauft und absolv(i)irt sind auf die Verheißung der Vergebung der Sünden. Wo aber noch übrige Sünden in uns sind, dieselben werden uns doch nicht zugerechnet. Desgleichen sind wir auch berufen auf die gewisse Hoffnung des ewigen Lebens und des Reiches, darauf wir warten, und auf Verachtung des Todes und des Teufels. Lasset uns nur etwas vollkommenlicher trunken werden von dem Troste und Freude der Gnade und Wohlthaten Christi Jesu, um welcher Ursache willen der Esel an den Weinstock und edlen Reben gehangen ist, das ist, an den allerkräftigsten und besten Weinstock, daß er trinke guten Mostkeller und \unsicher{Korseifer}.

211. Auf diese Weise verstehen wir nun dieses Stück auch, daß er sagt: „Er wird sein Kleid in Wein waschen“, nämlich, rothes; denn der Wein im Lande Canaan ist gemeiniglich roth gewesen; und daß er dann weiter sagt: „seinen Mantel in Weinbeerblut“. Denn durch das Blut versteht die Schrift die rothe Farbe, wie 5 Mos. 32, 14. steht: „Und tränkte ihn mit gutem Traubenblut“, das ist, mit dem allerbesten rothen Weine. Das ist die Eigenschaft und Art dieser Sprache; Traubenblut ist rother Wein. Und es ist das eine Wiederholung (Tautologie), wie droben auch mit dem Anbinden des Esels, da er einerlei Meinung auch zweimal wiederholt und mit andern Worten ausdrückt.

212. Was aber dies für ein Waschen sei, davon habe ich gesagt, daß die Ausleger darüber nicht einig sind, und was meine Meinung sei, habe ich angezeigt. Die Glossa ordinaria, wie man dieselbe nennt, legt es geistlich aus von dem Leiden Christi; welches keine böse Auslegung ist, aber doch ist sie etwas fremdartig. Das muß man aber sagen und fleißig behalten, daß wir alle getauft werden in den Tod und Blut Christi Jesu, damit wir gewaschen werden; wie 1 Cor. 6, 11. gesagt wird: „Und solche sind euer etliche gewesen; aber ihr seid abgewaschen, ihr seid geheiliget, ihr seid gerecht worden durch den Namen unsers Herrn Jesu, und durch den Geist unsers Gottes.“ Man muß das Blut Christi auch mit zählen unter die Gaben des Heiligen Geistes, wie gesagt wird 1 Petr. 1, 2.: „Durch die Heiligung des Geistes, zum Gehorsam, und zur Besprengung des Blutes Jesu Christi.“ Denn der Heilige Geist ist nicht gegeben worden, denn nur durch das Blut unsers Herrn Jesu Christi, welcher die Gaben des Heiligen Geistes erlangt hat, und diesen Reichthum des Weins und Traubenbluts, und dazu, daß wir so wunderbarlich erneuert worden sind, daß ist alles erworben durch das Blut Christi. Darum ist es nicht übel geredet, daß wir getauft und absolv(i)irt werden im Blute oder in der Kraft des Bluts des Sohnes Gottes: es muß Christus mit seinem Blute dabei sein.

213. Unter den Papisten ist das Wort noch geblieben, daß sie gesagt haben: Die Sacramente sind aus den Seiten Christi geflossen; denn sie haben ihre Kraft aus den Wunden und Blute Christi; darum ist es ein sehr feines, gutes, und gottseeliges Wort. Und das Evangelium, so wir gehört, und das Anfangen, da man dem Schiloh gehorcht und Gehör gibt, ist eben dasselbe Blut Christi, wie es St. Petrus nennt, 1 Petr. 1, 2., „zum Gehorsam und zur Besprengung des Bluts Jesu Christi.“ Die das Wort des Evangeliums von Christo hören, die werden mit dem Blute Christi, des Sohnes Gottes, besprenget. Es ist alles das schöne rothe Blut Christi. Welches also mag verstanden werden, daß alles durch die Kraft des Blutes Christi erworben sei, was wir haben durch den Heiligen Geist im Wort, in der Taufe, im Abendmahle des Herrn, in der Absolution, im Troste; und alles, was da zur Vergebung der Sünden und dem ewigen Leben gehört, das alles ist das gebened(e)ite Blut des Sohnes Gottes.

214. Derhalben, die diesen Text auf solche Weise auslegen, wiewohl es nicht so gar eigentlich geredet ist, so ist es doch an sich die Wahrheit. Denn das Evangelium, das Abendmahl und die Taufe müssen vom Blute Christi nicht abgesondert werden, dieweil es alles in demselben ist; wie Paulus etlichemal Eph. 1, 7. Cap. 2, 16. sagt: Er hat alles versöhnt und gereiniget durch sein Blut etc. Darnach, eigentlich zu reden, werden wir täglich gewaschen durch das Wort, und daß wir dem Schiloh Gehör geben, ihm gehorchen und anhangen. Wie Paulus ermahnt Col. 3, 9. 10. da er sagt: „Ziehet den alten Menschen mit seinen Werken aus; und ziehet den neuen an“; item, 1 Cor. 15, 53.: „Denn dies Vergängliche muß anziehen das Unvergängliche, und dies Sterbliche muß anziehen die Unsterblichkeit.“ Dieses arme, erbärmliche Leben ist wie ein Kleid, und wie es St. Petrus nennt, eine Hütte, darüber die Heiligen oder Gläubigen hart klagen, 2 Cor. 5, 4.: „Denn dieweil wir in der Hütte sind, seufzen wir uns, und sind beschweret; sintemal wir wollten lieber nicht entkleidet, sondern überkleidet werden“; wir wollten gerne den alten Rock behalten und den Oberrock anziehen.

215. So wird nun dies Kleid täglich gewaschen in rothem Weine, dem allerkräftigsten Safte, oder in dem großen Ueberflusse und Reichthume der Verheißung und Segens vom Samen des Weibes; welcher Segen durch diese Figur oder Gleichniß verstanden wird.

216. Wir werden erstlich trunken gemacht, wie der Esel und Sohn der Eselin vom Weine in diesem Weingarten, das ist, von göttlicher Verheißung im Heiligen Geiste, welche der Weinstock, die Traube und der Wein ist, so uns aufrichtet und stolz macht, und daß wir Verächter des Teufels und Todes werden. Was aber noch übrig ist am alten Menschen, das wäscht er und reinigt es, bis wir wiederum auferstehen werden unverweslich; da werden wir gar rein sein. Inzwischen halten wir uns aufrecht mit Vergebung der Sünden, und haben gewißlich das ewige Leben in der Hoffnung und sind voll Weins oder Heiliges Geistes, und werden abgewaschen in dem Bade, darin der alte Mensch getödtet und gewaschen, und von Tage zu Tage immer erneuert wird.

\quellseite{512}
217. Auf diese Weise haben die heiligen Väter vom Reiche Christi geredet mit vollem Munde und großen Freuden aus der Fülle des Heiligen Geistes. Und daraus kann man leichtlich abnehmen, was ihre Reden und Predigten gewesen sind. Es wird wahrlich Jakob dieses nicht also gelehrt und gepredigt haben allein in dieser letzten Stunde oder Tage seines Lebens, sondern was er von Abraham und Isaak empfangen, dasselbe hat er seine Kinder und Kindeskinder fleißig gelehrt und ihnen eingeprägt, nämlich, wie das Reich Christi sehr herrlich und reich sein würde, und daß bei denen, so geistlich trunken sind, ein Abwaschen und Reinigung von Sünden sein werde.

218. Wie Agatha zu dem Henker soll gesagt haben: Wo du meinen Leib nicht gar zerstören wirst, so werde ich nicht können in das Leben eingehen; also, sagt sie, soll man mit mir umgehen, so werde ich in das Paradies eingehen. Dergleichen hat Vincentius auch die Marter für Spiel und Scherz gehalten. In denselben hat der Geist das versuchte und widerwärtige Fleisch getödtet. Und handelt mit uns auch also, wenn er uns den Türken vorwirft, daß sie uns zerreißen und erwürgen, oder wenn wir sonst mit andern Jammer und Unfall beschweret werden; da wäscht er unsere Kleider im Blute mitten im Feuer, Marter und Pein. Und alsdann sollen wir denken, daß Christus diese Worte zu uns rede: Ich will dich waschen nicht allein mit Wasser, sondern mit dem allerbesten Weine, sofern du nur kannst stille und geduldig sein; halte mir nur stille. Daß man nun diesen Willen Gottes verstehe auch mitten im Tode, heißt das nicht trunken sein?

219. Wir wollen aber nun auf das letzte Stück kommen, da er sagt: „Seine Augen sind röthlicher, denn Wein“. Der ganze Text redet, wie nun oft gesagt ist, von der Trunkenheit und Reichthum der Gaben Christi, und daß wir so wunderbahrlich verändert werden, daß da wir zuvor uns vor dem jüngsten Tage gefürchtet, und den Tod, Leiden und Plage geflohen haben, wir jetzt der Zukunft Christi begehren und schreien: Komm, lieber Herr, komm. Wenn wir geplagt werden, so danken wir GOtt dafür, und werden wie in einem Bade gewaschen mit dem besten Weine. Also behalten wir an diesem Orte eben dieselbe Materie und Meinung, nämlich, daß Jakob hier beschreibe die Fülle, Ueberfluß und Trunkenheit, welche wir empfangen aus göttlicher Verheißung und Gaben des Heiligen Geistes.

220. In der Auslegung des Worts chaklili folge ich Salomo nach, der mich in der hebräischen Sprache dünkt sehr fein zu sein, wie wir droben auch einen Spruch angezogen haben aus Spr. 23, 29., da das Wort chakiluth die Röthe der Augen bedeutet, so aus übermäßigen Weinsaufen herkommt; welches eine sehr wunderliche Mißgestaltung des Angesichts und der Augen ist. Und ich kann es an diesen Orte nicht anders verstehen, denn von gar zu großer Trunkenheit und Vollheit, daher solches Gebrechen der Augen kommt. Denn wie die, so immerdar trunken sind, leiblicher Weise diesen Fehler oder Gebrechen haben, und über dies auch in den Weinhäusern schreien, GOtt lästern, fluchen, Lärm anrichten; wie denn solches an unsern Bürgern und Bauern wohl zu sehen ist, die hin und wieder in den Häusern, da man Bier schenkt, auch etwa unter der Predigt, bei dem Biere weiblich heulen und tören: also macht auch nun wiederum die geistliche Trunkenheit die Leute fröhlich, muthig und auch stolz, daß sie GOtt loben, danken und seine göttliche Gnade und Wohlthaten preisen und rühmen.

221. Derhalben will Jakob hier anzeigen, daß zur Zeit Christi er selbst in allen den Seinen werde rothe Augen haben und einem vollen und trunkenen Menschen gleich sein: er wird aussehen wie ein toller und voller Mensch, nicht vor GOtt, sondern vor der Welt, welche die, so Christo angehören, für Narren und unsinnig hält. Denn wer von dem zukünftigen Leben redet, und alles verachtet, daß sich die Welt hoch verwundert, denselben halten die Leute für toll und unsinnig. Wie die Italiener, wenn sie wollen sehr klug und weise sein, und hören, daß Einer an Christum glaubt und auf das zukünftige und ewige Leben hofft, nennen sie denselben Bon Christian.

222. Also sind wir Christen vor den Augen der Römischen Cardinalle eitel Klöße und Stöcke, tolle, volle Leute. Sehet, sagen sie, wie roth sehen sie aus, die nichts Anderes können, denn daß sie nur von ihrem Christo predigen. Wir werden von ihnen für trunkenen und schändliche Leute, ja, für Narren, Esel und Kinder gehalten, die ihre Kleider mit Wein begossen haben, und trinken noch von dem gestrigen Fraß und Trunkenheit, davon die Augen so roth sind.

Solches wird uns zwar vom Pabste vorgeworfen; und die Märtyrer haben auch dasselbe also hören müssen, als ob sie unsinnig und toll wären, daß sie das Römische Reich also sich unterstünden zu verachten, des Gefängnisses und Marter, so man ihnen angelegt, zu spotten und zu verlachen. Sie halten uns für trunkene Leute, die nicht mehr so scharf sehen können, wie Andere, die viel klüger, weiser und gelehrter in weltlichen Sachen sind denn wir, die nicht so blinzen, wie diejenigen, welchen die Völlerei und tägliches Schlemmen das Gesicht geschwächt und verderbt hat, die Tag und Nacht im Luder liegen und die Augen vertrincken.

223. Die weißen Zähne bedeuten auch Narrheit, Milchsinder, Milchsresser. Sie werden Milch essen, will er sagen, wie die kleinen Kinder pflegen. Aber vor GOtt werden wir mit reicher Milch und Honig gespeiset werden. Denn er meint das nicht, davon sich die Juden träumen lassen, nämlich, daß in dem Stamme so überschwänglich viel Milch sein werde, daß seine Zähne davon, daß sie immerdar so viel Milch essen oder gebrauchen, weiß werden sollen; denn auf die Weise könnten beide Türken und Pabst auch weiß werden: sondern es bedeutet ein stetiges Anhalten (wie er auch haben will, daß das Anbinden an den Weinstock immerdar für und für währen soll), daß wir nimmer aufhören, weiße Zähne und rothe Augen zu haben, oder die Kleider in Wein zu waschen.

224. Denn die Grammatici lassen es sich hier sehr sauer werden und zerbrechen sich heftig den Kopf wegen des Buchstabens j, der überflüssig ist. Aber es bedeutet eine entelechia, wie sie es in den Schulen auf Griechisch nennen, das ist, nicht sowohl eine Geschicklichkeit, als vielmehr eine solche Wirkung, die für und für währet, also, daß das Christenthum ein beständiges Waschen, Binden, und solche weiße Milchfarbe sein soll, die für und für währet, daß wir nimmermehr aufhören Milch zu essen.

225. Das ganze zukünftige Leben ist nicht eine schläfrige müßige Geschicklichkeit; wie die Philosophen einen Unterschied machen zwischen actus und habitus, das ist, zwischen dem Werke und der Geschicklichkeit; und wie sie in den Schulen den Glauben also definirt oder beschrieben haben, daß es eine Geschicklichkeit sei, und daß doch gleichwohl ein solcher geschickter Glaube nichts mehr thue, denn daß er allein das Werk des Glaubens hervorbringe. Das sind nur lauter ungeheure und seltsame Worte, welche sie auch selbst nicht verstanden haben. Und mit diesen ihren Träumen haben sie den Glauben verdunkelt, und haben also die Theologie und Philosophie des Aristoteles untereinander gemengt. Derhalben sollen wir sie immer fahren lassen mit ihren Geschicklichkeiten und Unterschied zwischen dem Glauben und dem Werke des Glaubens.

226. Und wir sollen wissen, daß der Patriarch Jakob an diesem Ort hat ausdrücken und anzeigen wollen, daß dieses Binden und Anderes, davon er sagt, immerdar währen und für und für im Schwange gehen werde. Er beschreibt nicht einen solchen Schilo, der nur einmal binde und darnach bald wieder davon gehe. Wie die Geselsköpfe zu Köln vom Glauben unnützlich plaudern und waschen, und sagen, daß an keinem Ort in der heiligen Schrift gelehrt werde, daß der Heilige Geist durch den Glauben gegeben werde; so doch alle Bücher der Propheten und Apostel dasselbe klärlich sagen.

227. Dieses Gehören und Waschen der Kleider, und die weißen Zähne sind eigentlich ein solch Ding, das immerdar für und für währen soll und muß. Denn haben wir nicht Vergebung der Sünden um Christi JESU willen? Und das ist ein Christ, der allezeit, er schlafe, wache, stehe oder gehe, und es sei gleich in welchem Augenblick, daß es sein kann, in seinem Herzen also gesinnet ist, daß er sagt: Ich glaube an GOtt, den Vater unsers HErrn JESU Christi, welcher sein Blut für mich vergossen hat zur Vergebung der Sünden. Das glaubt er immer ohn Unterlaß, und er könnte oder möchte kein Christ sein, wo er nicht alle Augenblick also glaubte. Diesen Glauben verleugnen die Sophisten zu Zeiten, daß sie verleugnen die Substanz und das rechte Wesen des christlichen Glaubens und Lebens.

228. Derhalben halte ich dafür, daß ich den rechten und eigentlichen Verstand dieses Textes erlangt und getroffen habe. Der Juden ungeschickte Deutung aber, und die Eitelkeit, so sie hier einführen, will ich übergehen und will die Grammatici den Hebräern befohlen haben; denn sie sagen, daß die Buchstaben j und h dreimal überflüssig seien. Ich schreibe mir das nicht zu, daß ich der hebräischen Sprache so gar vollkommen Verstand habe; und ist jetzt zur Zeit niemand und meines Erachtens ist auch bei der Apostel Zeiten niemand gewesen, der diese Sprache recht und vollkommen gekonnt hätte. Es ist genug, daß wir damit umgehen, dieselbige Sprache zu studiren, und darin immerdar wachsen und zunehmen, dieweil wir darin zu keiner Vollkommenheit kommen mögen, so daß wir so reden und urteilen könnten, wie Samuel, David und Jesaja geredet haben.

\quellseite{513}
229. Es scheint aber doch, als habe er eines sonderlichen Worts oder Namens gebrauchen wollen, der von dem Namen der Vorfahren oder eines Volkes herkomme, oder der dieselben bedeuten soll. Wie diese Worte „Eber“, „Gom“ und „Moab“ auch Namen der Personen sind: und wenn du den Buchstaben j dazusetzest, so werden daraus patronymica, gentilia oder possessiva, damit die Väter oder Völker verstanden werden. Eben dasselbe wollte ich wohl hier auch gern sagen, aber ich unterstehe mich nicht, es zu thun, daß das Wort osri so viel heiße als, vinculator, Einer, der immer anbindet; also daß es nicht ein solch Werk sei, das vergehe und aufhöre, sondern daß dieses Binden also verstanden werde, das immer von Einem auf den Andern erbe, welches Binden nie gehe aufgehört, sondern das gleichsam erblich und wesen sei zu allen Zeiten der Apostel und Väter, als Augustinus und der andern.

230. Also verstehe ich auch das Wort chaklili, daß damit angezeigt wird, daß es immer für und für, auch von Einem auf den Andern mit den roten Augen währen soll. Desgleichen, daß in dem Sohne der Geselin Kindesfinder begriffen werden. Aber ich kann davon nichts Gewisses sagen.

231. Eben also ist der Buchstabe h ein solcher Buchstabe und Endung, der das weibliche Geschlecht bedeutet, vornehmlich aber ein Relativ und Affix, wie man es in der hebräischen Grammatik heißt, wie hier steht roch, serech etc. Darum wollte ich es gern also auslegen, daß Christus und die Kirche, welche sein Füllen ist, Eins sind, das ist, der Bräutigam und die Braut sind Ein Leib, wie man sagen möchte, suus et sua, pullus et pulla, sind Eins und einerlei in der Kirche. Mein Geliebter ist mir, das ist, er ist mein, und ich bin sein. Derhalben sind wir Eins. Ich bin Christi und Christus ist meine Gerechtigkeit und Weisheit und Rechtfertigung. Ich bin sein Füllen. Wiederum ist der meine Sünde, der da sagt: Ich bin dein Sünder; und ich sage zu ihm auch also: Du hast meine Sünde getragen; du bist das Lamm Gottes, der du trägst die Sünde der Welt. Hier werden Bräutigam und Braut mit einander verbunden in Einem Geist, und wird Alles unter ihnen gemein, was sie haben.

232. Daher kommt es, daß Christus so sehr klagt im 40. Psalm V. 13., da er sagt: „Denn es hat mich umgeben Leiden ohne Zahl; es haben mich meine Sünden ergriffen, daß ich nicht sehen kann; ihrer ist mehr, denn Haare auf meinem Haupt“ etc.; item im 41. Psalm V. 5.: „HErr, sei mir gnädig, heile meine Seele; denn ich habe an dir gesündiget.“ Da wird er einfach ein Füllen und Esel oder die Kirche; Christus wird sein Leib, und sagt: Ich habe an dir gesündiget, so er doch nie etwas gesündigt hat; er legt die Sünden, so unser sind, auf sich und schreibt sie sich zu. Warum thut er das? Antwort: Er thut es darum, dieweil er also schließet: Es ist meine Braut, ich bin ihr Bräutigam; darum ist ihre Sünde mein und meine Gerechtigkeit ist ihr. Dieses sage ich als Einer, der es also errathen, und nicht, daß ich es also gewiß wollte abgeredet haben, und begehre lieber, daß es mit diesen Buchstaben die Meinung haben möchte, als daß ich es also gar gewiß sein sollte.

233. Wir wollen aber diesen großen, reichen und leiblichen Segen Judas fleißig betrachten, und denselben allen wunderlichen und seltsamen Träumen der Papisten und Scholastiker entgegensetzen, welche nicht allein den Text der Bibel und das schöne helle Licht der Verheißungen mit erschrecklicher Dunkelheit verwirret, daß niemand hat verstehen oder wissen können die Lehre von Christo und seinem Reich, oder was Gutes wir von ihm zu erwarten hätten; sondern wollen jetzt noch mit Gewalt und Schwert das Licht des Evangeliums, so wiederum aufgeht, gar auslöschen und vertilgen. Und sie haben schon eine lange Zeit in der Kirche und ihrem Reich das Schwert und Wort an einander gesetzt und zusammen gemengt, ja, sie haben endlich das Wort gar hingelegt und allein das Schwert behalten; wiewohl hier klärlich das Gegenteil gelehrt wird. Denn in dieser ganzen Rede Jakobs ist gar nichts von äußerlicher Gewalt.

234. Er hat gesagt, daß der Schiloh allein durch das Wort regieren und über seine Kirche herrschen wird; und wir haben droben Erinnerung gethan, daß man solches fleißig merken soll. Dagegen aber rühmt sich der Papst unverschämt und mit großer Hoffart und sagt: Des Kaisers Schwert ist in meiner Hand, und wiewohl es der Kaiser in seiner Scheide stecken hat, darf er es doch nicht zücken, es sei denn, daß ich es ihm erlaube und zugebe. Also haben sie das Reich Christi, welches geistlich ist, in ein äußerlich Reich, das mit dem Schwert umgeht, verkehrt und verwandelt.

235. Darum sollen wir lernen, daß das Reich Christi darin besteht, daß ihm die Völker gehorchen in seinem Wort. Denn Christus hat es also beschlossen, daß er wolle gegenwärtig und kräftig sein durch das Wort; wo dasselbe gehört wird, alsdann will er das Wort bekräftigen durch mitfolgende Zeichen, wie Marc. 16, 20. gesagt wird. Das hat er verheissen, und hält es auch ohne Schwert, ohne Todtschlag, ohne Wunden, Blut und Krieg.

236. Ich werde getauft und allein mit dem Wort unterrichtet, dazu auch absolvirt, ich gebrauche das Abendmahl des Herrn; aber mit dem Wort und durch das Wort ist dabei der Heilige Geist und die ganze heilige Dreifaltigkeit und wirkt also die Seligkeit der Menschen, wie die Worte von der Taufe lauten.

237. Auf diese Weise malt Jakob hier ab ein wunderbarlich Reich, und daß den weltlichen Regimenten gar ungleich ist, welche durch Gesetz und äußerliche Wehr oder Waffen verwaltet werden, und sind auch darum Reiche der Gebote, der Gesetze und äußerlicher Waffen. Dies aber ist ein Reich der Verheißung, da allein Gott gegenwärtig ist und wirket Alles durch das Wort; denn also lehren alle Verheißungen und Propheten.

238. Darum soll man immer fleißig merken den Unterschied der weltlichen Obrigkeit und des Kirchenregiments. Es sollen zwar die Waffen und weltlichen Könige dazu dienen, auf daß im Reich Christi möge Friede sein, das Evangelium zu lehren und auszubreiten: aber dies Reich soll nicht durch Gesetze verwaltet oder regiert werden; denn Gesetze machen keine Christen, sondern das Wort und die heiligen Sacramente, als, Abendmahl und Taufe etc., die machen und bauen das Reich Christi.

239. Das soll man lernen und mit höchstem Fleiß behalten, daß wir dem Papst, Mönchen und Sophisten nicht gleich werden, welche die greulichsten Feinde der Kirche und Gemeinde Gottes sind; wie jetzt die Sophisten zu Rom, Paris und Köln nichts Anderes thun, denn daß sie die Könige und Fürsten wider uns hetzen und rüsten. Und reimt sich eben recht auf sie, was Sprüchw. 1, 16. steht: „Ihre Füße eilen Blut zu vergießen“ etc. Sie lehren nicht, verstehen auch nicht, was das Reich Christi sei oder was darin geschehe. Sonst strafe sie und wehre ihnen.

240. So sagt nun Jakob klärlich: „Die Völker werden ihm anhangen“, und gehorchen. Es wird mit der Verheißung zugehen, und der Schiloh will dabei sein, und durch unsere Zunge und Mund wirken und kräftig sein. Wir aber lehren und reichen die heiligen Sacramente ohne alle äußerliche Gewalt und Schwert, und werden doch durch das Wort, die Taufe und Communion oder Abendmahl Kinder des ewigen Reichs geboren. Und eben dasselbe sagt auch Jesaia klärlich Cap. 11, 6. 9.: „Die Wölfe werden bei den Lämmern wohnen, und der Pardel bei den Böcken liegen. Ein kleiner Knabe wird Kälber und junge Löwen und Mastvieh mit einander treiben“ etc. „Man wird nirgend leben noch verderben auf meinem heiligen Berge; denn das Land ist voll Erkenntniß des HErrn, wie mit Wasser des Meers bedeckt.“ Die Christen werden niemand verletzen oder Krieg führen auf diesem ganzen Berge, nämlich also, daß sie auf die Weise sollten Christen werden, wie sich die Türken und Mahometisten träumen lassen; sondern die Erkenntniß des HErrn Christi wird reichlich im Schwange gehen, wie das Meer übergeht und ausbricht über die ganze weite Welt. Diese Erkenntniß Christi wird zuwege bringen, daß man die Wehre und Waffen ablegen, und die Schwerter zu Pflugschaaren machen wird etc., Jes. 2, 4., nämlich, so viel das Reich Christi angeht.

241. Von dieser Beschreibung des Reichs Christ ist noch der Pabst weit, welcher ein Tyrann und Feind der Kirche und Gemeinde Gottes ist; und dieweil er mit dem Bann nichts ausrichten kann, greift er nur zum Feuer und Schwert. Wie seine Sophisten zu Löwen, die groben ungelehrten Esel, nur damit umgehen, daß man uns unbilliger Weise mit grausamer Tyrannei und Schwert zusetzen soll. Aber wir sind von Gottes Gnaden nun verständig, gelehrt und mächtig im Worte, zu lehren, zu vermahnen, zu überwinden und zu strafen die Widerspenstigen. Aber nach dem allen fragt der Papst gar nicht, ja, die ihm widersprechen, tödtet er mit dem Schwerdt und verbrennt sie mit Feuer. Das haben die Apostel nicht gethan, deren Nachfolger und Stuhlerben er sich doch rühmt; sondern sie lehren, daß man den Widersacher strafen soll durch das Wort, welches der heilsamen gesunden Lehre gemäß ist.

\quellseite{514}
242. Eben dasselbe haben wir auch gehört, daß in den andern Stücken dieses Segens Jakobs eigentlich also gelehrt werde, darin er mit dieser ungemeinen Rede beschreibt den Reichthum und unermeßliche Verheißung des Heiligen Geistes, so uns durch Christum gebracht worden ist. Denn diese Gabe des Heiligen Geistes bringt nicht mit sich äußerliche Waffen oder Feuer, deren jetzt der Pabst gebraucht, sondern sie bringt mit sich ein solches Gut, das da trunken macht. Und der Heilige Geist mit seinen Gaben ist dieser Wein und Weinberg, daran Christus sein Füllen bindet, das ist, die Kirche und ihre Nachkommen, das Füllen und aber das Füllen. Denn dieses Reich wird nicht aufhören und wird für und für mit Reichthum und mit den Gaben des Heiligen Geistes gespeiset werden.

243. Dies, habe ich gesagt, sei der rechte und einige Verstand und Auslegung dieser Worte. Denn Jakob redet nicht von leiblichem Weingarten und Esel oder Füllen, sondern von den Eseln und Füllen der andern Völker bis an das Ende der Welt. Der Heilige Geist ist nimmer von der Kirche, sondern findet allezeit Etliche, die er trunken macht.

244. Zur Zeit des Johannes Huß schien die Kirche gar vertilget, und Gott hat doch beide ihn und seinen Gesellen Hieronymus von Prag erfüllt mit dem Heiligen Geist und geistlicher Trunkenheit, daß ein oder zwei Männer alle Nationen der Deutschen, Böhmen und Spanier verachtet haben. Da hat Ein Mensch dem Concil und der ganzen Welt widerstanden. Das heißt recht trunken werden vom Heiligen Geist und vor den Weisen und Heiligen dieser Welt zum Narren und Esel, ja, Erzketzer werden. Aber der Heilige Geist fragt nach ihrem Urtheil nicht, und Gott nimmt seinen Heiligen oder sein Füllen weg und rafft ihn auf, und läßt jene dieweil Episcopische Säue sein, und Klöße und Stöcke bleiben, die gar von keinen heiligen oder geistlichen Dingen wissen.

245. Das sind aber rechte Gaben des Heiligen Geistes, so die Leute trunken machen mit der sehr reichen Erkenntniß des Sohnes Gottes. Wie Johannes Huß, da er zur Marter geführt worden, mit großem Muth und unverzagtem Herzen also gebetet hat: „JEsu, du Sohn Gottes, der du für uns gelitten hast, erbarme dich meiner!“ Das wird nicht gesagt ohne den Heiligen Geist. Es hätte niemand anders also reden können, denn der von diesem Wein trunken gewesen wäre. Der ist nun Einer gewesen von den Füllen, so an den edlen Weinstock gebunden sind.

246. Wir werden aber nicht allein trunken vom Geist, und wissen, was da sei der Geist der Weisheit, des Verstandes, des Raths und der Stärke, Jes. 11, 2., sondern wir werden auch täglich gewachsen. Denn die Sünde setzt noch in unserm Fleisch, und der Geist und das Fleisch sind wider einander. Aber die Verheißung des Heiligen Geistes wirkt in uns, und wächst immer die ganze Zeit unsers Lebens sein Kleid, nämlich, uns. Denn er fegt immer die Sünde aus, die noch übrig in uns ist, und wir werden also gewaschen und gereinigt, nicht mit Wasser, sondern mit Wein und Traubenblut, mit röthlichem Wein, das ist, mit den vollen Gaben des Heiligen Geistes.

247. Denn die Sünde wird nicht getilgt, wie sich die Sophisten träumen lassen, mit unsern Reuen, Bereitungen und Genugthuungen; welches die groben Esel zu Löwen neulich in ihren Artikeln fälschlich gesetzt, nämlich, daß nach der Taufe und Absolution gar keine Sünde an dem Menschen übrig bleibe. Denn wenn dieselbe närrische und lose Theologie bestehen sollte, daß nach der Taufe keine Sünde mehr übrig in uns wäre, was bedürfte man denn des Waschens? Sagt nicht Christus Joh. 15, 1. 2.: „Mein Vater ist ein Weingärtner. Einen jeglichen Reben an mir, der Frucht bringet, wird er reinigen.“ Der Rebe ist rein, und wird doch gereiniget, wie geschrieben steht Offenb. 22, 11.: „Wer fromm ist, der werde immer frömmer.“ Wir sind nicht im vollkommenen Wesen, sondern müssen immer noch fromm werden und also fortgehen, oder uns von einer Tugend zur andern wenden, von Tag zu Tag, von einer Reinigung zur andern, und die wir geheiliget sind, müssen noch immerhin geheiliget werden.

248. Das verstehen die Sophisten nicht, die das immer wollen zum Voraus gesetzt haben: Es soll keine Sünde gereiniget werden, sondern man soll allein viel Verdienste sammeln, welche wir auch Andern verkaufen können. Wie ich, da ich noch ein Mönch war, täglich pflegte zu beichten, zu lesen, zu fasten, zu beten, Opfermesse zu halten, zu dem Ende und um der Ursache willen, daß ich von den Vigilien, Messen und übrigen Werken den Laien etwas mittheilen und verkaufen könnte. Dafür nahmen die Mönche Geld, Getreide und Wein; niemals noch jetzt viel Briefe und Siegel hat von solchem Kaufen und Verkaufen der Mönche und Pfaffen, darin sie sich gegen die Leute also verschrieben haben: Wir wollen dir N. hierin mitgetheilt haben für einen Scheffel Getreide unser Fasten, Wachen, Beten, Disciplin, Messe 2c., alles vermöge dieses ordentlichen Contracts und Kaufbriefs 2c. Also pflegten wir in unsern Briefen zu schreiben und verkaufen also unsere guten Werke. Denn wir waren, unsers Erachtens, ganz rein nach der Taufe und hatten gar keine Sünde mehr.

249. Und eben dasselbe lehren jetzt noch ganz und gar also die Sophisten zu Löwen und Paris. Aber wenn die Anfechtung daher ging, da waren wir alle erschrocken, wurden gar bleich davor, und war nichts, das verwirrter wäre, und sich mehr entsetzt oder gefürchtet hätte denn unser Gewissen. Niemand fürchtete sich mehr vor dem Tod, vor dem Gerichte Gottes und vor der Hölle denn eben die Mönche. Ich bin mir es wohl bewußt und weiß mich schuldig, wie greulich ich mich gefürchtet habe vor Gottes Gericht und vor dem Angesichte Christi, als eines gestrengen Richters und Tyrannen.

250. Und die Papisten behalten heutiges Tages gleichwohl noch diese höchst gottlose Meinung, und wenn der Tod daher gehet, zittern und erschrecken sie noch eben so sehr. Sind es Mönche, so halten sie den Ihrigen vor und prägen ihnen ein ihre Patrone, als, Augustinus, Franciscus, Dominicus 2c. Ist aber das nicht ein erbärmlich und erschrecklich Ding? Sind es aber Laien, so werden sie Mönche und ziehen noch in der Stunde des Todes Kappen an, und dürfen noch solch gottlos Wesen mit Gewalt und unbilliger Grausamkeit vertheidigen, welche Gewalt sie an denen üben, die ihnen nicht beifallen wollen.

251. Derhalben sollen wir Gott danken, die wir Christum haben, den Bischof unserer Seelen, zu dem wir mit gutem und fröhlichem Gewissen unsere Augen und Herzen erheben mögen und sagen: „Herr Jesu Christe, du Sohn Gottes, erbarme dich unser!“ welches wir vor dieser Zeit nimmermehr hätten thun können. Darum hütet euch vor solchen Theologen, und glaubt, daß wir immer von Tag zu Tag müssen gewaschen werden und daß der alte Sauerteig auch müsse ausgefegt werden. Denn wir sind nicht so rein, wie das Gesetz von uns fordert, sind auch so vollkommen und feurig nicht im Glauben und in der Liebe. Darum müssen wir immer beten und seufzen, daß wir von Tag zu Tag in der Reinigkeit wachsen mögen.

252. Und das meint hier Jakob, da er sagt: „Er wird sein Kleid waschen“, das ist, die Kirche wird er waschen in diesem Leben, welche nach dem Fleisch noch immer beflecket und befleckt ist. Aber in dem andern und zukünftigen Leben wird er die Kirche oder Gemeinde Gottes darstellen, daß sie keine Runzel oder Flecken habe, sondern daß sie rein sei. Inzwischen aber thut uns noth, daß wir immer beten: „Dein Name werde geheiliget, vergib uns unsere Schulden, und führe uns nicht in Versuchung“ 2c., Matth. 6, 9. 12. 13.

253. Wir wollen die Sophisten immer fahren lassen mit ihrem eitlen Ruhm, damit sie sich träumen lassen, daß sie sagen dürfen: Ich bin heilig, ich habe mich gereiniget durch meine Reue und Beichte; alles, was ich nun Gutes thun werde, ist überflüssig, dasselbe will ich den Bauern und Laien verkaufen, daß ich dafür meinen Leib nur wohl speisen und demselben gütlich thun möge 2c.

254. Aber ein recht christlich Leben ist ein solches Leben, da man täglich an den Weinstock gebunden wird, das ist, an das Wort, und da man trunken wird von den Gaben des Heiligen Geistes oder vom Worte. Zum Andern, daß man nicht allein also trunken werde vom Geist und mit Freudigkeit erfüllt, welches dem neuen Menschen die heilsamste Trunkenheit ist; sondern daß man auch nach dem alten Menschen im Wein gewaschen werde. Der alte Schalk muß noch gewaschen werden, nicht mit Seifenwasser, sondern durch das Wort, durch das Blut des Sohnes Gottes, welches unter uns ausgeprengt wird durch das Predigtamt. Denn wenn wir lehren und predigen, thun wir Anderes nichts, denn daß wir aussprechen und austheilen die Kraft des Blutes Jesu Christi unter das Volk; wie droben aus 1 Pet. 1, 2. gesagt worden ist. Man soll in der Kirche nicht Aristoteles lehren oder des Papstes Decretalen, sondern vom Blut Christi, des Sohnes Gottes, soll man uns predigen, daß uns dasselbe von Tag zu Tag reinigen möge, bis ich vollkommen rein bin, daß ich dem Heiland Christo, wenn er kommen wird, begegnen und unter die Augen gehen könne.

\quellseite{515}
255. Hierher gehört auch die Röthe der Augen, oder die blöden, vollen, vertrunkenen Augen Christi in seinen Völkern und Gläubigen. Welche Röthe nicht allein verstanden werden muß, wenn wir selbst trunken sind; sondern wenn wir den Heiligen Geist haben, und täglich gesegt und gereinigt werden vom alten Sauerteige. Danach sollen wir für Andere auch sorgen, und dieselben lehren und regieren. Denn das Auge bedeutet in der heiligen Schrift einen Regenten oder Lehrer, und durch die Zähne werden auch die Lehrer verstanden, welche beißen und die Laster und Sünden strafen sollen. Dies ist eine rechte und gute Allegorie, ja, es ist vielmehr der geistliche und rechte eigentliche Verstand des Textes.

256. Derhalben, wenn ich nun also trunken bin und den Heiligen Geist in seinen Gaben habe, als, im Glauben, in Erkenntniß Christi, der Taufe und des Worts, welches die allergrößten Gaben sind und die allerköstlichsten, die uns zum Leben führen, und auch das Bad für den alten Menschen habe: alsdann stößt mich Gott der Herr in seine Ernte und so setzt mich unser Herr Gott auf das Rad und auf die Scheibe.

257. Denn man muß Kirchendiener haben, die das Wort lehren, und muß das Predigtamt da sein, deß kann man nicht entbehren. Sie können nicht alle der heiligen Schrift warten. Die Nothdurft dieses Lebens erfordert es, daß allerlei Handwerksleute sein müssen, als, Schmiede, Zimmerleute, Töpfer etc., wie Sirach Cap. 38, 26. ff. zeugt; dieser aller kann man in einer Stadt nicht entbehren. Sie sollen nicht alle den Ackerbau, die Haushaltung, die Verwaltung der Polizei und andere Aemter oder Werke dieses gemeinen Lebens verlassen. Darum sind sonderliche gewisse Tage bestimmt und geordnet zur Predigt, an welchen Tagen die Laien oder das gemeine Volk zusammen kommen sollen, Gottes Wort zu hören; da müssen wahrlich die Augen roth und die Zähne weiß sein.

258. Aber diese Röthe der Augen ist schändlich, steht nicht wohl, ist nicht schön, wie am goldgelben Regenbogen, und in den braunen Augen, welche auf Griechisch χαροποί, das ist, Liebesaugen genannt werden; es sind auch nicht die Augen, so da glänzen wie Feuer von Liebe oder Zorn, oder sonst von andern heftigen Affecten oder Begierden, wie sie an den Hunden und Wölfen gesehen werden; sondern hier sind die Winkel an den Augen mit dieser Röthe geschändet oder befleckt, das ist, die Diener des Worts, welche in der Kirche die Augen sind, durch welcher Amt Gott das Anbinden an den Weinstock und die Reinigung der Reben ausrichtet, die sind nicht schön, sondern verächtlich und schändlich, ohne alle schöne Gestalt, und werden angesehen, wenn sie reden, wie Narren und trunkene Leute, als ob sie gestern zuviel gesoffen hätten; haben Augen, als wären sie im Salz gelegen wie Pökelfleisch.

259. Also sind die Diener Christi verschmäht und verachtet um dieser Trunkenheit willen des Heiligen Geistes. Denn sie reden nicht, was man gerne hört, das den Leuten gefalle, das schön, lieblich und den Menschen angenehm sei; sie haben nicht so hübsche schöne Liebesaugen wie die Mädchen, sondern haben traurige Augen, und die schändlich roth sind. Denn fleischliche Menschen haben dazu keine Lust, daß sie Gottes Wort hören; wie Paulus sagt 1 Cor. 2, 14.: „Der natürliche Mensch vernimmt nichts vom Geist Gottes; es ist ihm eine Thorheit.“ Man spricht bald: Es ist ein voller Narr, ein feindseliger Mensch, ein Ketzer etc.

260. Ein solch Amt ist das Predigtamt nach dem Urtheile der ganzen Welt, welches doch vor Gott ganz rein und schön ist; denn da sind es nicht rothe Augen, sondern Engelaugen. Dieweil es aber so eine feindselige und verächtliche Lehre ist, nicht lieblich und lustig anzuhören, ja, sie straft vielmehr und erzürnt die Zuhörer, und redet von solchen Dingen, die sich mit der Vernunft gar nicht reimen, und die ganz verächtlich und närrisch sind: darum schreien die Leute und sagen: Die Apostel und andere Diener des Evangeliums „sind voll süßes Weins“, Ap. Gesch. 2, 13. Und doch sind diese Augen Gott sehr angenehm, die mit dem Wort die armen betrübten Herzen regieren, lehren, aufrichten und trösten.

261. Danach sind die Zähne auch „weißer denn Milch“; denn sie sind voll Proselyten. Und wiewohl sie beißen, halten sie doch den Leuten das vor, das sehr gut und lieblich ist. Das Strafen, so durch das Predigtamt geschieht, tödtet nicht, sondern heilt vielmehr und hilft den armen erschrockenen Gewissen. Und wiewohl wir sehr hart strafen und die Zähne hineinschlagen mit Strafen, Schelten und Bannen, sind doch die Zähne weiß wie Milch, das ist, die Strafe der Diener Christi ist heilsam und lieblich.

262. Auf diese Weise, habe ich gesagt, verstehe ich diesen Text von den Gaben des Heiligen Geistes und vom Predigtamt, durch welche der Heilige Geist unter das Volk ausgetheilt, und also das Reich Christi bestellt und ausgerichtet wird. Die Kirche hört die Trunknen, und die rothen Augen und die weißen Zähne, und werden alle täglich gewaschen durch den Geist des Rechts und durch den Geist der Hize, wie Jes. 4, 4. steht, zum ewigen Leben. Denn also hat der Patriarch mit ungereimten Worten die Kirche abgemalt und beschrieben; und in denselben Worten soll man einen geistlichen Verstand suchen und keinen allegorischen. Denn es sind eigentliche und deutliche Worte, und ist darin keine Allegorie. Aber man muß dazu andere Augen, Ohren und Sinne haben, denn die fleischlichen Menschen haben.

263. Christus thut und richtet Alles aus allein mit dem Wort und will keines Schwerts brauchen. Davon weiß die Welt gar nichts; denn sie versteht und glaubt nicht, daß das Wort so große Dinge ausrichten könne. Wie Christus Joh. 14, 17. vom Heiligen Geist sagt, „welchen die Welt nicht kann empfahen; denn sie siehet ihn nicht und kennet ihn nicht.“ Sie hört ja das Evangelium wohl, sieht, daß wir getauft werden, daß wir des Abendmahls des Herrn gebrauchen: sie sieht aber nicht, daß der Heilige Geist durch das Wort und unser Amt kräftig ist, das, sage ich, weiß sie nicht, und ist doch dabei, von man predigt und die Sacramente austheilt, ist aber dabei wie im Traum.

264. Aber die glauben und die der Heilige Geist rührt durch das Wort, die kennen und empfahen den Heiligen Geist, wiewohl dasselbe vor der Welt nicht also scheint. Darum ist es ein ander Wort und weit eine andere Regierung denn eines Menschen, nämlich, Gottes Wort, welches er durch uns redet und dadurch er in der Kirche kräftig ist. Er richtet Alles aus allein mit dem Wort: er erleuchtet die Leute damit; er richtet sie auf und macht sie dadurch selig; darum ist es ein Wort der Verheißung, der Gnade, des ewigen Lebens und Seligkeit.

265. Von dieser Kraft des Worts weiß die Welt gar nichts und ist doch also, nämlich: wie Christus im Evangelium durch das Wort Teufel austreibt, also sagt der Diener Christi auch zum Sünder: „Ich absolviere dich von allen deinen Sünden. Und es geschieht auch also, denn es ist nicht ein Wort eines Menschen, vor dessen Stimme der Teufel mit nichten fliehen wird. Und wenn schon alle Fürsten und Philosophen alle ihre Bücher würden zusammentragen, würde doch nichts daraus werden. Aber auf die Absolution des Dieners Christi folgt gewißlich die Erlösung von Teufel und Sünde. Wenn dir der Heilige Geist Gnade gibt, daß du glaubst, da treibt er mit Einem Wort den Satan und den Tod aus.

266. Und in dem Theil hat sich der Papst auch über Christum erhoben, da er seine Decretalen und Decrete dem Worte Gottes gleich gemacht. Wir aber wollen den Papst in der Kirche nicht reden hören, sondern Gott allein, wie der 60. Psalm V. 8. sagt: „Gott redet in seinem Heiligthum, deß bin ich froh“, das ist, unter dem Haufen oder Versammlung derer, die da berufen und getauft sind, da soll Gott allein hören und reden. Wo aber das Gehör Gottes nicht ist, da fliehe immer hinweg, so fern und sehr als du kannst. Denn es ist von Gott also beschlossen, daß Christus allein in seiner Kirche reden und neben sich keinen andern Lehrer dulden oder leiden will; wie die Stimme des Vaters vom Himmel herab lautet, Matth. 17, 5.: „Den sollt ihr hören.“ Desgleichen sagt er auch Joh. 10, 27.: „Meine Schafe hören meine Stimme.“ Des Papstes Decretalen aber wollen mir auf das heimliche Gemach weisen.

\quellseite{516}
267. Denn dieweil mich der Seelen Heil und Seligkeit zwinget und dringet, so will ich an den Weinstock und an den edlen Reben gebunden sein wie ein Füllen; ich bin ein armer Esel, esse aber gern gute Trauben, trinke auch gern guten Wein. Ich will mich waschen lassen in dem Blut dieses Schiloh und mit Traubenblut oder Gaben des Heiligen Geistes; ich will auch meine chakilim hören, die da rothe Augen haben, die verachtet sind, und weiße Zähne in der Welt, denn sie strafen mich und lehren Christum: sie beißen mich mit ihren Zähnen, aber sie suchen damit nicht mein Verderben, sondern meine Seligkeit.

268. Die groben Esel aber von Paris und Löwen wollen wir unter die Füße treten, die Gott dahin gegeben hat in verkehrtem Sinn, daß sie keinen Unterschied wissen zwischen Gottes- und eines Menschen Wort. Ist aber das nicht eine erschreckliche Straße, daß sie so gar grob und unverständig sind, daß sie solches nicht wissen?

269. Derhalben lasset uns lernen, daß in der Kirche oder Gemeinde Gottes nur Gottes Wort gelehrt und gehört werden soll, und daß man sich dessen befleißige, daß wir ja desselben gewiß sein mögen, nach dem Befehl des ewigen Vaters, der da Matth. 17, 5. sagt: „Den sollt ihr hören“, und auch des Sohnes, der da sagt Joh. 10, 5.: „Meine Schafe kennen der fremden Stimme nicht.“ Also streitet die ganze Schrift wider die Lehre der Menschen, des Papstes, Cardinalen und Sophisten etc., und sammelt und behält uns in dem Schafstall, darin Christus allein der einige Hirte ist. Da gehen wir ein und aus; und wohin uns Christus führt, da finden wir allenthalben Weide, draußen und drinnen bindet er uns an den Weinstock, wäscht uns unser Kleid durch die rothen Augen und die weißen Zähne.

270. Bis daher haben wir nun den vornehmsten Theil dieser Verheißungen gehabt; die andern, so nun hernach folgen, sind schier alle zeitliche Verheißungen, damit geht Jakob kurz hindurch.

\section[Vierter Theil]{Vierter Theil. Von den Segen und Verheißungen über Sebulon, Jissachar, Dan, Gad, Asser, Naphthali}

\subsection{I}

\begin{quote}
\textbf{V. 13. Sebulon wird an der Anfurt des Meers wohnen, und an der Anfurt der Schiffe, und reichen an Sidon.}
\end{quote}

271. Der Segen Sebulons ist kurz, und scheint, daß der Patriarch Jacob aus der Bedeutung der Namen bei vielen Stämmen Prophezeiungen oder Weissagungen gemacht habe. Wie er denn von Juda gesagt hat: „Dich werden deine Brüder loben.“ Also sagt er hernach von Gad: „Gad, geduld, gerüstet, wird das Heer führen“; item, von Dan: „Dan, jadin, wird Richter sein“ etc. Also scheint es hier auch, daß er auf das Wort sabal gesehen habe, welches heißt, bei Einem wohnen. Daher auch droben Lea diesem Sohne den Namen gegeben, da sie sagt: „Nun wird sich mein Mann wieder zu mir thun“ etc., 1 Mos. 29, 34. Und daher hat auch die Fibel, des Königs Ahab Weib, ihren Namen.

272. Also sagt er aber nun: „Sebulon wird wohnen“, er wird ein rechter Sebulon werden, an der Anfurt der Schiffe“, also daß er mit der Seite reiche an Sidon. Jetzt wäre es aber nun vonnöthen, daß man eine gute Geographie hätte und eine Karte vom Gelobten Lande, die etwas besser wäre, denn die gemeine Tafel ist. Es haben sich zwar Viele über diese Beschreibung sehr hart bemüht und fleißig darin gearbeitet, und ich weiß nicht, ob sie es auch alles verstanden, oder erreicht und erlangt haben, darnach sie gearbeitet. Denn die Schrift ist an vielen Orten mit ihnen nicht eins. In den Tafeln, die man gemeiniglich gebraucht, ist Asser an dem Ort gemalt, an welchen hier Sebulon gesetzt wird. Denn demselben Asser werden die Häfen und Anfurte des Meers zugeschrieben, welches doch die Lage von Sebulon ist. Wie sich aber dasselbe reimen wolle, da mögen sie nachsehen; es wäre denn, daß du die Anfurt verstehen wolltest jenseit des Meers, gegen Mitternacht und nach Antiochien.

273. Sidon ist die äußerste Grenze des Gelobten Landes, desgleichen auch Joppe, Cäsarea, Ptolemais, Tyrus etc. Lies das 19. Capitel des Buchs Josua und beinahe das ganze Buch, so wirst du sehen, wie diese Stämme sind ausgetheilt worden, ob es sich mit diesem Texte reimen wolle und ob die Tafeln auch hiermit übereinstimmen. Denn ich lasse es unentschieden und befehle es den Grammatikern und Geographen.

274. Jakob segnet Sebulon und eignet ihm eine gewisse Wohnung zu vor den Andern, nämlich, daß Asser des Mittelmeers, das thut er sonst keinem Stamme mehr. Diesen Stamm aber segnet er und lobt ihn vom Orte der Wohnung, nämlich, daß er wohnen solle an einem hafenreichen Ufer, wo die Schiffe vor Anker liegen; wovon in der heiligen Schrift nichts steht. Denn es wird an seinem Orte gelesen, was Sebulon gethan habe, wie er dieser Anfurte gebraucht. Denn die Gewalt der zwei Städte Tyrus und Sidon ist geblieben, welche gleichsam die Königinnen im Meere gewesen sind; wie Tyrus in den Propheten die Stärke des Meers genannt wird, darum daß sie die allergemaltigste Stadt gewesen ist unter denen, so am Meere gelegen sind, und eine herrliche und hochberühmte Handelsstadt, deren alle Völker in ihren Schriften und Büchern gedenken und sie rühmen. Also daß der Stamm Sebulon diese Städte nie inne gehabt, wiewohl das Jakob hier sagt: Er wird wohnen im Lande, darin Tyrus und Sidon gewesen sind.

275. Es erhellt aber aus diesem Texte, daß zur Zeit des Patriarchen Jacob die Stadt Tyrus noch nicht gebaut gewesen, sondern Sidon, welche älter ist denn Tyrus. Darnach aber ist Tyrus aufgenommen, und viel gewaltiger und berühmter geworden als Sidon. Wie zu unserer Zeit Venedig im Adriatischen Meere ist, also ist zu der Zeit Tyrus, die Insel, auf einem sehr festen Felsen gelegen gewesen, so daß sie nicht ohne große Beschwerung, Mühe und Arbeit von Alexander dem Großen erobert worden ist. Aber zu Jakobs Zeiten ist sie entweder sehr unbekannt oder aber noch nicht gebaut gewesen, und in nachfolgenden Zeiten ist sie eine solche große Handelsstadt geworden, die über die ganze Welt berühmt gewesen ist.

276. Derhalben ist die gewisse Wohnung Sebulons bei Sidon und Tyrus, beinahe bis an den Berg Carmel und die Stadt Ptolemais. Dies ist die Summa und Inhalt dieses Textes, daß der Vater Jakob voll Heiligen Geistes und Glaubens verheißt und diese bestimmte gewisse Verheißung hinterläßt, wiewohl sie leiblich ist, nämlich, daß Sebulon diesen Ort einnehmen sollte.

277. Und wiewohl die Verheißungen leiblich sind, soll doch der Glaube dazu kommen, daß nothwendig Sebulon in dem Lande wohnen werde. Denn der Teufel ist nicht allein unsern Seelen feind und des zukünftigen Lebens willen, sondern ist auch unsrer Habe und Nahrung, unsrer Wohnung, unserm Leibe und andern Gütern dieses Lebens feind. Wie heftig er aber nun tobt und wüthet, und wie hart er die Kirche verfolgt, setzt doch hier der heilige Geist dem Teufel ein Ziel; als wollte er sagen: wiewohl du heftig tobst und wüthest durch das Meer, durch die Philister, durch die Syrer, so wirst du doch diesen Ort meinem Sohne Sebulon lassen müssen, den ich ihm zuvor ausgesehen und bestimmt habe.

278. Diesen Glauben soll man fleißig ansehen und betrachten in diesen leiblichen Verheißungen. Also ist Juda ein Löwe auch leiblich oder weltlich und hat das Regiment, bis der Schiloh kommt, welcher dann bleiben wird bis an das Ende der Welt, mitten unter seinen Feinden. Das Reich Juda muß dennoch bleiben bis auf den Schiloh, wenn der Teufel schon sehr zornig ist. Es wird ohne Zweifel der Teufel fleißig damit sein umzugehen, daß er beide das politische oder äußerliche und auch das geistliche Reich dieses Volks ausrotten möchte; und er hätte auch leichtlich alle und jede Stämme vertreiben können, wo nicht der Heilige Geist zuvor verkündigt hätte, daß diese Grenze sollte Sebulons sein und die Andern sollten ihre Grenze an andern Orten haben. Und Josua hat hernach das Loos gelegt, wie es hier vorgeschrieben und bestimmt ist, und hat einen jeglichen Stamm an seinen Ort gesetzt; welche Stämme auch also in ihren Grenzen umschrieben, bewahrt und erhalten worden sind.

279. Darnach ist das auch zu merken, alle leiblichen Verheißungen und der Glaube daran nothwendig auch in sich begreift den Glauben an die geistliche Verheißung. Denn wer Gott anrufen will, daß ist vonnöthen, daß derselbe in seinem Herzen schließe und es dafür halte, daß Gott mit ihm versöhnt, ihm günstig und gnädig sei und ihn erhören wolle. Wiewohl die Sophisten zu Romen diese Lehre hart aufsetzten und derselben feind sind, ist es doch also gewiß und die lautere Wahrheit. Du wirst Gott den Vater, allmächtigen Schöpfer, nimmermehr anrufen, der dich erschaffen hat, dich mit Brode speist und von allem Uebel erlöst, es sei denn, daß du glaubest und gewiß seiest, daß du an ihm einen gnädigen Gott habest.

\quellseite{517}
280. Also schreien wir heutiges Tages wider den Papst und den letzten grimmigen Zorn Satans, und beten für das weltliche Regiment, um Frieden für Fürsten und alle Nothdurft dieser Zeit, nicht allein für das Wort, sondern auch für den Bauch und das tägliche Brod. Aber wir könnten nicht ernstlich und recht beten, wenn wir nicht glaubten, daß Gott unser Vater wäre, und daß er uns erhöre und gnädig sein wollte. Also wird es alles auf den Einen Christum gezogen, wie Paulus Col. 1, 16. 17. sagt. Alle Creaturen, und alles, was im Himmel und auf Erden ist, das alles miteinander wird in diese Summa und auf den einigen Christum gebracht, also daß er der Mittler sein soll, wir bitten gleich um geistliche oder leibliche Güter.

281. Auf diese Weise redet der Heilige Geist durch Jakob von der leiblichen Verheißung, welche gar nichts ist, wo nicht der Glaube an die geistlichen Verheißungen dazu kommt. Denn daher kommt es, daß in den Historien dieses Volks oftmals gesagt wird, wie die Kinder Israel zu Gott dem Herrn gerufen und geschrieen haben, und sind doch nicht erhört oder erlöset worden. Denn sie wollten Gott versöhnen mit Opfern und dadurch, daß sie immer schlachteten und die Opfer häuften; ja, das noch mehr ist, sie opferten ihm auch ihre Kinder, und ihre eigenen Leiber, die sie mit Messern und Priestien zerhacken, bis sie gar blutig wurden, oder aber führten sonst ein hartes strenges Leben. Aber sie waren abgöttische Leute und ehrten Gott mit selbsterdichtetem Gottesdienste; darum sind sie auch nicht erhört worden.

282. Also wollte Ahas auch ein sehr heiliger und gottfeliger König sein, der alle seine Söhne Gott aufopferte, rief die Götter Syriens an und viel andere mehr; aber Gott hörte ihn nicht. Also thut der Papst auch, der häuft immer das Gebet und andere vermeinte Gottesdienste mit höchstem Fleiß: sie gehen barfuß, gebieten den Leuten zu fasten, und daß sie ihre Sünde beichten sollen und zum Sacrament gehen; aber mit dem allen wird Gott mehr erzürnt, denn versöhnt, oder erreicht, Hülfe zu erzeugen.

283. Wir aber haben die Verheißung Psalm 50, 15.: „Rufe mich an in der Noth, so will ich dich erretten“; item im 145. Psalm V. 18. 19.: „Er thut, was die Gottesfürchtigen begehren. Der Herr ist nahe allen, die ihn anrufen, allen, die ihn mit Ernst anrufen“; item im 32. Psalm V. 2.: „Wohl dem Menschen, dem der Herr die Missethat nicht zurechnet, in deß Geist kein Falsch ist.“ An diese Verheißung halten wir uns, und verlassen uns darauf, daß wir mit Gott versöhnet sind und daß er unser Seufzen und Beten erhöre.

284. Der Papst mit dem ganzen Haufen der Mönche können Gott nicht recht anrufen; denn sie treten vor Gottes Angesicht, und rühmen ihre Verdienste und selbsterdichteten Gottesdienste, wie die Baaliten thaten. Und dieses Aberglaubens und Vertrauens halben auf die Werke haben die Propheten immer mit dem Volke Israel streiten müssen, welche alle ihre vermeinten Gottesdienste und ihr ganzes Reich verdammten, als das von Gott vermaledeyet und verworfen ist.

285. Verhalben, sage ich, begreifen alle leiblichen Verheißungen in sich den Glauben an die geistlichen Verheißungen. Gott wird auch nicht einen Bissen Brods geben denen, die sich auch ganz heisch schreien, wenn sie keinen Glauben haben; es sei denn traun, daß er uns ohngefähr als den unwürdigen leibliche Güter zuwerfe, wie er seinen Feinden und den Gottlosen auch etliche Güter gibt. Aber daß er uns günstig sei, und uns aus väterlichem Herzen erhöre und schütze, dazu ist der Glaube nöthig auch in äußerlichen weltlichen Werken, und wenn man um etwas bittet, das auch sehr fleischlich ist. Ich muß ja allewege Gott den Vater ergreifen, der einen Sohn hat, in welchem er sich geoffenbart und verklärt hat.

286. Es ist nie etwas Heiligeres oder Gottfeligeres gewesen denn die äußerliche Zucht der falschen Propheten: die ließen ihre Kinder durch das Feuer gehen, und verbrannten ihre herzliebsten Kinder und Früchte ihres Leibes ohne alle Barmherzigkeit oder Mitleiden. Damit wollten sie Gott ehren und dienen, wie im Propheten Micha am 6. Cap. V. 7. gesagt wird: „Soll ich meinen ersten Sohn für meine Uebertretung geben? oder meines Leibes Frucht für die Sünde meiner Seele?“ Aber der Prophet antwortet daselbst und sagt also V. 8.: „Es ist dir gesagt, Mensch, was gut ist, und was der Herr von dir fordert, nämlich Gottes Wort halten, und Liebe üben, und demüthig sein vor deinem Gott.“ Da schrieen aber die Götzendiener darüber und sagten, solches wäre Ketzerei, ungöttlich und unrecht. Denn es ist allezeit bei den Leuten größerer Fleiß und Hitze, über dem Aberglauben und Abgötterei zu halten, denn die rechte Lehre und wahren Glauben anzunehmen und daran festzuhalten; wie das heutiges Tages im Papstthume wohl zu sehen ist und die Welt ist zu allen Zeiten der Exempel voll.

287. Ich würde jetzt die Arbeit, das Wachen und Kasteien des Leibes nicht dulden oder ertragen können, welches ich vorzeiten getragen und geduldet habe, da ich noch ein Mönch war. Denn da ich zu Erfurt in der Hochschule angefangen hatte in guten Künsten und in der Philosophie zu studiren, und darin so viel gefaßt und gelernt hatte, daß ich Magister geworden war, hätte ich daselbst, nach dem Exempel der Andern, die Jugend wiederum lehren und unterrichten können, oder aber hätte mögen fortfahren und weiter studiren. Aber ich verließ meine Eltern und Herrn Verwandten, und begab mich wider ihrer aller Willen in das Kloster und zog eine Kappe an. Denn ich hatte mich überreden lassen, daß ich glaubte, ich würde in demselben Stande und mit solcher harten sauren Arbeit Gott einen großen Dienst thun.

288. Aber die Schrift zeuget und sagt Röm. 14, 23.: „Alles, was nicht aus dem Glauben gehet, das ist Sünde“, und wiederum, was aus dem Glauben gehet, das ist Gerechtigkeit. Derhalben, wo man Haus und Hof hat und das weltliche Regiment verwaltet, und in diesem ganzen zeitlichen oder fleischlichen Leben sollen wir immer im Glauben an den Mittler Christum beten, und sollen es dafür gewiß halten, solch Gebet sei Gott angenehm und daß es auch erhört werde.

289. Also wird Sebulon von seinem Vater an diesen Ort des Landes gesetzt, welcher ist bei Sidon, dem Satan zum Spott und Hohn, der heftig dabinder getobt hat, ja, auch allen höllischen Pforten zu Spott. Und er ist daselbst geblieben, so lange er den Glauben behalten: aber da er denselben hat fahren lassen und da die Sünden überhand genommen, ist er nach Assyrien geführt worden; und war vergeblich, daß die Israeliten sich dagegen sträubten, schrieen und rühmten, daß sie Gottes Volk wären.

290. Denn dies hat Gott der Herr durch die Propheten fleißig getrieben und den Leuten immer lassen vorhalten: Das ist mein Volk, das da glaubt, und mich als einen solchen Gott anruft, wie ich mich geoffenbart habe; denn also will ich erkannt und gebetet sein und nicht anders. Wo aber dieser Glaube erlischt und dies Anrufen oder Gebet auch darnieder liegt, so ist es nicht mehr Gottes Volk; sondern ich will das mein Volk heißen, das da glaubt, daß ich der Vater Christi, des verheißenen Schilo, sei, dieselben will ich erhören und erhalten im Lande, das ich ihnen zu besitzen gegeben habe; die Andern will ich herausstoßen und verwerfen.

291. Womit aber dieser Stamm umgegangen sei und was er für Hantierung getrieben habe, drückt die Schrift nicht aus. Es kann wohl sein, daß sie von den Früchten des Landes Getreide, Oel, Wein und Vieh in die nächsten Städte Tyrus und Sidon geführt haben und wiederum derselben Städte Waaren in ihr Land mitgebracht haben. Und es scheint, als ob sie Gemeinschaft gehabt haben mit den Heiden, so allernächst umher gewohnt, mehr aus Geiz und um gottlosen Wesens willen, denn sonst irgend aus Noth und daß sie davon Nutzen gehabt. Denn sie sind bald dahingefallen zu der Abgötterei der Sidonier, Askaroth etc. Und da sie also besieget worden sind vom Exempel und Gemeinschaft der Heiden, haben sie darüber müssen umkommen und zu Trümmern gehen.

\quellseite{518}
\subsection{II}

\begin{quote}
\textbf{V. 14. 15. Issachar wird ein beinerner Esel sein, und sich lagern zwischen die Grenzen. Und er sahe die Ruhe, daß sie gut ist, und das Land, daß es lustig ist; er hat aber seine Schulter geneiget zu tragen, und ist ein zinsbarer Knecht worden.}
\end{quote}

292. Dieser Stamm wird in den gewöhnlichen Tafeln richtiger angegeben, denn Sebulon. Denn er hat mit den andern Zwei, Sebulon und Naphtali, gewohnt in den Grenzen des Berges Thabor, wie Josua Cap. 19, 10. ff. zeuget und Matth. 4, 15. auch steht; dazu wird dasselbe Jes. 9, 1, 2. auch angezeigt. Hieronymus sagt, daß der Berg Thabor gleich am Ende oder an der Grenze liege der vier Stämme Ephraim, Issachar, Sebulon und Naphtali; und Asser hat gegen Abend, Ephraim gegen Mittag, Naphtali gegen Mitternacht gelegen, Issachar aber ist beinahe in der Mitte des galiläischen Landes gelegen gewesen.

293. Warum braucht er aber diese Metapher oder dieses Gleichnis, daß er ihn einen „beinernen Esel“ nennt, das ist, der nicht könne müde werden, der da Stärke habe, der nicht allein von Fleisch und Adern stark sei, sondern ganz beinnern? Und was mag er damit meinen, das darnach folgt, robez, das ist, der sich lagert, wie die vierfüßigen Thiere, wenn sie die Beine an den Leib ziehen; desgleichen in dem Worte hamisch ketajim, das ist, Grenze? Den Grammatikern ist zwar damit auch ein groß Kreuz und Marter auferlegt; denn sie wollen haben, daß das Wort saphat mit dem Punkte zur linken Lippe heiße, aber wenn dasselbe Wort einen andern Punkt oben auf der rechten Seite hat, so heiße es, Sattel oder Nest, darauf die Esel ihre Last tragen, und zanken miteinander und hat noch niemand entscheiden können; so kann ich denselbigen Zank auch nicht entscheiden.

294. Es ist aber eben dasselbe Wort, das Ps. 68, 14. steht: „Wenn ihr zu Felde lieget“ etc. An demselben Ort murmeln sie sich sehr, und ziehet es der Eine hieher, der Andere dorthin. Aber ich lasse die Auslegung der neuern Grammatiker und der Rabbinen fahren, welche das Wort also auslegen, daß es Sattel, Nest oder Last der Esel oder Ochsen heiße, und behalte an dem Ort die gemeine lateinische und des heiligen Hieronymus Uebersetzung, die also lautet: Si dormieritis inter medios terminos, wie wirs auf Deutsch gegeben haben: „Wenn ihr zu Felde lieget zwischen den Grenzen.“ Denn die Hauptleute und Obersten im Krieg müssen die Heerlager schlagen zwischen die Grenzen und Enden der Länder.

295. Issachar ist ein guter Pater, ist ein solcher Stamm, der Lust hat zu guter Ruhe, der nichts so lieb und werth hat, als Friede und gut Gemach; wenn er schon darüber große Last und Beschwerung tragen sollte, leidet er es doch alles mit Geduld, daß er nur Frieden und Ruhe haben möge. Er ist ein fauler Pater, wie man sagt, er ließe Holz auf sich hauen, daß er nur Frieden hätte; aber er ist weise. Und er ist zwar an dem allerbesten Ort des galiläischen Landes gelegen gewesen, um den Berg Thabor an dem Jordan, zwischen Asser, Sebulon, Ephraim und Naphtali, hat umher an allen Orten seine Brüder und Verwandten zu Nachbarn gehabt, und ist mitten unter ihnen wie eine Grenze eingeschlossen und geschützt worden.

296. Also lege ich es aus, daß er nicht gewohnt habe an den äußersten Grenzen, sondern daß er allenthalben umgeben gewesen sei mit Hülfe seiner Brüder und habe also guten Frieden und Ruhe gehabt, und wo er darüber hat etwas leiden oder dulden sollen oder Geld ausgeben, oder aber Andern dienen, \unsicher{hat er sich derselben gewehret} nur um Friedens und guten Gemachs willen. Es sind wohl gute faule Patres gewesen, aber doch weise dabei.

297. Auf diese Weise hat der heilige Geist diesem Stamm seine Wohnung bestimmt, daß er innerhalb der Grenzen der andern Stämme soll sitzen bleiben, und eignet ihm ein reich und fruchtbar Land zu, darin er sich lagert wie ein beinnerner Esel, der mit Arbeit oder Geld erkauft, daß er möge Frieden und Ruhe haben. Da liegt er wie ein fauler Pater. Wenn ihm geboten wird, Geld oder Schoß zu geben, oder sonst Dienstbarkeit zu leisten, was es auch sein mag, weigert er sich nicht; denn er will seinen Frieden und gut Gemach ungehindert haben, dabei will er gern bleiben.

298. Er ist wohl ein Esel, aber ist nicht unverständig: er leidet viel, gibt viel, läßt es sich viel kosten, duldet und verträgt viel, daß er nur das gute Land besitzen mag, und ist gar keiner, kann Alles leiden; er liegt zwischen den Grenzen, und hat die Landschaft der andern Stämme, so an die Grenze der Heiden reichen, um sich herliegen wie eine Mauer und starke Festung.

299. Derhalben ist dies auch eine leibliche Verheißung, und ist ebenso zu verstehen, wie wir die droben von Sebulon ausgelegt haben. Und GOtt hat ihnen diese Ruhe also erhalten nach seiner Verheißung, so lange sie im Glauben und Furcht des HErrn gelebt haben.

300. Er setzt aber noch hinzu: „Und er sahe die Ruhe, daß sie gut ist“, das ist, er hat Lust dazu. Denn es ist ein fein Ding um die Ruhe; wie ein Bauer auf dem nächsten Dorf hier bei uns zu Dabrun soll gesagt haben: „Wer zwei Kühe habe, soll gern die eine darum geben, daß man nur möge Frieden behalten. Es ist besser, man habe Eine in gutem Frieden, denn zwei im Krieg. Es ist besser, eine Hand voll mit Ruhe, denn beide Fäuste voll mit Mühe und Jammer“, wie Salomo sagt im Prediger Cap. 4, 6.

301. Dieser Esel, sagt er, wird so viel leichter Alles dulden und ertragen; denn er wird nicht Lust haben zum Krieg, sondern wird sehen, daß die Ruhe ein fein lustig angenehm Ding ist in dem guten fruchtbaren Lande, und wird denken, es sei besser das Geld verlieren denn das Land. Darum hält und neigt er die Schultern her zu tragen, was ihm wird auferlegt werden, trägt die Last gern, so ihm auferlegt ist, und ist, sagt der Text, „ein zinsbarer Knecht worden“, nicht allein seiner Richter und Könige, sondern vielleicht auch wohl der Fremden, daß er nur mit Gold und Silber, oder womit er konnte, den Frieden kaufen möchte; und es ist wohl getauft. Denn man muß viel größere Unkosten und viel mehr darauf wenden, wenn man mit Krieg und Waffen um die Grenze streiten soll.

302. Also ist nun der leibliche Segen dieses Stammes bestätigt, welchen sie behalten haben, so lange sie den geistlichen Segen auch gehabt. Da aber derselbe mit ihnen ein Ende genommen, da haben sie auch ihre Aecker und das Land verloren, und sind nach Assyrien geführt und gesetzt worden.

303. Zuletzt soll man auch merken, daß er dies so fein auf den Namen deutet; denn Issachar hat den Namen vom Lohn, als wollte man sagen, ein Taglöhner. Daher auch Lea ihrem Sohn diesen Namen gegeben hat. Und Judas Ischariot wird also genannt vom Lohn, darum daß er Christum um Lohn verkauft hat. Dahin sieht er auch mit diesem Stamm, nämlich, daß er zu seinem Lohn Lust habe, und sich genügen lasse an seinem Theil und den andern Theil nicht begehre. Aber der Juden unnützes Gemäsch will ich fahren lassen, welche das Neigen der Schultern auf das Gesetz ziehen, daß er sich desselben befleißigen und damit abmühen werde.

\subsection{III}

\begin{quote}
\textbf{V. 16—18. Dan wird Richter sein in seinem Volk, wie ein andrer Geschlecht in Israel. Dan wird eine Schlange werden auf dem Wege, und eine Otter auf dem Steige, und das Pferd in die Ferse beißen, daß sein Reiter zurückfalle. HErr, ich warte auf dein Heil.}
\end{quote}

304. Hier ist abermal ein Wortspiel: dan jad(in), das ist: Dan wird Richter sein. Es scheint aber, als sei dies ingemein geredet, und daß diesem Stamm nichts mehr zugeeignet werde, als den andern, welche abgesondert und unterschieden waren, also daß ein jeglicher Stamm seine Richter, Amtleute und Hausväter für sich gehabt. Derhalben, was hier von Dan gesagt wird, gehört zu seinem Hausregiment, und nicht zum politischen oder weltlichen Regiment, damit er über Andere regieren und herrschen sollte; als wollte er sagen: Das Amt und Stand dieses Stammes wird mit seinem Namen übereinkommen; wie er genannt wird, also wird er auch ein Urtheil haben. Die Verheißung wird fest und unbeweglich sein. Dan wird nicht umkommen, sondern wird bleiben, daß er über seinen Stamm richte, wie ein jeglicher Stamm unter den andern seinen Richter, Regiment und sein Volk hat nach dem Gesetz Moses; wird aber doch seinen Stamm also regieren, daß er auch wohl die andern regieren könne. Welcher Segen in Simson ist erfüllt worden, wie wir sagen werden.

305. An welchem Orte aber sein Erbe gewesen, muß man aus den geographischen Tafeln ersehen, nämlich, im Erbtheil Benjamins am Meer, wie Simeons Erbe im Stamm Juda gewesen ist. Denn diese vier Stämme sind im Winkel gelegen gegen Mittag. Aber gegen Mitternacht und nach dem Stamm Sebulon und Ephraim zu haben sie die Stadt Laïs eingenommen, welche sie von ihres Vaters Namen Dan genannt haben, wie im Buch der Richter am 18. Cap. 23. 27. ff. steht.

\quellseite{519}
306. Was aber Jakob sagt von der Schlange, das zieht auch Lyra auf Simson, der den reißigen Zeug und das Fußvolk der Philister gebissen hat. Denn er hat niemals ein solch Heer gehabt, wie die andern Hauptleute oder Richter des Volks Israel; sondern er selbst einzig und allein hat ohne fremde Hilfe und Waffen, da ihn der Geist getrieben, die wichtigsten Dinge ausgerichtet; hat die Stricke und Bande zerrissen, damit er gebunden war, als wären es Faden oder Garn, so man zu wirken pflegt; hat große Haufen Streiter und Knechte niedergelegt, hat auf einmal mit einem Eselskinnbacken tausend Mann geschlagen, Richt. 15, 15., und viele haben vor ihm gezittert und sind geflohen; und endlich, da er hat sterben sollen, hat er mehr Leute getödtet und umgebracht, denn Andere, da sie gelebt haben. Richt. 16, 30.

307. Derhalben ist Simson ein wunderbarer Kriegsmann und Heiland gewesen vor den andern allen, der Alles durch Eingeben des Heiligen Geistes gethan hat, welcher ihn dazu getrieben, auch im Tode. Darum wird er von Mose also wunderbärlich abgemalt, wie sein ganzes Leben voll großer Wunderwerke ist. „Er wird“, sagt er, „eine Schlange werden“, nicht wider sein Volk, sondern wider seine Feinde und Widersacher. Und sagt Ein Ding zweimal: „Er wird eine Otter sein auf dem Steige“, welche, wenn sie mit Füßen getreten wird, bald um sich sticht, und sich mit einem giftigen Biß an dem rächt, der sie getreten hat. Also hat Simson auch gethan, er mochte gleich in Israel, oder bei den Philistern zu Gath oder Eshtaol sein, hat er Gott vertraut und mit einerlei Sturm den Feind angegriffen.

308. Und wird in seiner Historie, weder der Heiden noch auch der heiligen Schrift, dergleichen Exempel mehr gefunden eines solchen starken unüberwindlichen Kriegsmannes, der sich hätte dürfen wider ein ganzes Volk legen und der so viel herrliche Siege erlangt hätte ohne einiges Heer oder Kriegsrüstung; so er doch am Leibe und Gestalt Andern nicht ungleich gewesen, wiewohl er größere Stärke gehabt und sehr ungestüm und heftig gewesen.

309. Darum sagt Mose: Er wird niederlegen oder verstoßen, das ist, demütigen die Philister, also, daß seine Reiter zurück fallen, wie eine Schlange oder Otter. Er wird keine Wehre tragen, kein Schwert, seine Bogen oder Roß haben, sondern wird eines Esels Kinnbacken ergreifen, und damit eine große Schlacht thun, dergleichen Andere mit Gewalt und dem größten Geschütz nicht würden thun können: er wird heißen wie eine Schlange. Ja, er bindet die Fuchsschwänze und Feuerbrände zusammen, und zündet damit die Früchte der Philister an und verbrennt sie; und da er endlich in ihre Hände und Gewalt gekommen, haben sie ihn doch nicht unterdrücken, oder von aller Furcht und Gefahr erlösen können. Ein solcher Krieg ist niemals mit irgend einem Volk geführt worden. Darum wird Simson recht und billig mit einer Schlange verglichen. Denn wiewohl dieselbe keine Wehre oder Waffen hat, wirft sie doch beide Roß und Reiter darnieder mit ihrem vergifteten Stich. Ich halte aber dafür, daß sie dies nicht verstanden haben, bis es vollzogen und die Prophezeiung Jakobs erfüllt worden ist.

310. Die andern Ausleger folgen einer andern Meinung, und legen diesen Text also aus: Er wird das Volk richten wie einen Stamm, das ist, alles Volk. Welches nicht wahr ist. Denn das liest man nirgend, daß Ein Richter das ganze Volk Israel regiert habe. Darum ist es synecdocheisch geredet: „Er wird Richter sein in seinem Volk, wie Israel“, das ist, ein Theil wird er richten, so daß Gideon, Jephthah, Debora im Nominativo, nicht im Accusativo, wie man in der Grammatik redet, verstanden werden. Denn wie die andern Stämme richten, also wird dieser auch richten; wie der Herr zu David sagt 2 Sam. 7, 7.: „Habe ich auch je geredet mit irgend der Stämme Israel einem, denen ich befohlen habe, mein Volk Israel zu weiden?“

311. Es ist aber eine gemeine Meinung bei allen christlichen oder kirchlichen Schreibern, daß der Antichrist aus dem Stamm Dan kommen würde, welchen sie unter der Schlange verstehen, und ist solche Meinung wie ein vornehmer Artikel des Glaubens von jedermann angenommen und bewilligt worden. Aber die Meinung will sich mit diesem Text durchaus nicht reimen, und ist ganz falsch und unrecht. Wo sie hergekommen ist, weiß ich nicht; man findet sie aber auch bei den ältesten Schreibern, und ist vom Anfang der Kirche allezeit ein gemein Gerücht gewesen, daß der Antichrist kommen würde, wie St. Johannes 1. Ep. 2, 18. sagt: „Es sind viel Widerchristen.“

312. Sie haben gesagt, er wird geboren werden zu Babylon aus dem Stamm Dan, und diese Lüge ist durch alle Kirchen oder Gemeinden gegangen, so auch, daß ich oftmals gewünscht habe, daß unsere Vorfahren die Lehre von Christo mit solchem Ernst und Fleiß dem Volk vorgehalten und eingeprägt hätten.

313. Ob man aber von dem Namen Dan eine Allegorie oder heimliche Deutung nehmen soll, weiß ich zwar nicht; denn dies Wort „Dan“ heißt so viel als Gericht, wie 1 Mos. 6, 3. steht: „Die Menschen wollen sich meinen Geist nicht mehr strafen lassen“ etc., also daß es dahin gezogen werde, daß der Antichrist das Amt eines Richters und nicht eines Heilandes führen wird. Das sage ich, weiß ich nicht. Es hat wahrlich der Papst die Kirche voll Gesetze und erdichteter Menschensatzungen gemacht, und hat das Wort und den Glauben an die göttlichen Verheißungen gar ausgelöscht und vertilgt. Dan, das ist, der Papst kann nichts, denn nur eitel Decrete und Decretalen, und solche Danitische Gesetze schreiben. Aber nach dem Buchstaben wird es niemand recht verstehen können, daß der Antichrist zu Babel soll geboren und nach jüdischer Weise beschnitten werden.

314. Und ich halte dafür, daß der Teufel diese Fabel erdacht, und daß er diese Glosse erdichtet habe, auf daß er damit unsere Gedanken möchte abführen von dem rechten gegenwärtigen Antichrist. Denn unter allen hohen Schulen und papistischen Lehrern ist keiner, der den Papst für den Antichrist halte; sie sind alle der Meinung, er werde aus Babylon kommen. Inzwischen aber, dieweil sie sich von ihm also träumen lassen und auf ihn warten, werden sie von dem rechten wahren Antichrist, nämlich, dem Papst zu Rom, unterdrückt und verschlungen, welcher ein rechter Daniter ist und aus Babylon hergekommen, nicht aus dem, das in Assyrien gelegen ist, sondern aus Rom, welches das rechte Babylon ist, darin Dan, das ist, der Papst regiert.

315. Derhalben hat uns der Teufel mit seiner Gaufelei wollen betrügen und beides unsern Verstand und Augen verblenden, daß wir ja nicht auf den Antichrist, der in der heiligen Stadt sitzt, sehen könnten. Denn durch des Teufels Rath und Lüge ist es dazu gekommen, daß der Antichrist in der Kirche sitzen sollte, und doch nicht regieren mit Gottes Gesetz, Verheißung und Gnade (denn das hat er nicht thun können), sondern dem alten stracks zuwider. Wo Christus durch die Gabe des Heiligen Geistes die Sünde vergibt, die Gewissen von Sünden, Tod, Hölle und Gewalt des Teufels erlöset; da hat der Papst zu Rom dagegen mit unzähligen vielen Stricken und Lügen die Welt gefüllt, ehe wir dasselbe gesehen, oder davon urtheilen können. Denn was thut er doch Anderes, da er im Tempel Gottes sitzt, denn daß er Stricke legt, damit er nur die armen Gewissen fängt, und tödtet auch dieselben mit seinen närrischen unzähligen Gesetzen und unnöthigen Menschensatzungen vom Essen, Trinken, Kleidern und dergleichen Narrwerken mehr, und hält und donnert greulich wider diejenigen, so ihm nicht gehorchen wollen, und gibt dieselben gar dem Teufel?

316. Das heißt ja recht, wider Christum und sein Evangelium lehren. Denn Christus erlöset uns von der Gewalt des Gesetzes Gottes, von dem Anklagen des Gesetzes oder der Zehen Gebote und von der Verpflichtung der äußerlichen Ceremonien des Gesetzes. Aber der Papst hat gesagt, daß auch unsere Haare, Essen, Trinken, Haus, Weib etc., und alles, was schier lebt, der Sünde unterworfen und befleckt sei, wo es nicht von ihm erlaubt und in seinen Decretalen geheiligt ist.

317. Auf diese Weise hat dieser Daniter mit seinen Gesetzen das Evangelium ausgelöscht, und ist geboren und hergekommen aus dem rechten Babylon, nämlich aus Rom; wie Augustin auch Rom das andere Babylon nennt, welches angefangen hat und aufgekommen ist, da jenes Babylon abgenommen hat. Und Dan oder der Papst ist die rechte Schlange auf dem Wege Christi, der mit seinem Gift und Gewalt diejenigen tödtet, die darauf wandeln.

\quellseite{520}
318. Dies sage ich allein darum, daß ihr wisset das Böse zu verwerfen und das Gute zu erwählen. Wenn die Sophisten zu Paris und Löwen diese Auslegung höreten, würden sie rufen: Ich lehrete Ketzerei, und das wider die angenommene Meinung der Väter wäre, die nun fast tausend Jahre gelehrt und in der Kirche ausgebreitet worden ist. Aber es sei gleich der Kirche oder der Synagoge Meinung, ist sie doch wahrlich nicht recht und schickt sich eigentlich hierher nicht. Denn der Text redet vom Segen des Stammes Dan und nicht vom Fluch. Droben hat Jakob offenbar Levi und Simeon verflucht; aber hier wird das Gegenteil verheißen, nämlich, daß Simson seinen Stamm mit den andern, so allernächst dabei gelegen, erlösen werde.

319. So aber jemand die Allegorie gefällt, der mag es geistlich verstehen vom Königlichen Babylon, und von dem schändlichen heuchlerischen Ungeheuer, dem Papste zu Rom, welcher der rechte teuflische Dan ist; denn er hat die Kirche ausgelöscht und verschlungen, vornehmlich derjenigen, so erwachsen und zu Jahren gekommen sind, die sich auf ihr Klosterleben und Werke verlassen haben; derselbigen Seele, Leib und Gut hat er verderbet, und was Gott hat wollen gut, heilsam und frei haben, das alles hat er schädlich und verdammlich gemacht. Der verfluchte Tropf zu Rom. Man soll ihm billig fluchen.

320. Zuletzt wird gefragt: Warum Jakob dies Stück habe hinzugesetzt: „Herr, ich warte auf dein Heil“? Darauf antworten Etliche also: Dieweil Simson so ein großer trefflicher Mann hat werden sollen, daß er auch für den Messias selbst hätte mögen angenommen werden, so habe der Patriarch Jakob dies wie zur Warnung dazu gethan, daß, wie große Dinge Simson würde ausrichten, man doch noch auf den rechten wahren Messias warten sollte. Wie viel aber dasselbe gelten möge, da mögen Andere zusehen. Meines Verstandes ist es nicht die rechte Meinung dieses Textes. Denn im Volke Israel sind vor Simson Etliche gewesen, die großer Dinge halben, so sie ausgerichtet haben, keineswegs geringer gewesen sind denn Simson. Gideon hat mit einem kleinen Häuflein ein groß gewaltig Heer der Midianiter niedergelegt. Jephtah hat die Ammoniter niedergeschlagen, und sind andere starke und löbliche Hauptleute mehr in Israel gewesen, von welchen man auch dasselbe sagen möchte.

321. Mich dünkt nicht, daß diese Worte müssen auf Christum gezogen werden, von dem Jakob droben in dem Stamme Juda gesagt hat, sondern er redet insgemein vom Heil, nach der Weise der heiligen Schrift, darin oftmals das Heil gesetzt wird für den Sieg; also, 1 Sam. 11, 13. spricht Saul: „Es soll auf diesen Tag niemand sterben; denn der Herr hat heute Heil gegeben in Israel“, das ist, er hat Sieg gegeben; und David sagt in seinen letzten Worten 2 Sam. 23, 5.: „Mein Heil sproßt noch nicht ganz“, als wollte er sagen: Ich habe auch Heil und viele Siege genug gebracht; aber die sind nicht zu vergleichen mit Christo.

322. Simson aber ist fast der letzte Heiland, im Buche der Richter. Denn Samuel kommt bald nach ihm, und das Volk bittet um einen König. Da folgen nach einander Saul, David, Salomo, und die zwölf Stämme werden in Ein Reich und Ein Volk versammelt. Aber zuvor war es ein zerstreut und zerrissen Reich. Denn Gott hat bald aus diesem, bald aus einem andern Stamme Richter erweckt; wiederum die Stämme zerstreut und sind oftmals unter einander sehr uneinig gewesen; wie das im Gesange Deboras im Buche der Richter am 5. Cap. zu sehen ist, welche beinahe alle Stämme verklagt, wie sie von ihnen sei verlassen worden, da sie die beiden stille gewesen, oder aber auf dem Meere Kananiter getrieben haben. „Warum“, sagt sie nicht. 5, 16. 17., „bleibest du zwischen den Hürden, zu hören das Blöken der Herden, und hältst groß von dir, und sonderst dich von uns? Gilead blieb jenseit des Jordans. Und warum wohnete Dan unter den Schiffen?“ Die Ephraimiter haben sich empört wider Jephtah, da er wieder kam und gesiegt hatte wider die Ammoniter, da sie zu ihm sagten, Richt. 12, 1.: „Warum bist du in den Streit gezogen wider die Kinder Ammon, und hast uns nicht gerufen, daß wir mit dir zögen?“ Und zuvor im 8. Cap. 9. 1. des Buchs der Richter haben sich die von Ephraim auch heftig gezankt mit Gideon, daß er sie nicht hatte gerufen, da er wider die Midianiter in den Streit gezogen war.

323. Darum sind zur Zeit der Richter viel große wunderbare Dinge ausgerichtet worden; aber das Volk hat niemals einem einigen Hauptmanne gehorcht: Gott hat sie gleichsam stückweise dazumal regiert, also daß ein jeglicher Stamm seinen eigenen Hauptmann gehabt, allein, daß sich bisweilen die, so am nächsten bei einander gewohnt, zusammen gehalten, als Manasse und Sebulon; aber das ganze Israel ist nie unter Einem Haupte versammelt gewesen denn allein unter den Königen.

324. So verstehe ich nun dies Stück, daß dieser große Mann und Patriarch Jakob zuvor hat vorgesehen, daß viel schweres Unglück und Gefahr über dies Volk kommen würde, von welchem Christus sollte geboren werden, und daß es derhalben von Gott sonderlich müßte regiert und erhalten werden; und daßselbe werde nicht geschehen können ohne großen Kampf und Beschwerung unter den Heiden und Feinden, von denen er wußte, daß sie alle ihre Macht und höchstes Vermögen darauf wenden würden, daß sie diese neuen Gäste im Lande Kanaan vertilgen und ausrotten möchten.

325. Denn er wird also gedacht haben, wie Einer, der weissagt und zukünftige Dinge verkündigt: Es wird ein neu unerhörtes Ding sein, wenn sie werden anfangen die Kananiter zu vertreiben aus ihrem Sitz, wie ihnen Gott verheißen hat, daß sie thun sollen. Da werden die Kananiter ihre Nachbarn, Schwäger und Blutsfreunde zu sich rufen und sich zur Hülfe bewerben, damit sie ihre Grenzen schützen und bewahren mögen; alsdann werden große schwere Empörungen folgen, dadurch die Israeliten sehr hart müssen geplagt werden. Darum wird ihnen göttliche Hülfe noth thun, daß sie ihre Feinde schlagen und überwinden mögen, und das Land einnehmen. Du Herr, mein Gott, will er sagen, hast meinen Nachkommen dieses Land verheißen, bis auf die Zukunft Schilo’s; dieweil aber große Gefahr vorhanden ist, so ihnen von den Einwohnern und Nachbarn des Landes begegnen wird, bitte ich, du wollest ihnen zu Hülfe kommen und deinem Volke Sieg verleihen, wollest sie nicht zulassen, daß die Israeliten vertilget werden möchten, um deiner Verheißung und Worts willen.

326. Also ist dieses Stück recht hinzugesetzt, den Glauben damit zu erwecken, welcher nöthig ist auch in leiblichen Dingen. Wie Isaak erstlich hat müssen durch den Glauben empfangen werden, ehe er ist leiblich von Abraham gezeuget worden. Denn wo Abraham nicht geglaubt hätte, daß er gewißlich würde von ihm geboren werden, das ist, daß Gott in seiner Verheißung wahrhaftig wäre, hätte er ihn nie gezeuget.

327. Auf solche Weise ist auch der Glaube leiblicher Dinge, die man noch nicht vor Augen sieht, einerlei mit dem Glauben der Rechtfertigung und Vergebung der Sünden, dadurch wir schließen und gewiß dafür halten, daß uns Gott günstig und gnädig sei, und dasjenige, so er verheißt, treulich und gewiß geben wolle.

328. Derhalben wünscht und bittet Jakob hier für die Kinder Israel und für seine Nachkommen und sagt: Herr, ich will auf deinen Sieg warten. Wo aber? Antwort: In meinem Samen, und in diesem Lande, so ihm verheißen ist, bis der Schiloh komme. Und was ich diesem Dan sage, dessen will ich euch alle erinnert und vermahnt haben, daß ihr nicht ungläubig werdet, auch nicht verzaget wegen der zukünftigen Gefahr. Wenn dermaßen die Feinde, die Philister und Syrer, über euch herrschen wollen, und euch sehr hart plagen und unter die Füße treten werden, sollt ihr doch nicht kleinmüthig werden, sondern auf die Hülfe warten, so euch von Gott verheißen ist. Denn daran ist kein Zweifel, mein Same wird dies Land erben bis auf den Schiloh. Diese Verheißung will ich euch geben und lehren, und will haben, daß dieselbe soll von euch ausgebreitet werden, und daß ihr sie bewahren sollt und fest darauf halten. Gott wird euch nicht verlassen, sondern ihr werdet das Land, so euch verheißen ist, einnehmen, und wenn auch Alles über und über gehen wird und die ganze Welt darüber brechen sollte. Und wie ich des Sieges warte, also sollt ihr auch für euch auf Glück und Heil, und daß ihr werdet erhalten werden, hoffen und warten.

329. Derhalben ist dies Wort, das Jakob sagt: „Herr, ich warte auf dein Heil“, ein Wort des Glaubens auch in diesen leiblichen Dingen, welche man von Gott nicht bitten oder erwarten kann, es sei denn, daß wir es mit gewissem Glauben dafür halten, daß Gott für uns sorge, uns günstig sei, uns die Sünde vergebe, und uns helfen wolle in allerlei Noth und Gefahr, nicht allein die geistlich ist, sondern auch in leiblichen Nöthen. Denn also
\quellseite{521}
stärkt Jakob auch diesen Stamm Dan und die andern alle; als wollte er sagen: Seid unerschrocken, GOtt wird euch nicht lassen; wie im 27. Psalm V. 14. steht: „Harre des HErrn, sey getrost und unverzagt, und harre des HErrn.“

\subsection{IV}

\begin{quote}
\textbf{V. 19. Gad gerüstet, wird das Heer führen, und wieder herumführen.}
\end{quote}

330. Ueber den rechten Verstand dieser Worte haben sich die Grammatiker zermaartert. In den zwei Worten gad gedud, das ist, „Gad gerüstet“, spielet er auf den Namen an, wie auch im Vorhergehenden; aber was damit bedeutet sei, weiß ich nicht. Die Hebräer und alle Ausleger wissen es auch nicht. Denn was wir Christen nicht verstehen oder auslegen, das können sie auch nicht verstehen oder auslegen; denn sie haben den Verstand oder Sinn der heiligen Schrift nicht. Die heilige Schrift aber lesen ohne Glauben an Christum, ist nichts Anderes denn im Finsterniß wandeln; wie Christus sagt Joh. 8, 12.: „Ich bin das Licht der Welt.“ Und dieweil die Juden desselben Lichts beraubt sind, ist es unmöglich, daß sie auch nur Eine Stelle oder Spruch der Verheißung recht verstehen sollten. Wie ihnen das, so von Schiloh gesagt wird im Segen Judas, gar unbekannt ist; denn sie wandeln nicht im Lichte. Und wenn sie etwa über solche Sprüche kommen, rühmen sie von ihren Werken und Würdigkeit, und lassen sich dieweil träumen und dichten sich einen solchen Messias, wie Alexander der Große oder Julius Cäsar gewesen sind. Die mit solchem Herzen die heilige Schrift lesen, würden die Zeit viel besser anlegen, wenn sie Ovids Metamorphosen läsen.

331. Solche Leute sind auch die Sophisten zu Löwen und Paris, die gar ersoffen und verkehrt sind mit dem Wahn von ihren Werken, so sie bei sich gefaßt haben, und sind beraubt des Lichts des Glaubens und Christi. Wenn sie aber aufs beste lehren, schärfen sie den Leuten die Werke der Liebe und den äußerlichen Buchstaben des Gesetzes ein. Darum ist es nicht möglich, daß sie die heilige Schrift recht verstehen oder handeln können. Denn es eben so viel, als wenn Einer in der dunkeln Nacht den Tag oder das Licht suchen wollte.

332. So will ich nun ihre Meinungen aufzählen. Gad heißt bey den Hebräern so viel als, glückselig, daher auch dieser Sohn Jakobs den Namen bekommen hat; daher kommt aber dann gedud, welches ein gemein und bekannt Wort ist und so viel als Mörder heißt, und nicht ein rechtmäßig Heer, sondern einen Haufen oder Versammlung von Mördern oder Kriegsknechten, so auf ein Scharmützel, oder das Land zu verderben und zu plündern ausgesandt wird; wie 2 Kön. 24, 2. von Josaphat gesagt wird: „Und der HErr ließ auf ihn Kriegsknechte kommen aus Chaldäa, aus Syrien, aus Moab, und aus den Kindern Ammon.“ Und 2 Kön. 13, 20. steht also: „Da Elisa gestorben war, und man ihn begraben hatte, fielen die Kriegsleute der Moabiter ins Land desselben Jahres.“ Daraus können wir abnehmen, daß sie ein solcher Haufe Kriegsknechte sind, die hin und wieder streifen und das Land verderben, nicht aber die, so mit einem rechtmäßigen Heere streiten: der Vortrab, der verlorne Haufe, die streifende Rotte, deren der türkische Tyrann viele hat, die oftmals herausfallen, und die Nachbarn umher angreifen und beschädigen. Die Ostreicher nennen es den „Sackmann“.

333. In der Grafschaft Mansfeld, welches mein Vaterland ist, ist ein Bild wie ein großer Riese gehauen oder geschnitzt, welches „Gedud“ genannt wird; und die Einwohner des Landes sagen, es sei daselbst ein Krieg geführt worden zwischen den Bischöfen und Kaiser Heinrich dem Sechsten, und daß zu derselben Zeit dies Wort „Gedud“ soll gehört worden sein. Woher aber daselbst gekommen sei, weiß ich nicht; ob vielleicht an dem Orte Juden gewohnt haben, oder ob es vom Teufel hergekommen sei.

334. Derhalben sagt Jakob, daß eine Rotte Kriegsknechte oder andere Räuber und Mörder in den Stamm Gad fallen werden, um der Feinde willen, die ihre Nachbarn waren, als, die Moabiter, Ammoniter, Syrer und Andere mehr. Es wird dieser Stamm leiden müssen, daß die Kriegsknechte sie überfallen, berauben und plündern werden, und die nächsten Dörfer oder Flecken an den Grenzen werden sie sonderlich ihres Gefallens plündern, aber mit einem rechtmäßigen Heere sollen sie nicht bekriegt werden, wie es Hieronymus im Lateinischen verdolmetscht hat; denn er hat gewohnt an den Grenzen Israel. Darum haben beide Reiter und Fußvolk, Syrer, Ammoniter und Moabiter unter diesem Stamm sehr hart angegriffen mit Rauben und Plündern, und Verderben des Landes. Aber mit dieser Hoffnung hält sich Jakob noch immer aufrecht: „HErr, ich warte auf dein Heil“, als wollte er sagen: Wir werden das Land nicht ohne große Mühe und Arbeit besitzen; aber doch wollen wir es mit GOttes Gnade und Hülfe behalten. Die Kriegsknechte werden Gad überfallen, zwacken und plündern; aber er wird auch um sich greifen und rauben, und seine Feinde wiederum plagen und ängsten. Wie aber? Antwort: Darum daß ich auf dein Heil warten will.

335. Dies scheint der allereinfältigste Verstand zu sein. Den Hebräern und Grammatikern will ich ihr Urtheil frei bleiben lassen, daß sie etwas Besseres beibringen. Im 94. Psalm V. 21. steht auch dasselbe Wort, jagodu, da in der gemeinen lateinischen Uebersetzung stehet: Captabunt in animam justi, das ist: „Sie haschen nach der Seele des Gerechten.“ Dasselbe hätte man besser also geben sollen: Concursant, turmatim conveniunt, wie unser Deutsch lautet: „Sie rüsten sich“, sie laufen zusammen etc. Denn er klagt über die Hohenpriester, Pharisäer und falschen Propheten, die sich sammeln und zu Haufen werfen, und einander helfen und mit einander rathschlagen wider die rechten Lehrer, daß sie also das Volk der rechten Lehre berauben mögen. Wie heutiges Tages die Sophisten zu Löwen und Paris sich zusammenrotten wider uns, und tödten unsere Brüder und zerreißen die armen Kirchen greulich, sie zwingen und verwunden, welche sie nur können oder mögen. Das sind nun unsere Gedudin. Wir haben es im Psalm also gegeben: „Sie rüsten sich wider die Seele des Gerechten“, das ist, es sind Löwener, Pariser, Kölner.

336. Auf dieselbe Weise sagt Jakob hier auch: Die Kriegsknechte werden sich zusammenwerfen und einen Haufen Räuber versammeln und werden diesen Stamm angreifen, der also gelegen ist, daß ihn die Nachbarn leichtlich beleidigen können; da werden sie ihn berauben und Gültige fangen. Aber Dan wird ihnen wieder auf der Ferse sein, er wird hinter ihnen her sein, er wird mit göttlicher Hülfe seinen Mördern und Räubern von Moab und Ammon nacheilen. Wie er droben von Juda gesagt hat: „Deine Hand wird deinen Feinden auf dem Halse sein“, damit angezeigt wird, daß Juda über seine Feinde herrschen werde. Desgleichen steht auch 5. Mos. 27.: „Zu der Zeit wird seine Last von deiner Schulter weichen müssen, und sein Joch von deinem Halse.“ Auf diese Weise wird Gad seine Feinde nicht überwinden, sondern wird sich schützen und den Feinden auf dem Fuße nach eilen, und sie wiederum drücken und matreiben mit GOttes Hülfe. Darum heißt Jakob seine Nachkommen GOtt vertrauen, daß sie durch seine Hülfe und Beistand vor Feinden sicher sein mögen.

337. Nun wollte ich aber gern in diesen Stamm setzen Elia und Jehu, den König Israel; wie wir in der deutschen Bibel im Segen Moses 5 Mos. 33, 21. angezeigt haben. Und es scheint, daß Mose dies eigentlich also gemeint habe; vornehmlich aber halte ich dafür, daß man Elia daher setzen und zählen soll. Ich darf es aber nicht gewiß sagen. Der Text 1 Kön. 17, 1. sagt klarlich, daß Elia der Thisbiter aus den Bürgern Gilead gewesen sei, welches beinahe also lautet, als ob er zum Stamme Gad gehört. Aber es reimt sich gleichwohl nicht alles mit einander. Denn es sind zweierlei Gilead, im Stamme Gad und auch in Manasse. Wie in Deutschland zweierlei Pannonia sind: das eine heißt man Pannonia superior, das ist Oestreich, das andere Pannonia inferior, das ist Ungerland. Also ist Gilead auch getheilt in die zwei Stämme getheilt gewesen; einen Theil hat Gad, den anderen Manasse inne gehabt. Darum wollte ich gern Elia setzen in den Theil Gads, welcher das oberste Gilead ist. Jehu ist zu Ramoth zum Könige gesalbt worden; es folgt aber daraus nicht, daß er derhalben sollte im Stamm Gad geboren sein. Darum kann ich nichts Gewisses hieraus schließen; und diese Frage gehört auch mehr zum Segen Moses, 5 Mos. 33, 20. Aber den vorigen Verstand behalte ich, nämlich, daß Gad von den Nachbarn viel habe leiden müssen, welche er doch, da sie ihn also überfallen, mit GOttes Hülfe weggetrieben hat.

\subsection{V}

\begin{quote}
\textbf{V. 20. Von Asser kommt sein fett Brod, und er wird den Königen niedliche Speise geben.}
\end{quote}

338. Hier wird nichts gesagt von Kriege oder Jammer von den gedudim. Das fette Brod, halte ich, bedeutet das fruchtbare Land. Denn das hebräische Wort, lechem, heißt nicht allein Brod, sondern begreift ingemein alle Speise in sich, auch die Früchte der Bäume. Dasselbe wirst du nicht leichtlich in irgend einer andern Sprache also finden. Darum will Jakob so viel sagen: Asser sei der allerschönste Stamm wegen der Fruchtbarkeit des Landes, und sonst vom Ueberfluß oder Reichthum der Dinge, so zu nothdürftiger Unterhaltung dieses Lebens gehören, und daß daselbst sehr guter Wein wachse, dazu Oel und ander Getreide, als an einem fetten und fruchtbaren Orte, an der Grenze des Meeres und schier am allerbesten Orte des galiläischen Landes; wie wir droben den Ort beschrieben haben, da Isachar gewohnt hat, um den Berg Thabor und von Thabor bis an das Meer und, wie die Juden sagen, bis an Sidon; wiewohl ich daran etwas zweifele. Aber das ist gewiß, daß in Galiläa ein sehr gutes und fruchtbar Land ist, also daß der Stamm Asser recht und billig möchte Bethlehem genannt werden.

\quellseite{522}

339. Er macht aber diese Fruchtbarkeit und Ueberfluß sehr groß, da er sagt: „Er wird den Königen niedliche Speise geben.“ Das ist, sein Getreide und andere Früchte der Erde werden so gut und köstlich sein, daß sie auch auf der Könige und Fürsten Tische sollen getragen werden; wie von Nabal gesagt wird, 1 Sam. 25, 36.: „Da hatte er ein Mahl zugerichtet in seinem Hause, wie eines Königes Mahl.“ Es wird wahrlich der Berg Carmel sehr gerühmt wegen der großen Fruchtbarkeit, welcher nahe bei dem Stamm Asser gelegen hat. Derhalben verkündigt Jakob dem Asser nichts Widerwärtiges, sondern nur das Allerbeste, und sagt: Es soll ein gut Land sein, daß Könige davon essen möchten.

\subsection{VI}

\begin{quote}
\textbf{V. 21. Naphthali ist ein schneller Hirsch, und gibt schöne Rede.}
\end{quote}

340. Dieser Stamm ist an der äußersten Grenze des heiligen Landes gegen Mitternacht gelegen, allernächst den Syren, und reicht bis an Damascus und den Berg Libanon. Es sind aber die Worte des Patriarchen Jakob sehr dunkel und sind beinahe von niemand verstanden worden. Die Juden bringen nichts Anderes hervor, denn daß Naphtali ein schneller Hirsch sei, das ist: wie ein Hirsch das allerschnellste Thier sei, also soll das Land Naphthali auch seine Früchte schnell hervorbringen, welche in Tempel geopfert werden, da Gott mit sehr hübschen Reden gelobt und gerühmt wird. Als könnte man dasselbe nicht auch von den andern Stämmen sagen, als, von Juda und den andern. Darum sind es harte und jüdische Auslegungen, welche von den Verheißungen gar nichts wissen.

341. Es gefällt mir besser, daß die Schnelle der Hindin gezogen werde auf die Jagd und \unsicher{Erhöhung}, oder aber auf die Zärtlichkeit und schöne Gestalt; als wollte er sagen, daß dieser Stamm würde hoch berühmt und sehr zart werden, wie eine Hindin oder Rehbolin pflegt zu sein, nämlich, das allerschönste und lieblichste unter den wilden Thieren; wie Salomo sagt Spr. 5, 18. 19.: „Freue dich des Weibes deiner Jugend. Sie ist lieblich wie eine Hindin, und holdselig wie ein Rehe.“ Item, dieweil der Stamm Naphtali Syrien zunächst gelegen hat, nach Damascus und dem Libanon hin, darum haben ihn die Nachbarn können angreifen, Unrecht und Gewalt zu thun; wie die Hindin der Jagd unterworfen ist, daß sie damit gedrängt wird. Die Könige in Syrien haben daselbst viele Gelegenheiten gehabt, ihre Truppen und Jagd zu üben: aber es ist gleichwohl dieser Stamm durch den Glauben der Väter, das ist, der gegebenen Verheißungen, und durch den göttlichen Schutz ausgefallen und aus denselbigen Stricken errettet worden; ist allezeit den Hunden entlaufen und den Spießen der Jäger entkommen, daß er gleichwohl noch ganz und unverletzet geblieben ist.

342. Also verstehe ich es, wiewohl es die andern noch auf eine andere Weise verstehen wollen. Wie in der ganzen heiligen Schrift kein Ort ist, der so mancherlei Weise zerrissen, hin und her gezogen worden wäre wie dieses Capitel. Man pflegt im deutschen Sprüchwort zu sagen: Wo viel Wege sind, da ist keiner gut. Die Unwissenheit gebärt mancherlei Meinungen. Dieweil die Juden die Sachen an sich selbst nicht verstehen, darum dichten sie mancherlei Verstand oder Auslegung.

343. Mich dünkt dieser Verstand der nächste zu sein, und der sich mit der Sache am allerbesten reimet. Denn wir behalten diese beiden vollkommenen Sprüche, als, daß Jakob erstlich rede von der geistlichen Verheißung Christi und darnach von dem leiblichen Segen der andern Stämme. Denn das ist die Verheißung: Dies Land will ich deinem Samen geben nach dir; daraus werden den andern Stämmen die leiblichen Verheißungen vorgehalten, daß sie sollen geschützt und erhalten werden bis auf das babylonische Gefängniß oder auf die Zukunft Christi. Denn jenes Gefängniß hat nicht Alles so gar verwüstet und vertilgt, daß nicht Etliche noch sollten übergeblieben sein, die wiederum aus der ganzen Welt in das gelobte Land gekommen sind. Darum ist vonnöthen gewesen, daß ihnen dies Land nicht allein hat müssen gegeben, sondern auch erhalten werden, und Gott hat sie oftmals von den Feinden, ihren Nachbarn, errettet, wiewohl dieselben immer begehrten, die Israeliten zu verderben und gar auszutilgen, und sie selbst auch bisweilen mit einander sehr uneinig gewesen sind.

344. Der Teufel ist diesem Volk immerdar feind gewesen um der göttlichen Verheißung willen, und um der Kirche willen, so in dem Volk ist versammelt worden. Denn wiewohl immer viel abgöttische Leute unter ihnen gewesen sind, war doch daran kein Zweifel, daß Gott daselbst ist geehrt und erkannt worden; wie solches der 76. Psalm V. 2. zeuget, da er sagt: „Gott ist in Juda bekannt, in Israel ist sein Name herrlich“; und im 147. Psalm V. 20. steht also: „So thut er keinen Heiden, noch läßt sie wissen seine Rechte.“ Sind sie schon nicht alle rechtschaffen und fromm gewesen, so ist doch daselbst die rechte Kirche gewesen und man hat im Volk Israel die rechte Lehre von Gott gehört. Derhalben sind diese zwei Stücke, davon wir gesagt haben, gewiß, nämlich, von dem Schild und von dem zeitlichen Segen und Erhaltung aller Stämme insonderheit, wie wir von Gad, Asser, und Dan gehört haben.

345. Also sagt Jakob auch von Naphtali: Es werden ihn die Nachbarn jagen und aufhetzen, und ihm nachstellen, welche Nachbarn sich unterstehen werden, ihn zu fangen und gar zu verschlingen; aber er soll genennet werden eine Hindin, die ausgefallen, das ist, die errettet und nicht gefangen oder in den Stricken der Feinde behalten worden ist. Darum halte ich dafür, daß dieser Verstand der rechte natürliche Verstand sei von allen Stämmen: Ich will ihnen dieses Land geben und erhalten; wiewohl wir eines jeglichen Stammes Geschichte insonderheit nicht wissen. Als, vom Stamm Gad kann ich nicht sagen, wie er durch die Räubereien und Ausfälle der gedudim geplagt worden ist, und wie er sich solcher Gefahr erwehrt und sie dagegen geschützt habe. Also werden die Jagdhunde den Stamm Naphtali plagen und verfolgen, man wird ihm mit Spießen und Netzen nachstellen; aber er wird genennet werden und soll mir heißen eine ausgefallene Hindin, die errettet ist.

346. Es zeugen auch die Historien, daß Naphtali viel erlitten habe, und der König von Assyrien hat ihn zu allererst angegriffen; wie das im 9. Cap. Jesaias zu sehen ist. Wiewohl aber nun dies nicht übereinstimmt mit den Auslegungen der Andern, ist es doch der Sache an sich und den Verheißungen gar nicht zuwider. Denn wie die geistliche Verheißung in Juda erfüllt ist: also ist die geistliche in den andern Stämmen auch erfüllt worden.

347. Aber was meint er mit den folgenden Worten, da er sagt: „Und gibt schöne Rede“? Andere legen es aus vom Opfer; aber dasselbe gehört den andern Stämmen allen zu, und vornehmlich dem Benjamin, darin Jerusalem gelegen gewesen. Daher weiß ich nicht, was dieses für eine Rede sein möge; und wie es die Andern auslegen, gefällt mir es nicht und ich weiß doch nichts Anderes dazu zu sagen. Etliche ziehen es auf den Gesang Deboras, Richt. 5, 1., welche eine Prophetin gewesen ist aus dem Stamm Ephraim, und den Sohn Abinoams berufen hat vom Stamm Naphtali und ihn heißen Kriegsvolk sammeln wider Sisera. Und da sie denselben überwunden, hat sie den schönen Gesang gesungen, welchen Jakob an diesem Ort eine hübsche Rede hat nennen wollen. Dennselben mag zwar also sein, dieweil wir nichts Besseres haben, denn wir folgen mögen. Es liegt uns aber das Wort imre noch im Weg, das Jakob hier gebraucht hat. Denn es würde sich zu diesem Gesang besser reimen, wenn er gesagt hätte simre.

348. Wenn mir es aber gezienen wollte etwas zu dichten, wollte ich es lieber also auslegen: Dieweil Naphtali der äußerste Stamm gewesen im Lande Canaan, allernächst bei den Syren gelegen, ist er oftmals durch derselben Hinterlist und auch mit offenbarer Gewalt angegriffen worden, und ist auch durch so viel Widerwärtigkeit und Trübsal unterrichtet und gewöhnt worden, sich mit größerem Fleiß und Liebe um Gottes Wort anzunehmen; wie der Prophet Jesaia sagt Cap. 28, V. 19.: „Allein die Aufsuchung“ oder Unfall, „lehret aufs Wort merken“; und man sagt im gemeinen Sprüchwort: Sauer macht essen; das Kreuz und Verfolgung oder Trübsal macht, daß uns das Wort süß schmeckt. Denn wer gar müßig und sicher ist, läßt sich den Trost, so ihm von Gott im Wort angezeigt wird, nicht bewegen, fragt nach ehrlichen Dingen gar nichts, noch viel weniger nach dem, das göttlich ist; wie Jes. 26, 16. 17. steht: „Herr, wenn Trübsal da ist, so suchet man dich; wenn du sie züchtigest, so rufen sie ängstiglich.“ Gleichwie eine Schwangere, wenn sie schier gebären soll, so ist ihr angst, schreiet in ihrem Schmerzen: so gehts uns auch vor deinem Angesicht; und im 78. Psalm V. 34.: „Wenn er sie erwürgete, suchten sie ihn, und kehrten sich frühe zu Gott.“

\quellseite{523}

349. Wenn aber Glück vorhanden ist und daß es uns nach allem unserm Willen geht, so werden wir gar faul und verderben darüber; denn das Glück macht die Leute zu Narren; wie der Papst durch den Überfluß aller Dinge und durch Sicherheit bethört, und durch die Stricke des Satans erstickt ist. Es sagt ein Poet: Stultitiam patiuntur opes, das ist so viel gesagt: Die da reich sind, gute Tage und Ruhe haben, dieselben werden gar zu geil, daß sie in aller Wollust leben und sich keines ehrlichen oder ernsten Dinges befleißigen, sondern gehen nur immerdar mit Scherz und Narrenbethätigung um. Aber denen, die versucht, geplagt und gemartert werden, wird die Schlafsucht wohl vertrieben, daß sie lernen schreien und Gott um Hülfe und Errettung bitten.

350. Von diesem Verstand halte ich, daß er sich mit diesem Text nicht übel reime; und will gleichwohl damit Andern ihre Meinung nicht genommen haben, dieweil man des rechten eigentlichen Verstandes nicht gewiß ist, nämlich, daß wenn die Leute verjagt und verfolgt werden, sie zu Gott fliehen, ihn anrufen mit Beten und Dank sagen, den lieben Gott preisen, loben und rühmen; welches sehr schöne und hübsche Reden sind unsers Mundes, und ein sehr schönes Lob und vortrefflicher Gottesdienst, nämlich, daß wir erstlich unsere Leiber in der Angst und Noth, darnach aber das Opfer des Lobes Gottes opfern; wie der 50. Psalm V. 14. 15. lehrt: „Opfere Gott Dank“, es steht aber kurz vorher: „Rufe mich an in der Noth.“

351. Diese Auslegung fließt aus zweierlei Verheißungen, wie ich droben gesagt habe. Das Land ist dem Volk Israel gegeben worden; darnach ist ihnen auch verheißen, daß sie sollen erhalten werden wider des Teufels und der umliegenden Völker grimmigen Zorn, die sich unterstanden haben, dies Volk zu vertilgen, welches sich Gott erwählet hatte, daß der Heiland Christus aus demselben sollte geboren werden; um deß willen ist es Alles zu thun.

\section[Fünfter Theil]{Fünfter Theil. Von dem Segen und Verheißungen Gottes über Joseph und Benjamin}

\subsection{I}

\begin{quote}
\textbf{V. 22. Joseph wird wachsen, er wird wachsen wie an einer Quelle.}
\end{quote}

352. Nun folgt Rachel mit ihren Söhnen, die sie dem Patriarchen Jakob geboren hat, und ist der Segen dieser Söhne eben so dunkel, wie die vorigen auch gewesen sind. Das Wort, ben porath, heißt eigentlich Mehrung, und hier soll man ziehen, das droben Jakob vor Joseph angezeigt hat, 1 Mos. 48, 3. 4.: „Der allmächtige Gott erschien mir zu Luz, im Lande Canaan, und segnete mich, und sprach zu mir: pere urebèh: Siehe, ich will dich wachsen lassen, und mehren“ etc. Diese Erscheinung aber ist geschehen, da er vor dem Männern zu Sichem floh, da ihm Rachel gestorben, und Bilha von Ruben war geschwächt worden, deren er sich hat müssen enthalten, also daß gar keine Hoffnung mehr vorhanden gewesen, wie er möchte gemehret werden und wachsen. Und Gott hat doch zu ihm gesagt: „Ich will dich wachsen lassen“ etc. Derhalben hat Jakob diese Verheißung wiederholt, da er Ephraim gesehen hat, und gedacht, dies Wachsen und Gemehretwerden muß den Joseph angehen, der von dem vornehmsten Weibe geboren war und den er für den Erstgebornen hielt. Und dieweil es dem Heiligen Geist anders gefallen hat, der ihn gelehrt, daß Juda der Vater Christi sein sollte, hat er doch in den Nachkommen Josephs zweierlei Stämme gesetzt. Und an diesem Ort sieht er auf den Namen Ephraim, des Sohnes Josephs, da er sagt: ben porath, das ist, ein Sohn der Mehrung oder des Wachsens.

353. Und in dem Wort ben soll man die Art und Weise zu reden merken, so die hebräische Sprache hat. Denn wo dies Wort gesetzt wird zu einem solchen Wort, das ein Ding bedeutet, welches nicht lebt, so bedeutet es keinen Sohn, sondern heißt eben dasselbe, was mir auf Lateinisch sagen, res, ein Ding, oder ens; als wenn ich sage: Du bist ein Sohn des Todes, da ist der Tod nicht der Vater, sondern ein solch Ding, das kein Leben an sich hat. Dergleichen, wenn man sagt: filius gladii, filius vaginae, das ist, ein Sohn des Schwertes, ein Sohn der Scheide. Und im Propheten Jeremia, Klagl. 3, 13.: filii pharetrae, Söhne des Köchers, welches daselbst Pfeile heißt. Und im 80. Psalm V. 15. 16. steht: „Suche beim diesen Weinstock, und halte ihn im Bau, den deine Rechte gepflanzet hat, und den du dir festiglich erwählet hast.“ Da nennt er den Sohn einen Weinstock und redet von dem Volk Israel. Dergleichen steht auch Jes. 5, 1.: „Mein Lieber hat einen Weinberg an einem fetten Ort.“ Da steht im Lateinischen: in cornu filio olei, hoc est, in angulo, qui est res pinguis, locus fertilis, etc., an einem fetten oder fruchtbaren Ort.

354. Derhalben, wenn Joseph ben porath genannt wird, ein Sohn des Wachsens, so wird damit nicht verstanden ein natürlicher oder Menschensohn, sondern der Sohn eines solchen Dinges, das kein Leben hat, nämlich, ein Sohn der Mehrung. Als wollte er sagen: Er wird etwas Fruchtbares sein, und das sehr viel Frucht bringen wird, nicht allein was das Land belanget, sondern auch den ganzen Stamm Ephraim und Manasse, welcher viel Volks haben wird, und wird wachsen, daß eine große Menge Volks daraus werde. Denn er redet nicht von der Person Josephs, sondern von den Nachkommen des Stammes, und dieselben lobt er wegen der Fruchtbarkeit, und daß sie so sehr wachsen und zunehmen werden.

355. Die Andern verdolmetschen es: ein Sohn eines Zweiges; und ich will darüber nicht streiten, daß es ein Zweig oder ein kleines Bäumlein heißen möge; wie Jerusalem im 48. Psalm V. 3. ein schöner fruchtbarer Zweig genannt wird. Aber hier redet Jakob mehr ingemein, nämlich, von einem fruchtbaren Dinge, als von einem Baum oder Wald; denn er versteht damit die Menge dieses Stammes, und vornehmlich Ephraim, welcher sehr viel Volks gehabt hat.

356. Und setzt noch mehr dazu und sagt: „wie an einer Quelle“, als wollte er sagen: Er wird also wachsen und Frucht bringen, wie ein Baum, der bei einem Brunnen gepflanzt ist. Und es scheint, als habe David im 1. Psalm hierher gesehen, da er V. 3. sagt: „Der ist wie ein Baum, gepflanzet an den Wasserbächen“ etc. Denn wo es trocken und dürre ist, da pfleget es auch unfruchtbar zu sein. Darum sagt auch der 24. Psalm V. 2.: „Denn er hat ihn an die Meere gegründet, und an den Wassern bereitet.“ Denn ohne Wasser können wir nicht leben und sind gemeiniglich alle Dörfer und Städte an die Wasser gebaut.

357. Auf diese Weise lobt nun Jakob den Ort oder Sitz, da dieser Stamm gewohnt, und sagt: Es wird ein fruchtbar Ding sein mit diesem Stamm, so an einer Quelle und Wasserbächen gelegen ist. Damit zeigt er an, daß dies Volk wunderlicher Weise wachsen und gemehret werden soll, und daß es ein sehr zahlreicher Stamm werden würde; wie solches aus dem Buch der Richter und auch der Könige zu sehen ist. Ephraim hat ein schön Reich und viel treffliche Männer und vornehme Hauptleute gehabt, wiewohl sie gottlos gewesen. Dazu hat er auch treffliche Propheten gehabt. Von Manasse ist der streitbare Fürst Gideon hergekommen, desgleichen Jephthah und andere treffliche Männer mehr. Und also sagt im 60. Psalm V. 9.: „Mein ist Manasse, Ephraim ist die Macht meines Haupts“, ist die vornehmste Kraft und Stärke meines Volks, ist mein bestes Volk, das ich habe, „und Gilead ist mein“, das ist mir auch gut. Daran ist wohl zu sehen, daß David auch von Ephraim und Manasse viel gehalten hat. Und also ist der Segen Jakobs erfüllt worden, der da sagt, daß diese beiden Stämme also wie an einer Quelle, die immerdar springt, gelegen sind, schön daher wachsen und gemehret werden sollen.

\subsection{Die Töchter treten einher im Regiment}

358. Es erzählen schier alle Ausleger und selbst auch Hieronymus eine Fabel von der schönen Gestalt Josephs, der so überaus hübsch und schön soll gewesen sein, daß die Jungfrauen und Weiber mit großen Haufen von den Mauern, Thürmen und Fenstern nach ihm gesehen haben, wo er über die Gassen gegangen ist. Es sind aber jüdische Fabeln und Lügen. Denn Jakob redet nicht von der Person Josephs, sondern vom ganzen Stamm, so von ihm hergekommen ist.

\quellseite{524}

359. Darum sagen wir, daß das Wort schur herkomme von scharar, welches heißt, das Regiment oder Herrschaft haben, streiten und siegen. Daher Sarah auch den Namen hat, daß ist, Sarah soll eine Herrin, eine Ueberwinderin und rechte Hausmutter sein; denn sie wird viel Anfechtungen haben, dieselben muß sie überwinden und also siegen. Darum wollen wir es lieber also geben im Lateinischen, welches auch besser ist: Filiae incedunt super regimine, das ist, wie wir es auch verdeutscht haben: „Die Töchter treten einher im Regiment.“ Da dies auch zu merken ist, daß hier die Mehrzahl mit der Einzahl verbunden wird; wie das gemein ist in der heiligen Schrift, daß Objectiv und Substantiv verschiedenen Numerus (Zahl) haben.

360. Das Wort zaadah heißt nicht, umherlaufen, sondern, stattlich und herrlich einhergehen; wie Juno sagt im Virgili: Ast ego, quae divum incedo regina etc., daher schmücken und majestätisch einhertreten.

361. Darum, will Jakob sagen, wird mein Joseph der allerschönste Stamm werden, sein Reich wird wunderbarer Weise wachsen und reichlich zunehmen. Was wird aber denn mehr werden? Die Töchter, sagt er, werden eintreten im Regiment, das ist, die Städte und Töchter der Städte werden in eine feine Ordnung und Regiment gefaßt werden, darin gute Gesetze, Zucht und Sitten sein sollen.

362. Und das ist auch eine sonderliche Art und Weise zu reden in der heiligen Schrift, wenn man sagt: „Die Städte und Töchter der Städte“, welche im Buch Josua oftmals gebraucht wird; und Jerusalem wird Sach. 9, 9. eine „Tochter Zion“ genannt. Und daher hat der Papst dieselben Namen auch auf die Kirche gezogen, da er gesagt, daß die kleinen Kirchen der großen Pfarren Töchter sind; denn die obersten Kirchen hat er die Pfarre oder bischöflichen Sitz genannt.

363. Das ist nun der rechte Verstand dieser Worte, daß er sagen will: Dies wird ein schön Reich werden, wird sehr wohl bestellt und geordnet werden mit Gesetzen, mit Amtleuten und andern nothwendigen Ständen, wird nicht ein aufrührerischer loser Haufen Volks sein; wie Ahab gerühmt wird um des schönen und herrlichen Regiments willen, so er im Volke Israel gehalten hat, wiewohl er ein gottloser Baalischer König war. Denn Jesbel sagte zu ihm 1 Kön. 21, 7.: „Was wäre für ein Königreich in Israel, wenn du so säßest?“ Das sind zeitliche Güter, damit GOtt diesen Stamm gesegnet hat, daß er sollte groß und fruchtbar sein, und sollte eine Polizei oder Regiment haben, das mit Fürsten, obrigkeitlichen Personen und Kriegsheerten wohl bestellt wäre. Aber er wird das alles nicht mit Frieden und Ruhe haben oder gebrauchen. Denn also folgt weiter im Text:

\begin{quote}
\textbf{V. 23. 24. Und wiewohl ihn die Schützen erzürnen, und wider ihn kriegen und ihn verfolgen; so bleibt doch sein Bogen fest, und die Arme seiner Hände stark, durch die Hände des Mächtigen in Jakob.}
\end{quote}

364. Das ist eine hebräische Art zu reden: viri sagittarum, das ist, Männer der Pfeile, die gute Schützen sind; wie sie auch sagen, filius vaginae, das ist, ein Sohn der Scheide, und wie im 127. Psalm V. 4. steht: „Wie die Pfeile in der Hand eines Starken, also gerathen die jungen Knaben“ etc. Jakob erzählt, wie es mit diesem Stamm zugehen werde und was er werde für Glück haben. Er wird wohl Widersacher haben und viel leiden müssen, aber er wird doch obliegen.

365. Durch diese Schützen aber verstehe ich nicht den Stamm Juda, wie die Andern haben wollen, sondern die Syrer, welche dies Reich sehr geplagt haben und treffliche Schützen gewesen sind; wie Ahab von einem Syrer mit einem Pfeil geschossen und davon gestorben ist. Es sind auch zu derselben Zeit sehr gute Kriegsrechte und Schützen gewesen, und man hat ihrer im Streit mehr gebraucht, denn derjenigen, so da Schwerter und Schleudern gehabt; wie noch jetzt die Türken mit vielen Pfeilen schießen. Aber die, so mit großem Geschütz umgehen, überwinden diese noch alle, welche damit große Schlachten thun.

366. Derhalben werden diese Syrer des Reich Juda oftmals plagen und betrüben, mit Geschoß denselben zuzusetzen und es verfolgen mit bitterem Neid und Haß, und werden es gar wollen herunter und vertilgt haben um des Glücks willen, und daß sie also wachsen und zunehmen; vornehmlich aber dieweil man daselbst GOttes Stimme und sein Wort lehrt und hört, wie es die Propheten dem Volke vortragen, wird große Feindschaft sein, die nicht beigelegt werden kann, zwischen diesem Stamme und den Syrern, aber nicht dem Stamm Juda; wiewohl derselbe auch etlichemal Ephraim angegriffen hat. Aber die Syrer haben viel greulicher mit ihnen gezürnt und ihnen mehrmals mit heftigem Schießen zugesetzt, auf daß sie dies Volk gar vertilgen und ausrotten möchten; wie im 137. Psalm V. 7. steht: „Rein ab, rein ab, bis auf ihren Boden.“

367. Aber er wird in seiner Feste sitzen, laß sie zürnen, laß die gewaltigen Schützen aus Syrien herkommen; hier sind auch Leute, ihr Bogen kann auch schießen. „Die Arme seiner Hände sind stark“, das ist: Er hat auch eine Faust, Ephraim vermag auch etwas mit seiner Faust, Bogen und Pfeilen auszurichten. Aber das ist zwischen ihnen der Unterschied: dort bei den Syrern sind mehr Schützen, die stärker, gerüsteter und in allewege größer sind. Woher sie aber diese Stärke haben, zeigt er an, da er sagt: „durch die Hände des Mächtigen in Jakob.“ Das hebräische Wort aber, abi, heißt an etlichen Orten, Ochß; doch wird es zum öftern und auch gemeiniglich GOtt zugeeignet; als, im 132. Psalm V. 2. 5.: „David schwur dem HErrn und gelobte dem Mächtigen Jakobs.“ Wir haben es verdeutscht, mächtig.

368. Das ist wahrlich eine große Herrlichkeit, daß Jakob also reden darf: Der Bogen, die Pfeile, Fäuste und Arme sind fest, kräftig und stark. Woher aber kommt ihnen dasselbe? Antwort: „Durch die Hände des Mächtigen“, welcher ist der GOtt Jakobs. Dieweil es aber nun GOtt verheißen hat, darum wird er auch diesen Stamm erhalten, welcher ein fruchtbar Ding ist, das ist, gewaltig und mächtig von vielem Volk, daß „Töchter“, das ist, Städte, „treten einher im“ politischen „Regiment.“ Aber die Widersacher haben ihn sehr hart erbittert und mit ihm in den Haaren gelegen; denn der Teufel streitet wider die Gottlosen, nicht allein mit leiblichen, sondern auch mit geistlichen Waffen. Es ist aber doch kein Bogen fest, und kein Arm stark und unüberwindlich geblieben, nicht durch seine Kräfte, sondern durch die Hände des Mächtigen in Jakob.

369. Denn dieses ist die rechte christliche Lehre, die uns auch im 44. Psalm V. 7. vorgehalten wird, da er sagt: „Denn ich verlasse mich nicht auf meinen Bogen, und mein Schwert kann mir nicht helfen.“ Man muß zwar Waffen sein in Bereitschaft haben und alles, was sonst mehr dazu gehört, das Regiment damit zu verwalten, und wie man beides im Frieden und auch in Kriegsläuften regieren soll; man muß allerdings auch etwas thun, muß die Mittel brauchen, die vorhanden sind, muß rathschlagen und Alles wohl erwägen; aber wenn das nun alles bereit und wohl bestellt ist, soll man sich darauf keineswegs verlassen, sondern soll beten, wie im 70. Psalm V. 2. steht: „Eile, GOtt, mich zu erretten, HErr, mir zu helfen“; und im 124. Psalm V. 8.: „Unsere Hülfe stehet im Namen des HErrn“ etc. Denn die Widersacher sind mit Rath und Stärke uns weit überlegen.

\subsection{Aus ihnen sind kommen Hirten und Steine in Israel}

370. Dies Stück ist etwas dunkel um der sonderlichen Art willen zu reden in der hebräischen Sprache. Aber das Wort „Stein“ wird in der heiligen Schrift durch ein Gleichniß genommen für eine Ecke, für einen Obersten, ein Haupt oder König. Als, Hizkia ist eine Ecke oder Haupt des Volks gewesen. Und Jakob will so viel sagen, daß in diesem Stamm Ephraim Hirten und Steine sein werden. Denn er braucht die Einzahl für die Mehrzahl, das ist, Eines für Viele. Wie denn dieselbe Figur in allen Sprachen gemein ist, als, daß die Lateiner sagen: Romanus proelio victor: Der Römer hat in dem Streit den Sieg behalten, das ist, die Römer haben gesiegt etc.

371. Und wird dies Wort „Stein“ nicht allein genommen oder verstanden für einen König, sondern für alle Hirten und Regenten in diesem Stamm; denn derselbe hat vor David das Regiment inne gehabt mit Manasse; wie im 78. Psalm V. 67. 68. steht: „Er verwarf die Hütte Josephs, und erwählte nicht den Stamm Ephraim; sondern erwählte den Stamm Juda“ 2c. Aber nach David und Salomon hat Jerobeam regiert, wie 1 Kön. 11, 31. ff. steht, und der HErr hat das Reich Israel erweckt; daselbst sind viel treffliche Könige und Propheten gewesen, als, Ahia, Hananin, Jona, Micha und viel andere mehr, und noch vor ihnen Elia und Elisa; das sind die Hirten und Ecken gewesen. Und das Wort „weiden“, heißt in der heiligen Schrift beiderlei, nämlich, leiblich und geistlich regieren; als, im 2. Psalm V. 9.: „Du sollst sie mit einem eisernen Scepter zerschlagen.“

\quellseite{525}

372. Derhalben ist es eine große Gnade gewesen, daß dieser Stamm so viel herrliche Steine, Propheten und Könige gehabt hat; wiewohl viele abgöttisch gewesen, die doch noch vor ihrem Tode Buße gethan zu haben scheinen, sintemal die heilige Schrift von Vielen sagt: „Und er ist mit seinen Vätern entschlafen.“

373. Es hat übrigens GOtt dies Reich und Volk erhalten um der Auserwählten willen; wie er auch noch jetzt seine Sonne aufgehen läßt über die Bösen und über die Guten. Ja, der Himmel und alle Creaturen dienen viel mehr den Ungerechten oder Bösen, als den Gerechten und Frommen; und geht doch den Gottlosen wohl nur um der Frommen willen. Die Bösen genießen des Segens der Heiligen und Gläubigen, welchen sie doch Feind sind, so ihnen doch alles Gute widerfährt nur um der Gemeinschaft und Wohlfahrt willen der Frommen.

374. Derhalben, wiewohl der mehrere Theil der Könige im Volk Israel gottlos und auch falsche Propheten unter ihnen gewesen sind, sagt doch der HErr zu Elia, 1 Kön. 19, 18.: „Ich habe mir lassen überbleiben sieben tausend Mann“ 2c. Um desselben Haufens willen hat GOtt auch die Gottlosen und Abgöttischen erhalten. Desgleichen die hundert Propheten, welche Obadja verborgen hat, so von den Gottlosen unter die Füße getreten waren, da sie mußten in das Elend fliehen: die sind es gewesen, so das Land beschützt und erhalten haben.

375. Also haben hingegen Tages die Papisten die Regimente inne und haben Macht und Gewalt in der Welt; da scheint bei ihnen die Sonne, und steht Alles sehr wohl und haben Alles vollauf. Unsere Pfarrherren aber und Kirchendiener leiden Mangel und sind verachtet und werden verstossen; wie im Prediger Salomons steht Cap. 8, 14.: „Es sind Gerechte, denen gehet es, als hätten sie Werke der Gottlosen, und sind Gottlose, denen gehet es, als hätten sie Werke der Gerechten.“

376. Darum hat der Stamm Ephraim und Manasse etliche treffliche Männer gehabt, welche „Steine“ gewesen sind, das ist, Fürsten und Propheten. Denn GOtt ist ein vortrefflicher und weiser Regent, der auch durch die böselste Obrigkeit die Regimente und Polizeien verwalten, und Frieden und dazu das Gericht im Schwang erhalten kann, eine Zeit besser denn die andere. Sie haben auch in dem Stamm GOttes Wort gehabt; und wiewohl daselbige ihrer Wenige angenommen haben, sind doch etliche fromm und gottselig gewesen, und sind siebenhundert übrig geblieben, die nicht haben ihre Kniee gebeugt vor Baal, wie droben 1 Kön. 19, 18. gesagt ist.

\begin{quote}
\textbf{V. 25. Von deines Vaters GOtt ist dir geholfen.}
\end{quote}

377. Dies ist eine sehr große und reiche Verheißung, und Mose sagt recht an dem, daß 5 Mos. 4, 7. steht: „Wo ist so ein herrlich Volk, zu dem Götter also nahe sich thun, als der HErr, unser GOtt, so oft wir ihn anrufen?“ Denn das ist eine große Herrlichkeit, wo man GOtt zum Helfer hat in allem, was man thut. Darum ist das Königreich Israel nicht ein solches Reich gewesen, da alles ohngefähr geschehen sei, was darin gehandelt ist, wie es mit dem römischen oder persischen Reich zugegangen ist, sondern es ist ein gewisses Wort gefasset und begriffen gewesen.

378. GOtt, sagt Jacob, ist es, der diesem Reich und den Propheten hilft; um der siebentausend willen, das ist, um der Auserwählten willen, muß das Königreich Israel das Wort haben, daß es öffentlich gelehrt werde, und muß eine Polizei haben. Denn der Haufe der siebentausend Männer wird ohne Zweifel nicht zu den Götzen oder Bildern zu Dan oder Bethel gelaufen sein, sondern hat im Tempel zu Jerusalem GOtt gedient und angerufen. Hosea, der Prophet, hat achtzig Jahre gearbeitet, das Volk zu lehren und die Abgötterei zu strafen, und wird ohne Zweifel Viele bekehrt haben. Desgleichen Jona und Micha haben auch Viele bei der reinen gesunden Lehre und rechtem wahrem Gottesdienste erhalten; aber nicht ohne große Gefahr. Amos hat auch im Reiche Ephraim gelehrt und gepredigt.

379. Also hat GOtt den Kindern Israel viel Propheten erweckt, daraus man genugsam verstehen kann, daß er sei dieses Volkes Helfer gewesen. Es sagt aber Jacob: „von deines Vaters GOtt“, das ist, der deinem Vater die Verheißung gegeben hat, oder, mein GOtt, dessen Verheißung ich habe, der wird dich nicht verlassen, sondern wird dir in Allem helfen.

\begin{quote}
\textbf{Und von dem Allmächtigen bist du gesegnet, mit Segen oben vom Himmel herab, mit Segen von der Tiefe, die unten liegt, mit Segen an Brüsten und Bäuchen.}
\end{quote}

380. Hier bezeichnet er nur kürzlich den leiblichen Segen. Schaddai und, wie die Griechen reden, πολυμάστορος, ist der GOtt, der Alles ernährt, der ingemein Alles speist und erhält, der viel Brüste hat. Denn das hebräische Wort, schad, heißt Brust; wie er es denn einen Segen von Brüsten und Bäuchen nennt.

381. Und es ist wahrlich eine feine Beschreibung, die er mit kurzen und herrlichen Worten gibt. Der ganze Himmel, sagt er, wird dein sein, das Erdreich und Meer, und alles, was darin ist, soll dein sein. Beinahe auf dieselbe Weise segnet auch Isaac droben seinen Sohn Esau 1 Mos. 27, 28., da er zu ihm sagt: „GOtt gebe dir vom Thau des Himmels, und von der Fettigkeit der Erden, und Korn und Weins die Fülle.“ Also sagt an diesem Ort Jacob: Der Himmel wird dir dienen, und das Meer, und alles, was das Erdreich hat. Es ist wahrlich ein großer Segen, wenn GOtt gut Wetter, Regen früh und spät gibt, dazu die Sonne und bequeme Winde zu den Früchten, so aus der Erde wachsen. Denn wo es nicht regnet, da verdorrt Alles, und vergeht Alles, wenn es so überaus trocken ist. Wenn aber die Kraft des Himmels da ist, wie es denjenigen, so den Acker bauen, selbst wünschen und begehren mögen, alsdann bringt das Erdreich nicht Dornen oder Disteln, sondern gut Korn und andere Früchte, welche nicht verletzt oder verderbet werden von übergroßer Hitze, Mehlthau, Hagel, oder dergleichen andern Unfall, dadurch die Früchte des Feldes und Bäume beschädigt werden. Da sind wahrlich die Ackerleute über solche Fruchtbarkeit des Jahres, da Alles so reichlich und wohl gerathen ist, fröhlich und guter Dinge, daß sie dies alles gebrauchen mögen; welche Gabe zu unsern Zeiten sehr selten ist.

382. Die Tiefe aber, die unten liegt, bedeutet nicht allein das große Meer, das um die ganze Welt geht, sondern alle tiefen Wasser, als, die stehenden Wasser, Seen und Pfützen.

383. Endlich verheißt er ihm auch den Segen des Leibes und Bauches, damit er fürzlich geheiße die Fruchtbarkeit der Menschen und des Viehes, und alles, was zur Nothdurft der Menschen geschaffen ist. Das alles, sagt er, wird fruchtbar sein. Die Mütter sollen nicht allein fruchtbar sein, sondern sollen auch Milch genug haben, daß sie die Kindlein säugen und nähren können. Die Kühe sollen Kälber, Milch, Butter und Käse geben, desgleichen die Schafe und Ochsen sollen Fleisch geben.

384. Dieses sind nun wahrlich große leibliche Wohlthaten; aber diemeil sie so gemein sind und fast täglich geschehen, werden sie verachtet. Ich habe oftmals Rechnung gemacht, und gefunden, daß die Menschen mehr Milch verthun und verzehren denn Wein. Die Heerden der Schafe sind doch gleichsam wie Wolken, die nichts denn Milch regnen. Aber wir gehen an solchen Gaben vorüber, und verachten dieselben, dieweil sie nun von Alters wegen, und daß sie täglich gebraucht werden, so verächtlich geworden sind; und sind doch lauter Wunderwerke und unermessliche Wohlthaten GOttes.

385. Wiederum ist es auch ein sehr schwerer Fluch, wenn wir derselben beraubt werden; wie der Prophet Hosea dasselbe Unglück den Unbußfertigen droht, da er Hos. 9, 14. sagt: „Gib ihnen unfruchtbare Leiber, und versiegene Brüste.“ Die Mütter werden gelobt wegen der Fruchtbarkeit und Geburt, und sind von GOtt mit diesen Gaben vor den Männern geziert und begnadet. Die Frucht wird in Mutterleibe, und wenn sie schon auf die Welt gekommen ist, wunderbärlich von Blut und Milch genährt. Derhalben sieht man an den Weibern größeren Nutzen als an den Männern; welches die Hurere weder sehen, noch sich dessen verwundern, sondern fromme, gottselige Ehe-männer erkennen und rühmen solche wunderbare Werke GOttes. Und bis daher ist nun gesagt worden vom Segen der Menschen und des Viehes.

\quellseite{526}

\begin{quote}
\textbf{V. 26. Die Segen deines Vaters gehen stärker, denn die Segen meiner Voreltern, nach Wunsch der Höhen in der Welt; und sollen kommen auf das Haupt Joseph und auf die Scheitel des Nasiräers unter seinen Brüdern.}
\end{quote}

386. Dieser Text ist etwas dunkel; und ich habe nun oftmals gesagt, daß mir die Glosse der jüdischen Rabbinen, die diese Worte und andere mehr auf die Person Josephs ziehen, nicht gefalle. „Die Segen deines Vaters“, sagt Jakob, „gehen stärker“, das ist: Was mir, deinem Vater, ist verheißen worden, ist noch viel größer als der Segen meiner Voreltern. Du, mein erstgeborner Sohn, und deine Söhne, so ich zu Kindern angenommen, habt größere Gaben, denn die Gaben eurer Voreltern gewesen sind, deren sie in diesem Leben haben gebraucht haben. Ich habe zwölf Söhne gezeugt, damit Gott mein Haus gemehret hat, und habe gesehen die Kinder, so von diesen meinen Söhnen geboren sind, also daß ich mit siebenzig Seelen, Knechte und Mägde ausgenommen, herab nach Egypten gezogen bin: das sind wahrlich große herrliche Gaben.

387. Denn Abraham hat wohl im Lande Sichem einen sonderlichen Ort gekauft, dahin er seine Sarah begraben; aber Ephraim und Manasse werden im Lande Canaan zwei große Theile inne haben und besitzen, und werden wunderbarer Weise gemehret werden und wachsen, daß ein groß Volk aus ihnen werde. Derhalben, sagt Jakob, ist der Segen deines Vaters über dir größer und stärker, und ist vor mir unserer Voreltern keiner gewesen, der also stark und mächtig geworden wäre, wie du Joseph in deinen Söhnen gemehret und groß geworden bist. Abraham und Isaac haben nichts mehr, denn nur einen Fuß breit oder gar keinen Raum von diesem Lande zu Sichem gehabt. Derhalben sind die Verheißungen, so mir gegeben waren, in dir erfüllt worden.

388. Daß er noch hinzusetzt: „nach Wunsch der Höhen in der Welt“, das ziehen die Juden auf ihre Lust und Begierde. Denn sie verdolmetschen und setzen anstatt der Worte „nach Wunsch“, „Ende“, und thun das aus keiner andern Ursache, denn daß sie lügen mögen, wie der Stamm Ephraim werde die ganze Welt inne haben und besitzen. Wie sie sich sonst oftmals mehr solche Lügen träumen lassen wider die klare helle Schrift, und haben vornehmlich dies Capitel mit ihren aller schädlichsten Glossen greulich zerrissen.

389. Vom Stamm Juda sagt Jakob droben: „Dem Schiloh werden die Völker anhangen“ und gehorchen; und im 72. Psalm V. 8. steht geschrieben: „Er wird herrschen von einem Meer bis ans andere, und von dem Wasser an bis zur Welt Ende.“ Desgleichen steht im 19. Psalm V. 5.: „Ihre Schnur gehet aus in alle Lande, und ihre Rede an der Welt Ende“; und im 2. Psalm V. 8.: „Heische von mir, so will ich dir die Heiden zum Erbe geben.“ Da wird das Reich über die ganze Welt nicht Ephraim, sondern Juda verheißen, und dazu noch nicht dem leiblichen Juda, sondern dem Schiloh oder Messias. Warum schreiben denn die Juden das Regiment über die ganze Welt Ephraim zu, so sie doch unter dem König Hiskia gefangen und nach Assyrien weggeführt worden sind, daß sie gar keine Hoffnung gehabt, daß sie wiederum in das Land Israel kommen würden, und haben daselbst gar keine Prophezeiung noch irgend einen Propheten gehabt, die zu gleicher Zeit mit ihnen gen Babel gezogen sind, nämlich, Daniel, Hesekiel etc. Da ist gewisse Verheißung und Hoffnung gewesen auf die Erlösung. Aber die zehen Stämme haben gar kein Wort gehabt, und alle Hoffnung, wiederum in das Land zu kommen, ist ihnen abgeschnitten gewesen. Wie sie denn auch nimmermehr sind wiedergekommen, denn allein unter dem Herrn Christo, da die Prophezeiung von Schiloh ist erfüllt worden, dem die Völker anhangen und gehorchen sollten.

390. Darum verfälschen und lästern die Juden diesen Text offenbar, da sie denselben also deuten, daß der Stamm Ephraim sich die Herrschaft über die ganze Welt sollte anmaßen. Derhalben wird dies Stück also im Lateinischen gegeben: Usque ad desiderium collium in mundo, das ist: „nach Wunsch der Höhen in der Welt“; und wollen damit der Rabbinen Auslegung verworfen haben. Und steht an diesem Orte eben dasselbe hebräische Wort, das Mose braucht, 5 Mos. 5, 21., da er sagt: „Laß dich nicht gelüsten deines Nächsten Weib.“

391. Die „Hohen“ aber werden durch ein Gleichniß genommen und verstanden für Könige und Königreiche der Welt. Und das will so viel sagen: Ephraim, dein Reich ist wie ein Reich dieser Welt, damit es streiten wird; wie David mit den Syrern, Ammonitern und Moabitern oftmals gestritten und ihr Land eingenommen hat. Aber die weltlichen Herrschaften und Regimente haben keine Propheten, kein Wort, keinen Gottesdienst, von Gott selbst geordnet, gehabt, sondern sind ohngefähr also aufgekommen und versammelt worden; wie sonst Gott den Regen streut, und läßt seine Sonne aufgehen über die Guten und über die Bösen, Matth. 5, 45. Dies Reich aber ist mit gewissem Wort und gewissen Gesetzen gefaßt und verordnet, und hat viel herrliche tapfere Männer, Propheten und Fürsten. Wie aber die Hohen in der Welt wünschen und begehren, daß sie Reichthum, glückseliges Regiment, und ein fruchtbar Land haben mögen, mit aller Nothdurft, so zu diesem Leben gehört: also, sagt Jakob, hast du auch ein solches Reich, „nach Wunsch der Höhen in der Welt“. Joseph wird das haben, das die weltlichen Herrschaften wünschen und begehren. Das, sagt er, wird auf Josephs Haupt kommen und auf seinen andern Stamm. Denn Juda und die andern Stämme haben ihr eigen Land, das sie besitzen.

392. „Die sollen kommen auf das Haupt Joseph“ etc. Mit diesen Worten wiederholt er eben dasselbe, das er auch zuvor gesagt hat, allein daß es andere Worte sind. Wer aber die Nasiräer gewesen sind, kann man aus der Beschreibung Mosis 4 Mos. 7, 2. ff. abnehmen. Simson ist ein Nasiräer gewesen, daher Nazareth den Namen hat; und von Christo steht im 132. Psalm V. 18.: „Aber über ihm soll blühen seine Krone“, welches man im Lateinischen möchte geben, nazareitas ejus. Und Matth. 2, 23.: „Er soll Nazarenus heißen“, kommt her von diesem hebräischen Worte nasir, das da heißt, absonderlich und heilig sein. Denn die Nasiräer hatten ihre sonderlichen Ceremonien, trunken keinen Wein, noch alles, was trunken macht, beschoren auch ihr Haupt nicht.

393. Zu den Zeiten der Väter sind sie noch nicht gewesen, oder aber sind ja noch nicht so berühmt gewesen. Es ist auch ihre Religion und Leben in den Historien nicht beschrieben. Mose, wie es scheint, ist der Erste gewesen, der von dieser Weise und Gebrauch geschrieben hat. Es war aber eine solche Religion, die entweder eine bestimmte Zeit gewähret, wie Paulus sein Haupt beschoren zu Kenchrea, Ap. Gesch. 18, 18., und solche sind Nasiräer gewesen einen Monat oder zwei nach eines jeglichen Wohlgefallen; oder aber es war ein solches Gelübde, das für und für gewähret hat, so lang sie gelebet haben, wie Simons Gelübde war.

394. Darum macht nun Jakob den Segen Josephs groß, nicht allein damit, daß er alles Gute haben werde, das ein König dieser Welt begehren oder erlangen möchte; sondern daß ihm solches auch mit einem gewissen Worte Gottes widerfahren soll, und daß er wie ein Nasiräer soll gehalten werden, also daß dieser Segen, gegen andere Stämme zu rechnen, etwas Sonderliches und Treffliches sein soll, und vor den andern einen großen Vorzug haben, den Stamm Juda ausgenommen.

395. Und ich glaube, daß dieser Titel oder Name nasir die Ephraimiten bethört und etwas stolz gemacht habe, daß sie sich beredet und haben dünken lassen, daß nichts ohne ihren Rath und Macht könnte oder sollte ausgerichtet werden; wie sie denn derhalben Gideon schelten Richt. 8, 1. und sagen: „Warum hast du uns nicht gethan, daß du uns nicht riefest, da du in Streit zogest wider die Midianiter?“ und zankten sich mit ihm heftiglich. Von dieser fleischlichen Ehre sind, sage ich, die Nachkommen Josephs stolz und aufgeblasen geworden, vornehmlich, da die Herrlichkeit der Erstgeburt dazu gekommen ist, welche ihr Vater Joseph erlangt hatte, als der von dem vornehmsten Weibe seines Vaters Jakob geboren war.

\subsection{II}

\begin{quote}
\textbf{V. 27. Benjamin ist ein reißender Wolf; des Morgens wird er Raub fressen, aber des Abends wird er Beut austheilen.}
\end{quote}

396. Dies ist ein kurzer Segen, und der sehr fein auf den Apostel Paulus ist gedeutet worden, aber nach der Allegorie oder heimlichen Deutung. Niemals es beinahe Ein Verstand ist nach dem Buchstaben. Denn er hat den heiligen Stephanus verschlungen wie ein Wolf; darnach hat er den Raub ausgetheilt über die ganze Welt, da er das Apostel- und Kirchenamt geführt und verwaltet, und also die Lehre des Evangeliums hin und wider ausgebreitet hat.

\quellseite{527}

397. Viele ziehen es auf den König Saul, der die Amalekiter zerstört und geplündert, und darnach das Regiment verwaltet hat. Aber es reimt sich nicht wohl, daß der soll ein Wolf genannt werden, der auf Gottes Befehl Krieg führt. Der Wolf wartet nicht so lange, bis die Schafe zum Tode verurtheilt sind, oder auf den Befehl des Hirten; sondern tobt und wüthet wider die armen unschuldigen Schäflein ohne und wider des Herrn Willen.

398. Ich verstehe es einfältig also, daß dieser Stamm gewesen sei wie ein reißender Wolf um deswillen, daß jenem Leviten sein Weib geschwächt worden. Denn da hat das Schwerdt der Benjaminiten vierzigtausend Männer gefressen und aufgerieben aus allen Stämmen Israel in zwei Kämpfen. Darnach ist sein gottseeliger Stamm gewesen als Benjamin, nicht allein um St. Pauli willen, sondern auch darum, daß sie sich zum Stamme Juda gehalten wider den König Jerobeam und die andern Stämme Israel, da sie sich doch um billiger rechtmäßiger Ursachen willen hätten mögen zu den andern Stämmen halten; denn er war Josephs rechter leiblicher Bruder. Sein Sitz oder Wohnung ist gewesen zwischen Ephraim und Manasse, er hat Juda im Süden gehabt.

\section[Sechster Theil]{Sechster Theil. Von den Segen und Verheißungen, so Jakob auf seine Söhne gelegt; item, wie Jakob seinen Söhnen Befehl gibt, ihn in Canaan zu begraben, und darauf stirbt}

\subsection{I}

\begin{quote}
\textbf{V. 28. Das sind die zwölf Stämme Israel alle, und das ists, daß ihr Vater mit ihnen geredet hat, da er sie segnete, einen jeglichen mit einem sondern Segen.}
\end{quote}

399. Mose sagt insgemein von allen und jeglichen Söhnen Jakob, wie der Vater einen jeglichen mit einem sondern und eigenen Segen gesegnet habe. Aber bei Levi, Simeon und Ruben scheint es, daß er dieselben nicht gesegnet, sondern vielmehr verflucht habe. Wie sagt er denn hier, daß ein Jeder seinen sondern Segen vom Vater empfangen habe, so doch diese mit ausdrücklichen Worten verflucht worden sind? Denn er sagt von ihnen: „Verflucht sei ihr Zorn, daß er so heftig“ etc. Und: „Meine Ehre sei nicht in ihrer Kirche.“

400. Darauf antworte ich also: Der Patriarch Jakob hat das Exempel und Sünde des Mords, so sie an ihrem Bruder begangen, und der Blutschande Rubens verdammt, er hat aber Ruben nicht gar verflucht oder verworfen; er hat nur von ihm gesagt, er soll nicht wachsen, nicht der Oberste sein, und hat ihm die Erstgeburt, oder das Priesterthum und Regiment genommen, daß er nicht sollte zwei Theile haben, sondern soll sich genügen lassen an einem kleinen Theil, und sollte nicht sterben. Wie Mose sagt 5 Mos. 33, 6.: „Ruben lebe, und sterbe nicht, und sein Pöbel sei gering“; da hat er diese Worte Jakobs ausgelegt, da er sagt: „Du sollst nicht der Oberste sein“, das ist, dein Pöbel soll gering sein.

401. Levi wird sehr herrlich gesegnet im 5. Buch Mose. Aber Simeon ist der allerverächtlichste Stamm; denn derselbe wird von seinem Vater verflucht, und Mose schweiget sein gar und gedenkt seiner nicht mit einem Wort. Hieronymus, wie ich droben gesagt habe, ist der Meinung, daß Simeons darum nicht gedacht wird, weil Judas Ischarioth von demselben Stamm hat sollen geboren werden. Aber wer da will, der mag dasselbe also glauben. Jakob will nicht haben, daß sie gar sollen untergehen, sondern sollen erhalten werden bis auf die Zerstreuung in Jakob. Er gibt ihnen keinen sonderlichen Sitz, Wohnung oder Ehre, wie er Juda oder Joseph thut, und beschuldigt und straft sie wegen einer sehr großen Uebelthat und greulichen Lasters, und bittet, daß Gott ihn und die Seinen davor bewahren wolle. „Meine Seele“, sagt er, „komme nicht in ihren Rath“, das ist, meine Nachkommen, und „meine Ehre“, das ist, das rechte Israel. Derhalben hat Simeon wie auch Ruben ohne Ehre und gar verächtlich gelebt unter den andern Stämmen, ohne sonderliche Gaben; und wird doch um des Vaters Jakob willen mit unter die Nachkommen und die zwölf Stämme gezählt. Und bis daher ist nun geredet worden von eines Jeden Segen, wie es Mose nennt.

402. Hieraus sollten man nun unsrem Exempel nehmen und Vergleichungen anstellen zwischen unseren Offenbarungen und denen, so die heiligen Väter gehabt. Denn das ist ja eine treffliche und sehr klare Prophezeiung oder Weissagung gewesen von Schiloh oder Messias, dem Sohn der Jungfrau, in dem Segen Juda, da ich gesagt habe, daß Schiloh genannt wird gleichsam eines Weibes oder Jungfrauen Sohn. Denn dieselbe Geburt ist allein eine jungfräuliche Geburt gewesen, aus dem Leibe seiner Mutter (auf Hebräisch schiljah), ist nicht ein Sohn gewesen, so aus den Lenden hergekommen oder geboren wäre. Dazu kommt das auch noch, daß dieser Jungfrauen Sohn zugleich ein König und Herr der Welt werden sollte, der das Füllen wird an den Weinstock binden und sein Kleid in Wein waschen, das ist, er sollte Vergebung der Sünden bringen, dazu Errettung vom Tode, und uns alle selig machen.

403. Dies ist eine sehr hohe Offenbarung von dem Sohne, so aus dem Leibe herkommen und zugleich ein Herr sein soll im göttlichen Wesen und Kraft. Denn das Kleid waschen, Augen haben, die da roth sind, und Zähne weißer denn Milch, daß er die Sünde hinwegnehme, beiße und verschlinge, das gebühret nicht allein einer Jungfrauen Sohn, sondern der auch zugleich Gottes Sohn sei.

404. Derhalben haben die Väter ein groß Licht und Erkenntniß gehabt von den zwei Naturen in Christo. Und Jakob zeigt an diesem Orte an, was der heiligen Väter Reden und Predigten gewesen sind vom Samen des Weibes, nämlich eben dieselben oder denen gleich, welche wir von ihm gehört, die er von den andern Vätern, so vor ihm gewesen, empfangen hat, nämlich, von dem Sohn der Jungfrau und von seinem Reich, der die ganze Welt reinigen würde und die Gläubigen mit dem Heiligen Geist und geistlichen Gaben trunken machen; wie dasselbe St. Paulus zu den Römern und Galatern reichlich auslegt.

405. Darum werden uns nun an diesem Ort herrliche Weissagungen vorgehalten; aber wenn man sie mit dem Neuen Testament vergleichen wollte, so haben wir eine weit klarere Weissagung. Denn die Erfüllung, und daß dieser Heiland nun dargestellt und gesandt worden ist, das ist ja viel größer und trefflicher, denn jenes, da er erstlich verheißen ist; nur daß wir nicht Acht darauf geben und fragen nicht sonderlich viel nach solcher unaussprechlichen Majestät und unendlichem Reichthum des Heiligen Geistes, der uns gegeben ist; wie Paulus sagt Tit. 3, 6.: „Welchen er ausgegossen hat über uns reichlich durch Jesum Christ, unsern Heiland.“

406. Derhalben sollte man die Propheten und die ganze Bibel mit großem Fleiß lesen; und ich zwar, wenn ich dies also mit einander vergleiche, zürne ich mit mir selbst und schäme mich meiner und meines Lebens, und gereut mich, daß ich leben soll, darum daß wir, nachdem Christus nun dargestellt und gekommen ist, gleichwohl so gar kalt sind, und unsere Gaben so gering achten und dem Wort so gar schlechtlich glauben; so doch die Väter mit so großer Beständigkeit geglaubt und im Glauben der Verheißung gelebt haben, dadurch sie große Gefahr und mancherlei Beschwerung überwunden haben. Wenn man sie mit uns vergleichen will, sind sie die Besten gewesen, aber sie sind die Ersten geworden; und da wir die Ersten waren, sind wir nun die Letzten geworden. Denn sie haben so gar fest an den Verheißungen vom künftigen Schiloh gehangen, daß wenn wir mit ihnen vergleichen werden, wir gar kalt, schläfrig und faul sind bei diesem Reichthum und großer Herrlichkeit der Darstellung und Erscheinung unsers Messias.

407. Sollte sich doch Einer selbst Leid anthun, wenn es nicht Sünde wäre, daß wir die große Majestät der göttlichen Werke so wenig achten. Denn ist das nicht eine große Gabe und Herrlichkeit, daß auch ein Weib in der Noth taufen kann und sagen: Ich erlöse dich vom Tode, Teufel, Sünde und allem Unglück und gebe dir das ewige Leben; ich mache aus einem Kinde des Teufels ein Kind Gottes. Aber dieser Reichthum des Geistes ist bei uns gering geworden, daß wir dessen nicht so groß achten, dieweil es täglich also geschieht, und ist doch an sich ein wahrhaftig Ding. Ein Diener des Evangeliums, der da predigt und tauft, ist ein größerer Prophet, denn Jakob oder Mose gewesen sind: die Klarheit unserer Prophezeiung übertrifft weit jener Klarheit. Sie haben wohl einerlei geistliche Speise, Trank und Glauben gehabt: aber unsere Güter sind so reich, daß auch ein Kind absolvieren könnte und aus dem Reich des Teufels in das Reich Gottes versetzen; und dasselbe durch nichts Anderes, denn durch das Wort; wie Jesaia Cap. 11, V. 6. 8. sagt: „Ein kleiner Knabe wird Kälber und junge Löwen und Mastvieh mit einander treiben“ etc., „und ein Säugling wird seine Lust haben am Loch der Otter, und ein Entwöhnter wird seine Hand stecken in die Höhle des Basilisken.“

\quellseite{528}

408. Daher sagt auch St. Petrus 2. Ep. 1, 19.: „Wir haben ein festes prophetisches Wort, und ihr thut wohl, daß ihr darauf achtet.“ O Lieber, will er sagen, habt Acht darauf, haltet nur mit dem Glauben fest daran, denn es ist ein gewiß Ding. Wir haben die heilige Taufe, Erlösung vom Tode und von Sünden. Denn der Sohn Gottes sagt selbst Joh. 3, 16.: „Also hat Gott die Welt geliebet, daß er seinen eingebornen Sohn gab, auf daß alle, die an ihn glauben, nicht verloren werden, sondern das ewige Leben haben.“

409. Auf diese Weise sollten wir nun unsere Herrlichkeit und Reichthum, so wir im Neuen Testament haben, gegen die Offenbarung der Väter im Alten Testament halten, daß wir erweckt und munter würden durch dies große Licht, und also mit fröhlichem Herzen das Evangelium annehmen und behalten möchten.

410. Aber es ist ein schändlich Ding, und ist wohl zu erbarmen und zu beweinen, daß wir so gar faul sind und die Schlafseuche uns so gar eingenommen hat, daß wir solcher großen Gaben so gar wenig achten und uns dagegen so schläfrig halten. Die schändliche verdammte „Sünde, so uns immer anklebt und träge macht“, wie zu den Hebräern am 12. Cap. V. 1. steht, läßt nicht zu, daß wir GOtt vertrauen, loben und danken mit solchem Fleiß, Inbrunst und Verwunderung, wie es die Größe dieser Sache an sich fordert. Denn es wäre nicht Wunder, wenn es Einer glaubte, daß er vor Freuden stürbe.

411. Darum klage ich oftmals darüber und ist mir leid, daß ich nicht glauben kann, so ich doch wohl weiß, daß wahrhaftig ist, was uns im Wort vorgetragen wird; welches ich nicht allein aus der heiligen Schrift gelernt, sondern es hat mich auch die Erfahrung selbst dies in mancherlei Anfechtung gelehrt und mich dessen vergewissert. Deshalb wünsche ich oft, daß ich nur aus dem Unflat meines Fleisches fahren und von so vielen Hindernissen des Glaubens erlöset werden möchte, entweder durch den jüngsten Tag, darauf ich mit großer Begierde warte und darnach ich groß Verlangen habe, oder aber sonst auf irgend eine andere Weise, wie es auch geschehen möchte. Denn ich schäme mich und bin gar verdrossen des schändlichen Unglaubens bei dem großen Reichthum der vielen Verheißungen, damit wir überschüttet und gleichsam trunken gemacht sind, wenn ich betrachte und sehe, daß die heiligen Patriarchen, solchen großen Glauben gehabt haben an die Verheißungen, die noch nicht erfüllt oder ins Werk gesetzt waren. Und andere heilige gottselige Menschen klagen darüber auch, in welchen die Natur auch wider den Glauben streitet, und können den alten Menschen nicht so hart im Zaum halten und bezwingen, daß er nicht wider den Geist streiten sollte.

412. Da kommen denn die Rottengeister und des Papsts Scholastiker dazu, die da lehren und halten, daß man immer noch im Zweifel stehen müsse, als ob wir nicht allbereit von Natur wider die Verheißung stritten. Darum ist es eine verfluchte Unsinnigkeit an den Sophisten zu Löwen, die da unverschämt sagen dürfen, daß in der heiligen Schrift an keinem Orte gelehret werde, daß uns durch den Glauben an Christum Vergebung der Sünden gegeben und das ewige Leben geschenkt werde.

413. Da die heiligen Väter in der Verheißung gelebt, haben sie das geglaubt, das Jacob vom Schilo gesagt hat, der das Füllen würde an den Weinstock binden. Und alle Propheten schreien, daß wir durch den Glauben Vergebung der Sünden haben. Und wider so viele und so herrliche Zeugnisse des heiligen Geistes unterstehen sich diese Säue und Esel diese Lehre zu bestreiten. Ist es nicht einem scheußlichen und seltsamen Ungeheuer gleich, daß man in der Kirche das lehren und halten soll, daß wir doch sonst von uns selbst thun, dieweil wir durch die Sünde zerrüttet und verderbt sind?

414. Dies soll man aber oftmals wiederholen und fleißig treiben, daß wir das Predigtamt lernen in Ehren halten, und GOtt danken, daß wir noch trefflichere Propheten sind, als die Väter und Propheten im Alten Testament gewesen. Denn jetzt zur Zeit kann auch ein Kind oder Weib zu mir sagen: Sei getrost, mein lieber Sohn, ich verkündige dir Vergebung der Sünden, absolviere dich etc. Ist es nicht also: Wer das hört und glaubt, daß derselbe Vergebung der Sünden und das ewige Leben hat? Ist es denn nicht eine große greuliche Narrheit und Unsinnigkeit, daß man lehrt, man soll daran zweifeln und sagen, dies alles werde uns in der heiligen Schrift nicht vorgehalten; ja, darüber noch heftig streiten, daß die Schrift dieser Lehre stracks zuwider sei? O welche erschreckliche Zeit ist das und voll großer Gefahr, in welche wir armen Kinder gekommen und gerathen sind!

415. Daran ist wahrlich kein Zweifel, daß der, so die heiligen Sacramente reicht und austheilt, ein Prophet ist, und der sie empfängt und glaubt, ist auch ein Prophet. Denn das sind ja Worte eines großen Propheten, der da sagt, Matth. 26, 26. 28.: „Das ist mein Leib, das ist mein Blut.“ Christus ist der Prophet, welcher 5 Mos. 18, 15. 18. 19. ist verheißen worden, da Mose sagt: „Einen Propheten, wie mich, wird der HErr, dein GOtt, dir erwecken, aus dir und aus deinen Brüdern, dem sollt ihr gehorchen. Ich will“, spricht der HErr, „meine Worte in seinen Mund geben; der soll zu ihnen reden alles, was ich ihm gebieten werde. Und wer meine Worte nicht hören wird, die er in meinem Namen reden wird, von dem will ichs fordern.“ Sollen wir denn diesem Propheten nicht glauben, so doch zu dieser Verheißung eine schwere Drohung gesetzt ist, welche endlich den schwärmerischen und gottlosen Geist der Sophisten zu Löwen zu Schanden machen und verdammen wird?

416. Darum ist es vonnöthen, daß wir das liebe Predigtamt der Kirche, welches unsere Herrlichkeit ist, groß machen und hoch rühmen. Denn unser Glaube hält sich jetzt nicht allein an die Verheißungen, sondern auch an dasjenige, so verheißen war und uns nun gegeben ist. Wir haben jetzt den HErrn selbst, der mit uns redet, und uns viel herrlichere und klarere Verheißungen vorhält, als die gewesen sind, so den Vätern gegeben waren. In der Taufe sagt Christus: „Ich erlöse und errette dich aus der Gewalt des Teufels, und übergebe dich meinem himmlischen Vater. Da muß ich es wahrlich gewiß dafür halten, daß Christus nicht lüge; und je fester ich glaube, je heiliger und seliger ich bin. Und wir haben zwar keine Ursache überall, daß wir da zweifeln sollten, sintemal

417. Ich kann auf solche Lästerung nicht deuten oder davon etwas sagen, ich werde darüber in meinem Herzen heftig bewegt. Denn was thun doch die heillosen Leute Anderes, denn daß sie damit die Erbsünde bestätigen, welche sonst ihrer Art und Natur nach nicht gern glaubt? Um deswillen ich mit mir selbst oftmals hart zürne, daß da ich herzlich gern glauben wollte mit so großem Glauben, wie die Wahrheit, Gewißheit und Versicherung der Verheißung fordern und haben will, also daß ich durch einen ehern Berg brechen könnte, und alle Gefahr und Schaden, so mir vorkommen kann, unerschrocken verachten; dennoch das Gesetz in meinen Gliedern wider mich streitet und mich gefangen nimmt. Diese böse, verderbte und schändliche Natur loben und preisen nun die Sophisten zu Löwen, und lehren eben dasselbe, dazu doch sonst die Natur uns arme Menschen zwingt, daß wir es thun müssen. Es ist gleich, als wenn ich zu einem Hurer oder Dieb sagte: Du sündigest daran nicht, sondern thust recht daran, daß du also nach deiner Lust Unzucht treibst und andern Sünden und Schanden nachhängst; da doch die Natur von selbst das also mitbringt, daß die jungen Leute zur Unzucht geneigt sind; wie alle Menschen von Art und Natur auch zu andern Lastern und Sünden, zu Hoffart und Geiz, und dergleichen Sünden mehr, geneigt und beflissen sind.

418. Die Sophisten verstehen nicht, was Glaube oder was Unglaube sei. Darum verdammen sie Christum in seiner Rede, und verdammen die auch als Ketzer, die da lehren, daß man den Verheißungen Christi glauben solle. Solche greuliche und höllische Unsinnigkeit des Teufels wird niemand groß genug machen können, sie übertreffen wahrlich allen Verstand, und kann sie niemand mit Worten genugsam ausreden. Es kann doch nichts Scheußlicheres geredet oder erdacht werden, als das ist, daß sie die Leute zwingen und treiben zu sündigen. Es ist gut, sagen sie, daß du deines Glaubens nicht gewiß bist; Christus redet eben die Worte der Verheißung zu dir nicht, und wenn du an Vergebung der Sünden zweifelst, so bist du ein Christ.

\quellseite{529}

419. Das ist des Papstes und seiner Theologen Theologie, und ist der Antichrist selbst. Der Türke ist doch noch etwas bescheidener, und macht es ja nicht so gar grob, welcher bekennt, daß Christus ein Prophet sei und daß er zur linken Hand GOttes sitze. Aber die groben Esel zu Löwen verleugnen offenbar, daß Christus ein Prophet sei. Wiewohl sie dasselbe mit Worten sagen, widersprechen sie doch demselben mit der That, und sagen: Wir haben nicht Vergebung der Sünden und das ewige Leben durch Christum. Sie verleugnen und verwerfen, so viel an ihnen ist, die schönen hellen Sprüche, welche voll Trostes und Lebens sind, als der ist, da der Täufer Johannes sagt Joh. 1, 29.: „Siehe, das ist GOttes Lamm, welches der Welt Sünde trägt“; item, da St. Petrus sagt Ap. Gesch. 10, 43.: „Von dem zeugen alle Propheten, daß durch seinen Namen alle, die an ihn glauben, Vergebung der Sünden empfangen sollen.“ Desgleichen sagt er auch Ap. Gesch. 4, 12.: „Es ist kein anderer Name den Menschen gegeben, darinnen wir sollen selig werden.“

420. Dies sage ich darum, daß ich damit euch ermahnen und erwecken möge zu Betrachtung und hoher Bewunderung des Glaubens der lieben Väter, und unserer Herrlichkeit und Offenbarungen, so wir im Neuen Testament haben; daß wir unsere Schwachheit und Unglauben beweinen, welcher gar unwürdig ist der großen Güter, die uns im Evangelium angeboten und vorgetragen werden. Denn da uns jetzt das Heil viel näher ist, wie Paulus sagt Röm. 13, 11., glauben wir doch fast schläferlich, und geht schwach zu, daß wir dem Wort GOttes beifallen. Die Väter sind sehr brünstig im Glauben an den zukünftigen Christum gewesen und der noch hat sollen geoffenbart werden, in welchem sie sich gerühmt und den Teufel überwunden haben, und dazu auch allerlei Jammer, so ihnen begegnet ist. Wir aber werden nicht sonderlich fröhlich über die Verheißungen Christi, der jetzt nun gegenwärtig ist und regieret. Und über das fährt noch das ganze Papstthum zu und verdammt diese Lehre öffentlich.

421. Derhalben lasset uns Fleiß anwenden, daß wir unsere Herrlichkeit im Neuen Testament groß machen, und von Tag zu Tag immer wachsen und zunehmen mögen in Erkenntniß Christi, das ist, im Glauben, welcher nicht zweifele, sondern es gewiß dafür halte, daß ihm in Christo Vergebung der Sünden und das ewige Leben geschenkt sei. Denn wer das glaubt, der ist ein rechter Christ und ist mit in der Zahl der Kinder GOttes.

422. Die wüthenden, unsinnigen und gottelsterlichen Sophisten sind nicht Menschen, ja, das noch mehr ist, sie sind nicht werth, daß sie Säue heißen, sondern sollen billig Teufel genannt und unter die Teufel auch gerechnet werden, dieweil sie sich unterstehen vor der christlichen Kirche zu lehren wider den Heiligen Geist, daß man zweifeln solle an GOttes Gnade und Vergebung der Sünden um Christi willen. Der Türke, wie ich zuvor gesagt habe, ist noch so gar unerfahren nicht; denn wo er die heilige Schrift hätte, würde er wahrlich das nimmermehr anfechten.

423. Wir aber sollen lernen unsern reichen Segen groß machen und hoch heben, welcher ungleich größer und herrlicher ist, denn die Segen, so die Väter gehabt. Wir haben ja die Schlüssel, das Evangelium, die heiligen Sacramente und reiche Vergebung der Sünden, und das ewige Leben, so uns auf mancherlei Weise angeboten und vorgetragen wird. Unser lieber HErr GOtt helfe uns durch den Heiligen Geist und seine Gnade, daß wir dies alles mit festem Glauben annehmen und behalten mögen.

\subsection{II}

\begin{quote}
\textbf{V. 29—32. Und er gebot ihnen, und sprach zu ihnen: Ich werde versammelt zu meinem Volk, begrabet mich bei meine Väter in der Höhle auf dem Acker Ephron, des Hethiters, in der zweifachen Höhle, die gegen Mamre liegt, im Lande Canaan, die Abraham kaufte, sammt dem Acker, von Ephron, dem Hethiter, zum Erbbegräbniß. Daselbst haben sie Abraham begraben, und Sara, sein Weib. Daselbst haben sie auch Isaak begraben, und Rebekka, sein Weib. Daselbst hab ich auch Lea begraben; in dem Acker und der Höhle, die von den Kindern Heth gekauft ist.}
\end{quote}

424. Von dieser Höhle ist droben genug gesagt worden. Hier aber siehst du, daß der heilige Vater Jakob in sehr großem Glauben gestorben ist. Denn diese Segen zeigen genugsam an, was er bekannt und gelehrt, und auch wie er gelebt habe. In diesem Schloß hat er gelebt, in demselben stirbt er auch mit Frieden, und bietet seinen Söhnen, daß sie ihn begraben sollen, nicht in Egypten, sondern bei seinen Vätern im Lande Canaan; nicht allein um der Ursache willen, daß das Volk ein ewiges Gedächtniß ihrer Väter in dem Lande haben möchte, damit es nicht wiederum nach Egypten flöhe, oder sonst zu anderer Hülfe der Heiden, sondern auch darum, daß sie sich dabei wissen sollten zu erinnern, daß die Väter die Verheißung gehabt hätten; und also gedachtet: Die schlafen bei uns, mit denen GOtt geredet hat, von welchen wir das Wort und Verheißung empfangen haben, daß wir nicht fremden Göttern nachfolgen sollen. Dies ist eine große wichtige Ursache gewesen, warum Jakob hat wollen auf dem Acker begraben werden, nämlich, zum Zeugniß und Erhaltung der Predigt und rechten Lehre bei den Nachkommen. Es ist aber dies die Ursache nicht allein, sondern er hat damit auch angezeigt, daß er hoffe und warte auf die Auferstehung der Todten, so in der Auferstehung Christi geschehen sollte; wie denn Viele mit Christo auferweckt worden sind, Matth. 27, 52. 53., und sind Vielen erschienen.

425. So ist nun das die erste Ursache gewesen: damit daß Volk ein ewiges und lebendiges Gedächtniß ihres Vaters haben möchte, daß sie wüßten, sie wären ja nicht von den Heiden geboren, sondern von dem Vater, der die Verheißung gehabt, der gepredigt hat von der Jungfrau Sohn, an welchem Jakob selbst auch geglaubt hat; und haben eben denselben Glauben die Nachkommen auch also behalten und behalten ihn jetzt auch noch.

426. Und wiewohl ich schwach bin, zweifle ich doch nicht, sondern halte es mit dem Evangelium, wiewohl der Satan dawider ficht, und die Sünde, so der Natur anklebt; welches Uebel die Sophisten zu Löwen sich unterstehen auch noch größer zu machen, und reizen das Fleisch zum Aberglauben und gottlosen Wesen, welches doch sonst uns arme Menschen zum Zweifel treibt.

427. Die andere Ursache dieses Begräbnisses ist, daß ich ewiglich glaube, daß diese Väter, Abraham, Isaak, Jakob und Joseph unter denen gewesen sind, die sammt Christo auferstanden sind. Denn befiehlt und gebietet Jakob so fleißig, daß sie ihn sollten legen und begraben in der Väter Begräbniß. Und dies ist das Bekenntniß seines Glaubens, davon er hat haben wollen, daß es nicht allein sollte im Wort stehen, das er gelehrt hat, da er noch am Leben gewesen ist, sondern auch in seinem Tod und Begräbniß, daß dies Gedächtniß ihre Nachkommen immer möchten vor Augen haben, und daß sie die Väter höreten, die gleichsam aus ihren Gräbern ihnen eben die Lehre vorhielten, welche sie ihnen auch vorzeiten, da sie noch am Leben waren, vorgehalten hatten, als, von der leiblichen Verheißung des Landes und von dem gebenedeiten Samen. Da liegen sie, sollten sie sagen, die es uns gesagt und gelehrt haben: darum wollen wir auch nun also glauben, und darauf warten, daß uns verheissen ist.

428. Also haben wahrlich die Propheten geglaubt, wie ihre Schriften zeugen; wiewohl ihrer keiner gewesen ist, der nicht über seine eigene Schwachheit klagte, und begehrte, daß er hätte mögen stärker sein im Glauben; wie wir auch schwach sind, und wünschen doch und seufzen darnach, daß wir möchten einen vollkommenen Glauben haben, und in unserm Gewissen also versichert sein, daß wir an der Gnade und Wohlthaten Christi, so er uns bewiesen hat, gar nicht zweifeln möchten.

429. Derhalben sollen wir uns freuen, daß wir würdig sind dieser Anfechtung und Verfolgung, so wir jetzt leiden müssen um dieses Bekenntnisses willen; und lasset uns für unsere Feinde bitten, daß doch Etliche mögen bekehrt werden von denen, die jetzt vor unsern Ohren und zu unsern Lebzeiten an der Lehre und Glauben Christi so unverschämt zweifeln dürfen, und sich unterstehen, dieselben zu vertilgen. Der HErr wolle sich ihrer erbarmen, die noch zu bekehren sind, die Andern hole der Teufel weg.

\begin{quote}
\textbf{V. 33. Und da Jakob vollendet hatte die Gebote an seine Kinder, that er seine Füße zusammen aufs Bette, und verschied, und ward versammelt zu seinem Volk.}
\end{quote}

430. Ihr habt droben gehört diese Art und Weise zu reden, daß die heilige Schrift den Tod der Heiligen, ja, auch Himmels Tod nennt, „versammelt werden zum Volk seiner Väter.“ Also wird an diesem Ort nicht gesagt, daß Jakob gestorben sei, sondern er thut seine Füße zusammen auf das Bett, ist matt geworden und ward versammelt zu seinem Volk. Denn der Heilige Geist hat dies also mit sonderlichem Fleiß beschreiben wollen, daß er uns damit versichern möchte, daß die Väter im Glauben an die Verheißung und Hoffnung der zukünftigen Unsterblichkeit und ewigen Lebens gestorben seien. Denn wenn sie Papisten, Bischöfe oder Cardinalé gewesen wären, welche des zukünftigen Lebens beider der Gerechten und Gottlosen spotten und dasselbe nur für eine Fabel halten, so würden sie ihr Begräbniß nicht so fleißig an einem gewissen bestimmten Ort bestellt haben.

\quellseite{530}

431. Es sind wahrlich große prächtige volle Worte, da Mose saget: „Er ward versammelt zu seinem Volk.“ Und ist eine solche Weise zu reden, die gänzlich auf die Heiligen und Gottseeligen gehört. Und ich halte deshalb dafür, daß auch Israel sei selig geworden; denn es steht von ihm auch geschrieben 1 Mos. 25, 17.: „Er ward gesammlet zu seinem Volk.“ Darum soll uns diese Art und Weise zu reden lieb und angenehm sein; denn sie zeigt an, daß die Heiligen von Anfang der Welt im Glauben und Hoffnung der Auferstehung entschlafen sind.

432. Die davon disputiren, ob auch die Heiligen in den Schooß der Väter, limbus patrum, wie sie es geheißen haben, sind geführt worden, oder aber, ob sie in die Gemeinschaft der Engel aufgenommen sind, dieselben zanken darüber ohne Aufhören; und ist doch eitel vergeblich Ding, darum sie zanken. Es ist besser, daß wir bei dem gewissen Wort bleiben, wie wir droben etlichemal gesagt haben; welches Wort zeuget, daß nach diesem Leben ein Volk sei. Denn Jakob ist nicht aufgefahren gen Himmel, ist auch nicht in die Hölle gefahren. Wo ist er denn hingekommen? Gott hat einen sonderlichen Ort oder Behältniß, darin die Heiligen und Auserwählten ruhen ohne Tod, ohne Pein und ohne Hölle. Wie aber derselbe Ort genannt werde, oder was das für ein Ort sein möge, weiß niemand; das ist aber gewiß, daß ein Volk genannt wird und auch da ist.

433. Der Juden Rabbinen dichten allerlei, was sie können und mögen. Christus aber sagt allein, Luc. 16, 22., Lazarus sei gestorben und von den Engeln in Abrahams Schooß getragen worden. Diesem Ausleger und Lehrer soll man glauben; denn er ist Gott und Mensch. Der Schooß Abrahams aber ist der Ort, an welchen alle Gerechten oder Gläubigen nach dem Tod geführt und aufgenommen sind, und Abraham wird darum dabei angezeigt, weil er ist ein Vater der Verheißung gewesen. Es wird Noah hier nicht genannt, Adam auch nicht, die doch auch in demselben Schooß versammelt worden sind. Denn Adam ist nicht der Erste unter denen, die gestorben sind, sondern Abel, der von seinem Bruder Kain ist erwürget worden. Henoch ist lebendig vom Herrn weggenommen worden, daß er den Tod nicht sähe, und wird doch nicht der Schooß Adams oder Abels genannt, sondern Abrahams: entweder daß die Verheißung, so dem Adam geschehen, nicht so klar gewesen ist, oder aber sonst um einer andern Ursache willen.

434. Wie aber von uns in der Kirche auch gesagt wird, daß wir durch die Taufe, durch das Evangelium und Sacrament zu unserm Volk versammelt werden, und wissen, daß wir unter unserm Volk sind: also, wenn wir sterben in der Verheißung Christi, dem Abraham geschehen, werden wir auch durch die Engel in den Schooß Abrahams oder zu unserm Volk geführt und aufgenommen. Und wir haben oftmals gesagt, daß die Verheißungen der Väter zu der Auferstehung und zum künftigen Leben gehören, und nicht allein zu diesem zeitlichen Leben, sondern auch zu dem andern geistlichen und ewigen Leben, in welchem die Werke dieses natürlichen Lebens, wie wir genähret und gezeuget werden und dergleichen andere Werke mehr, aufhören werden. Wo aber dasselbe Volk sei, wissen wir nicht.

435. Der Papst dichtet, daß viererlei Ordnungen der Verstorbenen seien. Die erste ist die Hölle der Verdammten. Die andere das Fegefeuer. Und Thomas Aquinas sagt, daß die Hölle gleichsam der Mittelpunct sei, da das Fegefeuer umhergehe. Bei demselben aber ist der dritte Bezirk oder Ort der unmündigen Kinder, so nicht getauft sind. Der vierte Ort ist, sagen sie, limbus patrum, der Schooß der Väter, darin die Gottseeligen gewesen sind vor der Auferstehung Christi.

436. Dies sind lauter Träume und Menschengedichte. Petrus und Paulus sagen klarlich, Eph. 6, 12., daß die bösen Geister in der Luft schweben, und 2 Petr. 2, 4. wird auch gesagt, daß Gott der Engel, die gesündigt haben, nicht verschonet habe, sondern hat sie mit Ketten der Finsterniß zur Hölle verstoßen, und übergeben, daß sie zum Gericht behalten werden. Damit bin ich zwar wohl zufrieden, und will nicht etwas Höheres suchen, denn von den Aposteln gelehrt worden ist. Des Fegefeuers wird in der heiligen Schrift gar nicht gedacht, sondern es ist des Teufels Gedicht, daß die Papisten etliche Jahrmärkte hätten und dem Geld damit nachstellen möchten. Die Sophisten halten es mit dem Papst um Thomas willen, aber der Thomas geht uns nichts an. Augustin gedenkt an etlichen Orten des Fegefeuers, aber er redet sehr dunkel. Darum glaube ich nicht, daß diese vielerlei Ordnungen wahr sind. Denn die Schrift redet nicht also, sondern zeuget, daß die verstorbenen Heiligen versammelt werden zu ihrem Volk, oder zu denen, die an den Messias geglaubt und auf seine Zukunft gewartet haben; wie Adam mit allen seinen Nachkommen im Glauben Christi gestorben ist. Wie aber dieselbigen an sonderlichen bestimmten Orten bewahret werden, wissen wir nicht.

437. Wir haben daran genug, daß die Schrift zeuget, daß vom Anfang der Welt die nicht sind verloren worden, die an den Samen des Weibes geglaubt haben, und daß ihrer nicht ist vergessen worden, sondern daß sie versammelt worden sind zu ihrem Volk. Was aber dasselbe für ein Ort sei, kann niemand sagen.

438. Ich habe droben ein grob Gleichniß angezogen von dem Leben der Frucht im Mutterleibe. Denn niemand ist unter allen Menschen, die jetzt leben, der da wissen könne, wo er die ersten zwei Jahr gewesen sei, da er im Mutterleibe lebte, und da er auf die Welt geboren war und seiner Mutter Milch gesogen hat. Er weiß nicht, wie die Tage und die Nächte oder Zeiten gewesen sind, oder auch wer die gewesen sind, so ihn regiert und sein gewartet haben: und hat doch dazumal gelebt, und ist Leib und Seele mit einander vereinigt und zu allen natürlichen Werken tauglich und geschickt gewesen. Derhalben ist dies das gewisseste Argument und Anzeigen, daß Gott den Menschen wunderbarerlicher Weise erhalten wolle, die ihm, dem Menschen, selbst gar unbekannt ist. Siehe an den großen Jammer und die harten Anstöße derer, so mit der fallenden Seuche und Wahnsinn betrübt und beschweret sind, dadurch sie so gar von sich selbst kommen, daß es scheint, als ob sie aller Sinne gar beraubet sind. Ja, auch des Nachts, wenn wir liegen und schlafen, wissen wir nicht, an welchem Ort wir schlafen, ob wir bei Freunden oder Feinden sind.

439. Wie ist das denn so eine große Thorheit und Frevel, daß wir die göttliche Weisheit so einschließen und fassen wollen in die enge menschliche Vernunft! Können wir nicht denken, daß der die Seele bewahren und erhalten werde, daß ich eben nicht wissen kann, wo ich sei, und bin doch an einem gewissen bestimmten Ort, dieweil er dasselbe täglich thut, wie die Exempel anzeigen, so wir jetzt erzählt haben.

440. Darum sollen wir Gott die Ehre geben, und ihn damit ehren, daß wir es dafür halten, daß seine Weisheit unermäßlich groß sei, und daß er mehr und wunderbare Weise hat, uns zu erhalten, als wir mit unsern Sinnen und Verstand begreifen können. Er bewahrt und erhält uns anders in Mutterleibe, anders in der Wiege, auf eine andere Weise im Schlafe, und aber auf eine andere Weise, wenn wir krank sind. Wir haben daran genug, daß wir wissen, daß die Heiligen im Alten Testament gestorben sind im Glauben an den zukünftigen Christus, und daß die Gottseeligen im Neuen Testament, so im Glauben an denselben Christus, der nun gekommen ist, sterben, zu ihrem Volk versammelt werden.

441. Wie aber, oder was es für ein Ort sei, wissen wir nicht. Es ist aber gewiß, daß es kein böser Ort sei, und daß sie gar mit keinerlei Pein gemartert oder geplaget werden, sondern ruhen in Gottes Gnade. Wie sie in diesem Leben auch sanft schlafen und in Gottes und der lieben Engel Schutz und Schirm sind, und sich vor keinerlei Gefahr fürchten, wenn sie auch mitten unter den Teufeln stünden: also sind auch nach diesem Leben die Seelen der Gerechten in Gottes Hand, wie im Buch der Weisheit steht am 3. Cap. V. 1.

442. Was es aber mit der Ruhe für eine Weise habe, oder wie es damit beschaffen sei, hat niemand je gewußt, ausgenommen die mit Christo auferstanden sind, oder die er von den Todten auferweckt hat, als, Adam, Eva, Noah, Abraham und Andere mehr, desgleichen der Wittwe Sohn zu Nain und die Tochter des Obersten der Schule, des Jairus.

443. Daran ist aber kein Zweifel, die, so zu ihrem Volk versammelt sind, ruhen im Frieden. Ja, sagst du, wie geht aber das zu? Sage du mir, wie das zugehe, daß du schläfst, oder auf welche Weise die Frucht im Mutterleibe erhalten werde. Deine Mutter hat wohl gefühlt,
\quellseite{531}
daß du dich in ihrem Leibe gezeugt und bewegt, und gelebt, du weißt aber nicht, wie du gelebt, mit den Händen getastet und gesehen hast etc. Du hast auch deiner Mutter Brüste gesogen, da du gar nichts davon gewußt hast; und willst nun jetzt wissen und verstehen, wie dich GOtt im Frieden bewahre und erhalte nach dem Tode!

444. Derhalben sollen wir glauben, daß ein sonderlicher Ort sei der Auserwählten, da sie alle ruhen, ausgenommen die, welche nach GOttes Willen mit Christo auferstehen sollten. Denn es hat GOtt also wohlgefallen, daß er mit ihrer Auferstehung den Tag der Eröffnung, der Auferstehung Christi, hat ehren wollen. Und ich glaube, daß auch Joseph dazumal sei auferweckt worden, welche alle in den Gräbern gelegen haben bis auf die Zukunft Christi. Das sollen wir glauben, und nicht so vorwitzig und genau forschen nach der Weise dieser Ruhe. Laß dir viel lieber das befohlen und angelegen sein, daß du Einer von GOttes Volk seiest, und glaubest, daß Christus von den Todten auferstanden sei. GOtt wird dafür wohl sorgen, daß du auch bei dem Volk der Auserwählten sein werdest, der da genugsam und zuwege gebracht hat, daß du in Mutterleib, in der Wiege und im Schlaf gesund und lebendig wärest.

445. Und es ist möglich ein wunderbärlich Ding, daß uns GOtt mit Schlafen, Essen, Wachen und andern Dingen mehr den unnützen Thieren gleich macht. Denn die Seele des Menschen schläft also, daß alle Sinne gleichsam begraben sind, und unser Bett ist wie ein Grab, darin doch nichts Verdrießliches oder Beschwerliches ist. Also ist auch an diesem Ort, dahin die Todten geführet werden, gar keine Pein, sondern, wie man zu sagen pflegt, sie ruhen daselbst im Frieden; wie kein Anzeichen irgend eines Schmerzes oder Mühseligkeit gespüret wird an denen, so aus dem Schlaf auferweckt sind und jetzt nun wachen; wiewohl sie nicht wissen, was mittlerweile, dieweil sie geschlafen, geschehen sei. Eben dasselbe geschieht auch, wenn ein Mensch in seinem Gemüth entzückt wird, das sie ecstasis nennen; wie Augustin erzählt von einer Entzückung, so er selbst und auch seine Mutter soll gehabt haben, daß sie beide nichts davon gewußt haben, da sie wiederum zu sich selbst gekommen sind, an welchem Ort sie gewesen wären.

446. Darum ist der Tod der Heiligen ehrlich, und wie der 116. Psalm V. 15. sagt, ist er vor dem HErrn köstlich und werth gehalten. Die Ruhe und Ochsen haben keinen Ort, dahin sie versammelt werden, wie die Heiligen haben, wiewohl wir denselben in diesem Leben nicht wissen; wie wir denn sonst viel andere Dinge mehr auch nicht wissen. Ich habe oftmals versucht, den Augenblick der Zeit zu merken und zu behalten, darin ich entweder einschlafe oder aber wiederum aufwache; ich habe aber dasselbe nie ergreifen können, oder den vorkommen, daß mich der Schlaf darüber nicht übereilt und betreten habe, ehe denn ich es gedacht oder gemeint hätte. Gleich also wird auch unser Tod und Auferstehung sein. Wir fahren dahin und kommen am jüngsten Tag herwieder, ehe wir es gewahr werden, wissen auch nicht, wie lange wir aus gewesen sind.

447. Darum sagen wir, daß kein Fegefeuer, oder limbus patrum (Schooß der Väter) sei, darin sie sagen, daß noch gewisse Hoffnung sein soll, da man auf die Erlösung wartet; sondern dasselbe sind lauter Gedichte irgend eines groben ungelehrten Sophisten. Denn eben diese Hoffnung und Erwartung würde sie plagen; wie der Spruch lautet Spr. 13, 12.: „Die Hoffnung, die sich verzeucht, ängstet das Herz“; wie wir in der Hoffnung leiden, wiewohl wir darüber nicht zu Schaden werden. Aber jenes Volk GOttes nach diesem Leben leidet keine Angst von Verlangen der Erlösung, sondern schläft den sanftesten Schlaf und in guter stiller Ruhe; wie Christus Matth. 5, V. 39. sagt von der Tochter des Jairus: Das Mägdlein ist nicht todt, es lebt und ruht etc.

448. Zuletzt, daß Jakob seine Füße zusammen thut auf dem Bette, damit wird angezeigt, daß solches eine sonderliche Weise desselben Volks Gott gewiesen sein. Und zeigt Mose an, daß Jakob im Bette gesessen, da er mit seinen Kindern geredet habe; da er aber ausgeredet, hat er sein Haupt auf das Bette gelegt und auf demselben auch die Füße wieder nach sich gezogen. Zeigt auch damit zugleich an, daß er willig und mit Wissen den Geist aufgegeben habe im Glauben und Hoffnung auf den zukünftigen HErrn Christum. Er hat die Beine von der Erde auf das Bette gelegt, und gesagt: Ich fahre dahin im Namen Christi.

\chapter{Das fünfzigste Capitel}

\subsection[Von dem Leichengepränge und Verehrigung Jakobs; von der Furcht der …]{Von dem Leichengepränge und Verehrigung Jakobs; von der Furcht der Brüder Josephs, wie sie Joseph suchen zu versöhnen; wie sich Joseph gegen sie verhält; und wie endlich Joseph verlangt in Canaan begraben zu werden}

\subsection{I}

\begin{quote}
\textbf{V. 1—3. Da fiel Joseph auf seines Vaters Angesicht, und weinete über ihm, und küssete ihn. Und Joseph befahl seinen Knechten, den Aerzten, daß sie seinen Vater salbten. Und die Aerzte salbten Israel, bis daß vierzig Tage um waren; denn so lange währen die Salbetage. Und die Egypter beweinten ihn siebenzig Tage.}
\end{quote}

1. In diesem letzten Capitel wird beschrieben das Leichengepränge Jakobs. Und solche Ehrerbietung Josephs, so er seinem Vater erzeigt, kommt her aus dem Glauben und Hoffnung der Auferstehung; denn er hält es dafür, daß sein Vater lebe, sonst würde er ihn weder salben noch küssen. Es ist aber eine sehr schöne Religion, Gebrauch und Ceremonie gewesen, daß sie ihn vierzig Tage gesalbt und dreyßig Tage beweint haben. Denn wenn die beiden Zahlen zusammen kommen, so werden siebenzig Tage daraus.

2. Es ist wahrlich eine große Pracht gewesen. GOtt hat Jakob mit größerer Herrlichkeit begraben denn seinen Sohn Christum Jesum. Und Mose hat nicht vergeblich die Tage, darin man Jakob gesalbt und beweint, so fleißig gezählet, sondern hat also mit sonderlichem Fleiß die sehr ehrliche Leichenbestattung Jakobs rühmen wollen.

3. Es ist in allewege ein groß Ding, daß dieser Fremdling bei dem König Pharao und den Fürsten Egyptens in solchen Ehren gehalten worden und in so großen Gnaden gewesen ist, daß sie ihn so lange Zeit beweinten, und den letzten Dienst oder Wohlthat, so die Gottseligen einander zu erzeigen pflegen, dem Jakob, da er verstorben war, so freiwillig und mit solcher großen Ehrerbietung erzeigt und bewiesen haben.

4. Darnach kann man auch an diesem Salben sehen, daß die Egypter viel herrliche und treffliche Würze gehabt haben, als, Myrrhen, Balsam und Kasia, welche den Leib siebenzig Tage haben erhalten können. Denn diese Kraft schreibt man den Myrrhen zu, daß sie wehren und vertreiben sollen, daß der Leib nicht verfaule und stinke. Da Augustus nach Egypten gekommen ist, haben sie ihm den Leib Alexanders des Großen und des Ptolemäus gezeigt, welche viel Jahre lang in Gräbern verwahret waren, darüber er sich sehr verwundert hat, daß sie also haben verwahret werden können, daß noch die Haut und alle Glieder an ihnen gewesen. Eben dasselbe sagen sie auch vom Leichnam des Kaisers Titus. Derhalben ist es sehr köstliche Würze gewesen, damit sie die Leichname also zugerichtet, daß sie in so vielen Jahren nicht verfault oder verwest sind, so es doch ein sehr heiß Land ist, welche Hitze die Körper sehr weich macht.

5. Wir pflegen im Winter anstatt der Myrrhen mit dem Frost das Fleisch zu verwahren, daß es nicht faul werde, vornehmlich an den Orten, so gegen Norden gelegen sind, als, in Dänemark und an andern Orten mehr. Das ist unsere deutsche Myrrhe. Aber daß sie in Egypten in der heißen Sonne die Leichname haben bewahren können, daß sie nicht faul und stinkend geworden sind, ist sehr wunderlich. Es hat GOtt damit wollen anzeigen, daß die Todten in Christo wahrhaftig mit Myrrhen gesalbt sind. Christus ist unsere Myrrhe, wie ihm die Weisen auch Myrrhen geopfert haben, Matth. 2, 11. Denn so wir an ihn glauben, werden wir mit Myrrhen gesalbt, daß wir nicht verfaulen, sondern erhalten und bewahrt werden auf die zukünftige Auferstehung.

\begin{quote}
\textbf{V. 4—6. Da nun die Leichentage aus waren, redete Joseph mit Pharao Gesinde, und sprach: Habe ich Gnade vor euch funden, so redet mit Pharao, und sprechet: Mein Vater hat einen Eid von mir genommen und gesagt: Siehe, ich sterbe, begrabe mich in meinem Grabe, das ich mir im Lande Canaan gegraben habe. So will ich nun hinaufziehen, und meinen Vater begraben, und wiederkommen. Pharao sprach: Gehe hinauf, und begrabe deinen Vater, wie du ihm geschworen hast.}
\end{quote}

\quellseite{532}

6. Joseph bittet um Erlaubniß, daß er seinen Vater mit Pharaos Willen begraben möge, nicht allein, daß er damit hat anzeigen wollen seine Ehrerbietung gegen die königliche Majestät, welches denn Gottes Gebot also fordert und haben will, 1 Petr. 2, 17.: „Ehret den König“; sondern auch um der gemeinen Aemter willen, so er zu verwalten gehabt und die am Regiment desselben Königreichs nöthig gewesen sind, welche er ohne Pharaos Willen mit Ehren nicht hat können verlassen. Darnach hat er sich auch dessen besorgt, daß vielleicht die Egypter denken möchten, daß er damit umginge, daß er wollte davon fliehen und von ihnen also hinwegziehen, oder aber, daß sie dächten, Joseph verachtete sie darum, daß er es dafür hielt, es wäre besser, daß sein Vater im Lande Canaan begraben würde, denn in dem schönen herrlichen Königreich Egypten. Und haben auch zwar also denken mögen: Wir übertreffen die andern Völker weit über die ganze Welt an Adel, an Gewalt, dazu auch an Geld und Gut, warum sucht denn nun Joseph seinem Vater noch ein ander Grab denn in Egyptenland? Um der Ursache willen führt er den Eid an, den er seinem Vater gethan hatte, und verspricht ihnen, daß er wiederkommen wolle, ja, daß er nicht wollte einen Fuß aus dem Wege geben, wo er durch den Eid dazu nicht wäre genöthigt worden und wo er nicht vom König Erlaubniß erlangt hätte; als wollte er sagen: Ich bin diesem Lande nicht feind, begehre auch nicht daraus zu entfliehen, sondern ich muß meinem Vater Glauben halten. Derselben ist nun der König durch diese Ursache bewogen worden, daß er ihn mit Gnaden hat ziehen lassen, und hat ihm etliche Gefährten mitgegeben, die mit ihm gezogen sind, wie im Text nun weiter folgt.

\begin{quote}
\textbf{V. 7—14. Also zog Joseph hinauf, seinen Vater zu begraben. Und es zogen mit ihm alle Knechte Pharao, die Aeltesten seines Hauses, und alle Aeltesten des Landes Egypten; dazu das ganze Gesinde Josephs, und seine Brüder, und das Gesinde seines Vaters. Allein ihre Kinder, Schafe und Ochsen ließen sie im Lande Gosen. Und zogen auch mit ihm hinauf Wagen und Reiter, und waren ein fast großes Heer. Da sie nun an die Tenne Atad kamen, die jenseit des Jordans liegt, da hielten sie eine sehr große und bittere Klage; und er trug über seinem Vater Leide sieben Tage. Und da die Leute im Lande, die Canaaniter, die Klage bei der Tenne Atad sahen, sprachen sie: Die Egypter halten da große Klage. Daher heißt man den Ort der Egypter Klage, welcher liegt jenseit des Jordans. Und seine Kinder thaten, wie er ihnen befohlen hatte; und führten ihn ins Land Canaan, und begruben ihn in der zwiefachen Höhle des Ackers, die Abraham erkauft hatte mit dem Acker, zum Erbbegräbniß, von Ephron, dem Hethiter, gegen Mamre. Als sie ihn nun begraben hatten, zog Joseph wieder in Egypten mit seinen Brüdern, und mit allen, die mit ihm hinauf gezogen waren, seinen Vater zu begraben.}
\end{quote}

7. Pharao erlaubt Joseph nicht allein, daß er ins Land Canaan ziehe, sondern es ziehen mit ihm das ganze Haus des Königs, alle Aeltesten des Landes Egypten, seine Brüder mit ihren Kindern und Gesinde, und ließen allein die kleinen Kinder und das Vieh in Egypten bleiben; welches wahrlich ein sehr schön Gepränge und Leichenbestattung gewesen ist. Denn Joseph hat seinen Vater ehrlich begraben wollen, auf daß er seiner Zusage genug thun möchte, und daß über ganz Egypten und Canaan lautbar würde, daß Jacob gestorben wäre. Darum hat er noch über die siebenzig Tage, darin er Leid getragen und den Vater zuvor gesalbt und beweint hatte, noch sieben Tage im Lande Canaan verzogen, nicht aus Fürstlichkeit, sondern deswegen, daß er auch dieselben Tage über hat wollen Leid tragen.

8. Das Wort atad verstehen Etliche als ein Appellativum, das ist, ein Gattungswort, und deuten es Stechdornstrauch, oder eine andere Art von Dornen, damit derselbe Ort verzäunt gewesen; wie bei uns die Hecken oder Gärten auch mit Dornen umher verzäunt werden, mit Schlehenbäumen oder Hainbutten. Es ist ein Platz oder Tenne gewesen, mit Dornen verzäunt, darauf man das Getreide hat zu dreschen pflegen. Andere machen aus diesem Wort den Eigennamen eines Mannes, wie derselbigen auch in der deutschen Sprache viel sind: Dorn, Birnstiel etc.

9. Aber hier kommt eine Frage, die mich sehr bewegt und bekümmert: Wie doch Mose sagen konnte, daß Joseph aus Egypten über den Jordan gezogen sei. Denn es ist der gerade Weg aus Egypten durch die Wüste Etham, und ist eine solche Reise nach Hebron, die man innerhalb vierzig Tagen thun kann: der Jordan aber liegt gegen Aufgang der Sonne. Ich löse aber diese Frage also auf: Jakob ist nicht auf demselben Wege geführt worden über den Jordan, welchen die Kinder Israel gezogen sind; denn dieselbigen sind vom Morgen oder Aufgang der Sonne her über den Jordan gezogen. Das hebräische Wort begreift beiderlei in sich und heißt, diesseits und jenseits des Jordans. Hebron ist Egyptenland etwas näher oder aber je so nahe, als der Jordan ist; Joseph aber hat seinen Vater nicht über den Jordan geführt; denn er ist dahin nicht gekommen, dieweil Atad neben dem Jordan gelegen ist.

\subsection{II}

\begin{quote}
\textbf{V. 15. Die Brüder aber Josephs fürchteten sich, da ihr Vater gestorben war, und sprachen: Joseph wird uns gram sein, und vergelten alle Bosheit, die wir an ihm gethan haben.}
\end{quote}

10. Hier siehst du erstlich, welch ein erschrecklich Unglück die Sünde und ein böses Gewissen sei: daß ist beinahe eine solche Wunde, die nicht zu heilen ist. Die Söhne Jakobs hatten gehört, wie sie ihr Vater gesegnet; und wiewohl den zwei, Simeon und Levi, solche Strafen verkündigt worden, die traurig genug sind, war ihnen doch ihre Sünde und Missethat vergeben. Denn der Vater hat sie ja gesegnet, daß sie wachsen und gemehret werden sollten. Die andern Stämme aber waren alle gesegnet und reich gemacht worden, vornehmlich Juda, und haben noch bei ihres Vaters Leben mit Joseph siebenzehn Jahre in Egypten gelebt, und haben viel große Wohlthaten von ihm empfangen, und seine Gnade und Barmherzigkeit, als ihres Vaters, ohne Unterlaß mit der That erfahren. Sie hatten auch an ihm Liebe gesehen und empfunden die brüderliche Liebe, daß er seine Brüder herzlich lieb hatte und sich ganz freundlich gegen sie erzeigte. Wie er sie denn nicht allein mit sehr freundlichen Worten angeredet, sondern auch von sich hatte ziehen lassen, daß sie wiederum zu ihrem Vater kämen, da er sie mit den größten Geschenken beehret hatte.

11. Derhalben sieht man ja an Joseph Anderes nichts denn lauter barmherzige Barmherzigkeit, wenn man also reden möchte, inwendig und auswendig. Sein Herz ist gar brünstig gewesen und gleichsam entbrannt von väterlicher Liebe; seine Hände waren voll großer Wohlthaten gegen seine Brüder: und gleichwohl ist ihr Herz noch immer unruhig, und werden bei sich selbst des bösen Gewissens halben, daß sie solche schwere Sünde begangen hätten, und vom Stachel des Todes also geplagt und gemartert, daß sie selbst dem Menschen nicht vertrauen, noch sich Gutes zu ihm versehen können, der ihnen doch bis daher schon viel Gutes gethan und sich nicht anders denn nur auf das allerfreundlichste erzeigt hatte.

12. Ist denn nun die Sünde nicht ein erschrecklich greulich Ding? Sie wird wohl leichtlich begangen und vollbracht, sonderlich außerhalb der Anfechtung, da die Leute ohne alle Scheu sicher und frei dahin sündigen, denken, es hat keine Noth: aber wenn die Sünde offenbar wird und das Beißen und der Stachel des bösen Gewissens hernach folgt, alsdann ist keine Vergebung oder Trost stark genug, solch Beißen und denselben Stachel zu vertreiben oder demselben zu widerstehen.

13. Derhalben ist der Papst mit seinem Anhang billig wie der Teufel selbst zu verfluchen, welcher die Welt mit vielen unzähligen Sünden erfüllt hat, und gelehrt, daß man zweifeln solle an der Barmherzigkeit und Gnade Christi. Denn es geht nicht ohne große Beschwerung zu, daß ein armes Herz nach Gottes Gnade und daß ihm seine Sünden mögen verziehen werden, verlangt; vor welchem Gott es von Natur flieht und ihm feind ist; wenn es fühlt, daß er durch viel große und greuliche Sünden erzürnet worden ist, da läßt man sich schwerlich trösten. Die Vergebung der Sünde haftet im Herzen fast kümmerlich, wiewohl sie mit vielen Versicherungen und Zeichen der Gnade Gottes bekräftigt ist. Daraus abzunehmen ist, welch ein scharfes Gift die Sünde in der Natur des Menschen sei. Dasselbe kann man aber nicht bald sehen, wenn nur einfach das Gesetz dem Menschen vorgehalten wird, es sei denn, daß es zu seiner Kraft komme und daß die Sünde anfängt lebendig zu werden; da findet es sich denn, wie ein greulich vergiftetes höllisches Uebel und Unglück es ist, auch so, daß der Tod und das Blut des Sohnes Gottes die Sünde kaum bedecken und tilgen kann.

\quellseite{533}

14. Siehe doch, was Petrus und die andern Jünger gethan haben, wie hart, wie ungläubig und unverständig ihr Herz ist, die Auferstehung Christi zu glauben. Also hängt die Sünde diesen Brüdern Josephs auch auf das härteste an, daß sie bis daher nicht hat können gestillt oder weggenommen werden durch so große Wohlthaten, und daß Joseph so freundlich mit ihnen geredet, und Befehle gegeben an ihren Vater, die voll Freundlichkeit und guten geneigten Willens gewesen, und endlich, da sie so viel Jahre auf das allerfreundlichste mit einander umgegangen waren. Denn darum sind sie jetzt so erschrocken, da ihr Vater gestorben war, und fürchten sich, als die sich schuldig wissen, daß ihnen Joseph gram sein, und die Bosheit, so sie an ihm begangen, vergelten und Böses mit Bösem bezahlen möchte. Wer denkt oder redet das nun anders, denn allein das Gift der Sünde selbst, welche nun durch das Gesetz aufgeweckt ist? Dieselbe können sie nicht bezwingen oder heilen, da sie doch allebereit nun fast siebenzehn Jahre daher etwas geheilt und gelindert ist. Wie gar leichtlich und gelinde ist es aber zugegangen, da Simeon und Levi, und die andern Brüder gesündigt haben! aber wie schwer sie nun wiederum geheilt und zufrieden gestellt werden können, das bezeugt die Schrift an diesem Orte. Und daher kommt es auch, daß Viele, so das Wort der Gnaden nicht hören, dahin getrieben werden, daß sie gar verzweifeln und in das Wasser springen, oder aber sich selbst an einem Strick hängen und also dahin sterben. Denn sie können oder mögen die Kraft der Sünde, so nun geoffenbart und lebendig in ihnen geworden ist, nicht ertragen. Wenn die Sünde noch vor der Thüre schläft, wie zu Kain 1 Mos. 4, 7. gesagt wird, da achtet man ihrer nicht und wird immer noch mit andern augenblicklichen Sünden gehäuft; man fällt in eine Sünde über die andere. Aber wenn sie aufgeweckt und lebendig gemacht wird, alsdann muß das Blut des Sohnes Gottes darauf gewendet werden, damit sie abgeschafft und weggenommen werden möge. Ja, eine solche große theure Arznei und Hülfe gehört dazu, nämlich, die Gottheit, so Mensch geworden ist, und das Blut des Sohnes Gottes selbst.

15. Dieweil nun so gar viel Unglück und Jammer der Erbsünde ist, darin wir empfangen und geboren sind, so sind auch die schändlichen Päpisten so viel mehr zu verfluchen, die unsern Schaden immer noch mit den wirklichen Sünden und Gesetzen des Papstes mehren und schwerer machen.

16. Wer sich aber vor Sünden hüten kann, der gedenke, daß er es thue; und die, so gefallen sind, sollen lernen wiederum aufstehen und fest glauben. Und in solchem Kämpfen lernen wir verstehen, was Glaube sei; welches die Sophisten nicht wissen, und denken einfach, der Glaube sei nur ein solcher Wahn, so aus menschlichem Gutdünken gefaßt ist. Denn sie wissen vom Gesetz und von der Sünde nichts; und darum verstehen sie auch das Evangelium und Gottes Gnade nicht, meinen, es seien geringe Dinge, denen man leichtlich mit guten Werken und natürlichen Kräften rathen könne. Wenn sie es aber in Anfechtungen gelernt und erfahren hätten, wie es so schwer sei, daß man sich aufrichte und erwecke im Glauben an den Heiland Christum wider die Sünde und Tod, würden sie viel anders lehren. Aber dieweil sie davon nichts wissen, liegen sie und schnarchen, und lassen sich träumen, die Sünde sei nichts Anderes denn nur Lust des Fleisches. Aber vom Zweifel und Unglauben, und daß man Gott feind ist und vor ihm flieht, davon wissen sie nichts. Darum sind sie nicht werth, daß sie sollen Theologen heißen; denn sie wissen von den vornehmsten Hauptstücken der christlichen Lehre gar nichts.

17. Derhalben sollen wir lernen, daß die Sünde ein erschrecklich Unglück sei, nicht, wenn sie begangen wird; denn da gefällt sie uns über die Maßen wohl und wir haben dazu große Lust: sondern wenn sie durch das Gesetz aufgeweckt wird. Da sollen wir lernen, daß es die Hölle selbst ist, und ein viel gewaltiger Ding, denn Himmel und Erde sein kann, auch also, daß man die Gnade nicht ohne große Mühe und Arbeit ergreifen oder fassen kann, daß ein Herz, so vom Gesetze und der Sünde beschweret ist, bei sich also schließen könne: „Wiewohl ich schwer gesündigt habe, so hat doch Gott seinen Sohn in die Welt gesandt, nicht, daß er die Welt richte, sondern daß die Welt durch ihn selig werde, Joh. 3, 17. Wenn dieser Trost nicht da ist, da ist keine andere Hülfe und Rath wider die Sünde und ihren Stachel, wenn du auch schon tausend Kappen anzögest und bis an der Welt Ende wallen gingest.“

18. Wie vorzeiten ich auch, da ich noch ein Mönch war, die Hoffnung hatte, ich würde mein Gewissen können zufrieden stellen mit Fasten, Beten, und vielem Wachen, damit ich meinen Leib jämmerlich plagte und marterte. Aber je sauer ich mir es werden ließ, je weniger Ruhe und Frieden ich fühlte; denn das rechte Licht war von meinen Augen weggethan: ich war ohne Glauben, und rief die verstorbenen Heiligen und die Jungfrau Maria an, opferte ihnen Messen, bis wir nun jetzt wiederum aus großer Gnade Gottes aus solcher Finsterniß hervor gekommen sind, und Christum erkennen, welchen das scheußliche Ungeheuer, der Papst und seine Sophisten, gar begraben haben. Ich weiß nun, Gott Lob! daß er nicht allein nach der Historie gestorben ist vor fünfzehn hundert Jahren, sondern daß sein Tod währet vom Anfang bis an der Welt Ende, und daß er zu Hülfe kommt allen Heiligen durch die ganze Zeit der Welt, als, Adam, Eva und den Andern etc. Christus ist mir jetzt noch so neu, als ob er diese Stunde sein Blut vergossen hätte. Davon habe ich im Papstthum gar nichts gewußt; denn da meinte ich, der Tod Christi gehöre allein auf die Erbsünde, oder sonst die vergangenen Sünden, in meinen kindlichen Tagen begangen; aber mich ginge er jetzt nichts an.

19. Jetzt aber lasset uns nun lernen, und dasselbe fest behalten, daß die Gebeine der Sünder keine Ruhe noch Frieden haben können, denn allein durch den einigen Glauben, welcher die Schlange, so von Gott erhöhet ist, anschaut, Joh. 3, 14. 15.: „Wie Moses in der Wüste eine Schlange erhöhet hat, also muß des Menschen Sohn erhöhet werden, auf daß alle, die an ihn glauben, nicht verloren werden, sondern das ewige Leben haben.“ Der Glaube thut es, der macht, daß die Gewissen stille werden und Frieden haben, bis an der Welt Ende. Nicht daß der Sünder ganz und gar nicht sollte gedacht werden; denn Gott will haben, daß wir der unermäßlichen Wohlthaten gedenken und ihm dafür danken sollen, dadurch er uns aus dem Gefängnisse des Teufels geführt und erlöset hat; aber dasselbe also, daß das Gewissen nicht gemartert werde, sondern wisse, daß es Vergebung der Sünden habe und das ewige Leben durch Christum.

20. Nun folget die Rede dieser Brüder, die sie dahin richten, daß sie Joseph versöhnen mögen; und Stücke sagen, daß Ephraim ihr Gesandter und Redner gewesen sei.

\begin{quote}
\textbf{V. 16—18. Darum ließen sie ihm sagen: Dein Vater befahl vor seinem Tode und sprach: Also sollt ihr Joseph sagen: Lieber, vergib deinen Brüdern die Missethat und ihre Sünde, daß sie so übel an dir gethan haben. Lieber, so vergib nun die Missethat uns, den Dienern des Gottes deines Vaters. Aber Joseph weinete, da sie solches mit ihm redeten. Und seine Brüder gingen hin, und fielen vor ihm nieder, und sprachen: Siehe, wir sind deine Knechte.}
\end{quote}

21. Sie brauchen in dieser Rede eine sonderliche Rhetorik, nicht allein, daß sie wiederholen ihres Vaters Befehl, damit sie ihn wie durch ein Gesetz zwingen wollen, daß er ihnen ihre Sünde verzeihen soll; sondern daß sie sich auch damit schmücken und loben mögen, da sie sich Diener nennen des Gottes seines Vaters; als wollten sie sagen: Wir sind auch Gottes Diener, dienen demselben Gott, welchem du und dein Vater auch gedienet habet: darum vergib uns die Sünde um des gemeinsamen Gottes willen, den wir alle fürchten, ehren und vor Augen haben. Sie zweifeln an Gott nicht, daß er ihnen alle ihre Sünden vergeben habe; aber darüber sind sie ungewiß, wie ihr Bruder gegen sie gesinnet sei, was er für ein Herz zu ihnen trage, der doch nun siebenzehn Jahre lang mit vielen Wohlthaten bezeugt hatte, daß er schon mit ihnen versöhnt wäre und ihnen mit höchster Barmherzigkeit Alles verziehen hätte. Aber sie fühlen noch immer das Gegentheil, und halten es stracks dafür, daß er es ihnen noch nicht vergeben habe, und können mit so vielen Anzeichen des guten Willens und brüderlicher Liebe noch nicht zufrieden sein. Das ist nun, wie ich gesagt habe, die Art und Natur der Sünde, welche das Gewissen auch vor den Menschen greulich martret; aber vor GOtt beschuldigt und erschreckt sie dasselbe noch viel härter.

\quellseite{534}

\subsection{III}

\begin{quote}
\textbf{V. 19—23. Joseph sprach zu ihnen: Fürchtet euch nicht, denn ich bin unter GOtt. Ihr gedachtet es böse mit mir zu machen; aber GOtt gedachte es gut zu machen, daß er thäte, wie es jetzt am Tage ist, zu erhalten viel Volks. So fürchtet euch nun nicht; ich will euch versorgen und eure Kinder. Und er tröstete sie, und redete freundlich mit ihnen. Also redete Joseph in Egypten mit seines Vaters Hause, und lebte hundert und zehn Jahr. Und sahe Ephraims Kinder, bis ins dritte Glied. Desgleichen gleichen die Kinder Machir, Manasses Sohnes, zeugten auch Kinder auf Josephs Schooß.}
\end{quote}

22. Joseph antwortet seinen Brüdern sehr freundlich und friedlich: Ihr, sagt er, ziehet GOtt an, für welches Diener ihr euch ausgebt; das ich denn gern höre und glaube. Und ihr sollt eben dasselbe von mir auch halten; denn ihr habt von unserm Vater dasselbe auch gehört. Darum so erkenne ich, und nehme diese Gesandschaft oder Werbung von euch an, ganz freundlich und demüthig an, und erhebe mich nicht über GOtt, und will nicht haben, daß ihr euch zu mir etwas Böses versehen sollt, oder den Argwohn auf mich werfen, als ob ich etwas wider GOtt thun wollte. Haltet ihr mich für einen Diener GOttes, so denket nicht, daß ich etwas wider ihn thun wollte. So tut euch GOtt eure Sünde vergeben, habt ihr ein gut Gewissen, und seid solcher Verheißung und Vergebung, so euch von GOtt widerfahren, gewiß, was dürft ihr denn an mir zweifeln? Denn ich bin ja nicht über GOtt, sondern unter GOtt. Bleibet ihr auch mit mir unter GOtt, und lasset uns hoffen, daß er uns gnädig sei und daß wir mit ihm ausgesöhnt sind.

23. Diese Worte sind voll herzlicher Neigung und Liebe, so Joseph zu seinen Brüdern getragen; und dieweil er durch solche Freundlichkeit und Barmherzigkeit überwunden wird, hat er sich nicht enthalten können, er hat müssen weinen. Und jetzt, dieweil er hört, daß seine Brüder daran noch zweifeln, ob er ihnen vergeben habe, nachdem sie von GOtt und auch von Jakob, ihrem Vater, Vergebung erlangt hatten, sagt er zu ihnen also: So euch GOtt eure Sünde vergeben hat, warum sollte denn ich euch nicht auch alle Bosheit vergeben, so ihr an mir gethan habt? Und setzt einen sehr ernsten Spruch hinzu, und sagt: Es ist wohl wahr, „ihr gedachtet es böse mit mir zu machen“; aber GOtt ist in seinem Rathe wunderbarlich, der eure bösen Gedanken zu unser aller Nutzen und Guten gewandt hat.

24. Derhalben will Joseph nicht, daß sie die Sünde und ihre Pläne, ihn zu tödten, verleugnen, oder derselben vergessen sollen; sondern sehet, sagt er, welch ein wunderlich Werk GOttes das ist, der dies alles zur höchsten Wohlfahrt gewandt hat, auf daß ich vielen Völkern dienen und helfen möchte, welche Hungers gestorben wären, und GOtt auch nie erkannt, sein Wort auch nicht gehört hätten, wo ich nicht von euch nach Egypten wäre verkauft worden. Also hat GOtt aus dem, das so sehr böse war, etwas, das sehr gut ist, gemacht. Darum will ich euch gern verziehen und vergeben, dieweil ich sehe, daß GOtt durch euren sehr bösen Rath mir und unzähligen andern Leuten mehr so viel Gutes gethan hat. Und das ist es, das Paulus sagt Röm. 8, 28.: „Wir wissen, daß denen, die GOtt lieben, alle Dinge zum Besten dienen, die nach dem Vorsatz berufen sind“ etc. Und St. Augustin sagt auch an einem Orte: GOtt ist so gütig, daß er nichts Böses geschehen läßt, es sei denn, daß er daraus ein groß Gut machen könnte.

25. Daraus folgt aber das nicht, was die gottlosen Leute daraus schließen wollen und St. Paulus widerlegt hat, Röm. 3, 7. 8.: „Denn so die Wahrheit GOttes durch meine Lüge herrlicher wird zu seinem Preis, warum sollte ich denn noch als ein Sünder gerichtet werden? Und nicht vielmehr also thun, wie mir gelästert werden, und wie etliche sprechen, daß wir sagen sollen: Lasset uns Uebel thun, auf daß Gutes daraus komme?“ Denn Paulus setzt selbst hinzu und sagt: „Welcher Verdammniß ist ganz recht.“ „Es ist zwar das Gesetz neben eingekommen, auf daß die Sünde mächtiger würde“, das ist, durch das Gesetz wird angezeigt, wie groß die Sünde sei, und wird dadurch nicht allein nicht geheilt noch weggenommen, sondern wird größer und mächtiger; aber diese Größe macht, daß auch die Gnade mächtiger wird. Derhalben soll die Größe der Sünde niemand zur Verzweiflung treiben, sondern man soll so viel mehr rühmen und groß machen die Größe der Gnade, welche die Sünden verschlingt, wie groß sie auch sein mögen.

26. Das ist aber zumal ein ungeschicktes und gottloses Argument, wenn man also schließen will: Wenn durch unsere Sünde GOttes Gnade und Barmherzigkeit gepriesen wird, so sollen wir nur frei sündigen. Denn GOtt will das nicht haben oder zugeben; er schreibt dir auch keine Regel vor, nach welcher du frei sündigen mögest. Er sagt nicht, daß er um der Sünde willen jemand wolle Gutes thun. Es ist ein Anderes, Gutes thun um der Sünde willen, und ein Anderes, einem um der Sünde willen zu Hülfe kommen, daß er derselben möge los werden.

27. Denn GOtt ist gänzlich der Sünde feind und haßt dieselbe sehr heftig; und legt uns gleichwohl die Verheißung vor: Wer gefallen ist, soll nicht verzweifeln. Er heißt uns nicht Sünde thun, sondern verbietet solches, und will haben, wir sollen nicht sündigen. Wenn aber die Sünde begangen ist, und das Gesetz beschuldigt und erschreckt das Gewissen, „auf daß die Sünde überaus sündig werde durch das Gebot“, Röm. 7, 13., wie daselbst hier an den Brüdern Josephs zu sehen ist: da will traun GOtt nicht haben, daß der Tod herrschen soll; wie er das durch den Propheten bezeugt, Hos. 18, 23., da er sagt: „Meinest du, daß ich Gefallen habe am Tode des Gottlosen, spricht der HErr HErr, und nicht vielmehr, daß er sich bekehre von seinem Wesen und lebe?“ Denn wie er der Sünde feind ist, also will er auch nicht haben, daß der, so gefallen ist, in der Sünde bleiben soll, wenn sie lebendig geworden ist, und durch das Gesetz das Gewissen verwundet und peinigt; sondern er gibt alsdann eine Verheißung und Hülfe, dadurch ein solch Herz, so verwundet ist, erhalten werde, daß es nicht verzweifele. Ein Pharisäer oder Werkheiliger soll dadurch nicht selig werden, daß er so vermessen ist und sich auf seine Frömmigkeit verläßt; und David oder der Schächer am Kreuz soll auch nicht um der Sünde willen verzweifeln, als müßte er darum verloren werden.

28. Sondern man soll auf der rechten Mittelstraße bleiben und von Sünden sich hüten. Denn wiewohl GOtt Gnade und Vergebung verheißen hat, wie Augustin sagt, hat er doch das nicht verheißet, daß du eben so gewiß nach dem Fall werdest wieder kommen; wie Saul und Judas nicht wieder gekommen. Es steht nicht in unsrer Macht, die Gnade zu ergreifen, und du weißt nicht, ob du auch die Vergebung, so dir angeboten wird, annehmen könnest. Darum soll man GOtt fürchten, welcher beide der Barmherzigkeit und auch der Verzweiflung Feind ist. Wie ich keinen Gefallen habe, sagt er, am Sünder; also habe ich auch am Tode des Sünders keinen Gefallen: sondern wo du durch des Teufels Bosheit dich hast zu sündigen treiben lassen, und das Gesetz dich nun derhalben verdammt, und der Teufel dich mit seinen feurigen Pfeilen schießt, und damit umgeht, daß er dich möge in die Hölle stürzen; so will ich doch nicht, daß du sterben oder verzweifeln sollst; sondern du sollst straks zu dem Heiland Christo fliehen, welcher nicht Lust hat am Tode des Sünders, Hes. 18, 32., wie er auch am Sünder keinen Gefallen hat.

29. Auf diese Weise macht GOtt aus dem Bösen etwas Gutes: nicht daß er haben wollte, daß das Böse geschehen sollte; sondern seine Güte ist so groß, auch in unsrer Bosheit, daß er es nicht lassen kann, er muß die Sünde vergeben, wenn der Sünder reu’t und um Gnade bittet. Wenn das geschieht, so soll es vergeben sein.

30. Dieser große Trost soll den Herzen der Gottseeligen fleißig eingeprägt werden, den Glauben damit zu erwecken und zu stärken. Aber wir müssen wohl zusehen, daß wir derselben nicht mißbrauchen, wie Erasmus vermahnt am 5. Cap. V. 5. 7.: „Sei nicht so sicher, obwohl deine Sünde noch nicht gestraft ist, daß du darum für und für sündigen wolltest. Denke auch nicht: GOtt ist sehr barmherzig, er wird mich nicht strafen, ich sündige, wie viel ich will. Er kann bald also zornig werden, als gnädig er ist, und sein Zorn über die Gottlosen hat sein Aufhören.“ Wir pflegen auf Deutsch zu sagen: Du sollst nicht auf Barmherzigkeit sündigen. Und kurz zuvor, V. 4., sagt Sirach: „Denke nicht: Ich habe wohl mehr gesündiget, und ist mir nichts Böses widerfahren; denn der HErr ist wohl geduldig, aber er wird dich nicht ungestraft lassen.“

31. Derhalben, je größer die Gnade und Barmherzigkeit GOttes ist, je weniger man derselben mißbrauchen soll; sondern wer gefallen ist, der soll zu GOttes Gnade Zuflucht haben. Da wird er alsdann wohl sehen, wie schwer es sei, daß man sich aufrichten und trösten lasse. Wenn man aber den Trost ergriffen hat, so ist schon die Sünde geheilt. Wirst du aber dich auf GOttes Güte und Gnade verlassen, wissentlich und vorsätzlich seine Gebote zu übertreten dich unterstehen, ist große Gefahr dabei, daß dich die Sünde gar unterdrücken möchte, ehe du die Vergebung in deinem Sohne empfängest; wie Judas, Saul, Ahithophel, Absalom dieselbe nicht haben begehren können, welche zuvor so verstockt gewesen, daß sie auf GOttes Gnade und Barmherzigkeit gesündigt haben.

\quellseite{535}

32. Derhalben, das Sirach sagt: „Sei nicht so sicher, obwohl deine Sünde noch nicht gestraft ist“ \&c., das gehört auf die zukünftige Sünde, daß sich nicht Einer unterstehe, den Irrthum der Sophisten daraus zu beweisen, die da lehren, man solle nicht Vergebung der Sünden glauben, und daß sie ungewiß sei; sondern das ist der rechte Verstand der Worte Sirachs: Du sollst forthin nicht sündigen. Man soll sich nicht also verlassen auf GOttes Gnade und Barmherzigkeit, daß wir derselben mißbrauchen wollten; denn solche Leute sind, die da sagen: „Lasset uns Uebel thun, auf daß Gutes daraus komme“, Röm. 3, 8.

33. Wieweit GOtt die Sünde verboten hat, will er haben, daß du dieselbe fliehen und dich davor hüten sollst. Wenn sie aber nun begangen ist, und dich beißt und zur Verzweiflung treibt, wie Adam im Paradiese vor GOtt flieht und sich vor ihm verbirgt, nachdem er erkannte, daß er nacket war und nicht bestehen konnte vor der Stimme GOttes, da war die Sünde in ihrer rechten Wirkung und in erschrecklichem Gericht der Anklage und Verdammniß, und Adam war da in Wahrheit recht todt, nach dem Urtheil, so GOtt zuvor gesprochen hatte, 1 Mos. 2, 17.: „Welches Tages du davon issest, wirst du des Todes sterben“: alsdann sollst du zum Wort der Gnade eilen, wie Adam auch wiederum lebendig geworden ist, da er die Verheißung vom Samen des Weibes gehört hat, 1 Mos. 3, 15.: „Ich will Feindschaft setzen zwischen dir und dem Weibe, und zwischen deinem Samen und ihrem Samen. Derselbe soll dir den Kopf zertreten“ \&c. Das ist so viel gesagt: Adam und Eva sollen nicht sterben, sondern leben, und sollen einen Samen haben, welcher sie erhalten wird. Da sind sie beide wiederum lebendig und selig geworden.

34. Darum muß man das fleißig merken und wohl behalten, daß GOtt an der Sünde keinen Gefallen hat, und daß er auch nichts thue um der Sünde willen; denn er ist ihr feind und straft sie, wie geschrieben steht: Ich will einem jeglichen geben nach seinen Werken. Aber denen, in welchen die Sünde durch das Gesetz mächtig ist, bietet er Gnade und Seligkeit an. Zuvor, sagt er, habe ich verboten, du solltest nicht vermessen sein: du hast mir aber nicht gehorcht. Nun verbiete ich dir, daß du nicht verzweifeln sollst. Du mußt dich vor beiderlei Sünde hüten, und sollst nicht sagen: Die Sünde ist vergeben und bezahlt, ich will ohne Furcht dahin leben. In solcher Sicherheit sollst du nicht fortfahren; denn es kann wohl geschehen, daß du mit Verzweiflung überfallen werdest, ehe du die Barmherzigkeit ergreifst. Du weißt nicht die Stunde des Todes, weißt auch nicht, ob dich der Teufel nicht werde angreifen, ehe du anfängst Buße zu thun, oder ob du nicht möchtest von ihm übereilt werden, und nicht Raum oder Zeit haben, dich zu bekehren.

35. Denn die Sünde oder Tod ist nicht ein solch Ding, daß man sie so frei unsers Gefallens ablegen und wegwerfen könnte, wie man einen Rock auszieht und ablegt: sondern es dringt durch das Mark des Fleisches und Geistes; wie wir sehen, daß diejenigen, so in Verzweiflung stecken, in den schweren Anfechtungen, so sie haben, von schwerer Bekümmerniß des Herzens und von der Schwindsucht gar verzehret werden. Und David hat diese Kraft und Tyrannei der Sünden greulich gefühlt und empfunden, da er vom Propheten Nathan die harte Strafe gehört, daß er zu ihm sagt, 2 Sam. 12, 7.: „Du bist der Mann“, der die Sünde gethan hat \&c. Da schreit er mit lauter Stimme, V. 13.: „Ich habe gesündigt wider den HErrn“, und steht ihm vor seinen Augen und Herzen nichts Anderes, denn der Tod und ewiges Verderben. Und er hätte auch müssen in demselben Augenblick sterben, wo nicht der Prophet bald zu ihm gesagt hätte: „Du wirst nicht sterben, sondern der HErr hat deine Sünde weggenommen“. 2 Sam. 12, 13. Ein solch kräftig Ding ist es, sage ich, um die Verzweiflung und Sünde, daß es beide Leib und Seele verzehrt.

36. Derhalben siehe, daß du nicht leichtlich sündigest auf den Trost der Gnade GOttes; sondern man soll sich darauf allein verlassen und sich daran fest halten, wenn die Verzweiflung daher geht. Und da sollen die Pfarrherren und Kirchendiener fleißig helfen, und da soll man die Gnade und Barmherzigkeit GOttes, die viel mächtiger ist, als die Sünde immer sein kann, groß machen und rühmen; wie Nathan den David mit diesem Trost von dem Tode auferweckt hat, da er zu ihm sagt: „Der HErr hat deine Sünde weggenommen.“ Wie thust du? Willst du darum sterben in deinen Sünden? Er konnte es aber gleichwohl lange nicht überwinden; denn er hat darnach noch eine lange Zeit mit dem Stachel und Beißen der Sünde gekämpft, wiewohl sie ihm vergeben und verziehen war. Denn siehe, wie jämmerlich und elend er in den Psalmen darüber klagt, so wirst du sehen, was die Sünde für Schaden thut. Es ist eine unerträgliche Last und ein recht teuflisch Gift.

37. Derhalben sei nicht so sicher, Uebel zu thun, daß Gutes daraus komme. Denn wiewohl GOtt diese Sünde der Brüder Josephs dazu gebraucht, daß dadurch vielen Leuten ist geholfen worden, und sonst viel mehr Exempel von GOttes Gnade und Barmherzigkeit hin und wieder in der Schrift stehen, ist doch Gefahr dabei, daß diejenigen, so gar ohne Furcht sind, vom Tode möchten übereilt werden und in die Hölle fahren, ehe sie Zuflucht haben können zu der Barmherzigkeit GOttes.

38. Denn dasselbe bezeuget diese gegenwärtige Historie auch, und zeigt an, wie schwer es sei, daß man nach begangener Sünde wiederum Buße thun möge, und dieser Hülfe glauben, daß GOtt umsonst wolle gnädig sein um seines Sohnes Christi Jesu willen denen, die ihre Sünde bekennen. Und vornehmlich können die nicht ohne große Mühe und Arbeit aus der Verzweiflung kommen, so im Gesetz und Gefängniß des Papstes auferzogen und in den Banden menschlicher Satzungen bestrickt sind. Da wollen wir ja immer etliches Verdienst herbringen, und begehren, daß GOtt unsere Treue und Genugthuung ansehen wolle; welchen Irrthum noch jetzt die Sophisten zu Löwen unverwandt und ganz unchristlich bestätigen wollen.

39. Wir aber haben denselben verworfen und verdammt. Denn die Buße begreift in sich den Schrecken des Gewissens und den Glauben, das ist: sie lehrt, daß Christus umsonst und aus Gnaden dem armen Sünder hilft, mit welchem es jetzt an dem ist, daß er würde verzweifeln müssen, und der sich doch noch durch den Glauben aufrichtet und schreit: GOtt, erbarme dich mein um Christi Jesu willen, an dem ich glaube, der für mich den Tod gelitten hat! Und wo du in dem Glauben also verharrst, wirst du gewißlich selig werden.

40. Du wirst aber noch dabei zugleich fühlen einen erschrecklichen Kampf des Gesetzes, der Natur, der Gewohnheit oder langen Gebrauches, und endlich der ganzen Welt, so diesem Glauben und Vertrauen der Seligkeit widerstrebt. Darum können wir denselben mit unsern Kräften nicht machen; es ist kein selbst erworbener Glaube, sondern wie Paulus sagt: Es ist GOttes Gabe und kommt nicht aus uns selbst.

41. Und darum erzeigt sich Joseph hier so gar freundlich gegen seine Brüder und redet ihnen in das Herz, tröstet sie mit so freundlichen Worten, als er immer kann und mag, nur daß er damit den Zweifel aus ihren Herzen nehmen möge, daß sie sich zu ihm versehen könnten, daß er ihnen gnädig wäre, und daß er ihnen Muth machte, daß sie getrost werden, und es gewißlich dafür halten, daß aller ihrer Sünde und Bosheit bei ihm gar vergessen sei.

42. Es setzt aber Mose noch dazu, daß Joseph von seinen Söhnen Kinder gesehen habe bis ins dritte Glied, und gelebt habe hundert und zehn Jahre: zu der Zeit ist dies Volk sehr gemehret und groß geworden.

\subsection{IV}

\begin{quote}
\textbf{V. 24—26. Und Joseph sprach zu seinen Brüdern: Ich sterbe, und GOtt wird euch heimsuchen, und aus diesem Lande führen in das Land, das er Abraham, Isaak und Jakob geschworen hat. Darum nahm er einen Eid von den Kindern Israel, und sprach: Wenn euch GOtt heimsuchen wird, so führet meine Gebeine von dannen. Also starb Joseph, da er war hundert und zehn Jahr alt. Und sie salbeten ihn, und legten ihn in eine Lade in Egypten.}
\end{quote}

43. Joseph beweiset hier seinen Glauben, da er begehrt, daß er unter diejenigen möge gezählt werden, die mit Christo sollten auferweckt werden. Und ich glaube, daß er mit den andern Heiligen, deren Matth. 27, 52. 53. gedacht wird, wiederum sei lebendig geworden. Er will durchaus im Lande Canaan ruhen; wiewohl an der Stätte und Begräbniß wenig gelegen ist; denn er hätte auch mögen in Egypten begraben werden. Und wäre es Christo nicht schwerer gewesen, die Väter aufzuwecken aus Egypten oder anders woher; sondern seinen Glauben an Christum damit zu bezeugen, gebietet er, daß man seine Gebeine in das Land Canaan führen soll. Dergleichen auch darum, weil er gewußt hat, daß das Volk Israel in dem Lande, so den Vätern verheißen war, gemehret werden sollte.

\quellseite{536}

44. Darum begehrt er, daß sein Begräbniß allen Nachkommen möge vor Augen sein, daß die Kinder und Kindeskinder des Vaters und ihrer Vorfahren dabei gedenken möchten, und nach ihrem Exempel in demselbigen Glauben und Verheißung, darin er mit seinen Vätern entschlafen wäre, auch beharren möchten.

45. Das ist nun der liebe Genesis. Unser HErr GOtt gebe, daß es Andere nach mir besser machen! Ich kann nicht mehr; ich bin schwach. Bittet GOtt für mich, daß er mir ein gutes seliges Stündlein verleihe.

\backmatter
\chapter{I. Hebräisches Wort-Register}

\notensatz{facsimiles/register-536-right}{Hebräisches Wort-Register. Rechte Druckseite von Scan 536.}

\quellseite{537}
\notensatz{facsimiles/register-537-left}{Hebräisches Wort-Register. Linke Druckseite von Scan 537.}

\notensatz{facsimiles/register-537-right}{Hebräisches Wort-Register. Rechte Druckseite von Scan 537.}

\quellseite{538}
\notensatz{facsimiles/register-538-left}{Hebräisches Wort-Register. Linke Druckseite von Scan 538.}

\chapter{II. Spruch-Register}
\notensatz{facsimiles/register-538-right}{Spruch-Register. Rechte Druckseite von Scan 538.}

\quellseite{539}
\notensatz{facsimiles/register-539-left}{Spruch-Register. Linke Druckseite von Scan 539.}

\notensatz{facsimiles/register-539-right}{Spruch-Register. Rechte Druckseite von Scan 539.}

\quellseite{540}
\notensatz{facsimiles/register-540-left}{Spruch-Register. Linke Druckseite von Scan 540.}

\chapter{III. Sach-Register}
\notensatz{facsimiles/register-540-right}{Sach-Register. Rechte Druckseite von Scan 540.}

\quellseite{541}
\notensatz{facsimiles/register-541-left}{Sach-Register. Linke Druckseite von Scan 541.}

\notensatz{facsimiles/register-541-right}{Sach-Register. Rechte Druckseite von Scan 541.}

\quellseite{542}
\notensatz{facsimiles/register-542-left}{Sach-Register. Linke Druckseite von Scan 542.}

\notensatz{facsimiles/register-542-right}{Sach-Register. Rechte Druckseite von Scan 542.}

\quellseite{543}
\notensatz{facsimiles/register-543-left}{Sach-Register. Linke Druckseite von Scan 543.}

\notensatz{facsimiles/register-543-right}{Sach-Register. Rechte Druckseite von Scan 543.}

\quellseite{544}
\notensatz{facsimiles/register-544-left}{Sach-Register. Linke Druckseite von Scan 544.}

\notensatz{facsimiles/register-544-right}{Sach-Register. Rechte Druckseite von Scan 544.}

\quellseite{545}
\notensatz{facsimiles/register-545-left}{Sach-Register. Linke Druckseite von Scan 545.}

\notensatz{facsimiles/register-545-right}{Sach-Register. Rechte Druckseite von Scan 545.}

\quellseite{546}
\notensatz{facsimiles/register-546-left}{Sach-Register. Linke Druckseite von Scan 546.}

\notensatz{facsimiles/register-546-right}{Sach-Register. Rechte Druckseite von Scan 546.}

\quellseite{547}
\notensatz{facsimiles/register-547-left}{Sach-Register. Linke Druckseite von Scan 547.}

\notensatz{facsimiles/register-547-right}{Sach-Register. Rechte Druckseite von Scan 547.}

\quellseite{548}
\notensatz{facsimiles/register-548-left}{Sach-Register. Linke Druckseite von Scan 548.}

\notensatz{facsimiles/register-548-right}{Sach-Register. Rechte Druckseite von Scan 548.}

\quellseite{549}
\notensatz{facsimiles/register-549-left}{Sach-Register. Linke Druckseite von Scan 549.}

\notensatz{facsimiles/register-549-right}{Sach-Register. Rechte Druckseite von Scan 549.}

\quellseite{550}
\notensatz{facsimiles/register-550-left}{Sach-Register. Linke Druckseite von Scan 550.}

\notensatz{facsimiles/register-550-right}{Sach-Register. Rechte Druckseite von Scan 550.}

\quellseite{551}
\notensatz{facsimiles/register-551-left}{Sach-Register. Linke Druckseite von Scan 551.}

\notensatz{facsimiles/register-551-right}{Sach-Register. Rechte Druckseite von Scan 551.}

\quellseite{552}
\notensatz{facsimiles/register-552-left}{Sach-Register. Linke Druckseite von Scan 552.}

\notensatz{facsimiles/register-552-right}{Sach-Register. Rechte Druckseite von Scan 552.}

\quellseite{553}
\notensatz{facsimiles/register-553-left}{Sach-Register. Linke Druckseite von Scan 553.}

\notensatz{facsimiles/register-553-right}{Sach-Register. Rechte Druckseite von Scan 553.}

\quellseite{554}
\notensatz{facsimiles/register-554-left}{Sach-Register. Linke Druckseite von Scan 554.}

\notensatz{facsimiles/register-554-right}{Sach-Register. Rechte Druckseite von Scan 554.}

\quellseite{555}
\notensatz{facsimiles/register-555-left}{Sach-Register. Linke Druckseite von Scan 555.}

\notensatz{facsimiles/register-555-right}{Sach-Register. Rechte Druckseite von Scan 555.}

\quellseite{556}
\notensatz{facsimiles/register-556-left}{Sach-Register. Linke Druckseite von Scan 556.}

\notensatz{facsimiles/register-556-right}{Sach-Register. Rechte Druckseite von Scan 556.}

\quellseite{557}
\notensatz{facsimiles/register-557-left}{Sach-Register. Linke Druckseite von Scan 557.}

\notensatz{facsimiles/register-557-right}{Sach-Register. Rechte Druckseite von Scan 557.}

\quellseite{558}
\notensatz{facsimiles/register-558-left}{Sach-Register. Linke Druckseite von Scan 558.}

\notensatz{facsimiles/register-558-right}{Sach-Register. Rechte Druckseite von Scan 558.}

\quellseite{559}
\notensatz{facsimiles/register-559-left}{Sach-Register. Linke Druckseite von Scan 559.}

\notensatz{facsimiles/register-559-right}{Sach-Register. Rechte Druckseite von Scan 559.}

\quellseite{560}
\notensatz{facsimiles/register-560-left}{Sach-Register. Linke Druckseite von Scan 560.}

\notensatz{facsimiles/register-560-right}{Sach-Register. Rechte Druckseite von Scan 560.}

\quellseite{561}
\notensatz{facsimiles/register-561-left}{Sach-Register. Linke Druckseite von Scan 561.}

\notensatz{facsimiles/register-561-right}{Sach-Register. Rechte Druckseite von Scan 561.}

\quellseite{562}
\notensatz{facsimiles/register-562-left}{Sach-Register. Linke Druckseite von Scan 562.}

\notensatz{facsimiles/register-562-right}{Sach-Register. Rechte Druckseite von Scan 562.}

\quellseite{563}
\notensatz{facsimiles/register-563-left}{Sach-Register. Linke Druckseite von Scan 563.}

\notensatz{facsimiles/register-563-right}{Sach-Register. Rechte Druckseite von Scan 563.}

\quellseite{564}
\notensatz{facsimiles/register-564-left}{Sach-Register. Linke Druckseite von Scan 564.}

\notensatz{facsimiles/register-564-right}{Sach-Register. Rechte Druckseite von Scan 564.}

\quellseite{565}
\notensatz{facsimiles/register-565-left}{Sach-Register. Linke Druckseite von Scan 565.}

\notensatz{facsimiles/register-565-right}{Sach-Register. Rechte Druckseite von Scan 565.}

\quellseite{566}
\notensatz{facsimiles/register-566-left}{Sach-Register. Linke Druckseite von Scan 566.}

\notensatz{facsimiles/register-566-right}{Sach-Register. Rechte Druckseite von Scan 566.}

\quellseite{567}
\notensatz{facsimiles/register-567-left}{Sach-Register. Linke Druckseite von Scan 567.}

\notensatz{facsimiles/register-567-right}{Sach-Register. Rechte Druckseite von Scan 567.}

\quellseite{568}
\notensatz{facsimiles/register-568-left}{Sach-Register. Linke Druckseite von Scan 568.}

\notensatz{facsimiles/register-568-right}{Sach-Register. Rechte Druckseite von Scan 568.}

\quellseite{569}
\notensatz{facsimiles/register-569-left}{Sach-Register. Linke Druckseite von Scan 569.}

\notensatz{facsimiles/register-569-right}{Sach-Register. Rechte Druckseite von Scan 569.}

\quellseite{570}
\notensatz{facsimiles/register-570-left}{Sach-Register. Linke Druckseite von Scan 570.}

\notensatz{facsimiles/register-570-right}{Sach-Register. Rechte Druckseite von Scan 570.}

\quellseite{571}
\notensatz{facsimiles/register-571-left}{Sach-Register. Linke Druckseite von Scan 571.}

\notensatz{facsimiles/register-571-right}{Sach-Register. Rechte Druckseite von Scan 571.}

\quellseite{572}
\notensatz{facsimiles/register-572-left}{Sach-Register. Linke Druckseite von Scan 572.}

\notensatz{facsimiles/register-572-right}{Sach-Register. Rechte Druckseite von Scan 572.}

\quellseite{573}
\notensatz{facsimiles/register-573-left}{Sach-Register. Linke Druckseite von Scan 573.}

\notensatz{facsimiles/register-573-right}{Sach-Register. Rechte Druckseite von Scan 573.}

\quellseite{574}
\notensatz{facsimiles/register-574-left}{Sach-Register. Linke Druckseite von Scan 574.}

\notensatz{facsimiles/register-574-right}{Sach-Register. Rechte Druckseite von Scan 574.}

\quellseite{575}
\notensatz{facsimiles/register-575-left}{Sach-Register. Linke Druckseite von Scan 575.}

\notensatz{facsimiles/register-575-right}{Sach-Register. Rechte Druckseite von Scan 575.}

\quellseite{576}
\notensatz{facsimiles/register-576-left}{Sach-Register. Linke Druckseite von Scan 576.}

\notensatz{facsimiles/register-576-right}{Sach-Register. Rechte Druckseite von Scan 576.}

\quellseite{577}
\notensatz{facsimiles/register-577-left}{Sach-Register. Linke Druckseite von Scan 577.}

\notensatz{facsimiles/register-577-right}{Sach-Register. Rechte Druckseite von Scan 577.}

\quellseite{578}
\notensatz{facsimiles/register-578-left}{Sach-Register. Linke Druckseite von Scan 578.}

\notensatz{facsimiles/register-578-right}{Sach-Register. Rechte Druckseite von Scan 578.}

\quellseite{579}
\notensatz{facsimiles/register-579-left}{Sach-Register. Linke Druckseite von Scan 579.}

\notensatz{facsimiles/register-579-right}{Sach-Register. Rechte Druckseite von Scan 579.}

\quellseite{580}
\notensatz{facsimiles/register-580-left}{Sach-Register. Linke Druckseite von Scan 580.}
